\section*{Item 4}\label{it:4}

\textbf{Enunciado}: Simule, no LTspice, o circuito do conversor Flyback projetado no Exercício~1, utilizando o transistor IRFH5007. Acione o transistor com uma tensão de gate entre $0$ e $10$ V, bem como um resistor de gate de $10~\Omega$.

\medskip
\textbf{a)} Apresente as formas de onda da corrente no dreno do transistor, da tensão dreno--source do transistor e da tensão de saída para $K~L1~L2~1$ no LTspice.

\medskip
\textbf{Resolução}:
a Figura~\ref{fig:ex4_1} apresenta as formas de onda solicitadas no enunciado.
\begin{figure}[htb!]
	\centering
	%% Creator: Matplotlib, PGF backend
%%
%% To include the figure in your LaTeX document, write
%%   \input{<filename>.pgf}
%%
%% Make sure the required packages are loaded in your preamble
%%   \usepackage{pgf}
%%
%% Also ensure that all the required font packages are loaded; for instance,
%% the lmodern package is sometimes necessary when using math font.
%%   \usepackage{lmodern}
%%
%% Figures using additional raster images can only be included by \input if
%% they are in the same directory as the main LaTeX file. For loading figures
%% from other directories you can use the `import` package
%%   \usepackage{import}
%%
%% and then include the figures with
%%   \import{<path to file>}{<filename>.pgf}
%%
%% Matplotlib used the following preamble
%%   \def\mathdefault#1{#1}
%%   \everymath=\expandafter{\the\everymath\displaystyle}
%%   \IfFileExists{scrextend.sty}{
%%     \usepackage[fontsize=8.000000pt]{scrextend}
%%   }{
%%     \renewcommand{\normalsize}{\fontsize{8.000000}{9.600000}\selectfont}
%%     \normalsize
%%   }
%%   
%%           \usepackage{amsmath}
%%           \usepackage{amssymb}
%%           \usepackage{mathtools}
%%           \usepackage{dsfont}
%%           \providecommand{\mathdefault}[1]{#1}
%%       
%%   \makeatletter\@ifpackageloaded{underscore}{}{\usepackage[strings]{underscore}}\makeatother
%%
\begingroup%
\makeatletter%
\begin{pgfpicture}%
\pgfpathrectangle{\pgfpointorigin}{\pgfqpoint{3.393768in}{3.646235in}}%
\pgfusepath{use as bounding box, clip}%
\begin{pgfscope}%
\pgfsetbuttcap%
\pgfsetmiterjoin%
\definecolor{currentfill}{rgb}{1.000000,1.000000,1.000000}%
\pgfsetfillcolor{currentfill}%
\pgfsetlinewidth{0.000000pt}%
\definecolor{currentstroke}{rgb}{1.000000,1.000000,1.000000}%
\pgfsetstrokecolor{currentstroke}%
\pgfsetdash{}{0pt}%
\pgfpathmoveto{\pgfqpoint{0.000000in}{0.000000in}}%
\pgfpathlineto{\pgfqpoint{3.393768in}{0.000000in}}%
\pgfpathlineto{\pgfqpoint{3.393768in}{3.646235in}}%
\pgfpathlineto{\pgfqpoint{0.000000in}{3.646235in}}%
\pgfpathlineto{\pgfqpoint{0.000000in}{0.000000in}}%
\pgfpathclose%
\pgfusepath{fill}%
\end{pgfscope}%
\begin{pgfscope}%
\pgfsetbuttcap%
\pgfsetmiterjoin%
\definecolor{currentfill}{rgb}{1.000000,1.000000,1.000000}%
\pgfsetfillcolor{currentfill}%
\pgfsetlinewidth{0.000000pt}%
\definecolor{currentstroke}{rgb}{0.000000,0.000000,0.000000}%
\pgfsetstrokecolor{currentstroke}%
\pgfsetstrokeopacity{0.000000}%
\pgfsetdash{}{0pt}%
\pgfpathmoveto{\pgfqpoint{0.573768in}{2.601404in}}%
\pgfpathlineto{\pgfqpoint{3.293768in}{2.601404in}}%
\pgfpathlineto{\pgfqpoint{3.293768in}{3.507654in}}%
\pgfpathlineto{\pgfqpoint{0.573768in}{3.507654in}}%
\pgfpathlineto{\pgfqpoint{0.573768in}{2.601404in}}%
\pgfpathclose%
\pgfusepath{fill}%
\end{pgfscope}%
\begin{pgfscope}%
\pgfpathrectangle{\pgfqpoint{0.573768in}{2.601404in}}{\pgfqpoint{2.720000in}{0.906250in}}%
\pgfusepath{clip}%
\pgfsetbuttcap%
\pgfsetroundjoin%
\pgfsetlinewidth{0.501875pt}%
\definecolor{currentstroke}{rgb}{0.690196,0.690196,0.690196}%
\pgfsetstrokecolor{currentstroke}%
\pgfsetstrokeopacity{0.250000}%
\pgfsetdash{{1.850000pt}{0.800000pt}}{0.000000pt}%
\pgfpathmoveto{\pgfqpoint{0.697405in}{2.601404in}}%
\pgfpathlineto{\pgfqpoint{0.697405in}{3.507654in}}%
\pgfusepath{stroke}%
\end{pgfscope}%
\begin{pgfscope}%
\pgfsetbuttcap%
\pgfsetroundjoin%
\definecolor{currentfill}{rgb}{0.000000,0.000000,0.000000}%
\pgfsetfillcolor{currentfill}%
\pgfsetlinewidth{0.803000pt}%
\definecolor{currentstroke}{rgb}{0.000000,0.000000,0.000000}%
\pgfsetstrokecolor{currentstroke}%
\pgfsetdash{}{0pt}%
\pgfsys@defobject{currentmarker}{\pgfqpoint{0.000000in}{-0.048611in}}{\pgfqpoint{0.000000in}{0.000000in}}{%
\pgfpathmoveto{\pgfqpoint{0.000000in}{0.000000in}}%
\pgfpathlineto{\pgfqpoint{0.000000in}{-0.048611in}}%
\pgfusepath{stroke,fill}%
}%
\begin{pgfscope}%
\pgfsys@transformshift{0.697405in}{2.601404in}%
\pgfsys@useobject{currentmarker}{}%
\end{pgfscope}%
\end{pgfscope}%
\begin{pgfscope}%
\pgfpathrectangle{\pgfqpoint{0.573768in}{2.601404in}}{\pgfqpoint{2.720000in}{0.906250in}}%
\pgfusepath{clip}%
\pgfsetbuttcap%
\pgfsetroundjoin%
\pgfsetlinewidth{0.501875pt}%
\definecolor{currentstroke}{rgb}{0.690196,0.690196,0.690196}%
\pgfsetstrokecolor{currentstroke}%
\pgfsetstrokeopacity{0.250000}%
\pgfsetdash{{1.850000pt}{0.800000pt}}{0.000000pt}%
\pgfpathmoveto{\pgfqpoint{1.191950in}{2.601404in}}%
\pgfpathlineto{\pgfqpoint{1.191950in}{3.507654in}}%
\pgfusepath{stroke}%
\end{pgfscope}%
\begin{pgfscope}%
\pgfsetbuttcap%
\pgfsetroundjoin%
\definecolor{currentfill}{rgb}{0.000000,0.000000,0.000000}%
\pgfsetfillcolor{currentfill}%
\pgfsetlinewidth{0.803000pt}%
\definecolor{currentstroke}{rgb}{0.000000,0.000000,0.000000}%
\pgfsetstrokecolor{currentstroke}%
\pgfsetdash{}{0pt}%
\pgfsys@defobject{currentmarker}{\pgfqpoint{0.000000in}{-0.048611in}}{\pgfqpoint{0.000000in}{0.000000in}}{%
\pgfpathmoveto{\pgfqpoint{0.000000in}{0.000000in}}%
\pgfpathlineto{\pgfqpoint{0.000000in}{-0.048611in}}%
\pgfusepath{stroke,fill}%
}%
\begin{pgfscope}%
\pgfsys@transformshift{1.191950in}{2.601404in}%
\pgfsys@useobject{currentmarker}{}%
\end{pgfscope}%
\end{pgfscope}%
\begin{pgfscope}%
\pgfpathrectangle{\pgfqpoint{0.573768in}{2.601404in}}{\pgfqpoint{2.720000in}{0.906250in}}%
\pgfusepath{clip}%
\pgfsetbuttcap%
\pgfsetroundjoin%
\pgfsetlinewidth{0.501875pt}%
\definecolor{currentstroke}{rgb}{0.690196,0.690196,0.690196}%
\pgfsetstrokecolor{currentstroke}%
\pgfsetstrokeopacity{0.250000}%
\pgfsetdash{{1.850000pt}{0.800000pt}}{0.000000pt}%
\pgfpathmoveto{\pgfqpoint{1.686496in}{2.601404in}}%
\pgfpathlineto{\pgfqpoint{1.686496in}{3.507654in}}%
\pgfusepath{stroke}%
\end{pgfscope}%
\begin{pgfscope}%
\pgfsetbuttcap%
\pgfsetroundjoin%
\definecolor{currentfill}{rgb}{0.000000,0.000000,0.000000}%
\pgfsetfillcolor{currentfill}%
\pgfsetlinewidth{0.803000pt}%
\definecolor{currentstroke}{rgb}{0.000000,0.000000,0.000000}%
\pgfsetstrokecolor{currentstroke}%
\pgfsetdash{}{0pt}%
\pgfsys@defobject{currentmarker}{\pgfqpoint{0.000000in}{-0.048611in}}{\pgfqpoint{0.000000in}{0.000000in}}{%
\pgfpathmoveto{\pgfqpoint{0.000000in}{0.000000in}}%
\pgfpathlineto{\pgfqpoint{0.000000in}{-0.048611in}}%
\pgfusepath{stroke,fill}%
}%
\begin{pgfscope}%
\pgfsys@transformshift{1.686496in}{2.601404in}%
\pgfsys@useobject{currentmarker}{}%
\end{pgfscope}%
\end{pgfscope}%
\begin{pgfscope}%
\pgfpathrectangle{\pgfqpoint{0.573768in}{2.601404in}}{\pgfqpoint{2.720000in}{0.906250in}}%
\pgfusepath{clip}%
\pgfsetbuttcap%
\pgfsetroundjoin%
\pgfsetlinewidth{0.501875pt}%
\definecolor{currentstroke}{rgb}{0.690196,0.690196,0.690196}%
\pgfsetstrokecolor{currentstroke}%
\pgfsetstrokeopacity{0.250000}%
\pgfsetdash{{1.850000pt}{0.800000pt}}{0.000000pt}%
\pgfpathmoveto{\pgfqpoint{2.181041in}{2.601404in}}%
\pgfpathlineto{\pgfqpoint{2.181041in}{3.507654in}}%
\pgfusepath{stroke}%
\end{pgfscope}%
\begin{pgfscope}%
\pgfsetbuttcap%
\pgfsetroundjoin%
\definecolor{currentfill}{rgb}{0.000000,0.000000,0.000000}%
\pgfsetfillcolor{currentfill}%
\pgfsetlinewidth{0.803000pt}%
\definecolor{currentstroke}{rgb}{0.000000,0.000000,0.000000}%
\pgfsetstrokecolor{currentstroke}%
\pgfsetdash{}{0pt}%
\pgfsys@defobject{currentmarker}{\pgfqpoint{0.000000in}{-0.048611in}}{\pgfqpoint{0.000000in}{0.000000in}}{%
\pgfpathmoveto{\pgfqpoint{0.000000in}{0.000000in}}%
\pgfpathlineto{\pgfqpoint{0.000000in}{-0.048611in}}%
\pgfusepath{stroke,fill}%
}%
\begin{pgfscope}%
\pgfsys@transformshift{2.181041in}{2.601404in}%
\pgfsys@useobject{currentmarker}{}%
\end{pgfscope}%
\end{pgfscope}%
\begin{pgfscope}%
\pgfpathrectangle{\pgfqpoint{0.573768in}{2.601404in}}{\pgfqpoint{2.720000in}{0.906250in}}%
\pgfusepath{clip}%
\pgfsetbuttcap%
\pgfsetroundjoin%
\pgfsetlinewidth{0.501875pt}%
\definecolor{currentstroke}{rgb}{0.690196,0.690196,0.690196}%
\pgfsetstrokecolor{currentstroke}%
\pgfsetstrokeopacity{0.250000}%
\pgfsetdash{{1.850000pt}{0.800000pt}}{0.000000pt}%
\pgfpathmoveto{\pgfqpoint{2.675587in}{2.601404in}}%
\pgfpathlineto{\pgfqpoint{2.675587in}{3.507654in}}%
\pgfusepath{stroke}%
\end{pgfscope}%
\begin{pgfscope}%
\pgfsetbuttcap%
\pgfsetroundjoin%
\definecolor{currentfill}{rgb}{0.000000,0.000000,0.000000}%
\pgfsetfillcolor{currentfill}%
\pgfsetlinewidth{0.803000pt}%
\definecolor{currentstroke}{rgb}{0.000000,0.000000,0.000000}%
\pgfsetstrokecolor{currentstroke}%
\pgfsetdash{}{0pt}%
\pgfsys@defobject{currentmarker}{\pgfqpoint{0.000000in}{-0.048611in}}{\pgfqpoint{0.000000in}{0.000000in}}{%
\pgfpathmoveto{\pgfqpoint{0.000000in}{0.000000in}}%
\pgfpathlineto{\pgfqpoint{0.000000in}{-0.048611in}}%
\pgfusepath{stroke,fill}%
}%
\begin{pgfscope}%
\pgfsys@transformshift{2.675587in}{2.601404in}%
\pgfsys@useobject{currentmarker}{}%
\end{pgfscope}%
\end{pgfscope}%
\begin{pgfscope}%
\pgfpathrectangle{\pgfqpoint{0.573768in}{2.601404in}}{\pgfqpoint{2.720000in}{0.906250in}}%
\pgfusepath{clip}%
\pgfsetbuttcap%
\pgfsetroundjoin%
\pgfsetlinewidth{0.501875pt}%
\definecolor{currentstroke}{rgb}{0.690196,0.690196,0.690196}%
\pgfsetstrokecolor{currentstroke}%
\pgfsetstrokeopacity{0.250000}%
\pgfsetdash{{1.850000pt}{0.800000pt}}{0.000000pt}%
\pgfpathmoveto{\pgfqpoint{3.170132in}{2.601404in}}%
\pgfpathlineto{\pgfqpoint{3.170132in}{3.507654in}}%
\pgfusepath{stroke}%
\end{pgfscope}%
\begin{pgfscope}%
\pgfsetbuttcap%
\pgfsetroundjoin%
\definecolor{currentfill}{rgb}{0.000000,0.000000,0.000000}%
\pgfsetfillcolor{currentfill}%
\pgfsetlinewidth{0.803000pt}%
\definecolor{currentstroke}{rgb}{0.000000,0.000000,0.000000}%
\pgfsetstrokecolor{currentstroke}%
\pgfsetdash{}{0pt}%
\pgfsys@defobject{currentmarker}{\pgfqpoint{0.000000in}{-0.048611in}}{\pgfqpoint{0.000000in}{0.000000in}}{%
\pgfpathmoveto{\pgfqpoint{0.000000in}{0.000000in}}%
\pgfpathlineto{\pgfqpoint{0.000000in}{-0.048611in}}%
\pgfusepath{stroke,fill}%
}%
\begin{pgfscope}%
\pgfsys@transformshift{3.170132in}{2.601404in}%
\pgfsys@useobject{currentmarker}{}%
\end{pgfscope}%
\end{pgfscope}%
\begin{pgfscope}%
\pgfpathrectangle{\pgfqpoint{0.573768in}{2.601404in}}{\pgfqpoint{2.720000in}{0.906250in}}%
\pgfusepath{clip}%
\pgfsetbuttcap%
\pgfsetroundjoin%
\pgfsetlinewidth{0.501875pt}%
\definecolor{currentstroke}{rgb}{0.690196,0.690196,0.690196}%
\pgfsetstrokecolor{currentstroke}%
\pgfsetstrokeopacity{0.250000}%
\pgfsetdash{{1.850000pt}{0.800000pt}}{0.000000pt}%
\pgfpathmoveto{\pgfqpoint{0.573768in}{2.644559in}}%
\pgfpathlineto{\pgfqpoint{3.293768in}{2.644559in}}%
\pgfusepath{stroke}%
\end{pgfscope}%
\begin{pgfscope}%
\pgfsetbuttcap%
\pgfsetroundjoin%
\definecolor{currentfill}{rgb}{0.000000,0.000000,0.000000}%
\pgfsetfillcolor{currentfill}%
\pgfsetlinewidth{0.803000pt}%
\definecolor{currentstroke}{rgb}{0.000000,0.000000,0.000000}%
\pgfsetstrokecolor{currentstroke}%
\pgfsetdash{}{0pt}%
\pgfsys@defobject{currentmarker}{\pgfqpoint{-0.048611in}{0.000000in}}{\pgfqpoint{-0.000000in}{0.000000in}}{%
\pgfpathmoveto{\pgfqpoint{-0.000000in}{0.000000in}}%
\pgfpathlineto{\pgfqpoint{-0.048611in}{0.000000in}}%
\pgfusepath{stroke,fill}%
}%
\begin{pgfscope}%
\pgfsys@transformshift{0.573768in}{2.644559in}%
\pgfsys@useobject{currentmarker}{}%
\end{pgfscope}%
\end{pgfscope}%
\begin{pgfscope}%
\definecolor{textcolor}{rgb}{0.000000,0.000000,0.000000}%
\pgfsetstrokecolor{textcolor}%
\pgfsetfillcolor{textcolor}%
\pgftext[x=0.417518in, y=2.605979in, left, base]{\color{textcolor}{\rmfamily\fontsize{8.000000}{9.600000}\selectfont\catcode`\^=\active\def^{\ifmmode\sp\else\^{}\fi}\catcode`\%=\active\def%{\%}$\mathdefault{0}$}}%
\end{pgfscope}%
\begin{pgfscope}%
\pgfpathrectangle{\pgfqpoint{0.573768in}{2.601404in}}{\pgfqpoint{2.720000in}{0.906250in}}%
\pgfusepath{clip}%
\pgfsetbuttcap%
\pgfsetroundjoin%
\pgfsetlinewidth{0.501875pt}%
\definecolor{currentstroke}{rgb}{0.690196,0.690196,0.690196}%
\pgfsetstrokecolor{currentstroke}%
\pgfsetstrokeopacity{0.250000}%
\pgfsetdash{{1.850000pt}{0.800000pt}}{0.000000pt}%
\pgfpathmoveto{\pgfqpoint{0.573768in}{3.076107in}}%
\pgfpathlineto{\pgfqpoint{3.293768in}{3.076107in}}%
\pgfusepath{stroke}%
\end{pgfscope}%
\begin{pgfscope}%
\pgfsetbuttcap%
\pgfsetroundjoin%
\definecolor{currentfill}{rgb}{0.000000,0.000000,0.000000}%
\pgfsetfillcolor{currentfill}%
\pgfsetlinewidth{0.803000pt}%
\definecolor{currentstroke}{rgb}{0.000000,0.000000,0.000000}%
\pgfsetstrokecolor{currentstroke}%
\pgfsetdash{}{0pt}%
\pgfsys@defobject{currentmarker}{\pgfqpoint{-0.048611in}{0.000000in}}{\pgfqpoint{-0.000000in}{0.000000in}}{%
\pgfpathmoveto{\pgfqpoint{-0.000000in}{0.000000in}}%
\pgfpathlineto{\pgfqpoint{-0.048611in}{0.000000in}}%
\pgfusepath{stroke,fill}%
}%
\begin{pgfscope}%
\pgfsys@transformshift{0.573768in}{3.076107in}%
\pgfsys@useobject{currentmarker}{}%
\end{pgfscope}%
\end{pgfscope}%
\begin{pgfscope}%
\definecolor{textcolor}{rgb}{0.000000,0.000000,0.000000}%
\pgfsetstrokecolor{textcolor}%
\pgfsetfillcolor{textcolor}%
\pgftext[x=0.417518in, y=3.037527in, left, base]{\color{textcolor}{\rmfamily\fontsize{8.000000}{9.600000}\selectfont\catcode`\^=\active\def^{\ifmmode\sp\else\^{}\fi}\catcode`\%=\active\def%{\%}$\mathdefault{5}$}}%
\end{pgfscope}%
\begin{pgfscope}%
\pgfpathrectangle{\pgfqpoint{0.573768in}{2.601404in}}{\pgfqpoint{2.720000in}{0.906250in}}%
\pgfusepath{clip}%
\pgfsetbuttcap%
\pgfsetroundjoin%
\pgfsetlinewidth{0.501875pt}%
\definecolor{currentstroke}{rgb}{0.690196,0.690196,0.690196}%
\pgfsetstrokecolor{currentstroke}%
\pgfsetstrokeopacity{0.250000}%
\pgfsetdash{{1.850000pt}{0.800000pt}}{0.000000pt}%
\pgfpathmoveto{\pgfqpoint{0.573768in}{3.507654in}}%
\pgfpathlineto{\pgfqpoint{3.293768in}{3.507654in}}%
\pgfusepath{stroke}%
\end{pgfscope}%
\begin{pgfscope}%
\pgfsetbuttcap%
\pgfsetroundjoin%
\definecolor{currentfill}{rgb}{0.000000,0.000000,0.000000}%
\pgfsetfillcolor{currentfill}%
\pgfsetlinewidth{0.803000pt}%
\definecolor{currentstroke}{rgb}{0.000000,0.000000,0.000000}%
\pgfsetstrokecolor{currentstroke}%
\pgfsetdash{}{0pt}%
\pgfsys@defobject{currentmarker}{\pgfqpoint{-0.048611in}{0.000000in}}{\pgfqpoint{-0.000000in}{0.000000in}}{%
\pgfpathmoveto{\pgfqpoint{-0.000000in}{0.000000in}}%
\pgfpathlineto{\pgfqpoint{-0.048611in}{0.000000in}}%
\pgfusepath{stroke,fill}%
}%
\begin{pgfscope}%
\pgfsys@transformshift{0.573768in}{3.507654in}%
\pgfsys@useobject{currentmarker}{}%
\end{pgfscope}%
\end{pgfscope}%
\begin{pgfscope}%
\definecolor{textcolor}{rgb}{0.000000,0.000000,0.000000}%
\pgfsetstrokecolor{textcolor}%
\pgfsetfillcolor{textcolor}%
\pgftext[x=0.358489in, y=3.469074in, left, base]{\color{textcolor}{\rmfamily\fontsize{8.000000}{9.600000}\selectfont\catcode`\^=\active\def^{\ifmmode\sp\else\^{}\fi}\catcode`\%=\active\def%{\%}$\mathdefault{10}$}}%
\end{pgfscope}%
\begin{pgfscope}%
\definecolor{textcolor}{rgb}{0.000000,0.000000,0.000000}%
\pgfsetstrokecolor{textcolor}%
\pgfsetfillcolor{textcolor}%
\pgftext[x=0.302933in,y=3.054529in,,bottom,rotate=90.000000]{\color{textcolor}{\rmfamily\fontsize{8.000000}{9.600000}\selectfont\catcode`\^=\active\def^{\ifmmode\sp\else\^{}\fi}\catcode`\%=\active\def%{\%}$i_D$ (A)}}%
\end{pgfscope}%
\begin{pgfscope}%
\pgfpathrectangle{\pgfqpoint{0.573768in}{2.601404in}}{\pgfqpoint{2.720000in}{0.906250in}}%
\pgfusepath{clip}%
\pgfsetrectcap%
\pgfsetroundjoin%
\pgfsetlinewidth{1.505625pt}%
\definecolor{currentstroke}{rgb}{0.000000,0.447059,0.698039}%
\pgfsetstrokecolor{currentstroke}%
\pgfsetdash{}{0pt}%
\pgfpathmoveto{\pgfqpoint{0.697405in}{2.644559in}}%
\pgfpathlineto{\pgfqpoint{0.697409in}{2.644559in}}%
\pgfpathlineto{\pgfqpoint{0.697416in}{2.743950in}}%
\pgfpathlineto{\pgfqpoint{0.698332in}{2.702816in}}%
\pgfpathlineto{\pgfqpoint{0.699384in}{2.768855in}}%
\pgfpathlineto{\pgfqpoint{0.699411in}{2.755130in}}%
\pgfpathlineto{\pgfqpoint{0.700200in}{2.644559in}}%
\pgfpathlineto{\pgfqpoint{0.702356in}{2.644559in}}%
\pgfpathlineto{\pgfqpoint{0.702361in}{2.869774in}}%
\pgfpathlineto{\pgfqpoint{0.703089in}{2.844532in}}%
\pgfpathlineto{\pgfqpoint{0.703138in}{2.782560in}}%
\pgfpathlineto{\pgfqpoint{0.703781in}{2.887873in}}%
\pgfpathlineto{\pgfqpoint{0.703831in}{2.826191in}}%
\pgfpathlineto{\pgfqpoint{0.704329in}{2.922122in}}%
\pgfpathlineto{\pgfqpoint{0.704402in}{2.644559in}}%
\pgfpathlineto{\pgfqpoint{0.704449in}{2.644559in}}%
\pgfpathlineto{\pgfqpoint{0.707301in}{2.644559in}}%
\pgfpathlineto{\pgfqpoint{0.707312in}{3.006975in}}%
\pgfpathlineto{\pgfqpoint{0.708040in}{2.946658in}}%
\pgfpathlineto{\pgfqpoint{0.708090in}{2.912359in}}%
\pgfpathlineto{\pgfqpoint{0.708733in}{2.990041in}}%
\pgfpathlineto{\pgfqpoint{0.708782in}{2.955932in}}%
\pgfpathlineto{\pgfqpoint{0.709227in}{3.021028in}}%
\pgfpathlineto{\pgfqpoint{0.709329in}{2.920687in}}%
\pgfpathlineto{\pgfqpoint{0.710083in}{2.644559in}}%
\pgfpathlineto{\pgfqpoint{0.712248in}{2.644688in}}%
\pgfpathlineto{\pgfqpoint{0.712253in}{3.113946in}}%
\pgfpathlineto{\pgfqpoint{0.713002in}{3.023773in}}%
\pgfpathlineto{\pgfqpoint{0.713843in}{3.101363in}}%
\pgfpathlineto{\pgfqpoint{0.714220in}{3.125003in}}%
\pgfpathlineto{\pgfqpoint{0.715622in}{2.644559in}}%
\pgfpathlineto{\pgfqpoint{0.717193in}{2.644575in}}%
\pgfpathlineto{\pgfqpoint{0.717199in}{3.199080in}}%
\pgfpathlineto{\pgfqpoint{0.717926in}{3.152350in}}%
\pgfpathlineto{\pgfqpoint{0.717976in}{3.112374in}}%
\pgfpathlineto{\pgfqpoint{0.718618in}{3.195704in}}%
\pgfpathlineto{\pgfqpoint{0.718668in}{3.155947in}}%
\pgfpathlineto{\pgfqpoint{0.719165in}{3.229923in}}%
\pgfpathlineto{\pgfqpoint{0.719221in}{3.145377in}}%
\pgfpathlineto{\pgfqpoint{0.719990in}{2.644559in}}%
\pgfpathlineto{\pgfqpoint{0.722138in}{2.644559in}}%
\pgfpathlineto{\pgfqpoint{0.722144in}{3.269503in}}%
\pgfpathlineto{\pgfqpoint{0.722873in}{3.234143in}}%
\pgfpathlineto{\pgfqpoint{0.722922in}{3.194994in}}%
\pgfpathlineto{\pgfqpoint{0.723565in}{3.277488in}}%
\pgfpathlineto{\pgfqpoint{0.723615in}{3.238564in}}%
\pgfpathlineto{\pgfqpoint{0.724111in}{3.311606in}}%
\pgfpathlineto{\pgfqpoint{0.724167in}{3.223413in}}%
\pgfpathlineto{\pgfqpoint{0.724935in}{2.644559in}}%
\pgfpathlineto{\pgfqpoint{0.727084in}{2.644573in}}%
\pgfpathlineto{\pgfqpoint{0.727090in}{3.342230in}}%
\pgfpathlineto{\pgfqpoint{0.727830in}{3.253503in}}%
\pgfpathlineto{\pgfqpoint{0.728670in}{3.359050in}}%
\pgfpathlineto{\pgfqpoint{0.729056in}{3.383196in}}%
\pgfpathlineto{\pgfqpoint{0.730133in}{2.644559in}}%
\pgfpathlineto{\pgfqpoint{0.732030in}{2.644608in}}%
\pgfpathlineto{\pgfqpoint{0.732035in}{3.392488in}}%
\pgfpathlineto{\pgfqpoint{0.732803in}{3.328914in}}%
\pgfpathlineto{\pgfqpoint{0.734002in}{3.407003in}}%
\pgfpathlineto{\pgfqpoint{0.734002in}{3.381636in}}%
\pgfpathlineto{\pgfqpoint{0.734029in}{3.386623in}}%
\pgfpathlineto{\pgfqpoint{0.734783in}{2.644559in}}%
\pgfpathlineto{\pgfqpoint{0.736975in}{2.644560in}}%
\pgfpathlineto{\pgfqpoint{0.736981in}{3.413370in}}%
\pgfpathlineto{\pgfqpoint{0.737738in}{3.358724in}}%
\pgfpathlineto{\pgfqpoint{0.738947in}{3.436583in}}%
\pgfpathlineto{\pgfqpoint{0.738947in}{3.411600in}}%
\pgfpathlineto{\pgfqpoint{0.738965in}{3.416410in}}%
\pgfpathlineto{\pgfqpoint{0.739727in}{2.644559in}}%
\pgfpathlineto{\pgfqpoint{0.741920in}{2.644560in}}%
\pgfpathlineto{\pgfqpoint{0.741926in}{3.422089in}}%
\pgfpathlineto{\pgfqpoint{0.742665in}{3.351991in}}%
\pgfpathlineto{\pgfqpoint{0.743505in}{3.439161in}}%
\pgfpathlineto{\pgfqpoint{0.743893in}{3.463416in}}%
\pgfpathlineto{\pgfqpoint{0.743961in}{2.678676in}}%
\pgfpathlineto{\pgfqpoint{0.744745in}{2.644559in}}%
\pgfpathlineto{\pgfqpoint{0.746866in}{2.644565in}}%
\pgfpathlineto{\pgfqpoint{0.746871in}{3.413957in}}%
\pgfpathlineto{\pgfqpoint{0.747639in}{3.385071in}}%
\pgfpathlineto{\pgfqpoint{0.747688in}{3.342472in}}%
\pgfpathlineto{\pgfqpoint{0.748331in}{3.428429in}}%
\pgfpathlineto{\pgfqpoint{0.748381in}{3.386008in}}%
\pgfpathlineto{\pgfqpoint{0.748807in}{3.458228in}}%
\pgfpathlineto{\pgfqpoint{0.748894in}{3.361627in}}%
\pgfpathlineto{\pgfqpoint{0.749663in}{2.644559in}}%
\pgfpathlineto{\pgfqpoint{0.751811in}{2.644572in}}%
\pgfpathlineto{\pgfqpoint{0.751828in}{3.398484in}}%
\pgfpathlineto{\pgfqpoint{0.752582in}{3.344713in}}%
\pgfpathlineto{\pgfqpoint{0.752631in}{3.326921in}}%
\pgfpathlineto{\pgfqpoint{0.753175in}{3.381909in}}%
\pgfpathlineto{\pgfqpoint{0.753225in}{3.364209in}}%
\pgfpathlineto{\pgfqpoint{0.753751in}{3.418018in}}%
\pgfpathlineto{\pgfqpoint{0.753838in}{3.323043in}}%
\pgfpathlineto{\pgfqpoint{0.753851in}{3.392471in}}%
\pgfpathlineto{\pgfqpoint{0.754617in}{2.644559in}}%
\pgfpathlineto{\pgfqpoint{0.756757in}{2.644560in}}%
\pgfpathlineto{\pgfqpoint{0.756774in}{3.349929in}}%
\pgfpathlineto{\pgfqpoint{0.757497in}{3.265090in}}%
\pgfpathlineto{\pgfqpoint{0.758337in}{3.360993in}}%
\pgfpathlineto{\pgfqpoint{0.758729in}{3.385511in}}%
\pgfpathlineto{\pgfqpoint{0.759975in}{2.644559in}}%
\pgfpathlineto{\pgfqpoint{0.761702in}{2.644559in}}%
\pgfpathlineto{\pgfqpoint{0.761720in}{3.282674in}}%
\pgfpathlineto{\pgfqpoint{0.762454in}{3.209576in}}%
\pgfpathlineto{\pgfqpoint{0.763295in}{3.293749in}}%
\pgfpathlineto{\pgfqpoint{0.763674in}{3.317530in}}%
\pgfpathlineto{\pgfqpoint{0.765185in}{2.644559in}}%
\pgfpathlineto{\pgfqpoint{0.766648in}{2.644563in}}%
\pgfpathlineto{\pgfqpoint{0.766664in}{3.212899in}}%
\pgfpathlineto{\pgfqpoint{0.767412in}{3.122549in}}%
\pgfpathlineto{\pgfqpoint{0.768253in}{3.229054in}}%
\pgfpathlineto{\pgfqpoint{0.768620in}{3.252000in}}%
\pgfpathlineto{\pgfqpoint{0.769569in}{2.644559in}}%
\pgfpathlineto{\pgfqpoint{0.771593in}{2.644559in}}%
\pgfpathlineto{\pgfqpoint{0.771610in}{3.115951in}}%
\pgfpathlineto{\pgfqpoint{0.772365in}{3.028234in}}%
\pgfpathlineto{\pgfqpoint{0.773205in}{3.146329in}}%
\pgfpathlineto{\pgfqpoint{0.773565in}{3.168837in}}%
\pgfpathlineto{\pgfqpoint{0.774373in}{2.644559in}}%
\pgfpathlineto{\pgfqpoint{0.776538in}{2.644559in}}%
\pgfpathlineto{\pgfqpoint{0.776557in}{3.012697in}}%
\pgfpathlineto{\pgfqpoint{0.777277in}{2.949633in}}%
\pgfpathlineto{\pgfqpoint{0.778117in}{3.024018in}}%
\pgfpathlineto{\pgfqpoint{0.778511in}{3.048707in}}%
\pgfpathlineto{\pgfqpoint{0.779873in}{2.644559in}}%
\pgfpathlineto{\pgfqpoint{0.781483in}{2.644559in}}%
\pgfpathlineto{\pgfqpoint{0.781501in}{2.907630in}}%
\pgfpathlineto{\pgfqpoint{0.782215in}{2.841371in}}%
\pgfpathlineto{\pgfqpoint{0.783056in}{2.921002in}}%
\pgfpathlineto{\pgfqpoint{0.783429in}{2.944400in}}%
\pgfpathlineto{\pgfqpoint{0.784391in}{2.644559in}}%
\pgfpathlineto{\pgfqpoint{0.786428in}{2.644559in}}%
\pgfpathlineto{\pgfqpoint{0.786438in}{2.804446in}}%
\pgfpathlineto{\pgfqpoint{0.787158in}{2.730481in}}%
\pgfpathlineto{\pgfqpoint{0.787999in}{2.813494in}}%
\pgfpathlineto{\pgfqpoint{0.788373in}{2.836964in}}%
\pgfpathlineto{\pgfqpoint{0.789052in}{2.644559in}}%
\pgfpathlineto{\pgfqpoint{0.791373in}{2.644559in}}%
\pgfpathlineto{\pgfqpoint{0.791387in}{2.759141in}}%
\pgfpathlineto{\pgfqpoint{0.792119in}{2.696010in}}%
\pgfpathlineto{\pgfqpoint{0.792169in}{2.696003in}}%
\pgfpathlineto{\pgfqpoint{0.793302in}{2.770325in}}%
\pgfpathlineto{\pgfqpoint{0.793347in}{2.747895in}}%
\pgfpathlineto{\pgfqpoint{0.793372in}{2.757029in}}%
\pgfpathlineto{\pgfqpoint{0.794077in}{2.644559in}}%
\pgfpathlineto{\pgfqpoint{0.796319in}{2.644560in}}%
\pgfpathlineto{\pgfqpoint{0.796323in}{2.764980in}}%
\pgfpathlineto{\pgfqpoint{0.796416in}{2.641558in}}%
\pgfpathlineto{\pgfqpoint{0.797089in}{2.682583in}}%
\pgfpathlineto{\pgfqpoint{0.797929in}{2.769170in}}%
\pgfpathlineto{\pgfqpoint{0.798293in}{2.791933in}}%
\pgfpathlineto{\pgfqpoint{0.798880in}{2.644559in}}%
\pgfpathlineto{\pgfqpoint{0.801264in}{2.644559in}}%
\pgfpathlineto{\pgfqpoint{0.801268in}{2.758223in}}%
\pgfpathlineto{\pgfqpoint{0.801990in}{2.686330in}}%
\pgfpathlineto{\pgfqpoint{0.802831in}{2.759560in}}%
\pgfpathlineto{\pgfqpoint{0.803208in}{2.783189in}}%
\pgfpathlineto{\pgfqpoint{0.803927in}{2.644559in}}%
\pgfpathlineto{\pgfqpoint{0.806210in}{2.644559in}}%
\pgfpathlineto{\pgfqpoint{0.806213in}{2.751128in}}%
\pgfpathlineto{\pgfqpoint{0.806937in}{2.689556in}}%
\pgfpathlineto{\pgfqpoint{0.807877in}{2.762051in}}%
\pgfpathlineto{\pgfqpoint{0.808154in}{2.779451in}}%
\pgfpathlineto{\pgfqpoint{0.809174in}{2.644559in}}%
\pgfpathlineto{\pgfqpoint{0.811155in}{2.644559in}}%
\pgfpathlineto{\pgfqpoint{0.811160in}{2.750855in}}%
\pgfpathlineto{\pgfqpoint{0.811924in}{2.695793in}}%
\pgfpathlineto{\pgfqpoint{0.813129in}{2.777480in}}%
\pgfpathlineto{\pgfqpoint{0.813880in}{2.644559in}}%
\pgfpathlineto{\pgfqpoint{0.816101in}{2.644559in}}%
\pgfpathlineto{\pgfqpoint{0.816106in}{2.751060in}}%
\pgfpathlineto{\pgfqpoint{0.816845in}{2.700500in}}%
\pgfpathlineto{\pgfqpoint{0.816894in}{2.697805in}}%
\pgfpathlineto{\pgfqpoint{0.817142in}{2.719145in}}%
\pgfpathlineto{\pgfqpoint{0.817191in}{2.716442in}}%
\pgfpathlineto{\pgfqpoint{0.818028in}{2.774859in}}%
\pgfpathlineto{\pgfqpoint{0.818098in}{2.759596in}}%
\pgfpathlineto{\pgfqpoint{0.818899in}{2.644559in}}%
\pgfpathlineto{\pgfqpoint{0.821046in}{2.644559in}}%
\pgfpathlineto{\pgfqpoint{0.821051in}{2.750602in}}%
\pgfpathlineto{\pgfqpoint{0.821800in}{2.710562in}}%
\pgfpathlineto{\pgfqpoint{0.821849in}{2.687994in}}%
\pgfpathlineto{\pgfqpoint{0.822393in}{2.747810in}}%
\pgfpathlineto{\pgfqpoint{0.822443in}{2.725308in}}%
\pgfpathlineto{\pgfqpoint{0.822979in}{2.784546in}}%
\pgfpathlineto{\pgfqpoint{0.823074in}{2.699334in}}%
\pgfpathlineto{\pgfqpoint{0.823089in}{2.750220in}}%
\pgfpathlineto{\pgfqpoint{0.823869in}{2.644559in}}%
\pgfpathlineto{\pgfqpoint{0.825992in}{2.644559in}}%
\pgfpathlineto{\pgfqpoint{0.825996in}{2.747070in}}%
\pgfpathlineto{\pgfqpoint{0.826063in}{2.621730in}}%
\pgfpathlineto{\pgfqpoint{0.826747in}{2.661877in}}%
\pgfpathlineto{\pgfqpoint{0.827588in}{2.785487in}}%
\pgfpathlineto{\pgfqpoint{0.827965in}{2.809128in}}%
\pgfpathlineto{\pgfqpoint{0.828434in}{2.644559in}}%
\pgfpathlineto{\pgfqpoint{0.830936in}{2.644559in}}%
\pgfpathlineto{\pgfqpoint{0.830942in}{2.749364in}}%
\pgfpathlineto{\pgfqpoint{0.831002in}{2.643924in}}%
\pgfpathlineto{\pgfqpoint{0.831684in}{2.682845in}}%
\pgfpathlineto{\pgfqpoint{0.832524in}{2.762622in}}%
\pgfpathlineto{\pgfqpoint{0.832911in}{2.786863in}}%
\pgfpathlineto{\pgfqpoint{0.833505in}{2.644559in}}%
\pgfpathlineto{\pgfqpoint{0.835882in}{2.644559in}}%
\pgfpathlineto{\pgfqpoint{0.835887in}{2.749574in}}%
\pgfpathlineto{\pgfqpoint{0.835977in}{2.641158in}}%
\pgfpathlineto{\pgfqpoint{0.836614in}{2.711058in}}%
\pgfpathlineto{\pgfqpoint{0.836664in}{2.683215in}}%
\pgfpathlineto{\pgfqpoint{0.837307in}{2.754495in}}%
\pgfpathlineto{\pgfqpoint{0.837356in}{2.726767in}}%
\pgfpathlineto{\pgfqpoint{0.837856in}{2.788955in}}%
\pgfpathlineto{\pgfqpoint{0.837943in}{2.644559in}}%
\pgfpathlineto{\pgfqpoint{0.838005in}{2.644559in}}%
\pgfpathlineto{\pgfqpoint{0.840828in}{2.644559in}}%
\pgfpathlineto{\pgfqpoint{0.840833in}{2.743484in}}%
\pgfpathlineto{\pgfqpoint{0.840913in}{2.624979in}}%
\pgfpathlineto{\pgfqpoint{0.841598in}{2.666561in}}%
\pgfpathlineto{\pgfqpoint{0.842439in}{2.782464in}}%
\pgfpathlineto{\pgfqpoint{0.842802in}{2.805140in}}%
\pgfpathlineto{\pgfqpoint{0.843315in}{2.644559in}}%
\pgfpathlineto{\pgfqpoint{0.845773in}{2.644559in}}%
\pgfpathlineto{\pgfqpoint{0.845778in}{2.745911in}}%
\pgfpathlineto{\pgfqpoint{0.846533in}{2.696366in}}%
\pgfpathlineto{\pgfqpoint{0.847747in}{2.772386in}}%
\pgfpathlineto{\pgfqpoint{0.847790in}{2.742315in}}%
\pgfpathlineto{\pgfqpoint{0.848602in}{2.644559in}}%
\pgfpathlineto{\pgfqpoint{0.850720in}{2.644617in}}%
\pgfpathlineto{\pgfqpoint{0.850724in}{2.766831in}}%
\pgfpathlineto{\pgfqpoint{0.851486in}{2.689411in}}%
\pgfpathlineto{\pgfqpoint{0.852426in}{2.765265in}}%
\pgfpathlineto{\pgfqpoint{0.852693in}{2.781992in}}%
\pgfpathlineto{\pgfqpoint{0.853632in}{2.644559in}}%
\pgfpathlineto{\pgfqpoint{0.855664in}{2.644559in}}%
\pgfpathlineto{\pgfqpoint{0.855669in}{2.733449in}}%
\pgfpathlineto{\pgfqpoint{0.855733in}{2.643162in}}%
\pgfpathlineto{\pgfqpoint{0.856418in}{2.683084in}}%
\pgfpathlineto{\pgfqpoint{0.857259in}{2.760332in}}%
\pgfpathlineto{\pgfqpoint{0.857638in}{2.784128in}}%
\pgfpathlineto{\pgfqpoint{0.858278in}{2.644559in}}%
\pgfpathlineto{\pgfqpoint{0.860610in}{2.644559in}}%
\pgfpathlineto{\pgfqpoint{0.860615in}{2.745594in}}%
\pgfpathlineto{\pgfqpoint{0.860695in}{2.617510in}}%
\pgfpathlineto{\pgfqpoint{0.861380in}{2.659064in}}%
\pgfpathlineto{\pgfqpoint{0.862221in}{2.787488in}}%
\pgfpathlineto{\pgfqpoint{0.862584in}{2.810167in}}%
\pgfpathlineto{\pgfqpoint{0.863062in}{2.644559in}}%
\pgfpathlineto{\pgfqpoint{0.865557in}{2.644914in}}%
\pgfpathlineto{\pgfqpoint{0.865559in}{2.762175in}}%
\pgfpathlineto{\pgfqpoint{0.865656in}{2.611594in}}%
\pgfpathlineto{\pgfqpoint{0.866294in}{2.737940in}}%
\pgfpathlineto{\pgfqpoint{0.866343in}{2.654236in}}%
\pgfpathlineto{\pgfqpoint{0.866986in}{2.781221in}}%
\pgfpathlineto{\pgfqpoint{0.867036in}{2.697946in}}%
\pgfpathlineto{\pgfqpoint{0.867481in}{2.812134in}}%
\pgfpathlineto{\pgfqpoint{0.867615in}{2.644559in}}%
\pgfpathlineto{\pgfqpoint{0.867772in}{2.644559in}}%
\pgfpathlineto{\pgfqpoint{0.870501in}{2.644559in}}%
\pgfpathlineto{\pgfqpoint{0.870506in}{2.746472in}}%
\pgfpathlineto{\pgfqpoint{0.871257in}{2.694854in}}%
\pgfpathlineto{\pgfqpoint{0.872474in}{2.772760in}}%
\pgfpathlineto{\pgfqpoint{0.872475in}{2.748413in}}%
\pgfpathlineto{\pgfqpoint{0.872498in}{2.756862in}}%
\pgfpathlineto{\pgfqpoint{0.873223in}{2.644559in}}%
\pgfpathlineto{\pgfqpoint{0.875446in}{2.644559in}}%
\pgfpathlineto{\pgfqpoint{0.875451in}{2.737050in}}%
\pgfpathlineto{\pgfqpoint{0.875526in}{2.640923in}}%
\pgfpathlineto{\pgfqpoint{0.876212in}{2.682018in}}%
\pgfpathlineto{\pgfqpoint{0.877052in}{2.762914in}}%
\pgfpathlineto{\pgfqpoint{0.877420in}{2.785958in}}%
\pgfpathlineto{\pgfqpoint{0.878057in}{2.644559in}}%
\pgfpathlineto{\pgfqpoint{0.880393in}{2.644668in}}%
\pgfpathlineto{\pgfqpoint{0.880395in}{2.748682in}}%
\pgfpathlineto{\pgfqpoint{0.880493in}{2.623339in}}%
\pgfpathlineto{\pgfqpoint{0.881130in}{2.724441in}}%
\pgfpathlineto{\pgfqpoint{0.881180in}{2.665955in}}%
\pgfpathlineto{\pgfqpoint{0.881823in}{2.767770in}}%
\pgfpathlineto{\pgfqpoint{0.881872in}{2.709615in}}%
\pgfpathlineto{\pgfqpoint{0.882317in}{2.798718in}}%
\pgfpathlineto{\pgfqpoint{0.882451in}{2.644559in}}%
\pgfpathlineto{\pgfqpoint{0.882587in}{2.644559in}}%
\pgfpathlineto{\pgfqpoint{0.885337in}{2.644559in}}%
\pgfpathlineto{\pgfqpoint{0.885342in}{2.733234in}}%
\pgfpathlineto{\pgfqpoint{0.885405in}{2.644405in}}%
\pgfpathlineto{\pgfqpoint{0.886089in}{2.684113in}}%
\pgfpathlineto{\pgfqpoint{0.886929in}{2.758201in}}%
\pgfpathlineto{\pgfqpoint{0.887311in}{2.782119in}}%
\pgfpathlineto{\pgfqpoint{0.887995in}{2.644559in}}%
\pgfpathlineto{\pgfqpoint{0.890282in}{2.644559in}}%
\pgfpathlineto{\pgfqpoint{0.890287in}{2.735322in}}%
\pgfpathlineto{\pgfqpoint{0.890349in}{2.640294in}}%
\pgfpathlineto{\pgfqpoint{0.891031in}{2.679719in}}%
\pgfpathlineto{\pgfqpoint{0.891872in}{2.761040in}}%
\pgfpathlineto{\pgfqpoint{0.892256in}{2.785123in}}%
\pgfpathlineto{\pgfqpoint{0.892847in}{2.644559in}}%
\pgfpathlineto{\pgfqpoint{0.895227in}{2.644559in}}%
\pgfpathlineto{\pgfqpoint{0.895235in}{2.733513in}}%
\pgfpathlineto{\pgfqpoint{0.895291in}{2.632993in}}%
\pgfpathlineto{\pgfqpoint{0.895974in}{2.671994in}}%
\pgfpathlineto{\pgfqpoint{0.896815in}{2.767729in}}%
\pgfpathlineto{\pgfqpoint{0.897202in}{2.792008in}}%
\pgfpathlineto{\pgfqpoint{0.897728in}{2.644559in}}%
\pgfpathlineto{\pgfqpoint{0.900173in}{2.644559in}}%
\pgfpathlineto{\pgfqpoint{0.900178in}{2.742206in}}%
\pgfpathlineto{\pgfqpoint{0.900259in}{2.624479in}}%
\pgfpathlineto{\pgfqpoint{0.900945in}{2.666067in}}%
\pgfpathlineto{\pgfqpoint{0.901785in}{2.777105in}}%
\pgfpathlineto{\pgfqpoint{0.902147in}{2.799743in}}%
\pgfpathlineto{\pgfqpoint{0.902646in}{2.644559in}}%
\pgfpathlineto{\pgfqpoint{0.905121in}{2.644620in}}%
\pgfpathlineto{\pgfqpoint{0.905123in}{2.767079in}}%
\pgfpathlineto{\pgfqpoint{0.905178in}{2.644540in}}%
\pgfpathlineto{\pgfqpoint{0.905856in}{2.682583in}}%
\pgfpathlineto{\pgfqpoint{0.906796in}{2.760470in}}%
\pgfpathlineto{\pgfqpoint{0.907093in}{2.779078in}}%
\pgfpathlineto{\pgfqpoint{0.907884in}{2.644559in}}%
\pgfpathlineto{\pgfqpoint{0.910063in}{2.644559in}}%
\pgfpathlineto{\pgfqpoint{0.910069in}{2.743014in}}%
\pgfpathlineto{\pgfqpoint{0.910150in}{2.614354in}}%
\pgfpathlineto{\pgfqpoint{0.910835in}{2.656064in}}%
\pgfpathlineto{\pgfqpoint{0.911675in}{2.785647in}}%
\pgfpathlineto{\pgfqpoint{0.912038in}{2.808310in}}%
\pgfpathlineto{\pgfqpoint{0.912498in}{2.644559in}}%
\pgfpathlineto{\pgfqpoint{0.915011in}{2.644562in}}%
\pgfpathlineto{\pgfqpoint{0.915013in}{2.739676in}}%
\pgfpathlineto{\pgfqpoint{0.915072in}{2.619619in}}%
\pgfpathlineto{\pgfqpoint{0.915802in}{2.728926in}}%
\pgfpathlineto{\pgfqpoint{0.915851in}{2.664649in}}%
\pgfpathlineto{\pgfqpoint{0.916494in}{2.772282in}}%
\pgfpathlineto{\pgfqpoint{0.916544in}{2.708282in}}%
\pgfpathlineto{\pgfqpoint{0.916984in}{2.802920in}}%
\pgfpathlineto{\pgfqpoint{0.917072in}{2.644559in}}%
\pgfpathlineto{\pgfqpoint{0.917226in}{2.644559in}}%
\pgfpathlineto{\pgfqpoint{0.919954in}{2.644559in}}%
\pgfpathlineto{\pgfqpoint{0.919960in}{2.757087in}}%
\pgfpathlineto{\pgfqpoint{0.920012in}{2.640000in}}%
\pgfpathlineto{\pgfqpoint{0.920737in}{2.709627in}}%
\pgfpathlineto{\pgfqpoint{0.920786in}{2.683668in}}%
\pgfpathlineto{\pgfqpoint{0.921330in}{2.746842in}}%
\pgfpathlineto{\pgfqpoint{0.921380in}{2.721016in}}%
\pgfpathlineto{\pgfqpoint{0.921902in}{2.782664in}}%
\pgfpathlineto{\pgfqpoint{0.922021in}{2.644559in}}%
\pgfpathlineto{\pgfqpoint{0.922032in}{2.644559in}}%
\pgfpathlineto{\pgfqpoint{0.924903in}{2.644909in}}%
\pgfpathlineto{\pgfqpoint{0.924905in}{2.735136in}}%
\pgfpathlineto{\pgfqpoint{0.924968in}{2.616683in}}%
\pgfpathlineto{\pgfqpoint{0.925651in}{2.656338in}}%
\pgfpathlineto{\pgfqpoint{0.926492in}{2.779208in}}%
\pgfpathlineto{\pgfqpoint{0.926874in}{2.803171in}}%
\pgfpathlineto{\pgfqpoint{0.927362in}{2.644559in}}%
\pgfpathlineto{\pgfqpoint{0.929848in}{2.644563in}}%
\pgfpathlineto{\pgfqpoint{0.929850in}{2.741905in}}%
\pgfpathlineto{\pgfqpoint{0.929918in}{2.622433in}}%
\pgfpathlineto{\pgfqpoint{0.930653in}{2.726246in}}%
\pgfpathlineto{\pgfqpoint{0.930703in}{2.669164in}}%
\pgfpathlineto{\pgfqpoint{0.931346in}{2.769596in}}%
\pgfpathlineto{\pgfqpoint{0.931395in}{2.712803in}}%
\pgfpathlineto{\pgfqpoint{0.931820in}{2.799267in}}%
\pgfpathlineto{\pgfqpoint{0.931904in}{2.644559in}}%
\pgfpathlineto{\pgfqpoint{0.932111in}{2.644559in}}%
\pgfpathlineto{\pgfqpoint{0.934793in}{2.644562in}}%
\pgfpathlineto{\pgfqpoint{0.934795in}{2.744929in}}%
\pgfpathlineto{\pgfqpoint{0.934849in}{2.620886in}}%
\pgfpathlineto{\pgfqpoint{0.935575in}{2.725894in}}%
\pgfpathlineto{\pgfqpoint{0.935624in}{2.664740in}}%
\pgfpathlineto{\pgfqpoint{0.936267in}{2.769258in}}%
\pgfpathlineto{\pgfqpoint{0.936317in}{2.708366in}}%
\pgfpathlineto{\pgfqpoint{0.936739in}{2.798784in}}%
\pgfpathlineto{\pgfqpoint{0.936856in}{2.644559in}}%
\pgfpathlineto{\pgfqpoint{0.937009in}{2.644559in}}%
\pgfpathlineto{\pgfqpoint{0.939739in}{2.645030in}}%
\pgfpathlineto{\pgfqpoint{0.939741in}{2.728565in}}%
\pgfpathlineto{\pgfqpoint{0.939829in}{2.617413in}}%
\pgfpathlineto{\pgfqpoint{0.940465in}{2.723550in}}%
\pgfpathlineto{\pgfqpoint{0.940515in}{2.659514in}}%
\pgfpathlineto{\pgfqpoint{0.941158in}{2.766884in}}%
\pgfpathlineto{\pgfqpoint{0.941207in}{2.703170in}}%
\pgfpathlineto{\pgfqpoint{0.941711in}{2.801485in}}%
\pgfpathlineto{\pgfqpoint{0.941797in}{2.644559in}}%
\pgfpathlineto{\pgfqpoint{0.941930in}{2.644559in}}%
\pgfpathlineto{\pgfqpoint{0.944684in}{2.644565in}}%
\pgfpathlineto{\pgfqpoint{0.944686in}{2.753100in}}%
\pgfpathlineto{\pgfqpoint{0.944773in}{2.633063in}}%
\pgfpathlineto{\pgfqpoint{0.945459in}{2.674946in}}%
\pgfpathlineto{\pgfqpoint{0.946300in}{2.765257in}}%
\pgfpathlineto{\pgfqpoint{0.946656in}{2.787578in}}%
\pgfpathlineto{\pgfqpoint{0.947213in}{2.644559in}}%
\pgfpathlineto{\pgfqpoint{0.949627in}{2.644559in}}%
\pgfpathlineto{\pgfqpoint{0.949637in}{2.758503in}}%
\pgfpathlineto{\pgfqpoint{0.949746in}{2.630323in}}%
\pgfpathlineto{\pgfqpoint{0.950389in}{2.711679in}}%
\pgfpathlineto{\pgfqpoint{0.950438in}{2.673519in}}%
\pgfpathlineto{\pgfqpoint{0.951081in}{2.755108in}}%
\pgfpathlineto{\pgfqpoint{0.951131in}{2.717080in}}%
\pgfpathlineto{\pgfqpoint{0.951564in}{2.785385in}}%
\pgfpathlineto{\pgfqpoint{0.951693in}{2.644559in}}%
\pgfpathlineto{\pgfqpoint{0.951803in}{2.644559in}}%
\pgfpathlineto{\pgfqpoint{0.954576in}{2.645260in}}%
\pgfpathlineto{\pgfqpoint{0.954577in}{2.747677in}}%
\pgfpathlineto{\pgfqpoint{0.954699in}{2.642546in}}%
\pgfpathlineto{\pgfqpoint{0.955342in}{2.700955in}}%
\pgfpathlineto{\pgfqpoint{0.955391in}{2.685840in}}%
\pgfpathlineto{\pgfqpoint{0.955935in}{2.738192in}}%
\pgfpathlineto{\pgfqpoint{0.955985in}{2.723166in}}%
\pgfpathlineto{\pgfqpoint{0.956513in}{2.774450in}}%
\pgfpathlineto{\pgfqpoint{0.956601in}{2.714446in}}%
\pgfpathlineto{\pgfqpoint{0.956615in}{2.722576in}}%
\pgfpathlineto{\pgfqpoint{0.957357in}{2.644559in}}%
\pgfpathlineto{\pgfqpoint{0.959521in}{2.645263in}}%
\pgfpathlineto{\pgfqpoint{0.959523in}{2.739848in}}%
\pgfpathlineto{\pgfqpoint{0.959582in}{2.618074in}}%
\pgfpathlineto{\pgfqpoint{0.960262in}{2.656915in}}%
\pgfpathlineto{\pgfqpoint{0.961103in}{2.777869in}}%
\pgfpathlineto{\pgfqpoint{0.961493in}{2.802260in}}%
\pgfpathlineto{\pgfqpoint{0.961958in}{2.644559in}}%
\pgfpathlineto{\pgfqpoint{0.964464in}{2.644559in}}%
\pgfpathlineto{\pgfqpoint{0.964473in}{2.724527in}}%
\pgfpathlineto{\pgfqpoint{0.965231in}{2.689321in}}%
\pgfpathlineto{\pgfqpoint{0.966438in}{2.768101in}}%
\pgfpathlineto{\pgfqpoint{0.966438in}{2.742653in}}%
\pgfpathlineto{\pgfqpoint{0.966460in}{2.750701in}}%
\pgfpathlineto{\pgfqpoint{0.967175in}{2.644559in}}%
\pgfpathlineto{\pgfqpoint{0.969411in}{2.644562in}}%
\pgfpathlineto{\pgfqpoint{0.969414in}{2.743180in}}%
\pgfpathlineto{\pgfqpoint{0.970223in}{2.693915in}}%
\pgfpathlineto{\pgfqpoint{0.971347in}{2.767819in}}%
\pgfpathlineto{\pgfqpoint{0.971384in}{2.744727in}}%
\pgfpathlineto{\pgfqpoint{0.971411in}{2.754682in}}%
\pgfpathlineto{\pgfqpoint{0.972139in}{2.644559in}}%
\pgfpathlineto{\pgfqpoint{0.974357in}{2.644559in}}%
\pgfpathlineto{\pgfqpoint{0.974360in}{2.766714in}}%
\pgfpathlineto{\pgfqpoint{0.974458in}{2.604616in}}%
\pgfpathlineto{\pgfqpoint{0.975097in}{2.733490in}}%
\pgfpathlineto{\pgfqpoint{0.975146in}{2.647359in}}%
\pgfpathlineto{\pgfqpoint{0.975789in}{2.776757in}}%
\pgfpathlineto{\pgfqpoint{0.975839in}{2.691082in}}%
\pgfpathlineto{\pgfqpoint{0.976284in}{2.807661in}}%
\pgfpathlineto{\pgfqpoint{0.976416in}{2.644559in}}%
\pgfpathlineto{\pgfqpoint{0.976545in}{2.644559in}}%
\pgfpathlineto{\pgfqpoint{0.979303in}{2.645062in}}%
\pgfpathlineto{\pgfqpoint{0.979305in}{2.764208in}}%
\pgfpathlineto{\pgfqpoint{0.979368in}{2.635740in}}%
\pgfpathlineto{\pgfqpoint{0.980053in}{2.675429in}}%
\pgfpathlineto{\pgfqpoint{0.980894in}{2.759031in}}%
\pgfpathlineto{\pgfqpoint{0.981274in}{2.782872in}}%
\pgfpathlineto{\pgfqpoint{0.981836in}{2.644559in}}%
\pgfpathlineto{\pgfqpoint{0.984246in}{2.644559in}}%
\pgfpathlineto{\pgfqpoint{0.984252in}{2.726741in}}%
\pgfpathlineto{\pgfqpoint{0.984374in}{2.640691in}}%
\pgfpathlineto{\pgfqpoint{0.985017in}{2.699330in}}%
\pgfpathlineto{\pgfqpoint{0.985066in}{2.683989in}}%
\pgfpathlineto{\pgfqpoint{0.985610in}{2.736568in}}%
\pgfpathlineto{\pgfqpoint{0.985660in}{2.721315in}}%
\pgfpathlineto{\pgfqpoint{0.986187in}{2.772752in}}%
\pgfpathlineto{\pgfqpoint{0.986274in}{2.712459in}}%
\pgfpathlineto{\pgfqpoint{0.986288in}{2.720850in}}%
\pgfpathlineto{\pgfqpoint{0.987043in}{2.644559in}}%
\pgfpathlineto{\pgfqpoint{0.989191in}{2.644559in}}%
\pgfpathlineto{\pgfqpoint{0.989196in}{2.768163in}}%
\pgfpathlineto{\pgfqpoint{0.989253in}{2.636118in}}%
\pgfpathlineto{\pgfqpoint{0.990082in}{2.714945in}}%
\pgfpathlineto{\pgfqpoint{0.990131in}{2.686972in}}%
\pgfpathlineto{\pgfqpoint{0.990774in}{2.758374in}}%
\pgfpathlineto{\pgfqpoint{0.990824in}{2.730531in}}%
\pgfpathlineto{\pgfqpoint{0.991165in}{2.782914in}}%
\pgfpathlineto{\pgfqpoint{0.991255in}{2.644559in}}%
\pgfpathlineto{\pgfqpoint{0.991426in}{2.644559in}}%
\pgfpathlineto{\pgfqpoint{0.994137in}{2.644559in}}%
\pgfpathlineto{\pgfqpoint{0.994143in}{2.738298in}}%
\pgfpathlineto{\pgfqpoint{0.994228in}{2.624437in}}%
\pgfpathlineto{\pgfqpoint{0.994865in}{2.711779in}}%
\pgfpathlineto{\pgfqpoint{0.994914in}{2.666314in}}%
\pgfpathlineto{\pgfqpoint{0.995557in}{2.755145in}}%
\pgfpathlineto{\pgfqpoint{0.995606in}{2.709939in}}%
\pgfpathlineto{\pgfqpoint{0.996111in}{2.789807in}}%
\pgfpathlineto{\pgfqpoint{0.996198in}{2.644559in}}%
\pgfpathlineto{\pgfqpoint{0.996290in}{2.644559in}}%
\pgfpathlineto{\pgfqpoint{0.999084in}{2.644562in}}%
\pgfpathlineto{\pgfqpoint{0.999086in}{2.746349in}}%
\pgfpathlineto{\pgfqpoint{0.999178in}{2.630564in}}%
\pgfpathlineto{\pgfqpoint{0.999865in}{2.672735in}}%
\pgfpathlineto{\pgfqpoint{1.000706in}{2.763902in}}%
\pgfpathlineto{\pgfqpoint{1.001056in}{2.785870in}}%
\pgfpathlineto{\pgfqpoint{1.001618in}{2.644559in}}%
\pgfpathlineto{\pgfqpoint{1.004028in}{2.644559in}}%
\pgfpathlineto{\pgfqpoint{1.004034in}{2.737906in}}%
\pgfpathlineto{\pgfqpoint{1.004123in}{2.626188in}}%
\pgfpathlineto{\pgfqpoint{1.004760in}{2.709043in}}%
\pgfpathlineto{\pgfqpoint{1.004809in}{2.668292in}}%
\pgfpathlineto{\pgfqpoint{1.005452in}{2.752417in}}%
\pgfpathlineto{\pgfqpoint{1.005502in}{2.711909in}}%
\pgfpathlineto{\pgfqpoint{1.006002in}{2.786826in}}%
\pgfpathlineto{\pgfqpoint{1.006093in}{2.644559in}}%
\pgfpathlineto{\pgfqpoint{1.006196in}{2.644559in}}%
\pgfpathlineto{\pgfqpoint{1.008974in}{2.644559in}}%
\pgfpathlineto{\pgfqpoint{1.008980in}{2.735313in}}%
\pgfpathlineto{\pgfqpoint{1.009065in}{2.627799in}}%
\pgfpathlineto{\pgfqpoint{1.009701in}{2.709261in}}%
\pgfpathlineto{\pgfqpoint{1.009751in}{2.669675in}}%
\pgfpathlineto{\pgfqpoint{1.010394in}{2.752643in}}%
\pgfpathlineto{\pgfqpoint{1.010443in}{2.713283in}}%
\pgfpathlineto{\pgfqpoint{1.010947in}{2.787310in}}%
\pgfpathlineto{\pgfqpoint{1.011033in}{2.644559in}}%
\pgfpathlineto{\pgfqpoint{1.011117in}{2.644559in}}%
\pgfpathlineto{\pgfqpoint{1.013919in}{2.644559in}}%
\pgfpathlineto{\pgfqpoint{1.013925in}{2.738209in}}%
\pgfpathlineto{\pgfqpoint{1.014014in}{2.625208in}}%
\pgfpathlineto{\pgfqpoint{1.014651in}{2.710254in}}%
\pgfpathlineto{\pgfqpoint{1.014701in}{2.667340in}}%
\pgfpathlineto{\pgfqpoint{1.015343in}{2.753624in}}%
\pgfpathlineto{\pgfqpoint{1.015393in}{2.710961in}}%
\pgfpathlineto{\pgfqpoint{1.015893in}{2.788000in}}%
\pgfpathlineto{\pgfqpoint{1.015981in}{2.644559in}}%
\pgfpathlineto{\pgfqpoint{1.016068in}{2.644559in}}%
\pgfpathlineto{\pgfqpoint{1.018864in}{2.644559in}}%
\pgfpathlineto{\pgfqpoint{1.018870in}{2.734986in}}%
\pgfpathlineto{\pgfqpoint{1.018954in}{2.631909in}}%
\pgfpathlineto{\pgfqpoint{1.019591in}{2.703835in}}%
\pgfpathlineto{\pgfqpoint{1.019640in}{2.673696in}}%
\pgfpathlineto{\pgfqpoint{1.020283in}{2.747235in}}%
\pgfpathlineto{\pgfqpoint{1.020333in}{2.717286in}}%
\pgfpathlineto{\pgfqpoint{1.020838in}{2.781997in}}%
\pgfpathlineto{\pgfqpoint{1.020937in}{2.644559in}}%
\pgfpathlineto{\pgfqpoint{1.020976in}{2.644559in}}%
\pgfpathlineto{\pgfqpoint{1.023810in}{2.644559in}}%
\pgfpathlineto{\pgfqpoint{1.023816in}{2.736815in}}%
\pgfpathlineto{\pgfqpoint{1.023905in}{2.629440in}}%
\pgfpathlineto{\pgfqpoint{1.024542in}{2.704659in}}%
\pgfpathlineto{\pgfqpoint{1.024592in}{2.671576in}}%
\pgfpathlineto{\pgfqpoint{1.025235in}{2.748048in}}%
\pgfpathlineto{\pgfqpoint{1.025284in}{2.715177in}}%
\pgfpathlineto{\pgfqpoint{1.025784in}{2.782417in}}%
\pgfpathlineto{\pgfqpoint{1.025870in}{2.644559in}}%
\pgfpathlineto{\pgfqpoint{1.025939in}{2.644559in}}%
\pgfpathlineto{\pgfqpoint{1.028755in}{2.644559in}}%
\pgfpathlineto{\pgfqpoint{1.028761in}{2.730224in}}%
\pgfpathlineto{\pgfqpoint{1.028815in}{2.639665in}}%
\pgfpathlineto{\pgfqpoint{1.029492in}{2.677464in}}%
\pgfpathlineto{\pgfqpoint{1.030432in}{2.755835in}}%
\pgfpathlineto{\pgfqpoint{1.030704in}{2.772903in}}%
\pgfpathlineto{\pgfqpoint{1.031521in}{2.644559in}}%
\pgfpathlineto{\pgfqpoint{1.033701in}{2.644559in}}%
\pgfpathlineto{\pgfqpoint{1.033707in}{2.742584in}}%
\pgfpathlineto{\pgfqpoint{1.033761in}{2.618457in}}%
\pgfpathlineto{\pgfqpoint{1.034437in}{2.656149in}}%
\pgfpathlineto{\pgfqpoint{1.035278in}{2.768577in}}%
\pgfpathlineto{\pgfqpoint{1.035649in}{2.791820in}}%
\pgfpathlineto{\pgfqpoint{1.036116in}{2.644559in}}%
\pgfpathlineto{\pgfqpoint{1.038646in}{2.644559in}}%
\pgfpathlineto{\pgfqpoint{1.038652in}{2.741151in}}%
\pgfpathlineto{\pgfqpoint{1.038705in}{2.619448in}}%
\pgfpathlineto{\pgfqpoint{1.039381in}{2.656915in}}%
\pgfpathlineto{\pgfqpoint{1.040221in}{2.769220in}}%
\pgfpathlineto{\pgfqpoint{1.040594in}{2.792533in}}%
\pgfpathlineto{\pgfqpoint{1.041055in}{2.644559in}}%
\pgfpathlineto{\pgfqpoint{1.043592in}{2.644559in}}%
\pgfpathlineto{\pgfqpoint{1.043598in}{2.741153in}}%
\pgfpathlineto{\pgfqpoint{1.043650in}{2.620051in}}%
\pgfpathlineto{\pgfqpoint{1.044325in}{2.657366in}}%
\pgfpathlineto{\pgfqpoint{1.045166in}{2.768731in}}%
\pgfpathlineto{\pgfqpoint{1.045539in}{2.792084in}}%
\pgfpathlineto{\pgfqpoint{1.046018in}{2.644559in}}%
\pgfpathlineto{\pgfqpoint{1.048537in}{2.644559in}}%
\pgfpathlineto{\pgfqpoint{1.048543in}{2.744988in}}%
\pgfpathlineto{\pgfqpoint{1.048670in}{2.620312in}}%
\pgfpathlineto{\pgfqpoint{1.049312in}{2.718125in}}%
\pgfpathlineto{\pgfqpoint{1.049362in}{2.663703in}}%
\pgfpathlineto{\pgfqpoint{1.050005in}{2.761504in}}%
\pgfpathlineto{\pgfqpoint{1.050054in}{2.707314in}}%
\pgfpathlineto{\pgfqpoint{1.050480in}{2.791287in}}%
\pgfpathlineto{\pgfqpoint{1.050594in}{2.644559in}}%
\pgfpathlineto{\pgfqpoint{1.050733in}{2.644559in}}%
\pgfpathlineto{\pgfqpoint{1.053482in}{2.644559in}}%
\pgfpathlineto{\pgfqpoint{1.053489in}{2.746408in}}%
\pgfpathlineto{\pgfqpoint{1.053602in}{2.639423in}}%
\pgfpathlineto{\pgfqpoint{1.054244in}{2.697653in}}%
\pgfpathlineto{\pgfqpoint{1.054294in}{2.682573in}}%
\pgfpathlineto{\pgfqpoint{1.054838in}{2.734913in}}%
\pgfpathlineto{\pgfqpoint{1.054887in}{2.719877in}}%
\pgfpathlineto{\pgfqpoint{1.055419in}{2.771400in}}%
\pgfpathlineto{\pgfqpoint{1.055511in}{2.711531in}}%
\pgfpathlineto{\pgfqpoint{1.055525in}{2.720019in}}%
\pgfpathlineto{\pgfqpoint{1.056272in}{2.644559in}}%
\pgfpathlineto{\pgfqpoint{1.058428in}{2.644559in}}%
\pgfpathlineto{\pgfqpoint{1.058435in}{2.746158in}}%
\pgfpathlineto{\pgfqpoint{1.058563in}{2.620590in}}%
\pgfpathlineto{\pgfqpoint{1.059206in}{2.715276in}}%
\pgfpathlineto{\pgfqpoint{1.059255in}{2.663993in}}%
\pgfpathlineto{\pgfqpoint{1.059898in}{2.758660in}}%
\pgfpathlineto{\pgfqpoint{1.059947in}{2.707600in}}%
\pgfpathlineto{\pgfqpoint{1.060372in}{2.788379in}}%
\pgfpathlineto{\pgfqpoint{1.060486in}{2.644559in}}%
\pgfpathlineto{\pgfqpoint{1.060639in}{2.644559in}}%
\pgfpathlineto{\pgfqpoint{1.063373in}{2.644559in}}%
\pgfpathlineto{\pgfqpoint{1.063380in}{2.744403in}}%
\pgfpathlineto{\pgfqpoint{1.063495in}{2.627488in}}%
\pgfpathlineto{\pgfqpoint{1.064138in}{2.709722in}}%
\pgfpathlineto{\pgfqpoint{1.064187in}{2.670707in}}%
\pgfpathlineto{\pgfqpoint{1.064830in}{2.753134in}}%
\pgfpathlineto{\pgfqpoint{1.064879in}{2.714285in}}%
\pgfpathlineto{\pgfqpoint{1.065311in}{2.783296in}}%
\pgfpathlineto{\pgfqpoint{1.065440in}{2.644559in}}%
\pgfpathlineto{\pgfqpoint{1.065547in}{2.644559in}}%
\pgfpathlineto{\pgfqpoint{1.068319in}{2.644559in}}%
\pgfpathlineto{\pgfqpoint{1.068326in}{2.746691in}}%
\pgfpathlineto{\pgfqpoint{1.068395in}{2.644384in}}%
\pgfpathlineto{\pgfqpoint{1.069080in}{2.684451in}}%
\pgfpathlineto{\pgfqpoint{1.070293in}{2.766019in}}%
\pgfpathlineto{\pgfqpoint{1.071060in}{2.644559in}}%
\pgfpathlineto{\pgfqpoint{1.073264in}{2.644559in}}%
\pgfpathlineto{\pgfqpoint{1.073271in}{2.746159in}}%
\pgfpathlineto{\pgfqpoint{1.073399in}{2.620564in}}%
\pgfpathlineto{\pgfqpoint{1.074042in}{2.715362in}}%
\pgfpathlineto{\pgfqpoint{1.074091in}{2.663963in}}%
\pgfpathlineto{\pgfqpoint{1.074734in}{2.758745in}}%
\pgfpathlineto{\pgfqpoint{1.074783in}{2.707570in}}%
\pgfpathlineto{\pgfqpoint{1.075209in}{2.788474in}}%
\pgfpathlineto{\pgfqpoint{1.075322in}{2.644559in}}%
\pgfpathlineto{\pgfqpoint{1.075479in}{2.644559in}}%
\pgfpathlineto{\pgfqpoint{1.078210in}{2.644559in}}%
\pgfpathlineto{\pgfqpoint{1.078216in}{2.746941in}}%
\pgfpathlineto{\pgfqpoint{1.078955in}{2.684965in}}%
\pgfpathlineto{\pgfqpoint{1.080184in}{2.765823in}}%
\pgfpathlineto{\pgfqpoint{1.080184in}{2.753151in}}%
\pgfpathlineto{\pgfqpoint{1.080923in}{2.644559in}}%
\pgfpathlineto{\pgfqpoint{1.083155in}{2.644559in}}%
\pgfpathlineto{\pgfqpoint{1.083162in}{2.746180in}}%
\pgfpathlineto{\pgfqpoint{1.083227in}{2.632004in}}%
\pgfpathlineto{\pgfqpoint{1.083910in}{2.671462in}}%
\pgfpathlineto{\pgfqpoint{1.084750in}{2.755174in}}%
\pgfpathlineto{\pgfqpoint{1.085129in}{2.778901in}}%
\pgfpathlineto{\pgfqpoint{1.085694in}{2.644559in}}%
\pgfpathlineto{\pgfqpoint{1.088101in}{2.644559in}}%
\pgfpathlineto{\pgfqpoint{1.088108in}{2.746211in}}%
\pgfpathlineto{\pgfqpoint{1.088828in}{2.683676in}}%
\pgfpathlineto{\pgfqpoint{1.090074in}{2.762045in}}%
\pgfpathlineto{\pgfqpoint{1.090075in}{2.733876in}}%
\pgfpathlineto{\pgfqpoint{1.090095in}{2.745896in}}%
\pgfpathlineto{\pgfqpoint{1.090810in}{2.644559in}}%
\pgfpathlineto{\pgfqpoint{1.093046in}{2.644559in}}%
\pgfpathlineto{\pgfqpoint{1.093053in}{2.745822in}}%
\pgfpathlineto{\pgfqpoint{1.093109in}{2.642456in}}%
\pgfpathlineto{\pgfqpoint{1.093786in}{2.680052in}}%
\pgfpathlineto{\pgfqpoint{1.094825in}{2.755621in}}%
\pgfpathlineto{\pgfqpoint{1.095020in}{2.767862in}}%
\pgfpathlineto{\pgfqpoint{1.096324in}{2.644559in}}%
\pgfpathlineto{\pgfqpoint{1.097992in}{2.644559in}}%
\pgfpathlineto{\pgfqpoint{1.097999in}{2.745358in}}%
\pgfpathlineto{\pgfqpoint{1.098065in}{2.634516in}}%
\pgfpathlineto{\pgfqpoint{1.098748in}{2.674043in}}%
\pgfpathlineto{\pgfqpoint{1.099588in}{2.751484in}}%
\pgfpathlineto{\pgfqpoint{1.099965in}{2.775115in}}%
\pgfpathlineto{\pgfqpoint{1.100558in}{2.644559in}}%
\pgfpathlineto{\pgfqpoint{1.102937in}{2.644559in}}%
\pgfpathlineto{\pgfqpoint{1.102944in}{2.745851in}}%
\pgfpathlineto{\pgfqpoint{1.103010in}{2.631746in}}%
\pgfpathlineto{\pgfqpoint{1.103692in}{2.671276in}}%
\pgfpathlineto{\pgfqpoint{1.104533in}{2.755013in}}%
\pgfpathlineto{\pgfqpoint{1.104911in}{2.778678in}}%
\pgfpathlineto{\pgfqpoint{1.105478in}{2.644559in}}%
\pgfpathlineto{\pgfqpoint{1.107882in}{2.644559in}}%
\pgfpathlineto{\pgfqpoint{1.107890in}{2.742686in}}%
\pgfpathlineto{\pgfqpoint{1.107964in}{2.626168in}}%
\pgfpathlineto{\pgfqpoint{1.108649in}{2.666645in}}%
\pgfpathlineto{\pgfqpoint{1.109489in}{2.758782in}}%
\pgfpathlineto{\pgfqpoint{1.109856in}{2.781775in}}%
\pgfpathlineto{\pgfqpoint{1.110369in}{2.644559in}}%
\pgfpathlineto{\pgfqpoint{1.112828in}{2.644559in}}%
\pgfpathlineto{\pgfqpoint{1.112835in}{2.743351in}}%
\pgfpathlineto{\pgfqpoint{1.112909in}{2.623778in}}%
\pgfpathlineto{\pgfqpoint{1.113593in}{2.664275in}}%
\pgfpathlineto{\pgfqpoint{1.114434in}{2.762456in}}%
\pgfpathlineto{\pgfqpoint{1.114802in}{2.785477in}}%
\pgfpathlineto{\pgfqpoint{1.115318in}{2.644559in}}%
\pgfpathlineto{\pgfqpoint{1.117773in}{2.644559in}}%
\pgfpathlineto{\pgfqpoint{1.117781in}{2.742246in}}%
\pgfpathlineto{\pgfqpoint{1.117914in}{2.620781in}}%
\pgfpathlineto{\pgfqpoint{1.118507in}{2.657975in}}%
\pgfpathlineto{\pgfqpoint{1.119348in}{2.764617in}}%
\pgfpathlineto{\pgfqpoint{1.119721in}{2.787949in}}%
\pgfpathlineto{\pgfqpoint{1.120228in}{2.644559in}}%
\pgfpathlineto{\pgfqpoint{1.122719in}{2.644559in}}%
\pgfpathlineto{\pgfqpoint{1.122726in}{2.743086in}}%
\pgfpathlineto{\pgfqpoint{1.122789in}{2.632337in}}%
\pgfpathlineto{\pgfqpoint{1.123470in}{2.671139in}}%
\pgfpathlineto{\pgfqpoint{1.124311in}{2.752216in}}%
\pgfpathlineto{\pgfqpoint{1.124693in}{2.776170in}}%
\pgfpathlineto{\pgfqpoint{1.125254in}{2.644559in}}%
\pgfpathlineto{\pgfqpoint{1.127664in}{2.644559in}}%
\pgfpathlineto{\pgfqpoint{1.127671in}{2.740871in}}%
\pgfpathlineto{\pgfqpoint{1.127730in}{2.615274in}}%
\pgfpathlineto{\pgfqpoint{1.128407in}{2.652998in}}%
\pgfpathlineto{\pgfqpoint{1.129247in}{2.770940in}}%
\pgfpathlineto{\pgfqpoint{1.129638in}{2.795413in}}%
\pgfpathlineto{\pgfqpoint{1.130072in}{2.644559in}}%
\pgfpathlineto{\pgfqpoint{1.132610in}{2.644559in}}%
\pgfpathlineto{\pgfqpoint{1.132617in}{2.741548in}}%
\pgfpathlineto{\pgfqpoint{1.132683in}{2.612115in}}%
\pgfpathlineto{\pgfqpoint{1.133366in}{2.651520in}}%
\pgfpathlineto{\pgfqpoint{1.134206in}{2.772363in}}%
\pgfpathlineto{\pgfqpoint{1.134584in}{2.795973in}}%
\pgfpathlineto{\pgfqpoint{1.135047in}{2.644559in}}%
\pgfpathlineto{\pgfqpoint{1.137555in}{2.644559in}}%
\pgfpathlineto{\pgfqpoint{1.137562in}{2.743220in}}%
\pgfpathlineto{\pgfqpoint{1.137631in}{2.628119in}}%
\pgfpathlineto{\pgfqpoint{1.138316in}{2.668041in}}%
\pgfpathlineto{\pgfqpoint{1.139157in}{2.756940in}}%
\pgfpathlineto{\pgfqpoint{1.139529in}{2.780283in}}%
\pgfpathlineto{\pgfqpoint{1.140091in}{2.644559in}}%
\pgfpathlineto{\pgfqpoint{1.142501in}{2.644559in}}%
\pgfpathlineto{\pgfqpoint{1.142513in}{2.770218in}}%
\pgfpathlineto{\pgfqpoint{1.142623in}{2.639864in}}%
\pgfpathlineto{\pgfqpoint{1.143266in}{2.693368in}}%
\pgfpathlineto{\pgfqpoint{1.143315in}{2.682933in}}%
\pgfpathlineto{\pgfqpoint{1.143859in}{2.730615in}}%
\pgfpathlineto{\pgfqpoint{1.143909in}{2.720250in}}%
\pgfpathlineto{\pgfqpoint{1.144439in}{2.766993in}}%
\pgfpathlineto{\pgfqpoint{1.144538in}{2.695898in}}%
\pgfpathlineto{\pgfqpoint{1.145305in}{2.644559in}}%
\pgfpathlineto{\pgfqpoint{1.147446in}{2.644559in}}%
\pgfpathlineto{\pgfqpoint{1.147458in}{2.771005in}}%
\pgfpathlineto{\pgfqpoint{1.147526in}{2.635846in}}%
\pgfpathlineto{\pgfqpoint{1.148211in}{2.675867in}}%
\pgfpathlineto{\pgfqpoint{1.149052in}{2.749351in}}%
\pgfpathlineto{\pgfqpoint{1.149420in}{2.772430in}}%
\pgfpathlineto{\pgfqpoint{1.150080in}{2.644559in}}%
\pgfpathlineto{\pgfqpoint{1.152392in}{2.644559in}}%
\pgfpathlineto{\pgfqpoint{1.152403in}{2.768108in}}%
\pgfpathlineto{\pgfqpoint{1.152491in}{2.622805in}}%
\pgfpathlineto{\pgfqpoint{1.153127in}{2.708452in}}%
\pgfpathlineto{\pgfqpoint{1.153177in}{2.664663in}}%
\pgfpathlineto{\pgfqpoint{1.153820in}{2.751826in}}%
\pgfpathlineto{\pgfqpoint{1.153869in}{2.708280in}}%
\pgfpathlineto{\pgfqpoint{1.154365in}{2.785977in}}%
\pgfpathlineto{\pgfqpoint{1.154456in}{2.644559in}}%
\pgfpathlineto{\pgfqpoint{1.154552in}{2.644559in}}%
\pgfpathlineto{\pgfqpoint{1.157337in}{2.644559in}}%
\pgfpathlineto{\pgfqpoint{1.157349in}{2.768074in}}%
\pgfpathlineto{\pgfqpoint{1.157427in}{2.630160in}}%
\pgfpathlineto{\pgfqpoint{1.158063in}{2.700888in}}%
\pgfpathlineto{\pgfqpoint{1.158112in}{2.671254in}}%
\pgfpathlineto{\pgfqpoint{1.158755in}{2.744294in}}%
\pgfpathlineto{\pgfqpoint{1.158804in}{2.714839in}}%
\pgfpathlineto{\pgfqpoint{1.159311in}{2.779122in}}%
\pgfpathlineto{\pgfqpoint{1.159403in}{2.644559in}}%
\pgfpathlineto{\pgfqpoint{1.159476in}{2.644559in}}%
\pgfpathlineto{\pgfqpoint{1.162282in}{2.644559in}}%
\pgfpathlineto{\pgfqpoint{1.162294in}{2.764997in}}%
\pgfpathlineto{\pgfqpoint{1.162411in}{2.639759in}}%
\pgfpathlineto{\pgfqpoint{1.163054in}{2.694899in}}%
\pgfpathlineto{\pgfqpoint{1.163104in}{2.682922in}}%
\pgfpathlineto{\pgfqpoint{1.163648in}{2.732149in}}%
\pgfpathlineto{\pgfqpoint{1.163697in}{2.720235in}}%
\pgfpathlineto{\pgfqpoint{1.164224in}{2.768312in}}%
\pgfpathlineto{\pgfqpoint{1.164312in}{2.688800in}}%
\pgfpathlineto{\pgfqpoint{1.164327in}{2.738802in}}%
\pgfpathlineto{\pgfqpoint{1.165074in}{2.644559in}}%
\pgfpathlineto{\pgfqpoint{1.167228in}{2.644559in}}%
\pgfpathlineto{\pgfqpoint{1.167240in}{2.769675in}}%
\pgfpathlineto{\pgfqpoint{1.167326in}{2.622833in}}%
\pgfpathlineto{\pgfqpoint{1.167962in}{2.707507in}}%
\pgfpathlineto{\pgfqpoint{1.168011in}{2.664558in}}%
\pgfpathlineto{\pgfqpoint{1.168654in}{2.750885in}}%
\pgfpathlineto{\pgfqpoint{1.168704in}{2.708172in}}%
\pgfpathlineto{\pgfqpoint{1.169202in}{2.785160in}}%
\pgfpathlineto{\pgfqpoint{1.169291in}{2.644559in}}%
\pgfpathlineto{\pgfqpoint{1.169391in}{2.644559in}}%
\pgfpathlineto{\pgfqpoint{1.172173in}{2.644559in}}%
\pgfpathlineto{\pgfqpoint{1.172185in}{2.764866in}}%
\pgfpathlineto{\pgfqpoint{1.172247in}{2.632514in}}%
\pgfpathlineto{\pgfqpoint{1.172928in}{2.671199in}}%
\pgfpathlineto{\pgfqpoint{1.173768in}{2.753929in}}%
\pgfpathlineto{\pgfqpoint{1.174147in}{2.777677in}}%
\pgfpathlineto{\pgfqpoint{1.174707in}{2.644559in}}%
\pgfpathlineto{\pgfqpoint{1.177119in}{2.644559in}}%
\pgfpathlineto{\pgfqpoint{1.177130in}{2.761180in}}%
\pgfpathlineto{\pgfqpoint{1.177244in}{2.639435in}}%
\pgfpathlineto{\pgfqpoint{1.177886in}{2.695125in}}%
\pgfpathlineto{\pgfqpoint{1.177935in}{2.682485in}}%
\pgfpathlineto{\pgfqpoint{1.178479in}{2.732384in}}%
\pgfpathlineto{\pgfqpoint{1.178529in}{2.719790in}}%
\pgfpathlineto{\pgfqpoint{1.179058in}{2.768704in}}%
\pgfpathlineto{\pgfqpoint{1.179155in}{2.692730in}}%
\pgfpathlineto{\pgfqpoint{1.179923in}{2.644559in}}%
\pgfpathlineto{\pgfqpoint{1.182064in}{2.644559in}}%
\pgfpathlineto{\pgfqpoint{1.182075in}{2.758208in}}%
\pgfpathlineto{\pgfqpoint{1.182145in}{2.617348in}}%
\pgfpathlineto{\pgfqpoint{1.182830in}{2.657347in}}%
\pgfpathlineto{\pgfqpoint{1.183670in}{2.770991in}}%
\pgfpathlineto{\pgfqpoint{1.184038in}{2.794004in}}%
\pgfpathlineto{\pgfqpoint{1.184518in}{2.644559in}}%
\pgfpathlineto{\pgfqpoint{1.187010in}{2.644559in}}%
\pgfpathlineto{\pgfqpoint{1.187021in}{2.759706in}}%
\pgfpathlineto{\pgfqpoint{1.187124in}{2.644359in}}%
\pgfpathlineto{\pgfqpoint{1.187761in}{2.688758in}}%
\pgfpathlineto{\pgfqpoint{1.187810in}{2.686799in}}%
\pgfpathlineto{\pgfqpoint{1.188058in}{2.707389in}}%
\pgfpathlineto{\pgfqpoint{1.188107in}{2.705451in}}%
\pgfpathlineto{\pgfqpoint{1.188941in}{2.762846in}}%
\pgfpathlineto{\pgfqpoint{1.189011in}{2.749221in}}%
\pgfpathlineto{\pgfqpoint{1.189792in}{2.644559in}}%
\pgfpathlineto{\pgfqpoint{1.191955in}{2.644559in}}%
\pgfpathlineto{\pgfqpoint{1.191966in}{2.758299in}}%
\pgfpathlineto{\pgfqpoint{1.192095in}{2.637737in}}%
\pgfpathlineto{\pgfqpoint{1.192688in}{2.674815in}}%
\pgfpathlineto{\pgfqpoint{1.193529in}{2.748924in}}%
\pgfpathlineto{\pgfqpoint{1.193902in}{2.772334in}}%
\pgfpathlineto{\pgfqpoint{1.194557in}{2.644559in}}%
\pgfpathlineto{\pgfqpoint{1.196901in}{2.644559in}}%
\pgfpathlineto{\pgfqpoint{1.196912in}{2.758509in}}%
\pgfpathlineto{\pgfqpoint{1.197632in}{2.686529in}}%
\pgfpathlineto{\pgfqpoint{1.197781in}{2.693344in}}%
\pgfpathlineto{\pgfqpoint{1.198874in}{2.764573in}}%
\pgfpathlineto{\pgfqpoint{1.198875in}{2.739474in}}%
\pgfpathlineto{\pgfqpoint{1.198901in}{2.749626in}}%
\pgfpathlineto{\pgfqpoint{1.199618in}{2.644559in}}%
\pgfpathlineto{\pgfqpoint{1.201846in}{2.644559in}}%
\pgfpathlineto{\pgfqpoint{1.201857in}{2.757034in}}%
\pgfpathlineto{\pgfqpoint{1.201925in}{2.633704in}}%
\pgfpathlineto{\pgfqpoint{1.202610in}{2.673250in}}%
\pgfpathlineto{\pgfqpoint{1.203451in}{2.754602in}}%
\pgfpathlineto{\pgfqpoint{1.203820in}{2.777765in}}%
\pgfpathlineto{\pgfqpoint{1.204406in}{2.644559in}}%
\pgfpathlineto{\pgfqpoint{1.206792in}{2.644559in}}%
\pgfpathlineto{\pgfqpoint{1.206805in}{2.738487in}}%
\pgfpathlineto{\pgfqpoint{1.206873in}{2.619050in}}%
\pgfpathlineto{\pgfqpoint{1.207557in}{2.658828in}}%
\pgfpathlineto{\pgfqpoint{1.208397in}{2.770011in}}%
\pgfpathlineto{\pgfqpoint{1.208765in}{2.793039in}}%
\pgfpathlineto{\pgfqpoint{1.209244in}{2.644559in}}%
\pgfpathlineto{\pgfqpoint{1.211736in}{2.644559in}}%
\pgfpathlineto{\pgfqpoint{1.211746in}{2.782093in}}%
\pgfpathlineto{\pgfqpoint{1.212479in}{2.687010in}}%
\pgfpathlineto{\pgfqpoint{1.213711in}{2.765288in}}%
\pgfpathlineto{\pgfqpoint{1.213743in}{2.727428in}}%
\pgfpathlineto{\pgfqpoint{1.214776in}{2.644559in}}%
\pgfpathlineto{\pgfqpoint{1.216684in}{2.644562in}}%
\pgfpathlineto{\pgfqpoint{1.216698in}{2.761418in}}%
\pgfpathlineto{\pgfqpoint{1.217409in}{2.689632in}}%
\pgfpathlineto{\pgfqpoint{1.217458in}{2.687655in}}%
\pgfpathlineto{\pgfqpoint{1.217705in}{2.708281in}}%
\pgfpathlineto{\pgfqpoint{1.217755in}{2.706288in}}%
\pgfpathlineto{\pgfqpoint{1.218656in}{2.768046in}}%
\pgfpathlineto{\pgfqpoint{1.218686in}{2.741840in}}%
\pgfpathlineto{\pgfqpoint{1.219484in}{2.644559in}}%
\pgfpathlineto{\pgfqpoint{1.221627in}{2.644559in}}%
\pgfpathlineto{\pgfqpoint{1.221635in}{2.784300in}}%
\pgfpathlineto{\pgfqpoint{1.221703in}{2.622906in}}%
\pgfpathlineto{\pgfqpoint{1.222433in}{2.722038in}}%
\pgfpathlineto{\pgfqpoint{1.222483in}{2.667799in}}%
\pgfpathlineto{\pgfqpoint{1.223126in}{2.765403in}}%
\pgfpathlineto{\pgfqpoint{1.223175in}{2.711423in}}%
\pgfpathlineto{\pgfqpoint{1.223602in}{2.795214in}}%
\pgfpathlineto{\pgfqpoint{1.223694in}{2.644559in}}%
\pgfpathlineto{\pgfqpoint{1.223889in}{2.644559in}}%
\pgfpathlineto{\pgfqpoint{1.226573in}{2.644559in}}%
\pgfpathlineto{\pgfqpoint{1.226590in}{2.744581in}}%
\pgfpathlineto{\pgfqpoint{1.226646in}{2.621295in}}%
\pgfpathlineto{\pgfqpoint{1.227325in}{2.659381in}}%
\pgfpathlineto{\pgfqpoint{1.228166in}{2.775560in}}%
\pgfpathlineto{\pgfqpoint{1.228547in}{2.799447in}}%
\pgfpathlineto{\pgfqpoint{1.229031in}{2.644559in}}%
\pgfpathlineto{\pgfqpoint{1.231519in}{2.644559in}}%
\pgfpathlineto{\pgfqpoint{1.231533in}{2.747370in}}%
\pgfpathlineto{\pgfqpoint{1.231588in}{2.625528in}}%
\pgfpathlineto{\pgfqpoint{1.232264in}{2.662888in}}%
\pgfpathlineto{\pgfqpoint{1.233104in}{2.773774in}}%
\pgfpathlineto{\pgfqpoint{1.233493in}{2.798090in}}%
\pgfpathlineto{\pgfqpoint{1.233991in}{2.644559in}}%
\pgfpathlineto{\pgfqpoint{1.236465in}{2.644584in}}%
\pgfpathlineto{\pgfqpoint{1.236469in}{2.774346in}}%
\pgfpathlineto{\pgfqpoint{1.236534in}{2.623814in}}%
\pgfpathlineto{\pgfqpoint{1.237210in}{2.661517in}}%
\pgfpathlineto{\pgfqpoint{1.238051in}{2.778471in}}%
\pgfpathlineto{\pgfqpoint{1.238438in}{2.802694in}}%
\pgfpathlineto{\pgfqpoint{1.238920in}{2.644559in}}%
\pgfpathlineto{\pgfqpoint{1.241410in}{2.644559in}}%
\pgfpathlineto{\pgfqpoint{1.241414in}{2.747971in}}%
\pgfpathlineto{\pgfqpoint{1.241489in}{2.621587in}}%
\pgfpathlineto{\pgfqpoint{1.242172in}{2.661464in}}%
\pgfpathlineto{\pgfqpoint{1.243013in}{2.781761in}}%
\pgfpathlineto{\pgfqpoint{1.243384in}{2.804942in}}%
\pgfpathlineto{\pgfqpoint{1.243844in}{2.644559in}}%
\pgfpathlineto{\pgfqpoint{1.246355in}{2.644559in}}%
\pgfpathlineto{\pgfqpoint{1.246369in}{2.758299in}}%
\pgfpathlineto{\pgfqpoint{1.246429in}{2.623232in}}%
\pgfpathlineto{\pgfqpoint{1.247109in}{2.662090in}}%
\pgfpathlineto{\pgfqpoint{1.247950in}{2.782111in}}%
\pgfpathlineto{\pgfqpoint{1.248329in}{2.805830in}}%
\pgfpathlineto{\pgfqpoint{1.248819in}{2.644559in}}%
\pgfpathlineto{\pgfqpoint{1.251301in}{2.644559in}}%
\pgfpathlineto{\pgfqpoint{1.251305in}{2.753752in}}%
\pgfpathlineto{\pgfqpoint{1.251393in}{2.621757in}}%
\pgfpathlineto{\pgfqpoint{1.252029in}{2.727909in}}%
\pgfpathlineto{\pgfqpoint{1.252078in}{2.663144in}}%
\pgfpathlineto{\pgfqpoint{1.252721in}{2.771256in}}%
\pgfpathlineto{\pgfqpoint{1.252770in}{2.706787in}}%
\pgfpathlineto{\pgfqpoint{1.253274in}{2.805881in}}%
\pgfpathlineto{\pgfqpoint{1.253357in}{2.644559in}}%
\pgfpathlineto{\pgfqpoint{1.253481in}{2.644559in}}%
\pgfpathlineto{\pgfqpoint{1.256246in}{2.644559in}}%
\pgfpathlineto{\pgfqpoint{1.256251in}{2.759144in}}%
\pgfpathlineto{\pgfqpoint{1.256314in}{2.632035in}}%
\pgfpathlineto{\pgfqpoint{1.256990in}{2.669676in}}%
\pgfpathlineto{\pgfqpoint{1.257831in}{2.774533in}}%
\pgfpathlineto{\pgfqpoint{1.258220in}{2.798947in}}%
\pgfpathlineto{\pgfqpoint{1.258732in}{2.644559in}}%
\pgfpathlineto{\pgfqpoint{1.261192in}{2.644559in}}%
\pgfpathlineto{\pgfqpoint{1.261196in}{2.760403in}}%
\pgfpathlineto{\pgfqpoint{1.261263in}{2.623450in}}%
\pgfpathlineto{\pgfqpoint{1.261944in}{2.662613in}}%
\pgfpathlineto{\pgfqpoint{1.262785in}{2.783262in}}%
\pgfpathlineto{\pgfqpoint{1.263165in}{2.807054in}}%
\pgfpathlineto{\pgfqpoint{1.263645in}{2.644559in}}%
\pgfpathlineto{\pgfqpoint{1.266137in}{2.644559in}}%
\pgfpathlineto{\pgfqpoint{1.266141in}{2.760433in}}%
\pgfpathlineto{\pgfqpoint{1.266205in}{2.624171in}}%
\pgfpathlineto{\pgfqpoint{1.266885in}{2.662686in}}%
\pgfpathlineto{\pgfqpoint{1.267726in}{2.782565in}}%
\pgfpathlineto{\pgfqpoint{1.268111in}{2.806661in}}%
\pgfpathlineto{\pgfqpoint{1.268592in}{2.644559in}}%
\pgfpathlineto{\pgfqpoint{1.271082in}{2.644559in}}%
\pgfpathlineto{\pgfqpoint{1.271087in}{2.761791in}}%
\pgfpathlineto{\pgfqpoint{1.271151in}{2.624345in}}%
\pgfpathlineto{\pgfqpoint{1.271830in}{2.662842in}}%
\pgfpathlineto{\pgfqpoint{1.272671in}{2.782768in}}%
\pgfpathlineto{\pgfqpoint{1.273056in}{2.806877in}}%
\pgfpathlineto{\pgfqpoint{1.273537in}{2.644559in}}%
\pgfpathlineto{\pgfqpoint{1.276028in}{2.644559in}}%
\pgfpathlineto{\pgfqpoint{1.276032in}{2.760597in}}%
\pgfpathlineto{\pgfqpoint{1.276101in}{2.622439in}}%
\pgfpathlineto{\pgfqpoint{1.276784in}{2.661967in}}%
\pgfpathlineto{\pgfqpoint{1.277625in}{2.784342in}}%
\pgfpathlineto{\pgfqpoint{1.278002in}{2.807940in}}%
\pgfpathlineto{\pgfqpoint{1.278467in}{2.644559in}}%
\pgfpathlineto{\pgfqpoint{1.280973in}{2.644559in}}%
\pgfpathlineto{\pgfqpoint{1.280978in}{2.763191in}}%
\pgfpathlineto{\pgfqpoint{1.281107in}{2.633771in}}%
\pgfpathlineto{\pgfqpoint{1.281700in}{2.670876in}}%
\pgfpathlineto{\pgfqpoint{1.282541in}{2.772494in}}%
\pgfpathlineto{\pgfqpoint{1.282917in}{2.796073in}}%
\pgfpathlineto{\pgfqpoint{1.283435in}{2.644559in}}%
\pgfpathlineto{\pgfqpoint{1.285919in}{2.644559in}}%
\pgfpathlineto{\pgfqpoint{1.285923in}{2.757787in}}%
\pgfpathlineto{\pgfqpoint{1.285996in}{2.621139in}}%
\pgfpathlineto{\pgfqpoint{1.286680in}{2.661221in}}%
\pgfpathlineto{\pgfqpoint{1.287521in}{2.784950in}}%
\pgfpathlineto{\pgfqpoint{1.287893in}{2.808217in}}%
\pgfpathlineto{\pgfqpoint{1.288355in}{2.644559in}}%
\pgfpathlineto{\pgfqpoint{1.290864in}{2.644559in}}%
\pgfpathlineto{\pgfqpoint{1.290869in}{2.756286in}}%
\pgfpathlineto{\pgfqpoint{1.290933in}{2.623836in}}%
\pgfpathlineto{\pgfqpoint{1.291614in}{2.662480in}}%
\pgfpathlineto{\pgfqpoint{1.292454in}{2.781751in}}%
\pgfpathlineto{\pgfqpoint{1.292838in}{2.805771in}}%
\pgfpathlineto{\pgfqpoint{1.293321in}{2.644559in}}%
\pgfpathlineto{\pgfqpoint{1.295810in}{2.644559in}}%
\pgfpathlineto{\pgfqpoint{1.295814in}{2.762453in}}%
\pgfpathlineto{\pgfqpoint{1.295952in}{2.626265in}}%
\pgfpathlineto{\pgfqpoint{1.296545in}{2.663489in}}%
\pgfpathlineto{\pgfqpoint{1.297386in}{2.780770in}}%
\pgfpathlineto{\pgfqpoint{1.297758in}{2.804039in}}%
\pgfpathlineto{\pgfqpoint{1.298234in}{2.644559in}}%
\pgfpathlineto{\pgfqpoint{1.300755in}{2.644559in}}%
\pgfpathlineto{\pgfqpoint{1.300760in}{2.762559in}}%
\pgfpathlineto{\pgfqpoint{1.300834in}{2.621889in}}%
\pgfpathlineto{\pgfqpoint{1.301519in}{2.662241in}}%
\pgfpathlineto{\pgfqpoint{1.302359in}{2.785546in}}%
\pgfpathlineto{\pgfqpoint{1.302729in}{2.808679in}}%
\pgfpathlineto{\pgfqpoint{1.303215in}{2.644559in}}%
\pgfpathlineto{\pgfqpoint{1.305701in}{2.644559in}}%
\pgfpathlineto{\pgfqpoint{1.305705in}{2.762271in}}%
\pgfpathlineto{\pgfqpoint{1.305842in}{2.632209in}}%
\pgfpathlineto{\pgfqpoint{1.306435in}{2.669407in}}%
\pgfpathlineto{\pgfqpoint{1.307276in}{2.774734in}}%
\pgfpathlineto{\pgfqpoint{1.307648in}{2.798049in}}%
\pgfpathlineto{\pgfqpoint{1.308160in}{2.644559in}}%
\pgfpathlineto{\pgfqpoint{1.310646in}{2.644559in}}%
\pgfpathlineto{\pgfqpoint{1.310651in}{2.762818in}}%
\pgfpathlineto{\pgfqpoint{1.311415in}{2.703257in}}%
\pgfpathlineto{\pgfqpoint{1.311464in}{2.695764in}}%
\pgfpathlineto{\pgfqpoint{1.311909in}{2.734297in}}%
\pgfpathlineto{\pgfqpoint{1.311959in}{2.726861in}}%
\pgfpathlineto{\pgfqpoint{1.312586in}{2.776771in}}%
\pgfpathlineto{\pgfqpoint{1.312659in}{2.724986in}}%
\pgfpathlineto{\pgfqpoint{1.312670in}{2.744483in}}%
\pgfpathlineto{\pgfqpoint{1.313446in}{2.644559in}}%
\pgfpathlineto{\pgfqpoint{1.315591in}{2.644559in}}%
\pgfpathlineto{\pgfqpoint{1.315601in}{2.747265in}}%
\pgfpathlineto{\pgfqpoint{1.315658in}{2.627967in}}%
\pgfpathlineto{\pgfqpoint{1.316338in}{2.666407in}}%
\pgfpathlineto{\pgfqpoint{1.317179in}{2.777879in}}%
\pgfpathlineto{\pgfqpoint{1.317565in}{2.802089in}}%
\pgfpathlineto{\pgfqpoint{1.318023in}{2.644559in}}%
\pgfpathlineto{\pgfqpoint{1.320537in}{2.644559in}}%
\pgfpathlineto{\pgfqpoint{1.320541in}{2.762501in}}%
\pgfpathlineto{\pgfqpoint{1.320618in}{2.634896in}}%
\pgfpathlineto{\pgfqpoint{1.321303in}{2.675495in}}%
\pgfpathlineto{\pgfqpoint{1.322144in}{2.772641in}}%
\pgfpathlineto{\pgfqpoint{1.322511in}{2.795625in}}%
\pgfpathlineto{\pgfqpoint{1.323063in}{2.644559in}}%
\pgfpathlineto{\pgfqpoint{1.325482in}{2.644559in}}%
\pgfpathlineto{\pgfqpoint{1.325487in}{2.762927in}}%
\pgfpathlineto{\pgfqpoint{1.325562in}{2.635826in}}%
\pgfpathlineto{\pgfqpoint{1.326247in}{2.676249in}}%
\pgfpathlineto{\pgfqpoint{1.327087in}{2.771795in}}%
\pgfpathlineto{\pgfqpoint{1.327456in}{2.794898in}}%
\pgfpathlineto{\pgfqpoint{1.328013in}{2.644559in}}%
\pgfpathlineto{\pgfqpoint{1.330428in}{2.644559in}}%
\pgfpathlineto{\pgfqpoint{1.330432in}{2.761458in}}%
\pgfpathlineto{\pgfqpoint{1.331188in}{2.707974in}}%
\pgfpathlineto{\pgfqpoint{1.331238in}{2.689629in}}%
\pgfpathlineto{\pgfqpoint{1.331782in}{2.745226in}}%
\pgfpathlineto{\pgfqpoint{1.331831in}{2.726939in}}%
\pgfpathlineto{\pgfqpoint{1.332364in}{2.781760in}}%
\pgfpathlineto{\pgfqpoint{1.332457in}{2.693784in}}%
\pgfpathlineto{\pgfqpoint{1.332471in}{2.753785in}}%
\pgfpathlineto{\pgfqpoint{1.333243in}{2.644559in}}%
\pgfpathlineto{\pgfqpoint{1.335373in}{2.644559in}}%
\pgfpathlineto{\pgfqpoint{1.335378in}{2.760160in}}%
\pgfpathlineto{\pgfqpoint{1.335452in}{2.635748in}}%
\pgfpathlineto{\pgfqpoint{1.336137in}{2.676093in}}%
\pgfpathlineto{\pgfqpoint{1.336978in}{2.771102in}}%
\pgfpathlineto{\pgfqpoint{1.337347in}{2.794243in}}%
\pgfpathlineto{\pgfqpoint{1.337904in}{2.644559in}}%
\pgfpathlineto{\pgfqpoint{1.340319in}{2.644559in}}%
\pgfpathlineto{\pgfqpoint{1.340323in}{2.760828in}}%
\pgfpathlineto{\pgfqpoint{1.341072in}{2.697470in}}%
\pgfpathlineto{\pgfqpoint{1.341913in}{2.748853in}}%
\pgfpathlineto{\pgfqpoint{1.342293in}{2.772699in}}%
\pgfpathlineto{\pgfqpoint{1.342347in}{2.721495in}}%
\pgfpathlineto{\pgfqpoint{1.342361in}{2.727154in}}%
\pgfpathlineto{\pgfqpoint{1.343128in}{2.644559in}}%
\pgfpathlineto{\pgfqpoint{1.345264in}{2.644559in}}%
\pgfpathlineto{\pgfqpoint{1.345269in}{2.761304in}}%
\pgfpathlineto{\pgfqpoint{1.346028in}{2.701294in}}%
\pgfpathlineto{\pgfqpoint{1.346078in}{2.696761in}}%
\pgfpathlineto{\pgfqpoint{1.346424in}{2.726129in}}%
\pgfpathlineto{\pgfqpoint{1.346474in}{2.721634in}}%
\pgfpathlineto{\pgfqpoint{1.347202in}{2.774955in}}%
\pgfpathlineto{\pgfqpoint{1.347265in}{2.759721in}}%
\pgfpathlineto{\pgfqpoint{1.348018in}{2.644559in}}%
\pgfpathlineto{\pgfqpoint{1.350210in}{2.644559in}}%
\pgfpathlineto{\pgfqpoint{1.350214in}{2.763586in}}%
\pgfpathlineto{\pgfqpoint{1.350296in}{2.635713in}}%
\pgfpathlineto{\pgfqpoint{1.350981in}{2.676785in}}%
\pgfpathlineto{\pgfqpoint{1.351821in}{2.772294in}}%
\pgfpathlineto{\pgfqpoint{1.352184in}{2.794958in}}%
\pgfpathlineto{\pgfqpoint{1.352713in}{2.644559in}}%
\pgfpathlineto{\pgfqpoint{1.355155in}{2.644559in}}%
\pgfpathlineto{\pgfqpoint{1.355160in}{2.757509in}}%
\pgfpathlineto{\pgfqpoint{1.355238in}{2.640731in}}%
\pgfpathlineto{\pgfqpoint{1.355925in}{2.681581in}}%
\pgfpathlineto{\pgfqpoint{1.356765in}{2.765700in}}%
\pgfpathlineto{\pgfqpoint{1.357129in}{2.788511in}}%
\pgfpathlineto{\pgfqpoint{1.357744in}{2.644559in}}%
\pgfpathlineto{\pgfqpoint{1.360101in}{2.644559in}}%
\pgfpathlineto{\pgfqpoint{1.360105in}{2.763499in}}%
\pgfpathlineto{\pgfqpoint{1.360172in}{2.623432in}}%
\pgfpathlineto{\pgfqpoint{1.360854in}{2.662709in}}%
\pgfpathlineto{\pgfqpoint{1.361695in}{2.784133in}}%
\pgfpathlineto{\pgfqpoint{1.362074in}{2.807876in}}%
\pgfpathlineto{\pgfqpoint{1.362553in}{2.644559in}}%
\pgfpathlineto{\pgfqpoint{1.365046in}{2.644559in}}%
\pgfpathlineto{\pgfqpoint{1.365051in}{2.755194in}}%
\pgfpathlineto{\pgfqpoint{1.365119in}{2.624049in}}%
\pgfpathlineto{\pgfqpoint{1.365801in}{2.663655in}}%
\pgfpathlineto{\pgfqpoint{1.366642in}{2.781118in}}%
\pgfpathlineto{\pgfqpoint{1.367020in}{2.804763in}}%
\pgfpathlineto{\pgfqpoint{1.367497in}{2.644559in}}%
\pgfpathlineto{\pgfqpoint{1.369992in}{2.644559in}}%
\pgfpathlineto{\pgfqpoint{1.369996in}{2.756078in}}%
\pgfpathlineto{\pgfqpoint{1.370071in}{2.634347in}}%
\pgfpathlineto{\pgfqpoint{1.370756in}{2.674803in}}%
\pgfpathlineto{\pgfqpoint{1.371597in}{2.771389in}}%
\pgfpathlineto{\pgfqpoint{1.371965in}{2.794450in}}%
\pgfpathlineto{\pgfqpoint{1.372489in}{2.644559in}}%
\pgfpathlineto{\pgfqpoint{1.374937in}{2.644559in}}%
\pgfpathlineto{\pgfqpoint{1.374941in}{2.759628in}}%
\pgfpathlineto{\pgfqpoint{1.375014in}{2.637498in}}%
\pgfpathlineto{\pgfqpoint{1.375698in}{2.677575in}}%
\pgfpathlineto{\pgfqpoint{1.376539in}{2.769190in}}%
\pgfpathlineto{\pgfqpoint{1.376911in}{2.792483in}}%
\pgfpathlineto{\pgfqpoint{1.377463in}{2.644559in}}%
\pgfpathlineto{\pgfqpoint{1.379882in}{2.644559in}}%
\pgfpathlineto{\pgfqpoint{1.379887in}{2.761422in}}%
\pgfpathlineto{\pgfqpoint{1.380613in}{2.698027in}}%
\pgfpathlineto{\pgfqpoint{1.380662in}{2.695821in}}%
\pgfpathlineto{\pgfqpoint{1.380910in}{2.716668in}}%
\pgfpathlineto{\pgfqpoint{1.380959in}{2.714462in}}%
\pgfpathlineto{\pgfqpoint{1.381856in}{2.776135in}}%
\pgfpathlineto{\pgfqpoint{1.381883in}{2.759739in}}%
\pgfpathlineto{\pgfqpoint{1.382639in}{2.644559in}}%
\pgfpathlineto{\pgfqpoint{1.384828in}{2.644559in}}%
\pgfpathlineto{\pgfqpoint{1.384832in}{2.763387in}}%
\pgfpathlineto{\pgfqpoint{1.385572in}{2.685879in}}%
\pgfpathlineto{\pgfqpoint{1.386412in}{2.759796in}}%
\pgfpathlineto{\pgfqpoint{1.386802in}{2.784223in}}%
\pgfpathlineto{\pgfqpoint{1.387461in}{2.644559in}}%
\pgfpathlineto{\pgfqpoint{1.389773in}{2.644559in}}%
\pgfpathlineto{\pgfqpoint{1.389778in}{2.758387in}}%
\pgfpathlineto{\pgfqpoint{1.390521in}{2.689795in}}%
\pgfpathlineto{\pgfqpoint{1.391461in}{2.761325in}}%
\pgfpathlineto{\pgfqpoint{1.391747in}{2.779274in}}%
\pgfpathlineto{\pgfqpoint{1.392806in}{2.644559in}}%
\pgfpathlineto{\pgfqpoint{1.394719in}{2.644559in}}%
\pgfpathlineto{\pgfqpoint{1.394723in}{2.761658in}}%
\pgfpathlineto{\pgfqpoint{1.394804in}{2.622381in}}%
\pgfpathlineto{\pgfqpoint{1.395489in}{2.663546in}}%
\pgfpathlineto{\pgfqpoint{1.396329in}{2.784867in}}%
\pgfpathlineto{\pgfqpoint{1.396693in}{2.807587in}}%
\pgfpathlineto{\pgfqpoint{1.397204in}{2.644559in}}%
\pgfpathlineto{\pgfqpoint{1.399664in}{2.644559in}}%
\pgfpathlineto{\pgfqpoint{1.399669in}{2.758782in}}%
\pgfpathlineto{\pgfqpoint{1.399737in}{2.624203in}}%
\pgfpathlineto{\pgfqpoint{1.400420in}{2.663875in}}%
\pgfpathlineto{\pgfqpoint{1.401261in}{2.781959in}}%
\pgfpathlineto{\pgfqpoint{1.401638in}{2.805576in}}%
\pgfpathlineto{\pgfqpoint{1.402112in}{2.644559in}}%
\pgfpathlineto{\pgfqpoint{1.404610in}{2.644559in}}%
\pgfpathlineto{\pgfqpoint{1.404614in}{2.759333in}}%
\pgfpathlineto{\pgfqpoint{1.404678in}{2.624411in}}%
\pgfpathlineto{\pgfqpoint{1.405359in}{2.663333in}}%
\pgfpathlineto{\pgfqpoint{1.406200in}{2.781845in}}%
\pgfpathlineto{\pgfqpoint{1.406584in}{2.805871in}}%
\pgfpathlineto{\pgfqpoint{1.407062in}{2.644559in}}%
\pgfpathlineto{\pgfqpoint{1.409555in}{2.644559in}}%
\pgfpathlineto{\pgfqpoint{1.409560in}{2.758290in}}%
\pgfpathlineto{\pgfqpoint{1.409640in}{2.621250in}}%
\pgfpathlineto{\pgfqpoint{1.410325in}{2.662400in}}%
\pgfpathlineto{\pgfqpoint{1.411166in}{2.785058in}}%
\pgfpathlineto{\pgfqpoint{1.411529in}{2.807775in}}%
\pgfpathlineto{\pgfqpoint{1.411994in}{2.644559in}}%
\pgfpathlineto{\pgfqpoint{1.414501in}{2.644559in}}%
\pgfpathlineto{\pgfqpoint{1.414505in}{2.757759in}}%
\pgfpathlineto{\pgfqpoint{1.415269in}{2.709002in}}%
\pgfpathlineto{\pgfqpoint{1.415318in}{2.688550in}}%
\pgfpathlineto{\pgfqpoint{1.415862in}{2.746251in}}%
\pgfpathlineto{\pgfqpoint{1.415912in}{2.725864in}}%
\pgfpathlineto{\pgfqpoint{1.416440in}{2.782534in}}%
\pgfpathlineto{\pgfqpoint{1.416529in}{2.699213in}}%
\pgfpathlineto{\pgfqpoint{1.416544in}{2.747750in}}%
\pgfpathlineto{\pgfqpoint{1.417286in}{2.644559in}}%
\pgfpathlineto{\pgfqpoint{1.419446in}{2.644559in}}%
\pgfpathlineto{\pgfqpoint{1.419451in}{2.762309in}}%
\pgfpathlineto{\pgfqpoint{1.419517in}{2.643751in}}%
\pgfpathlineto{\pgfqpoint{1.420199in}{2.682891in}}%
\pgfpathlineto{\pgfqpoint{1.421039in}{2.763577in}}%
\pgfpathlineto{\pgfqpoint{1.421420in}{2.787434in}}%
\pgfpathlineto{\pgfqpoint{1.422030in}{2.644559in}}%
\pgfpathlineto{\pgfqpoint{1.424392in}{2.644559in}}%
\pgfpathlineto{\pgfqpoint{1.424397in}{2.749278in}}%
\pgfpathlineto{\pgfqpoint{1.424461in}{2.638753in}}%
\pgfpathlineto{\pgfqpoint{1.425144in}{2.678185in}}%
\pgfpathlineto{\pgfqpoint{1.425985in}{2.767171in}}%
\pgfpathlineto{\pgfqpoint{1.426365in}{2.791014in}}%
\pgfpathlineto{\pgfqpoint{1.426905in}{2.644559in}}%
\pgfpathlineto{\pgfqpoint{1.429337in}{2.644559in}}%
\pgfpathlineto{\pgfqpoint{1.429341in}{2.755975in}}%
\pgfpathlineto{\pgfqpoint{1.430098in}{2.697871in}}%
\pgfpathlineto{\pgfqpoint{1.430247in}{2.704562in}}%
\pgfpathlineto{\pgfqpoint{1.431273in}{2.771690in}}%
\pgfpathlineto{\pgfqpoint{1.431311in}{2.749002in}}%
\pgfpathlineto{\pgfqpoint{1.431335in}{2.758047in}}%
\pgfpathlineto{\pgfqpoint{1.432082in}{2.644559in}}%
\pgfpathlineto{\pgfqpoint{1.434282in}{2.644559in}}%
\pgfpathlineto{\pgfqpoint{1.434287in}{2.756838in}}%
\pgfpathlineto{\pgfqpoint{1.434348in}{2.624933in}}%
\pgfpathlineto{\pgfqpoint{1.435027in}{2.663258in}}%
\pgfpathlineto{\pgfqpoint{1.435868in}{2.780674in}}%
\pgfpathlineto{\pgfqpoint{1.436256in}{2.804978in}}%
\pgfpathlineto{\pgfqpoint{1.436737in}{2.644559in}}%
\pgfpathlineto{\pgfqpoint{1.439228in}{2.644559in}}%
\pgfpathlineto{\pgfqpoint{1.439232in}{2.755208in}}%
\pgfpathlineto{\pgfqpoint{1.439314in}{2.620415in}}%
\pgfpathlineto{\pgfqpoint{1.439999in}{2.661664in}}%
\pgfpathlineto{\pgfqpoint{1.440840in}{2.785072in}}%
\pgfpathlineto{\pgfqpoint{1.441202in}{2.807716in}}%
\pgfpathlineto{\pgfqpoint{1.441667in}{2.644559in}}%
\pgfpathlineto{\pgfqpoint{1.444173in}{2.644559in}}%
\pgfpathlineto{\pgfqpoint{1.444178in}{2.759285in}}%
\pgfpathlineto{\pgfqpoint{1.444247in}{2.639845in}}%
\pgfpathlineto{\pgfqpoint{1.444930in}{2.679464in}}%
\pgfpathlineto{\pgfqpoint{1.445771in}{2.766672in}}%
\pgfpathlineto{\pgfqpoint{1.446147in}{2.790247in}}%
\pgfpathlineto{\pgfqpoint{1.446711in}{2.644559in}}%
\pgfpathlineto{\pgfqpoint{1.449119in}{2.644559in}}%
\pgfpathlineto{\pgfqpoint{1.449123in}{2.759774in}}%
\pgfpathlineto{\pgfqpoint{1.449195in}{2.637108in}}%
\pgfpathlineto{\pgfqpoint{1.449880in}{2.677414in}}%
\pgfpathlineto{\pgfqpoint{1.450721in}{2.769522in}}%
\pgfpathlineto{\pgfqpoint{1.451093in}{2.792833in}}%
\pgfpathlineto{\pgfqpoint{1.451653in}{2.644559in}}%
\pgfpathlineto{\pgfqpoint{1.454064in}{2.644559in}}%
\pgfpathlineto{\pgfqpoint{1.454069in}{2.756554in}}%
\pgfpathlineto{\pgfqpoint{1.454812in}{2.699004in}}%
\pgfpathlineto{\pgfqpoint{1.454861in}{2.695722in}}%
\pgfpathlineto{\pgfqpoint{1.455207in}{2.723868in}}%
\pgfpathlineto{\pgfqpoint{1.455257in}{2.720567in}}%
\pgfpathlineto{\pgfqpoint{1.455994in}{2.773283in}}%
\pgfpathlineto{\pgfqpoint{1.456069in}{2.744738in}}%
\pgfpathlineto{\pgfqpoint{1.456863in}{2.644559in}}%
\pgfpathlineto{\pgfqpoint{1.459010in}{2.644559in}}%
\pgfpathlineto{\pgfqpoint{1.459014in}{2.757202in}}%
\pgfpathlineto{\pgfqpoint{1.459095in}{2.621741in}}%
\pgfpathlineto{\pgfqpoint{1.459780in}{2.662969in}}%
\pgfpathlineto{\pgfqpoint{1.460621in}{2.784292in}}%
\pgfpathlineto{\pgfqpoint{1.460984in}{2.806964in}}%
\pgfpathlineto{\pgfqpoint{1.461453in}{2.644559in}}%
\pgfpathlineto{\pgfqpoint{1.463955in}{2.644559in}}%
\pgfpathlineto{\pgfqpoint{1.463960in}{2.757819in}}%
\pgfpathlineto{\pgfqpoint{1.464041in}{2.621036in}}%
\pgfpathlineto{\pgfqpoint{1.464725in}{2.662215in}}%
\pgfpathlineto{\pgfqpoint{1.465566in}{2.785150in}}%
\pgfpathlineto{\pgfqpoint{1.465929in}{2.807848in}}%
\pgfpathlineto{\pgfqpoint{1.466407in}{2.644559in}}%
\pgfpathlineto{\pgfqpoint{1.468901in}{2.644559in}}%
\pgfpathlineto{\pgfqpoint{1.468905in}{2.756287in}}%
\pgfpathlineto{\pgfqpoint{1.468986in}{2.620718in}}%
\pgfpathlineto{\pgfqpoint{1.469671in}{2.661901in}}%
\pgfpathlineto{\pgfqpoint{1.470512in}{2.785048in}}%
\pgfpathlineto{\pgfqpoint{1.470874in}{2.807736in}}%
\pgfpathlineto{\pgfqpoint{1.471339in}{2.644559in}}%
\pgfpathlineto{\pgfqpoint{1.473846in}{2.644559in}}%
\pgfpathlineto{\pgfqpoint{1.473851in}{2.756960in}}%
\pgfpathlineto{\pgfqpoint{1.474576in}{2.692924in}}%
\pgfpathlineto{\pgfqpoint{1.475791in}{2.772850in}}%
\pgfpathlineto{\pgfqpoint{1.475820in}{2.761822in}}%
\pgfpathlineto{\pgfqpoint{1.476550in}{2.644559in}}%
\pgfpathlineto{\pgfqpoint{1.478792in}{2.644559in}}%
\pgfpathlineto{\pgfqpoint{1.478796in}{2.755240in}}%
\pgfpathlineto{\pgfqpoint{1.479530in}{2.685675in}}%
\pgfpathlineto{\pgfqpoint{1.480470in}{2.763364in}}%
\pgfpathlineto{\pgfqpoint{1.480765in}{2.781914in}}%
\pgfpathlineto{\pgfqpoint{1.481594in}{2.644559in}}%
\pgfpathlineto{\pgfqpoint{1.483737in}{2.644559in}}%
\pgfpathlineto{\pgfqpoint{1.483741in}{2.760270in}}%
\pgfpathlineto{\pgfqpoint{1.483808in}{2.643917in}}%
\pgfpathlineto{\pgfqpoint{1.484491in}{2.683335in}}%
\pgfpathlineto{\pgfqpoint{1.485332in}{2.762809in}}%
\pgfpathlineto{\pgfqpoint{1.485711in}{2.786593in}}%
\pgfpathlineto{\pgfqpoint{1.486321in}{2.644559in}}%
\pgfpathlineto{\pgfqpoint{1.488682in}{2.644559in}}%
\pgfpathlineto{\pgfqpoint{1.488688in}{2.749722in}}%
\pgfpathlineto{\pgfqpoint{1.489408in}{2.697545in}}%
\pgfpathlineto{\pgfqpoint{1.489458in}{2.694324in}}%
\pgfpathlineto{\pgfqpoint{1.489705in}{2.716192in}}%
\pgfpathlineto{\pgfqpoint{1.489755in}{2.712958in}}%
\pgfpathlineto{\pgfqpoint{1.490656in}{2.775966in}}%
\pgfpathlineto{\pgfqpoint{1.490657in}{2.748823in}}%
\pgfpathlineto{\pgfqpoint{1.490678in}{2.756234in}}%
\pgfpathlineto{\pgfqpoint{1.491391in}{2.644559in}}%
\pgfpathlineto{\pgfqpoint{1.493628in}{2.644559in}}%
\pgfpathlineto{\pgfqpoint{1.493632in}{2.759846in}}%
\pgfpathlineto{\pgfqpoint{1.493752in}{2.635362in}}%
\pgfpathlineto{\pgfqpoint{1.494395in}{2.719424in}}%
\pgfpathlineto{\pgfqpoint{1.494444in}{2.678583in}}%
\pgfpathlineto{\pgfqpoint{1.495087in}{2.762831in}}%
\pgfpathlineto{\pgfqpoint{1.495137in}{2.722164in}}%
\pgfpathlineto{\pgfqpoint{1.495567in}{2.792906in}}%
\pgfpathlineto{\pgfqpoint{1.495691in}{2.644559in}}%
\pgfpathlineto{\pgfqpoint{1.495790in}{2.644559in}}%
\pgfpathlineto{\pgfqpoint{1.498573in}{2.644559in}}%
\pgfpathlineto{\pgfqpoint{1.498578in}{2.755630in}}%
\pgfpathlineto{\pgfqpoint{1.499301in}{2.690235in}}%
\pgfpathlineto{\pgfqpoint{1.500340in}{2.763912in}}%
\pgfpathlineto{\pgfqpoint{1.500518in}{2.775095in}}%
\pgfpathlineto{\pgfqpoint{1.500610in}{2.708491in}}%
\pgfpathlineto{\pgfqpoint{1.501382in}{2.644559in}}%
\pgfpathlineto{\pgfqpoint{1.503519in}{2.644559in}}%
\pgfpathlineto{\pgfqpoint{1.503523in}{2.755144in}}%
\pgfpathlineto{\pgfqpoint{1.504248in}{2.686687in}}%
\pgfpathlineto{\pgfqpoint{1.505187in}{2.761132in}}%
\pgfpathlineto{\pgfqpoint{1.505464in}{2.778466in}}%
\pgfpathlineto{\pgfqpoint{1.506426in}{2.644559in}}%
\pgfpathlineto{\pgfqpoint{1.508464in}{2.644559in}}%
\pgfpathlineto{\pgfqpoint{1.508469in}{2.760100in}}%
\pgfpathlineto{\pgfqpoint{1.508599in}{2.631465in}}%
\pgfpathlineto{\pgfqpoint{1.509192in}{2.668615in}}%
\pgfpathlineto{\pgfqpoint{1.510033in}{2.774170in}}%
\pgfpathlineto{\pgfqpoint{1.510409in}{2.797692in}}%
\pgfpathlineto{\pgfqpoint{1.510922in}{2.644559in}}%
\pgfpathlineto{\pgfqpoint{1.513410in}{2.644559in}}%
\pgfpathlineto{\pgfqpoint{1.513414in}{2.754895in}}%
\pgfpathlineto{\pgfqpoint{1.513535in}{2.632230in}}%
\pgfpathlineto{\pgfqpoint{1.514178in}{2.721425in}}%
\pgfpathlineto{\pgfqpoint{1.514228in}{2.675498in}}%
\pgfpathlineto{\pgfqpoint{1.514870in}{2.764846in}}%
\pgfpathlineto{\pgfqpoint{1.514920in}{2.719066in}}%
\pgfpathlineto{\pgfqpoint{1.515350in}{2.794885in}}%
\pgfpathlineto{\pgfqpoint{1.515476in}{2.644559in}}%
\pgfpathlineto{\pgfqpoint{1.515574in}{2.644559in}}%
\pgfpathlineto{\pgfqpoint{1.518355in}{2.644559in}}%
\pgfpathlineto{\pgfqpoint{1.518360in}{2.755729in}}%
\pgfpathlineto{\pgfqpoint{1.519123in}{2.706934in}}%
\pgfpathlineto{\pgfqpoint{1.519172in}{2.690050in}}%
\pgfpathlineto{\pgfqpoint{1.519716in}{2.744175in}}%
\pgfpathlineto{\pgfqpoint{1.519766in}{2.727372in}}%
\pgfpathlineto{\pgfqpoint{1.520295in}{2.780464in}}%
\pgfpathlineto{\pgfqpoint{1.520392in}{2.706918in}}%
\pgfpathlineto{\pgfqpoint{1.521159in}{2.644559in}}%
\pgfpathlineto{\pgfqpoint{1.523301in}{2.644559in}}%
\pgfpathlineto{\pgfqpoint{1.523305in}{2.758698in}}%
\pgfpathlineto{\pgfqpoint{1.523379in}{2.630711in}}%
\pgfpathlineto{\pgfqpoint{1.524064in}{2.671277in}}%
\pgfpathlineto{\pgfqpoint{1.524905in}{2.775580in}}%
\pgfpathlineto{\pgfqpoint{1.525274in}{2.798741in}}%
\pgfpathlineto{\pgfqpoint{1.525765in}{2.644559in}}%
\pgfpathlineto{\pgfqpoint{1.528246in}{2.644559in}}%
\pgfpathlineto{\pgfqpoint{1.528251in}{2.754506in}}%
\pgfpathlineto{\pgfqpoint{1.528999in}{2.704478in}}%
\pgfpathlineto{\pgfqpoint{1.529048in}{2.690286in}}%
\pgfpathlineto{\pgfqpoint{1.529592in}{2.741724in}}%
\pgfpathlineto{\pgfqpoint{1.529642in}{2.727602in}}%
\pgfpathlineto{\pgfqpoint{1.530178in}{2.778490in}}%
\pgfpathlineto{\pgfqpoint{1.530275in}{2.697403in}}%
\pgfpathlineto{\pgfqpoint{1.530290in}{2.749221in}}%
\pgfpathlineto{\pgfqpoint{1.531055in}{2.644559in}}%
\pgfpathlineto{\pgfqpoint{1.533192in}{2.644559in}}%
\pgfpathlineto{\pgfqpoint{1.533196in}{2.759497in}}%
\pgfpathlineto{\pgfqpoint{1.533263in}{2.624050in}}%
\pgfpathlineto{\pgfqpoint{1.533945in}{2.663441in}}%
\pgfpathlineto{\pgfqpoint{1.534785in}{2.782269in}}%
\pgfpathlineto{\pgfqpoint{1.535165in}{2.806048in}}%
\pgfpathlineto{\pgfqpoint{1.535642in}{2.644559in}}%
\pgfpathlineto{\pgfqpoint{1.538137in}{2.644559in}}%
\pgfpathlineto{\pgfqpoint{1.538143in}{2.749344in}}%
\pgfpathlineto{\pgfqpoint{1.538876in}{2.692219in}}%
\pgfpathlineto{\pgfqpoint{1.540013in}{2.769638in}}%
\pgfpathlineto{\pgfqpoint{1.540144in}{2.733881in}}%
\pgfpathlineto{\pgfqpoint{1.540111in}{2.775776in}}%
\pgfpathlineto{\pgfqpoint{1.540154in}{2.734311in}}%
\pgfpathlineto{\pgfqpoint{1.540969in}{2.644559in}}%
\pgfpathlineto{\pgfqpoint{1.543082in}{2.644559in}}%
\pgfpathlineto{\pgfqpoint{1.543087in}{2.752230in}}%
\pgfpathlineto{\pgfqpoint{1.543816in}{2.693171in}}%
\pgfpathlineto{\pgfqpoint{1.545029in}{2.771602in}}%
\pgfpathlineto{\pgfqpoint{1.545057in}{2.750384in}}%
\pgfpathlineto{\pgfqpoint{1.545078in}{2.755799in}}%
\pgfpathlineto{\pgfqpoint{1.545784in}{2.644559in}}%
\pgfpathlineto{\pgfqpoint{1.548028in}{2.644559in}}%
\pgfpathlineto{\pgfqpoint{1.548032in}{2.755564in}}%
\pgfpathlineto{\pgfqpoint{1.548761in}{2.686265in}}%
\pgfpathlineto{\pgfqpoint{1.549701in}{2.762179in}}%
\pgfpathlineto{\pgfqpoint{1.549975in}{2.779384in}}%
\pgfpathlineto{\pgfqpoint{1.550876in}{2.644559in}}%
\pgfpathlineto{\pgfqpoint{1.552973in}{2.644559in}}%
\pgfpathlineto{\pgfqpoint{1.552978in}{2.755836in}}%
\pgfpathlineto{\pgfqpoint{1.553700in}{2.695537in}}%
\pgfpathlineto{\pgfqpoint{1.553947in}{2.708727in}}%
\pgfpathlineto{\pgfqpoint{1.554947in}{2.773968in}}%
\pgfpathlineto{\pgfqpoint{1.554977in}{2.749421in}}%
\pgfpathlineto{\pgfqpoint{1.555771in}{2.644559in}}%
\pgfpathlineto{\pgfqpoint{1.557919in}{2.644559in}}%
\pgfpathlineto{\pgfqpoint{1.557923in}{2.756530in}}%
\pgfpathlineto{\pgfqpoint{1.558665in}{2.698953in}}%
\pgfpathlineto{\pgfqpoint{1.558714in}{2.695592in}}%
\pgfpathlineto{\pgfqpoint{1.559060in}{2.723818in}}%
\pgfpathlineto{\pgfqpoint{1.559110in}{2.720436in}}%
\pgfpathlineto{\pgfqpoint{1.559847in}{2.773275in}}%
\pgfpathlineto{\pgfqpoint{1.559923in}{2.744483in}}%
\pgfpathlineto{\pgfqpoint{1.560718in}{2.644559in}}%
\pgfpathlineto{\pgfqpoint{1.562864in}{2.644559in}}%
\pgfpathlineto{\pgfqpoint{1.562870in}{2.749295in}}%
\pgfpathlineto{\pgfqpoint{1.562939in}{2.637215in}}%
\pgfpathlineto{\pgfqpoint{1.563624in}{2.677443in}}%
\pgfpathlineto{\pgfqpoint{1.564465in}{2.768978in}}%
\pgfpathlineto{\pgfqpoint{1.564838in}{2.792360in}}%
\pgfpathlineto{\pgfqpoint{1.565398in}{2.644559in}}%
\pgfpathlineto{\pgfqpoint{1.567810in}{2.644559in}}%
\pgfpathlineto{\pgfqpoint{1.567815in}{2.749702in}}%
\pgfpathlineto{\pgfqpoint{1.568573in}{2.700361in}}%
\pgfpathlineto{\pgfqpoint{1.568622in}{2.695424in}}%
\pgfpathlineto{\pgfqpoint{1.568968in}{2.725220in}}%
\pgfpathlineto{\pgfqpoint{1.569018in}{2.720274in}}%
\pgfpathlineto{\pgfqpoint{1.569747in}{2.774129in}}%
\pgfpathlineto{\pgfqpoint{1.569826in}{2.735462in}}%
\pgfpathlineto{\pgfqpoint{1.570657in}{2.644559in}}%
\pgfpathlineto{\pgfqpoint{1.572755in}{2.644559in}}%
\pgfpathlineto{\pgfqpoint{1.572760in}{2.756060in}}%
\pgfpathlineto{\pgfqpoint{1.572833in}{2.630875in}}%
\pgfpathlineto{\pgfqpoint{1.573518in}{2.671407in}}%
\pgfpathlineto{\pgfqpoint{1.574359in}{2.774677in}}%
\pgfpathlineto{\pgfqpoint{1.574729in}{2.797856in}}%
\pgfpathlineto{\pgfqpoint{1.575252in}{2.644559in}}%
\pgfpathlineto{\pgfqpoint{1.577701in}{2.644559in}}%
\pgfpathlineto{\pgfqpoint{1.577705in}{2.759775in}}%
\pgfpathlineto{\pgfqpoint{1.578471in}{2.691321in}}%
\pgfpathlineto{\pgfqpoint{1.579411in}{2.763068in}}%
\pgfpathlineto{\pgfqpoint{1.579674in}{2.779611in}}%
\pgfpathlineto{\pgfqpoint{1.580794in}{2.644559in}}%
\pgfpathlineto{\pgfqpoint{1.582646in}{2.644559in}}%
\pgfpathlineto{\pgfqpoint{1.582651in}{2.753840in}}%
\pgfpathlineto{\pgfqpoint{1.582765in}{2.632078in}}%
\pgfpathlineto{\pgfqpoint{1.583407in}{2.720498in}}%
\pgfpathlineto{\pgfqpoint{1.583457in}{2.675245in}}%
\pgfpathlineto{\pgfqpoint{1.584100in}{2.763921in}}%
\pgfpathlineto{\pgfqpoint{1.584149in}{2.718812in}}%
\pgfpathlineto{\pgfqpoint{1.584582in}{2.794182in}}%
\pgfpathlineto{\pgfqpoint{1.584712in}{2.644559in}}%
\pgfpathlineto{\pgfqpoint{1.584806in}{2.644559in}}%
\pgfpathlineto{\pgfqpoint{1.587592in}{2.644559in}}%
\pgfpathlineto{\pgfqpoint{1.587596in}{2.758652in}}%
\pgfpathlineto{\pgfqpoint{1.588348in}{2.705072in}}%
\pgfpathlineto{\pgfqpoint{1.588398in}{2.691419in}}%
\pgfpathlineto{\pgfqpoint{1.588942in}{2.742356in}}%
\pgfpathlineto{\pgfqpoint{1.588991in}{2.728697in}}%
\pgfpathlineto{\pgfqpoint{1.589526in}{2.779026in}}%
\pgfpathlineto{\pgfqpoint{1.589620in}{2.716148in}}%
\pgfpathlineto{\pgfqpoint{1.589634in}{2.732045in}}%
\pgfpathlineto{\pgfqpoint{1.590391in}{2.644559in}}%
\pgfpathlineto{\pgfqpoint{1.592537in}{2.644559in}}%
\pgfpathlineto{\pgfqpoint{1.592542in}{2.749468in}}%
\pgfpathlineto{\pgfqpoint{1.593264in}{2.686906in}}%
\pgfpathlineto{\pgfqpoint{1.594204in}{2.760684in}}%
\pgfpathlineto{\pgfqpoint{1.594481in}{2.778098in}}%
\pgfpathlineto{\pgfqpoint{1.595423in}{2.644559in}}%
\pgfpathlineto{\pgfqpoint{1.597482in}{2.644559in}}%
\pgfpathlineto{\pgfqpoint{1.597487in}{2.758580in}}%
\pgfpathlineto{\pgfqpoint{1.598230in}{2.699714in}}%
\pgfpathlineto{\pgfqpoint{1.598280in}{2.695633in}}%
\pgfpathlineto{\pgfqpoint{1.598626in}{2.724579in}}%
\pgfpathlineto{\pgfqpoint{1.598675in}{2.720477in}}%
\pgfpathlineto{\pgfqpoint{1.599412in}{2.773980in}}%
\pgfpathlineto{\pgfqpoint{1.599487in}{2.746261in}}%
\pgfpathlineto{\pgfqpoint{1.600281in}{2.644559in}}%
\pgfpathlineto{\pgfqpoint{1.602428in}{2.644559in}}%
\pgfpathlineto{\pgfqpoint{1.602433in}{2.749795in}}%
\pgfpathlineto{\pgfqpoint{1.603158in}{2.690980in}}%
\pgfpathlineto{\pgfqpoint{1.604296in}{2.769893in}}%
\pgfpathlineto{\pgfqpoint{1.604373in}{2.774776in}}%
\pgfpathlineto{\pgfqpoint{1.604465in}{2.709343in}}%
\pgfpathlineto{\pgfqpoint{1.605233in}{2.644559in}}%
\pgfpathlineto{\pgfqpoint{1.607373in}{2.644559in}}%
\pgfpathlineto{\pgfqpoint{1.607379in}{2.749437in}}%
\pgfpathlineto{\pgfqpoint{1.608108in}{2.693268in}}%
\pgfpathlineto{\pgfqpoint{1.609321in}{2.773018in}}%
\pgfpathlineto{\pgfqpoint{1.609347in}{2.762534in}}%
\pgfpathlineto{\pgfqpoint{1.610077in}{2.644559in}}%
\pgfpathlineto{\pgfqpoint{1.612319in}{2.644559in}}%
\pgfpathlineto{\pgfqpoint{1.612323in}{2.752918in}}%
\pgfpathlineto{\pgfqpoint{1.612456in}{2.632371in}}%
\pgfpathlineto{\pgfqpoint{1.613049in}{2.669540in}}%
\pgfpathlineto{\pgfqpoint{1.613890in}{2.771691in}}%
\pgfpathlineto{\pgfqpoint{1.614264in}{2.795159in}}%
\pgfpathlineto{\pgfqpoint{1.614775in}{2.644559in}}%
\pgfpathlineto{\pgfqpoint{1.617264in}{2.644559in}}%
\pgfpathlineto{\pgfqpoint{1.617269in}{2.754363in}}%
\pgfpathlineto{\pgfqpoint{1.618032in}{2.693215in}}%
\pgfpathlineto{\pgfqpoint{1.619170in}{2.771854in}}%
\pgfpathlineto{\pgfqpoint{1.619238in}{2.776149in}}%
\pgfpathlineto{\pgfqpoint{1.620757in}{2.644559in}}%
\pgfpathlineto{\pgfqpoint{1.622210in}{2.644559in}}%
\pgfpathlineto{\pgfqpoint{1.622214in}{2.755193in}}%
\pgfpathlineto{\pgfqpoint{1.622958in}{2.684425in}}%
\pgfpathlineto{\pgfqpoint{1.623799in}{2.759664in}}%
\pgfpathlineto{\pgfqpoint{1.624184in}{2.783773in}}%
\pgfpathlineto{\pgfqpoint{1.624844in}{2.644559in}}%
\pgfpathlineto{\pgfqpoint{1.627155in}{2.644559in}}%
\pgfpathlineto{\pgfqpoint{1.627160in}{2.758141in}}%
\pgfpathlineto{\pgfqpoint{1.627906in}{2.699511in}}%
\pgfpathlineto{\pgfqpoint{1.627955in}{2.696060in}}%
\pgfpathlineto{\pgfqpoint{1.628301in}{2.724376in}}%
\pgfpathlineto{\pgfqpoint{1.628351in}{2.720904in}}%
\pgfpathlineto{\pgfqpoint{1.629086in}{2.773691in}}%
\pgfpathlineto{\pgfqpoint{1.629159in}{2.748412in}}%
\pgfpathlineto{\pgfqpoint{1.629953in}{2.644559in}}%
\pgfpathlineto{\pgfqpoint{1.632101in}{2.644559in}}%
\pgfpathlineto{\pgfqpoint{1.632106in}{2.749587in}}%
\pgfpathlineto{\pgfqpoint{1.632863in}{2.693211in}}%
\pgfpathlineto{\pgfqpoint{1.634001in}{2.771285in}}%
\pgfpathlineto{\pgfqpoint{1.634129in}{2.720626in}}%
\pgfpathlineto{\pgfqpoint{1.634074in}{2.775917in}}%
\pgfpathlineto{\pgfqpoint{1.634142in}{2.726436in}}%
\pgfpathlineto{\pgfqpoint{1.634915in}{2.644559in}}%
\pgfpathlineto{\pgfqpoint{1.637046in}{2.644559in}}%
\pgfpathlineto{\pgfqpoint{1.637052in}{2.749275in}}%
\pgfpathlineto{\pgfqpoint{1.637776in}{2.694316in}}%
\pgfpathlineto{\pgfqpoint{1.639020in}{2.772499in}}%
\pgfpathlineto{\pgfqpoint{1.639050in}{2.745664in}}%
\pgfpathlineto{\pgfqpoint{1.639845in}{2.644559in}}%
\pgfpathlineto{\pgfqpoint{1.641992in}{2.644559in}}%
\pgfpathlineto{\pgfqpoint{1.641997in}{2.749692in}}%
\pgfpathlineto{\pgfqpoint{1.642746in}{2.700947in}}%
\pgfpathlineto{\pgfqpoint{1.642795in}{2.694698in}}%
\pgfpathlineto{\pgfqpoint{1.643240in}{2.732023in}}%
\pgfpathlineto{\pgfqpoint{1.643290in}{2.725758in}}%
\pgfpathlineto{\pgfqpoint{1.643924in}{2.775000in}}%
\pgfpathlineto{\pgfqpoint{1.644016in}{2.738102in}}%
\pgfpathlineto{\pgfqpoint{1.644794in}{2.644559in}}%
\pgfpathlineto{\pgfqpoint{1.646937in}{2.644559in}}%
\pgfpathlineto{\pgfqpoint{1.646942in}{2.749550in}}%
\pgfpathlineto{\pgfqpoint{1.647696in}{2.691894in}}%
\pgfpathlineto{\pgfqpoint{1.648734in}{2.765389in}}%
\pgfpathlineto{\pgfqpoint{1.648965in}{2.719559in}}%
\pgfpathlineto{\pgfqpoint{1.648911in}{2.776466in}}%
\pgfpathlineto{\pgfqpoint{1.648979in}{2.726807in}}%
\pgfpathlineto{\pgfqpoint{1.649741in}{2.644559in}}%
\pgfpathlineto{\pgfqpoint{1.651882in}{2.644559in}}%
\pgfpathlineto{\pgfqpoint{1.651887in}{2.753406in}}%
\pgfpathlineto{\pgfqpoint{1.652628in}{2.698604in}}%
\pgfpathlineto{\pgfqpoint{1.652677in}{2.694957in}}%
\pgfpathlineto{\pgfqpoint{1.653023in}{2.723469in}}%
\pgfpathlineto{\pgfqpoint{1.653073in}{2.719801in}}%
\pgfpathlineto{\pgfqpoint{1.653811in}{2.772950in}}%
\pgfpathlineto{\pgfqpoint{1.653887in}{2.746185in}}%
\pgfpathlineto{\pgfqpoint{1.654681in}{2.644559in}}%
\pgfpathlineto{\pgfqpoint{1.656828in}{2.644559in}}%
\pgfpathlineto{\pgfqpoint{1.656833in}{2.749535in}}%
\pgfpathlineto{\pgfqpoint{1.657594in}{2.692784in}}%
\pgfpathlineto{\pgfqpoint{1.658632in}{2.766329in}}%
\pgfpathlineto{\pgfqpoint{1.658802in}{2.776983in}}%
\pgfpathlineto{\pgfqpoint{1.658865in}{2.699345in}}%
\pgfpathlineto{\pgfqpoint{1.659626in}{2.644559in}}%
\pgfpathlineto{\pgfqpoint{1.661773in}{2.644559in}}%
\pgfpathlineto{\pgfqpoint{1.661779in}{2.749531in}}%
\pgfpathlineto{\pgfqpoint{1.662498in}{2.690638in}}%
\pgfpathlineto{\pgfqpoint{1.663636in}{2.769195in}}%
\pgfpathlineto{\pgfqpoint{1.663801in}{2.719959in}}%
\pgfpathlineto{\pgfqpoint{1.663716in}{2.774247in}}%
\pgfpathlineto{\pgfqpoint{1.663815in}{2.727096in}}%
\pgfpathlineto{\pgfqpoint{1.664584in}{2.644559in}}%
\pgfpathlineto{\pgfqpoint{1.666719in}{2.644559in}}%
\pgfpathlineto{\pgfqpoint{1.666724in}{2.749297in}}%
\pgfpathlineto{\pgfqpoint{1.667464in}{2.688446in}}%
\pgfpathlineto{\pgfqpoint{1.668403in}{2.760993in}}%
\pgfpathlineto{\pgfqpoint{1.668693in}{2.779175in}}%
\pgfpathlineto{\pgfqpoint{1.669720in}{2.644559in}}%
\pgfpathlineto{\pgfqpoint{1.671664in}{2.644559in}}%
\pgfpathlineto{\pgfqpoint{1.671669in}{2.753412in}}%
\pgfpathlineto{\pgfqpoint{1.671783in}{2.632008in}}%
\pgfpathlineto{\pgfqpoint{1.672426in}{2.720503in}}%
\pgfpathlineto{\pgfqpoint{1.672476in}{2.675182in}}%
\pgfpathlineto{\pgfqpoint{1.673119in}{2.763926in}}%
\pgfpathlineto{\pgfqpoint{1.673168in}{2.718749in}}%
\pgfpathlineto{\pgfqpoint{1.673601in}{2.794171in}}%
\pgfpathlineto{\pgfqpoint{1.673730in}{2.644559in}}%
\pgfpathlineto{\pgfqpoint{1.673824in}{2.644559in}}%
\pgfpathlineto{\pgfqpoint{1.676610in}{2.644559in}}%
\pgfpathlineto{\pgfqpoint{1.676615in}{2.749561in}}%
\pgfpathlineto{\pgfqpoint{1.677344in}{2.693256in}}%
\pgfpathlineto{\pgfqpoint{1.678557in}{2.773288in}}%
\pgfpathlineto{\pgfqpoint{1.678584in}{2.762678in}}%
\pgfpathlineto{\pgfqpoint{1.679315in}{2.644559in}}%
\pgfpathlineto{\pgfqpoint{1.681555in}{2.644559in}}%
\pgfpathlineto{\pgfqpoint{1.681561in}{2.749613in}}%
\pgfpathlineto{\pgfqpoint{1.682324in}{2.693291in}}%
\pgfpathlineto{\pgfqpoint{1.683461in}{2.771233in}}%
\pgfpathlineto{\pgfqpoint{1.683583in}{2.718638in}}%
\pgfpathlineto{\pgfqpoint{1.683529in}{2.775493in}}%
\pgfpathlineto{\pgfqpoint{1.683597in}{2.727556in}}%
\pgfpathlineto{\pgfqpoint{1.684358in}{2.644559in}}%
\pgfpathlineto{\pgfqpoint{1.686501in}{2.644559in}}%
\pgfpathlineto{\pgfqpoint{1.686505in}{2.753083in}}%
\pgfpathlineto{\pgfqpoint{1.686639in}{2.635315in}}%
\pgfpathlineto{\pgfqpoint{1.687232in}{2.672487in}}%
\pgfpathlineto{\pgfqpoint{1.688073in}{2.768936in}}%
\pgfpathlineto{\pgfqpoint{1.688447in}{2.792378in}}%
\pgfpathlineto{\pgfqpoint{1.688958in}{2.644559in}}%
\pgfpathlineto{\pgfqpoint{1.691446in}{2.644559in}}%
\pgfpathlineto{\pgfqpoint{1.691452in}{2.749537in}}%
\pgfpathlineto{\pgfqpoint{1.692178in}{2.694005in}}%
\pgfpathlineto{\pgfqpoint{1.693393in}{2.772197in}}%
\pgfpathlineto{\pgfqpoint{1.693420in}{2.744216in}}%
\pgfpathlineto{\pgfqpoint{1.693441in}{2.756344in}}%
\pgfpathlineto{\pgfqpoint{1.694189in}{2.644559in}}%
\pgfpathlineto{\pgfqpoint{1.696391in}{2.644559in}}%
\pgfpathlineto{\pgfqpoint{1.696397in}{2.750416in}}%
\pgfpathlineto{\pgfqpoint{1.696476in}{2.634360in}}%
\pgfpathlineto{\pgfqpoint{1.697163in}{2.675641in}}%
\pgfpathlineto{\pgfqpoint{1.698004in}{2.771190in}}%
\pgfpathlineto{\pgfqpoint{1.698365in}{2.793868in}}%
\pgfpathlineto{\pgfqpoint{1.698927in}{2.644559in}}%
\pgfpathlineto{\pgfqpoint{1.701337in}{2.644559in}}%
\pgfpathlineto{\pgfqpoint{1.701342in}{2.749617in}}%
\pgfpathlineto{\pgfqpoint{1.702067in}{2.697024in}}%
\pgfpathlineto{\pgfqpoint{1.702116in}{2.694121in}}%
\pgfpathlineto{\pgfqpoint{1.702364in}{2.715667in}}%
\pgfpathlineto{\pgfqpoint{1.702413in}{2.712759in}}%
\pgfpathlineto{\pgfqpoint{1.703311in}{2.775169in}}%
\pgfpathlineto{\pgfqpoint{1.703311in}{2.748312in}}%
\pgfpathlineto{\pgfqpoint{1.703336in}{2.757613in}}%
\pgfpathlineto{\pgfqpoint{1.704046in}{2.644559in}}%
\pgfpathlineto{\pgfqpoint{1.706282in}{2.644559in}}%
\pgfpathlineto{\pgfqpoint{1.706288in}{2.749427in}}%
\pgfpathlineto{\pgfqpoint{1.707046in}{2.691788in}}%
\pgfpathlineto{\pgfqpoint{1.708084in}{2.766675in}}%
\pgfpathlineto{\pgfqpoint{1.708256in}{2.777470in}}%
\pgfpathlineto{\pgfqpoint{1.709743in}{2.644559in}}%
\pgfpathlineto{\pgfqpoint{1.711228in}{2.644559in}}%
\pgfpathlineto{\pgfqpoint{1.711232in}{2.754264in}}%
\pgfpathlineto{\pgfqpoint{1.711306in}{2.634882in}}%
\pgfpathlineto{\pgfqpoint{1.711991in}{2.675374in}}%
\pgfpathlineto{\pgfqpoint{1.712831in}{2.770213in}}%
\pgfpathlineto{\pgfqpoint{1.713202in}{2.793416in}}%
\pgfpathlineto{\pgfqpoint{1.713761in}{2.644559in}}%
\pgfpathlineto{\pgfqpoint{1.716173in}{2.644559in}}%
\pgfpathlineto{\pgfqpoint{1.716179in}{2.749326in}}%
\pgfpathlineto{\pgfqpoint{1.716924in}{2.691674in}}%
\pgfpathlineto{\pgfqpoint{1.718061in}{2.769798in}}%
\pgfpathlineto{\pgfqpoint{1.718202in}{2.718303in}}%
\pgfpathlineto{\pgfqpoint{1.718147in}{2.775195in}}%
\pgfpathlineto{\pgfqpoint{1.718215in}{2.727117in}}%
\pgfpathlineto{\pgfqpoint{1.718977in}{2.644559in}}%
\pgfpathlineto{\pgfqpoint{1.721119in}{2.644559in}}%
\pgfpathlineto{\pgfqpoint{1.721123in}{2.758941in}}%
\pgfpathlineto{\pgfqpoint{1.721860in}{2.686460in}}%
\pgfpathlineto{\pgfqpoint{1.722800in}{2.764101in}}%
\pgfpathlineto{\pgfqpoint{1.723093in}{2.782487in}}%
\pgfpathlineto{\pgfqpoint{1.723937in}{2.644559in}}%
\pgfpathlineto{\pgfqpoint{1.726064in}{2.644559in}}%
\pgfpathlineto{\pgfqpoint{1.726070in}{2.749522in}}%
\pgfpathlineto{\pgfqpoint{1.726179in}{2.642577in}}%
\pgfpathlineto{\pgfqpoint{1.726820in}{2.708573in}}%
\pgfpathlineto{\pgfqpoint{1.726869in}{2.685558in}}%
\pgfpathlineto{\pgfqpoint{1.727413in}{2.745821in}}%
\pgfpathlineto{\pgfqpoint{1.727463in}{2.722873in}}%
\pgfpathlineto{\pgfqpoint{1.727998in}{2.782504in}}%
\pgfpathlineto{\pgfqpoint{1.728092in}{2.698211in}}%
\pgfpathlineto{\pgfqpoint{1.728107in}{2.746926in}}%
\pgfpathlineto{\pgfqpoint{1.728860in}{2.644559in}}%
\pgfpathlineto{\pgfqpoint{1.731010in}{2.644559in}}%
\pgfpathlineto{\pgfqpoint{1.731014in}{2.756847in}}%
\pgfpathlineto{\pgfqpoint{1.731751in}{2.703160in}}%
\pgfpathlineto{\pgfqpoint{1.731800in}{2.690802in}}%
\pgfpathlineto{\pgfqpoint{1.732344in}{2.740411in}}%
\pgfpathlineto{\pgfqpoint{1.732394in}{2.728114in}}%
\pgfpathlineto{\pgfqpoint{1.732938in}{2.777659in}}%
\pgfpathlineto{\pgfqpoint{1.733039in}{2.697806in}}%
\pgfpathlineto{\pgfqpoint{1.733054in}{2.749451in}}%
\pgfpathlineto{\pgfqpoint{1.733818in}{2.644559in}}%
\pgfpathlineto{\pgfqpoint{1.735955in}{2.644559in}}%
\pgfpathlineto{\pgfqpoint{1.735960in}{2.752236in}}%
\pgfpathlineto{\pgfqpoint{1.736694in}{2.691890in}}%
\pgfpathlineto{\pgfqpoint{1.737929in}{2.774894in}}%
\pgfpathlineto{\pgfqpoint{1.738688in}{2.644559in}}%
\pgfpathlineto{\pgfqpoint{1.740901in}{2.644559in}}%
\pgfpathlineto{\pgfqpoint{1.740905in}{2.749174in}}%
\pgfpathlineto{\pgfqpoint{1.740996in}{2.626837in}}%
\pgfpathlineto{\pgfqpoint{1.741634in}{2.722006in}}%
\pgfpathlineto{\pgfqpoint{1.741683in}{2.668899in}}%
\pgfpathlineto{\pgfqpoint{1.742326in}{2.765403in}}%
\pgfpathlineto{\pgfqpoint{1.742375in}{2.712491in}}%
\pgfpathlineto{\pgfqpoint{1.742874in}{2.799762in}}%
\pgfpathlineto{\pgfqpoint{1.742965in}{2.644559in}}%
\pgfpathlineto{\pgfqpoint{1.743088in}{2.644559in}}%
\pgfpathlineto{\pgfqpoint{1.745846in}{2.644559in}}%
\pgfpathlineto{\pgfqpoint{1.745852in}{2.749312in}}%
\pgfpathlineto{\pgfqpoint{1.746600in}{2.698815in}}%
\pgfpathlineto{\pgfqpoint{1.746649in}{2.695285in}}%
\pgfpathlineto{\pgfqpoint{1.746995in}{2.723677in}}%
\pgfpathlineto{\pgfqpoint{1.747045in}{2.720131in}}%
\pgfpathlineto{\pgfqpoint{1.747778in}{2.772894in}}%
\pgfpathlineto{\pgfqpoint{1.747862in}{2.731953in}}%
\pgfpathlineto{\pgfqpoint{1.748711in}{2.644559in}}%
\pgfpathlineto{\pgfqpoint{1.750792in}{2.644559in}}%
\pgfpathlineto{\pgfqpoint{1.750797in}{2.749401in}}%
\pgfpathlineto{\pgfqpoint{1.751545in}{2.699606in}}%
\pgfpathlineto{\pgfqpoint{1.751594in}{2.695011in}}%
\pgfpathlineto{\pgfqpoint{1.751940in}{2.724468in}}%
\pgfpathlineto{\pgfqpoint{1.751990in}{2.719858in}}%
\pgfpathlineto{\pgfqpoint{1.752724in}{2.773693in}}%
\pgfpathlineto{\pgfqpoint{1.752808in}{2.733082in}}%
\pgfpathlineto{\pgfqpoint{1.753607in}{2.644559in}}%
\pgfpathlineto{\pgfqpoint{1.755737in}{2.644559in}}%
\pgfpathlineto{\pgfqpoint{1.755742in}{2.749677in}}%
\pgfpathlineto{\pgfqpoint{1.756473in}{2.698080in}}%
\pgfpathlineto{\pgfqpoint{1.756523in}{2.694285in}}%
\pgfpathlineto{\pgfqpoint{1.756869in}{2.722941in}}%
\pgfpathlineto{\pgfqpoint{1.756918in}{2.719134in}}%
\pgfpathlineto{\pgfqpoint{1.757711in}{2.775840in}}%
\pgfpathlineto{\pgfqpoint{1.757753in}{2.732808in}}%
\pgfpathlineto{\pgfqpoint{1.758591in}{2.644559in}}%
\pgfpathlineto{\pgfqpoint{1.760682in}{2.644559in}}%
\pgfpathlineto{\pgfqpoint{1.760688in}{2.749509in}}%
\pgfpathlineto{\pgfqpoint{1.761442in}{2.691858in}}%
\pgfpathlineto{\pgfqpoint{1.762480in}{2.765343in}}%
\pgfpathlineto{\pgfqpoint{1.762711in}{2.719425in}}%
\pgfpathlineto{\pgfqpoint{1.762656in}{2.776394in}}%
\pgfpathlineto{\pgfqpoint{1.762724in}{2.726807in}}%
\pgfpathlineto{\pgfqpoint{1.763487in}{2.644559in}}%
\pgfpathlineto{\pgfqpoint{1.765628in}{2.644559in}}%
\pgfpathlineto{\pgfqpoint{1.765633in}{2.749409in}}%
\pgfpathlineto{\pgfqpoint{1.766378in}{2.698683in}}%
\pgfpathlineto{\pgfqpoint{1.766428in}{2.695182in}}%
\pgfpathlineto{\pgfqpoint{1.766774in}{2.723545in}}%
\pgfpathlineto{\pgfqpoint{1.766823in}{2.720028in}}%
\pgfpathlineto{\pgfqpoint{1.767559in}{2.772857in}}%
\pgfpathlineto{\pgfqpoint{1.767644in}{2.733110in}}%
\pgfpathlineto{\pgfqpoint{1.768444in}{2.644559in}}%
\pgfpathlineto{\pgfqpoint{1.770573in}{2.644559in}}%
\pgfpathlineto{\pgfqpoint{1.770577in}{2.752332in}}%
\pgfpathlineto{\pgfqpoint{1.770663in}{2.631685in}}%
\pgfpathlineto{\pgfqpoint{1.771299in}{2.717786in}}%
\pgfpathlineto{\pgfqpoint{1.771348in}{2.673356in}}%
\pgfpathlineto{\pgfqpoint{1.771991in}{2.761163in}}%
\pgfpathlineto{\pgfqpoint{1.772041in}{2.716968in}}%
\pgfpathlineto{\pgfqpoint{1.772547in}{2.795982in}}%
\pgfpathlineto{\pgfqpoint{1.772633in}{2.644559in}}%
\pgfpathlineto{\pgfqpoint{1.772710in}{2.644559in}}%
\pgfpathlineto{\pgfqpoint{1.775519in}{2.644560in}}%
\pgfpathlineto{\pgfqpoint{1.775522in}{2.750932in}}%
\pgfpathlineto{\pgfqpoint{1.775586in}{2.622044in}}%
\pgfpathlineto{\pgfqpoint{1.776317in}{2.732508in}}%
\pgfpathlineto{\pgfqpoint{1.776366in}{2.667480in}}%
\pgfpathlineto{\pgfqpoint{1.777009in}{2.775859in}}%
\pgfpathlineto{\pgfqpoint{1.777059in}{2.711118in}}%
\pgfpathlineto{\pgfqpoint{1.777493in}{2.806107in}}%
\pgfpathlineto{\pgfqpoint{1.777578in}{2.644559in}}%
\pgfpathlineto{\pgfqpoint{1.777757in}{2.644559in}}%
\pgfpathlineto{\pgfqpoint{1.780464in}{2.644559in}}%
\pgfpathlineto{\pgfqpoint{1.780470in}{2.749832in}}%
\pgfpathlineto{\pgfqpoint{1.781202in}{2.691349in}}%
\pgfpathlineto{\pgfqpoint{1.782339in}{2.770359in}}%
\pgfpathlineto{\pgfqpoint{1.782438in}{2.776572in}}%
\pgfpathlineto{\pgfqpoint{1.782502in}{2.714142in}}%
\pgfpathlineto{\pgfqpoint{1.783276in}{2.644559in}}%
\pgfpathlineto{\pgfqpoint{1.785410in}{2.644559in}}%
\pgfpathlineto{\pgfqpoint{1.785414in}{2.751442in}}%
\pgfpathlineto{\pgfqpoint{1.786183in}{2.694368in}}%
\pgfpathlineto{\pgfqpoint{1.787384in}{2.774496in}}%
\pgfpathlineto{\pgfqpoint{1.787384in}{2.768148in}}%
\pgfpathlineto{\pgfqpoint{1.788120in}{2.644559in}}%
\pgfpathlineto{\pgfqpoint{1.790355in}{2.644559in}}%
\pgfpathlineto{\pgfqpoint{1.790360in}{2.753140in}}%
\pgfpathlineto{\pgfqpoint{1.791124in}{2.690723in}}%
\pgfpathlineto{\pgfqpoint{1.792064in}{2.761619in}}%
\pgfpathlineto{\pgfqpoint{1.792329in}{2.778249in}}%
\pgfpathlineto{\pgfqpoint{1.793499in}{2.644559in}}%
\pgfpathlineto{\pgfqpoint{1.795301in}{2.644559in}}%
\pgfpathlineto{\pgfqpoint{1.795306in}{2.749495in}}%
\pgfpathlineto{\pgfqpoint{1.796059in}{2.699263in}}%
\pgfpathlineto{\pgfqpoint{1.796108in}{2.695764in}}%
\pgfpathlineto{\pgfqpoint{1.796454in}{2.724126in}}%
\pgfpathlineto{\pgfqpoint{1.796504in}{2.720610in}}%
\pgfpathlineto{\pgfqpoint{1.797235in}{2.773206in}}%
\pgfpathlineto{\pgfqpoint{1.797317in}{2.733196in}}%
\pgfpathlineto{\pgfqpoint{1.798117in}{2.644559in}}%
\pgfpathlineto{\pgfqpoint{1.800246in}{2.644559in}}%
\pgfpathlineto{\pgfqpoint{1.800251in}{2.749709in}}%
\pgfpathlineto{\pgfqpoint{1.801006in}{2.692798in}}%
\pgfpathlineto{\pgfqpoint{1.802144in}{2.770839in}}%
\pgfpathlineto{\pgfqpoint{1.802275in}{2.716474in}}%
\pgfpathlineto{\pgfqpoint{1.802220in}{2.775636in}}%
\pgfpathlineto{\pgfqpoint{1.802289in}{2.730028in}}%
\pgfpathlineto{\pgfqpoint{1.803046in}{2.644559in}}%
\pgfpathlineto{\pgfqpoint{1.805192in}{2.644559in}}%
\pgfpathlineto{\pgfqpoint{1.805197in}{2.749808in}}%
\pgfpathlineto{\pgfqpoint{1.805952in}{2.692990in}}%
\pgfpathlineto{\pgfqpoint{1.807090in}{2.771066in}}%
\pgfpathlineto{\pgfqpoint{1.807220in}{2.716524in}}%
\pgfpathlineto{\pgfqpoint{1.807165in}{2.775806in}}%
\pgfpathlineto{\pgfqpoint{1.807234in}{2.730289in}}%
\pgfpathlineto{\pgfqpoint{1.807992in}{2.644559in}}%
\pgfpathlineto{\pgfqpoint{1.810137in}{2.644559in}}%
\pgfpathlineto{\pgfqpoint{1.810142in}{2.749807in}}%
\pgfpathlineto{\pgfqpoint{1.810871in}{2.697529in}}%
\pgfpathlineto{\pgfqpoint{1.810920in}{2.694717in}}%
\pgfpathlineto{\pgfqpoint{1.811168in}{2.716172in}}%
\pgfpathlineto{\pgfqpoint{1.811217in}{2.713355in}}%
\pgfpathlineto{\pgfqpoint{1.812111in}{2.775430in}}%
\pgfpathlineto{\pgfqpoint{1.812137in}{2.758417in}}%
\pgfpathlineto{\pgfqpoint{1.812906in}{2.644559in}}%
\pgfpathlineto{\pgfqpoint{1.815082in}{2.644559in}}%
\pgfpathlineto{\pgfqpoint{1.815088in}{2.749704in}}%
\pgfpathlineto{\pgfqpoint{1.815154in}{2.643260in}}%
\pgfpathlineto{\pgfqpoint{1.815842in}{2.683263in}}%
\pgfpathlineto{\pgfqpoint{1.816683in}{2.761598in}}%
\pgfpathlineto{\pgfqpoint{1.817056in}{2.785028in}}%
\pgfpathlineto{\pgfqpoint{1.817681in}{2.644559in}}%
\pgfpathlineto{\pgfqpoint{1.820028in}{2.644559in}}%
\pgfpathlineto{\pgfqpoint{1.820033in}{2.749582in}}%
\pgfpathlineto{\pgfqpoint{1.820787in}{2.691964in}}%
\pgfpathlineto{\pgfqpoint{1.821826in}{2.765461in}}%
\pgfpathlineto{\pgfqpoint{1.822056in}{2.719647in}}%
\pgfpathlineto{\pgfqpoint{1.822002in}{2.776517in}}%
\pgfpathlineto{\pgfqpoint{1.822070in}{2.726818in}}%
\pgfpathlineto{\pgfqpoint{1.822833in}{2.644559in}}%
\pgfpathlineto{\pgfqpoint{1.824973in}{2.644559in}}%
\pgfpathlineto{\pgfqpoint{1.824979in}{2.749599in}}%
\pgfpathlineto{\pgfqpoint{1.825704in}{2.690750in}}%
\pgfpathlineto{\pgfqpoint{1.826841in}{2.768845in}}%
\pgfpathlineto{\pgfqpoint{1.827003in}{2.713095in}}%
\pgfpathlineto{\pgfqpoint{1.826919in}{2.773719in}}%
\pgfpathlineto{\pgfqpoint{1.827017in}{2.733547in}}%
\pgfpathlineto{\pgfqpoint{1.827783in}{2.644559in}}%
\pgfpathlineto{\pgfqpoint{1.829919in}{2.644559in}}%
\pgfpathlineto{\pgfqpoint{1.829924in}{2.749628in}}%
\pgfpathlineto{\pgfqpoint{1.830660in}{2.698369in}}%
\pgfpathlineto{\pgfqpoint{1.830709in}{2.694379in}}%
\pgfpathlineto{\pgfqpoint{1.831056in}{2.723228in}}%
\pgfpathlineto{\pgfqpoint{1.831105in}{2.719229in}}%
\pgfpathlineto{\pgfqpoint{1.831847in}{2.772944in}}%
\pgfpathlineto{\pgfqpoint{1.831935in}{2.735652in}}%
\pgfpathlineto{\pgfqpoint{1.832763in}{2.644559in}}%
\pgfpathlineto{\pgfqpoint{1.834864in}{2.644559in}}%
\pgfpathlineto{\pgfqpoint{1.834870in}{2.749841in}}%
\pgfpathlineto{\pgfqpoint{1.835590in}{2.697622in}}%
\pgfpathlineto{\pgfqpoint{1.835639in}{2.693820in}}%
\pgfpathlineto{\pgfqpoint{1.835986in}{2.722482in}}%
\pgfpathlineto{\pgfqpoint{1.836035in}{2.718669in}}%
\pgfpathlineto{\pgfqpoint{1.836838in}{2.776046in}}%
\pgfpathlineto{\pgfqpoint{1.836880in}{2.734654in}}%
\pgfpathlineto{\pgfqpoint{1.837719in}{2.644559in}}%
\pgfpathlineto{\pgfqpoint{1.839810in}{2.644559in}}%
\pgfpathlineto{\pgfqpoint{1.839815in}{2.749671in}}%
\pgfpathlineto{\pgfqpoint{1.839917in}{2.642283in}}%
\pgfpathlineto{\pgfqpoint{1.840557in}{2.708623in}}%
\pgfpathlineto{\pgfqpoint{1.840607in}{2.685005in}}%
\pgfpathlineto{\pgfqpoint{1.841151in}{2.745871in}}%
\pgfpathlineto{\pgfqpoint{1.841200in}{2.722320in}}%
\pgfpathlineto{\pgfqpoint{1.841739in}{2.782795in}}%
\pgfpathlineto{\pgfqpoint{1.841838in}{2.698876in}}%
\pgfpathlineto{\pgfqpoint{1.841853in}{2.746794in}}%
\pgfpathlineto{\pgfqpoint{1.842611in}{2.644559in}}%
\pgfpathlineto{\pgfqpoint{1.844755in}{2.644559in}}%
\pgfpathlineto{\pgfqpoint{1.844760in}{2.753243in}}%
\pgfpathlineto{\pgfqpoint{1.845529in}{2.694891in}}%
\pgfpathlineto{\pgfqpoint{1.846729in}{2.774518in}}%
\pgfpathlineto{\pgfqpoint{1.846729in}{2.761787in}}%
\pgfpathlineto{\pgfqpoint{1.847466in}{2.644559in}}%
\pgfpathlineto{\pgfqpoint{1.849701in}{2.644559in}}%
\pgfpathlineto{\pgfqpoint{1.849705in}{2.751436in}}%
\pgfpathlineto{\pgfqpoint{1.850439in}{2.692038in}}%
\pgfpathlineto{\pgfqpoint{1.851674in}{2.774486in}}%
\pgfpathlineto{\pgfqpoint{1.851674in}{2.769295in}}%
\pgfpathlineto{\pgfqpoint{1.852433in}{2.644559in}}%
\pgfpathlineto{\pgfqpoint{1.854646in}{2.644559in}}%
\pgfpathlineto{\pgfqpoint{1.854650in}{2.748918in}}%
\pgfpathlineto{\pgfqpoint{1.854724in}{2.640181in}}%
\pgfpathlineto{\pgfqpoint{1.855409in}{2.680870in}}%
\pgfpathlineto{\pgfqpoint{1.856250in}{2.763444in}}%
\pgfpathlineto{\pgfqpoint{1.856620in}{2.786621in}}%
\pgfpathlineto{\pgfqpoint{1.857230in}{2.644559in}}%
\pgfpathlineto{\pgfqpoint{1.859592in}{2.644559in}}%
\pgfpathlineto{\pgfqpoint{1.859597in}{2.749587in}}%
\pgfpathlineto{\pgfqpoint{1.859709in}{2.642957in}}%
\pgfpathlineto{\pgfqpoint{1.860352in}{2.708903in}}%
\pgfpathlineto{\pgfqpoint{1.860401in}{2.686101in}}%
\pgfpathlineto{\pgfqpoint{1.860945in}{2.746151in}}%
\pgfpathlineto{\pgfqpoint{1.860995in}{2.723415in}}%
\pgfpathlineto{\pgfqpoint{1.861527in}{2.782681in}}%
\pgfpathlineto{\pgfqpoint{1.861620in}{2.698766in}}%
\pgfpathlineto{\pgfqpoint{1.861635in}{2.746702in}}%
\pgfpathlineto{\pgfqpoint{1.862393in}{2.644559in}}%
\pgfpathlineto{\pgfqpoint{1.864537in}{2.644559in}}%
\pgfpathlineto{\pgfqpoint{1.864541in}{2.753186in}}%
\pgfpathlineto{\pgfqpoint{1.865264in}{2.695189in}}%
\pgfpathlineto{\pgfqpoint{1.865511in}{2.708368in}}%
\pgfpathlineto{\pgfqpoint{1.866511in}{2.773599in}}%
\pgfpathlineto{\pgfqpoint{1.866538in}{2.758639in}}%
\pgfpathlineto{\pgfqpoint{1.867298in}{2.644559in}}%
\pgfpathlineto{\pgfqpoint{1.869482in}{2.644559in}}%
\pgfpathlineto{\pgfqpoint{1.869487in}{2.751080in}}%
\pgfpathlineto{\pgfqpoint{1.870216in}{2.693108in}}%
\pgfpathlineto{\pgfqpoint{1.871430in}{2.771349in}}%
\pgfpathlineto{\pgfqpoint{1.871457in}{2.743403in}}%
\pgfpathlineto{\pgfqpoint{1.871478in}{2.755597in}}%
\pgfpathlineto{\pgfqpoint{1.872185in}{2.644559in}}%
\pgfpathlineto{\pgfqpoint{1.874428in}{2.644559in}}%
\pgfpathlineto{\pgfqpoint{1.874433in}{2.749636in}}%
\pgfpathlineto{\pgfqpoint{1.875174in}{2.691691in}}%
\pgfpathlineto{\pgfqpoint{1.876311in}{2.769974in}}%
\pgfpathlineto{\pgfqpoint{1.876456in}{2.717029in}}%
\pgfpathlineto{\pgfqpoint{1.876402in}{2.775670in}}%
\pgfpathlineto{\pgfqpoint{1.876470in}{2.729244in}}%
\pgfpathlineto{\pgfqpoint{1.877233in}{2.644559in}}%
\pgfpathlineto{\pgfqpoint{1.879373in}{2.644559in}}%
\pgfpathlineto{\pgfqpoint{1.879379in}{2.749712in}}%
\pgfpathlineto{\pgfqpoint{1.880104in}{2.689293in}}%
\pgfpathlineto{\pgfqpoint{1.881143in}{2.764370in}}%
\pgfpathlineto{\pgfqpoint{1.881319in}{2.775454in}}%
\pgfpathlineto{\pgfqpoint{1.882777in}{2.644559in}}%
\pgfpathlineto{\pgfqpoint{1.884319in}{2.644559in}}%
\pgfpathlineto{\pgfqpoint{1.884323in}{2.750887in}}%
\pgfpathlineto{\pgfqpoint{1.885054in}{2.692637in}}%
\pgfpathlineto{\pgfqpoint{1.886267in}{2.771941in}}%
\pgfpathlineto{\pgfqpoint{1.886293in}{2.761660in}}%
\pgfpathlineto{\pgfqpoint{1.887025in}{2.644559in}}%
\pgfpathlineto{\pgfqpoint{1.889264in}{2.644559in}}%
\pgfpathlineto{\pgfqpoint{1.889270in}{2.749620in}}%
\pgfpathlineto{\pgfqpoint{1.889989in}{2.690404in}}%
\pgfpathlineto{\pgfqpoint{1.891127in}{2.768580in}}%
\pgfpathlineto{\pgfqpoint{1.891293in}{2.714715in}}%
\pgfpathlineto{\pgfqpoint{1.891207in}{2.773630in}}%
\pgfpathlineto{\pgfqpoint{1.891307in}{2.731706in}}%
\pgfpathlineto{\pgfqpoint{1.892065in}{2.644559in}}%
\pgfpathlineto{\pgfqpoint{1.894210in}{2.644559in}}%
\pgfpathlineto{\pgfqpoint{1.894215in}{2.749742in}}%
\pgfpathlineto{\pgfqpoint{1.894280in}{2.643731in}}%
\pgfpathlineto{\pgfqpoint{1.894966in}{2.683450in}}%
\pgfpathlineto{\pgfqpoint{1.895807in}{2.761164in}}%
\pgfpathlineto{\pgfqpoint{1.896184in}{2.784787in}}%
\pgfpathlineto{\pgfqpoint{1.896809in}{2.644559in}}%
\pgfpathlineto{\pgfqpoint{1.899155in}{2.644559in}}%
\pgfpathlineto{\pgfqpoint{1.899161in}{2.749300in}}%
\pgfpathlineto{\pgfqpoint{1.899920in}{2.699625in}}%
\pgfpathlineto{\pgfqpoint{1.899969in}{2.695430in}}%
\pgfpathlineto{\pgfqpoint{1.900315in}{2.724486in}}%
\pgfpathlineto{\pgfqpoint{1.900365in}{2.720278in}}%
\pgfpathlineto{\pgfqpoint{1.901093in}{2.773357in}}%
\pgfpathlineto{\pgfqpoint{1.901179in}{2.737723in}}%
\pgfpathlineto{\pgfqpoint{1.901958in}{2.644559in}}%
\pgfpathlineto{\pgfqpoint{1.904101in}{2.644559in}}%
\pgfpathlineto{\pgfqpoint{1.904106in}{2.749634in}}%
\pgfpathlineto{\pgfqpoint{1.904846in}{2.691781in}}%
\pgfpathlineto{\pgfqpoint{1.905983in}{2.769692in}}%
\pgfpathlineto{\pgfqpoint{1.906131in}{2.708513in}}%
\pgfpathlineto{\pgfqpoint{1.906074in}{2.775414in}}%
\pgfpathlineto{\pgfqpoint{1.906146in}{2.738967in}}%
\pgfpathlineto{\pgfqpoint{1.906911in}{2.644559in}}%
\pgfpathlineto{\pgfqpoint{1.909046in}{2.644559in}}%
\pgfpathlineto{\pgfqpoint{1.909051in}{2.749559in}}%
\pgfpathlineto{\pgfqpoint{1.909790in}{2.697754in}}%
\pgfpathlineto{\pgfqpoint{1.909840in}{2.695036in}}%
\pgfpathlineto{\pgfqpoint{1.910087in}{2.716398in}}%
\pgfpathlineto{\pgfqpoint{1.910137in}{2.713674in}}%
\pgfpathlineto{\pgfqpoint{1.910974in}{2.772111in}}%
\pgfpathlineto{\pgfqpoint{1.911045in}{2.757590in}}%
\pgfpathlineto{\pgfqpoint{1.911809in}{2.644559in}}%
\pgfpathlineto{\pgfqpoint{1.913992in}{2.644559in}}%
\pgfpathlineto{\pgfqpoint{1.913997in}{2.749671in}}%
\pgfpathlineto{\pgfqpoint{1.914092in}{2.642122in}}%
\pgfpathlineto{\pgfqpoint{1.914732in}{2.708197in}}%
\pgfpathlineto{\pgfqpoint{1.914782in}{2.684553in}}%
\pgfpathlineto{\pgfqpoint{1.915326in}{2.745445in}}%
\pgfpathlineto{\pgfqpoint{1.915375in}{2.721868in}}%
\pgfpathlineto{\pgfqpoint{1.915919in}{2.782691in}}%
\pgfpathlineto{\pgfqpoint{1.916020in}{2.698082in}}%
\pgfpathlineto{\pgfqpoint{1.916034in}{2.747441in}}%
\pgfpathlineto{\pgfqpoint{1.916785in}{2.644559in}}%
\pgfpathlineto{\pgfqpoint{1.918937in}{2.644559in}}%
\pgfpathlineto{\pgfqpoint{1.918942in}{2.749462in}}%
\pgfpathlineto{\pgfqpoint{1.919680in}{2.688547in}}%
\pgfpathlineto{\pgfqpoint{1.920620in}{2.761137in}}%
\pgfpathlineto{\pgfqpoint{1.920911in}{2.779398in}}%
\pgfpathlineto{\pgfqpoint{1.921940in}{2.644559in}}%
\pgfpathlineto{\pgfqpoint{1.923882in}{2.644559in}}%
\pgfpathlineto{\pgfqpoint{1.923888in}{2.749261in}}%
\pgfpathlineto{\pgfqpoint{1.923968in}{2.632626in}}%
\pgfpathlineto{\pgfqpoint{1.924654in}{2.674060in}}%
\pgfpathlineto{\pgfqpoint{1.925494in}{2.771723in}}%
\pgfpathlineto{\pgfqpoint{1.925856in}{2.794392in}}%
\pgfpathlineto{\pgfqpoint{1.926401in}{2.644559in}}%
\pgfpathlineto{\pgfqpoint{1.928829in}{2.644586in}}%
\pgfpathlineto{\pgfqpoint{1.928831in}{2.778691in}}%
\pgfpathlineto{\pgfqpoint{1.929577in}{2.694497in}}%
\pgfpathlineto{\pgfqpoint{1.930802in}{2.773260in}}%
\pgfpathlineto{\pgfqpoint{1.930802in}{2.748709in}}%
\pgfpathlineto{\pgfqpoint{1.930826in}{2.757478in}}%
\pgfpathlineto{\pgfqpoint{1.931530in}{2.644559in}}%
\pgfpathlineto{\pgfqpoint{1.933774in}{2.644595in}}%
\pgfpathlineto{\pgfqpoint{1.933779in}{2.772138in}}%
\pgfpathlineto{\pgfqpoint{1.933895in}{2.637012in}}%
\pgfpathlineto{\pgfqpoint{1.934538in}{2.715327in}}%
\pgfpathlineto{\pgfqpoint{1.934587in}{2.680236in}}%
\pgfpathlineto{\pgfqpoint{1.935230in}{2.758743in}}%
\pgfpathlineto{\pgfqpoint{1.935280in}{2.723809in}}%
\pgfpathlineto{\pgfqpoint{1.935711in}{2.788905in}}%
\pgfpathlineto{\pgfqpoint{1.935836in}{2.644559in}}%
\pgfpathlineto{\pgfqpoint{1.935933in}{2.644559in}}%
\pgfpathlineto{\pgfqpoint{1.938719in}{2.644559in}}%
\pgfpathlineto{\pgfqpoint{1.938724in}{2.749371in}}%
\pgfpathlineto{\pgfqpoint{1.938785in}{2.642136in}}%
\pgfpathlineto{\pgfqpoint{1.939469in}{2.681267in}}%
\pgfpathlineto{\pgfqpoint{1.940310in}{2.762153in}}%
\pgfpathlineto{\pgfqpoint{1.940693in}{2.786172in}}%
\pgfpathlineto{\pgfqpoint{1.941277in}{2.644559in}}%
\pgfpathlineto{\pgfqpoint{1.943664in}{2.644559in}}%
\pgfpathlineto{\pgfqpoint{1.943669in}{2.751212in}}%
\pgfpathlineto{\pgfqpoint{1.944426in}{2.699796in}}%
\pgfpathlineto{\pgfqpoint{1.944475in}{2.695174in}}%
\pgfpathlineto{\pgfqpoint{1.944821in}{2.724659in}}%
\pgfpathlineto{\pgfqpoint{1.944871in}{2.720019in}}%
\pgfpathlineto{\pgfqpoint{1.945601in}{2.773629in}}%
\pgfpathlineto{\pgfqpoint{1.945680in}{2.729445in}}%
\pgfpathlineto{\pgfqpoint{1.946568in}{2.644559in}}%
\pgfpathlineto{\pgfqpoint{1.948610in}{2.644559in}}%
\pgfpathlineto{\pgfqpoint{1.948615in}{2.749624in}}%
\pgfpathlineto{\pgfqpoint{1.948694in}{2.642263in}}%
\pgfpathlineto{\pgfqpoint{1.949383in}{2.683600in}}%
\pgfpathlineto{\pgfqpoint{1.950223in}{2.762673in}}%
\pgfpathlineto{\pgfqpoint{1.950584in}{2.785263in}}%
\pgfpathlineto{\pgfqpoint{1.951207in}{2.644559in}}%
\pgfpathlineto{\pgfqpoint{1.953555in}{2.644559in}}%
\pgfpathlineto{\pgfqpoint{1.953561in}{2.749582in}}%
\pgfpathlineto{\pgfqpoint{1.954305in}{2.692353in}}%
\pgfpathlineto{\pgfqpoint{1.955529in}{2.774993in}}%
\pgfpathlineto{\pgfqpoint{1.956265in}{2.644559in}}%
\pgfpathlineto{\pgfqpoint{1.958501in}{2.644559in}}%
\pgfpathlineto{\pgfqpoint{1.958506in}{2.749572in}}%
\pgfpathlineto{\pgfqpoint{1.958566in}{2.643567in}}%
\pgfpathlineto{\pgfqpoint{1.959246in}{2.682196in}}%
\pgfpathlineto{\pgfqpoint{1.960087in}{2.760480in}}%
\pgfpathlineto{\pgfqpoint{1.960474in}{2.784783in}}%
\pgfpathlineto{\pgfqpoint{1.961100in}{2.644559in}}%
\pgfpathlineto{\pgfqpoint{1.963446in}{2.644559in}}%
\pgfpathlineto{\pgfqpoint{1.963452in}{2.749322in}}%
\pgfpathlineto{\pgfqpoint{1.964210in}{2.699612in}}%
\pgfpathlineto{\pgfqpoint{1.964259in}{2.695476in}}%
\pgfpathlineto{\pgfqpoint{1.964605in}{2.724473in}}%
\pgfpathlineto{\pgfqpoint{1.964655in}{2.720324in}}%
\pgfpathlineto{\pgfqpoint{1.965384in}{2.773367in}}%
\pgfpathlineto{\pgfqpoint{1.965470in}{2.737741in}}%
\pgfpathlineto{\pgfqpoint{1.966248in}{2.644559in}}%
\pgfpathlineto{\pgfqpoint{1.968392in}{2.644559in}}%
\pgfpathlineto{\pgfqpoint{1.968397in}{2.749606in}}%
\pgfpathlineto{\pgfqpoint{1.968524in}{2.643189in}}%
\pgfpathlineto{\pgfqpoint{1.969117in}{2.680296in}}%
\pgfpathlineto{\pgfqpoint{1.969958in}{2.760161in}}%
\pgfpathlineto{\pgfqpoint{1.970335in}{2.783808in}}%
\pgfpathlineto{\pgfqpoint{1.970955in}{2.644559in}}%
\pgfpathlineto{\pgfqpoint{1.973337in}{2.644559in}}%
\pgfpathlineto{\pgfqpoint{1.973341in}{2.750084in}}%
\pgfpathlineto{\pgfqpoint{1.973356in}{2.643822in}}%
\pgfpathlineto{\pgfqpoint{1.974068in}{2.688094in}}%
\pgfpathlineto{\pgfqpoint{1.975008in}{2.758657in}}%
\pgfpathlineto{\pgfqpoint{1.975283in}{2.775951in}}%
\pgfpathlineto{\pgfqpoint{1.976441in}{2.644559in}}%
\pgfpathlineto{\pgfqpoint{1.978282in}{2.644559in}}%
\pgfpathlineto{\pgfqpoint{1.978288in}{2.749686in}}%
\pgfpathlineto{\pgfqpoint{1.979008in}{2.697416in}}%
\pgfpathlineto{\pgfqpoint{1.979058in}{2.693636in}}%
\pgfpathlineto{\pgfqpoint{1.979404in}{2.722276in}}%
\pgfpathlineto{\pgfqpoint{1.979454in}{2.718486in}}%
\pgfpathlineto{\pgfqpoint{1.980256in}{2.775817in}}%
\pgfpathlineto{\pgfqpoint{1.980298in}{2.733839in}}%
\pgfpathlineto{\pgfqpoint{1.981137in}{2.644559in}}%
\pgfpathlineto{\pgfqpoint{1.983228in}{2.644559in}}%
\pgfpathlineto{\pgfqpoint{1.983233in}{2.749307in}}%
\pgfpathlineto{\pgfqpoint{1.983302in}{2.640833in}}%
\pgfpathlineto{\pgfqpoint{1.983989in}{2.681168in}}%
\pgfpathlineto{\pgfqpoint{1.984830in}{2.763536in}}%
\pgfpathlineto{\pgfqpoint{1.985202in}{2.786864in}}%
\pgfpathlineto{\pgfqpoint{1.985779in}{2.644559in}}%
\pgfpathlineto{\pgfqpoint{1.988173in}{2.644559in}}%
\pgfpathlineto{\pgfqpoint{1.988178in}{2.750024in}}%
\pgfpathlineto{\pgfqpoint{1.988270in}{2.626880in}}%
\pgfpathlineto{\pgfqpoint{1.988908in}{2.722314in}}%
\pgfpathlineto{\pgfqpoint{1.988957in}{2.669034in}}%
\pgfpathlineto{\pgfqpoint{1.989600in}{2.765711in}}%
\pgfpathlineto{\pgfqpoint{1.989650in}{2.712627in}}%
\pgfpathlineto{\pgfqpoint{1.990147in}{2.799968in}}%
\pgfpathlineto{\pgfqpoint{1.990238in}{2.644559in}}%
\pgfpathlineto{\pgfqpoint{1.990360in}{2.644559in}}%
\pgfpathlineto{\pgfqpoint{1.993119in}{2.644559in}}%
\pgfpathlineto{\pgfqpoint{1.993124in}{2.749204in}}%
\pgfpathlineto{\pgfqpoint{1.993181in}{2.641697in}}%
\pgfpathlineto{\pgfqpoint{1.993858in}{2.679640in}}%
\pgfpathlineto{\pgfqpoint{1.994699in}{2.761235in}}%
\pgfpathlineto{\pgfqpoint{1.995093in}{2.785929in}}%
\pgfpathlineto{\pgfqpoint{1.995678in}{2.644559in}}%
\pgfpathlineto{\pgfqpoint{1.998064in}{2.644559in}}%
\pgfpathlineto{\pgfqpoint{1.998069in}{2.752115in}}%
\pgfpathlineto{\pgfqpoint{1.998142in}{2.630910in}}%
\pgfpathlineto{\pgfqpoint{1.998827in}{2.671385in}}%
\pgfpathlineto{\pgfqpoint{1.999668in}{2.773534in}}%
\pgfpathlineto{\pgfqpoint{2.000038in}{2.796744in}}%
\pgfpathlineto{\pgfqpoint{2.000547in}{2.644559in}}%
\pgfpathlineto{\pgfqpoint{2.003010in}{2.644559in}}%
\pgfpathlineto{\pgfqpoint{2.003015in}{2.749837in}}%
\pgfpathlineto{\pgfqpoint{2.003735in}{2.690027in}}%
\pgfpathlineto{\pgfqpoint{2.004773in}{2.763541in}}%
\pgfpathlineto{\pgfqpoint{2.005038in}{2.716480in}}%
\pgfpathlineto{\pgfqpoint{2.004953in}{2.774807in}}%
\pgfpathlineto{\pgfqpoint{2.005052in}{2.730515in}}%
\pgfpathlineto{\pgfqpoint{2.005808in}{2.644559in}}%
\pgfpathlineto{\pgfqpoint{2.007955in}{2.644559in}}%
\pgfpathlineto{\pgfqpoint{2.007961in}{2.749389in}}%
\pgfpathlineto{\pgfqpoint{2.008014in}{2.642717in}}%
\pgfpathlineto{\pgfqpoint{2.008689in}{2.680007in}}%
\pgfpathlineto{\pgfqpoint{2.009530in}{2.760775in}}%
\pgfpathlineto{\pgfqpoint{2.009903in}{2.784155in}}%
\pgfpathlineto{\pgfqpoint{2.010519in}{2.644559in}}%
\pgfpathlineto{\pgfqpoint{2.012901in}{2.644559in}}%
\pgfpathlineto{\pgfqpoint{2.012905in}{2.750786in}}%
\pgfpathlineto{\pgfqpoint{2.013639in}{2.691855in}}%
\pgfpathlineto{\pgfqpoint{2.014874in}{2.774499in}}%
\pgfpathlineto{\pgfqpoint{2.014874in}{2.769209in}}%
\pgfpathlineto{\pgfqpoint{2.015633in}{2.644559in}}%
\pgfpathlineto{\pgfqpoint{2.017846in}{2.644559in}}%
\pgfpathlineto{\pgfqpoint{2.017851in}{2.749082in}}%
\pgfpathlineto{\pgfqpoint{2.017915in}{2.640979in}}%
\pgfpathlineto{\pgfqpoint{2.018599in}{2.680437in}}%
\pgfpathlineto{\pgfqpoint{2.019440in}{2.762275in}}%
\pgfpathlineto{\pgfqpoint{2.019820in}{2.786099in}}%
\pgfpathlineto{\pgfqpoint{2.020396in}{2.644559in}}%
\pgfpathlineto{\pgfqpoint{2.022792in}{2.644559in}}%
\pgfpathlineto{\pgfqpoint{2.022796in}{2.751601in}}%
\pgfpathlineto{\pgfqpoint{2.022870in}{2.629622in}}%
\pgfpathlineto{\pgfqpoint{2.023556in}{2.670296in}}%
\pgfpathlineto{\pgfqpoint{2.024397in}{2.774883in}}%
\pgfpathlineto{\pgfqpoint{2.024765in}{2.797958in}}%
\pgfpathlineto{\pgfqpoint{2.025298in}{2.644559in}}%
\pgfpathlineto{\pgfqpoint{2.027737in}{2.644559in}}%
\pgfpathlineto{\pgfqpoint{2.027742in}{2.749508in}}%
\pgfpathlineto{\pgfqpoint{2.028497in}{2.699560in}}%
\pgfpathlineto{\pgfqpoint{2.028546in}{2.695893in}}%
\pgfpathlineto{\pgfqpoint{2.028892in}{2.724422in}}%
\pgfpathlineto{\pgfqpoint{2.028942in}{2.720739in}}%
\pgfpathlineto{\pgfqpoint{2.029673in}{2.773441in}}%
\pgfpathlineto{\pgfqpoint{2.029753in}{2.732517in}}%
\pgfpathlineto{\pgfqpoint{2.030601in}{2.644559in}}%
\pgfpathlineto{\pgfqpoint{2.032682in}{2.644559in}}%
\pgfpathlineto{\pgfqpoint{2.032688in}{2.748933in}}%
\pgfpathlineto{\pgfqpoint{2.032762in}{2.637437in}}%
\pgfpathlineto{\pgfqpoint{2.033447in}{2.678324in}}%
\pgfpathlineto{\pgfqpoint{2.034288in}{2.766358in}}%
\pgfpathlineto{\pgfqpoint{2.034656in}{2.789443in}}%
\pgfpathlineto{\pgfqpoint{2.035215in}{2.644559in}}%
\pgfpathlineto{\pgfqpoint{2.037628in}{2.644559in}}%
\pgfpathlineto{\pgfqpoint{2.037633in}{2.749055in}}%
\pgfpathlineto{\pgfqpoint{2.037715in}{2.631547in}}%
\pgfpathlineto{\pgfqpoint{2.038400in}{2.673069in}}%
\pgfpathlineto{\pgfqpoint{2.039241in}{2.771933in}}%
\pgfpathlineto{\pgfqpoint{2.039602in}{2.794517in}}%
\pgfpathlineto{\pgfqpoint{2.040114in}{2.644559in}}%
\pgfpathlineto{\pgfqpoint{2.042573in}{2.644559in}}%
\pgfpathlineto{\pgfqpoint{2.042579in}{2.749319in}}%
\pgfpathlineto{\pgfqpoint{2.042675in}{2.641581in}}%
\pgfpathlineto{\pgfqpoint{2.043315in}{2.707677in}}%
\pgfpathlineto{\pgfqpoint{2.043364in}{2.684043in}}%
\pgfpathlineto{\pgfqpoint{2.043908in}{2.744925in}}%
\pgfpathlineto{\pgfqpoint{2.043958in}{2.721358in}}%
\pgfpathlineto{\pgfqpoint{2.044502in}{2.782170in}}%
\pgfpathlineto{\pgfqpoint{2.044602in}{2.697709in}}%
\pgfpathlineto{\pgfqpoint{2.044616in}{2.746718in}}%
\pgfpathlineto{\pgfqpoint{2.045371in}{2.644559in}}%
\pgfpathlineto{\pgfqpoint{2.047519in}{2.644559in}}%
\pgfpathlineto{\pgfqpoint{2.047523in}{2.754174in}}%
\pgfpathlineto{\pgfqpoint{2.047593in}{2.633705in}}%
\pgfpathlineto{\pgfqpoint{2.048281in}{2.673895in}}%
\pgfpathlineto{\pgfqpoint{2.049122in}{2.771696in}}%
\pgfpathlineto{\pgfqpoint{2.049493in}{2.794937in}}%
\pgfpathlineto{\pgfqpoint{2.050041in}{2.644559in}}%
\pgfpathlineto{\pgfqpoint{2.052464in}{2.644559in}}%
\pgfpathlineto{\pgfqpoint{2.052470in}{2.749623in}}%
\pgfpathlineto{\pgfqpoint{2.053195in}{2.697654in}}%
\pgfpathlineto{\pgfqpoint{2.053245in}{2.693893in}}%
\pgfpathlineto{\pgfqpoint{2.053591in}{2.722515in}}%
\pgfpathlineto{\pgfqpoint{2.053640in}{2.718742in}}%
\pgfpathlineto{\pgfqpoint{2.054438in}{2.775745in}}%
\pgfpathlineto{\pgfqpoint{2.054480in}{2.732699in}}%
\pgfpathlineto{\pgfqpoint{2.055319in}{2.644559in}}%
\pgfpathlineto{\pgfqpoint{2.057410in}{2.644559in}}%
\pgfpathlineto{\pgfqpoint{2.057415in}{2.749096in}}%
\pgfpathlineto{\pgfqpoint{2.057492in}{2.635104in}}%
\pgfpathlineto{\pgfqpoint{2.058178in}{2.676288in}}%
\pgfpathlineto{\pgfqpoint{2.059018in}{2.769032in}}%
\pgfpathlineto{\pgfqpoint{2.059384in}{2.791900in}}%
\pgfpathlineto{\pgfqpoint{2.059945in}{2.644559in}}%
\pgfpathlineto{\pgfqpoint{2.062355in}{2.644559in}}%
\pgfpathlineto{\pgfqpoint{2.062361in}{2.749352in}}%
\pgfpathlineto{\pgfqpoint{2.062437in}{2.641962in}}%
\pgfpathlineto{\pgfqpoint{2.063126in}{2.683102in}}%
\pgfpathlineto{\pgfqpoint{2.063967in}{2.762046in}}%
\pgfpathlineto{\pgfqpoint{2.064329in}{2.784770in}}%
\pgfpathlineto{\pgfqpoint{2.064953in}{2.644559in}}%
\pgfpathlineto{\pgfqpoint{2.067301in}{2.644559in}}%
\pgfpathlineto{\pgfqpoint{2.067306in}{2.749560in}}%
\pgfpathlineto{\pgfqpoint{2.068037in}{2.697931in}}%
\pgfpathlineto{\pgfqpoint{2.068086in}{2.694001in}}%
\pgfpathlineto{\pgfqpoint{2.068432in}{2.722791in}}%
\pgfpathlineto{\pgfqpoint{2.068482in}{2.718850in}}%
\pgfpathlineto{\pgfqpoint{2.069274in}{2.775696in}}%
\pgfpathlineto{\pgfqpoint{2.069317in}{2.734216in}}%
\pgfpathlineto{\pgfqpoint{2.070107in}{2.644559in}}%
\pgfpathlineto{\pgfqpoint{2.072246in}{2.644559in}}%
\pgfpathlineto{\pgfqpoint{2.072251in}{2.749276in}}%
\pgfpathlineto{\pgfqpoint{2.072332in}{2.632531in}}%
\pgfpathlineto{\pgfqpoint{2.073017in}{2.673975in}}%
\pgfpathlineto{\pgfqpoint{2.073858in}{2.771865in}}%
\pgfpathlineto{\pgfqpoint{2.074220in}{2.794526in}}%
\pgfpathlineto{\pgfqpoint{2.074769in}{2.644559in}}%
\pgfpathlineto{\pgfqpoint{2.077192in}{2.644559in}}%
\pgfpathlineto{\pgfqpoint{2.077197in}{2.749721in}}%
\pgfpathlineto{\pgfqpoint{2.077927in}{2.691706in}}%
\pgfpathlineto{\pgfqpoint{2.079139in}{2.773508in}}%
\pgfpathlineto{\pgfqpoint{2.080743in}{2.644559in}}%
\pgfpathlineto{\pgfqpoint{2.082137in}{2.644559in}}%
\pgfpathlineto{\pgfqpoint{2.082142in}{2.748844in}}%
\pgfpathlineto{\pgfqpoint{2.082215in}{2.638257in}}%
\pgfpathlineto{\pgfqpoint{2.082901in}{2.679048in}}%
\pgfpathlineto{\pgfqpoint{2.083741in}{2.765338in}}%
\pgfpathlineto{\pgfqpoint{2.084111in}{2.788487in}}%
\pgfpathlineto{\pgfqpoint{2.084670in}{2.644559in}}%
\pgfpathlineto{\pgfqpoint{2.087082in}{2.644559in}}%
\pgfpathlineto{\pgfqpoint{2.087088in}{2.748769in}}%
\pgfpathlineto{\pgfqpoint{2.087161in}{2.638617in}}%
\pgfpathlineto{\pgfqpoint{2.087846in}{2.679360in}}%
\pgfpathlineto{\pgfqpoint{2.088686in}{2.764741in}}%
\pgfpathlineto{\pgfqpoint{2.089056in}{2.787919in}}%
\pgfpathlineto{\pgfqpoint{2.089616in}{2.644559in}}%
\pgfpathlineto{\pgfqpoint{2.092028in}{2.644559in}}%
\pgfpathlineto{\pgfqpoint{2.092033in}{2.749063in}}%
\pgfpathlineto{\pgfqpoint{2.092101in}{2.639999in}}%
\pgfpathlineto{\pgfqpoint{2.092788in}{2.680166in}}%
\pgfpathlineto{\pgfqpoint{2.093629in}{2.763336in}}%
\pgfpathlineto{\pgfqpoint{2.094002in}{2.786742in}}%
\pgfpathlineto{\pgfqpoint{2.094579in}{2.644559in}}%
\pgfpathlineto{\pgfqpoint{2.096973in}{2.644559in}}%
\pgfpathlineto{\pgfqpoint{2.096978in}{2.751592in}}%
\pgfpathlineto{\pgfqpoint{2.097059in}{2.628593in}}%
\pgfpathlineto{\pgfqpoint{2.097745in}{2.669906in}}%
\pgfpathlineto{\pgfqpoint{2.098586in}{2.776123in}}%
\pgfpathlineto{\pgfqpoint{2.098947in}{2.798768in}}%
\pgfpathlineto{\pgfqpoint{2.099458in}{2.644559in}}%
\pgfpathlineto{\pgfqpoint{2.101919in}{2.644559in}}%
\pgfpathlineto{\pgfqpoint{2.101924in}{2.748220in}}%
\pgfpathlineto{\pgfqpoint{2.101985in}{2.641721in}}%
\pgfpathlineto{\pgfqpoint{2.102667in}{2.680866in}}%
\pgfpathlineto{\pgfqpoint{2.103508in}{2.760912in}}%
\pgfpathlineto{\pgfqpoint{2.103893in}{2.785032in}}%
\pgfpathlineto{\pgfqpoint{2.104503in}{2.644559in}}%
\pgfpathlineto{\pgfqpoint{2.106864in}{2.644559in}}%
\pgfpathlineto{\pgfqpoint{2.106870in}{2.749288in}}%
\pgfpathlineto{\pgfqpoint{2.107598in}{2.690678in}}%
\pgfpathlineto{\pgfqpoint{2.108736in}{2.768791in}}%
\pgfpathlineto{\pgfqpoint{2.108812in}{2.773561in}}%
\pgfpathlineto{\pgfqpoint{2.108901in}{2.706234in}}%
\pgfpathlineto{\pgfqpoint{2.109668in}{2.644559in}}%
\pgfpathlineto{\pgfqpoint{2.111810in}{2.644559in}}%
\pgfpathlineto{\pgfqpoint{2.111815in}{2.749505in}}%
\pgfpathlineto{\pgfqpoint{2.111949in}{2.643334in}}%
\pgfpathlineto{\pgfqpoint{2.112542in}{2.680501in}}%
\pgfpathlineto{\pgfqpoint{2.113383in}{2.760523in}}%
\pgfpathlineto{\pgfqpoint{2.113756in}{2.783951in}}%
\pgfpathlineto{\pgfqpoint{2.114374in}{2.644559in}}%
\pgfpathlineto{\pgfqpoint{2.116755in}{2.644559in}}%
\pgfpathlineto{\pgfqpoint{2.116760in}{2.749220in}}%
\pgfpathlineto{\pgfqpoint{2.116859in}{2.631402in}}%
\pgfpathlineto{\pgfqpoint{2.117498in}{2.718230in}}%
\pgfpathlineto{\pgfqpoint{2.117548in}{2.673907in}}%
\pgfpathlineto{\pgfqpoint{2.118191in}{2.761657in}}%
\pgfpathlineto{\pgfqpoint{2.118240in}{2.717470in}}%
\pgfpathlineto{\pgfqpoint{2.118682in}{2.792491in}}%
\pgfpathlineto{\pgfqpoint{2.118813in}{2.644559in}}%
\pgfpathlineto{\pgfqpoint{2.118925in}{2.644559in}}%
\pgfpathlineto{\pgfqpoint{2.121701in}{2.644559in}}%
\pgfpathlineto{\pgfqpoint{2.121706in}{2.749540in}}%
\pgfpathlineto{\pgfqpoint{2.121814in}{2.642612in}}%
\pgfpathlineto{\pgfqpoint{2.122455in}{2.708531in}}%
\pgfpathlineto{\pgfqpoint{2.122505in}{2.685567in}}%
\pgfpathlineto{\pgfqpoint{2.123049in}{2.745780in}}%
\pgfpathlineto{\pgfqpoint{2.123098in}{2.722882in}}%
\pgfpathlineto{\pgfqpoint{2.123634in}{2.782488in}}%
\pgfpathlineto{\pgfqpoint{2.123729in}{2.698022in}}%
\pgfpathlineto{\pgfqpoint{2.123744in}{2.747147in}}%
\pgfpathlineto{\pgfqpoint{2.124498in}{2.644559in}}%
\pgfpathlineto{\pgfqpoint{2.126646in}{2.644559in}}%
\pgfpathlineto{\pgfqpoint{2.126652in}{2.749524in}}%
\pgfpathlineto{\pgfqpoint{2.127406in}{2.699532in}}%
\pgfpathlineto{\pgfqpoint{2.127455in}{2.695766in}}%
\pgfpathlineto{\pgfqpoint{2.127801in}{2.724395in}}%
\pgfpathlineto{\pgfqpoint{2.127851in}{2.720613in}}%
\pgfpathlineto{\pgfqpoint{2.128581in}{2.773421in}}%
\pgfpathlineto{\pgfqpoint{2.128651in}{2.741317in}}%
\pgfpathlineto{\pgfqpoint{2.129448in}{2.644559in}}%
\pgfpathlineto{\pgfqpoint{2.131592in}{2.644559in}}%
\pgfpathlineto{\pgfqpoint{2.131597in}{2.749423in}}%
\pgfpathlineto{\pgfqpoint{2.132362in}{2.692963in}}%
\pgfpathlineto{\pgfqpoint{2.133499in}{2.770989in}}%
\pgfpathlineto{\pgfqpoint{2.133623in}{2.712090in}}%
\pgfpathlineto{\pgfqpoint{2.133565in}{2.775155in}}%
\pgfpathlineto{\pgfqpoint{2.133637in}{2.734972in}}%
\pgfpathlineto{\pgfqpoint{2.134400in}{2.644559in}}%
\pgfpathlineto{\pgfqpoint{2.136537in}{2.644559in}}%
\pgfpathlineto{\pgfqpoint{2.136542in}{2.749370in}}%
\pgfpathlineto{\pgfqpoint{2.137262in}{2.697049in}}%
\pgfpathlineto{\pgfqpoint{2.137311in}{2.692677in}}%
\pgfpathlineto{\pgfqpoint{2.137657in}{2.721905in}}%
\pgfpathlineto{\pgfqpoint{2.137707in}{2.717530in}}%
\pgfpathlineto{\pgfqpoint{2.138511in}{2.775521in}}%
\pgfpathlineto{\pgfqpoint{2.138537in}{2.757840in}}%
\pgfpathlineto{\pgfqpoint{2.139290in}{2.644559in}}%
\pgfpathlineto{\pgfqpoint{2.141482in}{2.644559in}}%
\pgfpathlineto{\pgfqpoint{2.141488in}{2.749668in}}%
\pgfpathlineto{\pgfqpoint{2.142225in}{2.684040in}}%
\pgfpathlineto{\pgfqpoint{2.143065in}{2.758493in}}%
\pgfpathlineto{\pgfqpoint{2.143456in}{2.783026in}}%
\pgfpathlineto{\pgfqpoint{2.144137in}{2.644559in}}%
\pgfpathlineto{\pgfqpoint{2.146428in}{2.644559in}}%
\pgfpathlineto{\pgfqpoint{2.146433in}{2.748948in}}%
\pgfpathlineto{\pgfqpoint{2.146510in}{2.635479in}}%
\pgfpathlineto{\pgfqpoint{2.147195in}{2.676618in}}%
\pgfpathlineto{\pgfqpoint{2.148036in}{2.768125in}}%
\pgfpathlineto{\pgfqpoint{2.148402in}{2.791023in}}%
\pgfpathlineto{\pgfqpoint{2.148964in}{2.644559in}}%
\pgfpathlineto{\pgfqpoint{2.151373in}{2.644559in}}%
\pgfpathlineto{\pgfqpoint{2.151379in}{2.749448in}}%
\pgfpathlineto{\pgfqpoint{2.151433in}{2.642935in}}%
\pgfpathlineto{\pgfqpoint{2.152108in}{2.680339in}}%
\pgfpathlineto{\pgfqpoint{2.152949in}{2.760788in}}%
\pgfpathlineto{\pgfqpoint{2.153321in}{2.784140in}}%
\pgfpathlineto{\pgfqpoint{2.153938in}{2.644559in}}%
\pgfpathlineto{\pgfqpoint{2.156319in}{2.644559in}}%
\pgfpathlineto{\pgfqpoint{2.156323in}{2.749205in}}%
\pgfpathlineto{\pgfqpoint{2.156411in}{2.631555in}}%
\pgfpathlineto{\pgfqpoint{2.157051in}{2.717283in}}%
\pgfpathlineto{\pgfqpoint{2.157100in}{2.673454in}}%
\pgfpathlineto{\pgfqpoint{2.157743in}{2.760711in}}%
\pgfpathlineto{\pgfqpoint{2.157793in}{2.717015in}}%
\pgfpathlineto{\pgfqpoint{2.158293in}{2.795152in}}%
\pgfpathlineto{\pgfqpoint{2.158379in}{2.644559in}}%
\pgfpathlineto{\pgfqpoint{2.158458in}{2.644559in}}%
\pgfpathlineto{\pgfqpoint{2.161264in}{2.644559in}}%
\pgfpathlineto{\pgfqpoint{2.161270in}{2.748985in}}%
\pgfpathlineto{\pgfqpoint{2.161349in}{2.633054in}}%
\pgfpathlineto{\pgfqpoint{2.162035in}{2.674418in}}%
\pgfpathlineto{\pgfqpoint{2.162875in}{2.770358in}}%
\pgfpathlineto{\pgfqpoint{2.163238in}{2.793075in}}%
\pgfpathlineto{\pgfqpoint{2.163771in}{2.644559in}}%
\pgfpathlineto{\pgfqpoint{2.166210in}{2.644559in}}%
\pgfpathlineto{\pgfqpoint{2.166215in}{2.749760in}}%
\pgfpathlineto{\pgfqpoint{2.166949in}{2.698508in}}%
\pgfpathlineto{\pgfqpoint{2.166998in}{2.694728in}}%
\pgfpathlineto{\pgfqpoint{2.167344in}{2.723370in}}%
\pgfpathlineto{\pgfqpoint{2.167394in}{2.719575in}}%
\pgfpathlineto{\pgfqpoint{2.168136in}{2.773091in}}%
\pgfpathlineto{\pgfqpoint{2.168214in}{2.742226in}}%
\pgfpathlineto{\pgfqpoint{2.169017in}{2.644559in}}%
\pgfpathlineto{\pgfqpoint{2.171155in}{2.644559in}}%
\pgfpathlineto{\pgfqpoint{2.171161in}{2.749161in}}%
\pgfpathlineto{\pgfqpoint{2.171282in}{2.641191in}}%
\pgfpathlineto{\pgfqpoint{2.171925in}{2.710299in}}%
\pgfpathlineto{\pgfqpoint{2.171974in}{2.684459in}}%
\pgfpathlineto{\pgfqpoint{2.172518in}{2.747543in}}%
\pgfpathlineto{\pgfqpoint{2.172568in}{2.721778in}}%
\pgfpathlineto{\pgfqpoint{2.173096in}{2.783775in}}%
\pgfpathlineto{\pgfqpoint{2.173212in}{2.644559in}}%
\pgfpathlineto{\pgfqpoint{2.176101in}{2.644559in}}%
\pgfpathlineto{\pgfqpoint{2.176105in}{2.751262in}}%
\pgfpathlineto{\pgfqpoint{2.176838in}{2.692199in}}%
\pgfpathlineto{\pgfqpoint{2.178074in}{2.774237in}}%
\pgfpathlineto{\pgfqpoint{2.178074in}{2.769255in}}%
\pgfpathlineto{\pgfqpoint{2.178832in}{2.644559in}}%
\pgfpathlineto{\pgfqpoint{2.181046in}{2.644559in}}%
\pgfpathlineto{\pgfqpoint{2.181051in}{2.748436in}}%
\pgfpathlineto{\pgfqpoint{2.181118in}{2.644351in}}%
\pgfpathlineto{\pgfqpoint{2.181803in}{2.684355in}}%
\pgfpathlineto{\pgfqpoint{2.182643in}{2.758863in}}%
\pgfpathlineto{\pgfqpoint{2.183020in}{2.782476in}}%
\pgfpathlineto{\pgfqpoint{2.183680in}{2.644559in}}%
\pgfpathlineto{\pgfqpoint{2.185992in}{2.644559in}}%
\pgfpathlineto{\pgfqpoint{2.185997in}{2.748749in}}%
\pgfpathlineto{\pgfqpoint{2.186069in}{2.639144in}}%
\pgfpathlineto{\pgfqpoint{2.186754in}{2.679826in}}%
\pgfpathlineto{\pgfqpoint{2.187595in}{2.764262in}}%
\pgfpathlineto{\pgfqpoint{2.187965in}{2.787479in}}%
\pgfpathlineto{\pgfqpoint{2.188575in}{2.644559in}}%
\pgfpathlineto{\pgfqpoint{2.190937in}{2.644559in}}%
\pgfpathlineto{\pgfqpoint{2.190942in}{2.748610in}}%
\pgfpathlineto{\pgfqpoint{2.191013in}{2.640462in}}%
\pgfpathlineto{\pgfqpoint{2.191698in}{2.680972in}}%
\pgfpathlineto{\pgfqpoint{2.192539in}{2.762680in}}%
\pgfpathlineto{\pgfqpoint{2.192911in}{2.786003in}}%
\pgfpathlineto{\pgfqpoint{2.193493in}{2.644559in}}%
\pgfpathlineto{\pgfqpoint{2.195882in}{2.644559in}}%
\pgfpathlineto{\pgfqpoint{2.195888in}{2.749578in}}%
\pgfpathlineto{\pgfqpoint{2.196651in}{2.693605in}}%
\pgfpathlineto{\pgfqpoint{2.197788in}{2.771559in}}%
\pgfpathlineto{\pgfqpoint{2.197910in}{2.720594in}}%
\pgfpathlineto{\pgfqpoint{2.197856in}{2.775828in}}%
\pgfpathlineto{\pgfqpoint{2.197924in}{2.726410in}}%
\pgfpathlineto{\pgfqpoint{2.198701in}{2.644559in}}%
\pgfpathlineto{\pgfqpoint{2.200828in}{2.644559in}}%
\pgfpathlineto{\pgfqpoint{2.200833in}{2.748453in}}%
\pgfpathlineto{\pgfqpoint{2.200900in}{2.643502in}}%
\pgfpathlineto{\pgfqpoint{2.201586in}{2.683639in}}%
\pgfpathlineto{\pgfqpoint{2.202426in}{2.759543in}}%
\pgfpathlineto{\pgfqpoint{2.202802in}{2.783076in}}%
\pgfpathlineto{\pgfqpoint{2.203461in}{2.644559in}}%
\pgfpathlineto{\pgfqpoint{2.205773in}{2.644559in}}%
\pgfpathlineto{\pgfqpoint{2.205779in}{2.748702in}}%
\pgfpathlineto{\pgfqpoint{2.205850in}{2.640252in}}%
\pgfpathlineto{\pgfqpoint{2.206535in}{2.680802in}}%
\pgfpathlineto{\pgfqpoint{2.207375in}{2.763248in}}%
\pgfpathlineto{\pgfqpoint{2.207747in}{2.786550in}}%
\pgfpathlineto{\pgfqpoint{2.208357in}{2.644559in}}%
\pgfpathlineto{\pgfqpoint{2.210719in}{2.644559in}}%
\pgfpathlineto{\pgfqpoint{2.210724in}{2.748665in}}%
\pgfpathlineto{\pgfqpoint{2.210796in}{2.639318in}}%
\pgfpathlineto{\pgfqpoint{2.211481in}{2.679969in}}%
\pgfpathlineto{\pgfqpoint{2.212322in}{2.763769in}}%
\pgfpathlineto{\pgfqpoint{2.212693in}{2.787004in}}%
\pgfpathlineto{\pgfqpoint{2.213302in}{2.644559in}}%
\pgfpathlineto{\pgfqpoint{2.215664in}{2.644559in}}%
\pgfpathlineto{\pgfqpoint{2.215670in}{2.748719in}}%
\pgfpathlineto{\pgfqpoint{2.215741in}{2.639554in}}%
\pgfpathlineto{\pgfqpoint{2.216426in}{2.680186in}}%
\pgfpathlineto{\pgfqpoint{2.217267in}{2.763824in}}%
\pgfpathlineto{\pgfqpoint{2.217638in}{2.787073in}}%
\pgfpathlineto{\pgfqpoint{2.218248in}{2.644559in}}%
\pgfpathlineto{\pgfqpoint{2.220610in}{2.644559in}}%
\pgfpathlineto{\pgfqpoint{2.220615in}{2.748373in}}%
\pgfpathlineto{\pgfqpoint{2.221363in}{2.685731in}}%
\pgfpathlineto{\pgfqpoint{2.222303in}{2.763302in}}%
\pgfpathlineto{\pgfqpoint{2.222584in}{2.780916in}}%
\pgfpathlineto{\pgfqpoint{2.223441in}{2.644559in}}%
\pgfpathlineto{\pgfqpoint{2.225555in}{2.644559in}}%
\pgfpathlineto{\pgfqpoint{2.225561in}{2.748948in}}%
\pgfpathlineto{\pgfqpoint{2.225640in}{2.633348in}}%
\pgfpathlineto{\pgfqpoint{2.226325in}{2.674683in}}%
\pgfpathlineto{\pgfqpoint{2.227166in}{2.769959in}}%
\pgfpathlineto{\pgfqpoint{2.227529in}{2.792699in}}%
\pgfpathlineto{\pgfqpoint{2.228092in}{2.644559in}}%
\pgfpathlineto{\pgfqpoint{2.230501in}{2.644559in}}%
\pgfpathlineto{\pgfqpoint{2.230506in}{2.749797in}}%
\pgfpathlineto{\pgfqpoint{2.231243in}{2.698969in}}%
\pgfpathlineto{\pgfqpoint{2.231293in}{2.694792in}}%
\pgfpathlineto{\pgfqpoint{2.231639in}{2.723831in}}%
\pgfpathlineto{\pgfqpoint{2.231688in}{2.719640in}}%
\pgfpathlineto{\pgfqpoint{2.232428in}{2.773385in}}%
\pgfpathlineto{\pgfqpoint{2.232505in}{2.742886in}}%
\pgfpathlineto{\pgfqpoint{2.233303in}{2.644559in}}%
\pgfpathlineto{\pgfqpoint{2.235446in}{2.644559in}}%
\pgfpathlineto{\pgfqpoint{2.235451in}{2.748728in}}%
\pgfpathlineto{\pgfqpoint{2.235529in}{2.634727in}}%
\pgfpathlineto{\pgfqpoint{2.236214in}{2.675920in}}%
\pgfpathlineto{\pgfqpoint{2.237055in}{2.767958in}}%
\pgfpathlineto{\pgfqpoint{2.237420in}{2.790805in}}%
\pgfpathlineto{\pgfqpoint{2.237981in}{2.644559in}}%
\pgfpathlineto{\pgfqpoint{2.240392in}{2.644559in}}%
\pgfpathlineto{\pgfqpoint{2.240397in}{2.749410in}}%
\pgfpathlineto{\pgfqpoint{2.241162in}{2.699824in}}%
\pgfpathlineto{\pgfqpoint{2.241211in}{2.695901in}}%
\pgfpathlineto{\pgfqpoint{2.241557in}{2.724683in}}%
\pgfpathlineto{\pgfqpoint{2.241607in}{2.720751in}}%
\pgfpathlineto{\pgfqpoint{2.242332in}{2.773374in}}%
\pgfpathlineto{\pgfqpoint{2.242407in}{2.735397in}}%
\pgfpathlineto{\pgfqpoint{2.243238in}{2.644559in}}%
\pgfpathlineto{\pgfqpoint{2.245337in}{2.644559in}}%
\pgfpathlineto{\pgfqpoint{2.245342in}{2.749501in}}%
\pgfpathlineto{\pgfqpoint{2.245441in}{2.642186in}}%
\pgfpathlineto{\pgfqpoint{2.246081in}{2.707862in}}%
\pgfpathlineto{\pgfqpoint{2.246131in}{2.684761in}}%
\pgfpathlineto{\pgfqpoint{2.246675in}{2.745111in}}%
\pgfpathlineto{\pgfqpoint{2.246724in}{2.722075in}}%
\pgfpathlineto{\pgfqpoint{2.247265in}{2.782148in}}%
\pgfpathlineto{\pgfqpoint{2.247365in}{2.697100in}}%
\pgfpathlineto{\pgfqpoint{2.247380in}{2.747780in}}%
\pgfpathlineto{\pgfqpoint{2.248122in}{2.644559in}}%
\pgfpathlineto{\pgfqpoint{2.250282in}{2.644559in}}%
\pgfpathlineto{\pgfqpoint{2.250288in}{2.749782in}}%
\pgfpathlineto{\pgfqpoint{2.251030in}{2.698933in}}%
\pgfpathlineto{\pgfqpoint{2.251079in}{2.695121in}}%
\pgfpathlineto{\pgfqpoint{2.251425in}{2.723793in}}%
\pgfpathlineto{\pgfqpoint{2.251475in}{2.719970in}}%
\pgfpathlineto{\pgfqpoint{2.252212in}{2.773202in}}%
\pgfpathlineto{\pgfqpoint{2.252296in}{2.730497in}}%
\pgfpathlineto{\pgfqpoint{2.252308in}{2.731897in}}%
\pgfpathlineto{\pgfqpoint{2.252349in}{2.644559in}}%
\pgfpathlineto{\pgfqpoint{2.252381in}{2.644559in}}%
\pgfpathlineto{\pgfqpoint{2.255228in}{2.644559in}}%
\pgfpathlineto{\pgfqpoint{2.255233in}{2.749188in}}%
\pgfpathlineto{\pgfqpoint{2.255310in}{2.641361in}}%
\pgfpathlineto{\pgfqpoint{2.256000in}{2.682576in}}%
\pgfpathlineto{\pgfqpoint{2.256840in}{2.762195in}}%
\pgfpathlineto{\pgfqpoint{2.257202in}{2.784860in}}%
\pgfpathlineto{\pgfqpoint{2.257822in}{2.644559in}}%
\pgfpathlineto{\pgfqpoint{2.260173in}{2.644559in}}%
\pgfpathlineto{\pgfqpoint{2.260178in}{2.750674in}}%
\pgfpathlineto{\pgfqpoint{2.260266in}{2.627548in}}%
\pgfpathlineto{\pgfqpoint{2.260903in}{2.721507in}}%
\pgfpathlineto{\pgfqpoint{2.260953in}{2.669416in}}%
\pgfpathlineto{\pgfqpoint{2.261595in}{2.764906in}}%
\pgfpathlineto{\pgfqpoint{2.261645in}{2.713006in}}%
\pgfpathlineto{\pgfqpoint{2.262147in}{2.799471in}}%
\pgfpathlineto{\pgfqpoint{2.262238in}{2.644559in}}%
\pgfpathlineto{\pgfqpoint{2.262311in}{2.644559in}}%
\pgfpathlineto{\pgfqpoint{2.265119in}{2.644559in}}%
\pgfpathlineto{\pgfqpoint{2.265124in}{2.748070in}}%
\pgfpathlineto{\pgfqpoint{2.265179in}{2.643566in}}%
\pgfpathlineto{\pgfqpoint{2.265857in}{2.681440in}}%
\pgfpathlineto{\pgfqpoint{2.266698in}{2.759220in}}%
\pgfpathlineto{\pgfqpoint{2.267093in}{2.784001in}}%
\pgfpathlineto{\pgfqpoint{2.267702in}{2.644559in}}%
\pgfpathlineto{\pgfqpoint{2.270064in}{2.644559in}}%
\pgfpathlineto{\pgfqpoint{2.270070in}{2.748784in}}%
\pgfpathlineto{\pgfqpoint{2.270143in}{2.637595in}}%
\pgfpathlineto{\pgfqpoint{2.270829in}{2.678448in}}%
\pgfpathlineto{\pgfqpoint{2.271669in}{2.765597in}}%
\pgfpathlineto{\pgfqpoint{2.272038in}{2.788699in}}%
\pgfpathlineto{\pgfqpoint{2.272597in}{2.644559in}}%
\pgfpathlineto{\pgfqpoint{2.275010in}{2.644559in}}%
\pgfpathlineto{\pgfqpoint{2.275015in}{2.748828in}}%
\pgfpathlineto{\pgfqpoint{2.275093in}{2.634545in}}%
\pgfpathlineto{\pgfqpoint{2.275778in}{2.675761in}}%
\pgfpathlineto{\pgfqpoint{2.276619in}{2.768467in}}%
\pgfpathlineto{\pgfqpoint{2.276984in}{2.791299in}}%
\pgfpathlineto{\pgfqpoint{2.277544in}{2.644559in}}%
\pgfpathlineto{\pgfqpoint{2.279955in}{2.644559in}}%
\pgfpathlineto{\pgfqpoint{2.279961in}{2.749258in}}%
\pgfpathlineto{\pgfqpoint{2.280712in}{2.698512in}}%
\pgfpathlineto{\pgfqpoint{2.280761in}{2.695064in}}%
\pgfpathlineto{\pgfqpoint{2.281107in}{2.723371in}}%
\pgfpathlineto{\pgfqpoint{2.281157in}{2.719913in}}%
\pgfpathlineto{\pgfqpoint{2.281889in}{2.772483in}}%
\pgfpathlineto{\pgfqpoint{2.281971in}{2.735933in}}%
\pgfpathlineto{\pgfqpoint{2.282756in}{2.644559in}}%
\pgfpathlineto{\pgfqpoint{2.284901in}{2.644559in}}%
\pgfpathlineto{\pgfqpoint{2.284906in}{2.749540in}}%
\pgfpathlineto{\pgfqpoint{2.285026in}{2.643208in}}%
\pgfpathlineto{\pgfqpoint{2.285669in}{2.709476in}}%
\pgfpathlineto{\pgfqpoint{2.285719in}{2.686464in}}%
\pgfpathlineto{\pgfqpoint{2.286263in}{2.746723in}}%
\pgfpathlineto{\pgfqpoint{2.286312in}{2.723779in}}%
\pgfpathlineto{\pgfqpoint{2.286840in}{2.782992in}}%
\pgfpathlineto{\pgfqpoint{2.286930in}{2.697165in}}%
\pgfpathlineto{\pgfqpoint{2.286944in}{2.748638in}}%
\pgfpathlineto{\pgfqpoint{2.287713in}{2.644559in}}%
\pgfpathlineto{\pgfqpoint{2.289846in}{2.644559in}}%
\pgfpathlineto{\pgfqpoint{2.289851in}{2.748016in}}%
\pgfpathlineto{\pgfqpoint{2.289921in}{2.628543in}}%
\pgfpathlineto{\pgfqpoint{2.290608in}{2.668882in}}%
\pgfpathlineto{\pgfqpoint{2.291449in}{2.774917in}}%
\pgfpathlineto{\pgfqpoint{2.291820in}{2.798166in}}%
\pgfpathlineto{\pgfqpoint{2.292303in}{2.644559in}}%
\pgfpathlineto{\pgfqpoint{2.294792in}{2.644559in}}%
\pgfpathlineto{\pgfqpoint{2.294797in}{2.749451in}}%
\pgfpathlineto{\pgfqpoint{2.295538in}{2.692024in}}%
\pgfpathlineto{\pgfqpoint{2.296676in}{2.770109in}}%
\pgfpathlineto{\pgfqpoint{2.296820in}{2.720333in}}%
\pgfpathlineto{\pgfqpoint{2.296765in}{2.775744in}}%
\pgfpathlineto{\pgfqpoint{2.296833in}{2.726387in}}%
\pgfpathlineto{\pgfqpoint{2.297610in}{2.644559in}}%
\pgfpathlineto{\pgfqpoint{2.299737in}{2.644559in}}%
\pgfpathlineto{\pgfqpoint{2.299742in}{2.748476in}}%
\pgfpathlineto{\pgfqpoint{2.299810in}{2.642909in}}%
\pgfpathlineto{\pgfqpoint{2.300495in}{2.683123in}}%
\pgfpathlineto{\pgfqpoint{2.301336in}{2.760118in}}%
\pgfpathlineto{\pgfqpoint{2.301711in}{2.783609in}}%
\pgfpathlineto{\pgfqpoint{2.302369in}{2.644559in}}%
\pgfpathlineto{\pgfqpoint{2.304682in}{2.644559in}}%
\pgfpathlineto{\pgfqpoint{2.304688in}{2.748936in}}%
\pgfpathlineto{\pgfqpoint{2.304765in}{2.634891in}}%
\pgfpathlineto{\pgfqpoint{2.305451in}{2.676083in}}%
\pgfpathlineto{\pgfqpoint{2.306291in}{2.768567in}}%
\pgfpathlineto{\pgfqpoint{2.306656in}{2.791423in}}%
\pgfpathlineto{\pgfqpoint{2.307217in}{2.644559in}}%
\pgfpathlineto{\pgfqpoint{2.309628in}{2.644559in}}%
\pgfpathlineto{\pgfqpoint{2.309633in}{2.749297in}}%
\pgfpathlineto{\pgfqpoint{2.309747in}{2.641395in}}%
\pgfpathlineto{\pgfqpoint{2.310390in}{2.709678in}}%
\pgfpathlineto{\pgfqpoint{2.310440in}{2.684572in}}%
\pgfpathlineto{\pgfqpoint{2.310984in}{2.746924in}}%
\pgfpathlineto{\pgfqpoint{2.311033in}{2.721889in}}%
\pgfpathlineto{\pgfqpoint{2.311565in}{2.783384in}}%
\pgfpathlineto{\pgfqpoint{2.311657in}{2.696723in}}%
\pgfpathlineto{\pgfqpoint{2.311672in}{2.748184in}}%
\pgfpathlineto{\pgfqpoint{2.312440in}{2.644559in}}%
\pgfpathlineto{\pgfqpoint{2.314573in}{2.644559in}}%
\pgfpathlineto{\pgfqpoint{2.314578in}{2.749332in}}%
\pgfpathlineto{\pgfqpoint{2.314593in}{2.643189in}}%
\pgfpathlineto{\pgfqpoint{2.315341in}{2.705890in}}%
\pgfpathlineto{\pgfqpoint{2.315390in}{2.689329in}}%
\pgfpathlineto{\pgfqpoint{2.315934in}{2.743165in}}%
\pgfpathlineto{\pgfqpoint{2.315984in}{2.726616in}}%
\pgfpathlineto{\pgfqpoint{2.316513in}{2.779490in}}%
\pgfpathlineto{\pgfqpoint{2.316602in}{2.697059in}}%
\pgfpathlineto{\pgfqpoint{2.316617in}{2.748140in}}%
\pgfpathlineto{\pgfqpoint{2.317386in}{2.644559in}}%
\pgfpathlineto{\pgfqpoint{2.319519in}{2.644559in}}%
\pgfpathlineto{\pgfqpoint{2.319523in}{2.748688in}}%
\pgfpathlineto{\pgfqpoint{2.319624in}{2.631353in}}%
\pgfpathlineto{\pgfqpoint{2.320264in}{2.718291in}}%
\pgfpathlineto{\pgfqpoint{2.320313in}{2.673945in}}%
\pgfpathlineto{\pgfqpoint{2.320956in}{2.761717in}}%
\pgfpathlineto{\pgfqpoint{2.321006in}{2.717509in}}%
\pgfpathlineto{\pgfqpoint{2.321447in}{2.792489in}}%
\pgfpathlineto{\pgfqpoint{2.321576in}{2.644559in}}%
\pgfpathlineto{\pgfqpoint{2.321690in}{2.644559in}}%
\pgfpathlineto{\pgfqpoint{2.324464in}{2.644559in}}%
\pgfpathlineto{\pgfqpoint{2.324470in}{2.749201in}}%
\pgfpathlineto{\pgfqpoint{2.324550in}{2.632777in}}%
\pgfpathlineto{\pgfqpoint{2.325235in}{2.674193in}}%
\pgfpathlineto{\pgfqpoint{2.326076in}{2.771370in}}%
\pgfpathlineto{\pgfqpoint{2.326438in}{2.794052in}}%
\pgfpathlineto{\pgfqpoint{2.326978in}{2.644559in}}%
\pgfpathlineto{\pgfqpoint{2.329410in}{2.644559in}}%
\pgfpathlineto{\pgfqpoint{2.329415in}{2.749432in}}%
\pgfpathlineto{\pgfqpoint{2.329529in}{2.641840in}}%
\pgfpathlineto{\pgfqpoint{2.330172in}{2.709701in}}%
\pgfpathlineto{\pgfqpoint{2.330222in}{2.685020in}}%
\pgfpathlineto{\pgfqpoint{2.330766in}{2.746947in}}%
\pgfpathlineto{\pgfqpoint{2.330815in}{2.722337in}}%
\pgfpathlineto{\pgfqpoint{2.331347in}{2.783400in}}%
\pgfpathlineto{\pgfqpoint{2.331439in}{2.696930in}}%
\pgfpathlineto{\pgfqpoint{2.331454in}{2.748406in}}%
\pgfpathlineto{\pgfqpoint{2.332222in}{2.644559in}}%
\pgfpathlineto{\pgfqpoint{2.334355in}{2.644559in}}%
\pgfpathlineto{\pgfqpoint{2.334360in}{2.748497in}}%
\pgfpathlineto{\pgfqpoint{2.334481in}{2.631559in}}%
\pgfpathlineto{\pgfqpoint{2.335124in}{2.720370in}}%
\pgfpathlineto{\pgfqpoint{2.335174in}{2.674835in}}%
\pgfpathlineto{\pgfqpoint{2.335817in}{2.763791in}}%
\pgfpathlineto{\pgfqpoint{2.335866in}{2.718404in}}%
\pgfpathlineto{\pgfqpoint{2.336295in}{2.793805in}}%
\pgfpathlineto{\pgfqpoint{2.336421in}{2.644559in}}%
\pgfpathlineto{\pgfqpoint{2.336566in}{2.644559in}}%
\pgfpathlineto{\pgfqpoint{2.339301in}{2.644559in}}%
\pgfpathlineto{\pgfqpoint{2.339306in}{2.749304in}}%
\pgfpathlineto{\pgfqpoint{2.340069in}{2.692808in}}%
\pgfpathlineto{\pgfqpoint{2.341206in}{2.771004in}}%
\pgfpathlineto{\pgfqpoint{2.341329in}{2.719325in}}%
\pgfpathlineto{\pgfqpoint{2.341274in}{2.775296in}}%
\pgfpathlineto{\pgfqpoint{2.341342in}{2.726225in}}%
\pgfpathlineto{\pgfqpoint{2.342108in}{2.644559in}}%
\pgfpathlineto{\pgfqpoint{2.344246in}{2.644559in}}%
\pgfpathlineto{\pgfqpoint{2.344251in}{2.748870in}}%
\pgfpathlineto{\pgfqpoint{2.344324in}{2.638369in}}%
\pgfpathlineto{\pgfqpoint{2.345010in}{2.679152in}}%
\pgfpathlineto{\pgfqpoint{2.345850in}{2.765388in}}%
\pgfpathlineto{\pgfqpoint{2.346220in}{2.788543in}}%
\pgfpathlineto{\pgfqpoint{2.346779in}{2.644559in}}%
\pgfpathlineto{\pgfqpoint{2.349192in}{2.644559in}}%
\pgfpathlineto{\pgfqpoint{2.349197in}{2.747674in}}%
\pgfpathlineto{\pgfqpoint{2.349959in}{2.704971in}}%
\pgfpathlineto{\pgfqpoint{2.350008in}{2.689501in}}%
\pgfpathlineto{\pgfqpoint{2.350552in}{2.742210in}}%
\pgfpathlineto{\pgfqpoint{2.350602in}{2.726825in}}%
\pgfpathlineto{\pgfqpoint{2.351131in}{2.778508in}}%
\pgfpathlineto{\pgfqpoint{2.351219in}{2.720051in}}%
\pgfpathlineto{\pgfqpoint{2.351233in}{2.724529in}}%
\pgfpathlineto{\pgfqpoint{2.351997in}{2.644559in}}%
\pgfpathlineto{\pgfqpoint{2.354137in}{2.644559in}}%
\pgfpathlineto{\pgfqpoint{2.354142in}{2.749519in}}%
\pgfpathlineto{\pgfqpoint{2.354265in}{2.643281in}}%
\pgfpathlineto{\pgfqpoint{2.354908in}{2.709642in}}%
\pgfpathlineto{\pgfqpoint{2.354957in}{2.686565in}}%
\pgfpathlineto{\pgfqpoint{2.355501in}{2.746889in}}%
\pgfpathlineto{\pgfqpoint{2.355551in}{2.723881in}}%
\pgfpathlineto{\pgfqpoint{2.356078in}{2.783078in}}%
\pgfpathlineto{\pgfqpoint{2.356166in}{2.697134in}}%
\pgfpathlineto{\pgfqpoint{2.356181in}{2.748624in}}%
\pgfpathlineto{\pgfqpoint{2.356949in}{2.644559in}}%
\pgfpathlineto{\pgfqpoint{2.359082in}{2.644559in}}%
\pgfpathlineto{\pgfqpoint{2.359088in}{2.747738in}}%
\pgfpathlineto{\pgfqpoint{2.359164in}{2.627814in}}%
\pgfpathlineto{\pgfqpoint{2.359850in}{2.668774in}}%
\pgfpathlineto{\pgfqpoint{2.360691in}{2.775677in}}%
\pgfpathlineto{\pgfqpoint{2.361056in}{2.798559in}}%
\pgfpathlineto{\pgfqpoint{2.361567in}{2.644559in}}%
\pgfpathlineto{\pgfqpoint{2.364028in}{2.644559in}}%
\pgfpathlineto{\pgfqpoint{2.364033in}{2.748366in}}%
\pgfpathlineto{\pgfqpoint{2.364783in}{2.684925in}}%
\pgfpathlineto{\pgfqpoint{2.365624in}{2.757943in}}%
\pgfpathlineto{\pgfqpoint{2.366002in}{2.781648in}}%
\pgfpathlineto{\pgfqpoint{2.366713in}{2.644559in}}%
\pgfpathlineto{\pgfqpoint{2.368973in}{2.644559in}}%
\pgfpathlineto{\pgfqpoint{2.368979in}{2.748521in}}%
\pgfpathlineto{\pgfqpoint{2.369048in}{2.641656in}}%
\pgfpathlineto{\pgfqpoint{2.369733in}{2.682022in}}%
\pgfpathlineto{\pgfqpoint{2.370574in}{2.761324in}}%
\pgfpathlineto{\pgfqpoint{2.370947in}{2.784730in}}%
\pgfpathlineto{\pgfqpoint{2.371557in}{2.644559in}}%
\pgfpathlineto{\pgfqpoint{2.373919in}{2.644559in}}%
\pgfpathlineto{\pgfqpoint{2.373924in}{2.749295in}}%
\pgfpathlineto{\pgfqpoint{2.374683in}{2.699267in}}%
\pgfpathlineto{\pgfqpoint{2.374732in}{2.695189in}}%
\pgfpathlineto{\pgfqpoint{2.375079in}{2.724123in}}%
\pgfpathlineto{\pgfqpoint{2.375128in}{2.720042in}}%
\pgfpathlineto{\pgfqpoint{2.375856in}{2.772997in}}%
\pgfpathlineto{\pgfqpoint{2.375920in}{2.757935in}}%
\pgfpathlineto{\pgfqpoint{2.376688in}{2.644559in}}%
\pgfpathlineto{\pgfqpoint{2.378864in}{2.644559in}}%
\pgfpathlineto{\pgfqpoint{2.378869in}{2.750415in}}%
\pgfpathlineto{\pgfqpoint{2.378952in}{2.628043in}}%
\pgfpathlineto{\pgfqpoint{2.379589in}{2.720663in}}%
\pgfpathlineto{\pgfqpoint{2.379639in}{2.669566in}}%
\pgfpathlineto{\pgfqpoint{2.380281in}{2.764065in}}%
\pgfpathlineto{\pgfqpoint{2.380331in}{2.713153in}}%
\pgfpathlineto{\pgfqpoint{2.380838in}{2.798942in}}%
\pgfpathlineto{\pgfqpoint{2.380928in}{2.644559in}}%
\pgfpathlineto{\pgfqpoint{2.381002in}{2.644559in}}%
\pgfpathlineto{\pgfqpoint{2.383810in}{2.644559in}}%
\pgfpathlineto{\pgfqpoint{2.383815in}{2.747735in}}%
\pgfpathlineto{\pgfqpoint{2.384583in}{2.706481in}}%
\pgfpathlineto{\pgfqpoint{2.384633in}{2.688518in}}%
\pgfpathlineto{\pgfqpoint{2.385177in}{2.743715in}}%
\pgfpathlineto{\pgfqpoint{2.385226in}{2.725847in}}%
\pgfpathlineto{\pgfqpoint{2.385752in}{2.779803in}}%
\pgfpathlineto{\pgfqpoint{2.385838in}{2.717708in}}%
\pgfpathlineto{\pgfqpoint{2.385852in}{2.726831in}}%
\pgfpathlineto{\pgfqpoint{2.386625in}{2.644559in}}%
\pgfpathlineto{\pgfqpoint{2.388755in}{2.644559in}}%
\pgfpathlineto{\pgfqpoint{2.388761in}{2.749563in}}%
\pgfpathlineto{\pgfqpoint{2.388859in}{2.642419in}}%
\pgfpathlineto{\pgfqpoint{2.389499in}{2.707848in}}%
\pgfpathlineto{\pgfqpoint{2.389549in}{2.685002in}}%
\pgfpathlineto{\pgfqpoint{2.390093in}{2.745097in}}%
\pgfpathlineto{\pgfqpoint{2.390142in}{2.722316in}}%
\pgfpathlineto{\pgfqpoint{2.390683in}{2.782129in}}%
\pgfpathlineto{\pgfqpoint{2.390784in}{2.696674in}}%
\pgfpathlineto{\pgfqpoint{2.390798in}{2.748338in}}%
\pgfpathlineto{\pgfqpoint{2.391534in}{2.644559in}}%
\pgfpathlineto{\pgfqpoint{2.393701in}{2.644559in}}%
\pgfpathlineto{\pgfqpoint{2.393706in}{2.749611in}}%
\pgfpathlineto{\pgfqpoint{2.394452in}{2.691923in}}%
\pgfpathlineto{\pgfqpoint{2.395589in}{2.770374in}}%
\pgfpathlineto{\pgfqpoint{2.395729in}{2.717165in}}%
\pgfpathlineto{\pgfqpoint{2.395674in}{2.775726in}}%
\pgfpathlineto{\pgfqpoint{2.395743in}{2.729039in}}%
\pgfpathlineto{\pgfqpoint{2.396507in}{2.644559in}}%
\pgfpathlineto{\pgfqpoint{2.398646in}{2.644559in}}%
\pgfpathlineto{\pgfqpoint{2.398651in}{2.748688in}}%
\pgfpathlineto{\pgfqpoint{2.398723in}{2.639748in}}%
\pgfpathlineto{\pgfqpoint{2.399408in}{2.680353in}}%
\pgfpathlineto{\pgfqpoint{2.400249in}{2.763538in}}%
\pgfpathlineto{\pgfqpoint{2.400620in}{2.786803in}}%
\pgfpathlineto{\pgfqpoint{2.401229in}{2.644559in}}%
\pgfpathlineto{\pgfqpoint{2.403592in}{2.644559in}}%
\pgfpathlineto{\pgfqpoint{2.403597in}{2.748216in}}%
\pgfpathlineto{\pgfqpoint{2.403657in}{2.641902in}}%
\pgfpathlineto{\pgfqpoint{2.404339in}{2.680931in}}%
\pgfpathlineto{\pgfqpoint{2.405180in}{2.760806in}}%
\pgfpathlineto{\pgfqpoint{2.405565in}{2.784989in}}%
\pgfpathlineto{\pgfqpoint{2.406176in}{2.644559in}}%
\pgfpathlineto{\pgfqpoint{2.408537in}{2.644559in}}%
\pgfpathlineto{\pgfqpoint{2.408542in}{2.749328in}}%
\pgfpathlineto{\pgfqpoint{2.408655in}{2.641453in}}%
\pgfpathlineto{\pgfqpoint{2.409298in}{2.709540in}}%
\pgfpathlineto{\pgfqpoint{2.409347in}{2.684608in}}%
\pgfpathlineto{\pgfqpoint{2.409891in}{2.746786in}}%
\pgfpathlineto{\pgfqpoint{2.409941in}{2.721925in}}%
\pgfpathlineto{\pgfqpoint{2.410473in}{2.783296in}}%
\pgfpathlineto{\pgfqpoint{2.410566in}{2.696775in}}%
\pgfpathlineto{\pgfqpoint{2.410581in}{2.748224in}}%
\pgfpathlineto{\pgfqpoint{2.411350in}{2.644559in}}%
\pgfpathlineto{\pgfqpoint{2.413482in}{2.644559in}}%
\pgfpathlineto{\pgfqpoint{2.413487in}{2.748627in}}%
\pgfpathlineto{\pgfqpoint{2.413595in}{2.631425in}}%
\pgfpathlineto{\pgfqpoint{2.414235in}{2.718829in}}%
\pgfpathlineto{\pgfqpoint{2.414284in}{2.674288in}}%
\pgfpathlineto{\pgfqpoint{2.414927in}{2.762254in}}%
\pgfpathlineto{\pgfqpoint{2.414977in}{2.717853in}}%
\pgfpathlineto{\pgfqpoint{2.415414in}{2.792801in}}%
\pgfpathlineto{\pgfqpoint{2.415540in}{2.644559in}}%
\pgfpathlineto{\pgfqpoint{2.415652in}{2.644559in}}%
\pgfpathlineto{\pgfqpoint{2.418428in}{2.644559in}}%
\pgfpathlineto{\pgfqpoint{2.418433in}{2.749042in}}%
\pgfpathlineto{\pgfqpoint{2.418510in}{2.635323in}}%
\pgfpathlineto{\pgfqpoint{2.419196in}{2.676483in}}%
\pgfpathlineto{\pgfqpoint{2.420036in}{2.768639in}}%
\pgfpathlineto{\pgfqpoint{2.420402in}{2.791524in}}%
\pgfpathlineto{\pgfqpoint{2.420963in}{2.644559in}}%
\pgfpathlineto{\pgfqpoint{2.423373in}{2.644559in}}%
\pgfpathlineto{\pgfqpoint{2.423379in}{2.748743in}}%
\pgfpathlineto{\pgfqpoint{2.423452in}{2.637932in}}%
\pgfpathlineto{\pgfqpoint{2.424137in}{2.678745in}}%
\pgfpathlineto{\pgfqpoint{2.424978in}{2.765166in}}%
\pgfpathlineto{\pgfqpoint{2.425347in}{2.788295in}}%
\pgfpathlineto{\pgfqpoint{2.425907in}{2.644559in}}%
\pgfpathlineto{\pgfqpoint{2.428319in}{2.644559in}}%
\pgfpathlineto{\pgfqpoint{2.428324in}{2.748385in}}%
\pgfpathlineto{\pgfqpoint{2.428391in}{2.643866in}}%
\pgfpathlineto{\pgfqpoint{2.429076in}{2.683935in}}%
\pgfpathlineto{\pgfqpoint{2.429917in}{2.759020in}}%
\pgfpathlineto{\pgfqpoint{2.430293in}{2.782589in}}%
\pgfpathlineto{\pgfqpoint{2.430952in}{2.644559in}}%
\pgfpathlineto{\pgfqpoint{2.433264in}{2.644559in}}%
\pgfpathlineto{\pgfqpoint{2.433270in}{2.748858in}}%
\pgfpathlineto{\pgfqpoint{2.433347in}{2.634746in}}%
\pgfpathlineto{\pgfqpoint{2.434033in}{2.675946in}}%
\pgfpathlineto{\pgfqpoint{2.434873in}{2.768398in}}%
\pgfpathlineto{\pgfqpoint{2.435238in}{2.791245in}}%
\pgfpathlineto{\pgfqpoint{2.435799in}{2.644559in}}%
\pgfpathlineto{\pgfqpoint{2.438210in}{2.644559in}}%
\pgfpathlineto{\pgfqpoint{2.438215in}{2.749272in}}%
\pgfpathlineto{\pgfqpoint{2.438305in}{2.641350in}}%
\pgfpathlineto{\pgfqpoint{2.438945in}{2.707268in}}%
\pgfpathlineto{\pgfqpoint{2.438994in}{2.683472in}}%
\pgfpathlineto{\pgfqpoint{2.439538in}{2.744516in}}%
\pgfpathlineto{\pgfqpoint{2.439588in}{2.720786in}}%
\pgfpathlineto{\pgfqpoint{2.440184in}{2.785006in}}%
\pgfpathlineto{\pgfqpoint{2.440239in}{2.694064in}}%
\pgfpathlineto{\pgfqpoint{2.440253in}{2.750287in}}%
\pgfpathlineto{\pgfqpoint{2.440996in}{2.644559in}}%
\pgfpathlineto{\pgfqpoint{2.443155in}{2.644559in}}%
\pgfpathlineto{\pgfqpoint{2.443161in}{2.749449in}}%
\pgfpathlineto{\pgfqpoint{2.443914in}{2.692787in}}%
\pgfpathlineto{\pgfqpoint{2.445052in}{2.770851in}}%
\pgfpathlineto{\pgfqpoint{2.445183in}{2.720218in}}%
\pgfpathlineto{\pgfqpoint{2.445129in}{2.775709in}}%
\pgfpathlineto{\pgfqpoint{2.445197in}{2.726459in}}%
\pgfpathlineto{\pgfqpoint{2.445975in}{2.644559in}}%
\pgfpathlineto{\pgfqpoint{2.448101in}{2.644559in}}%
\pgfpathlineto{\pgfqpoint{2.448106in}{2.748247in}}%
\pgfpathlineto{\pgfqpoint{2.448169in}{2.641366in}}%
\pgfpathlineto{\pgfqpoint{2.448851in}{2.680762in}}%
\pgfpathlineto{\pgfqpoint{2.449692in}{2.761240in}}%
\pgfpathlineto{\pgfqpoint{2.450074in}{2.785221in}}%
\pgfpathlineto{\pgfqpoint{2.450683in}{2.644559in}}%
\pgfpathlineto{\pgfqpoint{2.453046in}{2.644559in}}%
\pgfpathlineto{\pgfqpoint{2.453052in}{2.749396in}}%
\pgfpathlineto{\pgfqpoint{2.453156in}{2.643987in}}%
\pgfpathlineto{\pgfqpoint{2.453799in}{2.706486in}}%
\pgfpathlineto{\pgfqpoint{2.453849in}{2.686882in}}%
\pgfpathlineto{\pgfqpoint{2.454393in}{2.743744in}}%
\pgfpathlineto{\pgfqpoint{2.454442in}{2.724187in}}%
\pgfpathlineto{\pgfqpoint{2.454978in}{2.780507in}}%
\pgfpathlineto{\pgfqpoint{2.455075in}{2.715710in}}%
\pgfpathlineto{\pgfqpoint{2.455089in}{2.730488in}}%
\pgfpathlineto{\pgfqpoint{2.455871in}{2.644559in}}%
\pgfpathlineto{\pgfqpoint{2.457992in}{2.644559in}}%
\pgfpathlineto{\pgfqpoint{2.457997in}{2.749234in}}%
\pgfpathlineto{\pgfqpoint{2.458053in}{2.641825in}}%
\pgfpathlineto{\pgfqpoint{2.458731in}{2.679762in}}%
\pgfpathlineto{\pgfqpoint{2.459572in}{2.761206in}}%
\pgfpathlineto{\pgfqpoint{2.459965in}{2.785905in}}%
\pgfpathlineto{\pgfqpoint{2.460553in}{2.644559in}}%
\pgfpathlineto{\pgfqpoint{2.462937in}{2.644559in}}%
\pgfpathlineto{\pgfqpoint{2.462941in}{2.749197in}}%
\pgfpathlineto{\pgfqpoint{2.463063in}{2.631638in}}%
\pgfpathlineto{\pgfqpoint{2.463706in}{2.720451in}}%
\pgfpathlineto{\pgfqpoint{2.463755in}{2.674910in}}%
\pgfpathlineto{\pgfqpoint{2.464398in}{2.763872in}}%
\pgfpathlineto{\pgfqpoint{2.464448in}{2.718479in}}%
\pgfpathlineto{\pgfqpoint{2.464877in}{2.793896in}}%
\pgfpathlineto{\pgfqpoint{2.465003in}{2.644559in}}%
\pgfpathlineto{\pgfqpoint{2.465148in}{2.644559in}}%
\pgfpathlineto{\pgfqpoint{2.467882in}{2.644559in}}%
\pgfpathlineto{\pgfqpoint{2.467888in}{2.749245in}}%
\pgfpathlineto{\pgfqpoint{2.468652in}{2.692836in}}%
\pgfpathlineto{\pgfqpoint{2.469789in}{2.770783in}}%
\pgfpathlineto{\pgfqpoint{2.469911in}{2.718224in}}%
\pgfpathlineto{\pgfqpoint{2.469856in}{2.774989in}}%
\pgfpathlineto{\pgfqpoint{2.469924in}{2.726961in}}%
\pgfpathlineto{\pgfqpoint{2.470689in}{2.644559in}}%
\pgfpathlineto{\pgfqpoint{2.472828in}{2.644559in}}%
\pgfpathlineto{\pgfqpoint{2.472833in}{2.749021in}}%
\pgfpathlineto{\pgfqpoint{2.472910in}{2.635869in}}%
\pgfpathlineto{\pgfqpoint{2.473595in}{2.676979in}}%
\pgfpathlineto{\pgfqpoint{2.474436in}{2.768116in}}%
\pgfpathlineto{\pgfqpoint{2.474802in}{2.791039in}}%
\pgfpathlineto{\pgfqpoint{2.475365in}{2.644559in}}%
\pgfpathlineto{\pgfqpoint{2.477773in}{2.644559in}}%
\pgfpathlineto{\pgfqpoint{2.477779in}{2.748536in}}%
\pgfpathlineto{\pgfqpoint{2.477847in}{2.642393in}}%
\pgfpathlineto{\pgfqpoint{2.478532in}{2.682680in}}%
\pgfpathlineto{\pgfqpoint{2.479373in}{2.760797in}}%
\pgfpathlineto{\pgfqpoint{2.479747in}{2.784249in}}%
\pgfpathlineto{\pgfqpoint{2.480360in}{2.644559in}}%
\pgfpathlineto{\pgfqpoint{2.482719in}{2.644559in}}%
\pgfpathlineto{\pgfqpoint{2.482724in}{2.748951in}}%
\pgfpathlineto{\pgfqpoint{2.482800in}{2.636362in}}%
\pgfpathlineto{\pgfqpoint{2.483485in}{2.677424in}}%
\pgfpathlineto{\pgfqpoint{2.484326in}{2.767424in}}%
\pgfpathlineto{\pgfqpoint{2.484693in}{2.790381in}}%
\pgfpathlineto{\pgfqpoint{2.485251in}{2.644559in}}%
\pgfpathlineto{\pgfqpoint{2.487664in}{2.644559in}}%
\pgfpathlineto{\pgfqpoint{2.487670in}{2.749050in}}%
\pgfpathlineto{\pgfqpoint{2.487733in}{2.640839in}}%
\pgfpathlineto{\pgfqpoint{2.488418in}{2.680364in}}%
\pgfpathlineto{\pgfqpoint{2.489259in}{2.762318in}}%
\pgfpathlineto{\pgfqpoint{2.489638in}{2.786104in}}%
\pgfpathlineto{\pgfqpoint{2.490215in}{2.644559in}}%
\pgfpathlineto{\pgfqpoint{2.492610in}{2.644559in}}%
\pgfpathlineto{\pgfqpoint{2.492615in}{2.749198in}}%
\pgfpathlineto{\pgfqpoint{2.493376in}{2.692671in}}%
\pgfpathlineto{\pgfqpoint{2.494513in}{2.770514in}}%
\pgfpathlineto{\pgfqpoint{2.494638in}{2.718850in}}%
\pgfpathlineto{\pgfqpoint{2.494584in}{2.774922in}}%
\pgfpathlineto{\pgfqpoint{2.494652in}{2.726303in}}%
\pgfpathlineto{\pgfqpoint{2.495417in}{2.644559in}}%
\pgfpathlineto{\pgfqpoint{2.497555in}{2.644559in}}%
\pgfpathlineto{\pgfqpoint{2.497560in}{2.748946in}}%
\pgfpathlineto{\pgfqpoint{2.497635in}{2.637374in}}%
\pgfpathlineto{\pgfqpoint{2.498320in}{2.678268in}}%
\pgfpathlineto{\pgfqpoint{2.499161in}{2.766462in}}%
\pgfpathlineto{\pgfqpoint{2.499529in}{2.789542in}}%
\pgfpathlineto{\pgfqpoint{2.500088in}{2.644559in}}%
\pgfpathlineto{\pgfqpoint{2.502501in}{2.644559in}}%
\pgfpathlineto{\pgfqpoint{2.502506in}{2.747885in}}%
\pgfpathlineto{\pgfqpoint{2.503230in}{2.681902in}}%
\pgfpathlineto{\pgfqpoint{2.504071in}{2.757485in}}%
\pgfpathlineto{\pgfqpoint{2.504446in}{2.780992in}}%
\pgfpathlineto{\pgfqpoint{2.505105in}{2.644559in}}%
\pgfpathlineto{\pgfqpoint{2.507446in}{2.644559in}}%
\pgfpathlineto{\pgfqpoint{2.507452in}{2.749441in}}%
\pgfpathlineto{\pgfqpoint{2.508209in}{2.692768in}}%
\pgfpathlineto{\pgfqpoint{2.509346in}{2.770682in}}%
\pgfpathlineto{\pgfqpoint{2.509474in}{2.719280in}}%
\pgfpathlineto{\pgfqpoint{2.509420in}{2.775301in}}%
\pgfpathlineto{\pgfqpoint{2.509488in}{2.726556in}}%
\pgfpathlineto{\pgfqpoint{2.510253in}{2.644559in}}%
\pgfpathlineto{\pgfqpoint{2.512392in}{2.644559in}}%
\pgfpathlineto{\pgfqpoint{2.512397in}{2.748519in}}%
\pgfpathlineto{\pgfqpoint{2.512465in}{2.643073in}}%
\pgfpathlineto{\pgfqpoint{2.513150in}{2.683274in}}%
\pgfpathlineto{\pgfqpoint{2.513991in}{2.760293in}}%
\pgfpathlineto{\pgfqpoint{2.514365in}{2.783795in}}%
\pgfpathlineto{\pgfqpoint{2.515024in}{2.644559in}}%
\pgfpathlineto{\pgfqpoint{2.517337in}{2.644559in}}%
\pgfpathlineto{\pgfqpoint{2.517342in}{2.748151in}}%
\pgfpathlineto{\pgfqpoint{2.517400in}{2.642647in}}%
\pgfpathlineto{\pgfqpoint{2.518080in}{2.681132in}}%
\pgfpathlineto{\pgfqpoint{2.518921in}{2.760087in}}%
\pgfpathlineto{\pgfqpoint{2.519311in}{2.784547in}}%
\pgfpathlineto{\pgfqpoint{2.519919in}{2.644559in}}%
\pgfpathlineto{\pgfqpoint{2.522282in}{2.644559in}}%
\pgfpathlineto{\pgfqpoint{2.522288in}{2.749340in}}%
\pgfpathlineto{\pgfqpoint{2.522345in}{2.642520in}}%
\pgfpathlineto{\pgfqpoint{2.523023in}{2.680649in}}%
\pgfpathlineto{\pgfqpoint{2.523864in}{2.760962in}}%
\pgfpathlineto{\pgfqpoint{2.524256in}{2.785562in}}%
\pgfpathlineto{\pgfqpoint{2.524844in}{2.644559in}}%
\pgfpathlineto{\pgfqpoint{2.527228in}{2.644559in}}%
\pgfpathlineto{\pgfqpoint{2.527232in}{2.748536in}}%
\pgfpathlineto{\pgfqpoint{2.527309in}{2.627966in}}%
\pgfpathlineto{\pgfqpoint{2.527995in}{2.668893in}}%
\pgfpathlineto{\pgfqpoint{2.528836in}{2.775739in}}%
\pgfpathlineto{\pgfqpoint{2.529202in}{2.798644in}}%
\pgfpathlineto{\pgfqpoint{2.529712in}{2.644559in}}%
\pgfpathlineto{\pgfqpoint{2.532173in}{2.644559in}}%
\pgfpathlineto{\pgfqpoint{2.532179in}{2.747767in}}%
\pgfpathlineto{\pgfqpoint{2.532946in}{2.706453in}}%
\pgfpathlineto{\pgfqpoint{2.532996in}{2.688765in}}%
\pgfpathlineto{\pgfqpoint{2.533540in}{2.743688in}}%
\pgfpathlineto{\pgfqpoint{2.533589in}{2.726093in}}%
\pgfpathlineto{\pgfqpoint{2.534115in}{2.779803in}}%
\pgfpathlineto{\pgfqpoint{2.534201in}{2.718163in}}%
\pgfpathlineto{\pgfqpoint{2.534215in}{2.726660in}}%
\pgfpathlineto{\pgfqpoint{2.534987in}{2.644559in}}%
\pgfpathlineto{\pgfqpoint{2.537119in}{2.644559in}}%
\pgfpathlineto{\pgfqpoint{2.537124in}{2.749262in}}%
\pgfpathlineto{\pgfqpoint{2.537189in}{2.641531in}}%
\pgfpathlineto{\pgfqpoint{2.537874in}{2.681239in}}%
\pgfpathlineto{\pgfqpoint{2.538715in}{2.762514in}}%
\pgfpathlineto{\pgfqpoint{2.539093in}{2.786206in}}%
\pgfpathlineto{\pgfqpoint{2.539669in}{2.644559in}}%
\pgfpathlineto{\pgfqpoint{2.542064in}{2.644559in}}%
\pgfpathlineto{\pgfqpoint{2.542070in}{2.749137in}}%
\pgfpathlineto{\pgfqpoint{2.542126in}{2.644231in}}%
\pgfpathlineto{\pgfqpoint{2.542802in}{2.681951in}}%
\pgfpathlineto{\pgfqpoint{2.543643in}{2.758430in}}%
\pgfpathlineto{\pgfqpoint{2.544038in}{2.783237in}}%
\pgfpathlineto{\pgfqpoint{2.544680in}{2.644559in}}%
\pgfpathlineto{\pgfqpoint{2.547010in}{2.644559in}}%
\pgfpathlineto{\pgfqpoint{2.547015in}{2.749276in}}%
\pgfpathlineto{\pgfqpoint{2.547083in}{2.640731in}}%
\pgfpathlineto{\pgfqpoint{2.547770in}{2.681012in}}%
\pgfpathlineto{\pgfqpoint{2.548611in}{2.763476in}}%
\pgfpathlineto{\pgfqpoint{2.548984in}{2.786831in}}%
\pgfpathlineto{\pgfqpoint{2.549560in}{2.644559in}}%
\pgfpathlineto{\pgfqpoint{2.551955in}{2.644559in}}%
\pgfpathlineto{\pgfqpoint{2.551961in}{2.749149in}}%
\pgfpathlineto{\pgfqpoint{2.552689in}{2.689264in}}%
\pgfpathlineto{\pgfqpoint{2.553728in}{2.763146in}}%
\pgfpathlineto{\pgfqpoint{2.553984in}{2.696905in}}%
\pgfpathlineto{\pgfqpoint{2.553902in}{2.774129in}}%
\pgfpathlineto{\pgfqpoint{2.553999in}{2.747560in}}%
\pgfpathlineto{\pgfqpoint{2.554768in}{2.644559in}}%
\pgfpathlineto{\pgfqpoint{2.556901in}{2.644559in}}%
\pgfpathlineto{\pgfqpoint{2.556905in}{2.749390in}}%
\pgfpathlineto{\pgfqpoint{2.556977in}{2.632582in}}%
\pgfpathlineto{\pgfqpoint{2.557666in}{2.673095in}}%
\pgfpathlineto{\pgfqpoint{2.558507in}{2.771542in}}%
\pgfpathlineto{\pgfqpoint{2.558874in}{2.794575in}}%
\pgfpathlineto{\pgfqpoint{2.559411in}{2.644559in}}%
\pgfpathlineto{\pgfqpoint{2.561846in}{2.644559in}}%
\pgfpathlineto{\pgfqpoint{2.561852in}{2.749554in}}%
\pgfpathlineto{\pgfqpoint{2.562618in}{2.693501in}}%
\pgfpathlineto{\pgfqpoint{2.563756in}{2.771453in}}%
\pgfpathlineto{\pgfqpoint{2.563820in}{2.775477in}}%
\pgfpathlineto{\pgfqpoint{2.565395in}{2.644559in}}%
\pgfpathlineto{\pgfqpoint{2.566792in}{2.644559in}}%
\pgfpathlineto{\pgfqpoint{2.566797in}{2.748223in}}%
\pgfpathlineto{\pgfqpoint{2.566856in}{2.642424in}}%
\pgfpathlineto{\pgfqpoint{2.567536in}{2.681145in}}%
\pgfpathlineto{\pgfqpoint{2.568377in}{2.760507in}}%
\pgfpathlineto{\pgfqpoint{2.568765in}{2.784852in}}%
\pgfpathlineto{\pgfqpoint{2.569378in}{2.644559in}}%
\pgfpathlineto{\pgfqpoint{2.571737in}{2.644559in}}%
\pgfpathlineto{\pgfqpoint{2.571742in}{2.748894in}}%
\pgfpathlineto{\pgfqpoint{2.571816in}{2.637713in}}%
\pgfpathlineto{\pgfqpoint{2.572501in}{2.678567in}}%
\pgfpathlineto{\pgfqpoint{2.573342in}{2.765988in}}%
\pgfpathlineto{\pgfqpoint{2.573711in}{2.789094in}}%
\pgfpathlineto{\pgfqpoint{2.574270in}{2.644559in}}%
\pgfpathlineto{\pgfqpoint{2.576682in}{2.644559in}}%
\pgfpathlineto{\pgfqpoint{2.576688in}{2.747658in}}%
\pgfpathlineto{\pgfqpoint{2.577452in}{2.705427in}}%
\pgfpathlineto{\pgfqpoint{2.577501in}{2.689082in}}%
\pgfpathlineto{\pgfqpoint{2.578045in}{2.742664in}}%
\pgfpathlineto{\pgfqpoint{2.578095in}{2.726408in}}%
\pgfpathlineto{\pgfqpoint{2.578623in}{2.778900in}}%
\pgfpathlineto{\pgfqpoint{2.578710in}{2.719150in}}%
\pgfpathlineto{\pgfqpoint{2.578724in}{2.725314in}}%
\pgfpathlineto{\pgfqpoint{2.579491in}{2.644559in}}%
\pgfpathlineto{\pgfqpoint{2.581628in}{2.644559in}}%
\pgfpathlineto{\pgfqpoint{2.581633in}{2.749364in}}%
\pgfpathlineto{\pgfqpoint{2.581695in}{2.641208in}}%
\pgfpathlineto{\pgfqpoint{2.582379in}{2.680432in}}%
\pgfpathlineto{\pgfqpoint{2.583220in}{2.762901in}}%
\pgfpathlineto{\pgfqpoint{2.583602in}{2.786859in}}%
\pgfpathlineto{\pgfqpoint{2.584178in}{2.644559in}}%
\pgfpathlineto{\pgfqpoint{2.586573in}{2.644559in}}%
\pgfpathlineto{\pgfqpoint{2.586579in}{2.746890in}}%
\pgfpathlineto{\pgfqpoint{2.586695in}{2.625685in}}%
\pgfpathlineto{\pgfqpoint{2.587338in}{2.724940in}}%
\pgfpathlineto{\pgfqpoint{2.587388in}{2.668939in}}%
\pgfpathlineto{\pgfqpoint{2.588031in}{2.768323in}}%
\pgfpathlineto{\pgfqpoint{2.588080in}{2.712545in}}%
\pgfpathlineto{\pgfqpoint{2.588511in}{2.798441in}}%
\pgfpathlineto{\pgfqpoint{2.588636in}{2.644559in}}%
\pgfpathlineto{\pgfqpoint{2.588782in}{2.644559in}}%
\pgfpathlineto{\pgfqpoint{2.591519in}{2.644559in}}%
\pgfpathlineto{\pgfqpoint{2.591524in}{2.749458in}}%
\pgfpathlineto{\pgfqpoint{2.591579in}{2.643212in}}%
\pgfpathlineto{\pgfqpoint{2.592256in}{2.680884in}}%
\pgfpathlineto{\pgfqpoint{2.593097in}{2.760638in}}%
\pgfpathlineto{\pgfqpoint{2.593468in}{2.783922in}}%
\pgfpathlineto{\pgfqpoint{2.594084in}{2.644559in}}%
\pgfpathlineto{\pgfqpoint{2.596464in}{2.644559in}}%
\pgfpathlineto{\pgfqpoint{2.596470in}{2.747643in}}%
\pgfpathlineto{\pgfqpoint{2.596558in}{2.626693in}}%
\pgfpathlineto{\pgfqpoint{2.597195in}{2.721440in}}%
\pgfpathlineto{\pgfqpoint{2.597245in}{2.668640in}}%
\pgfpathlineto{\pgfqpoint{2.597888in}{2.764839in}}%
\pgfpathlineto{\pgfqpoint{2.597937in}{2.712231in}}%
\pgfpathlineto{\pgfqpoint{2.598438in}{2.799316in}}%
\pgfpathlineto{\pgfqpoint{2.598528in}{2.644559in}}%
\pgfpathlineto{\pgfqpoint{2.598651in}{2.644559in}}%
\pgfpathlineto{\pgfqpoint{2.601410in}{2.644559in}}%
\pgfpathlineto{\pgfqpoint{2.601415in}{2.748343in}}%
\pgfpathlineto{\pgfqpoint{2.601481in}{2.644366in}}%
\pgfpathlineto{\pgfqpoint{2.602166in}{2.684336in}}%
\pgfpathlineto{\pgfqpoint{2.603007in}{2.758407in}}%
\pgfpathlineto{\pgfqpoint{2.603384in}{2.782037in}}%
\pgfpathlineto{\pgfqpoint{2.604044in}{2.644559in}}%
\pgfpathlineto{\pgfqpoint{2.606355in}{2.644559in}}%
\pgfpathlineto{\pgfqpoint{2.606360in}{2.748487in}}%
\pgfpathlineto{\pgfqpoint{2.606429in}{2.642761in}}%
\pgfpathlineto{\pgfqpoint{2.607114in}{2.682994in}}%
\pgfpathlineto{\pgfqpoint{2.607955in}{2.760277in}}%
\pgfpathlineto{\pgfqpoint{2.608329in}{2.783758in}}%
\pgfpathlineto{\pgfqpoint{2.608987in}{2.644559in}}%
\pgfpathlineto{\pgfqpoint{2.611301in}{2.644559in}}%
\pgfpathlineto{\pgfqpoint{2.611306in}{2.748753in}}%
\pgfpathlineto{\pgfqpoint{2.611378in}{2.638989in}}%
\pgfpathlineto{\pgfqpoint{2.612063in}{2.679689in}}%
\pgfpathlineto{\pgfqpoint{2.612904in}{2.764394in}}%
\pgfpathlineto{\pgfqpoint{2.613274in}{2.787600in}}%
\pgfpathlineto{\pgfqpoint{2.613834in}{2.644559in}}%
\pgfpathlineto{\pgfqpoint{2.616246in}{2.644559in}}%
\pgfpathlineto{\pgfqpoint{2.616251in}{2.747815in}}%
\pgfpathlineto{\pgfqpoint{2.616973in}{2.682041in}}%
\pgfpathlineto{\pgfqpoint{2.617814in}{2.756889in}}%
\pgfpathlineto{\pgfqpoint{2.618190in}{2.780470in}}%
\pgfpathlineto{\pgfqpoint{2.618844in}{2.644559in}}%
\pgfpathlineto{\pgfqpoint{2.621192in}{2.644559in}}%
\pgfpathlineto{\pgfqpoint{2.621197in}{2.749644in}}%
\pgfpathlineto{\pgfqpoint{2.621931in}{2.691403in}}%
\pgfpathlineto{\pgfqpoint{2.623068in}{2.769428in}}%
\pgfpathlineto{\pgfqpoint{2.623220in}{2.715651in}}%
\pgfpathlineto{\pgfqpoint{2.623165in}{2.775519in}}%
\pgfpathlineto{\pgfqpoint{2.623234in}{2.730713in}}%
\pgfpathlineto{\pgfqpoint{2.623996in}{2.644559in}}%
\pgfpathlineto{\pgfqpoint{2.626137in}{2.644559in}}%
\pgfpathlineto{\pgfqpoint{2.626142in}{2.748569in}}%
\pgfpathlineto{\pgfqpoint{2.626212in}{2.640746in}}%
\pgfpathlineto{\pgfqpoint{2.626898in}{2.681217in}}%
\pgfpathlineto{\pgfqpoint{2.627738in}{2.762293in}}%
\pgfpathlineto{\pgfqpoint{2.628111in}{2.785639in}}%
\pgfpathlineto{\pgfqpoint{2.628709in}{2.644559in}}%
\pgfpathlineto{\pgfqpoint{2.631082in}{2.644559in}}%
\pgfpathlineto{\pgfqpoint{2.631088in}{2.749590in}}%
\pgfpathlineto{\pgfqpoint{2.631198in}{2.642255in}}%
\pgfpathlineto{\pgfqpoint{2.631840in}{2.709305in}}%
\pgfpathlineto{\pgfqpoint{2.631890in}{2.685314in}}%
\pgfpathlineto{\pgfqpoint{2.632434in}{2.746552in}}%
\pgfpathlineto{\pgfqpoint{2.632483in}{2.722630in}}%
\pgfpathlineto{\pgfqpoint{2.633017in}{2.783161in}}%
\pgfpathlineto{\pgfqpoint{2.633111in}{2.697191in}}%
\pgfpathlineto{\pgfqpoint{2.633126in}{2.748649in}}%
\pgfpathlineto{\pgfqpoint{2.633895in}{2.644559in}}%
\pgfpathlineto{\pgfqpoint{2.636028in}{2.644559in}}%
\pgfpathlineto{\pgfqpoint{2.636032in}{2.750091in}}%
\pgfpathlineto{\pgfqpoint{2.636088in}{2.631456in}}%
\pgfpathlineto{\pgfqpoint{2.636764in}{2.668892in}}%
\pgfpathlineto{\pgfqpoint{2.637604in}{2.771296in}}%
\pgfpathlineto{\pgfqpoint{2.637976in}{2.794602in}}%
\pgfpathlineto{\pgfqpoint{2.638483in}{2.644559in}}%
\pgfpathlineto{\pgfqpoint{2.640973in}{2.644559in}}%
\pgfpathlineto{\pgfqpoint{2.640979in}{2.749436in}}%
\pgfpathlineto{\pgfqpoint{2.641053in}{2.644223in}}%
\pgfpathlineto{\pgfqpoint{2.641745in}{2.685273in}}%
\pgfpathlineto{\pgfqpoint{2.642586in}{2.760227in}}%
\pgfpathlineto{\pgfqpoint{2.642947in}{2.782911in}}%
\pgfpathlineto{\pgfqpoint{2.643627in}{2.644559in}}%
\pgfpathlineto{\pgfqpoint{2.645918in}{2.644559in}}%
\pgfpathlineto{\pgfqpoint{2.645924in}{2.753376in}}%
\pgfpathlineto{\pgfqpoint{2.646685in}{2.703611in}}%
\pgfpathlineto{\pgfqpoint{2.646734in}{2.691068in}}%
\pgfpathlineto{\pgfqpoint{2.647278in}{2.740852in}}%
\pgfpathlineto{\pgfqpoint{2.647328in}{2.728390in}}%
\pgfpathlineto{\pgfqpoint{2.647858in}{2.777187in}}%
\pgfpathlineto{\pgfqpoint{2.647946in}{2.720884in}}%
\pgfpathlineto{\pgfqpoint{2.647960in}{2.723969in}}%
\pgfpathlineto{\pgfqpoint{2.648722in}{2.644559in}}%
\pgfpathlineto{\pgfqpoint{2.650864in}{2.644559in}}%
\pgfpathlineto{\pgfqpoint{2.650870in}{2.749074in}}%
\pgfpathlineto{\pgfqpoint{2.650947in}{2.634892in}}%
\pgfpathlineto{\pgfqpoint{2.651633in}{2.676095in}}%
\pgfpathlineto{\pgfqpoint{2.652473in}{2.769121in}}%
\pgfpathlineto{\pgfqpoint{2.652838in}{2.791974in}}%
\pgfpathlineto{\pgfqpoint{2.653398in}{2.644559in}}%
\pgfpathlineto{\pgfqpoint{2.655810in}{2.644559in}}%
\pgfpathlineto{\pgfqpoint{2.655815in}{2.746374in}}%
\pgfpathlineto{\pgfqpoint{2.656552in}{2.698720in}}%
\pgfpathlineto{\pgfqpoint{2.656601in}{2.692446in}}%
\pgfpathlineto{\pgfqpoint{2.657046in}{2.729766in}}%
\pgfpathlineto{\pgfqpoint{2.657096in}{2.723536in}}%
\pgfpathlineto{\pgfqpoint{2.657739in}{2.773229in}}%
\pgfpathlineto{\pgfqpoint{2.657819in}{2.715725in}}%
\pgfpathlineto{\pgfqpoint{2.657828in}{2.749037in}}%
\pgfpathlineto{\pgfqpoint{2.658615in}{2.644559in}}%
\pgfpathlineto{\pgfqpoint{2.660755in}{2.644559in}}%
\pgfpathlineto{\pgfqpoint{2.660760in}{2.750260in}}%
\pgfpathlineto{\pgfqpoint{2.660872in}{2.631620in}}%
\pgfpathlineto{\pgfqpoint{2.661515in}{2.719765in}}%
\pgfpathlineto{\pgfqpoint{2.661564in}{2.674731in}}%
\pgfpathlineto{\pgfqpoint{2.662207in}{2.763188in}}%
\pgfpathlineto{\pgfqpoint{2.662257in}{2.718297in}}%
\pgfpathlineto{\pgfqpoint{2.662691in}{2.793506in}}%
\pgfpathlineto{\pgfqpoint{2.662821in}{2.644559in}}%
\pgfpathlineto{\pgfqpoint{2.662909in}{2.644559in}}%
\pgfpathlineto{\pgfqpoint{2.665700in}{2.644559in}}%
\pgfpathlineto{\pgfqpoint{2.665706in}{2.748789in}}%
\pgfpathlineto{\pgfqpoint{2.665762in}{2.643441in}}%
\pgfpathlineto{\pgfqpoint{2.666441in}{2.681585in}}%
\pgfpathlineto{\pgfqpoint{2.667282in}{2.759904in}}%
\pgfpathlineto{\pgfqpoint{2.667674in}{2.784509in}}%
\pgfpathlineto{\pgfqpoint{2.668282in}{2.644559in}}%
\pgfpathlineto{\pgfqpoint{2.670646in}{2.644559in}}%
\pgfpathlineto{\pgfqpoint{2.670651in}{2.748999in}}%
\pgfpathlineto{\pgfqpoint{2.670732in}{2.632289in}}%
\pgfpathlineto{\pgfqpoint{2.671418in}{2.673738in}}%
\pgfpathlineto{\pgfqpoint{2.672258in}{2.771057in}}%
\pgfpathlineto{\pgfqpoint{2.672620in}{2.793703in}}%
\pgfpathlineto{\pgfqpoint{2.673164in}{2.644559in}}%
\pgfpathlineto{\pgfqpoint{2.675592in}{2.644568in}}%
\pgfpathlineto{\pgfqpoint{2.675595in}{2.772815in}}%
\pgfpathlineto{\pgfqpoint{2.675714in}{2.633225in}}%
\pgfpathlineto{\pgfqpoint{2.676357in}{2.719215in}}%
\pgfpathlineto{\pgfqpoint{2.676407in}{2.676482in}}%
\pgfpathlineto{\pgfqpoint{2.677050in}{2.762617in}}%
\pgfpathlineto{\pgfqpoint{2.677099in}{2.720070in}}%
\pgfpathlineto{\pgfqpoint{2.677530in}{2.792724in}}%
\pgfpathlineto{\pgfqpoint{2.677654in}{2.644559in}}%
\pgfpathlineto{\pgfqpoint{2.677753in}{2.644559in}}%
\pgfpathlineto{\pgfqpoint{2.680537in}{2.644559in}}%
\pgfpathlineto{\pgfqpoint{2.680542in}{2.747874in}}%
\pgfpathlineto{\pgfqpoint{2.681267in}{2.681838in}}%
\pgfpathlineto{\pgfqpoint{2.682108in}{2.757489in}}%
\pgfpathlineto{\pgfqpoint{2.682482in}{2.780988in}}%
\pgfpathlineto{\pgfqpoint{2.683141in}{2.644559in}}%
\pgfpathlineto{\pgfqpoint{2.685481in}{2.644559in}}%
\pgfpathlineto{\pgfqpoint{2.685488in}{2.753268in}}%
\pgfpathlineto{\pgfqpoint{2.685576in}{2.634059in}}%
\pgfpathlineto{\pgfqpoint{2.686213in}{2.714963in}}%
\pgfpathlineto{\pgfqpoint{2.686262in}{2.676030in}}%
\pgfpathlineto{\pgfqpoint{2.686905in}{2.758345in}}%
\pgfpathlineto{\pgfqpoint{2.686955in}{2.719638in}}%
\pgfpathlineto{\pgfqpoint{2.687456in}{2.792852in}}%
\pgfpathlineto{\pgfqpoint{2.687543in}{2.644559in}}%
\pgfpathlineto{\pgfqpoint{2.687621in}{2.644559in}}%
\pgfpathlineto{\pgfqpoint{2.690429in}{2.644562in}}%
\pgfpathlineto{\pgfqpoint{2.690433in}{2.788585in}}%
\pgfpathlineto{\pgfqpoint{2.690509in}{2.621489in}}%
\pgfpathlineto{\pgfqpoint{2.691194in}{2.662655in}}%
\pgfpathlineto{\pgfqpoint{2.692035in}{2.781003in}}%
\pgfpathlineto{\pgfqpoint{2.692402in}{2.803939in}}%
\pgfpathlineto{\pgfqpoint{2.692882in}{2.644559in}}%
\pgfpathlineto{\pgfqpoint{2.695372in}{2.644559in}}%
\pgfpathlineto{\pgfqpoint{2.695378in}{2.752060in}}%
\pgfpathlineto{\pgfqpoint{2.695457in}{2.637860in}}%
\pgfpathlineto{\pgfqpoint{2.696142in}{2.679075in}}%
\pgfpathlineto{\pgfqpoint{2.696983in}{2.766647in}}%
\pgfpathlineto{\pgfqpoint{2.697347in}{2.789467in}}%
\pgfpathlineto{\pgfqpoint{2.697905in}{2.644559in}}%
\pgfpathlineto{\pgfqpoint{2.700319in}{2.644559in}}%
\pgfpathlineto{\pgfqpoint{2.700323in}{2.745506in}}%
\pgfpathlineto{\pgfqpoint{2.701084in}{2.693409in}}%
\pgfpathlineto{\pgfqpoint{2.702293in}{2.773360in}}%
\pgfpathlineto{\pgfqpoint{2.702293in}{2.760207in}}%
\pgfpathlineto{\pgfqpoint{2.703025in}{2.644559in}}%
\pgfpathlineto{\pgfqpoint{2.705264in}{2.644559in}}%
\pgfpathlineto{\pgfqpoint{2.705270in}{2.749356in}}%
\pgfpathlineto{\pgfqpoint{2.706033in}{2.693094in}}%
\pgfpathlineto{\pgfqpoint{2.707171in}{2.771045in}}%
\pgfpathlineto{\pgfqpoint{2.707292in}{2.719688in}}%
\pgfpathlineto{\pgfqpoint{2.707238in}{2.775267in}}%
\pgfpathlineto{\pgfqpoint{2.707306in}{2.726057in}}%
\pgfpathlineto{\pgfqpoint{2.708073in}{2.644559in}}%
\pgfpathlineto{\pgfqpoint{2.710210in}{2.644560in}}%
\pgfpathlineto{\pgfqpoint{2.710214in}{2.761596in}}%
\pgfpathlineto{\pgfqpoint{2.710281in}{2.622888in}}%
\pgfpathlineto{\pgfqpoint{2.710966in}{2.662991in}}%
\pgfpathlineto{\pgfqpoint{2.711807in}{2.780190in}}%
\pgfpathlineto{\pgfqpoint{2.712184in}{2.803773in}}%
\pgfpathlineto{\pgfqpoint{2.712660in}{2.644559in}}%
\pgfpathlineto{\pgfqpoint{2.715155in}{2.644559in}}%
\pgfpathlineto{\pgfqpoint{2.715161in}{2.749090in}}%
\pgfpathlineto{\pgfqpoint{2.715242in}{2.631732in}}%
\pgfpathlineto{\pgfqpoint{2.715927in}{2.673238in}}%
\pgfpathlineto{\pgfqpoint{2.716768in}{2.771887in}}%
\pgfpathlineto{\pgfqpoint{2.717129in}{2.794485in}}%
\pgfpathlineto{\pgfqpoint{2.717642in}{2.644559in}}%
\pgfpathlineto{\pgfqpoint{2.720101in}{2.644559in}}%
\pgfpathlineto{\pgfqpoint{2.720106in}{2.744901in}}%
\pgfpathlineto{\pgfqpoint{2.720844in}{2.699121in}}%
\pgfpathlineto{\pgfqpoint{2.720894in}{2.692320in}}%
\pgfpathlineto{\pgfqpoint{2.721339in}{2.730181in}}%
\pgfpathlineto{\pgfqpoint{2.721388in}{2.723396in}}%
\pgfpathlineto{\pgfqpoint{2.722028in}{2.773470in}}%
\pgfpathlineto{\pgfqpoint{2.722110in}{2.717107in}}%
\pgfpathlineto{\pgfqpoint{2.722119in}{2.747894in}}%
\pgfpathlineto{\pgfqpoint{2.722898in}{2.644559in}}%
\pgfpathlineto{\pgfqpoint{2.725046in}{2.644559in}}%
\pgfpathlineto{\pgfqpoint{2.725051in}{2.749570in}}%
\pgfpathlineto{\pgfqpoint{2.725142in}{2.626905in}}%
\pgfpathlineto{\pgfqpoint{2.725779in}{2.722054in}}%
\pgfpathlineto{\pgfqpoint{2.725829in}{2.668972in}}%
\pgfpathlineto{\pgfqpoint{2.726471in}{2.765451in}}%
\pgfpathlineto{\pgfqpoint{2.726521in}{2.712564in}}%
\pgfpathlineto{\pgfqpoint{2.727020in}{2.799805in}}%
\pgfpathlineto{\pgfqpoint{2.727110in}{2.644559in}}%
\pgfpathlineto{\pgfqpoint{2.727233in}{2.644559in}}%
\pgfpathlineto{\pgfqpoint{2.729991in}{2.644559in}}%
\pgfpathlineto{\pgfqpoint{2.729996in}{2.743295in}}%
\pgfpathlineto{\pgfqpoint{2.730761in}{2.705569in}}%
\pgfpathlineto{\pgfqpoint{2.730810in}{2.688964in}}%
\pgfpathlineto{\pgfqpoint{2.731354in}{2.742823in}}%
\pgfpathlineto{\pgfqpoint{2.731404in}{2.726273in}}%
\pgfpathlineto{\pgfqpoint{2.731932in}{2.779075in}}%
\pgfpathlineto{\pgfqpoint{2.732021in}{2.691389in}}%
\pgfpathlineto{\pgfqpoint{2.732035in}{2.752176in}}%
\pgfpathlineto{\pgfqpoint{2.732789in}{2.644559in}}%
\pgfpathlineto{\pgfqpoint{2.734937in}{2.644559in}}%
\pgfpathlineto{\pgfqpoint{2.734941in}{2.748452in}}%
\pgfpathlineto{\pgfqpoint{2.735071in}{2.631503in}}%
\pgfpathlineto{\pgfqpoint{2.735664in}{2.668638in}}%
\pgfpathlineto{\pgfqpoint{2.736505in}{2.770879in}}%
\pgfpathlineto{\pgfqpoint{2.736881in}{2.794457in}}%
\pgfpathlineto{\pgfqpoint{2.737398in}{2.644559in}}%
\pgfpathlineto{\pgfqpoint{2.739881in}{2.644559in}}%
\pgfpathlineto{\pgfqpoint{2.739887in}{2.757529in}}%
\pgfpathlineto{\pgfqpoint{2.740634in}{2.694683in}}%
\pgfpathlineto{\pgfqpoint{2.741856in}{2.772075in}}%
\pgfpathlineto{\pgfqpoint{2.741889in}{2.736130in}}%
\pgfpathlineto{\pgfqpoint{2.742800in}{2.644559in}}%
\pgfpathlineto{\pgfqpoint{2.744828in}{2.644559in}}%
\pgfpathlineto{\pgfqpoint{2.744833in}{2.748961in}}%
\pgfpathlineto{\pgfqpoint{2.744910in}{2.635281in}}%
\pgfpathlineto{\pgfqpoint{2.745596in}{2.676438in}}%
\pgfpathlineto{\pgfqpoint{2.746436in}{2.768347in}}%
\pgfpathlineto{\pgfqpoint{2.746802in}{2.791230in}}%
\pgfpathlineto{\pgfqpoint{2.747363in}{2.644559in}}%
\pgfpathlineto{\pgfqpoint{2.749773in}{2.644559in}}%
\pgfpathlineto{\pgfqpoint{2.749779in}{2.746610in}}%
\pgfpathlineto{\pgfqpoint{2.750504in}{2.696913in}}%
\pgfpathlineto{\pgfqpoint{2.750553in}{2.693261in}}%
\pgfpathlineto{\pgfqpoint{2.750900in}{2.721753in}}%
\pgfpathlineto{\pgfqpoint{2.750949in}{2.718130in}}%
\pgfpathlineto{\pgfqpoint{2.751747in}{2.774957in}}%
\pgfpathlineto{\pgfqpoint{2.751777in}{2.745992in}}%
\pgfpathlineto{\pgfqpoint{2.752571in}{2.644559in}}%
\pgfpathlineto{\pgfqpoint{2.754719in}{2.644559in}}%
\pgfpathlineto{\pgfqpoint{2.754724in}{2.749165in}}%
\pgfpathlineto{\pgfqpoint{2.755483in}{2.692464in}}%
\pgfpathlineto{\pgfqpoint{2.756621in}{2.770291in}}%
\pgfpathlineto{\pgfqpoint{2.756747in}{2.718102in}}%
\pgfpathlineto{\pgfqpoint{2.756693in}{2.774814in}}%
\pgfpathlineto{\pgfqpoint{2.756761in}{2.726856in}}%
\pgfpathlineto{\pgfqpoint{2.757525in}{2.644559in}}%
\pgfpathlineto{\pgfqpoint{2.759665in}{2.644569in}}%
\pgfpathlineto{\pgfqpoint{2.759668in}{2.774997in}}%
\pgfpathlineto{\pgfqpoint{2.759739in}{2.627373in}}%
\pgfpathlineto{\pgfqpoint{2.760424in}{2.667804in}}%
\pgfpathlineto{\pgfqpoint{2.761265in}{2.776272in}}%
\pgfpathlineto{\pgfqpoint{2.761638in}{2.799656in}}%
\pgfpathlineto{\pgfqpoint{2.762150in}{2.644559in}}%
\pgfpathlineto{\pgfqpoint{2.764611in}{2.644562in}}%
\pgfpathlineto{\pgfqpoint{2.764615in}{2.790021in}}%
\pgfpathlineto{\pgfqpoint{2.764684in}{2.618978in}}%
\pgfpathlineto{\pgfqpoint{2.765369in}{2.659461in}}%
\pgfpathlineto{\pgfqpoint{2.766210in}{2.783250in}}%
\pgfpathlineto{\pgfqpoint{2.766584in}{2.806630in}}%
\pgfpathlineto{\pgfqpoint{2.767078in}{2.644559in}}%
\pgfpathlineto{\pgfqpoint{2.769556in}{2.644586in}}%
\pgfpathlineto{\pgfqpoint{2.769559in}{2.776027in}}%
\pgfpathlineto{\pgfqpoint{2.769625in}{2.642141in}}%
\pgfpathlineto{\pgfqpoint{2.770307in}{2.681462in}}%
\pgfpathlineto{\pgfqpoint{2.771148in}{2.761351in}}%
\pgfpathlineto{\pgfqpoint{2.771529in}{2.785256in}}%
\pgfpathlineto{\pgfqpoint{2.772158in}{2.644559in}}%
\pgfpathlineto{\pgfqpoint{2.774500in}{2.644559in}}%
\pgfpathlineto{\pgfqpoint{2.774506in}{2.752272in}}%
\pgfpathlineto{\pgfqpoint{2.774576in}{2.639812in}}%
\pgfpathlineto{\pgfqpoint{2.775266in}{2.680249in}}%
\pgfpathlineto{\pgfqpoint{2.776106in}{2.764681in}}%
\pgfpathlineto{\pgfqpoint{2.776474in}{2.787780in}}%
\pgfpathlineto{\pgfqpoint{2.777060in}{2.644559in}}%
\pgfpathlineto{\pgfqpoint{2.779445in}{2.644559in}}%
\pgfpathlineto{\pgfqpoint{2.779451in}{2.756592in}}%
\pgfpathlineto{\pgfqpoint{2.779515in}{2.640087in}}%
\pgfpathlineto{\pgfqpoint{2.780197in}{2.679412in}}%
\pgfpathlineto{\pgfqpoint{2.781037in}{2.763369in}}%
\pgfpathlineto{\pgfqpoint{2.781420in}{2.787342in}}%
\pgfpathlineto{\pgfqpoint{2.781978in}{2.644559in}}%
\pgfpathlineto{\pgfqpoint{2.784392in}{2.644559in}}%
\pgfpathlineto{\pgfqpoint{2.784396in}{2.745875in}}%
\pgfpathlineto{\pgfqpoint{2.785122in}{2.696833in}}%
\pgfpathlineto{\pgfqpoint{2.785171in}{2.693039in}}%
\pgfpathlineto{\pgfqpoint{2.785517in}{2.721673in}}%
\pgfpathlineto{\pgfqpoint{2.785567in}{2.717908in}}%
\pgfpathlineto{\pgfqpoint{2.786365in}{2.774910in}}%
\pgfpathlineto{\pgfqpoint{2.786396in}{2.745900in}}%
\pgfpathlineto{\pgfqpoint{2.787189in}{2.644559in}}%
\pgfpathlineto{\pgfqpoint{2.789337in}{2.644559in}}%
\pgfpathlineto{\pgfqpoint{2.789343in}{2.749204in}}%
\pgfpathlineto{\pgfqpoint{2.790066in}{2.690222in}}%
\pgfpathlineto{\pgfqpoint{2.791203in}{2.768233in}}%
\pgfpathlineto{\pgfqpoint{2.791282in}{2.773164in}}%
\pgfpathlineto{\pgfqpoint{2.791375in}{2.711808in}}%
\pgfpathlineto{\pgfqpoint{2.792134in}{2.644559in}}%
\pgfpathlineto{\pgfqpoint{2.794282in}{2.644559in}}%
\pgfpathlineto{\pgfqpoint{2.794288in}{2.749740in}}%
\pgfpathlineto{\pgfqpoint{2.794358in}{2.642481in}}%
\pgfpathlineto{\pgfqpoint{2.795044in}{2.682953in}}%
\pgfpathlineto{\pgfqpoint{2.795884in}{2.761317in}}%
\pgfpathlineto{\pgfqpoint{2.796256in}{2.784635in}}%
\pgfpathlineto{\pgfqpoint{2.796914in}{2.644559in}}%
\pgfpathlineto{\pgfqpoint{2.799228in}{2.644559in}}%
\pgfpathlineto{\pgfqpoint{2.799233in}{2.748766in}}%
\pgfpathlineto{\pgfqpoint{2.799305in}{2.639379in}}%
\pgfpathlineto{\pgfqpoint{2.799990in}{2.680038in}}%
\pgfpathlineto{\pgfqpoint{2.800831in}{2.764164in}}%
\pgfpathlineto{\pgfqpoint{2.801202in}{2.787398in}}%
\pgfpathlineto{\pgfqpoint{2.801811in}{2.644559in}}%
\pgfpathlineto{\pgfqpoint{2.804173in}{2.644559in}}%
\pgfpathlineto{\pgfqpoint{2.804179in}{2.743558in}}%
\pgfpathlineto{\pgfqpoint{2.804946in}{2.706052in}}%
\pgfpathlineto{\pgfqpoint{2.804996in}{2.688933in}}%
\pgfpathlineto{\pgfqpoint{2.805540in}{2.743305in}}%
\pgfpathlineto{\pgfqpoint{2.805589in}{2.726244in}}%
\pgfpathlineto{\pgfqpoint{2.806115in}{2.779440in}}%
\pgfpathlineto{\pgfqpoint{2.806202in}{2.692767in}}%
\pgfpathlineto{\pgfqpoint{2.806216in}{2.750585in}}%
\pgfpathlineto{\pgfqpoint{2.806968in}{2.644559in}}%
\pgfpathlineto{\pgfqpoint{2.809119in}{2.644559in}}%
\pgfpathlineto{\pgfqpoint{2.809123in}{2.748567in}}%
\pgfpathlineto{\pgfqpoint{2.809252in}{2.631546in}}%
\pgfpathlineto{\pgfqpoint{2.809845in}{2.668672in}}%
\pgfpathlineto{\pgfqpoint{2.810686in}{2.770775in}}%
\pgfpathlineto{\pgfqpoint{2.811062in}{2.794379in}}%
\pgfpathlineto{\pgfqpoint{2.811579in}{2.644559in}}%
\pgfpathlineto{\pgfqpoint{2.814063in}{2.644559in}}%
\pgfpathlineto{\pgfqpoint{2.814069in}{2.757293in}}%
\pgfpathlineto{\pgfqpoint{2.814830in}{2.698470in}}%
\pgfpathlineto{\pgfqpoint{2.814879in}{2.696263in}}%
\pgfpathlineto{\pgfqpoint{2.815126in}{2.717098in}}%
\pgfpathlineto{\pgfqpoint{2.815176in}{2.714917in}}%
\pgfpathlineto{\pgfqpoint{2.816003in}{2.772102in}}%
\pgfpathlineto{\pgfqpoint{2.816065in}{2.758007in}}%
\pgfpathlineto{\pgfqpoint{2.816828in}{2.644559in}}%
\pgfpathlineto{\pgfqpoint{2.819010in}{2.644559in}}%
\pgfpathlineto{\pgfqpoint{2.819014in}{2.747904in}}%
\pgfpathlineto{\pgfqpoint{2.819074in}{2.630356in}}%
\pgfpathlineto{\pgfqpoint{2.819753in}{2.668639in}}%
\pgfpathlineto{\pgfqpoint{2.820594in}{2.772779in}}%
\pgfpathlineto{\pgfqpoint{2.820984in}{2.797206in}}%
\pgfpathlineto{\pgfqpoint{2.821497in}{2.644559in}}%
\pgfpathlineto{\pgfqpoint{2.823955in}{2.644559in}}%
\pgfpathlineto{\pgfqpoint{2.823960in}{2.746837in}}%
\pgfpathlineto{\pgfqpoint{2.824687in}{2.697263in}}%
\pgfpathlineto{\pgfqpoint{2.824736in}{2.693381in}}%
\pgfpathlineto{\pgfqpoint{2.825082in}{2.722103in}}%
\pgfpathlineto{\pgfqpoint{2.825132in}{2.718250in}}%
\pgfpathlineto{\pgfqpoint{2.825929in}{2.775258in}}%
\pgfpathlineto{\pgfqpoint{2.825959in}{2.747062in}}%
\pgfpathlineto{\pgfqpoint{2.826753in}{2.644559in}}%
\pgfpathlineto{\pgfqpoint{2.828901in}{2.644559in}}%
\pgfpathlineto{\pgfqpoint{2.828906in}{2.749146in}}%
\pgfpathlineto{\pgfqpoint{2.829644in}{2.691122in}}%
\pgfpathlineto{\pgfqpoint{2.830781in}{2.768693in}}%
\pgfpathlineto{\pgfqpoint{2.830930in}{2.714457in}}%
\pgfpathlineto{\pgfqpoint{2.830874in}{2.774557in}}%
\pgfpathlineto{\pgfqpoint{2.830944in}{2.730534in}}%
\pgfpathlineto{\pgfqpoint{2.831703in}{2.644559in}}%
\pgfpathlineto{\pgfqpoint{2.833845in}{2.644559in}}%
\pgfpathlineto{\pgfqpoint{2.833851in}{2.750781in}}%
\pgfpathlineto{\pgfqpoint{2.834619in}{2.706703in}}%
\pgfpathlineto{\pgfqpoint{2.834669in}{2.689680in}}%
\pgfpathlineto{\pgfqpoint{2.835213in}{2.743937in}}%
\pgfpathlineto{\pgfqpoint{2.835262in}{2.727008in}}%
\pgfpathlineto{\pgfqpoint{2.835788in}{2.780045in}}%
\pgfpathlineto{\pgfqpoint{2.835874in}{2.718455in}}%
\pgfpathlineto{\pgfqpoint{2.835888in}{2.727525in}}%
\pgfpathlineto{\pgfqpoint{2.836662in}{2.644559in}}%
\pgfpathlineto{\pgfqpoint{2.838790in}{2.644559in}}%
\pgfpathlineto{\pgfqpoint{2.838796in}{2.755437in}}%
\pgfpathlineto{\pgfqpoint{2.838915in}{2.629794in}}%
\pgfpathlineto{\pgfqpoint{2.839557in}{2.721613in}}%
\pgfpathlineto{\pgfqpoint{2.839607in}{2.673072in}}%
\pgfpathlineto{\pgfqpoint{2.840250in}{2.765006in}}%
\pgfpathlineto{\pgfqpoint{2.840299in}{2.716668in}}%
\pgfpathlineto{\pgfqpoint{2.840730in}{2.795103in}}%
\pgfpathlineto{\pgfqpoint{2.840854in}{2.644559in}}%
\pgfpathlineto{\pgfqpoint{2.841002in}{2.644559in}}%
\pgfpathlineto{\pgfqpoint{2.843736in}{2.644559in}}%
\pgfpathlineto{\pgfqpoint{2.843742in}{2.753823in}}%
\pgfpathlineto{\pgfqpoint{2.843808in}{2.640839in}}%
\pgfpathlineto{\pgfqpoint{2.844495in}{2.680880in}}%
\pgfpathlineto{\pgfqpoint{2.845335in}{2.762639in}}%
\pgfpathlineto{\pgfqpoint{2.845711in}{2.786191in}}%
\pgfpathlineto{\pgfqpoint{2.846294in}{2.644559in}}%
\pgfpathlineto{\pgfqpoint{2.848681in}{2.644559in}}%
\pgfpathlineto{\pgfqpoint{2.848687in}{2.755301in}}%
\pgfpathlineto{\pgfqpoint{2.848756in}{2.636813in}}%
\pgfpathlineto{\pgfqpoint{2.849441in}{2.677046in}}%
\pgfpathlineto{\pgfqpoint{2.850282in}{2.766424in}}%
\pgfpathlineto{\pgfqpoint{2.850656in}{2.789869in}}%
\pgfpathlineto{\pgfqpoint{2.851202in}{2.644559in}}%
\pgfpathlineto{\pgfqpoint{2.853629in}{2.644565in}}%
\pgfpathlineto{\pgfqpoint{2.853631in}{2.767913in}}%
\pgfpathlineto{\pgfqpoint{2.853742in}{2.635393in}}%
\pgfpathlineto{\pgfqpoint{2.854385in}{2.716319in}}%
\pgfpathlineto{\pgfqpoint{2.854435in}{2.678498in}}%
\pgfpathlineto{\pgfqpoint{2.855078in}{2.759731in}}%
\pgfpathlineto{\pgfqpoint{2.855127in}{2.722075in}}%
\pgfpathlineto{\pgfqpoint{2.855562in}{2.790111in}}%
\pgfpathlineto{\pgfqpoint{2.855691in}{2.644559in}}%
\pgfpathlineto{\pgfqpoint{2.855787in}{2.644559in}}%
\pgfpathlineto{\pgfqpoint{2.858572in}{2.644559in}}%
\pgfpathlineto{\pgfqpoint{2.858578in}{2.753762in}}%
\pgfpathlineto{\pgfqpoint{2.858660in}{2.637059in}}%
\pgfpathlineto{\pgfqpoint{2.859346in}{2.678578in}}%
\pgfpathlineto{\pgfqpoint{2.860187in}{2.765947in}}%
\pgfpathlineto{\pgfqpoint{2.860547in}{2.788523in}}%
\pgfpathlineto{\pgfqpoint{2.861109in}{2.644559in}}%
\pgfpathlineto{\pgfqpoint{2.863519in}{2.644559in}}%
\pgfpathlineto{\pgfqpoint{2.863524in}{2.748776in}}%
\pgfpathlineto{\pgfqpoint{2.863600in}{2.636017in}}%
\pgfpathlineto{\pgfqpoint{2.864286in}{2.677093in}}%
\pgfpathlineto{\pgfqpoint{2.865127in}{2.767006in}}%
\pgfpathlineto{\pgfqpoint{2.865493in}{2.789944in}}%
\pgfpathlineto{\pgfqpoint{2.866051in}{2.644559in}}%
\pgfpathlineto{\pgfqpoint{2.868464in}{2.644559in}}%
\pgfpathlineto{\pgfqpoint{2.868469in}{2.747867in}}%
\pgfpathlineto{\pgfqpoint{2.869190in}{2.682781in}}%
\pgfpathlineto{\pgfqpoint{2.870031in}{2.756866in}}%
\pgfpathlineto{\pgfqpoint{2.870408in}{2.780489in}}%
\pgfpathlineto{\pgfqpoint{2.871079in}{2.644559in}}%
\pgfpathlineto{\pgfqpoint{2.873410in}{2.644559in}}%
\pgfpathlineto{\pgfqpoint{2.873415in}{2.748692in}}%
\pgfpathlineto{\pgfqpoint{2.873486in}{2.640128in}}%
\pgfpathlineto{\pgfqpoint{2.874171in}{2.680691in}}%
\pgfpathlineto{\pgfqpoint{2.875012in}{2.763298in}}%
\pgfpathlineto{\pgfqpoint{2.875384in}{2.786592in}}%
\pgfpathlineto{\pgfqpoint{2.875993in}{2.644559in}}%
\pgfpathlineto{\pgfqpoint{2.878355in}{2.644559in}}%
\pgfpathlineto{\pgfqpoint{2.878361in}{2.745357in}}%
\pgfpathlineto{\pgfqpoint{2.879114in}{2.694204in}}%
\pgfpathlineto{\pgfqpoint{2.880329in}{2.772065in}}%
\pgfpathlineto{\pgfqpoint{2.880329in}{2.747456in}}%
\pgfpathlineto{\pgfqpoint{2.880356in}{2.757580in}}%
\pgfpathlineto{\pgfqpoint{2.881076in}{2.644559in}}%
\pgfpathlineto{\pgfqpoint{2.883301in}{2.644559in}}%
\pgfpathlineto{\pgfqpoint{2.883305in}{2.751126in}}%
\pgfpathlineto{\pgfqpoint{2.883382in}{2.628860in}}%
\pgfpathlineto{\pgfqpoint{2.884068in}{2.669820in}}%
\pgfpathlineto{\pgfqpoint{2.884909in}{2.775587in}}%
\pgfpathlineto{\pgfqpoint{2.885274in}{2.798480in}}%
\pgfpathlineto{\pgfqpoint{2.885785in}{2.644559in}}%
\pgfpathlineto{\pgfqpoint{2.888246in}{2.644559in}}%
\pgfpathlineto{\pgfqpoint{2.888251in}{2.743364in}}%
\pgfpathlineto{\pgfqpoint{2.888325in}{2.640002in}}%
\pgfpathlineto{\pgfqpoint{2.889012in}{2.680928in}}%
\pgfpathlineto{\pgfqpoint{2.889852in}{2.762338in}}%
\pgfpathlineto{\pgfqpoint{2.890220in}{2.785388in}}%
\pgfpathlineto{\pgfqpoint{2.890848in}{2.644559in}}%
\pgfpathlineto{\pgfqpoint{2.893192in}{2.644559in}}%
\pgfpathlineto{\pgfqpoint{2.893197in}{2.749707in}}%
\pgfpathlineto{\pgfqpoint{2.893933in}{2.691531in}}%
\pgfpathlineto{\pgfqpoint{2.895070in}{2.769794in}}%
\pgfpathlineto{\pgfqpoint{2.895220in}{2.716783in}}%
\pgfpathlineto{\pgfqpoint{2.895165in}{2.775757in}}%
\pgfpathlineto{\pgfqpoint{2.895234in}{2.729711in}}%
\pgfpathlineto{\pgfqpoint{2.895993in}{2.644559in}}%
\pgfpathlineto{\pgfqpoint{2.898136in}{2.644559in}}%
\pgfpathlineto{\pgfqpoint{2.898142in}{2.753273in}}%
\pgfpathlineto{\pgfqpoint{2.898206in}{2.640904in}}%
\pgfpathlineto{\pgfqpoint{2.898891in}{2.680552in}}%
\pgfpathlineto{\pgfqpoint{2.899732in}{2.761890in}}%
\pgfpathlineto{\pgfqpoint{2.900111in}{2.785670in}}%
\pgfpathlineto{\pgfqpoint{2.900688in}{2.644559in}}%
\pgfpathlineto{\pgfqpoint{2.903082in}{2.644559in}}%
\pgfpathlineto{\pgfqpoint{2.903087in}{2.748244in}}%
\pgfpathlineto{\pgfqpoint{2.903190in}{2.626298in}}%
\pgfpathlineto{\pgfqpoint{2.903828in}{2.723297in}}%
\pgfpathlineto{\pgfqpoint{2.903878in}{2.668974in}}%
\pgfpathlineto{\pgfqpoint{2.904521in}{2.766690in}}%
\pgfpathlineto{\pgfqpoint{2.904570in}{2.712570in}}%
\pgfpathlineto{\pgfqpoint{2.905011in}{2.797421in}}%
\pgfpathlineto{\pgfqpoint{2.905147in}{2.644559in}}%
\pgfpathlineto{\pgfqpoint{2.905241in}{2.644559in}}%
\pgfpathlineto{\pgfqpoint{2.908027in}{2.644559in}}%
\pgfpathlineto{\pgfqpoint{2.908033in}{2.752703in}}%
\pgfpathlineto{\pgfqpoint{2.908102in}{2.639386in}}%
\pgfpathlineto{\pgfqpoint{2.908787in}{2.679655in}}%
\pgfpathlineto{\pgfqpoint{2.909628in}{2.764187in}}%
\pgfpathlineto{\pgfqpoint{2.910002in}{2.787636in}}%
\pgfpathlineto{\pgfqpoint{2.910561in}{2.644559in}}%
\pgfpathlineto{\pgfqpoint{2.912973in}{2.644559in}}%
\pgfpathlineto{\pgfqpoint{2.912979in}{2.747135in}}%
\pgfpathlineto{\pgfqpoint{2.913735in}{2.694022in}}%
\pgfpathlineto{\pgfqpoint{2.914947in}{2.772478in}}%
\pgfpathlineto{\pgfqpoint{2.914947in}{2.747707in}}%
\pgfpathlineto{\pgfqpoint{2.914969in}{2.755156in}}%
\pgfpathlineto{\pgfqpoint{2.915683in}{2.644559in}}%
\pgfpathlineto{\pgfqpoint{2.917919in}{2.644560in}}%
\pgfpathlineto{\pgfqpoint{2.917923in}{2.761034in}}%
\pgfpathlineto{\pgfqpoint{2.918032in}{2.630629in}}%
\pgfpathlineto{\pgfqpoint{2.918674in}{2.720480in}}%
\pgfpathlineto{\pgfqpoint{2.918724in}{2.673683in}}%
\pgfpathlineto{\pgfqpoint{2.919367in}{2.763880in}}%
\pgfpathlineto{\pgfqpoint{2.919416in}{2.717272in}}%
\pgfpathlineto{\pgfqpoint{2.919852in}{2.794313in}}%
\pgfpathlineto{\pgfqpoint{2.919981in}{2.644559in}}%
\pgfpathlineto{\pgfqpoint{2.920081in}{2.644559in}}%
\pgfpathlineto{\pgfqpoint{2.922863in}{2.644559in}}%
\pgfpathlineto{\pgfqpoint{2.922869in}{2.755444in}}%
\pgfpathlineto{\pgfqpoint{2.922981in}{2.634981in}}%
\pgfpathlineto{\pgfqpoint{2.923624in}{2.715393in}}%
\pgfpathlineto{\pgfqpoint{2.923674in}{2.678146in}}%
\pgfpathlineto{\pgfqpoint{2.924316in}{2.758809in}}%
\pgfpathlineto{\pgfqpoint{2.924366in}{2.721720in}}%
\pgfpathlineto{\pgfqpoint{2.924800in}{2.789114in}}%
\pgfpathlineto{\pgfqpoint{2.924927in}{2.644559in}}%
\pgfpathlineto{\pgfqpoint{2.925023in}{2.644559in}}%
\pgfpathlineto{\pgfqpoint{2.927809in}{2.644559in}}%
\pgfpathlineto{\pgfqpoint{2.927815in}{2.752921in}}%
\pgfpathlineto{\pgfqpoint{2.928539in}{2.682837in}}%
\pgfpathlineto{\pgfqpoint{2.929380in}{2.757390in}}%
\pgfpathlineto{\pgfqpoint{2.929755in}{2.780900in}}%
\pgfpathlineto{\pgfqpoint{2.930416in}{2.644559in}}%
\pgfpathlineto{\pgfqpoint{2.932754in}{2.644559in}}%
\pgfpathlineto{\pgfqpoint{2.932760in}{2.755136in}}%
\pgfpathlineto{\pgfqpoint{2.932830in}{2.636370in}}%
\pgfpathlineto{\pgfqpoint{2.933515in}{2.676678in}}%
\pgfpathlineto{\pgfqpoint{2.934355in}{2.766688in}}%
\pgfpathlineto{\pgfqpoint{2.934729in}{2.790088in}}%
\pgfpathlineto{\pgfqpoint{2.935280in}{2.644559in}}%
\pgfpathlineto{\pgfqpoint{2.937701in}{2.644560in}}%
\pgfpathlineto{\pgfqpoint{2.937703in}{2.750670in}}%
\pgfpathlineto{\pgfqpoint{2.937768in}{2.622027in}}%
\pgfpathlineto{\pgfqpoint{2.938499in}{2.732803in}}%
\pgfpathlineto{\pgfqpoint{2.938549in}{2.667495in}}%
\pgfpathlineto{\pgfqpoint{2.939191in}{2.776153in}}%
\pgfpathlineto{\pgfqpoint{2.939241in}{2.711134in}}%
\pgfpathlineto{\pgfqpoint{2.939674in}{2.806385in}}%
\pgfpathlineto{\pgfqpoint{2.939759in}{2.644559in}}%
\pgfpathlineto{\pgfqpoint{2.939932in}{2.644559in}}%
\pgfpathlineto{\pgfqpoint{2.942646in}{2.644559in}}%
\pgfpathlineto{\pgfqpoint{2.942650in}{2.753989in}}%
\pgfpathlineto{\pgfqpoint{2.942706in}{2.643904in}}%
\pgfpathlineto{\pgfqpoint{2.943383in}{2.681537in}}%
\pgfpathlineto{\pgfqpoint{2.944223in}{2.758609in}}%
\pgfpathlineto{\pgfqpoint{2.944595in}{2.781894in}}%
\pgfpathlineto{\pgfqpoint{2.945255in}{2.644559in}}%
\pgfpathlineto{\pgfqpoint{2.947592in}{2.644559in}}%
\pgfpathlineto{\pgfqpoint{2.947597in}{2.749477in}}%
\pgfpathlineto{\pgfqpoint{2.948347in}{2.698644in}}%
\pgfpathlineto{\pgfqpoint{2.948396in}{2.695306in}}%
\pgfpathlineto{\pgfqpoint{2.948742in}{2.723502in}}%
\pgfpathlineto{\pgfqpoint{2.948792in}{2.720156in}}%
\pgfpathlineto{\pgfqpoint{2.949525in}{2.772656in}}%
\pgfpathlineto{\pgfqpoint{2.949607in}{2.737246in}}%
\pgfpathlineto{\pgfqpoint{2.950433in}{2.644559in}}%
\pgfpathlineto{\pgfqpoint{2.952536in}{2.644559in}}%
\pgfpathlineto{\pgfqpoint{2.952542in}{2.753555in}}%
\pgfpathlineto{\pgfqpoint{2.953272in}{2.682151in}}%
\pgfpathlineto{\pgfqpoint{2.954113in}{2.758441in}}%
\pgfpathlineto{\pgfqpoint{2.954485in}{2.781765in}}%
\pgfpathlineto{\pgfqpoint{2.955141in}{2.644559in}}%
\pgfpathlineto{\pgfqpoint{2.957482in}{2.644559in}}%
\pgfpathlineto{\pgfqpoint{2.957488in}{2.749301in}}%
\pgfpathlineto{\pgfqpoint{2.957561in}{2.641805in}}%
\pgfpathlineto{\pgfqpoint{2.958250in}{2.682664in}}%
\pgfpathlineto{\pgfqpoint{2.959091in}{2.761962in}}%
\pgfpathlineto{\pgfqpoint{2.959456in}{2.784867in}}%
\pgfpathlineto{\pgfqpoint{2.960079in}{2.644559in}}%
\pgfpathlineto{\pgfqpoint{2.962427in}{2.644559in}}%
\pgfpathlineto{\pgfqpoint{2.962434in}{2.748422in}}%
\pgfpathlineto{\pgfqpoint{2.962502in}{2.643580in}}%
\pgfpathlineto{\pgfqpoint{2.963187in}{2.683803in}}%
\pgfpathlineto{\pgfqpoint{2.964028in}{2.759625in}}%
\pgfpathlineto{\pgfqpoint{2.964402in}{2.783075in}}%
\pgfpathlineto{\pgfqpoint{2.965060in}{2.644559in}}%
\pgfpathlineto{\pgfqpoint{2.967373in}{2.644559in}}%
\pgfpathlineto{\pgfqpoint{2.967379in}{2.748170in}}%
\pgfpathlineto{\pgfqpoint{2.967438in}{2.642209in}}%
\pgfpathlineto{\pgfqpoint{2.968119in}{2.680985in}}%
\pgfpathlineto{\pgfqpoint{2.968959in}{2.760467in}}%
\pgfpathlineto{\pgfqpoint{2.969347in}{2.784782in}}%
\pgfpathlineto{\pgfqpoint{2.969959in}{2.644559in}}%
\pgfpathlineto{\pgfqpoint{2.972319in}{2.644559in}}%
\pgfpathlineto{\pgfqpoint{2.972324in}{2.747079in}}%
\pgfpathlineto{\pgfqpoint{2.973072in}{2.701579in}}%
\pgfpathlineto{\pgfqpoint{2.973121in}{2.691888in}}%
\pgfpathlineto{\pgfqpoint{2.973566in}{2.732620in}}%
\pgfpathlineto{\pgfqpoint{2.973616in}{2.722983in}}%
\pgfpathlineto{\pgfqpoint{2.974251in}{2.775589in}}%
\pgfpathlineto{\pgfqpoint{2.974328in}{2.715609in}}%
\pgfpathlineto{\pgfqpoint{2.974337in}{2.750221in}}%
\pgfpathlineto{\pgfqpoint{2.975082in}{2.644559in}}%
\pgfpathlineto{\pgfqpoint{2.977264in}{2.644559in}}%
\pgfpathlineto{\pgfqpoint{2.977270in}{2.747742in}}%
\pgfpathlineto{\pgfqpoint{2.977342in}{2.628114in}}%
\pgfpathlineto{\pgfqpoint{2.978029in}{2.668764in}}%
\pgfpathlineto{\pgfqpoint{2.978870in}{2.775214in}}%
\pgfpathlineto{\pgfqpoint{2.979238in}{2.798288in}}%
\pgfpathlineto{\pgfqpoint{2.979767in}{2.644559in}}%
\pgfpathlineto{\pgfqpoint{2.982209in}{2.644559in}}%
\pgfpathlineto{\pgfqpoint{2.982215in}{2.751538in}}%
\pgfpathlineto{\pgfqpoint{2.982288in}{2.639070in}}%
\pgfpathlineto{\pgfqpoint{2.982977in}{2.679820in}}%
\pgfpathlineto{\pgfqpoint{2.983818in}{2.765800in}}%
\pgfpathlineto{\pgfqpoint{2.984184in}{2.788727in}}%
\pgfpathlineto{\pgfqpoint{2.984771in}{2.644559in}}%
\pgfpathlineto{\pgfqpoint{2.987155in}{2.644559in}}%
\pgfpathlineto{\pgfqpoint{2.987161in}{2.749156in}}%
\pgfpathlineto{\pgfqpoint{2.987890in}{2.690388in}}%
\pgfpathlineto{\pgfqpoint{2.989028in}{2.768622in}}%
\pgfpathlineto{\pgfqpoint{2.989184in}{2.716601in}}%
\pgfpathlineto{\pgfqpoint{2.989103in}{2.773354in}}%
\pgfpathlineto{\pgfqpoint{2.989197in}{2.728221in}}%
\pgfpathlineto{\pgfqpoint{2.989956in}{2.644559in}}%
\pgfpathlineto{\pgfqpoint{2.992100in}{2.644559in}}%
\pgfpathlineto{\pgfqpoint{2.992106in}{2.747720in}}%
\pgfpathlineto{\pgfqpoint{2.992161in}{2.644062in}}%
\pgfpathlineto{\pgfqpoint{2.992839in}{2.681958in}}%
\pgfpathlineto{\pgfqpoint{2.993680in}{2.759686in}}%
\pgfpathlineto{\pgfqpoint{2.994074in}{2.784439in}}%
\pgfpathlineto{\pgfqpoint{2.994683in}{2.644559in}}%
\pgfpathlineto{\pgfqpoint{2.997046in}{2.644559in}}%
\pgfpathlineto{\pgfqpoint{2.997051in}{2.747254in}}%
\pgfpathlineto{\pgfqpoint{2.997793in}{2.700042in}}%
\pgfpathlineto{\pgfqpoint{2.997843in}{2.691539in}}%
\pgfpathlineto{\pgfqpoint{2.998288in}{2.731086in}}%
\pgfpathlineto{\pgfqpoint{2.998337in}{2.722631in}}%
\pgfpathlineto{\pgfqpoint{2.998975in}{2.774243in}}%
\pgfpathlineto{\pgfqpoint{2.999056in}{2.715092in}}%
\pgfpathlineto{\pgfqpoint{2.999065in}{2.749594in}}%
\pgfpathlineto{\pgfqpoint{2.999810in}{2.644559in}}%
\pgfpathlineto{\pgfqpoint{3.001992in}{2.644559in}}%
\pgfpathlineto{\pgfqpoint{3.001996in}{2.756905in}}%
\pgfpathlineto{\pgfqpoint{3.002730in}{2.691003in}}%
\pgfpathlineto{\pgfqpoint{3.003867in}{2.768903in}}%
\pgfpathlineto{\pgfqpoint{3.004022in}{2.712683in}}%
\pgfpathlineto{\pgfqpoint{3.003965in}{2.775051in}}%
\pgfpathlineto{\pgfqpoint{3.004037in}{2.734954in}}%
\pgfpathlineto{\pgfqpoint{3.004825in}{2.644559in}}%
\pgfpathlineto{\pgfqpoint{3.006937in}{2.644559in}}%
\pgfpathlineto{\pgfqpoint{3.006942in}{2.745285in}}%
\pgfpathlineto{\pgfqpoint{3.007706in}{2.705538in}}%
\pgfpathlineto{\pgfqpoint{3.007755in}{2.688832in}}%
\pgfpathlineto{\pgfqpoint{3.008299in}{2.742791in}}%
\pgfpathlineto{\pgfqpoint{3.008349in}{2.726141in}}%
\pgfpathlineto{\pgfqpoint{3.008877in}{2.779057in}}%
\pgfpathlineto{\pgfqpoint{3.008967in}{2.691552in}}%
\pgfpathlineto{\pgfqpoint{3.008980in}{2.751830in}}%
\pgfpathlineto{\pgfqpoint{3.009730in}{2.644559in}}%
\pgfpathlineto{\pgfqpoint{3.011882in}{2.644559in}}%
\pgfpathlineto{\pgfqpoint{3.011888in}{2.748491in}}%
\pgfpathlineto{\pgfqpoint{3.011998in}{2.625262in}}%
\pgfpathlineto{\pgfqpoint{3.012641in}{2.725514in}}%
\pgfpathlineto{\pgfqpoint{3.012691in}{2.668404in}}%
\pgfpathlineto{\pgfqpoint{3.013334in}{2.768899in}}%
\pgfpathlineto{\pgfqpoint{3.013383in}{2.712008in}}%
\pgfpathlineto{\pgfqpoint{3.013818in}{2.799212in}}%
\pgfpathlineto{\pgfqpoint{3.013945in}{2.644559in}}%
\pgfpathlineto{\pgfqpoint{3.014092in}{2.644559in}}%
\pgfpathlineto{\pgfqpoint{3.016827in}{2.644559in}}%
\pgfpathlineto{\pgfqpoint{3.016833in}{2.753407in}}%
\pgfpathlineto{\pgfqpoint{3.017559in}{2.695159in}}%
\pgfpathlineto{\pgfqpoint{3.017608in}{2.695183in}}%
\pgfpathlineto{\pgfqpoint{3.018802in}{2.773266in}}%
\pgfpathlineto{\pgfqpoint{3.018802in}{2.760854in}}%
\pgfpathlineto{\pgfqpoint{3.019548in}{2.644559in}}%
\pgfpathlineto{\pgfqpoint{3.021773in}{2.644559in}}%
\pgfpathlineto{\pgfqpoint{3.021779in}{2.748765in}}%
\pgfpathlineto{\pgfqpoint{3.021855in}{2.636108in}}%
\pgfpathlineto{\pgfqpoint{3.022540in}{2.677178in}}%
\pgfpathlineto{\pgfqpoint{3.023381in}{2.766890in}}%
\pgfpathlineto{\pgfqpoint{3.023747in}{2.789832in}}%
\pgfpathlineto{\pgfqpoint{3.024305in}{2.644559in}}%
\pgfpathlineto{\pgfqpoint{3.026719in}{2.644559in}}%
\pgfpathlineto{\pgfqpoint{3.026723in}{2.747713in}}%
\pgfpathlineto{\pgfqpoint{3.027447in}{2.682559in}}%
\pgfpathlineto{\pgfqpoint{3.028287in}{2.757241in}}%
\pgfpathlineto{\pgfqpoint{3.028663in}{2.780805in}}%
\pgfpathlineto{\pgfqpoint{3.029316in}{2.644559in}}%
\pgfpathlineto{\pgfqpoint{3.031664in}{2.644559in}}%
\pgfpathlineto{\pgfqpoint{3.031670in}{2.749449in}}%
\pgfpathlineto{\pgfqpoint{3.031787in}{2.642688in}}%
\pgfpathlineto{\pgfqpoint{3.032430in}{2.709313in}}%
\pgfpathlineto{\pgfqpoint{3.032479in}{2.685908in}}%
\pgfpathlineto{\pgfqpoint{3.033023in}{2.746560in}}%
\pgfpathlineto{\pgfqpoint{3.033073in}{2.723224in}}%
\pgfpathlineto{\pgfqpoint{3.033603in}{2.782919in}}%
\pgfpathlineto{\pgfqpoint{3.033693in}{2.697005in}}%
\pgfpathlineto{\pgfqpoint{3.033708in}{2.748445in}}%
\pgfpathlineto{\pgfqpoint{3.034477in}{2.644559in}}%
\pgfpathlineto{\pgfqpoint{3.036610in}{2.644559in}}%
\pgfpathlineto{\pgfqpoint{3.036614in}{2.759252in}}%
\pgfpathlineto{\pgfqpoint{3.036746in}{2.631449in}}%
\pgfpathlineto{\pgfqpoint{3.037339in}{2.668609in}}%
\pgfpathlineto{\pgfqpoint{3.038180in}{2.770837in}}%
\pgfpathlineto{\pgfqpoint{3.038555in}{2.794344in}}%
\pgfpathlineto{\pgfqpoint{3.039071in}{2.644559in}}%
\pgfpathlineto{\pgfqpoint{3.041555in}{2.644559in}}%
\pgfpathlineto{\pgfqpoint{3.041561in}{2.749426in}}%
\pgfpathlineto{\pgfqpoint{3.041631in}{2.642260in}}%
\pgfpathlineto{\pgfqpoint{3.042320in}{2.682796in}}%
\pgfpathlineto{\pgfqpoint{3.043161in}{2.761825in}}%
\pgfpathlineto{\pgfqpoint{3.043529in}{2.784925in}}%
\pgfpathlineto{\pgfqpoint{3.044153in}{2.644559in}}%
\pgfpathlineto{\pgfqpoint{3.046500in}{2.644559in}}%
\pgfpathlineto{\pgfqpoint{3.046506in}{2.748948in}}%
\pgfpathlineto{\pgfqpoint{3.046577in}{2.641887in}}%
\pgfpathlineto{\pgfqpoint{3.047262in}{2.682392in}}%
\pgfpathlineto{\pgfqpoint{3.048103in}{2.761150in}}%
\pgfpathlineto{\pgfqpoint{3.048474in}{2.784438in}}%
\pgfpathlineto{\pgfqpoint{3.049079in}{2.644559in}}%
\pgfpathlineto{\pgfqpoint{3.051446in}{2.644559in}}%
\pgfpathlineto{\pgfqpoint{3.051451in}{2.749581in}}%
\pgfpathlineto{\pgfqpoint{3.051541in}{2.642380in}}%
\pgfpathlineto{\pgfqpoint{3.052181in}{2.707247in}}%
\pgfpathlineto{\pgfqpoint{3.052231in}{2.684524in}}%
\pgfpathlineto{\pgfqpoint{3.052775in}{2.744496in}}%
\pgfpathlineto{\pgfqpoint{3.052824in}{2.721837in}}%
\pgfpathlineto{\pgfqpoint{3.053420in}{2.784974in}}%
\pgfpathlineto{\pgfqpoint{3.053479in}{2.693232in}}%
\pgfpathlineto{\pgfqpoint{3.053494in}{2.754550in}}%
\pgfpathlineto{\pgfqpoint{3.054240in}{2.644559in}}%
\pgfpathlineto{\pgfqpoint{3.056391in}{2.644559in}}%
\pgfpathlineto{\pgfqpoint{3.056397in}{2.748166in}}%
\pgfpathlineto{\pgfqpoint{3.056466in}{2.643135in}}%
\pgfpathlineto{\pgfqpoint{3.057151in}{2.683400in}}%
\pgfpathlineto{\pgfqpoint{3.057992in}{2.759275in}}%
\pgfpathlineto{\pgfqpoint{3.058365in}{2.782694in}}%
\pgfpathlineto{\pgfqpoint{3.059023in}{2.644559in}}%
\pgfpathlineto{\pgfqpoint{3.061337in}{2.644559in}}%
\pgfpathlineto{\pgfqpoint{3.061342in}{2.749121in}}%
\pgfpathlineto{\pgfqpoint{3.061421in}{2.633493in}}%
\pgfpathlineto{\pgfqpoint{3.062107in}{2.674826in}}%
\pgfpathlineto{\pgfqpoint{3.062948in}{2.770490in}}%
\pgfpathlineto{\pgfqpoint{3.063311in}{2.793239in}}%
\pgfpathlineto{\pgfqpoint{3.063872in}{2.644559in}}%
\pgfpathlineto{\pgfqpoint{3.066282in}{2.644559in}}%
\pgfpathlineto{\pgfqpoint{3.066287in}{2.747815in}}%
\pgfpathlineto{\pgfqpoint{3.066351in}{2.642249in}}%
\pgfpathlineto{\pgfqpoint{3.067034in}{2.681705in}}%
\pgfpathlineto{\pgfqpoint{3.067875in}{2.762147in}}%
\pgfpathlineto{\pgfqpoint{3.068256in}{2.786082in}}%
\pgfpathlineto{\pgfqpoint{3.068865in}{2.644559in}}%
\pgfpathlineto{\pgfqpoint{3.071228in}{2.644559in}}%
\pgfpathlineto{\pgfqpoint{3.071233in}{2.748312in}}%
\pgfpathlineto{\pgfqpoint{3.071982in}{2.685172in}}%
\pgfpathlineto{\pgfqpoint{3.072922in}{2.763622in}}%
\pgfpathlineto{\pgfqpoint{3.073202in}{2.781167in}}%
\pgfpathlineto{\pgfqpoint{3.074011in}{2.644559in}}%
\pgfpathlineto{\pgfqpoint{3.076173in}{2.644559in}}%
\pgfpathlineto{\pgfqpoint{3.076179in}{2.745172in}}%
\pgfpathlineto{\pgfqpoint{3.076899in}{2.691991in}}%
\pgfpathlineto{\pgfqpoint{3.078116in}{2.770503in}}%
\pgfpathlineto{\pgfqpoint{3.078147in}{2.747667in}}%
\pgfpathlineto{\pgfqpoint{3.078174in}{2.757690in}}%
\pgfpathlineto{\pgfqpoint{3.078888in}{2.644559in}}%
\pgfpathlineto{\pgfqpoint{3.081118in}{2.644559in}}%
\pgfpathlineto{\pgfqpoint{3.081124in}{2.750162in}}%
\pgfpathlineto{\pgfqpoint{3.081198in}{2.634322in}}%
\pgfpathlineto{\pgfqpoint{3.081885in}{2.675118in}}%
\pgfpathlineto{\pgfqpoint{3.082725in}{2.771208in}}%
\pgfpathlineto{\pgfqpoint{3.083093in}{2.794230in}}%
\pgfpathlineto{\pgfqpoint{3.083604in}{2.644559in}}%
\pgfpathlineto{\pgfqpoint{3.086064in}{2.644559in}}%
\pgfpathlineto{\pgfqpoint{3.086070in}{2.745820in}}%
\pgfpathlineto{\pgfqpoint{3.086825in}{2.703748in}}%
\pgfpathlineto{\pgfqpoint{3.086874in}{2.687972in}}%
\pgfpathlineto{\pgfqpoint{3.087418in}{2.741003in}}%
\pgfpathlineto{\pgfqpoint{3.087467in}{2.725280in}}%
\pgfpathlineto{\pgfqpoint{3.088000in}{2.777537in}}%
\pgfpathlineto{\pgfqpoint{3.088100in}{2.696110in}}%
\pgfpathlineto{\pgfqpoint{3.088861in}{2.644559in}}%
\pgfpathlineto{\pgfqpoint{3.091009in}{2.644559in}}%
\pgfpathlineto{\pgfqpoint{3.091015in}{2.749899in}}%
\pgfpathlineto{\pgfqpoint{3.091085in}{2.634601in}}%
\pgfpathlineto{\pgfqpoint{3.091772in}{2.674972in}}%
\pgfpathlineto{\pgfqpoint{3.092613in}{2.770282in}}%
\pgfpathlineto{\pgfqpoint{3.092984in}{2.793534in}}%
\pgfpathlineto{\pgfqpoint{3.093496in}{2.644559in}}%
\pgfpathlineto{\pgfqpoint{3.095955in}{2.644559in}}%
\pgfpathlineto{\pgfqpoint{3.095960in}{2.745298in}}%
\pgfpathlineto{\pgfqpoint{3.096016in}{2.643921in}}%
\pgfpathlineto{\pgfqpoint{3.096693in}{2.681836in}}%
\pgfpathlineto{\pgfqpoint{3.097534in}{2.757419in}}%
\pgfpathlineto{\pgfqpoint{3.097929in}{2.782207in}}%
\pgfpathlineto{\pgfqpoint{3.098568in}{2.644559in}}%
\pgfpathlineto{\pgfqpoint{3.100899in}{2.644559in}}%
\pgfpathlineto{\pgfqpoint{3.100905in}{2.762905in}}%
\pgfpathlineto{\pgfqpoint{3.100985in}{2.627661in}}%
\pgfpathlineto{\pgfqpoint{3.101670in}{2.668850in}}%
\pgfpathlineto{\pgfqpoint{3.102511in}{2.776594in}}%
\pgfpathlineto{\pgfqpoint{3.102874in}{2.799356in}}%
\pgfpathlineto{\pgfqpoint{3.103403in}{2.644559in}}%
\pgfpathlineto{\pgfqpoint{3.105846in}{2.644559in}}%
\pgfpathlineto{\pgfqpoint{3.105852in}{2.746836in}}%
\pgfpathlineto{\pgfqpoint{3.105969in}{2.625713in}}%
\pgfpathlineto{\pgfqpoint{3.106612in}{2.724965in}}%
\pgfpathlineto{\pgfqpoint{3.106662in}{2.668983in}}%
\pgfpathlineto{\pgfqpoint{3.107305in}{2.768349in}}%
\pgfpathlineto{\pgfqpoint{3.107354in}{2.712589in}}%
\pgfpathlineto{\pgfqpoint{3.107785in}{2.798432in}}%
\pgfpathlineto{\pgfqpoint{3.107909in}{2.644559in}}%
\pgfpathlineto{\pgfqpoint{3.108054in}{2.644559in}}%
\pgfpathlineto{\pgfqpoint{3.110790in}{2.644559in}}%
\pgfpathlineto{\pgfqpoint{3.110797in}{2.752733in}}%
\pgfpathlineto{\pgfqpoint{3.110857in}{2.642255in}}%
\pgfpathlineto{\pgfqpoint{3.111537in}{2.681083in}}%
\pgfpathlineto{\pgfqpoint{3.112378in}{2.761392in}}%
\pgfpathlineto{\pgfqpoint{3.112765in}{2.785669in}}%
\pgfpathlineto{\pgfqpoint{3.113365in}{2.644559in}}%
\pgfpathlineto{\pgfqpoint{3.115737in}{2.644559in}}%
\pgfpathlineto{\pgfqpoint{3.115742in}{2.749106in}}%
\pgfpathlineto{\pgfqpoint{3.115822in}{2.633073in}}%
\pgfpathlineto{\pgfqpoint{3.116507in}{2.674443in}}%
\pgfpathlineto{\pgfqpoint{3.117348in}{2.770785in}}%
\pgfpathlineto{\pgfqpoint{3.117711in}{2.793502in}}%
\pgfpathlineto{\pgfqpoint{3.118244in}{2.644559in}}%
\pgfpathlineto{\pgfqpoint{3.120681in}{2.644559in}}%
\pgfpathlineto{\pgfqpoint{3.120687in}{2.754416in}}%
\pgfpathlineto{\pgfqpoint{3.120772in}{2.639238in}}%
\pgfpathlineto{\pgfqpoint{3.121409in}{2.709509in}}%
\pgfpathlineto{\pgfqpoint{3.121459in}{2.680958in}}%
\pgfpathlineto{\pgfqpoint{3.122101in}{2.752913in}}%
\pgfpathlineto{\pgfqpoint{3.122151in}{2.724544in}}%
\pgfpathlineto{\pgfqpoint{3.122656in}{2.787677in}}%
\pgfpathlineto{\pgfqpoint{3.122748in}{2.644559in}}%
\pgfpathlineto{\pgfqpoint{3.122789in}{2.644559in}}%
\pgfpathlineto{\pgfqpoint{3.125628in}{2.644559in}}%
\pgfpathlineto{\pgfqpoint{3.125633in}{2.748779in}}%
\pgfpathlineto{\pgfqpoint{3.125711in}{2.634925in}}%
\pgfpathlineto{\pgfqpoint{3.126396in}{2.676102in}}%
\pgfpathlineto{\pgfqpoint{3.127237in}{2.767958in}}%
\pgfpathlineto{\pgfqpoint{3.127602in}{2.790819in}}%
\pgfpathlineto{\pgfqpoint{3.128163in}{2.644559in}}%
\pgfpathlineto{\pgfqpoint{3.130573in}{2.644559in}}%
\pgfpathlineto{\pgfqpoint{3.130579in}{2.747289in}}%
\pgfpathlineto{\pgfqpoint{3.130638in}{2.642287in}}%
\pgfpathlineto{\pgfqpoint{3.131320in}{2.681122in}}%
\pgfpathlineto{\pgfqpoint{3.132160in}{2.760797in}}%
\pgfpathlineto{\pgfqpoint{3.132547in}{2.785060in}}%
\pgfpathlineto{\pgfqpoint{3.133157in}{2.644559in}}%
\pgfpathlineto{\pgfqpoint{3.135519in}{2.644559in}}%
\pgfpathlineto{\pgfqpoint{3.135524in}{2.748978in}}%
\pgfpathlineto{\pgfqpoint{3.135606in}{2.631276in}}%
\pgfpathlineto{\pgfqpoint{3.136292in}{2.672819in}}%
\pgfpathlineto{\pgfqpoint{3.137132in}{2.771928in}}%
\pgfpathlineto{\pgfqpoint{3.137493in}{2.794489in}}%
\pgfpathlineto{\pgfqpoint{3.138004in}{2.644559in}}%
\pgfpathlineto{\pgfqpoint{3.140464in}{2.644559in}}%
\pgfpathlineto{\pgfqpoint{3.140469in}{2.745828in}}%
\pgfpathlineto{\pgfqpoint{3.141222in}{2.693913in}}%
\pgfpathlineto{\pgfqpoint{3.142438in}{2.772475in}}%
\pgfpathlineto{\pgfqpoint{3.142438in}{2.747544in}}%
\pgfpathlineto{\pgfqpoint{3.142465in}{2.757661in}}%
\pgfpathlineto{\pgfqpoint{3.143180in}{2.644559in}}%
\pgfpathlineto{\pgfqpoint{3.145410in}{2.644559in}}%
\pgfpathlineto{\pgfqpoint{3.145414in}{2.749290in}}%
\pgfpathlineto{\pgfqpoint{3.145529in}{2.631508in}}%
\pgfpathlineto{\pgfqpoint{3.146171in}{2.719821in}}%
\pgfpathlineto{\pgfqpoint{3.146221in}{2.674679in}}%
\pgfpathlineto{\pgfqpoint{3.146864in}{2.763244in}}%
\pgfpathlineto{\pgfqpoint{3.146913in}{2.718245in}}%
\pgfpathlineto{\pgfqpoint{3.147346in}{2.793493in}}%
\pgfpathlineto{\pgfqpoint{3.147476in}{2.644559in}}%
\pgfpathlineto{\pgfqpoint{3.147564in}{2.644559in}}%
\pgfpathlineto{\pgfqpoint{3.150354in}{2.644559in}}%
\pgfpathlineto{\pgfqpoint{3.150361in}{2.750020in}}%
\pgfpathlineto{\pgfqpoint{3.150430in}{2.643552in}}%
\pgfpathlineto{\pgfqpoint{3.151115in}{2.683880in}}%
\pgfpathlineto{\pgfqpoint{3.151956in}{2.759943in}}%
\pgfpathlineto{\pgfqpoint{3.152329in}{2.783344in}}%
\pgfpathlineto{\pgfqpoint{3.152986in}{2.644559in}}%
\pgfpathlineto{\pgfqpoint{3.155301in}{2.644559in}}%
\pgfpathlineto{\pgfqpoint{3.155305in}{2.746320in}}%
\pgfpathlineto{\pgfqpoint{3.156060in}{2.702702in}}%
\pgfpathlineto{\pgfqpoint{3.156109in}{2.690884in}}%
\pgfpathlineto{\pgfqpoint{3.156653in}{2.739947in}}%
\pgfpathlineto{\pgfqpoint{3.156703in}{2.728201in}}%
\pgfpathlineto{\pgfqpoint{3.157236in}{2.776511in}}%
\pgfpathlineto{\pgfqpoint{3.157337in}{2.701523in}}%
\pgfpathlineto{\pgfqpoint{3.158099in}{2.644559in}}%
\pgfpathlineto{\pgfqpoint{3.160245in}{2.644559in}}%
\pgfpathlineto{\pgfqpoint{3.160251in}{2.751044in}}%
\pgfpathlineto{\pgfqpoint{3.160338in}{2.631997in}}%
\pgfpathlineto{\pgfqpoint{3.160974in}{2.716084in}}%
\pgfpathlineto{\pgfqpoint{3.161024in}{2.673838in}}%
\pgfpathlineto{\pgfqpoint{3.161667in}{2.759461in}}%
\pgfpathlineto{\pgfqpoint{3.161716in}{2.717450in}}%
\pgfpathlineto{\pgfqpoint{3.162220in}{2.794102in}}%
\pgfpathlineto{\pgfqpoint{3.162305in}{2.644559in}}%
\pgfpathlineto{\pgfqpoint{3.162382in}{2.644559in}}%
\pgfpathlineto{\pgfqpoint{3.165193in}{2.644734in}}%
\pgfpathlineto{\pgfqpoint{3.165195in}{2.775206in}}%
\pgfpathlineto{\pgfqpoint{3.165953in}{2.694473in}}%
\pgfpathlineto{\pgfqpoint{3.167165in}{2.772650in}}%
\pgfpathlineto{\pgfqpoint{3.167166in}{2.747808in}}%
\pgfpathlineto{\pgfqpoint{3.167192in}{2.757917in}}%
\pgfpathlineto{\pgfqpoint{3.167907in}{2.644559in}}%
\pgfpathlineto{\pgfqpoint{3.170132in}{2.644559in}}%
\pgfpathlineto{\pgfqpoint{3.170132in}{2.644559in}}%
\pgfusepath{stroke}%
\end{pgfscope}%
\begin{pgfscope}%
\pgfsetrectcap%
\pgfsetmiterjoin%
\pgfsetlinewidth{0.803000pt}%
\definecolor{currentstroke}{rgb}{0.000000,0.000000,0.000000}%
\pgfsetstrokecolor{currentstroke}%
\pgfsetdash{}{0pt}%
\pgfpathmoveto{\pgfqpoint{0.573768in}{2.601404in}}%
\pgfpathlineto{\pgfqpoint{0.573768in}{3.507654in}}%
\pgfusepath{stroke}%
\end{pgfscope}%
\begin{pgfscope}%
\pgfsetrectcap%
\pgfsetmiterjoin%
\pgfsetlinewidth{0.803000pt}%
\definecolor{currentstroke}{rgb}{0.000000,0.000000,0.000000}%
\pgfsetstrokecolor{currentstroke}%
\pgfsetdash{}{0pt}%
\pgfpathmoveto{\pgfqpoint{3.293768in}{2.601404in}}%
\pgfpathlineto{\pgfqpoint{3.293768in}{3.507654in}}%
\pgfusepath{stroke}%
\end{pgfscope}%
\begin{pgfscope}%
\pgfsetrectcap%
\pgfsetmiterjoin%
\pgfsetlinewidth{0.803000pt}%
\definecolor{currentstroke}{rgb}{0.000000,0.000000,0.000000}%
\pgfsetstrokecolor{currentstroke}%
\pgfsetdash{}{0pt}%
\pgfpathmoveto{\pgfqpoint{0.573768in}{2.601404in}}%
\pgfpathlineto{\pgfqpoint{3.293768in}{2.601404in}}%
\pgfusepath{stroke}%
\end{pgfscope}%
\begin{pgfscope}%
\pgfsetrectcap%
\pgfsetmiterjoin%
\pgfsetlinewidth{0.803000pt}%
\definecolor{currentstroke}{rgb}{0.000000,0.000000,0.000000}%
\pgfsetstrokecolor{currentstroke}%
\pgfsetdash{}{0pt}%
\pgfpathmoveto{\pgfqpoint{0.573768in}{3.507654in}}%
\pgfpathlineto{\pgfqpoint{3.293768in}{3.507654in}}%
\pgfusepath{stroke}%
\end{pgfscope}%
\begin{pgfscope}%
\definecolor{textcolor}{rgb}{0.000000,0.000000,0.000000}%
\pgfsetstrokecolor{textcolor}%
\pgfsetfillcolor{textcolor}%
\pgftext[x=0.628168in,y=3.444217in,left,top]{\color{textcolor}{\rmfamily\fontsize{8.000000}{9.600000}\selectfont\catcode`\^=\active\def^{\ifmmode\sp\else\^{}\fi}\catcode`\%=\active\def%{\%}(a)}}%
\end{pgfscope}%
\begin{pgfscope}%
\pgfsetbuttcap%
\pgfsetmiterjoin%
\definecolor{currentfill}{rgb}{1.000000,1.000000,1.000000}%
\pgfsetfillcolor{currentfill}%
\pgfsetlinewidth{0.000000pt}%
\definecolor{currentstroke}{rgb}{0.000000,0.000000,0.000000}%
\pgfsetstrokecolor{currentstroke}%
\pgfsetstrokeopacity{0.000000}%
\pgfsetdash{}{0pt}%
\pgfpathmoveto{\pgfqpoint{0.573768in}{1.532029in}}%
\pgfpathlineto{\pgfqpoint{3.293768in}{1.532029in}}%
\pgfpathlineto{\pgfqpoint{3.293768in}{2.438279in}}%
\pgfpathlineto{\pgfqpoint{0.573768in}{2.438279in}}%
\pgfpathlineto{\pgfqpoint{0.573768in}{1.532029in}}%
\pgfpathclose%
\pgfusepath{fill}%
\end{pgfscope}%
\begin{pgfscope}%
\pgfpathrectangle{\pgfqpoint{0.573768in}{1.532029in}}{\pgfqpoint{2.720000in}{0.906250in}}%
\pgfusepath{clip}%
\pgfsetbuttcap%
\pgfsetroundjoin%
\pgfsetlinewidth{0.501875pt}%
\definecolor{currentstroke}{rgb}{0.690196,0.690196,0.690196}%
\pgfsetstrokecolor{currentstroke}%
\pgfsetstrokeopacity{0.250000}%
\pgfsetdash{{1.850000pt}{0.800000pt}}{0.000000pt}%
\pgfpathmoveto{\pgfqpoint{0.697405in}{1.532029in}}%
\pgfpathlineto{\pgfqpoint{0.697405in}{2.438279in}}%
\pgfusepath{stroke}%
\end{pgfscope}%
\begin{pgfscope}%
\pgfsetbuttcap%
\pgfsetroundjoin%
\definecolor{currentfill}{rgb}{0.000000,0.000000,0.000000}%
\pgfsetfillcolor{currentfill}%
\pgfsetlinewidth{0.803000pt}%
\definecolor{currentstroke}{rgb}{0.000000,0.000000,0.000000}%
\pgfsetstrokecolor{currentstroke}%
\pgfsetdash{}{0pt}%
\pgfsys@defobject{currentmarker}{\pgfqpoint{0.000000in}{-0.048611in}}{\pgfqpoint{0.000000in}{0.000000in}}{%
\pgfpathmoveto{\pgfqpoint{0.000000in}{0.000000in}}%
\pgfpathlineto{\pgfqpoint{0.000000in}{-0.048611in}}%
\pgfusepath{stroke,fill}%
}%
\begin{pgfscope}%
\pgfsys@transformshift{0.697405in}{1.532029in}%
\pgfsys@useobject{currentmarker}{}%
\end{pgfscope}%
\end{pgfscope}%
\begin{pgfscope}%
\pgfpathrectangle{\pgfqpoint{0.573768in}{1.532029in}}{\pgfqpoint{2.720000in}{0.906250in}}%
\pgfusepath{clip}%
\pgfsetbuttcap%
\pgfsetroundjoin%
\pgfsetlinewidth{0.501875pt}%
\definecolor{currentstroke}{rgb}{0.690196,0.690196,0.690196}%
\pgfsetstrokecolor{currentstroke}%
\pgfsetstrokeopacity{0.250000}%
\pgfsetdash{{1.850000pt}{0.800000pt}}{0.000000pt}%
\pgfpathmoveto{\pgfqpoint{1.191950in}{1.532029in}}%
\pgfpathlineto{\pgfqpoint{1.191950in}{2.438279in}}%
\pgfusepath{stroke}%
\end{pgfscope}%
\begin{pgfscope}%
\pgfsetbuttcap%
\pgfsetroundjoin%
\definecolor{currentfill}{rgb}{0.000000,0.000000,0.000000}%
\pgfsetfillcolor{currentfill}%
\pgfsetlinewidth{0.803000pt}%
\definecolor{currentstroke}{rgb}{0.000000,0.000000,0.000000}%
\pgfsetstrokecolor{currentstroke}%
\pgfsetdash{}{0pt}%
\pgfsys@defobject{currentmarker}{\pgfqpoint{0.000000in}{-0.048611in}}{\pgfqpoint{0.000000in}{0.000000in}}{%
\pgfpathmoveto{\pgfqpoint{0.000000in}{0.000000in}}%
\pgfpathlineto{\pgfqpoint{0.000000in}{-0.048611in}}%
\pgfusepath{stroke,fill}%
}%
\begin{pgfscope}%
\pgfsys@transformshift{1.191950in}{1.532029in}%
\pgfsys@useobject{currentmarker}{}%
\end{pgfscope}%
\end{pgfscope}%
\begin{pgfscope}%
\pgfpathrectangle{\pgfqpoint{0.573768in}{1.532029in}}{\pgfqpoint{2.720000in}{0.906250in}}%
\pgfusepath{clip}%
\pgfsetbuttcap%
\pgfsetroundjoin%
\pgfsetlinewidth{0.501875pt}%
\definecolor{currentstroke}{rgb}{0.690196,0.690196,0.690196}%
\pgfsetstrokecolor{currentstroke}%
\pgfsetstrokeopacity{0.250000}%
\pgfsetdash{{1.850000pt}{0.800000pt}}{0.000000pt}%
\pgfpathmoveto{\pgfqpoint{1.686496in}{1.532029in}}%
\pgfpathlineto{\pgfqpoint{1.686496in}{2.438279in}}%
\pgfusepath{stroke}%
\end{pgfscope}%
\begin{pgfscope}%
\pgfsetbuttcap%
\pgfsetroundjoin%
\definecolor{currentfill}{rgb}{0.000000,0.000000,0.000000}%
\pgfsetfillcolor{currentfill}%
\pgfsetlinewidth{0.803000pt}%
\definecolor{currentstroke}{rgb}{0.000000,0.000000,0.000000}%
\pgfsetstrokecolor{currentstroke}%
\pgfsetdash{}{0pt}%
\pgfsys@defobject{currentmarker}{\pgfqpoint{0.000000in}{-0.048611in}}{\pgfqpoint{0.000000in}{0.000000in}}{%
\pgfpathmoveto{\pgfqpoint{0.000000in}{0.000000in}}%
\pgfpathlineto{\pgfqpoint{0.000000in}{-0.048611in}}%
\pgfusepath{stroke,fill}%
}%
\begin{pgfscope}%
\pgfsys@transformshift{1.686496in}{1.532029in}%
\pgfsys@useobject{currentmarker}{}%
\end{pgfscope}%
\end{pgfscope}%
\begin{pgfscope}%
\pgfpathrectangle{\pgfqpoint{0.573768in}{1.532029in}}{\pgfqpoint{2.720000in}{0.906250in}}%
\pgfusepath{clip}%
\pgfsetbuttcap%
\pgfsetroundjoin%
\pgfsetlinewidth{0.501875pt}%
\definecolor{currentstroke}{rgb}{0.690196,0.690196,0.690196}%
\pgfsetstrokecolor{currentstroke}%
\pgfsetstrokeopacity{0.250000}%
\pgfsetdash{{1.850000pt}{0.800000pt}}{0.000000pt}%
\pgfpathmoveto{\pgfqpoint{2.181041in}{1.532029in}}%
\pgfpathlineto{\pgfqpoint{2.181041in}{2.438279in}}%
\pgfusepath{stroke}%
\end{pgfscope}%
\begin{pgfscope}%
\pgfsetbuttcap%
\pgfsetroundjoin%
\definecolor{currentfill}{rgb}{0.000000,0.000000,0.000000}%
\pgfsetfillcolor{currentfill}%
\pgfsetlinewidth{0.803000pt}%
\definecolor{currentstroke}{rgb}{0.000000,0.000000,0.000000}%
\pgfsetstrokecolor{currentstroke}%
\pgfsetdash{}{0pt}%
\pgfsys@defobject{currentmarker}{\pgfqpoint{0.000000in}{-0.048611in}}{\pgfqpoint{0.000000in}{0.000000in}}{%
\pgfpathmoveto{\pgfqpoint{0.000000in}{0.000000in}}%
\pgfpathlineto{\pgfqpoint{0.000000in}{-0.048611in}}%
\pgfusepath{stroke,fill}%
}%
\begin{pgfscope}%
\pgfsys@transformshift{2.181041in}{1.532029in}%
\pgfsys@useobject{currentmarker}{}%
\end{pgfscope}%
\end{pgfscope}%
\begin{pgfscope}%
\pgfpathrectangle{\pgfqpoint{0.573768in}{1.532029in}}{\pgfqpoint{2.720000in}{0.906250in}}%
\pgfusepath{clip}%
\pgfsetbuttcap%
\pgfsetroundjoin%
\pgfsetlinewidth{0.501875pt}%
\definecolor{currentstroke}{rgb}{0.690196,0.690196,0.690196}%
\pgfsetstrokecolor{currentstroke}%
\pgfsetstrokeopacity{0.250000}%
\pgfsetdash{{1.850000pt}{0.800000pt}}{0.000000pt}%
\pgfpathmoveto{\pgfqpoint{2.675587in}{1.532029in}}%
\pgfpathlineto{\pgfqpoint{2.675587in}{2.438279in}}%
\pgfusepath{stroke}%
\end{pgfscope}%
\begin{pgfscope}%
\pgfsetbuttcap%
\pgfsetroundjoin%
\definecolor{currentfill}{rgb}{0.000000,0.000000,0.000000}%
\pgfsetfillcolor{currentfill}%
\pgfsetlinewidth{0.803000pt}%
\definecolor{currentstroke}{rgb}{0.000000,0.000000,0.000000}%
\pgfsetstrokecolor{currentstroke}%
\pgfsetdash{}{0pt}%
\pgfsys@defobject{currentmarker}{\pgfqpoint{0.000000in}{-0.048611in}}{\pgfqpoint{0.000000in}{0.000000in}}{%
\pgfpathmoveto{\pgfqpoint{0.000000in}{0.000000in}}%
\pgfpathlineto{\pgfqpoint{0.000000in}{-0.048611in}}%
\pgfusepath{stroke,fill}%
}%
\begin{pgfscope}%
\pgfsys@transformshift{2.675587in}{1.532029in}%
\pgfsys@useobject{currentmarker}{}%
\end{pgfscope}%
\end{pgfscope}%
\begin{pgfscope}%
\pgfpathrectangle{\pgfqpoint{0.573768in}{1.532029in}}{\pgfqpoint{2.720000in}{0.906250in}}%
\pgfusepath{clip}%
\pgfsetbuttcap%
\pgfsetroundjoin%
\pgfsetlinewidth{0.501875pt}%
\definecolor{currentstroke}{rgb}{0.690196,0.690196,0.690196}%
\pgfsetstrokecolor{currentstroke}%
\pgfsetstrokeopacity{0.250000}%
\pgfsetdash{{1.850000pt}{0.800000pt}}{0.000000pt}%
\pgfpathmoveto{\pgfqpoint{3.170132in}{1.532029in}}%
\pgfpathlineto{\pgfqpoint{3.170132in}{2.438279in}}%
\pgfusepath{stroke}%
\end{pgfscope}%
\begin{pgfscope}%
\pgfsetbuttcap%
\pgfsetroundjoin%
\definecolor{currentfill}{rgb}{0.000000,0.000000,0.000000}%
\pgfsetfillcolor{currentfill}%
\pgfsetlinewidth{0.803000pt}%
\definecolor{currentstroke}{rgb}{0.000000,0.000000,0.000000}%
\pgfsetstrokecolor{currentstroke}%
\pgfsetdash{}{0pt}%
\pgfsys@defobject{currentmarker}{\pgfqpoint{0.000000in}{-0.048611in}}{\pgfqpoint{0.000000in}{0.000000in}}{%
\pgfpathmoveto{\pgfqpoint{0.000000in}{0.000000in}}%
\pgfpathlineto{\pgfqpoint{0.000000in}{-0.048611in}}%
\pgfusepath{stroke,fill}%
}%
\begin{pgfscope}%
\pgfsys@transformshift{3.170132in}{1.532029in}%
\pgfsys@useobject{currentmarker}{}%
\end{pgfscope}%
\end{pgfscope}%
\begin{pgfscope}%
\pgfpathrectangle{\pgfqpoint{0.573768in}{1.532029in}}{\pgfqpoint{2.720000in}{0.906250in}}%
\pgfusepath{clip}%
\pgfsetbuttcap%
\pgfsetroundjoin%
\pgfsetlinewidth{0.501875pt}%
\definecolor{currentstroke}{rgb}{0.690196,0.690196,0.690196}%
\pgfsetstrokecolor{currentstroke}%
\pgfsetstrokeopacity{0.250000}%
\pgfsetdash{{1.850000pt}{0.800000pt}}{0.000000pt}%
\pgfpathmoveto{\pgfqpoint{0.573768in}{1.532029in}}%
\pgfpathlineto{\pgfqpoint{3.293768in}{1.532029in}}%
\pgfusepath{stroke}%
\end{pgfscope}%
\begin{pgfscope}%
\pgfsetbuttcap%
\pgfsetroundjoin%
\definecolor{currentfill}{rgb}{0.000000,0.000000,0.000000}%
\pgfsetfillcolor{currentfill}%
\pgfsetlinewidth{0.803000pt}%
\definecolor{currentstroke}{rgb}{0.000000,0.000000,0.000000}%
\pgfsetstrokecolor{currentstroke}%
\pgfsetdash{}{0pt}%
\pgfsys@defobject{currentmarker}{\pgfqpoint{-0.048611in}{0.000000in}}{\pgfqpoint{-0.000000in}{0.000000in}}{%
\pgfpathmoveto{\pgfqpoint{-0.000000in}{0.000000in}}%
\pgfpathlineto{\pgfqpoint{-0.048611in}{0.000000in}}%
\pgfusepath{stroke,fill}%
}%
\begin{pgfscope}%
\pgfsys@transformshift{0.573768in}{1.532029in}%
\pgfsys@useobject{currentmarker}{}%
\end{pgfscope}%
\end{pgfscope}%
\begin{pgfscope}%
\definecolor{textcolor}{rgb}{0.000000,0.000000,0.000000}%
\pgfsetstrokecolor{textcolor}%
\pgfsetfillcolor{textcolor}%
\pgftext[x=0.417518in, y=1.493449in, left, base]{\color{textcolor}{\rmfamily\fontsize{8.000000}{9.600000}\selectfont\catcode`\^=\active\def^{\ifmmode\sp\else\^{}\fi}\catcode`\%=\active\def%{\%}$\mathdefault{0}$}}%
\end{pgfscope}%
\begin{pgfscope}%
\pgfpathrectangle{\pgfqpoint{0.573768in}{1.532029in}}{\pgfqpoint{2.720000in}{0.906250in}}%
\pgfusepath{clip}%
\pgfsetbuttcap%
\pgfsetroundjoin%
\pgfsetlinewidth{0.501875pt}%
\definecolor{currentstroke}{rgb}{0.690196,0.690196,0.690196}%
\pgfsetstrokecolor{currentstroke}%
\pgfsetstrokeopacity{0.250000}%
\pgfsetdash{{1.850000pt}{0.800000pt}}{0.000000pt}%
\pgfpathmoveto{\pgfqpoint{0.573768in}{1.783766in}}%
\pgfpathlineto{\pgfqpoint{3.293768in}{1.783766in}}%
\pgfusepath{stroke}%
\end{pgfscope}%
\begin{pgfscope}%
\pgfsetbuttcap%
\pgfsetroundjoin%
\definecolor{currentfill}{rgb}{0.000000,0.000000,0.000000}%
\pgfsetfillcolor{currentfill}%
\pgfsetlinewidth{0.803000pt}%
\definecolor{currentstroke}{rgb}{0.000000,0.000000,0.000000}%
\pgfsetstrokecolor{currentstroke}%
\pgfsetdash{}{0pt}%
\pgfsys@defobject{currentmarker}{\pgfqpoint{-0.048611in}{0.000000in}}{\pgfqpoint{-0.000000in}{0.000000in}}{%
\pgfpathmoveto{\pgfqpoint{-0.000000in}{0.000000in}}%
\pgfpathlineto{\pgfqpoint{-0.048611in}{0.000000in}}%
\pgfusepath{stroke,fill}%
}%
\begin{pgfscope}%
\pgfsys@transformshift{0.573768in}{1.783766in}%
\pgfsys@useobject{currentmarker}{}%
\end{pgfscope}%
\end{pgfscope}%
\begin{pgfscope}%
\definecolor{textcolor}{rgb}{0.000000,0.000000,0.000000}%
\pgfsetstrokecolor{textcolor}%
\pgfsetfillcolor{textcolor}%
\pgftext[x=0.358489in, y=1.745185in, left, base]{\color{textcolor}{\rmfamily\fontsize{8.000000}{9.600000}\selectfont\catcode`\^=\active\def^{\ifmmode\sp\else\^{}\fi}\catcode`\%=\active\def%{\%}$\mathdefault{25}$}}%
\end{pgfscope}%
\begin{pgfscope}%
\pgfpathrectangle{\pgfqpoint{0.573768in}{1.532029in}}{\pgfqpoint{2.720000in}{0.906250in}}%
\pgfusepath{clip}%
\pgfsetbuttcap%
\pgfsetroundjoin%
\pgfsetlinewidth{0.501875pt}%
\definecolor{currentstroke}{rgb}{0.690196,0.690196,0.690196}%
\pgfsetstrokecolor{currentstroke}%
\pgfsetstrokeopacity{0.250000}%
\pgfsetdash{{1.850000pt}{0.800000pt}}{0.000000pt}%
\pgfpathmoveto{\pgfqpoint{0.573768in}{2.035502in}}%
\pgfpathlineto{\pgfqpoint{3.293768in}{2.035502in}}%
\pgfusepath{stroke}%
\end{pgfscope}%
\begin{pgfscope}%
\pgfsetbuttcap%
\pgfsetroundjoin%
\definecolor{currentfill}{rgb}{0.000000,0.000000,0.000000}%
\pgfsetfillcolor{currentfill}%
\pgfsetlinewidth{0.803000pt}%
\definecolor{currentstroke}{rgb}{0.000000,0.000000,0.000000}%
\pgfsetstrokecolor{currentstroke}%
\pgfsetdash{}{0pt}%
\pgfsys@defobject{currentmarker}{\pgfqpoint{-0.048611in}{0.000000in}}{\pgfqpoint{-0.000000in}{0.000000in}}{%
\pgfpathmoveto{\pgfqpoint{-0.000000in}{0.000000in}}%
\pgfpathlineto{\pgfqpoint{-0.048611in}{0.000000in}}%
\pgfusepath{stroke,fill}%
}%
\begin{pgfscope}%
\pgfsys@transformshift{0.573768in}{2.035502in}%
\pgfsys@useobject{currentmarker}{}%
\end{pgfscope}%
\end{pgfscope}%
\begin{pgfscope}%
\definecolor{textcolor}{rgb}{0.000000,0.000000,0.000000}%
\pgfsetstrokecolor{textcolor}%
\pgfsetfillcolor{textcolor}%
\pgftext[x=0.358489in, y=1.996921in, left, base]{\color{textcolor}{\rmfamily\fontsize{8.000000}{9.600000}\selectfont\catcode`\^=\active\def^{\ifmmode\sp\else\^{}\fi}\catcode`\%=\active\def%{\%}$\mathdefault{50}$}}%
\end{pgfscope}%
\begin{pgfscope}%
\pgfpathrectangle{\pgfqpoint{0.573768in}{1.532029in}}{\pgfqpoint{2.720000in}{0.906250in}}%
\pgfusepath{clip}%
\pgfsetbuttcap%
\pgfsetroundjoin%
\pgfsetlinewidth{0.501875pt}%
\definecolor{currentstroke}{rgb}{0.690196,0.690196,0.690196}%
\pgfsetstrokecolor{currentstroke}%
\pgfsetstrokeopacity{0.250000}%
\pgfsetdash{{1.850000pt}{0.800000pt}}{0.000000pt}%
\pgfpathmoveto{\pgfqpoint{0.573768in}{2.287238in}}%
\pgfpathlineto{\pgfqpoint{3.293768in}{2.287238in}}%
\pgfusepath{stroke}%
\end{pgfscope}%
\begin{pgfscope}%
\pgfsetbuttcap%
\pgfsetroundjoin%
\definecolor{currentfill}{rgb}{0.000000,0.000000,0.000000}%
\pgfsetfillcolor{currentfill}%
\pgfsetlinewidth{0.803000pt}%
\definecolor{currentstroke}{rgb}{0.000000,0.000000,0.000000}%
\pgfsetstrokecolor{currentstroke}%
\pgfsetdash{}{0pt}%
\pgfsys@defobject{currentmarker}{\pgfqpoint{-0.048611in}{0.000000in}}{\pgfqpoint{-0.000000in}{0.000000in}}{%
\pgfpathmoveto{\pgfqpoint{-0.000000in}{0.000000in}}%
\pgfpathlineto{\pgfqpoint{-0.048611in}{0.000000in}}%
\pgfusepath{stroke,fill}%
}%
\begin{pgfscope}%
\pgfsys@transformshift{0.573768in}{2.287238in}%
\pgfsys@useobject{currentmarker}{}%
\end{pgfscope}%
\end{pgfscope}%
\begin{pgfscope}%
\definecolor{textcolor}{rgb}{0.000000,0.000000,0.000000}%
\pgfsetstrokecolor{textcolor}%
\pgfsetfillcolor{textcolor}%
\pgftext[x=0.358489in, y=2.248658in, left, base]{\color{textcolor}{\rmfamily\fontsize{8.000000}{9.600000}\selectfont\catcode`\^=\active\def^{\ifmmode\sp\else\^{}\fi}\catcode`\%=\active\def%{\%}$\mathdefault{75}$}}%
\end{pgfscope}%
\begin{pgfscope}%
\definecolor{textcolor}{rgb}{0.000000,0.000000,0.000000}%
\pgfsetstrokecolor{textcolor}%
\pgfsetfillcolor{textcolor}%
\pgftext[x=0.302933in,y=1.985154in,,bottom,rotate=90.000000]{\color{textcolor}{\rmfamily\fontsize{8.000000}{9.600000}\selectfont\catcode`\^=\active\def^{\ifmmode\sp\else\^{}\fi}\catcode`\%=\active\def%{\%}$v_{DS}$ (V)}}%
\end{pgfscope}%
\begin{pgfscope}%
\pgfpathrectangle{\pgfqpoint{0.573768in}{1.532029in}}{\pgfqpoint{2.720000in}{0.906250in}}%
\pgfusepath{clip}%
\pgfsetrectcap%
\pgfsetroundjoin%
\pgfsetlinewidth{1.505625pt}%
\definecolor{currentstroke}{rgb}{0.835294,0.368627,0.000000}%
\pgfsetstrokecolor{currentstroke}%
\pgfsetdash{}{0pt}%
\pgfpathmoveto{\pgfqpoint{0.697405in}{1.532030in}}%
\pgfpathlineto{\pgfqpoint{0.697409in}{1.544233in}}%
\pgfpathlineto{\pgfqpoint{0.697413in}{1.555394in}}%
\pgfpathlineto{\pgfqpoint{0.697444in}{1.532028in}}%
\pgfpathlineto{\pgfqpoint{0.698134in}{1.532057in}}%
\pgfpathlineto{\pgfqpoint{0.699411in}{1.532476in}}%
\pgfpathlineto{\pgfqpoint{0.700250in}{1.904272in}}%
\pgfpathlineto{\pgfqpoint{0.700299in}{1.904266in}}%
\pgfpathlineto{\pgfqpoint{0.701832in}{1.907975in}}%
\pgfpathlineto{\pgfqpoint{0.702178in}{1.908645in}}%
\pgfpathlineto{\pgfqpoint{0.702356in}{1.909137in}}%
\pgfpathlineto{\pgfqpoint{0.702455in}{1.532093in}}%
\pgfpathlineto{\pgfqpoint{0.703138in}{1.532118in}}%
\pgfpathlineto{\pgfqpoint{0.704353in}{1.532451in}}%
\pgfpathlineto{\pgfqpoint{0.705270in}{1.913651in}}%
\pgfpathlineto{\pgfqpoint{0.705320in}{1.913599in}}%
\pgfpathlineto{\pgfqpoint{0.706754in}{1.920405in}}%
\pgfpathlineto{\pgfqpoint{0.706803in}{1.920343in}}%
\pgfpathlineto{\pgfqpoint{0.706952in}{1.921291in}}%
\pgfpathlineto{\pgfqpoint{0.707001in}{1.921227in}}%
\pgfpathlineto{\pgfqpoint{0.707301in}{1.922731in}}%
\pgfpathlineto{\pgfqpoint{0.707406in}{1.532167in}}%
\pgfpathlineto{\pgfqpoint{0.708090in}{1.532191in}}%
\pgfpathlineto{\pgfqpoint{0.709301in}{1.533449in}}%
\pgfpathlineto{\pgfqpoint{0.710181in}{1.928645in}}%
\pgfpathlineto{\pgfqpoint{0.712248in}{1.942212in}}%
\pgfpathlineto{\pgfqpoint{0.712409in}{1.532235in}}%
\pgfpathlineto{\pgfqpoint{0.713002in}{1.532257in}}%
\pgfpathlineto{\pgfqpoint{0.714242in}{1.532649in}}%
\pgfpathlineto{\pgfqpoint{0.715276in}{1.950893in}}%
\pgfpathlineto{\pgfqpoint{0.717193in}{1.967029in}}%
\pgfpathlineto{\pgfqpoint{0.717382in}{1.532288in}}%
\pgfpathlineto{\pgfqpoint{0.717976in}{1.532310in}}%
\pgfpathlineto{\pgfqpoint{0.719186in}{1.532732in}}%
\pgfpathlineto{\pgfqpoint{0.720237in}{1.977462in}}%
\pgfpathlineto{\pgfqpoint{0.722138in}{1.996431in}}%
\pgfpathlineto{\pgfqpoint{0.722329in}{1.532337in}}%
\pgfpathlineto{\pgfqpoint{0.722922in}{1.532359in}}%
\pgfpathlineto{\pgfqpoint{0.724133in}{1.532921in}}%
\pgfpathlineto{\pgfqpoint{0.725132in}{2.007584in}}%
\pgfpathlineto{\pgfqpoint{0.727084in}{2.029585in}}%
\pgfpathlineto{\pgfqpoint{0.727236in}{1.532376in}}%
\pgfpathlineto{\pgfqpoint{0.727830in}{1.532398in}}%
\pgfpathlineto{\pgfqpoint{0.729078in}{1.532951in}}%
\pgfpathlineto{\pgfqpoint{0.730182in}{2.043050in}}%
\pgfpathlineto{\pgfqpoint{0.732030in}{2.065590in}}%
\pgfpathlineto{\pgfqpoint{0.732162in}{1.532415in}}%
\pgfpathlineto{\pgfqpoint{0.732803in}{1.532439in}}%
\pgfpathlineto{\pgfqpoint{0.734020in}{1.532790in}}%
\pgfpathlineto{\pgfqpoint{0.734047in}{1.611087in}}%
\pgfpathlineto{\pgfqpoint{0.734832in}{2.075687in}}%
\pgfpathlineto{\pgfqpoint{0.736975in}{2.103348in}}%
\pgfpathlineto{\pgfqpoint{0.737103in}{1.532432in}}%
\pgfpathlineto{\pgfqpoint{0.737738in}{1.532457in}}%
\pgfpathlineto{\pgfqpoint{0.738965in}{1.532821in}}%
\pgfpathlineto{\pgfqpoint{0.738992in}{1.611990in}}%
\pgfpathlineto{\pgfqpoint{0.739776in}{2.113588in}}%
\pgfpathlineto{\pgfqpoint{0.741920in}{2.141827in}}%
\pgfpathlineto{\pgfqpoint{0.742071in}{1.532433in}}%
\pgfpathlineto{\pgfqpoint{0.742714in}{1.532469in}}%
\pgfpathlineto{\pgfqpoint{0.743893in}{1.532502in}}%
\pgfpathlineto{\pgfqpoint{0.743909in}{1.532743in}}%
\pgfpathlineto{\pgfqpoint{0.743935in}{1.598418in}}%
\pgfpathlineto{\pgfqpoint{0.744696in}{2.151624in}}%
\pgfpathlineto{\pgfqpoint{0.746866in}{2.179973in}}%
\pgfpathlineto{\pgfqpoint{0.747001in}{1.532426in}}%
\pgfpathlineto{\pgfqpoint{0.747639in}{1.532465in}}%
\pgfpathlineto{\pgfqpoint{0.748838in}{1.532493in}}%
\pgfpathlineto{\pgfqpoint{0.748852in}{1.532684in}}%
\pgfpathlineto{\pgfqpoint{0.748871in}{1.554010in}}%
\pgfpathlineto{\pgfqpoint{0.749762in}{2.191012in}}%
\pgfpathlineto{\pgfqpoint{0.751811in}{2.216685in}}%
\pgfpathlineto{\pgfqpoint{0.751949in}{1.532413in}}%
\pgfpathlineto{\pgfqpoint{0.752582in}{1.532445in}}%
\pgfpathlineto{\pgfqpoint{0.753784in}{1.532481in}}%
\pgfpathlineto{\pgfqpoint{0.753801in}{1.532755in}}%
\pgfpathlineto{\pgfqpoint{0.753828in}{1.604464in}}%
\pgfpathlineto{\pgfqpoint{0.754617in}{2.225683in}}%
\pgfpathlineto{\pgfqpoint{0.756757in}{2.250856in}}%
\pgfpathlineto{\pgfqpoint{0.756906in}{1.532382in}}%
\pgfpathlineto{\pgfqpoint{0.757546in}{1.532422in}}%
\pgfpathlineto{\pgfqpoint{0.758729in}{1.532452in}}%
\pgfpathlineto{\pgfqpoint{0.758743in}{1.532609in}}%
\pgfpathlineto{\pgfqpoint{0.758762in}{1.552726in}}%
\pgfpathlineto{\pgfqpoint{0.759678in}{2.260475in}}%
\pgfpathlineto{\pgfqpoint{0.761702in}{2.281825in}}%
\pgfpathlineto{\pgfqpoint{0.761873in}{1.532345in}}%
\pgfpathlineto{\pgfqpoint{0.762454in}{1.532366in}}%
\pgfpathlineto{\pgfqpoint{0.763696in}{1.532866in}}%
\pgfpathlineto{\pgfqpoint{0.764987in}{2.293834in}}%
\pgfpathlineto{\pgfqpoint{0.766648in}{2.308589in}}%
\pgfpathlineto{\pgfqpoint{0.766819in}{1.532297in}}%
\pgfpathlineto{\pgfqpoint{0.767412in}{1.532319in}}%
\pgfpathlineto{\pgfqpoint{0.768641in}{1.532757in}}%
\pgfpathlineto{\pgfqpoint{0.769915in}{2.318319in}}%
\pgfpathlineto{\pgfqpoint{0.771593in}{2.330454in}}%
\pgfpathlineto{\pgfqpoint{0.771771in}{1.532242in}}%
\pgfpathlineto{\pgfqpoint{0.772365in}{1.532265in}}%
\pgfpathlineto{\pgfqpoint{0.773584in}{1.532517in}}%
\pgfpathlineto{\pgfqpoint{0.773610in}{1.602373in}}%
\pgfpathlineto{\pgfqpoint{0.774422in}{2.335002in}}%
\pgfpathlineto{\pgfqpoint{0.776538in}{2.346966in}}%
\pgfpathlineto{\pgfqpoint{0.776688in}{1.532194in}}%
\pgfpathlineto{\pgfqpoint{0.777326in}{1.532222in}}%
\pgfpathlineto{\pgfqpoint{0.778529in}{1.532417in}}%
\pgfpathlineto{\pgfqpoint{0.778555in}{1.601009in}}%
\pgfpathlineto{\pgfqpoint{0.779329in}{2.349666in}}%
\pgfpathlineto{\pgfqpoint{0.780961in}{2.356465in}}%
\pgfpathlineto{\pgfqpoint{0.781010in}{2.356431in}}%
\pgfpathlineto{\pgfqpoint{0.781478in}{2.357649in}}%
\pgfpathlineto{\pgfqpoint{0.781483in}{2.357647in}}%
\pgfpathlineto{\pgfqpoint{0.781634in}{1.532128in}}%
\pgfpathlineto{\pgfqpoint{0.782265in}{1.532162in}}%
\pgfpathlineto{\pgfqpoint{0.783467in}{1.532232in}}%
\pgfpathlineto{\pgfqpoint{0.783477in}{1.532361in}}%
\pgfpathlineto{\pgfqpoint{0.783519in}{1.970979in}}%
\pgfpathlineto{\pgfqpoint{0.784292in}{2.358992in}}%
\pgfpathlineto{\pgfqpoint{0.784639in}{2.359876in}}%
\pgfpathlineto{\pgfqpoint{0.785776in}{2.362046in}}%
\pgfpathlineto{\pgfqpoint{0.786396in}{2.361703in}}%
\pgfpathlineto{\pgfqpoint{0.786172in}{2.362240in}}%
\pgfpathlineto{\pgfqpoint{0.786428in}{2.362050in}}%
\pgfpathlineto{\pgfqpoint{0.786509in}{1.532060in}}%
\pgfpathlineto{\pgfqpoint{0.787208in}{1.532099in}}%
\pgfpathlineto{\pgfqpoint{0.788413in}{1.532156in}}%
\pgfpathlineto{\pgfqpoint{0.788424in}{1.532266in}}%
\pgfpathlineto{\pgfqpoint{0.789794in}{2.362243in}}%
\pgfpathlineto{\pgfqpoint{0.790634in}{2.360527in}}%
\pgfpathlineto{\pgfqpoint{0.790684in}{2.360788in}}%
\pgfpathlineto{\pgfqpoint{0.791460in}{1.532034in}}%
\pgfpathlineto{\pgfqpoint{0.792119in}{1.532058in}}%
\pgfpathlineto{\pgfqpoint{0.793372in}{1.532262in}}%
\pgfpathlineto{\pgfqpoint{0.794126in}{2.360232in}}%
\pgfpathlineto{\pgfqpoint{0.794621in}{2.359911in}}%
\pgfpathlineto{\pgfqpoint{0.794868in}{2.358853in}}%
\pgfpathlineto{\pgfqpoint{0.794917in}{2.359281in}}%
\pgfpathlineto{\pgfqpoint{0.795143in}{2.104259in}}%
\pgfpathlineto{\pgfqpoint{0.795391in}{1.568262in}}%
\pgfpathlineto{\pgfqpoint{0.795737in}{2.304676in}}%
\pgfpathlineto{\pgfqpoint{0.796033in}{1.615510in}}%
\pgfpathlineto{\pgfqpoint{0.796416in}{1.532028in}}%
\pgfpathlineto{\pgfqpoint{0.796319in}{1.961458in}}%
\pgfpathlineto{\pgfqpoint{0.796891in}{1.532046in}}%
\pgfpathlineto{\pgfqpoint{0.798316in}{1.532229in}}%
\pgfpathlineto{\pgfqpoint{0.799127in}{2.357709in}}%
\pgfpathlineto{\pgfqpoint{0.799572in}{2.357205in}}%
\pgfpathlineto{\pgfqpoint{0.799622in}{2.357396in}}%
\pgfpathlineto{\pgfqpoint{0.799869in}{2.356121in}}%
\pgfpathlineto{\pgfqpoint{0.799919in}{2.356793in}}%
\pgfpathlineto{\pgfqpoint{0.800413in}{1.575326in}}%
\pgfpathlineto{\pgfqpoint{0.800710in}{2.263063in}}%
\pgfpathlineto{\pgfqpoint{0.800759in}{2.298966in}}%
\pgfpathlineto{\pgfqpoint{0.800908in}{1.938131in}}%
\pgfpathlineto{\pgfqpoint{0.801323in}{1.532032in}}%
\pgfpathlineto{\pgfqpoint{0.801792in}{1.532048in}}%
\pgfpathlineto{\pgfqpoint{0.803259in}{1.532180in}}%
\pgfpathlineto{\pgfqpoint{0.803302in}{1.934392in}}%
\pgfpathlineto{\pgfqpoint{0.804075in}{2.355301in}}%
\pgfpathlineto{\pgfqpoint{0.804718in}{2.353578in}}%
\pgfpathlineto{\pgfqpoint{0.804768in}{2.354757in}}%
\pgfpathlineto{\pgfqpoint{0.805361in}{1.577540in}}%
\pgfpathlineto{\pgfqpoint{0.805608in}{2.105116in}}%
\pgfpathlineto{\pgfqpoint{0.805707in}{2.291075in}}%
\pgfpathlineto{\pgfqpoint{0.806053in}{1.583665in}}%
\pgfpathlineto{\pgfqpoint{0.806214in}{1.675004in}}%
\pgfpathlineto{\pgfqpoint{0.806307in}{1.532036in}}%
\pgfpathlineto{\pgfqpoint{0.806987in}{1.532061in}}%
\pgfpathlineto{\pgfqpoint{0.808205in}{1.532191in}}%
\pgfpathlineto{\pgfqpoint{0.808262in}{2.306227in}}%
\pgfpathlineto{\pgfqpoint{0.809025in}{2.352863in}}%
\pgfpathlineto{\pgfqpoint{0.809668in}{2.351197in}}%
\pgfpathlineto{\pgfqpoint{0.809718in}{2.352316in}}%
\pgfpathlineto{\pgfqpoint{0.810311in}{1.576452in}}%
\pgfpathlineto{\pgfqpoint{0.810558in}{2.104802in}}%
\pgfpathlineto{\pgfqpoint{0.810657in}{2.293075in}}%
\pgfpathlineto{\pgfqpoint{0.811004in}{1.582692in}}%
\pgfpathlineto{\pgfqpoint{0.811159in}{1.665303in}}%
\pgfpathlineto{\pgfqpoint{0.811209in}{1.532035in}}%
\pgfpathlineto{\pgfqpoint{0.811924in}{1.532063in}}%
\pgfpathlineto{\pgfqpoint{0.813151in}{1.532203in}}%
\pgfpathlineto{\pgfqpoint{0.813930in}{2.350381in}}%
\pgfpathlineto{\pgfqpoint{0.814523in}{2.349953in}}%
\pgfpathlineto{\pgfqpoint{0.814771in}{2.348935in}}%
\pgfpathlineto{\pgfqpoint{0.814820in}{2.349082in}}%
\pgfpathlineto{\pgfqpoint{0.815038in}{1.950061in}}%
\pgfpathlineto{\pgfqpoint{0.815236in}{1.563928in}}%
\pgfpathlineto{\pgfqpoint{0.815582in}{2.308417in}}%
\pgfpathlineto{\pgfqpoint{0.815879in}{1.620868in}}%
\pgfpathlineto{\pgfqpoint{0.816166in}{1.532036in}}%
\pgfpathlineto{\pgfqpoint{0.816101in}{1.726852in}}%
\pgfpathlineto{\pgfqpoint{0.816746in}{1.532056in}}%
\pgfpathlineto{\pgfqpoint{0.818098in}{1.532226in}}%
\pgfpathlineto{\pgfqpoint{0.818850in}{2.348039in}}%
\pgfpathlineto{\pgfqpoint{0.819394in}{2.347587in}}%
\pgfpathlineto{\pgfqpoint{0.819443in}{2.347696in}}%
\pgfpathlineto{\pgfqpoint{0.819690in}{2.346655in}}%
\pgfpathlineto{\pgfqpoint{0.819740in}{2.347028in}}%
\pgfpathlineto{\pgfqpoint{0.819967in}{2.055007in}}%
\pgfpathlineto{\pgfqpoint{0.820214in}{1.570484in}}%
\pgfpathlineto{\pgfqpoint{0.820560in}{2.306652in}}%
\pgfpathlineto{\pgfqpoint{0.820857in}{1.601953in}}%
\pgfpathlineto{\pgfqpoint{0.821093in}{1.532033in}}%
\pgfpathlineto{\pgfqpoint{0.821046in}{1.711231in}}%
\pgfpathlineto{\pgfqpoint{0.821750in}{1.532056in}}%
\pgfpathlineto{\pgfqpoint{0.823047in}{1.532532in}}%
\pgfpathlineto{\pgfqpoint{0.823869in}{2.345817in}}%
\pgfpathlineto{\pgfqpoint{0.824165in}{2.345745in}}%
\pgfpathlineto{\pgfqpoint{0.824413in}{2.344461in}}%
\pgfpathlineto{\pgfqpoint{0.824462in}{2.345485in}}%
\pgfpathlineto{\pgfqpoint{0.824949in}{1.974891in}}%
\pgfpathlineto{\pgfqpoint{0.825147in}{1.567738in}}%
\pgfpathlineto{\pgfqpoint{0.825542in}{2.307355in}}%
\pgfpathlineto{\pgfqpoint{0.825789in}{1.632425in}}%
\pgfpathlineto{\pgfqpoint{0.826008in}{1.532003in}}%
\pgfpathlineto{\pgfqpoint{0.825992in}{1.672815in}}%
\pgfpathlineto{\pgfqpoint{0.826648in}{1.532043in}}%
\pgfpathlineto{\pgfqpoint{0.827993in}{1.532751in}}%
\pgfpathlineto{\pgfqpoint{0.828830in}{2.343551in}}%
\pgfpathlineto{\pgfqpoint{0.829077in}{2.342800in}}%
\pgfpathlineto{\pgfqpoint{0.829126in}{2.343483in}}%
\pgfpathlineto{\pgfqpoint{0.829275in}{2.342135in}}%
\pgfpathlineto{\pgfqpoint{0.829324in}{2.343341in}}%
\pgfpathlineto{\pgfqpoint{0.829979in}{1.736241in}}%
\pgfpathlineto{\pgfqpoint{0.830127in}{1.599303in}}%
\pgfpathlineto{\pgfqpoint{0.830424in}{2.225252in}}%
\pgfpathlineto{\pgfqpoint{0.830474in}{2.254409in}}%
\pgfpathlineto{\pgfqpoint{0.830622in}{1.930812in}}%
\pgfpathlineto{\pgfqpoint{0.831002in}{1.532029in}}%
\pgfpathlineto{\pgfqpoint{0.831486in}{1.532046in}}%
\pgfpathlineto{\pgfqpoint{0.832934in}{1.532231in}}%
\pgfpathlineto{\pgfqpoint{0.833753in}{2.341152in}}%
\pgfpathlineto{\pgfqpoint{0.834148in}{2.340991in}}%
\pgfpathlineto{\pgfqpoint{0.834495in}{2.339486in}}%
\pgfpathlineto{\pgfqpoint{0.834544in}{2.340390in}}%
\pgfpathlineto{\pgfqpoint{0.835039in}{1.586706in}}%
\pgfpathlineto{\pgfqpoint{0.835385in}{2.262078in}}%
\pgfpathlineto{\pgfqpoint{0.835434in}{2.270180in}}%
\pgfpathlineto{\pgfqpoint{0.835484in}{2.192442in}}%
\pgfpathlineto{\pgfqpoint{0.835922in}{1.532031in}}%
\pgfpathlineto{\pgfqpoint{0.836466in}{1.532049in}}%
\pgfpathlineto{\pgfqpoint{0.837878in}{1.532188in}}%
\pgfpathlineto{\pgfqpoint{0.838696in}{2.338824in}}%
\pgfpathlineto{\pgfqpoint{0.839290in}{2.338359in}}%
\pgfpathlineto{\pgfqpoint{0.839586in}{2.337322in}}%
\pgfpathlineto{\pgfqpoint{0.839685in}{2.301933in}}%
\pgfpathlineto{\pgfqpoint{0.839982in}{1.563630in}}%
\pgfpathlineto{\pgfqpoint{0.840378in}{2.323942in}}%
\pgfpathlineto{\pgfqpoint{0.840724in}{1.564946in}}%
\pgfpathlineto{\pgfqpoint{0.840828in}{1.654076in}}%
\pgfpathlineto{\pgfqpoint{0.840845in}{1.532014in}}%
\pgfpathlineto{\pgfqpoint{0.841351in}{1.532066in}}%
\pgfpathlineto{\pgfqpoint{0.842811in}{1.532115in}}%
\pgfpathlineto{\pgfqpoint{0.842821in}{1.532157in}}%
\pgfpathlineto{\pgfqpoint{0.842848in}{1.598985in}}%
\pgfpathlineto{\pgfqpoint{0.843661in}{2.336780in}}%
\pgfpathlineto{\pgfqpoint{0.844205in}{2.334080in}}%
\pgfpathlineto{\pgfqpoint{0.844254in}{2.336500in}}%
\pgfpathlineto{\pgfqpoint{0.844967in}{1.595434in}}%
\pgfpathlineto{\pgfqpoint{0.845165in}{1.897658in}}%
\pgfpathlineto{\pgfqpoint{0.845313in}{2.253845in}}%
\pgfpathlineto{\pgfqpoint{0.845820in}{1.532037in}}%
\pgfpathlineto{\pgfqpoint{0.847763in}{1.532137in}}%
\pgfpathlineto{\pgfqpoint{0.847772in}{1.532290in}}%
\pgfpathlineto{\pgfqpoint{0.848552in}{2.334417in}}%
\pgfpathlineto{\pgfqpoint{0.848948in}{2.334271in}}%
\pgfpathlineto{\pgfqpoint{0.849492in}{2.331236in}}%
\pgfpathlineto{\pgfqpoint{0.849541in}{2.332614in}}%
\pgfpathlineto{\pgfqpoint{0.849890in}{1.580561in}}%
\pgfpathlineto{\pgfqpoint{0.850673in}{1.613481in}}%
\pgfpathlineto{\pgfqpoint{0.850720in}{1.694893in}}%
\pgfpathlineto{\pgfqpoint{0.850797in}{1.532034in}}%
\pgfpathlineto{\pgfqpoint{0.851338in}{1.532058in}}%
\pgfpathlineto{\pgfqpoint{0.852714in}{1.532184in}}%
\pgfpathlineto{\pgfqpoint{0.852756in}{1.913737in}}%
\pgfpathlineto{\pgfqpoint{0.853533in}{2.332336in}}%
\pgfpathlineto{\pgfqpoint{0.854176in}{2.330916in}}%
\pgfpathlineto{\pgfqpoint{0.854226in}{2.331925in}}%
\pgfpathlineto{\pgfqpoint{0.854720in}{1.791947in}}%
\pgfpathlineto{\pgfqpoint{0.854869in}{1.577729in}}%
\pgfpathlineto{\pgfqpoint{0.855215in}{2.282137in}}%
\pgfpathlineto{\pgfqpoint{0.855561in}{1.589256in}}%
\pgfpathlineto{\pgfqpoint{0.855733in}{1.532028in}}%
\pgfpathlineto{\pgfqpoint{0.855664in}{1.626727in}}%
\pgfpathlineto{\pgfqpoint{0.856764in}{1.532079in}}%
\pgfpathlineto{\pgfqpoint{0.857650in}{1.532119in}}%
\pgfpathlineto{\pgfqpoint{0.857661in}{1.532221in}}%
\pgfpathlineto{\pgfqpoint{0.858476in}{2.330121in}}%
\pgfpathlineto{\pgfqpoint{0.858921in}{2.329639in}}%
\pgfpathlineto{\pgfqpoint{0.858971in}{2.329889in}}%
\pgfpathlineto{\pgfqpoint{0.859218in}{2.328739in}}%
\pgfpathlineto{\pgfqpoint{0.859267in}{2.329420in}}%
\pgfpathlineto{\pgfqpoint{0.859612in}{1.923721in}}%
\pgfpathlineto{\pgfqpoint{0.859809in}{1.592332in}}%
\pgfpathlineto{\pgfqpoint{0.860156in}{2.261788in}}%
\pgfpathlineto{\pgfqpoint{0.860452in}{1.632876in}}%
\pgfpathlineto{\pgfqpoint{0.860627in}{1.531997in}}%
\pgfpathlineto{\pgfqpoint{0.860610in}{1.667466in}}%
\pgfpathlineto{\pgfqpoint{0.861331in}{1.532075in}}%
\pgfpathlineto{\pgfqpoint{0.862588in}{1.532105in}}%
\pgfpathlineto{\pgfqpoint{0.862601in}{1.532144in}}%
\pgfpathlineto{\pgfqpoint{0.862620in}{1.554712in}}%
\pgfpathlineto{\pgfqpoint{0.863458in}{2.327993in}}%
\pgfpathlineto{\pgfqpoint{0.863507in}{2.327798in}}%
\pgfpathlineto{\pgfqpoint{0.863557in}{2.327984in}}%
\pgfpathlineto{\pgfqpoint{0.864002in}{2.327246in}}%
\pgfpathlineto{\pgfqpoint{0.864051in}{2.327625in}}%
\pgfpathlineto{\pgfqpoint{0.864200in}{2.326286in}}%
\pgfpathlineto{\pgfqpoint{0.864249in}{2.327268in}}%
\pgfpathlineto{\pgfqpoint{0.864778in}{1.568813in}}%
\pgfpathlineto{\pgfqpoint{0.865074in}{2.254835in}}%
\pgfpathlineto{\pgfqpoint{0.865124in}{2.309771in}}%
\pgfpathlineto{\pgfqpoint{0.865322in}{1.803438in}}%
\pgfpathlineto{\pgfqpoint{0.865571in}{1.531986in}}%
\pgfpathlineto{\pgfqpoint{0.866195in}{1.532076in}}%
\pgfpathlineto{\pgfqpoint{0.867431in}{1.532082in}}%
\pgfpathlineto{\pgfqpoint{0.867554in}{1.532285in}}%
\pgfpathlineto{\pgfqpoint{0.868415in}{2.325850in}}%
\pgfpathlineto{\pgfqpoint{0.868761in}{2.325563in}}%
\pgfpathlineto{\pgfqpoint{0.868811in}{2.325666in}}%
\pgfpathlineto{\pgfqpoint{0.869157in}{2.324665in}}%
\pgfpathlineto{\pgfqpoint{0.869206in}{2.324988in}}%
\pgfpathlineto{\pgfqpoint{0.869256in}{2.323607in}}%
\pgfpathlineto{\pgfqpoint{0.869305in}{2.324608in}}%
\pgfpathlineto{\pgfqpoint{0.869694in}{1.611019in}}%
\pgfpathlineto{\pgfqpoint{0.870466in}{1.629678in}}%
\pgfpathlineto{\pgfqpoint{0.870501in}{1.676178in}}%
\pgfpathlineto{\pgfqpoint{0.870572in}{1.532034in}}%
\pgfpathlineto{\pgfqpoint{0.871108in}{1.532056in}}%
\pgfpathlineto{\pgfqpoint{0.872498in}{1.532223in}}%
\pgfpathlineto{\pgfqpoint{0.873372in}{2.323752in}}%
\pgfpathlineto{\pgfqpoint{0.873767in}{2.323549in}}%
\pgfpathlineto{\pgfqpoint{0.874213in}{2.322133in}}%
\pgfpathlineto{\pgfqpoint{0.874262in}{2.322294in}}%
\pgfpathlineto{\pgfqpoint{0.875526in}{1.532027in}}%
\pgfpathlineto{\pgfqpoint{0.876014in}{1.532045in}}%
\pgfpathlineto{\pgfqpoint{0.877443in}{1.532215in}}%
\pgfpathlineto{\pgfqpoint{0.878353in}{2.321747in}}%
\pgfpathlineto{\pgfqpoint{0.878700in}{2.321322in}}%
\pgfpathlineto{\pgfqpoint{0.878749in}{2.321552in}}%
\pgfpathlineto{\pgfqpoint{0.879095in}{2.319670in}}%
\pgfpathlineto{\pgfqpoint{0.879145in}{2.320889in}}%
\pgfpathlineto{\pgfqpoint{0.879634in}{1.575088in}}%
\pgfpathlineto{\pgfqpoint{0.879930in}{2.247368in}}%
\pgfpathlineto{\pgfqpoint{0.879980in}{2.297231in}}%
\pgfpathlineto{\pgfqpoint{0.880178in}{1.801164in}}%
\pgfpathlineto{\pgfqpoint{0.880409in}{1.532011in}}%
\pgfpathlineto{\pgfqpoint{0.881032in}{1.532071in}}%
\pgfpathlineto{\pgfqpoint{0.882365in}{1.532089in}}%
\pgfpathlineto{\pgfqpoint{0.882383in}{1.532143in}}%
\pgfpathlineto{\pgfqpoint{0.882399in}{1.547589in}}%
\pgfpathlineto{\pgfqpoint{0.883279in}{2.319702in}}%
\pgfpathlineto{\pgfqpoint{0.883328in}{2.319549in}}%
\pgfpathlineto{\pgfqpoint{0.883477in}{2.319652in}}%
\pgfpathlineto{\pgfqpoint{0.883823in}{2.319001in}}%
\pgfpathlineto{\pgfqpoint{0.883872in}{2.319275in}}%
\pgfpathlineto{\pgfqpoint{0.884021in}{2.318159in}}%
\pgfpathlineto{\pgfqpoint{0.884070in}{2.318864in}}%
\pgfpathlineto{\pgfqpoint{0.884570in}{1.584315in}}%
\pgfpathlineto{\pgfqpoint{0.884867in}{2.230743in}}%
\pgfpathlineto{\pgfqpoint{0.884917in}{2.279101in}}%
\pgfpathlineto{\pgfqpoint{0.885114in}{1.807278in}}%
\pgfpathlineto{\pgfqpoint{0.885405in}{1.532029in}}%
\pgfpathlineto{\pgfqpoint{0.885990in}{1.532050in}}%
\pgfpathlineto{\pgfqpoint{0.887334in}{1.532207in}}%
\pgfpathlineto{\pgfqpoint{0.888242in}{2.317690in}}%
\pgfpathlineto{\pgfqpoint{0.888589in}{2.317314in}}%
\pgfpathlineto{\pgfqpoint{0.888638in}{2.317498in}}%
\pgfpathlineto{\pgfqpoint{0.888984in}{2.316085in}}%
\pgfpathlineto{\pgfqpoint{0.889034in}{2.316836in}}%
\pgfpathlineto{\pgfqpoint{0.889528in}{1.596713in}}%
\pgfpathlineto{\pgfqpoint{0.889825in}{2.214183in}}%
\pgfpathlineto{\pgfqpoint{0.889874in}{2.257141in}}%
\pgfpathlineto{\pgfqpoint{0.890072in}{1.808729in}}%
\pgfpathlineto{\pgfqpoint{0.890349in}{1.532027in}}%
\pgfpathlineto{\pgfqpoint{0.890932in}{1.532048in}}%
\pgfpathlineto{\pgfqpoint{0.892279in}{1.532213in}}%
\pgfpathlineto{\pgfqpoint{0.893193in}{2.315715in}}%
\pgfpathlineto{\pgfqpoint{0.893539in}{2.315269in}}%
\pgfpathlineto{\pgfqpoint{0.893588in}{2.315530in}}%
\pgfpathlineto{\pgfqpoint{0.893935in}{2.313376in}}%
\pgfpathlineto{\pgfqpoint{0.893984in}{2.314902in}}%
\pgfpathlineto{\pgfqpoint{0.894470in}{1.577171in}}%
\pgfpathlineto{\pgfqpoint{0.894767in}{2.220617in}}%
\pgfpathlineto{\pgfqpoint{0.894817in}{2.288884in}}%
\pgfpathlineto{\pgfqpoint{0.895064in}{1.702069in}}%
\pgfpathlineto{\pgfqpoint{0.895259in}{1.532027in}}%
\pgfpathlineto{\pgfqpoint{0.895924in}{1.532062in}}%
\pgfpathlineto{\pgfqpoint{0.897213in}{1.532116in}}%
\pgfpathlineto{\pgfqpoint{0.897224in}{1.532196in}}%
\pgfpathlineto{\pgfqpoint{0.898124in}{2.313715in}}%
\pgfpathlineto{\pgfqpoint{0.898519in}{2.313529in}}%
\pgfpathlineto{\pgfqpoint{0.898865in}{2.312353in}}%
\pgfpathlineto{\pgfqpoint{0.898915in}{2.312873in}}%
\pgfpathlineto{\pgfqpoint{0.899162in}{2.074048in}}%
\pgfpathlineto{\pgfqpoint{0.899409in}{1.617435in}}%
\pgfpathlineto{\pgfqpoint{0.899756in}{2.220896in}}%
\pgfpathlineto{\pgfqpoint{0.900052in}{1.656472in}}%
\pgfpathlineto{\pgfqpoint{0.900191in}{1.532012in}}%
\pgfpathlineto{\pgfqpoint{0.900994in}{1.532075in}}%
\pgfpathlineto{\pgfqpoint{0.902157in}{1.532112in}}%
\pgfpathlineto{\pgfqpoint{0.902166in}{1.532153in}}%
\pgfpathlineto{\pgfqpoint{0.902184in}{1.554608in}}%
\pgfpathlineto{\pgfqpoint{0.903091in}{2.311884in}}%
\pgfpathlineto{\pgfqpoint{0.903635in}{2.310058in}}%
\pgfpathlineto{\pgfqpoint{0.903684in}{2.311559in}}%
\pgfpathlineto{\pgfqpoint{0.904393in}{1.594441in}}%
\pgfpathlineto{\pgfqpoint{0.904591in}{1.957746in}}%
\pgfpathlineto{\pgfqpoint{0.904739in}{2.261701in}}%
\pgfpathlineto{\pgfqpoint{0.905178in}{1.532027in}}%
\pgfpathlineto{\pgfqpoint{0.905218in}{1.532040in}}%
\pgfpathlineto{\pgfqpoint{0.907104in}{1.532117in}}%
\pgfpathlineto{\pgfqpoint{0.907116in}{1.532213in}}%
\pgfpathlineto{\pgfqpoint{0.908033in}{2.309874in}}%
\pgfpathlineto{\pgfqpoint{0.908379in}{2.309396in}}%
\pgfpathlineto{\pgfqpoint{0.908428in}{2.309697in}}%
\pgfpathlineto{\pgfqpoint{0.908676in}{2.308586in}}%
\pgfpathlineto{\pgfqpoint{0.908725in}{2.309311in}}%
\pgfpathlineto{\pgfqpoint{0.908774in}{2.305608in}}%
\pgfpathlineto{\pgfqpoint{0.908824in}{2.309094in}}%
\pgfpathlineto{\pgfqpoint{0.909297in}{1.576872in}}%
\pgfpathlineto{\pgfqpoint{0.909594in}{2.178001in}}%
\pgfpathlineto{\pgfqpoint{0.909693in}{2.290222in}}%
\pgfpathlineto{\pgfqpoint{0.910025in}{1.579627in}}%
\pgfpathlineto{\pgfqpoint{0.910069in}{1.580468in}}%
\pgfpathlineto{\pgfqpoint{0.910080in}{1.531991in}}%
\pgfpathlineto{\pgfqpoint{0.910835in}{1.532043in}}%
\pgfpathlineto{\pgfqpoint{0.912066in}{1.532624in}}%
\pgfpathlineto{\pgfqpoint{0.912943in}{2.307954in}}%
\pgfpathlineto{\pgfqpoint{0.913141in}{2.307922in}}%
\pgfpathlineto{\pgfqpoint{0.913685in}{2.306468in}}%
\pgfpathlineto{\pgfqpoint{0.913734in}{2.307276in}}%
\pgfpathlineto{\pgfqpoint{0.914267in}{1.593774in}}%
\pgfpathlineto{\pgfqpoint{0.914563in}{2.180396in}}%
\pgfpathlineto{\pgfqpoint{0.914613in}{2.255209in}}%
\pgfpathlineto{\pgfqpoint{0.914910in}{1.644969in}}%
\pgfpathlineto{\pgfqpoint{0.915025in}{1.531995in}}%
\pgfpathlineto{\pgfqpoint{0.915802in}{1.532072in}}%
\pgfpathlineto{\pgfqpoint{0.916994in}{1.532113in}}%
\pgfpathlineto{\pgfqpoint{0.917004in}{1.532161in}}%
\pgfpathlineto{\pgfqpoint{0.917044in}{1.824372in}}%
\pgfpathlineto{\pgfqpoint{0.917819in}{2.305991in}}%
\pgfpathlineto{\pgfqpoint{0.918709in}{2.305008in}}%
\pgfpathlineto{\pgfqpoint{0.918808in}{2.304471in}}%
\pgfpathlineto{\pgfqpoint{0.919056in}{1.792077in}}%
\pgfpathlineto{\pgfqpoint{0.919204in}{1.588309in}}%
\pgfpathlineto{\pgfqpoint{0.919550in}{2.265521in}}%
\pgfpathlineto{\pgfqpoint{0.919896in}{1.598196in}}%
\pgfpathlineto{\pgfqpoint{0.920012in}{1.532026in}}%
\pgfpathlineto{\pgfqpoint{0.919957in}{1.614673in}}%
\pgfpathlineto{\pgfqpoint{0.920984in}{1.532061in}}%
\pgfpathlineto{\pgfqpoint{0.921950in}{1.532177in}}%
\pgfpathlineto{\pgfqpoint{0.921991in}{1.866609in}}%
\pgfpathlineto{\pgfqpoint{0.922776in}{2.304261in}}%
\pgfpathlineto{\pgfqpoint{0.923419in}{2.302756in}}%
\pgfpathlineto{\pgfqpoint{0.923469in}{2.303968in}}%
\pgfpathlineto{\pgfqpoint{0.924210in}{1.607615in}}%
\pgfpathlineto{\pgfqpoint{0.924458in}{2.118216in}}%
\pgfpathlineto{\pgfqpoint{0.924557in}{2.240544in}}%
\pgfpathlineto{\pgfqpoint{0.924917in}{1.531989in}}%
\pgfpathlineto{\pgfqpoint{0.924930in}{1.532015in}}%
\pgfpathlineto{\pgfqpoint{0.925700in}{1.532073in}}%
\pgfpathlineto{\pgfqpoint{0.925750in}{1.532047in}}%
\pgfpathlineto{\pgfqpoint{0.926902in}{1.532702in}}%
\pgfpathlineto{\pgfqpoint{0.927807in}{2.302402in}}%
\pgfpathlineto{\pgfqpoint{0.927956in}{2.301974in}}%
\pgfpathlineto{\pgfqpoint{0.928005in}{2.302371in}}%
\pgfpathlineto{\pgfqpoint{0.928351in}{2.301138in}}%
\pgfpathlineto{\pgfqpoint{0.928401in}{2.302080in}}%
\pgfpathlineto{\pgfqpoint{0.928450in}{2.299550in}}%
\pgfpathlineto{\pgfqpoint{0.928500in}{2.301946in}}%
\pgfpathlineto{\pgfqpoint{0.929127in}{1.583394in}}%
\pgfpathlineto{\pgfqpoint{0.929325in}{1.962610in}}%
\pgfpathlineto{\pgfqpoint{0.929473in}{2.282058in}}%
\pgfpathlineto{\pgfqpoint{0.929861in}{1.532005in}}%
\pgfpathlineto{\pgfqpoint{0.929961in}{1.532048in}}%
\pgfpathlineto{\pgfqpoint{0.931494in}{1.532078in}}%
\pgfpathlineto{\pgfqpoint{0.931848in}{1.532775in}}%
\pgfpathlineto{\pgfqpoint{0.932753in}{2.300584in}}%
\pgfpathlineto{\pgfqpoint{0.932902in}{2.300156in}}%
\pgfpathlineto{\pgfqpoint{0.932951in}{2.300562in}}%
\pgfpathlineto{\pgfqpoint{0.933396in}{2.298292in}}%
\pgfpathlineto{\pgfqpoint{0.933446in}{2.300170in}}%
\pgfpathlineto{\pgfqpoint{0.934075in}{1.612452in}}%
\pgfpathlineto{\pgfqpoint{0.934323in}{2.054153in}}%
\pgfpathlineto{\pgfqpoint{0.934422in}{2.227967in}}%
\pgfpathlineto{\pgfqpoint{0.934807in}{1.531995in}}%
\pgfpathlineto{\pgfqpoint{0.934849in}{1.532020in}}%
\pgfpathlineto{\pgfqpoint{0.936366in}{1.532100in}}%
\pgfpathlineto{\pgfqpoint{0.936776in}{1.532112in}}%
\pgfpathlineto{\pgfqpoint{0.936786in}{1.532165in}}%
\pgfpathlineto{\pgfqpoint{0.936828in}{1.873721in}}%
\pgfpathlineto{\pgfqpoint{0.937602in}{2.298738in}}%
\pgfpathlineto{\pgfqpoint{0.938344in}{2.297279in}}%
\pgfpathlineto{\pgfqpoint{0.938394in}{2.298295in}}%
\pgfpathlineto{\pgfqpoint{0.939037in}{1.585838in}}%
\pgfpathlineto{\pgfqpoint{0.939284in}{2.108252in}}%
\pgfpathlineto{\pgfqpoint{0.939383in}{2.276417in}}%
\pgfpathlineto{\pgfqpoint{0.939753in}{1.531999in}}%
\pgfpathlineto{\pgfqpoint{0.939791in}{1.532049in}}%
\pgfpathlineto{\pgfqpoint{0.940910in}{1.532058in}}%
\pgfpathlineto{\pgfqpoint{0.941738in}{1.532517in}}%
\pgfpathlineto{\pgfqpoint{0.942672in}{2.296894in}}%
\pgfpathlineto{\pgfqpoint{0.942820in}{2.296704in}}%
\pgfpathlineto{\pgfqpoint{0.942870in}{2.296843in}}%
\pgfpathlineto{\pgfqpoint{0.943315in}{2.295953in}}%
\pgfpathlineto{\pgfqpoint{0.943364in}{2.296307in}}%
\pgfpathlineto{\pgfqpoint{0.943414in}{2.295467in}}%
\pgfpathlineto{\pgfqpoint{0.943463in}{2.296080in}}%
\pgfpathlineto{\pgfqpoint{0.943809in}{1.809392in}}%
\pgfpathlineto{\pgfqpoint{0.943958in}{1.621059in}}%
\pgfpathlineto{\pgfqpoint{0.944304in}{2.214569in}}%
\pgfpathlineto{\pgfqpoint{0.944639in}{1.632940in}}%
\pgfpathlineto{\pgfqpoint{0.944719in}{1.532026in}}%
\pgfpathlineto{\pgfqpoint{0.944685in}{1.642492in}}%
\pgfpathlineto{\pgfqpoint{0.945954in}{1.532071in}}%
\pgfpathlineto{\pgfqpoint{0.946680in}{1.532215in}}%
\pgfpathlineto{\pgfqpoint{0.947608in}{2.295196in}}%
\pgfpathlineto{\pgfqpoint{0.947905in}{2.295126in}}%
\pgfpathlineto{\pgfqpoint{0.948251in}{2.293718in}}%
\pgfpathlineto{\pgfqpoint{0.948301in}{2.294749in}}%
\pgfpathlineto{\pgfqpoint{0.948944in}{1.585921in}}%
\pgfpathlineto{\pgfqpoint{0.949191in}{2.099517in}}%
\pgfpathlineto{\pgfqpoint{0.949290in}{2.275419in}}%
\pgfpathlineto{\pgfqpoint{0.949646in}{1.532024in}}%
\pgfpathlineto{\pgfqpoint{0.949702in}{1.532038in}}%
\pgfpathlineto{\pgfqpoint{0.951613in}{1.532114in}}%
\pgfpathlineto{\pgfqpoint{0.951623in}{1.532176in}}%
\pgfpathlineto{\pgfqpoint{0.951665in}{1.929914in}}%
\pgfpathlineto{\pgfqpoint{0.952446in}{2.293478in}}%
\pgfpathlineto{\pgfqpoint{0.953089in}{2.290993in}}%
\pgfpathlineto{\pgfqpoint{0.952545in}{2.293491in}}%
\pgfpathlineto{\pgfqpoint{0.953138in}{2.293218in}}%
\pgfpathlineto{\pgfqpoint{0.953879in}{1.581885in}}%
\pgfpathlineto{\pgfqpoint{0.954077in}{1.929180in}}%
\pgfpathlineto{\pgfqpoint{0.954225in}{2.281922in}}%
\pgfpathlineto{\pgfqpoint{0.954624in}{1.532027in}}%
\pgfpathlineto{\pgfqpoint{0.954699in}{1.532029in}}%
\pgfpathlineto{\pgfqpoint{0.956569in}{1.532190in}}%
\pgfpathlineto{\pgfqpoint{0.957505in}{2.291607in}}%
\pgfpathlineto{\pgfqpoint{0.957950in}{2.291231in}}%
\pgfpathlineto{\pgfqpoint{0.958000in}{2.291317in}}%
\pgfpathlineto{\pgfqpoint{0.958247in}{2.290538in}}%
\pgfpathlineto{\pgfqpoint{0.958297in}{2.290757in}}%
\pgfpathlineto{\pgfqpoint{0.958346in}{2.289842in}}%
\pgfpathlineto{\pgfqpoint{0.958395in}{2.290380in}}%
\pgfpathlineto{\pgfqpoint{0.959535in}{1.531991in}}%
\pgfpathlineto{\pgfqpoint{0.959817in}{1.532053in}}%
\pgfpathlineto{\pgfqpoint{0.961446in}{1.532089in}}%
\pgfpathlineto{\pgfqpoint{0.961512in}{1.532154in}}%
\pgfpathlineto{\pgfqpoint{0.961539in}{1.603002in}}%
\pgfpathlineto{\pgfqpoint{0.962255in}{2.289948in}}%
\pgfpathlineto{\pgfqpoint{0.962304in}{2.289754in}}%
\pgfpathlineto{\pgfqpoint{0.962552in}{2.289978in}}%
\pgfpathlineto{\pgfqpoint{0.962898in}{2.289354in}}%
\pgfpathlineto{\pgfqpoint{0.962947in}{2.289772in}}%
\pgfpathlineto{\pgfqpoint{0.963195in}{2.287613in}}%
\pgfpathlineto{\pgfqpoint{0.963244in}{2.289362in}}%
\pgfpathlineto{\pgfqpoint{0.963795in}{1.581751in}}%
\pgfpathlineto{\pgfqpoint{0.964042in}{2.113274in}}%
\pgfpathlineto{\pgfqpoint{0.964141in}{2.284103in}}%
\pgfpathlineto{\pgfqpoint{0.964514in}{1.532030in}}%
\pgfpathlineto{\pgfqpoint{0.964544in}{1.532033in}}%
\pgfpathlineto{\pgfqpoint{0.966460in}{1.532189in}}%
\pgfpathlineto{\pgfqpoint{0.967323in}{2.288232in}}%
\pgfpathlineto{\pgfqpoint{0.967818in}{2.288034in}}%
\pgfpathlineto{\pgfqpoint{0.968164in}{2.286874in}}%
\pgfpathlineto{\pgfqpoint{0.968214in}{2.287365in}}%
\pgfpathlineto{\pgfqpoint{0.968452in}{2.067206in}}%
\pgfpathlineto{\pgfqpoint{0.968677in}{1.633610in}}%
\pgfpathlineto{\pgfqpoint{0.969072in}{2.200080in}}%
\pgfpathlineto{\pgfqpoint{0.969362in}{1.649081in}}%
\pgfpathlineto{\pgfqpoint{0.969487in}{1.532032in}}%
\pgfpathlineto{\pgfqpoint{0.969412in}{1.649970in}}%
\pgfpathlineto{\pgfqpoint{0.970420in}{1.532068in}}%
\pgfpathlineto{\pgfqpoint{0.971411in}{1.532470in}}%
\pgfpathlineto{\pgfqpoint{0.972386in}{2.286609in}}%
\pgfpathlineto{\pgfqpoint{0.972485in}{2.286595in}}%
\pgfpathlineto{\pgfqpoint{0.973128in}{2.284959in}}%
\pgfpathlineto{\pgfqpoint{0.973177in}{2.285865in}}%
\pgfpathlineto{\pgfqpoint{0.973662in}{1.594564in}}%
\pgfpathlineto{\pgfqpoint{0.973959in}{2.152885in}}%
\pgfpathlineto{\pgfqpoint{0.974058in}{2.260982in}}%
\pgfpathlineto{\pgfqpoint{0.974372in}{1.531971in}}%
\pgfpathlineto{\pgfqpoint{0.974395in}{1.532006in}}%
\pgfpathlineto{\pgfqpoint{0.975196in}{1.532081in}}%
\pgfpathlineto{\pgfqpoint{0.975245in}{1.532040in}}%
\pgfpathlineto{\pgfqpoint{0.976354in}{1.532257in}}%
\pgfpathlineto{\pgfqpoint{0.977287in}{2.284914in}}%
\pgfpathlineto{\pgfqpoint{0.977534in}{2.284742in}}%
\pgfpathlineto{\pgfqpoint{0.977584in}{2.284813in}}%
\pgfpathlineto{\pgfqpoint{0.977930in}{2.284203in}}%
\pgfpathlineto{\pgfqpoint{0.977979in}{2.284335in}}%
\pgfpathlineto{\pgfqpoint{0.978128in}{2.283309in}}%
\pgfpathlineto{\pgfqpoint{0.978177in}{2.283753in}}%
\pgfpathlineto{\pgfqpoint{0.979368in}{1.532024in}}%
\pgfpathlineto{\pgfqpoint{0.979856in}{1.532041in}}%
\pgfpathlineto{\pgfqpoint{0.981298in}{1.532216in}}%
\pgfpathlineto{\pgfqpoint{0.982330in}{2.283391in}}%
\pgfpathlineto{\pgfqpoint{0.982479in}{2.282899in}}%
\pgfpathlineto{\pgfqpoint{0.982528in}{2.283342in}}%
\pgfpathlineto{\pgfqpoint{0.982874in}{2.281664in}}%
\pgfpathlineto{\pgfqpoint{0.982924in}{2.283016in}}%
\pgfpathlineto{\pgfqpoint{0.983616in}{1.605201in}}%
\pgfpathlineto{\pgfqpoint{0.983814in}{1.995672in}}%
\pgfpathlineto{\pgfqpoint{0.983962in}{2.245094in}}%
\pgfpathlineto{\pgfqpoint{0.984298in}{1.532026in}}%
\pgfpathlineto{\pgfqpoint{0.984374in}{1.532028in}}%
\pgfpathlineto{\pgfqpoint{0.986242in}{1.532189in}}%
\pgfpathlineto{\pgfqpoint{0.987192in}{2.281672in}}%
\pgfpathlineto{\pgfqpoint{0.987587in}{2.281496in}}%
\pgfpathlineto{\pgfqpoint{0.988032in}{2.279646in}}%
\pgfpathlineto{\pgfqpoint{0.988082in}{2.280541in}}%
\pgfpathlineto{\pgfqpoint{0.989253in}{1.532026in}}%
\pgfpathlineto{\pgfqpoint{0.989389in}{1.532042in}}%
\pgfpathlineto{\pgfqpoint{0.991177in}{1.532115in}}%
\pgfpathlineto{\pgfqpoint{0.991188in}{1.532209in}}%
\pgfpathlineto{\pgfqpoint{0.992217in}{2.280171in}}%
\pgfpathlineto{\pgfqpoint{0.992415in}{2.280125in}}%
\pgfpathlineto{\pgfqpoint{0.992860in}{2.269632in}}%
\pgfpathlineto{\pgfqpoint{0.992910in}{2.279650in}}%
\pgfpathlineto{\pgfqpoint{0.993503in}{1.632542in}}%
\pgfpathlineto{\pgfqpoint{0.993750in}{2.083089in}}%
\pgfpathlineto{\pgfqpoint{0.993849in}{2.201000in}}%
\pgfpathlineto{\pgfqpoint{0.994156in}{1.532012in}}%
\pgfpathlineto{\pgfqpoint{0.994192in}{1.532045in}}%
\pgfpathlineto{\pgfqpoint{0.995409in}{1.532063in}}%
\pgfpathlineto{\pgfqpoint{0.996131in}{1.532154in}}%
\pgfpathlineto{\pgfqpoint{0.996159in}{1.608686in}}%
\pgfpathlineto{\pgfqpoint{0.996925in}{2.278479in}}%
\pgfpathlineto{\pgfqpoint{0.997964in}{2.276662in}}%
\pgfpathlineto{\pgfqpoint{0.998013in}{2.277132in}}%
\pgfpathlineto{\pgfqpoint{0.999121in}{1.532024in}}%
\pgfpathlineto{\pgfqpoint{0.999420in}{1.532047in}}%
\pgfpathlineto{\pgfqpoint{1.001068in}{1.532115in}}%
\pgfpathlineto{\pgfqpoint{1.001080in}{1.532213in}}%
\pgfpathlineto{\pgfqpoint{1.002113in}{2.277058in}}%
\pgfpathlineto{\pgfqpoint{1.002310in}{2.277015in}}%
\pgfpathlineto{\pgfqpoint{1.002657in}{2.275415in}}%
\pgfpathlineto{\pgfqpoint{1.002706in}{2.276710in}}%
\pgfpathlineto{\pgfqpoint{1.003398in}{1.617291in}}%
\pgfpathlineto{\pgfqpoint{1.003646in}{2.079074in}}%
\pgfpathlineto{\pgfqpoint{1.003745in}{2.223139in}}%
\pgfpathlineto{\pgfqpoint{1.004048in}{1.532015in}}%
\pgfpathlineto{\pgfqpoint{1.004123in}{1.532019in}}%
\pgfpathlineto{\pgfqpoint{1.005749in}{1.532103in}}%
\pgfpathlineto{\pgfqpoint{1.006012in}{1.532110in}}%
\pgfpathlineto{\pgfqpoint{1.006022in}{1.532161in}}%
\pgfpathlineto{\pgfqpoint{1.006064in}{1.854875in}}%
\pgfpathlineto{\pgfqpoint{1.006838in}{2.275476in}}%
\pgfpathlineto{\pgfqpoint{1.007679in}{2.272809in}}%
\pgfpathlineto{\pgfqpoint{1.007729in}{2.274996in}}%
\pgfpathlineto{\pgfqpoint{1.008322in}{1.637136in}}%
\pgfpathlineto{\pgfqpoint{1.008569in}{2.058190in}}%
\pgfpathlineto{\pgfqpoint{1.008668in}{2.192937in}}%
\pgfpathlineto{\pgfqpoint{1.008993in}{1.532018in}}%
\pgfpathlineto{\pgfqpoint{1.009065in}{1.532021in}}%
\pgfpathlineto{\pgfqpoint{1.010789in}{1.532106in}}%
\pgfpathlineto{\pgfqpoint{1.010958in}{1.532112in}}%
\pgfpathlineto{\pgfqpoint{1.010969in}{1.532181in}}%
\pgfpathlineto{\pgfqpoint{1.011950in}{2.273929in}}%
\pgfpathlineto{\pgfqpoint{1.012346in}{2.273778in}}%
\pgfpathlineto{\pgfqpoint{1.012692in}{2.272551in}}%
\pgfpathlineto{\pgfqpoint{1.012741in}{2.273240in}}%
\pgfpathlineto{\pgfqpoint{1.013071in}{1.964900in}}%
\pgfpathlineto{\pgfqpoint{1.013939in}{1.532013in}}%
\pgfpathlineto{\pgfqpoint{1.013615in}{2.211669in}}%
\pgfpathlineto{\pgfqpoint{1.013977in}{1.532044in}}%
\pgfpathlineto{\pgfqpoint{1.015294in}{1.532067in}}%
\pgfpathlineto{\pgfqpoint{1.015913in}{1.532156in}}%
\pgfpathlineto{\pgfqpoint{1.015953in}{1.822716in}}%
\pgfpathlineto{\pgfqpoint{1.016703in}{2.272365in}}%
\pgfpathlineto{\pgfqpoint{1.017741in}{2.270606in}}%
\pgfpathlineto{\pgfqpoint{1.017791in}{2.271252in}}%
\pgfpathlineto{\pgfqpoint{1.018902in}{1.532021in}}%
\pgfpathlineto{\pgfqpoint{1.019245in}{1.532033in}}%
\pgfpathlineto{\pgfqpoint{1.020862in}{1.532242in}}%
\pgfpathlineto{\pgfqpoint{1.021797in}{2.270960in}}%
\pgfpathlineto{\pgfqpoint{1.021995in}{2.270941in}}%
\pgfpathlineto{\pgfqpoint{1.022637in}{2.268325in}}%
\pgfpathlineto{\pgfqpoint{1.022687in}{2.270197in}}%
\pgfpathlineto{\pgfqpoint{1.023181in}{1.600443in}}%
\pgfpathlineto{\pgfqpoint{1.023478in}{2.193060in}}%
\pgfpathlineto{\pgfqpoint{1.023528in}{2.248419in}}%
\pgfpathlineto{\pgfqpoint{1.023765in}{1.733848in}}%
\pgfpathlineto{\pgfqpoint{1.023849in}{1.532019in}}%
\pgfpathlineto{\pgfqpoint{1.024641in}{1.532068in}}%
\pgfpathlineto{\pgfqpoint{1.025795in}{1.532112in}}%
\pgfpathlineto{\pgfqpoint{1.025806in}{1.532186in}}%
\pgfpathlineto{\pgfqpoint{1.026814in}{2.269487in}}%
\pgfpathlineto{\pgfqpoint{1.027110in}{2.269394in}}%
\pgfpathlineto{\pgfqpoint{1.027555in}{2.266574in}}%
\pgfpathlineto{\pgfqpoint{1.027605in}{2.268784in}}%
\pgfpathlineto{\pgfqpoint{1.028134in}{1.595988in}}%
\pgfpathlineto{\pgfqpoint{1.028382in}{2.128209in}}%
\pgfpathlineto{\pgfqpoint{1.028480in}{2.260302in}}%
\pgfpathlineto{\pgfqpoint{1.028815in}{1.532024in}}%
\pgfpathlineto{\pgfqpoint{1.028856in}{1.532037in}}%
\pgfpathlineto{\pgfqpoint{1.030741in}{1.532113in}}%
\pgfpathlineto{\pgfqpoint{1.030750in}{1.532172in}}%
\pgfpathlineto{\pgfqpoint{1.030792in}{1.893794in}}%
\pgfpathlineto{\pgfqpoint{1.031570in}{2.268068in}}%
\pgfpathlineto{\pgfqpoint{1.032312in}{2.265558in}}%
\pgfpathlineto{\pgfqpoint{1.031669in}{2.268092in}}%
\pgfpathlineto{\pgfqpoint{1.032362in}{2.267815in}}%
\pgfpathlineto{\pgfqpoint{1.033103in}{1.599051in}}%
\pgfpathlineto{\pgfqpoint{1.033301in}{1.964506in}}%
\pgfpathlineto{\pgfqpoint{1.033450in}{2.254285in}}%
\pgfpathlineto{\pgfqpoint{1.033719in}{1.531989in}}%
\pgfpathlineto{\pgfqpoint{1.033893in}{1.532045in}}%
\pgfpathlineto{\pgfqpoint{1.035685in}{1.532109in}}%
\pgfpathlineto{\pgfqpoint{1.035695in}{1.532165in}}%
\pgfpathlineto{\pgfqpoint{1.035738in}{1.908961in}}%
\pgfpathlineto{\pgfqpoint{1.036511in}{2.266603in}}%
\pgfpathlineto{\pgfqpoint{1.037253in}{2.264808in}}%
\pgfpathlineto{\pgfqpoint{1.036610in}{2.266624in}}%
\pgfpathlineto{\pgfqpoint{1.037303in}{2.266330in}}%
\pgfpathlineto{\pgfqpoint{1.038044in}{1.616959in}}%
\pgfpathlineto{\pgfqpoint{1.038242in}{1.986686in}}%
\pgfpathlineto{\pgfqpoint{1.038391in}{2.225746in}}%
\pgfpathlineto{\pgfqpoint{1.038664in}{1.531992in}}%
\pgfpathlineto{\pgfqpoint{1.038837in}{1.532045in}}%
\pgfpathlineto{\pgfqpoint{1.040631in}{1.532111in}}%
\pgfpathlineto{\pgfqpoint{1.040641in}{1.532171in}}%
\pgfpathlineto{\pgfqpoint{1.041648in}{2.265079in}}%
\pgfpathlineto{\pgfqpoint{1.042094in}{2.264707in}}%
\pgfpathlineto{\pgfqpoint{1.042143in}{2.264803in}}%
\pgfpathlineto{\pgfqpoint{1.042390in}{2.264046in}}%
\pgfpathlineto{\pgfqpoint{1.042440in}{2.264280in}}%
\pgfpathlineto{\pgfqpoint{1.042489in}{2.263424in}}%
\pgfpathlineto{\pgfqpoint{1.042539in}{2.263942in}}%
\pgfpathlineto{\pgfqpoint{1.043609in}{1.531993in}}%
\pgfpathlineto{\pgfqpoint{1.044127in}{1.532036in}}%
\pgfpathlineto{\pgfqpoint{1.045587in}{1.532176in}}%
\pgfpathlineto{\pgfqpoint{1.046513in}{2.263691in}}%
\pgfpathlineto{\pgfqpoint{1.047007in}{2.263511in}}%
\pgfpathlineto{\pgfqpoint{1.047453in}{2.261653in}}%
\pgfpathlineto{\pgfqpoint{1.047502in}{2.262559in}}%
\pgfpathlineto{\pgfqpoint{1.048555in}{1.531993in}}%
\pgfpathlineto{\pgfqpoint{1.048818in}{1.532048in}}%
\pgfpathlineto{\pgfqpoint{1.050522in}{1.532111in}}%
\pgfpathlineto{\pgfqpoint{1.050533in}{1.532181in}}%
\pgfpathlineto{\pgfqpoint{1.051475in}{2.262328in}}%
\pgfpathlineto{\pgfqpoint{1.051920in}{2.261993in}}%
\pgfpathlineto{\pgfqpoint{1.051970in}{2.262144in}}%
\pgfpathlineto{\pgfqpoint{1.052316in}{2.261038in}}%
\pgfpathlineto{\pgfqpoint{1.052365in}{2.261530in}}%
\pgfpathlineto{\pgfqpoint{1.052415in}{2.257802in}}%
\pgfpathlineto{\pgfqpoint{1.052464in}{2.261193in}}%
\pgfpathlineto{\pgfqpoint{1.053532in}{1.532028in}}%
\pgfpathlineto{\pgfqpoint{1.053651in}{1.532036in}}%
\pgfpathlineto{\pgfqpoint{1.055478in}{1.532183in}}%
\pgfpathlineto{\pgfqpoint{1.056470in}{2.261007in}}%
\pgfpathlineto{\pgfqpoint{1.056816in}{2.260636in}}%
\pgfpathlineto{\pgfqpoint{1.056866in}{2.260883in}}%
\pgfpathlineto{\pgfqpoint{1.057212in}{2.259691in}}%
\pgfpathlineto{\pgfqpoint{1.057262in}{2.260409in}}%
\pgfpathlineto{\pgfqpoint{1.057565in}{2.138184in}}%
\pgfpathlineto{\pgfqpoint{1.058447in}{1.531994in}}%
\pgfpathlineto{\pgfqpoint{1.058207in}{2.243431in}}%
\pgfpathlineto{\pgfqpoint{1.058520in}{1.532040in}}%
\pgfpathlineto{\pgfqpoint{1.060244in}{1.532084in}}%
\pgfpathlineto{\pgfqpoint{1.060424in}{1.532185in}}%
\pgfpathlineto{\pgfqpoint{1.061381in}{2.259621in}}%
\pgfpathlineto{\pgfqpoint{1.061776in}{2.259499in}}%
\pgfpathlineto{\pgfqpoint{1.062221in}{2.258049in}}%
\pgfpathlineto{\pgfqpoint{1.062271in}{2.258787in}}%
\pgfpathlineto{\pgfqpoint{1.062733in}{1.634100in}}%
\pgfpathlineto{\pgfqpoint{1.063079in}{2.189528in}}%
\pgfpathlineto{\pgfqpoint{1.063128in}{2.198421in}}%
\pgfpathlineto{\pgfqpoint{1.063178in}{2.139212in}}%
\pgfpathlineto{\pgfqpoint{1.063392in}{1.532015in}}%
\pgfpathlineto{\pgfqpoint{1.064187in}{1.532050in}}%
\pgfpathlineto{\pgfqpoint{1.065368in}{1.532170in}}%
\pgfpathlineto{\pgfqpoint{1.065409in}{1.862017in}}%
\pgfpathlineto{\pgfqpoint{1.066189in}{2.258339in}}%
\pgfpathlineto{\pgfqpoint{1.066288in}{2.258365in}}%
\pgfpathlineto{\pgfqpoint{1.067030in}{2.256959in}}%
\pgfpathlineto{\pgfqpoint{1.067080in}{2.257988in}}%
\pgfpathlineto{\pgfqpoint{1.067574in}{1.850058in}}%
\pgfpathlineto{\pgfqpoint{1.068368in}{1.532029in}}%
\pgfpathlineto{\pgfqpoint{1.068118in}{2.198428in}}%
\pgfpathlineto{\pgfqpoint{1.068486in}{1.532034in}}%
\pgfpathlineto{\pgfqpoint{1.070315in}{1.532181in}}%
\pgfpathlineto{\pgfqpoint{1.071307in}{2.257019in}}%
\pgfpathlineto{\pgfqpoint{1.071654in}{2.256633in}}%
\pgfpathlineto{\pgfqpoint{1.071703in}{2.256894in}}%
\pgfpathlineto{\pgfqpoint{1.072049in}{2.255646in}}%
\pgfpathlineto{\pgfqpoint{1.072099in}{2.256425in}}%
\pgfpathlineto{\pgfqpoint{1.072454in}{1.982895in}}%
\pgfpathlineto{\pgfqpoint{1.073283in}{1.531993in}}%
\pgfpathlineto{\pgfqpoint{1.073047in}{2.238603in}}%
\pgfpathlineto{\pgfqpoint{1.073356in}{1.532040in}}%
\pgfpathlineto{\pgfqpoint{1.075080in}{1.532084in}}%
\pgfpathlineto{\pgfqpoint{1.075260in}{1.532185in}}%
\pgfpathlineto{\pgfqpoint{1.076221in}{2.255683in}}%
\pgfpathlineto{\pgfqpoint{1.076616in}{2.255575in}}%
\pgfpathlineto{\pgfqpoint{1.077062in}{2.254172in}}%
\pgfpathlineto{\pgfqpoint{1.077111in}{2.254913in}}%
\pgfpathlineto{\pgfqpoint{1.078317in}{1.532030in}}%
\pgfpathlineto{\pgfqpoint{1.079054in}{1.532059in}}%
\pgfpathlineto{\pgfqpoint{1.080207in}{1.532219in}}%
\pgfpathlineto{\pgfqpoint{1.081269in}{2.254402in}}%
\pgfpathlineto{\pgfqpoint{1.081417in}{2.254203in}}%
\pgfpathlineto{\pgfqpoint{1.081467in}{2.254354in}}%
\pgfpathlineto{\pgfqpoint{1.081912in}{2.253495in}}%
\pgfpathlineto{\pgfqpoint{1.081961in}{2.253865in}}%
\pgfpathlineto{\pgfqpoint{1.082110in}{2.242853in}}%
\pgfpathlineto{\pgfqpoint{1.082159in}{2.253324in}}%
\pgfpathlineto{\pgfqpoint{1.083190in}{1.532020in}}%
\pgfpathlineto{\pgfqpoint{1.083415in}{1.532027in}}%
\pgfpathlineto{\pgfqpoint{1.085153in}{1.532218in}}%
\pgfpathlineto{\pgfqpoint{1.086188in}{2.253222in}}%
\pgfpathlineto{\pgfqpoint{1.086287in}{2.253217in}}%
\pgfpathlineto{\pgfqpoint{1.086831in}{2.251758in}}%
\pgfpathlineto{\pgfqpoint{1.086881in}{2.252817in}}%
\pgfpathlineto{\pgfqpoint{1.087524in}{1.621567in}}%
\pgfpathlineto{\pgfqpoint{1.087820in}{2.148818in}}%
\pgfpathlineto{\pgfqpoint{1.087870in}{2.212098in}}%
\pgfpathlineto{\pgfqpoint{1.088116in}{1.549741in}}%
\pgfpathlineto{\pgfqpoint{1.088158in}{1.532026in}}%
\pgfpathlineto{\pgfqpoint{1.088877in}{1.532059in}}%
\pgfpathlineto{\pgfqpoint{1.090095in}{1.532160in}}%
\pgfpathlineto{\pgfqpoint{1.090137in}{1.881141in}}%
\pgfpathlineto{\pgfqpoint{1.090909in}{2.251950in}}%
\pgfpathlineto{\pgfqpoint{1.091008in}{2.251978in}}%
\pgfpathlineto{\pgfqpoint{1.091750in}{2.249943in}}%
\pgfpathlineto{\pgfqpoint{1.091799in}{2.251650in}}%
\pgfpathlineto{\pgfqpoint{1.092491in}{1.619390in}}%
\pgfpathlineto{\pgfqpoint{1.092689in}{1.942521in}}%
\pgfpathlineto{\pgfqpoint{1.092838in}{2.219946in}}%
\pgfpathlineto{\pgfqpoint{1.093109in}{1.532024in}}%
\pgfpathlineto{\pgfqpoint{1.093292in}{1.532030in}}%
\pgfpathlineto{\pgfqpoint{1.095043in}{1.532207in}}%
\pgfpathlineto{\pgfqpoint{1.096077in}{2.250764in}}%
\pgfpathlineto{\pgfqpoint{1.096226in}{2.250185in}}%
\pgfpathlineto{\pgfqpoint{1.096275in}{2.250747in}}%
\pgfpathlineto{\pgfqpoint{1.096621in}{2.248804in}}%
\pgfpathlineto{\pgfqpoint{1.096671in}{2.250522in}}%
\pgfpathlineto{\pgfqpoint{1.097412in}{1.631172in}}%
\pgfpathlineto{\pgfqpoint{1.097610in}{1.869272in}}%
\pgfpathlineto{\pgfqpoint{1.097808in}{2.200065in}}%
\pgfpathlineto{\pgfqpoint{1.098065in}{1.532021in}}%
\pgfpathlineto{\pgfqpoint{1.098253in}{1.532027in}}%
\pgfpathlineto{\pgfqpoint{1.099990in}{1.532248in}}%
\pgfpathlineto{\pgfqpoint{1.101003in}{2.249423in}}%
\pgfpathlineto{\pgfqpoint{1.101102in}{2.249416in}}%
\pgfpathlineto{\pgfqpoint{1.101844in}{2.247960in}}%
\pgfpathlineto{\pgfqpoint{1.101893in}{2.248543in}}%
\pgfpathlineto{\pgfqpoint{1.102289in}{1.689311in}}%
\pgfpathlineto{\pgfqpoint{1.102388in}{1.641065in}}%
\pgfpathlineto{\pgfqpoint{1.102586in}{1.968952in}}%
\pgfpathlineto{\pgfqpoint{1.102734in}{2.188942in}}%
\pgfpathlineto{\pgfqpoint{1.102972in}{1.532019in}}%
\pgfpathlineto{\pgfqpoint{1.103148in}{1.532044in}}%
\pgfpathlineto{\pgfqpoint{1.104923in}{1.532112in}}%
\pgfpathlineto{\pgfqpoint{1.104934in}{1.532204in}}%
\pgfpathlineto{\pgfqpoint{1.105973in}{2.248345in}}%
\pgfpathlineto{\pgfqpoint{1.106121in}{2.247881in}}%
\pgfpathlineto{\pgfqpoint{1.106171in}{2.248329in}}%
\pgfpathlineto{\pgfqpoint{1.106616in}{2.246004in}}%
\pgfpathlineto{\pgfqpoint{1.106665in}{2.247999in}}%
\pgfpathlineto{\pgfqpoint{1.107358in}{1.616451in}}%
\pgfpathlineto{\pgfqpoint{1.107555in}{1.952953in}}%
\pgfpathlineto{\pgfqpoint{1.107704in}{2.225165in}}%
\pgfpathlineto{\pgfqpoint{1.107907in}{1.532015in}}%
\pgfpathlineto{\pgfqpoint{1.108154in}{1.532027in}}%
\pgfpathlineto{\pgfqpoint{1.109880in}{1.532208in}}%
\pgfpathlineto{\pgfqpoint{1.111012in}{2.247188in}}%
\pgfpathlineto{\pgfqpoint{1.111061in}{2.246498in}}%
\pgfpathlineto{\pgfqpoint{1.111111in}{2.247176in}}%
\pgfpathlineto{\pgfqpoint{1.111358in}{2.245655in}}%
\pgfpathlineto{\pgfqpoint{1.111407in}{2.247048in}}%
\pgfpathlineto{\pgfqpoint{1.112298in}{1.640451in}}%
\pgfpathlineto{\pgfqpoint{1.112495in}{1.956763in}}%
\pgfpathlineto{\pgfqpoint{1.112644in}{2.189418in}}%
\pgfpathlineto{\pgfqpoint{1.112852in}{1.532010in}}%
\pgfpathlineto{\pgfqpoint{1.113049in}{1.532044in}}%
\pgfpathlineto{\pgfqpoint{1.114814in}{1.532112in}}%
\pgfpathlineto{\pgfqpoint{1.114825in}{1.532209in}}%
\pgfpathlineto{\pgfqpoint{1.115862in}{2.245963in}}%
\pgfpathlineto{\pgfqpoint{1.116010in}{2.245562in}}%
\pgfpathlineto{\pgfqpoint{1.116060in}{2.245947in}}%
\pgfpathlineto{\pgfqpoint{1.116505in}{2.244425in}}%
\pgfpathlineto{\pgfqpoint{1.116554in}{2.245613in}}%
\pgfpathlineto{\pgfqpoint{1.117246in}{1.640292in}}%
\pgfpathlineto{\pgfqpoint{1.117494in}{2.064803in}}%
\pgfpathlineto{\pgfqpoint{1.117593in}{2.188979in}}%
\pgfpathlineto{\pgfqpoint{1.117796in}{1.531993in}}%
\pgfpathlineto{\pgfqpoint{1.117963in}{1.532044in}}%
\pgfpathlineto{\pgfqpoint{1.119758in}{1.532110in}}%
\pgfpathlineto{\pgfqpoint{1.119769in}{1.532174in}}%
\pgfpathlineto{\pgfqpoint{1.120822in}{2.244749in}}%
\pgfpathlineto{\pgfqpoint{1.121168in}{2.244410in}}%
\pgfpathlineto{\pgfqpoint{1.121217in}{2.244614in}}%
\pgfpathlineto{\pgfqpoint{1.121563in}{2.243563in}}%
\pgfpathlineto{\pgfqpoint{1.121613in}{2.244124in}}%
\pgfpathlineto{\pgfqpoint{1.121662in}{2.242537in}}%
\pgfpathlineto{\pgfqpoint{1.121712in}{2.243892in}}%
\pgfpathlineto{\pgfqpoint{1.122169in}{1.611456in}}%
\pgfpathlineto{\pgfqpoint{1.122466in}{2.157057in}}%
\pgfpathlineto{\pgfqpoint{1.122515in}{2.226665in}}%
\pgfpathlineto{\pgfqpoint{1.122755in}{1.532022in}}%
\pgfpathlineto{\pgfqpoint{1.122789in}{1.532022in}}%
\pgfpathlineto{\pgfqpoint{1.124716in}{1.532216in}}%
\pgfpathlineto{\pgfqpoint{1.125848in}{2.243723in}}%
\pgfpathlineto{\pgfqpoint{1.125897in}{2.243097in}}%
\pgfpathlineto{\pgfqpoint{1.125946in}{2.243713in}}%
\pgfpathlineto{\pgfqpoint{1.126293in}{2.241386in}}%
\pgfpathlineto{\pgfqpoint{1.126342in}{2.243516in}}%
\pgfpathlineto{\pgfqpoint{1.127133in}{1.650389in}}%
\pgfpathlineto{\pgfqpoint{1.127331in}{1.910512in}}%
\pgfpathlineto{\pgfqpoint{1.127480in}{2.172555in}}%
\pgfpathlineto{\pgfqpoint{1.127688in}{1.531979in}}%
\pgfpathlineto{\pgfqpoint{1.127912in}{1.532023in}}%
\pgfpathlineto{\pgfqpoint{1.129658in}{1.532150in}}%
\pgfpathlineto{\pgfqpoint{1.129686in}{1.606525in}}%
\pgfpathlineto{\pgfqpoint{1.130468in}{2.242375in}}%
\pgfpathlineto{\pgfqpoint{1.131407in}{2.241851in}}%
\pgfpathlineto{\pgfqpoint{1.131704in}{2.240293in}}%
\pgfpathlineto{\pgfqpoint{1.131852in}{2.021299in}}%
\pgfpathlineto{\pgfqpoint{1.132633in}{1.531975in}}%
\pgfpathlineto{\pgfqpoint{1.132446in}{2.225002in}}%
\pgfpathlineto{\pgfqpoint{1.132772in}{1.532018in}}%
\pgfpathlineto{\pgfqpoint{1.134404in}{1.532106in}}%
\pgfpathlineto{\pgfqpoint{1.134593in}{1.532105in}}%
\pgfpathlineto{\pgfqpoint{1.134609in}{1.532308in}}%
\pgfpathlineto{\pgfqpoint{1.135591in}{2.241315in}}%
\pgfpathlineto{\pgfqpoint{1.135739in}{2.241266in}}%
\pgfpathlineto{\pgfqpoint{1.136283in}{2.240899in}}%
\pgfpathlineto{\pgfqpoint{1.136629in}{2.238364in}}%
\pgfpathlineto{\pgfqpoint{1.136679in}{2.239048in}}%
\pgfpathlineto{\pgfqpoint{1.137579in}{1.532018in}}%
\pgfpathlineto{\pgfqpoint{1.137871in}{1.532045in}}%
\pgfpathlineto{\pgfqpoint{1.139542in}{1.532115in}}%
\pgfpathlineto{\pgfqpoint{1.139552in}{1.532208in}}%
\pgfpathlineto{\pgfqpoint{1.140684in}{2.240385in}}%
\pgfpathlineto{\pgfqpoint{1.140733in}{2.239792in}}%
\pgfpathlineto{\pgfqpoint{1.140783in}{2.240378in}}%
\pgfpathlineto{\pgfqpoint{1.141129in}{2.238632in}}%
\pgfpathlineto{\pgfqpoint{1.141179in}{2.240196in}}%
\pgfpathlineto{\pgfqpoint{1.142019in}{1.635277in}}%
\pgfpathlineto{\pgfqpoint{1.142217in}{1.973941in}}%
\pgfpathlineto{\pgfqpoint{1.142365in}{2.197529in}}%
\pgfpathlineto{\pgfqpoint{1.142554in}{1.532025in}}%
\pgfpathlineto{\pgfqpoint{1.142771in}{1.532040in}}%
\pgfpathlineto{\pgfqpoint{1.144495in}{1.532163in}}%
\pgfpathlineto{\pgfqpoint{1.144538in}{1.911854in}}%
\pgfpathlineto{\pgfqpoint{1.145305in}{2.239213in}}%
\pgfpathlineto{\pgfqpoint{1.145503in}{2.239270in}}%
\pgfpathlineto{\pgfqpoint{1.146146in}{2.238085in}}%
\pgfpathlineto{\pgfqpoint{1.146195in}{2.239000in}}%
\pgfpathlineto{\pgfqpoint{1.146245in}{2.237179in}}%
\pgfpathlineto{\pgfqpoint{1.146294in}{2.238884in}}%
\pgfpathlineto{\pgfqpoint{1.146937in}{1.619529in}}%
\pgfpathlineto{\pgfqpoint{1.147185in}{2.026282in}}%
\pgfpathlineto{\pgfqpoint{1.147283in}{2.214195in}}%
\pgfpathlineto{\pgfqpoint{1.147526in}{1.532023in}}%
\pgfpathlineto{\pgfqpoint{1.147717in}{1.532030in}}%
\pgfpathlineto{\pgfqpoint{1.149443in}{1.532209in}}%
\pgfpathlineto{\pgfqpoint{1.150575in}{2.238274in}}%
\pgfpathlineto{\pgfqpoint{1.150624in}{2.237623in}}%
\pgfpathlineto{\pgfqpoint{1.150674in}{2.238269in}}%
\pgfpathlineto{\pgfqpoint{1.151020in}{2.235907in}}%
\pgfpathlineto{\pgfqpoint{1.151070in}{2.238099in}}%
\pgfpathlineto{\pgfqpoint{1.151910in}{1.638720in}}%
\pgfpathlineto{\pgfqpoint{1.152059in}{1.815259in}}%
\pgfpathlineto{\pgfqpoint{1.152256in}{2.190988in}}%
\pgfpathlineto{\pgfqpoint{1.152419in}{1.532010in}}%
\pgfpathlineto{\pgfqpoint{1.152682in}{1.532025in}}%
\pgfpathlineto{\pgfqpoint{1.154386in}{1.532156in}}%
\pgfpathlineto{\pgfqpoint{1.154427in}{1.844182in}}%
\pgfpathlineto{\pgfqpoint{1.155195in}{2.237100in}}%
\pgfpathlineto{\pgfqpoint{1.155294in}{2.237135in}}%
\pgfpathlineto{\pgfqpoint{1.156036in}{2.236289in}}%
\pgfpathlineto{\pgfqpoint{1.156085in}{2.236883in}}%
\pgfpathlineto{\pgfqpoint{1.156234in}{2.234654in}}%
\pgfpathlineto{\pgfqpoint{1.156283in}{2.236618in}}%
\pgfpathlineto{\pgfqpoint{1.156827in}{1.649154in}}%
\pgfpathlineto{\pgfqpoint{1.157124in}{2.106241in}}%
\pgfpathlineto{\pgfqpoint{1.157223in}{2.173494in}}%
\pgfpathlineto{\pgfqpoint{1.157380in}{1.532020in}}%
\pgfpathlineto{\pgfqpoint{1.157427in}{1.532021in}}%
\pgfpathlineto{\pgfqpoint{1.159334in}{1.532210in}}%
\pgfpathlineto{\pgfqpoint{1.160465in}{2.236251in}}%
\pgfpathlineto{\pgfqpoint{1.160515in}{2.235682in}}%
\pgfpathlineto{\pgfqpoint{1.160564in}{2.236248in}}%
\pgfpathlineto{\pgfqpoint{1.160911in}{2.234781in}}%
\pgfpathlineto{\pgfqpoint{1.160960in}{2.236089in}}%
\pgfpathlineto{\pgfqpoint{1.161801in}{1.634371in}}%
\pgfpathlineto{\pgfqpoint{1.162048in}{2.024748in}}%
\pgfpathlineto{\pgfqpoint{1.162147in}{2.193453in}}%
\pgfpathlineto{\pgfqpoint{1.162340in}{1.532026in}}%
\pgfpathlineto{\pgfqpoint{1.162560in}{1.532040in}}%
\pgfpathlineto{\pgfqpoint{1.164274in}{1.532139in}}%
\pgfpathlineto{\pgfqpoint{1.164283in}{1.532456in}}%
\pgfpathlineto{\pgfqpoint{1.165272in}{2.235068in}}%
\pgfpathlineto{\pgfqpoint{1.165321in}{2.235067in}}%
\pgfpathlineto{\pgfqpoint{1.166112in}{2.234409in}}%
\pgfpathlineto{\pgfqpoint{1.166360in}{2.233144in}}%
\pgfpathlineto{\pgfqpoint{1.166558in}{1.883237in}}%
\pgfpathlineto{\pgfqpoint{1.167256in}{1.532010in}}%
\pgfpathlineto{\pgfqpoint{1.167102in}{2.212478in}}%
\pgfpathlineto{\pgfqpoint{1.167467in}{1.532046in}}%
\pgfpathlineto{\pgfqpoint{1.169212in}{1.532108in}}%
\pgfpathlineto{\pgfqpoint{1.169222in}{1.532160in}}%
\pgfpathlineto{\pgfqpoint{1.169264in}{1.867956in}}%
\pgfpathlineto{\pgfqpoint{1.170034in}{2.234165in}}%
\pgfpathlineto{\pgfqpoint{1.170232in}{2.234228in}}%
\pgfpathlineto{\pgfqpoint{1.170874in}{2.233066in}}%
\pgfpathlineto{\pgfqpoint{1.170924in}{2.233994in}}%
\pgfpathlineto{\pgfqpoint{1.170973in}{2.232251in}}%
\pgfpathlineto{\pgfqpoint{1.171023in}{2.233889in}}%
\pgfpathlineto{\pgfqpoint{1.171715in}{1.630990in}}%
\pgfpathlineto{\pgfqpoint{1.171913in}{1.953784in}}%
\pgfpathlineto{\pgfqpoint{1.172061in}{2.203130in}}%
\pgfpathlineto{\pgfqpoint{1.172214in}{1.532023in}}%
\pgfpathlineto{\pgfqpoint{1.172483in}{1.532044in}}%
\pgfpathlineto{\pgfqpoint{1.174159in}{1.532113in}}%
\pgfpathlineto{\pgfqpoint{1.174171in}{1.532222in}}%
\pgfpathlineto{\pgfqpoint{1.175300in}{2.233312in}}%
\pgfpathlineto{\pgfqpoint{1.175844in}{2.231756in}}%
\pgfpathlineto{\pgfqpoint{1.175894in}{2.233067in}}%
\pgfpathlineto{\pgfqpoint{1.176685in}{1.631806in}}%
\pgfpathlineto{\pgfqpoint{1.176883in}{1.983971in}}%
\pgfpathlineto{\pgfqpoint{1.177031in}{2.202915in}}%
\pgfpathlineto{\pgfqpoint{1.177175in}{1.532027in}}%
\pgfpathlineto{\pgfqpoint{1.177441in}{1.532036in}}%
\pgfpathlineto{\pgfqpoint{1.179113in}{1.532159in}}%
\pgfpathlineto{\pgfqpoint{1.179155in}{1.880404in}}%
\pgfpathlineto{\pgfqpoint{1.179923in}{2.232280in}}%
\pgfpathlineto{\pgfqpoint{1.180121in}{2.232348in}}%
\pgfpathlineto{\pgfqpoint{1.180764in}{2.231363in}}%
\pgfpathlineto{\pgfqpoint{1.180813in}{2.232136in}}%
\pgfpathlineto{\pgfqpoint{1.180862in}{2.230891in}}%
\pgfpathlineto{\pgfqpoint{1.180912in}{2.232036in}}%
\pgfpathlineto{\pgfqpoint{1.181604in}{1.653844in}}%
\pgfpathlineto{\pgfqpoint{1.181901in}{2.112274in}}%
\pgfpathlineto{\pgfqpoint{1.181950in}{2.166301in}}%
\pgfpathlineto{\pgfqpoint{1.182092in}{1.531993in}}%
\pgfpathlineto{\pgfqpoint{1.182119in}{1.532046in}}%
\pgfpathlineto{\pgfqpoint{1.183126in}{1.532054in}}%
\pgfpathlineto{\pgfqpoint{1.184066in}{1.532855in}}%
\pgfpathlineto{\pgfqpoint{1.185013in}{2.231393in}}%
\pgfpathlineto{\pgfqpoint{1.185062in}{2.231224in}}%
\pgfpathlineto{\pgfqpoint{1.185408in}{2.231388in}}%
\pgfpathlineto{\pgfqpoint{1.185656in}{2.230874in}}%
\pgfpathlineto{\pgfqpoint{1.185705in}{2.231216in}}%
\pgfpathlineto{\pgfqpoint{1.186201in}{2.228664in}}%
\pgfpathlineto{\pgfqpoint{1.187061in}{1.532028in}}%
\pgfpathlineto{\pgfqpoint{1.187959in}{1.532064in}}%
\pgfpathlineto{\pgfqpoint{1.189003in}{1.532145in}}%
\pgfpathlineto{\pgfqpoint{1.189011in}{1.532706in}}%
\pgfpathlineto{\pgfqpoint{1.189990in}{2.230498in}}%
\pgfpathlineto{\pgfqpoint{1.190039in}{2.230343in}}%
\pgfpathlineto{\pgfqpoint{1.190287in}{2.230496in}}%
\pgfpathlineto{\pgfqpoint{1.190633in}{2.229955in}}%
\pgfpathlineto{\pgfqpoint{1.190682in}{2.230265in}}%
\pgfpathlineto{\pgfqpoint{1.190929in}{2.228881in}}%
\pgfpathlineto{\pgfqpoint{1.190979in}{2.229833in}}%
\pgfpathlineto{\pgfqpoint{1.191514in}{1.634788in}}%
\pgfpathlineto{\pgfqpoint{1.191810in}{2.158983in}}%
\pgfpathlineto{\pgfqpoint{1.191860in}{2.196649in}}%
\pgfpathlineto{\pgfqpoint{1.191968in}{1.716644in}}%
\pgfpathlineto{\pgfqpoint{1.192019in}{1.532026in}}%
\pgfpathlineto{\pgfqpoint{1.192738in}{1.532059in}}%
\pgfpathlineto{\pgfqpoint{1.193947in}{1.532139in}}%
\pgfpathlineto{\pgfqpoint{1.193956in}{1.532457in}}%
\pgfpathlineto{\pgfqpoint{1.195002in}{2.229560in}}%
\pgfpathlineto{\pgfqpoint{1.195892in}{2.228771in}}%
\pgfpathlineto{\pgfqpoint{1.196090in}{2.227538in}}%
\pgfpathlineto{\pgfqpoint{1.196964in}{1.532030in}}%
\pgfpathlineto{\pgfqpoint{1.197731in}{1.532057in}}%
\pgfpathlineto{\pgfqpoint{1.198901in}{1.532426in}}%
\pgfpathlineto{\pgfqpoint{1.199964in}{2.228725in}}%
\pgfpathlineto{\pgfqpoint{1.200904in}{2.227322in}}%
\pgfpathlineto{\pgfqpoint{1.200953in}{2.227702in}}%
\pgfpathlineto{\pgfqpoint{1.201201in}{1.944118in}}%
\pgfpathlineto{\pgfqpoint{1.201890in}{1.532025in}}%
\pgfpathlineto{\pgfqpoint{1.201745in}{2.192876in}}%
\pgfpathlineto{\pgfqpoint{1.202115in}{1.532032in}}%
\pgfpathlineto{\pgfqpoint{1.203841in}{1.532167in}}%
\pgfpathlineto{\pgfqpoint{1.203884in}{1.942400in}}%
\pgfpathlineto{\pgfqpoint{1.204653in}{2.227882in}}%
\pgfpathlineto{\pgfqpoint{1.204752in}{2.227923in}}%
\pgfpathlineto{\pgfqpoint{1.205494in}{2.227283in}}%
\pgfpathlineto{\pgfqpoint{1.205543in}{2.227743in}}%
\pgfpathlineto{\pgfqpoint{1.205691in}{2.226634in}}%
\pgfpathlineto{\pgfqpoint{1.205741in}{2.227514in}}%
\pgfpathlineto{\pgfqpoint{1.206087in}{2.156917in}}%
\pgfpathlineto{\pgfqpoint{1.206821in}{1.531997in}}%
\pgfpathlineto{\pgfqpoint{1.206724in}{2.219348in}}%
\pgfpathlineto{\pgfqpoint{1.207112in}{1.532051in}}%
\pgfpathlineto{\pgfqpoint{1.208776in}{1.532109in}}%
\pgfpathlineto{\pgfqpoint{1.208795in}{1.535523in}}%
\pgfpathlineto{\pgfqpoint{1.209788in}{2.227140in}}%
\pgfpathlineto{\pgfqpoint{1.210628in}{2.225842in}}%
\pgfpathlineto{\pgfqpoint{1.210678in}{2.226750in}}%
\pgfpathlineto{\pgfqpoint{1.210727in}{2.221286in}}%
\pgfpathlineto{\pgfqpoint{1.210777in}{2.226591in}}%
\pgfpathlineto{\pgfqpoint{1.211321in}{1.645711in}}%
\pgfpathlineto{\pgfqpoint{1.211568in}{2.028023in}}%
\pgfpathlineto{\pgfqpoint{1.211667in}{2.178595in}}%
\pgfpathlineto{\pgfqpoint{1.211844in}{1.532030in}}%
\pgfpathlineto{\pgfqpoint{1.212084in}{1.532042in}}%
\pgfpathlineto{\pgfqpoint{1.213733in}{1.532185in}}%
\pgfpathlineto{\pgfqpoint{1.214925in}{2.226396in}}%
\pgfpathlineto{\pgfqpoint{1.215024in}{2.226390in}}%
\pgfpathlineto{\pgfqpoint{1.215666in}{2.224485in}}%
\pgfpathlineto{\pgfqpoint{1.215716in}{2.225860in}}%
\pgfpathlineto{\pgfqpoint{1.216281in}{1.636593in}}%
\pgfpathlineto{\pgfqpoint{1.216528in}{2.024327in}}%
\pgfpathlineto{\pgfqpoint{1.216652in}{2.194676in}}%
\pgfpathlineto{\pgfqpoint{1.216744in}{1.532030in}}%
\pgfpathlineto{\pgfqpoint{1.217062in}{1.532044in}}%
\pgfpathlineto{\pgfqpoint{1.218676in}{1.532157in}}%
\pgfpathlineto{\pgfqpoint{1.218719in}{1.859558in}}%
\pgfpathlineto{\pgfqpoint{1.219484in}{2.225581in}}%
\pgfpathlineto{\pgfqpoint{1.219781in}{2.225692in}}%
\pgfpathlineto{\pgfqpoint{1.220226in}{2.225136in}}%
\pgfpathlineto{\pgfqpoint{1.220275in}{2.225604in}}%
\pgfpathlineto{\pgfqpoint{1.220523in}{2.224316in}}%
\pgfpathlineto{\pgfqpoint{1.220572in}{2.225352in}}%
\pgfpathlineto{\pgfqpoint{1.221067in}{1.924467in}}%
\pgfpathlineto{\pgfqpoint{1.221657in}{1.532003in}}%
\pgfpathlineto{\pgfqpoint{1.221630in}{2.170839in}}%
\pgfpathlineto{\pgfqpoint{1.221988in}{1.532030in}}%
\pgfpathlineto{\pgfqpoint{1.223622in}{1.532161in}}%
\pgfpathlineto{\pgfqpoint{1.223663in}{1.846168in}}%
\pgfpathlineto{\pgfqpoint{1.224433in}{2.224895in}}%
\pgfpathlineto{\pgfqpoint{1.224532in}{2.224946in}}%
\pgfpathlineto{\pgfqpoint{1.225274in}{2.224519in}}%
\pgfpathlineto{\pgfqpoint{1.225323in}{2.224851in}}%
\pgfpathlineto{\pgfqpoint{1.225669in}{2.222067in}}%
\pgfpathlineto{\pgfqpoint{1.225719in}{2.224332in}}%
\pgfpathlineto{\pgfqpoint{1.226213in}{1.628387in}}%
\pgfpathlineto{\pgfqpoint{1.226510in}{2.129140in}}%
\pgfpathlineto{\pgfqpoint{1.226573in}{2.222511in}}%
\pgfpathlineto{\pgfqpoint{1.226603in}{1.531998in}}%
\pgfpathlineto{\pgfqpoint{1.226929in}{1.532030in}}%
\pgfpathlineto{\pgfqpoint{1.228568in}{1.532163in}}%
\pgfpathlineto{\pgfqpoint{1.228609in}{1.848750in}}%
\pgfpathlineto{\pgfqpoint{1.229377in}{2.224288in}}%
\pgfpathlineto{\pgfqpoint{1.229476in}{2.224342in}}%
\pgfpathlineto{\pgfqpoint{1.230218in}{2.223992in}}%
\pgfpathlineto{\pgfqpoint{1.230267in}{2.224271in}}%
\pgfpathlineto{\pgfqpoint{1.230613in}{2.222758in}}%
\pgfpathlineto{\pgfqpoint{1.230663in}{2.223773in}}%
\pgfpathlineto{\pgfqpoint{1.231207in}{1.633419in}}%
\pgfpathlineto{\pgfqpoint{1.231484in}{2.130568in}}%
\pgfpathlineto{\pgfqpoint{1.231519in}{2.196978in}}%
\pgfpathlineto{\pgfqpoint{1.231548in}{1.532007in}}%
\pgfpathlineto{\pgfqpoint{1.231819in}{1.532052in}}%
\pgfpathlineto{\pgfqpoint{1.233504in}{1.532115in}}%
\pgfpathlineto{\pgfqpoint{1.233515in}{1.532188in}}%
\pgfpathlineto{\pgfqpoint{1.234832in}{2.223948in}}%
\pgfpathlineto{\pgfqpoint{1.235475in}{2.220561in}}%
\pgfpathlineto{\pgfqpoint{1.235524in}{2.223557in}}%
\pgfpathlineto{\pgfqpoint{1.236172in}{1.637304in}}%
\pgfpathlineto{\pgfqpoint{1.236414in}{2.023066in}}%
\pgfpathlineto{\pgfqpoint{1.236465in}{2.150766in}}%
\pgfpathlineto{\pgfqpoint{1.236492in}{1.532004in}}%
\pgfpathlineto{\pgfqpoint{1.236963in}{1.532061in}}%
\pgfpathlineto{\pgfqpoint{1.238449in}{1.532115in}}%
\pgfpathlineto{\pgfqpoint{1.238459in}{1.532174in}}%
\pgfpathlineto{\pgfqpoint{1.238501in}{1.902272in}}%
\pgfpathlineto{\pgfqpoint{1.239266in}{2.223292in}}%
\pgfpathlineto{\pgfqpoint{1.239563in}{2.223442in}}%
\pgfpathlineto{\pgfqpoint{1.240008in}{2.223018in}}%
\pgfpathlineto{\pgfqpoint{1.240057in}{2.223431in}}%
\pgfpathlineto{\pgfqpoint{1.240403in}{2.222096in}}%
\pgfpathlineto{\pgfqpoint{1.240453in}{2.223139in}}%
\pgfpathlineto{\pgfqpoint{1.240947in}{1.944937in}}%
\pgfpathlineto{\pgfqpoint{1.241436in}{1.532001in}}%
\pgfpathlineto{\pgfqpoint{1.241410in}{2.073303in}}%
\pgfpathlineto{\pgfqpoint{1.241876in}{1.532035in}}%
\pgfpathlineto{\pgfqpoint{1.243412in}{1.532904in}}%
\pgfpathlineto{\pgfqpoint{1.244438in}{2.222949in}}%
\pgfpathlineto{\pgfqpoint{1.245377in}{2.222090in}}%
\pgfpathlineto{\pgfqpoint{1.245427in}{2.222586in}}%
\pgfpathlineto{\pgfqpoint{1.245476in}{2.221668in}}%
\pgfpathlineto{\pgfqpoint{1.245526in}{2.222431in}}%
\pgfpathlineto{\pgfqpoint{1.245864in}{2.046310in}}%
\pgfpathlineto{\pgfqpoint{1.246382in}{1.532003in}}%
\pgfpathlineto{\pgfqpoint{1.246355in}{2.082932in}}%
\pgfpathlineto{\pgfqpoint{1.246813in}{1.532036in}}%
\pgfpathlineto{\pgfqpoint{1.248349in}{1.532160in}}%
\pgfpathlineto{\pgfqpoint{1.248376in}{1.606135in}}%
\pgfpathlineto{\pgfqpoint{1.249165in}{2.222362in}}%
\pgfpathlineto{\pgfqpoint{1.250104in}{2.222389in}}%
\pgfpathlineto{\pgfqpoint{1.250648in}{2.220495in}}%
\pgfpathlineto{\pgfqpoint{1.250698in}{2.220810in}}%
\pgfpathlineto{\pgfqpoint{1.251327in}{1.532007in}}%
\pgfpathlineto{\pgfqpoint{1.252029in}{1.532072in}}%
\pgfpathlineto{\pgfqpoint{1.253284in}{1.532114in}}%
\pgfpathlineto{\pgfqpoint{1.253294in}{1.532156in}}%
\pgfpathlineto{\pgfqpoint{1.253320in}{1.599688in}}%
\pgfpathlineto{\pgfqpoint{1.254072in}{2.222040in}}%
\pgfpathlineto{\pgfqpoint{1.254121in}{2.221955in}}%
\pgfpathlineto{\pgfqpoint{1.254863in}{2.222216in}}%
\pgfpathlineto{\pgfqpoint{1.254962in}{2.222173in}}%
\pgfpathlineto{\pgfqpoint{1.255407in}{2.220848in}}%
\pgfpathlineto{\pgfqpoint{1.255457in}{2.221634in}}%
\pgfpathlineto{\pgfqpoint{1.255852in}{1.839607in}}%
\pgfpathlineto{\pgfqpoint{1.256271in}{1.532022in}}%
\pgfpathlineto{\pgfqpoint{1.256246in}{2.026636in}}%
\pgfpathlineto{\pgfqpoint{1.256743in}{1.532057in}}%
\pgfpathlineto{\pgfqpoint{1.258232in}{1.532119in}}%
\pgfpathlineto{\pgfqpoint{1.258243in}{1.532217in}}%
\pgfpathlineto{\pgfqpoint{1.259374in}{2.221996in}}%
\pgfpathlineto{\pgfqpoint{1.259424in}{2.221494in}}%
\pgfpathlineto{\pgfqpoint{1.259770in}{2.222047in}}%
\pgfpathlineto{\pgfqpoint{1.260017in}{2.220431in}}%
\pgfpathlineto{\pgfqpoint{1.260067in}{2.221944in}}%
\pgfpathlineto{\pgfqpoint{1.261214in}{1.532006in}}%
\pgfpathlineto{\pgfqpoint{1.261994in}{1.532075in}}%
\pgfpathlineto{\pgfqpoint{1.263176in}{1.532115in}}%
\pgfpathlineto{\pgfqpoint{1.263185in}{1.532166in}}%
\pgfpathlineto{\pgfqpoint{1.263226in}{1.831867in}}%
\pgfpathlineto{\pgfqpoint{1.263991in}{2.221434in}}%
\pgfpathlineto{\pgfqpoint{1.264387in}{2.221636in}}%
\pgfpathlineto{\pgfqpoint{1.264733in}{2.221265in}}%
\pgfpathlineto{\pgfqpoint{1.264782in}{2.221632in}}%
\pgfpathlineto{\pgfqpoint{1.265227in}{2.220074in}}%
\pgfpathlineto{\pgfqpoint{1.265277in}{2.221268in}}%
\pgfpathlineto{\pgfqpoint{1.266160in}{1.532006in}}%
\pgfpathlineto{\pgfqpoint{1.267083in}{1.532055in}}%
\pgfpathlineto{\pgfqpoint{1.268131in}{1.532169in}}%
\pgfpathlineto{\pgfqpoint{1.268172in}{1.849117in}}%
\pgfpathlineto{\pgfqpoint{1.268938in}{2.221118in}}%
\pgfpathlineto{\pgfqpoint{1.269037in}{2.221185in}}%
\pgfpathlineto{\pgfqpoint{1.269779in}{2.221027in}}%
\pgfpathlineto{\pgfqpoint{1.269829in}{2.221247in}}%
\pgfpathlineto{\pgfqpoint{1.270274in}{2.219934in}}%
\pgfpathlineto{\pgfqpoint{1.270323in}{2.220703in}}%
\pgfpathlineto{\pgfqpoint{1.270669in}{1.967532in}}%
\pgfpathlineto{\pgfqpoint{1.271105in}{1.532006in}}%
\pgfpathlineto{\pgfqpoint{1.271583in}{1.532064in}}%
\pgfpathlineto{\pgfqpoint{1.273067in}{1.532116in}}%
\pgfpathlineto{\pgfqpoint{1.273077in}{1.532170in}}%
\pgfpathlineto{\pgfqpoint{1.273118in}{1.853321in}}%
\pgfpathlineto{\pgfqpoint{1.273883in}{2.220874in}}%
\pgfpathlineto{\pgfqpoint{1.274279in}{2.221081in}}%
\pgfpathlineto{\pgfqpoint{1.274625in}{2.220695in}}%
\pgfpathlineto{\pgfqpoint{1.274675in}{2.221081in}}%
\pgfpathlineto{\pgfqpoint{1.275120in}{2.219431in}}%
\pgfpathlineto{\pgfqpoint{1.275169in}{2.220730in}}%
\pgfpathlineto{\pgfqpoint{1.276051in}{1.532004in}}%
\pgfpathlineto{\pgfqpoint{1.276833in}{1.532076in}}%
\pgfpathlineto{\pgfqpoint{1.278012in}{1.532115in}}%
\pgfpathlineto{\pgfqpoint{1.278021in}{1.532162in}}%
\pgfpathlineto{\pgfqpoint{1.278050in}{1.609188in}}%
\pgfpathlineto{\pgfqpoint{1.278813in}{2.220535in}}%
\pgfpathlineto{\pgfqpoint{1.278912in}{2.220604in}}%
\pgfpathlineto{\pgfqpoint{1.279555in}{2.220531in}}%
\pgfpathlineto{\pgfqpoint{1.279604in}{2.220726in}}%
\pgfpathlineto{\pgfqpoint{1.280049in}{2.219877in}}%
\pgfpathlineto{\pgfqpoint{1.280099in}{2.220327in}}%
\pgfpathlineto{\pgfqpoint{1.280247in}{2.217544in}}%
\pgfpathlineto{\pgfqpoint{1.280297in}{2.219951in}}%
\pgfpathlineto{\pgfqpoint{1.281031in}{1.532026in}}%
\pgfpathlineto{\pgfqpoint{1.281502in}{1.532043in}}%
\pgfpathlineto{\pgfqpoint{1.282968in}{1.532182in}}%
\pgfpathlineto{\pgfqpoint{1.283008in}{1.850743in}}%
\pgfpathlineto{\pgfqpoint{1.283781in}{2.220318in}}%
\pgfpathlineto{\pgfqpoint{1.284078in}{2.220489in}}%
\pgfpathlineto{\pgfqpoint{1.284523in}{2.220191in}}%
\pgfpathlineto{\pgfqpoint{1.284572in}{2.220516in}}%
\pgfpathlineto{\pgfqpoint{1.285017in}{2.219218in}}%
\pgfpathlineto{\pgfqpoint{1.285067in}{2.220148in}}%
\pgfpathlineto{\pgfqpoint{1.285215in}{2.150573in}}%
\pgfpathlineto{\pgfqpoint{1.285413in}{2.152226in}}%
\pgfpathlineto{\pgfqpoint{1.285942in}{1.532002in}}%
\pgfpathlineto{\pgfqpoint{1.286581in}{1.532043in}}%
\pgfpathlineto{\pgfqpoint{1.287921in}{1.533359in}}%
\pgfpathlineto{\pgfqpoint{1.288949in}{2.220133in}}%
\pgfpathlineto{\pgfqpoint{1.290086in}{2.218820in}}%
\pgfpathlineto{\pgfqpoint{1.290136in}{2.219473in}}%
\pgfpathlineto{\pgfqpoint{1.290432in}{1.997802in}}%
\pgfpathlineto{\pgfqpoint{1.290888in}{1.532005in}}%
\pgfpathlineto{\pgfqpoint{1.291366in}{1.532063in}}%
\pgfpathlineto{\pgfqpoint{1.292848in}{1.532115in}}%
\pgfpathlineto{\pgfqpoint{1.292858in}{1.532168in}}%
\pgfpathlineto{\pgfqpoint{1.292900in}{1.845845in}}%
\pgfpathlineto{\pgfqpoint{1.293667in}{2.219738in}}%
\pgfpathlineto{\pgfqpoint{1.293865in}{2.219859in}}%
\pgfpathlineto{\pgfqpoint{1.294409in}{2.219686in}}%
\pgfpathlineto{\pgfqpoint{1.294458in}{2.219910in}}%
\pgfpathlineto{\pgfqpoint{1.294903in}{2.218933in}}%
\pgfpathlineto{\pgfqpoint{1.294953in}{2.219499in}}%
\pgfpathlineto{\pgfqpoint{1.295002in}{2.218391in}}%
\pgfpathlineto{\pgfqpoint{1.295052in}{2.219336in}}%
\pgfpathlineto{\pgfqpoint{1.295833in}{1.532008in}}%
\pgfpathlineto{\pgfqpoint{1.296941in}{1.532063in}}%
\pgfpathlineto{\pgfqpoint{1.297805in}{1.532180in}}%
\pgfpathlineto{\pgfqpoint{1.299173in}{2.219705in}}%
\pgfpathlineto{\pgfqpoint{1.299223in}{2.219450in}}%
\pgfpathlineto{\pgfqpoint{1.299371in}{2.219684in}}%
\pgfpathlineto{\pgfqpoint{1.299816in}{2.218683in}}%
\pgfpathlineto{\pgfqpoint{1.299866in}{2.219341in}}%
\pgfpathlineto{\pgfqpoint{1.299915in}{2.218148in}}%
\pgfpathlineto{\pgfqpoint{1.299965in}{2.219204in}}%
\pgfpathlineto{\pgfqpoint{1.300778in}{1.532005in}}%
\pgfpathlineto{\pgfqpoint{1.301865in}{1.532087in}}%
\pgfpathlineto{\pgfqpoint{1.302739in}{1.532115in}}%
\pgfpathlineto{\pgfqpoint{1.302748in}{1.532161in}}%
\pgfpathlineto{\pgfqpoint{1.302776in}{1.605465in}}%
\pgfpathlineto{\pgfqpoint{1.303561in}{2.219141in}}%
\pgfpathlineto{\pgfqpoint{1.304600in}{2.219131in}}%
\pgfpathlineto{\pgfqpoint{1.305094in}{2.217935in}}%
\pgfpathlineto{\pgfqpoint{1.305242in}{2.096819in}}%
\pgfpathlineto{\pgfqpoint{1.305735in}{1.532025in}}%
\pgfpathlineto{\pgfqpoint{1.306188in}{1.532060in}}%
\pgfpathlineto{\pgfqpoint{1.307687in}{1.532121in}}%
\pgfpathlineto{\pgfqpoint{1.307696in}{1.532182in}}%
\pgfpathlineto{\pgfqpoint{1.307736in}{1.871156in}}%
\pgfpathlineto{\pgfqpoint{1.308506in}{2.219020in}}%
\pgfpathlineto{\pgfqpoint{1.308902in}{2.219233in}}%
\pgfpathlineto{\pgfqpoint{1.309248in}{2.218779in}}%
\pgfpathlineto{\pgfqpoint{1.309298in}{2.219244in}}%
\pgfpathlineto{\pgfqpoint{1.309743in}{2.216213in}}%
\pgfpathlineto{\pgfqpoint{1.309792in}{2.218919in}}%
\pgfpathlineto{\pgfqpoint{1.310699in}{1.532033in}}%
\pgfpathlineto{\pgfqpoint{1.311019in}{1.532050in}}%
\pgfpathlineto{\pgfqpoint{1.312640in}{1.532170in}}%
\pgfpathlineto{\pgfqpoint{1.312682in}{1.871527in}}%
\pgfpathlineto{\pgfqpoint{1.313446in}{2.218771in}}%
\pgfpathlineto{\pgfqpoint{1.313842in}{2.218984in}}%
\pgfpathlineto{\pgfqpoint{1.314188in}{2.218626in}}%
\pgfpathlineto{\pgfqpoint{1.314237in}{2.218993in}}%
\pgfpathlineto{\pgfqpoint{1.314682in}{2.217576in}}%
\pgfpathlineto{\pgfqpoint{1.314732in}{2.218657in}}%
\pgfpathlineto{\pgfqpoint{1.315276in}{1.800435in}}%
\pgfpathlineto{\pgfqpoint{1.315614in}{1.532016in}}%
\pgfpathlineto{\pgfqpoint{1.315591in}{1.889269in}}%
\pgfpathlineto{\pgfqpoint{1.316190in}{1.532066in}}%
\pgfpathlineto{\pgfqpoint{1.317577in}{1.532118in}}%
\pgfpathlineto{\pgfqpoint{1.317588in}{1.532197in}}%
\pgfpathlineto{\pgfqpoint{1.318815in}{2.218769in}}%
\pgfpathlineto{\pgfqpoint{1.318864in}{2.218390in}}%
\pgfpathlineto{\pgfqpoint{1.319111in}{2.218786in}}%
\pgfpathlineto{\pgfqpoint{1.319457in}{2.217638in}}%
\pgfpathlineto{\pgfqpoint{1.319507in}{2.218612in}}%
\pgfpathlineto{\pgfqpoint{1.319556in}{2.216939in}}%
\pgfpathlineto{\pgfqpoint{1.319606in}{2.218527in}}%
\pgfpathlineto{\pgfqpoint{1.320575in}{1.532024in}}%
\pgfpathlineto{\pgfqpoint{1.321056in}{1.532064in}}%
\pgfpathlineto{\pgfqpoint{1.322522in}{1.532117in}}%
\pgfpathlineto{\pgfqpoint{1.322532in}{1.532182in}}%
\pgfpathlineto{\pgfqpoint{1.323903in}{2.218495in}}%
\pgfpathlineto{\pgfqpoint{1.324645in}{2.216936in}}%
\pgfpathlineto{\pgfqpoint{1.324695in}{2.218009in}}%
\pgfpathlineto{\pgfqpoint{1.325562in}{1.532025in}}%
\pgfpathlineto{\pgfqpoint{1.326593in}{1.532085in}}%
\pgfpathlineto{\pgfqpoint{1.327468in}{1.532119in}}%
\pgfpathlineto{\pgfqpoint{1.327478in}{1.532198in}}%
\pgfpathlineto{\pgfqpoint{1.328705in}{2.218297in}}%
\pgfpathlineto{\pgfqpoint{1.328754in}{2.217928in}}%
\pgfpathlineto{\pgfqpoint{1.329002in}{2.218321in}}%
\pgfpathlineto{\pgfqpoint{1.329348in}{2.217241in}}%
\pgfpathlineto{\pgfqpoint{1.329397in}{2.218158in}}%
\pgfpathlineto{\pgfqpoint{1.329447in}{2.216693in}}%
\pgfpathlineto{\pgfqpoint{1.329496in}{2.218076in}}%
\pgfpathlineto{\pgfqpoint{1.330478in}{1.532033in}}%
\pgfpathlineto{\pgfqpoint{1.331287in}{1.532067in}}%
\pgfpathlineto{\pgfqpoint{1.332421in}{1.532158in}}%
\pgfpathlineto{\pgfqpoint{1.332447in}{1.598324in}}%
\pgfpathlineto{\pgfqpoint{1.333194in}{2.217755in}}%
\pgfpathlineto{\pgfqpoint{1.333243in}{2.217679in}}%
\pgfpathlineto{\pgfqpoint{1.333985in}{2.217987in}}%
\pgfpathlineto{\pgfqpoint{1.334084in}{2.217953in}}%
\pgfpathlineto{\pgfqpoint{1.334628in}{2.216240in}}%
\pgfpathlineto{\pgfqpoint{1.334677in}{2.217302in}}%
\pgfpathlineto{\pgfqpoint{1.335452in}{1.532024in}}%
\pgfpathlineto{\pgfqpoint{1.336186in}{1.532074in}}%
\pgfpathlineto{\pgfqpoint{1.337359in}{1.532118in}}%
\pgfpathlineto{\pgfqpoint{1.337369in}{1.532198in}}%
\pgfpathlineto{\pgfqpoint{1.338597in}{2.217840in}}%
\pgfpathlineto{\pgfqpoint{1.338646in}{2.217460in}}%
\pgfpathlineto{\pgfqpoint{1.338893in}{2.217862in}}%
\pgfpathlineto{\pgfqpoint{1.339240in}{2.216735in}}%
\pgfpathlineto{\pgfqpoint{1.339289in}{2.217698in}}%
\pgfpathlineto{\pgfqpoint{1.339339in}{2.216100in}}%
\pgfpathlineto{\pgfqpoint{1.339388in}{2.217615in}}%
\pgfpathlineto{\pgfqpoint{1.340432in}{1.532035in}}%
\pgfpathlineto{\pgfqpoint{1.340924in}{1.532053in}}%
\pgfpathlineto{\pgfqpoint{1.342314in}{1.532184in}}%
\pgfpathlineto{\pgfqpoint{1.343622in}{2.217608in}}%
\pgfpathlineto{\pgfqpoint{1.343672in}{2.217272in}}%
\pgfpathlineto{\pgfqpoint{1.343820in}{2.217612in}}%
\pgfpathlineto{\pgfqpoint{1.344265in}{2.216481in}}%
\pgfpathlineto{\pgfqpoint{1.344315in}{2.217366in}}%
\pgfpathlineto{\pgfqpoint{1.344364in}{2.215781in}}%
\pgfpathlineto{\pgfqpoint{1.344414in}{2.217261in}}%
\pgfpathlineto{\pgfqpoint{1.345316in}{1.532033in}}%
\pgfpathlineto{\pgfqpoint{1.345880in}{1.532052in}}%
\pgfpathlineto{\pgfqpoint{1.347265in}{1.532498in}}%
\pgfpathlineto{\pgfqpoint{1.348315in}{2.217202in}}%
\pgfpathlineto{\pgfqpoint{1.349502in}{2.216390in}}%
\pgfpathlineto{\pgfqpoint{1.349650in}{2.208724in}}%
\pgfpathlineto{\pgfqpoint{1.349700in}{2.204943in}}%
\pgfpathlineto{\pgfqpoint{1.350249in}{1.532024in}}%
\pgfpathlineto{\pgfqpoint{1.350882in}{1.532046in}}%
\pgfpathlineto{\pgfqpoint{1.352205in}{1.532187in}}%
\pgfpathlineto{\pgfqpoint{1.353504in}{2.217169in}}%
\pgfpathlineto{\pgfqpoint{1.353554in}{2.216870in}}%
\pgfpathlineto{\pgfqpoint{1.353702in}{2.217177in}}%
\pgfpathlineto{\pgfqpoint{1.354147in}{2.216211in}}%
\pgfpathlineto{\pgfqpoint{1.354197in}{2.216938in}}%
\pgfpathlineto{\pgfqpoint{1.354345in}{2.203239in}}%
\pgfpathlineto{\pgfqpoint{1.354395in}{2.216681in}}%
\pgfpathlineto{\pgfqpoint{1.355238in}{1.532031in}}%
\pgfpathlineto{\pgfqpoint{1.355677in}{1.532057in}}%
\pgfpathlineto{\pgfqpoint{1.357140in}{1.532117in}}%
\pgfpathlineto{\pgfqpoint{1.357150in}{1.532183in}}%
\pgfpathlineto{\pgfqpoint{1.358486in}{2.216841in}}%
\pgfpathlineto{\pgfqpoint{1.358635in}{2.216805in}}%
\pgfpathlineto{\pgfqpoint{1.359277in}{2.216219in}}%
\pgfpathlineto{\pgfqpoint{1.359574in}{2.212565in}}%
\pgfpathlineto{\pgfqpoint{1.360123in}{1.532006in}}%
\pgfpathlineto{\pgfqpoint{1.360854in}{1.532048in}}%
\pgfpathlineto{\pgfqpoint{1.362094in}{1.532166in}}%
\pgfpathlineto{\pgfqpoint{1.362135in}{1.828477in}}%
\pgfpathlineto{\pgfqpoint{1.362899in}{2.216509in}}%
\pgfpathlineto{\pgfqpoint{1.363295in}{2.216729in}}%
\pgfpathlineto{\pgfqpoint{1.363641in}{2.216386in}}%
\pgfpathlineto{\pgfqpoint{1.363691in}{2.216747in}}%
\pgfpathlineto{\pgfqpoint{1.364136in}{2.215435in}}%
\pgfpathlineto{\pgfqpoint{1.364185in}{2.216428in}}%
\pgfpathlineto{\pgfqpoint{1.364581in}{2.130856in}}%
\pgfpathlineto{\pgfqpoint{1.365068in}{1.532007in}}%
\pgfpathlineto{\pgfqpoint{1.365752in}{1.532071in}}%
\pgfpathlineto{\pgfqpoint{1.367030in}{1.532115in}}%
\pgfpathlineto{\pgfqpoint{1.367040in}{1.532163in}}%
\pgfpathlineto{\pgfqpoint{1.367068in}{1.612872in}}%
\pgfpathlineto{\pgfqpoint{1.367843in}{2.216289in}}%
\pgfpathlineto{\pgfqpoint{1.368239in}{2.216500in}}%
\pgfpathlineto{\pgfqpoint{1.368585in}{2.216157in}}%
\pgfpathlineto{\pgfqpoint{1.368635in}{2.216506in}}%
\pgfpathlineto{\pgfqpoint{1.369080in}{2.215177in}}%
\pgfpathlineto{\pgfqpoint{1.369129in}{2.216166in}}%
\pgfpathlineto{\pgfqpoint{1.369574in}{2.017108in}}%
\pgfpathlineto{\pgfqpoint{1.370029in}{1.532024in}}%
\pgfpathlineto{\pgfqpoint{1.370558in}{1.532042in}}%
\pgfpathlineto{\pgfqpoint{1.371987in}{1.532183in}}%
\pgfpathlineto{\pgfqpoint{1.373280in}{2.216293in}}%
\pgfpathlineto{\pgfqpoint{1.373329in}{2.216013in}}%
\pgfpathlineto{\pgfqpoint{1.373478in}{2.216298in}}%
\pgfpathlineto{\pgfqpoint{1.373923in}{2.215379in}}%
\pgfpathlineto{\pgfqpoint{1.373972in}{2.216048in}}%
\pgfpathlineto{\pgfqpoint{1.374121in}{2.212231in}}%
\pgfpathlineto{\pgfqpoint{1.374170in}{2.215806in}}%
\pgfpathlineto{\pgfqpoint{1.375014in}{1.532025in}}%
\pgfpathlineto{\pgfqpoint{1.375401in}{1.532039in}}%
\pgfpathlineto{\pgfqpoint{1.376934in}{1.532225in}}%
\pgfpathlineto{\pgfqpoint{1.378057in}{2.216020in}}%
\pgfpathlineto{\pgfqpoint{1.378996in}{2.215070in}}%
\pgfpathlineto{\pgfqpoint{1.379046in}{2.215648in}}%
\pgfpathlineto{\pgfqpoint{1.379095in}{2.214554in}}%
\pgfpathlineto{\pgfqpoint{1.379145in}{2.215494in}}%
\pgfpathlineto{\pgfqpoint{1.379942in}{1.532036in}}%
\pgfpathlineto{\pgfqpoint{1.381206in}{1.532082in}}%
\pgfpathlineto{\pgfqpoint{1.381883in}{1.532495in}}%
\pgfpathlineto{\pgfqpoint{1.382935in}{2.215740in}}%
\pgfpathlineto{\pgfqpoint{1.384122in}{2.214976in}}%
\pgfpathlineto{\pgfqpoint{1.384271in}{2.207193in}}%
\pgfpathlineto{\pgfqpoint{1.384320in}{2.212927in}}%
\pgfpathlineto{\pgfqpoint{1.384893in}{1.532030in}}%
\pgfpathlineto{\pgfqpoint{1.385572in}{1.532054in}}%
\pgfpathlineto{\pgfqpoint{1.386825in}{1.532222in}}%
\pgfpathlineto{\pgfqpoint{1.387956in}{2.215753in}}%
\pgfpathlineto{\pgfqpoint{1.388005in}{2.215275in}}%
\pgfpathlineto{\pgfqpoint{1.388351in}{2.215830in}}%
\pgfpathlineto{\pgfqpoint{1.388599in}{2.214526in}}%
\pgfpathlineto{\pgfqpoint{1.388648in}{2.215753in}}%
\pgfpathlineto{\pgfqpoint{1.388698in}{2.213473in}}%
\pgfpathlineto{\pgfqpoint{1.388747in}{2.215698in}}%
\pgfpathlineto{\pgfqpoint{1.389841in}{1.532029in}}%
\pgfpathlineto{\pgfqpoint{1.390027in}{1.532035in}}%
\pgfpathlineto{\pgfqpoint{1.391771in}{1.532235in}}%
\pgfpathlineto{\pgfqpoint{1.392905in}{2.215535in}}%
\pgfpathlineto{\pgfqpoint{1.393102in}{2.215588in}}%
\pgfpathlineto{\pgfqpoint{1.393745in}{2.214167in}}%
\pgfpathlineto{\pgfqpoint{1.393795in}{2.215343in}}%
\pgfpathlineto{\pgfqpoint{1.393844in}{2.200407in}}%
\pgfpathlineto{\pgfqpoint{1.393894in}{2.215218in}}%
\pgfpathlineto{\pgfqpoint{1.394740in}{1.532008in}}%
\pgfpathlineto{\pgfqpoint{1.395192in}{1.532036in}}%
\pgfpathlineto{\pgfqpoint{1.396721in}{1.532974in}}%
\pgfpathlineto{\pgfqpoint{1.397699in}{2.215212in}}%
\pgfpathlineto{\pgfqpoint{1.397748in}{2.215114in}}%
\pgfpathlineto{\pgfqpoint{1.398292in}{2.215309in}}%
\pgfpathlineto{\pgfqpoint{1.398490in}{2.215225in}}%
\pgfpathlineto{\pgfqpoint{1.398935in}{2.213689in}}%
\pgfpathlineto{\pgfqpoint{1.398984in}{2.214625in}}%
\pgfpathlineto{\pgfqpoint{1.399686in}{1.532008in}}%
\pgfpathlineto{\pgfqpoint{1.400667in}{1.532083in}}%
\pgfpathlineto{\pgfqpoint{1.401648in}{1.532115in}}%
\pgfpathlineto{\pgfqpoint{1.401658in}{1.532163in}}%
\pgfpathlineto{\pgfqpoint{1.401686in}{1.610545in}}%
\pgfpathlineto{\pgfqpoint{1.402458in}{2.214921in}}%
\pgfpathlineto{\pgfqpoint{1.402755in}{2.215101in}}%
\pgfpathlineto{\pgfqpoint{1.403200in}{2.214870in}}%
\pgfpathlineto{\pgfqpoint{1.403250in}{2.215149in}}%
\pgfpathlineto{\pgfqpoint{1.403695in}{2.214143in}}%
\pgfpathlineto{\pgfqpoint{1.403744in}{2.214811in}}%
\pgfpathlineto{\pgfqpoint{1.403794in}{2.213606in}}%
\pgfpathlineto{\pgfqpoint{1.403843in}{2.214676in}}%
\pgfpathlineto{\pgfqpoint{1.404631in}{1.532006in}}%
\pgfpathlineto{\pgfqpoint{1.405755in}{1.532062in}}%
\pgfpathlineto{\pgfqpoint{1.406604in}{1.532167in}}%
\pgfpathlineto{\pgfqpoint{1.406644in}{1.835371in}}%
\pgfpathlineto{\pgfqpoint{1.407409in}{2.214757in}}%
\pgfpathlineto{\pgfqpoint{1.407705in}{2.214937in}}%
\pgfpathlineto{\pgfqpoint{1.408150in}{2.214694in}}%
\pgfpathlineto{\pgfqpoint{1.408200in}{2.214985in}}%
\pgfpathlineto{\pgfqpoint{1.408645in}{2.213939in}}%
\pgfpathlineto{\pgfqpoint{1.408694in}{2.214651in}}%
\pgfpathlineto{\pgfqpoint{1.408744in}{2.213338in}}%
\pgfpathlineto{\pgfqpoint{1.408793in}{2.214517in}}%
\pgfpathlineto{\pgfqpoint{1.409577in}{1.532005in}}%
\pgfpathlineto{\pgfqpoint{1.410473in}{1.532081in}}%
\pgfpathlineto{\pgfqpoint{1.411539in}{1.532114in}}%
\pgfpathlineto{\pgfqpoint{1.411548in}{1.532152in}}%
\pgfpathlineto{\pgfqpoint{1.411565in}{1.553297in}}%
\pgfpathlineto{\pgfqpoint{1.412439in}{2.214608in}}%
\pgfpathlineto{\pgfqpoint{1.413379in}{2.214373in}}%
\pgfpathlineto{\pgfqpoint{1.413428in}{2.214593in}}%
\pgfpathlineto{\pgfqpoint{1.413775in}{2.213228in}}%
\pgfpathlineto{\pgfqpoint{1.413824in}{2.214034in}}%
\pgfpathlineto{\pgfqpoint{1.414554in}{1.532032in}}%
\pgfpathlineto{\pgfqpoint{1.415714in}{1.532074in}}%
\pgfpathlineto{\pgfqpoint{1.416501in}{1.532489in}}%
\pgfpathlineto{\pgfqpoint{1.417534in}{2.214456in}}%
\pgfpathlineto{\pgfqpoint{1.418572in}{2.214000in}}%
\pgfpathlineto{\pgfqpoint{1.418869in}{2.212788in}}%
\pgfpathlineto{\pgfqpoint{1.418918in}{2.211939in}}%
\pgfpathlineto{\pgfqpoint{1.419517in}{1.532028in}}%
\pgfpathlineto{\pgfqpoint{1.420446in}{1.532077in}}%
\pgfpathlineto{\pgfqpoint{1.421432in}{1.532122in}}%
\pgfpathlineto{\pgfqpoint{1.421443in}{1.532219in}}%
\pgfpathlineto{\pgfqpoint{1.422574in}{2.214501in}}%
\pgfpathlineto{\pgfqpoint{1.422623in}{2.214004in}}%
\pgfpathlineto{\pgfqpoint{1.422970in}{2.214582in}}%
\pgfpathlineto{\pgfqpoint{1.423217in}{2.213198in}}%
\pgfpathlineto{\pgfqpoint{1.423266in}{2.214509in}}%
\pgfpathlineto{\pgfqpoint{1.423316in}{2.211180in}}%
\pgfpathlineto{\pgfqpoint{1.423365in}{2.214454in}}%
\pgfpathlineto{\pgfqpoint{1.424461in}{1.532028in}}%
\pgfpathlineto{\pgfqpoint{1.424600in}{1.532047in}}%
\pgfpathlineto{\pgfqpoint{1.426378in}{1.532121in}}%
\pgfpathlineto{\pgfqpoint{1.426390in}{1.532251in}}%
\pgfpathlineto{\pgfqpoint{1.427498in}{2.214194in}}%
\pgfpathlineto{\pgfqpoint{1.428635in}{2.212875in}}%
\pgfpathlineto{\pgfqpoint{1.428685in}{2.213387in}}%
\pgfpathlineto{\pgfqpoint{1.428932in}{2.017576in}}%
\pgfpathlineto{\pgfqpoint{1.429414in}{1.532035in}}%
\pgfpathlineto{\pgfqpoint{1.429901in}{1.532052in}}%
\pgfpathlineto{\pgfqpoint{1.431335in}{1.532258in}}%
\pgfpathlineto{\pgfqpoint{1.432428in}{2.214068in}}%
\pgfpathlineto{\pgfqpoint{1.432477in}{2.213850in}}%
\pgfpathlineto{\pgfqpoint{1.432923in}{2.214121in}}%
\pgfpathlineto{\pgfqpoint{1.433170in}{2.213471in}}%
\pgfpathlineto{\pgfqpoint{1.433219in}{2.213982in}}%
\pgfpathlineto{\pgfqpoint{1.433467in}{2.211838in}}%
\pgfpathlineto{\pgfqpoint{1.433516in}{2.213671in}}%
\pgfpathlineto{\pgfqpoint{1.434304in}{1.532006in}}%
\pgfpathlineto{\pgfqpoint{1.434830in}{1.532040in}}%
\pgfpathlineto{\pgfqpoint{1.436277in}{1.532172in}}%
\pgfpathlineto{\pgfqpoint{1.436318in}{1.864174in}}%
\pgfpathlineto{\pgfqpoint{1.437083in}{2.213729in}}%
\pgfpathlineto{\pgfqpoint{1.437478in}{2.213949in}}%
\pgfpathlineto{\pgfqpoint{1.437825in}{2.213593in}}%
\pgfpathlineto{\pgfqpoint{1.437874in}{2.213968in}}%
\pgfpathlineto{\pgfqpoint{1.438319in}{2.212607in}}%
\pgfpathlineto{\pgfqpoint{1.438369in}{2.213656in}}%
\pgfpathlineto{\pgfqpoint{1.438863in}{1.966785in}}%
\pgfpathlineto{\pgfqpoint{1.439249in}{1.532003in}}%
\pgfpathlineto{\pgfqpoint{1.439801in}{1.532039in}}%
\pgfpathlineto{\pgfqpoint{1.441228in}{1.532429in}}%
\pgfpathlineto{\pgfqpoint{1.442260in}{2.213622in}}%
\pgfpathlineto{\pgfqpoint{1.442309in}{2.213590in}}%
\pgfpathlineto{\pgfqpoint{1.442952in}{2.213598in}}%
\pgfpathlineto{\pgfqpoint{1.443051in}{2.213529in}}%
\pgfpathlineto{\pgfqpoint{1.443496in}{2.212350in}}%
\pgfpathlineto{\pgfqpoint{1.443546in}{2.212669in}}%
\pgfpathlineto{\pgfqpoint{1.443744in}{2.061378in}}%
\pgfpathlineto{\pgfqpoint{1.444247in}{1.532026in}}%
\pgfpathlineto{\pgfqpoint{1.444732in}{1.532044in}}%
\pgfpathlineto{\pgfqpoint{1.446171in}{1.532226in}}%
\pgfpathlineto{\pgfqpoint{1.447304in}{2.213612in}}%
\pgfpathlineto{\pgfqpoint{1.447502in}{2.213669in}}%
\pgfpathlineto{\pgfqpoint{1.448145in}{2.212337in}}%
\pgfpathlineto{\pgfqpoint{1.448194in}{2.213441in}}%
\pgfpathlineto{\pgfqpoint{1.448244in}{2.208662in}}%
\pgfpathlineto{\pgfqpoint{1.448293in}{2.213340in}}%
\pgfpathlineto{\pgfqpoint{1.449195in}{1.532027in}}%
\pgfpathlineto{\pgfqpoint{1.449534in}{1.532056in}}%
\pgfpathlineto{\pgfqpoint{1.451105in}{1.532122in}}%
\pgfpathlineto{\pgfqpoint{1.451116in}{1.532220in}}%
\pgfpathlineto{\pgfqpoint{1.452246in}{2.213511in}}%
\pgfpathlineto{\pgfqpoint{1.452296in}{2.213035in}}%
\pgfpathlineto{\pgfqpoint{1.452642in}{2.213591in}}%
\pgfpathlineto{\pgfqpoint{1.452889in}{2.212314in}}%
\pgfpathlineto{\pgfqpoint{1.452939in}{2.213518in}}%
\pgfpathlineto{\pgfqpoint{1.452988in}{2.211415in}}%
\pgfpathlineto{\pgfqpoint{1.453038in}{2.213464in}}%
\pgfpathlineto{\pgfqpoint{1.454131in}{1.532034in}}%
\pgfpathlineto{\pgfqpoint{1.454367in}{1.532045in}}%
\pgfpathlineto{\pgfqpoint{1.456059in}{1.532174in}}%
\pgfpathlineto{\pgfqpoint{1.456101in}{1.907592in}}%
\pgfpathlineto{\pgfqpoint{1.456863in}{2.213090in}}%
\pgfpathlineto{\pgfqpoint{1.457259in}{2.213308in}}%
\pgfpathlineto{\pgfqpoint{1.457605in}{2.213015in}}%
\pgfpathlineto{\pgfqpoint{1.457655in}{2.213324in}}%
\pgfpathlineto{\pgfqpoint{1.458100in}{2.212236in}}%
\pgfpathlineto{\pgfqpoint{1.458149in}{2.213000in}}%
\pgfpathlineto{\pgfqpoint{1.458199in}{2.211558in}}%
\pgfpathlineto{\pgfqpoint{1.458248in}{2.212870in}}%
\pgfpathlineto{\pgfqpoint{1.459031in}{1.532007in}}%
\pgfpathlineto{\pgfqpoint{1.459780in}{1.532047in}}%
\pgfpathlineto{\pgfqpoint{1.461012in}{1.532941in}}%
\pgfpathlineto{\pgfqpoint{1.461997in}{2.213023in}}%
\pgfpathlineto{\pgfqpoint{1.462937in}{2.212583in}}%
\pgfpathlineto{\pgfqpoint{1.462986in}{2.212869in}}%
\pgfpathlineto{\pgfqpoint{1.463234in}{2.211533in}}%
\pgfpathlineto{\pgfqpoint{1.463283in}{2.212418in}}%
\pgfpathlineto{\pgfqpoint{1.463977in}{1.532005in}}%
\pgfpathlineto{\pgfqpoint{1.465072in}{1.532088in}}%
\pgfpathlineto{\pgfqpoint{1.465939in}{1.532114in}}%
\pgfpathlineto{\pgfqpoint{1.465948in}{1.532152in}}%
\pgfpathlineto{\pgfqpoint{1.465965in}{1.553123in}}%
\pgfpathlineto{\pgfqpoint{1.466852in}{2.212804in}}%
\pgfpathlineto{\pgfqpoint{1.467792in}{2.212581in}}%
\pgfpathlineto{\pgfqpoint{1.467841in}{2.212781in}}%
\pgfpathlineto{\pgfqpoint{1.468188in}{2.211502in}}%
\pgfpathlineto{\pgfqpoint{1.468237in}{2.212210in}}%
\pgfpathlineto{\pgfqpoint{1.468912in}{1.571144in}}%
\pgfpathlineto{\pgfqpoint{1.468922in}{1.532004in}}%
\pgfpathlineto{\pgfqpoint{1.469671in}{1.532047in}}%
\pgfpathlineto{\pgfqpoint{1.470901in}{1.532431in}}%
\pgfpathlineto{\pgfqpoint{1.471932in}{2.212698in}}%
\pgfpathlineto{\pgfqpoint{1.471982in}{2.212668in}}%
\pgfpathlineto{\pgfqpoint{1.472723in}{2.212616in}}%
\pgfpathlineto{\pgfqpoint{1.472822in}{2.212528in}}%
\pgfpathlineto{\pgfqpoint{1.473267in}{2.209149in}}%
\pgfpathlineto{\pgfqpoint{1.473317in}{2.211245in}}%
\pgfpathlineto{\pgfqpoint{1.473940in}{1.532035in}}%
\pgfpathlineto{\pgfqpoint{1.474526in}{1.532056in}}%
\pgfpathlineto{\pgfqpoint{1.475844in}{1.532255in}}%
\pgfpathlineto{\pgfqpoint{1.476946in}{2.212655in}}%
\pgfpathlineto{\pgfqpoint{1.476995in}{2.212406in}}%
\pgfpathlineto{\pgfqpoint{1.477441in}{2.212714in}}%
\pgfpathlineto{\pgfqpoint{1.477688in}{2.211988in}}%
\pgfpathlineto{\pgfqpoint{1.477737in}{2.212583in}}%
\pgfpathlineto{\pgfqpoint{1.477985in}{2.198005in}}%
\pgfpathlineto{\pgfqpoint{1.478034in}{2.212261in}}%
\pgfpathlineto{\pgfqpoint{1.478854in}{1.532031in}}%
\pgfpathlineto{\pgfqpoint{1.479332in}{1.532048in}}%
\pgfpathlineto{\pgfqpoint{1.480786in}{1.532176in}}%
\pgfpathlineto{\pgfqpoint{1.480829in}{1.933861in}}%
\pgfpathlineto{\pgfqpoint{1.481594in}{2.212300in}}%
\pgfpathlineto{\pgfqpoint{1.481890in}{2.212479in}}%
\pgfpathlineto{\pgfqpoint{1.482335in}{2.212273in}}%
\pgfpathlineto{\pgfqpoint{1.482385in}{2.212522in}}%
\pgfpathlineto{\pgfqpoint{1.482830in}{2.211603in}}%
\pgfpathlineto{\pgfqpoint{1.482879in}{2.212177in}}%
\pgfpathlineto{\pgfqpoint{1.483028in}{2.200240in}}%
\pgfpathlineto{\pgfqpoint{1.483077in}{2.211826in}}%
\pgfpathlineto{\pgfqpoint{1.483808in}{1.532029in}}%
\pgfpathlineto{\pgfqpoint{1.484392in}{1.532051in}}%
\pgfpathlineto{\pgfqpoint{1.485734in}{1.532217in}}%
\pgfpathlineto{\pgfqpoint{1.486865in}{2.212450in}}%
\pgfpathlineto{\pgfqpoint{1.486915in}{2.211950in}}%
\pgfpathlineto{\pgfqpoint{1.487261in}{2.212536in}}%
\pgfpathlineto{\pgfqpoint{1.487508in}{2.211156in}}%
\pgfpathlineto{\pgfqpoint{1.487557in}{2.212468in}}%
\pgfpathlineto{\pgfqpoint{1.487607in}{2.209366in}}%
\pgfpathlineto{\pgfqpoint{1.487656in}{2.212415in}}%
\pgfpathlineto{\pgfqpoint{1.488738in}{1.532034in}}%
\pgfpathlineto{\pgfqpoint{1.488914in}{1.532039in}}%
\pgfpathlineto{\pgfqpoint{1.490678in}{1.532182in}}%
\pgfpathlineto{\pgfqpoint{1.492034in}{2.212197in}}%
\pgfpathlineto{\pgfqpoint{1.492182in}{2.212130in}}%
\pgfpathlineto{\pgfqpoint{1.492627in}{2.211918in}}%
\pgfpathlineto{\pgfqpoint{1.492726in}{2.211794in}}%
\pgfpathlineto{\pgfqpoint{1.493072in}{2.207779in}}%
\pgfpathlineto{\pgfqpoint{1.493122in}{2.210276in}}%
\pgfpathlineto{\pgfqpoint{1.493681in}{1.532027in}}%
\pgfpathlineto{\pgfqpoint{1.494345in}{1.532051in}}%
\pgfpathlineto{\pgfqpoint{1.495623in}{1.532182in}}%
\pgfpathlineto{\pgfqpoint{1.495664in}{1.869063in}}%
\pgfpathlineto{\pgfqpoint{1.496433in}{2.211939in}}%
\pgfpathlineto{\pgfqpoint{1.496927in}{2.212194in}}%
\pgfpathlineto{\pgfqpoint{1.497174in}{2.211766in}}%
\pgfpathlineto{\pgfqpoint{1.497224in}{2.212199in}}%
\pgfpathlineto{\pgfqpoint{1.497669in}{2.210587in}}%
\pgfpathlineto{\pgfqpoint{1.497718in}{2.211915in}}%
\pgfpathlineto{\pgfqpoint{1.498663in}{1.532034in}}%
\pgfpathlineto{\pgfqpoint{1.499499in}{1.532068in}}%
\pgfpathlineto{\pgfqpoint{1.500569in}{1.532183in}}%
\pgfpathlineto{\pgfqpoint{1.501975in}{2.212056in}}%
\pgfpathlineto{\pgfqpoint{1.502074in}{2.212054in}}%
\pgfpathlineto{\pgfqpoint{1.502717in}{2.209954in}}%
\pgfpathlineto{\pgfqpoint{1.502766in}{2.211588in}}%
\pgfpathlineto{\pgfqpoint{1.503576in}{1.532032in}}%
\pgfpathlineto{\pgfqpoint{1.504099in}{1.532057in}}%
\pgfpathlineto{\pgfqpoint{1.505511in}{1.532148in}}%
\pgfpathlineto{\pgfqpoint{1.505520in}{1.532496in}}%
\pgfpathlineto{\pgfqpoint{1.506574in}{2.211738in}}%
\pgfpathlineto{\pgfqpoint{1.507811in}{2.210491in}}%
\pgfpathlineto{\pgfqpoint{1.507860in}{2.210705in}}%
\pgfpathlineto{\pgfqpoint{1.508008in}{2.143006in}}%
\pgfpathlineto{\pgfqpoint{1.508485in}{1.532023in}}%
\pgfpathlineto{\pgfqpoint{1.509044in}{1.532064in}}%
\pgfpathlineto{\pgfqpoint{1.510450in}{1.532121in}}%
\pgfpathlineto{\pgfqpoint{1.510459in}{1.532182in}}%
\pgfpathlineto{\pgfqpoint{1.510500in}{1.871306in}}%
\pgfpathlineto{\pgfqpoint{1.511269in}{2.211552in}}%
\pgfpathlineto{\pgfqpoint{1.511763in}{2.211809in}}%
\pgfpathlineto{\pgfqpoint{1.512010in}{2.211389in}}%
\pgfpathlineto{\pgfqpoint{1.512060in}{2.211815in}}%
\pgfpathlineto{\pgfqpoint{1.512505in}{2.210274in}}%
\pgfpathlineto{\pgfqpoint{1.512554in}{2.211533in}}%
\pgfpathlineto{\pgfqpoint{1.513429in}{1.532022in}}%
\pgfpathlineto{\pgfqpoint{1.514574in}{1.532082in}}%
\pgfpathlineto{\pgfqpoint{1.515395in}{1.532119in}}%
\pgfpathlineto{\pgfqpoint{1.515405in}{1.532183in}}%
\pgfpathlineto{\pgfqpoint{1.516810in}{2.211674in}}%
\pgfpathlineto{\pgfqpoint{1.516909in}{2.211672in}}%
\pgfpathlineto{\pgfqpoint{1.517552in}{2.209759in}}%
\pgfpathlineto{\pgfqpoint{1.517602in}{2.211205in}}%
\pgfpathlineto{\pgfqpoint{1.518408in}{1.532031in}}%
\pgfpathlineto{\pgfqpoint{1.519024in}{1.532061in}}%
\pgfpathlineto{\pgfqpoint{1.520350in}{1.532176in}}%
\pgfpathlineto{\pgfqpoint{1.520392in}{1.924078in}}%
\pgfpathlineto{\pgfqpoint{1.521159in}{2.211285in}}%
\pgfpathlineto{\pgfqpoint{1.521456in}{2.211466in}}%
\pgfpathlineto{\pgfqpoint{1.521901in}{2.211246in}}%
\pgfpathlineto{\pgfqpoint{1.521950in}{2.211517in}}%
\pgfpathlineto{\pgfqpoint{1.522396in}{2.210555in}}%
\pgfpathlineto{\pgfqpoint{1.522445in}{2.211187in}}%
\pgfpathlineto{\pgfqpoint{1.522593in}{2.191397in}}%
\pgfpathlineto{\pgfqpoint{1.522643in}{2.210824in}}%
\pgfpathlineto{\pgfqpoint{1.523336in}{1.532024in}}%
\pgfpathlineto{\pgfqpoint{1.524014in}{1.532069in}}%
\pgfpathlineto{\pgfqpoint{1.525287in}{1.532120in}}%
\pgfpathlineto{\pgfqpoint{1.525298in}{1.532233in}}%
\pgfpathlineto{\pgfqpoint{1.526407in}{2.211325in}}%
\pgfpathlineto{\pgfqpoint{1.527545in}{2.209966in}}%
\pgfpathlineto{\pgfqpoint{1.527594in}{2.210633in}}%
\pgfpathlineto{\pgfqpoint{1.527941in}{1.861581in}}%
\pgfpathlineto{\pgfqpoint{1.528292in}{1.532032in}}%
\pgfpathlineto{\pgfqpoint{1.528850in}{1.532051in}}%
\pgfpathlineto{\pgfqpoint{1.530247in}{1.532496in}}%
\pgfpathlineto{\pgfqpoint{1.531302in}{2.211143in}}%
\pgfpathlineto{\pgfqpoint{1.532538in}{2.209895in}}%
\pgfpathlineto{\pgfqpoint{1.532588in}{2.210129in}}%
\pgfpathlineto{\pgfqpoint{1.532736in}{2.142894in}}%
\pgfpathlineto{\pgfqpoint{1.533213in}{1.532007in}}%
\pgfpathlineto{\pgfqpoint{1.533796in}{1.532068in}}%
\pgfpathlineto{\pgfqpoint{1.535175in}{1.532115in}}%
\pgfpathlineto{\pgfqpoint{1.535185in}{1.532163in}}%
\pgfpathlineto{\pgfqpoint{1.535214in}{1.611638in}}%
\pgfpathlineto{\pgfqpoint{1.535988in}{2.210917in}}%
\pgfpathlineto{\pgfqpoint{1.536383in}{2.211147in}}%
\pgfpathlineto{\pgfqpoint{1.536730in}{2.210847in}}%
\pgfpathlineto{\pgfqpoint{1.536779in}{2.211179in}}%
\pgfpathlineto{\pgfqpoint{1.537224in}{2.210089in}}%
\pgfpathlineto{\pgfqpoint{1.537274in}{2.210885in}}%
\pgfpathlineto{\pgfqpoint{1.537323in}{2.209441in}}%
\pgfpathlineto{\pgfqpoint{1.537372in}{2.210766in}}%
\pgfpathlineto{\pgfqpoint{1.538198in}{1.532031in}}%
\pgfpathlineto{\pgfqpoint{1.538925in}{1.532066in}}%
\pgfpathlineto{\pgfqpoint{1.540122in}{1.532118in}}%
\pgfpathlineto{\pgfqpoint{1.540133in}{1.532196in}}%
\pgfpathlineto{\pgfqpoint{1.541364in}{2.211099in}}%
\pgfpathlineto{\pgfqpoint{1.542106in}{2.208606in}}%
\pgfpathlineto{\pgfqpoint{1.541562in}{2.211138in}}%
\pgfpathlineto{\pgfqpoint{1.542156in}{2.210935in}}%
\pgfpathlineto{\pgfqpoint{1.543142in}{1.532031in}}%
\pgfpathlineto{\pgfqpoint{1.543420in}{1.532040in}}%
\pgfpathlineto{\pgfqpoint{1.545078in}{1.532183in}}%
\pgfpathlineto{\pgfqpoint{1.546377in}{2.210934in}}%
\pgfpathlineto{\pgfqpoint{1.546427in}{2.210624in}}%
\pgfpathlineto{\pgfqpoint{1.546674in}{2.210932in}}%
\pgfpathlineto{\pgfqpoint{1.547020in}{2.209981in}}%
\pgfpathlineto{\pgfqpoint{1.547070in}{2.210726in}}%
\pgfpathlineto{\pgfqpoint{1.547218in}{2.199476in}}%
\pgfpathlineto{\pgfqpoint{1.547267in}{2.210490in}}%
\pgfpathlineto{\pgfqpoint{1.548087in}{1.532031in}}%
\pgfpathlineto{\pgfqpoint{1.548563in}{1.532048in}}%
\pgfpathlineto{\pgfqpoint{1.550022in}{1.532169in}}%
\pgfpathlineto{\pgfqpoint{1.550064in}{1.876114in}}%
\pgfpathlineto{\pgfqpoint{1.550827in}{2.210554in}}%
\pgfpathlineto{\pgfqpoint{1.551222in}{2.210780in}}%
\pgfpathlineto{\pgfqpoint{1.551569in}{2.210473in}}%
\pgfpathlineto{\pgfqpoint{1.551618in}{2.210806in}}%
\pgfpathlineto{\pgfqpoint{1.552063in}{2.209685in}}%
\pgfpathlineto{\pgfqpoint{1.552113in}{2.210506in}}%
\pgfpathlineto{\pgfqpoint{1.552162in}{2.208956in}}%
\pgfpathlineto{\pgfqpoint{1.552212in}{2.210384in}}%
\pgfpathlineto{\pgfqpoint{1.553029in}{1.532032in}}%
\pgfpathlineto{\pgfqpoint{1.553700in}{1.532056in}}%
\pgfpathlineto{\pgfqpoint{1.554967in}{1.532165in}}%
\pgfpathlineto{\pgfqpoint{1.555009in}{1.855338in}}%
\pgfpathlineto{\pgfqpoint{1.555771in}{2.210445in}}%
\pgfpathlineto{\pgfqpoint{1.556167in}{2.210672in}}%
\pgfpathlineto{\pgfqpoint{1.556513in}{2.210377in}}%
\pgfpathlineto{\pgfqpoint{1.556563in}{2.210699in}}%
\pgfpathlineto{\pgfqpoint{1.557008in}{2.209624in}}%
\pgfpathlineto{\pgfqpoint{1.557057in}{2.210400in}}%
\pgfpathlineto{\pgfqpoint{1.557107in}{2.208995in}}%
\pgfpathlineto{\pgfqpoint{1.557156in}{2.210278in}}%
\pgfpathlineto{\pgfqpoint{1.557984in}{1.532034in}}%
\pgfpathlineto{\pgfqpoint{1.558764in}{1.532062in}}%
\pgfpathlineto{\pgfqpoint{1.559913in}{1.532175in}}%
\pgfpathlineto{\pgfqpoint{1.559956in}{1.912107in}}%
\pgfpathlineto{\pgfqpoint{1.560718in}{2.210341in}}%
\pgfpathlineto{\pgfqpoint{1.561113in}{2.210565in}}%
\pgfpathlineto{\pgfqpoint{1.561459in}{2.210303in}}%
\pgfpathlineto{\pgfqpoint{1.561509in}{2.210590in}}%
\pgfpathlineto{\pgfqpoint{1.561954in}{2.209619in}}%
\pgfpathlineto{\pgfqpoint{1.562003in}{2.210280in}}%
\pgfpathlineto{\pgfqpoint{1.562152in}{2.190437in}}%
\pgfpathlineto{\pgfqpoint{1.562201in}{2.209940in}}%
\pgfpathlineto{\pgfqpoint{1.562939in}{1.532027in}}%
\pgfpathlineto{\pgfqpoint{1.563575in}{1.532066in}}%
\pgfpathlineto{\pgfqpoint{1.564851in}{1.532122in}}%
\pgfpathlineto{\pgfqpoint{1.564861in}{1.532218in}}%
\pgfpathlineto{\pgfqpoint{1.565992in}{2.210530in}}%
\pgfpathlineto{\pgfqpoint{1.566041in}{2.210053in}}%
\pgfpathlineto{\pgfqpoint{1.566486in}{2.210610in}}%
\pgfpathlineto{\pgfqpoint{1.566635in}{2.209374in}}%
\pgfpathlineto{\pgfqpoint{1.566684in}{2.210553in}}%
\pgfpathlineto{\pgfqpoint{1.566734in}{2.208620in}}%
\pgfpathlineto{\pgfqpoint{1.566783in}{2.210502in}}%
\pgfpathlineto{\pgfqpoint{1.567887in}{1.532034in}}%
\pgfpathlineto{\pgfqpoint{1.568227in}{1.532048in}}%
\pgfpathlineto{\pgfqpoint{1.569807in}{1.532217in}}%
\pgfpathlineto{\pgfqpoint{1.571003in}{2.210361in}}%
\pgfpathlineto{\pgfqpoint{1.571052in}{2.210119in}}%
\pgfpathlineto{\pgfqpoint{1.571399in}{2.210385in}}%
\pgfpathlineto{\pgfqpoint{1.571646in}{2.209757in}}%
\pgfpathlineto{\pgfqpoint{1.571695in}{2.210255in}}%
\pgfpathlineto{\pgfqpoint{1.571943in}{2.208482in}}%
\pgfpathlineto{\pgfqpoint{1.571992in}{2.209958in}}%
\pgfpathlineto{\pgfqpoint{1.572791in}{1.532024in}}%
\pgfpathlineto{\pgfqpoint{1.573419in}{1.532047in}}%
\pgfpathlineto{\pgfqpoint{1.574753in}{1.532242in}}%
\pgfpathlineto{\pgfqpoint{1.575895in}{2.210133in}}%
\pgfpathlineto{\pgfqpoint{1.576834in}{2.209707in}}%
\pgfpathlineto{\pgfqpoint{1.577181in}{2.208102in}}%
\pgfpathlineto{\pgfqpoint{1.577378in}{1.879946in}}%
\pgfpathlineto{\pgfqpoint{1.577785in}{1.532033in}}%
\pgfpathlineto{\pgfqpoint{1.578273in}{1.532051in}}%
\pgfpathlineto{\pgfqpoint{1.579694in}{1.532164in}}%
\pgfpathlineto{\pgfqpoint{1.579736in}{1.848258in}}%
\pgfpathlineto{\pgfqpoint{1.580497in}{2.209904in}}%
\pgfpathlineto{\pgfqpoint{1.581436in}{2.209901in}}%
\pgfpathlineto{\pgfqpoint{1.581486in}{2.210063in}}%
\pgfpathlineto{\pgfqpoint{1.581931in}{2.208934in}}%
\pgfpathlineto{\pgfqpoint{1.581980in}{2.209470in}}%
\pgfpathlineto{\pgfqpoint{1.582030in}{2.207647in}}%
\pgfpathlineto{\pgfqpoint{1.582079in}{2.209208in}}%
\pgfpathlineto{\pgfqpoint{1.582665in}{1.532022in}}%
\pgfpathlineto{\pgfqpoint{1.583358in}{1.532052in}}%
\pgfpathlineto{\pgfqpoint{1.584641in}{1.532183in}}%
\pgfpathlineto{\pgfqpoint{1.586043in}{2.210111in}}%
\pgfpathlineto{\pgfqpoint{1.586142in}{2.210109in}}%
\pgfpathlineto{\pgfqpoint{1.586784in}{2.208700in}}%
\pgfpathlineto{\pgfqpoint{1.586834in}{2.209641in}}%
\pgfpathlineto{\pgfqpoint{1.587279in}{1.876294in}}%
\pgfpathlineto{\pgfqpoint{1.587664in}{1.532035in}}%
\pgfpathlineto{\pgfqpoint{1.588151in}{1.532052in}}%
\pgfpathlineto{\pgfqpoint{1.589589in}{1.532234in}}%
\pgfpathlineto{\pgfqpoint{1.590787in}{2.209997in}}%
\pgfpathlineto{\pgfqpoint{1.590886in}{2.210025in}}%
\pgfpathlineto{\pgfqpoint{1.591628in}{2.208870in}}%
\pgfpathlineto{\pgfqpoint{1.591677in}{2.209746in}}%
\pgfpathlineto{\pgfqpoint{1.591727in}{2.208046in}}%
\pgfpathlineto{\pgfqpoint{1.591776in}{2.209629in}}%
\pgfpathlineto{\pgfqpoint{1.592628in}{1.532035in}}%
\pgfpathlineto{\pgfqpoint{1.593165in}{1.532056in}}%
\pgfpathlineto{\pgfqpoint{1.594535in}{1.532239in}}%
\pgfpathlineto{\pgfqpoint{1.595720in}{2.209882in}}%
\pgfpathlineto{\pgfqpoint{1.595819in}{2.209910in}}%
\pgfpathlineto{\pgfqpoint{1.596560in}{2.208889in}}%
\pgfpathlineto{\pgfqpoint{1.596610in}{2.209624in}}%
\pgfpathlineto{\pgfqpoint{1.596659in}{2.208354in}}%
\pgfpathlineto{\pgfqpoint{1.596709in}{2.209506in}}%
\pgfpathlineto{\pgfqpoint{1.597550in}{1.532034in}}%
\pgfpathlineto{\pgfqpoint{1.598576in}{1.532075in}}%
\pgfpathlineto{\pgfqpoint{1.599477in}{1.532172in}}%
\pgfpathlineto{\pgfqpoint{1.599519in}{1.888804in}}%
\pgfpathlineto{\pgfqpoint{1.600281in}{2.209576in}}%
\pgfpathlineto{\pgfqpoint{1.600677in}{2.209807in}}%
\pgfpathlineto{\pgfqpoint{1.601023in}{2.209525in}}%
\pgfpathlineto{\pgfqpoint{1.601072in}{2.209840in}}%
\pgfpathlineto{\pgfqpoint{1.601518in}{2.208818in}}%
\pgfpathlineto{\pgfqpoint{1.601567in}{2.209550in}}%
\pgfpathlineto{\pgfqpoint{1.601616in}{2.208288in}}%
\pgfpathlineto{\pgfqpoint{1.601666in}{2.209432in}}%
\pgfpathlineto{\pgfqpoint{1.602522in}{1.532034in}}%
\pgfpathlineto{\pgfqpoint{1.603603in}{1.532074in}}%
\pgfpathlineto{\pgfqpoint{1.604423in}{1.532184in}}%
\pgfpathlineto{\pgfqpoint{1.605827in}{2.209763in}}%
\pgfpathlineto{\pgfqpoint{1.605925in}{2.209763in}}%
\pgfpathlineto{\pgfqpoint{1.606568in}{2.208260in}}%
\pgfpathlineto{\pgfqpoint{1.606618in}{2.209315in}}%
\pgfpathlineto{\pgfqpoint{1.607433in}{1.532031in}}%
\pgfpathlineto{\pgfqpoint{1.608751in}{1.532088in}}%
\pgfpathlineto{\pgfqpoint{1.609369in}{1.532186in}}%
\pgfpathlineto{\pgfqpoint{1.610670in}{2.209685in}}%
\pgfpathlineto{\pgfqpoint{1.610720in}{2.209372in}}%
\pgfpathlineto{\pgfqpoint{1.610967in}{2.209691in}}%
\pgfpathlineto{\pgfqpoint{1.611313in}{2.208763in}}%
\pgfpathlineto{\pgfqpoint{1.611363in}{2.209500in}}%
\pgfpathlineto{\pgfqpoint{1.611511in}{2.205504in}}%
\pgfpathlineto{\pgfqpoint{1.611561in}{2.209293in}}%
\pgfpathlineto{\pgfqpoint{1.612338in}{1.532021in}}%
\pgfpathlineto{\pgfqpoint{1.612802in}{1.532056in}}%
\pgfpathlineto{\pgfqpoint{1.614304in}{1.532119in}}%
\pgfpathlineto{\pgfqpoint{1.614314in}{1.532182in}}%
\pgfpathlineto{\pgfqpoint{1.615813in}{2.209597in}}%
\pgfpathlineto{\pgfqpoint{1.615863in}{2.209231in}}%
\pgfpathlineto{\pgfqpoint{1.615912in}{2.209584in}}%
\pgfpathlineto{\pgfqpoint{1.616357in}{2.208393in}}%
\pgfpathlineto{\pgfqpoint{1.616407in}{2.209287in}}%
\pgfpathlineto{\pgfqpoint{1.616456in}{2.207443in}}%
\pgfpathlineto{\pgfqpoint{1.616505in}{2.209166in}}%
\pgfpathlineto{\pgfqpoint{1.617316in}{1.532034in}}%
\pgfpathlineto{\pgfqpoint{1.617834in}{1.532056in}}%
\pgfpathlineto{\pgfqpoint{1.619257in}{1.532153in}}%
\pgfpathlineto{\pgfqpoint{1.619266in}{1.532617in}}%
\pgfpathlineto{\pgfqpoint{1.620312in}{2.209315in}}%
\pgfpathlineto{\pgfqpoint{1.621548in}{2.208052in}}%
\pgfpathlineto{\pgfqpoint{1.621598in}{2.208463in}}%
\pgfpathlineto{\pgfqpoint{1.621795in}{2.081887in}}%
\pgfpathlineto{\pgfqpoint{1.622277in}{1.532030in}}%
\pgfpathlineto{\pgfqpoint{1.622810in}{1.532060in}}%
\pgfpathlineto{\pgfqpoint{1.624195in}{1.532118in}}%
\pgfpathlineto{\pgfqpoint{1.624206in}{1.532197in}}%
\pgfpathlineto{\pgfqpoint{1.625438in}{2.209414in}}%
\pgfpathlineto{\pgfqpoint{1.626179in}{2.205844in}}%
\pgfpathlineto{\pgfqpoint{1.625635in}{2.209455in}}%
\pgfpathlineto{\pgfqpoint{1.626229in}{2.209262in}}%
\pgfpathlineto{\pgfqpoint{1.627224in}{1.532034in}}%
\pgfpathlineto{\pgfqpoint{1.627461in}{1.532046in}}%
\pgfpathlineto{\pgfqpoint{1.629149in}{1.532167in}}%
\pgfpathlineto{\pgfqpoint{1.629192in}{1.869172in}}%
\pgfpathlineto{\pgfqpoint{1.629953in}{2.209028in}}%
\pgfpathlineto{\pgfqpoint{1.630349in}{2.209260in}}%
\pgfpathlineto{\pgfqpoint{1.630695in}{2.208992in}}%
\pgfpathlineto{\pgfqpoint{1.630745in}{2.209294in}}%
\pgfpathlineto{\pgfqpoint{1.631289in}{2.207853in}}%
\pgfpathlineto{\pgfqpoint{1.631338in}{2.208889in}}%
\pgfpathlineto{\pgfqpoint{1.631882in}{1.723014in}}%
\pgfpathlineto{\pgfqpoint{1.632150in}{1.532035in}}%
\pgfpathlineto{\pgfqpoint{1.632101in}{1.756088in}}%
\pgfpathlineto{\pgfqpoint{1.632814in}{1.532057in}}%
\pgfpathlineto{\pgfqpoint{1.634096in}{1.532184in}}%
\pgfpathlineto{\pgfqpoint{1.635409in}{2.209239in}}%
\pgfpathlineto{\pgfqpoint{1.636151in}{2.207537in}}%
\pgfpathlineto{\pgfqpoint{1.635508in}{2.209253in}}%
\pgfpathlineto{\pgfqpoint{1.636201in}{2.208958in}}%
\pgfpathlineto{\pgfqpoint{1.637104in}{1.532032in}}%
\pgfpathlineto{\pgfqpoint{1.637826in}{1.532065in}}%
\pgfpathlineto{\pgfqpoint{1.639041in}{1.532172in}}%
\pgfpathlineto{\pgfqpoint{1.639083in}{1.889493in}}%
\pgfpathlineto{\pgfqpoint{1.639845in}{2.208870in}}%
\pgfpathlineto{\pgfqpoint{1.640240in}{2.209098in}}%
\pgfpathlineto{\pgfqpoint{1.640587in}{2.208814in}}%
\pgfpathlineto{\pgfqpoint{1.640636in}{2.209129in}}%
\pgfpathlineto{\pgfqpoint{1.641081in}{2.208096in}}%
\pgfpathlineto{\pgfqpoint{1.641131in}{2.208835in}}%
\pgfpathlineto{\pgfqpoint{1.641180in}{2.207548in}}%
\pgfpathlineto{\pgfqpoint{1.641229in}{2.208715in}}%
\pgfpathlineto{\pgfqpoint{1.642061in}{1.532034in}}%
\pgfpathlineto{\pgfqpoint{1.643092in}{1.532075in}}%
\pgfpathlineto{\pgfqpoint{1.643987in}{1.532182in}}%
\pgfpathlineto{\pgfqpoint{1.645486in}{2.209094in}}%
\pgfpathlineto{\pgfqpoint{1.645536in}{2.208704in}}%
\pgfpathlineto{\pgfqpoint{1.645585in}{2.209084in}}%
\pgfpathlineto{\pgfqpoint{1.646030in}{2.207830in}}%
\pgfpathlineto{\pgfqpoint{1.646080in}{2.208807in}}%
\pgfpathlineto{\pgfqpoint{1.646129in}{2.206536in}}%
\pgfpathlineto{\pgfqpoint{1.646179in}{2.208692in}}%
\pgfpathlineto{\pgfqpoint{1.647053in}{1.532035in}}%
\pgfpathlineto{\pgfqpoint{1.647449in}{1.532049in}}%
\pgfpathlineto{\pgfqpoint{1.648933in}{1.532193in}}%
\pgfpathlineto{\pgfqpoint{1.650235in}{2.208984in}}%
\pgfpathlineto{\pgfqpoint{1.651076in}{2.201520in}}%
\pgfpathlineto{\pgfqpoint{1.651126in}{2.208574in}}%
\pgfpathlineto{\pgfqpoint{1.651175in}{2.143305in}}%
\pgfpathlineto{\pgfqpoint{1.651422in}{2.164243in}}%
\pgfpathlineto{\pgfqpoint{1.651948in}{1.532033in}}%
\pgfpathlineto{\pgfqpoint{1.653023in}{1.532072in}}%
\pgfpathlineto{\pgfqpoint{1.653877in}{1.532170in}}%
\pgfpathlineto{\pgfqpoint{1.653919in}{1.880095in}}%
\pgfpathlineto{\pgfqpoint{1.654681in}{2.208628in}}%
\pgfpathlineto{\pgfqpoint{1.655077in}{2.208855in}}%
\pgfpathlineto{\pgfqpoint{1.655423in}{2.208562in}}%
\pgfpathlineto{\pgfqpoint{1.655472in}{2.208885in}}%
\pgfpathlineto{\pgfqpoint{1.655918in}{2.207823in}}%
\pgfpathlineto{\pgfqpoint{1.655967in}{2.208592in}}%
\pgfpathlineto{\pgfqpoint{1.656016in}{2.207224in}}%
\pgfpathlineto{\pgfqpoint{1.656066in}{2.208473in}}%
\pgfpathlineto{\pgfqpoint{1.656879in}{1.532034in}}%
\pgfpathlineto{\pgfqpoint{1.657742in}{1.532064in}}%
\pgfpathlineto{\pgfqpoint{1.658822in}{1.532166in}}%
\pgfpathlineto{\pgfqpoint{1.658865in}{1.870333in}}%
\pgfpathlineto{\pgfqpoint{1.659626in}{2.208549in}}%
\pgfpathlineto{\pgfqpoint{1.660022in}{2.208780in}}%
\pgfpathlineto{\pgfqpoint{1.660368in}{2.208516in}}%
\pgfpathlineto{\pgfqpoint{1.660417in}{2.208813in}}%
\pgfpathlineto{\pgfqpoint{1.660961in}{2.207400in}}%
\pgfpathlineto{\pgfqpoint{1.661011in}{2.208404in}}%
\pgfpathlineto{\pgfqpoint{1.661505in}{1.815456in}}%
\pgfpathlineto{\pgfqpoint{1.661862in}{1.532034in}}%
\pgfpathlineto{\pgfqpoint{1.662399in}{1.532057in}}%
\pgfpathlineto{\pgfqpoint{1.663769in}{1.532191in}}%
\pgfpathlineto{\pgfqpoint{1.665078in}{2.208761in}}%
\pgfpathlineto{\pgfqpoint{1.665820in}{2.207278in}}%
\pgfpathlineto{\pgfqpoint{1.665870in}{2.208484in}}%
\pgfpathlineto{\pgfqpoint{1.666826in}{1.532035in}}%
\pgfpathlineto{\pgfqpoint{1.668156in}{1.532086in}}%
\pgfpathlineto{\pgfqpoint{1.668715in}{1.532207in}}%
\pgfpathlineto{\pgfqpoint{1.669917in}{2.208629in}}%
\pgfpathlineto{\pgfqpoint{1.669967in}{2.208364in}}%
\pgfpathlineto{\pgfqpoint{1.670313in}{2.208660in}}%
\pgfpathlineto{\pgfqpoint{1.670560in}{2.207985in}}%
\pgfpathlineto{\pgfqpoint{1.670610in}{2.208538in}}%
\pgfpathlineto{\pgfqpoint{1.670857in}{2.206264in}}%
\pgfpathlineto{\pgfqpoint{1.670906in}{2.208254in}}%
\pgfpathlineto{\pgfqpoint{1.671683in}{1.532022in}}%
\pgfpathlineto{\pgfqpoint{1.672179in}{1.532045in}}%
\pgfpathlineto{\pgfqpoint{1.673660in}{1.532183in}}%
\pgfpathlineto{\pgfqpoint{1.675061in}{2.208561in}}%
\pgfpathlineto{\pgfqpoint{1.675160in}{2.208561in}}%
\pgfpathlineto{\pgfqpoint{1.675901in}{2.200911in}}%
\pgfpathlineto{\pgfqpoint{1.675951in}{2.207929in}}%
\pgfpathlineto{\pgfqpoint{1.676669in}{1.532031in}}%
\pgfpathlineto{\pgfqpoint{1.677196in}{1.532059in}}%
\pgfpathlineto{\pgfqpoint{1.678605in}{1.532188in}}%
\pgfpathlineto{\pgfqpoint{1.679908in}{2.208514in}}%
\pgfpathlineto{\pgfqpoint{1.679958in}{2.208198in}}%
\pgfpathlineto{\pgfqpoint{1.680205in}{2.208524in}}%
\pgfpathlineto{\pgfqpoint{1.680551in}{2.207603in}}%
\pgfpathlineto{\pgfqpoint{1.680600in}{2.208341in}}%
\pgfpathlineto{\pgfqpoint{1.680749in}{2.205355in}}%
\pgfpathlineto{\pgfqpoint{1.680798in}{2.208138in}}%
\pgfpathlineto{\pgfqpoint{1.681607in}{1.532035in}}%
\pgfpathlineto{\pgfqpoint{1.682027in}{1.532051in}}%
\pgfpathlineto{\pgfqpoint{1.683551in}{1.532200in}}%
\pgfpathlineto{\pgfqpoint{1.684754in}{2.208397in}}%
\pgfpathlineto{\pgfqpoint{1.684803in}{2.208125in}}%
\pgfpathlineto{\pgfqpoint{1.685150in}{2.208428in}}%
\pgfpathlineto{\pgfqpoint{1.685397in}{2.207734in}}%
\pgfpathlineto{\pgfqpoint{1.685446in}{2.208305in}}%
\pgfpathlineto{\pgfqpoint{1.685694in}{2.205474in}}%
\pgfpathlineto{\pgfqpoint{1.685743in}{2.208022in}}%
\pgfpathlineto{\pgfqpoint{1.686520in}{1.532026in}}%
\pgfpathlineto{\pgfqpoint{1.686985in}{1.532055in}}%
\pgfpathlineto{\pgfqpoint{1.688495in}{1.532176in}}%
\pgfpathlineto{\pgfqpoint{1.688538in}{1.920315in}}%
\pgfpathlineto{\pgfqpoint{1.689304in}{2.208067in}}%
\pgfpathlineto{\pgfqpoint{1.689601in}{2.208252in}}%
\pgfpathlineto{\pgfqpoint{1.690046in}{2.208048in}}%
\pgfpathlineto{\pgfqpoint{1.690095in}{2.208313in}}%
\pgfpathlineto{\pgfqpoint{1.690540in}{2.207410in}}%
\pgfpathlineto{\pgfqpoint{1.690590in}{2.207999in}}%
\pgfpathlineto{\pgfqpoint{1.690738in}{2.202856in}}%
\pgfpathlineto{\pgfqpoint{1.690788in}{2.207703in}}%
\pgfpathlineto{\pgfqpoint{1.691505in}{1.532032in}}%
\pgfpathlineto{\pgfqpoint{1.692030in}{1.532058in}}%
\pgfpathlineto{\pgfqpoint{1.693441in}{1.532179in}}%
\pgfpathlineto{\pgfqpoint{1.694931in}{2.208245in}}%
\pgfpathlineto{\pgfqpoint{1.694980in}{2.208049in}}%
\pgfpathlineto{\pgfqpoint{1.695030in}{2.208231in}}%
\pgfpathlineto{\pgfqpoint{1.695574in}{2.207302in}}%
\pgfpathlineto{\pgfqpoint{1.695623in}{2.207782in}}%
\pgfpathlineto{\pgfqpoint{1.695772in}{2.193396in}}%
\pgfpathlineto{\pgfqpoint{1.695821in}{2.207331in}}%
\pgfpathlineto{\pgfqpoint{1.696430in}{1.532029in}}%
\pgfpathlineto{\pgfqpoint{1.697113in}{1.532065in}}%
\pgfpathlineto{\pgfqpoint{1.698377in}{1.532118in}}%
\pgfpathlineto{\pgfqpoint{1.698388in}{1.532206in}}%
\pgfpathlineto{\pgfqpoint{1.699521in}{2.208224in}}%
\pgfpathlineto{\pgfqpoint{1.699570in}{2.207752in}}%
\pgfpathlineto{\pgfqpoint{1.700015in}{2.208306in}}%
\pgfpathlineto{\pgfqpoint{1.700164in}{2.207104in}}%
\pgfpathlineto{\pgfqpoint{1.700213in}{2.208251in}}%
\pgfpathlineto{\pgfqpoint{1.700262in}{2.206451in}}%
\pgfpathlineto{\pgfqpoint{1.700312in}{2.208200in}}%
\pgfpathlineto{\pgfqpoint{1.701395in}{1.532035in}}%
\pgfpathlineto{\pgfqpoint{1.701820in}{1.532050in}}%
\pgfpathlineto{\pgfqpoint{1.703336in}{1.532268in}}%
\pgfpathlineto{\pgfqpoint{1.704442in}{2.208085in}}%
\pgfpathlineto{\pgfqpoint{1.704491in}{2.207757in}}%
\pgfpathlineto{\pgfqpoint{1.704936in}{2.208163in}}%
\pgfpathlineto{\pgfqpoint{1.705183in}{2.207223in}}%
\pgfpathlineto{\pgfqpoint{1.705233in}{2.208053in}}%
\pgfpathlineto{\pgfqpoint{1.705381in}{2.205385in}}%
\pgfpathlineto{\pgfqpoint{1.705431in}{2.207899in}}%
\pgfpathlineto{\pgfqpoint{1.706403in}{1.532035in}}%
\pgfpathlineto{\pgfqpoint{1.706700in}{1.532046in}}%
\pgfpathlineto{\pgfqpoint{1.708278in}{1.532180in}}%
\pgfpathlineto{\pgfqpoint{1.709644in}{2.207970in}}%
\pgfpathlineto{\pgfqpoint{1.709792in}{2.207873in}}%
\pgfpathlineto{\pgfqpoint{1.710039in}{2.207879in}}%
\pgfpathlineto{\pgfqpoint{1.710287in}{2.207487in}}%
\pgfpathlineto{\pgfqpoint{1.710336in}{2.207617in}}%
\pgfpathlineto{\pgfqpoint{1.710583in}{2.206602in}}%
\pgfpathlineto{\pgfqpoint{1.710633in}{2.207021in}}%
\pgfpathlineto{\pgfqpoint{1.710880in}{1.972069in}}%
\pgfpathlineto{\pgfqpoint{1.711263in}{1.532026in}}%
\pgfpathlineto{\pgfqpoint{1.711842in}{1.532063in}}%
\pgfpathlineto{\pgfqpoint{1.713215in}{1.532123in}}%
\pgfpathlineto{\pgfqpoint{1.713225in}{1.532225in}}%
\pgfpathlineto{\pgfqpoint{1.714355in}{2.207988in}}%
\pgfpathlineto{\pgfqpoint{1.714404in}{2.207563in}}%
\pgfpathlineto{\pgfqpoint{1.714849in}{2.208068in}}%
\pgfpathlineto{\pgfqpoint{1.714997in}{2.207042in}}%
\pgfpathlineto{\pgfqpoint{1.715047in}{2.208010in}}%
\pgfpathlineto{\pgfqpoint{1.715195in}{2.198330in}}%
\pgfpathlineto{\pgfqpoint{1.715245in}{2.207894in}}%
\pgfpathlineto{\pgfqpoint{1.716284in}{1.532034in}}%
\pgfpathlineto{\pgfqpoint{1.716528in}{1.532046in}}%
\pgfpathlineto{\pgfqpoint{1.718170in}{1.532199in}}%
\pgfpathlineto{\pgfqpoint{1.719372in}{2.207872in}}%
\pgfpathlineto{\pgfqpoint{1.719422in}{2.207596in}}%
\pgfpathlineto{\pgfqpoint{1.719768in}{2.207900in}}%
\pgfpathlineto{\pgfqpoint{1.720015in}{2.207197in}}%
\pgfpathlineto{\pgfqpoint{1.720065in}{2.207777in}}%
\pgfpathlineto{\pgfqpoint{1.720312in}{2.203974in}}%
\pgfpathlineto{\pgfqpoint{1.720361in}{2.207493in}}%
\pgfpathlineto{\pgfqpoint{1.721182in}{1.532036in}}%
\pgfpathlineto{\pgfqpoint{1.721613in}{1.532050in}}%
\pgfpathlineto{\pgfqpoint{1.723119in}{1.532349in}}%
\pgfpathlineto{\pgfqpoint{1.724134in}{2.207745in}}%
\pgfpathlineto{\pgfqpoint{1.724184in}{2.207422in}}%
\pgfpathlineto{\pgfqpoint{1.724728in}{2.207887in}}%
\pgfpathlineto{\pgfqpoint{1.724876in}{2.207083in}}%
\pgfpathlineto{\pgfqpoint{1.724926in}{2.207833in}}%
\pgfpathlineto{\pgfqpoint{1.725173in}{2.190074in}}%
\pgfpathlineto{\pgfqpoint{1.725222in}{2.207621in}}%
\pgfpathlineto{\pgfqpoint{1.726111in}{1.532031in}}%
\pgfpathlineto{\pgfqpoint{1.726572in}{1.532047in}}%
\pgfpathlineto{\pgfqpoint{1.728065in}{1.532491in}}%
\pgfpathlineto{\pgfqpoint{1.729107in}{2.207558in}}%
\pgfpathlineto{\pgfqpoint{1.730294in}{2.206907in}}%
\pgfpathlineto{\pgfqpoint{1.730442in}{2.205967in}}%
\pgfpathlineto{\pgfqpoint{1.730492in}{2.206005in}}%
\pgfpathlineto{\pgfqpoint{1.731113in}{1.532033in}}%
\pgfpathlineto{\pgfqpoint{1.732146in}{1.532078in}}%
\pgfpathlineto{\pgfqpoint{1.733002in}{1.532149in}}%
\pgfpathlineto{\pgfqpoint{1.733011in}{1.532496in}}%
\pgfpathlineto{\pgfqpoint{1.734065in}{2.207502in}}%
\pgfpathlineto{\pgfqpoint{1.735302in}{2.206388in}}%
\pgfpathlineto{\pgfqpoint{1.735351in}{2.206660in}}%
\pgfpathlineto{\pgfqpoint{1.735400in}{2.205253in}}%
\pgfpathlineto{\pgfqpoint{1.735450in}{2.206222in}}%
\pgfpathlineto{\pgfqpoint{1.736017in}{1.532030in}}%
\pgfpathlineto{\pgfqpoint{1.736743in}{1.532066in}}%
\pgfpathlineto{\pgfqpoint{1.737940in}{1.532117in}}%
\pgfpathlineto{\pgfqpoint{1.737951in}{1.532197in}}%
\pgfpathlineto{\pgfqpoint{1.739183in}{2.207635in}}%
\pgfpathlineto{\pgfqpoint{1.739925in}{2.204261in}}%
\pgfpathlineto{\pgfqpoint{1.739381in}{2.207677in}}%
\pgfpathlineto{\pgfqpoint{1.739974in}{2.207491in}}%
\pgfpathlineto{\pgfqpoint{1.740920in}{1.532012in}}%
\pgfpathlineto{\pgfqpoint{1.741238in}{1.532053in}}%
\pgfpathlineto{\pgfqpoint{1.742888in}{1.532123in}}%
\pgfpathlineto{\pgfqpoint{1.742898in}{1.532224in}}%
\pgfpathlineto{\pgfqpoint{1.744027in}{2.207523in}}%
\pgfpathlineto{\pgfqpoint{1.744077in}{2.207079in}}%
\pgfpathlineto{\pgfqpoint{1.744522in}{2.207598in}}%
\pgfpathlineto{\pgfqpoint{1.744670in}{2.206498in}}%
\pgfpathlineto{\pgfqpoint{1.744720in}{2.207540in}}%
\pgfpathlineto{\pgfqpoint{1.744868in}{2.181450in}}%
\pgfpathlineto{\pgfqpoint{1.744918in}{2.207398in}}%
\pgfpathlineto{\pgfqpoint{1.745916in}{1.532033in}}%
\pgfpathlineto{\pgfqpoint{1.746352in}{1.532052in}}%
\pgfpathlineto{\pgfqpoint{1.747842in}{1.532196in}}%
\pgfpathlineto{\pgfqpoint{1.749156in}{2.207449in}}%
\pgfpathlineto{\pgfqpoint{1.749898in}{2.205848in}}%
\pgfpathlineto{\pgfqpoint{1.749948in}{2.207170in}}%
\pgfpathlineto{\pgfqpoint{1.750861in}{1.532034in}}%
\pgfpathlineto{\pgfqpoint{1.751792in}{1.532071in}}%
\pgfpathlineto{\pgfqpoint{1.752788in}{1.532202in}}%
\pgfpathlineto{\pgfqpoint{1.754002in}{2.207352in}}%
\pgfpathlineto{\pgfqpoint{1.754052in}{2.207035in}}%
\pgfpathlineto{\pgfqpoint{1.754398in}{2.207386in}}%
\pgfpathlineto{\pgfqpoint{1.754645in}{2.206581in}}%
\pgfpathlineto{\pgfqpoint{1.754695in}{2.207270in}}%
\pgfpathlineto{\pgfqpoint{1.754843in}{2.205759in}}%
\pgfpathlineto{\pgfqpoint{1.754893in}{2.207111in}}%
\pgfpathlineto{\pgfqpoint{1.755797in}{1.532034in}}%
\pgfpathlineto{\pgfqpoint{1.756671in}{1.532065in}}%
\pgfpathlineto{\pgfqpoint{1.757733in}{1.532197in}}%
\pgfpathlineto{\pgfqpoint{1.758937in}{2.207268in}}%
\pgfpathlineto{\pgfqpoint{1.758986in}{2.206991in}}%
\pgfpathlineto{\pgfqpoint{1.759332in}{2.207304in}}%
\pgfpathlineto{\pgfqpoint{1.759580in}{2.206608in}}%
\pgfpathlineto{\pgfqpoint{1.759629in}{2.207187in}}%
\pgfpathlineto{\pgfqpoint{1.759876in}{2.204537in}}%
\pgfpathlineto{\pgfqpoint{1.759926in}{2.206913in}}%
\pgfpathlineto{\pgfqpoint{1.760799in}{1.532035in}}%
\pgfpathlineto{\pgfqpoint{1.761195in}{1.532049in}}%
\pgfpathlineto{\pgfqpoint{1.762678in}{1.532193in}}%
\pgfpathlineto{\pgfqpoint{1.763982in}{2.207228in}}%
\pgfpathlineto{\pgfqpoint{1.764081in}{2.207243in}}%
\pgfpathlineto{\pgfqpoint{1.764822in}{2.202970in}}%
\pgfpathlineto{\pgfqpoint{1.764872in}{2.206845in}}%
\pgfpathlineto{\pgfqpoint{1.765696in}{1.532034in}}%
\pgfpathlineto{\pgfqpoint{1.766131in}{1.532052in}}%
\pgfpathlineto{\pgfqpoint{1.767624in}{1.532204in}}%
\pgfpathlineto{\pgfqpoint{1.768839in}{2.207142in}}%
\pgfpathlineto{\pgfqpoint{1.768889in}{2.206822in}}%
\pgfpathlineto{\pgfqpoint{1.769235in}{2.207175in}}%
\pgfpathlineto{\pgfqpoint{1.769482in}{2.206358in}}%
\pgfpathlineto{\pgfqpoint{1.769532in}{2.207059in}}%
\pgfpathlineto{\pgfqpoint{1.769680in}{2.205492in}}%
\pgfpathlineto{\pgfqpoint{1.769730in}{2.206899in}}%
\pgfpathlineto{\pgfqpoint{1.770611in}{1.532023in}}%
\pgfpathlineto{\pgfqpoint{1.771398in}{1.532074in}}%
\pgfpathlineto{\pgfqpoint{1.772559in}{1.532118in}}%
\pgfpathlineto{\pgfqpoint{1.772569in}{1.532196in}}%
\pgfpathlineto{\pgfqpoint{1.773797in}{2.207100in}}%
\pgfpathlineto{\pgfqpoint{1.773846in}{2.206706in}}%
\pgfpathlineto{\pgfqpoint{1.774192in}{2.207137in}}%
\pgfpathlineto{\pgfqpoint{1.774439in}{2.206083in}}%
\pgfpathlineto{\pgfqpoint{1.774489in}{2.207025in}}%
\pgfpathlineto{\pgfqpoint{1.774637in}{2.183645in}}%
\pgfpathlineto{\pgfqpoint{1.774687in}{2.206843in}}%
\pgfpathlineto{\pgfqpoint{1.775537in}{1.532002in}}%
\pgfpathlineto{\pgfqpoint{1.776119in}{1.532067in}}%
\pgfpathlineto{\pgfqpoint{1.777503in}{1.532114in}}%
\pgfpathlineto{\pgfqpoint{1.777512in}{1.532160in}}%
\pgfpathlineto{\pgfqpoint{1.777540in}{1.607610in}}%
\pgfpathlineto{\pgfqpoint{1.778301in}{2.206696in}}%
\pgfpathlineto{\pgfqpoint{1.778400in}{2.206771in}}%
\pgfpathlineto{\pgfqpoint{1.779043in}{2.206770in}}%
\pgfpathlineto{\pgfqpoint{1.779092in}{2.206957in}}%
\pgfpathlineto{\pgfqpoint{1.779636in}{2.206032in}}%
\pgfpathlineto{\pgfqpoint{1.779686in}{2.206521in}}%
\pgfpathlineto{\pgfqpoint{1.779834in}{2.193836in}}%
\pgfpathlineto{\pgfqpoint{1.779883in}{2.206098in}}%
\pgfpathlineto{\pgfqpoint{1.780565in}{1.532034in}}%
\pgfpathlineto{\pgfqpoint{1.781202in}{1.532061in}}%
\pgfpathlineto{\pgfqpoint{1.782460in}{1.532188in}}%
\pgfpathlineto{\pgfqpoint{1.783820in}{2.206883in}}%
\pgfpathlineto{\pgfqpoint{1.783870in}{2.206814in}}%
\pgfpathlineto{\pgfqpoint{1.784315in}{2.206735in}}%
\pgfpathlineto{\pgfqpoint{1.784513in}{2.206543in}}%
\pgfpathlineto{\pgfqpoint{1.784859in}{2.204466in}}%
\pgfpathlineto{\pgfqpoint{1.784908in}{2.205530in}}%
\pgfpathlineto{\pgfqpoint{1.785464in}{1.532033in}}%
\pgfpathlineto{\pgfqpoint{1.786183in}{1.532063in}}%
\pgfpathlineto{\pgfqpoint{1.787402in}{1.532148in}}%
\pgfpathlineto{\pgfqpoint{1.787411in}{1.532493in}}%
\pgfpathlineto{\pgfqpoint{1.788466in}{2.206734in}}%
\pgfpathlineto{\pgfqpoint{1.789703in}{2.205540in}}%
\pgfpathlineto{\pgfqpoint{1.789752in}{2.205871in}}%
\pgfpathlineto{\pgfqpoint{1.789801in}{2.203312in}}%
\pgfpathlineto{\pgfqpoint{1.789851in}{2.205403in}}%
\pgfpathlineto{\pgfqpoint{1.790438in}{1.532033in}}%
\pgfpathlineto{\pgfqpoint{1.791075in}{1.532061in}}%
\pgfpathlineto{\pgfqpoint{1.792349in}{1.532165in}}%
\pgfpathlineto{\pgfqpoint{1.792391in}{1.860954in}}%
\pgfpathlineto{\pgfqpoint{1.793153in}{2.206572in}}%
\pgfpathlineto{\pgfqpoint{1.793548in}{2.206805in}}%
\pgfpathlineto{\pgfqpoint{1.793895in}{2.206535in}}%
\pgfpathlineto{\pgfqpoint{1.793944in}{2.206841in}}%
\pgfpathlineto{\pgfqpoint{1.794488in}{2.205408in}}%
\pgfpathlineto{\pgfqpoint{1.794538in}{2.206444in}}%
\pgfpathlineto{\pgfqpoint{1.795082in}{1.737130in}}%
\pgfpathlineto{\pgfqpoint{1.795374in}{1.532034in}}%
\pgfpathlineto{\pgfqpoint{1.795301in}{1.754034in}}%
\pgfpathlineto{\pgfqpoint{1.796009in}{1.532060in}}%
\pgfpathlineto{\pgfqpoint{1.797297in}{1.532204in}}%
\pgfpathlineto{\pgfqpoint{1.798512in}{2.206786in}}%
\pgfpathlineto{\pgfqpoint{1.798562in}{2.206466in}}%
\pgfpathlineto{\pgfqpoint{1.798908in}{2.206821in}}%
\pgfpathlineto{\pgfqpoint{1.799155in}{2.206010in}}%
\pgfpathlineto{\pgfqpoint{1.799205in}{2.206707in}}%
\pgfpathlineto{\pgfqpoint{1.799353in}{2.205172in}}%
\pgfpathlineto{\pgfqpoint{1.799402in}{2.206549in}}%
\pgfpathlineto{\pgfqpoint{1.800363in}{1.532035in}}%
\pgfpathlineto{\pgfqpoint{1.801154in}{1.532065in}}%
\pgfpathlineto{\pgfqpoint{1.802243in}{1.532221in}}%
\pgfpathlineto{\pgfqpoint{1.803442in}{2.206703in}}%
\pgfpathlineto{\pgfqpoint{1.803540in}{2.206732in}}%
\pgfpathlineto{\pgfqpoint{1.804282in}{2.205565in}}%
\pgfpathlineto{\pgfqpoint{1.804332in}{2.206470in}}%
\pgfpathlineto{\pgfqpoint{1.804381in}{2.204708in}}%
\pgfpathlineto{\pgfqpoint{1.804431in}{2.206357in}}%
\pgfpathlineto{\pgfqpoint{1.805310in}{1.532035in}}%
\pgfpathlineto{\pgfqpoint{1.805804in}{1.532054in}}%
\pgfpathlineto{\pgfqpoint{1.807189in}{1.532222in}}%
\pgfpathlineto{\pgfqpoint{1.808387in}{2.206649in}}%
\pgfpathlineto{\pgfqpoint{1.808486in}{2.206678in}}%
\pgfpathlineto{\pgfqpoint{1.809228in}{2.205521in}}%
\pgfpathlineto{\pgfqpoint{1.809278in}{2.206419in}}%
\pgfpathlineto{\pgfqpoint{1.809327in}{2.204699in}}%
\pgfpathlineto{\pgfqpoint{1.809376in}{2.206307in}}%
\pgfpathlineto{\pgfqpoint{1.810196in}{1.532035in}}%
\pgfpathlineto{\pgfqpoint{1.810772in}{1.532055in}}%
\pgfpathlineto{\pgfqpoint{1.812137in}{1.532329in}}%
\pgfpathlineto{\pgfqpoint{1.813153in}{2.206552in}}%
\pgfpathlineto{\pgfqpoint{1.813203in}{2.206266in}}%
\pgfpathlineto{\pgfqpoint{1.813747in}{2.206683in}}%
\pgfpathlineto{\pgfqpoint{1.813895in}{2.205978in}}%
\pgfpathlineto{\pgfqpoint{1.813945in}{2.206625in}}%
\pgfpathlineto{\pgfqpoint{1.814192in}{2.204557in}}%
\pgfpathlineto{\pgfqpoint{1.814241in}{2.206421in}}%
\pgfpathlineto{\pgfqpoint{1.815154in}{1.532032in}}%
\pgfpathlineto{\pgfqpoint{1.815595in}{1.532054in}}%
\pgfpathlineto{\pgfqpoint{1.817076in}{1.532162in}}%
\pgfpathlineto{\pgfqpoint{1.817118in}{1.844831in}}%
\pgfpathlineto{\pgfqpoint{1.817879in}{2.206288in}}%
\pgfpathlineto{\pgfqpoint{1.818720in}{2.206342in}}%
\pgfpathlineto{\pgfqpoint{1.818769in}{2.206504in}}%
\pgfpathlineto{\pgfqpoint{1.819214in}{2.205689in}}%
\pgfpathlineto{\pgfqpoint{1.819264in}{2.206067in}}%
\pgfpathlineto{\pgfqpoint{1.819412in}{2.203943in}}%
\pgfpathlineto{\pgfqpoint{1.819461in}{2.205646in}}%
\pgfpathlineto{\pgfqpoint{1.820144in}{1.532035in}}%
\pgfpathlineto{\pgfqpoint{1.820738in}{1.532057in}}%
\pgfpathlineto{\pgfqpoint{1.822024in}{1.532192in}}%
\pgfpathlineto{\pgfqpoint{1.823328in}{2.206559in}}%
\pgfpathlineto{\pgfqpoint{1.823427in}{2.206576in}}%
\pgfpathlineto{\pgfqpoint{1.824169in}{2.203003in}}%
\pgfpathlineto{\pgfqpoint{1.824218in}{2.206187in}}%
\pgfpathlineto{\pgfqpoint{1.825066in}{1.532034in}}%
\pgfpathlineto{\pgfqpoint{1.825457in}{1.532049in}}%
\pgfpathlineto{\pgfqpoint{1.826972in}{1.532258in}}%
\pgfpathlineto{\pgfqpoint{1.828080in}{2.206436in}}%
\pgfpathlineto{\pgfqpoint{1.828130in}{2.206164in}}%
\pgfpathlineto{\pgfqpoint{1.828575in}{2.206512in}}%
\pgfpathlineto{\pgfqpoint{1.828822in}{2.205753in}}%
\pgfpathlineto{\pgfqpoint{1.828871in}{2.206398in}}%
\pgfpathlineto{\pgfqpoint{1.829020in}{2.205059in}}%
\pgfpathlineto{\pgfqpoint{1.829069in}{2.206241in}}%
\pgfpathlineto{\pgfqpoint{1.829711in}{1.726495in}}%
\pgfpathlineto{\pgfqpoint{1.829981in}{1.532034in}}%
\pgfpathlineto{\pgfqpoint{1.829919in}{1.741348in}}%
\pgfpathlineto{\pgfqpoint{1.830611in}{1.532059in}}%
\pgfpathlineto{\pgfqpoint{1.831916in}{1.532220in}}%
\pgfpathlineto{\pgfqpoint{1.833110in}{2.206396in}}%
\pgfpathlineto{\pgfqpoint{1.833209in}{2.206425in}}%
\pgfpathlineto{\pgfqpoint{1.833950in}{2.205340in}}%
\pgfpathlineto{\pgfqpoint{1.834000in}{2.206164in}}%
\pgfpathlineto{\pgfqpoint{1.834049in}{2.204699in}}%
\pgfpathlineto{\pgfqpoint{1.834099in}{2.206052in}}%
\pgfpathlineto{\pgfqpoint{1.834919in}{1.532035in}}%
\pgfpathlineto{\pgfqpoint{1.835689in}{1.532062in}}%
\pgfpathlineto{\pgfqpoint{1.836861in}{1.532209in}}%
\pgfpathlineto{\pgfqpoint{1.838065in}{2.206359in}}%
\pgfpathlineto{\pgfqpoint{1.838114in}{2.206089in}}%
\pgfpathlineto{\pgfqpoint{1.838461in}{2.206399in}}%
\pgfpathlineto{\pgfqpoint{1.838708in}{2.205728in}}%
\pgfpathlineto{\pgfqpoint{1.838757in}{2.206286in}}%
\pgfpathlineto{\pgfqpoint{1.839005in}{2.204251in}}%
\pgfpathlineto{\pgfqpoint{1.839054in}{2.206018in}}%
\pgfpathlineto{\pgfqpoint{1.839855in}{1.532032in}}%
\pgfpathlineto{\pgfqpoint{1.840409in}{1.532051in}}%
\pgfpathlineto{\pgfqpoint{1.841811in}{1.532479in}}%
\pgfpathlineto{\pgfqpoint{1.842858in}{2.206152in}}%
\pgfpathlineto{\pgfqpoint{1.844095in}{2.205140in}}%
\pgfpathlineto{\pgfqpoint{1.844144in}{2.205297in}}%
\pgfpathlineto{\pgfqpoint{1.844194in}{2.204390in}}%
\pgfpathlineto{\pgfqpoint{1.844243in}{2.204839in}}%
\pgfpathlineto{\pgfqpoint{1.844810in}{1.532033in}}%
\pgfpathlineto{\pgfqpoint{1.845727in}{1.532071in}}%
\pgfpathlineto{\pgfqpoint{1.846747in}{1.532148in}}%
\pgfpathlineto{\pgfqpoint{1.846756in}{1.532494in}}%
\pgfpathlineto{\pgfqpoint{1.847812in}{2.206108in}}%
\pgfpathlineto{\pgfqpoint{1.849048in}{2.204957in}}%
\pgfpathlineto{\pgfqpoint{1.849098in}{2.205282in}}%
\pgfpathlineto{\pgfqpoint{1.849147in}{2.203535in}}%
\pgfpathlineto{\pgfqpoint{1.849197in}{2.204859in}}%
\pgfpathlineto{\pgfqpoint{1.849762in}{1.532030in}}%
\pgfpathlineto{\pgfqpoint{1.850439in}{1.532054in}}%
\pgfpathlineto{\pgfqpoint{1.851697in}{1.532193in}}%
\pgfpathlineto{\pgfqpoint{1.852928in}{2.206250in}}%
\pgfpathlineto{\pgfqpoint{1.852977in}{2.205807in}}%
\pgfpathlineto{\pgfqpoint{1.853323in}{2.206289in}}%
\pgfpathlineto{\pgfqpoint{1.853570in}{2.205037in}}%
\pgfpathlineto{\pgfqpoint{1.853620in}{2.206181in}}%
\pgfpathlineto{\pgfqpoint{1.853669in}{2.204101in}}%
\pgfpathlineto{\pgfqpoint{1.853719in}{2.206114in}}%
\pgfpathlineto{\pgfqpoint{1.854724in}{1.532027in}}%
\pgfpathlineto{\pgfqpoint{1.855063in}{1.532054in}}%
\pgfpathlineto{\pgfqpoint{1.856632in}{1.532120in}}%
\pgfpathlineto{\pgfqpoint{1.856643in}{1.532215in}}%
\pgfpathlineto{\pgfqpoint{1.857774in}{2.206170in}}%
\pgfpathlineto{\pgfqpoint{1.857823in}{2.205683in}}%
\pgfpathlineto{\pgfqpoint{1.858268in}{2.206255in}}%
\pgfpathlineto{\pgfqpoint{1.858417in}{2.204991in}}%
\pgfpathlineto{\pgfqpoint{1.858466in}{2.206201in}}%
\pgfpathlineto{\pgfqpoint{1.858515in}{2.204176in}}%
\pgfpathlineto{\pgfqpoint{1.858565in}{2.206154in}}%
\pgfpathlineto{\pgfqpoint{1.859640in}{1.532032in}}%
\pgfpathlineto{\pgfqpoint{1.859956in}{1.532049in}}%
\pgfpathlineto{\pgfqpoint{1.861584in}{1.532148in}}%
\pgfpathlineto{\pgfqpoint{1.861592in}{1.532479in}}%
\pgfpathlineto{\pgfqpoint{1.862640in}{2.205912in}}%
\pgfpathlineto{\pgfqpoint{1.863877in}{2.204897in}}%
\pgfpathlineto{\pgfqpoint{1.863926in}{2.205058in}}%
\pgfpathlineto{\pgfqpoint{1.863976in}{2.204140in}}%
\pgfpathlineto{\pgfqpoint{1.864025in}{2.204601in}}%
\pgfpathlineto{\pgfqpoint{1.864593in}{1.532032in}}%
\pgfpathlineto{\pgfqpoint{1.865461in}{1.532063in}}%
\pgfpathlineto{\pgfqpoint{1.866538in}{1.532495in}}%
\pgfpathlineto{\pgfqpoint{1.867594in}{2.205871in}}%
\pgfpathlineto{\pgfqpoint{1.868831in}{2.204719in}}%
\pgfpathlineto{\pgfqpoint{1.868880in}{2.205053in}}%
\pgfpathlineto{\pgfqpoint{1.868930in}{2.203274in}}%
\pgfpathlineto{\pgfqpoint{1.868979in}{2.204637in}}%
\pgfpathlineto{\pgfqpoint{1.869542in}{1.532031in}}%
\pgfpathlineto{\pgfqpoint{1.870265in}{1.532065in}}%
\pgfpathlineto{\pgfqpoint{1.871478in}{1.532182in}}%
\pgfpathlineto{\pgfqpoint{1.872779in}{2.205995in}}%
\pgfpathlineto{\pgfqpoint{1.872828in}{2.205676in}}%
\pgfpathlineto{\pgfqpoint{1.873076in}{2.206006in}}%
\pgfpathlineto{\pgfqpoint{1.873422in}{2.205079in}}%
\pgfpathlineto{\pgfqpoint{1.873471in}{2.205825in}}%
\pgfpathlineto{\pgfqpoint{1.873620in}{2.202612in}}%
\pgfpathlineto{\pgfqpoint{1.873669in}{2.205626in}}%
\pgfpathlineto{\pgfqpoint{1.874535in}{1.532035in}}%
\pgfpathlineto{\pgfqpoint{1.874926in}{1.532049in}}%
\pgfpathlineto{\pgfqpoint{1.876425in}{1.532215in}}%
\pgfpathlineto{\pgfqpoint{1.877629in}{2.205910in}}%
\pgfpathlineto{\pgfqpoint{1.877727in}{2.205938in}}%
\pgfpathlineto{\pgfqpoint{1.878469in}{2.204751in}}%
\pgfpathlineto{\pgfqpoint{1.878519in}{2.205676in}}%
\pgfpathlineto{\pgfqpoint{1.878568in}{2.203816in}}%
\pgfpathlineto{\pgfqpoint{1.878618in}{2.205563in}}%
\pgfpathlineto{\pgfqpoint{1.879461in}{1.532035in}}%
\pgfpathlineto{\pgfqpoint{1.879956in}{1.532053in}}%
\pgfpathlineto{\pgfqpoint{1.881374in}{1.532488in}}%
\pgfpathlineto{\pgfqpoint{1.882431in}{2.205720in}}%
\pgfpathlineto{\pgfqpoint{1.883667in}{2.204548in}}%
\pgfpathlineto{\pgfqpoint{1.883717in}{2.204909in}}%
\pgfpathlineto{\pgfqpoint{1.883766in}{2.202818in}}%
\pgfpathlineto{\pgfqpoint{1.883816in}{2.204497in}}%
\pgfpathlineto{\pgfqpoint{1.884379in}{1.532031in}}%
\pgfpathlineto{\pgfqpoint{1.885054in}{1.532055in}}%
\pgfpathlineto{\pgfqpoint{1.886314in}{1.532187in}}%
\pgfpathlineto{\pgfqpoint{1.887619in}{2.205846in}}%
\pgfpathlineto{\pgfqpoint{1.887718in}{2.205863in}}%
\pgfpathlineto{\pgfqpoint{1.888459in}{2.197649in}}%
\pgfpathlineto{\pgfqpoint{1.888509in}{2.205469in}}%
\pgfpathlineto{\pgfqpoint{1.888558in}{2.141685in}}%
\pgfpathlineto{\pgfqpoint{1.888858in}{2.153261in}}%
\pgfpathlineto{\pgfqpoint{1.889352in}{1.532034in}}%
\pgfpathlineto{\pgfqpoint{1.889989in}{1.532060in}}%
\pgfpathlineto{\pgfqpoint{1.891262in}{1.532239in}}%
\pgfpathlineto{\pgfqpoint{1.892461in}{2.205758in}}%
\pgfpathlineto{\pgfqpoint{1.892560in}{2.205787in}}%
\pgfpathlineto{\pgfqpoint{1.893301in}{2.204603in}}%
\pgfpathlineto{\pgfqpoint{1.893351in}{2.205533in}}%
\pgfpathlineto{\pgfqpoint{1.893400in}{2.203681in}}%
\pgfpathlineto{\pgfqpoint{1.893450in}{2.205422in}}%
\pgfpathlineto{\pgfqpoint{1.894280in}{1.532032in}}%
\pgfpathlineto{\pgfqpoint{1.894818in}{1.532057in}}%
\pgfpathlineto{\pgfqpoint{1.896204in}{1.532165in}}%
\pgfpathlineto{\pgfqpoint{1.896246in}{1.860446in}}%
\pgfpathlineto{\pgfqpoint{1.897007in}{2.205476in}}%
\pgfpathlineto{\pgfqpoint{1.897403in}{2.205712in}}%
\pgfpathlineto{\pgfqpoint{1.897749in}{2.205449in}}%
\pgfpathlineto{\pgfqpoint{1.897798in}{2.205752in}}%
\pgfpathlineto{\pgfqpoint{1.898342in}{2.204386in}}%
\pgfpathlineto{\pgfqpoint{1.898392in}{2.205365in}}%
\pgfpathlineto{\pgfqpoint{1.898837in}{1.968934in}}%
\pgfpathlineto{\pgfqpoint{1.899234in}{1.532034in}}%
\pgfpathlineto{\pgfqpoint{1.899771in}{1.532056in}}%
\pgfpathlineto{\pgfqpoint{1.901150in}{1.532182in}}%
\pgfpathlineto{\pgfqpoint{1.902651in}{2.205755in}}%
\pgfpathlineto{\pgfqpoint{1.902700in}{2.205351in}}%
\pgfpathlineto{\pgfqpoint{1.902750in}{2.205745in}}%
\pgfpathlineto{\pgfqpoint{1.903195in}{2.204459in}}%
\pgfpathlineto{\pgfqpoint{1.903244in}{2.205479in}}%
\pgfpathlineto{\pgfqpoint{1.903294in}{2.202834in}}%
\pgfpathlineto{\pgfqpoint{1.903343in}{2.205368in}}%
\pgfpathlineto{\pgfqpoint{1.904206in}{1.532035in}}%
\pgfpathlineto{\pgfqpoint{1.904599in}{1.532050in}}%
\pgfpathlineto{\pgfqpoint{1.906100in}{1.532295in}}%
\pgfpathlineto{\pgfqpoint{1.907208in}{2.205638in}}%
\pgfpathlineto{\pgfqpoint{1.907505in}{2.205727in}}%
\pgfpathlineto{\pgfqpoint{1.907950in}{2.204812in}}%
\pgfpathlineto{\pgfqpoint{1.907999in}{2.205622in}}%
\pgfpathlineto{\pgfqpoint{1.908148in}{2.203502in}}%
\pgfpathlineto{\pgfqpoint{1.908197in}{2.205473in}}%
\pgfpathlineto{\pgfqpoint{1.909110in}{1.532034in}}%
\pgfpathlineto{\pgfqpoint{1.909543in}{1.532051in}}%
\pgfpathlineto{\pgfqpoint{1.911045in}{1.532270in}}%
\pgfpathlineto{\pgfqpoint{1.912156in}{2.205582in}}%
\pgfpathlineto{\pgfqpoint{1.912205in}{2.205280in}}%
\pgfpathlineto{\pgfqpoint{1.912650in}{2.205664in}}%
\pgfpathlineto{\pgfqpoint{1.912897in}{2.204820in}}%
\pgfpathlineto{\pgfqpoint{1.912947in}{2.205556in}}%
\pgfpathlineto{\pgfqpoint{1.913095in}{2.203892in}}%
\pgfpathlineto{\pgfqpoint{1.913145in}{2.205403in}}%
\pgfpathlineto{\pgfqpoint{1.914035in}{1.532032in}}%
\pgfpathlineto{\pgfqpoint{1.914683in}{1.532054in}}%
\pgfpathlineto{\pgfqpoint{1.915992in}{1.532502in}}%
\pgfpathlineto{\pgfqpoint{1.917032in}{2.205385in}}%
\pgfpathlineto{\pgfqpoint{1.918219in}{2.204790in}}%
\pgfpathlineto{\pgfqpoint{1.918417in}{2.204008in}}%
\pgfpathlineto{\pgfqpoint{1.918615in}{1.945882in}}%
\pgfpathlineto{\pgfqpoint{1.919043in}{1.532035in}}%
\pgfpathlineto{\pgfqpoint{1.919532in}{1.532053in}}%
\pgfpathlineto{\pgfqpoint{1.920934in}{1.532209in}}%
\pgfpathlineto{\pgfqpoint{1.922138in}{2.205524in}}%
\pgfpathlineto{\pgfqpoint{1.922187in}{2.205253in}}%
\pgfpathlineto{\pgfqpoint{1.922533in}{2.205568in}}%
\pgfpathlineto{\pgfqpoint{1.922781in}{2.204901in}}%
\pgfpathlineto{\pgfqpoint{1.922830in}{2.205458in}}%
\pgfpathlineto{\pgfqpoint{1.923077in}{2.203511in}}%
\pgfpathlineto{\pgfqpoint{1.923127in}{2.205196in}}%
\pgfpathlineto{\pgfqpoint{1.923918in}{1.532024in}}%
\pgfpathlineto{\pgfqpoint{1.924505in}{1.532065in}}%
\pgfpathlineto{\pgfqpoint{1.925868in}{1.532119in}}%
\pgfpathlineto{\pgfqpoint{1.925880in}{1.532221in}}%
\pgfpathlineto{\pgfqpoint{1.926994in}{2.205419in}}%
\pgfpathlineto{\pgfqpoint{1.927934in}{2.204726in}}%
\pgfpathlineto{\pgfqpoint{1.927983in}{2.205144in}}%
\pgfpathlineto{\pgfqpoint{1.928132in}{2.203826in}}%
\pgfpathlineto{\pgfqpoint{1.928181in}{2.204841in}}%
\pgfpathlineto{\pgfqpoint{1.928895in}{1.532034in}}%
\pgfpathlineto{\pgfqpoint{1.929873in}{1.532070in}}%
\pgfpathlineto{\pgfqpoint{1.930826in}{1.532241in}}%
\pgfpathlineto{\pgfqpoint{1.932025in}{2.205461in}}%
\pgfpathlineto{\pgfqpoint{1.932124in}{2.205490in}}%
\pgfpathlineto{\pgfqpoint{1.932865in}{2.204316in}}%
\pgfpathlineto{\pgfqpoint{1.932915in}{2.205243in}}%
\pgfpathlineto{\pgfqpoint{1.932964in}{2.203432in}}%
\pgfpathlineto{\pgfqpoint{1.933014in}{2.205134in}}%
\pgfpathlineto{\pgfqpoint{1.933823in}{1.532027in}}%
\pgfpathlineto{\pgfqpoint{1.934389in}{1.532048in}}%
\pgfpathlineto{\pgfqpoint{1.935769in}{1.532181in}}%
\pgfpathlineto{\pgfqpoint{1.935809in}{1.879200in}}%
\pgfpathlineto{\pgfqpoint{1.936576in}{2.205205in}}%
\pgfpathlineto{\pgfqpoint{1.937070in}{2.205471in}}%
\pgfpathlineto{\pgfqpoint{1.937318in}{2.205096in}}%
\pgfpathlineto{\pgfqpoint{1.937367in}{2.205484in}}%
\pgfpathlineto{\pgfqpoint{1.937812in}{2.204248in}}%
\pgfpathlineto{\pgfqpoint{1.937862in}{2.205224in}}%
\pgfpathlineto{\pgfqpoint{1.937911in}{2.203135in}}%
\pgfpathlineto{\pgfqpoint{1.937961in}{2.205114in}}%
\pgfpathlineto{\pgfqpoint{1.938785in}{1.532031in}}%
\pgfpathlineto{\pgfqpoint{1.939271in}{1.532048in}}%
\pgfpathlineto{\pgfqpoint{1.940720in}{1.532471in}}%
\pgfpathlineto{\pgfqpoint{1.941772in}{2.205250in}}%
\pgfpathlineto{\pgfqpoint{1.943008in}{2.204163in}}%
\pgfpathlineto{\pgfqpoint{1.943057in}{2.204439in}}%
\pgfpathlineto{\pgfqpoint{1.943107in}{2.203159in}}%
\pgfpathlineto{\pgfqpoint{1.943156in}{2.204029in}}%
\pgfpathlineto{\pgfqpoint{1.943684in}{1.532032in}}%
\pgfpathlineto{\pgfqpoint{1.944475in}{1.532064in}}%
\pgfpathlineto{\pgfqpoint{1.945659in}{1.532181in}}%
\pgfpathlineto{\pgfqpoint{1.946963in}{2.205393in}}%
\pgfpathlineto{\pgfqpoint{1.947013in}{2.205062in}}%
\pgfpathlineto{\pgfqpoint{1.947260in}{2.205405in}}%
\pgfpathlineto{\pgfqpoint{1.947606in}{2.204446in}}%
\pgfpathlineto{\pgfqpoint{1.947656in}{2.205229in}}%
\pgfpathlineto{\pgfqpoint{1.947804in}{2.196791in}}%
\pgfpathlineto{\pgfqpoint{1.947854in}{2.205025in}}%
\pgfpathlineto{\pgfqpoint{1.948694in}{1.532032in}}%
\pgfpathlineto{\pgfqpoint{1.949135in}{1.532054in}}%
\pgfpathlineto{\pgfqpoint{1.950603in}{1.532156in}}%
\pgfpathlineto{\pgfqpoint{1.950632in}{1.611225in}}%
\pgfpathlineto{\pgfqpoint{1.951405in}{2.205062in}}%
\pgfpathlineto{\pgfqpoint{1.951899in}{2.205330in}}%
\pgfpathlineto{\pgfqpoint{1.952146in}{2.204988in}}%
\pgfpathlineto{\pgfqpoint{1.952196in}{2.205345in}}%
\pgfpathlineto{\pgfqpoint{1.952641in}{2.204241in}}%
\pgfpathlineto{\pgfqpoint{1.952690in}{2.205084in}}%
\pgfpathlineto{\pgfqpoint{1.952740in}{2.203592in}}%
\pgfpathlineto{\pgfqpoint{1.952789in}{2.204975in}}%
\pgfpathlineto{\pgfqpoint{1.953666in}{1.532035in}}%
\pgfpathlineto{\pgfqpoint{1.954355in}{1.532061in}}%
\pgfpathlineto{\pgfqpoint{1.955553in}{1.532257in}}%
\pgfpathlineto{\pgfqpoint{1.956661in}{2.205249in}}%
\pgfpathlineto{\pgfqpoint{1.956710in}{2.204939in}}%
\pgfpathlineto{\pgfqpoint{1.957155in}{2.205336in}}%
\pgfpathlineto{\pgfqpoint{1.957403in}{2.204476in}}%
\pgfpathlineto{\pgfqpoint{1.957452in}{2.205232in}}%
\pgfpathlineto{\pgfqpoint{1.957600in}{2.203496in}}%
\pgfpathlineto{\pgfqpoint{1.957650in}{2.205083in}}%
\pgfpathlineto{\pgfqpoint{1.958566in}{1.532032in}}%
\pgfpathlineto{\pgfqpoint{1.959148in}{1.532053in}}%
\pgfpathlineto{\pgfqpoint{1.960494in}{1.532165in}}%
\pgfpathlineto{\pgfqpoint{1.960537in}{1.858883in}}%
\pgfpathlineto{\pgfqpoint{1.961298in}{2.204984in}}%
\pgfpathlineto{\pgfqpoint{1.961694in}{2.205219in}}%
\pgfpathlineto{\pgfqpoint{1.962040in}{2.204954in}}%
\pgfpathlineto{\pgfqpoint{1.962089in}{2.205258in}}%
\pgfpathlineto{\pgfqpoint{1.962633in}{2.203879in}}%
\pgfpathlineto{\pgfqpoint{1.962683in}{2.204870in}}%
\pgfpathlineto{\pgfqpoint{1.963128in}{1.966313in}}%
\pgfpathlineto{\pgfqpoint{1.963524in}{1.532034in}}%
\pgfpathlineto{\pgfqpoint{1.964061in}{1.532056in}}%
\pgfpathlineto{\pgfqpoint{1.965441in}{1.532181in}}%
\pgfpathlineto{\pgfqpoint{1.966941in}{2.205265in}}%
\pgfpathlineto{\pgfqpoint{1.966990in}{2.204877in}}%
\pgfpathlineto{\pgfqpoint{1.967040in}{2.205257in}}%
\pgfpathlineto{\pgfqpoint{1.967485in}{2.204047in}}%
\pgfpathlineto{\pgfqpoint{1.967534in}{2.204992in}}%
\pgfpathlineto{\pgfqpoint{1.967584in}{2.203073in}}%
\pgfpathlineto{\pgfqpoint{1.967633in}{2.204883in}}%
\pgfpathlineto{\pgfqpoint{1.968447in}{1.532031in}}%
\pgfpathlineto{\pgfqpoint{1.968969in}{1.532057in}}%
\pgfpathlineto{\pgfqpoint{1.970384in}{1.532148in}}%
\pgfpathlineto{\pgfqpoint{1.970392in}{1.532494in}}%
\pgfpathlineto{\pgfqpoint{1.971450in}{2.205034in}}%
\pgfpathlineto{\pgfqpoint{1.972686in}{2.203885in}}%
\pgfpathlineto{\pgfqpoint{1.972735in}{2.204237in}}%
\pgfpathlineto{\pgfqpoint{1.972785in}{2.202464in}}%
\pgfpathlineto{\pgfqpoint{1.972834in}{2.203840in}}%
\pgfpathlineto{\pgfqpoint{1.973356in}{1.532027in}}%
\pgfpathlineto{\pgfqpoint{1.974117in}{1.532059in}}%
\pgfpathlineto{\pgfqpoint{1.975333in}{1.532192in}}%
\pgfpathlineto{\pgfqpoint{1.976737in}{2.205171in}}%
\pgfpathlineto{\pgfqpoint{1.977479in}{2.203687in}}%
\pgfpathlineto{\pgfqpoint{1.977529in}{2.204757in}}%
\pgfpathlineto{\pgfqpoint{1.978338in}{1.532035in}}%
\pgfpathlineto{\pgfqpoint{1.979800in}{1.532087in}}%
\pgfpathlineto{\pgfqpoint{1.980279in}{1.532205in}}%
\pgfpathlineto{\pgfqpoint{1.981483in}{2.205115in}}%
\pgfpathlineto{\pgfqpoint{1.981533in}{2.204840in}}%
\pgfpathlineto{\pgfqpoint{1.981879in}{2.205158in}}%
\pgfpathlineto{\pgfqpoint{1.982126in}{2.204480in}}%
\pgfpathlineto{\pgfqpoint{1.982176in}{2.205049in}}%
\pgfpathlineto{\pgfqpoint{1.982423in}{2.202960in}}%
\pgfpathlineto{\pgfqpoint{1.982472in}{2.204787in}}%
\pgfpathlineto{\pgfqpoint{1.983302in}{1.532030in}}%
\pgfpathlineto{\pgfqpoint{1.983791in}{1.532048in}}%
\pgfpathlineto{\pgfqpoint{1.985222in}{1.532171in}}%
\pgfpathlineto{\pgfqpoint{1.985265in}{1.894245in}}%
\pgfpathlineto{\pgfqpoint{1.986026in}{2.204837in}}%
\pgfpathlineto{\pgfqpoint{1.986422in}{2.205070in}}%
\pgfpathlineto{\pgfqpoint{1.986768in}{2.204821in}}%
\pgfpathlineto{\pgfqpoint{1.986817in}{2.205109in}}%
\pgfpathlineto{\pgfqpoint{1.987361in}{2.203812in}}%
\pgfpathlineto{\pgfqpoint{1.987411in}{2.204715in}}%
\pgfpathlineto{\pgfqpoint{1.987460in}{2.201344in}}%
\pgfpathlineto{\pgfqpoint{1.987510in}{2.204565in}}%
\pgfpathlineto{\pgfqpoint{1.988192in}{1.532012in}}%
\pgfpathlineto{\pgfqpoint{1.988760in}{1.532045in}}%
\pgfpathlineto{\pgfqpoint{1.990171in}{1.532224in}}%
\pgfpathlineto{\pgfqpoint{1.991300in}{2.205092in}}%
\pgfpathlineto{\pgfqpoint{1.991349in}{2.204674in}}%
\pgfpathlineto{\pgfqpoint{1.991794in}{2.205177in}}%
\pgfpathlineto{\pgfqpoint{1.991943in}{2.204190in}}%
\pgfpathlineto{\pgfqpoint{1.991992in}{2.205122in}}%
\pgfpathlineto{\pgfqpoint{1.992141in}{2.202303in}}%
\pgfpathlineto{\pgfqpoint{1.992190in}{2.205009in}}%
\pgfpathlineto{\pgfqpoint{1.993181in}{1.532030in}}%
\pgfpathlineto{\pgfqpoint{1.993463in}{1.532040in}}%
\pgfpathlineto{\pgfqpoint{1.995120in}{1.532469in}}%
\pgfpathlineto{\pgfqpoint{1.996173in}{2.204870in}}%
\pgfpathlineto{\pgfqpoint{1.997409in}{2.203713in}}%
\pgfpathlineto{\pgfqpoint{1.997458in}{2.204034in}}%
\pgfpathlineto{\pgfqpoint{1.997508in}{2.202234in}}%
\pgfpathlineto{\pgfqpoint{1.997557in}{2.203602in}}%
\pgfpathlineto{\pgfqpoint{1.998100in}{1.532024in}}%
\pgfpathlineto{\pgfqpoint{1.998827in}{1.532050in}}%
\pgfpathlineto{\pgfqpoint{2.000062in}{1.532252in}}%
\pgfpathlineto{\pgfqpoint{2.001189in}{2.204855in}}%
\pgfpathlineto{\pgfqpoint{2.002376in}{2.204044in}}%
\pgfpathlineto{\pgfqpoint{2.002426in}{2.203555in}}%
\pgfpathlineto{\pgfqpoint{2.002475in}{2.203648in}}%
\pgfpathlineto{\pgfqpoint{2.002673in}{2.002990in}}%
\pgfpathlineto{\pgfqpoint{2.003098in}{1.532035in}}%
\pgfpathlineto{\pgfqpoint{2.003586in}{1.532052in}}%
\pgfpathlineto{\pgfqpoint{2.005007in}{1.532226in}}%
\pgfpathlineto{\pgfqpoint{2.006204in}{2.204958in}}%
\pgfpathlineto{\pgfqpoint{2.006303in}{2.204989in}}%
\pgfpathlineto{\pgfqpoint{2.007044in}{2.203918in}}%
\pgfpathlineto{\pgfqpoint{2.007094in}{2.204750in}}%
\pgfpathlineto{\pgfqpoint{2.007143in}{2.203306in}}%
\pgfpathlineto{\pgfqpoint{2.007193in}{2.204643in}}%
\pgfpathlineto{\pgfqpoint{2.008014in}{1.532031in}}%
\pgfpathlineto{\pgfqpoint{2.008837in}{1.532068in}}%
\pgfpathlineto{\pgfqpoint{2.009947in}{1.532148in}}%
\pgfpathlineto{\pgfqpoint{2.009956in}{1.532493in}}%
\pgfpathlineto{\pgfqpoint{2.011014in}{2.204789in}}%
\pgfpathlineto{\pgfqpoint{2.012250in}{2.203594in}}%
\pgfpathlineto{\pgfqpoint{2.012300in}{2.203978in}}%
\pgfpathlineto{\pgfqpoint{2.012349in}{2.201182in}}%
\pgfpathlineto{\pgfqpoint{2.012399in}{2.203561in}}%
\pgfpathlineto{\pgfqpoint{2.012962in}{1.532030in}}%
\pgfpathlineto{\pgfqpoint{2.013639in}{1.532054in}}%
\pgfpathlineto{\pgfqpoint{2.014897in}{1.532194in}}%
\pgfpathlineto{\pgfqpoint{2.016128in}{2.204946in}}%
\pgfpathlineto{\pgfqpoint{2.016325in}{2.204991in}}%
\pgfpathlineto{\pgfqpoint{2.016869in}{2.202748in}}%
\pgfpathlineto{\pgfqpoint{2.016919in}{2.204820in}}%
\pgfpathlineto{\pgfqpoint{2.017915in}{1.532030in}}%
\pgfpathlineto{\pgfqpoint{2.018253in}{1.532051in}}%
\pgfpathlineto{\pgfqpoint{2.019840in}{1.532165in}}%
\pgfpathlineto{\pgfqpoint{2.019883in}{1.872437in}}%
\pgfpathlineto{\pgfqpoint{2.020644in}{2.204597in}}%
\pgfpathlineto{\pgfqpoint{2.021039in}{2.204827in}}%
\pgfpathlineto{\pgfqpoint{2.021386in}{2.204579in}}%
\pgfpathlineto{\pgfqpoint{2.021435in}{2.204862in}}%
\pgfpathlineto{\pgfqpoint{2.021979in}{2.203567in}}%
\pgfpathlineto{\pgfqpoint{2.022028in}{2.204459in}}%
\pgfpathlineto{\pgfqpoint{2.022078in}{2.201094in}}%
\pgfpathlineto{\pgfqpoint{2.022127in}{2.204306in}}%
\pgfpathlineto{\pgfqpoint{2.022811in}{1.532018in}}%
\pgfpathlineto{\pgfqpoint{2.023359in}{1.532045in}}%
\pgfpathlineto{\pgfqpoint{2.024790in}{1.532248in}}%
\pgfpathlineto{\pgfqpoint{2.025892in}{2.204704in}}%
\pgfpathlineto{\pgfqpoint{2.027128in}{2.203384in}}%
\pgfpathlineto{\pgfqpoint{2.027177in}{2.203728in}}%
\pgfpathlineto{\pgfqpoint{2.027425in}{1.933880in}}%
\pgfpathlineto{\pgfqpoint{2.027812in}{1.532034in}}%
\pgfpathlineto{\pgfqpoint{2.028348in}{1.532056in}}%
\pgfpathlineto{\pgfqpoint{2.029733in}{1.532199in}}%
\pgfpathlineto{\pgfqpoint{2.030947in}{2.204828in}}%
\pgfpathlineto{\pgfqpoint{2.030996in}{2.204512in}}%
\pgfpathlineto{\pgfqpoint{2.031343in}{2.204872in}}%
\pgfpathlineto{\pgfqpoint{2.031590in}{2.204090in}}%
\pgfpathlineto{\pgfqpoint{2.031639in}{2.204766in}}%
\pgfpathlineto{\pgfqpoint{2.031788in}{2.203366in}}%
\pgfpathlineto{\pgfqpoint{2.031837in}{2.204616in}}%
\pgfpathlineto{\pgfqpoint{2.032762in}{1.532026in}}%
\pgfpathlineto{\pgfqpoint{2.034090in}{1.532092in}}%
\pgfpathlineto{\pgfqpoint{2.034669in}{1.532122in}}%
\pgfpathlineto{\pgfqpoint{2.034680in}{1.532224in}}%
\pgfpathlineto{\pgfqpoint{2.035809in}{2.204812in}}%
\pgfpathlineto{\pgfqpoint{2.035858in}{2.204402in}}%
\pgfpathlineto{\pgfqpoint{2.036303in}{2.204897in}}%
\pgfpathlineto{\pgfqpoint{2.036452in}{2.203935in}}%
\pgfpathlineto{\pgfqpoint{2.036501in}{2.204842in}}%
\pgfpathlineto{\pgfqpoint{2.036650in}{2.202510in}}%
\pgfpathlineto{\pgfqpoint{2.036699in}{2.204730in}}%
\pgfpathlineto{\pgfqpoint{2.037664in}{1.532023in}}%
\pgfpathlineto{\pgfqpoint{2.038005in}{1.532035in}}%
\pgfpathlineto{\pgfqpoint{2.039624in}{1.532206in}}%
\pgfpathlineto{\pgfqpoint{2.040856in}{2.204810in}}%
\pgfpathlineto{\pgfqpoint{2.041499in}{2.203374in}}%
\pgfpathlineto{\pgfqpoint{2.041053in}{2.204855in}}%
\pgfpathlineto{\pgfqpoint{2.041548in}{2.204749in}}%
\pgfpathlineto{\pgfqpoint{2.042389in}{1.720184in}}%
\pgfpathlineto{\pgfqpoint{2.042617in}{1.532031in}}%
\pgfpathlineto{\pgfqpoint{2.042573in}{1.728715in}}%
\pgfpathlineto{\pgfqpoint{2.043315in}{1.532062in}}%
\pgfpathlineto{\pgfqpoint{2.044565in}{1.532148in}}%
\pgfpathlineto{\pgfqpoint{2.044574in}{1.532495in}}%
\pgfpathlineto{\pgfqpoint{2.045618in}{2.204526in}}%
\pgfpathlineto{\pgfqpoint{2.046805in}{2.203924in}}%
\pgfpathlineto{\pgfqpoint{2.046954in}{2.203043in}}%
\pgfpathlineto{\pgfqpoint{2.047003in}{2.203149in}}%
\pgfpathlineto{\pgfqpoint{2.047538in}{1.532025in}}%
\pgfpathlineto{\pgfqpoint{2.048875in}{1.532077in}}%
\pgfpathlineto{\pgfqpoint{2.049514in}{1.532182in}}%
\pgfpathlineto{\pgfqpoint{2.050882in}{2.204590in}}%
\pgfpathlineto{\pgfqpoint{2.050981in}{2.204589in}}%
\pgfpathlineto{\pgfqpoint{2.051871in}{2.203560in}}%
\pgfpathlineto{\pgfqpoint{2.051970in}{2.202954in}}%
\pgfpathlineto{\pgfqpoint{2.052217in}{1.794523in}}%
\pgfpathlineto{\pgfqpoint{2.052522in}{1.532034in}}%
\pgfpathlineto{\pgfqpoint{2.053096in}{1.532054in}}%
\pgfpathlineto{\pgfqpoint{2.054460in}{1.532197in}}%
\pgfpathlineto{\pgfqpoint{2.055666in}{2.204658in}}%
\pgfpathlineto{\pgfqpoint{2.055715in}{2.204375in}}%
\pgfpathlineto{\pgfqpoint{2.056061in}{2.204702in}}%
\pgfpathlineto{\pgfqpoint{2.056308in}{2.204008in}}%
\pgfpathlineto{\pgfqpoint{2.056358in}{2.204595in}}%
\pgfpathlineto{\pgfqpoint{2.056605in}{2.202255in}}%
\pgfpathlineto{\pgfqpoint{2.056655in}{2.204335in}}%
\pgfpathlineto{\pgfqpoint{2.057445in}{1.532025in}}%
\pgfpathlineto{\pgfqpoint{2.057930in}{1.532060in}}%
\pgfpathlineto{\pgfqpoint{2.059396in}{1.532121in}}%
\pgfpathlineto{\pgfqpoint{2.059407in}{1.532225in}}%
\pgfpathlineto{\pgfqpoint{2.060538in}{2.204619in}}%
\pgfpathlineto{\pgfqpoint{2.060736in}{2.204683in}}%
\pgfpathlineto{\pgfqpoint{2.061379in}{2.203685in}}%
\pgfpathlineto{\pgfqpoint{2.061428in}{2.204497in}}%
\pgfpathlineto{\pgfqpoint{2.061478in}{2.203222in}}%
\pgfpathlineto{\pgfqpoint{2.061527in}{2.204406in}}%
\pgfpathlineto{\pgfqpoint{2.062363in}{1.585059in}}%
\pgfpathlineto{\pgfqpoint{2.062437in}{1.532032in}}%
\pgfpathlineto{\pgfqpoint{2.063126in}{1.532057in}}%
\pgfpathlineto{\pgfqpoint{2.064348in}{1.532156in}}%
\pgfpathlineto{\pgfqpoint{2.064377in}{1.612579in}}%
\pgfpathlineto{\pgfqpoint{2.065150in}{2.204358in}}%
\pgfpathlineto{\pgfqpoint{2.065546in}{2.204593in}}%
\pgfpathlineto{\pgfqpoint{2.065892in}{2.204295in}}%
\pgfpathlineto{\pgfqpoint{2.065942in}{2.204634in}}%
\pgfpathlineto{\pgfqpoint{2.066486in}{2.203023in}}%
\pgfpathlineto{\pgfqpoint{2.066535in}{2.204256in}}%
\pgfpathlineto{\pgfqpoint{2.067361in}{1.532034in}}%
\pgfpathlineto{\pgfqpoint{2.068334in}{1.532069in}}%
\pgfpathlineto{\pgfqpoint{2.069297in}{1.532208in}}%
\pgfpathlineto{\pgfqpoint{2.070502in}{2.204576in}}%
\pgfpathlineto{\pgfqpoint{2.070552in}{2.204302in}}%
\pgfpathlineto{\pgfqpoint{2.070898in}{2.204619in}}%
\pgfpathlineto{\pgfqpoint{2.071145in}{2.203941in}}%
\pgfpathlineto{\pgfqpoint{2.071195in}{2.204510in}}%
\pgfpathlineto{\pgfqpoint{2.071442in}{2.202424in}}%
\pgfpathlineto{\pgfqpoint{2.071491in}{2.204248in}}%
\pgfpathlineto{\pgfqpoint{2.072282in}{1.532024in}}%
\pgfpathlineto{\pgfqpoint{2.072819in}{1.532043in}}%
\pgfpathlineto{\pgfqpoint{2.074243in}{1.532220in}}%
\pgfpathlineto{\pgfqpoint{2.075363in}{2.204486in}}%
\pgfpathlineto{\pgfqpoint{2.075462in}{2.204523in}}%
\pgfpathlineto{\pgfqpoint{2.076204in}{2.203969in}}%
\pgfpathlineto{\pgfqpoint{2.076253in}{2.204324in}}%
\pgfpathlineto{\pgfqpoint{2.076500in}{2.202807in}}%
\pgfpathlineto{\pgfqpoint{2.076550in}{2.203924in}}%
\pgfpathlineto{\pgfqpoint{2.077288in}{1.532035in}}%
\pgfpathlineto{\pgfqpoint{2.078124in}{1.532067in}}%
\pgfpathlineto{\pgfqpoint{2.079190in}{1.532264in}}%
\pgfpathlineto{\pgfqpoint{2.080298in}{2.204520in}}%
\pgfpathlineto{\pgfqpoint{2.080347in}{2.204217in}}%
\pgfpathlineto{\pgfqpoint{2.080792in}{2.204613in}}%
\pgfpathlineto{\pgfqpoint{2.081039in}{2.203786in}}%
\pgfpathlineto{\pgfqpoint{2.081089in}{2.204513in}}%
\pgfpathlineto{\pgfqpoint{2.081237in}{2.202954in}}%
\pgfpathlineto{\pgfqpoint{2.081287in}{2.204368in}}%
\pgfpathlineto{\pgfqpoint{2.082215in}{1.532027in}}%
\pgfpathlineto{\pgfqpoint{2.083000in}{1.532056in}}%
\pgfpathlineto{\pgfqpoint{2.084134in}{1.532224in}}%
\pgfpathlineto{\pgfqpoint{2.085264in}{2.204551in}}%
\pgfpathlineto{\pgfqpoint{2.085313in}{2.204090in}}%
\pgfpathlineto{\pgfqpoint{2.085758in}{2.204641in}}%
\pgfpathlineto{\pgfqpoint{2.085907in}{2.203510in}}%
\pgfpathlineto{\pgfqpoint{2.085956in}{2.204590in}}%
\pgfpathlineto{\pgfqpoint{2.086104in}{2.175752in}}%
\pgfpathlineto{\pgfqpoint{2.086154in}{2.204456in}}%
\pgfpathlineto{\pgfqpoint{2.087161in}{1.532027in}}%
\pgfpathlineto{\pgfqpoint{2.087648in}{1.532045in}}%
\pgfpathlineto{\pgfqpoint{2.089080in}{1.532218in}}%
\pgfpathlineto{\pgfqpoint{2.090209in}{2.204516in}}%
\pgfpathlineto{\pgfqpoint{2.090259in}{2.204079in}}%
\pgfpathlineto{\pgfqpoint{2.090704in}{2.204602in}}%
\pgfpathlineto{\pgfqpoint{2.090852in}{2.203553in}}%
\pgfpathlineto{\pgfqpoint{2.090902in}{2.204548in}}%
\pgfpathlineto{\pgfqpoint{2.091050in}{2.191730in}}%
\pgfpathlineto{\pgfqpoint{2.091100in}{2.204433in}}%
\pgfpathlineto{\pgfqpoint{2.092101in}{1.532030in}}%
\pgfpathlineto{\pgfqpoint{2.092392in}{1.532040in}}%
\pgfpathlineto{\pgfqpoint{2.094022in}{1.532171in}}%
\pgfpathlineto{\pgfqpoint{2.094065in}{1.901924in}}%
\pgfpathlineto{\pgfqpoint{2.094826in}{2.204200in}}%
\pgfpathlineto{\pgfqpoint{2.095222in}{2.204430in}}%
\pgfpathlineto{\pgfqpoint{2.095568in}{2.204187in}}%
\pgfpathlineto{\pgfqpoint{2.095617in}{2.204464in}}%
\pgfpathlineto{\pgfqpoint{2.096161in}{2.203198in}}%
\pgfpathlineto{\pgfqpoint{2.096211in}{2.204059in}}%
\pgfpathlineto{\pgfqpoint{2.096260in}{2.201558in}}%
\pgfpathlineto{\pgfqpoint{2.096310in}{2.203907in}}%
\pgfpathlineto{\pgfqpoint{2.096992in}{1.532017in}}%
\pgfpathlineto{\pgfqpoint{2.097547in}{1.532045in}}%
\pgfpathlineto{\pgfqpoint{2.098970in}{1.532217in}}%
\pgfpathlineto{\pgfqpoint{2.100101in}{2.204468in}}%
\pgfpathlineto{\pgfqpoint{2.100150in}{2.204023in}}%
\pgfpathlineto{\pgfqpoint{2.100595in}{2.204561in}}%
\pgfpathlineto{\pgfqpoint{2.100743in}{2.203490in}}%
\pgfpathlineto{\pgfqpoint{2.100793in}{2.204511in}}%
\pgfpathlineto{\pgfqpoint{2.100941in}{2.188258in}}%
\pgfpathlineto{\pgfqpoint{2.100991in}{2.204395in}}%
\pgfpathlineto{\pgfqpoint{2.101985in}{1.532028in}}%
\pgfpathlineto{\pgfqpoint{2.102321in}{1.532053in}}%
\pgfpathlineto{\pgfqpoint{2.103905in}{1.532119in}}%
\pgfpathlineto{\pgfqpoint{2.103916in}{1.532233in}}%
\pgfpathlineto{\pgfqpoint{2.105047in}{2.204389in}}%
\pgfpathlineto{\pgfqpoint{2.105245in}{2.204451in}}%
\pgfpathlineto{\pgfqpoint{2.105888in}{2.203466in}}%
\pgfpathlineto{\pgfqpoint{2.105937in}{2.204252in}}%
\pgfpathlineto{\pgfqpoint{2.105987in}{2.203017in}}%
\pgfpathlineto{\pgfqpoint{2.106036in}{2.204158in}}%
\pgfpathlineto{\pgfqpoint{2.106630in}{1.789458in}}%
\pgfpathlineto{\pgfqpoint{2.106961in}{1.532034in}}%
\pgfpathlineto{\pgfqpoint{2.107549in}{1.532056in}}%
\pgfpathlineto{\pgfqpoint{2.108860in}{1.532185in}}%
\pgfpathlineto{\pgfqpoint{2.110360in}{2.204438in}}%
\pgfpathlineto{\pgfqpoint{2.111003in}{2.202442in}}%
\pgfpathlineto{\pgfqpoint{2.111052in}{2.204057in}}%
\pgfpathlineto{\pgfqpoint{2.111868in}{1.532031in}}%
\pgfpathlineto{\pgfqpoint{2.112443in}{1.532052in}}%
\pgfpathlineto{\pgfqpoint{2.113811in}{1.532493in}}%
\pgfpathlineto{\pgfqpoint{2.114869in}{2.204220in}}%
\pgfpathlineto{\pgfqpoint{2.116105in}{2.203065in}}%
\pgfpathlineto{\pgfqpoint{2.116154in}{2.203440in}}%
\pgfpathlineto{\pgfqpoint{2.116204in}{2.201559in}}%
\pgfpathlineto{\pgfqpoint{2.116253in}{2.203058in}}%
\pgfpathlineto{\pgfqpoint{2.116775in}{1.532021in}}%
\pgfpathlineto{\pgfqpoint{2.117498in}{1.532065in}}%
\pgfpathlineto{\pgfqpoint{2.118740in}{1.532117in}}%
\pgfpathlineto{\pgfqpoint{2.118751in}{1.532197in}}%
\pgfpathlineto{\pgfqpoint{2.120063in}{2.204370in}}%
\pgfpathlineto{\pgfqpoint{2.120804in}{2.202929in}}%
\pgfpathlineto{\pgfqpoint{2.120162in}{2.204386in}}%
\pgfpathlineto{\pgfqpoint{2.120854in}{2.204116in}}%
\pgfpathlineto{\pgfqpoint{2.121492in}{1.737492in}}%
\pgfpathlineto{\pgfqpoint{2.121747in}{1.532032in}}%
\pgfpathlineto{\pgfqpoint{2.122406in}{1.532055in}}%
\pgfpathlineto{\pgfqpoint{2.123702in}{1.532497in}}%
\pgfpathlineto{\pgfqpoint{2.124746in}{2.204132in}}%
\pgfpathlineto{\pgfqpoint{2.126031in}{2.203290in}}%
\pgfpathlineto{\pgfqpoint{2.126081in}{2.202749in}}%
\pgfpathlineto{\pgfqpoint{2.126130in}{2.202831in}}%
\pgfpathlineto{\pgfqpoint{2.126427in}{1.735684in}}%
\pgfpathlineto{\pgfqpoint{2.126720in}{1.532034in}}%
\pgfpathlineto{\pgfqpoint{2.126646in}{1.753436in}}%
\pgfpathlineto{\pgfqpoint{2.127356in}{1.532060in}}%
\pgfpathlineto{\pgfqpoint{2.128641in}{1.532183in}}%
\pgfpathlineto{\pgfqpoint{2.130140in}{2.204306in}}%
\pgfpathlineto{\pgfqpoint{2.130882in}{2.202431in}}%
\pgfpathlineto{\pgfqpoint{2.130931in}{2.203743in}}%
\pgfpathlineto{\pgfqpoint{2.131645in}{1.532035in}}%
\pgfpathlineto{\pgfqpoint{2.132411in}{1.532061in}}%
\pgfpathlineto{\pgfqpoint{2.133590in}{1.532268in}}%
\pgfpathlineto{\pgfqpoint{2.134697in}{2.204258in}}%
\pgfpathlineto{\pgfqpoint{2.134746in}{2.203920in}}%
\pgfpathlineto{\pgfqpoint{2.135191in}{2.204350in}}%
\pgfpathlineto{\pgfqpoint{2.135439in}{2.203407in}}%
\pgfpathlineto{\pgfqpoint{2.135488in}{2.204252in}}%
\pgfpathlineto{\pgfqpoint{2.135637in}{2.201794in}}%
\pgfpathlineto{\pgfqpoint{2.135686in}{2.204108in}}%
\pgfpathlineto{\pgfqpoint{2.136592in}{1.532035in}}%
\pgfpathlineto{\pgfqpoint{2.136965in}{1.532047in}}%
\pgfpathlineto{\pgfqpoint{2.138537in}{1.532396in}}%
\pgfpathlineto{\pgfqpoint{2.139538in}{2.204127in}}%
\pgfpathlineto{\pgfqpoint{2.139587in}{2.203924in}}%
\pgfpathlineto{\pgfqpoint{2.140131in}{2.204258in}}%
\pgfpathlineto{\pgfqpoint{2.140279in}{2.203758in}}%
\pgfpathlineto{\pgfqpoint{2.140329in}{2.204198in}}%
\pgfpathlineto{\pgfqpoint{2.140675in}{2.202495in}}%
\pgfpathlineto{\pgfqpoint{2.140724in}{2.203875in}}%
\pgfpathlineto{\pgfqpoint{2.141548in}{1.532033in}}%
\pgfpathlineto{\pgfqpoint{2.142274in}{1.532063in}}%
\pgfpathlineto{\pgfqpoint{2.143478in}{1.532183in}}%
\pgfpathlineto{\pgfqpoint{2.144978in}{2.204276in}}%
\pgfpathlineto{\pgfqpoint{2.145027in}{2.203882in}}%
\pgfpathlineto{\pgfqpoint{2.145077in}{2.204269in}}%
\pgfpathlineto{\pgfqpoint{2.145522in}{2.203069in}}%
\pgfpathlineto{\pgfqpoint{2.145571in}{2.204016in}}%
\pgfpathlineto{\pgfqpoint{2.145621in}{2.202158in}}%
\pgfpathlineto{\pgfqpoint{2.145670in}{2.203909in}}%
\pgfpathlineto{\pgfqpoint{2.146463in}{1.532025in}}%
\pgfpathlineto{\pgfqpoint{2.147047in}{1.532064in}}%
\pgfpathlineto{\pgfqpoint{2.148414in}{1.532121in}}%
\pgfpathlineto{\pgfqpoint{2.148425in}{1.532224in}}%
\pgfpathlineto{\pgfqpoint{2.149558in}{2.204199in}}%
\pgfpathlineto{\pgfqpoint{2.149756in}{2.204263in}}%
\pgfpathlineto{\pgfqpoint{2.150300in}{2.203418in}}%
\pgfpathlineto{\pgfqpoint{2.150349in}{2.204152in}}%
\pgfpathlineto{\pgfqpoint{2.150497in}{2.202423in}}%
\pgfpathlineto{\pgfqpoint{2.150547in}{2.203993in}}%
\pgfpathlineto{\pgfqpoint{2.151433in}{1.532031in}}%
\pgfpathlineto{\pgfqpoint{2.152059in}{1.532061in}}%
\pgfpathlineto{\pgfqpoint{2.153365in}{1.532148in}}%
\pgfpathlineto{\pgfqpoint{2.153374in}{1.532493in}}%
\pgfpathlineto{\pgfqpoint{2.154433in}{2.204039in}}%
\pgfpathlineto{\pgfqpoint{2.155669in}{2.202873in}}%
\pgfpathlineto{\pgfqpoint{2.155718in}{2.203259in}}%
\pgfpathlineto{\pgfqpoint{2.155768in}{2.201223in}}%
\pgfpathlineto{\pgfqpoint{2.155817in}{2.202875in}}%
\pgfpathlineto{\pgfqpoint{2.156338in}{1.532021in}}%
\pgfpathlineto{\pgfqpoint{2.157051in}{1.532064in}}%
\pgfpathlineto{\pgfqpoint{2.158304in}{1.532117in}}%
\pgfpathlineto{\pgfqpoint{2.158315in}{1.532190in}}%
\pgfpathlineto{\pgfqpoint{2.159546in}{2.204206in}}%
\pgfpathlineto{\pgfqpoint{2.159595in}{2.203760in}}%
\pgfpathlineto{\pgfqpoint{2.159941in}{2.204249in}}%
\pgfpathlineto{\pgfqpoint{2.160189in}{2.203004in}}%
\pgfpathlineto{\pgfqpoint{2.160238in}{2.204147in}}%
\pgfpathlineto{\pgfqpoint{2.160287in}{2.202125in}}%
\pgfpathlineto{\pgfqpoint{2.160337in}{2.204082in}}%
\pgfpathlineto{\pgfqpoint{2.161300in}{1.532024in}}%
\pgfpathlineto{\pgfqpoint{2.161688in}{1.532057in}}%
\pgfpathlineto{\pgfqpoint{2.163250in}{1.532119in}}%
\pgfpathlineto{\pgfqpoint{2.163262in}{1.532235in}}%
\pgfpathlineto{\pgfqpoint{2.164364in}{2.204002in}}%
\pgfpathlineto{\pgfqpoint{2.165600in}{2.202529in}}%
\pgfpathlineto{\pgfqpoint{2.165650in}{2.203048in}}%
\pgfpathlineto{\pgfqpoint{2.166271in}{1.532034in}}%
\pgfpathlineto{\pgfqpoint{2.167592in}{1.532085in}}%
\pgfpathlineto{\pgfqpoint{2.168205in}{1.532183in}}%
\pgfpathlineto{\pgfqpoint{2.169710in}{2.204165in}}%
\pgfpathlineto{\pgfqpoint{2.170352in}{2.202724in}}%
\pgfpathlineto{\pgfqpoint{2.170402in}{2.203782in}}%
\pgfpathlineto{\pgfqpoint{2.170946in}{1.771227in}}%
\pgfpathlineto{\pgfqpoint{2.171208in}{1.532030in}}%
\pgfpathlineto{\pgfqpoint{2.171875in}{1.532054in}}%
\pgfpathlineto{\pgfqpoint{2.173156in}{1.532491in}}%
\pgfpathlineto{\pgfqpoint{2.174215in}{2.203960in}}%
\pgfpathlineto{\pgfqpoint{2.175451in}{2.202722in}}%
\pgfpathlineto{\pgfqpoint{2.175501in}{2.203150in}}%
\pgfpathlineto{\pgfqpoint{2.175698in}{2.089586in}}%
\pgfpathlineto{\pgfqpoint{2.176162in}{1.532030in}}%
\pgfpathlineto{\pgfqpoint{2.176739in}{1.532051in}}%
\pgfpathlineto{\pgfqpoint{2.178096in}{1.532191in}}%
\pgfpathlineto{\pgfqpoint{2.179327in}{2.204121in}}%
\pgfpathlineto{\pgfqpoint{2.179376in}{2.203690in}}%
\pgfpathlineto{\pgfqpoint{2.179722in}{2.204166in}}%
\pgfpathlineto{\pgfqpoint{2.179970in}{2.203001in}}%
\pgfpathlineto{\pgfqpoint{2.180019in}{2.204065in}}%
\pgfpathlineto{\pgfqpoint{2.180069in}{2.202356in}}%
\pgfpathlineto{\pgfqpoint{2.180118in}{2.204001in}}%
\pgfpathlineto{\pgfqpoint{2.181118in}{1.532029in}}%
\pgfpathlineto{\pgfqpoint{2.181654in}{1.532060in}}%
\pgfpathlineto{\pgfqpoint{2.183032in}{1.532118in}}%
\pgfpathlineto{\pgfqpoint{2.183043in}{1.532208in}}%
\pgfpathlineto{\pgfqpoint{2.184174in}{2.204074in}}%
\pgfpathlineto{\pgfqpoint{2.184224in}{2.203580in}}%
\pgfpathlineto{\pgfqpoint{2.184669in}{2.204165in}}%
\pgfpathlineto{\pgfqpoint{2.184817in}{2.202900in}}%
\pgfpathlineto{\pgfqpoint{2.184867in}{2.204116in}}%
\pgfpathlineto{\pgfqpoint{2.184916in}{2.202105in}}%
\pgfpathlineto{\pgfqpoint{2.184966in}{2.204071in}}%
\pgfpathlineto{\pgfqpoint{2.186069in}{1.532027in}}%
\pgfpathlineto{\pgfqpoint{2.186359in}{1.532037in}}%
\pgfpathlineto{\pgfqpoint{2.187989in}{1.532217in}}%
\pgfpathlineto{\pgfqpoint{2.189119in}{2.204039in}}%
\pgfpathlineto{\pgfqpoint{2.189168in}{2.203583in}}%
\pgfpathlineto{\pgfqpoint{2.189613in}{2.204129in}}%
\pgfpathlineto{\pgfqpoint{2.189762in}{2.203019in}}%
\pgfpathlineto{\pgfqpoint{2.189811in}{2.204079in}}%
\pgfpathlineto{\pgfqpoint{2.189959in}{2.179479in}}%
\pgfpathlineto{\pgfqpoint{2.190009in}{2.203950in}}%
\pgfpathlineto{\pgfqpoint{2.191013in}{1.532027in}}%
\pgfpathlineto{\pgfqpoint{2.191451in}{1.532058in}}%
\pgfpathlineto{\pgfqpoint{2.192923in}{1.532120in}}%
\pgfpathlineto{\pgfqpoint{2.192935in}{1.532249in}}%
\pgfpathlineto{\pgfqpoint{2.194037in}{2.203858in}}%
\pgfpathlineto{\pgfqpoint{2.195273in}{2.202505in}}%
\pgfpathlineto{\pgfqpoint{2.195323in}{2.202871in}}%
\pgfpathlineto{\pgfqpoint{2.195570in}{1.928593in}}%
\pgfpathlineto{\pgfqpoint{2.195935in}{1.532034in}}%
\pgfpathlineto{\pgfqpoint{2.196503in}{1.532053in}}%
\pgfpathlineto{\pgfqpoint{2.197878in}{1.532184in}}%
\pgfpathlineto{\pgfqpoint{2.199195in}{2.204029in}}%
\pgfpathlineto{\pgfqpoint{2.199937in}{2.202455in}}%
\pgfpathlineto{\pgfqpoint{2.199294in}{2.204047in}}%
\pgfpathlineto{\pgfqpoint{2.199987in}{2.203797in}}%
\pgfpathlineto{\pgfqpoint{2.200900in}{1.532029in}}%
\pgfpathlineto{\pgfqpoint{2.201882in}{1.532065in}}%
\pgfpathlineto{\pgfqpoint{2.202825in}{1.532215in}}%
\pgfpathlineto{\pgfqpoint{2.203955in}{2.203996in}}%
\pgfpathlineto{\pgfqpoint{2.204005in}{2.203513in}}%
\pgfpathlineto{\pgfqpoint{2.204450in}{2.204087in}}%
\pgfpathlineto{\pgfqpoint{2.204598in}{2.202866in}}%
\pgfpathlineto{\pgfqpoint{2.204648in}{2.204037in}}%
\pgfpathlineto{\pgfqpoint{2.204697in}{2.202194in}}%
\pgfpathlineto{\pgfqpoint{2.204747in}{2.203989in}}%
\pgfpathlineto{\pgfqpoint{2.205850in}{1.532027in}}%
\pgfpathlineto{\pgfqpoint{2.206238in}{1.532042in}}%
\pgfpathlineto{\pgfqpoint{2.207770in}{1.532214in}}%
\pgfpathlineto{\pgfqpoint{2.208901in}{2.203968in}}%
\pgfpathlineto{\pgfqpoint{2.208950in}{2.203471in}}%
\pgfpathlineto{\pgfqpoint{2.209396in}{2.204063in}}%
\pgfpathlineto{\pgfqpoint{2.209544in}{2.202785in}}%
\pgfpathlineto{\pgfqpoint{2.209593in}{2.204014in}}%
\pgfpathlineto{\pgfqpoint{2.209643in}{2.201958in}}%
\pgfpathlineto{\pgfqpoint{2.209692in}{2.203969in}}%
\pgfpathlineto{\pgfqpoint{2.210796in}{1.532027in}}%
\pgfpathlineto{\pgfqpoint{2.211086in}{1.532037in}}%
\pgfpathlineto{\pgfqpoint{2.212716in}{1.532216in}}%
\pgfpathlineto{\pgfqpoint{2.213846in}{2.203934in}}%
\pgfpathlineto{\pgfqpoint{2.213895in}{2.203468in}}%
\pgfpathlineto{\pgfqpoint{2.214341in}{2.204024in}}%
\pgfpathlineto{\pgfqpoint{2.214489in}{2.202870in}}%
\pgfpathlineto{\pgfqpoint{2.214538in}{2.203973in}}%
\pgfpathlineto{\pgfqpoint{2.214588in}{2.202328in}}%
\pgfpathlineto{\pgfqpoint{2.214637in}{2.203926in}}%
\pgfpathlineto{\pgfqpoint{2.215741in}{1.532027in}}%
\pgfpathlineto{\pgfqpoint{2.216377in}{1.532065in}}%
\pgfpathlineto{\pgfqpoint{2.217651in}{1.532120in}}%
\pgfpathlineto{\pgfqpoint{2.217661in}{1.532216in}}%
\pgfpathlineto{\pgfqpoint{2.218792in}{2.203909in}}%
\pgfpathlineto{\pgfqpoint{2.218841in}{2.203439in}}%
\pgfpathlineto{\pgfqpoint{2.219286in}{2.204001in}}%
\pgfpathlineto{\pgfqpoint{2.219434in}{2.202838in}}%
\pgfpathlineto{\pgfqpoint{2.219484in}{2.203951in}}%
\pgfpathlineto{\pgfqpoint{2.219533in}{2.202285in}}%
\pgfpathlineto{\pgfqpoint{2.219583in}{2.203904in}}%
\pgfpathlineto{\pgfqpoint{2.220679in}{1.532029in}}%
\pgfpathlineto{\pgfqpoint{2.221314in}{1.532063in}}%
\pgfpathlineto{\pgfqpoint{2.222595in}{1.532117in}}%
\pgfpathlineto{\pgfqpoint{2.222606in}{1.532192in}}%
\pgfpathlineto{\pgfqpoint{2.223837in}{2.203901in}}%
\pgfpathlineto{\pgfqpoint{2.223886in}{2.203450in}}%
\pgfpathlineto{\pgfqpoint{2.224232in}{2.203944in}}%
\pgfpathlineto{\pgfqpoint{2.224480in}{2.202678in}}%
\pgfpathlineto{\pgfqpoint{2.224529in}{2.203843in}}%
\pgfpathlineto{\pgfqpoint{2.224578in}{2.201707in}}%
\pgfpathlineto{\pgfqpoint{2.224628in}{2.203777in}}%
\pgfpathlineto{\pgfqpoint{2.225591in}{1.532024in}}%
\pgfpathlineto{\pgfqpoint{2.225979in}{1.532057in}}%
\pgfpathlineto{\pgfqpoint{2.227541in}{1.532120in}}%
\pgfpathlineto{\pgfqpoint{2.227553in}{1.532242in}}%
\pgfpathlineto{\pgfqpoint{2.228685in}{2.203687in}}%
\pgfpathlineto{\pgfqpoint{2.229823in}{2.202850in}}%
\pgfpathlineto{\pgfqpoint{2.229971in}{2.202372in}}%
\pgfpathlineto{\pgfqpoint{2.230119in}{2.101937in}}%
\pgfpathlineto{\pgfqpoint{2.230564in}{1.532034in}}%
\pgfpathlineto{\pgfqpoint{2.231095in}{1.532056in}}%
\pgfpathlineto{\pgfqpoint{2.232496in}{1.532182in}}%
\pgfpathlineto{\pgfqpoint{2.233995in}{2.203883in}}%
\pgfpathlineto{\pgfqpoint{2.234045in}{2.203516in}}%
\pgfpathlineto{\pgfqpoint{2.234094in}{2.203877in}}%
\pgfpathlineto{\pgfqpoint{2.234638in}{2.202204in}}%
\pgfpathlineto{\pgfqpoint{2.234687in}{2.203525in}}%
\pgfpathlineto{\pgfqpoint{2.235482in}{1.532025in}}%
\pgfpathlineto{\pgfqpoint{2.236363in}{1.532075in}}%
\pgfpathlineto{\pgfqpoint{2.237433in}{1.532121in}}%
\pgfpathlineto{\pgfqpoint{2.237443in}{1.532223in}}%
\pgfpathlineto{\pgfqpoint{2.238574in}{2.203800in}}%
\pgfpathlineto{\pgfqpoint{2.238772in}{2.203862in}}%
\pgfpathlineto{\pgfqpoint{2.239415in}{2.202840in}}%
\pgfpathlineto{\pgfqpoint{2.239465in}{2.203668in}}%
\pgfpathlineto{\pgfqpoint{2.239514in}{2.202347in}}%
\pgfpathlineto{\pgfqpoint{2.239564in}{2.203575in}}%
\pgfpathlineto{\pgfqpoint{2.240475in}{1.532034in}}%
\pgfpathlineto{\pgfqpoint{2.241706in}{1.532082in}}%
\pgfpathlineto{\pgfqpoint{2.242389in}{1.532220in}}%
\pgfpathlineto{\pgfqpoint{2.243584in}{2.203785in}}%
\pgfpathlineto{\pgfqpoint{2.243683in}{2.203815in}}%
\pgfpathlineto{\pgfqpoint{2.244425in}{2.202736in}}%
\pgfpathlineto{\pgfqpoint{2.244474in}{2.203573in}}%
\pgfpathlineto{\pgfqpoint{2.244523in}{2.202111in}}%
\pgfpathlineto{\pgfqpoint{2.244573in}{2.203466in}}%
\pgfpathlineto{\pgfqpoint{2.245381in}{1.532032in}}%
\pgfpathlineto{\pgfqpoint{2.246180in}{1.532066in}}%
\pgfpathlineto{\pgfqpoint{2.247329in}{1.532150in}}%
\pgfpathlineto{\pgfqpoint{2.247338in}{1.532526in}}%
\pgfpathlineto{\pgfqpoint{2.248370in}{2.203589in}}%
\pgfpathlineto{\pgfqpoint{2.249408in}{2.203270in}}%
\pgfpathlineto{\pgfqpoint{2.249754in}{2.202219in}}%
\pgfpathlineto{\pgfqpoint{2.249804in}{2.201371in}}%
\pgfpathlineto{\pgfqpoint{2.250002in}{1.860603in}}%
\pgfpathlineto{\pgfqpoint{2.250348in}{1.532034in}}%
\pgfpathlineto{\pgfqpoint{2.250881in}{1.532056in}}%
\pgfpathlineto{\pgfqpoint{2.252278in}{1.532194in}}%
\pgfpathlineto{\pgfqpoint{2.253677in}{2.203798in}}%
\pgfpathlineto{\pgfqpoint{2.254517in}{2.201777in}}%
\pgfpathlineto{\pgfqpoint{2.254567in}{2.203264in}}%
\pgfpathlineto{\pgfqpoint{2.255310in}{1.532031in}}%
\pgfpathlineto{\pgfqpoint{2.255950in}{1.532061in}}%
\pgfpathlineto{\pgfqpoint{2.257221in}{1.532154in}}%
\pgfpathlineto{\pgfqpoint{2.257250in}{1.608991in}}%
\pgfpathlineto{\pgfqpoint{2.258020in}{2.203489in}}%
\pgfpathlineto{\pgfqpoint{2.258416in}{2.203721in}}%
\pgfpathlineto{\pgfqpoint{2.258762in}{2.203487in}}%
\pgfpathlineto{\pgfqpoint{2.258812in}{2.203759in}}%
\pgfpathlineto{\pgfqpoint{2.259356in}{2.202548in}}%
\pgfpathlineto{\pgfqpoint{2.259405in}{2.203362in}}%
\pgfpathlineto{\pgfqpoint{2.259454in}{2.201457in}}%
\pgfpathlineto{\pgfqpoint{2.259504in}{2.203214in}}%
\pgfpathlineto{\pgfqpoint{2.260192in}{1.532014in}}%
\pgfpathlineto{\pgfqpoint{2.260804in}{1.532063in}}%
\pgfpathlineto{\pgfqpoint{2.262160in}{1.532121in}}%
\pgfpathlineto{\pgfqpoint{2.262170in}{1.532220in}}%
\pgfpathlineto{\pgfqpoint{2.263300in}{2.203762in}}%
\pgfpathlineto{\pgfqpoint{2.263349in}{2.203364in}}%
\pgfpathlineto{\pgfqpoint{2.263794in}{2.203851in}}%
\pgfpathlineto{\pgfqpoint{2.264042in}{2.202649in}}%
\pgfpathlineto{\pgfqpoint{2.264091in}{2.203752in}}%
\pgfpathlineto{\pgfqpoint{2.264141in}{2.201925in}}%
\pgfpathlineto{\pgfqpoint{2.264190in}{2.203689in}}%
\pgfpathlineto{\pgfqpoint{2.265179in}{1.532028in}}%
\pgfpathlineto{\pgfqpoint{2.265659in}{1.532045in}}%
\pgfpathlineto{\pgfqpoint{2.267116in}{1.532218in}}%
\pgfpathlineto{\pgfqpoint{2.268246in}{2.203759in}}%
\pgfpathlineto{\pgfqpoint{2.268295in}{2.203309in}}%
\pgfpathlineto{\pgfqpoint{2.268741in}{2.203848in}}%
\pgfpathlineto{\pgfqpoint{2.268889in}{2.202755in}}%
\pgfpathlineto{\pgfqpoint{2.268938in}{2.203796in}}%
\pgfpathlineto{\pgfqpoint{2.269087in}{2.182470in}}%
\pgfpathlineto{\pgfqpoint{2.269136in}{2.203672in}}%
\pgfpathlineto{\pgfqpoint{2.270143in}{1.532026in}}%
\pgfpathlineto{\pgfqpoint{2.270532in}{1.532041in}}%
\pgfpathlineto{\pgfqpoint{2.272062in}{1.532223in}}%
\pgfpathlineto{\pgfqpoint{2.273191in}{2.203732in}}%
\pgfpathlineto{\pgfqpoint{2.273240in}{2.203305in}}%
\pgfpathlineto{\pgfqpoint{2.273685in}{2.203820in}}%
\pgfpathlineto{\pgfqpoint{2.273834in}{2.202808in}}%
\pgfpathlineto{\pgfqpoint{2.273883in}{2.203767in}}%
\pgfpathlineto{\pgfqpoint{2.274032in}{2.199475in}}%
\pgfpathlineto{\pgfqpoint{2.274081in}{2.203654in}}%
\pgfpathlineto{\pgfqpoint{2.275045in}{1.532025in}}%
\pgfpathlineto{\pgfqpoint{2.275333in}{1.532053in}}%
\pgfpathlineto{\pgfqpoint{2.276996in}{1.532121in}}%
\pgfpathlineto{\pgfqpoint{2.277007in}{1.532224in}}%
\pgfpathlineto{\pgfqpoint{2.278137in}{2.203659in}}%
\pgfpathlineto{\pgfqpoint{2.278335in}{2.203722in}}%
\pgfpathlineto{\pgfqpoint{2.278978in}{2.202777in}}%
\pgfpathlineto{\pgfqpoint{2.279028in}{2.203527in}}%
\pgfpathlineto{\pgfqpoint{2.279176in}{2.195185in}}%
\pgfpathlineto{\pgfqpoint{2.279225in}{2.203307in}}%
\pgfpathlineto{\pgfqpoint{2.280026in}{1.532034in}}%
\pgfpathlineto{\pgfqpoint{2.280514in}{1.532051in}}%
\pgfpathlineto{\pgfqpoint{2.281953in}{1.532226in}}%
\pgfpathlineto{\pgfqpoint{2.283151in}{2.203654in}}%
\pgfpathlineto{\pgfqpoint{2.283250in}{2.203683in}}%
\pgfpathlineto{\pgfqpoint{2.283992in}{2.202536in}}%
\pgfpathlineto{\pgfqpoint{2.284041in}{2.203436in}}%
\pgfpathlineto{\pgfqpoint{2.284091in}{2.201746in}}%
\pgfpathlineto{\pgfqpoint{2.284140in}{2.203327in}}%
\pgfpathlineto{\pgfqpoint{2.284953in}{1.532031in}}%
\pgfpathlineto{\pgfqpoint{2.285570in}{1.532060in}}%
\pgfpathlineto{\pgfqpoint{2.286893in}{1.532148in}}%
\pgfpathlineto{\pgfqpoint{2.286901in}{1.532492in}}%
\pgfpathlineto{\pgfqpoint{2.287960in}{2.203497in}}%
\pgfpathlineto{\pgfqpoint{2.289196in}{2.202361in}}%
\pgfpathlineto{\pgfqpoint{2.289246in}{2.202740in}}%
\pgfpathlineto{\pgfqpoint{2.289295in}{2.201044in}}%
\pgfpathlineto{\pgfqpoint{2.289345in}{2.202377in}}%
\pgfpathlineto{\pgfqpoint{2.289865in}{1.532015in}}%
\pgfpathlineto{\pgfqpoint{2.290608in}{1.532052in}}%
\pgfpathlineto{\pgfqpoint{2.291844in}{1.532247in}}%
\pgfpathlineto{\pgfqpoint{2.292946in}{2.203499in}}%
\pgfpathlineto{\pgfqpoint{2.294182in}{2.202191in}}%
\pgfpathlineto{\pgfqpoint{2.294232in}{2.202518in}}%
\pgfpathlineto{\pgfqpoint{2.294429in}{2.052237in}}%
\pgfpathlineto{\pgfqpoint{2.294900in}{1.532034in}}%
\pgfpathlineto{\pgfqpoint{2.295390in}{1.532052in}}%
\pgfpathlineto{\pgfqpoint{2.296787in}{1.532184in}}%
\pgfpathlineto{\pgfqpoint{2.298105in}{2.203675in}}%
\pgfpathlineto{\pgfqpoint{2.298847in}{2.202081in}}%
\pgfpathlineto{\pgfqpoint{2.298204in}{2.203693in}}%
\pgfpathlineto{\pgfqpoint{2.298896in}{2.203442in}}%
\pgfpathlineto{\pgfqpoint{2.299810in}{1.532028in}}%
\pgfpathlineto{\pgfqpoint{2.300743in}{1.532075in}}%
\pgfpathlineto{\pgfqpoint{2.301723in}{1.532119in}}%
\pgfpathlineto{\pgfqpoint{2.301734in}{1.532221in}}%
\pgfpathlineto{\pgfqpoint{2.302863in}{2.203633in}}%
\pgfpathlineto{\pgfqpoint{2.302913in}{2.203237in}}%
\pgfpathlineto{\pgfqpoint{2.303358in}{2.203717in}}%
\pgfpathlineto{\pgfqpoint{2.303605in}{2.202508in}}%
\pgfpathlineto{\pgfqpoint{2.303655in}{2.203615in}}%
\pgfpathlineto{\pgfqpoint{2.303704in}{2.201748in}}%
\pgfpathlineto{\pgfqpoint{2.303754in}{2.203550in}}%
\pgfpathlineto{\pgfqpoint{2.304718in}{1.532025in}}%
\pgfpathlineto{\pgfqpoint{2.305203in}{1.532060in}}%
\pgfpathlineto{\pgfqpoint{2.306669in}{1.532121in}}%
\pgfpathlineto{\pgfqpoint{2.306680in}{1.532224in}}%
\pgfpathlineto{\pgfqpoint{2.307811in}{2.203580in}}%
\pgfpathlineto{\pgfqpoint{2.308008in}{2.203645in}}%
\pgfpathlineto{\pgfqpoint{2.308651in}{2.202665in}}%
\pgfpathlineto{\pgfqpoint{2.308701in}{2.203459in}}%
\pgfpathlineto{\pgfqpoint{2.308750in}{2.202225in}}%
\pgfpathlineto{\pgfqpoint{2.308800in}{2.203368in}}%
\pgfpathlineto{\pgfqpoint{2.309393in}{1.808405in}}%
\pgfpathlineto{\pgfqpoint{2.309677in}{1.532031in}}%
\pgfpathlineto{\pgfqpoint{2.310291in}{1.532060in}}%
\pgfpathlineto{\pgfqpoint{2.311620in}{1.532147in}}%
\pgfpathlineto{\pgfqpoint{2.311629in}{1.532491in}}%
\pgfpathlineto{\pgfqpoint{2.312688in}{2.203437in}}%
\pgfpathlineto{\pgfqpoint{2.313924in}{2.202240in}}%
\pgfpathlineto{\pgfqpoint{2.313973in}{2.202655in}}%
\pgfpathlineto{\pgfqpoint{2.314023in}{2.199668in}}%
\pgfpathlineto{\pgfqpoint{2.314072in}{2.202264in}}%
\pgfpathlineto{\pgfqpoint{2.314593in}{1.532025in}}%
\pgfpathlineto{\pgfqpoint{2.315292in}{1.532058in}}%
\pgfpathlineto{\pgfqpoint{2.316574in}{1.532487in}}%
\pgfpathlineto{\pgfqpoint{2.317633in}{2.203410in}}%
\pgfpathlineto{\pgfqpoint{2.318870in}{2.202212in}}%
\pgfpathlineto{\pgfqpoint{2.318919in}{2.202643in}}%
\pgfpathlineto{\pgfqpoint{2.318969in}{2.199494in}}%
\pgfpathlineto{\pgfqpoint{2.319018in}{2.202262in}}%
\pgfpathlineto{\pgfqpoint{2.319538in}{1.532021in}}%
\pgfpathlineto{\pgfqpoint{2.320264in}{1.532065in}}%
\pgfpathlineto{\pgfqpoint{2.321504in}{1.532117in}}%
\pgfpathlineto{\pgfqpoint{2.321515in}{1.532197in}}%
\pgfpathlineto{\pgfqpoint{2.322729in}{2.203539in}}%
\pgfpathlineto{\pgfqpoint{2.322778in}{2.203230in}}%
\pgfpathlineto{\pgfqpoint{2.323124in}{2.203580in}}%
\pgfpathlineto{\pgfqpoint{2.323372in}{2.202813in}}%
\pgfpathlineto{\pgfqpoint{2.323421in}{2.203471in}}%
\pgfpathlineto{\pgfqpoint{2.323668in}{2.185020in}}%
\pgfpathlineto{\pgfqpoint{2.323718in}{2.203182in}}%
\pgfpathlineto{\pgfqpoint{2.324500in}{1.532024in}}%
\pgfpathlineto{\pgfqpoint{2.325087in}{1.532065in}}%
\pgfpathlineto{\pgfqpoint{2.326450in}{1.532119in}}%
\pgfpathlineto{\pgfqpoint{2.326462in}{1.532225in}}%
\pgfpathlineto{\pgfqpoint{2.327572in}{2.203426in}}%
\pgfpathlineto{\pgfqpoint{2.328610in}{2.202616in}}%
\pgfpathlineto{\pgfqpoint{2.328660in}{2.203026in}}%
\pgfpathlineto{\pgfqpoint{2.328808in}{2.200168in}}%
\pgfpathlineto{\pgfqpoint{2.328858in}{2.202614in}}%
\pgfpathlineto{\pgfqpoint{2.329459in}{1.532031in}}%
\pgfpathlineto{\pgfqpoint{2.330123in}{1.532055in}}%
\pgfpathlineto{\pgfqpoint{2.331411in}{1.532492in}}%
\pgfpathlineto{\pgfqpoint{2.332469in}{2.203365in}}%
\pgfpathlineto{\pgfqpoint{2.333706in}{2.202200in}}%
\pgfpathlineto{\pgfqpoint{2.333755in}{2.202599in}}%
\pgfpathlineto{\pgfqpoint{2.333805in}{2.200543in}}%
\pgfpathlineto{\pgfqpoint{2.333854in}{2.202227in}}%
\pgfpathlineto{\pgfqpoint{2.334375in}{1.532020in}}%
\pgfpathlineto{\pgfqpoint{2.335124in}{1.532066in}}%
\pgfpathlineto{\pgfqpoint{2.336341in}{1.532118in}}%
\pgfpathlineto{\pgfqpoint{2.336350in}{1.532182in}}%
\pgfpathlineto{\pgfqpoint{2.337852in}{2.203519in}}%
\pgfpathlineto{\pgfqpoint{2.338594in}{2.198078in}}%
\pgfpathlineto{\pgfqpoint{2.338643in}{2.202956in}}%
\pgfpathlineto{\pgfqpoint{2.339353in}{1.532034in}}%
\pgfpathlineto{\pgfqpoint{2.339871in}{1.532054in}}%
\pgfpathlineto{\pgfqpoint{2.341296in}{1.532191in}}%
\pgfpathlineto{\pgfqpoint{2.342603in}{2.203519in}}%
\pgfpathlineto{\pgfqpoint{2.342702in}{2.203536in}}%
\pgfpathlineto{\pgfqpoint{2.343443in}{2.199310in}}%
\pgfpathlineto{\pgfqpoint{2.343493in}{2.203174in}}%
\pgfpathlineto{\pgfqpoint{2.344324in}{1.532027in}}%
\pgfpathlineto{\pgfqpoint{2.344713in}{1.532041in}}%
\pgfpathlineto{\pgfqpoint{2.346243in}{1.532224in}}%
\pgfpathlineto{\pgfqpoint{2.347373in}{2.203504in}}%
\pgfpathlineto{\pgfqpoint{2.347422in}{2.203046in}}%
\pgfpathlineto{\pgfqpoint{2.347867in}{2.203599in}}%
\pgfpathlineto{\pgfqpoint{2.348016in}{2.202492in}}%
\pgfpathlineto{\pgfqpoint{2.348065in}{2.203550in}}%
\pgfpathlineto{\pgfqpoint{2.348214in}{2.180596in}}%
\pgfpathlineto{\pgfqpoint{2.348263in}{2.203428in}}%
\pgfpathlineto{\pgfqpoint{2.349242in}{1.532030in}}%
\pgfpathlineto{\pgfqpoint{2.349662in}{1.532054in}}%
\pgfpathlineto{\pgfqpoint{2.351176in}{1.532116in}}%
\pgfpathlineto{\pgfqpoint{2.351187in}{1.532185in}}%
\pgfpathlineto{\pgfqpoint{2.352491in}{2.203484in}}%
\pgfpathlineto{\pgfqpoint{2.352541in}{2.203153in}}%
\pgfpathlineto{\pgfqpoint{2.352788in}{2.203497in}}%
\pgfpathlineto{\pgfqpoint{2.353134in}{2.202542in}}%
\pgfpathlineto{\pgfqpoint{2.353184in}{2.203322in}}%
\pgfpathlineto{\pgfqpoint{2.353332in}{2.196236in}}%
\pgfpathlineto{\pgfqpoint{2.353381in}{2.203123in}}%
\pgfpathlineto{\pgfqpoint{2.354190in}{1.532031in}}%
\pgfpathlineto{\pgfqpoint{2.354661in}{1.532048in}}%
\pgfpathlineto{\pgfqpoint{2.356138in}{1.532492in}}%
\pgfpathlineto{\pgfqpoint{2.357196in}{2.203286in}}%
\pgfpathlineto{\pgfqpoint{2.358433in}{2.202151in}}%
\pgfpathlineto{\pgfqpoint{2.358482in}{2.202535in}}%
\pgfpathlineto{\pgfqpoint{2.358532in}{2.200832in}}%
\pgfpathlineto{\pgfqpoint{2.358581in}{2.202176in}}%
\pgfpathlineto{\pgfqpoint{2.359102in}{1.532014in}}%
\pgfpathlineto{\pgfqpoint{2.359850in}{1.532052in}}%
\pgfpathlineto{\pgfqpoint{2.361079in}{1.532215in}}%
\pgfpathlineto{\pgfqpoint{2.362210in}{2.203448in}}%
\pgfpathlineto{\pgfqpoint{2.362259in}{2.202989in}}%
\pgfpathlineto{\pgfqpoint{2.362704in}{2.203539in}}%
\pgfpathlineto{\pgfqpoint{2.362853in}{2.202423in}}%
\pgfpathlineto{\pgfqpoint{2.362902in}{2.203489in}}%
\pgfpathlineto{\pgfqpoint{2.363050in}{2.178040in}}%
\pgfpathlineto{\pgfqpoint{2.363100in}{2.203361in}}%
\pgfpathlineto{\pgfqpoint{2.364099in}{1.532029in}}%
\pgfpathlineto{\pgfqpoint{2.364536in}{1.532056in}}%
\pgfpathlineto{\pgfqpoint{2.366013in}{1.532117in}}%
\pgfpathlineto{\pgfqpoint{2.366024in}{1.532200in}}%
\pgfpathlineto{\pgfqpoint{2.367257in}{2.203459in}}%
\pgfpathlineto{\pgfqpoint{2.367899in}{2.201927in}}%
\pgfpathlineto{\pgfqpoint{2.367454in}{2.203507in}}%
\pgfpathlineto{\pgfqpoint{2.367949in}{2.203409in}}%
\pgfpathlineto{\pgfqpoint{2.369048in}{1.532028in}}%
\pgfpathlineto{\pgfqpoint{2.370079in}{1.532079in}}%
\pgfpathlineto{\pgfqpoint{2.370959in}{1.532119in}}%
\pgfpathlineto{\pgfqpoint{2.370971in}{1.532234in}}%
\pgfpathlineto{\pgfqpoint{2.372101in}{2.203339in}}%
\pgfpathlineto{\pgfqpoint{2.372298in}{2.203403in}}%
\pgfpathlineto{\pgfqpoint{2.372941in}{2.202517in}}%
\pgfpathlineto{\pgfqpoint{2.372991in}{2.203207in}}%
\pgfpathlineto{\pgfqpoint{2.373139in}{2.200959in}}%
\pgfpathlineto{\pgfqpoint{2.373189in}{2.202997in}}%
\pgfpathlineto{\pgfqpoint{2.373994in}{1.532034in}}%
\pgfpathlineto{\pgfqpoint{2.374485in}{1.532052in}}%
\pgfpathlineto{\pgfqpoint{2.375920in}{1.532465in}}%
\pgfpathlineto{\pgfqpoint{2.376985in}{2.203221in}}%
\pgfpathlineto{\pgfqpoint{2.378221in}{2.201767in}}%
\pgfpathlineto{\pgfqpoint{2.378270in}{2.202482in}}%
\pgfpathlineto{\pgfqpoint{2.378883in}{1.532015in}}%
\pgfpathlineto{\pgfqpoint{2.380232in}{1.532075in}}%
\pgfpathlineto{\pgfqpoint{2.380861in}{1.532218in}}%
\pgfpathlineto{\pgfqpoint{2.381991in}{2.203366in}}%
\pgfpathlineto{\pgfqpoint{2.382041in}{2.202941in}}%
\pgfpathlineto{\pgfqpoint{2.382486in}{2.203459in}}%
\pgfpathlineto{\pgfqpoint{2.382634in}{2.202465in}}%
\pgfpathlineto{\pgfqpoint{2.382683in}{2.203408in}}%
\pgfpathlineto{\pgfqpoint{2.382832in}{2.200641in}}%
\pgfpathlineto{\pgfqpoint{2.382881in}{2.203301in}}%
\pgfpathlineto{\pgfqpoint{2.383863in}{1.532029in}}%
\pgfpathlineto{\pgfqpoint{2.384138in}{1.532038in}}%
\pgfpathlineto{\pgfqpoint{2.385806in}{1.532197in}}%
\pgfpathlineto{\pgfqpoint{2.387021in}{2.203332in}}%
\pgfpathlineto{\pgfqpoint{2.387070in}{2.203018in}}%
\pgfpathlineto{\pgfqpoint{2.387417in}{2.203370in}}%
\pgfpathlineto{\pgfqpoint{2.387664in}{2.202583in}}%
\pgfpathlineto{\pgfqpoint{2.387713in}{2.203258in}}%
\pgfpathlineto{\pgfqpoint{2.387862in}{2.201832in}}%
\pgfpathlineto{\pgfqpoint{2.387911in}{2.203104in}}%
\pgfpathlineto{\pgfqpoint{2.388800in}{1.532032in}}%
\pgfpathlineto{\pgfqpoint{2.389994in}{1.532081in}}%
\pgfpathlineto{\pgfqpoint{2.390747in}{1.532150in}}%
\pgfpathlineto{\pgfqpoint{2.390756in}{1.532546in}}%
\pgfpathlineto{\pgfqpoint{2.391781in}{2.203104in}}%
\pgfpathlineto{\pgfqpoint{2.392721in}{2.202956in}}%
\pgfpathlineto{\pgfqpoint{2.392770in}{2.202868in}}%
\pgfpathlineto{\pgfqpoint{2.393166in}{2.201718in}}%
\pgfpathlineto{\pgfqpoint{2.393215in}{2.201474in}}%
\pgfpathlineto{\pgfqpoint{2.393708in}{1.588179in}}%
\pgfpathlineto{\pgfqpoint{2.393812in}{1.532035in}}%
\pgfpathlineto{\pgfqpoint{2.394452in}{1.532061in}}%
\pgfpathlineto{\pgfqpoint{2.395698in}{1.532213in}}%
\pgfpathlineto{\pgfqpoint{2.396903in}{2.203298in}}%
\pgfpathlineto{\pgfqpoint{2.397101in}{2.203347in}}%
\pgfpathlineto{\pgfqpoint{2.397744in}{2.202171in}}%
\pgfpathlineto{\pgfqpoint{2.397793in}{2.203092in}}%
\pgfpathlineto{\pgfqpoint{2.397842in}{2.201349in}}%
\pgfpathlineto{\pgfqpoint{2.397892in}{2.202985in}}%
\pgfpathlineto{\pgfqpoint{2.398723in}{1.532027in}}%
\pgfpathlineto{\pgfqpoint{2.399309in}{1.532049in}}%
\pgfpathlineto{\pgfqpoint{2.400643in}{1.532215in}}%
\pgfpathlineto{\pgfqpoint{2.401773in}{2.203327in}}%
\pgfpathlineto{\pgfqpoint{2.401823in}{2.202849in}}%
\pgfpathlineto{\pgfqpoint{2.402268in}{2.203421in}}%
\pgfpathlineto{\pgfqpoint{2.402416in}{2.202233in}}%
\pgfpathlineto{\pgfqpoint{2.402466in}{2.203373in}}%
\pgfpathlineto{\pgfqpoint{2.402515in}{2.201640in}}%
\pgfpathlineto{\pgfqpoint{2.402565in}{2.203327in}}%
\pgfpathlineto{\pgfqpoint{2.403657in}{1.532028in}}%
\pgfpathlineto{\pgfqpoint{2.404191in}{1.532060in}}%
\pgfpathlineto{\pgfqpoint{2.405577in}{1.532119in}}%
\pgfpathlineto{\pgfqpoint{2.405589in}{1.532231in}}%
\pgfpathlineto{\pgfqpoint{2.406720in}{2.203258in}}%
\pgfpathlineto{\pgfqpoint{2.406918in}{2.203322in}}%
\pgfpathlineto{\pgfqpoint{2.407561in}{2.202322in}}%
\pgfpathlineto{\pgfqpoint{2.407610in}{2.203136in}}%
\pgfpathlineto{\pgfqpoint{2.407660in}{2.201857in}}%
\pgfpathlineto{\pgfqpoint{2.407709in}{2.203046in}}%
\pgfpathlineto{\pgfqpoint{2.408544in}{1.584740in}}%
\pgfpathlineto{\pgfqpoint{2.408585in}{1.532031in}}%
\pgfpathlineto{\pgfqpoint{2.409298in}{1.532063in}}%
\pgfpathlineto{\pgfqpoint{2.410529in}{1.532148in}}%
\pgfpathlineto{\pgfqpoint{2.410538in}{1.532491in}}%
\pgfpathlineto{\pgfqpoint{2.411597in}{2.203117in}}%
\pgfpathlineto{\pgfqpoint{2.412833in}{2.201934in}}%
\pgfpathlineto{\pgfqpoint{2.412883in}{2.202349in}}%
\pgfpathlineto{\pgfqpoint{2.412932in}{2.199878in}}%
\pgfpathlineto{\pgfqpoint{2.412982in}{2.201973in}}%
\pgfpathlineto{\pgfqpoint{2.413502in}{1.532020in}}%
\pgfpathlineto{\pgfqpoint{2.414235in}{1.532065in}}%
\pgfpathlineto{\pgfqpoint{2.415468in}{1.532117in}}%
\pgfpathlineto{\pgfqpoint{2.415479in}{1.532200in}}%
\pgfpathlineto{\pgfqpoint{2.416690in}{2.203251in}}%
\pgfpathlineto{\pgfqpoint{2.416740in}{2.202940in}}%
\pgfpathlineto{\pgfqpoint{2.417086in}{2.203295in}}%
\pgfpathlineto{\pgfqpoint{2.417333in}{2.202524in}}%
\pgfpathlineto{\pgfqpoint{2.417383in}{2.203189in}}%
\pgfpathlineto{\pgfqpoint{2.417531in}{2.201827in}}%
\pgfpathlineto{\pgfqpoint{2.417581in}{2.203039in}}%
\pgfpathlineto{\pgfqpoint{2.418434in}{1.599281in}}%
\pgfpathlineto{\pgfqpoint{2.418463in}{1.532025in}}%
\pgfpathlineto{\pgfqpoint{2.419196in}{1.532051in}}%
\pgfpathlineto{\pgfqpoint{2.420425in}{1.532224in}}%
\pgfpathlineto{\pgfqpoint{2.421557in}{2.203220in}}%
\pgfpathlineto{\pgfqpoint{2.421755in}{2.203287in}}%
\pgfpathlineto{\pgfqpoint{2.422299in}{2.202518in}}%
\pgfpathlineto{\pgfqpoint{2.422348in}{2.203181in}}%
\pgfpathlineto{\pgfqpoint{2.422496in}{2.201786in}}%
\pgfpathlineto{\pgfqpoint{2.422546in}{2.203025in}}%
\pgfpathlineto{\pgfqpoint{2.423452in}{1.532026in}}%
\pgfpathlineto{\pgfqpoint{2.424731in}{1.532074in}}%
\pgfpathlineto{\pgfqpoint{2.425371in}{1.532223in}}%
\pgfpathlineto{\pgfqpoint{2.426500in}{2.203261in}}%
\pgfpathlineto{\pgfqpoint{2.426550in}{2.202814in}}%
\pgfpathlineto{\pgfqpoint{2.426995in}{2.203352in}}%
\pgfpathlineto{\pgfqpoint{2.427143in}{2.202277in}}%
\pgfpathlineto{\pgfqpoint{2.427192in}{2.203302in}}%
\pgfpathlineto{\pgfqpoint{2.427341in}{2.186181in}}%
\pgfpathlineto{\pgfqpoint{2.427390in}{2.203184in}}%
\pgfpathlineto{\pgfqpoint{2.428391in}{1.532029in}}%
\pgfpathlineto{\pgfqpoint{2.428730in}{1.532052in}}%
\pgfpathlineto{\pgfqpoint{2.430304in}{1.532118in}}%
\pgfpathlineto{\pgfqpoint{2.430316in}{1.532211in}}%
\pgfpathlineto{\pgfqpoint{2.431447in}{2.203248in}}%
\pgfpathlineto{\pgfqpoint{2.431496in}{2.202731in}}%
\pgfpathlineto{\pgfqpoint{2.431942in}{2.203343in}}%
\pgfpathlineto{\pgfqpoint{2.432090in}{2.201969in}}%
\pgfpathlineto{\pgfqpoint{2.432139in}{2.203296in}}%
\pgfpathlineto{\pgfqpoint{2.432189in}{2.200525in}}%
\pgfpathlineto{\pgfqpoint{2.432238in}{2.203249in}}%
\pgfpathlineto{\pgfqpoint{2.433300in}{1.532025in}}%
\pgfpathlineto{\pgfqpoint{2.433489in}{1.532049in}}%
\pgfpathlineto{\pgfqpoint{2.435251in}{1.532121in}}%
\pgfpathlineto{\pgfqpoint{2.435262in}{1.532224in}}%
\pgfpathlineto{\pgfqpoint{2.436392in}{2.203165in}}%
\pgfpathlineto{\pgfqpoint{2.436590in}{2.203229in}}%
\pgfpathlineto{\pgfqpoint{2.437233in}{2.202269in}}%
\pgfpathlineto{\pgfqpoint{2.437282in}{2.203042in}}%
\pgfpathlineto{\pgfqpoint{2.437332in}{2.201848in}}%
\pgfpathlineto{\pgfqpoint{2.437381in}{2.202951in}}%
\pgfpathlineto{\pgfqpoint{2.437925in}{1.921329in}}%
\pgfpathlineto{\pgfqpoint{2.438252in}{1.532031in}}%
\pgfpathlineto{\pgfqpoint{2.438846in}{1.532058in}}%
\pgfpathlineto{\pgfqpoint{2.440202in}{1.532151in}}%
\pgfpathlineto{\pgfqpoint{2.440211in}{1.532622in}}%
\pgfpathlineto{\pgfqpoint{2.441243in}{2.202975in}}%
\pgfpathlineto{\pgfqpoint{2.442183in}{2.202812in}}%
\pgfpathlineto{\pgfqpoint{2.442678in}{2.201211in}}%
\pgfpathlineto{\pgfqpoint{2.443271in}{1.532034in}}%
\pgfpathlineto{\pgfqpoint{2.444260in}{1.532071in}}%
\pgfpathlineto{\pgfqpoint{2.445151in}{1.532185in}}%
\pgfpathlineto{\pgfqpoint{2.446469in}{2.203220in}}%
\pgfpathlineto{\pgfqpoint{2.447211in}{2.201628in}}%
\pgfpathlineto{\pgfqpoint{2.446568in}{2.203239in}}%
\pgfpathlineto{\pgfqpoint{2.447261in}{2.202991in}}%
\pgfpathlineto{\pgfqpoint{2.448169in}{1.532028in}}%
\pgfpathlineto{\pgfqpoint{2.449099in}{1.532075in}}%
\pgfpathlineto{\pgfqpoint{2.450087in}{1.532119in}}%
\pgfpathlineto{\pgfqpoint{2.450098in}{1.532235in}}%
\pgfpathlineto{\pgfqpoint{2.451227in}{2.203129in}}%
\pgfpathlineto{\pgfqpoint{2.451425in}{2.203192in}}%
\pgfpathlineto{\pgfqpoint{2.452068in}{2.202351in}}%
\pgfpathlineto{\pgfqpoint{2.452117in}{2.202992in}}%
\pgfpathlineto{\pgfqpoint{2.452266in}{2.201200in}}%
\pgfpathlineto{\pgfqpoint{2.452315in}{2.202779in}}%
\pgfpathlineto{\pgfqpoint{2.453096in}{1.532032in}}%
\pgfpathlineto{\pgfqpoint{2.453700in}{1.532058in}}%
\pgfpathlineto{\pgfqpoint{2.455043in}{1.532204in}}%
\pgfpathlineto{\pgfqpoint{2.456266in}{2.203188in}}%
\pgfpathlineto{\pgfqpoint{2.456464in}{2.203235in}}%
\pgfpathlineto{\pgfqpoint{2.457107in}{2.198808in}}%
\pgfpathlineto{\pgfqpoint{2.457157in}{2.202979in}}%
\pgfpathlineto{\pgfqpoint{2.458053in}{1.532030in}}%
\pgfpathlineto{\pgfqpoint{2.458385in}{1.532050in}}%
\pgfpathlineto{\pgfqpoint{2.459984in}{1.532147in}}%
\pgfpathlineto{\pgfqpoint{2.459992in}{1.532470in}}%
\pgfpathlineto{\pgfqpoint{2.461047in}{2.202985in}}%
\pgfpathlineto{\pgfqpoint{2.462284in}{2.201867in}}%
\pgfpathlineto{\pgfqpoint{2.462333in}{2.202212in}}%
\pgfpathlineto{\pgfqpoint{2.462382in}{2.200688in}}%
\pgfpathlineto{\pgfqpoint{2.462432in}{2.201836in}}%
\pgfpathlineto{\pgfqpoint{2.462956in}{1.532020in}}%
\pgfpathlineto{\pgfqpoint{2.463706in}{1.532066in}}%
\pgfpathlineto{\pgfqpoint{2.464922in}{1.532118in}}%
\pgfpathlineto{\pgfqpoint{2.464932in}{1.532182in}}%
\pgfpathlineto{\pgfqpoint{2.466434in}{2.203148in}}%
\pgfpathlineto{\pgfqpoint{2.467175in}{2.200196in}}%
\pgfpathlineto{\pgfqpoint{2.467225in}{2.202598in}}%
\pgfpathlineto{\pgfqpoint{2.467935in}{1.532034in}}%
\pgfpathlineto{\pgfqpoint{2.468504in}{1.532053in}}%
\pgfpathlineto{\pgfqpoint{2.469879in}{1.532198in}}%
\pgfpathlineto{\pgfqpoint{2.471084in}{2.203118in}}%
\pgfpathlineto{\pgfqpoint{2.471134in}{2.202833in}}%
\pgfpathlineto{\pgfqpoint{2.471480in}{2.203162in}}%
\pgfpathlineto{\pgfqpoint{2.471727in}{2.202462in}}%
\pgfpathlineto{\pgfqpoint{2.471777in}{2.203056in}}%
\pgfpathlineto{\pgfqpoint{2.472024in}{2.200620in}}%
\pgfpathlineto{\pgfqpoint{2.472073in}{2.202798in}}%
\pgfpathlineto{\pgfqpoint{2.472863in}{1.532026in}}%
\pgfpathlineto{\pgfqpoint{2.473348in}{1.532060in}}%
\pgfpathlineto{\pgfqpoint{2.474815in}{1.532122in}}%
\pgfpathlineto{\pgfqpoint{2.474825in}{1.532224in}}%
\pgfpathlineto{\pgfqpoint{2.475958in}{2.203111in}}%
\pgfpathlineto{\pgfqpoint{2.476255in}{2.203195in}}%
\pgfpathlineto{\pgfqpoint{2.476700in}{2.202322in}}%
\pgfpathlineto{\pgfqpoint{2.476750in}{2.203081in}}%
\pgfpathlineto{\pgfqpoint{2.476898in}{2.201278in}}%
\pgfpathlineto{\pgfqpoint{2.476947in}{2.202927in}}%
\pgfpathlineto{\pgfqpoint{2.477847in}{1.532028in}}%
\pgfpathlineto{\pgfqpoint{2.478384in}{1.532061in}}%
\pgfpathlineto{\pgfqpoint{2.479759in}{1.532119in}}%
\pgfpathlineto{\pgfqpoint{2.479771in}{1.532227in}}%
\pgfpathlineto{\pgfqpoint{2.480904in}{2.203118in}}%
\pgfpathlineto{\pgfqpoint{2.481102in}{2.203184in}}%
\pgfpathlineto{\pgfqpoint{2.481646in}{2.202289in}}%
\pgfpathlineto{\pgfqpoint{2.481696in}{2.203082in}}%
\pgfpathlineto{\pgfqpoint{2.481844in}{2.200993in}}%
\pgfpathlineto{\pgfqpoint{2.481893in}{2.202927in}}%
\pgfpathlineto{\pgfqpoint{2.482754in}{1.532026in}}%
\pgfpathlineto{\pgfqpoint{2.483238in}{1.532060in}}%
\pgfpathlineto{\pgfqpoint{2.484706in}{1.532122in}}%
\pgfpathlineto{\pgfqpoint{2.484716in}{1.532224in}}%
\pgfpathlineto{\pgfqpoint{2.485844in}{2.203127in}}%
\pgfpathlineto{\pgfqpoint{2.486141in}{2.203214in}}%
\pgfpathlineto{\pgfqpoint{2.486586in}{2.202242in}}%
\pgfpathlineto{\pgfqpoint{2.486636in}{2.203108in}}%
\pgfpathlineto{\pgfqpoint{2.486784in}{2.199415in}}%
\pgfpathlineto{\pgfqpoint{2.486834in}{2.202958in}}%
\pgfpathlineto{\pgfqpoint{2.487733in}{1.532030in}}%
\pgfpathlineto{\pgfqpoint{2.488072in}{1.532051in}}%
\pgfpathlineto{\pgfqpoint{2.489658in}{1.532166in}}%
\pgfpathlineto{\pgfqpoint{2.489700in}{1.869509in}}%
\pgfpathlineto{\pgfqpoint{2.490462in}{2.202871in}}%
\pgfpathlineto{\pgfqpoint{2.490858in}{2.203106in}}%
\pgfpathlineto{\pgfqpoint{2.491204in}{2.202843in}}%
\pgfpathlineto{\pgfqpoint{2.491253in}{2.203148in}}%
\pgfpathlineto{\pgfqpoint{2.491797in}{2.201795in}}%
\pgfpathlineto{\pgfqpoint{2.491847in}{2.202767in}}%
\pgfpathlineto{\pgfqpoint{2.492242in}{2.078973in}}%
\pgfpathlineto{\pgfqpoint{2.492661in}{1.532034in}}%
\pgfpathlineto{\pgfqpoint{2.493277in}{1.532058in}}%
\pgfpathlineto{\pgfqpoint{2.494606in}{1.532193in}}%
\pgfpathlineto{\pgfqpoint{2.495911in}{2.203150in}}%
\pgfpathlineto{\pgfqpoint{2.496010in}{2.203168in}}%
\pgfpathlineto{\pgfqpoint{2.496752in}{2.199795in}}%
\pgfpathlineto{\pgfqpoint{2.496801in}{2.202800in}}%
\pgfpathlineto{\pgfqpoint{2.497635in}{1.532026in}}%
\pgfpathlineto{\pgfqpoint{2.498073in}{1.532059in}}%
\pgfpathlineto{\pgfqpoint{2.499542in}{1.532122in}}%
\pgfpathlineto{\pgfqpoint{2.499552in}{1.532224in}}%
\pgfpathlineto{\pgfqpoint{2.500681in}{2.203132in}}%
\pgfpathlineto{\pgfqpoint{2.500978in}{2.203222in}}%
\pgfpathlineto{\pgfqpoint{2.501423in}{2.202033in}}%
\pgfpathlineto{\pgfqpoint{2.501473in}{2.203126in}}%
\pgfpathlineto{\pgfqpoint{2.501522in}{2.201338in}}%
\pgfpathlineto{\pgfqpoint{2.501572in}{2.203063in}}%
\pgfpathlineto{\pgfqpoint{2.502557in}{1.532029in}}%
\pgfpathlineto{\pgfqpoint{2.503082in}{1.532059in}}%
\pgfpathlineto{\pgfqpoint{2.504486in}{1.532118in}}%
\pgfpathlineto{\pgfqpoint{2.504496in}{1.532182in}}%
\pgfpathlineto{\pgfqpoint{2.505995in}{2.203122in}}%
\pgfpathlineto{\pgfqpoint{2.506737in}{2.200973in}}%
\pgfpathlineto{\pgfqpoint{2.506786in}{2.202551in}}%
\pgfpathlineto{\pgfqpoint{2.507495in}{1.532035in}}%
\pgfpathlineto{\pgfqpoint{2.508110in}{1.532058in}}%
\pgfpathlineto{\pgfqpoint{2.509442in}{1.532193in}}%
\pgfpathlineto{\pgfqpoint{2.510748in}{2.203140in}}%
\pgfpathlineto{\pgfqpoint{2.510847in}{2.203158in}}%
\pgfpathlineto{\pgfqpoint{2.511589in}{2.200290in}}%
\pgfpathlineto{\pgfqpoint{2.511638in}{2.202801in}}%
\pgfpathlineto{\pgfqpoint{2.512465in}{1.532028in}}%
\pgfpathlineto{\pgfqpoint{2.512903in}{1.532057in}}%
\pgfpathlineto{\pgfqpoint{2.514377in}{1.532119in}}%
\pgfpathlineto{\pgfqpoint{2.514389in}{1.532221in}}%
\pgfpathlineto{\pgfqpoint{2.515518in}{2.203123in}}%
\pgfpathlineto{\pgfqpoint{2.515568in}{2.202721in}}%
\pgfpathlineto{\pgfqpoint{2.516013in}{2.203212in}}%
\pgfpathlineto{\pgfqpoint{2.516260in}{2.201988in}}%
\pgfpathlineto{\pgfqpoint{2.516309in}{2.203114in}}%
\pgfpathlineto{\pgfqpoint{2.516359in}{2.201199in}}%
\pgfpathlineto{\pgfqpoint{2.516408in}{2.203050in}}%
\pgfpathlineto{\pgfqpoint{2.517400in}{1.532028in}}%
\pgfpathlineto{\pgfqpoint{2.517833in}{1.532056in}}%
\pgfpathlineto{\pgfqpoint{2.519323in}{1.532119in}}%
\pgfpathlineto{\pgfqpoint{2.519334in}{1.532225in}}%
\pgfpathlineto{\pgfqpoint{2.520463in}{2.203102in}}%
\pgfpathlineto{\pgfqpoint{2.520759in}{2.203186in}}%
\pgfpathlineto{\pgfqpoint{2.521205in}{2.202198in}}%
\pgfpathlineto{\pgfqpoint{2.521254in}{2.203075in}}%
\pgfpathlineto{\pgfqpoint{2.521402in}{2.195779in}}%
\pgfpathlineto{\pgfqpoint{2.521452in}{2.202924in}}%
\pgfpathlineto{\pgfqpoint{2.522345in}{1.532031in}}%
\pgfpathlineto{\pgfqpoint{2.522727in}{1.532044in}}%
\pgfpathlineto{\pgfqpoint{2.524283in}{1.532471in}}%
\pgfpathlineto{\pgfqpoint{2.525338in}{2.202941in}}%
\pgfpathlineto{\pgfqpoint{2.526574in}{2.201849in}}%
\pgfpathlineto{\pgfqpoint{2.526624in}{2.202181in}}%
\pgfpathlineto{\pgfqpoint{2.526673in}{2.200816in}}%
\pgfpathlineto{\pgfqpoint{2.526723in}{2.201816in}}%
\pgfpathlineto{\pgfqpoint{2.527247in}{1.532014in}}%
\pgfpathlineto{\pgfqpoint{2.528045in}{1.532070in}}%
\pgfpathlineto{\pgfqpoint{2.529214in}{1.532120in}}%
\pgfpathlineto{\pgfqpoint{2.529225in}{1.532215in}}%
\pgfpathlineto{\pgfqpoint{2.530355in}{2.203111in}}%
\pgfpathlineto{\pgfqpoint{2.530405in}{2.202655in}}%
\pgfpathlineto{\pgfqpoint{2.530850in}{2.203205in}}%
\pgfpathlineto{\pgfqpoint{2.530998in}{2.202102in}}%
\pgfpathlineto{\pgfqpoint{2.531048in}{2.203156in}}%
\pgfpathlineto{\pgfqpoint{2.531196in}{2.180845in}}%
\pgfpathlineto{\pgfqpoint{2.531245in}{2.203033in}}%
\pgfpathlineto{\pgfqpoint{2.532227in}{1.532029in}}%
\pgfpathlineto{\pgfqpoint{2.532650in}{1.532054in}}%
\pgfpathlineto{\pgfqpoint{2.534159in}{1.532117in}}%
\pgfpathlineto{\pgfqpoint{2.534169in}{1.532196in}}%
\pgfpathlineto{\pgfqpoint{2.535482in}{2.203102in}}%
\pgfpathlineto{\pgfqpoint{2.536224in}{2.201635in}}%
\pgfpathlineto{\pgfqpoint{2.535581in}{2.203118in}}%
\pgfpathlineto{\pgfqpoint{2.536273in}{2.202853in}}%
\pgfpathlineto{\pgfqpoint{2.537131in}{1.549345in}}%
\pgfpathlineto{\pgfqpoint{2.537189in}{1.532030in}}%
\pgfpathlineto{\pgfqpoint{2.537874in}{1.532056in}}%
\pgfpathlineto{\pgfqpoint{2.539113in}{1.532165in}}%
\pgfpathlineto{\pgfqpoint{2.539155in}{1.869852in}}%
\pgfpathlineto{\pgfqpoint{2.539916in}{2.202796in}}%
\pgfpathlineto{\pgfqpoint{2.540312in}{2.203035in}}%
\pgfpathlineto{\pgfqpoint{2.540658in}{2.202802in}}%
\pgfpathlineto{\pgfqpoint{2.540707in}{2.203079in}}%
\pgfpathlineto{\pgfqpoint{2.541251in}{2.201893in}}%
\pgfpathlineto{\pgfqpoint{2.541301in}{2.202701in}}%
\pgfpathlineto{\pgfqpoint{2.541350in}{2.200945in}}%
\pgfpathlineto{\pgfqpoint{2.541400in}{2.202560in}}%
\pgfpathlineto{\pgfqpoint{2.542126in}{1.532032in}}%
\pgfpathlineto{\pgfqpoint{2.542753in}{1.532059in}}%
\pgfpathlineto{\pgfqpoint{2.544060in}{1.532197in}}%
\pgfpathlineto{\pgfqpoint{2.545274in}{2.203068in}}%
\pgfpathlineto{\pgfqpoint{2.545323in}{2.202747in}}%
\pgfpathlineto{\pgfqpoint{2.545669in}{2.203109in}}%
\pgfpathlineto{\pgfqpoint{2.545917in}{2.202311in}}%
\pgfpathlineto{\pgfqpoint{2.545966in}{2.203001in}}%
\pgfpathlineto{\pgfqpoint{2.546114in}{2.201542in}}%
\pgfpathlineto{\pgfqpoint{2.546164in}{2.202850in}}%
\pgfpathlineto{\pgfqpoint{2.547083in}{1.532030in}}%
\pgfpathlineto{\pgfqpoint{2.548117in}{1.532077in}}%
\pgfpathlineto{\pgfqpoint{2.549004in}{1.532171in}}%
\pgfpathlineto{\pgfqpoint{2.549046in}{1.893926in}}%
\pgfpathlineto{\pgfqpoint{2.549808in}{2.202789in}}%
\pgfpathlineto{\pgfqpoint{2.550203in}{2.203026in}}%
\pgfpathlineto{\pgfqpoint{2.550549in}{2.202791in}}%
\pgfpathlineto{\pgfqpoint{2.550599in}{2.203071in}}%
\pgfpathlineto{\pgfqpoint{2.551143in}{2.201871in}}%
\pgfpathlineto{\pgfqpoint{2.551192in}{2.202691in}}%
\pgfpathlineto{\pgfqpoint{2.551242in}{2.200860in}}%
\pgfpathlineto{\pgfqpoint{2.551291in}{2.202551in}}%
\pgfpathlineto{\pgfqpoint{2.552048in}{1.532035in}}%
\pgfpathlineto{\pgfqpoint{2.552640in}{1.532056in}}%
\pgfpathlineto{\pgfqpoint{2.553956in}{1.532483in}}%
\pgfpathlineto{\pgfqpoint{2.555015in}{2.202901in}}%
\pgfpathlineto{\pgfqpoint{2.556252in}{2.201670in}}%
\pgfpathlineto{\pgfqpoint{2.556301in}{2.202129in}}%
\pgfpathlineto{\pgfqpoint{2.556499in}{2.098815in}}%
\pgfpathlineto{\pgfqpoint{2.556920in}{1.532023in}}%
\pgfpathlineto{\pgfqpoint{2.557567in}{1.532051in}}%
\pgfpathlineto{\pgfqpoint{2.558896in}{1.532182in}}%
\pgfpathlineto{\pgfqpoint{2.560351in}{2.202953in}}%
\pgfpathlineto{\pgfqpoint{2.561340in}{2.201519in}}%
\pgfpathlineto{\pgfqpoint{2.561636in}{1.730374in}}%
\pgfpathlineto{\pgfqpoint{2.561900in}{1.532034in}}%
\pgfpathlineto{\pgfqpoint{2.561846in}{1.746236in}}%
\pgfpathlineto{\pgfqpoint{2.562569in}{1.532057in}}%
\pgfpathlineto{\pgfqpoint{2.563842in}{1.532193in}}%
\pgfpathlineto{\pgfqpoint{2.565148in}{2.203068in}}%
\pgfpathlineto{\pgfqpoint{2.565247in}{2.203087in}}%
\pgfpathlineto{\pgfqpoint{2.565989in}{2.200448in}}%
\pgfpathlineto{\pgfqpoint{2.566038in}{2.202734in}}%
\pgfpathlineto{\pgfqpoint{2.566856in}{1.532028in}}%
\pgfpathlineto{\pgfqpoint{2.567289in}{1.532056in}}%
\pgfpathlineto{\pgfqpoint{2.568777in}{1.532119in}}%
\pgfpathlineto{\pgfqpoint{2.568789in}{1.532228in}}%
\pgfpathlineto{\pgfqpoint{2.569922in}{2.203029in}}%
\pgfpathlineto{\pgfqpoint{2.570120in}{2.203096in}}%
\pgfpathlineto{\pgfqpoint{2.570664in}{2.202215in}}%
\pgfpathlineto{\pgfqpoint{2.570713in}{2.202993in}}%
\pgfpathlineto{\pgfqpoint{2.570862in}{2.201029in}}%
\pgfpathlineto{\pgfqpoint{2.570911in}{2.202838in}}%
\pgfpathlineto{\pgfqpoint{2.571816in}{1.532026in}}%
\pgfpathlineto{\pgfqpoint{2.572304in}{1.532044in}}%
\pgfpathlineto{\pgfqpoint{2.573734in}{1.532224in}}%
\pgfpathlineto{\pgfqpoint{2.574863in}{2.203060in}}%
\pgfpathlineto{\pgfqpoint{2.574913in}{2.202641in}}%
\pgfpathlineto{\pgfqpoint{2.575358in}{2.203153in}}%
\pgfpathlineto{\pgfqpoint{2.575506in}{2.202176in}}%
\pgfpathlineto{\pgfqpoint{2.575556in}{2.203102in}}%
\pgfpathlineto{\pgfqpoint{2.575704in}{2.200675in}}%
\pgfpathlineto{\pgfqpoint{2.575754in}{2.202995in}}%
\pgfpathlineto{\pgfqpoint{2.576734in}{1.532029in}}%
\pgfpathlineto{\pgfqpoint{2.577056in}{1.532050in}}%
\pgfpathlineto{\pgfqpoint{2.578667in}{1.532116in}}%
\pgfpathlineto{\pgfqpoint{2.578678in}{1.532189in}}%
\pgfpathlineto{\pgfqpoint{2.579986in}{2.203058in}}%
\pgfpathlineto{\pgfqpoint{2.580728in}{2.201659in}}%
\pgfpathlineto{\pgfqpoint{2.580085in}{2.203075in}}%
\pgfpathlineto{\pgfqpoint{2.580777in}{2.202811in}}%
\pgfpathlineto{\pgfqpoint{2.581345in}{1.904734in}}%
\pgfpathlineto{\pgfqpoint{2.581695in}{1.532030in}}%
\pgfpathlineto{\pgfqpoint{2.582231in}{1.532058in}}%
\pgfpathlineto{\pgfqpoint{2.583622in}{1.532166in}}%
\pgfpathlineto{\pgfqpoint{2.583664in}{1.869508in}}%
\pgfpathlineto{\pgfqpoint{2.584425in}{2.202761in}}%
\pgfpathlineto{\pgfqpoint{2.584821in}{2.203001in}}%
\pgfpathlineto{\pgfqpoint{2.585167in}{2.202747in}}%
\pgfpathlineto{\pgfqpoint{2.585217in}{2.203048in}}%
\pgfpathlineto{\pgfqpoint{2.585761in}{2.201757in}}%
\pgfpathlineto{\pgfqpoint{2.585810in}{2.202678in}}%
\pgfpathlineto{\pgfqpoint{2.585860in}{2.199765in}}%
\pgfpathlineto{\pgfqpoint{2.585909in}{2.202538in}}%
\pgfpathlineto{\pgfqpoint{2.586592in}{1.532009in}}%
\pgfpathlineto{\pgfqpoint{2.587141in}{1.532062in}}%
\pgfpathlineto{\pgfqpoint{2.588559in}{1.532119in}}%
\pgfpathlineto{\pgfqpoint{2.588568in}{1.532179in}}%
\pgfpathlineto{\pgfqpoint{2.588609in}{1.878962in}}%
\pgfpathlineto{\pgfqpoint{2.589375in}{2.202787in}}%
\pgfpathlineto{\pgfqpoint{2.589870in}{2.203051in}}%
\pgfpathlineto{\pgfqpoint{2.590117in}{2.202684in}}%
\pgfpathlineto{\pgfqpoint{2.590167in}{2.203065in}}%
\pgfpathlineto{\pgfqpoint{2.590612in}{2.201867in}}%
\pgfpathlineto{\pgfqpoint{2.590661in}{2.202805in}}%
\pgfpathlineto{\pgfqpoint{2.590711in}{2.200951in}}%
\pgfpathlineto{\pgfqpoint{2.590760in}{2.202696in}}%
\pgfpathlineto{\pgfqpoint{2.591579in}{1.532031in}}%
\pgfpathlineto{\pgfqpoint{2.592107in}{1.532057in}}%
\pgfpathlineto{\pgfqpoint{2.593511in}{1.532148in}}%
\pgfpathlineto{\pgfqpoint{2.593520in}{1.532493in}}%
\pgfpathlineto{\pgfqpoint{2.594579in}{2.202870in}}%
\pgfpathlineto{\pgfqpoint{2.595815in}{2.201733in}}%
\pgfpathlineto{\pgfqpoint{2.595865in}{2.202129in}}%
\pgfpathlineto{\pgfqpoint{2.595914in}{2.200401in}}%
\pgfpathlineto{\pgfqpoint{2.595964in}{2.201777in}}%
\pgfpathlineto{\pgfqpoint{2.596483in}{1.532012in}}%
\pgfpathlineto{\pgfqpoint{2.597245in}{1.532052in}}%
\pgfpathlineto{\pgfqpoint{2.598462in}{1.532224in}}%
\pgfpathlineto{\pgfqpoint{2.599591in}{2.203033in}}%
\pgfpathlineto{\pgfqpoint{2.599640in}{2.202590in}}%
\pgfpathlineto{\pgfqpoint{2.600085in}{2.203124in}}%
\pgfpathlineto{\pgfqpoint{2.600234in}{2.202065in}}%
\pgfpathlineto{\pgfqpoint{2.600283in}{2.203073in}}%
\pgfpathlineto{\pgfqpoint{2.600432in}{2.189543in}}%
\pgfpathlineto{\pgfqpoint{2.600481in}{2.202960in}}%
\pgfpathlineto{\pgfqpoint{2.601481in}{1.532029in}}%
\pgfpathlineto{\pgfqpoint{2.601820in}{1.532052in}}%
\pgfpathlineto{\pgfqpoint{2.603395in}{1.532118in}}%
\pgfpathlineto{\pgfqpoint{2.603406in}{1.532206in}}%
\pgfpathlineto{\pgfqpoint{2.604538in}{2.203017in}}%
\pgfpathlineto{\pgfqpoint{2.604588in}{2.202529in}}%
\pgfpathlineto{\pgfqpoint{2.605033in}{2.203110in}}%
\pgfpathlineto{\pgfqpoint{2.605181in}{2.201872in}}%
\pgfpathlineto{\pgfqpoint{2.605231in}{2.203061in}}%
\pgfpathlineto{\pgfqpoint{2.605280in}{2.201164in}}%
\pgfpathlineto{\pgfqpoint{2.605330in}{2.203015in}}%
\pgfpathlineto{\pgfqpoint{2.606429in}{1.532028in}}%
\pgfpathlineto{\pgfqpoint{2.606817in}{1.532042in}}%
\pgfpathlineto{\pgfqpoint{2.608352in}{1.532223in}}%
\pgfpathlineto{\pgfqpoint{2.609481in}{2.202980in}}%
\pgfpathlineto{\pgfqpoint{2.609531in}{2.202609in}}%
\pgfpathlineto{\pgfqpoint{2.609976in}{2.203064in}}%
\pgfpathlineto{\pgfqpoint{2.610223in}{2.201977in}}%
\pgfpathlineto{\pgfqpoint{2.610273in}{2.202961in}}%
\pgfpathlineto{\pgfqpoint{2.610421in}{2.175288in}}%
\pgfpathlineto{\pgfqpoint{2.610470in}{2.202781in}}%
\pgfpathlineto{\pgfqpoint{2.611378in}{1.532027in}}%
\pgfpathlineto{\pgfqpoint{2.611965in}{1.532048in}}%
\pgfpathlineto{\pgfqpoint{2.613298in}{1.532217in}}%
\pgfpathlineto{\pgfqpoint{2.614428in}{2.202996in}}%
\pgfpathlineto{\pgfqpoint{2.614477in}{2.202552in}}%
\pgfpathlineto{\pgfqpoint{2.614922in}{2.203089in}}%
\pgfpathlineto{\pgfqpoint{2.615071in}{2.202032in}}%
\pgfpathlineto{\pgfqpoint{2.615120in}{2.203040in}}%
\pgfpathlineto{\pgfqpoint{2.615268in}{2.190251in}}%
\pgfpathlineto{\pgfqpoint{2.615318in}{2.202929in}}%
\pgfpathlineto{\pgfqpoint{2.616301in}{1.532029in}}%
\pgfpathlineto{\pgfqpoint{2.616627in}{1.532051in}}%
\pgfpathlineto{\pgfqpoint{2.618231in}{1.532117in}}%
\pgfpathlineto{\pgfqpoint{2.618241in}{1.532182in}}%
\pgfpathlineto{\pgfqpoint{2.619734in}{2.202956in}}%
\pgfpathlineto{\pgfqpoint{2.620476in}{2.201356in}}%
\pgfpathlineto{\pgfqpoint{2.620525in}{2.202373in}}%
\pgfpathlineto{\pgfqpoint{2.621293in}{1.532035in}}%
\pgfpathlineto{\pgfqpoint{2.622327in}{1.532075in}}%
\pgfpathlineto{\pgfqpoint{2.623189in}{1.532228in}}%
\pgfpathlineto{\pgfqpoint{2.624391in}{2.202954in}}%
\pgfpathlineto{\pgfqpoint{2.624589in}{2.203004in}}%
\pgfpathlineto{\pgfqpoint{2.625232in}{2.201838in}}%
\pgfpathlineto{\pgfqpoint{2.625282in}{2.202755in}}%
\pgfpathlineto{\pgfqpoint{2.625331in}{2.201051in}}%
\pgfpathlineto{\pgfqpoint{2.625381in}{2.202650in}}%
\pgfpathlineto{\pgfqpoint{2.626212in}{1.532027in}}%
\pgfpathlineto{\pgfqpoint{2.626799in}{1.532049in}}%
\pgfpathlineto{\pgfqpoint{2.628135in}{1.532245in}}%
\pgfpathlineto{\pgfqpoint{2.629253in}{2.202880in}}%
\pgfpathlineto{\pgfqpoint{2.629352in}{2.202917in}}%
\pgfpathlineto{\pgfqpoint{2.630094in}{2.202375in}}%
\pgfpathlineto{\pgfqpoint{2.630143in}{2.202720in}}%
\pgfpathlineto{\pgfqpoint{2.630391in}{2.201284in}}%
\pgfpathlineto{\pgfqpoint{2.630440in}{2.202321in}}%
\pgfpathlineto{\pgfqpoint{2.631130in}{1.532031in}}%
\pgfpathlineto{\pgfqpoint{2.632137in}{1.532074in}}%
\pgfpathlineto{\pgfqpoint{2.633074in}{1.532148in}}%
\pgfpathlineto{\pgfqpoint{2.633083in}{1.532492in}}%
\pgfpathlineto{\pgfqpoint{2.634142in}{2.202812in}}%
\pgfpathlineto{\pgfqpoint{2.635477in}{2.200478in}}%
\pgfpathlineto{\pgfqpoint{2.635527in}{2.201730in}}%
\pgfpathlineto{\pgfqpoint{2.636047in}{1.532018in}}%
\pgfpathlineto{\pgfqpoint{2.636813in}{1.532067in}}%
\pgfpathlineto{\pgfqpoint{2.638014in}{1.532119in}}%
\pgfpathlineto{\pgfqpoint{2.638023in}{1.532181in}}%
\pgfpathlineto{\pgfqpoint{2.639522in}{2.202994in}}%
\pgfpathlineto{\pgfqpoint{2.639571in}{2.202646in}}%
\pgfpathlineto{\pgfqpoint{2.639621in}{2.202985in}}%
\pgfpathlineto{\pgfqpoint{2.640165in}{2.201386in}}%
\pgfpathlineto{\pgfqpoint{2.640214in}{2.202611in}}%
\pgfpathlineto{\pgfqpoint{2.641053in}{1.532033in}}%
\pgfpathlineto{\pgfqpoint{2.642042in}{1.532069in}}%
\pgfpathlineto{\pgfqpoint{2.642969in}{1.532181in}}%
\pgfpathlineto{\pgfqpoint{2.644468in}{2.203012in}}%
\pgfpathlineto{\pgfqpoint{2.644517in}{2.202634in}}%
\pgfpathlineto{\pgfqpoint{2.644567in}{2.203005in}}%
\pgfpathlineto{\pgfqpoint{2.645012in}{2.201873in}}%
\pgfpathlineto{\pgfqpoint{2.645061in}{2.202752in}}%
\pgfpathlineto{\pgfqpoint{2.645111in}{2.201173in}}%
\pgfpathlineto{\pgfqpoint{2.645160in}{2.202646in}}%
\pgfpathlineto{\pgfqpoint{2.645969in}{1.532030in}}%
\pgfpathlineto{\pgfqpoint{2.646685in}{1.532064in}}%
\pgfpathlineto{\pgfqpoint{2.647914in}{1.532182in}}%
\pgfpathlineto{\pgfqpoint{2.649217in}{2.202987in}}%
\pgfpathlineto{\pgfqpoint{2.649266in}{2.202661in}}%
\pgfpathlineto{\pgfqpoint{2.649513in}{2.203003in}}%
\pgfpathlineto{\pgfqpoint{2.649859in}{2.202076in}}%
\pgfpathlineto{\pgfqpoint{2.649909in}{2.202834in}}%
\pgfpathlineto{\pgfqpoint{2.650057in}{2.199984in}}%
\pgfpathlineto{\pgfqpoint{2.650107in}{2.202644in}}%
\pgfpathlineto{\pgfqpoint{2.650900in}{1.532025in}}%
\pgfpathlineto{\pgfqpoint{2.651385in}{1.532060in}}%
\pgfpathlineto{\pgfqpoint{2.652851in}{1.532121in}}%
\pgfpathlineto{\pgfqpoint{2.652862in}{1.532225in}}%
\pgfpathlineto{\pgfqpoint{2.653992in}{2.202929in}}%
\pgfpathlineto{\pgfqpoint{2.654190in}{2.202995in}}%
\pgfpathlineto{\pgfqpoint{2.654833in}{2.202087in}}%
\pgfpathlineto{\pgfqpoint{2.654882in}{2.202816in}}%
\pgfpathlineto{\pgfqpoint{2.655030in}{2.200073in}}%
\pgfpathlineto{\pgfqpoint{2.655080in}{2.202615in}}%
\pgfpathlineto{\pgfqpoint{2.655914in}{1.532032in}}%
\pgfpathlineto{\pgfqpoint{2.656354in}{1.532054in}}%
\pgfpathlineto{\pgfqpoint{2.657802in}{1.532147in}}%
\pgfpathlineto{\pgfqpoint{2.657811in}{1.532474in}}%
\pgfpathlineto{\pgfqpoint{2.658862in}{2.202791in}}%
\pgfpathlineto{\pgfqpoint{2.660099in}{2.201723in}}%
\pgfpathlineto{\pgfqpoint{2.660148in}{2.201998in}}%
\pgfpathlineto{\pgfqpoint{2.660198in}{2.200800in}}%
\pgfpathlineto{\pgfqpoint{2.660247in}{2.201606in}}%
\pgfpathlineto{\pgfqpoint{2.660775in}{1.532021in}}%
\pgfpathlineto{\pgfqpoint{2.661614in}{1.532070in}}%
\pgfpathlineto{\pgfqpoint{2.662741in}{1.532118in}}%
\pgfpathlineto{\pgfqpoint{2.662750in}{1.532182in}}%
\pgfpathlineto{\pgfqpoint{2.664244in}{2.202946in}}%
\pgfpathlineto{\pgfqpoint{2.664986in}{2.201340in}}%
\pgfpathlineto{\pgfqpoint{2.665035in}{2.202387in}}%
\pgfpathlineto{\pgfqpoint{2.665762in}{1.532028in}}%
\pgfpathlineto{\pgfqpoint{2.666837in}{1.532068in}}%
\pgfpathlineto{\pgfqpoint{2.667698in}{1.532225in}}%
\pgfpathlineto{\pgfqpoint{2.668826in}{2.202950in}}%
\pgfpathlineto{\pgfqpoint{2.668876in}{2.202604in}}%
\pgfpathlineto{\pgfqpoint{2.669321in}{2.203035in}}%
\pgfpathlineto{\pgfqpoint{2.669568in}{2.202045in}}%
\pgfpathlineto{\pgfqpoint{2.669618in}{2.202931in}}%
\pgfpathlineto{\pgfqpoint{2.669766in}{2.195258in}}%
\pgfpathlineto{\pgfqpoint{2.669816in}{2.202783in}}%
\pgfpathlineto{\pgfqpoint{2.670682in}{1.532023in}}%
\pgfpathlineto{\pgfqpoint{2.671071in}{1.532057in}}%
\pgfpathlineto{\pgfqpoint{2.672632in}{1.532118in}}%
\pgfpathlineto{\pgfqpoint{2.672643in}{1.532219in}}%
\pgfpathlineto{\pgfqpoint{2.673758in}{2.202869in}}%
\pgfpathlineto{\pgfqpoint{2.673857in}{2.202906in}}%
\pgfpathlineto{\pgfqpoint{2.674599in}{2.202372in}}%
\pgfpathlineto{\pgfqpoint{2.674648in}{2.202715in}}%
\pgfpathlineto{\pgfqpoint{2.674895in}{2.201318in}}%
\pgfpathlineto{\pgfqpoint{2.674945in}{2.202320in}}%
\pgfpathlineto{\pgfqpoint{2.675642in}{1.532026in}}%
\pgfpathlineto{\pgfqpoint{2.676703in}{1.532065in}}%
\pgfpathlineto{\pgfqpoint{2.677587in}{1.532183in}}%
\pgfpathlineto{\pgfqpoint{2.677627in}{1.858420in}}%
\pgfpathlineto{\pgfqpoint{2.678396in}{2.202731in}}%
\pgfpathlineto{\pgfqpoint{2.678891in}{2.203006in}}%
\pgfpathlineto{\pgfqpoint{2.679138in}{2.202600in}}%
\pgfpathlineto{\pgfqpoint{2.679188in}{2.203028in}}%
\pgfpathlineto{\pgfqpoint{2.679633in}{2.201681in}}%
\pgfpathlineto{\pgfqpoint{2.679682in}{2.202790in}}%
\pgfpathlineto{\pgfqpoint{2.680177in}{2.133080in}}%
\pgfpathlineto{\pgfqpoint{2.680593in}{1.532029in}}%
\pgfpathlineto{\pgfqpoint{2.681217in}{1.532062in}}%
\pgfpathlineto{\pgfqpoint{2.682522in}{1.532118in}}%
\pgfpathlineto{\pgfqpoint{2.682532in}{1.532181in}}%
\pgfpathlineto{\pgfqpoint{2.684032in}{2.202977in}}%
\pgfpathlineto{\pgfqpoint{2.684773in}{2.200765in}}%
\pgfpathlineto{\pgfqpoint{2.684823in}{2.202407in}}%
\pgfpathlineto{\pgfqpoint{2.685521in}{1.532024in}}%
\pgfpathlineto{\pgfqpoint{2.686163in}{1.532046in}}%
\pgfpathlineto{\pgfqpoint{2.687479in}{1.532203in}}%
\pgfpathlineto{\pgfqpoint{2.688709in}{2.203015in}}%
\pgfpathlineto{\pgfqpoint{2.688907in}{2.203064in}}%
\pgfpathlineto{\pgfqpoint{2.689451in}{2.200883in}}%
\pgfpathlineto{\pgfqpoint{2.689501in}{2.202908in}}%
\pgfpathlineto{\pgfqpoint{2.690446in}{1.532007in}}%
\pgfpathlineto{\pgfqpoint{2.690848in}{1.532059in}}%
\pgfpathlineto{\pgfqpoint{2.692412in}{1.532114in}}%
\pgfpathlineto{\pgfqpoint{2.692422in}{1.532165in}}%
\pgfpathlineto{\pgfqpoint{2.692462in}{1.835859in}}%
\pgfpathlineto{\pgfqpoint{2.693229in}{2.202682in}}%
\pgfpathlineto{\pgfqpoint{2.693525in}{2.202870in}}%
\pgfpathlineto{\pgfqpoint{2.693970in}{2.202715in}}%
\pgfpathlineto{\pgfqpoint{2.694020in}{2.202940in}}%
\pgfpathlineto{\pgfqpoint{2.694564in}{2.201890in}}%
\pgfpathlineto{\pgfqpoint{2.694613in}{2.202519in}}%
\pgfpathlineto{\pgfqpoint{2.694663in}{2.201284in}}%
\pgfpathlineto{\pgfqpoint{2.694712in}{2.202364in}}%
\pgfpathlineto{\pgfqpoint{2.695457in}{1.532026in}}%
\pgfpathlineto{\pgfqpoint{2.696389in}{1.532078in}}%
\pgfpathlineto{\pgfqpoint{2.697360in}{1.532122in}}%
\pgfpathlineto{\pgfqpoint{2.697371in}{1.532224in}}%
\pgfpathlineto{\pgfqpoint{2.698499in}{2.202959in}}%
\pgfpathlineto{\pgfqpoint{2.698796in}{2.203047in}}%
\pgfpathlineto{\pgfqpoint{2.699241in}{2.202134in}}%
\pgfpathlineto{\pgfqpoint{2.699290in}{2.202944in}}%
\pgfpathlineto{\pgfqpoint{2.699438in}{2.200859in}}%
\pgfpathlineto{\pgfqpoint{2.699488in}{2.202797in}}%
\pgfpathlineto{\pgfqpoint{2.700399in}{1.532033in}}%
\pgfpathlineto{\pgfqpoint{2.700836in}{1.532052in}}%
\pgfpathlineto{\pgfqpoint{2.702314in}{1.532181in}}%
\pgfpathlineto{\pgfqpoint{2.703816in}{2.202986in}}%
\pgfpathlineto{\pgfqpoint{2.704558in}{2.195221in}}%
\pgfpathlineto{\pgfqpoint{2.704607in}{2.202413in}}%
\pgfpathlineto{\pgfqpoint{2.705317in}{1.532034in}}%
\pgfpathlineto{\pgfqpoint{2.705885in}{1.532053in}}%
\pgfpathlineto{\pgfqpoint{2.707260in}{1.532189in}}%
\pgfpathlineto{\pgfqpoint{2.708568in}{2.202997in}}%
\pgfpathlineto{\pgfqpoint{2.708666in}{2.203015in}}%
\pgfpathlineto{\pgfqpoint{2.709408in}{2.196770in}}%
\pgfpathlineto{\pgfqpoint{2.709458in}{2.202657in}}%
\pgfpathlineto{\pgfqpoint{2.709507in}{2.137790in}}%
\pgfpathlineto{\pgfqpoint{2.709656in}{2.202119in}}%
\pgfpathlineto{\pgfqpoint{2.710228in}{1.532006in}}%
\pgfpathlineto{\pgfqpoint{2.710917in}{1.532071in}}%
\pgfpathlineto{\pgfqpoint{2.712194in}{1.532114in}}%
\pgfpathlineto{\pgfqpoint{2.712203in}{1.532162in}}%
\pgfpathlineto{\pgfqpoint{2.712232in}{1.612001in}}%
\pgfpathlineto{\pgfqpoint{2.713006in}{2.202699in}}%
\pgfpathlineto{\pgfqpoint{2.713501in}{2.202971in}}%
\pgfpathlineto{\pgfqpoint{2.713748in}{2.202609in}}%
\pgfpathlineto{\pgfqpoint{2.713797in}{2.202990in}}%
\pgfpathlineto{\pgfqpoint{2.714242in}{2.201829in}}%
\pgfpathlineto{\pgfqpoint{2.714292in}{2.202742in}}%
\pgfpathlineto{\pgfqpoint{2.714341in}{2.201057in}}%
\pgfpathlineto{\pgfqpoint{2.714391in}{2.202637in}}%
\pgfpathlineto{\pgfqpoint{2.715192in}{1.532023in}}%
\pgfpathlineto{\pgfqpoint{2.715828in}{1.532046in}}%
\pgfpathlineto{\pgfqpoint{2.717152in}{1.532209in}}%
\pgfpathlineto{\pgfqpoint{2.718284in}{2.202958in}}%
\pgfpathlineto{\pgfqpoint{2.718334in}{2.202506in}}%
\pgfpathlineto{\pgfqpoint{2.718779in}{2.203050in}}%
\pgfpathlineto{\pgfqpoint{2.718927in}{2.201957in}}%
\pgfpathlineto{\pgfqpoint{2.718977in}{2.203001in}}%
\pgfpathlineto{\pgfqpoint{2.719125in}{2.182395in}}%
\pgfpathlineto{\pgfqpoint{2.719175in}{2.202880in}}%
\pgfpathlineto{\pgfqpoint{2.720151in}{1.532034in}}%
\pgfpathlineto{\pgfqpoint{2.720547in}{1.532049in}}%
\pgfpathlineto{\pgfqpoint{2.722101in}{1.532413in}}%
\pgfpathlineto{\pgfqpoint{2.723096in}{2.202781in}}%
\pgfpathlineto{\pgfqpoint{2.723145in}{2.202617in}}%
\pgfpathlineto{\pgfqpoint{2.723689in}{2.202912in}}%
\pgfpathlineto{\pgfqpoint{2.723887in}{2.202850in}}%
\pgfpathlineto{\pgfqpoint{2.724332in}{2.199606in}}%
\pgfpathlineto{\pgfqpoint{2.724381in}{2.202372in}}%
\pgfpathlineto{\pgfqpoint{2.725065in}{1.532012in}}%
\pgfpathlineto{\pgfqpoint{2.725631in}{1.532045in}}%
\pgfpathlineto{\pgfqpoint{2.727043in}{1.532224in}}%
\pgfpathlineto{\pgfqpoint{2.728172in}{2.202921in}}%
\pgfpathlineto{\pgfqpoint{2.728469in}{2.203013in}}%
\pgfpathlineto{\pgfqpoint{2.728914in}{2.201724in}}%
\pgfpathlineto{\pgfqpoint{2.728964in}{2.202920in}}%
\pgfpathlineto{\pgfqpoint{2.729013in}{2.200665in}}%
\pgfpathlineto{\pgfqpoint{2.729063in}{2.202858in}}%
\pgfpathlineto{\pgfqpoint{2.730043in}{1.532032in}}%
\pgfpathlineto{\pgfqpoint{2.730365in}{1.532048in}}%
\pgfpathlineto{\pgfqpoint{2.731985in}{1.532157in}}%
\pgfpathlineto{\pgfqpoint{2.732011in}{1.599492in}}%
\pgfpathlineto{\pgfqpoint{2.732740in}{2.202541in}}%
\pgfpathlineto{\pgfqpoint{2.732789in}{2.202510in}}%
\pgfpathlineto{\pgfqpoint{2.733729in}{2.202771in}}%
\pgfpathlineto{\pgfqpoint{2.734273in}{2.201522in}}%
\pgfpathlineto{\pgfqpoint{2.734322in}{2.202050in}}%
\pgfpathlineto{\pgfqpoint{2.734520in}{2.121066in}}%
\pgfpathlineto{\pgfqpoint{2.734956in}{1.532019in}}%
\pgfpathlineto{\pgfqpoint{2.735714in}{1.532067in}}%
\pgfpathlineto{\pgfqpoint{2.736923in}{1.532118in}}%
\pgfpathlineto{\pgfqpoint{2.736932in}{1.532181in}}%
\pgfpathlineto{\pgfqpoint{2.738436in}{2.202911in}}%
\pgfpathlineto{\pgfqpoint{2.738486in}{2.202586in}}%
\pgfpathlineto{\pgfqpoint{2.738535in}{2.202902in}}%
\pgfpathlineto{\pgfqpoint{2.739079in}{2.201455in}}%
\pgfpathlineto{\pgfqpoint{2.739129in}{2.202517in}}%
\pgfpathlineto{\pgfqpoint{2.739673in}{1.756638in}}%
\pgfpathlineto{\pgfqpoint{2.739993in}{1.532034in}}%
\pgfpathlineto{\pgfqpoint{2.740584in}{1.532056in}}%
\pgfpathlineto{\pgfqpoint{2.741878in}{1.532183in}}%
\pgfpathlineto{\pgfqpoint{2.743196in}{2.202923in}}%
\pgfpathlineto{\pgfqpoint{2.743938in}{2.201233in}}%
\pgfpathlineto{\pgfqpoint{2.743295in}{2.202941in}}%
\pgfpathlineto{\pgfqpoint{2.743987in}{2.202684in}}%
\pgfpathlineto{\pgfqpoint{2.744863in}{1.532025in}}%
\pgfpathlineto{\pgfqpoint{2.745645in}{1.532071in}}%
\pgfpathlineto{\pgfqpoint{2.746814in}{1.532121in}}%
\pgfpathlineto{\pgfqpoint{2.746825in}{1.532224in}}%
\pgfpathlineto{\pgfqpoint{2.747957in}{2.202858in}}%
\pgfpathlineto{\pgfqpoint{2.748155in}{2.202924in}}%
\pgfpathlineto{\pgfqpoint{2.748699in}{2.202149in}}%
\pgfpathlineto{\pgfqpoint{2.748748in}{2.202818in}}%
\pgfpathlineto{\pgfqpoint{2.748896in}{2.201401in}}%
\pgfpathlineto{\pgfqpoint{2.748946in}{2.202662in}}%
\pgfpathlineto{\pgfqpoint{2.749867in}{1.532033in}}%
\pgfpathlineto{\pgfqpoint{2.751048in}{1.532077in}}%
\pgfpathlineto{\pgfqpoint{2.751768in}{1.532168in}}%
\pgfpathlineto{\pgfqpoint{2.751810in}{1.873718in}}%
\pgfpathlineto{\pgfqpoint{2.752571in}{2.202623in}}%
\pgfpathlineto{\pgfqpoint{2.752967in}{2.202859in}}%
\pgfpathlineto{\pgfqpoint{2.753313in}{2.202596in}}%
\pgfpathlineto{\pgfqpoint{2.753362in}{2.202902in}}%
\pgfpathlineto{\pgfqpoint{2.753906in}{2.201554in}}%
\pgfpathlineto{\pgfqpoint{2.753956in}{2.202525in}}%
\pgfpathlineto{\pgfqpoint{2.754352in}{2.081528in}}%
\pgfpathlineto{\pgfqpoint{2.754769in}{1.532034in}}%
\pgfpathlineto{\pgfqpoint{2.755384in}{1.532057in}}%
\pgfpathlineto{\pgfqpoint{2.756715in}{1.532198in}}%
\pgfpathlineto{\pgfqpoint{2.757921in}{2.202873in}}%
\pgfpathlineto{\pgfqpoint{2.757970in}{2.202586in}}%
\pgfpathlineto{\pgfqpoint{2.758316in}{2.202917in}}%
\pgfpathlineto{\pgfqpoint{2.758564in}{2.202215in}}%
\pgfpathlineto{\pgfqpoint{2.758613in}{2.202811in}}%
\pgfpathlineto{\pgfqpoint{2.758860in}{2.200311in}}%
\pgfpathlineto{\pgfqpoint{2.758910in}{2.202552in}}%
\pgfpathlineto{\pgfqpoint{2.759683in}{1.532018in}}%
\pgfpathlineto{\pgfqpoint{2.760177in}{1.532061in}}%
\pgfpathlineto{\pgfqpoint{2.761650in}{1.532118in}}%
\pgfpathlineto{\pgfqpoint{2.761661in}{1.532203in}}%
\pgfpathlineto{\pgfqpoint{2.762891in}{2.202921in}}%
\pgfpathlineto{\pgfqpoint{2.763089in}{2.202971in}}%
\pgfpathlineto{\pgfqpoint{2.763633in}{2.200650in}}%
\pgfpathlineto{\pgfqpoint{2.763683in}{2.202816in}}%
\pgfpathlineto{\pgfqpoint{2.764627in}{1.531998in}}%
\pgfpathlineto{\pgfqpoint{2.765023in}{1.532060in}}%
\pgfpathlineto{\pgfqpoint{2.766593in}{1.532113in}}%
\pgfpathlineto{\pgfqpoint{2.766612in}{1.533241in}}%
\pgfpathlineto{\pgfqpoint{2.767572in}{2.202646in}}%
\pgfpathlineto{\pgfqpoint{2.767622in}{2.202611in}}%
\pgfpathlineto{\pgfqpoint{2.768364in}{2.202703in}}%
\pgfpathlineto{\pgfqpoint{2.768462in}{2.202642in}}%
\pgfpathlineto{\pgfqpoint{2.769006in}{2.199543in}}%
\pgfpathlineto{\pgfqpoint{2.769056in}{2.201510in}}%
\pgfpathlineto{\pgfqpoint{2.769625in}{1.532027in}}%
\pgfpathlineto{\pgfqpoint{2.770307in}{1.532052in}}%
\pgfpathlineto{\pgfqpoint{2.771553in}{1.532251in}}%
\pgfpathlineto{\pgfqpoint{2.772653in}{2.202701in}}%
\pgfpathlineto{\pgfqpoint{2.773889in}{2.201529in}}%
\pgfpathlineto{\pgfqpoint{2.773939in}{2.201738in}}%
\pgfpathlineto{\pgfqpoint{2.773988in}{2.199998in}}%
\pgfpathlineto{\pgfqpoint{2.774037in}{2.201149in}}%
\pgfpathlineto{\pgfqpoint{2.774576in}{1.532031in}}%
\pgfpathlineto{\pgfqpoint{2.775315in}{1.532065in}}%
\pgfpathlineto{\pgfqpoint{2.776493in}{1.532149in}}%
\pgfpathlineto{\pgfqpoint{2.776502in}{1.532526in}}%
\pgfpathlineto{\pgfqpoint{2.777554in}{2.202679in}}%
\pgfpathlineto{\pgfqpoint{2.778791in}{2.201724in}}%
\pgfpathlineto{\pgfqpoint{2.778840in}{2.201899in}}%
\pgfpathlineto{\pgfqpoint{2.778889in}{2.201088in}}%
\pgfpathlineto{\pgfqpoint{2.778939in}{2.201509in}}%
\pgfpathlineto{\pgfqpoint{2.779515in}{1.532026in}}%
\pgfpathlineto{\pgfqpoint{2.780592in}{1.532066in}}%
\pgfpathlineto{\pgfqpoint{2.781443in}{1.532223in}}%
\pgfpathlineto{\pgfqpoint{2.782572in}{2.202850in}}%
\pgfpathlineto{\pgfqpoint{2.782869in}{2.202939in}}%
\pgfpathlineto{\pgfqpoint{2.783314in}{2.201905in}}%
\pgfpathlineto{\pgfqpoint{2.783363in}{2.202836in}}%
\pgfpathlineto{\pgfqpoint{2.783413in}{2.201472in}}%
\pgfpathlineto{\pgfqpoint{2.783462in}{2.202773in}}%
\pgfpathlineto{\pgfqpoint{2.784204in}{1.734406in}}%
\pgfpathlineto{\pgfqpoint{2.784485in}{1.532033in}}%
\pgfpathlineto{\pgfqpoint{2.785122in}{1.532061in}}%
\pgfpathlineto{\pgfqpoint{2.786386in}{1.532168in}}%
\pgfpathlineto{\pgfqpoint{2.786428in}{1.873441in}}%
\pgfpathlineto{\pgfqpoint{2.787189in}{2.202592in}}%
\pgfpathlineto{\pgfqpoint{2.787585in}{2.202827in}}%
\pgfpathlineto{\pgfqpoint{2.787931in}{2.202562in}}%
\pgfpathlineto{\pgfqpoint{2.787981in}{2.202869in}}%
\pgfpathlineto{\pgfqpoint{2.788525in}{2.201508in}}%
\pgfpathlineto{\pgfqpoint{2.788574in}{2.202490in}}%
\pgfpathlineto{\pgfqpoint{2.789019in}{1.979854in}}%
\pgfpathlineto{\pgfqpoint{2.789429in}{1.532034in}}%
\pgfpathlineto{\pgfqpoint{2.789967in}{1.532056in}}%
\pgfpathlineto{\pgfqpoint{2.791333in}{1.532195in}}%
\pgfpathlineto{\pgfqpoint{2.792728in}{2.202851in}}%
\pgfpathlineto{\pgfqpoint{2.793569in}{2.201164in}}%
\pgfpathlineto{\pgfqpoint{2.793618in}{2.202289in}}%
\pgfpathlineto{\pgfqpoint{2.794358in}{1.532028in}}%
\pgfpathlineto{\pgfqpoint{2.795241in}{1.532061in}}%
\pgfpathlineto{\pgfqpoint{2.796280in}{1.532237in}}%
\pgfpathlineto{\pgfqpoint{2.797408in}{2.202813in}}%
\pgfpathlineto{\pgfqpoint{2.797606in}{2.202878in}}%
\pgfpathlineto{\pgfqpoint{2.798249in}{2.202101in}}%
\pgfpathlineto{\pgfqpoint{2.798299in}{2.202686in}}%
\pgfpathlineto{\pgfqpoint{2.798447in}{2.201238in}}%
\pgfpathlineto{\pgfqpoint{2.798496in}{2.202477in}}%
\pgfpathlineto{\pgfqpoint{2.799305in}{1.532027in}}%
\pgfpathlineto{\pgfqpoint{2.800238in}{1.532077in}}%
\pgfpathlineto{\pgfqpoint{2.801214in}{1.532121in}}%
\pgfpathlineto{\pgfqpoint{2.801225in}{1.532216in}}%
\pgfpathlineto{\pgfqpoint{2.802355in}{2.202893in}}%
\pgfpathlineto{\pgfqpoint{2.802404in}{2.202438in}}%
\pgfpathlineto{\pgfqpoint{2.802850in}{2.202989in}}%
\pgfpathlineto{\pgfqpoint{2.802998in}{2.201895in}}%
\pgfpathlineto{\pgfqpoint{2.803047in}{2.202941in}}%
\pgfpathlineto{\pgfqpoint{2.803196in}{2.182768in}}%
\pgfpathlineto{\pgfqpoint{2.803245in}{2.202823in}}%
\pgfpathlineto{\pgfqpoint{2.804227in}{1.532031in}}%
\pgfpathlineto{\pgfqpoint{2.804650in}{1.532052in}}%
\pgfpathlineto{\pgfqpoint{2.806166in}{1.532154in}}%
\pgfpathlineto{\pgfqpoint{2.806184in}{1.554306in}}%
\pgfpathlineto{\pgfqpoint{2.807017in}{2.202576in}}%
\pgfpathlineto{\pgfqpoint{2.807067in}{2.202541in}}%
\pgfpathlineto{\pgfqpoint{2.807907in}{2.202728in}}%
\pgfpathlineto{\pgfqpoint{2.808550in}{2.199967in}}%
\pgfpathlineto{\pgfqpoint{2.808600in}{2.201677in}}%
\pgfpathlineto{\pgfqpoint{2.809138in}{1.532019in}}%
\pgfpathlineto{\pgfqpoint{2.809845in}{1.532052in}}%
\pgfpathlineto{\pgfqpoint{2.811114in}{1.532181in}}%
\pgfpathlineto{\pgfqpoint{2.812618in}{2.202872in}}%
\pgfpathlineto{\pgfqpoint{2.812667in}{2.202550in}}%
\pgfpathlineto{\pgfqpoint{2.812717in}{2.202862in}}%
\pgfpathlineto{\pgfqpoint{2.813261in}{2.201441in}}%
\pgfpathlineto{\pgfqpoint{2.813310in}{2.202478in}}%
\pgfpathlineto{\pgfqpoint{2.813805in}{1.856008in}}%
\pgfpathlineto{\pgfqpoint{2.814114in}{1.532032in}}%
\pgfpathlineto{\pgfqpoint{2.814731in}{1.532059in}}%
\pgfpathlineto{\pgfqpoint{2.816056in}{1.532148in}}%
\pgfpathlineto{\pgfqpoint{2.816065in}{1.532491in}}%
\pgfpathlineto{\pgfqpoint{2.817124in}{2.202689in}}%
\pgfpathlineto{\pgfqpoint{2.818361in}{2.201516in}}%
\pgfpathlineto{\pgfqpoint{2.818410in}{2.201934in}}%
\pgfpathlineto{\pgfqpoint{2.818460in}{2.199704in}}%
\pgfpathlineto{\pgfqpoint{2.818509in}{2.201570in}}%
\pgfpathlineto{\pgfqpoint{2.819029in}{1.532017in}}%
\pgfpathlineto{\pgfqpoint{2.819753in}{1.532052in}}%
\pgfpathlineto{\pgfqpoint{2.821007in}{1.532228in}}%
\pgfpathlineto{\pgfqpoint{2.822140in}{2.202816in}}%
\pgfpathlineto{\pgfqpoint{2.822337in}{2.202882in}}%
\pgfpathlineto{\pgfqpoint{2.822881in}{2.202053in}}%
\pgfpathlineto{\pgfqpoint{2.822931in}{2.202778in}}%
\pgfpathlineto{\pgfqpoint{2.823079in}{2.201136in}}%
\pgfpathlineto{\pgfqpoint{2.823129in}{2.202622in}}%
\pgfpathlineto{\pgfqpoint{2.824050in}{1.532033in}}%
\pgfpathlineto{\pgfqpoint{2.824736in}{1.532058in}}%
\pgfpathlineto{\pgfqpoint{2.825949in}{1.532166in}}%
\pgfpathlineto{\pgfqpoint{2.825992in}{1.870026in}}%
\pgfpathlineto{\pgfqpoint{2.826753in}{2.202571in}}%
\pgfpathlineto{\pgfqpoint{2.827148in}{2.202808in}}%
\pgfpathlineto{\pgfqpoint{2.827494in}{2.202570in}}%
\pgfpathlineto{\pgfqpoint{2.827544in}{2.202852in}}%
\pgfpathlineto{\pgfqpoint{2.828088in}{2.201637in}}%
\pgfpathlineto{\pgfqpoint{2.828137in}{2.202473in}}%
\pgfpathlineto{\pgfqpoint{2.828187in}{2.200544in}}%
\pgfpathlineto{\pgfqpoint{2.828236in}{2.202331in}}%
\pgfpathlineto{\pgfqpoint{2.829005in}{1.532034in}}%
\pgfpathlineto{\pgfqpoint{2.829545in}{1.532056in}}%
\pgfpathlineto{\pgfqpoint{2.830898in}{1.532232in}}%
\pgfpathlineto{\pgfqpoint{2.832099in}{2.202826in}}%
\pgfpathlineto{\pgfqpoint{2.832296in}{2.202873in}}%
\pgfpathlineto{\pgfqpoint{2.832939in}{2.201635in}}%
\pgfpathlineto{\pgfqpoint{2.832989in}{2.202615in}}%
\pgfpathlineto{\pgfqpoint{2.833038in}{2.200546in}}%
\pgfpathlineto{\pgfqpoint{2.833088in}{2.202508in}}%
\pgfpathlineto{\pgfqpoint{2.833900in}{1.532030in}}%
\pgfpathlineto{\pgfqpoint{2.834421in}{1.532058in}}%
\pgfpathlineto{\pgfqpoint{2.835831in}{1.532117in}}%
\pgfpathlineto{\pgfqpoint{2.835842in}{1.532198in}}%
\pgfpathlineto{\pgfqpoint{2.837058in}{2.202835in}}%
\pgfpathlineto{\pgfqpoint{2.837107in}{2.202524in}}%
\pgfpathlineto{\pgfqpoint{2.837454in}{2.202881in}}%
\pgfpathlineto{\pgfqpoint{2.837701in}{2.202117in}}%
\pgfpathlineto{\pgfqpoint{2.837750in}{2.202777in}}%
\pgfpathlineto{\pgfqpoint{2.837998in}{2.187029in}}%
\pgfpathlineto{\pgfqpoint{2.838047in}{2.202502in}}%
\pgfpathlineto{\pgfqpoint{2.838809in}{1.532020in}}%
\pgfpathlineto{\pgfqpoint{2.839409in}{1.532045in}}%
\pgfpathlineto{\pgfqpoint{2.840787in}{1.532179in}}%
\pgfpathlineto{\pgfqpoint{2.840827in}{1.854420in}}%
\pgfpathlineto{\pgfqpoint{2.841595in}{2.202596in}}%
\pgfpathlineto{\pgfqpoint{2.842090in}{2.202866in}}%
\pgfpathlineto{\pgfqpoint{2.842337in}{2.202467in}}%
\pgfpathlineto{\pgfqpoint{2.842386in}{2.202884in}}%
\pgfpathlineto{\pgfqpoint{2.842831in}{2.201561in}}%
\pgfpathlineto{\pgfqpoint{2.842881in}{2.202638in}}%
\pgfpathlineto{\pgfqpoint{2.843375in}{2.123273in}}%
\pgfpathlineto{\pgfqpoint{2.843808in}{1.532030in}}%
\pgfpathlineto{\pgfqpoint{2.844396in}{1.532052in}}%
\pgfpathlineto{\pgfqpoint{2.845731in}{1.532171in}}%
\pgfpathlineto{\pgfqpoint{2.845774in}{1.901326in}}%
\pgfpathlineto{\pgfqpoint{2.846541in}{2.202573in}}%
\pgfpathlineto{\pgfqpoint{2.846937in}{2.202805in}}%
\pgfpathlineto{\pgfqpoint{2.847283in}{2.202572in}}%
\pgfpathlineto{\pgfqpoint{2.847333in}{2.202843in}}%
\pgfpathlineto{\pgfqpoint{2.847877in}{2.201639in}}%
\pgfpathlineto{\pgfqpoint{2.847926in}{2.202447in}}%
\pgfpathlineto{\pgfqpoint{2.847975in}{2.200591in}}%
\pgfpathlineto{\pgfqpoint{2.848025in}{2.202301in}}%
\pgfpathlineto{\pgfqpoint{2.848756in}{1.532025in}}%
\pgfpathlineto{\pgfqpoint{2.849342in}{1.532047in}}%
\pgfpathlineto{\pgfqpoint{2.850681in}{1.532268in}}%
\pgfpathlineto{\pgfqpoint{2.851795in}{2.202719in}}%
\pgfpathlineto{\pgfqpoint{2.853032in}{2.201295in}}%
\pgfpathlineto{\pgfqpoint{2.853081in}{2.201746in}}%
\pgfpathlineto{\pgfqpoint{2.853635in}{1.595667in}}%
\pgfpathlineto{\pgfqpoint{2.853674in}{1.532027in}}%
\pgfpathlineto{\pgfqpoint{2.854385in}{1.532067in}}%
\pgfpathlineto{\pgfqpoint{2.855614in}{1.532120in}}%
\pgfpathlineto{\pgfqpoint{2.855623in}{1.532180in}}%
\pgfpathlineto{\pgfqpoint{2.855664in}{1.874111in}}%
\pgfpathlineto{\pgfqpoint{2.856430in}{2.202618in}}%
\pgfpathlineto{\pgfqpoint{2.856925in}{2.202893in}}%
\pgfpathlineto{\pgfqpoint{2.857172in}{2.202530in}}%
\pgfpathlineto{\pgfqpoint{2.857222in}{2.202915in}}%
\pgfpathlineto{\pgfqpoint{2.857667in}{2.201758in}}%
\pgfpathlineto{\pgfqpoint{2.857716in}{2.202672in}}%
\pgfpathlineto{\pgfqpoint{2.857766in}{2.201006in}}%
\pgfpathlineto{\pgfqpoint{2.857815in}{2.202568in}}%
\pgfpathlineto{\pgfqpoint{2.858611in}{1.532029in}}%
\pgfpathlineto{\pgfqpoint{2.859247in}{1.532051in}}%
\pgfpathlineto{\pgfqpoint{2.860569in}{1.532196in}}%
\pgfpathlineto{\pgfqpoint{2.861802in}{2.202916in}}%
\pgfpathlineto{\pgfqpoint{2.862000in}{2.202963in}}%
\pgfpathlineto{\pgfqpoint{2.862544in}{2.188453in}}%
\pgfpathlineto{\pgfqpoint{2.862593in}{2.202793in}}%
\pgfpathlineto{\pgfqpoint{2.863137in}{2.144357in}}%
\pgfpathlineto{\pgfqpoint{2.863554in}{1.532026in}}%
\pgfpathlineto{\pgfqpoint{2.864236in}{1.532067in}}%
\pgfpathlineto{\pgfqpoint{2.865505in}{1.532121in}}%
\pgfpathlineto{\pgfqpoint{2.865516in}{1.532223in}}%
\pgfpathlineto{\pgfqpoint{2.866644in}{2.202836in}}%
\pgfpathlineto{\pgfqpoint{2.866941in}{2.202921in}}%
\pgfpathlineto{\pgfqpoint{2.867386in}{2.201964in}}%
\pgfpathlineto{\pgfqpoint{2.867436in}{2.202810in}}%
\pgfpathlineto{\pgfqpoint{2.867584in}{2.199940in}}%
\pgfpathlineto{\pgfqpoint{2.867633in}{2.202659in}}%
\pgfpathlineto{\pgfqpoint{2.868519in}{1.532029in}}%
\pgfpathlineto{\pgfqpoint{2.868894in}{1.532042in}}%
\pgfpathlineto{\pgfqpoint{2.870461in}{1.532200in}}%
\pgfpathlineto{\pgfqpoint{2.871673in}{2.202845in}}%
\pgfpathlineto{\pgfqpoint{2.871722in}{2.202531in}}%
\pgfpathlineto{\pgfqpoint{2.872068in}{2.202891in}}%
\pgfpathlineto{\pgfqpoint{2.872316in}{2.202119in}}%
\pgfpathlineto{\pgfqpoint{2.872365in}{2.202787in}}%
\pgfpathlineto{\pgfqpoint{2.872612in}{2.184629in}}%
\pgfpathlineto{\pgfqpoint{2.872662in}{2.202506in}}%
\pgfpathlineto{\pgfqpoint{2.873486in}{1.532027in}}%
\pgfpathlineto{\pgfqpoint{2.874023in}{1.532062in}}%
\pgfpathlineto{\pgfqpoint{2.875396in}{1.532120in}}%
\pgfpathlineto{\pgfqpoint{2.875407in}{1.532214in}}%
\pgfpathlineto{\pgfqpoint{2.876537in}{2.202864in}}%
\pgfpathlineto{\pgfqpoint{2.876587in}{2.202375in}}%
\pgfpathlineto{\pgfqpoint{2.877032in}{2.202962in}}%
\pgfpathlineto{\pgfqpoint{2.877180in}{2.201734in}}%
\pgfpathlineto{\pgfqpoint{2.877230in}{2.202915in}}%
\pgfpathlineto{\pgfqpoint{2.877279in}{2.201064in}}%
\pgfpathlineto{\pgfqpoint{2.877329in}{2.202870in}}%
\pgfpathlineto{\pgfqpoint{2.878401in}{1.532035in}}%
\pgfpathlineto{\pgfqpoint{2.878817in}{1.532049in}}%
\pgfpathlineto{\pgfqpoint{2.880356in}{1.532474in}}%
\pgfpathlineto{\pgfqpoint{2.881323in}{2.202623in}}%
\pgfpathlineto{\pgfqpoint{2.881372in}{2.202561in}}%
\pgfpathlineto{\pgfqpoint{2.882114in}{2.202684in}}%
\pgfpathlineto{\pgfqpoint{2.882213in}{2.202625in}}%
\pgfpathlineto{\pgfqpoint{2.882658in}{2.201185in}}%
\pgfpathlineto{\pgfqpoint{2.882708in}{2.201927in}}%
\pgfpathlineto{\pgfqpoint{2.883320in}{1.532017in}}%
\pgfpathlineto{\pgfqpoint{2.884563in}{1.532071in}}%
\pgfpathlineto{\pgfqpoint{2.885298in}{1.532216in}}%
\pgfpathlineto{\pgfqpoint{2.886428in}{2.202831in}}%
\pgfpathlineto{\pgfqpoint{2.886477in}{2.202375in}}%
\pgfpathlineto{\pgfqpoint{2.886922in}{2.202930in}}%
\pgfpathlineto{\pgfqpoint{2.887071in}{2.201840in}}%
\pgfpathlineto{\pgfqpoint{2.887120in}{2.202884in}}%
\pgfpathlineto{\pgfqpoint{2.887269in}{2.184118in}}%
\pgfpathlineto{\pgfqpoint{2.887318in}{2.202770in}}%
\pgfpathlineto{\pgfqpoint{2.888325in}{1.532030in}}%
\pgfpathlineto{\pgfqpoint{2.888715in}{1.532044in}}%
\pgfpathlineto{\pgfqpoint{2.890241in}{1.532174in}}%
\pgfpathlineto{\pgfqpoint{2.890283in}{1.927696in}}%
\pgfpathlineto{\pgfqpoint{2.891045in}{2.202523in}}%
\pgfpathlineto{\pgfqpoint{2.891243in}{2.202659in}}%
\pgfpathlineto{\pgfqpoint{2.891787in}{2.202558in}}%
\pgfpathlineto{\pgfqpoint{2.891837in}{2.202777in}}%
\pgfpathlineto{\pgfqpoint{2.892381in}{2.201739in}}%
\pgfpathlineto{\pgfqpoint{2.892430in}{2.202350in}}%
\pgfpathlineto{\pgfqpoint{2.892480in}{2.201150in}}%
\pgfpathlineto{\pgfqpoint{2.892529in}{2.202193in}}%
\pgfpathlineto{\pgfqpoint{2.893295in}{1.532035in}}%
\pgfpathlineto{\pgfqpoint{2.894279in}{1.532071in}}%
\pgfpathlineto{\pgfqpoint{2.895189in}{1.532218in}}%
\pgfpathlineto{\pgfqpoint{2.896389in}{2.202791in}}%
\pgfpathlineto{\pgfqpoint{2.896587in}{2.202843in}}%
\pgfpathlineto{\pgfqpoint{2.897230in}{2.201707in}}%
\pgfpathlineto{\pgfqpoint{2.897279in}{2.202600in}}%
\pgfpathlineto{\pgfqpoint{2.897329in}{2.201003in}}%
\pgfpathlineto{\pgfqpoint{2.897378in}{2.202497in}}%
\pgfpathlineto{\pgfqpoint{2.898206in}{1.532030in}}%
\pgfpathlineto{\pgfqpoint{2.898891in}{1.532055in}}%
\pgfpathlineto{\pgfqpoint{2.900131in}{1.532169in}}%
\pgfpathlineto{\pgfqpoint{2.900174in}{1.885949in}}%
\pgfpathlineto{\pgfqpoint{2.900935in}{2.202542in}}%
\pgfpathlineto{\pgfqpoint{2.901331in}{2.202776in}}%
\pgfpathlineto{\pgfqpoint{2.901677in}{2.202525in}}%
\pgfpathlineto{\pgfqpoint{2.901726in}{2.202816in}}%
\pgfpathlineto{\pgfqpoint{2.902270in}{2.201527in}}%
\pgfpathlineto{\pgfqpoint{2.902320in}{2.202431in}}%
\pgfpathlineto{\pgfqpoint{2.902369in}{2.199448in}}%
\pgfpathlineto{\pgfqpoint{2.902419in}{2.202285in}}%
\pgfpathlineto{\pgfqpoint{2.903101in}{1.532011in}}%
\pgfpathlineto{\pgfqpoint{2.903680in}{1.532045in}}%
\pgfpathlineto{\pgfqpoint{2.905078in}{1.532179in}}%
\pgfpathlineto{\pgfqpoint{2.905117in}{1.849756in}}%
\pgfpathlineto{\pgfqpoint{2.905884in}{2.202535in}}%
\pgfpathlineto{\pgfqpoint{2.906180in}{2.202727in}}%
\pgfpathlineto{\pgfqpoint{2.906625in}{2.202569in}}%
\pgfpathlineto{\pgfqpoint{2.906675in}{2.202803in}}%
\pgfpathlineto{\pgfqpoint{2.907219in}{2.201749in}}%
\pgfpathlineto{\pgfqpoint{2.907268in}{2.202396in}}%
\pgfpathlineto{\pgfqpoint{2.907318in}{2.201142in}}%
\pgfpathlineto{\pgfqpoint{2.907367in}{2.202245in}}%
\pgfpathlineto{\pgfqpoint{2.908102in}{1.532027in}}%
\pgfpathlineto{\pgfqpoint{2.909034in}{1.532077in}}%
\pgfpathlineto{\pgfqpoint{2.910015in}{1.532123in}}%
\pgfpathlineto{\pgfqpoint{2.910025in}{1.532224in}}%
\pgfpathlineto{\pgfqpoint{2.911155in}{2.202841in}}%
\pgfpathlineto{\pgfqpoint{2.911204in}{2.202393in}}%
\pgfpathlineto{\pgfqpoint{2.911649in}{2.202937in}}%
\pgfpathlineto{\pgfqpoint{2.911798in}{2.201868in}}%
\pgfpathlineto{\pgfqpoint{2.911847in}{2.202890in}}%
\pgfpathlineto{\pgfqpoint{2.911995in}{2.187979in}}%
\pgfpathlineto{\pgfqpoint{2.912045in}{2.202778in}}%
\pgfpathlineto{\pgfqpoint{2.913050in}{1.532033in}}%
\pgfpathlineto{\pgfqpoint{2.913389in}{1.532048in}}%
\pgfpathlineto{\pgfqpoint{2.914969in}{1.532188in}}%
\pgfpathlineto{\pgfqpoint{2.916376in}{2.202841in}}%
\pgfpathlineto{\pgfqpoint{2.916425in}{2.202536in}}%
\pgfpathlineto{\pgfqpoint{2.916573in}{2.202834in}}%
\pgfpathlineto{\pgfqpoint{2.917018in}{2.201773in}}%
\pgfpathlineto{\pgfqpoint{2.917068in}{2.202559in}}%
\pgfpathlineto{\pgfqpoint{2.917117in}{2.201196in}}%
\pgfpathlineto{\pgfqpoint{2.917167in}{2.202446in}}%
\pgfpathlineto{\pgfqpoint{2.917937in}{1.532022in}}%
\pgfpathlineto{\pgfqpoint{2.918921in}{1.532061in}}%
\pgfpathlineto{\pgfqpoint{2.919914in}{1.532183in}}%
\pgfpathlineto{\pgfqpoint{2.919954in}{1.860611in}}%
\pgfpathlineto{\pgfqpoint{2.920724in}{2.202586in}}%
\pgfpathlineto{\pgfqpoint{2.921218in}{2.202860in}}%
\pgfpathlineto{\pgfqpoint{2.921466in}{2.202454in}}%
\pgfpathlineto{\pgfqpoint{2.921515in}{2.202882in}}%
\pgfpathlineto{\pgfqpoint{2.921960in}{2.201531in}}%
\pgfpathlineto{\pgfqpoint{2.922010in}{2.202643in}}%
\pgfpathlineto{\pgfqpoint{2.922504in}{2.131155in}}%
\pgfpathlineto{\pgfqpoint{2.922911in}{1.532027in}}%
\pgfpathlineto{\pgfqpoint{2.923525in}{1.532063in}}%
\pgfpathlineto{\pgfqpoint{2.924850in}{1.532119in}}%
\pgfpathlineto{\pgfqpoint{2.924859in}{1.532180in}}%
\pgfpathlineto{\pgfqpoint{2.924900in}{1.879481in}}%
\pgfpathlineto{\pgfqpoint{2.925666in}{2.202586in}}%
\pgfpathlineto{\pgfqpoint{2.926161in}{2.202852in}}%
\pgfpathlineto{\pgfqpoint{2.926408in}{2.202487in}}%
\pgfpathlineto{\pgfqpoint{2.926458in}{2.202867in}}%
\pgfpathlineto{\pgfqpoint{2.926903in}{2.201686in}}%
\pgfpathlineto{\pgfqpoint{2.926952in}{2.202611in}}%
\pgfpathlineto{\pgfqpoint{2.927002in}{2.200842in}}%
\pgfpathlineto{\pgfqpoint{2.927051in}{2.202503in}}%
\pgfpathlineto{\pgfqpoint{2.927866in}{1.532029in}}%
\pgfpathlineto{\pgfqpoint{2.928440in}{1.532049in}}%
\pgfpathlineto{\pgfqpoint{2.929805in}{1.532182in}}%
\pgfpathlineto{\pgfqpoint{2.931306in}{2.202852in}}%
\pgfpathlineto{\pgfqpoint{2.932048in}{2.200367in}}%
\pgfpathlineto{\pgfqpoint{2.932097in}{2.202308in}}%
\pgfpathlineto{\pgfqpoint{2.932830in}{1.532025in}}%
\pgfpathlineto{\pgfqpoint{2.933366in}{1.532064in}}%
\pgfpathlineto{\pgfqpoint{2.934741in}{1.532120in}}%
\pgfpathlineto{\pgfqpoint{2.934753in}{1.532258in}}%
\pgfpathlineto{\pgfqpoint{2.935874in}{2.202708in}}%
\pgfpathlineto{\pgfqpoint{2.937110in}{2.201364in}}%
\pgfpathlineto{\pgfqpoint{2.937160in}{2.201644in}}%
\pgfpathlineto{\pgfqpoint{2.937407in}{1.907114in}}%
\pgfpathlineto{\pgfqpoint{2.937719in}{1.532002in}}%
\pgfpathlineto{\pgfqpoint{2.938301in}{1.532067in}}%
\pgfpathlineto{\pgfqpoint{2.939684in}{1.532114in}}%
\pgfpathlineto{\pgfqpoint{2.939694in}{1.532160in}}%
\pgfpathlineto{\pgfqpoint{2.939722in}{1.606351in}}%
\pgfpathlineto{\pgfqpoint{2.940476in}{2.202516in}}%
\pgfpathlineto{\pgfqpoint{2.940525in}{2.202468in}}%
\pgfpathlineto{\pgfqpoint{2.941366in}{2.202783in}}%
\pgfpathlineto{\pgfqpoint{2.942009in}{2.201256in}}%
\pgfpathlineto{\pgfqpoint{2.942058in}{2.202043in}}%
\pgfpathlineto{\pgfqpoint{2.942706in}{1.532028in}}%
\pgfpathlineto{\pgfqpoint{2.943976in}{1.532075in}}%
\pgfpathlineto{\pgfqpoint{2.944641in}{1.532181in}}%
\pgfpathlineto{\pgfqpoint{2.946145in}{2.202874in}}%
\pgfpathlineto{\pgfqpoint{2.946788in}{2.201394in}}%
\pgfpathlineto{\pgfqpoint{2.946837in}{2.202473in}}%
\pgfpathlineto{\pgfqpoint{2.947602in}{1.560933in}}%
\pgfpathlineto{\pgfqpoint{2.947662in}{1.532034in}}%
\pgfpathlineto{\pgfqpoint{2.948347in}{1.532059in}}%
\pgfpathlineto{\pgfqpoint{2.949589in}{1.532235in}}%
\pgfpathlineto{\pgfqpoint{2.950779in}{2.202822in}}%
\pgfpathlineto{\pgfqpoint{2.950878in}{2.202853in}}%
\pgfpathlineto{\pgfqpoint{2.951620in}{2.201871in}}%
\pgfpathlineto{\pgfqpoint{2.951669in}{2.202623in}}%
\pgfpathlineto{\pgfqpoint{2.951718in}{2.201396in}}%
\pgfpathlineto{\pgfqpoint{2.951768in}{2.202518in}}%
\pgfpathlineto{\pgfqpoint{2.952547in}{1.562314in}}%
\pgfpathlineto{\pgfqpoint{2.952596in}{1.532028in}}%
\pgfpathlineto{\pgfqpoint{2.953322in}{1.532067in}}%
\pgfpathlineto{\pgfqpoint{2.954523in}{1.532119in}}%
\pgfpathlineto{\pgfqpoint{2.954532in}{1.532181in}}%
\pgfpathlineto{\pgfqpoint{2.956031in}{2.202886in}}%
\pgfpathlineto{\pgfqpoint{2.956080in}{2.202540in}}%
\pgfpathlineto{\pgfqpoint{2.956130in}{2.202878in}}%
\pgfpathlineto{\pgfqpoint{2.956674in}{2.201324in}}%
\pgfpathlineto{\pgfqpoint{2.956723in}{2.202508in}}%
\pgfpathlineto{\pgfqpoint{2.957561in}{1.532031in}}%
\pgfpathlineto{\pgfqpoint{2.958696in}{1.532080in}}%
\pgfpathlineto{\pgfqpoint{2.959475in}{1.532156in}}%
\pgfpathlineto{\pgfqpoint{2.959505in}{1.612219in}}%
\pgfpathlineto{\pgfqpoint{2.960277in}{2.202584in}}%
\pgfpathlineto{\pgfqpoint{2.960673in}{2.202823in}}%
\pgfpathlineto{\pgfqpoint{2.961019in}{2.202535in}}%
\pgfpathlineto{\pgfqpoint{2.961069in}{2.202870in}}%
\pgfpathlineto{\pgfqpoint{2.961613in}{2.201371in}}%
\pgfpathlineto{\pgfqpoint{2.961662in}{2.202506in}}%
\pgfpathlineto{\pgfqpoint{2.962502in}{1.532028in}}%
\pgfpathlineto{\pgfqpoint{2.963879in}{1.532079in}}%
\pgfpathlineto{\pgfqpoint{2.964425in}{1.532221in}}%
\pgfpathlineto{\pgfqpoint{2.965554in}{2.202857in}}%
\pgfpathlineto{\pgfqpoint{2.965604in}{2.202452in}}%
\pgfpathlineto{\pgfqpoint{2.966049in}{2.202946in}}%
\pgfpathlineto{\pgfqpoint{2.966296in}{2.201707in}}%
\pgfpathlineto{\pgfqpoint{2.966346in}{2.202848in}}%
\pgfpathlineto{\pgfqpoint{2.966395in}{2.200869in}}%
\pgfpathlineto{\pgfqpoint{2.966445in}{2.202786in}}%
\pgfpathlineto{\pgfqpoint{2.967438in}{1.532028in}}%
\pgfpathlineto{\pgfqpoint{2.967822in}{1.532042in}}%
\pgfpathlineto{\pgfqpoint{2.969371in}{1.532229in}}%
\pgfpathlineto{\pgfqpoint{2.970503in}{2.202825in}}%
\pgfpathlineto{\pgfqpoint{2.970701in}{2.202891in}}%
\pgfpathlineto{\pgfqpoint{2.971245in}{2.202053in}}%
\pgfpathlineto{\pgfqpoint{2.971295in}{2.202784in}}%
\pgfpathlineto{\pgfqpoint{2.971443in}{2.201097in}}%
\pgfpathlineto{\pgfqpoint{2.971492in}{2.202628in}}%
\pgfpathlineto{\pgfqpoint{2.972363in}{1.532032in}}%
\pgfpathlineto{\pgfqpoint{2.973022in}{1.532054in}}%
\pgfpathlineto{\pgfqpoint{2.974320in}{1.532493in}}%
\pgfpathlineto{\pgfqpoint{2.975379in}{2.202686in}}%
\pgfpathlineto{\pgfqpoint{2.976615in}{2.201527in}}%
\pgfpathlineto{\pgfqpoint{2.976665in}{2.201940in}}%
\pgfpathlineto{\pgfqpoint{2.976714in}{2.199946in}}%
\pgfpathlineto{\pgfqpoint{2.976764in}{2.201583in}}%
\pgfpathlineto{\pgfqpoint{2.977283in}{1.532014in}}%
\pgfpathlineto{\pgfqpoint{2.978029in}{1.532052in}}%
\pgfpathlineto{\pgfqpoint{2.979262in}{1.532250in}}%
\pgfpathlineto{\pgfqpoint{2.980361in}{2.202687in}}%
\pgfpathlineto{\pgfqpoint{2.981597in}{2.201532in}}%
\pgfpathlineto{\pgfqpoint{2.981647in}{2.201690in}}%
\pgfpathlineto{\pgfqpoint{2.981696in}{2.200200in}}%
\pgfpathlineto{\pgfqpoint{2.981746in}{2.201044in}}%
\pgfpathlineto{\pgfqpoint{2.982288in}{1.532031in}}%
\pgfpathlineto{\pgfqpoint{2.983027in}{1.532066in}}%
\pgfpathlineto{\pgfqpoint{2.984202in}{1.532148in}}%
\pgfpathlineto{\pgfqpoint{2.984210in}{1.532473in}}%
\pgfpathlineto{\pgfqpoint{2.985266in}{2.202677in}}%
\pgfpathlineto{\pgfqpoint{2.986601in}{2.200750in}}%
\pgfpathlineto{\pgfqpoint{2.986650in}{2.201596in}}%
\pgfpathlineto{\pgfqpoint{2.987253in}{1.532034in}}%
\pgfpathlineto{\pgfqpoint{2.988039in}{1.532063in}}%
\pgfpathlineto{\pgfqpoint{2.989152in}{1.532215in}}%
\pgfpathlineto{\pgfqpoint{2.990352in}{2.202809in}}%
\pgfpathlineto{\pgfqpoint{2.990451in}{2.202839in}}%
\pgfpathlineto{\pgfqpoint{2.991193in}{2.201742in}}%
\pgfpathlineto{\pgfqpoint{2.991242in}{2.202593in}}%
\pgfpathlineto{\pgfqpoint{2.991292in}{2.201080in}}%
\pgfpathlineto{\pgfqpoint{2.991341in}{2.202485in}}%
\pgfpathlineto{\pgfqpoint{2.992161in}{1.532028in}}%
\pgfpathlineto{\pgfqpoint{2.992889in}{1.532067in}}%
\pgfpathlineto{\pgfqpoint{2.994086in}{1.532119in}}%
\pgfpathlineto{\pgfqpoint{2.994098in}{1.532220in}}%
\pgfpathlineto{\pgfqpoint{2.995227in}{2.202844in}}%
\pgfpathlineto{\pgfqpoint{2.995277in}{2.202434in}}%
\pgfpathlineto{\pgfqpoint{2.995722in}{2.202938in}}%
\pgfpathlineto{\pgfqpoint{2.995969in}{2.201693in}}%
\pgfpathlineto{\pgfqpoint{2.996019in}{2.202844in}}%
\pgfpathlineto{\pgfqpoint{2.996068in}{2.200848in}}%
\pgfpathlineto{\pgfqpoint{2.996117in}{2.202783in}}%
\pgfpathlineto{\pgfqpoint{2.997088in}{1.532032in}}%
\pgfpathlineto{\pgfqpoint{2.997496in}{1.532051in}}%
\pgfpathlineto{\pgfqpoint{2.999038in}{1.532147in}}%
\pgfpathlineto{\pgfqpoint{2.999047in}{1.532490in}}%
\pgfpathlineto{\pgfqpoint{3.000107in}{2.202668in}}%
\pgfpathlineto{\pgfqpoint{3.001343in}{2.201434in}}%
\pgfpathlineto{\pgfqpoint{3.001392in}{2.201890in}}%
\pgfpathlineto{\pgfqpoint{3.001590in}{2.096942in}}%
\pgfpathlineto{\pgfqpoint{3.002053in}{1.532030in}}%
\pgfpathlineto{\pgfqpoint{3.002680in}{1.532062in}}%
\pgfpathlineto{\pgfqpoint{3.003977in}{1.532118in}}%
\pgfpathlineto{\pgfqpoint{3.003988in}{1.532203in}}%
\pgfpathlineto{\pgfqpoint{3.005220in}{2.202864in}}%
\pgfpathlineto{\pgfqpoint{3.005863in}{2.201305in}}%
\pgfpathlineto{\pgfqpoint{3.005418in}{2.202914in}}%
\pgfpathlineto{\pgfqpoint{3.005913in}{2.202823in}}%
\pgfpathlineto{\pgfqpoint{3.006988in}{1.532031in}}%
\pgfpathlineto{\pgfqpoint{3.007953in}{1.532066in}}%
\pgfpathlineto{\pgfqpoint{3.008939in}{1.532964in}}%
\pgfpathlineto{\pgfqpoint{3.009879in}{2.202572in}}%
\pgfpathlineto{\pgfqpoint{3.009928in}{2.202519in}}%
\pgfpathlineto{\pgfqpoint{3.010670in}{2.202661in}}%
\pgfpathlineto{\pgfqpoint{3.010769in}{2.202607in}}%
\pgfpathlineto{\pgfqpoint{3.011313in}{2.196626in}}%
\pgfpathlineto{\pgfqpoint{3.011362in}{2.201620in}}%
\pgfpathlineto{\pgfqpoint{3.011901in}{1.532008in}}%
\pgfpathlineto{\pgfqpoint{3.012592in}{1.532048in}}%
\pgfpathlineto{\pgfqpoint{3.013878in}{1.532179in}}%
\pgfpathlineto{\pgfqpoint{3.013918in}{1.868208in}}%
\pgfpathlineto{\pgfqpoint{3.014685in}{2.202531in}}%
\pgfpathlineto{\pgfqpoint{3.015180in}{2.202801in}}%
\pgfpathlineto{\pgfqpoint{3.015427in}{2.202423in}}%
\pgfpathlineto{\pgfqpoint{3.015476in}{2.202819in}}%
\pgfpathlineto{\pgfqpoint{3.015922in}{2.201590in}}%
\pgfpathlineto{\pgfqpoint{3.015971in}{2.202571in}}%
\pgfpathlineto{\pgfqpoint{3.016020in}{2.200553in}}%
\pgfpathlineto{\pgfqpoint{3.016070in}{2.202466in}}%
\pgfpathlineto{\pgfqpoint{3.016884in}{1.532034in}}%
\pgfpathlineto{\pgfqpoint{3.017410in}{1.532054in}}%
\pgfpathlineto{\pgfqpoint{3.018828in}{1.532346in}}%
\pgfpathlineto{\pgfqpoint{3.019844in}{2.202705in}}%
\pgfpathlineto{\pgfqpoint{3.019894in}{2.202359in}}%
\pgfpathlineto{\pgfqpoint{3.020438in}{2.202854in}}%
\pgfpathlineto{\pgfqpoint{3.020586in}{2.201993in}}%
\pgfpathlineto{\pgfqpoint{3.020636in}{2.202805in}}%
\pgfpathlineto{\pgfqpoint{3.020784in}{2.201209in}}%
\pgfpathlineto{\pgfqpoint{3.020834in}{2.202698in}}%
\pgfpathlineto{\pgfqpoint{3.021808in}{1.532026in}}%
\pgfpathlineto{\pgfqpoint{3.022639in}{1.532055in}}%
\pgfpathlineto{\pgfqpoint{3.023771in}{1.532223in}}%
\pgfpathlineto{\pgfqpoint{3.024899in}{2.202752in}}%
\pgfpathlineto{\pgfqpoint{3.025196in}{2.202836in}}%
\pgfpathlineto{\pgfqpoint{3.025641in}{2.201860in}}%
\pgfpathlineto{\pgfqpoint{3.025690in}{2.202727in}}%
\pgfpathlineto{\pgfqpoint{3.025839in}{2.198248in}}%
\pgfpathlineto{\pgfqpoint{3.025888in}{2.202573in}}%
\pgfpathlineto{\pgfqpoint{3.026774in}{1.532029in}}%
\pgfpathlineto{\pgfqpoint{3.027150in}{1.532042in}}%
\pgfpathlineto{\pgfqpoint{3.028714in}{1.532182in}}%
\pgfpathlineto{\pgfqpoint{3.030206in}{2.202755in}}%
\pgfpathlineto{\pgfqpoint{3.030948in}{2.201262in}}%
\pgfpathlineto{\pgfqpoint{3.030997in}{2.202186in}}%
\pgfpathlineto{\pgfqpoint{3.031680in}{1.535066in}}%
\pgfpathlineto{\pgfqpoint{3.031715in}{1.532031in}}%
\pgfpathlineto{\pgfqpoint{3.032430in}{1.532064in}}%
\pgfpathlineto{\pgfqpoint{3.033656in}{1.532148in}}%
\pgfpathlineto{\pgfqpoint{3.033665in}{1.532492in}}%
\pgfpathlineto{\pgfqpoint{3.034724in}{2.202610in}}%
\pgfpathlineto{\pgfqpoint{3.035961in}{2.201468in}}%
\pgfpathlineto{\pgfqpoint{3.036010in}{2.201870in}}%
\pgfpathlineto{\pgfqpoint{3.036060in}{2.200076in}}%
\pgfpathlineto{\pgfqpoint{3.036109in}{2.201519in}}%
\pgfpathlineto{\pgfqpoint{3.036629in}{1.532018in}}%
\pgfpathlineto{\pgfqpoint{3.037388in}{1.532067in}}%
\pgfpathlineto{\pgfqpoint{3.038595in}{1.532118in}}%
\pgfpathlineto{\pgfqpoint{3.038605in}{1.532181in}}%
\pgfpathlineto{\pgfqpoint{3.040110in}{2.202789in}}%
\pgfpathlineto{\pgfqpoint{3.040159in}{2.202450in}}%
\pgfpathlineto{\pgfqpoint{3.040209in}{2.202779in}}%
\pgfpathlineto{\pgfqpoint{3.040753in}{2.201229in}}%
\pgfpathlineto{\pgfqpoint{3.040802in}{2.202396in}}%
\pgfpathlineto{\pgfqpoint{3.041631in}{1.532032in}}%
\pgfpathlineto{\pgfqpoint{3.042765in}{1.532080in}}%
\pgfpathlineto{\pgfqpoint{3.043548in}{1.532158in}}%
\pgfpathlineto{\pgfqpoint{3.043578in}{1.615378in}}%
\pgfpathlineto{\pgfqpoint{3.044350in}{2.202492in}}%
\pgfpathlineto{\pgfqpoint{3.044746in}{2.202729in}}%
\pgfpathlineto{\pgfqpoint{3.045092in}{2.202499in}}%
\pgfpathlineto{\pgfqpoint{3.045142in}{2.202774in}}%
\pgfpathlineto{\pgfqpoint{3.045686in}{2.201597in}}%
\pgfpathlineto{\pgfqpoint{3.045735in}{2.202393in}}%
\pgfpathlineto{\pgfqpoint{3.045785in}{2.200687in}}%
\pgfpathlineto{\pgfqpoint{3.045834in}{2.202251in}}%
\pgfpathlineto{\pgfqpoint{3.046577in}{1.532028in}}%
\pgfpathlineto{\pgfqpoint{3.047213in}{1.532065in}}%
\pgfpathlineto{\pgfqpoint{3.048487in}{1.532119in}}%
\pgfpathlineto{\pgfqpoint{3.048498in}{1.532240in}}%
\pgfpathlineto{\pgfqpoint{3.049623in}{2.202707in}}%
\pgfpathlineto{\pgfqpoint{3.049722in}{2.202744in}}%
\pgfpathlineto{\pgfqpoint{3.050464in}{2.202077in}}%
\pgfpathlineto{\pgfqpoint{3.050513in}{2.202562in}}%
\pgfpathlineto{\pgfqpoint{3.050662in}{2.201416in}}%
\pgfpathlineto{\pgfqpoint{3.050711in}{2.202343in}}%
\pgfpathlineto{\pgfqpoint{3.051107in}{2.040738in}}%
\pgfpathlineto{\pgfqpoint{3.051488in}{1.532032in}}%
\pgfpathlineto{\pgfqpoint{3.052033in}{1.532050in}}%
\pgfpathlineto{\pgfqpoint{3.053439in}{1.532155in}}%
\pgfpathlineto{\pgfqpoint{3.053468in}{1.609584in}}%
\pgfpathlineto{\pgfqpoint{3.054240in}{2.202493in}}%
\pgfpathlineto{\pgfqpoint{3.054635in}{2.202735in}}%
\pgfpathlineto{\pgfqpoint{3.054981in}{2.202480in}}%
\pgfpathlineto{\pgfqpoint{3.055031in}{2.202785in}}%
\pgfpathlineto{\pgfqpoint{3.055575in}{2.201498in}}%
\pgfpathlineto{\pgfqpoint{3.055624in}{2.202423in}}%
\pgfpathlineto{\pgfqpoint{3.055674in}{2.199663in}}%
\pgfpathlineto{\pgfqpoint{3.055723in}{2.202285in}}%
\pgfpathlineto{\pgfqpoint{3.056466in}{1.532028in}}%
\pgfpathlineto{\pgfqpoint{3.056953in}{1.532046in}}%
\pgfpathlineto{\pgfqpoint{3.058389in}{1.532222in}}%
\pgfpathlineto{\pgfqpoint{3.059518in}{2.202769in}}%
\pgfpathlineto{\pgfqpoint{3.059567in}{2.202393in}}%
\pgfpathlineto{\pgfqpoint{3.060012in}{2.202850in}}%
\pgfpathlineto{\pgfqpoint{3.060260in}{2.201731in}}%
\pgfpathlineto{\pgfqpoint{3.060309in}{2.202746in}}%
\pgfpathlineto{\pgfqpoint{3.060358in}{2.201167in}}%
\pgfpathlineto{\pgfqpoint{3.060408in}{2.202681in}}%
\pgfpathlineto{\pgfqpoint{3.061373in}{1.532024in}}%
\pgfpathlineto{\pgfqpoint{3.062156in}{1.532072in}}%
\pgfpathlineto{\pgfqpoint{3.063323in}{1.532120in}}%
\pgfpathlineto{\pgfqpoint{3.063335in}{1.532242in}}%
\pgfpathlineto{\pgfqpoint{3.064466in}{2.202628in}}%
\pgfpathlineto{\pgfqpoint{3.065653in}{2.201827in}}%
\pgfpathlineto{\pgfqpoint{3.065850in}{2.199791in}}%
\pgfpathlineto{\pgfqpoint{3.066351in}{1.532028in}}%
\pgfpathlineto{\pgfqpoint{3.067083in}{1.532068in}}%
\pgfpathlineto{\pgfqpoint{3.068268in}{1.532120in}}%
\pgfpathlineto{\pgfqpoint{3.068280in}{1.532237in}}%
\pgfpathlineto{\pgfqpoint{3.069409in}{2.202745in}}%
\pgfpathlineto{\pgfqpoint{3.069606in}{2.202812in}}%
\pgfpathlineto{\pgfqpoint{3.070249in}{2.202035in}}%
\pgfpathlineto{\pgfqpoint{3.070299in}{2.202631in}}%
\pgfpathlineto{\pgfqpoint{3.070447in}{2.201174in}}%
\pgfpathlineto{\pgfqpoint{3.070497in}{2.202427in}}%
\pgfpathlineto{\pgfqpoint{3.071298in}{1.532029in}}%
\pgfpathlineto{\pgfqpoint{3.072230in}{1.532074in}}%
\pgfpathlineto{\pgfqpoint{3.073213in}{1.532117in}}%
\pgfpathlineto{\pgfqpoint{3.073224in}{1.532197in}}%
\pgfpathlineto{\pgfqpoint{3.074456in}{2.202850in}}%
\pgfpathlineto{\pgfqpoint{3.074654in}{2.202898in}}%
\pgfpathlineto{\pgfqpoint{3.075198in}{2.192139in}}%
\pgfpathlineto{\pgfqpoint{3.075247in}{2.202736in}}%
\pgfpathlineto{\pgfqpoint{3.075297in}{2.142228in}}%
\pgfpathlineto{\pgfqpoint{3.075742in}{2.183082in}}%
\pgfpathlineto{\pgfqpoint{3.076263in}{1.532034in}}%
\pgfpathlineto{\pgfqpoint{3.077542in}{1.532082in}}%
\pgfpathlineto{\pgfqpoint{3.078174in}{1.532482in}}%
\pgfpathlineto{\pgfqpoint{3.079234in}{2.202620in}}%
\pgfpathlineto{\pgfqpoint{3.080633in}{2.200590in}}%
\pgfpathlineto{\pgfqpoint{3.080714in}{2.146894in}}%
\pgfpathlineto{\pgfqpoint{3.081155in}{1.532029in}}%
\pgfpathlineto{\pgfqpoint{3.081786in}{1.532051in}}%
\pgfpathlineto{\pgfqpoint{3.083116in}{1.532211in}}%
\pgfpathlineto{\pgfqpoint{3.084247in}{2.202804in}}%
\pgfpathlineto{\pgfqpoint{3.084296in}{2.202297in}}%
\pgfpathlineto{\pgfqpoint{3.084742in}{2.202914in}}%
\pgfpathlineto{\pgfqpoint{3.084890in}{2.201635in}}%
\pgfpathlineto{\pgfqpoint{3.084939in}{2.202873in}}%
\pgfpathlineto{\pgfqpoint{3.084989in}{2.200856in}}%
\pgfpathlineto{\pgfqpoint{3.085038in}{2.202833in}}%
\pgfpathlineto{\pgfqpoint{3.086112in}{1.532031in}}%
\pgfpathlineto{\pgfqpoint{3.086429in}{1.532047in}}%
\pgfpathlineto{\pgfqpoint{3.088058in}{1.532162in}}%
\pgfpathlineto{\pgfqpoint{3.088100in}{1.858276in}}%
\pgfpathlineto{\pgfqpoint{3.088861in}{2.202500in}}%
\pgfpathlineto{\pgfqpoint{3.089257in}{2.202726in}}%
\pgfpathlineto{\pgfqpoint{3.089603in}{2.202460in}}%
\pgfpathlineto{\pgfqpoint{3.089653in}{2.202759in}}%
\pgfpathlineto{\pgfqpoint{3.090197in}{2.201364in}}%
\pgfpathlineto{\pgfqpoint{3.090246in}{2.202358in}}%
\pgfpathlineto{\pgfqpoint{3.090691in}{1.946565in}}%
\pgfpathlineto{\pgfqpoint{3.091046in}{1.532029in}}%
\pgfpathlineto{\pgfqpoint{3.091624in}{1.532061in}}%
\pgfpathlineto{\pgfqpoint{3.092995in}{1.532118in}}%
\pgfpathlineto{\pgfqpoint{3.093006in}{1.532206in}}%
\pgfpathlineto{\pgfqpoint{3.094139in}{2.202776in}}%
\pgfpathlineto{\pgfqpoint{3.094188in}{2.202295in}}%
\pgfpathlineto{\pgfqpoint{3.094633in}{2.202880in}}%
\pgfpathlineto{\pgfqpoint{3.094781in}{2.201701in}}%
\pgfpathlineto{\pgfqpoint{3.094831in}{2.202835in}}%
\pgfpathlineto{\pgfqpoint{3.094880in}{2.201144in}}%
\pgfpathlineto{\pgfqpoint{3.094930in}{2.202792in}}%
\pgfpathlineto{\pgfqpoint{3.096016in}{1.532030in}}%
\pgfpathlineto{\pgfqpoint{3.096644in}{1.532060in}}%
\pgfpathlineto{\pgfqpoint{3.097947in}{1.532147in}}%
\pgfpathlineto{\pgfqpoint{3.097956in}{1.532467in}}%
\pgfpathlineto{\pgfqpoint{3.099013in}{2.202575in}}%
\pgfpathlineto{\pgfqpoint{3.100249in}{2.201349in}}%
\pgfpathlineto{\pgfqpoint{3.100299in}{2.201767in}}%
\pgfpathlineto{\pgfqpoint{3.100447in}{2.135928in}}%
\pgfpathlineto{\pgfqpoint{3.100921in}{1.532018in}}%
\pgfpathlineto{\pgfqpoint{3.101818in}{1.532078in}}%
\pgfpathlineto{\pgfqpoint{3.102885in}{1.532116in}}%
\pgfpathlineto{\pgfqpoint{3.102895in}{1.532175in}}%
\pgfpathlineto{\pgfqpoint{3.102937in}{1.891719in}}%
\pgfpathlineto{\pgfqpoint{3.103699in}{2.202454in}}%
\pgfpathlineto{\pgfqpoint{3.104095in}{2.202693in}}%
\pgfpathlineto{\pgfqpoint{3.104441in}{2.202437in}}%
\pgfpathlineto{\pgfqpoint{3.104491in}{2.202740in}}%
\pgfpathlineto{\pgfqpoint{3.105035in}{2.201440in}}%
\pgfpathlineto{\pgfqpoint{3.105084in}{2.202371in}}%
\pgfpathlineto{\pgfqpoint{3.105134in}{2.199080in}}%
\pgfpathlineto{\pgfqpoint{3.105183in}{2.202230in}}%
\pgfpathlineto{\pgfqpoint{3.105865in}{1.532009in}}%
\pgfpathlineto{\pgfqpoint{3.106414in}{1.532062in}}%
\pgfpathlineto{\pgfqpoint{3.107832in}{1.532119in}}%
\pgfpathlineto{\pgfqpoint{3.107841in}{1.532179in}}%
\pgfpathlineto{\pgfqpoint{3.107882in}{1.881403in}}%
\pgfpathlineto{\pgfqpoint{3.108648in}{2.202481in}}%
\pgfpathlineto{\pgfqpoint{3.109142in}{2.202745in}}%
\pgfpathlineto{\pgfqpoint{3.109390in}{2.202384in}}%
\pgfpathlineto{\pgfqpoint{3.109439in}{2.202759in}}%
\pgfpathlineto{\pgfqpoint{3.109884in}{2.201589in}}%
\pgfpathlineto{\pgfqpoint{3.109934in}{2.202499in}}%
\pgfpathlineto{\pgfqpoint{3.109983in}{2.200776in}}%
\pgfpathlineto{\pgfqpoint{3.110033in}{2.202391in}}%
\pgfpathlineto{\pgfqpoint{3.110857in}{1.532028in}}%
\pgfpathlineto{\pgfqpoint{3.111439in}{1.532049in}}%
\pgfpathlineto{\pgfqpoint{3.112789in}{1.532244in}}%
\pgfpathlineto{\pgfqpoint{3.113909in}{2.202652in}}%
\pgfpathlineto{\pgfqpoint{3.114008in}{2.202690in}}%
\pgfpathlineto{\pgfqpoint{3.114750in}{2.202134in}}%
\pgfpathlineto{\pgfqpoint{3.114800in}{2.202509in}}%
\pgfpathlineto{\pgfqpoint{3.115047in}{2.200970in}}%
\pgfpathlineto{\pgfqpoint{3.115096in}{2.202126in}}%
\pgfpathlineto{\pgfqpoint{3.115773in}{1.532024in}}%
\pgfpathlineto{\pgfqpoint{3.116656in}{1.532076in}}%
\pgfpathlineto{\pgfqpoint{3.117723in}{1.532119in}}%
\pgfpathlineto{\pgfqpoint{3.117735in}{1.532233in}}%
\pgfpathlineto{\pgfqpoint{3.118837in}{2.202620in}}%
\pgfpathlineto{\pgfqpoint{3.120074in}{2.201212in}}%
\pgfpathlineto{\pgfqpoint{3.120123in}{2.201740in}}%
\pgfpathlineto{\pgfqpoint{3.120688in}{1.613519in}}%
\pgfpathlineto{\pgfqpoint{3.120720in}{1.532026in}}%
\pgfpathlineto{\pgfqpoint{3.121459in}{1.532052in}}%
\pgfpathlineto{\pgfqpoint{3.122680in}{1.532224in}}%
\pgfpathlineto{\pgfqpoint{3.123808in}{2.202755in}}%
\pgfpathlineto{\pgfqpoint{3.124105in}{2.202843in}}%
\pgfpathlineto{\pgfqpoint{3.124550in}{2.201915in}}%
\pgfpathlineto{\pgfqpoint{3.124599in}{2.202739in}}%
\pgfpathlineto{\pgfqpoint{3.124748in}{2.200482in}}%
\pgfpathlineto{\pgfqpoint{3.124797in}{2.202593in}}%
\pgfpathlineto{\pgfqpoint{3.125663in}{1.532025in}}%
\pgfpathlineto{\pgfqpoint{3.126099in}{1.532040in}}%
\pgfpathlineto{\pgfqpoint{3.127625in}{1.532223in}}%
\pgfpathlineto{\pgfqpoint{3.128756in}{2.202742in}}%
\pgfpathlineto{\pgfqpoint{3.128954in}{2.202807in}}%
\pgfpathlineto{\pgfqpoint{3.129597in}{2.201796in}}%
\pgfpathlineto{\pgfqpoint{3.129647in}{2.202624in}}%
\pgfpathlineto{\pgfqpoint{3.129696in}{2.201320in}}%
\pgfpathlineto{\pgfqpoint{3.129745in}{2.202534in}}%
\pgfpathlineto{\pgfqpoint{3.130638in}{1.532028in}}%
\pgfpathlineto{\pgfqpoint{3.132012in}{1.532079in}}%
\pgfpathlineto{\pgfqpoint{3.132571in}{1.532233in}}%
\pgfpathlineto{\pgfqpoint{3.133701in}{2.202734in}}%
\pgfpathlineto{\pgfqpoint{3.133899in}{2.202799in}}%
\pgfpathlineto{\pgfqpoint{3.134542in}{2.201897in}}%
\pgfpathlineto{\pgfqpoint{3.134591in}{2.202615in}}%
\pgfpathlineto{\pgfqpoint{3.134740in}{2.200043in}}%
\pgfpathlineto{\pgfqpoint{3.134789in}{2.202409in}}%
\pgfpathlineto{\pgfqpoint{3.135555in}{1.532023in}}%
\pgfpathlineto{\pgfqpoint{3.136044in}{1.532061in}}%
\pgfpathlineto{\pgfqpoint{3.137504in}{1.532117in}}%
\pgfpathlineto{\pgfqpoint{3.137515in}{1.532203in}}%
\pgfpathlineto{\pgfqpoint{3.138746in}{2.202812in}}%
\pgfpathlineto{\pgfqpoint{3.138944in}{2.202859in}}%
\pgfpathlineto{\pgfqpoint{3.139488in}{2.199630in}}%
\pgfpathlineto{\pgfqpoint{3.139537in}{2.202696in}}%
\pgfpathlineto{\pgfqpoint{3.140510in}{1.532035in}}%
\pgfpathlineto{\pgfqpoint{3.140777in}{1.532042in}}%
\pgfpathlineto{\pgfqpoint{3.142465in}{1.532489in}}%
\pgfpathlineto{\pgfqpoint{3.143427in}{2.202534in}}%
\pgfpathlineto{\pgfqpoint{3.143477in}{2.202489in}}%
\pgfpathlineto{\pgfqpoint{3.144218in}{2.202595in}}%
\pgfpathlineto{\pgfqpoint{3.144317in}{2.202535in}}%
\pgfpathlineto{\pgfqpoint{3.144762in}{2.201327in}}%
\pgfpathlineto{\pgfqpoint{3.144812in}{2.201820in}}%
\pgfpathlineto{\pgfqpoint{3.145059in}{2.016780in}}%
\pgfpathlineto{\pgfqpoint{3.145429in}{1.532020in}}%
\pgfpathlineto{\pgfqpoint{3.146023in}{1.532048in}}%
\pgfpathlineto{\pgfqpoint{3.147405in}{1.532182in}}%
\pgfpathlineto{\pgfqpoint{3.148900in}{2.202732in}}%
\pgfpathlineto{\pgfqpoint{3.149641in}{2.201012in}}%
\pgfpathlineto{\pgfqpoint{3.149691in}{2.202176in}}%
\pgfpathlineto{\pgfqpoint{3.150430in}{1.532028in}}%
\pgfpathlineto{\pgfqpoint{3.151313in}{1.532061in}}%
\pgfpathlineto{\pgfqpoint{3.152352in}{1.532225in}}%
\pgfpathlineto{\pgfqpoint{3.153481in}{2.202734in}}%
\pgfpathlineto{\pgfqpoint{3.153530in}{2.202386in}}%
\pgfpathlineto{\pgfqpoint{3.153976in}{2.202820in}}%
\pgfpathlineto{\pgfqpoint{3.154223in}{2.201825in}}%
\pgfpathlineto{\pgfqpoint{3.154272in}{2.202717in}}%
\pgfpathlineto{\pgfqpoint{3.154421in}{2.194068in}}%
\pgfpathlineto{\pgfqpoint{3.154470in}{2.202568in}}%
\pgfpathlineto{\pgfqpoint{3.155347in}{1.532031in}}%
\pgfpathlineto{\pgfqpoint{3.155763in}{1.532053in}}%
\pgfpathlineto{\pgfqpoint{3.157295in}{1.532170in}}%
\pgfpathlineto{\pgfqpoint{3.157337in}{1.886143in}}%
\pgfpathlineto{\pgfqpoint{3.158099in}{2.202479in}}%
\pgfpathlineto{\pgfqpoint{3.158494in}{2.202713in}}%
\pgfpathlineto{\pgfqpoint{3.158840in}{2.202459in}}%
\pgfpathlineto{\pgfqpoint{3.158890in}{2.202755in}}%
\pgfpathlineto{\pgfqpoint{3.159434in}{2.201453in}}%
\pgfpathlineto{\pgfqpoint{3.159483in}{2.202373in}}%
\pgfpathlineto{\pgfqpoint{3.159533in}{2.198857in}}%
\pgfpathlineto{\pgfqpoint{3.159582in}{2.202227in}}%
\pgfpathlineto{\pgfqpoint{3.160284in}{1.532023in}}%
\pgfpathlineto{\pgfqpoint{3.160826in}{1.532042in}}%
\pgfpathlineto{\pgfqpoint{3.162242in}{1.532192in}}%
\pgfpathlineto{\pgfqpoint{3.163467in}{2.202766in}}%
\pgfpathlineto{\pgfqpoint{3.163516in}{2.202383in}}%
\pgfpathlineto{\pgfqpoint{3.163862in}{2.202810in}}%
\pgfpathlineto{\pgfqpoint{3.164110in}{2.201829in}}%
\pgfpathlineto{\pgfqpoint{3.164159in}{2.202707in}}%
\pgfpathlineto{\pgfqpoint{3.164308in}{2.197208in}}%
\pgfpathlineto{\pgfqpoint{3.164357in}{2.202555in}}%
\pgfpathlineto{\pgfqpoint{3.165239in}{1.532035in}}%
\pgfpathlineto{\pgfqpoint{3.165607in}{1.532046in}}%
\pgfpathlineto{\pgfqpoint{3.167192in}{1.532490in}}%
\pgfpathlineto{\pgfqpoint{3.168154in}{2.202521in}}%
\pgfpathlineto{\pgfqpoint{3.168204in}{2.202477in}}%
\pgfpathlineto{\pgfqpoint{3.168945in}{2.202588in}}%
\pgfpathlineto{\pgfqpoint{3.169044in}{2.202529in}}%
\pgfpathlineto{\pgfqpoint{3.169489in}{2.201356in}}%
\pgfpathlineto{\pgfqpoint{3.169539in}{2.201824in}}%
\pgfpathlineto{\pgfqpoint{3.169737in}{2.104057in}}%
\pgfpathlineto{\pgfqpoint{3.170033in}{1.681539in}}%
\pgfpathlineto{\pgfqpoint{3.170132in}{1.748829in}}%
\pgfusepath{stroke}%
\end{pgfscope}%
\begin{pgfscope}%
\pgfsetrectcap%
\pgfsetmiterjoin%
\pgfsetlinewidth{0.803000pt}%
\definecolor{currentstroke}{rgb}{0.000000,0.000000,0.000000}%
\pgfsetstrokecolor{currentstroke}%
\pgfsetdash{}{0pt}%
\pgfpathmoveto{\pgfqpoint{0.573768in}{1.532029in}}%
\pgfpathlineto{\pgfqpoint{0.573768in}{2.438279in}}%
\pgfusepath{stroke}%
\end{pgfscope}%
\begin{pgfscope}%
\pgfsetrectcap%
\pgfsetmiterjoin%
\pgfsetlinewidth{0.803000pt}%
\definecolor{currentstroke}{rgb}{0.000000,0.000000,0.000000}%
\pgfsetstrokecolor{currentstroke}%
\pgfsetdash{}{0pt}%
\pgfpathmoveto{\pgfqpoint{3.293768in}{1.532029in}}%
\pgfpathlineto{\pgfqpoint{3.293768in}{2.438279in}}%
\pgfusepath{stroke}%
\end{pgfscope}%
\begin{pgfscope}%
\pgfsetrectcap%
\pgfsetmiterjoin%
\pgfsetlinewidth{0.803000pt}%
\definecolor{currentstroke}{rgb}{0.000000,0.000000,0.000000}%
\pgfsetstrokecolor{currentstroke}%
\pgfsetdash{}{0pt}%
\pgfpathmoveto{\pgfqpoint{0.573768in}{1.532029in}}%
\pgfpathlineto{\pgfqpoint{3.293768in}{1.532029in}}%
\pgfusepath{stroke}%
\end{pgfscope}%
\begin{pgfscope}%
\pgfsetrectcap%
\pgfsetmiterjoin%
\pgfsetlinewidth{0.803000pt}%
\definecolor{currentstroke}{rgb}{0.000000,0.000000,0.000000}%
\pgfsetstrokecolor{currentstroke}%
\pgfsetdash{}{0pt}%
\pgfpathmoveto{\pgfqpoint{0.573768in}{2.438279in}}%
\pgfpathlineto{\pgfqpoint{3.293768in}{2.438279in}}%
\pgfusepath{stroke}%
\end{pgfscope}%
\begin{pgfscope}%
\definecolor{textcolor}{rgb}{0.000000,0.000000,0.000000}%
\pgfsetstrokecolor{textcolor}%
\pgfsetfillcolor{textcolor}%
\pgftext[x=0.628168in,y=2.374842in,left,top]{\color{textcolor}{\rmfamily\fontsize{8.000000}{9.600000}\selectfont\catcode`\^=\active\def^{\ifmmode\sp\else\^{}\fi}\catcode`\%=\active\def%{\%}(b)}}%
\end{pgfscope}%
\begin{pgfscope}%
\pgfsetbuttcap%
\pgfsetmiterjoin%
\definecolor{currentfill}{rgb}{1.000000,1.000000,1.000000}%
\pgfsetfillcolor{currentfill}%
\pgfsetlinewidth{0.000000pt}%
\definecolor{currentstroke}{rgb}{0.000000,0.000000,0.000000}%
\pgfsetstrokecolor{currentstroke}%
\pgfsetstrokeopacity{0.000000}%
\pgfsetdash{}{0pt}%
\pgfpathmoveto{\pgfqpoint{0.573768in}{0.462654in}}%
\pgfpathlineto{\pgfqpoint{3.293768in}{0.462654in}}%
\pgfpathlineto{\pgfqpoint{3.293768in}{1.368904in}}%
\pgfpathlineto{\pgfqpoint{0.573768in}{1.368904in}}%
\pgfpathlineto{\pgfqpoint{0.573768in}{0.462654in}}%
\pgfpathclose%
\pgfusepath{fill}%
\end{pgfscope}%
\begin{pgfscope}%
\pgfpathrectangle{\pgfqpoint{0.573768in}{0.462654in}}{\pgfqpoint{2.720000in}{0.906250in}}%
\pgfusepath{clip}%
\pgfsetbuttcap%
\pgfsetroundjoin%
\pgfsetlinewidth{0.501875pt}%
\definecolor{currentstroke}{rgb}{0.690196,0.690196,0.690196}%
\pgfsetstrokecolor{currentstroke}%
\pgfsetstrokeopacity{0.250000}%
\pgfsetdash{{1.850000pt}{0.800000pt}}{0.000000pt}%
\pgfpathmoveto{\pgfqpoint{0.697405in}{0.462654in}}%
\pgfpathlineto{\pgfqpoint{0.697405in}{1.368904in}}%
\pgfusepath{stroke}%
\end{pgfscope}%
\begin{pgfscope}%
\pgfsetbuttcap%
\pgfsetroundjoin%
\definecolor{currentfill}{rgb}{0.000000,0.000000,0.000000}%
\pgfsetfillcolor{currentfill}%
\pgfsetlinewidth{0.803000pt}%
\definecolor{currentstroke}{rgb}{0.000000,0.000000,0.000000}%
\pgfsetstrokecolor{currentstroke}%
\pgfsetdash{}{0pt}%
\pgfsys@defobject{currentmarker}{\pgfqpoint{0.000000in}{-0.048611in}}{\pgfqpoint{0.000000in}{0.000000in}}{%
\pgfpathmoveto{\pgfqpoint{0.000000in}{0.000000in}}%
\pgfpathlineto{\pgfqpoint{0.000000in}{-0.048611in}}%
\pgfusepath{stroke,fill}%
}%
\begin{pgfscope}%
\pgfsys@transformshift{0.697405in}{0.462654in}%
\pgfsys@useobject{currentmarker}{}%
\end{pgfscope}%
\end{pgfscope}%
\begin{pgfscope}%
\definecolor{textcolor}{rgb}{0.000000,0.000000,0.000000}%
\pgfsetstrokecolor{textcolor}%
\pgfsetfillcolor{textcolor}%
\pgftext[x=0.697405in,y=0.365432in,,top]{\color{textcolor}{\rmfamily\fontsize{8.000000}{9.600000}\selectfont\catcode`\^=\active\def^{\ifmmode\sp\else\^{}\fi}\catcode`\%=\active\def%{\%}$\mathdefault{0}$}}%
\end{pgfscope}%
\begin{pgfscope}%
\pgfpathrectangle{\pgfqpoint{0.573768in}{0.462654in}}{\pgfqpoint{2.720000in}{0.906250in}}%
\pgfusepath{clip}%
\pgfsetbuttcap%
\pgfsetroundjoin%
\pgfsetlinewidth{0.501875pt}%
\definecolor{currentstroke}{rgb}{0.690196,0.690196,0.690196}%
\pgfsetstrokecolor{currentstroke}%
\pgfsetstrokeopacity{0.250000}%
\pgfsetdash{{1.850000pt}{0.800000pt}}{0.000000pt}%
\pgfpathmoveto{\pgfqpoint{1.191950in}{0.462654in}}%
\pgfpathlineto{\pgfqpoint{1.191950in}{1.368904in}}%
\pgfusepath{stroke}%
\end{pgfscope}%
\begin{pgfscope}%
\pgfsetbuttcap%
\pgfsetroundjoin%
\definecolor{currentfill}{rgb}{0.000000,0.000000,0.000000}%
\pgfsetfillcolor{currentfill}%
\pgfsetlinewidth{0.803000pt}%
\definecolor{currentstroke}{rgb}{0.000000,0.000000,0.000000}%
\pgfsetstrokecolor{currentstroke}%
\pgfsetdash{}{0pt}%
\pgfsys@defobject{currentmarker}{\pgfqpoint{0.000000in}{-0.048611in}}{\pgfqpoint{0.000000in}{0.000000in}}{%
\pgfpathmoveto{\pgfqpoint{0.000000in}{0.000000in}}%
\pgfpathlineto{\pgfqpoint{0.000000in}{-0.048611in}}%
\pgfusepath{stroke,fill}%
}%
\begin{pgfscope}%
\pgfsys@transformshift{1.191950in}{0.462654in}%
\pgfsys@useobject{currentmarker}{}%
\end{pgfscope}%
\end{pgfscope}%
\begin{pgfscope}%
\definecolor{textcolor}{rgb}{0.000000,0.000000,0.000000}%
\pgfsetstrokecolor{textcolor}%
\pgfsetfillcolor{textcolor}%
\pgftext[x=1.191950in,y=0.365432in,,top]{\color{textcolor}{\rmfamily\fontsize{8.000000}{9.600000}\selectfont\catcode`\^=\active\def^{\ifmmode\sp\else\^{}\fi}\catcode`\%=\active\def%{\%}$\mathdefault{1}$}}%
\end{pgfscope}%
\begin{pgfscope}%
\pgfpathrectangle{\pgfqpoint{0.573768in}{0.462654in}}{\pgfqpoint{2.720000in}{0.906250in}}%
\pgfusepath{clip}%
\pgfsetbuttcap%
\pgfsetroundjoin%
\pgfsetlinewidth{0.501875pt}%
\definecolor{currentstroke}{rgb}{0.690196,0.690196,0.690196}%
\pgfsetstrokecolor{currentstroke}%
\pgfsetstrokeopacity{0.250000}%
\pgfsetdash{{1.850000pt}{0.800000pt}}{0.000000pt}%
\pgfpathmoveto{\pgfqpoint{1.686496in}{0.462654in}}%
\pgfpathlineto{\pgfqpoint{1.686496in}{1.368904in}}%
\pgfusepath{stroke}%
\end{pgfscope}%
\begin{pgfscope}%
\pgfsetbuttcap%
\pgfsetroundjoin%
\definecolor{currentfill}{rgb}{0.000000,0.000000,0.000000}%
\pgfsetfillcolor{currentfill}%
\pgfsetlinewidth{0.803000pt}%
\definecolor{currentstroke}{rgb}{0.000000,0.000000,0.000000}%
\pgfsetstrokecolor{currentstroke}%
\pgfsetdash{}{0pt}%
\pgfsys@defobject{currentmarker}{\pgfqpoint{0.000000in}{-0.048611in}}{\pgfqpoint{0.000000in}{0.000000in}}{%
\pgfpathmoveto{\pgfqpoint{0.000000in}{0.000000in}}%
\pgfpathlineto{\pgfqpoint{0.000000in}{-0.048611in}}%
\pgfusepath{stroke,fill}%
}%
\begin{pgfscope}%
\pgfsys@transformshift{1.686496in}{0.462654in}%
\pgfsys@useobject{currentmarker}{}%
\end{pgfscope}%
\end{pgfscope}%
\begin{pgfscope}%
\definecolor{textcolor}{rgb}{0.000000,0.000000,0.000000}%
\pgfsetstrokecolor{textcolor}%
\pgfsetfillcolor{textcolor}%
\pgftext[x=1.686496in,y=0.365432in,,top]{\color{textcolor}{\rmfamily\fontsize{8.000000}{9.600000}\selectfont\catcode`\^=\active\def^{\ifmmode\sp\else\^{}\fi}\catcode`\%=\active\def%{\%}$\mathdefault{2}$}}%
\end{pgfscope}%
\begin{pgfscope}%
\pgfpathrectangle{\pgfqpoint{0.573768in}{0.462654in}}{\pgfqpoint{2.720000in}{0.906250in}}%
\pgfusepath{clip}%
\pgfsetbuttcap%
\pgfsetroundjoin%
\pgfsetlinewidth{0.501875pt}%
\definecolor{currentstroke}{rgb}{0.690196,0.690196,0.690196}%
\pgfsetstrokecolor{currentstroke}%
\pgfsetstrokeopacity{0.250000}%
\pgfsetdash{{1.850000pt}{0.800000pt}}{0.000000pt}%
\pgfpathmoveto{\pgfqpoint{2.181041in}{0.462654in}}%
\pgfpathlineto{\pgfqpoint{2.181041in}{1.368904in}}%
\pgfusepath{stroke}%
\end{pgfscope}%
\begin{pgfscope}%
\pgfsetbuttcap%
\pgfsetroundjoin%
\definecolor{currentfill}{rgb}{0.000000,0.000000,0.000000}%
\pgfsetfillcolor{currentfill}%
\pgfsetlinewidth{0.803000pt}%
\definecolor{currentstroke}{rgb}{0.000000,0.000000,0.000000}%
\pgfsetstrokecolor{currentstroke}%
\pgfsetdash{}{0pt}%
\pgfsys@defobject{currentmarker}{\pgfqpoint{0.000000in}{-0.048611in}}{\pgfqpoint{0.000000in}{0.000000in}}{%
\pgfpathmoveto{\pgfqpoint{0.000000in}{0.000000in}}%
\pgfpathlineto{\pgfqpoint{0.000000in}{-0.048611in}}%
\pgfusepath{stroke,fill}%
}%
\begin{pgfscope}%
\pgfsys@transformshift{2.181041in}{0.462654in}%
\pgfsys@useobject{currentmarker}{}%
\end{pgfscope}%
\end{pgfscope}%
\begin{pgfscope}%
\definecolor{textcolor}{rgb}{0.000000,0.000000,0.000000}%
\pgfsetstrokecolor{textcolor}%
\pgfsetfillcolor{textcolor}%
\pgftext[x=2.181041in,y=0.365432in,,top]{\color{textcolor}{\rmfamily\fontsize{8.000000}{9.600000}\selectfont\catcode`\^=\active\def^{\ifmmode\sp\else\^{}\fi}\catcode`\%=\active\def%{\%}$\mathdefault{3}$}}%
\end{pgfscope}%
\begin{pgfscope}%
\pgfpathrectangle{\pgfqpoint{0.573768in}{0.462654in}}{\pgfqpoint{2.720000in}{0.906250in}}%
\pgfusepath{clip}%
\pgfsetbuttcap%
\pgfsetroundjoin%
\pgfsetlinewidth{0.501875pt}%
\definecolor{currentstroke}{rgb}{0.690196,0.690196,0.690196}%
\pgfsetstrokecolor{currentstroke}%
\pgfsetstrokeopacity{0.250000}%
\pgfsetdash{{1.850000pt}{0.800000pt}}{0.000000pt}%
\pgfpathmoveto{\pgfqpoint{2.675587in}{0.462654in}}%
\pgfpathlineto{\pgfqpoint{2.675587in}{1.368904in}}%
\pgfusepath{stroke}%
\end{pgfscope}%
\begin{pgfscope}%
\pgfsetbuttcap%
\pgfsetroundjoin%
\definecolor{currentfill}{rgb}{0.000000,0.000000,0.000000}%
\pgfsetfillcolor{currentfill}%
\pgfsetlinewidth{0.803000pt}%
\definecolor{currentstroke}{rgb}{0.000000,0.000000,0.000000}%
\pgfsetstrokecolor{currentstroke}%
\pgfsetdash{}{0pt}%
\pgfsys@defobject{currentmarker}{\pgfqpoint{0.000000in}{-0.048611in}}{\pgfqpoint{0.000000in}{0.000000in}}{%
\pgfpathmoveto{\pgfqpoint{0.000000in}{0.000000in}}%
\pgfpathlineto{\pgfqpoint{0.000000in}{-0.048611in}}%
\pgfusepath{stroke,fill}%
}%
\begin{pgfscope}%
\pgfsys@transformshift{2.675587in}{0.462654in}%
\pgfsys@useobject{currentmarker}{}%
\end{pgfscope}%
\end{pgfscope}%
\begin{pgfscope}%
\definecolor{textcolor}{rgb}{0.000000,0.000000,0.000000}%
\pgfsetstrokecolor{textcolor}%
\pgfsetfillcolor{textcolor}%
\pgftext[x=2.675587in,y=0.365432in,,top]{\color{textcolor}{\rmfamily\fontsize{8.000000}{9.600000}\selectfont\catcode`\^=\active\def^{\ifmmode\sp\else\^{}\fi}\catcode`\%=\active\def%{\%}$\mathdefault{4}$}}%
\end{pgfscope}%
\begin{pgfscope}%
\pgfpathrectangle{\pgfqpoint{0.573768in}{0.462654in}}{\pgfqpoint{2.720000in}{0.906250in}}%
\pgfusepath{clip}%
\pgfsetbuttcap%
\pgfsetroundjoin%
\pgfsetlinewidth{0.501875pt}%
\definecolor{currentstroke}{rgb}{0.690196,0.690196,0.690196}%
\pgfsetstrokecolor{currentstroke}%
\pgfsetstrokeopacity{0.250000}%
\pgfsetdash{{1.850000pt}{0.800000pt}}{0.000000pt}%
\pgfpathmoveto{\pgfqpoint{3.170132in}{0.462654in}}%
\pgfpathlineto{\pgfqpoint{3.170132in}{1.368904in}}%
\pgfusepath{stroke}%
\end{pgfscope}%
\begin{pgfscope}%
\pgfsetbuttcap%
\pgfsetroundjoin%
\definecolor{currentfill}{rgb}{0.000000,0.000000,0.000000}%
\pgfsetfillcolor{currentfill}%
\pgfsetlinewidth{0.803000pt}%
\definecolor{currentstroke}{rgb}{0.000000,0.000000,0.000000}%
\pgfsetstrokecolor{currentstroke}%
\pgfsetdash{}{0pt}%
\pgfsys@defobject{currentmarker}{\pgfqpoint{0.000000in}{-0.048611in}}{\pgfqpoint{0.000000in}{0.000000in}}{%
\pgfpathmoveto{\pgfqpoint{0.000000in}{0.000000in}}%
\pgfpathlineto{\pgfqpoint{0.000000in}{-0.048611in}}%
\pgfusepath{stroke,fill}%
}%
\begin{pgfscope}%
\pgfsys@transformshift{3.170132in}{0.462654in}%
\pgfsys@useobject{currentmarker}{}%
\end{pgfscope}%
\end{pgfscope}%
\begin{pgfscope}%
\definecolor{textcolor}{rgb}{0.000000,0.000000,0.000000}%
\pgfsetstrokecolor{textcolor}%
\pgfsetfillcolor{textcolor}%
\pgftext[x=3.170132in,y=0.365432in,,top]{\color{textcolor}{\rmfamily\fontsize{8.000000}{9.600000}\selectfont\catcode`\^=\active\def^{\ifmmode\sp\else\^{}\fi}\catcode`\%=\active\def%{\%}$\mathdefault{5}$}}%
\end{pgfscope}%
\begin{pgfscope}%
\definecolor{textcolor}{rgb}{0.000000,0.000000,0.000000}%
\pgfsetstrokecolor{textcolor}%
\pgfsetfillcolor{textcolor}%
\pgftext[x=1.933768in,y=0.211111in,,top]{\color{textcolor}{\rmfamily\fontsize{8.000000}{9.600000}\selectfont\catcode`\^=\active\def^{\ifmmode\sp\else\^{}\fi}\catcode`\%=\active\def%{\%}Time (ms)}}%
\end{pgfscope}%
\begin{pgfscope}%
\pgfpathrectangle{\pgfqpoint{0.573768in}{0.462654in}}{\pgfqpoint{2.720000in}{0.906250in}}%
\pgfusepath{clip}%
\pgfsetbuttcap%
\pgfsetroundjoin%
\pgfsetlinewidth{0.501875pt}%
\definecolor{currentstroke}{rgb}{0.690196,0.690196,0.690196}%
\pgfsetstrokecolor{currentstroke}%
\pgfsetstrokeopacity{0.250000}%
\pgfsetdash{{1.850000pt}{0.800000pt}}{0.000000pt}%
\pgfpathmoveto{\pgfqpoint{0.573768in}{0.462654in}}%
\pgfpathlineto{\pgfqpoint{3.293768in}{0.462654in}}%
\pgfusepath{stroke}%
\end{pgfscope}%
\begin{pgfscope}%
\pgfsetbuttcap%
\pgfsetroundjoin%
\definecolor{currentfill}{rgb}{0.000000,0.000000,0.000000}%
\pgfsetfillcolor{currentfill}%
\pgfsetlinewidth{0.803000pt}%
\definecolor{currentstroke}{rgb}{0.000000,0.000000,0.000000}%
\pgfsetstrokecolor{currentstroke}%
\pgfsetdash{}{0pt}%
\pgfsys@defobject{currentmarker}{\pgfqpoint{-0.048611in}{0.000000in}}{\pgfqpoint{-0.000000in}{0.000000in}}{%
\pgfpathmoveto{\pgfqpoint{-0.000000in}{0.000000in}}%
\pgfpathlineto{\pgfqpoint{-0.048611in}{0.000000in}}%
\pgfusepath{stroke,fill}%
}%
\begin{pgfscope}%
\pgfsys@transformshift{0.573768in}{0.462654in}%
\pgfsys@useobject{currentmarker}{}%
\end{pgfscope}%
\end{pgfscope}%
\begin{pgfscope}%
\definecolor{textcolor}{rgb}{0.000000,0.000000,0.000000}%
\pgfsetstrokecolor{textcolor}%
\pgfsetfillcolor{textcolor}%
\pgftext[x=0.325695in, y=0.424074in, left, base]{\color{textcolor}{\rmfamily\fontsize{8.000000}{9.600000}\selectfont\catcode`\^=\active\def^{\ifmmode\sp\else\^{}\fi}\catcode`\%=\active\def%{\%}0.0}}%
\end{pgfscope}%
\begin{pgfscope}%
\pgfpathrectangle{\pgfqpoint{0.573768in}{0.462654in}}{\pgfqpoint{2.720000in}{0.906250in}}%
\pgfusepath{clip}%
\pgfsetbuttcap%
\pgfsetroundjoin%
\pgfsetlinewidth{0.501875pt}%
\definecolor{currentstroke}{rgb}{0.690196,0.690196,0.690196}%
\pgfsetstrokecolor{currentstroke}%
\pgfsetstrokeopacity{0.250000}%
\pgfsetdash{{1.850000pt}{0.800000pt}}{0.000000pt}%
\pgfpathmoveto{\pgfqpoint{0.573768in}{0.689217in}}%
\pgfpathlineto{\pgfqpoint{3.293768in}{0.689217in}}%
\pgfusepath{stroke}%
\end{pgfscope}%
\begin{pgfscope}%
\pgfsetbuttcap%
\pgfsetroundjoin%
\definecolor{currentfill}{rgb}{0.000000,0.000000,0.000000}%
\pgfsetfillcolor{currentfill}%
\pgfsetlinewidth{0.803000pt}%
\definecolor{currentstroke}{rgb}{0.000000,0.000000,0.000000}%
\pgfsetstrokecolor{currentstroke}%
\pgfsetdash{}{0pt}%
\pgfsys@defobject{currentmarker}{\pgfqpoint{-0.048611in}{0.000000in}}{\pgfqpoint{-0.000000in}{0.000000in}}{%
\pgfpathmoveto{\pgfqpoint{-0.000000in}{0.000000in}}%
\pgfpathlineto{\pgfqpoint{-0.048611in}{0.000000in}}%
\pgfusepath{stroke,fill}%
}%
\begin{pgfscope}%
\pgfsys@transformshift{0.573768in}{0.689217in}%
\pgfsys@useobject{currentmarker}{}%
\end{pgfscope}%
\end{pgfscope}%
\begin{pgfscope}%
\definecolor{textcolor}{rgb}{0.000000,0.000000,0.000000}%
\pgfsetstrokecolor{textcolor}%
\pgfsetfillcolor{textcolor}%
\pgftext[x=0.325695in, y=0.650637in, left, base]{\color{textcolor}{\rmfamily\fontsize{8.000000}{9.600000}\selectfont\catcode`\^=\active\def^{\ifmmode\sp\else\^{}\fi}\catcode`\%=\active\def%{\%}6.2}}%
\end{pgfscope}%
\begin{pgfscope}%
\pgfpathrectangle{\pgfqpoint{0.573768in}{0.462654in}}{\pgfqpoint{2.720000in}{0.906250in}}%
\pgfusepath{clip}%
\pgfsetbuttcap%
\pgfsetroundjoin%
\pgfsetlinewidth{0.501875pt}%
\definecolor{currentstroke}{rgb}{0.690196,0.690196,0.690196}%
\pgfsetstrokecolor{currentstroke}%
\pgfsetstrokeopacity{0.250000}%
\pgfsetdash{{1.850000pt}{0.800000pt}}{0.000000pt}%
\pgfpathmoveto{\pgfqpoint{0.573768in}{0.915779in}}%
\pgfpathlineto{\pgfqpoint{3.293768in}{0.915779in}}%
\pgfusepath{stroke}%
\end{pgfscope}%
\begin{pgfscope}%
\pgfsetbuttcap%
\pgfsetroundjoin%
\definecolor{currentfill}{rgb}{0.000000,0.000000,0.000000}%
\pgfsetfillcolor{currentfill}%
\pgfsetlinewidth{0.803000pt}%
\definecolor{currentstroke}{rgb}{0.000000,0.000000,0.000000}%
\pgfsetstrokecolor{currentstroke}%
\pgfsetdash{}{0pt}%
\pgfsys@defobject{currentmarker}{\pgfqpoint{-0.048611in}{0.000000in}}{\pgfqpoint{-0.000000in}{0.000000in}}{%
\pgfpathmoveto{\pgfqpoint{-0.000000in}{0.000000in}}%
\pgfpathlineto{\pgfqpoint{-0.048611in}{0.000000in}}%
\pgfusepath{stroke,fill}%
}%
\begin{pgfscope}%
\pgfsys@transformshift{0.573768in}{0.915779in}%
\pgfsys@useobject{currentmarker}{}%
\end{pgfscope}%
\end{pgfscope}%
\begin{pgfscope}%
\definecolor{textcolor}{rgb}{0.000000,0.000000,0.000000}%
\pgfsetstrokecolor{textcolor}%
\pgfsetfillcolor{textcolor}%
\pgftext[x=0.266667in, y=0.877199in, left, base]{\color{textcolor}{\rmfamily\fontsize{8.000000}{9.600000}\selectfont\catcode`\^=\active\def^{\ifmmode\sp\else\^{}\fi}\catcode`\%=\active\def%{\%}12.5}}%
\end{pgfscope}%
\begin{pgfscope}%
\pgfpathrectangle{\pgfqpoint{0.573768in}{0.462654in}}{\pgfqpoint{2.720000in}{0.906250in}}%
\pgfusepath{clip}%
\pgfsetbuttcap%
\pgfsetroundjoin%
\pgfsetlinewidth{0.501875pt}%
\definecolor{currentstroke}{rgb}{0.690196,0.690196,0.690196}%
\pgfsetstrokecolor{currentstroke}%
\pgfsetstrokeopacity{0.250000}%
\pgfsetdash{{1.850000pt}{0.800000pt}}{0.000000pt}%
\pgfpathmoveto{\pgfqpoint{0.573768in}{1.142342in}}%
\pgfpathlineto{\pgfqpoint{3.293768in}{1.142342in}}%
\pgfusepath{stroke}%
\end{pgfscope}%
\begin{pgfscope}%
\pgfsetbuttcap%
\pgfsetroundjoin%
\definecolor{currentfill}{rgb}{0.000000,0.000000,0.000000}%
\pgfsetfillcolor{currentfill}%
\pgfsetlinewidth{0.803000pt}%
\definecolor{currentstroke}{rgb}{0.000000,0.000000,0.000000}%
\pgfsetstrokecolor{currentstroke}%
\pgfsetdash{}{0pt}%
\pgfsys@defobject{currentmarker}{\pgfqpoint{-0.048611in}{0.000000in}}{\pgfqpoint{-0.000000in}{0.000000in}}{%
\pgfpathmoveto{\pgfqpoint{-0.000000in}{0.000000in}}%
\pgfpathlineto{\pgfqpoint{-0.048611in}{0.000000in}}%
\pgfusepath{stroke,fill}%
}%
\begin{pgfscope}%
\pgfsys@transformshift{0.573768in}{1.142342in}%
\pgfsys@useobject{currentmarker}{}%
\end{pgfscope}%
\end{pgfscope}%
\begin{pgfscope}%
\definecolor{textcolor}{rgb}{0.000000,0.000000,0.000000}%
\pgfsetstrokecolor{textcolor}%
\pgfsetfillcolor{textcolor}%
\pgftext[x=0.266667in, y=1.103762in, left, base]{\color{textcolor}{\rmfamily\fontsize{8.000000}{9.600000}\selectfont\catcode`\^=\active\def^{\ifmmode\sp\else\^{}\fi}\catcode`\%=\active\def%{\%}18.8}}%
\end{pgfscope}%
\begin{pgfscope}%
\pgfpathrectangle{\pgfqpoint{0.573768in}{0.462654in}}{\pgfqpoint{2.720000in}{0.906250in}}%
\pgfusepath{clip}%
\pgfsetbuttcap%
\pgfsetroundjoin%
\pgfsetlinewidth{0.501875pt}%
\definecolor{currentstroke}{rgb}{0.690196,0.690196,0.690196}%
\pgfsetstrokecolor{currentstroke}%
\pgfsetstrokeopacity{0.250000}%
\pgfsetdash{{1.850000pt}{0.800000pt}}{0.000000pt}%
\pgfpathmoveto{\pgfqpoint{0.573768in}{1.368904in}}%
\pgfpathlineto{\pgfqpoint{3.293768in}{1.368904in}}%
\pgfusepath{stroke}%
\end{pgfscope}%
\begin{pgfscope}%
\pgfsetbuttcap%
\pgfsetroundjoin%
\definecolor{currentfill}{rgb}{0.000000,0.000000,0.000000}%
\pgfsetfillcolor{currentfill}%
\pgfsetlinewidth{0.803000pt}%
\definecolor{currentstroke}{rgb}{0.000000,0.000000,0.000000}%
\pgfsetstrokecolor{currentstroke}%
\pgfsetdash{}{0pt}%
\pgfsys@defobject{currentmarker}{\pgfqpoint{-0.048611in}{0.000000in}}{\pgfqpoint{-0.000000in}{0.000000in}}{%
\pgfpathmoveto{\pgfqpoint{-0.000000in}{0.000000in}}%
\pgfpathlineto{\pgfqpoint{-0.048611in}{0.000000in}}%
\pgfusepath{stroke,fill}%
}%
\begin{pgfscope}%
\pgfsys@transformshift{0.573768in}{1.368904in}%
\pgfsys@useobject{currentmarker}{}%
\end{pgfscope}%
\end{pgfscope}%
\begin{pgfscope}%
\definecolor{textcolor}{rgb}{0.000000,0.000000,0.000000}%
\pgfsetstrokecolor{textcolor}%
\pgfsetfillcolor{textcolor}%
\pgftext[x=0.266667in, y=1.330324in, left, base]{\color{textcolor}{\rmfamily\fontsize{8.000000}{9.600000}\selectfont\catcode`\^=\active\def^{\ifmmode\sp\else\^{}\fi}\catcode`\%=\active\def%{\%}25.0}}%
\end{pgfscope}%
\begin{pgfscope}%
\definecolor{textcolor}{rgb}{0.000000,0.000000,0.000000}%
\pgfsetstrokecolor{textcolor}%
\pgfsetfillcolor{textcolor}%
\pgftext[x=0.211111in,y=0.915779in,,bottom,rotate=90.000000]{\color{textcolor}{\rmfamily\fontsize{8.000000}{9.600000}\selectfont\catcode`\^=\active\def^{\ifmmode\sp\else\^{}\fi}\catcode`\%=\active\def%{\%}$v_o$ (V)}}%
\end{pgfscope}%
\begin{pgfscope}%
\pgfpathrectangle{\pgfqpoint{0.573768in}{0.462654in}}{\pgfqpoint{2.720000in}{0.906250in}}%
\pgfusepath{clip}%
\pgfsetrectcap%
\pgfsetroundjoin%
\pgfsetlinewidth{1.505625pt}%
\definecolor{currentstroke}{rgb}{0.000000,0.619608,0.450980}%
\pgfsetstrokecolor{currentstroke}%
\pgfsetdash{}{0pt}%
\pgfpathmoveto{\pgfqpoint{0.697405in}{0.462654in}}%
\pgfpathlineto{\pgfqpoint{0.699607in}{0.463359in}}%
\pgfpathlineto{\pgfqpoint{0.703040in}{0.474979in}}%
\pgfpathlineto{\pgfqpoint{0.704479in}{0.475654in}}%
\pgfpathlineto{\pgfqpoint{0.707991in}{0.498813in}}%
\pgfpathlineto{\pgfqpoint{0.709394in}{0.499451in}}%
\pgfpathlineto{\pgfqpoint{0.712253in}{0.533597in}}%
\pgfpathlineto{\pgfqpoint{0.712903in}{0.533504in}}%
\pgfpathlineto{\pgfqpoint{0.714272in}{0.533308in}}%
\pgfpathlineto{\pgfqpoint{0.714322in}{0.533986in}}%
\pgfpathlineto{\pgfqpoint{0.717199in}{0.578050in}}%
\pgfpathlineto{\pgfqpoint{0.717827in}{0.577903in}}%
\pgfpathlineto{\pgfqpoint{0.719221in}{0.577579in}}%
\pgfpathlineto{\pgfqpoint{0.719264in}{0.578296in}}%
\pgfpathlineto{\pgfqpoint{0.722144in}{0.630825in}}%
\pgfpathlineto{\pgfqpoint{0.722774in}{0.630611in}}%
\pgfpathlineto{\pgfqpoint{0.724167in}{0.630138in}}%
\pgfpathlineto{\pgfqpoint{0.724202in}{0.630757in}}%
\pgfpathlineto{\pgfqpoint{0.727090in}{0.690393in}}%
\pgfpathlineto{\pgfqpoint{0.727681in}{0.690120in}}%
\pgfpathlineto{\pgfqpoint{0.729108in}{0.689465in}}%
\pgfpathlineto{\pgfqpoint{0.729149in}{0.690255in}}%
\pgfpathlineto{\pgfqpoint{0.732035in}{0.755132in}}%
\pgfpathlineto{\pgfqpoint{0.732606in}{0.754795in}}%
\pgfpathlineto{\pgfqpoint{0.734057in}{0.753938in}}%
\pgfpathlineto{\pgfqpoint{0.734097in}{0.754794in}}%
\pgfpathlineto{\pgfqpoint{0.736981in}{0.823061in}}%
\pgfpathlineto{\pgfqpoint{0.737540in}{0.822653in}}%
\pgfpathlineto{\pgfqpoint{0.739003in}{0.821590in}}%
\pgfpathlineto{\pgfqpoint{0.739045in}{0.822506in}}%
\pgfpathlineto{\pgfqpoint{0.741926in}{0.892309in}}%
\pgfpathlineto{\pgfqpoint{0.742417in}{0.891882in}}%
\pgfpathlineto{\pgfqpoint{0.743946in}{0.890557in}}%
\pgfpathlineto{\pgfqpoint{0.743985in}{0.891359in}}%
\pgfpathlineto{\pgfqpoint{0.746871in}{0.960977in}}%
\pgfpathlineto{\pgfqpoint{0.747342in}{0.960502in}}%
\pgfpathlineto{\pgfqpoint{0.748894in}{0.958943in}}%
\pgfpathlineto{\pgfqpoint{0.748935in}{0.959782in}}%
\pgfpathlineto{\pgfqpoint{0.751817in}{1.027093in}}%
\pgfpathlineto{\pgfqpoint{0.752235in}{1.026616in}}%
\pgfpathlineto{\pgfqpoint{0.753851in}{1.024776in}}%
\pgfpathlineto{\pgfqpoint{0.753900in}{1.025829in}}%
\pgfpathlineto{\pgfqpoint{0.756762in}{1.088670in}}%
\pgfpathlineto{\pgfqpoint{0.757151in}{1.088178in}}%
\pgfpathlineto{\pgfqpoint{0.758787in}{1.086113in}}%
\pgfpathlineto{\pgfqpoint{0.758844in}{1.087220in}}%
\pgfpathlineto{\pgfqpoint{0.761707in}{1.144508in}}%
\pgfpathlineto{\pgfqpoint{0.762058in}{1.144024in}}%
\pgfpathlineto{\pgfqpoint{0.763739in}{1.141712in}}%
\pgfpathlineto{\pgfqpoint{0.763800in}{1.142819in}}%
\pgfpathlineto{\pgfqpoint{0.766653in}{1.192817in}}%
\pgfpathlineto{\pgfqpoint{0.766918in}{1.192426in}}%
\pgfpathlineto{\pgfqpoint{0.768690in}{1.189815in}}%
\pgfpathlineto{\pgfqpoint{0.768757in}{1.190878in}}%
\pgfpathlineto{\pgfqpoint{0.771598in}{1.232366in}}%
\pgfpathlineto{\pgfqpoint{0.771771in}{1.232097in}}%
\pgfpathlineto{\pgfqpoint{0.773633in}{1.229206in}}%
\pgfpathlineto{\pgfqpoint{0.773706in}{1.230193in}}%
\pgfpathlineto{\pgfqpoint{0.776543in}{1.262363in}}%
\pgfpathlineto{\pgfqpoint{0.776603in}{1.262266in}}%
\pgfpathlineto{\pgfqpoint{0.778579in}{1.259078in}}%
\pgfpathlineto{\pgfqpoint{0.778686in}{1.260231in}}%
\pgfpathlineto{\pgfqpoint{0.781258in}{1.281057in}}%
\pgfpathlineto{\pgfqpoint{0.781488in}{1.282030in}}%
\pgfpathlineto{\pgfqpoint{0.782017in}{1.281153in}}%
\pgfpathlineto{\pgfqpoint{0.783533in}{1.278656in}}%
\pgfpathlineto{\pgfqpoint{0.783649in}{1.279540in}}%
\pgfpathlineto{\pgfqpoint{0.785776in}{1.290143in}}%
\pgfpathlineto{\pgfqpoint{0.786428in}{1.290900in}}%
\pgfpathlineto{\pgfqpoint{0.786713in}{1.290427in}}%
\pgfpathlineto{\pgfqpoint{0.788470in}{1.287492in}}%
\pgfpathlineto{\pgfqpoint{0.788706in}{1.288427in}}%
\pgfpathlineto{\pgfqpoint{0.790189in}{1.291153in}}%
\pgfpathlineto{\pgfqpoint{0.791042in}{1.290122in}}%
\pgfpathlineto{\pgfqpoint{0.793435in}{1.286163in}}%
\pgfpathlineto{\pgfqpoint{0.793780in}{1.286936in}}%
\pgfpathlineto{\pgfqpoint{0.794621in}{1.287455in}}%
\pgfpathlineto{\pgfqpoint{0.794818in}{1.287295in}}%
\pgfpathlineto{\pgfqpoint{0.797187in}{1.283421in}}%
\pgfpathlineto{\pgfqpoint{0.798368in}{1.281464in}}%
\pgfpathlineto{\pgfqpoint{0.798732in}{1.282296in}}%
\pgfpathlineto{\pgfqpoint{0.799622in}{1.282895in}}%
\pgfpathlineto{\pgfqpoint{0.799820in}{1.282734in}}%
\pgfpathlineto{\pgfqpoint{0.802237in}{1.278800in}}%
\pgfpathlineto{\pgfqpoint{0.803316in}{1.277022in}}%
\pgfpathlineto{\pgfqpoint{0.803680in}{1.277860in}}%
\pgfpathlineto{\pgfqpoint{0.804570in}{1.278465in}}%
\pgfpathlineto{\pgfqpoint{0.804768in}{1.278307in}}%
\pgfpathlineto{\pgfqpoint{0.806789in}{1.275051in}}%
\pgfpathlineto{\pgfqpoint{0.808262in}{1.272636in}}%
\pgfpathlineto{\pgfqpoint{0.808630in}{1.273482in}}%
\pgfpathlineto{\pgfqpoint{0.809520in}{1.274084in}}%
\pgfpathlineto{\pgfqpoint{0.809718in}{1.273928in}}%
\pgfpathlineto{\pgfqpoint{0.811726in}{1.270709in}}%
\pgfpathlineto{\pgfqpoint{0.813215in}{1.268318in}}%
\pgfpathlineto{\pgfqpoint{0.813534in}{1.269090in}}%
\pgfpathlineto{\pgfqpoint{0.814474in}{1.269778in}}%
\pgfpathlineto{\pgfqpoint{0.814672in}{1.269621in}}%
\pgfpathlineto{\pgfqpoint{0.817092in}{1.265743in}}%
\pgfpathlineto{\pgfqpoint{0.818161in}{1.264047in}}%
\pgfpathlineto{\pgfqpoint{0.818503in}{1.264874in}}%
\pgfpathlineto{\pgfqpoint{0.819443in}{1.265535in}}%
\pgfpathlineto{\pgfqpoint{0.819591in}{1.265424in}}%
\pgfpathlineto{\pgfqpoint{0.821503in}{1.262414in}}%
\pgfpathlineto{\pgfqpoint{0.823105in}{1.259857in}}%
\pgfpathlineto{\pgfqpoint{0.823424in}{1.260634in}}%
\pgfpathlineto{\pgfqpoint{0.824363in}{1.261361in}}%
\pgfpathlineto{\pgfqpoint{0.824561in}{1.261216in}}%
\pgfpathlineto{\pgfqpoint{0.826252in}{1.258569in}}%
\pgfpathlineto{\pgfqpoint{0.828050in}{1.255707in}}%
\pgfpathlineto{\pgfqpoint{0.828385in}{1.256517in}}%
\pgfpathlineto{\pgfqpoint{0.829324in}{1.257210in}}%
\pgfpathlineto{\pgfqpoint{0.829522in}{1.257061in}}%
\pgfpathlineto{\pgfqpoint{0.831585in}{1.253830in}}%
\pgfpathlineto{\pgfqpoint{0.832986in}{1.251592in}}%
\pgfpathlineto{\pgfqpoint{0.833357in}{1.252446in}}%
\pgfpathlineto{\pgfqpoint{0.834297in}{1.253075in}}%
\pgfpathlineto{\pgfqpoint{0.834445in}{1.252962in}}%
\pgfpathlineto{\pgfqpoint{0.836466in}{1.249823in}}%
\pgfpathlineto{\pgfqpoint{0.837943in}{1.247514in}}%
\pgfpathlineto{\pgfqpoint{0.838251in}{1.248272in}}%
\pgfpathlineto{\pgfqpoint{0.839240in}{1.249024in}}%
\pgfpathlineto{\pgfqpoint{0.839388in}{1.248917in}}%
\pgfpathlineto{\pgfqpoint{0.841302in}{1.245966in}}%
\pgfpathlineto{\pgfqpoint{0.842886in}{1.243486in}}%
\pgfpathlineto{\pgfqpoint{0.843216in}{1.244294in}}%
\pgfpathlineto{\pgfqpoint{0.844205in}{1.245026in}}%
\pgfpathlineto{\pgfqpoint{0.844353in}{1.244918in}}%
\pgfpathlineto{\pgfqpoint{0.846187in}{1.242115in}}%
\pgfpathlineto{\pgfqpoint{0.847835in}{1.239559in}}%
\pgfpathlineto{\pgfqpoint{0.848157in}{1.240339in}}%
\pgfpathlineto{\pgfqpoint{0.849146in}{1.241059in}}%
\pgfpathlineto{\pgfqpoint{0.849294in}{1.240949in}}%
\pgfpathlineto{\pgfqpoint{0.851338in}{1.237819in}}%
\pgfpathlineto{\pgfqpoint{0.852770in}{1.235578in}}%
\pgfpathlineto{\pgfqpoint{0.853088in}{1.236348in}}%
\pgfpathlineto{\pgfqpoint{0.854077in}{1.237157in}}%
\pgfpathlineto{\pgfqpoint{0.854275in}{1.237018in}}%
\pgfpathlineto{\pgfqpoint{0.856566in}{1.233514in}}%
\pgfpathlineto{\pgfqpoint{0.857712in}{1.231730in}}%
\pgfpathlineto{\pgfqpoint{0.858080in}{1.232581in}}%
\pgfpathlineto{\pgfqpoint{0.859020in}{1.233261in}}%
\pgfpathlineto{\pgfqpoint{0.859218in}{1.233115in}}%
\pgfpathlineto{\pgfqpoint{0.861578in}{1.229516in}}%
\pgfpathlineto{\pgfqpoint{0.862672in}{1.227853in}}%
\pgfpathlineto{\pgfqpoint{0.863013in}{1.228683in}}%
\pgfpathlineto{\pgfqpoint{0.864002in}{1.229409in}}%
\pgfpathlineto{\pgfqpoint{0.864150in}{1.229303in}}%
\pgfpathlineto{\pgfqpoint{0.866096in}{1.226376in}}%
\pgfpathlineto{\pgfqpoint{0.867615in}{1.224068in}}%
\pgfpathlineto{\pgfqpoint{0.867920in}{1.224824in}}%
\pgfpathlineto{\pgfqpoint{0.868959in}{1.225624in}}%
\pgfpathlineto{\pgfqpoint{0.869107in}{1.225515in}}%
\pgfpathlineto{\pgfqpoint{0.871207in}{1.222363in}}%
\pgfpathlineto{\pgfqpoint{0.872560in}{1.220323in}}%
\pgfpathlineto{\pgfqpoint{0.872877in}{1.221106in}}%
\pgfpathlineto{\pgfqpoint{0.873916in}{1.221891in}}%
\pgfpathlineto{\pgfqpoint{0.874064in}{1.221780in}}%
\pgfpathlineto{\pgfqpoint{0.876212in}{1.218565in}}%
\pgfpathlineto{\pgfqpoint{0.877493in}{1.216609in}}%
\pgfpathlineto{\pgfqpoint{0.877809in}{1.217368in}}%
\pgfpathlineto{\pgfqpoint{0.878848in}{1.218203in}}%
\pgfpathlineto{\pgfqpoint{0.878996in}{1.218101in}}%
\pgfpathlineto{\pgfqpoint{0.880933in}{1.215236in}}%
\pgfpathlineto{\pgfqpoint{0.882451in}{1.212963in}}%
\pgfpathlineto{\pgfqpoint{0.882784in}{1.213777in}}%
\pgfpathlineto{\pgfqpoint{0.883774in}{1.214552in}}%
\pgfpathlineto{\pgfqpoint{0.883971in}{1.214414in}}%
\pgfpathlineto{\pgfqpoint{0.886237in}{1.211046in}}%
\pgfpathlineto{\pgfqpoint{0.887383in}{1.209313in}}%
\pgfpathlineto{\pgfqpoint{0.887698in}{1.210072in}}%
\pgfpathlineto{\pgfqpoint{0.888737in}{1.210931in}}%
\pgfpathlineto{\pgfqpoint{0.888885in}{1.210834in}}%
\pgfpathlineto{\pgfqpoint{0.890735in}{1.208138in}}%
\pgfpathlineto{\pgfqpoint{0.892330in}{1.205738in}}%
\pgfpathlineto{\pgfqpoint{0.892649in}{1.206501in}}%
\pgfpathlineto{\pgfqpoint{0.893687in}{1.207345in}}%
\pgfpathlineto{\pgfqpoint{0.893836in}{1.207247in}}%
\pgfpathlineto{\pgfqpoint{0.895677in}{1.204571in}}%
\pgfpathlineto{\pgfqpoint{0.897273in}{1.202181in}}%
\pgfpathlineto{\pgfqpoint{0.897629in}{1.203020in}}%
\pgfpathlineto{\pgfqpoint{0.898668in}{1.203779in}}%
\pgfpathlineto{\pgfqpoint{0.898816in}{1.203670in}}%
\pgfpathlineto{\pgfqpoint{0.900994in}{1.200488in}}%
\pgfpathlineto{\pgfqpoint{0.902232in}{1.198671in}}%
\pgfpathlineto{\pgfqpoint{0.902547in}{1.199454in}}%
\pgfpathlineto{\pgfqpoint{0.903585in}{1.200312in}}%
\pgfpathlineto{\pgfqpoint{0.903734in}{1.200222in}}%
\pgfpathlineto{\pgfqpoint{0.905757in}{1.197308in}}%
\pgfpathlineto{\pgfqpoint{0.907166in}{1.195219in}}%
\pgfpathlineto{\pgfqpoint{0.907489in}{1.195985in}}%
\pgfpathlineto{\pgfqpoint{0.908527in}{1.196828in}}%
\pgfpathlineto{\pgfqpoint{0.908676in}{1.196733in}}%
\pgfpathlineto{\pgfqpoint{0.910538in}{1.194066in}}%
\pgfpathlineto{\pgfqpoint{0.912126in}{1.191753in}}%
\pgfpathlineto{\pgfqpoint{0.912448in}{1.192549in}}%
\pgfpathlineto{\pgfqpoint{0.913487in}{1.193391in}}%
\pgfpathlineto{\pgfqpoint{0.913635in}{1.193298in}}%
\pgfpathlineto{\pgfqpoint{0.915456in}{1.190708in}}%
\pgfpathlineto{\pgfqpoint{0.917072in}{1.188361in}}%
\pgfpathlineto{\pgfqpoint{0.917374in}{1.189109in}}%
\pgfpathlineto{\pgfqpoint{0.918462in}{1.189974in}}%
\pgfpathlineto{\pgfqpoint{0.918561in}{1.189910in}}%
\pgfpathlineto{\pgfqpoint{0.920242in}{1.187547in}}%
\pgfpathlineto{\pgfqpoint{0.922007in}{1.184965in}}%
\pgfpathlineto{\pgfqpoint{0.922331in}{1.185751in}}%
\pgfpathlineto{\pgfqpoint{0.923419in}{1.186620in}}%
\pgfpathlineto{\pgfqpoint{0.923518in}{1.186560in}}%
\pgfpathlineto{\pgfqpoint{0.925305in}{1.184064in}}%
\pgfpathlineto{\pgfqpoint{0.926959in}{1.181682in}}%
\pgfpathlineto{\pgfqpoint{0.927263in}{1.182430in}}%
\pgfpathlineto{\pgfqpoint{0.928351in}{1.183288in}}%
\pgfpathlineto{\pgfqpoint{0.928450in}{1.183225in}}%
\pgfpathlineto{\pgfqpoint{0.930159in}{1.180848in}}%
\pgfpathlineto{\pgfqpoint{0.931904in}{1.178345in}}%
\pgfpathlineto{\pgfqpoint{0.932209in}{1.179108in}}%
\pgfpathlineto{\pgfqpoint{0.933297in}{1.180022in}}%
\pgfpathlineto{\pgfqpoint{0.933446in}{1.179928in}}%
\pgfpathlineto{\pgfqpoint{0.935031in}{1.177727in}}%
\pgfpathlineto{\pgfqpoint{0.936843in}{1.175111in}}%
\pgfpathlineto{\pgfqpoint{0.937157in}{1.175869in}}%
\pgfpathlineto{\pgfqpoint{0.938245in}{1.176766in}}%
\pgfpathlineto{\pgfqpoint{0.938394in}{1.176670in}}%
\pgfpathlineto{\pgfqpoint{0.940218in}{1.174127in}}%
\pgfpathlineto{\pgfqpoint{0.941797in}{1.171887in}}%
\pgfpathlineto{\pgfqpoint{0.942128in}{1.172695in}}%
\pgfpathlineto{\pgfqpoint{0.943216in}{1.173528in}}%
\pgfpathlineto{\pgfqpoint{0.943315in}{1.173465in}}%
\pgfpathlineto{\pgfqpoint{0.945064in}{1.171058in}}%
\pgfpathlineto{\pgfqpoint{0.946730in}{1.168675in}}%
\pgfpathlineto{\pgfqpoint{0.947064in}{1.169480in}}%
\pgfpathlineto{\pgfqpoint{0.948152in}{1.170376in}}%
\pgfpathlineto{\pgfqpoint{0.948301in}{1.170282in}}%
\pgfpathlineto{\pgfqpoint{0.950092in}{1.167810in}}%
\pgfpathlineto{\pgfqpoint{0.951679in}{1.165553in}}%
\pgfpathlineto{\pgfqpoint{0.952001in}{1.166321in}}%
\pgfpathlineto{\pgfqpoint{0.953089in}{1.167196in}}%
\pgfpathlineto{\pgfqpoint{0.953237in}{1.167101in}}%
\pgfpathlineto{\pgfqpoint{0.955095in}{1.164548in}}%
\pgfpathlineto{\pgfqpoint{0.956615in}{1.162393in}}%
\pgfpathlineto{\pgfqpoint{0.956961in}{1.163212in}}%
\pgfpathlineto{\pgfqpoint{0.958049in}{1.164075in}}%
\pgfpathlineto{\pgfqpoint{0.958198in}{1.163979in}}%
\pgfpathlineto{\pgfqpoint{0.960262in}{1.161135in}}%
\pgfpathlineto{\pgfqpoint{0.961579in}{1.159307in}}%
\pgfpathlineto{\pgfqpoint{0.961909in}{1.160127in}}%
\pgfpathlineto{\pgfqpoint{0.963046in}{1.161013in}}%
\pgfpathlineto{\pgfqpoint{0.963145in}{1.160947in}}%
\pgfpathlineto{\pgfqpoint{0.964984in}{1.158445in}}%
\pgfpathlineto{\pgfqpoint{0.966506in}{1.156307in}}%
\pgfpathlineto{\pgfqpoint{0.966829in}{1.157071in}}%
\pgfpathlineto{\pgfqpoint{0.967966in}{1.157967in}}%
\pgfpathlineto{\pgfqpoint{0.968065in}{1.157903in}}%
\pgfpathlineto{\pgfqpoint{0.969926in}{1.155387in}}%
\pgfpathlineto{\pgfqpoint{0.971468in}{1.153258in}}%
\pgfpathlineto{\pgfqpoint{0.971793in}{1.154072in}}%
\pgfpathlineto{\pgfqpoint{0.972930in}{1.155011in}}%
\pgfpathlineto{\pgfqpoint{0.973029in}{1.154951in}}%
\pgfpathlineto{\pgfqpoint{0.974751in}{1.152648in}}%
\pgfpathlineto{\pgfqpoint{0.976399in}{1.150352in}}%
\pgfpathlineto{\pgfqpoint{0.976743in}{1.151163in}}%
\pgfpathlineto{\pgfqpoint{0.977881in}{1.152022in}}%
\pgfpathlineto{\pgfqpoint{0.977979in}{1.151956in}}%
\pgfpathlineto{\pgfqpoint{0.979905in}{1.149361in}}%
\pgfpathlineto{\pgfqpoint{0.981349in}{1.147359in}}%
\pgfpathlineto{\pgfqpoint{0.981687in}{1.148177in}}%
\pgfpathlineto{\pgfqpoint{0.982825in}{1.149104in}}%
\pgfpathlineto{\pgfqpoint{0.982924in}{1.149044in}}%
\pgfpathlineto{\pgfqpoint{0.984473in}{1.147005in}}%
\pgfpathlineto{\pgfqpoint{0.986288in}{1.144498in}}%
\pgfpathlineto{\pgfqpoint{0.986598in}{1.145237in}}%
\pgfpathlineto{\pgfqpoint{0.987736in}{1.146195in}}%
\pgfpathlineto{\pgfqpoint{0.987834in}{1.146139in}}%
\pgfpathlineto{\pgfqpoint{0.989488in}{1.143970in}}%
\pgfpathlineto{\pgfqpoint{0.991239in}{1.141563in}}%
\pgfpathlineto{\pgfqpoint{0.991574in}{1.142379in}}%
\pgfpathlineto{\pgfqpoint{0.992712in}{1.143343in}}%
\pgfpathlineto{\pgfqpoint{0.992811in}{1.143288in}}%
\pgfpathlineto{\pgfqpoint{0.994370in}{1.141265in}}%
\pgfpathlineto{\pgfqpoint{0.996198in}{1.138789in}}%
\pgfpathlineto{\pgfqpoint{0.996529in}{1.139601in}}%
\pgfpathlineto{\pgfqpoint{0.997667in}{1.140498in}}%
\pgfpathlineto{\pgfqpoint{0.997766in}{1.140437in}}%
\pgfpathlineto{\pgfqpoint{0.999568in}{1.138068in}}%
\pgfpathlineto{\pgfqpoint{1.001130in}{1.135938in}}%
\pgfpathlineto{\pgfqpoint{1.001470in}{1.136760in}}%
\pgfpathlineto{\pgfqpoint{1.002607in}{1.137718in}}%
\pgfpathlineto{\pgfqpoint{1.002706in}{1.137664in}}%
\pgfpathlineto{\pgfqpoint{1.004166in}{1.135794in}}%
\pgfpathlineto{\pgfqpoint{1.006080in}{1.133207in}}%
\pgfpathlineto{\pgfqpoint{1.006393in}{1.133971in}}%
\pgfpathlineto{\pgfqpoint{1.007580in}{1.134923in}}%
\pgfpathlineto{\pgfqpoint{1.007630in}{1.134894in}}%
\pgfpathlineto{\pgfqpoint{1.009207in}{1.132879in}}%
\pgfpathlineto{\pgfqpoint{1.011016in}{1.130432in}}%
\pgfpathlineto{\pgfqpoint{1.011357in}{1.131254in}}%
\pgfpathlineto{\pgfqpoint{1.012544in}{1.132206in}}%
\pgfpathlineto{\pgfqpoint{1.012642in}{1.132144in}}%
\pgfpathlineto{\pgfqpoint{1.014552in}{1.129651in}}%
\pgfpathlineto{\pgfqpoint{1.015968in}{1.127745in}}%
\pgfpathlineto{\pgfqpoint{1.016307in}{1.128559in}}%
\pgfpathlineto{\pgfqpoint{1.017494in}{1.129485in}}%
\pgfpathlineto{\pgfqpoint{1.017543in}{1.129455in}}%
\pgfpathlineto{\pgfqpoint{1.019146in}{1.127411in}}%
\pgfpathlineto{\pgfqpoint{1.020918in}{1.125064in}}%
\pgfpathlineto{\pgfqpoint{1.021203in}{1.125789in}}%
\pgfpathlineto{\pgfqpoint{1.022440in}{1.126849in}}%
\pgfpathlineto{\pgfqpoint{1.022489in}{1.126821in}}%
\pgfpathlineto{\pgfqpoint{1.024048in}{1.124851in}}%
\pgfpathlineto{\pgfqpoint{1.025853in}{1.122438in}}%
\pgfpathlineto{\pgfqpoint{1.026171in}{1.123199in}}%
\pgfpathlineto{\pgfqpoint{1.027358in}{1.124191in}}%
\pgfpathlineto{\pgfqpoint{1.027457in}{1.124134in}}%
\pgfpathlineto{\pgfqpoint{1.029245in}{1.121843in}}%
\pgfpathlineto{\pgfqpoint{1.030806in}{1.119769in}}%
\pgfpathlineto{\pgfqpoint{1.031125in}{1.120548in}}%
\pgfpathlineto{\pgfqpoint{1.032312in}{1.121551in}}%
\pgfpathlineto{\pgfqpoint{1.032411in}{1.121500in}}%
\pgfpathlineto{\pgfqpoint{1.034437in}{1.118909in}}%
\pgfpathlineto{\pgfqpoint{1.035751in}{1.117179in}}%
\pgfpathlineto{\pgfqpoint{1.036066in}{1.117946in}}%
\pgfpathlineto{\pgfqpoint{1.037253in}{1.118931in}}%
\pgfpathlineto{\pgfqpoint{1.037352in}{1.118877in}}%
\pgfpathlineto{\pgfqpoint{1.039331in}{1.116357in}}%
\pgfpathlineto{\pgfqpoint{1.040687in}{1.114568in}}%
\pgfpathlineto{\pgfqpoint{1.041006in}{1.115339in}}%
\pgfpathlineto{\pgfqpoint{1.042242in}{1.116361in}}%
\pgfpathlineto{\pgfqpoint{1.042291in}{1.116332in}}%
\pgfpathlineto{\pgfqpoint{1.043929in}{1.114283in}}%
\pgfpathlineto{\pgfqpoint{1.045632in}{1.112042in}}%
\pgfpathlineto{\pgfqpoint{1.045969in}{1.112855in}}%
\pgfpathlineto{\pgfqpoint{1.047156in}{1.113870in}}%
\pgfpathlineto{\pgfqpoint{1.047255in}{1.113817in}}%
\pgfpathlineto{\pgfqpoint{1.048966in}{1.111674in}}%
\pgfpathlineto{\pgfqpoint{1.050578in}{1.109562in}}%
\pgfpathlineto{\pgfqpoint{1.050882in}{1.110304in}}%
\pgfpathlineto{\pgfqpoint{1.052118in}{1.111398in}}%
\pgfpathlineto{\pgfqpoint{1.052168in}{1.111373in}}%
\pgfpathlineto{\pgfqpoint{1.053651in}{1.109566in}}%
\pgfpathlineto{\pgfqpoint{1.055525in}{1.107120in}}%
\pgfpathlineto{\pgfqpoint{1.055827in}{1.107862in}}%
\pgfpathlineto{\pgfqpoint{1.057064in}{1.108974in}}%
\pgfpathlineto{\pgfqpoint{1.057113in}{1.108950in}}%
\pgfpathlineto{\pgfqpoint{1.058563in}{1.107196in}}%
\pgfpathlineto{\pgfqpoint{1.060469in}{1.104716in}}%
\pgfpathlineto{\pgfqpoint{1.060787in}{1.105479in}}%
\pgfpathlineto{\pgfqpoint{1.062024in}{1.106506in}}%
\pgfpathlineto{\pgfqpoint{1.062073in}{1.106478in}}%
\pgfpathlineto{\pgfqpoint{1.063643in}{1.104555in}}%
\pgfpathlineto{\pgfqpoint{1.065425in}{1.102273in}}%
\pgfpathlineto{\pgfqpoint{1.065744in}{1.103080in}}%
\pgfpathlineto{\pgfqpoint{1.066981in}{1.104159in}}%
\pgfpathlineto{\pgfqpoint{1.067030in}{1.104134in}}%
\pgfpathlineto{\pgfqpoint{1.068536in}{1.102320in}}%
\pgfpathlineto{\pgfqpoint{1.070361in}{1.099963in}}%
\pgfpathlineto{\pgfqpoint{1.070665in}{1.100700in}}%
\pgfpathlineto{\pgfqpoint{1.071901in}{1.101797in}}%
\pgfpathlineto{\pgfqpoint{1.071950in}{1.101773in}}%
\pgfpathlineto{\pgfqpoint{1.073399in}{1.100039in}}%
\pgfpathlineto{\pgfqpoint{1.075306in}{1.097586in}}%
\pgfpathlineto{\pgfqpoint{1.075627in}{1.098364in}}%
\pgfpathlineto{\pgfqpoint{1.076864in}{1.099427in}}%
\pgfpathlineto{\pgfqpoint{1.076913in}{1.099402in}}%
\pgfpathlineto{\pgfqpoint{1.078411in}{1.097606in}}%
\pgfpathlineto{\pgfqpoint{1.080253in}{1.095246in}}%
\pgfpathlineto{\pgfqpoint{1.080576in}{1.096038in}}%
\pgfpathlineto{\pgfqpoint{1.081813in}{1.097139in}}%
\pgfpathlineto{\pgfqpoint{1.081912in}{1.097087in}}%
\pgfpathlineto{\pgfqpoint{1.083613in}{1.095014in}}%
\pgfpathlineto{\pgfqpoint{1.085207in}{1.093007in}}%
\pgfpathlineto{\pgfqpoint{1.085496in}{1.093744in}}%
\pgfpathlineto{\pgfqpoint{1.086782in}{1.094890in}}%
\pgfpathlineto{\pgfqpoint{1.086831in}{1.094865in}}%
\pgfpathlineto{\pgfqpoint{1.088383in}{1.093009in}}%
\pgfpathlineto{\pgfqpoint{1.090151in}{1.090769in}}%
\pgfpathlineto{\pgfqpoint{1.090464in}{1.091543in}}%
\pgfpathlineto{\pgfqpoint{1.091750in}{1.092618in}}%
\pgfpathlineto{\pgfqpoint{1.091799in}{1.092590in}}%
\pgfpathlineto{\pgfqpoint{1.093242in}{1.090865in}}%
\pgfpathlineto{\pgfqpoint{1.095094in}{1.088517in}}%
\pgfpathlineto{\pgfqpoint{1.095385in}{1.089230in}}%
\pgfpathlineto{\pgfqpoint{1.096720in}{1.090349in}}%
\pgfpathlineto{\pgfqpoint{1.098451in}{1.088285in}}%
\pgfpathlineto{\pgfqpoint{1.100045in}{1.086303in}}%
\pgfpathlineto{\pgfqpoint{1.100360in}{1.087097in}}%
\pgfpathlineto{\pgfqpoint{1.101646in}{1.088187in}}%
\pgfpathlineto{\pgfqpoint{1.101695in}{1.088160in}}%
\pgfpathlineto{\pgfqpoint{1.103346in}{1.086185in}}%
\pgfpathlineto{\pgfqpoint{1.104988in}{1.084144in}}%
\pgfpathlineto{\pgfqpoint{1.105280in}{1.084893in}}%
\pgfpathlineto{\pgfqpoint{1.106566in}{1.086077in}}%
\pgfpathlineto{\pgfqpoint{1.106616in}{1.086054in}}%
\pgfpathlineto{\pgfqpoint{1.108105in}{1.084320in}}%
\pgfpathlineto{\pgfqpoint{1.109931in}{1.082028in}}%
\pgfpathlineto{\pgfqpoint{1.110221in}{1.082737in}}%
\pgfpathlineto{\pgfqpoint{1.111556in}{1.083879in}}%
\pgfpathlineto{\pgfqpoint{1.113198in}{1.081959in}}%
\pgfpathlineto{\pgfqpoint{1.114879in}{1.079885in}}%
\pgfpathlineto{\pgfqpoint{1.115169in}{1.080628in}}%
\pgfpathlineto{\pgfqpoint{1.116505in}{1.081806in}}%
\pgfpathlineto{\pgfqpoint{1.117963in}{1.080129in}}%
\pgfpathlineto{\pgfqpoint{1.119814in}{1.077822in}}%
\pgfpathlineto{\pgfqpoint{1.120129in}{1.078597in}}%
\pgfpathlineto{\pgfqpoint{1.121415in}{1.079766in}}%
\pgfpathlineto{\pgfqpoint{1.121464in}{1.079743in}}%
\pgfpathlineto{\pgfqpoint{1.122926in}{1.078053in}}%
\pgfpathlineto{\pgfqpoint{1.124768in}{1.075775in}}%
\pgfpathlineto{\pgfqpoint{1.125056in}{1.076495in}}%
\pgfpathlineto{\pgfqpoint{1.126392in}{1.077678in}}%
\pgfpathlineto{\pgfqpoint{1.128011in}{1.075818in}}%
\pgfpathlineto{\pgfqpoint{1.129724in}{1.073724in}}%
\pgfpathlineto{\pgfqpoint{1.130023in}{1.074489in}}%
\pgfpathlineto{\pgfqpoint{1.131358in}{1.075672in}}%
\pgfpathlineto{\pgfqpoint{1.132772in}{1.074062in}}%
\pgfpathlineto{\pgfqpoint{1.134652in}{1.071742in}}%
\pgfpathlineto{\pgfqpoint{1.134948in}{1.072471in}}%
\pgfpathlineto{\pgfqpoint{1.136283in}{1.073679in}}%
\pgfpathlineto{\pgfqpoint{1.137674in}{1.072113in}}%
\pgfpathlineto{\pgfqpoint{1.139603in}{1.069740in}}%
\pgfpathlineto{\pgfqpoint{1.139893in}{1.070460in}}%
\pgfpathlineto{\pgfqpoint{1.141228in}{1.071692in}}%
\pgfpathlineto{\pgfqpoint{1.141277in}{1.071666in}}%
\pgfpathlineto{\pgfqpoint{1.142672in}{1.070083in}}%
\pgfpathlineto{\pgfqpoint{1.144551in}{1.067800in}}%
\pgfpathlineto{\pgfqpoint{1.144860in}{1.068590in}}%
\pgfpathlineto{\pgfqpoint{1.146195in}{1.069788in}}%
\pgfpathlineto{\pgfqpoint{1.147568in}{1.068257in}}%
\pgfpathlineto{\pgfqpoint{1.149494in}{1.065903in}}%
\pgfpathlineto{\pgfqpoint{1.149784in}{1.066624in}}%
\pgfpathlineto{\pgfqpoint{1.151119in}{1.067875in}}%
\pgfpathlineto{\pgfqpoint{1.151168in}{1.067849in}}%
\pgfpathlineto{\pgfqpoint{1.152583in}{1.066260in}}%
\pgfpathlineto{\pgfqpoint{1.154443in}{1.064018in}}%
\pgfpathlineto{\pgfqpoint{1.154750in}{1.064809in}}%
\pgfpathlineto{\pgfqpoint{1.156085in}{1.066051in}}%
\pgfpathlineto{\pgfqpoint{1.156135in}{1.066028in}}%
\pgfpathlineto{\pgfqpoint{1.157667in}{1.064291in}}%
\pgfpathlineto{\pgfqpoint{1.159386in}{1.062219in}}%
\pgfpathlineto{\pgfqpoint{1.159674in}{1.062962in}}%
\pgfpathlineto{\pgfqpoint{1.161059in}{1.064250in}}%
\pgfpathlineto{\pgfqpoint{1.162319in}{1.062875in}}%
\pgfpathlineto{\pgfqpoint{1.164340in}{1.060454in}}%
\pgfpathlineto{\pgfqpoint{1.164629in}{1.061203in}}%
\pgfpathlineto{\pgfqpoint{1.166014in}{1.062477in}}%
\pgfpathlineto{\pgfqpoint{1.167418in}{1.060916in}}%
\pgfpathlineto{\pgfqpoint{1.169278in}{1.058684in}}%
\pgfpathlineto{\pgfqpoint{1.169589in}{1.059476in}}%
\pgfpathlineto{\pgfqpoint{1.170973in}{1.060707in}}%
\pgfpathlineto{\pgfqpoint{1.172483in}{1.059021in}}%
\pgfpathlineto{\pgfqpoint{1.174223in}{1.056949in}}%
\pgfpathlineto{\pgfqpoint{1.174509in}{1.057696in}}%
\pgfpathlineto{\pgfqpoint{1.175894in}{1.059022in}}%
\pgfpathlineto{\pgfqpoint{1.177200in}{1.057613in}}%
\pgfpathlineto{\pgfqpoint{1.179169in}{1.055267in}}%
\pgfpathlineto{\pgfqpoint{1.179478in}{1.056068in}}%
\pgfpathlineto{\pgfqpoint{1.180862in}{1.057361in}}%
\pgfpathlineto{\pgfqpoint{1.182286in}{1.055807in}}%
\pgfpathlineto{\pgfqpoint{1.184121in}{1.053634in}}%
\pgfpathlineto{\pgfqpoint{1.184419in}{1.054422in}}%
\pgfpathlineto{\pgfqpoint{1.185804in}{1.055766in}}%
\pgfpathlineto{\pgfqpoint{1.187124in}{1.054346in}}%
\pgfpathlineto{\pgfqpoint{1.189066in}{1.052052in}}%
\pgfpathlineto{\pgfqpoint{1.189347in}{1.052791in}}%
\pgfpathlineto{\pgfqpoint{1.190732in}{1.054161in}}%
\pgfpathlineto{\pgfqpoint{1.192051in}{1.052757in}}%
\pgfpathlineto{\pgfqpoint{1.193999in}{1.050437in}}%
\pgfpathlineto{\pgfqpoint{1.194309in}{1.051229in}}%
\pgfpathlineto{\pgfqpoint{1.195694in}{1.052581in}}%
\pgfpathlineto{\pgfqpoint{1.196995in}{1.051196in}}%
\pgfpathlineto{\pgfqpoint{1.198945in}{1.048879in}}%
\pgfpathlineto{\pgfqpoint{1.199222in}{1.049594in}}%
\pgfpathlineto{\pgfqpoint{1.200657in}{1.051022in}}%
\pgfpathlineto{\pgfqpoint{1.202016in}{1.049560in}}%
\pgfpathlineto{\pgfqpoint{1.203897in}{1.047356in}}%
\pgfpathlineto{\pgfqpoint{1.204158in}{1.048061in}}%
\pgfpathlineto{\pgfqpoint{1.205593in}{1.049541in}}%
\pgfpathlineto{\pgfqpoint{1.206915in}{1.048144in}}%
\pgfpathlineto{\pgfqpoint{1.208840in}{1.045870in}}%
\pgfpathlineto{\pgfqpoint{1.209145in}{1.046663in}}%
\pgfpathlineto{\pgfqpoint{1.210579in}{1.048066in}}%
\pgfpathlineto{\pgfqpoint{1.211985in}{1.046556in}}%
\pgfpathlineto{\pgfqpoint{1.213779in}{1.044442in}}%
\pgfpathlineto{\pgfqpoint{1.214084in}{1.045242in}}%
\pgfpathlineto{\pgfqpoint{1.215518in}{1.046705in}}%
\pgfpathlineto{\pgfqpoint{1.216815in}{1.045343in}}%
\pgfpathlineto{\pgfqpoint{1.218733in}{1.043105in}}%
\pgfpathlineto{\pgfqpoint{1.218989in}{1.043810in}}%
\pgfpathlineto{\pgfqpoint{1.220424in}{1.045437in}}%
\pgfpathlineto{\pgfqpoint{1.221592in}{1.044285in}}%
\pgfpathlineto{\pgfqpoint{1.223679in}{1.041865in}}%
\pgfpathlineto{\pgfqpoint{1.223938in}{1.042596in}}%
\pgfpathlineto{\pgfqpoint{1.225422in}{1.044272in}}%
\pgfpathlineto{\pgfqpoint{1.226646in}{1.043029in}}%
\pgfpathlineto{\pgfqpoint{1.228625in}{1.040741in}}%
\pgfpathlineto{\pgfqpoint{1.228882in}{1.041479in}}%
\pgfpathlineto{\pgfqpoint{1.230366in}{1.043228in}}%
\pgfpathlineto{\pgfqpoint{1.231552in}{1.042054in}}%
\pgfpathlineto{\pgfqpoint{1.233561in}{1.039705in}}%
\pgfpathlineto{\pgfqpoint{1.233843in}{1.040491in}}%
\pgfpathlineto{\pgfqpoint{1.235327in}{1.042253in}}%
\pgfpathlineto{\pgfqpoint{1.236504in}{1.041103in}}%
\pgfpathlineto{\pgfqpoint{1.238514in}{1.038779in}}%
\pgfpathlineto{\pgfqpoint{1.238771in}{1.039540in}}%
\pgfpathlineto{\pgfqpoint{1.240255in}{1.041447in}}%
\pgfpathlineto{\pgfqpoint{1.241343in}{1.040451in}}%
\pgfpathlineto{\pgfqpoint{1.243466in}{1.038003in}}%
\pgfpathlineto{\pgfqpoint{1.243745in}{1.038829in}}%
\pgfpathlineto{\pgfqpoint{1.245229in}{1.040686in}}%
\pgfpathlineto{\pgfqpoint{1.246305in}{1.039684in}}%
\pgfpathlineto{\pgfqpoint{1.248401in}{1.037243in}}%
\pgfpathlineto{\pgfqpoint{1.248670in}{1.038029in}}%
\pgfpathlineto{\pgfqpoint{1.250203in}{1.039998in}}%
\pgfpathlineto{\pgfqpoint{1.251331in}{1.038929in}}%
\pgfpathlineto{\pgfqpoint{1.253357in}{1.036599in}}%
\pgfpathlineto{\pgfqpoint{1.253627in}{1.037407in}}%
\pgfpathlineto{\pgfqpoint{1.255160in}{1.039338in}}%
\pgfpathlineto{\pgfqpoint{1.256295in}{1.038250in}}%
\pgfpathlineto{\pgfqpoint{1.258293in}{1.035929in}}%
\pgfpathlineto{\pgfqpoint{1.258534in}{1.036639in}}%
\pgfpathlineto{\pgfqpoint{1.260067in}{1.038722in}}%
\pgfpathlineto{\pgfqpoint{1.261105in}{1.037816in}}%
\pgfpathlineto{\pgfqpoint{1.263241in}{1.035353in}}%
\pgfpathlineto{\pgfqpoint{1.263497in}{1.036125in}}%
\pgfpathlineto{\pgfqpoint{1.265030in}{1.038189in}}%
\pgfpathlineto{\pgfqpoint{1.266100in}{1.037227in}}%
\pgfpathlineto{\pgfqpoint{1.268188in}{1.034838in}}%
\pgfpathlineto{\pgfqpoint{1.268444in}{1.035621in}}%
\pgfpathlineto{\pgfqpoint{1.269977in}{1.037687in}}%
\pgfpathlineto{\pgfqpoint{1.271047in}{1.036726in}}%
\pgfpathlineto{\pgfqpoint{1.273133in}{1.034325in}}%
\pgfpathlineto{\pgfqpoint{1.273389in}{1.035104in}}%
\pgfpathlineto{\pgfqpoint{1.274922in}{1.037184in}}%
\pgfpathlineto{\pgfqpoint{1.275992in}{1.036233in}}%
\pgfpathlineto{\pgfqpoint{1.278075in}{1.033822in}}%
\pgfpathlineto{\pgfqpoint{1.278319in}{1.034547in}}%
\pgfpathlineto{\pgfqpoint{1.279852in}{1.036664in}}%
\pgfpathlineto{\pgfqpoint{1.280890in}{1.035759in}}%
\pgfpathlineto{\pgfqpoint{1.283024in}{1.033322in}}%
\pgfpathlineto{\pgfqpoint{1.283286in}{1.034128in}}%
\pgfpathlineto{\pgfqpoint{1.284820in}{1.036205in}}%
\pgfpathlineto{\pgfqpoint{1.285886in}{1.035259in}}%
\pgfpathlineto{\pgfqpoint{1.287963in}{1.032860in}}%
\pgfpathlineto{\pgfqpoint{1.288207in}{1.033580in}}%
\pgfpathlineto{\pgfqpoint{1.289740in}{1.035692in}}%
\pgfpathlineto{\pgfqpoint{1.290778in}{1.034790in}}%
\pgfpathlineto{\pgfqpoint{1.292915in}{1.032354in}}%
\pgfpathlineto{\pgfqpoint{1.293172in}{1.033137in}}%
\pgfpathlineto{\pgfqpoint{1.294705in}{1.035189in}}%
\pgfpathlineto{\pgfqpoint{1.295774in}{1.034231in}}%
\pgfpathlineto{\pgfqpoint{1.297850in}{1.031836in}}%
\pgfpathlineto{\pgfqpoint{1.298135in}{1.032678in}}%
\pgfpathlineto{\pgfqpoint{1.299668in}{1.034716in}}%
\pgfpathlineto{\pgfqpoint{1.300760in}{1.033728in}}%
\pgfpathlineto{\pgfqpoint{1.302801in}{1.031376in}}%
\pgfpathlineto{\pgfqpoint{1.303066in}{1.032166in}}%
\pgfpathlineto{\pgfqpoint{1.304600in}{1.034257in}}%
\pgfpathlineto{\pgfqpoint{1.305638in}{1.033349in}}%
\pgfpathlineto{\pgfqpoint{1.307750in}{1.030937in}}%
\pgfpathlineto{\pgfqpoint{1.308012in}{1.031738in}}%
\pgfpathlineto{\pgfqpoint{1.309545in}{1.033832in}}%
\pgfpathlineto{\pgfqpoint{1.310612in}{1.032899in}}%
\pgfpathlineto{\pgfqpoint{1.312696in}{1.030524in}}%
\pgfpathlineto{\pgfqpoint{1.312951in}{1.031308in}}%
\pgfpathlineto{\pgfqpoint{1.314485in}{1.033430in}}%
\pgfpathlineto{\pgfqpoint{1.315523in}{1.032543in}}%
\pgfpathlineto{\pgfqpoint{1.317635in}{1.030114in}}%
\pgfpathlineto{\pgfqpoint{1.317875in}{1.030824in}}%
\pgfpathlineto{\pgfqpoint{1.319408in}{1.032963in}}%
\pgfpathlineto{\pgfqpoint{1.320437in}{1.032095in}}%
\pgfpathlineto{\pgfqpoint{1.322578in}{1.029636in}}%
\pgfpathlineto{\pgfqpoint{1.322815in}{1.030346in}}%
\pgfpathlineto{\pgfqpoint{1.324398in}{1.032537in}}%
\pgfpathlineto{\pgfqpoint{1.325482in}{1.031568in}}%
\pgfpathlineto{\pgfqpoint{1.327526in}{1.029222in}}%
\pgfpathlineto{\pgfqpoint{1.327765in}{1.029938in}}%
\pgfpathlineto{\pgfqpoint{1.329348in}{1.032121in}}%
\pgfpathlineto{\pgfqpoint{1.330444in}{1.031145in}}%
\pgfpathlineto{\pgfqpoint{1.332471in}{1.028819in}}%
\pgfpathlineto{\pgfqpoint{1.332749in}{1.029642in}}%
\pgfpathlineto{\pgfqpoint{1.334282in}{1.031717in}}%
\pgfpathlineto{\pgfqpoint{1.335343in}{1.030786in}}%
\pgfpathlineto{\pgfqpoint{1.337417in}{1.028409in}}%
\pgfpathlineto{\pgfqpoint{1.337657in}{1.029124in}}%
\pgfpathlineto{\pgfqpoint{1.339240in}{1.031291in}}%
\pgfpathlineto{\pgfqpoint{1.340344in}{1.030302in}}%
\pgfpathlineto{\pgfqpoint{1.342361in}{1.027993in}}%
\pgfpathlineto{\pgfqpoint{1.342633in}{1.028802in}}%
\pgfpathlineto{\pgfqpoint{1.344166in}{1.030892in}}%
\pgfpathlineto{\pgfqpoint{1.345214in}{1.029989in}}%
\pgfpathlineto{\pgfqpoint{1.347308in}{1.027592in}}%
\pgfpathlineto{\pgfqpoint{1.347573in}{1.028380in}}%
\pgfpathlineto{\pgfqpoint{1.349106in}{1.030496in}}%
\pgfpathlineto{\pgfqpoint{1.350145in}{1.029607in}}%
\pgfpathlineto{\pgfqpoint{1.352251in}{1.027197in}}%
\pgfpathlineto{\pgfqpoint{1.352515in}{1.027986in}}%
\pgfpathlineto{\pgfqpoint{1.354048in}{1.030128in}}%
\pgfpathlineto{\pgfqpoint{1.355097in}{1.029245in}}%
\pgfpathlineto{\pgfqpoint{1.357194in}{1.026849in}}%
\pgfpathlineto{\pgfqpoint{1.357448in}{1.027603in}}%
\pgfpathlineto{\pgfqpoint{1.358981in}{1.029745in}}%
\pgfpathlineto{\pgfqpoint{1.360019in}{1.028867in}}%
\pgfpathlineto{\pgfqpoint{1.362150in}{1.026447in}}%
\pgfpathlineto{\pgfqpoint{1.362405in}{1.027225in}}%
\pgfpathlineto{\pgfqpoint{1.363987in}{1.029381in}}%
\pgfpathlineto{\pgfqpoint{1.365081in}{1.028406in}}%
\pgfpathlineto{\pgfqpoint{1.367094in}{1.026107in}}%
\pgfpathlineto{\pgfqpoint{1.367349in}{1.026863in}}%
\pgfpathlineto{\pgfqpoint{1.368882in}{1.028971in}}%
\pgfpathlineto{\pgfqpoint{1.369920in}{1.028091in}}%
\pgfpathlineto{\pgfqpoint{1.372033in}{1.025681in}}%
\pgfpathlineto{\pgfqpoint{1.372291in}{1.026446in}}%
\pgfpathlineto{\pgfqpoint{1.373824in}{1.028564in}}%
\pgfpathlineto{\pgfqpoint{1.374853in}{1.027698in}}%
\pgfpathlineto{\pgfqpoint{1.376988in}{1.025296in}}%
\pgfpathlineto{\pgfqpoint{1.377216in}{1.026003in}}%
\pgfpathlineto{\pgfqpoint{1.378798in}{1.028214in}}%
\pgfpathlineto{\pgfqpoint{1.379833in}{1.027323in}}%
\pgfpathlineto{\pgfqpoint{1.381926in}{1.024937in}}%
\pgfpathlineto{\pgfqpoint{1.382193in}{1.025733in}}%
\pgfpathlineto{\pgfqpoint{1.383727in}{1.027866in}}%
\pgfpathlineto{\pgfqpoint{1.384765in}{1.026990in}}%
\pgfpathlineto{\pgfqpoint{1.386875in}{1.024595in}}%
\pgfpathlineto{\pgfqpoint{1.387115in}{1.025323in}}%
\pgfpathlineto{\pgfqpoint{1.388698in}{1.027536in}}%
\pgfpathlineto{\pgfqpoint{1.389788in}{1.026598in}}%
\pgfpathlineto{\pgfqpoint{1.391823in}{1.024310in}}%
\pgfpathlineto{\pgfqpoint{1.392064in}{1.025056in}}%
\pgfpathlineto{\pgfqpoint{1.393646in}{1.027229in}}%
\pgfpathlineto{\pgfqpoint{1.394744in}{1.026260in}}%
\pgfpathlineto{\pgfqpoint{1.396762in}{1.023966in}}%
\pgfpathlineto{\pgfqpoint{1.397006in}{1.024697in}}%
\pgfpathlineto{\pgfqpoint{1.398589in}{1.026908in}}%
\pgfpathlineto{\pgfqpoint{1.399664in}{1.025974in}}%
\pgfpathlineto{\pgfqpoint{1.401712in}{1.023647in}}%
\pgfpathlineto{\pgfqpoint{1.401964in}{1.024401in}}%
\pgfpathlineto{\pgfqpoint{1.403546in}{1.026560in}}%
\pgfpathlineto{\pgfqpoint{1.404643in}{1.025586in}}%
\pgfpathlineto{\pgfqpoint{1.406660in}{1.023321in}}%
\pgfpathlineto{\pgfqpoint{1.406914in}{1.024104in}}%
\pgfpathlineto{\pgfqpoint{1.408497in}{1.026259in}}%
\pgfpathlineto{\pgfqpoint{1.409593in}{1.025286in}}%
\pgfpathlineto{\pgfqpoint{1.411598in}{1.023009in}}%
\pgfpathlineto{\pgfqpoint{1.411846in}{1.023747in}}%
\pgfpathlineto{\pgfqpoint{1.413428in}{1.025930in}}%
\pgfpathlineto{\pgfqpoint{1.414514in}{1.024977in}}%
\pgfpathlineto{\pgfqpoint{1.416544in}{1.022675in}}%
\pgfpathlineto{\pgfqpoint{1.416792in}{1.023414in}}%
\pgfpathlineto{\pgfqpoint{1.418374in}{1.025591in}}%
\pgfpathlineto{\pgfqpoint{1.419460in}{1.024636in}}%
\pgfpathlineto{\pgfqpoint{1.421493in}{1.022332in}}%
\pgfpathlineto{\pgfqpoint{1.421733in}{1.023056in}}%
\pgfpathlineto{\pgfqpoint{1.423316in}{1.025279in}}%
\pgfpathlineto{\pgfqpoint{1.424404in}{1.024355in}}%
\pgfpathlineto{\pgfqpoint{1.426428in}{1.022061in}}%
\pgfpathlineto{\pgfqpoint{1.426707in}{1.022884in}}%
\pgfpathlineto{\pgfqpoint{1.428240in}{1.025011in}}%
\pgfpathlineto{\pgfqpoint{1.429278in}{1.024145in}}%
\pgfpathlineto{\pgfqpoint{1.431381in}{1.021764in}}%
\pgfpathlineto{\pgfqpoint{1.431637in}{1.022524in}}%
\pgfpathlineto{\pgfqpoint{1.433219in}{1.024674in}}%
\pgfpathlineto{\pgfqpoint{1.434316in}{1.023706in}}%
\pgfpathlineto{\pgfqpoint{1.436332in}{1.021445in}}%
\pgfpathlineto{\pgfqpoint{1.436588in}{1.022229in}}%
\pgfpathlineto{\pgfqpoint{1.438171in}{1.024377in}}%
\pgfpathlineto{\pgfqpoint{1.439266in}{1.023412in}}%
\pgfpathlineto{\pgfqpoint{1.441270in}{1.021144in}}%
\pgfpathlineto{\pgfqpoint{1.441518in}{1.021879in}}%
\pgfpathlineto{\pgfqpoint{1.443101in}{1.024054in}}%
\pgfpathlineto{\pgfqpoint{1.444189in}{1.023099in}}%
\pgfpathlineto{\pgfqpoint{1.446223in}{1.020828in}}%
\pgfpathlineto{\pgfqpoint{1.446463in}{1.021574in}}%
\pgfpathlineto{\pgfqpoint{1.448046in}{1.023783in}}%
\pgfpathlineto{\pgfqpoint{1.449128in}{1.022855in}}%
\pgfpathlineto{\pgfqpoint{1.451166in}{1.020560in}}%
\pgfpathlineto{\pgfqpoint{1.451406in}{1.021287in}}%
\pgfpathlineto{\pgfqpoint{1.452988in}{1.023510in}}%
\pgfpathlineto{\pgfqpoint{1.454076in}{1.022592in}}%
\pgfpathlineto{\pgfqpoint{1.456114in}{1.020314in}}%
\pgfpathlineto{\pgfqpoint{1.456369in}{1.021099in}}%
\pgfpathlineto{\pgfqpoint{1.457951in}{1.023258in}}%
\pgfpathlineto{\pgfqpoint{1.459047in}{1.022296in}}%
\pgfpathlineto{\pgfqpoint{1.461053in}{1.020032in}}%
\pgfpathlineto{\pgfqpoint{1.461305in}{1.020782in}}%
\pgfpathlineto{\pgfqpoint{1.462887in}{1.022971in}}%
\pgfpathlineto{\pgfqpoint{1.463972in}{1.022031in}}%
\pgfpathlineto{\pgfqpoint{1.465998in}{1.019744in}}%
\pgfpathlineto{\pgfqpoint{1.466259in}{1.020521in}}%
\pgfpathlineto{\pgfqpoint{1.467841in}{1.022689in}}%
\pgfpathlineto{\pgfqpoint{1.468930in}{1.021738in}}%
\pgfpathlineto{\pgfqpoint{1.470943in}{1.019468in}}%
\pgfpathlineto{\pgfqpoint{1.471190in}{1.020204in}}%
\pgfpathlineto{\pgfqpoint{1.472773in}{1.022402in}}%
\pgfpathlineto{\pgfqpoint{1.473851in}{1.021474in}}%
\pgfpathlineto{\pgfqpoint{1.475890in}{1.019175in}}%
\pgfpathlineto{\pgfqpoint{1.476155in}{1.019963in}}%
\pgfpathlineto{\pgfqpoint{1.477737in}{1.022115in}}%
\pgfpathlineto{\pgfqpoint{1.478834in}{1.021156in}}%
\pgfpathlineto{\pgfqpoint{1.480842in}{1.018918in}}%
\pgfpathlineto{\pgfqpoint{1.481099in}{1.019710in}}%
\pgfpathlineto{\pgfqpoint{1.482682in}{1.021858in}}%
\pgfpathlineto{\pgfqpoint{1.483771in}{1.020901in}}%
\pgfpathlineto{\pgfqpoint{1.485784in}{1.018634in}}%
\pgfpathlineto{\pgfqpoint{1.486024in}{1.019357in}}%
\pgfpathlineto{\pgfqpoint{1.487607in}{1.021597in}}%
\pgfpathlineto{\pgfqpoint{1.488690in}{1.020696in}}%
\pgfpathlineto{\pgfqpoint{1.490721in}{1.018410in}}%
\pgfpathlineto{\pgfqpoint{1.490995in}{1.019227in}}%
\pgfpathlineto{\pgfqpoint{1.492578in}{1.021368in}}%
\pgfpathlineto{\pgfqpoint{1.493681in}{1.020393in}}%
\pgfpathlineto{\pgfqpoint{1.495677in}{1.018168in}}%
\pgfpathlineto{\pgfqpoint{1.495938in}{1.018972in}}%
\pgfpathlineto{\pgfqpoint{1.497521in}{1.021156in}}%
\pgfpathlineto{\pgfqpoint{1.498605in}{1.020225in}}%
\pgfpathlineto{\pgfqpoint{1.500623in}{1.017979in}}%
\pgfpathlineto{\pgfqpoint{1.500887in}{1.018790in}}%
\pgfpathlineto{\pgfqpoint{1.502470in}{1.020933in}}%
\pgfpathlineto{\pgfqpoint{1.503576in}{1.019962in}}%
\pgfpathlineto{\pgfqpoint{1.505563in}{1.017728in}}%
\pgfpathlineto{\pgfqpoint{1.505832in}{1.018529in}}%
\pgfpathlineto{\pgfqpoint{1.507415in}{1.020666in}}%
\pgfpathlineto{\pgfqpoint{1.508504in}{1.019707in}}%
\pgfpathlineto{\pgfqpoint{1.510514in}{1.017470in}}%
\pgfpathlineto{\pgfqpoint{1.510774in}{1.018275in}}%
\pgfpathlineto{\pgfqpoint{1.512357in}{1.020468in}}%
\pgfpathlineto{\pgfqpoint{1.513441in}{1.019542in}}%
\pgfpathlineto{\pgfqpoint{1.515460in}{1.017298in}}%
\pgfpathlineto{\pgfqpoint{1.515722in}{1.018106in}}%
\pgfpathlineto{\pgfqpoint{1.517305in}{1.020253in}}%
\pgfpathlineto{\pgfqpoint{1.518408in}{1.019289in}}%
\pgfpathlineto{\pgfqpoint{1.520405in}{1.017070in}}%
\pgfpathlineto{\pgfqpoint{1.520665in}{1.017869in}}%
\pgfpathlineto{\pgfqpoint{1.522247in}{1.020033in}}%
\pgfpathlineto{\pgfqpoint{1.523336in}{1.019088in}}%
\pgfpathlineto{\pgfqpoint{1.525336in}{1.016842in}}%
\pgfpathlineto{\pgfqpoint{1.525616in}{1.017673in}}%
\pgfpathlineto{\pgfqpoint{1.527199in}{1.019844in}}%
\pgfpathlineto{\pgfqpoint{1.528281in}{1.018910in}}%
\pgfpathlineto{\pgfqpoint{1.530290in}{1.016654in}}%
\pgfpathlineto{\pgfqpoint{1.530560in}{1.017457in}}%
\pgfpathlineto{\pgfqpoint{1.532143in}{1.019596in}}%
\pgfpathlineto{\pgfqpoint{1.533238in}{1.018635in}}%
\pgfpathlineto{\pgfqpoint{1.535239in}{1.016390in}}%
\pgfpathlineto{\pgfqpoint{1.535493in}{1.017155in}}%
\pgfpathlineto{\pgfqpoint{1.537076in}{1.019375in}}%
\pgfpathlineto{\pgfqpoint{1.538148in}{1.018472in}}%
\pgfpathlineto{\pgfqpoint{1.540181in}{1.016192in}}%
\pgfpathlineto{\pgfqpoint{1.540425in}{1.016921in}}%
\pgfpathlineto{\pgfqpoint{1.542007in}{1.019155in}}%
\pgfpathlineto{\pgfqpoint{1.543087in}{1.018259in}}%
\pgfpathlineto{\pgfqpoint{1.545124in}{1.015976in}}%
\pgfpathlineto{\pgfqpoint{1.545388in}{1.016760in}}%
\pgfpathlineto{\pgfqpoint{1.546971in}{1.018911in}}%
\pgfpathlineto{\pgfqpoint{1.548069in}{1.017959in}}%
\pgfpathlineto{\pgfqpoint{1.550078in}{1.015732in}}%
\pgfpathlineto{\pgfqpoint{1.550332in}{1.016516in}}%
\pgfpathlineto{\pgfqpoint{1.551915in}{1.018709in}}%
\pgfpathlineto{\pgfqpoint{1.552997in}{1.017789in}}%
\pgfpathlineto{\pgfqpoint{1.555023in}{1.015540in}}%
\pgfpathlineto{\pgfqpoint{1.555277in}{1.016319in}}%
\pgfpathlineto{\pgfqpoint{1.556859in}{1.018522in}}%
\pgfpathlineto{\pgfqpoint{1.557938in}{1.017608in}}%
\pgfpathlineto{\pgfqpoint{1.559969in}{1.015359in}}%
\pgfpathlineto{\pgfqpoint{1.560223in}{1.016146in}}%
\pgfpathlineto{\pgfqpoint{1.561806in}{1.018349in}}%
\pgfpathlineto{\pgfqpoint{1.562883in}{1.017435in}}%
\pgfpathlineto{\pgfqpoint{1.564912in}{1.015173in}}%
\pgfpathlineto{\pgfqpoint{1.565151in}{1.015902in}}%
\pgfpathlineto{\pgfqpoint{1.566734in}{1.018161in}}%
\pgfpathlineto{\pgfqpoint{1.567808in}{1.017287in}}%
\pgfpathlineto{\pgfqpoint{1.569852in}{1.014999in}}%
\pgfpathlineto{\pgfqpoint{1.570113in}{1.015775in}}%
\pgfpathlineto{\pgfqpoint{1.571695in}{1.017954in}}%
\pgfpathlineto{\pgfqpoint{1.572783in}{1.017024in}}%
\pgfpathlineto{\pgfqpoint{1.574792in}{1.014776in}}%
\pgfpathlineto{\pgfqpoint{1.575054in}{1.015555in}}%
\pgfpathlineto{\pgfqpoint{1.576637in}{1.017775in}}%
\pgfpathlineto{\pgfqpoint{1.577705in}{1.016877in}}%
\pgfpathlineto{\pgfqpoint{1.579736in}{1.014605in}}%
\pgfpathlineto{\pgfqpoint{1.580002in}{1.015401in}}%
\pgfpathlineto{\pgfqpoint{1.581585in}{1.017647in}}%
\pgfpathlineto{\pgfqpoint{1.582646in}{1.016769in}}%
\pgfpathlineto{\pgfqpoint{1.584696in}{1.014501in}}%
\pgfpathlineto{\pgfqpoint{1.584955in}{1.015298in}}%
\pgfpathlineto{\pgfqpoint{1.586537in}{1.017472in}}%
\pgfpathlineto{\pgfqpoint{1.587626in}{1.016538in}}%
\pgfpathlineto{\pgfqpoint{1.589634in}{1.014293in}}%
\pgfpathlineto{\pgfqpoint{1.589897in}{1.015080in}}%
\pgfpathlineto{\pgfqpoint{1.591479in}{1.017299in}}%
\pgfpathlineto{\pgfqpoint{1.592551in}{1.016403in}}%
\pgfpathlineto{\pgfqpoint{1.594580in}{1.014136in}}%
\pgfpathlineto{\pgfqpoint{1.594829in}{1.014882in}}%
\pgfpathlineto{\pgfqpoint{1.596412in}{1.017112in}}%
\pgfpathlineto{\pgfqpoint{1.597479in}{1.016224in}}%
\pgfpathlineto{\pgfqpoint{1.599532in}{1.013955in}}%
\pgfpathlineto{\pgfqpoint{1.599787in}{1.014745in}}%
\pgfpathlineto{\pgfqpoint{1.601369in}{1.016977in}}%
\pgfpathlineto{\pgfqpoint{1.602435in}{1.016092in}}%
\pgfpathlineto{\pgfqpoint{1.604478in}{1.013836in}}%
\pgfpathlineto{\pgfqpoint{1.604739in}{1.014642in}}%
\pgfpathlineto{\pgfqpoint{1.606321in}{1.016834in}}%
\pgfpathlineto{\pgfqpoint{1.607399in}{1.015921in}}%
\pgfpathlineto{\pgfqpoint{1.609415in}{1.013670in}}%
\pgfpathlineto{\pgfqpoint{1.609681in}{1.014468in}}%
\pgfpathlineto{\pgfqpoint{1.611264in}{1.016668in}}%
\pgfpathlineto{\pgfqpoint{1.612344in}{1.015761in}}%
\pgfpathlineto{\pgfqpoint{1.614369in}{1.013525in}}%
\pgfpathlineto{\pgfqpoint{1.614626in}{1.014319in}}%
\pgfpathlineto{\pgfqpoint{1.616209in}{1.016500in}}%
\pgfpathlineto{\pgfqpoint{1.617293in}{1.015581in}}%
\pgfpathlineto{\pgfqpoint{1.619307in}{1.013334in}}%
\pgfpathlineto{\pgfqpoint{1.619570in}{1.014117in}}%
\pgfpathlineto{\pgfqpoint{1.621152in}{1.016310in}}%
\pgfpathlineto{\pgfqpoint{1.622233in}{1.015396in}}%
\pgfpathlineto{\pgfqpoint{1.624254in}{1.013142in}}%
\pgfpathlineto{\pgfqpoint{1.624498in}{1.013871in}}%
\pgfpathlineto{\pgfqpoint{1.626081in}{1.016118in}}%
\pgfpathlineto{\pgfqpoint{1.627119in}{1.015275in}}%
\pgfpathlineto{\pgfqpoint{1.629206in}{1.012973in}}%
\pgfpathlineto{\pgfqpoint{1.629459in}{1.013758in}}%
\pgfpathlineto{\pgfqpoint{1.631041in}{1.016002in}}%
\pgfpathlineto{\pgfqpoint{1.632105in}{1.015127in}}%
\pgfpathlineto{\pgfqpoint{1.634142in}{1.012856in}}%
\pgfpathlineto{\pgfqpoint{1.634420in}{1.013683in}}%
\pgfpathlineto{\pgfqpoint{1.636003in}{1.015839in}}%
\pgfpathlineto{\pgfqpoint{1.637104in}{1.014895in}}%
\pgfpathlineto{\pgfqpoint{1.639096in}{1.012700in}}%
\pgfpathlineto{\pgfqpoint{1.639350in}{1.013486in}}%
\pgfpathlineto{\pgfqpoint{1.640933in}{1.015701in}}%
\pgfpathlineto{\pgfqpoint{1.642006in}{1.014805in}}%
\pgfpathlineto{\pgfqpoint{1.644041in}{1.012561in}}%
\pgfpathlineto{\pgfqpoint{1.644300in}{1.013362in}}%
\pgfpathlineto{\pgfqpoint{1.645882in}{1.015579in}}%
\pgfpathlineto{\pgfqpoint{1.646957in}{1.014686in}}%
\pgfpathlineto{\pgfqpoint{1.648979in}{1.012434in}}%
\pgfpathlineto{\pgfqpoint{1.649246in}{1.013230in}}%
\pgfpathlineto{\pgfqpoint{1.650829in}{1.015407in}}%
\pgfpathlineto{\pgfqpoint{1.651915in}{1.014488in}}%
\pgfpathlineto{\pgfqpoint{1.653932in}{1.012266in}}%
\pgfpathlineto{\pgfqpoint{1.654187in}{1.013050in}}%
\pgfpathlineto{\pgfqpoint{1.655769in}{1.015258in}}%
\pgfpathlineto{\pgfqpoint{1.656845in}{1.014359in}}%
\pgfpathlineto{\pgfqpoint{1.658878in}{1.012119in}}%
\pgfpathlineto{\pgfqpoint{1.659131in}{1.012903in}}%
\pgfpathlineto{\pgfqpoint{1.660714in}{1.015141in}}%
\pgfpathlineto{\pgfqpoint{1.661779in}{1.014265in}}%
\pgfpathlineto{\pgfqpoint{1.663815in}{1.011998in}}%
\pgfpathlineto{\pgfqpoint{1.664089in}{1.012815in}}%
\pgfpathlineto{\pgfqpoint{1.665672in}{1.014992in}}%
\pgfpathlineto{\pgfqpoint{1.666762in}{1.014070in}}%
\pgfpathlineto{\pgfqpoint{1.668761in}{1.011845in}}%
\pgfpathlineto{\pgfqpoint{1.669027in}{1.012638in}}%
\pgfpathlineto{\pgfqpoint{1.670610in}{1.014829in}}%
\pgfpathlineto{\pgfqpoint{1.671695in}{1.013917in}}%
\pgfpathlineto{\pgfqpoint{1.673714in}{1.011695in}}%
\pgfpathlineto{\pgfqpoint{1.673973in}{1.012492in}}%
\pgfpathlineto{\pgfqpoint{1.675555in}{1.014688in}}%
\pgfpathlineto{\pgfqpoint{1.676633in}{1.013782in}}%
\pgfpathlineto{\pgfqpoint{1.678651in}{1.011538in}}%
\pgfpathlineto{\pgfqpoint{1.678919in}{1.012342in}}%
\pgfpathlineto{\pgfqpoint{1.680502in}{1.014564in}}%
\pgfpathlineto{\pgfqpoint{1.681574in}{1.013679in}}%
\pgfpathlineto{\pgfqpoint{1.683597in}{1.011429in}}%
\pgfpathlineto{\pgfqpoint{1.683864in}{1.012222in}}%
\pgfpathlineto{\pgfqpoint{1.685446in}{1.014405in}}%
\pgfpathlineto{\pgfqpoint{1.686531in}{1.013493in}}%
\pgfpathlineto{\pgfqpoint{1.688551in}{1.011273in}}%
\pgfpathlineto{\pgfqpoint{1.688810in}{1.012072in}}%
\pgfpathlineto{\pgfqpoint{1.690392in}{1.014269in}}%
\pgfpathlineto{\pgfqpoint{1.691469in}{1.013366in}}%
\pgfpathlineto{\pgfqpoint{1.693484in}{1.011127in}}%
\pgfpathlineto{\pgfqpoint{1.693744in}{1.011910in}}%
\pgfpathlineto{\pgfqpoint{1.695326in}{1.014172in}}%
\pgfpathlineto{\pgfqpoint{1.696351in}{1.013350in}}%
\pgfpathlineto{\pgfqpoint{1.698437in}{1.011032in}}%
\pgfpathlineto{\pgfqpoint{1.698680in}{1.011759in}}%
\pgfpathlineto{\pgfqpoint{1.700262in}{1.014020in}}%
\pgfpathlineto{\pgfqpoint{1.701335in}{1.013162in}}%
\pgfpathlineto{\pgfqpoint{1.703383in}{1.010887in}}%
\pgfpathlineto{\pgfqpoint{1.703650in}{1.011685in}}%
\pgfpathlineto{\pgfqpoint{1.705233in}{1.013859in}}%
\pgfpathlineto{\pgfqpoint{1.706318in}{1.012947in}}%
\pgfpathlineto{\pgfqpoint{1.708320in}{1.010723in}}%
\pgfpathlineto{\pgfqpoint{1.708556in}{1.011431in}}%
\pgfpathlineto{\pgfqpoint{1.710138in}{1.013748in}}%
\pgfpathlineto{\pgfqpoint{1.711177in}{1.012933in}}%
\pgfpathlineto{\pgfqpoint{1.713276in}{1.010620in}}%
\pgfpathlineto{\pgfqpoint{1.713514in}{1.011353in}}%
\pgfpathlineto{\pgfqpoint{1.715096in}{1.013623in}}%
\pgfpathlineto{\pgfqpoint{1.716127in}{1.012814in}}%
\pgfpathlineto{\pgfqpoint{1.718215in}{1.010496in}}%
\pgfpathlineto{\pgfqpoint{1.718482in}{1.011287in}}%
\pgfpathlineto{\pgfqpoint{1.720065in}{1.013457in}}%
\pgfpathlineto{\pgfqpoint{1.721152in}{1.012537in}}%
\pgfpathlineto{\pgfqpoint{1.723165in}{1.010303in}}%
\pgfpathlineto{\pgfqpoint{1.723442in}{1.011135in}}%
\pgfpathlineto{\pgfqpoint{1.725025in}{1.013334in}}%
\pgfpathlineto{\pgfqpoint{1.726100in}{1.012446in}}%
\pgfpathlineto{\pgfqpoint{1.728107in}{1.010218in}}%
\pgfpathlineto{\pgfqpoint{1.728365in}{1.010985in}}%
\pgfpathlineto{\pgfqpoint{1.729948in}{1.013192in}}%
\pgfpathlineto{\pgfqpoint{1.731021in}{1.012298in}}%
\pgfpathlineto{\pgfqpoint{1.733054in}{1.010044in}}%
\pgfpathlineto{\pgfqpoint{1.733323in}{1.010853in}}%
\pgfpathlineto{\pgfqpoint{1.734906in}{1.013069in}}%
\pgfpathlineto{\pgfqpoint{1.735978in}{1.012181in}}%
\pgfpathlineto{\pgfqpoint{1.737999in}{1.009940in}}%
\pgfpathlineto{\pgfqpoint{1.738243in}{1.010668in}}%
\pgfpathlineto{\pgfqpoint{1.739826in}{1.012921in}}%
\pgfpathlineto{\pgfqpoint{1.740898in}{1.012060in}}%
\pgfpathlineto{\pgfqpoint{1.742948in}{1.009804in}}%
\pgfpathlineto{\pgfqpoint{1.743187in}{1.010530in}}%
\pgfpathlineto{\pgfqpoint{1.744769in}{1.012771in}}%
\pgfpathlineto{\pgfqpoint{1.745846in}{1.011903in}}%
\pgfpathlineto{\pgfqpoint{1.747888in}{1.009639in}}%
\pgfpathlineto{\pgfqpoint{1.748167in}{1.010468in}}%
\pgfpathlineto{\pgfqpoint{1.749750in}{1.012625in}}%
\pgfpathlineto{\pgfqpoint{1.750838in}{1.011706in}}%
\pgfpathlineto{\pgfqpoint{1.752834in}{1.009494in}}%
\pgfpathlineto{\pgfqpoint{1.753112in}{1.010323in}}%
\pgfpathlineto{\pgfqpoint{1.754695in}{1.012497in}}%
\pgfpathlineto{\pgfqpoint{1.755778in}{1.011589in}}%
\pgfpathlineto{\pgfqpoint{1.757779in}{1.009373in}}%
\pgfpathlineto{\pgfqpoint{1.758047in}{1.010172in}}%
\pgfpathlineto{\pgfqpoint{1.759629in}{1.012377in}}%
\pgfpathlineto{\pgfqpoint{1.760707in}{1.011486in}}%
\pgfpathlineto{\pgfqpoint{1.762724in}{1.009251in}}%
\pgfpathlineto{\pgfqpoint{1.762993in}{1.010051in}}%
\pgfpathlineto{\pgfqpoint{1.764575in}{1.012250in}}%
\pgfpathlineto{\pgfqpoint{1.765655in}{1.011355in}}%
\pgfpathlineto{\pgfqpoint{1.767670in}{1.009123in}}%
\pgfpathlineto{\pgfqpoint{1.767949in}{1.009953in}}%
\pgfpathlineto{\pgfqpoint{1.769532in}{1.012117in}}%
\pgfpathlineto{\pgfqpoint{1.770628in}{1.011193in}}%
\pgfpathlineto{\pgfqpoint{1.772617in}{1.008992in}}%
\pgfpathlineto{\pgfqpoint{1.772857in}{1.009712in}}%
\pgfpathlineto{\pgfqpoint{1.774439in}{1.011998in}}%
\pgfpathlineto{\pgfqpoint{1.775468in}{1.011197in}}%
\pgfpathlineto{\pgfqpoint{1.777565in}{1.008875in}}%
\pgfpathlineto{\pgfqpoint{1.777806in}{1.009600in}}%
\pgfpathlineto{\pgfqpoint{1.779389in}{1.011875in}}%
\pgfpathlineto{\pgfqpoint{1.780419in}{1.011059in}}%
\pgfpathlineto{\pgfqpoint{1.782502in}{1.008753in}}%
\pgfpathlineto{\pgfqpoint{1.782732in}{1.009447in}}%
\pgfpathlineto{\pgfqpoint{1.784315in}{1.011803in}}%
\pgfpathlineto{\pgfqpoint{1.785304in}{1.011064in}}%
\pgfpathlineto{\pgfqpoint{1.787454in}{1.008686in}}%
\pgfpathlineto{\pgfqpoint{1.787724in}{1.009492in}}%
\pgfpathlineto{\pgfqpoint{1.789307in}{1.011676in}}%
\pgfpathlineto{\pgfqpoint{1.790385in}{1.010775in}}%
\pgfpathlineto{\pgfqpoint{1.792405in}{1.008564in}}%
\pgfpathlineto{\pgfqpoint{1.792658in}{1.009345in}}%
\pgfpathlineto{\pgfqpoint{1.794241in}{1.011588in}}%
\pgfpathlineto{\pgfqpoint{1.795305in}{1.010726in}}%
\pgfpathlineto{\pgfqpoint{1.797343in}{1.008472in}}%
\pgfpathlineto{\pgfqpoint{1.797622in}{1.009303in}}%
\pgfpathlineto{\pgfqpoint{1.799205in}{1.011479in}}%
\pgfpathlineto{\pgfqpoint{1.800294in}{1.010569in}}%
\pgfpathlineto{\pgfqpoint{1.802289in}{1.008363in}}%
\pgfpathlineto{\pgfqpoint{1.802551in}{1.009148in}}%
\pgfpathlineto{\pgfqpoint{1.804134in}{1.011374in}}%
\pgfpathlineto{\pgfqpoint{1.805202in}{1.010502in}}%
\pgfpathlineto{\pgfqpoint{1.807234in}{1.008256in}}%
\pgfpathlineto{\pgfqpoint{1.807497in}{1.009043in}}%
\pgfpathlineto{\pgfqpoint{1.809080in}{1.011277in}}%
\pgfpathlineto{\pgfqpoint{1.810144in}{1.010413in}}%
\pgfpathlineto{\pgfqpoint{1.812185in}{1.008183in}}%
\pgfpathlineto{\pgfqpoint{1.812412in}{1.008886in}}%
\pgfpathlineto{\pgfqpoint{1.813994in}{1.011208in}}%
\pgfpathlineto{\pgfqpoint{1.815033in}{1.010409in}}%
\pgfpathlineto{\pgfqpoint{1.817118in}{1.008105in}}%
\pgfpathlineto{\pgfqpoint{1.817384in}{1.008898in}}%
\pgfpathlineto{\pgfqpoint{1.818967in}{1.011146in}}%
\pgfpathlineto{\pgfqpoint{1.820026in}{1.010290in}}%
\pgfpathlineto{\pgfqpoint{1.822070in}{1.008032in}}%
\pgfpathlineto{\pgfqpoint{1.822339in}{1.008835in}}%
\pgfpathlineto{\pgfqpoint{1.823921in}{1.011049in}}%
\pgfpathlineto{\pgfqpoint{1.824998in}{1.010166in}}%
\pgfpathlineto{\pgfqpoint{1.827017in}{1.007935in}}%
\pgfpathlineto{\pgfqpoint{1.827289in}{1.008744in}}%
\pgfpathlineto{\pgfqpoint{1.828871in}{1.010935in}}%
\pgfpathlineto{\pgfqpoint{1.829951in}{1.010043in}}%
\pgfpathlineto{\pgfqpoint{1.831961in}{1.007822in}}%
\pgfpathlineto{\pgfqpoint{1.832219in}{1.008594in}}%
\pgfpathlineto{\pgfqpoint{1.833802in}{1.010834in}}%
\pgfpathlineto{\pgfqpoint{1.834864in}{1.009975in}}%
\pgfpathlineto{\pgfqpoint{1.836906in}{1.007719in}}%
\pgfpathlineto{\pgfqpoint{1.837175in}{1.008523in}}%
\pgfpathlineto{\pgfqpoint{1.838757in}{1.010751in}}%
\pgfpathlineto{\pgfqpoint{1.839826in}{1.009883in}}%
\pgfpathlineto{\pgfqpoint{1.841853in}{1.007644in}}%
\pgfpathlineto{\pgfqpoint{1.842117in}{1.008431in}}%
\pgfpathlineto{\pgfqpoint{1.843699in}{1.010656in}}%
\pgfpathlineto{\pgfqpoint{1.844767in}{1.009783in}}%
\pgfpathlineto{\pgfqpoint{1.846799in}{1.007539in}}%
\pgfpathlineto{\pgfqpoint{1.847070in}{1.008350in}}%
\pgfpathlineto{\pgfqpoint{1.848653in}{1.010556in}}%
\pgfpathlineto{\pgfqpoint{1.849727in}{1.009671in}}%
\pgfpathlineto{\pgfqpoint{1.851745in}{1.007444in}}%
\pgfpathlineto{\pgfqpoint{1.851988in}{1.008172in}}%
\pgfpathlineto{\pgfqpoint{1.853570in}{1.010442in}}%
\pgfpathlineto{\pgfqpoint{1.854599in}{1.009643in}}%
\pgfpathlineto{\pgfqpoint{1.856693in}{1.007340in}}%
\pgfpathlineto{\pgfqpoint{1.856933in}{1.008064in}}%
\pgfpathlineto{\pgfqpoint{1.858515in}{1.010322in}}%
\pgfpathlineto{\pgfqpoint{1.859589in}{1.009472in}}%
\pgfpathlineto{\pgfqpoint{1.861635in}{1.007215in}}%
\pgfpathlineto{\pgfqpoint{1.861898in}{1.008002in}}%
\pgfpathlineto{\pgfqpoint{1.863481in}{1.010225in}}%
\pgfpathlineto{\pgfqpoint{1.864549in}{1.009353in}}%
\pgfpathlineto{\pgfqpoint{1.866581in}{1.007111in}}%
\pgfpathlineto{\pgfqpoint{1.866853in}{1.007923in}}%
\pgfpathlineto{\pgfqpoint{1.868435in}{1.010131in}}%
\pgfpathlineto{\pgfqpoint{1.869509in}{1.009249in}}%
\pgfpathlineto{\pgfqpoint{1.871524in}{1.007026in}}%
\pgfpathlineto{\pgfqpoint{1.871790in}{1.007820in}}%
\pgfpathlineto{\pgfqpoint{1.873372in}{1.010034in}}%
\pgfpathlineto{\pgfqpoint{1.874448in}{1.009155in}}%
\pgfpathlineto{\pgfqpoint{1.876470in}{1.006925in}}%
\pgfpathlineto{\pgfqpoint{1.876738in}{1.007726in}}%
\pgfpathlineto{\pgfqpoint{1.878321in}{1.009944in}}%
\pgfpathlineto{\pgfqpoint{1.879394in}{1.009069in}}%
\pgfpathlineto{\pgfqpoint{1.881417in}{1.006838in}}%
\pgfpathlineto{\pgfqpoint{1.881689in}{1.007651in}}%
\pgfpathlineto{\pgfqpoint{1.883272in}{1.009857in}}%
\pgfpathlineto{\pgfqpoint{1.884345in}{1.008975in}}%
\pgfpathlineto{\pgfqpoint{1.886360in}{1.006754in}}%
\pgfpathlineto{\pgfqpoint{1.886630in}{1.007556in}}%
\pgfpathlineto{\pgfqpoint{1.888212in}{1.009762in}}%
\pgfpathlineto{\pgfqpoint{1.889291in}{1.008878in}}%
\pgfpathlineto{\pgfqpoint{1.891307in}{1.006655in}}%
\pgfpathlineto{\pgfqpoint{1.891570in}{1.007441in}}%
\pgfpathlineto{\pgfqpoint{1.893153in}{1.009673in}}%
\pgfpathlineto{\pgfqpoint{1.894220in}{1.008810in}}%
\pgfpathlineto{\pgfqpoint{1.896260in}{1.006585in}}%
\pgfpathlineto{\pgfqpoint{1.896513in}{1.007367in}}%
\pgfpathlineto{\pgfqpoint{1.898095in}{1.009629in}}%
\pgfpathlineto{\pgfqpoint{1.899153in}{1.008788in}}%
\pgfpathlineto{\pgfqpoint{1.901205in}{1.006551in}}%
\pgfpathlineto{\pgfqpoint{1.901464in}{1.007349in}}%
\pgfpathlineto{\pgfqpoint{1.903046in}{1.009567in}}%
\pgfpathlineto{\pgfqpoint{1.904119in}{1.008694in}}%
\pgfpathlineto{\pgfqpoint{1.906146in}{1.006460in}}%
\pgfpathlineto{\pgfqpoint{1.906417in}{1.007270in}}%
\pgfpathlineto{\pgfqpoint{1.907999in}{1.009472in}}%
\pgfpathlineto{\pgfqpoint{1.909079in}{1.008590in}}%
\pgfpathlineto{\pgfqpoint{1.911091in}{1.006374in}}%
\pgfpathlineto{\pgfqpoint{1.911364in}{1.007189in}}%
\pgfpathlineto{\pgfqpoint{1.912947in}{1.009382in}}%
\pgfpathlineto{\pgfqpoint{1.914026in}{1.008496in}}%
\pgfpathlineto{\pgfqpoint{1.916034in}{1.006283in}}%
\pgfpathlineto{\pgfqpoint{1.916290in}{1.007049in}}%
\pgfpathlineto{\pgfqpoint{1.917873in}{1.009308in}}%
\pgfpathlineto{\pgfqpoint{1.918897in}{1.008498in}}%
\pgfpathlineto{\pgfqpoint{1.920979in}{1.006205in}}%
\pgfpathlineto{\pgfqpoint{1.921248in}{1.007010in}}%
\pgfpathlineto{\pgfqpoint{1.922830in}{1.009253in}}%
\pgfpathlineto{\pgfqpoint{1.923893in}{1.008400in}}%
\pgfpathlineto{\pgfqpoint{1.925933in}{1.006186in}}%
\pgfpathlineto{\pgfqpoint{1.926153in}{1.006872in}}%
\pgfpathlineto{\pgfqpoint{1.927736in}{1.009209in}}%
\pgfpathlineto{\pgfqpoint{1.928725in}{1.008469in}}%
\pgfpathlineto{\pgfqpoint{1.930871in}{1.006106in}}%
\pgfpathlineto{\pgfqpoint{1.931135in}{1.006895in}}%
\pgfpathlineto{\pgfqpoint{1.932717in}{1.009139in}}%
\pgfpathlineto{\pgfqpoint{1.933781in}{1.008286in}}%
\pgfpathlineto{\pgfqpoint{1.935823in}{1.006061in}}%
\pgfpathlineto{\pgfqpoint{1.936081in}{1.006860in}}%
\pgfpathlineto{\pgfqpoint{1.937664in}{1.009098in}}%
\pgfpathlineto{\pgfqpoint{1.938729in}{1.008241in}}%
\pgfpathlineto{\pgfqpoint{1.940762in}{1.006003in}}%
\pgfpathlineto{\pgfqpoint{1.941030in}{1.006802in}}%
\pgfpathlineto{\pgfqpoint{1.942612in}{1.009024in}}%
\pgfpathlineto{\pgfqpoint{1.943682in}{1.008153in}}%
\pgfpathlineto{\pgfqpoint{1.945706in}{1.005926in}}%
\pgfpathlineto{\pgfqpoint{1.945974in}{1.006728in}}%
\pgfpathlineto{\pgfqpoint{1.947557in}{1.008945in}}%
\pgfpathlineto{\pgfqpoint{1.948631in}{1.008073in}}%
\pgfpathlineto{\pgfqpoint{1.950658in}{1.005843in}}%
\pgfpathlineto{\pgfqpoint{1.950910in}{1.006604in}}%
\pgfpathlineto{\pgfqpoint{1.952493in}{1.008863in}}%
\pgfpathlineto{\pgfqpoint{1.953553in}{1.008020in}}%
\pgfpathlineto{\pgfqpoint{1.955600in}{1.005767in}}%
\pgfpathlineto{\pgfqpoint{1.955869in}{1.006573in}}%
\pgfpathlineto{\pgfqpoint{1.957452in}{1.008783in}}%
\pgfpathlineto{\pgfqpoint{1.958531in}{1.007906in}}%
\pgfpathlineto{\pgfqpoint{1.960551in}{1.005707in}}%
\pgfpathlineto{\pgfqpoint{1.960803in}{1.006487in}}%
\pgfpathlineto{\pgfqpoint{1.962386in}{1.008741in}}%
\pgfpathlineto{\pgfqpoint{1.963446in}{1.007897in}}%
\pgfpathlineto{\pgfqpoint{1.965496in}{1.005666in}}%
\pgfpathlineto{\pgfqpoint{1.965754in}{1.006463in}}%
\pgfpathlineto{\pgfqpoint{1.967336in}{1.008694in}}%
\pgfpathlineto{\pgfqpoint{1.968404in}{1.007833in}}%
\pgfpathlineto{\pgfqpoint{1.970435in}{1.005598in}}%
\pgfpathlineto{\pgfqpoint{1.970708in}{1.006412in}}%
\pgfpathlineto{\pgfqpoint{1.972290in}{1.008626in}}%
\pgfpathlineto{\pgfqpoint{1.973364in}{1.007751in}}%
\pgfpathlineto{\pgfqpoint{1.975387in}{1.005550in}}%
\pgfpathlineto{\pgfqpoint{1.975649in}{1.006358in}}%
\pgfpathlineto{\pgfqpoint{1.977232in}{1.008571in}}%
\pgfpathlineto{\pgfqpoint{1.978306in}{1.007695in}}%
\pgfpathlineto{\pgfqpoint{1.980325in}{1.005475in}}%
\pgfpathlineto{\pgfqpoint{1.980593in}{1.006279in}}%
\pgfpathlineto{\pgfqpoint{1.982176in}{1.008515in}}%
\pgfpathlineto{\pgfqpoint{1.983240in}{1.007659in}}%
\pgfpathlineto{\pgfqpoint{1.985278in}{1.005444in}}%
\pgfpathlineto{\pgfqpoint{1.985531in}{1.006229in}}%
\pgfpathlineto{\pgfqpoint{1.987114in}{1.008483in}}%
\pgfpathlineto{\pgfqpoint{1.988173in}{1.007641in}}%
\pgfpathlineto{\pgfqpoint{1.990221in}{1.005408in}}%
\pgfpathlineto{\pgfqpoint{1.990459in}{1.006140in}}%
\pgfpathlineto{\pgfqpoint{1.992042in}{1.008426in}}%
\pgfpathlineto{\pgfqpoint{1.993072in}{1.007637in}}%
\pgfpathlineto{\pgfqpoint{1.995162in}{1.005339in}}%
\pgfpathlineto{\pgfqpoint{1.995431in}{1.006136in}}%
\pgfpathlineto{\pgfqpoint{1.997013in}{1.008328in}}%
\pgfpathlineto{\pgfqpoint{1.998093in}{1.007440in}}%
\pgfpathlineto{\pgfqpoint{2.000101in}{1.005231in}}%
\pgfpathlineto{\pgfqpoint{2.000349in}{1.005970in}}%
\pgfpathlineto{\pgfqpoint{2.001931in}{1.008285in}}%
\pgfpathlineto{\pgfqpoint{2.002963in}{1.007491in}}%
\pgfpathlineto{\pgfqpoint{2.005052in}{1.005194in}}%
\pgfpathlineto{\pgfqpoint{2.005313in}{1.005979in}}%
\pgfpathlineto{\pgfqpoint{2.006896in}{1.008251in}}%
\pgfpathlineto{\pgfqpoint{2.007952in}{1.007421in}}%
\pgfpathlineto{\pgfqpoint{2.009999in}{1.005170in}}%
\pgfpathlineto{\pgfqpoint{2.010272in}{1.005983in}}%
\pgfpathlineto{\pgfqpoint{2.011855in}{1.008180in}}%
\pgfpathlineto{\pgfqpoint{2.012932in}{1.007295in}}%
\pgfpathlineto{\pgfqpoint{2.014945in}{1.005083in}}%
\pgfpathlineto{\pgfqpoint{2.015188in}{1.005811in}}%
\pgfpathlineto{\pgfqpoint{2.016771in}{1.008097in}}%
\pgfpathlineto{\pgfqpoint{2.017799in}{1.007311in}}%
\pgfpathlineto{\pgfqpoint{2.019897in}{1.005030in}}%
\pgfpathlineto{\pgfqpoint{2.020149in}{1.005807in}}%
\pgfpathlineto{\pgfqpoint{2.021732in}{1.008045in}}%
\pgfpathlineto{\pgfqpoint{2.022797in}{1.007190in}}%
\pgfpathlineto{\pgfqpoint{2.024828in}{1.004958in}}%
\pgfpathlineto{\pgfqpoint{2.025100in}{1.005767in}}%
\pgfpathlineto{\pgfqpoint{2.026683in}{1.008011in}}%
\pgfpathlineto{\pgfqpoint{2.027742in}{1.007161in}}%
\pgfpathlineto{\pgfqpoint{2.029779in}{1.004923in}}%
\pgfpathlineto{\pgfqpoint{2.030057in}{1.005754in}}%
\pgfpathlineto{\pgfqpoint{2.031639in}{1.007970in}}%
\pgfpathlineto{\pgfqpoint{2.032709in}{1.007105in}}%
\pgfpathlineto{\pgfqpoint{2.034730in}{1.004905in}}%
\pgfpathlineto{\pgfqpoint{2.034968in}{1.005638in}}%
\pgfpathlineto{\pgfqpoint{2.036551in}{1.007927in}}%
\pgfpathlineto{\pgfqpoint{2.037581in}{1.007140in}}%
\pgfpathlineto{\pgfqpoint{2.039673in}{1.004843in}}%
\pgfpathlineto{\pgfqpoint{2.039916in}{1.005567in}}%
\pgfpathlineto{\pgfqpoint{2.041499in}{1.007832in}}%
\pgfpathlineto{\pgfqpoint{2.042569in}{1.006994in}}%
\pgfpathlineto{\pgfqpoint{2.044616in}{1.004746in}}%
\pgfpathlineto{\pgfqpoint{2.044877in}{1.005520in}}%
\pgfpathlineto{\pgfqpoint{2.046459in}{1.007749in}}%
\pgfpathlineto{\pgfqpoint{2.047525in}{1.006888in}}%
\pgfpathlineto{\pgfqpoint{2.049557in}{1.004656in}}%
\pgfpathlineto{\pgfqpoint{2.049794in}{1.005370in}}%
\pgfpathlineto{\pgfqpoint{2.051376in}{1.007725in}}%
\pgfpathlineto{\pgfqpoint{2.052365in}{1.006997in}}%
\pgfpathlineto{\pgfqpoint{2.054506in}{1.004646in}}%
\pgfpathlineto{\pgfqpoint{2.054775in}{1.005451in}}%
\pgfpathlineto{\pgfqpoint{2.056358in}{1.007689in}}%
\pgfpathlineto{\pgfqpoint{2.057423in}{1.006836in}}%
\pgfpathlineto{\pgfqpoint{2.059459in}{1.004629in}}%
\pgfpathlineto{\pgfqpoint{2.059697in}{1.005370in}}%
\pgfpathlineto{\pgfqpoint{2.061280in}{1.007665in}}%
\pgfpathlineto{\pgfqpoint{2.062310in}{1.006875in}}%
\pgfpathlineto{\pgfqpoint{2.064404in}{1.004587in}}%
\pgfpathlineto{\pgfqpoint{2.064656in}{1.005351in}}%
\pgfpathlineto{\pgfqpoint{2.066238in}{1.007598in}}%
\pgfpathlineto{\pgfqpoint{2.067301in}{1.006752in}}%
\pgfpathlineto{\pgfqpoint{2.069343in}{1.004510in}}%
\pgfpathlineto{\pgfqpoint{2.069612in}{1.005314in}}%
\pgfpathlineto{\pgfqpoint{2.071195in}{1.007546in}}%
\pgfpathlineto{\pgfqpoint{2.072263in}{1.006688in}}%
\pgfpathlineto{\pgfqpoint{2.074297in}{1.004487in}}%
\pgfpathlineto{\pgfqpoint{2.074522in}{1.005190in}}%
\pgfpathlineto{\pgfqpoint{2.076105in}{1.007527in}}%
\pgfpathlineto{\pgfqpoint{2.077134in}{1.006749in}}%
\pgfpathlineto{\pgfqpoint{2.079237in}{1.004440in}}%
\pgfpathlineto{\pgfqpoint{2.079506in}{1.005248in}}%
\pgfpathlineto{\pgfqpoint{2.081089in}{1.007485in}}%
\pgfpathlineto{\pgfqpoint{2.082158in}{1.006631in}}%
\pgfpathlineto{\pgfqpoint{2.084184in}{1.004419in}}%
\pgfpathlineto{\pgfqpoint{2.084423in}{1.005145in}}%
\pgfpathlineto{\pgfqpoint{2.086006in}{1.007432in}}%
\pgfpathlineto{\pgfqpoint{2.087036in}{1.006649in}}%
\pgfpathlineto{\pgfqpoint{2.089130in}{1.004368in}}%
\pgfpathlineto{\pgfqpoint{2.089369in}{1.005099in}}%
\pgfpathlineto{\pgfqpoint{2.090951in}{1.007379in}}%
\pgfpathlineto{\pgfqpoint{2.091982in}{1.006593in}}%
\pgfpathlineto{\pgfqpoint{2.094078in}{1.004317in}}%
\pgfpathlineto{\pgfqpoint{2.094331in}{1.005099in}}%
\pgfpathlineto{\pgfqpoint{2.095914in}{1.007334in}}%
\pgfpathlineto{\pgfqpoint{2.096981in}{1.006479in}}%
\pgfpathlineto{\pgfqpoint{2.099021in}{1.004256in}}%
\pgfpathlineto{\pgfqpoint{2.099260in}{1.004990in}}%
\pgfpathlineto{\pgfqpoint{2.100842in}{1.007294in}}%
\pgfpathlineto{\pgfqpoint{2.101873in}{1.006517in}}%
\pgfpathlineto{\pgfqpoint{2.103968in}{1.004247in}}%
\pgfpathlineto{\pgfqpoint{2.104206in}{1.004982in}}%
\pgfpathlineto{\pgfqpoint{2.105789in}{1.007251in}}%
\pgfpathlineto{\pgfqpoint{2.106819in}{1.006453in}}%
\pgfpathlineto{\pgfqpoint{2.108914in}{1.004178in}}%
\pgfpathlineto{\pgfqpoint{2.109173in}{1.004978in}}%
\pgfpathlineto{\pgfqpoint{2.110756in}{1.007212in}}%
\pgfpathlineto{\pgfqpoint{2.111822in}{1.006358in}}%
\pgfpathlineto{\pgfqpoint{2.113854in}{1.004129in}}%
\pgfpathlineto{\pgfqpoint{2.114127in}{1.004945in}}%
\pgfpathlineto{\pgfqpoint{2.115709in}{1.007161in}}%
\pgfpathlineto{\pgfqpoint{2.116782in}{1.006292in}}%
\pgfpathlineto{\pgfqpoint{2.118797in}{1.004082in}}%
\pgfpathlineto{\pgfqpoint{2.119074in}{1.004904in}}%
\pgfpathlineto{\pgfqpoint{2.120656in}{1.007099in}}%
\pgfpathlineto{\pgfqpoint{2.121736in}{1.006218in}}%
\pgfpathlineto{\pgfqpoint{2.123744in}{1.004016in}}%
\pgfpathlineto{\pgfqpoint{2.124004in}{1.004793in}}%
\pgfpathlineto{\pgfqpoint{2.125586in}{1.007049in}}%
\pgfpathlineto{\pgfqpoint{2.126643in}{1.006210in}}%
\pgfpathlineto{\pgfqpoint{2.128696in}{1.003983in}}%
\pgfpathlineto{\pgfqpoint{2.128953in}{1.004781in}}%
\pgfpathlineto{\pgfqpoint{2.130536in}{1.007045in}}%
\pgfpathlineto{\pgfqpoint{2.131592in}{1.006213in}}%
\pgfpathlineto{\pgfqpoint{2.133637in}{1.003969in}}%
\pgfpathlineto{\pgfqpoint{2.133906in}{1.004772in}}%
\pgfpathlineto{\pgfqpoint{2.135488in}{1.006987in}}%
\pgfpathlineto{\pgfqpoint{2.136567in}{1.006120in}}%
\pgfpathlineto{\pgfqpoint{2.138582in}{1.003910in}}%
\pgfpathlineto{\pgfqpoint{2.138845in}{1.004694in}}%
\pgfpathlineto{\pgfqpoint{2.140428in}{1.006918in}}%
\pgfpathlineto{\pgfqpoint{2.141499in}{1.006057in}}%
\pgfpathlineto{\pgfqpoint{2.143532in}{1.003851in}}%
\pgfpathlineto{\pgfqpoint{2.143791in}{1.004653in}}%
\pgfpathlineto{\pgfqpoint{2.145374in}{1.006910in}}%
\pgfpathlineto{\pgfqpoint{2.146434in}{1.006073in}}%
\pgfpathlineto{\pgfqpoint{2.148477in}{1.003861in}}%
\pgfpathlineto{\pgfqpoint{2.148717in}{1.004605in}}%
\pgfpathlineto{\pgfqpoint{2.150300in}{1.006887in}}%
\pgfpathlineto{\pgfqpoint{2.151329in}{1.006096in}}%
\pgfpathlineto{\pgfqpoint{2.153417in}{1.003806in}}%
\pgfpathlineto{\pgfqpoint{2.153691in}{1.004622in}}%
\pgfpathlineto{\pgfqpoint{2.155273in}{1.006835in}}%
\pgfpathlineto{\pgfqpoint{2.156346in}{1.005966in}}%
\pgfpathlineto{\pgfqpoint{2.158363in}{1.003755in}}%
\pgfpathlineto{\pgfqpoint{2.158606in}{1.004482in}}%
\pgfpathlineto{\pgfqpoint{2.160189in}{1.006769in}}%
\pgfpathlineto{\pgfqpoint{2.161216in}{1.005987in}}%
\pgfpathlineto{\pgfqpoint{2.163300in}{1.003703in}}%
\pgfpathlineto{\pgfqpoint{2.163573in}{1.004509in}}%
\pgfpathlineto{\pgfqpoint{2.165155in}{1.006735in}}%
\pgfpathlineto{\pgfqpoint{2.166221in}{1.005875in}}%
\pgfpathlineto{\pgfqpoint{2.168259in}{1.003665in}}%
\pgfpathlineto{\pgfqpoint{2.168523in}{1.004483in}}%
\pgfpathlineto{\pgfqpoint{2.170105in}{1.006749in}}%
\pgfpathlineto{\pgfqpoint{2.171163in}{1.005919in}}%
\pgfpathlineto{\pgfqpoint{2.173199in}{1.003686in}}%
\pgfpathlineto{\pgfqpoint{2.173473in}{1.004498in}}%
\pgfpathlineto{\pgfqpoint{2.175055in}{1.006683in}}%
\pgfpathlineto{\pgfqpoint{2.176132in}{1.005799in}}%
\pgfpathlineto{\pgfqpoint{2.178144in}{1.003594in}}%
\pgfpathlineto{\pgfqpoint{2.178387in}{1.004322in}}%
\pgfpathlineto{\pgfqpoint{2.179970in}{1.006627in}}%
\pgfpathlineto{\pgfqpoint{2.180994in}{1.005855in}}%
\pgfpathlineto{\pgfqpoint{2.183092in}{1.003555in}}%
\pgfpathlineto{\pgfqpoint{2.183334in}{1.004277in}}%
\pgfpathlineto{\pgfqpoint{2.184916in}{1.006556in}}%
\pgfpathlineto{\pgfqpoint{2.185946in}{1.005775in}}%
\pgfpathlineto{\pgfqpoint{2.188039in}{1.003496in}}%
\pgfpathlineto{\pgfqpoint{2.188278in}{1.004226in}}%
\pgfpathlineto{\pgfqpoint{2.189861in}{1.006515in}}%
\pgfpathlineto{\pgfqpoint{2.190891in}{1.005735in}}%
\pgfpathlineto{\pgfqpoint{2.192974in}{1.003452in}}%
\pgfpathlineto{\pgfqpoint{2.193246in}{1.004258in}}%
\pgfpathlineto{\pgfqpoint{2.194828in}{1.006486in}}%
\pgfpathlineto{\pgfqpoint{2.195894in}{1.005628in}}%
\pgfpathlineto{\pgfqpoint{2.197924in}{1.003403in}}%
\pgfpathlineto{\pgfqpoint{2.198157in}{1.004106in}}%
\pgfpathlineto{\pgfqpoint{2.199789in}{1.006473in}}%
\pgfpathlineto{\pgfqpoint{2.200861in}{1.005615in}}%
\pgfpathlineto{\pgfqpoint{2.202875in}{1.003417in}}%
\pgfpathlineto{\pgfqpoint{2.203115in}{1.004143in}}%
\pgfpathlineto{\pgfqpoint{2.204697in}{1.006422in}}%
\pgfpathlineto{\pgfqpoint{2.205727in}{1.005640in}}%
\pgfpathlineto{\pgfqpoint{2.207820in}{1.003352in}}%
\pgfpathlineto{\pgfqpoint{2.208060in}{1.004077in}}%
\pgfpathlineto{\pgfqpoint{2.209643in}{1.006367in}}%
\pgfpathlineto{\pgfqpoint{2.210673in}{1.005591in}}%
\pgfpathlineto{\pgfqpoint{2.212766in}{1.003311in}}%
\pgfpathlineto{\pgfqpoint{2.213005in}{1.004038in}}%
\pgfpathlineto{\pgfqpoint{2.214588in}{1.006320in}}%
\pgfpathlineto{\pgfqpoint{2.215618in}{1.005538in}}%
\pgfpathlineto{\pgfqpoint{2.217712in}{1.003256in}}%
\pgfpathlineto{\pgfqpoint{2.217951in}{1.003984in}}%
\pgfpathlineto{\pgfqpoint{2.219533in}{1.006274in}}%
\pgfpathlineto{\pgfqpoint{2.220564in}{1.005495in}}%
\pgfpathlineto{\pgfqpoint{2.222654in}{1.003206in}}%
\pgfpathlineto{\pgfqpoint{2.222897in}{1.003933in}}%
\pgfpathlineto{\pgfqpoint{2.224480in}{1.006218in}}%
\pgfpathlineto{\pgfqpoint{2.225508in}{1.005437in}}%
\pgfpathlineto{\pgfqpoint{2.227592in}{1.003154in}}%
\pgfpathlineto{\pgfqpoint{2.227844in}{1.003902in}}%
\pgfpathlineto{\pgfqpoint{2.229427in}{1.006188in}}%
\pgfpathlineto{\pgfqpoint{2.230456in}{1.005393in}}%
\pgfpathlineto{\pgfqpoint{2.232550in}{1.003123in}}%
\pgfpathlineto{\pgfqpoint{2.232808in}{1.003927in}}%
\pgfpathlineto{\pgfqpoint{2.234391in}{1.006216in}}%
\pgfpathlineto{\pgfqpoint{2.235443in}{1.005403in}}%
\pgfpathlineto{\pgfqpoint{2.237495in}{1.003183in}}%
\pgfpathlineto{\pgfqpoint{2.237734in}{1.003919in}}%
\pgfpathlineto{\pgfqpoint{2.239316in}{1.006186in}}%
\pgfpathlineto{\pgfqpoint{2.240346in}{1.005391in}}%
\pgfpathlineto{\pgfqpoint{2.242434in}{1.003104in}}%
\pgfpathlineto{\pgfqpoint{2.242694in}{1.003881in}}%
\pgfpathlineto{\pgfqpoint{2.244276in}{1.006140in}}%
\pgfpathlineto{\pgfqpoint{2.245335in}{1.005307in}}%
\pgfpathlineto{\pgfqpoint{2.247380in}{1.003068in}}%
\pgfpathlineto{\pgfqpoint{2.247628in}{1.003810in}}%
\pgfpathlineto{\pgfqpoint{2.249210in}{1.006106in}}%
\pgfpathlineto{\pgfqpoint{2.250238in}{1.005315in}}%
\pgfpathlineto{\pgfqpoint{2.252333in}{1.003046in}}%
\pgfpathlineto{\pgfqpoint{2.252589in}{1.003843in}}%
\pgfpathlineto{\pgfqpoint{2.254171in}{1.006126in}}%
\pgfpathlineto{\pgfqpoint{2.255224in}{1.005309in}}%
\pgfpathlineto{\pgfqpoint{2.257275in}{1.003063in}}%
\pgfpathlineto{\pgfqpoint{2.257526in}{1.003813in}}%
\pgfpathlineto{\pgfqpoint{2.259108in}{1.006066in}}%
\pgfpathlineto{\pgfqpoint{2.260169in}{1.005228in}}%
\pgfpathlineto{\pgfqpoint{2.262221in}{1.003004in}}%
\pgfpathlineto{\pgfqpoint{2.262459in}{1.003740in}}%
\pgfpathlineto{\pgfqpoint{2.264042in}{1.006051in}}%
\pgfpathlineto{\pgfqpoint{2.265072in}{1.005277in}}%
\pgfpathlineto{\pgfqpoint{2.267166in}{1.003000in}}%
\pgfpathlineto{\pgfqpoint{2.267405in}{1.003730in}}%
\pgfpathlineto{\pgfqpoint{2.268988in}{1.006013in}}%
\pgfpathlineto{\pgfqpoint{2.270018in}{1.005232in}}%
\pgfpathlineto{\pgfqpoint{2.272112in}{1.002959in}}%
\pgfpathlineto{\pgfqpoint{2.272350in}{1.003689in}}%
\pgfpathlineto{\pgfqpoint{2.273933in}{1.005978in}}%
\pgfpathlineto{\pgfqpoint{2.274963in}{1.005198in}}%
\pgfpathlineto{\pgfqpoint{2.277059in}{1.002933in}}%
\pgfpathlineto{\pgfqpoint{2.277297in}{1.003667in}}%
\pgfpathlineto{\pgfqpoint{2.278879in}{1.005948in}}%
\pgfpathlineto{\pgfqpoint{2.279909in}{1.005158in}}%
\pgfpathlineto{\pgfqpoint{2.281998in}{1.002871in}}%
\pgfpathlineto{\pgfqpoint{2.282261in}{1.003656in}}%
\pgfpathlineto{\pgfqpoint{2.283844in}{1.005894in}}%
\pgfpathlineto{\pgfqpoint{2.284908in}{1.005047in}}%
\pgfpathlineto{\pgfqpoint{2.286944in}{1.002818in}}%
\pgfpathlineto{\pgfqpoint{2.287218in}{1.003636in}}%
\pgfpathlineto{\pgfqpoint{2.288801in}{1.005864in}}%
\pgfpathlineto{\pgfqpoint{2.289870in}{1.005007in}}%
\pgfpathlineto{\pgfqpoint{2.291883in}{1.002804in}}%
\pgfpathlineto{\pgfqpoint{2.292155in}{1.003609in}}%
\pgfpathlineto{\pgfqpoint{2.293737in}{1.005844in}}%
\pgfpathlineto{\pgfqpoint{2.294799in}{1.004995in}}%
\pgfpathlineto{\pgfqpoint{2.296833in}{1.002769in}}%
\pgfpathlineto{\pgfqpoint{2.297066in}{1.003473in}}%
\pgfpathlineto{\pgfqpoint{2.298698in}{1.005836in}}%
\pgfpathlineto{\pgfqpoint{2.299770in}{1.004978in}}%
\pgfpathlineto{\pgfqpoint{2.301785in}{1.002796in}}%
\pgfpathlineto{\pgfqpoint{2.302023in}{1.003528in}}%
\pgfpathlineto{\pgfqpoint{2.303605in}{1.005817in}}%
\pgfpathlineto{\pgfqpoint{2.304636in}{1.005035in}}%
\pgfpathlineto{\pgfqpoint{2.306732in}{1.002770in}}%
\pgfpathlineto{\pgfqpoint{2.306970in}{1.003508in}}%
\pgfpathlineto{\pgfqpoint{2.308552in}{1.005801in}}%
\pgfpathlineto{\pgfqpoint{2.309582in}{1.005016in}}%
\pgfpathlineto{\pgfqpoint{2.311672in}{1.002729in}}%
\pgfpathlineto{\pgfqpoint{2.311946in}{1.003545in}}%
\pgfpathlineto{\pgfqpoint{2.313528in}{1.005747in}}%
\pgfpathlineto{\pgfqpoint{2.314606in}{1.004872in}}%
\pgfpathlineto{\pgfqpoint{2.316617in}{1.002672in}}%
\pgfpathlineto{\pgfqpoint{2.316892in}{1.003490in}}%
\pgfpathlineto{\pgfqpoint{2.318474in}{1.005700in}}%
\pgfpathlineto{\pgfqpoint{2.319549in}{1.004831in}}%
\pgfpathlineto{\pgfqpoint{2.321561in}{1.002629in}}%
\pgfpathlineto{\pgfqpoint{2.321839in}{1.003456in}}%
\pgfpathlineto{\pgfqpoint{2.323421in}{1.005655in}}%
\pgfpathlineto{\pgfqpoint{2.324500in}{1.004781in}}%
\pgfpathlineto{\pgfqpoint{2.326499in}{1.002594in}}%
\pgfpathlineto{\pgfqpoint{2.326731in}{1.003284in}}%
\pgfpathlineto{\pgfqpoint{2.328363in}{1.005658in}}%
\pgfpathlineto{\pgfqpoint{2.329428in}{1.004806in}}%
\pgfpathlineto{\pgfqpoint{2.331454in}{1.002590in}}%
\pgfpathlineto{\pgfqpoint{2.331727in}{1.003407in}}%
\pgfpathlineto{\pgfqpoint{2.333310in}{1.005623in}}%
\pgfpathlineto{\pgfqpoint{2.334383in}{1.004759in}}%
\pgfpathlineto{\pgfqpoint{2.336405in}{1.002571in}}%
\pgfpathlineto{\pgfqpoint{2.336665in}{1.003374in}}%
\pgfpathlineto{\pgfqpoint{2.338248in}{1.005607in}}%
\pgfpathlineto{\pgfqpoint{2.339314in}{1.004758in}}%
\pgfpathlineto{\pgfqpoint{2.341342in}{1.002539in}}%
\pgfpathlineto{\pgfqpoint{2.341614in}{1.003349in}}%
\pgfpathlineto{\pgfqpoint{2.343196in}{1.005578in}}%
\pgfpathlineto{\pgfqpoint{2.344265in}{1.004725in}}%
\pgfpathlineto{\pgfqpoint{2.346293in}{1.002519in}}%
\pgfpathlineto{\pgfqpoint{2.346532in}{1.003248in}}%
\pgfpathlineto{\pgfqpoint{2.348115in}{1.005556in}}%
\pgfpathlineto{\pgfqpoint{2.349145in}{1.004785in}}%
\pgfpathlineto{\pgfqpoint{2.351233in}{1.002501in}}%
\pgfpathlineto{\pgfqpoint{2.351502in}{1.003301in}}%
\pgfpathlineto{\pgfqpoint{2.353085in}{1.005514in}}%
\pgfpathlineto{\pgfqpoint{2.354159in}{1.004649in}}%
\pgfpathlineto{\pgfqpoint{2.356181in}{1.002438in}}%
\pgfpathlineto{\pgfqpoint{2.356455in}{1.003257in}}%
\pgfpathlineto{\pgfqpoint{2.358037in}{1.005486in}}%
\pgfpathlineto{\pgfqpoint{2.359104in}{1.004633in}}%
\pgfpathlineto{\pgfqpoint{2.361130in}{1.002432in}}%
\pgfpathlineto{\pgfqpoint{2.361369in}{1.003161in}}%
\pgfpathlineto{\pgfqpoint{2.362951in}{1.005451in}}%
\pgfpathlineto{\pgfqpoint{2.363982in}{1.004675in}}%
\pgfpathlineto{\pgfqpoint{2.366073in}{1.002388in}}%
\pgfpathlineto{\pgfqpoint{2.366317in}{1.003116in}}%
\pgfpathlineto{\pgfqpoint{2.367899in}{1.005398in}}%
\pgfpathlineto{\pgfqpoint{2.368928in}{1.004619in}}%
\pgfpathlineto{\pgfqpoint{2.371023in}{1.002358in}}%
\pgfpathlineto{\pgfqpoint{2.371260in}{1.003092in}}%
\pgfpathlineto{\pgfqpoint{2.372842in}{1.005385in}}%
\pgfpathlineto{\pgfqpoint{2.373873in}{1.004599in}}%
\pgfpathlineto{\pgfqpoint{2.375963in}{1.002314in}}%
\pgfpathlineto{\pgfqpoint{2.376243in}{1.003146in}}%
\pgfpathlineto{\pgfqpoint{2.377825in}{1.005334in}}%
\pgfpathlineto{\pgfqpoint{2.378903in}{1.004455in}}%
\pgfpathlineto{\pgfqpoint{2.380912in}{1.002277in}}%
\pgfpathlineto{\pgfqpoint{2.381150in}{1.003012in}}%
\pgfpathlineto{\pgfqpoint{2.382733in}{1.005326in}}%
\pgfpathlineto{\pgfqpoint{2.383763in}{1.004557in}}%
\pgfpathlineto{\pgfqpoint{2.385852in}{1.002274in}}%
\pgfpathlineto{\pgfqpoint{2.386131in}{1.003100in}}%
\pgfpathlineto{\pgfqpoint{2.387713in}{1.005277in}}%
\pgfpathlineto{\pgfqpoint{2.388800in}{1.004388in}}%
\pgfpathlineto{\pgfqpoint{2.390798in}{1.002204in}}%
\pgfpathlineto{\pgfqpoint{2.391039in}{1.002928in}}%
\pgfpathlineto{\pgfqpoint{2.392622in}{1.005256in}}%
\pgfpathlineto{\pgfqpoint{2.393653in}{1.004477in}}%
\pgfpathlineto{\pgfqpoint{2.395743in}{1.002192in}}%
\pgfpathlineto{\pgfqpoint{2.396013in}{1.003001in}}%
\pgfpathlineto{\pgfqpoint{2.397595in}{1.005254in}}%
\pgfpathlineto{\pgfqpoint{2.398657in}{1.004418in}}%
\pgfpathlineto{\pgfqpoint{2.400693in}{1.002203in}}%
\pgfpathlineto{\pgfqpoint{2.400933in}{1.002930in}}%
\pgfpathlineto{\pgfqpoint{2.402515in}{1.005226in}}%
\pgfpathlineto{\pgfqpoint{2.403546in}{1.004454in}}%
\pgfpathlineto{\pgfqpoint{2.405641in}{1.002192in}}%
\pgfpathlineto{\pgfqpoint{2.405879in}{1.002930in}}%
\pgfpathlineto{\pgfqpoint{2.407462in}{1.005218in}}%
\pgfpathlineto{\pgfqpoint{2.408492in}{1.004433in}}%
\pgfpathlineto{\pgfqpoint{2.410581in}{1.002149in}}%
\pgfpathlineto{\pgfqpoint{2.410855in}{1.002965in}}%
\pgfpathlineto{\pgfqpoint{2.412438in}{1.005175in}}%
\pgfpathlineto{\pgfqpoint{2.413513in}{1.004307in}}%
\pgfpathlineto{\pgfqpoint{2.415524in}{1.002108in}}%
\pgfpathlineto{\pgfqpoint{2.415800in}{1.002929in}}%
\pgfpathlineto{\pgfqpoint{2.417383in}{1.005138in}}%
\pgfpathlineto{\pgfqpoint{2.418455in}{1.004276in}}%
\pgfpathlineto{\pgfqpoint{2.420477in}{1.002095in}}%
\pgfpathlineto{\pgfqpoint{2.420716in}{1.002838in}}%
\pgfpathlineto{\pgfqpoint{2.422299in}{1.005149in}}%
\pgfpathlineto{\pgfqpoint{2.423328in}{1.004373in}}%
\pgfpathlineto{\pgfqpoint{2.425421in}{1.002100in}}%
\pgfpathlineto{\pgfqpoint{2.425659in}{1.002829in}}%
\pgfpathlineto{\pgfqpoint{2.427242in}{1.005122in}}%
\pgfpathlineto{\pgfqpoint{2.428272in}{1.004347in}}%
\pgfpathlineto{\pgfqpoint{2.430365in}{1.002059in}}%
\pgfpathlineto{\pgfqpoint{2.430606in}{1.002779in}}%
\pgfpathlineto{\pgfqpoint{2.432189in}{1.005063in}}%
\pgfpathlineto{\pgfqpoint{2.433260in}{1.004245in}}%
\pgfpathlineto{\pgfqpoint{2.435314in}{1.002029in}}%
\pgfpathlineto{\pgfqpoint{2.435551in}{1.002766in}}%
\pgfpathlineto{\pgfqpoint{2.437134in}{1.005057in}}%
\pgfpathlineto{\pgfqpoint{2.438164in}{1.004274in}}%
\pgfpathlineto{\pgfqpoint{2.440253in}{1.001991in}}%
\pgfpathlineto{\pgfqpoint{2.440502in}{1.002732in}}%
\pgfpathlineto{\pgfqpoint{2.442084in}{1.005015in}}%
\pgfpathlineto{\pgfqpoint{2.443112in}{1.004225in}}%
\pgfpathlineto{\pgfqpoint{2.445197in}{1.001946in}}%
\pgfpathlineto{\pgfqpoint{2.445431in}{1.002652in}}%
\pgfpathlineto{\pgfqpoint{2.447063in}{1.005019in}}%
\pgfpathlineto{\pgfqpoint{2.448132in}{1.004168in}}%
\pgfpathlineto{\pgfqpoint{2.450150in}{1.001992in}}%
\pgfpathlineto{\pgfqpoint{2.450386in}{1.002723in}}%
\pgfpathlineto{\pgfqpoint{2.451969in}{1.005015in}}%
\pgfpathlineto{\pgfqpoint{2.453000in}{1.004230in}}%
\pgfpathlineto{\pgfqpoint{2.455089in}{1.001947in}}%
\pgfpathlineto{\pgfqpoint{2.455327in}{1.002658in}}%
\pgfpathlineto{\pgfqpoint{2.456909in}{1.004979in}}%
\pgfpathlineto{\pgfqpoint{2.457940in}{1.004211in}}%
\pgfpathlineto{\pgfqpoint{2.460035in}{1.001922in}}%
\pgfpathlineto{\pgfqpoint{2.460305in}{1.002727in}}%
\pgfpathlineto{\pgfqpoint{2.461888in}{1.004946in}}%
\pgfpathlineto{\pgfqpoint{2.462959in}{1.004086in}}%
\pgfpathlineto{\pgfqpoint{2.464987in}{1.001895in}}%
\pgfpathlineto{\pgfqpoint{2.465247in}{1.002698in}}%
\pgfpathlineto{\pgfqpoint{2.466829in}{1.004944in}}%
\pgfpathlineto{\pgfqpoint{2.467894in}{1.004103in}}%
\pgfpathlineto{\pgfqpoint{2.469924in}{1.001884in}}%
\pgfpathlineto{\pgfqpoint{2.470194in}{1.002688in}}%
\pgfpathlineto{\pgfqpoint{2.471777in}{1.004918in}}%
\pgfpathlineto{\pgfqpoint{2.472845in}{1.004068in}}%
\pgfpathlineto{\pgfqpoint{2.474877in}{1.001876in}}%
\pgfpathlineto{\pgfqpoint{2.475118in}{1.002624in}}%
\pgfpathlineto{\pgfqpoint{2.476700in}{1.004933in}}%
\pgfpathlineto{\pgfqpoint{2.477729in}{1.004158in}}%
\pgfpathlineto{\pgfqpoint{2.479822in}{1.001898in}}%
\pgfpathlineto{\pgfqpoint{2.480064in}{1.002646in}}%
\pgfpathlineto{\pgfqpoint{2.481646in}{1.004935in}}%
\pgfpathlineto{\pgfqpoint{2.482675in}{1.004154in}}%
\pgfpathlineto{\pgfqpoint{2.484767in}{1.001893in}}%
\pgfpathlineto{\pgfqpoint{2.485004in}{1.002628in}}%
\pgfpathlineto{\pgfqpoint{2.486586in}{1.004944in}}%
\pgfpathlineto{\pgfqpoint{2.487617in}{1.004173in}}%
\pgfpathlineto{\pgfqpoint{2.489714in}{1.001906in}}%
\pgfpathlineto{\pgfqpoint{2.489967in}{1.002687in}}%
\pgfpathlineto{\pgfqpoint{2.491550in}{1.004945in}}%
\pgfpathlineto{\pgfqpoint{2.492610in}{1.004115in}}%
\pgfpathlineto{\pgfqpoint{2.494652in}{1.001884in}}%
\pgfpathlineto{\pgfqpoint{2.494922in}{1.002690in}}%
\pgfpathlineto{\pgfqpoint{2.496505in}{1.004917in}}%
\pgfpathlineto{\pgfqpoint{2.497574in}{1.004064in}}%
\pgfpathlineto{\pgfqpoint{2.499603in}{1.001870in}}%
\pgfpathlineto{\pgfqpoint{2.499841in}{1.002604in}}%
\pgfpathlineto{\pgfqpoint{2.501423in}{1.004921in}}%
\pgfpathlineto{\pgfqpoint{2.502454in}{1.004153in}}%
\pgfpathlineto{\pgfqpoint{2.504550in}{1.001886in}}%
\pgfpathlineto{\pgfqpoint{2.504808in}{1.002680in}}%
\pgfpathlineto{\pgfqpoint{2.506391in}{1.004913in}}%
\pgfpathlineto{\pgfqpoint{2.507457in}{1.004064in}}%
\pgfpathlineto{\pgfqpoint{2.509488in}{1.001846in}}%
\pgfpathlineto{\pgfqpoint{2.509759in}{1.002657in}}%
\pgfpathlineto{\pgfqpoint{2.511341in}{1.004900in}}%
\pgfpathlineto{\pgfqpoint{2.512405in}{1.004061in}}%
\pgfpathlineto{\pgfqpoint{2.514439in}{1.001860in}}%
\pgfpathlineto{\pgfqpoint{2.514677in}{1.002595in}}%
\pgfpathlineto{\pgfqpoint{2.516260in}{1.004900in}}%
\pgfpathlineto{\pgfqpoint{2.517291in}{1.004127in}}%
\pgfpathlineto{\pgfqpoint{2.519385in}{1.001864in}}%
\pgfpathlineto{\pgfqpoint{2.519622in}{1.002596in}}%
\pgfpathlineto{\pgfqpoint{2.521205in}{1.004894in}}%
\pgfpathlineto{\pgfqpoint{2.522236in}{1.004115in}}%
\pgfpathlineto{\pgfqpoint{2.524326in}{1.001831in}}%
\pgfpathlineto{\pgfqpoint{2.524596in}{1.002638in}}%
\pgfpathlineto{\pgfqpoint{2.526179in}{1.004870in}}%
\pgfpathlineto{\pgfqpoint{2.527245in}{1.004020in}}%
\pgfpathlineto{\pgfqpoint{2.529275in}{1.001817in}}%
\pgfpathlineto{\pgfqpoint{2.529514in}{1.002548in}}%
\pgfpathlineto{\pgfqpoint{2.531097in}{1.004849in}}%
\pgfpathlineto{\pgfqpoint{2.532127in}{1.004079in}}%
\pgfpathlineto{\pgfqpoint{2.534215in}{1.001798in}}%
\pgfpathlineto{\pgfqpoint{2.534493in}{1.002622in}}%
\pgfpathlineto{\pgfqpoint{2.536075in}{1.004816in}}%
\pgfpathlineto{\pgfqpoint{2.537153in}{1.003944in}}%
\pgfpathlineto{\pgfqpoint{2.539169in}{1.001766in}}%
\pgfpathlineto{\pgfqpoint{2.539422in}{1.002549in}}%
\pgfpathlineto{\pgfqpoint{2.541004in}{1.004836in}}%
\pgfpathlineto{\pgfqpoint{2.542026in}{1.004058in}}%
\pgfpathlineto{\pgfqpoint{2.544106in}{1.001785in}}%
\pgfpathlineto{\pgfqpoint{2.544383in}{1.002608in}}%
\pgfpathlineto{\pgfqpoint{2.545966in}{1.004800in}}%
\pgfpathlineto{\pgfqpoint{2.547045in}{1.003926in}}%
\pgfpathlineto{\pgfqpoint{2.549060in}{1.001751in}}%
\pgfpathlineto{\pgfqpoint{2.549313in}{1.002537in}}%
\pgfpathlineto{\pgfqpoint{2.550896in}{1.004819in}}%
\pgfpathlineto{\pgfqpoint{2.551917in}{1.004040in}}%
\pgfpathlineto{\pgfqpoint{2.553999in}{1.001766in}}%
\pgfpathlineto{\pgfqpoint{2.554274in}{1.002582in}}%
\pgfpathlineto{\pgfqpoint{2.555856in}{1.004780in}}%
\pgfpathlineto{\pgfqpoint{2.556933in}{1.003907in}}%
\pgfpathlineto{\pgfqpoint{2.558939in}{1.001716in}}%
\pgfpathlineto{\pgfqpoint{2.559213in}{1.002535in}}%
\pgfpathlineto{\pgfqpoint{2.560796in}{1.004770in}}%
\pgfpathlineto{\pgfqpoint{2.561860in}{1.003923in}}%
\pgfpathlineto{\pgfqpoint{2.563888in}{1.001707in}}%
\pgfpathlineto{\pgfqpoint{2.564159in}{1.002520in}}%
\pgfpathlineto{\pgfqpoint{2.565742in}{1.004773in}}%
\pgfpathlineto{\pgfqpoint{2.566803in}{1.003940in}}%
\pgfpathlineto{\pgfqpoint{2.568840in}{1.001742in}}%
\pgfpathlineto{\pgfqpoint{2.569081in}{1.002489in}}%
\pgfpathlineto{\pgfqpoint{2.570664in}{1.004779in}}%
\pgfpathlineto{\pgfqpoint{2.571693in}{1.003998in}}%
\pgfpathlineto{\pgfqpoint{2.573785in}{1.001732in}}%
\pgfpathlineto{\pgfqpoint{2.574023in}{1.002465in}}%
\pgfpathlineto{\pgfqpoint{2.575605in}{1.004779in}}%
\pgfpathlineto{\pgfqpoint{2.576636in}{1.004011in}}%
\pgfpathlineto{\pgfqpoint{2.578724in}{1.001730in}}%
\pgfpathlineto{\pgfqpoint{2.578997in}{1.002539in}}%
\pgfpathlineto{\pgfqpoint{2.580579in}{1.004743in}}%
\pgfpathlineto{\pgfqpoint{2.581656in}{1.003875in}}%
\pgfpathlineto{\pgfqpoint{2.583678in}{1.001692in}}%
\pgfpathlineto{\pgfqpoint{2.583931in}{1.002477in}}%
\pgfpathlineto{\pgfqpoint{2.585513in}{1.004763in}}%
\pgfpathlineto{\pgfqpoint{2.586536in}{1.003986in}}%
\pgfpathlineto{\pgfqpoint{2.588623in}{1.001730in}}%
\pgfpathlineto{\pgfqpoint{2.588881in}{1.002522in}}%
\pgfpathlineto{\pgfqpoint{2.590463in}{1.004752in}}%
\pgfpathlineto{\pgfqpoint{2.591531in}{1.003902in}}%
\pgfpathlineto{\pgfqpoint{2.593563in}{1.001684in}}%
\pgfpathlineto{\pgfqpoint{2.593837in}{1.002504in}}%
\pgfpathlineto{\pgfqpoint{2.595420in}{1.004737in}}%
\pgfpathlineto{\pgfqpoint{2.596486in}{1.003887in}}%
\pgfpathlineto{\pgfqpoint{2.598512in}{1.001691in}}%
\pgfpathlineto{\pgfqpoint{2.598750in}{1.002420in}}%
\pgfpathlineto{\pgfqpoint{2.600333in}{1.004714in}}%
\pgfpathlineto{\pgfqpoint{2.601363in}{1.003941in}}%
\pgfpathlineto{\pgfqpoint{2.603456in}{1.001655in}}%
\pgfpathlineto{\pgfqpoint{2.603698in}{1.002378in}}%
\pgfpathlineto{\pgfqpoint{2.605280in}{1.004661in}}%
\pgfpathlineto{\pgfqpoint{2.606310in}{1.003887in}}%
\pgfpathlineto{\pgfqpoint{2.608403in}{1.001625in}}%
\pgfpathlineto{\pgfqpoint{2.608641in}{1.002359in}}%
\pgfpathlineto{\pgfqpoint{2.610223in}{1.004659in}}%
\pgfpathlineto{\pgfqpoint{2.611254in}{1.003883in}}%
\pgfpathlineto{\pgfqpoint{2.613348in}{1.001613in}}%
\pgfpathlineto{\pgfqpoint{2.613587in}{1.002345in}}%
\pgfpathlineto{\pgfqpoint{2.615169in}{1.004649in}}%
\pgfpathlineto{\pgfqpoint{2.616200in}{1.003880in}}%
\pgfpathlineto{\pgfqpoint{2.618296in}{1.001615in}}%
\pgfpathlineto{\pgfqpoint{2.618547in}{1.002388in}}%
\pgfpathlineto{\pgfqpoint{2.620129in}{1.004638in}}%
\pgfpathlineto{\pgfqpoint{2.621189in}{1.003802in}}%
\pgfpathlineto{\pgfqpoint{2.623234in}{1.001569in}}%
\pgfpathlineto{\pgfqpoint{2.623501in}{1.002370in}}%
\pgfpathlineto{\pgfqpoint{2.625084in}{1.004636in}}%
\pgfpathlineto{\pgfqpoint{2.626143in}{1.003811in}}%
\pgfpathlineto{\pgfqpoint{2.628188in}{1.001610in}}%
\pgfpathlineto{\pgfqpoint{2.628412in}{1.002308in}}%
\pgfpathlineto{\pgfqpoint{2.629995in}{1.004643in}}%
\pgfpathlineto{\pgfqpoint{2.631031in}{1.003865in}}%
\pgfpathlineto{\pgfqpoint{2.633126in}{1.001577in}}%
\pgfpathlineto{\pgfqpoint{2.633400in}{1.002397in}}%
\pgfpathlineto{\pgfqpoint{2.634983in}{1.004636in}}%
\pgfpathlineto{\pgfqpoint{2.636048in}{1.003790in}}%
\pgfpathlineto{\pgfqpoint{2.638077in}{1.001598in}}%
\pgfpathlineto{\pgfqpoint{2.638335in}{1.002391in}}%
\pgfpathlineto{\pgfqpoint{2.639917in}{1.004633in}}%
\pgfpathlineto{\pgfqpoint{2.640981in}{1.003794in}}%
\pgfpathlineto{\pgfqpoint{2.643023in}{1.001588in}}%
\pgfpathlineto{\pgfqpoint{2.643281in}{1.002386in}}%
\pgfpathlineto{\pgfqpoint{2.644863in}{1.004645in}}%
\pgfpathlineto{\pgfqpoint{2.645924in}{1.003816in}}%
\pgfpathlineto{\pgfqpoint{2.647960in}{1.001593in}}%
\pgfpathlineto{\pgfqpoint{2.648227in}{1.002391in}}%
\pgfpathlineto{\pgfqpoint{2.649810in}{1.004625in}}%
\pgfpathlineto{\pgfqpoint{2.650878in}{1.003779in}}%
\pgfpathlineto{\pgfqpoint{2.652914in}{1.001585in}}%
\pgfpathlineto{\pgfqpoint{2.653151in}{1.002324in}}%
\pgfpathlineto{\pgfqpoint{2.654734in}{1.004644in}}%
\pgfpathlineto{\pgfqpoint{2.655764in}{1.003872in}}%
\pgfpathlineto{\pgfqpoint{2.657853in}{1.001591in}}%
\pgfpathlineto{\pgfqpoint{2.658121in}{1.002386in}}%
\pgfpathlineto{\pgfqpoint{2.659703in}{1.004604in}}%
\pgfpathlineto{\pgfqpoint{2.660775in}{1.003744in}}%
\pgfpathlineto{\pgfqpoint{2.662805in}{1.001552in}}%
\pgfpathlineto{\pgfqpoint{2.663057in}{1.002334in}}%
\pgfpathlineto{\pgfqpoint{2.664640in}{1.004611in}}%
\pgfpathlineto{\pgfqpoint{2.665662in}{1.003827in}}%
\pgfpathlineto{\pgfqpoint{2.667749in}{1.001574in}}%
\pgfpathlineto{\pgfqpoint{2.667986in}{1.002309in}}%
\pgfpathlineto{\pgfqpoint{2.669568in}{1.004623in}}%
\pgfpathlineto{\pgfqpoint{2.670599in}{1.003852in}}%
\pgfpathlineto{\pgfqpoint{2.672697in}{1.001594in}}%
\pgfpathlineto{\pgfqpoint{2.672917in}{1.002280in}}%
\pgfpathlineto{\pgfqpoint{2.674549in}{1.004626in}}%
\pgfpathlineto{\pgfqpoint{2.675619in}{1.003762in}}%
\pgfpathlineto{\pgfqpoint{2.677641in}{1.001577in}}%
\pgfpathlineto{\pgfqpoint{2.677902in}{1.002384in}}%
\pgfpathlineto{\pgfqpoint{2.679484in}{1.004653in}}%
\pgfpathlineto{\pgfqpoint{2.680543in}{1.003830in}}%
\pgfpathlineto{\pgfqpoint{2.682587in}{1.001623in}}%
\pgfpathlineto{\pgfqpoint{2.682845in}{1.002417in}}%
\pgfpathlineto{\pgfqpoint{2.684427in}{1.004649in}}%
\pgfpathlineto{\pgfqpoint{2.685494in}{1.003801in}}%
\pgfpathlineto{\pgfqpoint{2.687527in}{1.001581in}}%
\pgfpathlineto{\pgfqpoint{2.687770in}{1.002310in}}%
\pgfpathlineto{\pgfqpoint{2.689352in}{1.004628in}}%
\pgfpathlineto{\pgfqpoint{2.690380in}{1.003864in}}%
\pgfpathlineto{\pgfqpoint{2.692478in}{1.001602in}}%
\pgfpathlineto{\pgfqpoint{2.692734in}{1.002390in}}%
\pgfpathlineto{\pgfqpoint{2.694317in}{1.004629in}}%
\pgfpathlineto{\pgfqpoint{2.695380in}{1.003785in}}%
\pgfpathlineto{\pgfqpoint{2.697422in}{1.001582in}}%
\pgfpathlineto{\pgfqpoint{2.697658in}{1.002318in}}%
\pgfpathlineto{\pgfqpoint{2.699241in}{1.004654in}}%
\pgfpathlineto{\pgfqpoint{2.700272in}{1.003891in}}%
\pgfpathlineto{\pgfqpoint{2.702369in}{1.001625in}}%
\pgfpathlineto{\pgfqpoint{2.702629in}{1.002426in}}%
\pgfpathlineto{\pgfqpoint{2.704212in}{1.004649in}}%
\pgfpathlineto{\pgfqpoint{2.705283in}{1.003792in}}%
\pgfpathlineto{\pgfqpoint{2.707306in}{1.001583in}}%
\pgfpathlineto{\pgfqpoint{2.707578in}{1.002398in}}%
\pgfpathlineto{\pgfqpoint{2.709161in}{1.004637in}}%
\pgfpathlineto{\pgfqpoint{2.710228in}{1.003794in}}%
\pgfpathlineto{\pgfqpoint{2.712258in}{1.001579in}}%
\pgfpathlineto{\pgfqpoint{2.712512in}{1.002344in}}%
\pgfpathlineto{\pgfqpoint{2.714094in}{1.004612in}}%
\pgfpathlineto{\pgfqpoint{2.715152in}{1.003789in}}%
\pgfpathlineto{\pgfqpoint{2.717200in}{1.001552in}}%
\pgfpathlineto{\pgfqpoint{2.717444in}{1.002280in}}%
\pgfpathlineto{\pgfqpoint{2.719026in}{1.004575in}}%
\pgfpathlineto{\pgfqpoint{2.720055in}{1.003805in}}%
\pgfpathlineto{\pgfqpoint{2.722146in}{1.001522in}}%
\pgfpathlineto{\pgfqpoint{2.722403in}{1.002291in}}%
\pgfpathlineto{\pgfqpoint{2.723986in}{1.004534in}}%
\pgfpathlineto{\pgfqpoint{2.725051in}{1.003694in}}%
\pgfpathlineto{\pgfqpoint{2.727094in}{1.001483in}}%
\pgfpathlineto{\pgfqpoint{2.727332in}{1.002216in}}%
\pgfpathlineto{\pgfqpoint{2.728914in}{1.004532in}}%
\pgfpathlineto{\pgfqpoint{2.729945in}{1.003766in}}%
\pgfpathlineto{\pgfqpoint{2.732035in}{1.001484in}}%
\pgfpathlineto{\pgfqpoint{2.732295in}{1.002256in}}%
\pgfpathlineto{\pgfqpoint{2.733877in}{1.004500in}}%
\pgfpathlineto{\pgfqpoint{2.734941in}{1.003659in}}%
\pgfpathlineto{\pgfqpoint{2.736987in}{1.001450in}}%
\pgfpathlineto{\pgfqpoint{2.737249in}{1.002260in}}%
\pgfpathlineto{\pgfqpoint{2.738832in}{1.004498in}}%
\pgfpathlineto{\pgfqpoint{2.739898in}{1.003654in}}%
\pgfpathlineto{\pgfqpoint{2.741924in}{1.001442in}}%
\pgfpathlineto{\pgfqpoint{2.742158in}{1.002142in}}%
\pgfpathlineto{\pgfqpoint{2.743740in}{1.004480in}}%
\pgfpathlineto{\pgfqpoint{2.744772in}{1.003716in}}%
\pgfpathlineto{\pgfqpoint{2.746877in}{1.001446in}}%
\pgfpathlineto{\pgfqpoint{2.747116in}{1.002188in}}%
\pgfpathlineto{\pgfqpoint{2.748699in}{1.004495in}}%
\pgfpathlineto{\pgfqpoint{2.749728in}{1.003719in}}%
\pgfpathlineto{\pgfqpoint{2.751823in}{1.001457in}}%
\pgfpathlineto{\pgfqpoint{2.752077in}{1.002239in}}%
\pgfpathlineto{\pgfqpoint{2.753659in}{1.004502in}}%
\pgfpathlineto{\pgfqpoint{2.754719in}{1.003675in}}%
\pgfpathlineto{\pgfqpoint{2.756761in}{1.001446in}}%
\pgfpathlineto{\pgfqpoint{2.757031in}{1.002249in}}%
\pgfpathlineto{\pgfqpoint{2.758613in}{1.004477in}}%
\pgfpathlineto{\pgfqpoint{2.759683in}{1.003626in}}%
\pgfpathlineto{\pgfqpoint{2.761709in}{1.001414in}}%
\pgfpathlineto{\pgfqpoint{2.761952in}{1.002143in}}%
\pgfpathlineto{\pgfqpoint{2.763534in}{1.004457in}}%
\pgfpathlineto{\pgfqpoint{2.764562in}{1.003693in}}%
\pgfpathlineto{\pgfqpoint{2.766653in}{1.001410in}}%
\pgfpathlineto{\pgfqpoint{2.766929in}{1.002228in}}%
\pgfpathlineto{\pgfqpoint{2.768512in}{1.004422in}}%
\pgfpathlineto{\pgfqpoint{2.769588in}{1.003549in}}%
\pgfpathlineto{\pgfqpoint{2.771592in}{1.001362in}}%
\pgfpathlineto{\pgfqpoint{2.771861in}{1.002163in}}%
\pgfpathlineto{\pgfqpoint{2.773444in}{1.004426in}}%
\pgfpathlineto{\pgfqpoint{2.774500in}{1.003599in}}%
\pgfpathlineto{\pgfqpoint{2.776544in}{1.001367in}}%
\pgfpathlineto{\pgfqpoint{2.776812in}{1.002169in}}%
\pgfpathlineto{\pgfqpoint{2.778395in}{1.004420in}}%
\pgfpathlineto{\pgfqpoint{2.779453in}{1.003587in}}%
\pgfpathlineto{\pgfqpoint{2.781494in}{1.001384in}}%
\pgfpathlineto{\pgfqpoint{2.781731in}{1.002118in}}%
\pgfpathlineto{\pgfqpoint{2.783314in}{1.004437in}}%
\pgfpathlineto{\pgfqpoint{2.784345in}{1.003669in}}%
\pgfpathlineto{\pgfqpoint{2.786441in}{1.001404in}}%
\pgfpathlineto{\pgfqpoint{2.786695in}{1.002186in}}%
\pgfpathlineto{\pgfqpoint{2.788277in}{1.004442in}}%
\pgfpathlineto{\pgfqpoint{2.789337in}{1.003614in}}%
\pgfpathlineto{\pgfqpoint{2.791387in}{1.001400in}}%
\pgfpathlineto{\pgfqpoint{2.791640in}{1.002181in}}%
\pgfpathlineto{\pgfqpoint{2.793222in}{1.004444in}}%
\pgfpathlineto{\pgfqpoint{2.794279in}{1.003617in}}%
\pgfpathlineto{\pgfqpoint{2.796332in}{1.001406in}}%
\pgfpathlineto{\pgfqpoint{2.796568in}{1.002139in}}%
\pgfpathlineto{\pgfqpoint{2.798150in}{1.004464in}}%
\pgfpathlineto{\pgfqpoint{2.799181in}{1.003691in}}%
\pgfpathlineto{\pgfqpoint{2.801275in}{1.001419in}}%
\pgfpathlineto{\pgfqpoint{2.801514in}{1.002150in}}%
\pgfpathlineto{\pgfqpoint{2.803097in}{1.004460in}}%
\pgfpathlineto{\pgfqpoint{2.804127in}{1.003694in}}%
\pgfpathlineto{\pgfqpoint{2.806216in}{1.001414in}}%
\pgfpathlineto{\pgfqpoint{2.806473in}{1.002178in}}%
\pgfpathlineto{\pgfqpoint{2.808056in}{1.004430in}}%
\pgfpathlineto{\pgfqpoint{2.809117in}{1.003595in}}%
\pgfpathlineto{\pgfqpoint{2.811168in}{1.001380in}}%
\pgfpathlineto{\pgfqpoint{2.811431in}{1.002189in}}%
\pgfpathlineto{\pgfqpoint{2.813013in}{1.004429in}}%
\pgfpathlineto{\pgfqpoint{2.814080in}{1.003586in}}%
\pgfpathlineto{\pgfqpoint{2.816108in}{1.001372in}}%
\pgfpathlineto{\pgfqpoint{2.816383in}{1.002190in}}%
\pgfpathlineto{\pgfqpoint{2.817965in}{1.004406in}}%
\pgfpathlineto{\pgfqpoint{2.819038in}{1.003545in}}%
\pgfpathlineto{\pgfqpoint{2.821059in}{1.001368in}}%
\pgfpathlineto{\pgfqpoint{2.821299in}{1.002112in}}%
\pgfpathlineto{\pgfqpoint{2.822881in}{1.004407in}}%
\pgfpathlineto{\pgfqpoint{2.823910in}{1.003628in}}%
\pgfpathlineto{\pgfqpoint{2.826006in}{1.001366in}}%
\pgfpathlineto{\pgfqpoint{2.826258in}{1.002147in}}%
\pgfpathlineto{\pgfqpoint{2.827841in}{1.004425in}}%
\pgfpathlineto{\pgfqpoint{2.828863in}{1.003645in}}%
\pgfpathlineto{\pgfqpoint{2.830944in}{1.001374in}}%
\pgfpathlineto{\pgfqpoint{2.831208in}{1.002162in}}%
\pgfpathlineto{\pgfqpoint{2.832791in}{1.004396in}}%
\pgfpathlineto{\pgfqpoint{2.833858in}{1.003551in}}%
\pgfpathlineto{\pgfqpoint{2.835888in}{1.001335in}}%
\pgfpathlineto{\pgfqpoint{2.836118in}{1.002028in}}%
\pgfpathlineto{\pgfqpoint{2.837750in}{1.004393in}}%
\pgfpathlineto{\pgfqpoint{2.838818in}{1.003543in}}%
\pgfpathlineto{\pgfqpoint{2.840841in}{1.001357in}}%
\pgfpathlineto{\pgfqpoint{2.841101in}{1.002156in}}%
\pgfpathlineto{\pgfqpoint{2.842683in}{1.004404in}}%
\pgfpathlineto{\pgfqpoint{2.843746in}{1.003568in}}%
\pgfpathlineto{\pgfqpoint{2.845787in}{1.001366in}}%
\pgfpathlineto{\pgfqpoint{2.846047in}{1.002168in}}%
\pgfpathlineto{\pgfqpoint{2.847629in}{1.004419in}}%
\pgfpathlineto{\pgfqpoint{2.848690in}{1.003586in}}%
\pgfpathlineto{\pgfqpoint{2.850719in}{1.001371in}}%
\pgfpathlineto{\pgfqpoint{2.850954in}{1.002073in}}%
\pgfpathlineto{\pgfqpoint{2.852586in}{1.004431in}}%
\pgfpathlineto{\pgfqpoint{2.853655in}{1.003571in}}%
\pgfpathlineto{\pgfqpoint{2.855677in}{1.001387in}}%
\pgfpathlineto{\pgfqpoint{2.855936in}{1.002189in}}%
\pgfpathlineto{\pgfqpoint{2.857518in}{1.004472in}}%
\pgfpathlineto{\pgfqpoint{2.858572in}{1.003660in}}%
\pgfpathlineto{\pgfqpoint{2.860618in}{1.001426in}}%
\pgfpathlineto{\pgfqpoint{2.860862in}{1.002153in}}%
\pgfpathlineto{\pgfqpoint{2.862445in}{1.004429in}}%
\pgfpathlineto{\pgfqpoint{2.863473in}{1.003651in}}%
\pgfpathlineto{\pgfqpoint{2.865567in}{1.001391in}}%
\pgfpathlineto{\pgfqpoint{2.865804in}{1.002122in}}%
\pgfpathlineto{\pgfqpoint{2.867386in}{1.004422in}}%
\pgfpathlineto{\pgfqpoint{2.868417in}{1.003645in}}%
\pgfpathlineto{\pgfqpoint{2.870506in}{1.001365in}}%
\pgfpathlineto{\pgfqpoint{2.870783in}{1.002189in}}%
\pgfpathlineto{\pgfqpoint{2.872365in}{1.004408in}}%
\pgfpathlineto{\pgfqpoint{2.873436in}{1.003554in}}%
\pgfpathlineto{\pgfqpoint{2.875457in}{1.001357in}}%
\pgfpathlineto{\pgfqpoint{2.875697in}{1.002084in}}%
\pgfpathlineto{\pgfqpoint{2.877279in}{1.004390in}}%
\pgfpathlineto{\pgfqpoint{2.878309in}{1.003626in}}%
\pgfpathlineto{\pgfqpoint{2.880399in}{1.001344in}}%
\pgfpathlineto{\pgfqpoint{2.880631in}{1.002035in}}%
\pgfpathlineto{\pgfqpoint{2.882213in}{1.004352in}}%
\pgfpathlineto{\pgfqpoint{2.883249in}{1.003569in}}%
\pgfpathlineto{\pgfqpoint{2.885348in}{1.001292in}}%
\pgfpathlineto{\pgfqpoint{2.885587in}{1.002026in}}%
\pgfpathlineto{\pgfqpoint{2.887170in}{1.004352in}}%
\pgfpathlineto{\pgfqpoint{2.888200in}{1.003592in}}%
\pgfpathlineto{\pgfqpoint{2.890296in}{1.001329in}}%
\pgfpathlineto{\pgfqpoint{2.890551in}{1.002112in}}%
\pgfpathlineto{\pgfqpoint{2.892133in}{1.004341in}}%
\pgfpathlineto{\pgfqpoint{2.893199in}{1.003492in}}%
\pgfpathlineto{\pgfqpoint{2.895234in}{1.001271in}}%
\pgfpathlineto{\pgfqpoint{2.895499in}{1.002067in}}%
\pgfpathlineto{\pgfqpoint{2.897081in}{1.004349in}}%
\pgfpathlineto{\pgfqpoint{2.898136in}{1.003536in}}%
\pgfpathlineto{\pgfqpoint{2.900187in}{1.001322in}}%
\pgfpathlineto{\pgfqpoint{2.900440in}{1.002104in}}%
\pgfpathlineto{\pgfqpoint{2.902023in}{1.004358in}}%
\pgfpathlineto{\pgfqpoint{2.903082in}{1.003528in}}%
\pgfpathlineto{\pgfqpoint{2.905133in}{1.001321in}}%
\pgfpathlineto{\pgfqpoint{2.905389in}{1.002113in}}%
\pgfpathlineto{\pgfqpoint{2.906972in}{1.004373in}}%
\pgfpathlineto{\pgfqpoint{2.908027in}{1.003547in}}%
\pgfpathlineto{\pgfqpoint{2.910075in}{1.001326in}}%
\pgfpathlineto{\pgfqpoint{2.910314in}{1.002057in}}%
\pgfpathlineto{\pgfqpoint{2.911896in}{1.004371in}}%
\pgfpathlineto{\pgfqpoint{2.912927in}{1.003607in}}%
\pgfpathlineto{\pgfqpoint{2.915023in}{1.001343in}}%
\pgfpathlineto{\pgfqpoint{2.915288in}{1.002154in}}%
\pgfpathlineto{\pgfqpoint{2.916870in}{1.004367in}}%
\pgfpathlineto{\pgfqpoint{2.917942in}{1.003508in}}%
\pgfpathlineto{\pgfqpoint{2.919968in}{1.001319in}}%
\pgfpathlineto{\pgfqpoint{2.920229in}{1.002126in}}%
\pgfpathlineto{\pgfqpoint{2.921812in}{1.004391in}}%
\pgfpathlineto{\pgfqpoint{2.922871in}{1.003566in}}%
\pgfpathlineto{\pgfqpoint{2.924914in}{1.001361in}}%
\pgfpathlineto{\pgfqpoint{2.925172in}{1.002155in}}%
\pgfpathlineto{\pgfqpoint{2.926754in}{1.004396in}}%
\pgfpathlineto{\pgfqpoint{2.927819in}{1.003556in}}%
\pgfpathlineto{\pgfqpoint{2.929859in}{1.001353in}}%
\pgfpathlineto{\pgfqpoint{2.930119in}{1.002156in}}%
\pgfpathlineto{\pgfqpoint{2.931702in}{1.004412in}}%
\pgfpathlineto{\pgfqpoint{2.932762in}{1.003581in}}%
\pgfpathlineto{\pgfqpoint{2.934792in}{1.001366in}}%
\pgfpathlineto{\pgfqpoint{2.935033in}{1.002083in}}%
\pgfpathlineto{\pgfqpoint{2.936616in}{1.004419in}}%
\pgfpathlineto{\pgfqpoint{2.937650in}{1.003641in}}%
\pgfpathlineto{\pgfqpoint{2.939747in}{1.001353in}}%
\pgfpathlineto{\pgfqpoint{2.939981in}{1.002064in}}%
\pgfpathlineto{\pgfqpoint{2.941613in}{1.004431in}}%
\pgfpathlineto{\pgfqpoint{2.942676in}{1.003583in}}%
\pgfpathlineto{\pgfqpoint{2.944696in}{1.001403in}}%
\pgfpathlineto{\pgfqpoint{2.944958in}{1.002210in}}%
\pgfpathlineto{\pgfqpoint{2.946541in}{1.004432in}}%
\pgfpathlineto{\pgfqpoint{2.947612in}{1.003577in}}%
\pgfpathlineto{\pgfqpoint{2.949634in}{1.001370in}}%
\pgfpathlineto{\pgfqpoint{2.949889in}{1.002133in}}%
\pgfpathlineto{\pgfqpoint{2.951471in}{1.004424in}}%
\pgfpathlineto{\pgfqpoint{2.952502in}{1.003640in}}%
\pgfpathlineto{\pgfqpoint{2.954587in}{1.001388in}}%
\pgfpathlineto{\pgfqpoint{2.954844in}{1.002183in}}%
\pgfpathlineto{\pgfqpoint{2.956426in}{1.004438in}}%
\pgfpathlineto{\pgfqpoint{2.957488in}{1.003607in}}%
\pgfpathlineto{\pgfqpoint{2.959531in}{1.001388in}}%
\pgfpathlineto{\pgfqpoint{2.959783in}{1.002153in}}%
\pgfpathlineto{\pgfqpoint{2.961365in}{1.004424in}}%
\pgfpathlineto{\pgfqpoint{2.962423in}{1.003603in}}%
\pgfpathlineto{\pgfqpoint{2.964476in}{1.001383in}}%
\pgfpathlineto{\pgfqpoint{2.964714in}{1.002117in}}%
\pgfpathlineto{\pgfqpoint{2.966296in}{1.004420in}}%
\pgfpathlineto{\pgfqpoint{2.967327in}{1.003648in}}%
\pgfpathlineto{\pgfqpoint{2.969422in}{1.001388in}}%
\pgfpathlineto{\pgfqpoint{2.969663in}{1.002132in}}%
\pgfpathlineto{\pgfqpoint{2.971245in}{1.004421in}}%
\pgfpathlineto{\pgfqpoint{2.972274in}{1.003640in}}%
\pgfpathlineto{\pgfqpoint{2.974363in}{1.001360in}}%
\pgfpathlineto{\pgfqpoint{2.974637in}{1.002179in}}%
\pgfpathlineto{\pgfqpoint{2.976220in}{1.004403in}}%
\pgfpathlineto{\pgfqpoint{2.977290in}{1.003547in}}%
\pgfpathlineto{\pgfqpoint{2.979301in}{1.001353in}}%
\pgfpathlineto{\pgfqpoint{2.979570in}{1.002149in}}%
\pgfpathlineto{\pgfqpoint{2.981152in}{1.004402in}}%
\pgfpathlineto{\pgfqpoint{2.982209in}{1.003570in}}%
\pgfpathlineto{\pgfqpoint{2.984253in}{1.001339in}}%
\pgfpathlineto{\pgfqpoint{2.984524in}{1.002151in}}%
\pgfpathlineto{\pgfqpoint{2.986106in}{1.004404in}}%
\pgfpathlineto{\pgfqpoint{2.987166in}{1.003570in}}%
\pgfpathlineto{\pgfqpoint{2.989197in}{1.001353in}}%
\pgfpathlineto{\pgfqpoint{2.989462in}{1.002140in}}%
\pgfpathlineto{\pgfqpoint{2.991044in}{1.004383in}}%
\pgfpathlineto{\pgfqpoint{2.992108in}{1.003543in}}%
\pgfpathlineto{\pgfqpoint{2.994148in}{1.001337in}}%
\pgfpathlineto{\pgfqpoint{2.994387in}{1.002074in}}%
\pgfpathlineto{\pgfqpoint{2.995969in}{1.004398in}}%
\pgfpathlineto{\pgfqpoint{2.997000in}{1.003635in}}%
\pgfpathlineto{\pgfqpoint{2.999090in}{1.001353in}}%
\pgfpathlineto{\pgfqpoint{2.999365in}{1.002169in}}%
\pgfpathlineto{\pgfqpoint{3.000947in}{1.004361in}}%
\pgfpathlineto{\pgfqpoint{3.002023in}{1.003489in}}%
\pgfpathlineto{\pgfqpoint{3.004037in}{1.001291in}}%
\pgfpathlineto{\pgfqpoint{3.004281in}{1.002021in}}%
\pgfpathlineto{\pgfqpoint{3.005863in}{1.004326in}}%
\pgfpathlineto{\pgfqpoint{3.006892in}{1.003560in}}%
\pgfpathlineto{\pgfqpoint{3.008980in}{1.001281in}}%
\pgfpathlineto{\pgfqpoint{3.009236in}{1.002041in}}%
\pgfpathlineto{\pgfqpoint{3.010818in}{1.004296in}}%
\pgfpathlineto{\pgfqpoint{3.011878in}{1.003464in}}%
\pgfpathlineto{\pgfqpoint{3.013932in}{1.001246in}}%
\pgfpathlineto{\pgfqpoint{3.014191in}{1.002044in}}%
\pgfpathlineto{\pgfqpoint{3.015773in}{1.004298in}}%
\pgfpathlineto{\pgfqpoint{3.016835in}{1.003467in}}%
\pgfpathlineto{\pgfqpoint{3.018874in}{1.001241in}}%
\pgfpathlineto{\pgfqpoint{3.019152in}{1.002068in}}%
\pgfpathlineto{\pgfqpoint{3.020735in}{1.004264in}}%
\pgfpathlineto{\pgfqpoint{3.021808in}{1.003406in}}%
\pgfpathlineto{\pgfqpoint{3.023822in}{1.001235in}}%
\pgfpathlineto{\pgfqpoint{3.024058in}{1.001966in}}%
\pgfpathlineto{\pgfqpoint{3.025641in}{1.004266in}}%
\pgfpathlineto{\pgfqpoint{3.026672in}{1.003491in}}%
\pgfpathlineto{\pgfqpoint{3.028769in}{1.001227in}}%
\pgfpathlineto{\pgfqpoint{3.029019in}{1.002003in}}%
\pgfpathlineto{\pgfqpoint{3.030602in}{1.004278in}}%
\pgfpathlineto{\pgfqpoint{3.031625in}{1.003493in}}%
\pgfpathlineto{\pgfqpoint{3.033708in}{1.001220in}}%
\pgfpathlineto{\pgfqpoint{3.033982in}{1.002039in}}%
\pgfpathlineto{\pgfqpoint{3.035565in}{1.004269in}}%
\pgfpathlineto{\pgfqpoint{3.036633in}{1.003419in}}%
\pgfpathlineto{\pgfqpoint{3.038659in}{1.001232in}}%
\pgfpathlineto{\pgfqpoint{3.038923in}{1.002043in}}%
\pgfpathlineto{\pgfqpoint{3.040505in}{1.004271in}}%
\pgfpathlineto{\pgfqpoint{3.041575in}{1.003420in}}%
\pgfpathlineto{\pgfqpoint{3.043604in}{1.001230in}}%
\pgfpathlineto{\pgfqpoint{3.043856in}{1.002006in}}%
\pgfpathlineto{\pgfqpoint{3.045438in}{1.004288in}}%
\pgfpathlineto{\pgfqpoint{3.046462in}{1.003509in}}%
\pgfpathlineto{\pgfqpoint{3.048551in}{1.001260in}}%
\pgfpathlineto{\pgfqpoint{3.048782in}{1.001978in}}%
\pgfpathlineto{\pgfqpoint{3.050365in}{1.004294in}}%
\pgfpathlineto{\pgfqpoint{3.051398in}{1.003515in}}%
\pgfpathlineto{\pgfqpoint{3.053494in}{1.001228in}}%
\pgfpathlineto{\pgfqpoint{3.053745in}{1.001987in}}%
\pgfpathlineto{\pgfqpoint{3.055328in}{1.004286in}}%
\pgfpathlineto{\pgfqpoint{3.056352in}{1.003514in}}%
\pgfpathlineto{\pgfqpoint{3.058440in}{1.001259in}}%
\pgfpathlineto{\pgfqpoint{3.058677in}{1.001990in}}%
\pgfpathlineto{\pgfqpoint{3.060260in}{1.004273in}}%
\pgfpathlineto{\pgfqpoint{3.061307in}{1.003474in}}%
\pgfpathlineto{\pgfqpoint{3.063374in}{1.001218in}}%
\pgfpathlineto{\pgfqpoint{3.063625in}{1.001967in}}%
\pgfpathlineto{\pgfqpoint{3.065207in}{1.004292in}}%
\pgfpathlineto{\pgfqpoint{3.066237in}{1.003516in}}%
\pgfpathlineto{\pgfqpoint{3.068332in}{1.001259in}}%
\pgfpathlineto{\pgfqpoint{3.068568in}{1.001997in}}%
\pgfpathlineto{\pgfqpoint{3.070200in}{1.004344in}}%
\pgfpathlineto{\pgfqpoint{3.071273in}{1.003480in}}%
\pgfpathlineto{\pgfqpoint{3.073272in}{1.001297in}}%
\pgfpathlineto{\pgfqpoint{3.073517in}{1.002026in}}%
\pgfpathlineto{\pgfqpoint{3.075099in}{1.004313in}}%
\pgfpathlineto{\pgfqpoint{3.076127in}{1.003541in}}%
\pgfpathlineto{\pgfqpoint{3.078217in}{1.001259in}}%
\pgfpathlineto{\pgfqpoint{3.078492in}{1.002075in}}%
\pgfpathlineto{\pgfqpoint{3.080074in}{1.004275in}}%
\pgfpathlineto{\pgfqpoint{3.081148in}{1.003403in}}%
\pgfpathlineto{\pgfqpoint{3.083165in}{1.001202in}}%
\pgfpathlineto{\pgfqpoint{3.083406in}{1.001932in}}%
\pgfpathlineto{\pgfqpoint{3.085038in}{1.004284in}}%
\pgfpathlineto{\pgfqpoint{3.086064in}{1.003492in}}%
\pgfpathlineto{\pgfqpoint{3.088114in}{1.001278in}}%
\pgfpathlineto{\pgfqpoint{3.088367in}{1.002047in}}%
\pgfpathlineto{\pgfqpoint{3.089949in}{1.004257in}}%
\pgfpathlineto{\pgfqpoint{3.091023in}{1.003395in}}%
\pgfpathlineto{\pgfqpoint{3.093056in}{1.001177in}}%
\pgfpathlineto{\pgfqpoint{3.093298in}{1.001908in}}%
\pgfpathlineto{\pgfqpoint{3.094880in}{1.004242in}}%
\pgfpathlineto{\pgfqpoint{3.095910in}{1.003490in}}%
\pgfpathlineto{\pgfqpoint{3.097998in}{1.001211in}}%
\pgfpathlineto{\pgfqpoint{3.098271in}{1.002016in}}%
\pgfpathlineto{\pgfqpoint{3.099854in}{1.004194in}}%
\pgfpathlineto{\pgfqpoint{3.100937in}{1.003307in}}%
\pgfpathlineto{\pgfqpoint{3.102950in}{1.001136in}}%
\pgfpathlineto{\pgfqpoint{3.103205in}{1.001926in}}%
\pgfpathlineto{\pgfqpoint{3.104787in}{1.004209in}}%
\pgfpathlineto{\pgfqpoint{3.105842in}{1.003397in}}%
\pgfpathlineto{\pgfqpoint{3.107896in}{1.001180in}}%
\pgfpathlineto{\pgfqpoint{3.108153in}{1.001972in}}%
\pgfpathlineto{\pgfqpoint{3.109736in}{1.004205in}}%
\pgfpathlineto{\pgfqpoint{3.110803in}{1.003360in}}%
\pgfpathlineto{\pgfqpoint{3.112842in}{1.001167in}}%
\pgfpathlineto{\pgfqpoint{3.113069in}{1.001874in}}%
\pgfpathlineto{\pgfqpoint{3.114701in}{1.004230in}}%
\pgfpathlineto{\pgfqpoint{3.115773in}{1.003368in}}%
\pgfpathlineto{\pgfqpoint{3.117773in}{1.001186in}}%
\pgfpathlineto{\pgfqpoint{3.118046in}{1.001998in}}%
\pgfpathlineto{\pgfqpoint{3.119629in}{1.004256in}}%
\pgfpathlineto{\pgfqpoint{3.120686in}{1.003426in}}%
\pgfpathlineto{\pgfqpoint{3.122731in}{1.001221in}}%
\pgfpathlineto{\pgfqpoint{3.122967in}{1.001956in}}%
\pgfpathlineto{\pgfqpoint{3.124550in}{1.004285in}}%
\pgfpathlineto{\pgfqpoint{3.125581in}{1.003520in}}%
\pgfpathlineto{\pgfqpoint{3.127677in}{1.001261in}}%
\pgfpathlineto{\pgfqpoint{3.127916in}{1.001999in}}%
\pgfpathlineto{\pgfqpoint{3.129498in}{1.004288in}}%
\pgfpathlineto{\pgfqpoint{3.130528in}{1.003506in}}%
\pgfpathlineto{\pgfqpoint{3.132623in}{1.001249in}}%
\pgfpathlineto{\pgfqpoint{3.132860in}{1.001987in}}%
\pgfpathlineto{\pgfqpoint{3.134443in}{1.004294in}}%
\pgfpathlineto{\pgfqpoint{3.135473in}{1.003519in}}%
\pgfpathlineto{\pgfqpoint{3.137563in}{1.001237in}}%
\pgfpathlineto{\pgfqpoint{3.137807in}{1.001963in}}%
\pgfpathlineto{\pgfqpoint{3.139389in}{1.004252in}}%
\pgfpathlineto{\pgfqpoint{3.140417in}{1.003480in}}%
\pgfpathlineto{\pgfqpoint{3.142508in}{1.001198in}}%
\pgfpathlineto{\pgfqpoint{3.142784in}{1.002017in}}%
\pgfpathlineto{\pgfqpoint{3.144367in}{1.004212in}}%
\pgfpathlineto{\pgfqpoint{3.145447in}{1.003335in}}%
\pgfpathlineto{\pgfqpoint{3.147460in}{1.001164in}}%
\pgfpathlineto{\pgfqpoint{3.147713in}{1.001948in}}%
\pgfpathlineto{\pgfqpoint{3.149295in}{1.004219in}}%
\pgfpathlineto{\pgfqpoint{3.150350in}{1.003399in}}%
\pgfpathlineto{\pgfqpoint{3.152403in}{1.001184in}}%
\pgfpathlineto{\pgfqpoint{3.152640in}{1.001918in}}%
\pgfpathlineto{\pgfqpoint{3.154223in}{1.004234in}}%
\pgfpathlineto{\pgfqpoint{3.155254in}{1.003465in}}%
\pgfpathlineto{\pgfqpoint{3.157350in}{1.001202in}}%
\pgfpathlineto{\pgfqpoint{3.157604in}{1.001985in}}%
\pgfpathlineto{\pgfqpoint{3.159187in}{1.004244in}}%
\pgfpathlineto{\pgfqpoint{3.160245in}{1.003418in}}%
\pgfpathlineto{\pgfqpoint{3.162289in}{1.001188in}}%
\pgfpathlineto{\pgfqpoint{3.162527in}{1.001902in}}%
\pgfpathlineto{\pgfqpoint{3.164110in}{1.004219in}}%
\pgfpathlineto{\pgfqpoint{3.165137in}{1.003455in}}%
\pgfpathlineto{\pgfqpoint{3.167235in}{1.001165in}}%
\pgfpathlineto{\pgfqpoint{3.167511in}{1.001986in}}%
\pgfpathlineto{\pgfqpoint{3.169094in}{1.004196in}}%
\pgfpathlineto{\pgfqpoint{3.170132in}{1.003371in}}%
\pgfpathlineto{\pgfqpoint{3.170132in}{1.003371in}}%
\pgfusepath{stroke}%
\end{pgfscope}%
\begin{pgfscope}%
\pgfsetrectcap%
\pgfsetmiterjoin%
\pgfsetlinewidth{0.803000pt}%
\definecolor{currentstroke}{rgb}{0.000000,0.000000,0.000000}%
\pgfsetstrokecolor{currentstroke}%
\pgfsetdash{}{0pt}%
\pgfpathmoveto{\pgfqpoint{0.573768in}{0.462654in}}%
\pgfpathlineto{\pgfqpoint{0.573768in}{1.368904in}}%
\pgfusepath{stroke}%
\end{pgfscope}%
\begin{pgfscope}%
\pgfsetrectcap%
\pgfsetmiterjoin%
\pgfsetlinewidth{0.803000pt}%
\definecolor{currentstroke}{rgb}{0.000000,0.000000,0.000000}%
\pgfsetstrokecolor{currentstroke}%
\pgfsetdash{}{0pt}%
\pgfpathmoveto{\pgfqpoint{3.293768in}{0.462654in}}%
\pgfpathlineto{\pgfqpoint{3.293768in}{1.368904in}}%
\pgfusepath{stroke}%
\end{pgfscope}%
\begin{pgfscope}%
\pgfsetrectcap%
\pgfsetmiterjoin%
\pgfsetlinewidth{0.803000pt}%
\definecolor{currentstroke}{rgb}{0.000000,0.000000,0.000000}%
\pgfsetstrokecolor{currentstroke}%
\pgfsetdash{}{0pt}%
\pgfpathmoveto{\pgfqpoint{0.573768in}{0.462654in}}%
\pgfpathlineto{\pgfqpoint{3.293768in}{0.462654in}}%
\pgfusepath{stroke}%
\end{pgfscope}%
\begin{pgfscope}%
\pgfsetrectcap%
\pgfsetmiterjoin%
\pgfsetlinewidth{0.803000pt}%
\definecolor{currentstroke}{rgb}{0.000000,0.000000,0.000000}%
\pgfsetstrokecolor{currentstroke}%
\pgfsetdash{}{0pt}%
\pgfpathmoveto{\pgfqpoint{0.573768in}{1.368904in}}%
\pgfpathlineto{\pgfqpoint{3.293768in}{1.368904in}}%
\pgfusepath{stroke}%
\end{pgfscope}%
\begin{pgfscope}%
\definecolor{textcolor}{rgb}{0.000000,0.000000,0.000000}%
\pgfsetstrokecolor{textcolor}%
\pgfsetfillcolor{textcolor}%
\pgftext[x=0.628168in,y=1.305467in,left,top]{\color{textcolor}{\rmfamily\fontsize{8.000000}{9.600000}\selectfont\catcode`\^=\active\def^{\ifmmode\sp\else\^{}\fi}\catcode`\%=\active\def%{\%}(c)}}%
\end{pgfscope}%
\begin{pgfscope}%
\pgfsetbuttcap%
\pgfsetmiterjoin%
\definecolor{currentfill}{rgb}{1.000000,1.000000,1.000000}%
\pgfsetfillcolor{currentfill}%
\pgfsetlinewidth{0.000000pt}%
\definecolor{currentstroke}{rgb}{0.000000,0.000000,0.000000}%
\pgfsetstrokecolor{currentstroke}%
\pgfsetstrokeopacity{0.000000}%
\pgfsetdash{}{0pt}%
\pgfpathmoveto{\pgfqpoint{2.197902in}{3.064738in}}%
\pgfpathlineto{\pgfqpoint{3.177102in}{3.064738in}}%
\pgfpathlineto{\pgfqpoint{3.177102in}{3.390988in}}%
\pgfpathlineto{\pgfqpoint{2.197902in}{3.390988in}}%
\pgfpathlineto{\pgfqpoint{2.197902in}{3.064738in}}%
\pgfpathclose%
\pgfusepath{fill}%
\end{pgfscope}%
\begin{pgfscope}%
\pgfpathrectangle{\pgfqpoint{2.197902in}{3.064738in}}{\pgfqpoint{0.979200in}{0.326250in}}%
\pgfusepath{clip}%
\pgfsetbuttcap%
\pgfsetroundjoin%
\pgfsetlinewidth{0.401500pt}%
\definecolor{currentstroke}{rgb}{0.690196,0.690196,0.690196}%
\pgfsetstrokecolor{currentstroke}%
\pgfsetstrokeopacity{0.250000}%
\pgfsetdash{{1.480000pt}{0.640000pt}}{0.000000pt}%
\pgfpathmoveto{\pgfqpoint{2.197902in}{3.064738in}}%
\pgfpathlineto{\pgfqpoint{2.197902in}{3.390988in}}%
\pgfusepath{stroke}%
\end{pgfscope}%
\begin{pgfscope}%
\pgfsetbuttcap%
\pgfsetroundjoin%
\definecolor{currentfill}{rgb}{0.000000,0.000000,0.000000}%
\pgfsetfillcolor{currentfill}%
\pgfsetlinewidth{0.803000pt}%
\definecolor{currentstroke}{rgb}{0.000000,0.000000,0.000000}%
\pgfsetstrokecolor{currentstroke}%
\pgfsetdash{}{0pt}%
\pgfsys@defobject{currentmarker}{\pgfqpoint{0.000000in}{0.000000in}}{\pgfqpoint{0.000000in}{0.027778in}}{%
\pgfpathmoveto{\pgfqpoint{0.000000in}{0.000000in}}%
\pgfpathlineto{\pgfqpoint{0.000000in}{0.027778in}}%
\pgfusepath{stroke,fill}%
}%
\begin{pgfscope}%
\pgfsys@transformshift{2.197902in}{3.064738in}%
\pgfsys@useobject{currentmarker}{}%
\end{pgfscope}%
\end{pgfscope}%
\begin{pgfscope}%
\definecolor{textcolor}{rgb}{0.000000,0.000000,0.000000}%
\pgfsetstrokecolor{textcolor}%
\pgfsetfillcolor{textcolor}%
\pgftext[x=2.197902in,y=3.050849in,,top]{\color{textcolor}{\rmfamily\fontsize{5.000000}{6.000000}\selectfont\catcode`\^=\active\def^{\ifmmode\sp\else\^{}\fi}\catcode`\%=\active\def%{\%}4.980}}%
\end{pgfscope}%
\begin{pgfscope}%
\pgfpathrectangle{\pgfqpoint{2.197902in}{3.064738in}}{\pgfqpoint{0.979200in}{0.326250in}}%
\pgfusepath{clip}%
\pgfsetbuttcap%
\pgfsetroundjoin%
\pgfsetlinewidth{0.401500pt}%
\definecolor{currentstroke}{rgb}{0.690196,0.690196,0.690196}%
\pgfsetstrokecolor{currentstroke}%
\pgfsetstrokeopacity{0.250000}%
\pgfsetdash{{1.480000pt}{0.640000pt}}{0.000000pt}%
\pgfpathmoveto{\pgfqpoint{2.442702in}{3.064738in}}%
\pgfpathlineto{\pgfqpoint{2.442702in}{3.390988in}}%
\pgfusepath{stroke}%
\end{pgfscope}%
\begin{pgfscope}%
\pgfsetbuttcap%
\pgfsetroundjoin%
\definecolor{currentfill}{rgb}{0.000000,0.000000,0.000000}%
\pgfsetfillcolor{currentfill}%
\pgfsetlinewidth{0.803000pt}%
\definecolor{currentstroke}{rgb}{0.000000,0.000000,0.000000}%
\pgfsetstrokecolor{currentstroke}%
\pgfsetdash{}{0pt}%
\pgfsys@defobject{currentmarker}{\pgfqpoint{0.000000in}{0.000000in}}{\pgfqpoint{0.000000in}{0.027778in}}{%
\pgfpathmoveto{\pgfqpoint{0.000000in}{0.000000in}}%
\pgfpathlineto{\pgfqpoint{0.000000in}{0.027778in}}%
\pgfusepath{stroke,fill}%
}%
\begin{pgfscope}%
\pgfsys@transformshift{2.442702in}{3.064738in}%
\pgfsys@useobject{currentmarker}{}%
\end{pgfscope}%
\end{pgfscope}%
\begin{pgfscope}%
\definecolor{textcolor}{rgb}{0.000000,0.000000,0.000000}%
\pgfsetstrokecolor{textcolor}%
\pgfsetfillcolor{textcolor}%
\pgftext[x=2.442702in,y=3.050849in,,top]{\color{textcolor}{\rmfamily\fontsize{5.000000}{6.000000}\selectfont\catcode`\^=\active\def^{\ifmmode\sp\else\^{}\fi}\catcode`\%=\active\def%{\%}4.985}}%
\end{pgfscope}%
\begin{pgfscope}%
\pgfpathrectangle{\pgfqpoint{2.197902in}{3.064738in}}{\pgfqpoint{0.979200in}{0.326250in}}%
\pgfusepath{clip}%
\pgfsetbuttcap%
\pgfsetroundjoin%
\pgfsetlinewidth{0.401500pt}%
\definecolor{currentstroke}{rgb}{0.690196,0.690196,0.690196}%
\pgfsetstrokecolor{currentstroke}%
\pgfsetstrokeopacity{0.250000}%
\pgfsetdash{{1.480000pt}{0.640000pt}}{0.000000pt}%
\pgfpathmoveto{\pgfqpoint{2.687502in}{3.064738in}}%
\pgfpathlineto{\pgfqpoint{2.687502in}{3.390988in}}%
\pgfusepath{stroke}%
\end{pgfscope}%
\begin{pgfscope}%
\pgfsetbuttcap%
\pgfsetroundjoin%
\definecolor{currentfill}{rgb}{0.000000,0.000000,0.000000}%
\pgfsetfillcolor{currentfill}%
\pgfsetlinewidth{0.803000pt}%
\definecolor{currentstroke}{rgb}{0.000000,0.000000,0.000000}%
\pgfsetstrokecolor{currentstroke}%
\pgfsetdash{}{0pt}%
\pgfsys@defobject{currentmarker}{\pgfqpoint{0.000000in}{0.000000in}}{\pgfqpoint{0.000000in}{0.027778in}}{%
\pgfpathmoveto{\pgfqpoint{0.000000in}{0.000000in}}%
\pgfpathlineto{\pgfqpoint{0.000000in}{0.027778in}}%
\pgfusepath{stroke,fill}%
}%
\begin{pgfscope}%
\pgfsys@transformshift{2.687502in}{3.064738in}%
\pgfsys@useobject{currentmarker}{}%
\end{pgfscope}%
\end{pgfscope}%
\begin{pgfscope}%
\definecolor{textcolor}{rgb}{0.000000,0.000000,0.000000}%
\pgfsetstrokecolor{textcolor}%
\pgfsetfillcolor{textcolor}%
\pgftext[x=2.687502in,y=3.050849in,,top]{\color{textcolor}{\rmfamily\fontsize{5.000000}{6.000000}\selectfont\catcode`\^=\active\def^{\ifmmode\sp\else\^{}\fi}\catcode`\%=\active\def%{\%}4.990}}%
\end{pgfscope}%
\begin{pgfscope}%
\pgfpathrectangle{\pgfqpoint{2.197902in}{3.064738in}}{\pgfqpoint{0.979200in}{0.326250in}}%
\pgfusepath{clip}%
\pgfsetbuttcap%
\pgfsetroundjoin%
\pgfsetlinewidth{0.401500pt}%
\definecolor{currentstroke}{rgb}{0.690196,0.690196,0.690196}%
\pgfsetstrokecolor{currentstroke}%
\pgfsetstrokeopacity{0.250000}%
\pgfsetdash{{1.480000pt}{0.640000pt}}{0.000000pt}%
\pgfpathmoveto{\pgfqpoint{2.932302in}{3.064738in}}%
\pgfpathlineto{\pgfqpoint{2.932302in}{3.390988in}}%
\pgfusepath{stroke}%
\end{pgfscope}%
\begin{pgfscope}%
\pgfsetbuttcap%
\pgfsetroundjoin%
\definecolor{currentfill}{rgb}{0.000000,0.000000,0.000000}%
\pgfsetfillcolor{currentfill}%
\pgfsetlinewidth{0.803000pt}%
\definecolor{currentstroke}{rgb}{0.000000,0.000000,0.000000}%
\pgfsetstrokecolor{currentstroke}%
\pgfsetdash{}{0pt}%
\pgfsys@defobject{currentmarker}{\pgfqpoint{0.000000in}{0.000000in}}{\pgfqpoint{0.000000in}{0.027778in}}{%
\pgfpathmoveto{\pgfqpoint{0.000000in}{0.000000in}}%
\pgfpathlineto{\pgfqpoint{0.000000in}{0.027778in}}%
\pgfusepath{stroke,fill}%
}%
\begin{pgfscope}%
\pgfsys@transformshift{2.932302in}{3.064738in}%
\pgfsys@useobject{currentmarker}{}%
\end{pgfscope}%
\end{pgfscope}%
\begin{pgfscope}%
\definecolor{textcolor}{rgb}{0.000000,0.000000,0.000000}%
\pgfsetstrokecolor{textcolor}%
\pgfsetfillcolor{textcolor}%
\pgftext[x=2.932302in,y=3.050849in,,top]{\color{textcolor}{\rmfamily\fontsize{5.000000}{6.000000}\selectfont\catcode`\^=\active\def^{\ifmmode\sp\else\^{}\fi}\catcode`\%=\active\def%{\%}4.995}}%
\end{pgfscope}%
\begin{pgfscope}%
\pgfpathrectangle{\pgfqpoint{2.197902in}{3.064738in}}{\pgfqpoint{0.979200in}{0.326250in}}%
\pgfusepath{clip}%
\pgfsetbuttcap%
\pgfsetroundjoin%
\pgfsetlinewidth{0.401500pt}%
\definecolor{currentstroke}{rgb}{0.690196,0.690196,0.690196}%
\pgfsetstrokecolor{currentstroke}%
\pgfsetstrokeopacity{0.250000}%
\pgfsetdash{{1.480000pt}{0.640000pt}}{0.000000pt}%
\pgfpathmoveto{\pgfqpoint{3.177102in}{3.064738in}}%
\pgfpathlineto{\pgfqpoint{3.177102in}{3.390988in}}%
\pgfusepath{stroke}%
\end{pgfscope}%
\begin{pgfscope}%
\pgfsetbuttcap%
\pgfsetroundjoin%
\definecolor{currentfill}{rgb}{0.000000,0.000000,0.000000}%
\pgfsetfillcolor{currentfill}%
\pgfsetlinewidth{0.803000pt}%
\definecolor{currentstroke}{rgb}{0.000000,0.000000,0.000000}%
\pgfsetstrokecolor{currentstroke}%
\pgfsetdash{}{0pt}%
\pgfsys@defobject{currentmarker}{\pgfqpoint{0.000000in}{0.000000in}}{\pgfqpoint{0.000000in}{0.027778in}}{%
\pgfpathmoveto{\pgfqpoint{0.000000in}{0.000000in}}%
\pgfpathlineto{\pgfqpoint{0.000000in}{0.027778in}}%
\pgfusepath{stroke,fill}%
}%
\begin{pgfscope}%
\pgfsys@transformshift{3.177102in}{3.064738in}%
\pgfsys@useobject{currentmarker}{}%
\end{pgfscope}%
\end{pgfscope}%
\begin{pgfscope}%
\definecolor{textcolor}{rgb}{0.000000,0.000000,0.000000}%
\pgfsetstrokecolor{textcolor}%
\pgfsetfillcolor{textcolor}%
\pgftext[x=3.177102in,y=3.050849in,,top]{\color{textcolor}{\rmfamily\fontsize{5.000000}{6.000000}\selectfont\catcode`\^=\active\def^{\ifmmode\sp\else\^{}\fi}\catcode`\%=\active\def%{\%}5.000}}%
\end{pgfscope}%
\begin{pgfscope}%
\pgfpathrectangle{\pgfqpoint{2.197902in}{3.064738in}}{\pgfqpoint{0.979200in}{0.326250in}}%
\pgfusepath{clip}%
\pgfsetbuttcap%
\pgfsetroundjoin%
\pgfsetlinewidth{0.401500pt}%
\definecolor{currentstroke}{rgb}{0.690196,0.690196,0.690196}%
\pgfsetstrokecolor{currentstroke}%
\pgfsetstrokeopacity{0.250000}%
\pgfsetdash{{1.480000pt}{0.640000pt}}{0.000000pt}%
\pgfpathmoveto{\pgfqpoint{2.197902in}{3.064738in}}%
\pgfpathlineto{\pgfqpoint{3.177102in}{3.064738in}}%
\pgfusepath{stroke}%
\end{pgfscope}%
\begin{pgfscope}%
\pgfsetbuttcap%
\pgfsetroundjoin%
\definecolor{currentfill}{rgb}{0.000000,0.000000,0.000000}%
\pgfsetfillcolor{currentfill}%
\pgfsetlinewidth{0.803000pt}%
\definecolor{currentstroke}{rgb}{0.000000,0.000000,0.000000}%
\pgfsetstrokecolor{currentstroke}%
\pgfsetdash{}{0pt}%
\pgfsys@defobject{currentmarker}{\pgfqpoint{0.000000in}{0.000000in}}{\pgfqpoint{0.027778in}{0.000000in}}{%
\pgfpathmoveto{\pgfqpoint{0.000000in}{0.000000in}}%
\pgfpathlineto{\pgfqpoint{0.027778in}{0.000000in}}%
\pgfusepath{stroke,fill}%
}%
\begin{pgfscope}%
\pgfsys@transformshift{2.197902in}{3.064738in}%
\pgfsys@useobject{currentmarker}{}%
\end{pgfscope}%
\end{pgfscope}%
\begin{pgfscope}%
\definecolor{textcolor}{rgb}{0.000000,0.000000,0.000000}%
\pgfsetstrokecolor{textcolor}%
\pgfsetfillcolor{textcolor}%
\pgftext[x=1.981463in, y=3.040625in, left, base]{\color{textcolor}{\rmfamily\fontsize{5.000000}{6.000000}\selectfont\catcode`\^=\active\def^{\ifmmode\sp\else\^{}\fi}\catcode`\%=\active\def%{\%}-0.50}}%
\end{pgfscope}%
\begin{pgfscope}%
\pgfpathrectangle{\pgfqpoint{2.197902in}{3.064738in}}{\pgfqpoint{0.979200in}{0.326250in}}%
\pgfusepath{clip}%
\pgfsetbuttcap%
\pgfsetroundjoin%
\pgfsetlinewidth{0.401500pt}%
\definecolor{currentstroke}{rgb}{0.690196,0.690196,0.690196}%
\pgfsetstrokecolor{currentstroke}%
\pgfsetstrokeopacity{0.250000}%
\pgfsetdash{{1.480000pt}{0.640000pt}}{0.000000pt}%
\pgfpathmoveto{\pgfqpoint{2.197902in}{3.227863in}}%
\pgfpathlineto{\pgfqpoint{3.177102in}{3.227863in}}%
\pgfusepath{stroke}%
\end{pgfscope}%
\begin{pgfscope}%
\pgfsetbuttcap%
\pgfsetroundjoin%
\definecolor{currentfill}{rgb}{0.000000,0.000000,0.000000}%
\pgfsetfillcolor{currentfill}%
\pgfsetlinewidth{0.803000pt}%
\definecolor{currentstroke}{rgb}{0.000000,0.000000,0.000000}%
\pgfsetstrokecolor{currentstroke}%
\pgfsetdash{}{0pt}%
\pgfsys@defobject{currentmarker}{\pgfqpoint{0.000000in}{0.000000in}}{\pgfqpoint{0.027778in}{0.000000in}}{%
\pgfpathmoveto{\pgfqpoint{0.000000in}{0.000000in}}%
\pgfpathlineto{\pgfqpoint{0.027778in}{0.000000in}}%
\pgfusepath{stroke,fill}%
}%
\begin{pgfscope}%
\pgfsys@transformshift{2.197902in}{3.227863in}%
\pgfsys@useobject{currentmarker}{}%
\end{pgfscope}%
\end{pgfscope}%
\begin{pgfscope}%
\definecolor{textcolor}{rgb}{0.000000,0.000000,0.000000}%
\pgfsetstrokecolor{textcolor}%
\pgfsetfillcolor{textcolor}%
\pgftext[x=2.014257in, y=3.203750in, left, base]{\color{textcolor}{\rmfamily\fontsize{5.000000}{6.000000}\selectfont\catcode`\^=\active\def^{\ifmmode\sp\else\^{}\fi}\catcode`\%=\active\def%{\%}0.73}}%
\end{pgfscope}%
\begin{pgfscope}%
\pgfpathrectangle{\pgfqpoint{2.197902in}{3.064738in}}{\pgfqpoint{0.979200in}{0.326250in}}%
\pgfusepath{clip}%
\pgfsetbuttcap%
\pgfsetroundjoin%
\pgfsetlinewidth{0.401500pt}%
\definecolor{currentstroke}{rgb}{0.690196,0.690196,0.690196}%
\pgfsetstrokecolor{currentstroke}%
\pgfsetstrokeopacity{0.250000}%
\pgfsetdash{{1.480000pt}{0.640000pt}}{0.000000pt}%
\pgfpathmoveto{\pgfqpoint{2.197902in}{3.390988in}}%
\pgfpathlineto{\pgfqpoint{3.177102in}{3.390988in}}%
\pgfusepath{stroke}%
\end{pgfscope}%
\begin{pgfscope}%
\pgfsetbuttcap%
\pgfsetroundjoin%
\definecolor{currentfill}{rgb}{0.000000,0.000000,0.000000}%
\pgfsetfillcolor{currentfill}%
\pgfsetlinewidth{0.803000pt}%
\definecolor{currentstroke}{rgb}{0.000000,0.000000,0.000000}%
\pgfsetstrokecolor{currentstroke}%
\pgfsetdash{}{0pt}%
\pgfsys@defobject{currentmarker}{\pgfqpoint{0.000000in}{0.000000in}}{\pgfqpoint{0.027778in}{0.000000in}}{%
\pgfpathmoveto{\pgfqpoint{0.000000in}{0.000000in}}%
\pgfpathlineto{\pgfqpoint{0.027778in}{0.000000in}}%
\pgfusepath{stroke,fill}%
}%
\begin{pgfscope}%
\pgfsys@transformshift{2.197902in}{3.390988in}%
\pgfsys@useobject{currentmarker}{}%
\end{pgfscope}%
\end{pgfscope}%
\begin{pgfscope}%
\definecolor{textcolor}{rgb}{0.000000,0.000000,0.000000}%
\pgfsetstrokecolor{textcolor}%
\pgfsetfillcolor{textcolor}%
\pgftext[x=2.014257in, y=3.366875in, left, base]{\color{textcolor}{\rmfamily\fontsize{5.000000}{6.000000}\selectfont\catcode`\^=\active\def^{\ifmmode\sp\else\^{}\fi}\catcode`\%=\active\def%{\%}1.95}}%
\end{pgfscope}%
\begin{pgfscope}%
\pgfpathrectangle{\pgfqpoint{2.197902in}{3.064738in}}{\pgfqpoint{0.979200in}{0.326250in}}%
\pgfusepath{clip}%
\pgfsetrectcap%
\pgfsetroundjoin%
\pgfsetlinewidth{1.505625pt}%
\definecolor{currentstroke}{rgb}{0.000000,0.447059,0.698039}%
\pgfsetstrokecolor{currentstroke}%
\pgfsetdash{}{0pt}%
\pgfpathmoveto{\pgfqpoint{2.194661in}{3.131319in}}%
\pgfpathlineto{\pgfqpoint{2.198295in}{3.131319in}}%
\pgfpathlineto{\pgfqpoint{2.198917in}{3.295610in}}%
\pgfpathlineto{\pgfqpoint{2.199349in}{3.228735in}}%
\pgfpathlineto{\pgfqpoint{2.199548in}{3.240896in}}%
\pgfpathlineto{\pgfqpoint{2.199823in}{3.236207in}}%
\pgfpathlineto{\pgfqpoint{2.200147in}{3.212297in}}%
\pgfpathlineto{\pgfqpoint{2.200858in}{3.129984in}}%
\pgfpathlineto{\pgfqpoint{2.201299in}{3.197472in}}%
\pgfpathlineto{\pgfqpoint{2.202180in}{3.121558in}}%
\pgfpathlineto{\pgfqpoint{2.203943in}{3.185677in}}%
\pgfpathlineto{\pgfqpoint{2.207470in}{3.111937in}}%
\pgfpathlineto{\pgfqpoint{2.211725in}{3.184612in}}%
\pgfpathlineto{\pgfqpoint{2.216621in}{3.118837in}}%
\pgfpathlineto{\pgfqpoint{2.221517in}{3.193868in}}%
\pgfpathlineto{\pgfqpoint{2.226413in}{3.128428in}}%
\pgfpathlineto{\pgfqpoint{2.231309in}{3.203421in}}%
\pgfpathlineto{\pgfqpoint{2.236205in}{3.138041in}}%
\pgfpathlineto{\pgfqpoint{2.241101in}{3.212988in}}%
\pgfpathlineto{\pgfqpoint{2.245997in}{3.147654in}}%
\pgfpathlineto{\pgfqpoint{2.250893in}{3.222549in}}%
\pgfpathlineto{\pgfqpoint{2.255789in}{3.157267in}}%
\pgfpathlineto{\pgfqpoint{2.260685in}{3.232111in}}%
\pgfpathlineto{\pgfqpoint{2.265581in}{3.166880in}}%
\pgfpathlineto{\pgfqpoint{2.270477in}{3.241672in}}%
\pgfpathlineto{\pgfqpoint{2.275373in}{3.176493in}}%
\pgfpathlineto{\pgfqpoint{2.280269in}{3.251233in}}%
\pgfpathlineto{\pgfqpoint{2.285165in}{3.186106in}}%
\pgfpathlineto{\pgfqpoint{2.290061in}{3.260794in}}%
\pgfpathlineto{\pgfqpoint{2.294957in}{3.195718in}}%
\pgfpathlineto{\pgfqpoint{2.299853in}{3.270354in}}%
\pgfpathlineto{\pgfqpoint{2.304749in}{3.205331in}}%
\pgfpathlineto{\pgfqpoint{2.309645in}{3.279915in}}%
\pgfpathlineto{\pgfqpoint{2.314541in}{3.214943in}}%
\pgfpathlineto{\pgfqpoint{2.319437in}{3.289476in}}%
\pgfpathlineto{\pgfqpoint{2.324333in}{3.224556in}}%
\pgfpathlineto{\pgfqpoint{2.329229in}{3.299037in}}%
\pgfpathlineto{\pgfqpoint{2.334125in}{3.234168in}}%
\pgfpathlineto{\pgfqpoint{2.339021in}{3.308597in}}%
\pgfpathlineto{\pgfqpoint{2.343917in}{3.243780in}}%
\pgfpathlineto{\pgfqpoint{2.348813in}{3.318158in}}%
\pgfpathlineto{\pgfqpoint{2.353709in}{3.253392in}}%
\pgfpathlineto{\pgfqpoint{2.358605in}{3.327718in}}%
\pgfpathlineto{\pgfqpoint{2.363501in}{3.263004in}}%
\pgfpathlineto{\pgfqpoint{2.368397in}{3.337278in}}%
\pgfpathlineto{\pgfqpoint{2.373293in}{3.272615in}}%
\pgfpathlineto{\pgfqpoint{2.378189in}{3.346839in}}%
\pgfpathlineto{\pgfqpoint{2.383085in}{3.282227in}}%
\pgfpathlineto{\pgfqpoint{2.387981in}{3.356399in}}%
\pgfpathlineto{\pgfqpoint{2.390886in}{3.289898in}}%
\pgfpathlineto{\pgfqpoint{2.393791in}{3.362043in}}%
\pgfpathlineto{\pgfqpoint{2.393846in}{3.270712in}}%
\pgfpathlineto{\pgfqpoint{2.394372in}{3.280865in}}%
\pgfpathlineto{\pgfqpoint{2.395976in}{3.302400in}}%
\pgfpathlineto{\pgfqpoint{2.397029in}{3.269420in}}%
\pgfpathlineto{\pgfqpoint{2.398004in}{3.265754in}}%
\pgfpathlineto{\pgfqpoint{2.399228in}{3.243219in}}%
\pgfpathlineto{\pgfqpoint{2.400619in}{3.261065in}}%
\pgfpathlineto{\pgfqpoint{2.402197in}{3.131319in}}%
\pgfpathlineto{\pgfqpoint{2.688108in}{3.131589in}}%
\pgfpathlineto{\pgfqpoint{2.688320in}{3.332888in}}%
\pgfpathlineto{\pgfqpoint{2.688992in}{3.233110in}}%
\pgfpathlineto{\pgfqpoint{2.689487in}{3.243743in}}%
\pgfpathlineto{\pgfqpoint{2.690201in}{3.165797in}}%
\pgfpathlineto{\pgfqpoint{2.690562in}{3.165720in}}%
\pgfpathlineto{\pgfqpoint{2.691282in}{3.157617in}}%
\pgfpathlineto{\pgfqpoint{2.695603in}{3.145900in}}%
\pgfpathlineto{\pgfqpoint{2.699699in}{3.149726in}}%
\pgfpathlineto{\pgfqpoint{2.704595in}{3.150861in}}%
\pgfpathlineto{\pgfqpoint{2.709491in}{3.158680in}}%
\pgfpathlineto{\pgfqpoint{2.714387in}{3.160406in}}%
\pgfpathlineto{\pgfqpoint{2.719283in}{3.168264in}}%
\pgfpathlineto{\pgfqpoint{2.724179in}{3.169985in}}%
\pgfpathlineto{\pgfqpoint{2.729075in}{3.177853in}}%
\pgfpathlineto{\pgfqpoint{2.733971in}{3.179571in}}%
\pgfpathlineto{\pgfqpoint{2.738867in}{3.187441in}}%
\pgfpathlineto{\pgfqpoint{2.743763in}{3.189158in}}%
\pgfpathlineto{\pgfqpoint{2.748659in}{3.197029in}}%
\pgfpathlineto{\pgfqpoint{2.753555in}{3.198744in}}%
\pgfpathlineto{\pgfqpoint{2.758451in}{3.206617in}}%
\pgfpathlineto{\pgfqpoint{2.763347in}{3.208330in}}%
\pgfpathlineto{\pgfqpoint{2.768243in}{3.216205in}}%
\pgfpathlineto{\pgfqpoint{2.773139in}{3.217916in}}%
\pgfpathlineto{\pgfqpoint{2.778035in}{3.225793in}}%
\pgfpathlineto{\pgfqpoint{2.782931in}{3.227502in}}%
\pgfpathlineto{\pgfqpoint{2.787827in}{3.235381in}}%
\pgfpathlineto{\pgfqpoint{2.792723in}{3.237087in}}%
\pgfpathlineto{\pgfqpoint{2.797619in}{3.244968in}}%
\pgfpathlineto{\pgfqpoint{2.802515in}{3.246673in}}%
\pgfpathlineto{\pgfqpoint{2.807411in}{3.254556in}}%
\pgfpathlineto{\pgfqpoint{2.812307in}{3.256259in}}%
\pgfpathlineto{\pgfqpoint{2.817203in}{3.264143in}}%
\pgfpathlineto{\pgfqpoint{2.822099in}{3.265844in}}%
\pgfpathlineto{\pgfqpoint{2.826995in}{3.273730in}}%
\pgfpathlineto{\pgfqpoint{2.831891in}{3.275430in}}%
\pgfpathlineto{\pgfqpoint{2.836787in}{3.283318in}}%
\pgfpathlineto{\pgfqpoint{2.841683in}{3.285015in}}%
\pgfpathlineto{\pgfqpoint{2.846579in}{3.292905in}}%
\pgfpathlineto{\pgfqpoint{2.851475in}{3.294600in}}%
\pgfpathlineto{\pgfqpoint{2.856371in}{3.302492in}}%
\pgfpathlineto{\pgfqpoint{2.861267in}{3.304185in}}%
\pgfpathlineto{\pgfqpoint{2.866163in}{3.312079in}}%
\pgfpathlineto{\pgfqpoint{2.871059in}{3.313770in}}%
\pgfpathlineto{\pgfqpoint{2.875955in}{3.321665in}}%
\pgfpathlineto{\pgfqpoint{2.879673in}{3.322201in}}%
\pgfpathlineto{\pgfqpoint{2.883391in}{3.328945in}}%
\pgfpathlineto{\pgfqpoint{2.883443in}{3.270861in}}%
\pgfpathlineto{\pgfqpoint{2.883875in}{3.279242in}}%
\pgfpathlineto{\pgfqpoint{2.886070in}{3.306215in}}%
\pgfpathlineto{\pgfqpoint{2.886935in}{3.240673in}}%
\pgfpathlineto{\pgfqpoint{2.887823in}{3.293898in}}%
\pgfpathlineto{\pgfqpoint{2.888866in}{3.211713in}}%
\pgfpathlineto{\pgfqpoint{2.890336in}{3.291252in}}%
\pgfpathlineto{\pgfqpoint{2.891643in}{3.131319in}}%
\pgfpathlineto{\pgfqpoint{3.177102in}{3.131319in}}%
\pgfpathlineto{\pgfqpoint{3.177102in}{3.131319in}}%
\pgfusepath{stroke}%
\end{pgfscope}%
\begin{pgfscope}%
\pgfsetrectcap%
\pgfsetmiterjoin%
\pgfsetlinewidth{0.803000pt}%
\definecolor{currentstroke}{rgb}{0.000000,0.000000,0.000000}%
\pgfsetstrokecolor{currentstroke}%
\pgfsetdash{}{0pt}%
\pgfpathmoveto{\pgfqpoint{2.197902in}{3.064738in}}%
\pgfpathlineto{\pgfqpoint{2.197902in}{3.390988in}}%
\pgfusepath{stroke}%
\end{pgfscope}%
\begin{pgfscope}%
\pgfsetrectcap%
\pgfsetmiterjoin%
\pgfsetlinewidth{0.803000pt}%
\definecolor{currentstroke}{rgb}{0.000000,0.000000,0.000000}%
\pgfsetstrokecolor{currentstroke}%
\pgfsetdash{}{0pt}%
\pgfpathmoveto{\pgfqpoint{3.177102in}{3.064738in}}%
\pgfpathlineto{\pgfqpoint{3.177102in}{3.390988in}}%
\pgfusepath{stroke}%
\end{pgfscope}%
\begin{pgfscope}%
\pgfsetrectcap%
\pgfsetmiterjoin%
\pgfsetlinewidth{0.803000pt}%
\definecolor{currentstroke}{rgb}{0.000000,0.000000,0.000000}%
\pgfsetstrokecolor{currentstroke}%
\pgfsetdash{}{0pt}%
\pgfpathmoveto{\pgfqpoint{2.197902in}{3.064738in}}%
\pgfpathlineto{\pgfqpoint{3.177102in}{3.064738in}}%
\pgfusepath{stroke}%
\end{pgfscope}%
\begin{pgfscope}%
\pgfsetrectcap%
\pgfsetmiterjoin%
\pgfsetlinewidth{0.803000pt}%
\definecolor{currentstroke}{rgb}{0.000000,0.000000,0.000000}%
\pgfsetstrokecolor{currentstroke}%
\pgfsetdash{}{0pt}%
\pgfpathmoveto{\pgfqpoint{2.197902in}{3.390988in}}%
\pgfpathlineto{\pgfqpoint{3.177102in}{3.390988in}}%
\pgfusepath{stroke}%
\end{pgfscope}%
\begin{pgfscope}%
\pgfsetbuttcap%
\pgfsetmiterjoin%
\definecolor{currentfill}{rgb}{1.000000,1.000000,1.000000}%
\pgfsetfillcolor{currentfill}%
\pgfsetlinewidth{0.000000pt}%
\definecolor{currentstroke}{rgb}{0.000000,0.000000,0.000000}%
\pgfsetstrokecolor{currentstroke}%
\pgfsetstrokeopacity{0.000000}%
\pgfsetdash{}{0pt}%
\pgfpathmoveto{\pgfqpoint{2.197902in}{1.995363in}}%
\pgfpathlineto{\pgfqpoint{3.177102in}{1.995363in}}%
\pgfpathlineto{\pgfqpoint{3.177102in}{2.321613in}}%
\pgfpathlineto{\pgfqpoint{2.197902in}{2.321613in}}%
\pgfpathlineto{\pgfqpoint{2.197902in}{1.995363in}}%
\pgfpathclose%
\pgfusepath{fill}%
\end{pgfscope}%
\begin{pgfscope}%
\pgfpathrectangle{\pgfqpoint{2.197902in}{1.995363in}}{\pgfqpoint{0.979200in}{0.326250in}}%
\pgfusepath{clip}%
\pgfsetbuttcap%
\pgfsetroundjoin%
\pgfsetlinewidth{0.401500pt}%
\definecolor{currentstroke}{rgb}{0.690196,0.690196,0.690196}%
\pgfsetstrokecolor{currentstroke}%
\pgfsetstrokeopacity{0.250000}%
\pgfsetdash{{1.480000pt}{0.640000pt}}{0.000000pt}%
\pgfpathmoveto{\pgfqpoint{2.197902in}{1.995363in}}%
\pgfpathlineto{\pgfqpoint{2.197902in}{2.321613in}}%
\pgfusepath{stroke}%
\end{pgfscope}%
\begin{pgfscope}%
\pgfsetbuttcap%
\pgfsetroundjoin%
\definecolor{currentfill}{rgb}{0.000000,0.000000,0.000000}%
\pgfsetfillcolor{currentfill}%
\pgfsetlinewidth{0.803000pt}%
\definecolor{currentstroke}{rgb}{0.000000,0.000000,0.000000}%
\pgfsetstrokecolor{currentstroke}%
\pgfsetdash{}{0pt}%
\pgfsys@defobject{currentmarker}{\pgfqpoint{0.000000in}{0.000000in}}{\pgfqpoint{0.000000in}{0.027778in}}{%
\pgfpathmoveto{\pgfqpoint{0.000000in}{0.000000in}}%
\pgfpathlineto{\pgfqpoint{0.000000in}{0.027778in}}%
\pgfusepath{stroke,fill}%
}%
\begin{pgfscope}%
\pgfsys@transformshift{2.197902in}{1.995363in}%
\pgfsys@useobject{currentmarker}{}%
\end{pgfscope}%
\end{pgfscope}%
\begin{pgfscope}%
\definecolor{textcolor}{rgb}{0.000000,0.000000,0.000000}%
\pgfsetstrokecolor{textcolor}%
\pgfsetfillcolor{textcolor}%
\pgftext[x=2.197902in,y=1.981474in,,top]{\color{textcolor}{\rmfamily\fontsize{5.000000}{6.000000}\selectfont\catcode`\^=\active\def^{\ifmmode\sp\else\^{}\fi}\catcode`\%=\active\def%{\%}4.980}}%
\end{pgfscope}%
\begin{pgfscope}%
\pgfpathrectangle{\pgfqpoint{2.197902in}{1.995363in}}{\pgfqpoint{0.979200in}{0.326250in}}%
\pgfusepath{clip}%
\pgfsetbuttcap%
\pgfsetroundjoin%
\pgfsetlinewidth{0.401500pt}%
\definecolor{currentstroke}{rgb}{0.690196,0.690196,0.690196}%
\pgfsetstrokecolor{currentstroke}%
\pgfsetstrokeopacity{0.250000}%
\pgfsetdash{{1.480000pt}{0.640000pt}}{0.000000pt}%
\pgfpathmoveto{\pgfqpoint{2.442702in}{1.995363in}}%
\pgfpathlineto{\pgfqpoint{2.442702in}{2.321613in}}%
\pgfusepath{stroke}%
\end{pgfscope}%
\begin{pgfscope}%
\pgfsetbuttcap%
\pgfsetroundjoin%
\definecolor{currentfill}{rgb}{0.000000,0.000000,0.000000}%
\pgfsetfillcolor{currentfill}%
\pgfsetlinewidth{0.803000pt}%
\definecolor{currentstroke}{rgb}{0.000000,0.000000,0.000000}%
\pgfsetstrokecolor{currentstroke}%
\pgfsetdash{}{0pt}%
\pgfsys@defobject{currentmarker}{\pgfqpoint{0.000000in}{0.000000in}}{\pgfqpoint{0.000000in}{0.027778in}}{%
\pgfpathmoveto{\pgfqpoint{0.000000in}{0.000000in}}%
\pgfpathlineto{\pgfqpoint{0.000000in}{0.027778in}}%
\pgfusepath{stroke,fill}%
}%
\begin{pgfscope}%
\pgfsys@transformshift{2.442702in}{1.995363in}%
\pgfsys@useobject{currentmarker}{}%
\end{pgfscope}%
\end{pgfscope}%
\begin{pgfscope}%
\definecolor{textcolor}{rgb}{0.000000,0.000000,0.000000}%
\pgfsetstrokecolor{textcolor}%
\pgfsetfillcolor{textcolor}%
\pgftext[x=2.442702in,y=1.981474in,,top]{\color{textcolor}{\rmfamily\fontsize{5.000000}{6.000000}\selectfont\catcode`\^=\active\def^{\ifmmode\sp\else\^{}\fi}\catcode`\%=\active\def%{\%}4.985}}%
\end{pgfscope}%
\begin{pgfscope}%
\pgfpathrectangle{\pgfqpoint{2.197902in}{1.995363in}}{\pgfqpoint{0.979200in}{0.326250in}}%
\pgfusepath{clip}%
\pgfsetbuttcap%
\pgfsetroundjoin%
\pgfsetlinewidth{0.401500pt}%
\definecolor{currentstroke}{rgb}{0.690196,0.690196,0.690196}%
\pgfsetstrokecolor{currentstroke}%
\pgfsetstrokeopacity{0.250000}%
\pgfsetdash{{1.480000pt}{0.640000pt}}{0.000000pt}%
\pgfpathmoveto{\pgfqpoint{2.687502in}{1.995363in}}%
\pgfpathlineto{\pgfqpoint{2.687502in}{2.321613in}}%
\pgfusepath{stroke}%
\end{pgfscope}%
\begin{pgfscope}%
\pgfsetbuttcap%
\pgfsetroundjoin%
\definecolor{currentfill}{rgb}{0.000000,0.000000,0.000000}%
\pgfsetfillcolor{currentfill}%
\pgfsetlinewidth{0.803000pt}%
\definecolor{currentstroke}{rgb}{0.000000,0.000000,0.000000}%
\pgfsetstrokecolor{currentstroke}%
\pgfsetdash{}{0pt}%
\pgfsys@defobject{currentmarker}{\pgfqpoint{0.000000in}{0.000000in}}{\pgfqpoint{0.000000in}{0.027778in}}{%
\pgfpathmoveto{\pgfqpoint{0.000000in}{0.000000in}}%
\pgfpathlineto{\pgfqpoint{0.000000in}{0.027778in}}%
\pgfusepath{stroke,fill}%
}%
\begin{pgfscope}%
\pgfsys@transformshift{2.687502in}{1.995363in}%
\pgfsys@useobject{currentmarker}{}%
\end{pgfscope}%
\end{pgfscope}%
\begin{pgfscope}%
\definecolor{textcolor}{rgb}{0.000000,0.000000,0.000000}%
\pgfsetstrokecolor{textcolor}%
\pgfsetfillcolor{textcolor}%
\pgftext[x=2.687502in,y=1.981474in,,top]{\color{textcolor}{\rmfamily\fontsize{5.000000}{6.000000}\selectfont\catcode`\^=\active\def^{\ifmmode\sp\else\^{}\fi}\catcode`\%=\active\def%{\%}4.990}}%
\end{pgfscope}%
\begin{pgfscope}%
\pgfpathrectangle{\pgfqpoint{2.197902in}{1.995363in}}{\pgfqpoint{0.979200in}{0.326250in}}%
\pgfusepath{clip}%
\pgfsetbuttcap%
\pgfsetroundjoin%
\pgfsetlinewidth{0.401500pt}%
\definecolor{currentstroke}{rgb}{0.690196,0.690196,0.690196}%
\pgfsetstrokecolor{currentstroke}%
\pgfsetstrokeopacity{0.250000}%
\pgfsetdash{{1.480000pt}{0.640000pt}}{0.000000pt}%
\pgfpathmoveto{\pgfqpoint{2.932302in}{1.995363in}}%
\pgfpathlineto{\pgfqpoint{2.932302in}{2.321613in}}%
\pgfusepath{stroke}%
\end{pgfscope}%
\begin{pgfscope}%
\pgfsetbuttcap%
\pgfsetroundjoin%
\definecolor{currentfill}{rgb}{0.000000,0.000000,0.000000}%
\pgfsetfillcolor{currentfill}%
\pgfsetlinewidth{0.803000pt}%
\definecolor{currentstroke}{rgb}{0.000000,0.000000,0.000000}%
\pgfsetstrokecolor{currentstroke}%
\pgfsetdash{}{0pt}%
\pgfsys@defobject{currentmarker}{\pgfqpoint{0.000000in}{0.000000in}}{\pgfqpoint{0.000000in}{0.027778in}}{%
\pgfpathmoveto{\pgfqpoint{0.000000in}{0.000000in}}%
\pgfpathlineto{\pgfqpoint{0.000000in}{0.027778in}}%
\pgfusepath{stroke,fill}%
}%
\begin{pgfscope}%
\pgfsys@transformshift{2.932302in}{1.995363in}%
\pgfsys@useobject{currentmarker}{}%
\end{pgfscope}%
\end{pgfscope}%
\begin{pgfscope}%
\definecolor{textcolor}{rgb}{0.000000,0.000000,0.000000}%
\pgfsetstrokecolor{textcolor}%
\pgfsetfillcolor{textcolor}%
\pgftext[x=2.932302in,y=1.981474in,,top]{\color{textcolor}{\rmfamily\fontsize{5.000000}{6.000000}\selectfont\catcode`\^=\active\def^{\ifmmode\sp\else\^{}\fi}\catcode`\%=\active\def%{\%}4.995}}%
\end{pgfscope}%
\begin{pgfscope}%
\pgfpathrectangle{\pgfqpoint{2.197902in}{1.995363in}}{\pgfqpoint{0.979200in}{0.326250in}}%
\pgfusepath{clip}%
\pgfsetbuttcap%
\pgfsetroundjoin%
\pgfsetlinewidth{0.401500pt}%
\definecolor{currentstroke}{rgb}{0.690196,0.690196,0.690196}%
\pgfsetstrokecolor{currentstroke}%
\pgfsetstrokeopacity{0.250000}%
\pgfsetdash{{1.480000pt}{0.640000pt}}{0.000000pt}%
\pgfpathmoveto{\pgfqpoint{3.177102in}{1.995363in}}%
\pgfpathlineto{\pgfqpoint{3.177102in}{2.321613in}}%
\pgfusepath{stroke}%
\end{pgfscope}%
\begin{pgfscope}%
\pgfsetbuttcap%
\pgfsetroundjoin%
\definecolor{currentfill}{rgb}{0.000000,0.000000,0.000000}%
\pgfsetfillcolor{currentfill}%
\pgfsetlinewidth{0.803000pt}%
\definecolor{currentstroke}{rgb}{0.000000,0.000000,0.000000}%
\pgfsetstrokecolor{currentstroke}%
\pgfsetdash{}{0pt}%
\pgfsys@defobject{currentmarker}{\pgfqpoint{0.000000in}{0.000000in}}{\pgfqpoint{0.000000in}{0.027778in}}{%
\pgfpathmoveto{\pgfqpoint{0.000000in}{0.000000in}}%
\pgfpathlineto{\pgfqpoint{0.000000in}{0.027778in}}%
\pgfusepath{stroke,fill}%
}%
\begin{pgfscope}%
\pgfsys@transformshift{3.177102in}{1.995363in}%
\pgfsys@useobject{currentmarker}{}%
\end{pgfscope}%
\end{pgfscope}%
\begin{pgfscope}%
\definecolor{textcolor}{rgb}{0.000000,0.000000,0.000000}%
\pgfsetstrokecolor{textcolor}%
\pgfsetfillcolor{textcolor}%
\pgftext[x=3.177102in,y=1.981474in,,top]{\color{textcolor}{\rmfamily\fontsize{5.000000}{6.000000}\selectfont\catcode`\^=\active\def^{\ifmmode\sp\else\^{}\fi}\catcode`\%=\active\def%{\%}5.000}}%
\end{pgfscope}%
\begin{pgfscope}%
\pgfpathrectangle{\pgfqpoint{2.197902in}{1.995363in}}{\pgfqpoint{0.979200in}{0.326250in}}%
\pgfusepath{clip}%
\pgfsetbuttcap%
\pgfsetroundjoin%
\pgfsetlinewidth{0.401500pt}%
\definecolor{currentstroke}{rgb}{0.690196,0.690196,0.690196}%
\pgfsetstrokecolor{currentstroke}%
\pgfsetstrokeopacity{0.250000}%
\pgfsetdash{{1.480000pt}{0.640000pt}}{0.000000pt}%
\pgfpathmoveto{\pgfqpoint{2.197902in}{1.995363in}}%
\pgfpathlineto{\pgfqpoint{3.177102in}{1.995363in}}%
\pgfusepath{stroke}%
\end{pgfscope}%
\begin{pgfscope}%
\pgfsetbuttcap%
\pgfsetroundjoin%
\definecolor{currentfill}{rgb}{0.000000,0.000000,0.000000}%
\pgfsetfillcolor{currentfill}%
\pgfsetlinewidth{0.803000pt}%
\definecolor{currentstroke}{rgb}{0.000000,0.000000,0.000000}%
\pgfsetstrokecolor{currentstroke}%
\pgfsetdash{}{0pt}%
\pgfsys@defobject{currentmarker}{\pgfqpoint{0.000000in}{0.000000in}}{\pgfqpoint{0.027778in}{0.000000in}}{%
\pgfpathmoveto{\pgfqpoint{0.000000in}{0.000000in}}%
\pgfpathlineto{\pgfqpoint{0.027778in}{0.000000in}}%
\pgfusepath{stroke,fill}%
}%
\begin{pgfscope}%
\pgfsys@transformshift{2.197902in}{1.995363in}%
\pgfsys@useobject{currentmarker}{}%
\end{pgfscope}%
\end{pgfscope}%
\begin{pgfscope}%
\definecolor{textcolor}{rgb}{0.000000,0.000000,0.000000}%
\pgfsetstrokecolor{textcolor}%
\pgfsetfillcolor{textcolor}%
\pgftext[x=1.981463in, y=1.971250in, left, base]{\color{textcolor}{\rmfamily\fontsize{5.000000}{6.000000}\selectfont\catcode`\^=\active\def^{\ifmmode\sp\else\^{}\fi}\catcode`\%=\active\def%{\%}-6.66}}%
\end{pgfscope}%
\begin{pgfscope}%
\pgfpathrectangle{\pgfqpoint{2.197902in}{1.995363in}}{\pgfqpoint{0.979200in}{0.326250in}}%
\pgfusepath{clip}%
\pgfsetbuttcap%
\pgfsetroundjoin%
\pgfsetlinewidth{0.401500pt}%
\definecolor{currentstroke}{rgb}{0.690196,0.690196,0.690196}%
\pgfsetstrokecolor{currentstroke}%
\pgfsetstrokeopacity{0.250000}%
\pgfsetdash{{1.480000pt}{0.640000pt}}{0.000000pt}%
\pgfpathmoveto{\pgfqpoint{2.197902in}{2.158488in}}%
\pgfpathlineto{\pgfqpoint{3.177102in}{2.158488in}}%
\pgfusepath{stroke}%
\end{pgfscope}%
\begin{pgfscope}%
\pgfsetbuttcap%
\pgfsetroundjoin%
\definecolor{currentfill}{rgb}{0.000000,0.000000,0.000000}%
\pgfsetfillcolor{currentfill}%
\pgfsetlinewidth{0.803000pt}%
\definecolor{currentstroke}{rgb}{0.000000,0.000000,0.000000}%
\pgfsetstrokecolor{currentstroke}%
\pgfsetdash{}{0pt}%
\pgfsys@defobject{currentmarker}{\pgfqpoint{0.000000in}{0.000000in}}{\pgfqpoint{0.027778in}{0.000000in}}{%
\pgfpathmoveto{\pgfqpoint{0.000000in}{0.000000in}}%
\pgfpathlineto{\pgfqpoint{0.027778in}{0.000000in}}%
\pgfusepath{stroke,fill}%
}%
\begin{pgfscope}%
\pgfsys@transformshift{2.197902in}{2.158488in}%
\pgfsys@useobject{currentmarker}{}%
\end{pgfscope}%
\end{pgfscope}%
\begin{pgfscope}%
\definecolor{textcolor}{rgb}{0.000000,0.000000,0.000000}%
\pgfsetstrokecolor{textcolor}%
\pgfsetfillcolor{textcolor}%
\pgftext[x=1.966996in, y=2.134375in, left, base]{\color{textcolor}{\rmfamily\fontsize{5.000000}{6.000000}\selectfont\catcode`\^=\active\def^{\ifmmode\sp\else\^{}\fi}\catcode`\%=\active\def%{\%}33.31}}%
\end{pgfscope}%
\begin{pgfscope}%
\pgfpathrectangle{\pgfqpoint{2.197902in}{1.995363in}}{\pgfqpoint{0.979200in}{0.326250in}}%
\pgfusepath{clip}%
\pgfsetbuttcap%
\pgfsetroundjoin%
\pgfsetlinewidth{0.401500pt}%
\definecolor{currentstroke}{rgb}{0.690196,0.690196,0.690196}%
\pgfsetstrokecolor{currentstroke}%
\pgfsetstrokeopacity{0.250000}%
\pgfsetdash{{1.480000pt}{0.640000pt}}{0.000000pt}%
\pgfpathmoveto{\pgfqpoint{2.197902in}{2.321613in}}%
\pgfpathlineto{\pgfqpoint{3.177102in}{2.321613in}}%
\pgfusepath{stroke}%
\end{pgfscope}%
\begin{pgfscope}%
\pgfsetbuttcap%
\pgfsetroundjoin%
\definecolor{currentfill}{rgb}{0.000000,0.000000,0.000000}%
\pgfsetfillcolor{currentfill}%
\pgfsetlinewidth{0.803000pt}%
\definecolor{currentstroke}{rgb}{0.000000,0.000000,0.000000}%
\pgfsetstrokecolor{currentstroke}%
\pgfsetdash{}{0pt}%
\pgfsys@defobject{currentmarker}{\pgfqpoint{0.000000in}{0.000000in}}{\pgfqpoint{0.027778in}{0.000000in}}{%
\pgfpathmoveto{\pgfqpoint{0.000000in}{0.000000in}}%
\pgfpathlineto{\pgfqpoint{0.027778in}{0.000000in}}%
\pgfusepath{stroke,fill}%
}%
\begin{pgfscope}%
\pgfsys@transformshift{2.197902in}{2.321613in}%
\pgfsys@useobject{currentmarker}{}%
\end{pgfscope}%
\end{pgfscope}%
\begin{pgfscope}%
\definecolor{textcolor}{rgb}{0.000000,0.000000,0.000000}%
\pgfsetstrokecolor{textcolor}%
\pgfsetfillcolor{textcolor}%
\pgftext[x=1.966996in, y=2.297500in, left, base]{\color{textcolor}{\rmfamily\fontsize{5.000000}{6.000000}\selectfont\catcode`\^=\active\def^{\ifmmode\sp\else\^{}\fi}\catcode`\%=\active\def%{\%}73.28}}%
\end{pgfscope}%
\begin{pgfscope}%
\pgfpathrectangle{\pgfqpoint{2.197902in}{1.995363in}}{\pgfqpoint{0.979200in}{0.326250in}}%
\pgfusepath{clip}%
\pgfsetrectcap%
\pgfsetroundjoin%
\pgfsetlinewidth{1.505625pt}%
\definecolor{currentstroke}{rgb}{0.835294,0.368627,0.000000}%
\pgfsetstrokecolor{currentstroke}%
\pgfsetdash{}{0pt}%
\pgfpathmoveto{\pgfqpoint{2.194661in}{2.089224in}}%
\pgfpathlineto{\pgfqpoint{2.198070in}{2.106287in}}%
\pgfpathlineto{\pgfqpoint{2.198295in}{2.109984in}}%
\pgfpathlineto{\pgfqpoint{2.200403in}{2.022551in}}%
\pgfpathlineto{\pgfqpoint{2.200638in}{2.022572in}}%
\pgfpathlineto{\pgfqpoint{2.207470in}{2.022551in}}%
\pgfpathlineto{\pgfqpoint{2.395976in}{2.022619in}}%
\pgfpathlineto{\pgfqpoint{2.397029in}{2.028059in}}%
\pgfpathlineto{\pgfqpoint{2.398004in}{2.045161in}}%
\pgfpathlineto{\pgfqpoint{2.399228in}{2.114549in}}%
\pgfpathlineto{\pgfqpoint{2.400619in}{2.257547in}}%
\pgfpathlineto{\pgfqpoint{2.402197in}{2.293955in}}%
\pgfpathlineto{\pgfqpoint{2.409841in}{2.293976in}}%
\pgfpathlineto{\pgfqpoint{2.453592in}{2.294114in}}%
\pgfpathlineto{\pgfqpoint{2.487864in}{2.294341in}}%
\pgfpathlineto{\pgfqpoint{2.522136in}{2.294249in}}%
\pgfpathlineto{\pgfqpoint{2.546616in}{2.294425in}}%
\pgfpathlineto{\pgfqpoint{2.561304in}{2.294166in}}%
\pgfpathlineto{\pgfqpoint{2.575992in}{2.294400in}}%
\pgfpathlineto{\pgfqpoint{2.590680in}{2.293870in}}%
\pgfpathlineto{\pgfqpoint{2.595576in}{2.294355in}}%
\pgfpathlineto{\pgfqpoint{2.600472in}{2.292152in}}%
\pgfpathlineto{\pgfqpoint{2.605368in}{2.294319in}}%
\pgfpathlineto{\pgfqpoint{2.610264in}{2.267842in}}%
\pgfpathlineto{\pgfqpoint{2.615160in}{2.294246in}}%
\pgfpathlineto{\pgfqpoint{2.620056in}{2.290715in}}%
\pgfpathlineto{\pgfqpoint{2.624952in}{2.294157in}}%
\pgfpathlineto{\pgfqpoint{2.629848in}{2.268712in}}%
\pgfpathlineto{\pgfqpoint{2.634744in}{2.293940in}}%
\pgfpathlineto{\pgfqpoint{2.639640in}{2.289328in}}%
\pgfpathlineto{\pgfqpoint{2.644317in}{2.293627in}}%
\pgfpathlineto{\pgfqpoint{2.648698in}{2.273441in}}%
\pgfpathlineto{\pgfqpoint{2.653392in}{2.243428in}}%
\pgfpathlineto{\pgfqpoint{2.667898in}{2.111526in}}%
\pgfpathlineto{\pgfqpoint{2.672794in}{2.084792in}}%
\pgfpathlineto{\pgfqpoint{2.677690in}{2.072920in}}%
\pgfpathlineto{\pgfqpoint{2.682586in}{2.075180in}}%
\pgfpathlineto{\pgfqpoint{2.685044in}{2.081737in}}%
\pgfpathlineto{\pgfqpoint{2.687683in}{2.095127in}}%
\pgfpathlineto{\pgfqpoint{2.688108in}{2.101605in}}%
\pgfpathlineto{\pgfqpoint{2.688178in}{2.100664in}}%
\pgfpathlineto{\pgfqpoint{2.689487in}{2.024609in}}%
\pgfpathlineto{\pgfqpoint{2.690201in}{2.022561in}}%
\pgfpathlineto{\pgfqpoint{2.787827in}{2.022569in}}%
\pgfpathlineto{\pgfqpoint{2.886070in}{2.022740in}}%
\pgfpathlineto{\pgfqpoint{2.886935in}{2.030970in}}%
\pgfpathlineto{\pgfqpoint{2.887823in}{2.047717in}}%
\pgfpathlineto{\pgfqpoint{2.888866in}{2.112227in}}%
\pgfpathlineto{\pgfqpoint{2.890336in}{2.264960in}}%
\pgfpathlineto{\pgfqpoint{2.891643in}{2.293893in}}%
\pgfpathlineto{\pgfqpoint{2.932320in}{2.294126in}}%
\pgfpathlineto{\pgfqpoint{3.118368in}{2.294023in}}%
\pgfpathlineto{\pgfqpoint{3.123264in}{2.291904in}}%
\pgfpathlineto{\pgfqpoint{3.128160in}{2.293865in}}%
\pgfpathlineto{\pgfqpoint{3.133056in}{2.267301in}}%
\pgfpathlineto{\pgfqpoint{3.137952in}{2.254397in}}%
\pgfpathlineto{\pgfqpoint{3.142848in}{2.222710in}}%
\pgfpathlineto{\pgfqpoint{3.152640in}{2.141159in}}%
\pgfpathlineto{\pgfqpoint{3.157536in}{2.108990in}}%
\pgfpathlineto{\pgfqpoint{3.162432in}{2.089373in}}%
\pgfpathlineto{\pgfqpoint{3.167328in}{2.083150in}}%
\pgfpathlineto{\pgfqpoint{3.172224in}{2.090171in}}%
\pgfpathlineto{\pgfqpoint{3.177102in}{2.110423in}}%
\pgfpathlineto{\pgfqpoint{3.177102in}{2.110423in}}%
\pgfusepath{stroke}%
\end{pgfscope}%
\begin{pgfscope}%
\pgfsetrectcap%
\pgfsetmiterjoin%
\pgfsetlinewidth{0.803000pt}%
\definecolor{currentstroke}{rgb}{0.000000,0.000000,0.000000}%
\pgfsetstrokecolor{currentstroke}%
\pgfsetdash{}{0pt}%
\pgfpathmoveto{\pgfqpoint{2.197902in}{1.995363in}}%
\pgfpathlineto{\pgfqpoint{2.197902in}{2.321613in}}%
\pgfusepath{stroke}%
\end{pgfscope}%
\begin{pgfscope}%
\pgfsetrectcap%
\pgfsetmiterjoin%
\pgfsetlinewidth{0.803000pt}%
\definecolor{currentstroke}{rgb}{0.000000,0.000000,0.000000}%
\pgfsetstrokecolor{currentstroke}%
\pgfsetdash{}{0pt}%
\pgfpathmoveto{\pgfqpoint{3.177102in}{1.995363in}}%
\pgfpathlineto{\pgfqpoint{3.177102in}{2.321613in}}%
\pgfusepath{stroke}%
\end{pgfscope}%
\begin{pgfscope}%
\pgfsetrectcap%
\pgfsetmiterjoin%
\pgfsetlinewidth{0.803000pt}%
\definecolor{currentstroke}{rgb}{0.000000,0.000000,0.000000}%
\pgfsetstrokecolor{currentstroke}%
\pgfsetdash{}{0pt}%
\pgfpathmoveto{\pgfqpoint{2.197902in}{1.995363in}}%
\pgfpathlineto{\pgfqpoint{3.177102in}{1.995363in}}%
\pgfusepath{stroke}%
\end{pgfscope}%
\begin{pgfscope}%
\pgfsetrectcap%
\pgfsetmiterjoin%
\pgfsetlinewidth{0.803000pt}%
\definecolor{currentstroke}{rgb}{0.000000,0.000000,0.000000}%
\pgfsetstrokecolor{currentstroke}%
\pgfsetdash{}{0pt}%
\pgfpathmoveto{\pgfqpoint{2.197902in}{2.321613in}}%
\pgfpathlineto{\pgfqpoint{3.177102in}{2.321613in}}%
\pgfusepath{stroke}%
\end{pgfscope}%
\begin{pgfscope}%
\pgfsetbuttcap%
\pgfsetmiterjoin%
\definecolor{currentfill}{rgb}{1.000000,1.000000,1.000000}%
\pgfsetfillcolor{currentfill}%
\pgfsetlinewidth{0.000000pt}%
\definecolor{currentstroke}{rgb}{0.000000,0.000000,0.000000}%
\pgfsetstrokecolor{currentstroke}%
\pgfsetstrokeopacity{0.000000}%
\pgfsetdash{}{0pt}%
\pgfpathmoveto{\pgfqpoint{2.197902in}{0.579321in}}%
\pgfpathlineto{\pgfqpoint{3.177102in}{0.579321in}}%
\pgfpathlineto{\pgfqpoint{3.177102in}{0.905571in}}%
\pgfpathlineto{\pgfqpoint{2.197902in}{0.905571in}}%
\pgfpathlineto{\pgfqpoint{2.197902in}{0.579321in}}%
\pgfpathclose%
\pgfusepath{fill}%
\end{pgfscope}%
\begin{pgfscope}%
\pgfpathrectangle{\pgfqpoint{2.197902in}{0.579321in}}{\pgfqpoint{0.979200in}{0.326250in}}%
\pgfusepath{clip}%
\pgfsetbuttcap%
\pgfsetroundjoin%
\pgfsetlinewidth{0.401500pt}%
\definecolor{currentstroke}{rgb}{0.690196,0.690196,0.690196}%
\pgfsetstrokecolor{currentstroke}%
\pgfsetstrokeopacity{0.250000}%
\pgfsetdash{{1.480000pt}{0.640000pt}}{0.000000pt}%
\pgfpathmoveto{\pgfqpoint{2.197902in}{0.579321in}}%
\pgfpathlineto{\pgfqpoint{2.197902in}{0.905571in}}%
\pgfusepath{stroke}%
\end{pgfscope}%
\begin{pgfscope}%
\pgfsetbuttcap%
\pgfsetroundjoin%
\definecolor{currentfill}{rgb}{0.000000,0.000000,0.000000}%
\pgfsetfillcolor{currentfill}%
\pgfsetlinewidth{0.803000pt}%
\definecolor{currentstroke}{rgb}{0.000000,0.000000,0.000000}%
\pgfsetstrokecolor{currentstroke}%
\pgfsetdash{}{0pt}%
\pgfsys@defobject{currentmarker}{\pgfqpoint{0.000000in}{0.000000in}}{\pgfqpoint{0.000000in}{0.027778in}}{%
\pgfpathmoveto{\pgfqpoint{0.000000in}{0.000000in}}%
\pgfpathlineto{\pgfqpoint{0.000000in}{0.027778in}}%
\pgfusepath{stroke,fill}%
}%
\begin{pgfscope}%
\pgfsys@transformshift{2.197902in}{0.579321in}%
\pgfsys@useobject{currentmarker}{}%
\end{pgfscope}%
\end{pgfscope}%
\begin{pgfscope}%
\definecolor{textcolor}{rgb}{0.000000,0.000000,0.000000}%
\pgfsetstrokecolor{textcolor}%
\pgfsetfillcolor{textcolor}%
\pgftext[x=2.197902in,y=0.565432in,,top]{\color{textcolor}{\rmfamily\fontsize{5.000000}{6.000000}\selectfont\catcode`\^=\active\def^{\ifmmode\sp\else\^{}\fi}\catcode`\%=\active\def%{\%}4.980}}%
\end{pgfscope}%
\begin{pgfscope}%
\pgfpathrectangle{\pgfqpoint{2.197902in}{0.579321in}}{\pgfqpoint{0.979200in}{0.326250in}}%
\pgfusepath{clip}%
\pgfsetbuttcap%
\pgfsetroundjoin%
\pgfsetlinewidth{0.401500pt}%
\definecolor{currentstroke}{rgb}{0.690196,0.690196,0.690196}%
\pgfsetstrokecolor{currentstroke}%
\pgfsetstrokeopacity{0.250000}%
\pgfsetdash{{1.480000pt}{0.640000pt}}{0.000000pt}%
\pgfpathmoveto{\pgfqpoint{2.442702in}{0.579321in}}%
\pgfpathlineto{\pgfqpoint{2.442702in}{0.905571in}}%
\pgfusepath{stroke}%
\end{pgfscope}%
\begin{pgfscope}%
\pgfsetbuttcap%
\pgfsetroundjoin%
\definecolor{currentfill}{rgb}{0.000000,0.000000,0.000000}%
\pgfsetfillcolor{currentfill}%
\pgfsetlinewidth{0.803000pt}%
\definecolor{currentstroke}{rgb}{0.000000,0.000000,0.000000}%
\pgfsetstrokecolor{currentstroke}%
\pgfsetdash{}{0pt}%
\pgfsys@defobject{currentmarker}{\pgfqpoint{0.000000in}{0.000000in}}{\pgfqpoint{0.000000in}{0.027778in}}{%
\pgfpathmoveto{\pgfqpoint{0.000000in}{0.000000in}}%
\pgfpathlineto{\pgfqpoint{0.000000in}{0.027778in}}%
\pgfusepath{stroke,fill}%
}%
\begin{pgfscope}%
\pgfsys@transformshift{2.442702in}{0.579321in}%
\pgfsys@useobject{currentmarker}{}%
\end{pgfscope}%
\end{pgfscope}%
\begin{pgfscope}%
\definecolor{textcolor}{rgb}{0.000000,0.000000,0.000000}%
\pgfsetstrokecolor{textcolor}%
\pgfsetfillcolor{textcolor}%
\pgftext[x=2.442702in,y=0.565432in,,top]{\color{textcolor}{\rmfamily\fontsize{5.000000}{6.000000}\selectfont\catcode`\^=\active\def^{\ifmmode\sp\else\^{}\fi}\catcode`\%=\active\def%{\%}4.985}}%
\end{pgfscope}%
\begin{pgfscope}%
\pgfpathrectangle{\pgfqpoint{2.197902in}{0.579321in}}{\pgfqpoint{0.979200in}{0.326250in}}%
\pgfusepath{clip}%
\pgfsetbuttcap%
\pgfsetroundjoin%
\pgfsetlinewidth{0.401500pt}%
\definecolor{currentstroke}{rgb}{0.690196,0.690196,0.690196}%
\pgfsetstrokecolor{currentstroke}%
\pgfsetstrokeopacity{0.250000}%
\pgfsetdash{{1.480000pt}{0.640000pt}}{0.000000pt}%
\pgfpathmoveto{\pgfqpoint{2.687502in}{0.579321in}}%
\pgfpathlineto{\pgfqpoint{2.687502in}{0.905571in}}%
\pgfusepath{stroke}%
\end{pgfscope}%
\begin{pgfscope}%
\pgfsetbuttcap%
\pgfsetroundjoin%
\definecolor{currentfill}{rgb}{0.000000,0.000000,0.000000}%
\pgfsetfillcolor{currentfill}%
\pgfsetlinewidth{0.803000pt}%
\definecolor{currentstroke}{rgb}{0.000000,0.000000,0.000000}%
\pgfsetstrokecolor{currentstroke}%
\pgfsetdash{}{0pt}%
\pgfsys@defobject{currentmarker}{\pgfqpoint{0.000000in}{0.000000in}}{\pgfqpoint{0.000000in}{0.027778in}}{%
\pgfpathmoveto{\pgfqpoint{0.000000in}{0.000000in}}%
\pgfpathlineto{\pgfqpoint{0.000000in}{0.027778in}}%
\pgfusepath{stroke,fill}%
}%
\begin{pgfscope}%
\pgfsys@transformshift{2.687502in}{0.579321in}%
\pgfsys@useobject{currentmarker}{}%
\end{pgfscope}%
\end{pgfscope}%
\begin{pgfscope}%
\definecolor{textcolor}{rgb}{0.000000,0.000000,0.000000}%
\pgfsetstrokecolor{textcolor}%
\pgfsetfillcolor{textcolor}%
\pgftext[x=2.687502in,y=0.565432in,,top]{\color{textcolor}{\rmfamily\fontsize{5.000000}{6.000000}\selectfont\catcode`\^=\active\def^{\ifmmode\sp\else\^{}\fi}\catcode`\%=\active\def%{\%}4.990}}%
\end{pgfscope}%
\begin{pgfscope}%
\pgfpathrectangle{\pgfqpoint{2.197902in}{0.579321in}}{\pgfqpoint{0.979200in}{0.326250in}}%
\pgfusepath{clip}%
\pgfsetbuttcap%
\pgfsetroundjoin%
\pgfsetlinewidth{0.401500pt}%
\definecolor{currentstroke}{rgb}{0.690196,0.690196,0.690196}%
\pgfsetstrokecolor{currentstroke}%
\pgfsetstrokeopacity{0.250000}%
\pgfsetdash{{1.480000pt}{0.640000pt}}{0.000000pt}%
\pgfpathmoveto{\pgfqpoint{2.932302in}{0.579321in}}%
\pgfpathlineto{\pgfqpoint{2.932302in}{0.905571in}}%
\pgfusepath{stroke}%
\end{pgfscope}%
\begin{pgfscope}%
\pgfsetbuttcap%
\pgfsetroundjoin%
\definecolor{currentfill}{rgb}{0.000000,0.000000,0.000000}%
\pgfsetfillcolor{currentfill}%
\pgfsetlinewidth{0.803000pt}%
\definecolor{currentstroke}{rgb}{0.000000,0.000000,0.000000}%
\pgfsetstrokecolor{currentstroke}%
\pgfsetdash{}{0pt}%
\pgfsys@defobject{currentmarker}{\pgfqpoint{0.000000in}{0.000000in}}{\pgfqpoint{0.000000in}{0.027778in}}{%
\pgfpathmoveto{\pgfqpoint{0.000000in}{0.000000in}}%
\pgfpathlineto{\pgfqpoint{0.000000in}{0.027778in}}%
\pgfusepath{stroke,fill}%
}%
\begin{pgfscope}%
\pgfsys@transformshift{2.932302in}{0.579321in}%
\pgfsys@useobject{currentmarker}{}%
\end{pgfscope}%
\end{pgfscope}%
\begin{pgfscope}%
\definecolor{textcolor}{rgb}{0.000000,0.000000,0.000000}%
\pgfsetstrokecolor{textcolor}%
\pgfsetfillcolor{textcolor}%
\pgftext[x=2.932302in,y=0.565432in,,top]{\color{textcolor}{\rmfamily\fontsize{5.000000}{6.000000}\selectfont\catcode`\^=\active\def^{\ifmmode\sp\else\^{}\fi}\catcode`\%=\active\def%{\%}4.995}}%
\end{pgfscope}%
\begin{pgfscope}%
\pgfpathrectangle{\pgfqpoint{2.197902in}{0.579321in}}{\pgfqpoint{0.979200in}{0.326250in}}%
\pgfusepath{clip}%
\pgfsetbuttcap%
\pgfsetroundjoin%
\pgfsetlinewidth{0.401500pt}%
\definecolor{currentstroke}{rgb}{0.690196,0.690196,0.690196}%
\pgfsetstrokecolor{currentstroke}%
\pgfsetstrokeopacity{0.250000}%
\pgfsetdash{{1.480000pt}{0.640000pt}}{0.000000pt}%
\pgfpathmoveto{\pgfqpoint{3.177102in}{0.579321in}}%
\pgfpathlineto{\pgfqpoint{3.177102in}{0.905571in}}%
\pgfusepath{stroke}%
\end{pgfscope}%
\begin{pgfscope}%
\pgfsetbuttcap%
\pgfsetroundjoin%
\definecolor{currentfill}{rgb}{0.000000,0.000000,0.000000}%
\pgfsetfillcolor{currentfill}%
\pgfsetlinewidth{0.803000pt}%
\definecolor{currentstroke}{rgb}{0.000000,0.000000,0.000000}%
\pgfsetstrokecolor{currentstroke}%
\pgfsetdash{}{0pt}%
\pgfsys@defobject{currentmarker}{\pgfqpoint{0.000000in}{0.000000in}}{\pgfqpoint{0.000000in}{0.027778in}}{%
\pgfpathmoveto{\pgfqpoint{0.000000in}{0.000000in}}%
\pgfpathlineto{\pgfqpoint{0.000000in}{0.027778in}}%
\pgfusepath{stroke,fill}%
}%
\begin{pgfscope}%
\pgfsys@transformshift{3.177102in}{0.579321in}%
\pgfsys@useobject{currentmarker}{}%
\end{pgfscope}%
\end{pgfscope}%
\begin{pgfscope}%
\definecolor{textcolor}{rgb}{0.000000,0.000000,0.000000}%
\pgfsetstrokecolor{textcolor}%
\pgfsetfillcolor{textcolor}%
\pgftext[x=3.177102in,y=0.565432in,,top]{\color{textcolor}{\rmfamily\fontsize{5.000000}{6.000000}\selectfont\catcode`\^=\active\def^{\ifmmode\sp\else\^{}\fi}\catcode`\%=\active\def%{\%}5.000}}%
\end{pgfscope}%
\begin{pgfscope}%
\pgfpathrectangle{\pgfqpoint{2.197902in}{0.579321in}}{\pgfqpoint{0.979200in}{0.326250in}}%
\pgfusepath{clip}%
\pgfsetbuttcap%
\pgfsetroundjoin%
\pgfsetlinewidth{0.401500pt}%
\definecolor{currentstroke}{rgb}{0.690196,0.690196,0.690196}%
\pgfsetstrokecolor{currentstroke}%
\pgfsetstrokeopacity{0.250000}%
\pgfsetdash{{1.480000pt}{0.640000pt}}{0.000000pt}%
\pgfpathmoveto{\pgfqpoint{2.197902in}{0.579321in}}%
\pgfpathlineto{\pgfqpoint{3.177102in}{0.579321in}}%
\pgfusepath{stroke}%
\end{pgfscope}%
\begin{pgfscope}%
\pgfsetbuttcap%
\pgfsetroundjoin%
\definecolor{currentfill}{rgb}{0.000000,0.000000,0.000000}%
\pgfsetfillcolor{currentfill}%
\pgfsetlinewidth{0.803000pt}%
\definecolor{currentstroke}{rgb}{0.000000,0.000000,0.000000}%
\pgfsetstrokecolor{currentstroke}%
\pgfsetdash{}{0pt}%
\pgfsys@defobject{currentmarker}{\pgfqpoint{0.000000in}{0.000000in}}{\pgfqpoint{0.027778in}{0.000000in}}{%
\pgfpathmoveto{\pgfqpoint{0.000000in}{0.000000in}}%
\pgfpathlineto{\pgfqpoint{0.027778in}{0.000000in}}%
\pgfusepath{stroke,fill}%
}%
\begin{pgfscope}%
\pgfsys@transformshift{2.197902in}{0.579321in}%
\pgfsys@useobject{currentmarker}{}%
\end{pgfscope}%
\end{pgfscope}%
\begin{pgfscope}%
\definecolor{textcolor}{rgb}{0.000000,0.000000,0.000000}%
\pgfsetstrokecolor{textcolor}%
\pgfsetfillcolor{textcolor}%
\pgftext[x=1.966996in, y=0.555209in, left, base]{\color{textcolor}{\rmfamily\fontsize{5.000000}{6.000000}\selectfont\catcode`\^=\active\def^{\ifmmode\sp\else\^{}\fi}\catcode`\%=\active\def%{\%}14.85}}%
\end{pgfscope}%
\begin{pgfscope}%
\pgfpathrectangle{\pgfqpoint{2.197902in}{0.579321in}}{\pgfqpoint{0.979200in}{0.326250in}}%
\pgfusepath{clip}%
\pgfsetbuttcap%
\pgfsetroundjoin%
\pgfsetlinewidth{0.401500pt}%
\definecolor{currentstroke}{rgb}{0.690196,0.690196,0.690196}%
\pgfsetstrokecolor{currentstroke}%
\pgfsetstrokeopacity{0.250000}%
\pgfsetdash{{1.480000pt}{0.640000pt}}{0.000000pt}%
\pgfpathmoveto{\pgfqpoint{2.197902in}{0.742446in}}%
\pgfpathlineto{\pgfqpoint{3.177102in}{0.742446in}}%
\pgfusepath{stroke}%
\end{pgfscope}%
\begin{pgfscope}%
\pgfsetbuttcap%
\pgfsetroundjoin%
\definecolor{currentfill}{rgb}{0.000000,0.000000,0.000000}%
\pgfsetfillcolor{currentfill}%
\pgfsetlinewidth{0.803000pt}%
\definecolor{currentstroke}{rgb}{0.000000,0.000000,0.000000}%
\pgfsetstrokecolor{currentstroke}%
\pgfsetdash{}{0pt}%
\pgfsys@defobject{currentmarker}{\pgfqpoint{0.000000in}{0.000000in}}{\pgfqpoint{0.027778in}{0.000000in}}{%
\pgfpathmoveto{\pgfqpoint{0.000000in}{0.000000in}}%
\pgfpathlineto{\pgfqpoint{0.027778in}{0.000000in}}%
\pgfusepath{stroke,fill}%
}%
\begin{pgfscope}%
\pgfsys@transformshift{2.197902in}{0.742446in}%
\pgfsys@useobject{currentmarker}{}%
\end{pgfscope}%
\end{pgfscope}%
\begin{pgfscope}%
\definecolor{textcolor}{rgb}{0.000000,0.000000,0.000000}%
\pgfsetstrokecolor{textcolor}%
\pgfsetfillcolor{textcolor}%
\pgftext[x=1.966996in, y=0.718334in, left, base]{\color{textcolor}{\rmfamily\fontsize{5.000000}{6.000000}\selectfont\catcode`\^=\active\def^{\ifmmode\sp\else\^{}\fi}\catcode`\%=\active\def%{\%}14.90}}%
\end{pgfscope}%
\begin{pgfscope}%
\pgfpathrectangle{\pgfqpoint{2.197902in}{0.579321in}}{\pgfqpoint{0.979200in}{0.326250in}}%
\pgfusepath{clip}%
\pgfsetbuttcap%
\pgfsetroundjoin%
\pgfsetlinewidth{0.401500pt}%
\definecolor{currentstroke}{rgb}{0.690196,0.690196,0.690196}%
\pgfsetstrokecolor{currentstroke}%
\pgfsetstrokeopacity{0.250000}%
\pgfsetdash{{1.480000pt}{0.640000pt}}{0.000000pt}%
\pgfpathmoveto{\pgfqpoint{2.197902in}{0.905571in}}%
\pgfpathlineto{\pgfqpoint{3.177102in}{0.905571in}}%
\pgfusepath{stroke}%
\end{pgfscope}%
\begin{pgfscope}%
\pgfsetbuttcap%
\pgfsetroundjoin%
\definecolor{currentfill}{rgb}{0.000000,0.000000,0.000000}%
\pgfsetfillcolor{currentfill}%
\pgfsetlinewidth{0.803000pt}%
\definecolor{currentstroke}{rgb}{0.000000,0.000000,0.000000}%
\pgfsetstrokecolor{currentstroke}%
\pgfsetdash{}{0pt}%
\pgfsys@defobject{currentmarker}{\pgfqpoint{0.000000in}{0.000000in}}{\pgfqpoint{0.027778in}{0.000000in}}{%
\pgfpathmoveto{\pgfqpoint{0.000000in}{0.000000in}}%
\pgfpathlineto{\pgfqpoint{0.027778in}{0.000000in}}%
\pgfusepath{stroke,fill}%
}%
\begin{pgfscope}%
\pgfsys@transformshift{2.197902in}{0.905571in}%
\pgfsys@useobject{currentmarker}{}%
\end{pgfscope}%
\end{pgfscope}%
\begin{pgfscope}%
\definecolor{textcolor}{rgb}{0.000000,0.000000,0.000000}%
\pgfsetstrokecolor{textcolor}%
\pgfsetfillcolor{textcolor}%
\pgftext[x=1.966996in, y=0.881459in, left, base]{\color{textcolor}{\rmfamily\fontsize{5.000000}{6.000000}\selectfont\catcode`\^=\active\def^{\ifmmode\sp\else\^{}\fi}\catcode`\%=\active\def%{\%}14.95}}%
\end{pgfscope}%
\begin{pgfscope}%
\pgfpathrectangle{\pgfqpoint{2.197902in}{0.579321in}}{\pgfqpoint{0.979200in}{0.326250in}}%
\pgfusepath{clip}%
\pgfsetrectcap%
\pgfsetroundjoin%
\pgfsetlinewidth{1.505625pt}%
\definecolor{currentstroke}{rgb}{0.000000,0.619608,0.450980}%
\pgfsetstrokecolor{currentstroke}%
\pgfsetdash{}{0pt}%
\pgfpathmoveto{\pgfqpoint{2.194661in}{0.810669in}}%
\pgfpathlineto{\pgfqpoint{2.400619in}{0.608576in}}%
\pgfpathlineto{\pgfqpoint{2.402197in}{0.610705in}}%
\pgfpathlineto{\pgfqpoint{2.419320in}{0.659113in}}%
\pgfpathlineto{\pgfqpoint{2.434008in}{0.696839in}}%
\pgfpathlineto{\pgfqpoint{2.448696in}{0.731004in}}%
\pgfpathlineto{\pgfqpoint{2.463384in}{0.761614in}}%
\pgfpathlineto{\pgfqpoint{2.478072in}{0.788663in}}%
\pgfpathlineto{\pgfqpoint{2.492760in}{0.812155in}}%
\pgfpathlineto{\pgfqpoint{2.507448in}{0.832085in}}%
\pgfpathlineto{\pgfqpoint{2.522136in}{0.848459in}}%
\pgfpathlineto{\pgfqpoint{2.531928in}{0.857398in}}%
\pgfpathlineto{\pgfqpoint{2.541720in}{0.864755in}}%
\pgfpathlineto{\pgfqpoint{2.551512in}{0.870532in}}%
\pgfpathlineto{\pgfqpoint{2.561304in}{0.874728in}}%
\pgfpathlineto{\pgfqpoint{2.571096in}{0.877345in}}%
\pgfpathlineto{\pgfqpoint{2.580888in}{0.878384in}}%
\pgfpathlineto{\pgfqpoint{2.590680in}{0.877845in}}%
\pgfpathlineto{\pgfqpoint{2.600472in}{0.875751in}}%
\pgfpathlineto{\pgfqpoint{2.610264in}{0.872405in}}%
\pgfpathlineto{\pgfqpoint{2.624952in}{0.863821in}}%
\pgfpathlineto{\pgfqpoint{2.629848in}{0.860651in}}%
\pgfpathlineto{\pgfqpoint{2.648698in}{0.843688in}}%
\pgfpathlineto{\pgfqpoint{2.890336in}{0.606509in}}%
\pgfpathlineto{\pgfqpoint{2.891643in}{0.607898in}}%
\pgfpathlineto{\pgfqpoint{2.907840in}{0.653869in}}%
\pgfpathlineto{\pgfqpoint{2.922528in}{0.691879in}}%
\pgfpathlineto{\pgfqpoint{2.937216in}{0.726332in}}%
\pgfpathlineto{\pgfqpoint{2.951904in}{0.757225in}}%
\pgfpathlineto{\pgfqpoint{2.966592in}{0.784561in}}%
\pgfpathlineto{\pgfqpoint{2.981280in}{0.808337in}}%
\pgfpathlineto{\pgfqpoint{2.995968in}{0.828554in}}%
\pgfpathlineto{\pgfqpoint{3.010656in}{0.845212in}}%
\pgfpathlineto{\pgfqpoint{3.020448in}{0.854340in}}%
\pgfpathlineto{\pgfqpoint{3.030240in}{0.861887in}}%
\pgfpathlineto{\pgfqpoint{3.040032in}{0.867853in}}%
\pgfpathlineto{\pgfqpoint{3.049824in}{0.872238in}}%
\pgfpathlineto{\pgfqpoint{3.059616in}{0.875043in}}%
\pgfpathlineto{\pgfqpoint{3.069408in}{0.876269in}}%
\pgfpathlineto{\pgfqpoint{3.079200in}{0.875916in}}%
\pgfpathlineto{\pgfqpoint{3.088992in}{0.873985in}}%
\pgfpathlineto{\pgfqpoint{3.098784in}{0.870477in}}%
\pgfpathlineto{\pgfqpoint{3.108576in}{0.865393in}}%
\pgfpathlineto{\pgfqpoint{3.118368in}{0.858735in}}%
\pgfpathlineto{\pgfqpoint{3.133056in}{0.846179in}}%
\pgfpathlineto{\pgfqpoint{3.177102in}{0.802857in}}%
\pgfpathlineto{\pgfqpoint{3.177102in}{0.802857in}}%
\pgfusepath{stroke}%
\end{pgfscope}%
\begin{pgfscope}%
\pgfsetrectcap%
\pgfsetmiterjoin%
\pgfsetlinewidth{0.803000pt}%
\definecolor{currentstroke}{rgb}{0.000000,0.000000,0.000000}%
\pgfsetstrokecolor{currentstroke}%
\pgfsetdash{}{0pt}%
\pgfpathmoveto{\pgfqpoint{2.197902in}{0.579321in}}%
\pgfpathlineto{\pgfqpoint{2.197902in}{0.905571in}}%
\pgfusepath{stroke}%
\end{pgfscope}%
\begin{pgfscope}%
\pgfsetrectcap%
\pgfsetmiterjoin%
\pgfsetlinewidth{0.803000pt}%
\definecolor{currentstroke}{rgb}{0.000000,0.000000,0.000000}%
\pgfsetstrokecolor{currentstroke}%
\pgfsetdash{}{0pt}%
\pgfpathmoveto{\pgfqpoint{3.177102in}{0.579321in}}%
\pgfpathlineto{\pgfqpoint{3.177102in}{0.905571in}}%
\pgfusepath{stroke}%
\end{pgfscope}%
\begin{pgfscope}%
\pgfsetrectcap%
\pgfsetmiterjoin%
\pgfsetlinewidth{0.803000pt}%
\definecolor{currentstroke}{rgb}{0.000000,0.000000,0.000000}%
\pgfsetstrokecolor{currentstroke}%
\pgfsetdash{}{0pt}%
\pgfpathmoveto{\pgfqpoint{2.197902in}{0.579321in}}%
\pgfpathlineto{\pgfqpoint{3.177102in}{0.579321in}}%
\pgfusepath{stroke}%
\end{pgfscope}%
\begin{pgfscope}%
\pgfsetrectcap%
\pgfsetmiterjoin%
\pgfsetlinewidth{0.803000pt}%
\definecolor{currentstroke}{rgb}{0.000000,0.000000,0.000000}%
\pgfsetstrokecolor{currentstroke}%
\pgfsetdash{}{0pt}%
\pgfpathmoveto{\pgfqpoint{2.197902in}{0.905571in}}%
\pgfpathlineto{\pgfqpoint{3.177102in}{0.905571in}}%
\pgfusepath{stroke}%
\end{pgfscope}%
\end{pgfpicture}%
\makeatother%
\endgroup%

	\caption{(a) corrente no dreno do transistor, (b) tensão dreno--source do transistor e (c) tensão de saída do circuito Flyback para $K~L1~L2~1$ .}
	\label{fig:ex4_1}
\end{figure}

\medskip
\textbf{b)} Altere o acoplamento para $K~L1~L2~0.99$ e apresente as mesmas formas de onda.

\medskip
\textbf{Resolução}:
A Figura~\ref{fig:ex4_2} apresenta as formas de onda solicitadas no enunciado. Pode-se observar o surgimento de oscilação na tensão $V_{DS}$.
\begin{figure}[htb!]
	\centering
	%% Creator: Matplotlib, PGF backend
%%
%% To include the figure in your LaTeX document, write
%%   \input{<filename>.pgf}
%%
%% Make sure the required packages are loaded in your preamble
%%   \usepackage{pgf}
%%
%% Also ensure that all the required font packages are loaded; for instance,
%% the lmodern package is sometimes necessary when using math font.
%%   \usepackage{lmodern}
%%
%% Figures using additional raster images can only be included by \input if
%% they are in the same directory as the main LaTeX file. For loading figures
%% from other directories you can use the `import` package
%%   \usepackage{import}
%%
%% and then include the figures with
%%   \import{<path to file>}{<filename>.pgf}
%%
%% Matplotlib used the following preamble
%%   \def\mathdefault#1{#1}
%%   \everymath=\expandafter{\the\everymath\displaystyle}
%%   \IfFileExists{scrextend.sty}{
%%     \usepackage[fontsize=8.000000pt]{scrextend}
%%   }{
%%     \renewcommand{\normalsize}{\fontsize{8.000000}{9.600000}\selectfont}
%%     \normalsize
%%   }
%%   
%%           \usepackage{amsmath}
%%           \usepackage{amssymb}
%%           \usepackage{mathtools}
%%           \usepackage{dsfont}
%%           \providecommand{\mathdefault}[1]{#1}
%%       
%%   \makeatletter\@ifpackageloaded{underscore}{}{\usepackage[strings]{underscore}}\makeatother
%%
\begingroup%
\makeatletter%
\begin{pgfpicture}%
\pgfpathrectangle{\pgfpointorigin}{\pgfqpoint{3.393768in}{3.646235in}}%
\pgfusepath{use as bounding box, clip}%
\begin{pgfscope}%
\pgfsetbuttcap%
\pgfsetmiterjoin%
\definecolor{currentfill}{rgb}{1.000000,1.000000,1.000000}%
\pgfsetfillcolor{currentfill}%
\pgfsetlinewidth{0.000000pt}%
\definecolor{currentstroke}{rgb}{1.000000,1.000000,1.000000}%
\pgfsetstrokecolor{currentstroke}%
\pgfsetdash{}{0pt}%
\pgfpathmoveto{\pgfqpoint{0.000000in}{0.000000in}}%
\pgfpathlineto{\pgfqpoint{3.393768in}{0.000000in}}%
\pgfpathlineto{\pgfqpoint{3.393768in}{3.646235in}}%
\pgfpathlineto{\pgfqpoint{0.000000in}{3.646235in}}%
\pgfpathlineto{\pgfqpoint{0.000000in}{0.000000in}}%
\pgfpathclose%
\pgfusepath{fill}%
\end{pgfscope}%
\begin{pgfscope}%
\pgfsetbuttcap%
\pgfsetmiterjoin%
\definecolor{currentfill}{rgb}{1.000000,1.000000,1.000000}%
\pgfsetfillcolor{currentfill}%
\pgfsetlinewidth{0.000000pt}%
\definecolor{currentstroke}{rgb}{0.000000,0.000000,0.000000}%
\pgfsetstrokecolor{currentstroke}%
\pgfsetstrokeopacity{0.000000}%
\pgfsetdash{}{0pt}%
\pgfpathmoveto{\pgfqpoint{0.573768in}{2.601404in}}%
\pgfpathlineto{\pgfqpoint{3.293768in}{2.601404in}}%
\pgfpathlineto{\pgfqpoint{3.293768in}{3.507654in}}%
\pgfpathlineto{\pgfqpoint{0.573768in}{3.507654in}}%
\pgfpathlineto{\pgfqpoint{0.573768in}{2.601404in}}%
\pgfpathclose%
\pgfusepath{fill}%
\end{pgfscope}%
\begin{pgfscope}%
\pgfpathrectangle{\pgfqpoint{0.573768in}{2.601404in}}{\pgfqpoint{2.720000in}{0.906250in}}%
\pgfusepath{clip}%
\pgfsetbuttcap%
\pgfsetroundjoin%
\pgfsetlinewidth{0.501875pt}%
\definecolor{currentstroke}{rgb}{0.690196,0.690196,0.690196}%
\pgfsetstrokecolor{currentstroke}%
\pgfsetstrokeopacity{0.250000}%
\pgfsetdash{{1.850000pt}{0.800000pt}}{0.000000pt}%
\pgfpathmoveto{\pgfqpoint{0.697405in}{2.601404in}}%
\pgfpathlineto{\pgfqpoint{0.697405in}{3.507654in}}%
\pgfusepath{stroke}%
\end{pgfscope}%
\begin{pgfscope}%
\pgfsetbuttcap%
\pgfsetroundjoin%
\definecolor{currentfill}{rgb}{0.000000,0.000000,0.000000}%
\pgfsetfillcolor{currentfill}%
\pgfsetlinewidth{0.803000pt}%
\definecolor{currentstroke}{rgb}{0.000000,0.000000,0.000000}%
\pgfsetstrokecolor{currentstroke}%
\pgfsetdash{}{0pt}%
\pgfsys@defobject{currentmarker}{\pgfqpoint{0.000000in}{-0.048611in}}{\pgfqpoint{0.000000in}{0.000000in}}{%
\pgfpathmoveto{\pgfqpoint{0.000000in}{0.000000in}}%
\pgfpathlineto{\pgfqpoint{0.000000in}{-0.048611in}}%
\pgfusepath{stroke,fill}%
}%
\begin{pgfscope}%
\pgfsys@transformshift{0.697405in}{2.601404in}%
\pgfsys@useobject{currentmarker}{}%
\end{pgfscope}%
\end{pgfscope}%
\begin{pgfscope}%
\pgfpathrectangle{\pgfqpoint{0.573768in}{2.601404in}}{\pgfqpoint{2.720000in}{0.906250in}}%
\pgfusepath{clip}%
\pgfsetbuttcap%
\pgfsetroundjoin%
\pgfsetlinewidth{0.501875pt}%
\definecolor{currentstroke}{rgb}{0.690196,0.690196,0.690196}%
\pgfsetstrokecolor{currentstroke}%
\pgfsetstrokeopacity{0.250000}%
\pgfsetdash{{1.850000pt}{0.800000pt}}{0.000000pt}%
\pgfpathmoveto{\pgfqpoint{1.191950in}{2.601404in}}%
\pgfpathlineto{\pgfqpoint{1.191950in}{3.507654in}}%
\pgfusepath{stroke}%
\end{pgfscope}%
\begin{pgfscope}%
\pgfsetbuttcap%
\pgfsetroundjoin%
\definecolor{currentfill}{rgb}{0.000000,0.000000,0.000000}%
\pgfsetfillcolor{currentfill}%
\pgfsetlinewidth{0.803000pt}%
\definecolor{currentstroke}{rgb}{0.000000,0.000000,0.000000}%
\pgfsetstrokecolor{currentstroke}%
\pgfsetdash{}{0pt}%
\pgfsys@defobject{currentmarker}{\pgfqpoint{0.000000in}{-0.048611in}}{\pgfqpoint{0.000000in}{0.000000in}}{%
\pgfpathmoveto{\pgfqpoint{0.000000in}{0.000000in}}%
\pgfpathlineto{\pgfqpoint{0.000000in}{-0.048611in}}%
\pgfusepath{stroke,fill}%
}%
\begin{pgfscope}%
\pgfsys@transformshift{1.191950in}{2.601404in}%
\pgfsys@useobject{currentmarker}{}%
\end{pgfscope}%
\end{pgfscope}%
\begin{pgfscope}%
\pgfpathrectangle{\pgfqpoint{0.573768in}{2.601404in}}{\pgfqpoint{2.720000in}{0.906250in}}%
\pgfusepath{clip}%
\pgfsetbuttcap%
\pgfsetroundjoin%
\pgfsetlinewidth{0.501875pt}%
\definecolor{currentstroke}{rgb}{0.690196,0.690196,0.690196}%
\pgfsetstrokecolor{currentstroke}%
\pgfsetstrokeopacity{0.250000}%
\pgfsetdash{{1.850000pt}{0.800000pt}}{0.000000pt}%
\pgfpathmoveto{\pgfqpoint{1.686496in}{2.601404in}}%
\pgfpathlineto{\pgfqpoint{1.686496in}{3.507654in}}%
\pgfusepath{stroke}%
\end{pgfscope}%
\begin{pgfscope}%
\pgfsetbuttcap%
\pgfsetroundjoin%
\definecolor{currentfill}{rgb}{0.000000,0.000000,0.000000}%
\pgfsetfillcolor{currentfill}%
\pgfsetlinewidth{0.803000pt}%
\definecolor{currentstroke}{rgb}{0.000000,0.000000,0.000000}%
\pgfsetstrokecolor{currentstroke}%
\pgfsetdash{}{0pt}%
\pgfsys@defobject{currentmarker}{\pgfqpoint{0.000000in}{-0.048611in}}{\pgfqpoint{0.000000in}{0.000000in}}{%
\pgfpathmoveto{\pgfqpoint{0.000000in}{0.000000in}}%
\pgfpathlineto{\pgfqpoint{0.000000in}{-0.048611in}}%
\pgfusepath{stroke,fill}%
}%
\begin{pgfscope}%
\pgfsys@transformshift{1.686496in}{2.601404in}%
\pgfsys@useobject{currentmarker}{}%
\end{pgfscope}%
\end{pgfscope}%
\begin{pgfscope}%
\pgfpathrectangle{\pgfqpoint{0.573768in}{2.601404in}}{\pgfqpoint{2.720000in}{0.906250in}}%
\pgfusepath{clip}%
\pgfsetbuttcap%
\pgfsetroundjoin%
\pgfsetlinewidth{0.501875pt}%
\definecolor{currentstroke}{rgb}{0.690196,0.690196,0.690196}%
\pgfsetstrokecolor{currentstroke}%
\pgfsetstrokeopacity{0.250000}%
\pgfsetdash{{1.850000pt}{0.800000pt}}{0.000000pt}%
\pgfpathmoveto{\pgfqpoint{2.181041in}{2.601404in}}%
\pgfpathlineto{\pgfqpoint{2.181041in}{3.507654in}}%
\pgfusepath{stroke}%
\end{pgfscope}%
\begin{pgfscope}%
\pgfsetbuttcap%
\pgfsetroundjoin%
\definecolor{currentfill}{rgb}{0.000000,0.000000,0.000000}%
\pgfsetfillcolor{currentfill}%
\pgfsetlinewidth{0.803000pt}%
\definecolor{currentstroke}{rgb}{0.000000,0.000000,0.000000}%
\pgfsetstrokecolor{currentstroke}%
\pgfsetdash{}{0pt}%
\pgfsys@defobject{currentmarker}{\pgfqpoint{0.000000in}{-0.048611in}}{\pgfqpoint{0.000000in}{0.000000in}}{%
\pgfpathmoveto{\pgfqpoint{0.000000in}{0.000000in}}%
\pgfpathlineto{\pgfqpoint{0.000000in}{-0.048611in}}%
\pgfusepath{stroke,fill}%
}%
\begin{pgfscope}%
\pgfsys@transformshift{2.181041in}{2.601404in}%
\pgfsys@useobject{currentmarker}{}%
\end{pgfscope}%
\end{pgfscope}%
\begin{pgfscope}%
\pgfpathrectangle{\pgfqpoint{0.573768in}{2.601404in}}{\pgfqpoint{2.720000in}{0.906250in}}%
\pgfusepath{clip}%
\pgfsetbuttcap%
\pgfsetroundjoin%
\pgfsetlinewidth{0.501875pt}%
\definecolor{currentstroke}{rgb}{0.690196,0.690196,0.690196}%
\pgfsetstrokecolor{currentstroke}%
\pgfsetstrokeopacity{0.250000}%
\pgfsetdash{{1.850000pt}{0.800000pt}}{0.000000pt}%
\pgfpathmoveto{\pgfqpoint{2.675587in}{2.601404in}}%
\pgfpathlineto{\pgfqpoint{2.675587in}{3.507654in}}%
\pgfusepath{stroke}%
\end{pgfscope}%
\begin{pgfscope}%
\pgfsetbuttcap%
\pgfsetroundjoin%
\definecolor{currentfill}{rgb}{0.000000,0.000000,0.000000}%
\pgfsetfillcolor{currentfill}%
\pgfsetlinewidth{0.803000pt}%
\definecolor{currentstroke}{rgb}{0.000000,0.000000,0.000000}%
\pgfsetstrokecolor{currentstroke}%
\pgfsetdash{}{0pt}%
\pgfsys@defobject{currentmarker}{\pgfqpoint{0.000000in}{-0.048611in}}{\pgfqpoint{0.000000in}{0.000000in}}{%
\pgfpathmoveto{\pgfqpoint{0.000000in}{0.000000in}}%
\pgfpathlineto{\pgfqpoint{0.000000in}{-0.048611in}}%
\pgfusepath{stroke,fill}%
}%
\begin{pgfscope}%
\pgfsys@transformshift{2.675587in}{2.601404in}%
\pgfsys@useobject{currentmarker}{}%
\end{pgfscope}%
\end{pgfscope}%
\begin{pgfscope}%
\pgfpathrectangle{\pgfqpoint{0.573768in}{2.601404in}}{\pgfqpoint{2.720000in}{0.906250in}}%
\pgfusepath{clip}%
\pgfsetbuttcap%
\pgfsetroundjoin%
\pgfsetlinewidth{0.501875pt}%
\definecolor{currentstroke}{rgb}{0.690196,0.690196,0.690196}%
\pgfsetstrokecolor{currentstroke}%
\pgfsetstrokeopacity{0.250000}%
\pgfsetdash{{1.850000pt}{0.800000pt}}{0.000000pt}%
\pgfpathmoveto{\pgfqpoint{3.170132in}{2.601404in}}%
\pgfpathlineto{\pgfqpoint{3.170132in}{3.507654in}}%
\pgfusepath{stroke}%
\end{pgfscope}%
\begin{pgfscope}%
\pgfsetbuttcap%
\pgfsetroundjoin%
\definecolor{currentfill}{rgb}{0.000000,0.000000,0.000000}%
\pgfsetfillcolor{currentfill}%
\pgfsetlinewidth{0.803000pt}%
\definecolor{currentstroke}{rgb}{0.000000,0.000000,0.000000}%
\pgfsetstrokecolor{currentstroke}%
\pgfsetdash{}{0pt}%
\pgfsys@defobject{currentmarker}{\pgfqpoint{0.000000in}{-0.048611in}}{\pgfqpoint{0.000000in}{0.000000in}}{%
\pgfpathmoveto{\pgfqpoint{0.000000in}{0.000000in}}%
\pgfpathlineto{\pgfqpoint{0.000000in}{-0.048611in}}%
\pgfusepath{stroke,fill}%
}%
\begin{pgfscope}%
\pgfsys@transformshift{3.170132in}{2.601404in}%
\pgfsys@useobject{currentmarker}{}%
\end{pgfscope}%
\end{pgfscope}%
\begin{pgfscope}%
\pgfpathrectangle{\pgfqpoint{0.573768in}{2.601404in}}{\pgfqpoint{2.720000in}{0.906250in}}%
\pgfusepath{clip}%
\pgfsetbuttcap%
\pgfsetroundjoin%
\pgfsetlinewidth{0.501875pt}%
\definecolor{currentstroke}{rgb}{0.690196,0.690196,0.690196}%
\pgfsetstrokecolor{currentstroke}%
\pgfsetstrokeopacity{0.250000}%
\pgfsetdash{{1.850000pt}{0.800000pt}}{0.000000pt}%
\pgfpathmoveto{\pgfqpoint{0.573768in}{2.644559in}}%
\pgfpathlineto{\pgfqpoint{3.293768in}{2.644559in}}%
\pgfusepath{stroke}%
\end{pgfscope}%
\begin{pgfscope}%
\pgfsetbuttcap%
\pgfsetroundjoin%
\definecolor{currentfill}{rgb}{0.000000,0.000000,0.000000}%
\pgfsetfillcolor{currentfill}%
\pgfsetlinewidth{0.803000pt}%
\definecolor{currentstroke}{rgb}{0.000000,0.000000,0.000000}%
\pgfsetstrokecolor{currentstroke}%
\pgfsetdash{}{0pt}%
\pgfsys@defobject{currentmarker}{\pgfqpoint{-0.048611in}{0.000000in}}{\pgfqpoint{-0.000000in}{0.000000in}}{%
\pgfpathmoveto{\pgfqpoint{-0.000000in}{0.000000in}}%
\pgfpathlineto{\pgfqpoint{-0.048611in}{0.000000in}}%
\pgfusepath{stroke,fill}%
}%
\begin{pgfscope}%
\pgfsys@transformshift{0.573768in}{2.644559in}%
\pgfsys@useobject{currentmarker}{}%
\end{pgfscope}%
\end{pgfscope}%
\begin{pgfscope}%
\definecolor{textcolor}{rgb}{0.000000,0.000000,0.000000}%
\pgfsetstrokecolor{textcolor}%
\pgfsetfillcolor{textcolor}%
\pgftext[x=0.417518in, y=2.605979in, left, base]{\color{textcolor}{\rmfamily\fontsize{8.000000}{9.600000}\selectfont\catcode`\^=\active\def^{\ifmmode\sp\else\^{}\fi}\catcode`\%=\active\def%{\%}$\mathdefault{0}$}}%
\end{pgfscope}%
\begin{pgfscope}%
\pgfpathrectangle{\pgfqpoint{0.573768in}{2.601404in}}{\pgfqpoint{2.720000in}{0.906250in}}%
\pgfusepath{clip}%
\pgfsetbuttcap%
\pgfsetroundjoin%
\pgfsetlinewidth{0.501875pt}%
\definecolor{currentstroke}{rgb}{0.690196,0.690196,0.690196}%
\pgfsetstrokecolor{currentstroke}%
\pgfsetstrokeopacity{0.250000}%
\pgfsetdash{{1.850000pt}{0.800000pt}}{0.000000pt}%
\pgfpathmoveto{\pgfqpoint{0.573768in}{3.076107in}}%
\pgfpathlineto{\pgfqpoint{3.293768in}{3.076107in}}%
\pgfusepath{stroke}%
\end{pgfscope}%
\begin{pgfscope}%
\pgfsetbuttcap%
\pgfsetroundjoin%
\definecolor{currentfill}{rgb}{0.000000,0.000000,0.000000}%
\pgfsetfillcolor{currentfill}%
\pgfsetlinewidth{0.803000pt}%
\definecolor{currentstroke}{rgb}{0.000000,0.000000,0.000000}%
\pgfsetstrokecolor{currentstroke}%
\pgfsetdash{}{0pt}%
\pgfsys@defobject{currentmarker}{\pgfqpoint{-0.048611in}{0.000000in}}{\pgfqpoint{-0.000000in}{0.000000in}}{%
\pgfpathmoveto{\pgfqpoint{-0.000000in}{0.000000in}}%
\pgfpathlineto{\pgfqpoint{-0.048611in}{0.000000in}}%
\pgfusepath{stroke,fill}%
}%
\begin{pgfscope}%
\pgfsys@transformshift{0.573768in}{3.076107in}%
\pgfsys@useobject{currentmarker}{}%
\end{pgfscope}%
\end{pgfscope}%
\begin{pgfscope}%
\definecolor{textcolor}{rgb}{0.000000,0.000000,0.000000}%
\pgfsetstrokecolor{textcolor}%
\pgfsetfillcolor{textcolor}%
\pgftext[x=0.417518in, y=3.037527in, left, base]{\color{textcolor}{\rmfamily\fontsize{8.000000}{9.600000}\selectfont\catcode`\^=\active\def^{\ifmmode\sp\else\^{}\fi}\catcode`\%=\active\def%{\%}$\mathdefault{5}$}}%
\end{pgfscope}%
\begin{pgfscope}%
\pgfpathrectangle{\pgfqpoint{0.573768in}{2.601404in}}{\pgfqpoint{2.720000in}{0.906250in}}%
\pgfusepath{clip}%
\pgfsetbuttcap%
\pgfsetroundjoin%
\pgfsetlinewidth{0.501875pt}%
\definecolor{currentstroke}{rgb}{0.690196,0.690196,0.690196}%
\pgfsetstrokecolor{currentstroke}%
\pgfsetstrokeopacity{0.250000}%
\pgfsetdash{{1.850000pt}{0.800000pt}}{0.000000pt}%
\pgfpathmoveto{\pgfqpoint{0.573768in}{3.507654in}}%
\pgfpathlineto{\pgfqpoint{3.293768in}{3.507654in}}%
\pgfusepath{stroke}%
\end{pgfscope}%
\begin{pgfscope}%
\pgfsetbuttcap%
\pgfsetroundjoin%
\definecolor{currentfill}{rgb}{0.000000,0.000000,0.000000}%
\pgfsetfillcolor{currentfill}%
\pgfsetlinewidth{0.803000pt}%
\definecolor{currentstroke}{rgb}{0.000000,0.000000,0.000000}%
\pgfsetstrokecolor{currentstroke}%
\pgfsetdash{}{0pt}%
\pgfsys@defobject{currentmarker}{\pgfqpoint{-0.048611in}{0.000000in}}{\pgfqpoint{-0.000000in}{0.000000in}}{%
\pgfpathmoveto{\pgfqpoint{-0.000000in}{0.000000in}}%
\pgfpathlineto{\pgfqpoint{-0.048611in}{0.000000in}}%
\pgfusepath{stroke,fill}%
}%
\begin{pgfscope}%
\pgfsys@transformshift{0.573768in}{3.507654in}%
\pgfsys@useobject{currentmarker}{}%
\end{pgfscope}%
\end{pgfscope}%
\begin{pgfscope}%
\definecolor{textcolor}{rgb}{0.000000,0.000000,0.000000}%
\pgfsetstrokecolor{textcolor}%
\pgfsetfillcolor{textcolor}%
\pgftext[x=0.358489in, y=3.469074in, left, base]{\color{textcolor}{\rmfamily\fontsize{8.000000}{9.600000}\selectfont\catcode`\^=\active\def^{\ifmmode\sp\else\^{}\fi}\catcode`\%=\active\def%{\%}$\mathdefault{10}$}}%
\end{pgfscope}%
\begin{pgfscope}%
\definecolor{textcolor}{rgb}{0.000000,0.000000,0.000000}%
\pgfsetstrokecolor{textcolor}%
\pgfsetfillcolor{textcolor}%
\pgftext[x=0.302933in,y=3.054529in,,bottom,rotate=90.000000]{\color{textcolor}{\rmfamily\fontsize{8.000000}{9.600000}\selectfont\catcode`\^=\active\def^{\ifmmode\sp\else\^{}\fi}\catcode`\%=\active\def%{\%}$i_D$ (A)}}%
\end{pgfscope}%
\begin{pgfscope}%
\pgfpathrectangle{\pgfqpoint{0.573768in}{2.601404in}}{\pgfqpoint{2.720000in}{0.906250in}}%
\pgfusepath{clip}%
\pgfsetrectcap%
\pgfsetroundjoin%
\pgfsetlinewidth{1.505625pt}%
\definecolor{currentstroke}{rgb}{0.000000,0.447059,0.698039}%
\pgfsetstrokecolor{currentstroke}%
\pgfsetdash{}{0pt}%
\pgfpathmoveto{\pgfqpoint{0.697405in}{2.644559in}}%
\pgfpathlineto{\pgfqpoint{0.697409in}{2.644559in}}%
\pgfpathlineto{\pgfqpoint{0.697416in}{2.743950in}}%
\pgfpathlineto{\pgfqpoint{0.698332in}{2.702816in}}%
\pgfpathlineto{\pgfqpoint{0.699384in}{2.768855in}}%
\pgfpathlineto{\pgfqpoint{0.699411in}{2.755130in}}%
\pgfpathlineto{\pgfqpoint{0.700203in}{2.644559in}}%
\pgfpathlineto{\pgfqpoint{0.702355in}{2.644559in}}%
\pgfpathlineto{\pgfqpoint{0.702366in}{2.872219in}}%
\pgfpathlineto{\pgfqpoint{0.703136in}{2.801147in}}%
\pgfpathlineto{\pgfqpoint{0.703977in}{2.880938in}}%
\pgfpathlineto{\pgfqpoint{0.704329in}{2.902987in}}%
\pgfpathlineto{\pgfqpoint{0.705237in}{2.644559in}}%
\pgfpathlineto{\pgfqpoint{0.707302in}{2.644561in}}%
\pgfpathlineto{\pgfqpoint{0.707319in}{2.953752in}}%
\pgfpathlineto{\pgfqpoint{0.708079in}{2.904437in}}%
\pgfpathlineto{\pgfqpoint{0.708920in}{3.007676in}}%
\pgfpathlineto{\pgfqpoint{0.709274in}{3.029873in}}%
\pgfpathlineto{\pgfqpoint{0.710070in}{2.644559in}}%
\pgfpathlineto{\pgfqpoint{0.712248in}{2.644580in}}%
\pgfpathlineto{\pgfqpoint{0.712269in}{3.057923in}}%
\pgfpathlineto{\pgfqpoint{0.713019in}{3.057248in}}%
\pgfpathlineto{\pgfqpoint{0.713069in}{3.017924in}}%
\pgfpathlineto{\pgfqpoint{0.713712in}{3.100638in}}%
\pgfpathlineto{\pgfqpoint{0.713761in}{3.061474in}}%
\pgfpathlineto{\pgfqpoint{0.714188in}{3.130499in}}%
\pgfpathlineto{\pgfqpoint{0.714286in}{2.976064in}}%
\pgfpathlineto{\pgfqpoint{0.715069in}{2.644559in}}%
\pgfpathlineto{\pgfqpoint{0.717191in}{2.644559in}}%
\pgfpathlineto{\pgfqpoint{0.717214in}{3.150438in}}%
\pgfpathlineto{\pgfqpoint{0.718077in}{3.136816in}}%
\pgfpathlineto{\pgfqpoint{0.719165in}{3.207938in}}%
\pgfpathlineto{\pgfqpoint{0.719166in}{3.182740in}}%
\pgfpathlineto{\pgfqpoint{0.719192in}{3.189227in}}%
\pgfpathlineto{\pgfqpoint{0.719941in}{2.644559in}}%
\pgfpathlineto{\pgfqpoint{0.722137in}{2.644559in}}%
\pgfpathlineto{\pgfqpoint{0.722965in}{3.246776in}}%
\pgfpathlineto{\pgfqpoint{0.723014in}{3.188196in}}%
\pgfpathlineto{\pgfqpoint{0.723657in}{3.290106in}}%
\pgfpathlineto{\pgfqpoint{0.723707in}{3.231781in}}%
\pgfpathlineto{\pgfqpoint{0.724111in}{3.318476in}}%
\pgfpathlineto{\pgfqpoint{0.724182in}{3.096169in}}%
\pgfpathlineto{\pgfqpoint{0.724968in}{2.644559in}}%
\pgfpathlineto{\pgfqpoint{0.727082in}{2.644559in}}%
\pgfpathlineto{\pgfqpoint{0.727799in}{3.302398in}}%
\pgfpathlineto{\pgfqpoint{0.727848in}{3.249287in}}%
\pgfpathlineto{\pgfqpoint{0.728689in}{3.358112in}}%
\pgfpathlineto{\pgfqpoint{0.729056in}{3.381073in}}%
\pgfpathlineto{\pgfqpoint{0.730127in}{2.644559in}}%
\pgfpathlineto{\pgfqpoint{0.732029in}{2.644632in}}%
\pgfpathlineto{\pgfqpoint{0.732743in}{3.343412in}}%
\pgfpathlineto{\pgfqpoint{0.732793in}{3.303641in}}%
\pgfpathlineto{\pgfqpoint{0.733633in}{3.399157in}}%
\pgfpathlineto{\pgfqpoint{0.734002in}{3.422218in}}%
\pgfpathlineto{\pgfqpoint{0.735346in}{2.644559in}}%
\pgfpathlineto{\pgfqpoint{0.736975in}{2.644564in}}%
\pgfpathlineto{\pgfqpoint{0.737000in}{3.377927in}}%
\pgfpathlineto{\pgfqpoint{0.737730in}{3.372579in}}%
\pgfpathlineto{\pgfqpoint{0.737780in}{3.338286in}}%
\pgfpathlineto{\pgfqpoint{0.738423in}{3.415956in}}%
\pgfpathlineto{\pgfqpoint{0.738472in}{3.381804in}}%
\pgfpathlineto{\pgfqpoint{0.738908in}{3.446321in}}%
\pgfpathlineto{\pgfqpoint{0.739016in}{3.347183in}}%
\pgfpathlineto{\pgfqpoint{0.739811in}{2.644559in}}%
\pgfpathlineto{\pgfqpoint{0.741918in}{2.644559in}}%
\pgfpathlineto{\pgfqpoint{0.742890in}{3.391782in}}%
\pgfpathlineto{\pgfqpoint{0.742939in}{3.365946in}}%
\pgfpathlineto{\pgfqpoint{0.743780in}{3.447546in}}%
\pgfpathlineto{\pgfqpoint{0.743948in}{3.356647in}}%
\pgfpathlineto{\pgfqpoint{0.743861in}{3.452624in}}%
\pgfpathlineto{\pgfqpoint{0.743963in}{3.359704in}}%
\pgfpathlineto{\pgfqpoint{0.744757in}{2.644559in}}%
\pgfpathlineto{\pgfqpoint{0.746863in}{2.644559in}}%
\pgfpathlineto{\pgfqpoint{0.746877in}{2.845413in}}%
\pgfpathlineto{\pgfqpoint{0.747562in}{3.377124in}}%
\pgfpathlineto{\pgfqpoint{0.747612in}{3.328099in}}%
\pgfpathlineto{\pgfqpoint{0.748452in}{3.432849in}}%
\pgfpathlineto{\pgfqpoint{0.748838in}{3.456982in}}%
\pgfpathlineto{\pgfqpoint{0.750022in}{2.644559in}}%
\pgfpathlineto{\pgfqpoint{0.751809in}{2.644559in}}%
\pgfpathlineto{\pgfqpoint{0.752932in}{3.383945in}}%
\pgfpathlineto{\pgfqpoint{0.752981in}{3.318862in}}%
\pgfpathlineto{\pgfqpoint{0.753624in}{3.427246in}}%
\pgfpathlineto{\pgfqpoint{0.753674in}{3.362457in}}%
\pgfpathlineto{\pgfqpoint{0.753784in}{3.437181in}}%
\pgfpathlineto{\pgfqpoint{0.753893in}{2.644559in}}%
\pgfpathlineto{\pgfqpoint{0.753941in}{2.644559in}}%
\pgfpathlineto{\pgfqpoint{0.756757in}{2.644561in}}%
\pgfpathlineto{\pgfqpoint{0.757478in}{3.308123in}}%
\pgfpathlineto{\pgfqpoint{0.757527in}{3.253494in}}%
\pgfpathlineto{\pgfqpoint{0.758368in}{3.363830in}}%
\pgfpathlineto{\pgfqpoint{0.758729in}{3.386388in}}%
\pgfpathlineto{\pgfqpoint{0.759792in}{2.644559in}}%
\pgfpathlineto{\pgfqpoint{0.761700in}{2.644559in}}%
\pgfpathlineto{\pgfqpoint{0.761722in}{3.264345in}}%
\pgfpathlineto{\pgfqpoint{0.762509in}{3.229631in}}%
\pgfpathlineto{\pgfqpoint{0.762558in}{3.219118in}}%
\pgfpathlineto{\pgfqpoint{0.763102in}{3.266850in}}%
\pgfpathlineto{\pgfqpoint{0.763152in}{3.256397in}}%
\pgfpathlineto{\pgfqpoint{0.763674in}{3.302741in}}%
\pgfpathlineto{\pgfqpoint{0.763730in}{3.208137in}}%
\pgfpathlineto{\pgfqpoint{0.763743in}{3.280774in}}%
\pgfpathlineto{\pgfqpoint{0.764553in}{2.644559in}}%
\pgfpathlineto{\pgfqpoint{0.766645in}{2.644559in}}%
\pgfpathlineto{\pgfqpoint{0.766667in}{3.190810in}}%
\pgfpathlineto{\pgfqpoint{0.767456in}{3.126962in}}%
\pgfpathlineto{\pgfqpoint{0.768296in}{3.219962in}}%
\pgfpathlineto{\pgfqpoint{0.768582in}{3.237845in}}%
\pgfpathlineto{\pgfqpoint{0.769834in}{2.644559in}}%
\pgfpathlineto{\pgfqpoint{0.771593in}{2.644705in}}%
\pgfpathlineto{\pgfqpoint{0.771611in}{3.107364in}}%
\pgfpathlineto{\pgfqpoint{0.772369in}{3.067340in}}%
\pgfpathlineto{\pgfqpoint{0.772418in}{3.050303in}}%
\pgfpathlineto{\pgfqpoint{0.772962in}{3.104554in}}%
\pgfpathlineto{\pgfqpoint{0.773012in}{3.087607in}}%
\pgfpathlineto{\pgfqpoint{0.773536in}{3.140518in}}%
\pgfpathlineto{\pgfqpoint{0.773623in}{3.064439in}}%
\pgfpathlineto{\pgfqpoint{0.773638in}{3.101340in}}%
\pgfpathlineto{\pgfqpoint{0.774431in}{2.644559in}}%
\pgfpathlineto{\pgfqpoint{0.776537in}{2.644559in}}%
\pgfpathlineto{\pgfqpoint{0.776557in}{3.011689in}}%
\pgfpathlineto{\pgfqpoint{0.777292in}{2.948975in}}%
\pgfpathlineto{\pgfqpoint{0.778231in}{3.027423in}}%
\pgfpathlineto{\pgfqpoint{0.778511in}{3.044968in}}%
\pgfpathlineto{\pgfqpoint{0.778593in}{2.852145in}}%
\pgfpathlineto{\pgfqpoint{0.779389in}{2.644559in}}%
\pgfpathlineto{\pgfqpoint{0.781481in}{2.644559in}}%
\pgfpathlineto{\pgfqpoint{0.781495in}{2.838651in}}%
\pgfpathlineto{\pgfqpoint{0.781500in}{2.943672in}}%
\pgfpathlineto{\pgfqpoint{0.781595in}{2.791194in}}%
\pgfpathlineto{\pgfqpoint{0.782221in}{2.878675in}}%
\pgfpathlineto{\pgfqpoint{0.782270in}{2.832566in}}%
\pgfpathlineto{\pgfqpoint{0.782913in}{2.922036in}}%
\pgfpathlineto{\pgfqpoint{0.782962in}{2.876172in}}%
\pgfpathlineto{\pgfqpoint{0.783407in}{2.953007in}}%
\pgfpathlineto{\pgfqpoint{0.783511in}{2.872213in}}%
\pgfpathlineto{\pgfqpoint{0.783526in}{2.891716in}}%
\pgfpathlineto{\pgfqpoint{0.784316in}{2.644559in}}%
\pgfpathlineto{\pgfqpoint{0.786427in}{2.644559in}}%
\pgfpathlineto{\pgfqpoint{0.786439in}{2.816302in}}%
\pgfpathlineto{\pgfqpoint{0.787206in}{2.742957in}}%
\pgfpathlineto{\pgfqpoint{0.788146in}{2.814908in}}%
\pgfpathlineto{\pgfqpoint{0.788402in}{2.830955in}}%
\pgfpathlineto{\pgfqpoint{0.789776in}{2.644559in}}%
\pgfpathlineto{\pgfqpoint{0.791373in}{2.644561in}}%
\pgfpathlineto{\pgfqpoint{0.791386in}{2.764795in}}%
\pgfpathlineto{\pgfqpoint{0.791501in}{2.631418in}}%
\pgfpathlineto{\pgfqpoint{0.792139in}{2.725114in}}%
\pgfpathlineto{\pgfqpoint{0.792188in}{2.674255in}}%
\pgfpathlineto{\pgfqpoint{0.792831in}{2.768502in}}%
\pgfpathlineto{\pgfqpoint{0.792881in}{2.717855in}}%
\pgfpathlineto{\pgfqpoint{0.793312in}{2.798611in}}%
\pgfpathlineto{\pgfqpoint{0.793440in}{2.644559in}}%
\pgfpathlineto{\pgfqpoint{0.793588in}{2.644559in}}%
\pgfpathlineto{\pgfqpoint{0.796319in}{2.644559in}}%
\pgfpathlineto{\pgfqpoint{0.796323in}{2.760880in}}%
\pgfpathlineto{\pgfqpoint{0.797050in}{2.686732in}}%
\pgfpathlineto{\pgfqpoint{0.797891in}{2.761121in}}%
\pgfpathlineto{\pgfqpoint{0.798265in}{2.784591in}}%
\pgfpathlineto{\pgfqpoint{0.798972in}{2.644559in}}%
\pgfpathlineto{\pgfqpoint{0.801264in}{2.644559in}}%
\pgfpathlineto{\pgfqpoint{0.801269in}{2.750716in}}%
\pgfpathlineto{\pgfqpoint{0.801327in}{2.643175in}}%
\pgfpathlineto{\pgfqpoint{0.802005in}{2.681453in}}%
\pgfpathlineto{\pgfqpoint{0.802845in}{2.766290in}}%
\pgfpathlineto{\pgfqpoint{0.803238in}{2.790935in}}%
\pgfpathlineto{\pgfqpoint{0.803824in}{2.644559in}}%
\pgfpathlineto{\pgfqpoint{0.806210in}{2.644559in}}%
\pgfpathlineto{\pgfqpoint{0.806215in}{2.750587in}}%
\pgfpathlineto{\pgfqpoint{0.806273in}{2.643849in}}%
\pgfpathlineto{\pgfqpoint{0.806954in}{2.682427in}}%
\pgfpathlineto{\pgfqpoint{0.807795in}{2.765552in}}%
\pgfpathlineto{\pgfqpoint{0.808184in}{2.789952in}}%
\pgfpathlineto{\pgfqpoint{0.808785in}{2.644559in}}%
\pgfpathlineto{\pgfqpoint{0.811155in}{2.644559in}}%
\pgfpathlineto{\pgfqpoint{0.811160in}{2.749652in}}%
\pgfpathlineto{\pgfqpoint{0.811265in}{2.640412in}}%
\pgfpathlineto{\pgfqpoint{0.811905in}{2.715097in}}%
\pgfpathlineto{\pgfqpoint{0.811954in}{2.683258in}}%
\pgfpathlineto{\pgfqpoint{0.812597in}{2.758525in}}%
\pgfpathlineto{\pgfqpoint{0.812647in}{2.726819in}}%
\pgfpathlineto{\pgfqpoint{0.813086in}{2.789162in}}%
\pgfpathlineto{\pgfqpoint{0.813226in}{2.644559in}}%
\pgfpathlineto{\pgfqpoint{0.813308in}{2.644559in}}%
\pgfpathlineto{\pgfqpoint{0.816101in}{2.644559in}}%
\pgfpathlineto{\pgfqpoint{0.816105in}{2.749034in}}%
\pgfpathlineto{\pgfqpoint{0.816232in}{2.636493in}}%
\pgfpathlineto{\pgfqpoint{0.816875in}{2.721606in}}%
\pgfpathlineto{\pgfqpoint{0.816924in}{2.679865in}}%
\pgfpathlineto{\pgfqpoint{0.817567in}{2.765010in}}%
\pgfpathlineto{\pgfqpoint{0.817617in}{2.723450in}}%
\pgfpathlineto{\pgfqpoint{0.818043in}{2.794851in}}%
\pgfpathlineto{\pgfqpoint{0.818168in}{2.644559in}}%
\pgfpathlineto{\pgfqpoint{0.818305in}{2.644559in}}%
\pgfpathlineto{\pgfqpoint{0.821046in}{2.644559in}}%
\pgfpathlineto{\pgfqpoint{0.821051in}{2.747115in}}%
\pgfpathlineto{\pgfqpoint{0.821119in}{2.621430in}}%
\pgfpathlineto{\pgfqpoint{0.821803in}{2.661762in}}%
\pgfpathlineto{\pgfqpoint{0.822644in}{2.786249in}}%
\pgfpathlineto{\pgfqpoint{0.823020in}{2.809754in}}%
\pgfpathlineto{\pgfqpoint{0.823517in}{2.644559in}}%
\pgfpathlineto{\pgfqpoint{0.825992in}{2.644559in}}%
\pgfpathlineto{\pgfqpoint{0.825996in}{2.746546in}}%
\pgfpathlineto{\pgfqpoint{0.826069in}{2.620157in}}%
\pgfpathlineto{\pgfqpoint{0.826754in}{2.661048in}}%
\pgfpathlineto{\pgfqpoint{0.827595in}{2.787529in}}%
\pgfpathlineto{\pgfqpoint{0.827965in}{2.810711in}}%
\pgfpathlineto{\pgfqpoint{0.828467in}{2.644559in}}%
\pgfpathlineto{\pgfqpoint{0.830938in}{2.644562in}}%
\pgfpathlineto{\pgfqpoint{0.830942in}{2.788099in}}%
\pgfpathlineto{\pgfqpoint{0.831035in}{2.627248in}}%
\pgfpathlineto{\pgfqpoint{0.831674in}{2.725397in}}%
\pgfpathlineto{\pgfqpoint{0.831723in}{2.669715in}}%
\pgfpathlineto{\pgfqpoint{0.832366in}{2.768737in}}%
\pgfpathlineto{\pgfqpoint{0.832416in}{2.713364in}}%
\pgfpathlineto{\pgfqpoint{0.832911in}{2.802794in}}%
\pgfpathlineto{\pgfqpoint{0.833001in}{2.644559in}}%
\pgfpathlineto{\pgfqpoint{0.833128in}{2.644559in}}%
\pgfpathlineto{\pgfqpoint{0.835882in}{2.644559in}}%
\pgfpathlineto{\pgfqpoint{0.835887in}{2.741667in}}%
\pgfpathlineto{\pgfqpoint{0.835965in}{2.628489in}}%
\pgfpathlineto{\pgfqpoint{0.836651in}{2.669879in}}%
\pgfpathlineto{\pgfqpoint{0.837491in}{2.778373in}}%
\pgfpathlineto{\pgfqpoint{0.837856in}{2.801215in}}%
\pgfpathlineto{\pgfqpoint{0.838387in}{2.644559in}}%
\pgfpathlineto{\pgfqpoint{0.840828in}{2.644559in}}%
\pgfpathlineto{\pgfqpoint{0.840833in}{2.739332in}}%
\pgfpathlineto{\pgfqpoint{0.840912in}{2.633938in}}%
\pgfpathlineto{\pgfqpoint{0.841598in}{2.675432in}}%
\pgfpathlineto{\pgfqpoint{0.842438in}{2.772232in}}%
\pgfpathlineto{\pgfqpoint{0.842802in}{2.794977in}}%
\pgfpathlineto{\pgfqpoint{0.843362in}{2.644559in}}%
\pgfpathlineto{\pgfqpoint{0.845773in}{2.644559in}}%
\pgfpathlineto{\pgfqpoint{0.845778in}{2.739290in}}%
\pgfpathlineto{\pgfqpoint{0.845857in}{2.634294in}}%
\pgfpathlineto{\pgfqpoint{0.846543in}{2.675765in}}%
\pgfpathlineto{\pgfqpoint{0.847384in}{2.771984in}}%
\pgfpathlineto{\pgfqpoint{0.847747in}{2.794748in}}%
\pgfpathlineto{\pgfqpoint{0.848312in}{2.644559in}}%
\pgfpathlineto{\pgfqpoint{0.850720in}{2.644692in}}%
\pgfpathlineto{\pgfqpoint{0.850723in}{2.742607in}}%
\pgfpathlineto{\pgfqpoint{0.850842in}{2.633582in}}%
\pgfpathlineto{\pgfqpoint{0.851485in}{2.719972in}}%
\pgfpathlineto{\pgfqpoint{0.851535in}{2.676885in}}%
\pgfpathlineto{\pgfqpoint{0.852177in}{2.763372in}}%
\pgfpathlineto{\pgfqpoint{0.852227in}{2.720474in}}%
\pgfpathlineto{\pgfqpoint{0.852658in}{2.793460in}}%
\pgfpathlineto{\pgfqpoint{0.852785in}{2.644559in}}%
\pgfpathlineto{\pgfqpoint{0.852918in}{2.644559in}}%
\pgfpathlineto{\pgfqpoint{0.855664in}{2.644559in}}%
\pgfpathlineto{\pgfqpoint{0.855669in}{2.733425in}}%
\pgfpathlineto{\pgfqpoint{0.855734in}{2.641999in}}%
\pgfpathlineto{\pgfqpoint{0.856419in}{2.682053in}}%
\pgfpathlineto{\pgfqpoint{0.857260in}{2.762530in}}%
\pgfpathlineto{\pgfqpoint{0.857638in}{2.786252in}}%
\pgfpathlineto{\pgfqpoint{0.858269in}{2.644559in}}%
\pgfpathlineto{\pgfqpoint{0.860611in}{2.645049in}}%
\pgfpathlineto{\pgfqpoint{0.860614in}{2.766235in}}%
\pgfpathlineto{\pgfqpoint{0.860701in}{2.619193in}}%
\pgfpathlineto{\pgfqpoint{0.861338in}{2.729998in}}%
\pgfpathlineto{\pgfqpoint{0.861387in}{2.661275in}}%
\pgfpathlineto{\pgfqpoint{0.862030in}{2.773311in}}%
\pgfpathlineto{\pgfqpoint{0.862080in}{2.704951in}}%
\pgfpathlineto{\pgfqpoint{0.862584in}{2.807905in}}%
\pgfpathlineto{\pgfqpoint{0.862690in}{2.644559in}}%
\pgfpathlineto{\pgfqpoint{0.862809in}{2.644559in}}%
\pgfpathlineto{\pgfqpoint{0.865557in}{2.644835in}}%
\pgfpathlineto{\pgfqpoint{0.865559in}{2.758721in}}%
\pgfpathlineto{\pgfqpoint{0.865655in}{2.611732in}}%
\pgfpathlineto{\pgfqpoint{0.866292in}{2.737370in}}%
\pgfpathlineto{\pgfqpoint{0.866342in}{2.654322in}}%
\pgfpathlineto{\pgfqpoint{0.866985in}{2.780650in}}%
\pgfpathlineto{\pgfqpoint{0.867034in}{2.698032in}}%
\pgfpathlineto{\pgfqpoint{0.867529in}{2.814617in}}%
\pgfpathlineto{\pgfqpoint{0.867621in}{2.644559in}}%
\pgfpathlineto{\pgfqpoint{0.867774in}{2.644559in}}%
\pgfpathlineto{\pgfqpoint{0.870502in}{2.644836in}}%
\pgfpathlineto{\pgfqpoint{0.870505in}{2.759888in}}%
\pgfpathlineto{\pgfqpoint{0.870600in}{2.610986in}}%
\pgfpathlineto{\pgfqpoint{0.871238in}{2.737989in}}%
\pgfpathlineto{\pgfqpoint{0.871288in}{2.653584in}}%
\pgfpathlineto{\pgfqpoint{0.871931in}{2.781268in}}%
\pgfpathlineto{\pgfqpoint{0.871980in}{2.697294in}}%
\pgfpathlineto{\pgfqpoint{0.872425in}{2.812180in}}%
\pgfpathlineto{\pgfqpoint{0.872584in}{2.644559in}}%
\pgfpathlineto{\pgfqpoint{0.872722in}{2.644559in}}%
\pgfpathlineto{\pgfqpoint{0.875448in}{2.644702in}}%
\pgfpathlineto{\pgfqpoint{0.875450in}{2.750957in}}%
\pgfpathlineto{\pgfqpoint{0.875547in}{2.620689in}}%
\pgfpathlineto{\pgfqpoint{0.876185in}{2.727572in}}%
\pgfpathlineto{\pgfqpoint{0.876234in}{2.663298in}}%
\pgfpathlineto{\pgfqpoint{0.876877in}{2.770889in}}%
\pgfpathlineto{\pgfqpoint{0.876926in}{2.706970in}}%
\pgfpathlineto{\pgfqpoint{0.877371in}{2.801829in}}%
\pgfpathlineto{\pgfqpoint{0.877508in}{2.644559in}}%
\pgfpathlineto{\pgfqpoint{0.877651in}{2.644559in}}%
\pgfpathlineto{\pgfqpoint{0.880393in}{2.644682in}}%
\pgfpathlineto{\pgfqpoint{0.880395in}{2.749107in}}%
\pgfpathlineto{\pgfqpoint{0.880493in}{2.624077in}}%
\pgfpathlineto{\pgfqpoint{0.881131in}{2.724074in}}%
\pgfpathlineto{\pgfqpoint{0.881180in}{2.666715in}}%
\pgfpathlineto{\pgfqpoint{0.881823in}{2.767405in}}%
\pgfpathlineto{\pgfqpoint{0.881873in}{2.710373in}}%
\pgfpathlineto{\pgfqpoint{0.882318in}{2.798354in}}%
\pgfpathlineto{\pgfqpoint{0.882456in}{2.644559in}}%
\pgfpathlineto{\pgfqpoint{0.882592in}{2.644559in}}%
\pgfpathlineto{\pgfqpoint{0.885337in}{2.644559in}}%
\pgfpathlineto{\pgfqpoint{0.885344in}{2.732255in}}%
\pgfpathlineto{\pgfqpoint{0.885405in}{2.628468in}}%
\pgfpathlineto{\pgfqpoint{0.886090in}{2.668296in}}%
\pgfpathlineto{\pgfqpoint{0.886931in}{2.773118in}}%
\pgfpathlineto{\pgfqpoint{0.887311in}{2.796930in}}%
\pgfpathlineto{\pgfqpoint{0.887832in}{2.644559in}}%
\pgfpathlineto{\pgfqpoint{0.890282in}{2.644559in}}%
\pgfpathlineto{\pgfqpoint{0.890290in}{2.733778in}}%
\pgfpathlineto{\pgfqpoint{0.890351in}{2.629786in}}%
\pgfpathlineto{\pgfqpoint{0.891037in}{2.669766in}}%
\pgfpathlineto{\pgfqpoint{0.891878in}{2.771561in}}%
\pgfpathlineto{\pgfqpoint{0.892256in}{2.795277in}}%
\pgfpathlineto{\pgfqpoint{0.892784in}{2.644559in}}%
\pgfpathlineto{\pgfqpoint{0.895227in}{2.644559in}}%
\pgfpathlineto{\pgfqpoint{0.895235in}{2.734443in}}%
\pgfpathlineto{\pgfqpoint{0.895295in}{2.632146in}}%
\pgfpathlineto{\pgfqpoint{0.895982in}{2.671969in}}%
\pgfpathlineto{\pgfqpoint{0.896822in}{2.768718in}}%
\pgfpathlineto{\pgfqpoint{0.897202in}{2.792509in}}%
\pgfpathlineto{\pgfqpoint{0.897742in}{2.644559in}}%
\pgfpathlineto{\pgfqpoint{0.900173in}{2.644559in}}%
\pgfpathlineto{\pgfqpoint{0.900181in}{2.750374in}}%
\pgfpathlineto{\pgfqpoint{0.900251in}{2.625981in}}%
\pgfpathlineto{\pgfqpoint{0.900938in}{2.666972in}}%
\pgfpathlineto{\pgfqpoint{0.901778in}{2.774507in}}%
\pgfpathlineto{\pgfqpoint{0.902147in}{2.797612in}}%
\pgfpathlineto{\pgfqpoint{0.902673in}{2.644559in}}%
\pgfpathlineto{\pgfqpoint{0.905118in}{2.644559in}}%
\pgfpathlineto{\pgfqpoint{0.905123in}{2.748057in}}%
\pgfpathlineto{\pgfqpoint{0.905183in}{2.626121in}}%
\pgfpathlineto{\pgfqpoint{0.905916in}{2.723757in}}%
\pgfpathlineto{\pgfqpoint{0.905965in}{2.671730in}}%
\pgfpathlineto{\pgfqpoint{0.906608in}{2.767123in}}%
\pgfpathlineto{\pgfqpoint{0.906658in}{2.715354in}}%
\pgfpathlineto{\pgfqpoint{0.907093in}{2.797450in}}%
\pgfpathlineto{\pgfqpoint{0.907187in}{2.644559in}}%
\pgfpathlineto{\pgfqpoint{0.907369in}{2.644559in}}%
\pgfpathlineto{\pgfqpoint{0.910066in}{2.644632in}}%
\pgfpathlineto{\pgfqpoint{0.910068in}{2.725789in}}%
\pgfpathlineto{\pgfqpoint{0.910812in}{2.695825in}}%
\pgfpathlineto{\pgfqpoint{0.910862in}{2.692866in}}%
\pgfpathlineto{\pgfqpoint{0.911109in}{2.714459in}}%
\pgfpathlineto{\pgfqpoint{0.911159in}{2.711514in}}%
\pgfpathlineto{\pgfqpoint{0.911994in}{2.770022in}}%
\pgfpathlineto{\pgfqpoint{0.912061in}{2.754047in}}%
\pgfpathlineto{\pgfqpoint{0.912879in}{2.644559in}}%
\pgfpathlineto{\pgfqpoint{0.915011in}{2.644563in}}%
\pgfpathlineto{\pgfqpoint{0.915013in}{2.739540in}}%
\pgfpathlineto{\pgfqpoint{0.915100in}{2.613263in}}%
\pgfpathlineto{\pgfqpoint{0.915834in}{2.736555in}}%
\pgfpathlineto{\pgfqpoint{0.915884in}{2.661504in}}%
\pgfpathlineto{\pgfqpoint{0.916527in}{2.779865in}}%
\pgfpathlineto{\pgfqpoint{0.916576in}{2.705184in}}%
\pgfpathlineto{\pgfqpoint{0.916984in}{2.808422in}}%
\pgfpathlineto{\pgfqpoint{0.917089in}{2.644559in}}%
\pgfpathlineto{\pgfqpoint{0.917310in}{2.644559in}}%
\pgfpathlineto{\pgfqpoint{0.919957in}{2.644562in}}%
\pgfpathlineto{\pgfqpoint{0.919959in}{2.738000in}}%
\pgfpathlineto{\pgfqpoint{0.920042in}{2.613748in}}%
\pgfpathlineto{\pgfqpoint{0.920776in}{2.735129in}}%
\pgfpathlineto{\pgfqpoint{0.920826in}{2.661765in}}%
\pgfpathlineto{\pgfqpoint{0.921469in}{2.778443in}}%
\pgfpathlineto{\pgfqpoint{0.921518in}{2.705440in}}%
\pgfpathlineto{\pgfqpoint{0.921929in}{2.807212in}}%
\pgfpathlineto{\pgfqpoint{0.922035in}{2.644559in}}%
\pgfpathlineto{\pgfqpoint{0.922248in}{2.644559in}}%
\pgfpathlineto{\pgfqpoint{0.924902in}{2.644562in}}%
\pgfpathlineto{\pgfqpoint{0.924904in}{2.737625in}}%
\pgfpathlineto{\pgfqpoint{0.924984in}{2.614581in}}%
\pgfpathlineto{\pgfqpoint{0.925719in}{2.733906in}}%
\pgfpathlineto{\pgfqpoint{0.925768in}{2.662367in}}%
\pgfpathlineto{\pgfqpoint{0.926411in}{2.777225in}}%
\pgfpathlineto{\pgfqpoint{0.926461in}{2.706037in}}%
\pgfpathlineto{\pgfqpoint{0.926874in}{2.806193in}}%
\pgfpathlineto{\pgfqpoint{0.926980in}{2.644559in}}%
\pgfpathlineto{\pgfqpoint{0.927190in}{2.644559in}}%
\pgfpathlineto{\pgfqpoint{0.929848in}{2.644951in}}%
\pgfpathlineto{\pgfqpoint{0.929850in}{2.725975in}}%
\pgfpathlineto{\pgfqpoint{0.929938in}{2.621329in}}%
\pgfpathlineto{\pgfqpoint{0.930574in}{2.720955in}}%
\pgfpathlineto{\pgfqpoint{0.930624in}{2.663430in}}%
\pgfpathlineto{\pgfqpoint{0.931267in}{2.764302in}}%
\pgfpathlineto{\pgfqpoint{0.931316in}{2.707073in}}%
\pgfpathlineto{\pgfqpoint{0.931820in}{2.798905in}}%
\pgfpathlineto{\pgfqpoint{0.931909in}{2.644559in}}%
\pgfpathlineto{\pgfqpoint{0.932037in}{2.644559in}}%
\pgfpathlineto{\pgfqpoint{0.934794in}{2.644997in}}%
\pgfpathlineto{\pgfqpoint{0.934795in}{2.726624in}}%
\pgfpathlineto{\pgfqpoint{0.934882in}{2.622677in}}%
\pgfpathlineto{\pgfqpoint{0.935518in}{2.719128in}}%
\pgfpathlineto{\pgfqpoint{0.935567in}{2.664656in}}%
\pgfpathlineto{\pgfqpoint{0.936210in}{2.762480in}}%
\pgfpathlineto{\pgfqpoint{0.936260in}{2.708295in}}%
\pgfpathlineto{\pgfqpoint{0.936765in}{2.797221in}}%
\pgfpathlineto{\pgfqpoint{0.936870in}{2.644559in}}%
\pgfpathlineto{\pgfqpoint{0.936976in}{2.644559in}}%
\pgfpathlineto{\pgfqpoint{0.939739in}{2.644976in}}%
\pgfpathlineto{\pgfqpoint{0.939741in}{2.726886in}}%
\pgfpathlineto{\pgfqpoint{0.939829in}{2.619675in}}%
\pgfpathlineto{\pgfqpoint{0.940465in}{2.721729in}}%
\pgfpathlineto{\pgfqpoint{0.940515in}{2.661774in}}%
\pgfpathlineto{\pgfqpoint{0.941158in}{2.765071in}}%
\pgfpathlineto{\pgfqpoint{0.941207in}{2.705422in}}%
\pgfpathlineto{\pgfqpoint{0.941711in}{2.799673in}}%
\pgfpathlineto{\pgfqpoint{0.941818in}{2.644559in}}%
\pgfpathlineto{\pgfqpoint{0.941932in}{2.644559in}}%
\pgfpathlineto{\pgfqpoint{0.944685in}{2.644803in}}%
\pgfpathlineto{\pgfqpoint{0.944686in}{2.719680in}}%
\pgfpathlineto{\pgfqpoint{0.944772in}{2.631445in}}%
\pgfpathlineto{\pgfqpoint{0.945408in}{2.709351in}}%
\pgfpathlineto{\pgfqpoint{0.945457in}{2.673324in}}%
\pgfpathlineto{\pgfqpoint{0.946100in}{2.752740in}}%
\pgfpathlineto{\pgfqpoint{0.946150in}{2.716925in}}%
\pgfpathlineto{\pgfqpoint{0.946656in}{2.787560in}}%
\pgfpathlineto{\pgfqpoint{0.946746in}{2.644559in}}%
\pgfpathlineto{\pgfqpoint{0.946833in}{2.644559in}}%
\pgfpathlineto{\pgfqpoint{0.949630in}{2.644816in}}%
\pgfpathlineto{\pgfqpoint{0.949632in}{2.720751in}}%
\pgfpathlineto{\pgfqpoint{0.949717in}{2.628589in}}%
\pgfpathlineto{\pgfqpoint{0.950353in}{2.711878in}}%
\pgfpathlineto{\pgfqpoint{0.950403in}{2.670478in}}%
\pgfpathlineto{\pgfqpoint{0.951046in}{2.755257in}}%
\pgfpathlineto{\pgfqpoint{0.951095in}{2.714089in}}%
\pgfpathlineto{\pgfqpoint{0.951602in}{2.790067in}}%
\pgfpathlineto{\pgfqpoint{0.951689in}{2.644559in}}%
\pgfpathlineto{\pgfqpoint{0.951794in}{2.644559in}}%
\pgfpathlineto{\pgfqpoint{0.954575in}{2.644625in}}%
\pgfpathlineto{\pgfqpoint{0.954577in}{2.726934in}}%
\pgfpathlineto{\pgfqpoint{0.954695in}{2.636743in}}%
\pgfpathlineto{\pgfqpoint{0.955338in}{2.706642in}}%
\pgfpathlineto{\pgfqpoint{0.955387in}{2.679992in}}%
\pgfpathlineto{\pgfqpoint{0.955931in}{2.743868in}}%
\pgfpathlineto{\pgfqpoint{0.955980in}{2.717330in}}%
\pgfpathlineto{\pgfqpoint{0.956511in}{2.780248in}}%
\pgfpathlineto{\pgfqpoint{0.956610in}{2.701783in}}%
\pgfpathlineto{\pgfqpoint{0.956626in}{2.723168in}}%
\pgfpathlineto{\pgfqpoint{0.957413in}{2.644559in}}%
\pgfpathlineto{\pgfqpoint{0.959518in}{2.644559in}}%
\pgfpathlineto{\pgfqpoint{0.959523in}{2.752626in}}%
\pgfpathlineto{\pgfqpoint{0.960382in}{2.703176in}}%
\pgfpathlineto{\pgfqpoint{0.960431in}{2.695372in}}%
\pgfpathlineto{\pgfqpoint{0.960876in}{2.734217in}}%
\pgfpathlineto{\pgfqpoint{0.960926in}{2.726467in}}%
\pgfpathlineto{\pgfqpoint{0.961456in}{2.770632in}}%
\pgfpathlineto{\pgfqpoint{0.961555in}{2.694132in}}%
\pgfpathlineto{\pgfqpoint{0.961572in}{2.732802in}}%
\pgfpathlineto{\pgfqpoint{0.962357in}{2.644559in}}%
\pgfpathlineto{\pgfqpoint{0.964464in}{2.644559in}}%
\pgfpathlineto{\pgfqpoint{0.964469in}{2.743841in}}%
\pgfpathlineto{\pgfqpoint{0.964551in}{2.616315in}}%
\pgfpathlineto{\pgfqpoint{0.965236in}{2.657993in}}%
\pgfpathlineto{\pgfqpoint{0.966076in}{2.778143in}}%
\pgfpathlineto{\pgfqpoint{0.966438in}{2.800756in}}%
\pgfpathlineto{\pgfqpoint{0.966948in}{2.644559in}}%
\pgfpathlineto{\pgfqpoint{0.969409in}{2.644559in}}%
\pgfpathlineto{\pgfqpoint{0.969419in}{2.754688in}}%
\pgfpathlineto{\pgfqpoint{0.969476in}{2.626341in}}%
\pgfpathlineto{\pgfqpoint{0.970160in}{2.665611in}}%
\pgfpathlineto{\pgfqpoint{0.971000in}{2.767562in}}%
\pgfpathlineto{\pgfqpoint{0.971384in}{2.791589in}}%
\pgfpathlineto{\pgfqpoint{0.971901in}{2.644559in}}%
\pgfpathlineto{\pgfqpoint{0.974358in}{2.645002in}}%
\pgfpathlineto{\pgfqpoint{0.974360in}{2.759140in}}%
\pgfpathlineto{\pgfqpoint{0.974423in}{2.626336in}}%
\pgfpathlineto{\pgfqpoint{0.975108in}{2.666190in}}%
\pgfpathlineto{\pgfqpoint{0.975949in}{2.767008in}}%
\pgfpathlineto{\pgfqpoint{0.976329in}{2.790792in}}%
\pgfpathlineto{\pgfqpoint{0.976853in}{2.644559in}}%
\pgfpathlineto{\pgfqpoint{0.979301in}{2.644559in}}%
\pgfpathlineto{\pgfqpoint{0.979309in}{2.721863in}}%
\pgfpathlineto{\pgfqpoint{0.980059in}{2.691129in}}%
\pgfpathlineto{\pgfqpoint{0.980504in}{2.717064in}}%
\pgfpathlineto{\pgfqpoint{0.981274in}{2.765417in}}%
\pgfpathlineto{\pgfqpoint{0.981304in}{2.744834in}}%
\pgfpathlineto{\pgfqpoint{0.981351in}{2.747413in}}%
\pgfpathlineto{\pgfqpoint{0.982104in}{2.644559in}}%
\pgfpathlineto{\pgfqpoint{0.984246in}{2.644559in}}%
\pgfpathlineto{\pgfqpoint{0.984255in}{2.723370in}}%
\pgfpathlineto{\pgfqpoint{0.984994in}{2.690367in}}%
\pgfpathlineto{\pgfqpoint{0.985637in}{2.729057in}}%
\pgfpathlineto{\pgfqpoint{0.986220in}{2.765666in}}%
\pgfpathlineto{\pgfqpoint{0.986258in}{2.714429in}}%
\pgfpathlineto{\pgfqpoint{0.986270in}{2.741854in}}%
\pgfpathlineto{\pgfqpoint{0.986297in}{2.747502in}}%
\pgfpathlineto{\pgfqpoint{0.987062in}{2.644559in}}%
\pgfpathlineto{\pgfqpoint{0.989191in}{2.644559in}}%
\pgfpathlineto{\pgfqpoint{0.989196in}{2.751520in}}%
\pgfpathlineto{\pgfqpoint{0.989321in}{2.641825in}}%
\pgfpathlineto{\pgfqpoint{0.989964in}{2.699182in}}%
\pgfpathlineto{\pgfqpoint{0.990014in}{2.685151in}}%
\pgfpathlineto{\pgfqpoint{0.990558in}{2.736421in}}%
\pgfpathlineto{\pgfqpoint{0.990607in}{2.722475in}}%
\pgfpathlineto{\pgfqpoint{0.991133in}{2.772552in}}%
\pgfpathlineto{\pgfqpoint{0.991219in}{2.713896in}}%
\pgfpathlineto{\pgfqpoint{0.991233in}{2.720083in}}%
\pgfpathlineto{\pgfqpoint{0.992021in}{2.644559in}}%
\pgfpathlineto{\pgfqpoint{0.994139in}{2.644615in}}%
\pgfpathlineto{\pgfqpoint{0.994142in}{2.735518in}}%
\pgfpathlineto{\pgfqpoint{0.994267in}{2.627841in}}%
\pgfpathlineto{\pgfqpoint{0.994910in}{2.713445in}}%
\pgfpathlineto{\pgfqpoint{0.994959in}{2.671215in}}%
\pgfpathlineto{\pgfqpoint{0.995602in}{2.756843in}}%
\pgfpathlineto{\pgfqpoint{0.995652in}{2.714808in}}%
\pgfpathlineto{\pgfqpoint{0.996079in}{2.786718in}}%
\pgfpathlineto{\pgfqpoint{0.996204in}{2.644559in}}%
\pgfpathlineto{\pgfqpoint{0.996351in}{2.644559in}}%
\pgfpathlineto{\pgfqpoint{0.999082in}{2.644559in}}%
\pgfpathlineto{\pgfqpoint{0.999088in}{2.723947in}}%
\pgfpathlineto{\pgfqpoint{0.999851in}{2.687535in}}%
\pgfpathlineto{\pgfqpoint{1.001056in}{2.768172in}}%
\pgfpathlineto{\pgfqpoint{1.001056in}{2.762480in}}%
\pgfpathlineto{\pgfqpoint{1.001828in}{2.644559in}}%
\pgfpathlineto{\pgfqpoint{1.004028in}{2.644559in}}%
\pgfpathlineto{\pgfqpoint{1.004037in}{2.721259in}}%
\pgfpathlineto{\pgfqpoint{1.004784in}{2.689886in}}%
\pgfpathlineto{\pgfqpoint{1.006002in}{2.765148in}}%
\pgfpathlineto{\pgfqpoint{1.006030in}{2.751922in}}%
\pgfpathlineto{\pgfqpoint{1.006819in}{2.644559in}}%
\pgfpathlineto{\pgfqpoint{1.008973in}{2.644559in}}%
\pgfpathlineto{\pgfqpoint{1.008979in}{2.726937in}}%
\pgfpathlineto{\pgfqpoint{1.009098in}{2.641178in}}%
\pgfpathlineto{\pgfqpoint{1.009741in}{2.697905in}}%
\pgfpathlineto{\pgfqpoint{1.009791in}{2.684437in}}%
\pgfpathlineto{\pgfqpoint{1.010335in}{2.735145in}}%
\pgfpathlineto{\pgfqpoint{1.010384in}{2.721760in}}%
\pgfpathlineto{\pgfqpoint{1.010913in}{2.771423in}}%
\pgfpathlineto{\pgfqpoint{1.011001in}{2.713254in}}%
\pgfpathlineto{\pgfqpoint{1.011015in}{2.719331in}}%
\pgfpathlineto{\pgfqpoint{1.011803in}{2.644559in}}%
\pgfpathlineto{\pgfqpoint{1.013919in}{2.644559in}}%
\pgfpathlineto{\pgfqpoint{1.013925in}{2.726339in}}%
\pgfpathlineto{\pgfqpoint{1.014038in}{2.641918in}}%
\pgfpathlineto{\pgfqpoint{1.014681in}{2.696487in}}%
\pgfpathlineto{\pgfqpoint{1.014730in}{2.685090in}}%
\pgfpathlineto{\pgfqpoint{1.015274in}{2.733732in}}%
\pgfpathlineto{\pgfqpoint{1.015324in}{2.722409in}}%
\pgfpathlineto{\pgfqpoint{1.015855in}{2.770197in}}%
\pgfpathlineto{\pgfqpoint{1.015946in}{2.715603in}}%
\pgfpathlineto{\pgfqpoint{1.015960in}{2.716715in}}%
\pgfpathlineto{\pgfqpoint{1.016749in}{2.644559in}}%
\pgfpathlineto{\pgfqpoint{1.018864in}{2.644559in}}%
\pgfpathlineto{\pgfqpoint{1.018870in}{2.729729in}}%
\pgfpathlineto{\pgfqpoint{1.018925in}{2.636875in}}%
\pgfpathlineto{\pgfqpoint{1.019603in}{2.674919in}}%
\pgfpathlineto{\pgfqpoint{1.020444in}{2.752726in}}%
\pgfpathlineto{\pgfqpoint{1.020838in}{2.777439in}}%
\pgfpathlineto{\pgfqpoint{1.021459in}{2.644559in}}%
\pgfpathlineto{\pgfqpoint{1.023810in}{2.644559in}}%
\pgfpathlineto{\pgfqpoint{1.023816in}{2.731140in}}%
\pgfpathlineto{\pgfqpoint{1.023876in}{2.634863in}}%
\pgfpathlineto{\pgfqpoint{1.024558in}{2.674097in}}%
\pgfpathlineto{\pgfqpoint{1.025399in}{2.754645in}}%
\pgfpathlineto{\pgfqpoint{1.025784in}{2.778751in}}%
\pgfpathlineto{\pgfqpoint{1.026392in}{2.644559in}}%
\pgfpathlineto{\pgfqpoint{1.028755in}{2.644559in}}%
\pgfpathlineto{\pgfqpoint{1.028761in}{2.734737in}}%
\pgfpathlineto{\pgfqpoint{1.028845in}{2.633833in}}%
\pgfpathlineto{\pgfqpoint{1.029482in}{2.700262in}}%
\pgfpathlineto{\pgfqpoint{1.029531in}{2.675606in}}%
\pgfpathlineto{\pgfqpoint{1.030075in}{2.737470in}}%
\pgfpathlineto{\pgfqpoint{1.030125in}{2.712961in}}%
\pgfpathlineto{\pgfqpoint{1.030729in}{2.778437in}}%
\pgfpathlineto{\pgfqpoint{1.030788in}{2.701422in}}%
\pgfpathlineto{\pgfqpoint{1.030804in}{2.733576in}}%
\pgfpathlineto{\pgfqpoint{1.031590in}{2.644559in}}%
\pgfpathlineto{\pgfqpoint{1.033701in}{2.644559in}}%
\pgfpathlineto{\pgfqpoint{1.033707in}{2.735898in}}%
\pgfpathlineto{\pgfqpoint{1.033793in}{2.631089in}}%
\pgfpathlineto{\pgfqpoint{1.034430in}{2.702817in}}%
\pgfpathlineto{\pgfqpoint{1.034479in}{2.673029in}}%
\pgfpathlineto{\pgfqpoint{1.035122in}{2.746214in}}%
\pgfpathlineto{\pgfqpoint{1.035172in}{2.716623in}}%
\pgfpathlineto{\pgfqpoint{1.035674in}{2.780805in}}%
\pgfpathlineto{\pgfqpoint{1.035768in}{2.644559in}}%
\pgfpathlineto{\pgfqpoint{1.035842in}{2.644559in}}%
\pgfpathlineto{\pgfqpoint{1.038646in}{2.644559in}}%
\pgfpathlineto{\pgfqpoint{1.038652in}{2.739644in}}%
\pgfpathlineto{\pgfqpoint{1.038715in}{2.612965in}}%
\pgfpathlineto{\pgfqpoint{1.039398in}{2.652590in}}%
\pgfpathlineto{\pgfqpoint{1.040239in}{2.775542in}}%
\pgfpathlineto{\pgfqpoint{1.040620in}{2.799394in}}%
\pgfpathlineto{\pgfqpoint{1.041105in}{2.644559in}}%
\pgfpathlineto{\pgfqpoint{1.043592in}{2.644559in}}%
\pgfpathlineto{\pgfqpoint{1.043598in}{2.741222in}}%
\pgfpathlineto{\pgfqpoint{1.043655in}{2.614816in}}%
\pgfpathlineto{\pgfqpoint{1.044336in}{2.653458in}}%
\pgfpathlineto{\pgfqpoint{1.045176in}{2.772913in}}%
\pgfpathlineto{\pgfqpoint{1.045565in}{2.797264in}}%
\pgfpathlineto{\pgfqpoint{1.046048in}{2.644559in}}%
\pgfpathlineto{\pgfqpoint{1.048537in}{2.644559in}}%
\pgfpathlineto{\pgfqpoint{1.048543in}{2.741973in}}%
\pgfpathlineto{\pgfqpoint{1.048597in}{2.618159in}}%
\pgfpathlineto{\pgfqpoint{1.049274in}{2.655929in}}%
\pgfpathlineto{\pgfqpoint{1.050115in}{2.769461in}}%
\pgfpathlineto{\pgfqpoint{1.050486in}{2.792687in}}%
\pgfpathlineto{\pgfqpoint{1.050992in}{2.644559in}}%
\pgfpathlineto{\pgfqpoint{1.053482in}{2.644559in}}%
\pgfpathlineto{\pgfqpoint{1.053489in}{2.744162in}}%
\pgfpathlineto{\pgfqpoint{1.053606in}{2.625586in}}%
\pgfpathlineto{\pgfqpoint{1.054249in}{2.710240in}}%
\pgfpathlineto{\pgfqpoint{1.054299in}{2.668849in}}%
\pgfpathlineto{\pgfqpoint{1.054942in}{2.753645in}}%
\pgfpathlineto{\pgfqpoint{1.054991in}{2.712434in}}%
\pgfpathlineto{\pgfqpoint{1.055422in}{2.783718in}}%
\pgfpathlineto{\pgfqpoint{1.055549in}{2.644559in}}%
\pgfpathlineto{\pgfqpoint{1.055691in}{2.644559in}}%
\pgfpathlineto{\pgfqpoint{1.058428in}{2.644559in}}%
\pgfpathlineto{\pgfqpoint{1.058434in}{2.744408in}}%
\pgfpathlineto{\pgfqpoint{1.058552in}{2.625556in}}%
\pgfpathlineto{\pgfqpoint{1.059195in}{2.710120in}}%
\pgfpathlineto{\pgfqpoint{1.059244in}{2.668820in}}%
\pgfpathlineto{\pgfqpoint{1.059887in}{2.753525in}}%
\pgfpathlineto{\pgfqpoint{1.059937in}{2.712405in}}%
\pgfpathlineto{\pgfqpoint{1.060367in}{2.783594in}}%
\pgfpathlineto{\pgfqpoint{1.060494in}{2.644559in}}%
\pgfpathlineto{\pgfqpoint{1.060636in}{2.644559in}}%
\pgfpathlineto{\pgfqpoint{1.063373in}{2.644559in}}%
\pgfpathlineto{\pgfqpoint{1.063380in}{2.745418in}}%
\pgfpathlineto{\pgfqpoint{1.063495in}{2.625587in}}%
\pgfpathlineto{\pgfqpoint{1.064138in}{2.709518in}}%
\pgfpathlineto{\pgfqpoint{1.064187in}{2.668807in}}%
\pgfpathlineto{\pgfqpoint{1.064830in}{2.752925in}}%
\pgfpathlineto{\pgfqpoint{1.064880in}{2.712391in}}%
\pgfpathlineto{\pgfqpoint{1.065311in}{2.783074in}}%
\pgfpathlineto{\pgfqpoint{1.065439in}{2.644559in}}%
\pgfpathlineto{\pgfqpoint{1.065577in}{2.644559in}}%
\pgfpathlineto{\pgfqpoint{1.068319in}{2.644559in}}%
\pgfpathlineto{\pgfqpoint{1.068325in}{2.746184in}}%
\pgfpathlineto{\pgfqpoint{1.068452in}{2.619886in}}%
\pgfpathlineto{\pgfqpoint{1.069095in}{2.716273in}}%
\pgfpathlineto{\pgfqpoint{1.069144in}{2.663275in}}%
\pgfpathlineto{\pgfqpoint{1.069787in}{2.759654in}}%
\pgfpathlineto{\pgfqpoint{1.069837in}{2.706884in}}%
\pgfpathlineto{\pgfqpoint{1.070262in}{2.789423in}}%
\pgfpathlineto{\pgfqpoint{1.070378in}{2.644559in}}%
\pgfpathlineto{\pgfqpoint{1.070555in}{2.644559in}}%
\pgfpathlineto{\pgfqpoint{1.073264in}{2.644559in}}%
\pgfpathlineto{\pgfqpoint{1.073271in}{2.746362in}}%
\pgfpathlineto{\pgfqpoint{1.073401in}{2.618275in}}%
\pgfpathlineto{\pgfqpoint{1.073994in}{2.655471in}}%
\pgfpathlineto{\pgfqpoint{1.074835in}{2.767457in}}%
\pgfpathlineto{\pgfqpoint{1.075209in}{2.790922in}}%
\pgfpathlineto{\pgfqpoint{1.075713in}{2.644559in}}%
\pgfpathlineto{\pgfqpoint{1.078210in}{2.644559in}}%
\pgfpathlineto{\pgfqpoint{1.078216in}{2.746886in}}%
\pgfpathlineto{\pgfqpoint{1.078340in}{2.615435in}}%
\pgfpathlineto{\pgfqpoint{1.078983in}{2.719840in}}%
\pgfpathlineto{\pgfqpoint{1.079032in}{2.658798in}}%
\pgfpathlineto{\pgfqpoint{1.079675in}{2.763208in}}%
\pgfpathlineto{\pgfqpoint{1.079725in}{2.702421in}}%
\pgfpathlineto{\pgfqpoint{1.080152in}{2.793059in}}%
\pgfpathlineto{\pgfqpoint{1.080277in}{2.644559in}}%
\pgfpathlineto{\pgfqpoint{1.080452in}{2.644559in}}%
\pgfpathlineto{\pgfqpoint{1.083155in}{2.644559in}}%
\pgfpathlineto{\pgfqpoint{1.083162in}{2.747031in}}%
\pgfpathlineto{\pgfqpoint{1.083275in}{2.616376in}}%
\pgfpathlineto{\pgfqpoint{1.083918in}{2.716719in}}%
\pgfpathlineto{\pgfqpoint{1.083967in}{2.659568in}}%
\pgfpathlineto{\pgfqpoint{1.084610in}{2.760101in}}%
\pgfpathlineto{\pgfqpoint{1.084660in}{2.703177in}}%
\pgfpathlineto{\pgfqpoint{1.085092in}{2.790289in}}%
\pgfpathlineto{\pgfqpoint{1.085215in}{2.644559in}}%
\pgfpathlineto{\pgfqpoint{1.085385in}{2.644559in}}%
\pgfpathlineto{\pgfqpoint{1.088101in}{2.644559in}}%
\pgfpathlineto{\pgfqpoint{1.088107in}{2.746895in}}%
\pgfpathlineto{\pgfqpoint{1.088235in}{2.615582in}}%
\pgfpathlineto{\pgfqpoint{1.088878in}{2.719959in}}%
\pgfpathlineto{\pgfqpoint{1.088928in}{2.658995in}}%
\pgfpathlineto{\pgfqpoint{1.089570in}{2.763325in}}%
\pgfpathlineto{\pgfqpoint{1.089620in}{2.702619in}}%
\pgfpathlineto{\pgfqpoint{1.090045in}{2.793036in}}%
\pgfpathlineto{\pgfqpoint{1.090168in}{2.644559in}}%
\pgfpathlineto{\pgfqpoint{1.090345in}{2.644559in}}%
\pgfpathlineto{\pgfqpoint{1.093046in}{2.644559in}}%
\pgfpathlineto{\pgfqpoint{1.093053in}{2.746439in}}%
\pgfpathlineto{\pgfqpoint{1.093133in}{2.643332in}}%
\pgfpathlineto{\pgfqpoint{1.093819in}{2.684534in}}%
\pgfpathlineto{\pgfqpoint{1.095020in}{2.765416in}}%
\pgfpathlineto{\pgfqpoint{1.095020in}{2.758936in}}%
\pgfpathlineto{\pgfqpoint{1.095791in}{2.644559in}}%
\pgfpathlineto{\pgfqpoint{1.097992in}{2.644559in}}%
\pgfpathlineto{\pgfqpoint{1.097998in}{2.746777in}}%
\pgfpathlineto{\pgfqpoint{1.098110in}{2.628697in}}%
\pgfpathlineto{\pgfqpoint{1.098752in}{2.704415in}}%
\pgfpathlineto{\pgfqpoint{1.098802in}{2.671802in}}%
\pgfpathlineto{\pgfqpoint{1.099445in}{2.747855in}}%
\pgfpathlineto{\pgfqpoint{1.099494in}{2.715352in}}%
\pgfpathlineto{\pgfqpoint{1.099928in}{2.778142in}}%
\pgfpathlineto{\pgfqpoint{1.100051in}{2.644559in}}%
\pgfpathlineto{\pgfqpoint{1.100173in}{2.644559in}}%
\pgfpathlineto{\pgfqpoint{1.102937in}{2.644559in}}%
\pgfpathlineto{\pgfqpoint{1.102944in}{2.746244in}}%
\pgfpathlineto{\pgfqpoint{1.103044in}{2.642912in}}%
\pgfpathlineto{\pgfqpoint{1.103681in}{2.688423in}}%
\pgfpathlineto{\pgfqpoint{1.103730in}{2.685299in}}%
\pgfpathlineto{\pgfqpoint{1.103978in}{2.707069in}}%
\pgfpathlineto{\pgfqpoint{1.104027in}{2.703935in}}%
\pgfpathlineto{\pgfqpoint{1.104865in}{2.762800in}}%
\pgfpathlineto{\pgfqpoint{1.104942in}{2.732411in}}%
\pgfpathlineto{\pgfqpoint{1.105768in}{2.644559in}}%
\pgfpathlineto{\pgfqpoint{1.107882in}{2.644559in}}%
\pgfpathlineto{\pgfqpoint{1.107890in}{2.745287in}}%
\pgfpathlineto{\pgfqpoint{1.108025in}{2.644152in}}%
\pgfpathlineto{\pgfqpoint{1.108619in}{2.681321in}}%
\pgfpathlineto{\pgfqpoint{1.109756in}{2.759616in}}%
\pgfpathlineto{\pgfqpoint{1.109919in}{2.697106in}}%
\pgfpathlineto{\pgfqpoint{1.109831in}{2.764295in}}%
\pgfpathlineto{\pgfqpoint{1.109935in}{2.698671in}}%
\pgfpathlineto{\pgfqpoint{1.110718in}{2.644559in}}%
\pgfpathlineto{\pgfqpoint{1.112828in}{2.644559in}}%
\pgfpathlineto{\pgfqpoint{1.112835in}{2.744898in}}%
\pgfpathlineto{\pgfqpoint{1.112902in}{2.636331in}}%
\pgfpathlineto{\pgfqpoint{1.113585in}{2.675935in}}%
\pgfpathlineto{\pgfqpoint{1.114426in}{2.749299in}}%
\pgfpathlineto{\pgfqpoint{1.114802in}{2.772861in}}%
\pgfpathlineto{\pgfqpoint{1.115496in}{2.644559in}}%
\pgfpathlineto{\pgfqpoint{1.117773in}{2.644559in}}%
\pgfpathlineto{\pgfqpoint{1.117781in}{2.744098in}}%
\pgfpathlineto{\pgfqpoint{1.117906in}{2.637413in}}%
\pgfpathlineto{\pgfqpoint{1.118549in}{2.696592in}}%
\pgfpathlineto{\pgfqpoint{1.118598in}{2.680694in}}%
\pgfpathlineto{\pgfqpoint{1.119142in}{2.733836in}}%
\pgfpathlineto{\pgfqpoint{1.119192in}{2.718014in}}%
\pgfpathlineto{\pgfqpoint{1.119717in}{2.769879in}}%
\pgfpathlineto{\pgfqpoint{1.119825in}{2.696353in}}%
\pgfpathlineto{\pgfqpoint{1.120607in}{2.644559in}}%
\pgfpathlineto{\pgfqpoint{1.122719in}{2.644559in}}%
\pgfpathlineto{\pgfqpoint{1.122726in}{2.743644in}}%
\pgfpathlineto{\pgfqpoint{1.122782in}{2.635934in}}%
\pgfpathlineto{\pgfqpoint{1.123456in}{2.673116in}}%
\pgfpathlineto{\pgfqpoint{1.124297in}{2.748771in}}%
\pgfpathlineto{\pgfqpoint{1.124693in}{2.773589in}}%
\pgfpathlineto{\pgfqpoint{1.125329in}{2.644559in}}%
\pgfpathlineto{\pgfqpoint{1.127664in}{2.644559in}}%
\pgfpathlineto{\pgfqpoint{1.127672in}{2.743403in}}%
\pgfpathlineto{\pgfqpoint{1.127734in}{2.632159in}}%
\pgfpathlineto{\pgfqpoint{1.128415in}{2.670913in}}%
\pgfpathlineto{\pgfqpoint{1.129256in}{2.752657in}}%
\pgfpathlineto{\pgfqpoint{1.129638in}{2.776645in}}%
\pgfpathlineto{\pgfqpoint{1.130232in}{2.644559in}}%
\pgfpathlineto{\pgfqpoint{1.132610in}{2.644559in}}%
\pgfpathlineto{\pgfqpoint{1.132617in}{2.742272in}}%
\pgfpathlineto{\pgfqpoint{1.132694in}{2.620861in}}%
\pgfpathlineto{\pgfqpoint{1.133379in}{2.661676in}}%
\pgfpathlineto{\pgfqpoint{1.134220in}{2.764158in}}%
\pgfpathlineto{\pgfqpoint{1.134584in}{2.786942in}}%
\pgfpathlineto{\pgfqpoint{1.135118in}{2.644559in}}%
\pgfpathlineto{\pgfqpoint{1.137555in}{2.644559in}}%
\pgfpathlineto{\pgfqpoint{1.137563in}{2.741337in}}%
\pgfpathlineto{\pgfqpoint{1.137629in}{2.612292in}}%
\pgfpathlineto{\pgfqpoint{1.138311in}{2.651719in}}%
\pgfpathlineto{\pgfqpoint{1.139152in}{2.772124in}}%
\pgfpathlineto{\pgfqpoint{1.139529in}{2.795714in}}%
\pgfpathlineto{\pgfqpoint{1.140018in}{2.644559in}}%
\pgfpathlineto{\pgfqpoint{1.142501in}{2.644559in}}%
\pgfpathlineto{\pgfqpoint{1.142508in}{2.740513in}}%
\pgfpathlineto{\pgfqpoint{1.142573in}{2.613761in}}%
\pgfpathlineto{\pgfqpoint{1.143255in}{2.652886in}}%
\pgfpathlineto{\pgfqpoint{1.144095in}{2.770646in}}%
\pgfpathlineto{\pgfqpoint{1.144474in}{2.794377in}}%
\pgfpathlineto{\pgfqpoint{1.144967in}{2.644559in}}%
\pgfpathlineto{\pgfqpoint{1.147446in}{2.644559in}}%
\pgfpathlineto{\pgfqpoint{1.147453in}{2.739531in}}%
\pgfpathlineto{\pgfqpoint{1.147513in}{2.616817in}}%
\pgfpathlineto{\pgfqpoint{1.148191in}{2.654745in}}%
\pgfpathlineto{\pgfqpoint{1.149032in}{2.767791in}}%
\pgfpathlineto{\pgfqpoint{1.149420in}{2.792106in}}%
\pgfpathlineto{\pgfqpoint{1.149910in}{2.644559in}}%
\pgfpathlineto{\pgfqpoint{1.152392in}{2.644559in}}%
\pgfpathlineto{\pgfqpoint{1.152404in}{2.772030in}}%
\pgfpathlineto{\pgfqpoint{1.153148in}{2.684058in}}%
\pgfpathlineto{\pgfqpoint{1.154365in}{2.763311in}}%
\pgfpathlineto{\pgfqpoint{1.154366in}{2.738092in}}%
\pgfpathlineto{\pgfqpoint{1.154392in}{2.748035in}}%
\pgfpathlineto{\pgfqpoint{1.155136in}{2.644559in}}%
\pgfpathlineto{\pgfqpoint{1.157337in}{2.644559in}}%
\pgfpathlineto{\pgfqpoint{1.157349in}{2.771627in}}%
\pgfpathlineto{\pgfqpoint{1.157420in}{2.630036in}}%
\pgfpathlineto{\pgfqpoint{1.158105in}{2.670418in}}%
\pgfpathlineto{\pgfqpoint{1.158945in}{2.754929in}}%
\pgfpathlineto{\pgfqpoint{1.159311in}{2.777835in}}%
\pgfpathlineto{\pgfqpoint{1.159908in}{2.644559in}}%
\pgfpathlineto{\pgfqpoint{1.162282in}{2.644559in}}%
\pgfpathlineto{\pgfqpoint{1.162295in}{2.771398in}}%
\pgfpathlineto{\pgfqpoint{1.162419in}{2.636985in}}%
\pgfpathlineto{\pgfqpoint{1.163013in}{2.674063in}}%
\pgfpathlineto{\pgfqpoint{1.163952in}{2.752903in}}%
\pgfpathlineto{\pgfqpoint{1.164228in}{2.770187in}}%
\pgfpathlineto{\pgfqpoint{1.165043in}{2.644559in}}%
\pgfpathlineto{\pgfqpoint{1.167228in}{2.644559in}}%
\pgfpathlineto{\pgfqpoint{1.167240in}{2.768612in}}%
\pgfpathlineto{\pgfqpoint{1.167316in}{2.630791in}}%
\pgfpathlineto{\pgfqpoint{1.168002in}{2.671752in}}%
\pgfpathlineto{\pgfqpoint{1.168842in}{2.755770in}}%
\pgfpathlineto{\pgfqpoint{1.169202in}{2.778290in}}%
\pgfpathlineto{\pgfqpoint{1.169808in}{2.644559in}}%
\pgfpathlineto{\pgfqpoint{1.172173in}{2.644559in}}%
\pgfpathlineto{\pgfqpoint{1.172185in}{2.766766in}}%
\pgfpathlineto{\pgfqpoint{1.172200in}{2.640272in}}%
\pgfpathlineto{\pgfqpoint{1.172945in}{2.681806in}}%
\pgfpathlineto{\pgfqpoint{1.173984in}{2.757467in}}%
\pgfpathlineto{\pgfqpoint{1.174147in}{2.767749in}}%
\pgfpathlineto{\pgfqpoint{1.175514in}{2.644559in}}%
\pgfpathlineto{\pgfqpoint{1.177119in}{2.644559in}}%
\pgfpathlineto{\pgfqpoint{1.177130in}{2.766094in}}%
\pgfpathlineto{\pgfqpoint{1.177262in}{2.637345in}}%
\pgfpathlineto{\pgfqpoint{1.177856in}{2.674474in}}%
\pgfpathlineto{\pgfqpoint{1.178696in}{2.747958in}}%
\pgfpathlineto{\pgfqpoint{1.179068in}{2.771242in}}%
\pgfpathlineto{\pgfqpoint{1.179760in}{2.644559in}}%
\pgfpathlineto{\pgfqpoint{1.182064in}{2.644559in}}%
\pgfpathlineto{\pgfqpoint{1.182076in}{2.765054in}}%
\pgfpathlineto{\pgfqpoint{1.182151in}{2.619745in}}%
\pgfpathlineto{\pgfqpoint{1.182836in}{2.660465in}}%
\pgfpathlineto{\pgfqpoint{1.183677in}{2.766369in}}%
\pgfpathlineto{\pgfqpoint{1.184038in}{2.788968in}}%
\pgfpathlineto{\pgfqpoint{1.184570in}{2.644559in}}%
\pgfpathlineto{\pgfqpoint{1.187010in}{2.644559in}}%
\pgfpathlineto{\pgfqpoint{1.187021in}{2.763265in}}%
\pgfpathlineto{\pgfqpoint{1.187086in}{2.637581in}}%
\pgfpathlineto{\pgfqpoint{1.187769in}{2.676700in}}%
\pgfpathlineto{\pgfqpoint{1.188709in}{2.755223in}}%
\pgfpathlineto{\pgfqpoint{1.188984in}{2.772456in}}%
\pgfpathlineto{\pgfqpoint{1.189815in}{2.644559in}}%
\pgfpathlineto{\pgfqpoint{1.191955in}{2.644559in}}%
\pgfpathlineto{\pgfqpoint{1.191966in}{2.760363in}}%
\pgfpathlineto{\pgfqpoint{1.192713in}{2.687651in}}%
\pgfpathlineto{\pgfqpoint{1.193059in}{2.707339in}}%
\pgfpathlineto{\pgfqpoint{1.193929in}{2.762000in}}%
\pgfpathlineto{\pgfqpoint{1.193956in}{2.749346in}}%
\pgfpathlineto{\pgfqpoint{1.194748in}{2.644559in}}%
\pgfpathlineto{\pgfqpoint{1.196901in}{2.644559in}}%
\pgfpathlineto{\pgfqpoint{1.196912in}{2.757255in}}%
\pgfpathlineto{\pgfqpoint{1.197642in}{2.684730in}}%
\pgfpathlineto{\pgfqpoint{1.198874in}{2.764522in}}%
\pgfpathlineto{\pgfqpoint{1.198875in}{2.739424in}}%
\pgfpathlineto{\pgfqpoint{1.198897in}{2.747424in}}%
\pgfpathlineto{\pgfqpoint{1.199645in}{2.644559in}}%
\pgfpathlineto{\pgfqpoint{1.201846in}{2.644559in}}%
\pgfpathlineto{\pgfqpoint{1.201857in}{2.753952in}}%
\pgfpathlineto{\pgfqpoint{1.201919in}{2.640046in}}%
\pgfpathlineto{\pgfqpoint{1.202597in}{2.678105in}}%
\pgfpathlineto{\pgfqpoint{1.203537in}{2.754476in}}%
\pgfpathlineto{\pgfqpoint{1.203820in}{2.772212in}}%
\pgfpathlineto{\pgfqpoint{1.204692in}{2.644559in}}%
\pgfpathlineto{\pgfqpoint{1.206794in}{2.644791in}}%
\pgfpathlineto{\pgfqpoint{1.206799in}{2.772076in}}%
\pgfpathlineto{\pgfqpoint{1.206863in}{2.623432in}}%
\pgfpathlineto{\pgfqpoint{1.207540in}{2.661192in}}%
\pgfpathlineto{\pgfqpoint{1.208381in}{2.765438in}}%
\pgfpathlineto{\pgfqpoint{1.208765in}{2.789524in}}%
\pgfpathlineto{\pgfqpoint{1.209273in}{2.644559in}}%
\pgfpathlineto{\pgfqpoint{1.211738in}{2.644560in}}%
\pgfpathlineto{\pgfqpoint{1.211754in}{2.756471in}}%
\pgfpathlineto{\pgfqpoint{1.211882in}{2.617207in}}%
\pgfpathlineto{\pgfqpoint{1.212476in}{2.654396in}}%
\pgfpathlineto{\pgfqpoint{1.213317in}{2.771886in}}%
\pgfpathlineto{\pgfqpoint{1.213711in}{2.796571in}}%
\pgfpathlineto{\pgfqpoint{1.214190in}{2.644559in}}%
\pgfpathlineto{\pgfqpoint{1.216684in}{2.644562in}}%
\pgfpathlineto{\pgfqpoint{1.216698in}{2.761830in}}%
\pgfpathlineto{\pgfqpoint{1.216822in}{2.643921in}}%
\pgfpathlineto{\pgfqpoint{1.217415in}{2.680944in}}%
\pgfpathlineto{\pgfqpoint{1.218355in}{2.753558in}}%
\pgfpathlineto{\pgfqpoint{1.218629in}{2.770786in}}%
\pgfpathlineto{\pgfqpoint{1.219639in}{2.644559in}}%
\pgfpathlineto{\pgfqpoint{1.221627in}{2.644559in}}%
\pgfpathlineto{\pgfqpoint{1.221635in}{2.777578in}}%
\pgfpathlineto{\pgfqpoint{1.222504in}{2.701268in}}%
\pgfpathlineto{\pgfqpoint{1.222553in}{2.696597in}}%
\pgfpathlineto{\pgfqpoint{1.222899in}{2.726131in}}%
\pgfpathlineto{\pgfqpoint{1.222949in}{2.721442in}}%
\pgfpathlineto{\pgfqpoint{1.223572in}{2.768400in}}%
\pgfpathlineto{\pgfqpoint{1.223641in}{2.719143in}}%
\pgfpathlineto{\pgfqpoint{1.223653in}{2.734442in}}%
\pgfpathlineto{\pgfqpoint{1.224452in}{2.644559in}}%
\pgfpathlineto{\pgfqpoint{1.226573in}{2.644559in}}%
\pgfpathlineto{\pgfqpoint{1.226583in}{2.750025in}}%
\pgfpathlineto{\pgfqpoint{1.226722in}{2.619946in}}%
\pgfpathlineto{\pgfqpoint{1.227316in}{2.657176in}}%
\pgfpathlineto{\pgfqpoint{1.228156in}{2.775404in}}%
\pgfpathlineto{\pgfqpoint{1.228547in}{2.799873in}}%
\pgfpathlineto{\pgfqpoint{1.229030in}{2.644559in}}%
\pgfpathlineto{\pgfqpoint{1.231519in}{2.644559in}}%
\pgfpathlineto{\pgfqpoint{1.231534in}{2.744407in}}%
\pgfpathlineto{\pgfqpoint{1.231590in}{2.621184in}}%
\pgfpathlineto{\pgfqpoint{1.232268in}{2.659008in}}%
\pgfpathlineto{\pgfqpoint{1.233108in}{2.776591in}}%
\pgfpathlineto{\pgfqpoint{1.233493in}{2.800652in}}%
\pgfpathlineto{\pgfqpoint{1.233984in}{2.644559in}}%
\pgfpathlineto{\pgfqpoint{1.236464in}{2.644559in}}%
\pgfpathlineto{\pgfqpoint{1.236478in}{2.748284in}}%
\pgfpathlineto{\pgfqpoint{1.236539in}{2.621126in}}%
\pgfpathlineto{\pgfqpoint{1.237219in}{2.659819in}}%
\pgfpathlineto{\pgfqpoint{1.238060in}{2.779437in}}%
\pgfpathlineto{\pgfqpoint{1.238438in}{2.803110in}}%
\pgfpathlineto{\pgfqpoint{1.238936in}{2.644559in}}%
\pgfpathlineto{\pgfqpoint{1.241410in}{2.644559in}}%
\pgfpathlineto{\pgfqpoint{1.241424in}{2.757030in}}%
\pgfpathlineto{\pgfqpoint{1.241507in}{2.635871in}}%
\pgfpathlineto{\pgfqpoint{1.242143in}{2.712359in}}%
\pgfpathlineto{\pgfqpoint{1.242192in}{2.677375in}}%
\pgfpathlineto{\pgfqpoint{1.242835in}{2.755744in}}%
\pgfpathlineto{\pgfqpoint{1.242885in}{2.720980in}}%
\pgfpathlineto{\pgfqpoint{1.243384in}{2.790087in}}%
\pgfpathlineto{\pgfqpoint{1.243471in}{2.644559in}}%
\pgfpathlineto{\pgfqpoint{1.243563in}{2.644559in}}%
\pgfpathlineto{\pgfqpoint{1.246355in}{2.644559in}}%
\pgfpathlineto{\pgfqpoint{1.246369in}{2.758562in}}%
\pgfpathlineto{\pgfqpoint{1.246456in}{2.640411in}}%
\pgfpathlineto{\pgfqpoint{1.247094in}{2.710031in}}%
\pgfpathlineto{\pgfqpoint{1.247143in}{2.682323in}}%
\pgfpathlineto{\pgfqpoint{1.247786in}{2.753468in}}%
\pgfpathlineto{\pgfqpoint{1.247836in}{2.725875in}}%
\pgfpathlineto{\pgfqpoint{1.248281in}{2.784493in}}%
\pgfpathlineto{\pgfqpoint{1.248415in}{2.644559in}}%
\pgfpathlineto{\pgfqpoint{1.248488in}{2.644559in}}%
\pgfpathlineto{\pgfqpoint{1.251301in}{2.644559in}}%
\pgfpathlineto{\pgfqpoint{1.251305in}{2.758363in}}%
\pgfpathlineto{\pgfqpoint{1.251391in}{2.622398in}}%
\pgfpathlineto{\pgfqpoint{1.252027in}{2.728481in}}%
\pgfpathlineto{\pgfqpoint{1.252076in}{2.663651in}}%
\pgfpathlineto{\pgfqpoint{1.252719in}{2.771829in}}%
\pgfpathlineto{\pgfqpoint{1.252768in}{2.707293in}}%
\pgfpathlineto{\pgfqpoint{1.253274in}{2.806578in}}%
\pgfpathlineto{\pgfqpoint{1.253360in}{2.644559in}}%
\pgfpathlineto{\pgfqpoint{1.253495in}{2.644559in}}%
\pgfpathlineto{\pgfqpoint{1.256246in}{2.644559in}}%
\pgfpathlineto{\pgfqpoint{1.256251in}{2.760964in}}%
\pgfpathlineto{\pgfqpoint{1.256335in}{2.637443in}}%
\pgfpathlineto{\pgfqpoint{1.256971in}{2.714216in}}%
\pgfpathlineto{\pgfqpoint{1.257020in}{2.678600in}}%
\pgfpathlineto{\pgfqpoint{1.257663in}{2.757603in}}%
\pgfpathlineto{\pgfqpoint{1.257713in}{2.722202in}}%
\pgfpathlineto{\pgfqpoint{1.258220in}{2.792476in}}%
\pgfpathlineto{\pgfqpoint{1.258314in}{2.644559in}}%
\pgfpathlineto{\pgfqpoint{1.258392in}{2.644559in}}%
\pgfpathlineto{\pgfqpoint{1.261192in}{2.644559in}}%
\pgfpathlineto{\pgfqpoint{1.261196in}{2.758738in}}%
\pgfpathlineto{\pgfqpoint{1.261305in}{2.634416in}}%
\pgfpathlineto{\pgfqpoint{1.261944in}{2.718565in}}%
\pgfpathlineto{\pgfqpoint{1.261993in}{2.677142in}}%
\pgfpathlineto{\pgfqpoint{1.262636in}{2.761975in}}%
\pgfpathlineto{\pgfqpoint{1.262685in}{2.720721in}}%
\pgfpathlineto{\pgfqpoint{1.263123in}{2.792520in}}%
\pgfpathlineto{\pgfqpoint{1.263260in}{2.644559in}}%
\pgfpathlineto{\pgfqpoint{1.263375in}{2.644559in}}%
\pgfpathlineto{\pgfqpoint{1.266137in}{2.644559in}}%
\pgfpathlineto{\pgfqpoint{1.266141in}{2.761025in}}%
\pgfpathlineto{\pgfqpoint{1.266271in}{2.626544in}}%
\pgfpathlineto{\pgfqpoint{1.266864in}{2.663676in}}%
\pgfpathlineto{\pgfqpoint{1.267705in}{2.779125in}}%
\pgfpathlineto{\pgfqpoint{1.268081in}{2.802664in}}%
\pgfpathlineto{\pgfqpoint{1.268586in}{2.644559in}}%
\pgfpathlineto{\pgfqpoint{1.271082in}{2.644559in}}%
\pgfpathlineto{\pgfqpoint{1.271087in}{2.762168in}}%
\pgfpathlineto{\pgfqpoint{1.271825in}{2.706919in}}%
\pgfpathlineto{\pgfqpoint{1.271875in}{2.688602in}}%
\pgfpathlineto{\pgfqpoint{1.272419in}{2.744172in}}%
\pgfpathlineto{\pgfqpoint{1.272468in}{2.725912in}}%
\pgfpathlineto{\pgfqpoint{1.273010in}{2.781249in}}%
\pgfpathlineto{\pgfqpoint{1.273118in}{2.698471in}}%
\pgfpathlineto{\pgfqpoint{1.273134in}{2.718830in}}%
\pgfpathlineto{\pgfqpoint{1.273919in}{2.644559in}}%
\pgfpathlineto{\pgfqpoint{1.276028in}{2.644559in}}%
\pgfpathlineto{\pgfqpoint{1.276032in}{2.761054in}}%
\pgfpathlineto{\pgfqpoint{1.276167in}{2.632831in}}%
\pgfpathlineto{\pgfqpoint{1.276761in}{2.670008in}}%
\pgfpathlineto{\pgfqpoint{1.277601in}{2.773556in}}%
\pgfpathlineto{\pgfqpoint{1.277975in}{2.796933in}}%
\pgfpathlineto{\pgfqpoint{1.278504in}{2.644559in}}%
\pgfpathlineto{\pgfqpoint{1.280973in}{2.644559in}}%
\pgfpathlineto{\pgfqpoint{1.280978in}{2.756650in}}%
\pgfpathlineto{\pgfqpoint{1.281037in}{2.630147in}}%
\pgfpathlineto{\pgfqpoint{1.281712in}{2.667491in}}%
\pgfpathlineto{\pgfqpoint{1.282553in}{2.775553in}}%
\pgfpathlineto{\pgfqpoint{1.282947in}{2.800261in}}%
\pgfpathlineto{\pgfqpoint{1.283447in}{2.644559in}}%
\pgfpathlineto{\pgfqpoint{1.285919in}{2.644559in}}%
\pgfpathlineto{\pgfqpoint{1.285923in}{2.762335in}}%
\pgfpathlineto{\pgfqpoint{1.286016in}{2.638945in}}%
\pgfpathlineto{\pgfqpoint{1.286653in}{2.713695in}}%
\pgfpathlineto{\pgfqpoint{1.286703in}{2.680813in}}%
\pgfpathlineto{\pgfqpoint{1.287346in}{2.757124in}}%
\pgfpathlineto{\pgfqpoint{1.287395in}{2.724373in}}%
\pgfpathlineto{\pgfqpoint{1.287893in}{2.791417in}}%
\pgfpathlineto{\pgfqpoint{1.287981in}{2.644559in}}%
\pgfpathlineto{\pgfqpoint{1.288063in}{2.644559in}}%
\pgfpathlineto{\pgfqpoint{1.290864in}{2.644559in}}%
\pgfpathlineto{\pgfqpoint{1.290869in}{2.762692in}}%
\pgfpathlineto{\pgfqpoint{1.291001in}{2.633932in}}%
\pgfpathlineto{\pgfqpoint{1.291595in}{2.671081in}}%
\pgfpathlineto{\pgfqpoint{1.292435in}{2.772625in}}%
\pgfpathlineto{\pgfqpoint{1.292810in}{2.796086in}}%
\pgfpathlineto{\pgfqpoint{1.293342in}{2.644559in}}%
\pgfpathlineto{\pgfqpoint{1.295810in}{2.644559in}}%
\pgfpathlineto{\pgfqpoint{1.295814in}{2.762265in}}%
\pgfpathlineto{\pgfqpoint{1.295892in}{2.624128in}}%
\pgfpathlineto{\pgfqpoint{1.296577in}{2.664882in}}%
\pgfpathlineto{\pgfqpoint{1.297418in}{2.783328in}}%
\pgfpathlineto{\pgfqpoint{1.297784in}{2.806222in}}%
\pgfpathlineto{\pgfqpoint{1.298296in}{2.644559in}}%
\pgfpathlineto{\pgfqpoint{1.300755in}{2.644559in}}%
\pgfpathlineto{\pgfqpoint{1.300760in}{2.758318in}}%
\pgfpathlineto{\pgfqpoint{1.300838in}{2.633827in}}%
\pgfpathlineto{\pgfqpoint{1.301523in}{2.674595in}}%
\pgfpathlineto{\pgfqpoint{1.302364in}{2.772600in}}%
\pgfpathlineto{\pgfqpoint{1.302729in}{2.795462in}}%
\pgfpathlineto{\pgfqpoint{1.303285in}{2.644559in}}%
\pgfpathlineto{\pgfqpoint{1.305701in}{2.644559in}}%
\pgfpathlineto{\pgfqpoint{1.305705in}{2.760444in}}%
\pgfpathlineto{\pgfqpoint{1.305844in}{2.631444in}}%
\pgfpathlineto{\pgfqpoint{1.306437in}{2.668660in}}%
\pgfpathlineto{\pgfqpoint{1.307278in}{2.775184in}}%
\pgfpathlineto{\pgfqpoint{1.307649in}{2.798442in}}%
\pgfpathlineto{\pgfqpoint{1.308176in}{2.644559in}}%
\pgfpathlineto{\pgfqpoint{1.310646in}{2.644559in}}%
\pgfpathlineto{\pgfqpoint{1.310651in}{2.757274in}}%
\pgfpathlineto{\pgfqpoint{1.310728in}{2.622036in}}%
\pgfpathlineto{\pgfqpoint{1.311413in}{2.662720in}}%
\pgfpathlineto{\pgfqpoint{1.312253in}{2.784009in}}%
\pgfpathlineto{\pgfqpoint{1.312620in}{2.806933in}}%
\pgfpathlineto{\pgfqpoint{1.313127in}{2.644559in}}%
\pgfpathlineto{\pgfqpoint{1.315590in}{2.644559in}}%
\pgfpathlineto{\pgfqpoint{1.315600in}{2.766788in}}%
\pgfpathlineto{\pgfqpoint{1.315670in}{2.639368in}}%
\pgfpathlineto{\pgfqpoint{1.316355in}{2.679650in}}%
\pgfpathlineto{\pgfqpoint{1.317196in}{2.768142in}}%
\pgfpathlineto{\pgfqpoint{1.317565in}{2.791284in}}%
\pgfpathlineto{\pgfqpoint{1.318155in}{2.644559in}}%
\pgfpathlineto{\pgfqpoint{1.320537in}{2.644559in}}%
\pgfpathlineto{\pgfqpoint{1.320541in}{2.762254in}}%
\pgfpathlineto{\pgfqpoint{1.320616in}{2.627023in}}%
\pgfpathlineto{\pgfqpoint{1.321301in}{2.667430in}}%
\pgfpathlineto{\pgfqpoint{1.322142in}{2.780393in}}%
\pgfpathlineto{\pgfqpoint{1.322511in}{2.803496in}}%
\pgfpathlineto{\pgfqpoint{1.323028in}{2.644559in}}%
\pgfpathlineto{\pgfqpoint{1.325482in}{2.644559in}}%
\pgfpathlineto{\pgfqpoint{1.325487in}{2.763192in}}%
\pgfpathlineto{\pgfqpoint{1.325564in}{2.634435in}}%
\pgfpathlineto{\pgfqpoint{1.326249in}{2.675096in}}%
\pgfpathlineto{\pgfqpoint{1.327090in}{2.773301in}}%
\pgfpathlineto{\pgfqpoint{1.327456in}{2.796248in}}%
\pgfpathlineto{\pgfqpoint{1.328011in}{2.644559in}}%
\pgfpathlineto{\pgfqpoint{1.330428in}{2.644559in}}%
\pgfpathlineto{\pgfqpoint{1.330432in}{2.761040in}}%
\pgfpathlineto{\pgfqpoint{1.330499in}{2.643134in}}%
\pgfpathlineto{\pgfqpoint{1.331182in}{2.682404in}}%
\pgfpathlineto{\pgfqpoint{1.332022in}{2.763852in}}%
\pgfpathlineto{\pgfqpoint{1.332402in}{2.787634in}}%
\pgfpathlineto{\pgfqpoint{1.333025in}{2.644559in}}%
\pgfpathlineto{\pgfqpoint{1.335373in}{2.644559in}}%
\pgfpathlineto{\pgfqpoint{1.335378in}{2.763222in}}%
\pgfpathlineto{\pgfqpoint{1.335455in}{2.634537in}}%
\pgfpathlineto{\pgfqpoint{1.336140in}{2.675236in}}%
\pgfpathlineto{\pgfqpoint{1.336981in}{2.773218in}}%
\pgfpathlineto{\pgfqpoint{1.337347in}{2.796141in}}%
\pgfpathlineto{\pgfqpoint{1.337902in}{2.644559in}}%
\pgfpathlineto{\pgfqpoint{1.340319in}{2.644559in}}%
\pgfpathlineto{\pgfqpoint{1.340323in}{2.760397in}}%
\pgfpathlineto{\pgfqpoint{1.340391in}{2.640953in}}%
\pgfpathlineto{\pgfqpoint{1.341074in}{2.680354in}}%
\pgfpathlineto{\pgfqpoint{1.341914in}{2.765846in}}%
\pgfpathlineto{\pgfqpoint{1.342293in}{2.789552in}}%
\pgfpathlineto{\pgfqpoint{1.342887in}{2.644559in}}%
\pgfpathlineto{\pgfqpoint{1.345264in}{2.644559in}}%
\pgfpathlineto{\pgfqpoint{1.345269in}{2.760895in}}%
\pgfpathlineto{\pgfqpoint{1.345343in}{2.636444in}}%
\pgfpathlineto{\pgfqpoint{1.346028in}{2.676773in}}%
\pgfpathlineto{\pgfqpoint{1.346869in}{2.770622in}}%
\pgfpathlineto{\pgfqpoint{1.347238in}{2.793761in}}%
\pgfpathlineto{\pgfqpoint{1.347805in}{2.644559in}}%
\pgfpathlineto{\pgfqpoint{1.350210in}{2.644559in}}%
\pgfpathlineto{\pgfqpoint{1.350214in}{2.762762in}}%
\pgfpathlineto{\pgfqpoint{1.350289in}{2.635719in}}%
\pgfpathlineto{\pgfqpoint{1.350974in}{2.676135in}}%
\pgfpathlineto{\pgfqpoint{1.351815in}{2.771856in}}%
\pgfpathlineto{\pgfqpoint{1.352184in}{2.794962in}}%
\pgfpathlineto{\pgfqpoint{1.352742in}{2.644559in}}%
\pgfpathlineto{\pgfqpoint{1.355155in}{2.644559in}}%
\pgfpathlineto{\pgfqpoint{1.355160in}{2.762969in}}%
\pgfpathlineto{\pgfqpoint{1.355882in}{2.687741in}}%
\pgfpathlineto{\pgfqpoint{1.356822in}{2.761960in}}%
\pgfpathlineto{\pgfqpoint{1.357099in}{2.779348in}}%
\pgfpathlineto{\pgfqpoint{1.358077in}{2.644559in}}%
\pgfpathlineto{\pgfqpoint{1.360101in}{2.644559in}}%
\pgfpathlineto{\pgfqpoint{1.360105in}{2.762932in}}%
\pgfpathlineto{\pgfqpoint{1.360851in}{2.697586in}}%
\pgfpathlineto{\pgfqpoint{1.361692in}{2.748934in}}%
\pgfpathlineto{\pgfqpoint{1.362074in}{2.772975in}}%
\pgfpathlineto{\pgfqpoint{1.362129in}{2.720547in}}%
\pgfpathlineto{\pgfqpoint{1.362143in}{2.728796in}}%
\pgfpathlineto{\pgfqpoint{1.362930in}{2.644559in}}%
\pgfpathlineto{\pgfqpoint{1.365046in}{2.644559in}}%
\pgfpathlineto{\pgfqpoint{1.365051in}{2.762816in}}%
\pgfpathlineto{\pgfqpoint{1.365799in}{2.698813in}}%
\pgfpathlineto{\pgfqpoint{1.365849in}{2.698281in}}%
\pgfpathlineto{\pgfqpoint{1.365898in}{2.705028in}}%
\pgfpathlineto{\pgfqpoint{1.365948in}{2.704494in}}%
\pgfpathlineto{\pgfqpoint{1.366978in}{2.772890in}}%
\pgfpathlineto{\pgfqpoint{1.367020in}{2.749939in}}%
\pgfpathlineto{\pgfqpoint{1.367046in}{2.759650in}}%
\pgfpathlineto{\pgfqpoint{1.367790in}{2.644559in}}%
\pgfpathlineto{\pgfqpoint{1.369992in}{2.644559in}}%
\pgfpathlineto{\pgfqpoint{1.369996in}{2.762668in}}%
\pgfpathlineto{\pgfqpoint{1.370751in}{2.696451in}}%
\pgfpathlineto{\pgfqpoint{1.371965in}{2.774602in}}%
\pgfpathlineto{\pgfqpoint{1.371966in}{2.749786in}}%
\pgfpathlineto{\pgfqpoint{1.371991in}{2.759357in}}%
\pgfpathlineto{\pgfqpoint{1.372739in}{2.644559in}}%
\pgfpathlineto{\pgfqpoint{1.374937in}{2.644559in}}%
\pgfpathlineto{\pgfqpoint{1.374941in}{2.762656in}}%
\pgfpathlineto{\pgfqpoint{1.375666in}{2.683760in}}%
\pgfpathlineto{\pgfqpoint{1.376507in}{2.759979in}}%
\pgfpathlineto{\pgfqpoint{1.376882in}{2.783503in}}%
\pgfpathlineto{\pgfqpoint{1.377558in}{2.644559in}}%
\pgfpathlineto{\pgfqpoint{1.379882in}{2.644559in}}%
\pgfpathlineto{\pgfqpoint{1.379887in}{2.761888in}}%
\pgfpathlineto{\pgfqpoint{1.379992in}{2.639090in}}%
\pgfpathlineto{\pgfqpoint{1.380630in}{2.714447in}}%
\pgfpathlineto{\pgfqpoint{1.380680in}{2.681723in}}%
\pgfpathlineto{\pgfqpoint{1.381323in}{2.757874in}}%
\pgfpathlineto{\pgfqpoint{1.381372in}{2.725286in}}%
\pgfpathlineto{\pgfqpoint{1.381812in}{2.788565in}}%
\pgfpathlineto{\pgfqpoint{1.381948in}{2.644559in}}%
\pgfpathlineto{\pgfqpoint{1.382037in}{2.644559in}}%
\pgfpathlineto{\pgfqpoint{1.384828in}{2.644559in}}%
\pgfpathlineto{\pgfqpoint{1.384832in}{2.759606in}}%
\pgfpathlineto{\pgfqpoint{1.385560in}{2.682778in}}%
\pgfpathlineto{\pgfqpoint{1.386400in}{2.760384in}}%
\pgfpathlineto{\pgfqpoint{1.386774in}{2.783837in}}%
\pgfpathlineto{\pgfqpoint{1.387433in}{2.644559in}}%
\pgfpathlineto{\pgfqpoint{1.389773in}{2.644559in}}%
\pgfpathlineto{\pgfqpoint{1.389778in}{2.762452in}}%
\pgfpathlineto{\pgfqpoint{1.390515in}{2.698123in}}%
\pgfpathlineto{\pgfqpoint{1.390564in}{2.697425in}}%
\pgfpathlineto{\pgfqpoint{1.390614in}{2.704337in}}%
\pgfpathlineto{\pgfqpoint{1.390663in}{2.703637in}}%
\pgfpathlineto{\pgfqpoint{1.391702in}{2.772693in}}%
\pgfpathlineto{\pgfqpoint{1.391747in}{2.749882in}}%
\pgfpathlineto{\pgfqpoint{1.391774in}{2.759769in}}%
\pgfpathlineto{\pgfqpoint{1.392516in}{2.644559in}}%
\pgfpathlineto{\pgfqpoint{1.394719in}{2.644559in}}%
\pgfpathlineto{\pgfqpoint{1.394723in}{2.761881in}}%
\pgfpathlineto{\pgfqpoint{1.394859in}{2.629795in}}%
\pgfpathlineto{\pgfqpoint{1.395453in}{2.667010in}}%
\pgfpathlineto{\pgfqpoint{1.396293in}{2.776984in}}%
\pgfpathlineto{\pgfqpoint{1.396666in}{2.800318in}}%
\pgfpathlineto{\pgfqpoint{1.397182in}{2.644559in}}%
\pgfpathlineto{\pgfqpoint{1.399664in}{2.644559in}}%
\pgfpathlineto{\pgfqpoint{1.399669in}{2.759443in}}%
\pgfpathlineto{\pgfqpoint{1.399730in}{2.625454in}}%
\pgfpathlineto{\pgfqpoint{1.400409in}{2.663805in}}%
\pgfpathlineto{\pgfqpoint{1.401250in}{2.780871in}}%
\pgfpathlineto{\pgfqpoint{1.401638in}{2.805169in}}%
\pgfpathlineto{\pgfqpoint{1.402127in}{2.644559in}}%
\pgfpathlineto{\pgfqpoint{1.404610in}{2.644559in}}%
\pgfpathlineto{\pgfqpoint{1.404614in}{2.758156in}}%
\pgfpathlineto{\pgfqpoint{1.404695in}{2.640436in}}%
\pgfpathlineto{\pgfqpoint{1.405382in}{2.681590in}}%
\pgfpathlineto{\pgfqpoint{1.406223in}{2.766280in}}%
\pgfpathlineto{\pgfqpoint{1.406584in}{2.788886in}}%
\pgfpathlineto{\pgfqpoint{1.407203in}{2.644559in}}%
\pgfpathlineto{\pgfqpoint{1.409555in}{2.644559in}}%
\pgfpathlineto{\pgfqpoint{1.409560in}{2.762822in}}%
\pgfpathlineto{\pgfqpoint{1.410303in}{2.696908in}}%
\pgfpathlineto{\pgfqpoint{1.411529in}{2.773443in}}%
\pgfpathlineto{\pgfqpoint{1.411581in}{2.731126in}}%
\pgfpathlineto{\pgfqpoint{1.412419in}{2.644559in}}%
\pgfpathlineto{\pgfqpoint{1.414501in}{2.644559in}}%
\pgfpathlineto{\pgfqpoint{1.414505in}{2.762552in}}%
\pgfpathlineto{\pgfqpoint{1.414585in}{2.622600in}}%
\pgfpathlineto{\pgfqpoint{1.415270in}{2.663761in}}%
\pgfpathlineto{\pgfqpoint{1.416111in}{2.784892in}}%
\pgfpathlineto{\pgfqpoint{1.416474in}{2.807621in}}%
\pgfpathlineto{\pgfqpoint{1.416989in}{2.644559in}}%
\pgfpathlineto{\pgfqpoint{1.419446in}{2.644559in}}%
\pgfpathlineto{\pgfqpoint{1.419451in}{2.761872in}}%
\pgfpathlineto{\pgfqpoint{1.419531in}{2.622402in}}%
\pgfpathlineto{\pgfqpoint{1.420216in}{2.663548in}}%
\pgfpathlineto{\pgfqpoint{1.421056in}{2.784893in}}%
\pgfpathlineto{\pgfqpoint{1.421420in}{2.807633in}}%
\pgfpathlineto{\pgfqpoint{1.421931in}{2.644559in}}%
\pgfpathlineto{\pgfqpoint{1.424392in}{2.644559in}}%
\pgfpathlineto{\pgfqpoint{1.424396in}{2.761689in}}%
\pgfpathlineto{\pgfqpoint{1.424476in}{2.622348in}}%
\pgfpathlineto{\pgfqpoint{1.425161in}{2.663496in}}%
\pgfpathlineto{\pgfqpoint{1.426002in}{2.784898in}}%
\pgfpathlineto{\pgfqpoint{1.426365in}{2.807633in}}%
\pgfpathlineto{\pgfqpoint{1.426876in}{2.644559in}}%
\pgfpathlineto{\pgfqpoint{1.429337in}{2.644559in}}%
\pgfpathlineto{\pgfqpoint{1.429341in}{2.761824in}}%
\pgfpathlineto{\pgfqpoint{1.429422in}{2.622199in}}%
\pgfpathlineto{\pgfqpoint{1.430107in}{2.663348in}}%
\pgfpathlineto{\pgfqpoint{1.430947in}{2.785082in}}%
\pgfpathlineto{\pgfqpoint{1.431311in}{2.807818in}}%
\pgfpathlineto{\pgfqpoint{1.431821in}{2.644559in}}%
\pgfpathlineto{\pgfqpoint{1.434282in}{2.644559in}}%
\pgfpathlineto{\pgfqpoint{1.434287in}{2.761616in}}%
\pgfpathlineto{\pgfqpoint{1.434354in}{2.643557in}}%
\pgfpathlineto{\pgfqpoint{1.435037in}{2.683095in}}%
\pgfpathlineto{\pgfqpoint{1.435878in}{2.763558in}}%
\pgfpathlineto{\pgfqpoint{1.436256in}{2.787276in}}%
\pgfpathlineto{\pgfqpoint{1.436888in}{2.644559in}}%
\pgfpathlineto{\pgfqpoint{1.439228in}{2.644559in}}%
\pgfpathlineto{\pgfqpoint{1.439232in}{2.761970in}}%
\pgfpathlineto{\pgfqpoint{1.439313in}{2.622461in}}%
\pgfpathlineto{\pgfqpoint{1.439998in}{2.663612in}}%
\pgfpathlineto{\pgfqpoint{1.440838in}{2.784866in}}%
\pgfpathlineto{\pgfqpoint{1.441202in}{2.807600in}}%
\pgfpathlineto{\pgfqpoint{1.441712in}{2.644559in}}%
\pgfpathlineto{\pgfqpoint{1.444173in}{2.644559in}}%
\pgfpathlineto{\pgfqpoint{1.444178in}{2.761319in}}%
\pgfpathlineto{\pgfqpoint{1.444920in}{2.684774in}}%
\pgfpathlineto{\pgfqpoint{1.445761in}{2.760788in}}%
\pgfpathlineto{\pgfqpoint{1.446147in}{2.785023in}}%
\pgfpathlineto{\pgfqpoint{1.446819in}{2.644559in}}%
\pgfpathlineto{\pgfqpoint{1.449119in}{2.644559in}}%
\pgfpathlineto{\pgfqpoint{1.449123in}{2.761699in}}%
\pgfpathlineto{\pgfqpoint{1.449204in}{2.622233in}}%
\pgfpathlineto{\pgfqpoint{1.449888in}{2.663381in}}%
\pgfpathlineto{\pgfqpoint{1.450729in}{2.785012in}}%
\pgfpathlineto{\pgfqpoint{1.451093in}{2.807749in}}%
\pgfpathlineto{\pgfqpoint{1.451603in}{2.644559in}}%
\pgfpathlineto{\pgfqpoint{1.454064in}{2.644559in}}%
\pgfpathlineto{\pgfqpoint{1.454069in}{2.761185in}}%
\pgfpathlineto{\pgfqpoint{1.454801in}{2.685861in}}%
\pgfpathlineto{\pgfqpoint{1.455741in}{2.764673in}}%
\pgfpathlineto{\pgfqpoint{1.456013in}{2.781746in}}%
\pgfpathlineto{\pgfqpoint{1.456867in}{2.644559in}}%
\pgfpathlineto{\pgfqpoint{1.459010in}{2.644559in}}%
\pgfpathlineto{\pgfqpoint{1.459014in}{2.759942in}}%
\pgfpathlineto{\pgfqpoint{1.459751in}{2.685357in}}%
\pgfpathlineto{\pgfqpoint{1.460592in}{2.759139in}}%
\pgfpathlineto{\pgfqpoint{1.460984in}{2.783722in}}%
\pgfpathlineto{\pgfqpoint{1.461680in}{2.644559in}}%
\pgfpathlineto{\pgfqpoint{1.463955in}{2.644559in}}%
\pgfpathlineto{\pgfqpoint{1.463960in}{2.758330in}}%
\pgfpathlineto{\pgfqpoint{1.464041in}{2.621330in}}%
\pgfpathlineto{\pgfqpoint{1.464725in}{2.662518in}}%
\pgfpathlineto{\pgfqpoint{1.465566in}{2.785001in}}%
\pgfpathlineto{\pgfqpoint{1.465929in}{2.807698in}}%
\pgfpathlineto{\pgfqpoint{1.466440in}{2.644559in}}%
\pgfpathlineto{\pgfqpoint{1.468901in}{2.644559in}}%
\pgfpathlineto{\pgfqpoint{1.468905in}{2.761523in}}%
\pgfpathlineto{\pgfqpoint{1.468986in}{2.622326in}}%
\pgfpathlineto{\pgfqpoint{1.469670in}{2.663487in}}%
\pgfpathlineto{\pgfqpoint{1.470511in}{2.784883in}}%
\pgfpathlineto{\pgfqpoint{1.470874in}{2.807606in}}%
\pgfpathlineto{\pgfqpoint{1.471385in}{2.644559in}}%
\pgfpathlineto{\pgfqpoint{1.473846in}{2.644559in}}%
\pgfpathlineto{\pgfqpoint{1.473851in}{2.760947in}}%
\pgfpathlineto{\pgfqpoint{1.474604in}{2.706312in}}%
\pgfpathlineto{\pgfqpoint{1.474654in}{2.690953in}}%
\pgfpathlineto{\pgfqpoint{1.475198in}{2.743557in}}%
\pgfpathlineto{\pgfqpoint{1.475247in}{2.728271in}}%
\pgfpathlineto{\pgfqpoint{1.475781in}{2.780150in}}%
\pgfpathlineto{\pgfqpoint{1.475882in}{2.699390in}}%
\pgfpathlineto{\pgfqpoint{1.475899in}{2.699770in}}%
\pgfpathlineto{\pgfqpoint{1.476683in}{2.644559in}}%
\pgfpathlineto{\pgfqpoint{1.478792in}{2.644559in}}%
\pgfpathlineto{\pgfqpoint{1.478796in}{2.754191in}}%
\pgfpathlineto{\pgfqpoint{1.479522in}{2.686547in}}%
\pgfpathlineto{\pgfqpoint{1.480462in}{2.761185in}}%
\pgfpathlineto{\pgfqpoint{1.480737in}{2.778474in}}%
\pgfpathlineto{\pgfqpoint{1.481692in}{2.644559in}}%
\pgfpathlineto{\pgfqpoint{1.483737in}{2.644559in}}%
\pgfpathlineto{\pgfqpoint{1.483741in}{2.760671in}}%
\pgfpathlineto{\pgfqpoint{1.483822in}{2.622424in}}%
\pgfpathlineto{\pgfqpoint{1.484507in}{2.663557in}}%
\pgfpathlineto{\pgfqpoint{1.485347in}{2.784548in}}%
\pgfpathlineto{\pgfqpoint{1.485711in}{2.807278in}}%
\pgfpathlineto{\pgfqpoint{1.486226in}{2.644559in}}%
\pgfpathlineto{\pgfqpoint{1.488682in}{2.644559in}}%
\pgfpathlineto{\pgfqpoint{1.488687in}{2.760495in}}%
\pgfpathlineto{\pgfqpoint{1.489442in}{2.706461in}}%
\pgfpathlineto{\pgfqpoint{1.489492in}{2.690891in}}%
\pgfpathlineto{\pgfqpoint{1.490036in}{2.743705in}}%
\pgfpathlineto{\pgfqpoint{1.490085in}{2.728209in}}%
\pgfpathlineto{\pgfqpoint{1.490618in}{2.780245in}}%
\pgfpathlineto{\pgfqpoint{1.490736in}{2.688172in}}%
\pgfpathlineto{\pgfqpoint{1.491520in}{2.644559in}}%
\pgfpathlineto{\pgfqpoint{1.493628in}{2.644559in}}%
\pgfpathlineto{\pgfqpoint{1.493632in}{2.759595in}}%
\pgfpathlineto{\pgfqpoint{1.494374in}{2.690503in}}%
\pgfpathlineto{\pgfqpoint{1.495313in}{2.760683in}}%
\pgfpathlineto{\pgfqpoint{1.495602in}{2.778774in}}%
\pgfpathlineto{\pgfqpoint{1.496804in}{2.644559in}}%
\pgfpathlineto{\pgfqpoint{1.498573in}{2.644559in}}%
\pgfpathlineto{\pgfqpoint{1.498578in}{2.759840in}}%
\pgfpathlineto{\pgfqpoint{1.499317in}{2.685125in}}%
\pgfpathlineto{\pgfqpoint{1.500157in}{2.759587in}}%
\pgfpathlineto{\pgfqpoint{1.500547in}{2.784047in}}%
\pgfpathlineto{\pgfqpoint{1.501235in}{2.644559in}}%
\pgfpathlineto{\pgfqpoint{1.503519in}{2.644559in}}%
\pgfpathlineto{\pgfqpoint{1.503523in}{2.760910in}}%
\pgfpathlineto{\pgfqpoint{1.504289in}{2.705117in}}%
\pgfpathlineto{\pgfqpoint{1.504339in}{2.693687in}}%
\pgfpathlineto{\pgfqpoint{1.504883in}{2.742367in}}%
\pgfpathlineto{\pgfqpoint{1.504932in}{2.731000in}}%
\pgfpathlineto{\pgfqpoint{1.505460in}{2.778576in}}%
\pgfpathlineto{\pgfqpoint{1.505548in}{2.698273in}}%
\pgfpathlineto{\pgfqpoint{1.505563in}{2.750139in}}%
\pgfpathlineto{\pgfqpoint{1.506349in}{2.644559in}}%
\pgfpathlineto{\pgfqpoint{1.508464in}{2.644559in}}%
\pgfpathlineto{\pgfqpoint{1.508469in}{2.759587in}}%
\pgfpathlineto{\pgfqpoint{1.508542in}{2.631059in}}%
\pgfpathlineto{\pgfqpoint{1.509227in}{2.671618in}}%
\pgfpathlineto{\pgfqpoint{1.510068in}{2.775487in}}%
\pgfpathlineto{\pgfqpoint{1.510438in}{2.798653in}}%
\pgfpathlineto{\pgfqpoint{1.510973in}{2.644559in}}%
\pgfpathlineto{\pgfqpoint{1.513410in}{2.644559in}}%
\pgfpathlineto{\pgfqpoint{1.513414in}{2.758963in}}%
\pgfpathlineto{\pgfqpoint{1.513488in}{2.630170in}}%
\pgfpathlineto{\pgfqpoint{1.514173in}{2.670774in}}%
\pgfpathlineto{\pgfqpoint{1.515014in}{2.776202in}}%
\pgfpathlineto{\pgfqpoint{1.515384in}{2.799340in}}%
\pgfpathlineto{\pgfqpoint{1.515914in}{2.644559in}}%
\pgfpathlineto{\pgfqpoint{1.518355in}{2.644559in}}%
\pgfpathlineto{\pgfqpoint{1.518360in}{2.759986in}}%
\pgfpathlineto{\pgfqpoint{1.519105in}{2.704998in}}%
\pgfpathlineto{\pgfqpoint{1.519154in}{2.690951in}}%
\pgfpathlineto{\pgfqpoint{1.519698in}{2.742245in}}%
\pgfpathlineto{\pgfqpoint{1.519748in}{2.728266in}}%
\pgfpathlineto{\pgfqpoint{1.520286in}{2.779104in}}%
\pgfpathlineto{\pgfqpoint{1.520384in}{2.698110in}}%
\pgfpathlineto{\pgfqpoint{1.520399in}{2.750066in}}%
\pgfpathlineto{\pgfqpoint{1.521186in}{2.644559in}}%
\pgfpathlineto{\pgfqpoint{1.523301in}{2.644559in}}%
\pgfpathlineto{\pgfqpoint{1.523305in}{2.759949in}}%
\pgfpathlineto{\pgfqpoint{1.523377in}{2.633827in}}%
\pgfpathlineto{\pgfqpoint{1.524062in}{2.674183in}}%
\pgfpathlineto{\pgfqpoint{1.524903in}{2.772818in}}%
\pgfpathlineto{\pgfqpoint{1.525274in}{2.796102in}}%
\pgfpathlineto{\pgfqpoint{1.525820in}{2.644559in}}%
\pgfpathlineto{\pgfqpoint{1.528246in}{2.644559in}}%
\pgfpathlineto{\pgfqpoint{1.528251in}{2.757917in}}%
\pgfpathlineto{\pgfqpoint{1.528998in}{2.704939in}}%
\pgfpathlineto{\pgfqpoint{1.529047in}{2.690676in}}%
\pgfpathlineto{\pgfqpoint{1.529591in}{2.742186in}}%
\pgfpathlineto{\pgfqpoint{1.529641in}{2.727992in}}%
\pgfpathlineto{\pgfqpoint{1.530178in}{2.778981in}}%
\pgfpathlineto{\pgfqpoint{1.530275in}{2.697825in}}%
\pgfpathlineto{\pgfqpoint{1.530290in}{2.749777in}}%
\pgfpathlineto{\pgfqpoint{1.531077in}{2.644559in}}%
\pgfpathlineto{\pgfqpoint{1.533192in}{2.644559in}}%
\pgfpathlineto{\pgfqpoint{1.533196in}{2.760080in}}%
\pgfpathlineto{\pgfqpoint{1.533956in}{2.707066in}}%
\pgfpathlineto{\pgfqpoint{1.534005in}{2.690694in}}%
\pgfpathlineto{\pgfqpoint{1.534549in}{2.744308in}}%
\pgfpathlineto{\pgfqpoint{1.534598in}{2.728015in}}%
\pgfpathlineto{\pgfqpoint{1.535129in}{2.780713in}}%
\pgfpathlineto{\pgfqpoint{1.535244in}{2.698763in}}%
\pgfpathlineto{\pgfqpoint{1.536027in}{2.644559in}}%
\pgfpathlineto{\pgfqpoint{1.538137in}{2.644559in}}%
\pgfpathlineto{\pgfqpoint{1.538141in}{2.756250in}}%
\pgfpathlineto{\pgfqpoint{1.538876in}{2.691973in}}%
\pgfpathlineto{\pgfqpoint{1.540013in}{2.769800in}}%
\pgfpathlineto{\pgfqpoint{1.540167in}{2.714119in}}%
\pgfpathlineto{\pgfqpoint{1.540111in}{2.775932in}}%
\pgfpathlineto{\pgfqpoint{1.540182in}{2.735161in}}%
\pgfpathlineto{\pgfqpoint{1.540967in}{2.644559in}}%
\pgfpathlineto{\pgfqpoint{1.543082in}{2.644559in}}%
\pgfpathlineto{\pgfqpoint{1.543087in}{2.760964in}}%
\pgfpathlineto{\pgfqpoint{1.543829in}{2.684701in}}%
\pgfpathlineto{\pgfqpoint{1.544669in}{2.760683in}}%
\pgfpathlineto{\pgfqpoint{1.545056in}{2.784959in}}%
\pgfpathlineto{\pgfqpoint{1.545727in}{2.644559in}}%
\pgfpathlineto{\pgfqpoint{1.548028in}{2.644559in}}%
\pgfpathlineto{\pgfqpoint{1.548032in}{2.759636in}}%
\pgfpathlineto{\pgfqpoint{1.548766in}{2.692270in}}%
\pgfpathlineto{\pgfqpoint{1.549904in}{2.770458in}}%
\pgfpathlineto{\pgfqpoint{1.550058in}{2.713996in}}%
\pgfpathlineto{\pgfqpoint{1.550002in}{2.776591in}}%
\pgfpathlineto{\pgfqpoint{1.550073in}{2.736413in}}%
\pgfpathlineto{\pgfqpoint{1.550860in}{2.644559in}}%
\pgfpathlineto{\pgfqpoint{1.552973in}{2.644559in}}%
\pgfpathlineto{\pgfqpoint{1.552978in}{2.759840in}}%
\pgfpathlineto{\pgfqpoint{1.553043in}{2.639522in}}%
\pgfpathlineto{\pgfqpoint{1.553726in}{2.678898in}}%
\pgfpathlineto{\pgfqpoint{1.554567in}{2.767014in}}%
\pgfpathlineto{\pgfqpoint{1.554947in}{2.790859in}}%
\pgfpathlineto{\pgfqpoint{1.555524in}{2.644559in}}%
\pgfpathlineto{\pgfqpoint{1.557919in}{2.644559in}}%
\pgfpathlineto{\pgfqpoint{1.557923in}{2.754869in}}%
\pgfpathlineto{\pgfqpoint{1.557994in}{2.637146in}}%
\pgfpathlineto{\pgfqpoint{1.558679in}{2.677323in}}%
\pgfpathlineto{\pgfqpoint{1.559520in}{2.768109in}}%
\pgfpathlineto{\pgfqpoint{1.559893in}{2.791487in}}%
\pgfpathlineto{\pgfqpoint{1.560464in}{2.644559in}}%
\pgfpathlineto{\pgfqpoint{1.562864in}{2.644559in}}%
\pgfpathlineto{\pgfqpoint{1.562869in}{2.761014in}}%
\pgfpathlineto{\pgfqpoint{1.563607in}{2.685223in}}%
\pgfpathlineto{\pgfqpoint{1.564447in}{2.759725in}}%
\pgfpathlineto{\pgfqpoint{1.564838in}{2.784232in}}%
\pgfpathlineto{\pgfqpoint{1.565525in}{2.644559in}}%
\pgfpathlineto{\pgfqpoint{1.567810in}{2.644559in}}%
\pgfpathlineto{\pgfqpoint{1.567814in}{2.759146in}}%
\pgfpathlineto{\pgfqpoint{1.568577in}{2.689468in}}%
\pgfpathlineto{\pgfqpoint{1.569517in}{2.764474in}}%
\pgfpathlineto{\pgfqpoint{1.569784in}{2.781209in}}%
\pgfpathlineto{\pgfqpoint{1.570748in}{2.644559in}}%
\pgfpathlineto{\pgfqpoint{1.572755in}{2.644559in}}%
\pgfpathlineto{\pgfqpoint{1.572760in}{2.759381in}}%
\pgfpathlineto{\pgfqpoint{1.573521in}{2.705221in}}%
\pgfpathlineto{\pgfqpoint{1.573571in}{2.692628in}}%
\pgfpathlineto{\pgfqpoint{1.574115in}{2.742506in}}%
\pgfpathlineto{\pgfqpoint{1.574164in}{2.729904in}}%
\pgfpathlineto{\pgfqpoint{1.574694in}{2.778892in}}%
\pgfpathlineto{\pgfqpoint{1.574783in}{2.719180in}}%
\pgfpathlineto{\pgfqpoint{1.574797in}{2.728861in}}%
\pgfpathlineto{\pgfqpoint{1.575584in}{2.644559in}}%
\pgfpathlineto{\pgfqpoint{1.577701in}{2.644559in}}%
\pgfpathlineto{\pgfqpoint{1.577705in}{2.759353in}}%
\pgfpathlineto{\pgfqpoint{1.578459in}{2.705071in}}%
\pgfpathlineto{\pgfqpoint{1.578508in}{2.691780in}}%
\pgfpathlineto{\pgfqpoint{1.579052in}{2.742355in}}%
\pgfpathlineto{\pgfqpoint{1.579102in}{2.729059in}}%
\pgfpathlineto{\pgfqpoint{1.579635in}{2.778985in}}%
\pgfpathlineto{\pgfqpoint{1.579730in}{2.716163in}}%
\pgfpathlineto{\pgfqpoint{1.579744in}{2.732242in}}%
\pgfpathlineto{\pgfqpoint{1.580531in}{2.644559in}}%
\pgfpathlineto{\pgfqpoint{1.582646in}{2.644559in}}%
\pgfpathlineto{\pgfqpoint{1.582651in}{2.759253in}}%
\pgfpathlineto{\pgfqpoint{1.583394in}{2.687170in}}%
\pgfpathlineto{\pgfqpoint{1.584334in}{2.764315in}}%
\pgfpathlineto{\pgfqpoint{1.584620in}{2.782283in}}%
\pgfpathlineto{\pgfqpoint{1.585497in}{2.644559in}}%
\pgfpathlineto{\pgfqpoint{1.587592in}{2.644559in}}%
\pgfpathlineto{\pgfqpoint{1.587596in}{2.759059in}}%
\pgfpathlineto{\pgfqpoint{1.588341in}{2.687230in}}%
\pgfpathlineto{\pgfqpoint{1.589280in}{2.764324in}}%
\pgfpathlineto{\pgfqpoint{1.589565in}{2.782233in}}%
\pgfpathlineto{\pgfqpoint{1.590444in}{2.644559in}}%
\pgfpathlineto{\pgfqpoint{1.592537in}{2.644559in}}%
\pgfpathlineto{\pgfqpoint{1.592541in}{2.759094in}}%
\pgfpathlineto{\pgfqpoint{1.593294in}{2.687694in}}%
\pgfpathlineto{\pgfqpoint{1.594233in}{2.764834in}}%
\pgfpathlineto{\pgfqpoint{1.594511in}{2.782257in}}%
\pgfpathlineto{\pgfqpoint{1.595399in}{2.644559in}}%
\pgfpathlineto{\pgfqpoint{1.597482in}{2.644559in}}%
\pgfpathlineto{\pgfqpoint{1.597487in}{2.759040in}}%
\pgfpathlineto{\pgfqpoint{1.598251in}{2.689454in}}%
\pgfpathlineto{\pgfqpoint{1.599190in}{2.764520in}}%
\pgfpathlineto{\pgfqpoint{1.599456in}{2.781222in}}%
\pgfpathlineto{\pgfqpoint{1.600418in}{2.644559in}}%
\pgfpathlineto{\pgfqpoint{1.602428in}{2.644559in}}%
\pgfpathlineto{\pgfqpoint{1.602432in}{2.758910in}}%
\pgfpathlineto{\pgfqpoint{1.603164in}{2.691045in}}%
\pgfpathlineto{\pgfqpoint{1.604202in}{2.765039in}}%
\pgfpathlineto{\pgfqpoint{1.604376in}{2.775970in}}%
\pgfpathlineto{\pgfqpoint{1.606021in}{2.644559in}}%
\pgfpathlineto{\pgfqpoint{1.607373in}{2.644559in}}%
\pgfpathlineto{\pgfqpoint{1.607378in}{2.758917in}}%
\pgfpathlineto{\pgfqpoint{1.608114in}{2.691181in}}%
\pgfpathlineto{\pgfqpoint{1.609153in}{2.765530in}}%
\pgfpathlineto{\pgfqpoint{1.609347in}{2.777756in}}%
\pgfpathlineto{\pgfqpoint{1.610916in}{2.644559in}}%
\pgfpathlineto{\pgfqpoint{1.612319in}{2.644559in}}%
\pgfpathlineto{\pgfqpoint{1.612323in}{2.755326in}}%
\pgfpathlineto{\pgfqpoint{1.613069in}{2.699693in}}%
\pgfpathlineto{\pgfqpoint{1.613118in}{2.695030in}}%
\pgfpathlineto{\pgfqpoint{1.613465in}{2.724556in}}%
\pgfpathlineto{\pgfqpoint{1.613514in}{2.719876in}}%
\pgfpathlineto{\pgfqpoint{1.614249in}{2.773881in}}%
\pgfpathlineto{\pgfqpoint{1.614335in}{2.732100in}}%
\pgfpathlineto{\pgfqpoint{1.615188in}{2.644559in}}%
\pgfpathlineto{\pgfqpoint{1.617264in}{2.644559in}}%
\pgfpathlineto{\pgfqpoint{1.617269in}{2.760024in}}%
\pgfpathlineto{\pgfqpoint{1.617339in}{2.635647in}}%
\pgfpathlineto{\pgfqpoint{1.618024in}{2.675805in}}%
\pgfpathlineto{\pgfqpoint{1.618865in}{2.771017in}}%
\pgfpathlineto{\pgfqpoint{1.619238in}{2.794409in}}%
\pgfpathlineto{\pgfqpoint{1.619796in}{2.644559in}}%
\pgfpathlineto{\pgfqpoint{1.622210in}{2.644559in}}%
\pgfpathlineto{\pgfqpoint{1.622215in}{2.749268in}}%
\pgfpathlineto{\pgfqpoint{1.622279in}{2.638805in}}%
\pgfpathlineto{\pgfqpoint{1.622962in}{2.678220in}}%
\pgfpathlineto{\pgfqpoint{1.623803in}{2.767084in}}%
\pgfpathlineto{\pgfqpoint{1.624184in}{2.790936in}}%
\pgfpathlineto{\pgfqpoint{1.624759in}{2.644559in}}%
\pgfpathlineto{\pgfqpoint{1.627155in}{2.644559in}}%
\pgfpathlineto{\pgfqpoint{1.627160in}{2.756238in}}%
\pgfpathlineto{\pgfqpoint{1.627896in}{2.691505in}}%
\pgfpathlineto{\pgfqpoint{1.628934in}{2.764332in}}%
\pgfpathlineto{\pgfqpoint{1.629186in}{2.712694in}}%
\pgfpathlineto{\pgfqpoint{1.629129in}{2.776533in}}%
\pgfpathlineto{\pgfqpoint{1.629201in}{2.737036in}}%
\pgfpathlineto{\pgfqpoint{1.629986in}{2.644559in}}%
\pgfpathlineto{\pgfqpoint{1.632101in}{2.644559in}}%
\pgfpathlineto{\pgfqpoint{1.632105in}{2.759216in}}%
\pgfpathlineto{\pgfqpoint{1.632849in}{2.689734in}}%
\pgfpathlineto{\pgfqpoint{1.633789in}{2.761664in}}%
\pgfpathlineto{\pgfqpoint{1.634074in}{2.779590in}}%
\pgfpathlineto{\pgfqpoint{1.635155in}{2.644559in}}%
\pgfpathlineto{\pgfqpoint{1.637046in}{2.644559in}}%
\pgfpathlineto{\pgfqpoint{1.637052in}{2.749574in}}%
\pgfpathlineto{\pgfqpoint{1.637782in}{2.693076in}}%
\pgfpathlineto{\pgfqpoint{1.638994in}{2.773671in}}%
\pgfpathlineto{\pgfqpoint{1.639020in}{2.768830in}}%
\pgfpathlineto{\pgfqpoint{1.639793in}{2.644559in}}%
\pgfpathlineto{\pgfqpoint{1.641992in}{2.644559in}}%
\pgfpathlineto{\pgfqpoint{1.641997in}{2.749276in}}%
\pgfpathlineto{\pgfqpoint{1.642068in}{2.635895in}}%
\pgfpathlineto{\pgfqpoint{1.642753in}{2.676310in}}%
\pgfpathlineto{\pgfqpoint{1.643594in}{2.770374in}}%
\pgfpathlineto{\pgfqpoint{1.643965in}{2.793652in}}%
\pgfpathlineto{\pgfqpoint{1.644529in}{2.644559in}}%
\pgfpathlineto{\pgfqpoint{1.646937in}{2.644559in}}%
\pgfpathlineto{\pgfqpoint{1.646943in}{2.749298in}}%
\pgfpathlineto{\pgfqpoint{1.647011in}{2.637625in}}%
\pgfpathlineto{\pgfqpoint{1.647696in}{2.677779in}}%
\pgfpathlineto{\pgfqpoint{1.648537in}{2.768537in}}%
\pgfpathlineto{\pgfqpoint{1.648911in}{2.791960in}}%
\pgfpathlineto{\pgfqpoint{1.649485in}{2.644559in}}%
\pgfpathlineto{\pgfqpoint{1.651882in}{2.644559in}}%
\pgfpathlineto{\pgfqpoint{1.651888in}{2.749253in}}%
\pgfpathlineto{\pgfqpoint{1.651958in}{2.636230in}}%
\pgfpathlineto{\pgfqpoint{1.652644in}{2.676593in}}%
\pgfpathlineto{\pgfqpoint{1.653484in}{2.769936in}}%
\pgfpathlineto{\pgfqpoint{1.653856in}{2.793244in}}%
\pgfpathlineto{\pgfqpoint{1.654419in}{2.644559in}}%
\pgfpathlineto{\pgfqpoint{1.656828in}{2.644559in}}%
\pgfpathlineto{\pgfqpoint{1.656833in}{2.749303in}}%
\pgfpathlineto{\pgfqpoint{1.656896in}{2.639541in}}%
\pgfpathlineto{\pgfqpoint{1.657579in}{2.678776in}}%
\pgfpathlineto{\pgfqpoint{1.658420in}{2.766382in}}%
\pgfpathlineto{\pgfqpoint{1.658802in}{2.790334in}}%
\pgfpathlineto{\pgfqpoint{1.659381in}{2.644559in}}%
\pgfpathlineto{\pgfqpoint{1.661773in}{2.644559in}}%
\pgfpathlineto{\pgfqpoint{1.661779in}{2.749406in}}%
\pgfpathlineto{\pgfqpoint{1.662514in}{2.691759in}}%
\pgfpathlineto{\pgfqpoint{1.663652in}{2.770640in}}%
\pgfpathlineto{\pgfqpoint{1.663804in}{2.714118in}}%
\pgfpathlineto{\pgfqpoint{1.663747in}{2.776639in}}%
\pgfpathlineto{\pgfqpoint{1.663818in}{2.735574in}}%
\pgfpathlineto{\pgfqpoint{1.664603in}{2.644559in}}%
\pgfpathlineto{\pgfqpoint{1.666719in}{2.644559in}}%
\pgfpathlineto{\pgfqpoint{1.666724in}{2.749249in}}%
\pgfpathlineto{\pgfqpoint{1.666796in}{2.635653in}}%
\pgfpathlineto{\pgfqpoint{1.667481in}{2.676119in}}%
\pgfpathlineto{\pgfqpoint{1.668322in}{2.770595in}}%
\pgfpathlineto{\pgfqpoint{1.668693in}{2.793845in}}%
\pgfpathlineto{\pgfqpoint{1.669251in}{2.644559in}}%
\pgfpathlineto{\pgfqpoint{1.671664in}{2.644559in}}%
\pgfpathlineto{\pgfqpoint{1.671670in}{2.749328in}}%
\pgfpathlineto{\pgfqpoint{1.672408in}{2.690995in}}%
\pgfpathlineto{\pgfqpoint{1.673447in}{2.765550in}}%
\pgfpathlineto{\pgfqpoint{1.673638in}{2.777558in}}%
\pgfpathlineto{\pgfqpoint{1.675194in}{2.644559in}}%
\pgfpathlineto{\pgfqpoint{1.676610in}{2.644559in}}%
\pgfpathlineto{\pgfqpoint{1.676614in}{2.758048in}}%
\pgfpathlineto{\pgfqpoint{1.677369in}{2.700089in}}%
\pgfpathlineto{\pgfqpoint{1.677418in}{2.696540in}}%
\pgfpathlineto{\pgfqpoint{1.677764in}{2.724953in}}%
\pgfpathlineto{\pgfqpoint{1.677814in}{2.721385in}}%
\pgfpathlineto{\pgfqpoint{1.678545in}{2.773997in}}%
\pgfpathlineto{\pgfqpoint{1.678614in}{2.743639in}}%
\pgfpathlineto{\pgfqpoint{1.679438in}{2.644559in}}%
\pgfpathlineto{\pgfqpoint{1.681555in}{2.644559in}}%
\pgfpathlineto{\pgfqpoint{1.681561in}{2.749533in}}%
\pgfpathlineto{\pgfqpoint{1.682315in}{2.690707in}}%
\pgfpathlineto{\pgfqpoint{1.683255in}{2.761344in}}%
\pgfpathlineto{\pgfqpoint{1.683529in}{2.778569in}}%
\pgfpathlineto{\pgfqpoint{1.684703in}{2.644559in}}%
\pgfpathlineto{\pgfqpoint{1.686501in}{2.644559in}}%
\pgfpathlineto{\pgfqpoint{1.686505in}{2.756045in}}%
\pgfpathlineto{\pgfqpoint{1.687273in}{2.701892in}}%
\pgfpathlineto{\pgfqpoint{1.687322in}{2.695771in}}%
\pgfpathlineto{\pgfqpoint{1.687767in}{2.732971in}}%
\pgfpathlineto{\pgfqpoint{1.687817in}{2.726828in}}%
\pgfpathlineto{\pgfqpoint{1.688442in}{2.775394in}}%
\pgfpathlineto{\pgfqpoint{1.688526in}{2.734694in}}%
\pgfpathlineto{\pgfqpoint{1.689340in}{2.644559in}}%
\pgfpathlineto{\pgfqpoint{1.691446in}{2.644559in}}%
\pgfpathlineto{\pgfqpoint{1.691452in}{2.749519in}}%
\pgfpathlineto{\pgfqpoint{1.692187in}{2.691754in}}%
\pgfpathlineto{\pgfqpoint{1.693226in}{2.764753in}}%
\pgfpathlineto{\pgfqpoint{1.693477in}{2.713903in}}%
\pgfpathlineto{\pgfqpoint{1.693420in}{2.776943in}}%
\pgfpathlineto{\pgfqpoint{1.693491in}{2.736181in}}%
\pgfpathlineto{\pgfqpoint{1.694278in}{2.644559in}}%
\pgfpathlineto{\pgfqpoint{1.696392in}{2.644564in}}%
\pgfpathlineto{\pgfqpoint{1.696396in}{2.781907in}}%
\pgfpathlineto{\pgfqpoint{1.697117in}{2.701155in}}%
\pgfpathlineto{\pgfqpoint{1.697166in}{2.690949in}}%
\pgfpathlineto{\pgfqpoint{1.697611in}{2.732205in}}%
\pgfpathlineto{\pgfqpoint{1.697661in}{2.722034in}}%
\pgfpathlineto{\pgfqpoint{1.698365in}{2.779553in}}%
\pgfpathlineto{\pgfqpoint{1.698420in}{2.718497in}}%
\pgfpathlineto{\pgfqpoint{1.698434in}{2.729397in}}%
\pgfpathlineto{\pgfqpoint{1.699221in}{2.644559in}}%
\pgfpathlineto{\pgfqpoint{1.701337in}{2.644559in}}%
\pgfpathlineto{\pgfqpoint{1.701342in}{2.749467in}}%
\pgfpathlineto{\pgfqpoint{1.702076in}{2.692192in}}%
\pgfpathlineto{\pgfqpoint{1.703214in}{2.770184in}}%
\pgfpathlineto{\pgfqpoint{1.703367in}{2.715322in}}%
\pgfpathlineto{\pgfqpoint{1.703311in}{2.776268in}}%
\pgfpathlineto{\pgfqpoint{1.703381in}{2.734176in}}%
\pgfpathlineto{\pgfqpoint{1.704168in}{2.644559in}}%
\pgfpathlineto{\pgfqpoint{1.706282in}{2.644559in}}%
\pgfpathlineto{\pgfqpoint{1.706288in}{2.749457in}}%
\pgfpathlineto{\pgfqpoint{1.707021in}{2.692337in}}%
\pgfpathlineto{\pgfqpoint{1.708158in}{2.769742in}}%
\pgfpathlineto{\pgfqpoint{1.708289in}{2.734065in}}%
\pgfpathlineto{\pgfqpoint{1.708256in}{2.775904in}}%
\pgfpathlineto{\pgfqpoint{1.708299in}{2.734392in}}%
\pgfpathlineto{\pgfqpoint{1.709133in}{2.644559in}}%
\pgfpathlineto{\pgfqpoint{1.711228in}{2.644559in}}%
\pgfpathlineto{\pgfqpoint{1.711233in}{2.749460in}}%
\pgfpathlineto{\pgfqpoint{1.711966in}{2.692533in}}%
\pgfpathlineto{\pgfqpoint{1.713202in}{2.775771in}}%
\pgfpathlineto{\pgfqpoint{1.713973in}{2.644559in}}%
\pgfpathlineto{\pgfqpoint{1.716173in}{2.644559in}}%
\pgfpathlineto{\pgfqpoint{1.716179in}{2.749518in}}%
\pgfpathlineto{\pgfqpoint{1.716907in}{2.693439in}}%
\pgfpathlineto{\pgfqpoint{1.718121in}{2.772903in}}%
\pgfpathlineto{\pgfqpoint{1.718147in}{2.762593in}}%
\pgfpathlineto{\pgfqpoint{1.718919in}{2.644559in}}%
\pgfpathlineto{\pgfqpoint{1.721119in}{2.644559in}}%
\pgfpathlineto{\pgfqpoint{1.721124in}{2.749569in}}%
\pgfpathlineto{\pgfqpoint{1.721849in}{2.694550in}}%
\pgfpathlineto{\pgfqpoint{1.723093in}{2.772756in}}%
\pgfpathlineto{\pgfqpoint{1.723123in}{2.745716in}}%
\pgfpathlineto{\pgfqpoint{1.723939in}{2.644559in}}%
\pgfpathlineto{\pgfqpoint{1.726064in}{2.644559in}}%
\pgfpathlineto{\pgfqpoint{1.726069in}{2.752672in}}%
\pgfpathlineto{\pgfqpoint{1.726808in}{2.698449in}}%
\pgfpathlineto{\pgfqpoint{1.726857in}{2.694676in}}%
\pgfpathlineto{\pgfqpoint{1.727203in}{2.723314in}}%
\pgfpathlineto{\pgfqpoint{1.727253in}{2.719520in}}%
\pgfpathlineto{\pgfqpoint{1.727992in}{2.772852in}}%
\pgfpathlineto{\pgfqpoint{1.728068in}{2.746722in}}%
\pgfpathlineto{\pgfqpoint{1.728879in}{2.644559in}}%
\pgfpathlineto{\pgfqpoint{1.731010in}{2.644559in}}%
\pgfpathlineto{\pgfqpoint{1.731014in}{2.754997in}}%
\pgfpathlineto{\pgfqpoint{1.731086in}{2.635773in}}%
\pgfpathlineto{\pgfqpoint{1.731772in}{2.676160in}}%
\pgfpathlineto{\pgfqpoint{1.732612in}{2.769520in}}%
\pgfpathlineto{\pgfqpoint{1.732984in}{2.792783in}}%
\pgfpathlineto{\pgfqpoint{1.733547in}{2.644559in}}%
\pgfpathlineto{\pgfqpoint{1.735955in}{2.644559in}}%
\pgfpathlineto{\pgfqpoint{1.735960in}{2.757843in}}%
\pgfpathlineto{\pgfqpoint{1.736724in}{2.694003in}}%
\pgfpathlineto{\pgfqpoint{1.737861in}{2.772105in}}%
\pgfpathlineto{\pgfqpoint{1.737929in}{2.776379in}}%
\pgfpathlineto{\pgfqpoint{1.738007in}{2.692889in}}%
\pgfpathlineto{\pgfqpoint{1.738792in}{2.644559in}}%
\pgfpathlineto{\pgfqpoint{1.740901in}{2.644559in}}%
\pgfpathlineto{\pgfqpoint{1.740906in}{2.749691in}}%
\pgfpathlineto{\pgfqpoint{1.741639in}{2.698231in}}%
\pgfpathlineto{\pgfqpoint{1.741688in}{2.695016in}}%
\pgfpathlineto{\pgfqpoint{1.741936in}{2.716878in}}%
\pgfpathlineto{\pgfqpoint{1.741985in}{2.713650in}}%
\pgfpathlineto{\pgfqpoint{1.742826in}{2.772817in}}%
\pgfpathlineto{\pgfqpoint{1.742875in}{2.749047in}}%
\pgfpathlineto{\pgfqpoint{1.742897in}{2.756577in}}%
\pgfpathlineto{\pgfqpoint{1.743646in}{2.644559in}}%
\pgfpathlineto{\pgfqpoint{1.745846in}{2.644559in}}%
\pgfpathlineto{\pgfqpoint{1.745852in}{2.749423in}}%
\pgfpathlineto{\pgfqpoint{1.746589in}{2.691182in}}%
\pgfpathlineto{\pgfqpoint{1.747628in}{2.765463in}}%
\pgfpathlineto{\pgfqpoint{1.747820in}{2.777515in}}%
\pgfpathlineto{\pgfqpoint{1.749418in}{2.644559in}}%
\pgfpathlineto{\pgfqpoint{1.750792in}{2.644559in}}%
\pgfpathlineto{\pgfqpoint{1.750797in}{2.749594in}}%
\pgfpathlineto{\pgfqpoint{1.751555in}{2.692062in}}%
\pgfpathlineto{\pgfqpoint{1.752594in}{2.766861in}}%
\pgfpathlineto{\pgfqpoint{1.752765in}{2.777657in}}%
\pgfpathlineto{\pgfqpoint{1.754295in}{2.644559in}}%
\pgfpathlineto{\pgfqpoint{1.755737in}{2.644559in}}%
\pgfpathlineto{\pgfqpoint{1.755741in}{2.754490in}}%
\pgfpathlineto{\pgfqpoint{1.755865in}{2.632194in}}%
\pgfpathlineto{\pgfqpoint{1.756508in}{2.721684in}}%
\pgfpathlineto{\pgfqpoint{1.756558in}{2.675498in}}%
\pgfpathlineto{\pgfqpoint{1.757201in}{2.765103in}}%
\pgfpathlineto{\pgfqpoint{1.757250in}{2.719068in}}%
\pgfpathlineto{\pgfqpoint{1.757678in}{2.795048in}}%
\pgfpathlineto{\pgfqpoint{1.757803in}{2.644559in}}%
\pgfpathlineto{\pgfqpoint{1.757951in}{2.644559in}}%
\pgfpathlineto{\pgfqpoint{1.760682in}{2.644559in}}%
\pgfpathlineto{\pgfqpoint{1.760687in}{2.754296in}}%
\pgfpathlineto{\pgfqpoint{1.761452in}{2.693813in}}%
\pgfpathlineto{\pgfqpoint{1.762589in}{2.771431in}}%
\pgfpathlineto{\pgfqpoint{1.762711in}{2.697564in}}%
\pgfpathlineto{\pgfqpoint{1.762656in}{2.775634in}}%
\pgfpathlineto{\pgfqpoint{1.762726in}{2.749003in}}%
\pgfpathlineto{\pgfqpoint{1.763513in}{2.644559in}}%
\pgfpathlineto{\pgfqpoint{1.765628in}{2.644559in}}%
\pgfpathlineto{\pgfqpoint{1.765633in}{2.749363in}}%
\pgfpathlineto{\pgfqpoint{1.766371in}{2.691217in}}%
\pgfpathlineto{\pgfqpoint{1.767409in}{2.765258in}}%
\pgfpathlineto{\pgfqpoint{1.767659in}{2.712529in}}%
\pgfpathlineto{\pgfqpoint{1.767602in}{2.777319in}}%
\pgfpathlineto{\pgfqpoint{1.767674in}{2.737573in}}%
\pgfpathlineto{\pgfqpoint{1.768460in}{2.644559in}}%
\pgfpathlineto{\pgfqpoint{1.770574in}{2.644560in}}%
\pgfpathlineto{\pgfqpoint{1.770578in}{2.773675in}}%
\pgfpathlineto{\pgfqpoint{1.771318in}{2.702389in}}%
\pgfpathlineto{\pgfqpoint{1.771368in}{2.692147in}}%
\pgfpathlineto{\pgfqpoint{1.771813in}{2.733438in}}%
\pgfpathlineto{\pgfqpoint{1.771862in}{2.723233in}}%
\pgfpathlineto{\pgfqpoint{1.772501in}{2.776666in}}%
\pgfpathlineto{\pgfqpoint{1.772599in}{2.732581in}}%
\pgfpathlineto{\pgfqpoint{1.773421in}{2.644559in}}%
\pgfpathlineto{\pgfqpoint{1.775519in}{2.644559in}}%
\pgfpathlineto{\pgfqpoint{1.775523in}{2.766637in}}%
\pgfpathlineto{\pgfqpoint{1.775625in}{2.636428in}}%
\pgfpathlineto{\pgfqpoint{1.776264in}{2.714973in}}%
\pgfpathlineto{\pgfqpoint{1.776313in}{2.679044in}}%
\pgfpathlineto{\pgfqpoint{1.776956in}{2.758395in}}%
\pgfpathlineto{\pgfqpoint{1.777005in}{2.722611in}}%
\pgfpathlineto{\pgfqpoint{1.777447in}{2.789177in}}%
\pgfpathlineto{\pgfqpoint{1.777585in}{2.644559in}}%
\pgfpathlineto{\pgfqpoint{1.777684in}{2.644559in}}%
\pgfpathlineto{\pgfqpoint{1.780464in}{2.644559in}}%
\pgfpathlineto{\pgfqpoint{1.780470in}{2.749794in}}%
\pgfpathlineto{\pgfqpoint{1.781192in}{2.697896in}}%
\pgfpathlineto{\pgfqpoint{1.781242in}{2.694405in}}%
\pgfpathlineto{\pgfqpoint{1.781588in}{2.722759in}}%
\pgfpathlineto{\pgfqpoint{1.781637in}{2.719251in}}%
\pgfpathlineto{\pgfqpoint{1.782438in}{2.776193in}}%
\pgfpathlineto{\pgfqpoint{1.782470in}{2.737107in}}%
\pgfpathlineto{\pgfqpoint{1.783391in}{2.644559in}}%
\pgfpathlineto{\pgfqpoint{1.785410in}{2.644559in}}%
\pgfpathlineto{\pgfqpoint{1.785415in}{2.749676in}}%
\pgfpathlineto{\pgfqpoint{1.786144in}{2.687984in}}%
\pgfpathlineto{\pgfqpoint{1.787084in}{2.761172in}}%
\pgfpathlineto{\pgfqpoint{1.787357in}{2.778345in}}%
\pgfpathlineto{\pgfqpoint{1.788377in}{2.644559in}}%
\pgfpathlineto{\pgfqpoint{1.790355in}{2.644559in}}%
\pgfpathlineto{\pgfqpoint{1.790361in}{2.749744in}}%
\pgfpathlineto{\pgfqpoint{1.791100in}{2.698813in}}%
\pgfpathlineto{\pgfqpoint{1.791150in}{2.695617in}}%
\pgfpathlineto{\pgfqpoint{1.791397in}{2.717460in}}%
\pgfpathlineto{\pgfqpoint{1.791447in}{2.714251in}}%
\pgfpathlineto{\pgfqpoint{1.792283in}{2.773157in}}%
\pgfpathlineto{\pgfqpoint{1.792351in}{2.756488in}}%
\pgfpathlineto{\pgfqpoint{1.793183in}{2.644559in}}%
\pgfpathlineto{\pgfqpoint{1.795301in}{2.644559in}}%
\pgfpathlineto{\pgfqpoint{1.795306in}{2.749693in}}%
\pgfpathlineto{\pgfqpoint{1.796058in}{2.701097in}}%
\pgfpathlineto{\pgfqpoint{1.796108in}{2.694985in}}%
\pgfpathlineto{\pgfqpoint{1.796553in}{2.732173in}}%
\pgfpathlineto{\pgfqpoint{1.796602in}{2.726045in}}%
\pgfpathlineto{\pgfqpoint{1.797235in}{2.775040in}}%
\pgfpathlineto{\pgfqpoint{1.797325in}{2.737829in}}%
\pgfpathlineto{\pgfqpoint{1.798133in}{2.644559in}}%
\pgfpathlineto{\pgfqpoint{1.800246in}{2.644559in}}%
\pgfpathlineto{\pgfqpoint{1.800251in}{2.749887in}}%
\pgfpathlineto{\pgfqpoint{1.800974in}{2.690262in}}%
\pgfpathlineto{\pgfqpoint{1.802012in}{2.763819in}}%
\pgfpathlineto{\pgfqpoint{1.802274in}{2.717864in}}%
\pgfpathlineto{\pgfqpoint{1.802190in}{2.775009in}}%
\pgfpathlineto{\pgfqpoint{1.802288in}{2.729184in}}%
\pgfpathlineto{\pgfqpoint{1.803077in}{2.644559in}}%
\pgfpathlineto{\pgfqpoint{1.805192in}{2.644559in}}%
\pgfpathlineto{\pgfqpoint{1.805197in}{2.749658in}}%
\pgfpathlineto{\pgfqpoint{1.805946in}{2.690037in}}%
\pgfpathlineto{\pgfqpoint{1.806886in}{2.761650in}}%
\pgfpathlineto{\pgfqpoint{1.807165in}{2.779204in}}%
\pgfpathlineto{\pgfqpoint{1.808266in}{2.644559in}}%
\pgfpathlineto{\pgfqpoint{1.810137in}{2.644559in}}%
\pgfpathlineto{\pgfqpoint{1.810142in}{2.749759in}}%
\pgfpathlineto{\pgfqpoint{1.810868in}{2.697916in}}%
\pgfpathlineto{\pgfqpoint{1.810918in}{2.694715in}}%
\pgfpathlineto{\pgfqpoint{1.811165in}{2.716563in}}%
\pgfpathlineto{\pgfqpoint{1.811214in}{2.713349in}}%
\pgfpathlineto{\pgfqpoint{1.812111in}{2.776001in}}%
\pgfpathlineto{\pgfqpoint{1.812111in}{2.748868in}}%
\pgfpathlineto{\pgfqpoint{1.812132in}{2.756289in}}%
\pgfpathlineto{\pgfqpoint{1.812881in}{2.644559in}}%
\pgfpathlineto{\pgfqpoint{1.815082in}{2.644559in}}%
\pgfpathlineto{\pgfqpoint{1.815087in}{2.753531in}}%
\pgfpathlineto{\pgfqpoint{1.815223in}{2.632566in}}%
\pgfpathlineto{\pgfqpoint{1.815816in}{2.669760in}}%
\pgfpathlineto{\pgfqpoint{1.816657in}{2.772031in}}%
\pgfpathlineto{\pgfqpoint{1.817030in}{2.795400in}}%
\pgfpathlineto{\pgfqpoint{1.817562in}{2.644559in}}%
\pgfpathlineto{\pgfqpoint{1.820028in}{2.644559in}}%
\pgfpathlineto{\pgfqpoint{1.820032in}{2.755188in}}%
\pgfpathlineto{\pgfqpoint{1.820101in}{2.634064in}}%
\pgfpathlineto{\pgfqpoint{1.820789in}{2.674083in}}%
\pgfpathlineto{\pgfqpoint{1.821629in}{2.771585in}}%
\pgfpathlineto{\pgfqpoint{1.822002in}{2.794934in}}%
\pgfpathlineto{\pgfqpoint{1.822552in}{2.644559in}}%
\pgfpathlineto{\pgfqpoint{1.824973in}{2.644559in}}%
\pgfpathlineto{\pgfqpoint{1.824979in}{2.749864in}}%
\pgfpathlineto{\pgfqpoint{1.825719in}{2.691565in}}%
\pgfpathlineto{\pgfqpoint{1.826757in}{2.765096in}}%
\pgfpathlineto{\pgfqpoint{1.826947in}{2.777036in}}%
\pgfpathlineto{\pgfqpoint{1.827026in}{2.662274in}}%
\pgfpathlineto{\pgfqpoint{1.827811in}{2.644559in}}%
\pgfpathlineto{\pgfqpoint{1.829919in}{2.644559in}}%
\pgfpathlineto{\pgfqpoint{1.829924in}{2.749388in}}%
\pgfpathlineto{\pgfqpoint{1.830664in}{2.688610in}}%
\pgfpathlineto{\pgfqpoint{1.831604in}{2.761108in}}%
\pgfpathlineto{\pgfqpoint{1.831893in}{2.779269in}}%
\pgfpathlineto{\pgfqpoint{1.832928in}{2.644559in}}%
\pgfpathlineto{\pgfqpoint{1.834864in}{2.644559in}}%
\pgfpathlineto{\pgfqpoint{1.834870in}{2.749577in}}%
\pgfpathlineto{\pgfqpoint{1.835590in}{2.695873in}}%
\pgfpathlineto{\pgfqpoint{1.835738in}{2.702486in}}%
\pgfpathlineto{\pgfqpoint{1.836838in}{2.774350in}}%
\pgfpathlineto{\pgfqpoint{1.836838in}{2.743001in}}%
\pgfpathlineto{\pgfqpoint{1.836865in}{2.759212in}}%
\pgfpathlineto{\pgfqpoint{1.837611in}{2.644559in}}%
\pgfpathlineto{\pgfqpoint{1.839810in}{2.644559in}}%
\pgfpathlineto{\pgfqpoint{1.839815in}{2.749716in}}%
\pgfpathlineto{\pgfqpoint{1.840550in}{2.691097in}}%
\pgfpathlineto{\pgfqpoint{1.841589in}{2.764568in}}%
\pgfpathlineto{\pgfqpoint{1.841784in}{2.776792in}}%
\pgfpathlineto{\pgfqpoint{1.841862in}{2.662267in}}%
\pgfpathlineto{\pgfqpoint{1.842647in}{2.644559in}}%
\pgfpathlineto{\pgfqpoint{1.844755in}{2.644559in}}%
\pgfpathlineto{\pgfqpoint{1.844761in}{2.749559in}}%
\pgfpathlineto{\pgfqpoint{1.845501in}{2.688894in}}%
\pgfpathlineto{\pgfqpoint{1.846440in}{2.761331in}}%
\pgfpathlineto{\pgfqpoint{1.846729in}{2.779471in}}%
\pgfpathlineto{\pgfqpoint{1.847772in}{2.644559in}}%
\pgfpathlineto{\pgfqpoint{1.849701in}{2.644559in}}%
\pgfpathlineto{\pgfqpoint{1.849706in}{2.749753in}}%
\pgfpathlineto{\pgfqpoint{1.850431in}{2.697853in}}%
\pgfpathlineto{\pgfqpoint{1.850481in}{2.694669in}}%
\pgfpathlineto{\pgfqpoint{1.850728in}{2.716501in}}%
\pgfpathlineto{\pgfqpoint{1.850777in}{2.713304in}}%
\pgfpathlineto{\pgfqpoint{1.851674in}{2.775986in}}%
\pgfpathlineto{\pgfqpoint{1.851675in}{2.748861in}}%
\pgfpathlineto{\pgfqpoint{1.851696in}{2.756272in}}%
\pgfpathlineto{\pgfqpoint{1.852445in}{2.644559in}}%
\pgfpathlineto{\pgfqpoint{1.854645in}{2.644559in}}%
\pgfpathlineto{\pgfqpoint{1.854651in}{2.754103in}}%
\pgfpathlineto{\pgfqpoint{1.854785in}{2.644054in}}%
\pgfpathlineto{\pgfqpoint{1.855379in}{2.681240in}}%
\pgfpathlineto{\pgfqpoint{1.856219in}{2.761061in}}%
\pgfpathlineto{\pgfqpoint{1.856593in}{2.784493in}}%
\pgfpathlineto{\pgfqpoint{1.857229in}{2.644559in}}%
\pgfpathlineto{\pgfqpoint{1.859592in}{2.644559in}}%
\pgfpathlineto{\pgfqpoint{1.859597in}{2.749965in}}%
\pgfpathlineto{\pgfqpoint{1.860341in}{2.699381in}}%
\pgfpathlineto{\pgfqpoint{1.860390in}{2.695576in}}%
\pgfpathlineto{\pgfqpoint{1.860736in}{2.724242in}}%
\pgfpathlineto{\pgfqpoint{1.860786in}{2.720424in}}%
\pgfpathlineto{\pgfqpoint{1.861522in}{2.773595in}}%
\pgfpathlineto{\pgfqpoint{1.861596in}{2.742697in}}%
\pgfpathlineto{\pgfqpoint{1.862424in}{2.644559in}}%
\pgfpathlineto{\pgfqpoint{1.864537in}{2.644559in}}%
\pgfpathlineto{\pgfqpoint{1.864542in}{2.749457in}}%
\pgfpathlineto{\pgfqpoint{1.865293in}{2.699088in}}%
\pgfpathlineto{\pgfqpoint{1.865342in}{2.695588in}}%
\pgfpathlineto{\pgfqpoint{1.865689in}{2.723950in}}%
\pgfpathlineto{\pgfqpoint{1.865738in}{2.720435in}}%
\pgfpathlineto{\pgfqpoint{1.866471in}{2.773089in}}%
\pgfpathlineto{\pgfqpoint{1.866542in}{2.740928in}}%
\pgfpathlineto{\pgfqpoint{1.867380in}{2.644559in}}%
\pgfpathlineto{\pgfqpoint{1.869482in}{2.644559in}}%
\pgfpathlineto{\pgfqpoint{1.869487in}{2.754094in}}%
\pgfpathlineto{\pgfqpoint{1.870239in}{2.691586in}}%
\pgfpathlineto{\pgfqpoint{1.871278in}{2.765758in}}%
\pgfpathlineto{\pgfqpoint{1.871456in}{2.776982in}}%
\pgfpathlineto{\pgfqpoint{1.873070in}{2.644559in}}%
\pgfpathlineto{\pgfqpoint{1.874428in}{2.644559in}}%
\pgfpathlineto{\pgfqpoint{1.874433in}{2.749599in}}%
\pgfpathlineto{\pgfqpoint{1.875192in}{2.700086in}}%
\pgfpathlineto{\pgfqpoint{1.875242in}{2.695999in}}%
\pgfpathlineto{\pgfqpoint{1.875588in}{2.724948in}}%
\pgfpathlineto{\pgfqpoint{1.875637in}{2.720846in}}%
\pgfpathlineto{\pgfqpoint{1.876366in}{2.773820in}}%
\pgfpathlineto{\pgfqpoint{1.876452in}{2.735783in}}%
\pgfpathlineto{\pgfqpoint{1.877263in}{2.644559in}}%
\pgfpathlineto{\pgfqpoint{1.879373in}{2.644559in}}%
\pgfpathlineto{\pgfqpoint{1.879378in}{2.753287in}}%
\pgfpathlineto{\pgfqpoint{1.879439in}{2.632341in}}%
\pgfpathlineto{\pgfqpoint{1.880121in}{2.671133in}}%
\pgfpathlineto{\pgfqpoint{1.880962in}{2.772395in}}%
\pgfpathlineto{\pgfqpoint{1.881347in}{2.796542in}}%
\pgfpathlineto{\pgfqpoint{1.881872in}{2.644559in}}%
\pgfpathlineto{\pgfqpoint{1.884319in}{2.644559in}}%
\pgfpathlineto{\pgfqpoint{1.884324in}{2.749742in}}%
\pgfpathlineto{\pgfqpoint{1.885062in}{2.700493in}}%
\pgfpathlineto{\pgfqpoint{1.885112in}{2.693955in}}%
\pgfpathlineto{\pgfqpoint{1.885557in}{2.731569in}}%
\pgfpathlineto{\pgfqpoint{1.885606in}{2.725015in}}%
\pgfpathlineto{\pgfqpoint{1.886246in}{2.774887in}}%
\pgfpathlineto{\pgfqpoint{1.886343in}{2.738644in}}%
\pgfpathlineto{\pgfqpoint{1.887146in}{2.644559in}}%
\pgfpathlineto{\pgfqpoint{1.889264in}{2.644559in}}%
\pgfpathlineto{\pgfqpoint{1.889270in}{2.749276in}}%
\pgfpathlineto{\pgfqpoint{1.890012in}{2.698953in}}%
\pgfpathlineto{\pgfqpoint{1.890061in}{2.694565in}}%
\pgfpathlineto{\pgfqpoint{1.890408in}{2.723815in}}%
\pgfpathlineto{\pgfqpoint{1.890457in}{2.719411in}}%
\pgfpathlineto{\pgfqpoint{1.891194in}{2.773213in}}%
\pgfpathlineto{\pgfqpoint{1.891280in}{2.733101in}}%
\pgfpathlineto{\pgfqpoint{1.892120in}{2.644559in}}%
\pgfpathlineto{\pgfqpoint{1.894210in}{2.644559in}}%
\pgfpathlineto{\pgfqpoint{1.894215in}{2.749466in}}%
\pgfpathlineto{\pgfqpoint{1.894947in}{2.692603in}}%
\pgfpathlineto{\pgfqpoint{1.896184in}{2.775645in}}%
\pgfpathlineto{\pgfqpoint{1.896955in}{2.644559in}}%
\pgfpathlineto{\pgfqpoint{1.899155in}{2.644559in}}%
\pgfpathlineto{\pgfqpoint{1.899161in}{2.749864in}}%
\pgfpathlineto{\pgfqpoint{1.899893in}{2.698492in}}%
\pgfpathlineto{\pgfqpoint{1.899943in}{2.694669in}}%
\pgfpathlineto{\pgfqpoint{1.900289in}{2.723352in}}%
\pgfpathlineto{\pgfqpoint{1.900338in}{2.719518in}}%
\pgfpathlineto{\pgfqpoint{1.901080in}{2.773070in}}%
\pgfpathlineto{\pgfqpoint{1.901161in}{2.737917in}}%
\pgfpathlineto{\pgfqpoint{1.902053in}{2.644559in}}%
\pgfpathlineto{\pgfqpoint{1.904101in}{2.644559in}}%
\pgfpathlineto{\pgfqpoint{1.904106in}{2.749840in}}%
\pgfpathlineto{\pgfqpoint{1.904842in}{2.698685in}}%
\pgfpathlineto{\pgfqpoint{1.904891in}{2.694881in}}%
\pgfpathlineto{\pgfqpoint{1.905238in}{2.723545in}}%
\pgfpathlineto{\pgfqpoint{1.905287in}{2.719729in}}%
\pgfpathlineto{\pgfqpoint{1.906029in}{2.773264in}}%
\pgfpathlineto{\pgfqpoint{1.906105in}{2.741615in}}%
\pgfpathlineto{\pgfqpoint{1.906942in}{2.644559in}}%
\pgfpathlineto{\pgfqpoint{1.909046in}{2.644559in}}%
\pgfpathlineto{\pgfqpoint{1.909052in}{2.749667in}}%
\pgfpathlineto{\pgfqpoint{1.909773in}{2.690486in}}%
\pgfpathlineto{\pgfqpoint{1.910910in}{2.769375in}}%
\pgfpathlineto{\pgfqpoint{1.910990in}{2.774370in}}%
\pgfpathlineto{\pgfqpoint{1.911098in}{2.692499in}}%
\pgfpathlineto{\pgfqpoint{1.911882in}{2.644559in}}%
\pgfpathlineto{\pgfqpoint{1.913992in}{2.644559in}}%
\pgfpathlineto{\pgfqpoint{1.913997in}{2.749940in}}%
\pgfpathlineto{\pgfqpoint{1.914734in}{2.698905in}}%
\pgfpathlineto{\pgfqpoint{1.914784in}{2.695066in}}%
\pgfpathlineto{\pgfqpoint{1.915130in}{2.723765in}}%
\pgfpathlineto{\pgfqpoint{1.915179in}{2.719915in}}%
\pgfpathlineto{\pgfqpoint{1.915919in}{2.773312in}}%
\pgfpathlineto{\pgfqpoint{1.916007in}{2.733828in}}%
\pgfpathlineto{\pgfqpoint{1.916852in}{2.644559in}}%
\pgfpathlineto{\pgfqpoint{1.918937in}{2.644559in}}%
\pgfpathlineto{\pgfqpoint{1.918942in}{2.749841in}}%
\pgfpathlineto{\pgfqpoint{1.919692in}{2.699776in}}%
\pgfpathlineto{\pgfqpoint{1.919742in}{2.695633in}}%
\pgfpathlineto{\pgfqpoint{1.920088in}{2.724637in}}%
\pgfpathlineto{\pgfqpoint{1.920137in}{2.720480in}}%
\pgfpathlineto{\pgfqpoint{1.920870in}{2.773803in}}%
\pgfpathlineto{\pgfqpoint{1.920942in}{2.743042in}}%
\pgfpathlineto{\pgfqpoint{1.921768in}{2.644559in}}%
\pgfpathlineto{\pgfqpoint{1.923882in}{2.644559in}}%
\pgfpathlineto{\pgfqpoint{1.923888in}{2.749828in}}%
\pgfpathlineto{\pgfqpoint{1.924628in}{2.698994in}}%
\pgfpathlineto{\pgfqpoint{1.924678in}{2.694951in}}%
\pgfpathlineto{\pgfqpoint{1.925024in}{2.723853in}}%
\pgfpathlineto{\pgfqpoint{1.925073in}{2.719801in}}%
\pgfpathlineto{\pgfqpoint{1.925811in}{2.773301in}}%
\pgfpathlineto{\pgfqpoint{1.925898in}{2.736348in}}%
\pgfpathlineto{\pgfqpoint{1.926726in}{2.644559in}}%
\pgfpathlineto{\pgfqpoint{1.928828in}{2.644559in}}%
\pgfpathlineto{\pgfqpoint{1.928833in}{2.749737in}}%
\pgfpathlineto{\pgfqpoint{1.928935in}{2.628016in}}%
\pgfpathlineto{\pgfqpoint{1.929573in}{2.723663in}}%
\pgfpathlineto{\pgfqpoint{1.929622in}{2.670673in}}%
\pgfpathlineto{\pgfqpoint{1.930265in}{2.767060in}}%
\pgfpathlineto{\pgfqpoint{1.930315in}{2.714265in}}%
\pgfpathlineto{\pgfqpoint{1.930756in}{2.797815in}}%
\pgfpathlineto{\pgfqpoint{1.930895in}{2.644559in}}%
\pgfpathlineto{\pgfqpoint{1.931027in}{2.644559in}}%
\pgfpathlineto{\pgfqpoint{1.933774in}{2.644559in}}%
\pgfpathlineto{\pgfqpoint{1.933778in}{2.764509in}}%
\pgfpathlineto{\pgfqpoint{1.933853in}{2.638906in}}%
\pgfpathlineto{\pgfqpoint{1.934538in}{2.679740in}}%
\pgfpathlineto{\pgfqpoint{1.935379in}{2.766118in}}%
\pgfpathlineto{\pgfqpoint{1.935747in}{2.789200in}}%
\pgfpathlineto{\pgfqpoint{1.936345in}{2.644559in}}%
\pgfpathlineto{\pgfqpoint{1.938719in}{2.644559in}}%
\pgfpathlineto{\pgfqpoint{1.938724in}{2.749744in}}%
\pgfpathlineto{\pgfqpoint{1.939488in}{2.693881in}}%
\pgfpathlineto{\pgfqpoint{1.940625in}{2.771852in}}%
\pgfpathlineto{\pgfqpoint{1.940693in}{2.776092in}}%
\pgfpathlineto{\pgfqpoint{1.940772in}{2.657206in}}%
\pgfpathlineto{\pgfqpoint{1.941555in}{2.644559in}}%
\pgfpathlineto{\pgfqpoint{1.943664in}{2.644559in}}%
\pgfpathlineto{\pgfqpoint{1.943670in}{2.749611in}}%
\pgfpathlineto{\pgfqpoint{1.944404in}{2.698302in}}%
\pgfpathlineto{\pgfqpoint{1.944453in}{2.694486in}}%
\pgfpathlineto{\pgfqpoint{1.944799in}{2.723164in}}%
\pgfpathlineto{\pgfqpoint{1.944849in}{2.719334in}}%
\pgfpathlineto{\pgfqpoint{1.945591in}{2.772884in}}%
\pgfpathlineto{\pgfqpoint{1.945669in}{2.742346in}}%
\pgfpathlineto{\pgfqpoint{1.946500in}{2.644559in}}%
\pgfpathlineto{\pgfqpoint{1.948610in}{2.644559in}}%
\pgfpathlineto{\pgfqpoint{1.948615in}{2.749442in}}%
\pgfpathlineto{\pgfqpoint{1.949346in}{2.687820in}}%
\pgfpathlineto{\pgfqpoint{1.950286in}{2.760908in}}%
\pgfpathlineto{\pgfqpoint{1.950558in}{2.778028in}}%
\pgfpathlineto{\pgfqpoint{1.951579in}{2.644559in}}%
\pgfpathlineto{\pgfqpoint{1.953555in}{2.644559in}}%
\pgfpathlineto{\pgfqpoint{1.953561in}{2.749802in}}%
\pgfpathlineto{\pgfqpoint{1.954292in}{2.691258in}}%
\pgfpathlineto{\pgfqpoint{1.955429in}{2.770221in}}%
\pgfpathlineto{\pgfqpoint{1.955504in}{2.774916in}}%
\pgfpathlineto{\pgfqpoint{1.955607in}{2.692825in}}%
\pgfpathlineto{\pgfqpoint{1.956392in}{2.644559in}}%
\pgfpathlineto{\pgfqpoint{1.958501in}{2.644559in}}%
\pgfpathlineto{\pgfqpoint{1.958506in}{2.749470in}}%
\pgfpathlineto{\pgfqpoint{1.959240in}{2.698020in}}%
\pgfpathlineto{\pgfqpoint{1.959289in}{2.694220in}}%
\pgfpathlineto{\pgfqpoint{1.959635in}{2.722881in}}%
\pgfpathlineto{\pgfqpoint{1.959685in}{2.719069in}}%
\pgfpathlineto{\pgfqpoint{1.960427in}{2.772599in}}%
\pgfpathlineto{\pgfqpoint{1.960505in}{2.741045in}}%
\pgfpathlineto{\pgfqpoint{1.961344in}{2.644559in}}%
\pgfpathlineto{\pgfqpoint{1.963446in}{2.644559in}}%
\pgfpathlineto{\pgfqpoint{1.963452in}{2.749556in}}%
\pgfpathlineto{\pgfqpoint{1.964177in}{2.699509in}}%
\pgfpathlineto{\pgfqpoint{1.964226in}{2.692811in}}%
\pgfpathlineto{\pgfqpoint{1.964671in}{2.730584in}}%
\pgfpathlineto{\pgfqpoint{1.964721in}{2.723871in}}%
\pgfpathlineto{\pgfqpoint{1.965420in}{2.777627in}}%
\pgfpathlineto{\pgfqpoint{1.965474in}{2.720298in}}%
\pgfpathlineto{\pgfqpoint{1.965488in}{2.726824in}}%
\pgfpathlineto{\pgfqpoint{1.966273in}{2.644559in}}%
\pgfpathlineto{\pgfqpoint{1.968392in}{2.644559in}}%
\pgfpathlineto{\pgfqpoint{1.968397in}{2.749841in}}%
\pgfpathlineto{\pgfqpoint{1.969129in}{2.698496in}}%
\pgfpathlineto{\pgfqpoint{1.969179in}{2.694681in}}%
\pgfpathlineto{\pgfqpoint{1.969525in}{2.723357in}}%
\pgfpathlineto{\pgfqpoint{1.969574in}{2.719529in}}%
\pgfpathlineto{\pgfqpoint{1.970316in}{2.773077in}}%
\pgfpathlineto{\pgfqpoint{1.970396in}{2.742915in}}%
\pgfpathlineto{\pgfqpoint{1.971222in}{2.644559in}}%
\pgfpathlineto{\pgfqpoint{1.973337in}{2.644559in}}%
\pgfpathlineto{\pgfqpoint{1.973341in}{2.750090in}}%
\pgfpathlineto{\pgfqpoint{1.973422in}{2.628166in}}%
\pgfpathlineto{\pgfqpoint{1.974108in}{2.669459in}}%
\pgfpathlineto{\pgfqpoint{1.974949in}{2.776120in}}%
\pgfpathlineto{\pgfqpoint{1.975311in}{2.798777in}}%
\pgfpathlineto{\pgfqpoint{1.975845in}{2.644559in}}%
\pgfpathlineto{\pgfqpoint{1.978282in}{2.644559in}}%
\pgfpathlineto{\pgfqpoint{1.978288in}{2.749476in}}%
\pgfpathlineto{\pgfqpoint{1.979019in}{2.692812in}}%
\pgfpathlineto{\pgfqpoint{1.980231in}{2.773804in}}%
\pgfpathlineto{\pgfqpoint{1.980256in}{2.768706in}}%
\pgfpathlineto{\pgfqpoint{1.981026in}{2.644559in}}%
\pgfpathlineto{\pgfqpoint{1.983228in}{2.644559in}}%
\pgfpathlineto{\pgfqpoint{1.983233in}{2.749834in}}%
\pgfpathlineto{\pgfqpoint{1.983997in}{2.700062in}}%
\pgfpathlineto{\pgfqpoint{1.984047in}{2.696741in}}%
\pgfpathlineto{\pgfqpoint{1.984393in}{2.724921in}}%
\pgfpathlineto{\pgfqpoint{1.984442in}{2.721591in}}%
\pgfpathlineto{\pgfqpoint{1.985168in}{2.773638in}}%
\pgfpathlineto{\pgfqpoint{1.985244in}{2.738803in}}%
\pgfpathlineto{\pgfqpoint{1.986065in}{2.644559in}}%
\pgfpathlineto{\pgfqpoint{1.988173in}{2.644559in}}%
\pgfpathlineto{\pgfqpoint{1.988179in}{2.749818in}}%
\pgfpathlineto{\pgfqpoint{1.988918in}{2.691377in}}%
\pgfpathlineto{\pgfqpoint{1.989956in}{2.764825in}}%
\pgfpathlineto{\pgfqpoint{1.990201in}{2.720017in}}%
\pgfpathlineto{\pgfqpoint{1.990147in}{2.776831in}}%
\pgfpathlineto{\pgfqpoint{1.990215in}{2.727122in}}%
\pgfpathlineto{\pgfqpoint{1.991004in}{2.644559in}}%
\pgfpathlineto{\pgfqpoint{1.993119in}{2.644559in}}%
\pgfpathlineto{\pgfqpoint{1.993124in}{2.749857in}}%
\pgfpathlineto{\pgfqpoint{1.993871in}{2.699498in}}%
\pgfpathlineto{\pgfqpoint{1.993920in}{2.695409in}}%
\pgfpathlineto{\pgfqpoint{1.994267in}{2.724359in}}%
\pgfpathlineto{\pgfqpoint{1.994316in}{2.720258in}}%
\pgfpathlineto{\pgfqpoint{1.995051in}{2.773614in}}%
\pgfpathlineto{\pgfqpoint{1.995133in}{2.730894in}}%
\pgfpathlineto{\pgfqpoint{1.995144in}{2.731047in}}%
\pgfpathlineto{\pgfqpoint{1.995972in}{2.644559in}}%
\pgfpathlineto{\pgfqpoint{1.998064in}{2.644559in}}%
\pgfpathlineto{\pgfqpoint{1.998070in}{2.749430in}}%
\pgfpathlineto{\pgfqpoint{1.998826in}{2.699642in}}%
\pgfpathlineto{\pgfqpoint{1.998875in}{2.695737in}}%
\pgfpathlineto{\pgfqpoint{1.999221in}{2.724505in}}%
\pgfpathlineto{\pgfqpoint{1.999271in}{2.720583in}}%
\pgfpathlineto{\pgfqpoint{2.000001in}{2.773471in}}%
\pgfpathlineto{\pgfqpoint{2.000080in}{2.733161in}}%
\pgfpathlineto{\pgfqpoint{2.000920in}{2.644559in}}%
\pgfpathlineto{\pgfqpoint{2.003010in}{2.644559in}}%
\pgfpathlineto{\pgfqpoint{2.003015in}{2.749716in}}%
\pgfpathlineto{\pgfqpoint{2.003778in}{2.693803in}}%
\pgfpathlineto{\pgfqpoint{2.004915in}{2.771883in}}%
\pgfpathlineto{\pgfqpoint{2.004984in}{2.776166in}}%
\pgfpathlineto{\pgfqpoint{2.005062in}{2.675999in}}%
\pgfpathlineto{\pgfqpoint{2.005845in}{2.644559in}}%
\pgfpathlineto{\pgfqpoint{2.007955in}{2.644559in}}%
\pgfpathlineto{\pgfqpoint{2.007960in}{2.749841in}}%
\pgfpathlineto{\pgfqpoint{2.008049in}{2.643682in}}%
\pgfpathlineto{\pgfqpoint{2.008690in}{2.706741in}}%
\pgfpathlineto{\pgfqpoint{2.008740in}{2.685782in}}%
\pgfpathlineto{\pgfqpoint{2.009284in}{2.743995in}}%
\pgfpathlineto{\pgfqpoint{2.009333in}{2.723091in}}%
\pgfpathlineto{\pgfqpoint{2.009929in}{2.784497in}}%
\pgfpathlineto{\pgfqpoint{2.009983in}{2.721321in}}%
\pgfpathlineto{\pgfqpoint{2.009997in}{2.725356in}}%
\pgfpathlineto{\pgfqpoint{2.010785in}{2.644559in}}%
\pgfpathlineto{\pgfqpoint{2.012901in}{2.644559in}}%
\pgfpathlineto{\pgfqpoint{2.012906in}{2.749820in}}%
\pgfpathlineto{\pgfqpoint{2.013644in}{2.691353in}}%
\pgfpathlineto{\pgfqpoint{2.014683in}{2.764798in}}%
\pgfpathlineto{\pgfqpoint{2.014929in}{2.720051in}}%
\pgfpathlineto{\pgfqpoint{2.014874in}{2.776836in}}%
\pgfpathlineto{\pgfqpoint{2.014942in}{2.727101in}}%
\pgfpathlineto{\pgfqpoint{2.015731in}{2.644559in}}%
\pgfpathlineto{\pgfqpoint{2.017846in}{2.644559in}}%
\pgfpathlineto{\pgfqpoint{2.017851in}{2.749852in}}%
\pgfpathlineto{\pgfqpoint{2.018607in}{2.700442in}}%
\pgfpathlineto{\pgfqpoint{2.018657in}{2.695548in}}%
\pgfpathlineto{\pgfqpoint{2.019003in}{2.725301in}}%
\pgfpathlineto{\pgfqpoint{2.019052in}{2.720398in}}%
\pgfpathlineto{\pgfqpoint{2.019782in}{2.774272in}}%
\pgfpathlineto{\pgfqpoint{2.019862in}{2.736498in}}%
\pgfpathlineto{\pgfqpoint{2.020688in}{2.644559in}}%
\pgfpathlineto{\pgfqpoint{2.022792in}{2.644559in}}%
\pgfpathlineto{\pgfqpoint{2.022797in}{2.749846in}}%
\pgfpathlineto{\pgfqpoint{2.023517in}{2.690093in}}%
\pgfpathlineto{\pgfqpoint{2.024556in}{2.763616in}}%
\pgfpathlineto{\pgfqpoint{2.024820in}{2.718042in}}%
\pgfpathlineto{\pgfqpoint{2.024735in}{2.774862in}}%
\pgfpathlineto{\pgfqpoint{2.024834in}{2.728847in}}%
\pgfpathlineto{\pgfqpoint{2.025622in}{2.644559in}}%
\pgfpathlineto{\pgfqpoint{2.027737in}{2.644559in}}%
\pgfpathlineto{\pgfqpoint{2.027742in}{2.749852in}}%
\pgfpathlineto{\pgfqpoint{2.028489in}{2.699473in}}%
\pgfpathlineto{\pgfqpoint{2.028538in}{2.695387in}}%
\pgfpathlineto{\pgfqpoint{2.028885in}{2.724333in}}%
\pgfpathlineto{\pgfqpoint{2.028934in}{2.720236in}}%
\pgfpathlineto{\pgfqpoint{2.029669in}{2.773595in}}%
\pgfpathlineto{\pgfqpoint{2.029751in}{2.730977in}}%
\pgfpathlineto{\pgfqpoint{2.029775in}{2.713420in}}%
\pgfpathlineto{\pgfqpoint{2.030572in}{2.644559in}}%
\pgfpathlineto{\pgfqpoint{2.032682in}{2.644559in}}%
\pgfpathlineto{\pgfqpoint{2.032688in}{2.749811in}}%
\pgfpathlineto{\pgfqpoint{2.033427in}{2.691963in}}%
\pgfpathlineto{\pgfqpoint{2.034565in}{2.770127in}}%
\pgfpathlineto{\pgfqpoint{2.034711in}{2.716945in}}%
\pgfpathlineto{\pgfqpoint{2.034656in}{2.775866in}}%
\pgfpathlineto{\pgfqpoint{2.034725in}{2.729862in}}%
\pgfpathlineto{\pgfqpoint{2.035512in}{2.644559in}}%
\pgfpathlineto{\pgfqpoint{2.037628in}{2.644559in}}%
\pgfpathlineto{\pgfqpoint{2.037633in}{2.749829in}}%
\pgfpathlineto{\pgfqpoint{2.038385in}{2.699931in}}%
\pgfpathlineto{\pgfqpoint{2.038435in}{2.695694in}}%
\pgfpathlineto{\pgfqpoint{2.038781in}{2.724792in}}%
\pgfpathlineto{\pgfqpoint{2.038830in}{2.720542in}}%
\pgfpathlineto{\pgfqpoint{2.039562in}{2.773892in}}%
\pgfpathlineto{\pgfqpoint{2.039633in}{2.742370in}}%
\pgfpathlineto{\pgfqpoint{2.040464in}{2.644559in}}%
\pgfpathlineto{\pgfqpoint{2.042573in}{2.644559in}}%
\pgfpathlineto{\pgfqpoint{2.042579in}{2.749483in}}%
\pgfpathlineto{\pgfqpoint{2.043317in}{2.698267in}}%
\pgfpathlineto{\pgfqpoint{2.043366in}{2.694463in}}%
\pgfpathlineto{\pgfqpoint{2.043712in}{2.723128in}}%
\pgfpathlineto{\pgfqpoint{2.043762in}{2.719312in}}%
\pgfpathlineto{\pgfqpoint{2.044501in}{2.772657in}}%
\pgfpathlineto{\pgfqpoint{2.044579in}{2.738765in}}%
\pgfpathlineto{\pgfqpoint{2.045442in}{2.644559in}}%
\pgfpathlineto{\pgfqpoint{2.047519in}{2.644559in}}%
\pgfpathlineto{\pgfqpoint{2.047524in}{2.749338in}}%
\pgfpathlineto{\pgfqpoint{2.048280in}{2.690840in}}%
\pgfpathlineto{\pgfqpoint{2.049319in}{2.767067in}}%
\pgfpathlineto{\pgfqpoint{2.049493in}{2.777999in}}%
\pgfpathlineto{\pgfqpoint{2.050855in}{2.644559in}}%
\pgfpathlineto{\pgfqpoint{2.052464in}{2.644559in}}%
\pgfpathlineto{\pgfqpoint{2.052470in}{2.749838in}}%
\pgfpathlineto{\pgfqpoint{2.053223in}{2.700092in}}%
\pgfpathlineto{\pgfqpoint{2.053273in}{2.695787in}}%
\pgfpathlineto{\pgfqpoint{2.053619in}{2.724953in}}%
\pgfpathlineto{\pgfqpoint{2.053668in}{2.720635in}}%
\pgfpathlineto{\pgfqpoint{2.054399in}{2.773994in}}%
\pgfpathlineto{\pgfqpoint{2.054488in}{2.737125in}}%
\pgfpathlineto{\pgfqpoint{2.055295in}{2.644559in}}%
\pgfpathlineto{\pgfqpoint{2.057410in}{2.644559in}}%
\pgfpathlineto{\pgfqpoint{2.057415in}{2.749835in}}%
\pgfpathlineto{\pgfqpoint{2.057517in}{2.643938in}}%
\pgfpathlineto{\pgfqpoint{2.058159in}{2.707623in}}%
\pgfpathlineto{\pgfqpoint{2.058209in}{2.686742in}}%
\pgfpathlineto{\pgfqpoint{2.058753in}{2.744877in}}%
\pgfpathlineto{\pgfqpoint{2.058802in}{2.724050in}}%
\pgfpathlineto{\pgfqpoint{2.059340in}{2.781743in}}%
\pgfpathlineto{\pgfqpoint{2.059461in}{2.691690in}}%
\pgfpathlineto{\pgfqpoint{2.060245in}{2.644559in}}%
\pgfpathlineto{\pgfqpoint{2.062355in}{2.644559in}}%
\pgfpathlineto{\pgfqpoint{2.062361in}{2.749832in}}%
\pgfpathlineto{\pgfqpoint{2.063080in}{2.697551in}}%
\pgfpathlineto{\pgfqpoint{2.063129in}{2.693756in}}%
\pgfpathlineto{\pgfqpoint{2.063476in}{2.722411in}}%
\pgfpathlineto{\pgfqpoint{2.063525in}{2.718605in}}%
\pgfpathlineto{\pgfqpoint{2.064329in}{2.776026in}}%
\pgfpathlineto{\pgfqpoint{2.064371in}{2.734656in}}%
\pgfpathlineto{\pgfqpoint{2.065208in}{2.644559in}}%
\pgfpathlineto{\pgfqpoint{2.067301in}{2.644559in}}%
\pgfpathlineto{\pgfqpoint{2.067306in}{2.749843in}}%
\pgfpathlineto{\pgfqpoint{2.068036in}{2.698300in}}%
\pgfpathlineto{\pgfqpoint{2.068086in}{2.694504in}}%
\pgfpathlineto{\pgfqpoint{2.068432in}{2.723160in}}%
\pgfpathlineto{\pgfqpoint{2.068481in}{2.719352in}}%
\pgfpathlineto{\pgfqpoint{2.069274in}{2.776107in}}%
\pgfpathlineto{\pgfqpoint{2.069317in}{2.733074in}}%
\pgfpathlineto{\pgfqpoint{2.070167in}{2.644559in}}%
\pgfpathlineto{\pgfqpoint{2.072246in}{2.644559in}}%
\pgfpathlineto{\pgfqpoint{2.072251in}{2.749849in}}%
\pgfpathlineto{\pgfqpoint{2.073005in}{2.700240in}}%
\pgfpathlineto{\pgfqpoint{2.073054in}{2.695418in}}%
\pgfpathlineto{\pgfqpoint{2.073401in}{2.725099in}}%
\pgfpathlineto{\pgfqpoint{2.073450in}{2.720268in}}%
\pgfpathlineto{\pgfqpoint{2.074181in}{2.774143in}}%
\pgfpathlineto{\pgfqpoint{2.074262in}{2.737333in}}%
\pgfpathlineto{\pgfqpoint{2.075085in}{2.644559in}}%
\pgfpathlineto{\pgfqpoint{2.077192in}{2.644559in}}%
\pgfpathlineto{\pgfqpoint{2.077197in}{2.749847in}}%
\pgfpathlineto{\pgfqpoint{2.077951in}{2.700265in}}%
\pgfpathlineto{\pgfqpoint{2.078000in}{2.695437in}}%
\pgfpathlineto{\pgfqpoint{2.078346in}{2.725124in}}%
\pgfpathlineto{\pgfqpoint{2.078396in}{2.720287in}}%
\pgfpathlineto{\pgfqpoint{2.079127in}{2.774157in}}%
\pgfpathlineto{\pgfqpoint{2.079208in}{2.737186in}}%
\pgfpathlineto{\pgfqpoint{2.080032in}{2.644559in}}%
\pgfpathlineto{\pgfqpoint{2.082137in}{2.644559in}}%
\pgfpathlineto{\pgfqpoint{2.082142in}{2.749847in}}%
\pgfpathlineto{\pgfqpoint{2.082872in}{2.698295in}}%
\pgfpathlineto{\pgfqpoint{2.082922in}{2.694432in}}%
\pgfpathlineto{\pgfqpoint{2.083268in}{2.723154in}}%
\pgfpathlineto{\pgfqpoint{2.083318in}{2.719281in}}%
\pgfpathlineto{\pgfqpoint{2.084111in}{2.776101in}}%
\pgfpathlineto{\pgfqpoint{2.084153in}{2.734260in}}%
\pgfpathlineto{\pgfqpoint{2.084991in}{2.644559in}}%
\pgfpathlineto{\pgfqpoint{2.087082in}{2.644559in}}%
\pgfpathlineto{\pgfqpoint{2.087088in}{2.749848in}}%
\pgfpathlineto{\pgfqpoint{2.087842in}{2.700322in}}%
\pgfpathlineto{\pgfqpoint{2.087892in}{2.695472in}}%
\pgfpathlineto{\pgfqpoint{2.088238in}{2.725181in}}%
\pgfpathlineto{\pgfqpoint{2.088287in}{2.720322in}}%
\pgfpathlineto{\pgfqpoint{2.089018in}{2.774194in}}%
\pgfpathlineto{\pgfqpoint{2.089098in}{2.736987in}}%
\pgfpathlineto{\pgfqpoint{2.089925in}{2.644559in}}%
\pgfpathlineto{\pgfqpoint{2.092028in}{2.644559in}}%
\pgfpathlineto{\pgfqpoint{2.092033in}{2.749847in}}%
\pgfpathlineto{\pgfqpoint{2.092752in}{2.697536in}}%
\pgfpathlineto{\pgfqpoint{2.092802in}{2.693748in}}%
\pgfpathlineto{\pgfqpoint{2.093148in}{2.722395in}}%
\pgfpathlineto{\pgfqpoint{2.093197in}{2.718597in}}%
\pgfpathlineto{\pgfqpoint{2.094002in}{2.776045in}}%
\pgfpathlineto{\pgfqpoint{2.094044in}{2.734725in}}%
\pgfpathlineto{\pgfqpoint{2.094879in}{2.644559in}}%
\pgfpathlineto{\pgfqpoint{2.096973in}{2.644559in}}%
\pgfpathlineto{\pgfqpoint{2.096979in}{2.749847in}}%
\pgfpathlineto{\pgfqpoint{2.097732in}{2.700236in}}%
\pgfpathlineto{\pgfqpoint{2.097782in}{2.695412in}}%
\pgfpathlineto{\pgfqpoint{2.098128in}{2.725095in}}%
\pgfpathlineto{\pgfqpoint{2.098177in}{2.720262in}}%
\pgfpathlineto{\pgfqpoint{2.098908in}{2.774140in}}%
\pgfpathlineto{\pgfqpoint{2.098989in}{2.737319in}}%
\pgfpathlineto{\pgfqpoint{2.099816in}{2.644559in}}%
\pgfpathlineto{\pgfqpoint{2.101919in}{2.644559in}}%
\pgfpathlineto{\pgfqpoint{2.101924in}{2.749845in}}%
\pgfpathlineto{\pgfqpoint{2.102645in}{2.697674in}}%
\pgfpathlineto{\pgfqpoint{2.102695in}{2.693867in}}%
\pgfpathlineto{\pgfqpoint{2.103041in}{2.722534in}}%
\pgfpathlineto{\pgfqpoint{2.103090in}{2.718716in}}%
\pgfpathlineto{\pgfqpoint{2.103893in}{2.776056in}}%
\pgfpathlineto{\pgfqpoint{2.103935in}{2.734636in}}%
\pgfpathlineto{\pgfqpoint{2.104773in}{2.644559in}}%
\pgfpathlineto{\pgfqpoint{2.106864in}{2.644559in}}%
\pgfpathlineto{\pgfqpoint{2.106870in}{2.749851in}}%
\pgfpathlineto{\pgfqpoint{2.107592in}{2.697751in}}%
\pgfpathlineto{\pgfqpoint{2.107641in}{2.693932in}}%
\pgfpathlineto{\pgfqpoint{2.107987in}{2.722610in}}%
\pgfpathlineto{\pgfqpoint{2.108037in}{2.718782in}}%
\pgfpathlineto{\pgfqpoint{2.108838in}{2.776069in}}%
\pgfpathlineto{\pgfqpoint{2.108880in}{2.734630in}}%
\pgfpathlineto{\pgfqpoint{2.109716in}{2.644559in}}%
\pgfpathlineto{\pgfqpoint{2.111810in}{2.644559in}}%
\pgfpathlineto{\pgfqpoint{2.111815in}{2.749848in}}%
\pgfpathlineto{\pgfqpoint{2.112535in}{2.697598in}}%
\pgfpathlineto{\pgfqpoint{2.112584in}{2.693800in}}%
\pgfpathlineto{\pgfqpoint{2.112931in}{2.722458in}}%
\pgfpathlineto{\pgfqpoint{2.112980in}{2.718650in}}%
\pgfpathlineto{\pgfqpoint{2.113784in}{2.776052in}}%
\pgfpathlineto{\pgfqpoint{2.113826in}{2.734696in}}%
\pgfpathlineto{\pgfqpoint{2.114663in}{2.644559in}}%
\pgfpathlineto{\pgfqpoint{2.116755in}{2.644559in}}%
\pgfpathlineto{\pgfqpoint{2.116761in}{2.749845in}}%
\pgfpathlineto{\pgfqpoint{2.117484in}{2.697856in}}%
\pgfpathlineto{\pgfqpoint{2.117534in}{2.694026in}}%
\pgfpathlineto{\pgfqpoint{2.117880in}{2.722716in}}%
\pgfpathlineto{\pgfqpoint{2.117929in}{2.718875in}}%
\pgfpathlineto{\pgfqpoint{2.118729in}{2.776070in}}%
\pgfpathlineto{\pgfqpoint{2.118771in}{2.734543in}}%
\pgfpathlineto{\pgfqpoint{2.119608in}{2.644559in}}%
\pgfpathlineto{\pgfqpoint{2.121701in}{2.644559in}}%
\pgfpathlineto{\pgfqpoint{2.121706in}{2.749850in}}%
\pgfpathlineto{\pgfqpoint{2.122426in}{2.697588in}}%
\pgfpathlineto{\pgfqpoint{2.122475in}{2.693794in}}%
\pgfpathlineto{\pgfqpoint{2.122821in}{2.722448in}}%
\pgfpathlineto{\pgfqpoint{2.122871in}{2.718644in}}%
\pgfpathlineto{\pgfqpoint{2.123674in}{2.776055in}}%
\pgfpathlineto{\pgfqpoint{2.123717in}{2.734701in}}%
\pgfpathlineto{\pgfqpoint{2.124554in}{2.644559in}}%
\pgfpathlineto{\pgfqpoint{2.126646in}{2.644559in}}%
\pgfpathlineto{\pgfqpoint{2.126651in}{2.749845in}}%
\pgfpathlineto{\pgfqpoint{2.127389in}{2.698824in}}%
\pgfpathlineto{\pgfqpoint{2.127439in}{2.694831in}}%
\pgfpathlineto{\pgfqpoint{2.127785in}{2.723683in}}%
\pgfpathlineto{\pgfqpoint{2.127834in}{2.719681in}}%
\pgfpathlineto{\pgfqpoint{2.128573in}{2.773214in}}%
\pgfpathlineto{\pgfqpoint{2.128662in}{2.736641in}}%
\pgfpathlineto{\pgfqpoint{2.129489in}{2.644559in}}%
\pgfpathlineto{\pgfqpoint{2.131592in}{2.644559in}}%
\pgfpathlineto{\pgfqpoint{2.131597in}{2.749842in}}%
\pgfpathlineto{\pgfqpoint{2.132336in}{2.698867in}}%
\pgfpathlineto{\pgfqpoint{2.132385in}{2.694878in}}%
\pgfpathlineto{\pgfqpoint{2.132731in}{2.723726in}}%
\pgfpathlineto{\pgfqpoint{2.132781in}{2.719728in}}%
\pgfpathlineto{\pgfqpoint{2.133519in}{2.773233in}}%
\pgfpathlineto{\pgfqpoint{2.133608in}{2.736518in}}%
\pgfpathlineto{\pgfqpoint{2.134437in}{2.644559in}}%
\pgfpathlineto{\pgfqpoint{2.136537in}{2.644559in}}%
\pgfpathlineto{\pgfqpoint{2.136542in}{2.749839in}}%
\pgfpathlineto{\pgfqpoint{2.137282in}{2.698902in}}%
\pgfpathlineto{\pgfqpoint{2.137331in}{2.694910in}}%
\pgfpathlineto{\pgfqpoint{2.137677in}{2.723761in}}%
\pgfpathlineto{\pgfqpoint{2.137727in}{2.719760in}}%
\pgfpathlineto{\pgfqpoint{2.138465in}{2.773250in}}%
\pgfpathlineto{\pgfqpoint{2.138553in}{2.736432in}}%
\pgfpathlineto{\pgfqpoint{2.139380in}{2.644559in}}%
\pgfpathlineto{\pgfqpoint{2.141482in}{2.644559in}}%
\pgfpathlineto{\pgfqpoint{2.141488in}{2.749837in}}%
\pgfpathlineto{\pgfqpoint{2.142210in}{2.698068in}}%
\pgfpathlineto{\pgfqpoint{2.142260in}{2.693637in}}%
\pgfpathlineto{\pgfqpoint{2.142606in}{2.722927in}}%
\pgfpathlineto{\pgfqpoint{2.142656in}{2.718487in}}%
\pgfpathlineto{\pgfqpoint{2.143456in}{2.776341in}}%
\pgfpathlineto{\pgfqpoint{2.143480in}{2.757334in}}%
\pgfpathlineto{\pgfqpoint{2.144289in}{2.644559in}}%
\pgfpathlineto{\pgfqpoint{2.146428in}{2.644559in}}%
\pgfpathlineto{\pgfqpoint{2.146433in}{2.749845in}}%
\pgfpathlineto{\pgfqpoint{2.147154in}{2.697909in}}%
\pgfpathlineto{\pgfqpoint{2.147203in}{2.693555in}}%
\pgfpathlineto{\pgfqpoint{2.147550in}{2.722768in}}%
\pgfpathlineto{\pgfqpoint{2.147599in}{2.718405in}}%
\pgfpathlineto{\pgfqpoint{2.148402in}{2.776309in}}%
\pgfpathlineto{\pgfqpoint{2.148425in}{2.757379in}}%
\pgfpathlineto{\pgfqpoint{2.149235in}{2.644559in}}%
\pgfpathlineto{\pgfqpoint{2.151373in}{2.644559in}}%
\pgfpathlineto{\pgfqpoint{2.151379in}{2.749844in}}%
\pgfpathlineto{\pgfqpoint{2.152119in}{2.692070in}}%
\pgfpathlineto{\pgfqpoint{2.153257in}{2.770211in}}%
\pgfpathlineto{\pgfqpoint{2.153402in}{2.716972in}}%
\pgfpathlineto{\pgfqpoint{2.153347in}{2.775902in}}%
\pgfpathlineto{\pgfqpoint{2.153416in}{2.729930in}}%
\pgfpathlineto{\pgfqpoint{2.154202in}{2.644559in}}%
\pgfpathlineto{\pgfqpoint{2.156319in}{2.644559in}}%
\pgfpathlineto{\pgfqpoint{2.156324in}{2.749844in}}%
\pgfpathlineto{\pgfqpoint{2.157044in}{2.697853in}}%
\pgfpathlineto{\pgfqpoint{2.157094in}{2.693523in}}%
\pgfpathlineto{\pgfqpoint{2.157440in}{2.722711in}}%
\pgfpathlineto{\pgfqpoint{2.157489in}{2.718374in}}%
\pgfpathlineto{\pgfqpoint{2.158293in}{2.776294in}}%
\pgfpathlineto{\pgfqpoint{2.158316in}{2.757386in}}%
\pgfpathlineto{\pgfqpoint{2.159126in}{2.644559in}}%
\pgfpathlineto{\pgfqpoint{2.161264in}{2.644559in}}%
\pgfpathlineto{\pgfqpoint{2.161270in}{2.749846in}}%
\pgfpathlineto{\pgfqpoint{2.162015in}{2.692321in}}%
\pgfpathlineto{\pgfqpoint{2.163152in}{2.770573in}}%
\pgfpathlineto{\pgfqpoint{2.163293in}{2.717078in}}%
\pgfpathlineto{\pgfqpoint{2.163238in}{2.775968in}}%
\pgfpathlineto{\pgfqpoint{2.163307in}{2.729839in}}%
\pgfpathlineto{\pgfqpoint{2.164093in}{2.644559in}}%
\pgfpathlineto{\pgfqpoint{2.166210in}{2.644559in}}%
\pgfpathlineto{\pgfqpoint{2.166215in}{2.749839in}}%
\pgfpathlineto{\pgfqpoint{2.166955in}{2.692011in}}%
\pgfpathlineto{\pgfqpoint{2.168092in}{2.770160in}}%
\pgfpathlineto{\pgfqpoint{2.168238in}{2.716938in}}%
\pgfpathlineto{\pgfqpoint{2.168184in}{2.775898in}}%
\pgfpathlineto{\pgfqpoint{2.168252in}{2.729950in}}%
\pgfpathlineto{\pgfqpoint{2.169041in}{2.644559in}}%
\pgfpathlineto{\pgfqpoint{2.171155in}{2.644559in}}%
\pgfpathlineto{\pgfqpoint{2.171161in}{2.749839in}}%
\pgfpathlineto{\pgfqpoint{2.171881in}{2.697907in}}%
\pgfpathlineto{\pgfqpoint{2.171931in}{2.693547in}}%
\pgfpathlineto{\pgfqpoint{2.172277in}{2.722765in}}%
\pgfpathlineto{\pgfqpoint{2.172326in}{2.718398in}}%
\pgfpathlineto{\pgfqpoint{2.173129in}{2.776304in}}%
\pgfpathlineto{\pgfqpoint{2.173153in}{2.757366in}}%
\pgfpathlineto{\pgfqpoint{2.173962in}{2.644559in}}%
\pgfpathlineto{\pgfqpoint{2.176101in}{2.644559in}}%
\pgfpathlineto{\pgfqpoint{2.176106in}{2.749842in}}%
\pgfpathlineto{\pgfqpoint{2.176840in}{2.691645in}}%
\pgfpathlineto{\pgfqpoint{2.177978in}{2.769857in}}%
\pgfpathlineto{\pgfqpoint{2.178129in}{2.716656in}}%
\pgfpathlineto{\pgfqpoint{2.178074in}{2.775931in}}%
\pgfpathlineto{\pgfqpoint{2.178143in}{2.730263in}}%
\pgfpathlineto{\pgfqpoint{2.178930in}{2.644559in}}%
\pgfpathlineto{\pgfqpoint{2.181046in}{2.644559in}}%
\pgfpathlineto{\pgfqpoint{2.181051in}{2.749840in}}%
\pgfpathlineto{\pgfqpoint{2.181797in}{2.692354in}}%
\pgfpathlineto{\pgfqpoint{2.182935in}{2.770622in}}%
\pgfpathlineto{\pgfqpoint{2.183075in}{2.717098in}}%
\pgfpathlineto{\pgfqpoint{2.183020in}{2.775966in}}%
\pgfpathlineto{\pgfqpoint{2.183089in}{2.729796in}}%
\pgfpathlineto{\pgfqpoint{2.183877in}{2.644559in}}%
\pgfpathlineto{\pgfqpoint{2.185992in}{2.644559in}}%
\pgfpathlineto{\pgfqpoint{2.185997in}{2.749838in}}%
\pgfpathlineto{\pgfqpoint{2.186754in}{2.692958in}}%
\pgfpathlineto{\pgfqpoint{2.187891in}{2.771376in}}%
\pgfpathlineto{\pgfqpoint{2.188020in}{2.716019in}}%
\pgfpathlineto{\pgfqpoint{2.187965in}{2.776024in}}%
\pgfpathlineto{\pgfqpoint{2.188034in}{2.730947in}}%
\pgfpathlineto{\pgfqpoint{2.188821in}{2.644559in}}%
\pgfpathlineto{\pgfqpoint{2.190937in}{2.644559in}}%
\pgfpathlineto{\pgfqpoint{2.190942in}{2.749844in}}%
\pgfpathlineto{\pgfqpoint{2.191683in}{2.692106in}}%
\pgfpathlineto{\pgfqpoint{2.192821in}{2.770267in}}%
\pgfpathlineto{\pgfqpoint{2.192966in}{2.716994in}}%
\pgfpathlineto{\pgfqpoint{2.192911in}{2.775916in}}%
\pgfpathlineto{\pgfqpoint{2.192979in}{2.729916in}}%
\pgfpathlineto{\pgfqpoint{2.193766in}{2.644559in}}%
\pgfpathlineto{\pgfqpoint{2.195882in}{2.644559in}}%
\pgfpathlineto{\pgfqpoint{2.195888in}{2.749838in}}%
\pgfpathlineto{\pgfqpoint{2.196642in}{2.693010in}}%
\pgfpathlineto{\pgfqpoint{2.197780in}{2.771004in}}%
\pgfpathlineto{\pgfqpoint{2.197911in}{2.716523in}}%
\pgfpathlineto{\pgfqpoint{2.197856in}{2.775809in}}%
\pgfpathlineto{\pgfqpoint{2.197925in}{2.730382in}}%
\pgfpathlineto{\pgfqpoint{2.198713in}{2.644559in}}%
\pgfpathlineto{\pgfqpoint{2.200828in}{2.644559in}}%
\pgfpathlineto{\pgfqpoint{2.200833in}{2.749832in}}%
\pgfpathlineto{\pgfqpoint{2.201571in}{2.691872in}}%
\pgfpathlineto{\pgfqpoint{2.202709in}{2.770047in}}%
\pgfpathlineto{\pgfqpoint{2.202856in}{2.716844in}}%
\pgfpathlineto{\pgfqpoint{2.202802in}{2.775899in}}%
\pgfpathlineto{\pgfqpoint{2.202870in}{2.730029in}}%
\pgfpathlineto{\pgfqpoint{2.203659in}{2.644559in}}%
\pgfpathlineto{\pgfqpoint{2.205773in}{2.644559in}}%
\pgfpathlineto{\pgfqpoint{2.205779in}{2.749836in}}%
\pgfpathlineto{\pgfqpoint{2.206521in}{2.692140in}}%
\pgfpathlineto{\pgfqpoint{2.207658in}{2.770326in}}%
\pgfpathlineto{\pgfqpoint{2.207802in}{2.717007in}}%
\pgfpathlineto{\pgfqpoint{2.207747in}{2.775917in}}%
\pgfpathlineto{\pgfqpoint{2.207816in}{2.729878in}}%
\pgfpathlineto{\pgfqpoint{2.208604in}{2.644559in}}%
\pgfpathlineto{\pgfqpoint{2.210719in}{2.644559in}}%
\pgfpathlineto{\pgfqpoint{2.210724in}{2.749835in}}%
\pgfpathlineto{\pgfqpoint{2.211480in}{2.693045in}}%
\pgfpathlineto{\pgfqpoint{2.212618in}{2.771145in}}%
\pgfpathlineto{\pgfqpoint{2.212747in}{2.716551in}}%
\pgfpathlineto{\pgfqpoint{2.212693in}{2.775861in}}%
\pgfpathlineto{\pgfqpoint{2.212761in}{2.730350in}}%
\pgfpathlineto{\pgfqpoint{2.213549in}{2.644559in}}%
\pgfpathlineto{\pgfqpoint{2.215664in}{2.644559in}}%
\pgfpathlineto{\pgfqpoint{2.215670in}{2.749837in}}%
\pgfpathlineto{\pgfqpoint{2.216426in}{2.693055in}}%
\pgfpathlineto{\pgfqpoint{2.217563in}{2.771169in}}%
\pgfpathlineto{\pgfqpoint{2.217693in}{2.716563in}}%
\pgfpathlineto{\pgfqpoint{2.217638in}{2.775873in}}%
\pgfpathlineto{\pgfqpoint{2.217707in}{2.730346in}}%
\pgfpathlineto{\pgfqpoint{2.218492in}{2.644559in}}%
\pgfpathlineto{\pgfqpoint{2.220610in}{2.644559in}}%
\pgfpathlineto{\pgfqpoint{2.220615in}{2.749834in}}%
\pgfpathlineto{\pgfqpoint{2.221368in}{2.692924in}}%
\pgfpathlineto{\pgfqpoint{2.222506in}{2.770908in}}%
\pgfpathlineto{\pgfqpoint{2.222638in}{2.716452in}}%
\pgfpathlineto{\pgfqpoint{2.222584in}{2.775798in}}%
\pgfpathlineto{\pgfqpoint{2.222652in}{2.730447in}}%
\pgfpathlineto{\pgfqpoint{2.223438in}{2.644559in}}%
\pgfpathlineto{\pgfqpoint{2.225555in}{2.644559in}}%
\pgfpathlineto{\pgfqpoint{2.225561in}{2.749838in}}%
\pgfpathlineto{\pgfqpoint{2.226315in}{2.693017in}}%
\pgfpathlineto{\pgfqpoint{2.227453in}{2.771031in}}%
\pgfpathlineto{\pgfqpoint{2.227584in}{2.716544in}}%
\pgfpathlineto{\pgfqpoint{2.227529in}{2.775821in}}%
\pgfpathlineto{\pgfqpoint{2.227598in}{2.730364in}}%
\pgfpathlineto{\pgfqpoint{2.228386in}{2.644559in}}%
\pgfpathlineto{\pgfqpoint{2.230501in}{2.644559in}}%
\pgfpathlineto{\pgfqpoint{2.230506in}{2.749832in}}%
\pgfpathlineto{\pgfqpoint{2.231259in}{2.692881in}}%
\pgfpathlineto{\pgfqpoint{2.232396in}{2.770865in}}%
\pgfpathlineto{\pgfqpoint{2.232529in}{2.716423in}}%
\pgfpathlineto{\pgfqpoint{2.232474in}{2.775795in}}%
\pgfpathlineto{\pgfqpoint{2.232543in}{2.730473in}}%
\pgfpathlineto{\pgfqpoint{2.233330in}{2.644559in}}%
\pgfpathlineto{\pgfqpoint{2.235446in}{2.644559in}}%
\pgfpathlineto{\pgfqpoint{2.235451in}{2.749835in}}%
\pgfpathlineto{\pgfqpoint{2.236214in}{2.699959in}}%
\pgfpathlineto{\pgfqpoint{2.236263in}{2.696646in}}%
\pgfpathlineto{\pgfqpoint{2.236609in}{2.724818in}}%
\pgfpathlineto{\pgfqpoint{2.236659in}{2.721496in}}%
\pgfpathlineto{\pgfqpoint{2.237386in}{2.773584in}}%
\pgfpathlineto{\pgfqpoint{2.237462in}{2.738931in}}%
\pgfpathlineto{\pgfqpoint{2.238280in}{2.644559in}}%
\pgfpathlineto{\pgfqpoint{2.240392in}{2.644559in}}%
\pgfpathlineto{\pgfqpoint{2.240397in}{2.749829in}}%
\pgfpathlineto{\pgfqpoint{2.241161in}{2.700103in}}%
\pgfpathlineto{\pgfqpoint{2.241211in}{2.696758in}}%
\pgfpathlineto{\pgfqpoint{2.241557in}{2.724962in}}%
\pgfpathlineto{\pgfqpoint{2.241606in}{2.721608in}}%
\pgfpathlineto{\pgfqpoint{2.242332in}{2.773660in}}%
\pgfpathlineto{\pgfqpoint{2.242408in}{2.738778in}}%
\pgfpathlineto{\pgfqpoint{2.243227in}{2.644559in}}%
\pgfpathlineto{\pgfqpoint{2.245337in}{2.644559in}}%
\pgfpathlineto{\pgfqpoint{2.245342in}{2.749830in}}%
\pgfpathlineto{\pgfqpoint{2.246096in}{2.699323in}}%
\pgfpathlineto{\pgfqpoint{2.246145in}{2.696139in}}%
\pgfpathlineto{\pgfqpoint{2.246392in}{2.717967in}}%
\pgfpathlineto{\pgfqpoint{2.246442in}{2.714777in}}%
\pgfpathlineto{\pgfqpoint{2.247272in}{2.773227in}}%
\pgfpathlineto{\pgfqpoint{2.247335in}{2.757823in}}%
\pgfpathlineto{\pgfqpoint{2.248139in}{2.644559in}}%
\pgfpathlineto{\pgfqpoint{2.250282in}{2.644559in}}%
\pgfpathlineto{\pgfqpoint{2.250288in}{2.749832in}}%
\pgfpathlineto{\pgfqpoint{2.251052in}{2.700087in}}%
\pgfpathlineto{\pgfqpoint{2.251101in}{2.696753in}}%
\pgfpathlineto{\pgfqpoint{2.251448in}{2.724946in}}%
\pgfpathlineto{\pgfqpoint{2.251497in}{2.721603in}}%
\pgfpathlineto{\pgfqpoint{2.252223in}{2.773654in}}%
\pgfpathlineto{\pgfqpoint{2.252299in}{2.738784in}}%
\pgfpathlineto{\pgfqpoint{2.253119in}{2.644559in}}%
\pgfpathlineto{\pgfqpoint{2.255228in}{2.644559in}}%
\pgfpathlineto{\pgfqpoint{2.255233in}{2.749826in}}%
\pgfpathlineto{\pgfqpoint{2.255984in}{2.699262in}}%
\pgfpathlineto{\pgfqpoint{2.256033in}{2.695847in}}%
\pgfpathlineto{\pgfqpoint{2.256380in}{2.724121in}}%
\pgfpathlineto{\pgfqpoint{2.256429in}{2.720698in}}%
\pgfpathlineto{\pgfqpoint{2.257162in}{2.773252in}}%
\pgfpathlineto{\pgfqpoint{2.257226in}{2.757946in}}%
\pgfpathlineto{\pgfqpoint{2.258029in}{2.644559in}}%
\pgfpathlineto{\pgfqpoint{2.260173in}{2.644559in}}%
\pgfpathlineto{\pgfqpoint{2.260179in}{2.749829in}}%
\pgfpathlineto{\pgfqpoint{2.260901in}{2.691003in}}%
\pgfpathlineto{\pgfqpoint{2.262039in}{2.768941in}}%
\pgfpathlineto{\pgfqpoint{2.262203in}{2.712690in}}%
\pgfpathlineto{\pgfqpoint{2.262118in}{2.773901in}}%
\pgfpathlineto{\pgfqpoint{2.262218in}{2.734904in}}%
\pgfpathlineto{\pgfqpoint{2.263004in}{2.644559in}}%
\pgfpathlineto{\pgfqpoint{2.265119in}{2.644559in}}%
\pgfpathlineto{\pgfqpoint{2.265124in}{2.749823in}}%
\pgfpathlineto{\pgfqpoint{2.265879in}{2.699505in}}%
\pgfpathlineto{\pgfqpoint{2.265928in}{2.696105in}}%
\pgfpathlineto{\pgfqpoint{2.266274in}{2.724364in}}%
\pgfpathlineto{\pgfqpoint{2.266324in}{2.720955in}}%
\pgfpathlineto{\pgfqpoint{2.267054in}{2.773369in}}%
\pgfpathlineto{\pgfqpoint{2.267135in}{2.739854in}}%
\pgfpathlineto{\pgfqpoint{2.267952in}{2.644559in}}%
\pgfpathlineto{\pgfqpoint{2.270064in}{2.644559in}}%
\pgfpathlineto{\pgfqpoint{2.270070in}{2.749820in}}%
\pgfpathlineto{\pgfqpoint{2.270789in}{2.690856in}}%
\pgfpathlineto{\pgfqpoint{2.271927in}{2.768711in}}%
\pgfpathlineto{\pgfqpoint{2.272094in}{2.712599in}}%
\pgfpathlineto{\pgfqpoint{2.272007in}{2.773757in}}%
\pgfpathlineto{\pgfqpoint{2.272109in}{2.734986in}}%
\pgfpathlineto{\pgfqpoint{2.272896in}{2.644559in}}%
\pgfpathlineto{\pgfqpoint{2.275010in}{2.644559in}}%
\pgfpathlineto{\pgfqpoint{2.275015in}{2.749825in}}%
\pgfpathlineto{\pgfqpoint{2.275767in}{2.699346in}}%
\pgfpathlineto{\pgfqpoint{2.275817in}{2.695935in}}%
\pgfpathlineto{\pgfqpoint{2.276163in}{2.724205in}}%
\pgfpathlineto{\pgfqpoint{2.276212in}{2.720786in}}%
\pgfpathlineto{\pgfqpoint{2.276944in}{2.773293in}}%
\pgfpathlineto{\pgfqpoint{2.277008in}{2.757927in}}%
\pgfpathlineto{\pgfqpoint{2.277810in}{2.644559in}}%
\pgfpathlineto{\pgfqpoint{2.279955in}{2.644559in}}%
\pgfpathlineto{\pgfqpoint{2.279961in}{2.749825in}}%
\pgfpathlineto{\pgfqpoint{2.280729in}{2.700277in}}%
\pgfpathlineto{\pgfqpoint{2.280778in}{2.697026in}}%
\pgfpathlineto{\pgfqpoint{2.281025in}{2.718921in}}%
\pgfpathlineto{\pgfqpoint{2.281075in}{2.715663in}}%
\pgfpathlineto{\pgfqpoint{2.281898in}{2.773718in}}%
\pgfpathlineto{\pgfqpoint{2.281954in}{2.757995in}}%
\pgfpathlineto{\pgfqpoint{2.282756in}{2.644559in}}%
\pgfpathlineto{\pgfqpoint{2.284901in}{2.644559in}}%
\pgfpathlineto{\pgfqpoint{2.284906in}{2.749823in}}%
\pgfpathlineto{\pgfqpoint{2.285674in}{2.700225in}}%
\pgfpathlineto{\pgfqpoint{2.285723in}{2.696993in}}%
\pgfpathlineto{\pgfqpoint{2.285970in}{2.718869in}}%
\pgfpathlineto{\pgfqpoint{2.286020in}{2.715631in}}%
\pgfpathlineto{\pgfqpoint{2.286843in}{2.773684in}}%
\pgfpathlineto{\pgfqpoint{2.286899in}{2.757992in}}%
\pgfpathlineto{\pgfqpoint{2.287699in}{2.644559in}}%
\pgfpathlineto{\pgfqpoint{2.289846in}{2.644559in}}%
\pgfpathlineto{\pgfqpoint{2.289851in}{2.749818in}}%
\pgfpathlineto{\pgfqpoint{2.290571in}{2.690846in}}%
\pgfpathlineto{\pgfqpoint{2.291709in}{2.768700in}}%
\pgfpathlineto{\pgfqpoint{2.291876in}{2.712626in}}%
\pgfpathlineto{\pgfqpoint{2.291789in}{2.773751in}}%
\pgfpathlineto{\pgfqpoint{2.291891in}{2.734951in}}%
\pgfpathlineto{\pgfqpoint{2.292679in}{2.644559in}}%
\pgfpathlineto{\pgfqpoint{2.294792in}{2.644559in}}%
\pgfpathlineto{\pgfqpoint{2.294797in}{2.749820in}}%
\pgfpathlineto{\pgfqpoint{2.295535in}{2.698133in}}%
\pgfpathlineto{\pgfqpoint{2.295584in}{2.695326in}}%
\pgfpathlineto{\pgfqpoint{2.295831in}{2.716777in}}%
\pgfpathlineto{\pgfqpoint{2.295881in}{2.713964in}}%
\pgfpathlineto{\pgfqpoint{2.296719in}{2.772527in}}%
\pgfpathlineto{\pgfqpoint{2.296790in}{2.757876in}}%
\pgfpathlineto{\pgfqpoint{2.297592in}{2.644559in}}%
\pgfpathlineto{\pgfqpoint{2.299737in}{2.644559in}}%
\pgfpathlineto{\pgfqpoint{2.299742in}{2.749826in}}%
\pgfpathlineto{\pgfqpoint{2.300508in}{2.700047in}}%
\pgfpathlineto{\pgfqpoint{2.300557in}{2.696894in}}%
\pgfpathlineto{\pgfqpoint{2.300804in}{2.718691in}}%
\pgfpathlineto{\pgfqpoint{2.300854in}{2.715532in}}%
\pgfpathlineto{\pgfqpoint{2.301678in}{2.773577in}}%
\pgfpathlineto{\pgfqpoint{2.301736in}{2.758004in}}%
\pgfpathlineto{\pgfqpoint{2.302537in}{2.644559in}}%
\pgfpathlineto{\pgfqpoint{2.304682in}{2.644559in}}%
\pgfpathlineto{\pgfqpoint{2.304688in}{2.749815in}}%
\pgfpathlineto{\pgfqpoint{2.305415in}{2.697453in}}%
\pgfpathlineto{\pgfqpoint{2.305464in}{2.694616in}}%
\pgfpathlineto{\pgfqpoint{2.305711in}{2.716097in}}%
\pgfpathlineto{\pgfqpoint{2.305761in}{2.713254in}}%
\pgfpathlineto{\pgfqpoint{2.306656in}{2.775462in}}%
\pgfpathlineto{\pgfqpoint{2.306657in}{2.748614in}}%
\pgfpathlineto{\pgfqpoint{2.306680in}{2.757479in}}%
\pgfpathlineto{\pgfqpoint{2.307427in}{2.644559in}}%
\pgfpathlineto{\pgfqpoint{2.309628in}{2.644559in}}%
\pgfpathlineto{\pgfqpoint{2.309633in}{2.749826in}}%
\pgfpathlineto{\pgfqpoint{2.310353in}{2.690852in}}%
\pgfpathlineto{\pgfqpoint{2.311490in}{2.768702in}}%
\pgfpathlineto{\pgfqpoint{2.311658in}{2.712635in}}%
\pgfpathlineto{\pgfqpoint{2.311571in}{2.773758in}}%
\pgfpathlineto{\pgfqpoint{2.311672in}{2.734969in}}%
\pgfpathlineto{\pgfqpoint{2.312459in}{2.644559in}}%
\pgfpathlineto{\pgfqpoint{2.314573in}{2.644559in}}%
\pgfpathlineto{\pgfqpoint{2.314579in}{2.749811in}}%
\pgfpathlineto{\pgfqpoint{2.315309in}{2.697668in}}%
\pgfpathlineto{\pgfqpoint{2.315359in}{2.694876in}}%
\pgfpathlineto{\pgfqpoint{2.315606in}{2.716312in}}%
\pgfpathlineto{\pgfqpoint{2.315656in}{2.713514in}}%
\pgfpathlineto{\pgfqpoint{2.316547in}{2.775436in}}%
\pgfpathlineto{\pgfqpoint{2.316571in}{2.757438in}}%
\pgfpathlineto{\pgfqpoint{2.317378in}{2.644559in}}%
\pgfpathlineto{\pgfqpoint{2.319519in}{2.644559in}}%
\pgfpathlineto{\pgfqpoint{2.319524in}{2.749824in}}%
\pgfpathlineto{\pgfqpoint{2.320292in}{2.700268in}}%
\pgfpathlineto{\pgfqpoint{2.320342in}{2.697019in}}%
\pgfpathlineto{\pgfqpoint{2.320589in}{2.718912in}}%
\pgfpathlineto{\pgfqpoint{2.320638in}{2.715657in}}%
\pgfpathlineto{\pgfqpoint{2.321461in}{2.773711in}}%
\pgfpathlineto{\pgfqpoint{2.321517in}{2.757992in}}%
\pgfpathlineto{\pgfqpoint{2.322320in}{2.644559in}}%
\pgfpathlineto{\pgfqpoint{2.324464in}{2.644559in}}%
\pgfpathlineto{\pgfqpoint{2.324470in}{2.749811in}}%
\pgfpathlineto{\pgfqpoint{2.325192in}{2.697194in}}%
\pgfpathlineto{\pgfqpoint{2.325241in}{2.694304in}}%
\pgfpathlineto{\pgfqpoint{2.325489in}{2.715838in}}%
\pgfpathlineto{\pgfqpoint{2.325538in}{2.712942in}}%
\pgfpathlineto{\pgfqpoint{2.326438in}{2.775484in}}%
\pgfpathlineto{\pgfqpoint{2.326438in}{2.748612in}}%
\pgfpathlineto{\pgfqpoint{2.326462in}{2.757511in}}%
\pgfpathlineto{\pgfqpoint{2.327210in}{2.644559in}}%
\pgfpathlineto{\pgfqpoint{2.329410in}{2.644559in}}%
\pgfpathlineto{\pgfqpoint{2.329415in}{2.749818in}}%
\pgfpathlineto{\pgfqpoint{2.330155in}{2.698280in}}%
\pgfpathlineto{\pgfqpoint{2.330205in}{2.695454in}}%
\pgfpathlineto{\pgfqpoint{2.330452in}{2.716924in}}%
\pgfpathlineto{\pgfqpoint{2.330501in}{2.714092in}}%
\pgfpathlineto{\pgfqpoint{2.331338in}{2.772605in}}%
\pgfpathlineto{\pgfqpoint{2.331408in}{2.757860in}}%
\pgfpathlineto{\pgfqpoint{2.332210in}{2.644559in}}%
\pgfpathlineto{\pgfqpoint{2.334355in}{2.644559in}}%
\pgfpathlineto{\pgfqpoint{2.334360in}{2.749808in}}%
\pgfpathlineto{\pgfqpoint{2.335118in}{2.693488in}}%
\pgfpathlineto{\pgfqpoint{2.336329in}{2.775343in}}%
\pgfpathlineto{\pgfqpoint{2.337099in}{2.644559in}}%
\pgfpathlineto{\pgfqpoint{2.339301in}{2.644559in}}%
\pgfpathlineto{\pgfqpoint{2.339306in}{2.749819in}}%
\pgfpathlineto{\pgfqpoint{2.340049in}{2.698456in}}%
\pgfpathlineto{\pgfqpoint{2.340098in}{2.695617in}}%
\pgfpathlineto{\pgfqpoint{2.340345in}{2.717101in}}%
\pgfpathlineto{\pgfqpoint{2.340395in}{2.714254in}}%
\pgfpathlineto{\pgfqpoint{2.341230in}{2.772698in}}%
\pgfpathlineto{\pgfqpoint{2.341299in}{2.757846in}}%
\pgfpathlineto{\pgfqpoint{2.342102in}{2.644559in}}%
\pgfpathlineto{\pgfqpoint{2.344246in}{2.644559in}}%
\pgfpathlineto{\pgfqpoint{2.344251in}{2.749809in}}%
\pgfpathlineto{\pgfqpoint{2.345000in}{2.698714in}}%
\pgfpathlineto{\pgfqpoint{2.345049in}{2.696000in}}%
\pgfpathlineto{\pgfqpoint{2.345296in}{2.717358in}}%
\pgfpathlineto{\pgfqpoint{2.345346in}{2.714638in}}%
\pgfpathlineto{\pgfqpoint{2.346178in}{2.772784in}}%
\pgfpathlineto{\pgfqpoint{2.346245in}{2.757998in}}%
\pgfpathlineto{\pgfqpoint{2.347045in}{2.644559in}}%
\pgfpathlineto{\pgfqpoint{2.349192in}{2.644559in}}%
\pgfpathlineto{\pgfqpoint{2.349197in}{2.749801in}}%
\pgfpathlineto{\pgfqpoint{2.349958in}{2.693751in}}%
\pgfpathlineto{\pgfqpoint{2.351165in}{2.775349in}}%
\pgfpathlineto{\pgfqpoint{2.351936in}{2.644559in}}%
\pgfpathlineto{\pgfqpoint{2.354137in}{2.644559in}}%
\pgfpathlineto{\pgfqpoint{2.354142in}{2.749820in}}%
\pgfpathlineto{\pgfqpoint{2.354881in}{2.698164in}}%
\pgfpathlineto{\pgfqpoint{2.354930in}{2.695351in}}%
\pgfpathlineto{\pgfqpoint{2.355177in}{2.716808in}}%
\pgfpathlineto{\pgfqpoint{2.355227in}{2.713989in}}%
\pgfpathlineto{\pgfqpoint{2.356064in}{2.772546in}}%
\pgfpathlineto{\pgfqpoint{2.356135in}{2.757875in}}%
\pgfpathlineto{\pgfqpoint{2.356937in}{2.644559in}}%
\pgfpathlineto{\pgfqpoint{2.359082in}{2.644559in}}%
\pgfpathlineto{\pgfqpoint{2.359088in}{2.749800in}}%
\pgfpathlineto{\pgfqpoint{2.359851in}{2.693946in}}%
\pgfpathlineto{\pgfqpoint{2.361056in}{2.775261in}}%
\pgfpathlineto{\pgfqpoint{2.361828in}{2.644559in}}%
\pgfpathlineto{\pgfqpoint{2.364028in}{2.644559in}}%
\pgfpathlineto{\pgfqpoint{2.364033in}{2.749818in}}%
\pgfpathlineto{\pgfqpoint{2.364776in}{2.698432in}}%
\pgfpathlineto{\pgfqpoint{2.364825in}{2.695593in}}%
\pgfpathlineto{\pgfqpoint{2.365072in}{2.717076in}}%
\pgfpathlineto{\pgfqpoint{2.365122in}{2.714230in}}%
\pgfpathlineto{\pgfqpoint{2.365957in}{2.772684in}}%
\pgfpathlineto{\pgfqpoint{2.366026in}{2.757846in}}%
\pgfpathlineto{\pgfqpoint{2.366828in}{2.644559in}}%
\pgfpathlineto{\pgfqpoint{2.368973in}{2.644559in}}%
\pgfpathlineto{\pgfqpoint{2.368979in}{2.749797in}}%
\pgfpathlineto{\pgfqpoint{2.369702in}{2.697124in}}%
\pgfpathlineto{\pgfqpoint{2.369751in}{2.694424in}}%
\pgfpathlineto{\pgfqpoint{2.369999in}{2.715768in}}%
\pgfpathlineto{\pgfqpoint{2.370048in}{2.713062in}}%
\pgfpathlineto{\pgfqpoint{2.370947in}{2.775358in}}%
\pgfpathlineto{\pgfqpoint{2.370947in}{2.748584in}}%
\pgfpathlineto{\pgfqpoint{2.370971in}{2.757625in}}%
\pgfpathlineto{\pgfqpoint{2.371717in}{2.644559in}}%
\pgfpathlineto{\pgfqpoint{2.373919in}{2.644559in}}%
\pgfpathlineto{\pgfqpoint{2.373924in}{2.749816in}}%
\pgfpathlineto{\pgfqpoint{2.374659in}{2.697909in}}%
\pgfpathlineto{\pgfqpoint{2.374708in}{2.695140in}}%
\pgfpathlineto{\pgfqpoint{2.374955in}{2.716553in}}%
\pgfpathlineto{\pgfqpoint{2.375005in}{2.713778in}}%
\pgfpathlineto{\pgfqpoint{2.375846in}{2.772483in}}%
\pgfpathlineto{\pgfqpoint{2.375917in}{2.757883in}}%
\pgfpathlineto{\pgfqpoint{2.376721in}{2.644559in}}%
\pgfpathlineto{\pgfqpoint{2.378864in}{2.644559in}}%
\pgfpathlineto{\pgfqpoint{2.378870in}{2.749803in}}%
\pgfpathlineto{\pgfqpoint{2.379633in}{2.693977in}}%
\pgfpathlineto{\pgfqpoint{2.380838in}{2.775273in}}%
\pgfpathlineto{\pgfqpoint{2.381608in}{2.644559in}}%
\pgfpathlineto{\pgfqpoint{2.383810in}{2.644559in}}%
\pgfpathlineto{\pgfqpoint{2.383815in}{2.749809in}}%
\pgfpathlineto{\pgfqpoint{2.384567in}{2.693138in}}%
\pgfpathlineto{\pgfqpoint{2.385784in}{2.775367in}}%
\pgfpathlineto{\pgfqpoint{2.386556in}{2.644559in}}%
\pgfpathlineto{\pgfqpoint{2.388755in}{2.644559in}}%
\pgfpathlineto{\pgfqpoint{2.388760in}{2.749805in}}%
\pgfpathlineto{\pgfqpoint{2.389518in}{2.693487in}}%
\pgfpathlineto{\pgfqpoint{2.390729in}{2.775342in}}%
\pgfpathlineto{\pgfqpoint{2.391501in}{2.644559in}}%
\pgfpathlineto{\pgfqpoint{2.393701in}{2.644559in}}%
\pgfpathlineto{\pgfqpoint{2.393706in}{2.749806in}}%
\pgfpathlineto{\pgfqpoint{2.394462in}{2.693433in}}%
\pgfpathlineto{\pgfqpoint{2.395674in}{2.775340in}}%
\pgfpathlineto{\pgfqpoint{2.396445in}{2.644559in}}%
\pgfpathlineto{\pgfqpoint{2.398646in}{2.644559in}}%
\pgfpathlineto{\pgfqpoint{2.398651in}{2.749804in}}%
\pgfpathlineto{\pgfqpoint{2.399372in}{2.696979in}}%
\pgfpathlineto{\pgfqpoint{2.399421in}{2.694253in}}%
\pgfpathlineto{\pgfqpoint{2.399669in}{2.715623in}}%
\pgfpathlineto{\pgfqpoint{2.399718in}{2.712891in}}%
\pgfpathlineto{\pgfqpoint{2.400620in}{2.775384in}}%
\pgfpathlineto{\pgfqpoint{2.400620in}{2.748598in}}%
\pgfpathlineto{\pgfqpoint{2.400644in}{2.757681in}}%
\pgfpathlineto{\pgfqpoint{2.401391in}{2.644559in}}%
\pgfpathlineto{\pgfqpoint{2.403592in}{2.644559in}}%
\pgfpathlineto{\pgfqpoint{2.403597in}{2.749808in}}%
\pgfpathlineto{\pgfqpoint{2.404360in}{2.693846in}}%
\pgfpathlineto{\pgfqpoint{2.405565in}{2.775361in}}%
\pgfpathlineto{\pgfqpoint{2.406336in}{2.644559in}}%
\pgfpathlineto{\pgfqpoint{2.408537in}{2.644559in}}%
\pgfpathlineto{\pgfqpoint{2.408542in}{2.749804in}}%
\pgfpathlineto{\pgfqpoint{2.409308in}{2.694096in}}%
\pgfpathlineto{\pgfqpoint{2.410511in}{2.775274in}}%
\pgfpathlineto{\pgfqpoint{2.411280in}{2.644559in}}%
\pgfpathlineto{\pgfqpoint{2.413482in}{2.644559in}}%
\pgfpathlineto{\pgfqpoint{2.413488in}{2.749804in}}%
\pgfpathlineto{\pgfqpoint{2.414211in}{2.697120in}}%
\pgfpathlineto{\pgfqpoint{2.414260in}{2.694416in}}%
\pgfpathlineto{\pgfqpoint{2.414507in}{2.715765in}}%
\pgfpathlineto{\pgfqpoint{2.414557in}{2.713053in}}%
\pgfpathlineto{\pgfqpoint{2.415456in}{2.775376in}}%
\pgfpathlineto{\pgfqpoint{2.415457in}{2.748599in}}%
\pgfpathlineto{\pgfqpoint{2.415480in}{2.757643in}}%
\pgfpathlineto{\pgfqpoint{2.416228in}{2.644559in}}%
\pgfpathlineto{\pgfqpoint{2.418428in}{2.644559in}}%
\pgfpathlineto{\pgfqpoint{2.418433in}{2.749801in}}%
\pgfpathlineto{\pgfqpoint{2.419198in}{2.694004in}}%
\pgfpathlineto{\pgfqpoint{2.420402in}{2.775266in}}%
\pgfpathlineto{\pgfqpoint{2.421173in}{2.644559in}}%
\pgfpathlineto{\pgfqpoint{2.423373in}{2.644559in}}%
\pgfpathlineto{\pgfqpoint{2.423379in}{2.749806in}}%
\pgfpathlineto{\pgfqpoint{2.424103in}{2.697211in}}%
\pgfpathlineto{\pgfqpoint{2.424153in}{2.694520in}}%
\pgfpathlineto{\pgfqpoint{2.424400in}{2.715855in}}%
\pgfpathlineto{\pgfqpoint{2.424449in}{2.713158in}}%
\pgfpathlineto{\pgfqpoint{2.425347in}{2.775372in}}%
\pgfpathlineto{\pgfqpoint{2.425347in}{2.748602in}}%
\pgfpathlineto{\pgfqpoint{2.425371in}{2.757623in}}%
\pgfpathlineto{\pgfqpoint{2.426118in}{2.644559in}}%
\pgfpathlineto{\pgfqpoint{2.428319in}{2.644559in}}%
\pgfpathlineto{\pgfqpoint{2.428324in}{2.749799in}}%
\pgfpathlineto{\pgfqpoint{2.429083in}{2.693661in}}%
\pgfpathlineto{\pgfqpoint{2.430293in}{2.775253in}}%
\pgfpathlineto{\pgfqpoint{2.431063in}{2.644559in}}%
\pgfpathlineto{\pgfqpoint{2.433264in}{2.644559in}}%
\pgfpathlineto{\pgfqpoint{2.433270in}{2.749802in}}%
\pgfpathlineto{\pgfqpoint{2.434036in}{2.694160in}}%
\pgfpathlineto{\pgfqpoint{2.435238in}{2.775269in}}%
\pgfpathlineto{\pgfqpoint{2.436009in}{2.644559in}}%
\pgfpathlineto{\pgfqpoint{2.438210in}{2.644559in}}%
\pgfpathlineto{\pgfqpoint{2.438215in}{2.749797in}}%
\pgfpathlineto{\pgfqpoint{2.438937in}{2.697075in}}%
\pgfpathlineto{\pgfqpoint{2.438987in}{2.694364in}}%
\pgfpathlineto{\pgfqpoint{2.439234in}{2.715719in}}%
\pgfpathlineto{\pgfqpoint{2.439284in}{2.713002in}}%
\pgfpathlineto{\pgfqpoint{2.440184in}{2.775366in}}%
\pgfpathlineto{\pgfqpoint{2.440184in}{2.748587in}}%
\pgfpathlineto{\pgfqpoint{2.440208in}{2.757640in}}%
\pgfpathlineto{\pgfqpoint{2.440955in}{2.644559in}}%
\pgfpathlineto{\pgfqpoint{2.443155in}{2.644559in}}%
\pgfpathlineto{\pgfqpoint{2.443161in}{2.749795in}}%
\pgfpathlineto{\pgfqpoint{2.443910in}{2.692918in}}%
\pgfpathlineto{\pgfqpoint{2.445129in}{2.775378in}}%
\pgfpathlineto{\pgfqpoint{2.445900in}{2.644559in}}%
\pgfpathlineto{\pgfqpoint{2.448101in}{2.644559in}}%
\pgfpathlineto{\pgfqpoint{2.448106in}{2.749802in}}%
\pgfpathlineto{\pgfqpoint{2.448870in}{2.693999in}}%
\pgfpathlineto{\pgfqpoint{2.450074in}{2.775271in}}%
\pgfpathlineto{\pgfqpoint{2.450846in}{2.644559in}}%
\pgfpathlineto{\pgfqpoint{2.453046in}{2.644559in}}%
\pgfpathlineto{\pgfqpoint{2.453051in}{2.749798in}}%
\pgfpathlineto{\pgfqpoint{2.453814in}{2.699729in}}%
\pgfpathlineto{\pgfqpoint{2.453863in}{2.696747in}}%
\pgfpathlineto{\pgfqpoint{2.454111in}{2.718372in}}%
\pgfpathlineto{\pgfqpoint{2.454160in}{2.715385in}}%
\pgfpathlineto{\pgfqpoint{2.454986in}{2.773347in}}%
\pgfpathlineto{\pgfqpoint{2.455045in}{2.757989in}}%
\pgfpathlineto{\pgfqpoint{2.455844in}{2.644559in}}%
\pgfpathlineto{\pgfqpoint{2.457992in}{2.644559in}}%
\pgfpathlineto{\pgfqpoint{2.457997in}{2.749775in}}%
\pgfpathlineto{\pgfqpoint{2.458761in}{2.692121in}}%
\pgfpathlineto{\pgfqpoint{2.459800in}{2.766626in}}%
\pgfpathlineto{\pgfqpoint{2.459965in}{2.777040in}}%
\pgfpathlineto{\pgfqpoint{2.461550in}{2.644559in}}%
\pgfpathlineto{\pgfqpoint{2.462937in}{2.644559in}}%
\pgfpathlineto{\pgfqpoint{2.462942in}{2.749729in}}%
\pgfpathlineto{\pgfqpoint{2.463698in}{2.699447in}}%
\pgfpathlineto{\pgfqpoint{2.463747in}{2.696004in}}%
\pgfpathlineto{\pgfqpoint{2.464094in}{2.724306in}}%
\pgfpathlineto{\pgfqpoint{2.464143in}{2.720854in}}%
\pgfpathlineto{\pgfqpoint{2.464873in}{2.773278in}}%
\pgfpathlineto{\pgfqpoint{2.464953in}{2.739281in}}%
\pgfpathlineto{\pgfqpoint{2.465771in}{2.644559in}}%
\pgfpathlineto{\pgfqpoint{2.467882in}{2.644559in}}%
\pgfpathlineto{\pgfqpoint{2.467888in}{2.749813in}}%
\pgfpathlineto{\pgfqpoint{2.468627in}{2.698239in}}%
\pgfpathlineto{\pgfqpoint{2.468677in}{2.695415in}}%
\pgfpathlineto{\pgfqpoint{2.468924in}{2.716884in}}%
\pgfpathlineto{\pgfqpoint{2.468973in}{2.714052in}}%
\pgfpathlineto{\pgfqpoint{2.469810in}{2.772581in}}%
\pgfpathlineto{\pgfqpoint{2.469881in}{2.757855in}}%
\pgfpathlineto{\pgfqpoint{2.470685in}{2.644559in}}%
\pgfpathlineto{\pgfqpoint{2.472828in}{2.644559in}}%
\pgfpathlineto{\pgfqpoint{2.472833in}{2.749773in}}%
\pgfpathlineto{\pgfqpoint{2.473573in}{2.699790in}}%
\pgfpathlineto{\pgfqpoint{2.473622in}{2.693700in}}%
\pgfpathlineto{\pgfqpoint{2.474067in}{2.730858in}}%
\pgfpathlineto{\pgfqpoint{2.474117in}{2.724768in}}%
\pgfpathlineto{\pgfqpoint{2.474756in}{2.774124in}}%
\pgfpathlineto{\pgfqpoint{2.474837in}{2.716301in}}%
\pgfpathlineto{\pgfqpoint{2.474846in}{2.750557in}}%
\pgfpathlineto{\pgfqpoint{2.475635in}{2.644559in}}%
\pgfpathlineto{\pgfqpoint{2.477773in}{2.644559in}}%
\pgfpathlineto{\pgfqpoint{2.477779in}{2.749778in}}%
\pgfpathlineto{\pgfqpoint{2.478506in}{2.697348in}}%
\pgfpathlineto{\pgfqpoint{2.478556in}{2.694693in}}%
\pgfpathlineto{\pgfqpoint{2.478803in}{2.715992in}}%
\pgfpathlineto{\pgfqpoint{2.478853in}{2.713331in}}%
\pgfpathlineto{\pgfqpoint{2.479747in}{2.775307in}}%
\pgfpathlineto{\pgfqpoint{2.479771in}{2.757529in}}%
\pgfpathlineto{\pgfqpoint{2.480579in}{2.644559in}}%
\pgfpathlineto{\pgfqpoint{2.482719in}{2.644559in}}%
\pgfpathlineto{\pgfqpoint{2.482724in}{2.749796in}}%
\pgfpathlineto{\pgfqpoint{2.483485in}{2.693790in}}%
\pgfpathlineto{\pgfqpoint{2.484693in}{2.775253in}}%
\pgfpathlineto{\pgfqpoint{2.485464in}{2.644559in}}%
\pgfpathlineto{\pgfqpoint{2.487664in}{2.644559in}}%
\pgfpathlineto{\pgfqpoint{2.487670in}{2.749786in}}%
\pgfpathlineto{\pgfqpoint{2.488416in}{2.698631in}}%
\pgfpathlineto{\pgfqpoint{2.488466in}{2.695797in}}%
\pgfpathlineto{\pgfqpoint{2.488713in}{2.717274in}}%
\pgfpathlineto{\pgfqpoint{2.488762in}{2.714435in}}%
\pgfpathlineto{\pgfqpoint{2.489596in}{2.772749in}}%
\pgfpathlineto{\pgfqpoint{2.489663in}{2.758140in}}%
\pgfpathlineto{\pgfqpoint{2.490463in}{2.644559in}}%
\pgfpathlineto{\pgfqpoint{2.492610in}{2.644559in}}%
\pgfpathlineto{\pgfqpoint{2.492615in}{2.749790in}}%
\pgfpathlineto{\pgfqpoint{2.493383in}{2.699897in}}%
\pgfpathlineto{\pgfqpoint{2.493432in}{2.697177in}}%
\pgfpathlineto{\pgfqpoint{2.493679in}{2.718540in}}%
\pgfpathlineto{\pgfqpoint{2.493729in}{2.715816in}}%
\pgfpathlineto{\pgfqpoint{2.494552in}{2.773358in}}%
\pgfpathlineto{\pgfqpoint{2.494608in}{2.758003in}}%
\pgfpathlineto{\pgfqpoint{2.495410in}{2.644559in}}%
\pgfpathlineto{\pgfqpoint{2.497555in}{2.644559in}}%
\pgfpathlineto{\pgfqpoint{2.497560in}{2.749768in}}%
\pgfpathlineto{\pgfqpoint{2.498290in}{2.684569in}}%
\pgfpathlineto{\pgfqpoint{2.499230in}{2.763546in}}%
\pgfpathlineto{\pgfqpoint{2.499503in}{2.780707in}}%
\pgfpathlineto{\pgfqpoint{2.500349in}{2.644559in}}%
\pgfpathlineto{\pgfqpoint{2.502501in}{2.644559in}}%
\pgfpathlineto{\pgfqpoint{2.502506in}{2.749801in}}%
\pgfpathlineto{\pgfqpoint{2.503263in}{2.693453in}}%
\pgfpathlineto{\pgfqpoint{2.504474in}{2.775335in}}%
\pgfpathlineto{\pgfqpoint{2.505245in}{2.644559in}}%
\pgfpathlineto{\pgfqpoint{2.507446in}{2.644559in}}%
\pgfpathlineto{\pgfqpoint{2.507451in}{2.749780in}}%
\pgfpathlineto{\pgfqpoint{2.508180in}{2.697506in}}%
\pgfpathlineto{\pgfqpoint{2.508230in}{2.694695in}}%
\pgfpathlineto{\pgfqpoint{2.508477in}{2.716149in}}%
\pgfpathlineto{\pgfqpoint{2.508526in}{2.713333in}}%
\pgfpathlineto{\pgfqpoint{2.509420in}{2.775386in}}%
\pgfpathlineto{\pgfqpoint{2.509446in}{2.758371in}}%
\pgfpathlineto{\pgfqpoint{2.510244in}{2.644559in}}%
\pgfpathlineto{\pgfqpoint{2.512392in}{2.644559in}}%
\pgfpathlineto{\pgfqpoint{2.512397in}{2.749798in}}%
\pgfpathlineto{\pgfqpoint{2.513158in}{2.699639in}}%
\pgfpathlineto{\pgfqpoint{2.513207in}{2.696619in}}%
\pgfpathlineto{\pgfqpoint{2.513454in}{2.718282in}}%
\pgfpathlineto{\pgfqpoint{2.513504in}{2.715258in}}%
\pgfpathlineto{\pgfqpoint{2.514330in}{2.773312in}}%
\pgfpathlineto{\pgfqpoint{2.514390in}{2.758017in}}%
\pgfpathlineto{\pgfqpoint{2.515190in}{2.644559in}}%
\pgfpathlineto{\pgfqpoint{2.517337in}{2.644559in}}%
\pgfpathlineto{\pgfqpoint{2.517342in}{2.749750in}}%
\pgfpathlineto{\pgfqpoint{2.517413in}{2.644337in}}%
\pgfpathlineto{\pgfqpoint{2.518105in}{2.685067in}}%
\pgfpathlineto{\pgfqpoint{2.518946in}{2.760947in}}%
\pgfpathlineto{\pgfqpoint{2.519311in}{2.783869in}}%
\pgfpathlineto{\pgfqpoint{2.520005in}{2.644559in}}%
\pgfpathlineto{\pgfqpoint{2.522282in}{2.644559in}}%
\pgfpathlineto{\pgfqpoint{2.522288in}{2.749813in}}%
\pgfpathlineto{\pgfqpoint{2.523026in}{2.698169in}}%
\pgfpathlineto{\pgfqpoint{2.523076in}{2.695352in}}%
\pgfpathlineto{\pgfqpoint{2.523323in}{2.716813in}}%
\pgfpathlineto{\pgfqpoint{2.523372in}{2.713990in}}%
\pgfpathlineto{\pgfqpoint{2.524210in}{2.772542in}}%
\pgfpathlineto{\pgfqpoint{2.524281in}{2.757860in}}%
\pgfpathlineto{\pgfqpoint{2.525084in}{2.644559in}}%
\pgfpathlineto{\pgfqpoint{2.527228in}{2.644559in}}%
\pgfpathlineto{\pgfqpoint{2.527233in}{2.749758in}}%
\pgfpathlineto{\pgfqpoint{2.527978in}{2.700502in}}%
\pgfpathlineto{\pgfqpoint{2.528027in}{2.693611in}}%
\pgfpathlineto{\pgfqpoint{2.528472in}{2.731569in}}%
\pgfpathlineto{\pgfqpoint{2.528522in}{2.724680in}}%
\pgfpathlineto{\pgfqpoint{2.529159in}{2.774665in}}%
\pgfpathlineto{\pgfqpoint{2.529237in}{2.716254in}}%
\pgfpathlineto{\pgfqpoint{2.529246in}{2.750564in}}%
\pgfpathlineto{\pgfqpoint{2.530036in}{2.644559in}}%
\pgfpathlineto{\pgfqpoint{2.532173in}{2.644559in}}%
\pgfpathlineto{\pgfqpoint{2.532179in}{2.749771in}}%
\pgfpathlineto{\pgfqpoint{2.532945in}{2.694087in}}%
\pgfpathlineto{\pgfqpoint{2.534147in}{2.775218in}}%
\pgfpathlineto{\pgfqpoint{2.534919in}{2.644559in}}%
\pgfpathlineto{\pgfqpoint{2.537119in}{2.644559in}}%
\pgfpathlineto{\pgfqpoint{2.537124in}{2.749792in}}%
\pgfpathlineto{\pgfqpoint{2.537874in}{2.698915in}}%
\pgfpathlineto{\pgfqpoint{2.537923in}{2.695914in}}%
\pgfpathlineto{\pgfqpoint{2.538170in}{2.717558in}}%
\pgfpathlineto{\pgfqpoint{2.538220in}{2.714553in}}%
\pgfpathlineto{\pgfqpoint{2.539052in}{2.772939in}}%
\pgfpathlineto{\pgfqpoint{2.539118in}{2.758035in}}%
\pgfpathlineto{\pgfqpoint{2.539917in}{2.644559in}}%
\pgfpathlineto{\pgfqpoint{2.542064in}{2.644559in}}%
\pgfpathlineto{\pgfqpoint{2.542070in}{2.749757in}}%
\pgfpathlineto{\pgfqpoint{2.542139in}{2.644214in}}%
\pgfpathlineto{\pgfqpoint{2.542829in}{2.684636in}}%
\pgfpathlineto{\pgfqpoint{2.543669in}{2.760947in}}%
\pgfpathlineto{\pgfqpoint{2.544038in}{2.784095in}}%
\pgfpathlineto{\pgfqpoint{2.544722in}{2.644559in}}%
\pgfpathlineto{\pgfqpoint{2.547010in}{2.644559in}}%
\pgfpathlineto{\pgfqpoint{2.547015in}{2.749806in}}%
\pgfpathlineto{\pgfqpoint{2.547740in}{2.697355in}}%
\pgfpathlineto{\pgfqpoint{2.547790in}{2.694497in}}%
\pgfpathlineto{\pgfqpoint{2.548037in}{2.715999in}}%
\pgfpathlineto{\pgfqpoint{2.548087in}{2.713135in}}%
\pgfpathlineto{\pgfqpoint{2.548984in}{2.775458in}}%
\pgfpathlineto{\pgfqpoint{2.548984in}{2.748601in}}%
\pgfpathlineto{\pgfqpoint{2.549007in}{2.757477in}}%
\pgfpathlineto{\pgfqpoint{2.549755in}{2.644559in}}%
\pgfpathlineto{\pgfqpoint{2.551955in}{2.644559in}}%
\pgfpathlineto{\pgfqpoint{2.551960in}{2.749763in}}%
\pgfpathlineto{\pgfqpoint{2.552681in}{2.682579in}}%
\pgfpathlineto{\pgfqpoint{2.553521in}{2.758136in}}%
\pgfpathlineto{\pgfqpoint{2.553898in}{2.781797in}}%
\pgfpathlineto{\pgfqpoint{2.554576in}{2.644559in}}%
\pgfpathlineto{\pgfqpoint{2.556901in}{2.644559in}}%
\pgfpathlineto{\pgfqpoint{2.556906in}{2.749799in}}%
\pgfpathlineto{\pgfqpoint{2.557669in}{2.693835in}}%
\pgfpathlineto{\pgfqpoint{2.558874in}{2.775347in}}%
\pgfpathlineto{\pgfqpoint{2.559645in}{2.644559in}}%
\pgfpathlineto{\pgfqpoint{2.561846in}{2.644559in}}%
\pgfpathlineto{\pgfqpoint{2.561851in}{2.749771in}}%
\pgfpathlineto{\pgfqpoint{2.562578in}{2.689604in}}%
\pgfpathlineto{\pgfqpoint{2.563617in}{2.764427in}}%
\pgfpathlineto{\pgfqpoint{2.563792in}{2.775470in}}%
\pgfpathlineto{\pgfqpoint{2.565303in}{2.644559in}}%
\pgfpathlineto{\pgfqpoint{2.566792in}{2.644559in}}%
\pgfpathlineto{\pgfqpoint{2.566797in}{2.749759in}}%
\pgfpathlineto{\pgfqpoint{2.567534in}{2.698169in}}%
\pgfpathlineto{\pgfqpoint{2.567584in}{2.695051in}}%
\pgfpathlineto{\pgfqpoint{2.567831in}{2.716813in}}%
\pgfpathlineto{\pgfqpoint{2.567881in}{2.713689in}}%
\pgfpathlineto{\pgfqpoint{2.568719in}{2.772569in}}%
\pgfpathlineto{\pgfqpoint{2.568791in}{2.758055in}}%
\pgfpathlineto{\pgfqpoint{2.569589in}{2.644559in}}%
\pgfpathlineto{\pgfqpoint{2.571737in}{2.644559in}}%
\pgfpathlineto{\pgfqpoint{2.571742in}{2.749784in}}%
\pgfpathlineto{\pgfqpoint{2.572488in}{2.698586in}}%
\pgfpathlineto{\pgfqpoint{2.572538in}{2.695751in}}%
\pgfpathlineto{\pgfqpoint{2.572785in}{2.717229in}}%
\pgfpathlineto{\pgfqpoint{2.572834in}{2.714389in}}%
\pgfpathlineto{\pgfqpoint{2.573668in}{2.772727in}}%
\pgfpathlineto{\pgfqpoint{2.573736in}{2.758140in}}%
\pgfpathlineto{\pgfqpoint{2.574533in}{2.644559in}}%
\pgfpathlineto{\pgfqpoint{2.576682in}{2.644559in}}%
\pgfpathlineto{\pgfqpoint{2.576688in}{2.749776in}}%
\pgfpathlineto{\pgfqpoint{2.577426in}{2.689565in}}%
\pgfpathlineto{\pgfqpoint{2.578366in}{2.759719in}}%
\pgfpathlineto{\pgfqpoint{2.578656in}{2.777970in}}%
\pgfpathlineto{\pgfqpoint{2.579844in}{2.644559in}}%
\pgfpathlineto{\pgfqpoint{2.581628in}{2.644559in}}%
\pgfpathlineto{\pgfqpoint{2.581633in}{2.749723in}}%
\pgfpathlineto{\pgfqpoint{2.582396in}{2.693721in}}%
\pgfpathlineto{\pgfqpoint{2.583602in}{2.775237in}}%
\pgfpathlineto{\pgfqpoint{2.584371in}{2.644559in}}%
\pgfpathlineto{\pgfqpoint{2.586573in}{2.644559in}}%
\pgfpathlineto{\pgfqpoint{2.586579in}{2.749802in}}%
\pgfpathlineto{\pgfqpoint{2.587304in}{2.697254in}}%
\pgfpathlineto{\pgfqpoint{2.587354in}{2.694575in}}%
\pgfpathlineto{\pgfqpoint{2.587601in}{2.715898in}}%
\pgfpathlineto{\pgfqpoint{2.587650in}{2.713213in}}%
\pgfpathlineto{\pgfqpoint{2.588547in}{2.775357in}}%
\pgfpathlineto{\pgfqpoint{2.588548in}{2.735948in}}%
\pgfpathlineto{\pgfqpoint{2.588571in}{2.757602in}}%
\pgfpathlineto{\pgfqpoint{2.589336in}{2.644559in}}%
\pgfpathlineto{\pgfqpoint{2.591519in}{2.644559in}}%
\pgfpathlineto{\pgfqpoint{2.591524in}{2.749761in}}%
\pgfpathlineto{\pgfqpoint{2.592249in}{2.682828in}}%
\pgfpathlineto{\pgfqpoint{2.593090in}{2.758513in}}%
\pgfpathlineto{\pgfqpoint{2.593464in}{2.782017in}}%
\pgfpathlineto{\pgfqpoint{2.594143in}{2.644559in}}%
\pgfpathlineto{\pgfqpoint{2.596464in}{2.644559in}}%
\pgfpathlineto{\pgfqpoint{2.596470in}{2.749793in}}%
\pgfpathlineto{\pgfqpoint{2.597196in}{2.697312in}}%
\pgfpathlineto{\pgfqpoint{2.597246in}{2.694647in}}%
\pgfpathlineto{\pgfqpoint{2.597493in}{2.715956in}}%
\pgfpathlineto{\pgfqpoint{2.597542in}{2.713285in}}%
\pgfpathlineto{\pgfqpoint{2.598438in}{2.775335in}}%
\pgfpathlineto{\pgfqpoint{2.598462in}{2.757569in}}%
\pgfpathlineto{\pgfqpoint{2.599268in}{2.644559in}}%
\pgfpathlineto{\pgfqpoint{2.601410in}{2.644559in}}%
\pgfpathlineto{\pgfqpoint{2.601415in}{2.749763in}}%
\pgfpathlineto{\pgfqpoint{2.602164in}{2.702047in}}%
\pgfpathlineto{\pgfqpoint{2.602213in}{2.692606in}}%
\pgfpathlineto{\pgfqpoint{2.602658in}{2.733110in}}%
\pgfpathlineto{\pgfqpoint{2.602708in}{2.723679in}}%
\pgfpathlineto{\pgfqpoint{2.603342in}{2.776072in}}%
\pgfpathlineto{\pgfqpoint{2.603419in}{2.716205in}}%
\pgfpathlineto{\pgfqpoint{2.603428in}{2.750647in}}%
\pgfpathlineto{\pgfqpoint{2.604221in}{2.644559in}}%
\pgfpathlineto{\pgfqpoint{2.606355in}{2.644559in}}%
\pgfpathlineto{\pgfqpoint{2.606360in}{2.749762in}}%
\pgfpathlineto{\pgfqpoint{2.607113in}{2.688355in}}%
\pgfpathlineto{\pgfqpoint{2.608053in}{2.762664in}}%
\pgfpathlineto{\pgfqpoint{2.608329in}{2.780014in}}%
\pgfpathlineto{\pgfqpoint{2.609303in}{2.644559in}}%
\pgfpathlineto{\pgfqpoint{2.611301in}{2.644559in}}%
\pgfpathlineto{\pgfqpoint{2.611306in}{2.749757in}}%
\pgfpathlineto{\pgfqpoint{2.612050in}{2.698441in}}%
\pgfpathlineto{\pgfqpoint{2.612099in}{2.695587in}}%
\pgfpathlineto{\pgfqpoint{2.612347in}{2.717084in}}%
\pgfpathlineto{\pgfqpoint{2.612396in}{2.714226in}}%
\pgfpathlineto{\pgfqpoint{2.613231in}{2.772635in}}%
\pgfpathlineto{\pgfqpoint{2.613300in}{2.758107in}}%
\pgfpathlineto{\pgfqpoint{2.614097in}{2.644559in}}%
\pgfpathlineto{\pgfqpoint{2.616246in}{2.644559in}}%
\pgfpathlineto{\pgfqpoint{2.616251in}{2.749783in}}%
\pgfpathlineto{\pgfqpoint{2.616998in}{2.698647in}}%
\pgfpathlineto{\pgfqpoint{2.617048in}{2.695809in}}%
\pgfpathlineto{\pgfqpoint{2.617295in}{2.717291in}}%
\pgfpathlineto{\pgfqpoint{2.617344in}{2.714448in}}%
\pgfpathlineto{\pgfqpoint{2.618178in}{2.772758in}}%
\pgfpathlineto{\pgfqpoint{2.618245in}{2.758136in}}%
\pgfpathlineto{\pgfqpoint{2.619042in}{2.644559in}}%
\pgfpathlineto{\pgfqpoint{2.621192in}{2.644559in}}%
\pgfpathlineto{\pgfqpoint{2.621197in}{2.749761in}}%
\pgfpathlineto{\pgfqpoint{2.621928in}{2.683120in}}%
\pgfpathlineto{\pgfqpoint{2.622768in}{2.758930in}}%
\pgfpathlineto{\pgfqpoint{2.623140in}{2.782254in}}%
\pgfpathlineto{\pgfqpoint{2.623822in}{2.644559in}}%
\pgfpathlineto{\pgfqpoint{2.626137in}{2.644559in}}%
\pgfpathlineto{\pgfqpoint{2.626142in}{2.749790in}}%
\pgfpathlineto{\pgfqpoint{2.626898in}{2.693394in}}%
\pgfpathlineto{\pgfqpoint{2.628111in}{2.775314in}}%
\pgfpathlineto{\pgfqpoint{2.628882in}{2.644559in}}%
\pgfpathlineto{\pgfqpoint{2.631082in}{2.644559in}}%
\pgfpathlineto{\pgfqpoint{2.631088in}{2.749760in}}%
\pgfpathlineto{\pgfqpoint{2.631818in}{2.683070in}}%
\pgfpathlineto{\pgfqpoint{2.632658in}{2.758864in}}%
\pgfpathlineto{\pgfqpoint{2.633030in}{2.782218in}}%
\pgfpathlineto{\pgfqpoint{2.633713in}{2.644559in}}%
\pgfpathlineto{\pgfqpoint{2.636027in}{2.644559in}}%
\pgfpathlineto{\pgfqpoint{2.636033in}{2.752984in}}%
\pgfpathlineto{\pgfqpoint{2.636111in}{2.638518in}}%
\pgfpathlineto{\pgfqpoint{2.636797in}{2.679737in}}%
\pgfpathlineto{\pgfqpoint{2.637637in}{2.766478in}}%
\pgfpathlineto{\pgfqpoint{2.638002in}{2.789298in}}%
\pgfpathlineto{\pgfqpoint{2.638604in}{2.644559in}}%
\pgfpathlineto{\pgfqpoint{2.640973in}{2.644559in}}%
\pgfpathlineto{\pgfqpoint{2.640979in}{2.749475in}}%
\pgfpathlineto{\pgfqpoint{2.641719in}{2.698386in}}%
\pgfpathlineto{\pgfqpoint{2.641768in}{2.694413in}}%
\pgfpathlineto{\pgfqpoint{2.642114in}{2.723246in}}%
\pgfpathlineto{\pgfqpoint{2.642164in}{2.719262in}}%
\pgfpathlineto{\pgfqpoint{2.642902in}{2.772718in}}%
\pgfpathlineto{\pgfqpoint{2.642989in}{2.734647in}}%
\pgfpathlineto{\pgfqpoint{2.643826in}{2.644559in}}%
\pgfpathlineto{\pgfqpoint{2.645919in}{2.644565in}}%
\pgfpathlineto{\pgfqpoint{2.645923in}{2.774677in}}%
\pgfpathlineto{\pgfqpoint{2.646016in}{2.628571in}}%
\pgfpathlineto{\pgfqpoint{2.646653in}{2.721829in}}%
\pgfpathlineto{\pgfqpoint{2.646703in}{2.670756in}}%
\pgfpathlineto{\pgfqpoint{2.647346in}{2.765228in}}%
\pgfpathlineto{\pgfqpoint{2.647395in}{2.714346in}}%
\pgfpathlineto{\pgfqpoint{2.647893in}{2.799494in}}%
\pgfpathlineto{\pgfqpoint{2.647983in}{2.644559in}}%
\pgfpathlineto{\pgfqpoint{2.648104in}{2.644559in}}%
\pgfpathlineto{\pgfqpoint{2.650864in}{2.644559in}}%
\pgfpathlineto{\pgfqpoint{2.650870in}{2.749684in}}%
\pgfpathlineto{\pgfqpoint{2.650953in}{2.642188in}}%
\pgfpathlineto{\pgfqpoint{2.651593in}{2.707320in}}%
\pgfpathlineto{\pgfqpoint{2.651642in}{2.683904in}}%
\pgfpathlineto{\pgfqpoint{2.652186in}{2.744569in}}%
\pgfpathlineto{\pgfqpoint{2.652236in}{2.721218in}}%
\pgfpathlineto{\pgfqpoint{2.652838in}{2.785475in}}%
\pgfpathlineto{\pgfqpoint{2.652894in}{2.693549in}}%
\pgfpathlineto{\pgfqpoint{2.652908in}{2.751948in}}%
\pgfpathlineto{\pgfqpoint{2.653691in}{2.644559in}}%
\pgfpathlineto{\pgfqpoint{2.655809in}{2.644559in}}%
\pgfpathlineto{\pgfqpoint{2.655815in}{2.750056in}}%
\pgfpathlineto{\pgfqpoint{2.655895in}{2.636104in}}%
\pgfpathlineto{\pgfqpoint{2.656581in}{2.677480in}}%
\pgfpathlineto{\pgfqpoint{2.657421in}{2.768748in}}%
\pgfpathlineto{\pgfqpoint{2.657784in}{2.791439in}}%
\pgfpathlineto{\pgfqpoint{2.658364in}{2.644559in}}%
\pgfpathlineto{\pgfqpoint{2.660755in}{2.644559in}}%
\pgfpathlineto{\pgfqpoint{2.660761in}{2.749365in}}%
\pgfpathlineto{\pgfqpoint{2.661525in}{2.700322in}}%
\pgfpathlineto{\pgfqpoint{2.661574in}{2.695550in}}%
\pgfpathlineto{\pgfqpoint{2.661920in}{2.725182in}}%
\pgfpathlineto{\pgfqpoint{2.661970in}{2.720399in}}%
\pgfpathlineto{\pgfqpoint{2.662696in}{2.773891in}}%
\pgfpathlineto{\pgfqpoint{2.662780in}{2.734667in}}%
\pgfpathlineto{\pgfqpoint{2.663593in}{2.644559in}}%
\pgfpathlineto{\pgfqpoint{2.665701in}{2.644559in}}%
\pgfpathlineto{\pgfqpoint{2.665705in}{2.758680in}}%
\pgfpathlineto{\pgfqpoint{2.666456in}{2.685424in}}%
\pgfpathlineto{\pgfqpoint{2.667297in}{2.759118in}}%
\pgfpathlineto{\pgfqpoint{2.667674in}{2.782799in}}%
\pgfpathlineto{\pgfqpoint{2.668388in}{2.644559in}}%
\pgfpathlineto{\pgfqpoint{2.670646in}{2.644559in}}%
\pgfpathlineto{\pgfqpoint{2.670651in}{2.749752in}}%
\pgfpathlineto{\pgfqpoint{2.671416in}{2.693846in}}%
\pgfpathlineto{\pgfqpoint{2.672620in}{2.775281in}}%
\pgfpathlineto{\pgfqpoint{2.673392in}{2.644559in}}%
\pgfpathlineto{\pgfqpoint{2.675590in}{2.644559in}}%
\pgfpathlineto{\pgfqpoint{2.675596in}{2.758041in}}%
\pgfpathlineto{\pgfqpoint{2.676355in}{2.699812in}}%
\pgfpathlineto{\pgfqpoint{2.676404in}{2.696031in}}%
\pgfpathlineto{\pgfqpoint{2.676750in}{2.724648in}}%
\pgfpathlineto{\pgfqpoint{2.676800in}{2.720905in}}%
\pgfpathlineto{\pgfqpoint{2.677529in}{2.773505in}}%
\pgfpathlineto{\pgfqpoint{2.677592in}{2.758685in}}%
\pgfpathlineto{\pgfqpoint{2.678387in}{2.644559in}}%
\pgfpathlineto{\pgfqpoint{2.680537in}{2.644559in}}%
\pgfpathlineto{\pgfqpoint{2.680542in}{2.749716in}}%
\pgfpathlineto{\pgfqpoint{2.681295in}{2.691394in}}%
\pgfpathlineto{\pgfqpoint{2.682334in}{2.765763in}}%
\pgfpathlineto{\pgfqpoint{2.682511in}{2.776879in}}%
\pgfpathlineto{\pgfqpoint{2.684100in}{2.644559in}}%
\pgfpathlineto{\pgfqpoint{2.685481in}{2.644559in}}%
\pgfpathlineto{\pgfqpoint{2.685487in}{2.753387in}}%
\pgfpathlineto{\pgfqpoint{2.685568in}{2.637207in}}%
\pgfpathlineto{\pgfqpoint{2.686254in}{2.678635in}}%
\pgfpathlineto{\pgfqpoint{2.687095in}{2.767681in}}%
\pgfpathlineto{\pgfqpoint{2.687456in}{2.790328in}}%
\pgfpathlineto{\pgfqpoint{2.688050in}{2.644559in}}%
\pgfpathlineto{\pgfqpoint{2.690428in}{2.644559in}}%
\pgfpathlineto{\pgfqpoint{2.690432in}{2.754592in}}%
\pgfpathlineto{\pgfqpoint{2.690506in}{2.639742in}}%
\pgfpathlineto{\pgfqpoint{2.691192in}{2.680494in}}%
\pgfpathlineto{\pgfqpoint{2.692032in}{2.764884in}}%
\pgfpathlineto{\pgfqpoint{2.692402in}{2.788026in}}%
\pgfpathlineto{\pgfqpoint{2.693011in}{2.644559in}}%
\pgfpathlineto{\pgfqpoint{2.695375in}{2.644920in}}%
\pgfpathlineto{\pgfqpoint{2.695379in}{2.750011in}}%
\pgfpathlineto{\pgfqpoint{2.695455in}{2.618651in}}%
\pgfpathlineto{\pgfqpoint{2.696140in}{2.659736in}}%
\pgfpathlineto{\pgfqpoint{2.696980in}{2.785899in}}%
\pgfpathlineto{\pgfqpoint{2.697347in}{2.808835in}}%
\pgfpathlineto{\pgfqpoint{2.697850in}{2.644559in}}%
\pgfpathlineto{\pgfqpoint{2.700318in}{2.644559in}}%
\pgfpathlineto{\pgfqpoint{2.700324in}{2.751264in}}%
\pgfpathlineto{\pgfqpoint{2.700400in}{2.639121in}}%
\pgfpathlineto{\pgfqpoint{2.701085in}{2.680124in}}%
\pgfpathlineto{\pgfqpoint{2.701926in}{2.765683in}}%
\pgfpathlineto{\pgfqpoint{2.702293in}{2.788667in}}%
\pgfpathlineto{\pgfqpoint{2.702899in}{2.644559in}}%
\pgfpathlineto{\pgfqpoint{2.705264in}{2.644559in}}%
\pgfpathlineto{\pgfqpoint{2.705270in}{2.749757in}}%
\pgfpathlineto{\pgfqpoint{2.706026in}{2.707523in}}%
\pgfpathlineto{\pgfqpoint{2.706075in}{2.688058in}}%
\pgfpathlineto{\pgfqpoint{2.706619in}{2.744780in}}%
\pgfpathlineto{\pgfqpoint{2.706669in}{2.725363in}}%
\pgfpathlineto{\pgfqpoint{2.707201in}{2.781275in}}%
\pgfpathlineto{\pgfqpoint{2.707293in}{2.716251in}}%
\pgfpathlineto{\pgfqpoint{2.707307in}{2.731077in}}%
\pgfpathlineto{\pgfqpoint{2.708093in}{2.644559in}}%
\pgfpathlineto{\pgfqpoint{2.710209in}{2.644559in}}%
\pgfpathlineto{\pgfqpoint{2.710215in}{2.749754in}}%
\pgfpathlineto{\pgfqpoint{2.710286in}{2.641928in}}%
\pgfpathlineto{\pgfqpoint{2.710972in}{2.682488in}}%
\pgfpathlineto{\pgfqpoint{2.711812in}{2.762835in}}%
\pgfpathlineto{\pgfqpoint{2.712184in}{2.786107in}}%
\pgfpathlineto{\pgfqpoint{2.712824in}{2.644559in}}%
\pgfpathlineto{\pgfqpoint{2.715155in}{2.644559in}}%
\pgfpathlineto{\pgfqpoint{2.715160in}{2.749579in}}%
\pgfpathlineto{\pgfqpoint{2.715270in}{2.642799in}}%
\pgfpathlineto{\pgfqpoint{2.715911in}{2.708578in}}%
\pgfpathlineto{\pgfqpoint{2.715960in}{2.685786in}}%
\pgfpathlineto{\pgfqpoint{2.716504in}{2.745826in}}%
\pgfpathlineto{\pgfqpoint{2.716554in}{2.723100in}}%
\pgfpathlineto{\pgfqpoint{2.717089in}{2.782503in}}%
\pgfpathlineto{\pgfqpoint{2.717183in}{2.697867in}}%
\pgfpathlineto{\pgfqpoint{2.717198in}{2.747408in}}%
\pgfpathlineto{\pgfqpoint{2.717983in}{2.644559in}}%
\pgfpathlineto{\pgfqpoint{2.720101in}{2.644559in}}%
\pgfpathlineto{\pgfqpoint{2.720105in}{2.756404in}}%
\pgfpathlineto{\pgfqpoint{2.720180in}{2.639123in}}%
\pgfpathlineto{\pgfqpoint{2.720865in}{2.679949in}}%
\pgfpathlineto{\pgfqpoint{2.721706in}{2.765669in}}%
\pgfpathlineto{\pgfqpoint{2.722074in}{2.788763in}}%
\pgfpathlineto{\pgfqpoint{2.722677in}{2.644559in}}%
\pgfpathlineto{\pgfqpoint{2.725046in}{2.644559in}}%
\pgfpathlineto{\pgfqpoint{2.725051in}{2.749764in}}%
\pgfpathlineto{\pgfqpoint{2.725813in}{2.686502in}}%
\pgfpathlineto{\pgfqpoint{2.726753in}{2.765678in}}%
\pgfpathlineto{\pgfqpoint{2.727020in}{2.782441in}}%
\pgfpathlineto{\pgfqpoint{2.727875in}{2.644559in}}%
\pgfpathlineto{\pgfqpoint{2.729990in}{2.644559in}}%
\pgfpathlineto{\pgfqpoint{2.729996in}{2.758244in}}%
\pgfpathlineto{\pgfqpoint{2.730750in}{2.697572in}}%
\pgfpathlineto{\pgfqpoint{2.730799in}{2.697578in}}%
\pgfpathlineto{\pgfqpoint{2.731926in}{2.771441in}}%
\pgfpathlineto{\pgfqpoint{2.731966in}{2.748322in}}%
\pgfpathlineto{\pgfqpoint{2.731994in}{2.758061in}}%
\pgfpathlineto{\pgfqpoint{2.732738in}{2.644559in}}%
\pgfpathlineto{\pgfqpoint{2.734937in}{2.644559in}}%
\pgfpathlineto{\pgfqpoint{2.734942in}{2.749713in}}%
\pgfpathlineto{\pgfqpoint{2.735008in}{2.643635in}}%
\pgfpathlineto{\pgfqpoint{2.735695in}{2.683444in}}%
\pgfpathlineto{\pgfqpoint{2.736535in}{2.761198in}}%
\pgfpathlineto{\pgfqpoint{2.736911in}{2.784761in}}%
\pgfpathlineto{\pgfqpoint{2.737570in}{2.644559in}}%
\pgfpathlineto{\pgfqpoint{2.739883in}{2.644570in}}%
\pgfpathlineto{\pgfqpoint{2.739886in}{2.773390in}}%
\pgfpathlineto{\pgfqpoint{2.740636in}{2.696079in}}%
\pgfpathlineto{\pgfqpoint{2.741477in}{2.747387in}}%
\pgfpathlineto{\pgfqpoint{2.741856in}{2.771198in}}%
\pgfpathlineto{\pgfqpoint{2.741892in}{2.715878in}}%
\pgfpathlineto{\pgfqpoint{2.741901in}{2.750084in}}%
\pgfpathlineto{\pgfqpoint{2.742693in}{2.644559in}}%
\pgfpathlineto{\pgfqpoint{2.744828in}{2.644559in}}%
\pgfpathlineto{\pgfqpoint{2.744833in}{2.749743in}}%
\pgfpathlineto{\pgfqpoint{2.745592in}{2.700556in}}%
\pgfpathlineto{\pgfqpoint{2.745642in}{2.695561in}}%
\pgfpathlineto{\pgfqpoint{2.745988in}{2.725416in}}%
\pgfpathlineto{\pgfqpoint{2.746037in}{2.720410in}}%
\pgfpathlineto{\pgfqpoint{2.746766in}{2.774289in}}%
\pgfpathlineto{\pgfqpoint{2.746844in}{2.734328in}}%
\pgfpathlineto{\pgfqpoint{2.747683in}{2.644559in}}%
\pgfpathlineto{\pgfqpoint{2.749773in}{2.644559in}}%
\pgfpathlineto{\pgfqpoint{2.749779in}{2.750038in}}%
\pgfpathlineto{\pgfqpoint{2.749850in}{2.641958in}}%
\pgfpathlineto{\pgfqpoint{2.750536in}{2.682572in}}%
\pgfpathlineto{\pgfqpoint{2.751376in}{2.762615in}}%
\pgfpathlineto{\pgfqpoint{2.751747in}{2.785859in}}%
\pgfpathlineto{\pgfqpoint{2.752389in}{2.644559in}}%
\pgfpathlineto{\pgfqpoint{2.754719in}{2.644559in}}%
\pgfpathlineto{\pgfqpoint{2.754724in}{2.749720in}}%
\pgfpathlineto{\pgfqpoint{2.755456in}{2.683089in}}%
\pgfpathlineto{\pgfqpoint{2.756296in}{2.758927in}}%
\pgfpathlineto{\pgfqpoint{2.756668in}{2.782227in}}%
\pgfpathlineto{\pgfqpoint{2.757349in}{2.644559in}}%
\pgfpathlineto{\pgfqpoint{2.759663in}{2.644559in}}%
\pgfpathlineto{\pgfqpoint{2.759669in}{2.755222in}}%
\pgfpathlineto{\pgfqpoint{2.760434in}{2.692396in}}%
\pgfpathlineto{\pgfqpoint{2.761473in}{2.766446in}}%
\pgfpathlineto{\pgfqpoint{2.761695in}{2.710916in}}%
\pgfpathlineto{\pgfqpoint{2.761638in}{2.776845in}}%
\pgfpathlineto{\pgfqpoint{2.761710in}{2.737132in}}%
\pgfpathlineto{\pgfqpoint{2.762496in}{2.644559in}}%
\pgfpathlineto{\pgfqpoint{2.764609in}{2.644559in}}%
\pgfpathlineto{\pgfqpoint{2.764615in}{2.748594in}}%
\pgfpathlineto{\pgfqpoint{2.764686in}{2.642189in}}%
\pgfpathlineto{\pgfqpoint{2.765371in}{2.682696in}}%
\pgfpathlineto{\pgfqpoint{2.766212in}{2.762397in}}%
\pgfpathlineto{\pgfqpoint{2.766584in}{2.785707in}}%
\pgfpathlineto{\pgfqpoint{2.767226in}{2.644559in}}%
\pgfpathlineto{\pgfqpoint{2.769555in}{2.644559in}}%
\pgfpathlineto{\pgfqpoint{2.769559in}{2.755150in}}%
\pgfpathlineto{\pgfqpoint{2.770311in}{2.685431in}}%
\pgfpathlineto{\pgfqpoint{2.771151in}{2.758698in}}%
\pgfpathlineto{\pgfqpoint{2.771529in}{2.782390in}}%
\pgfpathlineto{\pgfqpoint{2.772249in}{2.644559in}}%
\pgfpathlineto{\pgfqpoint{2.774501in}{2.644559in}}%
\pgfpathlineto{\pgfqpoint{2.774505in}{2.762378in}}%
\pgfpathlineto{\pgfqpoint{2.774576in}{2.642148in}}%
\pgfpathlineto{\pgfqpoint{2.775261in}{2.682568in}}%
\pgfpathlineto{\pgfqpoint{2.776102in}{2.762697in}}%
\pgfpathlineto{\pgfqpoint{2.776474in}{2.786042in}}%
\pgfpathlineto{\pgfqpoint{2.777113in}{2.644559in}}%
\pgfpathlineto{\pgfqpoint{2.779445in}{2.644559in}}%
\pgfpathlineto{\pgfqpoint{2.779451in}{2.748415in}}%
\pgfpathlineto{\pgfqpoint{2.780177in}{2.682626in}}%
\pgfpathlineto{\pgfqpoint{2.781018in}{2.758303in}}%
\pgfpathlineto{\pgfqpoint{2.781392in}{2.781769in}}%
\pgfpathlineto{\pgfqpoint{2.782074in}{2.644559in}}%
\pgfpathlineto{\pgfqpoint{2.784392in}{2.644583in}}%
\pgfpathlineto{\pgfqpoint{2.784396in}{2.782672in}}%
\pgfpathlineto{\pgfqpoint{2.784501in}{2.630335in}}%
\pgfpathlineto{\pgfqpoint{2.785141in}{2.720544in}}%
\pgfpathlineto{\pgfqpoint{2.785190in}{2.673150in}}%
\pgfpathlineto{\pgfqpoint{2.785833in}{2.763944in}}%
\pgfpathlineto{\pgfqpoint{2.785883in}{2.716739in}}%
\pgfpathlineto{\pgfqpoint{2.786322in}{2.794570in}}%
\pgfpathlineto{\pgfqpoint{2.786459in}{2.644559in}}%
\pgfpathlineto{\pgfqpoint{2.786586in}{2.644559in}}%
\pgfpathlineto{\pgfqpoint{2.789337in}{2.644559in}}%
\pgfpathlineto{\pgfqpoint{2.789342in}{2.749694in}}%
\pgfpathlineto{\pgfqpoint{2.789453in}{2.643273in}}%
\pgfpathlineto{\pgfqpoint{2.790094in}{2.708697in}}%
\pgfpathlineto{\pgfqpoint{2.790144in}{2.686319in}}%
\pgfpathlineto{\pgfqpoint{2.790688in}{2.745945in}}%
\pgfpathlineto{\pgfqpoint{2.790737in}{2.723633in}}%
\pgfpathlineto{\pgfqpoint{2.791271in}{2.782570in}}%
\pgfpathlineto{\pgfqpoint{2.791365in}{2.697520in}}%
\pgfpathlineto{\pgfqpoint{2.791380in}{2.748105in}}%
\pgfpathlineto{\pgfqpoint{2.792168in}{2.644559in}}%
\pgfpathlineto{\pgfqpoint{2.794284in}{2.644838in}}%
\pgfpathlineto{\pgfqpoint{2.794288in}{2.752040in}}%
\pgfpathlineto{\pgfqpoint{2.794371in}{2.619644in}}%
\pgfpathlineto{\pgfqpoint{2.795006in}{2.729408in}}%
\pgfpathlineto{\pgfqpoint{2.795056in}{2.661294in}}%
\pgfpathlineto{\pgfqpoint{2.795699in}{2.772745in}}%
\pgfpathlineto{\pgfqpoint{2.795748in}{2.704946in}}%
\pgfpathlineto{\pgfqpoint{2.796256in}{2.807621in}}%
\pgfpathlineto{\pgfqpoint{2.796342in}{2.644559in}}%
\pgfpathlineto{\pgfqpoint{2.796477in}{2.644559in}}%
\pgfpathlineto{\pgfqpoint{2.799228in}{2.644559in}}%
\pgfpathlineto{\pgfqpoint{2.799233in}{2.749779in}}%
\pgfpathlineto{\pgfqpoint{2.799994in}{2.699549in}}%
\pgfpathlineto{\pgfqpoint{2.800044in}{2.696693in}}%
\pgfpathlineto{\pgfqpoint{2.800291in}{2.718192in}}%
\pgfpathlineto{\pgfqpoint{2.800341in}{2.715332in}}%
\pgfpathlineto{\pgfqpoint{2.801167in}{2.773209in}}%
\pgfpathlineto{\pgfqpoint{2.801227in}{2.758082in}}%
\pgfpathlineto{\pgfqpoint{2.802026in}{2.644559in}}%
\pgfpathlineto{\pgfqpoint{2.804173in}{2.644559in}}%
\pgfpathlineto{\pgfqpoint{2.804178in}{2.751048in}}%
\pgfpathlineto{\pgfqpoint{2.804252in}{2.641344in}}%
\pgfpathlineto{\pgfqpoint{2.804937in}{2.682075in}}%
\pgfpathlineto{\pgfqpoint{2.805778in}{2.763279in}}%
\pgfpathlineto{\pgfqpoint{2.806147in}{2.786449in}}%
\pgfpathlineto{\pgfqpoint{2.806779in}{2.644559in}}%
\pgfpathlineto{\pgfqpoint{2.809119in}{2.644559in}}%
\pgfpathlineto{\pgfqpoint{2.809124in}{2.749787in}}%
\pgfpathlineto{\pgfqpoint{2.809871in}{2.698627in}}%
\pgfpathlineto{\pgfqpoint{2.809920in}{2.695792in}}%
\pgfpathlineto{\pgfqpoint{2.810167in}{2.717270in}}%
\pgfpathlineto{\pgfqpoint{2.810217in}{2.714431in}}%
\pgfpathlineto{\pgfqpoint{2.811050in}{2.772748in}}%
\pgfpathlineto{\pgfqpoint{2.811118in}{2.758142in}}%
\pgfpathlineto{\pgfqpoint{2.811917in}{2.644559in}}%
\pgfpathlineto{\pgfqpoint{2.814064in}{2.644559in}}%
\pgfpathlineto{\pgfqpoint{2.814069in}{2.750347in}}%
\pgfpathlineto{\pgfqpoint{2.814142in}{2.641047in}}%
\pgfpathlineto{\pgfqpoint{2.814828in}{2.681781in}}%
\pgfpathlineto{\pgfqpoint{2.815668in}{2.763574in}}%
\pgfpathlineto{\pgfqpoint{2.816038in}{2.786742in}}%
\pgfpathlineto{\pgfqpoint{2.816665in}{2.644559in}}%
\pgfpathlineto{\pgfqpoint{2.819010in}{2.644559in}}%
\pgfpathlineto{\pgfqpoint{2.819015in}{2.749769in}}%
\pgfpathlineto{\pgfqpoint{2.819748in}{2.697895in}}%
\pgfpathlineto{\pgfqpoint{2.819797in}{2.694719in}}%
\pgfpathlineto{\pgfqpoint{2.820044in}{2.716538in}}%
\pgfpathlineto{\pgfqpoint{2.820094in}{2.713357in}}%
\pgfpathlineto{\pgfqpoint{2.820934in}{2.772465in}}%
\pgfpathlineto{\pgfqpoint{2.820984in}{2.748648in}}%
\pgfpathlineto{\pgfqpoint{2.821009in}{2.758111in}}%
\pgfpathlineto{\pgfqpoint{2.821755in}{2.644559in}}%
\pgfpathlineto{\pgfqpoint{2.823955in}{2.644559in}}%
\pgfpathlineto{\pgfqpoint{2.823959in}{2.752258in}}%
\pgfpathlineto{\pgfqpoint{2.824067in}{2.626313in}}%
\pgfpathlineto{\pgfqpoint{2.824708in}{2.725159in}}%
\pgfpathlineto{\pgfqpoint{2.824758in}{2.669293in}}%
\pgfpathlineto{\pgfqpoint{2.825400in}{2.768545in}}%
\pgfpathlineto{\pgfqpoint{2.825450in}{2.712896in}}%
\pgfpathlineto{\pgfqpoint{2.825887in}{2.799046in}}%
\pgfpathlineto{\pgfqpoint{2.826022in}{2.644559in}}%
\pgfpathlineto{\pgfqpoint{2.826165in}{2.644559in}}%
\pgfpathlineto{\pgfqpoint{2.828901in}{2.644559in}}%
\pgfpathlineto{\pgfqpoint{2.828906in}{2.749507in}}%
\pgfpathlineto{\pgfqpoint{2.828961in}{2.643022in}}%
\pgfpathlineto{\pgfqpoint{2.829637in}{2.680622in}}%
\pgfpathlineto{\pgfqpoint{2.830478in}{2.760951in}}%
\pgfpathlineto{\pgfqpoint{2.830849in}{2.784254in}}%
\pgfpathlineto{\pgfqpoint{2.831483in}{2.644559in}}%
\pgfpathlineto{\pgfqpoint{2.833847in}{2.644562in}}%
\pgfpathlineto{\pgfqpoint{2.833850in}{2.769443in}}%
\pgfpathlineto{\pgfqpoint{2.833953in}{2.629533in}}%
\pgfpathlineto{\pgfqpoint{2.834591in}{2.721502in}}%
\pgfpathlineto{\pgfqpoint{2.834641in}{2.672212in}}%
\pgfpathlineto{\pgfqpoint{2.835284in}{2.764901in}}%
\pgfpathlineto{\pgfqpoint{2.835333in}{2.715803in}}%
\pgfpathlineto{\pgfqpoint{2.835774in}{2.795652in}}%
\pgfpathlineto{\pgfqpoint{2.835913in}{2.644559in}}%
\pgfpathlineto{\pgfqpoint{2.836038in}{2.644559in}}%
\pgfpathlineto{\pgfqpoint{2.838793in}{2.644723in}}%
\pgfpathlineto{\pgfqpoint{2.838797in}{2.756871in}}%
\pgfpathlineto{\pgfqpoint{2.838878in}{2.627307in}}%
\pgfpathlineto{\pgfqpoint{2.839564in}{2.668797in}}%
\pgfpathlineto{\pgfqpoint{2.840404in}{2.776285in}}%
\pgfpathlineto{\pgfqpoint{2.840765in}{2.798880in}}%
\pgfpathlineto{\pgfqpoint{2.841301in}{2.644559in}}%
\pgfpathlineto{\pgfqpoint{2.843738in}{2.644561in}}%
\pgfpathlineto{\pgfqpoint{2.843740in}{2.753318in}}%
\pgfpathlineto{\pgfqpoint{2.843805in}{2.621473in}}%
\pgfpathlineto{\pgfqpoint{2.844536in}{2.733212in}}%
\pgfpathlineto{\pgfqpoint{2.844585in}{2.666989in}}%
\pgfpathlineto{\pgfqpoint{2.845228in}{2.776560in}}%
\pgfpathlineto{\pgfqpoint{2.845278in}{2.710630in}}%
\pgfpathlineto{\pgfqpoint{2.845711in}{2.806761in}}%
\pgfpathlineto{\pgfqpoint{2.845797in}{2.644559in}}%
\pgfpathlineto{\pgfqpoint{2.846003in}{2.644559in}}%
\pgfpathlineto{\pgfqpoint{2.848682in}{2.644559in}}%
\pgfpathlineto{\pgfqpoint{2.848688in}{2.750861in}}%
\pgfpathlineto{\pgfqpoint{2.848764in}{2.638126in}}%
\pgfpathlineto{\pgfqpoint{2.849450in}{2.679220in}}%
\pgfpathlineto{\pgfqpoint{2.850291in}{2.766334in}}%
\pgfpathlineto{\pgfqpoint{2.850656in}{2.789248in}}%
\pgfpathlineto{\pgfqpoint{2.851254in}{2.644559in}}%
\pgfpathlineto{\pgfqpoint{2.853627in}{2.644559in}}%
\pgfpathlineto{\pgfqpoint{2.853632in}{2.756608in}}%
\pgfpathlineto{\pgfqpoint{2.853693in}{2.642580in}}%
\pgfpathlineto{\pgfqpoint{2.854370in}{2.680850in}}%
\pgfpathlineto{\pgfqpoint{2.855211in}{2.761814in}}%
\pgfpathlineto{\pgfqpoint{2.855602in}{2.786340in}}%
\pgfpathlineto{\pgfqpoint{2.856211in}{2.644559in}}%
\pgfpathlineto{\pgfqpoint{2.858574in}{2.644575in}}%
\pgfpathlineto{\pgfqpoint{2.858577in}{2.776645in}}%
\pgfpathlineto{\pgfqpoint{2.859353in}{2.706299in}}%
\pgfpathlineto{\pgfqpoint{2.859402in}{2.691566in}}%
\pgfpathlineto{\pgfqpoint{2.859946in}{2.743554in}}%
\pgfpathlineto{\pgfqpoint{2.859996in}{2.728874in}}%
\pgfpathlineto{\pgfqpoint{2.860519in}{2.779485in}}%
\pgfpathlineto{\pgfqpoint{2.860608in}{2.696221in}}%
\pgfpathlineto{\pgfqpoint{2.860624in}{2.725206in}}%
\pgfpathlineto{\pgfqpoint{2.861408in}{2.644559in}}%
\pgfpathlineto{\pgfqpoint{2.863519in}{2.644559in}}%
\pgfpathlineto{\pgfqpoint{2.863524in}{2.749334in}}%
\pgfpathlineto{\pgfqpoint{2.864267in}{2.689816in}}%
\pgfpathlineto{\pgfqpoint{2.865306in}{2.764895in}}%
\pgfpathlineto{\pgfqpoint{2.865493in}{2.776646in}}%
\pgfpathlineto{\pgfqpoint{2.866959in}{2.644559in}}%
\pgfpathlineto{\pgfqpoint{2.868465in}{2.644559in}}%
\pgfpathlineto{\pgfqpoint{2.868469in}{2.766653in}}%
\pgfpathlineto{\pgfqpoint{2.868550in}{2.634277in}}%
\pgfpathlineto{\pgfqpoint{2.869235in}{2.675649in}}%
\pgfpathlineto{\pgfqpoint{2.870076in}{2.771066in}}%
\pgfpathlineto{\pgfqpoint{2.870438in}{2.793743in}}%
\pgfpathlineto{\pgfqpoint{2.871006in}{2.644559in}}%
\pgfpathlineto{\pgfqpoint{2.873410in}{2.644559in}}%
\pgfpathlineto{\pgfqpoint{2.873415in}{2.749225in}}%
\pgfpathlineto{\pgfqpoint{2.873494in}{2.633249in}}%
\pgfpathlineto{\pgfqpoint{2.874180in}{2.674612in}}%
\pgfpathlineto{\pgfqpoint{2.875021in}{2.771108in}}%
\pgfpathlineto{\pgfqpoint{2.875384in}{2.793836in}}%
\pgfpathlineto{\pgfqpoint{2.875946in}{2.644559in}}%
\pgfpathlineto{\pgfqpoint{2.878354in}{2.644559in}}%
\pgfpathlineto{\pgfqpoint{2.878360in}{2.756597in}}%
\pgfpathlineto{\pgfqpoint{2.878428in}{2.637653in}}%
\pgfpathlineto{\pgfqpoint{2.879113in}{2.677754in}}%
\pgfpathlineto{\pgfqpoint{2.879954in}{2.766745in}}%
\pgfpathlineto{\pgfqpoint{2.880329in}{2.790266in}}%
\pgfpathlineto{\pgfqpoint{2.880905in}{2.644559in}}%
\pgfpathlineto{\pgfqpoint{2.883301in}{2.644559in}}%
\pgfpathlineto{\pgfqpoint{2.883306in}{2.749471in}}%
\pgfpathlineto{\pgfqpoint{2.884061in}{2.692379in}}%
\pgfpathlineto{\pgfqpoint{2.885199in}{2.770663in}}%
\pgfpathlineto{\pgfqpoint{2.885329in}{2.716775in}}%
\pgfpathlineto{\pgfqpoint{2.885274in}{2.775429in}}%
\pgfpathlineto{\pgfqpoint{2.885343in}{2.729020in}}%
\pgfpathlineto{\pgfqpoint{2.886132in}{2.644559in}}%
\pgfpathlineto{\pgfqpoint{2.888246in}{2.644559in}}%
\pgfpathlineto{\pgfqpoint{2.888250in}{2.752038in}}%
\pgfpathlineto{\pgfqpoint{2.888330in}{2.636086in}}%
\pgfpathlineto{\pgfqpoint{2.889015in}{2.677326in}}%
\pgfpathlineto{\pgfqpoint{2.889856in}{2.769131in}}%
\pgfpathlineto{\pgfqpoint{2.890220in}{2.791928in}}%
\pgfpathlineto{\pgfqpoint{2.890797in}{2.644559in}}%
\pgfpathlineto{\pgfqpoint{2.893192in}{2.644559in}}%
\pgfpathlineto{\pgfqpoint{2.893197in}{2.749269in}}%
\pgfpathlineto{\pgfqpoint{2.893957in}{2.699286in}}%
\pgfpathlineto{\pgfqpoint{2.894006in}{2.695448in}}%
\pgfpathlineto{\pgfqpoint{2.894352in}{2.724146in}}%
\pgfpathlineto{\pgfqpoint{2.894402in}{2.720297in}}%
\pgfpathlineto{\pgfqpoint{2.895130in}{2.772996in}}%
\pgfpathlineto{\pgfqpoint{2.895207in}{2.734190in}}%
\pgfpathlineto{\pgfqpoint{2.896043in}{2.644559in}}%
\pgfpathlineto{\pgfqpoint{2.898137in}{2.644559in}}%
\pgfpathlineto{\pgfqpoint{2.898142in}{2.748816in}}%
\pgfpathlineto{\pgfqpoint{2.898263in}{2.632718in}}%
\pgfpathlineto{\pgfqpoint{2.898906in}{2.721108in}}%
\pgfpathlineto{\pgfqpoint{2.898956in}{2.676005in}}%
\pgfpathlineto{\pgfqpoint{2.899599in}{2.764531in}}%
\pgfpathlineto{\pgfqpoint{2.899648in}{2.719571in}}%
\pgfpathlineto{\pgfqpoint{2.900077in}{2.794545in}}%
\pgfpathlineto{\pgfqpoint{2.900203in}{2.644559in}}%
\pgfpathlineto{\pgfqpoint{2.900346in}{2.644559in}}%
\pgfpathlineto{\pgfqpoint{2.903082in}{2.644559in}}%
\pgfpathlineto{\pgfqpoint{2.903088in}{2.749387in}}%
\pgfpathlineto{\pgfqpoint{2.903838in}{2.698402in}}%
\pgfpathlineto{\pgfqpoint{2.903887in}{2.695257in}}%
\pgfpathlineto{\pgfqpoint{2.904135in}{2.717045in}}%
\pgfpathlineto{\pgfqpoint{2.904184in}{2.713896in}}%
\pgfpathlineto{\pgfqpoint{2.905016in}{2.772408in}}%
\pgfpathlineto{\pgfqpoint{2.905081in}{2.757311in}}%
\pgfpathlineto{\pgfqpoint{2.905883in}{2.644559in}}%
\pgfpathlineto{\pgfqpoint{2.908028in}{2.644559in}}%
\pgfpathlineto{\pgfqpoint{2.908033in}{2.748943in}}%
\pgfpathlineto{\pgfqpoint{2.908148in}{2.638413in}}%
\pgfpathlineto{\pgfqpoint{2.908791in}{2.714573in}}%
\pgfpathlineto{\pgfqpoint{2.908840in}{2.681637in}}%
\pgfpathlineto{\pgfqpoint{2.909483in}{2.758020in}}%
\pgfpathlineto{\pgfqpoint{2.909533in}{2.725180in}}%
\pgfpathlineto{\pgfqpoint{2.909965in}{2.788243in}}%
\pgfpathlineto{\pgfqpoint{2.910088in}{2.644559in}}%
\pgfpathlineto{\pgfqpoint{2.910202in}{2.644559in}}%
\pgfpathlineto{\pgfqpoint{2.912973in}{2.644559in}}%
\pgfpathlineto{\pgfqpoint{2.912979in}{2.749560in}}%
\pgfpathlineto{\pgfqpoint{2.913101in}{2.643428in}}%
\pgfpathlineto{\pgfqpoint{2.913744in}{2.709657in}}%
\pgfpathlineto{\pgfqpoint{2.913794in}{2.686712in}}%
\pgfpathlineto{\pgfqpoint{2.914338in}{2.746904in}}%
\pgfpathlineto{\pgfqpoint{2.914387in}{2.724028in}}%
\pgfpathlineto{\pgfqpoint{2.914914in}{2.783095in}}%
\pgfpathlineto{\pgfqpoint{2.915002in}{2.697211in}}%
\pgfpathlineto{\pgfqpoint{2.915017in}{2.748709in}}%
\pgfpathlineto{\pgfqpoint{2.915804in}{2.644559in}}%
\pgfpathlineto{\pgfqpoint{2.917918in}{2.644559in}}%
\pgfpathlineto{\pgfqpoint{2.917924in}{2.753609in}}%
\pgfpathlineto{\pgfqpoint{2.918003in}{2.638387in}}%
\pgfpathlineto{\pgfqpoint{2.918689in}{2.679702in}}%
\pgfpathlineto{\pgfqpoint{2.919530in}{2.766577in}}%
\pgfpathlineto{\pgfqpoint{2.919893in}{2.789318in}}%
\pgfpathlineto{\pgfqpoint{2.920495in}{2.644559in}}%
\pgfpathlineto{\pgfqpoint{2.922864in}{2.644559in}}%
\pgfpathlineto{\pgfqpoint{2.922870in}{2.748301in}}%
\pgfpathlineto{\pgfqpoint{2.922946in}{2.637392in}}%
\pgfpathlineto{\pgfqpoint{2.923631in}{2.678454in}}%
\pgfpathlineto{\pgfqpoint{2.924472in}{2.766934in}}%
\pgfpathlineto{\pgfqpoint{2.924838in}{2.789867in}}%
\pgfpathlineto{\pgfqpoint{2.925428in}{2.644559in}}%
\pgfpathlineto{\pgfqpoint{2.927809in}{2.644559in}}%
\pgfpathlineto{\pgfqpoint{2.927815in}{2.749934in}}%
\pgfpathlineto{\pgfqpoint{2.927885in}{2.634331in}}%
\pgfpathlineto{\pgfqpoint{2.928572in}{2.674756in}}%
\pgfpathlineto{\pgfqpoint{2.929413in}{2.770068in}}%
\pgfpathlineto{\pgfqpoint{2.929784in}{2.793296in}}%
\pgfpathlineto{\pgfqpoint{2.930341in}{2.644559in}}%
\pgfpathlineto{\pgfqpoint{2.932756in}{2.644958in}}%
\pgfpathlineto{\pgfqpoint{2.932761in}{2.749455in}}%
\pgfpathlineto{\pgfqpoint{2.932835in}{2.618960in}}%
\pgfpathlineto{\pgfqpoint{2.933519in}{2.659842in}}%
\pgfpathlineto{\pgfqpoint{2.934360in}{2.785784in}}%
\pgfpathlineto{\pgfqpoint{2.934729in}{2.808855in}}%
\pgfpathlineto{\pgfqpoint{2.935228in}{2.644559in}}%
\pgfpathlineto{\pgfqpoint{2.937700in}{2.644559in}}%
\pgfpathlineto{\pgfqpoint{2.937706in}{2.749048in}}%
\pgfpathlineto{\pgfqpoint{2.937774in}{2.644100in}}%
\pgfpathlineto{\pgfqpoint{2.938460in}{2.684344in}}%
\pgfpathlineto{\pgfqpoint{2.939300in}{2.760331in}}%
\pgfpathlineto{\pgfqpoint{2.939674in}{2.783793in}}%
\pgfpathlineto{\pgfqpoint{2.940357in}{2.644559in}}%
\pgfpathlineto{\pgfqpoint{2.942646in}{2.644559in}}%
\pgfpathlineto{\pgfqpoint{2.942651in}{2.748419in}}%
\pgfpathlineto{\pgfqpoint{2.942728in}{2.637054in}}%
\pgfpathlineto{\pgfqpoint{2.943414in}{2.678158in}}%
\pgfpathlineto{\pgfqpoint{2.944254in}{2.767523in}}%
\pgfpathlineto{\pgfqpoint{2.944620in}{2.790427in}}%
\pgfpathlineto{\pgfqpoint{2.945206in}{2.644559in}}%
\pgfpathlineto{\pgfqpoint{2.947592in}{2.644559in}}%
\pgfpathlineto{\pgfqpoint{2.947597in}{2.749316in}}%
\pgfpathlineto{\pgfqpoint{2.948331in}{2.691085in}}%
\pgfpathlineto{\pgfqpoint{2.949468in}{2.769134in}}%
\pgfpathlineto{\pgfqpoint{2.949620in}{2.719200in}}%
\pgfpathlineto{\pgfqpoint{2.949565in}{2.775219in}}%
\pgfpathlineto{\pgfqpoint{2.949633in}{2.726341in}}%
\pgfpathlineto{\pgfqpoint{2.950418in}{2.644559in}}%
\pgfpathlineto{\pgfqpoint{2.952537in}{2.644559in}}%
\pgfpathlineto{\pgfqpoint{2.952542in}{2.750459in}}%
\pgfpathlineto{\pgfqpoint{2.952628in}{2.631853in}}%
\pgfpathlineto{\pgfqpoint{2.953264in}{2.718032in}}%
\pgfpathlineto{\pgfqpoint{2.953314in}{2.673648in}}%
\pgfpathlineto{\pgfqpoint{2.953956in}{2.761407in}}%
\pgfpathlineto{\pgfqpoint{2.954006in}{2.717262in}}%
\pgfpathlineto{\pgfqpoint{2.954511in}{2.796118in}}%
\pgfpathlineto{\pgfqpoint{2.954597in}{2.644559in}}%
\pgfpathlineto{\pgfqpoint{2.954703in}{2.644559in}}%
\pgfpathlineto{\pgfqpoint{2.957482in}{2.644559in}}%
\pgfpathlineto{\pgfqpoint{2.957488in}{2.749621in}}%
\pgfpathlineto{\pgfqpoint{2.957597in}{2.642938in}}%
\pgfpathlineto{\pgfqpoint{2.958238in}{2.708574in}}%
\pgfpathlineto{\pgfqpoint{2.958287in}{2.685919in}}%
\pgfpathlineto{\pgfqpoint{2.958831in}{2.745822in}}%
\pgfpathlineto{\pgfqpoint{2.958881in}{2.723234in}}%
\pgfpathlineto{\pgfqpoint{2.959416in}{2.782506in}}%
\pgfpathlineto{\pgfqpoint{2.959511in}{2.697619in}}%
\pgfpathlineto{\pgfqpoint{2.959525in}{2.747759in}}%
\pgfpathlineto{\pgfqpoint{2.960313in}{2.644559in}}%
\pgfpathlineto{\pgfqpoint{2.962427in}{2.644559in}}%
\pgfpathlineto{\pgfqpoint{2.962433in}{2.750760in}}%
\pgfpathlineto{\pgfqpoint{2.962509in}{2.638342in}}%
\pgfpathlineto{\pgfqpoint{2.963195in}{2.679389in}}%
\pgfpathlineto{\pgfqpoint{2.964035in}{2.766313in}}%
\pgfpathlineto{\pgfqpoint{2.964402in}{2.789260in}}%
\pgfpathlineto{\pgfqpoint{2.965000in}{2.644559in}}%
\pgfpathlineto{\pgfqpoint{2.967373in}{2.644559in}}%
\pgfpathlineto{\pgfqpoint{2.967379in}{2.749719in}}%
\pgfpathlineto{\pgfqpoint{2.967448in}{2.643011in}}%
\pgfpathlineto{\pgfqpoint{2.968137in}{2.683448in}}%
\pgfpathlineto{\pgfqpoint{2.968978in}{2.761994in}}%
\pgfpathlineto{\pgfqpoint{2.969347in}{2.785159in}}%
\pgfpathlineto{\pgfqpoint{2.970003in}{2.644559in}}%
\pgfpathlineto{\pgfqpoint{2.972317in}{2.644559in}}%
\pgfpathlineto{\pgfqpoint{2.972324in}{2.757201in}}%
\pgfpathlineto{\pgfqpoint{2.973086in}{2.702392in}}%
\pgfpathlineto{\pgfqpoint{2.973136in}{2.693546in}}%
\pgfpathlineto{\pgfqpoint{2.973581in}{2.733429in}}%
\pgfpathlineto{\pgfqpoint{2.973630in}{2.724645in}}%
\pgfpathlineto{\pgfqpoint{2.974258in}{2.775939in}}%
\pgfpathlineto{\pgfqpoint{2.974344in}{2.737030in}}%
\pgfpathlineto{\pgfqpoint{2.975148in}{2.644559in}}%
\pgfpathlineto{\pgfqpoint{2.977264in}{2.644559in}}%
\pgfpathlineto{\pgfqpoint{2.977270in}{2.749599in}}%
\pgfpathlineto{\pgfqpoint{2.977381in}{2.642334in}}%
\pgfpathlineto{\pgfqpoint{2.978024in}{2.709452in}}%
\pgfpathlineto{\pgfqpoint{2.978073in}{2.685459in}}%
\pgfpathlineto{\pgfqpoint{2.978617in}{2.746699in}}%
\pgfpathlineto{\pgfqpoint{2.978667in}{2.722775in}}%
\pgfpathlineto{\pgfqpoint{2.979200in}{2.783245in}}%
\pgfpathlineto{\pgfqpoint{2.979293in}{2.697207in}}%
\pgfpathlineto{\pgfqpoint{2.979308in}{2.748675in}}%
\pgfpathlineto{\pgfqpoint{2.980095in}{2.644559in}}%
\pgfpathlineto{\pgfqpoint{2.982210in}{2.644559in}}%
\pgfpathlineto{\pgfqpoint{2.982214in}{2.748343in}}%
\pgfpathlineto{\pgfqpoint{2.982297in}{2.629914in}}%
\pgfpathlineto{\pgfqpoint{2.982934in}{2.719549in}}%
\pgfpathlineto{\pgfqpoint{2.982984in}{2.671467in}}%
\pgfpathlineto{\pgfqpoint{2.983627in}{2.762958in}}%
\pgfpathlineto{\pgfqpoint{2.983676in}{2.715047in}}%
\pgfpathlineto{\pgfqpoint{2.984184in}{2.797858in}}%
\pgfpathlineto{\pgfqpoint{2.984278in}{2.644559in}}%
\pgfpathlineto{\pgfqpoint{2.984380in}{2.644559in}}%
\pgfpathlineto{\pgfqpoint{2.987155in}{2.644559in}}%
\pgfpathlineto{\pgfqpoint{2.987161in}{2.749562in}}%
\pgfpathlineto{\pgfqpoint{2.987242in}{2.643352in}}%
\pgfpathlineto{\pgfqpoint{2.987883in}{2.705779in}}%
\pgfpathlineto{\pgfqpoint{2.987933in}{2.684989in}}%
\pgfpathlineto{\pgfqpoint{2.988477in}{2.743034in}}%
\pgfpathlineto{\pgfqpoint{2.988526in}{2.722297in}}%
\pgfpathlineto{\pgfqpoint{2.989129in}{2.783959in}}%
\pgfpathlineto{\pgfqpoint{2.989183in}{2.720319in}}%
\pgfpathlineto{\pgfqpoint{2.989197in}{2.725514in}}%
\pgfpathlineto{\pgfqpoint{2.989984in}{2.644559in}}%
\pgfpathlineto{\pgfqpoint{2.992100in}{2.644559in}}%
\pgfpathlineto{\pgfqpoint{2.992106in}{2.749235in}}%
\pgfpathlineto{\pgfqpoint{2.992187in}{2.634956in}}%
\pgfpathlineto{\pgfqpoint{2.992872in}{2.676403in}}%
\pgfpathlineto{\pgfqpoint{2.993713in}{2.770300in}}%
\pgfpathlineto{\pgfqpoint{2.994074in}{2.792935in}}%
\pgfpathlineto{\pgfqpoint{2.994646in}{2.644559in}}%
\pgfpathlineto{\pgfqpoint{2.997046in}{2.644559in}}%
\pgfpathlineto{\pgfqpoint{2.997051in}{2.749444in}}%
\pgfpathlineto{\pgfqpoint{2.997183in}{2.643478in}}%
\pgfpathlineto{\pgfqpoint{2.997777in}{2.680629in}}%
\pgfpathlineto{\pgfqpoint{2.998617in}{2.759976in}}%
\pgfpathlineto{\pgfqpoint{2.998992in}{2.783463in}}%
\pgfpathlineto{\pgfqpoint{2.999632in}{2.644559in}}%
\pgfpathlineto{\pgfqpoint{3.001992in}{2.644559in}}%
\pgfpathlineto{\pgfqpoint{3.001996in}{2.765627in}}%
\pgfpathlineto{\pgfqpoint{3.002066in}{2.642734in}}%
\pgfpathlineto{\pgfqpoint{3.002751in}{2.683058in}}%
\pgfpathlineto{\pgfqpoint{3.003592in}{2.762037in}}%
\pgfpathlineto{\pgfqpoint{3.003965in}{2.785433in}}%
\pgfpathlineto{\pgfqpoint{3.004617in}{2.644559in}}%
\pgfpathlineto{\pgfqpoint{3.006936in}{2.644559in}}%
\pgfpathlineto{\pgfqpoint{3.006942in}{2.755728in}}%
\pgfpathlineto{\pgfqpoint{3.007020in}{2.636153in}}%
\pgfpathlineto{\pgfqpoint{3.007706in}{2.677325in}}%
\pgfpathlineto{\pgfqpoint{3.008546in}{2.767832in}}%
\pgfpathlineto{\pgfqpoint{3.008911in}{2.790655in}}%
\pgfpathlineto{\pgfqpoint{3.009493in}{2.644559in}}%
\pgfpathlineto{\pgfqpoint{3.011883in}{2.644560in}}%
\pgfpathlineto{\pgfqpoint{3.011885in}{2.750229in}}%
\pgfpathlineto{\pgfqpoint{3.011952in}{2.620923in}}%
\pgfpathlineto{\pgfqpoint{3.012684in}{2.733706in}}%
\pgfpathlineto{\pgfqpoint{3.012734in}{2.666747in}}%
\pgfpathlineto{\pgfqpoint{3.013377in}{2.777051in}}%
\pgfpathlineto{\pgfqpoint{3.013426in}{2.710390in}}%
\pgfpathlineto{\pgfqpoint{3.013856in}{2.807077in}}%
\pgfpathlineto{\pgfqpoint{3.013941in}{2.644559in}}%
\pgfpathlineto{\pgfqpoint{3.014153in}{2.644559in}}%
\pgfpathlineto{\pgfqpoint{3.016827in}{2.644559in}}%
\pgfpathlineto{\pgfqpoint{3.016833in}{2.750670in}}%
\pgfpathlineto{\pgfqpoint{3.016908in}{2.639718in}}%
\pgfpathlineto{\pgfqpoint{3.017593in}{2.680627in}}%
\pgfpathlineto{\pgfqpoint{3.018434in}{2.764803in}}%
\pgfpathlineto{\pgfqpoint{3.018802in}{2.787851in}}%
\pgfpathlineto{\pgfqpoint{3.019416in}{2.644559in}}%
\pgfpathlineto{\pgfqpoint{3.021773in}{2.644559in}}%
\pgfpathlineto{\pgfqpoint{3.021779in}{2.749761in}}%
\pgfpathlineto{\pgfqpoint{3.022510in}{2.697730in}}%
\pgfpathlineto{\pgfqpoint{3.022560in}{2.694748in}}%
\pgfpathlineto{\pgfqpoint{3.022807in}{2.716374in}}%
\pgfpathlineto{\pgfqpoint{3.022856in}{2.713386in}}%
\pgfpathlineto{\pgfqpoint{3.023747in}{2.775439in}}%
\pgfpathlineto{\pgfqpoint{3.023772in}{2.758163in}}%
\pgfpathlineto{\pgfqpoint{3.024572in}{2.644559in}}%
\pgfpathlineto{\pgfqpoint{3.026719in}{2.644559in}}%
\pgfpathlineto{\pgfqpoint{3.026723in}{2.759600in}}%
\pgfpathlineto{\pgfqpoint{3.027455in}{2.682466in}}%
\pgfpathlineto{\pgfqpoint{3.028295in}{2.759419in}}%
\pgfpathlineto{\pgfqpoint{3.028667in}{2.782733in}}%
\pgfpathlineto{\pgfqpoint{3.029337in}{2.644559in}}%
\pgfpathlineto{\pgfqpoint{3.031664in}{2.644559in}}%
\pgfpathlineto{\pgfqpoint{3.031670in}{2.749236in}}%
\pgfpathlineto{\pgfqpoint{3.032430in}{2.692508in}}%
\pgfpathlineto{\pgfqpoint{3.033567in}{2.770681in}}%
\pgfpathlineto{\pgfqpoint{3.033692in}{2.718788in}}%
\pgfpathlineto{\pgfqpoint{3.033638in}{2.775146in}}%
\pgfpathlineto{\pgfqpoint{3.033706in}{2.726481in}}%
\pgfpathlineto{\pgfqpoint{3.034491in}{2.644559in}}%
\pgfpathlineto{\pgfqpoint{3.036609in}{2.644559in}}%
\pgfpathlineto{\pgfqpoint{3.036615in}{2.749567in}}%
\pgfpathlineto{\pgfqpoint{3.036725in}{2.628163in}}%
\pgfpathlineto{\pgfqpoint{3.037368in}{2.724359in}}%
\pgfpathlineto{\pgfqpoint{3.037417in}{2.671293in}}%
\pgfpathlineto{\pgfqpoint{3.038060in}{2.767752in}}%
\pgfpathlineto{\pgfqpoint{3.038109in}{2.714890in}}%
\pgfpathlineto{\pgfqpoint{3.038544in}{2.798099in}}%
\pgfpathlineto{\pgfqpoint{3.038676in}{2.644559in}}%
\pgfpathlineto{\pgfqpoint{3.038821in}{2.644559in}}%
\pgfpathlineto{\pgfqpoint{3.041555in}{2.644559in}}%
\pgfpathlineto{\pgfqpoint{3.041561in}{2.749108in}}%
\pgfpathlineto{\pgfqpoint{3.041625in}{2.640965in}}%
\pgfpathlineto{\pgfqpoint{3.042310in}{2.680602in}}%
\pgfpathlineto{\pgfqpoint{3.043151in}{2.762437in}}%
\pgfpathlineto{\pgfqpoint{3.043529in}{2.786162in}}%
\pgfpathlineto{\pgfqpoint{3.044147in}{2.644559in}}%
\pgfpathlineto{\pgfqpoint{3.046500in}{2.644559in}}%
\pgfpathlineto{\pgfqpoint{3.046506in}{2.750073in}}%
\pgfpathlineto{\pgfqpoint{3.046592in}{2.631321in}}%
\pgfpathlineto{\pgfqpoint{3.047228in}{2.718065in}}%
\pgfpathlineto{\pgfqpoint{3.047278in}{2.673136in}}%
\pgfpathlineto{\pgfqpoint{3.047920in}{2.761439in}}%
\pgfpathlineto{\pgfqpoint{3.047970in}{2.716752in}}%
\pgfpathlineto{\pgfqpoint{3.048474in}{2.796123in}}%
\pgfpathlineto{\pgfqpoint{3.048560in}{2.644559in}}%
\pgfpathlineto{\pgfqpoint{3.048671in}{2.644559in}}%
\pgfpathlineto{\pgfqpoint{3.051446in}{2.644559in}}%
\pgfpathlineto{\pgfqpoint{3.051451in}{2.749745in}}%
\pgfpathlineto{\pgfqpoint{3.052218in}{2.707985in}}%
\pgfpathlineto{\pgfqpoint{3.052267in}{2.688826in}}%
\pgfpathlineto{\pgfqpoint{3.052811in}{2.745243in}}%
\pgfpathlineto{\pgfqpoint{3.052861in}{2.726130in}}%
\pgfpathlineto{\pgfqpoint{3.053388in}{2.781421in}}%
\pgfpathlineto{\pgfqpoint{3.053475in}{2.716033in}}%
\pgfpathlineto{\pgfqpoint{3.053489in}{2.731308in}}%
\pgfpathlineto{\pgfqpoint{3.054273in}{2.644559in}}%
\pgfpathlineto{\pgfqpoint{3.056391in}{2.644559in}}%
\pgfpathlineto{\pgfqpoint{3.056396in}{2.747727in}}%
\pgfpathlineto{\pgfqpoint{3.056467in}{2.642172in}}%
\pgfpathlineto{\pgfqpoint{3.057153in}{2.682647in}}%
\pgfpathlineto{\pgfqpoint{3.057993in}{2.762442in}}%
\pgfpathlineto{\pgfqpoint{3.058365in}{2.785766in}}%
\pgfpathlineto{\pgfqpoint{3.059006in}{2.644559in}}%
\pgfpathlineto{\pgfqpoint{3.061337in}{2.644559in}}%
\pgfpathlineto{\pgfqpoint{3.061342in}{2.749394in}}%
\pgfpathlineto{\pgfqpoint{3.062104in}{2.699045in}}%
\pgfpathlineto{\pgfqpoint{3.062153in}{2.696074in}}%
\pgfpathlineto{\pgfqpoint{3.062401in}{2.717688in}}%
\pgfpathlineto{\pgfqpoint{3.062450in}{2.714713in}}%
\pgfpathlineto{\pgfqpoint{3.063276in}{2.772688in}}%
\pgfpathlineto{\pgfqpoint{3.063336in}{2.757400in}}%
\pgfpathlineto{\pgfqpoint{3.064135in}{2.644559in}}%
\pgfpathlineto{\pgfqpoint{3.066282in}{2.644559in}}%
\pgfpathlineto{\pgfqpoint{3.066288in}{2.749923in}}%
\pgfpathlineto{\pgfqpoint{3.066367in}{2.636044in}}%
\pgfpathlineto{\pgfqpoint{3.067053in}{2.677375in}}%
\pgfpathlineto{\pgfqpoint{3.067893in}{2.769293in}}%
\pgfpathlineto{\pgfqpoint{3.068256in}{2.792018in}}%
\pgfpathlineto{\pgfqpoint{3.068837in}{2.644559in}}%
\pgfpathlineto{\pgfqpoint{3.071228in}{2.644559in}}%
\pgfpathlineto{\pgfqpoint{3.071233in}{2.749359in}}%
\pgfpathlineto{\pgfqpoint{3.071994in}{2.699400in}}%
\pgfpathlineto{\pgfqpoint{3.072044in}{2.695667in}}%
\pgfpathlineto{\pgfqpoint{3.072390in}{2.724259in}}%
\pgfpathlineto{\pgfqpoint{3.072439in}{2.720516in}}%
\pgfpathlineto{\pgfqpoint{3.073167in}{2.773070in}}%
\pgfpathlineto{\pgfqpoint{3.073244in}{2.736612in}}%
\pgfpathlineto{\pgfqpoint{3.074068in}{2.644559in}}%
\pgfpathlineto{\pgfqpoint{3.076175in}{2.645184in}}%
\pgfpathlineto{\pgfqpoint{3.076180in}{2.746174in}}%
\pgfpathlineto{\pgfqpoint{3.076246in}{2.620397in}}%
\pgfpathlineto{\pgfqpoint{3.076931in}{2.660528in}}%
\pgfpathlineto{\pgfqpoint{3.077771in}{2.784410in}}%
\pgfpathlineto{\pgfqpoint{3.078147in}{2.807917in}}%
\pgfpathlineto{\pgfqpoint{3.078644in}{2.644559in}}%
\pgfpathlineto{\pgfqpoint{3.081119in}{2.644559in}}%
\pgfpathlineto{\pgfqpoint{3.081124in}{2.746990in}}%
\pgfpathlineto{\pgfqpoint{3.081179in}{2.644324in}}%
\pgfpathlineto{\pgfqpoint{3.081857in}{2.682136in}}%
\pgfpathlineto{\pgfqpoint{3.082697in}{2.759652in}}%
\pgfpathlineto{\pgfqpoint{3.083093in}{2.784449in}}%
\pgfpathlineto{\pgfqpoint{3.083732in}{2.644559in}}%
\pgfpathlineto{\pgfqpoint{3.086063in}{2.644559in}}%
\pgfpathlineto{\pgfqpoint{3.086069in}{2.753207in}}%
\pgfpathlineto{\pgfqpoint{3.086150in}{2.637044in}}%
\pgfpathlineto{\pgfqpoint{3.086836in}{2.678456in}}%
\pgfpathlineto{\pgfqpoint{3.087676in}{2.768010in}}%
\pgfpathlineto{\pgfqpoint{3.088038in}{2.790672in}}%
\pgfpathlineto{\pgfqpoint{3.088626in}{2.644559in}}%
\pgfpathlineto{\pgfqpoint{3.091010in}{2.644559in}}%
\pgfpathlineto{\pgfqpoint{3.091014in}{2.750124in}}%
\pgfpathlineto{\pgfqpoint{3.091097in}{2.627616in}}%
\pgfpathlineto{\pgfqpoint{3.091734in}{2.721027in}}%
\pgfpathlineto{\pgfqpoint{3.091783in}{2.669135in}}%
\pgfpathlineto{\pgfqpoint{3.092426in}{2.764429in}}%
\pgfpathlineto{\pgfqpoint{3.092476in}{2.712724in}}%
\pgfpathlineto{\pgfqpoint{3.092984in}{2.799348in}}%
\pgfpathlineto{\pgfqpoint{3.093074in}{2.644559in}}%
\pgfpathlineto{\pgfqpoint{3.093186in}{2.644559in}}%
\pgfpathlineto{\pgfqpoint{3.095955in}{2.644559in}}%
\pgfpathlineto{\pgfqpoint{3.095959in}{2.756866in}}%
\pgfpathlineto{\pgfqpoint{3.096038in}{2.636520in}}%
\pgfpathlineto{\pgfqpoint{3.096724in}{2.677691in}}%
\pgfpathlineto{\pgfqpoint{3.097564in}{2.768403in}}%
\pgfpathlineto{\pgfqpoint{3.097929in}{2.791247in}}%
\pgfpathlineto{\pgfqpoint{3.098511in}{2.644559in}}%
\pgfpathlineto{\pgfqpoint{3.100900in}{2.644559in}}%
\pgfpathlineto{\pgfqpoint{3.100906in}{2.750466in}}%
\pgfpathlineto{\pgfqpoint{3.101661in}{2.685636in}}%
\pgfpathlineto{\pgfqpoint{3.102502in}{2.759174in}}%
\pgfpathlineto{\pgfqpoint{3.102874in}{2.782549in}}%
\pgfpathlineto{\pgfqpoint{3.103594in}{2.644559in}}%
\pgfpathlineto{\pgfqpoint{3.105846in}{2.644559in}}%
\pgfpathlineto{\pgfqpoint{3.105851in}{2.749790in}}%
\pgfpathlineto{\pgfqpoint{3.106618in}{2.694119in}}%
\pgfpathlineto{\pgfqpoint{3.107820in}{2.775252in}}%
\pgfpathlineto{\pgfqpoint{3.108591in}{2.644559in}}%
\pgfpathlineto{\pgfqpoint{3.110793in}{2.644898in}}%
\pgfpathlineto{\pgfqpoint{3.110797in}{2.748952in}}%
\pgfpathlineto{\pgfqpoint{3.110875in}{2.619346in}}%
\pgfpathlineto{\pgfqpoint{3.111560in}{2.660600in}}%
\pgfpathlineto{\pgfqpoint{3.112400in}{2.785058in}}%
\pgfpathlineto{\pgfqpoint{3.112765in}{2.807887in}}%
\pgfpathlineto{\pgfqpoint{3.113273in}{2.644559in}}%
\pgfpathlineto{\pgfqpoint{3.115737in}{2.644559in}}%
\pgfpathlineto{\pgfqpoint{3.115742in}{2.749753in}}%
\pgfpathlineto{\pgfqpoint{3.116490in}{2.691035in}}%
\pgfpathlineto{\pgfqpoint{3.117528in}{2.765547in}}%
\pgfpathlineto{\pgfqpoint{3.117711in}{2.777009in}}%
\pgfpathlineto{\pgfqpoint{3.119266in}{2.644559in}}%
\pgfpathlineto{\pgfqpoint{3.120683in}{2.644561in}}%
\pgfpathlineto{\pgfqpoint{3.120685in}{2.752230in}}%
\pgfpathlineto{\pgfqpoint{3.120749in}{2.621787in}}%
\pgfpathlineto{\pgfqpoint{3.121481in}{2.732615in}}%
\pgfpathlineto{\pgfqpoint{3.121530in}{2.667232in}}%
\pgfpathlineto{\pgfqpoint{3.122173in}{2.775965in}}%
\pgfpathlineto{\pgfqpoint{3.122223in}{2.710870in}}%
\pgfpathlineto{\pgfqpoint{3.122656in}{2.806211in}}%
\pgfpathlineto{\pgfqpoint{3.122744in}{2.644559in}}%
\pgfpathlineto{\pgfqpoint{3.122946in}{2.644559in}}%
\pgfpathlineto{\pgfqpoint{3.125628in}{2.644559in}}%
\pgfpathlineto{\pgfqpoint{3.125633in}{2.749607in}}%
\pgfpathlineto{\pgfqpoint{3.126368in}{2.697544in}}%
\pgfpathlineto{\pgfqpoint{3.126418in}{2.694899in}}%
\pgfpathlineto{\pgfqpoint{3.126665in}{2.716188in}}%
\pgfpathlineto{\pgfqpoint{3.126715in}{2.713536in}}%
\pgfpathlineto{\pgfqpoint{3.127555in}{2.772117in}}%
\pgfpathlineto{\pgfqpoint{3.127626in}{2.757158in}}%
\pgfpathlineto{\pgfqpoint{3.128433in}{2.644559in}}%
\pgfpathlineto{\pgfqpoint{3.130575in}{2.644984in}}%
\pgfpathlineto{\pgfqpoint{3.130579in}{2.747765in}}%
\pgfpathlineto{\pgfqpoint{3.130650in}{2.619186in}}%
\pgfpathlineto{\pgfqpoint{3.131335in}{2.659838in}}%
\pgfpathlineto{\pgfqpoint{3.132176in}{2.785311in}}%
\pgfpathlineto{\pgfqpoint{3.132547in}{2.808534in}}%
\pgfpathlineto{\pgfqpoint{3.133046in}{2.644559in}}%
\pgfpathlineto{\pgfqpoint{3.135519in}{2.644559in}}%
\pgfpathlineto{\pgfqpoint{3.135524in}{2.749686in}}%
\pgfpathlineto{\pgfqpoint{3.135620in}{2.642189in}}%
\pgfpathlineto{\pgfqpoint{3.136260in}{2.708202in}}%
\pgfpathlineto{\pgfqpoint{3.136309in}{2.684625in}}%
\pgfpathlineto{\pgfqpoint{3.136853in}{2.745450in}}%
\pgfpathlineto{\pgfqpoint{3.136903in}{2.721940in}}%
\pgfpathlineto{\pgfqpoint{3.137447in}{2.782696in}}%
\pgfpathlineto{\pgfqpoint{3.137547in}{2.697977in}}%
\pgfpathlineto{\pgfqpoint{3.137562in}{2.747588in}}%
\pgfpathlineto{\pgfqpoint{3.138347in}{2.644559in}}%
\pgfpathlineto{\pgfqpoint{3.140465in}{2.644798in}}%
\pgfpathlineto{\pgfqpoint{3.140470in}{2.753421in}}%
\pgfpathlineto{\pgfqpoint{3.140550in}{2.622259in}}%
\pgfpathlineto{\pgfqpoint{3.141235in}{2.663705in}}%
\pgfpathlineto{\pgfqpoint{3.142076in}{2.782356in}}%
\pgfpathlineto{\pgfqpoint{3.142438in}{2.805026in}}%
\pgfpathlineto{\pgfqpoint{3.142953in}{2.644559in}}%
\pgfpathlineto{\pgfqpoint{3.145410in}{2.644559in}}%
\pgfpathlineto{\pgfqpoint{3.145415in}{2.749677in}}%
\pgfpathlineto{\pgfqpoint{3.145508in}{2.642134in}}%
\pgfpathlineto{\pgfqpoint{3.146148in}{2.708052in}}%
\pgfpathlineto{\pgfqpoint{3.146198in}{2.684466in}}%
\pgfpathlineto{\pgfqpoint{3.146742in}{2.745300in}}%
\pgfpathlineto{\pgfqpoint{3.146791in}{2.721781in}}%
\pgfpathlineto{\pgfqpoint{3.147335in}{2.782546in}}%
\pgfpathlineto{\pgfqpoint{3.147438in}{2.697709in}}%
\pgfpathlineto{\pgfqpoint{3.147453in}{2.747778in}}%
\pgfpathlineto{\pgfqpoint{3.148240in}{2.644559in}}%
\pgfpathlineto{\pgfqpoint{3.150356in}{2.644783in}}%
\pgfpathlineto{\pgfqpoint{3.150361in}{2.754253in}}%
\pgfpathlineto{\pgfqpoint{3.150441in}{2.622879in}}%
\pgfpathlineto{\pgfqpoint{3.151126in}{2.664346in}}%
\pgfpathlineto{\pgfqpoint{3.151967in}{2.781780in}}%
\pgfpathlineto{\pgfqpoint{3.152329in}{2.804429in}}%
\pgfpathlineto{\pgfqpoint{3.152848in}{2.644559in}}%
\pgfpathlineto{\pgfqpoint{3.155302in}{2.644806in}}%
\pgfpathlineto{\pgfqpoint{3.155306in}{2.753069in}}%
\pgfpathlineto{\pgfqpoint{3.155389in}{2.620921in}}%
\pgfpathlineto{\pgfqpoint{3.156025in}{2.727964in}}%
\pgfpathlineto{\pgfqpoint{3.156074in}{2.662590in}}%
\pgfpathlineto{\pgfqpoint{3.156717in}{2.771306in}}%
\pgfpathlineto{\pgfqpoint{3.156767in}{2.706237in}}%
\pgfpathlineto{\pgfqpoint{3.157274in}{2.806161in}}%
\pgfpathlineto{\pgfqpoint{3.157359in}{2.644559in}}%
\pgfpathlineto{\pgfqpoint{3.157494in}{2.644559in}}%
\pgfpathlineto{\pgfqpoint{3.160246in}{2.644559in}}%
\pgfpathlineto{\pgfqpoint{3.160250in}{2.747596in}}%
\pgfpathlineto{\pgfqpoint{3.160322in}{2.642396in}}%
\pgfpathlineto{\pgfqpoint{3.161007in}{2.682843in}}%
\pgfpathlineto{\pgfqpoint{3.161848in}{2.762250in}}%
\pgfpathlineto{\pgfqpoint{3.162220in}{2.785591in}}%
\pgfpathlineto{\pgfqpoint{3.162868in}{2.644559in}}%
\pgfpathlineto{\pgfqpoint{3.165192in}{2.644559in}}%
\pgfpathlineto{\pgfqpoint{3.165196in}{2.754672in}}%
\pgfpathlineto{\pgfqpoint{3.165286in}{2.627579in}}%
\pgfpathlineto{\pgfqpoint{3.165923in}{2.720968in}}%
\pgfpathlineto{\pgfqpoint{3.165973in}{2.669612in}}%
\pgfpathlineto{\pgfqpoint{3.166616in}{2.764368in}}%
\pgfpathlineto{\pgfqpoint{3.166665in}{2.713202in}}%
\pgfpathlineto{\pgfqpoint{3.167165in}{2.798795in}}%
\pgfpathlineto{\pgfqpoint{3.167255in}{2.644559in}}%
\pgfpathlineto{\pgfqpoint{3.167378in}{2.644559in}}%
\pgfpathlineto{\pgfqpoint{3.170132in}{2.644559in}}%
\pgfpathlineto{\pgfqpoint{3.170132in}{2.644559in}}%
\pgfusepath{stroke}%
\end{pgfscope}%
\begin{pgfscope}%
\pgfsetrectcap%
\pgfsetmiterjoin%
\pgfsetlinewidth{0.803000pt}%
\definecolor{currentstroke}{rgb}{0.000000,0.000000,0.000000}%
\pgfsetstrokecolor{currentstroke}%
\pgfsetdash{}{0pt}%
\pgfpathmoveto{\pgfqpoint{0.573768in}{2.601404in}}%
\pgfpathlineto{\pgfqpoint{0.573768in}{3.507654in}}%
\pgfusepath{stroke}%
\end{pgfscope}%
\begin{pgfscope}%
\pgfsetrectcap%
\pgfsetmiterjoin%
\pgfsetlinewidth{0.803000pt}%
\definecolor{currentstroke}{rgb}{0.000000,0.000000,0.000000}%
\pgfsetstrokecolor{currentstroke}%
\pgfsetdash{}{0pt}%
\pgfpathmoveto{\pgfqpoint{3.293768in}{2.601404in}}%
\pgfpathlineto{\pgfqpoint{3.293768in}{3.507654in}}%
\pgfusepath{stroke}%
\end{pgfscope}%
\begin{pgfscope}%
\pgfsetrectcap%
\pgfsetmiterjoin%
\pgfsetlinewidth{0.803000pt}%
\definecolor{currentstroke}{rgb}{0.000000,0.000000,0.000000}%
\pgfsetstrokecolor{currentstroke}%
\pgfsetdash{}{0pt}%
\pgfpathmoveto{\pgfqpoint{0.573768in}{2.601404in}}%
\pgfpathlineto{\pgfqpoint{3.293768in}{2.601404in}}%
\pgfusepath{stroke}%
\end{pgfscope}%
\begin{pgfscope}%
\pgfsetrectcap%
\pgfsetmiterjoin%
\pgfsetlinewidth{0.803000pt}%
\definecolor{currentstroke}{rgb}{0.000000,0.000000,0.000000}%
\pgfsetstrokecolor{currentstroke}%
\pgfsetdash{}{0pt}%
\pgfpathmoveto{\pgfqpoint{0.573768in}{3.507654in}}%
\pgfpathlineto{\pgfqpoint{3.293768in}{3.507654in}}%
\pgfusepath{stroke}%
\end{pgfscope}%
\begin{pgfscope}%
\definecolor{textcolor}{rgb}{0.000000,0.000000,0.000000}%
\pgfsetstrokecolor{textcolor}%
\pgfsetfillcolor{textcolor}%
\pgftext[x=0.628168in,y=3.444217in,left,top]{\color{textcolor}{\rmfamily\fontsize{8.000000}{9.600000}\selectfont\catcode`\^=\active\def^{\ifmmode\sp\else\^{}\fi}\catcode`\%=\active\def%{\%}(a)}}%
\end{pgfscope}%
\begin{pgfscope}%
\pgfsetbuttcap%
\pgfsetmiterjoin%
\definecolor{currentfill}{rgb}{1.000000,1.000000,1.000000}%
\pgfsetfillcolor{currentfill}%
\pgfsetlinewidth{0.000000pt}%
\definecolor{currentstroke}{rgb}{0.000000,0.000000,0.000000}%
\pgfsetstrokecolor{currentstroke}%
\pgfsetstrokeopacity{0.000000}%
\pgfsetdash{}{0pt}%
\pgfpathmoveto{\pgfqpoint{0.573768in}{1.532029in}}%
\pgfpathlineto{\pgfqpoint{3.293768in}{1.532029in}}%
\pgfpathlineto{\pgfqpoint{3.293768in}{2.438279in}}%
\pgfpathlineto{\pgfqpoint{0.573768in}{2.438279in}}%
\pgfpathlineto{\pgfqpoint{0.573768in}{1.532029in}}%
\pgfpathclose%
\pgfusepath{fill}%
\end{pgfscope}%
\begin{pgfscope}%
\pgfpathrectangle{\pgfqpoint{0.573768in}{1.532029in}}{\pgfqpoint{2.720000in}{0.906250in}}%
\pgfusepath{clip}%
\pgfsetbuttcap%
\pgfsetroundjoin%
\pgfsetlinewidth{0.501875pt}%
\definecolor{currentstroke}{rgb}{0.690196,0.690196,0.690196}%
\pgfsetstrokecolor{currentstroke}%
\pgfsetstrokeopacity{0.250000}%
\pgfsetdash{{1.850000pt}{0.800000pt}}{0.000000pt}%
\pgfpathmoveto{\pgfqpoint{0.697405in}{1.532029in}}%
\pgfpathlineto{\pgfqpoint{0.697405in}{2.438279in}}%
\pgfusepath{stroke}%
\end{pgfscope}%
\begin{pgfscope}%
\pgfsetbuttcap%
\pgfsetroundjoin%
\definecolor{currentfill}{rgb}{0.000000,0.000000,0.000000}%
\pgfsetfillcolor{currentfill}%
\pgfsetlinewidth{0.803000pt}%
\definecolor{currentstroke}{rgb}{0.000000,0.000000,0.000000}%
\pgfsetstrokecolor{currentstroke}%
\pgfsetdash{}{0pt}%
\pgfsys@defobject{currentmarker}{\pgfqpoint{0.000000in}{-0.048611in}}{\pgfqpoint{0.000000in}{0.000000in}}{%
\pgfpathmoveto{\pgfqpoint{0.000000in}{0.000000in}}%
\pgfpathlineto{\pgfqpoint{0.000000in}{-0.048611in}}%
\pgfusepath{stroke,fill}%
}%
\begin{pgfscope}%
\pgfsys@transformshift{0.697405in}{1.532029in}%
\pgfsys@useobject{currentmarker}{}%
\end{pgfscope}%
\end{pgfscope}%
\begin{pgfscope}%
\pgfpathrectangle{\pgfqpoint{0.573768in}{1.532029in}}{\pgfqpoint{2.720000in}{0.906250in}}%
\pgfusepath{clip}%
\pgfsetbuttcap%
\pgfsetroundjoin%
\pgfsetlinewidth{0.501875pt}%
\definecolor{currentstroke}{rgb}{0.690196,0.690196,0.690196}%
\pgfsetstrokecolor{currentstroke}%
\pgfsetstrokeopacity{0.250000}%
\pgfsetdash{{1.850000pt}{0.800000pt}}{0.000000pt}%
\pgfpathmoveto{\pgfqpoint{1.191950in}{1.532029in}}%
\pgfpathlineto{\pgfqpoint{1.191950in}{2.438279in}}%
\pgfusepath{stroke}%
\end{pgfscope}%
\begin{pgfscope}%
\pgfsetbuttcap%
\pgfsetroundjoin%
\definecolor{currentfill}{rgb}{0.000000,0.000000,0.000000}%
\pgfsetfillcolor{currentfill}%
\pgfsetlinewidth{0.803000pt}%
\definecolor{currentstroke}{rgb}{0.000000,0.000000,0.000000}%
\pgfsetstrokecolor{currentstroke}%
\pgfsetdash{}{0pt}%
\pgfsys@defobject{currentmarker}{\pgfqpoint{0.000000in}{-0.048611in}}{\pgfqpoint{0.000000in}{0.000000in}}{%
\pgfpathmoveto{\pgfqpoint{0.000000in}{0.000000in}}%
\pgfpathlineto{\pgfqpoint{0.000000in}{-0.048611in}}%
\pgfusepath{stroke,fill}%
}%
\begin{pgfscope}%
\pgfsys@transformshift{1.191950in}{1.532029in}%
\pgfsys@useobject{currentmarker}{}%
\end{pgfscope}%
\end{pgfscope}%
\begin{pgfscope}%
\pgfpathrectangle{\pgfqpoint{0.573768in}{1.532029in}}{\pgfqpoint{2.720000in}{0.906250in}}%
\pgfusepath{clip}%
\pgfsetbuttcap%
\pgfsetroundjoin%
\pgfsetlinewidth{0.501875pt}%
\definecolor{currentstroke}{rgb}{0.690196,0.690196,0.690196}%
\pgfsetstrokecolor{currentstroke}%
\pgfsetstrokeopacity{0.250000}%
\pgfsetdash{{1.850000pt}{0.800000pt}}{0.000000pt}%
\pgfpathmoveto{\pgfqpoint{1.686496in}{1.532029in}}%
\pgfpathlineto{\pgfqpoint{1.686496in}{2.438279in}}%
\pgfusepath{stroke}%
\end{pgfscope}%
\begin{pgfscope}%
\pgfsetbuttcap%
\pgfsetroundjoin%
\definecolor{currentfill}{rgb}{0.000000,0.000000,0.000000}%
\pgfsetfillcolor{currentfill}%
\pgfsetlinewidth{0.803000pt}%
\definecolor{currentstroke}{rgb}{0.000000,0.000000,0.000000}%
\pgfsetstrokecolor{currentstroke}%
\pgfsetdash{}{0pt}%
\pgfsys@defobject{currentmarker}{\pgfqpoint{0.000000in}{-0.048611in}}{\pgfqpoint{0.000000in}{0.000000in}}{%
\pgfpathmoveto{\pgfqpoint{0.000000in}{0.000000in}}%
\pgfpathlineto{\pgfqpoint{0.000000in}{-0.048611in}}%
\pgfusepath{stroke,fill}%
}%
\begin{pgfscope}%
\pgfsys@transformshift{1.686496in}{1.532029in}%
\pgfsys@useobject{currentmarker}{}%
\end{pgfscope}%
\end{pgfscope}%
\begin{pgfscope}%
\pgfpathrectangle{\pgfqpoint{0.573768in}{1.532029in}}{\pgfqpoint{2.720000in}{0.906250in}}%
\pgfusepath{clip}%
\pgfsetbuttcap%
\pgfsetroundjoin%
\pgfsetlinewidth{0.501875pt}%
\definecolor{currentstroke}{rgb}{0.690196,0.690196,0.690196}%
\pgfsetstrokecolor{currentstroke}%
\pgfsetstrokeopacity{0.250000}%
\pgfsetdash{{1.850000pt}{0.800000pt}}{0.000000pt}%
\pgfpathmoveto{\pgfqpoint{2.181041in}{1.532029in}}%
\pgfpathlineto{\pgfqpoint{2.181041in}{2.438279in}}%
\pgfusepath{stroke}%
\end{pgfscope}%
\begin{pgfscope}%
\pgfsetbuttcap%
\pgfsetroundjoin%
\definecolor{currentfill}{rgb}{0.000000,0.000000,0.000000}%
\pgfsetfillcolor{currentfill}%
\pgfsetlinewidth{0.803000pt}%
\definecolor{currentstroke}{rgb}{0.000000,0.000000,0.000000}%
\pgfsetstrokecolor{currentstroke}%
\pgfsetdash{}{0pt}%
\pgfsys@defobject{currentmarker}{\pgfqpoint{0.000000in}{-0.048611in}}{\pgfqpoint{0.000000in}{0.000000in}}{%
\pgfpathmoveto{\pgfqpoint{0.000000in}{0.000000in}}%
\pgfpathlineto{\pgfqpoint{0.000000in}{-0.048611in}}%
\pgfusepath{stroke,fill}%
}%
\begin{pgfscope}%
\pgfsys@transformshift{2.181041in}{1.532029in}%
\pgfsys@useobject{currentmarker}{}%
\end{pgfscope}%
\end{pgfscope}%
\begin{pgfscope}%
\pgfpathrectangle{\pgfqpoint{0.573768in}{1.532029in}}{\pgfqpoint{2.720000in}{0.906250in}}%
\pgfusepath{clip}%
\pgfsetbuttcap%
\pgfsetroundjoin%
\pgfsetlinewidth{0.501875pt}%
\definecolor{currentstroke}{rgb}{0.690196,0.690196,0.690196}%
\pgfsetstrokecolor{currentstroke}%
\pgfsetstrokeopacity{0.250000}%
\pgfsetdash{{1.850000pt}{0.800000pt}}{0.000000pt}%
\pgfpathmoveto{\pgfqpoint{2.675587in}{1.532029in}}%
\pgfpathlineto{\pgfqpoint{2.675587in}{2.438279in}}%
\pgfusepath{stroke}%
\end{pgfscope}%
\begin{pgfscope}%
\pgfsetbuttcap%
\pgfsetroundjoin%
\definecolor{currentfill}{rgb}{0.000000,0.000000,0.000000}%
\pgfsetfillcolor{currentfill}%
\pgfsetlinewidth{0.803000pt}%
\definecolor{currentstroke}{rgb}{0.000000,0.000000,0.000000}%
\pgfsetstrokecolor{currentstroke}%
\pgfsetdash{}{0pt}%
\pgfsys@defobject{currentmarker}{\pgfqpoint{0.000000in}{-0.048611in}}{\pgfqpoint{0.000000in}{0.000000in}}{%
\pgfpathmoveto{\pgfqpoint{0.000000in}{0.000000in}}%
\pgfpathlineto{\pgfqpoint{0.000000in}{-0.048611in}}%
\pgfusepath{stroke,fill}%
}%
\begin{pgfscope}%
\pgfsys@transformshift{2.675587in}{1.532029in}%
\pgfsys@useobject{currentmarker}{}%
\end{pgfscope}%
\end{pgfscope}%
\begin{pgfscope}%
\pgfpathrectangle{\pgfqpoint{0.573768in}{1.532029in}}{\pgfqpoint{2.720000in}{0.906250in}}%
\pgfusepath{clip}%
\pgfsetbuttcap%
\pgfsetroundjoin%
\pgfsetlinewidth{0.501875pt}%
\definecolor{currentstroke}{rgb}{0.690196,0.690196,0.690196}%
\pgfsetstrokecolor{currentstroke}%
\pgfsetstrokeopacity{0.250000}%
\pgfsetdash{{1.850000pt}{0.800000pt}}{0.000000pt}%
\pgfpathmoveto{\pgfqpoint{3.170132in}{1.532029in}}%
\pgfpathlineto{\pgfqpoint{3.170132in}{2.438279in}}%
\pgfusepath{stroke}%
\end{pgfscope}%
\begin{pgfscope}%
\pgfsetbuttcap%
\pgfsetroundjoin%
\definecolor{currentfill}{rgb}{0.000000,0.000000,0.000000}%
\pgfsetfillcolor{currentfill}%
\pgfsetlinewidth{0.803000pt}%
\definecolor{currentstroke}{rgb}{0.000000,0.000000,0.000000}%
\pgfsetstrokecolor{currentstroke}%
\pgfsetdash{}{0pt}%
\pgfsys@defobject{currentmarker}{\pgfqpoint{0.000000in}{-0.048611in}}{\pgfqpoint{0.000000in}{0.000000in}}{%
\pgfpathmoveto{\pgfqpoint{0.000000in}{0.000000in}}%
\pgfpathlineto{\pgfqpoint{0.000000in}{-0.048611in}}%
\pgfusepath{stroke,fill}%
}%
\begin{pgfscope}%
\pgfsys@transformshift{3.170132in}{1.532029in}%
\pgfsys@useobject{currentmarker}{}%
\end{pgfscope}%
\end{pgfscope}%
\begin{pgfscope}%
\pgfpathrectangle{\pgfqpoint{0.573768in}{1.532029in}}{\pgfqpoint{2.720000in}{0.906250in}}%
\pgfusepath{clip}%
\pgfsetbuttcap%
\pgfsetroundjoin%
\pgfsetlinewidth{0.501875pt}%
\definecolor{currentstroke}{rgb}{0.690196,0.690196,0.690196}%
\pgfsetstrokecolor{currentstroke}%
\pgfsetstrokeopacity{0.250000}%
\pgfsetdash{{1.850000pt}{0.800000pt}}{0.000000pt}%
\pgfpathmoveto{\pgfqpoint{0.573768in}{1.532029in}}%
\pgfpathlineto{\pgfqpoint{3.293768in}{1.532029in}}%
\pgfusepath{stroke}%
\end{pgfscope}%
\begin{pgfscope}%
\pgfsetbuttcap%
\pgfsetroundjoin%
\definecolor{currentfill}{rgb}{0.000000,0.000000,0.000000}%
\pgfsetfillcolor{currentfill}%
\pgfsetlinewidth{0.803000pt}%
\definecolor{currentstroke}{rgb}{0.000000,0.000000,0.000000}%
\pgfsetstrokecolor{currentstroke}%
\pgfsetdash{}{0pt}%
\pgfsys@defobject{currentmarker}{\pgfqpoint{-0.048611in}{0.000000in}}{\pgfqpoint{-0.000000in}{0.000000in}}{%
\pgfpathmoveto{\pgfqpoint{-0.000000in}{0.000000in}}%
\pgfpathlineto{\pgfqpoint{-0.048611in}{0.000000in}}%
\pgfusepath{stroke,fill}%
}%
\begin{pgfscope}%
\pgfsys@transformshift{0.573768in}{1.532029in}%
\pgfsys@useobject{currentmarker}{}%
\end{pgfscope}%
\end{pgfscope}%
\begin{pgfscope}%
\definecolor{textcolor}{rgb}{0.000000,0.000000,0.000000}%
\pgfsetstrokecolor{textcolor}%
\pgfsetfillcolor{textcolor}%
\pgftext[x=0.417518in, y=1.493449in, left, base]{\color{textcolor}{\rmfamily\fontsize{8.000000}{9.600000}\selectfont\catcode`\^=\active\def^{\ifmmode\sp\else\^{}\fi}\catcode`\%=\active\def%{\%}$\mathdefault{0}$}}%
\end{pgfscope}%
\begin{pgfscope}%
\pgfpathrectangle{\pgfqpoint{0.573768in}{1.532029in}}{\pgfqpoint{2.720000in}{0.906250in}}%
\pgfusepath{clip}%
\pgfsetbuttcap%
\pgfsetroundjoin%
\pgfsetlinewidth{0.501875pt}%
\definecolor{currentstroke}{rgb}{0.690196,0.690196,0.690196}%
\pgfsetstrokecolor{currentstroke}%
\pgfsetstrokeopacity{0.250000}%
\pgfsetdash{{1.850000pt}{0.800000pt}}{0.000000pt}%
\pgfpathmoveto{\pgfqpoint{0.573768in}{1.834113in}}%
\pgfpathlineto{\pgfqpoint{3.293768in}{1.834113in}}%
\pgfusepath{stroke}%
\end{pgfscope}%
\begin{pgfscope}%
\pgfsetbuttcap%
\pgfsetroundjoin%
\definecolor{currentfill}{rgb}{0.000000,0.000000,0.000000}%
\pgfsetfillcolor{currentfill}%
\pgfsetlinewidth{0.803000pt}%
\definecolor{currentstroke}{rgb}{0.000000,0.000000,0.000000}%
\pgfsetstrokecolor{currentstroke}%
\pgfsetdash{}{0pt}%
\pgfsys@defobject{currentmarker}{\pgfqpoint{-0.048611in}{0.000000in}}{\pgfqpoint{-0.000000in}{0.000000in}}{%
\pgfpathmoveto{\pgfqpoint{-0.000000in}{0.000000in}}%
\pgfpathlineto{\pgfqpoint{-0.048611in}{0.000000in}}%
\pgfusepath{stroke,fill}%
}%
\begin{pgfscope}%
\pgfsys@transformshift{0.573768in}{1.834113in}%
\pgfsys@useobject{currentmarker}{}%
\end{pgfscope}%
\end{pgfscope}%
\begin{pgfscope}%
\definecolor{textcolor}{rgb}{0.000000,0.000000,0.000000}%
\pgfsetstrokecolor{textcolor}%
\pgfsetfillcolor{textcolor}%
\pgftext[x=0.358489in, y=1.795533in, left, base]{\color{textcolor}{\rmfamily\fontsize{8.000000}{9.600000}\selectfont\catcode`\^=\active\def^{\ifmmode\sp\else\^{}\fi}\catcode`\%=\active\def%{\%}$\mathdefault{50}$}}%
\end{pgfscope}%
\begin{pgfscope}%
\pgfpathrectangle{\pgfqpoint{0.573768in}{1.532029in}}{\pgfqpoint{2.720000in}{0.906250in}}%
\pgfusepath{clip}%
\pgfsetbuttcap%
\pgfsetroundjoin%
\pgfsetlinewidth{0.501875pt}%
\definecolor{currentstroke}{rgb}{0.690196,0.690196,0.690196}%
\pgfsetstrokecolor{currentstroke}%
\pgfsetstrokeopacity{0.250000}%
\pgfsetdash{{1.850000pt}{0.800000pt}}{0.000000pt}%
\pgfpathmoveto{\pgfqpoint{0.573768in}{2.136196in}}%
\pgfpathlineto{\pgfqpoint{3.293768in}{2.136196in}}%
\pgfusepath{stroke}%
\end{pgfscope}%
\begin{pgfscope}%
\pgfsetbuttcap%
\pgfsetroundjoin%
\definecolor{currentfill}{rgb}{0.000000,0.000000,0.000000}%
\pgfsetfillcolor{currentfill}%
\pgfsetlinewidth{0.803000pt}%
\definecolor{currentstroke}{rgb}{0.000000,0.000000,0.000000}%
\pgfsetstrokecolor{currentstroke}%
\pgfsetdash{}{0pt}%
\pgfsys@defobject{currentmarker}{\pgfqpoint{-0.048611in}{0.000000in}}{\pgfqpoint{-0.000000in}{0.000000in}}{%
\pgfpathmoveto{\pgfqpoint{-0.000000in}{0.000000in}}%
\pgfpathlineto{\pgfqpoint{-0.048611in}{0.000000in}}%
\pgfusepath{stroke,fill}%
}%
\begin{pgfscope}%
\pgfsys@transformshift{0.573768in}{2.136196in}%
\pgfsys@useobject{currentmarker}{}%
\end{pgfscope}%
\end{pgfscope}%
\begin{pgfscope}%
\definecolor{textcolor}{rgb}{0.000000,0.000000,0.000000}%
\pgfsetstrokecolor{textcolor}%
\pgfsetfillcolor{textcolor}%
\pgftext[x=0.299460in, y=2.097616in, left, base]{\color{textcolor}{\rmfamily\fontsize{8.000000}{9.600000}\selectfont\catcode`\^=\active\def^{\ifmmode\sp\else\^{}\fi}\catcode`\%=\active\def%{\%}$\mathdefault{100}$}}%
\end{pgfscope}%
\begin{pgfscope}%
\pgfpathrectangle{\pgfqpoint{0.573768in}{1.532029in}}{\pgfqpoint{2.720000in}{0.906250in}}%
\pgfusepath{clip}%
\pgfsetbuttcap%
\pgfsetroundjoin%
\pgfsetlinewidth{0.501875pt}%
\definecolor{currentstroke}{rgb}{0.690196,0.690196,0.690196}%
\pgfsetstrokecolor{currentstroke}%
\pgfsetstrokeopacity{0.250000}%
\pgfsetdash{{1.850000pt}{0.800000pt}}{0.000000pt}%
\pgfpathmoveto{\pgfqpoint{0.573768in}{2.438279in}}%
\pgfpathlineto{\pgfqpoint{3.293768in}{2.438279in}}%
\pgfusepath{stroke}%
\end{pgfscope}%
\begin{pgfscope}%
\pgfsetbuttcap%
\pgfsetroundjoin%
\definecolor{currentfill}{rgb}{0.000000,0.000000,0.000000}%
\pgfsetfillcolor{currentfill}%
\pgfsetlinewidth{0.803000pt}%
\definecolor{currentstroke}{rgb}{0.000000,0.000000,0.000000}%
\pgfsetstrokecolor{currentstroke}%
\pgfsetdash{}{0pt}%
\pgfsys@defobject{currentmarker}{\pgfqpoint{-0.048611in}{0.000000in}}{\pgfqpoint{-0.000000in}{0.000000in}}{%
\pgfpathmoveto{\pgfqpoint{-0.000000in}{0.000000in}}%
\pgfpathlineto{\pgfqpoint{-0.048611in}{0.000000in}}%
\pgfusepath{stroke,fill}%
}%
\begin{pgfscope}%
\pgfsys@transformshift{0.573768in}{2.438279in}%
\pgfsys@useobject{currentmarker}{}%
\end{pgfscope}%
\end{pgfscope}%
\begin{pgfscope}%
\definecolor{textcolor}{rgb}{0.000000,0.000000,0.000000}%
\pgfsetstrokecolor{textcolor}%
\pgfsetfillcolor{textcolor}%
\pgftext[x=0.299460in, y=2.399699in, left, base]{\color{textcolor}{\rmfamily\fontsize{8.000000}{9.600000}\selectfont\catcode`\^=\active\def^{\ifmmode\sp\else\^{}\fi}\catcode`\%=\active\def%{\%}$\mathdefault{150}$}}%
\end{pgfscope}%
\begin{pgfscope}%
\definecolor{textcolor}{rgb}{0.000000,0.000000,0.000000}%
\pgfsetstrokecolor{textcolor}%
\pgfsetfillcolor{textcolor}%
\pgftext[x=0.243905in,y=1.985154in,,bottom,rotate=90.000000]{\color{textcolor}{\rmfamily\fontsize{8.000000}{9.600000}\selectfont\catcode`\^=\active\def^{\ifmmode\sp\else\^{}\fi}\catcode`\%=\active\def%{\%}$v_{DS}$ (V)}}%
\end{pgfscope}%
\begin{pgfscope}%
\pgfpathrectangle{\pgfqpoint{0.573768in}{1.532029in}}{\pgfqpoint{2.720000in}{0.906250in}}%
\pgfusepath{clip}%
\pgfsetrectcap%
\pgfsetroundjoin%
\pgfsetlinewidth{1.505625pt}%
\definecolor{currentstroke}{rgb}{0.835294,0.368627,0.000000}%
\pgfsetstrokecolor{currentstroke}%
\pgfsetdash{}{0pt}%
\pgfpathmoveto{\pgfqpoint{0.697405in}{1.532030in}}%
\pgfpathlineto{\pgfqpoint{0.697409in}{1.539351in}}%
\pgfpathlineto{\pgfqpoint{0.697413in}{1.546048in}}%
\pgfpathlineto{\pgfqpoint{0.697444in}{1.532029in}}%
\pgfpathlineto{\pgfqpoint{0.698134in}{1.532046in}}%
\pgfpathlineto{\pgfqpoint{0.699411in}{1.532297in}}%
\pgfpathlineto{\pgfqpoint{0.699451in}{1.839162in}}%
\pgfpathlineto{\pgfqpoint{0.700332in}{1.748579in}}%
\pgfpathlineto{\pgfqpoint{0.700372in}{1.731944in}}%
\pgfpathlineto{\pgfqpoint{0.700386in}{1.777799in}}%
\pgfpathlineto{\pgfqpoint{0.701098in}{1.748370in}}%
\pgfpathlineto{\pgfqpoint{0.701145in}{1.765965in}}%
\pgfpathlineto{\pgfqpoint{0.701128in}{1.747071in}}%
\pgfpathlineto{\pgfqpoint{0.701873in}{1.759656in}}%
\pgfpathlineto{\pgfqpoint{0.702355in}{1.767393in}}%
\pgfpathlineto{\pgfqpoint{0.702359in}{1.746242in}}%
\pgfpathlineto{\pgfqpoint{0.702451in}{1.532071in}}%
\pgfpathlineto{\pgfqpoint{0.703136in}{1.532086in}}%
\pgfpathlineto{\pgfqpoint{0.704356in}{1.532568in}}%
\pgfpathlineto{\pgfqpoint{0.704397in}{1.866724in}}%
\pgfpathlineto{\pgfqpoint{0.705276in}{1.761145in}}%
\pgfpathlineto{\pgfqpoint{0.705284in}{1.729847in}}%
\pgfpathlineto{\pgfqpoint{0.705333in}{1.790447in}}%
\pgfpathlineto{\pgfqpoint{0.706050in}{1.751018in}}%
\pgfpathlineto{\pgfqpoint{0.706096in}{1.773908in}}%
\pgfpathlineto{\pgfqpoint{0.706832in}{1.761875in}}%
\pgfpathlineto{\pgfqpoint{0.707396in}{1.532111in}}%
\pgfpathlineto{\pgfqpoint{0.707302in}{1.770985in}}%
\pgfpathlineto{\pgfqpoint{0.708326in}{1.532144in}}%
\pgfpathlineto{\pgfqpoint{0.709285in}{1.532182in}}%
\pgfpathlineto{\pgfqpoint{0.709294in}{1.532271in}}%
\pgfpathlineto{\pgfqpoint{0.709337in}{1.784396in}}%
\pgfpathlineto{\pgfqpoint{0.709350in}{1.983826in}}%
\pgfpathlineto{\pgfqpoint{0.709372in}{1.623530in}}%
\pgfpathlineto{\pgfqpoint{0.710101in}{1.734060in}}%
\pgfpathlineto{\pgfqpoint{0.710105in}{1.710481in}}%
\pgfpathlineto{\pgfqpoint{0.710153in}{1.828258in}}%
\pgfpathlineto{\pgfqpoint{0.710868in}{1.749807in}}%
\pgfpathlineto{\pgfqpoint{0.710884in}{1.795311in}}%
\pgfpathlineto{\pgfqpoint{0.711674in}{1.783025in}}%
\pgfpathlineto{\pgfqpoint{0.712377in}{1.532151in}}%
\pgfpathlineto{\pgfqpoint{0.712248in}{1.786666in}}%
\pgfpathlineto{\pgfqpoint{0.712574in}{1.532155in}}%
\pgfpathlineto{\pgfqpoint{0.714241in}{1.532394in}}%
\pgfpathlineto{\pgfqpoint{0.714300in}{2.034694in}}%
\pgfpathlineto{\pgfqpoint{0.715239in}{1.729242in}}%
\pgfpathlineto{\pgfqpoint{0.715243in}{1.714433in}}%
\pgfpathlineto{\pgfqpoint{0.715257in}{1.851125in}}%
\pgfpathlineto{\pgfqpoint{0.715984in}{1.813823in}}%
\pgfpathlineto{\pgfqpoint{0.716035in}{1.761016in}}%
\pgfpathlineto{\pgfqpoint{0.716765in}{1.800796in}}%
\pgfpathlineto{\pgfqpoint{0.716839in}{1.802238in}}%
\pgfpathlineto{\pgfqpoint{0.717206in}{1.560996in}}%
\pgfpathlineto{\pgfqpoint{0.717335in}{1.532190in}}%
\pgfpathlineto{\pgfqpoint{0.717978in}{1.532203in}}%
\pgfpathlineto{\pgfqpoint{0.719184in}{1.532365in}}%
\pgfpathlineto{\pgfqpoint{0.719210in}{1.576239in}}%
\pgfpathlineto{\pgfqpoint{0.719248in}{2.084561in}}%
\pgfpathlineto{\pgfqpoint{0.719274in}{1.555173in}}%
\pgfpathlineto{\pgfqpoint{0.720000in}{1.778282in}}%
\pgfpathlineto{\pgfqpoint{0.720008in}{1.713936in}}%
\pgfpathlineto{\pgfqpoint{0.720023in}{1.886795in}}%
\pgfpathlineto{\pgfqpoint{0.720764in}{1.778344in}}%
\pgfpathlineto{\pgfqpoint{0.720847in}{1.838290in}}%
\pgfpathlineto{\pgfqpoint{0.720794in}{1.765477in}}%
\pgfpathlineto{\pgfqpoint{0.721509in}{1.782316in}}%
\pgfpathlineto{\pgfqpoint{0.721905in}{1.833790in}}%
\pgfpathlineto{\pgfqpoint{0.722141in}{1.780529in}}%
\pgfpathlineto{\pgfqpoint{0.722322in}{1.532212in}}%
\pgfpathlineto{\pgfqpoint{0.722915in}{1.532225in}}%
\pgfpathlineto{\pgfqpoint{0.724132in}{1.532527in}}%
\pgfpathlineto{\pgfqpoint{0.724198in}{2.074410in}}%
\pgfpathlineto{\pgfqpoint{0.725123in}{1.749675in}}%
\pgfpathlineto{\pgfqpoint{0.725133in}{1.810859in}}%
\pgfpathlineto{\pgfqpoint{0.725141in}{1.890951in}}%
\pgfpathlineto{\pgfqpoint{0.725157in}{1.748262in}}%
\pgfpathlineto{\pgfqpoint{0.725903in}{1.821454in}}%
\pgfpathlineto{\pgfqpoint{0.725911in}{1.849640in}}%
\pgfpathlineto{\pgfqpoint{0.725956in}{1.796920in}}%
\pgfpathlineto{\pgfqpoint{0.726675in}{1.835104in}}%
\pgfpathlineto{\pgfqpoint{0.727082in}{1.841722in}}%
\pgfpathlineto{\pgfqpoint{0.727094in}{1.631972in}}%
\pgfpathlineto{\pgfqpoint{0.727255in}{1.532236in}}%
\pgfpathlineto{\pgfqpoint{0.727848in}{1.532249in}}%
\pgfpathlineto{\pgfqpoint{0.729078in}{1.532580in}}%
\pgfpathlineto{\pgfqpoint{0.729145in}{2.150842in}}%
\pgfpathlineto{\pgfqpoint{0.730127in}{1.800527in}}%
\pgfpathlineto{\pgfqpoint{0.730131in}{1.780678in}}%
\pgfpathlineto{\pgfqpoint{0.730147in}{1.894998in}}%
\pgfpathlineto{\pgfqpoint{0.730892in}{1.821588in}}%
\pgfpathlineto{\pgfqpoint{0.730909in}{1.864374in}}%
\pgfpathlineto{\pgfqpoint{0.731674in}{1.851392in}}%
\pgfpathlineto{\pgfqpoint{0.732199in}{1.532254in}}%
\pgfpathlineto{\pgfqpoint{0.732029in}{1.863845in}}%
\pgfpathlineto{\pgfqpoint{0.732595in}{1.532263in}}%
\pgfpathlineto{\pgfqpoint{0.734025in}{1.532721in}}%
\pgfpathlineto{\pgfqpoint{0.734093in}{2.150517in}}%
\pgfpathlineto{\pgfqpoint{0.735032in}{1.839784in}}%
\pgfpathlineto{\pgfqpoint{0.735069in}{1.792736in}}%
\pgfpathlineto{\pgfqpoint{0.735054in}{1.929285in}}%
\pgfpathlineto{\pgfqpoint{0.735796in}{1.840816in}}%
\pgfpathlineto{\pgfqpoint{0.735810in}{1.891037in}}%
\pgfpathlineto{\pgfqpoint{0.736572in}{1.867845in}}%
\pgfpathlineto{\pgfqpoint{0.737098in}{1.532266in}}%
\pgfpathlineto{\pgfqpoint{0.736593in}{1.880513in}}%
\pgfpathlineto{\pgfqpoint{0.737533in}{1.532282in}}%
\pgfpathlineto{\pgfqpoint{0.738955in}{1.532350in}}%
\pgfpathlineto{\pgfqpoint{0.738963in}{1.532449in}}%
\pgfpathlineto{\pgfqpoint{0.738990in}{1.572831in}}%
\pgfpathlineto{\pgfqpoint{0.739032in}{2.185605in}}%
\pgfpathlineto{\pgfqpoint{0.739785in}{1.834351in}}%
\pgfpathlineto{\pgfqpoint{0.739789in}{1.797017in}}%
\pgfpathlineto{\pgfqpoint{0.739807in}{1.963958in}}%
\pgfpathlineto{\pgfqpoint{0.740548in}{1.862980in}}%
\pgfpathlineto{\pgfqpoint{0.740621in}{1.916667in}}%
\pgfpathlineto{\pgfqpoint{0.740574in}{1.855280in}}%
\pgfpathlineto{\pgfqpoint{0.741320in}{1.882858in}}%
\pgfpathlineto{\pgfqpoint{0.741332in}{1.905277in}}%
\pgfpathlineto{\pgfqpoint{0.741940in}{1.552607in}}%
\pgfpathlineto{\pgfqpoint{0.742057in}{1.532269in}}%
\pgfpathlineto{\pgfqpoint{0.742692in}{1.532291in}}%
\pgfpathlineto{\pgfqpoint{0.743900in}{1.532351in}}%
\pgfpathlineto{\pgfqpoint{0.743908in}{1.532437in}}%
\pgfpathlineto{\pgfqpoint{0.743936in}{1.575413in}}%
\pgfpathlineto{\pgfqpoint{0.743979in}{2.208342in}}%
\pgfpathlineto{\pgfqpoint{0.744727in}{1.869304in}}%
\pgfpathlineto{\pgfqpoint{0.744730in}{1.824091in}}%
\pgfpathlineto{\pgfqpoint{0.744749in}{1.986471in}}%
\pgfpathlineto{\pgfqpoint{0.745489in}{1.928732in}}%
\pgfpathlineto{\pgfqpoint{0.745493in}{1.941181in}}%
\pgfpathlineto{\pgfqpoint{0.745508in}{1.878029in}}%
\pgfpathlineto{\pgfqpoint{0.746250in}{1.912647in}}%
\pgfpathlineto{\pgfqpoint{0.747018in}{1.532264in}}%
\pgfpathlineto{\pgfqpoint{0.746302in}{1.926565in}}%
\pgfpathlineto{\pgfqpoint{0.747216in}{1.532268in}}%
\pgfpathlineto{\pgfqpoint{0.748860in}{1.532689in}}%
\pgfpathlineto{\pgfqpoint{0.748931in}{2.294030in}}%
\pgfpathlineto{\pgfqpoint{0.749941in}{1.914118in}}%
\pgfpathlineto{\pgfqpoint{0.749948in}{1.875643in}}%
\pgfpathlineto{\pgfqpoint{0.749961in}{1.979631in}}%
\pgfpathlineto{\pgfqpoint{0.750708in}{1.937349in}}%
\pgfpathlineto{\pgfqpoint{0.750715in}{1.914267in}}%
\pgfpathlineto{\pgfqpoint{0.750732in}{1.952224in}}%
\pgfpathlineto{\pgfqpoint{0.751477in}{1.942557in}}%
\pgfpathlineto{\pgfqpoint{0.751992in}{1.532252in}}%
\pgfpathlineto{\pgfqpoint{0.751809in}{1.947694in}}%
\pgfpathlineto{\pgfqpoint{0.752536in}{1.532281in}}%
\pgfpathlineto{\pgfqpoint{0.753784in}{1.532297in}}%
\pgfpathlineto{\pgfqpoint{0.753805in}{1.532622in}}%
\pgfpathlineto{\pgfqpoint{0.753874in}{2.285028in}}%
\pgfpathlineto{\pgfqpoint{0.754937in}{1.930539in}}%
\pgfpathlineto{\pgfqpoint{0.754973in}{1.887436in}}%
\pgfpathlineto{\pgfqpoint{0.754957in}{2.012302in}}%
\pgfpathlineto{\pgfqpoint{0.755709in}{1.931351in}}%
\pgfpathlineto{\pgfqpoint{0.755724in}{1.976675in}}%
\pgfpathlineto{\pgfqpoint{0.756491in}{1.962437in}}%
\pgfpathlineto{\pgfqpoint{0.756496in}{1.968071in}}%
\pgfpathlineto{\pgfqpoint{0.756857in}{1.532239in}}%
\pgfpathlineto{\pgfqpoint{0.758750in}{1.532587in}}%
\pgfpathlineto{\pgfqpoint{0.758819in}{2.285018in}}%
\pgfpathlineto{\pgfqpoint{0.759907in}{1.944720in}}%
\pgfpathlineto{\pgfqpoint{0.759911in}{1.915629in}}%
\pgfpathlineto{\pgfqpoint{0.759928in}{2.023441in}}%
\pgfpathlineto{\pgfqpoint{0.760675in}{1.953960in}}%
\pgfpathlineto{\pgfqpoint{0.760688in}{1.993383in}}%
\pgfpathlineto{\pgfqpoint{0.761453in}{1.980759in}}%
\pgfpathlineto{\pgfqpoint{0.761700in}{1.992431in}}%
\pgfpathlineto{\pgfqpoint{0.761728in}{1.545276in}}%
\pgfpathlineto{\pgfqpoint{0.761876in}{1.532220in}}%
\pgfpathlineto{\pgfqpoint{0.762459in}{1.532233in}}%
\pgfpathlineto{\pgfqpoint{0.763693in}{1.532419in}}%
\pgfpathlineto{\pgfqpoint{0.763719in}{1.577383in}}%
\pgfpathlineto{\pgfqpoint{0.763887in}{2.169150in}}%
\pgfpathlineto{\pgfqpoint{0.764529in}{2.061388in}}%
\pgfpathlineto{\pgfqpoint{0.764571in}{1.911395in}}%
\pgfpathlineto{\pgfqpoint{0.765307in}{1.988136in}}%
\pgfpathlineto{\pgfqpoint{0.765315in}{2.014937in}}%
\pgfpathlineto{\pgfqpoint{0.765359in}{1.964017in}}%
\pgfpathlineto{\pgfqpoint{0.766079in}{2.001124in}}%
\pgfpathlineto{\pgfqpoint{0.766772in}{1.532191in}}%
\pgfpathlineto{\pgfqpoint{0.766645in}{2.007386in}}%
\pgfpathlineto{\pgfqpoint{0.767109in}{1.532205in}}%
\pgfpathlineto{\pgfqpoint{0.768627in}{1.532256in}}%
\pgfpathlineto{\pgfqpoint{0.768643in}{1.532548in}}%
\pgfpathlineto{\pgfqpoint{0.768713in}{2.305408in}}%
\pgfpathlineto{\pgfqpoint{0.769794in}{1.998737in}}%
\pgfpathlineto{\pgfqpoint{0.769802in}{1.940363in}}%
\pgfpathlineto{\pgfqpoint{0.769846in}{2.057780in}}%
\pgfpathlineto{\pgfqpoint{0.770564in}{1.985877in}}%
\pgfpathlineto{\pgfqpoint{0.770577in}{2.024656in}}%
\pgfpathlineto{\pgfqpoint{0.770590in}{1.983062in}}%
\pgfpathlineto{\pgfqpoint{0.771337in}{2.008055in}}%
\pgfpathlineto{\pgfqpoint{0.771736in}{1.532160in}}%
\pgfpathlineto{\pgfqpoint{0.771593in}{2.022780in}}%
\pgfpathlineto{\pgfqpoint{0.772468in}{1.532182in}}%
\pgfpathlineto{\pgfqpoint{0.773579in}{1.532267in}}%
\pgfpathlineto{\pgfqpoint{0.773589in}{1.532512in}}%
\pgfpathlineto{\pgfqpoint{0.773660in}{2.291463in}}%
\pgfpathlineto{\pgfqpoint{0.774723in}{1.967674in}}%
\pgfpathlineto{\pgfqpoint{0.774726in}{1.950176in}}%
\pgfpathlineto{\pgfqpoint{0.774739in}{2.070069in}}%
\pgfpathlineto{\pgfqpoint{0.775482in}{2.012155in}}%
\pgfpathlineto{\pgfqpoint{0.775488in}{2.037681in}}%
\pgfpathlineto{\pgfqpoint{0.775539in}{1.993557in}}%
\pgfpathlineto{\pgfqpoint{0.776251in}{2.008252in}}%
\pgfpathlineto{\pgfqpoint{0.776534in}{2.029476in}}%
\pgfpathlineto{\pgfqpoint{0.776541in}{1.977840in}}%
\pgfpathlineto{\pgfqpoint{0.776698in}{1.532127in}}%
\pgfpathlineto{\pgfqpoint{0.777292in}{1.532141in}}%
\pgfpathlineto{\pgfqpoint{0.778538in}{1.532938in}}%
\pgfpathlineto{\pgfqpoint{0.778721in}{2.157316in}}%
\pgfpathlineto{\pgfqpoint{0.779595in}{2.020866in}}%
\pgfpathlineto{\pgfqpoint{0.779603in}{1.980847in}}%
\pgfpathlineto{\pgfqpoint{0.779646in}{2.058801in}}%
\pgfpathlineto{\pgfqpoint{0.780361in}{2.020998in}}%
\pgfpathlineto{\pgfqpoint{0.780369in}{2.036344in}}%
\pgfpathlineto{\pgfqpoint{0.780382in}{2.009935in}}%
\pgfpathlineto{\pgfqpoint{0.781132in}{2.026971in}}%
\pgfpathlineto{\pgfqpoint{0.781481in}{2.031197in}}%
\pgfpathlineto{\pgfqpoint{0.781490in}{1.911904in}}%
\pgfpathlineto{\pgfqpoint{0.781595in}{1.532084in}}%
\pgfpathlineto{\pgfqpoint{0.782270in}{1.532098in}}%
\pgfpathlineto{\pgfqpoint{0.783478in}{1.532260in}}%
\pgfpathlineto{\pgfqpoint{0.783642in}{2.198269in}}%
\pgfpathlineto{\pgfqpoint{0.784806in}{2.019204in}}%
\pgfpathlineto{\pgfqpoint{0.784814in}{1.994973in}}%
\pgfpathlineto{\pgfqpoint{0.784828in}{2.056018in}}%
\pgfpathlineto{\pgfqpoint{0.785576in}{2.015749in}}%
\pgfpathlineto{\pgfqpoint{0.785591in}{2.038106in}}%
\pgfpathlineto{\pgfqpoint{0.786358in}{2.031576in}}%
\pgfpathlineto{\pgfqpoint{0.786532in}{1.532051in}}%
\pgfpathlineto{\pgfqpoint{0.786427in}{2.038257in}}%
\pgfpathlineto{\pgfqpoint{0.787849in}{1.532083in}}%
\pgfpathlineto{\pgfqpoint{0.788423in}{1.532162in}}%
\pgfpathlineto{\pgfqpoint{0.788465in}{1.779166in}}%
\pgfpathlineto{\pgfqpoint{0.788498in}{2.151649in}}%
\pgfpathlineto{\pgfqpoint{0.789247in}{1.984885in}}%
\pgfpathlineto{\pgfqpoint{0.789262in}{2.067185in}}%
\pgfpathlineto{\pgfqpoint{0.790023in}{2.018540in}}%
\pgfpathlineto{\pgfqpoint{0.791125in}{1.546344in}}%
\pgfpathlineto{\pgfqpoint{0.790044in}{2.038538in}}%
\pgfpathlineto{\pgfqpoint{0.791174in}{1.559418in}}%
\pgfpathlineto{\pgfqpoint{0.791373in}{1.908626in}}%
\pgfpathlineto{\pgfqpoint{0.791401in}{1.532027in}}%
\pgfpathlineto{\pgfqpoint{0.791891in}{1.532037in}}%
\pgfpathlineto{\pgfqpoint{0.793368in}{1.532121in}}%
\pgfpathlineto{\pgfqpoint{0.793409in}{1.728270in}}%
\pgfpathlineto{\pgfqpoint{0.793440in}{2.121673in}}%
\pgfpathlineto{\pgfqpoint{0.794198in}{2.012517in}}%
\pgfpathlineto{\pgfqpoint{0.794208in}{2.056862in}}%
\pgfpathlineto{\pgfqpoint{0.794221in}{1.994662in}}%
\pgfpathlineto{\pgfqpoint{0.794968in}{2.028071in}}%
\pgfpathlineto{\pgfqpoint{0.794972in}{2.032885in}}%
\pgfpathlineto{\pgfqpoint{0.795297in}{1.643124in}}%
\pgfpathlineto{\pgfqpoint{0.795445in}{1.546270in}}%
\pgfpathlineto{\pgfqpoint{0.795742in}{1.972951in}}%
\pgfpathlineto{\pgfqpoint{0.795791in}{2.015480in}}%
\pgfpathlineto{\pgfqpoint{0.796038in}{1.617339in}}%
\pgfpathlineto{\pgfqpoint{0.796378in}{1.532031in}}%
\pgfpathlineto{\pgfqpoint{0.796319in}{1.681884in}}%
\pgfpathlineto{\pgfqpoint{0.796852in}{1.532041in}}%
\pgfpathlineto{\pgfqpoint{0.798314in}{1.532122in}}%
\pgfpathlineto{\pgfqpoint{0.798360in}{1.856958in}}%
\pgfpathlineto{\pgfqpoint{0.798380in}{2.102584in}}%
\pgfpathlineto{\pgfqpoint{0.799145in}{2.042445in}}%
\pgfpathlineto{\pgfqpoint{0.799149in}{2.051303in}}%
\pgfpathlineto{\pgfqpoint{0.799163in}{1.998958in}}%
\pgfpathlineto{\pgfqpoint{0.799906in}{2.023523in}}%
\pgfpathlineto{\pgfqpoint{0.800425in}{1.547120in}}%
\pgfpathlineto{\pgfqpoint{0.799928in}{2.029661in}}%
\pgfpathlineto{\pgfqpoint{0.800672in}{1.900362in}}%
\pgfpathlineto{\pgfqpoint{0.800771in}{2.016074in}}%
\pgfpathlineto{\pgfqpoint{0.801117in}{1.549422in}}%
\pgfpathlineto{\pgfqpoint{0.801268in}{1.609370in}}%
\pgfpathlineto{\pgfqpoint{0.801327in}{1.532031in}}%
\pgfpathlineto{\pgfqpoint{0.802005in}{1.532046in}}%
\pgfpathlineto{\pgfqpoint{0.803258in}{1.532114in}}%
\pgfpathlineto{\pgfqpoint{0.803301in}{1.739280in}}%
\pgfpathlineto{\pgfqpoint{0.803334in}{2.124034in}}%
\pgfpathlineto{\pgfqpoint{0.804087in}{1.995244in}}%
\pgfpathlineto{\pgfqpoint{0.804099in}{2.056018in}}%
\pgfpathlineto{\pgfqpoint{0.804115in}{1.992931in}}%
\pgfpathlineto{\pgfqpoint{0.804859in}{2.018501in}}%
\pgfpathlineto{\pgfqpoint{0.804871in}{2.030544in}}%
\pgfpathlineto{\pgfqpoint{0.805303in}{1.558124in}}%
\pgfpathlineto{\pgfqpoint{0.805353in}{1.547622in}}%
\pgfpathlineto{\pgfqpoint{0.805501in}{1.645051in}}%
\pgfpathlineto{\pgfqpoint{0.805699in}{2.007618in}}%
\pgfpathlineto{\pgfqpoint{0.806219in}{1.548894in}}%
\pgfpathlineto{\pgfqpoint{0.806273in}{1.532031in}}%
\pgfpathlineto{\pgfqpoint{0.806954in}{1.532046in}}%
\pgfpathlineto{\pgfqpoint{0.808204in}{1.532112in}}%
\pgfpathlineto{\pgfqpoint{0.808246in}{1.731478in}}%
\pgfpathlineto{\pgfqpoint{0.808322in}{2.118174in}}%
\pgfpathlineto{\pgfqpoint{0.809033in}{1.991090in}}%
\pgfpathlineto{\pgfqpoint{0.809045in}{2.054652in}}%
\pgfpathlineto{\pgfqpoint{0.809808in}{2.012978in}}%
\pgfpathlineto{\pgfqpoint{0.809820in}{2.029009in}}%
\pgfpathlineto{\pgfqpoint{0.810213in}{1.590873in}}%
\pgfpathlineto{\pgfqpoint{0.810312in}{1.549364in}}%
\pgfpathlineto{\pgfqpoint{0.810559in}{1.866427in}}%
\pgfpathlineto{\pgfqpoint{0.810658in}{2.006034in}}%
\pgfpathlineto{\pgfqpoint{0.811054in}{1.553546in}}%
\pgfpathlineto{\pgfqpoint{0.811162in}{1.560647in}}%
\pgfpathlineto{\pgfqpoint{0.811199in}{1.532030in}}%
\pgfpathlineto{\pgfqpoint{0.811905in}{1.532052in}}%
\pgfpathlineto{\pgfqpoint{0.813150in}{1.532123in}}%
\pgfpathlineto{\pgfqpoint{0.813192in}{1.772772in}}%
\pgfpathlineto{\pgfqpoint{0.813226in}{2.120108in}}%
\pgfpathlineto{\pgfqpoint{0.813977in}{2.002862in}}%
\pgfpathlineto{\pgfqpoint{0.814016in}{2.050178in}}%
\pgfpathlineto{\pgfqpoint{0.814000in}{1.990092in}}%
\pgfpathlineto{\pgfqpoint{0.814748in}{2.026814in}}%
\pgfpathlineto{\pgfqpoint{0.815281in}{1.549382in}}%
\pgfpathlineto{\pgfqpoint{0.815528in}{1.912901in}}%
\pgfpathlineto{\pgfqpoint{0.815627in}{2.011848in}}%
\pgfpathlineto{\pgfqpoint{0.815973in}{1.550401in}}%
\pgfpathlineto{\pgfqpoint{0.816101in}{1.624674in}}%
\pgfpathlineto{\pgfqpoint{0.816154in}{1.532028in}}%
\pgfpathlineto{\pgfqpoint{0.816875in}{1.532054in}}%
\pgfpathlineto{\pgfqpoint{0.818096in}{1.532122in}}%
\pgfpathlineto{\pgfqpoint{0.818136in}{1.731187in}}%
\pgfpathlineto{\pgfqpoint{0.818168in}{2.118093in}}%
\pgfpathlineto{\pgfqpoint{0.818922in}{1.986858in}}%
\pgfpathlineto{\pgfqpoint{0.818963in}{2.050146in}}%
\pgfpathlineto{\pgfqpoint{0.819698in}{2.021116in}}%
\pgfpathlineto{\pgfqpoint{0.819702in}{2.026472in}}%
\pgfpathlineto{\pgfqpoint{0.820092in}{1.619401in}}%
\pgfpathlineto{\pgfqpoint{0.820191in}{1.551821in}}%
\pgfpathlineto{\pgfqpoint{0.820488in}{1.918243in}}%
\pgfpathlineto{\pgfqpoint{0.820587in}{2.005701in}}%
\pgfpathlineto{\pgfqpoint{0.820933in}{1.552506in}}%
\pgfpathlineto{\pgfqpoint{0.821012in}{1.576599in}}%
\pgfpathlineto{\pgfqpoint{0.821046in}{1.616282in}}%
\pgfpathlineto{\pgfqpoint{0.821063in}{1.532013in}}%
\pgfpathlineto{\pgfqpoint{0.821705in}{1.532038in}}%
\pgfpathlineto{\pgfqpoint{0.823047in}{1.532281in}}%
\pgfpathlineto{\pgfqpoint{0.823190in}{2.099918in}}%
\pgfpathlineto{\pgfqpoint{0.824184in}{2.010560in}}%
\pgfpathlineto{\pgfqpoint{0.824204in}{2.034715in}}%
\pgfpathlineto{\pgfqpoint{0.825093in}{1.567069in}}%
\pgfpathlineto{\pgfqpoint{0.825143in}{1.550706in}}%
\pgfpathlineto{\pgfqpoint{0.825341in}{1.716330in}}%
\pgfpathlineto{\pgfqpoint{0.825538in}{2.004710in}}%
\pgfpathlineto{\pgfqpoint{0.826008in}{1.532012in}}%
\pgfpathlineto{\pgfqpoint{0.826024in}{1.532023in}}%
\pgfpathlineto{\pgfqpoint{0.826952in}{1.532044in}}%
\pgfpathlineto{\pgfqpoint{0.827991in}{1.532220in}}%
\pgfpathlineto{\pgfqpoint{0.828135in}{2.093087in}}%
\pgfpathlineto{\pgfqpoint{0.829180in}{2.021822in}}%
\pgfpathlineto{\pgfqpoint{0.829204in}{2.030609in}}%
\pgfpathlineto{\pgfqpoint{0.830063in}{1.559312in}}%
\pgfpathlineto{\pgfqpoint{0.830112in}{1.549412in}}%
\pgfpathlineto{\pgfqpoint{0.830261in}{1.651667in}}%
\pgfpathlineto{\pgfqpoint{0.830459in}{2.004705in}}%
\pgfpathlineto{\pgfqpoint{0.830954in}{1.532022in}}%
\pgfpathlineto{\pgfqpoint{0.830967in}{1.532048in}}%
\pgfpathlineto{\pgfqpoint{0.831773in}{1.532059in}}%
\pgfpathlineto{\pgfqpoint{0.832929in}{1.532101in}}%
\pgfpathlineto{\pgfqpoint{0.832937in}{1.532236in}}%
\pgfpathlineto{\pgfqpoint{0.833001in}{2.086234in}}%
\pgfpathlineto{\pgfqpoint{0.834099in}{2.014763in}}%
\pgfpathlineto{\pgfqpoint{0.834122in}{2.028426in}}%
\pgfpathlineto{\pgfqpoint{0.834983in}{1.577080in}}%
\pgfpathlineto{\pgfqpoint{0.835082in}{1.551050in}}%
\pgfpathlineto{\pgfqpoint{0.835280in}{1.793842in}}%
\pgfpathlineto{\pgfqpoint{0.835428in}{2.006919in}}%
\pgfpathlineto{\pgfqpoint{0.835899in}{1.532024in}}%
\pgfpathlineto{\pgfqpoint{0.835916in}{1.532025in}}%
\pgfpathlineto{\pgfqpoint{0.836997in}{1.532064in}}%
\pgfpathlineto{\pgfqpoint{0.837877in}{1.532114in}}%
\pgfpathlineto{\pgfqpoint{0.837919in}{1.728180in}}%
\pgfpathlineto{\pgfqpoint{0.837950in}{2.109607in}}%
\pgfpathlineto{\pgfqpoint{0.838705in}{1.980491in}}%
\pgfpathlineto{\pgfqpoint{0.838721in}{2.043578in}}%
\pgfpathlineto{\pgfqpoint{0.839490in}{2.019384in}}%
\pgfpathlineto{\pgfqpoint{0.839498in}{2.019270in}}%
\pgfpathlineto{\pgfqpoint{0.840033in}{1.553455in}}%
\pgfpathlineto{\pgfqpoint{0.839524in}{2.020224in}}%
\pgfpathlineto{\pgfqpoint{0.840280in}{1.893725in}}%
\pgfpathlineto{\pgfqpoint{0.840379in}{2.000988in}}%
\pgfpathlineto{\pgfqpoint{0.840834in}{1.554036in}}%
\pgfpathlineto{\pgfqpoint{0.840862in}{1.532027in}}%
\pgfpathlineto{\pgfqpoint{0.841598in}{1.532042in}}%
\pgfpathlineto{\pgfqpoint{0.842825in}{1.532148in}}%
\pgfpathlineto{\pgfqpoint{0.842896in}{2.110607in}}%
\pgfpathlineto{\pgfqpoint{0.844047in}{2.023739in}}%
\pgfpathlineto{\pgfqpoint{0.844073in}{2.028453in}}%
\pgfpathlineto{\pgfqpoint{0.844960in}{1.551953in}}%
\pgfpathlineto{\pgfqpoint{0.845108in}{1.654225in}}%
\pgfpathlineto{\pgfqpoint{0.845306in}{1.999158in}}%
\pgfpathlineto{\pgfqpoint{0.845808in}{1.532027in}}%
\pgfpathlineto{\pgfqpoint{0.847037in}{1.532053in}}%
\pgfpathlineto{\pgfqpoint{0.847771in}{1.532152in}}%
\pgfpathlineto{\pgfqpoint{0.847842in}{2.110931in}}%
\pgfpathlineto{\pgfqpoint{0.848990in}{2.026158in}}%
\pgfpathlineto{\pgfqpoint{0.849016in}{2.027322in}}%
\pgfpathlineto{\pgfqpoint{0.849882in}{1.557828in}}%
\pgfpathlineto{\pgfqpoint{0.849931in}{1.553287in}}%
\pgfpathlineto{\pgfqpoint{0.850030in}{1.607885in}}%
\pgfpathlineto{\pgfqpoint{0.850277in}{2.000953in}}%
\pgfpathlineto{\pgfqpoint{0.850744in}{1.532028in}}%
\pgfpathlineto{\pgfqpoint{0.851633in}{1.532046in}}%
\pgfpathlineto{\pgfqpoint{0.852714in}{1.532121in}}%
\pgfpathlineto{\pgfqpoint{0.852754in}{1.723901in}}%
\pgfpathlineto{\pgfqpoint{0.852785in}{2.106056in}}%
\pgfpathlineto{\pgfqpoint{0.853540in}{1.977775in}}%
\pgfpathlineto{\pgfqpoint{0.853555in}{2.040150in}}%
\pgfpathlineto{\pgfqpoint{0.854317in}{2.009595in}}%
\pgfpathlineto{\pgfqpoint{0.854354in}{2.016843in}}%
\pgfpathlineto{\pgfqpoint{0.854774in}{1.593212in}}%
\pgfpathlineto{\pgfqpoint{0.854872in}{1.553822in}}%
\pgfpathlineto{\pgfqpoint{0.855070in}{1.766363in}}%
\pgfpathlineto{\pgfqpoint{0.855219in}{1.997687in}}%
\pgfpathlineto{\pgfqpoint{0.855694in}{1.532029in}}%
\pgfpathlineto{\pgfqpoint{0.855708in}{1.532038in}}%
\pgfpathlineto{\pgfqpoint{0.857650in}{1.532084in}}%
\pgfpathlineto{\pgfqpoint{0.857662in}{1.532153in}}%
\pgfpathlineto{\pgfqpoint{0.857731in}{2.102431in}}%
\pgfpathlineto{\pgfqpoint{0.858886in}{2.004633in}}%
\pgfpathlineto{\pgfqpoint{0.858936in}{2.023150in}}%
\pgfpathlineto{\pgfqpoint{0.859827in}{1.553480in}}%
\pgfpathlineto{\pgfqpoint{0.860025in}{1.778023in}}%
\pgfpathlineto{\pgfqpoint{0.860173in}{1.999841in}}%
\pgfpathlineto{\pgfqpoint{0.860626in}{1.532012in}}%
\pgfpathlineto{\pgfqpoint{0.860664in}{1.532044in}}%
\pgfpathlineto{\pgfqpoint{0.862129in}{1.532075in}}%
\pgfpathlineto{\pgfqpoint{0.862602in}{1.532101in}}%
\pgfpathlineto{\pgfqpoint{0.862611in}{1.532466in}}%
\pgfpathlineto{\pgfqpoint{0.862751in}{2.101461in}}%
\pgfpathlineto{\pgfqpoint{0.863684in}{2.021023in}}%
\pgfpathlineto{\pgfqpoint{0.863696in}{1.987132in}}%
\pgfpathlineto{\pgfqpoint{0.863708in}{2.028885in}}%
\pgfpathlineto{\pgfqpoint{0.864443in}{1.993231in}}%
\pgfpathlineto{\pgfqpoint{0.865571in}{1.532005in}}%
\pgfpathlineto{\pgfqpoint{0.865119in}{1.997385in}}%
\pgfpathlineto{\pgfqpoint{0.865614in}{1.532046in}}%
\pgfpathlineto{\pgfqpoint{0.867034in}{1.532052in}}%
\pgfpathlineto{\pgfqpoint{0.867555in}{1.532235in}}%
\pgfpathlineto{\pgfqpoint{0.867621in}{2.092581in}}%
\pgfpathlineto{\pgfqpoint{0.868688in}{2.020574in}}%
\pgfpathlineto{\pgfqpoint{0.868713in}{2.022972in}}%
\pgfpathlineto{\pgfqpoint{0.869542in}{1.718485in}}%
\pgfpathlineto{\pgfqpoint{0.869739in}{1.557406in}}%
\pgfpathlineto{\pgfqpoint{0.870086in}{1.993775in}}%
\pgfpathlineto{\pgfqpoint{0.870382in}{1.569758in}}%
\pgfpathlineto{\pgfqpoint{0.870517in}{1.532003in}}%
\pgfpathlineto{\pgfqpoint{0.870502in}{1.586075in}}%
\pgfpathlineto{\pgfqpoint{0.871288in}{1.532036in}}%
\pgfpathlineto{\pgfqpoint{0.872499in}{1.532175in}}%
\pgfpathlineto{\pgfqpoint{0.872643in}{2.084998in}}%
\pgfpathlineto{\pgfqpoint{0.873722in}{2.001371in}}%
\pgfpathlineto{\pgfqpoint{0.873741in}{2.018240in}}%
\pgfpathlineto{\pgfqpoint{0.874628in}{1.562484in}}%
\pgfpathlineto{\pgfqpoint{0.874678in}{1.555321in}}%
\pgfpathlineto{\pgfqpoint{0.874776in}{1.603849in}}%
\pgfpathlineto{\pgfqpoint{0.875024in}{1.995178in}}%
\pgfpathlineto{\pgfqpoint{0.875463in}{1.532015in}}%
\pgfpathlineto{\pgfqpoint{0.875484in}{1.532021in}}%
\pgfpathlineto{\pgfqpoint{0.876580in}{1.532066in}}%
\pgfpathlineto{\pgfqpoint{0.877438in}{1.532099in}}%
\pgfpathlineto{\pgfqpoint{0.877446in}{1.532220in}}%
\pgfpathlineto{\pgfqpoint{0.877588in}{2.080893in}}%
\pgfpathlineto{\pgfqpoint{0.878610in}{1.995541in}}%
\pgfpathlineto{\pgfqpoint{0.878629in}{2.018501in}}%
\pgfpathlineto{\pgfqpoint{0.879617in}{1.555382in}}%
\pgfpathlineto{\pgfqpoint{0.879765in}{1.660237in}}%
\pgfpathlineto{\pgfqpoint{0.879963in}{1.992671in}}%
\pgfpathlineto{\pgfqpoint{0.880409in}{1.532019in}}%
\pgfpathlineto{\pgfqpoint{0.880451in}{1.532042in}}%
\pgfpathlineto{\pgfqpoint{0.882219in}{1.532081in}}%
\pgfpathlineto{\pgfqpoint{0.882383in}{1.532097in}}%
\pgfpathlineto{\pgfqpoint{0.882392in}{1.532233in}}%
\pgfpathlineto{\pgfqpoint{0.882456in}{2.075455in}}%
\pgfpathlineto{\pgfqpoint{0.883553in}{2.010123in}}%
\pgfpathlineto{\pgfqpoint{0.883577in}{2.016291in}}%
\pgfpathlineto{\pgfqpoint{0.884416in}{1.680290in}}%
\pgfpathlineto{\pgfqpoint{0.884564in}{1.556489in}}%
\pgfpathlineto{\pgfqpoint{0.884910in}{1.988604in}}%
\pgfpathlineto{\pgfqpoint{0.885256in}{1.562990in}}%
\pgfpathlineto{\pgfqpoint{0.885355in}{1.532026in}}%
\pgfpathlineto{\pgfqpoint{0.885337in}{1.574849in}}%
\pgfpathlineto{\pgfqpoint{0.886288in}{1.532047in}}%
\pgfpathlineto{\pgfqpoint{0.887335in}{1.532155in}}%
\pgfpathlineto{\pgfqpoint{0.887404in}{2.098279in}}%
\pgfpathlineto{\pgfqpoint{0.888534in}{1.986863in}}%
\pgfpathlineto{\pgfqpoint{0.889509in}{1.557324in}}%
\pgfpathlineto{\pgfqpoint{0.888550in}{2.017213in}}%
\pgfpathlineto{\pgfqpoint{0.889707in}{1.736515in}}%
\pgfpathlineto{\pgfqpoint{0.889855in}{1.985639in}}%
\pgfpathlineto{\pgfqpoint{0.890300in}{1.532027in}}%
\pgfpathlineto{\pgfqpoint{0.890351in}{1.532028in}}%
\pgfpathlineto{\pgfqpoint{0.892280in}{1.532147in}}%
\pgfpathlineto{\pgfqpoint{0.892348in}{2.093054in}}%
\pgfpathlineto{\pgfqpoint{0.893504in}{1.995570in}}%
\pgfpathlineto{\pgfqpoint{0.893527in}{2.014342in}}%
\pgfpathlineto{\pgfqpoint{0.894419in}{1.569798in}}%
\pgfpathlineto{\pgfqpoint{0.894469in}{1.559048in}}%
\pgfpathlineto{\pgfqpoint{0.894617in}{1.664319in}}%
\pgfpathlineto{\pgfqpoint{0.894815in}{1.985741in}}%
\pgfpathlineto{\pgfqpoint{0.895260in}{1.532028in}}%
\pgfpathlineto{\pgfqpoint{0.895295in}{1.532028in}}%
\pgfpathlineto{\pgfqpoint{0.897224in}{1.532133in}}%
\pgfpathlineto{\pgfqpoint{0.897293in}{2.081519in}}%
\pgfpathlineto{\pgfqpoint{0.898531in}{2.010195in}}%
\pgfpathlineto{\pgfqpoint{0.899398in}{1.561286in}}%
\pgfpathlineto{\pgfqpoint{0.899497in}{1.583354in}}%
\pgfpathlineto{\pgfqpoint{0.899794in}{1.985315in}}%
\pgfpathlineto{\pgfqpoint{0.900191in}{1.532024in}}%
\pgfpathlineto{\pgfqpoint{0.900207in}{1.532026in}}%
\pgfpathlineto{\pgfqpoint{0.901432in}{1.532054in}}%
\pgfpathlineto{\pgfqpoint{0.902170in}{1.532141in}}%
\pgfpathlineto{\pgfqpoint{0.902239in}{2.088417in}}%
\pgfpathlineto{\pgfqpoint{0.903423in}{1.999301in}}%
\pgfpathlineto{\pgfqpoint{0.903447in}{2.010353in}}%
\pgfpathlineto{\pgfqpoint{0.904328in}{1.566073in}}%
\pgfpathlineto{\pgfqpoint{0.904377in}{1.558693in}}%
\pgfpathlineto{\pgfqpoint{0.904476in}{1.608105in}}%
\pgfpathlineto{\pgfqpoint{0.904724in}{1.988485in}}%
\pgfpathlineto{\pgfqpoint{0.905134in}{1.532019in}}%
\pgfpathlineto{\pgfqpoint{0.905183in}{1.532024in}}%
\pgfpathlineto{\pgfqpoint{0.907113in}{1.532113in}}%
\pgfpathlineto{\pgfqpoint{0.907155in}{1.733122in}}%
\pgfpathlineto{\pgfqpoint{0.907187in}{2.094720in}}%
\pgfpathlineto{\pgfqpoint{0.907939in}{1.966274in}}%
\pgfpathlineto{\pgfqpoint{0.907955in}{2.028333in}}%
\pgfpathlineto{\pgfqpoint{0.908765in}{1.999235in}}%
\pgfpathlineto{\pgfqpoint{0.909341in}{1.561130in}}%
\pgfpathlineto{\pgfqpoint{0.908789in}{2.003090in}}%
\pgfpathlineto{\pgfqpoint{0.909589in}{1.895754in}}%
\pgfpathlineto{\pgfqpoint{0.909688in}{1.985899in}}%
\pgfpathlineto{\pgfqpoint{0.910088in}{1.532034in}}%
\pgfpathlineto{\pgfqpoint{0.910094in}{1.532035in}}%
\pgfpathlineto{\pgfqpoint{0.912051in}{1.532084in}}%
\pgfpathlineto{\pgfqpoint{0.912061in}{1.532133in}}%
\pgfpathlineto{\pgfqpoint{0.912125in}{2.075040in}}%
\pgfpathlineto{\pgfqpoint{0.913398in}{1.994399in}}%
\pgfpathlineto{\pgfqpoint{0.914273in}{1.558915in}}%
\pgfpathlineto{\pgfqpoint{0.913421in}{2.005756in}}%
\pgfpathlineto{\pgfqpoint{0.914323in}{1.569609in}}%
\pgfpathlineto{\pgfqpoint{0.914619in}{1.987301in}}%
\pgfpathlineto{\pgfqpoint{0.915025in}{1.532004in}}%
\pgfpathlineto{\pgfqpoint{0.915034in}{1.532053in}}%
\pgfpathlineto{\pgfqpoint{0.915785in}{1.532037in}}%
\pgfpathlineto{\pgfqpoint{0.917012in}{1.532504in}}%
\pgfpathlineto{\pgfqpoint{0.917149in}{2.087048in}}%
\pgfpathlineto{\pgfqpoint{0.918065in}{1.974085in}}%
\pgfpathlineto{\pgfqpoint{0.918083in}{2.015955in}}%
\pgfpathlineto{\pgfqpoint{0.918096in}{1.973185in}}%
\pgfpathlineto{\pgfqpoint{0.918867in}{1.991793in}}%
\pgfpathlineto{\pgfqpoint{0.919971in}{1.532004in}}%
\pgfpathlineto{\pgfqpoint{0.920727in}{1.532037in}}%
\pgfpathlineto{\pgfqpoint{0.921957in}{1.532450in}}%
\pgfpathlineto{\pgfqpoint{0.922094in}{2.084737in}}%
\pgfpathlineto{\pgfqpoint{0.923033in}{2.000309in}}%
\pgfpathlineto{\pgfqpoint{0.923040in}{1.972793in}}%
\pgfpathlineto{\pgfqpoint{0.923056in}{2.013404in}}%
\pgfpathlineto{\pgfqpoint{0.923801in}{1.993584in}}%
\pgfpathlineto{\pgfqpoint{0.923963in}{1.787496in}}%
\pgfpathlineto{\pgfqpoint{0.924161in}{1.563079in}}%
\pgfpathlineto{\pgfqpoint{0.924507in}{1.975656in}}%
\pgfpathlineto{\pgfqpoint{0.924804in}{1.595632in}}%
\pgfpathlineto{\pgfqpoint{0.924916in}{1.532005in}}%
\pgfpathlineto{\pgfqpoint{0.925669in}{1.532037in}}%
\pgfpathlineto{\pgfqpoint{0.926902in}{1.532406in}}%
\pgfpathlineto{\pgfqpoint{0.927039in}{2.082534in}}%
\pgfpathlineto{\pgfqpoint{0.927976in}{2.010188in}}%
\pgfpathlineto{\pgfqpoint{0.927989in}{1.971911in}}%
\pgfpathlineto{\pgfqpoint{0.928001in}{2.011821in}}%
\pgfpathlineto{\pgfqpoint{0.928751in}{1.992220in}}%
\pgfpathlineto{\pgfqpoint{0.928950in}{1.717890in}}%
\pgfpathlineto{\pgfqpoint{0.929148in}{1.565183in}}%
\pgfpathlineto{\pgfqpoint{0.929494in}{1.978807in}}%
\pgfpathlineto{\pgfqpoint{0.929791in}{1.575015in}}%
\pgfpathlineto{\pgfqpoint{0.929862in}{1.532016in}}%
\pgfpathlineto{\pgfqpoint{0.929848in}{1.575672in}}%
\pgfpathlineto{\pgfqpoint{0.930673in}{1.532056in}}%
\pgfpathlineto{\pgfqpoint{0.931838in}{1.532100in}}%
\pgfpathlineto{\pgfqpoint{0.931846in}{1.532251in}}%
\pgfpathlineto{\pgfqpoint{0.931909in}{2.064242in}}%
\pgfpathlineto{\pgfqpoint{0.932979in}{1.999209in}}%
\pgfpathlineto{\pgfqpoint{0.933018in}{1.976516in}}%
\pgfpathlineto{\pgfqpoint{0.933003in}{2.005820in}}%
\pgfpathlineto{\pgfqpoint{0.933745in}{1.978254in}}%
\pgfpathlineto{\pgfqpoint{0.934808in}{1.532017in}}%
\pgfpathlineto{\pgfqpoint{0.934437in}{1.982154in}}%
\pgfpathlineto{\pgfqpoint{0.934974in}{1.532025in}}%
\pgfpathlineto{\pgfqpoint{0.936793in}{1.532448in}}%
\pgfpathlineto{\pgfqpoint{0.936933in}{2.082347in}}%
\pgfpathlineto{\pgfqpoint{0.937871in}{1.993820in}}%
\pgfpathlineto{\pgfqpoint{0.937879in}{1.968615in}}%
\pgfpathlineto{\pgfqpoint{0.937895in}{2.009195in}}%
\pgfpathlineto{\pgfqpoint{0.938638in}{1.991722in}}%
\pgfpathlineto{\pgfqpoint{0.938750in}{1.909511in}}%
\pgfpathlineto{\pgfqpoint{0.939047in}{1.565914in}}%
\pgfpathlineto{\pgfqpoint{0.939393in}{1.977444in}}%
\pgfpathlineto{\pgfqpoint{0.939686in}{1.577149in}}%
\pgfpathlineto{\pgfqpoint{0.939753in}{1.532014in}}%
\pgfpathlineto{\pgfqpoint{0.940614in}{1.532041in}}%
\pgfpathlineto{\pgfqpoint{0.941737in}{1.532205in}}%
\pgfpathlineto{\pgfqpoint{0.941798in}{2.065494in}}%
\pgfpathlineto{\pgfqpoint{0.942901in}{1.984931in}}%
\pgfpathlineto{\pgfqpoint{0.942922in}{2.004265in}}%
\pgfpathlineto{\pgfqpoint{0.943700in}{1.910718in}}%
\pgfpathlineto{\pgfqpoint{0.943997in}{1.564654in}}%
\pgfpathlineto{\pgfqpoint{0.944343in}{1.979805in}}%
\pgfpathlineto{\pgfqpoint{0.944634in}{1.576905in}}%
\pgfpathlineto{\pgfqpoint{0.944717in}{1.532025in}}%
\pgfpathlineto{\pgfqpoint{0.945556in}{1.532043in}}%
\pgfpathlineto{\pgfqpoint{0.946679in}{1.532137in}}%
\pgfpathlineto{\pgfqpoint{0.946746in}{2.076314in}}%
\pgfpathlineto{\pgfqpoint{0.947933in}{1.987066in}}%
\pgfpathlineto{\pgfqpoint{0.947955in}{2.001792in}}%
\pgfpathlineto{\pgfqpoint{0.948851in}{1.584567in}}%
\pgfpathlineto{\pgfqpoint{0.948900in}{1.564914in}}%
\pgfpathlineto{\pgfqpoint{0.949098in}{1.721848in}}%
\pgfpathlineto{\pgfqpoint{0.949296in}{1.976928in}}%
\pgfpathlineto{\pgfqpoint{0.949645in}{1.532023in}}%
\pgfpathlineto{\pgfqpoint{0.949760in}{1.532039in}}%
\pgfpathlineto{\pgfqpoint{0.951623in}{1.532120in}}%
\pgfpathlineto{\pgfqpoint{0.951689in}{2.067322in}}%
\pgfpathlineto{\pgfqpoint{0.953029in}{1.977131in}}%
\pgfpathlineto{\pgfqpoint{0.953046in}{1.996368in}}%
\pgfpathlineto{\pgfqpoint{0.953562in}{1.961034in}}%
\pgfpathlineto{\pgfqpoint{0.954233in}{1.979679in}}%
\pgfpathlineto{\pgfqpoint{0.954622in}{1.532027in}}%
\pgfpathlineto{\pgfqpoint{0.956569in}{1.532117in}}%
\pgfpathlineto{\pgfqpoint{0.956610in}{1.752810in}}%
\pgfpathlineto{\pgfqpoint{0.956642in}{2.062982in}}%
\pgfpathlineto{\pgfqpoint{0.957391in}{1.962796in}}%
\pgfpathlineto{\pgfqpoint{0.957409in}{2.010379in}}%
\pgfpathlineto{\pgfqpoint{0.957422in}{1.962189in}}%
\pgfpathlineto{\pgfqpoint{0.958231in}{1.987169in}}%
\pgfpathlineto{\pgfqpoint{0.958824in}{1.566090in}}%
\pgfpathlineto{\pgfqpoint{0.958255in}{1.990633in}}%
\pgfpathlineto{\pgfqpoint{0.959071in}{1.861161in}}%
\pgfpathlineto{\pgfqpoint{0.959170in}{1.973812in}}%
\pgfpathlineto{\pgfqpoint{0.959567in}{1.532029in}}%
\pgfpathlineto{\pgfqpoint{0.959596in}{1.532033in}}%
\pgfpathlineto{\pgfqpoint{0.961513in}{1.532108in}}%
\pgfpathlineto{\pgfqpoint{0.961555in}{1.732619in}}%
\pgfpathlineto{\pgfqpoint{0.961588in}{2.064720in}}%
\pgfpathlineto{\pgfqpoint{0.962337in}{1.980858in}}%
\pgfpathlineto{\pgfqpoint{0.962374in}{1.959723in}}%
\pgfpathlineto{\pgfqpoint{0.962361in}{2.008664in}}%
\pgfpathlineto{\pgfqpoint{0.963107in}{1.978709in}}%
\pgfpathlineto{\pgfqpoint{0.963116in}{1.978815in}}%
\pgfpathlineto{\pgfqpoint{0.963161in}{1.992055in}}%
\pgfpathlineto{\pgfqpoint{0.963638in}{1.663235in}}%
\pgfpathlineto{\pgfqpoint{0.963786in}{1.571158in}}%
\pgfpathlineto{\pgfqpoint{0.964132in}{1.966786in}}%
\pgfpathlineto{\pgfqpoint{0.964481in}{1.532009in}}%
\pgfpathlineto{\pgfqpoint{0.965730in}{1.532048in}}%
\pgfpathlineto{\pgfqpoint{0.966465in}{1.532329in}}%
\pgfpathlineto{\pgfqpoint{0.966602in}{2.070133in}}%
\pgfpathlineto{\pgfqpoint{0.967544in}{1.987573in}}%
\pgfpathlineto{\pgfqpoint{0.967552in}{1.963782in}}%
\pgfpathlineto{\pgfqpoint{0.967597in}{2.002511in}}%
\pgfpathlineto{\pgfqpoint{0.968313in}{1.983930in}}%
\pgfpathlineto{\pgfqpoint{0.969429in}{1.532023in}}%
\pgfpathlineto{\pgfqpoint{0.968348in}{1.984821in}}%
\pgfpathlineto{\pgfqpoint{0.970011in}{1.532047in}}%
\pgfpathlineto{\pgfqpoint{0.971407in}{1.532150in}}%
\pgfpathlineto{\pgfqpoint{0.971476in}{2.082716in}}%
\pgfpathlineto{\pgfqpoint{0.972605in}{1.980305in}}%
\pgfpathlineto{\pgfqpoint{0.972655in}{1.999188in}}%
\pgfpathlineto{\pgfqpoint{0.973388in}{1.922746in}}%
\pgfpathlineto{\pgfqpoint{0.973685in}{1.567424in}}%
\pgfpathlineto{\pgfqpoint{0.974031in}{1.973673in}}%
\pgfpathlineto{\pgfqpoint{0.974354in}{1.578563in}}%
\pgfpathlineto{\pgfqpoint{0.974358in}{1.581964in}}%
\pgfpathlineto{\pgfqpoint{0.974371in}{1.532019in}}%
\pgfpathlineto{\pgfqpoint{0.975059in}{1.532052in}}%
\pgfpathlineto{\pgfqpoint{0.976350in}{1.532112in}}%
\pgfpathlineto{\pgfqpoint{0.976391in}{1.739102in}}%
\pgfpathlineto{\pgfqpoint{0.976424in}{2.041521in}}%
\pgfpathlineto{\pgfqpoint{0.977177in}{1.971197in}}%
\pgfpathlineto{\pgfqpoint{0.977182in}{1.961179in}}%
\pgfpathlineto{\pgfqpoint{0.977197in}{2.002291in}}%
\pgfpathlineto{\pgfqpoint{0.977940in}{1.986718in}}%
\pgfpathlineto{\pgfqpoint{0.978647in}{1.573574in}}%
\pgfpathlineto{\pgfqpoint{0.978053in}{1.988091in}}%
\pgfpathlineto{\pgfqpoint{0.978845in}{1.798519in}}%
\pgfpathlineto{\pgfqpoint{0.978993in}{1.963473in}}%
\pgfpathlineto{\pgfqpoint{0.979371in}{1.532029in}}%
\pgfpathlineto{\pgfqpoint{0.979417in}{1.532033in}}%
\pgfpathlineto{\pgfqpoint{0.981294in}{1.532105in}}%
\pgfpathlineto{\pgfqpoint{0.981336in}{1.717090in}}%
\pgfpathlineto{\pgfqpoint{0.981409in}{2.073963in}}%
\pgfpathlineto{\pgfqpoint{0.982121in}{1.958614in}}%
\pgfpathlineto{\pgfqpoint{0.982125in}{1.949604in}}%
\pgfpathlineto{\pgfqpoint{0.982141in}{2.008982in}}%
\pgfpathlineto{\pgfqpoint{0.982881in}{1.984052in}}%
\pgfpathlineto{\pgfqpoint{0.982885in}{1.989393in}}%
\pgfpathlineto{\pgfqpoint{0.983401in}{1.716635in}}%
\pgfpathlineto{\pgfqpoint{0.984279in}{1.532031in}}%
\pgfpathlineto{\pgfqpoint{0.983945in}{1.964834in}}%
\pgfpathlineto{\pgfqpoint{0.984290in}{1.532032in}}%
\pgfpathlineto{\pgfqpoint{0.986240in}{1.532105in}}%
\pgfpathlineto{\pgfqpoint{0.986281in}{1.717033in}}%
\pgfpathlineto{\pgfqpoint{0.986354in}{2.073191in}}%
\pgfpathlineto{\pgfqpoint{0.987066in}{1.958635in}}%
\pgfpathlineto{\pgfqpoint{0.987070in}{1.948435in}}%
\pgfpathlineto{\pgfqpoint{0.987086in}{2.008860in}}%
\pgfpathlineto{\pgfqpoint{0.987826in}{1.981855in}}%
\pgfpathlineto{\pgfqpoint{0.987863in}{1.988290in}}%
\pgfpathlineto{\pgfqpoint{0.988365in}{1.699282in}}%
\pgfpathlineto{\pgfqpoint{0.988859in}{1.964928in}}%
\pgfpathlineto{\pgfqpoint{0.989244in}{1.532027in}}%
\pgfpathlineto{\pgfqpoint{0.991187in}{1.532122in}}%
\pgfpathlineto{\pgfqpoint{0.991252in}{2.059870in}}%
\pgfpathlineto{\pgfqpoint{0.992584in}{1.987288in}}%
\pgfpathlineto{\pgfqpoint{0.993458in}{1.568308in}}%
\pgfpathlineto{\pgfqpoint{0.992611in}{1.988604in}}%
\pgfpathlineto{\pgfqpoint{0.993557in}{1.602626in}}%
\pgfpathlineto{\pgfqpoint{0.993804in}{1.966983in}}%
\pgfpathlineto{\pgfqpoint{0.994153in}{1.532020in}}%
\pgfpathlineto{\pgfqpoint{0.994267in}{1.532026in}}%
\pgfpathlineto{\pgfqpoint{0.996132in}{1.532115in}}%
\pgfpathlineto{\pgfqpoint{0.996172in}{1.726684in}}%
\pgfpathlineto{\pgfqpoint{0.996204in}{2.074633in}}%
\pgfpathlineto{\pgfqpoint{0.996956in}{1.951096in}}%
\pgfpathlineto{\pgfqpoint{0.996960in}{1.947814in}}%
\pgfpathlineto{\pgfqpoint{0.996974in}{2.009133in}}%
\pgfpathlineto{\pgfqpoint{0.997682in}{1.983148in}}%
\pgfpathlineto{\pgfqpoint{0.997687in}{1.987366in}}%
\pgfpathlineto{\pgfqpoint{0.998180in}{1.857619in}}%
\pgfpathlineto{\pgfqpoint{0.998774in}{1.966362in}}%
\pgfpathlineto{\pgfqpoint{0.999117in}{1.532030in}}%
\pgfpathlineto{\pgfqpoint{1.001078in}{1.532121in}}%
\pgfpathlineto{\pgfqpoint{1.001143in}{2.058403in}}%
\pgfpathlineto{\pgfqpoint{1.002455in}{1.967934in}}%
\pgfpathlineto{\pgfqpoint{1.003345in}{1.571575in}}%
\pgfpathlineto{\pgfqpoint{1.002471in}{1.987089in}}%
\pgfpathlineto{\pgfqpoint{1.003592in}{1.813840in}}%
\pgfpathlineto{\pgfqpoint{1.003741in}{1.966804in}}%
\pgfpathlineto{\pgfqpoint{1.004097in}{1.532029in}}%
\pgfpathlineto{\pgfqpoint{1.004191in}{1.532031in}}%
\pgfpathlineto{\pgfqpoint{1.006030in}{1.532430in}}%
\pgfpathlineto{\pgfqpoint{1.006136in}{2.067025in}}%
\pgfpathlineto{\pgfqpoint{1.007082in}{1.960981in}}%
\pgfpathlineto{\pgfqpoint{1.007086in}{1.954706in}}%
\pgfpathlineto{\pgfqpoint{1.007099in}{1.995518in}}%
\pgfpathlineto{\pgfqpoint{1.007841in}{1.970577in}}%
\pgfpathlineto{\pgfqpoint{1.007854in}{1.980570in}}%
\pgfpathlineto{\pgfqpoint{1.008198in}{1.648197in}}%
\pgfpathlineto{\pgfqpoint{1.008346in}{1.572645in}}%
\pgfpathlineto{\pgfqpoint{1.008643in}{1.952180in}}%
\pgfpathlineto{\pgfqpoint{1.008692in}{1.965347in}}%
\pgfpathlineto{\pgfqpoint{1.008791in}{1.854730in}}%
\pgfpathlineto{\pgfqpoint{1.009024in}{1.532027in}}%
\pgfpathlineto{\pgfqpoint{1.009692in}{1.532042in}}%
\pgfpathlineto{\pgfqpoint{1.010969in}{1.532121in}}%
\pgfpathlineto{\pgfqpoint{1.011034in}{2.056615in}}%
\pgfpathlineto{\pgfqpoint{1.012368in}{1.981634in}}%
\pgfpathlineto{\pgfqpoint{1.013251in}{1.571234in}}%
\pgfpathlineto{\pgfqpoint{1.012422in}{1.984550in}}%
\pgfpathlineto{\pgfqpoint{1.013300in}{1.573759in}}%
\pgfpathlineto{\pgfqpoint{1.013647in}{1.963781in}}%
\pgfpathlineto{\pgfqpoint{1.013967in}{1.532028in}}%
\pgfpathlineto{\pgfqpoint{1.015914in}{1.532114in}}%
\pgfpathlineto{\pgfqpoint{1.015978in}{2.056791in}}%
\pgfpathlineto{\pgfqpoint{1.017368in}{1.982151in}}%
\pgfpathlineto{\pgfqpoint{1.018219in}{1.570346in}}%
\pgfpathlineto{\pgfqpoint{1.017395in}{1.983109in}}%
\pgfpathlineto{\pgfqpoint{1.018318in}{1.612131in}}%
\pgfpathlineto{\pgfqpoint{1.018565in}{1.965650in}}%
\pgfpathlineto{\pgfqpoint{1.018925in}{1.532026in}}%
\pgfpathlineto{\pgfqpoint{1.019010in}{1.532027in}}%
\pgfpathlineto{\pgfqpoint{1.020861in}{1.532134in}}%
\pgfpathlineto{\pgfqpoint{1.020928in}{2.066238in}}%
\pgfpathlineto{\pgfqpoint{1.022089in}{1.970670in}}%
\pgfpathlineto{\pgfqpoint{1.022112in}{1.987529in}}%
\pgfpathlineto{\pgfqpoint{1.022865in}{1.943575in}}%
\pgfpathlineto{\pgfqpoint{1.023876in}{1.532025in}}%
\pgfpathlineto{\pgfqpoint{1.023545in}{1.961833in}}%
\pgfpathlineto{\pgfqpoint{1.023918in}{1.532035in}}%
\pgfpathlineto{\pgfqpoint{1.025807in}{1.532139in}}%
\pgfpathlineto{\pgfqpoint{1.025874in}{2.066295in}}%
\pgfpathlineto{\pgfqpoint{1.027031in}{1.984197in}}%
\pgfpathlineto{\pgfqpoint{1.027058in}{1.986840in}}%
\pgfpathlineto{\pgfqpoint{1.027808in}{1.950934in}}%
\pgfpathlineto{\pgfqpoint{1.028793in}{1.532025in}}%
\pgfpathlineto{\pgfqpoint{1.028481in}{1.963442in}}%
\pgfpathlineto{\pgfqpoint{1.028938in}{1.532027in}}%
\pgfpathlineto{\pgfqpoint{1.030752in}{1.532139in}}%
\pgfpathlineto{\pgfqpoint{1.030820in}{2.069479in}}%
\pgfpathlineto{\pgfqpoint{1.031952in}{1.967845in}}%
\pgfpathlineto{\pgfqpoint{1.031986in}{1.956082in}}%
\pgfpathlineto{\pgfqpoint{1.031974in}{1.988354in}}%
\pgfpathlineto{\pgfqpoint{1.032722in}{1.968871in}}%
\pgfpathlineto{\pgfqpoint{1.032806in}{1.888174in}}%
\pgfpathlineto{\pgfqpoint{1.033739in}{1.532024in}}%
\pgfpathlineto{\pgfqpoint{1.033449in}{1.959112in}}%
\pgfpathlineto{\pgfqpoint{1.033757in}{1.532036in}}%
\pgfpathlineto{\pgfqpoint{1.035686in}{1.532080in}}%
\pgfpathlineto{\pgfqpoint{1.035698in}{1.532141in}}%
\pgfpathlineto{\pgfqpoint{1.035768in}{2.045064in}}%
\pgfpathlineto{\pgfqpoint{1.036879in}{1.964128in}}%
\pgfpathlineto{\pgfqpoint{1.036884in}{1.956427in}}%
\pgfpathlineto{\pgfqpoint{1.036899in}{1.985291in}}%
\pgfpathlineto{\pgfqpoint{1.037631in}{1.973701in}}%
\pgfpathlineto{\pgfqpoint{1.037842in}{1.736189in}}%
\pgfpathlineto{\pgfqpoint{1.038664in}{1.532000in}}%
\pgfpathlineto{\pgfqpoint{1.038386in}{1.955189in}}%
\pgfpathlineto{\pgfqpoint{1.038715in}{1.532022in}}%
\pgfpathlineto{\pgfqpoint{1.040648in}{1.532491in}}%
\pgfpathlineto{\pgfqpoint{1.040753in}{2.060676in}}%
\pgfpathlineto{\pgfqpoint{1.041696in}{1.975796in}}%
\pgfpathlineto{\pgfqpoint{1.041705in}{1.949002in}}%
\pgfpathlineto{\pgfqpoint{1.041721in}{1.989599in}}%
\pgfpathlineto{\pgfqpoint{1.042468in}{1.971616in}}%
\pgfpathlineto{\pgfqpoint{1.042474in}{1.974796in}}%
\pgfpathlineto{\pgfqpoint{1.042756in}{1.805119in}}%
\pgfpathlineto{\pgfqpoint{1.043610in}{1.532001in}}%
\pgfpathlineto{\pgfqpoint{1.043350in}{1.956855in}}%
\pgfpathlineto{\pgfqpoint{1.043655in}{1.532022in}}%
\pgfpathlineto{\pgfqpoint{1.045594in}{1.533420in}}%
\pgfpathlineto{\pgfqpoint{1.045699in}{2.058694in}}%
\pgfpathlineto{\pgfqpoint{1.046620in}{1.949504in}}%
\pgfpathlineto{\pgfqpoint{1.046638in}{1.990530in}}%
\pgfpathlineto{\pgfqpoint{1.046651in}{1.948855in}}%
\pgfpathlineto{\pgfqpoint{1.047569in}{1.966461in}}%
\pgfpathlineto{\pgfqpoint{1.048555in}{1.532005in}}%
\pgfpathlineto{\pgfqpoint{1.049423in}{1.532054in}}%
\pgfpathlineto{\pgfqpoint{1.050532in}{1.532110in}}%
\pgfpathlineto{\pgfqpoint{1.050574in}{1.749281in}}%
\pgfpathlineto{\pgfqpoint{1.050603in}{2.037420in}}%
\pgfpathlineto{\pgfqpoint{1.051357in}{1.948389in}}%
\pgfpathlineto{\pgfqpoint{1.051362in}{1.948204in}}%
\pgfpathlineto{\pgfqpoint{1.051367in}{1.965044in}}%
\pgfpathlineto{\pgfqpoint{1.051376in}{1.990844in}}%
\pgfpathlineto{\pgfqpoint{1.051389in}{1.946940in}}%
\pgfpathlineto{\pgfqpoint{1.052133in}{1.961080in}}%
\pgfpathlineto{\pgfqpoint{1.052147in}{1.975390in}}%
\pgfpathlineto{\pgfqpoint{1.052670in}{1.781943in}}%
\pgfpathlineto{\pgfqpoint{1.053501in}{1.532017in}}%
\pgfpathlineto{\pgfqpoint{1.053214in}{1.955137in}}%
\pgfpathlineto{\pgfqpoint{1.053564in}{1.532035in}}%
\pgfpathlineto{\pgfqpoint{1.055478in}{1.532113in}}%
\pgfpathlineto{\pgfqpoint{1.055517in}{1.720626in}}%
\pgfpathlineto{\pgfqpoint{1.055549in}{2.065241in}}%
\pgfpathlineto{\pgfqpoint{1.056303in}{1.941141in}}%
\pgfpathlineto{\pgfqpoint{1.056307in}{1.936845in}}%
\pgfpathlineto{\pgfqpoint{1.056320in}{1.998876in}}%
\pgfpathlineto{\pgfqpoint{1.057051in}{1.958014in}}%
\pgfpathlineto{\pgfqpoint{1.057066in}{1.976739in}}%
\pgfpathlineto{\pgfqpoint{1.057626in}{1.764190in}}%
\pgfpathlineto{\pgfqpoint{1.058446in}{1.532017in}}%
\pgfpathlineto{\pgfqpoint{1.058170in}{1.958490in}}%
\pgfpathlineto{\pgfqpoint{1.058510in}{1.532035in}}%
\pgfpathlineto{\pgfqpoint{1.060423in}{1.532113in}}%
\pgfpathlineto{\pgfqpoint{1.060463in}{1.720613in}}%
\pgfpathlineto{\pgfqpoint{1.060494in}{2.063708in}}%
\pgfpathlineto{\pgfqpoint{1.061248in}{1.940295in}}%
\pgfpathlineto{\pgfqpoint{1.061252in}{1.936407in}}%
\pgfpathlineto{\pgfqpoint{1.061265in}{1.997912in}}%
\pgfpathlineto{\pgfqpoint{1.061996in}{1.957335in}}%
\pgfpathlineto{\pgfqpoint{1.062012in}{1.975809in}}%
\pgfpathlineto{\pgfqpoint{1.062565in}{1.787252in}}%
\pgfpathlineto{\pgfqpoint{1.063392in}{1.532017in}}%
\pgfpathlineto{\pgfqpoint{1.063109in}{1.953431in}}%
\pgfpathlineto{\pgfqpoint{1.063453in}{1.532035in}}%
\pgfpathlineto{\pgfqpoint{1.065369in}{1.532113in}}%
\pgfpathlineto{\pgfqpoint{1.065408in}{1.721094in}}%
\pgfpathlineto{\pgfqpoint{1.065439in}{2.061358in}}%
\pgfpathlineto{\pgfqpoint{1.066194in}{1.940242in}}%
\pgfpathlineto{\pgfqpoint{1.066197in}{1.936174in}}%
\pgfpathlineto{\pgfqpoint{1.066211in}{1.996397in}}%
\pgfpathlineto{\pgfqpoint{1.066942in}{1.956717in}}%
\pgfpathlineto{\pgfqpoint{1.066957in}{1.974861in}}%
\pgfpathlineto{\pgfqpoint{1.067510in}{1.796117in}}%
\pgfpathlineto{\pgfqpoint{1.068338in}{1.532007in}}%
\pgfpathlineto{\pgfqpoint{1.068054in}{1.949031in}}%
\pgfpathlineto{\pgfqpoint{1.068409in}{1.532036in}}%
\pgfpathlineto{\pgfqpoint{1.070304in}{1.532078in}}%
\pgfpathlineto{\pgfqpoint{1.070315in}{1.532119in}}%
\pgfpathlineto{\pgfqpoint{1.070378in}{2.047927in}}%
\pgfpathlineto{\pgfqpoint{1.071692in}{1.957283in}}%
\pgfpathlineto{\pgfqpoint{1.072664in}{1.575997in}}%
\pgfpathlineto{\pgfqpoint{1.071711in}{1.975673in}}%
\pgfpathlineto{\pgfqpoint{1.072812in}{1.670946in}}%
\pgfpathlineto{\pgfqpoint{1.073010in}{1.955451in}}%
\pgfpathlineto{\pgfqpoint{1.073283in}{1.532004in}}%
\pgfpathlineto{\pgfqpoint{1.073450in}{1.532038in}}%
\pgfpathlineto{\pgfqpoint{1.075259in}{1.532111in}}%
\pgfpathlineto{\pgfqpoint{1.075301in}{1.759525in}}%
\pgfpathlineto{\pgfqpoint{1.075364in}{2.027561in}}%
\pgfpathlineto{\pgfqpoint{1.076084in}{1.945394in}}%
\pgfpathlineto{\pgfqpoint{1.076103in}{1.984616in}}%
\pgfpathlineto{\pgfqpoint{1.076116in}{1.944082in}}%
\pgfpathlineto{\pgfqpoint{1.076961in}{1.962824in}}%
\pgfpathlineto{\pgfqpoint{1.077613in}{1.579699in}}%
\pgfpathlineto{\pgfqpoint{1.076987in}{1.969973in}}%
\pgfpathlineto{\pgfqpoint{1.077860in}{1.835433in}}%
\pgfpathlineto{\pgfqpoint{1.077959in}{1.949259in}}%
\pgfpathlineto{\pgfqpoint{1.078229in}{1.531999in}}%
\pgfpathlineto{\pgfqpoint{1.078389in}{1.532038in}}%
\pgfpathlineto{\pgfqpoint{1.080204in}{1.532106in}}%
\pgfpathlineto{\pgfqpoint{1.080246in}{1.731931in}}%
\pgfpathlineto{\pgfqpoint{1.080277in}{2.024292in}}%
\pgfpathlineto{\pgfqpoint{1.081030in}{1.973928in}}%
\pgfpathlineto{\pgfqpoint{1.081040in}{1.943009in}}%
\pgfpathlineto{\pgfqpoint{1.081055in}{1.984316in}}%
\pgfpathlineto{\pgfqpoint{1.081794in}{1.965606in}}%
\pgfpathlineto{\pgfqpoint{1.081804in}{1.970923in}}%
\pgfpathlineto{\pgfqpoint{1.082336in}{1.852999in}}%
\pgfpathlineto{\pgfqpoint{1.083175in}{1.532000in}}%
\pgfpathlineto{\pgfqpoint{1.082929in}{1.960732in}}%
\pgfpathlineto{\pgfqpoint{1.083275in}{1.532024in}}%
\pgfpathlineto{\pgfqpoint{1.085151in}{1.532122in}}%
\pgfpathlineto{\pgfqpoint{1.085215in}{2.045374in}}%
\pgfpathlineto{\pgfqpoint{1.086523in}{1.966107in}}%
\pgfpathlineto{\pgfqpoint{1.086546in}{1.973538in}}%
\pgfpathlineto{\pgfqpoint{1.087436in}{1.605496in}}%
\pgfpathlineto{\pgfqpoint{1.087535in}{1.578393in}}%
\pgfpathlineto{\pgfqpoint{1.087733in}{1.801129in}}%
\pgfpathlineto{\pgfqpoint{1.087881in}{1.955545in}}%
\pgfpathlineto{\pgfqpoint{1.088120in}{1.532000in}}%
\pgfpathlineto{\pgfqpoint{1.088334in}{1.532026in}}%
\pgfpathlineto{\pgfqpoint{1.090095in}{1.532104in}}%
\pgfpathlineto{\pgfqpoint{1.090136in}{1.722278in}}%
\pgfpathlineto{\pgfqpoint{1.090168in}{2.040987in}}%
\pgfpathlineto{\pgfqpoint{1.090921in}{1.962633in}}%
\pgfpathlineto{\pgfqpoint{1.090929in}{1.936426in}}%
\pgfpathlineto{\pgfqpoint{1.090942in}{1.986533in}}%
\pgfpathlineto{\pgfqpoint{1.091687in}{1.966858in}}%
\pgfpathlineto{\pgfqpoint{1.092476in}{1.579454in}}%
\pgfpathlineto{\pgfqpoint{1.091747in}{1.969628in}}%
\pgfpathlineto{\pgfqpoint{1.092624in}{1.679483in}}%
\pgfpathlineto{\pgfqpoint{1.092822in}{1.950962in}}%
\pgfpathlineto{\pgfqpoint{1.093102in}{1.532029in}}%
\pgfpathlineto{\pgfqpoint{1.093275in}{1.532032in}}%
\pgfpathlineto{\pgfqpoint{1.095042in}{1.532118in}}%
\pgfpathlineto{\pgfqpoint{1.095106in}{2.044691in}}%
\pgfpathlineto{\pgfqpoint{1.096411in}{1.969253in}}%
\pgfpathlineto{\pgfqpoint{1.096438in}{1.971838in}}%
\pgfpathlineto{\pgfqpoint{1.097320in}{1.624762in}}%
\pgfpathlineto{\pgfqpoint{1.097419in}{1.578248in}}%
\pgfpathlineto{\pgfqpoint{1.097666in}{1.851340in}}%
\pgfpathlineto{\pgfqpoint{1.097765in}{1.953329in}}%
\pgfpathlineto{\pgfqpoint{1.098012in}{1.532020in}}%
\pgfpathlineto{\pgfqpoint{1.098208in}{1.532029in}}%
\pgfpathlineto{\pgfqpoint{1.099987in}{1.532114in}}%
\pgfpathlineto{\pgfqpoint{1.100051in}{2.044420in}}%
\pgfpathlineto{\pgfqpoint{1.101392in}{1.955709in}}%
\pgfpathlineto{\pgfqpoint{1.102374in}{1.578437in}}%
\pgfpathlineto{\pgfqpoint{1.101411in}{1.970833in}}%
\pgfpathlineto{\pgfqpoint{1.102423in}{1.589194in}}%
\pgfpathlineto{\pgfqpoint{1.102720in}{1.953322in}}%
\pgfpathlineto{\pgfqpoint{1.103002in}{1.532028in}}%
\pgfpathlineto{\pgfqpoint{1.103087in}{1.532030in}}%
\pgfpathlineto{\pgfqpoint{1.104932in}{1.532114in}}%
\pgfpathlineto{\pgfqpoint{1.104974in}{1.757920in}}%
\pgfpathlineto{\pgfqpoint{1.105036in}{2.023404in}}%
\pgfpathlineto{\pgfqpoint{1.105756in}{1.942257in}}%
\pgfpathlineto{\pgfqpoint{1.105760in}{1.939517in}}%
\pgfpathlineto{\pgfqpoint{1.105772in}{1.981103in}}%
\pgfpathlineto{\pgfqpoint{1.106501in}{1.955936in}}%
\pgfpathlineto{\pgfqpoint{1.106547in}{1.966689in}}%
\pgfpathlineto{\pgfqpoint{1.107097in}{1.825634in}}%
\pgfpathlineto{\pgfqpoint{1.107944in}{1.532026in}}%
\pgfpathlineto{\pgfqpoint{1.107691in}{1.948246in}}%
\pgfpathlineto{\pgfqpoint{1.108025in}{1.532027in}}%
\pgfpathlineto{\pgfqpoint{1.109878in}{1.532113in}}%
\pgfpathlineto{\pgfqpoint{1.109919in}{1.755670in}}%
\pgfpathlineto{\pgfqpoint{1.109981in}{2.022068in}}%
\pgfpathlineto{\pgfqpoint{1.110702in}{1.943299in}}%
\pgfpathlineto{\pgfqpoint{1.110705in}{1.938181in}}%
\pgfpathlineto{\pgfqpoint{1.110718in}{1.979610in}}%
\pgfpathlineto{\pgfqpoint{1.111457in}{1.963023in}}%
\pgfpathlineto{\pgfqpoint{1.111531in}{1.965669in}}%
\pgfpathlineto{\pgfqpoint{1.111972in}{1.921951in}}%
\pgfpathlineto{\pgfqpoint{1.112902in}{1.532024in}}%
\pgfpathlineto{\pgfqpoint{1.112615in}{1.939484in}}%
\pgfpathlineto{\pgfqpoint{1.113041in}{1.532037in}}%
\pgfpathlineto{\pgfqpoint{1.114825in}{1.532137in}}%
\pgfpathlineto{\pgfqpoint{1.114893in}{2.058256in}}%
\pgfpathlineto{\pgfqpoint{1.116023in}{1.969306in}}%
\pgfpathlineto{\pgfqpoint{1.116035in}{1.942468in}}%
\pgfpathlineto{\pgfqpoint{1.116049in}{1.975883in}}%
\pgfpathlineto{\pgfqpoint{1.116791in}{1.955011in}}%
\pgfpathlineto{\pgfqpoint{1.116810in}{1.961066in}}%
\pgfpathlineto{\pgfqpoint{1.116943in}{1.908924in}}%
\pgfpathlineto{\pgfqpoint{1.117831in}{1.532027in}}%
\pgfpathlineto{\pgfqpoint{1.117585in}{1.948254in}}%
\pgfpathlineto{\pgfqpoint{1.117955in}{1.532034in}}%
\pgfpathlineto{\pgfqpoint{1.119768in}{1.532111in}}%
\pgfpathlineto{\pgfqpoint{1.119811in}{1.769354in}}%
\pgfpathlineto{\pgfqpoint{1.119871in}{2.019910in}}%
\pgfpathlineto{\pgfqpoint{1.120590in}{1.945476in}}%
\pgfpathlineto{\pgfqpoint{1.120594in}{1.937245in}}%
\pgfpathlineto{\pgfqpoint{1.120611in}{1.977292in}}%
\pgfpathlineto{\pgfqpoint{1.121353in}{1.964289in}}%
\pgfpathlineto{\pgfqpoint{1.122192in}{1.581193in}}%
\pgfpathlineto{\pgfqpoint{1.122340in}{1.708253in}}%
\pgfpathlineto{\pgfqpoint{1.122538in}{1.950690in}}%
\pgfpathlineto{\pgfqpoint{1.122782in}{1.532026in}}%
\pgfpathlineto{\pgfqpoint{1.122962in}{1.532030in}}%
\pgfpathlineto{\pgfqpoint{1.124715in}{1.532122in}}%
\pgfpathlineto{\pgfqpoint{1.124780in}{2.047723in}}%
\pgfpathlineto{\pgfqpoint{1.126000in}{1.966853in}}%
\pgfpathlineto{\pgfqpoint{1.126012in}{1.944274in}}%
\pgfpathlineto{\pgfqpoint{1.126027in}{1.969714in}}%
\pgfpathlineto{\pgfqpoint{1.126774in}{1.953456in}}%
\pgfpathlineto{\pgfqpoint{1.126916in}{1.787075in}}%
\pgfpathlineto{\pgfqpoint{1.127700in}{1.532025in}}%
\pgfpathlineto{\pgfqpoint{1.127460in}{1.945972in}}%
\pgfpathlineto{\pgfqpoint{1.127821in}{1.532027in}}%
\pgfpathlineto{\pgfqpoint{1.129662in}{1.532143in}}%
\pgfpathlineto{\pgfqpoint{1.129730in}{2.043910in}}%
\pgfpathlineto{\pgfqpoint{1.130834in}{1.961687in}}%
\pgfpathlineto{\pgfqpoint{1.130843in}{1.940330in}}%
\pgfpathlineto{\pgfqpoint{1.130886in}{1.971568in}}%
\pgfpathlineto{\pgfqpoint{1.131604in}{1.957239in}}%
\pgfpathlineto{\pgfqpoint{1.131615in}{1.959734in}}%
\pgfpathlineto{\pgfqpoint{1.131804in}{1.892125in}}%
\pgfpathlineto{\pgfqpoint{1.132634in}{1.532015in}}%
\pgfpathlineto{\pgfqpoint{1.132447in}{1.944025in}}%
\pgfpathlineto{\pgfqpoint{1.132785in}{1.532025in}}%
\pgfpathlineto{\pgfqpoint{1.134605in}{1.532116in}}%
\pgfpathlineto{\pgfqpoint{1.134670in}{2.040000in}}%
\pgfpathlineto{\pgfqpoint{1.135981in}{1.961049in}}%
\pgfpathlineto{\pgfqpoint{1.136005in}{1.966757in}}%
\pgfpathlineto{\pgfqpoint{1.136811in}{1.804272in}}%
\pgfpathlineto{\pgfqpoint{1.137579in}{1.531998in}}%
\pgfpathlineto{\pgfqpoint{1.137404in}{1.945211in}}%
\pgfpathlineto{\pgfqpoint{1.137718in}{1.532023in}}%
\pgfpathlineto{\pgfqpoint{1.139555in}{1.532201in}}%
\pgfpathlineto{\pgfqpoint{1.139615in}{2.035383in}}%
\pgfpathlineto{\pgfqpoint{1.140692in}{1.964744in}}%
\pgfpathlineto{\pgfqpoint{1.140704in}{1.939682in}}%
\pgfpathlineto{\pgfqpoint{1.140717in}{1.971633in}}%
\pgfpathlineto{\pgfqpoint{1.141466in}{1.951512in}}%
\pgfpathlineto{\pgfqpoint{1.141478in}{1.958004in}}%
\pgfpathlineto{\pgfqpoint{1.141760in}{1.806099in}}%
\pgfpathlineto{\pgfqpoint{1.142525in}{1.532000in}}%
\pgfpathlineto{\pgfqpoint{1.142354in}{1.945417in}}%
\pgfpathlineto{\pgfqpoint{1.142661in}{1.532023in}}%
\pgfpathlineto{\pgfqpoint{1.144502in}{1.532430in}}%
\pgfpathlineto{\pgfqpoint{1.144606in}{2.041583in}}%
\pgfpathlineto{\pgfqpoint{1.145554in}{1.952820in}}%
\pgfpathlineto{\pgfqpoint{1.145562in}{1.934939in}}%
\pgfpathlineto{\pgfqpoint{1.145578in}{1.973258in}}%
\pgfpathlineto{\pgfqpoint{1.146318in}{1.951088in}}%
\pgfpathlineto{\pgfqpoint{1.146330in}{1.959676in}}%
\pgfpathlineto{\pgfqpoint{1.146769in}{1.709674in}}%
\pgfpathlineto{\pgfqpoint{1.147471in}{1.532004in}}%
\pgfpathlineto{\pgfqpoint{1.147263in}{1.941667in}}%
\pgfpathlineto{\pgfqpoint{1.147647in}{1.532039in}}%
\pgfpathlineto{\pgfqpoint{1.149440in}{1.532106in}}%
\pgfpathlineto{\pgfqpoint{1.149482in}{1.727429in}}%
\pgfpathlineto{\pgfqpoint{1.149512in}{2.015733in}}%
\pgfpathlineto{\pgfqpoint{1.150267in}{1.965221in}}%
\pgfpathlineto{\pgfqpoint{1.150277in}{1.933127in}}%
\pgfpathlineto{\pgfqpoint{1.150292in}{1.974893in}}%
\pgfpathlineto{\pgfqpoint{1.151037in}{1.961762in}}%
\pgfpathlineto{\pgfqpoint{1.151897in}{1.579538in}}%
\pgfpathlineto{\pgfqpoint{1.152046in}{1.697634in}}%
\pgfpathlineto{\pgfqpoint{1.152244in}{1.952694in}}%
\pgfpathlineto{\pgfqpoint{1.152442in}{1.532029in}}%
\pgfpathlineto{\pgfqpoint{1.152702in}{1.532034in}}%
\pgfpathlineto{\pgfqpoint{1.154392in}{1.532223in}}%
\pgfpathlineto{\pgfqpoint{1.154455in}{2.039045in}}%
\pgfpathlineto{\pgfqpoint{1.155497in}{1.970675in}}%
\pgfpathlineto{\pgfqpoint{1.155511in}{1.936166in}}%
\pgfpathlineto{\pgfqpoint{1.156271in}{1.954608in}}%
\pgfpathlineto{\pgfqpoint{1.156276in}{1.957220in}}%
\pgfpathlineto{\pgfqpoint{1.156646in}{1.746205in}}%
\pgfpathlineto{\pgfqpoint{1.157378in}{1.532024in}}%
\pgfpathlineto{\pgfqpoint{1.157190in}{1.946260in}}%
\pgfpathlineto{\pgfqpoint{1.157511in}{1.532026in}}%
\pgfpathlineto{\pgfqpoint{1.159334in}{1.532136in}}%
\pgfpathlineto{\pgfqpoint{1.159402in}{2.049885in}}%
\pgfpathlineto{\pgfqpoint{1.160511in}{1.939690in}}%
\pgfpathlineto{\pgfqpoint{1.160515in}{1.935074in}}%
\pgfpathlineto{\pgfqpoint{1.160527in}{1.969119in}}%
\pgfpathlineto{\pgfqpoint{1.161273in}{1.947979in}}%
\pgfpathlineto{\pgfqpoint{1.161286in}{1.956093in}}%
\pgfpathlineto{\pgfqpoint{1.161600in}{1.745283in}}%
\pgfpathlineto{\pgfqpoint{1.162343in}{1.532026in}}%
\pgfpathlineto{\pgfqpoint{1.162144in}{1.946392in}}%
\pgfpathlineto{\pgfqpoint{1.162469in}{1.532034in}}%
\pgfpathlineto{\pgfqpoint{1.164278in}{1.532113in}}%
\pgfpathlineto{\pgfqpoint{1.164320in}{1.763434in}}%
\pgfpathlineto{\pgfqpoint{1.164380in}{2.013362in}}%
\pgfpathlineto{\pgfqpoint{1.165098in}{1.950916in}}%
\pgfpathlineto{\pgfqpoint{1.165135in}{1.932294in}}%
\pgfpathlineto{\pgfqpoint{1.165121in}{1.972193in}}%
\pgfpathlineto{\pgfqpoint{1.165867in}{1.958074in}}%
\pgfpathlineto{\pgfqpoint{1.166713in}{1.600218in}}%
\pgfpathlineto{\pgfqpoint{1.165896in}{1.958679in}}%
\pgfpathlineto{\pgfqpoint{1.166861in}{1.666650in}}%
\pgfpathlineto{\pgfqpoint{1.167109in}{1.922182in}}%
\pgfpathlineto{\pgfqpoint{1.167271in}{1.532024in}}%
\pgfpathlineto{\pgfqpoint{1.167507in}{1.532029in}}%
\pgfpathlineto{\pgfqpoint{1.169225in}{1.532137in}}%
\pgfpathlineto{\pgfqpoint{1.169293in}{2.047757in}}%
\pgfpathlineto{\pgfqpoint{1.170402in}{1.938791in}}%
\pgfpathlineto{\pgfqpoint{1.170406in}{1.934028in}}%
\pgfpathlineto{\pgfqpoint{1.170418in}{1.967915in}}%
\pgfpathlineto{\pgfqpoint{1.171167in}{1.947172in}}%
\pgfpathlineto{\pgfqpoint{1.171180in}{1.955560in}}%
\pgfpathlineto{\pgfqpoint{1.171544in}{1.699828in}}%
\pgfpathlineto{\pgfqpoint{1.172200in}{1.532020in}}%
\pgfpathlineto{\pgfqpoint{1.172039in}{1.939482in}}%
\pgfpathlineto{\pgfqpoint{1.172401in}{1.532031in}}%
\pgfpathlineto{\pgfqpoint{1.174175in}{1.532472in}}%
\pgfpathlineto{\pgfqpoint{1.174278in}{2.039750in}}%
\pgfpathlineto{\pgfqpoint{1.175225in}{1.962586in}}%
\pgfpathlineto{\pgfqpoint{1.175233in}{1.931920in}}%
\pgfpathlineto{\pgfqpoint{1.175250in}{1.971641in}}%
\pgfpathlineto{\pgfqpoint{1.175996in}{1.948349in}}%
\pgfpathlineto{\pgfqpoint{1.176008in}{1.956204in}}%
\pgfpathlineto{\pgfqpoint{1.176457in}{1.757571in}}%
\pgfpathlineto{\pgfqpoint{1.177184in}{1.532026in}}%
\pgfpathlineto{\pgfqpoint{1.177001in}{1.942800in}}%
\pgfpathlineto{\pgfqpoint{1.177361in}{1.532030in}}%
\pgfpathlineto{\pgfqpoint{1.179114in}{1.532117in}}%
\pgfpathlineto{\pgfqpoint{1.179178in}{2.036164in}}%
\pgfpathlineto{\pgfqpoint{1.180491in}{1.948573in}}%
\pgfpathlineto{\pgfqpoint{1.180499in}{1.939587in}}%
\pgfpathlineto{\pgfqpoint{1.180511in}{1.961261in}}%
\pgfpathlineto{\pgfqpoint{1.181261in}{1.946388in}}%
\pgfpathlineto{\pgfqpoint{1.182092in}{1.532013in}}%
\pgfpathlineto{\pgfqpoint{1.182935in}{1.532040in}}%
\pgfpathlineto{\pgfqpoint{1.184059in}{1.532111in}}%
\pgfpathlineto{\pgfqpoint{1.184121in}{2.036684in}}%
\pgfpathlineto{\pgfqpoint{1.185501in}{1.942084in}}%
\pgfpathlineto{\pgfqpoint{1.186543in}{1.586100in}}%
\pgfpathlineto{\pgfqpoint{1.185520in}{1.960023in}}%
\pgfpathlineto{\pgfqpoint{1.186692in}{1.682221in}}%
\pgfpathlineto{\pgfqpoint{1.186889in}{1.940678in}}%
\pgfpathlineto{\pgfqpoint{1.187086in}{1.532027in}}%
\pgfpathlineto{\pgfqpoint{1.187324in}{1.532036in}}%
\pgfpathlineto{\pgfqpoint{1.189002in}{1.532095in}}%
\pgfpathlineto{\pgfqpoint{1.189010in}{1.532273in}}%
\pgfpathlineto{\pgfqpoint{1.189070in}{2.029741in}}%
\pgfpathlineto{\pgfqpoint{1.190090in}{1.963226in}}%
\pgfpathlineto{\pgfqpoint{1.190102in}{1.931677in}}%
\pgfpathlineto{\pgfqpoint{1.190117in}{1.966854in}}%
\pgfpathlineto{\pgfqpoint{1.190860in}{1.952654in}}%
\pgfpathlineto{\pgfqpoint{1.190866in}{1.953709in}}%
\pgfpathlineto{\pgfqpoint{1.191204in}{1.917409in}}%
\pgfpathlineto{\pgfqpoint{1.192032in}{1.532029in}}%
\pgfpathlineto{\pgfqpoint{1.191838in}{1.936292in}}%
\pgfpathlineto{\pgfqpoint{1.192267in}{1.532035in}}%
\pgfpathlineto{\pgfqpoint{1.193956in}{1.532261in}}%
\pgfpathlineto{\pgfqpoint{1.194017in}{2.031810in}}%
\pgfpathlineto{\pgfqpoint{1.195037in}{1.957571in}}%
\pgfpathlineto{\pgfqpoint{1.195045in}{1.931139in}}%
\pgfpathlineto{\pgfqpoint{1.195062in}{1.967132in}}%
\pgfpathlineto{\pgfqpoint{1.195810in}{1.950865in}}%
\pgfpathlineto{\pgfqpoint{1.195814in}{1.953644in}}%
\pgfpathlineto{\pgfqpoint{1.196233in}{1.817133in}}%
\pgfpathlineto{\pgfqpoint{1.197006in}{1.532029in}}%
\pgfpathlineto{\pgfqpoint{1.196826in}{1.937565in}}%
\pgfpathlineto{\pgfqpoint{1.197147in}{1.532034in}}%
\pgfpathlineto{\pgfqpoint{1.198897in}{1.532123in}}%
\pgfpathlineto{\pgfqpoint{1.198960in}{2.035017in}}%
\pgfpathlineto{\pgfqpoint{1.200216in}{1.948279in}}%
\pgfpathlineto{\pgfqpoint{1.200223in}{1.935853in}}%
\pgfpathlineto{\pgfqpoint{1.200238in}{1.960727in}}%
\pgfpathlineto{\pgfqpoint{1.200988in}{1.945622in}}%
\pgfpathlineto{\pgfqpoint{1.201003in}{1.949503in}}%
\pgfpathlineto{\pgfqpoint{1.201155in}{1.870965in}}%
\pgfpathlineto{\pgfqpoint{1.201919in}{1.532026in}}%
\pgfpathlineto{\pgfqpoint{1.201748in}{1.938747in}}%
\pgfpathlineto{\pgfqpoint{1.202103in}{1.532030in}}%
\pgfpathlineto{\pgfqpoint{1.203844in}{1.532144in}}%
\pgfpathlineto{\pgfqpoint{1.203912in}{2.024464in}}%
\pgfpathlineto{\pgfqpoint{1.205021in}{1.956181in}}%
\pgfpathlineto{\pgfqpoint{1.205031in}{1.933085in}}%
\pgfpathlineto{\pgfqpoint{1.205046in}{1.962836in}}%
\pgfpathlineto{\pgfqpoint{1.205793in}{1.946593in}}%
\pgfpathlineto{\pgfqpoint{1.206374in}{1.587029in}}%
\pgfpathlineto{\pgfqpoint{1.205821in}{1.951851in}}%
\pgfpathlineto{\pgfqpoint{1.206621in}{1.854276in}}%
\pgfpathlineto{\pgfqpoint{1.206714in}{1.939432in}}%
\pgfpathlineto{\pgfqpoint{1.206821in}{1.532013in}}%
\pgfpathlineto{\pgfqpoint{1.207095in}{1.532041in}}%
\pgfpathlineto{\pgfqpoint{1.208787in}{1.532119in}}%
\pgfpathlineto{\pgfqpoint{1.208852in}{2.035618in}}%
\pgfpathlineto{\pgfqpoint{1.210136in}{1.953295in}}%
\pgfpathlineto{\pgfqpoint{1.210175in}{1.935854in}}%
\pgfpathlineto{\pgfqpoint{1.210160in}{1.959686in}}%
\pgfpathlineto{\pgfqpoint{1.210907in}{1.945722in}}%
\pgfpathlineto{\pgfqpoint{1.210926in}{1.948304in}}%
\pgfpathlineto{\pgfqpoint{1.211115in}{1.791680in}}%
\pgfpathlineto{\pgfqpoint{1.211767in}{1.532002in}}%
\pgfpathlineto{\pgfqpoint{1.211695in}{1.938452in}}%
\pgfpathlineto{\pgfqpoint{1.212031in}{1.532042in}}%
\pgfpathlineto{\pgfqpoint{1.213731in}{1.532108in}}%
\pgfpathlineto{\pgfqpoint{1.213773in}{1.730111in}}%
\pgfpathlineto{\pgfqpoint{1.213836in}{1.999636in}}%
\pgfpathlineto{\pgfqpoint{1.214558in}{1.966169in}}%
\pgfpathlineto{\pgfqpoint{1.214574in}{1.928968in}}%
\pgfpathlineto{\pgfqpoint{1.215337in}{1.942934in}}%
\pgfpathlineto{\pgfqpoint{1.216315in}{1.600077in}}%
\pgfpathlineto{\pgfqpoint{1.215363in}{1.953226in}}%
\pgfpathlineto{\pgfqpoint{1.216513in}{1.797840in}}%
\pgfpathlineto{\pgfqpoint{1.216684in}{1.925137in}}%
\pgfpathlineto{\pgfqpoint{1.216747in}{1.532031in}}%
\pgfpathlineto{\pgfqpoint{1.217069in}{1.532037in}}%
\pgfpathlineto{\pgfqpoint{1.218683in}{1.532289in}}%
\pgfpathlineto{\pgfqpoint{1.218743in}{2.030262in}}%
\pgfpathlineto{\pgfqpoint{1.219742in}{1.930752in}}%
\pgfpathlineto{\pgfqpoint{1.219746in}{1.926696in}}%
\pgfpathlineto{\pgfqpoint{1.219761in}{1.965453in}}%
\pgfpathlineto{\pgfqpoint{1.220502in}{1.941549in}}%
\pgfpathlineto{\pgfqpoint{1.220513in}{1.951956in}}%
\pgfpathlineto{\pgfqpoint{1.221068in}{1.755865in}}%
\pgfpathlineto{\pgfqpoint{1.221720in}{1.532030in}}%
\pgfpathlineto{\pgfqpoint{1.221630in}{1.947617in}}%
\pgfpathlineto{\pgfqpoint{1.221960in}{1.532037in}}%
\pgfpathlineto{\pgfqpoint{1.223622in}{1.532111in}}%
\pgfpathlineto{\pgfqpoint{1.223665in}{1.748474in}}%
\pgfpathlineto{\pgfqpoint{1.223725in}{2.003527in}}%
\pgfpathlineto{\pgfqpoint{1.224448in}{1.951636in}}%
\pgfpathlineto{\pgfqpoint{1.224456in}{1.926375in}}%
\pgfpathlineto{\pgfqpoint{1.224471in}{1.965793in}}%
\pgfpathlineto{\pgfqpoint{1.225218in}{1.946268in}}%
\pgfpathlineto{\pgfqpoint{1.226205in}{1.604310in}}%
\pgfpathlineto{\pgfqpoint{1.225245in}{1.953521in}}%
\pgfpathlineto{\pgfqpoint{1.226304in}{1.633729in}}%
\pgfpathlineto{\pgfqpoint{1.226573in}{1.925847in}}%
\pgfpathlineto{\pgfqpoint{1.226602in}{1.532006in}}%
\pgfpathlineto{\pgfqpoint{1.226969in}{1.532045in}}%
\pgfpathlineto{\pgfqpoint{1.228567in}{1.532109in}}%
\pgfpathlineto{\pgfqpoint{1.228608in}{1.718344in}}%
\pgfpathlineto{\pgfqpoint{1.228640in}{2.019525in}}%
\pgfpathlineto{\pgfqpoint{1.229390in}{1.968197in}}%
\pgfpathlineto{\pgfqpoint{1.229403in}{1.921653in}}%
\pgfpathlineto{\pgfqpoint{1.229416in}{1.969375in}}%
\pgfpathlineto{\pgfqpoint{1.230164in}{1.954449in}}%
\pgfpathlineto{\pgfqpoint{1.231190in}{1.595979in}}%
\pgfpathlineto{\pgfqpoint{1.231239in}{1.607609in}}%
\pgfpathlineto{\pgfqpoint{1.231519in}{1.933303in}}%
\pgfpathlineto{\pgfqpoint{1.231548in}{1.532009in}}%
\pgfpathlineto{\pgfqpoint{1.231872in}{1.532030in}}%
\pgfpathlineto{\pgfqpoint{1.233513in}{1.532109in}}%
\pgfpathlineto{\pgfqpoint{1.233554in}{1.718403in}}%
\pgfpathlineto{\pgfqpoint{1.233585in}{2.019573in}}%
\pgfpathlineto{\pgfqpoint{1.234335in}{1.968094in}}%
\pgfpathlineto{\pgfqpoint{1.234348in}{1.921056in}}%
\pgfpathlineto{\pgfqpoint{1.234362in}{1.969146in}}%
\pgfpathlineto{\pgfqpoint{1.235105in}{1.951351in}}%
\pgfpathlineto{\pgfqpoint{1.235142in}{1.953810in}}%
\pgfpathlineto{\pgfqpoint{1.235128in}{1.937524in}}%
\pgfpathlineto{\pgfqpoint{1.235717in}{1.949651in}}%
\pgfpathlineto{\pgfqpoint{1.235964in}{1.765347in}}%
\pgfpathlineto{\pgfqpoint{1.236493in}{1.532011in}}%
\pgfpathlineto{\pgfqpoint{1.236464in}{1.923092in}}%
\pgfpathlineto{\pgfqpoint{1.236873in}{1.532047in}}%
\pgfpathlineto{\pgfqpoint{1.238458in}{1.532108in}}%
\pgfpathlineto{\pgfqpoint{1.238486in}{1.579855in}}%
\pgfpathlineto{\pgfqpoint{1.238570in}{2.041726in}}%
\pgfpathlineto{\pgfqpoint{1.239288in}{1.920683in}}%
\pgfpathlineto{\pgfqpoint{1.239292in}{1.913166in}}%
\pgfpathlineto{\pgfqpoint{1.239308in}{1.976196in}}%
\pgfpathlineto{\pgfqpoint{1.240044in}{1.940397in}}%
\pgfpathlineto{\pgfqpoint{1.240054in}{1.955450in}}%
\pgfpathlineto{\pgfqpoint{1.240069in}{1.935505in}}%
\pgfpathlineto{\pgfqpoint{1.240798in}{1.941498in}}%
\pgfpathlineto{\pgfqpoint{1.241456in}{1.532026in}}%
\pgfpathlineto{\pgfqpoint{1.242192in}{1.532042in}}%
\pgfpathlineto{\pgfqpoint{1.243406in}{1.532135in}}%
\pgfpathlineto{\pgfqpoint{1.243471in}{2.045020in}}%
\pgfpathlineto{\pgfqpoint{1.244641in}{1.939240in}}%
\pgfpathlineto{\pgfqpoint{1.244676in}{1.928736in}}%
\pgfpathlineto{\pgfqpoint{1.244661in}{1.962045in}}%
\pgfpathlineto{\pgfqpoint{1.245411in}{1.942708in}}%
\pgfpathlineto{\pgfqpoint{1.246096in}{1.588508in}}%
\pgfpathlineto{\pgfqpoint{1.245432in}{1.948846in}}%
\pgfpathlineto{\pgfqpoint{1.246322in}{1.809038in}}%
\pgfpathlineto{\pgfqpoint{1.246355in}{1.870286in}}%
\pgfpathlineto{\pgfqpoint{1.246403in}{1.532030in}}%
\pgfpathlineto{\pgfqpoint{1.246896in}{1.532046in}}%
\pgfpathlineto{\pgfqpoint{1.248351in}{1.532126in}}%
\pgfpathlineto{\pgfqpoint{1.248415in}{2.037286in}}%
\pgfpathlineto{\pgfqpoint{1.249698in}{1.954659in}}%
\pgfpathlineto{\pgfqpoint{1.249709in}{1.931683in}}%
\pgfpathlineto{\pgfqpoint{1.249724in}{1.957679in}}%
\pgfpathlineto{\pgfqpoint{1.250474in}{1.941451in}}%
\pgfpathlineto{\pgfqpoint{1.250489in}{1.947306in}}%
\pgfpathlineto{\pgfqpoint{1.250834in}{1.824476in}}%
\pgfpathlineto{\pgfqpoint{1.251327in}{1.532017in}}%
\pgfpathlineto{\pgfqpoint{1.251301in}{1.839943in}}%
\pgfpathlineto{\pgfqpoint{1.251779in}{1.532034in}}%
\pgfpathlineto{\pgfqpoint{1.253303in}{1.532932in}}%
\pgfpathlineto{\pgfqpoint{1.253406in}{2.042193in}}%
\pgfpathlineto{\pgfqpoint{1.254325in}{1.950659in}}%
\pgfpathlineto{\pgfqpoint{1.254334in}{1.920427in}}%
\pgfpathlineto{\pgfqpoint{1.254350in}{1.967313in}}%
\pgfpathlineto{\pgfqpoint{1.255096in}{1.948935in}}%
\pgfpathlineto{\pgfqpoint{1.255102in}{1.951729in}}%
\pgfpathlineto{\pgfqpoint{1.255749in}{1.894274in}}%
\pgfpathlineto{\pgfqpoint{1.256335in}{1.532027in}}%
\pgfpathlineto{\pgfqpoint{1.256723in}{1.532036in}}%
\pgfpathlineto{\pgfqpoint{1.258243in}{1.532150in}}%
\pgfpathlineto{\pgfqpoint{1.258345in}{2.001766in}}%
\pgfpathlineto{\pgfqpoint{1.259406in}{1.951425in}}%
\pgfpathlineto{\pgfqpoint{1.259417in}{1.931782in}}%
\pgfpathlineto{\pgfqpoint{1.259433in}{1.956723in}}%
\pgfpathlineto{\pgfqpoint{1.260162in}{1.945719in}}%
\pgfpathlineto{\pgfqpoint{1.261240in}{1.532028in}}%
\pgfpathlineto{\pgfqpoint{1.260231in}{1.950168in}}%
\pgfpathlineto{\pgfqpoint{1.262586in}{1.532058in}}%
\pgfpathlineto{\pgfqpoint{1.263187in}{1.532121in}}%
\pgfpathlineto{\pgfqpoint{1.263227in}{1.736672in}}%
\pgfpathlineto{\pgfqpoint{1.263244in}{1.995934in}}%
\pgfpathlineto{\pgfqpoint{1.264011in}{1.955501in}}%
\pgfpathlineto{\pgfqpoint{1.264021in}{1.924814in}}%
\pgfpathlineto{\pgfqpoint{1.264037in}{1.963549in}}%
\pgfpathlineto{\pgfqpoint{1.264782in}{1.952097in}}%
\pgfpathlineto{\pgfqpoint{1.265551in}{1.927942in}}%
\pgfpathlineto{\pgfqpoint{1.265600in}{1.931906in}}%
\pgfpathlineto{\pgfqpoint{1.265649in}{1.893814in}}%
\pgfpathlineto{\pgfqpoint{1.266160in}{1.532018in}}%
\pgfpathlineto{\pgfqpoint{1.266667in}{1.532036in}}%
\pgfpathlineto{\pgfqpoint{1.268133in}{1.532126in}}%
\pgfpathlineto{\pgfqpoint{1.268196in}{2.037959in}}%
\pgfpathlineto{\pgfqpoint{1.269482in}{1.943210in}}%
\pgfpathlineto{\pgfqpoint{1.269489in}{1.930651in}}%
\pgfpathlineto{\pgfqpoint{1.269505in}{1.957031in}}%
\pgfpathlineto{\pgfqpoint{1.270251in}{1.941494in}}%
\pgfpathlineto{\pgfqpoint{1.270270in}{1.946752in}}%
\pgfpathlineto{\pgfqpoint{1.270703in}{1.728988in}}%
\pgfpathlineto{\pgfqpoint{1.271186in}{1.532032in}}%
\pgfpathlineto{\pgfqpoint{1.271082in}{1.792253in}}%
\pgfpathlineto{\pgfqpoint{1.271578in}{1.532041in}}%
\pgfpathlineto{\pgfqpoint{1.273076in}{1.532110in}}%
\pgfpathlineto{\pgfqpoint{1.273118in}{1.724879in}}%
\pgfpathlineto{\pgfqpoint{1.273148in}{2.010473in}}%
\pgfpathlineto{\pgfqpoint{1.273903in}{1.959232in}}%
\pgfpathlineto{\pgfqpoint{1.273914in}{1.921313in}}%
\pgfpathlineto{\pgfqpoint{1.273930in}{1.966441in}}%
\pgfpathlineto{\pgfqpoint{1.274674in}{1.935675in}}%
\pgfpathlineto{\pgfqpoint{1.274727in}{1.951987in}}%
\pgfpathlineto{\pgfqpoint{1.275421in}{1.948371in}}%
\pgfpathlineto{\pgfqpoint{1.275619in}{1.794068in}}%
\pgfpathlineto{\pgfqpoint{1.276089in}{1.532027in}}%
\pgfpathlineto{\pgfqpoint{1.276513in}{1.532048in}}%
\pgfpathlineto{\pgfqpoint{1.278023in}{1.532121in}}%
\pgfpathlineto{\pgfqpoint{1.278063in}{1.728607in}}%
\pgfpathlineto{\pgfqpoint{1.278093in}{2.008439in}}%
\pgfpathlineto{\pgfqpoint{1.278846in}{1.945520in}}%
\pgfpathlineto{\pgfqpoint{1.278854in}{1.921147in}}%
\pgfpathlineto{\pgfqpoint{1.278867in}{1.965467in}}%
\pgfpathlineto{\pgfqpoint{1.279615in}{1.951751in}}%
\pgfpathlineto{\pgfqpoint{1.280790in}{1.610352in}}%
\pgfpathlineto{\pgfqpoint{1.280839in}{1.626876in}}%
\pgfpathlineto{\pgfqpoint{1.280973in}{1.778955in}}%
\pgfpathlineto{\pgfqpoint{1.280996in}{1.532024in}}%
\pgfpathlineto{\pgfqpoint{1.281465in}{1.532048in}}%
\pgfpathlineto{\pgfqpoint{1.282959in}{1.532083in}}%
\pgfpathlineto{\pgfqpoint{1.282971in}{1.532145in}}%
\pgfpathlineto{\pgfqpoint{1.283076in}{2.002155in}}%
\pgfpathlineto{\pgfqpoint{1.284131in}{1.954734in}}%
\pgfpathlineto{\pgfqpoint{1.284142in}{1.930989in}}%
\pgfpathlineto{\pgfqpoint{1.284159in}{1.956845in}}%
\pgfpathlineto{\pgfqpoint{1.284894in}{1.940146in}}%
\pgfpathlineto{\pgfqpoint{1.285027in}{1.949507in}}%
\pgfpathlineto{\pgfqpoint{1.285324in}{1.930177in}}%
\pgfpathlineto{\pgfqpoint{1.285373in}{1.940720in}}%
\pgfpathlineto{\pgfqpoint{1.285963in}{1.532030in}}%
\pgfpathlineto{\pgfqpoint{1.286703in}{1.532045in}}%
\pgfpathlineto{\pgfqpoint{1.287915in}{1.532134in}}%
\pgfpathlineto{\pgfqpoint{1.287981in}{2.046253in}}%
\pgfpathlineto{\pgfqpoint{1.289151in}{1.933387in}}%
\pgfpathlineto{\pgfqpoint{1.289155in}{1.925811in}}%
\pgfpathlineto{\pgfqpoint{1.289171in}{1.960969in}}%
\pgfpathlineto{\pgfqpoint{1.289912in}{1.947170in}}%
\pgfpathlineto{\pgfqpoint{1.289916in}{1.948171in}}%
\pgfpathlineto{\pgfqpoint{1.290383in}{1.899488in}}%
\pgfpathlineto{\pgfqpoint{1.290924in}{1.532027in}}%
\pgfpathlineto{\pgfqpoint{1.291397in}{1.532038in}}%
\pgfpathlineto{\pgfqpoint{1.292859in}{1.532121in}}%
\pgfpathlineto{\pgfqpoint{1.292899in}{1.724263in}}%
\pgfpathlineto{\pgfqpoint{1.292930in}{2.012672in}}%
\pgfpathlineto{\pgfqpoint{1.293684in}{1.946057in}}%
\pgfpathlineto{\pgfqpoint{1.293693in}{1.920212in}}%
\pgfpathlineto{\pgfqpoint{1.293706in}{1.966397in}}%
\pgfpathlineto{\pgfqpoint{1.294456in}{1.937594in}}%
\pgfpathlineto{\pgfqpoint{1.294462in}{1.936196in}}%
\pgfpathlineto{\pgfqpoint{1.294476in}{1.951318in}}%
\pgfpathlineto{\pgfqpoint{1.295180in}{1.946750in}}%
\pgfpathlineto{\pgfqpoint{1.295377in}{1.844725in}}%
\pgfpathlineto{\pgfqpoint{1.295833in}{1.532019in}}%
\pgfpathlineto{\pgfqpoint{1.296330in}{1.532050in}}%
\pgfpathlineto{\pgfqpoint{1.297804in}{1.532114in}}%
\pgfpathlineto{\pgfqpoint{1.297846in}{1.728474in}}%
\pgfpathlineto{\pgfqpoint{1.297909in}{2.000957in}}%
\pgfpathlineto{\pgfqpoint{1.298628in}{1.964435in}}%
\pgfpathlineto{\pgfqpoint{1.298643in}{1.922564in}}%
\pgfpathlineto{\pgfqpoint{1.299407in}{1.947693in}}%
\pgfpathlineto{\pgfqpoint{1.299415in}{1.950697in}}%
\pgfpathlineto{\pgfqpoint{1.299433in}{1.935945in}}%
\pgfpathlineto{\pgfqpoint{1.300172in}{1.945963in}}%
\pgfpathlineto{\pgfqpoint{1.300794in}{1.532026in}}%
\pgfpathlineto{\pgfqpoint{1.301622in}{1.532043in}}%
\pgfpathlineto{\pgfqpoint{1.302750in}{1.532116in}}%
\pgfpathlineto{\pgfqpoint{1.302791in}{1.739235in}}%
\pgfpathlineto{\pgfqpoint{1.302854in}{2.006279in}}%
\pgfpathlineto{\pgfqpoint{1.303575in}{1.943646in}}%
\pgfpathlineto{\pgfqpoint{1.303583in}{1.921192in}}%
\pgfpathlineto{\pgfqpoint{1.303597in}{1.964434in}}%
\pgfpathlineto{\pgfqpoint{1.304341in}{1.940299in}}%
\pgfpathlineto{\pgfqpoint{1.304351in}{1.951009in}}%
\pgfpathlineto{\pgfqpoint{1.304373in}{1.935414in}}%
\pgfpathlineto{\pgfqpoint{1.305109in}{1.937029in}}%
\pgfpathlineto{\pgfqpoint{1.305159in}{1.942399in}}%
\pgfpathlineto{\pgfqpoint{1.305208in}{1.903218in}}%
\pgfpathlineto{\pgfqpoint{1.305723in}{1.532026in}}%
\pgfpathlineto{\pgfqpoint{1.306289in}{1.532050in}}%
\pgfpathlineto{\pgfqpoint{1.307696in}{1.532121in}}%
\pgfpathlineto{\pgfqpoint{1.307737in}{1.742840in}}%
\pgfpathlineto{\pgfqpoint{1.307799in}{2.005988in}}%
\pgfpathlineto{\pgfqpoint{1.308518in}{1.945053in}}%
\pgfpathlineto{\pgfqpoint{1.308526in}{1.920871in}}%
\pgfpathlineto{\pgfqpoint{1.308539in}{1.963970in}}%
\pgfpathlineto{\pgfqpoint{1.309288in}{1.950849in}}%
\pgfpathlineto{\pgfqpoint{1.309340in}{1.935736in}}%
\pgfpathlineto{\pgfqpoint{1.310074in}{1.939978in}}%
\pgfpathlineto{\pgfqpoint{1.310669in}{1.532016in}}%
\pgfpathlineto{\pgfqpoint{1.311512in}{1.532043in}}%
\pgfpathlineto{\pgfqpoint{1.312640in}{1.532109in}}%
\pgfpathlineto{\pgfqpoint{1.312668in}{1.578242in}}%
\pgfpathlineto{\pgfqpoint{1.312752in}{2.040154in}}%
\pgfpathlineto{\pgfqpoint{1.313470in}{1.917744in}}%
\pgfpathlineto{\pgfqpoint{1.313474in}{1.909813in}}%
\pgfpathlineto{\pgfqpoint{1.313490in}{1.974449in}}%
\pgfpathlineto{\pgfqpoint{1.314227in}{1.937517in}}%
\pgfpathlineto{\pgfqpoint{1.314237in}{1.953375in}}%
\pgfpathlineto{\pgfqpoint{1.314252in}{1.932631in}}%
\pgfpathlineto{\pgfqpoint{1.314986in}{1.940089in}}%
\pgfpathlineto{\pgfqpoint{1.315020in}{1.944269in}}%
\pgfpathlineto{\pgfqpoint{1.315090in}{1.918243in}}%
\pgfpathlineto{\pgfqpoint{1.315670in}{1.532028in}}%
\pgfpathlineto{\pgfqpoint{1.316157in}{1.532038in}}%
\pgfpathlineto{\pgfqpoint{1.317589in}{1.532148in}}%
\pgfpathlineto{\pgfqpoint{1.317657in}{2.025095in}}%
\pgfpathlineto{\pgfqpoint{1.318763in}{1.956526in}}%
\pgfpathlineto{\pgfqpoint{1.318804in}{1.927169in}}%
\pgfpathlineto{\pgfqpoint{1.318791in}{1.958634in}}%
\pgfpathlineto{\pgfqpoint{1.319529in}{1.941516in}}%
\pgfpathlineto{\pgfqpoint{1.319538in}{1.947837in}}%
\pgfpathlineto{\pgfqpoint{1.320117in}{1.828881in}}%
\pgfpathlineto{\pgfqpoint{1.320560in}{1.532022in}}%
\pgfpathlineto{\pgfqpoint{1.321054in}{1.532050in}}%
\pgfpathlineto{\pgfqpoint{1.322522in}{1.532083in}}%
\pgfpathlineto{\pgfqpoint{1.322533in}{1.532126in}}%
\pgfpathlineto{\pgfqpoint{1.322597in}{2.037685in}}%
\pgfpathlineto{\pgfqpoint{1.323857in}{1.938462in}}%
\pgfpathlineto{\pgfqpoint{1.323892in}{1.928817in}}%
\pgfpathlineto{\pgfqpoint{1.323877in}{1.956604in}}%
\pgfpathlineto{\pgfqpoint{1.324624in}{1.941691in}}%
\pgfpathlineto{\pgfqpoint{1.325521in}{1.532026in}}%
\pgfpathlineto{\pgfqpoint{1.324647in}{1.945498in}}%
\pgfpathlineto{\pgfqpoint{1.326101in}{1.532052in}}%
\pgfpathlineto{\pgfqpoint{1.327477in}{1.532119in}}%
\pgfpathlineto{\pgfqpoint{1.327519in}{1.757866in}}%
\pgfpathlineto{\pgfqpoint{1.327580in}{2.000426in}}%
\pgfpathlineto{\pgfqpoint{1.328302in}{1.936479in}}%
\pgfpathlineto{\pgfqpoint{1.328307in}{1.921977in}}%
\pgfpathlineto{\pgfqpoint{1.328322in}{1.962683in}}%
\pgfpathlineto{\pgfqpoint{1.329067in}{1.939974in}}%
\pgfpathlineto{\pgfqpoint{1.329077in}{1.950014in}}%
\pgfpathlineto{\pgfqpoint{1.329099in}{1.935144in}}%
\pgfpathlineto{\pgfqpoint{1.329807in}{1.938721in}}%
\pgfpathlineto{\pgfqpoint{1.329856in}{1.945584in}}%
\pgfpathlineto{\pgfqpoint{1.329955in}{1.904843in}}%
\pgfpathlineto{\pgfqpoint{1.330499in}{1.532028in}}%
\pgfpathlineto{\pgfqpoint{1.331033in}{1.532049in}}%
\pgfpathlineto{\pgfqpoint{1.332425in}{1.532144in}}%
\pgfpathlineto{\pgfqpoint{1.332492in}{2.046078in}}%
\pgfpathlineto{\pgfqpoint{1.333601in}{1.938322in}}%
\pgfpathlineto{\pgfqpoint{1.333608in}{1.923268in}}%
\pgfpathlineto{\pgfqpoint{1.333622in}{1.961769in}}%
\pgfpathlineto{\pgfqpoint{1.334366in}{1.945131in}}%
\pgfpathlineto{\pgfqpoint{1.334371in}{1.947830in}}%
\pgfpathlineto{\pgfqpoint{1.334910in}{1.893920in}}%
\pgfpathlineto{\pgfqpoint{1.335412in}{1.532026in}}%
\pgfpathlineto{\pgfqpoint{1.335893in}{1.532050in}}%
\pgfpathlineto{\pgfqpoint{1.337368in}{1.532119in}}%
\pgfpathlineto{\pgfqpoint{1.337410in}{1.760019in}}%
\pgfpathlineto{\pgfqpoint{1.337471in}{2.000379in}}%
\pgfpathlineto{\pgfqpoint{1.338191in}{1.941803in}}%
\pgfpathlineto{\pgfqpoint{1.338228in}{1.921797in}}%
\pgfpathlineto{\pgfqpoint{1.338215in}{1.962001in}}%
\pgfpathlineto{\pgfqpoint{1.338960in}{1.948547in}}%
\pgfpathlineto{\pgfqpoint{1.338982in}{1.934903in}}%
\pgfpathlineto{\pgfqpoint{1.339004in}{1.949484in}}%
\pgfpathlineto{\pgfqpoint{1.339706in}{1.936695in}}%
\pgfpathlineto{\pgfqpoint{1.339756in}{1.945036in}}%
\pgfpathlineto{\pgfqpoint{1.339855in}{1.893182in}}%
\pgfpathlineto{\pgfqpoint{1.340391in}{1.532028in}}%
\pgfpathlineto{\pgfqpoint{1.340876in}{1.532038in}}%
\pgfpathlineto{\pgfqpoint{1.342316in}{1.532147in}}%
\pgfpathlineto{\pgfqpoint{1.342384in}{2.030722in}}%
\pgfpathlineto{\pgfqpoint{1.343489in}{1.950019in}}%
\pgfpathlineto{\pgfqpoint{1.343498in}{1.925139in}}%
\pgfpathlineto{\pgfqpoint{1.343515in}{1.958627in}}%
\pgfpathlineto{\pgfqpoint{1.344259in}{1.938929in}}%
\pgfpathlineto{\pgfqpoint{1.345343in}{1.532026in}}%
\pgfpathlineto{\pgfqpoint{1.344279in}{1.947017in}}%
\pgfpathlineto{\pgfqpoint{1.346621in}{1.532055in}}%
\pgfpathlineto{\pgfqpoint{1.347260in}{1.532129in}}%
\pgfpathlineto{\pgfqpoint{1.347325in}{2.039818in}}%
\pgfpathlineto{\pgfqpoint{1.348554in}{1.942280in}}%
\pgfpathlineto{\pgfqpoint{1.348562in}{1.926470in}}%
\pgfpathlineto{\pgfqpoint{1.348606in}{1.956626in}}%
\pgfpathlineto{\pgfqpoint{1.349322in}{1.944316in}}%
\pgfpathlineto{\pgfqpoint{1.350289in}{1.532027in}}%
\pgfpathlineto{\pgfqpoint{1.349348in}{1.945551in}}%
\pgfpathlineto{\pgfqpoint{1.350578in}{1.532033in}}%
\pgfpathlineto{\pgfqpoint{1.352206in}{1.532130in}}%
\pgfpathlineto{\pgfqpoint{1.352270in}{2.040193in}}%
\pgfpathlineto{\pgfqpoint{1.353496in}{1.950290in}}%
\pgfpathlineto{\pgfqpoint{1.353536in}{1.927015in}}%
\pgfpathlineto{\pgfqpoint{1.353521in}{1.957202in}}%
\pgfpathlineto{\pgfqpoint{1.354271in}{1.940223in}}%
\pgfpathlineto{\pgfqpoint{1.355212in}{1.532031in}}%
\pgfpathlineto{\pgfqpoint{1.354291in}{1.945392in}}%
\pgfpathlineto{\pgfqpoint{1.355783in}{1.532042in}}%
\pgfpathlineto{\pgfqpoint{1.357151in}{1.532127in}}%
\pgfpathlineto{\pgfqpoint{1.357215in}{2.035878in}}%
\pgfpathlineto{\pgfqpoint{1.358477in}{1.930295in}}%
\pgfpathlineto{\pgfqpoint{1.358481in}{1.927707in}}%
\pgfpathlineto{\pgfqpoint{1.358493in}{1.954903in}}%
\pgfpathlineto{\pgfqpoint{1.359225in}{1.942348in}}%
\pgfpathlineto{\pgfqpoint{1.359233in}{1.945016in}}%
\pgfpathlineto{\pgfqpoint{1.359727in}{1.777764in}}%
\pgfpathlineto{\pgfqpoint{1.360212in}{1.532033in}}%
\pgfpathlineto{\pgfqpoint{1.360653in}{1.532045in}}%
\pgfpathlineto{\pgfqpoint{1.362096in}{1.532127in}}%
\pgfpathlineto{\pgfqpoint{1.362160in}{2.036870in}}%
\pgfpathlineto{\pgfqpoint{1.363420in}{1.937743in}}%
\pgfpathlineto{\pgfqpoint{1.363456in}{1.927879in}}%
\pgfpathlineto{\pgfqpoint{1.363440in}{1.955596in}}%
\pgfpathlineto{\pgfqpoint{1.364191in}{1.938181in}}%
\pgfpathlineto{\pgfqpoint{1.364207in}{1.944730in}}%
\pgfpathlineto{\pgfqpoint{1.364588in}{1.902691in}}%
\pgfpathlineto{\pgfqpoint{1.365117in}{1.532033in}}%
\pgfpathlineto{\pgfqpoint{1.365601in}{1.532043in}}%
\pgfpathlineto{\pgfqpoint{1.367046in}{1.532218in}}%
\pgfpathlineto{\pgfqpoint{1.367108in}{2.040053in}}%
\pgfpathlineto{\pgfqpoint{1.368156in}{1.952999in}}%
\pgfpathlineto{\pgfqpoint{1.368196in}{1.921518in}}%
\pgfpathlineto{\pgfqpoint{1.368180in}{1.962098in}}%
\pgfpathlineto{\pgfqpoint{1.368931in}{1.938984in}}%
\pgfpathlineto{\pgfqpoint{1.369818in}{1.591183in}}%
\pgfpathlineto{\pgfqpoint{1.368952in}{1.947045in}}%
\pgfpathlineto{\pgfqpoint{1.369952in}{1.668309in}}%
\pgfpathlineto{\pgfqpoint{1.369992in}{1.732118in}}%
\pgfpathlineto{\pgfqpoint{1.370109in}{1.532033in}}%
\pgfpathlineto{\pgfqpoint{1.370603in}{1.532045in}}%
\pgfpathlineto{\pgfqpoint{1.371991in}{1.532192in}}%
\pgfpathlineto{\pgfqpoint{1.372054in}{2.042010in}}%
\pgfpathlineto{\pgfqpoint{1.373107in}{1.926356in}}%
\pgfpathlineto{\pgfqpoint{1.373111in}{1.920316in}}%
\pgfpathlineto{\pgfqpoint{1.373127in}{1.961508in}}%
\pgfpathlineto{\pgfqpoint{1.373864in}{1.947105in}}%
\pgfpathlineto{\pgfqpoint{1.374759in}{1.591438in}}%
\pgfpathlineto{\pgfqpoint{1.374895in}{1.663662in}}%
\pgfpathlineto{\pgfqpoint{1.374937in}{1.729837in}}%
\pgfpathlineto{\pgfqpoint{1.374996in}{1.532031in}}%
\pgfpathlineto{\pgfqpoint{1.375568in}{1.532043in}}%
\pgfpathlineto{\pgfqpoint{1.376938in}{1.532311in}}%
\pgfpathlineto{\pgfqpoint{1.377071in}{2.031268in}}%
\pgfpathlineto{\pgfqpoint{1.377994in}{1.938379in}}%
\pgfpathlineto{\pgfqpoint{1.378001in}{1.919661in}}%
\pgfpathlineto{\pgfqpoint{1.378014in}{1.962858in}}%
\pgfpathlineto{\pgfqpoint{1.378759in}{1.944523in}}%
\pgfpathlineto{\pgfqpoint{1.378765in}{1.947917in}}%
\pgfpathlineto{\pgfqpoint{1.379425in}{1.903751in}}%
\pgfpathlineto{\pgfqpoint{1.379928in}{1.532029in}}%
\pgfpathlineto{\pgfqpoint{1.380482in}{1.532041in}}%
\pgfpathlineto{\pgfqpoint{1.381878in}{1.532122in}}%
\pgfpathlineto{\pgfqpoint{1.381919in}{1.763896in}}%
\pgfpathlineto{\pgfqpoint{1.382010in}{2.000229in}}%
\pgfpathlineto{\pgfqpoint{1.382698in}{1.939740in}}%
\pgfpathlineto{\pgfqpoint{1.382706in}{1.919427in}}%
\pgfpathlineto{\pgfqpoint{1.382719in}{1.961815in}}%
\pgfpathlineto{\pgfqpoint{1.383465in}{1.945346in}}%
\pgfpathlineto{\pgfqpoint{1.383513in}{1.934442in}}%
\pgfpathlineto{\pgfqpoint{1.383526in}{1.948220in}}%
\pgfpathlineto{\pgfqpoint{1.384212in}{1.936170in}}%
\pgfpathlineto{\pgfqpoint{1.384261in}{1.944424in}}%
\pgfpathlineto{\pgfqpoint{1.384409in}{1.849613in}}%
\pgfpathlineto{\pgfqpoint{1.384888in}{1.532031in}}%
\pgfpathlineto{\pgfqpoint{1.385362in}{1.532041in}}%
\pgfpathlineto{\pgfqpoint{1.386829in}{1.532310in}}%
\pgfpathlineto{\pgfqpoint{1.386962in}{2.030788in}}%
\pgfpathlineto{\pgfqpoint{1.387884in}{1.938500in}}%
\pgfpathlineto{\pgfqpoint{1.387892in}{1.919454in}}%
\pgfpathlineto{\pgfqpoint{1.387905in}{1.962477in}}%
\pgfpathlineto{\pgfqpoint{1.388650in}{1.944231in}}%
\pgfpathlineto{\pgfqpoint{1.388655in}{1.947642in}}%
\pgfpathlineto{\pgfqpoint{1.389327in}{1.886897in}}%
\pgfpathlineto{\pgfqpoint{1.389837in}{1.532033in}}%
\pgfpathlineto{\pgfqpoint{1.390317in}{1.532043in}}%
\pgfpathlineto{\pgfqpoint{1.391774in}{1.532248in}}%
\pgfpathlineto{\pgfqpoint{1.391836in}{2.034064in}}%
\pgfpathlineto{\pgfqpoint{1.392859in}{1.938266in}}%
\pgfpathlineto{\pgfqpoint{1.392867in}{1.920811in}}%
\pgfpathlineto{\pgfqpoint{1.392880in}{1.960941in}}%
\pgfpathlineto{\pgfqpoint{1.393628in}{1.947047in}}%
\pgfpathlineto{\pgfqpoint{1.394740in}{1.532023in}}%
\pgfpathlineto{\pgfqpoint{1.395156in}{1.532035in}}%
\pgfpathlineto{\pgfqpoint{1.396714in}{1.532122in}}%
\pgfpathlineto{\pgfqpoint{1.396756in}{1.766586in}}%
\pgfpathlineto{\pgfqpoint{1.396817in}{1.999925in}}%
\pgfpathlineto{\pgfqpoint{1.397535in}{1.942233in}}%
\pgfpathlineto{\pgfqpoint{1.397544in}{1.919881in}}%
\pgfpathlineto{\pgfqpoint{1.397557in}{1.962544in}}%
\pgfpathlineto{\pgfqpoint{1.398306in}{1.940581in}}%
\pgfpathlineto{\pgfqpoint{1.398314in}{1.933400in}}%
\pgfpathlineto{\pgfqpoint{1.398366in}{1.947715in}}%
\pgfpathlineto{\pgfqpoint{1.399037in}{1.936200in}}%
\pgfpathlineto{\pgfqpoint{1.399087in}{1.944396in}}%
\pgfpathlineto{\pgfqpoint{1.399235in}{1.869808in}}%
\pgfpathlineto{\pgfqpoint{1.399685in}{1.532016in}}%
\pgfpathlineto{\pgfqpoint{1.400211in}{1.532036in}}%
\pgfpathlineto{\pgfqpoint{1.401659in}{1.532116in}}%
\pgfpathlineto{\pgfqpoint{1.401700in}{1.734378in}}%
\pgfpathlineto{\pgfqpoint{1.401717in}{1.992348in}}%
\pgfpathlineto{\pgfqpoint{1.402485in}{1.959523in}}%
\pgfpathlineto{\pgfqpoint{1.402500in}{1.922522in}}%
\pgfpathlineto{\pgfqpoint{1.403265in}{1.934278in}}%
\pgfpathlineto{\pgfqpoint{1.403590in}{1.947525in}}%
\pgfpathlineto{\pgfqpoint{1.403340in}{1.933769in}}%
\pgfpathlineto{\pgfqpoint{1.404035in}{1.946445in}}%
\pgfpathlineto{\pgfqpoint{1.404282in}{1.720685in}}%
\pgfpathlineto{\pgfqpoint{1.404695in}{1.532030in}}%
\pgfpathlineto{\pgfqpoint{1.404610in}{1.731584in}}%
\pgfpathlineto{\pgfqpoint{1.405184in}{1.532041in}}%
\pgfpathlineto{\pgfqpoint{1.406605in}{1.532124in}}%
\pgfpathlineto{\pgfqpoint{1.406670in}{2.038375in}}%
\pgfpathlineto{\pgfqpoint{1.407928in}{1.944480in}}%
\pgfpathlineto{\pgfqpoint{1.407935in}{1.926399in}}%
\pgfpathlineto{\pgfqpoint{1.407952in}{1.955136in}}%
\pgfpathlineto{\pgfqpoint{1.408701in}{1.939785in}}%
\pgfpathlineto{\pgfqpoint{1.409664in}{1.532033in}}%
\pgfpathlineto{\pgfqpoint{1.408721in}{1.943795in}}%
\pgfpathlineto{\pgfqpoint{1.410155in}{1.532044in}}%
\pgfpathlineto{\pgfqpoint{1.411551in}{1.532125in}}%
\pgfpathlineto{\pgfqpoint{1.411608in}{2.019086in}}%
\pgfpathlineto{\pgfqpoint{1.413002in}{1.942914in}}%
\pgfpathlineto{\pgfqpoint{1.413012in}{1.930209in}}%
\pgfpathlineto{\pgfqpoint{1.413028in}{1.950854in}}%
\pgfpathlineto{\pgfqpoint{1.413767in}{1.936021in}}%
\pgfpathlineto{\pgfqpoint{1.413879in}{1.943565in}}%
\pgfpathlineto{\pgfqpoint{1.414077in}{1.880525in}}%
\pgfpathlineto{\pgfqpoint{1.414522in}{1.532017in}}%
\pgfpathlineto{\pgfqpoint{1.415073in}{1.532036in}}%
\pgfpathlineto{\pgfqpoint{1.416503in}{1.532664in}}%
\pgfpathlineto{\pgfqpoint{1.416605in}{2.038427in}}%
\pgfpathlineto{\pgfqpoint{1.417525in}{1.949094in}}%
\pgfpathlineto{\pgfqpoint{1.417534in}{1.916348in}}%
\pgfpathlineto{\pgfqpoint{1.417550in}{1.964175in}}%
\pgfpathlineto{\pgfqpoint{1.418294in}{1.944172in}}%
\pgfpathlineto{\pgfqpoint{1.418299in}{1.947940in}}%
\pgfpathlineto{\pgfqpoint{1.418973in}{1.928349in}}%
\pgfpathlineto{\pgfqpoint{1.419467in}{1.532017in}}%
\pgfpathlineto{\pgfqpoint{1.420166in}{1.532056in}}%
\pgfpathlineto{\pgfqpoint{1.421439in}{1.532106in}}%
\pgfpathlineto{\pgfqpoint{1.421466in}{1.572315in}}%
\pgfpathlineto{\pgfqpoint{1.421551in}{2.038160in}}%
\pgfpathlineto{\pgfqpoint{1.422269in}{1.920879in}}%
\pgfpathlineto{\pgfqpoint{1.422273in}{1.907229in}}%
\pgfpathlineto{\pgfqpoint{1.422289in}{1.972888in}}%
\pgfpathlineto{\pgfqpoint{1.423031in}{1.950849in}}%
\pgfpathlineto{\pgfqpoint{1.423045in}{1.930684in}}%
\pgfpathlineto{\pgfqpoint{1.423810in}{1.941926in}}%
\pgfpathlineto{\pgfqpoint{1.424413in}{1.532016in}}%
\pgfpathlineto{\pgfqpoint{1.425656in}{1.532051in}}%
\pgfpathlineto{\pgfqpoint{1.426394in}{1.532609in}}%
\pgfpathlineto{\pgfqpoint{1.426496in}{2.037850in}}%
\pgfpathlineto{\pgfqpoint{1.427416in}{1.949223in}}%
\pgfpathlineto{\pgfqpoint{1.427424in}{1.916359in}}%
\pgfpathlineto{\pgfqpoint{1.427441in}{1.963979in}}%
\pgfpathlineto{\pgfqpoint{1.428185in}{1.943951in}}%
\pgfpathlineto{\pgfqpoint{1.428190in}{1.947740in}}%
\pgfpathlineto{\pgfqpoint{1.428884in}{1.912038in}}%
\pgfpathlineto{\pgfqpoint{1.429358in}{1.532016in}}%
\pgfpathlineto{\pgfqpoint{1.429958in}{1.532053in}}%
\pgfpathlineto{\pgfqpoint{1.431330in}{1.532105in}}%
\pgfpathlineto{\pgfqpoint{1.431357in}{1.571758in}}%
\pgfpathlineto{\pgfqpoint{1.431441in}{2.037592in}}%
\pgfpathlineto{\pgfqpoint{1.432163in}{1.907171in}}%
\pgfpathlineto{\pgfqpoint{1.432179in}{1.972642in}}%
\pgfpathlineto{\pgfqpoint{1.432954in}{1.949120in}}%
\pgfpathlineto{\pgfqpoint{1.432968in}{1.931017in}}%
\pgfpathlineto{\pgfqpoint{1.432981in}{1.949710in}}%
\pgfpathlineto{\pgfqpoint{1.433725in}{1.940872in}}%
\pgfpathlineto{\pgfqpoint{1.433745in}{1.937422in}}%
\pgfpathlineto{\pgfqpoint{1.433778in}{1.938980in}}%
\pgfpathlineto{\pgfqpoint{1.433982in}{1.692501in}}%
\pgfpathlineto{\pgfqpoint{1.434354in}{1.532029in}}%
\pgfpathlineto{\pgfqpoint{1.434282in}{1.708399in}}%
\pgfpathlineto{\pgfqpoint{1.434889in}{1.532049in}}%
\pgfpathlineto{\pgfqpoint{1.436280in}{1.532145in}}%
\pgfpathlineto{\pgfqpoint{1.436346in}{2.042032in}}%
\pgfpathlineto{\pgfqpoint{1.437459in}{1.924486in}}%
\pgfpathlineto{\pgfqpoint{1.437463in}{1.921248in}}%
\pgfpathlineto{\pgfqpoint{1.437476in}{1.958800in}}%
\pgfpathlineto{\pgfqpoint{1.438215in}{1.937222in}}%
\pgfpathlineto{\pgfqpoint{1.438225in}{1.945664in}}%
\pgfpathlineto{\pgfqpoint{1.438843in}{1.820441in}}%
\pgfpathlineto{\pgfqpoint{1.439249in}{1.532017in}}%
\pgfpathlineto{\pgfqpoint{1.439750in}{1.532051in}}%
\pgfpathlineto{\pgfqpoint{1.441221in}{1.532106in}}%
\pgfpathlineto{\pgfqpoint{1.441248in}{1.572197in}}%
\pgfpathlineto{\pgfqpoint{1.441332in}{2.037757in}}%
\pgfpathlineto{\pgfqpoint{1.442050in}{1.920604in}}%
\pgfpathlineto{\pgfqpoint{1.442054in}{1.906922in}}%
\pgfpathlineto{\pgfqpoint{1.442070in}{1.972542in}}%
\pgfpathlineto{\pgfqpoint{1.442812in}{1.950261in}}%
\pgfpathlineto{\pgfqpoint{1.442827in}{1.930196in}}%
\pgfpathlineto{\pgfqpoint{1.443604in}{1.937190in}}%
\pgfpathlineto{\pgfqpoint{1.443625in}{1.940578in}}%
\pgfpathlineto{\pgfqpoint{1.443820in}{1.776611in}}%
\pgfpathlineto{\pgfqpoint{1.444239in}{1.532030in}}%
\pgfpathlineto{\pgfqpoint{1.444722in}{1.532040in}}%
\pgfpathlineto{\pgfqpoint{1.446169in}{1.532130in}}%
\pgfpathlineto{\pgfqpoint{1.446234in}{2.041251in}}%
\pgfpathlineto{\pgfqpoint{1.447435in}{1.930492in}}%
\pgfpathlineto{\pgfqpoint{1.447439in}{1.923220in}}%
\pgfpathlineto{\pgfqpoint{1.447455in}{1.956544in}}%
\pgfpathlineto{\pgfqpoint{1.448195in}{1.941892in}}%
\pgfpathlineto{\pgfqpoint{1.448201in}{1.944509in}}%
\pgfpathlineto{\pgfqpoint{1.448716in}{1.850601in}}%
\pgfpathlineto{\pgfqpoint{1.449140in}{1.532016in}}%
\pgfpathlineto{\pgfqpoint{1.449641in}{1.532051in}}%
\pgfpathlineto{\pgfqpoint{1.451112in}{1.532105in}}%
\pgfpathlineto{\pgfqpoint{1.451138in}{1.571898in}}%
\pgfpathlineto{\pgfqpoint{1.451223in}{2.037356in}}%
\pgfpathlineto{\pgfqpoint{1.451941in}{1.920521in}}%
\pgfpathlineto{\pgfqpoint{1.451945in}{1.906845in}}%
\pgfpathlineto{\pgfqpoint{1.451961in}{1.972328in}}%
\pgfpathlineto{\pgfqpoint{1.452703in}{1.950356in}}%
\pgfpathlineto{\pgfqpoint{1.452720in}{1.929900in}}%
\pgfpathlineto{\pgfqpoint{1.453486in}{1.937655in}}%
\pgfpathlineto{\pgfqpoint{1.453505in}{1.940710in}}%
\pgfpathlineto{\pgfqpoint{1.453593in}{1.928610in}}%
\pgfpathlineto{\pgfqpoint{1.454126in}{1.532030in}}%
\pgfpathlineto{\pgfqpoint{1.454801in}{1.532045in}}%
\pgfpathlineto{\pgfqpoint{1.456059in}{1.532120in}}%
\pgfpathlineto{\pgfqpoint{1.456102in}{1.783603in}}%
\pgfpathlineto{\pgfqpoint{1.456159in}{2.006189in}}%
\pgfpathlineto{\pgfqpoint{1.456880in}{1.933887in}}%
\pgfpathlineto{\pgfqpoint{1.456888in}{1.918075in}}%
\pgfpathlineto{\pgfqpoint{1.456900in}{1.962182in}}%
\pgfpathlineto{\pgfqpoint{1.457651in}{1.933462in}}%
\pgfpathlineto{\pgfqpoint{1.457739in}{1.946729in}}%
\pgfpathlineto{\pgfqpoint{1.457688in}{1.932379in}}%
\pgfpathlineto{\pgfqpoint{1.458431in}{1.942949in}}%
\pgfpathlineto{\pgfqpoint{1.458629in}{1.823550in}}%
\pgfpathlineto{\pgfqpoint{1.459073in}{1.532030in}}%
\pgfpathlineto{\pgfqpoint{1.459553in}{1.532040in}}%
\pgfpathlineto{\pgfqpoint{1.461005in}{1.532123in}}%
\pgfpathlineto{\pgfqpoint{1.461048in}{1.812273in}}%
\pgfpathlineto{\pgfqpoint{1.461108in}{2.034358in}}%
\pgfpathlineto{\pgfqpoint{1.461823in}{1.941744in}}%
\pgfpathlineto{\pgfqpoint{1.461832in}{1.908021in}}%
\pgfpathlineto{\pgfqpoint{1.461844in}{1.970079in}}%
\pgfpathlineto{\pgfqpoint{1.462593in}{1.945832in}}%
\pgfpathlineto{\pgfqpoint{1.462600in}{1.949923in}}%
\pgfpathlineto{\pgfqpoint{1.462615in}{1.930223in}}%
\pgfpathlineto{\pgfqpoint{1.463328in}{1.942546in}}%
\pgfpathlineto{\pgfqpoint{1.463568in}{1.823797in}}%
\pgfpathlineto{\pgfqpoint{1.463977in}{1.532015in}}%
\pgfpathlineto{\pgfqpoint{1.464528in}{1.532035in}}%
\pgfpathlineto{\pgfqpoint{1.465957in}{1.532472in}}%
\pgfpathlineto{\pgfqpoint{1.466059in}{2.035135in}}%
\pgfpathlineto{\pgfqpoint{1.466983in}{1.931002in}}%
\pgfpathlineto{\pgfqpoint{1.467019in}{1.917416in}}%
\pgfpathlineto{\pgfqpoint{1.467004in}{1.963045in}}%
\pgfpathlineto{\pgfqpoint{1.467748in}{1.946546in}}%
\pgfpathlineto{\pgfqpoint{1.468922in}{1.532016in}}%
\pgfpathlineto{\pgfqpoint{1.469522in}{1.532053in}}%
\pgfpathlineto{\pgfqpoint{1.470893in}{1.532105in}}%
\pgfpathlineto{\pgfqpoint{1.470920in}{1.571811in}}%
\pgfpathlineto{\pgfqpoint{1.471005in}{2.036960in}}%
\pgfpathlineto{\pgfqpoint{1.471723in}{1.920276in}}%
\pgfpathlineto{\pgfqpoint{1.471727in}{1.906577in}}%
\pgfpathlineto{\pgfqpoint{1.471743in}{1.971999in}}%
\pgfpathlineto{\pgfqpoint{1.472484in}{1.949709in}}%
\pgfpathlineto{\pgfqpoint{1.472499in}{1.929979in}}%
\pgfpathlineto{\pgfqpoint{1.473279in}{1.936539in}}%
\pgfpathlineto{\pgfqpoint{1.473300in}{1.940400in}}%
\pgfpathlineto{\pgfqpoint{1.473484in}{1.798363in}}%
\pgfpathlineto{\pgfqpoint{1.473894in}{1.532031in}}%
\pgfpathlineto{\pgfqpoint{1.474406in}{1.532046in}}%
\pgfpathlineto{\pgfqpoint{1.475840in}{1.532112in}}%
\pgfpathlineto{\pgfqpoint{1.475882in}{1.730597in}}%
\pgfpathlineto{\pgfqpoint{1.475944in}{1.992033in}}%
\pgfpathlineto{\pgfqpoint{1.476667in}{1.958862in}}%
\pgfpathlineto{\pgfqpoint{1.476683in}{1.920724in}}%
\pgfpathlineto{\pgfqpoint{1.477442in}{1.936863in}}%
\pgfpathlineto{\pgfqpoint{1.477668in}{1.946511in}}%
\pgfpathlineto{\pgfqpoint{1.478163in}{1.932298in}}%
\pgfpathlineto{\pgfqpoint{1.478212in}{1.944586in}}%
\pgfpathlineto{\pgfqpoint{1.478410in}{1.826685in}}%
\pgfpathlineto{\pgfqpoint{1.478850in}{1.532031in}}%
\pgfpathlineto{\pgfqpoint{1.479374in}{1.532046in}}%
\pgfpathlineto{\pgfqpoint{1.480785in}{1.532110in}}%
\pgfpathlineto{\pgfqpoint{1.480827in}{1.725510in}}%
\pgfpathlineto{\pgfqpoint{1.480857in}{1.999696in}}%
\pgfpathlineto{\pgfqpoint{1.481611in}{1.961249in}}%
\pgfpathlineto{\pgfqpoint{1.481626in}{1.918212in}}%
\pgfpathlineto{\pgfqpoint{1.482379in}{1.939516in}}%
\pgfpathlineto{\pgfqpoint{1.482389in}{1.947815in}}%
\pgfpathlineto{\pgfqpoint{1.482408in}{1.931690in}}%
\pgfpathlineto{\pgfqpoint{1.483129in}{1.934331in}}%
\pgfpathlineto{\pgfqpoint{1.483179in}{1.944099in}}%
\pgfpathlineto{\pgfqpoint{1.483327in}{1.851056in}}%
\pgfpathlineto{\pgfqpoint{1.483759in}{1.532016in}}%
\pgfpathlineto{\pgfqpoint{1.484309in}{1.532036in}}%
\pgfpathlineto{\pgfqpoint{1.485739in}{1.532560in}}%
\pgfpathlineto{\pgfqpoint{1.485841in}{2.036363in}}%
\pgfpathlineto{\pgfqpoint{1.486762in}{1.948318in}}%
\pgfpathlineto{\pgfqpoint{1.486770in}{1.915565in}}%
\pgfpathlineto{\pgfqpoint{1.486786in}{1.962934in}}%
\pgfpathlineto{\pgfqpoint{1.487530in}{1.942864in}}%
\pgfpathlineto{\pgfqpoint{1.487536in}{1.946763in}}%
\pgfpathlineto{\pgfqpoint{1.488220in}{1.923827in}}%
\pgfpathlineto{\pgfqpoint{1.488731in}{1.532031in}}%
\pgfpathlineto{\pgfqpoint{1.489393in}{1.532045in}}%
\pgfpathlineto{\pgfqpoint{1.490676in}{1.532113in}}%
\pgfpathlineto{\pgfqpoint{1.490719in}{1.734804in}}%
\pgfpathlineto{\pgfqpoint{1.490736in}{1.994878in}}%
\pgfpathlineto{\pgfqpoint{1.491505in}{1.957285in}}%
\pgfpathlineto{\pgfqpoint{1.491520in}{1.921789in}}%
\pgfpathlineto{\pgfqpoint{1.492278in}{1.946603in}}%
\pgfpathlineto{\pgfqpoint{1.492648in}{1.932428in}}%
\pgfpathlineto{\pgfqpoint{1.492387in}{1.946642in}}%
\pgfpathlineto{\pgfqpoint{1.493044in}{1.944495in}}%
\pgfpathlineto{\pgfqpoint{1.493341in}{1.695298in}}%
\pgfpathlineto{\pgfqpoint{1.493694in}{1.532030in}}%
\pgfpathlineto{\pgfqpoint{1.494225in}{1.532048in}}%
\pgfpathlineto{\pgfqpoint{1.495625in}{1.532149in}}%
\pgfpathlineto{\pgfqpoint{1.495693in}{2.026063in}}%
\pgfpathlineto{\pgfqpoint{1.496804in}{1.936188in}}%
\pgfpathlineto{\pgfqpoint{1.496811in}{1.922848in}}%
\pgfpathlineto{\pgfqpoint{1.496823in}{1.955439in}}%
\pgfpathlineto{\pgfqpoint{1.497572in}{1.939805in}}%
\pgfpathlineto{\pgfqpoint{1.498638in}{1.532030in}}%
\pgfpathlineto{\pgfqpoint{1.497595in}{1.943847in}}%
\pgfpathlineto{\pgfqpoint{1.498871in}{1.532041in}}%
\pgfpathlineto{\pgfqpoint{1.500569in}{1.532125in}}%
\pgfpathlineto{\pgfqpoint{1.500634in}{2.037920in}}%
\pgfpathlineto{\pgfqpoint{1.501866in}{1.933783in}}%
\pgfpathlineto{\pgfqpoint{1.501902in}{1.924249in}}%
\pgfpathlineto{\pgfqpoint{1.501886in}{1.954760in}}%
\pgfpathlineto{\pgfqpoint{1.502633in}{1.940944in}}%
\pgfpathlineto{\pgfqpoint{1.503572in}{1.532031in}}%
\pgfpathlineto{\pgfqpoint{1.502659in}{1.942761in}}%
\pgfpathlineto{\pgfqpoint{1.503894in}{1.532042in}}%
\pgfpathlineto{\pgfqpoint{1.505511in}{1.532101in}}%
\pgfpathlineto{\pgfqpoint{1.505520in}{1.532311in}}%
\pgfpathlineto{\pgfqpoint{1.505579in}{2.029052in}}%
\pgfpathlineto{\pgfqpoint{1.506576in}{1.936650in}}%
\pgfpathlineto{\pgfqpoint{1.506583in}{1.917395in}}%
\pgfpathlineto{\pgfqpoint{1.506596in}{1.960525in}}%
\pgfpathlineto{\pgfqpoint{1.507340in}{1.942209in}}%
\pgfpathlineto{\pgfqpoint{1.507346in}{1.945772in}}%
\pgfpathlineto{\pgfqpoint{1.508026in}{1.905235in}}%
\pgfpathlineto{\pgfqpoint{1.508500in}{1.532026in}}%
\pgfpathlineto{\pgfqpoint{1.509079in}{1.532051in}}%
\pgfpathlineto{\pgfqpoint{1.510450in}{1.532084in}}%
\pgfpathlineto{\pgfqpoint{1.510462in}{1.532155in}}%
\pgfpathlineto{\pgfqpoint{1.510518in}{2.005206in}}%
\pgfpathlineto{\pgfqpoint{1.511597in}{1.947982in}}%
\pgfpathlineto{\pgfqpoint{1.511608in}{1.924110in}}%
\pgfpathlineto{\pgfqpoint{1.511624in}{1.953764in}}%
\pgfpathlineto{\pgfqpoint{1.512350in}{1.942368in}}%
\pgfpathlineto{\pgfqpoint{1.512566in}{1.946033in}}%
\pgfpathlineto{\pgfqpoint{1.513011in}{1.858966in}}%
\pgfpathlineto{\pgfqpoint{1.513430in}{1.532026in}}%
\pgfpathlineto{\pgfqpoint{1.513975in}{1.532037in}}%
\pgfpathlineto{\pgfqpoint{1.515407in}{1.532145in}}%
\pgfpathlineto{\pgfqpoint{1.515508in}{1.997662in}}%
\pgfpathlineto{\pgfqpoint{1.516570in}{1.947554in}}%
\pgfpathlineto{\pgfqpoint{1.516581in}{1.925906in}}%
\pgfpathlineto{\pgfqpoint{1.516597in}{1.951998in}}%
\pgfpathlineto{\pgfqpoint{1.517334in}{1.940457in}}%
\pgfpathlineto{\pgfqpoint{1.517514in}{1.944711in}}%
\pgfpathlineto{\pgfqpoint{1.518107in}{1.642101in}}%
\pgfpathlineto{\pgfqpoint{1.518400in}{1.532031in}}%
\pgfpathlineto{\pgfqpoint{1.518355in}{1.698884in}}%
\pgfpathlineto{\pgfqpoint{1.519006in}{1.532048in}}%
\pgfpathlineto{\pgfqpoint{1.520347in}{1.532101in}}%
\pgfpathlineto{\pgfqpoint{1.520356in}{1.532311in}}%
\pgfpathlineto{\pgfqpoint{1.520415in}{2.028864in}}%
\pgfpathlineto{\pgfqpoint{1.521412in}{1.935819in}}%
\pgfpathlineto{\pgfqpoint{1.521420in}{1.917316in}}%
\pgfpathlineto{\pgfqpoint{1.521433in}{1.960613in}}%
\pgfpathlineto{\pgfqpoint{1.522178in}{1.942209in}}%
\pgfpathlineto{\pgfqpoint{1.522184in}{1.945604in}}%
\pgfpathlineto{\pgfqpoint{1.522873in}{1.896522in}}%
\pgfpathlineto{\pgfqpoint{1.523336in}{1.532027in}}%
\pgfpathlineto{\pgfqpoint{1.523914in}{1.532050in}}%
\pgfpathlineto{\pgfqpoint{1.525287in}{1.532085in}}%
\pgfpathlineto{\pgfqpoint{1.525298in}{1.532148in}}%
\pgfpathlineto{\pgfqpoint{1.525366in}{2.018477in}}%
\pgfpathlineto{\pgfqpoint{1.526471in}{1.953466in}}%
\pgfpathlineto{\pgfqpoint{1.526484in}{1.923120in}}%
\pgfpathlineto{\pgfqpoint{1.526497in}{1.954191in}}%
\pgfpathlineto{\pgfqpoint{1.527240in}{1.937130in}}%
\pgfpathlineto{\pgfqpoint{1.527288in}{1.943831in}}%
\pgfpathlineto{\pgfqpoint{1.527827in}{1.879982in}}%
\pgfpathlineto{\pgfqpoint{1.528292in}{1.532031in}}%
\pgfpathlineto{\pgfqpoint{1.528849in}{1.532042in}}%
\pgfpathlineto{\pgfqpoint{1.530247in}{1.532310in}}%
\pgfpathlineto{\pgfqpoint{1.530306in}{2.028600in}}%
\pgfpathlineto{\pgfqpoint{1.531303in}{1.936363in}}%
\pgfpathlineto{\pgfqpoint{1.531311in}{1.917108in}}%
\pgfpathlineto{\pgfqpoint{1.531324in}{1.960154in}}%
\pgfpathlineto{\pgfqpoint{1.532067in}{1.941886in}}%
\pgfpathlineto{\pgfqpoint{1.532073in}{1.945399in}}%
\pgfpathlineto{\pgfqpoint{1.532734in}{1.918643in}}%
\pgfpathlineto{\pgfqpoint{1.533242in}{1.532031in}}%
\pgfpathlineto{\pgfqpoint{1.533857in}{1.532049in}}%
\pgfpathlineto{\pgfqpoint{1.535186in}{1.532115in}}%
\pgfpathlineto{\pgfqpoint{1.535228in}{1.746358in}}%
\pgfpathlineto{\pgfqpoint{1.535289in}{1.995273in}}%
\pgfpathlineto{\pgfqpoint{1.536012in}{1.950047in}}%
\pgfpathlineto{\pgfqpoint{1.536022in}{1.918338in}}%
\pgfpathlineto{\pgfqpoint{1.536038in}{1.958837in}}%
\pgfpathlineto{\pgfqpoint{1.536784in}{1.938256in}}%
\pgfpathlineto{\pgfqpoint{1.536792in}{1.931024in}}%
\pgfpathlineto{\pgfqpoint{1.536810in}{1.946487in}}%
\pgfpathlineto{\pgfqpoint{1.537535in}{1.939862in}}%
\pgfpathlineto{\pgfqpoint{1.537683in}{1.914391in}}%
\pgfpathlineto{\pgfqpoint{1.538199in}{1.532030in}}%
\pgfpathlineto{\pgfqpoint{1.538876in}{1.532044in}}%
\pgfpathlineto{\pgfqpoint{1.540133in}{1.532133in}}%
\pgfpathlineto{\pgfqpoint{1.540199in}{2.042004in}}%
\pgfpathlineto{\pgfqpoint{1.541370in}{1.932552in}}%
\pgfpathlineto{\pgfqpoint{1.541405in}{1.921778in}}%
\pgfpathlineto{\pgfqpoint{1.541390in}{1.956291in}}%
\pgfpathlineto{\pgfqpoint{1.542135in}{1.943297in}}%
\pgfpathlineto{\pgfqpoint{1.543148in}{1.532030in}}%
\pgfpathlineto{\pgfqpoint{1.543532in}{1.532038in}}%
\pgfpathlineto{\pgfqpoint{1.545078in}{1.532129in}}%
\pgfpathlineto{\pgfqpoint{1.545143in}{2.040177in}}%
\pgfpathlineto{\pgfqpoint{1.546344in}{1.933946in}}%
\pgfpathlineto{\pgfqpoint{1.546351in}{1.922524in}}%
\pgfpathlineto{\pgfqpoint{1.546364in}{1.955013in}}%
\pgfpathlineto{\pgfqpoint{1.547109in}{1.941919in}}%
\pgfpathlineto{\pgfqpoint{1.548090in}{1.532030in}}%
\pgfpathlineto{\pgfqpoint{1.547135in}{1.942641in}}%
\pgfpathlineto{\pgfqpoint{1.548371in}{1.532036in}}%
\pgfpathlineto{\pgfqpoint{1.550024in}{1.532135in}}%
\pgfpathlineto{\pgfqpoint{1.550090in}{2.043054in}}%
\pgfpathlineto{\pgfqpoint{1.551257in}{1.943052in}}%
\pgfpathlineto{\pgfqpoint{1.551265in}{1.921014in}}%
\pgfpathlineto{\pgfqpoint{1.551282in}{1.956169in}}%
\pgfpathlineto{\pgfqpoint{1.552029in}{1.940214in}}%
\pgfpathlineto{\pgfqpoint{1.553043in}{1.532029in}}%
\pgfpathlineto{\pgfqpoint{1.552052in}{1.943081in}}%
\pgfpathlineto{\pgfqpoint{1.553182in}{1.532040in}}%
\pgfpathlineto{\pgfqpoint{1.554971in}{1.532160in}}%
\pgfpathlineto{\pgfqpoint{1.555075in}{1.996179in}}%
\pgfpathlineto{\pgfqpoint{1.556108in}{1.928823in}}%
\pgfpathlineto{\pgfqpoint{1.556114in}{1.925603in}}%
\pgfpathlineto{\pgfqpoint{1.556129in}{1.951090in}}%
\pgfpathlineto{\pgfqpoint{1.556857in}{1.932583in}}%
\pgfpathlineto{\pgfqpoint{1.557129in}{1.943639in}}%
\pgfpathlineto{\pgfqpoint{1.557476in}{1.901602in}}%
\pgfpathlineto{\pgfqpoint{1.557994in}{1.532028in}}%
\pgfpathlineto{\pgfqpoint{1.558629in}{1.532051in}}%
\pgfpathlineto{\pgfqpoint{1.559916in}{1.532142in}}%
\pgfpathlineto{\pgfqpoint{1.559982in}{2.041890in}}%
\pgfpathlineto{\pgfqpoint{1.561092in}{1.936423in}}%
\pgfpathlineto{\pgfqpoint{1.561100in}{1.919200in}}%
\pgfpathlineto{\pgfqpoint{1.561112in}{1.957297in}}%
\pgfpathlineto{\pgfqpoint{1.561858in}{1.941061in}}%
\pgfpathlineto{\pgfqpoint{1.561864in}{1.943981in}}%
\pgfpathlineto{\pgfqpoint{1.562431in}{1.894494in}}%
\pgfpathlineto{\pgfqpoint{1.562928in}{1.532030in}}%
\pgfpathlineto{\pgfqpoint{1.563458in}{1.532047in}}%
\pgfpathlineto{\pgfqpoint{1.564860in}{1.532125in}}%
\pgfpathlineto{\pgfqpoint{1.564924in}{2.037048in}}%
\pgfpathlineto{\pgfqpoint{1.566159in}{1.926643in}}%
\pgfpathlineto{\pgfqpoint{1.566163in}{1.922865in}}%
\pgfpathlineto{\pgfqpoint{1.566207in}{1.953064in}}%
\pgfpathlineto{\pgfqpoint{1.566913in}{1.939222in}}%
\pgfpathlineto{\pgfqpoint{1.566918in}{1.942288in}}%
\pgfpathlineto{\pgfqpoint{1.567415in}{1.867553in}}%
\pgfpathlineto{\pgfqpoint{1.567862in}{1.532033in}}%
\pgfpathlineto{\pgfqpoint{1.568380in}{1.532044in}}%
\pgfpathlineto{\pgfqpoint{1.569808in}{1.532164in}}%
\pgfpathlineto{\pgfqpoint{1.569871in}{2.035828in}}%
\pgfpathlineto{\pgfqpoint{1.570981in}{1.939636in}}%
\pgfpathlineto{\pgfqpoint{1.570989in}{1.919577in}}%
\pgfpathlineto{\pgfqpoint{1.571004in}{1.956728in}}%
\pgfpathlineto{\pgfqpoint{1.571755in}{1.933332in}}%
\pgfpathlineto{\pgfqpoint{1.571772in}{1.943123in}}%
\pgfpathlineto{\pgfqpoint{1.572302in}{1.925172in}}%
\pgfpathlineto{\pgfqpoint{1.572836in}{1.532033in}}%
\pgfpathlineto{\pgfqpoint{1.573521in}{1.532048in}}%
\pgfpathlineto{\pgfqpoint{1.574752in}{1.532137in}}%
\pgfpathlineto{\pgfqpoint{1.574814in}{2.033557in}}%
\pgfpathlineto{\pgfqpoint{1.576015in}{1.943893in}}%
\pgfpathlineto{\pgfqpoint{1.576023in}{1.922473in}}%
\pgfpathlineto{\pgfqpoint{1.576038in}{1.954142in}}%
\pgfpathlineto{\pgfqpoint{1.576787in}{1.935083in}}%
\pgfpathlineto{\pgfqpoint{1.577774in}{1.532033in}}%
\pgfpathlineto{\pgfqpoint{1.576807in}{1.942003in}}%
\pgfpathlineto{\pgfqpoint{1.578854in}{1.532056in}}%
\pgfpathlineto{\pgfqpoint{1.579698in}{1.532153in}}%
\pgfpathlineto{\pgfqpoint{1.579761in}{2.033117in}}%
\pgfpathlineto{\pgfqpoint{1.580903in}{1.931704in}}%
\pgfpathlineto{\pgfqpoint{1.580907in}{1.921768in}}%
\pgfpathlineto{\pgfqpoint{1.580924in}{1.955377in}}%
\pgfpathlineto{\pgfqpoint{1.581674in}{1.933502in}}%
\pgfpathlineto{\pgfqpoint{1.581689in}{1.942571in}}%
\pgfpathlineto{\pgfqpoint{1.582239in}{1.883013in}}%
\pgfpathlineto{\pgfqpoint{1.582756in}{1.532033in}}%
\pgfpathlineto{\pgfqpoint{1.583246in}{1.532044in}}%
\pgfpathlineto{\pgfqpoint{1.584645in}{1.532192in}}%
\pgfpathlineto{\pgfqpoint{1.584708in}{2.039147in}}%
\pgfpathlineto{\pgfqpoint{1.585763in}{1.922541in}}%
\pgfpathlineto{\pgfqpoint{1.585767in}{1.916976in}}%
\pgfpathlineto{\pgfqpoint{1.585782in}{1.958321in}}%
\pgfpathlineto{\pgfqpoint{1.586515in}{1.942841in}}%
\pgfpathlineto{\pgfqpoint{1.586519in}{1.944138in}}%
\pgfpathlineto{\pgfqpoint{1.587162in}{1.906987in}}%
\pgfpathlineto{\pgfqpoint{1.587659in}{1.532033in}}%
\pgfpathlineto{\pgfqpoint{1.588242in}{1.532046in}}%
\pgfpathlineto{\pgfqpoint{1.589591in}{1.532192in}}%
\pgfpathlineto{\pgfqpoint{1.589653in}{2.039095in}}%
\pgfpathlineto{\pgfqpoint{1.590706in}{1.929557in}}%
\pgfpathlineto{\pgfqpoint{1.590710in}{1.917953in}}%
\pgfpathlineto{\pgfqpoint{1.590727in}{1.959014in}}%
\pgfpathlineto{\pgfqpoint{1.591476in}{1.933302in}}%
\pgfpathlineto{\pgfqpoint{1.592652in}{1.532033in}}%
\pgfpathlineto{\pgfqpoint{1.591493in}{1.943615in}}%
\pgfpathlineto{\pgfqpoint{1.593689in}{1.532057in}}%
\pgfpathlineto{\pgfqpoint{1.594536in}{1.532185in}}%
\pgfpathlineto{\pgfqpoint{1.594599in}{2.038887in}}%
\pgfpathlineto{\pgfqpoint{1.595679in}{1.938976in}}%
\pgfpathlineto{\pgfqpoint{1.595686in}{1.917793in}}%
\pgfpathlineto{\pgfqpoint{1.595702in}{1.957671in}}%
\pgfpathlineto{\pgfqpoint{1.596450in}{1.933339in}}%
\pgfpathlineto{\pgfqpoint{1.596466in}{1.943094in}}%
\pgfpathlineto{\pgfqpoint{1.597205in}{1.689001in}}%
\pgfpathlineto{\pgfqpoint{1.597535in}{1.532033in}}%
\pgfpathlineto{\pgfqpoint{1.598102in}{1.532045in}}%
\pgfpathlineto{\pgfqpoint{1.599481in}{1.532165in}}%
\pgfpathlineto{\pgfqpoint{1.599544in}{2.036021in}}%
\pgfpathlineto{\pgfqpoint{1.600655in}{1.935687in}}%
\pgfpathlineto{\pgfqpoint{1.600662in}{1.919301in}}%
\pgfpathlineto{\pgfqpoint{1.600705in}{1.955812in}}%
\pgfpathlineto{\pgfqpoint{1.601425in}{1.934130in}}%
\pgfpathlineto{\pgfqpoint{1.602528in}{1.532032in}}%
\pgfpathlineto{\pgfqpoint{1.601473in}{1.942541in}}%
\pgfpathlineto{\pgfqpoint{1.603361in}{1.532053in}}%
\pgfpathlineto{\pgfqpoint{1.604423in}{1.532121in}}%
\pgfpathlineto{\pgfqpoint{1.604487in}{2.033419in}}%
\pgfpathlineto{\pgfqpoint{1.605802in}{1.946108in}}%
\pgfpathlineto{\pgfqpoint{1.605813in}{1.925099in}}%
\pgfpathlineto{\pgfqpoint{1.605825in}{1.950523in}}%
\pgfpathlineto{\pgfqpoint{1.606576in}{1.934738in}}%
\pgfpathlineto{\pgfqpoint{1.606593in}{1.940725in}}%
\pgfpathlineto{\pgfqpoint{1.606918in}{1.927985in}}%
\pgfpathlineto{\pgfqpoint{1.607478in}{1.532032in}}%
\pgfpathlineto{\pgfqpoint{1.608163in}{1.532048in}}%
\pgfpathlineto{\pgfqpoint{1.609368in}{1.532118in}}%
\pgfpathlineto{\pgfqpoint{1.609411in}{1.771156in}}%
\pgfpathlineto{\pgfqpoint{1.609467in}{1.994706in}}%
\pgfpathlineto{\pgfqpoint{1.610191in}{1.933631in}}%
\pgfpathlineto{\pgfqpoint{1.610196in}{1.917447in}}%
\pgfpathlineto{\pgfqpoint{1.610211in}{1.957487in}}%
\pgfpathlineto{\pgfqpoint{1.610954in}{1.944128in}}%
\pgfpathlineto{\pgfqpoint{1.610989in}{1.945154in}}%
\pgfpathlineto{\pgfqpoint{1.611926in}{1.867228in}}%
\pgfpathlineto{\pgfqpoint{1.612387in}{1.532032in}}%
\pgfpathlineto{\pgfqpoint{1.612921in}{1.532045in}}%
\pgfpathlineto{\pgfqpoint{1.614315in}{1.532129in}}%
\pgfpathlineto{\pgfqpoint{1.614379in}{2.032493in}}%
\pgfpathlineto{\pgfqpoint{1.615635in}{1.947245in}}%
\pgfpathlineto{\pgfqpoint{1.615646in}{1.923686in}}%
\pgfpathlineto{\pgfqpoint{1.615659in}{1.951175in}}%
\pgfpathlineto{\pgfqpoint{1.616409in}{1.935610in}}%
\pgfpathlineto{\pgfqpoint{1.617339in}{1.532028in}}%
\pgfpathlineto{\pgfqpoint{1.616428in}{1.941131in}}%
\pgfpathlineto{\pgfqpoint{1.618172in}{1.532057in}}%
\pgfpathlineto{\pgfqpoint{1.619251in}{1.532086in}}%
\pgfpathlineto{\pgfqpoint{1.619262in}{1.532147in}}%
\pgfpathlineto{\pgfqpoint{1.619328in}{2.032541in}}%
\pgfpathlineto{\pgfqpoint{1.620438in}{1.934271in}}%
\pgfpathlineto{\pgfqpoint{1.620445in}{1.920097in}}%
\pgfpathlineto{\pgfqpoint{1.620461in}{1.954510in}}%
\pgfpathlineto{\pgfqpoint{1.621203in}{1.933140in}}%
\pgfpathlineto{\pgfqpoint{1.621214in}{1.942837in}}%
\pgfpathlineto{\pgfqpoint{1.621823in}{1.861866in}}%
\pgfpathlineto{\pgfqpoint{1.622279in}{1.532028in}}%
\pgfpathlineto{\pgfqpoint{1.622814in}{1.532049in}}%
\pgfpathlineto{\pgfqpoint{1.624196in}{1.532084in}}%
\pgfpathlineto{\pgfqpoint{1.624208in}{1.532162in}}%
\pgfpathlineto{\pgfqpoint{1.624310in}{1.995089in}}%
\pgfpathlineto{\pgfqpoint{1.625344in}{1.934537in}}%
\pgfpathlineto{\pgfqpoint{1.625381in}{1.924624in}}%
\pgfpathlineto{\pgfqpoint{1.625364in}{1.950984in}}%
\pgfpathlineto{\pgfqpoint{1.626101in}{1.941307in}}%
\pgfpathlineto{\pgfqpoint{1.626175in}{1.943808in}}%
\pgfpathlineto{\pgfqpoint{1.626818in}{1.780672in}}%
\pgfpathlineto{\pgfqpoint{1.627218in}{1.532030in}}%
\pgfpathlineto{\pgfqpoint{1.627748in}{1.532047in}}%
\pgfpathlineto{\pgfqpoint{1.629152in}{1.532137in}}%
\pgfpathlineto{\pgfqpoint{1.629217in}{2.043464in}}%
\pgfpathlineto{\pgfqpoint{1.630360in}{1.923443in}}%
\pgfpathlineto{\pgfqpoint{1.630364in}{1.918563in}}%
\pgfpathlineto{\pgfqpoint{1.630409in}{1.955661in}}%
\pgfpathlineto{\pgfqpoint{1.631115in}{1.938723in}}%
\pgfpathlineto{\pgfqpoint{1.631123in}{1.942652in}}%
\pgfpathlineto{\pgfqpoint{1.631719in}{1.840743in}}%
\pgfpathlineto{\pgfqpoint{1.632168in}{1.532029in}}%
\pgfpathlineto{\pgfqpoint{1.632651in}{1.532040in}}%
\pgfpathlineto{\pgfqpoint{1.634098in}{1.532154in}}%
\pgfpathlineto{\pgfqpoint{1.634200in}{2.003120in}}%
\pgfpathlineto{\pgfqpoint{1.635252in}{1.945169in}}%
\pgfpathlineto{\pgfqpoint{1.635262in}{1.922860in}}%
\pgfpathlineto{\pgfqpoint{1.635277in}{1.952316in}}%
\pgfpathlineto{\pgfqpoint{1.636017in}{1.932420in}}%
\pgfpathlineto{\pgfqpoint{1.636034in}{1.942735in}}%
\pgfpathlineto{\pgfqpoint{1.636652in}{1.876493in}}%
\pgfpathlineto{\pgfqpoint{1.637106in}{1.532030in}}%
\pgfpathlineto{\pgfqpoint{1.637634in}{1.532047in}}%
\pgfpathlineto{\pgfqpoint{1.639042in}{1.532126in}}%
\pgfpathlineto{\pgfqpoint{1.639105in}{2.032479in}}%
\pgfpathlineto{\pgfqpoint{1.640369in}{1.927418in}}%
\pgfpathlineto{\pgfqpoint{1.640372in}{1.923252in}}%
\pgfpathlineto{\pgfqpoint{1.640388in}{1.951489in}}%
\pgfpathlineto{\pgfqpoint{1.641126in}{1.941138in}}%
\pgfpathlineto{\pgfqpoint{1.642068in}{1.532028in}}%
\pgfpathlineto{\pgfqpoint{1.642803in}{1.532054in}}%
\pgfpathlineto{\pgfqpoint{1.643979in}{1.532086in}}%
\pgfpathlineto{\pgfqpoint{1.643989in}{1.532147in}}%
\pgfpathlineto{\pgfqpoint{1.644055in}{2.038218in}}%
\pgfpathlineto{\pgfqpoint{1.645165in}{1.934721in}}%
\pgfpathlineto{\pgfqpoint{1.645173in}{1.918933in}}%
\pgfpathlineto{\pgfqpoint{1.645186in}{1.956296in}}%
\pgfpathlineto{\pgfqpoint{1.645931in}{1.940264in}}%
\pgfpathlineto{\pgfqpoint{1.645937in}{1.943055in}}%
\pgfpathlineto{\pgfqpoint{1.646522in}{1.896424in}}%
\pgfpathlineto{\pgfqpoint{1.647011in}{1.532028in}}%
\pgfpathlineto{\pgfqpoint{1.647548in}{1.532049in}}%
\pgfpathlineto{\pgfqpoint{1.648934in}{1.532142in}}%
\pgfpathlineto{\pgfqpoint{1.649000in}{2.040921in}}%
\pgfpathlineto{\pgfqpoint{1.650110in}{1.934859in}}%
\pgfpathlineto{\pgfqpoint{1.650118in}{1.918254in}}%
\pgfpathlineto{\pgfqpoint{1.650131in}{1.956630in}}%
\pgfpathlineto{\pgfqpoint{1.650875in}{1.940126in}}%
\pgfpathlineto{\pgfqpoint{1.650881in}{1.943167in}}%
\pgfpathlineto{\pgfqpoint{1.651457in}{1.906547in}}%
\pgfpathlineto{\pgfqpoint{1.651958in}{1.532028in}}%
\pgfpathlineto{\pgfqpoint{1.652495in}{1.532050in}}%
\pgfpathlineto{\pgfqpoint{1.653880in}{1.532144in}}%
\pgfpathlineto{\pgfqpoint{1.653946in}{2.033174in}}%
\pgfpathlineto{\pgfqpoint{1.655056in}{1.933790in}}%
\pgfpathlineto{\pgfqpoint{1.655064in}{1.919611in}}%
\pgfpathlineto{\pgfqpoint{1.655079in}{1.954332in}}%
\pgfpathlineto{\pgfqpoint{1.655821in}{1.932719in}}%
\pgfpathlineto{\pgfqpoint{1.655865in}{1.942541in}}%
\pgfpathlineto{\pgfqpoint{1.656431in}{1.875394in}}%
\pgfpathlineto{\pgfqpoint{1.656896in}{1.532029in}}%
\pgfpathlineto{\pgfqpoint{1.657431in}{1.532049in}}%
\pgfpathlineto{\pgfqpoint{1.658826in}{1.532157in}}%
\pgfpathlineto{\pgfqpoint{1.658929in}{1.998901in}}%
\pgfpathlineto{\pgfqpoint{1.659982in}{1.950712in}}%
\pgfpathlineto{\pgfqpoint{1.659997in}{1.924633in}}%
\pgfpathlineto{\pgfqpoint{1.660757in}{1.941764in}}%
\pgfpathlineto{\pgfqpoint{1.661836in}{1.532030in}}%
\pgfpathlineto{\pgfqpoint{1.660892in}{1.942356in}}%
\pgfpathlineto{\pgfqpoint{1.662020in}{1.532033in}}%
\pgfpathlineto{\pgfqpoint{1.663770in}{1.532134in}}%
\pgfpathlineto{\pgfqpoint{1.663835in}{2.041480in}}%
\pgfpathlineto{\pgfqpoint{1.665003in}{1.941408in}}%
\pgfpathlineto{\pgfqpoint{1.665010in}{1.919728in}}%
\pgfpathlineto{\pgfqpoint{1.665027in}{1.955102in}}%
\pgfpathlineto{\pgfqpoint{1.665773in}{1.940295in}}%
\pgfpathlineto{\pgfqpoint{1.666796in}{1.532028in}}%
\pgfpathlineto{\pgfqpoint{1.665798in}{1.941899in}}%
\pgfpathlineto{\pgfqpoint{1.666937in}{1.532041in}}%
\pgfpathlineto{\pgfqpoint{1.668706in}{1.532086in}}%
\pgfpathlineto{\pgfqpoint{1.668716in}{1.532147in}}%
\pgfpathlineto{\pgfqpoint{1.668783in}{2.035572in}}%
\pgfpathlineto{\pgfqpoint{1.669892in}{1.934480in}}%
\pgfpathlineto{\pgfqpoint{1.669900in}{1.919121in}}%
\pgfpathlineto{\pgfqpoint{1.669913in}{1.955778in}}%
\pgfpathlineto{\pgfqpoint{1.670660in}{1.940126in}}%
\pgfpathlineto{\pgfqpoint{1.670665in}{1.942761in}}%
\pgfpathlineto{\pgfqpoint{1.671236in}{1.909189in}}%
\pgfpathlineto{\pgfqpoint{1.671729in}{1.532030in}}%
\pgfpathlineto{\pgfqpoint{1.672309in}{1.532042in}}%
\pgfpathlineto{\pgfqpoint{1.673661in}{1.532139in}}%
\pgfpathlineto{\pgfqpoint{1.673727in}{2.044455in}}%
\pgfpathlineto{\pgfqpoint{1.674867in}{1.930638in}}%
\pgfpathlineto{\pgfqpoint{1.674903in}{1.918984in}}%
\pgfpathlineto{\pgfqpoint{1.674887in}{1.956408in}}%
\pgfpathlineto{\pgfqpoint{1.675631in}{1.942663in}}%
\pgfpathlineto{\pgfqpoint{1.676684in}{1.532032in}}%
\pgfpathlineto{\pgfqpoint{1.677567in}{1.532051in}}%
\pgfpathlineto{\pgfqpoint{1.678605in}{1.532118in}}%
\pgfpathlineto{\pgfqpoint{1.678647in}{1.773016in}}%
\pgfpathlineto{\pgfqpoint{1.678704in}{1.997447in}}%
\pgfpathlineto{\pgfqpoint{1.679425in}{1.948419in}}%
\pgfpathlineto{\pgfqpoint{1.679433in}{1.916122in}}%
\pgfpathlineto{\pgfqpoint{1.679451in}{1.957687in}}%
\pgfpathlineto{\pgfqpoint{1.680193in}{1.943598in}}%
\pgfpathlineto{\pgfqpoint{1.680227in}{1.944455in}}%
\pgfpathlineto{\pgfqpoint{1.681162in}{1.878996in}}%
\pgfpathlineto{\pgfqpoint{1.681672in}{1.532033in}}%
\pgfpathlineto{\pgfqpoint{1.682167in}{1.532044in}}%
\pgfpathlineto{\pgfqpoint{1.683551in}{1.532125in}}%
\pgfpathlineto{\pgfqpoint{1.683608in}{2.014988in}}%
\pgfpathlineto{\pgfqpoint{1.685002in}{1.938890in}}%
\pgfpathlineto{\pgfqpoint{1.685012in}{1.926959in}}%
\pgfpathlineto{\pgfqpoint{1.685027in}{1.947347in}}%
\pgfpathlineto{\pgfqpoint{1.685766in}{1.934714in}}%
\pgfpathlineto{\pgfqpoint{1.685783in}{1.940984in}}%
\pgfpathlineto{\pgfqpoint{1.686120in}{1.848371in}}%
\pgfpathlineto{\pgfqpoint{1.686588in}{1.532032in}}%
\pgfpathlineto{\pgfqpoint{1.687075in}{1.532043in}}%
\pgfpathlineto{\pgfqpoint{1.688496in}{1.532119in}}%
\pgfpathlineto{\pgfqpoint{1.688538in}{1.781603in}}%
\pgfpathlineto{\pgfqpoint{1.688595in}{2.002507in}}%
\pgfpathlineto{\pgfqpoint{1.689319in}{1.922726in}}%
\pgfpathlineto{\pgfqpoint{1.689323in}{1.914544in}}%
\pgfpathlineto{\pgfqpoint{1.689340in}{1.957677in}}%
\pgfpathlineto{\pgfqpoint{1.690079in}{1.942255in}}%
\pgfpathlineto{\pgfqpoint{1.690155in}{1.930457in}}%
\pgfpathlineto{\pgfqpoint{1.690107in}{1.944785in}}%
\pgfpathlineto{\pgfqpoint{1.690817in}{1.934725in}}%
\pgfpathlineto{\pgfqpoint{1.690966in}{1.939815in}}%
\pgfpathlineto{\pgfqpoint{1.691015in}{1.916010in}}%
\pgfpathlineto{\pgfqpoint{1.691509in}{1.532030in}}%
\pgfpathlineto{\pgfqpoint{1.692088in}{1.532042in}}%
\pgfpathlineto{\pgfqpoint{1.693442in}{1.532135in}}%
\pgfpathlineto{\pgfqpoint{1.693508in}{2.041912in}}%
\pgfpathlineto{\pgfqpoint{1.694677in}{1.939669in}}%
\pgfpathlineto{\pgfqpoint{1.694684in}{1.919216in}}%
\pgfpathlineto{\pgfqpoint{1.694700in}{1.954842in}}%
\pgfpathlineto{\pgfqpoint{1.695445in}{1.942100in}}%
\pgfpathlineto{\pgfqpoint{1.696479in}{1.532032in}}%
\pgfpathlineto{\pgfqpoint{1.697067in}{1.532045in}}%
\pgfpathlineto{\pgfqpoint{1.698388in}{1.532133in}}%
\pgfpathlineto{\pgfqpoint{1.698451in}{2.033054in}}%
\pgfpathlineto{\pgfqpoint{1.699681in}{1.938311in}}%
\pgfpathlineto{\pgfqpoint{1.699689in}{1.921912in}}%
\pgfpathlineto{\pgfqpoint{1.699704in}{1.951732in}}%
\pgfpathlineto{\pgfqpoint{1.700451in}{1.935391in}}%
\pgfpathlineto{\pgfqpoint{1.701399in}{1.532030in}}%
\pgfpathlineto{\pgfqpoint{1.700471in}{1.940776in}}%
\pgfpathlineto{\pgfqpoint{1.701879in}{1.532040in}}%
\pgfpathlineto{\pgfqpoint{1.703333in}{1.532131in}}%
\pgfpathlineto{\pgfqpoint{1.703398in}{2.039759in}}%
\pgfpathlineto{\pgfqpoint{1.704571in}{1.924311in}}%
\pgfpathlineto{\pgfqpoint{1.704575in}{1.919339in}}%
\pgfpathlineto{\pgfqpoint{1.704591in}{1.953802in}}%
\pgfpathlineto{\pgfqpoint{1.705327in}{1.941332in}}%
\pgfpathlineto{\pgfqpoint{1.705358in}{1.941378in}}%
\pgfpathlineto{\pgfqpoint{1.705827in}{1.929394in}}%
\pgfpathlineto{\pgfqpoint{1.706344in}{1.532030in}}%
\pgfpathlineto{\pgfqpoint{1.707268in}{1.532056in}}%
\pgfpathlineto{\pgfqpoint{1.708278in}{1.532129in}}%
\pgfpathlineto{\pgfqpoint{1.708343in}{2.038631in}}%
\pgfpathlineto{\pgfqpoint{1.709544in}{1.932826in}}%
\pgfpathlineto{\pgfqpoint{1.709552in}{1.920950in}}%
\pgfpathlineto{\pgfqpoint{1.709564in}{1.953110in}}%
\pgfpathlineto{\pgfqpoint{1.710313in}{1.937565in}}%
\pgfpathlineto{\pgfqpoint{1.711289in}{1.532030in}}%
\pgfpathlineto{\pgfqpoint{1.710334in}{1.940934in}}%
\pgfpathlineto{\pgfqpoint{1.711521in}{1.532041in}}%
\pgfpathlineto{\pgfqpoint{1.713224in}{1.532128in}}%
\pgfpathlineto{\pgfqpoint{1.713289in}{2.038135in}}%
\pgfpathlineto{\pgfqpoint{1.714492in}{1.927745in}}%
\pgfpathlineto{\pgfqpoint{1.714496in}{1.920372in}}%
\pgfpathlineto{\pgfqpoint{1.714512in}{1.953637in}}%
\pgfpathlineto{\pgfqpoint{1.715256in}{1.940374in}}%
\pgfpathlineto{\pgfqpoint{1.716233in}{1.532031in}}%
\pgfpathlineto{\pgfqpoint{1.715313in}{1.940873in}}%
\pgfpathlineto{\pgfqpoint{1.716512in}{1.532036in}}%
\pgfpathlineto{\pgfqpoint{1.718169in}{1.532123in}}%
\pgfpathlineto{\pgfqpoint{1.718232in}{2.032367in}}%
\pgfpathlineto{\pgfqpoint{1.719524in}{1.928630in}}%
\pgfpathlineto{\pgfqpoint{1.719528in}{1.923260in}}%
\pgfpathlineto{\pgfqpoint{1.719543in}{1.950586in}}%
\pgfpathlineto{\pgfqpoint{1.720290in}{1.935246in}}%
\pgfpathlineto{\pgfqpoint{1.721176in}{1.532031in}}%
\pgfpathlineto{\pgfqpoint{1.720310in}{1.940010in}}%
\pgfpathlineto{\pgfqpoint{1.721948in}{1.532047in}}%
\pgfpathlineto{\pgfqpoint{1.723113in}{1.532115in}}%
\pgfpathlineto{\pgfqpoint{1.723155in}{1.749220in}}%
\pgfpathlineto{\pgfqpoint{1.723216in}{1.992042in}}%
\pgfpathlineto{\pgfqpoint{1.723939in}{1.948352in}}%
\pgfpathlineto{\pgfqpoint{1.723950in}{1.916608in}}%
\pgfpathlineto{\pgfqpoint{1.723965in}{1.956877in}}%
\pgfpathlineto{\pgfqpoint{1.724714in}{1.933892in}}%
\pgfpathlineto{\pgfqpoint{1.724721in}{1.929386in}}%
\pgfpathlineto{\pgfqpoint{1.724738in}{1.944591in}}%
\pgfpathlineto{\pgfqpoint{1.725482in}{1.933802in}}%
\pgfpathlineto{\pgfqpoint{1.725581in}{1.936935in}}%
\pgfpathlineto{\pgfqpoint{1.726128in}{1.532032in}}%
\pgfpathlineto{\pgfqpoint{1.726808in}{1.532046in}}%
\pgfpathlineto{\pgfqpoint{1.728058in}{1.532113in}}%
\pgfpathlineto{\pgfqpoint{1.728100in}{1.736750in}}%
\pgfpathlineto{\pgfqpoint{1.728117in}{1.996591in}}%
\pgfpathlineto{\pgfqpoint{1.728884in}{1.956482in}}%
\pgfpathlineto{\pgfqpoint{1.728900in}{1.918078in}}%
\pgfpathlineto{\pgfqpoint{1.729660in}{1.933093in}}%
\pgfpathlineto{\pgfqpoint{1.730030in}{1.944112in}}%
\pgfpathlineto{\pgfqpoint{1.730425in}{1.929418in}}%
\pgfpathlineto{\pgfqpoint{1.730475in}{1.942965in}}%
\pgfpathlineto{\pgfqpoint{1.730623in}{1.847564in}}%
\pgfpathlineto{\pgfqpoint{1.731086in}{1.532028in}}%
\pgfpathlineto{\pgfqpoint{1.731623in}{1.532050in}}%
\pgfpathlineto{\pgfqpoint{1.732997in}{1.532086in}}%
\pgfpathlineto{\pgfqpoint{1.733007in}{1.532147in}}%
\pgfpathlineto{\pgfqpoint{1.733073in}{2.037777in}}%
\pgfpathlineto{\pgfqpoint{1.734181in}{1.943849in}}%
\pgfpathlineto{\pgfqpoint{1.734189in}{1.917956in}}%
\pgfpathlineto{\pgfqpoint{1.734205in}{1.955476in}}%
\pgfpathlineto{\pgfqpoint{1.734950in}{1.939568in}}%
\pgfpathlineto{\pgfqpoint{1.734955in}{1.942340in}}%
\pgfpathlineto{\pgfqpoint{1.735524in}{1.909825in}}%
\pgfpathlineto{\pgfqpoint{1.736008in}{1.532032in}}%
\pgfpathlineto{\pgfqpoint{1.736575in}{1.532044in}}%
\pgfpathlineto{\pgfqpoint{1.737950in}{1.532117in}}%
\pgfpathlineto{\pgfqpoint{1.737992in}{1.761959in}}%
\pgfpathlineto{\pgfqpoint{1.738053in}{1.992348in}}%
\pgfpathlineto{\pgfqpoint{1.738771in}{1.951530in}}%
\pgfpathlineto{\pgfqpoint{1.738782in}{1.916070in}}%
\pgfpathlineto{\pgfqpoint{1.738797in}{1.957427in}}%
\pgfpathlineto{\pgfqpoint{1.739538in}{1.929156in}}%
\pgfpathlineto{\pgfqpoint{1.739560in}{1.944451in}}%
\pgfpathlineto{\pgfqpoint{1.740295in}{1.940054in}}%
\pgfpathlineto{\pgfqpoint{1.740542in}{1.839244in}}%
\pgfpathlineto{\pgfqpoint{1.740962in}{1.532032in}}%
\pgfpathlineto{\pgfqpoint{1.741490in}{1.532045in}}%
\pgfpathlineto{\pgfqpoint{1.742897in}{1.532128in}}%
\pgfpathlineto{\pgfqpoint{1.742960in}{2.031710in}}%
\pgfpathlineto{\pgfqpoint{1.744217in}{1.945563in}}%
\pgfpathlineto{\pgfqpoint{1.744229in}{1.922920in}}%
\pgfpathlineto{\pgfqpoint{1.744241in}{1.950400in}}%
\pgfpathlineto{\pgfqpoint{1.744992in}{1.933894in}}%
\pgfpathlineto{\pgfqpoint{1.745011in}{1.939994in}}%
\pgfpathlineto{\pgfqpoint{1.745622in}{1.635933in}}%
\pgfpathlineto{\pgfqpoint{1.745910in}{1.532030in}}%
\pgfpathlineto{\pgfqpoint{1.745846in}{1.676027in}}%
\pgfpathlineto{\pgfqpoint{1.746540in}{1.532050in}}%
\pgfpathlineto{\pgfqpoint{1.747843in}{1.532139in}}%
\pgfpathlineto{\pgfqpoint{1.747908in}{2.043657in}}%
\pgfpathlineto{\pgfqpoint{1.749048in}{1.936674in}}%
\pgfpathlineto{\pgfqpoint{1.749056in}{1.917495in}}%
\pgfpathlineto{\pgfqpoint{1.749100in}{1.954787in}}%
\pgfpathlineto{\pgfqpoint{1.749815in}{1.942101in}}%
\pgfpathlineto{\pgfqpoint{1.750912in}{1.532033in}}%
\pgfpathlineto{\pgfqpoint{1.751308in}{1.532041in}}%
\pgfpathlineto{\pgfqpoint{1.752787in}{1.532120in}}%
\pgfpathlineto{\pgfqpoint{1.752829in}{1.786899in}}%
\pgfpathlineto{\pgfqpoint{1.752844in}{2.006291in}}%
\pgfpathlineto{\pgfqpoint{1.753607in}{1.927613in}}%
\pgfpathlineto{\pgfqpoint{1.753611in}{1.913529in}}%
\pgfpathlineto{\pgfqpoint{1.753627in}{1.959683in}}%
\pgfpathlineto{\pgfqpoint{1.754374in}{1.931909in}}%
\pgfpathlineto{\pgfqpoint{1.754379in}{1.928660in}}%
\pgfpathlineto{\pgfqpoint{1.754457in}{1.943754in}}%
\pgfpathlineto{\pgfqpoint{1.755130in}{1.936852in}}%
\pgfpathlineto{\pgfqpoint{1.755229in}{1.939046in}}%
\pgfpathlineto{\pgfqpoint{1.755328in}{1.889297in}}%
\pgfpathlineto{\pgfqpoint{1.755756in}{1.532025in}}%
\pgfpathlineto{\pgfqpoint{1.756360in}{1.532041in}}%
\pgfpathlineto{\pgfqpoint{1.757732in}{1.532121in}}%
\pgfpathlineto{\pgfqpoint{1.757774in}{1.763369in}}%
\pgfpathlineto{\pgfqpoint{1.757831in}{1.993506in}}%
\pgfpathlineto{\pgfqpoint{1.758555in}{1.932063in}}%
\pgfpathlineto{\pgfqpoint{1.758560in}{1.915985in}}%
\pgfpathlineto{\pgfqpoint{1.758575in}{1.956558in}}%
\pgfpathlineto{\pgfqpoint{1.759320in}{1.943904in}}%
\pgfpathlineto{\pgfqpoint{1.759442in}{1.930031in}}%
\pgfpathlineto{\pgfqpoint{1.760087in}{1.935880in}}%
\pgfpathlineto{\pgfqpoint{1.760186in}{1.938591in}}%
\pgfpathlineto{\pgfqpoint{1.760235in}{1.906894in}}%
\pgfpathlineto{\pgfqpoint{1.760735in}{1.532032in}}%
\pgfpathlineto{\pgfqpoint{1.761452in}{1.532050in}}%
\pgfpathlineto{\pgfqpoint{1.762683in}{1.532308in}}%
\pgfpathlineto{\pgfqpoint{1.762743in}{2.026300in}}%
\pgfpathlineto{\pgfqpoint{1.763740in}{1.933963in}}%
\pgfpathlineto{\pgfqpoint{1.763748in}{1.915111in}}%
\pgfpathlineto{\pgfqpoint{1.763761in}{1.957920in}}%
\pgfpathlineto{\pgfqpoint{1.764507in}{1.939797in}}%
\pgfpathlineto{\pgfqpoint{1.764512in}{1.943242in}}%
\pgfpathlineto{\pgfqpoint{1.765194in}{1.914510in}}%
\pgfpathlineto{\pgfqpoint{1.765692in}{1.532030in}}%
\pgfpathlineto{\pgfqpoint{1.766272in}{1.532042in}}%
\pgfpathlineto{\pgfqpoint{1.767625in}{1.532138in}}%
\pgfpathlineto{\pgfqpoint{1.767690in}{2.043180in}}%
\pgfpathlineto{\pgfqpoint{1.768831in}{1.930238in}}%
\pgfpathlineto{\pgfqpoint{1.768867in}{1.918308in}}%
\pgfpathlineto{\pgfqpoint{1.768851in}{1.955565in}}%
\pgfpathlineto{\pgfqpoint{1.769596in}{1.941645in}}%
\pgfpathlineto{\pgfqpoint{1.770680in}{1.532032in}}%
\pgfpathlineto{\pgfqpoint{1.770873in}{1.532036in}}%
\pgfpathlineto{\pgfqpoint{1.772568in}{1.532120in}}%
\pgfpathlineto{\pgfqpoint{1.772611in}{1.793417in}}%
\pgfpathlineto{\pgfqpoint{1.772626in}{2.013801in}}%
\pgfpathlineto{\pgfqpoint{1.773391in}{1.915794in}}%
\pgfpathlineto{\pgfqpoint{1.773395in}{1.910307in}}%
\pgfpathlineto{\pgfqpoint{1.773408in}{1.961720in}}%
\pgfpathlineto{\pgfqpoint{1.774137in}{1.943462in}}%
\pgfpathlineto{\pgfqpoint{1.774142in}{1.945374in}}%
\pgfpathlineto{\pgfqpoint{1.774153in}{1.929108in}}%
\pgfpathlineto{\pgfqpoint{1.774884in}{1.938250in}}%
\pgfpathlineto{\pgfqpoint{1.774983in}{1.939821in}}%
\pgfpathlineto{\pgfqpoint{1.775132in}{1.871915in}}%
\pgfpathlineto{\pgfqpoint{1.775562in}{1.532029in}}%
\pgfpathlineto{\pgfqpoint{1.776115in}{1.532040in}}%
\pgfpathlineto{\pgfqpoint{1.777514in}{1.532121in}}%
\pgfpathlineto{\pgfqpoint{1.777555in}{1.754735in}}%
\pgfpathlineto{\pgfqpoint{1.777617in}{1.992532in}}%
\pgfpathlineto{\pgfqpoint{1.778338in}{1.933113in}}%
\pgfpathlineto{\pgfqpoint{1.778343in}{1.916615in}}%
\pgfpathlineto{\pgfqpoint{1.778358in}{1.955701in}}%
\pgfpathlineto{\pgfqpoint{1.779101in}{1.937945in}}%
\pgfpathlineto{\pgfqpoint{1.779107in}{1.944081in}}%
\pgfpathlineto{\pgfqpoint{1.779161in}{1.929485in}}%
\pgfpathlineto{\pgfqpoint{1.779865in}{1.940371in}}%
\pgfpathlineto{\pgfqpoint{1.780112in}{1.827290in}}%
\pgfpathlineto{\pgfqpoint{1.780521in}{1.532032in}}%
\pgfpathlineto{\pgfqpoint{1.781044in}{1.532045in}}%
\pgfpathlineto{\pgfqpoint{1.782460in}{1.532123in}}%
\pgfpathlineto{\pgfqpoint{1.782524in}{2.031491in}}%
\pgfpathlineto{\pgfqpoint{1.783788in}{1.925227in}}%
\pgfpathlineto{\pgfqpoint{1.783792in}{1.922313in}}%
\pgfpathlineto{\pgfqpoint{1.783837in}{1.949923in}}%
\pgfpathlineto{\pgfqpoint{1.784539in}{1.937490in}}%
\pgfpathlineto{\pgfqpoint{1.784546in}{1.939849in}}%
\pgfpathlineto{\pgfqpoint{1.785040in}{1.849088in}}%
\pgfpathlineto{\pgfqpoint{1.785508in}{1.532033in}}%
\pgfpathlineto{\pgfqpoint{1.785996in}{1.532044in}}%
\pgfpathlineto{\pgfqpoint{1.787407in}{1.532153in}}%
\pgfpathlineto{\pgfqpoint{1.787470in}{2.031634in}}%
\pgfpathlineto{\pgfqpoint{1.788611in}{1.936848in}}%
\pgfpathlineto{\pgfqpoint{1.788619in}{1.919158in}}%
\pgfpathlineto{\pgfqpoint{1.788634in}{1.953232in}}%
\pgfpathlineto{\pgfqpoint{1.789382in}{1.932222in}}%
\pgfpathlineto{\pgfqpoint{1.789399in}{1.940933in}}%
\pgfpathlineto{\pgfqpoint{1.789977in}{1.864459in}}%
\pgfpathlineto{\pgfqpoint{1.790420in}{1.532032in}}%
\pgfpathlineto{\pgfqpoint{1.790952in}{1.532046in}}%
\pgfpathlineto{\pgfqpoint{1.792351in}{1.532124in}}%
\pgfpathlineto{\pgfqpoint{1.792415in}{2.031634in}}%
\pgfpathlineto{\pgfqpoint{1.793705in}{1.937272in}}%
\pgfpathlineto{\pgfqpoint{1.793712in}{1.922772in}}%
\pgfpathlineto{\pgfqpoint{1.793728in}{1.949828in}}%
\pgfpathlineto{\pgfqpoint{1.794478in}{1.933175in}}%
\pgfpathlineto{\pgfqpoint{1.794494in}{1.939517in}}%
\pgfpathlineto{\pgfqpoint{1.794883in}{1.908201in}}%
\pgfpathlineto{\pgfqpoint{1.795372in}{1.532032in}}%
\pgfpathlineto{\pgfqpoint{1.795959in}{1.532045in}}%
\pgfpathlineto{\pgfqpoint{1.797296in}{1.532121in}}%
\pgfpathlineto{\pgfqpoint{1.797337in}{1.750159in}}%
\pgfpathlineto{\pgfqpoint{1.797399in}{1.992201in}}%
\pgfpathlineto{\pgfqpoint{1.798118in}{1.950682in}}%
\pgfpathlineto{\pgfqpoint{1.798128in}{1.915590in}}%
\pgfpathlineto{\pgfqpoint{1.798144in}{1.957150in}}%
\pgfpathlineto{\pgfqpoint{1.798889in}{1.933638in}}%
\pgfpathlineto{\pgfqpoint{1.798899in}{1.944063in}}%
\pgfpathlineto{\pgfqpoint{1.798920in}{1.928934in}}%
\pgfpathlineto{\pgfqpoint{1.799651in}{1.935488in}}%
\pgfpathlineto{\pgfqpoint{1.799701in}{1.937705in}}%
\pgfpathlineto{\pgfqpoint{1.799750in}{1.926939in}}%
\pgfpathlineto{\pgfqpoint{1.799800in}{1.932293in}}%
\pgfpathlineto{\pgfqpoint{1.800337in}{1.532033in}}%
\pgfpathlineto{\pgfqpoint{1.801072in}{1.532050in}}%
\pgfpathlineto{\pgfqpoint{1.802243in}{1.532141in}}%
\pgfpathlineto{\pgfqpoint{1.802305in}{2.031724in}}%
\pgfpathlineto{\pgfqpoint{1.803481in}{1.927985in}}%
\pgfpathlineto{\pgfqpoint{1.803485in}{1.919942in}}%
\pgfpathlineto{\pgfqpoint{1.803500in}{1.953111in}}%
\pgfpathlineto{\pgfqpoint{1.804247in}{1.934527in}}%
\pgfpathlineto{\pgfqpoint{1.805305in}{1.532033in}}%
\pgfpathlineto{\pgfqpoint{1.804268in}{1.940568in}}%
\pgfpathlineto{\pgfqpoint{1.805650in}{1.532042in}}%
\pgfpathlineto{\pgfqpoint{1.807188in}{1.532130in}}%
\pgfpathlineto{\pgfqpoint{1.807251in}{2.031739in}}%
\pgfpathlineto{\pgfqpoint{1.808480in}{1.940521in}}%
\pgfpathlineto{\pgfqpoint{1.808488in}{1.921342in}}%
\pgfpathlineto{\pgfqpoint{1.808503in}{1.951447in}}%
\pgfpathlineto{\pgfqpoint{1.809254in}{1.932332in}}%
\pgfpathlineto{\pgfqpoint{1.809269in}{1.939950in}}%
\pgfpathlineto{\pgfqpoint{1.809749in}{1.877835in}}%
\pgfpathlineto{\pgfqpoint{1.810195in}{1.532032in}}%
\pgfpathlineto{\pgfqpoint{1.810720in}{1.532045in}}%
\pgfpathlineto{\pgfqpoint{1.812132in}{1.532121in}}%
\pgfpathlineto{\pgfqpoint{1.812175in}{1.798679in}}%
\pgfpathlineto{\pgfqpoint{1.812190in}{2.019208in}}%
\pgfpathlineto{\pgfqpoint{1.812954in}{1.913888in}}%
\pgfpathlineto{\pgfqpoint{1.812958in}{1.908191in}}%
\pgfpathlineto{\pgfqpoint{1.813001in}{1.963221in}}%
\pgfpathlineto{\pgfqpoint{1.813700in}{1.946214in}}%
\pgfpathlineto{\pgfqpoint{1.813716in}{1.926877in}}%
\pgfpathlineto{\pgfqpoint{1.814449in}{1.937790in}}%
\pgfpathlineto{\pgfqpoint{1.814581in}{1.938378in}}%
\pgfpathlineto{\pgfqpoint{1.814630in}{1.914617in}}%
\pgfpathlineto{\pgfqpoint{1.815102in}{1.532024in}}%
\pgfpathlineto{\pgfqpoint{1.816014in}{1.532048in}}%
\pgfpathlineto{\pgfqpoint{1.817078in}{1.532121in}}%
\pgfpathlineto{\pgfqpoint{1.817119in}{1.754878in}}%
\pgfpathlineto{\pgfqpoint{1.817181in}{1.992240in}}%
\pgfpathlineto{\pgfqpoint{1.817897in}{1.953096in}}%
\pgfpathlineto{\pgfqpoint{1.817908in}{1.915878in}}%
\pgfpathlineto{\pgfqpoint{1.817923in}{1.956260in}}%
\pgfpathlineto{\pgfqpoint{1.818670in}{1.943743in}}%
\pgfpathlineto{\pgfqpoint{1.818726in}{1.929236in}}%
\pgfpathlineto{\pgfqpoint{1.819459in}{1.933756in}}%
\pgfpathlineto{\pgfqpoint{1.819558in}{1.935954in}}%
\pgfpathlineto{\pgfqpoint{1.819607in}{1.897491in}}%
\pgfpathlineto{\pgfqpoint{1.820047in}{1.532027in}}%
\pgfpathlineto{\pgfqpoint{1.820690in}{1.532043in}}%
\pgfpathlineto{\pgfqpoint{1.822023in}{1.532121in}}%
\pgfpathlineto{\pgfqpoint{1.822066in}{1.798143in}}%
\pgfpathlineto{\pgfqpoint{1.822081in}{2.018661in}}%
\pgfpathlineto{\pgfqpoint{1.822844in}{1.918669in}}%
\pgfpathlineto{\pgfqpoint{1.822848in}{1.907772in}}%
\pgfpathlineto{\pgfqpoint{1.822891in}{1.961705in}}%
\pgfpathlineto{\pgfqpoint{1.823606in}{1.926883in}}%
\pgfpathlineto{\pgfqpoint{1.823622in}{1.945676in}}%
\pgfpathlineto{\pgfqpoint{1.824381in}{1.940085in}}%
\pgfpathlineto{\pgfqpoint{1.824628in}{1.818018in}}%
\pgfpathlineto{\pgfqpoint{1.825081in}{1.532033in}}%
\pgfpathlineto{\pgfqpoint{1.825570in}{1.532043in}}%
\pgfpathlineto{\pgfqpoint{1.826969in}{1.532125in}}%
\pgfpathlineto{\pgfqpoint{1.827026in}{2.015360in}}%
\pgfpathlineto{\pgfqpoint{1.828395in}{1.929756in}}%
\pgfpathlineto{\pgfqpoint{1.828401in}{1.925079in}}%
\pgfpathlineto{\pgfqpoint{1.828417in}{1.947202in}}%
\pgfpathlineto{\pgfqpoint{1.829154in}{1.937164in}}%
\pgfpathlineto{\pgfqpoint{1.829401in}{1.939524in}}%
\pgfpathlineto{\pgfqpoint{1.829500in}{1.902273in}}%
\pgfpathlineto{\pgfqpoint{1.830026in}{1.532033in}}%
\pgfpathlineto{\pgfqpoint{1.830614in}{1.532046in}}%
\pgfpathlineto{\pgfqpoint{1.831915in}{1.532136in}}%
\pgfpathlineto{\pgfqpoint{1.831978in}{2.031315in}}%
\pgfpathlineto{\pgfqpoint{1.833179in}{1.937283in}}%
\pgfpathlineto{\pgfqpoint{1.833187in}{1.920403in}}%
\pgfpathlineto{\pgfqpoint{1.833202in}{1.951754in}}%
\pgfpathlineto{\pgfqpoint{1.833952in}{1.932075in}}%
\pgfpathlineto{\pgfqpoint{1.833967in}{1.940129in}}%
\pgfpathlineto{\pgfqpoint{1.834437in}{1.913877in}}%
\pgfpathlineto{\pgfqpoint{1.834920in}{1.532031in}}%
\pgfpathlineto{\pgfqpoint{1.835540in}{1.532048in}}%
\pgfpathlineto{\pgfqpoint{1.836856in}{1.532101in}}%
\pgfpathlineto{\pgfqpoint{1.836865in}{1.532310in}}%
\pgfpathlineto{\pgfqpoint{1.836924in}{2.026501in}}%
\pgfpathlineto{\pgfqpoint{1.837922in}{1.933731in}}%
\pgfpathlineto{\pgfqpoint{1.837929in}{1.914614in}}%
\pgfpathlineto{\pgfqpoint{1.837942in}{1.957556in}}%
\pgfpathlineto{\pgfqpoint{1.838687in}{1.939401in}}%
\pgfpathlineto{\pgfqpoint{1.838692in}{1.942870in}}%
\pgfpathlineto{\pgfqpoint{1.839395in}{1.907268in}}%
\pgfpathlineto{\pgfqpoint{1.839913in}{1.532032in}}%
\pgfpathlineto{\pgfqpoint{1.840452in}{1.532046in}}%
\pgfpathlineto{\pgfqpoint{1.841805in}{1.532125in}}%
\pgfpathlineto{\pgfqpoint{1.841862in}{2.015143in}}%
\pgfpathlineto{\pgfqpoint{1.843231in}{1.929748in}}%
\pgfpathlineto{\pgfqpoint{1.843237in}{1.924972in}}%
\pgfpathlineto{\pgfqpoint{1.843253in}{1.947158in}}%
\pgfpathlineto{\pgfqpoint{1.843984in}{1.939240in}}%
\pgfpathlineto{\pgfqpoint{1.844086in}{1.939721in}}%
\pgfpathlineto{\pgfqpoint{1.844440in}{1.766942in}}%
\pgfpathlineto{\pgfqpoint{1.844863in}{1.532033in}}%
\pgfpathlineto{\pgfqpoint{1.845352in}{1.532044in}}%
\pgfpathlineto{\pgfqpoint{1.846752in}{1.532135in}}%
\pgfpathlineto{\pgfqpoint{1.846815in}{2.031414in}}%
\pgfpathlineto{\pgfqpoint{1.848016in}{1.934204in}}%
\pgfpathlineto{\pgfqpoint{1.848023in}{1.920429in}}%
\pgfpathlineto{\pgfqpoint{1.848068in}{1.951068in}}%
\pgfpathlineto{\pgfqpoint{1.848788in}{1.932114in}}%
\pgfpathlineto{\pgfqpoint{1.848804in}{1.940071in}}%
\pgfpathlineto{\pgfqpoint{1.849282in}{1.911014in}}%
\pgfpathlineto{\pgfqpoint{1.849758in}{1.532032in}}%
\pgfpathlineto{\pgfqpoint{1.850382in}{1.532047in}}%
\pgfpathlineto{\pgfqpoint{1.851696in}{1.532121in}}%
\pgfpathlineto{\pgfqpoint{1.851739in}{1.797725in}}%
\pgfpathlineto{\pgfqpoint{1.851753in}{2.018161in}}%
\pgfpathlineto{\pgfqpoint{1.852516in}{1.918556in}}%
\pgfpathlineto{\pgfqpoint{1.852520in}{1.907720in}}%
\pgfpathlineto{\pgfqpoint{1.852564in}{1.961465in}}%
\pgfpathlineto{\pgfqpoint{1.853280in}{1.926730in}}%
\pgfpathlineto{\pgfqpoint{1.853296in}{1.945588in}}%
\pgfpathlineto{\pgfqpoint{1.854043in}{1.936008in}}%
\pgfpathlineto{\pgfqpoint{1.854137in}{1.932246in}}%
\pgfpathlineto{\pgfqpoint{1.854187in}{1.938123in}}%
\pgfpathlineto{\pgfqpoint{1.854706in}{1.532031in}}%
\pgfpathlineto{\pgfqpoint{1.855478in}{1.532048in}}%
\pgfpathlineto{\pgfqpoint{1.856641in}{1.532119in}}%
\pgfpathlineto{\pgfqpoint{1.856684in}{1.786411in}}%
\pgfpathlineto{\pgfqpoint{1.856698in}{2.005484in}}%
\pgfpathlineto{\pgfqpoint{1.857460in}{1.932845in}}%
\pgfpathlineto{\pgfqpoint{1.857497in}{1.913014in}}%
\pgfpathlineto{\pgfqpoint{1.857484in}{1.958479in}}%
\pgfpathlineto{\pgfqpoint{1.858227in}{1.934289in}}%
\pgfpathlineto{\pgfqpoint{1.858235in}{1.928069in}}%
\pgfpathlineto{\pgfqpoint{1.858281in}{1.943324in}}%
\pgfpathlineto{\pgfqpoint{1.858958in}{1.932551in}}%
\pgfpathlineto{\pgfqpoint{1.859008in}{1.939528in}}%
\pgfpathlineto{\pgfqpoint{1.859205in}{1.882059in}}%
\pgfpathlineto{\pgfqpoint{1.859659in}{1.532032in}}%
\pgfpathlineto{\pgfqpoint{1.860192in}{1.532045in}}%
\pgfpathlineto{\pgfqpoint{1.861587in}{1.532122in}}%
\pgfpathlineto{\pgfqpoint{1.861628in}{1.744749in}}%
\pgfpathlineto{\pgfqpoint{1.861690in}{1.994324in}}%
\pgfpathlineto{\pgfqpoint{1.862411in}{1.940825in}}%
\pgfpathlineto{\pgfqpoint{1.862420in}{1.914891in}}%
\pgfpathlineto{\pgfqpoint{1.862433in}{1.956381in}}%
\pgfpathlineto{\pgfqpoint{1.863175in}{1.932640in}}%
\pgfpathlineto{\pgfqpoint{1.863226in}{1.943529in}}%
\pgfpathlineto{\pgfqpoint{1.863204in}{1.928799in}}%
\pgfpathlineto{\pgfqpoint{1.863942in}{1.930984in}}%
\pgfpathlineto{\pgfqpoint{1.863991in}{1.940342in}}%
\pgfpathlineto{\pgfqpoint{1.864139in}{1.886945in}}%
\pgfpathlineto{\pgfqpoint{1.864609in}{1.532032in}}%
\pgfpathlineto{\pgfqpoint{1.865194in}{1.532044in}}%
\pgfpathlineto{\pgfqpoint{1.866532in}{1.532122in}}%
\pgfpathlineto{\pgfqpoint{1.866575in}{1.780913in}}%
\pgfpathlineto{\pgfqpoint{1.866632in}{2.002588in}}%
\pgfpathlineto{\pgfqpoint{1.867352in}{1.933223in}}%
\pgfpathlineto{\pgfqpoint{1.867388in}{1.914043in}}%
\pgfpathlineto{\pgfqpoint{1.867375in}{1.957644in}}%
\pgfpathlineto{\pgfqpoint{1.868123in}{1.929575in}}%
\pgfpathlineto{\pgfqpoint{1.868142in}{1.943807in}}%
\pgfpathlineto{\pgfqpoint{1.868188in}{1.929562in}}%
\pgfpathlineto{\pgfqpoint{1.868944in}{1.938878in}}%
\pgfpathlineto{\pgfqpoint{1.869093in}{1.859834in}}%
\pgfpathlineto{\pgfqpoint{1.869597in}{1.532032in}}%
\pgfpathlineto{\pgfqpoint{1.870091in}{1.532043in}}%
\pgfpathlineto{\pgfqpoint{1.871477in}{1.532114in}}%
\pgfpathlineto{\pgfqpoint{1.871519in}{1.745329in}}%
\pgfpathlineto{\pgfqpoint{1.871580in}{1.991475in}}%
\pgfpathlineto{\pgfqpoint{1.872304in}{1.946614in}}%
\pgfpathlineto{\pgfqpoint{1.872314in}{1.915945in}}%
\pgfpathlineto{\pgfqpoint{1.872329in}{1.955867in}}%
\pgfpathlineto{\pgfqpoint{1.873070in}{1.944082in}}%
\pgfpathlineto{\pgfqpoint{1.873088in}{1.928390in}}%
\pgfpathlineto{\pgfqpoint{1.873837in}{1.937511in}}%
\pgfpathlineto{\pgfqpoint{1.873936in}{1.917907in}}%
\pgfpathlineto{\pgfqpoint{1.873986in}{1.920572in}}%
\pgfpathlineto{\pgfqpoint{1.874507in}{1.532032in}}%
\pgfpathlineto{\pgfqpoint{1.875242in}{1.532050in}}%
\pgfpathlineto{\pgfqpoint{1.876423in}{1.532121in}}%
\pgfpathlineto{\pgfqpoint{1.876465in}{1.762286in}}%
\pgfpathlineto{\pgfqpoint{1.876522in}{1.993450in}}%
\pgfpathlineto{\pgfqpoint{1.877244in}{1.942814in}}%
\pgfpathlineto{\pgfqpoint{1.877253in}{1.914317in}}%
\pgfpathlineto{\pgfqpoint{1.877297in}{1.955924in}}%
\pgfpathlineto{\pgfqpoint{1.878011in}{1.942591in}}%
\pgfpathlineto{\pgfqpoint{1.878033in}{1.929224in}}%
\pgfpathlineto{\pgfqpoint{1.878047in}{1.943464in}}%
\pgfpathlineto{\pgfqpoint{1.878756in}{1.930895in}}%
\pgfpathlineto{\pgfqpoint{1.878806in}{1.939797in}}%
\pgfpathlineto{\pgfqpoint{1.879004in}{1.848061in}}%
\pgfpathlineto{\pgfqpoint{1.879393in}{1.532026in}}%
\pgfpathlineto{\pgfqpoint{1.879973in}{1.532048in}}%
\pgfpathlineto{\pgfqpoint{1.881371in}{1.532147in}}%
\pgfpathlineto{\pgfqpoint{1.881438in}{2.018570in}}%
\pgfpathlineto{\pgfqpoint{1.882549in}{1.934575in}}%
\pgfpathlineto{\pgfqpoint{1.882586in}{1.920402in}}%
\pgfpathlineto{\pgfqpoint{1.882573in}{1.951276in}}%
\pgfpathlineto{\pgfqpoint{1.883313in}{1.938751in}}%
\pgfpathlineto{\pgfqpoint{1.883388in}{1.940835in}}%
\pgfpathlineto{\pgfqpoint{1.883904in}{1.905958in}}%
\pgfpathlineto{\pgfqpoint{1.884338in}{1.532032in}}%
\pgfpathlineto{\pgfqpoint{1.884963in}{1.532045in}}%
\pgfpathlineto{\pgfqpoint{1.886314in}{1.532121in}}%
\pgfpathlineto{\pgfqpoint{1.886355in}{1.744455in}}%
\pgfpathlineto{\pgfqpoint{1.886417in}{1.994259in}}%
\pgfpathlineto{\pgfqpoint{1.887137in}{1.943140in}}%
\pgfpathlineto{\pgfqpoint{1.887146in}{1.914673in}}%
\pgfpathlineto{\pgfqpoint{1.887191in}{1.955913in}}%
\pgfpathlineto{\pgfqpoint{1.887906in}{1.933476in}}%
\pgfpathlineto{\pgfqpoint{1.887917in}{1.943725in}}%
\pgfpathlineto{\pgfqpoint{1.887938in}{1.928605in}}%
\pgfpathlineto{\pgfqpoint{1.888659in}{1.939862in}}%
\pgfpathlineto{\pgfqpoint{1.888955in}{1.770158in}}%
\pgfpathlineto{\pgfqpoint{1.889330in}{1.532032in}}%
\pgfpathlineto{\pgfqpoint{1.889864in}{1.532045in}}%
\pgfpathlineto{\pgfqpoint{1.891261in}{1.532134in}}%
\pgfpathlineto{\pgfqpoint{1.891324in}{2.032018in}}%
\pgfpathlineto{\pgfqpoint{1.892531in}{1.922447in}}%
\pgfpathlineto{\pgfqpoint{1.892535in}{1.920607in}}%
\pgfpathlineto{\pgfqpoint{1.892547in}{1.951181in}}%
\pgfpathlineto{\pgfqpoint{1.893277in}{1.934194in}}%
\pgfpathlineto{\pgfqpoint{1.893285in}{1.940125in}}%
\pgfpathlineto{\pgfqpoint{1.893838in}{1.844484in}}%
\pgfpathlineto{\pgfqpoint{1.894271in}{1.532030in}}%
\pgfpathlineto{\pgfqpoint{1.894799in}{1.532047in}}%
\pgfpathlineto{\pgfqpoint{1.896206in}{1.532128in}}%
\pgfpathlineto{\pgfqpoint{1.896270in}{2.036999in}}%
\pgfpathlineto{\pgfqpoint{1.897472in}{1.931676in}}%
\pgfpathlineto{\pgfqpoint{1.897479in}{1.919991in}}%
\pgfpathlineto{\pgfqpoint{1.897491in}{1.952009in}}%
\pgfpathlineto{\pgfqpoint{1.898240in}{1.935648in}}%
\pgfpathlineto{\pgfqpoint{1.899216in}{1.532032in}}%
\pgfpathlineto{\pgfqpoint{1.898291in}{1.939787in}}%
\pgfpathlineto{\pgfqpoint{1.899448in}{1.532038in}}%
\pgfpathlineto{\pgfqpoint{1.901151in}{1.532131in}}%
\pgfpathlineto{\pgfqpoint{1.901251in}{2.008877in}}%
\pgfpathlineto{\pgfqpoint{1.902538in}{1.946862in}}%
\pgfpathlineto{\pgfqpoint{1.902553in}{1.925493in}}%
\pgfpathlineto{\pgfqpoint{1.903312in}{1.939987in}}%
\pgfpathlineto{\pgfqpoint{1.904163in}{1.532032in}}%
\pgfpathlineto{\pgfqpoint{1.904990in}{1.532052in}}%
\pgfpathlineto{\pgfqpoint{1.906096in}{1.532124in}}%
\pgfpathlineto{\pgfqpoint{1.906137in}{1.751557in}}%
\pgfpathlineto{\pgfqpoint{1.906195in}{1.991311in}}%
\pgfpathlineto{\pgfqpoint{1.906917in}{1.952841in}}%
\pgfpathlineto{\pgfqpoint{1.906927in}{1.916047in}}%
\pgfpathlineto{\pgfqpoint{1.906975in}{1.956078in}}%
\pgfpathlineto{\pgfqpoint{1.907688in}{1.933213in}}%
\pgfpathlineto{\pgfqpoint{1.907698in}{1.943647in}}%
\pgfpathlineto{\pgfqpoint{1.907720in}{1.928448in}}%
\pgfpathlineto{\pgfqpoint{1.908449in}{1.940049in}}%
\pgfpathlineto{\pgfqpoint{1.908745in}{1.754738in}}%
\pgfpathlineto{\pgfqpoint{1.909137in}{1.532032in}}%
\pgfpathlineto{\pgfqpoint{1.909674in}{1.532046in}}%
\pgfpathlineto{\pgfqpoint{1.911041in}{1.532121in}}%
\pgfpathlineto{\pgfqpoint{1.911083in}{1.751818in}}%
\pgfpathlineto{\pgfqpoint{1.911144in}{1.991248in}}%
\pgfpathlineto{\pgfqpoint{1.911862in}{1.952351in}}%
\pgfpathlineto{\pgfqpoint{1.911872in}{1.915963in}}%
\pgfpathlineto{\pgfqpoint{1.911888in}{1.955314in}}%
\pgfpathlineto{\pgfqpoint{1.912633in}{1.933091in}}%
\pgfpathlineto{\pgfqpoint{1.912643in}{1.943509in}}%
\pgfpathlineto{\pgfqpoint{1.912664in}{1.928402in}}%
\pgfpathlineto{\pgfqpoint{1.913395in}{1.934725in}}%
\pgfpathlineto{\pgfqpoint{1.913444in}{1.937385in}}%
\pgfpathlineto{\pgfqpoint{1.913494in}{1.928118in}}%
\pgfpathlineto{\pgfqpoint{1.913543in}{1.935124in}}%
\pgfpathlineto{\pgfqpoint{1.914055in}{1.532032in}}%
\pgfpathlineto{\pgfqpoint{1.914833in}{1.532049in}}%
\pgfpathlineto{\pgfqpoint{1.915988in}{1.532133in}}%
\pgfpathlineto{\pgfqpoint{1.916050in}{2.031673in}}%
\pgfpathlineto{\pgfqpoint{1.917256in}{1.924604in}}%
\pgfpathlineto{\pgfqpoint{1.917260in}{1.919920in}}%
\pgfpathlineto{\pgfqpoint{1.917275in}{1.951234in}}%
\pgfpathlineto{\pgfqpoint{1.918014in}{1.939882in}}%
\pgfpathlineto{\pgfqpoint{1.919008in}{1.532032in}}%
\pgfpathlineto{\pgfqpoint{1.919445in}{1.532043in}}%
\pgfpathlineto{\pgfqpoint{1.920932in}{1.532121in}}%
\pgfpathlineto{\pgfqpoint{1.920973in}{1.741499in}}%
\pgfpathlineto{\pgfqpoint{1.921034in}{1.996890in}}%
\pgfpathlineto{\pgfqpoint{1.921752in}{1.952059in}}%
\pgfpathlineto{\pgfqpoint{1.921763in}{1.914905in}}%
\pgfpathlineto{\pgfqpoint{1.921778in}{1.956578in}}%
\pgfpathlineto{\pgfqpoint{1.922525in}{1.931375in}}%
\pgfpathlineto{\pgfqpoint{1.922616in}{1.942657in}}%
\pgfpathlineto{\pgfqpoint{1.922593in}{1.928486in}}%
\pgfpathlineto{\pgfqpoint{1.923263in}{1.930958in}}%
\pgfpathlineto{\pgfqpoint{1.923312in}{1.940357in}}%
\pgfpathlineto{\pgfqpoint{1.923510in}{1.869033in}}%
\pgfpathlineto{\pgfqpoint{1.923947in}{1.532032in}}%
\pgfpathlineto{\pgfqpoint{1.924480in}{1.532045in}}%
\pgfpathlineto{\pgfqpoint{1.925880in}{1.532146in}}%
\pgfpathlineto{\pgfqpoint{1.925942in}{2.031065in}}%
\pgfpathlineto{\pgfqpoint{1.927111in}{1.947095in}}%
\pgfpathlineto{\pgfqpoint{1.927122in}{1.919283in}}%
\pgfpathlineto{\pgfqpoint{1.927165in}{1.951623in}}%
\pgfpathlineto{\pgfqpoint{1.927884in}{1.932589in}}%
\pgfpathlineto{\pgfqpoint{1.928846in}{1.532022in}}%
\pgfpathlineto{\pgfqpoint{1.927903in}{1.939886in}}%
\pgfpathlineto{\pgfqpoint{1.929919in}{1.532050in}}%
\pgfpathlineto{\pgfqpoint{1.930823in}{1.532123in}}%
\pgfpathlineto{\pgfqpoint{1.930863in}{1.732862in}}%
\pgfpathlineto{\pgfqpoint{1.930926in}{1.992590in}}%
\pgfpathlineto{\pgfqpoint{1.931647in}{1.956563in}}%
\pgfpathlineto{\pgfqpoint{1.931662in}{1.915545in}}%
\pgfpathlineto{\pgfqpoint{1.932415in}{1.935869in}}%
\pgfpathlineto{\pgfqpoint{1.932426in}{1.943845in}}%
\pgfpathlineto{\pgfqpoint{1.932444in}{1.928281in}}%
\pgfpathlineto{\pgfqpoint{1.933188in}{1.940282in}}%
\pgfpathlineto{\pgfqpoint{1.933484in}{1.741619in}}%
\pgfpathlineto{\pgfqpoint{1.933853in}{1.532028in}}%
\pgfpathlineto{\pgfqpoint{1.934390in}{1.532050in}}%
\pgfpathlineto{\pgfqpoint{1.935760in}{1.532086in}}%
\pgfpathlineto{\pgfqpoint{1.935771in}{1.532147in}}%
\pgfpathlineto{\pgfqpoint{1.935837in}{2.035636in}}%
\pgfpathlineto{\pgfqpoint{1.936947in}{1.932859in}}%
\pgfpathlineto{\pgfqpoint{1.936955in}{1.917353in}}%
\pgfpathlineto{\pgfqpoint{1.936968in}{1.954280in}}%
\pgfpathlineto{\pgfqpoint{1.937714in}{1.938383in}}%
\pgfpathlineto{\pgfqpoint{1.937719in}{1.941259in}}%
\pgfpathlineto{\pgfqpoint{1.938317in}{1.894331in}}%
\pgfpathlineto{\pgfqpoint{1.938771in}{1.532032in}}%
\pgfpathlineto{\pgfqpoint{1.939339in}{1.532044in}}%
\pgfpathlineto{\pgfqpoint{1.940714in}{1.532121in}}%
\pgfpathlineto{\pgfqpoint{1.940757in}{1.799749in}}%
\pgfpathlineto{\pgfqpoint{1.940772in}{2.019982in}}%
\pgfpathlineto{\pgfqpoint{1.941534in}{1.921879in}}%
\pgfpathlineto{\pgfqpoint{1.941538in}{1.907655in}}%
\pgfpathlineto{\pgfqpoint{1.941555in}{1.964008in}}%
\pgfpathlineto{\pgfqpoint{1.942300in}{1.926103in}}%
\pgfpathlineto{\pgfqpoint{1.942316in}{1.945270in}}%
\pgfpathlineto{\pgfqpoint{1.943046in}{1.937204in}}%
\pgfpathlineto{\pgfqpoint{1.943130in}{1.939226in}}%
\pgfpathlineto{\pgfqpoint{1.943328in}{1.812686in}}%
\pgfpathlineto{\pgfqpoint{1.943726in}{1.532032in}}%
\pgfpathlineto{\pgfqpoint{1.944255in}{1.532045in}}%
\pgfpathlineto{\pgfqpoint{1.945659in}{1.532121in}}%
\pgfpathlineto{\pgfqpoint{1.945701in}{1.750909in}}%
\pgfpathlineto{\pgfqpoint{1.945763in}{1.991092in}}%
\pgfpathlineto{\pgfqpoint{1.946480in}{1.952811in}}%
\pgfpathlineto{\pgfqpoint{1.946524in}{1.916146in}}%
\pgfpathlineto{\pgfqpoint{1.946537in}{1.955763in}}%
\pgfpathlineto{\pgfqpoint{1.947251in}{1.933298in}}%
\pgfpathlineto{\pgfqpoint{1.947262in}{1.943474in}}%
\pgfpathlineto{\pgfqpoint{1.947283in}{1.928405in}}%
\pgfpathlineto{\pgfqpoint{1.948008in}{1.939913in}}%
\pgfpathlineto{\pgfqpoint{1.948304in}{1.755673in}}%
\pgfpathlineto{\pgfqpoint{1.948709in}{1.532033in}}%
\pgfpathlineto{\pgfqpoint{1.949198in}{1.532044in}}%
\pgfpathlineto{\pgfqpoint{1.950607in}{1.532151in}}%
\pgfpathlineto{\pgfqpoint{1.950670in}{2.030897in}}%
\pgfpathlineto{\pgfqpoint{1.951814in}{1.926499in}}%
\pgfpathlineto{\pgfqpoint{1.951818in}{1.918563in}}%
\pgfpathlineto{\pgfqpoint{1.951833in}{1.952951in}}%
\pgfpathlineto{\pgfqpoint{1.952579in}{1.933303in}}%
\pgfpathlineto{\pgfqpoint{1.953655in}{1.532032in}}%
\pgfpathlineto{\pgfqpoint{1.952598in}{1.940190in}}%
\pgfpathlineto{\pgfqpoint{1.954044in}{1.532041in}}%
\pgfpathlineto{\pgfqpoint{1.955550in}{1.532121in}}%
\pgfpathlineto{\pgfqpoint{1.955592in}{1.750360in}}%
\pgfpathlineto{\pgfqpoint{1.955654in}{1.991137in}}%
\pgfpathlineto{\pgfqpoint{1.956372in}{1.951574in}}%
\pgfpathlineto{\pgfqpoint{1.956382in}{1.915403in}}%
\pgfpathlineto{\pgfqpoint{1.956397in}{1.955727in}}%
\pgfpathlineto{\pgfqpoint{1.957141in}{1.931301in}}%
\pgfpathlineto{\pgfqpoint{1.957222in}{1.942950in}}%
\pgfpathlineto{\pgfqpoint{1.957238in}{1.928829in}}%
\pgfpathlineto{\pgfqpoint{1.957882in}{1.932587in}}%
\pgfpathlineto{\pgfqpoint{1.957931in}{1.936005in}}%
\pgfpathlineto{\pgfqpoint{1.958080in}{1.907260in}}%
\pgfpathlineto{\pgfqpoint{1.958562in}{1.532032in}}%
\pgfpathlineto{\pgfqpoint{1.959240in}{1.532046in}}%
\pgfpathlineto{\pgfqpoint{1.960496in}{1.532123in}}%
\pgfpathlineto{\pgfqpoint{1.960537in}{1.751911in}}%
\pgfpathlineto{\pgfqpoint{1.960595in}{1.991327in}}%
\pgfpathlineto{\pgfqpoint{1.961321in}{1.937299in}}%
\pgfpathlineto{\pgfqpoint{1.961330in}{1.915448in}}%
\pgfpathlineto{\pgfqpoint{1.961344in}{1.956848in}}%
\pgfpathlineto{\pgfqpoint{1.962084in}{1.927898in}}%
\pgfpathlineto{\pgfqpoint{1.962106in}{1.943494in}}%
\pgfpathlineto{\pgfqpoint{1.962836in}{1.938633in}}%
\pgfpathlineto{\pgfqpoint{1.963083in}{1.841453in}}%
\pgfpathlineto{\pgfqpoint{1.963465in}{1.532032in}}%
\pgfpathlineto{\pgfqpoint{1.964028in}{1.532045in}}%
\pgfpathlineto{\pgfqpoint{1.965442in}{1.532126in}}%
\pgfpathlineto{\pgfqpoint{1.965505in}{2.031130in}}%
\pgfpathlineto{\pgfqpoint{1.966766in}{1.937905in}}%
\pgfpathlineto{\pgfqpoint{1.966774in}{1.921216in}}%
\pgfpathlineto{\pgfqpoint{1.966790in}{1.949926in}}%
\pgfpathlineto{\pgfqpoint{1.967538in}{1.932788in}}%
\pgfpathlineto{\pgfqpoint{1.967558in}{1.938786in}}%
\pgfpathlineto{\pgfqpoint{1.968177in}{1.640616in}}%
\pgfpathlineto{\pgfqpoint{1.968452in}{1.532032in}}%
\pgfpathlineto{\pgfqpoint{1.968392in}{1.660451in}}%
\pgfpathlineto{\pgfqpoint{1.969080in}{1.532047in}}%
\pgfpathlineto{\pgfqpoint{1.970387in}{1.532121in}}%
\pgfpathlineto{\pgfqpoint{1.970428in}{1.742197in}}%
\pgfpathlineto{\pgfqpoint{1.970489in}{1.995433in}}%
\pgfpathlineto{\pgfqpoint{1.971207in}{1.952101in}}%
\pgfpathlineto{\pgfqpoint{1.971217in}{1.914747in}}%
\pgfpathlineto{\pgfqpoint{1.971233in}{1.956325in}}%
\pgfpathlineto{\pgfqpoint{1.971981in}{1.934008in}}%
\pgfpathlineto{\pgfqpoint{1.971992in}{1.943578in}}%
\pgfpathlineto{\pgfqpoint{1.972012in}{1.928200in}}%
\pgfpathlineto{\pgfqpoint{1.972741in}{1.932233in}}%
\pgfpathlineto{\pgfqpoint{1.972791in}{1.936184in}}%
\pgfpathlineto{\pgfqpoint{1.972890in}{1.905579in}}%
\pgfpathlineto{\pgfqpoint{1.973356in}{1.532021in}}%
\pgfpathlineto{\pgfqpoint{1.974356in}{1.532058in}}%
\pgfpathlineto{\pgfqpoint{1.975334in}{1.532142in}}%
\pgfpathlineto{\pgfqpoint{1.975401in}{2.032678in}}%
\pgfpathlineto{\pgfqpoint{1.976511in}{1.930306in}}%
\pgfpathlineto{\pgfqpoint{1.976519in}{1.918245in}}%
\pgfpathlineto{\pgfqpoint{1.976533in}{1.953852in}}%
\pgfpathlineto{\pgfqpoint{1.977277in}{1.935001in}}%
\pgfpathlineto{\pgfqpoint{1.977287in}{1.940684in}}%
\pgfpathlineto{\pgfqpoint{1.977889in}{1.874125in}}%
\pgfpathlineto{\pgfqpoint{1.978343in}{1.532030in}}%
\pgfpathlineto{\pgfqpoint{1.978871in}{1.532047in}}%
\pgfpathlineto{\pgfqpoint{1.980278in}{1.532127in}}%
\pgfpathlineto{\pgfqpoint{1.980342in}{2.030780in}}%
\pgfpathlineto{\pgfqpoint{1.981604in}{1.930877in}}%
\pgfpathlineto{\pgfqpoint{1.981640in}{1.921928in}}%
\pgfpathlineto{\pgfqpoint{1.981625in}{1.949911in}}%
\pgfpathlineto{\pgfqpoint{1.982375in}{1.932589in}}%
\pgfpathlineto{\pgfqpoint{1.982394in}{1.938898in}}%
\pgfpathlineto{\pgfqpoint{1.982944in}{1.737614in}}%
\pgfpathlineto{\pgfqpoint{1.983310in}{1.532032in}}%
\pgfpathlineto{\pgfqpoint{1.983849in}{1.532045in}}%
\pgfpathlineto{\pgfqpoint{1.985226in}{1.532161in}}%
\pgfpathlineto{\pgfqpoint{1.985289in}{2.031046in}}%
\pgfpathlineto{\pgfqpoint{1.986402in}{1.927472in}}%
\pgfpathlineto{\pgfqpoint{1.986406in}{1.917898in}}%
\pgfpathlineto{\pgfqpoint{1.986421in}{1.953677in}}%
\pgfpathlineto{\pgfqpoint{1.987172in}{1.930865in}}%
\pgfpathlineto{\pgfqpoint{1.987189in}{1.940352in}}%
\pgfpathlineto{\pgfqpoint{1.987741in}{1.923461in}}%
\pgfpathlineto{\pgfqpoint{1.988280in}{1.532033in}}%
\pgfpathlineto{\pgfqpoint{1.989016in}{1.532051in}}%
\pgfpathlineto{\pgfqpoint{1.990169in}{1.532127in}}%
\pgfpathlineto{\pgfqpoint{1.990233in}{2.031051in}}%
\pgfpathlineto{\pgfqpoint{1.991467in}{1.924479in}}%
\pgfpathlineto{\pgfqpoint{1.991470in}{1.920532in}}%
\pgfpathlineto{\pgfqpoint{1.991486in}{1.950048in}}%
\pgfpathlineto{\pgfqpoint{1.992222in}{1.939460in}}%
\pgfpathlineto{\pgfqpoint{1.993188in}{1.532032in}}%
\pgfpathlineto{\pgfqpoint{1.993871in}{1.532047in}}%
\pgfpathlineto{\pgfqpoint{1.995115in}{1.532130in}}%
\pgfpathlineto{\pgfqpoint{1.995215in}{2.004288in}}%
\pgfpathlineto{\pgfqpoint{1.996510in}{1.936728in}}%
\pgfpathlineto{\pgfqpoint{1.996520in}{1.926011in}}%
\pgfpathlineto{\pgfqpoint{1.996535in}{1.944995in}}%
\pgfpathlineto{\pgfqpoint{1.997265in}{1.939662in}}%
\pgfpathlineto{\pgfqpoint{1.998141in}{1.532032in}}%
\pgfpathlineto{\pgfqpoint{1.999370in}{1.532061in}}%
\pgfpathlineto{\pgfqpoint{2.000061in}{1.532134in}}%
\pgfpathlineto{\pgfqpoint{2.000124in}{2.031728in}}%
\pgfpathlineto{\pgfqpoint{2.001329in}{1.924237in}}%
\pgfpathlineto{\pgfqpoint{2.001333in}{1.919600in}}%
\pgfpathlineto{\pgfqpoint{2.001348in}{1.950894in}}%
\pgfpathlineto{\pgfqpoint{2.002087in}{1.939563in}}%
\pgfpathlineto{\pgfqpoint{2.003062in}{1.532032in}}%
\pgfpathlineto{\pgfqpoint{2.003481in}{1.532043in}}%
\pgfpathlineto{\pgfqpoint{2.005005in}{1.532119in}}%
\pgfpathlineto{\pgfqpoint{2.005048in}{1.785642in}}%
\pgfpathlineto{\pgfqpoint{2.005105in}{2.004595in}}%
\pgfpathlineto{\pgfqpoint{2.005825in}{1.931060in}}%
\pgfpathlineto{\pgfqpoint{2.005833in}{1.912211in}}%
\pgfpathlineto{\pgfqpoint{2.005845in}{1.958246in}}%
\pgfpathlineto{\pgfqpoint{2.006592in}{1.935084in}}%
\pgfpathlineto{\pgfqpoint{2.006599in}{1.927450in}}%
\pgfpathlineto{\pgfqpoint{2.006648in}{1.943193in}}%
\pgfpathlineto{\pgfqpoint{2.007314in}{1.932891in}}%
\pgfpathlineto{\pgfqpoint{2.007463in}{1.928087in}}%
\pgfpathlineto{\pgfqpoint{2.007512in}{1.936563in}}%
\pgfpathlineto{\pgfqpoint{2.007999in}{1.532031in}}%
\pgfpathlineto{\pgfqpoint{2.008839in}{1.532049in}}%
\pgfpathlineto{\pgfqpoint{2.009951in}{1.532122in}}%
\pgfpathlineto{\pgfqpoint{2.010014in}{2.030960in}}%
\pgfpathlineto{\pgfqpoint{2.011304in}{1.934174in}}%
\pgfpathlineto{\pgfqpoint{2.011311in}{1.921968in}}%
\pgfpathlineto{\pgfqpoint{2.011354in}{1.948445in}}%
\pgfpathlineto{\pgfqpoint{2.012074in}{1.932149in}}%
\pgfpathlineto{\pgfqpoint{2.012091in}{1.938562in}}%
\pgfpathlineto{\pgfqpoint{2.012471in}{1.921477in}}%
\pgfpathlineto{\pgfqpoint{2.013007in}{1.532033in}}%
\pgfpathlineto{\pgfqpoint{2.013694in}{1.532048in}}%
\pgfpathlineto{\pgfqpoint{2.014896in}{1.532127in}}%
\pgfpathlineto{\pgfqpoint{2.014960in}{2.031001in}}%
\pgfpathlineto{\pgfqpoint{2.016194in}{1.924525in}}%
\pgfpathlineto{\pgfqpoint{2.016197in}{1.920474in}}%
\pgfpathlineto{\pgfqpoint{2.016213in}{1.950027in}}%
\pgfpathlineto{\pgfqpoint{2.016953in}{1.937399in}}%
\pgfpathlineto{\pgfqpoint{2.017921in}{1.532032in}}%
\pgfpathlineto{\pgfqpoint{2.016978in}{1.939146in}}%
\pgfpathlineto{\pgfqpoint{2.018212in}{1.532038in}}%
\pgfpathlineto{\pgfqpoint{2.019843in}{1.532146in}}%
\pgfpathlineto{\pgfqpoint{2.019906in}{2.030855in}}%
\pgfpathlineto{\pgfqpoint{2.021075in}{1.945888in}}%
\pgfpathlineto{\pgfqpoint{2.021086in}{1.919140in}}%
\pgfpathlineto{\pgfqpoint{2.021098in}{1.951484in}}%
\pgfpathlineto{\pgfqpoint{2.021849in}{1.932063in}}%
\pgfpathlineto{\pgfqpoint{2.022881in}{1.532032in}}%
\pgfpathlineto{\pgfqpoint{2.021897in}{1.939418in}}%
\pgfpathlineto{\pgfqpoint{2.024160in}{1.532061in}}%
\pgfpathlineto{\pgfqpoint{2.024788in}{1.532139in}}%
\pgfpathlineto{\pgfqpoint{2.024851in}{2.031031in}}%
\pgfpathlineto{\pgfqpoint{2.026028in}{1.921203in}}%
\pgfpathlineto{\pgfqpoint{2.026032in}{1.919142in}}%
\pgfpathlineto{\pgfqpoint{2.026044in}{1.951795in}}%
\pgfpathlineto{\pgfqpoint{2.026774in}{1.933612in}}%
\pgfpathlineto{\pgfqpoint{2.026782in}{1.939771in}}%
\pgfpathlineto{\pgfqpoint{2.027371in}{1.856873in}}%
\pgfpathlineto{\pgfqpoint{2.027806in}{1.532032in}}%
\pgfpathlineto{\pgfqpoint{2.028341in}{1.532045in}}%
\pgfpathlineto{\pgfqpoint{2.029733in}{1.532130in}}%
\pgfpathlineto{\pgfqpoint{2.029833in}{2.005057in}}%
\pgfpathlineto{\pgfqpoint{2.031127in}{1.936779in}}%
\pgfpathlineto{\pgfqpoint{2.031137in}{1.925712in}}%
\pgfpathlineto{\pgfqpoint{2.031153in}{1.945124in}}%
\pgfpathlineto{\pgfqpoint{2.031892in}{1.931022in}}%
\pgfpathlineto{\pgfqpoint{2.031923in}{1.938890in}}%
\pgfpathlineto{\pgfqpoint{2.032318in}{1.859492in}}%
\pgfpathlineto{\pgfqpoint{2.032789in}{1.532033in}}%
\pgfpathlineto{\pgfqpoint{2.033279in}{1.532044in}}%
\pgfpathlineto{\pgfqpoint{2.034680in}{1.532143in}}%
\pgfpathlineto{\pgfqpoint{2.034742in}{2.030638in}}%
\pgfpathlineto{\pgfqpoint{2.035914in}{1.936418in}}%
\pgfpathlineto{\pgfqpoint{2.035922in}{1.918817in}}%
\pgfpathlineto{\pgfqpoint{2.035937in}{1.951691in}}%
\pgfpathlineto{\pgfqpoint{2.036684in}{1.933071in}}%
\pgfpathlineto{\pgfqpoint{2.037700in}{1.532032in}}%
\pgfpathlineto{\pgfqpoint{2.036704in}{1.939631in}}%
\pgfpathlineto{\pgfqpoint{2.038187in}{1.532043in}}%
\pgfpathlineto{\pgfqpoint{2.039623in}{1.532122in}}%
\pgfpathlineto{\pgfqpoint{2.039664in}{1.745908in}}%
\pgfpathlineto{\pgfqpoint{2.039726in}{1.992464in}}%
\pgfpathlineto{\pgfqpoint{2.040444in}{1.950396in}}%
\pgfpathlineto{\pgfqpoint{2.040454in}{1.914654in}}%
\pgfpathlineto{\pgfqpoint{2.040469in}{1.956023in}}%
\pgfpathlineto{\pgfqpoint{2.041212in}{1.932002in}}%
\pgfpathlineto{\pgfqpoint{2.041263in}{1.942850in}}%
\pgfpathlineto{\pgfqpoint{2.041242in}{1.928061in}}%
\pgfpathlineto{\pgfqpoint{2.041966in}{1.934386in}}%
\pgfpathlineto{\pgfqpoint{2.042016in}{1.933462in}}%
\pgfpathlineto{\pgfqpoint{2.042065in}{1.937448in}}%
\pgfpathlineto{\pgfqpoint{2.042637in}{1.532032in}}%
\pgfpathlineto{\pgfqpoint{2.043762in}{1.532058in}}%
\pgfpathlineto{\pgfqpoint{2.044569in}{1.532128in}}%
\pgfpathlineto{\pgfqpoint{2.044670in}{2.001668in}}%
\pgfpathlineto{\pgfqpoint{2.045996in}{1.943043in}}%
\pgfpathlineto{\pgfqpoint{2.046008in}{1.925989in}}%
\pgfpathlineto{\pgfqpoint{2.046024in}{1.945007in}}%
\pgfpathlineto{\pgfqpoint{2.046750in}{1.937748in}}%
\pgfpathlineto{\pgfqpoint{2.046966in}{1.938938in}}%
\pgfpathlineto{\pgfqpoint{2.047164in}{1.828701in}}%
\pgfpathlineto{\pgfqpoint{2.047637in}{1.532033in}}%
\pgfpathlineto{\pgfqpoint{2.048132in}{1.532044in}}%
\pgfpathlineto{\pgfqpoint{2.049515in}{1.532134in}}%
\pgfpathlineto{\pgfqpoint{2.049579in}{2.032451in}}%
\pgfpathlineto{\pgfqpoint{2.050782in}{1.929889in}}%
\pgfpathlineto{\pgfqpoint{2.050786in}{1.920394in}}%
\pgfpathlineto{\pgfqpoint{2.050804in}{1.951074in}}%
\pgfpathlineto{\pgfqpoint{2.051551in}{1.933722in}}%
\pgfpathlineto{\pgfqpoint{2.052538in}{1.532032in}}%
\pgfpathlineto{\pgfqpoint{2.051571in}{1.939442in}}%
\pgfpathlineto{\pgfqpoint{2.052927in}{1.532040in}}%
\pgfpathlineto{\pgfqpoint{2.054460in}{1.532123in}}%
\pgfpathlineto{\pgfqpoint{2.054501in}{1.748923in}}%
\pgfpathlineto{\pgfqpoint{2.054563in}{1.990844in}}%
\pgfpathlineto{\pgfqpoint{2.055281in}{1.952651in}}%
\pgfpathlineto{\pgfqpoint{2.055324in}{1.915536in}}%
\pgfpathlineto{\pgfqpoint{2.055337in}{1.955427in}}%
\pgfpathlineto{\pgfqpoint{2.056047in}{1.927674in}}%
\pgfpathlineto{\pgfqpoint{2.056067in}{1.943258in}}%
\pgfpathlineto{\pgfqpoint{2.056808in}{1.932202in}}%
\pgfpathlineto{\pgfqpoint{2.056858in}{1.938786in}}%
\pgfpathlineto{\pgfqpoint{2.057006in}{1.912539in}}%
\pgfpathlineto{\pgfqpoint{2.057457in}{1.532031in}}%
\pgfpathlineto{\pgfqpoint{2.058061in}{1.532047in}}%
\pgfpathlineto{\pgfqpoint{2.059404in}{1.532116in}}%
\pgfpathlineto{\pgfqpoint{2.059447in}{1.762666in}}%
\pgfpathlineto{\pgfqpoint{2.059507in}{1.990408in}}%
\pgfpathlineto{\pgfqpoint{2.060228in}{1.936515in}}%
\pgfpathlineto{\pgfqpoint{2.060236in}{1.914402in}}%
\pgfpathlineto{\pgfqpoint{2.060250in}{1.955720in}}%
\pgfpathlineto{\pgfqpoint{2.060997in}{1.931331in}}%
\pgfpathlineto{\pgfqpoint{2.061044in}{1.942686in}}%
\pgfpathlineto{\pgfqpoint{2.061067in}{1.928327in}}%
\pgfpathlineto{\pgfqpoint{2.061768in}{1.938931in}}%
\pgfpathlineto{\pgfqpoint{2.061867in}{1.939283in}}%
\pgfpathlineto{\pgfqpoint{2.062065in}{1.743565in}}%
\pgfpathlineto{\pgfqpoint{2.062410in}{1.532033in}}%
\pgfpathlineto{\pgfqpoint{2.062981in}{1.532044in}}%
\pgfpathlineto{\pgfqpoint{2.064352in}{1.532137in}}%
\pgfpathlineto{\pgfqpoint{2.064415in}{2.030648in}}%
\pgfpathlineto{\pgfqpoint{2.065613in}{1.941658in}}%
\pgfpathlineto{\pgfqpoint{2.065620in}{1.920080in}}%
\pgfpathlineto{\pgfqpoint{2.065636in}{1.951077in}}%
\pgfpathlineto{\pgfqpoint{2.066385in}{1.931825in}}%
\pgfpathlineto{\pgfqpoint{2.066402in}{1.939196in}}%
\pgfpathlineto{\pgfqpoint{2.067021in}{1.724669in}}%
\pgfpathlineto{\pgfqpoint{2.067361in}{1.532032in}}%
\pgfpathlineto{\pgfqpoint{2.067937in}{1.532044in}}%
\pgfpathlineto{\pgfqpoint{2.069297in}{1.532130in}}%
\pgfpathlineto{\pgfqpoint{2.069360in}{2.030796in}}%
\pgfpathlineto{\pgfqpoint{2.070591in}{1.932228in}}%
\pgfpathlineto{\pgfqpoint{2.070599in}{1.920650in}}%
\pgfpathlineto{\pgfqpoint{2.070611in}{1.950033in}}%
\pgfpathlineto{\pgfqpoint{2.071363in}{1.931414in}}%
\pgfpathlineto{\pgfqpoint{2.071378in}{1.938981in}}%
\pgfpathlineto{\pgfqpoint{2.071835in}{1.910916in}}%
\pgfpathlineto{\pgfqpoint{2.072319in}{1.532032in}}%
\pgfpathlineto{\pgfqpoint{2.072955in}{1.532047in}}%
\pgfpathlineto{\pgfqpoint{2.074244in}{1.532151in}}%
\pgfpathlineto{\pgfqpoint{2.074306in}{2.030632in}}%
\pgfpathlineto{\pgfqpoint{2.075451in}{1.928797in}}%
\pgfpathlineto{\pgfqpoint{2.075455in}{1.918874in}}%
\pgfpathlineto{\pgfqpoint{2.075471in}{1.952605in}}%
\pgfpathlineto{\pgfqpoint{2.076219in}{1.931485in}}%
\pgfpathlineto{\pgfqpoint{2.076235in}{1.939567in}}%
\pgfpathlineto{\pgfqpoint{2.076977in}{1.644087in}}%
\pgfpathlineto{\pgfqpoint{2.077265in}{1.532032in}}%
\pgfpathlineto{\pgfqpoint{2.077192in}{1.657175in}}%
\pgfpathlineto{\pgfqpoint{2.077901in}{1.532047in}}%
\pgfpathlineto{\pgfqpoint{2.079189in}{1.532150in}}%
\pgfpathlineto{\pgfqpoint{2.079252in}{2.030632in}}%
\pgfpathlineto{\pgfqpoint{2.080393in}{1.932288in}}%
\pgfpathlineto{\pgfqpoint{2.080401in}{1.918117in}}%
\pgfpathlineto{\pgfqpoint{2.080413in}{1.951963in}}%
\pgfpathlineto{\pgfqpoint{2.081164in}{1.931915in}}%
\pgfpathlineto{\pgfqpoint{2.082197in}{1.532032in}}%
\pgfpathlineto{\pgfqpoint{2.081183in}{1.939758in}}%
\pgfpathlineto{\pgfqpoint{2.083070in}{1.532051in}}%
\pgfpathlineto{\pgfqpoint{2.084134in}{1.532135in}}%
\pgfpathlineto{\pgfqpoint{2.084197in}{2.030666in}}%
\pgfpathlineto{\pgfqpoint{2.085398in}{1.934585in}}%
\pgfpathlineto{\pgfqpoint{2.085405in}{1.919475in}}%
\pgfpathlineto{\pgfqpoint{2.085421in}{1.950626in}}%
\pgfpathlineto{\pgfqpoint{2.086169in}{1.931696in}}%
\pgfpathlineto{\pgfqpoint{2.086185in}{1.939117in}}%
\pgfpathlineto{\pgfqpoint{2.086769in}{1.778869in}}%
\pgfpathlineto{\pgfqpoint{2.087156in}{1.532032in}}%
\pgfpathlineto{\pgfqpoint{2.087694in}{1.532045in}}%
\pgfpathlineto{\pgfqpoint{2.089080in}{1.532149in}}%
\pgfpathlineto{\pgfqpoint{2.089142in}{2.030632in}}%
\pgfpathlineto{\pgfqpoint{2.090286in}{1.923939in}}%
\pgfpathlineto{\pgfqpoint{2.090290in}{1.917798in}}%
\pgfpathlineto{\pgfqpoint{2.090305in}{1.952737in}}%
\pgfpathlineto{\pgfqpoint{2.091052in}{1.933477in}}%
\pgfpathlineto{\pgfqpoint{2.092082in}{1.532033in}}%
\pgfpathlineto{\pgfqpoint{2.091073in}{1.939776in}}%
\pgfpathlineto{\pgfqpoint{2.092406in}{1.532040in}}%
\pgfpathlineto{\pgfqpoint{2.094025in}{1.532138in}}%
\pgfpathlineto{\pgfqpoint{2.094088in}{2.030605in}}%
\pgfpathlineto{\pgfqpoint{2.095288in}{1.934664in}}%
\pgfpathlineto{\pgfqpoint{2.095295in}{1.919333in}}%
\pgfpathlineto{\pgfqpoint{2.095310in}{1.950763in}}%
\pgfpathlineto{\pgfqpoint{2.096057in}{1.932539in}}%
\pgfpathlineto{\pgfqpoint{2.097046in}{1.532032in}}%
\pgfpathlineto{\pgfqpoint{2.096076in}{1.939247in}}%
\pgfpathlineto{\pgfqpoint{2.097782in}{1.532050in}}%
\pgfpathlineto{\pgfqpoint{2.098971in}{1.532151in}}%
\pgfpathlineto{\pgfqpoint{2.099033in}{2.030587in}}%
\pgfpathlineto{\pgfqpoint{2.100176in}{1.932620in}}%
\pgfpathlineto{\pgfqpoint{2.100183in}{1.918118in}}%
\pgfpathlineto{\pgfqpoint{2.100196in}{1.951578in}}%
\pgfpathlineto{\pgfqpoint{2.100946in}{1.932006in}}%
\pgfpathlineto{\pgfqpoint{2.101974in}{1.532033in}}%
\pgfpathlineto{\pgfqpoint{2.100966in}{1.939712in}}%
\pgfpathlineto{\pgfqpoint{2.102744in}{1.532049in}}%
\pgfpathlineto{\pgfqpoint{2.103915in}{1.532137in}}%
\pgfpathlineto{\pgfqpoint{2.103978in}{2.030595in}}%
\pgfpathlineto{\pgfqpoint{2.105178in}{1.936297in}}%
\pgfpathlineto{\pgfqpoint{2.105186in}{1.919380in}}%
\pgfpathlineto{\pgfqpoint{2.105201in}{1.950954in}}%
\pgfpathlineto{\pgfqpoint{2.105951in}{1.931052in}}%
\pgfpathlineto{\pgfqpoint{2.105966in}{1.939164in}}%
\pgfpathlineto{\pgfqpoint{2.106478in}{1.886383in}}%
\pgfpathlineto{\pgfqpoint{2.106920in}{1.532033in}}%
\pgfpathlineto{\pgfqpoint{2.107493in}{1.532044in}}%
\pgfpathlineto{\pgfqpoint{2.108861in}{1.532137in}}%
\pgfpathlineto{\pgfqpoint{2.108924in}{2.030592in}}%
\pgfpathlineto{\pgfqpoint{2.110124in}{1.933804in}}%
\pgfpathlineto{\pgfqpoint{2.110132in}{1.919568in}}%
\pgfpathlineto{\pgfqpoint{2.110147in}{1.950238in}}%
\pgfpathlineto{\pgfqpoint{2.110896in}{1.931333in}}%
\pgfpathlineto{\pgfqpoint{2.110913in}{1.939142in}}%
\pgfpathlineto{\pgfqpoint{2.111378in}{1.925814in}}%
\pgfpathlineto{\pgfqpoint{2.111865in}{1.532033in}}%
\pgfpathlineto{\pgfqpoint{2.112683in}{1.532051in}}%
\pgfpathlineto{\pgfqpoint{2.113806in}{1.532137in}}%
\pgfpathlineto{\pgfqpoint{2.113869in}{2.030575in}}%
\pgfpathlineto{\pgfqpoint{2.115067in}{1.944560in}}%
\pgfpathlineto{\pgfqpoint{2.115078in}{1.919878in}}%
\pgfpathlineto{\pgfqpoint{2.115091in}{1.950753in}}%
\pgfpathlineto{\pgfqpoint{2.115840in}{1.932404in}}%
\pgfpathlineto{\pgfqpoint{2.116812in}{1.532033in}}%
\pgfpathlineto{\pgfqpoint{2.115859in}{1.939187in}}%
\pgfpathlineto{\pgfqpoint{2.117583in}{1.532049in}}%
\pgfpathlineto{\pgfqpoint{2.118752in}{1.532137in}}%
\pgfpathlineto{\pgfqpoint{2.118815in}{2.030578in}}%
\pgfpathlineto{\pgfqpoint{2.120013in}{1.941488in}}%
\pgfpathlineto{\pgfqpoint{2.120021in}{1.919976in}}%
\pgfpathlineto{\pgfqpoint{2.120036in}{1.951001in}}%
\pgfpathlineto{\pgfqpoint{2.120785in}{1.931961in}}%
\pgfpathlineto{\pgfqpoint{2.121756in}{1.532033in}}%
\pgfpathlineto{\pgfqpoint{2.120804in}{1.938998in}}%
\pgfpathlineto{\pgfqpoint{2.123118in}{1.532062in}}%
\pgfpathlineto{\pgfqpoint{2.123697in}{1.532137in}}%
\pgfpathlineto{\pgfqpoint{2.123760in}{2.030562in}}%
\pgfpathlineto{\pgfqpoint{2.124959in}{1.940279in}}%
\pgfpathlineto{\pgfqpoint{2.124966in}{1.919694in}}%
\pgfpathlineto{\pgfqpoint{2.124982in}{1.950887in}}%
\pgfpathlineto{\pgfqpoint{2.125733in}{1.931229in}}%
\pgfpathlineto{\pgfqpoint{2.125750in}{1.939157in}}%
\pgfpathlineto{\pgfqpoint{2.126216in}{1.925028in}}%
\pgfpathlineto{\pgfqpoint{2.126710in}{1.532032in}}%
\pgfpathlineto{\pgfqpoint{2.127488in}{1.532049in}}%
\pgfpathlineto{\pgfqpoint{2.128643in}{1.532147in}}%
\pgfpathlineto{\pgfqpoint{2.128706in}{2.030605in}}%
\pgfpathlineto{\pgfqpoint{2.129851in}{1.922993in}}%
\pgfpathlineto{\pgfqpoint{2.129855in}{1.917755in}}%
\pgfpathlineto{\pgfqpoint{2.129871in}{1.952285in}}%
\pgfpathlineto{\pgfqpoint{2.130609in}{1.939683in}}%
\pgfpathlineto{\pgfqpoint{2.131655in}{1.532032in}}%
\pgfpathlineto{\pgfqpoint{2.131891in}{1.532038in}}%
\pgfpathlineto{\pgfqpoint{2.133589in}{1.532147in}}%
\pgfpathlineto{\pgfqpoint{2.133651in}{2.030607in}}%
\pgfpathlineto{\pgfqpoint{2.134799in}{1.920762in}}%
\pgfpathlineto{\pgfqpoint{2.134802in}{1.918064in}}%
\pgfpathlineto{\pgfqpoint{2.134815in}{1.951840in}}%
\pgfpathlineto{\pgfqpoint{2.135546in}{1.934081in}}%
\pgfpathlineto{\pgfqpoint{2.135553in}{1.939943in}}%
\pgfpathlineto{\pgfqpoint{2.136174in}{1.857192in}}%
\pgfpathlineto{\pgfqpoint{2.136601in}{1.532032in}}%
\pgfpathlineto{\pgfqpoint{2.137133in}{1.532045in}}%
\pgfpathlineto{\pgfqpoint{2.138534in}{1.532146in}}%
\pgfpathlineto{\pgfqpoint{2.138597in}{2.030607in}}%
\pgfpathlineto{\pgfqpoint{2.139765in}{1.947438in}}%
\pgfpathlineto{\pgfqpoint{2.139776in}{1.918492in}}%
\pgfpathlineto{\pgfqpoint{2.139791in}{1.951485in}}%
\pgfpathlineto{\pgfqpoint{2.140540in}{1.931748in}}%
\pgfpathlineto{\pgfqpoint{2.141539in}{1.532033in}}%
\pgfpathlineto{\pgfqpoint{2.140588in}{1.939166in}}%
\pgfpathlineto{\pgfqpoint{2.142903in}{1.532062in}}%
\pgfpathlineto{\pgfqpoint{2.143480in}{1.532146in}}%
\pgfpathlineto{\pgfqpoint{2.143542in}{2.030452in}}%
\pgfpathlineto{\pgfqpoint{2.144710in}{1.946835in}}%
\pgfpathlineto{\pgfqpoint{2.144721in}{1.918566in}}%
\pgfpathlineto{\pgfqpoint{2.144736in}{1.951111in}}%
\pgfpathlineto{\pgfqpoint{2.145482in}{1.932305in}}%
\pgfpathlineto{\pgfqpoint{2.146483in}{1.532033in}}%
\pgfpathlineto{\pgfqpoint{2.145501in}{1.939432in}}%
\pgfpathlineto{\pgfqpoint{2.147154in}{1.532047in}}%
\pgfpathlineto{\pgfqpoint{2.148425in}{1.532147in}}%
\pgfpathlineto{\pgfqpoint{2.148488in}{2.030441in}}%
\pgfpathlineto{\pgfqpoint{2.149633in}{1.921717in}}%
\pgfpathlineto{\pgfqpoint{2.149636in}{1.917763in}}%
\pgfpathlineto{\pgfqpoint{2.149680in}{1.951783in}}%
\pgfpathlineto{\pgfqpoint{2.150385in}{1.938482in}}%
\pgfpathlineto{\pgfqpoint{2.150389in}{1.940015in}}%
\pgfpathlineto{\pgfqpoint{2.150981in}{1.893599in}}%
\pgfpathlineto{\pgfqpoint{2.151481in}{1.532033in}}%
\pgfpathlineto{\pgfqpoint{2.152020in}{1.532046in}}%
\pgfpathlineto{\pgfqpoint{2.153370in}{1.532143in}}%
\pgfpathlineto{\pgfqpoint{2.153433in}{2.030457in}}%
\pgfpathlineto{\pgfqpoint{2.154605in}{1.933236in}}%
\pgfpathlineto{\pgfqpoint{2.154613in}{1.918651in}}%
\pgfpathlineto{\pgfqpoint{2.154656in}{1.950960in}}%
\pgfpathlineto{\pgfqpoint{2.155374in}{1.931851in}}%
\pgfpathlineto{\pgfqpoint{2.156374in}{1.532033in}}%
\pgfpathlineto{\pgfqpoint{2.155394in}{1.939268in}}%
\pgfpathlineto{\pgfqpoint{2.157440in}{1.532055in}}%
\pgfpathlineto{\pgfqpoint{2.158316in}{1.532147in}}%
\pgfpathlineto{\pgfqpoint{2.158379in}{2.030423in}}%
\pgfpathlineto{\pgfqpoint{2.159522in}{1.926173in}}%
\pgfpathlineto{\pgfqpoint{2.159526in}{1.918025in}}%
\pgfpathlineto{\pgfqpoint{2.159541in}{1.952386in}}%
\pgfpathlineto{\pgfqpoint{2.160288in}{1.931733in}}%
\pgfpathlineto{\pgfqpoint{2.161375in}{1.532033in}}%
\pgfpathlineto{\pgfqpoint{2.160308in}{1.939558in}}%
\pgfpathlineto{\pgfqpoint{2.162262in}{1.532053in}}%
\pgfpathlineto{\pgfqpoint{2.163261in}{1.532142in}}%
\pgfpathlineto{\pgfqpoint{2.163324in}{2.030455in}}%
\pgfpathlineto{\pgfqpoint{2.164498in}{1.933202in}}%
\pgfpathlineto{\pgfqpoint{2.164505in}{1.918554in}}%
\pgfpathlineto{\pgfqpoint{2.164549in}{1.950896in}}%
\pgfpathlineto{\pgfqpoint{2.165269in}{1.930928in}}%
\pgfpathlineto{\pgfqpoint{2.165287in}{1.939286in}}%
\pgfpathlineto{\pgfqpoint{2.165777in}{1.926715in}}%
\pgfpathlineto{\pgfqpoint{2.166316in}{1.532033in}}%
\pgfpathlineto{\pgfqpoint{2.167153in}{1.532053in}}%
\pgfpathlineto{\pgfqpoint{2.168207in}{1.532143in}}%
\pgfpathlineto{\pgfqpoint{2.168270in}{2.030429in}}%
\pgfpathlineto{\pgfqpoint{2.169442in}{1.931499in}}%
\pgfpathlineto{\pgfqpoint{2.169450in}{1.918978in}}%
\pgfpathlineto{\pgfqpoint{2.169462in}{1.951336in}}%
\pgfpathlineto{\pgfqpoint{2.170212in}{1.931866in}}%
\pgfpathlineto{\pgfqpoint{2.171211in}{1.532033in}}%
\pgfpathlineto{\pgfqpoint{2.170232in}{1.939246in}}%
\pgfpathlineto{\pgfqpoint{2.172079in}{1.532051in}}%
\pgfpathlineto{\pgfqpoint{2.173153in}{1.532147in}}%
\pgfpathlineto{\pgfqpoint{2.173215in}{2.030401in}}%
\pgfpathlineto{\pgfqpoint{2.174382in}{1.945621in}}%
\pgfpathlineto{\pgfqpoint{2.174393in}{1.918767in}}%
\pgfpathlineto{\pgfqpoint{2.174406in}{1.951030in}}%
\pgfpathlineto{\pgfqpoint{2.175155in}{1.931459in}}%
\pgfpathlineto{\pgfqpoint{2.175171in}{1.939134in}}%
\pgfpathlineto{\pgfqpoint{2.175862in}{1.672356in}}%
\pgfpathlineto{\pgfqpoint{2.176202in}{1.532033in}}%
\pgfpathlineto{\pgfqpoint{2.176741in}{1.532046in}}%
\pgfpathlineto{\pgfqpoint{2.178098in}{1.532145in}}%
\pgfpathlineto{\pgfqpoint{2.178161in}{2.030401in}}%
\pgfpathlineto{\pgfqpoint{2.179329in}{1.941016in}}%
\pgfpathlineto{\pgfqpoint{2.179337in}{1.918874in}}%
\pgfpathlineto{\pgfqpoint{2.179353in}{1.951590in}}%
\pgfpathlineto{\pgfqpoint{2.180101in}{1.932671in}}%
\pgfpathlineto{\pgfqpoint{2.181158in}{1.532033in}}%
\pgfpathlineto{\pgfqpoint{2.180120in}{1.939352in}}%
\pgfpathlineto{\pgfqpoint{2.181649in}{1.532044in}}%
\pgfpathlineto{\pgfqpoint{2.183043in}{1.532142in}}%
\pgfpathlineto{\pgfqpoint{2.183106in}{2.030422in}}%
\pgfpathlineto{\pgfqpoint{2.184277in}{1.940751in}}%
\pgfpathlineto{\pgfqpoint{2.184285in}{1.918848in}}%
\pgfpathlineto{\pgfqpoint{2.184300in}{1.951524in}}%
\pgfpathlineto{\pgfqpoint{2.185050in}{1.931258in}}%
\pgfpathlineto{\pgfqpoint{2.185066in}{1.939180in}}%
\pgfpathlineto{\pgfqpoint{2.185643in}{1.837288in}}%
\pgfpathlineto{\pgfqpoint{2.186041in}{1.532033in}}%
\pgfpathlineto{\pgfqpoint{2.186606in}{1.532044in}}%
\pgfpathlineto{\pgfqpoint{2.187989in}{1.532149in}}%
\pgfpathlineto{\pgfqpoint{2.188052in}{2.030357in}}%
\pgfpathlineto{\pgfqpoint{2.189194in}{1.926922in}}%
\pgfpathlineto{\pgfqpoint{2.189198in}{1.918173in}}%
\pgfpathlineto{\pgfqpoint{2.189214in}{1.952333in}}%
\pgfpathlineto{\pgfqpoint{2.189964in}{1.930416in}}%
\pgfpathlineto{\pgfqpoint{2.189979in}{1.939584in}}%
\pgfpathlineto{\pgfqpoint{2.190544in}{1.894977in}}%
\pgfpathlineto{\pgfqpoint{2.191045in}{1.532033in}}%
\pgfpathlineto{\pgfqpoint{2.191585in}{1.532046in}}%
\pgfpathlineto{\pgfqpoint{2.192934in}{1.532143in}}%
\pgfpathlineto{\pgfqpoint{2.192997in}{2.030405in}}%
\pgfpathlineto{\pgfqpoint{2.194166in}{1.941442in}}%
\pgfpathlineto{\pgfqpoint{2.194174in}{1.919208in}}%
\pgfpathlineto{\pgfqpoint{2.194190in}{1.951490in}}%
\pgfpathlineto{\pgfqpoint{2.194939in}{1.931450in}}%
\pgfpathlineto{\pgfqpoint{2.194955in}{1.939014in}}%
\pgfpathlineto{\pgfqpoint{2.195640in}{1.677783in}}%
\pgfpathlineto{\pgfqpoint{2.195999in}{1.532033in}}%
\pgfpathlineto{\pgfqpoint{2.196543in}{1.532046in}}%
\pgfpathlineto{\pgfqpoint{2.197880in}{1.532145in}}%
\pgfpathlineto{\pgfqpoint{2.197942in}{2.030361in}}%
\pgfpathlineto{\pgfqpoint{2.199114in}{1.936603in}}%
\pgfpathlineto{\pgfqpoint{2.199122in}{1.918472in}}%
\pgfpathlineto{\pgfqpoint{2.199138in}{1.951269in}}%
\pgfpathlineto{\pgfqpoint{2.199886in}{1.931776in}}%
\pgfpathlineto{\pgfqpoint{2.200933in}{1.532033in}}%
\pgfpathlineto{\pgfqpoint{2.199906in}{1.939113in}}%
\pgfpathlineto{\pgfqpoint{2.201917in}{1.532055in}}%
\pgfpathlineto{\pgfqpoint{2.202825in}{1.532144in}}%
\pgfpathlineto{\pgfqpoint{2.202888in}{2.030370in}}%
\pgfpathlineto{\pgfqpoint{2.204058in}{1.943262in}}%
\pgfpathlineto{\pgfqpoint{2.204098in}{1.919454in}}%
\pgfpathlineto{\pgfqpoint{2.204083in}{1.951676in}}%
\pgfpathlineto{\pgfqpoint{2.204833in}{1.931140in}}%
\pgfpathlineto{\pgfqpoint{2.204849in}{1.939199in}}%
\pgfpathlineto{\pgfqpoint{2.205470in}{1.766193in}}%
\pgfpathlineto{\pgfqpoint{2.205882in}{1.532033in}}%
\pgfpathlineto{\pgfqpoint{2.206372in}{1.532044in}}%
\pgfpathlineto{\pgfqpoint{2.207770in}{1.532143in}}%
\pgfpathlineto{\pgfqpoint{2.207833in}{2.030379in}}%
\pgfpathlineto{\pgfqpoint{2.209006in}{1.929059in}}%
\pgfpathlineto{\pgfqpoint{2.209010in}{1.919322in}}%
\pgfpathlineto{\pgfqpoint{2.209027in}{1.951484in}}%
\pgfpathlineto{\pgfqpoint{2.209775in}{1.932255in}}%
\pgfpathlineto{\pgfqpoint{2.210837in}{1.532033in}}%
\pgfpathlineto{\pgfqpoint{2.209794in}{1.939327in}}%
\pgfpathlineto{\pgfqpoint{2.211431in}{1.532046in}}%
\pgfpathlineto{\pgfqpoint{2.212716in}{1.532145in}}%
\pgfpathlineto{\pgfqpoint{2.212779in}{2.030342in}}%
\pgfpathlineto{\pgfqpoint{2.213949in}{1.940415in}}%
\pgfpathlineto{\pgfqpoint{2.213956in}{1.918966in}}%
\pgfpathlineto{\pgfqpoint{2.213972in}{1.951467in}}%
\pgfpathlineto{\pgfqpoint{2.214720in}{1.931782in}}%
\pgfpathlineto{\pgfqpoint{2.215783in}{1.532033in}}%
\pgfpathlineto{\pgfqpoint{2.214740in}{1.939077in}}%
\pgfpathlineto{\pgfqpoint{2.216821in}{1.532057in}}%
\pgfpathlineto{\pgfqpoint{2.217661in}{1.532145in}}%
\pgfpathlineto{\pgfqpoint{2.217724in}{2.030339in}}%
\pgfpathlineto{\pgfqpoint{2.218895in}{1.934483in}}%
\pgfpathlineto{\pgfqpoint{2.218903in}{1.918652in}}%
\pgfpathlineto{\pgfqpoint{2.218946in}{1.950783in}}%
\pgfpathlineto{\pgfqpoint{2.219666in}{1.931738in}}%
\pgfpathlineto{\pgfqpoint{2.220726in}{1.532033in}}%
\pgfpathlineto{\pgfqpoint{2.219686in}{1.939072in}}%
\pgfpathlineto{\pgfqpoint{2.221616in}{1.532053in}}%
\pgfpathlineto{\pgfqpoint{2.222607in}{1.532146in}}%
\pgfpathlineto{\pgfqpoint{2.222670in}{2.030322in}}%
\pgfpathlineto{\pgfqpoint{2.223839in}{1.939683in}}%
\pgfpathlineto{\pgfqpoint{2.223847in}{1.918666in}}%
\pgfpathlineto{\pgfqpoint{2.223862in}{1.951471in}}%
\pgfpathlineto{\pgfqpoint{2.224612in}{1.930805in}}%
\pgfpathlineto{\pgfqpoint{2.224630in}{1.939182in}}%
\pgfpathlineto{\pgfqpoint{2.225133in}{1.921491in}}%
\pgfpathlineto{\pgfqpoint{2.225672in}{1.532033in}}%
\pgfpathlineto{\pgfqpoint{2.226365in}{1.532048in}}%
\pgfpathlineto{\pgfqpoint{2.227552in}{1.532145in}}%
\pgfpathlineto{\pgfqpoint{2.227615in}{2.030325in}}%
\pgfpathlineto{\pgfqpoint{2.228787in}{1.935926in}}%
\pgfpathlineto{\pgfqpoint{2.228795in}{1.918378in}}%
\pgfpathlineto{\pgfqpoint{2.228810in}{1.951343in}}%
\pgfpathlineto{\pgfqpoint{2.229556in}{1.932892in}}%
\pgfpathlineto{\pgfqpoint{2.230616in}{1.532033in}}%
\pgfpathlineto{\pgfqpoint{2.229576in}{1.939242in}}%
\pgfpathlineto{\pgfqpoint{2.231061in}{1.532044in}}%
\pgfpathlineto{\pgfqpoint{2.232498in}{1.532146in}}%
\pgfpathlineto{\pgfqpoint{2.232561in}{2.030307in}}%
\pgfpathlineto{\pgfqpoint{2.233726in}{1.950489in}}%
\pgfpathlineto{\pgfqpoint{2.233739in}{1.918549in}}%
\pgfpathlineto{\pgfqpoint{2.233755in}{1.950904in}}%
\pgfpathlineto{\pgfqpoint{2.234502in}{1.931335in}}%
\pgfpathlineto{\pgfqpoint{2.234518in}{1.938994in}}%
\pgfpathlineto{\pgfqpoint{2.235200in}{1.682699in}}%
\pgfpathlineto{\pgfqpoint{2.235527in}{1.532032in}}%
\pgfpathlineto{\pgfqpoint{2.236065in}{1.532045in}}%
\pgfpathlineto{\pgfqpoint{2.237444in}{1.532161in}}%
\pgfpathlineto{\pgfqpoint{2.237507in}{2.030651in}}%
\pgfpathlineto{\pgfqpoint{2.238622in}{1.925437in}}%
\pgfpathlineto{\pgfqpoint{2.238626in}{1.916943in}}%
\pgfpathlineto{\pgfqpoint{2.238641in}{1.953149in}}%
\pgfpathlineto{\pgfqpoint{2.239388in}{1.933572in}}%
\pgfpathlineto{\pgfqpoint{2.240475in}{1.532032in}}%
\pgfpathlineto{\pgfqpoint{2.239408in}{1.939753in}}%
\pgfpathlineto{\pgfqpoint{2.240617in}{1.532037in}}%
\pgfpathlineto{\pgfqpoint{2.242390in}{1.532160in}}%
\pgfpathlineto{\pgfqpoint{2.242452in}{2.030552in}}%
\pgfpathlineto{\pgfqpoint{2.243567in}{1.926090in}}%
\pgfpathlineto{\pgfqpoint{2.243571in}{1.917069in}}%
\pgfpathlineto{\pgfqpoint{2.243586in}{1.953098in}}%
\pgfpathlineto{\pgfqpoint{2.244337in}{1.930453in}}%
\pgfpathlineto{\pgfqpoint{2.244354in}{1.939814in}}%
\pgfpathlineto{\pgfqpoint{2.244946in}{1.894297in}}%
\pgfpathlineto{\pgfqpoint{2.245409in}{1.532032in}}%
\pgfpathlineto{\pgfqpoint{2.245947in}{1.532045in}}%
\pgfpathlineto{\pgfqpoint{2.247335in}{1.532165in}}%
\pgfpathlineto{\pgfqpoint{2.247398in}{2.030924in}}%
\pgfpathlineto{\pgfqpoint{2.248508in}{1.941293in}}%
\pgfpathlineto{\pgfqpoint{2.248516in}{1.917175in}}%
\pgfpathlineto{\pgfqpoint{2.248532in}{1.953197in}}%
\pgfpathlineto{\pgfqpoint{2.249281in}{1.931132in}}%
\pgfpathlineto{\pgfqpoint{2.250365in}{1.532032in}}%
\pgfpathlineto{\pgfqpoint{2.249328in}{1.939524in}}%
\pgfpathlineto{\pgfqpoint{2.251546in}{1.532059in}}%
\pgfpathlineto{\pgfqpoint{2.252280in}{1.532160in}}%
\pgfpathlineto{\pgfqpoint{2.252343in}{2.030542in}}%
\pgfpathlineto{\pgfqpoint{2.253456in}{1.932073in}}%
\pgfpathlineto{\pgfqpoint{2.253464in}{1.916973in}}%
\pgfpathlineto{\pgfqpoint{2.253476in}{1.952112in}}%
\pgfpathlineto{\pgfqpoint{2.254225in}{1.933021in}}%
\pgfpathlineto{\pgfqpoint{2.255299in}{1.532032in}}%
\pgfpathlineto{\pgfqpoint{2.254245in}{1.939814in}}%
\pgfpathlineto{\pgfqpoint{2.255539in}{1.532038in}}%
\pgfpathlineto{\pgfqpoint{2.257226in}{1.532171in}}%
\pgfpathlineto{\pgfqpoint{2.257289in}{2.031408in}}%
\pgfpathlineto{\pgfqpoint{2.258396in}{1.948694in}}%
\pgfpathlineto{\pgfqpoint{2.258407in}{1.916480in}}%
\pgfpathlineto{\pgfqpoint{2.258422in}{1.952994in}}%
\pgfpathlineto{\pgfqpoint{2.259170in}{1.931073in}}%
\pgfpathlineto{\pgfqpoint{2.260264in}{1.532033in}}%
\pgfpathlineto{\pgfqpoint{2.259189in}{1.939702in}}%
\pgfpathlineto{\pgfqpoint{2.261198in}{1.532054in}}%
\pgfpathlineto{\pgfqpoint{2.262172in}{1.532173in}}%
\pgfpathlineto{\pgfqpoint{2.262235in}{2.031680in}}%
\pgfpathlineto{\pgfqpoint{2.263319in}{1.920385in}}%
\pgfpathlineto{\pgfqpoint{2.263323in}{1.915741in}}%
\pgfpathlineto{\pgfqpoint{2.263339in}{1.953413in}}%
\pgfpathlineto{\pgfqpoint{2.264072in}{1.938819in}}%
\pgfpathlineto{\pgfqpoint{2.264075in}{1.940409in}}%
\pgfpathlineto{\pgfqpoint{2.264717in}{1.905989in}}%
\pgfpathlineto{\pgfqpoint{2.265192in}{1.532032in}}%
\pgfpathlineto{\pgfqpoint{2.265780in}{1.532045in}}%
\pgfpathlineto{\pgfqpoint{2.267117in}{1.532168in}}%
\pgfpathlineto{\pgfqpoint{2.267180in}{2.031208in}}%
\pgfpathlineto{\pgfqpoint{2.268289in}{1.943403in}}%
\pgfpathlineto{\pgfqpoint{2.268297in}{1.917841in}}%
\pgfpathlineto{\pgfqpoint{2.268315in}{1.953053in}}%
\pgfpathlineto{\pgfqpoint{2.269063in}{1.931266in}}%
\pgfpathlineto{\pgfqpoint{2.270152in}{1.532033in}}%
\pgfpathlineto{\pgfqpoint{2.269083in}{1.939497in}}%
\pgfpathlineto{\pgfqpoint{2.270987in}{1.532052in}}%
\pgfpathlineto{\pgfqpoint{2.272063in}{1.532174in}}%
\pgfpathlineto{\pgfqpoint{2.272126in}{2.031713in}}%
\pgfpathlineto{\pgfqpoint{2.273209in}{1.927306in}}%
\pgfpathlineto{\pgfqpoint{2.273213in}{1.916524in}}%
\pgfpathlineto{\pgfqpoint{2.273230in}{1.954170in}}%
\pgfpathlineto{\pgfqpoint{2.273977in}{1.932818in}}%
\pgfpathlineto{\pgfqpoint{2.275081in}{1.532032in}}%
\pgfpathlineto{\pgfqpoint{2.273997in}{1.940132in}}%
\pgfpathlineto{\pgfqpoint{2.275273in}{1.532036in}}%
\pgfpathlineto{\pgfqpoint{2.277008in}{1.532170in}}%
\pgfpathlineto{\pgfqpoint{2.277071in}{2.031361in}}%
\pgfpathlineto{\pgfqpoint{2.278177in}{1.947611in}}%
\pgfpathlineto{\pgfqpoint{2.278189in}{1.916739in}}%
\pgfpathlineto{\pgfqpoint{2.278204in}{1.952288in}}%
\pgfpathlineto{\pgfqpoint{2.278953in}{1.930602in}}%
\pgfpathlineto{\pgfqpoint{2.278969in}{1.939636in}}%
\pgfpathlineto{\pgfqpoint{2.279652in}{1.768229in}}%
\pgfpathlineto{\pgfqpoint{2.280042in}{1.532033in}}%
\pgfpathlineto{\pgfqpoint{2.280580in}{1.532045in}}%
\pgfpathlineto{\pgfqpoint{2.281954in}{1.532174in}}%
\pgfpathlineto{\pgfqpoint{2.282017in}{2.031701in}}%
\pgfpathlineto{\pgfqpoint{2.283103in}{1.919631in}}%
\pgfpathlineto{\pgfqpoint{2.283107in}{1.915887in}}%
\pgfpathlineto{\pgfqpoint{2.283119in}{1.952810in}}%
\pgfpathlineto{\pgfqpoint{2.283852in}{1.935138in}}%
\pgfpathlineto{\pgfqpoint{2.283860in}{1.940267in}}%
\pgfpathlineto{\pgfqpoint{2.284508in}{1.897183in}}%
\pgfpathlineto{\pgfqpoint{2.284987in}{1.532032in}}%
\pgfpathlineto{\pgfqpoint{2.285525in}{1.532045in}}%
\pgfpathlineto{\pgfqpoint{2.286899in}{1.532174in}}%
\pgfpathlineto{\pgfqpoint{2.286962in}{2.031700in}}%
\pgfpathlineto{\pgfqpoint{2.288049in}{1.919525in}}%
\pgfpathlineto{\pgfqpoint{2.288052in}{1.915971in}}%
\pgfpathlineto{\pgfqpoint{2.288065in}{1.952766in}}%
\pgfpathlineto{\pgfqpoint{2.288797in}{1.934837in}}%
\pgfpathlineto{\pgfqpoint{2.288804in}{1.940326in}}%
\pgfpathlineto{\pgfqpoint{2.289476in}{1.869488in}}%
\pgfpathlineto{\pgfqpoint{2.289934in}{1.532033in}}%
\pgfpathlineto{\pgfqpoint{2.290472in}{1.532045in}}%
\pgfpathlineto{\pgfqpoint{2.291845in}{1.532173in}}%
\pgfpathlineto{\pgfqpoint{2.291908in}{2.031663in}}%
\pgfpathlineto{\pgfqpoint{2.292993in}{1.919876in}}%
\pgfpathlineto{\pgfqpoint{2.292997in}{1.915746in}}%
\pgfpathlineto{\pgfqpoint{2.293040in}{1.953269in}}%
\pgfpathlineto{\pgfqpoint{2.293744in}{1.937551in}}%
\pgfpathlineto{\pgfqpoint{2.293747in}{1.940263in}}%
\pgfpathlineto{\pgfqpoint{2.294406in}{1.889275in}}%
\pgfpathlineto{\pgfqpoint{2.294855in}{1.532032in}}%
\pgfpathlineto{\pgfqpoint{2.295436in}{1.532045in}}%
\pgfpathlineto{\pgfqpoint{2.296790in}{1.532168in}}%
\pgfpathlineto{\pgfqpoint{2.296853in}{2.031174in}}%
\pgfpathlineto{\pgfqpoint{2.297962in}{1.936544in}}%
\pgfpathlineto{\pgfqpoint{2.297969in}{1.916540in}}%
\pgfpathlineto{\pgfqpoint{2.297984in}{1.953224in}}%
\pgfpathlineto{\pgfqpoint{2.298733in}{1.930735in}}%
\pgfpathlineto{\pgfqpoint{2.298749in}{1.939521in}}%
\pgfpathlineto{\pgfqpoint{2.299527in}{1.643513in}}%
\pgfpathlineto{\pgfqpoint{2.299821in}{1.532032in}}%
\pgfpathlineto{\pgfqpoint{2.299737in}{1.653434in}}%
\pgfpathlineto{\pgfqpoint{2.300409in}{1.532045in}}%
\pgfpathlineto{\pgfqpoint{2.301736in}{1.532174in}}%
\pgfpathlineto{\pgfqpoint{2.301799in}{2.031779in}}%
\pgfpathlineto{\pgfqpoint{2.302884in}{1.920851in}}%
\pgfpathlineto{\pgfqpoint{2.302887in}{1.915634in}}%
\pgfpathlineto{\pgfqpoint{2.302903in}{1.953360in}}%
\pgfpathlineto{\pgfqpoint{2.303641in}{1.940153in}}%
\pgfpathlineto{\pgfqpoint{2.304741in}{1.532033in}}%
\pgfpathlineto{\pgfqpoint{2.304920in}{1.532036in}}%
\pgfpathlineto{\pgfqpoint{2.306680in}{1.532153in}}%
\pgfpathlineto{\pgfqpoint{2.306743in}{2.030130in}}%
\pgfpathlineto{\pgfqpoint{2.307882in}{1.942065in}}%
\pgfpathlineto{\pgfqpoint{2.307890in}{1.918280in}}%
\pgfpathlineto{\pgfqpoint{2.307906in}{1.952145in}}%
\pgfpathlineto{\pgfqpoint{2.308657in}{1.930441in}}%
\pgfpathlineto{\pgfqpoint{2.308674in}{1.939297in}}%
\pgfpathlineto{\pgfqpoint{2.309207in}{1.920799in}}%
\pgfpathlineto{\pgfqpoint{2.309715in}{1.532033in}}%
\pgfpathlineto{\pgfqpoint{2.310402in}{1.532048in}}%
\pgfpathlineto{\pgfqpoint{2.311626in}{1.532173in}}%
\pgfpathlineto{\pgfqpoint{2.311690in}{2.031648in}}%
\pgfpathlineto{\pgfqpoint{2.312774in}{1.921156in}}%
\pgfpathlineto{\pgfqpoint{2.312778in}{1.915650in}}%
\pgfpathlineto{\pgfqpoint{2.312794in}{1.953545in}}%
\pgfpathlineto{\pgfqpoint{2.313528in}{1.939280in}}%
\pgfpathlineto{\pgfqpoint{2.313533in}{1.940191in}}%
\pgfpathlineto{\pgfqpoint{2.314151in}{1.922485in}}%
\pgfpathlineto{\pgfqpoint{2.314634in}{1.532033in}}%
\pgfpathlineto{\pgfqpoint{2.315408in}{1.532049in}}%
\pgfpathlineto{\pgfqpoint{2.316571in}{1.532152in}}%
\pgfpathlineto{\pgfqpoint{2.316634in}{2.030113in}}%
\pgfpathlineto{\pgfqpoint{2.317776in}{1.932226in}}%
\pgfpathlineto{\pgfqpoint{2.317783in}{1.917791in}}%
\pgfpathlineto{\pgfqpoint{2.317796in}{1.951083in}}%
\pgfpathlineto{\pgfqpoint{2.318547in}{1.931049in}}%
\pgfpathlineto{\pgfqpoint{2.318563in}{1.939036in}}%
\pgfpathlineto{\pgfqpoint{2.319526in}{1.565041in}}%
\pgfpathlineto{\pgfqpoint{2.319605in}{1.532033in}}%
\pgfpathlineto{\pgfqpoint{2.320292in}{1.532048in}}%
\pgfpathlineto{\pgfqpoint{2.321517in}{1.532174in}}%
\pgfpathlineto{\pgfqpoint{2.321581in}{2.031647in}}%
\pgfpathlineto{\pgfqpoint{2.322668in}{1.916571in}}%
\pgfpathlineto{\pgfqpoint{2.322684in}{1.954019in}}%
\pgfpathlineto{\pgfqpoint{2.323491in}{1.932397in}}%
\pgfpathlineto{\pgfqpoint{2.324521in}{1.532033in}}%
\pgfpathlineto{\pgfqpoint{2.323511in}{1.939435in}}%
\pgfpathlineto{\pgfqpoint{2.324895in}{1.532040in}}%
\pgfpathlineto{\pgfqpoint{2.326462in}{1.532154in}}%
\pgfpathlineto{\pgfqpoint{2.326525in}{2.030117in}}%
\pgfpathlineto{\pgfqpoint{2.327664in}{1.944824in}}%
\pgfpathlineto{\pgfqpoint{2.327675in}{1.918142in}}%
\pgfpathlineto{\pgfqpoint{2.327688in}{1.951837in}}%
\pgfpathlineto{\pgfqpoint{2.328438in}{1.930769in}}%
\pgfpathlineto{\pgfqpoint{2.328454in}{1.939241in}}%
\pgfpathlineto{\pgfqpoint{2.329077in}{1.816038in}}%
\pgfpathlineto{\pgfqpoint{2.329474in}{1.532032in}}%
\pgfpathlineto{\pgfqpoint{2.330007in}{1.532045in}}%
\pgfpathlineto{\pgfqpoint{2.331408in}{1.532168in}}%
\pgfpathlineto{\pgfqpoint{2.331471in}{2.031066in}}%
\pgfpathlineto{\pgfqpoint{2.332582in}{1.941825in}}%
\pgfpathlineto{\pgfqpoint{2.332590in}{1.916868in}}%
\pgfpathlineto{\pgfqpoint{2.332605in}{1.953118in}}%
\pgfpathlineto{\pgfqpoint{2.333355in}{1.931566in}}%
\pgfpathlineto{\pgfqpoint{2.334475in}{1.532033in}}%
\pgfpathlineto{\pgfqpoint{2.333373in}{1.939710in}}%
\pgfpathlineto{\pgfqpoint{2.335019in}{1.532046in}}%
\pgfpathlineto{\pgfqpoint{2.336353in}{1.532155in}}%
\pgfpathlineto{\pgfqpoint{2.336416in}{2.030110in}}%
\pgfpathlineto{\pgfqpoint{2.337554in}{1.944353in}}%
\pgfpathlineto{\pgfqpoint{2.337594in}{1.918263in}}%
\pgfpathlineto{\pgfqpoint{2.337579in}{1.952184in}}%
\pgfpathlineto{\pgfqpoint{2.338327in}{1.930888in}}%
\pgfpathlineto{\pgfqpoint{2.338343in}{1.939137in}}%
\pgfpathlineto{\pgfqpoint{2.339024in}{1.727282in}}%
\pgfpathlineto{\pgfqpoint{2.339366in}{1.532032in}}%
\pgfpathlineto{\pgfqpoint{2.339900in}{1.532045in}}%
\pgfpathlineto{\pgfqpoint{2.341299in}{1.532167in}}%
\pgfpathlineto{\pgfqpoint{2.341362in}{2.030982in}}%
\pgfpathlineto{\pgfqpoint{2.342474in}{1.935756in}}%
\pgfpathlineto{\pgfqpoint{2.342482in}{1.916473in}}%
\pgfpathlineto{\pgfqpoint{2.342497in}{1.952736in}}%
\pgfpathlineto{\pgfqpoint{2.343245in}{1.931341in}}%
\pgfpathlineto{\pgfqpoint{2.344315in}{1.532032in}}%
\pgfpathlineto{\pgfqpoint{2.343265in}{1.939551in}}%
\pgfpathlineto{\pgfqpoint{2.344950in}{1.532047in}}%
\pgfpathlineto{\pgfqpoint{2.346245in}{1.532176in}}%
\pgfpathlineto{\pgfqpoint{2.346310in}{2.033591in}}%
\pgfpathlineto{\pgfqpoint{2.347391in}{1.933110in}}%
\pgfpathlineto{\pgfqpoint{2.347399in}{1.915114in}}%
\pgfpathlineto{\pgfqpoint{2.347414in}{1.953878in}}%
\pgfpathlineto{\pgfqpoint{2.348163in}{1.929865in}}%
\pgfpathlineto{\pgfqpoint{2.348180in}{1.940114in}}%
\pgfpathlineto{\pgfqpoint{2.348794in}{1.902127in}}%
\pgfpathlineto{\pgfqpoint{2.349243in}{1.532033in}}%
\pgfpathlineto{\pgfqpoint{2.349859in}{1.532046in}}%
\pgfpathlineto{\pgfqpoint{2.351189in}{1.532154in}}%
\pgfpathlineto{\pgfqpoint{2.351252in}{2.030078in}}%
\pgfpathlineto{\pgfqpoint{2.352391in}{1.942245in}}%
\pgfpathlineto{\pgfqpoint{2.352399in}{1.918187in}}%
\pgfpathlineto{\pgfqpoint{2.352415in}{1.952167in}}%
\pgfpathlineto{\pgfqpoint{2.353163in}{1.932267in}}%
\pgfpathlineto{\pgfqpoint{2.354201in}{1.532032in}}%
\pgfpathlineto{\pgfqpoint{2.353182in}{1.939308in}}%
\pgfpathlineto{\pgfqpoint{2.354633in}{1.532043in}}%
\pgfpathlineto{\pgfqpoint{2.356135in}{1.532168in}}%
\pgfpathlineto{\pgfqpoint{2.356198in}{2.031083in}}%
\pgfpathlineto{\pgfqpoint{2.357308in}{1.944133in}}%
\pgfpathlineto{\pgfqpoint{2.357348in}{1.917534in}}%
\pgfpathlineto{\pgfqpoint{2.357333in}{1.952887in}}%
\pgfpathlineto{\pgfqpoint{2.358082in}{1.931247in}}%
\pgfpathlineto{\pgfqpoint{2.359135in}{1.532033in}}%
\pgfpathlineto{\pgfqpoint{2.358102in}{1.939447in}}%
\pgfpathlineto{\pgfqpoint{2.359901in}{1.532049in}}%
\pgfpathlineto{\pgfqpoint{2.361080in}{1.532159in}}%
\pgfpathlineto{\pgfqpoint{2.361143in}{2.031386in}}%
\pgfpathlineto{\pgfqpoint{2.362256in}{1.930986in}}%
\pgfpathlineto{\pgfqpoint{2.362263in}{1.916901in}}%
\pgfpathlineto{\pgfqpoint{2.362276in}{1.952628in}}%
\pgfpathlineto{\pgfqpoint{2.363025in}{1.931075in}}%
\pgfpathlineto{\pgfqpoint{2.364093in}{1.532032in}}%
\pgfpathlineto{\pgfqpoint{2.363044in}{1.939499in}}%
\pgfpathlineto{\pgfqpoint{2.364874in}{1.532049in}}%
\pgfpathlineto{\pgfqpoint{2.366026in}{1.532167in}}%
\pgfpathlineto{\pgfqpoint{2.366089in}{2.030959in}}%
\pgfpathlineto{\pgfqpoint{2.367202in}{1.931698in}}%
\pgfpathlineto{\pgfqpoint{2.367209in}{1.916880in}}%
\pgfpathlineto{\pgfqpoint{2.367222in}{1.951854in}}%
\pgfpathlineto{\pgfqpoint{2.367972in}{1.932429in}}%
\pgfpathlineto{\pgfqpoint{2.369030in}{1.532033in}}%
\pgfpathlineto{\pgfqpoint{2.367992in}{1.939589in}}%
\pgfpathlineto{\pgfqpoint{2.369356in}{1.532040in}}%
\pgfpathlineto{\pgfqpoint{2.370971in}{1.532160in}}%
\pgfpathlineto{\pgfqpoint{2.371034in}{2.031600in}}%
\pgfpathlineto{\pgfqpoint{2.372147in}{1.924409in}}%
\pgfpathlineto{\pgfqpoint{2.372151in}{1.916358in}}%
\pgfpathlineto{\pgfqpoint{2.372166in}{1.953223in}}%
\pgfpathlineto{\pgfqpoint{2.372913in}{1.933262in}}%
\pgfpathlineto{\pgfqpoint{2.373981in}{1.532032in}}%
\pgfpathlineto{\pgfqpoint{2.372933in}{1.939627in}}%
\pgfpathlineto{\pgfqpoint{2.374164in}{1.532036in}}%
\pgfpathlineto{\pgfqpoint{2.375917in}{1.532169in}}%
\pgfpathlineto{\pgfqpoint{2.375980in}{2.031126in}}%
\pgfpathlineto{\pgfqpoint{2.377090in}{1.936117in}}%
\pgfpathlineto{\pgfqpoint{2.377098in}{1.916389in}}%
\pgfpathlineto{\pgfqpoint{2.377113in}{1.952999in}}%
\pgfpathlineto{\pgfqpoint{2.377863in}{1.929791in}}%
\pgfpathlineto{\pgfqpoint{2.377878in}{1.939585in}}%
\pgfpathlineto{\pgfqpoint{2.378471in}{1.899241in}}%
\pgfpathlineto{\pgfqpoint{2.378917in}{1.532033in}}%
\pgfpathlineto{\pgfqpoint{2.379535in}{1.532046in}}%
\pgfpathlineto{\pgfqpoint{2.380862in}{1.532159in}}%
\pgfpathlineto{\pgfqpoint{2.380925in}{2.031236in}}%
\pgfpathlineto{\pgfqpoint{2.382039in}{1.921647in}}%
\pgfpathlineto{\pgfqpoint{2.382043in}{1.916293in}}%
\pgfpathlineto{\pgfqpoint{2.382059in}{1.952731in}}%
\pgfpathlineto{\pgfqpoint{2.382800in}{1.937526in}}%
\pgfpathlineto{\pgfqpoint{2.382824in}{1.939592in}}%
\pgfpathlineto{\pgfqpoint{2.383925in}{1.532033in}}%
\pgfpathlineto{\pgfqpoint{2.385808in}{1.532158in}}%
\pgfpathlineto{\pgfqpoint{2.385870in}{2.030687in}}%
\pgfpathlineto{\pgfqpoint{2.386986in}{1.921399in}}%
\pgfpathlineto{\pgfqpoint{2.386990in}{1.916432in}}%
\pgfpathlineto{\pgfqpoint{2.387006in}{1.952563in}}%
\pgfpathlineto{\pgfqpoint{2.387743in}{1.939935in}}%
\pgfpathlineto{\pgfqpoint{2.388875in}{1.532033in}}%
\pgfpathlineto{\pgfqpoint{2.389320in}{1.532044in}}%
\pgfpathlineto{\pgfqpoint{2.390753in}{1.532155in}}%
\pgfpathlineto{\pgfqpoint{2.390816in}{2.030054in}}%
\pgfpathlineto{\pgfqpoint{2.391956in}{1.934269in}}%
\pgfpathlineto{\pgfqpoint{2.391964in}{1.917476in}}%
\pgfpathlineto{\pgfqpoint{2.391979in}{1.951574in}}%
\pgfpathlineto{\pgfqpoint{2.392729in}{1.930561in}}%
\pgfpathlineto{\pgfqpoint{2.392746in}{1.939262in}}%
\pgfpathlineto{\pgfqpoint{2.393329in}{1.873118in}}%
\pgfpathlineto{\pgfqpoint{2.393749in}{1.532033in}}%
\pgfpathlineto{\pgfqpoint{2.394314in}{1.532044in}}%
\pgfpathlineto{\pgfqpoint{2.395698in}{1.532156in}}%
\pgfpathlineto{\pgfqpoint{2.395761in}{2.030051in}}%
\pgfpathlineto{\pgfqpoint{2.396898in}{1.946762in}}%
\pgfpathlineto{\pgfqpoint{2.396909in}{1.917445in}}%
\pgfpathlineto{\pgfqpoint{2.396925in}{1.951276in}}%
\pgfpathlineto{\pgfqpoint{2.397673in}{1.930580in}}%
\pgfpathlineto{\pgfqpoint{2.397690in}{1.939241in}}%
\pgfpathlineto{\pgfqpoint{2.398269in}{1.880799in}}%
\pgfpathlineto{\pgfqpoint{2.398701in}{1.532033in}}%
\pgfpathlineto{\pgfqpoint{2.399273in}{1.532044in}}%
\pgfpathlineto{\pgfqpoint{2.400644in}{1.532161in}}%
\pgfpathlineto{\pgfqpoint{2.400707in}{2.032262in}}%
\pgfpathlineto{\pgfqpoint{2.401820in}{1.925145in}}%
\pgfpathlineto{\pgfqpoint{2.401824in}{1.916427in}}%
\pgfpathlineto{\pgfqpoint{2.401840in}{1.953253in}}%
\pgfpathlineto{\pgfqpoint{2.402588in}{1.931035in}}%
\pgfpathlineto{\pgfqpoint{2.403643in}{1.532033in}}%
\pgfpathlineto{\pgfqpoint{2.402608in}{1.939476in}}%
\pgfpathlineto{\pgfqpoint{2.404508in}{1.532051in}}%
\pgfpathlineto{\pgfqpoint{2.405589in}{1.532154in}}%
\pgfpathlineto{\pgfqpoint{2.405652in}{2.030031in}}%
\pgfpathlineto{\pgfqpoint{2.406791in}{1.942166in}}%
\pgfpathlineto{\pgfqpoint{2.406799in}{1.918142in}}%
\pgfpathlineto{\pgfqpoint{2.406815in}{1.952059in}}%
\pgfpathlineto{\pgfqpoint{2.407565in}{1.930412in}}%
\pgfpathlineto{\pgfqpoint{2.407583in}{1.939288in}}%
\pgfpathlineto{\pgfqpoint{2.408144in}{1.898862in}}%
\pgfpathlineto{\pgfqpoint{2.408590in}{1.532033in}}%
\pgfpathlineto{\pgfqpoint{2.409160in}{1.532044in}}%
\pgfpathlineto{\pgfqpoint{2.410535in}{1.532158in}}%
\pgfpathlineto{\pgfqpoint{2.410598in}{2.030829in}}%
\pgfpathlineto{\pgfqpoint{2.411711in}{1.927764in}}%
\pgfpathlineto{\pgfqpoint{2.411715in}{1.917290in}}%
\pgfpathlineto{\pgfqpoint{2.411732in}{1.953048in}}%
\pgfpathlineto{\pgfqpoint{2.412480in}{1.931137in}}%
\pgfpathlineto{\pgfqpoint{2.413539in}{1.532033in}}%
\pgfpathlineto{\pgfqpoint{2.412499in}{1.939464in}}%
\pgfpathlineto{\pgfqpoint{2.414310in}{1.532049in}}%
\pgfpathlineto{\pgfqpoint{2.415480in}{1.532160in}}%
\pgfpathlineto{\pgfqpoint{2.415543in}{2.031613in}}%
\pgfpathlineto{\pgfqpoint{2.416656in}{1.927927in}}%
\pgfpathlineto{\pgfqpoint{2.416660in}{1.917289in}}%
\pgfpathlineto{\pgfqpoint{2.416677in}{1.952841in}}%
\pgfpathlineto{\pgfqpoint{2.417427in}{1.930352in}}%
\pgfpathlineto{\pgfqpoint{2.417443in}{1.939477in}}%
\pgfpathlineto{\pgfqpoint{2.418141in}{1.745256in}}%
\pgfpathlineto{\pgfqpoint{2.418481in}{1.532033in}}%
\pgfpathlineto{\pgfqpoint{2.419049in}{1.532044in}}%
\pgfpathlineto{\pgfqpoint{2.420426in}{1.532159in}}%
\pgfpathlineto{\pgfqpoint{2.420489in}{2.031132in}}%
\pgfpathlineto{\pgfqpoint{2.421603in}{1.921374in}}%
\pgfpathlineto{\pgfqpoint{2.421607in}{1.916270in}}%
\pgfpathlineto{\pgfqpoint{2.421622in}{1.952558in}}%
\pgfpathlineto{\pgfqpoint{2.422359in}{1.939869in}}%
\pgfpathlineto{\pgfqpoint{2.423431in}{1.532033in}}%
\pgfpathlineto{\pgfqpoint{2.423856in}{1.532042in}}%
\pgfpathlineto{\pgfqpoint{2.425371in}{1.532159in}}%
\pgfpathlineto{\pgfqpoint{2.425434in}{2.031201in}}%
\pgfpathlineto{\pgfqpoint{2.426548in}{1.923455in}}%
\pgfpathlineto{\pgfqpoint{2.426552in}{1.916294in}}%
\pgfpathlineto{\pgfqpoint{2.426567in}{1.953087in}}%
\pgfpathlineto{\pgfqpoint{2.427316in}{1.930078in}}%
\pgfpathlineto{\pgfqpoint{2.427320in}{1.930052in}}%
\pgfpathlineto{\pgfqpoint{2.427325in}{1.934234in}}%
\pgfpathlineto{\pgfqpoint{2.427333in}{1.939595in}}%
\pgfpathlineto{\pgfqpoint{2.427944in}{1.877552in}}%
\pgfpathlineto{\pgfqpoint{2.428369in}{1.532033in}}%
\pgfpathlineto{\pgfqpoint{2.428935in}{1.532044in}}%
\pgfpathlineto{\pgfqpoint{2.430317in}{1.532161in}}%
\pgfpathlineto{\pgfqpoint{2.430380in}{2.032160in}}%
\pgfpathlineto{\pgfqpoint{2.431492in}{1.925483in}}%
\pgfpathlineto{\pgfqpoint{2.431496in}{1.916426in}}%
\pgfpathlineto{\pgfqpoint{2.431512in}{1.953154in}}%
\pgfpathlineto{\pgfqpoint{2.432258in}{1.933357in}}%
\pgfpathlineto{\pgfqpoint{2.433318in}{1.532033in}}%
\pgfpathlineto{\pgfqpoint{2.432278in}{1.939503in}}%
\pgfpathlineto{\pgfqpoint{2.433492in}{1.532035in}}%
\pgfpathlineto{\pgfqpoint{2.435262in}{1.532158in}}%
\pgfpathlineto{\pgfqpoint{2.435325in}{2.030634in}}%
\pgfpathlineto{\pgfqpoint{2.436441in}{1.922361in}}%
\pgfpathlineto{\pgfqpoint{2.436445in}{1.916359in}}%
\pgfpathlineto{\pgfqpoint{2.436461in}{1.952725in}}%
\pgfpathlineto{\pgfqpoint{2.437200in}{1.939595in}}%
\pgfpathlineto{\pgfqpoint{2.438266in}{1.532033in}}%
\pgfpathlineto{\pgfqpoint{2.438492in}{1.532038in}}%
\pgfpathlineto{\pgfqpoint{2.440208in}{1.532160in}}%
\pgfpathlineto{\pgfqpoint{2.440271in}{2.031725in}}%
\pgfpathlineto{\pgfqpoint{2.441385in}{1.928349in}}%
\pgfpathlineto{\pgfqpoint{2.441389in}{1.917583in}}%
\pgfpathlineto{\pgfqpoint{2.441406in}{1.952839in}}%
\pgfpathlineto{\pgfqpoint{2.442155in}{1.930905in}}%
\pgfpathlineto{\pgfqpoint{2.443269in}{1.532033in}}%
\pgfpathlineto{\pgfqpoint{2.442203in}{1.939265in}}%
\pgfpathlineto{\pgfqpoint{2.444157in}{1.532053in}}%
\pgfpathlineto{\pgfqpoint{2.445153in}{1.532159in}}%
\pgfpathlineto{\pgfqpoint{2.445216in}{2.031311in}}%
\pgfpathlineto{\pgfqpoint{2.446330in}{1.920768in}}%
\pgfpathlineto{\pgfqpoint{2.446334in}{1.916357in}}%
\pgfpathlineto{\pgfqpoint{2.446350in}{1.952321in}}%
\pgfpathlineto{\pgfqpoint{2.447085in}{1.939723in}}%
\pgfpathlineto{\pgfqpoint{2.447797in}{1.771430in}}%
\pgfpathlineto{\pgfqpoint{2.448153in}{1.532033in}}%
\pgfpathlineto{\pgfqpoint{2.448722in}{1.532044in}}%
\pgfpathlineto{\pgfqpoint{2.450099in}{1.532159in}}%
\pgfpathlineto{\pgfqpoint{2.450161in}{2.031084in}}%
\pgfpathlineto{\pgfqpoint{2.451276in}{1.921259in}}%
\pgfpathlineto{\pgfqpoint{2.451280in}{1.916278in}}%
\pgfpathlineto{\pgfqpoint{2.451295in}{1.952416in}}%
\pgfpathlineto{\pgfqpoint{2.452033in}{1.939721in}}%
\pgfpathlineto{\pgfqpoint{2.453126in}{1.532033in}}%
\pgfpathlineto{\pgfqpoint{2.453369in}{1.532039in}}%
\pgfpathlineto{\pgfqpoint{2.455045in}{1.532177in}}%
\pgfpathlineto{\pgfqpoint{2.455107in}{2.032153in}}%
\pgfpathlineto{\pgfqpoint{2.456191in}{1.933185in}}%
\pgfpathlineto{\pgfqpoint{2.456198in}{1.915287in}}%
\pgfpathlineto{\pgfqpoint{2.456213in}{1.953415in}}%
\pgfpathlineto{\pgfqpoint{2.456963in}{1.929976in}}%
\pgfpathlineto{\pgfqpoint{2.456979in}{1.939793in}}%
\pgfpathlineto{\pgfqpoint{2.457670in}{1.801752in}}%
\pgfpathlineto{\pgfqpoint{2.458069in}{1.532033in}}%
\pgfpathlineto{\pgfqpoint{2.458613in}{1.532045in}}%
\pgfpathlineto{\pgfqpoint{2.459992in}{1.532307in}}%
\pgfpathlineto{\pgfqpoint{2.460052in}{2.024909in}}%
\pgfpathlineto{\pgfqpoint{2.461049in}{1.932446in}}%
\pgfpathlineto{\pgfqpoint{2.461057in}{1.913165in}}%
\pgfpathlineto{\pgfqpoint{2.461070in}{1.955846in}}%
\pgfpathlineto{\pgfqpoint{2.461816in}{1.937815in}}%
\pgfpathlineto{\pgfqpoint{2.461822in}{1.941350in}}%
\pgfpathlineto{\pgfqpoint{2.462516in}{1.917249in}}%
\pgfpathlineto{\pgfqpoint{2.463011in}{1.532032in}}%
\pgfpathlineto{\pgfqpoint{2.463698in}{1.532047in}}%
\pgfpathlineto{\pgfqpoint{2.464935in}{1.532165in}}%
\pgfpathlineto{\pgfqpoint{2.464998in}{2.030601in}}%
\pgfpathlineto{\pgfqpoint{2.466111in}{1.933027in}}%
\pgfpathlineto{\pgfqpoint{2.466118in}{1.916376in}}%
\pgfpathlineto{\pgfqpoint{2.466133in}{1.952271in}}%
\pgfpathlineto{\pgfqpoint{2.466882in}{1.930432in}}%
\pgfpathlineto{\pgfqpoint{2.466898in}{1.939284in}}%
\pgfpathlineto{\pgfqpoint{2.467621in}{1.706911in}}%
\pgfpathlineto{\pgfqpoint{2.467947in}{1.532032in}}%
\pgfpathlineto{\pgfqpoint{2.468528in}{1.532045in}}%
\pgfpathlineto{\pgfqpoint{2.469881in}{1.532168in}}%
\pgfpathlineto{\pgfqpoint{2.469944in}{2.030924in}}%
\pgfpathlineto{\pgfqpoint{2.471053in}{1.946149in}}%
\pgfpathlineto{\pgfqpoint{2.471064in}{1.916703in}}%
\pgfpathlineto{\pgfqpoint{2.471076in}{1.952180in}}%
\pgfpathlineto{\pgfqpoint{2.471826in}{1.932345in}}%
\pgfpathlineto{\pgfqpoint{2.472893in}{1.532033in}}%
\pgfpathlineto{\pgfqpoint{2.471847in}{1.939530in}}%
\pgfpathlineto{\pgfqpoint{2.473177in}{1.532038in}}%
\pgfpathlineto{\pgfqpoint{2.474829in}{1.532304in}}%
\pgfpathlineto{\pgfqpoint{2.474888in}{2.024682in}}%
\pgfpathlineto{\pgfqpoint{2.475886in}{1.931758in}}%
\pgfpathlineto{\pgfqpoint{2.475893in}{1.913224in}}%
\pgfpathlineto{\pgfqpoint{2.475906in}{1.955925in}}%
\pgfpathlineto{\pgfqpoint{2.476651in}{1.937791in}}%
\pgfpathlineto{\pgfqpoint{2.476657in}{1.941334in}}%
\pgfpathlineto{\pgfqpoint{2.477358in}{1.916946in}}%
\pgfpathlineto{\pgfqpoint{2.477832in}{1.532033in}}%
\pgfpathlineto{\pgfqpoint{2.478506in}{1.532047in}}%
\pgfpathlineto{\pgfqpoint{2.479771in}{1.532157in}}%
\pgfpathlineto{\pgfqpoint{2.479834in}{2.030319in}}%
\pgfpathlineto{\pgfqpoint{2.480949in}{1.925631in}}%
\pgfpathlineto{\pgfqpoint{2.480953in}{1.916548in}}%
\pgfpathlineto{\pgfqpoint{2.480968in}{1.952938in}}%
\pgfpathlineto{\pgfqpoint{2.481716in}{1.932899in}}%
\pgfpathlineto{\pgfqpoint{2.482770in}{1.532033in}}%
\pgfpathlineto{\pgfqpoint{2.481738in}{1.939452in}}%
\pgfpathlineto{\pgfqpoint{2.482991in}{1.532037in}}%
\pgfpathlineto{\pgfqpoint{2.484717in}{1.532160in}}%
\pgfpathlineto{\pgfqpoint{2.484780in}{2.031712in}}%
\pgfpathlineto{\pgfqpoint{2.485892in}{1.930101in}}%
\pgfpathlineto{\pgfqpoint{2.485928in}{1.916970in}}%
\pgfpathlineto{\pgfqpoint{2.485913in}{1.953122in}}%
\pgfpathlineto{\pgfqpoint{2.486660in}{1.932534in}}%
\pgfpathlineto{\pgfqpoint{2.487732in}{1.532033in}}%
\pgfpathlineto{\pgfqpoint{2.486681in}{1.939509in}}%
\pgfpathlineto{\pgfqpoint{2.487971in}{1.532038in}}%
\pgfpathlineto{\pgfqpoint{2.489663in}{1.532189in}}%
\pgfpathlineto{\pgfqpoint{2.489725in}{2.033311in}}%
\pgfpathlineto{\pgfqpoint{2.490804in}{1.942639in}}%
\pgfpathlineto{\pgfqpoint{2.490812in}{1.915653in}}%
\pgfpathlineto{\pgfqpoint{2.490828in}{1.954313in}}%
\pgfpathlineto{\pgfqpoint{2.491576in}{1.930319in}}%
\pgfpathlineto{\pgfqpoint{2.492695in}{1.532033in}}%
\pgfpathlineto{\pgfqpoint{2.491592in}{1.939559in}}%
\pgfpathlineto{\pgfqpoint{2.493778in}{1.532057in}}%
\pgfpathlineto{\pgfqpoint{2.494608in}{1.532179in}}%
\pgfpathlineto{\pgfqpoint{2.494671in}{2.031913in}}%
\pgfpathlineto{\pgfqpoint{2.495755in}{1.932455in}}%
\pgfpathlineto{\pgfqpoint{2.495762in}{1.915238in}}%
\pgfpathlineto{\pgfqpoint{2.495806in}{1.953100in}}%
\pgfpathlineto{\pgfqpoint{2.496526in}{1.930687in}}%
\pgfpathlineto{\pgfqpoint{2.497621in}{1.532032in}}%
\pgfpathlineto{\pgfqpoint{2.496545in}{1.939522in}}%
\pgfpathlineto{\pgfqpoint{2.498438in}{1.532052in}}%
\pgfpathlineto{\pgfqpoint{2.499552in}{1.532140in}}%
\pgfpathlineto{\pgfqpoint{2.499615in}{2.029969in}}%
\pgfpathlineto{\pgfqpoint{2.500791in}{1.922778in}}%
\pgfpathlineto{\pgfqpoint{2.500794in}{1.918190in}}%
\pgfpathlineto{\pgfqpoint{2.500810in}{1.950749in}}%
\pgfpathlineto{\pgfqpoint{2.501546in}{1.938977in}}%
\pgfpathlineto{\pgfqpoint{2.502228in}{1.723771in}}%
\pgfpathlineto{\pgfqpoint{2.502549in}{1.532033in}}%
\pgfpathlineto{\pgfqpoint{2.503114in}{1.532044in}}%
\pgfpathlineto{\pgfqpoint{2.504498in}{1.532155in}}%
\pgfpathlineto{\pgfqpoint{2.504561in}{2.029958in}}%
\pgfpathlineto{\pgfqpoint{2.505702in}{1.936923in}}%
\pgfpathlineto{\pgfqpoint{2.505710in}{1.917285in}}%
\pgfpathlineto{\pgfqpoint{2.505725in}{1.951880in}}%
\pgfpathlineto{\pgfqpoint{2.506472in}{1.931539in}}%
\pgfpathlineto{\pgfqpoint{2.507506in}{1.532033in}}%
\pgfpathlineto{\pgfqpoint{2.506491in}{1.939127in}}%
\pgfpathlineto{\pgfqpoint{2.508131in}{1.532047in}}%
\pgfpathlineto{\pgfqpoint{2.509446in}{1.532209in}}%
\pgfpathlineto{\pgfqpoint{2.509509in}{2.021584in}}%
\pgfpathlineto{\pgfqpoint{2.510535in}{1.918647in}}%
\pgfpathlineto{\pgfqpoint{2.510539in}{1.916103in}}%
\pgfpathlineto{\pgfqpoint{2.510582in}{1.952448in}}%
\pgfpathlineto{\pgfqpoint{2.511271in}{1.934317in}}%
\pgfpathlineto{\pgfqpoint{2.511281in}{1.940345in}}%
\pgfpathlineto{\pgfqpoint{2.511990in}{1.906659in}}%
\pgfpathlineto{\pgfqpoint{2.512470in}{1.532033in}}%
\pgfpathlineto{\pgfqpoint{2.513059in}{1.532045in}}%
\pgfpathlineto{\pgfqpoint{2.514390in}{1.532179in}}%
\pgfpathlineto{\pgfqpoint{2.514453in}{2.031927in}}%
\pgfpathlineto{\pgfqpoint{2.515536in}{1.932888in}}%
\pgfpathlineto{\pgfqpoint{2.515543in}{1.915204in}}%
\pgfpathlineto{\pgfqpoint{2.515559in}{1.953316in}}%
\pgfpathlineto{\pgfqpoint{2.516308in}{1.929788in}}%
\pgfpathlineto{\pgfqpoint{2.516325in}{1.939817in}}%
\pgfpathlineto{\pgfqpoint{2.516947in}{1.897417in}}%
\pgfpathlineto{\pgfqpoint{2.517413in}{1.532031in}}%
\pgfpathlineto{\pgfqpoint{2.518006in}{1.532044in}}%
\pgfpathlineto{\pgfqpoint{2.519333in}{1.532130in}}%
\pgfpathlineto{\pgfqpoint{2.519397in}{2.029954in}}%
\pgfpathlineto{\pgfqpoint{2.520629in}{1.929408in}}%
\pgfpathlineto{\pgfqpoint{2.520664in}{1.920363in}}%
\pgfpathlineto{\pgfqpoint{2.520649in}{1.949586in}}%
\pgfpathlineto{\pgfqpoint{2.521398in}{1.931768in}}%
\pgfpathlineto{\pgfqpoint{2.522346in}{1.532032in}}%
\pgfpathlineto{\pgfqpoint{2.521418in}{1.938130in}}%
\pgfpathlineto{\pgfqpoint{2.523521in}{1.532058in}}%
\pgfpathlineto{\pgfqpoint{2.524281in}{1.532168in}}%
\pgfpathlineto{\pgfqpoint{2.524344in}{2.030912in}}%
\pgfpathlineto{\pgfqpoint{2.525456in}{1.935373in}}%
\pgfpathlineto{\pgfqpoint{2.525464in}{1.916358in}}%
\pgfpathlineto{\pgfqpoint{2.525480in}{1.952239in}}%
\pgfpathlineto{\pgfqpoint{2.526226in}{1.932924in}}%
\pgfpathlineto{\pgfqpoint{2.527296in}{1.532033in}}%
\pgfpathlineto{\pgfqpoint{2.526246in}{1.939472in}}%
\pgfpathlineto{\pgfqpoint{2.527483in}{1.532036in}}%
\pgfpathlineto{\pgfqpoint{2.529229in}{1.532305in}}%
\pgfpathlineto{\pgfqpoint{2.529288in}{2.024679in}}%
\pgfpathlineto{\pgfqpoint{2.530286in}{1.930747in}}%
\pgfpathlineto{\pgfqpoint{2.530294in}{1.913299in}}%
\pgfpathlineto{\pgfqpoint{2.530307in}{1.956113in}}%
\pgfpathlineto{\pgfqpoint{2.531051in}{1.937809in}}%
\pgfpathlineto{\pgfqpoint{2.531057in}{1.941277in}}%
\pgfpathlineto{\pgfqpoint{2.531759in}{1.916616in}}%
\pgfpathlineto{\pgfqpoint{2.532227in}{1.532033in}}%
\pgfpathlineto{\pgfqpoint{2.532896in}{1.532046in}}%
\pgfpathlineto{\pgfqpoint{2.534171in}{1.532158in}}%
\pgfpathlineto{\pgfqpoint{2.534234in}{2.030546in}}%
\pgfpathlineto{\pgfqpoint{2.535348in}{1.926541in}}%
\pgfpathlineto{\pgfqpoint{2.535352in}{1.916822in}}%
\pgfpathlineto{\pgfqpoint{2.535368in}{1.952973in}}%
\pgfpathlineto{\pgfqpoint{2.536114in}{1.933392in}}%
\pgfpathlineto{\pgfqpoint{2.537188in}{1.532033in}}%
\pgfpathlineto{\pgfqpoint{2.536134in}{1.939279in}}%
\pgfpathlineto{\pgfqpoint{2.537330in}{1.532036in}}%
\pgfpathlineto{\pgfqpoint{2.539118in}{1.532181in}}%
\pgfpathlineto{\pgfqpoint{2.539180in}{2.031551in}}%
\pgfpathlineto{\pgfqpoint{2.540261in}{1.940087in}}%
\pgfpathlineto{\pgfqpoint{2.540269in}{1.915418in}}%
\pgfpathlineto{\pgfqpoint{2.540284in}{1.953981in}}%
\pgfpathlineto{\pgfqpoint{2.541035in}{1.929586in}}%
\pgfpathlineto{\pgfqpoint{2.541053in}{1.939750in}}%
\pgfpathlineto{\pgfqpoint{2.541637in}{1.926513in}}%
\pgfpathlineto{\pgfqpoint{2.542139in}{1.532031in}}%
\pgfpathlineto{\pgfqpoint{2.542977in}{1.532053in}}%
\pgfpathlineto{\pgfqpoint{2.544060in}{1.532128in}}%
\pgfpathlineto{\pgfqpoint{2.544124in}{2.029640in}}%
\pgfpathlineto{\pgfqpoint{2.545358in}{1.924222in}}%
\pgfpathlineto{\pgfqpoint{2.545362in}{1.919688in}}%
\pgfpathlineto{\pgfqpoint{2.545378in}{1.949203in}}%
\pgfpathlineto{\pgfqpoint{2.546116in}{1.938490in}}%
\pgfpathlineto{\pgfqpoint{2.547067in}{1.532033in}}%
\pgfpathlineto{\pgfqpoint{2.547889in}{1.532051in}}%
\pgfpathlineto{\pgfqpoint{2.549007in}{1.532153in}}%
\pgfpathlineto{\pgfqpoint{2.549070in}{2.029915in}}%
\pgfpathlineto{\pgfqpoint{2.550211in}{1.932976in}}%
\pgfpathlineto{\pgfqpoint{2.550219in}{1.917449in}}%
\pgfpathlineto{\pgfqpoint{2.550234in}{1.950897in}}%
\pgfpathlineto{\pgfqpoint{2.550982in}{1.930412in}}%
\pgfpathlineto{\pgfqpoint{2.550999in}{1.939070in}}%
\pgfpathlineto{\pgfqpoint{2.551574in}{1.887980in}}%
\pgfpathlineto{\pgfqpoint{2.552011in}{1.532031in}}%
\pgfpathlineto{\pgfqpoint{2.552582in}{1.532044in}}%
\pgfpathlineto{\pgfqpoint{2.553951in}{1.532132in}}%
\pgfpathlineto{\pgfqpoint{2.554015in}{2.031927in}}%
\pgfpathlineto{\pgfqpoint{2.555221in}{1.923210in}}%
\pgfpathlineto{\pgfqpoint{2.555225in}{1.918417in}}%
\pgfpathlineto{\pgfqpoint{2.555241in}{1.950055in}}%
\pgfpathlineto{\pgfqpoint{2.555980in}{1.938954in}}%
\pgfpathlineto{\pgfqpoint{2.556953in}{1.532033in}}%
\pgfpathlineto{\pgfqpoint{2.557669in}{1.532049in}}%
\pgfpathlineto{\pgfqpoint{2.558898in}{1.532154in}}%
\pgfpathlineto{\pgfqpoint{2.558961in}{2.029907in}}%
\pgfpathlineto{\pgfqpoint{2.560101in}{1.936293in}}%
\pgfpathlineto{\pgfqpoint{2.560109in}{1.917281in}}%
\pgfpathlineto{\pgfqpoint{2.560124in}{1.951617in}}%
\pgfpathlineto{\pgfqpoint{2.560872in}{1.931363in}}%
\pgfpathlineto{\pgfqpoint{2.561935in}{1.532033in}}%
\pgfpathlineto{\pgfqpoint{2.560892in}{1.939050in}}%
\pgfpathlineto{\pgfqpoint{2.562726in}{1.532050in}}%
\pgfpathlineto{\pgfqpoint{2.563847in}{1.532305in}}%
\pgfpathlineto{\pgfqpoint{2.563906in}{2.024652in}}%
\pgfpathlineto{\pgfqpoint{2.564904in}{1.932069in}}%
\pgfpathlineto{\pgfqpoint{2.564912in}{1.913195in}}%
\pgfpathlineto{\pgfqpoint{2.564924in}{1.955798in}}%
\pgfpathlineto{\pgfqpoint{2.565671in}{1.937789in}}%
\pgfpathlineto{\pgfqpoint{2.565676in}{1.941215in}}%
\pgfpathlineto{\pgfqpoint{2.566382in}{1.913475in}}%
\pgfpathlineto{\pgfqpoint{2.566855in}{1.532033in}}%
\pgfpathlineto{\pgfqpoint{2.567485in}{1.532047in}}%
\pgfpathlineto{\pgfqpoint{2.568791in}{1.532186in}}%
\pgfpathlineto{\pgfqpoint{2.568853in}{2.031740in}}%
\pgfpathlineto{\pgfqpoint{2.569933in}{1.938447in}}%
\pgfpathlineto{\pgfqpoint{2.569941in}{1.915299in}}%
\pgfpathlineto{\pgfqpoint{2.569956in}{1.953779in}}%
\pgfpathlineto{\pgfqpoint{2.570705in}{1.930646in}}%
\pgfpathlineto{\pgfqpoint{2.571804in}{1.532033in}}%
\pgfpathlineto{\pgfqpoint{2.570724in}{1.939693in}}%
\pgfpathlineto{\pgfqpoint{2.572587in}{1.532049in}}%
\pgfpathlineto{\pgfqpoint{2.573736in}{1.532189in}}%
\pgfpathlineto{\pgfqpoint{2.573798in}{2.033235in}}%
\pgfpathlineto{\pgfqpoint{2.574875in}{1.948773in}}%
\pgfpathlineto{\pgfqpoint{2.574887in}{1.914625in}}%
\pgfpathlineto{\pgfqpoint{2.574902in}{1.953802in}}%
\pgfpathlineto{\pgfqpoint{2.575650in}{1.929911in}}%
\pgfpathlineto{\pgfqpoint{2.575666in}{1.939786in}}%
\pgfpathlineto{\pgfqpoint{2.576427in}{1.700852in}}%
\pgfpathlineto{\pgfqpoint{2.576748in}{1.532033in}}%
\pgfpathlineto{\pgfqpoint{2.577327in}{1.532045in}}%
\pgfpathlineto{\pgfqpoint{2.578683in}{1.532306in}}%
\pgfpathlineto{\pgfqpoint{2.578743in}{2.024739in}}%
\pgfpathlineto{\pgfqpoint{2.579738in}{1.944732in}}%
\pgfpathlineto{\pgfqpoint{2.579748in}{1.913360in}}%
\pgfpathlineto{\pgfqpoint{2.579762in}{1.956383in}}%
\pgfpathlineto{\pgfqpoint{2.580507in}{1.937826in}}%
\pgfpathlineto{\pgfqpoint{2.580512in}{1.941311in}}%
\pgfpathlineto{\pgfqpoint{2.581210in}{1.916566in}}%
\pgfpathlineto{\pgfqpoint{2.581680in}{1.532033in}}%
\pgfpathlineto{\pgfqpoint{2.582396in}{1.532049in}}%
\pgfpathlineto{\pgfqpoint{2.583626in}{1.532153in}}%
\pgfpathlineto{\pgfqpoint{2.583688in}{2.029811in}}%
\pgfpathlineto{\pgfqpoint{2.584828in}{1.939226in}}%
\pgfpathlineto{\pgfqpoint{2.584836in}{1.917551in}}%
\pgfpathlineto{\pgfqpoint{2.584851in}{1.951663in}}%
\pgfpathlineto{\pgfqpoint{2.585599in}{1.932157in}}%
\pgfpathlineto{\pgfqpoint{2.586631in}{1.532033in}}%
\pgfpathlineto{\pgfqpoint{2.585618in}{1.939042in}}%
\pgfpathlineto{\pgfqpoint{2.587106in}{1.532042in}}%
\pgfpathlineto{\pgfqpoint{2.588571in}{1.532158in}}%
\pgfpathlineto{\pgfqpoint{2.588634in}{2.030867in}}%
\pgfpathlineto{\pgfqpoint{2.589749in}{1.921850in}}%
\pgfpathlineto{\pgfqpoint{2.589753in}{1.916027in}}%
\pgfpathlineto{\pgfqpoint{2.589768in}{1.952790in}}%
\pgfpathlineto{\pgfqpoint{2.590513in}{1.933274in}}%
\pgfpathlineto{\pgfqpoint{2.591577in}{1.532031in}}%
\pgfpathlineto{\pgfqpoint{2.590534in}{1.939478in}}%
\pgfpathlineto{\pgfqpoint{2.591755in}{1.532035in}}%
\pgfpathlineto{\pgfqpoint{2.593515in}{1.532130in}}%
\pgfpathlineto{\pgfqpoint{2.593579in}{2.031056in}}%
\pgfpathlineto{\pgfqpoint{2.594809in}{1.935197in}}%
\pgfpathlineto{\pgfqpoint{2.594817in}{1.919373in}}%
\pgfpathlineto{\pgfqpoint{2.594832in}{1.949164in}}%
\pgfpathlineto{\pgfqpoint{2.595580in}{1.932223in}}%
\pgfpathlineto{\pgfqpoint{2.596523in}{1.532033in}}%
\pgfpathlineto{\pgfqpoint{2.595600in}{1.938250in}}%
\pgfpathlineto{\pgfqpoint{2.597196in}{1.532047in}}%
\pgfpathlineto{\pgfqpoint{2.598462in}{1.532158in}}%
\pgfpathlineto{\pgfqpoint{2.598525in}{2.030545in}}%
\pgfpathlineto{\pgfqpoint{2.599642in}{1.918112in}}%
\pgfpathlineto{\pgfqpoint{2.599674in}{1.916999in}}%
\pgfpathlineto{\pgfqpoint{2.599659in}{1.952879in}}%
\pgfpathlineto{\pgfqpoint{2.600334in}{1.938743in}}%
\pgfpathlineto{\pgfqpoint{2.600338in}{1.940335in}}%
\pgfpathlineto{\pgfqpoint{2.601005in}{1.911402in}}%
\pgfpathlineto{\pgfqpoint{2.601464in}{1.532033in}}%
\pgfpathlineto{\pgfqpoint{2.602114in}{1.532046in}}%
\pgfpathlineto{\pgfqpoint{2.603411in}{1.532306in}}%
\pgfpathlineto{\pgfqpoint{2.603470in}{2.024755in}}%
\pgfpathlineto{\pgfqpoint{2.604468in}{1.932014in}}%
\pgfpathlineto{\pgfqpoint{2.604475in}{1.912980in}}%
\pgfpathlineto{\pgfqpoint{2.604488in}{1.955770in}}%
\pgfpathlineto{\pgfqpoint{2.605233in}{1.937742in}}%
\pgfpathlineto{\pgfqpoint{2.605238in}{1.941253in}}%
\pgfpathlineto{\pgfqpoint{2.605933in}{1.923217in}}%
\pgfpathlineto{\pgfqpoint{2.606400in}{1.532032in}}%
\pgfpathlineto{\pgfqpoint{2.607162in}{1.532049in}}%
\pgfpathlineto{\pgfqpoint{2.608356in}{1.532303in}}%
\pgfpathlineto{\pgfqpoint{2.608415in}{2.024495in}}%
\pgfpathlineto{\pgfqpoint{2.609413in}{1.933318in}}%
\pgfpathlineto{\pgfqpoint{2.609420in}{1.912977in}}%
\pgfpathlineto{\pgfqpoint{2.609433in}{1.955031in}}%
\pgfpathlineto{\pgfqpoint{2.610182in}{1.941123in}}%
\pgfpathlineto{\pgfqpoint{2.611367in}{1.532033in}}%
\pgfpathlineto{\pgfqpoint{2.611654in}{1.532038in}}%
\pgfpathlineto{\pgfqpoint{2.613300in}{1.532190in}}%
\pgfpathlineto{\pgfqpoint{2.613362in}{2.033272in}}%
\pgfpathlineto{\pgfqpoint{2.614417in}{1.920258in}}%
\pgfpathlineto{\pgfqpoint{2.614421in}{1.913633in}}%
\pgfpathlineto{\pgfqpoint{2.614436in}{1.955094in}}%
\pgfpathlineto{\pgfqpoint{2.615181in}{1.933984in}}%
\pgfpathlineto{\pgfqpoint{2.616133in}{1.595376in}}%
\pgfpathlineto{\pgfqpoint{2.615203in}{1.939976in}}%
\pgfpathlineto{\pgfqpoint{2.616212in}{1.616000in}}%
\pgfpathlineto{\pgfqpoint{2.616246in}{1.651206in}}%
\pgfpathlineto{\pgfqpoint{2.616314in}{1.532033in}}%
\pgfpathlineto{\pgfqpoint{2.616850in}{1.532045in}}%
\pgfpathlineto{\pgfqpoint{2.618245in}{1.532189in}}%
\pgfpathlineto{\pgfqpoint{2.618307in}{2.033122in}}%
\pgfpathlineto{\pgfqpoint{2.619385in}{1.948791in}}%
\pgfpathlineto{\pgfqpoint{2.619396in}{1.914778in}}%
\pgfpathlineto{\pgfqpoint{2.619411in}{1.953679in}}%
\pgfpathlineto{\pgfqpoint{2.620158in}{1.930772in}}%
\pgfpathlineto{\pgfqpoint{2.621252in}{1.532031in}}%
\pgfpathlineto{\pgfqpoint{2.620178in}{1.939753in}}%
\pgfpathlineto{\pgfqpoint{2.621829in}{1.532044in}}%
\pgfpathlineto{\pgfqpoint{2.623187in}{1.532128in}}%
\pgfpathlineto{\pgfqpoint{2.623251in}{2.030736in}}%
\pgfpathlineto{\pgfqpoint{2.624482in}{1.932094in}}%
\pgfpathlineto{\pgfqpoint{2.624489in}{1.919572in}}%
\pgfpathlineto{\pgfqpoint{2.624501in}{1.948632in}}%
\pgfpathlineto{\pgfqpoint{2.625254in}{1.931339in}}%
\pgfpathlineto{\pgfqpoint{2.625273in}{1.938098in}}%
\pgfpathlineto{\pgfqpoint{2.625913in}{1.659884in}}%
\pgfpathlineto{\pgfqpoint{2.626255in}{1.532033in}}%
\pgfpathlineto{\pgfqpoint{2.626799in}{1.532046in}}%
\pgfpathlineto{\pgfqpoint{2.628135in}{1.532155in}}%
\pgfpathlineto{\pgfqpoint{2.628198in}{2.029869in}}%
\pgfpathlineto{\pgfqpoint{2.629335in}{1.947074in}}%
\pgfpathlineto{\pgfqpoint{2.629347in}{1.917176in}}%
\pgfpathlineto{\pgfqpoint{2.629362in}{1.951380in}}%
\pgfpathlineto{\pgfqpoint{2.630112in}{1.929819in}}%
\pgfpathlineto{\pgfqpoint{2.630127in}{1.938942in}}%
\pgfpathlineto{\pgfqpoint{2.630695in}{1.894951in}}%
\pgfpathlineto{\pgfqpoint{2.631143in}{1.532031in}}%
\pgfpathlineto{\pgfqpoint{2.631719in}{1.532044in}}%
\pgfpathlineto{\pgfqpoint{2.633078in}{1.532128in}}%
\pgfpathlineto{\pgfqpoint{2.633142in}{2.030775in}}%
\pgfpathlineto{\pgfqpoint{2.634375in}{1.928953in}}%
\pgfpathlineto{\pgfqpoint{2.634379in}{1.920350in}}%
\pgfpathlineto{\pgfqpoint{2.634396in}{1.949365in}}%
\pgfpathlineto{\pgfqpoint{2.635142in}{1.934020in}}%
\pgfpathlineto{\pgfqpoint{2.636111in}{1.532028in}}%
\pgfpathlineto{\pgfqpoint{2.635162in}{1.938166in}}%
\pgfpathlineto{\pgfqpoint{2.636451in}{1.532046in}}%
\pgfpathlineto{\pgfqpoint{2.638015in}{1.532085in}}%
\pgfpathlineto{\pgfqpoint{2.638025in}{1.532146in}}%
\pgfpathlineto{\pgfqpoint{2.638093in}{2.016224in}}%
\pgfpathlineto{\pgfqpoint{2.639200in}{1.944889in}}%
\pgfpathlineto{\pgfqpoint{2.639209in}{1.919705in}}%
\pgfpathlineto{\pgfqpoint{2.639226in}{1.950532in}}%
\pgfpathlineto{\pgfqpoint{2.639974in}{1.937465in}}%
\pgfpathlineto{\pgfqpoint{2.641038in}{1.532032in}}%
\pgfpathlineto{\pgfqpoint{2.640035in}{1.939238in}}%
\pgfpathlineto{\pgfqpoint{2.641224in}{1.532036in}}%
\pgfpathlineto{\pgfqpoint{2.642970in}{1.532139in}}%
\pgfpathlineto{\pgfqpoint{2.643033in}{2.029923in}}%
\pgfpathlineto{\pgfqpoint{2.644209in}{1.922797in}}%
\pgfpathlineto{\pgfqpoint{2.644213in}{1.917757in}}%
\pgfpathlineto{\pgfqpoint{2.644229in}{1.950862in}}%
\pgfpathlineto{\pgfqpoint{2.644967in}{1.938864in}}%
\pgfpathlineto{\pgfqpoint{2.645937in}{1.532022in}}%
\pgfpathlineto{\pgfqpoint{2.646357in}{1.532046in}}%
\pgfpathlineto{\pgfqpoint{2.647906in}{1.532086in}}%
\pgfpathlineto{\pgfqpoint{2.647916in}{1.532147in}}%
\pgfpathlineto{\pgfqpoint{2.647983in}{2.029988in}}%
\pgfpathlineto{\pgfqpoint{2.649090in}{1.941983in}}%
\pgfpathlineto{\pgfqpoint{2.649101in}{1.917207in}}%
\pgfpathlineto{\pgfqpoint{2.649115in}{1.952654in}}%
\pgfpathlineto{\pgfqpoint{2.649860in}{1.933810in}}%
\pgfpathlineto{\pgfqpoint{2.649870in}{1.939544in}}%
\pgfpathlineto{\pgfqpoint{2.650482in}{1.890837in}}%
\pgfpathlineto{\pgfqpoint{2.650953in}{1.532031in}}%
\pgfpathlineto{\pgfqpoint{2.651494in}{1.532047in}}%
\pgfpathlineto{\pgfqpoint{2.652857in}{1.532104in}}%
\pgfpathlineto{\pgfqpoint{2.652866in}{1.532448in}}%
\pgfpathlineto{\pgfqpoint{2.652968in}{2.030248in}}%
\pgfpathlineto{\pgfqpoint{2.653894in}{1.922651in}}%
\pgfpathlineto{\pgfqpoint{2.653898in}{1.911501in}}%
\pgfpathlineto{\pgfqpoint{2.653914in}{1.957828in}}%
\pgfpathlineto{\pgfqpoint{2.654658in}{1.937982in}}%
\pgfpathlineto{\pgfqpoint{2.654663in}{1.941751in}}%
\pgfpathlineto{\pgfqpoint{2.655387in}{1.918458in}}%
\pgfpathlineto{\pgfqpoint{2.655846in}{1.532027in}}%
\pgfpathlineto{\pgfqpoint{2.656581in}{1.532042in}}%
\pgfpathlineto{\pgfqpoint{2.657807in}{1.532147in}}%
\pgfpathlineto{\pgfqpoint{2.657910in}{2.001856in}}%
\pgfpathlineto{\pgfqpoint{2.658967in}{1.927482in}}%
\pgfpathlineto{\pgfqpoint{2.658972in}{1.920009in}}%
\pgfpathlineto{\pgfqpoint{2.658987in}{1.948949in}}%
\pgfpathlineto{\pgfqpoint{2.659730in}{1.929210in}}%
\pgfpathlineto{\pgfqpoint{2.659747in}{1.939741in}}%
\pgfpathlineto{\pgfqpoint{2.660385in}{1.860788in}}%
\pgfpathlineto{\pgfqpoint{2.660839in}{1.532032in}}%
\pgfpathlineto{\pgfqpoint{2.661376in}{1.532045in}}%
\pgfpathlineto{\pgfqpoint{2.662751in}{1.532125in}}%
\pgfpathlineto{\pgfqpoint{2.662851in}{1.995592in}}%
\pgfpathlineto{\pgfqpoint{2.664220in}{1.940826in}}%
\pgfpathlineto{\pgfqpoint{2.664234in}{1.926070in}}%
\pgfpathlineto{\pgfqpoint{2.664247in}{1.943091in}}%
\pgfpathlineto{\pgfqpoint{2.664968in}{1.930077in}}%
\pgfpathlineto{\pgfqpoint{2.665017in}{1.937009in}}%
\pgfpathlineto{\pgfqpoint{2.665314in}{1.892306in}}%
\pgfpathlineto{\pgfqpoint{2.665772in}{1.532029in}}%
\pgfpathlineto{\pgfqpoint{2.666308in}{1.532048in}}%
\pgfpathlineto{\pgfqpoint{2.667697in}{1.532134in}}%
\pgfpathlineto{\pgfqpoint{2.667762in}{2.039620in}}%
\pgfpathlineto{\pgfqpoint{2.668932in}{1.934588in}}%
\pgfpathlineto{\pgfqpoint{2.668940in}{1.916516in}}%
\pgfpathlineto{\pgfqpoint{2.668984in}{1.951439in}}%
\pgfpathlineto{\pgfqpoint{2.669700in}{1.938960in}}%
\pgfpathlineto{\pgfqpoint{2.670699in}{1.532033in}}%
\pgfpathlineto{\pgfqpoint{2.669727in}{1.939105in}}%
\pgfpathlineto{\pgfqpoint{2.670921in}{1.532038in}}%
\pgfpathlineto{\pgfqpoint{2.672644in}{1.532153in}}%
\pgfpathlineto{\pgfqpoint{2.672706in}{2.029781in}}%
\pgfpathlineto{\pgfqpoint{2.673847in}{1.937670in}}%
\pgfpathlineto{\pgfqpoint{2.673855in}{1.917203in}}%
\pgfpathlineto{\pgfqpoint{2.673870in}{1.951726in}}%
\pgfpathlineto{\pgfqpoint{2.674620in}{1.930378in}}%
\pgfpathlineto{\pgfqpoint{2.674636in}{1.938964in}}%
\pgfpathlineto{\pgfqpoint{2.675235in}{1.855763in}}%
\pgfpathlineto{\pgfqpoint{2.675640in}{1.532031in}}%
\pgfpathlineto{\pgfqpoint{2.676206in}{1.532043in}}%
\pgfpathlineto{\pgfqpoint{2.677592in}{1.532308in}}%
\pgfpathlineto{\pgfqpoint{2.677652in}{2.024666in}}%
\pgfpathlineto{\pgfqpoint{2.678649in}{1.932217in}}%
\pgfpathlineto{\pgfqpoint{2.678657in}{1.912952in}}%
\pgfpathlineto{\pgfqpoint{2.678670in}{1.955678in}}%
\pgfpathlineto{\pgfqpoint{2.679415in}{1.937728in}}%
\pgfpathlineto{\pgfqpoint{2.679420in}{1.941168in}}%
\pgfpathlineto{\pgfqpoint{2.680134in}{1.908583in}}%
\pgfpathlineto{\pgfqpoint{2.680609in}{1.532033in}}%
\pgfpathlineto{\pgfqpoint{2.681196in}{1.532046in}}%
\pgfpathlineto{\pgfqpoint{2.682538in}{1.532307in}}%
\pgfpathlineto{\pgfqpoint{2.682597in}{2.024688in}}%
\pgfpathlineto{\pgfqpoint{2.683595in}{1.931343in}}%
\pgfpathlineto{\pgfqpoint{2.683602in}{1.913003in}}%
\pgfpathlineto{\pgfqpoint{2.683616in}{1.955882in}}%
\pgfpathlineto{\pgfqpoint{2.684359in}{1.937892in}}%
\pgfpathlineto{\pgfqpoint{2.684364in}{1.941117in}}%
\pgfpathlineto{\pgfqpoint{2.685075in}{1.908260in}}%
\pgfpathlineto{\pgfqpoint{2.685519in}{1.532027in}}%
\pgfpathlineto{\pgfqpoint{2.686155in}{1.532040in}}%
\pgfpathlineto{\pgfqpoint{2.687480in}{1.532147in}}%
\pgfpathlineto{\pgfqpoint{2.687581in}{2.000612in}}%
\pgfpathlineto{\pgfqpoint{2.688639in}{1.927418in}}%
\pgfpathlineto{\pgfqpoint{2.688644in}{1.919717in}}%
\pgfpathlineto{\pgfqpoint{2.688660in}{1.949182in}}%
\pgfpathlineto{\pgfqpoint{2.689397in}{1.929204in}}%
\pgfpathlineto{\pgfqpoint{2.689416in}{1.939489in}}%
\pgfpathlineto{\pgfqpoint{2.690043in}{1.886218in}}%
\pgfpathlineto{\pgfqpoint{2.690506in}{1.532028in}}%
\pgfpathlineto{\pgfqpoint{2.691043in}{1.532049in}}%
\pgfpathlineto{\pgfqpoint{2.692425in}{1.532143in}}%
\pgfpathlineto{\pgfqpoint{2.692492in}{2.031156in}}%
\pgfpathlineto{\pgfqpoint{2.693603in}{1.928237in}}%
\pgfpathlineto{\pgfqpoint{2.693637in}{1.916983in}}%
\pgfpathlineto{\pgfqpoint{2.693625in}{1.952120in}}%
\pgfpathlineto{\pgfqpoint{2.694369in}{1.933970in}}%
\pgfpathlineto{\pgfqpoint{2.694379in}{1.939543in}}%
\pgfpathlineto{\pgfqpoint{2.695027in}{1.837364in}}%
\pgfpathlineto{\pgfqpoint{2.695392in}{1.532011in}}%
\pgfpathlineto{\pgfqpoint{2.695991in}{1.532054in}}%
\pgfpathlineto{\pgfqpoint{2.697365in}{1.532100in}}%
\pgfpathlineto{\pgfqpoint{2.697373in}{1.532213in}}%
\pgfpathlineto{\pgfqpoint{2.697433in}{2.026512in}}%
\pgfpathlineto{\pgfqpoint{2.698489in}{1.928804in}}%
\pgfpathlineto{\pgfqpoint{2.698525in}{1.915708in}}%
\pgfpathlineto{\pgfqpoint{2.698510in}{1.954184in}}%
\pgfpathlineto{\pgfqpoint{2.699259in}{1.930329in}}%
\pgfpathlineto{\pgfqpoint{2.700400in}{1.532028in}}%
\pgfpathlineto{\pgfqpoint{2.699308in}{1.939536in}}%
\pgfpathlineto{\pgfqpoint{2.701233in}{1.532056in}}%
\pgfpathlineto{\pgfqpoint{2.702306in}{1.532086in}}%
\pgfpathlineto{\pgfqpoint{2.702316in}{1.532146in}}%
\pgfpathlineto{\pgfqpoint{2.702383in}{2.027600in}}%
\pgfpathlineto{\pgfqpoint{2.703488in}{1.947901in}}%
\pgfpathlineto{\pgfqpoint{2.703501in}{1.917060in}}%
\pgfpathlineto{\pgfqpoint{2.703516in}{1.950935in}}%
\pgfpathlineto{\pgfqpoint{2.704259in}{1.929918in}}%
\pgfpathlineto{\pgfqpoint{2.704270in}{1.939516in}}%
\pgfpathlineto{\pgfqpoint{2.704896in}{1.871765in}}%
\pgfpathlineto{\pgfqpoint{2.705316in}{1.532031in}}%
\pgfpathlineto{\pgfqpoint{2.705878in}{1.532043in}}%
\pgfpathlineto{\pgfqpoint{2.707261in}{1.532135in}}%
\pgfpathlineto{\pgfqpoint{2.707324in}{2.033444in}}%
\pgfpathlineto{\pgfqpoint{2.708525in}{1.934761in}}%
\pgfpathlineto{\pgfqpoint{2.708532in}{1.918003in}}%
\pgfpathlineto{\pgfqpoint{2.708548in}{1.950406in}}%
\pgfpathlineto{\pgfqpoint{2.709296in}{1.931779in}}%
\pgfpathlineto{\pgfqpoint{2.710286in}{1.532029in}}%
\pgfpathlineto{\pgfqpoint{2.709316in}{1.938471in}}%
\pgfpathlineto{\pgfqpoint{2.710922in}{1.532051in}}%
\pgfpathlineto{\pgfqpoint{2.712208in}{1.532160in}}%
\pgfpathlineto{\pgfqpoint{2.712308in}{1.991774in}}%
\pgfpathlineto{\pgfqpoint{2.713339in}{1.944929in}}%
\pgfpathlineto{\pgfqpoint{2.713351in}{1.921269in}}%
\pgfpathlineto{\pgfqpoint{2.713367in}{1.947757in}}%
\pgfpathlineto{\pgfqpoint{2.714115in}{1.929416in}}%
\pgfpathlineto{\pgfqpoint{2.714144in}{1.940731in}}%
\pgfpathlineto{\pgfqpoint{2.714736in}{1.921690in}}%
\pgfpathlineto{\pgfqpoint{2.715202in}{1.532031in}}%
\pgfpathlineto{\pgfqpoint{2.715911in}{1.532050in}}%
\pgfpathlineto{\pgfqpoint{2.717147in}{1.532101in}}%
\pgfpathlineto{\pgfqpoint{2.717156in}{1.532315in}}%
\pgfpathlineto{\pgfqpoint{2.717289in}{2.022252in}}%
\pgfpathlineto{\pgfqpoint{2.718210in}{1.941506in}}%
\pgfpathlineto{\pgfqpoint{2.718219in}{1.913001in}}%
\pgfpathlineto{\pgfqpoint{2.718235in}{1.955452in}}%
\pgfpathlineto{\pgfqpoint{2.718978in}{1.937606in}}%
\pgfpathlineto{\pgfqpoint{2.718984in}{1.941002in}}%
\pgfpathlineto{\pgfqpoint{2.719672in}{1.921196in}}%
\pgfpathlineto{\pgfqpoint{2.720180in}{1.532028in}}%
\pgfpathlineto{\pgfqpoint{2.720915in}{1.532054in}}%
\pgfpathlineto{\pgfqpoint{2.722088in}{1.532086in}}%
\pgfpathlineto{\pgfqpoint{2.722098in}{1.532147in}}%
\pgfpathlineto{\pgfqpoint{2.722164in}{2.034441in}}%
\pgfpathlineto{\pgfqpoint{2.723274in}{1.934173in}}%
\pgfpathlineto{\pgfqpoint{2.723282in}{1.915706in}}%
\pgfpathlineto{\pgfqpoint{2.723327in}{1.952055in}}%
\pgfpathlineto{\pgfqpoint{2.724043in}{1.939700in}}%
\pgfpathlineto{\pgfqpoint{2.725128in}{1.532032in}}%
\pgfpathlineto{\pgfqpoint{2.726258in}{1.532059in}}%
\pgfpathlineto{\pgfqpoint{2.727043in}{1.532141in}}%
\pgfpathlineto{\pgfqpoint{2.727105in}{2.030090in}}%
\pgfpathlineto{\pgfqpoint{2.728280in}{1.925170in}}%
\pgfpathlineto{\pgfqpoint{2.728284in}{1.917862in}}%
\pgfpathlineto{\pgfqpoint{2.728299in}{1.951067in}}%
\pgfpathlineto{\pgfqpoint{2.729045in}{1.933296in}}%
\pgfpathlineto{\pgfqpoint{2.730038in}{1.532032in}}%
\pgfpathlineto{\pgfqpoint{2.729065in}{1.938576in}}%
\pgfpathlineto{\pgfqpoint{2.730305in}{1.532036in}}%
\pgfpathlineto{\pgfqpoint{2.731994in}{1.532597in}}%
\pgfpathlineto{\pgfqpoint{2.732095in}{2.031596in}}%
\pgfpathlineto{\pgfqpoint{2.733016in}{1.946525in}}%
\pgfpathlineto{\pgfqpoint{2.733057in}{1.911842in}}%
\pgfpathlineto{\pgfqpoint{2.733042in}{1.958172in}}%
\pgfpathlineto{\pgfqpoint{2.733787in}{1.941763in}}%
\pgfpathlineto{\pgfqpoint{2.735008in}{1.532031in}}%
\pgfpathlineto{\pgfqpoint{2.735200in}{1.532035in}}%
\pgfpathlineto{\pgfqpoint{2.736931in}{1.532111in}}%
\pgfpathlineto{\pgfqpoint{2.736973in}{1.727895in}}%
\pgfpathlineto{\pgfqpoint{2.737035in}{1.987536in}}%
\pgfpathlineto{\pgfqpoint{2.737759in}{1.953975in}}%
\pgfpathlineto{\pgfqpoint{2.737774in}{1.915438in}}%
\pgfpathlineto{\pgfqpoint{2.738538in}{1.927600in}}%
\pgfpathlineto{\pgfqpoint{2.738939in}{1.941336in}}%
\pgfpathlineto{\pgfqpoint{2.739285in}{1.939011in}}%
\pgfpathlineto{\pgfqpoint{2.739384in}{1.940233in}}%
\pgfpathlineto{\pgfqpoint{2.739631in}{1.706455in}}%
\pgfpathlineto{\pgfqpoint{2.739953in}{1.532033in}}%
\pgfpathlineto{\pgfqpoint{2.740538in}{1.532045in}}%
\pgfpathlineto{\pgfqpoint{2.741883in}{1.532303in}}%
\pgfpathlineto{\pgfqpoint{2.741943in}{2.024220in}}%
\pgfpathlineto{\pgfqpoint{2.742940in}{1.931933in}}%
\pgfpathlineto{\pgfqpoint{2.742948in}{1.912990in}}%
\pgfpathlineto{\pgfqpoint{2.742961in}{1.955513in}}%
\pgfpathlineto{\pgfqpoint{2.743706in}{1.937480in}}%
\pgfpathlineto{\pgfqpoint{2.743711in}{1.941086in}}%
\pgfpathlineto{\pgfqpoint{2.744395in}{1.921008in}}%
\pgfpathlineto{\pgfqpoint{2.744906in}{1.532032in}}%
\pgfpathlineto{\pgfqpoint{2.745642in}{1.532050in}}%
\pgfpathlineto{\pgfqpoint{2.746825in}{1.532136in}}%
\pgfpathlineto{\pgfqpoint{2.746887in}{2.030290in}}%
\pgfpathlineto{\pgfqpoint{2.748087in}{1.939423in}}%
\pgfpathlineto{\pgfqpoint{2.748095in}{1.918606in}}%
\pgfpathlineto{\pgfqpoint{2.748111in}{1.950378in}}%
\pgfpathlineto{\pgfqpoint{2.748859in}{1.932520in}}%
\pgfpathlineto{\pgfqpoint{2.749850in}{1.532029in}}%
\pgfpathlineto{\pgfqpoint{2.748879in}{1.938412in}}%
\pgfpathlineto{\pgfqpoint{2.750288in}{1.532047in}}%
\pgfpathlineto{\pgfqpoint{2.751771in}{1.532161in}}%
\pgfpathlineto{\pgfqpoint{2.751871in}{1.991421in}}%
\pgfpathlineto{\pgfqpoint{2.752909in}{1.929539in}}%
\pgfpathlineto{\pgfqpoint{2.752913in}{1.921760in}}%
\pgfpathlineto{\pgfqpoint{2.752930in}{1.947631in}}%
\pgfpathlineto{\pgfqpoint{2.753677in}{1.934067in}}%
\pgfpathlineto{\pgfqpoint{2.753793in}{1.940497in}}%
\pgfpathlineto{\pgfqpoint{2.754288in}{1.927093in}}%
\pgfpathlineto{\pgfqpoint{2.754780in}{1.532031in}}%
\pgfpathlineto{\pgfqpoint{2.755555in}{1.532048in}}%
\pgfpathlineto{\pgfqpoint{2.756715in}{1.532128in}}%
\pgfpathlineto{\pgfqpoint{2.756779in}{2.030680in}}%
\pgfpathlineto{\pgfqpoint{2.758011in}{1.928725in}}%
\pgfpathlineto{\pgfqpoint{2.758015in}{1.920166in}}%
\pgfpathlineto{\pgfqpoint{2.758033in}{1.949284in}}%
\pgfpathlineto{\pgfqpoint{2.758778in}{1.933868in}}%
\pgfpathlineto{\pgfqpoint{2.759717in}{1.532033in}}%
\pgfpathlineto{\pgfqpoint{2.758799in}{1.938146in}}%
\pgfpathlineto{\pgfqpoint{2.760088in}{1.532040in}}%
\pgfpathlineto{\pgfqpoint{2.761663in}{1.532180in}}%
\pgfpathlineto{\pgfqpoint{2.761727in}{2.034114in}}%
\pgfpathlineto{\pgfqpoint{2.762807in}{1.931760in}}%
\pgfpathlineto{\pgfqpoint{2.762815in}{1.914834in}}%
\pgfpathlineto{\pgfqpoint{2.762859in}{1.952898in}}%
\pgfpathlineto{\pgfqpoint{2.763579in}{1.929404in}}%
\pgfpathlineto{\pgfqpoint{2.763596in}{1.939692in}}%
\pgfpathlineto{\pgfqpoint{2.764172in}{1.930874in}}%
\pgfpathlineto{\pgfqpoint{2.764686in}{1.532029in}}%
\pgfpathlineto{\pgfqpoint{2.766014in}{1.532066in}}%
\pgfpathlineto{\pgfqpoint{2.766608in}{1.532157in}}%
\pgfpathlineto{\pgfqpoint{2.766711in}{1.992666in}}%
\pgfpathlineto{\pgfqpoint{2.767745in}{1.924552in}}%
\pgfpathlineto{\pgfqpoint{2.767751in}{1.921622in}}%
\pgfpathlineto{\pgfqpoint{2.767766in}{1.946963in}}%
\pgfpathlineto{\pgfqpoint{2.768498in}{1.930031in}}%
\pgfpathlineto{\pgfqpoint{2.768872in}{1.938956in}}%
\pgfpathlineto{\pgfqpoint{2.769169in}{1.892124in}}%
\pgfpathlineto{\pgfqpoint{2.769626in}{1.532029in}}%
\pgfpathlineto{\pgfqpoint{2.770212in}{1.532042in}}%
\pgfpathlineto{\pgfqpoint{2.771552in}{1.532133in}}%
\pgfpathlineto{\pgfqpoint{2.771617in}{2.039050in}}%
\pgfpathlineto{\pgfqpoint{2.772784in}{1.944519in}}%
\pgfpathlineto{\pgfqpoint{2.772824in}{1.917171in}}%
\pgfpathlineto{\pgfqpoint{2.772809in}{1.952311in}}%
\pgfpathlineto{\pgfqpoint{2.773554in}{1.938425in}}%
\pgfpathlineto{\pgfqpoint{2.774576in}{1.532029in}}%
\pgfpathlineto{\pgfqpoint{2.773584in}{1.938638in}}%
\pgfpathlineto{\pgfqpoint{2.774767in}{1.532033in}}%
\pgfpathlineto{\pgfqpoint{2.776498in}{1.532155in}}%
\pgfpathlineto{\pgfqpoint{2.776602in}{1.996810in}}%
\pgfpathlineto{\pgfqpoint{2.777658in}{1.939585in}}%
\pgfpathlineto{\pgfqpoint{2.777669in}{1.921265in}}%
\pgfpathlineto{\pgfqpoint{2.777685in}{1.947246in}}%
\pgfpathlineto{\pgfqpoint{2.778423in}{1.939421in}}%
\pgfpathlineto{\pgfqpoint{2.779503in}{1.532029in}}%
\pgfpathlineto{\pgfqpoint{2.779682in}{1.532033in}}%
\pgfpathlineto{\pgfqpoint{2.781441in}{1.532121in}}%
\pgfpathlineto{\pgfqpoint{2.781483in}{1.768092in}}%
\pgfpathlineto{\pgfqpoint{2.781541in}{1.995982in}}%
\pgfpathlineto{\pgfqpoint{2.782263in}{1.941467in}}%
\pgfpathlineto{\pgfqpoint{2.782272in}{1.912455in}}%
\pgfpathlineto{\pgfqpoint{2.782315in}{1.954389in}}%
\pgfpathlineto{\pgfqpoint{2.783027in}{1.941292in}}%
\pgfpathlineto{\pgfqpoint{2.783082in}{1.927306in}}%
\pgfpathlineto{\pgfqpoint{2.783096in}{1.941396in}}%
\pgfpathlineto{\pgfqpoint{2.783793in}{1.937061in}}%
\pgfpathlineto{\pgfqpoint{2.784040in}{1.839073in}}%
\pgfpathlineto{\pgfqpoint{2.784410in}{1.532025in}}%
\pgfpathlineto{\pgfqpoint{2.784992in}{1.532039in}}%
\pgfpathlineto{\pgfqpoint{2.786387in}{1.532123in}}%
\pgfpathlineto{\pgfqpoint{2.786427in}{1.734442in}}%
\pgfpathlineto{\pgfqpoint{2.786487in}{1.990961in}}%
\pgfpathlineto{\pgfqpoint{2.787209in}{1.955439in}}%
\pgfpathlineto{\pgfqpoint{2.787224in}{1.913964in}}%
\pgfpathlineto{\pgfqpoint{2.787977in}{1.934005in}}%
\pgfpathlineto{\pgfqpoint{2.787987in}{1.942623in}}%
\pgfpathlineto{\pgfqpoint{2.788006in}{1.926708in}}%
\pgfpathlineto{\pgfqpoint{2.788743in}{1.932039in}}%
\pgfpathlineto{\pgfqpoint{2.788891in}{1.938401in}}%
\pgfpathlineto{\pgfqpoint{2.788941in}{1.908168in}}%
\pgfpathlineto{\pgfqpoint{2.789384in}{1.532031in}}%
\pgfpathlineto{\pgfqpoint{2.789995in}{1.532048in}}%
\pgfpathlineto{\pgfqpoint{2.791329in}{1.532102in}}%
\pgfpathlineto{\pgfqpoint{2.791338in}{1.532326in}}%
\pgfpathlineto{\pgfqpoint{2.791471in}{2.022831in}}%
\pgfpathlineto{\pgfqpoint{2.792394in}{1.934252in}}%
\pgfpathlineto{\pgfqpoint{2.792402in}{1.912946in}}%
\pgfpathlineto{\pgfqpoint{2.792447in}{1.954548in}}%
\pgfpathlineto{\pgfqpoint{2.793163in}{1.940945in}}%
\pgfpathlineto{\pgfqpoint{2.794301in}{1.532013in}}%
\pgfpathlineto{\pgfqpoint{2.794660in}{1.532031in}}%
\pgfpathlineto{\pgfqpoint{2.796283in}{1.532240in}}%
\pgfpathlineto{\pgfqpoint{2.796342in}{2.025878in}}%
\pgfpathlineto{\pgfqpoint{2.797371in}{1.919810in}}%
\pgfpathlineto{\pgfqpoint{2.797375in}{1.913884in}}%
\pgfpathlineto{\pgfqpoint{2.797391in}{1.954400in}}%
\pgfpathlineto{\pgfqpoint{2.798126in}{1.940033in}}%
\pgfpathlineto{\pgfqpoint{2.798857in}{1.873376in}}%
\pgfpathlineto{\pgfqpoint{2.799307in}{1.532033in}}%
\pgfpathlineto{\pgfqpoint{2.799846in}{1.532045in}}%
\pgfpathlineto{\pgfqpoint{2.801227in}{1.532185in}}%
\pgfpathlineto{\pgfqpoint{2.801289in}{2.031421in}}%
\pgfpathlineto{\pgfqpoint{2.802370in}{1.942822in}}%
\pgfpathlineto{\pgfqpoint{2.802378in}{1.916039in}}%
\pgfpathlineto{\pgfqpoint{2.802395in}{1.953677in}}%
\pgfpathlineto{\pgfqpoint{2.803144in}{1.930588in}}%
\pgfpathlineto{\pgfqpoint{2.804252in}{1.532028in}}%
\pgfpathlineto{\pgfqpoint{2.803163in}{1.939442in}}%
\pgfpathlineto{\pgfqpoint{2.804887in}{1.532051in}}%
\pgfpathlineto{\pgfqpoint{2.806170in}{1.532141in}}%
\pgfpathlineto{\pgfqpoint{2.806237in}{2.036566in}}%
\pgfpathlineto{\pgfqpoint{2.807347in}{1.934041in}}%
\pgfpathlineto{\pgfqpoint{2.807355in}{1.915294in}}%
\pgfpathlineto{\pgfqpoint{2.807399in}{1.952365in}}%
\pgfpathlineto{\pgfqpoint{2.808115in}{1.939942in}}%
\pgfpathlineto{\pgfqpoint{2.809186in}{1.532033in}}%
\pgfpathlineto{\pgfqpoint{2.809772in}{1.532045in}}%
\pgfpathlineto{\pgfqpoint{2.811118in}{1.532189in}}%
\pgfpathlineto{\pgfqpoint{2.811180in}{2.033107in}}%
\pgfpathlineto{\pgfqpoint{2.812258in}{1.945647in}}%
\pgfpathlineto{\pgfqpoint{2.812298in}{1.915471in}}%
\pgfpathlineto{\pgfqpoint{2.812283in}{1.954087in}}%
\pgfpathlineto{\pgfqpoint{2.813031in}{1.929898in}}%
\pgfpathlineto{\pgfqpoint{2.814142in}{1.532028in}}%
\pgfpathlineto{\pgfqpoint{2.813048in}{1.939501in}}%
\pgfpathlineto{\pgfqpoint{2.815570in}{1.532069in}}%
\pgfpathlineto{\pgfqpoint{2.816061in}{1.532141in}}%
\pgfpathlineto{\pgfqpoint{2.816128in}{2.035419in}}%
\pgfpathlineto{\pgfqpoint{2.817238in}{1.932122in}}%
\pgfpathlineto{\pgfqpoint{2.817246in}{1.915688in}}%
\pgfpathlineto{\pgfqpoint{2.817259in}{1.953069in}}%
\pgfpathlineto{\pgfqpoint{2.818005in}{1.937117in}}%
\pgfpathlineto{\pgfqpoint{2.818010in}{1.939967in}}%
\pgfpathlineto{\pgfqpoint{2.818616in}{1.899090in}}%
\pgfpathlineto{\pgfqpoint{2.819071in}{1.532033in}}%
\pgfpathlineto{\pgfqpoint{2.819649in}{1.532045in}}%
\pgfpathlineto{\pgfqpoint{2.821009in}{1.532189in}}%
\pgfpathlineto{\pgfqpoint{2.821071in}{2.033077in}}%
\pgfpathlineto{\pgfqpoint{2.822128in}{1.916556in}}%
\pgfpathlineto{\pgfqpoint{2.822132in}{1.914302in}}%
\pgfpathlineto{\pgfqpoint{2.822144in}{1.954825in}}%
\pgfpathlineto{\pgfqpoint{2.822869in}{1.928900in}}%
\pgfpathlineto{\pgfqpoint{2.822883in}{1.940341in}}%
\pgfpathlineto{\pgfqpoint{2.823569in}{1.894865in}}%
\pgfpathlineto{\pgfqpoint{2.823973in}{1.532018in}}%
\pgfpathlineto{\pgfqpoint{2.824609in}{1.532051in}}%
\pgfpathlineto{\pgfqpoint{2.825950in}{1.532120in}}%
\pgfpathlineto{\pgfqpoint{2.825991in}{1.729901in}}%
\pgfpathlineto{\pgfqpoint{2.826051in}{1.985401in}}%
\pgfpathlineto{\pgfqpoint{2.826774in}{1.954624in}}%
\pgfpathlineto{\pgfqpoint{2.826790in}{1.914685in}}%
\pgfpathlineto{\pgfqpoint{2.827550in}{1.942272in}}%
\pgfpathlineto{\pgfqpoint{2.827569in}{1.927074in}}%
\pgfpathlineto{\pgfqpoint{2.828283in}{1.928417in}}%
\pgfpathlineto{\pgfqpoint{2.828332in}{1.939677in}}%
\pgfpathlineto{\pgfqpoint{2.828530in}{1.882886in}}%
\pgfpathlineto{\pgfqpoint{2.828961in}{1.532030in}}%
\pgfpathlineto{\pgfqpoint{2.829538in}{1.532043in}}%
\pgfpathlineto{\pgfqpoint{2.830901in}{1.532308in}}%
\pgfpathlineto{\pgfqpoint{2.830961in}{2.024467in}}%
\pgfpathlineto{\pgfqpoint{2.831958in}{1.932752in}}%
\pgfpathlineto{\pgfqpoint{2.831966in}{1.912869in}}%
\pgfpathlineto{\pgfqpoint{2.831979in}{1.955133in}}%
\pgfpathlineto{\pgfqpoint{2.832728in}{1.940979in}}%
\pgfpathlineto{\pgfqpoint{2.833864in}{1.532024in}}%
\pgfpathlineto{\pgfqpoint{2.834146in}{1.532032in}}%
\pgfpathlineto{\pgfqpoint{2.835841in}{1.532122in}}%
\pgfpathlineto{\pgfqpoint{2.835881in}{1.729262in}}%
\pgfpathlineto{\pgfqpoint{2.835942in}{1.986211in}}%
\pgfpathlineto{\pgfqpoint{2.836666in}{1.947805in}}%
\pgfpathlineto{\pgfqpoint{2.836679in}{1.914402in}}%
\pgfpathlineto{\pgfqpoint{2.836693in}{1.954294in}}%
\pgfpathlineto{\pgfqpoint{2.837441in}{1.932716in}}%
\pgfpathlineto{\pgfqpoint{2.837448in}{1.926796in}}%
\pgfpathlineto{\pgfqpoint{2.837466in}{1.942117in}}%
\pgfpathlineto{\pgfqpoint{2.838197in}{1.932638in}}%
\pgfpathlineto{\pgfqpoint{2.838296in}{1.935927in}}%
\pgfpathlineto{\pgfqpoint{2.838811in}{1.532023in}}%
\pgfpathlineto{\pgfqpoint{2.839860in}{1.532047in}}%
\pgfpathlineto{\pgfqpoint{2.840786in}{1.532115in}}%
\pgfpathlineto{\pgfqpoint{2.840828in}{1.738984in}}%
\pgfpathlineto{\pgfqpoint{2.840889in}{1.992906in}}%
\pgfpathlineto{\pgfqpoint{2.841609in}{1.954991in}}%
\pgfpathlineto{\pgfqpoint{2.841622in}{1.913810in}}%
\pgfpathlineto{\pgfqpoint{2.842385in}{1.935485in}}%
\pgfpathlineto{\pgfqpoint{2.842395in}{1.926621in}}%
\pgfpathlineto{\pgfqpoint{2.842413in}{1.942247in}}%
\pgfpathlineto{\pgfqpoint{2.843146in}{1.936217in}}%
\pgfpathlineto{\pgfqpoint{2.843344in}{1.894182in}}%
\pgfpathlineto{\pgfqpoint{2.843755in}{1.532012in}}%
\pgfpathlineto{\pgfqpoint{2.844388in}{1.532037in}}%
\pgfpathlineto{\pgfqpoint{2.845739in}{1.533056in}}%
\pgfpathlineto{\pgfqpoint{2.845841in}{2.031654in}}%
\pgfpathlineto{\pgfqpoint{2.846759in}{1.958188in}}%
\pgfpathlineto{\pgfqpoint{2.846772in}{1.910362in}}%
\pgfpathlineto{\pgfqpoint{2.847534in}{1.938631in}}%
\pgfpathlineto{\pgfqpoint{2.847540in}{1.941761in}}%
\pgfpathlineto{\pgfqpoint{2.848264in}{1.916147in}}%
\pgfpathlineto{\pgfqpoint{2.848764in}{1.532028in}}%
\pgfpathlineto{\pgfqpoint{2.849450in}{1.532043in}}%
\pgfpathlineto{\pgfqpoint{2.850680in}{1.532146in}}%
\pgfpathlineto{\pgfqpoint{2.850747in}{2.018596in}}%
\pgfpathlineto{\pgfqpoint{2.851857in}{1.940067in}}%
\pgfpathlineto{\pgfqpoint{2.851866in}{1.917799in}}%
\pgfpathlineto{\pgfqpoint{2.851883in}{1.949542in}}%
\pgfpathlineto{\pgfqpoint{2.852629in}{1.934347in}}%
\pgfpathlineto{\pgfqpoint{2.853693in}{1.532031in}}%
\pgfpathlineto{\pgfqpoint{2.852724in}{1.938910in}}%
\pgfpathlineto{\pgfqpoint{2.853975in}{1.532037in}}%
\pgfpathlineto{\pgfqpoint{2.855623in}{1.532120in}}%
\pgfpathlineto{\pgfqpoint{2.855687in}{2.030307in}}%
\pgfpathlineto{\pgfqpoint{2.856982in}{1.925425in}}%
\pgfpathlineto{\pgfqpoint{2.856986in}{1.920536in}}%
\pgfpathlineto{\pgfqpoint{2.857001in}{1.948111in}}%
\pgfpathlineto{\pgfqpoint{2.857745in}{1.934389in}}%
\pgfpathlineto{\pgfqpoint{2.858631in}{1.532031in}}%
\pgfpathlineto{\pgfqpoint{2.857769in}{1.937433in}}%
\pgfpathlineto{\pgfqpoint{2.859106in}{1.532041in}}%
\pgfpathlineto{\pgfqpoint{2.860567in}{1.532108in}}%
\pgfpathlineto{\pgfqpoint{2.860608in}{1.717552in}}%
\pgfpathlineto{\pgfqpoint{2.860637in}{2.007025in}}%
\pgfpathlineto{\pgfqpoint{2.861392in}{1.953678in}}%
\pgfpathlineto{\pgfqpoint{2.861403in}{1.910894in}}%
\pgfpathlineto{\pgfqpoint{2.861418in}{1.957365in}}%
\pgfpathlineto{\pgfqpoint{2.862165in}{1.942682in}}%
\pgfpathlineto{\pgfqpoint{2.862257in}{1.926523in}}%
\pgfpathlineto{\pgfqpoint{2.862913in}{1.931553in}}%
\pgfpathlineto{\pgfqpoint{2.863012in}{1.935196in}}%
\pgfpathlineto{\pgfqpoint{2.863062in}{1.911183in}}%
\pgfpathlineto{\pgfqpoint{2.863210in}{1.795275in}}%
\pgfpathlineto{\pgfqpoint{2.863586in}{1.532033in}}%
\pgfpathlineto{\pgfqpoint{2.864168in}{1.532045in}}%
\pgfpathlineto{\pgfqpoint{2.865520in}{1.532305in}}%
\pgfpathlineto{\pgfqpoint{2.865579in}{2.024244in}}%
\pgfpathlineto{\pgfqpoint{2.866577in}{1.931292in}}%
\pgfpathlineto{\pgfqpoint{2.866584in}{1.913007in}}%
\pgfpathlineto{\pgfqpoint{2.866597in}{1.955641in}}%
\pgfpathlineto{\pgfqpoint{2.867341in}{1.937670in}}%
\pgfpathlineto{\pgfqpoint{2.867346in}{1.940954in}}%
\pgfpathlineto{\pgfqpoint{2.868033in}{1.917800in}}%
\pgfpathlineto{\pgfqpoint{2.868501in}{1.532027in}}%
\pgfpathlineto{\pgfqpoint{2.869235in}{1.532042in}}%
\pgfpathlineto{\pgfqpoint{2.870462in}{1.532159in}}%
\pgfpathlineto{\pgfqpoint{2.870518in}{2.004632in}}%
\pgfpathlineto{\pgfqpoint{2.871594in}{1.948160in}}%
\pgfpathlineto{\pgfqpoint{2.871637in}{1.919458in}}%
\pgfpathlineto{\pgfqpoint{2.871621in}{1.949465in}}%
\pgfpathlineto{\pgfqpoint{2.872364in}{1.927574in}}%
\pgfpathlineto{\pgfqpoint{2.872624in}{1.940559in}}%
\pgfpathlineto{\pgfqpoint{2.873020in}{1.907453in}}%
\pgfpathlineto{\pgfqpoint{2.873446in}{1.532026in}}%
\pgfpathlineto{\pgfqpoint{2.874081in}{1.532040in}}%
\pgfpathlineto{\pgfqpoint{2.875407in}{1.532153in}}%
\pgfpathlineto{\pgfqpoint{2.875462in}{1.991467in}}%
\pgfpathlineto{\pgfqpoint{2.876542in}{1.941074in}}%
\pgfpathlineto{\pgfqpoint{2.876554in}{1.921022in}}%
\pgfpathlineto{\pgfqpoint{2.876570in}{1.946974in}}%
\pgfpathlineto{\pgfqpoint{2.877309in}{1.928811in}}%
\pgfpathlineto{\pgfqpoint{2.877552in}{1.940552in}}%
\pgfpathlineto{\pgfqpoint{2.877948in}{1.905395in}}%
\pgfpathlineto{\pgfqpoint{2.878428in}{1.532027in}}%
\pgfpathlineto{\pgfqpoint{2.879063in}{1.532052in}}%
\pgfpathlineto{\pgfqpoint{2.880342in}{1.532084in}}%
\pgfpathlineto{\pgfqpoint{2.880354in}{1.532175in}}%
\pgfpathlineto{\pgfqpoint{2.880409in}{2.007874in}}%
\pgfpathlineto{\pgfqpoint{2.881456in}{1.947015in}}%
\pgfpathlineto{\pgfqpoint{2.881467in}{1.917547in}}%
\pgfpathlineto{\pgfqpoint{2.881483in}{1.951136in}}%
\pgfpathlineto{\pgfqpoint{2.882224in}{1.939558in}}%
\pgfpathlineto{\pgfqpoint{2.882350in}{1.941695in}}%
\pgfpathlineto{\pgfqpoint{2.883418in}{1.532033in}}%
\pgfpathlineto{\pgfqpoint{2.885298in}{1.532141in}}%
\pgfpathlineto{\pgfqpoint{2.885360in}{2.029562in}}%
\pgfpathlineto{\pgfqpoint{2.886533in}{1.929408in}}%
\pgfpathlineto{\pgfqpoint{2.886569in}{1.918551in}}%
\pgfpathlineto{\pgfqpoint{2.886554in}{1.951054in}}%
\pgfpathlineto{\pgfqpoint{2.887303in}{1.930872in}}%
\pgfpathlineto{\pgfqpoint{2.888330in}{1.532027in}}%
\pgfpathlineto{\pgfqpoint{2.887322in}{1.938393in}}%
\pgfpathlineto{\pgfqpoint{2.889757in}{1.532070in}}%
\pgfpathlineto{\pgfqpoint{2.890233in}{1.532085in}}%
\pgfpathlineto{\pgfqpoint{2.890243in}{1.532147in}}%
\pgfpathlineto{\pgfqpoint{2.890312in}{2.001983in}}%
\pgfpathlineto{\pgfqpoint{2.891400in}{1.927344in}}%
\pgfpathlineto{\pgfqpoint{2.891405in}{1.919414in}}%
\pgfpathlineto{\pgfqpoint{2.891421in}{1.949269in}}%
\pgfpathlineto{\pgfqpoint{2.892158in}{1.936632in}}%
\pgfpathlineto{\pgfqpoint{2.892165in}{1.939827in}}%
\pgfpathlineto{\pgfqpoint{2.892764in}{1.904644in}}%
\pgfpathlineto{\pgfqpoint{2.893270in}{1.532032in}}%
\pgfpathlineto{\pgfqpoint{2.894105in}{1.532052in}}%
\pgfpathlineto{\pgfqpoint{2.895188in}{1.532138in}}%
\pgfpathlineto{\pgfqpoint{2.895251in}{2.029666in}}%
\pgfpathlineto{\pgfqpoint{2.896428in}{1.919240in}}%
\pgfpathlineto{\pgfqpoint{2.896460in}{1.918752in}}%
\pgfpathlineto{\pgfqpoint{2.896445in}{1.950844in}}%
\pgfpathlineto{\pgfqpoint{2.896914in}{1.940127in}}%
\pgfpathlineto{\pgfqpoint{2.896918in}{1.941893in}}%
\pgfpathlineto{\pgfqpoint{2.896961in}{1.927525in}}%
\pgfpathlineto{\pgfqpoint{2.897668in}{1.932529in}}%
\pgfpathlineto{\pgfqpoint{2.897745in}{1.887459in}}%
\pgfpathlineto{\pgfqpoint{2.898155in}{1.532026in}}%
\pgfpathlineto{\pgfqpoint{2.898758in}{1.532041in}}%
\pgfpathlineto{\pgfqpoint{2.900132in}{1.532122in}}%
\pgfpathlineto{\pgfqpoint{2.900174in}{1.772988in}}%
\pgfpathlineto{\pgfqpoint{2.900232in}{1.997949in}}%
\pgfpathlineto{\pgfqpoint{2.900955in}{1.928618in}}%
\pgfpathlineto{\pgfqpoint{2.900962in}{1.913294in}}%
\pgfpathlineto{\pgfqpoint{2.900975in}{1.956073in}}%
\pgfpathlineto{\pgfqpoint{2.901720in}{1.939358in}}%
\pgfpathlineto{\pgfqpoint{2.901800in}{1.927829in}}%
\pgfpathlineto{\pgfqpoint{2.901748in}{1.941957in}}%
\pgfpathlineto{\pgfqpoint{2.902476in}{1.937713in}}%
\pgfpathlineto{\pgfqpoint{2.902575in}{1.937721in}}%
\pgfpathlineto{\pgfqpoint{2.902773in}{1.794400in}}%
\pgfpathlineto{\pgfqpoint{2.903152in}{1.532032in}}%
\pgfpathlineto{\pgfqpoint{2.903690in}{1.532045in}}%
\pgfpathlineto{\pgfqpoint{2.905081in}{1.532174in}}%
\pgfpathlineto{\pgfqpoint{2.905144in}{2.031839in}}%
\pgfpathlineto{\pgfqpoint{2.906228in}{1.929323in}}%
\pgfpathlineto{\pgfqpoint{2.906236in}{1.915766in}}%
\pgfpathlineto{\pgfqpoint{2.906248in}{1.953446in}}%
\pgfpathlineto{\pgfqpoint{2.906999in}{1.929457in}}%
\pgfpathlineto{\pgfqpoint{2.907003in}{1.929432in}}%
\pgfpathlineto{\pgfqpoint{2.907008in}{1.933725in}}%
\pgfpathlineto{\pgfqpoint{2.907016in}{1.939571in}}%
\pgfpathlineto{\pgfqpoint{2.907676in}{1.836919in}}%
\pgfpathlineto{\pgfqpoint{2.908078in}{1.532030in}}%
\pgfpathlineto{\pgfqpoint{2.908643in}{1.532042in}}%
\pgfpathlineto{\pgfqpoint{2.910029in}{1.532309in}}%
\pgfpathlineto{\pgfqpoint{2.910088in}{2.024701in}}%
\pgfpathlineto{\pgfqpoint{2.911083in}{1.945093in}}%
\pgfpathlineto{\pgfqpoint{2.911094in}{1.912961in}}%
\pgfpathlineto{\pgfqpoint{2.911107in}{1.956052in}}%
\pgfpathlineto{\pgfqpoint{2.911851in}{1.937489in}}%
\pgfpathlineto{\pgfqpoint{2.911857in}{1.941008in}}%
\pgfpathlineto{\pgfqpoint{2.912564in}{1.918591in}}%
\pgfpathlineto{\pgfqpoint{2.913027in}{1.532031in}}%
\pgfpathlineto{\pgfqpoint{2.913695in}{1.532045in}}%
\pgfpathlineto{\pgfqpoint{2.914974in}{1.532307in}}%
\pgfpathlineto{\pgfqpoint{2.915034in}{2.024397in}}%
\pgfpathlineto{\pgfqpoint{2.916031in}{1.930959in}}%
\pgfpathlineto{\pgfqpoint{2.916039in}{1.912901in}}%
\pgfpathlineto{\pgfqpoint{2.916052in}{1.955971in}}%
\pgfpathlineto{\pgfqpoint{2.916798in}{1.937629in}}%
\pgfpathlineto{\pgfqpoint{2.916804in}{1.941062in}}%
\pgfpathlineto{\pgfqpoint{2.917509in}{1.911913in}}%
\pgfpathlineto{\pgfqpoint{2.918003in}{1.532028in}}%
\pgfpathlineto{\pgfqpoint{2.918590in}{1.532041in}}%
\pgfpathlineto{\pgfqpoint{2.919916in}{1.532146in}}%
\pgfpathlineto{\pgfqpoint{2.919984in}{2.011563in}}%
\pgfpathlineto{\pgfqpoint{2.921072in}{1.919025in}}%
\pgfpathlineto{\pgfqpoint{2.921087in}{1.949305in}}%
\pgfpathlineto{\pgfqpoint{2.921874in}{1.939556in}}%
\pgfpathlineto{\pgfqpoint{2.922946in}{1.532027in}}%
\pgfpathlineto{\pgfqpoint{2.923434in}{1.532038in}}%
\pgfpathlineto{\pgfqpoint{2.924862in}{1.532146in}}%
\pgfpathlineto{\pgfqpoint{2.924929in}{2.015989in}}%
\pgfpathlineto{\pgfqpoint{2.926036in}{1.947105in}}%
\pgfpathlineto{\pgfqpoint{2.926046in}{1.919219in}}%
\pgfpathlineto{\pgfqpoint{2.926093in}{1.949726in}}%
\pgfpathlineto{\pgfqpoint{2.926806in}{1.938895in}}%
\pgfpathlineto{\pgfqpoint{2.926844in}{1.939306in}}%
\pgfpathlineto{\pgfqpoint{2.927484in}{1.800191in}}%
\pgfpathlineto{\pgfqpoint{2.927846in}{1.532029in}}%
\pgfpathlineto{\pgfqpoint{2.928424in}{1.532048in}}%
\pgfpathlineto{\pgfqpoint{2.929806in}{1.532135in}}%
\pgfpathlineto{\pgfqpoint{2.929872in}{2.040314in}}%
\pgfpathlineto{\pgfqpoint{2.931016in}{1.920209in}}%
\pgfpathlineto{\pgfqpoint{2.931020in}{1.915244in}}%
\pgfpathlineto{\pgfqpoint{2.931065in}{1.952273in}}%
\pgfpathlineto{\pgfqpoint{2.931774in}{1.936416in}}%
\pgfpathlineto{\pgfqpoint{2.931779in}{1.939747in}}%
\pgfpathlineto{\pgfqpoint{2.932360in}{1.897531in}}%
\pgfpathlineto{\pgfqpoint{2.932774in}{1.532011in}}%
\pgfpathlineto{\pgfqpoint{2.933371in}{1.532053in}}%
\pgfpathlineto{\pgfqpoint{2.934747in}{1.532100in}}%
\pgfpathlineto{\pgfqpoint{2.934755in}{1.532221in}}%
\pgfpathlineto{\pgfqpoint{2.934815in}{2.026300in}}%
\pgfpathlineto{\pgfqpoint{2.935867in}{1.939710in}}%
\pgfpathlineto{\pgfqpoint{2.935875in}{1.914782in}}%
\pgfpathlineto{\pgfqpoint{2.935890in}{1.953902in}}%
\pgfpathlineto{\pgfqpoint{2.936641in}{1.929273in}}%
\pgfpathlineto{\pgfqpoint{2.936658in}{1.939580in}}%
\pgfpathlineto{\pgfqpoint{2.937278in}{1.925636in}}%
\pgfpathlineto{\pgfqpoint{2.937774in}{1.532029in}}%
\pgfpathlineto{\pgfqpoint{2.938608in}{1.532055in}}%
\pgfpathlineto{\pgfqpoint{2.939698in}{1.532145in}}%
\pgfpathlineto{\pgfqpoint{2.939765in}{2.024784in}}%
\pgfpathlineto{\pgfqpoint{2.940874in}{1.938290in}}%
\pgfpathlineto{\pgfqpoint{2.940882in}{1.917025in}}%
\pgfpathlineto{\pgfqpoint{2.940926in}{1.950335in}}%
\pgfpathlineto{\pgfqpoint{2.941646in}{1.929115in}}%
\pgfpathlineto{\pgfqpoint{2.941726in}{1.939070in}}%
\pgfpathlineto{\pgfqpoint{2.942266in}{1.882063in}}%
\pgfpathlineto{\pgfqpoint{2.942728in}{1.532027in}}%
\pgfpathlineto{\pgfqpoint{2.943265in}{1.532050in}}%
\pgfpathlineto{\pgfqpoint{2.944633in}{1.532085in}}%
\pgfpathlineto{\pgfqpoint{2.944643in}{1.532146in}}%
\pgfpathlineto{\pgfqpoint{2.944711in}{2.011825in}}%
\pgfpathlineto{\pgfqpoint{2.945799in}{1.919053in}}%
\pgfpathlineto{\pgfqpoint{2.945848in}{1.949223in}}%
\pgfpathlineto{\pgfqpoint{2.946596in}{1.939239in}}%
\pgfpathlineto{\pgfqpoint{2.947693in}{1.532032in}}%
\pgfpathlineto{\pgfqpoint{2.946636in}{1.939275in}}%
\pgfpathlineto{\pgfqpoint{2.948479in}{1.532050in}}%
\pgfpathlineto{\pgfqpoint{2.949587in}{1.532127in}}%
\pgfpathlineto{\pgfqpoint{2.949651in}{2.029632in}}%
\pgfpathlineto{\pgfqpoint{2.950885in}{1.922797in}}%
\pgfpathlineto{\pgfqpoint{2.950889in}{1.919327in}}%
\pgfpathlineto{\pgfqpoint{2.950932in}{1.948642in}}%
\pgfpathlineto{\pgfqpoint{2.951640in}{1.938251in}}%
\pgfpathlineto{\pgfqpoint{2.952575in}{1.532026in}}%
\pgfpathlineto{\pgfqpoint{2.953561in}{1.532061in}}%
\pgfpathlineto{\pgfqpoint{2.954522in}{1.532082in}}%
\pgfpathlineto{\pgfqpoint{2.954533in}{1.532126in}}%
\pgfpathlineto{\pgfqpoint{2.954597in}{2.031501in}}%
\pgfpathlineto{\pgfqpoint{2.955831in}{1.927388in}}%
\pgfpathlineto{\pgfqpoint{2.955835in}{1.919581in}}%
\pgfpathlineto{\pgfqpoint{2.955852in}{1.949560in}}%
\pgfpathlineto{\pgfqpoint{2.956599in}{1.933029in}}%
\pgfpathlineto{\pgfqpoint{2.957529in}{1.532031in}}%
\pgfpathlineto{\pgfqpoint{2.956619in}{1.938015in}}%
\pgfpathlineto{\pgfqpoint{2.957941in}{1.532043in}}%
\pgfpathlineto{\pgfqpoint{2.959475in}{1.532102in}}%
\pgfpathlineto{\pgfqpoint{2.959483in}{1.532321in}}%
\pgfpathlineto{\pgfqpoint{2.959616in}{2.022439in}}%
\pgfpathlineto{\pgfqpoint{2.960540in}{1.932067in}}%
\pgfpathlineto{\pgfqpoint{2.960548in}{1.913068in}}%
\pgfpathlineto{\pgfqpoint{2.960561in}{1.955098in}}%
\pgfpathlineto{\pgfqpoint{2.961309in}{1.940871in}}%
\pgfpathlineto{\pgfqpoint{2.962509in}{1.532028in}}%
\pgfpathlineto{\pgfqpoint{2.962700in}{1.532032in}}%
\pgfpathlineto{\pgfqpoint{2.964425in}{1.532146in}}%
\pgfpathlineto{\pgfqpoint{2.964492in}{2.021890in}}%
\pgfpathlineto{\pgfqpoint{2.965599in}{1.948256in}}%
\pgfpathlineto{\pgfqpoint{2.965613in}{1.918744in}}%
\pgfpathlineto{\pgfqpoint{2.965626in}{1.951331in}}%
\pgfpathlineto{\pgfqpoint{2.966371in}{1.930742in}}%
\pgfpathlineto{\pgfqpoint{2.967448in}{1.532031in}}%
\pgfpathlineto{\pgfqpoint{2.966425in}{1.939207in}}%
\pgfpathlineto{\pgfqpoint{2.968483in}{1.532058in}}%
\pgfpathlineto{\pgfqpoint{2.969367in}{1.532108in}}%
\pgfpathlineto{\pgfqpoint{2.969408in}{1.714793in}}%
\pgfpathlineto{\pgfqpoint{2.969439in}{2.006662in}}%
\pgfpathlineto{\pgfqpoint{2.970194in}{1.950134in}}%
\pgfpathlineto{\pgfqpoint{2.970233in}{1.911719in}}%
\pgfpathlineto{\pgfqpoint{2.970220in}{1.957887in}}%
\pgfpathlineto{\pgfqpoint{2.970960in}{1.937380in}}%
\pgfpathlineto{\pgfqpoint{2.970968in}{1.942943in}}%
\pgfpathlineto{\pgfqpoint{2.970981in}{1.926387in}}%
\pgfpathlineto{\pgfqpoint{2.971714in}{1.932377in}}%
\pgfpathlineto{\pgfqpoint{2.971862in}{1.938069in}}%
\pgfpathlineto{\pgfqpoint{2.971911in}{1.913391in}}%
\pgfpathlineto{\pgfqpoint{2.972370in}{1.532030in}}%
\pgfpathlineto{\pgfqpoint{2.973037in}{1.532044in}}%
\pgfpathlineto{\pgfqpoint{2.974314in}{1.532117in}}%
\pgfpathlineto{\pgfqpoint{2.974356in}{1.765839in}}%
\pgfpathlineto{\pgfqpoint{2.974413in}{1.989917in}}%
\pgfpathlineto{\pgfqpoint{2.975133in}{1.950797in}}%
\pgfpathlineto{\pgfqpoint{2.975143in}{1.913651in}}%
\pgfpathlineto{\pgfqpoint{2.975159in}{1.954680in}}%
\pgfpathlineto{\pgfqpoint{2.975902in}{1.930141in}}%
\pgfpathlineto{\pgfqpoint{2.975991in}{1.941290in}}%
\pgfpathlineto{\pgfqpoint{2.975969in}{1.927124in}}%
\pgfpathlineto{\pgfqpoint{2.976635in}{1.931128in}}%
\pgfpathlineto{\pgfqpoint{2.976685in}{1.934853in}}%
\pgfpathlineto{\pgfqpoint{2.976784in}{1.921069in}}%
\pgfpathlineto{\pgfqpoint{2.976833in}{1.926472in}}%
\pgfpathlineto{\pgfqpoint{2.977312in}{1.532031in}}%
\pgfpathlineto{\pgfqpoint{2.978123in}{1.532052in}}%
\pgfpathlineto{\pgfqpoint{2.979256in}{1.532101in}}%
\pgfpathlineto{\pgfqpoint{2.979265in}{1.532307in}}%
\pgfpathlineto{\pgfqpoint{2.979324in}{2.024342in}}%
\pgfpathlineto{\pgfqpoint{2.980322in}{1.930873in}}%
\pgfpathlineto{\pgfqpoint{2.980330in}{1.912959in}}%
\pgfpathlineto{\pgfqpoint{2.980343in}{1.955839in}}%
\pgfpathlineto{\pgfqpoint{2.981089in}{1.937576in}}%
\pgfpathlineto{\pgfqpoint{2.981094in}{1.940985in}}%
\pgfpathlineto{\pgfqpoint{2.981796in}{1.915885in}}%
\pgfpathlineto{\pgfqpoint{2.982228in}{1.532025in}}%
\pgfpathlineto{\pgfqpoint{2.982984in}{1.532044in}}%
\pgfpathlineto{\pgfqpoint{2.984208in}{1.532160in}}%
\pgfpathlineto{\pgfqpoint{2.984308in}{1.991342in}}%
\pgfpathlineto{\pgfqpoint{2.985345in}{1.929038in}}%
\pgfpathlineto{\pgfqpoint{2.985350in}{1.921507in}}%
\pgfpathlineto{\pgfqpoint{2.985367in}{1.947503in}}%
\pgfpathlineto{\pgfqpoint{2.986092in}{1.928771in}}%
\pgfpathlineto{\pgfqpoint{2.986322in}{1.939994in}}%
\pgfpathlineto{\pgfqpoint{2.986718in}{1.922108in}}%
\pgfpathlineto{\pgfqpoint{2.987242in}{1.532031in}}%
\pgfpathlineto{\pgfqpoint{2.988131in}{1.532051in}}%
\pgfpathlineto{\pgfqpoint{2.989151in}{1.532124in}}%
\pgfpathlineto{\pgfqpoint{2.989214in}{2.029553in}}%
\pgfpathlineto{\pgfqpoint{2.990480in}{1.921176in}}%
\pgfpathlineto{\pgfqpoint{2.990484in}{1.921105in}}%
\pgfpathlineto{\pgfqpoint{2.990489in}{1.933078in}}%
\pgfpathlineto{\pgfqpoint{2.990497in}{1.948577in}}%
\pgfpathlineto{\pgfqpoint{2.990512in}{1.920729in}}%
\pgfpathlineto{\pgfqpoint{2.991260in}{1.936771in}}%
\pgfpathlineto{\pgfqpoint{2.991264in}{1.937661in}}%
\pgfpathlineto{\pgfqpoint{2.991723in}{1.877593in}}%
\pgfpathlineto{\pgfqpoint{2.992137in}{1.532027in}}%
\pgfpathlineto{\pgfqpoint{2.992724in}{1.532051in}}%
\pgfpathlineto{\pgfqpoint{2.994087in}{1.532084in}}%
\pgfpathlineto{\pgfqpoint{2.994099in}{1.532162in}}%
\pgfpathlineto{\pgfqpoint{2.994154in}{2.005171in}}%
\pgfpathlineto{\pgfqpoint{2.995230in}{1.947975in}}%
\pgfpathlineto{\pgfqpoint{2.995273in}{1.919276in}}%
\pgfpathlineto{\pgfqpoint{2.995257in}{1.949721in}}%
\pgfpathlineto{\pgfqpoint{2.995995in}{1.932815in}}%
\pgfpathlineto{\pgfqpoint{2.996028in}{1.941126in}}%
\pgfpathlineto{\pgfqpoint{2.997059in}{1.542363in}}%
\pgfpathlineto{\pgfqpoint{2.997104in}{1.532030in}}%
\pgfpathlineto{\pgfqpoint{2.997826in}{1.532050in}}%
\pgfpathlineto{\pgfqpoint{2.999038in}{1.532100in}}%
\pgfpathlineto{\pgfqpoint{2.999047in}{1.532307in}}%
\pgfpathlineto{\pgfqpoint{2.999106in}{2.024291in}}%
\pgfpathlineto{\pgfqpoint{3.000104in}{1.931746in}}%
\pgfpathlineto{\pgfqpoint{3.000112in}{1.912970in}}%
\pgfpathlineto{\pgfqpoint{3.000125in}{1.955492in}}%
\pgfpathlineto{\pgfqpoint{3.000871in}{1.937518in}}%
\pgfpathlineto{\pgfqpoint{3.000876in}{1.940915in}}%
\pgfpathlineto{\pgfqpoint{3.001560in}{1.922489in}}%
\pgfpathlineto{\pgfqpoint{3.002066in}{1.532029in}}%
\pgfpathlineto{\pgfqpoint{3.002801in}{1.532053in}}%
\pgfpathlineto{\pgfqpoint{3.003989in}{1.532151in}}%
\pgfpathlineto{\pgfqpoint{3.004057in}{2.002804in}}%
\pgfpathlineto{\pgfqpoint{3.005145in}{1.927374in}}%
\pgfpathlineto{\pgfqpoint{3.005151in}{1.919267in}}%
\pgfpathlineto{\pgfqpoint{3.005166in}{1.949285in}}%
\pgfpathlineto{\pgfqpoint{3.005909in}{1.939916in}}%
\pgfpathlineto{\pgfqpoint{3.006973in}{1.532027in}}%
\pgfpathlineto{\pgfqpoint{3.007261in}{1.532044in}}%
\pgfpathlineto{\pgfqpoint{3.008923in}{1.532083in}}%
\pgfpathlineto{\pgfqpoint{3.008934in}{1.532148in}}%
\pgfpathlineto{\pgfqpoint{3.009035in}{1.992099in}}%
\pgfpathlineto{\pgfqpoint{3.010077in}{1.924272in}}%
\pgfpathlineto{\pgfqpoint{3.010112in}{1.921509in}}%
\pgfpathlineto{\pgfqpoint{3.010093in}{1.945931in}}%
\pgfpathlineto{\pgfqpoint{3.010813in}{1.935858in}}%
\pgfpathlineto{\pgfqpoint{3.010995in}{1.940809in}}%
\pgfpathlineto{\pgfqpoint{3.011687in}{1.635568in}}%
\pgfpathlineto{\pgfqpoint{3.011901in}{1.532011in}}%
\pgfpathlineto{\pgfqpoint{3.011884in}{1.659813in}}%
\pgfpathlineto{\pgfqpoint{3.012635in}{1.532040in}}%
\pgfpathlineto{\pgfqpoint{3.013885in}{1.532670in}}%
\pgfpathlineto{\pgfqpoint{3.013986in}{2.031200in}}%
\pgfpathlineto{\pgfqpoint{3.014908in}{1.944011in}}%
\pgfpathlineto{\pgfqpoint{3.014916in}{1.910963in}}%
\pgfpathlineto{\pgfqpoint{3.014933in}{1.957778in}}%
\pgfpathlineto{\pgfqpoint{3.015677in}{1.937840in}}%
\pgfpathlineto{\pgfqpoint{3.015682in}{1.941645in}}%
\pgfpathlineto{\pgfqpoint{3.016414in}{1.916531in}}%
\pgfpathlineto{\pgfqpoint{3.016908in}{1.532028in}}%
\pgfpathlineto{\pgfqpoint{3.017544in}{1.532052in}}%
\pgfpathlineto{\pgfqpoint{3.018815in}{1.532086in}}%
\pgfpathlineto{\pgfqpoint{3.018825in}{1.532147in}}%
\pgfpathlineto{\pgfqpoint{3.018891in}{2.034094in}}%
\pgfpathlineto{\pgfqpoint{3.020002in}{1.932322in}}%
\pgfpathlineto{\pgfqpoint{3.020010in}{1.915618in}}%
\pgfpathlineto{\pgfqpoint{3.020022in}{1.952651in}}%
\pgfpathlineto{\pgfqpoint{3.020767in}{1.936884in}}%
\pgfpathlineto{\pgfqpoint{3.020773in}{1.939844in}}%
\pgfpathlineto{\pgfqpoint{3.021385in}{1.893138in}}%
\pgfpathlineto{\pgfqpoint{3.021834in}{1.532033in}}%
\pgfpathlineto{\pgfqpoint{3.022411in}{1.532045in}}%
\pgfpathlineto{\pgfqpoint{3.023772in}{1.532193in}}%
\pgfpathlineto{\pgfqpoint{3.023836in}{2.036901in}}%
\pgfpathlineto{\pgfqpoint{3.024888in}{1.935184in}}%
\pgfpathlineto{\pgfqpoint{3.024895in}{1.912970in}}%
\pgfpathlineto{\pgfqpoint{3.024911in}{1.954466in}}%
\pgfpathlineto{\pgfqpoint{3.025657in}{1.935101in}}%
\pgfpathlineto{\pgfqpoint{3.026628in}{1.595548in}}%
\pgfpathlineto{\pgfqpoint{3.025679in}{1.939855in}}%
\pgfpathlineto{\pgfqpoint{3.026671in}{1.607397in}}%
\pgfpathlineto{\pgfqpoint{3.026719in}{1.649614in}}%
\pgfpathlineto{\pgfqpoint{3.026778in}{1.532029in}}%
\pgfpathlineto{\pgfqpoint{3.027306in}{1.532047in}}%
\pgfpathlineto{\pgfqpoint{3.028714in}{1.532121in}}%
\pgfpathlineto{\pgfqpoint{3.028756in}{1.763457in}}%
\pgfpathlineto{\pgfqpoint{3.028814in}{1.994404in}}%
\pgfpathlineto{\pgfqpoint{3.029534in}{1.949838in}}%
\pgfpathlineto{\pgfqpoint{3.029544in}{1.913345in}}%
\pgfpathlineto{\pgfqpoint{3.029559in}{1.954811in}}%
\pgfpathlineto{\pgfqpoint{3.030307in}{1.941990in}}%
\pgfpathlineto{\pgfqpoint{3.030363in}{1.927026in}}%
\pgfpathlineto{\pgfqpoint{3.031090in}{1.932213in}}%
\pgfpathlineto{\pgfqpoint{3.031189in}{1.935197in}}%
\pgfpathlineto{\pgfqpoint{3.031715in}{1.532032in}}%
\pgfpathlineto{\pgfqpoint{3.032528in}{1.532051in}}%
\pgfpathlineto{\pgfqpoint{3.033660in}{1.532128in}}%
\pgfpathlineto{\pgfqpoint{3.033724in}{2.029514in}}%
\pgfpathlineto{\pgfqpoint{3.034955in}{1.930701in}}%
\pgfpathlineto{\pgfqpoint{3.034963in}{1.919715in}}%
\pgfpathlineto{\pgfqpoint{3.034975in}{1.949201in}}%
\pgfpathlineto{\pgfqpoint{3.035725in}{1.930439in}}%
\pgfpathlineto{\pgfqpoint{3.035741in}{1.937969in}}%
\pgfpathlineto{\pgfqpoint{3.036179in}{1.922489in}}%
\pgfpathlineto{\pgfqpoint{3.036628in}{1.532022in}}%
\pgfpathlineto{\pgfqpoint{3.037417in}{1.532044in}}%
\pgfpathlineto{\pgfqpoint{3.038605in}{1.532120in}}%
\pgfpathlineto{\pgfqpoint{3.038645in}{1.722790in}}%
\pgfpathlineto{\pgfqpoint{3.038676in}{1.996342in}}%
\pgfpathlineto{\pgfqpoint{3.039428in}{1.950921in}}%
\pgfpathlineto{\pgfqpoint{3.039438in}{1.913142in}}%
\pgfpathlineto{\pgfqpoint{3.039453in}{1.954991in}}%
\pgfpathlineto{\pgfqpoint{3.040193in}{1.930624in}}%
\pgfpathlineto{\pgfqpoint{3.040202in}{1.941489in}}%
\pgfpathlineto{\pgfqpoint{3.040334in}{1.927614in}}%
\pgfpathlineto{\pgfqpoint{3.040946in}{1.929361in}}%
\pgfpathlineto{\pgfqpoint{3.040995in}{1.937153in}}%
\pgfpathlineto{\pgfqpoint{3.041193in}{1.875320in}}%
\pgfpathlineto{\pgfqpoint{3.041625in}{1.532030in}}%
\pgfpathlineto{\pgfqpoint{3.042211in}{1.532043in}}%
\pgfpathlineto{\pgfqpoint{3.043549in}{1.532111in}}%
\pgfpathlineto{\pgfqpoint{3.043591in}{1.734894in}}%
\pgfpathlineto{\pgfqpoint{3.043608in}{1.992390in}}%
\pgfpathlineto{\pgfqpoint{3.044372in}{1.945172in}}%
\pgfpathlineto{\pgfqpoint{3.044377in}{1.953461in}}%
\pgfpathlineto{\pgfqpoint{3.044393in}{1.915971in}}%
\pgfpathlineto{\pgfqpoint{3.045137in}{1.936186in}}%
\pgfpathlineto{\pgfqpoint{3.045147in}{1.941852in}}%
\pgfpathlineto{\pgfqpoint{3.045166in}{1.927246in}}%
\pgfpathlineto{\pgfqpoint{3.045878in}{1.936757in}}%
\pgfpathlineto{\pgfqpoint{3.045977in}{1.939598in}}%
\pgfpathlineto{\pgfqpoint{3.046175in}{1.801669in}}%
\pgfpathlineto{\pgfqpoint{3.046539in}{1.532025in}}%
\pgfpathlineto{\pgfqpoint{3.047080in}{1.532037in}}%
\pgfpathlineto{\pgfqpoint{3.048496in}{1.532124in}}%
\pgfpathlineto{\pgfqpoint{3.048560in}{2.029429in}}%
\pgfpathlineto{\pgfqpoint{3.049820in}{1.935054in}}%
\pgfpathlineto{\pgfqpoint{3.049828in}{1.920001in}}%
\pgfpathlineto{\pgfqpoint{3.049843in}{1.948188in}}%
\pgfpathlineto{\pgfqpoint{3.050590in}{1.932828in}}%
\pgfpathlineto{\pgfqpoint{3.051500in}{1.532031in}}%
\pgfpathlineto{\pgfqpoint{3.050610in}{1.937708in}}%
\pgfpathlineto{\pgfqpoint{3.052069in}{1.532043in}}%
\pgfpathlineto{\pgfqpoint{3.053443in}{1.532135in}}%
\pgfpathlineto{\pgfqpoint{3.053506in}{2.033543in}}%
\pgfpathlineto{\pgfqpoint{3.054703in}{1.944838in}}%
\pgfpathlineto{\pgfqpoint{3.054714in}{1.918302in}}%
\pgfpathlineto{\pgfqpoint{3.054726in}{1.949712in}}%
\pgfpathlineto{\pgfqpoint{3.055475in}{1.933258in}}%
\pgfpathlineto{\pgfqpoint{3.056467in}{1.532029in}}%
\pgfpathlineto{\pgfqpoint{3.055498in}{1.938113in}}%
\pgfpathlineto{\pgfqpoint{3.056757in}{1.532035in}}%
\pgfpathlineto{\pgfqpoint{3.058389in}{1.532156in}}%
\pgfpathlineto{\pgfqpoint{3.058493in}{1.994879in}}%
\pgfpathlineto{\pgfqpoint{3.059525in}{1.925170in}}%
\pgfpathlineto{\pgfqpoint{3.059531in}{1.920991in}}%
\pgfpathlineto{\pgfqpoint{3.059547in}{1.947420in}}%
\pgfpathlineto{\pgfqpoint{3.060274in}{1.930513in}}%
\pgfpathlineto{\pgfqpoint{3.060285in}{1.939481in}}%
\pgfpathlineto{\pgfqpoint{3.060922in}{1.907520in}}%
\pgfpathlineto{\pgfqpoint{3.061416in}{1.532032in}}%
\pgfpathlineto{\pgfqpoint{3.062104in}{1.532047in}}%
\pgfpathlineto{\pgfqpoint{3.063336in}{1.532179in}}%
\pgfpathlineto{\pgfqpoint{3.063398in}{2.031217in}}%
\pgfpathlineto{\pgfqpoint{3.064483in}{1.925798in}}%
\pgfpathlineto{\pgfqpoint{3.064488in}{1.915333in}}%
\pgfpathlineto{\pgfqpoint{3.064503in}{1.953671in}}%
\pgfpathlineto{\pgfqpoint{3.065253in}{1.930272in}}%
\pgfpathlineto{\pgfqpoint{3.066319in}{1.532027in}}%
\pgfpathlineto{\pgfqpoint{3.065273in}{1.939299in}}%
\pgfpathlineto{\pgfqpoint{3.067152in}{1.532045in}}%
\pgfpathlineto{\pgfqpoint{3.068280in}{1.532147in}}%
\pgfpathlineto{\pgfqpoint{3.068383in}{2.001720in}}%
\pgfpathlineto{\pgfqpoint{3.069439in}{1.927324in}}%
\pgfpathlineto{\pgfqpoint{3.069445in}{1.919875in}}%
\pgfpathlineto{\pgfqpoint{3.069460in}{1.948633in}}%
\pgfpathlineto{\pgfqpoint{3.070202in}{1.929037in}}%
\pgfpathlineto{\pgfqpoint{3.070220in}{1.939473in}}%
\pgfpathlineto{\pgfqpoint{3.070867in}{1.857586in}}%
\pgfpathlineto{\pgfqpoint{3.071307in}{1.532032in}}%
\pgfpathlineto{\pgfqpoint{3.071846in}{1.532045in}}%
\pgfpathlineto{\pgfqpoint{3.073225in}{1.532151in}}%
\pgfpathlineto{\pgfqpoint{3.073288in}{2.029379in}}%
\pgfpathlineto{\pgfqpoint{3.074431in}{1.929425in}}%
\pgfpathlineto{\pgfqpoint{3.074467in}{1.917793in}}%
\pgfpathlineto{\pgfqpoint{3.074452in}{1.951568in}}%
\pgfpathlineto{\pgfqpoint{3.075201in}{1.930200in}}%
\pgfpathlineto{\pgfqpoint{3.075217in}{1.938688in}}%
\pgfpathlineto{\pgfqpoint{3.075827in}{1.829907in}}%
\pgfpathlineto{\pgfqpoint{3.076192in}{1.532011in}}%
\pgfpathlineto{\pgfqpoint{3.076782in}{1.532053in}}%
\pgfpathlineto{\pgfqpoint{3.078166in}{1.532103in}}%
\pgfpathlineto{\pgfqpoint{3.078183in}{1.544949in}}%
\pgfpathlineto{\pgfqpoint{3.078277in}{2.029160in}}%
\pgfpathlineto{\pgfqpoint{3.079085in}{1.918945in}}%
\pgfpathlineto{\pgfqpoint{3.079089in}{1.906356in}}%
\pgfpathlineto{\pgfqpoint{3.079105in}{1.961891in}}%
\pgfpathlineto{\pgfqpoint{3.079847in}{1.942768in}}%
\pgfpathlineto{\pgfqpoint{3.079862in}{1.926239in}}%
\pgfpathlineto{\pgfqpoint{3.080646in}{1.931365in}}%
\pgfpathlineto{\pgfqpoint{3.080660in}{1.933405in}}%
\pgfpathlineto{\pgfqpoint{3.080738in}{1.894560in}}%
\pgfpathlineto{\pgfqpoint{3.081179in}{1.532029in}}%
\pgfpathlineto{\pgfqpoint{3.081758in}{1.532041in}}%
\pgfpathlineto{\pgfqpoint{3.083116in}{1.532143in}}%
\pgfpathlineto{\pgfqpoint{3.083182in}{2.031440in}}%
\pgfpathlineto{\pgfqpoint{3.084294in}{1.927625in}}%
\pgfpathlineto{\pgfqpoint{3.084329in}{1.916588in}}%
\pgfpathlineto{\pgfqpoint{3.084316in}{1.951583in}}%
\pgfpathlineto{\pgfqpoint{3.085062in}{1.933756in}}%
\pgfpathlineto{\pgfqpoint{3.085071in}{1.939369in}}%
\pgfpathlineto{\pgfqpoint{3.085702in}{1.859921in}}%
\pgfpathlineto{\pgfqpoint{3.086101in}{1.532027in}}%
\pgfpathlineto{\pgfqpoint{3.086687in}{1.532050in}}%
\pgfpathlineto{\pgfqpoint{3.088051in}{1.532085in}}%
\pgfpathlineto{\pgfqpoint{3.088062in}{1.532147in}}%
\pgfpathlineto{\pgfqpoint{3.088163in}{2.000300in}}%
\pgfpathlineto{\pgfqpoint{3.089221in}{1.927290in}}%
\pgfpathlineto{\pgfqpoint{3.089226in}{1.919673in}}%
\pgfpathlineto{\pgfqpoint{3.089242in}{1.948817in}}%
\pgfpathlineto{\pgfqpoint{3.089982in}{1.928976in}}%
\pgfpathlineto{\pgfqpoint{3.090000in}{1.939444in}}%
\pgfpathlineto{\pgfqpoint{3.090608in}{1.897446in}}%
\pgfpathlineto{\pgfqpoint{3.091028in}{1.532020in}}%
\pgfpathlineto{\pgfqpoint{3.091684in}{1.532041in}}%
\pgfpathlineto{\pgfqpoint{3.093007in}{1.532143in}}%
\pgfpathlineto{\pgfqpoint{3.093074in}{2.026507in}}%
\pgfpathlineto{\pgfqpoint{3.094185in}{1.927125in}}%
\pgfpathlineto{\pgfqpoint{3.094189in}{1.917754in}}%
\pgfpathlineto{\pgfqpoint{3.094206in}{1.951885in}}%
\pgfpathlineto{\pgfqpoint{3.094952in}{1.930196in}}%
\pgfpathlineto{\pgfqpoint{3.094964in}{1.939537in}}%
\pgfpathlineto{\pgfqpoint{3.095570in}{1.883310in}}%
\pgfpathlineto{\pgfqpoint{3.095991in}{1.532027in}}%
\pgfpathlineto{\pgfqpoint{3.096575in}{1.532050in}}%
\pgfpathlineto{\pgfqpoint{3.097942in}{1.532085in}}%
\pgfpathlineto{\pgfqpoint{3.097952in}{1.532147in}}%
\pgfpathlineto{\pgfqpoint{3.098020in}{2.008202in}}%
\pgfpathlineto{\pgfqpoint{3.099105in}{1.929845in}}%
\pgfpathlineto{\pgfqpoint{3.099110in}{1.918941in}}%
\pgfpathlineto{\pgfqpoint{3.099125in}{1.949665in}}%
\pgfpathlineto{\pgfqpoint{3.099871in}{1.928947in}}%
\pgfpathlineto{\pgfqpoint{3.099893in}{1.939703in}}%
\pgfpathlineto{\pgfqpoint{3.100528in}{1.883408in}}%
\pgfpathlineto{\pgfqpoint{3.100976in}{1.532029in}}%
\pgfpathlineto{\pgfqpoint{3.101513in}{1.532048in}}%
\pgfpathlineto{\pgfqpoint{3.102897in}{1.532138in}}%
\pgfpathlineto{\pgfqpoint{3.102963in}{2.041497in}}%
\pgfpathlineto{\pgfqpoint{3.104103in}{1.937632in}}%
\pgfpathlineto{\pgfqpoint{3.104110in}{1.915283in}}%
\pgfpathlineto{\pgfqpoint{3.104126in}{1.953228in}}%
\pgfpathlineto{\pgfqpoint{3.104874in}{1.938129in}}%
\pgfpathlineto{\pgfqpoint{3.105900in}{1.532033in}}%
\pgfpathlineto{\pgfqpoint{3.104901in}{1.939391in}}%
\pgfpathlineto{\pgfqpoint{3.106024in}{1.532035in}}%
\pgfpathlineto{\pgfqpoint{3.107844in}{1.532158in}}%
\pgfpathlineto{\pgfqpoint{3.107907in}{2.030276in}}%
\pgfpathlineto{\pgfqpoint{3.109023in}{1.920406in}}%
\pgfpathlineto{\pgfqpoint{3.109027in}{1.915874in}}%
\pgfpathlineto{\pgfqpoint{3.109042in}{1.951927in}}%
\pgfpathlineto{\pgfqpoint{3.109775in}{1.938008in}}%
\pgfpathlineto{\pgfqpoint{3.109780in}{1.939403in}}%
\pgfpathlineto{\pgfqpoint{3.110380in}{1.919104in}}%
\pgfpathlineto{\pgfqpoint{3.110810in}{1.532012in}}%
\pgfpathlineto{\pgfqpoint{3.111560in}{1.532039in}}%
\pgfpathlineto{\pgfqpoint{3.112792in}{1.532261in}}%
\pgfpathlineto{\pgfqpoint{3.112851in}{2.024282in}}%
\pgfpathlineto{\pgfqpoint{3.113878in}{1.931627in}}%
\pgfpathlineto{\pgfqpoint{3.113886in}{1.914051in}}%
\pgfpathlineto{\pgfqpoint{3.113930in}{1.953372in}}%
\pgfpathlineto{\pgfqpoint{3.114649in}{1.930142in}}%
\pgfpathlineto{\pgfqpoint{3.115806in}{1.532033in}}%
\pgfpathlineto{\pgfqpoint{3.114669in}{1.939584in}}%
\pgfpathlineto{\pgfqpoint{3.116292in}{1.532043in}}%
\pgfpathlineto{\pgfqpoint{3.117738in}{1.532307in}}%
\pgfpathlineto{\pgfqpoint{3.117797in}{2.024540in}}%
\pgfpathlineto{\pgfqpoint{3.118795in}{1.930970in}}%
\pgfpathlineto{\pgfqpoint{3.118803in}{1.912757in}}%
\pgfpathlineto{\pgfqpoint{3.118816in}{1.955891in}}%
\pgfpathlineto{\pgfqpoint{3.119562in}{1.937594in}}%
\pgfpathlineto{\pgfqpoint{3.119567in}{1.941009in}}%
\pgfpathlineto{\pgfqpoint{3.120256in}{1.923220in}}%
\pgfpathlineto{\pgfqpoint{3.120701in}{1.532012in}}%
\pgfpathlineto{\pgfqpoint{3.121580in}{1.532059in}}%
\pgfpathlineto{\pgfqpoint{3.122676in}{1.532108in}}%
\pgfpathlineto{\pgfqpoint{3.122704in}{1.576889in}}%
\pgfpathlineto{\pgfqpoint{3.122787in}{2.031168in}}%
\pgfpathlineto{\pgfqpoint{3.123504in}{1.930592in}}%
\pgfpathlineto{\pgfqpoint{3.123512in}{1.901371in}}%
\pgfpathlineto{\pgfqpoint{3.123527in}{1.966014in}}%
\pgfpathlineto{\pgfqpoint{3.124269in}{1.938482in}}%
\pgfpathlineto{\pgfqpoint{3.124274in}{1.944819in}}%
\pgfpathlineto{\pgfqpoint{3.124290in}{1.924127in}}%
\pgfpathlineto{\pgfqpoint{3.125035in}{1.936050in}}%
\pgfpathlineto{\pgfqpoint{3.125690in}{1.532032in}}%
\pgfpathlineto{\pgfqpoint{3.125069in}{1.936650in}}%
\pgfpathlineto{\pgfqpoint{3.126665in}{1.532053in}}%
\pgfpathlineto{\pgfqpoint{3.127626in}{1.532153in}}%
\pgfpathlineto{\pgfqpoint{3.127687in}{2.030025in}}%
\pgfpathlineto{\pgfqpoint{3.128831in}{1.935382in}}%
\pgfpathlineto{\pgfqpoint{3.128839in}{1.916637in}}%
\pgfpathlineto{\pgfqpoint{3.128854in}{1.951599in}}%
\pgfpathlineto{\pgfqpoint{3.129603in}{1.930502in}}%
\pgfpathlineto{\pgfqpoint{3.130592in}{1.532011in}}%
\pgfpathlineto{\pgfqpoint{3.129622in}{1.938548in}}%
\pgfpathlineto{\pgfqpoint{3.131879in}{1.532069in}}%
\pgfpathlineto{\pgfqpoint{3.132565in}{1.532101in}}%
\pgfpathlineto{\pgfqpoint{3.132574in}{1.532244in}}%
\pgfpathlineto{\pgfqpoint{3.132633in}{2.025772in}}%
\pgfpathlineto{\pgfqpoint{3.133662in}{1.920660in}}%
\pgfpathlineto{\pgfqpoint{3.133666in}{1.913598in}}%
\pgfpathlineto{\pgfqpoint{3.133681in}{1.954565in}}%
\pgfpathlineto{\pgfqpoint{3.134420in}{1.940029in}}%
\pgfpathlineto{\pgfqpoint{3.135562in}{1.532031in}}%
\pgfpathlineto{\pgfqpoint{3.135666in}{1.532036in}}%
\pgfpathlineto{\pgfqpoint{3.137511in}{1.532101in}}%
\pgfpathlineto{\pgfqpoint{3.137520in}{1.532315in}}%
\pgfpathlineto{\pgfqpoint{3.137652in}{2.022186in}}%
\pgfpathlineto{\pgfqpoint{3.138574in}{1.942742in}}%
\pgfpathlineto{\pgfqpoint{3.138582in}{1.912974in}}%
\pgfpathlineto{\pgfqpoint{3.138598in}{1.955544in}}%
\pgfpathlineto{\pgfqpoint{3.139343in}{1.937678in}}%
\pgfpathlineto{\pgfqpoint{3.139348in}{1.940825in}}%
\pgfpathlineto{\pgfqpoint{3.140042in}{1.923239in}}%
\pgfpathlineto{\pgfqpoint{3.140483in}{1.532016in}}%
\pgfpathlineto{\pgfqpoint{3.141334in}{1.532042in}}%
\pgfpathlineto{\pgfqpoint{3.142466in}{1.532697in}}%
\pgfpathlineto{\pgfqpoint{3.142568in}{2.031246in}}%
\pgfpathlineto{\pgfqpoint{3.143490in}{1.944874in}}%
\pgfpathlineto{\pgfqpoint{3.143498in}{1.911154in}}%
\pgfpathlineto{\pgfqpoint{3.143515in}{1.957898in}}%
\pgfpathlineto{\pgfqpoint{3.144257in}{1.937750in}}%
\pgfpathlineto{\pgfqpoint{3.144262in}{1.941588in}}%
\pgfpathlineto{\pgfqpoint{3.144277in}{1.927219in}}%
\pgfpathlineto{\pgfqpoint{3.144968in}{1.931040in}}%
\pgfpathlineto{\pgfqpoint{3.145453in}{1.532031in}}%
\pgfpathlineto{\pgfqpoint{3.146544in}{1.532058in}}%
\pgfpathlineto{\pgfqpoint{3.147402in}{1.532102in}}%
\pgfpathlineto{\pgfqpoint{3.147411in}{1.532320in}}%
\pgfpathlineto{\pgfqpoint{3.147543in}{2.022376in}}%
\pgfpathlineto{\pgfqpoint{3.148467in}{1.933919in}}%
\pgfpathlineto{\pgfqpoint{3.148475in}{1.912892in}}%
\pgfpathlineto{\pgfqpoint{3.148520in}{1.954265in}}%
\pgfpathlineto{\pgfqpoint{3.149236in}{1.940771in}}%
\pgfpathlineto{\pgfqpoint{3.150374in}{1.532017in}}%
\pgfpathlineto{\pgfqpoint{3.150582in}{1.532044in}}%
\pgfpathlineto{\pgfqpoint{3.152348in}{1.532106in}}%
\pgfpathlineto{\pgfqpoint{3.152375in}{1.573819in}}%
\pgfpathlineto{\pgfqpoint{3.152459in}{2.031508in}}%
\pgfpathlineto{\pgfqpoint{3.153179in}{1.914996in}}%
\pgfpathlineto{\pgfqpoint{3.153183in}{1.901452in}}%
\pgfpathlineto{\pgfqpoint{3.153199in}{1.966388in}}%
\pgfpathlineto{\pgfqpoint{3.153941in}{1.943589in}}%
\pgfpathlineto{\pgfqpoint{3.153950in}{1.935721in}}%
\pgfpathlineto{\pgfqpoint{3.153960in}{1.924656in}}%
\pgfpathlineto{\pgfqpoint{3.153973in}{1.944134in}}%
\pgfpathlineto{\pgfqpoint{3.154714in}{1.936412in}}%
\pgfpathlineto{\pgfqpoint{3.155320in}{1.532015in}}%
\pgfpathlineto{\pgfqpoint{3.156470in}{1.532049in}}%
\pgfpathlineto{\pgfqpoint{3.157302in}{1.532438in}}%
\pgfpathlineto{\pgfqpoint{3.157404in}{2.028331in}}%
\pgfpathlineto{\pgfqpoint{3.158329in}{1.926020in}}%
\pgfpathlineto{\pgfqpoint{3.158365in}{1.912146in}}%
\pgfpathlineto{\pgfqpoint{3.158350in}{1.957370in}}%
\pgfpathlineto{\pgfqpoint{3.159094in}{1.941181in}}%
\pgfpathlineto{\pgfqpoint{3.160322in}{1.532029in}}%
\pgfpathlineto{\pgfqpoint{3.160760in}{1.532046in}}%
\pgfpathlineto{\pgfqpoint{3.162244in}{1.532155in}}%
\pgfpathlineto{\pgfqpoint{3.162347in}{1.997517in}}%
\pgfpathlineto{\pgfqpoint{3.163404in}{1.939456in}}%
\pgfpathlineto{\pgfqpoint{3.163414in}{1.920857in}}%
\pgfpathlineto{\pgfqpoint{3.163430in}{1.947299in}}%
\pgfpathlineto{\pgfqpoint{3.164170in}{1.939081in}}%
\pgfpathlineto{\pgfqpoint{3.165210in}{1.532020in}}%
\pgfpathlineto{\pgfqpoint{3.164240in}{1.939450in}}%
\pgfpathlineto{\pgfqpoint{3.165379in}{1.532030in}}%
\pgfpathlineto{\pgfqpoint{3.167189in}{1.532146in}}%
\pgfpathlineto{\pgfqpoint{3.167255in}{2.031075in}}%
\pgfpathlineto{\pgfqpoint{3.168362in}{1.943730in}}%
\pgfpathlineto{\pgfqpoint{3.168373in}{1.916392in}}%
\pgfpathlineto{\pgfqpoint{3.168387in}{1.952584in}}%
\pgfpathlineto{\pgfqpoint{3.169134in}{1.937295in}}%
\pgfpathlineto{\pgfqpoint{3.169140in}{1.939467in}}%
\pgfpathlineto{\pgfqpoint{3.169731in}{1.901883in}}%
\pgfpathlineto{\pgfqpoint{3.170018in}{1.596566in}}%
\pgfpathlineto{\pgfqpoint{3.170132in}{1.648553in}}%
\pgfusepath{stroke}%
\end{pgfscope}%
\begin{pgfscope}%
\pgfsetrectcap%
\pgfsetmiterjoin%
\pgfsetlinewidth{0.803000pt}%
\definecolor{currentstroke}{rgb}{0.000000,0.000000,0.000000}%
\pgfsetstrokecolor{currentstroke}%
\pgfsetdash{}{0pt}%
\pgfpathmoveto{\pgfqpoint{0.573768in}{1.532029in}}%
\pgfpathlineto{\pgfqpoint{0.573768in}{2.438279in}}%
\pgfusepath{stroke}%
\end{pgfscope}%
\begin{pgfscope}%
\pgfsetrectcap%
\pgfsetmiterjoin%
\pgfsetlinewidth{0.803000pt}%
\definecolor{currentstroke}{rgb}{0.000000,0.000000,0.000000}%
\pgfsetstrokecolor{currentstroke}%
\pgfsetdash{}{0pt}%
\pgfpathmoveto{\pgfqpoint{3.293768in}{1.532029in}}%
\pgfpathlineto{\pgfqpoint{3.293768in}{2.438279in}}%
\pgfusepath{stroke}%
\end{pgfscope}%
\begin{pgfscope}%
\pgfsetrectcap%
\pgfsetmiterjoin%
\pgfsetlinewidth{0.803000pt}%
\definecolor{currentstroke}{rgb}{0.000000,0.000000,0.000000}%
\pgfsetstrokecolor{currentstroke}%
\pgfsetdash{}{0pt}%
\pgfpathmoveto{\pgfqpoint{0.573768in}{1.532029in}}%
\pgfpathlineto{\pgfqpoint{3.293768in}{1.532029in}}%
\pgfusepath{stroke}%
\end{pgfscope}%
\begin{pgfscope}%
\pgfsetrectcap%
\pgfsetmiterjoin%
\pgfsetlinewidth{0.803000pt}%
\definecolor{currentstroke}{rgb}{0.000000,0.000000,0.000000}%
\pgfsetstrokecolor{currentstroke}%
\pgfsetdash{}{0pt}%
\pgfpathmoveto{\pgfqpoint{0.573768in}{2.438279in}}%
\pgfpathlineto{\pgfqpoint{3.293768in}{2.438279in}}%
\pgfusepath{stroke}%
\end{pgfscope}%
\begin{pgfscope}%
\definecolor{textcolor}{rgb}{0.000000,0.000000,0.000000}%
\pgfsetstrokecolor{textcolor}%
\pgfsetfillcolor{textcolor}%
\pgftext[x=0.628168in,y=2.374842in,left,top]{\color{textcolor}{\rmfamily\fontsize{8.000000}{9.600000}\selectfont\catcode`\^=\active\def^{\ifmmode\sp\else\^{}\fi}\catcode`\%=\active\def%{\%}(b)}}%
\end{pgfscope}%
\begin{pgfscope}%
\pgfsetbuttcap%
\pgfsetmiterjoin%
\definecolor{currentfill}{rgb}{1.000000,1.000000,1.000000}%
\pgfsetfillcolor{currentfill}%
\pgfsetlinewidth{0.000000pt}%
\definecolor{currentstroke}{rgb}{0.000000,0.000000,0.000000}%
\pgfsetstrokecolor{currentstroke}%
\pgfsetstrokeopacity{0.000000}%
\pgfsetdash{}{0pt}%
\pgfpathmoveto{\pgfqpoint{0.573768in}{0.462654in}}%
\pgfpathlineto{\pgfqpoint{3.293768in}{0.462654in}}%
\pgfpathlineto{\pgfqpoint{3.293768in}{1.368904in}}%
\pgfpathlineto{\pgfqpoint{0.573768in}{1.368904in}}%
\pgfpathlineto{\pgfqpoint{0.573768in}{0.462654in}}%
\pgfpathclose%
\pgfusepath{fill}%
\end{pgfscope}%
\begin{pgfscope}%
\pgfpathrectangle{\pgfqpoint{0.573768in}{0.462654in}}{\pgfqpoint{2.720000in}{0.906250in}}%
\pgfusepath{clip}%
\pgfsetbuttcap%
\pgfsetroundjoin%
\pgfsetlinewidth{0.501875pt}%
\definecolor{currentstroke}{rgb}{0.690196,0.690196,0.690196}%
\pgfsetstrokecolor{currentstroke}%
\pgfsetstrokeopacity{0.250000}%
\pgfsetdash{{1.850000pt}{0.800000pt}}{0.000000pt}%
\pgfpathmoveto{\pgfqpoint{0.697405in}{0.462654in}}%
\pgfpathlineto{\pgfqpoint{0.697405in}{1.368904in}}%
\pgfusepath{stroke}%
\end{pgfscope}%
\begin{pgfscope}%
\pgfsetbuttcap%
\pgfsetroundjoin%
\definecolor{currentfill}{rgb}{0.000000,0.000000,0.000000}%
\pgfsetfillcolor{currentfill}%
\pgfsetlinewidth{0.803000pt}%
\definecolor{currentstroke}{rgb}{0.000000,0.000000,0.000000}%
\pgfsetstrokecolor{currentstroke}%
\pgfsetdash{}{0pt}%
\pgfsys@defobject{currentmarker}{\pgfqpoint{0.000000in}{-0.048611in}}{\pgfqpoint{0.000000in}{0.000000in}}{%
\pgfpathmoveto{\pgfqpoint{0.000000in}{0.000000in}}%
\pgfpathlineto{\pgfqpoint{0.000000in}{-0.048611in}}%
\pgfusepath{stroke,fill}%
}%
\begin{pgfscope}%
\pgfsys@transformshift{0.697405in}{0.462654in}%
\pgfsys@useobject{currentmarker}{}%
\end{pgfscope}%
\end{pgfscope}%
\begin{pgfscope}%
\definecolor{textcolor}{rgb}{0.000000,0.000000,0.000000}%
\pgfsetstrokecolor{textcolor}%
\pgfsetfillcolor{textcolor}%
\pgftext[x=0.697405in,y=0.365432in,,top]{\color{textcolor}{\rmfamily\fontsize{8.000000}{9.600000}\selectfont\catcode`\^=\active\def^{\ifmmode\sp\else\^{}\fi}\catcode`\%=\active\def%{\%}$\mathdefault{0}$}}%
\end{pgfscope}%
\begin{pgfscope}%
\pgfpathrectangle{\pgfqpoint{0.573768in}{0.462654in}}{\pgfqpoint{2.720000in}{0.906250in}}%
\pgfusepath{clip}%
\pgfsetbuttcap%
\pgfsetroundjoin%
\pgfsetlinewidth{0.501875pt}%
\definecolor{currentstroke}{rgb}{0.690196,0.690196,0.690196}%
\pgfsetstrokecolor{currentstroke}%
\pgfsetstrokeopacity{0.250000}%
\pgfsetdash{{1.850000pt}{0.800000pt}}{0.000000pt}%
\pgfpathmoveto{\pgfqpoint{1.191950in}{0.462654in}}%
\pgfpathlineto{\pgfqpoint{1.191950in}{1.368904in}}%
\pgfusepath{stroke}%
\end{pgfscope}%
\begin{pgfscope}%
\pgfsetbuttcap%
\pgfsetroundjoin%
\definecolor{currentfill}{rgb}{0.000000,0.000000,0.000000}%
\pgfsetfillcolor{currentfill}%
\pgfsetlinewidth{0.803000pt}%
\definecolor{currentstroke}{rgb}{0.000000,0.000000,0.000000}%
\pgfsetstrokecolor{currentstroke}%
\pgfsetdash{}{0pt}%
\pgfsys@defobject{currentmarker}{\pgfqpoint{0.000000in}{-0.048611in}}{\pgfqpoint{0.000000in}{0.000000in}}{%
\pgfpathmoveto{\pgfqpoint{0.000000in}{0.000000in}}%
\pgfpathlineto{\pgfqpoint{0.000000in}{-0.048611in}}%
\pgfusepath{stroke,fill}%
}%
\begin{pgfscope}%
\pgfsys@transformshift{1.191950in}{0.462654in}%
\pgfsys@useobject{currentmarker}{}%
\end{pgfscope}%
\end{pgfscope}%
\begin{pgfscope}%
\definecolor{textcolor}{rgb}{0.000000,0.000000,0.000000}%
\pgfsetstrokecolor{textcolor}%
\pgfsetfillcolor{textcolor}%
\pgftext[x=1.191950in,y=0.365432in,,top]{\color{textcolor}{\rmfamily\fontsize{8.000000}{9.600000}\selectfont\catcode`\^=\active\def^{\ifmmode\sp\else\^{}\fi}\catcode`\%=\active\def%{\%}$\mathdefault{1}$}}%
\end{pgfscope}%
\begin{pgfscope}%
\pgfpathrectangle{\pgfqpoint{0.573768in}{0.462654in}}{\pgfqpoint{2.720000in}{0.906250in}}%
\pgfusepath{clip}%
\pgfsetbuttcap%
\pgfsetroundjoin%
\pgfsetlinewidth{0.501875pt}%
\definecolor{currentstroke}{rgb}{0.690196,0.690196,0.690196}%
\pgfsetstrokecolor{currentstroke}%
\pgfsetstrokeopacity{0.250000}%
\pgfsetdash{{1.850000pt}{0.800000pt}}{0.000000pt}%
\pgfpathmoveto{\pgfqpoint{1.686496in}{0.462654in}}%
\pgfpathlineto{\pgfqpoint{1.686496in}{1.368904in}}%
\pgfusepath{stroke}%
\end{pgfscope}%
\begin{pgfscope}%
\pgfsetbuttcap%
\pgfsetroundjoin%
\definecolor{currentfill}{rgb}{0.000000,0.000000,0.000000}%
\pgfsetfillcolor{currentfill}%
\pgfsetlinewidth{0.803000pt}%
\definecolor{currentstroke}{rgb}{0.000000,0.000000,0.000000}%
\pgfsetstrokecolor{currentstroke}%
\pgfsetdash{}{0pt}%
\pgfsys@defobject{currentmarker}{\pgfqpoint{0.000000in}{-0.048611in}}{\pgfqpoint{0.000000in}{0.000000in}}{%
\pgfpathmoveto{\pgfqpoint{0.000000in}{0.000000in}}%
\pgfpathlineto{\pgfqpoint{0.000000in}{-0.048611in}}%
\pgfusepath{stroke,fill}%
}%
\begin{pgfscope}%
\pgfsys@transformshift{1.686496in}{0.462654in}%
\pgfsys@useobject{currentmarker}{}%
\end{pgfscope}%
\end{pgfscope}%
\begin{pgfscope}%
\definecolor{textcolor}{rgb}{0.000000,0.000000,0.000000}%
\pgfsetstrokecolor{textcolor}%
\pgfsetfillcolor{textcolor}%
\pgftext[x=1.686496in,y=0.365432in,,top]{\color{textcolor}{\rmfamily\fontsize{8.000000}{9.600000}\selectfont\catcode`\^=\active\def^{\ifmmode\sp\else\^{}\fi}\catcode`\%=\active\def%{\%}$\mathdefault{2}$}}%
\end{pgfscope}%
\begin{pgfscope}%
\pgfpathrectangle{\pgfqpoint{0.573768in}{0.462654in}}{\pgfqpoint{2.720000in}{0.906250in}}%
\pgfusepath{clip}%
\pgfsetbuttcap%
\pgfsetroundjoin%
\pgfsetlinewidth{0.501875pt}%
\definecolor{currentstroke}{rgb}{0.690196,0.690196,0.690196}%
\pgfsetstrokecolor{currentstroke}%
\pgfsetstrokeopacity{0.250000}%
\pgfsetdash{{1.850000pt}{0.800000pt}}{0.000000pt}%
\pgfpathmoveto{\pgfqpoint{2.181041in}{0.462654in}}%
\pgfpathlineto{\pgfqpoint{2.181041in}{1.368904in}}%
\pgfusepath{stroke}%
\end{pgfscope}%
\begin{pgfscope}%
\pgfsetbuttcap%
\pgfsetroundjoin%
\definecolor{currentfill}{rgb}{0.000000,0.000000,0.000000}%
\pgfsetfillcolor{currentfill}%
\pgfsetlinewidth{0.803000pt}%
\definecolor{currentstroke}{rgb}{0.000000,0.000000,0.000000}%
\pgfsetstrokecolor{currentstroke}%
\pgfsetdash{}{0pt}%
\pgfsys@defobject{currentmarker}{\pgfqpoint{0.000000in}{-0.048611in}}{\pgfqpoint{0.000000in}{0.000000in}}{%
\pgfpathmoveto{\pgfqpoint{0.000000in}{0.000000in}}%
\pgfpathlineto{\pgfqpoint{0.000000in}{-0.048611in}}%
\pgfusepath{stroke,fill}%
}%
\begin{pgfscope}%
\pgfsys@transformshift{2.181041in}{0.462654in}%
\pgfsys@useobject{currentmarker}{}%
\end{pgfscope}%
\end{pgfscope}%
\begin{pgfscope}%
\definecolor{textcolor}{rgb}{0.000000,0.000000,0.000000}%
\pgfsetstrokecolor{textcolor}%
\pgfsetfillcolor{textcolor}%
\pgftext[x=2.181041in,y=0.365432in,,top]{\color{textcolor}{\rmfamily\fontsize{8.000000}{9.600000}\selectfont\catcode`\^=\active\def^{\ifmmode\sp\else\^{}\fi}\catcode`\%=\active\def%{\%}$\mathdefault{3}$}}%
\end{pgfscope}%
\begin{pgfscope}%
\pgfpathrectangle{\pgfqpoint{0.573768in}{0.462654in}}{\pgfqpoint{2.720000in}{0.906250in}}%
\pgfusepath{clip}%
\pgfsetbuttcap%
\pgfsetroundjoin%
\pgfsetlinewidth{0.501875pt}%
\definecolor{currentstroke}{rgb}{0.690196,0.690196,0.690196}%
\pgfsetstrokecolor{currentstroke}%
\pgfsetstrokeopacity{0.250000}%
\pgfsetdash{{1.850000pt}{0.800000pt}}{0.000000pt}%
\pgfpathmoveto{\pgfqpoint{2.675587in}{0.462654in}}%
\pgfpathlineto{\pgfqpoint{2.675587in}{1.368904in}}%
\pgfusepath{stroke}%
\end{pgfscope}%
\begin{pgfscope}%
\pgfsetbuttcap%
\pgfsetroundjoin%
\definecolor{currentfill}{rgb}{0.000000,0.000000,0.000000}%
\pgfsetfillcolor{currentfill}%
\pgfsetlinewidth{0.803000pt}%
\definecolor{currentstroke}{rgb}{0.000000,0.000000,0.000000}%
\pgfsetstrokecolor{currentstroke}%
\pgfsetdash{}{0pt}%
\pgfsys@defobject{currentmarker}{\pgfqpoint{0.000000in}{-0.048611in}}{\pgfqpoint{0.000000in}{0.000000in}}{%
\pgfpathmoveto{\pgfqpoint{0.000000in}{0.000000in}}%
\pgfpathlineto{\pgfqpoint{0.000000in}{-0.048611in}}%
\pgfusepath{stroke,fill}%
}%
\begin{pgfscope}%
\pgfsys@transformshift{2.675587in}{0.462654in}%
\pgfsys@useobject{currentmarker}{}%
\end{pgfscope}%
\end{pgfscope}%
\begin{pgfscope}%
\definecolor{textcolor}{rgb}{0.000000,0.000000,0.000000}%
\pgfsetstrokecolor{textcolor}%
\pgfsetfillcolor{textcolor}%
\pgftext[x=2.675587in,y=0.365432in,,top]{\color{textcolor}{\rmfamily\fontsize{8.000000}{9.600000}\selectfont\catcode`\^=\active\def^{\ifmmode\sp\else\^{}\fi}\catcode`\%=\active\def%{\%}$\mathdefault{4}$}}%
\end{pgfscope}%
\begin{pgfscope}%
\pgfpathrectangle{\pgfqpoint{0.573768in}{0.462654in}}{\pgfqpoint{2.720000in}{0.906250in}}%
\pgfusepath{clip}%
\pgfsetbuttcap%
\pgfsetroundjoin%
\pgfsetlinewidth{0.501875pt}%
\definecolor{currentstroke}{rgb}{0.690196,0.690196,0.690196}%
\pgfsetstrokecolor{currentstroke}%
\pgfsetstrokeopacity{0.250000}%
\pgfsetdash{{1.850000pt}{0.800000pt}}{0.000000pt}%
\pgfpathmoveto{\pgfqpoint{3.170132in}{0.462654in}}%
\pgfpathlineto{\pgfqpoint{3.170132in}{1.368904in}}%
\pgfusepath{stroke}%
\end{pgfscope}%
\begin{pgfscope}%
\pgfsetbuttcap%
\pgfsetroundjoin%
\definecolor{currentfill}{rgb}{0.000000,0.000000,0.000000}%
\pgfsetfillcolor{currentfill}%
\pgfsetlinewidth{0.803000pt}%
\definecolor{currentstroke}{rgb}{0.000000,0.000000,0.000000}%
\pgfsetstrokecolor{currentstroke}%
\pgfsetdash{}{0pt}%
\pgfsys@defobject{currentmarker}{\pgfqpoint{0.000000in}{-0.048611in}}{\pgfqpoint{0.000000in}{0.000000in}}{%
\pgfpathmoveto{\pgfqpoint{0.000000in}{0.000000in}}%
\pgfpathlineto{\pgfqpoint{0.000000in}{-0.048611in}}%
\pgfusepath{stroke,fill}%
}%
\begin{pgfscope}%
\pgfsys@transformshift{3.170132in}{0.462654in}%
\pgfsys@useobject{currentmarker}{}%
\end{pgfscope}%
\end{pgfscope}%
\begin{pgfscope}%
\definecolor{textcolor}{rgb}{0.000000,0.000000,0.000000}%
\pgfsetstrokecolor{textcolor}%
\pgfsetfillcolor{textcolor}%
\pgftext[x=3.170132in,y=0.365432in,,top]{\color{textcolor}{\rmfamily\fontsize{8.000000}{9.600000}\selectfont\catcode`\^=\active\def^{\ifmmode\sp\else\^{}\fi}\catcode`\%=\active\def%{\%}$\mathdefault{5}$}}%
\end{pgfscope}%
\begin{pgfscope}%
\definecolor{textcolor}{rgb}{0.000000,0.000000,0.000000}%
\pgfsetstrokecolor{textcolor}%
\pgfsetfillcolor{textcolor}%
\pgftext[x=1.933768in,y=0.211111in,,top]{\color{textcolor}{\rmfamily\fontsize{8.000000}{9.600000}\selectfont\catcode`\^=\active\def^{\ifmmode\sp\else\^{}\fi}\catcode`\%=\active\def%{\%}Time (ms)}}%
\end{pgfscope}%
\begin{pgfscope}%
\pgfpathrectangle{\pgfqpoint{0.573768in}{0.462654in}}{\pgfqpoint{2.720000in}{0.906250in}}%
\pgfusepath{clip}%
\pgfsetbuttcap%
\pgfsetroundjoin%
\pgfsetlinewidth{0.501875pt}%
\definecolor{currentstroke}{rgb}{0.690196,0.690196,0.690196}%
\pgfsetstrokecolor{currentstroke}%
\pgfsetstrokeopacity{0.250000}%
\pgfsetdash{{1.850000pt}{0.800000pt}}{0.000000pt}%
\pgfpathmoveto{\pgfqpoint{0.573768in}{0.462654in}}%
\pgfpathlineto{\pgfqpoint{3.293768in}{0.462654in}}%
\pgfusepath{stroke}%
\end{pgfscope}%
\begin{pgfscope}%
\pgfsetbuttcap%
\pgfsetroundjoin%
\definecolor{currentfill}{rgb}{0.000000,0.000000,0.000000}%
\pgfsetfillcolor{currentfill}%
\pgfsetlinewidth{0.803000pt}%
\definecolor{currentstroke}{rgb}{0.000000,0.000000,0.000000}%
\pgfsetstrokecolor{currentstroke}%
\pgfsetdash{}{0pt}%
\pgfsys@defobject{currentmarker}{\pgfqpoint{-0.048611in}{0.000000in}}{\pgfqpoint{-0.000000in}{0.000000in}}{%
\pgfpathmoveto{\pgfqpoint{-0.000000in}{0.000000in}}%
\pgfpathlineto{\pgfqpoint{-0.048611in}{0.000000in}}%
\pgfusepath{stroke,fill}%
}%
\begin{pgfscope}%
\pgfsys@transformshift{0.573768in}{0.462654in}%
\pgfsys@useobject{currentmarker}{}%
\end{pgfscope}%
\end{pgfscope}%
\begin{pgfscope}%
\definecolor{textcolor}{rgb}{0.000000,0.000000,0.000000}%
\pgfsetstrokecolor{textcolor}%
\pgfsetfillcolor{textcolor}%
\pgftext[x=0.325695in, y=0.424074in, left, base]{\color{textcolor}{\rmfamily\fontsize{8.000000}{9.600000}\selectfont\catcode`\^=\active\def^{\ifmmode\sp\else\^{}\fi}\catcode`\%=\active\def%{\%}0.0}}%
\end{pgfscope}%
\begin{pgfscope}%
\pgfpathrectangle{\pgfqpoint{0.573768in}{0.462654in}}{\pgfqpoint{2.720000in}{0.906250in}}%
\pgfusepath{clip}%
\pgfsetbuttcap%
\pgfsetroundjoin%
\pgfsetlinewidth{0.501875pt}%
\definecolor{currentstroke}{rgb}{0.690196,0.690196,0.690196}%
\pgfsetstrokecolor{currentstroke}%
\pgfsetstrokeopacity{0.250000}%
\pgfsetdash{{1.850000pt}{0.800000pt}}{0.000000pt}%
\pgfpathmoveto{\pgfqpoint{0.573768in}{0.689217in}}%
\pgfpathlineto{\pgfqpoint{3.293768in}{0.689217in}}%
\pgfusepath{stroke}%
\end{pgfscope}%
\begin{pgfscope}%
\pgfsetbuttcap%
\pgfsetroundjoin%
\definecolor{currentfill}{rgb}{0.000000,0.000000,0.000000}%
\pgfsetfillcolor{currentfill}%
\pgfsetlinewidth{0.803000pt}%
\definecolor{currentstroke}{rgb}{0.000000,0.000000,0.000000}%
\pgfsetstrokecolor{currentstroke}%
\pgfsetdash{}{0pt}%
\pgfsys@defobject{currentmarker}{\pgfqpoint{-0.048611in}{0.000000in}}{\pgfqpoint{-0.000000in}{0.000000in}}{%
\pgfpathmoveto{\pgfqpoint{-0.000000in}{0.000000in}}%
\pgfpathlineto{\pgfqpoint{-0.048611in}{0.000000in}}%
\pgfusepath{stroke,fill}%
}%
\begin{pgfscope}%
\pgfsys@transformshift{0.573768in}{0.689217in}%
\pgfsys@useobject{currentmarker}{}%
\end{pgfscope}%
\end{pgfscope}%
\begin{pgfscope}%
\definecolor{textcolor}{rgb}{0.000000,0.000000,0.000000}%
\pgfsetstrokecolor{textcolor}%
\pgfsetfillcolor{textcolor}%
\pgftext[x=0.325695in, y=0.650637in, left, base]{\color{textcolor}{\rmfamily\fontsize{8.000000}{9.600000}\selectfont\catcode`\^=\active\def^{\ifmmode\sp\else\^{}\fi}\catcode`\%=\active\def%{\%}6.2}}%
\end{pgfscope}%
\begin{pgfscope}%
\pgfpathrectangle{\pgfqpoint{0.573768in}{0.462654in}}{\pgfqpoint{2.720000in}{0.906250in}}%
\pgfusepath{clip}%
\pgfsetbuttcap%
\pgfsetroundjoin%
\pgfsetlinewidth{0.501875pt}%
\definecolor{currentstroke}{rgb}{0.690196,0.690196,0.690196}%
\pgfsetstrokecolor{currentstroke}%
\pgfsetstrokeopacity{0.250000}%
\pgfsetdash{{1.850000pt}{0.800000pt}}{0.000000pt}%
\pgfpathmoveto{\pgfqpoint{0.573768in}{0.915779in}}%
\pgfpathlineto{\pgfqpoint{3.293768in}{0.915779in}}%
\pgfusepath{stroke}%
\end{pgfscope}%
\begin{pgfscope}%
\pgfsetbuttcap%
\pgfsetroundjoin%
\definecolor{currentfill}{rgb}{0.000000,0.000000,0.000000}%
\pgfsetfillcolor{currentfill}%
\pgfsetlinewidth{0.803000pt}%
\definecolor{currentstroke}{rgb}{0.000000,0.000000,0.000000}%
\pgfsetstrokecolor{currentstroke}%
\pgfsetdash{}{0pt}%
\pgfsys@defobject{currentmarker}{\pgfqpoint{-0.048611in}{0.000000in}}{\pgfqpoint{-0.000000in}{0.000000in}}{%
\pgfpathmoveto{\pgfqpoint{-0.000000in}{0.000000in}}%
\pgfpathlineto{\pgfqpoint{-0.048611in}{0.000000in}}%
\pgfusepath{stroke,fill}%
}%
\begin{pgfscope}%
\pgfsys@transformshift{0.573768in}{0.915779in}%
\pgfsys@useobject{currentmarker}{}%
\end{pgfscope}%
\end{pgfscope}%
\begin{pgfscope}%
\definecolor{textcolor}{rgb}{0.000000,0.000000,0.000000}%
\pgfsetstrokecolor{textcolor}%
\pgfsetfillcolor{textcolor}%
\pgftext[x=0.266667in, y=0.877199in, left, base]{\color{textcolor}{\rmfamily\fontsize{8.000000}{9.600000}\selectfont\catcode`\^=\active\def^{\ifmmode\sp\else\^{}\fi}\catcode`\%=\active\def%{\%}12.5}}%
\end{pgfscope}%
\begin{pgfscope}%
\pgfpathrectangle{\pgfqpoint{0.573768in}{0.462654in}}{\pgfqpoint{2.720000in}{0.906250in}}%
\pgfusepath{clip}%
\pgfsetbuttcap%
\pgfsetroundjoin%
\pgfsetlinewidth{0.501875pt}%
\definecolor{currentstroke}{rgb}{0.690196,0.690196,0.690196}%
\pgfsetstrokecolor{currentstroke}%
\pgfsetstrokeopacity{0.250000}%
\pgfsetdash{{1.850000pt}{0.800000pt}}{0.000000pt}%
\pgfpathmoveto{\pgfqpoint{0.573768in}{1.142342in}}%
\pgfpathlineto{\pgfqpoint{3.293768in}{1.142342in}}%
\pgfusepath{stroke}%
\end{pgfscope}%
\begin{pgfscope}%
\pgfsetbuttcap%
\pgfsetroundjoin%
\definecolor{currentfill}{rgb}{0.000000,0.000000,0.000000}%
\pgfsetfillcolor{currentfill}%
\pgfsetlinewidth{0.803000pt}%
\definecolor{currentstroke}{rgb}{0.000000,0.000000,0.000000}%
\pgfsetstrokecolor{currentstroke}%
\pgfsetdash{}{0pt}%
\pgfsys@defobject{currentmarker}{\pgfqpoint{-0.048611in}{0.000000in}}{\pgfqpoint{-0.000000in}{0.000000in}}{%
\pgfpathmoveto{\pgfqpoint{-0.000000in}{0.000000in}}%
\pgfpathlineto{\pgfqpoint{-0.048611in}{0.000000in}}%
\pgfusepath{stroke,fill}%
}%
\begin{pgfscope}%
\pgfsys@transformshift{0.573768in}{1.142342in}%
\pgfsys@useobject{currentmarker}{}%
\end{pgfscope}%
\end{pgfscope}%
\begin{pgfscope}%
\definecolor{textcolor}{rgb}{0.000000,0.000000,0.000000}%
\pgfsetstrokecolor{textcolor}%
\pgfsetfillcolor{textcolor}%
\pgftext[x=0.266667in, y=1.103762in, left, base]{\color{textcolor}{\rmfamily\fontsize{8.000000}{9.600000}\selectfont\catcode`\^=\active\def^{\ifmmode\sp\else\^{}\fi}\catcode`\%=\active\def%{\%}18.8}}%
\end{pgfscope}%
\begin{pgfscope}%
\pgfpathrectangle{\pgfqpoint{0.573768in}{0.462654in}}{\pgfqpoint{2.720000in}{0.906250in}}%
\pgfusepath{clip}%
\pgfsetbuttcap%
\pgfsetroundjoin%
\pgfsetlinewidth{0.501875pt}%
\definecolor{currentstroke}{rgb}{0.690196,0.690196,0.690196}%
\pgfsetstrokecolor{currentstroke}%
\pgfsetstrokeopacity{0.250000}%
\pgfsetdash{{1.850000pt}{0.800000pt}}{0.000000pt}%
\pgfpathmoveto{\pgfqpoint{0.573768in}{1.368904in}}%
\pgfpathlineto{\pgfqpoint{3.293768in}{1.368904in}}%
\pgfusepath{stroke}%
\end{pgfscope}%
\begin{pgfscope}%
\pgfsetbuttcap%
\pgfsetroundjoin%
\definecolor{currentfill}{rgb}{0.000000,0.000000,0.000000}%
\pgfsetfillcolor{currentfill}%
\pgfsetlinewidth{0.803000pt}%
\definecolor{currentstroke}{rgb}{0.000000,0.000000,0.000000}%
\pgfsetstrokecolor{currentstroke}%
\pgfsetdash{}{0pt}%
\pgfsys@defobject{currentmarker}{\pgfqpoint{-0.048611in}{0.000000in}}{\pgfqpoint{-0.000000in}{0.000000in}}{%
\pgfpathmoveto{\pgfqpoint{-0.000000in}{0.000000in}}%
\pgfpathlineto{\pgfqpoint{-0.048611in}{0.000000in}}%
\pgfusepath{stroke,fill}%
}%
\begin{pgfscope}%
\pgfsys@transformshift{0.573768in}{1.368904in}%
\pgfsys@useobject{currentmarker}{}%
\end{pgfscope}%
\end{pgfscope}%
\begin{pgfscope}%
\definecolor{textcolor}{rgb}{0.000000,0.000000,0.000000}%
\pgfsetstrokecolor{textcolor}%
\pgfsetfillcolor{textcolor}%
\pgftext[x=0.266667in, y=1.330324in, left, base]{\color{textcolor}{\rmfamily\fontsize{8.000000}{9.600000}\selectfont\catcode`\^=\active\def^{\ifmmode\sp\else\^{}\fi}\catcode`\%=\active\def%{\%}25.0}}%
\end{pgfscope}%
\begin{pgfscope}%
\definecolor{textcolor}{rgb}{0.000000,0.000000,0.000000}%
\pgfsetstrokecolor{textcolor}%
\pgfsetfillcolor{textcolor}%
\pgftext[x=0.211111in,y=0.915779in,,bottom,rotate=90.000000]{\color{textcolor}{\rmfamily\fontsize{8.000000}{9.600000}\selectfont\catcode`\^=\active\def^{\ifmmode\sp\else\^{}\fi}\catcode`\%=\active\def%{\%}$v_o$ (V)}}%
\end{pgfscope}%
\begin{pgfscope}%
\pgfpathrectangle{\pgfqpoint{0.573768in}{0.462654in}}{\pgfqpoint{2.720000in}{0.906250in}}%
\pgfusepath{clip}%
\pgfsetrectcap%
\pgfsetroundjoin%
\pgfsetlinewidth{1.505625pt}%
\definecolor{currentstroke}{rgb}{0.000000,0.619608,0.450980}%
\pgfsetstrokecolor{currentstroke}%
\pgfsetdash{}{0pt}%
\pgfpathmoveto{\pgfqpoint{0.697405in}{0.462654in}}%
\pgfpathlineto{\pgfqpoint{0.699624in}{0.463415in}}%
\pgfpathlineto{\pgfqpoint{0.702417in}{0.474990in}}%
\pgfpathlineto{\pgfqpoint{0.704484in}{0.475690in}}%
\pgfpathlineto{\pgfqpoint{0.707683in}{0.498852in}}%
\pgfpathlineto{\pgfqpoint{0.709394in}{0.499316in}}%
\pgfpathlineto{\pgfqpoint{0.709403in}{0.499454in}}%
\pgfpathlineto{\pgfqpoint{0.712266in}{0.533467in}}%
\pgfpathlineto{\pgfqpoint{0.712772in}{0.533394in}}%
\pgfpathlineto{\pgfqpoint{0.714286in}{0.533199in}}%
\pgfpathlineto{\pgfqpoint{0.714337in}{0.533907in}}%
\pgfpathlineto{\pgfqpoint{0.717214in}{0.577726in}}%
\pgfpathlineto{\pgfqpoint{0.718027in}{0.577537in}}%
\pgfpathlineto{\pgfqpoint{0.719233in}{0.577279in}}%
\pgfpathlineto{\pgfqpoint{0.719274in}{0.577989in}}%
\pgfpathlineto{\pgfqpoint{0.722162in}{0.630249in}}%
\pgfpathlineto{\pgfqpoint{0.723509in}{0.629793in}}%
\pgfpathlineto{\pgfqpoint{0.724182in}{0.629607in}}%
\pgfpathlineto{\pgfqpoint{0.724216in}{0.630256in}}%
\pgfpathlineto{\pgfqpoint{0.727109in}{0.689466in}}%
\pgfpathlineto{\pgfqpoint{0.727651in}{0.689218in}}%
\pgfpathlineto{\pgfqpoint{0.729123in}{0.688562in}}%
\pgfpathlineto{\pgfqpoint{0.729165in}{0.689362in}}%
\pgfpathlineto{\pgfqpoint{0.732052in}{0.753631in}}%
\pgfpathlineto{\pgfqpoint{0.732644in}{0.753283in}}%
\pgfpathlineto{\pgfqpoint{0.734059in}{0.752452in}}%
\pgfpathlineto{\pgfqpoint{0.734093in}{0.752752in}}%
\pgfpathlineto{\pgfqpoint{0.736997in}{0.820905in}}%
\pgfpathlineto{\pgfqpoint{0.737335in}{0.820661in}}%
\pgfpathlineto{\pgfqpoint{0.739016in}{0.819461in}}%
\pgfpathlineto{\pgfqpoint{0.739049in}{0.819995in}}%
\pgfpathlineto{\pgfqpoint{0.741942in}{0.889408in}}%
\pgfpathlineto{\pgfqpoint{0.742247in}{0.889146in}}%
\pgfpathlineto{\pgfqpoint{0.743963in}{0.887685in}}%
\pgfpathlineto{\pgfqpoint{0.744011in}{0.888630in}}%
\pgfpathlineto{\pgfqpoint{0.746888in}{0.957213in}}%
\pgfpathlineto{\pgfqpoint{0.747661in}{0.956440in}}%
\pgfpathlineto{\pgfqpoint{0.748911in}{0.955215in}}%
\pgfpathlineto{\pgfqpoint{0.748951in}{0.956028in}}%
\pgfpathlineto{\pgfqpoint{0.751831in}{1.022482in}}%
\pgfpathlineto{\pgfqpoint{0.752240in}{1.022023in}}%
\pgfpathlineto{\pgfqpoint{0.753855in}{1.020209in}}%
\pgfpathlineto{\pgfqpoint{0.753907in}{1.021210in}}%
\pgfpathlineto{\pgfqpoint{0.756776in}{1.083442in}}%
\pgfpathlineto{\pgfqpoint{0.757527in}{1.082500in}}%
\pgfpathlineto{\pgfqpoint{0.758800in}{1.080910in}}%
\pgfpathlineto{\pgfqpoint{0.758855in}{1.081833in}}%
\pgfpathlineto{\pgfqpoint{0.761721in}{1.138440in}}%
\pgfpathlineto{\pgfqpoint{0.762360in}{1.137567in}}%
\pgfpathlineto{\pgfqpoint{0.763743in}{1.135683in}}%
\pgfpathlineto{\pgfqpoint{0.763817in}{1.136801in}}%
\pgfpathlineto{\pgfqpoint{0.766667in}{1.186070in}}%
\pgfpathlineto{\pgfqpoint{0.767010in}{1.185567in}}%
\pgfpathlineto{\pgfqpoint{0.768693in}{1.183113in}}%
\pgfpathlineto{\pgfqpoint{0.768774in}{1.184335in}}%
\pgfpathlineto{\pgfqpoint{0.771265in}{1.221380in}}%
\pgfpathlineto{\pgfqpoint{0.771607in}{1.225201in}}%
\pgfpathlineto{\pgfqpoint{0.772270in}{1.224181in}}%
\pgfpathlineto{\pgfqpoint{0.773638in}{1.222076in}}%
\pgfpathlineto{\pgfqpoint{0.773729in}{1.223276in}}%
\pgfpathlineto{\pgfqpoint{0.776118in}{1.251388in}}%
\pgfpathlineto{\pgfqpoint{0.776555in}{1.254856in}}%
\pgfpathlineto{\pgfqpoint{0.777143in}{1.253913in}}%
\pgfpathlineto{\pgfqpoint{0.778593in}{1.251624in}}%
\pgfpathlineto{\pgfqpoint{0.778706in}{1.252862in}}%
\pgfpathlineto{\pgfqpoint{0.781495in}{1.274303in}}%
\pgfpathlineto{\pgfqpoint{0.781513in}{1.274274in}}%
\pgfpathlineto{\pgfqpoint{0.783526in}{1.270977in}}%
\pgfpathlineto{\pgfqpoint{0.783685in}{1.272166in}}%
\pgfpathlineto{\pgfqpoint{0.786087in}{1.282999in}}%
\pgfpathlineto{\pgfqpoint{0.786439in}{1.283259in}}%
\pgfpathlineto{\pgfqpoint{0.786761in}{1.282724in}}%
\pgfpathlineto{\pgfqpoint{0.788481in}{1.279880in}}%
\pgfpathlineto{\pgfqpoint{0.788741in}{1.280918in}}%
\pgfpathlineto{\pgfqpoint{0.790131in}{1.283607in}}%
\pgfpathlineto{\pgfqpoint{0.790828in}{1.282977in}}%
\pgfpathlineto{\pgfqpoint{0.793440in}{1.278684in}}%
\pgfpathlineto{\pgfqpoint{0.793807in}{1.279531in}}%
\pgfpathlineto{\pgfqpoint{0.794599in}{1.280111in}}%
\pgfpathlineto{\pgfqpoint{0.794779in}{1.280004in}}%
\pgfpathlineto{\pgfqpoint{0.796456in}{1.277357in}}%
\pgfpathlineto{\pgfqpoint{0.798380in}{1.274232in}}%
\pgfpathlineto{\pgfqpoint{0.798733in}{1.275087in}}%
\pgfpathlineto{\pgfqpoint{0.799740in}{1.275642in}}%
\pgfpathlineto{\pgfqpoint{0.799970in}{1.275385in}}%
\pgfpathlineto{\pgfqpoint{0.803334in}{1.269915in}}%
\pgfpathlineto{\pgfqpoint{0.803855in}{1.271044in}}%
\pgfpathlineto{\pgfqpoint{0.804949in}{1.271005in}}%
\pgfpathlineto{\pgfqpoint{0.805010in}{1.270909in}}%
\pgfpathlineto{\pgfqpoint{0.808278in}{1.265614in}}%
\pgfpathlineto{\pgfqpoint{0.808632in}{1.266464in}}%
\pgfpathlineto{\pgfqpoint{0.809834in}{1.266821in}}%
\pgfpathlineto{\pgfqpoint{0.809982in}{1.266600in}}%
\pgfpathlineto{\pgfqpoint{0.813226in}{1.261383in}}%
\pgfpathlineto{\pgfqpoint{0.813578in}{1.262222in}}%
\pgfpathlineto{\pgfqpoint{0.814526in}{1.262827in}}%
\pgfpathlineto{\pgfqpoint{0.814748in}{1.262622in}}%
\pgfpathlineto{\pgfqpoint{0.817073in}{1.258901in}}%
\pgfpathlineto{\pgfqpoint{0.818168in}{1.257169in}}%
\pgfpathlineto{\pgfqpoint{0.818522in}{1.258024in}}%
\pgfpathlineto{\pgfqpoint{0.819748in}{1.258389in}}%
\pgfpathlineto{\pgfqpoint{0.819901in}{1.258156in}}%
\pgfpathlineto{\pgfqpoint{0.823107in}{1.253048in}}%
\pgfpathlineto{\pgfqpoint{0.823456in}{1.253892in}}%
\pgfpathlineto{\pgfqpoint{0.824306in}{1.254595in}}%
\pgfpathlineto{\pgfqpoint{0.824469in}{1.254510in}}%
\pgfpathlineto{\pgfqpoint{0.825390in}{1.253184in}}%
\pgfpathlineto{\pgfqpoint{0.828054in}{1.248972in}}%
\pgfpathlineto{\pgfqpoint{0.828401in}{1.249817in}}%
\pgfpathlineto{\pgfqpoint{0.829274in}{1.250530in}}%
\pgfpathlineto{\pgfqpoint{0.829429in}{1.250446in}}%
\pgfpathlineto{\pgfqpoint{0.831080in}{1.247961in}}%
\pgfpathlineto{\pgfqpoint{0.833001in}{1.244956in}}%
\pgfpathlineto{\pgfqpoint{0.833348in}{1.245799in}}%
\pgfpathlineto{\pgfqpoint{0.834249in}{1.246502in}}%
\pgfpathlineto{\pgfqpoint{0.834395in}{1.246417in}}%
\pgfpathlineto{\pgfqpoint{0.835851in}{1.244239in}}%
\pgfpathlineto{\pgfqpoint{0.837950in}{1.240956in}}%
\pgfpathlineto{\pgfqpoint{0.838299in}{1.241809in}}%
\pgfpathlineto{\pgfqpoint{0.839415in}{1.242370in}}%
\pgfpathlineto{\pgfqpoint{0.839631in}{1.242105in}}%
\pgfpathlineto{\pgfqpoint{0.842896in}{1.237012in}}%
\pgfpathlineto{\pgfqpoint{0.843338in}{1.238023in}}%
\pgfpathlineto{\pgfqpoint{0.844327in}{1.238448in}}%
\pgfpathlineto{\pgfqpoint{0.844519in}{1.238236in}}%
\pgfpathlineto{\pgfqpoint{0.847842in}{1.233087in}}%
\pgfpathlineto{\pgfqpoint{0.848351in}{1.234213in}}%
\pgfpathlineto{\pgfqpoint{0.849121in}{1.234644in}}%
\pgfpathlineto{\pgfqpoint{0.849307in}{1.234519in}}%
\pgfpathlineto{\pgfqpoint{0.851930in}{1.230517in}}%
\pgfpathlineto{\pgfqpoint{0.852785in}{1.229219in}}%
\pgfpathlineto{\pgfqpoint{0.853131in}{1.230069in}}%
\pgfpathlineto{\pgfqpoint{0.854097in}{1.230777in}}%
\pgfpathlineto{\pgfqpoint{0.854286in}{1.230637in}}%
\pgfpathlineto{\pgfqpoint{0.857359in}{1.225945in}}%
\pgfpathlineto{\pgfqpoint{0.857731in}{1.225404in}}%
\pgfpathlineto{\pgfqpoint{0.858075in}{1.226248in}}%
\pgfpathlineto{\pgfqpoint{0.859040in}{1.226965in}}%
\pgfpathlineto{\pgfqpoint{0.859220in}{1.226838in}}%
\pgfpathlineto{\pgfqpoint{0.861486in}{1.223416in}}%
\pgfpathlineto{\pgfqpoint{0.862671in}{1.221619in}}%
\pgfpathlineto{\pgfqpoint{0.863015in}{1.222462in}}%
\pgfpathlineto{\pgfqpoint{0.863929in}{1.223235in}}%
\pgfpathlineto{\pgfqpoint{0.864045in}{1.223180in}}%
\pgfpathlineto{\pgfqpoint{0.865020in}{1.221853in}}%
\pgfpathlineto{\pgfqpoint{0.867621in}{1.217902in}}%
\pgfpathlineto{\pgfqpoint{0.867964in}{1.218743in}}%
\pgfpathlineto{\pgfqpoint{0.868908in}{1.219495in}}%
\pgfpathlineto{\pgfqpoint{0.869029in}{1.219428in}}%
\pgfpathlineto{\pgfqpoint{0.870645in}{1.217098in}}%
\pgfpathlineto{\pgfqpoint{0.872562in}{1.214210in}}%
\pgfpathlineto{\pgfqpoint{0.872907in}{1.215054in}}%
\pgfpathlineto{\pgfqpoint{0.873844in}{1.215836in}}%
\pgfpathlineto{\pgfqpoint{0.873955in}{1.215782in}}%
\pgfpathlineto{\pgfqpoint{0.874974in}{1.214394in}}%
\pgfpathlineto{\pgfqpoint{0.877508in}{1.210586in}}%
\pgfpathlineto{\pgfqpoint{0.877849in}{1.211418in}}%
\pgfpathlineto{\pgfqpoint{0.878817in}{1.212200in}}%
\pgfpathlineto{\pgfqpoint{0.878948in}{1.212123in}}%
\pgfpathlineto{\pgfqpoint{0.880161in}{1.210412in}}%
\pgfpathlineto{\pgfqpoint{0.882456in}{1.206989in}}%
\pgfpathlineto{\pgfqpoint{0.882796in}{1.207821in}}%
\pgfpathlineto{\pgfqpoint{0.883740in}{1.208620in}}%
\pgfpathlineto{\pgfqpoint{0.883877in}{1.208551in}}%
\pgfpathlineto{\pgfqpoint{0.884910in}{1.207135in}}%
\pgfpathlineto{\pgfqpoint{0.887404in}{1.203424in}}%
\pgfpathlineto{\pgfqpoint{0.887744in}{1.204259in}}%
\pgfpathlineto{\pgfqpoint{0.888688in}{1.205048in}}%
\pgfpathlineto{\pgfqpoint{0.888809in}{1.204990in}}%
\pgfpathlineto{\pgfqpoint{0.890543in}{1.202546in}}%
\pgfpathlineto{\pgfqpoint{0.892348in}{1.199883in}}%
\pgfpathlineto{\pgfqpoint{0.892686in}{1.200710in}}%
\pgfpathlineto{\pgfqpoint{0.893630in}{1.201524in}}%
\pgfpathlineto{\pgfqpoint{0.893737in}{1.201476in}}%
\pgfpathlineto{\pgfqpoint{0.894617in}{1.200338in}}%
\pgfpathlineto{\pgfqpoint{0.897293in}{1.196397in}}%
\pgfpathlineto{\pgfqpoint{0.897632in}{1.197227in}}%
\pgfpathlineto{\pgfqpoint{0.898603in}{1.198037in}}%
\pgfpathlineto{\pgfqpoint{0.898698in}{1.197992in}}%
\pgfpathlineto{\pgfqpoint{0.900251in}{1.195850in}}%
\pgfpathlineto{\pgfqpoint{0.902239in}{1.192944in}}%
\pgfpathlineto{\pgfqpoint{0.902576in}{1.193772in}}%
\pgfpathlineto{\pgfqpoint{0.903546in}{1.194593in}}%
\pgfpathlineto{\pgfqpoint{0.903635in}{1.194553in}}%
\pgfpathlineto{\pgfqpoint{0.904773in}{1.193046in}}%
\pgfpathlineto{\pgfqpoint{0.907187in}{1.189521in}}%
\pgfpathlineto{\pgfqpoint{0.907523in}{1.190345in}}%
\pgfpathlineto{\pgfqpoint{0.908502in}{1.191173in}}%
\pgfpathlineto{\pgfqpoint{0.908577in}{1.191138in}}%
\pgfpathlineto{\pgfqpoint{0.909539in}{1.189907in}}%
\pgfpathlineto{\pgfqpoint{0.912125in}{1.186149in}}%
\pgfpathlineto{\pgfqpoint{0.912460in}{1.186973in}}%
\pgfpathlineto{\pgfqpoint{0.913521in}{1.187787in}}%
\pgfpathlineto{\pgfqpoint{0.913687in}{1.187667in}}%
\pgfpathlineto{\pgfqpoint{0.915241in}{1.185456in}}%
\pgfpathlineto{\pgfqpoint{0.917070in}{1.182809in}}%
\pgfpathlineto{\pgfqpoint{0.917411in}{1.183644in}}%
\pgfpathlineto{\pgfqpoint{0.918426in}{1.184485in}}%
\pgfpathlineto{\pgfqpoint{0.918523in}{1.184434in}}%
\pgfpathlineto{\pgfqpoint{0.920331in}{1.181942in}}%
\pgfpathlineto{\pgfqpoint{0.922016in}{1.179518in}}%
\pgfpathlineto{\pgfqpoint{0.922356in}{1.180351in}}%
\pgfpathlineto{\pgfqpoint{0.923370in}{1.181196in}}%
\pgfpathlineto{\pgfqpoint{0.923471in}{1.181143in}}%
\pgfpathlineto{\pgfqpoint{0.924850in}{1.179285in}}%
\pgfpathlineto{\pgfqpoint{0.926962in}{1.176255in}}%
\pgfpathlineto{\pgfqpoint{0.927301in}{1.177089in}}%
\pgfpathlineto{\pgfqpoint{0.928317in}{1.177944in}}%
\pgfpathlineto{\pgfqpoint{0.928424in}{1.177888in}}%
\pgfpathlineto{\pgfqpoint{0.929444in}{1.176556in}}%
\pgfpathlineto{\pgfqpoint{0.931909in}{1.173044in}}%
\pgfpathlineto{\pgfqpoint{0.932244in}{1.173867in}}%
\pgfpathlineto{\pgfqpoint{0.933369in}{1.174677in}}%
\pgfpathlineto{\pgfqpoint{0.933532in}{1.174544in}}%
\pgfpathlineto{\pgfqpoint{0.935617in}{1.171590in}}%
\pgfpathlineto{\pgfqpoint{0.936850in}{1.169846in}}%
\pgfpathlineto{\pgfqpoint{0.937186in}{1.170681in}}%
\pgfpathlineto{\pgfqpoint{0.938559in}{1.171286in}}%
\pgfpathlineto{\pgfqpoint{0.938750in}{1.171027in}}%
\pgfpathlineto{\pgfqpoint{0.941798in}{1.166706in}}%
\pgfpathlineto{\pgfqpoint{0.942130in}{1.167531in}}%
\pgfpathlineto{\pgfqpoint{0.943430in}{1.168233in}}%
\pgfpathlineto{\pgfqpoint{0.943626in}{1.167993in}}%
\pgfpathlineto{\pgfqpoint{0.946746in}{1.163591in}}%
\pgfpathlineto{\pgfqpoint{0.947077in}{1.164406in}}%
\pgfpathlineto{\pgfqpoint{0.948205in}{1.165249in}}%
\pgfpathlineto{\pgfqpoint{0.948366in}{1.165124in}}%
\pgfpathlineto{\pgfqpoint{0.951194in}{1.161166in}}%
\pgfpathlineto{\pgfqpoint{0.951689in}{1.160504in}}%
\pgfpathlineto{\pgfqpoint{0.952023in}{1.161328in}}%
\pgfpathlineto{\pgfqpoint{0.953177in}{1.162165in}}%
\pgfpathlineto{\pgfqpoint{0.953357in}{1.162015in}}%
\pgfpathlineto{\pgfqpoint{0.956626in}{1.157442in}}%
\pgfpathlineto{\pgfqpoint{0.957246in}{1.158754in}}%
\pgfpathlineto{\pgfqpoint{0.958203in}{1.159053in}}%
\pgfpathlineto{\pgfqpoint{0.958342in}{1.158921in}}%
\pgfpathlineto{\pgfqpoint{0.961588in}{1.154425in}}%
\pgfpathlineto{\pgfqpoint{0.962234in}{1.155779in}}%
\pgfpathlineto{\pgfqpoint{0.963289in}{1.155903in}}%
\pgfpathlineto{\pgfqpoint{0.963380in}{1.155793in}}%
\pgfpathlineto{\pgfqpoint{0.966526in}{1.151426in}}%
\pgfpathlineto{\pgfqpoint{0.966952in}{1.152430in}}%
\pgfpathlineto{\pgfqpoint{0.967943in}{1.153148in}}%
\pgfpathlineto{\pgfqpoint{0.968057in}{1.153084in}}%
\pgfpathlineto{\pgfqpoint{0.969215in}{1.151580in}}%
\pgfpathlineto{\pgfqpoint{0.971476in}{1.148469in}}%
\pgfpathlineto{\pgfqpoint{0.971809in}{1.149287in}}%
\pgfpathlineto{\pgfqpoint{0.972847in}{1.150200in}}%
\pgfpathlineto{\pgfqpoint{0.972928in}{1.150165in}}%
\pgfpathlineto{\pgfqpoint{0.973833in}{1.149081in}}%
\pgfpathlineto{\pgfqpoint{0.976409in}{1.145526in}}%
\pgfpathlineto{\pgfqpoint{0.976748in}{1.146349in}}%
\pgfpathlineto{\pgfqpoint{0.977912in}{1.147236in}}%
\pgfpathlineto{\pgfqpoint{0.978053in}{1.147132in}}%
\pgfpathlineto{\pgfqpoint{0.980010in}{1.144485in}}%
\pgfpathlineto{\pgfqpoint{0.981367in}{1.142644in}}%
\pgfpathlineto{\pgfqpoint{0.981701in}{1.143464in}}%
\pgfpathlineto{\pgfqpoint{0.982781in}{1.144374in}}%
\pgfpathlineto{\pgfqpoint{0.982848in}{1.144341in}}%
\pgfpathlineto{\pgfqpoint{0.984499in}{1.142223in}}%
\pgfpathlineto{\pgfqpoint{0.986312in}{1.139764in}}%
\pgfpathlineto{\pgfqpoint{0.986643in}{1.140581in}}%
\pgfpathlineto{\pgfqpoint{0.987814in}{1.141475in}}%
\pgfpathlineto{\pgfqpoint{0.987966in}{1.141362in}}%
\pgfpathlineto{\pgfqpoint{0.989915in}{1.138740in}}%
\pgfpathlineto{\pgfqpoint{0.991252in}{1.136950in}}%
\pgfpathlineto{\pgfqpoint{0.991581in}{1.137767in}}%
\pgfpathlineto{\pgfqpoint{0.992918in}{1.138542in}}%
\pgfpathlineto{\pgfqpoint{0.993130in}{1.138295in}}%
\pgfpathlineto{\pgfqpoint{0.996204in}{1.134140in}}%
\pgfpathlineto{\pgfqpoint{0.996529in}{1.134950in}}%
\pgfpathlineto{\pgfqpoint{0.997852in}{1.135763in}}%
\pgfpathlineto{\pgfqpoint{0.998054in}{1.135541in}}%
\pgfpathlineto{\pgfqpoint{1.001143in}{1.131387in}}%
\pgfpathlineto{\pgfqpoint{1.001508in}{1.132273in}}%
\pgfpathlineto{\pgfqpoint{1.002635in}{1.133126in}}%
\pgfpathlineto{\pgfqpoint{1.002766in}{1.133039in}}%
\pgfpathlineto{\pgfqpoint{1.004389in}{1.130912in}}%
\pgfpathlineto{\pgfqpoint{1.006087in}{1.128643in}}%
\pgfpathlineto{\pgfqpoint{1.006415in}{1.129466in}}%
\pgfpathlineto{\pgfqpoint{1.007743in}{1.130301in}}%
\pgfpathlineto{\pgfqpoint{1.007950in}{1.130078in}}%
\pgfpathlineto{\pgfqpoint{1.011034in}{1.125965in}}%
\pgfpathlineto{\pgfqpoint{1.011357in}{1.126775in}}%
\pgfpathlineto{\pgfqpoint{1.012554in}{1.127708in}}%
\pgfpathlineto{\pgfqpoint{1.012706in}{1.127599in}}%
\pgfpathlineto{\pgfqpoint{1.015855in}{1.123427in}}%
\pgfpathlineto{\pgfqpoint{1.015978in}{1.123296in}}%
\pgfpathlineto{\pgfqpoint{1.016301in}{1.124109in}}%
\pgfpathlineto{\pgfqpoint{1.017555in}{1.125027in}}%
\pgfpathlineto{\pgfqpoint{1.017703in}{1.124906in}}%
\pgfpathlineto{\pgfqpoint{1.020928in}{1.120663in}}%
\pgfpathlineto{\pgfqpoint{1.021558in}{1.122008in}}%
\pgfpathlineto{\pgfqpoint{1.022420in}{1.122427in}}%
\pgfpathlineto{\pgfqpoint{1.022568in}{1.122334in}}%
\pgfpathlineto{\pgfqpoint{1.024262in}{1.120146in}}%
\pgfpathlineto{\pgfqpoint{1.025874in}{1.118038in}}%
\pgfpathlineto{\pgfqpoint{1.026200in}{1.118851in}}%
\pgfpathlineto{\pgfqpoint{1.027581in}{1.119670in}}%
\pgfpathlineto{\pgfqpoint{1.027775in}{1.119450in}}%
\pgfpathlineto{\pgfqpoint{1.030820in}{1.115450in}}%
\pgfpathlineto{\pgfqpoint{1.031142in}{1.116256in}}%
\pgfpathlineto{\pgfqpoint{1.032591in}{1.117022in}}%
\pgfpathlineto{\pgfqpoint{1.032773in}{1.116797in}}%
\pgfpathlineto{\pgfqpoint{1.035752in}{1.112875in}}%
\pgfpathlineto{\pgfqpoint{1.036087in}{1.113690in}}%
\pgfpathlineto{\pgfqpoint{1.037429in}{1.114571in}}%
\pgfpathlineto{\pgfqpoint{1.037631in}{1.114364in}}%
\pgfpathlineto{\pgfqpoint{1.040705in}{1.110351in}}%
\pgfpathlineto{\pgfqpoint{1.041093in}{1.111299in}}%
\pgfpathlineto{\pgfqpoint{1.042494in}{1.111963in}}%
\pgfpathlineto{\pgfqpoint{1.042665in}{1.111756in}}%
\pgfpathlineto{\pgfqpoint{1.045653in}{1.107869in}}%
\pgfpathlineto{\pgfqpoint{1.045975in}{1.108679in}}%
\pgfpathlineto{\pgfqpoint{1.047196in}{1.109664in}}%
\pgfpathlineto{\pgfqpoint{1.047297in}{1.109601in}}%
\pgfpathlineto{\pgfqpoint{1.049225in}{1.107164in}}%
\pgfpathlineto{\pgfqpoint{1.050589in}{1.105396in}}%
\pgfpathlineto{\pgfqpoint{1.050924in}{1.106215in}}%
\pgfpathlineto{\pgfqpoint{1.052241in}{1.107152in}}%
\pgfpathlineto{\pgfqpoint{1.052422in}{1.106989in}}%
\pgfpathlineto{\pgfqpoint{1.055549in}{1.102968in}}%
\pgfpathlineto{\pgfqpoint{1.056020in}{1.104066in}}%
\pgfpathlineto{\pgfqpoint{1.057406in}{1.104506in}}%
\pgfpathlineto{\pgfqpoint{1.057496in}{1.104395in}}%
\pgfpathlineto{\pgfqpoint{1.060494in}{1.100545in}}%
\pgfpathlineto{\pgfqpoint{1.060812in}{1.101348in}}%
\pgfpathlineto{\pgfqpoint{1.062116in}{1.102314in}}%
\pgfpathlineto{\pgfqpoint{1.062265in}{1.102193in}}%
\pgfpathlineto{\pgfqpoint{1.064880in}{1.098846in}}%
\pgfpathlineto{\pgfqpoint{1.065439in}{1.098158in}}%
\pgfpathlineto{\pgfqpoint{1.065761in}{1.098966in}}%
\pgfpathlineto{\pgfqpoint{1.067227in}{1.099801in}}%
\pgfpathlineto{\pgfqpoint{1.067419in}{1.099576in}}%
\pgfpathlineto{\pgfqpoint{1.070378in}{1.095809in}}%
\pgfpathlineto{\pgfqpoint{1.070696in}{1.096612in}}%
\pgfpathlineto{\pgfqpoint{1.071870in}{1.097676in}}%
\pgfpathlineto{\pgfqpoint{1.071900in}{1.097665in}}%
\pgfpathlineto{\pgfqpoint{1.073158in}{1.096235in}}%
\pgfpathlineto{\pgfqpoint{1.075317in}{1.093486in}}%
\pgfpathlineto{\pgfqpoint{1.075646in}{1.094300in}}%
\pgfpathlineto{\pgfqpoint{1.077005in}{1.095268in}}%
\pgfpathlineto{\pgfqpoint{1.077159in}{1.095132in}}%
\pgfpathlineto{\pgfqpoint{1.080263in}{1.091198in}}%
\pgfpathlineto{\pgfqpoint{1.080752in}{1.092324in}}%
\pgfpathlineto{\pgfqpoint{1.081893in}{1.093034in}}%
\pgfpathlineto{\pgfqpoint{1.082006in}{1.092955in}}%
\pgfpathlineto{\pgfqpoint{1.084264in}{1.090133in}}%
\pgfpathlineto{\pgfqpoint{1.085215in}{1.088961in}}%
\pgfpathlineto{\pgfqpoint{1.085535in}{1.089765in}}%
\pgfpathlineto{\pgfqpoint{1.086764in}{1.090811in}}%
\pgfpathlineto{\pgfqpoint{1.086791in}{1.090798in}}%
\pgfpathlineto{\pgfqpoint{1.087980in}{1.089433in}}%
\pgfpathlineto{\pgfqpoint{1.090152in}{1.086690in}}%
\pgfpathlineto{\pgfqpoint{1.090481in}{1.087498in}}%
\pgfpathlineto{\pgfqpoint{1.091733in}{1.088572in}}%
\pgfpathlineto{\pgfqpoint{1.091747in}{1.088563in}}%
\pgfpathlineto{\pgfqpoint{1.092921in}{1.087230in}}%
\pgfpathlineto{\pgfqpoint{1.095106in}{1.084509in}}%
\pgfpathlineto{\pgfqpoint{1.095423in}{1.085310in}}%
\pgfpathlineto{\pgfqpoint{1.096622in}{1.086403in}}%
\pgfpathlineto{\pgfqpoint{1.097567in}{1.085417in}}%
\pgfpathlineto{\pgfqpoint{1.100051in}{1.082331in}}%
\pgfpathlineto{\pgfqpoint{1.100364in}{1.083120in}}%
\pgfpathlineto{\pgfqpoint{1.101518in}{1.084257in}}%
\pgfpathlineto{\pgfqpoint{1.101522in}{1.084256in}}%
\pgfpathlineto{\pgfqpoint{1.102720in}{1.083009in}}%
\pgfpathlineto{\pgfqpoint{1.104989in}{1.080179in}}%
\pgfpathlineto{\pgfqpoint{1.105314in}{1.080989in}}%
\pgfpathlineto{\pgfqpoint{1.106598in}{1.082073in}}%
\pgfpathlineto{\pgfqpoint{1.106708in}{1.082008in}}%
\pgfpathlineto{\pgfqpoint{1.108965in}{1.079259in}}%
\pgfpathlineto{\pgfqpoint{1.109935in}{1.078062in}}%
\pgfpathlineto{\pgfqpoint{1.110260in}{1.078875in}}%
\pgfpathlineto{\pgfqpoint{1.111589in}{1.079947in}}%
\pgfpathlineto{\pgfqpoint{1.111677in}{1.079888in}}%
\pgfpathlineto{\pgfqpoint{1.113437in}{1.077760in}}%
\pgfpathlineto{\pgfqpoint{1.114893in}{1.075983in}}%
\pgfpathlineto{\pgfqpoint{1.115207in}{1.076784in}}%
\pgfpathlineto{\pgfqpoint{1.116452in}{1.077904in}}%
\pgfpathlineto{\pgfqpoint{1.117726in}{1.076504in}}%
\pgfpathlineto{\pgfqpoint{1.119825in}{1.073914in}}%
\pgfpathlineto{\pgfqpoint{1.120146in}{1.074717in}}%
\pgfpathlineto{\pgfqpoint{1.121399in}{1.075859in}}%
\pgfpathlineto{\pgfqpoint{1.121409in}{1.075854in}}%
\pgfpathlineto{\pgfqpoint{1.122753in}{1.074366in}}%
\pgfpathlineto{\pgfqpoint{1.124780in}{1.071896in}}%
\pgfpathlineto{\pgfqpoint{1.125094in}{1.072691in}}%
\pgfpathlineto{\pgfqpoint{1.126265in}{1.073852in}}%
\pgfpathlineto{\pgfqpoint{1.127114in}{1.073065in}}%
\pgfpathlineto{\pgfqpoint{1.129714in}{1.069868in}}%
\pgfpathlineto{\pgfqpoint{1.130039in}{1.070668in}}%
\pgfpathlineto{\pgfqpoint{1.131213in}{1.071845in}}%
\pgfpathlineto{\pgfqpoint{1.132447in}{1.070592in}}%
\pgfpathlineto{\pgfqpoint{1.134670in}{1.067900in}}%
\pgfpathlineto{\pgfqpoint{1.134983in}{1.068696in}}%
\pgfpathlineto{\pgfqpoint{1.136157in}{1.069890in}}%
\pgfpathlineto{\pgfqpoint{1.137009in}{1.069118in}}%
\pgfpathlineto{\pgfqpoint{1.139615in}{1.065960in}}%
\pgfpathlineto{\pgfqpoint{1.139929in}{1.066758in}}%
\pgfpathlineto{\pgfqpoint{1.141105in}{1.067959in}}%
\pgfpathlineto{\pgfqpoint{1.141958in}{1.067190in}}%
\pgfpathlineto{\pgfqpoint{1.144559in}{1.064042in}}%
\pgfpathlineto{\pgfqpoint{1.144872in}{1.064834in}}%
\pgfpathlineto{\pgfqpoint{1.146074in}{1.066050in}}%
\pgfpathlineto{\pgfqpoint{1.146085in}{1.066047in}}%
\pgfpathlineto{\pgfqpoint{1.147016in}{1.065157in}}%
\pgfpathlineto{\pgfqpoint{1.149499in}{1.062150in}}%
\pgfpathlineto{\pgfqpoint{1.149819in}{1.062954in}}%
\pgfpathlineto{\pgfqpoint{1.151047in}{1.064178in}}%
\pgfpathlineto{\pgfqpoint{1.152046in}{1.063202in}}%
\pgfpathlineto{\pgfqpoint{1.154455in}{1.060316in}}%
\pgfpathlineto{\pgfqpoint{1.154767in}{1.061109in}}%
\pgfpathlineto{\pgfqpoint{1.155941in}{1.062328in}}%
\pgfpathlineto{\pgfqpoint{1.156893in}{1.061463in}}%
\pgfpathlineto{\pgfqpoint{1.159402in}{1.058467in}}%
\pgfpathlineto{\pgfqpoint{1.159712in}{1.059262in}}%
\pgfpathlineto{\pgfqpoint{1.160879in}{1.060498in}}%
\pgfpathlineto{\pgfqpoint{1.160882in}{1.060497in}}%
\pgfpathlineto{\pgfqpoint{1.161946in}{1.059513in}}%
\pgfpathlineto{\pgfqpoint{1.164334in}{1.056648in}}%
\pgfpathlineto{\pgfqpoint{1.164651in}{1.057450in}}%
\pgfpathlineto{\pgfqpoint{1.165881in}{1.058713in}}%
\pgfpathlineto{\pgfqpoint{1.167010in}{1.057609in}}%
\pgfpathlineto{\pgfqpoint{1.169276in}{1.054889in}}%
\pgfpathlineto{\pgfqpoint{1.169599in}{1.055694in}}%
\pgfpathlineto{\pgfqpoint{1.170792in}{1.056992in}}%
\pgfpathlineto{\pgfqpoint{1.171742in}{1.056155in}}%
\pgfpathlineto{\pgfqpoint{1.174232in}{1.053197in}}%
\pgfpathlineto{\pgfqpoint{1.174540in}{1.053994in}}%
\pgfpathlineto{\pgfqpoint{1.175931in}{1.055249in}}%
\pgfpathlineto{\pgfqpoint{1.175970in}{1.055228in}}%
\pgfpathlineto{\pgfqpoint{1.177707in}{1.053254in}}%
\pgfpathlineto{\pgfqpoint{1.179178in}{1.051534in}}%
\pgfpathlineto{\pgfqpoint{1.179485in}{1.052330in}}%
\pgfpathlineto{\pgfqpoint{1.180851in}{1.053606in}}%
\pgfpathlineto{\pgfqpoint{1.180861in}{1.053601in}}%
\pgfpathlineto{\pgfqpoint{1.182243in}{1.052097in}}%
\pgfpathlineto{\pgfqpoint{1.184121in}{1.049892in}}%
\pgfpathlineto{\pgfqpoint{1.184427in}{1.050690in}}%
\pgfpathlineto{\pgfqpoint{1.185830in}{1.051971in}}%
\pgfpathlineto{\pgfqpoint{1.185884in}{1.051943in}}%
\pgfpathlineto{\pgfqpoint{1.187275in}{1.050384in}}%
\pgfpathlineto{\pgfqpoint{1.189070in}{1.048282in}}%
\pgfpathlineto{\pgfqpoint{1.189377in}{1.049081in}}%
\pgfpathlineto{\pgfqpoint{1.190788in}{1.050370in}}%
\pgfpathlineto{\pgfqpoint{1.190829in}{1.050348in}}%
\pgfpathlineto{\pgfqpoint{1.192515in}{1.048450in}}%
\pgfpathlineto{\pgfqpoint{1.194017in}{1.046703in}}%
\pgfpathlineto{\pgfqpoint{1.194322in}{1.047504in}}%
\pgfpathlineto{\pgfqpoint{1.195740in}{1.048816in}}%
\pgfpathlineto{\pgfqpoint{1.195777in}{1.048797in}}%
\pgfpathlineto{\pgfqpoint{1.197592in}{1.046756in}}%
\pgfpathlineto{\pgfqpoint{1.198960in}{1.045175in}}%
\pgfpathlineto{\pgfqpoint{1.199262in}{1.045978in}}%
\pgfpathlineto{\pgfqpoint{1.200870in}{1.047207in}}%
\pgfpathlineto{\pgfqpoint{1.201067in}{1.047010in}}%
\pgfpathlineto{\pgfqpoint{1.203896in}{1.043683in}}%
\pgfpathlineto{\pgfqpoint{1.204209in}{1.044491in}}%
\pgfpathlineto{\pgfqpoint{1.205692in}{1.045831in}}%
\pgfpathlineto{\pgfqpoint{1.205751in}{1.045795in}}%
\pgfpathlineto{\pgfqpoint{1.207392in}{1.043933in}}%
\pgfpathlineto{\pgfqpoint{1.208852in}{1.042255in}}%
\pgfpathlineto{\pgfqpoint{1.209149in}{1.043058in}}%
\pgfpathlineto{\pgfqpoint{1.210763in}{1.044349in}}%
\pgfpathlineto{\pgfqpoint{1.210946in}{1.044177in}}%
\pgfpathlineto{\pgfqpoint{1.213790in}{1.040857in}}%
\pgfpathlineto{\pgfqpoint{1.214112in}{1.041712in}}%
\pgfpathlineto{\pgfqpoint{1.215414in}{1.043150in}}%
\pgfpathlineto{\pgfqpoint{1.216513in}{1.042129in}}%
\pgfpathlineto{\pgfqpoint{1.218743in}{1.039549in}}%
\pgfpathlineto{\pgfqpoint{1.219034in}{1.040356in}}%
\pgfpathlineto{\pgfqpoint{1.220382in}{1.041900in}}%
\pgfpathlineto{\pgfqpoint{1.221593in}{1.040748in}}%
\pgfpathlineto{\pgfqpoint{1.223680in}{1.038326in}}%
\pgfpathlineto{\pgfqpoint{1.223975in}{1.039140in}}%
\pgfpathlineto{\pgfqpoint{1.225447in}{1.040712in}}%
\pgfpathlineto{\pgfqpoint{1.226821in}{1.039283in}}%
\pgfpathlineto{\pgfqpoint{1.228624in}{1.037190in}}%
\pgfpathlineto{\pgfqpoint{1.228923in}{1.038012in}}%
\pgfpathlineto{\pgfqpoint{1.230291in}{1.039659in}}%
\pgfpathlineto{\pgfqpoint{1.231388in}{1.038668in}}%
\pgfpathlineto{\pgfqpoint{1.233570in}{1.036138in}}%
\pgfpathlineto{\pgfqpoint{1.233863in}{1.036956in}}%
\pgfpathlineto{\pgfqpoint{1.235149in}{1.038676in}}%
\pgfpathlineto{\pgfqpoint{1.236014in}{1.038084in}}%
\pgfpathlineto{\pgfqpoint{1.238528in}{1.035194in}}%
\pgfpathlineto{\pgfqpoint{1.238809in}{1.036012in}}%
\pgfpathlineto{\pgfqpoint{1.240075in}{1.037798in}}%
\pgfpathlineto{\pgfqpoint{1.241118in}{1.037060in}}%
\pgfpathlineto{\pgfqpoint{1.243455in}{1.034352in}}%
\pgfpathlineto{\pgfqpoint{1.243748in}{1.035186in}}%
\pgfpathlineto{\pgfqpoint{1.245112in}{1.037059in}}%
\pgfpathlineto{\pgfqpoint{1.245948in}{1.036476in}}%
\pgfpathlineto{\pgfqpoint{1.248397in}{1.033641in}}%
\pgfpathlineto{\pgfqpoint{1.248685in}{1.034477in}}%
\pgfpathlineto{\pgfqpoint{1.250470in}{1.036291in}}%
\pgfpathlineto{\pgfqpoint{1.250698in}{1.036083in}}%
\pgfpathlineto{\pgfqpoint{1.253360in}{1.033030in}}%
\pgfpathlineto{\pgfqpoint{1.253656in}{1.033930in}}%
\pgfpathlineto{\pgfqpoint{1.255305in}{1.035789in}}%
\pgfpathlineto{\pgfqpoint{1.255365in}{1.035758in}}%
\pgfpathlineto{\pgfqpoint{1.257020in}{1.033935in}}%
\pgfpathlineto{\pgfqpoint{1.258298in}{1.032476in}}%
\pgfpathlineto{\pgfqpoint{1.258576in}{1.033311in}}%
\pgfpathlineto{\pgfqpoint{1.260280in}{1.035261in}}%
\pgfpathlineto{\pgfqpoint{1.260329in}{1.035233in}}%
\pgfpathlineto{\pgfqpoint{1.262141in}{1.033228in}}%
\pgfpathlineto{\pgfqpoint{1.263244in}{1.031975in}}%
\pgfpathlineto{\pgfqpoint{1.263519in}{1.032803in}}%
\pgfpathlineto{\pgfqpoint{1.265208in}{1.034765in}}%
\pgfpathlineto{\pgfqpoint{1.265254in}{1.034741in}}%
\pgfpathlineto{\pgfqpoint{1.266914in}{1.032924in}}%
\pgfpathlineto{\pgfqpoint{1.268178in}{1.031468in}}%
\pgfpathlineto{\pgfqpoint{1.268464in}{1.032311in}}%
\pgfpathlineto{\pgfqpoint{1.270226in}{1.034241in}}%
\pgfpathlineto{\pgfqpoint{1.270403in}{1.034107in}}%
\pgfpathlineto{\pgfqpoint{1.273134in}{1.030984in}}%
\pgfpathlineto{\pgfqpoint{1.273542in}{1.032164in}}%
\pgfpathlineto{\pgfqpoint{1.274871in}{1.033860in}}%
\pgfpathlineto{\pgfqpoint{1.275767in}{1.033180in}}%
\pgfpathlineto{\pgfqpoint{1.278079in}{1.030530in}}%
\pgfpathlineto{\pgfqpoint{1.278359in}{1.031368in}}%
\pgfpathlineto{\pgfqpoint{1.280023in}{1.033362in}}%
\pgfpathlineto{\pgfqpoint{1.281217in}{1.032141in}}%
\pgfpathlineto{\pgfqpoint{1.283024in}{1.030077in}}%
\pgfpathlineto{\pgfqpoint{1.283301in}{1.030906in}}%
\pgfpathlineto{\pgfqpoint{1.285027in}{1.032838in}}%
\pgfpathlineto{\pgfqpoint{1.285077in}{1.032811in}}%
\pgfpathlineto{\pgfqpoint{1.287197in}{1.030446in}}%
\pgfpathlineto{\pgfqpoint{1.287964in}{1.029566in}}%
\pgfpathlineto{\pgfqpoint{1.288249in}{1.030410in}}%
\pgfpathlineto{\pgfqpoint{1.290078in}{1.032315in}}%
\pgfpathlineto{\pgfqpoint{1.290324in}{1.032086in}}%
\pgfpathlineto{\pgfqpoint{1.292915in}{1.029119in}}%
\pgfpathlineto{\pgfqpoint{1.293195in}{1.029957in}}%
\pgfpathlineto{\pgfqpoint{1.294926in}{1.031948in}}%
\pgfpathlineto{\pgfqpoint{1.295028in}{1.031887in}}%
\pgfpathlineto{\pgfqpoint{1.297418in}{1.029198in}}%
\pgfpathlineto{\pgfqpoint{1.297862in}{1.028702in}}%
\pgfpathlineto{\pgfqpoint{1.298137in}{1.029536in}}%
\pgfpathlineto{\pgfqpoint{1.299875in}{1.031540in}}%
\pgfpathlineto{\pgfqpoint{1.299925in}{1.031511in}}%
\pgfpathlineto{\pgfqpoint{1.301770in}{1.029477in}}%
\pgfpathlineto{\pgfqpoint{1.302807in}{1.028301in}}%
\pgfpathlineto{\pgfqpoint{1.303087in}{1.029139in}}%
\pgfpathlineto{\pgfqpoint{1.304678in}{1.031155in}}%
\pgfpathlineto{\pgfqpoint{1.305710in}{1.030185in}}%
\pgfpathlineto{\pgfqpoint{1.307752in}{1.027859in}}%
\pgfpathlineto{\pgfqpoint{1.308026in}{1.028687in}}%
\pgfpathlineto{\pgfqpoint{1.309768in}{1.030681in}}%
\pgfpathlineto{\pgfqpoint{1.309876in}{1.030615in}}%
\pgfpathlineto{\pgfqpoint{1.312709in}{1.027451in}}%
\pgfpathlineto{\pgfqpoint{1.313248in}{1.028953in}}%
\pgfpathlineto{\pgfqpoint{1.314726in}{1.030220in}}%
\pgfpathlineto{\pgfqpoint{1.314806in}{1.030171in}}%
\pgfpathlineto{\pgfqpoint{1.317529in}{1.027107in}}%
\pgfpathlineto{\pgfqpoint{1.317641in}{1.026986in}}%
\pgfpathlineto{\pgfqpoint{1.317921in}{1.027825in}}%
\pgfpathlineto{\pgfqpoint{1.319702in}{1.029797in}}%
\pgfpathlineto{\pgfqpoint{1.319857in}{1.029684in}}%
\pgfpathlineto{\pgfqpoint{1.322597in}{1.026608in}}%
\pgfpathlineto{\pgfqpoint{1.323032in}{1.027879in}}%
\pgfpathlineto{\pgfqpoint{1.324327in}{1.029505in}}%
\pgfpathlineto{\pgfqpoint{1.325204in}{1.028864in}}%
\pgfpathlineto{\pgfqpoint{1.327534in}{1.026216in}}%
\pgfpathlineto{\pgfqpoint{1.327808in}{1.027050in}}%
\pgfpathlineto{\pgfqpoint{1.329577in}{1.029070in}}%
\pgfpathlineto{\pgfqpoint{1.329614in}{1.029047in}}%
\pgfpathlineto{\pgfqpoint{1.331528in}{1.026938in}}%
\pgfpathlineto{\pgfqpoint{1.332476in}{1.025859in}}%
\pgfpathlineto{\pgfqpoint{1.332759in}{1.026701in}}%
\pgfpathlineto{\pgfqpoint{1.334576in}{1.028644in}}%
\pgfpathlineto{\pgfqpoint{1.334824in}{1.028426in}}%
\pgfpathlineto{\pgfqpoint{1.337425in}{1.025471in}}%
\pgfpathlineto{\pgfqpoint{1.337698in}{1.026304in}}%
\pgfpathlineto{\pgfqpoint{1.339465in}{1.028333in}}%
\pgfpathlineto{\pgfqpoint{1.339509in}{1.028309in}}%
\pgfpathlineto{\pgfqpoint{1.341568in}{1.026036in}}%
\pgfpathlineto{\pgfqpoint{1.342367in}{1.025130in}}%
\pgfpathlineto{\pgfqpoint{1.342649in}{1.025970in}}%
\pgfpathlineto{\pgfqpoint{1.344437in}{1.027940in}}%
\pgfpathlineto{\pgfqpoint{1.344619in}{1.027801in}}%
\pgfpathlineto{\pgfqpoint{1.347325in}{1.024766in}}%
\pgfpathlineto{\pgfqpoint{1.347705in}{1.025903in}}%
\pgfpathlineto{\pgfqpoint{1.349230in}{1.027642in}}%
\pgfpathlineto{\pgfqpoint{1.350289in}{1.026624in}}%
\pgfpathlineto{\pgfqpoint{1.352253in}{1.024388in}}%
\pgfpathlineto{\pgfqpoint{1.352532in}{1.025224in}}%
\pgfpathlineto{\pgfqpoint{1.354085in}{1.027319in}}%
\pgfpathlineto{\pgfqpoint{1.355014in}{1.026539in}}%
\pgfpathlineto{\pgfqpoint{1.357197in}{1.024056in}}%
\pgfpathlineto{\pgfqpoint{1.357477in}{1.024893in}}%
\pgfpathlineto{\pgfqpoint{1.358939in}{1.026998in}}%
\pgfpathlineto{\pgfqpoint{1.359826in}{1.026372in}}%
\pgfpathlineto{\pgfqpoint{1.362143in}{1.023737in}}%
\pgfpathlineto{\pgfqpoint{1.362422in}{1.024571in}}%
\pgfpathlineto{\pgfqpoint{1.363863in}{1.026680in}}%
\pgfpathlineto{\pgfqpoint{1.364725in}{1.026109in}}%
\pgfpathlineto{\pgfqpoint{1.367092in}{1.023419in}}%
\pgfpathlineto{\pgfqpoint{1.367371in}{1.024255in}}%
\pgfpathlineto{\pgfqpoint{1.369089in}{1.026309in}}%
\pgfpathlineto{\pgfqpoint{1.369107in}{1.026301in}}%
\pgfpathlineto{\pgfqpoint{1.370356in}{1.025004in}}%
\pgfpathlineto{\pgfqpoint{1.372037in}{1.023095in}}%
\pgfpathlineto{\pgfqpoint{1.372317in}{1.023936in}}%
\pgfpathlineto{\pgfqpoint{1.374088in}{1.025966in}}%
\pgfpathlineto{\pgfqpoint{1.374212in}{1.025887in}}%
\pgfpathlineto{\pgfqpoint{1.376997in}{1.022793in}}%
\pgfpathlineto{\pgfqpoint{1.377494in}{1.024220in}}%
\pgfpathlineto{\pgfqpoint{1.378805in}{1.025727in}}%
\pgfpathlineto{\pgfqpoint{1.379771in}{1.024923in}}%
\pgfpathlineto{\pgfqpoint{1.381934in}{1.022482in}}%
\pgfpathlineto{\pgfqpoint{1.382206in}{1.023314in}}%
\pgfpathlineto{\pgfqpoint{1.383959in}{1.025373in}}%
\pgfpathlineto{\pgfqpoint{1.384073in}{1.025306in}}%
\pgfpathlineto{\pgfqpoint{1.386450in}{1.022664in}}%
\pgfpathlineto{\pgfqpoint{1.386888in}{1.022197in}}%
\pgfpathlineto{\pgfqpoint{1.387153in}{1.023021in}}%
\pgfpathlineto{\pgfqpoint{1.388779in}{1.025095in}}%
\pgfpathlineto{\pgfqpoint{1.389816in}{1.024126in}}%
\pgfpathlineto{\pgfqpoint{1.391836in}{1.021873in}}%
\pgfpathlineto{\pgfqpoint{1.392097in}{1.022690in}}%
\pgfpathlineto{\pgfqpoint{1.393684in}{1.024798in}}%
\pgfpathlineto{\pgfqpoint{1.394817in}{1.023765in}}%
\pgfpathlineto{\pgfqpoint{1.396770in}{1.021566in}}%
\pgfpathlineto{\pgfqpoint{1.397042in}{1.022398in}}%
\pgfpathlineto{\pgfqpoint{1.398839in}{1.024444in}}%
\pgfpathlineto{\pgfqpoint{1.398901in}{1.024406in}}%
\pgfpathlineto{\pgfqpoint{1.401151in}{1.021917in}}%
\pgfpathlineto{\pgfqpoint{1.401717in}{1.021295in}}%
\pgfpathlineto{\pgfqpoint{1.401987in}{1.022122in}}%
\pgfpathlineto{\pgfqpoint{1.403738in}{1.024182in}}%
\pgfpathlineto{\pgfqpoint{1.405333in}{1.022495in}}%
\pgfpathlineto{\pgfqpoint{1.406653in}{1.021003in}}%
\pgfpathlineto{\pgfqpoint{1.406933in}{1.021838in}}%
\pgfpathlineto{\pgfqpoint{1.408553in}{1.023915in}}%
\pgfpathlineto{\pgfqpoint{1.409600in}{1.022943in}}%
\pgfpathlineto{\pgfqpoint{1.411608in}{1.020693in}}%
\pgfpathlineto{\pgfqpoint{1.411874in}{1.021522in}}%
\pgfpathlineto{\pgfqpoint{1.413649in}{1.023604in}}%
\pgfpathlineto{\pgfqpoint{1.413750in}{1.023545in}}%
\pgfpathlineto{\pgfqpoint{1.415765in}{1.021319in}}%
\pgfpathlineto{\pgfqpoint{1.416559in}{1.020447in}}%
\pgfpathlineto{\pgfqpoint{1.416824in}{1.021276in}}%
\pgfpathlineto{\pgfqpoint{1.418345in}{1.023413in}}%
\pgfpathlineto{\pgfqpoint{1.419402in}{1.022545in}}%
\pgfpathlineto{\pgfqpoint{1.421505in}{1.020194in}}%
\pgfpathlineto{\pgfqpoint{1.421769in}{1.021023in}}%
\pgfpathlineto{\pgfqpoint{1.423291in}{1.023157in}}%
\pgfpathlineto{\pgfqpoint{1.424352in}{1.022283in}}%
\pgfpathlineto{\pgfqpoint{1.426450in}{1.019939in}}%
\pgfpathlineto{\pgfqpoint{1.426714in}{1.020768in}}%
\pgfpathlineto{\pgfqpoint{1.428236in}{1.022902in}}%
\pgfpathlineto{\pgfqpoint{1.429301in}{1.022025in}}%
\pgfpathlineto{\pgfqpoint{1.431395in}{1.019687in}}%
\pgfpathlineto{\pgfqpoint{1.431660in}{1.020516in}}%
\pgfpathlineto{\pgfqpoint{1.433180in}{1.022653in}}%
\pgfpathlineto{\pgfqpoint{1.434229in}{1.021799in}}%
\pgfpathlineto{\pgfqpoint{1.436330in}{1.019429in}}%
\pgfpathlineto{\pgfqpoint{1.436608in}{1.020267in}}%
\pgfpathlineto{\pgfqpoint{1.438423in}{1.022302in}}%
\pgfpathlineto{\pgfqpoint{1.438574in}{1.022193in}}%
\pgfpathlineto{\pgfqpoint{1.441286in}{1.019177in}}%
\pgfpathlineto{\pgfqpoint{1.441764in}{1.020567in}}%
\pgfpathlineto{\pgfqpoint{1.443418in}{1.022033in}}%
\pgfpathlineto{\pgfqpoint{1.443594in}{1.021887in}}%
\pgfpathlineto{\pgfqpoint{1.446234in}{1.018951in}}%
\pgfpathlineto{\pgfqpoint{1.446584in}{1.020013in}}%
\pgfpathlineto{\pgfqpoint{1.448185in}{1.021875in}}%
\pgfpathlineto{\pgfqpoint{1.449493in}{1.020568in}}%
\pgfpathlineto{\pgfqpoint{1.451177in}{1.018696in}}%
\pgfpathlineto{\pgfqpoint{1.451442in}{1.019525in}}%
\pgfpathlineto{\pgfqpoint{1.452963in}{1.021670in}}%
\pgfpathlineto{\pgfqpoint{1.454019in}{1.020811in}}%
\pgfpathlineto{\pgfqpoint{1.456116in}{1.018466in}}%
\pgfpathlineto{\pgfqpoint{1.456386in}{1.019296in}}%
\pgfpathlineto{\pgfqpoint{1.458200in}{1.021365in}}%
\pgfpathlineto{\pgfqpoint{1.458288in}{1.021309in}}%
\pgfpathlineto{\pgfqpoint{1.461064in}{1.018254in}}%
\pgfpathlineto{\pgfqpoint{1.461587in}{1.019749in}}%
\pgfpathlineto{\pgfqpoint{1.462856in}{1.021228in}}%
\pgfpathlineto{\pgfqpoint{1.463765in}{1.020532in}}%
\pgfpathlineto{\pgfqpoint{1.466013in}{1.018027in}}%
\pgfpathlineto{\pgfqpoint{1.466278in}{1.018854in}}%
\pgfpathlineto{\pgfqpoint{1.467826in}{1.020976in}}%
\pgfpathlineto{\pgfqpoint{1.468912in}{1.020050in}}%
\pgfpathlineto{\pgfqpoint{1.470959in}{1.017772in}}%
\pgfpathlineto{\pgfqpoint{1.471223in}{1.018602in}}%
\pgfpathlineto{\pgfqpoint{1.472745in}{1.020753in}}%
\pgfpathlineto{\pgfqpoint{1.473781in}{1.019922in}}%
\pgfpathlineto{\pgfqpoint{1.475899in}{1.017557in}}%
\pgfpathlineto{\pgfqpoint{1.476169in}{1.018388in}}%
\pgfpathlineto{\pgfqpoint{1.477915in}{1.020504in}}%
\pgfpathlineto{\pgfqpoint{1.479374in}{1.019001in}}%
\pgfpathlineto{\pgfqpoint{1.480843in}{1.017362in}}%
\pgfpathlineto{\pgfqpoint{1.481119in}{1.018194in}}%
\pgfpathlineto{\pgfqpoint{1.482882in}{1.020234in}}%
\pgfpathlineto{\pgfqpoint{1.482931in}{1.020207in}}%
\pgfpathlineto{\pgfqpoint{1.484853in}{1.018126in}}%
\pgfpathlineto{\pgfqpoint{1.485795in}{1.017093in}}%
\pgfpathlineto{\pgfqpoint{1.486060in}{1.017922in}}%
\pgfpathlineto{\pgfqpoint{1.487582in}{1.020072in}}%
\pgfpathlineto{\pgfqpoint{1.488637in}{1.019221in}}%
\pgfpathlineto{\pgfqpoint{1.490736in}{1.016884in}}%
\pgfpathlineto{\pgfqpoint{1.491004in}{1.017711in}}%
\pgfpathlineto{\pgfqpoint{1.492747in}{1.019839in}}%
\pgfpathlineto{\pgfqpoint{1.494225in}{1.018316in}}%
\pgfpathlineto{\pgfqpoint{1.495676in}{1.016693in}}%
\pgfpathlineto{\pgfqpoint{1.495952in}{1.017528in}}%
\pgfpathlineto{\pgfqpoint{1.497760in}{1.019586in}}%
\pgfpathlineto{\pgfqpoint{1.497872in}{1.019513in}}%
\pgfpathlineto{\pgfqpoint{1.500634in}{1.016486in}}%
\pgfpathlineto{\pgfqpoint{1.501113in}{1.017873in}}%
\pgfpathlineto{\pgfqpoint{1.502763in}{1.019339in}}%
\pgfpathlineto{\pgfqpoint{1.502881in}{1.019249in}}%
\pgfpathlineto{\pgfqpoint{1.505579in}{1.016268in}}%
\pgfpathlineto{\pgfqpoint{1.506017in}{1.017558in}}%
\pgfpathlineto{\pgfqpoint{1.507366in}{1.019249in}}%
\pgfpathlineto{\pgfqpoint{1.508487in}{1.018328in}}%
\pgfpathlineto{\pgfqpoint{1.510518in}{1.016072in}}%
\pgfpathlineto{\pgfqpoint{1.510786in}{1.016901in}}%
\pgfpathlineto{\pgfqpoint{1.512350in}{1.019061in}}%
\pgfpathlineto{\pgfqpoint{1.513357in}{1.018227in}}%
\pgfpathlineto{\pgfqpoint{1.515461in}{1.015882in}}%
\pgfpathlineto{\pgfqpoint{1.515730in}{1.016708in}}%
\pgfpathlineto{\pgfqpoint{1.517334in}{1.018858in}}%
\pgfpathlineto{\pgfqpoint{1.518400in}{1.017916in}}%
\pgfpathlineto{\pgfqpoint{1.520415in}{1.015688in}}%
\pgfpathlineto{\pgfqpoint{1.520678in}{1.016511in}}%
\pgfpathlineto{\pgfqpoint{1.522184in}{1.018665in}}%
\pgfpathlineto{\pgfqpoint{1.523069in}{1.018032in}}%
\pgfpathlineto{\pgfqpoint{1.525349in}{1.015482in}}%
\pgfpathlineto{\pgfqpoint{1.525626in}{1.016320in}}%
\pgfpathlineto{\pgfqpoint{1.527415in}{1.018403in}}%
\pgfpathlineto{\pgfqpoint{1.527541in}{1.018325in}}%
\pgfpathlineto{\pgfqpoint{1.530306in}{1.015299in}}%
\pgfpathlineto{\pgfqpoint{1.530816in}{1.016756in}}%
\pgfpathlineto{\pgfqpoint{1.532429in}{1.018156in}}%
\pgfpathlineto{\pgfqpoint{1.532550in}{1.018065in}}%
\pgfpathlineto{\pgfqpoint{1.535244in}{1.015082in}}%
\pgfpathlineto{\pgfqpoint{1.535677in}{1.016348in}}%
\pgfpathlineto{\pgfqpoint{1.537048in}{1.018088in}}%
\pgfpathlineto{\pgfqpoint{1.538141in}{1.017197in}}%
\pgfpathlineto{\pgfqpoint{1.540182in}{1.014914in}}%
\pgfpathlineto{\pgfqpoint{1.540461in}{1.015749in}}%
\pgfpathlineto{\pgfqpoint{1.542209in}{1.017825in}}%
\pgfpathlineto{\pgfqpoint{1.542225in}{1.017817in}}%
\pgfpathlineto{\pgfqpoint{1.543581in}{1.016414in}}%
\pgfpathlineto{\pgfqpoint{1.545127in}{1.014686in}}%
\pgfpathlineto{\pgfqpoint{1.545404in}{1.015523in}}%
\pgfpathlineto{\pgfqpoint{1.547002in}{1.017687in}}%
\pgfpathlineto{\pgfqpoint{1.548130in}{1.016688in}}%
\pgfpathlineto{\pgfqpoint{1.550073in}{1.014515in}}%
\pgfpathlineto{\pgfqpoint{1.550352in}{1.015355in}}%
\pgfpathlineto{\pgfqpoint{1.552131in}{1.017448in}}%
\pgfpathlineto{\pgfqpoint{1.552136in}{1.017445in}}%
\pgfpathlineto{\pgfqpoint{1.553924in}{1.015554in}}%
\pgfpathlineto{\pgfqpoint{1.555024in}{1.014338in}}%
\pgfpathlineto{\pgfqpoint{1.555295in}{1.015169in}}%
\pgfpathlineto{\pgfqpoint{1.556871in}{1.017339in}}%
\pgfpathlineto{\pgfqpoint{1.557924in}{1.016452in}}%
\pgfpathlineto{\pgfqpoint{1.559966in}{1.014171in}}%
\pgfpathlineto{\pgfqpoint{1.560240in}{1.014993in}}%
\pgfpathlineto{\pgfqpoint{1.561808in}{1.017119in}}%
\pgfpathlineto{\pgfqpoint{1.562724in}{1.016378in}}%
\pgfpathlineto{\pgfqpoint{1.564907in}{1.013939in}}%
\pgfpathlineto{\pgfqpoint{1.565183in}{1.014767in}}%
\pgfpathlineto{\pgfqpoint{1.566540in}{1.016924in}}%
\pgfpathlineto{\pgfqpoint{1.567415in}{1.016513in}}%
\pgfpathlineto{\pgfqpoint{1.569854in}{1.013788in}}%
\pgfpathlineto{\pgfqpoint{1.570132in}{1.014624in}}%
\pgfpathlineto{\pgfqpoint{1.571683in}{1.016790in}}%
\pgfpathlineto{\pgfqpoint{1.572605in}{1.016073in}}%
\pgfpathlineto{\pgfqpoint{1.574797in}{1.013626in}}%
\pgfpathlineto{\pgfqpoint{1.575075in}{1.014458in}}%
\pgfpathlineto{\pgfqpoint{1.576441in}{1.016612in}}%
\pgfpathlineto{\pgfqpoint{1.577443in}{1.016042in}}%
\pgfpathlineto{\pgfqpoint{1.579744in}{1.013474in}}%
\pgfpathlineto{\pgfqpoint{1.580021in}{1.014305in}}%
\pgfpathlineto{\pgfqpoint{1.581415in}{1.016466in}}%
\pgfpathlineto{\pgfqpoint{1.582435in}{1.015837in}}%
\pgfpathlineto{\pgfqpoint{1.584692in}{1.013319in}}%
\pgfpathlineto{\pgfqpoint{1.584970in}{1.014155in}}%
\pgfpathlineto{\pgfqpoint{1.586600in}{1.016312in}}%
\pgfpathlineto{\pgfqpoint{1.587797in}{1.015210in}}%
\pgfpathlineto{\pgfqpoint{1.589637in}{1.013158in}}%
\pgfpathlineto{\pgfqpoint{1.589915in}{1.013994in}}%
\pgfpathlineto{\pgfqpoint{1.591550in}{1.016150in}}%
\pgfpathlineto{\pgfqpoint{1.592652in}{1.015150in}}%
\pgfpathlineto{\pgfqpoint{1.594582in}{1.012998in}}%
\pgfpathlineto{\pgfqpoint{1.594862in}{1.013838in}}%
\pgfpathlineto{\pgfqpoint{1.596631in}{1.015951in}}%
\pgfpathlineto{\pgfqpoint{1.596651in}{1.015941in}}%
\pgfpathlineto{\pgfqpoint{1.598003in}{1.014541in}}%
\pgfpathlineto{\pgfqpoint{1.599527in}{1.012844in}}%
\pgfpathlineto{\pgfqpoint{1.599803in}{1.013675in}}%
\pgfpathlineto{\pgfqpoint{1.601272in}{1.015849in}}%
\pgfpathlineto{\pgfqpoint{1.602121in}{1.015317in}}%
\pgfpathlineto{\pgfqpoint{1.604469in}{1.012699in}}%
\pgfpathlineto{\pgfqpoint{1.604747in}{1.013528in}}%
\pgfpathlineto{\pgfqpoint{1.606096in}{1.015682in}}%
\pgfpathlineto{\pgfqpoint{1.606974in}{1.015281in}}%
\pgfpathlineto{\pgfqpoint{1.609425in}{1.012562in}}%
\pgfpathlineto{\pgfqpoint{1.609692in}{1.013387in}}%
\pgfpathlineto{\pgfqpoint{1.611415in}{1.015549in}}%
\pgfpathlineto{\pgfqpoint{1.612772in}{1.014199in}}%
\pgfpathlineto{\pgfqpoint{1.614361in}{1.012431in}}%
\pgfpathlineto{\pgfqpoint{1.614641in}{1.013268in}}%
\pgfpathlineto{\pgfqpoint{1.616211in}{1.015414in}}%
\pgfpathlineto{\pgfqpoint{1.617339in}{1.014444in}}%
\pgfpathlineto{\pgfqpoint{1.619312in}{1.012250in}}%
\pgfpathlineto{\pgfqpoint{1.619586in}{1.013080in}}%
\pgfpathlineto{\pgfqpoint{1.621187in}{1.015260in}}%
\pgfpathlineto{\pgfqpoint{1.622120in}{1.014491in}}%
\pgfpathlineto{\pgfqpoint{1.624262in}{1.012122in}}%
\pgfpathlineto{\pgfqpoint{1.624530in}{1.012947in}}%
\pgfpathlineto{\pgfqpoint{1.626144in}{1.015126in}}%
\pgfpathlineto{\pgfqpoint{1.627259in}{1.014139in}}%
\pgfpathlineto{\pgfqpoint{1.629201in}{1.011978in}}%
\pgfpathlineto{\pgfqpoint{1.629480in}{1.012817in}}%
\pgfpathlineto{\pgfqpoint{1.631320in}{1.014872in}}%
\pgfpathlineto{\pgfqpoint{1.631460in}{1.014772in}}%
\pgfpathlineto{\pgfqpoint{1.634150in}{1.011808in}}%
\pgfpathlineto{\pgfqpoint{1.634627in}{1.013181in}}%
\pgfpathlineto{\pgfqpoint{1.636078in}{1.014812in}}%
\pgfpathlineto{\pgfqpoint{1.637146in}{1.013841in}}%
\pgfpathlineto{\pgfqpoint{1.639088in}{1.011681in}}%
\pgfpathlineto{\pgfqpoint{1.639365in}{1.012509in}}%
\pgfpathlineto{\pgfqpoint{1.640707in}{1.014663in}}%
\pgfpathlineto{\pgfqpoint{1.641681in}{1.014170in}}%
\pgfpathlineto{\pgfqpoint{1.644039in}{1.011548in}}%
\pgfpathlineto{\pgfqpoint{1.644314in}{1.012381in}}%
\pgfpathlineto{\pgfqpoint{1.645978in}{1.014542in}}%
\pgfpathlineto{\pgfqpoint{1.647053in}{1.013553in}}%
\pgfpathlineto{\pgfqpoint{1.648984in}{1.011407in}}%
\pgfpathlineto{\pgfqpoint{1.649262in}{1.012247in}}%
\pgfpathlineto{\pgfqpoint{1.651108in}{1.014323in}}%
\pgfpathlineto{\pgfqpoint{1.651255in}{1.014218in}}%
\pgfpathlineto{\pgfqpoint{1.653930in}{1.011265in}}%
\pgfpathlineto{\pgfqpoint{1.654381in}{1.012567in}}%
\pgfpathlineto{\pgfqpoint{1.655893in}{1.014252in}}%
\pgfpathlineto{\pgfqpoint{1.657084in}{1.013115in}}%
\pgfpathlineto{\pgfqpoint{1.658878in}{1.011133in}}%
\pgfpathlineto{\pgfqpoint{1.659151in}{1.011967in}}%
\pgfpathlineto{\pgfqpoint{1.660855in}{1.014122in}}%
\pgfpathlineto{\pgfqpoint{1.662069in}{1.012948in}}%
\pgfpathlineto{\pgfqpoint{1.663818in}{1.011004in}}%
\pgfpathlineto{\pgfqpoint{1.664094in}{1.011833in}}%
\pgfpathlineto{\pgfqpoint{1.665561in}{1.014005in}}%
\pgfpathlineto{\pgfqpoint{1.666457in}{1.013428in}}%
\pgfpathlineto{\pgfqpoint{1.668766in}{1.010864in}}%
\pgfpathlineto{\pgfqpoint{1.669040in}{1.011692in}}%
\pgfpathlineto{\pgfqpoint{1.670604in}{1.013877in}}%
\pgfpathlineto{\pgfqpoint{1.671533in}{1.013154in}}%
\pgfpathlineto{\pgfqpoint{1.673711in}{1.010736in}}%
\pgfpathlineto{\pgfqpoint{1.673986in}{1.011565in}}%
\pgfpathlineto{\pgfqpoint{1.675535in}{1.013740in}}%
\pgfpathlineto{\pgfqpoint{1.676644in}{1.012829in}}%
\pgfpathlineto{\pgfqpoint{1.678661in}{1.010603in}}%
\pgfpathlineto{\pgfqpoint{1.678928in}{1.011429in}}%
\pgfpathlineto{\pgfqpoint{1.680684in}{1.013590in}}%
\pgfpathlineto{\pgfqpoint{1.682018in}{1.012250in}}%
\pgfpathlineto{\pgfqpoint{1.683608in}{1.010507in}}%
\pgfpathlineto{\pgfqpoint{1.683874in}{1.011334in}}%
\pgfpathlineto{\pgfqpoint{1.685682in}{1.013452in}}%
\pgfpathlineto{\pgfqpoint{1.685766in}{1.013403in}}%
\pgfpathlineto{\pgfqpoint{1.688163in}{1.010800in}}%
\pgfpathlineto{\pgfqpoint{1.688552in}{1.010384in}}%
\pgfpathlineto{\pgfqpoint{1.688823in}{1.011216in}}%
\pgfpathlineto{\pgfqpoint{1.690645in}{1.013321in}}%
\pgfpathlineto{\pgfqpoint{1.690768in}{1.013243in}}%
\pgfpathlineto{\pgfqpoint{1.693508in}{1.010272in}}%
\pgfpathlineto{\pgfqpoint{1.693991in}{1.011667in}}%
\pgfpathlineto{\pgfqpoint{1.695354in}{1.013274in}}%
\pgfpathlineto{\pgfqpoint{1.696297in}{1.012516in}}%
\pgfpathlineto{\pgfqpoint{1.698434in}{1.010145in}}%
\pgfpathlineto{\pgfqpoint{1.698712in}{1.010977in}}%
\pgfpathlineto{\pgfqpoint{1.700128in}{1.013151in}}%
\pgfpathlineto{\pgfqpoint{1.701009in}{1.012662in}}%
\pgfpathlineto{\pgfqpoint{1.703381in}{1.010029in}}%
\pgfpathlineto{\pgfqpoint{1.703658in}{1.010864in}}%
\pgfpathlineto{\pgfqpoint{1.705226in}{1.013051in}}%
\pgfpathlineto{\pgfqpoint{1.706296in}{1.012169in}}%
\pgfpathlineto{\pgfqpoint{1.708326in}{1.009918in}}%
\pgfpathlineto{\pgfqpoint{1.708603in}{1.010749in}}%
\pgfpathlineto{\pgfqpoint{1.710075in}{1.012932in}}%
\pgfpathlineto{\pgfqpoint{1.710965in}{1.012362in}}%
\pgfpathlineto{\pgfqpoint{1.713272in}{1.009804in}}%
\pgfpathlineto{\pgfqpoint{1.713549in}{1.010637in}}%
\pgfpathlineto{\pgfqpoint{1.715053in}{1.012825in}}%
\pgfpathlineto{\pgfqpoint{1.715957in}{1.012201in}}%
\pgfpathlineto{\pgfqpoint{1.718215in}{1.009698in}}%
\pgfpathlineto{\pgfqpoint{1.718491in}{1.010522in}}%
\pgfpathlineto{\pgfqpoint{1.719805in}{1.012678in}}%
\pgfpathlineto{\pgfqpoint{1.720788in}{1.012231in}}%
\pgfpathlineto{\pgfqpoint{1.723171in}{1.009604in}}%
\pgfpathlineto{\pgfqpoint{1.723438in}{1.010428in}}%
\pgfpathlineto{\pgfqpoint{1.725102in}{1.012627in}}%
\pgfpathlineto{\pgfqpoint{1.726264in}{1.011560in}}%
\pgfpathlineto{\pgfqpoint{1.728117in}{1.009528in}}%
\pgfpathlineto{\pgfqpoint{1.728383in}{1.010345in}}%
\pgfpathlineto{\pgfqpoint{1.730030in}{1.012514in}}%
\pgfpathlineto{\pgfqpoint{1.731129in}{1.011524in}}%
\pgfpathlineto{\pgfqpoint{1.733057in}{1.009389in}}%
\pgfpathlineto{\pgfqpoint{1.733335in}{1.010225in}}%
\pgfpathlineto{\pgfqpoint{1.735158in}{1.012309in}}%
\pgfpathlineto{\pgfqpoint{1.735280in}{1.012230in}}%
\pgfpathlineto{\pgfqpoint{1.737969in}{1.009285in}}%
\pgfpathlineto{\pgfqpoint{1.738007in}{1.009257in}}%
\pgfpathlineto{\pgfqpoint{1.738275in}{1.010085in}}%
\pgfpathlineto{\pgfqpoint{1.740047in}{1.012252in}}%
\pgfpathlineto{\pgfqpoint{1.741540in}{1.010731in}}%
\pgfpathlineto{\pgfqpoint{1.742943in}{1.009179in}}%
\pgfpathlineto{\pgfqpoint{1.743221in}{1.010012in}}%
\pgfpathlineto{\pgfqpoint{1.744653in}{1.012182in}}%
\pgfpathlineto{\pgfqpoint{1.745622in}{1.011575in}}%
\pgfpathlineto{\pgfqpoint{1.747892in}{1.009061in}}%
\pgfpathlineto{\pgfqpoint{1.748169in}{1.009897in}}%
\pgfpathlineto{\pgfqpoint{1.749860in}{1.012067in}}%
\pgfpathlineto{\pgfqpoint{1.751110in}{1.010873in}}%
\pgfpathlineto{\pgfqpoint{1.752844in}{1.008972in}}%
\pgfpathlineto{\pgfqpoint{1.753111in}{1.009801in}}%
\pgfpathlineto{\pgfqpoint{1.754911in}{1.011944in}}%
\pgfpathlineto{\pgfqpoint{1.755004in}{1.011892in}}%
\pgfpathlineto{\pgfqpoint{1.757102in}{1.009631in}}%
\pgfpathlineto{\pgfqpoint{1.757788in}{1.008886in}}%
\pgfpathlineto{\pgfqpoint{1.758060in}{1.009719in}}%
\pgfpathlineto{\pgfqpoint{1.759889in}{1.011824in}}%
\pgfpathlineto{\pgfqpoint{1.759939in}{1.011794in}}%
\pgfpathlineto{\pgfqpoint{1.762095in}{1.009473in}}%
\pgfpathlineto{\pgfqpoint{1.762743in}{1.008787in}}%
\pgfpathlineto{\pgfqpoint{1.763006in}{1.009610in}}%
\pgfpathlineto{\pgfqpoint{1.764573in}{1.011778in}}%
\pgfpathlineto{\pgfqpoint{1.765537in}{1.011013in}}%
\pgfpathlineto{\pgfqpoint{1.767674in}{1.008649in}}%
\pgfpathlineto{\pgfqpoint{1.767948in}{1.009476in}}%
\pgfpathlineto{\pgfqpoint{1.769432in}{1.011671in}}%
\pgfpathlineto{\pgfqpoint{1.770481in}{1.010916in}}%
\pgfpathlineto{\pgfqpoint{1.772626in}{1.008564in}}%
\pgfpathlineto{\pgfqpoint{1.772892in}{1.009394in}}%
\pgfpathlineto{\pgfqpoint{1.774715in}{1.011533in}}%
\pgfpathlineto{\pgfqpoint{1.774786in}{1.011493in}}%
\pgfpathlineto{\pgfqpoint{1.776857in}{1.009263in}}%
\pgfpathlineto{\pgfqpoint{1.777570in}{1.008489in}}%
\pgfpathlineto{\pgfqpoint{1.777839in}{1.009315in}}%
\pgfpathlineto{\pgfqpoint{1.779598in}{1.011472in}}%
\pgfpathlineto{\pgfqpoint{1.780994in}{1.010066in}}%
\pgfpathlineto{\pgfqpoint{1.782506in}{1.008395in}}%
\pgfpathlineto{\pgfqpoint{1.782782in}{1.009220in}}%
\pgfpathlineto{\pgfqpoint{1.784113in}{1.011384in}}%
\pgfpathlineto{\pgfqpoint{1.785040in}{1.010978in}}%
\pgfpathlineto{\pgfqpoint{1.787453in}{1.008310in}}%
\pgfpathlineto{\pgfqpoint{1.787730in}{1.009139in}}%
\pgfpathlineto{\pgfqpoint{1.789109in}{1.011319in}}%
\pgfpathlineto{\pgfqpoint{1.789977in}{1.010909in}}%
\pgfpathlineto{\pgfqpoint{1.792397in}{1.008233in}}%
\pgfpathlineto{\pgfqpoint{1.792673in}{1.009058in}}%
\pgfpathlineto{\pgfqpoint{1.794002in}{1.011222in}}%
\pgfpathlineto{\pgfqpoint{1.794911in}{1.010840in}}%
\pgfpathlineto{\pgfqpoint{1.797352in}{1.008155in}}%
\pgfpathlineto{\pgfqpoint{1.797621in}{1.008984in}}%
\pgfpathlineto{\pgfqpoint{1.799404in}{1.011148in}}%
\pgfpathlineto{\pgfqpoint{1.800974in}{1.009538in}}%
\pgfpathlineto{\pgfqpoint{1.802305in}{1.008101in}}%
\pgfpathlineto{\pgfqpoint{1.802568in}{1.008923in}}%
\pgfpathlineto{\pgfqpoint{1.804018in}{1.011106in}}%
\pgfpathlineto{\pgfqpoint{1.805004in}{1.010463in}}%
\pgfpathlineto{\pgfqpoint{1.807233in}{1.007999in}}%
\pgfpathlineto{\pgfqpoint{1.807511in}{1.008829in}}%
\pgfpathlineto{\pgfqpoint{1.808866in}{1.011004in}}%
\pgfpathlineto{\pgfqpoint{1.809749in}{1.010613in}}%
\pgfpathlineto{\pgfqpoint{1.812190in}{1.007937in}}%
\pgfpathlineto{\pgfqpoint{1.812451in}{1.008757in}}%
\pgfpathlineto{\pgfqpoint{1.814043in}{1.010974in}}%
\pgfpathlineto{\pgfqpoint{1.815026in}{1.010184in}}%
\pgfpathlineto{\pgfqpoint{1.817134in}{1.007869in}}%
\pgfpathlineto{\pgfqpoint{1.817402in}{1.008693in}}%
\pgfpathlineto{\pgfqpoint{1.819179in}{1.010837in}}%
\pgfpathlineto{\pgfqpoint{1.820690in}{1.009294in}}%
\pgfpathlineto{\pgfqpoint{1.822081in}{1.007781in}}%
\pgfpathlineto{\pgfqpoint{1.822346in}{1.008609in}}%
\pgfpathlineto{\pgfqpoint{1.824163in}{1.010746in}}%
\pgfpathlineto{\pgfqpoint{1.824264in}{1.010687in}}%
\pgfpathlineto{\pgfqpoint{1.826807in}{1.007928in}}%
\pgfpathlineto{\pgfqpoint{1.827026in}{1.007707in}}%
\pgfpathlineto{\pgfqpoint{1.827292in}{1.008538in}}%
\pgfpathlineto{\pgfqpoint{1.829139in}{1.010673in}}%
\pgfpathlineto{\pgfqpoint{1.829187in}{1.010644in}}%
\pgfpathlineto{\pgfqpoint{1.831505in}{1.008147in}}%
\pgfpathlineto{\pgfqpoint{1.831961in}{1.007644in}}%
\pgfpathlineto{\pgfqpoint{1.832240in}{1.008479in}}%
\pgfpathlineto{\pgfqpoint{1.833692in}{1.010657in}}%
\pgfpathlineto{\pgfqpoint{1.834756in}{1.009924in}}%
\pgfpathlineto{\pgfqpoint{1.836924in}{1.007559in}}%
\pgfpathlineto{\pgfqpoint{1.837184in}{1.008373in}}%
\pgfpathlineto{\pgfqpoint{1.838543in}{1.010552in}}%
\pgfpathlineto{\pgfqpoint{1.839445in}{1.010141in}}%
\pgfpathlineto{\pgfqpoint{1.841862in}{1.007491in}}%
\pgfpathlineto{\pgfqpoint{1.842128in}{1.008320in}}%
\pgfpathlineto{\pgfqpoint{1.843965in}{1.010450in}}%
\pgfpathlineto{\pgfqpoint{1.844030in}{1.010411in}}%
\pgfpathlineto{\pgfqpoint{1.846144in}{1.008135in}}%
\pgfpathlineto{\pgfqpoint{1.846797in}{1.007415in}}%
\pgfpathlineto{\pgfqpoint{1.847074in}{1.008240in}}%
\pgfpathlineto{\pgfqpoint{1.848402in}{1.010406in}}%
\pgfpathlineto{\pgfqpoint{1.849314in}{1.010025in}}%
\pgfpathlineto{\pgfqpoint{1.851753in}{1.007353in}}%
\pgfpathlineto{\pgfqpoint{1.852015in}{1.008172in}}%
\pgfpathlineto{\pgfqpoint{1.853593in}{1.010395in}}%
\pgfpathlineto{\pgfqpoint{1.854533in}{1.009670in}}%
\pgfpathlineto{\pgfqpoint{1.856698in}{1.007297in}}%
\pgfpathlineto{\pgfqpoint{1.856963in}{1.008120in}}%
\pgfpathlineto{\pgfqpoint{1.858731in}{1.010295in}}%
\pgfpathlineto{\pgfqpoint{1.860192in}{1.008819in}}%
\pgfpathlineto{\pgfqpoint{1.861643in}{1.007232in}}%
\pgfpathlineto{\pgfqpoint{1.861912in}{1.008061in}}%
\pgfpathlineto{\pgfqpoint{1.863701in}{1.010231in}}%
\pgfpathlineto{\pgfqpoint{1.865244in}{1.008654in}}%
\pgfpathlineto{\pgfqpoint{1.866589in}{1.007186in}}%
\pgfpathlineto{\pgfqpoint{1.866859in}{1.008016in}}%
\pgfpathlineto{\pgfqpoint{1.868711in}{1.010109in}}%
\pgfpathlineto{\pgfqpoint{1.868796in}{1.010055in}}%
\pgfpathlineto{\pgfqpoint{1.871534in}{1.007089in}}%
\pgfpathlineto{\pgfqpoint{1.872027in}{1.008500in}}%
\pgfpathlineto{\pgfqpoint{1.873414in}{1.010114in}}%
\pgfpathlineto{\pgfqpoint{1.874435in}{1.009257in}}%
\pgfpathlineto{\pgfqpoint{1.876479in}{1.007016in}}%
\pgfpathlineto{\pgfqpoint{1.876748in}{1.007843in}}%
\pgfpathlineto{\pgfqpoint{1.878538in}{1.009988in}}%
\pgfpathlineto{\pgfqpoint{1.878569in}{1.009973in}}%
\pgfpathlineto{\pgfqpoint{1.880270in}{1.008197in}}%
\pgfpathlineto{\pgfqpoint{1.881422in}{1.006933in}}%
\pgfpathlineto{\pgfqpoint{1.881698in}{1.007768in}}%
\pgfpathlineto{\pgfqpoint{1.883515in}{1.009873in}}%
\pgfpathlineto{\pgfqpoint{1.883640in}{1.009795in}}%
\pgfpathlineto{\pgfqpoint{1.886200in}{1.007009in}}%
\pgfpathlineto{\pgfqpoint{1.886370in}{1.006835in}}%
\pgfpathlineto{\pgfqpoint{1.886639in}{1.007665in}}%
\pgfpathlineto{\pgfqpoint{1.888425in}{1.009847in}}%
\pgfpathlineto{\pgfqpoint{1.888461in}{1.009831in}}%
\pgfpathlineto{\pgfqpoint{1.890061in}{1.008173in}}%
\pgfpathlineto{\pgfqpoint{1.891307in}{1.006801in}}%
\pgfpathlineto{\pgfqpoint{1.891587in}{1.007638in}}%
\pgfpathlineto{\pgfqpoint{1.893196in}{1.009803in}}%
\pgfpathlineto{\pgfqpoint{1.894251in}{1.008901in}}%
\pgfpathlineto{\pgfqpoint{1.896253in}{1.006693in}}%
\pgfpathlineto{\pgfqpoint{1.896530in}{1.007526in}}%
\pgfpathlineto{\pgfqpoint{1.897985in}{1.009735in}}%
\pgfpathlineto{\pgfqpoint{1.899053in}{1.009010in}}%
\pgfpathlineto{\pgfqpoint{1.901207in}{1.006652in}}%
\pgfpathlineto{\pgfqpoint{1.901472in}{1.007474in}}%
\pgfpathlineto{\pgfqpoint{1.903159in}{1.009677in}}%
\pgfpathlineto{\pgfqpoint{1.904347in}{1.008570in}}%
\pgfpathlineto{\pgfqpoint{1.906152in}{1.006595in}}%
\pgfpathlineto{\pgfqpoint{1.906421in}{1.007423in}}%
\pgfpathlineto{\pgfqpoint{1.908212in}{1.009586in}}%
\pgfpathlineto{\pgfqpoint{1.909723in}{1.008043in}}%
\pgfpathlineto{\pgfqpoint{1.911098in}{1.006544in}}%
\pgfpathlineto{\pgfqpoint{1.911366in}{1.007371in}}%
\pgfpathlineto{\pgfqpoint{1.913148in}{1.009531in}}%
\pgfpathlineto{\pgfqpoint{1.914685in}{1.007962in}}%
\pgfpathlineto{\pgfqpoint{1.916050in}{1.006493in}}%
\pgfpathlineto{\pgfqpoint{1.916311in}{1.007308in}}%
\pgfpathlineto{\pgfqpoint{1.917690in}{1.009495in}}%
\pgfpathlineto{\pgfqpoint{1.918564in}{1.009089in}}%
\pgfpathlineto{\pgfqpoint{1.920988in}{1.006429in}}%
\pgfpathlineto{\pgfqpoint{1.921257in}{1.007257in}}%
\pgfpathlineto{\pgfqpoint{1.923065in}{1.009418in}}%
\pgfpathlineto{\pgfqpoint{1.924628in}{1.007807in}}%
\pgfpathlineto{\pgfqpoint{1.925925in}{1.006380in}}%
\pgfpathlineto{\pgfqpoint{1.926203in}{1.007212in}}%
\pgfpathlineto{\pgfqpoint{1.927614in}{1.009395in}}%
\pgfpathlineto{\pgfqpoint{1.928501in}{1.008922in}}%
\pgfpathlineto{\pgfqpoint{1.930880in}{1.006318in}}%
\pgfpathlineto{\pgfqpoint{1.931148in}{1.007147in}}%
\pgfpathlineto{\pgfqpoint{1.932797in}{1.009366in}}%
\pgfpathlineto{\pgfqpoint{1.933945in}{1.008341in}}%
\pgfpathlineto{\pgfqpoint{1.935821in}{1.006277in}}%
\pgfpathlineto{\pgfqpoint{1.936095in}{1.007103in}}%
\pgfpathlineto{\pgfqpoint{1.937653in}{1.009295in}}%
\pgfpathlineto{\pgfqpoint{1.938726in}{1.008434in}}%
\pgfpathlineto{\pgfqpoint{1.940772in}{1.006204in}}%
\pgfpathlineto{\pgfqpoint{1.941033in}{1.007024in}}%
\pgfpathlineto{\pgfqpoint{1.942622in}{1.009257in}}%
\pgfpathlineto{\pgfqpoint{1.943576in}{1.008512in}}%
\pgfpathlineto{\pgfqpoint{1.945716in}{1.006169in}}%
\pgfpathlineto{\pgfqpoint{1.945985in}{1.006996in}}%
\pgfpathlineto{\pgfqpoint{1.947774in}{1.009147in}}%
\pgfpathlineto{\pgfqpoint{1.949297in}{1.007592in}}%
\pgfpathlineto{\pgfqpoint{1.950653in}{1.006100in}}%
\pgfpathlineto{\pgfqpoint{1.950930in}{1.006929in}}%
\pgfpathlineto{\pgfqpoint{1.952327in}{1.009114in}}%
\pgfpathlineto{\pgfqpoint{1.953317in}{1.008550in}}%
\pgfpathlineto{\pgfqpoint{1.955607in}{1.006043in}}%
\pgfpathlineto{\pgfqpoint{1.955876in}{1.006872in}}%
\pgfpathlineto{\pgfqpoint{1.957684in}{1.009038in}}%
\pgfpathlineto{\pgfqpoint{1.959339in}{1.007330in}}%
\pgfpathlineto{\pgfqpoint{1.960552in}{1.006009in}}%
\pgfpathlineto{\pgfqpoint{1.960821in}{1.006833in}}%
\pgfpathlineto{\pgfqpoint{1.962588in}{1.008980in}}%
\pgfpathlineto{\pgfqpoint{1.964078in}{1.007473in}}%
\pgfpathlineto{\pgfqpoint{1.965488in}{1.005922in}}%
\pgfpathlineto{\pgfqpoint{1.965765in}{1.006751in}}%
\pgfpathlineto{\pgfqpoint{1.967115in}{1.008929in}}%
\pgfpathlineto{\pgfqpoint{1.967991in}{1.008565in}}%
\pgfpathlineto{\pgfqpoint{1.970443in}{1.005878in}}%
\pgfpathlineto{\pgfqpoint{1.970712in}{1.006706in}}%
\pgfpathlineto{\pgfqpoint{1.972462in}{1.008895in}}%
\pgfpathlineto{\pgfqpoint{1.973812in}{1.007565in}}%
\pgfpathlineto{\pgfqpoint{1.975385in}{1.005837in}}%
\pgfpathlineto{\pgfqpoint{1.975662in}{1.006670in}}%
\pgfpathlineto{\pgfqpoint{1.977520in}{1.008732in}}%
\pgfpathlineto{\pgfqpoint{1.977651in}{1.008639in}}%
\pgfpathlineto{\pgfqpoint{1.980342in}{1.005736in}}%
\pgfpathlineto{\pgfqpoint{1.980810in}{1.007108in}}%
\pgfpathlineto{\pgfqpoint{1.982200in}{1.008783in}}%
\pgfpathlineto{\pgfqpoint{1.983254in}{1.007916in}}%
\pgfpathlineto{\pgfqpoint{1.985271in}{1.005696in}}%
\pgfpathlineto{\pgfqpoint{1.985549in}{1.006527in}}%
\pgfpathlineto{\pgfqpoint{1.986965in}{1.008715in}}%
\pgfpathlineto{\pgfqpoint{1.987892in}{1.008195in}}%
\pgfpathlineto{\pgfqpoint{1.990215in}{1.005638in}}%
\pgfpathlineto{\pgfqpoint{1.990493in}{1.006467in}}%
\pgfpathlineto{\pgfqpoint{1.991846in}{1.008650in}}%
\pgfpathlineto{\pgfqpoint{1.992740in}{1.008261in}}%
\pgfpathlineto{\pgfqpoint{1.995171in}{1.005601in}}%
\pgfpathlineto{\pgfqpoint{1.995436in}{1.006424in}}%
\pgfpathlineto{\pgfqpoint{1.997144in}{1.008629in}}%
\pgfpathlineto{\pgfqpoint{1.998232in}{1.007622in}}%
\pgfpathlineto{\pgfqpoint{2.000107in}{1.005561in}}%
\pgfpathlineto{\pgfqpoint{2.000383in}{1.006384in}}%
\pgfpathlineto{\pgfqpoint{2.001752in}{1.008553in}}%
\pgfpathlineto{\pgfqpoint{2.002783in}{1.007981in}}%
\pgfpathlineto{\pgfqpoint{2.005062in}{1.005491in}}%
\pgfpathlineto{\pgfqpoint{2.005327in}{1.006316in}}%
\pgfpathlineto{\pgfqpoint{2.007091in}{1.008516in}}%
\pgfpathlineto{\pgfqpoint{2.008344in}{1.007284in}}%
\pgfpathlineto{\pgfqpoint{2.009997in}{1.005467in}}%
\pgfpathlineto{\pgfqpoint{2.010273in}{1.006292in}}%
\pgfpathlineto{\pgfqpoint{2.011600in}{1.008459in}}%
\pgfpathlineto{\pgfqpoint{2.012495in}{1.008107in}}%
\pgfpathlineto{\pgfqpoint{2.014942in}{1.005415in}}%
\pgfpathlineto{\pgfqpoint{2.015220in}{1.006244in}}%
\pgfpathlineto{\pgfqpoint{2.016576in}{1.008430in}}%
\pgfpathlineto{\pgfqpoint{2.017492in}{1.008015in}}%
\pgfpathlineto{\pgfqpoint{2.019889in}{1.005378in}}%
\pgfpathlineto{\pgfqpoint{2.020165in}{1.006203in}}%
\pgfpathlineto{\pgfqpoint{2.021491in}{1.008376in}}%
\pgfpathlineto{\pgfqpoint{2.022383in}{1.008032in}}%
\pgfpathlineto{\pgfqpoint{2.024834in}{1.005336in}}%
\pgfpathlineto{\pgfqpoint{2.025111in}{1.006167in}}%
\pgfpathlineto{\pgfqpoint{2.026477in}{1.008352in}}%
\pgfpathlineto{\pgfqpoint{2.027520in}{1.007781in}}%
\pgfpathlineto{\pgfqpoint{2.029789in}{1.005301in}}%
\pgfpathlineto{\pgfqpoint{2.030054in}{1.006124in}}%
\pgfpathlineto{\pgfqpoint{2.031764in}{1.008332in}}%
\pgfpathlineto{\pgfqpoint{2.032883in}{1.007289in}}%
\pgfpathlineto{\pgfqpoint{2.034725in}{1.005265in}}%
\pgfpathlineto{\pgfqpoint{2.035002in}{1.006091in}}%
\pgfpathlineto{\pgfqpoint{2.036354in}{1.008266in}}%
\pgfpathlineto{\pgfqpoint{2.037231in}{1.007897in}}%
\pgfpathlineto{\pgfqpoint{2.039680in}{1.005218in}}%
\pgfpathlineto{\pgfqpoint{2.039948in}{1.006045in}}%
\pgfpathlineto{\pgfqpoint{2.041688in}{1.008246in}}%
\pgfpathlineto{\pgfqpoint{2.043020in}{1.006946in}}%
\pgfpathlineto{\pgfqpoint{2.044625in}{1.005196in}}%
\pgfpathlineto{\pgfqpoint{2.044893in}{1.006020in}}%
\pgfpathlineto{\pgfqpoint{2.046664in}{1.008171in}}%
\pgfpathlineto{\pgfqpoint{2.048082in}{1.006744in}}%
\pgfpathlineto{\pgfqpoint{2.049561in}{1.005120in}}%
\pgfpathlineto{\pgfqpoint{2.049839in}{1.005952in}}%
\pgfpathlineto{\pgfqpoint{2.051275in}{1.008144in}}%
\pgfpathlineto{\pgfqpoint{2.052205in}{1.007597in}}%
\pgfpathlineto{\pgfqpoint{2.054516in}{1.005071in}}%
\pgfpathlineto{\pgfqpoint{2.054785in}{1.005902in}}%
\pgfpathlineto{\pgfqpoint{2.056563in}{1.008088in}}%
\pgfpathlineto{\pgfqpoint{2.058061in}{1.006576in}}%
\pgfpathlineto{\pgfqpoint{2.059461in}{1.005052in}}%
\pgfpathlineto{\pgfqpoint{2.059730in}{1.005879in}}%
\pgfpathlineto{\pgfqpoint{2.061521in}{1.008049in}}%
\pgfpathlineto{\pgfqpoint{2.063031in}{1.006514in}}%
\pgfpathlineto{\pgfqpoint{2.064397in}{1.005013in}}%
\pgfpathlineto{\pgfqpoint{2.064675in}{1.005843in}}%
\pgfpathlineto{\pgfqpoint{2.066060in}{1.008032in}}%
\pgfpathlineto{\pgfqpoint{2.066971in}{1.007579in}}%
\pgfpathlineto{\pgfqpoint{2.069343in}{1.004973in}}%
\pgfpathlineto{\pgfqpoint{2.069620in}{1.005801in}}%
\pgfpathlineto{\pgfqpoint{2.070967in}{1.007982in}}%
\pgfpathlineto{\pgfqpoint{2.071866in}{1.007601in}}%
\pgfpathlineto{\pgfqpoint{2.074289in}{1.004938in}}%
\pgfpathlineto{\pgfqpoint{2.074566in}{1.005769in}}%
\pgfpathlineto{\pgfqpoint{2.075969in}{1.007964in}}%
\pgfpathlineto{\pgfqpoint{2.076928in}{1.007436in}}%
\pgfpathlineto{\pgfqpoint{2.079234in}{1.004901in}}%
\pgfpathlineto{\pgfqpoint{2.079512in}{1.005732in}}%
\pgfpathlineto{\pgfqpoint{2.080913in}{1.007927in}}%
\pgfpathlineto{\pgfqpoint{2.081875in}{1.007398in}}%
\pgfpathlineto{\pgfqpoint{2.084179in}{1.004866in}}%
\pgfpathlineto{\pgfqpoint{2.084457in}{1.005698in}}%
\pgfpathlineto{\pgfqpoint{2.085863in}{1.007893in}}%
\pgfpathlineto{\pgfqpoint{2.086769in}{1.007419in}}%
\pgfpathlineto{\pgfqpoint{2.089125in}{1.004830in}}%
\pgfpathlineto{\pgfqpoint{2.089402in}{1.005661in}}%
\pgfpathlineto{\pgfqpoint{2.090797in}{1.007856in}}%
\pgfpathlineto{\pgfqpoint{2.091844in}{1.007243in}}%
\pgfpathlineto{\pgfqpoint{2.094070in}{1.004796in}}%
\pgfpathlineto{\pgfqpoint{2.094347in}{1.005625in}}%
\pgfpathlineto{\pgfqpoint{2.095721in}{1.007814in}}%
\pgfpathlineto{\pgfqpoint{2.096642in}{1.007369in}}%
\pgfpathlineto{\pgfqpoint{2.099016in}{1.004761in}}%
\pgfpathlineto{\pgfqpoint{2.099294in}{1.005592in}}%
\pgfpathlineto{\pgfqpoint{2.100696in}{1.007788in}}%
\pgfpathlineto{\pgfqpoint{2.101666in}{1.007249in}}%
\pgfpathlineto{\pgfqpoint{2.103961in}{1.004728in}}%
\pgfpathlineto{\pgfqpoint{2.104239in}{1.005559in}}%
\pgfpathlineto{\pgfqpoint{2.105628in}{1.007752in}}%
\pgfpathlineto{\pgfqpoint{2.106622in}{1.007203in}}%
\pgfpathlineto{\pgfqpoint{2.108906in}{1.004694in}}%
\pgfpathlineto{\pgfqpoint{2.109184in}{1.005525in}}%
\pgfpathlineto{\pgfqpoint{2.110578in}{1.007720in}}%
\pgfpathlineto{\pgfqpoint{2.111621in}{1.007112in}}%
\pgfpathlineto{\pgfqpoint{2.113852in}{1.004662in}}%
\pgfpathlineto{\pgfqpoint{2.114129in}{1.005491in}}%
\pgfpathlineto{\pgfqpoint{2.115509in}{1.007682in}}%
\pgfpathlineto{\pgfqpoint{2.116380in}{1.007285in}}%
\pgfpathlineto{\pgfqpoint{2.118797in}{1.004629in}}%
\pgfpathlineto{\pgfqpoint{2.119075in}{1.005463in}}%
\pgfpathlineto{\pgfqpoint{2.120487in}{1.007661in}}%
\pgfpathlineto{\pgfqpoint{2.121345in}{1.007232in}}%
\pgfpathlineto{\pgfqpoint{2.123743in}{1.004598in}}%
\pgfpathlineto{\pgfqpoint{2.124020in}{1.005427in}}%
\pgfpathlineto{\pgfqpoint{2.125401in}{1.007620in}}%
\pgfpathlineto{\pgfqpoint{2.126266in}{1.007228in}}%
\pgfpathlineto{\pgfqpoint{2.128689in}{1.004567in}}%
\pgfpathlineto{\pgfqpoint{2.128963in}{1.005386in}}%
\pgfpathlineto{\pgfqpoint{2.130279in}{1.007561in}}%
\pgfpathlineto{\pgfqpoint{2.131268in}{1.007134in}}%
\pgfpathlineto{\pgfqpoint{2.133634in}{1.004536in}}%
\pgfpathlineto{\pgfqpoint{2.133910in}{1.005360in}}%
\pgfpathlineto{\pgfqpoint{2.135235in}{1.007535in}}%
\pgfpathlineto{\pgfqpoint{2.136139in}{1.007186in}}%
\pgfpathlineto{\pgfqpoint{2.138579in}{1.004505in}}%
\pgfpathlineto{\pgfqpoint{2.138855in}{1.005329in}}%
\pgfpathlineto{\pgfqpoint{2.140177in}{1.007503in}}%
\pgfpathlineto{\pgfqpoint{2.141076in}{1.007164in}}%
\pgfpathlineto{\pgfqpoint{2.143525in}{1.004474in}}%
\pgfpathlineto{\pgfqpoint{2.143802in}{1.005303in}}%
\pgfpathlineto{\pgfqpoint{2.145181in}{1.007494in}}%
\pgfpathlineto{\pgfqpoint{2.146045in}{1.007107in}}%
\pgfpathlineto{\pgfqpoint{2.148471in}{1.004443in}}%
\pgfpathlineto{\pgfqpoint{2.148748in}{1.005273in}}%
\pgfpathlineto{\pgfqpoint{2.150138in}{1.007468in}}%
\pgfpathlineto{\pgfqpoint{2.151123in}{1.006931in}}%
\pgfpathlineto{\pgfqpoint{2.153416in}{1.004413in}}%
\pgfpathlineto{\pgfqpoint{2.153691in}{1.005236in}}%
\pgfpathlineto{\pgfqpoint{2.155018in}{1.007413in}}%
\pgfpathlineto{\pgfqpoint{2.155934in}{1.007048in}}%
\pgfpathlineto{\pgfqpoint{2.158361in}{1.004383in}}%
\pgfpathlineto{\pgfqpoint{2.158639in}{1.005212in}}%
\pgfpathlineto{\pgfqpoint{2.160027in}{1.007408in}}%
\pgfpathlineto{\pgfqpoint{2.161002in}{1.006884in}}%
\pgfpathlineto{\pgfqpoint{2.163307in}{1.004353in}}%
\pgfpathlineto{\pgfqpoint{2.163583in}{1.005180in}}%
\pgfpathlineto{\pgfqpoint{2.164922in}{1.007361in}}%
\pgfpathlineto{\pgfqpoint{2.165960in}{1.006841in}}%
\pgfpathlineto{\pgfqpoint{2.168252in}{1.004325in}}%
\pgfpathlineto{\pgfqpoint{2.168528in}{1.005147in}}%
\pgfpathlineto{\pgfqpoint{2.169849in}{1.007322in}}%
\pgfpathlineto{\pgfqpoint{2.170750in}{1.006983in}}%
\pgfpathlineto{\pgfqpoint{2.173198in}{1.004296in}}%
\pgfpathlineto{\pgfqpoint{2.173475in}{1.005125in}}%
\pgfpathlineto{\pgfqpoint{2.174864in}{1.007321in}}%
\pgfpathlineto{\pgfqpoint{2.175862in}{1.006772in}}%
\pgfpathlineto{\pgfqpoint{2.178143in}{1.004268in}}%
\pgfpathlineto{\pgfqpoint{2.178420in}{1.005096in}}%
\pgfpathlineto{\pgfqpoint{2.179799in}{1.007289in}}%
\pgfpathlineto{\pgfqpoint{2.180674in}{1.006890in}}%
\pgfpathlineto{\pgfqpoint{2.183089in}{1.004240in}}%
\pgfpathlineto{\pgfqpoint{2.183365in}{1.005065in}}%
\pgfpathlineto{\pgfqpoint{2.184695in}{1.007244in}}%
\pgfpathlineto{\pgfqpoint{2.185643in}{1.006838in}}%
\pgfpathlineto{\pgfqpoint{2.188034in}{1.004213in}}%
\pgfpathlineto{\pgfqpoint{2.188312in}{1.005043in}}%
\pgfpathlineto{\pgfqpoint{2.189708in}{1.007239in}}%
\pgfpathlineto{\pgfqpoint{2.190733in}{1.006650in}}%
\pgfpathlineto{\pgfqpoint{2.192979in}{1.004184in}}%
\pgfpathlineto{\pgfqpoint{2.193255in}{1.005007in}}%
\pgfpathlineto{\pgfqpoint{2.194581in}{1.007186in}}%
\pgfpathlineto{\pgfqpoint{2.195470in}{1.006853in}}%
\pgfpathlineto{\pgfqpoint{2.197925in}{1.004159in}}%
\pgfpathlineto{\pgfqpoint{2.198202in}{1.004987in}}%
\pgfpathlineto{\pgfqpoint{2.199575in}{1.007178in}}%
\pgfpathlineto{\pgfqpoint{2.200537in}{1.006693in}}%
\pgfpathlineto{\pgfqpoint{2.202870in}{1.004132in}}%
\pgfpathlineto{\pgfqpoint{2.203149in}{1.004966in}}%
\pgfpathlineto{\pgfqpoint{2.204575in}{1.007168in}}%
\pgfpathlineto{\pgfqpoint{2.205618in}{1.006516in}}%
\pgfpathlineto{\pgfqpoint{2.207816in}{1.004105in}}%
\pgfpathlineto{\pgfqpoint{2.208091in}{1.004927in}}%
\pgfpathlineto{\pgfqpoint{2.209410in}{1.007103in}}%
\pgfpathlineto{\pgfqpoint{2.210343in}{1.006734in}}%
\pgfpathlineto{\pgfqpoint{2.212761in}{1.004080in}}%
\pgfpathlineto{\pgfqpoint{2.213038in}{1.004908in}}%
\pgfpathlineto{\pgfqpoint{2.214410in}{1.007099in}}%
\pgfpathlineto{\pgfqpoint{2.215398in}{1.006588in}}%
\pgfpathlineto{\pgfqpoint{2.217707in}{1.004055in}}%
\pgfpathlineto{\pgfqpoint{2.217983in}{1.004882in}}%
\pgfpathlineto{\pgfqpoint{2.219356in}{1.007074in}}%
\pgfpathlineto{\pgfqpoint{2.220325in}{1.006583in}}%
\pgfpathlineto{\pgfqpoint{2.222652in}{1.004030in}}%
\pgfpathlineto{\pgfqpoint{2.222929in}{1.004858in}}%
\pgfpathlineto{\pgfqpoint{2.224311in}{1.007053in}}%
\pgfpathlineto{\pgfqpoint{2.225191in}{1.006646in}}%
\pgfpathlineto{\pgfqpoint{2.227598in}{1.004004in}}%
\pgfpathlineto{\pgfqpoint{2.227874in}{1.004832in}}%
\pgfpathlineto{\pgfqpoint{2.229247in}{1.007025in}}%
\pgfpathlineto{\pgfqpoint{2.230225in}{1.006524in}}%
\pgfpathlineto{\pgfqpoint{2.232543in}{1.003980in}}%
\pgfpathlineto{\pgfqpoint{2.232820in}{1.004810in}}%
\pgfpathlineto{\pgfqpoint{2.234204in}{1.007004in}}%
\pgfpathlineto{\pgfqpoint{2.235101in}{1.006575in}}%
\pgfpathlineto{\pgfqpoint{2.237490in}{1.003955in}}%
\pgfpathlineto{\pgfqpoint{2.237767in}{1.004788in}}%
\pgfpathlineto{\pgfqpoint{2.239198in}{1.006990in}}%
\pgfpathlineto{\pgfqpoint{2.240158in}{1.006426in}}%
\pgfpathlineto{\pgfqpoint{2.242435in}{1.003928in}}%
\pgfpathlineto{\pgfqpoint{2.242713in}{1.004761in}}%
\pgfpathlineto{\pgfqpoint{2.244139in}{1.006963in}}%
\pgfpathlineto{\pgfqpoint{2.245184in}{1.006311in}}%
\pgfpathlineto{\pgfqpoint{2.247381in}{1.003902in}}%
\pgfpathlineto{\pgfqpoint{2.247657in}{1.004727in}}%
\pgfpathlineto{\pgfqpoint{2.249026in}{1.006918in}}%
\pgfpathlineto{\pgfqpoint{2.250025in}{1.006398in}}%
\pgfpathlineto{\pgfqpoint{2.252326in}{1.003874in}}%
\pgfpathlineto{\pgfqpoint{2.252604in}{1.004707in}}%
\pgfpathlineto{\pgfqpoint{2.254031in}{1.006910in}}%
\pgfpathlineto{\pgfqpoint{2.255041in}{1.006296in}}%
\pgfpathlineto{\pgfqpoint{2.257272in}{1.003849in}}%
\pgfpathlineto{\pgfqpoint{2.257549in}{1.004680in}}%
\pgfpathlineto{\pgfqpoint{2.259030in}{1.006893in}}%
\pgfpathlineto{\pgfqpoint{2.260084in}{1.006160in}}%
\pgfpathlineto{\pgfqpoint{2.262218in}{1.003820in}}%
\pgfpathlineto{\pgfqpoint{2.262494in}{1.004647in}}%
\pgfpathlineto{\pgfqpoint{2.263905in}{1.006848in}}%
\pgfpathlineto{\pgfqpoint{2.264796in}{1.006388in}}%
\pgfpathlineto{\pgfqpoint{2.267163in}{1.003792in}}%
\pgfpathlineto{\pgfqpoint{2.267438in}{1.004614in}}%
\pgfpathlineto{\pgfqpoint{2.268769in}{1.006794in}}%
\pgfpathlineto{\pgfqpoint{2.269715in}{1.006391in}}%
\pgfpathlineto{\pgfqpoint{2.272109in}{1.003765in}}%
\pgfpathlineto{\pgfqpoint{2.272385in}{1.004592in}}%
\pgfpathlineto{\pgfqpoint{2.273795in}{1.006792in}}%
\pgfpathlineto{\pgfqpoint{2.274680in}{1.006340in}}%
\pgfpathlineto{\pgfqpoint{2.277054in}{1.003737in}}%
\pgfpathlineto{\pgfqpoint{2.277331in}{1.004567in}}%
\pgfpathlineto{\pgfqpoint{2.278804in}{1.006780in}}%
\pgfpathlineto{\pgfqpoint{2.279702in}{1.006230in}}%
\pgfpathlineto{\pgfqpoint{2.282000in}{1.003710in}}%
\pgfpathlineto{\pgfqpoint{2.282276in}{1.004538in}}%
\pgfpathlineto{\pgfqpoint{2.283692in}{1.006741in}}%
\pgfpathlineto{\pgfqpoint{2.284595in}{1.006261in}}%
\pgfpathlineto{\pgfqpoint{2.286945in}{1.003684in}}%
\pgfpathlineto{\pgfqpoint{2.287221in}{1.004510in}}%
\pgfpathlineto{\pgfqpoint{2.288629in}{1.006711in}}%
\pgfpathlineto{\pgfqpoint{2.289526in}{1.006250in}}%
\pgfpathlineto{\pgfqpoint{2.291891in}{1.003657in}}%
\pgfpathlineto{\pgfqpoint{2.292167in}{1.004485in}}%
\pgfpathlineto{\pgfqpoint{2.293588in}{1.006689in}}%
\pgfpathlineto{\pgfqpoint{2.294595in}{1.006087in}}%
\pgfpathlineto{\pgfqpoint{2.296836in}{1.003631in}}%
\pgfpathlineto{\pgfqpoint{2.297112in}{1.004461in}}%
\pgfpathlineto{\pgfqpoint{2.298580in}{1.006673in}}%
\pgfpathlineto{\pgfqpoint{2.299428in}{1.006186in}}%
\pgfpathlineto{\pgfqpoint{2.301781in}{1.003606in}}%
\pgfpathlineto{\pgfqpoint{2.302056in}{1.004429in}}%
\pgfpathlineto{\pgfqpoint{2.303435in}{1.006624in}}%
\pgfpathlineto{\pgfqpoint{2.304301in}{1.006240in}}%
\pgfpathlineto{\pgfqpoint{2.306725in}{1.003582in}}%
\pgfpathlineto{\pgfqpoint{2.307003in}{1.004413in}}%
\pgfpathlineto{\pgfqpoint{2.308410in}{1.006612in}}%
\pgfpathlineto{\pgfqpoint{2.309317in}{1.006141in}}%
\pgfpathlineto{\pgfqpoint{2.311672in}{1.003559in}}%
\pgfpathlineto{\pgfqpoint{2.311949in}{1.004387in}}%
\pgfpathlineto{\pgfqpoint{2.313371in}{1.006593in}}%
\pgfpathlineto{\pgfqpoint{2.314403in}{1.005963in}}%
\pgfpathlineto{\pgfqpoint{2.316616in}{1.003537in}}%
\pgfpathlineto{\pgfqpoint{2.316894in}{1.004368in}}%
\pgfpathlineto{\pgfqpoint{2.318300in}{1.006567in}}%
\pgfpathlineto{\pgfqpoint{2.319164in}{1.006145in}}%
\pgfpathlineto{\pgfqpoint{2.321563in}{1.003514in}}%
\pgfpathlineto{\pgfqpoint{2.321839in}{1.004341in}}%
\pgfpathlineto{\pgfqpoint{2.323250in}{1.006544in}}%
\pgfpathlineto{\pgfqpoint{2.324111in}{1.006120in}}%
\pgfpathlineto{\pgfqpoint{2.326507in}{1.003493in}}%
\pgfpathlineto{\pgfqpoint{2.326785in}{1.004324in}}%
\pgfpathlineto{\pgfqpoint{2.328197in}{1.006525in}}%
\pgfpathlineto{\pgfqpoint{2.329077in}{1.006077in}}%
\pgfpathlineto{\pgfqpoint{2.331454in}{1.003471in}}%
\pgfpathlineto{\pgfqpoint{2.331729in}{1.004294in}}%
\pgfpathlineto{\pgfqpoint{2.333063in}{1.006476in}}%
\pgfpathlineto{\pgfqpoint{2.334056in}{1.006018in}}%
\pgfpathlineto{\pgfqpoint{2.336398in}{1.003450in}}%
\pgfpathlineto{\pgfqpoint{2.336676in}{1.004281in}}%
\pgfpathlineto{\pgfqpoint{2.338089in}{1.006483in}}%
\pgfpathlineto{\pgfqpoint{2.338975in}{1.006026in}}%
\pgfpathlineto{\pgfqpoint{2.341345in}{1.003429in}}%
\pgfpathlineto{\pgfqpoint{2.341622in}{1.004262in}}%
\pgfpathlineto{\pgfqpoint{2.343068in}{1.006470in}}%
\pgfpathlineto{\pgfqpoint{2.343976in}{1.005944in}}%
\pgfpathlineto{\pgfqpoint{2.346292in}{1.003406in}}%
\pgfpathlineto{\pgfqpoint{2.346568in}{1.004234in}}%
\pgfpathlineto{\pgfqpoint{2.347975in}{1.006431in}}%
\pgfpathlineto{\pgfqpoint{2.348880in}{1.005963in}}%
\pgfpathlineto{\pgfqpoint{2.351235in}{1.003382in}}%
\pgfpathlineto{\pgfqpoint{2.351512in}{1.004213in}}%
\pgfpathlineto{\pgfqpoint{2.352927in}{1.006415in}}%
\pgfpathlineto{\pgfqpoint{2.353833in}{1.005935in}}%
\pgfpathlineto{\pgfqpoint{2.356181in}{1.003361in}}%
\pgfpathlineto{\pgfqpoint{2.356456in}{1.004184in}}%
\pgfpathlineto{\pgfqpoint{2.357788in}{1.006366in}}%
\pgfpathlineto{\pgfqpoint{2.358737in}{1.005961in}}%
\pgfpathlineto{\pgfqpoint{2.361126in}{1.003342in}}%
\pgfpathlineto{\pgfqpoint{2.361403in}{1.004174in}}%
\pgfpathlineto{\pgfqpoint{2.362887in}{1.006388in}}%
\pgfpathlineto{\pgfqpoint{2.363952in}{1.005642in}}%
\pgfpathlineto{\pgfqpoint{2.366072in}{1.003320in}}%
\pgfpathlineto{\pgfqpoint{2.366350in}{1.004154in}}%
\pgfpathlineto{\pgfqpoint{2.367802in}{1.006364in}}%
\pgfpathlineto{\pgfqpoint{2.368830in}{1.005698in}}%
\pgfpathlineto{\pgfqpoint{2.371017in}{1.003302in}}%
\pgfpathlineto{\pgfqpoint{2.371294in}{1.004135in}}%
\pgfpathlineto{\pgfqpoint{2.372791in}{1.006349in}}%
\pgfpathlineto{\pgfqpoint{2.373642in}{1.005824in}}%
\pgfpathlineto{\pgfqpoint{2.375963in}{1.003280in}}%
\pgfpathlineto{\pgfqpoint{2.376239in}{1.004110in}}%
\pgfpathlineto{\pgfqpoint{2.377694in}{1.006323in}}%
\pgfpathlineto{\pgfqpoint{2.378751in}{1.005625in}}%
\pgfpathlineto{\pgfqpoint{2.380908in}{1.003262in}}%
\pgfpathlineto{\pgfqpoint{2.381185in}{1.004095in}}%
\pgfpathlineto{\pgfqpoint{2.382673in}{1.006310in}}%
\pgfpathlineto{\pgfqpoint{2.383663in}{1.005642in}}%
\pgfpathlineto{\pgfqpoint{2.385853in}{1.003243in}}%
\pgfpathlineto{\pgfqpoint{2.386131in}{1.004075in}}%
\pgfpathlineto{\pgfqpoint{2.387561in}{1.006281in}}%
\pgfpathlineto{\pgfqpoint{2.388529in}{1.005713in}}%
\pgfpathlineto{\pgfqpoint{2.390798in}{1.003226in}}%
\pgfpathlineto{\pgfqpoint{2.391076in}{1.004057in}}%
\pgfpathlineto{\pgfqpoint{2.392492in}{1.006261in}}%
\pgfpathlineto{\pgfqpoint{2.393428in}{1.005747in}}%
\pgfpathlineto{\pgfqpoint{2.395744in}{1.003209in}}%
\pgfpathlineto{\pgfqpoint{2.396021in}{1.004041in}}%
\pgfpathlineto{\pgfqpoint{2.397436in}{1.006245in}}%
\pgfpathlineto{\pgfqpoint{2.398367in}{1.005738in}}%
\pgfpathlineto{\pgfqpoint{2.400690in}{1.003192in}}%
\pgfpathlineto{\pgfqpoint{2.400966in}{1.004022in}}%
\pgfpathlineto{\pgfqpoint{2.402425in}{1.006235in}}%
\pgfpathlineto{\pgfqpoint{2.403422in}{1.005598in}}%
\pgfpathlineto{\pgfqpoint{2.405635in}{1.003174in}}%
\pgfpathlineto{\pgfqpoint{2.405912in}{1.004005in}}%
\pgfpathlineto{\pgfqpoint{2.407326in}{1.006210in}}%
\pgfpathlineto{\pgfqpoint{2.408231in}{1.005733in}}%
\pgfpathlineto{\pgfqpoint{2.410580in}{1.003159in}}%
\pgfpathlineto{\pgfqpoint{2.410857in}{1.003987in}}%
\pgfpathlineto{\pgfqpoint{2.412248in}{1.006187in}}%
\pgfpathlineto{\pgfqpoint{2.413220in}{1.005668in}}%
\pgfpathlineto{\pgfqpoint{2.415526in}{1.003142in}}%
\pgfpathlineto{\pgfqpoint{2.415800in}{1.003962in}}%
\pgfpathlineto{\pgfqpoint{2.417149in}{1.006152in}}%
\pgfpathlineto{\pgfqpoint{2.418046in}{1.005782in}}%
\pgfpathlineto{\pgfqpoint{2.420471in}{1.003125in}}%
\pgfpathlineto{\pgfqpoint{2.420749in}{1.003957in}}%
\pgfpathlineto{\pgfqpoint{2.422232in}{1.006173in}}%
\pgfpathlineto{\pgfqpoint{2.423269in}{1.005460in}}%
\pgfpathlineto{\pgfqpoint{2.425417in}{1.003108in}}%
\pgfpathlineto{\pgfqpoint{2.425694in}{1.003941in}}%
\pgfpathlineto{\pgfqpoint{2.427186in}{1.006157in}}%
\pgfpathlineto{\pgfqpoint{2.428093in}{1.005578in}}%
\pgfpathlineto{\pgfqpoint{2.430363in}{1.003092in}}%
\pgfpathlineto{\pgfqpoint{2.430640in}{1.003923in}}%
\pgfpathlineto{\pgfqpoint{2.432118in}{1.006139in}}%
\pgfpathlineto{\pgfqpoint{2.433130in}{1.005459in}}%
\pgfpathlineto{\pgfqpoint{2.435308in}{1.003074in}}%
\pgfpathlineto{\pgfqpoint{2.435585in}{1.003906in}}%
\pgfpathlineto{\pgfqpoint{2.437010in}{1.006112in}}%
\pgfpathlineto{\pgfqpoint{2.438042in}{1.005481in}}%
\pgfpathlineto{\pgfqpoint{2.440253in}{1.003059in}}%
\pgfpathlineto{\pgfqpoint{2.440531in}{1.003891in}}%
\pgfpathlineto{\pgfqpoint{2.441985in}{1.006102in}}%
\pgfpathlineto{\pgfqpoint{2.443041in}{1.005406in}}%
\pgfpathlineto{\pgfqpoint{2.445199in}{1.003042in}}%
\pgfpathlineto{\pgfqpoint{2.445476in}{1.003874in}}%
\pgfpathlineto{\pgfqpoint{2.446951in}{1.006090in}}%
\pgfpathlineto{\pgfqpoint{2.447896in}{1.005488in}}%
\pgfpathlineto{\pgfqpoint{2.450144in}{1.003026in}}%
\pgfpathlineto{\pgfqpoint{2.450421in}{1.003858in}}%
\pgfpathlineto{\pgfqpoint{2.451901in}{1.006076in}}%
\pgfpathlineto{\pgfqpoint{2.452934in}{1.005372in}}%
\pgfpathlineto{\pgfqpoint{2.455091in}{1.003011in}}%
\pgfpathlineto{\pgfqpoint{2.455368in}{1.003842in}}%
\pgfpathlineto{\pgfqpoint{2.456782in}{1.006046in}}%
\pgfpathlineto{\pgfqpoint{2.457670in}{1.005588in}}%
\pgfpathlineto{\pgfqpoint{2.460052in}{1.003009in}}%
\pgfpathlineto{\pgfqpoint{2.460312in}{1.003823in}}%
\pgfpathlineto{\pgfqpoint{2.461699in}{1.006017in}}%
\pgfpathlineto{\pgfqpoint{2.462613in}{1.005569in}}%
\pgfpathlineto{\pgfqpoint{2.464981in}{1.002976in}}%
\pgfpathlineto{\pgfqpoint{2.465258in}{1.003805in}}%
\pgfpathlineto{\pgfqpoint{2.466665in}{1.006003in}}%
\pgfpathlineto{\pgfqpoint{2.467571in}{1.005534in}}%
\pgfpathlineto{\pgfqpoint{2.469926in}{1.002954in}}%
\pgfpathlineto{\pgfqpoint{2.470202in}{1.003777in}}%
\pgfpathlineto{\pgfqpoint{2.471532in}{1.005961in}}%
\pgfpathlineto{\pgfqpoint{2.472455in}{1.005589in}}%
\pgfpathlineto{\pgfqpoint{2.474888in}{1.002954in}}%
\pgfpathlineto{\pgfqpoint{2.475148in}{1.003765in}}%
\pgfpathlineto{\pgfqpoint{2.476507in}{1.005953in}}%
\pgfpathlineto{\pgfqpoint{2.477388in}{1.005581in}}%
\pgfpathlineto{\pgfqpoint{2.479817in}{1.002922in}}%
\pgfpathlineto{\pgfqpoint{2.480092in}{1.003742in}}%
\pgfpathlineto{\pgfqpoint{2.481414in}{1.005921in}}%
\pgfpathlineto{\pgfqpoint{2.482311in}{1.005592in}}%
\pgfpathlineto{\pgfqpoint{2.484763in}{1.002907in}}%
\pgfpathlineto{\pgfqpoint{2.485039in}{1.003735in}}%
\pgfpathlineto{\pgfqpoint{2.486455in}{1.005942in}}%
\pgfpathlineto{\pgfqpoint{2.487365in}{1.005458in}}%
\pgfpathlineto{\pgfqpoint{2.489725in}{1.002906in}}%
\pgfpathlineto{\pgfqpoint{2.489986in}{1.003724in}}%
\pgfpathlineto{\pgfqpoint{2.491480in}{1.005940in}}%
\pgfpathlineto{\pgfqpoint{2.492369in}{1.005377in}}%
\pgfpathlineto{\pgfqpoint{2.494654in}{1.002875in}}%
\pgfpathlineto{\pgfqpoint{2.494932in}{1.003707in}}%
\pgfpathlineto{\pgfqpoint{2.496361in}{1.005914in}}%
\pgfpathlineto{\pgfqpoint{2.497371in}{1.005302in}}%
\pgfpathlineto{\pgfqpoint{2.499598in}{1.002864in}}%
\pgfpathlineto{\pgfqpoint{2.499874in}{1.003691in}}%
\pgfpathlineto{\pgfqpoint{2.501223in}{1.005877in}}%
\pgfpathlineto{\pgfqpoint{2.502129in}{1.005495in}}%
\pgfpathlineto{\pgfqpoint{2.504544in}{1.002851in}}%
\pgfpathlineto{\pgfqpoint{2.504821in}{1.003683in}}%
\pgfpathlineto{\pgfqpoint{2.506241in}{1.005890in}}%
\pgfpathlineto{\pgfqpoint{2.507225in}{1.005320in}}%
\pgfpathlineto{\pgfqpoint{2.509494in}{1.002841in}}%
\pgfpathlineto{\pgfqpoint{2.509767in}{1.003670in}}%
\pgfpathlineto{\pgfqpoint{2.511384in}{1.005886in}}%
\pgfpathlineto{\pgfqpoint{2.512410in}{1.005036in}}%
\pgfpathlineto{\pgfqpoint{2.514436in}{1.002818in}}%
\pgfpathlineto{\pgfqpoint{2.514714in}{1.003651in}}%
\pgfpathlineto{\pgfqpoint{2.516139in}{1.005858in}}%
\pgfpathlineto{\pgfqpoint{2.517184in}{1.005214in}}%
\pgfpathlineto{\pgfqpoint{2.519379in}{1.002811in}}%
\pgfpathlineto{\pgfqpoint{2.519655in}{1.003636in}}%
\pgfpathlineto{\pgfqpoint{2.521008in}{1.005822in}}%
\pgfpathlineto{\pgfqpoint{2.521901in}{1.005449in}}%
\pgfpathlineto{\pgfqpoint{2.524326in}{1.002793in}}%
\pgfpathlineto{\pgfqpoint{2.524602in}{1.003616in}}%
\pgfpathlineto{\pgfqpoint{2.525938in}{1.005804in}}%
\pgfpathlineto{\pgfqpoint{2.526945in}{1.005331in}}%
\pgfpathlineto{\pgfqpoint{2.529288in}{1.002796in}}%
\pgfpathlineto{\pgfqpoint{2.529547in}{1.003606in}}%
\pgfpathlineto{\pgfqpoint{2.530899in}{1.005792in}}%
\pgfpathlineto{\pgfqpoint{2.531888in}{1.005314in}}%
\pgfpathlineto{\pgfqpoint{2.534217in}{1.002765in}}%
\pgfpathlineto{\pgfqpoint{2.534494in}{1.003596in}}%
\pgfpathlineto{\pgfqpoint{2.535909in}{1.005800in}}%
\pgfpathlineto{\pgfqpoint{2.536817in}{1.005320in}}%
\pgfpathlineto{\pgfqpoint{2.539163in}{1.002751in}}%
\pgfpathlineto{\pgfqpoint{2.539441in}{1.003585in}}%
\pgfpathlineto{\pgfqpoint{2.540887in}{1.005796in}}%
\pgfpathlineto{\pgfqpoint{2.541828in}{1.005237in}}%
\pgfpathlineto{\pgfqpoint{2.544106in}{1.002744in}}%
\pgfpathlineto{\pgfqpoint{2.544382in}{1.003568in}}%
\pgfpathlineto{\pgfqpoint{2.545721in}{1.005752in}}%
\pgfpathlineto{\pgfqpoint{2.546735in}{1.005267in}}%
\pgfpathlineto{\pgfqpoint{2.549053in}{1.002730in}}%
\pgfpathlineto{\pgfqpoint{2.549330in}{1.003562in}}%
\pgfpathlineto{\pgfqpoint{2.550745in}{1.005769in}}%
\pgfpathlineto{\pgfqpoint{2.551666in}{1.005277in}}%
\pgfpathlineto{\pgfqpoint{2.553998in}{1.002724in}}%
\pgfpathlineto{\pgfqpoint{2.554276in}{1.003559in}}%
\pgfpathlineto{\pgfqpoint{2.555730in}{1.005768in}}%
\pgfpathlineto{\pgfqpoint{2.556773in}{1.005085in}}%
\pgfpathlineto{\pgfqpoint{2.558944in}{1.002709in}}%
\pgfpathlineto{\pgfqpoint{2.559221in}{1.003541in}}%
\pgfpathlineto{\pgfqpoint{2.560641in}{1.005750in}}%
\pgfpathlineto{\pgfqpoint{2.561637in}{1.005167in}}%
\pgfpathlineto{\pgfqpoint{2.563906in}{1.002713in}}%
\pgfpathlineto{\pgfqpoint{2.564166in}{1.003524in}}%
\pgfpathlineto{\pgfqpoint{2.565523in}{1.005713in}}%
\pgfpathlineto{\pgfqpoint{2.566415in}{1.005335in}}%
\pgfpathlineto{\pgfqpoint{2.568836in}{1.002685in}}%
\pgfpathlineto{\pgfqpoint{2.569110in}{1.003504in}}%
\pgfpathlineto{\pgfqpoint{2.570461in}{1.005694in}}%
\pgfpathlineto{\pgfqpoint{2.571342in}{1.005340in}}%
\pgfpathlineto{\pgfqpoint{2.573798in}{1.002684in}}%
\pgfpathlineto{\pgfqpoint{2.574058in}{1.003502in}}%
\pgfpathlineto{\pgfqpoint{2.575553in}{1.005719in}}%
\pgfpathlineto{\pgfqpoint{2.576427in}{1.005174in}}%
\pgfpathlineto{\pgfqpoint{2.578743in}{1.002669in}}%
\pgfpathlineto{\pgfqpoint{2.579002in}{1.003478in}}%
\pgfpathlineto{\pgfqpoint{2.580349in}{1.005666in}}%
\pgfpathlineto{\pgfqpoint{2.581396in}{1.005134in}}%
\pgfpathlineto{\pgfqpoint{2.583671in}{1.002644in}}%
\pgfpathlineto{\pgfqpoint{2.583949in}{1.003475in}}%
\pgfpathlineto{\pgfqpoint{2.585368in}{1.005678in}}%
\pgfpathlineto{\pgfqpoint{2.586344in}{1.005115in}}%
\pgfpathlineto{\pgfqpoint{2.588617in}{1.002628in}}%
\pgfpathlineto{\pgfqpoint{2.588893in}{1.003453in}}%
\pgfpathlineto{\pgfqpoint{2.590260in}{1.005650in}}%
\pgfpathlineto{\pgfqpoint{2.591306in}{1.005091in}}%
\pgfpathlineto{\pgfqpoint{2.593561in}{1.002622in}}%
\pgfpathlineto{\pgfqpoint{2.593837in}{1.003448in}}%
\pgfpathlineto{\pgfqpoint{2.595208in}{1.005643in}}%
\pgfpathlineto{\pgfqpoint{2.596220in}{1.005113in}}%
\pgfpathlineto{\pgfqpoint{2.598508in}{1.002610in}}%
\pgfpathlineto{\pgfqpoint{2.598785in}{1.003441in}}%
\pgfpathlineto{\pgfqpoint{2.600197in}{1.005647in}}%
\pgfpathlineto{\pgfqpoint{2.601085in}{1.005194in}}%
\pgfpathlineto{\pgfqpoint{2.603470in}{1.002614in}}%
\pgfpathlineto{\pgfqpoint{2.603730in}{1.003426in}}%
\pgfpathlineto{\pgfqpoint{2.605095in}{1.005617in}}%
\pgfpathlineto{\pgfqpoint{2.606007in}{1.005205in}}%
\pgfpathlineto{\pgfqpoint{2.608415in}{1.002599in}}%
\pgfpathlineto{\pgfqpoint{2.608676in}{1.003413in}}%
\pgfpathlineto{\pgfqpoint{2.610069in}{1.005611in}}%
\pgfpathlineto{\pgfqpoint{2.610970in}{1.005169in}}%
\pgfpathlineto{\pgfqpoint{2.613362in}{1.002585in}}%
\pgfpathlineto{\pgfqpoint{2.613621in}{1.003400in}}%
\pgfpathlineto{\pgfqpoint{2.615081in}{1.005613in}}%
\pgfpathlineto{\pgfqpoint{2.616034in}{1.005025in}}%
\pgfpathlineto{\pgfqpoint{2.618307in}{1.002570in}}%
\pgfpathlineto{\pgfqpoint{2.618565in}{1.003377in}}%
\pgfpathlineto{\pgfqpoint{2.619916in}{1.005569in}}%
\pgfpathlineto{\pgfqpoint{2.620813in}{1.005197in}}%
\pgfpathlineto{\pgfqpoint{2.623234in}{1.002548in}}%
\pgfpathlineto{\pgfqpoint{2.623511in}{1.003379in}}%
\pgfpathlineto{\pgfqpoint{2.624921in}{1.005582in}}%
\pgfpathlineto{\pgfqpoint{2.625814in}{1.005126in}}%
\pgfpathlineto{\pgfqpoint{2.628180in}{1.002537in}}%
\pgfpathlineto{\pgfqpoint{2.628458in}{1.003369in}}%
\pgfpathlineto{\pgfqpoint{2.629880in}{1.005578in}}%
\pgfpathlineto{\pgfqpoint{2.630926in}{1.004939in}}%
\pgfpathlineto{\pgfqpoint{2.633125in}{1.002533in}}%
\pgfpathlineto{\pgfqpoint{2.633402in}{1.003363in}}%
\pgfpathlineto{\pgfqpoint{2.634803in}{1.005564in}}%
\pgfpathlineto{\pgfqpoint{2.635814in}{1.004993in}}%
\pgfpathlineto{\pgfqpoint{2.638076in}{1.002522in}}%
\pgfpathlineto{\pgfqpoint{2.638349in}{1.003350in}}%
\pgfpathlineto{\pgfqpoint{2.639974in}{1.005571in}}%
\pgfpathlineto{\pgfqpoint{2.640924in}{1.004800in}}%
\pgfpathlineto{\pgfqpoint{2.643033in}{1.002528in}}%
\pgfpathlineto{\pgfqpoint{2.643291in}{1.003331in}}%
\pgfpathlineto{\pgfqpoint{2.644609in}{1.005495in}}%
\pgfpathlineto{\pgfqpoint{2.645564in}{1.005101in}}%
\pgfpathlineto{\pgfqpoint{2.647966in}{1.002475in}}%
\pgfpathlineto{\pgfqpoint{2.648240in}{1.003305in}}%
\pgfpathlineto{\pgfqpoint{2.649839in}{1.005532in}}%
\pgfpathlineto{\pgfqpoint{2.650779in}{1.004800in}}%
\pgfpathlineto{\pgfqpoint{2.652923in}{1.002480in}}%
\pgfpathlineto{\pgfqpoint{2.653185in}{1.003298in}}%
\pgfpathlineto{\pgfqpoint{2.654577in}{1.005494in}}%
\pgfpathlineto{\pgfqpoint{2.655434in}{1.005102in}}%
\pgfpathlineto{\pgfqpoint{2.657859in}{1.002458in}}%
\pgfpathlineto{\pgfqpoint{2.658129in}{1.003280in}}%
\pgfpathlineto{\pgfqpoint{2.659647in}{1.005506in}}%
\pgfpathlineto{\pgfqpoint{2.660583in}{1.004868in}}%
\pgfpathlineto{\pgfqpoint{2.662807in}{1.002449in}}%
\pgfpathlineto{\pgfqpoint{2.663075in}{1.003274in}}%
\pgfpathlineto{\pgfqpoint{2.664797in}{1.005467in}}%
\pgfpathlineto{\pgfqpoint{2.666061in}{1.004262in}}%
\pgfpathlineto{\pgfqpoint{2.667746in}{1.002419in}}%
\pgfpathlineto{\pgfqpoint{2.668020in}{1.003245in}}%
\pgfpathlineto{\pgfqpoint{2.669483in}{1.005459in}}%
\pgfpathlineto{\pgfqpoint{2.670482in}{1.004819in}}%
\pgfpathlineto{\pgfqpoint{2.672689in}{1.002404in}}%
\pgfpathlineto{\pgfqpoint{2.672967in}{1.003236in}}%
\pgfpathlineto{\pgfqpoint{2.674391in}{1.005444in}}%
\pgfpathlineto{\pgfqpoint{2.675433in}{1.004806in}}%
\pgfpathlineto{\pgfqpoint{2.677652in}{1.002408in}}%
\pgfpathlineto{\pgfqpoint{2.677912in}{1.003222in}}%
\pgfpathlineto{\pgfqpoint{2.679302in}{1.005422in}}%
\pgfpathlineto{\pgfqpoint{2.680183in}{1.005010in}}%
\pgfpathlineto{\pgfqpoint{2.682597in}{1.002399in}}%
\pgfpathlineto{\pgfqpoint{2.682857in}{1.003212in}}%
\pgfpathlineto{\pgfqpoint{2.684246in}{1.005407in}}%
\pgfpathlineto{\pgfqpoint{2.685120in}{1.005002in}}%
\pgfpathlineto{\pgfqpoint{2.687532in}{1.002373in}}%
\pgfpathlineto{\pgfqpoint{2.687801in}{1.003193in}}%
\pgfpathlineto{\pgfqpoint{2.689307in}{1.005419in}}%
\pgfpathlineto{\pgfqpoint{2.690191in}{1.004853in}}%
\pgfpathlineto{\pgfqpoint{2.692475in}{1.002356in}}%
\pgfpathlineto{\pgfqpoint{2.692749in}{1.003184in}}%
\pgfpathlineto{\pgfqpoint{2.694363in}{1.005398in}}%
\pgfpathlineto{\pgfqpoint{2.695497in}{1.004433in}}%
\pgfpathlineto{\pgfqpoint{2.697433in}{1.002349in}}%
\pgfpathlineto{\pgfqpoint{2.697692in}{1.003156in}}%
\pgfpathlineto{\pgfqpoint{2.699013in}{1.005337in}}%
\pgfpathlineto{\pgfqpoint{2.699920in}{1.005001in}}%
\pgfpathlineto{\pgfqpoint{2.702367in}{1.002327in}}%
\pgfpathlineto{\pgfqpoint{2.702640in}{1.003156in}}%
\pgfpathlineto{\pgfqpoint{2.704259in}{1.005374in}}%
\pgfpathlineto{\pgfqpoint{2.705383in}{1.004416in}}%
\pgfpathlineto{\pgfqpoint{2.707307in}{1.002312in}}%
\pgfpathlineto{\pgfqpoint{2.707584in}{1.003143in}}%
\pgfpathlineto{\pgfqpoint{2.709056in}{1.005360in}}%
\pgfpathlineto{\pgfqpoint{2.709942in}{1.004831in}}%
\pgfpathlineto{\pgfqpoint{2.712261in}{1.002309in}}%
\pgfpathlineto{\pgfqpoint{2.712529in}{1.003132in}}%
\pgfpathlineto{\pgfqpoint{2.714072in}{1.005369in}}%
\pgfpathlineto{\pgfqpoint{2.715033in}{1.004677in}}%
\pgfpathlineto{\pgfqpoint{2.717215in}{1.002321in}}%
\pgfpathlineto{\pgfqpoint{2.717476in}{1.003138in}}%
\pgfpathlineto{\pgfqpoint{2.718912in}{1.005337in}}%
\pgfpathlineto{\pgfqpoint{2.719962in}{1.004669in}}%
\pgfpathlineto{\pgfqpoint{2.722148in}{1.002281in}}%
\pgfpathlineto{\pgfqpoint{2.722422in}{1.003110in}}%
\pgfpathlineto{\pgfqpoint{2.723996in}{1.005330in}}%
\pgfpathlineto{\pgfqpoint{2.725009in}{1.004540in}}%
\pgfpathlineto{\pgfqpoint{2.727105in}{1.002282in}}%
\pgfpathlineto{\pgfqpoint{2.727365in}{1.003092in}}%
\pgfpathlineto{\pgfqpoint{2.728689in}{1.005275in}}%
\pgfpathlineto{\pgfqpoint{2.729601in}{1.004928in}}%
\pgfpathlineto{\pgfqpoint{2.732050in}{1.002274in}}%
\pgfpathlineto{\pgfqpoint{2.732312in}{1.003094in}}%
\pgfpathlineto{\pgfqpoint{2.733698in}{1.005296in}}%
\pgfpathlineto{\pgfqpoint{2.734580in}{1.004890in}}%
\pgfpathlineto{\pgfqpoint{2.736989in}{1.002270in}}%
\pgfpathlineto{\pgfqpoint{2.737256in}{1.003091in}}%
\pgfpathlineto{\pgfqpoint{2.738840in}{1.005334in}}%
\pgfpathlineto{\pgfqpoint{2.739878in}{1.004516in}}%
\pgfpathlineto{\pgfqpoint{2.741943in}{1.002289in}}%
\pgfpathlineto{\pgfqpoint{2.742206in}{1.003110in}}%
\pgfpathlineto{\pgfqpoint{2.743762in}{1.005305in}}%
\pgfpathlineto{\pgfqpoint{2.744878in}{1.004413in}}%
\pgfpathlineto{\pgfqpoint{2.746870in}{1.002235in}}%
\pgfpathlineto{\pgfqpoint{2.747145in}{1.003056in}}%
\pgfpathlineto{\pgfqpoint{2.748464in}{1.005245in}}%
\pgfpathlineto{\pgfqpoint{2.749401in}{1.004886in}}%
\pgfpathlineto{\pgfqpoint{2.751824in}{1.002250in}}%
\pgfpathlineto{\pgfqpoint{2.752094in}{1.003076in}}%
\pgfpathlineto{\pgfqpoint{2.753703in}{1.005304in}}%
\pgfpathlineto{\pgfqpoint{2.754756in}{1.004437in}}%
\pgfpathlineto{\pgfqpoint{2.756761in}{1.002245in}}%
\pgfpathlineto{\pgfqpoint{2.757038in}{1.003075in}}%
\pgfpathlineto{\pgfqpoint{2.758443in}{1.005275in}}%
\pgfpathlineto{\pgfqpoint{2.759411in}{1.004746in}}%
\pgfpathlineto{\pgfqpoint{2.761710in}{1.002232in}}%
\pgfpathlineto{\pgfqpoint{2.761985in}{1.003060in}}%
\pgfpathlineto{\pgfqpoint{2.763439in}{1.005277in}}%
\pgfpathlineto{\pgfqpoint{2.764478in}{1.004606in}}%
\pgfpathlineto{\pgfqpoint{2.766660in}{1.002232in}}%
\pgfpathlineto{\pgfqpoint{2.766930in}{1.003058in}}%
\pgfpathlineto{\pgfqpoint{2.768479in}{1.005287in}}%
\pgfpathlineto{\pgfqpoint{2.769416in}{1.004612in}}%
\pgfpathlineto{\pgfqpoint{2.771600in}{1.002224in}}%
\pgfpathlineto{\pgfqpoint{2.771876in}{1.003054in}}%
\pgfpathlineto{\pgfqpoint{2.773452in}{1.005264in}}%
\pgfpathlineto{\pgfqpoint{2.774576in}{1.004348in}}%
\pgfpathlineto{\pgfqpoint{2.776551in}{1.002201in}}%
\pgfpathlineto{\pgfqpoint{2.776820in}{1.003025in}}%
\pgfpathlineto{\pgfqpoint{2.778326in}{1.005259in}}%
\pgfpathlineto{\pgfqpoint{2.779276in}{1.004623in}}%
\pgfpathlineto{\pgfqpoint{2.781497in}{1.002208in}}%
\pgfpathlineto{\pgfqpoint{2.781764in}{1.003029in}}%
\pgfpathlineto{\pgfqpoint{2.783440in}{1.005247in}}%
\pgfpathlineto{\pgfqpoint{2.784696in}{1.004095in}}%
\pgfpathlineto{\pgfqpoint{2.786444in}{1.002203in}}%
\pgfpathlineto{\pgfqpoint{2.786710in}{1.003026in}}%
\pgfpathlineto{\pgfqpoint{2.788365in}{1.005253in}}%
\pgfpathlineto{\pgfqpoint{2.789453in}{1.004306in}}%
\pgfpathlineto{\pgfqpoint{2.791396in}{1.002210in}}%
\pgfpathlineto{\pgfqpoint{2.791658in}{1.003027in}}%
\pgfpathlineto{\pgfqpoint{2.793052in}{1.005229in}}%
\pgfpathlineto{\pgfqpoint{2.793936in}{1.004807in}}%
\pgfpathlineto{\pgfqpoint{2.796342in}{1.002208in}}%
\pgfpathlineto{\pgfqpoint{2.796601in}{1.003014in}}%
\pgfpathlineto{\pgfqpoint{2.797925in}{1.005193in}}%
\pgfpathlineto{\pgfqpoint{2.798826in}{1.004858in}}%
\pgfpathlineto{\pgfqpoint{2.801272in}{1.002182in}}%
\pgfpathlineto{\pgfqpoint{2.801548in}{1.003005in}}%
\pgfpathlineto{\pgfqpoint{2.802882in}{1.005193in}}%
\pgfpathlineto{\pgfqpoint{2.803853in}{1.004766in}}%
\pgfpathlineto{\pgfqpoint{2.806220in}{1.002178in}}%
\pgfpathlineto{\pgfqpoint{2.806495in}{1.003008in}}%
\pgfpathlineto{\pgfqpoint{2.808110in}{1.005220in}}%
\pgfpathlineto{\pgfqpoint{2.809230in}{1.004270in}}%
\pgfpathlineto{\pgfqpoint{2.811180in}{1.002171in}}%
\pgfpathlineto{\pgfqpoint{2.811440in}{1.002990in}}%
\pgfpathlineto{\pgfqpoint{2.812933in}{1.005211in}}%
\pgfpathlineto{\pgfqpoint{2.813839in}{1.004636in}}%
\pgfpathlineto{\pgfqpoint{2.816111in}{1.002153in}}%
\pgfpathlineto{\pgfqpoint{2.816385in}{1.002981in}}%
\pgfpathlineto{\pgfqpoint{2.817949in}{1.005198in}}%
\pgfpathlineto{\pgfqpoint{2.818909in}{1.004479in}}%
\pgfpathlineto{\pgfqpoint{2.821071in}{1.002148in}}%
\pgfpathlineto{\pgfqpoint{2.821331in}{1.002965in}}%
\pgfpathlineto{\pgfqpoint{2.822823in}{1.005186in}}%
\pgfpathlineto{\pgfqpoint{2.823752in}{1.004586in}}%
\pgfpathlineto{\pgfqpoint{2.826007in}{1.002136in}}%
\pgfpathlineto{\pgfqpoint{2.826274in}{1.002958in}}%
\pgfpathlineto{\pgfqpoint{2.827859in}{1.005201in}}%
\pgfpathlineto{\pgfqpoint{2.828933in}{1.004341in}}%
\pgfpathlineto{\pgfqpoint{2.830961in}{1.002154in}}%
\pgfpathlineto{\pgfqpoint{2.831220in}{1.002962in}}%
\pgfpathlineto{\pgfqpoint{2.832578in}{1.005143in}}%
\pgfpathlineto{\pgfqpoint{2.833448in}{1.004785in}}%
\pgfpathlineto{\pgfqpoint{2.835898in}{1.002121in}}%
\pgfpathlineto{\pgfqpoint{2.836163in}{1.002939in}}%
\pgfpathlineto{\pgfqpoint{2.837683in}{1.005189in}}%
\pgfpathlineto{\pgfqpoint{2.838593in}{1.004590in}}%
\pgfpathlineto{\pgfqpoint{2.840843in}{1.002143in}}%
\pgfpathlineto{\pgfqpoint{2.841113in}{1.002966in}}%
\pgfpathlineto{\pgfqpoint{2.842899in}{1.005120in}}%
\pgfpathlineto{\pgfqpoint{2.844486in}{1.003510in}}%
\pgfpathlineto{\pgfqpoint{2.845797in}{1.002104in}}%
\pgfpathlineto{\pgfqpoint{2.846060in}{1.002929in}}%
\pgfpathlineto{\pgfqpoint{2.847581in}{1.005150in}}%
\pgfpathlineto{\pgfqpoint{2.848509in}{1.004516in}}%
\pgfpathlineto{\pgfqpoint{2.850731in}{1.002091in}}%
\pgfpathlineto{\pgfqpoint{2.851003in}{1.002917in}}%
\pgfpathlineto{\pgfqpoint{2.852617in}{1.005130in}}%
\pgfpathlineto{\pgfqpoint{2.853693in}{1.004230in}}%
\pgfpathlineto{\pgfqpoint{2.855669in}{1.002070in}}%
\pgfpathlineto{\pgfqpoint{2.855945in}{1.002895in}}%
\pgfpathlineto{\pgfqpoint{2.857271in}{1.005076in}}%
\pgfpathlineto{\pgfqpoint{2.858160in}{1.004752in}}%
\pgfpathlineto{\pgfqpoint{2.860624in}{1.002065in}}%
\pgfpathlineto{\pgfqpoint{2.860892in}{1.002884in}}%
\pgfpathlineto{\pgfqpoint{2.862365in}{1.005121in}}%
\pgfpathlineto{\pgfqpoint{2.863359in}{1.004484in}}%
\pgfpathlineto{\pgfqpoint{2.865579in}{1.002086in}}%
\pgfpathlineto{\pgfqpoint{2.865842in}{1.002905in}}%
\pgfpathlineto{\pgfqpoint{2.867397in}{1.005090in}}%
\pgfpathlineto{\pgfqpoint{2.868517in}{1.004191in}}%
\pgfpathlineto{\pgfqpoint{2.870518in}{1.002028in}}%
\pgfpathlineto{\pgfqpoint{2.870782in}{1.002848in}}%
\pgfpathlineto{\pgfqpoint{2.872192in}{1.005082in}}%
\pgfpathlineto{\pgfqpoint{2.873168in}{1.004551in}}%
\pgfpathlineto{\pgfqpoint{2.875462in}{1.002062in}}%
\pgfpathlineto{\pgfqpoint{2.875729in}{1.002886in}}%
\pgfpathlineto{\pgfqpoint{2.877338in}{1.005109in}}%
\pgfpathlineto{\pgfqpoint{2.878401in}{1.004230in}}%
\pgfpathlineto{\pgfqpoint{2.880409in}{1.002063in}}%
\pgfpathlineto{\pgfqpoint{2.880671in}{1.002876in}}%
\pgfpathlineto{\pgfqpoint{2.882078in}{1.005103in}}%
\pgfpathlineto{\pgfqpoint{2.882993in}{1.004640in}}%
\pgfpathlineto{\pgfqpoint{2.885343in}{1.002071in}}%
\pgfpathlineto{\pgfqpoint{2.885618in}{1.002889in}}%
\pgfpathlineto{\pgfqpoint{2.886938in}{1.005055in}}%
\pgfpathlineto{\pgfqpoint{2.887864in}{1.004693in}}%
\pgfpathlineto{\pgfqpoint{2.890295in}{1.002043in}}%
\pgfpathlineto{\pgfqpoint{2.890565in}{1.002866in}}%
\pgfpathlineto{\pgfqpoint{2.892083in}{1.005107in}}%
\pgfpathlineto{\pgfqpoint{2.893110in}{1.004377in}}%
\pgfpathlineto{\pgfqpoint{2.895251in}{1.002071in}}%
\pgfpathlineto{\pgfqpoint{2.895511in}{1.002882in}}%
\pgfpathlineto{\pgfqpoint{2.896899in}{1.005055in}}%
\pgfpathlineto{\pgfqpoint{2.897786in}{1.004625in}}%
\pgfpathlineto{\pgfqpoint{2.900188in}{1.002011in}}%
\pgfpathlineto{\pgfqpoint{2.900455in}{1.002840in}}%
\pgfpathlineto{\pgfqpoint{2.902209in}{1.005065in}}%
\pgfpathlineto{\pgfqpoint{2.903492in}{1.003827in}}%
\pgfpathlineto{\pgfqpoint{2.905127in}{1.002041in}}%
\pgfpathlineto{\pgfqpoint{2.905404in}{1.002867in}}%
\pgfpathlineto{\pgfqpoint{2.906814in}{1.005046in}}%
\pgfpathlineto{\pgfqpoint{2.907676in}{1.004614in}}%
\pgfpathlineto{\pgfqpoint{2.910088in}{1.002007in}}%
\pgfpathlineto{\pgfqpoint{2.910348in}{1.002823in}}%
\pgfpathlineto{\pgfqpoint{2.911736in}{1.005034in}}%
\pgfpathlineto{\pgfqpoint{2.912610in}{1.004641in}}%
\pgfpathlineto{\pgfqpoint{2.915034in}{1.002021in}}%
\pgfpathlineto{\pgfqpoint{2.915293in}{1.002829in}}%
\pgfpathlineto{\pgfqpoint{2.916645in}{1.005012in}}%
\pgfpathlineto{\pgfqpoint{2.917656in}{1.004511in}}%
\pgfpathlineto{\pgfqpoint{2.919967in}{1.001990in}}%
\pgfpathlineto{\pgfqpoint{2.920240in}{1.002818in}}%
\pgfpathlineto{\pgfqpoint{2.921874in}{1.005041in}}%
\pgfpathlineto{\pgfqpoint{2.922900in}{1.004180in}}%
\pgfpathlineto{\pgfqpoint{2.924913in}{1.001986in}}%
\pgfpathlineto{\pgfqpoint{2.925185in}{1.002811in}}%
\pgfpathlineto{\pgfqpoint{2.926806in}{1.005021in}}%
\pgfpathlineto{\pgfqpoint{2.927859in}{1.004140in}}%
\pgfpathlineto{\pgfqpoint{2.929856in}{1.001958in}}%
\pgfpathlineto{\pgfqpoint{2.930131in}{1.002787in}}%
\pgfpathlineto{\pgfqpoint{2.931685in}{1.004993in}}%
\pgfpathlineto{\pgfqpoint{2.932706in}{1.004213in}}%
\pgfpathlineto{\pgfqpoint{2.934815in}{1.001941in}}%
\pgfpathlineto{\pgfqpoint{2.935073in}{1.002747in}}%
\pgfpathlineto{\pgfqpoint{2.936393in}{1.004938in}}%
\pgfpathlineto{\pgfqpoint{2.937298in}{1.004613in}}%
\pgfpathlineto{\pgfqpoint{2.939748in}{1.001938in}}%
\pgfpathlineto{\pgfqpoint{2.940021in}{1.002764in}}%
\pgfpathlineto{\pgfqpoint{2.941626in}{1.004982in}}%
\pgfpathlineto{\pgfqpoint{2.942658in}{1.004141in}}%
\pgfpathlineto{\pgfqpoint{2.944695in}{1.001922in}}%
\pgfpathlineto{\pgfqpoint{2.944967in}{1.002748in}}%
\pgfpathlineto{\pgfqpoint{2.946605in}{1.004964in}}%
\pgfpathlineto{\pgfqpoint{2.947624in}{1.004107in}}%
\pgfpathlineto{\pgfqpoint{2.949633in}{1.001911in}}%
\pgfpathlineto{\pgfqpoint{2.949910in}{1.002735in}}%
\pgfpathlineto{\pgfqpoint{2.951260in}{1.004905in}}%
\pgfpathlineto{\pgfqpoint{2.952163in}{1.004518in}}%
\pgfpathlineto{\pgfqpoint{2.954579in}{1.001877in}}%
\pgfpathlineto{\pgfqpoint{2.954856in}{1.002707in}}%
\pgfpathlineto{\pgfqpoint{2.956235in}{1.004917in}}%
\pgfpathlineto{\pgfqpoint{2.957171in}{1.004469in}}%
\pgfpathlineto{\pgfqpoint{2.959542in}{1.001908in}}%
\pgfpathlineto{\pgfqpoint{2.959802in}{1.002720in}}%
\pgfpathlineto{\pgfqpoint{2.961162in}{1.004908in}}%
\pgfpathlineto{\pgfqpoint{2.962062in}{1.004518in}}%
\pgfpathlineto{\pgfqpoint{2.964476in}{1.001884in}}%
\pgfpathlineto{\pgfqpoint{2.964749in}{1.002711in}}%
\pgfpathlineto{\pgfqpoint{2.966357in}{1.004931in}}%
\pgfpathlineto{\pgfqpoint{2.967384in}{1.004093in}}%
\pgfpathlineto{\pgfqpoint{2.969423in}{1.001871in}}%
\pgfpathlineto{\pgfqpoint{2.969694in}{1.002692in}}%
\pgfpathlineto{\pgfqpoint{2.971229in}{1.004928in}}%
\pgfpathlineto{\pgfqpoint{2.972159in}{1.004281in}}%
\pgfpathlineto{\pgfqpoint{2.974370in}{1.001878in}}%
\pgfpathlineto{\pgfqpoint{2.974638in}{1.002703in}}%
\pgfpathlineto{\pgfqpoint{2.976395in}{1.004901in}}%
\pgfpathlineto{\pgfqpoint{2.977727in}{1.003600in}}%
\pgfpathlineto{\pgfqpoint{2.979324in}{1.001885in}}%
\pgfpathlineto{\pgfqpoint{2.979587in}{1.002708in}}%
\pgfpathlineto{\pgfqpoint{2.981135in}{1.004918in}}%
\pgfpathlineto{\pgfqpoint{2.982056in}{1.004256in}}%
\pgfpathlineto{\pgfqpoint{2.984261in}{1.001862in}}%
\pgfpathlineto{\pgfqpoint{2.984530in}{1.002690in}}%
\pgfpathlineto{\pgfqpoint{2.986174in}{1.004923in}}%
\pgfpathlineto{\pgfqpoint{2.987290in}{1.003957in}}%
\pgfpathlineto{\pgfqpoint{2.989197in}{1.001874in}}%
\pgfpathlineto{\pgfqpoint{2.989474in}{1.002700in}}%
\pgfpathlineto{\pgfqpoint{2.990828in}{1.004879in}}%
\pgfpathlineto{\pgfqpoint{2.991723in}{1.004498in}}%
\pgfpathlineto{\pgfqpoint{2.994154in}{1.001865in}}%
\pgfpathlineto{\pgfqpoint{2.994418in}{1.002684in}}%
\pgfpathlineto{\pgfqpoint{2.995828in}{1.004917in}}%
\pgfpathlineto{\pgfqpoint{2.996820in}{1.004370in}}%
\pgfpathlineto{\pgfqpoint{2.999106in}{1.001901in}}%
\pgfpathlineto{\pgfqpoint{2.999366in}{1.002710in}}%
\pgfpathlineto{\pgfqpoint{3.000728in}{1.004888in}}%
\pgfpathlineto{\pgfqpoint{3.001602in}{1.004517in}}%
\pgfpathlineto{\pgfqpoint{3.004041in}{1.001860in}}%
\pgfpathlineto{\pgfqpoint{3.004309in}{1.002679in}}%
\pgfpathlineto{\pgfqpoint{3.005789in}{1.004913in}}%
\pgfpathlineto{\pgfqpoint{3.006775in}{1.004274in}}%
\pgfpathlineto{\pgfqpoint{3.008988in}{1.001872in}}%
\pgfpathlineto{\pgfqpoint{3.009255in}{1.002690in}}%
\pgfpathlineto{\pgfqpoint{3.010788in}{1.004908in}}%
\pgfpathlineto{\pgfqpoint{3.011786in}{1.004182in}}%
\pgfpathlineto{\pgfqpoint{3.013941in}{1.001852in}}%
\pgfpathlineto{\pgfqpoint{3.014203in}{1.002671in}}%
\pgfpathlineto{\pgfqpoint{3.015590in}{1.004871in}}%
\pgfpathlineto{\pgfqpoint{3.016505in}{1.004429in}}%
\pgfpathlineto{\pgfqpoint{3.018875in}{1.001840in}}%
\pgfpathlineto{\pgfqpoint{3.019149in}{1.002668in}}%
\pgfpathlineto{\pgfqpoint{3.020716in}{1.004884in}}%
\pgfpathlineto{\pgfqpoint{3.021633in}{1.004209in}}%
\pgfpathlineto{\pgfqpoint{3.023821in}{1.001819in}}%
\pgfpathlineto{\pgfqpoint{3.024096in}{1.002650in}}%
\pgfpathlineto{\pgfqpoint{3.025683in}{1.004869in}}%
\pgfpathlineto{\pgfqpoint{3.026578in}{1.004197in}}%
\pgfpathlineto{\pgfqpoint{3.028770in}{1.001815in}}%
\pgfpathlineto{\pgfqpoint{3.029039in}{1.002644in}}%
\pgfpathlineto{\pgfqpoint{3.030843in}{1.004831in}}%
\pgfpathlineto{\pgfqpoint{3.032331in}{1.003330in}}%
\pgfpathlineto{\pgfqpoint{3.033706in}{1.001828in}}%
\pgfpathlineto{\pgfqpoint{3.033986in}{1.002661in}}%
\pgfpathlineto{\pgfqpoint{3.035443in}{1.004846in}}%
\pgfpathlineto{\pgfqpoint{3.036431in}{1.004213in}}%
\pgfpathlineto{\pgfqpoint{3.038660in}{1.001786in}}%
\pgfpathlineto{\pgfqpoint{3.038928in}{1.002608in}}%
\pgfpathlineto{\pgfqpoint{3.040495in}{1.004868in}}%
\pgfpathlineto{\pgfqpoint{3.041391in}{1.004232in}}%
\pgfpathlineto{\pgfqpoint{3.043608in}{1.001830in}}%
\pgfpathlineto{\pgfqpoint{3.043876in}{1.002653in}}%
\pgfpathlineto{\pgfqpoint{3.045582in}{1.004843in}}%
\pgfpathlineto{\pgfqpoint{3.046882in}{1.003610in}}%
\pgfpathlineto{\pgfqpoint{3.048542in}{1.001797in}}%
\pgfpathlineto{\pgfqpoint{3.048820in}{1.002626in}}%
\pgfpathlineto{\pgfqpoint{3.050210in}{1.004830in}}%
\pgfpathlineto{\pgfqpoint{3.051168in}{1.004335in}}%
\pgfpathlineto{\pgfqpoint{3.053489in}{1.001797in}}%
\pgfpathlineto{\pgfqpoint{3.053764in}{1.002623in}}%
\pgfpathlineto{\pgfqpoint{3.055154in}{1.004828in}}%
\pgfpathlineto{\pgfqpoint{3.056108in}{1.004339in}}%
\pgfpathlineto{\pgfqpoint{3.058442in}{1.001801in}}%
\pgfpathlineto{\pgfqpoint{3.058711in}{1.002625in}}%
\pgfpathlineto{\pgfqpoint{3.060225in}{1.004858in}}%
\pgfpathlineto{\pgfqpoint{3.061170in}{1.004220in}}%
\pgfpathlineto{\pgfqpoint{3.063381in}{1.001803in}}%
\pgfpathlineto{\pgfqpoint{3.063660in}{1.002633in}}%
\pgfpathlineto{\pgfqpoint{3.065111in}{1.004822in}}%
\pgfpathlineto{\pgfqpoint{3.066113in}{1.004180in}}%
\pgfpathlineto{\pgfqpoint{3.068332in}{1.001766in}}%
\pgfpathlineto{\pgfqpoint{3.068602in}{1.002590in}}%
\pgfpathlineto{\pgfqpoint{3.070120in}{1.004835in}}%
\pgfpathlineto{\pgfqpoint{3.071015in}{1.004251in}}%
\pgfpathlineto{\pgfqpoint{3.073271in}{1.001786in}}%
\pgfpathlineto{\pgfqpoint{3.073549in}{1.002616in}}%
\pgfpathlineto{\pgfqpoint{3.074992in}{1.004806in}}%
\pgfpathlineto{\pgfqpoint{3.075877in}{1.004305in}}%
\pgfpathlineto{\pgfqpoint{3.078232in}{1.001758in}}%
\pgfpathlineto{\pgfqpoint{3.078492in}{1.002575in}}%
\pgfpathlineto{\pgfqpoint{3.079842in}{1.004779in}}%
\pgfpathlineto{\pgfqpoint{3.080835in}{1.004314in}}%
\pgfpathlineto{\pgfqpoint{3.083166in}{1.001769in}}%
\pgfpathlineto{\pgfqpoint{3.083440in}{1.002596in}}%
\pgfpathlineto{\pgfqpoint{3.085050in}{1.004806in}}%
\pgfpathlineto{\pgfqpoint{3.086076in}{1.003964in}}%
\pgfpathlineto{\pgfqpoint{3.088114in}{1.001747in}}%
\pgfpathlineto{\pgfqpoint{3.088383in}{1.002568in}}%
\pgfpathlineto{\pgfqpoint{3.089891in}{1.004804in}}%
\pgfpathlineto{\pgfqpoint{3.090806in}{1.004208in}}%
\pgfpathlineto{\pgfqpoint{3.093057in}{1.001752in}}%
\pgfpathlineto{\pgfqpoint{3.093330in}{1.002576in}}%
\pgfpathlineto{\pgfqpoint{3.094910in}{1.004783in}}%
\pgfpathlineto{\pgfqpoint{3.095817in}{1.004101in}}%
\pgfpathlineto{\pgfqpoint{3.098004in}{1.001718in}}%
\pgfpathlineto{\pgfqpoint{3.098276in}{1.002547in}}%
\pgfpathlineto{\pgfqpoint{3.099934in}{1.004770in}}%
\pgfpathlineto{\pgfqpoint{3.100934in}{1.003916in}}%
\pgfpathlineto{\pgfqpoint{3.102947in}{1.001717in}}%
\pgfpathlineto{\pgfqpoint{3.103222in}{1.002547in}}%
\pgfpathlineto{\pgfqpoint{3.104798in}{1.004761in}}%
\pgfpathlineto{\pgfqpoint{3.105900in}{1.003872in}}%
\pgfpathlineto{\pgfqpoint{3.107889in}{1.001699in}}%
\pgfpathlineto{\pgfqpoint{3.108167in}{1.002530in}}%
\pgfpathlineto{\pgfqpoint{3.109578in}{1.004742in}}%
\pgfpathlineto{\pgfqpoint{3.110449in}{1.004316in}}%
\pgfpathlineto{\pgfqpoint{3.112851in}{1.001724in}}%
\pgfpathlineto{\pgfqpoint{3.113111in}{1.002534in}}%
\pgfpathlineto{\pgfqpoint{3.114452in}{1.004718in}}%
\pgfpathlineto{\pgfqpoint{3.115465in}{1.004234in}}%
\pgfpathlineto{\pgfqpoint{3.117797in}{1.001716in}}%
\pgfpathlineto{\pgfqpoint{3.118056in}{1.002526in}}%
\pgfpathlineto{\pgfqpoint{3.119410in}{1.004720in}}%
\pgfpathlineto{\pgfqpoint{3.120387in}{1.004260in}}%
\pgfpathlineto{\pgfqpoint{3.122744in}{1.001713in}}%
\pgfpathlineto{\pgfqpoint{3.123006in}{1.002534in}}%
\pgfpathlineto{\pgfqpoint{3.124556in}{1.004750in}}%
\pgfpathlineto{\pgfqpoint{3.125444in}{1.004125in}}%
\pgfpathlineto{\pgfqpoint{3.127687in}{1.001708in}}%
\pgfpathlineto{\pgfqpoint{3.127947in}{1.002516in}}%
\pgfpathlineto{\pgfqpoint{3.129268in}{1.004690in}}%
\pgfpathlineto{\pgfqpoint{3.130186in}{1.004341in}}%
\pgfpathlineto{\pgfqpoint{3.132633in}{1.001699in}}%
\pgfpathlineto{\pgfqpoint{3.132892in}{1.002508in}}%
\pgfpathlineto{\pgfqpoint{3.134223in}{1.004697in}}%
\pgfpathlineto{\pgfqpoint{3.135164in}{1.004312in}}%
\pgfpathlineto{\pgfqpoint{3.137578in}{1.001703in}}%
\pgfpathlineto{\pgfqpoint{3.137840in}{1.002522in}}%
\pgfpathlineto{\pgfqpoint{3.139282in}{1.004734in}}%
\pgfpathlineto{\pgfqpoint{3.140266in}{1.004138in}}%
\pgfpathlineto{\pgfqpoint{3.142523in}{1.001698in}}%
\pgfpathlineto{\pgfqpoint{3.142784in}{1.002513in}}%
\pgfpathlineto{\pgfqpoint{3.144134in}{1.004705in}}%
\pgfpathlineto{\pgfqpoint{3.145140in}{1.004218in}}%
\pgfpathlineto{\pgfqpoint{3.147469in}{1.001703in}}%
\pgfpathlineto{\pgfqpoint{3.147730in}{1.002519in}}%
\pgfpathlineto{\pgfqpoint{3.149122in}{1.004719in}}%
\pgfpathlineto{\pgfqpoint{3.150014in}{1.004293in}}%
\pgfpathlineto{\pgfqpoint{3.152414in}{1.001696in}}%
\pgfpathlineto{\pgfqpoint{3.152674in}{1.002507in}}%
\pgfpathlineto{\pgfqpoint{3.154009in}{1.004695in}}%
\pgfpathlineto{\pgfqpoint{3.154909in}{1.004349in}}%
\pgfpathlineto{\pgfqpoint{3.157359in}{1.001699in}}%
\pgfpathlineto{\pgfqpoint{3.157620in}{1.002513in}}%
\pgfpathlineto{\pgfqpoint{3.158973in}{1.004701in}}%
\pgfpathlineto{\pgfqpoint{3.159914in}{1.004279in}}%
\pgfpathlineto{\pgfqpoint{3.162296in}{1.001687in}}%
\pgfpathlineto{\pgfqpoint{3.162565in}{1.002509in}}%
\pgfpathlineto{\pgfqpoint{3.164051in}{1.004741in}}%
\pgfpathlineto{\pgfqpoint{3.165037in}{1.004094in}}%
\pgfpathlineto{\pgfqpoint{3.167239in}{1.001691in}}%
\pgfpathlineto{\pgfqpoint{3.167516in}{1.002524in}}%
\pgfpathlineto{\pgfqpoint{3.169372in}{1.004621in}}%
\pgfpathlineto{\pgfqpoint{3.169467in}{1.004560in}}%
\pgfpathlineto{\pgfqpoint{3.170132in}{1.003871in}}%
\pgfpathlineto{\pgfqpoint{3.170132in}{1.003871in}}%
\pgfusepath{stroke}%
\end{pgfscope}%
\begin{pgfscope}%
\pgfsetrectcap%
\pgfsetmiterjoin%
\pgfsetlinewidth{0.803000pt}%
\definecolor{currentstroke}{rgb}{0.000000,0.000000,0.000000}%
\pgfsetstrokecolor{currentstroke}%
\pgfsetdash{}{0pt}%
\pgfpathmoveto{\pgfqpoint{0.573768in}{0.462654in}}%
\pgfpathlineto{\pgfqpoint{0.573768in}{1.368904in}}%
\pgfusepath{stroke}%
\end{pgfscope}%
\begin{pgfscope}%
\pgfsetrectcap%
\pgfsetmiterjoin%
\pgfsetlinewidth{0.803000pt}%
\definecolor{currentstroke}{rgb}{0.000000,0.000000,0.000000}%
\pgfsetstrokecolor{currentstroke}%
\pgfsetdash{}{0pt}%
\pgfpathmoveto{\pgfqpoint{3.293768in}{0.462654in}}%
\pgfpathlineto{\pgfqpoint{3.293768in}{1.368904in}}%
\pgfusepath{stroke}%
\end{pgfscope}%
\begin{pgfscope}%
\pgfsetrectcap%
\pgfsetmiterjoin%
\pgfsetlinewidth{0.803000pt}%
\definecolor{currentstroke}{rgb}{0.000000,0.000000,0.000000}%
\pgfsetstrokecolor{currentstroke}%
\pgfsetdash{}{0pt}%
\pgfpathmoveto{\pgfqpoint{0.573768in}{0.462654in}}%
\pgfpathlineto{\pgfqpoint{3.293768in}{0.462654in}}%
\pgfusepath{stroke}%
\end{pgfscope}%
\begin{pgfscope}%
\pgfsetrectcap%
\pgfsetmiterjoin%
\pgfsetlinewidth{0.803000pt}%
\definecolor{currentstroke}{rgb}{0.000000,0.000000,0.000000}%
\pgfsetstrokecolor{currentstroke}%
\pgfsetdash{}{0pt}%
\pgfpathmoveto{\pgfqpoint{0.573768in}{1.368904in}}%
\pgfpathlineto{\pgfqpoint{3.293768in}{1.368904in}}%
\pgfusepath{stroke}%
\end{pgfscope}%
\begin{pgfscope}%
\definecolor{textcolor}{rgb}{0.000000,0.000000,0.000000}%
\pgfsetstrokecolor{textcolor}%
\pgfsetfillcolor{textcolor}%
\pgftext[x=0.628168in,y=1.305467in,left,top]{\color{textcolor}{\rmfamily\fontsize{8.000000}{9.600000}\selectfont\catcode`\^=\active\def^{\ifmmode\sp\else\^{}\fi}\catcode`\%=\active\def%{\%}(c)}}%
\end{pgfscope}%
\begin{pgfscope}%
\pgfsetbuttcap%
\pgfsetmiterjoin%
\definecolor{currentfill}{rgb}{1.000000,1.000000,1.000000}%
\pgfsetfillcolor{currentfill}%
\pgfsetlinewidth{0.000000pt}%
\definecolor{currentstroke}{rgb}{0.000000,0.000000,0.000000}%
\pgfsetstrokecolor{currentstroke}%
\pgfsetstrokeopacity{0.000000}%
\pgfsetdash{}{0pt}%
\pgfpathmoveto{\pgfqpoint{2.197902in}{3.064738in}}%
\pgfpathlineto{\pgfqpoint{3.177102in}{3.064738in}}%
\pgfpathlineto{\pgfqpoint{3.177102in}{3.390988in}}%
\pgfpathlineto{\pgfqpoint{2.197902in}{3.390988in}}%
\pgfpathlineto{\pgfqpoint{2.197902in}{3.064738in}}%
\pgfpathclose%
\pgfusepath{fill}%
\end{pgfscope}%
\begin{pgfscope}%
\pgfpathrectangle{\pgfqpoint{2.197902in}{3.064738in}}{\pgfqpoint{0.979200in}{0.326250in}}%
\pgfusepath{clip}%
\pgfsetbuttcap%
\pgfsetroundjoin%
\pgfsetlinewidth{0.401500pt}%
\definecolor{currentstroke}{rgb}{0.690196,0.690196,0.690196}%
\pgfsetstrokecolor{currentstroke}%
\pgfsetstrokeopacity{0.250000}%
\pgfsetdash{{1.480000pt}{0.640000pt}}{0.000000pt}%
\pgfpathmoveto{\pgfqpoint{2.197902in}{3.064738in}}%
\pgfpathlineto{\pgfqpoint{2.197902in}{3.390988in}}%
\pgfusepath{stroke}%
\end{pgfscope}%
\begin{pgfscope}%
\pgfsetbuttcap%
\pgfsetroundjoin%
\definecolor{currentfill}{rgb}{0.000000,0.000000,0.000000}%
\pgfsetfillcolor{currentfill}%
\pgfsetlinewidth{0.803000pt}%
\definecolor{currentstroke}{rgb}{0.000000,0.000000,0.000000}%
\pgfsetstrokecolor{currentstroke}%
\pgfsetdash{}{0pt}%
\pgfsys@defobject{currentmarker}{\pgfqpoint{0.000000in}{0.000000in}}{\pgfqpoint{0.000000in}{0.027778in}}{%
\pgfpathmoveto{\pgfqpoint{0.000000in}{0.000000in}}%
\pgfpathlineto{\pgfqpoint{0.000000in}{0.027778in}}%
\pgfusepath{stroke,fill}%
}%
\begin{pgfscope}%
\pgfsys@transformshift{2.197902in}{3.064738in}%
\pgfsys@useobject{currentmarker}{}%
\end{pgfscope}%
\end{pgfscope}%
\begin{pgfscope}%
\definecolor{textcolor}{rgb}{0.000000,0.000000,0.000000}%
\pgfsetstrokecolor{textcolor}%
\pgfsetfillcolor{textcolor}%
\pgftext[x=2.197902in,y=3.050849in,,top]{\color{textcolor}{\rmfamily\fontsize{5.000000}{6.000000}\selectfont\catcode`\^=\active\def^{\ifmmode\sp\else\^{}\fi}\catcode`\%=\active\def%{\%}4.980}}%
\end{pgfscope}%
\begin{pgfscope}%
\pgfpathrectangle{\pgfqpoint{2.197902in}{3.064738in}}{\pgfqpoint{0.979200in}{0.326250in}}%
\pgfusepath{clip}%
\pgfsetbuttcap%
\pgfsetroundjoin%
\pgfsetlinewidth{0.401500pt}%
\definecolor{currentstroke}{rgb}{0.690196,0.690196,0.690196}%
\pgfsetstrokecolor{currentstroke}%
\pgfsetstrokeopacity{0.250000}%
\pgfsetdash{{1.480000pt}{0.640000pt}}{0.000000pt}%
\pgfpathmoveto{\pgfqpoint{2.442702in}{3.064738in}}%
\pgfpathlineto{\pgfqpoint{2.442702in}{3.390988in}}%
\pgfusepath{stroke}%
\end{pgfscope}%
\begin{pgfscope}%
\pgfsetbuttcap%
\pgfsetroundjoin%
\definecolor{currentfill}{rgb}{0.000000,0.000000,0.000000}%
\pgfsetfillcolor{currentfill}%
\pgfsetlinewidth{0.803000pt}%
\definecolor{currentstroke}{rgb}{0.000000,0.000000,0.000000}%
\pgfsetstrokecolor{currentstroke}%
\pgfsetdash{}{0pt}%
\pgfsys@defobject{currentmarker}{\pgfqpoint{0.000000in}{0.000000in}}{\pgfqpoint{0.000000in}{0.027778in}}{%
\pgfpathmoveto{\pgfqpoint{0.000000in}{0.000000in}}%
\pgfpathlineto{\pgfqpoint{0.000000in}{0.027778in}}%
\pgfusepath{stroke,fill}%
}%
\begin{pgfscope}%
\pgfsys@transformshift{2.442702in}{3.064738in}%
\pgfsys@useobject{currentmarker}{}%
\end{pgfscope}%
\end{pgfscope}%
\begin{pgfscope}%
\definecolor{textcolor}{rgb}{0.000000,0.000000,0.000000}%
\pgfsetstrokecolor{textcolor}%
\pgfsetfillcolor{textcolor}%
\pgftext[x=2.442702in,y=3.050849in,,top]{\color{textcolor}{\rmfamily\fontsize{5.000000}{6.000000}\selectfont\catcode`\^=\active\def^{\ifmmode\sp\else\^{}\fi}\catcode`\%=\active\def%{\%}4.985}}%
\end{pgfscope}%
\begin{pgfscope}%
\pgfpathrectangle{\pgfqpoint{2.197902in}{3.064738in}}{\pgfqpoint{0.979200in}{0.326250in}}%
\pgfusepath{clip}%
\pgfsetbuttcap%
\pgfsetroundjoin%
\pgfsetlinewidth{0.401500pt}%
\definecolor{currentstroke}{rgb}{0.690196,0.690196,0.690196}%
\pgfsetstrokecolor{currentstroke}%
\pgfsetstrokeopacity{0.250000}%
\pgfsetdash{{1.480000pt}{0.640000pt}}{0.000000pt}%
\pgfpathmoveto{\pgfqpoint{2.687502in}{3.064738in}}%
\pgfpathlineto{\pgfqpoint{2.687502in}{3.390988in}}%
\pgfusepath{stroke}%
\end{pgfscope}%
\begin{pgfscope}%
\pgfsetbuttcap%
\pgfsetroundjoin%
\definecolor{currentfill}{rgb}{0.000000,0.000000,0.000000}%
\pgfsetfillcolor{currentfill}%
\pgfsetlinewidth{0.803000pt}%
\definecolor{currentstroke}{rgb}{0.000000,0.000000,0.000000}%
\pgfsetstrokecolor{currentstroke}%
\pgfsetdash{}{0pt}%
\pgfsys@defobject{currentmarker}{\pgfqpoint{0.000000in}{0.000000in}}{\pgfqpoint{0.000000in}{0.027778in}}{%
\pgfpathmoveto{\pgfqpoint{0.000000in}{0.000000in}}%
\pgfpathlineto{\pgfqpoint{0.000000in}{0.027778in}}%
\pgfusepath{stroke,fill}%
}%
\begin{pgfscope}%
\pgfsys@transformshift{2.687502in}{3.064738in}%
\pgfsys@useobject{currentmarker}{}%
\end{pgfscope}%
\end{pgfscope}%
\begin{pgfscope}%
\definecolor{textcolor}{rgb}{0.000000,0.000000,0.000000}%
\pgfsetstrokecolor{textcolor}%
\pgfsetfillcolor{textcolor}%
\pgftext[x=2.687502in,y=3.050849in,,top]{\color{textcolor}{\rmfamily\fontsize{5.000000}{6.000000}\selectfont\catcode`\^=\active\def^{\ifmmode\sp\else\^{}\fi}\catcode`\%=\active\def%{\%}4.990}}%
\end{pgfscope}%
\begin{pgfscope}%
\pgfpathrectangle{\pgfqpoint{2.197902in}{3.064738in}}{\pgfqpoint{0.979200in}{0.326250in}}%
\pgfusepath{clip}%
\pgfsetbuttcap%
\pgfsetroundjoin%
\pgfsetlinewidth{0.401500pt}%
\definecolor{currentstroke}{rgb}{0.690196,0.690196,0.690196}%
\pgfsetstrokecolor{currentstroke}%
\pgfsetstrokeopacity{0.250000}%
\pgfsetdash{{1.480000pt}{0.640000pt}}{0.000000pt}%
\pgfpathmoveto{\pgfqpoint{2.932302in}{3.064738in}}%
\pgfpathlineto{\pgfqpoint{2.932302in}{3.390988in}}%
\pgfusepath{stroke}%
\end{pgfscope}%
\begin{pgfscope}%
\pgfsetbuttcap%
\pgfsetroundjoin%
\definecolor{currentfill}{rgb}{0.000000,0.000000,0.000000}%
\pgfsetfillcolor{currentfill}%
\pgfsetlinewidth{0.803000pt}%
\definecolor{currentstroke}{rgb}{0.000000,0.000000,0.000000}%
\pgfsetstrokecolor{currentstroke}%
\pgfsetdash{}{0pt}%
\pgfsys@defobject{currentmarker}{\pgfqpoint{0.000000in}{0.000000in}}{\pgfqpoint{0.000000in}{0.027778in}}{%
\pgfpathmoveto{\pgfqpoint{0.000000in}{0.000000in}}%
\pgfpathlineto{\pgfqpoint{0.000000in}{0.027778in}}%
\pgfusepath{stroke,fill}%
}%
\begin{pgfscope}%
\pgfsys@transformshift{2.932302in}{3.064738in}%
\pgfsys@useobject{currentmarker}{}%
\end{pgfscope}%
\end{pgfscope}%
\begin{pgfscope}%
\definecolor{textcolor}{rgb}{0.000000,0.000000,0.000000}%
\pgfsetstrokecolor{textcolor}%
\pgfsetfillcolor{textcolor}%
\pgftext[x=2.932302in,y=3.050849in,,top]{\color{textcolor}{\rmfamily\fontsize{5.000000}{6.000000}\selectfont\catcode`\^=\active\def^{\ifmmode\sp\else\^{}\fi}\catcode`\%=\active\def%{\%}4.995}}%
\end{pgfscope}%
\begin{pgfscope}%
\pgfpathrectangle{\pgfqpoint{2.197902in}{3.064738in}}{\pgfqpoint{0.979200in}{0.326250in}}%
\pgfusepath{clip}%
\pgfsetbuttcap%
\pgfsetroundjoin%
\pgfsetlinewidth{0.401500pt}%
\definecolor{currentstroke}{rgb}{0.690196,0.690196,0.690196}%
\pgfsetstrokecolor{currentstroke}%
\pgfsetstrokeopacity{0.250000}%
\pgfsetdash{{1.480000pt}{0.640000pt}}{0.000000pt}%
\pgfpathmoveto{\pgfqpoint{3.177102in}{3.064738in}}%
\pgfpathlineto{\pgfqpoint{3.177102in}{3.390988in}}%
\pgfusepath{stroke}%
\end{pgfscope}%
\begin{pgfscope}%
\pgfsetbuttcap%
\pgfsetroundjoin%
\definecolor{currentfill}{rgb}{0.000000,0.000000,0.000000}%
\pgfsetfillcolor{currentfill}%
\pgfsetlinewidth{0.803000pt}%
\definecolor{currentstroke}{rgb}{0.000000,0.000000,0.000000}%
\pgfsetstrokecolor{currentstroke}%
\pgfsetdash{}{0pt}%
\pgfsys@defobject{currentmarker}{\pgfqpoint{0.000000in}{0.000000in}}{\pgfqpoint{0.000000in}{0.027778in}}{%
\pgfpathmoveto{\pgfqpoint{0.000000in}{0.000000in}}%
\pgfpathlineto{\pgfqpoint{0.000000in}{0.027778in}}%
\pgfusepath{stroke,fill}%
}%
\begin{pgfscope}%
\pgfsys@transformshift{3.177102in}{3.064738in}%
\pgfsys@useobject{currentmarker}{}%
\end{pgfscope}%
\end{pgfscope}%
\begin{pgfscope}%
\definecolor{textcolor}{rgb}{0.000000,0.000000,0.000000}%
\pgfsetstrokecolor{textcolor}%
\pgfsetfillcolor{textcolor}%
\pgftext[x=3.177102in,y=3.050849in,,top]{\color{textcolor}{\rmfamily\fontsize{5.000000}{6.000000}\selectfont\catcode`\^=\active\def^{\ifmmode\sp\else\^{}\fi}\catcode`\%=\active\def%{\%}5.000}}%
\end{pgfscope}%
\begin{pgfscope}%
\pgfpathrectangle{\pgfqpoint{2.197902in}{3.064738in}}{\pgfqpoint{0.979200in}{0.326250in}}%
\pgfusepath{clip}%
\pgfsetbuttcap%
\pgfsetroundjoin%
\pgfsetlinewidth{0.401500pt}%
\definecolor{currentstroke}{rgb}{0.690196,0.690196,0.690196}%
\pgfsetstrokecolor{currentstroke}%
\pgfsetstrokeopacity{0.250000}%
\pgfsetdash{{1.480000pt}{0.640000pt}}{0.000000pt}%
\pgfpathmoveto{\pgfqpoint{2.197902in}{3.064738in}}%
\pgfpathlineto{\pgfqpoint{3.177102in}{3.064738in}}%
\pgfusepath{stroke}%
\end{pgfscope}%
\begin{pgfscope}%
\pgfsetbuttcap%
\pgfsetroundjoin%
\definecolor{currentfill}{rgb}{0.000000,0.000000,0.000000}%
\pgfsetfillcolor{currentfill}%
\pgfsetlinewidth{0.803000pt}%
\definecolor{currentstroke}{rgb}{0.000000,0.000000,0.000000}%
\pgfsetstrokecolor{currentstroke}%
\pgfsetdash{}{0pt}%
\pgfsys@defobject{currentmarker}{\pgfqpoint{0.000000in}{0.000000in}}{\pgfqpoint{0.027778in}{0.000000in}}{%
\pgfpathmoveto{\pgfqpoint{0.000000in}{0.000000in}}%
\pgfpathlineto{\pgfqpoint{0.027778in}{0.000000in}}%
\pgfusepath{stroke,fill}%
}%
\begin{pgfscope}%
\pgfsys@transformshift{2.197902in}{3.064738in}%
\pgfsys@useobject{currentmarker}{}%
\end{pgfscope}%
\end{pgfscope}%
\begin{pgfscope}%
\definecolor{textcolor}{rgb}{0.000000,0.000000,0.000000}%
\pgfsetstrokecolor{textcolor}%
\pgfsetfillcolor{textcolor}%
\pgftext[x=1.981463in, y=3.040625in, left, base]{\color{textcolor}{\rmfamily\fontsize{5.000000}{6.000000}\selectfont\catcode`\^=\active\def^{\ifmmode\sp\else\^{}\fi}\catcode`\%=\active\def%{\%}-0.50}}%
\end{pgfscope}%
\begin{pgfscope}%
\pgfpathrectangle{\pgfqpoint{2.197902in}{3.064738in}}{\pgfqpoint{0.979200in}{0.326250in}}%
\pgfusepath{clip}%
\pgfsetbuttcap%
\pgfsetroundjoin%
\pgfsetlinewidth{0.401500pt}%
\definecolor{currentstroke}{rgb}{0.690196,0.690196,0.690196}%
\pgfsetstrokecolor{currentstroke}%
\pgfsetstrokeopacity{0.250000}%
\pgfsetdash{{1.480000pt}{0.640000pt}}{0.000000pt}%
\pgfpathmoveto{\pgfqpoint{2.197902in}{3.227863in}}%
\pgfpathlineto{\pgfqpoint{3.177102in}{3.227863in}}%
\pgfusepath{stroke}%
\end{pgfscope}%
\begin{pgfscope}%
\pgfsetbuttcap%
\pgfsetroundjoin%
\definecolor{currentfill}{rgb}{0.000000,0.000000,0.000000}%
\pgfsetfillcolor{currentfill}%
\pgfsetlinewidth{0.803000pt}%
\definecolor{currentstroke}{rgb}{0.000000,0.000000,0.000000}%
\pgfsetstrokecolor{currentstroke}%
\pgfsetdash{}{0pt}%
\pgfsys@defobject{currentmarker}{\pgfqpoint{0.000000in}{0.000000in}}{\pgfqpoint{0.027778in}{0.000000in}}{%
\pgfpathmoveto{\pgfqpoint{0.000000in}{0.000000in}}%
\pgfpathlineto{\pgfqpoint{0.027778in}{0.000000in}}%
\pgfusepath{stroke,fill}%
}%
\begin{pgfscope}%
\pgfsys@transformshift{2.197902in}{3.227863in}%
\pgfsys@useobject{currentmarker}{}%
\end{pgfscope}%
\end{pgfscope}%
\begin{pgfscope}%
\definecolor{textcolor}{rgb}{0.000000,0.000000,0.000000}%
\pgfsetstrokecolor{textcolor}%
\pgfsetfillcolor{textcolor}%
\pgftext[x=2.014257in, y=3.203750in, left, base]{\color{textcolor}{\rmfamily\fontsize{5.000000}{6.000000}\selectfont\catcode`\^=\active\def^{\ifmmode\sp\else\^{}\fi}\catcode`\%=\active\def%{\%}0.73}}%
\end{pgfscope}%
\begin{pgfscope}%
\pgfpathrectangle{\pgfqpoint{2.197902in}{3.064738in}}{\pgfqpoint{0.979200in}{0.326250in}}%
\pgfusepath{clip}%
\pgfsetbuttcap%
\pgfsetroundjoin%
\pgfsetlinewidth{0.401500pt}%
\definecolor{currentstroke}{rgb}{0.690196,0.690196,0.690196}%
\pgfsetstrokecolor{currentstroke}%
\pgfsetstrokeopacity{0.250000}%
\pgfsetdash{{1.480000pt}{0.640000pt}}{0.000000pt}%
\pgfpathmoveto{\pgfqpoint{2.197902in}{3.390988in}}%
\pgfpathlineto{\pgfqpoint{3.177102in}{3.390988in}}%
\pgfusepath{stroke}%
\end{pgfscope}%
\begin{pgfscope}%
\pgfsetbuttcap%
\pgfsetroundjoin%
\definecolor{currentfill}{rgb}{0.000000,0.000000,0.000000}%
\pgfsetfillcolor{currentfill}%
\pgfsetlinewidth{0.803000pt}%
\definecolor{currentstroke}{rgb}{0.000000,0.000000,0.000000}%
\pgfsetstrokecolor{currentstroke}%
\pgfsetdash{}{0pt}%
\pgfsys@defobject{currentmarker}{\pgfqpoint{0.000000in}{0.000000in}}{\pgfqpoint{0.027778in}{0.000000in}}{%
\pgfpathmoveto{\pgfqpoint{0.000000in}{0.000000in}}%
\pgfpathlineto{\pgfqpoint{0.027778in}{0.000000in}}%
\pgfusepath{stroke,fill}%
}%
\begin{pgfscope}%
\pgfsys@transformshift{2.197902in}{3.390988in}%
\pgfsys@useobject{currentmarker}{}%
\end{pgfscope}%
\end{pgfscope}%
\begin{pgfscope}%
\definecolor{textcolor}{rgb}{0.000000,0.000000,0.000000}%
\pgfsetstrokecolor{textcolor}%
\pgfsetfillcolor{textcolor}%
\pgftext[x=2.014257in, y=3.366875in, left, base]{\color{textcolor}{\rmfamily\fontsize{5.000000}{6.000000}\selectfont\catcode`\^=\active\def^{\ifmmode\sp\else\^{}\fi}\catcode`\%=\active\def%{\%}1.95}}%
\end{pgfscope}%
\begin{pgfscope}%
\pgfpathrectangle{\pgfqpoint{2.197902in}{3.064738in}}{\pgfqpoint{0.979200in}{0.326250in}}%
\pgfusepath{clip}%
\pgfsetrectcap%
\pgfsetroundjoin%
\pgfsetlinewidth{1.505625pt}%
\definecolor{currentstroke}{rgb}{0.000000,0.447059,0.698039}%
\pgfsetstrokecolor{currentstroke}%
\pgfsetdash{}{0pt}%
\pgfpathmoveto{\pgfqpoint{2.193930in}{3.131319in}}%
\pgfpathlineto{\pgfqpoint{2.198379in}{3.131319in}}%
\pgfpathlineto{\pgfqpoint{2.198813in}{3.290290in}}%
\pgfpathlineto{\pgfqpoint{2.199311in}{3.232250in}}%
\pgfpathlineto{\pgfqpoint{2.199513in}{3.239868in}}%
\pgfpathlineto{\pgfqpoint{2.199777in}{3.239768in}}%
\pgfpathlineto{\pgfqpoint{2.200683in}{3.146508in}}%
\pgfpathlineto{\pgfqpoint{2.201028in}{3.185527in}}%
\pgfpathlineto{\pgfqpoint{2.201719in}{3.139239in}}%
\pgfpathlineto{\pgfqpoint{2.203101in}{3.174439in}}%
\pgfpathlineto{\pgfqpoint{2.205865in}{3.127982in}}%
\pgfpathlineto{\pgfqpoint{2.210055in}{3.169964in}}%
\pgfpathlineto{\pgfqpoint{2.214951in}{3.132816in}}%
\pgfpathlineto{\pgfqpoint{2.219847in}{3.178908in}}%
\pgfpathlineto{\pgfqpoint{2.224743in}{3.142372in}}%
\pgfpathlineto{\pgfqpoint{2.229639in}{3.188469in}}%
\pgfpathlineto{\pgfqpoint{2.234535in}{3.151975in}}%
\pgfpathlineto{\pgfqpoint{2.239431in}{3.198047in}}%
\pgfpathlineto{\pgfqpoint{2.244327in}{3.161578in}}%
\pgfpathlineto{\pgfqpoint{2.249223in}{3.207619in}}%
\pgfpathlineto{\pgfqpoint{2.254119in}{3.171181in}}%
\pgfpathlineto{\pgfqpoint{2.259015in}{3.217190in}}%
\pgfpathlineto{\pgfqpoint{2.263911in}{3.180784in}}%
\pgfpathlineto{\pgfqpoint{2.268807in}{3.226761in}}%
\pgfpathlineto{\pgfqpoint{2.273703in}{3.190387in}}%
\pgfpathlineto{\pgfqpoint{2.278599in}{3.236332in}}%
\pgfpathlineto{\pgfqpoint{2.283495in}{3.199989in}}%
\pgfpathlineto{\pgfqpoint{2.288391in}{3.245904in}}%
\pgfpathlineto{\pgfqpoint{2.293287in}{3.209592in}}%
\pgfpathlineto{\pgfqpoint{2.298183in}{3.255475in}}%
\pgfpathlineto{\pgfqpoint{2.303079in}{3.219194in}}%
\pgfpathlineto{\pgfqpoint{2.307975in}{3.265045in}}%
\pgfpathlineto{\pgfqpoint{2.312871in}{3.228797in}}%
\pgfpathlineto{\pgfqpoint{2.317767in}{3.274616in}}%
\pgfpathlineto{\pgfqpoint{2.322663in}{3.238399in}}%
\pgfpathlineto{\pgfqpoint{2.327559in}{3.284187in}}%
\pgfpathlineto{\pgfqpoint{2.332455in}{3.248001in}}%
\pgfpathlineto{\pgfqpoint{2.337351in}{3.293758in}}%
\pgfpathlineto{\pgfqpoint{2.342247in}{3.257603in}}%
\pgfpathlineto{\pgfqpoint{2.347143in}{3.303328in}}%
\pgfpathlineto{\pgfqpoint{2.352039in}{3.267205in}}%
\pgfpathlineto{\pgfqpoint{2.356935in}{3.312899in}}%
\pgfpathlineto{\pgfqpoint{2.361831in}{3.276807in}}%
\pgfpathlineto{\pgfqpoint{2.366727in}{3.322469in}}%
\pgfpathlineto{\pgfqpoint{2.371623in}{3.286408in}}%
\pgfpathlineto{\pgfqpoint{2.376519in}{3.332039in}}%
\pgfpathlineto{\pgfqpoint{2.381415in}{3.296010in}}%
\pgfpathlineto{\pgfqpoint{2.386311in}{3.341610in}}%
\pgfpathlineto{\pgfqpoint{2.390051in}{3.304482in}}%
\pgfpathlineto{\pgfqpoint{2.393791in}{3.348911in}}%
\pgfpathlineto{\pgfqpoint{2.393847in}{3.271917in}}%
\pgfpathlineto{\pgfqpoint{2.394417in}{3.282840in}}%
\pgfpathlineto{\pgfqpoint{2.396158in}{3.305458in}}%
\pgfpathlineto{\pgfqpoint{2.397224in}{3.257621in}}%
\pgfpathlineto{\pgfqpoint{2.398252in}{3.280377in}}%
\pgfpathlineto{\pgfqpoint{2.399755in}{3.228467in}}%
\pgfpathlineto{\pgfqpoint{2.401327in}{3.223569in}}%
\pgfpathlineto{\pgfqpoint{2.402972in}{3.131319in}}%
\pgfpathlineto{\pgfqpoint{2.688022in}{3.131320in}}%
\pgfpathlineto{\pgfqpoint{2.688385in}{3.301208in}}%
\pgfpathlineto{\pgfqpoint{2.688722in}{3.265440in}}%
\pgfpathlineto{\pgfqpoint{2.689809in}{3.123227in}}%
\pgfpathlineto{\pgfqpoint{2.689953in}{3.216709in}}%
\pgfpathlineto{\pgfqpoint{2.690562in}{3.122262in}}%
\pgfpathlineto{\pgfqpoint{2.691014in}{3.203733in}}%
\pgfpathlineto{\pgfqpoint{2.691919in}{3.114015in}}%
\pgfpathlineto{\pgfqpoint{2.693729in}{3.192400in}}%
\pgfpathlineto{\pgfqpoint{2.697347in}{3.105121in}}%
\pgfpathlineto{\pgfqpoint{2.701692in}{3.192063in}}%
\pgfpathlineto{\pgfqpoint{2.706588in}{3.112343in}}%
\pgfpathlineto{\pgfqpoint{2.711484in}{3.201379in}}%
\pgfpathlineto{\pgfqpoint{2.716380in}{3.121932in}}%
\pgfpathlineto{\pgfqpoint{2.721276in}{3.210937in}}%
\pgfpathlineto{\pgfqpoint{2.726172in}{3.131541in}}%
\pgfpathlineto{\pgfqpoint{2.731068in}{3.220509in}}%
\pgfpathlineto{\pgfqpoint{2.735964in}{3.141148in}}%
\pgfpathlineto{\pgfqpoint{2.740860in}{3.230075in}}%
\pgfpathlineto{\pgfqpoint{2.745756in}{3.150757in}}%
\pgfpathlineto{\pgfqpoint{2.750652in}{3.239641in}}%
\pgfpathlineto{\pgfqpoint{2.755548in}{3.160365in}}%
\pgfpathlineto{\pgfqpoint{2.760444in}{3.249207in}}%
\pgfpathlineto{\pgfqpoint{2.765340in}{3.169973in}}%
\pgfpathlineto{\pgfqpoint{2.770236in}{3.258773in}}%
\pgfpathlineto{\pgfqpoint{2.775132in}{3.179581in}}%
\pgfpathlineto{\pgfqpoint{2.780028in}{3.268339in}}%
\pgfpathlineto{\pgfqpoint{2.784924in}{3.189189in}}%
\pgfpathlineto{\pgfqpoint{2.789820in}{3.277905in}}%
\pgfpathlineto{\pgfqpoint{2.794716in}{3.198796in}}%
\pgfpathlineto{\pgfqpoint{2.799612in}{3.287470in}}%
\pgfpathlineto{\pgfqpoint{2.804508in}{3.208404in}}%
\pgfpathlineto{\pgfqpoint{2.809404in}{3.297036in}}%
\pgfpathlineto{\pgfqpoint{2.814300in}{3.218011in}}%
\pgfpathlineto{\pgfqpoint{2.819196in}{3.306601in}}%
\pgfpathlineto{\pgfqpoint{2.824092in}{3.227619in}}%
\pgfpathlineto{\pgfqpoint{2.828988in}{3.316167in}}%
\pgfpathlineto{\pgfqpoint{2.833884in}{3.237226in}}%
\pgfpathlineto{\pgfqpoint{2.838780in}{3.325732in}}%
\pgfpathlineto{\pgfqpoint{2.843676in}{3.246833in}}%
\pgfpathlineto{\pgfqpoint{2.848572in}{3.335297in}}%
\pgfpathlineto{\pgfqpoint{2.853468in}{3.256440in}}%
\pgfpathlineto{\pgfqpoint{2.858364in}{3.344862in}}%
\pgfpathlineto{\pgfqpoint{2.863260in}{3.266047in}}%
\pgfpathlineto{\pgfqpoint{2.868156in}{3.354428in}}%
\pgfpathlineto{\pgfqpoint{2.873052in}{3.275654in}}%
\pgfpathlineto{\pgfqpoint{2.877948in}{3.363993in}}%
\pgfpathlineto{\pgfqpoint{2.880670in}{3.283140in}}%
\pgfpathlineto{\pgfqpoint{2.883391in}{3.369283in}}%
\pgfpathlineto{\pgfqpoint{2.883447in}{3.270930in}}%
\pgfpathlineto{\pgfqpoint{2.884072in}{3.282808in}}%
\pgfpathlineto{\pgfqpoint{2.885721in}{3.304077in}}%
\pgfpathlineto{\pgfqpoint{2.886791in}{3.259238in}}%
\pgfpathlineto{\pgfqpoint{2.887770in}{3.276615in}}%
\pgfpathlineto{\pgfqpoint{2.889174in}{3.230987in}}%
\pgfpathlineto{\pgfqpoint{2.890691in}{3.265245in}}%
\pgfpathlineto{\pgfqpoint{2.892294in}{3.131319in}}%
\pgfpathlineto{\pgfqpoint{3.177102in}{3.131319in}}%
\pgfpathlineto{\pgfqpoint{3.177102in}{3.131319in}}%
\pgfusepath{stroke}%
\end{pgfscope}%
\begin{pgfscope}%
\pgfsetrectcap%
\pgfsetmiterjoin%
\pgfsetlinewidth{0.803000pt}%
\definecolor{currentstroke}{rgb}{0.000000,0.000000,0.000000}%
\pgfsetstrokecolor{currentstroke}%
\pgfsetdash{}{0pt}%
\pgfpathmoveto{\pgfqpoint{2.197902in}{3.064738in}}%
\pgfpathlineto{\pgfqpoint{2.197902in}{3.390988in}}%
\pgfusepath{stroke}%
\end{pgfscope}%
\begin{pgfscope}%
\pgfsetrectcap%
\pgfsetmiterjoin%
\pgfsetlinewidth{0.803000pt}%
\definecolor{currentstroke}{rgb}{0.000000,0.000000,0.000000}%
\pgfsetstrokecolor{currentstroke}%
\pgfsetdash{}{0pt}%
\pgfpathmoveto{\pgfqpoint{3.177102in}{3.064738in}}%
\pgfpathlineto{\pgfqpoint{3.177102in}{3.390988in}}%
\pgfusepath{stroke}%
\end{pgfscope}%
\begin{pgfscope}%
\pgfsetrectcap%
\pgfsetmiterjoin%
\pgfsetlinewidth{0.803000pt}%
\definecolor{currentstroke}{rgb}{0.000000,0.000000,0.000000}%
\pgfsetstrokecolor{currentstroke}%
\pgfsetdash{}{0pt}%
\pgfpathmoveto{\pgfqpoint{2.197902in}{3.064738in}}%
\pgfpathlineto{\pgfqpoint{3.177102in}{3.064738in}}%
\pgfusepath{stroke}%
\end{pgfscope}%
\begin{pgfscope}%
\pgfsetrectcap%
\pgfsetmiterjoin%
\pgfsetlinewidth{0.803000pt}%
\definecolor{currentstroke}{rgb}{0.000000,0.000000,0.000000}%
\pgfsetstrokecolor{currentstroke}%
\pgfsetdash{}{0pt}%
\pgfpathmoveto{\pgfqpoint{2.197902in}{3.390988in}}%
\pgfpathlineto{\pgfqpoint{3.177102in}{3.390988in}}%
\pgfusepath{stroke}%
\end{pgfscope}%
\begin{pgfscope}%
\pgfsetbuttcap%
\pgfsetmiterjoin%
\definecolor{currentfill}{rgb}{1.000000,1.000000,1.000000}%
\pgfsetfillcolor{currentfill}%
\pgfsetlinewidth{0.000000pt}%
\definecolor{currentstroke}{rgb}{0.000000,0.000000,0.000000}%
\pgfsetstrokecolor{currentstroke}%
\pgfsetstrokeopacity{0.000000}%
\pgfsetdash{}{0pt}%
\pgfpathmoveto{\pgfqpoint{2.197902in}{1.995363in}}%
\pgfpathlineto{\pgfqpoint{3.177102in}{1.995363in}}%
\pgfpathlineto{\pgfqpoint{3.177102in}{2.321613in}}%
\pgfpathlineto{\pgfqpoint{2.197902in}{2.321613in}}%
\pgfpathlineto{\pgfqpoint{2.197902in}{1.995363in}}%
\pgfpathclose%
\pgfusepath{fill}%
\end{pgfscope}%
\begin{pgfscope}%
\pgfpathrectangle{\pgfqpoint{2.197902in}{1.995363in}}{\pgfqpoint{0.979200in}{0.326250in}}%
\pgfusepath{clip}%
\pgfsetbuttcap%
\pgfsetroundjoin%
\pgfsetlinewidth{0.401500pt}%
\definecolor{currentstroke}{rgb}{0.690196,0.690196,0.690196}%
\pgfsetstrokecolor{currentstroke}%
\pgfsetstrokeopacity{0.250000}%
\pgfsetdash{{1.480000pt}{0.640000pt}}{0.000000pt}%
\pgfpathmoveto{\pgfqpoint{2.197902in}{1.995363in}}%
\pgfpathlineto{\pgfqpoint{2.197902in}{2.321613in}}%
\pgfusepath{stroke}%
\end{pgfscope}%
\begin{pgfscope}%
\pgfsetbuttcap%
\pgfsetroundjoin%
\definecolor{currentfill}{rgb}{0.000000,0.000000,0.000000}%
\pgfsetfillcolor{currentfill}%
\pgfsetlinewidth{0.803000pt}%
\definecolor{currentstroke}{rgb}{0.000000,0.000000,0.000000}%
\pgfsetstrokecolor{currentstroke}%
\pgfsetdash{}{0pt}%
\pgfsys@defobject{currentmarker}{\pgfqpoint{0.000000in}{0.000000in}}{\pgfqpoint{0.000000in}{0.027778in}}{%
\pgfpathmoveto{\pgfqpoint{0.000000in}{0.000000in}}%
\pgfpathlineto{\pgfqpoint{0.000000in}{0.027778in}}%
\pgfusepath{stroke,fill}%
}%
\begin{pgfscope}%
\pgfsys@transformshift{2.197902in}{1.995363in}%
\pgfsys@useobject{currentmarker}{}%
\end{pgfscope}%
\end{pgfscope}%
\begin{pgfscope}%
\definecolor{textcolor}{rgb}{0.000000,0.000000,0.000000}%
\pgfsetstrokecolor{textcolor}%
\pgfsetfillcolor{textcolor}%
\pgftext[x=2.197902in,y=1.981474in,,top]{\color{textcolor}{\rmfamily\fontsize{5.000000}{6.000000}\selectfont\catcode`\^=\active\def^{\ifmmode\sp\else\^{}\fi}\catcode`\%=\active\def%{\%}4.980}}%
\end{pgfscope}%
\begin{pgfscope}%
\pgfpathrectangle{\pgfqpoint{2.197902in}{1.995363in}}{\pgfqpoint{0.979200in}{0.326250in}}%
\pgfusepath{clip}%
\pgfsetbuttcap%
\pgfsetroundjoin%
\pgfsetlinewidth{0.401500pt}%
\definecolor{currentstroke}{rgb}{0.690196,0.690196,0.690196}%
\pgfsetstrokecolor{currentstroke}%
\pgfsetstrokeopacity{0.250000}%
\pgfsetdash{{1.480000pt}{0.640000pt}}{0.000000pt}%
\pgfpathmoveto{\pgfqpoint{2.442702in}{1.995363in}}%
\pgfpathlineto{\pgfqpoint{2.442702in}{2.321613in}}%
\pgfusepath{stroke}%
\end{pgfscope}%
\begin{pgfscope}%
\pgfsetbuttcap%
\pgfsetroundjoin%
\definecolor{currentfill}{rgb}{0.000000,0.000000,0.000000}%
\pgfsetfillcolor{currentfill}%
\pgfsetlinewidth{0.803000pt}%
\definecolor{currentstroke}{rgb}{0.000000,0.000000,0.000000}%
\pgfsetstrokecolor{currentstroke}%
\pgfsetdash{}{0pt}%
\pgfsys@defobject{currentmarker}{\pgfqpoint{0.000000in}{0.000000in}}{\pgfqpoint{0.000000in}{0.027778in}}{%
\pgfpathmoveto{\pgfqpoint{0.000000in}{0.000000in}}%
\pgfpathlineto{\pgfqpoint{0.000000in}{0.027778in}}%
\pgfusepath{stroke,fill}%
}%
\begin{pgfscope}%
\pgfsys@transformshift{2.442702in}{1.995363in}%
\pgfsys@useobject{currentmarker}{}%
\end{pgfscope}%
\end{pgfscope}%
\begin{pgfscope}%
\definecolor{textcolor}{rgb}{0.000000,0.000000,0.000000}%
\pgfsetstrokecolor{textcolor}%
\pgfsetfillcolor{textcolor}%
\pgftext[x=2.442702in,y=1.981474in,,top]{\color{textcolor}{\rmfamily\fontsize{5.000000}{6.000000}\selectfont\catcode`\^=\active\def^{\ifmmode\sp\else\^{}\fi}\catcode`\%=\active\def%{\%}4.985}}%
\end{pgfscope}%
\begin{pgfscope}%
\pgfpathrectangle{\pgfqpoint{2.197902in}{1.995363in}}{\pgfqpoint{0.979200in}{0.326250in}}%
\pgfusepath{clip}%
\pgfsetbuttcap%
\pgfsetroundjoin%
\pgfsetlinewidth{0.401500pt}%
\definecolor{currentstroke}{rgb}{0.690196,0.690196,0.690196}%
\pgfsetstrokecolor{currentstroke}%
\pgfsetstrokeopacity{0.250000}%
\pgfsetdash{{1.480000pt}{0.640000pt}}{0.000000pt}%
\pgfpathmoveto{\pgfqpoint{2.687502in}{1.995363in}}%
\pgfpathlineto{\pgfqpoint{2.687502in}{2.321613in}}%
\pgfusepath{stroke}%
\end{pgfscope}%
\begin{pgfscope}%
\pgfsetbuttcap%
\pgfsetroundjoin%
\definecolor{currentfill}{rgb}{0.000000,0.000000,0.000000}%
\pgfsetfillcolor{currentfill}%
\pgfsetlinewidth{0.803000pt}%
\definecolor{currentstroke}{rgb}{0.000000,0.000000,0.000000}%
\pgfsetstrokecolor{currentstroke}%
\pgfsetdash{}{0pt}%
\pgfsys@defobject{currentmarker}{\pgfqpoint{0.000000in}{0.000000in}}{\pgfqpoint{0.000000in}{0.027778in}}{%
\pgfpathmoveto{\pgfqpoint{0.000000in}{0.000000in}}%
\pgfpathlineto{\pgfqpoint{0.000000in}{0.027778in}}%
\pgfusepath{stroke,fill}%
}%
\begin{pgfscope}%
\pgfsys@transformshift{2.687502in}{1.995363in}%
\pgfsys@useobject{currentmarker}{}%
\end{pgfscope}%
\end{pgfscope}%
\begin{pgfscope}%
\definecolor{textcolor}{rgb}{0.000000,0.000000,0.000000}%
\pgfsetstrokecolor{textcolor}%
\pgfsetfillcolor{textcolor}%
\pgftext[x=2.687502in,y=1.981474in,,top]{\color{textcolor}{\rmfamily\fontsize{5.000000}{6.000000}\selectfont\catcode`\^=\active\def^{\ifmmode\sp\else\^{}\fi}\catcode`\%=\active\def%{\%}4.990}}%
\end{pgfscope}%
\begin{pgfscope}%
\pgfpathrectangle{\pgfqpoint{2.197902in}{1.995363in}}{\pgfqpoint{0.979200in}{0.326250in}}%
\pgfusepath{clip}%
\pgfsetbuttcap%
\pgfsetroundjoin%
\pgfsetlinewidth{0.401500pt}%
\definecolor{currentstroke}{rgb}{0.690196,0.690196,0.690196}%
\pgfsetstrokecolor{currentstroke}%
\pgfsetstrokeopacity{0.250000}%
\pgfsetdash{{1.480000pt}{0.640000pt}}{0.000000pt}%
\pgfpathmoveto{\pgfqpoint{2.932302in}{1.995363in}}%
\pgfpathlineto{\pgfqpoint{2.932302in}{2.321613in}}%
\pgfusepath{stroke}%
\end{pgfscope}%
\begin{pgfscope}%
\pgfsetbuttcap%
\pgfsetroundjoin%
\definecolor{currentfill}{rgb}{0.000000,0.000000,0.000000}%
\pgfsetfillcolor{currentfill}%
\pgfsetlinewidth{0.803000pt}%
\definecolor{currentstroke}{rgb}{0.000000,0.000000,0.000000}%
\pgfsetstrokecolor{currentstroke}%
\pgfsetdash{}{0pt}%
\pgfsys@defobject{currentmarker}{\pgfqpoint{0.000000in}{0.000000in}}{\pgfqpoint{0.000000in}{0.027778in}}{%
\pgfpathmoveto{\pgfqpoint{0.000000in}{0.000000in}}%
\pgfpathlineto{\pgfqpoint{0.000000in}{0.027778in}}%
\pgfusepath{stroke,fill}%
}%
\begin{pgfscope}%
\pgfsys@transformshift{2.932302in}{1.995363in}%
\pgfsys@useobject{currentmarker}{}%
\end{pgfscope}%
\end{pgfscope}%
\begin{pgfscope}%
\definecolor{textcolor}{rgb}{0.000000,0.000000,0.000000}%
\pgfsetstrokecolor{textcolor}%
\pgfsetfillcolor{textcolor}%
\pgftext[x=2.932302in,y=1.981474in,,top]{\color{textcolor}{\rmfamily\fontsize{5.000000}{6.000000}\selectfont\catcode`\^=\active\def^{\ifmmode\sp\else\^{}\fi}\catcode`\%=\active\def%{\%}4.995}}%
\end{pgfscope}%
\begin{pgfscope}%
\pgfpathrectangle{\pgfqpoint{2.197902in}{1.995363in}}{\pgfqpoint{0.979200in}{0.326250in}}%
\pgfusepath{clip}%
\pgfsetbuttcap%
\pgfsetroundjoin%
\pgfsetlinewidth{0.401500pt}%
\definecolor{currentstroke}{rgb}{0.690196,0.690196,0.690196}%
\pgfsetstrokecolor{currentstroke}%
\pgfsetstrokeopacity{0.250000}%
\pgfsetdash{{1.480000pt}{0.640000pt}}{0.000000pt}%
\pgfpathmoveto{\pgfqpoint{3.177102in}{1.995363in}}%
\pgfpathlineto{\pgfqpoint{3.177102in}{2.321613in}}%
\pgfusepath{stroke}%
\end{pgfscope}%
\begin{pgfscope}%
\pgfsetbuttcap%
\pgfsetroundjoin%
\definecolor{currentfill}{rgb}{0.000000,0.000000,0.000000}%
\pgfsetfillcolor{currentfill}%
\pgfsetlinewidth{0.803000pt}%
\definecolor{currentstroke}{rgb}{0.000000,0.000000,0.000000}%
\pgfsetstrokecolor{currentstroke}%
\pgfsetdash{}{0pt}%
\pgfsys@defobject{currentmarker}{\pgfqpoint{0.000000in}{0.000000in}}{\pgfqpoint{0.000000in}{0.027778in}}{%
\pgfpathmoveto{\pgfqpoint{0.000000in}{0.000000in}}%
\pgfpathlineto{\pgfqpoint{0.000000in}{0.027778in}}%
\pgfusepath{stroke,fill}%
}%
\begin{pgfscope}%
\pgfsys@transformshift{3.177102in}{1.995363in}%
\pgfsys@useobject{currentmarker}{}%
\end{pgfscope}%
\end{pgfscope}%
\begin{pgfscope}%
\definecolor{textcolor}{rgb}{0.000000,0.000000,0.000000}%
\pgfsetstrokecolor{textcolor}%
\pgfsetfillcolor{textcolor}%
\pgftext[x=3.177102in,y=1.981474in,,top]{\color{textcolor}{\rmfamily\fontsize{5.000000}{6.000000}\selectfont\catcode`\^=\active\def^{\ifmmode\sp\else\^{}\fi}\catcode`\%=\active\def%{\%}5.000}}%
\end{pgfscope}%
\begin{pgfscope}%
\pgfpathrectangle{\pgfqpoint{2.197902in}{1.995363in}}{\pgfqpoint{0.979200in}{0.326250in}}%
\pgfusepath{clip}%
\pgfsetbuttcap%
\pgfsetroundjoin%
\pgfsetlinewidth{0.401500pt}%
\definecolor{currentstroke}{rgb}{0.690196,0.690196,0.690196}%
\pgfsetstrokecolor{currentstroke}%
\pgfsetstrokeopacity{0.250000}%
\pgfsetdash{{1.480000pt}{0.640000pt}}{0.000000pt}%
\pgfpathmoveto{\pgfqpoint{2.197902in}{1.995363in}}%
\pgfpathlineto{\pgfqpoint{3.177102in}{1.995363in}}%
\pgfusepath{stroke}%
\end{pgfscope}%
\begin{pgfscope}%
\pgfsetbuttcap%
\pgfsetroundjoin%
\definecolor{currentfill}{rgb}{0.000000,0.000000,0.000000}%
\pgfsetfillcolor{currentfill}%
\pgfsetlinewidth{0.803000pt}%
\definecolor{currentstroke}{rgb}{0.000000,0.000000,0.000000}%
\pgfsetstrokecolor{currentstroke}%
\pgfsetdash{}{0pt}%
\pgfsys@defobject{currentmarker}{\pgfqpoint{0.000000in}{0.000000in}}{\pgfqpoint{0.027778in}{0.000000in}}{%
\pgfpathmoveto{\pgfqpoint{0.000000in}{0.000000in}}%
\pgfpathlineto{\pgfqpoint{0.027778in}{0.000000in}}%
\pgfusepath{stroke,fill}%
}%
\begin{pgfscope}%
\pgfsys@transformshift{2.197902in}{1.995363in}%
\pgfsys@useobject{currentmarker}{}%
\end{pgfscope}%
\end{pgfscope}%
\begin{pgfscope}%
\definecolor{textcolor}{rgb}{0.000000,0.000000,0.000000}%
\pgfsetstrokecolor{textcolor}%
\pgfsetfillcolor{textcolor}%
\pgftext[x=1.966996in, y=1.971250in, left, base]{\color{textcolor}{\rmfamily\fontsize{5.000000}{6.000000}\selectfont\catcode`\^=\active\def^{\ifmmode\sp\else\^{}\fi}\catcode`\%=\active\def%{\%}45.00}}%
\end{pgfscope}%
\begin{pgfscope}%
\pgfpathrectangle{\pgfqpoint{2.197902in}{1.995363in}}{\pgfqpoint{0.979200in}{0.326250in}}%
\pgfusepath{clip}%
\pgfsetbuttcap%
\pgfsetroundjoin%
\pgfsetlinewidth{0.401500pt}%
\definecolor{currentstroke}{rgb}{0.690196,0.690196,0.690196}%
\pgfsetstrokecolor{currentstroke}%
\pgfsetstrokeopacity{0.250000}%
\pgfsetdash{{1.480000pt}{0.640000pt}}{0.000000pt}%
\pgfpathmoveto{\pgfqpoint{2.197902in}{2.158488in}}%
\pgfpathlineto{\pgfqpoint{3.177102in}{2.158488in}}%
\pgfusepath{stroke}%
\end{pgfscope}%
\begin{pgfscope}%
\pgfsetbuttcap%
\pgfsetroundjoin%
\definecolor{currentfill}{rgb}{0.000000,0.000000,0.000000}%
\pgfsetfillcolor{currentfill}%
\pgfsetlinewidth{0.803000pt}%
\definecolor{currentstroke}{rgb}{0.000000,0.000000,0.000000}%
\pgfsetstrokecolor{currentstroke}%
\pgfsetdash{}{0pt}%
\pgfsys@defobject{currentmarker}{\pgfqpoint{0.000000in}{0.000000in}}{\pgfqpoint{0.027778in}{0.000000in}}{%
\pgfpathmoveto{\pgfqpoint{0.000000in}{0.000000in}}%
\pgfpathlineto{\pgfqpoint{0.027778in}{0.000000in}}%
\pgfusepath{stroke,fill}%
}%
\begin{pgfscope}%
\pgfsys@transformshift{2.197902in}{2.158488in}%
\pgfsys@useobject{currentmarker}{}%
\end{pgfscope}%
\end{pgfscope}%
\begin{pgfscope}%
\definecolor{textcolor}{rgb}{0.000000,0.000000,0.000000}%
\pgfsetstrokecolor{textcolor}%
\pgfsetfillcolor{textcolor}%
\pgftext[x=1.966996in, y=2.134375in, left, base]{\color{textcolor}{\rmfamily\fontsize{5.000000}{6.000000}\selectfont\catcode`\^=\active\def^{\ifmmode\sp\else\^{}\fi}\catcode`\%=\active\def%{\%}65.00}}%
\end{pgfscope}%
\begin{pgfscope}%
\pgfpathrectangle{\pgfqpoint{2.197902in}{1.995363in}}{\pgfqpoint{0.979200in}{0.326250in}}%
\pgfusepath{clip}%
\pgfsetbuttcap%
\pgfsetroundjoin%
\pgfsetlinewidth{0.401500pt}%
\definecolor{currentstroke}{rgb}{0.690196,0.690196,0.690196}%
\pgfsetstrokecolor{currentstroke}%
\pgfsetstrokeopacity{0.250000}%
\pgfsetdash{{1.480000pt}{0.640000pt}}{0.000000pt}%
\pgfpathmoveto{\pgfqpoint{2.197902in}{2.321613in}}%
\pgfpathlineto{\pgfqpoint{3.177102in}{2.321613in}}%
\pgfusepath{stroke}%
\end{pgfscope}%
\begin{pgfscope}%
\pgfsetbuttcap%
\pgfsetroundjoin%
\definecolor{currentfill}{rgb}{0.000000,0.000000,0.000000}%
\pgfsetfillcolor{currentfill}%
\pgfsetlinewidth{0.803000pt}%
\definecolor{currentstroke}{rgb}{0.000000,0.000000,0.000000}%
\pgfsetstrokecolor{currentstroke}%
\pgfsetdash{}{0pt}%
\pgfsys@defobject{currentmarker}{\pgfqpoint{0.000000in}{0.000000in}}{\pgfqpoint{0.027778in}{0.000000in}}{%
\pgfpathmoveto{\pgfqpoint{0.000000in}{0.000000in}}%
\pgfpathlineto{\pgfqpoint{0.027778in}{0.000000in}}%
\pgfusepath{stroke,fill}%
}%
\begin{pgfscope}%
\pgfsys@transformshift{2.197902in}{2.321613in}%
\pgfsys@useobject{currentmarker}{}%
\end{pgfscope}%
\end{pgfscope}%
\begin{pgfscope}%
\definecolor{textcolor}{rgb}{0.000000,0.000000,0.000000}%
\pgfsetstrokecolor{textcolor}%
\pgfsetfillcolor{textcolor}%
\pgftext[x=1.966996in, y=2.297500in, left, base]{\color{textcolor}{\rmfamily\fontsize{5.000000}{6.000000}\selectfont\catcode`\^=\active\def^{\ifmmode\sp\else\^{}\fi}\catcode`\%=\active\def%{\%}85.00}}%
\end{pgfscope}%
\begin{pgfscope}%
\pgfpathrectangle{\pgfqpoint{2.197902in}{1.995363in}}{\pgfqpoint{0.979200in}{0.326250in}}%
\pgfusepath{clip}%
\pgfsetrectcap%
\pgfsetroundjoin%
\pgfsetlinewidth{1.505625pt}%
\definecolor{currentstroke}{rgb}{0.835294,0.368627,0.000000}%
\pgfsetstrokecolor{currentstroke}%
\pgfsetdash{}{0pt}%
\pgfpathmoveto{\pgfqpoint{2.400236in}{1.988418in}}%
\pgfpathlineto{\pgfqpoint{2.401327in}{2.218286in}}%
\pgfpathlineto{\pgfqpoint{2.402972in}{2.246409in}}%
\pgfpathlineto{\pgfqpoint{2.404480in}{2.095469in}}%
\pgfpathlineto{\pgfqpoint{2.405235in}{2.117177in}}%
\pgfpathlineto{\pgfqpoint{2.406356in}{2.256740in}}%
\pgfpathlineto{\pgfqpoint{2.407951in}{2.092094in}}%
\pgfpathlineto{\pgfqpoint{2.409264in}{2.253053in}}%
\pgfpathlineto{\pgfqpoint{2.409676in}{2.225118in}}%
\pgfpathlineto{\pgfqpoint{2.410849in}{2.093508in}}%
\pgfpathlineto{\pgfqpoint{2.412357in}{2.244998in}}%
\pgfpathlineto{\pgfqpoint{2.413679in}{2.094906in}}%
\pgfpathlineto{\pgfqpoint{2.414088in}{2.121334in}}%
\pgfpathlineto{\pgfqpoint{2.415300in}{2.242156in}}%
\pgfpathlineto{\pgfqpoint{2.416770in}{2.103026in}}%
\pgfpathlineto{\pgfqpoint{2.418215in}{2.240970in}}%
\pgfpathlineto{\pgfqpoint{2.419438in}{2.103604in}}%
\pgfpathlineto{\pgfqpoint{2.419839in}{2.112572in}}%
\pgfpathlineto{\pgfqpoint{2.421097in}{2.239614in}}%
\pgfpathlineto{\pgfqpoint{2.422590in}{2.105655in}}%
\pgfpathlineto{\pgfqpoint{2.423057in}{2.145970in}}%
\pgfpathlineto{\pgfqpoint{2.423889in}{2.236338in}}%
\pgfpathlineto{\pgfqpoint{2.424304in}{2.220254in}}%
\pgfpathlineto{\pgfqpoint{2.425510in}{2.107801in}}%
\pgfpathlineto{\pgfqpoint{2.427029in}{2.231363in}}%
\pgfpathlineto{\pgfqpoint{2.428332in}{2.109741in}}%
\pgfpathlineto{\pgfqpoint{2.428744in}{2.128521in}}%
\pgfpathlineto{\pgfqpoint{2.429969in}{2.229057in}}%
\pgfpathlineto{\pgfqpoint{2.431434in}{2.115366in}}%
\pgfpathlineto{\pgfqpoint{2.432834in}{2.228754in}}%
\pgfpathlineto{\pgfqpoint{2.434057in}{2.118867in}}%
\pgfpathlineto{\pgfqpoint{2.434460in}{2.120713in}}%
\pgfpathlineto{\pgfqpoint{2.435686in}{2.226630in}}%
\pgfpathlineto{\pgfqpoint{2.436077in}{2.210087in}}%
\pgfpathlineto{\pgfqpoint{2.437302in}{2.119146in}}%
\pgfpathlineto{\pgfqpoint{2.438581in}{2.223799in}}%
\pgfpathlineto{\pgfqpoint{2.438997in}{2.209969in}}%
\pgfpathlineto{\pgfqpoint{2.440221in}{2.120720in}}%
\pgfpathlineto{\pgfqpoint{2.441761in}{2.217964in}}%
\pgfpathlineto{\pgfqpoint{2.443118in}{2.122013in}}%
\pgfpathlineto{\pgfqpoint{2.444364in}{2.216875in}}%
\pgfpathlineto{\pgfqpoint{2.444778in}{2.213138in}}%
\pgfpathlineto{\pgfqpoint{2.446049in}{2.123853in}}%
\pgfpathlineto{\pgfqpoint{2.447677in}{2.213460in}}%
\pgfpathlineto{\pgfqpoint{2.449002in}{2.125917in}}%
\pgfpathlineto{\pgfqpoint{2.450256in}{2.213236in}}%
\pgfpathlineto{\pgfqpoint{2.450673in}{2.209654in}}%
\pgfpathlineto{\pgfqpoint{2.451960in}{2.127844in}}%
\pgfpathlineto{\pgfqpoint{2.453169in}{2.210365in}}%
\pgfpathlineto{\pgfqpoint{2.453595in}{2.209190in}}%
\pgfpathlineto{\pgfqpoint{2.454959in}{2.130413in}}%
\pgfpathlineto{\pgfqpoint{2.456228in}{2.211501in}}%
\pgfpathlineto{\pgfqpoint{2.456649in}{2.203017in}}%
\pgfpathlineto{\pgfqpoint{2.457921in}{2.132312in}}%
\pgfpathlineto{\pgfqpoint{2.459107in}{2.208414in}}%
\pgfpathlineto{\pgfqpoint{2.459523in}{2.204940in}}%
\pgfpathlineto{\pgfqpoint{2.460556in}{2.136673in}}%
\pgfpathlineto{\pgfqpoint{2.461037in}{2.138398in}}%
\pgfpathlineto{\pgfqpoint{2.462375in}{2.208006in}}%
\pgfpathlineto{\pgfqpoint{2.463653in}{2.134926in}}%
\pgfpathlineto{\pgfqpoint{2.464074in}{2.143515in}}%
\pgfpathlineto{\pgfqpoint{2.465365in}{2.205664in}}%
\pgfpathlineto{\pgfqpoint{2.466564in}{2.137042in}}%
\pgfpathlineto{\pgfqpoint{2.466981in}{2.143021in}}%
\pgfpathlineto{\pgfqpoint{2.468045in}{2.204424in}}%
\pgfpathlineto{\pgfqpoint{2.468539in}{2.197240in}}%
\pgfpathlineto{\pgfqpoint{2.469493in}{2.140292in}}%
\pgfpathlineto{\pgfqpoint{2.469962in}{2.142811in}}%
\pgfpathlineto{\pgfqpoint{2.471364in}{2.202202in}}%
\pgfpathlineto{\pgfqpoint{2.472678in}{2.138879in}}%
\pgfpathlineto{\pgfqpoint{2.473108in}{2.152107in}}%
\pgfpathlineto{\pgfqpoint{2.473975in}{2.201145in}}%
\pgfpathlineto{\pgfqpoint{2.474404in}{2.198477in}}%
\pgfpathlineto{\pgfqpoint{2.475745in}{2.141029in}}%
\pgfpathlineto{\pgfqpoint{2.477005in}{2.201172in}}%
\pgfpathlineto{\pgfqpoint{2.477449in}{2.194286in}}%
\pgfpathlineto{\pgfqpoint{2.478562in}{2.141870in}}%
\pgfpathlineto{\pgfqpoint{2.479066in}{2.152585in}}%
\pgfpathlineto{\pgfqpoint{2.480101in}{2.200560in}}%
\pgfpathlineto{\pgfqpoint{2.480590in}{2.189130in}}%
\pgfpathlineto{\pgfqpoint{2.481732in}{2.142869in}}%
\pgfpathlineto{\pgfqpoint{2.483139in}{2.199405in}}%
\pgfpathlineto{\pgfqpoint{2.483602in}{2.189023in}}%
\pgfpathlineto{\pgfqpoint{2.484620in}{2.144220in}}%
\pgfpathlineto{\pgfqpoint{2.485111in}{2.153952in}}%
\pgfpathlineto{\pgfqpoint{2.486231in}{2.198360in}}%
\pgfpathlineto{\pgfqpoint{2.487702in}{2.145094in}}%
\pgfpathlineto{\pgfqpoint{2.488181in}{2.156309in}}%
\pgfpathlineto{\pgfqpoint{2.489181in}{2.197178in}}%
\pgfpathlineto{\pgfqpoint{2.490219in}{2.159605in}}%
\pgfpathlineto{\pgfqpoint{2.490732in}{2.146178in}}%
\pgfpathlineto{\pgfqpoint{2.492250in}{2.196204in}}%
\pgfpathlineto{\pgfqpoint{2.492749in}{2.184169in}}%
\pgfpathlineto{\pgfqpoint{2.493743in}{2.147250in}}%
\pgfpathlineto{\pgfqpoint{2.495264in}{2.195178in}}%
\pgfpathlineto{\pgfqpoint{2.495772in}{2.183583in}}%
\pgfpathlineto{\pgfqpoint{2.496786in}{2.148248in}}%
\pgfpathlineto{\pgfqpoint{2.498294in}{2.194207in}}%
\pgfpathlineto{\pgfqpoint{2.498801in}{2.182759in}}%
\pgfpathlineto{\pgfqpoint{2.499806in}{2.149224in}}%
\pgfpathlineto{\pgfqpoint{2.501326in}{2.193263in}}%
\pgfpathlineto{\pgfqpoint{2.501840in}{2.181699in}}%
\pgfpathlineto{\pgfqpoint{2.502840in}{2.150170in}}%
\pgfpathlineto{\pgfqpoint{2.504354in}{2.192354in}}%
\pgfpathlineto{\pgfqpoint{2.504871in}{2.180966in}}%
\pgfpathlineto{\pgfqpoint{2.505870in}{2.151081in}}%
\pgfpathlineto{\pgfqpoint{2.507395in}{2.191454in}}%
\pgfpathlineto{\pgfqpoint{2.507918in}{2.179768in}}%
\pgfpathlineto{\pgfqpoint{2.508915in}{2.151989in}}%
\pgfpathlineto{\pgfqpoint{2.510447in}{2.190536in}}%
\pgfpathlineto{\pgfqpoint{2.512031in}{2.153249in}}%
\pgfpathlineto{\pgfqpoint{2.513598in}{2.188946in}}%
\pgfpathlineto{\pgfqpoint{2.515117in}{2.154470in}}%
\pgfpathlineto{\pgfqpoint{2.516648in}{2.188040in}}%
\pgfpathlineto{\pgfqpoint{2.518173in}{2.155421in}}%
\pgfpathlineto{\pgfqpoint{2.519716in}{2.186974in}}%
\pgfpathlineto{\pgfqpoint{2.521255in}{2.156620in}}%
\pgfpathlineto{\pgfqpoint{2.522821in}{2.185466in}}%
\pgfpathlineto{\pgfqpoint{2.523869in}{2.157443in}}%
\pgfpathlineto{\pgfqpoint{2.524396in}{2.158630in}}%
\pgfpathlineto{\pgfqpoint{2.525498in}{2.186304in}}%
\pgfpathlineto{\pgfqpoint{2.526013in}{2.182584in}}%
\pgfpathlineto{\pgfqpoint{2.527192in}{2.156405in}}%
\pgfpathlineto{\pgfqpoint{2.528792in}{2.186272in}}%
\pgfpathlineto{\pgfqpoint{2.529251in}{2.178714in}}%
\pgfpathlineto{\pgfqpoint{2.530358in}{2.157241in}}%
\pgfpathlineto{\pgfqpoint{2.531957in}{2.185054in}}%
\pgfpathlineto{\pgfqpoint{2.533079in}{2.159798in}}%
\pgfpathlineto{\pgfqpoint{2.533649in}{2.160006in}}%
\pgfpathlineto{\pgfqpoint{2.534913in}{2.185114in}}%
\pgfpathlineto{\pgfqpoint{2.536195in}{2.160929in}}%
\pgfpathlineto{\pgfqpoint{2.536915in}{2.161538in}}%
\pgfpathlineto{\pgfqpoint{2.538194in}{2.184539in}}%
\pgfpathlineto{\pgfqpoint{2.539431in}{2.160878in}}%
\pgfpathlineto{\pgfqpoint{2.540033in}{2.161073in}}%
\pgfpathlineto{\pgfqpoint{2.541358in}{2.184045in}}%
\pgfpathlineto{\pgfqpoint{2.542636in}{2.161020in}}%
\pgfpathlineto{\pgfqpoint{2.543348in}{2.163112in}}%
\pgfpathlineto{\pgfqpoint{2.544683in}{2.183150in}}%
\pgfpathlineto{\pgfqpoint{2.545946in}{2.160563in}}%
\pgfpathlineto{\pgfqpoint{2.546644in}{2.164771in}}%
\pgfpathlineto{\pgfqpoint{2.547390in}{2.179889in}}%
\pgfpathlineto{\pgfqpoint{2.548335in}{2.178784in}}%
\pgfpathlineto{\pgfqpoint{2.549852in}{2.162096in}}%
\pgfpathlineto{\pgfqpoint{2.551300in}{2.182401in}}%
\pgfpathlineto{\pgfqpoint{2.552653in}{2.161031in}}%
\pgfpathlineto{\pgfqpoint{2.553354in}{2.166144in}}%
\pgfpathlineto{\pgfqpoint{2.554311in}{2.181309in}}%
\pgfpathlineto{\pgfqpoint{2.555228in}{2.175331in}}%
\pgfpathlineto{\pgfqpoint{2.556071in}{2.161774in}}%
\pgfpathlineto{\pgfqpoint{2.558221in}{2.181862in}}%
\pgfpathlineto{\pgfqpoint{2.560189in}{2.161512in}}%
\pgfpathlineto{\pgfqpoint{2.562032in}{2.180986in}}%
\pgfpathlineto{\pgfqpoint{2.563728in}{2.161956in}}%
\pgfpathlineto{\pgfqpoint{2.565384in}{2.181026in}}%
\pgfpathlineto{\pgfqpoint{2.566989in}{2.161987in}}%
\pgfpathlineto{\pgfqpoint{2.568748in}{2.180740in}}%
\pgfpathlineto{\pgfqpoint{2.570510in}{2.162314in}}%
\pgfpathlineto{\pgfqpoint{2.572422in}{2.180273in}}%
\pgfpathlineto{\pgfqpoint{2.574281in}{2.162966in}}%
\pgfpathlineto{\pgfqpoint{2.576074in}{2.179575in}}%
\pgfpathlineto{\pgfqpoint{2.577791in}{2.163277in}}%
\pgfpathlineto{\pgfqpoint{2.579604in}{2.179204in}}%
\pgfpathlineto{\pgfqpoint{2.581334in}{2.163644in}}%
\pgfpathlineto{\pgfqpoint{2.583207in}{2.178585in}}%
\pgfpathlineto{\pgfqpoint{2.585009in}{2.164471in}}%
\pgfpathlineto{\pgfqpoint{2.586813in}{2.177851in}}%
\pgfpathlineto{\pgfqpoint{2.588644in}{2.165331in}}%
\pgfpathlineto{\pgfqpoint{2.589474in}{2.175497in}}%
\pgfpathlineto{\pgfqpoint{2.590563in}{2.176319in}}%
\pgfpathlineto{\pgfqpoint{2.591608in}{2.165167in}}%
\pgfpathlineto{\pgfqpoint{2.592735in}{2.169029in}}%
\pgfpathlineto{\pgfqpoint{2.593788in}{2.178349in}}%
\pgfpathlineto{\pgfqpoint{2.595169in}{2.169058in}}%
\pgfpathlineto{\pgfqpoint{2.596838in}{2.167021in}}%
\pgfpathlineto{\pgfqpoint{2.598406in}{2.178134in}}%
\pgfpathlineto{\pgfqpoint{2.601114in}{2.165983in}}%
\pgfpathlineto{\pgfqpoint{2.603616in}{2.172429in}}%
\pgfpathlineto{\pgfqpoint{2.606029in}{2.174431in}}%
\pgfpathlineto{\pgfqpoint{2.609040in}{2.165170in}}%
\pgfpathlineto{\pgfqpoint{2.613936in}{2.177670in}}%
\pgfpathlineto{\pgfqpoint{2.618832in}{2.164379in}}%
\pgfpathlineto{\pgfqpoint{2.623728in}{2.176442in}}%
\pgfpathlineto{\pgfqpoint{2.628624in}{2.166665in}}%
\pgfpathlineto{\pgfqpoint{2.633520in}{2.172254in}}%
\pgfpathlineto{\pgfqpoint{2.638416in}{2.170652in}}%
\pgfpathlineto{\pgfqpoint{2.643312in}{2.138064in}}%
\pgfpathlineto{\pgfqpoint{2.648208in}{2.118517in}}%
\pgfpathlineto{\pgfqpoint{2.653104in}{2.055213in}}%
\pgfpathlineto{\pgfqpoint{2.656826in}{1.988418in}}%
\pgfpathmoveto{\pgfqpoint{2.889786in}{1.988418in}}%
\pgfpathlineto{\pgfqpoint{2.890691in}{2.179729in}}%
\pgfpathlineto{\pgfqpoint{2.892294in}{2.302043in}}%
\pgfpathlineto{\pgfqpoint{2.894741in}{2.063207in}}%
\pgfpathlineto{\pgfqpoint{2.895970in}{2.297296in}}%
\pgfpathlineto{\pgfqpoint{2.896452in}{2.253194in}}%
\pgfpathlineto{\pgfqpoint{2.897696in}{2.056729in}}%
\pgfpathlineto{\pgfqpoint{2.898965in}{2.293228in}}%
\pgfpathlineto{\pgfqpoint{2.899363in}{2.256179in}}%
\pgfpathlineto{\pgfqpoint{2.900469in}{2.054169in}}%
\pgfpathlineto{\pgfqpoint{2.902012in}{2.285129in}}%
\pgfpathlineto{\pgfqpoint{2.902438in}{2.218707in}}%
\pgfpathlineto{\pgfqpoint{2.903240in}{2.063075in}}%
\pgfpathlineto{\pgfqpoint{2.903614in}{2.073472in}}%
\pgfpathlineto{\pgfqpoint{2.904795in}{2.282992in}}%
\pgfpathlineto{\pgfqpoint{2.905179in}{2.251805in}}%
\pgfpathlineto{\pgfqpoint{2.906239in}{2.063975in}}%
\pgfpathlineto{\pgfqpoint{2.907091in}{2.186772in}}%
\pgfpathlineto{\pgfqpoint{2.907839in}{2.275801in}}%
\pgfpathlineto{\pgfqpoint{2.909342in}{2.073459in}}%
\pgfpathlineto{\pgfqpoint{2.910684in}{2.273748in}}%
\pgfpathlineto{\pgfqpoint{2.911088in}{2.235288in}}%
\pgfpathlineto{\pgfqpoint{2.912209in}{2.073940in}}%
\pgfpathlineto{\pgfqpoint{2.913779in}{2.261071in}}%
\pgfpathlineto{\pgfqpoint{2.915066in}{2.076654in}}%
\pgfpathlineto{\pgfqpoint{2.915459in}{2.107587in}}%
\pgfpathlineto{\pgfqpoint{2.916619in}{2.263128in}}%
\pgfpathlineto{\pgfqpoint{2.918127in}{2.085367in}}%
\pgfpathlineto{\pgfqpoint{2.919470in}{2.261142in}}%
\pgfpathlineto{\pgfqpoint{2.919874in}{2.227182in}}%
\pgfpathlineto{\pgfqpoint{2.921001in}{2.085986in}}%
\pgfpathlineto{\pgfqpoint{2.922183in}{2.249776in}}%
\pgfpathlineto{\pgfqpoint{2.922568in}{2.249728in}}%
\pgfpathlineto{\pgfqpoint{2.923872in}{2.088199in}}%
\pgfpathlineto{\pgfqpoint{2.924265in}{2.117960in}}%
\pgfpathlineto{\pgfqpoint{2.925428in}{2.250895in}}%
\pgfpathlineto{\pgfqpoint{2.926578in}{2.097538in}}%
\pgfpathlineto{\pgfqpoint{2.926950in}{2.097940in}}%
\pgfpathlineto{\pgfqpoint{2.928335in}{2.249043in}}%
\pgfpathlineto{\pgfqpoint{2.929883in}{2.099994in}}%
\pgfpathlineto{\pgfqpoint{2.931165in}{2.246765in}}%
\pgfpathlineto{\pgfqpoint{2.931542in}{2.225056in}}%
\pgfpathlineto{\pgfqpoint{2.932614in}{2.098307in}}%
\pgfpathlineto{\pgfqpoint{2.933468in}{2.185097in}}%
\pgfpathlineto{\pgfqpoint{2.934216in}{2.241051in}}%
\pgfpathlineto{\pgfqpoint{2.935368in}{2.106430in}}%
\pgfpathlineto{\pgfqpoint{2.935738in}{2.106963in}}%
\pgfpathlineto{\pgfqpoint{2.937125in}{2.239381in}}%
\pgfpathlineto{\pgfqpoint{2.938681in}{2.109255in}}%
\pgfpathlineto{\pgfqpoint{2.939965in}{2.237527in}}%
\pgfpathlineto{\pgfqpoint{2.940341in}{2.217440in}}%
\pgfpathlineto{\pgfqpoint{2.941413in}{2.107209in}}%
\pgfpathlineto{\pgfqpoint{2.942271in}{2.185468in}}%
\pgfpathlineto{\pgfqpoint{2.943021in}{2.231829in}}%
\pgfpathlineto{\pgfqpoint{2.944175in}{2.113573in}}%
\pgfpathlineto{\pgfqpoint{2.944550in}{2.116066in}}%
\pgfpathlineto{\pgfqpoint{2.945934in}{2.230359in}}%
\pgfpathlineto{\pgfqpoint{2.947089in}{2.117292in}}%
\pgfpathlineto{\pgfqpoint{2.947476in}{2.117536in}}%
\pgfpathlineto{\pgfqpoint{2.948770in}{2.229420in}}%
\pgfpathlineto{\pgfqpoint{2.949145in}{2.209970in}}%
\pgfpathlineto{\pgfqpoint{2.950209in}{2.115045in}}%
\pgfpathlineto{\pgfqpoint{2.951075in}{2.186212in}}%
\pgfpathlineto{\pgfqpoint{2.951829in}{2.223517in}}%
\pgfpathlineto{\pgfqpoint{2.952987in}{2.119714in}}%
\pgfpathlineto{\pgfqpoint{2.953369in}{2.124626in}}%
\pgfpathlineto{\pgfqpoint{2.954688in}{2.223596in}}%
\pgfpathlineto{\pgfqpoint{2.956233in}{2.123398in}}%
\pgfpathlineto{\pgfqpoint{2.957486in}{2.222011in}}%
\pgfpathlineto{\pgfqpoint{2.957863in}{2.209168in}}%
\pgfpathlineto{\pgfqpoint{2.958995in}{2.122065in}}%
\pgfpathlineto{\pgfqpoint{2.960597in}{2.217196in}}%
\pgfpathlineto{\pgfqpoint{2.961045in}{2.184821in}}%
\pgfpathlineto{\pgfqpoint{2.961760in}{2.126161in}}%
\pgfpathlineto{\pgfqpoint{2.962136in}{2.130192in}}%
\pgfpathlineto{\pgfqpoint{2.963477in}{2.216839in}}%
\pgfpathlineto{\pgfqpoint{2.964644in}{2.129770in}}%
\pgfpathlineto{\pgfqpoint{2.965035in}{2.130314in}}%
\pgfpathlineto{\pgfqpoint{2.966316in}{2.215874in}}%
\pgfpathlineto{\pgfqpoint{2.966692in}{2.201327in}}%
\pgfpathlineto{\pgfqpoint{2.967791in}{2.128324in}}%
\pgfpathlineto{\pgfqpoint{2.969417in}{2.210187in}}%
\pgfpathlineto{\pgfqpoint{2.970603in}{2.130674in}}%
\pgfpathlineto{\pgfqpoint{2.970984in}{2.138327in}}%
\pgfpathlineto{\pgfqpoint{2.972210in}{2.211706in}}%
\pgfpathlineto{\pgfqpoint{2.973782in}{2.134706in}}%
\pgfpathlineto{\pgfqpoint{2.975007in}{2.209925in}}%
\pgfpathlineto{\pgfqpoint{2.975381in}{2.201795in}}%
\pgfpathlineto{\pgfqpoint{2.976583in}{2.133828in}}%
\pgfpathlineto{\pgfqpoint{2.978188in}{2.205354in}}%
\pgfpathlineto{\pgfqpoint{2.979375in}{2.135781in}}%
\pgfpathlineto{\pgfqpoint{2.979752in}{2.142377in}}%
\pgfpathlineto{\pgfqpoint{2.981001in}{2.206440in}}%
\pgfpathlineto{\pgfqpoint{2.982187in}{2.139743in}}%
\pgfpathlineto{\pgfqpoint{2.982581in}{2.140060in}}%
\pgfpathlineto{\pgfqpoint{2.983836in}{2.205383in}}%
\pgfpathlineto{\pgfqpoint{2.984211in}{2.195544in}}%
\pgfpathlineto{\pgfqpoint{2.985375in}{2.138770in}}%
\pgfpathlineto{\pgfqpoint{2.986602in}{2.201745in}}%
\pgfpathlineto{\pgfqpoint{2.987005in}{2.199863in}}%
\pgfpathlineto{\pgfqpoint{2.988252in}{2.139856in}}%
\pgfpathlineto{\pgfqpoint{2.988635in}{2.150391in}}%
\pgfpathlineto{\pgfqpoint{2.989796in}{2.201776in}}%
\pgfpathlineto{\pgfqpoint{2.991032in}{2.142440in}}%
\pgfpathlineto{\pgfqpoint{2.991434in}{2.146184in}}%
\pgfpathlineto{\pgfqpoint{2.992681in}{2.200945in}}%
\pgfpathlineto{\pgfqpoint{2.994178in}{2.143222in}}%
\pgfpathlineto{\pgfqpoint{2.994630in}{2.159812in}}%
\pgfpathlineto{\pgfqpoint{2.995453in}{2.199093in}}%
\pgfpathlineto{\pgfqpoint{2.995862in}{2.193829in}}%
\pgfpathlineto{\pgfqpoint{2.997083in}{2.144130in}}%
\pgfpathlineto{\pgfqpoint{2.998595in}{2.197600in}}%
\pgfpathlineto{\pgfqpoint{2.999054in}{2.180567in}}%
\pgfpathlineto{\pgfqpoint{2.999882in}{2.145414in}}%
\pgfpathlineto{\pgfqpoint{3.000291in}{2.151868in}}%
\pgfpathlineto{\pgfqpoint{3.001508in}{2.196778in}}%
\pgfpathlineto{\pgfqpoint{3.002987in}{2.147221in}}%
\pgfpathlineto{\pgfqpoint{3.004309in}{2.196080in}}%
\pgfpathlineto{\pgfqpoint{3.004724in}{2.188368in}}%
\pgfpathlineto{\pgfqpoint{3.005930in}{2.148267in}}%
\pgfpathlineto{\pgfqpoint{3.007429in}{2.193499in}}%
\pgfpathlineto{\pgfqpoint{3.008790in}{2.148647in}}%
\pgfpathlineto{\pgfqpoint{3.010022in}{2.192014in}}%
\pgfpathlineto{\pgfqpoint{3.010430in}{2.191619in}}%
\pgfpathlineto{\pgfqpoint{3.011634in}{2.149744in}}%
\pgfpathlineto{\pgfqpoint{3.012007in}{2.154767in}}%
\pgfpathlineto{\pgfqpoint{3.013316in}{2.191637in}}%
\pgfpathlineto{\pgfqpoint{3.014577in}{2.150731in}}%
\pgfpathlineto{\pgfqpoint{3.014990in}{2.156643in}}%
\pgfpathlineto{\pgfqpoint{3.016224in}{2.191149in}}%
\pgfpathlineto{\pgfqpoint{3.017725in}{2.153027in}}%
\pgfpathlineto{\pgfqpoint{3.019129in}{2.190663in}}%
\pgfpathlineto{\pgfqpoint{3.020374in}{2.153264in}}%
\pgfpathlineto{\pgfqpoint{3.020784in}{2.155688in}}%
\pgfpathlineto{\pgfqpoint{3.022078in}{2.189778in}}%
\pgfpathlineto{\pgfqpoint{3.023645in}{2.154988in}}%
\pgfpathlineto{\pgfqpoint{3.025018in}{2.188975in}}%
\pgfpathlineto{\pgfqpoint{3.026269in}{2.154868in}}%
\pgfpathlineto{\pgfqpoint{3.026682in}{2.157169in}}%
\pgfpathlineto{\pgfqpoint{3.027985in}{2.188071in}}%
\pgfpathlineto{\pgfqpoint{3.029155in}{2.156279in}}%
\pgfpathlineto{\pgfqpoint{3.029569in}{2.156852in}}%
\pgfpathlineto{\pgfqpoint{3.030995in}{2.186797in}}%
\pgfpathlineto{\pgfqpoint{3.032279in}{2.155760in}}%
\pgfpathlineto{\pgfqpoint{3.032700in}{2.160730in}}%
\pgfpathlineto{\pgfqpoint{3.033964in}{2.185847in}}%
\pgfpathlineto{\pgfqpoint{3.035216in}{2.156532in}}%
\pgfpathlineto{\pgfqpoint{3.035627in}{2.160805in}}%
\pgfpathlineto{\pgfqpoint{3.036553in}{2.184975in}}%
\pgfpathlineto{\pgfqpoint{3.037011in}{2.183938in}}%
\pgfpathlineto{\pgfqpoint{3.038390in}{2.157915in}}%
\pgfpathlineto{\pgfqpoint{3.039688in}{2.185467in}}%
\pgfpathlineto{\pgfqpoint{3.040115in}{2.180682in}}%
\pgfpathlineto{\pgfqpoint{3.041390in}{2.159026in}}%
\pgfpathlineto{\pgfqpoint{3.042693in}{2.184811in}}%
\pgfpathlineto{\pgfqpoint{3.043991in}{2.159165in}}%
\pgfpathlineto{\pgfqpoint{3.044442in}{2.160884in}}%
\pgfpathlineto{\pgfqpoint{3.045468in}{2.183236in}}%
\pgfpathlineto{\pgfqpoint{3.045962in}{2.181667in}}%
\pgfpathlineto{\pgfqpoint{3.046922in}{2.160500in}}%
\pgfpathlineto{\pgfqpoint{3.047405in}{2.160706in}}%
\pgfpathlineto{\pgfqpoint{3.048773in}{2.183149in}}%
\pgfpathlineto{\pgfqpoint{3.050080in}{2.159694in}}%
\pgfpathlineto{\pgfqpoint{3.050511in}{2.163345in}}%
\pgfpathlineto{\pgfqpoint{3.051820in}{2.181868in}}%
\pgfpathlineto{\pgfqpoint{3.053160in}{2.160295in}}%
\pgfpathlineto{\pgfqpoint{3.054485in}{2.182195in}}%
\pgfpathlineto{\pgfqpoint{3.054950in}{2.179234in}}%
\pgfpathlineto{\pgfqpoint{3.056067in}{2.160754in}}%
\pgfpathlineto{\pgfqpoint{3.057640in}{2.181883in}}%
\pgfpathlineto{\pgfqpoint{3.058138in}{2.176536in}}%
\pgfpathlineto{\pgfqpoint{3.059097in}{2.161239in}}%
\pgfpathlineto{\pgfqpoint{3.060099in}{2.175665in}}%
\pgfpathlineto{\pgfqpoint{3.060608in}{2.181343in}}%
\pgfpathlineto{\pgfqpoint{3.062154in}{2.161691in}}%
\pgfpathlineto{\pgfqpoint{3.062658in}{2.166482in}}%
\pgfpathlineto{\pgfqpoint{3.063658in}{2.180865in}}%
\pgfpathlineto{\pgfqpoint{3.065161in}{2.162153in}}%
\pgfpathlineto{\pgfqpoint{3.065663in}{2.166511in}}%
\pgfpathlineto{\pgfqpoint{3.066686in}{2.180382in}}%
\pgfpathlineto{\pgfqpoint{3.068204in}{2.162587in}}%
\pgfpathlineto{\pgfqpoint{3.068713in}{2.167117in}}%
\pgfpathlineto{\pgfqpoint{3.069713in}{2.179913in}}%
\pgfpathlineto{\pgfqpoint{3.071231in}{2.163008in}}%
\pgfpathlineto{\pgfqpoint{3.071746in}{2.167464in}}%
\pgfpathlineto{\pgfqpoint{3.072750in}{2.179453in}}%
\pgfpathlineto{\pgfqpoint{3.074268in}{2.163417in}}%
\pgfpathlineto{\pgfqpoint{3.074789in}{2.167889in}}%
\pgfpathlineto{\pgfqpoint{3.075797in}{2.178995in}}%
\pgfpathlineto{\pgfqpoint{3.077327in}{2.163839in}}%
\pgfpathlineto{\pgfqpoint{3.078921in}{2.178372in}}%
\pgfpathlineto{\pgfqpoint{3.080481in}{2.164517in}}%
\pgfpathlineto{\pgfqpoint{3.082013in}{2.177764in}}%
\pgfpathlineto{\pgfqpoint{3.083543in}{2.164966in}}%
\pgfpathlineto{\pgfqpoint{3.085087in}{2.177213in}}%
\pgfpathlineto{\pgfqpoint{3.086641in}{2.165556in}}%
\pgfpathlineto{\pgfqpoint{3.087682in}{2.176561in}}%
\pgfpathlineto{\pgfqpoint{3.088221in}{2.176357in}}%
\pgfpathlineto{\pgfqpoint{3.089309in}{2.165512in}}%
\pgfpathlineto{\pgfqpoint{3.089829in}{2.166669in}}%
\pgfpathlineto{\pgfqpoint{3.090999in}{2.177002in}}%
\pgfpathlineto{\pgfqpoint{3.092629in}{2.165466in}}%
\pgfpathlineto{\pgfqpoint{3.093099in}{2.168428in}}%
\pgfpathlineto{\pgfqpoint{3.094246in}{2.176471in}}%
\pgfpathlineto{\pgfqpoint{3.095917in}{2.166409in}}%
\pgfpathlineto{\pgfqpoint{3.097152in}{2.176203in}}%
\pgfpathlineto{\pgfqpoint{3.098528in}{2.167026in}}%
\pgfpathlineto{\pgfqpoint{3.099290in}{2.167459in}}%
\pgfpathlineto{\pgfqpoint{3.100608in}{2.175845in}}%
\pgfpathlineto{\pgfqpoint{3.101867in}{2.166711in}}%
\pgfpathlineto{\pgfqpoint{3.102565in}{2.168007in}}%
\pgfpathlineto{\pgfqpoint{3.104097in}{2.174705in}}%
\pgfpathlineto{\pgfqpoint{3.105442in}{2.166472in}}%
\pgfpathlineto{\pgfqpoint{3.106824in}{2.174961in}}%
\pgfpathlineto{\pgfqpoint{3.107587in}{2.173018in}}%
\pgfpathlineto{\pgfqpoint{3.108387in}{2.167384in}}%
\pgfpathlineto{\pgfqpoint{3.109345in}{2.168821in}}%
\pgfpathlineto{\pgfqpoint{3.110309in}{2.174613in}}%
\pgfpathlineto{\pgfqpoint{3.111288in}{2.171984in}}%
\pgfpathlineto{\pgfqpoint{3.112170in}{2.166965in}}%
\pgfpathlineto{\pgfqpoint{3.114383in}{2.174523in}}%
\pgfpathlineto{\pgfqpoint{3.116358in}{2.167249in}}%
\pgfpathlineto{\pgfqpoint{3.118127in}{2.173871in}}%
\pgfpathlineto{\pgfqpoint{3.119762in}{2.167268in}}%
\pgfpathlineto{\pgfqpoint{3.121493in}{2.173645in}}%
\pgfpathlineto{\pgfqpoint{3.123300in}{2.167445in}}%
\pgfpathlineto{\pgfqpoint{3.125266in}{2.172785in}}%
\pgfpathlineto{\pgfqpoint{3.126305in}{2.168424in}}%
\pgfpathlineto{\pgfqpoint{3.127329in}{2.168279in}}%
\pgfpathlineto{\pgfqpoint{3.128303in}{2.171974in}}%
\pgfpathlineto{\pgfqpoint{3.129930in}{2.169889in}}%
\pgfpathlineto{\pgfqpoint{3.133426in}{2.162533in}}%
\pgfpathlineto{\pgfqpoint{3.137408in}{2.127634in}}%
\pgfpathlineto{\pgfqpoint{3.141370in}{2.063555in}}%
\pgfpathlineto{\pgfqpoint{3.145023in}{1.988418in}}%
\pgfpathlineto{\pgfqpoint{3.145023in}{1.988418in}}%
\pgfusepath{stroke}%
\end{pgfscope}%
\begin{pgfscope}%
\pgfsetrectcap%
\pgfsetmiterjoin%
\pgfsetlinewidth{0.803000pt}%
\definecolor{currentstroke}{rgb}{0.000000,0.000000,0.000000}%
\pgfsetstrokecolor{currentstroke}%
\pgfsetdash{}{0pt}%
\pgfpathmoveto{\pgfqpoint{2.197902in}{1.995363in}}%
\pgfpathlineto{\pgfqpoint{2.197902in}{2.321613in}}%
\pgfusepath{stroke}%
\end{pgfscope}%
\begin{pgfscope}%
\pgfsetrectcap%
\pgfsetmiterjoin%
\pgfsetlinewidth{0.803000pt}%
\definecolor{currentstroke}{rgb}{0.000000,0.000000,0.000000}%
\pgfsetstrokecolor{currentstroke}%
\pgfsetdash{}{0pt}%
\pgfpathmoveto{\pgfqpoint{3.177102in}{1.995363in}}%
\pgfpathlineto{\pgfqpoint{3.177102in}{2.321613in}}%
\pgfusepath{stroke}%
\end{pgfscope}%
\begin{pgfscope}%
\pgfsetrectcap%
\pgfsetmiterjoin%
\pgfsetlinewidth{0.803000pt}%
\definecolor{currentstroke}{rgb}{0.000000,0.000000,0.000000}%
\pgfsetstrokecolor{currentstroke}%
\pgfsetdash{}{0pt}%
\pgfpathmoveto{\pgfqpoint{2.197902in}{1.995363in}}%
\pgfpathlineto{\pgfqpoint{3.177102in}{1.995363in}}%
\pgfusepath{stroke}%
\end{pgfscope}%
\begin{pgfscope}%
\pgfsetrectcap%
\pgfsetmiterjoin%
\pgfsetlinewidth{0.803000pt}%
\definecolor{currentstroke}{rgb}{0.000000,0.000000,0.000000}%
\pgfsetstrokecolor{currentstroke}%
\pgfsetdash{}{0pt}%
\pgfpathmoveto{\pgfqpoint{2.197902in}{2.321613in}}%
\pgfpathlineto{\pgfqpoint{3.177102in}{2.321613in}}%
\pgfusepath{stroke}%
\end{pgfscope}%
\begin{pgfscope}%
\pgfsetbuttcap%
\pgfsetmiterjoin%
\definecolor{currentfill}{rgb}{1.000000,1.000000,1.000000}%
\pgfsetfillcolor{currentfill}%
\pgfsetlinewidth{0.000000pt}%
\definecolor{currentstroke}{rgb}{0.000000,0.000000,0.000000}%
\pgfsetstrokecolor{currentstroke}%
\pgfsetstrokeopacity{0.000000}%
\pgfsetdash{}{0pt}%
\pgfpathmoveto{\pgfqpoint{2.197902in}{0.579321in}}%
\pgfpathlineto{\pgfqpoint{3.177102in}{0.579321in}}%
\pgfpathlineto{\pgfqpoint{3.177102in}{0.905571in}}%
\pgfpathlineto{\pgfqpoint{2.197902in}{0.905571in}}%
\pgfpathlineto{\pgfqpoint{2.197902in}{0.579321in}}%
\pgfpathclose%
\pgfusepath{fill}%
\end{pgfscope}%
\begin{pgfscope}%
\pgfpathrectangle{\pgfqpoint{2.197902in}{0.579321in}}{\pgfqpoint{0.979200in}{0.326250in}}%
\pgfusepath{clip}%
\pgfsetbuttcap%
\pgfsetroundjoin%
\pgfsetlinewidth{0.401500pt}%
\definecolor{currentstroke}{rgb}{0.690196,0.690196,0.690196}%
\pgfsetstrokecolor{currentstroke}%
\pgfsetstrokeopacity{0.250000}%
\pgfsetdash{{1.480000pt}{0.640000pt}}{0.000000pt}%
\pgfpathmoveto{\pgfqpoint{2.197902in}{0.579321in}}%
\pgfpathlineto{\pgfqpoint{2.197902in}{0.905571in}}%
\pgfusepath{stroke}%
\end{pgfscope}%
\begin{pgfscope}%
\pgfsetbuttcap%
\pgfsetroundjoin%
\definecolor{currentfill}{rgb}{0.000000,0.000000,0.000000}%
\pgfsetfillcolor{currentfill}%
\pgfsetlinewidth{0.803000pt}%
\definecolor{currentstroke}{rgb}{0.000000,0.000000,0.000000}%
\pgfsetstrokecolor{currentstroke}%
\pgfsetdash{}{0pt}%
\pgfsys@defobject{currentmarker}{\pgfqpoint{0.000000in}{0.000000in}}{\pgfqpoint{0.000000in}{0.027778in}}{%
\pgfpathmoveto{\pgfqpoint{0.000000in}{0.000000in}}%
\pgfpathlineto{\pgfqpoint{0.000000in}{0.027778in}}%
\pgfusepath{stroke,fill}%
}%
\begin{pgfscope}%
\pgfsys@transformshift{2.197902in}{0.579321in}%
\pgfsys@useobject{currentmarker}{}%
\end{pgfscope}%
\end{pgfscope}%
\begin{pgfscope}%
\definecolor{textcolor}{rgb}{0.000000,0.000000,0.000000}%
\pgfsetstrokecolor{textcolor}%
\pgfsetfillcolor{textcolor}%
\pgftext[x=2.197902in,y=0.565432in,,top]{\color{textcolor}{\rmfamily\fontsize{5.000000}{6.000000}\selectfont\catcode`\^=\active\def^{\ifmmode\sp\else\^{}\fi}\catcode`\%=\active\def%{\%}4.980}}%
\end{pgfscope}%
\begin{pgfscope}%
\pgfpathrectangle{\pgfqpoint{2.197902in}{0.579321in}}{\pgfqpoint{0.979200in}{0.326250in}}%
\pgfusepath{clip}%
\pgfsetbuttcap%
\pgfsetroundjoin%
\pgfsetlinewidth{0.401500pt}%
\definecolor{currentstroke}{rgb}{0.690196,0.690196,0.690196}%
\pgfsetstrokecolor{currentstroke}%
\pgfsetstrokeopacity{0.250000}%
\pgfsetdash{{1.480000pt}{0.640000pt}}{0.000000pt}%
\pgfpathmoveto{\pgfqpoint{2.442702in}{0.579321in}}%
\pgfpathlineto{\pgfqpoint{2.442702in}{0.905571in}}%
\pgfusepath{stroke}%
\end{pgfscope}%
\begin{pgfscope}%
\pgfsetbuttcap%
\pgfsetroundjoin%
\definecolor{currentfill}{rgb}{0.000000,0.000000,0.000000}%
\pgfsetfillcolor{currentfill}%
\pgfsetlinewidth{0.803000pt}%
\definecolor{currentstroke}{rgb}{0.000000,0.000000,0.000000}%
\pgfsetstrokecolor{currentstroke}%
\pgfsetdash{}{0pt}%
\pgfsys@defobject{currentmarker}{\pgfqpoint{0.000000in}{0.000000in}}{\pgfqpoint{0.000000in}{0.027778in}}{%
\pgfpathmoveto{\pgfqpoint{0.000000in}{0.000000in}}%
\pgfpathlineto{\pgfqpoint{0.000000in}{0.027778in}}%
\pgfusepath{stroke,fill}%
}%
\begin{pgfscope}%
\pgfsys@transformshift{2.442702in}{0.579321in}%
\pgfsys@useobject{currentmarker}{}%
\end{pgfscope}%
\end{pgfscope}%
\begin{pgfscope}%
\definecolor{textcolor}{rgb}{0.000000,0.000000,0.000000}%
\pgfsetstrokecolor{textcolor}%
\pgfsetfillcolor{textcolor}%
\pgftext[x=2.442702in,y=0.565432in,,top]{\color{textcolor}{\rmfamily\fontsize{5.000000}{6.000000}\selectfont\catcode`\^=\active\def^{\ifmmode\sp\else\^{}\fi}\catcode`\%=\active\def%{\%}4.985}}%
\end{pgfscope}%
\begin{pgfscope}%
\pgfpathrectangle{\pgfqpoint{2.197902in}{0.579321in}}{\pgfqpoint{0.979200in}{0.326250in}}%
\pgfusepath{clip}%
\pgfsetbuttcap%
\pgfsetroundjoin%
\pgfsetlinewidth{0.401500pt}%
\definecolor{currentstroke}{rgb}{0.690196,0.690196,0.690196}%
\pgfsetstrokecolor{currentstroke}%
\pgfsetstrokeopacity{0.250000}%
\pgfsetdash{{1.480000pt}{0.640000pt}}{0.000000pt}%
\pgfpathmoveto{\pgfqpoint{2.687502in}{0.579321in}}%
\pgfpathlineto{\pgfqpoint{2.687502in}{0.905571in}}%
\pgfusepath{stroke}%
\end{pgfscope}%
\begin{pgfscope}%
\pgfsetbuttcap%
\pgfsetroundjoin%
\definecolor{currentfill}{rgb}{0.000000,0.000000,0.000000}%
\pgfsetfillcolor{currentfill}%
\pgfsetlinewidth{0.803000pt}%
\definecolor{currentstroke}{rgb}{0.000000,0.000000,0.000000}%
\pgfsetstrokecolor{currentstroke}%
\pgfsetdash{}{0pt}%
\pgfsys@defobject{currentmarker}{\pgfqpoint{0.000000in}{0.000000in}}{\pgfqpoint{0.000000in}{0.027778in}}{%
\pgfpathmoveto{\pgfqpoint{0.000000in}{0.000000in}}%
\pgfpathlineto{\pgfqpoint{0.000000in}{0.027778in}}%
\pgfusepath{stroke,fill}%
}%
\begin{pgfscope}%
\pgfsys@transformshift{2.687502in}{0.579321in}%
\pgfsys@useobject{currentmarker}{}%
\end{pgfscope}%
\end{pgfscope}%
\begin{pgfscope}%
\definecolor{textcolor}{rgb}{0.000000,0.000000,0.000000}%
\pgfsetstrokecolor{textcolor}%
\pgfsetfillcolor{textcolor}%
\pgftext[x=2.687502in,y=0.565432in,,top]{\color{textcolor}{\rmfamily\fontsize{5.000000}{6.000000}\selectfont\catcode`\^=\active\def^{\ifmmode\sp\else\^{}\fi}\catcode`\%=\active\def%{\%}4.990}}%
\end{pgfscope}%
\begin{pgfscope}%
\pgfpathrectangle{\pgfqpoint{2.197902in}{0.579321in}}{\pgfqpoint{0.979200in}{0.326250in}}%
\pgfusepath{clip}%
\pgfsetbuttcap%
\pgfsetroundjoin%
\pgfsetlinewidth{0.401500pt}%
\definecolor{currentstroke}{rgb}{0.690196,0.690196,0.690196}%
\pgfsetstrokecolor{currentstroke}%
\pgfsetstrokeopacity{0.250000}%
\pgfsetdash{{1.480000pt}{0.640000pt}}{0.000000pt}%
\pgfpathmoveto{\pgfqpoint{2.932302in}{0.579321in}}%
\pgfpathlineto{\pgfqpoint{2.932302in}{0.905571in}}%
\pgfusepath{stroke}%
\end{pgfscope}%
\begin{pgfscope}%
\pgfsetbuttcap%
\pgfsetroundjoin%
\definecolor{currentfill}{rgb}{0.000000,0.000000,0.000000}%
\pgfsetfillcolor{currentfill}%
\pgfsetlinewidth{0.803000pt}%
\definecolor{currentstroke}{rgb}{0.000000,0.000000,0.000000}%
\pgfsetstrokecolor{currentstroke}%
\pgfsetdash{}{0pt}%
\pgfsys@defobject{currentmarker}{\pgfqpoint{0.000000in}{0.000000in}}{\pgfqpoint{0.000000in}{0.027778in}}{%
\pgfpathmoveto{\pgfqpoint{0.000000in}{0.000000in}}%
\pgfpathlineto{\pgfqpoint{0.000000in}{0.027778in}}%
\pgfusepath{stroke,fill}%
}%
\begin{pgfscope}%
\pgfsys@transformshift{2.932302in}{0.579321in}%
\pgfsys@useobject{currentmarker}{}%
\end{pgfscope}%
\end{pgfscope}%
\begin{pgfscope}%
\definecolor{textcolor}{rgb}{0.000000,0.000000,0.000000}%
\pgfsetstrokecolor{textcolor}%
\pgfsetfillcolor{textcolor}%
\pgftext[x=2.932302in,y=0.565432in,,top]{\color{textcolor}{\rmfamily\fontsize{5.000000}{6.000000}\selectfont\catcode`\^=\active\def^{\ifmmode\sp\else\^{}\fi}\catcode`\%=\active\def%{\%}4.995}}%
\end{pgfscope}%
\begin{pgfscope}%
\pgfpathrectangle{\pgfqpoint{2.197902in}{0.579321in}}{\pgfqpoint{0.979200in}{0.326250in}}%
\pgfusepath{clip}%
\pgfsetbuttcap%
\pgfsetroundjoin%
\pgfsetlinewidth{0.401500pt}%
\definecolor{currentstroke}{rgb}{0.690196,0.690196,0.690196}%
\pgfsetstrokecolor{currentstroke}%
\pgfsetstrokeopacity{0.250000}%
\pgfsetdash{{1.480000pt}{0.640000pt}}{0.000000pt}%
\pgfpathmoveto{\pgfqpoint{3.177102in}{0.579321in}}%
\pgfpathlineto{\pgfqpoint{3.177102in}{0.905571in}}%
\pgfusepath{stroke}%
\end{pgfscope}%
\begin{pgfscope}%
\pgfsetbuttcap%
\pgfsetroundjoin%
\definecolor{currentfill}{rgb}{0.000000,0.000000,0.000000}%
\pgfsetfillcolor{currentfill}%
\pgfsetlinewidth{0.803000pt}%
\definecolor{currentstroke}{rgb}{0.000000,0.000000,0.000000}%
\pgfsetstrokecolor{currentstroke}%
\pgfsetdash{}{0pt}%
\pgfsys@defobject{currentmarker}{\pgfqpoint{0.000000in}{0.000000in}}{\pgfqpoint{0.000000in}{0.027778in}}{%
\pgfpathmoveto{\pgfqpoint{0.000000in}{0.000000in}}%
\pgfpathlineto{\pgfqpoint{0.000000in}{0.027778in}}%
\pgfusepath{stroke,fill}%
}%
\begin{pgfscope}%
\pgfsys@transformshift{3.177102in}{0.579321in}%
\pgfsys@useobject{currentmarker}{}%
\end{pgfscope}%
\end{pgfscope}%
\begin{pgfscope}%
\definecolor{textcolor}{rgb}{0.000000,0.000000,0.000000}%
\pgfsetstrokecolor{textcolor}%
\pgfsetfillcolor{textcolor}%
\pgftext[x=3.177102in,y=0.565432in,,top]{\color{textcolor}{\rmfamily\fontsize{5.000000}{6.000000}\selectfont\catcode`\^=\active\def^{\ifmmode\sp\else\^{}\fi}\catcode`\%=\active\def%{\%}5.000}}%
\end{pgfscope}%
\begin{pgfscope}%
\pgfpathrectangle{\pgfqpoint{2.197902in}{0.579321in}}{\pgfqpoint{0.979200in}{0.326250in}}%
\pgfusepath{clip}%
\pgfsetbuttcap%
\pgfsetroundjoin%
\pgfsetlinewidth{0.401500pt}%
\definecolor{currentstroke}{rgb}{0.690196,0.690196,0.690196}%
\pgfsetstrokecolor{currentstroke}%
\pgfsetstrokeopacity{0.250000}%
\pgfsetdash{{1.480000pt}{0.640000pt}}{0.000000pt}%
\pgfpathmoveto{\pgfqpoint{2.197902in}{0.579321in}}%
\pgfpathlineto{\pgfqpoint{3.177102in}{0.579321in}}%
\pgfusepath{stroke}%
\end{pgfscope}%
\begin{pgfscope}%
\pgfsetbuttcap%
\pgfsetroundjoin%
\definecolor{currentfill}{rgb}{0.000000,0.000000,0.000000}%
\pgfsetfillcolor{currentfill}%
\pgfsetlinewidth{0.803000pt}%
\definecolor{currentstroke}{rgb}{0.000000,0.000000,0.000000}%
\pgfsetstrokecolor{currentstroke}%
\pgfsetdash{}{0pt}%
\pgfsys@defobject{currentmarker}{\pgfqpoint{0.000000in}{0.000000in}}{\pgfqpoint{0.027778in}{0.000000in}}{%
\pgfpathmoveto{\pgfqpoint{0.000000in}{0.000000in}}%
\pgfpathlineto{\pgfqpoint{0.027778in}{0.000000in}}%
\pgfusepath{stroke,fill}%
}%
\begin{pgfscope}%
\pgfsys@transformshift{2.197902in}{0.579321in}%
\pgfsys@useobject{currentmarker}{}%
\end{pgfscope}%
\end{pgfscope}%
\begin{pgfscope}%
\definecolor{textcolor}{rgb}{0.000000,0.000000,0.000000}%
\pgfsetstrokecolor{textcolor}%
\pgfsetfillcolor{textcolor}%
\pgftext[x=1.966996in, y=0.555209in, left, base]{\color{textcolor}{\rmfamily\fontsize{5.000000}{6.000000}\selectfont\catcode`\^=\active\def^{\ifmmode\sp\else\^{}\fi}\catcode`\%=\active\def%{\%}14.86}}%
\end{pgfscope}%
\begin{pgfscope}%
\pgfpathrectangle{\pgfqpoint{2.197902in}{0.579321in}}{\pgfqpoint{0.979200in}{0.326250in}}%
\pgfusepath{clip}%
\pgfsetbuttcap%
\pgfsetroundjoin%
\pgfsetlinewidth{0.401500pt}%
\definecolor{currentstroke}{rgb}{0.690196,0.690196,0.690196}%
\pgfsetstrokecolor{currentstroke}%
\pgfsetstrokeopacity{0.250000}%
\pgfsetdash{{1.480000pt}{0.640000pt}}{0.000000pt}%
\pgfpathmoveto{\pgfqpoint{2.197902in}{0.742446in}}%
\pgfpathlineto{\pgfqpoint{3.177102in}{0.742446in}}%
\pgfusepath{stroke}%
\end{pgfscope}%
\begin{pgfscope}%
\pgfsetbuttcap%
\pgfsetroundjoin%
\definecolor{currentfill}{rgb}{0.000000,0.000000,0.000000}%
\pgfsetfillcolor{currentfill}%
\pgfsetlinewidth{0.803000pt}%
\definecolor{currentstroke}{rgb}{0.000000,0.000000,0.000000}%
\pgfsetstrokecolor{currentstroke}%
\pgfsetdash{}{0pt}%
\pgfsys@defobject{currentmarker}{\pgfqpoint{0.000000in}{0.000000in}}{\pgfqpoint{0.027778in}{0.000000in}}{%
\pgfpathmoveto{\pgfqpoint{0.000000in}{0.000000in}}%
\pgfpathlineto{\pgfqpoint{0.027778in}{0.000000in}}%
\pgfusepath{stroke,fill}%
}%
\begin{pgfscope}%
\pgfsys@transformshift{2.197902in}{0.742446in}%
\pgfsys@useobject{currentmarker}{}%
\end{pgfscope}%
\end{pgfscope}%
\begin{pgfscope}%
\definecolor{textcolor}{rgb}{0.000000,0.000000,0.000000}%
\pgfsetstrokecolor{textcolor}%
\pgfsetfillcolor{textcolor}%
\pgftext[x=1.966996in, y=0.718334in, left, base]{\color{textcolor}{\rmfamily\fontsize{5.000000}{6.000000}\selectfont\catcode`\^=\active\def^{\ifmmode\sp\else\^{}\fi}\catcode`\%=\active\def%{\%}14.91}}%
\end{pgfscope}%
\begin{pgfscope}%
\pgfpathrectangle{\pgfqpoint{2.197902in}{0.579321in}}{\pgfqpoint{0.979200in}{0.326250in}}%
\pgfusepath{clip}%
\pgfsetbuttcap%
\pgfsetroundjoin%
\pgfsetlinewidth{0.401500pt}%
\definecolor{currentstroke}{rgb}{0.690196,0.690196,0.690196}%
\pgfsetstrokecolor{currentstroke}%
\pgfsetstrokeopacity{0.250000}%
\pgfsetdash{{1.480000pt}{0.640000pt}}{0.000000pt}%
\pgfpathmoveto{\pgfqpoint{2.197902in}{0.905571in}}%
\pgfpathlineto{\pgfqpoint{3.177102in}{0.905571in}}%
\pgfusepath{stroke}%
\end{pgfscope}%
\begin{pgfscope}%
\pgfsetbuttcap%
\pgfsetroundjoin%
\definecolor{currentfill}{rgb}{0.000000,0.000000,0.000000}%
\pgfsetfillcolor{currentfill}%
\pgfsetlinewidth{0.803000pt}%
\definecolor{currentstroke}{rgb}{0.000000,0.000000,0.000000}%
\pgfsetstrokecolor{currentstroke}%
\pgfsetdash{}{0pt}%
\pgfsys@defobject{currentmarker}{\pgfqpoint{0.000000in}{0.000000in}}{\pgfqpoint{0.027778in}{0.000000in}}{%
\pgfpathmoveto{\pgfqpoint{0.000000in}{0.000000in}}%
\pgfpathlineto{\pgfqpoint{0.027778in}{0.000000in}}%
\pgfusepath{stroke,fill}%
}%
\begin{pgfscope}%
\pgfsys@transformshift{2.197902in}{0.905571in}%
\pgfsys@useobject{currentmarker}{}%
\end{pgfscope}%
\end{pgfscope}%
\begin{pgfscope}%
\definecolor{textcolor}{rgb}{0.000000,0.000000,0.000000}%
\pgfsetstrokecolor{textcolor}%
\pgfsetfillcolor{textcolor}%
\pgftext[x=1.966996in, y=0.881459in, left, base]{\color{textcolor}{\rmfamily\fontsize{5.000000}{6.000000}\selectfont\catcode`\^=\active\def^{\ifmmode\sp\else\^{}\fi}\catcode`\%=\active\def%{\%}14.96}}%
\end{pgfscope}%
\begin{pgfscope}%
\pgfpathrectangle{\pgfqpoint{2.197902in}{0.579321in}}{\pgfqpoint{0.979200in}{0.326250in}}%
\pgfusepath{clip}%
\pgfsetrectcap%
\pgfsetroundjoin%
\pgfsetlinewidth{1.505625pt}%
\definecolor{currentstroke}{rgb}{0.000000,0.619608,0.450980}%
\pgfsetstrokecolor{currentstroke}%
\pgfsetdash{}{0pt}%
\pgfpathmoveto{\pgfqpoint{2.193930in}{0.808770in}}%
\pgfpathlineto{\pgfqpoint{2.399755in}{0.607075in}}%
\pgfpathlineto{\pgfqpoint{2.401327in}{0.606509in}}%
\pgfpathlineto{\pgfqpoint{2.402972in}{0.609667in}}%
\pgfpathlineto{\pgfqpoint{2.407951in}{0.625215in}}%
\pgfpathlineto{\pgfqpoint{2.410460in}{0.632320in}}%
\pgfpathlineto{\pgfqpoint{2.417357in}{0.651185in}}%
\pgfpathlineto{\pgfqpoint{2.419067in}{0.656029in}}%
\pgfpathlineto{\pgfqpoint{2.423057in}{0.666655in}}%
\pgfpathlineto{\pgfqpoint{2.424771in}{0.671050in}}%
\pgfpathlineto{\pgfqpoint{2.428332in}{0.680588in}}%
\pgfpathlineto{\pgfqpoint{2.439476in}{0.707705in}}%
\pgfpathlineto{\pgfqpoint{2.443118in}{0.716488in}}%
\pgfpathlineto{\pgfqpoint{2.450673in}{0.732932in}}%
\pgfpathlineto{\pgfqpoint{2.455381in}{0.743405in}}%
\pgfpathlineto{\pgfqpoint{2.457921in}{0.748836in}}%
\pgfpathlineto{\pgfqpoint{2.463289in}{0.759583in}}%
\pgfpathlineto{\pgfqpoint{2.468045in}{0.768475in}}%
\pgfpathlineto{\pgfqpoint{2.479066in}{0.788749in}}%
\pgfpathlineto{\pgfqpoint{2.483602in}{0.796147in}}%
\pgfpathlineto{\pgfqpoint{2.488181in}{0.803713in}}%
\pgfpathlineto{\pgfqpoint{2.492749in}{0.810531in}}%
\pgfpathlineto{\pgfqpoint{2.497291in}{0.817301in}}%
\pgfpathlineto{\pgfqpoint{2.501840in}{0.823445in}}%
\pgfpathlineto{\pgfqpoint{2.506389in}{0.829513in}}%
\pgfpathlineto{\pgfqpoint{2.510977in}{0.835058in}}%
\pgfpathlineto{\pgfqpoint{2.515661in}{0.840530in}}%
\pgfpathlineto{\pgfqpoint{2.520703in}{0.846011in}}%
\pgfpathlineto{\pgfqpoint{2.530815in}{0.855579in}}%
\pgfpathlineto{\pgfqpoint{2.555228in}{0.871769in}}%
\pgfpathlineto{\pgfqpoint{2.561037in}{0.874206in}}%
\pgfpathlineto{\pgfqpoint{2.574281in}{0.877717in}}%
\pgfpathlineto{\pgfqpoint{2.589474in}{0.878038in}}%
\pgfpathlineto{\pgfqpoint{2.598406in}{0.876485in}}%
\pgfpathlineto{\pgfqpoint{2.606029in}{0.874160in}}%
\pgfpathlineto{\pgfqpoint{2.613936in}{0.870696in}}%
\pgfpathlineto{\pgfqpoint{2.623728in}{0.865017in}}%
\pgfpathlineto{\pgfqpoint{2.633520in}{0.857785in}}%
\pgfpathlineto{\pgfqpoint{2.648208in}{0.844378in}}%
\pgfpathlineto{\pgfqpoint{2.890691in}{0.606860in}}%
\pgfpathlineto{\pgfqpoint{2.892294in}{0.608662in}}%
\pgfpathlineto{\pgfqpoint{2.895970in}{0.619422in}}%
\pgfpathlineto{\pgfqpoint{2.901272in}{0.635190in}}%
\pgfpathlineto{\pgfqpoint{2.902438in}{0.638263in}}%
\pgfpathlineto{\pgfqpoint{2.904410in}{0.643623in}}%
\pgfpathlineto{\pgfqpoint{2.905530in}{0.647073in}}%
\pgfpathlineto{\pgfqpoint{2.907444in}{0.651798in}}%
\pgfpathlineto{\pgfqpoint{2.911485in}{0.663217in}}%
\pgfpathlineto{\pgfqpoint{2.913394in}{0.667582in}}%
\pgfpathlineto{\pgfqpoint{2.933045in}{0.716495in}}%
\pgfpathlineto{\pgfqpoint{2.934638in}{0.719629in}}%
\pgfpathlineto{\pgfqpoint{2.936740in}{0.724184in}}%
\pgfpathlineto{\pgfqpoint{2.943800in}{0.739953in}}%
\pgfpathlineto{\pgfqpoint{2.947476in}{0.747780in}}%
\pgfpathlineto{\pgfqpoint{2.949848in}{0.752496in}}%
\pgfpathlineto{\pgfqpoint{2.954272in}{0.760717in}}%
\pgfpathlineto{\pgfqpoint{2.964279in}{0.779594in}}%
\pgfpathlineto{\pgfqpoint{2.968223in}{0.786500in}}%
\pgfpathlineto{\pgfqpoint{2.970222in}{0.789860in}}%
\pgfpathlineto{\pgfqpoint{2.974273in}{0.796403in}}%
\pgfpathlineto{\pgfqpoint{2.976188in}{0.799575in}}%
\pgfpathlineto{\pgfqpoint{2.982581in}{0.809320in}}%
\pgfpathlineto{\pgfqpoint{2.987005in}{0.815275in}}%
\pgfpathlineto{\pgfqpoint{2.991434in}{0.821630in}}%
\pgfpathlineto{\pgfqpoint{2.995862in}{0.827023in}}%
\pgfpathlineto{\pgfqpoint{3.000291in}{0.832658in}}%
\pgfpathlineto{\pgfqpoint{3.004724in}{0.837480in}}%
\pgfpathlineto{\pgfqpoint{3.009204in}{0.842446in}}%
\pgfpathlineto{\pgfqpoint{3.013316in}{0.846331in}}%
\pgfpathlineto{\pgfqpoint{3.018262in}{0.851025in}}%
\pgfpathlineto{\pgfqpoint{3.022822in}{0.854878in}}%
\pgfpathlineto{\pgfqpoint{3.027217in}{0.858223in}}%
\pgfpathlineto{\pgfqpoint{3.031909in}{0.861553in}}%
\pgfpathlineto{\pgfqpoint{3.038390in}{0.865479in}}%
\pgfpathlineto{\pgfqpoint{3.059097in}{0.873295in}}%
\pgfpathlineto{\pgfqpoint{3.072220in}{0.874521in}}%
\pgfpathlineto{\pgfqpoint{3.086091in}{0.872889in}}%
\pgfpathlineto{\pgfqpoint{3.095364in}{0.869993in}}%
\pgfpathlineto{\pgfqpoint{3.104869in}{0.865551in}}%
\pgfpathlineto{\pgfqpoint{3.113203in}{0.860438in}}%
\pgfpathlineto{\pgfqpoint{3.125266in}{0.851019in}}%
\pgfpathlineto{\pgfqpoint{3.133426in}{0.843376in}}%
\pgfpathlineto{\pgfqpoint{3.177102in}{0.800474in}}%
\pgfpathlineto{\pgfqpoint{3.177102in}{0.800474in}}%
\pgfusepath{stroke}%
\end{pgfscope}%
\begin{pgfscope}%
\pgfsetrectcap%
\pgfsetmiterjoin%
\pgfsetlinewidth{0.803000pt}%
\definecolor{currentstroke}{rgb}{0.000000,0.000000,0.000000}%
\pgfsetstrokecolor{currentstroke}%
\pgfsetdash{}{0pt}%
\pgfpathmoveto{\pgfqpoint{2.197902in}{0.579321in}}%
\pgfpathlineto{\pgfqpoint{2.197902in}{0.905571in}}%
\pgfusepath{stroke}%
\end{pgfscope}%
\begin{pgfscope}%
\pgfsetrectcap%
\pgfsetmiterjoin%
\pgfsetlinewidth{0.803000pt}%
\definecolor{currentstroke}{rgb}{0.000000,0.000000,0.000000}%
\pgfsetstrokecolor{currentstroke}%
\pgfsetdash{}{0pt}%
\pgfpathmoveto{\pgfqpoint{3.177102in}{0.579321in}}%
\pgfpathlineto{\pgfqpoint{3.177102in}{0.905571in}}%
\pgfusepath{stroke}%
\end{pgfscope}%
\begin{pgfscope}%
\pgfsetrectcap%
\pgfsetmiterjoin%
\pgfsetlinewidth{0.803000pt}%
\definecolor{currentstroke}{rgb}{0.000000,0.000000,0.000000}%
\pgfsetstrokecolor{currentstroke}%
\pgfsetdash{}{0pt}%
\pgfpathmoveto{\pgfqpoint{2.197902in}{0.579321in}}%
\pgfpathlineto{\pgfqpoint{3.177102in}{0.579321in}}%
\pgfusepath{stroke}%
\end{pgfscope}%
\begin{pgfscope}%
\pgfsetrectcap%
\pgfsetmiterjoin%
\pgfsetlinewidth{0.803000pt}%
\definecolor{currentstroke}{rgb}{0.000000,0.000000,0.000000}%
\pgfsetstrokecolor{currentstroke}%
\pgfsetdash{}{0pt}%
\pgfpathmoveto{\pgfqpoint{2.197902in}{0.905571in}}%
\pgfpathlineto{\pgfqpoint{3.177102in}{0.905571in}}%
\pgfusepath{stroke}%
\end{pgfscope}%
\end{pgfpicture}%
\makeatother%
\endgroup%

	\caption{(a) corrente no dreno do transistor, (b) tensão dreno--source do transistor e (c) tensão de saída do circuito Flyback para $K~L1~L2~0.99$ .}
	\label{fig:ex4_2}
\end{figure}

\medskip
\textbf{c)} Projete o circuito \textit{clamper} RCD, conforme apresentado em aula, para o circuito e apresente as novas formas de onda.

\medskip
\textbf{Resolução}: a tensão $V_{DS}$ máxima, assim como a corrente de pico máxima, foram obtidos por inspeção dos gráficos gerados em simulação, com o auxílio de ferramentas de ampliação. Os valores obtidos para o circuito \textit{clamper} RCD foram~\cite{aula6_isolados}: $R_{clamp} = 2{,}6~\text{k}\Omega$, $C_{clamp} = 19{,}2$ nF, que resulta em uma perda $P_{loss,clamper} = 0{,}8352$ W.

A Figura~\ref{fig:ex4_3} apresenta as formas de onda solicitadas no enunciado. Pode-se observar a diminuição da amplitude da oscilação na tensão $V_{DS}$.
\begin{figure}[htb!]
	\centering
	%% Creator: Matplotlib, PGF backend
%%
%% To include the figure in your LaTeX document, write
%%   \input{<filename>.pgf}
%%
%% Make sure the required packages are loaded in your preamble
%%   \usepackage{pgf}
%%
%% Also ensure that all the required font packages are loaded; for instance,
%% the lmodern package is sometimes necessary when using math font.
%%   \usepackage{lmodern}
%%
%% Figures using additional raster images can only be included by \input if
%% they are in the same directory as the main LaTeX file. For loading figures
%% from other directories you can use the `import` package
%%   \usepackage{import}
%%
%% and then include the figures with
%%   \import{<path to file>}{<filename>.pgf}
%%
%% Matplotlib used the following preamble
%%   \def\mathdefault#1{#1}
%%   \everymath=\expandafter{\the\everymath\displaystyle}
%%   \IfFileExists{scrextend.sty}{
%%     \usepackage[fontsize=8.000000pt]{scrextend}
%%   }{
%%     \renewcommand{\normalsize}{\fontsize{8.000000}{9.600000}\selectfont}
%%     \normalsize
%%   }
%%   
%%           \usepackage{amsmath}
%%           \usepackage{amssymb}
%%           \usepackage{mathtools}
%%           \usepackage{dsfont}
%%           \providecommand{\mathdefault}[1]{#1}
%%       
%%   \makeatletter\@ifpackageloaded{underscore}{}{\usepackage[strings]{underscore}}\makeatother
%%
\begingroup%
\makeatletter%
\begin{pgfpicture}%
\pgfpathrectangle{\pgfpointorigin}{\pgfqpoint{3.393768in}{3.646235in}}%
\pgfusepath{use as bounding box, clip}%
\begin{pgfscope}%
\pgfsetbuttcap%
\pgfsetmiterjoin%
\definecolor{currentfill}{rgb}{1.000000,1.000000,1.000000}%
\pgfsetfillcolor{currentfill}%
\pgfsetlinewidth{0.000000pt}%
\definecolor{currentstroke}{rgb}{1.000000,1.000000,1.000000}%
\pgfsetstrokecolor{currentstroke}%
\pgfsetdash{}{0pt}%
\pgfpathmoveto{\pgfqpoint{0.000000in}{0.000000in}}%
\pgfpathlineto{\pgfqpoint{3.393768in}{0.000000in}}%
\pgfpathlineto{\pgfqpoint{3.393768in}{3.646235in}}%
\pgfpathlineto{\pgfqpoint{0.000000in}{3.646235in}}%
\pgfpathlineto{\pgfqpoint{0.000000in}{0.000000in}}%
\pgfpathclose%
\pgfusepath{fill}%
\end{pgfscope}%
\begin{pgfscope}%
\pgfsetbuttcap%
\pgfsetmiterjoin%
\definecolor{currentfill}{rgb}{1.000000,1.000000,1.000000}%
\pgfsetfillcolor{currentfill}%
\pgfsetlinewidth{0.000000pt}%
\definecolor{currentstroke}{rgb}{0.000000,0.000000,0.000000}%
\pgfsetstrokecolor{currentstroke}%
\pgfsetstrokeopacity{0.000000}%
\pgfsetdash{}{0pt}%
\pgfpathmoveto{\pgfqpoint{0.573768in}{2.601404in}}%
\pgfpathlineto{\pgfqpoint{3.293768in}{2.601404in}}%
\pgfpathlineto{\pgfqpoint{3.293768in}{3.507654in}}%
\pgfpathlineto{\pgfqpoint{0.573768in}{3.507654in}}%
\pgfpathlineto{\pgfqpoint{0.573768in}{2.601404in}}%
\pgfpathclose%
\pgfusepath{fill}%
\end{pgfscope}%
\begin{pgfscope}%
\pgfpathrectangle{\pgfqpoint{0.573768in}{2.601404in}}{\pgfqpoint{2.720000in}{0.906250in}}%
\pgfusepath{clip}%
\pgfsetbuttcap%
\pgfsetroundjoin%
\pgfsetlinewidth{0.501875pt}%
\definecolor{currentstroke}{rgb}{0.690196,0.690196,0.690196}%
\pgfsetstrokecolor{currentstroke}%
\pgfsetstrokeopacity{0.250000}%
\pgfsetdash{{1.850000pt}{0.800000pt}}{0.000000pt}%
\pgfpathmoveto{\pgfqpoint{0.697405in}{2.601404in}}%
\pgfpathlineto{\pgfqpoint{0.697405in}{3.507654in}}%
\pgfusepath{stroke}%
\end{pgfscope}%
\begin{pgfscope}%
\pgfsetbuttcap%
\pgfsetroundjoin%
\definecolor{currentfill}{rgb}{0.000000,0.000000,0.000000}%
\pgfsetfillcolor{currentfill}%
\pgfsetlinewidth{0.803000pt}%
\definecolor{currentstroke}{rgb}{0.000000,0.000000,0.000000}%
\pgfsetstrokecolor{currentstroke}%
\pgfsetdash{}{0pt}%
\pgfsys@defobject{currentmarker}{\pgfqpoint{0.000000in}{-0.048611in}}{\pgfqpoint{0.000000in}{0.000000in}}{%
\pgfpathmoveto{\pgfqpoint{0.000000in}{0.000000in}}%
\pgfpathlineto{\pgfqpoint{0.000000in}{-0.048611in}}%
\pgfusepath{stroke,fill}%
}%
\begin{pgfscope}%
\pgfsys@transformshift{0.697405in}{2.601404in}%
\pgfsys@useobject{currentmarker}{}%
\end{pgfscope}%
\end{pgfscope}%
\begin{pgfscope}%
\pgfpathrectangle{\pgfqpoint{0.573768in}{2.601404in}}{\pgfqpoint{2.720000in}{0.906250in}}%
\pgfusepath{clip}%
\pgfsetbuttcap%
\pgfsetroundjoin%
\pgfsetlinewidth{0.501875pt}%
\definecolor{currentstroke}{rgb}{0.690196,0.690196,0.690196}%
\pgfsetstrokecolor{currentstroke}%
\pgfsetstrokeopacity{0.250000}%
\pgfsetdash{{1.850000pt}{0.800000pt}}{0.000000pt}%
\pgfpathmoveto{\pgfqpoint{1.191950in}{2.601404in}}%
\pgfpathlineto{\pgfqpoint{1.191950in}{3.507654in}}%
\pgfusepath{stroke}%
\end{pgfscope}%
\begin{pgfscope}%
\pgfsetbuttcap%
\pgfsetroundjoin%
\definecolor{currentfill}{rgb}{0.000000,0.000000,0.000000}%
\pgfsetfillcolor{currentfill}%
\pgfsetlinewidth{0.803000pt}%
\definecolor{currentstroke}{rgb}{0.000000,0.000000,0.000000}%
\pgfsetstrokecolor{currentstroke}%
\pgfsetdash{}{0pt}%
\pgfsys@defobject{currentmarker}{\pgfqpoint{0.000000in}{-0.048611in}}{\pgfqpoint{0.000000in}{0.000000in}}{%
\pgfpathmoveto{\pgfqpoint{0.000000in}{0.000000in}}%
\pgfpathlineto{\pgfqpoint{0.000000in}{-0.048611in}}%
\pgfusepath{stroke,fill}%
}%
\begin{pgfscope}%
\pgfsys@transformshift{1.191950in}{2.601404in}%
\pgfsys@useobject{currentmarker}{}%
\end{pgfscope}%
\end{pgfscope}%
\begin{pgfscope}%
\pgfpathrectangle{\pgfqpoint{0.573768in}{2.601404in}}{\pgfqpoint{2.720000in}{0.906250in}}%
\pgfusepath{clip}%
\pgfsetbuttcap%
\pgfsetroundjoin%
\pgfsetlinewidth{0.501875pt}%
\definecolor{currentstroke}{rgb}{0.690196,0.690196,0.690196}%
\pgfsetstrokecolor{currentstroke}%
\pgfsetstrokeopacity{0.250000}%
\pgfsetdash{{1.850000pt}{0.800000pt}}{0.000000pt}%
\pgfpathmoveto{\pgfqpoint{1.686496in}{2.601404in}}%
\pgfpathlineto{\pgfqpoint{1.686496in}{3.507654in}}%
\pgfusepath{stroke}%
\end{pgfscope}%
\begin{pgfscope}%
\pgfsetbuttcap%
\pgfsetroundjoin%
\definecolor{currentfill}{rgb}{0.000000,0.000000,0.000000}%
\pgfsetfillcolor{currentfill}%
\pgfsetlinewidth{0.803000pt}%
\definecolor{currentstroke}{rgb}{0.000000,0.000000,0.000000}%
\pgfsetstrokecolor{currentstroke}%
\pgfsetdash{}{0pt}%
\pgfsys@defobject{currentmarker}{\pgfqpoint{0.000000in}{-0.048611in}}{\pgfqpoint{0.000000in}{0.000000in}}{%
\pgfpathmoveto{\pgfqpoint{0.000000in}{0.000000in}}%
\pgfpathlineto{\pgfqpoint{0.000000in}{-0.048611in}}%
\pgfusepath{stroke,fill}%
}%
\begin{pgfscope}%
\pgfsys@transformshift{1.686496in}{2.601404in}%
\pgfsys@useobject{currentmarker}{}%
\end{pgfscope}%
\end{pgfscope}%
\begin{pgfscope}%
\pgfpathrectangle{\pgfqpoint{0.573768in}{2.601404in}}{\pgfqpoint{2.720000in}{0.906250in}}%
\pgfusepath{clip}%
\pgfsetbuttcap%
\pgfsetroundjoin%
\pgfsetlinewidth{0.501875pt}%
\definecolor{currentstroke}{rgb}{0.690196,0.690196,0.690196}%
\pgfsetstrokecolor{currentstroke}%
\pgfsetstrokeopacity{0.250000}%
\pgfsetdash{{1.850000pt}{0.800000pt}}{0.000000pt}%
\pgfpathmoveto{\pgfqpoint{2.181041in}{2.601404in}}%
\pgfpathlineto{\pgfqpoint{2.181041in}{3.507654in}}%
\pgfusepath{stroke}%
\end{pgfscope}%
\begin{pgfscope}%
\pgfsetbuttcap%
\pgfsetroundjoin%
\definecolor{currentfill}{rgb}{0.000000,0.000000,0.000000}%
\pgfsetfillcolor{currentfill}%
\pgfsetlinewidth{0.803000pt}%
\definecolor{currentstroke}{rgb}{0.000000,0.000000,0.000000}%
\pgfsetstrokecolor{currentstroke}%
\pgfsetdash{}{0pt}%
\pgfsys@defobject{currentmarker}{\pgfqpoint{0.000000in}{-0.048611in}}{\pgfqpoint{0.000000in}{0.000000in}}{%
\pgfpathmoveto{\pgfqpoint{0.000000in}{0.000000in}}%
\pgfpathlineto{\pgfqpoint{0.000000in}{-0.048611in}}%
\pgfusepath{stroke,fill}%
}%
\begin{pgfscope}%
\pgfsys@transformshift{2.181041in}{2.601404in}%
\pgfsys@useobject{currentmarker}{}%
\end{pgfscope}%
\end{pgfscope}%
\begin{pgfscope}%
\pgfpathrectangle{\pgfqpoint{0.573768in}{2.601404in}}{\pgfqpoint{2.720000in}{0.906250in}}%
\pgfusepath{clip}%
\pgfsetbuttcap%
\pgfsetroundjoin%
\pgfsetlinewidth{0.501875pt}%
\definecolor{currentstroke}{rgb}{0.690196,0.690196,0.690196}%
\pgfsetstrokecolor{currentstroke}%
\pgfsetstrokeopacity{0.250000}%
\pgfsetdash{{1.850000pt}{0.800000pt}}{0.000000pt}%
\pgfpathmoveto{\pgfqpoint{2.675587in}{2.601404in}}%
\pgfpathlineto{\pgfqpoint{2.675587in}{3.507654in}}%
\pgfusepath{stroke}%
\end{pgfscope}%
\begin{pgfscope}%
\pgfsetbuttcap%
\pgfsetroundjoin%
\definecolor{currentfill}{rgb}{0.000000,0.000000,0.000000}%
\pgfsetfillcolor{currentfill}%
\pgfsetlinewidth{0.803000pt}%
\definecolor{currentstroke}{rgb}{0.000000,0.000000,0.000000}%
\pgfsetstrokecolor{currentstroke}%
\pgfsetdash{}{0pt}%
\pgfsys@defobject{currentmarker}{\pgfqpoint{0.000000in}{-0.048611in}}{\pgfqpoint{0.000000in}{0.000000in}}{%
\pgfpathmoveto{\pgfqpoint{0.000000in}{0.000000in}}%
\pgfpathlineto{\pgfqpoint{0.000000in}{-0.048611in}}%
\pgfusepath{stroke,fill}%
}%
\begin{pgfscope}%
\pgfsys@transformshift{2.675587in}{2.601404in}%
\pgfsys@useobject{currentmarker}{}%
\end{pgfscope}%
\end{pgfscope}%
\begin{pgfscope}%
\pgfpathrectangle{\pgfqpoint{0.573768in}{2.601404in}}{\pgfqpoint{2.720000in}{0.906250in}}%
\pgfusepath{clip}%
\pgfsetbuttcap%
\pgfsetroundjoin%
\pgfsetlinewidth{0.501875pt}%
\definecolor{currentstroke}{rgb}{0.690196,0.690196,0.690196}%
\pgfsetstrokecolor{currentstroke}%
\pgfsetstrokeopacity{0.250000}%
\pgfsetdash{{1.850000pt}{0.800000pt}}{0.000000pt}%
\pgfpathmoveto{\pgfqpoint{3.170132in}{2.601404in}}%
\pgfpathlineto{\pgfqpoint{3.170132in}{3.507654in}}%
\pgfusepath{stroke}%
\end{pgfscope}%
\begin{pgfscope}%
\pgfsetbuttcap%
\pgfsetroundjoin%
\definecolor{currentfill}{rgb}{0.000000,0.000000,0.000000}%
\pgfsetfillcolor{currentfill}%
\pgfsetlinewidth{0.803000pt}%
\definecolor{currentstroke}{rgb}{0.000000,0.000000,0.000000}%
\pgfsetstrokecolor{currentstroke}%
\pgfsetdash{}{0pt}%
\pgfsys@defobject{currentmarker}{\pgfqpoint{0.000000in}{-0.048611in}}{\pgfqpoint{0.000000in}{0.000000in}}{%
\pgfpathmoveto{\pgfqpoint{0.000000in}{0.000000in}}%
\pgfpathlineto{\pgfqpoint{0.000000in}{-0.048611in}}%
\pgfusepath{stroke,fill}%
}%
\begin{pgfscope}%
\pgfsys@transformshift{3.170132in}{2.601404in}%
\pgfsys@useobject{currentmarker}{}%
\end{pgfscope}%
\end{pgfscope}%
\begin{pgfscope}%
\pgfpathrectangle{\pgfqpoint{0.573768in}{2.601404in}}{\pgfqpoint{2.720000in}{0.906250in}}%
\pgfusepath{clip}%
\pgfsetbuttcap%
\pgfsetroundjoin%
\pgfsetlinewidth{0.501875pt}%
\definecolor{currentstroke}{rgb}{0.690196,0.690196,0.690196}%
\pgfsetstrokecolor{currentstroke}%
\pgfsetstrokeopacity{0.250000}%
\pgfsetdash{{1.850000pt}{0.800000pt}}{0.000000pt}%
\pgfpathmoveto{\pgfqpoint{0.573768in}{2.644559in}}%
\pgfpathlineto{\pgfqpoint{3.293768in}{2.644559in}}%
\pgfusepath{stroke}%
\end{pgfscope}%
\begin{pgfscope}%
\pgfsetbuttcap%
\pgfsetroundjoin%
\definecolor{currentfill}{rgb}{0.000000,0.000000,0.000000}%
\pgfsetfillcolor{currentfill}%
\pgfsetlinewidth{0.803000pt}%
\definecolor{currentstroke}{rgb}{0.000000,0.000000,0.000000}%
\pgfsetstrokecolor{currentstroke}%
\pgfsetdash{}{0pt}%
\pgfsys@defobject{currentmarker}{\pgfqpoint{-0.048611in}{0.000000in}}{\pgfqpoint{-0.000000in}{0.000000in}}{%
\pgfpathmoveto{\pgfqpoint{-0.000000in}{0.000000in}}%
\pgfpathlineto{\pgfqpoint{-0.048611in}{0.000000in}}%
\pgfusepath{stroke,fill}%
}%
\begin{pgfscope}%
\pgfsys@transformshift{0.573768in}{2.644559in}%
\pgfsys@useobject{currentmarker}{}%
\end{pgfscope}%
\end{pgfscope}%
\begin{pgfscope}%
\definecolor{textcolor}{rgb}{0.000000,0.000000,0.000000}%
\pgfsetstrokecolor{textcolor}%
\pgfsetfillcolor{textcolor}%
\pgftext[x=0.417518in, y=2.605979in, left, base]{\color{textcolor}{\rmfamily\fontsize{8.000000}{9.600000}\selectfont\catcode`\^=\active\def^{\ifmmode\sp\else\^{}\fi}\catcode`\%=\active\def%{\%}$\mathdefault{0}$}}%
\end{pgfscope}%
\begin{pgfscope}%
\pgfpathrectangle{\pgfqpoint{0.573768in}{2.601404in}}{\pgfqpoint{2.720000in}{0.906250in}}%
\pgfusepath{clip}%
\pgfsetbuttcap%
\pgfsetroundjoin%
\pgfsetlinewidth{0.501875pt}%
\definecolor{currentstroke}{rgb}{0.690196,0.690196,0.690196}%
\pgfsetstrokecolor{currentstroke}%
\pgfsetstrokeopacity{0.250000}%
\pgfsetdash{{1.850000pt}{0.800000pt}}{0.000000pt}%
\pgfpathmoveto{\pgfqpoint{0.573768in}{3.076107in}}%
\pgfpathlineto{\pgfqpoint{3.293768in}{3.076107in}}%
\pgfusepath{stroke}%
\end{pgfscope}%
\begin{pgfscope}%
\pgfsetbuttcap%
\pgfsetroundjoin%
\definecolor{currentfill}{rgb}{0.000000,0.000000,0.000000}%
\pgfsetfillcolor{currentfill}%
\pgfsetlinewidth{0.803000pt}%
\definecolor{currentstroke}{rgb}{0.000000,0.000000,0.000000}%
\pgfsetstrokecolor{currentstroke}%
\pgfsetdash{}{0pt}%
\pgfsys@defobject{currentmarker}{\pgfqpoint{-0.048611in}{0.000000in}}{\pgfqpoint{-0.000000in}{0.000000in}}{%
\pgfpathmoveto{\pgfqpoint{-0.000000in}{0.000000in}}%
\pgfpathlineto{\pgfqpoint{-0.048611in}{0.000000in}}%
\pgfusepath{stroke,fill}%
}%
\begin{pgfscope}%
\pgfsys@transformshift{0.573768in}{3.076107in}%
\pgfsys@useobject{currentmarker}{}%
\end{pgfscope}%
\end{pgfscope}%
\begin{pgfscope}%
\definecolor{textcolor}{rgb}{0.000000,0.000000,0.000000}%
\pgfsetstrokecolor{textcolor}%
\pgfsetfillcolor{textcolor}%
\pgftext[x=0.417518in, y=3.037527in, left, base]{\color{textcolor}{\rmfamily\fontsize{8.000000}{9.600000}\selectfont\catcode`\^=\active\def^{\ifmmode\sp\else\^{}\fi}\catcode`\%=\active\def%{\%}$\mathdefault{5}$}}%
\end{pgfscope}%
\begin{pgfscope}%
\pgfpathrectangle{\pgfqpoint{0.573768in}{2.601404in}}{\pgfqpoint{2.720000in}{0.906250in}}%
\pgfusepath{clip}%
\pgfsetbuttcap%
\pgfsetroundjoin%
\pgfsetlinewidth{0.501875pt}%
\definecolor{currentstroke}{rgb}{0.690196,0.690196,0.690196}%
\pgfsetstrokecolor{currentstroke}%
\pgfsetstrokeopacity{0.250000}%
\pgfsetdash{{1.850000pt}{0.800000pt}}{0.000000pt}%
\pgfpathmoveto{\pgfqpoint{0.573768in}{3.507654in}}%
\pgfpathlineto{\pgfqpoint{3.293768in}{3.507654in}}%
\pgfusepath{stroke}%
\end{pgfscope}%
\begin{pgfscope}%
\pgfsetbuttcap%
\pgfsetroundjoin%
\definecolor{currentfill}{rgb}{0.000000,0.000000,0.000000}%
\pgfsetfillcolor{currentfill}%
\pgfsetlinewidth{0.803000pt}%
\definecolor{currentstroke}{rgb}{0.000000,0.000000,0.000000}%
\pgfsetstrokecolor{currentstroke}%
\pgfsetdash{}{0pt}%
\pgfsys@defobject{currentmarker}{\pgfqpoint{-0.048611in}{0.000000in}}{\pgfqpoint{-0.000000in}{0.000000in}}{%
\pgfpathmoveto{\pgfqpoint{-0.000000in}{0.000000in}}%
\pgfpathlineto{\pgfqpoint{-0.048611in}{0.000000in}}%
\pgfusepath{stroke,fill}%
}%
\begin{pgfscope}%
\pgfsys@transformshift{0.573768in}{3.507654in}%
\pgfsys@useobject{currentmarker}{}%
\end{pgfscope}%
\end{pgfscope}%
\begin{pgfscope}%
\definecolor{textcolor}{rgb}{0.000000,0.000000,0.000000}%
\pgfsetstrokecolor{textcolor}%
\pgfsetfillcolor{textcolor}%
\pgftext[x=0.358489in, y=3.469074in, left, base]{\color{textcolor}{\rmfamily\fontsize{8.000000}{9.600000}\selectfont\catcode`\^=\active\def^{\ifmmode\sp\else\^{}\fi}\catcode`\%=\active\def%{\%}$\mathdefault{10}$}}%
\end{pgfscope}%
\begin{pgfscope}%
\definecolor{textcolor}{rgb}{0.000000,0.000000,0.000000}%
\pgfsetstrokecolor{textcolor}%
\pgfsetfillcolor{textcolor}%
\pgftext[x=0.302933in,y=3.054529in,,bottom,rotate=90.000000]{\color{textcolor}{\rmfamily\fontsize{8.000000}{9.600000}\selectfont\catcode`\^=\active\def^{\ifmmode\sp\else\^{}\fi}\catcode`\%=\active\def%{\%}$i_D$ (A)}}%
\end{pgfscope}%
\begin{pgfscope}%
\pgfpathrectangle{\pgfqpoint{0.573768in}{2.601404in}}{\pgfqpoint{2.720000in}{0.906250in}}%
\pgfusepath{clip}%
\pgfsetrectcap%
\pgfsetroundjoin%
\pgfsetlinewidth{1.505625pt}%
\definecolor{currentstroke}{rgb}{0.000000,0.447059,0.698039}%
\pgfsetstrokecolor{currentstroke}%
\pgfsetdash{}{0pt}%
\pgfpathmoveto{\pgfqpoint{0.697405in}{2.644559in}}%
\pgfpathlineto{\pgfqpoint{0.697409in}{2.644559in}}%
\pgfpathlineto{\pgfqpoint{0.697416in}{2.743970in}}%
\pgfpathlineto{\pgfqpoint{0.698332in}{2.702815in}}%
\pgfpathlineto{\pgfqpoint{0.699384in}{2.768855in}}%
\pgfpathlineto{\pgfqpoint{0.699411in}{2.755130in}}%
\pgfpathlineto{\pgfqpoint{0.700178in}{2.644559in}}%
\pgfpathlineto{\pgfqpoint{0.702357in}{2.644568in}}%
\pgfpathlineto{\pgfqpoint{0.702368in}{2.853116in}}%
\pgfpathlineto{\pgfqpoint{0.703119in}{2.792204in}}%
\pgfpathlineto{\pgfqpoint{0.703959in}{2.887751in}}%
\pgfpathlineto{\pgfqpoint{0.704329in}{2.910890in}}%
\pgfpathlineto{\pgfqpoint{0.705033in}{2.644559in}}%
\pgfpathlineto{\pgfqpoint{0.707302in}{2.644561in}}%
\pgfpathlineto{\pgfqpoint{0.707319in}{2.954175in}}%
\pgfpathlineto{\pgfqpoint{0.708080in}{2.904377in}}%
\pgfpathlineto{\pgfqpoint{0.708921in}{3.008063in}}%
\pgfpathlineto{\pgfqpoint{0.709274in}{3.030160in}}%
\pgfpathlineto{\pgfqpoint{0.710069in}{2.644559in}}%
\pgfpathlineto{\pgfqpoint{0.712248in}{2.644566in}}%
\pgfpathlineto{\pgfqpoint{0.712953in}{3.061959in}}%
\pgfpathlineto{\pgfqpoint{0.713002in}{3.005227in}}%
\pgfpathlineto{\pgfqpoint{0.713843in}{3.117697in}}%
\pgfpathlineto{\pgfqpoint{0.714220in}{3.141274in}}%
\pgfpathlineto{\pgfqpoint{0.715028in}{2.644559in}}%
\pgfpathlineto{\pgfqpoint{0.717191in}{2.644559in}}%
\pgfpathlineto{\pgfqpoint{0.717205in}{2.894072in}}%
\pgfpathlineto{\pgfqpoint{0.717935in}{3.146751in}}%
\pgfpathlineto{\pgfqpoint{0.717985in}{3.115879in}}%
\pgfpathlineto{\pgfqpoint{0.718628in}{3.190128in}}%
\pgfpathlineto{\pgfqpoint{0.718677in}{3.159428in}}%
\pgfpathlineto{\pgfqpoint{0.719119in}{3.220915in}}%
\pgfpathlineto{\pgfqpoint{0.719218in}{3.148642in}}%
\pgfpathlineto{\pgfqpoint{0.720019in}{2.644559in}}%
\pgfpathlineto{\pgfqpoint{0.722139in}{2.644565in}}%
\pgfpathlineto{\pgfqpoint{0.722160in}{3.225175in}}%
\pgfpathlineto{\pgfqpoint{0.722905in}{3.197416in}}%
\pgfpathlineto{\pgfqpoint{0.723745in}{3.280757in}}%
\pgfpathlineto{\pgfqpoint{0.724111in}{3.303638in}}%
\pgfpathlineto{\pgfqpoint{0.725665in}{2.644559in}}%
\pgfpathlineto{\pgfqpoint{0.727084in}{2.644579in}}%
\pgfpathlineto{\pgfqpoint{0.727108in}{3.293614in}}%
\pgfpathlineto{\pgfqpoint{0.727870in}{3.279853in}}%
\pgfpathlineto{\pgfqpoint{0.729056in}{3.354250in}}%
\pgfpathlineto{\pgfqpoint{0.729083in}{3.335688in}}%
\pgfpathlineto{\pgfqpoint{0.729993in}{2.644559in}}%
\pgfpathlineto{\pgfqpoint{0.732027in}{2.644559in}}%
\pgfpathlineto{\pgfqpoint{0.732053in}{3.345080in}}%
\pgfpathlineto{\pgfqpoint{0.732881in}{3.337423in}}%
\pgfpathlineto{\pgfqpoint{0.732930in}{3.329362in}}%
\pgfpathlineto{\pgfqpoint{0.733375in}{3.368448in}}%
\pgfpathlineto{\pgfqpoint{0.733425in}{3.360408in}}%
\pgfpathlineto{\pgfqpoint{0.733960in}{3.405162in}}%
\pgfpathlineto{\pgfqpoint{0.734037in}{3.341296in}}%
\pgfpathlineto{\pgfqpoint{0.734047in}{3.381344in}}%
\pgfpathlineto{\pgfqpoint{0.734933in}{2.644559in}}%
\pgfpathlineto{\pgfqpoint{0.736973in}{2.644559in}}%
\pgfpathlineto{\pgfqpoint{0.737777in}{3.374339in}}%
\pgfpathlineto{\pgfqpoint{0.737826in}{3.345371in}}%
\pgfpathlineto{\pgfqpoint{0.738469in}{3.417706in}}%
\pgfpathlineto{\pgfqpoint{0.738518in}{3.388898in}}%
\pgfpathlineto{\pgfqpoint{0.738947in}{3.447654in}}%
\pgfpathlineto{\pgfqpoint{0.739019in}{3.217556in}}%
\pgfpathlineto{\pgfqpoint{0.739804in}{2.644559in}}%
\pgfpathlineto{\pgfqpoint{0.741920in}{2.644570in}}%
\pgfpathlineto{\pgfqpoint{0.742715in}{3.388178in}}%
\pgfpathlineto{\pgfqpoint{0.742765in}{3.351440in}}%
\pgfpathlineto{\pgfqpoint{0.743408in}{3.431532in}}%
\pgfpathlineto{\pgfqpoint{0.743457in}{3.394979in}}%
\pgfpathlineto{\pgfqpoint{0.743893in}{3.461893in}}%
\pgfpathlineto{\pgfqpoint{0.743950in}{3.369822in}}%
\pgfpathlineto{\pgfqpoint{0.744760in}{2.644559in}}%
\pgfpathlineto{\pgfqpoint{0.746866in}{2.644674in}}%
\pgfpathlineto{\pgfqpoint{0.746893in}{3.375133in}}%
\pgfpathlineto{\pgfqpoint{0.747613in}{3.363644in}}%
\pgfpathlineto{\pgfqpoint{0.747663in}{3.352368in}}%
\pgfpathlineto{\pgfqpoint{0.748207in}{3.400866in}}%
\pgfpathlineto{\pgfqpoint{0.748256in}{3.389627in}}%
\pgfpathlineto{\pgfqpoint{0.748794in}{3.437727in}}%
\pgfpathlineto{\pgfqpoint{0.748894in}{3.345202in}}%
\pgfpathlineto{\pgfqpoint{0.748907in}{3.408882in}}%
\pgfpathlineto{\pgfqpoint{0.749691in}{2.644559in}}%
\pgfpathlineto{\pgfqpoint{0.751811in}{2.644561in}}%
\pgfpathlineto{\pgfqpoint{0.751835in}{3.345966in}}%
\pgfpathlineto{\pgfqpoint{0.752568in}{3.341105in}}%
\pgfpathlineto{\pgfqpoint{0.752617in}{3.321062in}}%
\pgfpathlineto{\pgfqpoint{0.753161in}{3.378305in}}%
\pgfpathlineto{\pgfqpoint{0.753211in}{3.358347in}}%
\pgfpathlineto{\pgfqpoint{0.753744in}{3.414850in}}%
\pgfpathlineto{\pgfqpoint{0.753839in}{3.317717in}}%
\pgfpathlineto{\pgfqpoint{0.753853in}{3.390937in}}%
\pgfpathlineto{\pgfqpoint{0.754634in}{2.644559in}}%
\pgfpathlineto{\pgfqpoint{0.756757in}{2.644610in}}%
\pgfpathlineto{\pgfqpoint{0.756777in}{3.305053in}}%
\pgfpathlineto{\pgfqpoint{0.757518in}{3.293563in}}%
\pgfpathlineto{\pgfqpoint{0.757568in}{3.278909in}}%
\pgfpathlineto{\pgfqpoint{0.758112in}{3.330786in}}%
\pgfpathlineto{\pgfqpoint{0.758161in}{3.316177in}}%
\pgfpathlineto{\pgfqpoint{0.758692in}{3.367204in}}%
\pgfpathlineto{\pgfqpoint{0.758785in}{3.272608in}}%
\pgfpathlineto{\pgfqpoint{0.758798in}{3.345737in}}%
\pgfpathlineto{\pgfqpoint{0.759580in}{2.644559in}}%
\pgfpathlineto{\pgfqpoint{0.761702in}{2.644610in}}%
\pgfpathlineto{\pgfqpoint{0.761721in}{3.264419in}}%
\pgfpathlineto{\pgfqpoint{0.762459in}{3.241524in}}%
\pgfpathlineto{\pgfqpoint{0.762508in}{3.207731in}}%
\pgfpathlineto{\pgfqpoint{0.763151in}{3.284913in}}%
\pgfpathlineto{\pgfqpoint{0.763201in}{3.251255in}}%
\pgfpathlineto{\pgfqpoint{0.763635in}{3.315252in}}%
\pgfpathlineto{\pgfqpoint{0.763730in}{3.217377in}}%
\pgfpathlineto{\pgfqpoint{0.763745in}{3.279865in}}%
\pgfpathlineto{\pgfqpoint{0.764532in}{2.644559in}}%
\pgfpathlineto{\pgfqpoint{0.766648in}{2.644595in}}%
\pgfpathlineto{\pgfqpoint{0.766666in}{3.196741in}}%
\pgfpathlineto{\pgfqpoint{0.767434in}{3.149853in}}%
\pgfpathlineto{\pgfqpoint{0.768620in}{3.224335in}}%
\pgfpathlineto{\pgfqpoint{0.768647in}{3.206226in}}%
\pgfpathlineto{\pgfqpoint{0.768688in}{3.208367in}}%
\pgfpathlineto{\pgfqpoint{0.769526in}{2.644559in}}%
\pgfpathlineto{\pgfqpoint{0.771593in}{2.644593in}}%
\pgfpathlineto{\pgfqpoint{0.771610in}{3.110004in}}%
\pgfpathlineto{\pgfqpoint{0.772372in}{3.077388in}}%
\pgfpathlineto{\pgfqpoint{0.772421in}{3.048184in}}%
\pgfpathlineto{\pgfqpoint{0.773064in}{3.120809in}}%
\pgfpathlineto{\pgfqpoint{0.773114in}{3.091699in}}%
\pgfpathlineto{\pgfqpoint{0.773537in}{3.150471in}}%
\pgfpathlineto{\pgfqpoint{0.773620in}{3.057206in}}%
\pgfpathlineto{\pgfqpoint{0.773635in}{3.113570in}}%
\pgfpathlineto{\pgfqpoint{0.774420in}{2.644559in}}%
\pgfpathlineto{\pgfqpoint{0.776536in}{2.644559in}}%
\pgfpathlineto{\pgfqpoint{0.776557in}{3.015540in}}%
\pgfpathlineto{\pgfqpoint{0.777326in}{2.990705in}}%
\pgfpathlineto{\pgfqpoint{0.777375in}{2.941749in}}%
\pgfpathlineto{\pgfqpoint{0.778018in}{3.034078in}}%
\pgfpathlineto{\pgfqpoint{0.778067in}{2.985326in}}%
\pgfpathlineto{\pgfqpoint{0.778511in}{3.064950in}}%
\pgfpathlineto{\pgfqpoint{0.778572in}{2.959690in}}%
\pgfpathlineto{\pgfqpoint{0.778588in}{3.021857in}}%
\pgfpathlineto{\pgfqpoint{0.779372in}{2.644559in}}%
\pgfpathlineto{\pgfqpoint{0.781484in}{2.644615in}}%
\pgfpathlineto{\pgfqpoint{0.781498in}{2.937176in}}%
\pgfpathlineto{\pgfqpoint{0.782234in}{2.882083in}}%
\pgfpathlineto{\pgfqpoint{0.782284in}{2.838934in}}%
\pgfpathlineto{\pgfqpoint{0.782926in}{2.925488in}}%
\pgfpathlineto{\pgfqpoint{0.782976in}{2.882495in}}%
\pgfpathlineto{\pgfqpoint{0.783414in}{2.956042in}}%
\pgfpathlineto{\pgfqpoint{0.783512in}{2.850547in}}%
\pgfpathlineto{\pgfqpoint{0.783525in}{2.920174in}}%
\pgfpathlineto{\pgfqpoint{0.784306in}{2.644559in}}%
\pgfpathlineto{\pgfqpoint{0.786427in}{2.644559in}}%
\pgfpathlineto{\pgfqpoint{0.786445in}{2.838203in}}%
\pgfpathlineto{\pgfqpoint{0.787248in}{2.788759in}}%
\pgfpathlineto{\pgfqpoint{0.787298in}{2.726608in}}%
\pgfpathlineto{\pgfqpoint{0.787940in}{2.832116in}}%
\pgfpathlineto{\pgfqpoint{0.787990in}{2.770232in}}%
\pgfpathlineto{\pgfqpoint{0.788402in}{2.860983in}}%
\pgfpathlineto{\pgfqpoint{0.788492in}{2.644559in}}%
\pgfpathlineto{\pgfqpoint{0.788662in}{2.644559in}}%
\pgfpathlineto{\pgfqpoint{0.791373in}{2.644560in}}%
\pgfpathlineto{\pgfqpoint{0.791377in}{2.765561in}}%
\pgfpathlineto{\pgfqpoint{0.792115in}{2.701514in}}%
\pgfpathlineto{\pgfqpoint{0.792164in}{2.697290in}}%
\pgfpathlineto{\pgfqpoint{0.792510in}{2.726376in}}%
\pgfpathlineto{\pgfqpoint{0.792560in}{2.722137in}}%
\pgfpathlineto{\pgfqpoint{0.793302in}{2.776098in}}%
\pgfpathlineto{\pgfqpoint{0.793390in}{2.736225in}}%
\pgfpathlineto{\pgfqpoint{0.794230in}{2.644559in}}%
\pgfpathlineto{\pgfqpoint{0.796319in}{2.644559in}}%
\pgfpathlineto{\pgfqpoint{0.796322in}{2.752932in}}%
\pgfpathlineto{\pgfqpoint{0.797057in}{2.705242in}}%
\pgfpathlineto{\pgfqpoint{0.797107in}{2.693499in}}%
\pgfpathlineto{\pgfqpoint{0.797651in}{2.742524in}}%
\pgfpathlineto{\pgfqpoint{0.797700in}{2.730779in}}%
\pgfpathlineto{\pgfqpoint{0.798244in}{2.779804in}}%
\pgfpathlineto{\pgfqpoint{0.798348in}{2.707563in}}%
\pgfpathlineto{\pgfqpoint{0.798365in}{2.746157in}}%
\pgfpathlineto{\pgfqpoint{0.799152in}{2.644559in}}%
\pgfpathlineto{\pgfqpoint{0.801264in}{2.644559in}}%
\pgfpathlineto{\pgfqpoint{0.801269in}{2.751114in}}%
\pgfpathlineto{\pgfqpoint{0.802006in}{2.700517in}}%
\pgfpathlineto{\pgfqpoint{0.802056in}{2.697663in}}%
\pgfpathlineto{\pgfqpoint{0.802303in}{2.719162in}}%
\pgfpathlineto{\pgfqpoint{0.802352in}{2.716300in}}%
\pgfpathlineto{\pgfqpoint{0.803193in}{2.775093in}}%
\pgfpathlineto{\pgfqpoint{0.803261in}{2.759590in}}%
\pgfpathlineto{\pgfqpoint{0.804074in}{2.644559in}}%
\pgfpathlineto{\pgfqpoint{0.806210in}{2.644559in}}%
\pgfpathlineto{\pgfqpoint{0.806215in}{2.750665in}}%
\pgfpathlineto{\pgfqpoint{0.806270in}{2.644218in}}%
\pgfpathlineto{\pgfqpoint{0.806947in}{2.682047in}}%
\pgfpathlineto{\pgfqpoint{0.807788in}{2.765233in}}%
\pgfpathlineto{\pgfqpoint{0.808184in}{2.790045in}}%
\pgfpathlineto{\pgfqpoint{0.808779in}{2.644559in}}%
\pgfpathlineto{\pgfqpoint{0.811155in}{2.644559in}}%
\pgfpathlineto{\pgfqpoint{0.811160in}{2.750509in}}%
\pgfpathlineto{\pgfqpoint{0.811887in}{2.682147in}}%
\pgfpathlineto{\pgfqpoint{0.812728in}{2.763696in}}%
\pgfpathlineto{\pgfqpoint{0.813102in}{2.787133in}}%
\pgfpathlineto{\pgfqpoint{0.813729in}{2.644559in}}%
\pgfpathlineto{\pgfqpoint{0.816101in}{2.644559in}}%
\pgfpathlineto{\pgfqpoint{0.816106in}{2.750723in}}%
\pgfpathlineto{\pgfqpoint{0.816849in}{2.709510in}}%
\pgfpathlineto{\pgfqpoint{0.816899in}{2.688062in}}%
\pgfpathlineto{\pgfqpoint{0.817443in}{2.746765in}}%
\pgfpathlineto{\pgfqpoint{0.817492in}{2.725369in}}%
\pgfpathlineto{\pgfqpoint{0.818031in}{2.783659in}}%
\pgfpathlineto{\pgfqpoint{0.818129in}{2.720539in}}%
\pgfpathlineto{\pgfqpoint{0.818143in}{2.729850in}}%
\pgfpathlineto{\pgfqpoint{0.818932in}{2.644559in}}%
\pgfpathlineto{\pgfqpoint{0.821046in}{2.644559in}}%
\pgfpathlineto{\pgfqpoint{0.821051in}{2.749314in}}%
\pgfpathlineto{\pgfqpoint{0.821163in}{2.639473in}}%
\pgfpathlineto{\pgfqpoint{0.821806in}{2.716393in}}%
\pgfpathlineto{\pgfqpoint{0.821856in}{2.682638in}}%
\pgfpathlineto{\pgfqpoint{0.822498in}{2.759816in}}%
\pgfpathlineto{\pgfqpoint{0.822548in}{2.726205in}}%
\pgfpathlineto{\pgfqpoint{0.822982in}{2.790120in}}%
\pgfpathlineto{\pgfqpoint{0.823111in}{2.644559in}}%
\pgfpathlineto{\pgfqpoint{0.823176in}{2.644559in}}%
\pgfpathlineto{\pgfqpoint{0.825992in}{2.644559in}}%
\pgfpathlineto{\pgfqpoint{0.825996in}{2.746969in}}%
\pgfpathlineto{\pgfqpoint{0.826064in}{2.621312in}}%
\pgfpathlineto{\pgfqpoint{0.826749in}{2.661636in}}%
\pgfpathlineto{\pgfqpoint{0.827590in}{2.785684in}}%
\pgfpathlineto{\pgfqpoint{0.827965in}{2.809187in}}%
\pgfpathlineto{\pgfqpoint{0.828462in}{2.644559in}}%
\pgfpathlineto{\pgfqpoint{0.830938in}{2.644595in}}%
\pgfpathlineto{\pgfqpoint{0.830940in}{2.752626in}}%
\pgfpathlineto{\pgfqpoint{0.831009in}{2.642392in}}%
\pgfpathlineto{\pgfqpoint{0.831694in}{2.682591in}}%
\pgfpathlineto{\pgfqpoint{0.832534in}{2.764692in}}%
\pgfpathlineto{\pgfqpoint{0.832911in}{2.788296in}}%
\pgfpathlineto{\pgfqpoint{0.833532in}{2.644559in}}%
\pgfpathlineto{\pgfqpoint{0.835882in}{2.644559in}}%
\pgfpathlineto{\pgfqpoint{0.835887in}{2.743279in}}%
\pgfpathlineto{\pgfqpoint{0.835968in}{2.624361in}}%
\pgfpathlineto{\pgfqpoint{0.836654in}{2.665986in}}%
\pgfpathlineto{\pgfqpoint{0.837494in}{2.782298in}}%
\pgfpathlineto{\pgfqpoint{0.837856in}{2.804931in}}%
\pgfpathlineto{\pgfqpoint{0.838377in}{2.644559in}}%
\pgfpathlineto{\pgfqpoint{0.840828in}{2.644559in}}%
\pgfpathlineto{\pgfqpoint{0.840833in}{2.741825in}}%
\pgfpathlineto{\pgfqpoint{0.840911in}{2.627253in}}%
\pgfpathlineto{\pgfqpoint{0.841597in}{2.668682in}}%
\pgfpathlineto{\pgfqpoint{0.842437in}{2.778583in}}%
\pgfpathlineto{\pgfqpoint{0.842802in}{2.801387in}}%
\pgfpathlineto{\pgfqpoint{0.843329in}{2.644559in}}%
\pgfpathlineto{\pgfqpoint{0.845773in}{2.644559in}}%
\pgfpathlineto{\pgfqpoint{0.845778in}{2.744889in}}%
\pgfpathlineto{\pgfqpoint{0.845861in}{2.618540in}}%
\pgfpathlineto{\pgfqpoint{0.846497in}{2.731194in}}%
\pgfpathlineto{\pgfqpoint{0.846546in}{2.660278in}}%
\pgfpathlineto{\pgfqpoint{0.847189in}{2.774520in}}%
\pgfpathlineto{\pgfqpoint{0.847238in}{2.703942in}}%
\pgfpathlineto{\pgfqpoint{0.847747in}{2.809422in}}%
\pgfpathlineto{\pgfqpoint{0.847835in}{2.644559in}}%
\pgfpathlineto{\pgfqpoint{0.847952in}{2.644559in}}%
\pgfpathlineto{\pgfqpoint{0.850718in}{2.644559in}}%
\pgfpathlineto{\pgfqpoint{0.850724in}{2.737867in}}%
\pgfpathlineto{\pgfqpoint{0.850790in}{2.639790in}}%
\pgfpathlineto{\pgfqpoint{0.851475in}{2.680011in}}%
\pgfpathlineto{\pgfqpoint{0.852316in}{2.765202in}}%
\pgfpathlineto{\pgfqpoint{0.852693in}{2.788825in}}%
\pgfpathlineto{\pgfqpoint{0.853289in}{2.644559in}}%
\pgfpathlineto{\pgfqpoint{0.855664in}{2.644559in}}%
\pgfpathlineto{\pgfqpoint{0.855669in}{2.734613in}}%
\pgfpathlineto{\pgfqpoint{0.855738in}{2.639877in}}%
\pgfpathlineto{\pgfqpoint{0.856423in}{2.680350in}}%
\pgfpathlineto{\pgfqpoint{0.857264in}{2.764347in}}%
\pgfpathlineto{\pgfqpoint{0.857638in}{2.787827in}}%
\pgfpathlineto{\pgfqpoint{0.858242in}{2.644559in}}%
\pgfpathlineto{\pgfqpoint{0.860611in}{2.645005in}}%
\pgfpathlineto{\pgfqpoint{0.860614in}{2.765768in}}%
\pgfpathlineto{\pgfqpoint{0.860701in}{2.618820in}}%
\pgfpathlineto{\pgfqpoint{0.861337in}{2.730108in}}%
\pgfpathlineto{\pgfqpoint{0.861387in}{2.660880in}}%
\pgfpathlineto{\pgfqpoint{0.862030in}{2.773421in}}%
\pgfpathlineto{\pgfqpoint{0.862079in}{2.704557in}}%
\pgfpathlineto{\pgfqpoint{0.862584in}{2.808036in}}%
\pgfpathlineto{\pgfqpoint{0.862671in}{2.644559in}}%
\pgfpathlineto{\pgfqpoint{0.862811in}{2.644559in}}%
\pgfpathlineto{\pgfqpoint{0.865555in}{2.644559in}}%
\pgfpathlineto{\pgfqpoint{0.865560in}{2.736405in}}%
\pgfpathlineto{\pgfqpoint{0.865632in}{2.643127in}}%
\pgfpathlineto{\pgfqpoint{0.866318in}{2.683973in}}%
\pgfpathlineto{\pgfqpoint{0.867159in}{2.760397in}}%
\pgfpathlineto{\pgfqpoint{0.867529in}{2.783612in}}%
\pgfpathlineto{\pgfqpoint{0.868209in}{2.644559in}}%
\pgfpathlineto{\pgfqpoint{0.870502in}{2.644711in}}%
\pgfpathlineto{\pgfqpoint{0.870505in}{2.751492in}}%
\pgfpathlineto{\pgfqpoint{0.870601in}{2.620424in}}%
\pgfpathlineto{\pgfqpoint{0.871239in}{2.727964in}}%
\pgfpathlineto{\pgfqpoint{0.871288in}{2.663032in}}%
\pgfpathlineto{\pgfqpoint{0.871931in}{2.771279in}}%
\pgfpathlineto{\pgfqpoint{0.871981in}{2.706706in}}%
\pgfpathlineto{\pgfqpoint{0.872426in}{2.802218in}}%
\pgfpathlineto{\pgfqpoint{0.872560in}{2.644559in}}%
\pgfpathlineto{\pgfqpoint{0.872708in}{2.644559in}}%
\pgfpathlineto{\pgfqpoint{0.875448in}{2.644688in}}%
\pgfpathlineto{\pgfqpoint{0.875450in}{2.749681in}}%
\pgfpathlineto{\pgfqpoint{0.875547in}{2.623143in}}%
\pgfpathlineto{\pgfqpoint{0.876185in}{2.725050in}}%
\pgfpathlineto{\pgfqpoint{0.876235in}{2.665773in}}%
\pgfpathlineto{\pgfqpoint{0.876878in}{2.768377in}}%
\pgfpathlineto{\pgfqpoint{0.876927in}{2.709435in}}%
\pgfpathlineto{\pgfqpoint{0.877372in}{2.799324in}}%
\pgfpathlineto{\pgfqpoint{0.877505in}{2.644559in}}%
\pgfpathlineto{\pgfqpoint{0.877647in}{2.644559in}}%
\pgfpathlineto{\pgfqpoint{0.880391in}{2.644559in}}%
\pgfpathlineto{\pgfqpoint{0.880399in}{2.732436in}}%
\pgfpathlineto{\pgfqpoint{0.880457in}{2.629901in}}%
\pgfpathlineto{\pgfqpoint{0.881141in}{2.669337in}}%
\pgfpathlineto{\pgfqpoint{0.881982in}{2.771617in}}%
\pgfpathlineto{\pgfqpoint{0.882365in}{2.795653in}}%
\pgfpathlineto{\pgfqpoint{0.882885in}{2.644559in}}%
\pgfpathlineto{\pgfqpoint{0.885337in}{2.644559in}}%
\pgfpathlineto{\pgfqpoint{0.885344in}{2.733984in}}%
\pgfpathlineto{\pgfqpoint{0.885420in}{2.621850in}}%
\pgfpathlineto{\pgfqpoint{0.886106in}{2.663293in}}%
\pgfpathlineto{\pgfqpoint{0.886947in}{2.779818in}}%
\pgfpathlineto{\pgfqpoint{0.887311in}{2.802590in}}%
\pgfpathlineto{\pgfqpoint{0.887827in}{2.644559in}}%
\pgfpathlineto{\pgfqpoint{0.890282in}{2.644559in}}%
\pgfpathlineto{\pgfqpoint{0.890291in}{2.752583in}}%
\pgfpathlineto{\pgfqpoint{0.890407in}{2.634210in}}%
\pgfpathlineto{\pgfqpoint{0.891050in}{2.714870in}}%
\pgfpathlineto{\pgfqpoint{0.891100in}{2.677511in}}%
\pgfpathlineto{\pgfqpoint{0.891743in}{2.758302in}}%
\pgfpathlineto{\pgfqpoint{0.891792in}{2.721068in}}%
\pgfpathlineto{\pgfqpoint{0.892222in}{2.788371in}}%
\pgfpathlineto{\pgfqpoint{0.892341in}{2.644559in}}%
\pgfpathlineto{\pgfqpoint{0.892473in}{2.644559in}}%
\pgfpathlineto{\pgfqpoint{0.895227in}{2.644559in}}%
\pgfpathlineto{\pgfqpoint{0.895236in}{2.743901in}}%
\pgfpathlineto{\pgfqpoint{0.895314in}{2.627824in}}%
\pgfpathlineto{\pgfqpoint{0.896001in}{2.669556in}}%
\pgfpathlineto{\pgfqpoint{0.896842in}{2.773367in}}%
\pgfpathlineto{\pgfqpoint{0.897202in}{2.795896in}}%
\pgfpathlineto{\pgfqpoint{0.897745in}{2.644559in}}%
\pgfpathlineto{\pgfqpoint{0.900172in}{2.644559in}}%
\pgfpathlineto{\pgfqpoint{0.900178in}{2.748789in}}%
\pgfpathlineto{\pgfqpoint{0.900252in}{2.615185in}}%
\pgfpathlineto{\pgfqpoint{0.900937in}{2.656295in}}%
\pgfpathlineto{\pgfqpoint{0.901778in}{2.784373in}}%
\pgfpathlineto{\pgfqpoint{0.902147in}{2.807471in}}%
\pgfpathlineto{\pgfqpoint{0.902641in}{2.644559in}}%
\pgfpathlineto{\pgfqpoint{0.905118in}{2.644559in}}%
\pgfpathlineto{\pgfqpoint{0.905123in}{2.759431in}}%
\pgfpathlineto{\pgfqpoint{0.905180in}{2.640211in}}%
\pgfpathlineto{\pgfqpoint{0.905910in}{2.709474in}}%
\pgfpathlineto{\pgfqpoint{0.905960in}{2.685015in}}%
\pgfpathlineto{\pgfqpoint{0.906504in}{2.746691in}}%
\pgfpathlineto{\pgfqpoint{0.906553in}{2.722360in}}%
\pgfpathlineto{\pgfqpoint{0.907093in}{2.783604in}}%
\pgfpathlineto{\pgfqpoint{0.907182in}{2.644559in}}%
\pgfpathlineto{\pgfqpoint{0.907195in}{2.644559in}}%
\pgfpathlineto{\pgfqpoint{0.910066in}{2.644577in}}%
\pgfpathlineto{\pgfqpoint{0.910068in}{2.760523in}}%
\pgfpathlineto{\pgfqpoint{0.910195in}{2.627627in}}%
\pgfpathlineto{\pgfqpoint{0.910838in}{2.720556in}}%
\pgfpathlineto{\pgfqpoint{0.910887in}{2.671016in}}%
\pgfpathlineto{\pgfqpoint{0.911530in}{2.763948in}}%
\pgfpathlineto{\pgfqpoint{0.911579in}{2.714613in}}%
\pgfpathlineto{\pgfqpoint{0.912007in}{2.793810in}}%
\pgfpathlineto{\pgfqpoint{0.912129in}{2.644559in}}%
\pgfpathlineto{\pgfqpoint{0.912288in}{2.644559in}}%
\pgfpathlineto{\pgfqpoint{0.915011in}{2.644566in}}%
\pgfpathlineto{\pgfqpoint{0.915013in}{2.752430in}}%
\pgfpathlineto{\pgfqpoint{0.915127in}{2.632826in}}%
\pgfpathlineto{\pgfqpoint{0.915770in}{2.713248in}}%
\pgfpathlineto{\pgfqpoint{0.915820in}{2.676030in}}%
\pgfpathlineto{\pgfqpoint{0.916463in}{2.756664in}}%
\pgfpathlineto{\pgfqpoint{0.916512in}{2.719604in}}%
\pgfpathlineto{\pgfqpoint{0.916946in}{2.786942in}}%
\pgfpathlineto{\pgfqpoint{0.917074in}{2.644559in}}%
\pgfpathlineto{\pgfqpoint{0.917182in}{2.644559in}}%
\pgfpathlineto{\pgfqpoint{0.919957in}{2.644564in}}%
\pgfpathlineto{\pgfqpoint{0.919959in}{2.746722in}}%
\pgfpathlineto{\pgfqpoint{0.920016in}{2.619806in}}%
\pgfpathlineto{\pgfqpoint{0.920746in}{2.728233in}}%
\pgfpathlineto{\pgfqpoint{0.920796in}{2.664669in}}%
\pgfpathlineto{\pgfqpoint{0.921438in}{2.771593in}}%
\pgfpathlineto{\pgfqpoint{0.921488in}{2.708298in}}%
\pgfpathlineto{\pgfqpoint{0.921929in}{2.802310in}}%
\pgfpathlineto{\pgfqpoint{0.922020in}{2.644559in}}%
\pgfpathlineto{\pgfqpoint{0.922171in}{2.644559in}}%
\pgfpathlineto{\pgfqpoint{0.924903in}{2.645071in}}%
\pgfpathlineto{\pgfqpoint{0.924905in}{2.737969in}}%
\pgfpathlineto{\pgfqpoint{0.924966in}{2.618245in}}%
\pgfpathlineto{\pgfqpoint{0.925648in}{2.657554in}}%
\pgfpathlineto{\pgfqpoint{0.926489in}{2.778284in}}%
\pgfpathlineto{\pgfqpoint{0.926874in}{2.802434in}}%
\pgfpathlineto{\pgfqpoint{0.927360in}{2.644559in}}%
\pgfpathlineto{\pgfqpoint{0.929848in}{2.644624in}}%
\pgfpathlineto{\pgfqpoint{0.929851in}{2.752510in}}%
\pgfpathlineto{\pgfqpoint{0.929926in}{2.617003in}}%
\pgfpathlineto{\pgfqpoint{0.930611in}{2.658250in}}%
\pgfpathlineto{\pgfqpoint{0.931452in}{2.779447in}}%
\pgfpathlineto{\pgfqpoint{0.931820in}{2.802473in}}%
\pgfpathlineto{\pgfqpoint{0.932324in}{2.644559in}}%
\pgfpathlineto{\pgfqpoint{0.934794in}{2.644833in}}%
\pgfpathlineto{\pgfqpoint{0.934795in}{2.721030in}}%
\pgfpathlineto{\pgfqpoint{0.934881in}{2.629626in}}%
\pgfpathlineto{\pgfqpoint{0.935517in}{2.711351in}}%
\pgfpathlineto{\pgfqpoint{0.935567in}{2.671538in}}%
\pgfpathlineto{\pgfqpoint{0.936210in}{2.754732in}}%
\pgfpathlineto{\pgfqpoint{0.936259in}{2.715147in}}%
\pgfpathlineto{\pgfqpoint{0.936765in}{2.789524in}}%
\pgfpathlineto{\pgfqpoint{0.936851in}{2.644559in}}%
\pgfpathlineto{\pgfqpoint{0.936922in}{2.644559in}}%
\pgfpathlineto{\pgfqpoint{0.939739in}{2.644743in}}%
\pgfpathlineto{\pgfqpoint{0.939743in}{2.719332in}}%
\pgfpathlineto{\pgfqpoint{0.939825in}{2.636388in}}%
\pgfpathlineto{\pgfqpoint{0.940511in}{2.678143in}}%
\pgfpathlineto{\pgfqpoint{0.941351in}{2.759872in}}%
\pgfpathlineto{\pgfqpoint{0.941711in}{2.782402in}}%
\pgfpathlineto{\pgfqpoint{0.942344in}{2.644559in}}%
\pgfpathlineto{\pgfqpoint{0.944685in}{2.645063in}}%
\pgfpathlineto{\pgfqpoint{0.944687in}{2.735475in}}%
\pgfpathlineto{\pgfqpoint{0.944749in}{2.616317in}}%
\pgfpathlineto{\pgfqpoint{0.945432in}{2.655939in}}%
\pgfpathlineto{\pgfqpoint{0.946273in}{2.779013in}}%
\pgfpathlineto{\pgfqpoint{0.946656in}{2.803003in}}%
\pgfpathlineto{\pgfqpoint{0.947143in}{2.644559in}}%
\pgfpathlineto{\pgfqpoint{0.949630in}{2.644860in}}%
\pgfpathlineto{\pgfqpoint{0.949633in}{2.765085in}}%
\pgfpathlineto{\pgfqpoint{0.949693in}{2.637955in}}%
\pgfpathlineto{\pgfqpoint{0.950377in}{2.677283in}}%
\pgfpathlineto{\pgfqpoint{0.951218in}{2.756882in}}%
\pgfpathlineto{\pgfqpoint{0.951602in}{2.780930in}}%
\pgfpathlineto{\pgfqpoint{0.952223in}{2.644559in}}%
\pgfpathlineto{\pgfqpoint{0.954572in}{2.644559in}}%
\pgfpathlineto{\pgfqpoint{0.954578in}{2.758101in}}%
\pgfpathlineto{\pgfqpoint{0.954642in}{2.633082in}}%
\pgfpathlineto{\pgfqpoint{0.955377in}{2.711899in}}%
\pgfpathlineto{\pgfqpoint{0.955427in}{2.679171in}}%
\pgfpathlineto{\pgfqpoint{0.956069in}{2.755300in}}%
\pgfpathlineto{\pgfqpoint{0.956119in}{2.722760in}}%
\pgfpathlineto{\pgfqpoint{0.956547in}{2.785239in}}%
\pgfpathlineto{\pgfqpoint{0.956639in}{2.644559in}}%
\pgfpathlineto{\pgfqpoint{0.956763in}{2.644559in}}%
\pgfpathlineto{\pgfqpoint{0.959518in}{2.644559in}}%
\pgfpathlineto{\pgfqpoint{0.959523in}{2.755111in}}%
\pgfpathlineto{\pgfqpoint{0.960285in}{2.697439in}}%
\pgfpathlineto{\pgfqpoint{0.960335in}{2.688586in}}%
\pgfpathlineto{\pgfqpoint{0.960780in}{2.728480in}}%
\pgfpathlineto{\pgfqpoint{0.960829in}{2.719681in}}%
\pgfpathlineto{\pgfqpoint{0.961458in}{2.771017in}}%
\pgfpathlineto{\pgfqpoint{0.961532in}{2.718037in}}%
\pgfpathlineto{\pgfqpoint{0.961544in}{2.739552in}}%
\pgfpathlineto{\pgfqpoint{0.961571in}{2.746656in}}%
\pgfpathlineto{\pgfqpoint{0.962339in}{2.644559in}}%
\pgfpathlineto{\pgfqpoint{0.964464in}{2.644559in}}%
\pgfpathlineto{\pgfqpoint{0.964473in}{2.723117in}}%
\pgfpathlineto{\pgfqpoint{0.965196in}{2.693776in}}%
\pgfpathlineto{\pgfqpoint{0.965245in}{2.686616in}}%
\pgfpathlineto{\pgfqpoint{0.965690in}{2.724847in}}%
\pgfpathlineto{\pgfqpoint{0.965740in}{2.717682in}}%
\pgfpathlineto{\pgfqpoint{0.966438in}{2.771817in}}%
\pgfpathlineto{\pgfqpoint{0.966481in}{2.736198in}}%
\pgfpathlineto{\pgfqpoint{0.967294in}{2.644559in}}%
\pgfpathlineto{\pgfqpoint{0.969410in}{2.644559in}}%
\pgfpathlineto{\pgfqpoint{0.969418in}{2.724318in}}%
\pgfpathlineto{\pgfqpoint{0.970150in}{2.689574in}}%
\pgfpathlineto{\pgfqpoint{0.971384in}{2.766388in}}%
\pgfpathlineto{\pgfqpoint{0.971433in}{2.741811in}}%
\pgfpathlineto{\pgfqpoint{0.971460in}{2.745037in}}%
\pgfpathlineto{\pgfqpoint{0.972222in}{2.644559in}}%
\pgfpathlineto{\pgfqpoint{0.974355in}{2.644559in}}%
\pgfpathlineto{\pgfqpoint{0.974364in}{2.722587in}}%
\pgfpathlineto{\pgfqpoint{0.975095in}{2.692994in}}%
\pgfpathlineto{\pgfqpoint{0.975145in}{2.687789in}}%
\pgfpathlineto{\pgfqpoint{0.975491in}{2.717854in}}%
\pgfpathlineto{\pgfqpoint{0.975540in}{2.712639in}}%
\pgfpathlineto{\pgfqpoint{0.976282in}{2.767569in}}%
\pgfpathlineto{\pgfqpoint{0.976371in}{2.728388in}}%
\pgfpathlineto{\pgfqpoint{0.977206in}{2.644559in}}%
\pgfpathlineto{\pgfqpoint{0.979302in}{2.644561in}}%
\pgfpathlineto{\pgfqpoint{0.979305in}{2.742307in}}%
\pgfpathlineto{\pgfqpoint{0.979408in}{2.643642in}}%
\pgfpathlineto{\pgfqpoint{0.980047in}{2.694993in}}%
\pgfpathlineto{\pgfqpoint{0.980096in}{2.686358in}}%
\pgfpathlineto{\pgfqpoint{0.980541in}{2.726041in}}%
\pgfpathlineto{\pgfqpoint{0.980591in}{2.717447in}}%
\pgfpathlineto{\pgfqpoint{0.981229in}{2.769233in}}%
\pgfpathlineto{\pgfqpoint{0.981311in}{2.707290in}}%
\pgfpathlineto{\pgfqpoint{0.981320in}{2.747284in}}%
\pgfpathlineto{\pgfqpoint{0.982105in}{2.644559in}}%
\pgfpathlineto{\pgfqpoint{0.984246in}{2.644559in}}%
\pgfpathlineto{\pgfqpoint{0.984252in}{2.723598in}}%
\pgfpathlineto{\pgfqpoint{0.985004in}{2.687530in}}%
\pgfpathlineto{\pgfqpoint{0.986220in}{2.767858in}}%
\pgfpathlineto{\pgfqpoint{0.986220in}{2.754752in}}%
\pgfpathlineto{\pgfqpoint{0.986988in}{2.644559in}}%
\pgfpathlineto{\pgfqpoint{0.989191in}{2.644559in}}%
\pgfpathlineto{\pgfqpoint{0.989197in}{2.753043in}}%
\pgfpathlineto{\pgfqpoint{0.989303in}{2.631453in}}%
\pgfpathlineto{\pgfqpoint{0.989944in}{2.707272in}}%
\pgfpathlineto{\pgfqpoint{0.989994in}{2.674418in}}%
\pgfpathlineto{\pgfqpoint{0.990637in}{2.750698in}}%
\pgfpathlineto{\pgfqpoint{0.990686in}{2.717982in}}%
\pgfpathlineto{\pgfqpoint{0.991123in}{2.781237in}}%
\pgfpathlineto{\pgfqpoint{0.991256in}{2.644559in}}%
\pgfpathlineto{\pgfqpoint{0.991361in}{2.644559in}}%
\pgfpathlineto{\pgfqpoint{0.994137in}{2.644559in}}%
\pgfpathlineto{\pgfqpoint{0.994143in}{2.729384in}}%
\pgfpathlineto{\pgfqpoint{0.994201in}{2.635730in}}%
\pgfpathlineto{\pgfqpoint{0.994881in}{2.674450in}}%
\pgfpathlineto{\pgfqpoint{0.995722in}{2.755042in}}%
\pgfpathlineto{\pgfqpoint{0.996111in}{2.779428in}}%
\pgfpathlineto{\pgfqpoint{0.996714in}{2.644559in}}%
\pgfpathlineto{\pgfqpoint{0.999082in}{2.644559in}}%
\pgfpathlineto{\pgfqpoint{0.999089in}{2.727454in}}%
\pgfpathlineto{\pgfqpoint{0.999212in}{2.640429in}}%
\pgfpathlineto{\pgfqpoint{0.999855in}{2.699154in}}%
\pgfpathlineto{\pgfqpoint{0.999905in}{2.683749in}}%
\pgfpathlineto{\pgfqpoint{1.000449in}{2.736390in}}%
\pgfpathlineto{\pgfqpoint{1.000498in}{2.721076in}}%
\pgfpathlineto{\pgfqpoint{1.001024in}{2.772513in}}%
\pgfpathlineto{\pgfqpoint{1.001120in}{2.702355in}}%
\pgfpathlineto{\pgfqpoint{1.001907in}{2.644559in}}%
\pgfpathlineto{\pgfqpoint{1.004028in}{2.644559in}}%
\pgfpathlineto{\pgfqpoint{1.004034in}{2.729557in}}%
\pgfpathlineto{\pgfqpoint{1.004087in}{2.637500in}}%
\pgfpathlineto{\pgfqpoint{1.004765in}{2.675279in}}%
\pgfpathlineto{\pgfqpoint{1.005605in}{2.751992in}}%
\pgfpathlineto{\pgfqpoint{1.005977in}{2.775270in}}%
\pgfpathlineto{\pgfqpoint{1.006632in}{2.644559in}}%
\pgfpathlineto{\pgfqpoint{1.008973in}{2.644559in}}%
\pgfpathlineto{\pgfqpoint{1.008980in}{2.738312in}}%
\pgfpathlineto{\pgfqpoint{1.009069in}{2.625416in}}%
\pgfpathlineto{\pgfqpoint{1.009706in}{2.710067in}}%
\pgfpathlineto{\pgfqpoint{1.009755in}{2.667547in}}%
\pgfpathlineto{\pgfqpoint{1.010398in}{2.753437in}}%
\pgfpathlineto{\pgfqpoint{1.010447in}{2.711167in}}%
\pgfpathlineto{\pgfqpoint{1.010947in}{2.787815in}}%
\pgfpathlineto{\pgfqpoint{1.011036in}{2.644559in}}%
\pgfpathlineto{\pgfqpoint{1.011131in}{2.644559in}}%
\pgfpathlineto{\pgfqpoint{1.013919in}{2.644559in}}%
\pgfpathlineto{\pgfqpoint{1.013925in}{2.736626in}}%
\pgfpathlineto{\pgfqpoint{1.014014in}{2.629617in}}%
\pgfpathlineto{\pgfqpoint{1.014651in}{2.704370in}}%
\pgfpathlineto{\pgfqpoint{1.014700in}{2.671726in}}%
\pgfpathlineto{\pgfqpoint{1.015343in}{2.747760in}}%
\pgfpathlineto{\pgfqpoint{1.015393in}{2.715327in}}%
\pgfpathlineto{\pgfqpoint{1.015893in}{2.782160in}}%
\pgfpathlineto{\pgfqpoint{1.015979in}{2.644559in}}%
\pgfpathlineto{\pgfqpoint{1.016057in}{2.644559in}}%
\pgfpathlineto{\pgfqpoint{1.018864in}{2.644559in}}%
\pgfpathlineto{\pgfqpoint{1.018870in}{2.734598in}}%
\pgfpathlineto{\pgfqpoint{1.018953in}{2.634523in}}%
\pgfpathlineto{\pgfqpoint{1.019589in}{2.699578in}}%
\pgfpathlineto{\pgfqpoint{1.019639in}{2.676189in}}%
\pgfpathlineto{\pgfqpoint{1.020183in}{2.736789in}}%
\pgfpathlineto{\pgfqpoint{1.020232in}{2.713542in}}%
\pgfpathlineto{\pgfqpoint{1.020838in}{2.777860in}}%
\pgfpathlineto{\pgfqpoint{1.020897in}{2.701849in}}%
\pgfpathlineto{\pgfqpoint{1.020913in}{2.733003in}}%
\pgfpathlineto{\pgfqpoint{1.021696in}{2.644559in}}%
\pgfpathlineto{\pgfqpoint{1.023810in}{2.644559in}}%
\pgfpathlineto{\pgfqpoint{1.023816in}{2.740276in}}%
\pgfpathlineto{\pgfqpoint{1.023878in}{2.612892in}}%
\pgfpathlineto{\pgfqpoint{1.024560in}{2.652316in}}%
\pgfpathlineto{\pgfqpoint{1.025401in}{2.775072in}}%
\pgfpathlineto{\pgfqpoint{1.025784in}{2.799030in}}%
\pgfpathlineto{\pgfqpoint{1.026267in}{2.644559in}}%
\pgfpathlineto{\pgfqpoint{1.028755in}{2.644559in}}%
\pgfpathlineto{\pgfqpoint{1.028762in}{2.740837in}}%
\pgfpathlineto{\pgfqpoint{1.028821in}{2.613833in}}%
\pgfpathlineto{\pgfqpoint{1.029502in}{2.652788in}}%
\pgfpathlineto{\pgfqpoint{1.030342in}{2.774035in}}%
\pgfpathlineto{\pgfqpoint{1.030729in}{2.798236in}}%
\pgfpathlineto{\pgfqpoint{1.031208in}{2.644559in}}%
\pgfpathlineto{\pgfqpoint{1.033701in}{2.644559in}}%
\pgfpathlineto{\pgfqpoint{1.033707in}{2.744338in}}%
\pgfpathlineto{\pgfqpoint{1.033826in}{2.626792in}}%
\pgfpathlineto{\pgfqpoint{1.034469in}{2.710423in}}%
\pgfpathlineto{\pgfqpoint{1.034519in}{2.670075in}}%
\pgfpathlineto{\pgfqpoint{1.035162in}{2.753832in}}%
\pgfpathlineto{\pgfqpoint{1.035211in}{2.713657in}}%
\pgfpathlineto{\pgfqpoint{1.035641in}{2.783855in}}%
\pgfpathlineto{\pgfqpoint{1.035764in}{2.644559in}}%
\pgfpathlineto{\pgfqpoint{1.035870in}{2.644559in}}%
\pgfpathlineto{\pgfqpoint{1.038646in}{2.644559in}}%
\pgfpathlineto{\pgfqpoint{1.038653in}{2.746119in}}%
\pgfpathlineto{\pgfqpoint{1.038707in}{2.615115in}}%
\pgfpathlineto{\pgfqpoint{1.039383in}{2.652730in}}%
\pgfpathlineto{\pgfqpoint{1.040224in}{2.772091in}}%
\pgfpathlineto{\pgfqpoint{1.040595in}{2.795323in}}%
\pgfpathlineto{\pgfqpoint{1.041090in}{2.644559in}}%
\pgfpathlineto{\pgfqpoint{1.043592in}{2.644559in}}%
\pgfpathlineto{\pgfqpoint{1.043598in}{2.746387in}}%
\pgfpathlineto{\pgfqpoint{1.043681in}{2.628296in}}%
\pgfpathlineto{\pgfqpoint{1.044318in}{2.705037in}}%
\pgfpathlineto{\pgfqpoint{1.044367in}{2.669900in}}%
\pgfpathlineto{\pgfqpoint{1.045010in}{2.748462in}}%
\pgfpathlineto{\pgfqpoint{1.045060in}{2.713466in}}%
\pgfpathlineto{\pgfqpoint{1.045565in}{2.783264in}}%
\pgfpathlineto{\pgfqpoint{1.045656in}{2.644559in}}%
\pgfpathlineto{\pgfqpoint{1.045731in}{2.644559in}}%
\pgfpathlineto{\pgfqpoint{1.048537in}{2.644559in}}%
\pgfpathlineto{\pgfqpoint{1.048544in}{2.746654in}}%
\pgfpathlineto{\pgfqpoint{1.048676in}{2.616792in}}%
\pgfpathlineto{\pgfqpoint{1.049270in}{2.654022in}}%
\pgfpathlineto{\pgfqpoint{1.050111in}{2.769519in}}%
\pgfpathlineto{\pgfqpoint{1.050484in}{2.792880in}}%
\pgfpathlineto{\pgfqpoint{1.050982in}{2.644559in}}%
\pgfpathlineto{\pgfqpoint{1.053482in}{2.644559in}}%
\pgfpathlineto{\pgfqpoint{1.053489in}{2.746591in}}%
\pgfpathlineto{\pgfqpoint{1.053590in}{2.624485in}}%
\pgfpathlineto{\pgfqpoint{1.054228in}{2.709540in}}%
\pgfpathlineto{\pgfqpoint{1.054277in}{2.667115in}}%
\pgfpathlineto{\pgfqpoint{1.054920in}{2.752949in}}%
\pgfpathlineto{\pgfqpoint{1.054970in}{2.710696in}}%
\pgfpathlineto{\pgfqpoint{1.055411in}{2.783703in}}%
\pgfpathlineto{\pgfqpoint{1.055549in}{2.644559in}}%
\pgfpathlineto{\pgfqpoint{1.055655in}{2.644559in}}%
\pgfpathlineto{\pgfqpoint{1.058428in}{2.644559in}}%
\pgfpathlineto{\pgfqpoint{1.058435in}{2.746782in}}%
\pgfpathlineto{\pgfqpoint{1.058528in}{2.642806in}}%
\pgfpathlineto{\pgfqpoint{1.059165in}{2.689093in}}%
\pgfpathlineto{\pgfqpoint{1.059214in}{2.684899in}}%
\pgfpathlineto{\pgfqpoint{1.059560in}{2.713952in}}%
\pgfpathlineto{\pgfqpoint{1.059610in}{2.709751in}}%
\pgfpathlineto{\pgfqpoint{1.060402in}{2.766803in}}%
\pgfpathlineto{\pgfqpoint{1.060444in}{2.727662in}}%
\pgfpathlineto{\pgfqpoint{1.061266in}{2.644559in}}%
\pgfpathlineto{\pgfqpoint{1.063373in}{2.644559in}}%
\pgfpathlineto{\pgfqpoint{1.063380in}{2.746818in}}%
\pgfpathlineto{\pgfqpoint{1.063514in}{2.618886in}}%
\pgfpathlineto{\pgfqpoint{1.064107in}{2.656116in}}%
\pgfpathlineto{\pgfqpoint{1.064948in}{2.766986in}}%
\pgfpathlineto{\pgfqpoint{1.065321in}{2.790317in}}%
\pgfpathlineto{\pgfqpoint{1.065825in}{2.644559in}}%
\pgfpathlineto{\pgfqpoint{1.068319in}{2.644559in}}%
\pgfpathlineto{\pgfqpoint{1.068326in}{2.746946in}}%
\pgfpathlineto{\pgfqpoint{1.068411in}{2.622072in}}%
\pgfpathlineto{\pgfqpoint{1.069048in}{2.708597in}}%
\pgfpathlineto{\pgfqpoint{1.069098in}{2.663819in}}%
\pgfpathlineto{\pgfqpoint{1.069741in}{2.752008in}}%
\pgfpathlineto{\pgfqpoint{1.069790in}{2.707399in}}%
\pgfpathlineto{\pgfqpoint{1.070293in}{2.786603in}}%
\pgfpathlineto{\pgfqpoint{1.070383in}{2.644559in}}%
\pgfpathlineto{\pgfqpoint{1.070458in}{2.644559in}}%
\pgfpathlineto{\pgfqpoint{1.073264in}{2.644559in}}%
\pgfpathlineto{\pgfqpoint{1.073271in}{2.746890in}}%
\pgfpathlineto{\pgfqpoint{1.073377in}{2.644380in}}%
\pgfpathlineto{\pgfqpoint{1.074017in}{2.689018in}}%
\pgfpathlineto{\pgfqpoint{1.074067in}{2.687157in}}%
\pgfpathlineto{\pgfqpoint{1.074215in}{2.701447in}}%
\pgfpathlineto{\pgfqpoint{1.074265in}{2.699582in}}%
\pgfpathlineto{\pgfqpoint{1.075196in}{2.763101in}}%
\pgfpathlineto{\pgfqpoint{1.075238in}{2.739500in}}%
\pgfpathlineto{\pgfqpoint{1.075262in}{2.748390in}}%
\pgfpathlineto{\pgfqpoint{1.076010in}{2.644559in}}%
\pgfpathlineto{\pgfqpoint{1.078210in}{2.644559in}}%
\pgfpathlineto{\pgfqpoint{1.078217in}{2.746213in}}%
\pgfpathlineto{\pgfqpoint{1.078278in}{2.643669in}}%
\pgfpathlineto{\pgfqpoint{1.078959in}{2.682333in}}%
\pgfpathlineto{\pgfqpoint{1.080096in}{2.760308in}}%
\pgfpathlineto{\pgfqpoint{1.080238in}{2.710980in}}%
\pgfpathlineto{\pgfqpoint{1.080184in}{2.765789in}}%
\pgfpathlineto{\pgfqpoint{1.080251in}{2.715976in}}%
\pgfpathlineto{\pgfqpoint{1.081037in}{2.644559in}}%
\pgfpathlineto{\pgfqpoint{1.083155in}{2.644559in}}%
\pgfpathlineto{\pgfqpoint{1.083162in}{2.746498in}}%
\pgfpathlineto{\pgfqpoint{1.083273in}{2.643596in}}%
\pgfpathlineto{\pgfqpoint{1.083915in}{2.689349in}}%
\pgfpathlineto{\pgfqpoint{1.083964in}{2.686572in}}%
\pgfpathlineto{\pgfqpoint{1.084211in}{2.707995in}}%
\pgfpathlineto{\pgfqpoint{1.084261in}{2.705209in}}%
\pgfpathlineto{\pgfqpoint{1.085091in}{2.763236in}}%
\pgfpathlineto{\pgfqpoint{1.085160in}{2.729839in}}%
\pgfpathlineto{\pgfqpoint{1.086008in}{2.644559in}}%
\pgfpathlineto{\pgfqpoint{1.088101in}{2.644559in}}%
\pgfpathlineto{\pgfqpoint{1.088108in}{2.745339in}}%
\pgfpathlineto{\pgfqpoint{1.088188in}{2.643795in}}%
\pgfpathlineto{\pgfqpoint{1.088873in}{2.684871in}}%
\pgfpathlineto{\pgfqpoint{1.090074in}{2.764340in}}%
\pgfpathlineto{\pgfqpoint{1.090075in}{2.751721in}}%
\pgfpathlineto{\pgfqpoint{1.090845in}{2.644559in}}%
\pgfpathlineto{\pgfqpoint{1.093046in}{2.644559in}}%
\pgfpathlineto{\pgfqpoint{1.093053in}{2.744939in}}%
\pgfpathlineto{\pgfqpoint{1.093117in}{2.638294in}}%
\pgfpathlineto{\pgfqpoint{1.093798in}{2.677320in}}%
\pgfpathlineto{\pgfqpoint{1.094738in}{2.753295in}}%
\pgfpathlineto{\pgfqpoint{1.095020in}{2.770978in}}%
\pgfpathlineto{\pgfqpoint{1.095897in}{2.644559in}}%
\pgfpathlineto{\pgfqpoint{1.097992in}{2.644559in}}%
\pgfpathlineto{\pgfqpoint{1.097999in}{2.743891in}}%
\pgfpathlineto{\pgfqpoint{1.098121in}{2.637941in}}%
\pgfpathlineto{\pgfqpoint{1.098764in}{2.695462in}}%
\pgfpathlineto{\pgfqpoint{1.098813in}{2.681179in}}%
\pgfpathlineto{\pgfqpoint{1.099357in}{2.732708in}}%
\pgfpathlineto{\pgfqpoint{1.099407in}{2.718496in}}%
\pgfpathlineto{\pgfqpoint{1.099933in}{2.768853in}}%
\pgfpathlineto{\pgfqpoint{1.100042in}{2.678980in}}%
\pgfpathlineto{\pgfqpoint{1.100821in}{2.644559in}}%
\pgfpathlineto{\pgfqpoint{1.102937in}{2.644559in}}%
\pgfpathlineto{\pgfqpoint{1.102944in}{2.742775in}}%
\pgfpathlineto{\pgfqpoint{1.103022in}{2.618150in}}%
\pgfpathlineto{\pgfqpoint{1.103707in}{2.659073in}}%
\pgfpathlineto{\pgfqpoint{1.104547in}{2.768084in}}%
\pgfpathlineto{\pgfqpoint{1.104911in}{2.790828in}}%
\pgfpathlineto{\pgfqpoint{1.105432in}{2.644559in}}%
\pgfpathlineto{\pgfqpoint{1.107882in}{2.644559in}}%
\pgfpathlineto{\pgfqpoint{1.107890in}{2.740466in}}%
\pgfpathlineto{\pgfqpoint{1.107956in}{2.613740in}}%
\pgfpathlineto{\pgfqpoint{1.108638in}{2.653044in}}%
\pgfpathlineto{\pgfqpoint{1.109479in}{2.770379in}}%
\pgfpathlineto{\pgfqpoint{1.109856in}{2.794008in}}%
\pgfpathlineto{\pgfqpoint{1.110351in}{2.644559in}}%
\pgfpathlineto{\pgfqpoint{1.112828in}{2.644559in}}%
\pgfpathlineto{\pgfqpoint{1.112835in}{2.742240in}}%
\pgfpathlineto{\pgfqpoint{1.112913in}{2.619670in}}%
\pgfpathlineto{\pgfqpoint{1.113598in}{2.660570in}}%
\pgfpathlineto{\pgfqpoint{1.114439in}{2.765492in}}%
\pgfpathlineto{\pgfqpoint{1.114802in}{2.788224in}}%
\pgfpathlineto{\pgfqpoint{1.115332in}{2.644559in}}%
\pgfpathlineto{\pgfqpoint{1.117773in}{2.644559in}}%
\pgfpathlineto{\pgfqpoint{1.117781in}{2.740031in}}%
\pgfpathlineto{\pgfqpoint{1.117917in}{2.616499in}}%
\pgfpathlineto{\pgfqpoint{1.118511in}{2.653725in}}%
\pgfpathlineto{\pgfqpoint{1.119351in}{2.768993in}}%
\pgfpathlineto{\pgfqpoint{1.119722in}{2.792220in}}%
\pgfpathlineto{\pgfqpoint{1.120225in}{2.644559in}}%
\pgfpathlineto{\pgfqpoint{1.122719in}{2.644559in}}%
\pgfpathlineto{\pgfqpoint{1.122731in}{2.771935in}}%
\pgfpathlineto{\pgfqpoint{1.123470in}{2.683754in}}%
\pgfpathlineto{\pgfqpoint{1.124693in}{2.763356in}}%
\pgfpathlineto{\pgfqpoint{1.124693in}{2.738127in}}%
\pgfpathlineto{\pgfqpoint{1.124719in}{2.748160in}}%
\pgfpathlineto{\pgfqpoint{1.125463in}{2.644559in}}%
\pgfpathlineto{\pgfqpoint{1.127664in}{2.644559in}}%
\pgfpathlineto{\pgfqpoint{1.127676in}{2.770272in}}%
\pgfpathlineto{\pgfqpoint{1.127806in}{2.637033in}}%
\pgfpathlineto{\pgfqpoint{1.128399in}{2.674170in}}%
\pgfpathlineto{\pgfqpoint{1.129240in}{2.747919in}}%
\pgfpathlineto{\pgfqpoint{1.129612in}{2.771260in}}%
\pgfpathlineto{\pgfqpoint{1.130297in}{2.644559in}}%
\pgfpathlineto{\pgfqpoint{1.132610in}{2.644559in}}%
\pgfpathlineto{\pgfqpoint{1.132622in}{2.770943in}}%
\pgfpathlineto{\pgfqpoint{1.132690in}{2.635703in}}%
\pgfpathlineto{\pgfqpoint{1.133375in}{2.675757in}}%
\pgfpathlineto{\pgfqpoint{1.134216in}{2.749554in}}%
\pgfpathlineto{\pgfqpoint{1.134584in}{2.772614in}}%
\pgfpathlineto{\pgfqpoint{1.135275in}{2.644559in}}%
\pgfpathlineto{\pgfqpoint{1.137555in}{2.644559in}}%
\pgfpathlineto{\pgfqpoint{1.137567in}{2.771033in}}%
\pgfpathlineto{\pgfqpoint{1.137691in}{2.637172in}}%
\pgfpathlineto{\pgfqpoint{1.138285in}{2.674243in}}%
\pgfpathlineto{\pgfqpoint{1.139225in}{2.752898in}}%
\pgfpathlineto{\pgfqpoint{1.139500in}{2.770200in}}%
\pgfpathlineto{\pgfqpoint{1.140320in}{2.644559in}}%
\pgfpathlineto{\pgfqpoint{1.142501in}{2.644559in}}%
\pgfpathlineto{\pgfqpoint{1.142512in}{2.769422in}}%
\pgfpathlineto{\pgfqpoint{1.142596in}{2.625358in}}%
\pgfpathlineto{\pgfqpoint{1.143232in}{2.705087in}}%
\pgfpathlineto{\pgfqpoint{1.143281in}{2.666882in}}%
\pgfpathlineto{\pgfqpoint{1.143924in}{2.748475in}}%
\pgfpathlineto{\pgfqpoint{1.143974in}{2.710485in}}%
\pgfpathlineto{\pgfqpoint{1.144474in}{2.782941in}}%
\pgfpathlineto{\pgfqpoint{1.144560in}{2.644559in}}%
\pgfpathlineto{\pgfqpoint{1.144635in}{2.644559in}}%
\pgfpathlineto{\pgfqpoint{1.147446in}{2.644559in}}%
\pgfpathlineto{\pgfqpoint{1.147458in}{2.767717in}}%
\pgfpathlineto{\pgfqpoint{1.147533in}{2.644313in}}%
\pgfpathlineto{\pgfqpoint{1.148218in}{2.684895in}}%
\pgfpathlineto{\pgfqpoint{1.149420in}{2.764528in}}%
\pgfpathlineto{\pgfqpoint{1.149420in}{2.751713in}}%
\pgfpathlineto{\pgfqpoint{1.150190in}{2.644559in}}%
\pgfpathlineto{\pgfqpoint{1.152392in}{2.644559in}}%
\pgfpathlineto{\pgfqpoint{1.152403in}{2.765307in}}%
\pgfpathlineto{\pgfqpoint{1.152477in}{2.626130in}}%
\pgfpathlineto{\pgfqpoint{1.153162in}{2.666617in}}%
\pgfpathlineto{\pgfqpoint{1.154002in}{2.760085in}}%
\pgfpathlineto{\pgfqpoint{1.154365in}{2.782826in}}%
\pgfpathlineto{\pgfqpoint{1.154926in}{2.644559in}}%
\pgfpathlineto{\pgfqpoint{1.157337in}{2.644559in}}%
\pgfpathlineto{\pgfqpoint{1.157348in}{2.763259in}}%
\pgfpathlineto{\pgfqpoint{1.157364in}{2.640282in}}%
\pgfpathlineto{\pgfqpoint{1.158106in}{2.682161in}}%
\pgfpathlineto{\pgfqpoint{1.159144in}{2.758268in}}%
\pgfpathlineto{\pgfqpoint{1.159311in}{2.768738in}}%
\pgfpathlineto{\pgfqpoint{1.160633in}{2.644559in}}%
\pgfpathlineto{\pgfqpoint{1.162282in}{2.644559in}}%
\pgfpathlineto{\pgfqpoint{1.162294in}{2.758660in}}%
\pgfpathlineto{\pgfqpoint{1.163038in}{2.688310in}}%
\pgfpathlineto{\pgfqpoint{1.163087in}{2.688047in}}%
\pgfpathlineto{\pgfqpoint{1.164216in}{2.762265in}}%
\pgfpathlineto{\pgfqpoint{1.164256in}{2.752524in}}%
\pgfpathlineto{\pgfqpoint{1.165024in}{2.644559in}}%
\pgfpathlineto{\pgfqpoint{1.167228in}{2.644559in}}%
\pgfpathlineto{\pgfqpoint{1.167239in}{2.757057in}}%
\pgfpathlineto{\pgfqpoint{1.167317in}{2.625372in}}%
\pgfpathlineto{\pgfqpoint{1.167953in}{2.708153in}}%
\pgfpathlineto{\pgfqpoint{1.168002in}{2.666283in}}%
\pgfpathlineto{\pgfqpoint{1.168645in}{2.751544in}}%
\pgfpathlineto{\pgfqpoint{1.168695in}{2.709882in}}%
\pgfpathlineto{\pgfqpoint{1.169202in}{2.786399in}}%
\pgfpathlineto{\pgfqpoint{1.169297in}{2.644559in}}%
\pgfpathlineto{\pgfqpoint{1.169361in}{2.644559in}}%
\pgfpathlineto{\pgfqpoint{1.172173in}{2.644559in}}%
\pgfpathlineto{\pgfqpoint{1.172184in}{2.751138in}}%
\pgfpathlineto{\pgfqpoint{1.172264in}{2.613665in}}%
\pgfpathlineto{\pgfqpoint{1.172899in}{2.720122in}}%
\pgfpathlineto{\pgfqpoint{1.172948in}{2.654519in}}%
\pgfpathlineto{\pgfqpoint{1.173591in}{2.763471in}}%
\pgfpathlineto{\pgfqpoint{1.173641in}{2.698161in}}%
\pgfpathlineto{\pgfqpoint{1.174147in}{2.798250in}}%
\pgfpathlineto{\pgfqpoint{1.174231in}{2.644559in}}%
\pgfpathlineto{\pgfqpoint{1.174374in}{2.644559in}}%
\pgfpathlineto{\pgfqpoint{1.177120in}{2.644559in}}%
\pgfpathlineto{\pgfqpoint{1.177136in}{2.747439in}}%
\pgfpathlineto{\pgfqpoint{1.177205in}{2.612423in}}%
\pgfpathlineto{\pgfqpoint{1.177890in}{2.652855in}}%
\pgfpathlineto{\pgfqpoint{1.178731in}{2.777432in}}%
\pgfpathlineto{\pgfqpoint{1.179093in}{2.800051in}}%
\pgfpathlineto{\pgfqpoint{1.179597in}{2.644559in}}%
\pgfpathlineto{\pgfqpoint{1.182063in}{2.644559in}}%
\pgfpathlineto{\pgfqpoint{1.182072in}{2.781517in}}%
\pgfpathlineto{\pgfqpoint{1.182103in}{2.594460in}}%
\pgfpathmoveto{\pgfqpoint{1.182103in}{2.594460in}}%
\pgfpathlineto{\pgfqpoint{1.182113in}{2.725321in}}%
\pgfpathlineto{\pgfqpoint{1.182122in}{2.594460in}}%
\pgfpathmoveto{\pgfqpoint{1.182125in}{2.594460in}}%
\pgfpathlineto{\pgfqpoint{1.182143in}{2.719256in}}%
\pgfpathlineto{\pgfqpoint{1.182178in}{2.594460in}}%
\pgfpathmoveto{\pgfqpoint{1.182187in}{2.594460in}}%
\pgfpathlineto{\pgfqpoint{1.182231in}{2.720671in}}%
\pgfpathlineto{\pgfqpoint{1.182277in}{2.594460in}}%
\pgfpathmoveto{\pgfqpoint{1.182283in}{2.594460in}}%
\pgfpathlineto{\pgfqpoint{1.182330in}{2.726773in}}%
\pgfpathlineto{\pgfqpoint{1.182378in}{2.594460in}}%
\pgfpathmoveto{\pgfqpoint{1.182380in}{2.594460in}}%
\pgfpathlineto{\pgfqpoint{1.183121in}{2.776110in}}%
\pgfpathlineto{\pgfqpoint{1.183170in}{2.641347in}}%
\pgfpathlineto{\pgfqpoint{1.183813in}{2.819282in}}%
\pgfpathlineto{\pgfqpoint{1.183863in}{2.685159in}}%
\pgfpathlineto{\pgfqpoint{1.184000in}{2.830905in}}%
\pgfpathlineto{\pgfqpoint{1.184122in}{2.644559in}}%
\pgfpathlineto{\pgfqpoint{1.184618in}{2.644559in}}%
\pgfpathlineto{\pgfqpoint{1.187011in}{2.644562in}}%
\pgfpathlineto{\pgfqpoint{1.187026in}{2.761478in}}%
\pgfpathlineto{\pgfqpoint{1.187096in}{2.631800in}}%
\pgfpathlineto{\pgfqpoint{1.187781in}{2.672030in}}%
\pgfpathlineto{\pgfqpoint{1.188621in}{2.760704in}}%
\pgfpathlineto{\pgfqpoint{1.188984in}{2.783380in}}%
\pgfpathlineto{\pgfqpoint{1.189566in}{2.644559in}}%
\pgfpathlineto{\pgfqpoint{1.191957in}{2.644572in}}%
\pgfpathlineto{\pgfqpoint{1.191971in}{2.764653in}}%
\pgfpathlineto{\pgfqpoint{1.192045in}{2.628242in}}%
\pgfpathlineto{\pgfqpoint{1.192681in}{2.709698in}}%
\pgfpathlineto{\pgfqpoint{1.192730in}{2.668919in}}%
\pgfpathlineto{\pgfqpoint{1.193373in}{2.753095in}}%
\pgfpathlineto{\pgfqpoint{1.193422in}{2.712512in}}%
\pgfpathlineto{\pgfqpoint{1.193929in}{2.787931in}}%
\pgfpathlineto{\pgfqpoint{1.194023in}{2.644559in}}%
\pgfpathlineto{\pgfqpoint{1.194122in}{2.644559in}}%
\pgfpathlineto{\pgfqpoint{1.196900in}{2.644559in}}%
\pgfpathlineto{\pgfqpoint{1.196910in}{2.762422in}}%
\pgfpathlineto{\pgfqpoint{1.197042in}{2.625522in}}%
\pgfpathlineto{\pgfqpoint{1.197636in}{2.662659in}}%
\pgfpathlineto{\pgfqpoint{1.198476in}{2.768619in}}%
\pgfpathlineto{\pgfqpoint{1.198849in}{2.791925in}}%
\pgfpathlineto{\pgfqpoint{1.199367in}{2.644559in}}%
\pgfpathlineto{\pgfqpoint{1.201846in}{2.644559in}}%
\pgfpathlineto{\pgfqpoint{1.201863in}{2.745168in}}%
\pgfpathlineto{\pgfqpoint{1.201952in}{2.631941in}}%
\pgfpathlineto{\pgfqpoint{1.202589in}{2.711500in}}%
\pgfpathlineto{\pgfqpoint{1.202639in}{2.673940in}}%
\pgfpathlineto{\pgfqpoint{1.203281in}{2.754920in}}%
\pgfpathlineto{\pgfqpoint{1.203331in}{2.717510in}}%
\pgfpathlineto{\pgfqpoint{1.203773in}{2.785756in}}%
\pgfpathlineto{\pgfqpoint{1.203910in}{2.644559in}}%
\pgfpathlineto{\pgfqpoint{1.204006in}{2.644559in}}%
\pgfpathlineto{\pgfqpoint{1.206792in}{2.644559in}}%
\pgfpathlineto{\pgfqpoint{1.206806in}{2.751956in}}%
\pgfpathlineto{\pgfqpoint{1.206863in}{2.625596in}}%
\pgfpathlineto{\pgfqpoint{1.207541in}{2.663493in}}%
\pgfpathlineto{\pgfqpoint{1.208381in}{2.775465in}}%
\pgfpathlineto{\pgfqpoint{1.208765in}{2.799504in}}%
\pgfpathlineto{\pgfqpoint{1.209265in}{2.644559in}}%
\pgfpathlineto{\pgfqpoint{1.211737in}{2.644559in}}%
\pgfpathlineto{\pgfqpoint{1.211751in}{2.756794in}}%
\pgfpathlineto{\pgfqpoint{1.211825in}{2.634513in}}%
\pgfpathlineto{\pgfqpoint{1.212510in}{2.675205in}}%
\pgfpathlineto{\pgfqpoint{1.213350in}{2.769182in}}%
\pgfpathlineto{\pgfqpoint{1.213711in}{2.791752in}}%
\pgfpathlineto{\pgfqpoint{1.214281in}{2.644559in}}%
\pgfpathlineto{\pgfqpoint{1.216682in}{2.644559in}}%
\pgfpathlineto{\pgfqpoint{1.216696in}{2.759500in}}%
\pgfpathlineto{\pgfqpoint{1.216830in}{2.628722in}}%
\pgfpathlineto{\pgfqpoint{1.217423in}{2.665953in}}%
\pgfpathlineto{\pgfqpoint{1.218264in}{2.776915in}}%
\pgfpathlineto{\pgfqpoint{1.218656in}{2.801484in}}%
\pgfpathlineto{\pgfqpoint{1.219152in}{2.644559in}}%
\pgfpathlineto{\pgfqpoint{1.221628in}{2.644559in}}%
\pgfpathlineto{\pgfqpoint{1.221632in}{2.760199in}}%
\pgfpathlineto{\pgfqpoint{1.221719in}{2.630027in}}%
\pgfpathlineto{\pgfqpoint{1.222356in}{2.721649in}}%
\pgfpathlineto{\pgfqpoint{1.222405in}{2.671326in}}%
\pgfpathlineto{\pgfqpoint{1.223048in}{2.765057in}}%
\pgfpathlineto{\pgfqpoint{1.223098in}{2.714907in}}%
\pgfpathlineto{\pgfqpoint{1.223602in}{2.799734in}}%
\pgfpathlineto{\pgfqpoint{1.223697in}{2.644559in}}%
\pgfpathlineto{\pgfqpoint{1.223791in}{2.644559in}}%
\pgfpathlineto{\pgfqpoint{1.226573in}{2.644559in}}%
\pgfpathlineto{\pgfqpoint{1.226578in}{2.759546in}}%
\pgfpathlineto{\pgfqpoint{1.226680in}{2.635830in}}%
\pgfpathlineto{\pgfqpoint{1.227317in}{2.716707in}}%
\pgfpathlineto{\pgfqpoint{1.227367in}{2.678225in}}%
\pgfpathlineto{\pgfqpoint{1.228010in}{2.760123in}}%
\pgfpathlineto{\pgfqpoint{1.228059in}{2.721798in}}%
\pgfpathlineto{\pgfqpoint{1.228501in}{2.790924in}}%
\pgfpathlineto{\pgfqpoint{1.228639in}{2.644559in}}%
\pgfpathlineto{\pgfqpoint{1.228740in}{2.644559in}}%
\pgfpathlineto{\pgfqpoint{1.231519in}{2.644559in}}%
\pgfpathlineto{\pgfqpoint{1.231523in}{2.759719in}}%
\pgfpathlineto{\pgfqpoint{1.232271in}{2.707751in}}%
\pgfpathlineto{\pgfqpoint{1.232321in}{2.688247in}}%
\pgfpathlineto{\pgfqpoint{1.232865in}{2.745002in}}%
\pgfpathlineto{\pgfqpoint{1.232914in}{2.725559in}}%
\pgfpathlineto{\pgfqpoint{1.233451in}{2.781785in}}%
\pgfpathlineto{\pgfqpoint{1.233548in}{2.693726in}}%
\pgfpathlineto{\pgfqpoint{1.233562in}{2.753215in}}%
\pgfpathlineto{\pgfqpoint{1.234344in}{2.644559in}}%
\pgfpathlineto{\pgfqpoint{1.236464in}{2.644559in}}%
\pgfpathlineto{\pgfqpoint{1.236468in}{2.753471in}}%
\pgfpathlineto{\pgfqpoint{1.236541in}{2.621325in}}%
\pgfpathlineto{\pgfqpoint{1.237226in}{2.661435in}}%
\pgfpathlineto{\pgfqpoint{1.238067in}{2.785354in}}%
\pgfpathlineto{\pgfqpoint{1.238438in}{2.808587in}}%
\pgfpathlineto{\pgfqpoint{1.238935in}{2.644559in}}%
\pgfpathlineto{\pgfqpoint{1.241410in}{2.644559in}}%
\pgfpathlineto{\pgfqpoint{1.241414in}{2.759901in}}%
\pgfpathlineto{\pgfqpoint{1.241476in}{2.624476in}}%
\pgfpathlineto{\pgfqpoint{1.242155in}{2.662642in}}%
\pgfpathlineto{\pgfqpoint{1.242996in}{2.782124in}}%
\pgfpathlineto{\pgfqpoint{1.243384in}{2.806387in}}%
\pgfpathlineto{\pgfqpoint{1.243874in}{2.644559in}}%
\pgfpathlineto{\pgfqpoint{1.246355in}{2.644559in}}%
\pgfpathlineto{\pgfqpoint{1.246360in}{2.757479in}}%
\pgfpathlineto{\pgfqpoint{1.246499in}{2.630836in}}%
\pgfpathlineto{\pgfqpoint{1.247093in}{2.668060in}}%
\pgfpathlineto{\pgfqpoint{1.247933in}{2.775076in}}%
\pgfpathlineto{\pgfqpoint{1.248329in}{2.799858in}}%
\pgfpathlineto{\pgfqpoint{1.248809in}{2.644559in}}%
\pgfpathlineto{\pgfqpoint{1.251301in}{2.644559in}}%
\pgfpathlineto{\pgfqpoint{1.251305in}{2.757203in}}%
\pgfpathlineto{\pgfqpoint{1.251365in}{2.629938in}}%
\pgfpathlineto{\pgfqpoint{1.252040in}{2.667398in}}%
\pgfpathlineto{\pgfqpoint{1.252881in}{2.775915in}}%
\pgfpathlineto{\pgfqpoint{1.253274in}{2.800567in}}%
\pgfpathlineto{\pgfqpoint{1.253764in}{2.644559in}}%
\pgfpathlineto{\pgfqpoint{1.256246in}{2.644559in}}%
\pgfpathlineto{\pgfqpoint{1.256251in}{2.759144in}}%
\pgfpathlineto{\pgfqpoint{1.256316in}{2.624460in}}%
\pgfpathlineto{\pgfqpoint{1.256997in}{2.663321in}}%
\pgfpathlineto{\pgfqpoint{1.257837in}{2.781913in}}%
\pgfpathlineto{\pgfqpoint{1.258220in}{2.805854in}}%
\pgfpathlineto{\pgfqpoint{1.258715in}{2.644559in}}%
\pgfpathlineto{\pgfqpoint{1.261192in}{2.644559in}}%
\pgfpathlineto{\pgfqpoint{1.261196in}{2.757619in}}%
\pgfpathlineto{\pgfqpoint{1.261274in}{2.622968in}}%
\pgfpathlineto{\pgfqpoint{1.261959in}{2.663762in}}%
\pgfpathlineto{\pgfqpoint{1.262800in}{2.783212in}}%
\pgfpathlineto{\pgfqpoint{1.263165in}{2.806068in}}%
\pgfpathlineto{\pgfqpoint{1.263679in}{2.644559in}}%
\pgfpathlineto{\pgfqpoint{1.266137in}{2.644559in}}%
\pgfpathlineto{\pgfqpoint{1.266141in}{2.759468in}}%
\pgfpathlineto{\pgfqpoint{1.266279in}{2.631936in}}%
\pgfpathlineto{\pgfqpoint{1.266872in}{2.669135in}}%
\pgfpathlineto{\pgfqpoint{1.267713in}{2.774246in}}%
\pgfpathlineto{\pgfqpoint{1.268085in}{2.797553in}}%
\pgfpathlineto{\pgfqpoint{1.268613in}{2.644559in}}%
\pgfpathlineto{\pgfqpoint{1.271082in}{2.644559in}}%
\pgfpathlineto{\pgfqpoint{1.271087in}{2.762494in}}%
\pgfpathlineto{\pgfqpoint{1.271165in}{2.624775in}}%
\pgfpathlineto{\pgfqpoint{1.271850in}{2.665565in}}%
\pgfpathlineto{\pgfqpoint{1.272691in}{2.782759in}}%
\pgfpathlineto{\pgfqpoint{1.273056in}{2.805632in}}%
\pgfpathlineto{\pgfqpoint{1.273572in}{2.644559in}}%
\pgfpathlineto{\pgfqpoint{1.276028in}{2.644559in}}%
\pgfpathlineto{\pgfqpoint{1.276032in}{2.760960in}}%
\pgfpathlineto{\pgfqpoint{1.276801in}{2.704185in}}%
\pgfpathlineto{\pgfqpoint{1.276850in}{2.694885in}}%
\pgfpathlineto{\pgfqpoint{1.277296in}{2.735221in}}%
\pgfpathlineto{\pgfqpoint{1.277345in}{2.725985in}}%
\pgfpathlineto{\pgfqpoint{1.277970in}{2.777548in}}%
\pgfpathlineto{\pgfqpoint{1.278053in}{2.735808in}}%
\pgfpathlineto{\pgfqpoint{1.278864in}{2.644559in}}%
\pgfpathlineto{\pgfqpoint{1.280973in}{2.644559in}}%
\pgfpathlineto{\pgfqpoint{1.280978in}{2.761291in}}%
\pgfpathlineto{\pgfqpoint{1.281055in}{2.634417in}}%
\pgfpathlineto{\pgfqpoint{1.281739in}{2.675034in}}%
\pgfpathlineto{\pgfqpoint{1.282580in}{2.772782in}}%
\pgfpathlineto{\pgfqpoint{1.282947in}{2.795760in}}%
\pgfpathlineto{\pgfqpoint{1.283500in}{2.644559in}}%
\pgfpathlineto{\pgfqpoint{1.285919in}{2.644559in}}%
\pgfpathlineto{\pgfqpoint{1.285923in}{2.762978in}}%
\pgfpathlineto{\pgfqpoint{1.285996in}{2.637232in}}%
\pgfpathlineto{\pgfqpoint{1.286681in}{2.677436in}}%
\pgfpathlineto{\pgfqpoint{1.287522in}{2.770388in}}%
\pgfpathlineto{\pgfqpoint{1.287893in}{2.793618in}}%
\pgfpathlineto{\pgfqpoint{1.288424in}{2.644559in}}%
\pgfpathlineto{\pgfqpoint{1.290864in}{2.644559in}}%
\pgfpathlineto{\pgfqpoint{1.290869in}{2.760593in}}%
\pgfpathlineto{\pgfqpoint{1.291604in}{2.686607in}}%
\pgfpathlineto{\pgfqpoint{1.292544in}{2.763997in}}%
\pgfpathlineto{\pgfqpoint{1.292838in}{2.782472in}}%
\pgfpathlineto{\pgfqpoint{1.293700in}{2.644559in}}%
\pgfpathlineto{\pgfqpoint{1.295810in}{2.644559in}}%
\pgfpathlineto{\pgfqpoint{1.295814in}{2.759135in}}%
\pgfpathlineto{\pgfqpoint{1.295888in}{2.636099in}}%
\pgfpathlineto{\pgfqpoint{1.296573in}{2.676367in}}%
\pgfpathlineto{\pgfqpoint{1.297413in}{2.770464in}}%
\pgfpathlineto{\pgfqpoint{1.297784in}{2.793646in}}%
\pgfpathlineto{\pgfqpoint{1.298348in}{2.644559in}}%
\pgfpathlineto{\pgfqpoint{1.300755in}{2.644559in}}%
\pgfpathlineto{\pgfqpoint{1.300760in}{2.760100in}}%
\pgfpathlineto{\pgfqpoint{1.300826in}{2.642063in}}%
\pgfpathlineto{\pgfqpoint{1.301508in}{2.681211in}}%
\pgfpathlineto{\pgfqpoint{1.302349in}{2.764639in}}%
\pgfpathlineto{\pgfqpoint{1.302729in}{2.788484in}}%
\pgfpathlineto{\pgfqpoint{1.303335in}{2.644559in}}%
\pgfpathlineto{\pgfqpoint{1.305701in}{2.644559in}}%
\pgfpathlineto{\pgfqpoint{1.305705in}{2.762574in}}%
\pgfpathlineto{\pgfqpoint{1.306451in}{2.697384in}}%
\pgfpathlineto{\pgfqpoint{1.307674in}{2.773082in}}%
\pgfpathlineto{\pgfqpoint{1.307743in}{2.729511in}}%
\pgfpathlineto{\pgfqpoint{1.308529in}{2.644559in}}%
\pgfpathlineto{\pgfqpoint{1.310646in}{2.644559in}}%
\pgfpathlineto{\pgfqpoint{1.310651in}{2.762574in}}%
\pgfpathlineto{\pgfqpoint{1.310781in}{2.631936in}}%
\pgfpathlineto{\pgfqpoint{1.311374in}{2.669088in}}%
\pgfpathlineto{\pgfqpoint{1.312215in}{2.774413in}}%
\pgfpathlineto{\pgfqpoint{1.312591in}{2.797929in}}%
\pgfpathlineto{\pgfqpoint{1.313114in}{2.644559in}}%
\pgfpathlineto{\pgfqpoint{1.315591in}{2.644559in}}%
\pgfpathlineto{\pgfqpoint{1.315596in}{2.762056in}}%
\pgfpathlineto{\pgfqpoint{1.316360in}{2.692216in}}%
\pgfpathlineto{\pgfqpoint{1.317398in}{2.767273in}}%
\pgfpathlineto{\pgfqpoint{1.317565in}{2.777750in}}%
\pgfpathlineto{\pgfqpoint{1.319057in}{2.644559in}}%
\pgfpathlineto{\pgfqpoint{1.320537in}{2.644559in}}%
\pgfpathlineto{\pgfqpoint{1.320541in}{2.762475in}}%
\pgfpathlineto{\pgfqpoint{1.321290in}{2.698174in}}%
\pgfpathlineto{\pgfqpoint{1.321438in}{2.705015in}}%
\pgfpathlineto{\pgfqpoint{1.322511in}{2.772419in}}%
\pgfpathlineto{\pgfqpoint{1.322511in}{2.749594in}}%
\pgfpathlineto{\pgfqpoint{1.322532in}{2.756729in}}%
\pgfpathlineto{\pgfqpoint{1.323281in}{2.644559in}}%
\pgfpathlineto{\pgfqpoint{1.325482in}{2.644559in}}%
\pgfpathlineto{\pgfqpoint{1.325487in}{2.761730in}}%
\pgfpathlineto{\pgfqpoint{1.325567in}{2.622437in}}%
\pgfpathlineto{\pgfqpoint{1.326252in}{2.663586in}}%
\pgfpathlineto{\pgfqpoint{1.327093in}{2.784820in}}%
\pgfpathlineto{\pgfqpoint{1.327456in}{2.807558in}}%
\pgfpathlineto{\pgfqpoint{1.327966in}{2.644559in}}%
\pgfpathlineto{\pgfqpoint{1.330428in}{2.644559in}}%
\pgfpathlineto{\pgfqpoint{1.330432in}{2.761979in}}%
\pgfpathlineto{\pgfqpoint{1.330513in}{2.622548in}}%
\pgfpathlineto{\pgfqpoint{1.331198in}{2.663733in}}%
\pgfpathlineto{\pgfqpoint{1.332039in}{2.784802in}}%
\pgfpathlineto{\pgfqpoint{1.332402in}{2.807504in}}%
\pgfpathlineto{\pgfqpoint{1.332917in}{2.644559in}}%
\pgfpathlineto{\pgfqpoint{1.335373in}{2.644559in}}%
\pgfpathlineto{\pgfqpoint{1.335378in}{2.760312in}}%
\pgfpathlineto{\pgfqpoint{1.335458in}{2.621867in}}%
\pgfpathlineto{\pgfqpoint{1.336143in}{2.663014in}}%
\pgfpathlineto{\pgfqpoint{1.336984in}{2.784998in}}%
\pgfpathlineto{\pgfqpoint{1.337347in}{2.807727in}}%
\pgfpathlineto{\pgfqpoint{1.337861in}{2.644559in}}%
\pgfpathlineto{\pgfqpoint{1.340319in}{2.644559in}}%
\pgfpathlineto{\pgfqpoint{1.340323in}{2.761854in}}%
\pgfpathlineto{\pgfqpoint{1.340404in}{2.622413in}}%
\pgfpathlineto{\pgfqpoint{1.341088in}{2.663566in}}%
\pgfpathlineto{\pgfqpoint{1.341929in}{2.784883in}}%
\pgfpathlineto{\pgfqpoint{1.342293in}{2.807615in}}%
\pgfpathlineto{\pgfqpoint{1.342806in}{2.644559in}}%
\pgfpathlineto{\pgfqpoint{1.345264in}{2.644559in}}%
\pgfpathlineto{\pgfqpoint{1.345269in}{2.755614in}}%
\pgfpathlineto{\pgfqpoint{1.345350in}{2.620478in}}%
\pgfpathlineto{\pgfqpoint{1.346035in}{2.661706in}}%
\pgfpathlineto{\pgfqpoint{1.346876in}{2.785117in}}%
\pgfpathlineto{\pgfqpoint{1.347238in}{2.807774in}}%
\pgfpathlineto{\pgfqpoint{1.347747in}{2.644559in}}%
\pgfpathlineto{\pgfqpoint{1.350210in}{2.644559in}}%
\pgfpathlineto{\pgfqpoint{1.350214in}{2.762207in}}%
\pgfpathlineto{\pgfqpoint{1.350309in}{2.640000in}}%
\pgfpathlineto{\pgfqpoint{1.350947in}{2.712788in}}%
\pgfpathlineto{\pgfqpoint{1.350996in}{2.682145in}}%
\pgfpathlineto{\pgfqpoint{1.351639in}{2.756220in}}%
\pgfpathlineto{\pgfqpoint{1.351689in}{2.725703in}}%
\pgfpathlineto{\pgfqpoint{1.352184in}{2.790342in}}%
\pgfpathlineto{\pgfqpoint{1.352271in}{2.644559in}}%
\pgfpathlineto{\pgfqpoint{1.352348in}{2.644559in}}%
\pgfpathlineto{\pgfqpoint{1.355155in}{2.644559in}}%
\pgfpathlineto{\pgfqpoint{1.355160in}{2.758530in}}%
\pgfpathlineto{\pgfqpoint{1.355240in}{2.621558in}}%
\pgfpathlineto{\pgfqpoint{1.355925in}{2.662743in}}%
\pgfpathlineto{\pgfqpoint{1.356766in}{2.784826in}}%
\pgfpathlineto{\pgfqpoint{1.357129in}{2.807529in}}%
\pgfpathlineto{\pgfqpoint{1.357623in}{2.644559in}}%
\pgfpathlineto{\pgfqpoint{1.360101in}{2.644559in}}%
\pgfpathlineto{\pgfqpoint{1.360105in}{2.754222in}}%
\pgfpathlineto{\pgfqpoint{1.360187in}{2.620125in}}%
\pgfpathlineto{\pgfqpoint{1.360872in}{2.661380in}}%
\pgfpathlineto{\pgfqpoint{1.361712in}{2.785099in}}%
\pgfpathlineto{\pgfqpoint{1.362074in}{2.807733in}}%
\pgfpathlineto{\pgfqpoint{1.362584in}{2.644559in}}%
\pgfpathlineto{\pgfqpoint{1.365046in}{2.644559in}}%
\pgfpathlineto{\pgfqpoint{1.365051in}{2.762530in}}%
\pgfpathlineto{\pgfqpoint{1.365179in}{2.632747in}}%
\pgfpathlineto{\pgfqpoint{1.365772in}{2.669872in}}%
\pgfpathlineto{\pgfqpoint{1.366613in}{2.773351in}}%
\pgfpathlineto{\pgfqpoint{1.366990in}{2.796938in}}%
\pgfpathlineto{\pgfqpoint{1.367518in}{2.644559in}}%
\pgfpathlineto{\pgfqpoint{1.369992in}{2.644559in}}%
\pgfpathlineto{\pgfqpoint{1.369996in}{2.756572in}}%
\pgfpathlineto{\pgfqpoint{1.370755in}{2.691139in}}%
\pgfpathlineto{\pgfqpoint{1.371695in}{2.761503in}}%
\pgfpathlineto{\pgfqpoint{1.371965in}{2.778471in}}%
\pgfpathlineto{\pgfqpoint{1.373168in}{2.644559in}}%
\pgfpathlineto{\pgfqpoint{1.374937in}{2.644559in}}%
\pgfpathlineto{\pgfqpoint{1.374941in}{2.762194in}}%
\pgfpathlineto{\pgfqpoint{1.375683in}{2.684862in}}%
\pgfpathlineto{\pgfqpoint{1.376524in}{2.760894in}}%
\pgfpathlineto{\pgfqpoint{1.376911in}{2.785156in}}%
\pgfpathlineto{\pgfqpoint{1.377582in}{2.644559in}}%
\pgfpathlineto{\pgfqpoint{1.379882in}{2.644559in}}%
\pgfpathlineto{\pgfqpoint{1.379887in}{2.757731in}}%
\pgfpathlineto{\pgfqpoint{1.379959in}{2.633406in}}%
\pgfpathlineto{\pgfqpoint{1.380644in}{2.673763in}}%
\pgfpathlineto{\pgfqpoint{1.381485in}{2.772618in}}%
\pgfpathlineto{\pgfqpoint{1.381856in}{2.795900in}}%
\pgfpathlineto{\pgfqpoint{1.382400in}{2.644559in}}%
\pgfpathlineto{\pgfqpoint{1.384828in}{2.644559in}}%
\pgfpathlineto{\pgfqpoint{1.384832in}{2.760542in}}%
\pgfpathlineto{\pgfqpoint{1.385591in}{2.706987in}}%
\pgfpathlineto{\pgfqpoint{1.385640in}{2.690782in}}%
\pgfpathlineto{\pgfqpoint{1.386184in}{2.744230in}}%
\pgfpathlineto{\pgfqpoint{1.386234in}{2.728102in}}%
\pgfpathlineto{\pgfqpoint{1.386765in}{2.780666in}}%
\pgfpathlineto{\pgfqpoint{1.386864in}{2.702792in}}%
\pgfpathlineto{\pgfqpoint{1.387646in}{2.644559in}}%
\pgfpathlineto{\pgfqpoint{1.389773in}{2.644559in}}%
\pgfpathlineto{\pgfqpoint{1.389778in}{2.756110in}}%
\pgfpathlineto{\pgfqpoint{1.390534in}{2.705907in}}%
\pgfpathlineto{\pgfqpoint{1.390584in}{2.690327in}}%
\pgfpathlineto{\pgfqpoint{1.391128in}{2.743150in}}%
\pgfpathlineto{\pgfqpoint{1.391177in}{2.727646in}}%
\pgfpathlineto{\pgfqpoint{1.391709in}{2.779660in}}%
\pgfpathlineto{\pgfqpoint{1.391810in}{2.701230in}}%
\pgfpathlineto{\pgfqpoint{1.392595in}{2.644559in}}%
\pgfpathlineto{\pgfqpoint{1.394719in}{2.644559in}}%
\pgfpathlineto{\pgfqpoint{1.394723in}{2.757798in}}%
\pgfpathlineto{\pgfqpoint{1.395482in}{2.706506in}}%
\pgfpathlineto{\pgfqpoint{1.395531in}{2.690448in}}%
\pgfpathlineto{\pgfqpoint{1.396075in}{2.743749in}}%
\pgfpathlineto{\pgfqpoint{1.396124in}{2.727767in}}%
\pgfpathlineto{\pgfqpoint{1.396656in}{2.780195in}}%
\pgfpathlineto{\pgfqpoint{1.396755in}{2.702667in}}%
\pgfpathlineto{\pgfqpoint{1.397537in}{2.644559in}}%
\pgfpathlineto{\pgfqpoint{1.399664in}{2.644559in}}%
\pgfpathlineto{\pgfqpoint{1.399669in}{2.755658in}}%
\pgfpathlineto{\pgfqpoint{1.400414in}{2.684211in}}%
\pgfpathlineto{\pgfqpoint{1.401255in}{2.760141in}}%
\pgfpathlineto{\pgfqpoint{1.401638in}{2.784183in}}%
\pgfpathlineto{\pgfqpoint{1.402312in}{2.644559in}}%
\pgfpathlineto{\pgfqpoint{1.404610in}{2.644559in}}%
\pgfpathlineto{\pgfqpoint{1.404614in}{2.759498in}}%
\pgfpathlineto{\pgfqpoint{1.404681in}{2.643544in}}%
\pgfpathlineto{\pgfqpoint{1.405365in}{2.683098in}}%
\pgfpathlineto{\pgfqpoint{1.406206in}{2.762982in}}%
\pgfpathlineto{\pgfqpoint{1.406584in}{2.786687in}}%
\pgfpathlineto{\pgfqpoint{1.407218in}{2.644559in}}%
\pgfpathlineto{\pgfqpoint{1.409555in}{2.644559in}}%
\pgfpathlineto{\pgfqpoint{1.409560in}{2.754927in}}%
\pgfpathlineto{\pgfqpoint{1.410320in}{2.702441in}}%
\pgfpathlineto{\pgfqpoint{1.410370in}{2.694014in}}%
\pgfpathlineto{\pgfqpoint{1.410815in}{2.733487in}}%
\pgfpathlineto{\pgfqpoint{1.410864in}{2.725104in}}%
\pgfpathlineto{\pgfqpoint{1.411493in}{2.776075in}}%
\pgfpathlineto{\pgfqpoint{1.411565in}{2.715125in}}%
\pgfpathlineto{\pgfqpoint{1.411574in}{2.752125in}}%
\pgfpathlineto{\pgfqpoint{1.412360in}{2.644559in}}%
\pgfpathlineto{\pgfqpoint{1.414501in}{2.644559in}}%
\pgfpathlineto{\pgfqpoint{1.414505in}{2.756609in}}%
\pgfpathlineto{\pgfqpoint{1.415232in}{2.686440in}}%
\pgfpathlineto{\pgfqpoint{1.416172in}{2.762103in}}%
\pgfpathlineto{\pgfqpoint{1.416447in}{2.779359in}}%
\pgfpathlineto{\pgfqpoint{1.417373in}{2.644559in}}%
\pgfpathlineto{\pgfqpoint{1.419446in}{2.644559in}}%
\pgfpathlineto{\pgfqpoint{1.419451in}{2.754335in}}%
\pgfpathlineto{\pgfqpoint{1.420209in}{2.706009in}}%
\pgfpathlineto{\pgfqpoint{1.420259in}{2.690020in}}%
\pgfpathlineto{\pgfqpoint{1.420803in}{2.743251in}}%
\pgfpathlineto{\pgfqpoint{1.420852in}{2.727340in}}%
\pgfpathlineto{\pgfqpoint{1.421383in}{2.779686in}}%
\pgfpathlineto{\pgfqpoint{1.421483in}{2.703058in}}%
\pgfpathlineto{\pgfqpoint{1.422267in}{2.644559in}}%
\pgfpathlineto{\pgfqpoint{1.424392in}{2.644559in}}%
\pgfpathlineto{\pgfqpoint{1.424396in}{2.755171in}}%
\pgfpathlineto{\pgfqpoint{1.425120in}{2.686729in}}%
\pgfpathlineto{\pgfqpoint{1.426060in}{2.761047in}}%
\pgfpathlineto{\pgfqpoint{1.426336in}{2.778395in}}%
\pgfpathlineto{\pgfqpoint{1.427305in}{2.644559in}}%
\pgfpathlineto{\pgfqpoint{1.429337in}{2.644559in}}%
\pgfpathlineto{\pgfqpoint{1.429341in}{2.755336in}}%
\pgfpathlineto{\pgfqpoint{1.430104in}{2.706792in}}%
\pgfpathlineto{\pgfqpoint{1.430153in}{2.690014in}}%
\pgfpathlineto{\pgfqpoint{1.430697in}{2.744033in}}%
\pgfpathlineto{\pgfqpoint{1.430747in}{2.727336in}}%
\pgfpathlineto{\pgfqpoint{1.431276in}{2.780341in}}%
\pgfpathlineto{\pgfqpoint{1.431374in}{2.706457in}}%
\pgfpathlineto{\pgfqpoint{1.432158in}{2.644559in}}%
\pgfpathlineto{\pgfqpoint{1.434282in}{2.644559in}}%
\pgfpathlineto{\pgfqpoint{1.434287in}{2.757431in}}%
\pgfpathlineto{\pgfqpoint{1.435027in}{2.690683in}}%
\pgfpathlineto{\pgfqpoint{1.436066in}{2.765954in}}%
\pgfpathlineto{\pgfqpoint{1.436256in}{2.777913in}}%
\pgfpathlineto{\pgfqpoint{1.437708in}{2.644559in}}%
\pgfpathlineto{\pgfqpoint{1.439228in}{2.644559in}}%
\pgfpathlineto{\pgfqpoint{1.439232in}{2.757448in}}%
\pgfpathlineto{\pgfqpoint{1.439980in}{2.700283in}}%
\pgfpathlineto{\pgfqpoint{1.440030in}{2.695294in}}%
\pgfpathlineto{\pgfqpoint{1.440376in}{2.725147in}}%
\pgfpathlineto{\pgfqpoint{1.440425in}{2.720139in}}%
\pgfpathlineto{\pgfqpoint{1.441160in}{2.774408in}}%
\pgfpathlineto{\pgfqpoint{1.441232in}{2.744337in}}%
\pgfpathlineto{\pgfqpoint{1.442053in}{2.644559in}}%
\pgfpathlineto{\pgfqpoint{1.444173in}{2.644559in}}%
\pgfpathlineto{\pgfqpoint{1.444178in}{2.756422in}}%
\pgfpathlineto{\pgfqpoint{1.444245in}{2.634485in}}%
\pgfpathlineto{\pgfqpoint{1.444932in}{2.674306in}}%
\pgfpathlineto{\pgfqpoint{1.445773in}{2.771454in}}%
\pgfpathlineto{\pgfqpoint{1.446147in}{2.794932in}}%
\pgfpathlineto{\pgfqpoint{1.446696in}{2.644559in}}%
\pgfpathlineto{\pgfqpoint{1.449119in}{2.644559in}}%
\pgfpathlineto{\pgfqpoint{1.449123in}{2.755735in}}%
\pgfpathlineto{\pgfqpoint{1.449185in}{2.641251in}}%
\pgfpathlineto{\pgfqpoint{1.449866in}{2.679981in}}%
\pgfpathlineto{\pgfqpoint{1.450707in}{2.764120in}}%
\pgfpathlineto{\pgfqpoint{1.451093in}{2.788304in}}%
\pgfpathlineto{\pgfqpoint{1.451685in}{2.644559in}}%
\pgfpathlineto{\pgfqpoint{1.454064in}{2.644559in}}%
\pgfpathlineto{\pgfqpoint{1.454069in}{2.754590in}}%
\pgfpathlineto{\pgfqpoint{1.454140in}{2.633686in}}%
\pgfpathlineto{\pgfqpoint{1.454825in}{2.673952in}}%
\pgfpathlineto{\pgfqpoint{1.455666in}{2.771456in}}%
\pgfpathlineto{\pgfqpoint{1.456038in}{2.794785in}}%
\pgfpathlineto{\pgfqpoint{1.456588in}{2.644559in}}%
\pgfpathlineto{\pgfqpoint{1.459010in}{2.644559in}}%
\pgfpathlineto{\pgfqpoint{1.459014in}{2.759101in}}%
\pgfpathlineto{\pgfqpoint{1.459755in}{2.704907in}}%
\pgfpathlineto{\pgfqpoint{1.459804in}{2.690248in}}%
\pgfpathlineto{\pgfqpoint{1.460348in}{2.742188in}}%
\pgfpathlineto{\pgfqpoint{1.460398in}{2.727530in}}%
\pgfpathlineto{\pgfqpoint{1.460938in}{2.779224in}}%
\pgfpathlineto{\pgfqpoint{1.461039in}{2.702282in}}%
\pgfpathlineto{\pgfqpoint{1.461055in}{2.746579in}}%
\pgfpathlineto{\pgfqpoint{1.461838in}{2.644559in}}%
\pgfpathlineto{\pgfqpoint{1.463955in}{2.644559in}}%
\pgfpathlineto{\pgfqpoint{1.463960in}{2.753448in}}%
\pgfpathlineto{\pgfqpoint{1.464697in}{2.698318in}}%
\pgfpathlineto{\pgfqpoint{1.464747in}{2.694833in}}%
\pgfpathlineto{\pgfqpoint{1.465093in}{2.723183in}}%
\pgfpathlineto{\pgfqpoint{1.465142in}{2.719677in}}%
\pgfpathlineto{\pgfqpoint{1.465884in}{2.772909in}}%
\pgfpathlineto{\pgfqpoint{1.465959in}{2.746213in}}%
\pgfpathlineto{\pgfqpoint{1.466771in}{2.644559in}}%
\pgfpathlineto{\pgfqpoint{1.468901in}{2.644559in}}%
\pgfpathlineto{\pgfqpoint{1.468905in}{2.756361in}}%
\pgfpathlineto{\pgfqpoint{1.468977in}{2.633372in}}%
\pgfpathlineto{\pgfqpoint{1.469662in}{2.673707in}}%
\pgfpathlineto{\pgfqpoint{1.470503in}{2.772266in}}%
\pgfpathlineto{\pgfqpoint{1.470874in}{2.795559in}}%
\pgfpathlineto{\pgfqpoint{1.471423in}{2.644559in}}%
\pgfpathlineto{\pgfqpoint{1.473846in}{2.644559in}}%
\pgfpathlineto{\pgfqpoint{1.473851in}{2.753181in}}%
\pgfpathlineto{\pgfqpoint{1.474580in}{2.693164in}}%
\pgfpathlineto{\pgfqpoint{1.475793in}{2.771936in}}%
\pgfpathlineto{\pgfqpoint{1.475820in}{2.750590in}}%
\pgfpathlineto{\pgfqpoint{1.475841in}{2.756000in}}%
\pgfpathlineto{\pgfqpoint{1.476592in}{2.644559in}}%
\pgfpathlineto{\pgfqpoint{1.478792in}{2.644559in}}%
\pgfpathlineto{\pgfqpoint{1.478796in}{2.758792in}}%
\pgfpathlineto{\pgfqpoint{1.479547in}{2.705037in}}%
\pgfpathlineto{\pgfqpoint{1.479597in}{2.691366in}}%
\pgfpathlineto{\pgfqpoint{1.480141in}{2.742321in}}%
\pgfpathlineto{\pgfqpoint{1.480190in}{2.728645in}}%
\pgfpathlineto{\pgfqpoint{1.480725in}{2.779023in}}%
\pgfpathlineto{\pgfqpoint{1.480820in}{2.715825in}}%
\pgfpathlineto{\pgfqpoint{1.480835in}{2.732458in}}%
\pgfpathlineto{\pgfqpoint{1.481619in}{2.644559in}}%
\pgfpathlineto{\pgfqpoint{1.483737in}{2.644559in}}%
\pgfpathlineto{\pgfqpoint{1.483742in}{2.749591in}}%
\pgfpathlineto{\pgfqpoint{1.484499in}{2.693221in}}%
\pgfpathlineto{\pgfqpoint{1.485637in}{2.771398in}}%
\pgfpathlineto{\pgfqpoint{1.485711in}{2.776041in}}%
\pgfpathlineto{\pgfqpoint{1.485775in}{2.711301in}}%
\pgfpathlineto{\pgfqpoint{1.486558in}{2.644559in}}%
\pgfpathlineto{\pgfqpoint{1.488682in}{2.644559in}}%
\pgfpathlineto{\pgfqpoint{1.488687in}{2.755246in}}%
\pgfpathlineto{\pgfqpoint{1.488789in}{2.632423in}}%
\pgfpathlineto{\pgfqpoint{1.489429in}{2.719179in}}%
\pgfpathlineto{\pgfqpoint{1.489478in}{2.675069in}}%
\pgfpathlineto{\pgfqpoint{1.490121in}{2.762606in}}%
\pgfpathlineto{\pgfqpoint{1.490171in}{2.718631in}}%
\pgfpathlineto{\pgfqpoint{1.490611in}{2.793343in}}%
\pgfpathlineto{\pgfqpoint{1.490740in}{2.644559in}}%
\pgfpathlineto{\pgfqpoint{1.490860in}{2.644559in}}%
\pgfpathlineto{\pgfqpoint{1.493628in}{2.644559in}}%
\pgfpathlineto{\pgfqpoint{1.493633in}{2.749506in}}%
\pgfpathlineto{\pgfqpoint{1.494362in}{2.693428in}}%
\pgfpathlineto{\pgfqpoint{1.495575in}{2.772887in}}%
\pgfpathlineto{\pgfqpoint{1.495602in}{2.762579in}}%
\pgfpathlineto{\pgfqpoint{1.496374in}{2.644559in}}%
\pgfpathlineto{\pgfqpoint{1.498573in}{2.644559in}}%
\pgfpathlineto{\pgfqpoint{1.498579in}{2.749682in}}%
\pgfpathlineto{\pgfqpoint{1.499316in}{2.692114in}}%
\pgfpathlineto{\pgfqpoint{1.500454in}{2.770373in}}%
\pgfpathlineto{\pgfqpoint{1.500547in}{2.776239in}}%
\pgfpathlineto{\pgfqpoint{1.500611in}{2.711641in}}%
\pgfpathlineto{\pgfqpoint{1.501390in}{2.644559in}}%
\pgfpathlineto{\pgfqpoint{1.503519in}{2.644559in}}%
\pgfpathlineto{\pgfqpoint{1.503524in}{2.749416in}}%
\pgfpathlineto{\pgfqpoint{1.504280in}{2.690844in}}%
\pgfpathlineto{\pgfqpoint{1.505219in}{2.761001in}}%
\pgfpathlineto{\pgfqpoint{1.505493in}{2.778170in}}%
\pgfpathlineto{\pgfqpoint{1.506703in}{2.644559in}}%
\pgfpathlineto{\pgfqpoint{1.508464in}{2.644559in}}%
\pgfpathlineto{\pgfqpoint{1.508470in}{2.749551in}}%
\pgfpathlineto{\pgfqpoint{1.509191in}{2.695669in}}%
\pgfpathlineto{\pgfqpoint{1.509537in}{2.715197in}}%
\pgfpathlineto{\pgfqpoint{1.510438in}{2.774091in}}%
\pgfpathlineto{\pgfqpoint{1.510465in}{2.759199in}}%
\pgfpathlineto{\pgfqpoint{1.511258in}{2.644559in}}%
\pgfpathlineto{\pgfqpoint{1.513410in}{2.644559in}}%
\pgfpathlineto{\pgfqpoint{1.513414in}{2.752269in}}%
\pgfpathlineto{\pgfqpoint{1.513540in}{2.631949in}}%
\pgfpathlineto{\pgfqpoint{1.514183in}{2.721480in}}%
\pgfpathlineto{\pgfqpoint{1.514232in}{2.675271in}}%
\pgfpathlineto{\pgfqpoint{1.514875in}{2.764900in}}%
\pgfpathlineto{\pgfqpoint{1.514924in}{2.718840in}}%
\pgfpathlineto{\pgfqpoint{1.515352in}{2.794796in}}%
\pgfpathlineto{\pgfqpoint{1.515475in}{2.644559in}}%
\pgfpathlineto{\pgfqpoint{1.515604in}{2.644559in}}%
\pgfpathlineto{\pgfqpoint{1.518355in}{2.644559in}}%
\pgfpathlineto{\pgfqpoint{1.518360in}{2.756034in}}%
\pgfpathlineto{\pgfqpoint{1.519082in}{2.690235in}}%
\pgfpathlineto{\pgfqpoint{1.520120in}{2.763903in}}%
\pgfpathlineto{\pgfqpoint{1.520299in}{2.775118in}}%
\pgfpathlineto{\pgfqpoint{1.520392in}{2.708264in}}%
\pgfpathlineto{\pgfqpoint{1.521173in}{2.644559in}}%
\pgfpathlineto{\pgfqpoint{1.523301in}{2.644559in}}%
\pgfpathlineto{\pgfqpoint{1.523306in}{2.749615in}}%
\pgfpathlineto{\pgfqpoint{1.524073in}{2.700766in}}%
\pgfpathlineto{\pgfqpoint{1.524122in}{2.696492in}}%
\pgfpathlineto{\pgfqpoint{1.524469in}{2.725628in}}%
\pgfpathlineto{\pgfqpoint{1.524518in}{2.721338in}}%
\pgfpathlineto{\pgfqpoint{1.525242in}{2.774252in}}%
\pgfpathlineto{\pgfqpoint{1.525325in}{2.734749in}}%
\pgfpathlineto{\pgfqpoint{1.526139in}{2.644559in}}%
\pgfpathlineto{\pgfqpoint{1.528246in}{2.644559in}}%
\pgfpathlineto{\pgfqpoint{1.528251in}{2.752468in}}%
\pgfpathlineto{\pgfqpoint{1.529004in}{2.700593in}}%
\pgfpathlineto{\pgfqpoint{1.529053in}{2.694291in}}%
\pgfpathlineto{\pgfqpoint{1.529498in}{2.731669in}}%
\pgfpathlineto{\pgfqpoint{1.529548in}{2.725351in}}%
\pgfpathlineto{\pgfqpoint{1.530181in}{2.774533in}}%
\pgfpathlineto{\pgfqpoint{1.530270in}{2.736390in}}%
\pgfpathlineto{\pgfqpoint{1.531077in}{2.644559in}}%
\pgfpathlineto{\pgfqpoint{1.533192in}{2.644559in}}%
\pgfpathlineto{\pgfqpoint{1.533197in}{2.749351in}}%
\pgfpathlineto{\pgfqpoint{1.533921in}{2.694708in}}%
\pgfpathlineto{\pgfqpoint{1.535165in}{2.772952in}}%
\pgfpathlineto{\pgfqpoint{1.535216in}{2.743428in}}%
\pgfpathlineto{\pgfqpoint{1.536013in}{2.644559in}}%
\pgfpathlineto{\pgfqpoint{1.538137in}{2.644559in}}%
\pgfpathlineto{\pgfqpoint{1.538141in}{2.751724in}}%
\pgfpathlineto{\pgfqpoint{1.538203in}{2.631365in}}%
\pgfpathlineto{\pgfqpoint{1.538885in}{2.670147in}}%
\pgfpathlineto{\pgfqpoint{1.539726in}{2.772926in}}%
\pgfpathlineto{\pgfqpoint{1.540111in}{2.797074in}}%
\pgfpathlineto{\pgfqpoint{1.540629in}{2.644559in}}%
\pgfpathlineto{\pgfqpoint{1.543082in}{2.644559in}}%
\pgfpathlineto{\pgfqpoint{1.543087in}{2.752770in}}%
\pgfpathlineto{\pgfqpoint{1.543102in}{2.644513in}}%
\pgfpathlineto{\pgfqpoint{1.543851in}{2.690251in}}%
\pgfpathlineto{\pgfqpoint{1.544790in}{2.761934in}}%
\pgfpathlineto{\pgfqpoint{1.545056in}{2.778651in}}%
\pgfpathlineto{\pgfqpoint{1.546166in}{2.644559in}}%
\pgfpathlineto{\pgfqpoint{1.548028in}{2.644559in}}%
\pgfpathlineto{\pgfqpoint{1.548033in}{2.749520in}}%
\pgfpathlineto{\pgfqpoint{1.548775in}{2.689011in}}%
\pgfpathlineto{\pgfqpoint{1.549715in}{2.761310in}}%
\pgfpathlineto{\pgfqpoint{1.550002in}{2.779350in}}%
\pgfpathlineto{\pgfqpoint{1.551053in}{2.644559in}}%
\pgfpathlineto{\pgfqpoint{1.552973in}{2.644559in}}%
\pgfpathlineto{\pgfqpoint{1.552979in}{2.749832in}}%
\pgfpathlineto{\pgfqpoint{1.553741in}{2.700673in}}%
\pgfpathlineto{\pgfqpoint{1.553790in}{2.696129in}}%
\pgfpathlineto{\pgfqpoint{1.554136in}{2.725533in}}%
\pgfpathlineto{\pgfqpoint{1.554186in}{2.720978in}}%
\pgfpathlineto{\pgfqpoint{1.554913in}{2.774310in}}%
\pgfpathlineto{\pgfqpoint{1.554999in}{2.730644in}}%
\pgfpathlineto{\pgfqpoint{1.555830in}{2.644559in}}%
\pgfpathlineto{\pgfqpoint{1.557919in}{2.644559in}}%
\pgfpathlineto{\pgfqpoint{1.557924in}{2.749895in}}%
\pgfpathlineto{\pgfqpoint{1.558673in}{2.699232in}}%
\pgfpathlineto{\pgfqpoint{1.558722in}{2.695842in}}%
\pgfpathlineto{\pgfqpoint{1.559068in}{2.724090in}}%
\pgfpathlineto{\pgfqpoint{1.559118in}{2.720693in}}%
\pgfpathlineto{\pgfqpoint{1.559851in}{2.773284in}}%
\pgfpathlineto{\pgfqpoint{1.559917in}{2.758103in}}%
\pgfpathlineto{\pgfqpoint{1.560719in}{2.644559in}}%
\pgfpathlineto{\pgfqpoint{1.562864in}{2.644559in}}%
\pgfpathlineto{\pgfqpoint{1.562870in}{2.749896in}}%
\pgfpathlineto{\pgfqpoint{1.563590in}{2.697974in}}%
\pgfpathlineto{\pgfqpoint{1.563640in}{2.693632in}}%
\pgfpathlineto{\pgfqpoint{1.563986in}{2.722833in}}%
\pgfpathlineto{\pgfqpoint{1.564035in}{2.718483in}}%
\pgfpathlineto{\pgfqpoint{1.564838in}{2.776382in}}%
\pgfpathlineto{\pgfqpoint{1.564862in}{2.757467in}}%
\pgfpathlineto{\pgfqpoint{1.565673in}{2.644559in}}%
\pgfpathlineto{\pgfqpoint{1.567810in}{2.644559in}}%
\pgfpathlineto{\pgfqpoint{1.567815in}{2.749706in}}%
\pgfpathlineto{\pgfqpoint{1.568567in}{2.699716in}}%
\pgfpathlineto{\pgfqpoint{1.568616in}{2.695447in}}%
\pgfpathlineto{\pgfqpoint{1.568963in}{2.724577in}}%
\pgfpathlineto{\pgfqpoint{1.569012in}{2.720296in}}%
\pgfpathlineto{\pgfqpoint{1.569744in}{2.773674in}}%
\pgfpathlineto{\pgfqpoint{1.569835in}{2.733154in}}%
\pgfpathlineto{\pgfqpoint{1.570650in}{2.644559in}}%
\pgfpathlineto{\pgfqpoint{1.572755in}{2.644559in}}%
\pgfpathlineto{\pgfqpoint{1.572761in}{2.749682in}}%
\pgfpathlineto{\pgfqpoint{1.573528in}{2.693693in}}%
\pgfpathlineto{\pgfqpoint{1.574665in}{2.771741in}}%
\pgfpathlineto{\pgfqpoint{1.574729in}{2.775750in}}%
\pgfpathlineto{\pgfqpoint{1.576312in}{2.644559in}}%
\pgfpathlineto{\pgfqpoint{1.577701in}{2.644559in}}%
\pgfpathlineto{\pgfqpoint{1.577706in}{2.749867in}}%
\pgfpathlineto{\pgfqpoint{1.578439in}{2.698502in}}%
\pgfpathlineto{\pgfqpoint{1.578488in}{2.694692in}}%
\pgfpathlineto{\pgfqpoint{1.578834in}{2.723362in}}%
\pgfpathlineto{\pgfqpoint{1.578884in}{2.719541in}}%
\pgfpathlineto{\pgfqpoint{1.579626in}{2.773080in}}%
\pgfpathlineto{\pgfqpoint{1.579706in}{2.738743in}}%
\pgfpathlineto{\pgfqpoint{1.580579in}{2.644559in}}%
\pgfpathlineto{\pgfqpoint{1.582646in}{2.644559in}}%
\pgfpathlineto{\pgfqpoint{1.582652in}{2.749563in}}%
\pgfpathlineto{\pgfqpoint{1.583394in}{2.698688in}}%
\pgfpathlineto{\pgfqpoint{1.583443in}{2.694890in}}%
\pgfpathlineto{\pgfqpoint{1.583789in}{2.723549in}}%
\pgfpathlineto{\pgfqpoint{1.583839in}{2.719738in}}%
\pgfpathlineto{\pgfqpoint{1.584576in}{2.772945in}}%
\pgfpathlineto{\pgfqpoint{1.584651in}{2.741118in}}%
\pgfpathlineto{\pgfqpoint{1.585488in}{2.644559in}}%
\pgfpathlineto{\pgfqpoint{1.587592in}{2.644559in}}%
\pgfpathlineto{\pgfqpoint{1.587597in}{2.749508in}}%
\pgfpathlineto{\pgfqpoint{1.588353in}{2.699664in}}%
\pgfpathlineto{\pgfqpoint{1.588402in}{2.695731in}}%
\pgfpathlineto{\pgfqpoint{1.588748in}{2.724527in}}%
\pgfpathlineto{\pgfqpoint{1.588798in}{2.720578in}}%
\pgfpathlineto{\pgfqpoint{1.589528in}{2.773503in}}%
\pgfpathlineto{\pgfqpoint{1.589596in}{2.741698in}}%
\pgfpathlineto{\pgfqpoint{1.590429in}{2.644559in}}%
\pgfpathlineto{\pgfqpoint{1.592537in}{2.644559in}}%
\pgfpathlineto{\pgfqpoint{1.592541in}{2.751356in}}%
\pgfpathlineto{\pgfqpoint{1.592602in}{2.631511in}}%
\pgfpathlineto{\pgfqpoint{1.593282in}{2.669928in}}%
\pgfpathlineto{\pgfqpoint{1.594122in}{2.772620in}}%
\pgfpathlineto{\pgfqpoint{1.594511in}{2.796978in}}%
\pgfpathlineto{\pgfqpoint{1.595027in}{2.644559in}}%
\pgfpathlineto{\pgfqpoint{1.597482in}{2.644559in}}%
\pgfpathlineto{\pgfqpoint{1.597488in}{2.749793in}}%
\pgfpathlineto{\pgfqpoint{1.598219in}{2.690927in}}%
\pgfpathlineto{\pgfqpoint{1.599258in}{2.764307in}}%
\pgfpathlineto{\pgfqpoint{1.599431in}{2.775211in}}%
\pgfpathlineto{\pgfqpoint{1.599519in}{2.705929in}}%
\pgfpathlineto{\pgfqpoint{1.600303in}{2.644559in}}%
\pgfpathlineto{\pgfqpoint{1.602428in}{2.644559in}}%
\pgfpathlineto{\pgfqpoint{1.602433in}{2.749548in}}%
\pgfpathlineto{\pgfqpoint{1.603186in}{2.698779in}}%
\pgfpathlineto{\pgfqpoint{1.603236in}{2.695742in}}%
\pgfpathlineto{\pgfqpoint{1.603483in}{2.717423in}}%
\pgfpathlineto{\pgfqpoint{1.603532in}{2.714380in}}%
\pgfpathlineto{\pgfqpoint{1.604363in}{2.772699in}}%
\pgfpathlineto{\pgfqpoint{1.604426in}{2.757538in}}%
\pgfpathlineto{\pgfqpoint{1.605227in}{2.644559in}}%
\pgfpathlineto{\pgfqpoint{1.607373in}{2.644559in}}%
\pgfpathlineto{\pgfqpoint{1.607379in}{2.749647in}}%
\pgfpathlineto{\pgfqpoint{1.608146in}{2.700722in}}%
\pgfpathlineto{\pgfqpoint{1.608195in}{2.696461in}}%
\pgfpathlineto{\pgfqpoint{1.608541in}{2.725584in}}%
\pgfpathlineto{\pgfqpoint{1.608591in}{2.721308in}}%
\pgfpathlineto{\pgfqpoint{1.609315in}{2.774206in}}%
\pgfpathlineto{\pgfqpoint{1.609398in}{2.736897in}}%
\pgfpathlineto{\pgfqpoint{1.610206in}{2.644559in}}%
\pgfpathlineto{\pgfqpoint{1.612319in}{2.644559in}}%
\pgfpathlineto{\pgfqpoint{1.612324in}{2.749366in}}%
\pgfpathlineto{\pgfqpoint{1.613045in}{2.696862in}}%
\pgfpathlineto{\pgfqpoint{1.613094in}{2.693486in}}%
\pgfpathlineto{\pgfqpoint{1.613441in}{2.721724in}}%
\pgfpathlineto{\pgfqpoint{1.613490in}{2.718333in}}%
\pgfpathlineto{\pgfqpoint{1.614293in}{2.775259in}}%
\pgfpathlineto{\pgfqpoint{1.614325in}{2.735193in}}%
\pgfpathlineto{\pgfqpoint{1.615318in}{2.644559in}}%
\pgfpathlineto{\pgfqpoint{1.617264in}{2.644559in}}%
\pgfpathlineto{\pgfqpoint{1.617270in}{2.749739in}}%
\pgfpathlineto{\pgfqpoint{1.617997in}{2.690924in}}%
\pgfpathlineto{\pgfqpoint{1.619134in}{2.769848in}}%
\pgfpathlineto{\pgfqpoint{1.619211in}{2.774669in}}%
\pgfpathlineto{\pgfqpoint{1.619301in}{2.706629in}}%
\pgfpathlineto{\pgfqpoint{1.620085in}{2.644559in}}%
\pgfpathlineto{\pgfqpoint{1.622210in}{2.644559in}}%
\pgfpathlineto{\pgfqpoint{1.622214in}{2.748134in}}%
\pgfpathlineto{\pgfqpoint{1.622307in}{2.626514in}}%
\pgfpathlineto{\pgfqpoint{1.622945in}{2.722179in}}%
\pgfpathlineto{\pgfqpoint{1.622994in}{2.668695in}}%
\pgfpathlineto{\pgfqpoint{1.623637in}{2.765576in}}%
\pgfpathlineto{\pgfqpoint{1.623687in}{2.712289in}}%
\pgfpathlineto{\pgfqpoint{1.624184in}{2.799800in}}%
\pgfpathlineto{\pgfqpoint{1.624274in}{2.644559in}}%
\pgfpathlineto{\pgfqpoint{1.624401in}{2.644559in}}%
\pgfpathlineto{\pgfqpoint{1.627155in}{2.644559in}}%
\pgfpathlineto{\pgfqpoint{1.627160in}{2.753656in}}%
\pgfpathlineto{\pgfqpoint{1.627255in}{2.632224in}}%
\pgfpathlineto{\pgfqpoint{1.627895in}{2.718376in}}%
\pgfpathlineto{\pgfqpoint{1.627944in}{2.674552in}}%
\pgfpathlineto{\pgfqpoint{1.628587in}{2.761803in}}%
\pgfpathlineto{\pgfqpoint{1.628636in}{2.718114in}}%
\pgfpathlineto{\pgfqpoint{1.629082in}{2.792822in}}%
\pgfpathlineto{\pgfqpoint{1.629212in}{2.644559in}}%
\pgfpathlineto{\pgfqpoint{1.629332in}{2.644559in}}%
\pgfpathlineto{\pgfqpoint{1.632101in}{2.644559in}}%
\pgfpathlineto{\pgfqpoint{1.632106in}{2.749721in}}%
\pgfpathlineto{\pgfqpoint{1.632836in}{2.697388in}}%
\pgfpathlineto{\pgfqpoint{1.632885in}{2.694747in}}%
\pgfpathlineto{\pgfqpoint{1.633133in}{2.716032in}}%
\pgfpathlineto{\pgfqpoint{1.633182in}{2.713385in}}%
\pgfpathlineto{\pgfqpoint{1.634074in}{2.775212in}}%
\pgfpathlineto{\pgfqpoint{1.634098in}{2.757405in}}%
\pgfpathlineto{\pgfqpoint{1.634904in}{2.644559in}}%
\pgfpathlineto{\pgfqpoint{1.637046in}{2.644559in}}%
\pgfpathlineto{\pgfqpoint{1.637051in}{2.749691in}}%
\pgfpathlineto{\pgfqpoint{1.637804in}{2.698899in}}%
\pgfpathlineto{\pgfqpoint{1.637853in}{2.695948in}}%
\pgfpathlineto{\pgfqpoint{1.638100in}{2.717544in}}%
\pgfpathlineto{\pgfqpoint{1.638150in}{2.714586in}}%
\pgfpathlineto{\pgfqpoint{1.638980in}{2.772843in}}%
\pgfpathlineto{\pgfqpoint{1.639044in}{2.757274in}}%
\pgfpathlineto{\pgfqpoint{1.639851in}{2.644559in}}%
\pgfpathlineto{\pgfqpoint{1.641992in}{2.644559in}}%
\pgfpathlineto{\pgfqpoint{1.641997in}{2.749780in}}%
\pgfpathlineto{\pgfqpoint{1.642742in}{2.698655in}}%
\pgfpathlineto{\pgfqpoint{1.642792in}{2.695643in}}%
\pgfpathlineto{\pgfqpoint{1.643039in}{2.717299in}}%
\pgfpathlineto{\pgfqpoint{1.643089in}{2.714281in}}%
\pgfpathlineto{\pgfqpoint{1.643923in}{2.772801in}}%
\pgfpathlineto{\pgfqpoint{1.643990in}{2.758030in}}%
\pgfpathlineto{\pgfqpoint{1.644791in}{2.644559in}}%
\pgfpathlineto{\pgfqpoint{1.646937in}{2.644559in}}%
\pgfpathlineto{\pgfqpoint{1.646942in}{2.749649in}}%
\pgfpathlineto{\pgfqpoint{1.647015in}{2.642931in}}%
\pgfpathlineto{\pgfqpoint{1.647705in}{2.683765in}}%
\pgfpathlineto{\pgfqpoint{1.648545in}{2.761920in}}%
\pgfpathlineto{\pgfqpoint{1.648911in}{2.784854in}}%
\pgfpathlineto{\pgfqpoint{1.649556in}{2.644559in}}%
\pgfpathlineto{\pgfqpoint{1.651882in}{2.644559in}}%
\pgfpathlineto{\pgfqpoint{1.651887in}{2.748343in}}%
\pgfpathlineto{\pgfqpoint{1.651961in}{2.628153in}}%
\pgfpathlineto{\pgfqpoint{1.652648in}{2.668878in}}%
\pgfpathlineto{\pgfqpoint{1.653489in}{2.775449in}}%
\pgfpathlineto{\pgfqpoint{1.653856in}{2.798481in}}%
\pgfpathlineto{\pgfqpoint{1.654388in}{2.644559in}}%
\pgfpathlineto{\pgfqpoint{1.656828in}{2.644559in}}%
\pgfpathlineto{\pgfqpoint{1.656833in}{2.749308in}}%
\pgfpathlineto{\pgfqpoint{1.657595in}{2.692603in}}%
\pgfpathlineto{\pgfqpoint{1.658634in}{2.765731in}}%
\pgfpathlineto{\pgfqpoint{1.658856in}{2.718827in}}%
\pgfpathlineto{\pgfqpoint{1.658802in}{2.776287in}}%
\pgfpathlineto{\pgfqpoint{1.658870in}{2.727862in}}%
\pgfpathlineto{\pgfqpoint{1.659654in}{2.644559in}}%
\pgfpathlineto{\pgfqpoint{1.661773in}{2.644559in}}%
\pgfpathlineto{\pgfqpoint{1.661779in}{2.749704in}}%
\pgfpathlineto{\pgfqpoint{1.661862in}{2.642852in}}%
\pgfpathlineto{\pgfqpoint{1.662502in}{2.706778in}}%
\pgfpathlineto{\pgfqpoint{1.662551in}{2.684583in}}%
\pgfpathlineto{\pgfqpoint{1.663095in}{2.744028in}}%
\pgfpathlineto{\pgfqpoint{1.663145in}{2.721895in}}%
\pgfpathlineto{\pgfqpoint{1.663747in}{2.784924in}}%
\pgfpathlineto{\pgfqpoint{1.663822in}{2.671483in}}%
\pgfpathlineto{\pgfqpoint{1.664603in}{2.644559in}}%
\pgfpathlineto{\pgfqpoint{1.666719in}{2.644559in}}%
\pgfpathlineto{\pgfqpoint{1.666724in}{2.749222in}}%
\pgfpathlineto{\pgfqpoint{1.666804in}{2.632878in}}%
\pgfpathlineto{\pgfqpoint{1.667489in}{2.674274in}}%
\pgfpathlineto{\pgfqpoint{1.668330in}{2.771412in}}%
\pgfpathlineto{\pgfqpoint{1.668693in}{2.794112in}}%
\pgfpathlineto{\pgfqpoint{1.669240in}{2.644559in}}%
\pgfpathlineto{\pgfqpoint{1.671664in}{2.644559in}}%
\pgfpathlineto{\pgfqpoint{1.671670in}{2.747668in}}%
\pgfpathlineto{\pgfqpoint{1.671741in}{2.628261in}}%
\pgfpathlineto{\pgfqpoint{1.672427in}{2.668715in}}%
\pgfpathlineto{\pgfqpoint{1.673268in}{2.774367in}}%
\pgfpathlineto{\pgfqpoint{1.673638in}{2.797546in}}%
\pgfpathlineto{\pgfqpoint{1.674159in}{2.644559in}}%
\pgfpathlineto{\pgfqpoint{1.676610in}{2.644559in}}%
\pgfpathlineto{\pgfqpoint{1.676615in}{2.749848in}}%
\pgfpathlineto{\pgfqpoint{1.677343in}{2.691188in}}%
\pgfpathlineto{\pgfqpoint{1.678481in}{2.770105in}}%
\pgfpathlineto{\pgfqpoint{1.678557in}{2.774891in}}%
\pgfpathlineto{\pgfqpoint{1.678646in}{2.707208in}}%
\pgfpathlineto{\pgfqpoint{1.679427in}{2.644559in}}%
\pgfpathlineto{\pgfqpoint{1.681555in}{2.644559in}}%
\pgfpathlineto{\pgfqpoint{1.681560in}{2.748871in}}%
\pgfpathlineto{\pgfqpoint{1.681637in}{2.636000in}}%
\pgfpathlineto{\pgfqpoint{1.682322in}{2.677085in}}%
\pgfpathlineto{\pgfqpoint{1.683163in}{2.767386in}}%
\pgfpathlineto{\pgfqpoint{1.683529in}{2.790321in}}%
\pgfpathlineto{\pgfqpoint{1.684111in}{2.644559in}}%
\pgfpathlineto{\pgfqpoint{1.686501in}{2.644559in}}%
\pgfpathlineto{\pgfqpoint{1.686506in}{2.749714in}}%
\pgfpathlineto{\pgfqpoint{1.687245in}{2.691775in}}%
\pgfpathlineto{\pgfqpoint{1.688383in}{2.769993in}}%
\pgfpathlineto{\pgfqpoint{1.688529in}{2.716950in}}%
\pgfpathlineto{\pgfqpoint{1.688474in}{2.775746in}}%
\pgfpathlineto{\pgfqpoint{1.688543in}{2.729553in}}%
\pgfpathlineto{\pgfqpoint{1.689327in}{2.644559in}}%
\pgfpathlineto{\pgfqpoint{1.691446in}{2.644559in}}%
\pgfpathlineto{\pgfqpoint{1.691451in}{2.749386in}}%
\pgfpathlineto{\pgfqpoint{1.691510in}{2.642471in}}%
\pgfpathlineto{\pgfqpoint{1.692190in}{2.680973in}}%
\pgfpathlineto{\pgfqpoint{1.693031in}{2.761648in}}%
\pgfpathlineto{\pgfqpoint{1.693420in}{2.786042in}}%
\pgfpathlineto{\pgfqpoint{1.694036in}{2.644559in}}%
\pgfpathlineto{\pgfqpoint{1.696393in}{2.645225in}}%
\pgfpathlineto{\pgfqpoint{1.696394in}{2.747204in}}%
\pgfpathlineto{\pgfqpoint{1.696462in}{2.620357in}}%
\pgfpathlineto{\pgfqpoint{1.697145in}{2.660106in}}%
\pgfpathlineto{\pgfqpoint{1.697986in}{2.782945in}}%
\pgfpathlineto{\pgfqpoint{1.698365in}{2.806681in}}%
\pgfpathlineto{\pgfqpoint{1.698860in}{2.644559in}}%
\pgfpathlineto{\pgfqpoint{1.701337in}{2.644559in}}%
\pgfpathlineto{\pgfqpoint{1.701342in}{2.749731in}}%
\pgfpathlineto{\pgfqpoint{1.701437in}{2.642992in}}%
\pgfpathlineto{\pgfqpoint{1.702077in}{2.707571in}}%
\pgfpathlineto{\pgfqpoint{1.702127in}{2.685407in}}%
\pgfpathlineto{\pgfqpoint{1.702671in}{2.744821in}}%
\pgfpathlineto{\pgfqpoint{1.702720in}{2.722720in}}%
\pgfpathlineto{\pgfqpoint{1.703264in}{2.782069in}}%
\pgfpathlineto{\pgfqpoint{1.703366in}{2.695112in}}%
\pgfpathlineto{\pgfqpoint{1.703380in}{2.750317in}}%
\pgfpathlineto{\pgfqpoint{1.704158in}{2.644559in}}%
\pgfpathlineto{\pgfqpoint{1.706282in}{2.644559in}}%
\pgfpathlineto{\pgfqpoint{1.706288in}{2.749196in}}%
\pgfpathlineto{\pgfqpoint{1.706368in}{2.632280in}}%
\pgfpathlineto{\pgfqpoint{1.707054in}{2.673742in}}%
\pgfpathlineto{\pgfqpoint{1.707895in}{2.771784in}}%
\pgfpathlineto{\pgfqpoint{1.708256in}{2.794427in}}%
\pgfpathlineto{\pgfqpoint{1.708803in}{2.644559in}}%
\pgfpathlineto{\pgfqpoint{1.711228in}{2.644559in}}%
\pgfpathlineto{\pgfqpoint{1.711233in}{2.749593in}}%
\pgfpathlineto{\pgfqpoint{1.711300in}{2.643096in}}%
\pgfpathlineto{\pgfqpoint{1.711988in}{2.683151in}}%
\pgfpathlineto{\pgfqpoint{1.712829in}{2.761423in}}%
\pgfpathlineto{\pgfqpoint{1.713202in}{2.784813in}}%
\pgfpathlineto{\pgfqpoint{1.713857in}{2.644559in}}%
\pgfpathlineto{\pgfqpoint{1.716173in}{2.644559in}}%
\pgfpathlineto{\pgfqpoint{1.716179in}{2.747211in}}%
\pgfpathlineto{\pgfqpoint{1.716288in}{2.626699in}}%
\pgfpathlineto{\pgfqpoint{1.716931in}{2.723197in}}%
\pgfpathlineto{\pgfqpoint{1.716981in}{2.669816in}}%
\pgfpathlineto{\pgfqpoint{1.717623in}{2.766589in}}%
\pgfpathlineto{\pgfqpoint{1.717673in}{2.713414in}}%
\pgfpathlineto{\pgfqpoint{1.718108in}{2.796940in}}%
\pgfpathlineto{\pgfqpoint{1.718238in}{2.644559in}}%
\pgfpathlineto{\pgfqpoint{1.718340in}{2.644559in}}%
\pgfpathlineto{\pgfqpoint{1.721119in}{2.644559in}}%
\pgfpathlineto{\pgfqpoint{1.721124in}{2.749602in}}%
\pgfpathlineto{\pgfqpoint{1.721871in}{2.699175in}}%
\pgfpathlineto{\pgfqpoint{1.721921in}{2.695082in}}%
\pgfpathlineto{\pgfqpoint{1.722267in}{2.724035in}}%
\pgfpathlineto{\pgfqpoint{1.722316in}{2.719931in}}%
\pgfpathlineto{\pgfqpoint{1.723051in}{2.773281in}}%
\pgfpathlineto{\pgfqpoint{1.723144in}{2.733192in}}%
\pgfpathlineto{\pgfqpoint{1.723961in}{2.644559in}}%
\pgfpathlineto{\pgfqpoint{1.726064in}{2.644559in}}%
\pgfpathlineto{\pgfqpoint{1.726070in}{2.749060in}}%
\pgfpathlineto{\pgfqpoint{1.726148in}{2.634410in}}%
\pgfpathlineto{\pgfqpoint{1.726833in}{2.675656in}}%
\pgfpathlineto{\pgfqpoint{1.727674in}{2.769468in}}%
\pgfpathlineto{\pgfqpoint{1.728038in}{2.792286in}}%
\pgfpathlineto{\pgfqpoint{1.728603in}{2.644559in}}%
\pgfpathlineto{\pgfqpoint{1.731010in}{2.644559in}}%
\pgfpathlineto{\pgfqpoint{1.731015in}{2.749521in}}%
\pgfpathlineto{\pgfqpoint{1.731114in}{2.642286in}}%
\pgfpathlineto{\pgfqpoint{1.731755in}{2.707896in}}%
\pgfpathlineto{\pgfqpoint{1.731804in}{2.684894in}}%
\pgfpathlineto{\pgfqpoint{1.732348in}{2.745145in}}%
\pgfpathlineto{\pgfqpoint{1.732398in}{2.722209in}}%
\pgfpathlineto{\pgfqpoint{1.732938in}{2.782157in}}%
\pgfpathlineto{\pgfqpoint{1.733038in}{2.697009in}}%
\pgfpathlineto{\pgfqpoint{1.733053in}{2.747921in}}%
\pgfpathlineto{\pgfqpoint{1.733836in}{2.644559in}}%
\pgfpathlineto{\pgfqpoint{1.735955in}{2.644559in}}%
\pgfpathlineto{\pgfqpoint{1.735961in}{2.749378in}}%
\pgfpathlineto{\pgfqpoint{1.736708in}{2.692100in}}%
\pgfpathlineto{\pgfqpoint{1.737846in}{2.769625in}}%
\pgfpathlineto{\pgfqpoint{1.737984in}{2.714314in}}%
\pgfpathlineto{\pgfqpoint{1.737929in}{2.774852in}}%
\pgfpathlineto{\pgfqpoint{1.737998in}{2.731365in}}%
\pgfpathlineto{\pgfqpoint{1.738785in}{2.644559in}}%
\pgfpathlineto{\pgfqpoint{1.740901in}{2.644559in}}%
\pgfpathlineto{\pgfqpoint{1.740906in}{2.749639in}}%
\pgfpathlineto{\pgfqpoint{1.740990in}{2.642635in}}%
\pgfpathlineto{\pgfqpoint{1.741630in}{2.706819in}}%
\pgfpathlineto{\pgfqpoint{1.741680in}{2.684421in}}%
\pgfpathlineto{\pgfqpoint{1.742224in}{2.744069in}}%
\pgfpathlineto{\pgfqpoint{1.742273in}{2.721734in}}%
\pgfpathlineto{\pgfqpoint{1.742874in}{2.784912in}}%
\pgfpathlineto{\pgfqpoint{1.742949in}{2.693563in}}%
\pgfpathlineto{\pgfqpoint{1.743732in}{2.644559in}}%
\pgfpathlineto{\pgfqpoint{1.745846in}{2.644559in}}%
\pgfpathlineto{\pgfqpoint{1.745852in}{2.749018in}}%
\pgfpathlineto{\pgfqpoint{1.745966in}{2.640876in}}%
\pgfpathlineto{\pgfqpoint{1.746609in}{2.709438in}}%
\pgfpathlineto{\pgfqpoint{1.746658in}{2.684058in}}%
\pgfpathlineto{\pgfqpoint{1.747202in}{2.746683in}}%
\pgfpathlineto{\pgfqpoint{1.747252in}{2.721376in}}%
\pgfpathlineto{\pgfqpoint{1.747783in}{2.783125in}}%
\pgfpathlineto{\pgfqpoint{1.747875in}{2.696344in}}%
\pgfpathlineto{\pgfqpoint{1.747890in}{2.747751in}}%
\pgfpathlineto{\pgfqpoint{1.748674in}{2.644559in}}%
\pgfpathlineto{\pgfqpoint{1.750792in}{2.644559in}}%
\pgfpathlineto{\pgfqpoint{1.750797in}{2.749596in}}%
\pgfpathlineto{\pgfqpoint{1.751558in}{2.693099in}}%
\pgfpathlineto{\pgfqpoint{1.752695in}{2.771099in}}%
\pgfpathlineto{\pgfqpoint{1.752820in}{2.718488in}}%
\pgfpathlineto{\pgfqpoint{1.752765in}{2.775495in}}%
\pgfpathlineto{\pgfqpoint{1.752834in}{2.727654in}}%
\pgfpathlineto{\pgfqpoint{1.753618in}{2.644559in}}%
\pgfpathlineto{\pgfqpoint{1.755737in}{2.644559in}}%
\pgfpathlineto{\pgfqpoint{1.755742in}{2.749166in}}%
\pgfpathlineto{\pgfqpoint{1.755823in}{2.632354in}}%
\pgfpathlineto{\pgfqpoint{1.756508in}{2.673807in}}%
\pgfpathlineto{\pgfqpoint{1.757349in}{2.771601in}}%
\pgfpathlineto{\pgfqpoint{1.757711in}{2.794250in}}%
\pgfpathlineto{\pgfqpoint{1.758264in}{2.644559in}}%
\pgfpathlineto{\pgfqpoint{1.760682in}{2.644559in}}%
\pgfpathlineto{\pgfqpoint{1.760688in}{2.749494in}}%
\pgfpathlineto{\pgfqpoint{1.760796in}{2.642436in}}%
\pgfpathlineto{\pgfqpoint{1.761437in}{2.708534in}}%
\pgfpathlineto{\pgfqpoint{1.761486in}{2.685389in}}%
\pgfpathlineto{\pgfqpoint{1.762030in}{2.745782in}}%
\pgfpathlineto{\pgfqpoint{1.762080in}{2.722704in}}%
\pgfpathlineto{\pgfqpoint{1.762615in}{2.782492in}}%
\pgfpathlineto{\pgfqpoint{1.762711in}{2.698355in}}%
\pgfpathlineto{\pgfqpoint{1.762725in}{2.746715in}}%
\pgfpathlineto{\pgfqpoint{1.763508in}{2.644559in}}%
\pgfpathlineto{\pgfqpoint{1.765628in}{2.644559in}}%
\pgfpathlineto{\pgfqpoint{1.765633in}{2.749331in}}%
\pgfpathlineto{\pgfqpoint{1.766382in}{2.698285in}}%
\pgfpathlineto{\pgfqpoint{1.766432in}{2.695106in}}%
\pgfpathlineto{\pgfqpoint{1.766679in}{2.716929in}}%
\pgfpathlineto{\pgfqpoint{1.766729in}{2.713744in}}%
\pgfpathlineto{\pgfqpoint{1.767561in}{2.772318in}}%
\pgfpathlineto{\pgfqpoint{1.767626in}{2.757210in}}%
\pgfpathlineto{\pgfqpoint{1.768426in}{2.644559in}}%
\pgfpathlineto{\pgfqpoint{1.770572in}{2.644559in}}%
\pgfpathlineto{\pgfqpoint{1.770578in}{2.751739in}}%
\pgfpathlineto{\pgfqpoint{1.770643in}{2.641292in}}%
\pgfpathlineto{\pgfqpoint{1.771326in}{2.680934in}}%
\pgfpathlineto{\pgfqpoint{1.772167in}{2.762618in}}%
\pgfpathlineto{\pgfqpoint{1.772547in}{2.786451in}}%
\pgfpathlineto{\pgfqpoint{1.773164in}{2.644559in}}%
\pgfpathlineto{\pgfqpoint{1.775519in}{2.644559in}}%
\pgfpathlineto{\pgfqpoint{1.775524in}{2.747101in}}%
\pgfpathlineto{\pgfqpoint{1.775596in}{2.640719in}}%
\pgfpathlineto{\pgfqpoint{1.776281in}{2.681330in}}%
\pgfpathlineto{\pgfqpoint{1.777122in}{2.762187in}}%
\pgfpathlineto{\pgfqpoint{1.777493in}{2.785410in}}%
\pgfpathlineto{\pgfqpoint{1.778127in}{2.644559in}}%
\pgfpathlineto{\pgfqpoint{1.780464in}{2.644559in}}%
\pgfpathlineto{\pgfqpoint{1.780470in}{2.749643in}}%
\pgfpathlineto{\pgfqpoint{1.780571in}{2.642779in}}%
\pgfpathlineto{\pgfqpoint{1.781211in}{2.707992in}}%
\pgfpathlineto{\pgfqpoint{1.781260in}{2.685462in}}%
\pgfpathlineto{\pgfqpoint{1.781804in}{2.745241in}}%
\pgfpathlineto{\pgfqpoint{1.781854in}{2.722776in}}%
\pgfpathlineto{\pgfqpoint{1.782393in}{2.782198in}}%
\pgfpathlineto{\pgfqpoint{1.782493in}{2.696387in}}%
\pgfpathlineto{\pgfqpoint{1.782507in}{2.748902in}}%
\pgfpathlineto{\pgfqpoint{1.783289in}{2.644559in}}%
\pgfpathlineto{\pgfqpoint{1.785410in}{2.644559in}}%
\pgfpathlineto{\pgfqpoint{1.785415in}{2.748937in}}%
\pgfpathlineto{\pgfqpoint{1.785493in}{2.634569in}}%
\pgfpathlineto{\pgfqpoint{1.786178in}{2.675791in}}%
\pgfpathlineto{\pgfqpoint{1.787019in}{2.768850in}}%
\pgfpathlineto{\pgfqpoint{1.787384in}{2.791682in}}%
\pgfpathlineto{\pgfqpoint{1.787956in}{2.644559in}}%
\pgfpathlineto{\pgfqpoint{1.790355in}{2.644559in}}%
\pgfpathlineto{\pgfqpoint{1.790360in}{2.749342in}}%
\pgfpathlineto{\pgfqpoint{1.790423in}{2.641116in}}%
\pgfpathlineto{\pgfqpoint{1.791107in}{2.680412in}}%
\pgfpathlineto{\pgfqpoint{1.791948in}{2.762931in}}%
\pgfpathlineto{\pgfqpoint{1.792329in}{2.786851in}}%
\pgfpathlineto{\pgfqpoint{1.792940in}{2.644559in}}%
\pgfpathlineto{\pgfqpoint{1.795301in}{2.644559in}}%
\pgfpathlineto{\pgfqpoint{1.795306in}{2.748787in}}%
\pgfpathlineto{\pgfqpoint{1.795380in}{2.637429in}}%
\pgfpathlineto{\pgfqpoint{1.796065in}{2.678299in}}%
\pgfpathlineto{\pgfqpoint{1.796906in}{2.765741in}}%
\pgfpathlineto{\pgfqpoint{1.797274in}{2.788831in}}%
\pgfpathlineto{\pgfqpoint{1.797865in}{2.644559in}}%
\pgfpathlineto{\pgfqpoint{1.800246in}{2.644559in}}%
\pgfpathlineto{\pgfqpoint{1.800252in}{2.749504in}}%
\pgfpathlineto{\pgfqpoint{1.800986in}{2.697364in}}%
\pgfpathlineto{\pgfqpoint{1.801035in}{2.694680in}}%
\pgfpathlineto{\pgfqpoint{1.801282in}{2.716009in}}%
\pgfpathlineto{\pgfqpoint{1.801332in}{2.713318in}}%
\pgfpathlineto{\pgfqpoint{1.802173in}{2.771938in}}%
\pgfpathlineto{\pgfqpoint{1.802245in}{2.757516in}}%
\pgfpathlineto{\pgfqpoint{1.803044in}{2.644559in}}%
\pgfpathlineto{\pgfqpoint{1.805192in}{2.644559in}}%
\pgfpathlineto{\pgfqpoint{1.805197in}{2.749070in}}%
\pgfpathlineto{\pgfqpoint{1.805276in}{2.633082in}}%
\pgfpathlineto{\pgfqpoint{1.805962in}{2.674450in}}%
\pgfpathlineto{\pgfqpoint{1.806803in}{2.770651in}}%
\pgfpathlineto{\pgfqpoint{1.807165in}{2.793370in}}%
\pgfpathlineto{\pgfqpoint{1.807729in}{2.644559in}}%
\pgfpathlineto{\pgfqpoint{1.810137in}{2.644559in}}%
\pgfpathlineto{\pgfqpoint{1.810142in}{2.749253in}}%
\pgfpathlineto{\pgfqpoint{1.810220in}{2.642849in}}%
\pgfpathlineto{\pgfqpoint{1.810862in}{2.705200in}}%
\pgfpathlineto{\pgfqpoint{1.810911in}{2.684174in}}%
\pgfpathlineto{\pgfqpoint{1.811455in}{2.742455in}}%
\pgfpathlineto{\pgfqpoint{1.811505in}{2.721482in}}%
\pgfpathlineto{\pgfqpoint{1.812111in}{2.783608in}}%
\pgfpathlineto{\pgfqpoint{1.812165in}{2.720009in}}%
\pgfpathlineto{\pgfqpoint{1.812178in}{2.724880in}}%
\pgfpathlineto{\pgfqpoint{1.812963in}{2.644559in}}%
\pgfpathlineto{\pgfqpoint{1.815082in}{2.644559in}}%
\pgfpathlineto{\pgfqpoint{1.815088in}{2.749537in}}%
\pgfpathlineto{\pgfqpoint{1.815203in}{2.642905in}}%
\pgfpathlineto{\pgfqpoint{1.815846in}{2.709117in}}%
\pgfpathlineto{\pgfqpoint{1.815895in}{2.686091in}}%
\pgfpathlineto{\pgfqpoint{1.816439in}{2.746364in}}%
\pgfpathlineto{\pgfqpoint{1.816488in}{2.723407in}}%
\pgfpathlineto{\pgfqpoint{1.817020in}{2.782802in}}%
\pgfpathlineto{\pgfqpoint{1.817111in}{2.697158in}}%
\pgfpathlineto{\pgfqpoint{1.817126in}{2.748584in}}%
\pgfpathlineto{\pgfqpoint{1.817905in}{2.644559in}}%
\pgfpathlineto{\pgfqpoint{1.820028in}{2.644559in}}%
\pgfpathlineto{\pgfqpoint{1.820033in}{2.749076in}}%
\pgfpathlineto{\pgfqpoint{1.820090in}{2.641064in}}%
\pgfpathlineto{\pgfqpoint{1.820769in}{2.679158in}}%
\pgfpathlineto{\pgfqpoint{1.821610in}{2.761469in}}%
\pgfpathlineto{\pgfqpoint{1.822002in}{2.786075in}}%
\pgfpathlineto{\pgfqpoint{1.822599in}{2.644559in}}%
\pgfpathlineto{\pgfqpoint{1.824973in}{2.644559in}}%
\pgfpathlineto{\pgfqpoint{1.824979in}{2.749377in}}%
\pgfpathlineto{\pgfqpoint{1.825726in}{2.691744in}}%
\pgfpathlineto{\pgfqpoint{1.826864in}{2.770082in}}%
\pgfpathlineto{\pgfqpoint{1.827002in}{2.716479in}}%
\pgfpathlineto{\pgfqpoint{1.826947in}{2.775318in}}%
\pgfpathlineto{\pgfqpoint{1.827016in}{2.729054in}}%
\pgfpathlineto{\pgfqpoint{1.827801in}{2.644559in}}%
\pgfpathlineto{\pgfqpoint{1.829919in}{2.644559in}}%
\pgfpathlineto{\pgfqpoint{1.829924in}{2.749068in}}%
\pgfpathlineto{\pgfqpoint{1.830004in}{2.633060in}}%
\pgfpathlineto{\pgfqpoint{1.830689in}{2.674429in}}%
\pgfpathlineto{\pgfqpoint{1.831530in}{2.770661in}}%
\pgfpathlineto{\pgfqpoint{1.831893in}{2.793378in}}%
\pgfpathlineto{\pgfqpoint{1.832454in}{2.644559in}}%
\pgfpathlineto{\pgfqpoint{1.834864in}{2.644559in}}%
\pgfpathlineto{\pgfqpoint{1.834870in}{2.748991in}}%
\pgfpathlineto{\pgfqpoint{1.834949in}{2.632787in}}%
\pgfpathlineto{\pgfqpoint{1.835635in}{2.674176in}}%
\pgfpathlineto{\pgfqpoint{1.836476in}{2.770626in}}%
\pgfpathlineto{\pgfqpoint{1.836838in}{2.793323in}}%
\pgfpathlineto{\pgfqpoint{1.837400in}{2.644559in}}%
\pgfpathlineto{\pgfqpoint{1.839810in}{2.644559in}}%
\pgfpathlineto{\pgfqpoint{1.839815in}{2.749307in}}%
\pgfpathlineto{\pgfqpoint{1.839904in}{2.642715in}}%
\pgfpathlineto{\pgfqpoint{1.840546in}{2.706120in}}%
\pgfpathlineto{\pgfqpoint{1.840595in}{2.684875in}}%
\pgfpathlineto{\pgfqpoint{1.841139in}{2.743375in}}%
\pgfpathlineto{\pgfqpoint{1.841189in}{2.722184in}}%
\pgfpathlineto{\pgfqpoint{1.841784in}{2.783796in}}%
\pgfpathlineto{\pgfqpoint{1.841837in}{2.721046in}}%
\pgfpathlineto{\pgfqpoint{1.841851in}{2.723905in}}%
\pgfpathlineto{\pgfqpoint{1.842635in}{2.644559in}}%
\pgfpathlineto{\pgfqpoint{1.844755in}{2.644559in}}%
\pgfpathlineto{\pgfqpoint{1.844760in}{2.749193in}}%
\pgfpathlineto{\pgfqpoint{1.844841in}{2.632649in}}%
\pgfpathlineto{\pgfqpoint{1.845526in}{2.674064in}}%
\pgfpathlineto{\pgfqpoint{1.846367in}{2.771496in}}%
\pgfpathlineto{\pgfqpoint{1.846729in}{2.794179in}}%
\pgfpathlineto{\pgfqpoint{1.847276in}{2.644559in}}%
\pgfpathlineto{\pgfqpoint{1.849701in}{2.644559in}}%
\pgfpathlineto{\pgfqpoint{1.849706in}{2.749030in}}%
\pgfpathlineto{\pgfqpoint{1.849788in}{2.631209in}}%
\pgfpathlineto{\pgfqpoint{1.850473in}{2.672761in}}%
\pgfpathlineto{\pgfqpoint{1.851314in}{2.772169in}}%
\pgfpathlineto{\pgfqpoint{1.851674in}{2.794725in}}%
\pgfpathlineto{\pgfqpoint{1.852230in}{2.644559in}}%
\pgfpathlineto{\pgfqpoint{1.854646in}{2.644559in}}%
\pgfpathlineto{\pgfqpoint{1.854651in}{2.749076in}}%
\pgfpathlineto{\pgfqpoint{1.854711in}{2.642620in}}%
\pgfpathlineto{\pgfqpoint{1.855391in}{2.681372in}}%
\pgfpathlineto{\pgfqpoint{1.856232in}{2.761091in}}%
\pgfpathlineto{\pgfqpoint{1.856620in}{2.785406in}}%
\pgfpathlineto{\pgfqpoint{1.857247in}{2.644559in}}%
\pgfpathlineto{\pgfqpoint{1.859592in}{2.644559in}}%
\pgfpathlineto{\pgfqpoint{1.859597in}{2.748773in}}%
\pgfpathlineto{\pgfqpoint{1.859671in}{2.637208in}}%
\pgfpathlineto{\pgfqpoint{1.860356in}{2.678100in}}%
\pgfpathlineto{\pgfqpoint{1.861197in}{2.765867in}}%
\pgfpathlineto{\pgfqpoint{1.861565in}{2.788941in}}%
\pgfpathlineto{\pgfqpoint{1.862156in}{2.644559in}}%
\pgfpathlineto{\pgfqpoint{1.864537in}{2.644559in}}%
\pgfpathlineto{\pgfqpoint{1.864542in}{2.748938in}}%
\pgfpathlineto{\pgfqpoint{1.864616in}{2.637585in}}%
\pgfpathlineto{\pgfqpoint{1.865302in}{2.678457in}}%
\pgfpathlineto{\pgfqpoint{1.866142in}{2.766281in}}%
\pgfpathlineto{\pgfqpoint{1.866511in}{2.789376in}}%
\pgfpathlineto{\pgfqpoint{1.867100in}{2.644559in}}%
\pgfpathlineto{\pgfqpoint{1.869482in}{2.644559in}}%
\pgfpathlineto{\pgfqpoint{1.869488in}{2.748635in}}%
\pgfpathlineto{\pgfqpoint{1.869558in}{2.640494in}}%
\pgfpathlineto{\pgfqpoint{1.870243in}{2.681004in}}%
\pgfpathlineto{\pgfqpoint{1.871084in}{2.762758in}}%
\pgfpathlineto{\pgfqpoint{1.871456in}{2.786081in}}%
\pgfpathlineto{\pgfqpoint{1.872078in}{2.644559in}}%
\pgfpathlineto{\pgfqpoint{1.874428in}{2.644559in}}%
\pgfpathlineto{\pgfqpoint{1.874433in}{2.748788in}}%
\pgfpathlineto{\pgfqpoint{1.874511in}{2.635030in}}%
\pgfpathlineto{\pgfqpoint{1.875196in}{2.676198in}}%
\pgfpathlineto{\pgfqpoint{1.876037in}{2.767905in}}%
\pgfpathlineto{\pgfqpoint{1.876402in}{2.790774in}}%
\pgfpathlineto{\pgfqpoint{1.876977in}{2.644559in}}%
\pgfpathlineto{\pgfqpoint{1.879373in}{2.644559in}}%
\pgfpathlineto{\pgfqpoint{1.879379in}{2.749159in}}%
\pgfpathlineto{\pgfqpoint{1.879459in}{2.632728in}}%
\pgfpathlineto{\pgfqpoint{1.880144in}{2.674145in}}%
\pgfpathlineto{\pgfqpoint{1.880985in}{2.771243in}}%
\pgfpathlineto{\pgfqpoint{1.881347in}{2.793922in}}%
\pgfpathlineto{\pgfqpoint{1.881906in}{2.644559in}}%
\pgfpathlineto{\pgfqpoint{1.884319in}{2.644559in}}%
\pgfpathlineto{\pgfqpoint{1.884324in}{2.749041in}}%
\pgfpathlineto{\pgfqpoint{1.884456in}{2.641768in}}%
\pgfpathlineto{\pgfqpoint{1.885049in}{2.678915in}}%
\pgfpathlineto{\pgfqpoint{1.885890in}{2.760225in}}%
\pgfpathlineto{\pgfqpoint{1.886264in}{2.783724in}}%
\pgfpathlineto{\pgfqpoint{1.886883in}{2.644559in}}%
\pgfpathlineto{\pgfqpoint{1.889264in}{2.644559in}}%
\pgfpathlineto{\pgfqpoint{1.889270in}{2.749443in}}%
\pgfpathlineto{\pgfqpoint{1.889388in}{2.642071in}}%
\pgfpathlineto{\pgfqpoint{1.890031in}{2.709968in}}%
\pgfpathlineto{\pgfqpoint{1.890080in}{2.685303in}}%
\pgfpathlineto{\pgfqpoint{1.890624in}{2.747213in}}%
\pgfpathlineto{\pgfqpoint{1.890674in}{2.722621in}}%
\pgfpathlineto{\pgfqpoint{1.891203in}{2.783541in}}%
\pgfpathlineto{\pgfqpoint{1.891293in}{2.696942in}}%
\pgfpathlineto{\pgfqpoint{1.891308in}{2.748450in}}%
\pgfpathlineto{\pgfqpoint{1.892089in}{2.644559in}}%
\pgfpathlineto{\pgfqpoint{1.894210in}{2.644559in}}%
\pgfpathlineto{\pgfqpoint{1.894215in}{2.748946in}}%
\pgfpathlineto{\pgfqpoint{1.894273in}{2.640273in}}%
\pgfpathlineto{\pgfqpoint{1.894952in}{2.678481in}}%
\pgfpathlineto{\pgfqpoint{1.895792in}{2.761861in}}%
\pgfpathlineto{\pgfqpoint{1.896184in}{2.786395in}}%
\pgfpathlineto{\pgfqpoint{1.896771in}{2.644559in}}%
\pgfpathlineto{\pgfqpoint{1.899155in}{2.644559in}}%
\pgfpathlineto{\pgfqpoint{1.899161in}{2.748990in}}%
\pgfpathlineto{\pgfqpoint{1.899241in}{2.632447in}}%
\pgfpathlineto{\pgfqpoint{1.899926in}{2.673880in}}%
\pgfpathlineto{\pgfqpoint{1.900767in}{2.770884in}}%
\pgfpathlineto{\pgfqpoint{1.901129in}{2.793542in}}%
\pgfpathlineto{\pgfqpoint{1.901687in}{2.644559in}}%
\pgfpathlineto{\pgfqpoint{1.904101in}{2.644559in}}%
\pgfpathlineto{\pgfqpoint{1.904106in}{2.748826in}}%
\pgfpathlineto{\pgfqpoint{1.904184in}{2.634755in}}%
\pgfpathlineto{\pgfqpoint{1.904869in}{2.675951in}}%
\pgfpathlineto{\pgfqpoint{1.905710in}{2.768273in}}%
\pgfpathlineto{\pgfqpoint{1.906074in}{2.791120in}}%
\pgfpathlineto{\pgfqpoint{1.906650in}{2.644559in}}%
\pgfpathlineto{\pgfqpoint{1.909046in}{2.644559in}}%
\pgfpathlineto{\pgfqpoint{1.909051in}{2.749114in}}%
\pgfpathlineto{\pgfqpoint{1.909115in}{2.641038in}}%
\pgfpathlineto{\pgfqpoint{1.909800in}{2.680589in}}%
\pgfpathlineto{\pgfqpoint{1.910641in}{2.762353in}}%
\pgfpathlineto{\pgfqpoint{1.911020in}{2.786127in}}%
\pgfpathlineto{\pgfqpoint{1.911637in}{2.644559in}}%
\pgfpathlineto{\pgfqpoint{1.913992in}{2.644559in}}%
\pgfpathlineto{\pgfqpoint{1.913997in}{2.748832in}}%
\pgfpathlineto{\pgfqpoint{1.914079in}{2.631148in}}%
\pgfpathlineto{\pgfqpoint{1.914764in}{2.672696in}}%
\pgfpathlineto{\pgfqpoint{1.915605in}{2.771584in}}%
\pgfpathlineto{\pgfqpoint{1.915965in}{2.794136in}}%
\pgfpathlineto{\pgfqpoint{1.916520in}{2.644559in}}%
\pgfpathlineto{\pgfqpoint{1.918937in}{2.644559in}}%
\pgfpathlineto{\pgfqpoint{1.918942in}{2.749271in}}%
\pgfpathlineto{\pgfqpoint{1.919009in}{2.640786in}}%
\pgfpathlineto{\pgfqpoint{1.919695in}{2.680808in}}%
\pgfpathlineto{\pgfqpoint{1.920536in}{2.763289in}}%
\pgfpathlineto{\pgfqpoint{1.920911in}{2.786794in}}%
\pgfpathlineto{\pgfqpoint{1.921524in}{2.644559in}}%
\pgfpathlineto{\pgfqpoint{1.923882in}{2.644559in}}%
\pgfpathlineto{\pgfqpoint{1.923888in}{2.748589in}}%
\pgfpathlineto{\pgfqpoint{1.923967in}{2.633484in}}%
\pgfpathlineto{\pgfqpoint{1.924652in}{2.674786in}}%
\pgfpathlineto{\pgfqpoint{1.925493in}{2.768641in}}%
\pgfpathlineto{\pgfqpoint{1.925856in}{2.791395in}}%
\pgfpathlineto{\pgfqpoint{1.926424in}{2.644559in}}%
\pgfpathlineto{\pgfqpoint{1.928828in}{2.644559in}}%
\pgfpathlineto{\pgfqpoint{1.928832in}{2.756004in}}%
\pgfpathlineto{\pgfqpoint{1.929581in}{2.685983in}}%
\pgfpathlineto{\pgfqpoint{1.930520in}{2.763137in}}%
\pgfpathlineto{\pgfqpoint{1.930802in}{2.780784in}}%
\pgfpathlineto{\pgfqpoint{1.931684in}{2.644559in}}%
\pgfpathlineto{\pgfqpoint{1.933773in}{2.644559in}}%
\pgfpathlineto{\pgfqpoint{1.933779in}{2.746994in}}%
\pgfpathlineto{\pgfqpoint{1.934547in}{2.706482in}}%
\pgfpathlineto{\pgfqpoint{1.934597in}{2.688994in}}%
\pgfpathlineto{\pgfqpoint{1.935141in}{2.743715in}}%
\pgfpathlineto{\pgfqpoint{1.935190in}{2.726323in}}%
\pgfpathlineto{\pgfqpoint{1.935716in}{2.779804in}}%
\pgfpathlineto{\pgfqpoint{1.935802in}{2.717930in}}%
\pgfpathlineto{\pgfqpoint{1.935815in}{2.727080in}}%
\pgfpathlineto{\pgfqpoint{1.936602in}{2.644559in}}%
\pgfpathlineto{\pgfqpoint{1.938719in}{2.644559in}}%
\pgfpathlineto{\pgfqpoint{1.938724in}{2.748948in}}%
\pgfpathlineto{\pgfqpoint{1.938803in}{2.633788in}}%
\pgfpathlineto{\pgfqpoint{1.939488in}{2.675083in}}%
\pgfpathlineto{\pgfqpoint{1.940329in}{2.769569in}}%
\pgfpathlineto{\pgfqpoint{1.940693in}{2.792343in}}%
\pgfpathlineto{\pgfqpoint{1.941259in}{2.644559in}}%
\pgfpathlineto{\pgfqpoint{1.943664in}{2.644559in}}%
\pgfpathlineto{\pgfqpoint{1.943670in}{2.749043in}}%
\pgfpathlineto{\pgfqpoint{1.943789in}{2.641192in}}%
\pgfpathlineto{\pgfqpoint{1.944432in}{2.709718in}}%
\pgfpathlineto{\pgfqpoint{1.944481in}{2.684435in}}%
\pgfpathlineto{\pgfqpoint{1.945025in}{2.746963in}}%
\pgfpathlineto{\pgfqpoint{1.945075in}{2.721753in}}%
\pgfpathlineto{\pgfqpoint{1.945604in}{2.783260in}}%
\pgfpathlineto{\pgfqpoint{1.945721in}{2.644559in}}%
\pgfpathlineto{\pgfqpoint{1.948610in}{2.644559in}}%
\pgfpathlineto{\pgfqpoint{1.948615in}{2.749114in}}%
\pgfpathlineto{\pgfqpoint{1.949359in}{2.691466in}}%
\pgfpathlineto{\pgfqpoint{1.950496in}{2.769027in}}%
\pgfpathlineto{\pgfqpoint{1.950639in}{2.714772in}}%
\pgfpathlineto{\pgfqpoint{1.950584in}{2.774512in}}%
\pgfpathlineto{\pgfqpoint{1.950653in}{2.730071in}}%
\pgfpathlineto{\pgfqpoint{1.951436in}{2.644559in}}%
\pgfpathlineto{\pgfqpoint{1.953555in}{2.644559in}}%
\pgfpathlineto{\pgfqpoint{1.953560in}{2.748804in}}%
\pgfpathlineto{\pgfqpoint{1.953634in}{2.638213in}}%
\pgfpathlineto{\pgfqpoint{1.954319in}{2.679003in}}%
\pgfpathlineto{\pgfqpoint{1.955160in}{2.765204in}}%
\pgfpathlineto{\pgfqpoint{1.955529in}{2.788351in}}%
\pgfpathlineto{\pgfqpoint{1.956127in}{2.644559in}}%
\pgfpathlineto{\pgfqpoint{1.958501in}{2.644559in}}%
\pgfpathlineto{\pgfqpoint{1.958506in}{2.748275in}}%
\pgfpathlineto{\pgfqpoint{1.959255in}{2.685309in}}%
\pgfpathlineto{\pgfqpoint{1.960194in}{2.763309in}}%
\pgfpathlineto{\pgfqpoint{1.960474in}{2.780882in}}%
\pgfpathlineto{\pgfqpoint{1.961337in}{2.644559in}}%
\pgfpathlineto{\pgfqpoint{1.963446in}{2.644559in}}%
\pgfpathlineto{\pgfqpoint{1.963451in}{2.748884in}}%
\pgfpathlineto{\pgfqpoint{1.963526in}{2.637195in}}%
\pgfpathlineto{\pgfqpoint{1.964211in}{2.678101in}}%
\pgfpathlineto{\pgfqpoint{1.965052in}{2.766345in}}%
\pgfpathlineto{\pgfqpoint{1.965420in}{2.789414in}}%
\pgfpathlineto{\pgfqpoint{1.966006in}{2.644559in}}%
\pgfpathlineto{\pgfqpoint{1.968392in}{2.644559in}}%
\pgfpathlineto{\pgfqpoint{1.968397in}{2.748704in}}%
\pgfpathlineto{\pgfqpoint{1.968476in}{2.633399in}}%
\pgfpathlineto{\pgfqpoint{1.969161in}{2.674715in}}%
\pgfpathlineto{\pgfqpoint{1.970002in}{2.769089in}}%
\pgfpathlineto{\pgfqpoint{1.970365in}{2.791836in}}%
\pgfpathlineto{\pgfqpoint{1.970928in}{2.644559in}}%
\pgfpathlineto{\pgfqpoint{1.973337in}{2.644559in}}%
\pgfpathlineto{\pgfqpoint{1.973342in}{2.746970in}}%
\pgfpathlineto{\pgfqpoint{1.973459in}{2.625620in}}%
\pgfpathlineto{\pgfqpoint{1.974102in}{2.725072in}}%
\pgfpathlineto{\pgfqpoint{1.974151in}{2.668869in}}%
\pgfpathlineto{\pgfqpoint{1.974794in}{2.768455in}}%
\pgfpathlineto{\pgfqpoint{1.974844in}{2.712475in}}%
\pgfpathlineto{\pgfqpoint{1.975275in}{2.798582in}}%
\pgfpathlineto{\pgfqpoint{1.975402in}{2.644559in}}%
\pgfpathlineto{\pgfqpoint{1.975543in}{2.644559in}}%
\pgfpathlineto{\pgfqpoint{1.978282in}{2.644559in}}%
\pgfpathlineto{\pgfqpoint{1.978288in}{2.748723in}}%
\pgfpathlineto{\pgfqpoint{1.978363in}{2.636668in}}%
\pgfpathlineto{\pgfqpoint{1.979048in}{2.677599in}}%
\pgfpathlineto{\pgfqpoint{1.979888in}{2.766182in}}%
\pgfpathlineto{\pgfqpoint{1.980256in}{2.789227in}}%
\pgfpathlineto{\pgfqpoint{1.980841in}{2.644559in}}%
\pgfpathlineto{\pgfqpoint{1.983228in}{2.644559in}}%
\pgfpathlineto{\pgfqpoint{1.983233in}{2.748811in}}%
\pgfpathlineto{\pgfqpoint{1.983308in}{2.637070in}}%
\pgfpathlineto{\pgfqpoint{1.983993in}{2.677980in}}%
\pgfpathlineto{\pgfqpoint{1.984834in}{2.766136in}}%
\pgfpathlineto{\pgfqpoint{1.985202in}{2.789199in}}%
\pgfpathlineto{\pgfqpoint{1.985792in}{2.644559in}}%
\pgfpathlineto{\pgfqpoint{1.988173in}{2.644559in}}%
\pgfpathlineto{\pgfqpoint{1.988179in}{2.748831in}}%
\pgfpathlineto{\pgfqpoint{1.988258in}{2.633466in}}%
\pgfpathlineto{\pgfqpoint{1.988943in}{2.674783in}}%
\pgfpathlineto{\pgfqpoint{1.989784in}{2.769450in}}%
\pgfpathlineto{\pgfqpoint{1.990147in}{2.792201in}}%
\pgfpathlineto{\pgfqpoint{1.990714in}{2.644559in}}%
\pgfpathlineto{\pgfqpoint{1.993119in}{2.644559in}}%
\pgfpathlineto{\pgfqpoint{1.993124in}{2.748719in}}%
\pgfpathlineto{\pgfqpoint{1.993198in}{2.637752in}}%
\pgfpathlineto{\pgfqpoint{1.993883in}{2.678581in}}%
\pgfpathlineto{\pgfqpoint{1.994724in}{2.765216in}}%
\pgfpathlineto{\pgfqpoint{1.995093in}{2.788332in}}%
\pgfpathlineto{\pgfqpoint{1.995689in}{2.644559in}}%
\pgfpathlineto{\pgfqpoint{1.998064in}{2.644559in}}%
\pgfpathlineto{\pgfqpoint{1.998070in}{2.748687in}}%
\pgfpathlineto{\pgfqpoint{1.998141in}{2.639629in}}%
\pgfpathlineto{\pgfqpoint{1.998826in}{2.680248in}}%
\pgfpathlineto{\pgfqpoint{1.999667in}{2.763638in}}%
\pgfpathlineto{\pgfqpoint{2.000038in}{2.786894in}}%
\pgfpathlineto{\pgfqpoint{2.000652in}{2.644559in}}%
\pgfpathlineto{\pgfqpoint{2.003010in}{2.644559in}}%
\pgfpathlineto{\pgfqpoint{2.003015in}{2.748948in}}%
\pgfpathlineto{\pgfqpoint{2.003073in}{2.641081in}}%
\pgfpathlineto{\pgfqpoint{2.003753in}{2.679415in}}%
\pgfpathlineto{\pgfqpoint{2.004594in}{2.761138in}}%
\pgfpathlineto{\pgfqpoint{2.004984in}{2.785607in}}%
\pgfpathlineto{\pgfqpoint{2.005587in}{2.644559in}}%
\pgfpathlineto{\pgfqpoint{2.007955in}{2.644559in}}%
\pgfpathlineto{\pgfqpoint{2.007961in}{2.749283in}}%
\pgfpathlineto{\pgfqpoint{2.008069in}{2.641645in}}%
\pgfpathlineto{\pgfqpoint{2.008710in}{2.708684in}}%
\pgfpathlineto{\pgfqpoint{2.008760in}{2.684612in}}%
\pgfpathlineto{\pgfqpoint{2.009304in}{2.745931in}}%
\pgfpathlineto{\pgfqpoint{2.009353in}{2.721928in}}%
\pgfpathlineto{\pgfqpoint{2.009888in}{2.782622in}}%
\pgfpathlineto{\pgfqpoint{2.009984in}{2.696746in}}%
\pgfpathlineto{\pgfqpoint{2.009999in}{2.748112in}}%
\pgfpathlineto{\pgfqpoint{2.010781in}{2.644559in}}%
\pgfpathlineto{\pgfqpoint{2.012901in}{2.644559in}}%
\pgfpathlineto{\pgfqpoint{2.012906in}{2.748621in}}%
\pgfpathlineto{\pgfqpoint{2.012979in}{2.638396in}}%
\pgfpathlineto{\pgfqpoint{2.013664in}{2.679144in}}%
\pgfpathlineto{\pgfqpoint{2.014505in}{2.764317in}}%
\pgfpathlineto{\pgfqpoint{2.014874in}{2.787484in}}%
\pgfpathlineto{\pgfqpoint{2.015479in}{2.644559in}}%
\pgfpathlineto{\pgfqpoint{2.017846in}{2.644559in}}%
\pgfpathlineto{\pgfqpoint{2.017851in}{2.748833in}}%
\pgfpathlineto{\pgfqpoint{2.017932in}{2.632574in}}%
\pgfpathlineto{\pgfqpoint{2.018617in}{2.673987in}}%
\pgfpathlineto{\pgfqpoint{2.019458in}{2.770235in}}%
\pgfpathlineto{\pgfqpoint{2.019820in}{2.792905in}}%
\pgfpathlineto{\pgfqpoint{2.020379in}{2.644559in}}%
\pgfpathlineto{\pgfqpoint{2.022792in}{2.644559in}}%
\pgfpathlineto{\pgfqpoint{2.022797in}{2.748830in}}%
\pgfpathlineto{\pgfqpoint{2.022874in}{2.634811in}}%
\pgfpathlineto{\pgfqpoint{2.023560in}{2.676003in}}%
\pgfpathlineto{\pgfqpoint{2.024401in}{2.768244in}}%
\pgfpathlineto{\pgfqpoint{2.024765in}{2.791096in}}%
\pgfpathlineto{\pgfqpoint{2.025338in}{2.644559in}}%
\pgfpathlineto{\pgfqpoint{2.027737in}{2.644559in}}%
\pgfpathlineto{\pgfqpoint{2.027742in}{2.748841in}}%
\pgfpathlineto{\pgfqpoint{2.027821in}{2.633859in}}%
\pgfpathlineto{\pgfqpoint{2.028506in}{2.675141in}}%
\pgfpathlineto{\pgfqpoint{2.029347in}{2.769125in}}%
\pgfpathlineto{\pgfqpoint{2.029711in}{2.791905in}}%
\pgfpathlineto{\pgfqpoint{2.030277in}{2.644559in}}%
\pgfpathlineto{\pgfqpoint{2.032682in}{2.644559in}}%
\pgfpathlineto{\pgfqpoint{2.032688in}{2.748817in}}%
\pgfpathlineto{\pgfqpoint{2.032771in}{2.630082in}}%
\pgfpathlineto{\pgfqpoint{2.033407in}{2.716530in}}%
\pgfpathlineto{\pgfqpoint{2.033457in}{2.671711in}}%
\pgfpathlineto{\pgfqpoint{2.034099in}{2.759907in}}%
\pgfpathlineto{\pgfqpoint{2.034149in}{2.715324in}}%
\pgfpathlineto{\pgfqpoint{2.034656in}{2.794774in}}%
\pgfpathlineto{\pgfqpoint{2.034741in}{2.644559in}}%
\pgfpathlineto{\pgfqpoint{2.034818in}{2.644559in}}%
\pgfpathlineto{\pgfqpoint{2.037628in}{2.644559in}}%
\pgfpathlineto{\pgfqpoint{2.037633in}{2.749177in}}%
\pgfpathlineto{\pgfqpoint{2.037714in}{2.632015in}}%
\pgfpathlineto{\pgfqpoint{2.038400in}{2.673500in}}%
\pgfpathlineto{\pgfqpoint{2.039240in}{2.771942in}}%
\pgfpathlineto{\pgfqpoint{2.039602in}{2.794562in}}%
\pgfpathlineto{\pgfqpoint{2.040160in}{2.644559in}}%
\pgfpathlineto{\pgfqpoint{2.042573in}{2.644559in}}%
\pgfpathlineto{\pgfqpoint{2.042579in}{2.748090in}}%
\pgfpathlineto{\pgfqpoint{2.042638in}{2.642028in}}%
\pgfpathlineto{\pgfqpoint{2.043319in}{2.680772in}}%
\pgfpathlineto{\pgfqpoint{2.044159in}{2.760242in}}%
\pgfpathlineto{\pgfqpoint{2.044547in}{2.784570in}}%
\pgfpathlineto{\pgfqpoint{2.045174in}{2.644559in}}%
\pgfpathlineto{\pgfqpoint{2.047519in}{2.644559in}}%
\pgfpathlineto{\pgfqpoint{2.047524in}{2.748570in}}%
\pgfpathlineto{\pgfqpoint{2.047595in}{2.640236in}}%
\pgfpathlineto{\pgfqpoint{2.048280in}{2.680766in}}%
\pgfpathlineto{\pgfqpoint{2.049121in}{2.762672in}}%
\pgfpathlineto{\pgfqpoint{2.049493in}{2.785980in}}%
\pgfpathlineto{\pgfqpoint{2.050115in}{2.644559in}}%
\pgfpathlineto{\pgfqpoint{2.052464in}{2.644559in}}%
\pgfpathlineto{\pgfqpoint{2.052469in}{2.747682in}}%
\pgfpathlineto{\pgfqpoint{2.053232in}{2.705027in}}%
\pgfpathlineto{\pgfqpoint{2.053281in}{2.689491in}}%
\pgfpathlineto{\pgfqpoint{2.053825in}{2.742265in}}%
\pgfpathlineto{\pgfqpoint{2.053875in}{2.726815in}}%
\pgfpathlineto{\pgfqpoint{2.054404in}{2.778558in}}%
\pgfpathlineto{\pgfqpoint{2.054492in}{2.720004in}}%
\pgfpathlineto{\pgfqpoint{2.054506in}{2.724604in}}%
\pgfpathlineto{\pgfqpoint{2.055292in}{2.644559in}}%
\pgfpathlineto{\pgfqpoint{2.057410in}{2.644559in}}%
\pgfpathlineto{\pgfqpoint{2.057415in}{2.748309in}}%
\pgfpathlineto{\pgfqpoint{2.057478in}{2.641500in}}%
\pgfpathlineto{\pgfqpoint{2.058161in}{2.680944in}}%
\pgfpathlineto{\pgfqpoint{2.059001in}{2.761414in}}%
\pgfpathlineto{\pgfqpoint{2.059384in}{2.785371in}}%
\pgfpathlineto{\pgfqpoint{2.059999in}{2.644559in}}%
\pgfpathlineto{\pgfqpoint{2.062355in}{2.644559in}}%
\pgfpathlineto{\pgfqpoint{2.062361in}{2.748417in}}%
\pgfpathlineto{\pgfqpoint{2.062433in}{2.639028in}}%
\pgfpathlineto{\pgfqpoint{2.063118in}{2.679680in}}%
\pgfpathlineto{\pgfqpoint{2.063958in}{2.763042in}}%
\pgfpathlineto{\pgfqpoint{2.064329in}{2.786265in}}%
\pgfpathlineto{\pgfqpoint{2.064944in}{2.644559in}}%
\pgfpathlineto{\pgfqpoint{2.067301in}{2.644559in}}%
\pgfpathlineto{\pgfqpoint{2.067306in}{2.749266in}}%
\pgfpathlineto{\pgfqpoint{2.067424in}{2.641963in}}%
\pgfpathlineto{\pgfqpoint{2.068067in}{2.709485in}}%
\pgfpathlineto{\pgfqpoint{2.068116in}{2.685190in}}%
\pgfpathlineto{\pgfqpoint{2.068660in}{2.746731in}}%
\pgfpathlineto{\pgfqpoint{2.068710in}{2.722507in}}%
\pgfpathlineto{\pgfqpoint{2.069239in}{2.783069in}}%
\pgfpathlineto{\pgfqpoint{2.069330in}{2.696701in}}%
\pgfpathlineto{\pgfqpoint{2.069345in}{2.748139in}}%
\pgfpathlineto{\pgfqpoint{2.070127in}{2.644559in}}%
\pgfpathlineto{\pgfqpoint{2.072246in}{2.644559in}}%
\pgfpathlineto{\pgfqpoint{2.072252in}{2.748883in}}%
\pgfpathlineto{\pgfqpoint{2.072312in}{2.640603in}}%
\pgfpathlineto{\pgfqpoint{2.072994in}{2.679473in}}%
\pgfpathlineto{\pgfqpoint{2.073835in}{2.761527in}}%
\pgfpathlineto{\pgfqpoint{2.074220in}{2.785682in}}%
\pgfpathlineto{\pgfqpoint{2.074828in}{2.644559in}}%
\pgfpathlineto{\pgfqpoint{2.077192in}{2.644559in}}%
\pgfpathlineto{\pgfqpoint{2.077197in}{2.748482in}}%
\pgfpathlineto{\pgfqpoint{2.077268in}{2.639834in}}%
\pgfpathlineto{\pgfqpoint{2.077953in}{2.680399in}}%
\pgfpathlineto{\pgfqpoint{2.078794in}{2.762626in}}%
\pgfpathlineto{\pgfqpoint{2.079165in}{2.785907in}}%
\pgfpathlineto{\pgfqpoint{2.079785in}{2.644559in}}%
\pgfpathlineto{\pgfqpoint{2.082137in}{2.644559in}}%
\pgfpathlineto{\pgfqpoint{2.082142in}{2.749018in}}%
\pgfpathlineto{\pgfqpoint{2.082274in}{2.641671in}}%
\pgfpathlineto{\pgfqpoint{2.082867in}{2.678819in}}%
\pgfpathlineto{\pgfqpoint{2.083708in}{2.760273in}}%
\pgfpathlineto{\pgfqpoint{2.084083in}{2.783768in}}%
\pgfpathlineto{\pgfqpoint{2.084702in}{2.644559in}}%
\pgfpathlineto{\pgfqpoint{2.087082in}{2.644559in}}%
\pgfpathlineto{\pgfqpoint{2.087088in}{2.748531in}}%
\pgfpathlineto{\pgfqpoint{2.087159in}{2.640068in}}%
\pgfpathlineto{\pgfqpoint{2.087844in}{2.680613in}}%
\pgfpathlineto{\pgfqpoint{2.088685in}{2.762636in}}%
\pgfpathlineto{\pgfqpoint{2.089056in}{2.785933in}}%
\pgfpathlineto{\pgfqpoint{2.089676in}{2.644559in}}%
\pgfpathlineto{\pgfqpoint{2.092028in}{2.644559in}}%
\pgfpathlineto{\pgfqpoint{2.092033in}{2.748995in}}%
\pgfpathlineto{\pgfqpoint{2.092089in}{2.640698in}}%
\pgfpathlineto{\pgfqpoint{2.092766in}{2.678406in}}%
\pgfpathlineto{\pgfqpoint{2.093606in}{2.761520in}}%
\pgfpathlineto{\pgfqpoint{2.094002in}{2.786329in}}%
\pgfpathlineto{\pgfqpoint{2.094588in}{2.644559in}}%
\pgfpathlineto{\pgfqpoint{2.096973in}{2.644559in}}%
\pgfpathlineto{\pgfqpoint{2.096979in}{2.748516in}}%
\pgfpathlineto{\pgfqpoint{2.097047in}{2.642223in}}%
\pgfpathlineto{\pgfqpoint{2.097733in}{2.682525in}}%
\pgfpathlineto{\pgfqpoint{2.098573in}{2.760846in}}%
\pgfpathlineto{\pgfqpoint{2.098947in}{2.784289in}}%
\pgfpathlineto{\pgfqpoint{2.099595in}{2.644559in}}%
\pgfpathlineto{\pgfqpoint{2.101919in}{2.644559in}}%
\pgfpathlineto{\pgfqpoint{2.101924in}{2.748076in}}%
\pgfpathlineto{\pgfqpoint{2.101983in}{2.641995in}}%
\pgfpathlineto{\pgfqpoint{2.102664in}{2.680734in}}%
\pgfpathlineto{\pgfqpoint{2.103505in}{2.760205in}}%
\pgfpathlineto{\pgfqpoint{2.103893in}{2.784534in}}%
\pgfpathlineto{\pgfqpoint{2.104513in}{2.644559in}}%
\pgfpathlineto{\pgfqpoint{2.106864in}{2.644559in}}%
\pgfpathlineto{\pgfqpoint{2.106870in}{2.748997in}}%
\pgfpathlineto{\pgfqpoint{2.106926in}{2.640642in}}%
\pgfpathlineto{\pgfqpoint{2.107603in}{2.678519in}}%
\pgfpathlineto{\pgfqpoint{2.108444in}{2.761598in}}%
\pgfpathlineto{\pgfqpoint{2.108838in}{2.786316in}}%
\pgfpathlineto{\pgfqpoint{2.109428in}{2.644559in}}%
\pgfpathlineto{\pgfqpoint{2.111810in}{2.644559in}}%
\pgfpathlineto{\pgfqpoint{2.111815in}{2.748821in}}%
\pgfpathlineto{\pgfqpoint{2.111897in}{2.631306in}}%
\pgfpathlineto{\pgfqpoint{2.112582in}{2.672838in}}%
\pgfpathlineto{\pgfqpoint{2.113423in}{2.771400in}}%
\pgfpathlineto{\pgfqpoint{2.113784in}{2.793965in}}%
\pgfpathlineto{\pgfqpoint{2.114342in}{2.644559in}}%
\pgfpathlineto{\pgfqpoint{2.116755in}{2.644559in}}%
\pgfpathlineto{\pgfqpoint{2.116761in}{2.748947in}}%
\pgfpathlineto{\pgfqpoint{2.116840in}{2.633081in}}%
\pgfpathlineto{\pgfqpoint{2.117525in}{2.674440in}}%
\pgfpathlineto{\pgfqpoint{2.118366in}{2.770203in}}%
\pgfpathlineto{\pgfqpoint{2.118729in}{2.792923in}}%
\pgfpathlineto{\pgfqpoint{2.119294in}{2.644559in}}%
\pgfpathlineto{\pgfqpoint{2.121701in}{2.644559in}}%
\pgfpathlineto{\pgfqpoint{2.121706in}{2.746589in}}%
\pgfpathlineto{\pgfqpoint{2.122459in}{2.694664in}}%
\pgfpathlineto{\pgfqpoint{2.123674in}{2.771524in}}%
\pgfpathlineto{\pgfqpoint{2.123699in}{2.756710in}}%
\pgfpathlineto{\pgfqpoint{2.124500in}{2.644559in}}%
\pgfpathlineto{\pgfqpoint{2.126646in}{2.644559in}}%
\pgfpathlineto{\pgfqpoint{2.126651in}{2.748212in}}%
\pgfpathlineto{\pgfqpoint{2.126712in}{2.641812in}}%
\pgfpathlineto{\pgfqpoint{2.127394in}{2.680890in}}%
\pgfpathlineto{\pgfqpoint{2.128235in}{2.760839in}}%
\pgfpathlineto{\pgfqpoint{2.128620in}{2.784996in}}%
\pgfpathlineto{\pgfqpoint{2.129243in}{2.644559in}}%
\pgfpathlineto{\pgfqpoint{2.131592in}{2.644559in}}%
\pgfpathlineto{\pgfqpoint{2.131597in}{2.748121in}}%
\pgfpathlineto{\pgfqpoint{2.131657in}{2.641630in}}%
\pgfpathlineto{\pgfqpoint{2.132339in}{2.680651in}}%
\pgfpathlineto{\pgfqpoint{2.133180in}{2.760578in}}%
\pgfpathlineto{\pgfqpoint{2.133565in}{2.784762in}}%
\pgfpathlineto{\pgfqpoint{2.134187in}{2.644559in}}%
\pgfpathlineto{\pgfqpoint{2.136537in}{2.644559in}}%
\pgfpathlineto{\pgfqpoint{2.136542in}{2.748982in}}%
\pgfpathlineto{\pgfqpoint{2.136672in}{2.641404in}}%
\pgfpathlineto{\pgfqpoint{2.137266in}{2.678539in}}%
\pgfpathlineto{\pgfqpoint{2.138106in}{2.760237in}}%
\pgfpathlineto{\pgfqpoint{2.138482in}{2.783782in}}%
\pgfpathlineto{\pgfqpoint{2.139099in}{2.644559in}}%
\pgfpathlineto{\pgfqpoint{2.141482in}{2.644559in}}%
\pgfpathlineto{\pgfqpoint{2.141488in}{2.748995in}}%
\pgfpathlineto{\pgfqpoint{2.141614in}{2.640879in}}%
\pgfpathlineto{\pgfqpoint{2.142257in}{2.710675in}}%
\pgfpathlineto{\pgfqpoint{2.142306in}{2.684199in}}%
\pgfpathlineto{\pgfqpoint{2.142850in}{2.747917in}}%
\pgfpathlineto{\pgfqpoint{2.142900in}{2.721519in}}%
\pgfpathlineto{\pgfqpoint{2.143425in}{2.783997in}}%
\pgfpathlineto{\pgfqpoint{2.143539in}{2.644559in}}%
\pgfpathlineto{\pgfqpoint{2.143550in}{2.644559in}}%
\pgfpathlineto{\pgfqpoint{2.146428in}{2.644559in}}%
\pgfpathlineto{\pgfqpoint{2.146433in}{2.748656in}}%
\pgfpathlineto{\pgfqpoint{2.146507in}{2.638003in}}%
\pgfpathlineto{\pgfqpoint{2.147192in}{2.678798in}}%
\pgfpathlineto{\pgfqpoint{2.148033in}{2.764771in}}%
\pgfpathlineto{\pgfqpoint{2.148402in}{2.787908in}}%
\pgfpathlineto{\pgfqpoint{2.149001in}{2.644559in}}%
\pgfpathlineto{\pgfqpoint{2.151373in}{2.644559in}}%
\pgfpathlineto{\pgfqpoint{2.151379in}{2.748462in}}%
\pgfpathlineto{\pgfqpoint{2.151452in}{2.638127in}}%
\pgfpathlineto{\pgfqpoint{2.152137in}{2.678887in}}%
\pgfpathlineto{\pgfqpoint{2.152978in}{2.763953in}}%
\pgfpathlineto{\pgfqpoint{2.153347in}{2.787105in}}%
\pgfpathlineto{\pgfqpoint{2.153948in}{2.644559in}}%
\pgfpathlineto{\pgfqpoint{2.156319in}{2.644559in}}%
\pgfpathlineto{\pgfqpoint{2.156324in}{2.749096in}}%
\pgfpathlineto{\pgfqpoint{2.156455in}{2.641946in}}%
\pgfpathlineto{\pgfqpoint{2.157048in}{2.679086in}}%
\pgfpathlineto{\pgfqpoint{2.157889in}{2.760108in}}%
\pgfpathlineto{\pgfqpoint{2.158264in}{2.783637in}}%
\pgfpathlineto{\pgfqpoint{2.158886in}{2.644559in}}%
\pgfpathlineto{\pgfqpoint{2.161264in}{2.644559in}}%
\pgfpathlineto{\pgfqpoint{2.161270in}{2.747636in}}%
\pgfpathlineto{\pgfqpoint{2.162034in}{2.705418in}}%
\pgfpathlineto{\pgfqpoint{2.162083in}{2.688963in}}%
\pgfpathlineto{\pgfqpoint{2.162627in}{2.742654in}}%
\pgfpathlineto{\pgfqpoint{2.162677in}{2.726289in}}%
\pgfpathlineto{\pgfqpoint{2.163205in}{2.778878in}}%
\pgfpathlineto{\pgfqpoint{2.163292in}{2.718932in}}%
\pgfpathlineto{\pgfqpoint{2.163306in}{2.725375in}}%
\pgfpathlineto{\pgfqpoint{2.164091in}{2.644559in}}%
\pgfpathlineto{\pgfqpoint{2.166210in}{2.644559in}}%
\pgfpathlineto{\pgfqpoint{2.166215in}{2.748160in}}%
\pgfpathlineto{\pgfqpoint{2.166274in}{2.642276in}}%
\pgfpathlineto{\pgfqpoint{2.166955in}{2.680999in}}%
\pgfpathlineto{\pgfqpoint{2.167795in}{2.760392in}}%
\pgfpathlineto{\pgfqpoint{2.168184in}{2.784733in}}%
\pgfpathlineto{\pgfqpoint{2.168811in}{2.644559in}}%
\pgfpathlineto{\pgfqpoint{2.171155in}{2.644559in}}%
\pgfpathlineto{\pgfqpoint{2.171161in}{2.748656in}}%
\pgfpathlineto{\pgfqpoint{2.171239in}{2.633683in}}%
\pgfpathlineto{\pgfqpoint{2.171925in}{2.674970in}}%
\pgfpathlineto{\pgfqpoint{2.172765in}{2.768668in}}%
\pgfpathlineto{\pgfqpoint{2.173129in}{2.791437in}}%
\pgfpathlineto{\pgfqpoint{2.173694in}{2.644559in}}%
\pgfpathlineto{\pgfqpoint{2.176101in}{2.644559in}}%
\pgfpathlineto{\pgfqpoint{2.176106in}{2.747951in}}%
\pgfpathlineto{\pgfqpoint{2.176159in}{2.644322in}}%
\pgfpathlineto{\pgfqpoint{2.176834in}{2.681617in}}%
\pgfpathlineto{\pgfqpoint{2.177675in}{2.758212in}}%
\pgfpathlineto{\pgfqpoint{2.178048in}{2.781605in}}%
\pgfpathlineto{\pgfqpoint{2.178718in}{2.644559in}}%
\pgfpathlineto{\pgfqpoint{2.181046in}{2.644559in}}%
\pgfpathlineto{\pgfqpoint{2.181051in}{2.747965in}}%
\pgfpathlineto{\pgfqpoint{2.181111in}{2.641402in}}%
\pgfpathlineto{\pgfqpoint{2.181793in}{2.680313in}}%
\pgfpathlineto{\pgfqpoint{2.182633in}{2.760149in}}%
\pgfpathlineto{\pgfqpoint{2.183020in}{2.784384in}}%
\pgfpathlineto{\pgfqpoint{2.183641in}{2.644559in}}%
\pgfpathlineto{\pgfqpoint{2.185992in}{2.644559in}}%
\pgfpathlineto{\pgfqpoint{2.185997in}{2.748893in}}%
\pgfpathlineto{\pgfqpoint{2.186075in}{2.634124in}}%
\pgfpathlineto{\pgfqpoint{2.186761in}{2.675384in}}%
\pgfpathlineto{\pgfqpoint{2.187601in}{2.769069in}}%
\pgfpathlineto{\pgfqpoint{2.187965in}{2.791868in}}%
\pgfpathlineto{\pgfqpoint{2.188535in}{2.644559in}}%
\pgfpathlineto{\pgfqpoint{2.190937in}{2.644559in}}%
\pgfpathlineto{\pgfqpoint{2.190942in}{2.747374in}}%
\pgfpathlineto{\pgfqpoint{2.191696in}{2.702281in}}%
\pgfpathlineto{\pgfqpoint{2.191745in}{2.690272in}}%
\pgfpathlineto{\pgfqpoint{2.192289in}{2.739525in}}%
\pgfpathlineto{\pgfqpoint{2.192339in}{2.727590in}}%
\pgfpathlineto{\pgfqpoint{2.192872in}{2.776099in}}%
\pgfpathlineto{\pgfqpoint{2.192974in}{2.701365in}}%
\pgfpathlineto{\pgfqpoint{2.193755in}{2.644559in}}%
\pgfpathlineto{\pgfqpoint{2.195882in}{2.644559in}}%
\pgfpathlineto{\pgfqpoint{2.195888in}{2.747874in}}%
\pgfpathlineto{\pgfqpoint{2.196612in}{2.681890in}}%
\pgfpathlineto{\pgfqpoint{2.197453in}{2.757430in}}%
\pgfpathlineto{\pgfqpoint{2.197828in}{2.780942in}}%
\pgfpathlineto{\pgfqpoint{2.198510in}{2.644559in}}%
\pgfpathlineto{\pgfqpoint{2.200828in}{2.644559in}}%
\pgfpathlineto{\pgfqpoint{2.200833in}{2.748565in}}%
\pgfpathlineto{\pgfqpoint{2.200908in}{2.636327in}}%
\pgfpathlineto{\pgfqpoint{2.201594in}{2.677278in}}%
\pgfpathlineto{\pgfqpoint{2.202434in}{2.765915in}}%
\pgfpathlineto{\pgfqpoint{2.202802in}{2.788939in}}%
\pgfpathlineto{\pgfqpoint{2.203390in}{2.644559in}}%
\pgfpathlineto{\pgfqpoint{2.205773in}{2.644559in}}%
\pgfpathlineto{\pgfqpoint{2.205779in}{2.748199in}}%
\pgfpathlineto{\pgfqpoint{2.205841in}{2.641210in}}%
\pgfpathlineto{\pgfqpoint{2.206524in}{2.680607in}}%
\pgfpathlineto{\pgfqpoint{2.207365in}{2.761130in}}%
\pgfpathlineto{\pgfqpoint{2.207747in}{2.785109in}}%
\pgfpathlineto{\pgfqpoint{2.208356in}{2.644559in}}%
\pgfpathlineto{\pgfqpoint{2.210719in}{2.644559in}}%
\pgfpathlineto{\pgfqpoint{2.210724in}{2.748615in}}%
\pgfpathlineto{\pgfqpoint{2.210800in}{2.636024in}}%
\pgfpathlineto{\pgfqpoint{2.211486in}{2.677090in}}%
\pgfpathlineto{\pgfqpoint{2.212326in}{2.766409in}}%
\pgfpathlineto{\pgfqpoint{2.212693in}{2.789349in}}%
\pgfpathlineto{\pgfqpoint{2.213274in}{2.644559in}}%
\pgfpathlineto{\pgfqpoint{2.215664in}{2.644559in}}%
\pgfpathlineto{\pgfqpoint{2.215670in}{2.748665in}}%
\pgfpathlineto{\pgfqpoint{2.215748in}{2.633952in}}%
\pgfpathlineto{\pgfqpoint{2.216434in}{2.675214in}}%
\pgfpathlineto{\pgfqpoint{2.217274in}{2.768450in}}%
\pgfpathlineto{\pgfqpoint{2.217638in}{2.791240in}}%
\pgfpathlineto{\pgfqpoint{2.218205in}{2.644559in}}%
\pgfpathlineto{\pgfqpoint{2.220610in}{2.644559in}}%
\pgfpathlineto{\pgfqpoint{2.220615in}{2.747248in}}%
\pgfpathlineto{\pgfqpoint{2.221356in}{2.699955in}}%
\pgfpathlineto{\pgfqpoint{2.221406in}{2.691594in}}%
\pgfpathlineto{\pgfqpoint{2.221851in}{2.730999in}}%
\pgfpathlineto{\pgfqpoint{2.221900in}{2.722687in}}%
\pgfpathlineto{\pgfqpoint{2.222539in}{2.774172in}}%
\pgfpathlineto{\pgfqpoint{2.222619in}{2.715109in}}%
\pgfpathlineto{\pgfqpoint{2.222628in}{2.749606in}}%
\pgfpathlineto{\pgfqpoint{2.223420in}{2.644559in}}%
\pgfpathlineto{\pgfqpoint{2.225555in}{2.644559in}}%
\pgfpathlineto{\pgfqpoint{2.225561in}{2.748862in}}%
\pgfpathlineto{\pgfqpoint{2.225643in}{2.630965in}}%
\pgfpathlineto{\pgfqpoint{2.226279in}{2.715759in}}%
\pgfpathlineto{\pgfqpoint{2.226328in}{2.672527in}}%
\pgfpathlineto{\pgfqpoint{2.226971in}{2.759140in}}%
\pgfpathlineto{\pgfqpoint{2.227021in}{2.716136in}}%
\pgfpathlineto{\pgfqpoint{2.227529in}{2.794070in}}%
\pgfpathlineto{\pgfqpoint{2.227616in}{2.644559in}}%
\pgfpathlineto{\pgfqpoint{2.227693in}{2.644559in}}%
\pgfpathlineto{\pgfqpoint{2.230501in}{2.644559in}}%
\pgfpathlineto{\pgfqpoint{2.230506in}{2.748118in}}%
\pgfpathlineto{\pgfqpoint{2.230562in}{2.643117in}}%
\pgfpathlineto{\pgfqpoint{2.231241in}{2.681317in}}%
\pgfpathlineto{\pgfqpoint{2.232082in}{2.759702in}}%
\pgfpathlineto{\pgfqpoint{2.232474in}{2.784303in}}%
\pgfpathlineto{\pgfqpoint{2.233105in}{2.644559in}}%
\pgfpathlineto{\pgfqpoint{2.235446in}{2.644559in}}%
\pgfpathlineto{\pgfqpoint{2.235451in}{2.747956in}}%
\pgfpathlineto{\pgfqpoint{2.235510in}{2.641879in}}%
\pgfpathlineto{\pgfqpoint{2.236190in}{2.680505in}}%
\pgfpathlineto{\pgfqpoint{2.237031in}{2.759821in}}%
\pgfpathlineto{\pgfqpoint{2.237420in}{2.784202in}}%
\pgfpathlineto{\pgfqpoint{2.238045in}{2.644559in}}%
\pgfpathlineto{\pgfqpoint{2.240392in}{2.644559in}}%
\pgfpathlineto{\pgfqpoint{2.240397in}{2.748103in}}%
\pgfpathlineto{\pgfqpoint{2.240453in}{2.643042in}}%
\pgfpathlineto{\pgfqpoint{2.241133in}{2.681257in}}%
\pgfpathlineto{\pgfqpoint{2.241973in}{2.759675in}}%
\pgfpathlineto{\pgfqpoint{2.242365in}{2.784267in}}%
\pgfpathlineto{\pgfqpoint{2.242998in}{2.644559in}}%
\pgfpathlineto{\pgfqpoint{2.245337in}{2.644559in}}%
\pgfpathlineto{\pgfqpoint{2.245342in}{2.747458in}}%
\pgfpathlineto{\pgfqpoint{2.246098in}{2.703087in}}%
\pgfpathlineto{\pgfqpoint{2.246148in}{2.689992in}}%
\pgfpathlineto{\pgfqpoint{2.246692in}{2.740330in}}%
\pgfpathlineto{\pgfqpoint{2.246741in}{2.727312in}}%
\pgfpathlineto{\pgfqpoint{2.247273in}{2.776820in}}%
\pgfpathlineto{\pgfqpoint{2.247374in}{2.705501in}}%
\pgfpathlineto{\pgfqpoint{2.248157in}{2.644559in}}%
\pgfpathlineto{\pgfqpoint{2.250282in}{2.644559in}}%
\pgfpathlineto{\pgfqpoint{2.250288in}{2.748187in}}%
\pgfpathlineto{\pgfqpoint{2.250349in}{2.641610in}}%
\pgfpathlineto{\pgfqpoint{2.251031in}{2.680762in}}%
\pgfpathlineto{\pgfqpoint{2.251872in}{2.760841in}}%
\pgfpathlineto{\pgfqpoint{2.252256in}{2.784957in}}%
\pgfpathlineto{\pgfqpoint{2.252871in}{2.644559in}}%
\pgfpathlineto{\pgfqpoint{2.255228in}{2.644559in}}%
\pgfpathlineto{\pgfqpoint{2.255234in}{2.746649in}}%
\pgfpathlineto{\pgfqpoint{2.255351in}{2.627323in}}%
\pgfpathlineto{\pgfqpoint{2.255994in}{2.722342in}}%
\pgfpathlineto{\pgfqpoint{2.256044in}{2.670590in}}%
\pgfpathlineto{\pgfqpoint{2.256687in}{2.765733in}}%
\pgfpathlineto{\pgfqpoint{2.256736in}{2.714189in}}%
\pgfpathlineto{\pgfqpoint{2.257167in}{2.795811in}}%
\pgfpathlineto{\pgfqpoint{2.257292in}{2.644559in}}%
\pgfpathlineto{\pgfqpoint{2.257447in}{2.644559in}}%
\pgfpathlineto{\pgfqpoint{2.260173in}{2.644559in}}%
\pgfpathlineto{\pgfqpoint{2.260179in}{2.748232in}}%
\pgfpathlineto{\pgfqpoint{2.260239in}{2.641963in}}%
\pgfpathlineto{\pgfqpoint{2.260921in}{2.680993in}}%
\pgfpathlineto{\pgfqpoint{2.261762in}{2.760852in}}%
\pgfpathlineto{\pgfqpoint{2.262147in}{2.785036in}}%
\pgfpathlineto{\pgfqpoint{2.262771in}{2.644559in}}%
\pgfpathlineto{\pgfqpoint{2.265119in}{2.644559in}}%
\pgfpathlineto{\pgfqpoint{2.265124in}{2.748333in}}%
\pgfpathlineto{\pgfqpoint{2.265200in}{2.635804in}}%
\pgfpathlineto{\pgfqpoint{2.265885in}{2.676792in}}%
\pgfpathlineto{\pgfqpoint{2.266726in}{2.765654in}}%
\pgfpathlineto{\pgfqpoint{2.267093in}{2.788644in}}%
\pgfpathlineto{\pgfqpoint{2.267679in}{2.644559in}}%
\pgfpathlineto{\pgfqpoint{2.270064in}{2.644559in}}%
\pgfpathlineto{\pgfqpoint{2.270070in}{2.748942in}}%
\pgfpathlineto{\pgfqpoint{2.270151in}{2.631787in}}%
\pgfpathlineto{\pgfqpoint{2.270836in}{2.673280in}}%
\pgfpathlineto{\pgfqpoint{2.271677in}{2.771329in}}%
\pgfpathlineto{\pgfqpoint{2.272038in}{2.793934in}}%
\pgfpathlineto{\pgfqpoint{2.272597in}{2.644559in}}%
\pgfpathlineto{\pgfqpoint{2.275010in}{2.644559in}}%
\pgfpathlineto{\pgfqpoint{2.275015in}{2.747174in}}%
\pgfpathlineto{\pgfqpoint{2.275772in}{2.694022in}}%
\pgfpathlineto{\pgfqpoint{2.276984in}{2.772519in}}%
\pgfpathlineto{\pgfqpoint{2.276984in}{2.747714in}}%
\pgfpathlineto{\pgfqpoint{2.277005in}{2.755094in}}%
\pgfpathlineto{\pgfqpoint{2.277727in}{2.644559in}}%
\pgfpathlineto{\pgfqpoint{2.279955in}{2.644559in}}%
\pgfpathlineto{\pgfqpoint{2.279960in}{2.747929in}}%
\pgfpathlineto{\pgfqpoint{2.280016in}{2.642779in}}%
\pgfpathlineto{\pgfqpoint{2.280695in}{2.680885in}}%
\pgfpathlineto{\pgfqpoint{2.281536in}{2.759155in}}%
\pgfpathlineto{\pgfqpoint{2.281929in}{2.783800in}}%
\pgfpathlineto{\pgfqpoint{2.282556in}{2.644559in}}%
\pgfpathlineto{\pgfqpoint{2.284901in}{2.644559in}}%
\pgfpathlineto{\pgfqpoint{2.284906in}{2.748842in}}%
\pgfpathlineto{\pgfqpoint{2.284964in}{2.640557in}}%
\pgfpathlineto{\pgfqpoint{2.285644in}{2.678937in}}%
\pgfpathlineto{\pgfqpoint{2.286485in}{2.761342in}}%
\pgfpathlineto{\pgfqpoint{2.286874in}{2.785780in}}%
\pgfpathlineto{\pgfqpoint{2.287475in}{2.644559in}}%
\pgfpathlineto{\pgfqpoint{2.289846in}{2.644559in}}%
\pgfpathlineto{\pgfqpoint{2.289851in}{2.747777in}}%
\pgfpathlineto{\pgfqpoint{2.289904in}{2.643988in}}%
\pgfpathlineto{\pgfqpoint{2.290579in}{2.681254in}}%
\pgfpathlineto{\pgfqpoint{2.291420in}{2.757719in}}%
\pgfpathlineto{\pgfqpoint{2.291793in}{2.781116in}}%
\pgfpathlineto{\pgfqpoint{2.292460in}{2.644559in}}%
\pgfpathlineto{\pgfqpoint{2.294792in}{2.644559in}}%
\pgfpathlineto{\pgfqpoint{2.294797in}{2.748818in}}%
\pgfpathlineto{\pgfqpoint{2.294879in}{2.631347in}}%
\pgfpathlineto{\pgfqpoint{2.295564in}{2.672876in}}%
\pgfpathlineto{\pgfqpoint{2.296405in}{2.771348in}}%
\pgfpathlineto{\pgfqpoint{2.296765in}{2.793917in}}%
\pgfpathlineto{\pgfqpoint{2.297322in}{2.644559in}}%
\pgfpathlineto{\pgfqpoint{2.299737in}{2.644559in}}%
\pgfpathlineto{\pgfqpoint{2.299742in}{2.748285in}}%
\pgfpathlineto{\pgfqpoint{2.299809in}{2.643887in}}%
\pgfpathlineto{\pgfqpoint{2.300494in}{2.683921in}}%
\pgfpathlineto{\pgfqpoint{2.301335in}{2.758549in}}%
\pgfpathlineto{\pgfqpoint{2.301711in}{2.782136in}}%
\pgfpathlineto{\pgfqpoint{2.302409in}{2.644559in}}%
\pgfpathlineto{\pgfqpoint{2.304682in}{2.644559in}}%
\pgfpathlineto{\pgfqpoint{2.304688in}{2.748213in}}%
\pgfpathlineto{\pgfqpoint{2.304755in}{2.643555in}}%
\pgfpathlineto{\pgfqpoint{2.305440in}{2.683631in}}%
\pgfpathlineto{\pgfqpoint{2.306281in}{2.758475in}}%
\pgfpathlineto{\pgfqpoint{2.306656in}{2.782030in}}%
\pgfpathlineto{\pgfqpoint{2.307346in}{2.644559in}}%
\pgfpathlineto{\pgfqpoint{2.309628in}{2.644559in}}%
\pgfpathlineto{\pgfqpoint{2.309633in}{2.748461in}}%
\pgfpathlineto{\pgfqpoint{2.309705in}{2.639245in}}%
\pgfpathlineto{\pgfqpoint{2.310390in}{2.679877in}}%
\pgfpathlineto{\pgfqpoint{2.311231in}{2.763018in}}%
\pgfpathlineto{\pgfqpoint{2.311602in}{2.786256in}}%
\pgfpathlineto{\pgfqpoint{2.312216in}{2.644559in}}%
\pgfpathlineto{\pgfqpoint{2.314573in}{2.644559in}}%
\pgfpathlineto{\pgfqpoint{2.314579in}{2.748734in}}%
\pgfpathlineto{\pgfqpoint{2.314658in}{2.633577in}}%
\pgfpathlineto{\pgfqpoint{2.315343in}{2.674878in}}%
\pgfpathlineto{\pgfqpoint{2.316184in}{2.769018in}}%
\pgfpathlineto{\pgfqpoint{2.316547in}{2.791778in}}%
\pgfpathlineto{\pgfqpoint{2.317116in}{2.644559in}}%
\pgfpathlineto{\pgfqpoint{2.319519in}{2.644559in}}%
\pgfpathlineto{\pgfqpoint{2.319524in}{2.747767in}}%
\pgfpathlineto{\pgfqpoint{2.319656in}{2.644312in}}%
\pgfpathlineto{\pgfqpoint{2.320249in}{2.681480in}}%
\pgfpathlineto{\pgfqpoint{2.321090in}{2.757346in}}%
\pgfpathlineto{\pgfqpoint{2.321465in}{2.780820in}}%
\pgfpathlineto{\pgfqpoint{2.322139in}{2.644559in}}%
\pgfpathlineto{\pgfqpoint{2.324464in}{2.644559in}}%
\pgfpathlineto{\pgfqpoint{2.324469in}{2.747728in}}%
\pgfpathlineto{\pgfqpoint{2.325237in}{2.706314in}}%
\pgfpathlineto{\pgfqpoint{2.325287in}{2.688632in}}%
\pgfpathlineto{\pgfqpoint{2.325831in}{2.743548in}}%
\pgfpathlineto{\pgfqpoint{2.325880in}{2.725960in}}%
\pgfpathlineto{\pgfqpoint{2.326406in}{2.779662in}}%
\pgfpathlineto{\pgfqpoint{2.326492in}{2.717984in}}%
\pgfpathlineto{\pgfqpoint{2.326506in}{2.726569in}}%
\pgfpathlineto{\pgfqpoint{2.327291in}{2.644559in}}%
\pgfpathlineto{\pgfqpoint{2.329410in}{2.644559in}}%
\pgfpathlineto{\pgfqpoint{2.329415in}{2.747885in}}%
\pgfpathlineto{\pgfqpoint{2.329469in}{2.643743in}}%
\pgfpathlineto{\pgfqpoint{2.330145in}{2.681262in}}%
\pgfpathlineto{\pgfqpoint{2.330986in}{2.758350in}}%
\pgfpathlineto{\pgfqpoint{2.331358in}{2.781683in}}%
\pgfpathlineto{\pgfqpoint{2.332024in}{2.644559in}}%
\pgfpathlineto{\pgfqpoint{2.334355in}{2.644559in}}%
\pgfpathlineto{\pgfqpoint{2.334361in}{2.748567in}}%
\pgfpathlineto{\pgfqpoint{2.334438in}{2.635074in}}%
\pgfpathlineto{\pgfqpoint{2.335123in}{2.676224in}}%
\pgfpathlineto{\pgfqpoint{2.335964in}{2.767106in}}%
\pgfpathlineto{\pgfqpoint{2.336329in}{2.789980in}}%
\pgfpathlineto{\pgfqpoint{2.336904in}{2.644559in}}%
\pgfpathlineto{\pgfqpoint{2.339301in}{2.644559in}}%
\pgfpathlineto{\pgfqpoint{2.339306in}{2.747613in}}%
\pgfpathlineto{\pgfqpoint{2.340066in}{2.704459in}}%
\pgfpathlineto{\pgfqpoint{2.340116in}{2.689623in}}%
\pgfpathlineto{\pgfqpoint{2.340660in}{2.741699in}}%
\pgfpathlineto{\pgfqpoint{2.340709in}{2.726946in}}%
\pgfpathlineto{\pgfqpoint{2.341239in}{2.778053in}}%
\pgfpathlineto{\pgfqpoint{2.341328in}{2.720535in}}%
\pgfpathlineto{\pgfqpoint{2.341342in}{2.723814in}}%
\pgfpathlineto{\pgfqpoint{2.342126in}{2.644559in}}%
\pgfpathlineto{\pgfqpoint{2.344246in}{2.644559in}}%
\pgfpathlineto{\pgfqpoint{2.344251in}{2.747719in}}%
\pgfpathlineto{\pgfqpoint{2.344971in}{2.682087in}}%
\pgfpathlineto{\pgfqpoint{2.345812in}{2.756241in}}%
\pgfpathlineto{\pgfqpoint{2.346189in}{2.779887in}}%
\pgfpathlineto{\pgfqpoint{2.346886in}{2.644559in}}%
\pgfpathlineto{\pgfqpoint{2.349192in}{2.644559in}}%
\pgfpathlineto{\pgfqpoint{2.349197in}{2.747451in}}%
\pgfpathlineto{\pgfqpoint{2.349950in}{2.702498in}}%
\pgfpathlineto{\pgfqpoint{2.349999in}{2.690621in}}%
\pgfpathlineto{\pgfqpoint{2.350543in}{2.739743in}}%
\pgfpathlineto{\pgfqpoint{2.350592in}{2.727939in}}%
\pgfpathlineto{\pgfqpoint{2.351126in}{2.776340in}}%
\pgfpathlineto{\pgfqpoint{2.351228in}{2.700694in}}%
\pgfpathlineto{\pgfqpoint{2.352011in}{2.644559in}}%
\pgfpathlineto{\pgfqpoint{2.354137in}{2.644559in}}%
\pgfpathlineto{\pgfqpoint{2.354142in}{2.748113in}}%
\pgfpathlineto{\pgfqpoint{2.354891in}{2.684915in}}%
\pgfpathlineto{\pgfqpoint{2.355831in}{2.762951in}}%
\pgfpathlineto{\pgfqpoint{2.356111in}{2.780502in}}%
\pgfpathlineto{\pgfqpoint{2.356969in}{2.644559in}}%
\pgfpathlineto{\pgfqpoint{2.359082in}{2.644559in}}%
\pgfpathlineto{\pgfqpoint{2.359088in}{2.748437in}}%
\pgfpathlineto{\pgfqpoint{2.359158in}{2.640872in}}%
\pgfpathlineto{\pgfqpoint{2.359843in}{2.681306in}}%
\pgfpathlineto{\pgfqpoint{2.360684in}{2.761635in}}%
\pgfpathlineto{\pgfqpoint{2.361056in}{2.784997in}}%
\pgfpathlineto{\pgfqpoint{2.361691in}{2.644559in}}%
\pgfpathlineto{\pgfqpoint{2.364028in}{2.644559in}}%
\pgfpathlineto{\pgfqpoint{2.364033in}{2.748552in}}%
\pgfpathlineto{\pgfqpoint{2.364113in}{2.633065in}}%
\pgfpathlineto{\pgfqpoint{2.364798in}{2.674417in}}%
\pgfpathlineto{\pgfqpoint{2.365639in}{2.768893in}}%
\pgfpathlineto{\pgfqpoint{2.366002in}{2.791604in}}%
\pgfpathlineto{\pgfqpoint{2.366564in}{2.644559in}}%
\pgfpathlineto{\pgfqpoint{2.368973in}{2.644559in}}%
\pgfpathlineto{\pgfqpoint{2.368979in}{2.748721in}}%
\pgfpathlineto{\pgfqpoint{2.369061in}{2.630907in}}%
\pgfpathlineto{\pgfqpoint{2.369697in}{2.715404in}}%
\pgfpathlineto{\pgfqpoint{2.369747in}{2.672469in}}%
\pgfpathlineto{\pgfqpoint{2.370389in}{2.758785in}}%
\pgfpathlineto{\pgfqpoint{2.370439in}{2.716078in}}%
\pgfpathlineto{\pgfqpoint{2.370947in}{2.793712in}}%
\pgfpathlineto{\pgfqpoint{2.371034in}{2.644559in}}%
\pgfpathlineto{\pgfqpoint{2.371111in}{2.644559in}}%
\pgfpathlineto{\pgfqpoint{2.373919in}{2.644559in}}%
\pgfpathlineto{\pgfqpoint{2.373924in}{2.748649in}}%
\pgfpathlineto{\pgfqpoint{2.373998in}{2.637378in}}%
\pgfpathlineto{\pgfqpoint{2.374683in}{2.678239in}}%
\pgfpathlineto{\pgfqpoint{2.375524in}{2.765256in}}%
\pgfpathlineto{\pgfqpoint{2.375893in}{2.788346in}}%
\pgfpathlineto{\pgfqpoint{2.376487in}{2.644559in}}%
\pgfpathlineto{\pgfqpoint{2.378864in}{2.644559in}}%
\pgfpathlineto{\pgfqpoint{2.378870in}{2.747210in}}%
\pgfpathlineto{\pgfqpoint{2.379612in}{2.699962in}}%
\pgfpathlineto{\pgfqpoint{2.379662in}{2.691267in}}%
\pgfpathlineto{\pgfqpoint{2.380107in}{2.731005in}}%
\pgfpathlineto{\pgfqpoint{2.380156in}{2.722361in}}%
\pgfpathlineto{\pgfqpoint{2.380794in}{2.774127in}}%
\pgfpathlineto{\pgfqpoint{2.380874in}{2.714850in}}%
\pgfpathlineto{\pgfqpoint{2.380883in}{2.749346in}}%
\pgfpathlineto{\pgfqpoint{2.381674in}{2.644559in}}%
\pgfpathlineto{\pgfqpoint{2.383810in}{2.644559in}}%
\pgfpathlineto{\pgfqpoint{2.383815in}{2.748821in}}%
\pgfpathlineto{\pgfqpoint{2.383882in}{2.639126in}}%
\pgfpathlineto{\pgfqpoint{2.384568in}{2.679104in}}%
\pgfpathlineto{\pgfqpoint{2.385409in}{2.762998in}}%
\pgfpathlineto{\pgfqpoint{2.385784in}{2.786498in}}%
\pgfpathlineto{\pgfqpoint{2.386387in}{2.644559in}}%
\pgfpathlineto{\pgfqpoint{2.388755in}{2.644559in}}%
\pgfpathlineto{\pgfqpoint{2.388761in}{2.748481in}}%
\pgfpathlineto{\pgfqpoint{2.388834in}{2.637479in}}%
\pgfpathlineto{\pgfqpoint{2.389520in}{2.678313in}}%
\pgfpathlineto{\pgfqpoint{2.390360in}{2.764576in}}%
\pgfpathlineto{\pgfqpoint{2.390729in}{2.787678in}}%
\pgfpathlineto{\pgfqpoint{2.391325in}{2.644559in}}%
\pgfpathlineto{\pgfqpoint{2.393701in}{2.644559in}}%
\pgfpathlineto{\pgfqpoint{2.393706in}{2.747951in}}%
\pgfpathlineto{\pgfqpoint{2.393760in}{2.643592in}}%
\pgfpathlineto{\pgfqpoint{2.394437in}{2.681295in}}%
\pgfpathlineto{\pgfqpoint{2.395278in}{2.758755in}}%
\pgfpathlineto{\pgfqpoint{2.395649in}{2.782040in}}%
\pgfpathlineto{\pgfqpoint{2.396308in}{2.644559in}}%
\pgfpathlineto{\pgfqpoint{2.398646in}{2.644559in}}%
\pgfpathlineto{\pgfqpoint{2.398651in}{2.747071in}}%
\pgfpathlineto{\pgfqpoint{2.399389in}{2.698594in}}%
\pgfpathlineto{\pgfqpoint{2.399438in}{2.691571in}}%
\pgfpathlineto{\pgfqpoint{2.399884in}{2.729639in}}%
\pgfpathlineto{\pgfqpoint{2.399933in}{2.722662in}}%
\pgfpathlineto{\pgfqpoint{2.400573in}{2.772929in}}%
\pgfpathlineto{\pgfqpoint{2.400656in}{2.714686in}}%
\pgfpathlineto{\pgfqpoint{2.400665in}{2.749097in}}%
\pgfpathlineto{\pgfqpoint{2.401455in}{2.644559in}}%
\pgfpathlineto{\pgfqpoint{2.403592in}{2.644559in}}%
\pgfpathlineto{\pgfqpoint{2.403597in}{2.748715in}}%
\pgfpathlineto{\pgfqpoint{2.403675in}{2.633956in}}%
\pgfpathlineto{\pgfqpoint{2.404361in}{2.675221in}}%
\pgfpathlineto{\pgfqpoint{2.405202in}{2.768611in}}%
\pgfpathlineto{\pgfqpoint{2.405565in}{2.791400in}}%
\pgfpathlineto{\pgfqpoint{2.406136in}{2.644559in}}%
\pgfpathlineto{\pgfqpoint{2.408537in}{2.644559in}}%
\pgfpathlineto{\pgfqpoint{2.408542in}{2.746994in}}%
\pgfpathlineto{\pgfqpoint{2.409298in}{2.694450in}}%
\pgfpathlineto{\pgfqpoint{2.410511in}{2.771464in}}%
\pgfpathlineto{\pgfqpoint{2.410534in}{2.755790in}}%
\pgfpathlineto{\pgfqpoint{2.411347in}{2.644559in}}%
\pgfpathlineto{\pgfqpoint{2.413482in}{2.644559in}}%
\pgfpathlineto{\pgfqpoint{2.413488in}{2.747958in}}%
\pgfpathlineto{\pgfqpoint{2.413543in}{2.643313in}}%
\pgfpathlineto{\pgfqpoint{2.414220in}{2.681170in}}%
\pgfpathlineto{\pgfqpoint{2.415061in}{2.758943in}}%
\pgfpathlineto{\pgfqpoint{2.415456in}{2.783728in}}%
\pgfpathlineto{\pgfqpoint{2.416088in}{2.644559in}}%
\pgfpathlineto{\pgfqpoint{2.418428in}{2.644559in}}%
\pgfpathlineto{\pgfqpoint{2.418433in}{2.748243in}}%
\pgfpathlineto{\pgfqpoint{2.418500in}{2.643403in}}%
\pgfpathlineto{\pgfqpoint{2.419186in}{2.683500in}}%
\pgfpathlineto{\pgfqpoint{2.420026in}{2.758799in}}%
\pgfpathlineto{\pgfqpoint{2.420402in}{2.782346in}}%
\pgfpathlineto{\pgfqpoint{2.421090in}{2.644559in}}%
\pgfpathlineto{\pgfqpoint{2.423373in}{2.644559in}}%
\pgfpathlineto{\pgfqpoint{2.423379in}{2.747319in}}%
\pgfpathlineto{\pgfqpoint{2.424125in}{2.700954in}}%
\pgfpathlineto{\pgfqpoint{2.424174in}{2.691177in}}%
\pgfpathlineto{\pgfqpoint{2.424619in}{2.731995in}}%
\pgfpathlineto{\pgfqpoint{2.424669in}{2.722272in}}%
\pgfpathlineto{\pgfqpoint{2.425305in}{2.775015in}}%
\pgfpathlineto{\pgfqpoint{2.425383in}{2.715071in}}%
\pgfpathlineto{\pgfqpoint{2.425392in}{2.749628in}}%
\pgfpathlineto{\pgfqpoint{2.426182in}{2.644559in}}%
\pgfpathlineto{\pgfqpoint{2.428319in}{2.644559in}}%
\pgfpathlineto{\pgfqpoint{2.428324in}{2.747836in}}%
\pgfpathlineto{\pgfqpoint{2.428378in}{2.643504in}}%
\pgfpathlineto{\pgfqpoint{2.429054in}{2.681079in}}%
\pgfpathlineto{\pgfqpoint{2.429895in}{2.758295in}}%
\pgfpathlineto{\pgfqpoint{2.430267in}{2.781611in}}%
\pgfpathlineto{\pgfqpoint{2.430928in}{2.644559in}}%
\pgfpathlineto{\pgfqpoint{2.433264in}{2.644559in}}%
\pgfpathlineto{\pgfqpoint{2.433270in}{2.747535in}}%
\pgfpathlineto{\pgfqpoint{2.434035in}{2.705233in}}%
\pgfpathlineto{\pgfqpoint{2.434084in}{2.688668in}}%
\pgfpathlineto{\pgfqpoint{2.434628in}{2.742469in}}%
\pgfpathlineto{\pgfqpoint{2.434677in}{2.725994in}}%
\pgfpathlineto{\pgfqpoint{2.435205in}{2.778670in}}%
\pgfpathlineto{\pgfqpoint{2.435292in}{2.718386in}}%
\pgfpathlineto{\pgfqpoint{2.435306in}{2.725389in}}%
\pgfpathlineto{\pgfqpoint{2.436090in}{2.644559in}}%
\pgfpathlineto{\pgfqpoint{2.438210in}{2.644559in}}%
\pgfpathlineto{\pgfqpoint{2.438215in}{2.748138in}}%
\pgfpathlineto{\pgfqpoint{2.438275in}{2.641761in}}%
\pgfpathlineto{\pgfqpoint{2.438957in}{2.680735in}}%
\pgfpathlineto{\pgfqpoint{2.439797in}{2.760576in}}%
\pgfpathlineto{\pgfqpoint{2.440184in}{2.784785in}}%
\pgfpathlineto{\pgfqpoint{2.440809in}{2.644559in}}%
\pgfpathlineto{\pgfqpoint{2.443155in}{2.644559in}}%
\pgfpathlineto{\pgfqpoint{2.443161in}{2.747862in}}%
\pgfpathlineto{\pgfqpoint{2.443216in}{2.642902in}}%
\pgfpathlineto{\pgfqpoint{2.443894in}{2.680840in}}%
\pgfpathlineto{\pgfqpoint{2.444735in}{2.758800in}}%
\pgfpathlineto{\pgfqpoint{2.445129in}{2.783535in}}%
\pgfpathlineto{\pgfqpoint{2.445760in}{2.644559in}}%
\pgfpathlineto{\pgfqpoint{2.448101in}{2.644559in}}%
\pgfpathlineto{\pgfqpoint{2.448106in}{2.748231in}}%
\pgfpathlineto{\pgfqpoint{2.448854in}{2.685316in}}%
\pgfpathlineto{\pgfqpoint{2.449794in}{2.763078in}}%
\pgfpathlineto{\pgfqpoint{2.450074in}{2.780662in}}%
\pgfpathlineto{\pgfqpoint{2.450940in}{2.644559in}}%
\pgfpathlineto{\pgfqpoint{2.453046in}{2.644559in}}%
\pgfpathlineto{\pgfqpoint{2.453051in}{2.747926in}}%
\pgfpathlineto{\pgfqpoint{2.453107in}{2.642886in}}%
\pgfpathlineto{\pgfqpoint{2.453786in}{2.680931in}}%
\pgfpathlineto{\pgfqpoint{2.454626in}{2.759083in}}%
\pgfpathlineto{\pgfqpoint{2.455020in}{2.783762in}}%
\pgfpathlineto{\pgfqpoint{2.455651in}{2.644559in}}%
\pgfpathlineto{\pgfqpoint{2.457992in}{2.644559in}}%
\pgfpathlineto{\pgfqpoint{2.457997in}{2.748166in}}%
\pgfpathlineto{\pgfqpoint{2.458063in}{2.643749in}}%
\pgfpathlineto{\pgfqpoint{2.458749in}{2.683773in}}%
\pgfpathlineto{\pgfqpoint{2.459589in}{2.758188in}}%
\pgfpathlineto{\pgfqpoint{2.459965in}{2.781776in}}%
\pgfpathlineto{\pgfqpoint{2.460663in}{2.644559in}}%
\pgfpathlineto{\pgfqpoint{2.462937in}{2.644559in}}%
\pgfpathlineto{\pgfqpoint{2.462942in}{2.747638in}}%
\pgfpathlineto{\pgfqpoint{2.463706in}{2.705286in}}%
\pgfpathlineto{\pgfqpoint{2.463755in}{2.689087in}}%
\pgfpathlineto{\pgfqpoint{2.464299in}{2.742524in}}%
\pgfpathlineto{\pgfqpoint{2.464349in}{2.726412in}}%
\pgfpathlineto{\pgfqpoint{2.464877in}{2.778772in}}%
\pgfpathlineto{\pgfqpoint{2.464965in}{2.719228in}}%
\pgfpathlineto{\pgfqpoint{2.464979in}{2.725137in}}%
\pgfpathlineto{\pgfqpoint{2.465766in}{2.644559in}}%
\pgfpathlineto{\pgfqpoint{2.467882in}{2.644559in}}%
\pgfpathlineto{\pgfqpoint{2.467888in}{2.747824in}}%
\pgfpathlineto{\pgfqpoint{2.467941in}{2.643557in}}%
\pgfpathlineto{\pgfqpoint{2.468618in}{2.681084in}}%
\pgfpathlineto{\pgfqpoint{2.469458in}{2.758201in}}%
\pgfpathlineto{\pgfqpoint{2.469830in}{2.781530in}}%
\pgfpathlineto{\pgfqpoint{2.470494in}{2.644559in}}%
\pgfpathlineto{\pgfqpoint{2.472828in}{2.644559in}}%
\pgfpathlineto{\pgfqpoint{2.472833in}{2.747823in}}%
\pgfpathlineto{\pgfqpoint{2.472886in}{2.644114in}}%
\pgfpathlineto{\pgfqpoint{2.473561in}{2.681357in}}%
\pgfpathlineto{\pgfqpoint{2.474401in}{2.757850in}}%
\pgfpathlineto{\pgfqpoint{2.474775in}{2.781254in}}%
\pgfpathlineto{\pgfqpoint{2.475445in}{2.644559in}}%
\pgfpathlineto{\pgfqpoint{2.477773in}{2.644559in}}%
\pgfpathlineto{\pgfqpoint{2.477779in}{2.746902in}}%
\pgfpathlineto{\pgfqpoint{2.478533in}{2.694146in}}%
\pgfpathlineto{\pgfqpoint{2.479747in}{2.771600in}}%
\pgfpathlineto{\pgfqpoint{2.479770in}{2.755626in}}%
\pgfpathlineto{\pgfqpoint{2.480582in}{2.644559in}}%
\pgfpathlineto{\pgfqpoint{2.482719in}{2.644559in}}%
\pgfpathlineto{\pgfqpoint{2.482724in}{2.747785in}}%
\pgfpathlineto{\pgfqpoint{2.483445in}{2.682049in}}%
\pgfpathlineto{\pgfqpoint{2.484286in}{2.756683in}}%
\pgfpathlineto{\pgfqpoint{2.484662in}{2.780285in}}%
\pgfpathlineto{\pgfqpoint{2.485355in}{2.644559in}}%
\pgfpathlineto{\pgfqpoint{2.487664in}{2.644559in}}%
\pgfpathlineto{\pgfqpoint{2.487669in}{2.747070in}}%
\pgfpathlineto{\pgfqpoint{2.488425in}{2.694048in}}%
\pgfpathlineto{\pgfqpoint{2.489638in}{2.772060in}}%
\pgfpathlineto{\pgfqpoint{2.489638in}{2.747521in}}%
\pgfpathlineto{\pgfqpoint{2.489660in}{2.755253in}}%
\pgfpathlineto{\pgfqpoint{2.490411in}{2.644559in}}%
\pgfpathlineto{\pgfqpoint{2.492610in}{2.644559in}}%
\pgfpathlineto{\pgfqpoint{2.492615in}{2.748357in}}%
\pgfpathlineto{\pgfqpoint{2.492686in}{2.640008in}}%
\pgfpathlineto{\pgfqpoint{2.493371in}{2.680535in}}%
\pgfpathlineto{\pgfqpoint{2.494212in}{2.762021in}}%
\pgfpathlineto{\pgfqpoint{2.494584in}{2.785321in}}%
\pgfpathlineto{\pgfqpoint{2.495210in}{2.644559in}}%
\pgfpathlineto{\pgfqpoint{2.497555in}{2.644559in}}%
\pgfpathlineto{\pgfqpoint{2.497561in}{2.748385in}}%
\pgfpathlineto{\pgfqpoint{2.497633in}{2.638531in}}%
\pgfpathlineto{\pgfqpoint{2.498318in}{2.679238in}}%
\pgfpathlineto{\pgfqpoint{2.499159in}{2.763343in}}%
\pgfpathlineto{\pgfqpoint{2.499529in}{2.786529in}}%
\pgfpathlineto{\pgfqpoint{2.500138in}{2.644559in}}%
\pgfpathlineto{\pgfqpoint{2.502501in}{2.644559in}}%
\pgfpathlineto{\pgfqpoint{2.502506in}{2.748255in}}%
\pgfpathlineto{\pgfqpoint{2.502573in}{2.643515in}}%
\pgfpathlineto{\pgfqpoint{2.503258in}{2.683605in}}%
\pgfpathlineto{\pgfqpoint{2.504099in}{2.758677in}}%
\pgfpathlineto{\pgfqpoint{2.504474in}{2.782227in}}%
\pgfpathlineto{\pgfqpoint{2.505162in}{2.644559in}}%
\pgfpathlineto{\pgfqpoint{2.507446in}{2.644559in}}%
\pgfpathlineto{\pgfqpoint{2.507451in}{2.747838in}}%
\pgfpathlineto{\pgfqpoint{2.507507in}{2.642903in}}%
\pgfpathlineto{\pgfqpoint{2.508185in}{2.680834in}}%
\pgfpathlineto{\pgfqpoint{2.509025in}{2.758692in}}%
\pgfpathlineto{\pgfqpoint{2.509420in}{2.783431in}}%
\pgfpathlineto{\pgfqpoint{2.510052in}{2.644559in}}%
\pgfpathlineto{\pgfqpoint{2.512392in}{2.644559in}}%
\pgfpathlineto{\pgfqpoint{2.512397in}{2.747842in}}%
\pgfpathlineto{\pgfqpoint{2.512450in}{2.643979in}}%
\pgfpathlineto{\pgfqpoint{2.513125in}{2.681313in}}%
\pgfpathlineto{\pgfqpoint{2.513966in}{2.758005in}}%
\pgfpathlineto{\pgfqpoint{2.514339in}{2.781386in}}%
\pgfpathlineto{\pgfqpoint{2.515006in}{2.644559in}}%
\pgfpathlineto{\pgfqpoint{2.517337in}{2.644559in}}%
\pgfpathlineto{\pgfqpoint{2.517342in}{2.748108in}}%
\pgfpathlineto{\pgfqpoint{2.518092in}{2.684610in}}%
\pgfpathlineto{\pgfqpoint{2.519032in}{2.763253in}}%
\pgfpathlineto{\pgfqpoint{2.519311in}{2.780761in}}%
\pgfpathlineto{\pgfqpoint{2.520157in}{2.644559in}}%
\pgfpathlineto{\pgfqpoint{2.522282in}{2.644559in}}%
\pgfpathlineto{\pgfqpoint{2.522288in}{2.747557in}}%
\pgfpathlineto{\pgfqpoint{2.523045in}{2.703816in}}%
\pgfpathlineto{\pgfqpoint{2.523095in}{2.689919in}}%
\pgfpathlineto{\pgfqpoint{2.523639in}{2.741058in}}%
\pgfpathlineto{\pgfqpoint{2.523688in}{2.727241in}}%
\pgfpathlineto{\pgfqpoint{2.524220in}{2.777492in}}%
\pgfpathlineto{\pgfqpoint{2.524320in}{2.708897in}}%
\pgfpathlineto{\pgfqpoint{2.525099in}{2.644559in}}%
\pgfpathlineto{\pgfqpoint{2.527228in}{2.644559in}}%
\pgfpathlineto{\pgfqpoint{2.527233in}{2.747749in}}%
\pgfpathlineto{\pgfqpoint{2.527288in}{2.643122in}}%
\pgfpathlineto{\pgfqpoint{2.527965in}{2.680809in}}%
\pgfpathlineto{\pgfqpoint{2.528805in}{2.758185in}}%
\pgfpathlineto{\pgfqpoint{2.529177in}{2.781469in}}%
\pgfpathlineto{\pgfqpoint{2.529836in}{2.644559in}}%
\pgfpathlineto{\pgfqpoint{2.532173in}{2.644559in}}%
\pgfpathlineto{\pgfqpoint{2.532179in}{2.748794in}}%
\pgfpathlineto{\pgfqpoint{2.532243in}{2.639108in}}%
\pgfpathlineto{\pgfqpoint{2.532928in}{2.678694in}}%
\pgfpathlineto{\pgfqpoint{2.533769in}{2.762849in}}%
\pgfpathlineto{\pgfqpoint{2.534147in}{2.786577in}}%
\pgfpathlineto{\pgfqpoint{2.534744in}{2.644559in}}%
\pgfpathlineto{\pgfqpoint{2.537119in}{2.644559in}}%
\pgfpathlineto{\pgfqpoint{2.537124in}{2.748052in}}%
\pgfpathlineto{\pgfqpoint{2.537186in}{2.641106in}}%
\pgfpathlineto{\pgfqpoint{2.537868in}{2.680346in}}%
\pgfpathlineto{\pgfqpoint{2.538709in}{2.760644in}}%
\pgfpathlineto{\pgfqpoint{2.539093in}{2.784704in}}%
\pgfpathlineto{\pgfqpoint{2.539673in}{2.644559in}}%
\pgfpathlineto{\pgfqpoint{2.542064in}{2.644559in}}%
\pgfpathlineto{\pgfqpoint{2.542070in}{2.748338in}}%
\pgfpathlineto{\pgfqpoint{2.542142in}{2.638952in}}%
\pgfpathlineto{\pgfqpoint{2.542827in}{2.679604in}}%
\pgfpathlineto{\pgfqpoint{2.543668in}{2.762830in}}%
\pgfpathlineto{\pgfqpoint{2.544038in}{2.786049in}}%
\pgfpathlineto{\pgfqpoint{2.544650in}{2.644559in}}%
\pgfpathlineto{\pgfqpoint{2.547010in}{2.644559in}}%
\pgfpathlineto{\pgfqpoint{2.547015in}{2.748435in}}%
\pgfpathlineto{\pgfqpoint{2.547084in}{2.641462in}}%
\pgfpathlineto{\pgfqpoint{2.547770in}{2.681833in}}%
\pgfpathlineto{\pgfqpoint{2.548610in}{2.761127in}}%
\pgfpathlineto{\pgfqpoint{2.548984in}{2.784526in}}%
\pgfpathlineto{\pgfqpoint{2.549594in}{2.644559in}}%
\pgfpathlineto{\pgfqpoint{2.551955in}{2.644559in}}%
\pgfpathlineto{\pgfqpoint{2.551961in}{2.747735in}}%
\pgfpathlineto{\pgfqpoint{2.552013in}{2.643863in}}%
\pgfpathlineto{\pgfqpoint{2.552688in}{2.681137in}}%
\pgfpathlineto{\pgfqpoint{2.553529in}{2.757619in}}%
\pgfpathlineto{\pgfqpoint{2.553902in}{2.781013in}}%
\pgfpathlineto{\pgfqpoint{2.554572in}{2.644559in}}%
\pgfpathlineto{\pgfqpoint{2.556901in}{2.644559in}}%
\pgfpathlineto{\pgfqpoint{2.556906in}{2.748450in}}%
\pgfpathlineto{\pgfqpoint{2.556984in}{2.634395in}}%
\pgfpathlineto{\pgfqpoint{2.557669in}{2.675602in}}%
\pgfpathlineto{\pgfqpoint{2.558510in}{2.767363in}}%
\pgfpathlineto{\pgfqpoint{2.558874in}{2.790188in}}%
\pgfpathlineto{\pgfqpoint{2.559423in}{2.644559in}}%
\pgfpathlineto{\pgfqpoint{2.561846in}{2.644559in}}%
\pgfpathlineto{\pgfqpoint{2.561852in}{2.748286in}}%
\pgfpathlineto{\pgfqpoint{2.561922in}{2.639791in}}%
\pgfpathlineto{\pgfqpoint{2.562608in}{2.680336in}}%
\pgfpathlineto{\pgfqpoint{2.563448in}{2.761942in}}%
\pgfpathlineto{\pgfqpoint{2.563820in}{2.785228in}}%
\pgfpathlineto{\pgfqpoint{2.564442in}{2.644559in}}%
\pgfpathlineto{\pgfqpoint{2.566792in}{2.644559in}}%
\pgfpathlineto{\pgfqpoint{2.566797in}{2.748739in}}%
\pgfpathlineto{\pgfqpoint{2.566876in}{2.633753in}}%
\pgfpathlineto{\pgfqpoint{2.567561in}{2.675038in}}%
\pgfpathlineto{\pgfqpoint{2.568402in}{2.768874in}}%
\pgfpathlineto{\pgfqpoint{2.568765in}{2.791648in}}%
\pgfpathlineto{\pgfqpoint{2.569332in}{2.644559in}}%
\pgfpathlineto{\pgfqpoint{2.571737in}{2.644559in}}%
\pgfpathlineto{\pgfqpoint{2.571742in}{2.748075in}}%
\pgfpathlineto{\pgfqpoint{2.571811in}{2.641600in}}%
\pgfpathlineto{\pgfqpoint{2.572496in}{2.681895in}}%
\pgfpathlineto{\pgfqpoint{2.573337in}{2.759670in}}%
\pgfpathlineto{\pgfqpoint{2.573711in}{2.783100in}}%
\pgfpathlineto{\pgfqpoint{2.574370in}{2.644559in}}%
\pgfpathlineto{\pgfqpoint{2.576682in}{2.644559in}}%
\pgfpathlineto{\pgfqpoint{2.576688in}{2.748446in}}%
\pgfpathlineto{\pgfqpoint{2.576756in}{2.642135in}}%
\pgfpathlineto{\pgfqpoint{2.577442in}{2.682433in}}%
\pgfpathlineto{\pgfqpoint{2.578282in}{2.760614in}}%
\pgfpathlineto{\pgfqpoint{2.578656in}{2.784056in}}%
\pgfpathlineto{\pgfqpoint{2.579303in}{2.644559in}}%
\pgfpathlineto{\pgfqpoint{2.581628in}{2.644559in}}%
\pgfpathlineto{\pgfqpoint{2.581633in}{2.747004in}}%
\pgfpathlineto{\pgfqpoint{2.582367in}{2.697820in}}%
\pgfpathlineto{\pgfqpoint{2.582417in}{2.691756in}}%
\pgfpathlineto{\pgfqpoint{2.582862in}{2.728867in}}%
\pgfpathlineto{\pgfqpoint{2.582911in}{2.722846in}}%
\pgfpathlineto{\pgfqpoint{2.583554in}{2.772330in}}%
\pgfpathlineto{\pgfqpoint{2.583637in}{2.714178in}}%
\pgfpathlineto{\pgfqpoint{2.583646in}{2.749343in}}%
\pgfpathlineto{\pgfqpoint{2.584438in}{2.644559in}}%
\pgfpathlineto{\pgfqpoint{2.586573in}{2.644559in}}%
\pgfpathlineto{\pgfqpoint{2.586579in}{2.748213in}}%
\pgfpathlineto{\pgfqpoint{2.586645in}{2.644382in}}%
\pgfpathlineto{\pgfqpoint{2.587329in}{2.684308in}}%
\pgfpathlineto{\pgfqpoint{2.588170in}{2.757821in}}%
\pgfpathlineto{\pgfqpoint{2.588547in}{2.781473in}}%
\pgfpathlineto{\pgfqpoint{2.589260in}{2.644559in}}%
\pgfpathlineto{\pgfqpoint{2.591519in}{2.644559in}}%
\pgfpathlineto{\pgfqpoint{2.591524in}{2.746996in}}%
\pgfpathlineto{\pgfqpoint{2.592279in}{2.694067in}}%
\pgfpathlineto{\pgfqpoint{2.593493in}{2.772080in}}%
\pgfpathlineto{\pgfqpoint{2.593493in}{2.747567in}}%
\pgfpathlineto{\pgfqpoint{2.593515in}{2.755392in}}%
\pgfpathlineto{\pgfqpoint{2.594264in}{2.644559in}}%
\pgfpathlineto{\pgfqpoint{2.596464in}{2.644559in}}%
\pgfpathlineto{\pgfqpoint{2.596470in}{2.747268in}}%
\pgfpathlineto{\pgfqpoint{2.597215in}{2.700554in}}%
\pgfpathlineto{\pgfqpoint{2.597264in}{2.691148in}}%
\pgfpathlineto{\pgfqpoint{2.597709in}{2.731596in}}%
\pgfpathlineto{\pgfqpoint{2.597759in}{2.722242in}}%
\pgfpathlineto{\pgfqpoint{2.598395in}{2.774644in}}%
\pgfpathlineto{\pgfqpoint{2.598474in}{2.714924in}}%
\pgfpathlineto{\pgfqpoint{2.598483in}{2.749456in}}%
\pgfpathlineto{\pgfqpoint{2.599274in}{2.644559in}}%
\pgfpathlineto{\pgfqpoint{2.601410in}{2.644559in}}%
\pgfpathlineto{\pgfqpoint{2.601415in}{2.748375in}}%
\pgfpathlineto{\pgfqpoint{2.601492in}{2.635420in}}%
\pgfpathlineto{\pgfqpoint{2.602177in}{2.676526in}}%
\pgfpathlineto{\pgfqpoint{2.603018in}{2.766185in}}%
\pgfpathlineto{\pgfqpoint{2.603384in}{2.789086in}}%
\pgfpathlineto{\pgfqpoint{2.603965in}{2.644559in}}%
\pgfpathlineto{\pgfqpoint{2.606355in}{2.644559in}}%
\pgfpathlineto{\pgfqpoint{2.606361in}{2.748587in}}%
\pgfpathlineto{\pgfqpoint{2.606439in}{2.633580in}}%
\pgfpathlineto{\pgfqpoint{2.607125in}{2.674872in}}%
\pgfpathlineto{\pgfqpoint{2.607966in}{2.768545in}}%
\pgfpathlineto{\pgfqpoint{2.608329in}{2.791307in}}%
\pgfpathlineto{\pgfqpoint{2.608893in}{2.644559in}}%
\pgfpathlineto{\pgfqpoint{2.611301in}{2.644559in}}%
\pgfpathlineto{\pgfqpoint{2.611306in}{2.747141in}}%
\pgfpathlineto{\pgfqpoint{2.612063in}{2.694043in}}%
\pgfpathlineto{\pgfqpoint{2.613274in}{2.772189in}}%
\pgfpathlineto{\pgfqpoint{2.613275in}{2.747563in}}%
\pgfpathlineto{\pgfqpoint{2.613297in}{2.755148in}}%
\pgfpathlineto{\pgfqpoint{2.614047in}{2.644559in}}%
\pgfpathlineto{\pgfqpoint{2.616246in}{2.644559in}}%
\pgfpathlineto{\pgfqpoint{2.616252in}{2.748283in}}%
\pgfpathlineto{\pgfqpoint{2.616323in}{2.639754in}}%
\pgfpathlineto{\pgfqpoint{2.617008in}{2.680303in}}%
\pgfpathlineto{\pgfqpoint{2.617848in}{2.761964in}}%
\pgfpathlineto{\pgfqpoint{2.618220in}{2.785248in}}%
\pgfpathlineto{\pgfqpoint{2.618843in}{2.644559in}}%
\pgfpathlineto{\pgfqpoint{2.621192in}{2.644559in}}%
\pgfpathlineto{\pgfqpoint{2.621197in}{2.748009in}}%
\pgfpathlineto{\pgfqpoint{2.621256in}{2.641784in}}%
\pgfpathlineto{\pgfqpoint{2.621937in}{2.680536in}}%
\pgfpathlineto{\pgfqpoint{2.622778in}{2.760067in}}%
\pgfpathlineto{\pgfqpoint{2.623165in}{2.784386in}}%
\pgfpathlineto{\pgfqpoint{2.623790in}{2.644559in}}%
\pgfpathlineto{\pgfqpoint{2.626137in}{2.644559in}}%
\pgfpathlineto{\pgfqpoint{2.626142in}{2.747956in}}%
\pgfpathlineto{\pgfqpoint{2.626201in}{2.641962in}}%
\pgfpathlineto{\pgfqpoint{2.626881in}{2.680544in}}%
\pgfpathlineto{\pgfqpoint{2.627722in}{2.759772in}}%
\pgfpathlineto{\pgfqpoint{2.628111in}{2.784176in}}%
\pgfpathlineto{\pgfqpoint{2.628735in}{2.644559in}}%
\pgfpathlineto{\pgfqpoint{2.631082in}{2.644559in}}%
\pgfpathlineto{\pgfqpoint{2.631088in}{2.748272in}}%
\pgfpathlineto{\pgfqpoint{2.631156in}{2.642187in}}%
\pgfpathlineto{\pgfqpoint{2.631841in}{2.682445in}}%
\pgfpathlineto{\pgfqpoint{2.632682in}{2.759862in}}%
\pgfpathlineto{\pgfqpoint{2.633056in}{2.783320in}}%
\pgfpathlineto{\pgfqpoint{2.633711in}{2.644559in}}%
\pgfpathlineto{\pgfqpoint{2.636029in}{2.644651in}}%
\pgfpathlineto{\pgfqpoint{2.636034in}{2.757286in}}%
\pgfpathlineto{\pgfqpoint{2.636769in}{2.688875in}}%
\pgfpathlineto{\pgfqpoint{2.637808in}{2.763544in}}%
\pgfpathlineto{\pgfqpoint{2.638002in}{2.775708in}}%
\pgfpathlineto{\pgfqpoint{2.639509in}{2.644559in}}%
\pgfpathlineto{\pgfqpoint{2.640973in}{2.644559in}}%
\pgfpathlineto{\pgfqpoint{2.640979in}{2.746902in}}%
\pgfpathlineto{\pgfqpoint{2.641732in}{2.694139in}}%
\pgfpathlineto{\pgfqpoint{2.642947in}{2.771991in}}%
\pgfpathlineto{\pgfqpoint{2.642947in}{2.747584in}}%
\pgfpathlineto{\pgfqpoint{2.642970in}{2.755591in}}%
\pgfpathlineto{\pgfqpoint{2.643717in}{2.644559in}}%
\pgfpathlineto{\pgfqpoint{2.645918in}{2.644559in}}%
\pgfpathlineto{\pgfqpoint{2.645924in}{2.747822in}}%
\pgfpathlineto{\pgfqpoint{2.646003in}{2.639007in}}%
\pgfpathlineto{\pgfqpoint{2.646690in}{2.680417in}}%
\pgfpathlineto{\pgfqpoint{2.647531in}{2.763537in}}%
\pgfpathlineto{\pgfqpoint{2.647893in}{2.786232in}}%
\pgfpathlineto{\pgfqpoint{2.648515in}{2.644559in}}%
\pgfpathlineto{\pgfqpoint{2.650864in}{2.644559in}}%
\pgfpathlineto{\pgfqpoint{2.650869in}{2.747046in}}%
\pgfpathlineto{\pgfqpoint{2.651625in}{2.694052in}}%
\pgfpathlineto{\pgfqpoint{2.652838in}{2.772057in}}%
\pgfpathlineto{\pgfqpoint{2.652838in}{2.747530in}}%
\pgfpathlineto{\pgfqpoint{2.652860in}{2.755298in}}%
\pgfpathlineto{\pgfqpoint{2.653610in}{2.644559in}}%
\pgfpathlineto{\pgfqpoint{2.655809in}{2.644559in}}%
\pgfpathlineto{\pgfqpoint{2.655815in}{2.744068in}}%
\pgfpathlineto{\pgfqpoint{2.655913in}{2.636974in}}%
\pgfpathlineto{\pgfqpoint{2.656551in}{2.710995in}}%
\pgfpathlineto{\pgfqpoint{2.656600in}{2.679517in}}%
\pgfpathlineto{\pgfqpoint{2.657243in}{2.754426in}}%
\pgfpathlineto{\pgfqpoint{2.657293in}{2.723076in}}%
\pgfpathlineto{\pgfqpoint{2.657738in}{2.785446in}}%
\pgfpathlineto{\pgfqpoint{2.657875in}{2.644559in}}%
\pgfpathlineto{\pgfqpoint{2.657951in}{2.644559in}}%
\pgfpathlineto{\pgfqpoint{2.660755in}{2.644559in}}%
\pgfpathlineto{\pgfqpoint{2.660760in}{2.746959in}}%
\pgfpathlineto{\pgfqpoint{2.661515in}{2.694091in}}%
\pgfpathlineto{\pgfqpoint{2.662729in}{2.771991in}}%
\pgfpathlineto{\pgfqpoint{2.662729in}{2.747542in}}%
\pgfpathlineto{\pgfqpoint{2.662752in}{2.755462in}}%
\pgfpathlineto{\pgfqpoint{2.663500in}{2.644559in}}%
\pgfpathlineto{\pgfqpoint{2.665701in}{2.644561in}}%
\pgfpathlineto{\pgfqpoint{2.665706in}{2.791036in}}%
\pgfpathlineto{\pgfqpoint{2.665767in}{2.621645in}}%
\pgfpathlineto{\pgfqpoint{2.666449in}{2.660899in}}%
\pgfpathlineto{\pgfqpoint{2.667290in}{2.780135in}}%
\pgfpathlineto{\pgfqpoint{2.667674in}{2.804212in}}%
\pgfpathlineto{\pgfqpoint{2.668164in}{2.644559in}}%
\pgfpathlineto{\pgfqpoint{2.670646in}{2.644559in}}%
\pgfpathlineto{\pgfqpoint{2.670652in}{2.748170in}}%
\pgfpathlineto{\pgfqpoint{2.670722in}{2.639837in}}%
\pgfpathlineto{\pgfqpoint{2.671408in}{2.680362in}}%
\pgfpathlineto{\pgfqpoint{2.672248in}{2.761510in}}%
\pgfpathlineto{\pgfqpoint{2.672620in}{2.784804in}}%
\pgfpathlineto{\pgfqpoint{2.673249in}{2.644559in}}%
\pgfpathlineto{\pgfqpoint{2.675590in}{2.644559in}}%
\pgfpathlineto{\pgfqpoint{2.675596in}{2.755221in}}%
\pgfpathlineto{\pgfqpoint{2.675719in}{2.633798in}}%
\pgfpathlineto{\pgfqpoint{2.676362in}{2.717458in}}%
\pgfpathlineto{\pgfqpoint{2.676411in}{2.677121in}}%
\pgfpathlineto{\pgfqpoint{2.677054in}{2.760865in}}%
\pgfpathlineto{\pgfqpoint{2.677103in}{2.720703in}}%
\pgfpathlineto{\pgfqpoint{2.677532in}{2.790839in}}%
\pgfpathlineto{\pgfqpoint{2.677656in}{2.644559in}}%
\pgfpathlineto{\pgfqpoint{2.677752in}{2.644559in}}%
\pgfpathlineto{\pgfqpoint{2.680537in}{2.644559in}}%
\pgfpathlineto{\pgfqpoint{2.680542in}{2.746383in}}%
\pgfpathlineto{\pgfqpoint{2.681299in}{2.695146in}}%
\pgfpathlineto{\pgfqpoint{2.682511in}{2.770456in}}%
\pgfpathlineto{\pgfqpoint{2.682554in}{2.742780in}}%
\pgfpathlineto{\pgfqpoint{2.683357in}{2.644559in}}%
\pgfpathlineto{\pgfqpoint{2.685482in}{2.644559in}}%
\pgfpathlineto{\pgfqpoint{2.685488in}{2.747165in}}%
\pgfpathlineto{\pgfqpoint{2.686220in}{2.697584in}}%
\pgfpathlineto{\pgfqpoint{2.686269in}{2.692647in}}%
\pgfpathlineto{\pgfqpoint{2.686615in}{2.722423in}}%
\pgfpathlineto{\pgfqpoint{2.686665in}{2.717518in}}%
\pgfpathlineto{\pgfqpoint{2.687456in}{2.775201in}}%
\pgfpathlineto{\pgfqpoint{2.687485in}{2.752253in}}%
\pgfpathlineto{\pgfqpoint{2.688282in}{2.644559in}}%
\pgfpathlineto{\pgfqpoint{2.690428in}{2.644559in}}%
\pgfpathlineto{\pgfqpoint{2.690433in}{2.746341in}}%
\pgfpathlineto{\pgfqpoint{2.691176in}{2.699134in}}%
\pgfpathlineto{\pgfqpoint{2.691225in}{2.691036in}}%
\pgfpathlineto{\pgfqpoint{2.691670in}{2.730177in}}%
\pgfpathlineto{\pgfqpoint{2.691720in}{2.722129in}}%
\pgfpathlineto{\pgfqpoint{2.692357in}{2.773308in}}%
\pgfpathlineto{\pgfqpoint{2.692437in}{2.714365in}}%
\pgfpathlineto{\pgfqpoint{2.692446in}{2.748807in}}%
\pgfpathlineto{\pgfqpoint{2.693238in}{2.644559in}}%
\pgfpathlineto{\pgfqpoint{2.695372in}{2.644559in}}%
\pgfpathlineto{\pgfqpoint{2.695379in}{2.750643in}}%
\pgfpathlineto{\pgfqpoint{2.696142in}{2.704514in}}%
\pgfpathlineto{\pgfqpoint{2.696191in}{2.689740in}}%
\pgfpathlineto{\pgfqpoint{2.696735in}{2.741752in}}%
\pgfpathlineto{\pgfqpoint{2.696785in}{2.727065in}}%
\pgfpathlineto{\pgfqpoint{2.697313in}{2.778009in}}%
\pgfpathlineto{\pgfqpoint{2.697401in}{2.719281in}}%
\pgfpathlineto{\pgfqpoint{2.697415in}{2.724985in}}%
\pgfpathlineto{\pgfqpoint{2.698200in}{2.644559in}}%
\pgfpathlineto{\pgfqpoint{2.700318in}{2.644559in}}%
\pgfpathlineto{\pgfqpoint{2.700324in}{2.748235in}}%
\pgfpathlineto{\pgfqpoint{2.700423in}{2.636755in}}%
\pgfpathlineto{\pgfqpoint{2.701061in}{2.711301in}}%
\pgfpathlineto{\pgfqpoint{2.701110in}{2.679335in}}%
\pgfpathlineto{\pgfqpoint{2.701753in}{2.754730in}}%
\pgfpathlineto{\pgfqpoint{2.701803in}{2.722896in}}%
\pgfpathlineto{\pgfqpoint{2.702248in}{2.785749in}}%
\pgfpathlineto{\pgfqpoint{2.702384in}{2.644559in}}%
\pgfpathlineto{\pgfqpoint{2.702465in}{2.644559in}}%
\pgfpathlineto{\pgfqpoint{2.705264in}{2.644559in}}%
\pgfpathlineto{\pgfqpoint{2.705270in}{2.747936in}}%
\pgfpathlineto{\pgfqpoint{2.705330in}{2.641342in}}%
\pgfpathlineto{\pgfqpoint{2.706011in}{2.680247in}}%
\pgfpathlineto{\pgfqpoint{2.706852in}{2.760086in}}%
\pgfpathlineto{\pgfqpoint{2.707238in}{2.784323in}}%
\pgfpathlineto{\pgfqpoint{2.707855in}{2.644559in}}%
\pgfpathlineto{\pgfqpoint{2.710208in}{2.644559in}}%
\pgfpathlineto{\pgfqpoint{2.710215in}{2.755787in}}%
\pgfpathlineto{\pgfqpoint{2.710336in}{2.633757in}}%
\pgfpathlineto{\pgfqpoint{2.710979in}{2.717139in}}%
\pgfpathlineto{\pgfqpoint{2.711028in}{2.677061in}}%
\pgfpathlineto{\pgfqpoint{2.711671in}{2.760548in}}%
\pgfpathlineto{\pgfqpoint{2.711720in}{2.720642in}}%
\pgfpathlineto{\pgfqpoint{2.712150in}{2.790561in}}%
\pgfpathlineto{\pgfqpoint{2.712274in}{2.644559in}}%
\pgfpathlineto{\pgfqpoint{2.712377in}{2.644559in}}%
\pgfpathlineto{\pgfqpoint{2.715155in}{2.644559in}}%
\pgfpathlineto{\pgfqpoint{2.715161in}{2.747289in}}%
\pgfpathlineto{\pgfqpoint{2.715917in}{2.702720in}}%
\pgfpathlineto{\pgfqpoint{2.715967in}{2.689505in}}%
\pgfpathlineto{\pgfqpoint{2.716511in}{2.739963in}}%
\pgfpathlineto{\pgfqpoint{2.716560in}{2.726826in}}%
\pgfpathlineto{\pgfqpoint{2.717092in}{2.776431in}}%
\pgfpathlineto{\pgfqpoint{2.717193in}{2.706508in}}%
\pgfpathlineto{\pgfqpoint{2.717973in}{2.644559in}}%
\pgfpathlineto{\pgfqpoint{2.720100in}{2.644559in}}%
\pgfpathlineto{\pgfqpoint{2.720106in}{2.746712in}}%
\pgfpathlineto{\pgfqpoint{2.720829in}{2.695697in}}%
\pgfpathlineto{\pgfqpoint{2.720878in}{2.693018in}}%
\pgfpathlineto{\pgfqpoint{2.721125in}{2.714328in}}%
\pgfpathlineto{\pgfqpoint{2.721175in}{2.711669in}}%
\pgfpathlineto{\pgfqpoint{2.722074in}{2.773908in}}%
\pgfpathlineto{\pgfqpoint{2.722075in}{2.747525in}}%
\pgfpathlineto{\pgfqpoint{2.722095in}{2.754354in}}%
\pgfpathlineto{\pgfqpoint{2.722843in}{2.644559in}}%
\pgfpathlineto{\pgfqpoint{2.725046in}{2.644559in}}%
\pgfpathlineto{\pgfqpoint{2.725051in}{2.747122in}}%
\pgfpathlineto{\pgfqpoint{2.725808in}{2.694049in}}%
\pgfpathlineto{\pgfqpoint{2.727020in}{2.772110in}}%
\pgfpathlineto{\pgfqpoint{2.727020in}{2.747530in}}%
\pgfpathlineto{\pgfqpoint{2.727042in}{2.755181in}}%
\pgfpathlineto{\pgfqpoint{2.727792in}{2.644559in}}%
\pgfpathlineto{\pgfqpoint{2.729991in}{2.644559in}}%
\pgfpathlineto{\pgfqpoint{2.729997in}{2.745806in}}%
\pgfpathlineto{\pgfqpoint{2.730754in}{2.694692in}}%
\pgfpathlineto{\pgfqpoint{2.731965in}{2.770935in}}%
\pgfpathlineto{\pgfqpoint{2.732008in}{2.737925in}}%
\pgfpathlineto{\pgfqpoint{2.732827in}{2.644559in}}%
\pgfpathlineto{\pgfqpoint{2.734937in}{2.644559in}}%
\pgfpathlineto{\pgfqpoint{2.734942in}{2.748032in}}%
\pgfpathlineto{\pgfqpoint{2.735000in}{2.642315in}}%
\pgfpathlineto{\pgfqpoint{2.735680in}{2.680817in}}%
\pgfpathlineto{\pgfqpoint{2.736521in}{2.759851in}}%
\pgfpathlineto{\pgfqpoint{2.736911in}{2.784298in}}%
\pgfpathlineto{\pgfqpoint{2.737539in}{2.644559in}}%
\pgfpathlineto{\pgfqpoint{2.739883in}{2.644563in}}%
\pgfpathlineto{\pgfqpoint{2.739888in}{2.784267in}}%
\pgfpathlineto{\pgfqpoint{2.739941in}{2.644039in}}%
\pgfpathlineto{\pgfqpoint{2.740618in}{2.681575in}}%
\pgfpathlineto{\pgfqpoint{2.741458in}{2.757283in}}%
\pgfpathlineto{\pgfqpoint{2.741830in}{2.780622in}}%
\pgfpathlineto{\pgfqpoint{2.742512in}{2.644559in}}%
\pgfpathlineto{\pgfqpoint{2.744828in}{2.644559in}}%
\pgfpathlineto{\pgfqpoint{2.744833in}{2.747805in}}%
\pgfpathlineto{\pgfqpoint{2.745555in}{2.682098in}}%
\pgfpathlineto{\pgfqpoint{2.746395in}{2.756781in}}%
\pgfpathlineto{\pgfqpoint{2.746772in}{2.780379in}}%
\pgfpathlineto{\pgfqpoint{2.747464in}{2.644559in}}%
\pgfpathlineto{\pgfqpoint{2.749774in}{2.644559in}}%
\pgfpathlineto{\pgfqpoint{2.749778in}{2.743498in}}%
\pgfpathlineto{\pgfqpoint{2.750537in}{2.692803in}}%
\pgfpathlineto{\pgfqpoint{2.751747in}{2.773127in}}%
\pgfpathlineto{\pgfqpoint{2.751747in}{2.759821in}}%
\pgfpathlineto{\pgfqpoint{2.752517in}{2.644559in}}%
\pgfpathlineto{\pgfqpoint{2.754719in}{2.644559in}}%
\pgfpathlineto{\pgfqpoint{2.754724in}{2.748507in}}%
\pgfpathlineto{\pgfqpoint{2.754801in}{2.635816in}}%
\pgfpathlineto{\pgfqpoint{2.755486in}{2.676895in}}%
\pgfpathlineto{\pgfqpoint{2.756327in}{2.766249in}}%
\pgfpathlineto{\pgfqpoint{2.756693in}{2.789175in}}%
\pgfpathlineto{\pgfqpoint{2.757273in}{2.644559in}}%
\pgfpathlineto{\pgfqpoint{2.759663in}{2.644559in}}%
\pgfpathlineto{\pgfqpoint{2.759670in}{2.747764in}}%
\pgfpathlineto{\pgfqpoint{2.760426in}{2.694007in}}%
\pgfpathlineto{\pgfqpoint{2.761638in}{2.772015in}}%
\pgfpathlineto{\pgfqpoint{2.761638in}{2.747459in}}%
\pgfpathlineto{\pgfqpoint{2.761660in}{2.755097in}}%
\pgfpathlineto{\pgfqpoint{2.762409in}{2.644559in}}%
\pgfpathlineto{\pgfqpoint{2.764609in}{2.644559in}}%
\pgfpathlineto{\pgfqpoint{2.764615in}{2.748923in}}%
\pgfpathlineto{\pgfqpoint{2.765337in}{2.701146in}}%
\pgfpathlineto{\pgfqpoint{2.765386in}{2.687537in}}%
\pgfpathlineto{\pgfqpoint{2.765930in}{2.738405in}}%
\pgfpathlineto{\pgfqpoint{2.765980in}{2.724841in}}%
\pgfpathlineto{\pgfqpoint{2.766584in}{2.779403in}}%
\pgfpathlineto{\pgfqpoint{2.766638in}{2.716712in}}%
\pgfpathlineto{\pgfqpoint{2.766652in}{2.727517in}}%
\pgfpathlineto{\pgfqpoint{2.767439in}{2.644559in}}%
\pgfpathlineto{\pgfqpoint{2.769554in}{2.644559in}}%
\pgfpathlineto{\pgfqpoint{2.769560in}{2.756913in}}%
\pgfpathlineto{\pgfqpoint{2.769619in}{2.640964in}}%
\pgfpathlineto{\pgfqpoint{2.770300in}{2.679608in}}%
\pgfpathlineto{\pgfqpoint{2.771141in}{2.760848in}}%
\pgfpathlineto{\pgfqpoint{2.771529in}{2.785204in}}%
\pgfpathlineto{\pgfqpoint{2.772137in}{2.644559in}}%
\pgfpathlineto{\pgfqpoint{2.774499in}{2.644559in}}%
\pgfpathlineto{\pgfqpoint{2.774506in}{2.757505in}}%
\pgfpathlineto{\pgfqpoint{2.774559in}{2.628213in}}%
\pgfpathlineto{\pgfqpoint{2.775236in}{2.665810in}}%
\pgfpathlineto{\pgfqpoint{2.776076in}{2.773154in}}%
\pgfpathlineto{\pgfqpoint{2.776449in}{2.796454in}}%
\pgfpathlineto{\pgfqpoint{2.776970in}{2.644559in}}%
\pgfpathlineto{\pgfqpoint{2.779446in}{2.644559in}}%
\pgfpathlineto{\pgfqpoint{2.779451in}{2.748353in}}%
\pgfpathlineto{\pgfqpoint{2.780184in}{2.697241in}}%
\pgfpathlineto{\pgfqpoint{2.780234in}{2.692693in}}%
\pgfpathlineto{\pgfqpoint{2.780580in}{2.722080in}}%
\pgfpathlineto{\pgfqpoint{2.780629in}{2.717563in}}%
\pgfpathlineto{\pgfqpoint{2.781371in}{2.771755in}}%
\pgfpathlineto{\pgfqpoint{2.781448in}{2.757171in}}%
\pgfpathlineto{\pgfqpoint{2.782235in}{2.644559in}}%
\pgfpathlineto{\pgfqpoint{2.784391in}{2.644559in}}%
\pgfpathlineto{\pgfqpoint{2.784397in}{2.746624in}}%
\pgfpathlineto{\pgfqpoint{2.785162in}{2.705403in}}%
\pgfpathlineto{\pgfqpoint{2.785211in}{2.688785in}}%
\pgfpathlineto{\pgfqpoint{2.785755in}{2.742640in}}%
\pgfpathlineto{\pgfqpoint{2.785805in}{2.726111in}}%
\pgfpathlineto{\pgfqpoint{2.786332in}{2.778847in}}%
\pgfpathlineto{\pgfqpoint{2.786419in}{2.718776in}}%
\pgfpathlineto{\pgfqpoint{2.786433in}{2.725278in}}%
\pgfpathlineto{\pgfqpoint{2.787219in}{2.644559in}}%
\pgfpathlineto{\pgfqpoint{2.789337in}{2.644559in}}%
\pgfpathlineto{\pgfqpoint{2.789342in}{2.747238in}}%
\pgfpathlineto{\pgfqpoint{2.790085in}{2.700137in}}%
\pgfpathlineto{\pgfqpoint{2.790135in}{2.691311in}}%
\pgfpathlineto{\pgfqpoint{2.790580in}{2.731180in}}%
\pgfpathlineto{\pgfqpoint{2.790629in}{2.722404in}}%
\pgfpathlineto{\pgfqpoint{2.791267in}{2.774299in}}%
\pgfpathlineto{\pgfqpoint{2.791347in}{2.714948in}}%
\pgfpathlineto{\pgfqpoint{2.791356in}{2.749454in}}%
\pgfpathlineto{\pgfqpoint{2.792143in}{2.644559in}}%
\pgfpathlineto{\pgfqpoint{2.794281in}{2.644559in}}%
\pgfpathlineto{\pgfqpoint{2.794287in}{2.768210in}}%
\pgfpathlineto{\pgfqpoint{2.794419in}{2.625666in}}%
\pgfpathlineto{\pgfqpoint{2.795062in}{2.726113in}}%
\pgfpathlineto{\pgfqpoint{2.795111in}{2.669111in}}%
\pgfpathlineto{\pgfqpoint{2.795754in}{2.769485in}}%
\pgfpathlineto{\pgfqpoint{2.795803in}{2.712729in}}%
\pgfpathlineto{\pgfqpoint{2.796228in}{2.799148in}}%
\pgfpathlineto{\pgfqpoint{2.796339in}{2.644559in}}%
\pgfpathlineto{\pgfqpoint{2.796503in}{2.644559in}}%
\pgfpathlineto{\pgfqpoint{2.799228in}{2.644559in}}%
\pgfpathlineto{\pgfqpoint{2.799233in}{2.747263in}}%
\pgfpathlineto{\pgfqpoint{2.799977in}{2.700316in}}%
\pgfpathlineto{\pgfqpoint{2.800026in}{2.691323in}}%
\pgfpathlineto{\pgfqpoint{2.800471in}{2.731359in}}%
\pgfpathlineto{\pgfqpoint{2.800521in}{2.722416in}}%
\pgfpathlineto{\pgfqpoint{2.801158in}{2.774462in}}%
\pgfpathlineto{\pgfqpoint{2.801237in}{2.715006in}}%
\pgfpathlineto{\pgfqpoint{2.801246in}{2.749524in}}%
\pgfpathlineto{\pgfqpoint{2.802035in}{2.644559in}}%
\pgfpathlineto{\pgfqpoint{2.804172in}{2.644559in}}%
\pgfpathlineto{\pgfqpoint{2.804178in}{2.759016in}}%
\pgfpathlineto{\pgfqpoint{2.804312in}{2.624721in}}%
\pgfpathlineto{\pgfqpoint{2.804905in}{2.661956in}}%
\pgfpathlineto{\pgfqpoint{2.805746in}{2.776185in}}%
\pgfpathlineto{\pgfqpoint{2.806120in}{2.799575in}}%
\pgfpathlineto{\pgfqpoint{2.806626in}{2.644559in}}%
\pgfpathlineto{\pgfqpoint{2.809119in}{2.644559in}}%
\pgfpathlineto{\pgfqpoint{2.809124in}{2.748052in}}%
\pgfpathlineto{\pgfqpoint{2.809180in}{2.643161in}}%
\pgfpathlineto{\pgfqpoint{2.809858in}{2.681221in}}%
\pgfpathlineto{\pgfqpoint{2.810699in}{2.759370in}}%
\pgfpathlineto{\pgfqpoint{2.811093in}{2.784047in}}%
\pgfpathlineto{\pgfqpoint{2.811724in}{2.644559in}}%
\pgfpathlineto{\pgfqpoint{2.814063in}{2.644559in}}%
\pgfpathlineto{\pgfqpoint{2.814069in}{2.757326in}}%
\pgfpathlineto{\pgfqpoint{2.814198in}{2.629797in}}%
\pgfpathlineto{\pgfqpoint{2.814841in}{2.721140in}}%
\pgfpathlineto{\pgfqpoint{2.814890in}{2.673199in}}%
\pgfpathlineto{\pgfqpoint{2.815533in}{2.764529in}}%
\pgfpathlineto{\pgfqpoint{2.815583in}{2.716800in}}%
\pgfpathlineto{\pgfqpoint{2.816008in}{2.794286in}}%
\pgfpathlineto{\pgfqpoint{2.816129in}{2.644559in}}%
\pgfpathlineto{\pgfqpoint{2.816283in}{2.644559in}}%
\pgfpathlineto{\pgfqpoint{2.819010in}{2.644559in}}%
\pgfpathlineto{\pgfqpoint{2.819015in}{2.748132in}}%
\pgfpathlineto{\pgfqpoint{2.819081in}{2.644321in}}%
\pgfpathlineto{\pgfqpoint{2.819766in}{2.684234in}}%
\pgfpathlineto{\pgfqpoint{2.820606in}{2.757554in}}%
\pgfpathlineto{\pgfqpoint{2.820984in}{2.781211in}}%
\pgfpathlineto{\pgfqpoint{2.821699in}{2.644559in}}%
\pgfpathlineto{\pgfqpoint{2.823956in}{2.644591in}}%
\pgfpathlineto{\pgfqpoint{2.823961in}{2.770059in}}%
\pgfpathlineto{\pgfqpoint{2.824020in}{2.622446in}}%
\pgfpathlineto{\pgfqpoint{2.824700in}{2.661185in}}%
\pgfpathlineto{\pgfqpoint{2.825541in}{2.778899in}}%
\pgfpathlineto{\pgfqpoint{2.825929in}{2.803200in}}%
\pgfpathlineto{\pgfqpoint{2.826418in}{2.644559in}}%
\pgfpathlineto{\pgfqpoint{2.828901in}{2.644559in}}%
\pgfpathlineto{\pgfqpoint{2.828906in}{2.748168in}}%
\pgfpathlineto{\pgfqpoint{2.828973in}{2.643431in}}%
\pgfpathlineto{\pgfqpoint{2.829658in}{2.683515in}}%
\pgfpathlineto{\pgfqpoint{2.830499in}{2.758409in}}%
\pgfpathlineto{\pgfqpoint{2.830874in}{2.781959in}}%
\pgfpathlineto{\pgfqpoint{2.831567in}{2.644559in}}%
\pgfpathlineto{\pgfqpoint{2.833846in}{2.644559in}}%
\pgfpathlineto{\pgfqpoint{2.833852in}{2.744934in}}%
\pgfpathlineto{\pgfqpoint{2.834599in}{2.695052in}}%
\pgfpathlineto{\pgfqpoint{2.835242in}{2.733789in}}%
\pgfpathlineto{\pgfqpoint{2.835820in}{2.770113in}}%
\pgfpathlineto{\pgfqpoint{2.835847in}{2.757229in}}%
\pgfpathlineto{\pgfqpoint{2.836638in}{2.644559in}}%
\pgfpathlineto{\pgfqpoint{2.838790in}{2.644559in}}%
\pgfpathlineto{\pgfqpoint{2.838797in}{2.753374in}}%
\pgfpathlineto{\pgfqpoint{2.839542in}{2.692143in}}%
\pgfpathlineto{\pgfqpoint{2.840765in}{2.772728in}}%
\pgfpathlineto{\pgfqpoint{2.840765in}{2.759996in}}%
\pgfpathlineto{\pgfqpoint{2.841538in}{2.644559in}}%
\pgfpathlineto{\pgfqpoint{2.843737in}{2.644559in}}%
\pgfpathlineto{\pgfqpoint{2.843742in}{2.745355in}}%
\pgfpathlineto{\pgfqpoint{2.844467in}{2.693559in}}%
\pgfpathlineto{\pgfqpoint{2.845307in}{2.744891in}}%
\pgfpathlineto{\pgfqpoint{2.845711in}{2.771732in}}%
\pgfpathlineto{\pgfqpoint{2.845746in}{2.716858in}}%
\pgfpathlineto{\pgfqpoint{2.845755in}{2.747214in}}%
\pgfpathlineto{\pgfqpoint{2.846550in}{2.644559in}}%
\pgfpathlineto{\pgfqpoint{2.848683in}{2.644574in}}%
\pgfpathlineto{\pgfqpoint{2.848688in}{2.773370in}}%
\pgfpathlineto{\pgfqpoint{2.849420in}{2.686291in}}%
\pgfpathlineto{\pgfqpoint{2.850359in}{2.759827in}}%
\pgfpathlineto{\pgfqpoint{2.850631in}{2.776911in}}%
\pgfpathlineto{\pgfqpoint{2.851630in}{2.644559in}}%
\pgfpathlineto{\pgfqpoint{2.853627in}{2.644559in}}%
\pgfpathlineto{\pgfqpoint{2.853633in}{2.754218in}}%
\pgfpathlineto{\pgfqpoint{2.853748in}{2.633636in}}%
\pgfpathlineto{\pgfqpoint{2.854391in}{2.716234in}}%
\pgfpathlineto{\pgfqpoint{2.854440in}{2.676861in}}%
\pgfpathlineto{\pgfqpoint{2.855083in}{2.759643in}}%
\pgfpathlineto{\pgfqpoint{2.855133in}{2.720441in}}%
\pgfpathlineto{\pgfqpoint{2.855565in}{2.789843in}}%
\pgfpathlineto{\pgfqpoint{2.855692in}{2.644559in}}%
\pgfpathlineto{\pgfqpoint{2.855798in}{2.644559in}}%
\pgfpathlineto{\pgfqpoint{2.858573in}{2.644559in}}%
\pgfpathlineto{\pgfqpoint{2.858579in}{2.747619in}}%
\pgfpathlineto{\pgfqpoint{2.859339in}{2.703482in}}%
\pgfpathlineto{\pgfqpoint{2.859389in}{2.689362in}}%
\pgfpathlineto{\pgfqpoint{2.859933in}{2.740722in}}%
\pgfpathlineto{\pgfqpoint{2.859982in}{2.726685in}}%
\pgfpathlineto{\pgfqpoint{2.860512in}{2.777061in}}%
\pgfpathlineto{\pgfqpoint{2.860601in}{2.719650in}}%
\pgfpathlineto{\pgfqpoint{2.860615in}{2.723441in}}%
\pgfpathlineto{\pgfqpoint{2.861402in}{2.644559in}}%
\pgfpathlineto{\pgfqpoint{2.863519in}{2.644559in}}%
\pgfpathlineto{\pgfqpoint{2.863524in}{2.748147in}}%
\pgfpathlineto{\pgfqpoint{2.863585in}{2.641659in}}%
\pgfpathlineto{\pgfqpoint{2.864266in}{2.680711in}}%
\pgfpathlineto{\pgfqpoint{2.865107in}{2.760669in}}%
\pgfpathlineto{\pgfqpoint{2.865493in}{2.784836in}}%
\pgfpathlineto{\pgfqpoint{2.866100in}{2.644559in}}%
\pgfpathlineto{\pgfqpoint{2.868464in}{2.644559in}}%
\pgfpathlineto{\pgfqpoint{2.868470in}{2.744447in}}%
\pgfpathlineto{\pgfqpoint{2.869231in}{2.696490in}}%
\pgfpathlineto{\pgfqpoint{2.869280in}{2.696412in}}%
\pgfpathlineto{\pgfqpoint{2.870403in}{2.770145in}}%
\pgfpathlineto{\pgfqpoint{2.870438in}{2.746968in}}%
\pgfpathlineto{\pgfqpoint{2.870465in}{2.757030in}}%
\pgfpathlineto{\pgfqpoint{2.871211in}{2.644559in}}%
\pgfpathlineto{\pgfqpoint{2.873410in}{2.644559in}}%
\pgfpathlineto{\pgfqpoint{2.873415in}{2.748660in}}%
\pgfpathlineto{\pgfqpoint{2.873496in}{2.632050in}}%
\pgfpathlineto{\pgfqpoint{2.874181in}{2.673504in}}%
\pgfpathlineto{\pgfqpoint{2.875022in}{2.770185in}}%
\pgfpathlineto{\pgfqpoint{2.875384in}{2.792814in}}%
\pgfpathlineto{\pgfqpoint{2.875940in}{2.644559in}}%
\pgfpathlineto{\pgfqpoint{2.878355in}{2.644559in}}%
\pgfpathlineto{\pgfqpoint{2.878361in}{2.749003in}}%
\pgfpathlineto{\pgfqpoint{2.878414in}{2.643792in}}%
\pgfpathlineto{\pgfqpoint{2.879089in}{2.681125in}}%
\pgfpathlineto{\pgfqpoint{2.879930in}{2.757480in}}%
\pgfpathlineto{\pgfqpoint{2.880303in}{2.780845in}}%
\pgfpathlineto{\pgfqpoint{2.880972in}{2.644559in}}%
\pgfpathlineto{\pgfqpoint{2.883301in}{2.644559in}}%
\pgfpathlineto{\pgfqpoint{2.883306in}{2.748027in}}%
\pgfpathlineto{\pgfqpoint{2.883369in}{2.640731in}}%
\pgfpathlineto{\pgfqpoint{2.884052in}{2.680152in}}%
\pgfpathlineto{\pgfqpoint{2.884892in}{2.760824in}}%
\pgfpathlineto{\pgfqpoint{2.885274in}{2.784782in}}%
\pgfpathlineto{\pgfqpoint{2.885870in}{2.644559in}}%
\pgfpathlineto{\pgfqpoint{2.888247in}{2.644641in}}%
\pgfpathlineto{\pgfqpoint{2.888252in}{2.755920in}}%
\pgfpathlineto{\pgfqpoint{2.889005in}{2.695649in}}%
\pgfpathlineto{\pgfqpoint{2.889549in}{2.727961in}}%
\pgfpathlineto{\pgfqpoint{2.890220in}{2.770123in}}%
\pgfpathlineto{\pgfqpoint{2.890247in}{2.757453in}}%
\pgfpathlineto{\pgfqpoint{2.891039in}{2.644559in}}%
\pgfpathlineto{\pgfqpoint{2.893192in}{2.644559in}}%
\pgfpathlineto{\pgfqpoint{2.893197in}{2.748136in}}%
\pgfpathlineto{\pgfqpoint{2.893267in}{2.640535in}}%
\pgfpathlineto{\pgfqpoint{2.893952in}{2.680966in}}%
\pgfpathlineto{\pgfqpoint{2.894793in}{2.760813in}}%
\pgfpathlineto{\pgfqpoint{2.895165in}{2.784164in}}%
\pgfpathlineto{\pgfqpoint{2.895800in}{2.644559in}}%
\pgfpathlineto{\pgfqpoint{2.898136in}{2.644559in}}%
\pgfpathlineto{\pgfqpoint{2.898143in}{2.746596in}}%
\pgfpathlineto{\pgfqpoint{2.898898in}{2.694155in}}%
\pgfpathlineto{\pgfqpoint{2.900111in}{2.771327in}}%
\pgfpathlineto{\pgfqpoint{2.900134in}{2.755507in}}%
\pgfpathlineto{\pgfqpoint{2.900946in}{2.644559in}}%
\pgfpathlineto{\pgfqpoint{2.903082in}{2.644559in}}%
\pgfpathlineto{\pgfqpoint{2.903088in}{2.747602in}}%
\pgfpathlineto{\pgfqpoint{2.903849in}{2.704628in}}%
\pgfpathlineto{\pgfqpoint{2.903899in}{2.689346in}}%
\pgfpathlineto{\pgfqpoint{2.904443in}{2.741866in}}%
\pgfpathlineto{\pgfqpoint{2.904492in}{2.726670in}}%
\pgfpathlineto{\pgfqpoint{2.905022in}{2.778177in}}%
\pgfpathlineto{\pgfqpoint{2.905110in}{2.719927in}}%
\pgfpathlineto{\pgfqpoint{2.905124in}{2.724196in}}%
\pgfpathlineto{\pgfqpoint{2.905910in}{2.644559in}}%
\pgfpathlineto{\pgfqpoint{2.908027in}{2.644559in}}%
\pgfpathlineto{\pgfqpoint{2.908033in}{2.747920in}}%
\pgfpathlineto{\pgfqpoint{2.908135in}{2.635976in}}%
\pgfpathlineto{\pgfqpoint{2.908775in}{2.712264in}}%
\pgfpathlineto{\pgfqpoint{2.908824in}{2.678728in}}%
\pgfpathlineto{\pgfqpoint{2.909467in}{2.755689in}}%
\pgfpathlineto{\pgfqpoint{2.909516in}{2.722292in}}%
\pgfpathlineto{\pgfqpoint{2.909957in}{2.786416in}}%
\pgfpathlineto{\pgfqpoint{2.910093in}{2.644559in}}%
\pgfpathlineto{\pgfqpoint{2.910182in}{2.644559in}}%
\pgfpathlineto{\pgfqpoint{2.912973in}{2.644559in}}%
\pgfpathlineto{\pgfqpoint{2.912979in}{2.747277in}}%
\pgfpathlineto{\pgfqpoint{2.913727in}{2.701088in}}%
\pgfpathlineto{\pgfqpoint{2.913777in}{2.690770in}}%
\pgfpathlineto{\pgfqpoint{2.914222in}{2.732128in}}%
\pgfpathlineto{\pgfqpoint{2.914271in}{2.721866in}}%
\pgfpathlineto{\pgfqpoint{2.914906in}{2.775070in}}%
\pgfpathlineto{\pgfqpoint{2.914986in}{2.720339in}}%
\pgfpathlineto{\pgfqpoint{2.914998in}{2.744811in}}%
\pgfpathlineto{\pgfqpoint{2.915789in}{2.644559in}}%
\pgfpathlineto{\pgfqpoint{2.917917in}{2.644559in}}%
\pgfpathlineto{\pgfqpoint{2.917924in}{2.757773in}}%
\pgfpathlineto{\pgfqpoint{2.917982in}{2.624763in}}%
\pgfpathlineto{\pgfqpoint{2.918662in}{2.663364in}}%
\pgfpathlineto{\pgfqpoint{2.919503in}{2.776372in}}%
\pgfpathlineto{\pgfqpoint{2.919893in}{2.800767in}}%
\pgfpathlineto{\pgfqpoint{2.920386in}{2.644559in}}%
\pgfpathlineto{\pgfqpoint{2.922865in}{2.644569in}}%
\pgfpathlineto{\pgfqpoint{2.922870in}{2.779378in}}%
\pgfpathlineto{\pgfqpoint{2.923629in}{2.696143in}}%
\pgfpathlineto{\pgfqpoint{2.924470in}{2.747499in}}%
\pgfpathlineto{\pgfqpoint{2.924838in}{2.770604in}}%
\pgfpathlineto{\pgfqpoint{2.924872in}{2.724463in}}%
\pgfpathlineto{\pgfqpoint{2.924880in}{2.740508in}}%
\pgfpathlineto{\pgfqpoint{2.925694in}{2.644559in}}%
\pgfpathlineto{\pgfqpoint{2.927811in}{2.644573in}}%
\pgfpathlineto{\pgfqpoint{2.927815in}{2.772235in}}%
\pgfpathlineto{\pgfqpoint{2.927877in}{2.642114in}}%
\pgfpathlineto{\pgfqpoint{2.928562in}{2.681603in}}%
\pgfpathlineto{\pgfqpoint{2.929403in}{2.759353in}}%
\pgfpathlineto{\pgfqpoint{2.929784in}{2.783224in}}%
\pgfpathlineto{\pgfqpoint{2.930431in}{2.644559in}}%
\pgfpathlineto{\pgfqpoint{2.932754in}{2.644559in}}%
\pgfpathlineto{\pgfqpoint{2.932760in}{2.754984in}}%
\pgfpathlineto{\pgfqpoint{2.932887in}{2.626031in}}%
\pgfpathlineto{\pgfqpoint{2.933530in}{2.725597in}}%
\pgfpathlineto{\pgfqpoint{2.933580in}{2.669436in}}%
\pgfpathlineto{\pgfqpoint{2.934223in}{2.768973in}}%
\pgfpathlineto{\pgfqpoint{2.934272in}{2.713049in}}%
\pgfpathlineto{\pgfqpoint{2.934698in}{2.798765in}}%
\pgfpathlineto{\pgfqpoint{2.934812in}{2.644559in}}%
\pgfpathlineto{\pgfqpoint{2.934984in}{2.644559in}}%
\pgfpathlineto{\pgfqpoint{2.937699in}{2.644559in}}%
\pgfpathlineto{\pgfqpoint{2.937705in}{2.755497in}}%
\pgfpathlineto{\pgfqpoint{2.937835in}{2.630231in}}%
\pgfpathlineto{\pgfqpoint{2.938478in}{2.721127in}}%
\pgfpathlineto{\pgfqpoint{2.938528in}{2.673649in}}%
\pgfpathlineto{\pgfqpoint{2.939171in}{2.764517in}}%
\pgfpathlineto{\pgfqpoint{2.939220in}{2.717247in}}%
\pgfpathlineto{\pgfqpoint{2.939645in}{2.794240in}}%
\pgfpathlineto{\pgfqpoint{2.939765in}{2.644559in}}%
\pgfpathlineto{\pgfqpoint{2.939902in}{2.644559in}}%
\pgfpathlineto{\pgfqpoint{2.942647in}{2.644571in}}%
\pgfpathlineto{\pgfqpoint{2.942652in}{2.773967in}}%
\pgfpathlineto{\pgfqpoint{2.943408in}{2.704134in}}%
\pgfpathlineto{\pgfqpoint{2.943457in}{2.688539in}}%
\pgfpathlineto{\pgfqpoint{2.944001in}{2.741390in}}%
\pgfpathlineto{\pgfqpoint{2.944051in}{2.725846in}}%
\pgfpathlineto{\pgfqpoint{2.944582in}{2.777886in}}%
\pgfpathlineto{\pgfqpoint{2.944682in}{2.700005in}}%
\pgfpathlineto{\pgfqpoint{2.945468in}{2.644559in}}%
\pgfpathlineto{\pgfqpoint{2.947592in}{2.644559in}}%
\pgfpathlineto{\pgfqpoint{2.947597in}{2.747901in}}%
\pgfpathlineto{\pgfqpoint{2.947656in}{2.641526in}}%
\pgfpathlineto{\pgfqpoint{2.948337in}{2.680310in}}%
\pgfpathlineto{\pgfqpoint{2.949178in}{2.759841in}}%
\pgfpathlineto{\pgfqpoint{2.949565in}{2.784140in}}%
\pgfpathlineto{\pgfqpoint{2.950175in}{2.644559in}}%
\pgfpathlineto{\pgfqpoint{2.952538in}{2.644567in}}%
\pgfpathlineto{\pgfqpoint{2.952542in}{2.781425in}}%
\pgfpathlineto{\pgfqpoint{2.953274in}{2.700790in}}%
\pgfpathlineto{\pgfqpoint{2.953324in}{2.689076in}}%
\pgfpathlineto{\pgfqpoint{2.953868in}{2.738042in}}%
\pgfpathlineto{\pgfqpoint{2.953917in}{2.726388in}}%
\pgfpathlineto{\pgfqpoint{2.954511in}{2.778405in}}%
\pgfpathlineto{\pgfqpoint{2.954565in}{2.697795in}}%
\pgfpathlineto{\pgfqpoint{2.954580in}{2.745593in}}%
\pgfpathlineto{\pgfqpoint{2.955362in}{2.644559in}}%
\pgfpathlineto{\pgfqpoint{2.957482in}{2.644559in}}%
\pgfpathlineto{\pgfqpoint{2.957488in}{2.747767in}}%
\pgfpathlineto{\pgfqpoint{2.957543in}{2.642703in}}%
\pgfpathlineto{\pgfqpoint{2.958221in}{2.680653in}}%
\pgfpathlineto{\pgfqpoint{2.959062in}{2.758562in}}%
\pgfpathlineto{\pgfqpoint{2.959456in}{2.783287in}}%
\pgfpathlineto{\pgfqpoint{2.960090in}{2.644559in}}%
\pgfpathlineto{\pgfqpoint{2.962427in}{2.644559in}}%
\pgfpathlineto{\pgfqpoint{2.962433in}{2.748962in}}%
\pgfpathlineto{\pgfqpoint{2.963176in}{2.693226in}}%
\pgfpathlineto{\pgfqpoint{2.964402in}{2.771817in}}%
\pgfpathlineto{\pgfqpoint{2.964402in}{2.747143in}}%
\pgfpathlineto{\pgfqpoint{2.964428in}{2.757068in}}%
\pgfpathlineto{\pgfqpoint{2.965172in}{2.644559in}}%
\pgfpathlineto{\pgfqpoint{2.967373in}{2.644559in}}%
\pgfpathlineto{\pgfqpoint{2.967379in}{2.748473in}}%
\pgfpathlineto{\pgfqpoint{2.967452in}{2.637659in}}%
\pgfpathlineto{\pgfqpoint{2.968138in}{2.678473in}}%
\pgfpathlineto{\pgfqpoint{2.968978in}{2.764392in}}%
\pgfpathlineto{\pgfqpoint{2.969347in}{2.787508in}}%
\pgfpathlineto{\pgfqpoint{2.969945in}{2.644559in}}%
\pgfpathlineto{\pgfqpoint{2.972318in}{2.644559in}}%
\pgfpathlineto{\pgfqpoint{2.972324in}{2.748164in}}%
\pgfpathlineto{\pgfqpoint{2.973075in}{2.695131in}}%
\pgfpathlineto{\pgfqpoint{2.973718in}{2.733814in}}%
\pgfpathlineto{\pgfqpoint{2.974293in}{2.769916in}}%
\pgfpathlineto{\pgfqpoint{2.974320in}{2.757146in}}%
\pgfpathlineto{\pgfqpoint{2.975109in}{2.644559in}}%
\pgfpathlineto{\pgfqpoint{2.977264in}{2.644559in}}%
\pgfpathlineto{\pgfqpoint{2.977270in}{2.748426in}}%
\pgfpathlineto{\pgfqpoint{2.977340in}{2.640395in}}%
\pgfpathlineto{\pgfqpoint{2.978025in}{2.680884in}}%
\pgfpathlineto{\pgfqpoint{2.978866in}{2.761974in}}%
\pgfpathlineto{\pgfqpoint{2.979238in}{2.785302in}}%
\pgfpathlineto{\pgfqpoint{2.979864in}{2.644559in}}%
\pgfpathlineto{\pgfqpoint{2.982211in}{2.644560in}}%
\pgfpathlineto{\pgfqpoint{2.982215in}{2.793070in}}%
\pgfpathlineto{\pgfqpoint{2.982347in}{2.630756in}}%
\pgfpathlineto{\pgfqpoint{2.982941in}{2.667963in}}%
\pgfpathlineto{\pgfqpoint{2.983782in}{2.770155in}}%
\pgfpathlineto{\pgfqpoint{2.984156in}{2.793592in}}%
\pgfpathlineto{\pgfqpoint{2.984686in}{2.644559in}}%
\pgfpathlineto{\pgfqpoint{2.987155in}{2.644559in}}%
\pgfpathlineto{\pgfqpoint{2.987161in}{2.747988in}}%
\pgfpathlineto{\pgfqpoint{2.987219in}{2.641812in}}%
\pgfpathlineto{\pgfqpoint{2.987900in}{2.680521in}}%
\pgfpathlineto{\pgfqpoint{2.988741in}{2.759981in}}%
\pgfpathlineto{\pgfqpoint{2.989129in}{2.784322in}}%
\pgfpathlineto{\pgfqpoint{2.989754in}{2.644559in}}%
\pgfpathlineto{\pgfqpoint{2.992100in}{2.644559in}}%
\pgfpathlineto{\pgfqpoint{2.992106in}{2.748694in}}%
\pgfpathlineto{\pgfqpoint{2.992862in}{2.702177in}}%
\pgfpathlineto{\pgfqpoint{2.992911in}{2.690280in}}%
\pgfpathlineto{\pgfqpoint{2.993455in}{2.739421in}}%
\pgfpathlineto{\pgfqpoint{2.993505in}{2.727599in}}%
\pgfpathlineto{\pgfqpoint{2.994037in}{2.775915in}}%
\pgfpathlineto{\pgfqpoint{2.994138in}{2.704312in}}%
\pgfpathlineto{\pgfqpoint{2.994918in}{2.644559in}}%
\pgfpathlineto{\pgfqpoint{2.997046in}{2.644559in}}%
\pgfpathlineto{\pgfqpoint{2.997051in}{2.747109in}}%
\pgfpathlineto{\pgfqpoint{2.997789in}{2.698692in}}%
\pgfpathlineto{\pgfqpoint{2.997838in}{2.691726in}}%
\pgfpathlineto{\pgfqpoint{2.998283in}{2.729737in}}%
\pgfpathlineto{\pgfqpoint{2.998333in}{2.722817in}}%
\pgfpathlineto{\pgfqpoint{2.998973in}{2.773038in}}%
\pgfpathlineto{\pgfqpoint{2.999056in}{2.714829in}}%
\pgfpathlineto{\pgfqpoint{2.999065in}{2.749248in}}%
\pgfpathlineto{\pgfqpoint{2.999854in}{2.644559in}}%
\pgfpathlineto{\pgfqpoint{3.001993in}{2.644899in}}%
\pgfpathlineto{\pgfqpoint{3.001998in}{2.744110in}}%
\pgfpathlineto{\pgfqpoint{3.002074in}{2.619587in}}%
\pgfpathlineto{\pgfqpoint{3.002759in}{2.660770in}}%
\pgfpathlineto{\pgfqpoint{3.003600in}{2.781745in}}%
\pgfpathlineto{\pgfqpoint{3.003965in}{2.804611in}}%
\pgfpathlineto{\pgfqpoint{3.004472in}{2.644559in}}%
\pgfpathlineto{\pgfqpoint{3.006937in}{2.644559in}}%
\pgfpathlineto{\pgfqpoint{3.006943in}{2.746192in}}%
\pgfpathlineto{\pgfqpoint{3.006998in}{2.642809in}}%
\pgfpathlineto{\pgfqpoint{3.007677in}{2.680835in}}%
\pgfpathlineto{\pgfqpoint{3.008518in}{2.758788in}}%
\pgfpathlineto{\pgfqpoint{3.008911in}{2.783444in}}%
\pgfpathlineto{\pgfqpoint{3.009544in}{2.644559in}}%
\pgfpathlineto{\pgfqpoint{3.011884in}{2.644744in}}%
\pgfpathlineto{\pgfqpoint{3.011888in}{2.751542in}}%
\pgfpathlineto{\pgfqpoint{3.011964in}{2.626692in}}%
\pgfpathlineto{\pgfqpoint{3.012649in}{2.667767in}}%
\pgfpathlineto{\pgfqpoint{3.013490in}{2.774704in}}%
\pgfpathlineto{\pgfqpoint{3.013856in}{2.797639in}}%
\pgfpathlineto{\pgfqpoint{3.014388in}{2.644559in}}%
\pgfpathlineto{\pgfqpoint{3.016827in}{2.644559in}}%
\pgfpathlineto{\pgfqpoint{3.016832in}{2.767701in}}%
\pgfpathlineto{\pgfqpoint{3.016958in}{2.624404in}}%
\pgfpathlineto{\pgfqpoint{3.017601in}{2.727070in}}%
\pgfpathlineto{\pgfqpoint{3.017651in}{2.667794in}}%
\pgfpathlineto{\pgfqpoint{3.018294in}{2.770441in}}%
\pgfpathlineto{\pgfqpoint{3.018343in}{2.711413in}}%
\pgfpathlineto{\pgfqpoint{3.018770in}{2.800283in}}%
\pgfpathlineto{\pgfqpoint{3.018884in}{2.644559in}}%
\pgfpathlineto{\pgfqpoint{3.019063in}{2.644559in}}%
\pgfpathlineto{\pgfqpoint{3.021773in}{2.644559in}}%
\pgfpathlineto{\pgfqpoint{3.021779in}{2.746580in}}%
\pgfpathlineto{\pgfqpoint{3.022534in}{2.694712in}}%
\pgfpathlineto{\pgfqpoint{3.023747in}{2.770993in}}%
\pgfpathlineto{\pgfqpoint{3.023789in}{2.736992in}}%
\pgfpathlineto{\pgfqpoint{3.024611in}{2.644559in}}%
\pgfpathlineto{\pgfqpoint{3.026719in}{2.644559in}}%
\pgfpathlineto{\pgfqpoint{3.026724in}{2.743907in}}%
\pgfpathlineto{\pgfqpoint{3.027453in}{2.682281in}}%
\pgfpathlineto{\pgfqpoint{3.028293in}{2.756962in}}%
\pgfpathlineto{\pgfqpoint{3.028666in}{2.780352in}}%
\pgfpathlineto{\pgfqpoint{3.029362in}{2.644559in}}%
\pgfpathlineto{\pgfqpoint{3.031664in}{2.644559in}}%
\pgfpathlineto{\pgfqpoint{3.031670in}{2.747543in}}%
\pgfpathlineto{\pgfqpoint{3.031798in}{2.644334in}}%
\pgfpathlineto{\pgfqpoint{3.032391in}{2.681468in}}%
\pgfpathlineto{\pgfqpoint{3.033232in}{2.756096in}}%
\pgfpathlineto{\pgfqpoint{3.033608in}{2.779682in}}%
\pgfpathlineto{\pgfqpoint{3.034295in}{2.644559in}}%
\pgfpathlineto{\pgfqpoint{3.036609in}{2.644559in}}%
\pgfpathlineto{\pgfqpoint{3.036615in}{2.760168in}}%
\pgfpathlineto{\pgfqpoint{3.036671in}{2.624563in}}%
\pgfpathlineto{\pgfqpoint{3.037349in}{2.662645in}}%
\pgfpathlineto{\pgfqpoint{3.038190in}{2.777082in}}%
\pgfpathlineto{\pgfqpoint{3.038584in}{2.801738in}}%
\pgfpathlineto{\pgfqpoint{3.039072in}{2.644559in}}%
\pgfpathlineto{\pgfqpoint{3.041555in}{2.644559in}}%
\pgfpathlineto{\pgfqpoint{3.041560in}{2.747066in}}%
\pgfpathlineto{\pgfqpoint{3.042316in}{2.694049in}}%
\pgfpathlineto{\pgfqpoint{3.043529in}{2.772041in}}%
\pgfpathlineto{\pgfqpoint{3.043529in}{2.747512in}}%
\pgfpathlineto{\pgfqpoint{3.043551in}{2.755257in}}%
\pgfpathlineto{\pgfqpoint{3.044298in}{2.644559in}}%
\pgfpathlineto{\pgfqpoint{3.046502in}{2.644607in}}%
\pgfpathlineto{\pgfqpoint{3.046504in}{2.769554in}}%
\pgfpathlineto{\pgfqpoint{3.046560in}{2.621966in}}%
\pgfpathlineto{\pgfqpoint{3.047236in}{2.659635in}}%
\pgfpathlineto{\pgfqpoint{3.048077in}{2.778870in}}%
\pgfpathlineto{\pgfqpoint{3.048449in}{2.802134in}}%
\pgfpathlineto{\pgfqpoint{3.048953in}{2.644559in}}%
\pgfpathlineto{\pgfqpoint{3.051446in}{2.644559in}}%
\pgfpathlineto{\pgfqpoint{3.051451in}{2.747917in}}%
\pgfpathlineto{\pgfqpoint{3.051509in}{2.641940in}}%
\pgfpathlineto{\pgfqpoint{3.052190in}{2.680481in}}%
\pgfpathlineto{\pgfqpoint{3.053030in}{2.759641in}}%
\pgfpathlineto{\pgfqpoint{3.053420in}{2.784064in}}%
\pgfpathlineto{\pgfqpoint{3.054046in}{2.644559in}}%
\pgfpathlineto{\pgfqpoint{3.056391in}{2.644559in}}%
\pgfpathlineto{\pgfqpoint{3.056397in}{2.746005in}}%
\pgfpathlineto{\pgfqpoint{3.057131in}{2.694232in}}%
\pgfpathlineto{\pgfqpoint{3.057873in}{2.739215in}}%
\pgfpathlineto{\pgfqpoint{3.058365in}{2.770167in}}%
\pgfpathlineto{\pgfqpoint{3.058366in}{2.734327in}}%
\pgfpathlineto{\pgfqpoint{3.058392in}{2.757368in}}%
\pgfpathlineto{\pgfqpoint{3.059186in}{2.644559in}}%
\pgfpathlineto{\pgfqpoint{3.061337in}{2.644559in}}%
\pgfpathlineto{\pgfqpoint{3.061342in}{2.748307in}}%
\pgfpathlineto{\pgfqpoint{3.061417in}{2.636214in}}%
\pgfpathlineto{\pgfqpoint{3.062103in}{2.677155in}}%
\pgfpathlineto{\pgfqpoint{3.062943in}{2.765200in}}%
\pgfpathlineto{\pgfqpoint{3.063311in}{2.788222in}}%
\pgfpathlineto{\pgfqpoint{3.063901in}{2.644559in}}%
\pgfpathlineto{\pgfqpoint{3.066281in}{2.644559in}}%
\pgfpathlineto{\pgfqpoint{3.066287in}{2.753047in}}%
\pgfpathlineto{\pgfqpoint{3.066391in}{2.633628in}}%
\pgfpathlineto{\pgfqpoint{3.067031in}{2.715280in}}%
\pgfpathlineto{\pgfqpoint{3.067081in}{2.676474in}}%
\pgfpathlineto{\pgfqpoint{3.067724in}{2.758694in}}%
\pgfpathlineto{\pgfqpoint{3.067773in}{2.720049in}}%
\pgfpathlineto{\pgfqpoint{3.068213in}{2.789336in}}%
\pgfpathlineto{\pgfqpoint{3.068347in}{2.644559in}}%
\pgfpathlineto{\pgfqpoint{3.068452in}{2.644559in}}%
\pgfpathlineto{\pgfqpoint{3.071228in}{2.644559in}}%
\pgfpathlineto{\pgfqpoint{3.071234in}{2.746779in}}%
\pgfpathlineto{\pgfqpoint{3.071294in}{2.640859in}}%
\pgfpathlineto{\pgfqpoint{3.071976in}{2.679880in}}%
\pgfpathlineto{\pgfqpoint{3.072816in}{2.759744in}}%
\pgfpathlineto{\pgfqpoint{3.073202in}{2.783903in}}%
\pgfpathlineto{\pgfqpoint{3.073804in}{2.644559in}}%
\pgfpathlineto{\pgfqpoint{3.076174in}{2.644559in}}%
\pgfpathlineto{\pgfqpoint{3.076178in}{2.752924in}}%
\pgfpathlineto{\pgfqpoint{3.076233in}{2.643611in}}%
\pgfpathlineto{\pgfqpoint{3.076910in}{2.681242in}}%
\pgfpathlineto{\pgfqpoint{3.077751in}{2.758313in}}%
\pgfpathlineto{\pgfqpoint{3.078122in}{2.781595in}}%
\pgfpathlineto{\pgfqpoint{3.078788in}{2.644559in}}%
\pgfpathlineto{\pgfqpoint{3.081120in}{2.644610in}}%
\pgfpathlineto{\pgfqpoint{3.081122in}{2.770062in}}%
\pgfpathlineto{\pgfqpoint{3.081178in}{2.622131in}}%
\pgfpathlineto{\pgfqpoint{3.081854in}{2.659796in}}%
\pgfpathlineto{\pgfqpoint{3.082695in}{2.779125in}}%
\pgfpathlineto{\pgfqpoint{3.083067in}{2.802392in}}%
\pgfpathlineto{\pgfqpoint{3.083567in}{2.644559in}}%
\pgfpathlineto{\pgfqpoint{3.086065in}{2.644587in}}%
\pgfpathlineto{\pgfqpoint{3.086068in}{2.765857in}}%
\pgfpathlineto{\pgfqpoint{3.086162in}{2.638144in}}%
\pgfpathlineto{\pgfqpoint{3.086799in}{2.709129in}}%
\pgfpathlineto{\pgfqpoint{3.086848in}{2.680396in}}%
\pgfpathlineto{\pgfqpoint{3.087491in}{2.752526in}}%
\pgfpathlineto{\pgfqpoint{3.087541in}{2.723989in}}%
\pgfpathlineto{\pgfqpoint{3.088038in}{2.786777in}}%
\pgfpathlineto{\pgfqpoint{3.088118in}{2.644559in}}%
\pgfpathlineto{\pgfqpoint{3.088198in}{2.644559in}}%
\pgfpathlineto{\pgfqpoint{3.091009in}{2.644559in}}%
\pgfpathlineto{\pgfqpoint{3.091015in}{2.743750in}}%
\pgfpathlineto{\pgfqpoint{3.091128in}{2.633223in}}%
\pgfpathlineto{\pgfqpoint{3.091771in}{2.716243in}}%
\pgfpathlineto{\pgfqpoint{3.091820in}{2.676422in}}%
\pgfpathlineto{\pgfqpoint{3.092463in}{2.759656in}}%
\pgfpathlineto{\pgfqpoint{3.092513in}{2.719999in}}%
\pgfpathlineto{\pgfqpoint{3.092946in}{2.789914in}}%
\pgfpathlineto{\pgfqpoint{3.093074in}{2.644559in}}%
\pgfpathlineto{\pgfqpoint{3.093169in}{2.644559in}}%
\pgfpathlineto{\pgfqpoint{3.095954in}{2.644559in}}%
\pgfpathlineto{\pgfqpoint{3.095960in}{2.751848in}}%
\pgfpathlineto{\pgfqpoint{3.096707in}{2.694753in}}%
\pgfpathlineto{\pgfqpoint{3.097251in}{2.727128in}}%
\pgfpathlineto{\pgfqpoint{3.097929in}{2.769736in}}%
\pgfpathlineto{\pgfqpoint{3.097955in}{2.756747in}}%
\pgfpathlineto{\pgfqpoint{3.098747in}{2.644559in}}%
\pgfpathlineto{\pgfqpoint{3.100902in}{2.644651in}}%
\pgfpathlineto{\pgfqpoint{3.100906in}{2.757073in}}%
\pgfpathlineto{\pgfqpoint{3.101668in}{2.692279in}}%
\pgfpathlineto{\pgfqpoint{3.102805in}{2.770206in}}%
\pgfpathlineto{\pgfqpoint{3.102930in}{2.713454in}}%
\pgfpathlineto{\pgfqpoint{3.102874in}{2.774561in}}%
\pgfpathlineto{\pgfqpoint{3.102944in}{2.731235in}}%
\pgfpathlineto{\pgfqpoint{3.103728in}{2.644559in}}%
\pgfpathlineto{\pgfqpoint{3.105846in}{2.644559in}}%
\pgfpathlineto{\pgfqpoint{3.105851in}{2.746455in}}%
\pgfpathlineto{\pgfqpoint{3.106610in}{2.695091in}}%
\pgfpathlineto{\pgfqpoint{3.107820in}{2.770507in}}%
\pgfpathlineto{\pgfqpoint{3.107862in}{2.740501in}}%
\pgfpathlineto{\pgfqpoint{3.108674in}{2.644559in}}%
\pgfpathlineto{\pgfqpoint{3.110792in}{2.644559in}}%
\pgfpathlineto{\pgfqpoint{3.110797in}{2.746327in}}%
\pgfpathlineto{\pgfqpoint{3.110853in}{2.642787in}}%
\pgfpathlineto{\pgfqpoint{3.111532in}{2.680992in}}%
\pgfpathlineto{\pgfqpoint{3.112373in}{2.759189in}}%
\pgfpathlineto{\pgfqpoint{3.112765in}{2.783800in}}%
\pgfpathlineto{\pgfqpoint{3.113399in}{2.644559in}}%
\pgfpathlineto{\pgfqpoint{3.115737in}{2.644559in}}%
\pgfpathlineto{\pgfqpoint{3.115742in}{2.748183in}}%
\pgfpathlineto{\pgfqpoint{3.115812in}{2.640601in}}%
\pgfpathlineto{\pgfqpoint{3.116498in}{2.681029in}}%
\pgfpathlineto{\pgfqpoint{3.117338in}{2.760914in}}%
\pgfpathlineto{\pgfqpoint{3.117711in}{2.784268in}}%
\pgfpathlineto{\pgfqpoint{3.118321in}{2.644559in}}%
\pgfpathlineto{\pgfqpoint{3.120684in}{2.644974in}}%
\pgfpathlineto{\pgfqpoint{3.120687in}{2.747596in}}%
\pgfpathlineto{\pgfqpoint{3.120753in}{2.633735in}}%
\pgfpathlineto{\pgfqpoint{3.121437in}{2.673565in}}%
\pgfpathlineto{\pgfqpoint{3.122278in}{2.767989in}}%
\pgfpathlineto{\pgfqpoint{3.122656in}{2.791726in}}%
\pgfpathlineto{\pgfqpoint{3.123205in}{2.644559in}}%
\pgfpathlineto{\pgfqpoint{3.125628in}{2.644559in}}%
\pgfpathlineto{\pgfqpoint{3.125633in}{2.748050in}}%
\pgfpathlineto{\pgfqpoint{3.125696in}{2.640764in}}%
\pgfpathlineto{\pgfqpoint{3.126379in}{2.680202in}}%
\pgfpathlineto{\pgfqpoint{3.127220in}{2.760877in}}%
\pgfpathlineto{\pgfqpoint{3.127602in}{2.784826in}}%
\pgfpathlineto{\pgfqpoint{3.128192in}{2.644559in}}%
\pgfpathlineto{\pgfqpoint{3.130572in}{2.644559in}}%
\pgfpathlineto{\pgfqpoint{3.130578in}{2.754069in}}%
\pgfpathlineto{\pgfqpoint{3.130633in}{2.644227in}}%
\pgfpathlineto{\pgfqpoint{3.131309in}{2.681746in}}%
\pgfpathlineto{\pgfqpoint{3.132150in}{2.757097in}}%
\pgfpathlineto{\pgfqpoint{3.132522in}{2.780404in}}%
\pgfpathlineto{\pgfqpoint{3.133205in}{2.644559in}}%
\pgfpathlineto{\pgfqpoint{3.135519in}{2.644559in}}%
\pgfpathlineto{\pgfqpoint{3.135524in}{2.747694in}}%
\pgfpathlineto{\pgfqpoint{3.136245in}{2.681973in}}%
\pgfpathlineto{\pgfqpoint{3.137085in}{2.756341in}}%
\pgfpathlineto{\pgfqpoint{3.137462in}{2.779966in}}%
\pgfpathlineto{\pgfqpoint{3.138157in}{2.644559in}}%
\pgfpathlineto{\pgfqpoint{3.140465in}{2.644563in}}%
\pgfpathlineto{\pgfqpoint{3.140470in}{2.786583in}}%
\pgfpathlineto{\pgfqpoint{3.140544in}{2.626381in}}%
\pgfpathlineto{\pgfqpoint{3.141229in}{2.667307in}}%
\pgfpathlineto{\pgfqpoint{3.142069in}{2.775285in}}%
\pgfpathlineto{\pgfqpoint{3.142438in}{2.798374in}}%
\pgfpathlineto{\pgfqpoint{3.142964in}{2.644559in}}%
\pgfpathlineto{\pgfqpoint{3.145410in}{2.644559in}}%
\pgfpathlineto{\pgfqpoint{3.145415in}{2.747983in}}%
\pgfpathlineto{\pgfqpoint{3.145472in}{2.642372in}}%
\pgfpathlineto{\pgfqpoint{3.146152in}{2.680777in}}%
\pgfpathlineto{\pgfqpoint{3.146993in}{2.759633in}}%
\pgfpathlineto{\pgfqpoint{3.147384in}{2.784125in}}%
\pgfpathlineto{\pgfqpoint{3.148013in}{2.644559in}}%
\pgfpathlineto{\pgfqpoint{3.150354in}{2.644559in}}%
\pgfpathlineto{\pgfqpoint{3.150360in}{2.755512in}}%
\pgfpathlineto{\pgfqpoint{3.150469in}{2.633551in}}%
\pgfpathlineto{\pgfqpoint{3.151112in}{2.715304in}}%
\pgfpathlineto{\pgfqpoint{3.151162in}{2.676665in}}%
\pgfpathlineto{\pgfqpoint{3.151804in}{2.758716in}}%
\pgfpathlineto{\pgfqpoint{3.151854in}{2.720243in}}%
\pgfpathlineto{\pgfqpoint{3.152289in}{2.789109in}}%
\pgfpathlineto{\pgfqpoint{3.152420in}{2.644559in}}%
\pgfpathlineto{\pgfqpoint{3.152521in}{2.644559in}}%
\pgfpathlineto{\pgfqpoint{3.155300in}{2.644559in}}%
\pgfpathlineto{\pgfqpoint{3.155306in}{2.748204in}}%
\pgfpathlineto{\pgfqpoint{3.156062in}{2.700953in}}%
\pgfpathlineto{\pgfqpoint{3.156111in}{2.691272in}}%
\pgfpathlineto{\pgfqpoint{3.156557in}{2.732009in}}%
\pgfpathlineto{\pgfqpoint{3.156606in}{2.722352in}}%
\pgfpathlineto{\pgfqpoint{3.157237in}{2.774734in}}%
\pgfpathlineto{\pgfqpoint{3.157330in}{2.713519in}}%
\pgfpathlineto{\pgfqpoint{3.157344in}{2.729983in}}%
\pgfpathlineto{\pgfqpoint{3.158128in}{2.644559in}}%
\pgfpathlineto{\pgfqpoint{3.160246in}{2.644559in}}%
\pgfpathlineto{\pgfqpoint{3.160251in}{2.748214in}}%
\pgfpathlineto{\pgfqpoint{3.160986in}{2.702510in}}%
\pgfpathlineto{\pgfqpoint{3.161036in}{2.688195in}}%
\pgfpathlineto{\pgfqpoint{3.161580in}{2.739768in}}%
\pgfpathlineto{\pgfqpoint{3.161629in}{2.725500in}}%
\pgfpathlineto{\pgfqpoint{3.162173in}{2.777024in}}%
\pgfpathlineto{\pgfqpoint{3.162284in}{2.709736in}}%
\pgfpathlineto{\pgfqpoint{3.163065in}{2.644559in}}%
\pgfpathlineto{\pgfqpoint{3.165191in}{2.644559in}}%
\pgfpathlineto{\pgfqpoint{3.165197in}{2.742045in}}%
\pgfpathlineto{\pgfqpoint{3.165279in}{2.627225in}}%
\pgfpathlineto{\pgfqpoint{3.165915in}{2.719012in}}%
\pgfpathlineto{\pgfqpoint{3.165965in}{2.668912in}}%
\pgfpathlineto{\pgfqpoint{3.166608in}{2.762373in}}%
\pgfpathlineto{\pgfqpoint{3.166657in}{2.712540in}}%
\pgfpathlineto{\pgfqpoint{3.167165in}{2.797282in}}%
\pgfpathlineto{\pgfqpoint{3.167255in}{2.644559in}}%
\pgfpathlineto{\pgfqpoint{3.167345in}{2.644559in}}%
\pgfpathlineto{\pgfqpoint{3.170132in}{2.644559in}}%
\pgfpathlineto{\pgfqpoint{3.170132in}{2.644559in}}%
\pgfusepath{stroke}%
\end{pgfscope}%
\begin{pgfscope}%
\pgfsetrectcap%
\pgfsetmiterjoin%
\pgfsetlinewidth{0.803000pt}%
\definecolor{currentstroke}{rgb}{0.000000,0.000000,0.000000}%
\pgfsetstrokecolor{currentstroke}%
\pgfsetdash{}{0pt}%
\pgfpathmoveto{\pgfqpoint{0.573768in}{2.601404in}}%
\pgfpathlineto{\pgfqpoint{0.573768in}{3.507654in}}%
\pgfusepath{stroke}%
\end{pgfscope}%
\begin{pgfscope}%
\pgfsetrectcap%
\pgfsetmiterjoin%
\pgfsetlinewidth{0.803000pt}%
\definecolor{currentstroke}{rgb}{0.000000,0.000000,0.000000}%
\pgfsetstrokecolor{currentstroke}%
\pgfsetdash{}{0pt}%
\pgfpathmoveto{\pgfqpoint{3.293768in}{2.601404in}}%
\pgfpathlineto{\pgfqpoint{3.293768in}{3.507654in}}%
\pgfusepath{stroke}%
\end{pgfscope}%
\begin{pgfscope}%
\pgfsetrectcap%
\pgfsetmiterjoin%
\pgfsetlinewidth{0.803000pt}%
\definecolor{currentstroke}{rgb}{0.000000,0.000000,0.000000}%
\pgfsetstrokecolor{currentstroke}%
\pgfsetdash{}{0pt}%
\pgfpathmoveto{\pgfqpoint{0.573768in}{2.601404in}}%
\pgfpathlineto{\pgfqpoint{3.293768in}{2.601404in}}%
\pgfusepath{stroke}%
\end{pgfscope}%
\begin{pgfscope}%
\pgfsetrectcap%
\pgfsetmiterjoin%
\pgfsetlinewidth{0.803000pt}%
\definecolor{currentstroke}{rgb}{0.000000,0.000000,0.000000}%
\pgfsetstrokecolor{currentstroke}%
\pgfsetdash{}{0pt}%
\pgfpathmoveto{\pgfqpoint{0.573768in}{3.507654in}}%
\pgfpathlineto{\pgfqpoint{3.293768in}{3.507654in}}%
\pgfusepath{stroke}%
\end{pgfscope}%
\begin{pgfscope}%
\definecolor{textcolor}{rgb}{0.000000,0.000000,0.000000}%
\pgfsetstrokecolor{textcolor}%
\pgfsetfillcolor{textcolor}%
\pgftext[x=0.628168in,y=3.444217in,left,top]{\color{textcolor}{\rmfamily\fontsize{8.000000}{9.600000}\selectfont\catcode`\^=\active\def^{\ifmmode\sp\else\^{}\fi}\catcode`\%=\active\def%{\%}(a)}}%
\end{pgfscope}%
\begin{pgfscope}%
\pgfsetbuttcap%
\pgfsetmiterjoin%
\definecolor{currentfill}{rgb}{1.000000,1.000000,1.000000}%
\pgfsetfillcolor{currentfill}%
\pgfsetlinewidth{0.000000pt}%
\definecolor{currentstroke}{rgb}{0.000000,0.000000,0.000000}%
\pgfsetstrokecolor{currentstroke}%
\pgfsetstrokeopacity{0.000000}%
\pgfsetdash{}{0pt}%
\pgfpathmoveto{\pgfqpoint{0.573768in}{1.532029in}}%
\pgfpathlineto{\pgfqpoint{3.293768in}{1.532029in}}%
\pgfpathlineto{\pgfqpoint{3.293768in}{2.438279in}}%
\pgfpathlineto{\pgfqpoint{0.573768in}{2.438279in}}%
\pgfpathlineto{\pgfqpoint{0.573768in}{1.532029in}}%
\pgfpathclose%
\pgfusepath{fill}%
\end{pgfscope}%
\begin{pgfscope}%
\pgfpathrectangle{\pgfqpoint{0.573768in}{1.532029in}}{\pgfqpoint{2.720000in}{0.906250in}}%
\pgfusepath{clip}%
\pgfsetbuttcap%
\pgfsetroundjoin%
\pgfsetlinewidth{0.501875pt}%
\definecolor{currentstroke}{rgb}{0.690196,0.690196,0.690196}%
\pgfsetstrokecolor{currentstroke}%
\pgfsetstrokeopacity{0.250000}%
\pgfsetdash{{1.850000pt}{0.800000pt}}{0.000000pt}%
\pgfpathmoveto{\pgfqpoint{0.697405in}{1.532029in}}%
\pgfpathlineto{\pgfqpoint{0.697405in}{2.438279in}}%
\pgfusepath{stroke}%
\end{pgfscope}%
\begin{pgfscope}%
\pgfsetbuttcap%
\pgfsetroundjoin%
\definecolor{currentfill}{rgb}{0.000000,0.000000,0.000000}%
\pgfsetfillcolor{currentfill}%
\pgfsetlinewidth{0.803000pt}%
\definecolor{currentstroke}{rgb}{0.000000,0.000000,0.000000}%
\pgfsetstrokecolor{currentstroke}%
\pgfsetdash{}{0pt}%
\pgfsys@defobject{currentmarker}{\pgfqpoint{0.000000in}{-0.048611in}}{\pgfqpoint{0.000000in}{0.000000in}}{%
\pgfpathmoveto{\pgfqpoint{0.000000in}{0.000000in}}%
\pgfpathlineto{\pgfqpoint{0.000000in}{-0.048611in}}%
\pgfusepath{stroke,fill}%
}%
\begin{pgfscope}%
\pgfsys@transformshift{0.697405in}{1.532029in}%
\pgfsys@useobject{currentmarker}{}%
\end{pgfscope}%
\end{pgfscope}%
\begin{pgfscope}%
\pgfpathrectangle{\pgfqpoint{0.573768in}{1.532029in}}{\pgfqpoint{2.720000in}{0.906250in}}%
\pgfusepath{clip}%
\pgfsetbuttcap%
\pgfsetroundjoin%
\pgfsetlinewidth{0.501875pt}%
\definecolor{currentstroke}{rgb}{0.690196,0.690196,0.690196}%
\pgfsetstrokecolor{currentstroke}%
\pgfsetstrokeopacity{0.250000}%
\pgfsetdash{{1.850000pt}{0.800000pt}}{0.000000pt}%
\pgfpathmoveto{\pgfqpoint{1.191950in}{1.532029in}}%
\pgfpathlineto{\pgfqpoint{1.191950in}{2.438279in}}%
\pgfusepath{stroke}%
\end{pgfscope}%
\begin{pgfscope}%
\pgfsetbuttcap%
\pgfsetroundjoin%
\definecolor{currentfill}{rgb}{0.000000,0.000000,0.000000}%
\pgfsetfillcolor{currentfill}%
\pgfsetlinewidth{0.803000pt}%
\definecolor{currentstroke}{rgb}{0.000000,0.000000,0.000000}%
\pgfsetstrokecolor{currentstroke}%
\pgfsetdash{}{0pt}%
\pgfsys@defobject{currentmarker}{\pgfqpoint{0.000000in}{-0.048611in}}{\pgfqpoint{0.000000in}{0.000000in}}{%
\pgfpathmoveto{\pgfqpoint{0.000000in}{0.000000in}}%
\pgfpathlineto{\pgfqpoint{0.000000in}{-0.048611in}}%
\pgfusepath{stroke,fill}%
}%
\begin{pgfscope}%
\pgfsys@transformshift{1.191950in}{1.532029in}%
\pgfsys@useobject{currentmarker}{}%
\end{pgfscope}%
\end{pgfscope}%
\begin{pgfscope}%
\pgfpathrectangle{\pgfqpoint{0.573768in}{1.532029in}}{\pgfqpoint{2.720000in}{0.906250in}}%
\pgfusepath{clip}%
\pgfsetbuttcap%
\pgfsetroundjoin%
\pgfsetlinewidth{0.501875pt}%
\definecolor{currentstroke}{rgb}{0.690196,0.690196,0.690196}%
\pgfsetstrokecolor{currentstroke}%
\pgfsetstrokeopacity{0.250000}%
\pgfsetdash{{1.850000pt}{0.800000pt}}{0.000000pt}%
\pgfpathmoveto{\pgfqpoint{1.686496in}{1.532029in}}%
\pgfpathlineto{\pgfqpoint{1.686496in}{2.438279in}}%
\pgfusepath{stroke}%
\end{pgfscope}%
\begin{pgfscope}%
\pgfsetbuttcap%
\pgfsetroundjoin%
\definecolor{currentfill}{rgb}{0.000000,0.000000,0.000000}%
\pgfsetfillcolor{currentfill}%
\pgfsetlinewidth{0.803000pt}%
\definecolor{currentstroke}{rgb}{0.000000,0.000000,0.000000}%
\pgfsetstrokecolor{currentstroke}%
\pgfsetdash{}{0pt}%
\pgfsys@defobject{currentmarker}{\pgfqpoint{0.000000in}{-0.048611in}}{\pgfqpoint{0.000000in}{0.000000in}}{%
\pgfpathmoveto{\pgfqpoint{0.000000in}{0.000000in}}%
\pgfpathlineto{\pgfqpoint{0.000000in}{-0.048611in}}%
\pgfusepath{stroke,fill}%
}%
\begin{pgfscope}%
\pgfsys@transformshift{1.686496in}{1.532029in}%
\pgfsys@useobject{currentmarker}{}%
\end{pgfscope}%
\end{pgfscope}%
\begin{pgfscope}%
\pgfpathrectangle{\pgfqpoint{0.573768in}{1.532029in}}{\pgfqpoint{2.720000in}{0.906250in}}%
\pgfusepath{clip}%
\pgfsetbuttcap%
\pgfsetroundjoin%
\pgfsetlinewidth{0.501875pt}%
\definecolor{currentstroke}{rgb}{0.690196,0.690196,0.690196}%
\pgfsetstrokecolor{currentstroke}%
\pgfsetstrokeopacity{0.250000}%
\pgfsetdash{{1.850000pt}{0.800000pt}}{0.000000pt}%
\pgfpathmoveto{\pgfqpoint{2.181041in}{1.532029in}}%
\pgfpathlineto{\pgfqpoint{2.181041in}{2.438279in}}%
\pgfusepath{stroke}%
\end{pgfscope}%
\begin{pgfscope}%
\pgfsetbuttcap%
\pgfsetroundjoin%
\definecolor{currentfill}{rgb}{0.000000,0.000000,0.000000}%
\pgfsetfillcolor{currentfill}%
\pgfsetlinewidth{0.803000pt}%
\definecolor{currentstroke}{rgb}{0.000000,0.000000,0.000000}%
\pgfsetstrokecolor{currentstroke}%
\pgfsetdash{}{0pt}%
\pgfsys@defobject{currentmarker}{\pgfqpoint{0.000000in}{-0.048611in}}{\pgfqpoint{0.000000in}{0.000000in}}{%
\pgfpathmoveto{\pgfqpoint{0.000000in}{0.000000in}}%
\pgfpathlineto{\pgfqpoint{0.000000in}{-0.048611in}}%
\pgfusepath{stroke,fill}%
}%
\begin{pgfscope}%
\pgfsys@transformshift{2.181041in}{1.532029in}%
\pgfsys@useobject{currentmarker}{}%
\end{pgfscope}%
\end{pgfscope}%
\begin{pgfscope}%
\pgfpathrectangle{\pgfqpoint{0.573768in}{1.532029in}}{\pgfqpoint{2.720000in}{0.906250in}}%
\pgfusepath{clip}%
\pgfsetbuttcap%
\pgfsetroundjoin%
\pgfsetlinewidth{0.501875pt}%
\definecolor{currentstroke}{rgb}{0.690196,0.690196,0.690196}%
\pgfsetstrokecolor{currentstroke}%
\pgfsetstrokeopacity{0.250000}%
\pgfsetdash{{1.850000pt}{0.800000pt}}{0.000000pt}%
\pgfpathmoveto{\pgfqpoint{2.675587in}{1.532029in}}%
\pgfpathlineto{\pgfqpoint{2.675587in}{2.438279in}}%
\pgfusepath{stroke}%
\end{pgfscope}%
\begin{pgfscope}%
\pgfsetbuttcap%
\pgfsetroundjoin%
\definecolor{currentfill}{rgb}{0.000000,0.000000,0.000000}%
\pgfsetfillcolor{currentfill}%
\pgfsetlinewidth{0.803000pt}%
\definecolor{currentstroke}{rgb}{0.000000,0.000000,0.000000}%
\pgfsetstrokecolor{currentstroke}%
\pgfsetdash{}{0pt}%
\pgfsys@defobject{currentmarker}{\pgfqpoint{0.000000in}{-0.048611in}}{\pgfqpoint{0.000000in}{0.000000in}}{%
\pgfpathmoveto{\pgfqpoint{0.000000in}{0.000000in}}%
\pgfpathlineto{\pgfqpoint{0.000000in}{-0.048611in}}%
\pgfusepath{stroke,fill}%
}%
\begin{pgfscope}%
\pgfsys@transformshift{2.675587in}{1.532029in}%
\pgfsys@useobject{currentmarker}{}%
\end{pgfscope}%
\end{pgfscope}%
\begin{pgfscope}%
\pgfpathrectangle{\pgfqpoint{0.573768in}{1.532029in}}{\pgfqpoint{2.720000in}{0.906250in}}%
\pgfusepath{clip}%
\pgfsetbuttcap%
\pgfsetroundjoin%
\pgfsetlinewidth{0.501875pt}%
\definecolor{currentstroke}{rgb}{0.690196,0.690196,0.690196}%
\pgfsetstrokecolor{currentstroke}%
\pgfsetstrokeopacity{0.250000}%
\pgfsetdash{{1.850000pt}{0.800000pt}}{0.000000pt}%
\pgfpathmoveto{\pgfqpoint{3.170132in}{1.532029in}}%
\pgfpathlineto{\pgfqpoint{3.170132in}{2.438279in}}%
\pgfusepath{stroke}%
\end{pgfscope}%
\begin{pgfscope}%
\pgfsetbuttcap%
\pgfsetroundjoin%
\definecolor{currentfill}{rgb}{0.000000,0.000000,0.000000}%
\pgfsetfillcolor{currentfill}%
\pgfsetlinewidth{0.803000pt}%
\definecolor{currentstroke}{rgb}{0.000000,0.000000,0.000000}%
\pgfsetstrokecolor{currentstroke}%
\pgfsetdash{}{0pt}%
\pgfsys@defobject{currentmarker}{\pgfqpoint{0.000000in}{-0.048611in}}{\pgfqpoint{0.000000in}{0.000000in}}{%
\pgfpathmoveto{\pgfqpoint{0.000000in}{0.000000in}}%
\pgfpathlineto{\pgfqpoint{0.000000in}{-0.048611in}}%
\pgfusepath{stroke,fill}%
}%
\begin{pgfscope}%
\pgfsys@transformshift{3.170132in}{1.532029in}%
\pgfsys@useobject{currentmarker}{}%
\end{pgfscope}%
\end{pgfscope}%
\begin{pgfscope}%
\pgfpathrectangle{\pgfqpoint{0.573768in}{1.532029in}}{\pgfqpoint{2.720000in}{0.906250in}}%
\pgfusepath{clip}%
\pgfsetbuttcap%
\pgfsetroundjoin%
\pgfsetlinewidth{0.501875pt}%
\definecolor{currentstroke}{rgb}{0.690196,0.690196,0.690196}%
\pgfsetstrokecolor{currentstroke}%
\pgfsetstrokeopacity{0.250000}%
\pgfsetdash{{1.850000pt}{0.800000pt}}{0.000000pt}%
\pgfpathmoveto{\pgfqpoint{0.573768in}{1.532029in}}%
\pgfpathlineto{\pgfqpoint{3.293768in}{1.532029in}}%
\pgfusepath{stroke}%
\end{pgfscope}%
\begin{pgfscope}%
\pgfsetbuttcap%
\pgfsetroundjoin%
\definecolor{currentfill}{rgb}{0.000000,0.000000,0.000000}%
\pgfsetfillcolor{currentfill}%
\pgfsetlinewidth{0.803000pt}%
\definecolor{currentstroke}{rgb}{0.000000,0.000000,0.000000}%
\pgfsetstrokecolor{currentstroke}%
\pgfsetdash{}{0pt}%
\pgfsys@defobject{currentmarker}{\pgfqpoint{-0.048611in}{0.000000in}}{\pgfqpoint{-0.000000in}{0.000000in}}{%
\pgfpathmoveto{\pgfqpoint{-0.000000in}{0.000000in}}%
\pgfpathlineto{\pgfqpoint{-0.048611in}{0.000000in}}%
\pgfusepath{stroke,fill}%
}%
\begin{pgfscope}%
\pgfsys@transformshift{0.573768in}{1.532029in}%
\pgfsys@useobject{currentmarker}{}%
\end{pgfscope}%
\end{pgfscope}%
\begin{pgfscope}%
\definecolor{textcolor}{rgb}{0.000000,0.000000,0.000000}%
\pgfsetstrokecolor{textcolor}%
\pgfsetfillcolor{textcolor}%
\pgftext[x=0.417518in, y=1.493449in, left, base]{\color{textcolor}{\rmfamily\fontsize{8.000000}{9.600000}\selectfont\catcode`\^=\active\def^{\ifmmode\sp\else\^{}\fi}\catcode`\%=\active\def%{\%}$\mathdefault{0}$}}%
\end{pgfscope}%
\begin{pgfscope}%
\pgfpathrectangle{\pgfqpoint{0.573768in}{1.532029in}}{\pgfqpoint{2.720000in}{0.906250in}}%
\pgfusepath{clip}%
\pgfsetbuttcap%
\pgfsetroundjoin%
\pgfsetlinewidth{0.501875pt}%
\definecolor{currentstroke}{rgb}{0.690196,0.690196,0.690196}%
\pgfsetstrokecolor{currentstroke}%
\pgfsetstrokeopacity{0.250000}%
\pgfsetdash{{1.850000pt}{0.800000pt}}{0.000000pt}%
\pgfpathmoveto{\pgfqpoint{0.573768in}{1.834113in}}%
\pgfpathlineto{\pgfqpoint{3.293768in}{1.834113in}}%
\pgfusepath{stroke}%
\end{pgfscope}%
\begin{pgfscope}%
\pgfsetbuttcap%
\pgfsetroundjoin%
\definecolor{currentfill}{rgb}{0.000000,0.000000,0.000000}%
\pgfsetfillcolor{currentfill}%
\pgfsetlinewidth{0.803000pt}%
\definecolor{currentstroke}{rgb}{0.000000,0.000000,0.000000}%
\pgfsetstrokecolor{currentstroke}%
\pgfsetdash{}{0pt}%
\pgfsys@defobject{currentmarker}{\pgfqpoint{-0.048611in}{0.000000in}}{\pgfqpoint{-0.000000in}{0.000000in}}{%
\pgfpathmoveto{\pgfqpoint{-0.000000in}{0.000000in}}%
\pgfpathlineto{\pgfqpoint{-0.048611in}{0.000000in}}%
\pgfusepath{stroke,fill}%
}%
\begin{pgfscope}%
\pgfsys@transformshift{0.573768in}{1.834113in}%
\pgfsys@useobject{currentmarker}{}%
\end{pgfscope}%
\end{pgfscope}%
\begin{pgfscope}%
\definecolor{textcolor}{rgb}{0.000000,0.000000,0.000000}%
\pgfsetstrokecolor{textcolor}%
\pgfsetfillcolor{textcolor}%
\pgftext[x=0.358489in, y=1.795533in, left, base]{\color{textcolor}{\rmfamily\fontsize{8.000000}{9.600000}\selectfont\catcode`\^=\active\def^{\ifmmode\sp\else\^{}\fi}\catcode`\%=\active\def%{\%}$\mathdefault{50}$}}%
\end{pgfscope}%
\begin{pgfscope}%
\pgfpathrectangle{\pgfqpoint{0.573768in}{1.532029in}}{\pgfqpoint{2.720000in}{0.906250in}}%
\pgfusepath{clip}%
\pgfsetbuttcap%
\pgfsetroundjoin%
\pgfsetlinewidth{0.501875pt}%
\definecolor{currentstroke}{rgb}{0.690196,0.690196,0.690196}%
\pgfsetstrokecolor{currentstroke}%
\pgfsetstrokeopacity{0.250000}%
\pgfsetdash{{1.850000pt}{0.800000pt}}{0.000000pt}%
\pgfpathmoveto{\pgfqpoint{0.573768in}{2.136196in}}%
\pgfpathlineto{\pgfqpoint{3.293768in}{2.136196in}}%
\pgfusepath{stroke}%
\end{pgfscope}%
\begin{pgfscope}%
\pgfsetbuttcap%
\pgfsetroundjoin%
\definecolor{currentfill}{rgb}{0.000000,0.000000,0.000000}%
\pgfsetfillcolor{currentfill}%
\pgfsetlinewidth{0.803000pt}%
\definecolor{currentstroke}{rgb}{0.000000,0.000000,0.000000}%
\pgfsetstrokecolor{currentstroke}%
\pgfsetdash{}{0pt}%
\pgfsys@defobject{currentmarker}{\pgfqpoint{-0.048611in}{0.000000in}}{\pgfqpoint{-0.000000in}{0.000000in}}{%
\pgfpathmoveto{\pgfqpoint{-0.000000in}{0.000000in}}%
\pgfpathlineto{\pgfqpoint{-0.048611in}{0.000000in}}%
\pgfusepath{stroke,fill}%
}%
\begin{pgfscope}%
\pgfsys@transformshift{0.573768in}{2.136196in}%
\pgfsys@useobject{currentmarker}{}%
\end{pgfscope}%
\end{pgfscope}%
\begin{pgfscope}%
\definecolor{textcolor}{rgb}{0.000000,0.000000,0.000000}%
\pgfsetstrokecolor{textcolor}%
\pgfsetfillcolor{textcolor}%
\pgftext[x=0.299460in, y=2.097616in, left, base]{\color{textcolor}{\rmfamily\fontsize{8.000000}{9.600000}\selectfont\catcode`\^=\active\def^{\ifmmode\sp\else\^{}\fi}\catcode`\%=\active\def%{\%}$\mathdefault{100}$}}%
\end{pgfscope}%
\begin{pgfscope}%
\pgfpathrectangle{\pgfqpoint{0.573768in}{1.532029in}}{\pgfqpoint{2.720000in}{0.906250in}}%
\pgfusepath{clip}%
\pgfsetbuttcap%
\pgfsetroundjoin%
\pgfsetlinewidth{0.501875pt}%
\definecolor{currentstroke}{rgb}{0.690196,0.690196,0.690196}%
\pgfsetstrokecolor{currentstroke}%
\pgfsetstrokeopacity{0.250000}%
\pgfsetdash{{1.850000pt}{0.800000pt}}{0.000000pt}%
\pgfpathmoveto{\pgfqpoint{0.573768in}{2.438279in}}%
\pgfpathlineto{\pgfqpoint{3.293768in}{2.438279in}}%
\pgfusepath{stroke}%
\end{pgfscope}%
\begin{pgfscope}%
\pgfsetbuttcap%
\pgfsetroundjoin%
\definecolor{currentfill}{rgb}{0.000000,0.000000,0.000000}%
\pgfsetfillcolor{currentfill}%
\pgfsetlinewidth{0.803000pt}%
\definecolor{currentstroke}{rgb}{0.000000,0.000000,0.000000}%
\pgfsetstrokecolor{currentstroke}%
\pgfsetdash{}{0pt}%
\pgfsys@defobject{currentmarker}{\pgfqpoint{-0.048611in}{0.000000in}}{\pgfqpoint{-0.000000in}{0.000000in}}{%
\pgfpathmoveto{\pgfqpoint{-0.000000in}{0.000000in}}%
\pgfpathlineto{\pgfqpoint{-0.048611in}{0.000000in}}%
\pgfusepath{stroke,fill}%
}%
\begin{pgfscope}%
\pgfsys@transformshift{0.573768in}{2.438279in}%
\pgfsys@useobject{currentmarker}{}%
\end{pgfscope}%
\end{pgfscope}%
\begin{pgfscope}%
\definecolor{textcolor}{rgb}{0.000000,0.000000,0.000000}%
\pgfsetstrokecolor{textcolor}%
\pgfsetfillcolor{textcolor}%
\pgftext[x=0.299460in, y=2.399699in, left, base]{\color{textcolor}{\rmfamily\fontsize{8.000000}{9.600000}\selectfont\catcode`\^=\active\def^{\ifmmode\sp\else\^{}\fi}\catcode`\%=\active\def%{\%}$\mathdefault{150}$}}%
\end{pgfscope}%
\begin{pgfscope}%
\definecolor{textcolor}{rgb}{0.000000,0.000000,0.000000}%
\pgfsetstrokecolor{textcolor}%
\pgfsetfillcolor{textcolor}%
\pgftext[x=0.243905in,y=1.985154in,,bottom,rotate=90.000000]{\color{textcolor}{\rmfamily\fontsize{8.000000}{9.600000}\selectfont\catcode`\^=\active\def^{\ifmmode\sp\else\^{}\fi}\catcode`\%=\active\def%{\%}$v_{DS}$ (V)}}%
\end{pgfscope}%
\begin{pgfscope}%
\pgfpathrectangle{\pgfqpoint{0.573768in}{1.532029in}}{\pgfqpoint{2.720000in}{0.906250in}}%
\pgfusepath{clip}%
\pgfsetrectcap%
\pgfsetroundjoin%
\pgfsetlinewidth{1.505625pt}%
\definecolor{currentstroke}{rgb}{0.835294,0.368627,0.000000}%
\pgfsetstrokecolor{currentstroke}%
\pgfsetdash{}{0pt}%
\pgfpathmoveto{\pgfqpoint{0.697405in}{1.532030in}}%
\pgfpathlineto{\pgfqpoint{0.697409in}{1.539352in}}%
\pgfpathlineto{\pgfqpoint{0.697413in}{1.546049in}}%
\pgfpathlineto{\pgfqpoint{0.697444in}{1.532029in}}%
\pgfpathlineto{\pgfqpoint{0.698134in}{1.532046in}}%
\pgfpathlineto{\pgfqpoint{0.699411in}{1.532297in}}%
\pgfpathlineto{\pgfqpoint{0.699491in}{1.781049in}}%
\pgfpathlineto{\pgfqpoint{0.700326in}{1.753906in}}%
\pgfpathlineto{\pgfqpoint{0.700821in}{1.757734in}}%
\pgfpathlineto{\pgfqpoint{0.700425in}{1.753893in}}%
\pgfpathlineto{\pgfqpoint{0.701117in}{1.755294in}}%
\pgfpathlineto{\pgfqpoint{0.701612in}{1.758777in}}%
\pgfpathlineto{\pgfqpoint{0.701216in}{1.755029in}}%
\pgfpathlineto{\pgfqpoint{0.701909in}{1.756704in}}%
\pgfpathlineto{\pgfqpoint{0.702008in}{1.756226in}}%
\pgfpathlineto{\pgfqpoint{0.702357in}{1.766581in}}%
\pgfpathlineto{\pgfqpoint{0.702436in}{1.532070in}}%
\pgfpathlineto{\pgfqpoint{0.703119in}{1.532084in}}%
\pgfpathlineto{\pgfqpoint{0.704351in}{1.532225in}}%
\pgfpathlineto{\pgfqpoint{0.704421in}{1.813112in}}%
\pgfpathlineto{\pgfqpoint{0.705377in}{1.754428in}}%
\pgfpathlineto{\pgfqpoint{0.705382in}{1.747718in}}%
\pgfpathlineto{\pgfqpoint{0.705397in}{1.773685in}}%
\pgfpathlineto{\pgfqpoint{0.706139in}{1.757740in}}%
\pgfpathlineto{\pgfqpoint{0.706329in}{1.768336in}}%
\pgfpathlineto{\pgfqpoint{0.706911in}{1.763379in}}%
\pgfpathlineto{\pgfqpoint{0.706954in}{1.762726in}}%
\pgfpathlineto{\pgfqpoint{0.707302in}{1.771380in}}%
\pgfpathlineto{\pgfqpoint{0.707397in}{1.532111in}}%
\pgfpathlineto{\pgfqpoint{0.708080in}{1.532125in}}%
\pgfpathlineto{\pgfqpoint{0.709295in}{1.532273in}}%
\pgfpathlineto{\pgfqpoint{0.709371in}{1.845326in}}%
\pgfpathlineto{\pgfqpoint{0.710492in}{1.763293in}}%
\pgfpathlineto{\pgfqpoint{0.710496in}{1.757142in}}%
\pgfpathlineto{\pgfqpoint{0.710514in}{1.784907in}}%
\pgfpathlineto{\pgfqpoint{0.711257in}{1.771050in}}%
\pgfpathlineto{\pgfqpoint{0.711263in}{1.768478in}}%
\pgfpathlineto{\pgfqpoint{0.711278in}{1.779481in}}%
\pgfpathlineto{\pgfqpoint{0.712019in}{1.774817in}}%
\pgfpathlineto{\pgfqpoint{0.712248in}{1.785180in}}%
\pgfpathlineto{\pgfqpoint{0.712409in}{1.532149in}}%
\pgfpathlineto{\pgfqpoint{0.713002in}{1.532162in}}%
\pgfpathlineto{\pgfqpoint{0.714239in}{1.532326in}}%
\pgfpathlineto{\pgfqpoint{0.714293in}{1.853612in}}%
\pgfpathlineto{\pgfqpoint{0.714318in}{1.895669in}}%
\pgfpathlineto{\pgfqpoint{0.714335in}{1.718502in}}%
\pgfpathlineto{\pgfqpoint{0.715028in}{1.796903in}}%
\pgfpathlineto{\pgfqpoint{0.715039in}{1.754977in}}%
\pgfpathlineto{\pgfqpoint{0.715055in}{1.808581in}}%
\pgfpathlineto{\pgfqpoint{0.715798in}{1.777382in}}%
\pgfpathlineto{\pgfqpoint{0.715805in}{1.774635in}}%
\pgfpathlineto{\pgfqpoint{0.715819in}{1.797631in}}%
\pgfpathlineto{\pgfqpoint{0.716510in}{1.797164in}}%
\pgfpathlineto{\pgfqpoint{0.717301in}{1.532185in}}%
\pgfpathlineto{\pgfqpoint{0.716955in}{1.799101in}}%
\pgfpathlineto{\pgfqpoint{0.718133in}{1.532212in}}%
\pgfpathlineto{\pgfqpoint{0.719172in}{1.532248in}}%
\pgfpathlineto{\pgfqpoint{0.719187in}{1.532475in}}%
\pgfpathlineto{\pgfqpoint{0.719259in}{1.918173in}}%
\pgfpathlineto{\pgfqpoint{0.720222in}{1.768720in}}%
\pgfpathlineto{\pgfqpoint{0.720226in}{1.760328in}}%
\pgfpathlineto{\pgfqpoint{0.720239in}{1.836140in}}%
\pgfpathlineto{\pgfqpoint{0.720969in}{1.813859in}}%
\pgfpathlineto{\pgfqpoint{0.720976in}{1.818939in}}%
\pgfpathlineto{\pgfqpoint{0.720992in}{1.788063in}}%
\pgfpathlineto{\pgfqpoint{0.721740in}{1.815406in}}%
\pgfpathlineto{\pgfqpoint{0.722311in}{1.532215in}}%
\pgfpathlineto{\pgfqpoint{0.722091in}{1.816531in}}%
\pgfpathlineto{\pgfqpoint{0.722707in}{1.532224in}}%
\pgfpathlineto{\pgfqpoint{0.724134in}{1.532620in}}%
\pgfpathlineto{\pgfqpoint{0.724202in}{1.963529in}}%
\pgfpathlineto{\pgfqpoint{0.725131in}{1.809635in}}%
\pgfpathlineto{\pgfqpoint{0.725138in}{1.789053in}}%
\pgfpathlineto{\pgfqpoint{0.725152in}{1.844919in}}%
\pgfpathlineto{\pgfqpoint{0.725900in}{1.817546in}}%
\pgfpathlineto{\pgfqpoint{0.725904in}{1.811705in}}%
\pgfpathlineto{\pgfqpoint{0.725922in}{1.832768in}}%
\pgfpathlineto{\pgfqpoint{0.726663in}{1.829374in}}%
\pgfpathlineto{\pgfqpoint{0.727228in}{1.532241in}}%
\pgfpathlineto{\pgfqpoint{0.727084in}{1.839355in}}%
\pgfpathlineto{\pgfqpoint{0.727771in}{1.532256in}}%
\pgfpathlineto{\pgfqpoint{0.729074in}{1.532454in}}%
\pgfpathlineto{\pgfqpoint{0.729101in}{1.578367in}}%
\pgfpathlineto{\pgfqpoint{0.729157in}{1.985400in}}%
\pgfpathlineto{\pgfqpoint{0.729897in}{1.881623in}}%
\pgfpathlineto{\pgfqpoint{0.729913in}{1.791265in}}%
\pgfpathlineto{\pgfqpoint{0.730684in}{1.836827in}}%
\pgfpathlineto{\pgfqpoint{0.730693in}{1.858346in}}%
\pgfpathlineto{\pgfqpoint{0.730710in}{1.826330in}}%
\pgfpathlineto{\pgfqpoint{0.731452in}{1.843657in}}%
\pgfpathlineto{\pgfqpoint{0.732148in}{1.532260in}}%
\pgfpathlineto{\pgfqpoint{0.732027in}{1.861685in}}%
\pgfpathlineto{\pgfqpoint{0.732732in}{1.532271in}}%
\pgfpathlineto{\pgfqpoint{0.734020in}{1.532484in}}%
\pgfpathlineto{\pgfqpoint{0.734047in}{1.579415in}}%
\pgfpathlineto{\pgfqpoint{0.734103in}{2.011338in}}%
\pgfpathlineto{\pgfqpoint{0.734845in}{1.898034in}}%
\pgfpathlineto{\pgfqpoint{0.734858in}{1.811610in}}%
\pgfpathlineto{\pgfqpoint{0.734903in}{1.900949in}}%
\pgfpathlineto{\pgfqpoint{0.735614in}{1.876201in}}%
\pgfpathlineto{\pgfqpoint{0.735620in}{1.881007in}}%
\pgfpathlineto{\pgfqpoint{0.735636in}{1.846249in}}%
\pgfpathlineto{\pgfqpoint{0.736367in}{1.877927in}}%
\pgfpathlineto{\pgfqpoint{0.737134in}{1.532266in}}%
\pgfpathlineto{\pgfqpoint{0.736973in}{1.887587in}}%
\pgfpathlineto{\pgfqpoint{0.737381in}{1.532280in}}%
\pgfpathlineto{\pgfqpoint{0.738951in}{1.532324in}}%
\pgfpathlineto{\pgfqpoint{0.738962in}{1.532425in}}%
\pgfpathlineto{\pgfqpoint{0.738990in}{1.572608in}}%
\pgfpathlineto{\pgfqpoint{0.739054in}{2.031987in}}%
\pgfpathlineto{\pgfqpoint{0.739787in}{1.865505in}}%
\pgfpathlineto{\pgfqpoint{0.739824in}{1.830440in}}%
\pgfpathlineto{\pgfqpoint{0.739808in}{1.930576in}}%
\pgfpathlineto{\pgfqpoint{0.740554in}{1.892688in}}%
\pgfpathlineto{\pgfqpoint{0.740590in}{1.904583in}}%
\pgfpathlineto{\pgfqpoint{0.740574in}{1.867336in}}%
\pgfpathlineto{\pgfqpoint{0.741320in}{1.891738in}}%
\pgfpathlineto{\pgfqpoint{0.742073in}{1.532269in}}%
\pgfpathlineto{\pgfqpoint{0.741920in}{1.904171in}}%
\pgfpathlineto{\pgfqpoint{0.742320in}{1.532284in}}%
\pgfpathlineto{\pgfqpoint{0.743896in}{1.532327in}}%
\pgfpathlineto{\pgfqpoint{0.743907in}{1.532418in}}%
\pgfpathlineto{\pgfqpoint{0.743926in}{1.544883in}}%
\pgfpathlineto{\pgfqpoint{0.743994in}{2.057301in}}%
\pgfpathlineto{\pgfqpoint{0.744818in}{1.942956in}}%
\pgfpathlineto{\pgfqpoint{0.744831in}{1.859459in}}%
\pgfpathlineto{\pgfqpoint{0.744847in}{1.944300in}}%
\pgfpathlineto{\pgfqpoint{0.745594in}{1.900584in}}%
\pgfpathlineto{\pgfqpoint{0.745604in}{1.926472in}}%
\pgfpathlineto{\pgfqpoint{0.745619in}{1.892998in}}%
\pgfpathlineto{\pgfqpoint{0.746363in}{1.908034in}}%
\pgfpathlineto{\pgfqpoint{0.746788in}{1.924087in}}%
\pgfpathlineto{\pgfqpoint{0.746876in}{1.694834in}}%
\pgfpathlineto{\pgfqpoint{0.746994in}{1.532270in}}%
\pgfpathlineto{\pgfqpoint{0.747613in}{1.532285in}}%
\pgfpathlineto{\pgfqpoint{0.748856in}{1.532503in}}%
\pgfpathlineto{\pgfqpoint{0.748883in}{1.579986in}}%
\pgfpathlineto{\pgfqpoint{0.748948in}{2.105453in}}%
\pgfpathlineto{\pgfqpoint{0.749679in}{1.894000in}}%
\pgfpathlineto{\pgfqpoint{0.749683in}{1.863327in}}%
\pgfpathlineto{\pgfqpoint{0.749699in}{1.988398in}}%
\pgfpathlineto{\pgfqpoint{0.750447in}{1.921335in}}%
\pgfpathlineto{\pgfqpoint{0.750450in}{1.908905in}}%
\pgfpathlineto{\pgfqpoint{0.750467in}{1.953621in}}%
\pgfpathlineto{\pgfqpoint{0.751214in}{1.934634in}}%
\pgfpathlineto{\pgfqpoint{0.751945in}{1.532259in}}%
\pgfpathlineto{\pgfqpoint{0.751811in}{1.945366in}}%
\pgfpathlineto{\pgfqpoint{0.752222in}{1.532264in}}%
\pgfpathlineto{\pgfqpoint{0.753802in}{1.532486in}}%
\pgfpathlineto{\pgfqpoint{0.753829in}{1.579720in}}%
\pgfpathlineto{\pgfqpoint{0.753885in}{2.128007in}}%
\pgfpathlineto{\pgfqpoint{0.754630in}{1.955682in}}%
\pgfpathlineto{\pgfqpoint{0.754640in}{1.924152in}}%
\pgfpathlineto{\pgfqpoint{0.754655in}{1.967966in}}%
\pgfpathlineto{\pgfqpoint{0.755398in}{1.953796in}}%
\pgfpathlineto{\pgfqpoint{0.755921in}{1.963524in}}%
\pgfpathlineto{\pgfqpoint{0.755426in}{1.942423in}}%
\pgfpathlineto{\pgfqpoint{0.756168in}{1.958363in}}%
\pgfpathlineto{\pgfqpoint{0.756757in}{1.970354in}}%
\pgfpathlineto{\pgfqpoint{0.756761in}{1.877534in}}%
\pgfpathlineto{\pgfqpoint{0.756889in}{1.532244in}}%
\pgfpathlineto{\pgfqpoint{0.757518in}{1.532259in}}%
\pgfpathlineto{\pgfqpoint{0.758747in}{1.532459in}}%
\pgfpathlineto{\pgfqpoint{0.758774in}{1.578816in}}%
\pgfpathlineto{\pgfqpoint{0.758834in}{2.148909in}}%
\pgfpathlineto{\pgfqpoint{0.759577in}{1.953708in}}%
\pgfpathlineto{\pgfqpoint{0.759584in}{1.904592in}}%
\pgfpathlineto{\pgfqpoint{0.759627in}{2.026023in}}%
\pgfpathlineto{\pgfqpoint{0.760347in}{1.953393in}}%
\pgfpathlineto{\pgfqpoint{0.760360in}{1.993077in}}%
\pgfpathlineto{\pgfqpoint{0.760375in}{1.950335in}}%
\pgfpathlineto{\pgfqpoint{0.761123in}{1.980275in}}%
\pgfpathlineto{\pgfqpoint{0.761702in}{1.990844in}}%
\pgfpathlineto{\pgfqpoint{0.761707in}{1.891151in}}%
\pgfpathlineto{\pgfqpoint{0.761826in}{1.532220in}}%
\pgfpathlineto{\pgfqpoint{0.762459in}{1.532240in}}%
\pgfpathlineto{\pgfqpoint{0.763682in}{1.532294in}}%
\pgfpathlineto{\pgfqpoint{0.763690in}{1.532365in}}%
\pgfpathlineto{\pgfqpoint{0.763718in}{1.574267in}}%
\pgfpathlineto{\pgfqpoint{0.763781in}{2.152168in}}%
\pgfpathlineto{\pgfqpoint{0.764529in}{2.014726in}}%
\pgfpathlineto{\pgfqpoint{0.764532in}{2.031666in}}%
\pgfpathlineto{\pgfqpoint{0.764547in}{1.937270in}}%
\pgfpathlineto{\pgfqpoint{0.765293in}{2.002202in}}%
\pgfpathlineto{\pgfqpoint{0.765307in}{1.972594in}}%
\pgfpathlineto{\pgfqpoint{0.765319in}{2.004612in}}%
\pgfpathlineto{\pgfqpoint{0.766066in}{1.995302in}}%
\pgfpathlineto{\pgfqpoint{0.766840in}{1.532196in}}%
\pgfpathlineto{\pgfqpoint{0.766648in}{2.006036in}}%
\pgfpathlineto{\pgfqpoint{0.767186in}{1.532205in}}%
\pgfpathlineto{\pgfqpoint{0.768638in}{1.532375in}}%
\pgfpathlineto{\pgfqpoint{0.768665in}{1.576576in}}%
\pgfpathlineto{\pgfqpoint{0.768816in}{2.146182in}}%
\pgfpathlineto{\pgfqpoint{0.769467in}{1.962799in}}%
\pgfpathlineto{\pgfqpoint{0.769471in}{1.943942in}}%
\pgfpathlineto{\pgfqpoint{0.769486in}{2.051929in}}%
\pgfpathlineto{\pgfqpoint{0.770234in}{1.983844in}}%
\pgfpathlineto{\pgfqpoint{0.770247in}{2.022409in}}%
\pgfpathlineto{\pgfqpoint{0.771011in}{2.012214in}}%
\pgfpathlineto{\pgfqpoint{0.771732in}{1.532162in}}%
\pgfpathlineto{\pgfqpoint{0.771593in}{2.017969in}}%
\pgfpathlineto{\pgfqpoint{0.772125in}{1.532170in}}%
\pgfpathlineto{\pgfqpoint{0.773591in}{1.532717in}}%
\pgfpathlineto{\pgfqpoint{0.773671in}{2.131086in}}%
\pgfpathlineto{\pgfqpoint{0.774696in}{2.014077in}}%
\pgfpathlineto{\pgfqpoint{0.774703in}{1.990691in}}%
\pgfpathlineto{\pgfqpoint{0.774716in}{2.032075in}}%
\pgfpathlineto{\pgfqpoint{0.775459in}{2.011346in}}%
\pgfpathlineto{\pgfqpoint{0.775598in}{2.022570in}}%
\pgfpathlineto{\pgfqpoint{0.775486in}{2.007032in}}%
\pgfpathlineto{\pgfqpoint{0.776216in}{2.015711in}}%
\pgfpathlineto{\pgfqpoint{0.776536in}{2.027099in}}%
\pgfpathlineto{\pgfqpoint{0.776540in}{2.008361in}}%
\pgfpathlineto{\pgfqpoint{0.776683in}{1.532124in}}%
\pgfpathlineto{\pgfqpoint{0.777326in}{1.532149in}}%
\pgfpathlineto{\pgfqpoint{0.778521in}{1.532194in}}%
\pgfpathlineto{\pgfqpoint{0.778530in}{1.532276in}}%
\pgfpathlineto{\pgfqpoint{0.778572in}{1.740981in}}%
\pgfpathlineto{\pgfqpoint{0.778627in}{2.116010in}}%
\pgfpathlineto{\pgfqpoint{0.779360in}{2.051577in}}%
\pgfpathlineto{\pgfqpoint{0.779372in}{1.988860in}}%
\pgfpathlineto{\pgfqpoint{0.780133in}{2.034072in}}%
\pgfpathlineto{\pgfqpoint{0.780148in}{2.010264in}}%
\pgfpathlineto{\pgfqpoint{0.780898in}{2.022202in}}%
\pgfpathlineto{\pgfqpoint{0.781484in}{2.035291in}}%
\pgfpathlineto{\pgfqpoint{0.781489in}{1.937784in}}%
\pgfpathlineto{\pgfqpoint{0.781601in}{1.532088in}}%
\pgfpathlineto{\pgfqpoint{0.782234in}{1.532110in}}%
\pgfpathlineto{\pgfqpoint{0.783475in}{1.532210in}}%
\pgfpathlineto{\pgfqpoint{0.783502in}{1.575596in}}%
\pgfpathlineto{\pgfqpoint{0.783560in}{2.087630in}}%
\pgfpathlineto{\pgfqpoint{0.784312in}{2.023623in}}%
\pgfpathlineto{\pgfqpoint{0.784348in}{2.011395in}}%
\pgfpathlineto{\pgfqpoint{0.784365in}{2.039810in}}%
\pgfpathlineto{\pgfqpoint{0.785078in}{2.022215in}}%
\pgfpathlineto{\pgfqpoint{0.785089in}{2.022125in}}%
\pgfpathlineto{\pgfqpoint{0.785122in}{2.019677in}}%
\pgfpathlineto{\pgfqpoint{0.785823in}{2.033723in}}%
\pgfpathlineto{\pgfqpoint{0.786430in}{2.040222in}}%
\pgfpathlineto{\pgfqpoint{0.786445in}{1.672120in}}%
\pgfpathlineto{\pgfqpoint{0.786522in}{1.532046in}}%
\pgfpathlineto{\pgfqpoint{0.787199in}{1.532061in}}%
\pgfpathlineto{\pgfqpoint{0.788425in}{1.532202in}}%
\pgfpathlineto{\pgfqpoint{0.788519in}{2.065444in}}%
\pgfpathlineto{\pgfqpoint{0.789653in}{2.031174in}}%
\pgfpathlineto{\pgfqpoint{0.790048in}{2.023116in}}%
\pgfpathlineto{\pgfqpoint{0.789999in}{2.031873in}}%
\pgfpathlineto{\pgfqpoint{0.790395in}{2.023445in}}%
\pgfpathlineto{\pgfqpoint{0.790444in}{2.031463in}}%
\pgfpathlineto{\pgfqpoint{0.790840in}{1.986745in}}%
\pgfpathlineto{\pgfqpoint{0.791441in}{1.532033in}}%
\pgfpathlineto{\pgfqpoint{0.791917in}{1.532043in}}%
\pgfpathlineto{\pgfqpoint{0.793370in}{1.532137in}}%
\pgfpathlineto{\pgfqpoint{0.793483in}{2.052931in}}%
\pgfpathlineto{\pgfqpoint{0.794772in}{2.028751in}}%
\pgfpathlineto{\pgfqpoint{0.795265in}{1.728734in}}%
\pgfpathlineto{\pgfqpoint{0.795463in}{1.545695in}}%
\pgfpathlineto{\pgfqpoint{0.795809in}{2.018690in}}%
\pgfpathlineto{\pgfqpoint{0.796106in}{1.564805in}}%
\pgfpathlineto{\pgfqpoint{0.796382in}{1.532034in}}%
\pgfpathlineto{\pgfqpoint{0.796319in}{1.657928in}}%
\pgfpathlineto{\pgfqpoint{0.796959in}{1.532046in}}%
\pgfpathlineto{\pgfqpoint{0.798319in}{1.532223in}}%
\pgfpathlineto{\pgfqpoint{0.798436in}{2.050169in}}%
\pgfpathlineto{\pgfqpoint{0.799492in}{2.021280in}}%
\pgfpathlineto{\pgfqpoint{0.799700in}{2.027857in}}%
\pgfpathlineto{\pgfqpoint{0.800081in}{1.999202in}}%
\pgfpathlineto{\pgfqpoint{0.800427in}{1.551007in}}%
\pgfpathlineto{\pgfqpoint{0.800774in}{2.007517in}}%
\pgfpathlineto{\pgfqpoint{0.801169in}{1.558448in}}%
\pgfpathlineto{\pgfqpoint{0.801264in}{1.649117in}}%
\pgfpathlineto{\pgfqpoint{0.801328in}{1.532033in}}%
\pgfpathlineto{\pgfqpoint{0.801858in}{1.532046in}}%
\pgfpathlineto{\pgfqpoint{0.803261in}{1.532145in}}%
\pgfpathlineto{\pgfqpoint{0.803390in}{2.050159in}}%
\pgfpathlineto{\pgfqpoint{0.804603in}{2.020832in}}%
\pgfpathlineto{\pgfqpoint{0.805365in}{1.547393in}}%
\pgfpathlineto{\pgfqpoint{0.804621in}{2.026630in}}%
\pgfpathlineto{\pgfqpoint{0.805612in}{1.885608in}}%
\pgfpathlineto{\pgfqpoint{0.805711in}{2.013489in}}%
\pgfpathlineto{\pgfqpoint{0.806057in}{1.551007in}}%
\pgfpathlineto{\pgfqpoint{0.806215in}{1.584102in}}%
\pgfpathlineto{\pgfqpoint{0.806270in}{1.532031in}}%
\pgfpathlineto{\pgfqpoint{0.806947in}{1.532046in}}%
\pgfpathlineto{\pgfqpoint{0.808210in}{1.532301in}}%
\pgfpathlineto{\pgfqpoint{0.808335in}{2.049313in}}%
\pgfpathlineto{\pgfqpoint{0.809333in}{2.025859in}}%
\pgfpathlineto{\pgfqpoint{0.810298in}{1.552396in}}%
\pgfpathlineto{\pgfqpoint{0.809402in}{2.026099in}}%
\pgfpathlineto{\pgfqpoint{0.810446in}{1.640671in}}%
\pgfpathlineto{\pgfqpoint{0.810693in}{1.996010in}}%
\pgfpathlineto{\pgfqpoint{0.811165in}{1.548978in}}%
\pgfpathlineto{\pgfqpoint{0.811213in}{1.532031in}}%
\pgfpathlineto{\pgfqpoint{0.811937in}{1.532051in}}%
\pgfpathlineto{\pgfqpoint{0.813147in}{1.532102in}}%
\pgfpathlineto{\pgfqpoint{0.813156in}{1.532314in}}%
\pgfpathlineto{\pgfqpoint{0.813280in}{2.047780in}}%
\pgfpathlineto{\pgfqpoint{0.814272in}{2.020261in}}%
\pgfpathlineto{\pgfqpoint{0.815276in}{1.561127in}}%
\pgfpathlineto{\pgfqpoint{0.814353in}{2.025011in}}%
\pgfpathlineto{\pgfqpoint{0.815473in}{1.795672in}}%
\pgfpathlineto{\pgfqpoint{0.815622in}{1.985724in}}%
\pgfpathlineto{\pgfqpoint{0.816113in}{1.541814in}}%
\pgfpathlineto{\pgfqpoint{0.816149in}{1.532032in}}%
\pgfpathlineto{\pgfqpoint{0.816849in}{1.532050in}}%
\pgfpathlineto{\pgfqpoint{0.818096in}{1.532129in}}%
\pgfpathlineto{\pgfqpoint{0.818214in}{2.044554in}}%
\pgfpathlineto{\pgfqpoint{0.819487in}{2.017696in}}%
\pgfpathlineto{\pgfqpoint{0.819784in}{2.019418in}}%
\pgfpathlineto{\pgfqpoint{0.819833in}{2.016062in}}%
\pgfpathlineto{\pgfqpoint{0.821093in}{1.532029in}}%
\pgfpathlineto{\pgfqpoint{0.821658in}{1.532041in}}%
\pgfpathlineto{\pgfqpoint{0.823041in}{1.532122in}}%
\pgfpathlineto{\pgfqpoint{0.823083in}{1.754804in}}%
\pgfpathlineto{\pgfqpoint{0.823269in}{2.050940in}}%
\pgfpathlineto{\pgfqpoint{0.823867in}{2.025102in}}%
\pgfpathlineto{\pgfqpoint{0.823872in}{2.033508in}}%
\pgfpathlineto{\pgfqpoint{0.823887in}{2.001544in}}%
\pgfpathlineto{\pgfqpoint{0.824632in}{2.019870in}}%
\pgfpathlineto{\pgfqpoint{0.824640in}{2.021281in}}%
\pgfpathlineto{\pgfqpoint{0.824855in}{1.976918in}}%
\pgfpathlineto{\pgfqpoint{0.825151in}{1.555212in}}%
\pgfpathlineto{\pgfqpoint{0.825547in}{1.992338in}}%
\pgfpathlineto{\pgfqpoint{0.825893in}{1.557283in}}%
\pgfpathlineto{\pgfqpoint{0.825992in}{1.616339in}}%
\pgfpathlineto{\pgfqpoint{0.826008in}{1.532013in}}%
\pgfpathlineto{\pgfqpoint{0.826551in}{1.532035in}}%
\pgfpathlineto{\pgfqpoint{0.827992in}{1.532279in}}%
\pgfpathlineto{\pgfqpoint{0.828176in}{2.042501in}}%
\pgfpathlineto{\pgfqpoint{0.829130in}{2.020224in}}%
\pgfpathlineto{\pgfqpoint{0.830108in}{1.550026in}}%
\pgfpathlineto{\pgfqpoint{0.829157in}{2.020920in}}%
\pgfpathlineto{\pgfqpoint{0.830256in}{1.647066in}}%
\pgfpathlineto{\pgfqpoint{0.830454in}{2.002074in}}%
\pgfpathlineto{\pgfqpoint{0.830968in}{1.532029in}}%
\pgfpathlineto{\pgfqpoint{0.832287in}{1.532057in}}%
\pgfpathlineto{\pgfqpoint{0.832934in}{1.532144in}}%
\pgfpathlineto{\pgfqpoint{0.833168in}{2.044050in}}%
\pgfpathlineto{\pgfqpoint{0.834202in}{2.018383in}}%
\pgfpathlineto{\pgfqpoint{0.835052in}{1.554889in}}%
\pgfpathlineto{\pgfqpoint{0.834262in}{2.020543in}}%
\pgfpathlineto{\pgfqpoint{0.835200in}{1.637268in}}%
\pgfpathlineto{\pgfqpoint{0.835447in}{1.993496in}}%
\pgfpathlineto{\pgfqpoint{0.835900in}{1.532019in}}%
\pgfpathlineto{\pgfqpoint{0.835908in}{1.532052in}}%
\pgfpathlineto{\pgfqpoint{0.836654in}{1.532040in}}%
\pgfpathlineto{\pgfqpoint{0.836703in}{1.532058in}}%
\pgfpathlineto{\pgfqpoint{0.837875in}{1.532104in}}%
\pgfpathlineto{\pgfqpoint{0.837893in}{1.545349in}}%
\pgfpathlineto{\pgfqpoint{0.838102in}{2.043053in}}%
\pgfpathlineto{\pgfqpoint{0.838829in}{2.021845in}}%
\pgfpathlineto{\pgfqpoint{0.838841in}{2.001030in}}%
\pgfpathlineto{\pgfqpoint{0.838857in}{2.024816in}}%
\pgfpathlineto{\pgfqpoint{0.839602in}{2.011410in}}%
\pgfpathlineto{\pgfqpoint{0.839894in}{1.633653in}}%
\pgfpathlineto{\pgfqpoint{0.840042in}{1.558172in}}%
\pgfpathlineto{\pgfqpoint{0.840339in}{1.973963in}}%
\pgfpathlineto{\pgfqpoint{0.840388in}{1.992243in}}%
\pgfpathlineto{\pgfqpoint{0.840537in}{1.769634in}}%
\pgfpathlineto{\pgfqpoint{0.840845in}{1.532022in}}%
\pgfpathlineto{\pgfqpoint{0.841399in}{1.532036in}}%
\pgfpathlineto{\pgfqpoint{0.842822in}{1.532111in}}%
\pgfpathlineto{\pgfqpoint{0.842862in}{1.707778in}}%
\pgfpathlineto{\pgfqpoint{0.842945in}{2.038441in}}%
\pgfpathlineto{\pgfqpoint{0.843646in}{2.015109in}}%
\pgfpathlineto{\pgfqpoint{0.843654in}{2.019012in}}%
\pgfpathlineto{\pgfqpoint{0.843672in}{2.004675in}}%
\pgfpathlineto{\pgfqpoint{0.844422in}{2.017242in}}%
\pgfpathlineto{\pgfqpoint{0.844768in}{1.775190in}}%
\pgfpathlineto{\pgfqpoint{0.844966in}{1.564711in}}%
\pgfpathlineto{\pgfqpoint{0.845312in}{1.974795in}}%
\pgfpathlineto{\pgfqpoint{0.845609in}{1.592590in}}%
\pgfpathlineto{\pgfqpoint{0.845791in}{1.532012in}}%
\pgfpathlineto{\pgfqpoint{0.845773in}{1.610546in}}%
\pgfpathlineto{\pgfqpoint{0.846497in}{1.532057in}}%
\pgfpathlineto{\pgfqpoint{0.847765in}{1.532098in}}%
\pgfpathlineto{\pgfqpoint{0.847774in}{1.532331in}}%
\pgfpathlineto{\pgfqpoint{0.848027in}{2.035181in}}%
\pgfpathlineto{\pgfqpoint{0.848872in}{2.018415in}}%
\pgfpathlineto{\pgfqpoint{0.849911in}{1.557650in}}%
\pgfpathlineto{\pgfqpoint{0.850059in}{1.654901in}}%
\pgfpathlineto{\pgfqpoint{0.850257in}{1.986682in}}%
\pgfpathlineto{\pgfqpoint{0.850749in}{1.532028in}}%
\pgfpathlineto{\pgfqpoint{0.850763in}{1.532039in}}%
\pgfpathlineto{\pgfqpoint{0.852653in}{1.532070in}}%
\pgfpathlineto{\pgfqpoint{0.852716in}{1.532147in}}%
\pgfpathlineto{\pgfqpoint{0.852950in}{2.038931in}}%
\pgfpathlineto{\pgfqpoint{0.853960in}{2.003397in}}%
\pgfpathlineto{\pgfqpoint{0.853979in}{2.015858in}}%
\pgfpathlineto{\pgfqpoint{0.854611in}{1.896007in}}%
\pgfpathlineto{\pgfqpoint{0.854858in}{1.557431in}}%
\pgfpathlineto{\pgfqpoint{0.855254in}{1.986832in}}%
\pgfpathlineto{\pgfqpoint{0.855501in}{1.593515in}}%
\pgfpathlineto{\pgfqpoint{0.855738in}{1.532028in}}%
\pgfpathlineto{\pgfqpoint{0.856373in}{1.532051in}}%
\pgfpathlineto{\pgfqpoint{0.857661in}{1.532142in}}%
\pgfpathlineto{\pgfqpoint{0.857991in}{2.035977in}}%
\pgfpathlineto{\pgfqpoint{0.858919in}{2.010524in}}%
\pgfpathlineto{\pgfqpoint{0.859828in}{1.555057in}}%
\pgfpathlineto{\pgfqpoint{0.858949in}{2.014090in}}%
\pgfpathlineto{\pgfqpoint{0.859976in}{1.676416in}}%
\pgfpathlineto{\pgfqpoint{0.860174in}{1.996114in}}%
\pgfpathlineto{\pgfqpoint{0.860626in}{1.532012in}}%
\pgfpathlineto{\pgfqpoint{0.860664in}{1.532044in}}%
\pgfpathlineto{\pgfqpoint{0.862129in}{1.532075in}}%
\pgfpathlineto{\pgfqpoint{0.862602in}{1.532100in}}%
\pgfpathlineto{\pgfqpoint{0.862611in}{1.532465in}}%
\pgfpathlineto{\pgfqpoint{0.862733in}{2.033154in}}%
\pgfpathlineto{\pgfqpoint{0.863685in}{2.007689in}}%
\pgfpathlineto{\pgfqpoint{0.864771in}{1.563293in}}%
\pgfpathlineto{\pgfqpoint{0.863712in}{2.014457in}}%
\pgfpathlineto{\pgfqpoint{0.864821in}{1.572858in}}%
\pgfpathlineto{\pgfqpoint{0.865117in}{1.978318in}}%
\pgfpathlineto{\pgfqpoint{0.865562in}{1.552948in}}%
\pgfpathlineto{\pgfqpoint{0.865632in}{1.532029in}}%
\pgfpathlineto{\pgfqpoint{0.866318in}{1.532044in}}%
\pgfpathlineto{\pgfqpoint{0.867553in}{1.532150in}}%
\pgfpathlineto{\pgfqpoint{0.867794in}{2.037001in}}%
\pgfpathlineto{\pgfqpoint{0.868779in}{2.004415in}}%
\pgfpathlineto{\pgfqpoint{0.869710in}{1.555353in}}%
\pgfpathlineto{\pgfqpoint{0.868838in}{2.012689in}}%
\pgfpathlineto{\pgfqpoint{0.869859in}{1.640340in}}%
\pgfpathlineto{\pgfqpoint{0.870106in}{1.991328in}}%
\pgfpathlineto{\pgfqpoint{0.870517in}{1.532015in}}%
\pgfpathlineto{\pgfqpoint{0.870560in}{1.532043in}}%
\pgfpathlineto{\pgfqpoint{0.872179in}{1.532058in}}%
\pgfpathlineto{\pgfqpoint{0.872501in}{1.532225in}}%
\pgfpathlineto{\pgfqpoint{0.872610in}{2.030100in}}%
\pgfpathlineto{\pgfqpoint{0.873673in}{2.008425in}}%
\pgfpathlineto{\pgfqpoint{0.874493in}{1.725823in}}%
\pgfpathlineto{\pgfqpoint{0.874690in}{1.556366in}}%
\pgfpathlineto{\pgfqpoint{0.875037in}{1.995388in}}%
\pgfpathlineto{\pgfqpoint{0.875333in}{1.571121in}}%
\pgfpathlineto{\pgfqpoint{0.875463in}{1.532018in}}%
\pgfpathlineto{\pgfqpoint{0.875448in}{1.580198in}}%
\pgfpathlineto{\pgfqpoint{0.876235in}{1.532039in}}%
\pgfpathlineto{\pgfqpoint{0.877445in}{1.532195in}}%
\pgfpathlineto{\pgfqpoint{0.877571in}{2.028616in}}%
\pgfpathlineto{\pgfqpoint{0.878638in}{1.998556in}}%
\pgfpathlineto{\pgfqpoint{0.879623in}{1.555919in}}%
\pgfpathlineto{\pgfqpoint{0.878654in}{2.006927in}}%
\pgfpathlineto{\pgfqpoint{0.879821in}{1.752315in}}%
\pgfpathlineto{\pgfqpoint{0.879969in}{1.991240in}}%
\pgfpathlineto{\pgfqpoint{0.880409in}{1.532027in}}%
\pgfpathlineto{\pgfqpoint{0.880457in}{1.532028in}}%
\pgfpathlineto{\pgfqpoint{0.882389in}{1.532147in}}%
\pgfpathlineto{\pgfqpoint{0.882630in}{2.033179in}}%
\pgfpathlineto{\pgfqpoint{0.883615in}{2.002322in}}%
\pgfpathlineto{\pgfqpoint{0.884555in}{1.561595in}}%
\pgfpathlineto{\pgfqpoint{0.883676in}{2.008852in}}%
\pgfpathlineto{\pgfqpoint{0.884654in}{1.586549in}}%
\pgfpathlineto{\pgfqpoint{0.884951in}{1.982531in}}%
\pgfpathlineto{\pgfqpoint{0.885355in}{1.532021in}}%
\pgfpathlineto{\pgfqpoint{0.885372in}{1.532025in}}%
\pgfpathlineto{\pgfqpoint{0.886552in}{1.532063in}}%
\pgfpathlineto{\pgfqpoint{0.887324in}{1.532084in}}%
\pgfpathlineto{\pgfqpoint{0.887334in}{1.532146in}}%
\pgfpathlineto{\pgfqpoint{0.887676in}{2.030671in}}%
\pgfpathlineto{\pgfqpoint{0.888572in}{2.007698in}}%
\pgfpathlineto{\pgfqpoint{0.889548in}{1.554851in}}%
\pgfpathlineto{\pgfqpoint{0.889647in}{1.617345in}}%
\pgfpathlineto{\pgfqpoint{0.889894in}{1.994618in}}%
\pgfpathlineto{\pgfqpoint{0.890333in}{1.532028in}}%
\pgfpathlineto{\pgfqpoint{0.890363in}{1.532036in}}%
\pgfpathlineto{\pgfqpoint{0.892278in}{1.532125in}}%
\pgfpathlineto{\pgfqpoint{0.892394in}{2.025187in}}%
\pgfpathlineto{\pgfqpoint{0.893694in}{1.996996in}}%
\pgfpathlineto{\pgfqpoint{0.894464in}{1.557860in}}%
\pgfpathlineto{\pgfqpoint{0.893714in}{2.001497in}}%
\pgfpathlineto{\pgfqpoint{0.894712in}{1.845756in}}%
\pgfpathlineto{\pgfqpoint{0.894810in}{1.985719in}}%
\pgfpathlineto{\pgfqpoint{0.895247in}{1.532026in}}%
\pgfpathlineto{\pgfqpoint{0.895281in}{1.532038in}}%
\pgfpathlineto{\pgfqpoint{0.897214in}{1.532083in}}%
\pgfpathlineto{\pgfqpoint{0.897226in}{1.532158in}}%
\pgfpathlineto{\pgfqpoint{0.897566in}{2.030923in}}%
\pgfpathlineto{\pgfqpoint{0.898426in}{2.007888in}}%
\pgfpathlineto{\pgfqpoint{0.899422in}{1.565375in}}%
\pgfpathlineto{\pgfqpoint{0.899521in}{1.603220in}}%
\pgfpathlineto{\pgfqpoint{0.899768in}{1.973088in}}%
\pgfpathlineto{\pgfqpoint{0.900189in}{1.532006in}}%
\pgfpathlineto{\pgfqpoint{0.900221in}{1.532043in}}%
\pgfpathlineto{\pgfqpoint{0.901530in}{1.532051in}}%
\pgfpathlineto{\pgfqpoint{0.902175in}{1.532380in}}%
\pgfpathlineto{\pgfqpoint{0.902284in}{2.022274in}}%
\pgfpathlineto{\pgfqpoint{0.903262in}{2.003410in}}%
\pgfpathlineto{\pgfqpoint{0.904360in}{1.569714in}}%
\pgfpathlineto{\pgfqpoint{0.904410in}{1.573265in}}%
\pgfpathlineto{\pgfqpoint{0.904756in}{1.964658in}}%
\pgfpathlineto{\pgfqpoint{0.905127in}{1.548665in}}%
\pgfpathlineto{\pgfqpoint{0.905180in}{1.532027in}}%
\pgfpathlineto{\pgfqpoint{0.905861in}{1.532042in}}%
\pgfpathlineto{\pgfqpoint{0.907116in}{1.532139in}}%
\pgfpathlineto{\pgfqpoint{0.907453in}{2.024786in}}%
\pgfpathlineto{\pgfqpoint{0.908381in}{2.001609in}}%
\pgfpathlineto{\pgfqpoint{0.908447in}{2.002472in}}%
\pgfpathlineto{\pgfqpoint{0.908399in}{1.986909in}}%
\pgfpathlineto{\pgfqpoint{0.908964in}{1.995325in}}%
\pgfpathlineto{\pgfqpoint{0.910080in}{1.532020in}}%
\pgfpathlineto{\pgfqpoint{0.910294in}{1.532029in}}%
\pgfpathlineto{\pgfqpoint{0.912059in}{1.532118in}}%
\pgfpathlineto{\pgfqpoint{0.912100in}{1.727635in}}%
\pgfpathlineto{\pgfqpoint{0.912288in}{2.026188in}}%
\pgfpathlineto{\pgfqpoint{0.912882in}{1.982142in}}%
\pgfpathlineto{\pgfqpoint{0.912925in}{2.007151in}}%
\pgfpathlineto{\pgfqpoint{0.912909in}{1.980562in}}%
\pgfpathlineto{\pgfqpoint{0.913635in}{1.996021in}}%
\pgfpathlineto{\pgfqpoint{0.913849in}{1.999190in}}%
\pgfpathlineto{\pgfqpoint{0.913948in}{1.960725in}}%
\pgfpathlineto{\pgfqpoint{0.915056in}{1.532027in}}%
\pgfpathlineto{\pgfqpoint{0.914640in}{1.960950in}}%
\pgfpathlineto{\pgfqpoint{0.915177in}{1.532038in}}%
\pgfpathlineto{\pgfqpoint{0.917005in}{1.532118in}}%
\pgfpathlineto{\pgfqpoint{0.917046in}{1.738740in}}%
\pgfpathlineto{\pgfqpoint{0.917231in}{2.023749in}}%
\pgfpathlineto{\pgfqpoint{0.917831in}{1.976293in}}%
\pgfpathlineto{\pgfqpoint{0.917847in}{2.008849in}}%
\pgfpathlineto{\pgfqpoint{0.918607in}{1.996021in}}%
\pgfpathlineto{\pgfqpoint{0.919228in}{1.569107in}}%
\pgfpathlineto{\pgfqpoint{0.918646in}{1.997628in}}%
\pgfpathlineto{\pgfqpoint{0.919476in}{1.859280in}}%
\pgfpathlineto{\pgfqpoint{0.919575in}{1.968564in}}%
\pgfpathlineto{\pgfqpoint{0.919971in}{1.532009in}}%
\pgfpathlineto{\pgfqpoint{0.920016in}{1.532024in}}%
\pgfpathlineto{\pgfqpoint{0.921949in}{1.532111in}}%
\pgfpathlineto{\pgfqpoint{0.921991in}{1.723376in}}%
\pgfpathlineto{\pgfqpoint{0.922171in}{2.021979in}}%
\pgfpathlineto{\pgfqpoint{0.922773in}{2.001976in}}%
\pgfpathlineto{\pgfqpoint{0.922795in}{1.980389in}}%
\pgfpathlineto{\pgfqpoint{0.922808in}{2.003592in}}%
\pgfpathlineto{\pgfqpoint{0.923584in}{1.985112in}}%
\pgfpathlineto{\pgfqpoint{0.923633in}{1.995439in}}%
\pgfpathlineto{\pgfqpoint{0.923880in}{1.936053in}}%
\pgfpathlineto{\pgfqpoint{0.924917in}{1.532008in}}%
\pgfpathlineto{\pgfqpoint{0.924523in}{1.951592in}}%
\pgfpathlineto{\pgfqpoint{0.924942in}{1.532042in}}%
\pgfpathlineto{\pgfqpoint{0.926142in}{1.532050in}}%
\pgfpathlineto{\pgfqpoint{0.926903in}{1.532753in}}%
\pgfpathlineto{\pgfqpoint{0.927009in}{2.015869in}}%
\pgfpathlineto{\pgfqpoint{0.927962in}{1.991342in}}%
\pgfpathlineto{\pgfqpoint{0.927972in}{1.997952in}}%
\pgfpathlineto{\pgfqpoint{0.927990in}{1.983056in}}%
\pgfpathlineto{\pgfqpoint{0.928679in}{1.993226in}}%
\pgfpathlineto{\pgfqpoint{0.928778in}{1.994005in}}%
\pgfpathlineto{\pgfqpoint{0.928877in}{1.877472in}}%
\pgfpathlineto{\pgfqpoint{0.929124in}{1.560067in}}%
\pgfpathlineto{\pgfqpoint{0.929470in}{1.982695in}}%
\pgfpathlineto{\pgfqpoint{0.929767in}{1.589373in}}%
\pgfpathlineto{\pgfqpoint{0.929862in}{1.532010in}}%
\pgfpathlineto{\pgfqpoint{0.930611in}{1.532038in}}%
\pgfpathlineto{\pgfqpoint{0.931845in}{1.532183in}}%
\pgfpathlineto{\pgfqpoint{0.931967in}{2.015265in}}%
\pgfpathlineto{\pgfqpoint{0.933040in}{1.984983in}}%
\pgfpathlineto{\pgfqpoint{0.933048in}{1.988768in}}%
\pgfpathlineto{\pgfqpoint{0.933056in}{1.993864in}}%
\pgfpathlineto{\pgfqpoint{0.933765in}{1.967742in}}%
\pgfpathlineto{\pgfqpoint{0.934827in}{1.532024in}}%
\pgfpathlineto{\pgfqpoint{0.934422in}{1.979043in}}%
\pgfpathlineto{\pgfqpoint{0.934881in}{1.532025in}}%
\pgfpathlineto{\pgfqpoint{0.936787in}{1.532123in}}%
\pgfpathlineto{\pgfqpoint{0.937005in}{2.015304in}}%
\pgfpathlineto{\pgfqpoint{0.938169in}{1.986392in}}%
\pgfpathlineto{\pgfqpoint{0.939008in}{1.564609in}}%
\pgfpathlineto{\pgfqpoint{0.938191in}{1.992180in}}%
\pgfpathlineto{\pgfqpoint{0.939206in}{1.719372in}}%
\pgfpathlineto{\pgfqpoint{0.939403in}{1.978294in}}%
\pgfpathlineto{\pgfqpoint{0.939772in}{1.532026in}}%
\pgfpathlineto{\pgfqpoint{0.939868in}{1.532037in}}%
\pgfpathlineto{\pgfqpoint{0.941734in}{1.532142in}}%
\pgfpathlineto{\pgfqpoint{0.942074in}{2.020450in}}%
\pgfpathlineto{\pgfqpoint{0.942958in}{1.991154in}}%
\pgfpathlineto{\pgfqpoint{0.942981in}{1.995414in}}%
\pgfpathlineto{\pgfqpoint{0.943916in}{1.589775in}}%
\pgfpathlineto{\pgfqpoint{0.943965in}{1.573234in}}%
\pgfpathlineto{\pgfqpoint{0.944163in}{1.740850in}}%
\pgfpathlineto{\pgfqpoint{0.944312in}{1.958426in}}%
\pgfpathlineto{\pgfqpoint{0.944698in}{1.532005in}}%
\pgfpathlineto{\pgfqpoint{0.944791in}{1.532040in}}%
\pgfpathlineto{\pgfqpoint{0.946675in}{1.532101in}}%
\pgfpathlineto{\pgfqpoint{0.946693in}{1.545108in}}%
\pgfpathlineto{\pgfqpoint{0.946790in}{2.011454in}}%
\pgfpathlineto{\pgfqpoint{0.947615in}{1.983480in}}%
\pgfpathlineto{\pgfqpoint{0.947623in}{1.976957in}}%
\pgfpathlineto{\pgfqpoint{0.947670in}{1.994644in}}%
\pgfpathlineto{\pgfqpoint{0.948341in}{1.988871in}}%
\pgfpathlineto{\pgfqpoint{0.948440in}{1.990411in}}%
\pgfpathlineto{\pgfqpoint{0.948687in}{1.856506in}}%
\pgfpathlineto{\pgfqpoint{0.948934in}{1.572575in}}%
\pgfpathlineto{\pgfqpoint{0.949280in}{1.962961in}}%
\pgfpathlineto{\pgfqpoint{0.949577in}{1.594992in}}%
\pgfpathlineto{\pgfqpoint{0.949693in}{1.532026in}}%
\pgfpathlineto{\pgfqpoint{0.950476in}{1.532044in}}%
\pgfpathlineto{\pgfqpoint{0.951625in}{1.532139in}}%
\pgfpathlineto{\pgfqpoint{0.951966in}{2.017186in}}%
\pgfpathlineto{\pgfqpoint{0.952853in}{1.978120in}}%
\pgfpathlineto{\pgfqpoint{0.953898in}{1.574688in}}%
\pgfpathlineto{\pgfqpoint{0.952874in}{1.993717in}}%
\pgfpathlineto{\pgfqpoint{0.954046in}{1.706518in}}%
\pgfpathlineto{\pgfqpoint{0.954244in}{1.961445in}}%
\pgfpathlineto{\pgfqpoint{0.954605in}{1.532025in}}%
\pgfpathlineto{\pgfqpoint{0.954686in}{1.532037in}}%
\pgfpathlineto{\pgfqpoint{0.956560in}{1.532082in}}%
\pgfpathlineto{\pgfqpoint{0.956571in}{1.532142in}}%
\pgfpathlineto{\pgfqpoint{0.956911in}{2.018315in}}%
\pgfpathlineto{\pgfqpoint{0.957775in}{1.976198in}}%
\pgfpathlineto{\pgfqpoint{0.958851in}{1.578001in}}%
\pgfpathlineto{\pgfqpoint{0.957797in}{1.992753in}}%
\pgfpathlineto{\pgfqpoint{0.959000in}{1.713699in}}%
\pgfpathlineto{\pgfqpoint{0.959198in}{1.956176in}}%
\pgfpathlineto{\pgfqpoint{0.959569in}{1.532029in}}%
\pgfpathlineto{\pgfqpoint{0.959643in}{1.532030in}}%
\pgfpathlineto{\pgfqpoint{0.961513in}{1.532108in}}%
\pgfpathlineto{\pgfqpoint{0.961555in}{1.733568in}}%
\pgfpathlineto{\pgfqpoint{0.961734in}{2.012245in}}%
\pgfpathlineto{\pgfqpoint{0.962339in}{1.972491in}}%
\pgfpathlineto{\pgfqpoint{0.962344in}{1.968066in}}%
\pgfpathlineto{\pgfqpoint{0.962390in}{1.996902in}}%
\pgfpathlineto{\pgfqpoint{0.963092in}{1.976126in}}%
\pgfpathlineto{\pgfqpoint{0.963344in}{1.988762in}}%
\pgfpathlineto{\pgfqpoint{0.963493in}{1.930907in}}%
\pgfpathlineto{\pgfqpoint{0.964136in}{1.960027in}}%
\pgfpathlineto{\pgfqpoint{0.964522in}{1.532030in}}%
\pgfpathlineto{\pgfqpoint{0.966463in}{1.532172in}}%
\pgfpathlineto{\pgfqpoint{0.966677in}{2.008876in}}%
\pgfpathlineto{\pgfqpoint{0.967646in}{1.978105in}}%
\pgfpathlineto{\pgfqpoint{0.968746in}{1.567176in}}%
\pgfpathlineto{\pgfqpoint{0.967667in}{1.989221in}}%
\pgfpathlineto{\pgfqpoint{0.968845in}{1.628707in}}%
\pgfpathlineto{\pgfqpoint{0.969093in}{1.974441in}}%
\pgfpathlineto{\pgfqpoint{0.969471in}{1.532029in}}%
\pgfpathlineto{\pgfqpoint{0.969511in}{1.532033in}}%
\pgfpathlineto{\pgfqpoint{0.971403in}{1.532105in}}%
\pgfpathlineto{\pgfqpoint{0.971445in}{1.718209in}}%
\pgfpathlineto{\pgfqpoint{0.971621in}{2.010782in}}%
\pgfpathlineto{\pgfqpoint{0.972228in}{1.968565in}}%
\pgfpathlineto{\pgfqpoint{0.972240in}{1.994067in}}%
\pgfpathlineto{\pgfqpoint{0.972989in}{1.986991in}}%
\pgfpathlineto{\pgfqpoint{0.973582in}{1.623503in}}%
\pgfpathlineto{\pgfqpoint{0.973681in}{1.572818in}}%
\pgfpathlineto{\pgfqpoint{0.973928in}{1.846952in}}%
\pgfpathlineto{\pgfqpoint{0.974027in}{1.961093in}}%
\pgfpathlineto{\pgfqpoint{0.974417in}{1.532030in}}%
\pgfpathlineto{\pgfqpoint{0.974454in}{1.532032in}}%
\pgfpathlineto{\pgfqpoint{0.976352in}{1.532132in}}%
\pgfpathlineto{\pgfqpoint{0.976467in}{2.005152in}}%
\pgfpathlineto{\pgfqpoint{0.977675in}{1.981849in}}%
\pgfpathlineto{\pgfqpoint{0.978653in}{1.569700in}}%
\pgfpathlineto{\pgfqpoint{0.977700in}{1.982091in}}%
\pgfpathlineto{\pgfqpoint{0.978801in}{1.711753in}}%
\pgfpathlineto{\pgfqpoint{0.978999in}{1.970369in}}%
\pgfpathlineto{\pgfqpoint{0.979342in}{1.532029in}}%
\pgfpathlineto{\pgfqpoint{0.979453in}{1.532034in}}%
\pgfpathlineto{\pgfqpoint{0.981294in}{1.532101in}}%
\pgfpathlineto{\pgfqpoint{0.981302in}{1.532462in}}%
\pgfpathlineto{\pgfqpoint{0.981421in}{2.004523in}}%
\pgfpathlineto{\pgfqpoint{0.982359in}{1.978777in}}%
\pgfpathlineto{\pgfqpoint{0.982370in}{1.971928in}}%
\pgfpathlineto{\pgfqpoint{0.982388in}{1.985217in}}%
\pgfpathlineto{\pgfqpoint{0.983125in}{1.975132in}}%
\pgfpathlineto{\pgfqpoint{0.983174in}{1.981960in}}%
\pgfpathlineto{\pgfqpoint{0.983273in}{1.951285in}}%
\pgfpathlineto{\pgfqpoint{0.984319in}{1.532029in}}%
\pgfpathlineto{\pgfqpoint{0.983965in}{1.959594in}}%
\pgfpathlineto{\pgfqpoint{0.984361in}{1.532032in}}%
\pgfpathlineto{\pgfqpoint{0.986241in}{1.532116in}}%
\pgfpathlineto{\pgfqpoint{0.986283in}{1.749065in}}%
\pgfpathlineto{\pgfqpoint{0.986464in}{2.009202in}}%
\pgfpathlineto{\pgfqpoint{0.987065in}{1.967431in}}%
\pgfpathlineto{\pgfqpoint{0.987077in}{1.980818in}}%
\pgfpathlineto{\pgfqpoint{0.987114in}{1.989028in}}%
\pgfpathlineto{\pgfqpoint{0.987097in}{1.966435in}}%
\pgfpathlineto{\pgfqpoint{0.987832in}{1.978476in}}%
\pgfpathlineto{\pgfqpoint{0.988527in}{1.583198in}}%
\pgfpathlineto{\pgfqpoint{0.987885in}{1.982661in}}%
\pgfpathlineto{\pgfqpoint{0.988823in}{1.907220in}}%
\pgfpathlineto{\pgfqpoint{0.988873in}{1.944329in}}%
\pgfpathlineto{\pgfqpoint{0.989153in}{1.622849in}}%
\pgfpathlineto{\pgfqpoint{0.989236in}{1.532026in}}%
\pgfpathlineto{\pgfqpoint{0.990043in}{1.532051in}}%
\pgfpathlineto{\pgfqpoint{0.991187in}{1.532116in}}%
\pgfpathlineto{\pgfqpoint{0.991228in}{1.749522in}}%
\pgfpathlineto{\pgfqpoint{0.991509in}{2.009860in}}%
\pgfpathlineto{\pgfqpoint{0.992009in}{1.959019in}}%
\pgfpathlineto{\pgfqpoint{0.992023in}{1.994517in}}%
\pgfpathlineto{\pgfqpoint{0.992783in}{1.977032in}}%
\pgfpathlineto{\pgfqpoint{0.992794in}{1.982509in}}%
\pgfpathlineto{\pgfqpoint{0.993235in}{1.863141in}}%
\pgfpathlineto{\pgfqpoint{0.994167in}{1.532026in}}%
\pgfpathlineto{\pgfqpoint{0.993829in}{1.955763in}}%
\pgfpathlineto{\pgfqpoint{0.994178in}{1.532036in}}%
\pgfpathlineto{\pgfqpoint{0.995969in}{1.532065in}}%
\pgfpathlineto{\pgfqpoint{0.996134in}{1.532136in}}%
\pgfpathlineto{\pgfqpoint{0.996474in}{2.006494in}}%
\pgfpathlineto{\pgfqpoint{0.997369in}{1.983496in}}%
\pgfpathlineto{\pgfqpoint{0.998436in}{1.571170in}}%
\pgfpathlineto{\pgfqpoint{0.998535in}{1.618117in}}%
\pgfpathlineto{\pgfqpoint{0.998782in}{1.965371in}}%
\pgfpathlineto{\pgfqpoint{0.999135in}{1.532027in}}%
\pgfpathlineto{\pgfqpoint{0.999212in}{1.532028in}}%
\pgfpathlineto{\pgfqpoint{1.001078in}{1.532116in}}%
\pgfpathlineto{\pgfqpoint{1.001133in}{1.969813in}}%
\pgfpathlineto{\pgfqpoint{1.001396in}{2.005577in}}%
\pgfpathlineto{\pgfqpoint{1.001346in}{1.943424in}}%
\pgfpathlineto{\pgfqpoint{1.001902in}{1.967964in}}%
\pgfpathlineto{\pgfqpoint{1.001907in}{1.960583in}}%
\pgfpathlineto{\pgfqpoint{1.001922in}{1.988815in}}%
\pgfpathlineto{\pgfqpoint{1.002659in}{1.976657in}}%
\pgfpathlineto{\pgfqpoint{1.002667in}{1.979511in}}%
\pgfpathlineto{\pgfqpoint{1.003107in}{1.914343in}}%
\pgfpathlineto{\pgfqpoint{1.003750in}{1.968294in}}%
\pgfpathlineto{\pgfqpoint{1.004087in}{1.532026in}}%
\pgfpathlineto{\pgfqpoint{1.006023in}{1.532115in}}%
\pgfpathlineto{\pgfqpoint{1.006064in}{1.737392in}}%
\pgfpathlineto{\pgfqpoint{1.006345in}{2.004763in}}%
\pgfpathlineto{\pgfqpoint{1.006849in}{1.965395in}}%
\pgfpathlineto{\pgfqpoint{1.006854in}{1.957090in}}%
\pgfpathlineto{\pgfqpoint{1.006870in}{1.990324in}}%
\pgfpathlineto{\pgfqpoint{1.007613in}{1.968808in}}%
\pgfpathlineto{\pgfqpoint{1.007785in}{1.979201in}}%
\pgfpathlineto{\pgfqpoint{1.008122in}{1.798294in}}%
\pgfpathlineto{\pgfqpoint{1.008993in}{1.532020in}}%
\pgfpathlineto{\pgfqpoint{1.008715in}{1.937924in}}%
\pgfpathlineto{\pgfqpoint{1.009031in}{1.532038in}}%
\pgfpathlineto{\pgfqpoint{1.010957in}{1.532078in}}%
\pgfpathlineto{\pgfqpoint{1.010967in}{1.532106in}}%
\pgfpathlineto{\pgfqpoint{1.011008in}{1.709526in}}%
\pgfpathlineto{\pgfqpoint{1.011180in}{1.999670in}}%
\pgfpathlineto{\pgfqpoint{1.011791in}{1.964335in}}%
\pgfpathlineto{\pgfqpoint{1.011805in}{1.984934in}}%
\pgfpathlineto{\pgfqpoint{1.011883in}{1.961747in}}%
\pgfpathlineto{\pgfqpoint{1.012526in}{1.966228in}}%
\pgfpathlineto{\pgfqpoint{1.012575in}{1.978135in}}%
\pgfpathlineto{\pgfqpoint{1.013020in}{1.900257in}}%
\pgfpathlineto{\pgfqpoint{1.013958in}{1.532023in}}%
\pgfpathlineto{\pgfqpoint{1.013656in}{1.963097in}}%
\pgfpathlineto{\pgfqpoint{1.013977in}{1.532037in}}%
\pgfpathlineto{\pgfqpoint{1.015904in}{1.532079in}}%
\pgfpathlineto{\pgfqpoint{1.015915in}{1.532125in}}%
\pgfpathlineto{\pgfqpoint{1.016251in}{2.000677in}}%
\pgfpathlineto{\pgfqpoint{1.017222in}{1.975190in}}%
\pgfpathlineto{\pgfqpoint{1.018232in}{1.573093in}}%
\pgfpathlineto{\pgfqpoint{1.017251in}{1.977299in}}%
\pgfpathlineto{\pgfqpoint{1.018331in}{1.620169in}}%
\pgfpathlineto{\pgfqpoint{1.018578in}{1.962190in}}%
\pgfpathlineto{\pgfqpoint{1.018901in}{1.532025in}}%
\pgfpathlineto{\pgfqpoint{1.018996in}{1.532036in}}%
\pgfpathlineto{\pgfqpoint{1.020862in}{1.532139in}}%
\pgfpathlineto{\pgfqpoint{1.021201in}{2.006038in}}%
\pgfpathlineto{\pgfqpoint{1.022062in}{1.962659in}}%
\pgfpathlineto{\pgfqpoint{1.023169in}{1.576307in}}%
\pgfpathlineto{\pgfqpoint{1.022110in}{1.980765in}}%
\pgfpathlineto{\pgfqpoint{1.023367in}{1.729020in}}%
\pgfpathlineto{\pgfqpoint{1.023564in}{1.956521in}}%
\pgfpathlineto{\pgfqpoint{1.023828in}{1.531999in}}%
\pgfpathlineto{\pgfqpoint{1.024016in}{1.532040in}}%
\pgfpathlineto{\pgfqpoint{1.025802in}{1.532099in}}%
\pgfpathlineto{\pgfqpoint{1.025820in}{1.545227in}}%
\pgfpathlineto{\pgfqpoint{1.025906in}{1.994483in}}%
\pgfpathlineto{\pgfqpoint{1.026736in}{1.977835in}}%
\pgfpathlineto{\pgfqpoint{1.026790in}{1.962942in}}%
\pgfpathlineto{\pgfqpoint{1.027475in}{1.970081in}}%
\pgfpathlineto{\pgfqpoint{1.027673in}{1.974537in}}%
\pgfpathlineto{\pgfqpoint{1.027821in}{1.948411in}}%
\pgfpathlineto{\pgfqpoint{1.028773in}{1.532000in}}%
\pgfpathlineto{\pgfqpoint{1.028513in}{1.955339in}}%
\pgfpathlineto{\pgfqpoint{1.028908in}{1.532024in}}%
\pgfpathlineto{\pgfqpoint{1.030757in}{1.532706in}}%
\pgfpathlineto{\pgfqpoint{1.030874in}{1.994668in}}%
\pgfpathlineto{\pgfqpoint{1.031798in}{1.975413in}}%
\pgfpathlineto{\pgfqpoint{1.031816in}{1.963406in}}%
\pgfpathlineto{\pgfqpoint{1.031907in}{1.975452in}}%
\pgfpathlineto{\pgfqpoint{1.032567in}{1.973432in}}%
\pgfpathlineto{\pgfqpoint{1.032666in}{1.974029in}}%
\pgfpathlineto{\pgfqpoint{1.032864in}{1.811200in}}%
\pgfpathlineto{\pgfqpoint{1.033719in}{1.532020in}}%
\pgfpathlineto{\pgfqpoint{1.033457in}{1.938433in}}%
\pgfpathlineto{\pgfqpoint{1.033784in}{1.532035in}}%
\pgfpathlineto{\pgfqpoint{1.035696in}{1.532117in}}%
\pgfpathlineto{\pgfqpoint{1.035736in}{1.731619in}}%
\pgfpathlineto{\pgfqpoint{1.035919in}{1.999557in}}%
\pgfpathlineto{\pgfqpoint{1.036520in}{1.979351in}}%
\pgfpathlineto{\pgfqpoint{1.036563in}{1.956337in}}%
\pgfpathlineto{\pgfqpoint{1.036547in}{1.980485in}}%
\pgfpathlineto{\pgfqpoint{1.037288in}{1.963324in}}%
\pgfpathlineto{\pgfqpoint{1.037392in}{1.973283in}}%
\pgfpathlineto{\pgfqpoint{1.037783in}{1.878632in}}%
\pgfpathlineto{\pgfqpoint{1.038665in}{1.531999in}}%
\pgfpathlineto{\pgfqpoint{1.038426in}{1.940305in}}%
\pgfpathlineto{\pgfqpoint{1.038747in}{1.532038in}}%
\pgfpathlineto{\pgfqpoint{1.040639in}{1.532101in}}%
\pgfpathlineto{\pgfqpoint{1.040668in}{1.578696in}}%
\pgfpathlineto{\pgfqpoint{1.040745in}{1.991457in}}%
\pgfpathlineto{\pgfqpoint{1.041473in}{1.959102in}}%
\pgfpathlineto{\pgfqpoint{1.041487in}{1.977477in}}%
\pgfpathlineto{\pgfqpoint{1.041501in}{1.957979in}}%
\pgfpathlineto{\pgfqpoint{1.042204in}{1.969983in}}%
\pgfpathlineto{\pgfqpoint{1.042452in}{1.973357in}}%
\pgfpathlineto{\pgfqpoint{1.042699in}{1.899685in}}%
\pgfpathlineto{\pgfqpoint{1.043610in}{1.532022in}}%
\pgfpathlineto{\pgfqpoint{1.043342in}{1.929481in}}%
\pgfpathlineto{\pgfqpoint{1.043724in}{1.532035in}}%
\pgfpathlineto{\pgfqpoint{1.045589in}{1.532140in}}%
\pgfpathlineto{\pgfqpoint{1.045929in}{2.002926in}}%
\pgfpathlineto{\pgfqpoint{1.046803in}{1.977882in}}%
\pgfpathlineto{\pgfqpoint{1.046819in}{1.955779in}}%
\pgfpathlineto{\pgfqpoint{1.047556in}{1.961745in}}%
\pgfpathlineto{\pgfqpoint{1.047592in}{1.967264in}}%
\pgfpathlineto{\pgfqpoint{1.047642in}{1.930111in}}%
\pgfpathlineto{\pgfqpoint{1.048557in}{1.532003in}}%
\pgfpathlineto{\pgfqpoint{1.048285in}{1.948993in}}%
\pgfpathlineto{\pgfqpoint{1.048676in}{1.532024in}}%
\pgfpathlineto{\pgfqpoint{1.050531in}{1.532104in}}%
\pgfpathlineto{\pgfqpoint{1.050572in}{1.720000in}}%
\pgfpathlineto{\pgfqpoint{1.050845in}{1.993661in}}%
\pgfpathlineto{\pgfqpoint{1.051355in}{1.972059in}}%
\pgfpathlineto{\pgfqpoint{1.051461in}{1.954409in}}%
\pgfpathlineto{\pgfqpoint{1.051383in}{1.978846in}}%
\pgfpathlineto{\pgfqpoint{1.052125in}{1.959583in}}%
\pgfpathlineto{\pgfqpoint{1.052174in}{1.972276in}}%
\pgfpathlineto{\pgfqpoint{1.052570in}{1.927136in}}%
\pgfpathlineto{\pgfqpoint{1.053502in}{1.532015in}}%
\pgfpathlineto{\pgfqpoint{1.053262in}{1.931674in}}%
\pgfpathlineto{\pgfqpoint{1.053684in}{1.532028in}}%
\pgfpathlineto{\pgfqpoint{1.055478in}{1.532113in}}%
\pgfpathlineto{\pgfqpoint{1.055518in}{1.728280in}}%
\pgfpathlineto{\pgfqpoint{1.055803in}{1.990750in}}%
\pgfpathlineto{\pgfqpoint{1.056306in}{1.968030in}}%
\pgfpathlineto{\pgfqpoint{1.056344in}{1.978233in}}%
\pgfpathlineto{\pgfqpoint{1.056327in}{1.951995in}}%
\pgfpathlineto{\pgfqpoint{1.057058in}{1.958651in}}%
\pgfpathlineto{\pgfqpoint{1.057090in}{1.970220in}}%
\pgfpathlineto{\pgfqpoint{1.057589in}{1.875830in}}%
\pgfpathlineto{\pgfqpoint{1.058489in}{1.532029in}}%
\pgfpathlineto{\pgfqpoint{1.058183in}{1.937541in}}%
\pgfpathlineto{\pgfqpoint{1.058571in}{1.532030in}}%
\pgfpathlineto{\pgfqpoint{1.060425in}{1.532141in}}%
\pgfpathlineto{\pgfqpoint{1.060627in}{1.988399in}}%
\pgfpathlineto{\pgfqpoint{1.061678in}{1.960652in}}%
\pgfpathlineto{\pgfqpoint{1.061719in}{1.967957in}}%
\pgfpathlineto{\pgfqpoint{1.061705in}{1.960255in}}%
\pgfpathlineto{\pgfqpoint{1.062444in}{1.961199in}}%
\pgfpathlineto{\pgfqpoint{1.062619in}{1.715739in}}%
\pgfpathlineto{\pgfqpoint{1.063393in}{1.532007in}}%
\pgfpathlineto{\pgfqpoint{1.063163in}{1.953425in}}%
\pgfpathlineto{\pgfqpoint{1.063514in}{1.532024in}}%
\pgfpathlineto{\pgfqpoint{1.065368in}{1.532110in}}%
\pgfpathlineto{\pgfqpoint{1.065410in}{1.755014in}}%
\pgfpathlineto{\pgfqpoint{1.065582in}{1.991244in}}%
\pgfpathlineto{\pgfqpoint{1.066195in}{1.967890in}}%
\pgfpathlineto{\pgfqpoint{1.066199in}{1.974902in}}%
\pgfpathlineto{\pgfqpoint{1.066216in}{1.951673in}}%
\pgfpathlineto{\pgfqpoint{1.066957in}{1.968477in}}%
\pgfpathlineto{\pgfqpoint{1.067722in}{1.580012in}}%
\pgfpathlineto{\pgfqpoint{1.067969in}{1.816507in}}%
\pgfpathlineto{\pgfqpoint{1.068118in}{1.951096in}}%
\pgfpathlineto{\pgfqpoint{1.068339in}{1.532010in}}%
\pgfpathlineto{\pgfqpoint{1.068554in}{1.532037in}}%
\pgfpathlineto{\pgfqpoint{1.070316in}{1.532134in}}%
\pgfpathlineto{\pgfqpoint{1.070656in}{1.997840in}}%
\pgfpathlineto{\pgfqpoint{1.071536in}{1.972787in}}%
\pgfpathlineto{\pgfqpoint{1.071551in}{1.952599in}}%
\pgfpathlineto{\pgfqpoint{1.072307in}{1.964930in}}%
\pgfpathlineto{\pgfqpoint{1.073334in}{1.532029in}}%
\pgfpathlineto{\pgfqpoint{1.073968in}{1.532044in}}%
\pgfpathlineto{\pgfqpoint{1.075262in}{1.532148in}}%
\pgfpathlineto{\pgfqpoint{1.075435in}{1.986069in}}%
\pgfpathlineto{\pgfqpoint{1.076499in}{1.966064in}}%
\pgfpathlineto{\pgfqpoint{1.076532in}{1.966087in}}%
\pgfpathlineto{\pgfqpoint{1.077398in}{1.859469in}}%
\pgfpathlineto{\pgfqpoint{1.077991in}{1.954109in}}%
\pgfpathlineto{\pgfqpoint{1.078320in}{1.532029in}}%
\pgfpathlineto{\pgfqpoint{1.080205in}{1.532116in}}%
\pgfpathlineto{\pgfqpoint{1.080418in}{1.988176in}}%
\pgfpathlineto{\pgfqpoint{1.081582in}{1.957164in}}%
\pgfpathlineto{\pgfqpoint{1.081597in}{1.964552in}}%
\pgfpathlineto{\pgfqpoint{1.082258in}{1.950874in}}%
\pgfpathlineto{\pgfqpoint{1.083229in}{1.532029in}}%
\pgfpathlineto{\pgfqpoint{1.083568in}{1.532037in}}%
\pgfpathlineto{\pgfqpoint{1.085151in}{1.532120in}}%
\pgfpathlineto{\pgfqpoint{1.085468in}{1.988956in}}%
\pgfpathlineto{\pgfqpoint{1.086649in}{1.965503in}}%
\pgfpathlineto{\pgfqpoint{1.087563in}{1.578358in}}%
\pgfpathlineto{\pgfqpoint{1.087712in}{1.710608in}}%
\pgfpathlineto{\pgfqpoint{1.087910in}{1.955154in}}%
\pgfpathlineto{\pgfqpoint{1.088156in}{1.532028in}}%
\pgfpathlineto{\pgfqpoint{1.088378in}{1.532035in}}%
\pgfpathlineto{\pgfqpoint{1.090093in}{1.532095in}}%
\pgfpathlineto{\pgfqpoint{1.090101in}{1.532284in}}%
\pgfpathlineto{\pgfqpoint{1.090212in}{1.984038in}}%
\pgfpathlineto{\pgfqpoint{1.091184in}{1.965730in}}%
\pgfpathlineto{\pgfqpoint{1.091242in}{1.952744in}}%
\pgfpathlineto{\pgfqpoint{1.091221in}{1.965992in}}%
\pgfpathlineto{\pgfqpoint{1.091971in}{1.956167in}}%
\pgfpathlineto{\pgfqpoint{1.092020in}{1.959968in}}%
\pgfpathlineto{\pgfqpoint{1.092169in}{1.940885in}}%
\pgfpathlineto{\pgfqpoint{1.093117in}{1.532025in}}%
\pgfpathlineto{\pgfqpoint{1.092861in}{1.942914in}}%
\pgfpathlineto{\pgfqpoint{1.093353in}{1.532039in}}%
\pgfpathlineto{\pgfqpoint{1.095043in}{1.532137in}}%
\pgfpathlineto{\pgfqpoint{1.095497in}{1.993395in}}%
\pgfpathlineto{\pgfqpoint{1.096251in}{1.962994in}}%
\pgfpathlineto{\pgfqpoint{1.096261in}{1.947325in}}%
\pgfpathlineto{\pgfqpoint{1.096277in}{1.969451in}}%
\pgfpathlineto{\pgfqpoint{1.097008in}{1.954216in}}%
\pgfpathlineto{\pgfqpoint{1.097031in}{1.959310in}}%
\pgfpathlineto{\pgfqpoint{1.097176in}{1.893902in}}%
\pgfpathlineto{\pgfqpoint{1.098047in}{1.532027in}}%
\pgfpathlineto{\pgfqpoint{1.097819in}{1.949481in}}%
\pgfpathlineto{\pgfqpoint{1.098121in}{1.532028in}}%
\pgfpathlineto{\pgfqpoint{1.099986in}{1.532107in}}%
\pgfpathlineto{\pgfqpoint{1.100028in}{1.740482in}}%
\pgfpathlineto{\pgfqpoint{1.100307in}{1.983223in}}%
\pgfpathlineto{\pgfqpoint{1.100814in}{1.969824in}}%
\pgfpathlineto{\pgfqpoint{1.100963in}{1.947588in}}%
\pgfpathlineto{\pgfqpoint{1.100849in}{1.969886in}}%
\pgfpathlineto{\pgfqpoint{1.101559in}{1.952909in}}%
\pgfpathlineto{\pgfqpoint{1.101848in}{1.963566in}}%
\pgfpathlineto{\pgfqpoint{1.102046in}{1.929126in}}%
\pgfpathlineto{\pgfqpoint{1.102961in}{1.532011in}}%
\pgfpathlineto{\pgfqpoint{1.103806in}{1.532040in}}%
\pgfpathlineto{\pgfqpoint{1.104931in}{1.532105in}}%
\pgfpathlineto{\pgfqpoint{1.104972in}{1.712608in}}%
\pgfpathlineto{\pgfqpoint{1.105145in}{1.985077in}}%
\pgfpathlineto{\pgfqpoint{1.105756in}{1.965648in}}%
\pgfpathlineto{\pgfqpoint{1.105770in}{1.945403in}}%
\pgfpathlineto{\pgfqpoint{1.105816in}{1.968921in}}%
\pgfpathlineto{\pgfqpoint{1.106531in}{1.951332in}}%
\pgfpathlineto{\pgfqpoint{1.106580in}{1.963553in}}%
\pgfpathlineto{\pgfqpoint{1.107075in}{1.915096in}}%
\pgfpathlineto{\pgfqpoint{1.107907in}{1.532001in}}%
\pgfpathlineto{\pgfqpoint{1.107718in}{1.954648in}}%
\pgfpathlineto{\pgfqpoint{1.108094in}{1.532040in}}%
\pgfpathlineto{\pgfqpoint{1.109875in}{1.532096in}}%
\pgfpathlineto{\pgfqpoint{1.109892in}{1.544674in}}%
\pgfpathlineto{\pgfqpoint{1.109988in}{1.980917in}}%
\pgfpathlineto{\pgfqpoint{1.110809in}{1.953831in}}%
\pgfpathlineto{\pgfqpoint{1.110817in}{1.948485in}}%
\pgfpathlineto{\pgfqpoint{1.110830in}{1.963900in}}%
\pgfpathlineto{\pgfqpoint{1.111580in}{1.951357in}}%
\pgfpathlineto{\pgfqpoint{1.111630in}{1.960579in}}%
\pgfpathlineto{\pgfqpoint{1.111976in}{1.937575in}}%
\pgfpathlineto{\pgfqpoint{1.112852in}{1.532013in}}%
\pgfpathlineto{\pgfqpoint{1.112668in}{1.940823in}}%
\pgfpathlineto{\pgfqpoint{1.113103in}{1.532027in}}%
\pgfpathlineto{\pgfqpoint{1.114823in}{1.532110in}}%
\pgfpathlineto{\pgfqpoint{1.114865in}{1.750003in}}%
\pgfpathlineto{\pgfqpoint{1.115036in}{1.982733in}}%
\pgfpathlineto{\pgfqpoint{1.115645in}{1.943975in}}%
\pgfpathlineto{\pgfqpoint{1.115662in}{1.968037in}}%
\pgfpathlineto{\pgfqpoint{1.115678in}{1.943688in}}%
\pgfpathlineto{\pgfqpoint{1.116432in}{1.958800in}}%
\pgfpathlineto{\pgfqpoint{1.116519in}{1.960747in}}%
\pgfpathlineto{\pgfqpoint{1.117014in}{1.837361in}}%
\pgfpathlineto{\pgfqpoint{1.117797in}{1.532002in}}%
\pgfpathlineto{\pgfqpoint{1.117607in}{1.926816in}}%
\pgfpathlineto{\pgfqpoint{1.117967in}{1.532039in}}%
\pgfpathlineto{\pgfqpoint{1.119767in}{1.532105in}}%
\pgfpathlineto{\pgfqpoint{1.119809in}{1.726108in}}%
\pgfpathlineto{\pgfqpoint{1.119980in}{1.981922in}}%
\pgfpathlineto{\pgfqpoint{1.120591in}{1.941517in}}%
\pgfpathlineto{\pgfqpoint{1.120740in}{1.966735in}}%
\pgfpathlineto{\pgfqpoint{1.121344in}{1.959799in}}%
\pgfpathlineto{\pgfqpoint{1.122227in}{1.580435in}}%
\pgfpathlineto{\pgfqpoint{1.121435in}{1.962160in}}%
\pgfpathlineto{\pgfqpoint{1.122425in}{1.784938in}}%
\pgfpathlineto{\pgfqpoint{1.122573in}{1.951480in}}%
\pgfpathlineto{\pgfqpoint{1.122768in}{1.532029in}}%
\pgfpathlineto{\pgfqpoint{1.123024in}{1.532034in}}%
\pgfpathlineto{\pgfqpoint{1.124719in}{1.532224in}}%
\pgfpathlineto{\pgfqpoint{1.125034in}{1.982100in}}%
\pgfpathlineto{\pgfqpoint{1.125823in}{1.960321in}}%
\pgfpathlineto{\pgfqpoint{1.125834in}{1.944754in}}%
\pgfpathlineto{\pgfqpoint{1.125850in}{1.963822in}}%
\pgfpathlineto{\pgfqpoint{1.126584in}{1.952421in}}%
\pgfpathlineto{\pgfqpoint{1.126717in}{1.957930in}}%
\pgfpathlineto{\pgfqpoint{1.126866in}{1.898069in}}%
\pgfpathlineto{\pgfqpoint{1.127727in}{1.532026in}}%
\pgfpathlineto{\pgfqpoint{1.127509in}{1.932769in}}%
\pgfpathlineto{\pgfqpoint{1.127905in}{1.532029in}}%
\pgfpathlineto{\pgfqpoint{1.129660in}{1.532113in}}%
\pgfpathlineto{\pgfqpoint{1.129729in}{1.960395in}}%
\pgfpathlineto{\pgfqpoint{1.129972in}{1.981814in}}%
\pgfpathlineto{\pgfqpoint{1.129923in}{1.923412in}}%
\pgfpathlineto{\pgfqpoint{1.130482in}{1.954740in}}%
\pgfpathlineto{\pgfqpoint{1.130492in}{1.942043in}}%
\pgfpathlineto{\pgfqpoint{1.130507in}{1.964968in}}%
\pgfpathlineto{\pgfqpoint{1.131244in}{1.957271in}}%
\pgfpathlineto{\pgfqpoint{1.131426in}{1.958092in}}%
\pgfpathlineto{\pgfqpoint{1.131871in}{1.846329in}}%
\pgfpathlineto{\pgfqpoint{1.132690in}{1.532026in}}%
\pgfpathlineto{\pgfqpoint{1.132464in}{1.944490in}}%
\pgfpathlineto{\pgfqpoint{1.132782in}{1.532028in}}%
\pgfpathlineto{\pgfqpoint{1.134607in}{1.532138in}}%
\pgfpathlineto{\pgfqpoint{1.134942in}{1.983229in}}%
\pgfpathlineto{\pgfqpoint{1.135804in}{1.961505in}}%
\pgfpathlineto{\pgfqpoint{1.135817in}{1.943514in}}%
\pgfpathlineto{\pgfqpoint{1.135866in}{1.961946in}}%
\pgfpathlineto{\pgfqpoint{1.136577in}{1.957205in}}%
\pgfpathlineto{\pgfqpoint{1.136874in}{1.743838in}}%
\pgfpathlineto{\pgfqpoint{1.137616in}{1.532026in}}%
\pgfpathlineto{\pgfqpoint{1.137418in}{1.941947in}}%
\pgfpathlineto{\pgfqpoint{1.137741in}{1.532034in}}%
\pgfpathlineto{\pgfqpoint{1.139550in}{1.532113in}}%
\pgfpathlineto{\pgfqpoint{1.139864in}{1.979498in}}%
\pgfpathlineto{\pgfqpoint{1.141128in}{1.950721in}}%
\pgfpathlineto{\pgfqpoint{1.141161in}{1.957140in}}%
\pgfpathlineto{\pgfqpoint{1.141762in}{1.852000in}}%
\pgfpathlineto{\pgfqpoint{1.142528in}{1.532020in}}%
\pgfpathlineto{\pgfqpoint{1.142355in}{1.935173in}}%
\pgfpathlineto{\pgfqpoint{1.142688in}{1.532025in}}%
\pgfpathlineto{\pgfqpoint{1.144497in}{1.532120in}}%
\pgfpathlineto{\pgfqpoint{1.144822in}{1.978692in}}%
\pgfpathlineto{\pgfqpoint{1.145823in}{1.956925in}}%
\pgfpathlineto{\pgfqpoint{1.145839in}{1.946753in}}%
\pgfpathlineto{\pgfqpoint{1.146604in}{1.948368in}}%
\pgfpathlineto{\pgfqpoint{1.147504in}{1.532029in}}%
\pgfpathlineto{\pgfqpoint{1.148564in}{1.532051in}}%
\pgfpathlineto{\pgfqpoint{1.149441in}{1.532111in}}%
\pgfpathlineto{\pgfqpoint{1.149497in}{1.945355in}}%
\pgfpathlineto{\pgfqpoint{1.149545in}{1.975089in}}%
\pgfpathlineto{\pgfqpoint{1.149691in}{1.927044in}}%
\pgfpathlineto{\pgfqpoint{1.150266in}{1.944378in}}%
\pgfpathlineto{\pgfqpoint{1.150272in}{1.939633in}}%
\pgfpathlineto{\pgfqpoint{1.150288in}{1.962204in}}%
\pgfpathlineto{\pgfqpoint{1.151027in}{1.948436in}}%
\pgfpathlineto{\pgfqpoint{1.151049in}{1.954980in}}%
\pgfpathlineto{\pgfqpoint{1.151649in}{1.889598in}}%
\pgfpathlineto{\pgfqpoint{1.152420in}{1.532021in}}%
\pgfpathlineto{\pgfqpoint{1.152292in}{1.939664in}}%
\pgfpathlineto{\pgfqpoint{1.152618in}{1.532038in}}%
\pgfpathlineto{\pgfqpoint{1.154389in}{1.532137in}}%
\pgfpathlineto{\pgfqpoint{1.154724in}{1.978833in}}%
\pgfpathlineto{\pgfqpoint{1.155582in}{1.957958in}}%
\pgfpathlineto{\pgfqpoint{1.155604in}{1.942515in}}%
\pgfpathlineto{\pgfqpoint{1.155624in}{1.958252in}}%
\pgfpathlineto{\pgfqpoint{1.156357in}{1.953546in}}%
\pgfpathlineto{\pgfqpoint{1.156605in}{1.865993in}}%
\pgfpathlineto{\pgfqpoint{1.157364in}{1.532021in}}%
\pgfpathlineto{\pgfqpoint{1.157248in}{1.913282in}}%
\pgfpathlineto{\pgfqpoint{1.157611in}{1.532035in}}%
\pgfpathlineto{\pgfqpoint{1.159330in}{1.532099in}}%
\pgfpathlineto{\pgfqpoint{1.159339in}{1.532436in}}%
\pgfpathlineto{\pgfqpoint{1.159441in}{1.974971in}}%
\pgfpathlineto{\pgfqpoint{1.160389in}{1.955093in}}%
\pgfpathlineto{\pgfqpoint{1.160434in}{1.942702in}}%
\pgfpathlineto{\pgfqpoint{1.160419in}{1.957190in}}%
\pgfpathlineto{\pgfqpoint{1.161159in}{1.948070in}}%
\pgfpathlineto{\pgfqpoint{1.161357in}{1.953707in}}%
\pgfpathlineto{\pgfqpoint{1.161505in}{1.944262in}}%
\pgfpathlineto{\pgfqpoint{1.162335in}{1.532029in}}%
\pgfpathlineto{\pgfqpoint{1.162692in}{1.532035in}}%
\pgfpathlineto{\pgfqpoint{1.164276in}{1.532102in}}%
\pgfpathlineto{\pgfqpoint{1.164317in}{1.714464in}}%
\pgfpathlineto{\pgfqpoint{1.164589in}{1.973225in}}%
\pgfpathlineto{\pgfqpoint{1.165103in}{1.962626in}}%
\pgfpathlineto{\pgfqpoint{1.165149in}{1.937795in}}%
\pgfpathlineto{\pgfqpoint{1.165877in}{1.942799in}}%
\pgfpathlineto{\pgfqpoint{1.166273in}{1.955908in}}%
\pgfpathlineto{\pgfqpoint{1.166223in}{1.941797in}}%
\pgfpathlineto{\pgfqpoint{1.166421in}{1.946059in}}%
\pgfpathlineto{\pgfqpoint{1.167257in}{1.532020in}}%
\pgfpathlineto{\pgfqpoint{1.167607in}{1.532031in}}%
\pgfpathlineto{\pgfqpoint{1.169225in}{1.532137in}}%
\pgfpathlineto{\pgfqpoint{1.169320in}{1.970690in}}%
\pgfpathlineto{\pgfqpoint{1.170403in}{1.951297in}}%
\pgfpathlineto{\pgfqpoint{1.170412in}{1.940140in}}%
\pgfpathlineto{\pgfqpoint{1.170434in}{1.957443in}}%
\pgfpathlineto{\pgfqpoint{1.171141in}{1.948138in}}%
\pgfpathlineto{\pgfqpoint{1.171289in}{1.953727in}}%
\pgfpathlineto{\pgfqpoint{1.171438in}{1.941106in}}%
\pgfpathlineto{\pgfqpoint{1.172203in}{1.532005in}}%
\pgfpathlineto{\pgfqpoint{1.172124in}{1.942165in}}%
\pgfpathlineto{\pgfqpoint{1.172553in}{1.532028in}}%
\pgfpathlineto{\pgfqpoint{1.174174in}{1.532239in}}%
\pgfpathlineto{\pgfqpoint{1.174290in}{1.973113in}}%
\pgfpathlineto{\pgfqpoint{1.175284in}{1.950997in}}%
\pgfpathlineto{\pgfqpoint{1.175289in}{1.954286in}}%
\pgfpathlineto{\pgfqpoint{1.175304in}{1.942533in}}%
\pgfpathlineto{\pgfqpoint{1.176045in}{1.946040in}}%
\pgfpathlineto{\pgfqpoint{1.176063in}{1.949883in}}%
\pgfpathlineto{\pgfqpoint{1.176446in}{1.871294in}}%
\pgfpathlineto{\pgfqpoint{1.177149in}{1.532000in}}%
\pgfpathlineto{\pgfqpoint{1.177077in}{1.937506in}}%
\pgfpathlineto{\pgfqpoint{1.177396in}{1.532025in}}%
\pgfpathlineto{\pgfqpoint{1.179120in}{1.532319in}}%
\pgfpathlineto{\pgfqpoint{1.179315in}{1.973932in}}%
\pgfpathlineto{\pgfqpoint{1.180178in}{1.956498in}}%
\pgfpathlineto{\pgfqpoint{1.180230in}{1.939573in}}%
\pgfpathlineto{\pgfqpoint{1.180958in}{1.943355in}}%
\pgfpathlineto{\pgfqpoint{1.181255in}{1.952219in}}%
\pgfpathlineto{\pgfqpoint{1.181354in}{1.918307in}}%
\pgfpathlineto{\pgfqpoint{1.182103in}{1.531999in}}%
\pgfpathlineto{\pgfqpoint{1.182028in}{1.931182in}}%
\pgfpathlineto{\pgfqpoint{1.182429in}{1.532054in}}%
\pgfpathlineto{\pgfqpoint{1.183764in}{1.532048in}}%
\pgfpathlineto{\pgfqpoint{1.184067in}{1.532715in}}%
\pgfpathlineto{\pgfqpoint{1.184178in}{1.972548in}}%
\pgfpathlineto{\pgfqpoint{1.185101in}{1.955396in}}%
\pgfpathlineto{\pgfqpoint{1.185120in}{1.939697in}}%
\pgfpathlineto{\pgfqpoint{1.185839in}{1.950170in}}%
\pgfpathlineto{\pgfqpoint{1.185938in}{1.952793in}}%
\pgfpathlineto{\pgfqpoint{1.186432in}{1.760015in}}%
\pgfpathlineto{\pgfqpoint{1.187055in}{1.532024in}}%
\pgfpathlineto{\pgfqpoint{1.187011in}{1.941143in}}%
\pgfpathlineto{\pgfqpoint{1.187336in}{1.532042in}}%
\pgfpathlineto{\pgfqpoint{1.189005in}{1.532115in}}%
\pgfpathlineto{\pgfqpoint{1.189211in}{1.974856in}}%
\pgfpathlineto{\pgfqpoint{1.190423in}{1.944835in}}%
\pgfpathlineto{\pgfqpoint{1.190427in}{1.943341in}}%
\pgfpathlineto{\pgfqpoint{1.190443in}{1.950175in}}%
\pgfpathlineto{\pgfqpoint{1.191185in}{1.945586in}}%
\pgfpathlineto{\pgfqpoint{1.191192in}{1.946223in}}%
\pgfpathlineto{\pgfqpoint{1.191309in}{1.888394in}}%
\pgfpathlineto{\pgfqpoint{1.191987in}{1.532023in}}%
\pgfpathlineto{\pgfqpoint{1.191957in}{1.947993in}}%
\pgfpathlineto{\pgfqpoint{1.192285in}{1.532041in}}%
\pgfpathlineto{\pgfqpoint{1.193942in}{1.532082in}}%
\pgfpathlineto{\pgfqpoint{1.193952in}{1.532141in}}%
\pgfpathlineto{\pgfqpoint{1.194048in}{1.970820in}}%
\pgfpathlineto{\pgfqpoint{1.195125in}{1.941448in}}%
\pgfpathlineto{\pgfqpoint{1.195231in}{1.951825in}}%
\pgfpathlineto{\pgfqpoint{1.195623in}{1.940830in}}%
\pgfpathlineto{\pgfqpoint{1.195870in}{1.945446in}}%
\pgfpathlineto{\pgfqpoint{1.196019in}{1.950691in}}%
\pgfpathlineto{\pgfqpoint{1.196068in}{1.940474in}}%
\pgfpathlineto{\pgfqpoint{1.196167in}{1.941278in}}%
\pgfpathlineto{\pgfqpoint{1.196930in}{1.532016in}}%
\pgfpathlineto{\pgfqpoint{1.197685in}{1.532052in}}%
\pgfpathlineto{\pgfqpoint{1.198896in}{1.532116in}}%
\pgfpathlineto{\pgfqpoint{1.198937in}{1.745827in}}%
\pgfpathlineto{\pgfqpoint{1.199207in}{1.974072in}}%
\pgfpathlineto{\pgfqpoint{1.199719in}{1.935326in}}%
\pgfpathlineto{\pgfqpoint{1.199784in}{1.933689in}}%
\pgfpathlineto{\pgfqpoint{1.199768in}{1.958189in}}%
\pgfpathlineto{\pgfqpoint{1.200436in}{1.942086in}}%
\pgfpathlineto{\pgfqpoint{1.200447in}{1.951294in}}%
\pgfpathlineto{\pgfqpoint{1.201124in}{1.939855in}}%
\pgfpathlineto{\pgfqpoint{1.201174in}{1.948860in}}%
\pgfpathlineto{\pgfqpoint{1.201897in}{1.532027in}}%
\pgfpathlineto{\pgfqpoint{1.201846in}{1.950302in}}%
\pgfpathlineto{\pgfqpoint{1.202490in}{1.532048in}}%
\pgfpathlineto{\pgfqpoint{1.203841in}{1.532117in}}%
\pgfpathlineto{\pgfqpoint{1.203882in}{1.744468in}}%
\pgfpathlineto{\pgfqpoint{1.204155in}{1.974462in}}%
\pgfpathlineto{\pgfqpoint{1.204666in}{1.933289in}}%
\pgfpathlineto{\pgfqpoint{1.204741in}{1.957353in}}%
\pgfpathlineto{\pgfqpoint{1.205432in}{1.946014in}}%
\pgfpathlineto{\pgfqpoint{1.206494in}{1.584468in}}%
\pgfpathlineto{\pgfqpoint{1.205518in}{1.950522in}}%
\pgfpathlineto{\pgfqpoint{1.206642in}{1.698279in}}%
\pgfpathlineto{\pgfqpoint{1.206792in}{1.931743in}}%
\pgfpathlineto{\pgfqpoint{1.206821in}{1.532017in}}%
\pgfpathlineto{\pgfqpoint{1.207293in}{1.532048in}}%
\pgfpathlineto{\pgfqpoint{1.208787in}{1.532122in}}%
\pgfpathlineto{\pgfqpoint{1.208879in}{1.968134in}}%
\pgfpathlineto{\pgfqpoint{1.210138in}{1.941217in}}%
\pgfpathlineto{\pgfqpoint{1.210154in}{1.949173in}}%
\pgfpathlineto{\pgfqpoint{1.210905in}{1.945359in}}%
\pgfpathlineto{\pgfqpoint{1.210948in}{1.945901in}}%
\pgfpathlineto{\pgfqpoint{1.211234in}{1.818794in}}%
\pgfpathlineto{\pgfqpoint{1.211781in}{1.532026in}}%
\pgfpathlineto{\pgfqpoint{1.211737in}{1.896269in}}%
\pgfpathlineto{\pgfqpoint{1.212163in}{1.532046in}}%
\pgfpathlineto{\pgfqpoint{1.213722in}{1.532082in}}%
\pgfpathlineto{\pgfqpoint{1.213733in}{1.532133in}}%
\pgfpathlineto{\pgfqpoint{1.214073in}{1.978643in}}%
\pgfpathlineto{\pgfqpoint{1.214986in}{1.938775in}}%
\pgfpathlineto{\pgfqpoint{1.214992in}{1.937115in}}%
\pgfpathlineto{\pgfqpoint{1.215007in}{1.952441in}}%
\pgfpathlineto{\pgfqpoint{1.215720in}{1.947138in}}%
\pgfpathlineto{\pgfqpoint{1.215962in}{1.947542in}}%
\pgfpathlineto{\pgfqpoint{1.216160in}{1.890193in}}%
\pgfpathlineto{\pgfqpoint{1.216709in}{1.532021in}}%
\pgfpathlineto{\pgfqpoint{1.217127in}{1.532034in}}%
\pgfpathlineto{\pgfqpoint{1.218678in}{1.532129in}}%
\pgfpathlineto{\pgfqpoint{1.218862in}{1.969891in}}%
\pgfpathlineto{\pgfqpoint{1.219981in}{1.947067in}}%
\pgfpathlineto{\pgfqpoint{1.219992in}{1.939643in}}%
\pgfpathlineto{\pgfqpoint{1.220009in}{1.949049in}}%
\pgfpathlineto{\pgfqpoint{1.220736in}{1.946528in}}%
\pgfpathlineto{\pgfqpoint{1.220835in}{1.946710in}}%
\pgfpathlineto{\pgfqpoint{1.221131in}{1.885679in}}%
\pgfpathlineto{\pgfqpoint{1.221653in}{1.532025in}}%
\pgfpathlineto{\pgfqpoint{1.222109in}{1.532037in}}%
\pgfpathlineto{\pgfqpoint{1.223626in}{1.532158in}}%
\pgfpathlineto{\pgfqpoint{1.223720in}{1.970557in}}%
\pgfpathlineto{\pgfqpoint{1.224764in}{1.944282in}}%
\pgfpathlineto{\pgfqpoint{1.225011in}{1.950706in}}%
\pgfpathlineto{\pgfqpoint{1.224962in}{1.938091in}}%
\pgfpathlineto{\pgfqpoint{1.225456in}{1.950364in}}%
\pgfpathlineto{\pgfqpoint{1.226620in}{1.532029in}}%
\pgfpathlineto{\pgfqpoint{1.228059in}{1.532060in}}%
\pgfpathlineto{\pgfqpoint{1.228569in}{1.532121in}}%
\pgfpathlineto{\pgfqpoint{1.228623in}{1.938911in}}%
\pgfpathlineto{\pgfqpoint{1.228888in}{1.972893in}}%
\pgfpathlineto{\pgfqpoint{1.228838in}{1.913819in}}%
\pgfpathlineto{\pgfqpoint{1.229394in}{1.949071in}}%
\pgfpathlineto{\pgfqpoint{1.229405in}{1.930218in}}%
\pgfpathlineto{\pgfqpoint{1.229421in}{1.957451in}}%
\pgfpathlineto{\pgfqpoint{1.230166in}{1.941545in}}%
\pgfpathlineto{\pgfqpoint{1.230174in}{1.938504in}}%
\pgfpathlineto{\pgfqpoint{1.230192in}{1.949669in}}%
\pgfpathlineto{\pgfqpoint{1.230902in}{1.940019in}}%
\pgfpathlineto{\pgfqpoint{1.230951in}{1.947105in}}%
\pgfpathlineto{\pgfqpoint{1.231050in}{1.863900in}}%
\pgfpathlineto{\pgfqpoint{1.231569in}{1.532032in}}%
\pgfpathlineto{\pgfqpoint{1.232024in}{1.532041in}}%
\pgfpathlineto{\pgfqpoint{1.233521in}{1.532531in}}%
\pgfpathlineto{\pgfqpoint{1.233622in}{1.970811in}}%
\pgfpathlineto{\pgfqpoint{1.234557in}{1.950813in}}%
\pgfpathlineto{\pgfqpoint{1.234571in}{1.935600in}}%
\pgfpathlineto{\pgfqpoint{1.234584in}{1.951882in}}%
\pgfpathlineto{\pgfqpoint{1.235302in}{1.938829in}}%
\pgfpathlineto{\pgfqpoint{1.235698in}{1.948446in}}%
\pgfpathlineto{\pgfqpoint{1.235945in}{1.918046in}}%
\pgfpathlineto{\pgfqpoint{1.236487in}{1.532013in}}%
\pgfpathlineto{\pgfqpoint{1.237127in}{1.532038in}}%
\pgfpathlineto{\pgfqpoint{1.238466in}{1.532683in}}%
\pgfpathlineto{\pgfqpoint{1.238568in}{1.970190in}}%
\pgfpathlineto{\pgfqpoint{1.239495in}{1.945923in}}%
\pgfpathlineto{\pgfqpoint{1.239503in}{1.951024in}}%
\pgfpathlineto{\pgfqpoint{1.239521in}{1.936608in}}%
\pgfpathlineto{\pgfqpoint{1.240269in}{1.948258in}}%
\pgfpathlineto{\pgfqpoint{1.240368in}{1.949276in}}%
\pgfpathlineto{\pgfqpoint{1.241110in}{1.639320in}}%
\pgfpathlineto{\pgfqpoint{1.241433in}{1.532015in}}%
\pgfpathlineto{\pgfqpoint{1.241410in}{1.785612in}}%
\pgfpathlineto{\pgfqpoint{1.242007in}{1.532052in}}%
\pgfpathlineto{\pgfqpoint{1.243404in}{1.532114in}}%
\pgfpathlineto{\pgfqpoint{1.243446in}{1.729555in}}%
\pgfpathlineto{\pgfqpoint{1.243712in}{1.970599in}}%
\pgfpathlineto{\pgfqpoint{1.244225in}{1.930377in}}%
\pgfpathlineto{\pgfqpoint{1.244277in}{1.956745in}}%
\pgfpathlineto{\pgfqpoint{1.245015in}{1.949719in}}%
\pgfpathlineto{\pgfqpoint{1.246134in}{1.609017in}}%
\pgfpathlineto{\pgfqpoint{1.245097in}{1.950037in}}%
\pgfpathlineto{\pgfqpoint{1.246317in}{1.715303in}}%
\pgfpathlineto{\pgfqpoint{1.246355in}{1.778060in}}%
\pgfpathlineto{\pgfqpoint{1.246378in}{1.532025in}}%
\pgfpathlineto{\pgfqpoint{1.246944in}{1.532050in}}%
\pgfpathlineto{\pgfqpoint{1.248341in}{1.532084in}}%
\pgfpathlineto{\pgfqpoint{1.248353in}{1.532152in}}%
\pgfpathlineto{\pgfqpoint{1.248445in}{1.968452in}}%
\pgfpathlineto{\pgfqpoint{1.249510in}{1.949846in}}%
\pgfpathlineto{\pgfqpoint{1.249554in}{1.934089in}}%
\pgfpathlineto{\pgfqpoint{1.249539in}{1.952961in}}%
\pgfpathlineto{\pgfqpoint{1.250247in}{1.940857in}}%
\pgfpathlineto{\pgfqpoint{1.250544in}{1.949144in}}%
\pgfpathlineto{\pgfqpoint{1.250495in}{1.936645in}}%
\pgfpathlineto{\pgfqpoint{1.250693in}{1.940463in}}%
\pgfpathlineto{\pgfqpoint{1.251324in}{1.532023in}}%
\pgfpathlineto{\pgfqpoint{1.252189in}{1.532057in}}%
\pgfpathlineto{\pgfqpoint{1.253286in}{1.532083in}}%
\pgfpathlineto{\pgfqpoint{1.253298in}{1.532143in}}%
\pgfpathlineto{\pgfqpoint{1.253392in}{1.967904in}}%
\pgfpathlineto{\pgfqpoint{1.254478in}{1.955195in}}%
\pgfpathlineto{\pgfqpoint{1.254527in}{1.928951in}}%
\pgfpathlineto{\pgfqpoint{1.254576in}{1.957812in}}%
\pgfpathlineto{\pgfqpoint{1.255245in}{1.952819in}}%
\pgfpathlineto{\pgfqpoint{1.256269in}{1.532016in}}%
\pgfpathlineto{\pgfqpoint{1.255579in}{1.953601in}}%
\pgfpathlineto{\pgfqpoint{1.257244in}{1.532061in}}%
\pgfpathlineto{\pgfqpoint{1.258240in}{1.532114in}}%
\pgfpathlineto{\pgfqpoint{1.258282in}{1.730509in}}%
\pgfpathlineto{\pgfqpoint{1.258548in}{1.970666in}}%
\pgfpathlineto{\pgfqpoint{1.259065in}{1.951201in}}%
\pgfpathlineto{\pgfqpoint{1.259110in}{1.930535in}}%
\pgfpathlineto{\pgfqpoint{1.259092in}{1.956373in}}%
\pgfpathlineto{\pgfqpoint{1.259814in}{1.949887in}}%
\pgfpathlineto{\pgfqpoint{1.260308in}{1.935517in}}%
\pgfpathlineto{\pgfqpoint{1.260555in}{1.940297in}}%
\pgfpathlineto{\pgfqpoint{1.260605in}{1.945792in}}%
\pgfpathlineto{\pgfqpoint{1.260704in}{1.894170in}}%
\pgfpathlineto{\pgfqpoint{1.261215in}{1.532017in}}%
\pgfpathlineto{\pgfqpoint{1.261811in}{1.532053in}}%
\pgfpathlineto{\pgfqpoint{1.263185in}{1.532111in}}%
\pgfpathlineto{\pgfqpoint{1.263226in}{1.712021in}}%
\pgfpathlineto{\pgfqpoint{1.263493in}{1.970390in}}%
\pgfpathlineto{\pgfqpoint{1.264009in}{1.950767in}}%
\pgfpathlineto{\pgfqpoint{1.264020in}{1.927967in}}%
\pgfpathlineto{\pgfqpoint{1.264036in}{1.956818in}}%
\pgfpathlineto{\pgfqpoint{1.264770in}{1.935500in}}%
\pgfpathlineto{\pgfqpoint{1.265113in}{1.949268in}}%
\pgfpathlineto{\pgfqpoint{1.265508in}{1.935183in}}%
\pgfpathlineto{\pgfqpoint{1.265558in}{1.947646in}}%
\pgfpathlineto{\pgfqpoint{1.265755in}{1.750341in}}%
\pgfpathlineto{\pgfqpoint{1.266160in}{1.532027in}}%
\pgfpathlineto{\pgfqpoint{1.266137in}{1.776801in}}%
\pgfpathlineto{\pgfqpoint{1.266674in}{1.532037in}}%
\pgfpathlineto{\pgfqpoint{1.268132in}{1.532121in}}%
\pgfpathlineto{\pgfqpoint{1.268173in}{1.736361in}}%
\pgfpathlineto{\pgfqpoint{1.268445in}{1.972405in}}%
\pgfpathlineto{\pgfqpoint{1.268957in}{1.929544in}}%
\pgfpathlineto{\pgfqpoint{1.268973in}{1.954926in}}%
\pgfpathlineto{\pgfqpoint{1.269729in}{1.937736in}}%
\pgfpathlineto{\pgfqpoint{1.269757in}{1.948132in}}%
\pgfpathlineto{\pgfqpoint{1.270462in}{1.936745in}}%
\pgfpathlineto{\pgfqpoint{1.270512in}{1.946763in}}%
\pgfpathlineto{\pgfqpoint{1.270710in}{1.750820in}}%
\pgfpathlineto{\pgfqpoint{1.271105in}{1.532020in}}%
\pgfpathlineto{\pgfqpoint{1.271082in}{1.760250in}}%
\pgfpathlineto{\pgfqpoint{1.271603in}{1.532050in}}%
\pgfpathlineto{\pgfqpoint{1.273077in}{1.532117in}}%
\pgfpathlineto{\pgfqpoint{1.273119in}{1.738934in}}%
\pgfpathlineto{\pgfqpoint{1.273394in}{1.968380in}}%
\pgfpathlineto{\pgfqpoint{1.273901in}{1.943847in}}%
\pgfpathlineto{\pgfqpoint{1.273910in}{1.928290in}}%
\pgfpathlineto{\pgfqpoint{1.273926in}{1.955820in}}%
\pgfpathlineto{\pgfqpoint{1.274665in}{1.947739in}}%
\pgfpathlineto{\pgfqpoint{1.276084in}{1.532031in}}%
\pgfpathlineto{\pgfqpoint{1.274944in}{1.948316in}}%
\pgfpathlineto{\pgfqpoint{1.276949in}{1.532049in}}%
\pgfpathlineto{\pgfqpoint{1.278023in}{1.532119in}}%
\pgfpathlineto{\pgfqpoint{1.278124in}{1.968441in}}%
\pgfpathlineto{\pgfqpoint{1.279603in}{1.946589in}}%
\pgfpathlineto{\pgfqpoint{1.280255in}{1.937723in}}%
\pgfpathlineto{\pgfqpoint{1.280354in}{1.939241in}}%
\pgfpathlineto{\pgfqpoint{1.280403in}{1.939299in}}%
\pgfpathlineto{\pgfqpoint{1.280453in}{1.941066in}}%
\pgfpathlineto{\pgfqpoint{1.281011in}{1.532026in}}%
\pgfpathlineto{\pgfqpoint{1.281690in}{1.532054in}}%
\pgfpathlineto{\pgfqpoint{1.282969in}{1.532122in}}%
\pgfpathlineto{\pgfqpoint{1.283273in}{1.970308in}}%
\pgfpathlineto{\pgfqpoint{1.284358in}{1.941317in}}%
\pgfpathlineto{\pgfqpoint{1.284366in}{1.937495in}}%
\pgfpathlineto{\pgfqpoint{1.284379in}{1.946625in}}%
\pgfpathlineto{\pgfqpoint{1.285126in}{1.942509in}}%
\pgfpathlineto{\pgfqpoint{1.285132in}{1.942741in}}%
\pgfpathlineto{\pgfqpoint{1.285419in}{1.930807in}}%
\pgfpathlineto{\pgfqpoint{1.285996in}{1.532027in}}%
\pgfpathlineto{\pgfqpoint{1.286681in}{1.532042in}}%
\pgfpathlineto{\pgfqpoint{1.287916in}{1.532143in}}%
\pgfpathlineto{\pgfqpoint{1.288010in}{1.968069in}}%
\pgfpathlineto{\pgfqpoint{1.289067in}{1.946541in}}%
\pgfpathlineto{\pgfqpoint{1.289413in}{1.948238in}}%
\pgfpathlineto{\pgfqpoint{1.289809in}{1.935697in}}%
\pgfpathlineto{\pgfqpoint{1.289858in}{1.947859in}}%
\pgfpathlineto{\pgfqpoint{1.290402in}{1.895249in}}%
\pgfpathlineto{\pgfqpoint{1.290927in}{1.532030in}}%
\pgfpathlineto{\pgfqpoint{1.291456in}{1.532048in}}%
\pgfpathlineto{\pgfqpoint{1.292861in}{1.532137in}}%
\pgfpathlineto{\pgfqpoint{1.293200in}{1.979591in}}%
\pgfpathlineto{\pgfqpoint{1.294097in}{1.946741in}}%
\pgfpathlineto{\pgfqpoint{1.294107in}{1.930563in}}%
\pgfpathlineto{\pgfqpoint{1.294123in}{1.952588in}}%
\pgfpathlineto{\pgfqpoint{1.294872in}{1.937797in}}%
\pgfpathlineto{\pgfqpoint{1.294923in}{1.945172in}}%
\pgfpathlineto{\pgfqpoint{1.295377in}{1.850732in}}%
\pgfpathlineto{\pgfqpoint{1.295888in}{1.532027in}}%
\pgfpathlineto{\pgfqpoint{1.296325in}{1.532049in}}%
\pgfpathlineto{\pgfqpoint{1.297795in}{1.532083in}}%
\pgfpathlineto{\pgfqpoint{1.297806in}{1.532133in}}%
\pgfpathlineto{\pgfqpoint{1.297896in}{1.966480in}}%
\pgfpathlineto{\pgfqpoint{1.299078in}{1.946086in}}%
\pgfpathlineto{\pgfqpoint{1.299091in}{1.935876in}}%
\pgfpathlineto{\pgfqpoint{1.299174in}{1.946538in}}%
\pgfpathlineto{\pgfqpoint{1.299831in}{1.942617in}}%
\pgfpathlineto{\pgfqpoint{1.299930in}{1.943919in}}%
\pgfpathlineto{\pgfqpoint{1.300325in}{1.852790in}}%
\pgfpathlineto{\pgfqpoint{1.300826in}{1.532028in}}%
\pgfpathlineto{\pgfqpoint{1.301310in}{1.532039in}}%
\pgfpathlineto{\pgfqpoint{1.302752in}{1.532147in}}%
\pgfpathlineto{\pgfqpoint{1.302949in}{1.969929in}}%
\pgfpathlineto{\pgfqpoint{1.303928in}{1.935065in}}%
\pgfpathlineto{\pgfqpoint{1.303946in}{1.947926in}}%
\pgfpathlineto{\pgfqpoint{1.304075in}{1.934823in}}%
\pgfpathlineto{\pgfqpoint{1.304727in}{1.945362in}}%
\pgfpathlineto{\pgfqpoint{1.305321in}{1.787969in}}%
\pgfpathlineto{\pgfqpoint{1.305812in}{1.532033in}}%
\pgfpathlineto{\pgfqpoint{1.306253in}{1.532045in}}%
\pgfpathlineto{\pgfqpoint{1.307697in}{1.532130in}}%
\pgfpathlineto{\pgfqpoint{1.308006in}{1.969026in}}%
\pgfpathlineto{\pgfqpoint{1.309007in}{1.943130in}}%
\pgfpathlineto{\pgfqpoint{1.309015in}{1.934182in}}%
\pgfpathlineto{\pgfqpoint{1.309059in}{1.947632in}}%
\pgfpathlineto{\pgfqpoint{1.309777in}{1.940756in}}%
\pgfpathlineto{\pgfqpoint{1.309789in}{1.943014in}}%
\pgfpathlineto{\pgfqpoint{1.310195in}{1.901719in}}%
\pgfpathlineto{\pgfqpoint{1.310667in}{1.532027in}}%
\pgfpathlineto{\pgfqpoint{1.311226in}{1.532050in}}%
\pgfpathlineto{\pgfqpoint{1.312641in}{1.532121in}}%
\pgfpathlineto{\pgfqpoint{1.312682in}{1.735025in}}%
\pgfpathlineto{\pgfqpoint{1.312953in}{1.970962in}}%
\pgfpathlineto{\pgfqpoint{1.313466in}{1.928002in}}%
\pgfpathlineto{\pgfqpoint{1.313482in}{1.953385in}}%
\pgfpathlineto{\pgfqpoint{1.314238in}{1.939704in}}%
\pgfpathlineto{\pgfqpoint{1.314477in}{1.946605in}}%
\pgfpathlineto{\pgfqpoint{1.314972in}{1.935239in}}%
\pgfpathlineto{\pgfqpoint{1.315021in}{1.943982in}}%
\pgfpathlineto{\pgfqpoint{1.315170in}{1.852594in}}%
\pgfpathlineto{\pgfqpoint{1.315674in}{1.532032in}}%
\pgfpathlineto{\pgfqpoint{1.316162in}{1.532042in}}%
\pgfpathlineto{\pgfqpoint{1.317586in}{1.532117in}}%
\pgfpathlineto{\pgfqpoint{1.317629in}{1.758246in}}%
\pgfpathlineto{\pgfqpoint{1.317677in}{1.966161in}}%
\pgfpathlineto{\pgfqpoint{1.318413in}{1.948184in}}%
\pgfpathlineto{\pgfqpoint{1.318427in}{1.931979in}}%
\pgfpathlineto{\pgfqpoint{1.318509in}{1.948868in}}%
\pgfpathlineto{\pgfqpoint{1.319156in}{1.935557in}}%
\pgfpathlineto{\pgfqpoint{1.319552in}{1.945972in}}%
\pgfpathlineto{\pgfqpoint{1.319601in}{1.935530in}}%
\pgfpathlineto{\pgfqpoint{1.319898in}{1.943872in}}%
\pgfpathlineto{\pgfqpoint{1.319997in}{1.944664in}}%
\pgfpathlineto{\pgfqpoint{1.320145in}{1.806541in}}%
\pgfpathlineto{\pgfqpoint{1.320650in}{1.532033in}}%
\pgfpathlineto{\pgfqpoint{1.321043in}{1.532041in}}%
\pgfpathlineto{\pgfqpoint{1.322532in}{1.532117in}}%
\pgfpathlineto{\pgfqpoint{1.322574in}{1.761686in}}%
\pgfpathlineto{\pgfqpoint{1.322636in}{1.967358in}}%
\pgfpathlineto{\pgfqpoint{1.323358in}{1.950241in}}%
\pgfpathlineto{\pgfqpoint{1.323364in}{1.953533in}}%
\pgfpathlineto{\pgfqpoint{1.323379in}{1.928117in}}%
\pgfpathlineto{\pgfqpoint{1.324106in}{1.946164in}}%
\pgfpathlineto{\pgfqpoint{1.324466in}{1.935442in}}%
\pgfpathlineto{\pgfqpoint{1.324212in}{1.946172in}}%
\pgfpathlineto{\pgfqpoint{1.324862in}{1.945083in}}%
\pgfpathlineto{\pgfqpoint{1.325109in}{1.805526in}}%
\pgfpathlineto{\pgfqpoint{1.325504in}{1.532017in}}%
\pgfpathlineto{\pgfqpoint{1.326005in}{1.532051in}}%
\pgfpathlineto{\pgfqpoint{1.327475in}{1.532106in}}%
\pgfpathlineto{\pgfqpoint{1.327502in}{1.572436in}}%
\pgfpathlineto{\pgfqpoint{1.327587in}{1.967472in}}%
\pgfpathlineto{\pgfqpoint{1.328312in}{1.942365in}}%
\pgfpathlineto{\pgfqpoint{1.328320in}{1.949434in}}%
\pgfpathlineto{\pgfqpoint{1.328334in}{1.932302in}}%
\pgfpathlineto{\pgfqpoint{1.329056in}{1.945925in}}%
\pgfpathlineto{\pgfqpoint{1.329551in}{1.935055in}}%
\pgfpathlineto{\pgfqpoint{1.329798in}{1.938170in}}%
\pgfpathlineto{\pgfqpoint{1.329874in}{1.943141in}}%
\pgfpathlineto{\pgfqpoint{1.329981in}{1.903783in}}%
\pgfpathlineto{\pgfqpoint{1.330449in}{1.532017in}}%
\pgfpathlineto{\pgfqpoint{1.331000in}{1.532036in}}%
\pgfpathlineto{\pgfqpoint{1.332430in}{1.532597in}}%
\pgfpathlineto{\pgfqpoint{1.332532in}{1.967150in}}%
\pgfpathlineto{\pgfqpoint{1.333456in}{1.938735in}}%
\pgfpathlineto{\pgfqpoint{1.333464in}{1.933206in}}%
\pgfpathlineto{\pgfqpoint{1.333481in}{1.947552in}}%
\pgfpathlineto{\pgfqpoint{1.334211in}{1.939025in}}%
\pgfpathlineto{\pgfqpoint{1.334508in}{1.945052in}}%
\pgfpathlineto{\pgfqpoint{1.334458in}{1.935139in}}%
\pgfpathlineto{\pgfqpoint{1.334755in}{1.942039in}}%
\pgfpathlineto{\pgfqpoint{1.334854in}{1.943781in}}%
\pgfpathlineto{\pgfqpoint{1.334953in}{1.869613in}}%
\pgfpathlineto{\pgfqpoint{1.335395in}{1.532016in}}%
\pgfpathlineto{\pgfqpoint{1.335945in}{1.532036in}}%
\pgfpathlineto{\pgfqpoint{1.337375in}{1.532522in}}%
\pgfpathlineto{\pgfqpoint{1.337477in}{1.966934in}}%
\pgfpathlineto{\pgfqpoint{1.338399in}{1.939668in}}%
\pgfpathlineto{\pgfqpoint{1.338409in}{1.933135in}}%
\pgfpathlineto{\pgfqpoint{1.338425in}{1.947405in}}%
\pgfpathlineto{\pgfqpoint{1.339145in}{1.941737in}}%
\pgfpathlineto{\pgfqpoint{1.339392in}{1.945213in}}%
\pgfpathlineto{\pgfqpoint{1.339837in}{1.931643in}}%
\pgfpathlineto{\pgfqpoint{1.340340in}{1.532017in}}%
\pgfpathlineto{\pgfqpoint{1.341039in}{1.532056in}}%
\pgfpathlineto{\pgfqpoint{1.342312in}{1.532105in}}%
\pgfpathlineto{\pgfqpoint{1.342338in}{1.572062in}}%
\pgfpathlineto{\pgfqpoint{1.342423in}{1.966885in}}%
\pgfpathlineto{\pgfqpoint{1.343143in}{1.932041in}}%
\pgfpathlineto{\pgfqpoint{1.343195in}{1.948390in}}%
\pgfpathlineto{\pgfqpoint{1.343896in}{1.939929in}}%
\pgfpathlineto{\pgfqpoint{1.344143in}{1.945434in}}%
\pgfpathlineto{\pgfqpoint{1.344538in}{1.934397in}}%
\pgfpathlineto{\pgfqpoint{1.344588in}{1.944740in}}%
\pgfpathlineto{\pgfqpoint{1.344934in}{1.743421in}}%
\pgfpathlineto{\pgfqpoint{1.345286in}{1.532014in}}%
\pgfpathlineto{\pgfqpoint{1.345837in}{1.532035in}}%
\pgfpathlineto{\pgfqpoint{1.347265in}{1.532266in}}%
\pgfpathlineto{\pgfqpoint{1.347421in}{1.966466in}}%
\pgfpathlineto{\pgfqpoint{1.348353in}{1.943149in}}%
\pgfpathlineto{\pgfqpoint{1.348363in}{1.933301in}}%
\pgfpathlineto{\pgfqpoint{1.348378in}{1.946569in}}%
\pgfpathlineto{\pgfqpoint{1.349122in}{1.940694in}}%
\pgfpathlineto{\pgfqpoint{1.349132in}{1.942184in}}%
\pgfpathlineto{\pgfqpoint{1.349725in}{1.926475in}}%
\pgfpathlineto{\pgfqpoint{1.350253in}{1.532030in}}%
\pgfpathlineto{\pgfqpoint{1.350897in}{1.532043in}}%
\pgfpathlineto{\pgfqpoint{1.352206in}{1.532132in}}%
\pgfpathlineto{\pgfqpoint{1.352546in}{1.976551in}}%
\pgfpathlineto{\pgfqpoint{1.353461in}{1.941083in}}%
\pgfpathlineto{\pgfqpoint{1.353471in}{1.930477in}}%
\pgfpathlineto{\pgfqpoint{1.353487in}{1.948945in}}%
\pgfpathlineto{\pgfqpoint{1.354230in}{1.938379in}}%
\pgfpathlineto{\pgfqpoint{1.354244in}{1.943678in}}%
\pgfpathlineto{\pgfqpoint{1.354708in}{1.896054in}}%
\pgfpathlineto{\pgfqpoint{1.355177in}{1.532015in}}%
\pgfpathlineto{\pgfqpoint{1.355777in}{1.532053in}}%
\pgfpathlineto{\pgfqpoint{1.357148in}{1.532104in}}%
\pgfpathlineto{\pgfqpoint{1.357165in}{1.545341in}}%
\pgfpathlineto{\pgfqpoint{1.357246in}{1.965919in}}%
\pgfpathlineto{\pgfqpoint{1.358068in}{1.947929in}}%
\pgfpathlineto{\pgfqpoint{1.358118in}{1.931706in}}%
\pgfpathlineto{\pgfqpoint{1.358810in}{1.935969in}}%
\pgfpathlineto{\pgfqpoint{1.358958in}{1.947505in}}%
\pgfpathlineto{\pgfqpoint{1.359453in}{1.932058in}}%
\pgfpathlineto{\pgfqpoint{1.359552in}{1.935239in}}%
\pgfpathlineto{\pgfqpoint{1.360122in}{1.532013in}}%
\pgfpathlineto{\pgfqpoint{1.360872in}{1.532040in}}%
\pgfpathlineto{\pgfqpoint{1.362101in}{1.532255in}}%
\pgfpathlineto{\pgfqpoint{1.362296in}{1.967185in}}%
\pgfpathlineto{\pgfqpoint{1.363193in}{1.936255in}}%
\pgfpathlineto{\pgfqpoint{1.363198in}{1.932385in}}%
\pgfpathlineto{\pgfqpoint{1.363214in}{1.946595in}}%
\pgfpathlineto{\pgfqpoint{1.363958in}{1.938966in}}%
\pgfpathlineto{\pgfqpoint{1.363984in}{1.941947in}}%
\pgfpathlineto{\pgfqpoint{1.365057in}{1.559679in}}%
\pgfpathlineto{\pgfqpoint{1.365103in}{1.532027in}}%
\pgfpathlineto{\pgfqpoint{1.365822in}{1.532054in}}%
\pgfpathlineto{\pgfqpoint{1.367041in}{1.532121in}}%
\pgfpathlineto{\pgfqpoint{1.367081in}{1.727018in}}%
\pgfpathlineto{\pgfqpoint{1.367353in}{1.969548in}}%
\pgfpathlineto{\pgfqpoint{1.367863in}{1.950060in}}%
\pgfpathlineto{\pgfqpoint{1.367874in}{1.926499in}}%
\pgfpathlineto{\pgfqpoint{1.367890in}{1.952310in}}%
\pgfpathlineto{\pgfqpoint{1.368638in}{1.937634in}}%
\pgfpathlineto{\pgfqpoint{1.368649in}{1.945396in}}%
\pgfpathlineto{\pgfqpoint{1.369373in}{1.933588in}}%
\pgfpathlineto{\pgfqpoint{1.369422in}{1.944065in}}%
\pgfpathlineto{\pgfqpoint{1.369570in}{1.882344in}}%
\pgfpathlineto{\pgfqpoint{1.370070in}{1.532032in}}%
\pgfpathlineto{\pgfqpoint{1.370607in}{1.532046in}}%
\pgfpathlineto{\pgfqpoint{1.371986in}{1.532117in}}%
\pgfpathlineto{\pgfqpoint{1.372029in}{1.758321in}}%
\pgfpathlineto{\pgfqpoint{1.372077in}{1.964442in}}%
\pgfpathlineto{\pgfqpoint{1.372810in}{1.931045in}}%
\pgfpathlineto{\pgfqpoint{1.372824in}{1.949247in}}%
\pgfpathlineto{\pgfqpoint{1.372906in}{1.930180in}}%
\pgfpathlineto{\pgfqpoint{1.373541in}{1.945108in}}%
\pgfpathlineto{\pgfqpoint{1.374035in}{1.933719in}}%
\pgfpathlineto{\pgfqpoint{1.374283in}{1.936068in}}%
\pgfpathlineto{\pgfqpoint{1.374431in}{1.942849in}}%
\pgfpathlineto{\pgfqpoint{1.374480in}{1.919998in}}%
\pgfpathlineto{\pgfqpoint{1.375003in}{1.532030in}}%
\pgfpathlineto{\pgfqpoint{1.375535in}{1.532048in}}%
\pgfpathlineto{\pgfqpoint{1.376933in}{1.532129in}}%
\pgfpathlineto{\pgfqpoint{1.377274in}{1.973882in}}%
\pgfpathlineto{\pgfqpoint{1.378212in}{1.941003in}}%
\pgfpathlineto{\pgfqpoint{1.378222in}{1.930862in}}%
\pgfpathlineto{\pgfqpoint{1.378238in}{1.947289in}}%
\pgfpathlineto{\pgfqpoint{1.378973in}{1.938399in}}%
\pgfpathlineto{\pgfqpoint{1.378984in}{1.942450in}}%
\pgfpathlineto{\pgfqpoint{1.379445in}{1.902512in}}%
\pgfpathlineto{\pgfqpoint{1.379917in}{1.532027in}}%
\pgfpathlineto{\pgfqpoint{1.380545in}{1.532040in}}%
\pgfpathlineto{\pgfqpoint{1.381880in}{1.532147in}}%
\pgfpathlineto{\pgfqpoint{1.381971in}{1.965378in}}%
\pgfpathlineto{\pgfqpoint{1.383053in}{1.944193in}}%
\pgfpathlineto{\pgfqpoint{1.383371in}{1.933926in}}%
\pgfpathlineto{\pgfqpoint{1.383125in}{1.944737in}}%
\pgfpathlineto{\pgfqpoint{1.383816in}{1.933960in}}%
\pgfpathlineto{\pgfqpoint{1.384212in}{1.943426in}}%
\pgfpathlineto{\pgfqpoint{1.384360in}{1.926645in}}%
\pgfpathlineto{\pgfqpoint{1.384878in}{1.532031in}}%
\pgfpathlineto{\pgfqpoint{1.385541in}{1.532045in}}%
\pgfpathlineto{\pgfqpoint{1.386822in}{1.532115in}}%
\pgfpathlineto{\pgfqpoint{1.386864in}{1.742146in}}%
\pgfpathlineto{\pgfqpoint{1.387130in}{1.966615in}}%
\pgfpathlineto{\pgfqpoint{1.387646in}{1.928148in}}%
\pgfpathlineto{\pgfqpoint{1.387652in}{1.925124in}}%
\pgfpathlineto{\pgfqpoint{1.387667in}{1.952060in}}%
\pgfpathlineto{\pgfqpoint{1.388397in}{1.932999in}}%
\pgfpathlineto{\pgfqpoint{1.388485in}{1.944949in}}%
\pgfpathlineto{\pgfqpoint{1.389137in}{1.943864in}}%
\pgfpathlineto{\pgfqpoint{1.389484in}{1.703262in}}%
\pgfpathlineto{\pgfqpoint{1.389822in}{1.532031in}}%
\pgfpathlineto{\pgfqpoint{1.390386in}{1.532042in}}%
\pgfpathlineto{\pgfqpoint{1.391768in}{1.532113in}}%
\pgfpathlineto{\pgfqpoint{1.391810in}{1.737529in}}%
\pgfpathlineto{\pgfqpoint{1.392071in}{1.967120in}}%
\pgfpathlineto{\pgfqpoint{1.392595in}{1.942011in}}%
\pgfpathlineto{\pgfqpoint{1.392632in}{1.952249in}}%
\pgfpathlineto{\pgfqpoint{1.392615in}{1.925343in}}%
\pgfpathlineto{\pgfqpoint{1.393355in}{1.945248in}}%
\pgfpathlineto{\pgfqpoint{1.394062in}{1.932379in}}%
\pgfpathlineto{\pgfqpoint{1.394112in}{1.943715in}}%
\pgfpathlineto{\pgfqpoint{1.394408in}{1.725837in}}%
\pgfpathlineto{\pgfqpoint{1.394769in}{1.532031in}}%
\pgfpathlineto{\pgfqpoint{1.395333in}{1.532042in}}%
\pgfpathlineto{\pgfqpoint{1.396713in}{1.532115in}}%
\pgfpathlineto{\pgfqpoint{1.396755in}{1.743346in}}%
\pgfpathlineto{\pgfqpoint{1.397020in}{1.966082in}}%
\pgfpathlineto{\pgfqpoint{1.397537in}{1.924769in}}%
\pgfpathlineto{\pgfqpoint{1.397553in}{1.951967in}}%
\pgfpathlineto{\pgfqpoint{1.398305in}{1.940834in}}%
\pgfpathlineto{\pgfqpoint{1.398315in}{1.932401in}}%
\pgfpathlineto{\pgfqpoint{1.398576in}{1.944682in}}%
\pgfpathlineto{\pgfqpoint{1.399071in}{1.932676in}}%
\pgfpathlineto{\pgfqpoint{1.399120in}{1.941409in}}%
\pgfpathlineto{\pgfqpoint{1.399268in}{1.840407in}}%
\pgfpathlineto{\pgfqpoint{1.399732in}{1.532030in}}%
\pgfpathlineto{\pgfqpoint{1.400266in}{1.532048in}}%
\pgfpathlineto{\pgfqpoint{1.401660in}{1.532132in}}%
\pgfpathlineto{\pgfqpoint{1.401995in}{1.972473in}}%
\pgfpathlineto{\pgfqpoint{1.402908in}{1.931672in}}%
\pgfpathlineto{\pgfqpoint{1.402913in}{1.929491in}}%
\pgfpathlineto{\pgfqpoint{1.402929in}{1.947371in}}%
\pgfpathlineto{\pgfqpoint{1.403664in}{1.934448in}}%
\pgfpathlineto{\pgfqpoint{1.403760in}{1.942236in}}%
\pgfpathlineto{\pgfqpoint{1.404150in}{1.911900in}}%
\pgfpathlineto{\pgfqpoint{1.404681in}{1.532029in}}%
\pgfpathlineto{\pgfqpoint{1.405266in}{1.532042in}}%
\pgfpathlineto{\pgfqpoint{1.406607in}{1.532143in}}%
\pgfpathlineto{\pgfqpoint{1.406818in}{1.966911in}}%
\pgfpathlineto{\pgfqpoint{1.407794in}{1.939496in}}%
\pgfpathlineto{\pgfqpoint{1.407804in}{1.930861in}}%
\pgfpathlineto{\pgfqpoint{1.407824in}{1.945767in}}%
\pgfpathlineto{\pgfqpoint{1.408550in}{1.940224in}}%
\pgfpathlineto{\pgfqpoint{1.408797in}{1.942853in}}%
\pgfpathlineto{\pgfqpoint{1.409094in}{1.901400in}}%
\pgfpathlineto{\pgfqpoint{1.409606in}{1.532031in}}%
\pgfpathlineto{\pgfqpoint{1.410320in}{1.532050in}}%
\pgfpathlineto{\pgfqpoint{1.411548in}{1.532103in}}%
\pgfpathlineto{\pgfqpoint{1.411556in}{1.532357in}}%
\pgfpathlineto{\pgfqpoint{1.411663in}{1.965174in}}%
\pgfpathlineto{\pgfqpoint{1.412605in}{1.938836in}}%
\pgfpathlineto{\pgfqpoint{1.412647in}{1.930982in}}%
\pgfpathlineto{\pgfqpoint{1.412660in}{1.945139in}}%
\pgfpathlineto{\pgfqpoint{1.413360in}{1.942237in}}%
\pgfpathlineto{\pgfqpoint{1.413706in}{1.942373in}}%
\pgfpathlineto{\pgfqpoint{1.414151in}{1.780469in}}%
\pgfpathlineto{\pgfqpoint{1.414559in}{1.532031in}}%
\pgfpathlineto{\pgfqpoint{1.415084in}{1.532046in}}%
\pgfpathlineto{\pgfqpoint{1.416495in}{1.532112in}}%
\pgfpathlineto{\pgfqpoint{1.416537in}{1.735929in}}%
\pgfpathlineto{\pgfqpoint{1.416598in}{1.964355in}}%
\pgfpathlineto{\pgfqpoint{1.417320in}{1.945250in}}%
\pgfpathlineto{\pgfqpoint{1.417332in}{1.923952in}}%
\pgfpathlineto{\pgfqpoint{1.417348in}{1.951300in}}%
\pgfpathlineto{\pgfqpoint{1.418095in}{1.934224in}}%
\pgfpathlineto{\pgfqpoint{1.418356in}{1.944818in}}%
\pgfpathlineto{\pgfqpoint{1.418751in}{1.930946in}}%
\pgfpathlineto{\pgfqpoint{1.418850in}{1.932422in}}%
\pgfpathlineto{\pgfqpoint{1.418900in}{1.936447in}}%
\pgfpathlineto{\pgfqpoint{1.418999in}{1.892492in}}%
\pgfpathlineto{\pgfqpoint{1.419496in}{1.532030in}}%
\pgfpathlineto{\pgfqpoint{1.420160in}{1.532044in}}%
\pgfpathlineto{\pgfqpoint{1.421441in}{1.532115in}}%
\pgfpathlineto{\pgfqpoint{1.421483in}{1.748092in}}%
\pgfpathlineto{\pgfqpoint{1.421745in}{1.965179in}}%
\pgfpathlineto{\pgfqpoint{1.422267in}{1.936686in}}%
\pgfpathlineto{\pgfqpoint{1.422276in}{1.951268in}}%
\pgfpathlineto{\pgfqpoint{1.422292in}{1.925181in}}%
\pgfpathlineto{\pgfqpoint{1.423028in}{1.940691in}}%
\pgfpathlineto{\pgfqpoint{1.423563in}{1.932065in}}%
\pgfpathlineto{\pgfqpoint{1.423177in}{1.944036in}}%
\pgfpathlineto{\pgfqpoint{1.423761in}{1.935755in}}%
\pgfpathlineto{\pgfqpoint{1.423860in}{1.939349in}}%
\pgfpathlineto{\pgfqpoint{1.424449in}{1.532031in}}%
\pgfpathlineto{\pgfqpoint{1.425219in}{1.532047in}}%
\pgfpathlineto{\pgfqpoint{1.426392in}{1.532310in}}%
\pgfpathlineto{\pgfqpoint{1.426489in}{1.963776in}}%
\pgfpathlineto{\pgfqpoint{1.427450in}{1.943684in}}%
\pgfpathlineto{\pgfqpoint{1.427541in}{1.930850in}}%
\pgfpathlineto{\pgfqpoint{1.427557in}{1.944893in}}%
\pgfpathlineto{\pgfqpoint{1.428263in}{1.933094in}}%
\pgfpathlineto{\pgfqpoint{1.428658in}{1.942275in}}%
\pgfpathlineto{\pgfqpoint{1.428856in}{1.909281in}}%
\pgfpathlineto{\pgfqpoint{1.429389in}{1.532030in}}%
\pgfpathlineto{\pgfqpoint{1.430351in}{1.532051in}}%
\pgfpathlineto{\pgfqpoint{1.431332in}{1.532117in}}%
\pgfpathlineto{\pgfqpoint{1.431374in}{1.765138in}}%
\pgfpathlineto{\pgfqpoint{1.431434in}{1.964132in}}%
\pgfpathlineto{\pgfqpoint{1.432154in}{1.938616in}}%
\pgfpathlineto{\pgfqpoint{1.432164in}{1.925800in}}%
\pgfpathlineto{\pgfqpoint{1.432239in}{1.949354in}}%
\pgfpathlineto{\pgfqpoint{1.432919in}{1.940878in}}%
\pgfpathlineto{\pgfqpoint{1.433065in}{1.943026in}}%
\pgfpathlineto{\pgfqpoint{1.433429in}{1.932642in}}%
\pgfpathlineto{\pgfqpoint{1.433627in}{1.935327in}}%
\pgfpathlineto{\pgfqpoint{1.433726in}{1.938541in}}%
\pgfpathlineto{\pgfqpoint{1.433824in}{1.918995in}}%
\pgfpathlineto{\pgfqpoint{1.434348in}{1.532030in}}%
\pgfpathlineto{\pgfqpoint{1.435027in}{1.532044in}}%
\pgfpathlineto{\pgfqpoint{1.436280in}{1.532145in}}%
\pgfpathlineto{\pgfqpoint{1.436491in}{1.966306in}}%
\pgfpathlineto{\pgfqpoint{1.437466in}{1.931116in}}%
\pgfpathlineto{\pgfqpoint{1.437518in}{1.945469in}}%
\pgfpathlineto{\pgfqpoint{1.437505in}{1.929972in}}%
\pgfpathlineto{\pgfqpoint{1.438276in}{1.939443in}}%
\pgfpathlineto{\pgfqpoint{1.438524in}{1.942210in}}%
\pgfpathlineto{\pgfqpoint{1.438771in}{1.924003in}}%
\pgfpathlineto{\pgfqpoint{1.439298in}{1.532032in}}%
\pgfpathlineto{\pgfqpoint{1.439980in}{1.532047in}}%
\pgfpathlineto{\pgfqpoint{1.441223in}{1.532117in}}%
\pgfpathlineto{\pgfqpoint{1.441265in}{1.763172in}}%
\pgfpathlineto{\pgfqpoint{1.441507in}{1.964987in}}%
\pgfpathlineto{\pgfqpoint{1.442047in}{1.927964in}}%
\pgfpathlineto{\pgfqpoint{1.442053in}{1.925555in}}%
\pgfpathlineto{\pgfqpoint{1.442128in}{1.949330in}}%
\pgfpathlineto{\pgfqpoint{1.442799in}{1.932890in}}%
\pgfpathlineto{\pgfqpoint{1.442897in}{1.942829in}}%
\pgfpathlineto{\pgfqpoint{1.442992in}{1.932418in}}%
\pgfpathlineto{\pgfqpoint{1.443572in}{1.937731in}}%
\pgfpathlineto{\pgfqpoint{1.443621in}{1.937844in}}%
\pgfpathlineto{\pgfqpoint{1.444192in}{1.532028in}}%
\pgfpathlineto{\pgfqpoint{1.445624in}{1.532060in}}%
\pgfpathlineto{\pgfqpoint{1.446168in}{1.532120in}}%
\pgfpathlineto{\pgfqpoint{1.446270in}{1.964407in}}%
\pgfpathlineto{\pgfqpoint{1.447681in}{1.941210in}}%
\pgfpathlineto{\pgfqpoint{1.448063in}{1.933930in}}%
\pgfpathlineto{\pgfqpoint{1.448458in}{1.940253in}}%
\pgfpathlineto{\pgfqpoint{1.448755in}{1.819360in}}%
\pgfpathlineto{\pgfqpoint{1.449185in}{1.532029in}}%
\pgfpathlineto{\pgfqpoint{1.449718in}{1.532048in}}%
\pgfpathlineto{\pgfqpoint{1.451116in}{1.532146in}}%
\pgfpathlineto{\pgfqpoint{1.451311in}{1.963193in}}%
\pgfpathlineto{\pgfqpoint{1.452297in}{1.932437in}}%
\pgfpathlineto{\pgfqpoint{1.452370in}{1.930501in}}%
\pgfpathlineto{\pgfqpoint{1.452318in}{1.944854in}}%
\pgfpathlineto{\pgfqpoint{1.452993in}{1.932677in}}%
\pgfpathlineto{\pgfqpoint{1.453389in}{1.941721in}}%
\pgfpathlineto{\pgfqpoint{1.453587in}{1.910653in}}%
\pgfpathlineto{\pgfqpoint{1.454099in}{1.532027in}}%
\pgfpathlineto{\pgfqpoint{1.455171in}{1.532061in}}%
\pgfpathlineto{\pgfqpoint{1.456051in}{1.532085in}}%
\pgfpathlineto{\pgfqpoint{1.456062in}{1.532147in}}%
\pgfpathlineto{\pgfqpoint{1.456153in}{1.962246in}}%
\pgfpathlineto{\pgfqpoint{1.457225in}{1.935670in}}%
\pgfpathlineto{\pgfqpoint{1.457233in}{1.931036in}}%
\pgfpathlineto{\pgfqpoint{1.457252in}{1.943363in}}%
\pgfpathlineto{\pgfqpoint{1.457954in}{1.936093in}}%
\pgfpathlineto{\pgfqpoint{1.458102in}{1.941869in}}%
\pgfpathlineto{\pgfqpoint{1.458498in}{1.930104in}}%
\pgfpathlineto{\pgfqpoint{1.458547in}{1.931162in}}%
\pgfpathlineto{\pgfqpoint{1.459075in}{1.532033in}}%
\pgfpathlineto{\pgfqpoint{1.459755in}{1.532048in}}%
\pgfpathlineto{\pgfqpoint{1.461010in}{1.532249in}}%
\pgfpathlineto{\pgfqpoint{1.461314in}{1.967887in}}%
\pgfpathlineto{\pgfqpoint{1.462095in}{1.947756in}}%
\pgfpathlineto{\pgfqpoint{1.462110in}{1.927066in}}%
\pgfpathlineto{\pgfqpoint{1.462864in}{1.939178in}}%
\pgfpathlineto{\pgfqpoint{1.463006in}{1.941322in}}%
\pgfpathlineto{\pgfqpoint{1.463353in}{1.932733in}}%
\pgfpathlineto{\pgfqpoint{1.463452in}{1.934097in}}%
\pgfpathlineto{\pgfqpoint{1.464019in}{1.532032in}}%
\pgfpathlineto{\pgfqpoint{1.464994in}{1.532053in}}%
\pgfpathlineto{\pgfqpoint{1.465949in}{1.532114in}}%
\pgfpathlineto{\pgfqpoint{1.465992in}{1.740699in}}%
\pgfpathlineto{\pgfqpoint{1.466259in}{1.963836in}}%
\pgfpathlineto{\pgfqpoint{1.466775in}{1.926294in}}%
\pgfpathlineto{\pgfqpoint{1.466781in}{1.923204in}}%
\pgfpathlineto{\pgfqpoint{1.466797in}{1.950451in}}%
\pgfpathlineto{\pgfqpoint{1.467526in}{1.931194in}}%
\pgfpathlineto{\pgfqpoint{1.467618in}{1.943334in}}%
\pgfpathlineto{\pgfqpoint{1.468275in}{1.941793in}}%
\pgfpathlineto{\pgfqpoint{1.468621in}{1.695196in}}%
\pgfpathlineto{\pgfqpoint{1.468936in}{1.532027in}}%
\pgfpathlineto{\pgfqpoint{1.469514in}{1.532050in}}%
\pgfpathlineto{\pgfqpoint{1.470887in}{1.532085in}}%
\pgfpathlineto{\pgfqpoint{1.470898in}{1.532147in}}%
\pgfpathlineto{\pgfqpoint{1.470992in}{1.963220in}}%
\pgfpathlineto{\pgfqpoint{1.472051in}{1.932514in}}%
\pgfpathlineto{\pgfqpoint{1.472289in}{1.942280in}}%
\pgfpathlineto{\pgfqpoint{1.472240in}{1.931943in}}%
\pgfpathlineto{\pgfqpoint{1.472833in}{1.940084in}}%
\pgfpathlineto{\pgfqpoint{1.473179in}{1.941360in}}%
\pgfpathlineto{\pgfqpoint{1.473476in}{1.833133in}}%
\pgfpathlineto{\pgfqpoint{1.473906in}{1.532030in}}%
\pgfpathlineto{\pgfqpoint{1.474432in}{1.532047in}}%
\pgfpathlineto{\pgfqpoint{1.475841in}{1.532123in}}%
\pgfpathlineto{\pgfqpoint{1.476061in}{1.964689in}}%
\pgfpathlineto{\pgfqpoint{1.477206in}{1.939104in}}%
\pgfpathlineto{\pgfqpoint{1.477215in}{1.931219in}}%
\pgfpathlineto{\pgfqpoint{1.477232in}{1.942468in}}%
\pgfpathlineto{\pgfqpoint{1.477981in}{1.935161in}}%
\pgfpathlineto{\pgfqpoint{1.477993in}{1.938208in}}%
\pgfpathlineto{\pgfqpoint{1.478390in}{1.873415in}}%
\pgfpathlineto{\pgfqpoint{1.478864in}{1.532033in}}%
\pgfpathlineto{\pgfqpoint{1.479350in}{1.532043in}}%
\pgfpathlineto{\pgfqpoint{1.480789in}{1.532155in}}%
\pgfpathlineto{\pgfqpoint{1.481106in}{1.966310in}}%
\pgfpathlineto{\pgfqpoint{1.481992in}{1.931258in}}%
\pgfpathlineto{\pgfqpoint{1.482008in}{1.942483in}}%
\pgfpathlineto{\pgfqpoint{1.482790in}{1.937780in}}%
\pgfpathlineto{\pgfqpoint{1.483857in}{1.532033in}}%
\pgfpathlineto{\pgfqpoint{1.482855in}{1.938275in}}%
\pgfpathlineto{\pgfqpoint{1.484846in}{1.532055in}}%
\pgfpathlineto{\pgfqpoint{1.485732in}{1.532120in}}%
\pgfpathlineto{\pgfqpoint{1.485833in}{1.963672in}}%
\pgfpathlineto{\pgfqpoint{1.487260in}{1.934018in}}%
\pgfpathlineto{\pgfqpoint{1.487377in}{1.940293in}}%
\pgfpathlineto{\pgfqpoint{1.487612in}{1.933198in}}%
\pgfpathlineto{\pgfqpoint{1.488008in}{1.939512in}}%
\pgfpathlineto{\pgfqpoint{1.488305in}{1.851599in}}%
\pgfpathlineto{\pgfqpoint{1.488702in}{1.532026in}}%
\pgfpathlineto{\pgfqpoint{1.489280in}{1.532041in}}%
\pgfpathlineto{\pgfqpoint{1.490678in}{1.532130in}}%
\pgfpathlineto{\pgfqpoint{1.490993in}{1.965674in}}%
\pgfpathlineto{\pgfqpoint{1.491994in}{1.935162in}}%
\pgfpathlineto{\pgfqpoint{1.491999in}{1.931099in}}%
\pgfpathlineto{\pgfqpoint{1.492014in}{1.942163in}}%
\pgfpathlineto{\pgfqpoint{1.492757in}{1.936308in}}%
\pgfpathlineto{\pgfqpoint{1.492765in}{1.938314in}}%
\pgfpathlineto{\pgfqpoint{1.493211in}{1.902916in}}%
\pgfpathlineto{\pgfqpoint{1.493687in}{1.532031in}}%
\pgfpathlineto{\pgfqpoint{1.494263in}{1.532043in}}%
\pgfpathlineto{\pgfqpoint{1.495623in}{1.532123in}}%
\pgfpathlineto{\pgfqpoint{1.495820in}{1.964068in}}%
\pgfpathlineto{\pgfqpoint{1.496987in}{1.937175in}}%
\pgfpathlineto{\pgfqpoint{1.496995in}{1.931718in}}%
\pgfpathlineto{\pgfqpoint{1.497039in}{1.941244in}}%
\pgfpathlineto{\pgfqpoint{1.497752in}{1.937257in}}%
\pgfpathlineto{\pgfqpoint{1.497781in}{1.937679in}}%
\pgfpathlineto{\pgfqpoint{1.498581in}{1.568715in}}%
\pgfpathlineto{\pgfqpoint{1.498679in}{1.532032in}}%
\pgfpathlineto{\pgfqpoint{1.499316in}{1.532049in}}%
\pgfpathlineto{\pgfqpoint{1.500568in}{1.532120in}}%
\pgfpathlineto{\pgfqpoint{1.500670in}{1.963183in}}%
\pgfpathlineto{\pgfqpoint{1.502080in}{1.940092in}}%
\pgfpathlineto{\pgfqpoint{1.502574in}{1.932886in}}%
\pgfpathlineto{\pgfqpoint{1.502822in}{1.934996in}}%
\pgfpathlineto{\pgfqpoint{1.502970in}{1.939079in}}%
\pgfpathlineto{\pgfqpoint{1.503069in}{1.920496in}}%
\pgfpathlineto{\pgfqpoint{1.503637in}{1.532033in}}%
\pgfpathlineto{\pgfqpoint{1.504379in}{1.532051in}}%
\pgfpathlineto{\pgfqpoint{1.505515in}{1.532136in}}%
\pgfpathlineto{\pgfqpoint{1.505845in}{1.968321in}}%
\pgfpathlineto{\pgfqpoint{1.506794in}{1.943465in}}%
\pgfpathlineto{\pgfqpoint{1.506810in}{1.929632in}}%
\pgfpathlineto{\pgfqpoint{1.507577in}{1.934120in}}%
\pgfpathlineto{\pgfqpoint{1.507631in}{1.938617in}}%
\pgfpathlineto{\pgfqpoint{1.508051in}{1.899796in}}%
\pgfpathlineto{\pgfqpoint{1.508520in}{1.532031in}}%
\pgfpathlineto{\pgfqpoint{1.509092in}{1.532043in}}%
\pgfpathlineto{\pgfqpoint{1.510465in}{1.532310in}}%
\pgfpathlineto{\pgfqpoint{1.510560in}{1.962492in}}%
\pgfpathlineto{\pgfqpoint{1.511530in}{1.943685in}}%
\pgfpathlineto{\pgfqpoint{1.511549in}{1.929033in}}%
\pgfpathlineto{\pgfqpoint{1.512275in}{1.932256in}}%
\pgfpathlineto{\pgfqpoint{1.512670in}{1.940583in}}%
\pgfpathlineto{\pgfqpoint{1.512621in}{1.931272in}}%
\pgfpathlineto{\pgfqpoint{1.512918in}{1.937522in}}%
\pgfpathlineto{\pgfqpoint{1.513429in}{1.532024in}}%
\pgfpathlineto{\pgfqpoint{1.514282in}{1.532054in}}%
\pgfpathlineto{\pgfqpoint{1.515405in}{1.532121in}}%
\pgfpathlineto{\pgfqpoint{1.515703in}{1.963156in}}%
\pgfpathlineto{\pgfqpoint{1.516945in}{1.933344in}}%
\pgfpathlineto{\pgfqpoint{1.517478in}{1.930980in}}%
\pgfpathlineto{\pgfqpoint{1.517007in}{1.941865in}}%
\pgfpathlineto{\pgfqpoint{1.517676in}{1.933799in}}%
\pgfpathlineto{\pgfqpoint{1.517825in}{1.931801in}}%
\pgfpathlineto{\pgfqpoint{1.517874in}{1.939287in}}%
\pgfpathlineto{\pgfqpoint{1.518444in}{1.532032in}}%
\pgfpathlineto{\pgfqpoint{1.519230in}{1.532050in}}%
\pgfpathlineto{\pgfqpoint{1.520350in}{1.532122in}}%
\pgfpathlineto{\pgfqpoint{1.520651in}{1.963159in}}%
\pgfpathlineto{\pgfqpoint{1.521890in}{1.942420in}}%
\pgfpathlineto{\pgfqpoint{1.521908in}{1.930183in}}%
\pgfpathlineto{\pgfqpoint{1.522667in}{1.937659in}}%
\pgfpathlineto{\pgfqpoint{1.522766in}{1.939402in}}%
\pgfpathlineto{\pgfqpoint{1.522915in}{1.876153in}}%
\pgfpathlineto{\pgfqpoint{1.523388in}{1.532032in}}%
\pgfpathlineto{\pgfqpoint{1.523925in}{1.532046in}}%
\pgfpathlineto{\pgfqpoint{1.525296in}{1.532122in}}%
\pgfpathlineto{\pgfqpoint{1.525398in}{1.962903in}}%
\pgfpathlineto{\pgfqpoint{1.526807in}{1.934077in}}%
\pgfpathlineto{\pgfqpoint{1.526814in}{1.930842in}}%
\pgfpathlineto{\pgfqpoint{1.526833in}{1.941254in}}%
\pgfpathlineto{\pgfqpoint{1.527562in}{1.931334in}}%
\pgfpathlineto{\pgfqpoint{1.527611in}{1.939631in}}%
\pgfpathlineto{\pgfqpoint{1.527858in}{1.859968in}}%
\pgfpathlineto{\pgfqpoint{1.528265in}{1.532030in}}%
\pgfpathlineto{\pgfqpoint{1.528856in}{1.532045in}}%
\pgfpathlineto{\pgfqpoint{1.530241in}{1.532121in}}%
\pgfpathlineto{\pgfqpoint{1.530546in}{1.963676in}}%
\pgfpathlineto{\pgfqpoint{1.531796in}{1.941622in}}%
\pgfpathlineto{\pgfqpoint{1.531814in}{1.930615in}}%
\pgfpathlineto{\pgfqpoint{1.532532in}{1.934973in}}%
\pgfpathlineto{\pgfqpoint{1.532680in}{1.939421in}}%
\pgfpathlineto{\pgfqpoint{1.532730in}{1.924402in}}%
\pgfpathlineto{\pgfqpoint{1.533249in}{1.532031in}}%
\pgfpathlineto{\pgfqpoint{1.534168in}{1.532055in}}%
\pgfpathlineto{\pgfqpoint{1.535186in}{1.532113in}}%
\pgfpathlineto{\pgfqpoint{1.535228in}{1.736134in}}%
\pgfpathlineto{\pgfqpoint{1.535487in}{1.964077in}}%
\pgfpathlineto{\pgfqpoint{1.536013in}{1.939224in}}%
\pgfpathlineto{\pgfqpoint{1.536050in}{1.949451in}}%
\pgfpathlineto{\pgfqpoint{1.536034in}{1.922403in}}%
\pgfpathlineto{\pgfqpoint{1.536772in}{1.942391in}}%
\pgfpathlineto{\pgfqpoint{1.537470in}{1.929752in}}%
\pgfpathlineto{\pgfqpoint{1.537520in}{1.940245in}}%
\pgfpathlineto{\pgfqpoint{1.537816in}{1.783407in}}%
\pgfpathlineto{\pgfqpoint{1.538156in}{1.532024in}}%
\pgfpathlineto{\pgfqpoint{1.538737in}{1.532048in}}%
\pgfpathlineto{\pgfqpoint{1.540134in}{1.532149in}}%
\pgfpathlineto{\pgfqpoint{1.540229in}{1.962335in}}%
\pgfpathlineto{\pgfqpoint{1.541288in}{1.934596in}}%
\pgfpathlineto{\pgfqpoint{1.541420in}{1.941624in}}%
\pgfpathlineto{\pgfqpoint{1.541469in}{1.930363in}}%
\pgfpathlineto{\pgfqpoint{1.542062in}{1.937125in}}%
\pgfpathlineto{\pgfqpoint{1.542359in}{1.930362in}}%
\pgfpathlineto{\pgfqpoint{1.542310in}{1.940959in}}%
\pgfpathlineto{\pgfqpoint{1.542557in}{1.935800in}}%
\pgfpathlineto{\pgfqpoint{1.543102in}{1.532029in}}%
\pgfpathlineto{\pgfqpoint{1.544048in}{1.532053in}}%
\pgfpathlineto{\pgfqpoint{1.545079in}{1.532140in}}%
\pgfpathlineto{\pgfqpoint{1.545394in}{1.965716in}}%
\pgfpathlineto{\pgfqpoint{1.546332in}{1.935472in}}%
\pgfpathlineto{\pgfqpoint{1.546340in}{1.930371in}}%
\pgfpathlineto{\pgfqpoint{1.546352in}{1.940879in}}%
\pgfpathlineto{\pgfqpoint{1.547104in}{1.935084in}}%
\pgfpathlineto{\pgfqpoint{1.547124in}{1.937124in}}%
\pgfpathlineto{\pgfqpoint{1.547732in}{1.737701in}}%
\pgfpathlineto{\pgfqpoint{1.548136in}{1.532033in}}%
\pgfpathlineto{\pgfqpoint{1.548627in}{1.532044in}}%
\pgfpathlineto{\pgfqpoint{1.550024in}{1.532134in}}%
\pgfpathlineto{\pgfqpoint{1.550237in}{1.962311in}}%
\pgfpathlineto{\pgfqpoint{1.551293in}{1.934785in}}%
\pgfpathlineto{\pgfqpoint{1.551329in}{1.929627in}}%
\pgfpathlineto{\pgfqpoint{1.551316in}{1.941880in}}%
\pgfpathlineto{\pgfqpoint{1.552066in}{1.933631in}}%
\pgfpathlineto{\pgfqpoint{1.552112in}{1.937264in}}%
\pgfpathlineto{\pgfqpoint{1.552557in}{1.912786in}}%
\pgfpathlineto{\pgfqpoint{1.553055in}{1.532032in}}%
\pgfpathlineto{\pgfqpoint{1.553642in}{1.532045in}}%
\pgfpathlineto{\pgfqpoint{1.554969in}{1.532130in}}%
\pgfpathlineto{\pgfqpoint{1.555070in}{1.962111in}}%
\pgfpathlineto{\pgfqpoint{1.556372in}{1.935038in}}%
\pgfpathlineto{\pgfqpoint{1.556383in}{1.930947in}}%
\pgfpathlineto{\pgfqpoint{1.556401in}{1.940347in}}%
\pgfpathlineto{\pgfqpoint{1.557131in}{1.938573in}}%
\pgfpathlineto{\pgfqpoint{1.557230in}{1.939051in}}%
\pgfpathlineto{\pgfqpoint{1.557576in}{1.828602in}}%
\pgfpathlineto{\pgfqpoint{1.557988in}{1.532032in}}%
\pgfpathlineto{\pgfqpoint{1.558524in}{1.532045in}}%
\pgfpathlineto{\pgfqpoint{1.559917in}{1.532174in}}%
\pgfpathlineto{\pgfqpoint{1.560224in}{1.969207in}}%
\pgfpathlineto{\pgfqpoint{1.561083in}{1.934881in}}%
\pgfpathlineto{\pgfqpoint{1.561119in}{1.926433in}}%
\pgfpathlineto{\pgfqpoint{1.561106in}{1.944886in}}%
\pgfpathlineto{\pgfqpoint{1.561851in}{1.937172in}}%
\pgfpathlineto{\pgfqpoint{1.562920in}{1.532033in}}%
\pgfpathlineto{\pgfqpoint{1.561914in}{1.938303in}}%
\pgfpathlineto{\pgfqpoint{1.563343in}{1.532042in}}%
\pgfpathlineto{\pgfqpoint{1.564862in}{1.532148in}}%
\pgfpathlineto{\pgfqpoint{1.565050in}{1.960451in}}%
\pgfpathlineto{\pgfqpoint{1.566095in}{1.940265in}}%
\pgfpathlineto{\pgfqpoint{1.566109in}{1.930358in}}%
\pgfpathlineto{\pgfqpoint{1.566122in}{1.940702in}}%
\pgfpathlineto{\pgfqpoint{1.566870in}{1.936297in}}%
\pgfpathlineto{\pgfqpoint{1.567882in}{1.532032in}}%
\pgfpathlineto{\pgfqpoint{1.566933in}{1.936853in}}%
\pgfpathlineto{\pgfqpoint{1.569160in}{1.532060in}}%
\pgfpathlineto{\pgfqpoint{1.569806in}{1.532127in}}%
\pgfpathlineto{\pgfqpoint{1.570099in}{1.962405in}}%
\pgfpathlineto{\pgfqpoint{1.571249in}{1.934185in}}%
\pgfpathlineto{\pgfqpoint{1.571260in}{1.929804in}}%
\pgfpathlineto{\pgfqpoint{1.571312in}{1.941237in}}%
\pgfpathlineto{\pgfqpoint{1.571975in}{1.937965in}}%
\pgfpathlineto{\pgfqpoint{1.572222in}{1.938922in}}%
\pgfpathlineto{\pgfqpoint{1.572371in}{1.878673in}}%
\pgfpathlineto{\pgfqpoint{1.572809in}{1.532032in}}%
\pgfpathlineto{\pgfqpoint{1.573379in}{1.532044in}}%
\pgfpathlineto{\pgfqpoint{1.574751in}{1.532126in}}%
\pgfpathlineto{\pgfqpoint{1.574970in}{1.963351in}}%
\pgfpathlineto{\pgfqpoint{1.576077in}{1.939970in}}%
\pgfpathlineto{\pgfqpoint{1.576119in}{1.930017in}}%
\pgfpathlineto{\pgfqpoint{1.576106in}{1.940554in}}%
\pgfpathlineto{\pgfqpoint{1.576851in}{1.936617in}}%
\pgfpathlineto{\pgfqpoint{1.576922in}{1.936638in}}%
\pgfpathlineto{\pgfqpoint{1.577305in}{1.894806in}}%
\pgfpathlineto{\pgfqpoint{1.577762in}{1.532032in}}%
\pgfpathlineto{\pgfqpoint{1.578340in}{1.532044in}}%
\pgfpathlineto{\pgfqpoint{1.579697in}{1.532129in}}%
\pgfpathlineto{\pgfqpoint{1.579798in}{1.962037in}}%
\pgfpathlineto{\pgfqpoint{1.581112in}{1.937411in}}%
\pgfpathlineto{\pgfqpoint{1.581121in}{1.930057in}}%
\pgfpathlineto{\pgfqpoint{1.581143in}{1.940955in}}%
\pgfpathlineto{\pgfqpoint{1.581883in}{1.931379in}}%
\pgfpathlineto{\pgfqpoint{1.581933in}{1.938000in}}%
\pgfpathlineto{\pgfqpoint{1.582230in}{1.906416in}}%
\pgfpathlineto{\pgfqpoint{1.582712in}{1.532032in}}%
\pgfpathlineto{\pgfqpoint{1.583344in}{1.532047in}}%
\pgfpathlineto{\pgfqpoint{1.584642in}{1.532123in}}%
\pgfpathlineto{\pgfqpoint{1.584950in}{1.963208in}}%
\pgfpathlineto{\pgfqpoint{1.586147in}{1.933675in}}%
\pgfpathlineto{\pgfqpoint{1.586156in}{1.929260in}}%
\pgfpathlineto{\pgfqpoint{1.586174in}{1.941273in}}%
\pgfpathlineto{\pgfqpoint{1.586903in}{1.933242in}}%
\pgfpathlineto{\pgfqpoint{1.587051in}{1.939277in}}%
\pgfpathlineto{\pgfqpoint{1.587150in}{1.916815in}}%
\pgfpathlineto{\pgfqpoint{1.587667in}{1.532032in}}%
\pgfpathlineto{\pgfqpoint{1.588452in}{1.532049in}}%
\pgfpathlineto{\pgfqpoint{1.589587in}{1.532121in}}%
\pgfpathlineto{\pgfqpoint{1.589677in}{1.959813in}}%
\pgfpathlineto{\pgfqpoint{1.591116in}{1.933887in}}%
\pgfpathlineto{\pgfqpoint{1.591127in}{1.930025in}}%
\pgfpathlineto{\pgfqpoint{1.591144in}{1.940228in}}%
\pgfpathlineto{\pgfqpoint{1.591867in}{1.930252in}}%
\pgfpathlineto{\pgfqpoint{1.591917in}{1.938260in}}%
\pgfpathlineto{\pgfqpoint{1.592114in}{1.896067in}}%
\pgfpathlineto{\pgfqpoint{1.592556in}{1.532024in}}%
\pgfpathlineto{\pgfqpoint{1.593282in}{1.532043in}}%
\pgfpathlineto{\pgfqpoint{1.594534in}{1.532148in}}%
\pgfpathlineto{\pgfqpoint{1.594628in}{1.960450in}}%
\pgfpathlineto{\pgfqpoint{1.595707in}{1.938599in}}%
\pgfpathlineto{\pgfqpoint{1.595832in}{1.929624in}}%
\pgfpathlineto{\pgfqpoint{1.595783in}{1.940728in}}%
\pgfpathlineto{\pgfqpoint{1.596475in}{1.935377in}}%
\pgfpathlineto{\pgfqpoint{1.596673in}{1.940147in}}%
\pgfpathlineto{\pgfqpoint{1.597090in}{1.886263in}}%
\pgfpathlineto{\pgfqpoint{1.597582in}{1.532032in}}%
\pgfpathlineto{\pgfqpoint{1.598071in}{1.532043in}}%
\pgfpathlineto{\pgfqpoint{1.599478in}{1.532122in}}%
\pgfpathlineto{\pgfqpoint{1.599800in}{1.962112in}}%
\pgfpathlineto{\pgfqpoint{1.601037in}{1.933815in}}%
\pgfpathlineto{\pgfqpoint{1.601048in}{1.929032in}}%
\pgfpathlineto{\pgfqpoint{1.601066in}{1.941033in}}%
\pgfpathlineto{\pgfqpoint{1.601759in}{1.933129in}}%
\pgfpathlineto{\pgfqpoint{1.601908in}{1.939284in}}%
\pgfpathlineto{\pgfqpoint{1.602007in}{1.913120in}}%
\pgfpathlineto{\pgfqpoint{1.602500in}{1.532032in}}%
\pgfpathlineto{\pgfqpoint{1.603236in}{1.532049in}}%
\pgfpathlineto{\pgfqpoint{1.604426in}{1.532172in}}%
\pgfpathlineto{\pgfqpoint{1.604748in}{1.966008in}}%
\pgfpathlineto{\pgfqpoint{1.605603in}{1.941452in}}%
\pgfpathlineto{\pgfqpoint{1.605617in}{1.928128in}}%
\pgfpathlineto{\pgfqpoint{1.605630in}{1.941764in}}%
\pgfpathlineto{\pgfqpoint{1.606379in}{1.932484in}}%
\pgfpathlineto{\pgfqpoint{1.606396in}{1.937215in}}%
\pgfpathlineto{\pgfqpoint{1.606970in}{1.896322in}}%
\pgfpathlineto{\pgfqpoint{1.607460in}{1.532032in}}%
\pgfpathlineto{\pgfqpoint{1.607997in}{1.532046in}}%
\pgfpathlineto{\pgfqpoint{1.609369in}{1.532121in}}%
\pgfpathlineto{\pgfqpoint{1.609459in}{1.959885in}}%
\pgfpathlineto{\pgfqpoint{1.610910in}{1.934082in}}%
\pgfpathlineto{\pgfqpoint{1.611047in}{1.939971in}}%
\pgfpathlineto{\pgfqpoint{1.611541in}{1.929683in}}%
\pgfpathlineto{\pgfqpoint{1.611690in}{1.936125in}}%
\pgfpathlineto{\pgfqpoint{1.611838in}{1.937045in}}%
\pgfpathlineto{\pgfqpoint{1.611888in}{1.908117in}}%
\pgfpathlineto{\pgfqpoint{1.612374in}{1.532032in}}%
\pgfpathlineto{\pgfqpoint{1.613193in}{1.532051in}}%
\pgfpathlineto{\pgfqpoint{1.614315in}{1.532126in}}%
\pgfpathlineto{\pgfqpoint{1.614625in}{1.963377in}}%
\pgfpathlineto{\pgfqpoint{1.615630in}{1.929564in}}%
\pgfpathlineto{\pgfqpoint{1.615634in}{1.928729in}}%
\pgfpathlineto{\pgfqpoint{1.615678in}{1.940907in}}%
\pgfpathlineto{\pgfqpoint{1.616369in}{1.933595in}}%
\pgfpathlineto{\pgfqpoint{1.616446in}{1.936334in}}%
\pgfpathlineto{\pgfqpoint{1.616882in}{1.874230in}}%
\pgfpathlineto{\pgfqpoint{1.617360in}{1.532032in}}%
\pgfpathlineto{\pgfqpoint{1.617848in}{1.532043in}}%
\pgfpathlineto{\pgfqpoint{1.619259in}{1.532121in}}%
\pgfpathlineto{\pgfqpoint{1.619579in}{1.959879in}}%
\pgfpathlineto{\pgfqpoint{1.620807in}{1.938701in}}%
\pgfpathlineto{\pgfqpoint{1.620901in}{1.928687in}}%
\pgfpathlineto{\pgfqpoint{1.620839in}{1.941223in}}%
\pgfpathlineto{\pgfqpoint{1.621543in}{1.928795in}}%
\pgfpathlineto{\pgfqpoint{1.621592in}{1.939159in}}%
\pgfpathlineto{\pgfqpoint{1.621840in}{1.857555in}}%
\pgfpathlineto{\pgfqpoint{1.622229in}{1.532019in}}%
\pgfpathlineto{\pgfqpoint{1.622846in}{1.532050in}}%
\pgfpathlineto{\pgfqpoint{1.624197in}{1.532085in}}%
\pgfpathlineto{\pgfqpoint{1.624207in}{1.532146in}}%
\pgfpathlineto{\pgfqpoint{1.624401in}{1.961318in}}%
\pgfpathlineto{\pgfqpoint{1.625377in}{1.932128in}}%
\pgfpathlineto{\pgfqpoint{1.625384in}{1.927304in}}%
\pgfpathlineto{\pgfqpoint{1.625398in}{1.941772in}}%
\pgfpathlineto{\pgfqpoint{1.626138in}{1.933322in}}%
\pgfpathlineto{\pgfqpoint{1.626435in}{1.938674in}}%
\pgfpathlineto{\pgfqpoint{1.626385in}{1.929939in}}%
\pgfpathlineto{\pgfqpoint{1.626682in}{1.935462in}}%
\pgfpathlineto{\pgfqpoint{1.627174in}{1.532026in}}%
\pgfpathlineto{\pgfqpoint{1.627944in}{1.532045in}}%
\pgfpathlineto{\pgfqpoint{1.629151in}{1.532128in}}%
\pgfpathlineto{\pgfqpoint{1.629378in}{1.962904in}}%
\pgfpathlineto{\pgfqpoint{1.630482in}{1.940619in}}%
\pgfpathlineto{\pgfqpoint{1.630495in}{1.928295in}}%
\pgfpathlineto{\pgfqpoint{1.630508in}{1.940677in}}%
\pgfpathlineto{\pgfqpoint{1.631253in}{1.936339in}}%
\pgfpathlineto{\pgfqpoint{1.632160in}{1.532032in}}%
\pgfpathlineto{\pgfqpoint{1.633479in}{1.532062in}}%
\pgfpathlineto{\pgfqpoint{1.634098in}{1.532156in}}%
\pgfpathlineto{\pgfqpoint{1.634414in}{1.963657in}}%
\pgfpathlineto{\pgfqpoint{1.635297in}{1.933459in}}%
\pgfpathlineto{\pgfqpoint{1.635302in}{1.928599in}}%
\pgfpathlineto{\pgfqpoint{1.635348in}{1.940581in}}%
\pgfpathlineto{\pgfqpoint{1.636060in}{1.936548in}}%
\pgfpathlineto{\pgfqpoint{1.637118in}{1.532032in}}%
\pgfpathlineto{\pgfqpoint{1.638446in}{1.532063in}}%
\pgfpathlineto{\pgfqpoint{1.639044in}{1.532152in}}%
\pgfpathlineto{\pgfqpoint{1.639342in}{1.962551in}}%
\pgfpathlineto{\pgfqpoint{1.640251in}{1.932942in}}%
\pgfpathlineto{\pgfqpoint{1.640259in}{1.928783in}}%
\pgfpathlineto{\pgfqpoint{1.640273in}{1.940561in}}%
\pgfpathlineto{\pgfqpoint{1.641018in}{1.934240in}}%
\pgfpathlineto{\pgfqpoint{1.641028in}{1.936067in}}%
\pgfpathlineto{\pgfqpoint{1.641582in}{1.914115in}}%
\pgfpathlineto{\pgfqpoint{1.642059in}{1.532033in}}%
\pgfpathlineto{\pgfqpoint{1.642693in}{1.532047in}}%
\pgfpathlineto{\pgfqpoint{1.643990in}{1.532182in}}%
\pgfpathlineto{\pgfqpoint{1.644308in}{1.965513in}}%
\pgfpathlineto{\pgfqpoint{1.645137in}{1.930703in}}%
\pgfpathlineto{\pgfqpoint{1.645141in}{1.927620in}}%
\pgfpathlineto{\pgfqpoint{1.645158in}{1.941127in}}%
\pgfpathlineto{\pgfqpoint{1.645900in}{1.935188in}}%
\pgfpathlineto{\pgfqpoint{1.647015in}{1.532031in}}%
\pgfpathlineto{\pgfqpoint{1.645928in}{1.936448in}}%
\pgfpathlineto{\pgfqpoint{1.647754in}{1.532051in}}%
\pgfpathlineto{\pgfqpoint{1.648930in}{1.532108in}}%
\pgfpathlineto{\pgfqpoint{1.648972in}{1.715348in}}%
\pgfpathlineto{\pgfqpoint{1.649021in}{1.958920in}}%
\pgfpathlineto{\pgfqpoint{1.649757in}{1.948211in}}%
\pgfpathlineto{\pgfqpoint{1.649822in}{1.949813in}}%
\pgfpathlineto{\pgfqpoint{1.649806in}{1.918789in}}%
\pgfpathlineto{\pgfqpoint{1.650342in}{1.931649in}}%
\pgfpathlineto{\pgfqpoint{1.650352in}{1.926377in}}%
\pgfpathlineto{\pgfqpoint{1.650369in}{1.942688in}}%
\pgfpathlineto{\pgfqpoint{1.651074in}{1.934230in}}%
\pgfpathlineto{\pgfqpoint{1.651272in}{1.940764in}}%
\pgfpathlineto{\pgfqpoint{1.651469in}{1.886967in}}%
\pgfpathlineto{\pgfqpoint{1.651902in}{1.532020in}}%
\pgfpathlineto{\pgfqpoint{1.652648in}{1.532043in}}%
\pgfpathlineto{\pgfqpoint{1.653881in}{1.532163in}}%
\pgfpathlineto{\pgfqpoint{1.653984in}{1.961112in}}%
\pgfpathlineto{\pgfqpoint{1.655019in}{1.938572in}}%
\pgfpathlineto{\pgfqpoint{1.655342in}{1.940525in}}%
\pgfpathlineto{\pgfqpoint{1.655083in}{1.927958in}}%
\pgfpathlineto{\pgfqpoint{1.655688in}{1.938203in}}%
\pgfpathlineto{\pgfqpoint{1.655787in}{1.940261in}}%
\pgfpathlineto{\pgfqpoint{1.656628in}{1.637760in}}%
\pgfpathlineto{\pgfqpoint{1.656880in}{1.532032in}}%
\pgfpathlineto{\pgfqpoint{1.656828in}{1.669673in}}%
\pgfpathlineto{\pgfqpoint{1.657595in}{1.532049in}}%
\pgfpathlineto{\pgfqpoint{1.658824in}{1.532127in}}%
\pgfpathlineto{\pgfqpoint{1.659145in}{1.964654in}}%
\pgfpathlineto{\pgfqpoint{1.660142in}{1.939337in}}%
\pgfpathlineto{\pgfqpoint{1.660152in}{1.928242in}}%
\pgfpathlineto{\pgfqpoint{1.660167in}{1.940210in}}%
\pgfpathlineto{\pgfqpoint{1.660908in}{1.931995in}}%
\pgfpathlineto{\pgfqpoint{1.660928in}{1.936161in}}%
\pgfpathlineto{\pgfqpoint{1.661392in}{1.891417in}}%
\pgfpathlineto{\pgfqpoint{1.661862in}{1.532031in}}%
\pgfpathlineto{\pgfqpoint{1.662403in}{1.532047in}}%
\pgfpathlineto{\pgfqpoint{1.663767in}{1.532107in}}%
\pgfpathlineto{\pgfqpoint{1.663796in}{1.581681in}}%
\pgfpathlineto{\pgfqpoint{1.664069in}{1.961764in}}%
\pgfpathlineto{\pgfqpoint{1.664599in}{1.943597in}}%
\pgfpathlineto{\pgfqpoint{1.664603in}{1.948246in}}%
\pgfpathlineto{\pgfqpoint{1.664619in}{1.920744in}}%
\pgfpathlineto{\pgfqpoint{1.665339in}{1.941155in}}%
\pgfpathlineto{\pgfqpoint{1.665986in}{1.927811in}}%
\pgfpathlineto{\pgfqpoint{1.666085in}{1.930549in}}%
\pgfpathlineto{\pgfqpoint{1.666184in}{1.934688in}}%
\pgfpathlineto{\pgfqpoint{1.666332in}{1.903771in}}%
\pgfpathlineto{\pgfqpoint{1.666755in}{1.532026in}}%
\pgfpathlineto{\pgfqpoint{1.667390in}{1.532040in}}%
\pgfpathlineto{\pgfqpoint{1.668716in}{1.532148in}}%
\pgfpathlineto{\pgfqpoint{1.668808in}{1.959715in}}%
\pgfpathlineto{\pgfqpoint{1.669839in}{1.947903in}}%
\pgfpathlineto{\pgfqpoint{1.670185in}{1.921203in}}%
\pgfpathlineto{\pgfqpoint{1.670631in}{1.921855in}}%
\pgfpathlineto{\pgfqpoint{1.671017in}{1.944670in}}%
\pgfpathlineto{\pgfqpoint{1.671264in}{1.874541in}}%
\pgfpathlineto{\pgfqpoint{1.671683in}{1.532020in}}%
\pgfpathlineto{\pgfqpoint{1.672427in}{1.532043in}}%
\pgfpathlineto{\pgfqpoint{1.673662in}{1.532158in}}%
\pgfpathlineto{\pgfqpoint{1.673754in}{1.959361in}}%
\pgfpathlineto{\pgfqpoint{1.674799in}{1.930388in}}%
\pgfpathlineto{\pgfqpoint{1.674807in}{1.925773in}}%
\pgfpathlineto{\pgfqpoint{1.674821in}{1.942850in}}%
\pgfpathlineto{\pgfqpoint{1.675550in}{1.931878in}}%
\pgfpathlineto{\pgfqpoint{1.675698in}{1.939997in}}%
\pgfpathlineto{\pgfqpoint{1.676094in}{1.927724in}}%
\pgfpathlineto{\pgfqpoint{1.676144in}{1.938193in}}%
\pgfpathlineto{\pgfqpoint{1.676707in}{1.532032in}}%
\pgfpathlineto{\pgfqpoint{1.677640in}{1.532055in}}%
\pgfpathlineto{\pgfqpoint{1.678605in}{1.532121in}}%
\pgfpathlineto{\pgfqpoint{1.678906in}{1.961591in}}%
\pgfpathlineto{\pgfqpoint{1.680152in}{1.939808in}}%
\pgfpathlineto{\pgfqpoint{1.680170in}{1.928657in}}%
\pgfpathlineto{\pgfqpoint{1.680882in}{1.932968in}}%
\pgfpathlineto{\pgfqpoint{1.681031in}{1.937839in}}%
\pgfpathlineto{\pgfqpoint{1.681130in}{1.924027in}}%
\pgfpathlineto{\pgfqpoint{1.681590in}{1.532027in}}%
\pgfpathlineto{\pgfqpoint{1.682569in}{1.532059in}}%
\pgfpathlineto{\pgfqpoint{1.683542in}{1.532085in}}%
\pgfpathlineto{\pgfqpoint{1.683552in}{1.532146in}}%
\pgfpathlineto{\pgfqpoint{1.683647in}{1.960334in}}%
\pgfpathlineto{\pgfqpoint{1.684718in}{1.939485in}}%
\pgfpathlineto{\pgfqpoint{1.684984in}{1.929002in}}%
\pgfpathlineto{\pgfqpoint{1.685478in}{1.938849in}}%
\pgfpathlineto{\pgfqpoint{1.686270in}{1.669447in}}%
\pgfpathlineto{\pgfqpoint{1.686607in}{1.532033in}}%
\pgfpathlineto{\pgfqpoint{1.687146in}{1.532046in}}%
\pgfpathlineto{\pgfqpoint{1.688498in}{1.532142in}}%
\pgfpathlineto{\pgfqpoint{1.688812in}{1.963878in}}%
\pgfpathlineto{\pgfqpoint{1.689746in}{1.937921in}}%
\pgfpathlineto{\pgfqpoint{1.689787in}{1.928706in}}%
\pgfpathlineto{\pgfqpoint{1.689774in}{1.939258in}}%
\pgfpathlineto{\pgfqpoint{1.690518in}{1.934504in}}%
\pgfpathlineto{\pgfqpoint{1.691510in}{1.532030in}}%
\pgfpathlineto{\pgfqpoint{1.690579in}{1.935393in}}%
\pgfpathlineto{\pgfqpoint{1.692734in}{1.532062in}}%
\pgfpathlineto{\pgfqpoint{1.693438in}{1.532101in}}%
\pgfpathlineto{\pgfqpoint{1.693447in}{1.532295in}}%
\pgfpathlineto{\pgfqpoint{1.693542in}{1.959785in}}%
\pgfpathlineto{\pgfqpoint{1.694511in}{1.930235in}}%
\pgfpathlineto{\pgfqpoint{1.694557in}{1.940778in}}%
\pgfpathlineto{\pgfqpoint{1.694571in}{1.927351in}}%
\pgfpathlineto{\pgfqpoint{1.695255in}{1.930872in}}%
\pgfpathlineto{\pgfqpoint{1.695403in}{1.938154in}}%
\pgfpathlineto{\pgfqpoint{1.695799in}{1.929095in}}%
\pgfpathlineto{\pgfqpoint{1.695898in}{1.929651in}}%
\pgfpathlineto{\pgfqpoint{1.696400in}{1.556242in}}%
\pgfpathlineto{\pgfqpoint{1.696410in}{1.532011in}}%
\pgfpathlineto{\pgfqpoint{1.697145in}{1.532039in}}%
\pgfpathlineto{\pgfqpoint{1.698393in}{1.532551in}}%
\pgfpathlineto{\pgfqpoint{1.698497in}{1.960552in}}%
\pgfpathlineto{\pgfqpoint{1.699421in}{1.930263in}}%
\pgfpathlineto{\pgfqpoint{1.699497in}{1.927226in}}%
\pgfpathlineto{\pgfqpoint{1.699480in}{1.940551in}}%
\pgfpathlineto{\pgfqpoint{1.700159in}{1.930136in}}%
\pgfpathlineto{\pgfqpoint{1.700456in}{1.938456in}}%
\pgfpathlineto{\pgfqpoint{1.700852in}{1.928107in}}%
\pgfpathlineto{\pgfqpoint{1.700901in}{1.933786in}}%
\pgfpathlineto{\pgfqpoint{1.701380in}{1.532031in}}%
\pgfpathlineto{\pgfqpoint{1.702176in}{1.532052in}}%
\pgfpathlineto{\pgfqpoint{1.703330in}{1.532104in}}%
\pgfpathlineto{\pgfqpoint{1.703338in}{1.532397in}}%
\pgfpathlineto{\pgfqpoint{1.703444in}{1.960702in}}%
\pgfpathlineto{\pgfqpoint{1.704374in}{1.927860in}}%
\pgfpathlineto{\pgfqpoint{1.704427in}{1.940686in}}%
\pgfpathlineto{\pgfqpoint{1.704412in}{1.926837in}}%
\pgfpathlineto{\pgfqpoint{1.705185in}{1.934578in}}%
\pgfpathlineto{\pgfqpoint{1.705234in}{1.934650in}}%
\pgfpathlineto{\pgfqpoint{1.705432in}{1.938281in}}%
\pgfpathlineto{\pgfqpoint{1.705927in}{1.868662in}}%
\pgfpathlineto{\pgfqpoint{1.706319in}{1.532026in}}%
\pgfpathlineto{\pgfqpoint{1.706906in}{1.532051in}}%
\pgfpathlineto{\pgfqpoint{1.708268in}{1.532083in}}%
\pgfpathlineto{\pgfqpoint{1.708279in}{1.532142in}}%
\pgfpathlineto{\pgfqpoint{1.708372in}{1.958442in}}%
\pgfpathlineto{\pgfqpoint{1.709434in}{1.925275in}}%
\pgfpathlineto{\pgfqpoint{1.709456in}{1.943268in}}%
\pgfpathlineto{\pgfqpoint{1.709476in}{1.924726in}}%
\pgfpathlineto{\pgfqpoint{1.710165in}{1.935609in}}%
\pgfpathlineto{\pgfqpoint{1.710363in}{1.939596in}}%
\pgfpathlineto{\pgfqpoint{1.710907in}{1.808738in}}%
\pgfpathlineto{\pgfqpoint{1.711300in}{1.532031in}}%
\pgfpathlineto{\pgfqpoint{1.711840in}{1.532046in}}%
\pgfpathlineto{\pgfqpoint{1.713222in}{1.532109in}}%
\pgfpathlineto{\pgfqpoint{1.713263in}{1.720297in}}%
\pgfpathlineto{\pgfqpoint{1.713325in}{1.960449in}}%
\pgfpathlineto{\pgfqpoint{1.714049in}{1.941816in}}%
\pgfpathlineto{\pgfqpoint{1.714106in}{1.923026in}}%
\pgfpathlineto{\pgfqpoint{1.714121in}{1.944137in}}%
\pgfpathlineto{\pgfqpoint{1.714850in}{1.935852in}}%
\pgfpathlineto{\pgfqpoint{1.715147in}{1.926853in}}%
\pgfpathlineto{\pgfqpoint{1.715098in}{1.940749in}}%
\pgfpathlineto{\pgfqpoint{1.715493in}{1.927086in}}%
\pgfpathlineto{\pgfqpoint{1.715543in}{1.940085in}}%
\pgfpathlineto{\pgfqpoint{1.715840in}{1.827823in}}%
\pgfpathlineto{\pgfqpoint{1.716192in}{1.532019in}}%
\pgfpathlineto{\pgfqpoint{1.716832in}{1.532051in}}%
\pgfpathlineto{\pgfqpoint{1.718168in}{1.532119in}}%
\pgfpathlineto{\pgfqpoint{1.718208in}{1.724480in}}%
\pgfpathlineto{\pgfqpoint{1.718488in}{1.962096in}}%
\pgfpathlineto{\pgfqpoint{1.718992in}{1.947237in}}%
\pgfpathlineto{\pgfqpoint{1.719007in}{1.920370in}}%
\pgfpathlineto{\pgfqpoint{1.719767in}{1.939860in}}%
\pgfpathlineto{\pgfqpoint{1.720364in}{1.927632in}}%
\pgfpathlineto{\pgfqpoint{1.720513in}{1.935937in}}%
\pgfpathlineto{\pgfqpoint{1.720711in}{1.901443in}}%
\pgfpathlineto{\pgfqpoint{1.721188in}{1.532032in}}%
\pgfpathlineto{\pgfqpoint{1.721822in}{1.532047in}}%
\pgfpathlineto{\pgfqpoint{1.723115in}{1.532127in}}%
\pgfpathlineto{\pgfqpoint{1.723205in}{1.958733in}}%
\pgfpathlineto{\pgfqpoint{1.724559in}{1.937829in}}%
\pgfpathlineto{\pgfqpoint{1.724611in}{1.928987in}}%
\pgfpathlineto{\pgfqpoint{1.724765in}{1.938215in}}%
\pgfpathlineto{\pgfqpoint{1.725309in}{1.937057in}}%
\pgfpathlineto{\pgfqpoint{1.725754in}{1.806630in}}%
\pgfpathlineto{\pgfqpoint{1.726100in}{1.532027in}}%
\pgfpathlineto{\pgfqpoint{1.726685in}{1.532050in}}%
\pgfpathlineto{\pgfqpoint{1.728051in}{1.532084in}}%
\pgfpathlineto{\pgfqpoint{1.728063in}{1.532174in}}%
\pgfpathlineto{\pgfqpoint{1.728165in}{1.960497in}}%
\pgfpathlineto{\pgfqpoint{1.729145in}{1.928595in}}%
\pgfpathlineto{\pgfqpoint{1.729187in}{1.939966in}}%
\pgfpathlineto{\pgfqpoint{1.729236in}{1.927051in}}%
\pgfpathlineto{\pgfqpoint{1.729879in}{1.933120in}}%
\pgfpathlineto{\pgfqpoint{1.730127in}{1.927255in}}%
\pgfpathlineto{\pgfqpoint{1.730077in}{1.939572in}}%
\pgfpathlineto{\pgfqpoint{1.730473in}{1.927645in}}%
\pgfpathlineto{\pgfqpoint{1.730522in}{1.938534in}}%
\pgfpathlineto{\pgfqpoint{1.730671in}{1.838171in}}%
\pgfpathlineto{\pgfqpoint{1.731054in}{1.532031in}}%
\pgfpathlineto{\pgfqpoint{1.731606in}{1.532042in}}%
\pgfpathlineto{\pgfqpoint{1.733011in}{1.532330in}}%
\pgfpathlineto{\pgfqpoint{1.733120in}{1.958169in}}%
\pgfpathlineto{\pgfqpoint{1.734065in}{1.932065in}}%
\pgfpathlineto{\pgfqpoint{1.734073in}{1.925767in}}%
\pgfpathlineto{\pgfqpoint{1.734123in}{1.940908in}}%
\pgfpathlineto{\pgfqpoint{1.734832in}{1.930415in}}%
\pgfpathlineto{\pgfqpoint{1.735129in}{1.937650in}}%
\pgfpathlineto{\pgfqpoint{1.735179in}{1.928941in}}%
\pgfpathlineto{\pgfqpoint{1.735475in}{1.935315in}}%
\pgfpathlineto{\pgfqpoint{1.736068in}{1.532033in}}%
\pgfpathlineto{\pgfqpoint{1.737054in}{1.532055in}}%
\pgfpathlineto{\pgfqpoint{1.737953in}{1.532154in}}%
\pgfpathlineto{\pgfqpoint{1.738278in}{1.965378in}}%
\pgfpathlineto{\pgfqpoint{1.739149in}{1.926974in}}%
\pgfpathlineto{\pgfqpoint{1.739166in}{1.940564in}}%
\pgfpathlineto{\pgfqpoint{1.739179in}{1.926946in}}%
\pgfpathlineto{\pgfqpoint{1.740063in}{1.932768in}}%
\pgfpathlineto{\pgfqpoint{1.740089in}{1.935097in}}%
\pgfpathlineto{\pgfqpoint{1.740593in}{1.787274in}}%
\pgfpathlineto{\pgfqpoint{1.740990in}{1.532031in}}%
\pgfpathlineto{\pgfqpoint{1.741531in}{1.532047in}}%
\pgfpathlineto{\pgfqpoint{1.742894in}{1.532106in}}%
\pgfpathlineto{\pgfqpoint{1.742923in}{1.580723in}}%
\pgfpathlineto{\pgfqpoint{1.743198in}{1.961520in}}%
\pgfpathlineto{\pgfqpoint{1.743723in}{1.923239in}}%
\pgfpathlineto{\pgfqpoint{1.743727in}{1.919672in}}%
\pgfpathlineto{\pgfqpoint{1.743744in}{1.945923in}}%
\pgfpathlineto{\pgfqpoint{1.744466in}{1.939899in}}%
\pgfpathlineto{\pgfqpoint{1.745152in}{1.926373in}}%
\pgfpathlineto{\pgfqpoint{1.744657in}{1.940075in}}%
\pgfpathlineto{\pgfqpoint{1.745251in}{1.928125in}}%
\pgfpathlineto{\pgfqpoint{1.745300in}{1.935474in}}%
\pgfpathlineto{\pgfqpoint{1.745498in}{1.853812in}}%
\pgfpathlineto{\pgfqpoint{1.745896in}{1.532030in}}%
\pgfpathlineto{\pgfqpoint{1.746510in}{1.532047in}}%
\pgfpathlineto{\pgfqpoint{1.747838in}{1.532100in}}%
\pgfpathlineto{\pgfqpoint{1.747847in}{1.532305in}}%
\pgfpathlineto{\pgfqpoint{1.747943in}{1.959264in}}%
\pgfpathlineto{\pgfqpoint{1.748910in}{1.939161in}}%
\pgfpathlineto{\pgfqpoint{1.748916in}{1.941400in}}%
\pgfpathlineto{\pgfqpoint{1.748929in}{1.925716in}}%
\pgfpathlineto{\pgfqpoint{1.749670in}{1.937553in}}%
\pgfpathlineto{\pgfqpoint{1.750017in}{1.937750in}}%
\pgfpathlineto{\pgfqpoint{1.750511in}{1.737723in}}%
\pgfpathlineto{\pgfqpoint{1.750843in}{1.532033in}}%
\pgfpathlineto{\pgfqpoint{1.751410in}{1.532044in}}%
\pgfpathlineto{\pgfqpoint{1.752788in}{1.532132in}}%
\pgfpathlineto{\pgfqpoint{1.753106in}{1.964301in}}%
\pgfpathlineto{\pgfqpoint{1.754069in}{1.938539in}}%
\pgfpathlineto{\pgfqpoint{1.754082in}{1.927748in}}%
\pgfpathlineto{\pgfqpoint{1.754126in}{1.938687in}}%
\pgfpathlineto{\pgfqpoint{1.754843in}{1.932867in}}%
\pgfpathlineto{\pgfqpoint{1.755773in}{1.532026in}}%
\pgfpathlineto{\pgfqpoint{1.754900in}{1.934740in}}%
\pgfpathlineto{\pgfqpoint{1.757201in}{1.532057in}}%
\pgfpathlineto{\pgfqpoint{1.757734in}{1.532143in}}%
\pgfpathlineto{\pgfqpoint{1.757827in}{1.959258in}}%
\pgfpathlineto{\pgfqpoint{1.758858in}{1.934985in}}%
\pgfpathlineto{\pgfqpoint{1.759204in}{1.935573in}}%
\pgfpathlineto{\pgfqpoint{1.759600in}{1.931201in}}%
\pgfpathlineto{\pgfqpoint{1.759649in}{1.935381in}}%
\pgfpathlineto{\pgfqpoint{1.760144in}{1.930826in}}%
\pgfpathlineto{\pgfqpoint{1.760243in}{1.931258in}}%
\pgfpathlineto{\pgfqpoint{1.760729in}{1.532031in}}%
\pgfpathlineto{\pgfqpoint{1.761585in}{1.532049in}}%
\pgfpathlineto{\pgfqpoint{1.762683in}{1.532303in}}%
\pgfpathlineto{\pgfqpoint{1.762843in}{1.958084in}}%
\pgfpathlineto{\pgfqpoint{1.763737in}{1.938823in}}%
\pgfpathlineto{\pgfqpoint{1.763751in}{1.925686in}}%
\pgfpathlineto{\pgfqpoint{1.763800in}{1.940569in}}%
\pgfpathlineto{\pgfqpoint{1.764477in}{1.935013in}}%
\pgfpathlineto{\pgfqpoint{1.764673in}{1.937238in}}%
\pgfpathlineto{\pgfqpoint{1.765316in}{1.798170in}}%
\pgfpathlineto{\pgfqpoint{1.765697in}{1.532032in}}%
\pgfpathlineto{\pgfqpoint{1.766234in}{1.532045in}}%
\pgfpathlineto{\pgfqpoint{1.767626in}{1.532172in}}%
\pgfpathlineto{\pgfqpoint{1.767946in}{1.967774in}}%
\pgfpathlineto{\pgfqpoint{1.768783in}{1.925267in}}%
\pgfpathlineto{\pgfqpoint{1.768788in}{1.923864in}}%
\pgfpathlineto{\pgfqpoint{1.768803in}{1.942216in}}%
\pgfpathlineto{\pgfqpoint{1.769519in}{1.934775in}}%
\pgfpathlineto{\pgfqpoint{1.769525in}{1.936173in}}%
\pgfpathlineto{\pgfqpoint{1.770162in}{1.917526in}}%
\pgfpathlineto{\pgfqpoint{1.770643in}{1.532028in}}%
\pgfpathlineto{\pgfqpoint{1.771326in}{1.532043in}}%
\pgfpathlineto{\pgfqpoint{1.772570in}{1.532140in}}%
\pgfpathlineto{\pgfqpoint{1.772747in}{1.961819in}}%
\pgfpathlineto{\pgfqpoint{1.773762in}{1.939742in}}%
\pgfpathlineto{\pgfqpoint{1.773779in}{1.926758in}}%
\pgfpathlineto{\pgfqpoint{1.774491in}{1.929282in}}%
\pgfpathlineto{\pgfqpoint{1.774788in}{1.937296in}}%
\pgfpathlineto{\pgfqpoint{1.774837in}{1.928073in}}%
\pgfpathlineto{\pgfqpoint{1.775035in}{1.931577in}}%
\pgfpathlineto{\pgfqpoint{1.775596in}{1.532028in}}%
\pgfpathlineto{\pgfqpoint{1.776529in}{1.532057in}}%
\pgfpathlineto{\pgfqpoint{1.777516in}{1.532140in}}%
\pgfpathlineto{\pgfqpoint{1.777723in}{1.962866in}}%
\pgfpathlineto{\pgfqpoint{1.778709in}{1.939662in}}%
\pgfpathlineto{\pgfqpoint{1.778728in}{1.926609in}}%
\pgfpathlineto{\pgfqpoint{1.778749in}{1.939972in}}%
\pgfpathlineto{\pgfqpoint{1.779489in}{1.929990in}}%
\pgfpathlineto{\pgfqpoint{1.779785in}{1.937183in}}%
\pgfpathlineto{\pgfqpoint{1.779983in}{1.924166in}}%
\pgfpathlineto{\pgfqpoint{1.780033in}{1.928976in}}%
\pgfpathlineto{\pgfqpoint{1.780509in}{1.532031in}}%
\pgfpathlineto{\pgfqpoint{1.781310in}{1.532052in}}%
\pgfpathlineto{\pgfqpoint{1.782457in}{1.532102in}}%
\pgfpathlineto{\pgfqpoint{1.782465in}{1.532350in}}%
\pgfpathlineto{\pgfqpoint{1.782557in}{1.958394in}}%
\pgfpathlineto{\pgfqpoint{1.783508in}{1.939425in}}%
\pgfpathlineto{\pgfqpoint{1.783559in}{1.926781in}}%
\pgfpathlineto{\pgfqpoint{1.784241in}{1.935923in}}%
\pgfpathlineto{\pgfqpoint{1.784587in}{1.937334in}}%
\pgfpathlineto{\pgfqpoint{1.785081in}{1.837509in}}%
\pgfpathlineto{\pgfqpoint{1.785445in}{1.532027in}}%
\pgfpathlineto{\pgfqpoint{1.786030in}{1.532050in}}%
\pgfpathlineto{\pgfqpoint{1.787396in}{1.532084in}}%
\pgfpathlineto{\pgfqpoint{1.787407in}{1.532146in}}%
\pgfpathlineto{\pgfqpoint{1.787500in}{1.958317in}}%
\pgfpathlineto{\pgfqpoint{1.788554in}{1.927943in}}%
\pgfpathlineto{\pgfqpoint{1.788603in}{1.938750in}}%
\pgfpathlineto{\pgfqpoint{1.788653in}{1.927615in}}%
\pgfpathlineto{\pgfqpoint{1.789296in}{1.933698in}}%
\pgfpathlineto{\pgfqpoint{1.789494in}{1.938229in}}%
\pgfpathlineto{\pgfqpoint{1.789988in}{1.879542in}}%
\pgfpathlineto{\pgfqpoint{1.790423in}{1.532030in}}%
\pgfpathlineto{\pgfqpoint{1.791008in}{1.532043in}}%
\pgfpathlineto{\pgfqpoint{1.792349in}{1.532112in}}%
\pgfpathlineto{\pgfqpoint{1.792391in}{1.735713in}}%
\pgfpathlineto{\pgfqpoint{1.792654in}{1.960152in}}%
\pgfpathlineto{\pgfqpoint{1.793172in}{1.921193in}}%
\pgfpathlineto{\pgfqpoint{1.793237in}{1.919396in}}%
\pgfpathlineto{\pgfqpoint{1.793221in}{1.946564in}}%
\pgfpathlineto{\pgfqpoint{1.793896in}{1.927068in}}%
\pgfpathlineto{\pgfqpoint{1.794017in}{1.939614in}}%
\pgfpathlineto{\pgfqpoint{1.794377in}{1.926520in}}%
\pgfpathlineto{\pgfqpoint{1.794674in}{1.935570in}}%
\pgfpathlineto{\pgfqpoint{1.794773in}{1.938078in}}%
\pgfpathlineto{\pgfqpoint{1.795020in}{1.770346in}}%
\pgfpathlineto{\pgfqpoint{1.795380in}{1.532028in}}%
\pgfpathlineto{\pgfqpoint{1.795917in}{1.532049in}}%
\pgfpathlineto{\pgfqpoint{1.797288in}{1.532085in}}%
\pgfpathlineto{\pgfqpoint{1.797298in}{1.532146in}}%
\pgfpathlineto{\pgfqpoint{1.797388in}{1.957938in}}%
\pgfpathlineto{\pgfqpoint{1.798477in}{1.939890in}}%
\pgfpathlineto{\pgfqpoint{1.798497in}{1.926209in}}%
\pgfpathlineto{\pgfqpoint{1.799248in}{1.933569in}}%
\pgfpathlineto{\pgfqpoint{1.799446in}{1.937511in}}%
\pgfpathlineto{\pgfqpoint{1.799891in}{1.853061in}}%
\pgfpathlineto{\pgfqpoint{1.800308in}{1.532032in}}%
\pgfpathlineto{\pgfqpoint{1.800887in}{1.532044in}}%
\pgfpathlineto{\pgfqpoint{1.802245in}{1.532174in}}%
\pgfpathlineto{\pgfqpoint{1.802559in}{1.967064in}}%
\pgfpathlineto{\pgfqpoint{1.803396in}{1.934248in}}%
\pgfpathlineto{\pgfqpoint{1.803406in}{1.923318in}}%
\pgfpathlineto{\pgfqpoint{1.803419in}{1.942414in}}%
\pgfpathlineto{\pgfqpoint{1.804168in}{1.934662in}}%
\pgfpathlineto{\pgfqpoint{1.805227in}{1.532026in}}%
\pgfpathlineto{\pgfqpoint{1.804229in}{1.935703in}}%
\pgfpathlineto{\pgfqpoint{1.805764in}{1.532037in}}%
\pgfpathlineto{\pgfqpoint{1.807189in}{1.532152in}}%
\pgfpathlineto{\pgfqpoint{1.807281in}{1.958979in}}%
\pgfpathlineto{\pgfqpoint{1.808306in}{1.934733in}}%
\pgfpathlineto{\pgfqpoint{1.808603in}{1.927002in}}%
\pgfpathlineto{\pgfqpoint{1.808554in}{1.939117in}}%
\pgfpathlineto{\pgfqpoint{1.808949in}{1.927411in}}%
\pgfpathlineto{\pgfqpoint{1.808999in}{1.938892in}}%
\pgfpathlineto{\pgfqpoint{1.809493in}{1.926991in}}%
\pgfpathlineto{\pgfqpoint{1.809691in}{1.932127in}}%
\pgfpathlineto{\pgfqpoint{1.810220in}{1.532031in}}%
\pgfpathlineto{\pgfqpoint{1.810961in}{1.532051in}}%
\pgfpathlineto{\pgfqpoint{1.812132in}{1.532123in}}%
\pgfpathlineto{\pgfqpoint{1.812445in}{1.962055in}}%
\pgfpathlineto{\pgfqpoint{1.813475in}{1.932866in}}%
\pgfpathlineto{\pgfqpoint{1.813484in}{1.926917in}}%
\pgfpathlineto{\pgfqpoint{1.813497in}{1.939146in}}%
\pgfpathlineto{\pgfqpoint{1.814244in}{1.931677in}}%
\pgfpathlineto{\pgfqpoint{1.814264in}{1.934363in}}%
\pgfpathlineto{\pgfqpoint{1.814769in}{1.800278in}}%
\pgfpathlineto{\pgfqpoint{1.815132in}{1.532031in}}%
\pgfpathlineto{\pgfqpoint{1.815697in}{1.532042in}}%
\pgfpathlineto{\pgfqpoint{1.817083in}{1.532307in}}%
\pgfpathlineto{\pgfqpoint{1.817179in}{1.958689in}}%
\pgfpathlineto{\pgfqpoint{1.818133in}{1.933640in}}%
\pgfpathlineto{\pgfqpoint{1.818176in}{1.925017in}}%
\pgfpathlineto{\pgfqpoint{1.818190in}{1.940416in}}%
\pgfpathlineto{\pgfqpoint{1.818907in}{1.928547in}}%
\pgfpathlineto{\pgfqpoint{1.819302in}{1.937294in}}%
\pgfpathlineto{\pgfqpoint{1.819253in}{1.927599in}}%
\pgfpathlineto{\pgfqpoint{1.819549in}{1.934143in}}%
\pgfpathlineto{\pgfqpoint{1.820090in}{1.532030in}}%
\pgfpathlineto{\pgfqpoint{1.821066in}{1.532051in}}%
\pgfpathlineto{\pgfqpoint{1.822029in}{1.532293in}}%
\pgfpathlineto{\pgfqpoint{1.822125in}{1.958031in}}%
\pgfpathlineto{\pgfqpoint{1.823083in}{1.927586in}}%
\pgfpathlineto{\pgfqpoint{1.823090in}{1.925548in}}%
\pgfpathlineto{\pgfqpoint{1.823102in}{1.940711in}}%
\pgfpathlineto{\pgfqpoint{1.823825in}{1.933498in}}%
\pgfpathlineto{\pgfqpoint{1.824009in}{1.936891in}}%
\pgfpathlineto{\pgfqpoint{1.824652in}{1.805229in}}%
\pgfpathlineto{\pgfqpoint{1.825086in}{1.532033in}}%
\pgfpathlineto{\pgfqpoint{1.825578in}{1.532044in}}%
\pgfpathlineto{\pgfqpoint{1.826970in}{1.532141in}}%
\pgfpathlineto{\pgfqpoint{1.827288in}{1.962852in}}%
\pgfpathlineto{\pgfqpoint{1.828218in}{1.938269in}}%
\pgfpathlineto{\pgfqpoint{1.828236in}{1.927265in}}%
\pgfpathlineto{\pgfqpoint{1.828249in}{1.938456in}}%
\pgfpathlineto{\pgfqpoint{1.829074in}{1.931534in}}%
\pgfpathlineto{\pgfqpoint{1.829119in}{1.933977in}}%
\pgfpathlineto{\pgfqpoint{1.829551in}{1.884287in}}%
\pgfpathlineto{\pgfqpoint{1.829955in}{1.532026in}}%
\pgfpathlineto{\pgfqpoint{1.830590in}{1.532040in}}%
\pgfpathlineto{\pgfqpoint{1.831916in}{1.532152in}}%
\pgfpathlineto{\pgfqpoint{1.832008in}{1.958256in}}%
\pgfpathlineto{\pgfqpoint{1.833041in}{1.922049in}}%
\pgfpathlineto{\pgfqpoint{1.833584in}{1.946407in}}%
\pgfpathlineto{\pgfqpoint{1.833221in}{1.918691in}}%
\pgfpathlineto{\pgfqpoint{1.833819in}{1.936377in}}%
\pgfpathlineto{\pgfqpoint{1.834900in}{1.532026in}}%
\pgfpathlineto{\pgfqpoint{1.834008in}{1.944340in}}%
\pgfpathlineto{\pgfqpoint{1.835882in}{1.532059in}}%
\pgfpathlineto{\pgfqpoint{1.836850in}{1.532083in}}%
\pgfpathlineto{\pgfqpoint{1.836862in}{1.532149in}}%
\pgfpathlineto{\pgfqpoint{1.836953in}{1.957518in}}%
\pgfpathlineto{\pgfqpoint{1.838005in}{1.939718in}}%
\pgfpathlineto{\pgfqpoint{1.838087in}{1.924019in}}%
\pgfpathlineto{\pgfqpoint{1.838034in}{1.941888in}}%
\pgfpathlineto{\pgfqpoint{1.838763in}{1.934458in}}%
\pgfpathlineto{\pgfqpoint{1.839011in}{1.938395in}}%
\pgfpathlineto{\pgfqpoint{1.839406in}{1.900521in}}%
\pgfpathlineto{\pgfqpoint{1.839854in}{1.532031in}}%
\pgfpathlineto{\pgfqpoint{1.840546in}{1.532049in}}%
\pgfpathlineto{\pgfqpoint{1.841805in}{1.532120in}}%
\pgfpathlineto{\pgfqpoint{1.842112in}{1.961512in}}%
\pgfpathlineto{\pgfqpoint{1.843173in}{1.931212in}}%
\pgfpathlineto{\pgfqpoint{1.843181in}{1.926685in}}%
\pgfpathlineto{\pgfqpoint{1.843224in}{1.938536in}}%
\pgfpathlineto{\pgfqpoint{1.843944in}{1.931553in}}%
\pgfpathlineto{\pgfqpoint{1.843956in}{1.934063in}}%
\pgfpathlineto{\pgfqpoint{1.844420in}{1.840275in}}%
\pgfpathlineto{\pgfqpoint{1.844791in}{1.532026in}}%
\pgfpathlineto{\pgfqpoint{1.845378in}{1.532051in}}%
\pgfpathlineto{\pgfqpoint{1.846741in}{1.532083in}}%
\pgfpathlineto{\pgfqpoint{1.846752in}{1.532146in}}%
\pgfpathlineto{\pgfqpoint{1.846844in}{1.958136in}}%
\pgfpathlineto{\pgfqpoint{1.847906in}{1.944844in}}%
\pgfpathlineto{\pgfqpoint{1.847956in}{1.920476in}}%
\pgfpathlineto{\pgfqpoint{1.848641in}{1.927414in}}%
\pgfpathlineto{\pgfqpoint{1.848790in}{1.944057in}}%
\pgfpathlineto{\pgfqpoint{1.849185in}{1.921419in}}%
\pgfpathlineto{\pgfqpoint{1.849284in}{1.921918in}}%
\pgfpathlineto{\pgfqpoint{1.849737in}{1.532026in}}%
\pgfpathlineto{\pgfqpoint{1.850523in}{1.532055in}}%
\pgfpathlineto{\pgfqpoint{1.851686in}{1.532082in}}%
\pgfpathlineto{\pgfqpoint{1.851697in}{1.532133in}}%
\pgfpathlineto{\pgfqpoint{1.852036in}{1.967425in}}%
\pgfpathlineto{\pgfqpoint{1.852943in}{1.936917in}}%
\pgfpathlineto{\pgfqpoint{1.852953in}{1.922700in}}%
\pgfpathlineto{\pgfqpoint{1.852968in}{1.942294in}}%
\pgfpathlineto{\pgfqpoint{1.853713in}{1.930388in}}%
\pgfpathlineto{\pgfqpoint{1.853740in}{1.935518in}}%
\pgfpathlineto{\pgfqpoint{1.854368in}{1.750693in}}%
\pgfpathlineto{\pgfqpoint{1.854711in}{1.532029in}}%
\pgfpathlineto{\pgfqpoint{1.855292in}{1.532041in}}%
\pgfpathlineto{\pgfqpoint{1.856644in}{1.532153in}}%
\pgfpathlineto{\pgfqpoint{1.856738in}{1.959029in}}%
\pgfpathlineto{\pgfqpoint{1.857796in}{1.930508in}}%
\pgfpathlineto{\pgfqpoint{1.857877in}{1.926701in}}%
\pgfpathlineto{\pgfqpoint{1.858135in}{1.938321in}}%
\pgfpathlineto{\pgfqpoint{1.858531in}{1.927005in}}%
\pgfpathlineto{\pgfqpoint{1.858580in}{1.938004in}}%
\pgfpathlineto{\pgfqpoint{1.858976in}{1.926724in}}%
\pgfpathlineto{\pgfqpoint{1.859174in}{1.928686in}}%
\pgfpathlineto{\pgfqpoint{1.859671in}{1.532027in}}%
\pgfpathlineto{\pgfqpoint{1.860406in}{1.532054in}}%
\pgfpathlineto{\pgfqpoint{1.861578in}{1.532085in}}%
\pgfpathlineto{\pgfqpoint{1.861589in}{1.532146in}}%
\pgfpathlineto{\pgfqpoint{1.861679in}{1.956769in}}%
\pgfpathlineto{\pgfqpoint{1.862757in}{1.934997in}}%
\pgfpathlineto{\pgfqpoint{1.862766in}{1.925883in}}%
\pgfpathlineto{\pgfqpoint{1.862787in}{1.939567in}}%
\pgfpathlineto{\pgfqpoint{1.863517in}{1.933976in}}%
\pgfpathlineto{\pgfqpoint{1.863764in}{1.936778in}}%
\pgfpathlineto{\pgfqpoint{1.864160in}{1.899230in}}%
\pgfpathlineto{\pgfqpoint{1.864616in}{1.532028in}}%
\pgfpathlineto{\pgfqpoint{1.865203in}{1.532041in}}%
\pgfpathlineto{\pgfqpoint{1.866534in}{1.532146in}}%
\pgfpathlineto{\pgfqpoint{1.866625in}{1.957085in}}%
\pgfpathlineto{\pgfqpoint{1.867702in}{1.925556in}}%
\pgfpathlineto{\pgfqpoint{1.867719in}{1.939280in}}%
\pgfpathlineto{\pgfqpoint{1.868490in}{1.935625in}}%
\pgfpathlineto{\pgfqpoint{1.868588in}{1.937136in}}%
\pgfpathlineto{\pgfqpoint{1.869231in}{1.724754in}}%
\pgfpathlineto{\pgfqpoint{1.869558in}{1.532028in}}%
\pgfpathlineto{\pgfqpoint{1.870145in}{1.532041in}}%
\pgfpathlineto{\pgfqpoint{1.871480in}{1.532161in}}%
\pgfpathlineto{\pgfqpoint{1.871585in}{1.959110in}}%
\pgfpathlineto{\pgfqpoint{1.872623in}{1.940721in}}%
\pgfpathlineto{\pgfqpoint{1.872640in}{1.924421in}}%
\pgfpathlineto{\pgfqpoint{1.873379in}{1.929659in}}%
\pgfpathlineto{\pgfqpoint{1.873676in}{1.938004in}}%
\pgfpathlineto{\pgfqpoint{1.873626in}{1.926617in}}%
\pgfpathlineto{\pgfqpoint{1.873923in}{1.933207in}}%
\pgfpathlineto{\pgfqpoint{1.874463in}{1.532027in}}%
\pgfpathlineto{\pgfqpoint{1.875641in}{1.532063in}}%
\pgfpathlineto{\pgfqpoint{1.876414in}{1.532084in}}%
\pgfpathlineto{\pgfqpoint{1.876425in}{1.532146in}}%
\pgfpathlineto{\pgfqpoint{1.876519in}{1.957546in}}%
\pgfpathlineto{\pgfqpoint{1.877555in}{1.926990in}}%
\pgfpathlineto{\pgfqpoint{1.877604in}{1.938108in}}%
\pgfpathlineto{\pgfqpoint{1.878296in}{1.934043in}}%
\pgfpathlineto{\pgfqpoint{1.878494in}{1.937551in}}%
\pgfpathlineto{\pgfqpoint{1.879038in}{1.841488in}}%
\pgfpathlineto{\pgfqpoint{1.879409in}{1.532026in}}%
\pgfpathlineto{\pgfqpoint{1.879996in}{1.532051in}}%
\pgfpathlineto{\pgfqpoint{1.881359in}{1.532083in}}%
\pgfpathlineto{\pgfqpoint{1.881371in}{1.532147in}}%
\pgfpathlineto{\pgfqpoint{1.881462in}{1.957323in}}%
\pgfpathlineto{\pgfqpoint{1.882523in}{1.937303in}}%
\pgfpathlineto{\pgfqpoint{1.882610in}{1.923704in}}%
\pgfpathlineto{\pgfqpoint{1.882590in}{1.941635in}}%
\pgfpathlineto{\pgfqpoint{1.883284in}{1.937061in}}%
\pgfpathlineto{\pgfqpoint{1.883630in}{1.937903in}}%
\pgfpathlineto{\pgfqpoint{1.884075in}{1.720962in}}%
\pgfpathlineto{\pgfqpoint{1.884376in}{1.532030in}}%
\pgfpathlineto{\pgfqpoint{1.885000in}{1.532048in}}%
\pgfpathlineto{\pgfqpoint{1.886311in}{1.532100in}}%
\pgfpathlineto{\pgfqpoint{1.886320in}{1.532306in}}%
\pgfpathlineto{\pgfqpoint{1.886513in}{1.958575in}}%
\pgfpathlineto{\pgfqpoint{1.887372in}{1.932707in}}%
\pgfpathlineto{\pgfqpoint{1.887447in}{1.924295in}}%
\pgfpathlineto{\pgfqpoint{1.887398in}{1.940834in}}%
\pgfpathlineto{\pgfqpoint{1.888137in}{1.937284in}}%
\pgfpathlineto{\pgfqpoint{1.889272in}{1.563403in}}%
\pgfpathlineto{\pgfqpoint{1.889315in}{1.532030in}}%
\pgfpathlineto{\pgfqpoint{1.890031in}{1.532050in}}%
\pgfpathlineto{\pgfqpoint{1.891256in}{1.532100in}}%
\pgfpathlineto{\pgfqpoint{1.891265in}{1.532307in}}%
\pgfpathlineto{\pgfqpoint{1.891458in}{1.958437in}}%
\pgfpathlineto{\pgfqpoint{1.892318in}{1.936745in}}%
\pgfpathlineto{\pgfqpoint{1.892362in}{1.923744in}}%
\pgfpathlineto{\pgfqpoint{1.892413in}{1.940486in}}%
\pgfpathlineto{\pgfqpoint{1.893047in}{1.937222in}}%
\pgfpathlineto{\pgfqpoint{1.894223in}{1.542167in}}%
\pgfpathlineto{\pgfqpoint{1.894273in}{1.532030in}}%
\pgfpathlineto{\pgfqpoint{1.894952in}{1.532045in}}%
\pgfpathlineto{\pgfqpoint{1.896204in}{1.532110in}}%
\pgfpathlineto{\pgfqpoint{1.896246in}{1.731544in}}%
\pgfpathlineto{\pgfqpoint{1.896307in}{1.958385in}}%
\pgfpathlineto{\pgfqpoint{1.897028in}{1.917343in}}%
\pgfpathlineto{\pgfqpoint{1.897044in}{1.946909in}}%
\pgfpathlineto{\pgfqpoint{1.897806in}{1.926050in}}%
\pgfpathlineto{\pgfqpoint{1.898043in}{1.939090in}}%
\pgfpathlineto{\pgfqpoint{1.898439in}{1.925400in}}%
\pgfpathlineto{\pgfqpoint{1.898587in}{1.935728in}}%
\pgfpathlineto{\pgfqpoint{1.898835in}{1.825692in}}%
\pgfpathlineto{\pgfqpoint{1.899191in}{1.532026in}}%
\pgfpathlineto{\pgfqpoint{1.899778in}{1.532051in}}%
\pgfpathlineto{\pgfqpoint{1.901141in}{1.532083in}}%
\pgfpathlineto{\pgfqpoint{1.901152in}{1.532145in}}%
\pgfpathlineto{\pgfqpoint{1.901244in}{1.957052in}}%
\pgfpathlineto{\pgfqpoint{1.902300in}{1.932028in}}%
\pgfpathlineto{\pgfqpoint{1.902306in}{1.923017in}}%
\pgfpathlineto{\pgfqpoint{1.902400in}{1.941051in}}%
\pgfpathlineto{\pgfqpoint{1.903060in}{1.926433in}}%
\pgfpathlineto{\pgfqpoint{1.903110in}{1.938052in}}%
\pgfpathlineto{\pgfqpoint{1.903703in}{1.908455in}}%
\pgfpathlineto{\pgfqpoint{1.904136in}{1.532027in}}%
\pgfpathlineto{\pgfqpoint{1.904869in}{1.532042in}}%
\pgfpathlineto{\pgfqpoint{1.906098in}{1.532146in}}%
\pgfpathlineto{\pgfqpoint{1.906191in}{1.957464in}}%
\pgfpathlineto{\pgfqpoint{1.907212in}{1.937863in}}%
\pgfpathlineto{\pgfqpoint{1.907261in}{1.926183in}}%
\pgfpathlineto{\pgfqpoint{1.907310in}{1.938159in}}%
\pgfpathlineto{\pgfqpoint{1.907953in}{1.931330in}}%
\pgfpathlineto{\pgfqpoint{1.908102in}{1.937570in}}%
\pgfpathlineto{\pgfqpoint{1.908596in}{1.925432in}}%
\pgfpathlineto{\pgfqpoint{1.909115in}{1.532030in}}%
\pgfpathlineto{\pgfqpoint{1.910493in}{1.532061in}}%
\pgfpathlineto{\pgfqpoint{1.911040in}{1.532111in}}%
\pgfpathlineto{\pgfqpoint{1.911082in}{1.733401in}}%
\pgfpathlineto{\pgfqpoint{1.911345in}{1.959329in}}%
\pgfpathlineto{\pgfqpoint{1.911864in}{1.919817in}}%
\pgfpathlineto{\pgfqpoint{1.911897in}{1.918567in}}%
\pgfpathlineto{\pgfqpoint{1.911914in}{1.946047in}}%
\pgfpathlineto{\pgfqpoint{1.912341in}{1.934729in}}%
\pgfpathlineto{\pgfqpoint{1.912384in}{1.940449in}}%
\pgfpathlineto{\pgfqpoint{1.912367in}{1.924158in}}%
\pgfpathlineto{\pgfqpoint{1.913073in}{1.935596in}}%
\pgfpathlineto{\pgfqpoint{1.913171in}{1.937908in}}%
\pgfpathlineto{\pgfqpoint{1.913765in}{1.703632in}}%
\pgfpathlineto{\pgfqpoint{1.914028in}{1.532025in}}%
\pgfpathlineto{\pgfqpoint{1.914666in}{1.532039in}}%
\pgfpathlineto{\pgfqpoint{1.915988in}{1.532133in}}%
\pgfpathlineto{\pgfqpoint{1.916320in}{1.964928in}}%
\pgfpathlineto{\pgfqpoint{1.917222in}{1.927112in}}%
\pgfpathlineto{\pgfqpoint{1.917227in}{1.922657in}}%
\pgfpathlineto{\pgfqpoint{1.917242in}{1.941570in}}%
\pgfpathlineto{\pgfqpoint{1.917989in}{1.930551in}}%
\pgfpathlineto{\pgfqpoint{1.918055in}{1.935570in}}%
\pgfpathlineto{\pgfqpoint{1.918682in}{1.717620in}}%
\pgfpathlineto{\pgfqpoint{1.919009in}{1.532030in}}%
\pgfpathlineto{\pgfqpoint{1.919596in}{1.532043in}}%
\pgfpathlineto{\pgfqpoint{1.920931in}{1.532113in}}%
\pgfpathlineto{\pgfqpoint{1.920974in}{1.744723in}}%
\pgfpathlineto{\pgfqpoint{1.921022in}{1.957474in}}%
\pgfpathlineto{\pgfqpoint{1.921722in}{1.932731in}}%
\pgfpathlineto{\pgfqpoint{1.921969in}{1.941993in}}%
\pgfpathlineto{\pgfqpoint{1.921920in}{1.922467in}}%
\pgfpathlineto{\pgfqpoint{1.922316in}{1.940517in}}%
\pgfpathlineto{\pgfqpoint{1.922365in}{1.922734in}}%
\pgfpathlineto{\pgfqpoint{1.922415in}{1.941757in}}%
\pgfpathlineto{\pgfqpoint{1.923057in}{1.931650in}}%
\pgfpathlineto{\pgfqpoint{1.923305in}{1.940383in}}%
\pgfpathlineto{\pgfqpoint{1.923404in}{1.922410in}}%
\pgfpathlineto{\pgfqpoint{1.923453in}{1.926825in}}%
\pgfpathlineto{\pgfqpoint{1.923918in}{1.532026in}}%
\pgfpathlineto{\pgfqpoint{1.924751in}{1.532044in}}%
\pgfpathlineto{\pgfqpoint{1.925881in}{1.532164in}}%
\pgfpathlineto{\pgfqpoint{1.925983in}{1.958838in}}%
\pgfpathlineto{\pgfqpoint{1.926973in}{1.925800in}}%
\pgfpathlineto{\pgfqpoint{1.927351in}{1.938753in}}%
\pgfpathlineto{\pgfqpoint{1.927302in}{1.925473in}}%
\pgfpathlineto{\pgfqpoint{1.927697in}{1.938113in}}%
\pgfpathlineto{\pgfqpoint{1.928897in}{1.532030in}}%
\pgfpathlineto{\pgfqpoint{1.927796in}{1.938386in}}%
\pgfpathlineto{\pgfqpoint{1.930026in}{1.532060in}}%
\pgfpathlineto{\pgfqpoint{1.930824in}{1.532126in}}%
\pgfpathlineto{\pgfqpoint{1.931158in}{1.966588in}}%
\pgfpathlineto{\pgfqpoint{1.932110in}{1.941277in}}%
\pgfpathlineto{\pgfqpoint{1.932126in}{1.923242in}}%
\pgfpathlineto{\pgfqpoint{1.932888in}{1.929520in}}%
\pgfpathlineto{\pgfqpoint{1.933193in}{1.934600in}}%
\pgfpathlineto{\pgfqpoint{1.933242in}{1.928433in}}%
\pgfpathlineto{\pgfqpoint{1.933341in}{1.929144in}}%
\pgfpathlineto{\pgfqpoint{1.933827in}{1.532029in}}%
\pgfpathlineto{\pgfqpoint{1.934695in}{1.532048in}}%
\pgfpathlineto{\pgfqpoint{1.935770in}{1.532130in}}%
\pgfpathlineto{\pgfqpoint{1.935859in}{1.956235in}}%
\pgfpathlineto{\pgfqpoint{1.937074in}{1.936240in}}%
\pgfpathlineto{\pgfqpoint{1.937090in}{1.928189in}}%
\pgfpathlineto{\pgfqpoint{1.937843in}{1.932649in}}%
\pgfpathlineto{\pgfqpoint{1.937935in}{1.933303in}}%
\pgfpathlineto{\pgfqpoint{1.938289in}{1.921135in}}%
\pgfpathlineto{\pgfqpoint{1.938755in}{1.532026in}}%
\pgfpathlineto{\pgfqpoint{1.939834in}{1.532061in}}%
\pgfpathlineto{\pgfqpoint{1.940705in}{1.532084in}}%
\pgfpathlineto{\pgfqpoint{1.940717in}{1.532164in}}%
\pgfpathlineto{\pgfqpoint{1.940819in}{1.958740in}}%
\pgfpathlineto{\pgfqpoint{1.941841in}{1.938376in}}%
\pgfpathlineto{\pgfqpoint{1.942240in}{1.925842in}}%
\pgfpathlineto{\pgfqpoint{1.942586in}{1.927150in}}%
\pgfpathlineto{\pgfqpoint{1.942635in}{1.937827in}}%
\pgfpathlineto{\pgfqpoint{1.943130in}{1.925564in}}%
\pgfpathlineto{\pgfqpoint{1.943229in}{1.927005in}}%
\pgfpathlineto{\pgfqpoint{1.943716in}{1.532030in}}%
\pgfpathlineto{\pgfqpoint{1.944481in}{1.532047in}}%
\pgfpathlineto{\pgfqpoint{1.945665in}{1.532306in}}%
\pgfpathlineto{\pgfqpoint{1.945859in}{1.958197in}}%
\pgfpathlineto{\pgfqpoint{1.946719in}{1.938064in}}%
\pgfpathlineto{\pgfqpoint{1.946733in}{1.923711in}}%
\pgfpathlineto{\pgfqpoint{1.946746in}{1.940545in}}%
\pgfpathlineto{\pgfqpoint{1.947452in}{1.937205in}}%
\pgfpathlineto{\pgfqpoint{1.948623in}{1.542377in}}%
\pgfpathlineto{\pgfqpoint{1.948719in}{1.532032in}}%
\pgfpathlineto{\pgfqpoint{1.949359in}{1.532048in}}%
\pgfpathlineto{\pgfqpoint{1.950607in}{1.532149in}}%
\pgfpathlineto{\pgfqpoint{1.950926in}{1.962750in}}%
\pgfpathlineto{\pgfqpoint{1.951831in}{1.938108in}}%
\pgfpathlineto{\pgfqpoint{1.951876in}{1.926303in}}%
\pgfpathlineto{\pgfqpoint{1.952610in}{1.932939in}}%
\pgfpathlineto{\pgfqpoint{1.953634in}{1.532028in}}%
\pgfpathlineto{\pgfqpoint{1.952671in}{1.933532in}}%
\pgfpathlineto{\pgfqpoint{1.954715in}{1.532052in}}%
\pgfpathlineto{\pgfqpoint{1.955552in}{1.532146in}}%
\pgfpathlineto{\pgfqpoint{1.955721in}{1.957514in}}%
\pgfpathlineto{\pgfqpoint{1.956722in}{1.930728in}}%
\pgfpathlineto{\pgfqpoint{1.956728in}{1.924653in}}%
\pgfpathlineto{\pgfqpoint{1.956781in}{1.938817in}}%
\pgfpathlineto{\pgfqpoint{1.957451in}{1.928915in}}%
\pgfpathlineto{\pgfqpoint{1.957747in}{1.936254in}}%
\pgfpathlineto{\pgfqpoint{1.958044in}{1.918935in}}%
\pgfpathlineto{\pgfqpoint{1.958094in}{1.921644in}}%
\pgfpathlineto{\pgfqpoint{1.958571in}{1.532029in}}%
\pgfpathlineto{\pgfqpoint{1.959354in}{1.532047in}}%
\pgfpathlineto{\pgfqpoint{1.960497in}{1.532129in}}%
\pgfpathlineto{\pgfqpoint{1.960830in}{1.965818in}}%
\pgfpathlineto{\pgfqpoint{1.961761in}{1.935804in}}%
\pgfpathlineto{\pgfqpoint{1.961771in}{1.922512in}}%
\pgfpathlineto{\pgfqpoint{1.961786in}{1.941183in}}%
\pgfpathlineto{\pgfqpoint{1.962533in}{1.929774in}}%
\pgfpathlineto{\pgfqpoint{1.962661in}{1.934863in}}%
\pgfpathlineto{\pgfqpoint{1.963075in}{1.892353in}}%
\pgfpathlineto{\pgfqpoint{1.963526in}{1.532027in}}%
\pgfpathlineto{\pgfqpoint{1.964112in}{1.532041in}}%
\pgfpathlineto{\pgfqpoint{1.965443in}{1.532146in}}%
\pgfpathlineto{\pgfqpoint{1.965534in}{1.957660in}}%
\pgfpathlineto{\pgfqpoint{1.966617in}{1.938547in}}%
\pgfpathlineto{\pgfqpoint{1.966635in}{1.925446in}}%
\pgfpathlineto{\pgfqpoint{1.967379in}{1.926851in}}%
\pgfpathlineto{\pgfqpoint{1.967428in}{1.936645in}}%
\pgfpathlineto{\pgfqpoint{1.967972in}{1.909023in}}%
\pgfpathlineto{\pgfqpoint{1.968428in}{1.532026in}}%
\pgfpathlineto{\pgfqpoint{1.969508in}{1.532061in}}%
\pgfpathlineto{\pgfqpoint{1.970378in}{1.532083in}}%
\pgfpathlineto{\pgfqpoint{1.970390in}{1.532161in}}%
\pgfpathlineto{\pgfqpoint{1.970492in}{1.958562in}}%
\pgfpathlineto{\pgfqpoint{1.971506in}{1.927382in}}%
\pgfpathlineto{\pgfqpoint{1.971547in}{1.925563in}}%
\pgfpathlineto{\pgfqpoint{1.971750in}{1.937974in}}%
\pgfpathlineto{\pgfqpoint{1.972146in}{1.926835in}}%
\pgfpathlineto{\pgfqpoint{1.972195in}{1.937792in}}%
\pgfpathlineto{\pgfqpoint{1.972690in}{1.925615in}}%
\pgfpathlineto{\pgfqpoint{1.972888in}{1.930610in}}%
\pgfpathlineto{\pgfqpoint{1.973356in}{1.532017in}}%
\pgfpathlineto{\pgfqpoint{1.974102in}{1.532054in}}%
\pgfpathlineto{\pgfqpoint{1.975332in}{1.532119in}}%
\pgfpathlineto{\pgfqpoint{1.975386in}{1.927796in}}%
\pgfpathlineto{\pgfqpoint{1.975642in}{1.959748in}}%
\pgfpathlineto{\pgfqpoint{1.975593in}{1.902796in}}%
\pgfpathlineto{\pgfqpoint{1.976149in}{1.933906in}}%
\pgfpathlineto{\pgfqpoint{1.976159in}{1.918252in}}%
\pgfpathlineto{\pgfqpoint{1.976174in}{1.944937in}}%
\pgfpathlineto{\pgfqpoint{1.976916in}{1.937317in}}%
\pgfpathlineto{\pgfqpoint{1.976993in}{1.937361in}}%
\pgfpathlineto{\pgfqpoint{1.977698in}{1.926370in}}%
\pgfpathlineto{\pgfqpoint{1.977747in}{1.936416in}}%
\pgfpathlineto{\pgfqpoint{1.977945in}{1.853924in}}%
\pgfpathlineto{\pgfqpoint{1.978363in}{1.532027in}}%
\pgfpathlineto{\pgfqpoint{1.978899in}{1.532049in}}%
\pgfpathlineto{\pgfqpoint{1.980269in}{1.532085in}}%
\pgfpathlineto{\pgfqpoint{1.980280in}{1.532145in}}%
\pgfpathlineto{\pgfqpoint{1.980370in}{1.957478in}}%
\pgfpathlineto{\pgfqpoint{1.981448in}{1.937423in}}%
\pgfpathlineto{\pgfqpoint{1.981465in}{1.926215in}}%
\pgfpathlineto{\pgfqpoint{1.982190in}{1.933594in}}%
\pgfpathlineto{\pgfqpoint{1.982388in}{1.936229in}}%
\pgfpathlineto{\pgfqpoint{1.982932in}{1.792786in}}%
\pgfpathlineto{\pgfqpoint{1.983308in}{1.532027in}}%
\pgfpathlineto{\pgfqpoint{1.983845in}{1.532049in}}%
\pgfpathlineto{\pgfqpoint{1.985215in}{1.532085in}}%
\pgfpathlineto{\pgfqpoint{1.985225in}{1.532146in}}%
\pgfpathlineto{\pgfqpoint{1.985316in}{1.956813in}}%
\pgfpathlineto{\pgfqpoint{1.986394in}{1.938406in}}%
\pgfpathlineto{\pgfqpoint{1.986413in}{1.925689in}}%
\pgfpathlineto{\pgfqpoint{1.987166in}{1.932770in}}%
\pgfpathlineto{\pgfqpoint{1.987364in}{1.936607in}}%
\pgfpathlineto{\pgfqpoint{1.987809in}{1.881103in}}%
\pgfpathlineto{\pgfqpoint{1.988209in}{1.532026in}}%
\pgfpathlineto{\pgfqpoint{1.988844in}{1.532040in}}%
\pgfpathlineto{\pgfqpoint{1.990171in}{1.532160in}}%
\pgfpathlineto{\pgfqpoint{1.990274in}{1.958497in}}%
\pgfpathlineto{\pgfqpoint{1.991297in}{1.932774in}}%
\pgfpathlineto{\pgfqpoint{1.991424in}{1.938525in}}%
\pgfpathlineto{\pgfqpoint{1.991329in}{1.925198in}}%
\pgfpathlineto{\pgfqpoint{1.992055in}{1.930918in}}%
\pgfpathlineto{\pgfqpoint{1.992302in}{1.937627in}}%
\pgfpathlineto{\pgfqpoint{1.992253in}{1.925227in}}%
\pgfpathlineto{\pgfqpoint{1.992648in}{1.936643in}}%
\pgfpathlineto{\pgfqpoint{1.993198in}{1.532028in}}%
\pgfpathlineto{\pgfqpoint{1.994625in}{1.532069in}}%
\pgfpathlineto{\pgfqpoint{1.995106in}{1.532085in}}%
\pgfpathlineto{\pgfqpoint{1.995116in}{1.532146in}}%
\pgfpathlineto{\pgfqpoint{1.995308in}{1.959096in}}%
\pgfpathlineto{\pgfqpoint{1.996297in}{1.937458in}}%
\pgfpathlineto{\pgfqpoint{1.996349in}{1.925190in}}%
\pgfpathlineto{\pgfqpoint{1.996335in}{1.939058in}}%
\pgfpathlineto{\pgfqpoint{1.997109in}{1.927176in}}%
\pgfpathlineto{\pgfqpoint{1.997159in}{1.935911in}}%
\pgfpathlineto{\pgfqpoint{1.997555in}{1.926682in}}%
\pgfpathlineto{\pgfqpoint{1.997653in}{1.926965in}}%
\pgfpathlineto{\pgfqpoint{1.998141in}{1.532028in}}%
\pgfpathlineto{\pgfqpoint{1.998876in}{1.532053in}}%
\pgfpathlineto{\pgfqpoint{2.000061in}{1.532141in}}%
\pgfpathlineto{\pgfqpoint{2.000256in}{1.960178in}}%
\pgfpathlineto{\pgfqpoint{2.001252in}{1.938643in}}%
\pgfpathlineto{\pgfqpoint{2.001374in}{1.925466in}}%
\pgfpathlineto{\pgfqpoint{2.002013in}{1.927246in}}%
\pgfpathlineto{\pgfqpoint{2.002409in}{1.935564in}}%
\pgfpathlineto{\pgfqpoint{2.002606in}{1.900149in}}%
\pgfpathlineto{\pgfqpoint{2.003073in}{1.532030in}}%
\pgfpathlineto{\pgfqpoint{2.003852in}{1.532047in}}%
\pgfpathlineto{\pgfqpoint{2.005010in}{1.532292in}}%
\pgfpathlineto{\pgfqpoint{2.005107in}{1.957477in}}%
\pgfpathlineto{\pgfqpoint{2.006066in}{1.936059in}}%
\pgfpathlineto{\pgfqpoint{2.006079in}{1.924231in}}%
\pgfpathlineto{\pgfqpoint{2.006091in}{1.939271in}}%
\pgfpathlineto{\pgfqpoint{2.006823in}{1.936077in}}%
\pgfpathlineto{\pgfqpoint{2.007713in}{1.707682in}}%
\pgfpathlineto{\pgfqpoint{2.008002in}{1.532030in}}%
\pgfpathlineto{\pgfqpoint{2.008611in}{1.532047in}}%
\pgfpathlineto{\pgfqpoint{2.009947in}{1.532100in}}%
\pgfpathlineto{\pgfqpoint{2.009956in}{1.532306in}}%
\pgfpathlineto{\pgfqpoint{2.010149in}{1.957941in}}%
\pgfpathlineto{\pgfqpoint{2.011009in}{1.936939in}}%
\pgfpathlineto{\pgfqpoint{2.011052in}{1.923578in}}%
\pgfpathlineto{\pgfqpoint{2.011039in}{1.940149in}}%
\pgfpathlineto{\pgfqpoint{2.011736in}{1.927016in}}%
\pgfpathlineto{\pgfqpoint{2.011785in}{1.936732in}}%
\pgfpathlineto{\pgfqpoint{2.012428in}{1.925690in}}%
\pgfpathlineto{\pgfqpoint{2.012478in}{1.930608in}}%
\pgfpathlineto{\pgfqpoint{2.012979in}{1.532028in}}%
\pgfpathlineto{\pgfqpoint{2.013763in}{1.532045in}}%
\pgfpathlineto{\pgfqpoint{2.014898in}{1.532142in}}%
\pgfpathlineto{\pgfqpoint{2.015088in}{1.957407in}}%
\pgfpathlineto{\pgfqpoint{2.016084in}{1.936194in}}%
\pgfpathlineto{\pgfqpoint{2.016122in}{1.938679in}}%
\pgfpathlineto{\pgfqpoint{2.016107in}{1.924661in}}%
\pgfpathlineto{\pgfqpoint{2.016767in}{1.936258in}}%
\pgfpathlineto{\pgfqpoint{2.017850in}{1.604042in}}%
\pgfpathlineto{\pgfqpoint{2.017882in}{1.532026in}}%
\pgfpathlineto{\pgfqpoint{2.018617in}{1.532042in}}%
\pgfpathlineto{\pgfqpoint{2.019843in}{1.532147in}}%
\pgfpathlineto{\pgfqpoint{2.019935in}{1.956520in}}%
\pgfpathlineto{\pgfqpoint{2.020997in}{1.940620in}}%
\pgfpathlineto{\pgfqpoint{2.021017in}{1.922508in}}%
\pgfpathlineto{\pgfqpoint{2.021036in}{1.940832in}}%
\pgfpathlineto{\pgfqpoint{2.021775in}{1.934080in}}%
\pgfpathlineto{\pgfqpoint{2.021973in}{1.937767in}}%
\pgfpathlineto{\pgfqpoint{2.022517in}{1.763038in}}%
\pgfpathlineto{\pgfqpoint{2.022827in}{1.532027in}}%
\pgfpathlineto{\pgfqpoint{2.023411in}{1.532050in}}%
\pgfpathlineto{\pgfqpoint{2.024778in}{1.532084in}}%
\pgfpathlineto{\pgfqpoint{2.024789in}{1.532146in}}%
\pgfpathlineto{\pgfqpoint{2.024882in}{1.956822in}}%
\pgfpathlineto{\pgfqpoint{2.025924in}{1.937518in}}%
\pgfpathlineto{\pgfqpoint{2.025974in}{1.926082in}}%
\pgfpathlineto{\pgfqpoint{2.026666in}{1.928811in}}%
\pgfpathlineto{\pgfqpoint{2.026815in}{1.937056in}}%
\pgfpathlineto{\pgfqpoint{2.027309in}{1.925581in}}%
\pgfpathlineto{\pgfqpoint{2.027773in}{1.532026in}}%
\pgfpathlineto{\pgfqpoint{2.028655in}{1.532057in}}%
\pgfpathlineto{\pgfqpoint{2.029723in}{1.532084in}}%
\pgfpathlineto{\pgfqpoint{2.029735in}{1.532167in}}%
\pgfpathlineto{\pgfqpoint{2.029837in}{1.958377in}}%
\pgfpathlineto{\pgfqpoint{2.030840in}{1.928934in}}%
\pgfpathlineto{\pgfqpoint{2.030875in}{1.938244in}}%
\pgfpathlineto{\pgfqpoint{2.031228in}{1.925212in}}%
\pgfpathlineto{\pgfqpoint{2.031277in}{1.937814in}}%
\pgfpathlineto{\pgfqpoint{2.031970in}{1.934497in}}%
\pgfpathlineto{\pgfqpoint{2.032069in}{1.936932in}}%
\pgfpathlineto{\pgfqpoint{2.032415in}{1.762014in}}%
\pgfpathlineto{\pgfqpoint{2.032720in}{1.532025in}}%
\pgfpathlineto{\pgfqpoint{2.033358in}{1.532039in}}%
\pgfpathlineto{\pgfqpoint{2.034678in}{1.532126in}}%
\pgfpathlineto{\pgfqpoint{2.035007in}{1.963960in}}%
\pgfpathlineto{\pgfqpoint{2.035977in}{1.936157in}}%
\pgfpathlineto{\pgfqpoint{2.035988in}{1.924245in}}%
\pgfpathlineto{\pgfqpoint{2.036003in}{1.938873in}}%
\pgfpathlineto{\pgfqpoint{2.036745in}{1.934022in}}%
\pgfpathlineto{\pgfqpoint{2.037664in}{1.532026in}}%
\pgfpathlineto{\pgfqpoint{2.038894in}{1.532053in}}%
\pgfpathlineto{\pgfqpoint{2.039625in}{1.532139in}}%
\pgfpathlineto{\pgfqpoint{2.039846in}{1.963444in}}%
\pgfpathlineto{\pgfqpoint{2.040839in}{1.939318in}}%
\pgfpathlineto{\pgfqpoint{2.040953in}{1.924830in}}%
\pgfpathlineto{\pgfqpoint{2.041612in}{1.927732in}}%
\pgfpathlineto{\pgfqpoint{2.041661in}{1.935490in}}%
\pgfpathlineto{\pgfqpoint{2.042156in}{1.926379in}}%
\pgfpathlineto{\pgfqpoint{2.042638in}{1.532028in}}%
\pgfpathlineto{\pgfqpoint{2.043566in}{1.532057in}}%
\pgfpathlineto{\pgfqpoint{2.044571in}{1.532149in}}%
\pgfpathlineto{\pgfqpoint{2.044665in}{1.957889in}}%
\pgfpathlineto{\pgfqpoint{2.045741in}{1.937782in}}%
\pgfpathlineto{\pgfqpoint{2.045758in}{1.925318in}}%
\pgfpathlineto{\pgfqpoint{2.046487in}{1.926985in}}%
\pgfpathlineto{\pgfqpoint{2.046536in}{1.936650in}}%
\pgfpathlineto{\pgfqpoint{2.047130in}{1.909025in}}%
\pgfpathlineto{\pgfqpoint{2.047595in}{1.532028in}}%
\pgfpathlineto{\pgfqpoint{2.048231in}{1.532051in}}%
\pgfpathlineto{\pgfqpoint{2.049517in}{1.532162in}}%
\pgfpathlineto{\pgfqpoint{2.049620in}{1.958324in}}%
\pgfpathlineto{\pgfqpoint{2.050649in}{1.929342in}}%
\pgfpathlineto{\pgfqpoint{2.050659in}{1.924718in}}%
\pgfpathlineto{\pgfqpoint{2.050711in}{1.938350in}}%
\pgfpathlineto{\pgfqpoint{2.051426in}{1.927009in}}%
\pgfpathlineto{\pgfqpoint{2.051574in}{1.937089in}}%
\pgfpathlineto{\pgfqpoint{2.051970in}{1.924950in}}%
\pgfpathlineto{\pgfqpoint{2.052069in}{1.925170in}}%
\pgfpathlineto{\pgfqpoint{2.052515in}{1.532030in}}%
\pgfpathlineto{\pgfqpoint{2.053232in}{1.532051in}}%
\pgfpathlineto{\pgfqpoint{2.054460in}{1.532123in}}%
\pgfpathlineto{\pgfqpoint{2.054686in}{1.960176in}}%
\pgfpathlineto{\pgfqpoint{2.055799in}{1.926505in}}%
\pgfpathlineto{\pgfqpoint{2.055802in}{1.925111in}}%
\pgfpathlineto{\pgfqpoint{2.055815in}{1.937568in}}%
\pgfpathlineto{\pgfqpoint{2.056561in}{1.929472in}}%
\pgfpathlineto{\pgfqpoint{2.056578in}{1.933037in}}%
\pgfpathlineto{\pgfqpoint{2.057052in}{1.885605in}}%
\pgfpathlineto{\pgfqpoint{2.057478in}{1.532028in}}%
\pgfpathlineto{\pgfqpoint{2.058062in}{1.532041in}}%
\pgfpathlineto{\pgfqpoint{2.059407in}{1.532153in}}%
\pgfpathlineto{\pgfqpoint{2.059501in}{1.957619in}}%
\pgfpathlineto{\pgfqpoint{2.060543in}{1.928936in}}%
\pgfpathlineto{\pgfqpoint{2.060840in}{1.935827in}}%
\pgfpathlineto{\pgfqpoint{2.060890in}{1.927431in}}%
\pgfpathlineto{\pgfqpoint{2.061285in}{1.935620in}}%
\pgfpathlineto{\pgfqpoint{2.062126in}{1.714166in}}%
\pgfpathlineto{\pgfqpoint{2.062433in}{1.532028in}}%
\pgfpathlineto{\pgfqpoint{2.063019in}{1.532041in}}%
\pgfpathlineto{\pgfqpoint{2.064352in}{1.532141in}}%
\pgfpathlineto{\pgfqpoint{2.064545in}{1.959446in}}%
\pgfpathlineto{\pgfqpoint{2.065540in}{1.939008in}}%
\pgfpathlineto{\pgfqpoint{2.065554in}{1.924693in}}%
\pgfpathlineto{\pgfqpoint{2.066319in}{1.929888in}}%
\pgfpathlineto{\pgfqpoint{2.066616in}{1.935584in}}%
\pgfpathlineto{\pgfqpoint{2.066566in}{1.926581in}}%
\pgfpathlineto{\pgfqpoint{2.066863in}{1.931932in}}%
\pgfpathlineto{\pgfqpoint{2.067352in}{1.532030in}}%
\pgfpathlineto{\pgfqpoint{2.068116in}{1.532047in}}%
\pgfpathlineto{\pgfqpoint{2.069301in}{1.532306in}}%
\pgfpathlineto{\pgfqpoint{2.069397in}{1.956790in}}%
\pgfpathlineto{\pgfqpoint{2.070357in}{1.939952in}}%
\pgfpathlineto{\pgfqpoint{2.070370in}{1.923391in}}%
\pgfpathlineto{\pgfqpoint{2.071117in}{1.935581in}}%
\pgfpathlineto{\pgfqpoint{2.071216in}{1.936221in}}%
\pgfpathlineto{\pgfqpoint{2.072251in}{1.606183in}}%
\pgfpathlineto{\pgfqpoint{2.072312in}{1.532030in}}%
\pgfpathlineto{\pgfqpoint{2.072994in}{1.532045in}}%
\pgfpathlineto{\pgfqpoint{2.074240in}{1.532110in}}%
\pgfpathlineto{\pgfqpoint{2.074282in}{1.731870in}}%
\pgfpathlineto{\pgfqpoint{2.074343in}{1.957558in}}%
\pgfpathlineto{\pgfqpoint{2.075065in}{1.916383in}}%
\pgfpathlineto{\pgfqpoint{2.075081in}{1.945827in}}%
\pgfpathlineto{\pgfqpoint{2.075831in}{1.935752in}}%
\pgfpathlineto{\pgfqpoint{2.075849in}{1.924394in}}%
\pgfpathlineto{\pgfqpoint{2.076205in}{1.938232in}}%
\pgfpathlineto{\pgfqpoint{2.076601in}{1.924488in}}%
\pgfpathlineto{\pgfqpoint{2.076650in}{1.937353in}}%
\pgfpathlineto{\pgfqpoint{2.076848in}{1.865313in}}%
\pgfpathlineto{\pgfqpoint{2.077268in}{1.532028in}}%
\pgfpathlineto{\pgfqpoint{2.077854in}{1.532041in}}%
\pgfpathlineto{\pgfqpoint{2.079188in}{1.532140in}}%
\pgfpathlineto{\pgfqpoint{2.079392in}{1.960426in}}%
\pgfpathlineto{\pgfqpoint{2.080378in}{1.936206in}}%
\pgfpathlineto{\pgfqpoint{2.080423in}{1.924797in}}%
\pgfpathlineto{\pgfqpoint{2.080408in}{1.938146in}}%
\pgfpathlineto{\pgfqpoint{2.081127in}{1.929495in}}%
\pgfpathlineto{\pgfqpoint{2.081424in}{1.935599in}}%
\pgfpathlineto{\pgfqpoint{2.081375in}{1.926663in}}%
\pgfpathlineto{\pgfqpoint{2.081671in}{1.931988in}}%
\pgfpathlineto{\pgfqpoint{2.082195in}{1.532030in}}%
\pgfpathlineto{\pgfqpoint{2.083115in}{1.532055in}}%
\pgfpathlineto{\pgfqpoint{2.084129in}{1.532100in}}%
\pgfpathlineto{\pgfqpoint{2.084138in}{1.532306in}}%
\pgfpathlineto{\pgfqpoint{2.084234in}{1.956970in}}%
\pgfpathlineto{\pgfqpoint{2.085184in}{1.931033in}}%
\pgfpathlineto{\pgfqpoint{2.085191in}{1.922785in}}%
\pgfpathlineto{\pgfqpoint{2.085243in}{1.939501in}}%
\pgfpathlineto{\pgfqpoint{2.085950in}{1.928820in}}%
\pgfpathlineto{\pgfqpoint{2.086098in}{1.936465in}}%
\pgfpathlineto{\pgfqpoint{2.086642in}{1.920356in}}%
\pgfpathlineto{\pgfqpoint{2.086840in}{1.721566in}}%
\pgfpathlineto{\pgfqpoint{2.087159in}{1.532028in}}%
\pgfpathlineto{\pgfqpoint{2.087745in}{1.532041in}}%
\pgfpathlineto{\pgfqpoint{2.089079in}{1.532140in}}%
\pgfpathlineto{\pgfqpoint{2.089287in}{1.960605in}}%
\pgfpathlineto{\pgfqpoint{2.090270in}{1.937881in}}%
\pgfpathlineto{\pgfqpoint{2.090292in}{1.924211in}}%
\pgfpathlineto{\pgfqpoint{2.090312in}{1.938322in}}%
\pgfpathlineto{\pgfqpoint{2.091042in}{1.928156in}}%
\pgfpathlineto{\pgfqpoint{2.091338in}{1.935680in}}%
\pgfpathlineto{\pgfqpoint{2.091388in}{1.926330in}}%
\pgfpathlineto{\pgfqpoint{2.091586in}{1.930542in}}%
\pgfpathlineto{\pgfqpoint{2.092089in}{1.532030in}}%
\pgfpathlineto{\pgfqpoint{2.092864in}{1.532047in}}%
\pgfpathlineto{\pgfqpoint{2.094029in}{1.532292in}}%
\pgfpathlineto{\pgfqpoint{2.094126in}{1.957327in}}%
\pgfpathlineto{\pgfqpoint{2.095088in}{1.935128in}}%
\pgfpathlineto{\pgfqpoint{2.095099in}{1.923334in}}%
\pgfpathlineto{\pgfqpoint{2.095116in}{1.938655in}}%
\pgfpathlineto{\pgfqpoint{2.095834in}{1.927611in}}%
\pgfpathlineto{\pgfqpoint{2.095867in}{1.935470in}}%
\pgfpathlineto{\pgfqpoint{2.096506in}{1.925317in}}%
\pgfpathlineto{\pgfqpoint{2.096555in}{1.930385in}}%
\pgfpathlineto{\pgfqpoint{2.097047in}{1.532029in}}%
\pgfpathlineto{\pgfqpoint{2.097831in}{1.532046in}}%
\pgfpathlineto{\pgfqpoint{2.098971in}{1.532149in}}%
\pgfpathlineto{\pgfqpoint{2.099064in}{1.956687in}}%
\pgfpathlineto{\pgfqpoint{2.100132in}{1.931763in}}%
\pgfpathlineto{\pgfqpoint{2.100362in}{1.936475in}}%
\pgfpathlineto{\pgfqpoint{2.100312in}{1.926055in}}%
\pgfpathlineto{\pgfqpoint{2.100807in}{1.936205in}}%
\pgfpathlineto{\pgfqpoint{2.101747in}{1.644816in}}%
\pgfpathlineto{\pgfqpoint{2.101983in}{1.532028in}}%
\pgfpathlineto{\pgfqpoint{2.102664in}{1.532043in}}%
\pgfpathlineto{\pgfqpoint{2.103916in}{1.532149in}}%
\pgfpathlineto{\pgfqpoint{2.104010in}{1.957094in}}%
\pgfpathlineto{\pgfqpoint{2.105082in}{1.933086in}}%
\pgfpathlineto{\pgfqpoint{2.105230in}{1.925760in}}%
\pgfpathlineto{\pgfqpoint{2.105280in}{1.937039in}}%
\pgfpathlineto{\pgfqpoint{2.105873in}{1.929437in}}%
\pgfpathlineto{\pgfqpoint{2.106170in}{1.936220in}}%
\pgfpathlineto{\pgfqpoint{2.106120in}{1.925684in}}%
\pgfpathlineto{\pgfqpoint{2.106417in}{1.932376in}}%
\pgfpathlineto{\pgfqpoint{2.106926in}{1.532030in}}%
\pgfpathlineto{\pgfqpoint{2.107702in}{1.532047in}}%
\pgfpathlineto{\pgfqpoint{2.108865in}{1.532292in}}%
\pgfpathlineto{\pgfqpoint{2.108962in}{1.956917in}}%
\pgfpathlineto{\pgfqpoint{2.109924in}{1.924729in}}%
\pgfpathlineto{\pgfqpoint{2.109944in}{1.939381in}}%
\pgfpathlineto{\pgfqpoint{2.110023in}{1.924225in}}%
\pgfpathlineto{\pgfqpoint{2.110744in}{1.935506in}}%
\pgfpathlineto{\pgfqpoint{2.111584in}{1.702402in}}%
\pgfpathlineto{\pgfqpoint{2.111846in}{1.532025in}}%
\pgfpathlineto{\pgfqpoint{2.112483in}{1.532039in}}%
\pgfpathlineto{\pgfqpoint{2.113806in}{1.532135in}}%
\pgfpathlineto{\pgfqpoint{2.114144in}{1.968127in}}%
\pgfpathlineto{\pgfqpoint{2.115033in}{1.925761in}}%
\pgfpathlineto{\pgfqpoint{2.115038in}{1.920161in}}%
\pgfpathlineto{\pgfqpoint{2.115054in}{1.942199in}}%
\pgfpathlineto{\pgfqpoint{2.115802in}{1.927383in}}%
\pgfpathlineto{\pgfqpoint{2.115820in}{1.934635in}}%
\pgfpathlineto{\pgfqpoint{2.116384in}{1.887657in}}%
\pgfpathlineto{\pgfqpoint{2.116791in}{1.532026in}}%
\pgfpathlineto{\pgfqpoint{2.117427in}{1.532040in}}%
\pgfpathlineto{\pgfqpoint{2.118753in}{1.532153in}}%
\pgfpathlineto{\pgfqpoint{2.118846in}{1.957524in}}%
\pgfpathlineto{\pgfqpoint{2.119883in}{1.927811in}}%
\pgfpathlineto{\pgfqpoint{2.119973in}{1.925192in}}%
\pgfpathlineto{\pgfqpoint{2.120068in}{1.937474in}}%
\pgfpathlineto{\pgfqpoint{2.120600in}{1.934494in}}%
\pgfpathlineto{\pgfqpoint{2.120847in}{1.936599in}}%
\pgfpathlineto{\pgfqpoint{2.121391in}{1.837865in}}%
\pgfpathlineto{\pgfqpoint{2.121774in}{1.532032in}}%
\pgfpathlineto{\pgfqpoint{2.122310in}{1.532045in}}%
\pgfpathlineto{\pgfqpoint{2.123699in}{1.532168in}}%
\pgfpathlineto{\pgfqpoint{2.124024in}{1.965297in}}%
\pgfpathlineto{\pgfqpoint{2.124863in}{1.928150in}}%
\pgfpathlineto{\pgfqpoint{2.124871in}{1.922975in}}%
\pgfpathlineto{\pgfqpoint{2.124884in}{1.940185in}}%
\pgfpathlineto{\pgfqpoint{2.125629in}{1.933638in}}%
\pgfpathlineto{\pgfqpoint{2.126712in}{1.532028in}}%
\pgfpathlineto{\pgfqpoint{2.125656in}{1.933760in}}%
\pgfpathlineto{\pgfqpoint{2.127147in}{1.532046in}}%
\pgfpathlineto{\pgfqpoint{2.128644in}{1.532151in}}%
\pgfpathlineto{\pgfqpoint{2.128738in}{1.957599in}}%
\pgfpathlineto{\pgfqpoint{2.129807in}{1.934963in}}%
\pgfpathlineto{\pgfqpoint{2.129815in}{1.937093in}}%
\pgfpathlineto{\pgfqpoint{2.129940in}{1.925306in}}%
\pgfpathlineto{\pgfqpoint{2.130567in}{1.933629in}}%
\pgfpathlineto{\pgfqpoint{2.130765in}{1.936280in}}%
\pgfpathlineto{\pgfqpoint{2.131309in}{1.786972in}}%
\pgfpathlineto{\pgfqpoint{2.131657in}{1.532028in}}%
\pgfpathlineto{\pgfqpoint{2.132240in}{1.532041in}}%
\pgfpathlineto{\pgfqpoint{2.133589in}{1.532151in}}%
\pgfpathlineto{\pgfqpoint{2.133682in}{1.956636in}}%
\pgfpathlineto{\pgfqpoint{2.134718in}{1.936346in}}%
\pgfpathlineto{\pgfqpoint{2.135113in}{1.925668in}}%
\pgfpathlineto{\pgfqpoint{2.135064in}{1.936896in}}%
\pgfpathlineto{\pgfqpoint{2.135459in}{1.926264in}}%
\pgfpathlineto{\pgfqpoint{2.135509in}{1.936650in}}%
\pgfpathlineto{\pgfqpoint{2.136003in}{1.925419in}}%
\pgfpathlineto{\pgfqpoint{2.136102in}{1.927187in}}%
\pgfpathlineto{\pgfqpoint{2.136594in}{1.532030in}}%
\pgfpathlineto{\pgfqpoint{2.137365in}{1.532047in}}%
\pgfpathlineto{\pgfqpoint{2.138538in}{1.532306in}}%
\pgfpathlineto{\pgfqpoint{2.138731in}{1.957062in}}%
\pgfpathlineto{\pgfqpoint{2.139594in}{1.939900in}}%
\pgfpathlineto{\pgfqpoint{2.139607in}{1.922947in}}%
\pgfpathlineto{\pgfqpoint{2.140349in}{1.935001in}}%
\pgfpathlineto{\pgfqpoint{2.140448in}{1.936251in}}%
\pgfpathlineto{\pgfqpoint{2.141338in}{1.631793in}}%
\pgfpathlineto{\pgfqpoint{2.141538in}{1.532030in}}%
\pgfpathlineto{\pgfqpoint{2.141482in}{1.651277in}}%
\pgfpathlineto{\pgfqpoint{2.142405in}{1.532049in}}%
\pgfpathlineto{\pgfqpoint{2.143483in}{1.532306in}}%
\pgfpathlineto{\pgfqpoint{2.143677in}{1.957558in}}%
\pgfpathlineto{\pgfqpoint{2.144538in}{1.937316in}}%
\pgfpathlineto{\pgfqpoint{2.144552in}{1.922594in}}%
\pgfpathlineto{\pgfqpoint{2.144564in}{1.939289in}}%
\pgfpathlineto{\pgfqpoint{2.145275in}{1.934267in}}%
\pgfpathlineto{\pgfqpoint{2.145621in}{1.935877in}}%
\pgfpathlineto{\pgfqpoint{2.146115in}{1.822059in}}%
\pgfpathlineto{\pgfqpoint{2.146507in}{1.532028in}}%
\pgfpathlineto{\pgfqpoint{2.147043in}{1.532049in}}%
\pgfpathlineto{\pgfqpoint{2.148415in}{1.532085in}}%
\pgfpathlineto{\pgfqpoint{2.148425in}{1.532146in}}%
\pgfpathlineto{\pgfqpoint{2.148619in}{1.959059in}}%
\pgfpathlineto{\pgfqpoint{2.149596in}{1.924667in}}%
\pgfpathlineto{\pgfqpoint{2.149649in}{1.938284in}}%
\pgfpathlineto{\pgfqpoint{2.150350in}{1.928558in}}%
\pgfpathlineto{\pgfqpoint{2.150647in}{1.935272in}}%
\pgfpathlineto{\pgfqpoint{2.150696in}{1.926587in}}%
\pgfpathlineto{\pgfqpoint{2.150894in}{1.931007in}}%
\pgfpathlineto{\pgfqpoint{2.151452in}{1.532028in}}%
\pgfpathlineto{\pgfqpoint{2.152632in}{1.532054in}}%
\pgfpathlineto{\pgfqpoint{2.153370in}{1.532142in}}%
\pgfpathlineto{\pgfqpoint{2.153461in}{1.955754in}}%
\pgfpathlineto{\pgfqpoint{2.154560in}{1.937834in}}%
\pgfpathlineto{\pgfqpoint{2.154611in}{1.924504in}}%
\pgfpathlineto{\pgfqpoint{2.155329in}{1.930138in}}%
\pgfpathlineto{\pgfqpoint{2.155379in}{1.930209in}}%
\pgfpathlineto{\pgfqpoint{2.155527in}{1.935451in}}%
\pgfpathlineto{\pgfqpoint{2.155577in}{1.926025in}}%
\pgfpathlineto{\pgfqpoint{2.155873in}{1.932540in}}%
\pgfpathlineto{\pgfqpoint{2.156376in}{1.532030in}}%
\pgfpathlineto{\pgfqpoint{2.157196in}{1.532052in}}%
\pgfpathlineto{\pgfqpoint{2.158311in}{1.532100in}}%
\pgfpathlineto{\pgfqpoint{2.158320in}{1.532306in}}%
\pgfpathlineto{\pgfqpoint{2.158512in}{1.956399in}}%
\pgfpathlineto{\pgfqpoint{2.159376in}{1.939638in}}%
\pgfpathlineto{\pgfqpoint{2.159389in}{1.922959in}}%
\pgfpathlineto{\pgfqpoint{2.160116in}{1.925910in}}%
\pgfpathlineto{\pgfqpoint{2.160511in}{1.935861in}}%
\pgfpathlineto{\pgfqpoint{2.160808in}{1.917709in}}%
\pgfpathlineto{\pgfqpoint{2.160857in}{1.924902in}}%
\pgfpathlineto{\pgfqpoint{2.161316in}{1.532029in}}%
\pgfpathlineto{\pgfqpoint{2.162182in}{1.532048in}}%
\pgfpathlineto{\pgfqpoint{2.163260in}{1.532126in}}%
\pgfpathlineto{\pgfqpoint{2.163489in}{1.959867in}}%
\pgfpathlineto{\pgfqpoint{2.164594in}{1.937364in}}%
\pgfpathlineto{\pgfqpoint{2.164607in}{1.924429in}}%
\pgfpathlineto{\pgfqpoint{2.165369in}{1.929440in}}%
\pgfpathlineto{\pgfqpoint{2.165388in}{1.932772in}}%
\pgfpathlineto{\pgfqpoint{2.165871in}{1.863354in}}%
\pgfpathlineto{\pgfqpoint{2.166274in}{1.532028in}}%
\pgfpathlineto{\pgfqpoint{2.166856in}{1.532041in}}%
\pgfpathlineto{\pgfqpoint{2.168207in}{1.532149in}}%
\pgfpathlineto{\pgfqpoint{2.168301in}{1.957323in}}%
\pgfpathlineto{\pgfqpoint{2.169349in}{1.936229in}}%
\pgfpathlineto{\pgfqpoint{2.169741in}{1.925839in}}%
\pgfpathlineto{\pgfqpoint{2.169691in}{1.936236in}}%
\pgfpathlineto{\pgfqpoint{2.170087in}{1.926813in}}%
\pgfpathlineto{\pgfqpoint{2.170136in}{1.936024in}}%
\pgfpathlineto{\pgfqpoint{2.170631in}{1.925504in}}%
\pgfpathlineto{\pgfqpoint{2.170730in}{1.926744in}}%
\pgfpathlineto{\pgfqpoint{2.171191in}{1.532026in}}%
\pgfpathlineto{\pgfqpoint{2.172024in}{1.532044in}}%
\pgfpathlineto{\pgfqpoint{2.173153in}{1.532166in}}%
\pgfpathlineto{\pgfqpoint{2.173256in}{1.957708in}}%
\pgfpathlineto{\pgfqpoint{2.174264in}{1.934790in}}%
\pgfpathlineto{\pgfqpoint{2.174714in}{1.924448in}}%
\pgfpathlineto{\pgfqpoint{2.174334in}{1.937881in}}%
\pgfpathlineto{\pgfqpoint{2.175060in}{1.926047in}}%
\pgfpathlineto{\pgfqpoint{2.175109in}{1.937114in}}%
\pgfpathlineto{\pgfqpoint{2.175604in}{1.924068in}}%
\pgfpathlineto{\pgfqpoint{2.175703in}{1.924475in}}%
\pgfpathlineto{\pgfqpoint{2.176159in}{1.532029in}}%
\pgfpathlineto{\pgfqpoint{2.176883in}{1.532052in}}%
\pgfpathlineto{\pgfqpoint{2.178096in}{1.532120in}}%
\pgfpathlineto{\pgfqpoint{2.178197in}{1.957296in}}%
\pgfpathlineto{\pgfqpoint{2.179617in}{1.929843in}}%
\pgfpathlineto{\pgfqpoint{2.179628in}{1.925576in}}%
\pgfpathlineto{\pgfqpoint{2.179646in}{1.936434in}}%
\pgfpathlineto{\pgfqpoint{2.180365in}{1.926132in}}%
\pgfpathlineto{\pgfqpoint{2.180414in}{1.934170in}}%
\pgfpathlineto{\pgfqpoint{2.180711in}{1.863847in}}%
\pgfpathlineto{\pgfqpoint{2.181111in}{1.532028in}}%
\pgfpathlineto{\pgfqpoint{2.181694in}{1.532041in}}%
\pgfpathlineto{\pgfqpoint{2.183044in}{1.532151in}}%
\pgfpathlineto{\pgfqpoint{2.183137in}{1.956311in}}%
\pgfpathlineto{\pgfqpoint{2.184165in}{1.927926in}}%
\pgfpathlineto{\pgfqpoint{2.184313in}{1.936468in}}%
\pgfpathlineto{\pgfqpoint{2.184264in}{1.925850in}}%
\pgfpathlineto{\pgfqpoint{2.184956in}{1.931090in}}%
\pgfpathlineto{\pgfqpoint{2.185203in}{1.935818in}}%
\pgfpathlineto{\pgfqpoint{2.185154in}{1.925838in}}%
\pgfpathlineto{\pgfqpoint{2.185550in}{1.934210in}}%
\pgfpathlineto{\pgfqpoint{2.186027in}{1.532026in}}%
\pgfpathlineto{\pgfqpoint{2.186958in}{1.532047in}}%
\pgfpathlineto{\pgfqpoint{2.187990in}{1.532172in}}%
\pgfpathlineto{\pgfqpoint{2.188092in}{1.957806in}}%
\pgfpathlineto{\pgfqpoint{2.189099in}{1.929348in}}%
\pgfpathlineto{\pgfqpoint{2.189216in}{1.937911in}}%
\pgfpathlineto{\pgfqpoint{2.189167in}{1.924434in}}%
\pgfpathlineto{\pgfqpoint{2.189859in}{1.931207in}}%
\pgfpathlineto{\pgfqpoint{2.190106in}{1.937240in}}%
\pgfpathlineto{\pgfqpoint{2.190057in}{1.924529in}}%
\pgfpathlineto{\pgfqpoint{2.190452in}{1.935749in}}%
\pgfpathlineto{\pgfqpoint{2.190700in}{1.730348in}}%
\pgfpathlineto{\pgfqpoint{2.190983in}{1.532030in}}%
\pgfpathlineto{\pgfqpoint{2.191597in}{1.532048in}}%
\pgfpathlineto{\pgfqpoint{2.192931in}{1.532114in}}%
\pgfpathlineto{\pgfqpoint{2.192974in}{1.746488in}}%
\pgfpathlineto{\pgfqpoint{2.193228in}{1.957456in}}%
\pgfpathlineto{\pgfqpoint{2.193755in}{1.941402in}}%
\pgfpathlineto{\pgfqpoint{2.193761in}{1.944226in}}%
\pgfpathlineto{\pgfqpoint{2.193776in}{1.918146in}}%
\pgfpathlineto{\pgfqpoint{2.194512in}{1.934533in}}%
\pgfpathlineto{\pgfqpoint{2.194572in}{1.925038in}}%
\pgfpathlineto{\pgfqpoint{2.194670in}{1.937085in}}%
\pgfpathlineto{\pgfqpoint{2.195264in}{1.927423in}}%
\pgfpathlineto{\pgfqpoint{2.195462in}{1.932569in}}%
\pgfpathlineto{\pgfqpoint{2.195511in}{1.903133in}}%
\pgfpathlineto{\pgfqpoint{2.195939in}{1.532029in}}%
\pgfpathlineto{\pgfqpoint{2.196562in}{1.532049in}}%
\pgfpathlineto{\pgfqpoint{2.197878in}{1.532121in}}%
\pgfpathlineto{\pgfqpoint{2.197979in}{1.957566in}}%
\pgfpathlineto{\pgfqpoint{2.199394in}{1.929562in}}%
\pgfpathlineto{\pgfqpoint{2.199404in}{1.926269in}}%
\pgfpathlineto{\pgfqpoint{2.199422in}{1.935669in}}%
\pgfpathlineto{\pgfqpoint{2.200109in}{1.929005in}}%
\pgfpathlineto{\pgfqpoint{2.200257in}{1.934715in}}%
\pgfpathlineto{\pgfqpoint{2.200455in}{1.890140in}}%
\pgfpathlineto{\pgfqpoint{2.200908in}{1.532027in}}%
\pgfpathlineto{\pgfqpoint{2.201544in}{1.532052in}}%
\pgfpathlineto{\pgfqpoint{2.202815in}{1.532085in}}%
\pgfpathlineto{\pgfqpoint{2.202825in}{1.532145in}}%
\pgfpathlineto{\pgfqpoint{2.202919in}{1.956318in}}%
\pgfpathlineto{\pgfqpoint{2.203989in}{1.936146in}}%
\pgfpathlineto{\pgfqpoint{2.204078in}{1.925910in}}%
\pgfpathlineto{\pgfqpoint{2.204730in}{1.926099in}}%
\pgfpathlineto{\pgfqpoint{2.204780in}{1.935416in}}%
\pgfpathlineto{\pgfqpoint{2.205373in}{1.917003in}}%
\pgfpathlineto{\pgfqpoint{2.205841in}{1.532028in}}%
\pgfpathlineto{\pgfqpoint{2.206623in}{1.532045in}}%
\pgfpathlineto{\pgfqpoint{2.207771in}{1.532153in}}%
\pgfpathlineto{\pgfqpoint{2.207864in}{1.956784in}}%
\pgfpathlineto{\pgfqpoint{2.208900in}{1.929037in}}%
\pgfpathlineto{\pgfqpoint{2.209048in}{1.937073in}}%
\pgfpathlineto{\pgfqpoint{2.209098in}{1.925063in}}%
\pgfpathlineto{\pgfqpoint{2.209691in}{1.932733in}}%
\pgfpathlineto{\pgfqpoint{2.209938in}{1.936496in}}%
\pgfpathlineto{\pgfqpoint{2.210334in}{1.898768in}}%
\pgfpathlineto{\pgfqpoint{2.210754in}{1.532027in}}%
\pgfpathlineto{\pgfqpoint{2.211486in}{1.532042in}}%
\pgfpathlineto{\pgfqpoint{2.212716in}{1.532145in}}%
\pgfpathlineto{\pgfqpoint{2.212810in}{1.956676in}}%
\pgfpathlineto{\pgfqpoint{2.213879in}{1.934978in}}%
\pgfpathlineto{\pgfqpoint{2.214006in}{1.925687in}}%
\pgfpathlineto{\pgfqpoint{2.214056in}{1.936164in}}%
\pgfpathlineto{\pgfqpoint{2.214649in}{1.930066in}}%
\pgfpathlineto{\pgfqpoint{2.214699in}{1.929971in}}%
\pgfpathlineto{\pgfqpoint{2.214847in}{1.935607in}}%
\pgfpathlineto{\pgfqpoint{2.214896in}{1.925688in}}%
\pgfpathlineto{\pgfqpoint{2.215193in}{1.932816in}}%
\pgfpathlineto{\pgfqpoint{2.215700in}{1.532026in}}%
\pgfpathlineto{\pgfqpoint{2.216780in}{1.532061in}}%
\pgfpathlineto{\pgfqpoint{2.217651in}{1.532084in}}%
\pgfpathlineto{\pgfqpoint{2.217663in}{1.532172in}}%
\pgfpathlineto{\pgfqpoint{2.217765in}{1.957599in}}%
\pgfpathlineto{\pgfqpoint{2.218760in}{1.924019in}}%
\pgfpathlineto{\pgfqpoint{2.218959in}{1.937970in}}%
\pgfpathlineto{\pgfqpoint{2.219503in}{1.935970in}}%
\pgfpathlineto{\pgfqpoint{2.219849in}{1.937243in}}%
\pgfpathlineto{\pgfqpoint{2.220393in}{1.693433in}}%
\pgfpathlineto{\pgfqpoint{2.220652in}{1.532031in}}%
\pgfpathlineto{\pgfqpoint{2.221307in}{1.532044in}}%
\pgfpathlineto{\pgfqpoint{2.222611in}{1.532306in}}%
\pgfpathlineto{\pgfqpoint{2.222804in}{1.957351in}}%
\pgfpathlineto{\pgfqpoint{2.223662in}{1.939646in}}%
\pgfpathlineto{\pgfqpoint{2.223708in}{1.922532in}}%
\pgfpathlineto{\pgfqpoint{2.224424in}{1.935973in}}%
\pgfpathlineto{\pgfqpoint{2.225562in}{1.563275in}}%
\pgfpathlineto{\pgfqpoint{2.225592in}{1.532025in}}%
\pgfpathlineto{\pgfqpoint{2.226328in}{1.532041in}}%
\pgfpathlineto{\pgfqpoint{2.227552in}{1.532132in}}%
\pgfpathlineto{\pgfqpoint{2.227891in}{1.966306in}}%
\pgfpathlineto{\pgfqpoint{2.228796in}{1.935398in}}%
\pgfpathlineto{\pgfqpoint{2.228806in}{1.920920in}}%
\pgfpathlineto{\pgfqpoint{2.228822in}{1.940840in}}%
\pgfpathlineto{\pgfqpoint{2.229568in}{1.928866in}}%
\pgfpathlineto{\pgfqpoint{2.229594in}{1.934281in}}%
\pgfpathlineto{\pgfqpoint{2.230351in}{1.618436in}}%
\pgfpathlineto{\pgfqpoint{2.230562in}{1.532029in}}%
\pgfpathlineto{\pgfqpoint{2.230501in}{1.630650in}}%
\pgfpathlineto{\pgfqpoint{2.231291in}{1.532052in}}%
\pgfpathlineto{\pgfqpoint{2.232498in}{1.532145in}}%
\pgfpathlineto{\pgfqpoint{2.232588in}{1.956794in}}%
\pgfpathlineto{\pgfqpoint{2.233675in}{1.937737in}}%
\pgfpathlineto{\pgfqpoint{2.233692in}{1.924232in}}%
\pgfpathlineto{\pgfqpoint{2.234433in}{1.931421in}}%
\pgfpathlineto{\pgfqpoint{2.234631in}{1.935723in}}%
\pgfpathlineto{\pgfqpoint{2.235076in}{1.897603in}}%
\pgfpathlineto{\pgfqpoint{2.235510in}{1.532028in}}%
\pgfpathlineto{\pgfqpoint{2.236190in}{1.532043in}}%
\pgfpathlineto{\pgfqpoint{2.237444in}{1.532149in}}%
\pgfpathlineto{\pgfqpoint{2.237537in}{1.956682in}}%
\pgfpathlineto{\pgfqpoint{2.238567in}{1.925508in}}%
\pgfpathlineto{\pgfqpoint{2.238616in}{1.936225in}}%
\pgfpathlineto{\pgfqpoint{2.239309in}{1.933364in}}%
\pgfpathlineto{\pgfqpoint{2.239408in}{1.935676in}}%
\pgfpathlineto{\pgfqpoint{2.240100in}{1.795333in}}%
\pgfpathlineto{\pgfqpoint{2.240453in}{1.532029in}}%
\pgfpathlineto{\pgfqpoint{2.241034in}{1.532041in}}%
\pgfpathlineto{\pgfqpoint{2.242389in}{1.532145in}}%
\pgfpathlineto{\pgfqpoint{2.242479in}{1.955965in}}%
\pgfpathlineto{\pgfqpoint{2.243572in}{1.937006in}}%
\pgfpathlineto{\pgfqpoint{2.243589in}{1.924945in}}%
\pgfpathlineto{\pgfqpoint{2.244322in}{1.926845in}}%
\pgfpathlineto{\pgfqpoint{2.244371in}{1.935817in}}%
\pgfpathlineto{\pgfqpoint{2.244964in}{1.906247in}}%
\pgfpathlineto{\pgfqpoint{2.245385in}{1.532030in}}%
\pgfpathlineto{\pgfqpoint{2.246049in}{1.532044in}}%
\pgfpathlineto{\pgfqpoint{2.247332in}{1.532117in}}%
\pgfpathlineto{\pgfqpoint{2.247374in}{1.767586in}}%
\pgfpathlineto{\pgfqpoint{2.247433in}{1.957357in}}%
\pgfpathlineto{\pgfqpoint{2.248157in}{1.920543in}}%
\pgfpathlineto{\pgfqpoint{2.248174in}{1.941525in}}%
\pgfpathlineto{\pgfqpoint{2.248190in}{1.920321in}}%
\pgfpathlineto{\pgfqpoint{2.248945in}{1.929167in}}%
\pgfpathlineto{\pgfqpoint{2.249064in}{1.935820in}}%
\pgfpathlineto{\pgfqpoint{2.249459in}{1.925840in}}%
\pgfpathlineto{\pgfqpoint{2.249707in}{1.929702in}}%
\pgfpathlineto{\pgfqpoint{2.249855in}{1.933945in}}%
\pgfpathlineto{\pgfqpoint{2.249904in}{1.904482in}}%
\pgfpathlineto{\pgfqpoint{2.250349in}{1.532028in}}%
\pgfpathlineto{\pgfqpoint{2.250981in}{1.532050in}}%
\pgfpathlineto{\pgfqpoint{2.252280in}{1.532152in}}%
\pgfpathlineto{\pgfqpoint{2.252373in}{1.956312in}}%
\pgfpathlineto{\pgfqpoint{2.253405in}{1.937682in}}%
\pgfpathlineto{\pgfqpoint{2.253701in}{1.922703in}}%
\pgfpathlineto{\pgfqpoint{2.253751in}{1.939654in}}%
\pgfpathlineto{\pgfqpoint{2.254146in}{1.922761in}}%
\pgfpathlineto{\pgfqpoint{2.254196in}{1.939280in}}%
\pgfpathlineto{\pgfqpoint{2.254592in}{1.922561in}}%
\pgfpathlineto{\pgfqpoint{2.254789in}{1.926597in}}%
\pgfpathlineto{\pgfqpoint{2.255246in}{1.532020in}}%
\pgfpathlineto{\pgfqpoint{2.255945in}{1.532041in}}%
\pgfpathlineto{\pgfqpoint{2.257223in}{1.532119in}}%
\pgfpathlineto{\pgfqpoint{2.257263in}{1.726896in}}%
\pgfpathlineto{\pgfqpoint{2.257546in}{1.958243in}}%
\pgfpathlineto{\pgfqpoint{2.258046in}{1.940539in}}%
\pgfpathlineto{\pgfqpoint{2.258052in}{1.944942in}}%
\pgfpathlineto{\pgfqpoint{2.258068in}{1.917057in}}%
\pgfpathlineto{\pgfqpoint{2.258792in}{1.936933in}}%
\pgfpathlineto{\pgfqpoint{2.258921in}{1.925159in}}%
\pgfpathlineto{\pgfqpoint{2.259572in}{1.925397in}}%
\pgfpathlineto{\pgfqpoint{2.259621in}{1.935602in}}%
\pgfpathlineto{\pgfqpoint{2.259819in}{1.882121in}}%
\pgfpathlineto{\pgfqpoint{2.260239in}{1.532028in}}%
\pgfpathlineto{\pgfqpoint{2.260822in}{1.532041in}}%
\pgfpathlineto{\pgfqpoint{2.262171in}{1.532150in}}%
\pgfpathlineto{\pgfqpoint{2.262264in}{1.957138in}}%
\pgfpathlineto{\pgfqpoint{2.263326in}{1.931122in}}%
\pgfpathlineto{\pgfqpoint{2.263574in}{1.925399in}}%
\pgfpathlineto{\pgfqpoint{2.263524in}{1.936627in}}%
\pgfpathlineto{\pgfqpoint{2.264019in}{1.925519in}}%
\pgfpathlineto{\pgfqpoint{2.264414in}{1.935947in}}%
\pgfpathlineto{\pgfqpoint{2.264464in}{1.925395in}}%
\pgfpathlineto{\pgfqpoint{2.264662in}{1.930784in}}%
\pgfpathlineto{\pgfqpoint{2.265200in}{1.532027in}}%
\pgfpathlineto{\pgfqpoint{2.266033in}{1.532056in}}%
\pgfpathlineto{\pgfqpoint{2.267106in}{1.532084in}}%
\pgfpathlineto{\pgfqpoint{2.267116in}{1.532145in}}%
\pgfpathlineto{\pgfqpoint{2.267210in}{1.956558in}}%
\pgfpathlineto{\pgfqpoint{2.268283in}{1.926984in}}%
\pgfpathlineto{\pgfqpoint{2.268299in}{1.936026in}}%
\pgfpathlineto{\pgfqpoint{2.268394in}{1.925791in}}%
\pgfpathlineto{\pgfqpoint{2.269034in}{1.930472in}}%
\pgfpathlineto{\pgfqpoint{2.269232in}{1.935417in}}%
\pgfpathlineto{\pgfqpoint{2.269628in}{1.914235in}}%
\pgfpathlineto{\pgfqpoint{2.270101in}{1.532026in}}%
\pgfpathlineto{\pgfqpoint{2.271529in}{1.532057in}}%
\pgfpathlineto{\pgfqpoint{2.272061in}{1.532138in}}%
\pgfpathlineto{\pgfqpoint{2.272398in}{1.968519in}}%
\pgfpathlineto{\pgfqpoint{2.273259in}{1.924248in}}%
\pgfpathlineto{\pgfqpoint{2.273265in}{1.918659in}}%
\pgfpathlineto{\pgfqpoint{2.273280in}{1.942822in}}%
\pgfpathlineto{\pgfqpoint{2.274020in}{1.934496in}}%
\pgfpathlineto{\pgfqpoint{2.275087in}{1.532032in}}%
\pgfpathlineto{\pgfqpoint{2.274052in}{1.935277in}}%
\pgfpathlineto{\pgfqpoint{2.275278in}{1.532036in}}%
\pgfpathlineto{\pgfqpoint{2.277005in}{1.532124in}}%
\pgfpathlineto{\pgfqpoint{2.277095in}{1.957098in}}%
\pgfpathlineto{\pgfqpoint{2.278469in}{1.930474in}}%
\pgfpathlineto{\pgfqpoint{2.278716in}{1.935867in}}%
\pgfpathlineto{\pgfqpoint{2.279112in}{1.925636in}}%
\pgfpathlineto{\pgfqpoint{2.279161in}{1.935407in}}%
\pgfpathlineto{\pgfqpoint{2.279705in}{1.742922in}}%
\pgfpathlineto{\pgfqpoint{2.280016in}{1.532028in}}%
\pgfpathlineto{\pgfqpoint{2.280596in}{1.532041in}}%
\pgfpathlineto{\pgfqpoint{2.281952in}{1.532144in}}%
\pgfpathlineto{\pgfqpoint{2.282046in}{1.956350in}}%
\pgfpathlineto{\pgfqpoint{2.283130in}{1.932452in}}%
\pgfpathlineto{\pgfqpoint{2.283140in}{1.936316in}}%
\pgfpathlineto{\pgfqpoint{2.283158in}{1.925539in}}%
\pgfpathlineto{\pgfqpoint{2.283851in}{1.935438in}}%
\pgfpathlineto{\pgfqpoint{2.284692in}{1.689390in}}%
\pgfpathlineto{\pgfqpoint{2.284964in}{1.532030in}}%
\pgfpathlineto{\pgfqpoint{2.285595in}{1.532048in}}%
\pgfpathlineto{\pgfqpoint{2.286894in}{1.532109in}}%
\pgfpathlineto{\pgfqpoint{2.286937in}{1.729639in}}%
\pgfpathlineto{\pgfqpoint{2.287195in}{1.957334in}}%
\pgfpathlineto{\pgfqpoint{2.287721in}{1.922650in}}%
\pgfpathlineto{\pgfqpoint{2.287732in}{1.945253in}}%
\pgfpathlineto{\pgfqpoint{2.287748in}{1.916772in}}%
\pgfpathlineto{\pgfqpoint{2.288478in}{1.937290in}}%
\pgfpathlineto{\pgfqpoint{2.289146in}{1.924172in}}%
\pgfpathlineto{\pgfqpoint{2.288651in}{1.937763in}}%
\pgfpathlineto{\pgfqpoint{2.289245in}{1.926961in}}%
\pgfpathlineto{\pgfqpoint{2.289294in}{1.931257in}}%
\pgfpathlineto{\pgfqpoint{2.289492in}{1.886127in}}%
\pgfpathlineto{\pgfqpoint{2.289904in}{1.532029in}}%
\pgfpathlineto{\pgfqpoint{2.290530in}{1.532049in}}%
\pgfpathlineto{\pgfqpoint{2.291841in}{1.532120in}}%
\pgfpathlineto{\pgfqpoint{2.292126in}{1.957436in}}%
\pgfpathlineto{\pgfqpoint{2.293363in}{1.928475in}}%
\pgfpathlineto{\pgfqpoint{2.293371in}{1.925172in}}%
\pgfpathlineto{\pgfqpoint{2.293390in}{1.936218in}}%
\pgfpathlineto{\pgfqpoint{2.294116in}{1.933637in}}%
\pgfpathlineto{\pgfqpoint{2.294215in}{1.935014in}}%
\pgfpathlineto{\pgfqpoint{2.294512in}{1.774562in}}%
\pgfpathlineto{\pgfqpoint{2.294828in}{1.532025in}}%
\pgfpathlineto{\pgfqpoint{2.295465in}{1.532039in}}%
\pgfpathlineto{\pgfqpoint{2.296788in}{1.532135in}}%
\pgfpathlineto{\pgfqpoint{2.297126in}{1.966588in}}%
\pgfpathlineto{\pgfqpoint{2.298010in}{1.924988in}}%
\pgfpathlineto{\pgfqpoint{2.298015in}{1.919870in}}%
\pgfpathlineto{\pgfqpoint{2.298031in}{1.941547in}}%
\pgfpathlineto{\pgfqpoint{2.298776in}{1.928411in}}%
\pgfpathlineto{\pgfqpoint{2.298937in}{1.934432in}}%
\pgfpathlineto{\pgfqpoint{2.299492in}{1.730407in}}%
\pgfpathlineto{\pgfqpoint{2.299809in}{1.532029in}}%
\pgfpathlineto{\pgfqpoint{2.300395in}{1.532042in}}%
\pgfpathlineto{\pgfqpoint{2.301734in}{1.532137in}}%
\pgfpathlineto{\pgfqpoint{2.301920in}{1.960276in}}%
\pgfpathlineto{\pgfqpoint{2.302944in}{1.934989in}}%
\pgfpathlineto{\pgfqpoint{2.302990in}{1.923328in}}%
\pgfpathlineto{\pgfqpoint{2.303004in}{1.937729in}}%
\pgfpathlineto{\pgfqpoint{2.303706in}{1.934922in}}%
\pgfpathlineto{\pgfqpoint{2.304448in}{1.721942in}}%
\pgfpathlineto{\pgfqpoint{2.304755in}{1.532029in}}%
\pgfpathlineto{\pgfqpoint{2.305341in}{1.532042in}}%
\pgfpathlineto{\pgfqpoint{2.306679in}{1.532138in}}%
\pgfpathlineto{\pgfqpoint{2.306886in}{1.960990in}}%
\pgfpathlineto{\pgfqpoint{2.307883in}{1.937092in}}%
\pgfpathlineto{\pgfqpoint{2.307902in}{1.924770in}}%
\pgfpathlineto{\pgfqpoint{2.308645in}{1.931547in}}%
\pgfpathlineto{\pgfqpoint{2.308744in}{1.928365in}}%
\pgfpathlineto{\pgfqpoint{2.308892in}{1.934958in}}%
\pgfpathlineto{\pgfqpoint{2.309238in}{1.904828in}}%
\pgfpathlineto{\pgfqpoint{2.309705in}{1.532028in}}%
\pgfpathlineto{\pgfqpoint{2.310390in}{1.532043in}}%
\pgfpathlineto{\pgfqpoint{2.311625in}{1.532141in}}%
\pgfpathlineto{\pgfqpoint{2.311795in}{1.956717in}}%
\pgfpathlineto{\pgfqpoint{2.312801in}{1.926100in}}%
\pgfpathlineto{\pgfqpoint{2.312839in}{1.923650in}}%
\pgfpathlineto{\pgfqpoint{2.312824in}{1.938174in}}%
\pgfpathlineto{\pgfqpoint{2.313500in}{1.926562in}}%
\pgfpathlineto{\pgfqpoint{2.313896in}{1.935025in}}%
\pgfpathlineto{\pgfqpoint{2.313846in}{1.925618in}}%
\pgfpathlineto{\pgfqpoint{2.314143in}{1.931607in}}%
\pgfpathlineto{\pgfqpoint{2.314609in}{1.532026in}}%
\pgfpathlineto{\pgfqpoint{2.315392in}{1.532055in}}%
\pgfpathlineto{\pgfqpoint{2.316560in}{1.532084in}}%
\pgfpathlineto{\pgfqpoint{2.316572in}{1.532164in}}%
\pgfpathlineto{\pgfqpoint{2.316674in}{1.957357in}}%
\pgfpathlineto{\pgfqpoint{2.317693in}{1.931658in}}%
\pgfpathlineto{\pgfqpoint{2.317821in}{1.937515in}}%
\pgfpathlineto{\pgfqpoint{2.317726in}{1.924125in}}%
\pgfpathlineto{\pgfqpoint{2.318453in}{1.929976in}}%
\pgfpathlineto{\pgfqpoint{2.318700in}{1.936652in}}%
\pgfpathlineto{\pgfqpoint{2.319095in}{1.922955in}}%
\pgfpathlineto{\pgfqpoint{2.319576in}{1.532029in}}%
\pgfpathlineto{\pgfqpoint{2.320843in}{1.532057in}}%
\pgfpathlineto{\pgfqpoint{2.321514in}{1.532120in}}%
\pgfpathlineto{\pgfqpoint{2.321615in}{1.956821in}}%
\pgfpathlineto{\pgfqpoint{2.323038in}{1.928428in}}%
\pgfpathlineto{\pgfqpoint{2.323218in}{1.925933in}}%
\pgfpathlineto{\pgfqpoint{2.323575in}{1.935306in}}%
\pgfpathlineto{\pgfqpoint{2.323773in}{1.931427in}}%
\pgfpathlineto{\pgfqpoint{2.323971in}{1.925644in}}%
\pgfpathlineto{\pgfqpoint{2.324020in}{1.934026in}}%
\pgfpathlineto{\pgfqpoint{2.324267in}{1.667174in}}%
\pgfpathlineto{\pgfqpoint{2.324518in}{1.532029in}}%
\pgfpathlineto{\pgfqpoint{2.325188in}{1.532044in}}%
\pgfpathlineto{\pgfqpoint{2.326460in}{1.532129in}}%
\pgfpathlineto{\pgfqpoint{2.326778in}{1.960742in}}%
\pgfpathlineto{\pgfqpoint{2.327772in}{1.936913in}}%
\pgfpathlineto{\pgfqpoint{2.327785in}{1.924245in}}%
\pgfpathlineto{\pgfqpoint{2.328546in}{1.932534in}}%
\pgfpathlineto{\pgfqpoint{2.329178in}{1.715136in}}%
\pgfpathlineto{\pgfqpoint{2.329469in}{1.532029in}}%
\pgfpathlineto{\pgfqpoint{2.330095in}{1.532049in}}%
\pgfpathlineto{\pgfqpoint{2.331405in}{1.532120in}}%
\pgfpathlineto{\pgfqpoint{2.331495in}{1.955162in}}%
\pgfpathlineto{\pgfqpoint{2.332941in}{1.927567in}}%
\pgfpathlineto{\pgfqpoint{2.333037in}{1.935898in}}%
\pgfpathlineto{\pgfqpoint{2.333633in}{1.925906in}}%
\pgfpathlineto{\pgfqpoint{2.333682in}{1.935130in}}%
\pgfpathlineto{\pgfqpoint{2.333979in}{1.890799in}}%
\pgfpathlineto{\pgfqpoint{2.334391in}{1.532027in}}%
\pgfpathlineto{\pgfqpoint{2.335123in}{1.532042in}}%
\pgfpathlineto{\pgfqpoint{2.336352in}{1.532145in}}%
\pgfpathlineto{\pgfqpoint{2.336446in}{1.956009in}}%
\pgfpathlineto{\pgfqpoint{2.337510in}{1.927479in}}%
\pgfpathlineto{\pgfqpoint{2.337657in}{1.936117in}}%
\pgfpathlineto{\pgfqpoint{2.337607in}{1.925467in}}%
\pgfpathlineto{\pgfqpoint{2.338300in}{1.930568in}}%
\pgfpathlineto{\pgfqpoint{2.338547in}{1.935418in}}%
\pgfpathlineto{\pgfqpoint{2.338497in}{1.925440in}}%
\pgfpathlineto{\pgfqpoint{2.338794in}{1.931966in}}%
\pgfpathlineto{\pgfqpoint{2.339041in}{1.752108in}}%
\pgfpathlineto{\pgfqpoint{2.339350in}{1.532030in}}%
\pgfpathlineto{\pgfqpoint{2.339967in}{1.532048in}}%
\pgfpathlineto{\pgfqpoint{2.341296in}{1.532121in}}%
\pgfpathlineto{\pgfqpoint{2.341508in}{1.959537in}}%
\pgfpathlineto{\pgfqpoint{2.342664in}{1.933200in}}%
\pgfpathlineto{\pgfqpoint{2.342673in}{1.925541in}}%
\pgfpathlineto{\pgfqpoint{2.342690in}{1.935632in}}%
\pgfpathlineto{\pgfqpoint{2.343437in}{1.930434in}}%
\pgfpathlineto{\pgfqpoint{2.343445in}{1.932158in}}%
\pgfpathlineto{\pgfqpoint{2.343865in}{1.910398in}}%
\pgfpathlineto{\pgfqpoint{2.344300in}{1.532029in}}%
\pgfpathlineto{\pgfqpoint{2.344971in}{1.532043in}}%
\pgfpathlineto{\pgfqpoint{2.346242in}{1.532131in}}%
\pgfpathlineto{\pgfqpoint{2.346456in}{1.959885in}}%
\pgfpathlineto{\pgfqpoint{2.347533in}{1.927688in}}%
\pgfpathlineto{\pgfqpoint{2.347538in}{1.924366in}}%
\pgfpathlineto{\pgfqpoint{2.347553in}{1.936976in}}%
\pgfpathlineto{\pgfqpoint{2.348292in}{1.931194in}}%
\pgfpathlineto{\pgfqpoint{2.348301in}{1.932733in}}%
\pgfpathlineto{\pgfqpoint{2.348801in}{1.920996in}}%
\pgfpathlineto{\pgfqpoint{2.349238in}{1.532030in}}%
\pgfpathlineto{\pgfqpoint{2.349950in}{1.532050in}}%
\pgfpathlineto{\pgfqpoint{2.351186in}{1.532113in}}%
\pgfpathlineto{\pgfqpoint{2.351228in}{1.741143in}}%
\pgfpathlineto{\pgfqpoint{2.351499in}{1.955723in}}%
\pgfpathlineto{\pgfqpoint{2.352011in}{1.945139in}}%
\pgfpathlineto{\pgfqpoint{2.352027in}{1.916475in}}%
\pgfpathlineto{\pgfqpoint{2.352782in}{1.935416in}}%
\pgfpathlineto{\pgfqpoint{2.353192in}{1.924597in}}%
\pgfpathlineto{\pgfqpoint{2.352831in}{1.936722in}}%
\pgfpathlineto{\pgfqpoint{2.353538in}{1.925014in}}%
\pgfpathlineto{\pgfqpoint{2.353587in}{1.935901in}}%
\pgfpathlineto{\pgfqpoint{2.353785in}{1.879384in}}%
\pgfpathlineto{\pgfqpoint{2.354207in}{1.532029in}}%
\pgfpathlineto{\pgfqpoint{2.354842in}{1.532049in}}%
\pgfpathlineto{\pgfqpoint{2.356133in}{1.532129in}}%
\pgfpathlineto{\pgfqpoint{2.356472in}{1.967027in}}%
\pgfpathlineto{\pgfqpoint{2.357381in}{1.935308in}}%
\pgfpathlineto{\pgfqpoint{2.357391in}{1.920156in}}%
\pgfpathlineto{\pgfqpoint{2.357406in}{1.941019in}}%
\pgfpathlineto{\pgfqpoint{2.358151in}{1.928475in}}%
\pgfpathlineto{\pgfqpoint{2.358178in}{1.934025in}}%
\pgfpathlineto{\pgfqpoint{2.358921in}{1.635680in}}%
\pgfpathlineto{\pgfqpoint{2.359158in}{1.532028in}}%
\pgfpathlineto{\pgfqpoint{2.359082in}{1.636600in}}%
\pgfpathlineto{\pgfqpoint{2.359843in}{1.532043in}}%
\pgfpathlineto{\pgfqpoint{2.361080in}{1.532156in}}%
\pgfpathlineto{\pgfqpoint{2.361174in}{1.956960in}}%
\pgfpathlineto{\pgfqpoint{2.362225in}{1.936749in}}%
\pgfpathlineto{\pgfqpoint{2.362361in}{1.924759in}}%
\pgfpathlineto{\pgfqpoint{2.362952in}{1.925122in}}%
\pgfpathlineto{\pgfqpoint{2.363348in}{1.935913in}}%
\pgfpathlineto{\pgfqpoint{2.363397in}{1.924988in}}%
\pgfpathlineto{\pgfqpoint{2.363595in}{1.930500in}}%
\pgfpathlineto{\pgfqpoint{2.364064in}{1.532026in}}%
\pgfpathlineto{\pgfqpoint{2.364848in}{1.532055in}}%
\pgfpathlineto{\pgfqpoint{2.366014in}{1.532083in}}%
\pgfpathlineto{\pgfqpoint{2.366026in}{1.532157in}}%
\pgfpathlineto{\pgfqpoint{2.366129in}{1.957325in}}%
\pgfpathlineto{\pgfqpoint{2.367143in}{1.933452in}}%
\pgfpathlineto{\pgfqpoint{2.367417in}{1.924602in}}%
\pgfpathlineto{\pgfqpoint{2.367368in}{1.937165in}}%
\pgfpathlineto{\pgfqpoint{2.367862in}{1.924768in}}%
\pgfpathlineto{\pgfqpoint{2.368258in}{1.936408in}}%
\pgfpathlineto{\pgfqpoint{2.368208in}{1.924606in}}%
\pgfpathlineto{\pgfqpoint{2.368505in}{1.931175in}}%
\pgfpathlineto{\pgfqpoint{2.369010in}{1.532025in}}%
\pgfpathlineto{\pgfqpoint{2.369845in}{1.532043in}}%
\pgfpathlineto{\pgfqpoint{2.370970in}{1.532132in}}%
\pgfpathlineto{\pgfqpoint{2.371309in}{1.965901in}}%
\pgfpathlineto{\pgfqpoint{2.372214in}{1.935161in}}%
\pgfpathlineto{\pgfqpoint{2.372224in}{1.920587in}}%
\pgfpathlineto{\pgfqpoint{2.372240in}{1.940534in}}%
\pgfpathlineto{\pgfqpoint{2.372985in}{1.928542in}}%
\pgfpathlineto{\pgfqpoint{2.373086in}{1.933915in}}%
\pgfpathlineto{\pgfqpoint{2.373700in}{1.692526in}}%
\pgfpathlineto{\pgfqpoint{2.373998in}{1.532027in}}%
\pgfpathlineto{\pgfqpoint{2.374585in}{1.532040in}}%
\pgfpathlineto{\pgfqpoint{2.375916in}{1.532145in}}%
\pgfpathlineto{\pgfqpoint{2.376006in}{1.956104in}}%
\pgfpathlineto{\pgfqpoint{2.377086in}{1.929442in}}%
\pgfpathlineto{\pgfqpoint{2.377095in}{1.923826in}}%
\pgfpathlineto{\pgfqpoint{2.377149in}{1.937119in}}%
\pgfpathlineto{\pgfqpoint{2.377825in}{1.926517in}}%
\pgfpathlineto{\pgfqpoint{2.378220in}{1.934601in}}%
\pgfpathlineto{\pgfqpoint{2.378418in}{1.917975in}}%
\pgfpathlineto{\pgfqpoint{2.378468in}{1.920940in}}%
\pgfpathlineto{\pgfqpoint{2.378907in}{1.532031in}}%
\pgfpathlineto{\pgfqpoint{2.379761in}{1.532048in}}%
\pgfpathlineto{\pgfqpoint{2.380865in}{1.532305in}}%
\pgfpathlineto{\pgfqpoint{2.381058in}{1.956364in}}%
\pgfpathlineto{\pgfqpoint{2.381923in}{1.936583in}}%
\pgfpathlineto{\pgfqpoint{2.381937in}{1.922556in}}%
\pgfpathlineto{\pgfqpoint{2.381951in}{1.939115in}}%
\pgfpathlineto{\pgfqpoint{2.382671in}{1.935778in}}%
\pgfpathlineto{\pgfqpoint{2.383823in}{1.542252in}}%
\pgfpathlineto{\pgfqpoint{2.383882in}{1.532029in}}%
\pgfpathlineto{\pgfqpoint{2.384568in}{1.532045in}}%
\pgfpathlineto{\pgfqpoint{2.385804in}{1.532114in}}%
\pgfpathlineto{\pgfqpoint{2.385847in}{1.753096in}}%
\pgfpathlineto{\pgfqpoint{2.385908in}{1.956881in}}%
\pgfpathlineto{\pgfqpoint{2.386630in}{1.942512in}}%
\pgfpathlineto{\pgfqpoint{2.386674in}{1.917715in}}%
\pgfpathlineto{\pgfqpoint{2.386690in}{1.943227in}}%
\pgfpathlineto{\pgfqpoint{2.387392in}{1.935366in}}%
\pgfpathlineto{\pgfqpoint{2.387410in}{1.924599in}}%
\pgfpathlineto{\pgfqpoint{2.387508in}{1.936659in}}%
\pgfpathlineto{\pgfqpoint{2.388150in}{1.931291in}}%
\pgfpathlineto{\pgfqpoint{2.388249in}{1.926862in}}%
\pgfpathlineto{\pgfqpoint{2.388299in}{1.933901in}}%
\pgfpathlineto{\pgfqpoint{2.388834in}{1.532027in}}%
\pgfpathlineto{\pgfqpoint{2.389816in}{1.532049in}}%
\pgfpathlineto{\pgfqpoint{2.390752in}{1.532145in}}%
\pgfpathlineto{\pgfqpoint{2.390941in}{1.956510in}}%
\pgfpathlineto{\pgfqpoint{2.391915in}{1.927897in}}%
\pgfpathlineto{\pgfqpoint{2.391922in}{1.923273in}}%
\pgfpathlineto{\pgfqpoint{2.391936in}{1.937541in}}%
\pgfpathlineto{\pgfqpoint{2.392650in}{1.926058in}}%
\pgfpathlineto{\pgfqpoint{2.392700in}{1.934804in}}%
\pgfpathlineto{\pgfqpoint{2.393095in}{1.925773in}}%
\pgfpathlineto{\pgfqpoint{2.393293in}{1.927585in}}%
\pgfpathlineto{\pgfqpoint{2.393760in}{1.532029in}}%
\pgfpathlineto{\pgfqpoint{2.394536in}{1.532046in}}%
\pgfpathlineto{\pgfqpoint{2.395696in}{1.532120in}}%
\pgfpathlineto{\pgfqpoint{2.395987in}{1.957852in}}%
\pgfpathlineto{\pgfqpoint{2.397238in}{1.928617in}}%
\pgfpathlineto{\pgfqpoint{2.397366in}{1.925272in}}%
\pgfpathlineto{\pgfqpoint{2.397464in}{1.935857in}}%
\pgfpathlineto{\pgfqpoint{2.397958in}{1.925286in}}%
\pgfpathlineto{\pgfqpoint{2.398008in}{1.934375in}}%
\pgfpathlineto{\pgfqpoint{2.398304in}{1.880640in}}%
\pgfpathlineto{\pgfqpoint{2.398751in}{1.532031in}}%
\pgfpathlineto{\pgfqpoint{2.399340in}{1.532044in}}%
\pgfpathlineto{\pgfqpoint{2.400647in}{1.532304in}}%
\pgfpathlineto{\pgfqpoint{2.400743in}{1.956083in}}%
\pgfpathlineto{\pgfqpoint{2.401694in}{1.927355in}}%
\pgfpathlineto{\pgfqpoint{2.401700in}{1.921838in}}%
\pgfpathlineto{\pgfqpoint{2.401783in}{1.938513in}}%
\pgfpathlineto{\pgfqpoint{2.402434in}{1.935031in}}%
\pgfpathlineto{\pgfqpoint{2.402776in}{1.935063in}}%
\pgfpathlineto{\pgfqpoint{2.403419in}{1.648842in}}%
\pgfpathlineto{\pgfqpoint{2.403627in}{1.532026in}}%
\pgfpathlineto{\pgfqpoint{2.404361in}{1.532042in}}%
\pgfpathlineto{\pgfqpoint{2.405590in}{1.532171in}}%
\pgfpathlineto{\pgfqpoint{2.405692in}{1.957313in}}%
\pgfpathlineto{\pgfqpoint{2.406681in}{1.929208in}}%
\pgfpathlineto{\pgfqpoint{2.406813in}{1.936991in}}%
\pgfpathlineto{\pgfqpoint{2.406862in}{1.924365in}}%
\pgfpathlineto{\pgfqpoint{2.407456in}{1.932098in}}%
\pgfpathlineto{\pgfqpoint{2.407703in}{1.936447in}}%
\pgfpathlineto{\pgfqpoint{2.408148in}{1.919415in}}%
\pgfpathlineto{\pgfqpoint{2.408614in}{1.532032in}}%
\pgfpathlineto{\pgfqpoint{2.409348in}{1.532049in}}%
\pgfpathlineto{\pgfqpoint{2.410534in}{1.532143in}}%
\pgfpathlineto{\pgfqpoint{2.410847in}{1.961318in}}%
\pgfpathlineto{\pgfqpoint{2.411760in}{1.925793in}}%
\pgfpathlineto{\pgfqpoint{2.411764in}{1.924401in}}%
\pgfpathlineto{\pgfqpoint{2.411776in}{1.936365in}}%
\pgfpathlineto{\pgfqpoint{2.412522in}{1.928767in}}%
\pgfpathlineto{\pgfqpoint{2.412542in}{1.932205in}}%
\pgfpathlineto{\pgfqpoint{2.413044in}{1.928392in}}%
\pgfpathlineto{\pgfqpoint{2.413488in}{1.569290in}}%
\pgfpathlineto{\pgfqpoint{2.413543in}{1.532029in}}%
\pgfpathlineto{\pgfqpoint{2.414220in}{1.532043in}}%
\pgfpathlineto{\pgfqpoint{2.415480in}{1.532142in}}%
\pgfpathlineto{\pgfqpoint{2.415570in}{1.955844in}}%
\pgfpathlineto{\pgfqpoint{2.416659in}{1.937566in}}%
\pgfpathlineto{\pgfqpoint{2.416675in}{1.923744in}}%
\pgfpathlineto{\pgfqpoint{2.417425in}{1.934901in}}%
\pgfpathlineto{\pgfqpoint{2.418265in}{1.640677in}}%
\pgfpathlineto{\pgfqpoint{2.418500in}{1.532029in}}%
\pgfpathlineto{\pgfqpoint{2.419186in}{1.532044in}}%
\pgfpathlineto{\pgfqpoint{2.420425in}{1.532139in}}%
\pgfpathlineto{\pgfqpoint{2.420623in}{1.959679in}}%
\pgfpathlineto{\pgfqpoint{2.421617in}{1.925876in}}%
\pgfpathlineto{\pgfqpoint{2.421654in}{1.924103in}}%
\pgfpathlineto{\pgfqpoint{2.421639in}{1.937156in}}%
\pgfpathlineto{\pgfqpoint{2.422330in}{1.926062in}}%
\pgfpathlineto{\pgfqpoint{2.422379in}{1.934779in}}%
\pgfpathlineto{\pgfqpoint{2.422923in}{1.924863in}}%
\pgfpathlineto{\pgfqpoint{2.422973in}{1.926617in}}%
\pgfpathlineto{\pgfqpoint{2.423417in}{1.532031in}}%
\pgfpathlineto{\pgfqpoint{2.424224in}{1.532052in}}%
\pgfpathlineto{\pgfqpoint{2.425365in}{1.532100in}}%
\pgfpathlineto{\pgfqpoint{2.425374in}{1.532306in}}%
\pgfpathlineto{\pgfqpoint{2.425470in}{1.956099in}}%
\pgfpathlineto{\pgfqpoint{2.426425in}{1.923259in}}%
\pgfpathlineto{\pgfqpoint{2.426445in}{1.939131in}}%
\pgfpathlineto{\pgfqpoint{2.426494in}{1.922411in}}%
\pgfpathlineto{\pgfqpoint{2.427215in}{1.935494in}}%
\pgfpathlineto{\pgfqpoint{2.428325in}{1.569333in}}%
\pgfpathlineto{\pgfqpoint{2.428378in}{1.532029in}}%
\pgfpathlineto{\pgfqpoint{2.429054in}{1.532043in}}%
\pgfpathlineto{\pgfqpoint{2.430314in}{1.532120in}}%
\pgfpathlineto{\pgfqpoint{2.430606in}{1.957099in}}%
\pgfpathlineto{\pgfqpoint{2.431853in}{1.928547in}}%
\pgfpathlineto{\pgfqpoint{2.431981in}{1.925194in}}%
\pgfpathlineto{\pgfqpoint{2.432078in}{1.935835in}}%
\pgfpathlineto{\pgfqpoint{2.432572in}{1.925224in}}%
\pgfpathlineto{\pgfqpoint{2.432622in}{1.934279in}}%
\pgfpathlineto{\pgfqpoint{2.432919in}{1.879753in}}%
\pgfpathlineto{\pgfqpoint{2.433316in}{1.532029in}}%
\pgfpathlineto{\pgfqpoint{2.433936in}{1.532048in}}%
\pgfpathlineto{\pgfqpoint{2.435260in}{1.532126in}}%
\pgfpathlineto{\pgfqpoint{2.435470in}{1.959390in}}%
\pgfpathlineto{\pgfqpoint{2.436572in}{1.927875in}}%
\pgfpathlineto{\pgfqpoint{2.436577in}{1.924810in}}%
\pgfpathlineto{\pgfqpoint{2.436593in}{1.936205in}}%
\pgfpathlineto{\pgfqpoint{2.437333in}{1.932210in}}%
\pgfpathlineto{\pgfqpoint{2.438275in}{1.532028in}}%
\pgfpathlineto{\pgfqpoint{2.439501in}{1.532063in}}%
\pgfpathlineto{\pgfqpoint{2.440207in}{1.532151in}}%
\pgfpathlineto{\pgfqpoint{2.440301in}{1.956606in}}%
\pgfpathlineto{\pgfqpoint{2.441336in}{1.927273in}}%
\pgfpathlineto{\pgfqpoint{2.441484in}{1.935937in}}%
\pgfpathlineto{\pgfqpoint{2.441434in}{1.925261in}}%
\pgfpathlineto{\pgfqpoint{2.442127in}{1.930446in}}%
\pgfpathlineto{\pgfqpoint{2.442374in}{1.935287in}}%
\pgfpathlineto{\pgfqpoint{2.442325in}{1.925265in}}%
\pgfpathlineto{\pgfqpoint{2.442720in}{1.933819in}}%
\pgfpathlineto{\pgfqpoint{2.443216in}{1.532028in}}%
\pgfpathlineto{\pgfqpoint{2.444240in}{1.532058in}}%
\pgfpathlineto{\pgfqpoint{2.445152in}{1.532143in}}%
\pgfpathlineto{\pgfqpoint{2.445243in}{1.955266in}}%
\pgfpathlineto{\pgfqpoint{2.446334in}{1.928561in}}%
\pgfpathlineto{\pgfqpoint{2.446342in}{1.923916in}}%
\pgfpathlineto{\pgfqpoint{2.446361in}{1.936835in}}%
\pgfpathlineto{\pgfqpoint{2.447093in}{1.934047in}}%
\pgfpathlineto{\pgfqpoint{2.447192in}{1.935408in}}%
\pgfpathlineto{\pgfqpoint{2.447885in}{1.689978in}}%
\pgfpathlineto{\pgfqpoint{2.448170in}{1.532029in}}%
\pgfpathlineto{\pgfqpoint{2.448805in}{1.532049in}}%
\pgfpathlineto{\pgfqpoint{2.450097in}{1.532128in}}%
\pgfpathlineto{\pgfqpoint{2.450430in}{1.964265in}}%
\pgfpathlineto{\pgfqpoint{2.451358in}{1.934432in}}%
\pgfpathlineto{\pgfqpoint{2.451368in}{1.921163in}}%
\pgfpathlineto{\pgfqpoint{2.451383in}{1.939690in}}%
\pgfpathlineto{\pgfqpoint{2.452129in}{1.927967in}}%
\pgfpathlineto{\pgfqpoint{2.452199in}{1.933120in}}%
\pgfpathlineto{\pgfqpoint{2.452687in}{1.889107in}}%
\pgfpathlineto{\pgfqpoint{2.453107in}{1.532028in}}%
\pgfpathlineto{\pgfqpoint{2.453736in}{1.532050in}}%
\pgfpathlineto{\pgfqpoint{2.455043in}{1.532144in}}%
\pgfpathlineto{\pgfqpoint{2.455134in}{1.956242in}}%
\pgfpathlineto{\pgfqpoint{2.456227in}{1.937271in}}%
\pgfpathlineto{\pgfqpoint{2.456246in}{1.924249in}}%
\pgfpathlineto{\pgfqpoint{2.456974in}{1.930319in}}%
\pgfpathlineto{\pgfqpoint{2.457171in}{1.935342in}}%
\pgfpathlineto{\pgfqpoint{2.457617in}{1.898333in}}%
\pgfpathlineto{\pgfqpoint{2.458063in}{1.532029in}}%
\pgfpathlineto{\pgfqpoint{2.458749in}{1.532044in}}%
\pgfpathlineto{\pgfqpoint{2.459988in}{1.532137in}}%
\pgfpathlineto{\pgfqpoint{2.460320in}{1.962512in}}%
\pgfpathlineto{\pgfqpoint{2.461205in}{1.936065in}}%
\pgfpathlineto{\pgfqpoint{2.461216in}{1.920526in}}%
\pgfpathlineto{\pgfqpoint{2.461232in}{1.940143in}}%
\pgfpathlineto{\pgfqpoint{2.461957in}{1.935104in}}%
\pgfpathlineto{\pgfqpoint{2.462942in}{1.578468in}}%
\pgfpathlineto{\pgfqpoint{2.462988in}{1.532029in}}%
\pgfpathlineto{\pgfqpoint{2.463706in}{1.532051in}}%
\pgfpathlineto{\pgfqpoint{2.464933in}{1.532125in}}%
\pgfpathlineto{\pgfqpoint{2.465023in}{1.954915in}}%
\pgfpathlineto{\pgfqpoint{2.466268in}{1.932344in}}%
\pgfpathlineto{\pgfqpoint{2.466278in}{1.926854in}}%
\pgfpathlineto{\pgfqpoint{2.466293in}{1.934113in}}%
\pgfpathlineto{\pgfqpoint{2.467037in}{1.929042in}}%
\pgfpathlineto{\pgfqpoint{2.467055in}{1.931484in}}%
\pgfpathlineto{\pgfqpoint{2.467497in}{1.909409in}}%
\pgfpathlineto{\pgfqpoint{2.467941in}{1.532029in}}%
\pgfpathlineto{\pgfqpoint{2.468667in}{1.532052in}}%
\pgfpathlineto{\pgfqpoint{2.469878in}{1.532120in}}%
\pgfpathlineto{\pgfqpoint{2.470175in}{1.956946in}}%
\pgfpathlineto{\pgfqpoint{2.471422in}{1.936417in}}%
\pgfpathlineto{\pgfqpoint{2.471440in}{1.924811in}}%
\pgfpathlineto{\pgfqpoint{2.472175in}{1.934309in}}%
\pgfpathlineto{\pgfqpoint{2.472571in}{1.753524in}}%
\pgfpathlineto{\pgfqpoint{2.472886in}{1.532029in}}%
\pgfpathlineto{\pgfqpoint{2.473462in}{1.532041in}}%
\pgfpathlineto{\pgfqpoint{2.474823in}{1.532120in}}%
\pgfpathlineto{\pgfqpoint{2.475105in}{1.957051in}}%
\pgfpathlineto{\pgfqpoint{2.476365in}{1.935860in}}%
\pgfpathlineto{\pgfqpoint{2.476383in}{1.925293in}}%
\pgfpathlineto{\pgfqpoint{2.477130in}{1.934890in}}%
\pgfpathlineto{\pgfqpoint{2.477526in}{1.747899in}}%
\pgfpathlineto{\pgfqpoint{2.477848in}{1.532032in}}%
\pgfpathlineto{\pgfqpoint{2.478434in}{1.532044in}}%
\pgfpathlineto{\pgfqpoint{2.479770in}{1.532140in}}%
\pgfpathlineto{\pgfqpoint{2.480081in}{1.960108in}}%
\pgfpathlineto{\pgfqpoint{2.481015in}{1.935666in}}%
\pgfpathlineto{\pgfqpoint{2.481030in}{1.924950in}}%
\pgfpathlineto{\pgfqpoint{2.481043in}{1.936442in}}%
\pgfpathlineto{\pgfqpoint{2.481792in}{1.929448in}}%
\pgfpathlineto{\pgfqpoint{2.481803in}{1.932045in}}%
\pgfpathlineto{\pgfqpoint{2.482330in}{1.916271in}}%
\pgfpathlineto{\pgfqpoint{2.482774in}{1.532029in}}%
\pgfpathlineto{\pgfqpoint{2.483495in}{1.532051in}}%
\pgfpathlineto{\pgfqpoint{2.484715in}{1.532133in}}%
\pgfpathlineto{\pgfqpoint{2.484908in}{1.959168in}}%
\pgfpathlineto{\pgfqpoint{2.485994in}{1.928523in}}%
\pgfpathlineto{\pgfqpoint{2.485999in}{1.924670in}}%
\pgfpathlineto{\pgfqpoint{2.486015in}{1.936260in}}%
\pgfpathlineto{\pgfqpoint{2.486756in}{1.929861in}}%
\pgfpathlineto{\pgfqpoint{2.486784in}{1.932151in}}%
\pgfpathlineto{\pgfqpoint{2.487383in}{1.794474in}}%
\pgfpathlineto{\pgfqpoint{2.487741in}{1.532032in}}%
\pgfpathlineto{\pgfqpoint{2.488326in}{1.532044in}}%
\pgfpathlineto{\pgfqpoint{2.489660in}{1.532129in}}%
\pgfpathlineto{\pgfqpoint{2.489761in}{1.956966in}}%
\pgfpathlineto{\pgfqpoint{2.491040in}{1.930875in}}%
\pgfpathlineto{\pgfqpoint{2.491050in}{1.925731in}}%
\pgfpathlineto{\pgfqpoint{2.491068in}{1.935193in}}%
\pgfpathlineto{\pgfqpoint{2.491792in}{1.928722in}}%
\pgfpathlineto{\pgfqpoint{2.491941in}{1.933596in}}%
\pgfpathlineto{\pgfqpoint{2.491990in}{1.926615in}}%
\pgfpathlineto{\pgfqpoint{2.492188in}{1.929727in}}%
\pgfpathlineto{\pgfqpoint{2.492686in}{1.532028in}}%
\pgfpathlineto{\pgfqpoint{2.493470in}{1.532045in}}%
\pgfpathlineto{\pgfqpoint{2.494608in}{1.532161in}}%
\pgfpathlineto{\pgfqpoint{2.494711in}{1.957130in}}%
\pgfpathlineto{\pgfqpoint{2.495735in}{1.937344in}}%
\pgfpathlineto{\pgfqpoint{2.495754in}{1.923983in}}%
\pgfpathlineto{\pgfqpoint{2.496482in}{1.925355in}}%
\pgfpathlineto{\pgfqpoint{2.496878in}{1.935716in}}%
\pgfpathlineto{\pgfqpoint{2.496829in}{1.924309in}}%
\pgfpathlineto{\pgfqpoint{2.497125in}{1.931362in}}%
\pgfpathlineto{\pgfqpoint{2.497633in}{1.532028in}}%
\pgfpathlineto{\pgfqpoint{2.498417in}{1.532045in}}%
\pgfpathlineto{\pgfqpoint{2.499552in}{1.532141in}}%
\pgfpathlineto{\pgfqpoint{2.499742in}{1.956675in}}%
\pgfpathlineto{\pgfqpoint{2.500738in}{1.937497in}}%
\pgfpathlineto{\pgfqpoint{2.500824in}{1.923765in}}%
\pgfpathlineto{\pgfqpoint{2.501486in}{1.926194in}}%
\pgfpathlineto{\pgfqpoint{2.501535in}{1.934901in}}%
\pgfpathlineto{\pgfqpoint{2.502079in}{1.917127in}}%
\pgfpathlineto{\pgfqpoint{2.502573in}{1.532029in}}%
\pgfpathlineto{\pgfqpoint{2.503753in}{1.532055in}}%
\pgfpathlineto{\pgfqpoint{2.504497in}{1.532139in}}%
\pgfpathlineto{\pgfqpoint{2.504703in}{1.960070in}}%
\pgfpathlineto{\pgfqpoint{2.505698in}{1.936986in}}%
\pgfpathlineto{\pgfqpoint{2.505715in}{1.923914in}}%
\pgfpathlineto{\pgfqpoint{2.505732in}{1.937062in}}%
\pgfpathlineto{\pgfqpoint{2.506464in}{1.932658in}}%
\pgfpathlineto{\pgfqpoint{2.506711in}{1.934443in}}%
\pgfpathlineto{\pgfqpoint{2.507106in}{1.870135in}}%
\pgfpathlineto{\pgfqpoint{2.507507in}{1.532028in}}%
\pgfpathlineto{\pgfqpoint{2.508135in}{1.532049in}}%
\pgfpathlineto{\pgfqpoint{2.509443in}{1.532143in}}%
\pgfpathlineto{\pgfqpoint{2.509534in}{1.955835in}}%
\pgfpathlineto{\pgfqpoint{2.510621in}{1.927843in}}%
\pgfpathlineto{\pgfqpoint{2.510629in}{1.924113in}}%
\pgfpathlineto{\pgfqpoint{2.510647in}{1.936652in}}%
\pgfpathlineto{\pgfqpoint{2.511351in}{1.928473in}}%
\pgfpathlineto{\pgfqpoint{2.511648in}{1.934816in}}%
\pgfpathlineto{\pgfqpoint{2.511599in}{1.925537in}}%
\pgfpathlineto{\pgfqpoint{2.511895in}{1.931239in}}%
\pgfpathlineto{\pgfqpoint{2.512242in}{1.629025in}}%
\pgfpathlineto{\pgfqpoint{2.512450in}{1.532029in}}%
\pgfpathlineto{\pgfqpoint{2.512392in}{1.631181in}}%
\pgfpathlineto{\pgfqpoint{2.513175in}{1.532051in}}%
\pgfpathlineto{\pgfqpoint{2.514387in}{1.532120in}}%
\pgfpathlineto{\pgfqpoint{2.514671in}{1.957095in}}%
\pgfpathlineto{\pgfqpoint{2.515909in}{1.928019in}}%
\pgfpathlineto{\pgfqpoint{2.515917in}{1.924777in}}%
\pgfpathlineto{\pgfqpoint{2.515936in}{1.935796in}}%
\pgfpathlineto{\pgfqpoint{2.516660in}{1.933848in}}%
\pgfpathlineto{\pgfqpoint{2.517056in}{1.790822in}}%
\pgfpathlineto{\pgfqpoint{2.517408in}{1.532029in}}%
\pgfpathlineto{\pgfqpoint{2.517993in}{1.532042in}}%
\pgfpathlineto{\pgfqpoint{2.519333in}{1.532131in}}%
\pgfpathlineto{\pgfqpoint{2.519671in}{1.965969in}}%
\pgfpathlineto{\pgfqpoint{2.520577in}{1.940892in}}%
\pgfpathlineto{\pgfqpoint{2.520593in}{1.920405in}}%
\pgfpathlineto{\pgfqpoint{2.521347in}{1.933430in}}%
\pgfpathlineto{\pgfqpoint{2.521527in}{1.933922in}}%
\pgfpathlineto{\pgfqpoint{2.521972in}{1.832134in}}%
\pgfpathlineto{\pgfqpoint{2.522331in}{1.532030in}}%
\pgfpathlineto{\pgfqpoint{2.522897in}{1.532042in}}%
\pgfpathlineto{\pgfqpoint{2.524277in}{1.532119in}}%
\pgfpathlineto{\pgfqpoint{2.524367in}{1.956007in}}%
\pgfpathlineto{\pgfqpoint{2.525813in}{1.934193in}}%
\pgfpathlineto{\pgfqpoint{2.526481in}{1.926454in}}%
\pgfpathlineto{\pgfqpoint{2.525987in}{1.934286in}}%
\pgfpathlineto{\pgfqpoint{2.526580in}{1.927774in}}%
\pgfpathlineto{\pgfqpoint{2.526778in}{1.932025in}}%
\pgfpathlineto{\pgfqpoint{2.526828in}{1.911806in}}%
\pgfpathlineto{\pgfqpoint{2.527288in}{1.532028in}}%
\pgfpathlineto{\pgfqpoint{2.528212in}{1.532056in}}%
\pgfpathlineto{\pgfqpoint{2.529223in}{1.532120in}}%
\pgfpathlineto{\pgfqpoint{2.529519in}{1.956916in}}%
\pgfpathlineto{\pgfqpoint{2.530763in}{1.935485in}}%
\pgfpathlineto{\pgfqpoint{2.530837in}{1.936041in}}%
\pgfpathlineto{\pgfqpoint{2.531534in}{1.924743in}}%
\pgfpathlineto{\pgfqpoint{2.531583in}{1.934620in}}%
\pgfpathlineto{\pgfqpoint{2.531831in}{1.853552in}}%
\pgfpathlineto{\pgfqpoint{2.532243in}{1.532029in}}%
\pgfpathlineto{\pgfqpoint{2.532829in}{1.532042in}}%
\pgfpathlineto{\pgfqpoint{2.534168in}{1.532113in}}%
\pgfpathlineto{\pgfqpoint{2.534210in}{1.749512in}}%
\pgfpathlineto{\pgfqpoint{2.534456in}{1.957248in}}%
\pgfpathlineto{\pgfqpoint{2.534992in}{1.930861in}}%
\pgfpathlineto{\pgfqpoint{2.535002in}{1.917486in}}%
\pgfpathlineto{\pgfqpoint{2.535078in}{1.942960in}}%
\pgfpathlineto{\pgfqpoint{2.535757in}{1.927308in}}%
\pgfpathlineto{\pgfqpoint{2.535769in}{1.936441in}}%
\pgfpathlineto{\pgfqpoint{2.535914in}{1.924736in}}%
\pgfpathlineto{\pgfqpoint{2.536501in}{1.930643in}}%
\pgfpathlineto{\pgfqpoint{2.536600in}{1.926500in}}%
\pgfpathlineto{\pgfqpoint{2.536649in}{1.933958in}}%
\pgfpathlineto{\pgfqpoint{2.537186in}{1.532028in}}%
\pgfpathlineto{\pgfqpoint{2.538511in}{1.532066in}}%
\pgfpathlineto{\pgfqpoint{2.539116in}{1.532153in}}%
\pgfpathlineto{\pgfqpoint{2.539210in}{1.956299in}}%
\pgfpathlineto{\pgfqpoint{2.540266in}{1.929043in}}%
\pgfpathlineto{\pgfqpoint{2.540563in}{1.934173in}}%
\pgfpathlineto{\pgfqpoint{2.540513in}{1.926702in}}%
\pgfpathlineto{\pgfqpoint{2.540909in}{1.933663in}}%
\pgfpathlineto{\pgfqpoint{2.541008in}{1.933944in}}%
\pgfpathlineto{\pgfqpoint{2.541849in}{1.702042in}}%
\pgfpathlineto{\pgfqpoint{2.542142in}{1.532028in}}%
\pgfpathlineto{\pgfqpoint{2.542777in}{1.532051in}}%
\pgfpathlineto{\pgfqpoint{2.544061in}{1.532141in}}%
\pgfpathlineto{\pgfqpoint{2.544251in}{1.957830in}}%
\pgfpathlineto{\pgfqpoint{2.545248in}{1.935251in}}%
\pgfpathlineto{\pgfqpoint{2.545291in}{1.923619in}}%
\pgfpathlineto{\pgfqpoint{2.545277in}{1.937856in}}%
\pgfpathlineto{\pgfqpoint{2.546005in}{1.934472in}}%
\pgfpathlineto{\pgfqpoint{2.546104in}{1.934707in}}%
\pgfpathlineto{\pgfqpoint{2.546796in}{1.686069in}}%
\pgfpathlineto{\pgfqpoint{2.547084in}{1.532028in}}%
\pgfpathlineto{\pgfqpoint{2.547720in}{1.532050in}}%
\pgfpathlineto{\pgfqpoint{2.549007in}{1.532153in}}%
\pgfpathlineto{\pgfqpoint{2.549100in}{1.956187in}}%
\pgfpathlineto{\pgfqpoint{2.550138in}{1.926890in}}%
\pgfpathlineto{\pgfqpoint{2.550287in}{1.936413in}}%
\pgfpathlineto{\pgfqpoint{2.550237in}{1.924577in}}%
\pgfpathlineto{\pgfqpoint{2.550930in}{1.930439in}}%
\pgfpathlineto{\pgfqpoint{2.551177in}{1.935730in}}%
\pgfpathlineto{\pgfqpoint{2.551127in}{1.924596in}}%
\pgfpathlineto{\pgfqpoint{2.551523in}{1.934124in}}%
\pgfpathlineto{\pgfqpoint{2.552013in}{1.532029in}}%
\pgfpathlineto{\pgfqpoint{2.552985in}{1.532050in}}%
\pgfpathlineto{\pgfqpoint{2.553950in}{1.532120in}}%
\pgfpathlineto{\pgfqpoint{2.554052in}{1.956722in}}%
\pgfpathlineto{\pgfqpoint{2.555471in}{1.928690in}}%
\pgfpathlineto{\pgfqpoint{2.555481in}{1.924899in}}%
\pgfpathlineto{\pgfqpoint{2.555499in}{1.935822in}}%
\pgfpathlineto{\pgfqpoint{2.556237in}{1.924975in}}%
\pgfpathlineto{\pgfqpoint{2.556286in}{1.934689in}}%
\pgfpathlineto{\pgfqpoint{2.556534in}{1.882660in}}%
\pgfpathlineto{\pgfqpoint{2.556936in}{1.532026in}}%
\pgfpathlineto{\pgfqpoint{2.557620in}{1.532052in}}%
\pgfpathlineto{\pgfqpoint{2.558887in}{1.532084in}}%
\pgfpathlineto{\pgfqpoint{2.558898in}{1.532145in}}%
\pgfpathlineto{\pgfqpoint{2.558991in}{1.955943in}}%
\pgfpathlineto{\pgfqpoint{2.560017in}{1.932217in}}%
\pgfpathlineto{\pgfqpoint{2.560165in}{1.922656in}}%
\pgfpathlineto{\pgfqpoint{2.560215in}{1.938488in}}%
\pgfpathlineto{\pgfqpoint{2.560808in}{1.927211in}}%
\pgfpathlineto{\pgfqpoint{2.561105in}{1.937620in}}%
\pgfpathlineto{\pgfqpoint{2.561055in}{1.922689in}}%
\pgfpathlineto{\pgfqpoint{2.561352in}{1.932058in}}%
\pgfpathlineto{\pgfqpoint{2.561922in}{1.532028in}}%
\pgfpathlineto{\pgfqpoint{2.562806in}{1.532048in}}%
\pgfpathlineto{\pgfqpoint{2.563843in}{1.532139in}}%
\pgfpathlineto{\pgfqpoint{2.564042in}{1.959246in}}%
\pgfpathlineto{\pgfqpoint{2.565028in}{1.930288in}}%
\pgfpathlineto{\pgfqpoint{2.565036in}{1.923571in}}%
\pgfpathlineto{\pgfqpoint{2.565090in}{1.936929in}}%
\pgfpathlineto{\pgfqpoint{2.565774in}{1.934018in}}%
\pgfpathlineto{\pgfqpoint{2.566120in}{1.934306in}}%
\pgfpathlineto{\pgfqpoint{2.566516in}{1.769866in}}%
\pgfpathlineto{\pgfqpoint{2.566827in}{1.532026in}}%
\pgfpathlineto{\pgfqpoint{2.567462in}{1.532040in}}%
\pgfpathlineto{\pgfqpoint{2.568790in}{1.532167in}}%
\pgfpathlineto{\pgfqpoint{2.568892in}{1.956849in}}%
\pgfpathlineto{\pgfqpoint{2.569898in}{1.937304in}}%
\pgfpathlineto{\pgfqpoint{2.569915in}{1.923534in}}%
\pgfpathlineto{\pgfqpoint{2.570653in}{1.926820in}}%
\pgfpathlineto{\pgfqpoint{2.570802in}{1.936315in}}%
\pgfpathlineto{\pgfqpoint{2.571197in}{1.924005in}}%
\pgfpathlineto{\pgfqpoint{2.571296in}{1.924347in}}%
\pgfpathlineto{\pgfqpoint{2.571811in}{1.532028in}}%
\pgfpathlineto{\pgfqpoint{2.572892in}{1.532052in}}%
\pgfpathlineto{\pgfqpoint{2.573734in}{1.532148in}}%
\pgfpathlineto{\pgfqpoint{2.573828in}{1.955516in}}%
\pgfpathlineto{\pgfqpoint{2.574886in}{1.935524in}}%
\pgfpathlineto{\pgfqpoint{2.574935in}{1.924972in}}%
\pgfpathlineto{\pgfqpoint{2.575628in}{1.929059in}}%
\pgfpathlineto{\pgfqpoint{2.575776in}{1.935185in}}%
\pgfpathlineto{\pgfqpoint{2.576271in}{1.920496in}}%
\pgfpathlineto{\pgfqpoint{2.576756in}{1.532029in}}%
\pgfpathlineto{\pgfqpoint{2.577590in}{1.532055in}}%
\pgfpathlineto{\pgfqpoint{2.578680in}{1.532149in}}%
\pgfpathlineto{\pgfqpoint{2.578773in}{1.956151in}}%
\pgfpathlineto{\pgfqpoint{2.579818in}{1.925179in}}%
\pgfpathlineto{\pgfqpoint{2.580214in}{1.935866in}}%
\pgfpathlineto{\pgfqpoint{2.580560in}{1.934675in}}%
\pgfpathlineto{\pgfqpoint{2.580659in}{1.935558in}}%
\pgfpathlineto{\pgfqpoint{2.581632in}{1.593834in}}%
\pgfpathlineto{\pgfqpoint{2.581730in}{1.532031in}}%
\pgfpathlineto{\pgfqpoint{2.582367in}{1.532049in}}%
\pgfpathlineto{\pgfqpoint{2.583620in}{1.532101in}}%
\pgfpathlineto{\pgfqpoint{2.583629in}{1.532322in}}%
\pgfpathlineto{\pgfqpoint{2.583722in}{1.955117in}}%
\pgfpathlineto{\pgfqpoint{2.584687in}{1.937655in}}%
\pgfpathlineto{\pgfqpoint{2.584701in}{1.923524in}}%
\pgfpathlineto{\pgfqpoint{2.585436in}{1.925839in}}%
\pgfpathlineto{\pgfqpoint{2.585832in}{1.934516in}}%
\pgfpathlineto{\pgfqpoint{2.586129in}{1.914928in}}%
\pgfpathlineto{\pgfqpoint{2.586228in}{1.871723in}}%
\pgfpathlineto{\pgfqpoint{2.586645in}{1.532029in}}%
\pgfpathlineto{\pgfqpoint{2.587230in}{1.532042in}}%
\pgfpathlineto{\pgfqpoint{2.588570in}{1.532134in}}%
\pgfpathlineto{\pgfqpoint{2.588908in}{1.969458in}}%
\pgfpathlineto{\pgfqpoint{2.589779in}{1.924285in}}%
\pgfpathlineto{\pgfqpoint{2.589784in}{1.917555in}}%
\pgfpathlineto{\pgfqpoint{2.589799in}{1.942915in}}%
\pgfpathlineto{\pgfqpoint{2.590535in}{1.925567in}}%
\pgfpathlineto{\pgfqpoint{2.590554in}{1.934686in}}%
\pgfpathlineto{\pgfqpoint{2.591182in}{1.873640in}}%
\pgfpathlineto{\pgfqpoint{2.591594in}{1.532032in}}%
\pgfpathlineto{\pgfqpoint{2.592180in}{1.532044in}}%
\pgfpathlineto{\pgfqpoint{2.593515in}{1.532131in}}%
\pgfpathlineto{\pgfqpoint{2.593605in}{1.956331in}}%
\pgfpathlineto{\pgfqpoint{2.594796in}{1.931445in}}%
\pgfpathlineto{\pgfqpoint{2.594801in}{1.932553in}}%
\pgfpathlineto{\pgfqpoint{2.594817in}{1.928244in}}%
\pgfpathlineto{\pgfqpoint{2.595564in}{1.930950in}}%
\pgfpathlineto{\pgfqpoint{2.596186in}{1.789388in}}%
\pgfpathlineto{\pgfqpoint{2.596508in}{1.532031in}}%
\pgfpathlineto{\pgfqpoint{2.597116in}{1.532047in}}%
\pgfpathlineto{\pgfqpoint{2.598456in}{1.532100in}}%
\pgfpathlineto{\pgfqpoint{2.598465in}{1.532306in}}%
\pgfpathlineto{\pgfqpoint{2.598658in}{1.956037in}}%
\pgfpathlineto{\pgfqpoint{2.599523in}{1.937959in}}%
\pgfpathlineto{\pgfqpoint{2.599536in}{1.921503in}}%
\pgfpathlineto{\pgfqpoint{2.599587in}{1.938595in}}%
\pgfpathlineto{\pgfqpoint{2.600260in}{1.935373in}}%
\pgfpathlineto{\pgfqpoint{2.601421in}{1.545553in}}%
\pgfpathlineto{\pgfqpoint{2.601445in}{1.532027in}}%
\pgfpathlineto{\pgfqpoint{2.602177in}{1.532042in}}%
\pgfpathlineto{\pgfqpoint{2.603407in}{1.532145in}}%
\pgfpathlineto{\pgfqpoint{2.603500in}{1.955592in}}%
\pgfpathlineto{\pgfqpoint{2.604562in}{1.930162in}}%
\pgfpathlineto{\pgfqpoint{2.604793in}{1.925209in}}%
\pgfpathlineto{\pgfqpoint{2.604744in}{1.935815in}}%
\pgfpathlineto{\pgfqpoint{2.605238in}{1.925336in}}%
\pgfpathlineto{\pgfqpoint{2.605634in}{1.935123in}}%
\pgfpathlineto{\pgfqpoint{2.605585in}{1.925153in}}%
\pgfpathlineto{\pgfqpoint{2.605881in}{1.930666in}}%
\pgfpathlineto{\pgfqpoint{2.606391in}{1.532026in}}%
\pgfpathlineto{\pgfqpoint{2.607372in}{1.532059in}}%
\pgfpathlineto{\pgfqpoint{2.608341in}{1.532084in}}%
\pgfpathlineto{\pgfqpoint{2.608353in}{1.532165in}}%
\pgfpathlineto{\pgfqpoint{2.608456in}{1.956926in}}%
\pgfpathlineto{\pgfqpoint{2.609467in}{1.929423in}}%
\pgfpathlineto{\pgfqpoint{2.609596in}{1.923893in}}%
\pgfpathlineto{\pgfqpoint{2.609500in}{1.936696in}}%
\pgfpathlineto{\pgfqpoint{2.610237in}{1.929788in}}%
\pgfpathlineto{\pgfqpoint{2.610484in}{1.924057in}}%
\pgfpathlineto{\pgfqpoint{2.610435in}{1.936096in}}%
\pgfpathlineto{\pgfqpoint{2.610831in}{1.924081in}}%
\pgfpathlineto{\pgfqpoint{2.610880in}{1.934600in}}%
\pgfpathlineto{\pgfqpoint{2.610979in}{1.853568in}}%
\pgfpathlineto{\pgfqpoint{2.611378in}{1.532032in}}%
\pgfpathlineto{\pgfqpoint{2.611964in}{1.532044in}}%
\pgfpathlineto{\pgfqpoint{2.613297in}{1.532127in}}%
\pgfpathlineto{\pgfqpoint{2.613398in}{1.956494in}}%
\pgfpathlineto{\pgfqpoint{2.614726in}{1.935857in}}%
\pgfpathlineto{\pgfqpoint{2.614776in}{1.924656in}}%
\pgfpathlineto{\pgfqpoint{2.614759in}{1.936169in}}%
\pgfpathlineto{\pgfqpoint{2.615475in}{1.933707in}}%
\pgfpathlineto{\pgfqpoint{2.615574in}{1.934194in}}%
\pgfpathlineto{\pgfqpoint{2.615970in}{1.783956in}}%
\pgfpathlineto{\pgfqpoint{2.616323in}{1.532028in}}%
\pgfpathlineto{\pgfqpoint{2.616909in}{1.532041in}}%
\pgfpathlineto{\pgfqpoint{2.618243in}{1.532139in}}%
\pgfpathlineto{\pgfqpoint{2.618447in}{1.959238in}}%
\pgfpathlineto{\pgfqpoint{2.619430in}{1.937176in}}%
\pgfpathlineto{\pgfqpoint{2.619450in}{1.923306in}}%
\pgfpathlineto{\pgfqpoint{2.620206in}{1.933893in}}%
\pgfpathlineto{\pgfqpoint{2.620552in}{1.934341in}}%
\pgfpathlineto{\pgfqpoint{2.620898in}{1.806915in}}%
\pgfpathlineto{\pgfqpoint{2.621256in}{1.532028in}}%
\pgfpathlineto{\pgfqpoint{2.621838in}{1.532041in}}%
\pgfpathlineto{\pgfqpoint{2.623189in}{1.532149in}}%
\pgfpathlineto{\pgfqpoint{2.623282in}{1.956398in}}%
\pgfpathlineto{\pgfqpoint{2.624341in}{1.934655in}}%
\pgfpathlineto{\pgfqpoint{2.624390in}{1.925152in}}%
\pgfpathlineto{\pgfqpoint{2.624440in}{1.935742in}}%
\pgfpathlineto{\pgfqpoint{2.625083in}{1.930613in}}%
\pgfpathlineto{\pgfqpoint{2.625330in}{1.935146in}}%
\pgfpathlineto{\pgfqpoint{2.625281in}{1.925147in}}%
\pgfpathlineto{\pgfqpoint{2.625676in}{1.933618in}}%
\pgfpathlineto{\pgfqpoint{2.626201in}{1.532028in}}%
\pgfpathlineto{\pgfqpoint{2.627573in}{1.532059in}}%
\pgfpathlineto{\pgfqpoint{2.628134in}{1.532148in}}%
\pgfpathlineto{\pgfqpoint{2.628228in}{1.956078in}}%
\pgfpathlineto{\pgfqpoint{2.629270in}{1.935681in}}%
\pgfpathlineto{\pgfqpoint{2.629319in}{1.925038in}}%
\pgfpathlineto{\pgfqpoint{2.630012in}{1.928106in}}%
\pgfpathlineto{\pgfqpoint{2.630160in}{1.935287in}}%
\pgfpathlineto{\pgfqpoint{2.630655in}{1.924498in}}%
\pgfpathlineto{\pgfqpoint{2.631156in}{1.532028in}}%
\pgfpathlineto{\pgfqpoint{2.632188in}{1.532059in}}%
\pgfpathlineto{\pgfqpoint{2.633080in}{1.532147in}}%
\pgfpathlineto{\pgfqpoint{2.633173in}{1.956139in}}%
\pgfpathlineto{\pgfqpoint{2.634235in}{1.930256in}}%
\pgfpathlineto{\pgfqpoint{2.634344in}{1.935810in}}%
\pgfpathlineto{\pgfqpoint{2.634393in}{1.925167in}}%
\pgfpathlineto{\pgfqpoint{2.634987in}{1.931100in}}%
\pgfpathlineto{\pgfqpoint{2.635036in}{1.931214in}}%
\pgfpathlineto{\pgfqpoint{2.635234in}{1.935228in}}%
\pgfpathlineto{\pgfqpoint{2.635679in}{1.874012in}}%
\pgfpathlineto{\pgfqpoint{2.636090in}{1.532031in}}%
\pgfpathlineto{\pgfqpoint{2.636720in}{1.532047in}}%
\pgfpathlineto{\pgfqpoint{2.638023in}{1.532121in}}%
\pgfpathlineto{\pgfqpoint{2.638124in}{1.956719in}}%
\pgfpathlineto{\pgfqpoint{2.639536in}{1.935415in}}%
\pgfpathlineto{\pgfqpoint{2.639553in}{1.925265in}}%
\pgfpathlineto{\pgfqpoint{2.640283in}{1.927940in}}%
\pgfpathlineto{\pgfqpoint{2.640431in}{1.933774in}}%
\pgfpathlineto{\pgfqpoint{2.640481in}{1.925701in}}%
\pgfpathlineto{\pgfqpoint{2.640579in}{1.925848in}}%
\pgfpathlineto{\pgfqpoint{2.641048in}{1.532032in}}%
\pgfpathlineto{\pgfqpoint{2.641782in}{1.532049in}}%
\pgfpathlineto{\pgfqpoint{2.642970in}{1.532136in}}%
\pgfpathlineto{\pgfqpoint{2.643282in}{1.961098in}}%
\pgfpathlineto{\pgfqpoint{2.644225in}{1.924470in}}%
\pgfpathlineto{\pgfqpoint{2.644242in}{1.936156in}}%
\pgfpathlineto{\pgfqpoint{2.645032in}{1.931452in}}%
\pgfpathlineto{\pgfqpoint{2.646003in}{1.532029in}}%
\pgfpathlineto{\pgfqpoint{2.645060in}{1.931727in}}%
\pgfpathlineto{\pgfqpoint{2.647185in}{1.532056in}}%
\pgfpathlineto{\pgfqpoint{2.647914in}{1.532121in}}%
\pgfpathlineto{\pgfqpoint{2.648017in}{1.956763in}}%
\pgfpathlineto{\pgfqpoint{2.649368in}{1.932318in}}%
\pgfpathlineto{\pgfqpoint{2.649374in}{1.933267in}}%
\pgfpathlineto{\pgfqpoint{2.649389in}{1.927305in}}%
\pgfpathlineto{\pgfqpoint{2.650124in}{1.931944in}}%
\pgfpathlineto{\pgfqpoint{2.650519in}{1.883286in}}%
\pgfpathlineto{\pgfqpoint{2.650940in}{1.532032in}}%
\pgfpathlineto{\pgfqpoint{2.651526in}{1.532044in}}%
\pgfpathlineto{\pgfqpoint{2.652860in}{1.532130in}}%
\pgfpathlineto{\pgfqpoint{2.652961in}{1.956177in}}%
\pgfpathlineto{\pgfqpoint{2.654230in}{1.933790in}}%
\pgfpathlineto{\pgfqpoint{2.654884in}{1.926759in}}%
\pgfpathlineto{\pgfqpoint{2.654347in}{1.933850in}}%
\pgfpathlineto{\pgfqpoint{2.654983in}{1.927783in}}%
\pgfpathlineto{\pgfqpoint{2.655279in}{1.932903in}}%
\pgfpathlineto{\pgfqpoint{2.655428in}{1.920188in}}%
\pgfpathlineto{\pgfqpoint{2.655851in}{1.532029in}}%
\pgfpathlineto{\pgfqpoint{2.656600in}{1.532045in}}%
\pgfpathlineto{\pgfqpoint{2.657805in}{1.532121in}}%
\pgfpathlineto{\pgfqpoint{2.657905in}{1.956316in}}%
\pgfpathlineto{\pgfqpoint{2.659323in}{1.934131in}}%
\pgfpathlineto{\pgfqpoint{2.659819in}{1.925608in}}%
\pgfpathlineto{\pgfqpoint{2.659435in}{1.935006in}}%
\pgfpathlineto{\pgfqpoint{2.660066in}{1.929409in}}%
\pgfpathlineto{\pgfqpoint{2.660215in}{1.933776in}}%
\pgfpathlineto{\pgfqpoint{2.660264in}{1.925093in}}%
\pgfpathlineto{\pgfqpoint{2.660363in}{1.925159in}}%
\pgfpathlineto{\pgfqpoint{2.660830in}{1.532032in}}%
\pgfpathlineto{\pgfqpoint{2.661564in}{1.532049in}}%
\pgfpathlineto{\pgfqpoint{2.662752in}{1.532133in}}%
\pgfpathlineto{\pgfqpoint{2.663063in}{1.960829in}}%
\pgfpathlineto{\pgfqpoint{2.664027in}{1.933587in}}%
\pgfpathlineto{\pgfqpoint{2.664036in}{1.924960in}}%
\pgfpathlineto{\pgfqpoint{2.664052in}{1.936037in}}%
\pgfpathlineto{\pgfqpoint{2.664801in}{1.930286in}}%
\pgfpathlineto{\pgfqpoint{2.664838in}{1.931646in}}%
\pgfpathlineto{\pgfqpoint{2.665328in}{1.908735in}}%
\pgfpathlineto{\pgfqpoint{2.665718in}{1.532013in}}%
\pgfpathlineto{\pgfqpoint{2.666449in}{1.532040in}}%
\pgfpathlineto{\pgfqpoint{2.667695in}{1.532112in}}%
\pgfpathlineto{\pgfqpoint{2.667736in}{1.724016in}}%
\pgfpathlineto{\pgfqpoint{2.668000in}{1.957196in}}%
\pgfpathlineto{\pgfqpoint{2.668517in}{1.938889in}}%
\pgfpathlineto{\pgfqpoint{2.668528in}{1.915665in}}%
\pgfpathlineto{\pgfqpoint{2.668544in}{1.944412in}}%
\pgfpathlineto{\pgfqpoint{2.669274in}{1.930508in}}%
\pgfpathlineto{\pgfqpoint{2.669520in}{1.936770in}}%
\pgfpathlineto{\pgfqpoint{2.669916in}{1.923487in}}%
\pgfpathlineto{\pgfqpoint{2.670015in}{1.924738in}}%
\pgfpathlineto{\pgfqpoint{2.670064in}{1.933137in}}%
\pgfpathlineto{\pgfqpoint{2.670311in}{1.859939in}}%
\pgfpathlineto{\pgfqpoint{2.670722in}{1.532028in}}%
\pgfpathlineto{\pgfqpoint{2.671309in}{1.532041in}}%
\pgfpathlineto{\pgfqpoint{2.672644in}{1.532161in}}%
\pgfpathlineto{\pgfqpoint{2.672736in}{1.955389in}}%
\pgfpathlineto{\pgfqpoint{2.673786in}{1.938080in}}%
\pgfpathlineto{\pgfqpoint{2.673802in}{1.922527in}}%
\pgfpathlineto{\pgfqpoint{2.674539in}{1.928343in}}%
\pgfpathlineto{\pgfqpoint{2.674836in}{1.935526in}}%
\pgfpathlineto{\pgfqpoint{2.674786in}{1.924461in}}%
\pgfpathlineto{\pgfqpoint{2.675083in}{1.931789in}}%
\pgfpathlineto{\pgfqpoint{2.675380in}{1.686243in}}%
\pgfpathlineto{\pgfqpoint{2.675643in}{1.532027in}}%
\pgfpathlineto{\pgfqpoint{2.676263in}{1.532050in}}%
\pgfpathlineto{\pgfqpoint{2.677587in}{1.532119in}}%
\pgfpathlineto{\pgfqpoint{2.677627in}{1.731751in}}%
\pgfpathlineto{\pgfqpoint{2.677899in}{1.958641in}}%
\pgfpathlineto{\pgfqpoint{2.678410in}{1.920191in}}%
\pgfpathlineto{\pgfqpoint{2.678415in}{1.917084in}}%
\pgfpathlineto{\pgfqpoint{2.678431in}{1.942839in}}%
\pgfpathlineto{\pgfqpoint{2.679162in}{1.925727in}}%
\pgfpathlineto{\pgfqpoint{2.679339in}{1.936035in}}%
\pgfpathlineto{\pgfqpoint{2.679242in}{1.924482in}}%
\pgfpathlineto{\pgfqpoint{2.679932in}{1.926650in}}%
\pgfpathlineto{\pgfqpoint{2.680130in}{1.931876in}}%
\pgfpathlineto{\pgfqpoint{2.680180in}{1.901098in}}%
\pgfpathlineto{\pgfqpoint{2.680615in}{1.532032in}}%
\pgfpathlineto{\pgfqpoint{2.681250in}{1.532047in}}%
\pgfpathlineto{\pgfqpoint{2.682536in}{1.532201in}}%
\pgfpathlineto{\pgfqpoint{2.682858in}{1.961342in}}%
\pgfpathlineto{\pgfqpoint{2.683648in}{1.925976in}}%
\pgfpathlineto{\pgfqpoint{2.683653in}{1.921446in}}%
\pgfpathlineto{\pgfqpoint{2.683668in}{1.938951in}}%
\pgfpathlineto{\pgfqpoint{2.684410in}{1.926954in}}%
\pgfpathlineto{\pgfqpoint{2.684428in}{1.933362in}}%
\pgfpathlineto{\pgfqpoint{2.685100in}{1.898603in}}%
\pgfpathlineto{\pgfqpoint{2.685583in}{1.532031in}}%
\pgfpathlineto{\pgfqpoint{2.686220in}{1.532049in}}%
\pgfpathlineto{\pgfqpoint{2.687476in}{1.532105in}}%
\pgfpathlineto{\pgfqpoint{2.687505in}{1.580658in}}%
\pgfpathlineto{\pgfqpoint{2.687776in}{1.957873in}}%
\pgfpathlineto{\pgfqpoint{2.688302in}{1.924416in}}%
\pgfpathlineto{\pgfqpoint{2.688314in}{1.943057in}}%
\pgfpathlineto{\pgfqpoint{2.688326in}{1.918568in}}%
\pgfpathlineto{\pgfqpoint{2.689070in}{1.929384in}}%
\pgfpathlineto{\pgfqpoint{2.689104in}{1.936734in}}%
\pgfpathlineto{\pgfqpoint{2.689785in}{1.923709in}}%
\pgfpathlineto{\pgfqpoint{2.689834in}{1.935670in}}%
\pgfpathlineto{\pgfqpoint{2.690082in}{1.863938in}}%
\pgfpathlineto{\pgfqpoint{2.690471in}{1.532031in}}%
\pgfpathlineto{\pgfqpoint{2.691126in}{1.532044in}}%
\pgfpathlineto{\pgfqpoint{2.692429in}{1.532304in}}%
\pgfpathlineto{\pgfqpoint{2.692525in}{1.955892in}}%
\pgfpathlineto{\pgfqpoint{2.693491in}{1.938426in}}%
\pgfpathlineto{\pgfqpoint{2.693609in}{1.922686in}}%
\pgfpathlineto{\pgfqpoint{2.693525in}{1.938776in}}%
\pgfpathlineto{\pgfqpoint{2.694265in}{1.927685in}}%
\pgfpathlineto{\pgfqpoint{2.694413in}{1.934770in}}%
\pgfpathlineto{\pgfqpoint{2.694809in}{1.925319in}}%
\pgfpathlineto{\pgfqpoint{2.694908in}{1.925387in}}%
\pgfpathlineto{\pgfqpoint{2.695425in}{1.532030in}}%
\pgfpathlineto{\pgfqpoint{2.696884in}{1.532061in}}%
\pgfpathlineto{\pgfqpoint{2.697369in}{1.532125in}}%
\pgfpathlineto{\pgfqpoint{2.697677in}{1.959243in}}%
\pgfpathlineto{\pgfqpoint{2.698703in}{1.936123in}}%
\pgfpathlineto{\pgfqpoint{2.698717in}{1.923753in}}%
\pgfpathlineto{\pgfqpoint{2.698759in}{1.936472in}}%
\pgfpathlineto{\pgfqpoint{2.699480in}{1.928666in}}%
\pgfpathlineto{\pgfqpoint{2.699494in}{1.931868in}}%
\pgfpathlineto{\pgfqpoint{2.700002in}{1.849414in}}%
\pgfpathlineto{\pgfqpoint{2.700360in}{1.532029in}}%
\pgfpathlineto{\pgfqpoint{2.700962in}{1.532048in}}%
\pgfpathlineto{\pgfqpoint{2.702314in}{1.532120in}}%
\pgfpathlineto{\pgfqpoint{2.702416in}{1.956702in}}%
\pgfpathlineto{\pgfqpoint{2.703829in}{1.928897in}}%
\pgfpathlineto{\pgfqpoint{2.703839in}{1.925324in}}%
\pgfpathlineto{\pgfqpoint{2.703857in}{1.935148in}}%
\pgfpathlineto{\pgfqpoint{2.704573in}{1.930539in}}%
\pgfpathlineto{\pgfqpoint{2.704771in}{1.933669in}}%
\pgfpathlineto{\pgfqpoint{2.704919in}{1.874271in}}%
\pgfpathlineto{\pgfqpoint{2.705330in}{1.532028in}}%
\pgfpathlineto{\pgfqpoint{2.705961in}{1.532050in}}%
\pgfpathlineto{\pgfqpoint{2.707262in}{1.532151in}}%
\pgfpathlineto{\pgfqpoint{2.707355in}{1.955563in}}%
\pgfpathlineto{\pgfqpoint{2.708415in}{1.924990in}}%
\pgfpathlineto{\pgfqpoint{2.708465in}{1.935728in}}%
\pgfpathlineto{\pgfqpoint{2.709157in}{1.931795in}}%
\pgfpathlineto{\pgfqpoint{2.709355in}{1.935172in}}%
\pgfpathlineto{\pgfqpoint{2.709899in}{1.833010in}}%
\pgfpathlineto{\pgfqpoint{2.710261in}{1.532027in}}%
\pgfpathlineto{\pgfqpoint{2.710830in}{1.532040in}}%
\pgfpathlineto{\pgfqpoint{2.712205in}{1.532119in}}%
\pgfpathlineto{\pgfqpoint{2.712245in}{1.732227in}}%
\pgfpathlineto{\pgfqpoint{2.712526in}{1.957376in}}%
\pgfpathlineto{\pgfqpoint{2.713027in}{1.938708in}}%
\pgfpathlineto{\pgfqpoint{2.713033in}{1.945154in}}%
\pgfpathlineto{\pgfqpoint{2.713049in}{1.915825in}}%
\pgfpathlineto{\pgfqpoint{2.713793in}{1.931882in}}%
\pgfpathlineto{\pgfqpoint{2.713803in}{1.936291in}}%
\pgfpathlineto{\pgfqpoint{2.713820in}{1.924403in}}%
\pgfpathlineto{\pgfqpoint{2.714533in}{1.934797in}}%
\pgfpathlineto{\pgfqpoint{2.714928in}{1.727725in}}%
\pgfpathlineto{\pgfqpoint{2.715203in}{1.532030in}}%
\pgfpathlineto{\pgfqpoint{2.715818in}{1.532048in}}%
\pgfpathlineto{\pgfqpoint{2.717150in}{1.532117in}}%
\pgfpathlineto{\pgfqpoint{2.717251in}{1.956690in}}%
\pgfpathlineto{\pgfqpoint{2.718717in}{1.935111in}}%
\pgfpathlineto{\pgfqpoint{2.719421in}{1.925456in}}%
\pgfpathlineto{\pgfqpoint{2.719470in}{1.933790in}}%
\pgfpathlineto{\pgfqpoint{2.719816in}{1.790420in}}%
\pgfpathlineto{\pgfqpoint{2.720192in}{1.532031in}}%
\pgfpathlineto{\pgfqpoint{2.720730in}{1.532046in}}%
\pgfpathlineto{\pgfqpoint{2.722095in}{1.532117in}}%
\pgfpathlineto{\pgfqpoint{2.722196in}{1.956126in}}%
\pgfpathlineto{\pgfqpoint{2.723659in}{1.925734in}}%
\pgfpathlineto{\pgfqpoint{2.723709in}{1.935220in}}%
\pgfpathlineto{\pgfqpoint{2.724203in}{1.925565in}}%
\pgfpathlineto{\pgfqpoint{2.724401in}{1.930698in}}%
\pgfpathlineto{\pgfqpoint{2.724599in}{1.933674in}}%
\pgfpathlineto{\pgfqpoint{2.724698in}{1.886026in}}%
\pgfpathlineto{\pgfqpoint{2.725123in}{1.532032in}}%
\pgfpathlineto{\pgfqpoint{2.725709in}{1.532044in}}%
\pgfpathlineto{\pgfqpoint{2.727042in}{1.532128in}}%
\pgfpathlineto{\pgfqpoint{2.727143in}{1.956359in}}%
\pgfpathlineto{\pgfqpoint{2.728436in}{1.928586in}}%
\pgfpathlineto{\pgfqpoint{2.728446in}{1.925405in}}%
\pgfpathlineto{\pgfqpoint{2.728463in}{1.935240in}}%
\pgfpathlineto{\pgfqpoint{2.729161in}{1.927195in}}%
\pgfpathlineto{\pgfqpoint{2.729310in}{1.933950in}}%
\pgfpathlineto{\pgfqpoint{2.729606in}{1.899172in}}%
\pgfpathlineto{\pgfqpoint{2.730070in}{1.532032in}}%
\pgfpathlineto{\pgfqpoint{2.730804in}{1.532049in}}%
\pgfpathlineto{\pgfqpoint{2.731990in}{1.532164in}}%
\pgfpathlineto{\pgfqpoint{2.732317in}{1.963204in}}%
\pgfpathlineto{\pgfqpoint{2.733158in}{1.925112in}}%
\pgfpathlineto{\pgfqpoint{2.733163in}{1.921813in}}%
\pgfpathlineto{\pgfqpoint{2.733176in}{1.938045in}}%
\pgfpathlineto{\pgfqpoint{2.733919in}{1.932625in}}%
\pgfpathlineto{\pgfqpoint{2.734654in}{1.788432in}}%
\pgfpathlineto{\pgfqpoint{2.735000in}{1.532028in}}%
\pgfpathlineto{\pgfqpoint{2.735581in}{1.532041in}}%
\pgfpathlineto{\pgfqpoint{2.736934in}{1.532147in}}%
\pgfpathlineto{\pgfqpoint{2.737028in}{1.956540in}}%
\pgfpathlineto{\pgfqpoint{2.738095in}{1.930225in}}%
\pgfpathlineto{\pgfqpoint{2.738105in}{1.923444in}}%
\pgfpathlineto{\pgfqpoint{2.738124in}{1.936832in}}%
\pgfpathlineto{\pgfqpoint{2.738844in}{1.927125in}}%
\pgfpathlineto{\pgfqpoint{2.739140in}{1.935452in}}%
\pgfpathlineto{\pgfqpoint{2.739437in}{1.915259in}}%
\pgfpathlineto{\pgfqpoint{2.739487in}{1.921449in}}%
\pgfpathlineto{\pgfqpoint{2.739941in}{1.532029in}}%
\pgfpathlineto{\pgfqpoint{2.740815in}{1.532048in}}%
\pgfpathlineto{\pgfqpoint{2.741878in}{1.532120in}}%
\pgfpathlineto{\pgfqpoint{2.742097in}{1.958847in}}%
\pgfpathlineto{\pgfqpoint{2.743237in}{1.925338in}}%
\pgfpathlineto{\pgfqpoint{2.743241in}{1.924826in}}%
\pgfpathlineto{\pgfqpoint{2.743254in}{1.935948in}}%
\pgfpathlineto{\pgfqpoint{2.743959in}{1.930936in}}%
\pgfpathlineto{\pgfqpoint{2.743965in}{1.931888in}}%
\pgfpathlineto{\pgfqpoint{2.744464in}{1.897024in}}%
\pgfpathlineto{\pgfqpoint{2.744883in}{1.532029in}}%
\pgfpathlineto{\pgfqpoint{2.745505in}{1.532049in}}%
\pgfpathlineto{\pgfqpoint{2.746824in}{1.532133in}}%
\pgfpathlineto{\pgfqpoint{2.747040in}{1.959495in}}%
\pgfpathlineto{\pgfqpoint{2.748115in}{1.933559in}}%
\pgfpathlineto{\pgfqpoint{2.748125in}{1.923809in}}%
\pgfpathlineto{\pgfqpoint{2.748140in}{1.936586in}}%
\pgfpathlineto{\pgfqpoint{2.748881in}{1.930510in}}%
\pgfpathlineto{\pgfqpoint{2.748891in}{1.932212in}}%
\pgfpathlineto{\pgfqpoint{2.749417in}{1.899968in}}%
\pgfpathlineto{\pgfqpoint{2.749852in}{1.532031in}}%
\pgfpathlineto{\pgfqpoint{2.750488in}{1.532047in}}%
\pgfpathlineto{\pgfqpoint{2.751769in}{1.532120in}}%
\pgfpathlineto{\pgfqpoint{2.752037in}{1.956761in}}%
\pgfpathlineto{\pgfqpoint{2.753286in}{1.928724in}}%
\pgfpathlineto{\pgfqpoint{2.753296in}{1.925363in}}%
\pgfpathlineto{\pgfqpoint{2.753390in}{1.935164in}}%
\pgfpathlineto{\pgfqpoint{2.754001in}{1.927626in}}%
\pgfpathlineto{\pgfqpoint{2.754150in}{1.934190in}}%
\pgfpathlineto{\pgfqpoint{2.754347in}{1.887976in}}%
\pgfpathlineto{\pgfqpoint{2.754754in}{1.532027in}}%
\pgfpathlineto{\pgfqpoint{2.755437in}{1.532052in}}%
\pgfpathlineto{\pgfqpoint{2.756705in}{1.532084in}}%
\pgfpathlineto{\pgfqpoint{2.756716in}{1.532145in}}%
\pgfpathlineto{\pgfqpoint{2.756810in}{1.955996in}}%
\pgfpathlineto{\pgfqpoint{2.757876in}{1.935575in}}%
\pgfpathlineto{\pgfqpoint{2.758256in}{1.924955in}}%
\pgfpathlineto{\pgfqpoint{2.758602in}{1.926047in}}%
\pgfpathlineto{\pgfqpoint{2.758651in}{1.935103in}}%
\pgfpathlineto{\pgfqpoint{2.759146in}{1.924595in}}%
\pgfpathlineto{\pgfqpoint{2.759245in}{1.925710in}}%
\pgfpathlineto{\pgfqpoint{2.759741in}{1.532032in}}%
\pgfpathlineto{\pgfqpoint{2.760525in}{1.532049in}}%
\pgfpathlineto{\pgfqpoint{2.761660in}{1.532128in}}%
\pgfpathlineto{\pgfqpoint{2.761761in}{1.956666in}}%
\pgfpathlineto{\pgfqpoint{2.763076in}{1.935780in}}%
\pgfpathlineto{\pgfqpoint{2.763091in}{1.924731in}}%
\pgfpathlineto{\pgfqpoint{2.763847in}{1.933873in}}%
\pgfpathlineto{\pgfqpoint{2.764341in}{1.772941in}}%
\pgfpathlineto{\pgfqpoint{2.764698in}{1.532031in}}%
\pgfpathlineto{\pgfqpoint{2.765238in}{1.532046in}}%
\pgfpathlineto{\pgfqpoint{2.766606in}{1.532132in}}%
\pgfpathlineto{\pgfqpoint{2.766931in}{1.961363in}}%
\pgfpathlineto{\pgfqpoint{2.767900in}{1.935246in}}%
\pgfpathlineto{\pgfqpoint{2.767911in}{1.923762in}}%
\pgfpathlineto{\pgfqpoint{2.767927in}{1.936673in}}%
\pgfpathlineto{\pgfqpoint{2.768675in}{1.927903in}}%
\pgfpathlineto{\pgfqpoint{2.768695in}{1.931969in}}%
\pgfpathlineto{\pgfqpoint{2.769198in}{1.886998in}}%
\pgfpathlineto{\pgfqpoint{2.769619in}{1.532030in}}%
\pgfpathlineto{\pgfqpoint{2.770251in}{1.532048in}}%
\pgfpathlineto{\pgfqpoint{2.771549in}{1.532111in}}%
\pgfpathlineto{\pgfqpoint{2.771592in}{1.736082in}}%
\pgfpathlineto{\pgfqpoint{2.771858in}{1.954385in}}%
\pgfpathlineto{\pgfqpoint{2.772372in}{1.935785in}}%
\pgfpathlineto{\pgfqpoint{2.772384in}{1.916790in}}%
\pgfpathlineto{\pgfqpoint{2.772396in}{1.943044in}}%
\pgfpathlineto{\pgfqpoint{2.773123in}{1.924825in}}%
\pgfpathlineto{\pgfqpoint{2.773152in}{1.937542in}}%
\pgfpathlineto{\pgfqpoint{2.773202in}{1.923010in}}%
\pgfpathlineto{\pgfqpoint{2.773894in}{1.926857in}}%
\pgfpathlineto{\pgfqpoint{2.774042in}{1.936068in}}%
\pgfpathlineto{\pgfqpoint{2.774141in}{1.890865in}}%
\pgfpathlineto{\pgfqpoint{2.774518in}{1.532021in}}%
\pgfpathlineto{\pgfqpoint{2.775186in}{1.532052in}}%
\pgfpathlineto{\pgfqpoint{2.776496in}{1.532120in}}%
\pgfpathlineto{\pgfqpoint{2.776597in}{1.956224in}}%
\pgfpathlineto{\pgfqpoint{2.778047in}{1.934467in}}%
\pgfpathlineto{\pgfqpoint{2.778721in}{1.924790in}}%
\pgfpathlineto{\pgfqpoint{2.778078in}{1.935639in}}%
\pgfpathlineto{\pgfqpoint{2.778820in}{1.926002in}}%
\pgfpathlineto{\pgfqpoint{2.778869in}{1.931975in}}%
\pgfpathlineto{\pgfqpoint{2.779067in}{1.894995in}}%
\pgfpathlineto{\pgfqpoint{2.779547in}{1.532031in}}%
\pgfpathlineto{\pgfqpoint{2.780234in}{1.532046in}}%
\pgfpathlineto{\pgfqpoint{2.781448in}{1.532437in}}%
\pgfpathlineto{\pgfqpoint{2.781553in}{1.956790in}}%
\pgfpathlineto{\pgfqpoint{2.782478in}{1.931602in}}%
\pgfpathlineto{\pgfqpoint{2.782486in}{1.922913in}}%
\pgfpathlineto{\pgfqpoint{2.782543in}{1.937323in}}%
\pgfpathlineto{\pgfqpoint{2.783246in}{1.927700in}}%
\pgfpathlineto{\pgfqpoint{2.783543in}{1.935014in}}%
\pgfpathlineto{\pgfqpoint{2.783939in}{1.924413in}}%
\pgfpathlineto{\pgfqpoint{2.783988in}{1.924383in}}%
\pgfpathlineto{\pgfqpoint{2.784443in}{1.532029in}}%
\pgfpathlineto{\pgfqpoint{2.785310in}{1.532048in}}%
\pgfpathlineto{\pgfqpoint{2.786387in}{1.532126in}}%
\pgfpathlineto{\pgfqpoint{2.786696in}{1.959277in}}%
\pgfpathlineto{\pgfqpoint{2.787702in}{1.924741in}}%
\pgfpathlineto{\pgfqpoint{2.787707in}{1.923282in}}%
\pgfpathlineto{\pgfqpoint{2.787750in}{1.936825in}}%
\pgfpathlineto{\pgfqpoint{2.788452in}{1.932177in}}%
\pgfpathlineto{\pgfqpoint{2.789157in}{1.655536in}}%
\pgfpathlineto{\pgfqpoint{2.789380in}{1.532031in}}%
\pgfpathlineto{\pgfqpoint{2.790036in}{1.532044in}}%
\pgfpathlineto{\pgfqpoint{2.791338in}{1.532305in}}%
\pgfpathlineto{\pgfqpoint{2.791434in}{1.955847in}}%
\pgfpathlineto{\pgfqpoint{2.792388in}{1.933856in}}%
\pgfpathlineto{\pgfqpoint{2.792430in}{1.922116in}}%
\pgfpathlineto{\pgfqpoint{2.792418in}{1.938254in}}%
\pgfpathlineto{\pgfqpoint{2.793136in}{1.931072in}}%
\pgfpathlineto{\pgfqpoint{2.793185in}{1.931008in}}%
\pgfpathlineto{\pgfqpoint{2.793383in}{1.934884in}}%
\pgfpathlineto{\pgfqpoint{2.793977in}{1.829359in}}%
\pgfpathlineto{\pgfqpoint{2.794300in}{1.532016in}}%
\pgfpathlineto{\pgfqpoint{2.794913in}{1.532038in}}%
\pgfpathlineto{\pgfqpoint{2.796278in}{1.532124in}}%
\pgfpathlineto{\pgfqpoint{2.796581in}{1.958381in}}%
\pgfpathlineto{\pgfqpoint{2.797616in}{1.925812in}}%
\pgfpathlineto{\pgfqpoint{2.797620in}{1.923894in}}%
\pgfpathlineto{\pgfqpoint{2.797637in}{1.936033in}}%
\pgfpathlineto{\pgfqpoint{2.798377in}{1.928360in}}%
\pgfpathlineto{\pgfqpoint{2.798393in}{1.931558in}}%
\pgfpathlineto{\pgfqpoint{2.798894in}{1.868295in}}%
\pgfpathlineto{\pgfqpoint{2.799271in}{1.532031in}}%
\pgfpathlineto{\pgfqpoint{2.799878in}{1.532047in}}%
\pgfpathlineto{\pgfqpoint{2.801220in}{1.532100in}}%
\pgfpathlineto{\pgfqpoint{2.801229in}{1.532306in}}%
\pgfpathlineto{\pgfqpoint{2.801324in}{1.955690in}}%
\pgfpathlineto{\pgfqpoint{2.802276in}{1.924208in}}%
\pgfpathlineto{\pgfqpoint{2.802282in}{1.921648in}}%
\pgfpathlineto{\pgfqpoint{2.802294in}{1.938277in}}%
\pgfpathlineto{\pgfqpoint{2.803024in}{1.932328in}}%
\pgfpathlineto{\pgfqpoint{2.803222in}{1.934715in}}%
\pgfpathlineto{\pgfqpoint{2.803914in}{1.767260in}}%
\pgfpathlineto{\pgfqpoint{2.804191in}{1.532014in}}%
\pgfpathlineto{\pgfqpoint{2.804856in}{1.532052in}}%
\pgfpathlineto{\pgfqpoint{2.806169in}{1.532119in}}%
\pgfpathlineto{\pgfqpoint{2.806386in}{1.957974in}}%
\pgfpathlineto{\pgfqpoint{2.807558in}{1.933251in}}%
\pgfpathlineto{\pgfqpoint{2.807598in}{1.924804in}}%
\pgfpathlineto{\pgfqpoint{2.807583in}{1.935780in}}%
\pgfpathlineto{\pgfqpoint{2.808329in}{1.931081in}}%
\pgfpathlineto{\pgfqpoint{2.809180in}{1.532029in}}%
\pgfpathlineto{\pgfqpoint{2.808355in}{1.931188in}}%
\pgfpathlineto{\pgfqpoint{2.810600in}{1.532067in}}%
\pgfpathlineto{\pgfqpoint{2.811116in}{1.532144in}}%
\pgfpathlineto{\pgfqpoint{2.811207in}{1.955318in}}%
\pgfpathlineto{\pgfqpoint{2.812288in}{1.936593in}}%
\pgfpathlineto{\pgfqpoint{2.812305in}{1.923777in}}%
\pgfpathlineto{\pgfqpoint{2.813061in}{1.935084in}}%
\pgfpathlineto{\pgfqpoint{2.814069in}{1.586877in}}%
\pgfpathlineto{\pgfqpoint{2.814082in}{1.532023in}}%
\pgfpathlineto{\pgfqpoint{2.814841in}{1.532053in}}%
\pgfpathlineto{\pgfqpoint{2.816059in}{1.532119in}}%
\pgfpathlineto{\pgfqpoint{2.816100in}{1.734461in}}%
\pgfpathlineto{\pgfqpoint{2.816382in}{1.956338in}}%
\pgfpathlineto{\pgfqpoint{2.816884in}{1.945331in}}%
\pgfpathlineto{\pgfqpoint{2.816900in}{1.915509in}}%
\pgfpathlineto{\pgfqpoint{2.817659in}{1.928310in}}%
\pgfpathlineto{\pgfqpoint{2.817840in}{1.924540in}}%
\pgfpathlineto{\pgfqpoint{2.817687in}{1.935508in}}%
\pgfpathlineto{\pgfqpoint{2.818395in}{1.930580in}}%
\pgfpathlineto{\pgfqpoint{2.818444in}{1.930597in}}%
\pgfpathlineto{\pgfqpoint{2.818543in}{1.933405in}}%
\pgfpathlineto{\pgfqpoint{2.818642in}{1.893139in}}%
\pgfpathlineto{\pgfqpoint{2.819081in}{1.532029in}}%
\pgfpathlineto{\pgfqpoint{2.819716in}{1.532049in}}%
\pgfpathlineto{\pgfqpoint{2.821006in}{1.532133in}}%
\pgfpathlineto{\pgfqpoint{2.821344in}{1.967439in}}%
\pgfpathlineto{\pgfqpoint{2.822230in}{1.934242in}}%
\pgfpathlineto{\pgfqpoint{2.822241in}{1.918805in}}%
\pgfpathlineto{\pgfqpoint{2.822257in}{1.941083in}}%
\pgfpathlineto{\pgfqpoint{2.822997in}{1.934475in}}%
\pgfpathlineto{\pgfqpoint{2.823973in}{1.532012in}}%
\pgfpathlineto{\pgfqpoint{2.824947in}{1.532060in}}%
\pgfpathlineto{\pgfqpoint{2.825949in}{1.532107in}}%
\pgfpathlineto{\pgfqpoint{2.825976in}{1.576710in}}%
\pgfpathlineto{\pgfqpoint{2.826050in}{1.956442in}}%
\pgfpathlineto{\pgfqpoint{2.826782in}{1.937399in}}%
\pgfpathlineto{\pgfqpoint{2.826793in}{1.916683in}}%
\pgfpathlineto{\pgfqpoint{2.826810in}{1.942921in}}%
\pgfpathlineto{\pgfqpoint{2.827540in}{1.924416in}}%
\pgfpathlineto{\pgfqpoint{2.827783in}{1.936578in}}%
\pgfpathlineto{\pgfqpoint{2.828179in}{1.923718in}}%
\pgfpathlineto{\pgfqpoint{2.828327in}{1.933700in}}%
\pgfpathlineto{\pgfqpoint{2.828575in}{1.850228in}}%
\pgfpathlineto{\pgfqpoint{2.828973in}{1.532029in}}%
\pgfpathlineto{\pgfqpoint{2.829559in}{1.532042in}}%
\pgfpathlineto{\pgfqpoint{2.830897in}{1.532139in}}%
\pgfpathlineto{\pgfqpoint{2.831108in}{1.959940in}}%
\pgfpathlineto{\pgfqpoint{2.832098in}{1.936972in}}%
\pgfpathlineto{\pgfqpoint{2.832116in}{1.923439in}}%
\pgfpathlineto{\pgfqpoint{2.832138in}{1.937236in}}%
\pgfpathlineto{\pgfqpoint{2.832867in}{1.931909in}}%
\pgfpathlineto{\pgfqpoint{2.833115in}{1.934328in}}%
\pgfpathlineto{\pgfqpoint{2.833461in}{1.906088in}}%
\pgfpathlineto{\pgfqpoint{2.833915in}{1.532032in}}%
\pgfpathlineto{\pgfqpoint{2.834648in}{1.532049in}}%
\pgfpathlineto{\pgfqpoint{2.835847in}{1.532300in}}%
\pgfpathlineto{\pgfqpoint{2.835943in}{1.955971in}}%
\pgfpathlineto{\pgfqpoint{2.836908in}{1.937162in}}%
\pgfpathlineto{\pgfqpoint{2.836928in}{1.921456in}}%
\pgfpathlineto{\pgfqpoint{2.837008in}{1.937950in}}%
\pgfpathlineto{\pgfqpoint{2.837933in}{1.928001in}}%
\pgfpathlineto{\pgfqpoint{2.838081in}{1.934389in}}%
\pgfpathlineto{\pgfqpoint{2.838378in}{1.904142in}}%
\pgfpathlineto{\pgfqpoint{2.838901in}{1.532032in}}%
\pgfpathlineto{\pgfqpoint{2.839987in}{1.532057in}}%
\pgfpathlineto{\pgfqpoint{2.840792in}{1.532233in}}%
\pgfpathlineto{\pgfqpoint{2.841109in}{1.964504in}}%
\pgfpathlineto{\pgfqpoint{2.841876in}{1.934717in}}%
\pgfpathlineto{\pgfqpoint{2.841886in}{1.917756in}}%
\pgfpathlineto{\pgfqpoint{2.841902in}{1.941963in}}%
\pgfpathlineto{\pgfqpoint{2.842641in}{1.934844in}}%
\pgfpathlineto{\pgfqpoint{2.843793in}{1.532032in}}%
\pgfpathlineto{\pgfqpoint{2.843923in}{1.532035in}}%
\pgfpathlineto{\pgfqpoint{2.845737in}{1.532255in}}%
\pgfpathlineto{\pgfqpoint{2.846047in}{1.958433in}}%
\pgfpathlineto{\pgfqpoint{2.846804in}{1.924279in}}%
\pgfpathlineto{\pgfqpoint{2.846843in}{1.921996in}}%
\pgfpathlineto{\pgfqpoint{2.846824in}{1.938898in}}%
\pgfpathlineto{\pgfqpoint{2.847542in}{1.927466in}}%
\pgfpathlineto{\pgfqpoint{2.847697in}{1.934506in}}%
\pgfpathlineto{\pgfqpoint{2.848192in}{1.924961in}}%
\pgfpathlineto{\pgfqpoint{2.848241in}{1.930229in}}%
\pgfpathlineto{\pgfqpoint{2.848340in}{1.880540in}}%
\pgfpathlineto{\pgfqpoint{2.848743in}{1.532031in}}%
\pgfpathlineto{\pgfqpoint{2.849321in}{1.532044in}}%
\pgfpathlineto{\pgfqpoint{2.850679in}{1.532129in}}%
\pgfpathlineto{\pgfqpoint{2.850852in}{1.955312in}}%
\pgfpathlineto{\pgfqpoint{2.851965in}{1.925835in}}%
\pgfpathlineto{\pgfqpoint{2.851997in}{1.925283in}}%
\pgfpathlineto{\pgfqpoint{2.851983in}{1.935194in}}%
\pgfpathlineto{\pgfqpoint{2.852700in}{1.928156in}}%
\pgfpathlineto{\pgfqpoint{2.852720in}{1.931663in}}%
\pgfpathlineto{\pgfqpoint{2.853267in}{1.903449in}}%
\pgfpathlineto{\pgfqpoint{2.853676in}{1.532027in}}%
\pgfpathlineto{\pgfqpoint{2.854341in}{1.532042in}}%
\pgfpathlineto{\pgfqpoint{2.855623in}{1.532119in}}%
\pgfpathlineto{\pgfqpoint{2.855663in}{1.732832in}}%
\pgfpathlineto{\pgfqpoint{2.855946in}{1.956392in}}%
\pgfpathlineto{\pgfqpoint{2.856448in}{1.945239in}}%
\pgfpathlineto{\pgfqpoint{2.856464in}{1.915447in}}%
\pgfpathlineto{\pgfqpoint{2.857222in}{1.928105in}}%
\pgfpathlineto{\pgfqpoint{2.857449in}{1.935671in}}%
\pgfpathlineto{\pgfqpoint{2.857944in}{1.924406in}}%
\pgfpathlineto{\pgfqpoint{2.857993in}{1.934013in}}%
\pgfpathlineto{\pgfqpoint{2.858241in}{1.852174in}}%
\pgfpathlineto{\pgfqpoint{2.858624in}{1.532029in}}%
\pgfpathlineto{\pgfqpoint{2.859240in}{1.532048in}}%
\pgfpathlineto{\pgfqpoint{2.860569in}{1.532121in}}%
\pgfpathlineto{\pgfqpoint{2.860774in}{1.958789in}}%
\pgfpathlineto{\pgfqpoint{2.861920in}{1.933389in}}%
\pgfpathlineto{\pgfqpoint{2.861959in}{1.925632in}}%
\pgfpathlineto{\pgfqpoint{2.861946in}{1.934950in}}%
\pgfpathlineto{\pgfqpoint{2.862687in}{1.929462in}}%
\pgfpathlineto{\pgfqpoint{2.862698in}{1.931420in}}%
\pgfpathlineto{\pgfqpoint{2.863165in}{1.882477in}}%
\pgfpathlineto{\pgfqpoint{2.863585in}{1.532028in}}%
\pgfpathlineto{\pgfqpoint{2.864168in}{1.532041in}}%
\pgfpathlineto{\pgfqpoint{2.865516in}{1.532151in}}%
\pgfpathlineto{\pgfqpoint{2.865610in}{1.956108in}}%
\pgfpathlineto{\pgfqpoint{2.866644in}{1.927835in}}%
\pgfpathlineto{\pgfqpoint{2.866792in}{1.935846in}}%
\pgfpathlineto{\pgfqpoint{2.866841in}{1.924597in}}%
\pgfpathlineto{\pgfqpoint{2.867435in}{1.931309in}}%
\pgfpathlineto{\pgfqpoint{2.867732in}{1.924700in}}%
\pgfpathlineto{\pgfqpoint{2.867682in}{1.935272in}}%
\pgfpathlineto{\pgfqpoint{2.868028in}{1.933070in}}%
\pgfpathlineto{\pgfqpoint{2.868542in}{1.532032in}}%
\pgfpathlineto{\pgfqpoint{2.869478in}{1.532053in}}%
\pgfpathlineto{\pgfqpoint{2.870465in}{1.532300in}}%
\pgfpathlineto{\pgfqpoint{2.870659in}{1.956154in}}%
\pgfpathlineto{\pgfqpoint{2.871519in}{1.938740in}}%
\pgfpathlineto{\pgfqpoint{2.871565in}{1.921921in}}%
\pgfpathlineto{\pgfqpoint{2.872274in}{1.932953in}}%
\pgfpathlineto{\pgfqpoint{2.872373in}{1.934755in}}%
\pgfpathlineto{\pgfqpoint{2.873164in}{1.734894in}}%
\pgfpathlineto{\pgfqpoint{2.873446in}{1.532026in}}%
\pgfpathlineto{\pgfqpoint{2.874083in}{1.532039in}}%
\pgfpathlineto{\pgfqpoint{2.875407in}{1.532143in}}%
\pgfpathlineto{\pgfqpoint{2.875498in}{1.954662in}}%
\pgfpathlineto{\pgfqpoint{2.876554in}{1.930302in}}%
\pgfpathlineto{\pgfqpoint{2.876560in}{1.920905in}}%
\pgfpathlineto{\pgfqpoint{2.876615in}{1.938915in}}%
\pgfpathlineto{\pgfqpoint{2.877285in}{1.924373in}}%
\pgfpathlineto{\pgfqpoint{2.877335in}{1.936154in}}%
\pgfpathlineto{\pgfqpoint{2.877829in}{1.923944in}}%
\pgfpathlineto{\pgfqpoint{2.877928in}{1.926267in}}%
\pgfpathlineto{\pgfqpoint{2.878414in}{1.532028in}}%
\pgfpathlineto{\pgfqpoint{2.879139in}{1.532051in}}%
\pgfpathlineto{\pgfqpoint{2.880350in}{1.532120in}}%
\pgfpathlineto{\pgfqpoint{2.880646in}{1.956694in}}%
\pgfpathlineto{\pgfqpoint{2.881892in}{1.935787in}}%
\pgfpathlineto{\pgfqpoint{2.881910in}{1.924668in}}%
\pgfpathlineto{\pgfqpoint{2.882634in}{1.926692in}}%
\pgfpathlineto{\pgfqpoint{2.882782in}{1.934182in}}%
\pgfpathlineto{\pgfqpoint{2.882931in}{1.895447in}}%
\pgfpathlineto{\pgfqpoint{2.883369in}{1.532028in}}%
\pgfpathlineto{\pgfqpoint{2.884052in}{1.532043in}}%
\pgfpathlineto{\pgfqpoint{2.885298in}{1.532154in}}%
\pgfpathlineto{\pgfqpoint{2.885391in}{1.955869in}}%
\pgfpathlineto{\pgfqpoint{2.886414in}{1.934021in}}%
\pgfpathlineto{\pgfqpoint{2.886463in}{1.925537in}}%
\pgfpathlineto{\pgfqpoint{2.886513in}{1.934793in}}%
\pgfpathlineto{\pgfqpoint{2.887156in}{1.930112in}}%
\pgfpathlineto{\pgfqpoint{2.887403in}{1.934267in}}%
\pgfpathlineto{\pgfqpoint{2.887353in}{1.925561in}}%
\pgfpathlineto{\pgfqpoint{2.887749in}{1.933131in}}%
\pgfpathlineto{\pgfqpoint{2.887947in}{1.820279in}}%
\pgfpathlineto{\pgfqpoint{2.888320in}{1.532032in}}%
\pgfpathlineto{\pgfqpoint{2.888906in}{1.532045in}}%
\pgfpathlineto{\pgfqpoint{2.890247in}{1.532300in}}%
\pgfpathlineto{\pgfqpoint{2.890441in}{1.956364in}}%
\pgfpathlineto{\pgfqpoint{2.891303in}{1.935818in}}%
\pgfpathlineto{\pgfqpoint{2.891318in}{1.922470in}}%
\pgfpathlineto{\pgfqpoint{2.891330in}{1.938706in}}%
\pgfpathlineto{\pgfqpoint{2.892033in}{1.934637in}}%
\pgfpathlineto{\pgfqpoint{2.893198in}{1.571623in}}%
\pgfpathlineto{\pgfqpoint{2.893267in}{1.532028in}}%
\pgfpathlineto{\pgfqpoint{2.893952in}{1.532043in}}%
\pgfpathlineto{\pgfqpoint{2.895189in}{1.532156in}}%
\pgfpathlineto{\pgfqpoint{2.895281in}{1.954854in}}%
\pgfpathlineto{\pgfqpoint{2.896329in}{1.921995in}}%
\pgfpathlineto{\pgfqpoint{2.896388in}{1.938478in}}%
\pgfpathlineto{\pgfqpoint{2.896368in}{1.921426in}}%
\pgfpathlineto{\pgfqpoint{2.897072in}{1.934707in}}%
\pgfpathlineto{\pgfqpoint{2.897419in}{1.935483in}}%
\pgfpathlineto{\pgfqpoint{2.897913in}{1.714478in}}%
\pgfpathlineto{\pgfqpoint{2.898213in}{1.532032in}}%
\pgfpathlineto{\pgfqpoint{2.898848in}{1.532047in}}%
\pgfpathlineto{\pgfqpoint{2.900134in}{1.532141in}}%
\pgfpathlineto{\pgfqpoint{2.900446in}{1.961034in}}%
\pgfpathlineto{\pgfqpoint{2.901379in}{1.935337in}}%
\pgfpathlineto{\pgfqpoint{2.901394in}{1.924441in}}%
\pgfpathlineto{\pgfqpoint{2.901406in}{1.936157in}}%
\pgfpathlineto{\pgfqpoint{2.902156in}{1.929033in}}%
\pgfpathlineto{\pgfqpoint{2.902167in}{1.931696in}}%
\pgfpathlineto{\pgfqpoint{2.902743in}{1.871882in}}%
\pgfpathlineto{\pgfqpoint{2.903133in}{1.532030in}}%
\pgfpathlineto{\pgfqpoint{2.903750in}{1.532048in}}%
\pgfpathlineto{\pgfqpoint{2.905078in}{1.532122in}}%
\pgfpathlineto{\pgfqpoint{2.905245in}{1.953505in}}%
\pgfpathlineto{\pgfqpoint{2.906420in}{1.928077in}}%
\pgfpathlineto{\pgfqpoint{2.906425in}{1.925758in}}%
\pgfpathlineto{\pgfqpoint{2.906440in}{1.934381in}}%
\pgfpathlineto{\pgfqpoint{2.907183in}{1.931360in}}%
\pgfpathlineto{\pgfqpoint{2.907902in}{1.615557in}}%
\pgfpathlineto{\pgfqpoint{2.908070in}{1.532028in}}%
\pgfpathlineto{\pgfqpoint{2.908027in}{1.623140in}}%
\pgfpathlineto{\pgfqpoint{2.908923in}{1.532047in}}%
\pgfpathlineto{\pgfqpoint{2.910023in}{1.532120in}}%
\pgfpathlineto{\pgfqpoint{2.910317in}{1.956770in}}%
\pgfpathlineto{\pgfqpoint{2.911566in}{1.936092in}}%
\pgfpathlineto{\pgfqpoint{2.911585in}{1.924472in}}%
\pgfpathlineto{\pgfqpoint{2.912316in}{1.928883in}}%
\pgfpathlineto{\pgfqpoint{2.912365in}{1.928663in}}%
\pgfpathlineto{\pgfqpoint{2.912514in}{1.933627in}}%
\pgfpathlineto{\pgfqpoint{2.912612in}{1.897660in}}%
\pgfpathlineto{\pgfqpoint{2.913018in}{1.532030in}}%
\pgfpathlineto{\pgfqpoint{2.913678in}{1.532044in}}%
\pgfpathlineto{\pgfqpoint{2.914967in}{1.532109in}}%
\pgfpathlineto{\pgfqpoint{2.915009in}{1.725129in}}%
\pgfpathlineto{\pgfqpoint{2.915268in}{1.956986in}}%
\pgfpathlineto{\pgfqpoint{2.915793in}{1.923496in}}%
\pgfpathlineto{\pgfqpoint{2.915805in}{1.943934in}}%
\pgfpathlineto{\pgfqpoint{2.915821in}{1.916839in}}%
\pgfpathlineto{\pgfqpoint{2.916553in}{1.928748in}}%
\pgfpathlineto{\pgfqpoint{2.916669in}{1.937145in}}%
\pgfpathlineto{\pgfqpoint{2.917065in}{1.923235in}}%
\pgfpathlineto{\pgfqpoint{2.917312in}{1.929865in}}%
\pgfpathlineto{\pgfqpoint{2.917461in}{1.934780in}}%
\pgfpathlineto{\pgfqpoint{2.917510in}{1.913596in}}%
\pgfpathlineto{\pgfqpoint{2.917936in}{1.532017in}}%
\pgfpathlineto{\pgfqpoint{2.919107in}{1.532064in}}%
\pgfpathlineto{\pgfqpoint{2.919914in}{1.532118in}}%
\pgfpathlineto{\pgfqpoint{2.920115in}{1.956941in}}%
\pgfpathlineto{\pgfqpoint{2.921317in}{1.932574in}}%
\pgfpathlineto{\pgfqpoint{2.921356in}{1.925782in}}%
\pgfpathlineto{\pgfqpoint{2.921341in}{1.934545in}}%
\pgfpathlineto{\pgfqpoint{2.922087in}{1.930640in}}%
\pgfpathlineto{\pgfqpoint{2.922111in}{1.930770in}}%
\pgfpathlineto{\pgfqpoint{2.922672in}{1.662239in}}%
\pgfpathlineto{\pgfqpoint{2.922945in}{1.532032in}}%
\pgfpathlineto{\pgfqpoint{2.923580in}{1.532047in}}%
\pgfpathlineto{\pgfqpoint{2.924863in}{1.532180in}}%
\pgfpathlineto{\pgfqpoint{2.925184in}{1.963210in}}%
\pgfpathlineto{\pgfqpoint{2.926013in}{1.938268in}}%
\pgfpathlineto{\pgfqpoint{2.926025in}{1.921407in}}%
\pgfpathlineto{\pgfqpoint{2.926790in}{1.927463in}}%
\pgfpathlineto{\pgfqpoint{2.926801in}{1.932366in}}%
\pgfpathlineto{\pgfqpoint{2.927465in}{1.888525in}}%
\pgfpathlineto{\pgfqpoint{2.927877in}{1.532030in}}%
\pgfpathlineto{\pgfqpoint{2.928513in}{1.532048in}}%
\pgfpathlineto{\pgfqpoint{2.929802in}{1.532100in}}%
\pgfpathlineto{\pgfqpoint{2.929810in}{1.532292in}}%
\pgfpathlineto{\pgfqpoint{2.929906in}{1.956089in}}%
\pgfpathlineto{\pgfqpoint{2.930866in}{1.934115in}}%
\pgfpathlineto{\pgfqpoint{2.930910in}{1.922714in}}%
\pgfpathlineto{\pgfqpoint{2.930897in}{1.937096in}}%
\pgfpathlineto{\pgfqpoint{2.931603in}{1.926106in}}%
\pgfpathlineto{\pgfqpoint{2.931652in}{1.934397in}}%
\pgfpathlineto{\pgfqpoint{2.932295in}{1.924849in}}%
\pgfpathlineto{\pgfqpoint{2.932345in}{1.928531in}}%
\pgfpathlineto{\pgfqpoint{2.932772in}{1.532017in}}%
\pgfpathlineto{\pgfqpoint{2.933629in}{1.532056in}}%
\pgfpathlineto{\pgfqpoint{2.934740in}{1.532082in}}%
\pgfpathlineto{\pgfqpoint{2.934751in}{1.532131in}}%
\pgfpathlineto{\pgfqpoint{2.935068in}{1.960082in}}%
\pgfpathlineto{\pgfqpoint{2.936043in}{1.924462in}}%
\pgfpathlineto{\pgfqpoint{2.936048in}{1.923861in}}%
\pgfpathlineto{\pgfqpoint{2.936061in}{1.936571in}}%
\pgfpathlineto{\pgfqpoint{2.936759in}{1.929362in}}%
\pgfpathlineto{\pgfqpoint{2.936802in}{1.931987in}}%
\pgfpathlineto{\pgfqpoint{2.937340in}{1.901160in}}%
\pgfpathlineto{\pgfqpoint{2.937718in}{1.532024in}}%
\pgfpathlineto{\pgfqpoint{2.938429in}{1.532041in}}%
\pgfpathlineto{\pgfqpoint{2.939696in}{1.532119in}}%
\pgfpathlineto{\pgfqpoint{2.939736in}{1.731181in}}%
\pgfpathlineto{\pgfqpoint{2.939785in}{1.954192in}}%
\pgfpathlineto{\pgfqpoint{2.940517in}{1.917591in}}%
\pgfpathlineto{\pgfqpoint{2.940533in}{1.942140in}}%
\pgfpathlineto{\pgfqpoint{2.941279in}{1.934682in}}%
\pgfpathlineto{\pgfqpoint{2.941854in}{1.924230in}}%
\pgfpathlineto{\pgfqpoint{2.941458in}{1.935789in}}%
\pgfpathlineto{\pgfqpoint{2.942052in}{1.927964in}}%
\pgfpathlineto{\pgfqpoint{2.942150in}{1.931540in}}%
\pgfpathlineto{\pgfqpoint{2.942200in}{1.917717in}}%
\pgfpathlineto{\pgfqpoint{2.942249in}{1.923770in}}%
\pgfpathlineto{\pgfqpoint{2.942694in}{1.532031in}}%
\pgfpathlineto{\pgfqpoint{2.943556in}{1.532049in}}%
\pgfpathlineto{\pgfqpoint{2.944640in}{1.532112in}}%
\pgfpathlineto{\pgfqpoint{2.944682in}{1.740498in}}%
\pgfpathlineto{\pgfqpoint{2.944939in}{1.956591in}}%
\pgfpathlineto{\pgfqpoint{2.945468in}{1.941457in}}%
\pgfpathlineto{\pgfqpoint{2.945533in}{1.943178in}}%
\pgfpathlineto{\pgfqpoint{2.945517in}{1.916759in}}%
\pgfpathlineto{\pgfqpoint{2.946163in}{1.926319in}}%
\pgfpathlineto{\pgfqpoint{2.946206in}{1.923630in}}%
\pgfpathlineto{\pgfqpoint{2.946346in}{1.936391in}}%
\pgfpathlineto{\pgfqpoint{2.946837in}{1.924136in}}%
\pgfpathlineto{\pgfqpoint{2.946887in}{1.934223in}}%
\pgfpathlineto{\pgfqpoint{2.947233in}{1.887683in}}%
\pgfpathlineto{\pgfqpoint{2.947656in}{1.532028in}}%
\pgfpathlineto{\pgfqpoint{2.948288in}{1.532050in}}%
\pgfpathlineto{\pgfqpoint{2.949589in}{1.532150in}}%
\pgfpathlineto{\pgfqpoint{2.949682in}{1.955949in}}%
\pgfpathlineto{\pgfqpoint{2.950712in}{1.935379in}}%
\pgfpathlineto{\pgfqpoint{2.950761in}{1.924380in}}%
\pgfpathlineto{\pgfqpoint{2.950811in}{1.935488in}}%
\pgfpathlineto{\pgfqpoint{2.951454in}{1.928945in}}%
\pgfpathlineto{\pgfqpoint{2.951602in}{1.935089in}}%
\pgfpathlineto{\pgfqpoint{2.952097in}{1.923785in}}%
\pgfpathlineto{\pgfqpoint{2.952543in}{1.572831in}}%
\pgfpathlineto{\pgfqpoint{2.952637in}{1.532031in}}%
\pgfpathlineto{\pgfqpoint{2.953274in}{1.532049in}}%
\pgfpathlineto{\pgfqpoint{2.954529in}{1.532100in}}%
\pgfpathlineto{\pgfqpoint{2.954538in}{1.532293in}}%
\pgfpathlineto{\pgfqpoint{2.954634in}{1.955561in}}%
\pgfpathlineto{\pgfqpoint{2.955596in}{1.923386in}}%
\pgfpathlineto{\pgfqpoint{2.955617in}{1.938227in}}%
\pgfpathlineto{\pgfqpoint{2.955696in}{1.922765in}}%
\pgfpathlineto{\pgfqpoint{2.956415in}{1.934367in}}%
\pgfpathlineto{\pgfqpoint{2.957305in}{1.663907in}}%
\pgfpathlineto{\pgfqpoint{2.957543in}{1.532028in}}%
\pgfpathlineto{\pgfqpoint{2.958221in}{1.532043in}}%
\pgfpathlineto{\pgfqpoint{2.959480in}{1.532144in}}%
\pgfpathlineto{\pgfqpoint{2.959573in}{1.955233in}}%
\pgfpathlineto{\pgfqpoint{2.960656in}{1.924461in}}%
\pgfpathlineto{\pgfqpoint{2.960674in}{1.935529in}}%
\pgfpathlineto{\pgfqpoint{2.961417in}{1.924891in}}%
\pgfpathlineto{\pgfqpoint{2.961466in}{1.935127in}}%
\pgfpathlineto{\pgfqpoint{2.962010in}{1.910817in}}%
\pgfpathlineto{\pgfqpoint{2.962536in}{1.532032in}}%
\pgfpathlineto{\pgfqpoint{2.963671in}{1.532059in}}%
\pgfpathlineto{\pgfqpoint{2.964428in}{1.532227in}}%
\pgfpathlineto{\pgfqpoint{2.964752in}{1.962390in}}%
\pgfpathlineto{\pgfqpoint{2.965507in}{1.921983in}}%
\pgfpathlineto{\pgfqpoint{2.965573in}{1.920449in}}%
\pgfpathlineto{\pgfqpoint{2.965589in}{1.939514in}}%
\pgfpathlineto{\pgfqpoint{2.966239in}{1.930790in}}%
\pgfpathlineto{\pgfqpoint{2.966367in}{1.934478in}}%
\pgfpathlineto{\pgfqpoint{2.966759in}{1.925306in}}%
\pgfpathlineto{\pgfqpoint{2.966957in}{1.927251in}}%
\pgfpathlineto{\pgfqpoint{2.967452in}{1.532027in}}%
\pgfpathlineto{\pgfqpoint{2.968187in}{1.532053in}}%
\pgfpathlineto{\pgfqpoint{2.969360in}{1.532085in}}%
\pgfpathlineto{\pgfqpoint{2.969371in}{1.532145in}}%
\pgfpathlineto{\pgfqpoint{2.969531in}{1.956211in}}%
\pgfpathlineto{\pgfqpoint{2.970536in}{1.923721in}}%
\pgfpathlineto{\pgfqpoint{2.970662in}{1.936135in}}%
\pgfpathlineto{\pgfqpoint{2.971306in}{1.926162in}}%
\pgfpathlineto{\pgfqpoint{2.971702in}{1.933856in}}%
\pgfpathlineto{\pgfqpoint{2.971900in}{1.906844in}}%
\pgfpathlineto{\pgfqpoint{2.972328in}{1.549244in}}%
\pgfpathlineto{\pgfqpoint{2.972390in}{1.532032in}}%
\pgfpathlineto{\pgfqpoint{2.973075in}{1.532047in}}%
\pgfpathlineto{\pgfqpoint{2.974320in}{1.532300in}}%
\pgfpathlineto{\pgfqpoint{2.974513in}{1.955944in}}%
\pgfpathlineto{\pgfqpoint{2.975375in}{1.938943in}}%
\pgfpathlineto{\pgfqpoint{2.975421in}{1.922010in}}%
\pgfpathlineto{\pgfqpoint{2.976111in}{1.929393in}}%
\pgfpathlineto{\pgfqpoint{2.976359in}{1.934500in}}%
\pgfpathlineto{\pgfqpoint{2.976754in}{1.924612in}}%
\pgfpathlineto{\pgfqpoint{2.976804in}{1.933144in}}%
\pgfpathlineto{\pgfqpoint{2.976952in}{1.827025in}}%
\pgfpathlineto{\pgfqpoint{2.977340in}{1.532028in}}%
\pgfpathlineto{\pgfqpoint{2.977926in}{1.532041in}}%
\pgfpathlineto{\pgfqpoint{2.979262in}{1.532159in}}%
\pgfpathlineto{\pgfqpoint{2.979353in}{1.954676in}}%
\pgfpathlineto{\pgfqpoint{2.980388in}{1.923526in}}%
\pgfpathlineto{\pgfqpoint{2.980394in}{1.920648in}}%
\pgfpathlineto{\pgfqpoint{2.980408in}{1.939538in}}%
\pgfpathlineto{\pgfqpoint{2.981128in}{1.935079in}}%
\pgfpathlineto{\pgfqpoint{2.982227in}{1.532025in}}%
\pgfpathlineto{\pgfqpoint{2.981227in}{1.935633in}}%
\pgfpathlineto{\pgfqpoint{2.983732in}{1.532059in}}%
\pgfpathlineto{\pgfqpoint{2.984205in}{1.532118in}}%
\pgfpathlineto{\pgfqpoint{2.984245in}{1.722104in}}%
\pgfpathlineto{\pgfqpoint{2.984521in}{1.958489in}}%
\pgfpathlineto{\pgfqpoint{2.985030in}{1.930491in}}%
\pgfpathlineto{\pgfqpoint{2.985040in}{1.917106in}}%
\pgfpathlineto{\pgfqpoint{2.985116in}{1.942614in}}%
\pgfpathlineto{\pgfqpoint{2.985795in}{1.933676in}}%
\pgfpathlineto{\pgfqpoint{2.985840in}{1.936031in}}%
\pgfpathlineto{\pgfqpoint{2.986104in}{1.924336in}}%
\pgfpathlineto{\pgfqpoint{2.986400in}{1.932153in}}%
\pgfpathlineto{\pgfqpoint{2.986499in}{1.934666in}}%
\pgfpathlineto{\pgfqpoint{2.986845in}{1.831094in}}%
\pgfpathlineto{\pgfqpoint{2.987219in}{1.532028in}}%
\pgfpathlineto{\pgfqpoint{2.987801in}{1.532041in}}%
\pgfpathlineto{\pgfqpoint{2.989153in}{1.532149in}}%
\pgfpathlineto{\pgfqpoint{2.989246in}{1.956296in}}%
\pgfpathlineto{\pgfqpoint{2.990297in}{1.924556in}}%
\pgfpathlineto{\pgfqpoint{2.990535in}{1.935653in}}%
\pgfpathlineto{\pgfqpoint{2.991079in}{1.934265in}}%
\pgfpathlineto{\pgfqpoint{2.991426in}{1.934921in}}%
\pgfpathlineto{\pgfqpoint{2.991871in}{1.721978in}}%
\pgfpathlineto{\pgfqpoint{2.992149in}{1.532030in}}%
\pgfpathlineto{\pgfqpoint{2.992812in}{1.532044in}}%
\pgfpathlineto{\pgfqpoint{2.994095in}{1.532116in}}%
\pgfpathlineto{\pgfqpoint{2.994138in}{1.762904in}}%
\pgfpathlineto{\pgfqpoint{2.994197in}{1.956280in}}%
\pgfpathlineto{\pgfqpoint{2.994918in}{1.918841in}}%
\pgfpathlineto{\pgfqpoint{2.994935in}{1.941359in}}%
\pgfpathlineto{\pgfqpoint{2.995704in}{1.935360in}}%
\pgfpathlineto{\pgfqpoint{2.996352in}{1.924982in}}%
\pgfpathlineto{\pgfqpoint{2.996451in}{1.927025in}}%
\pgfpathlineto{\pgfqpoint{2.996501in}{1.930218in}}%
\pgfpathlineto{\pgfqpoint{2.996600in}{1.916580in}}%
\pgfpathlineto{\pgfqpoint{2.996649in}{1.920924in}}%
\pgfpathlineto{\pgfqpoint{2.997151in}{1.532031in}}%
\pgfpathlineto{\pgfqpoint{2.997937in}{1.532048in}}%
\pgfpathlineto{\pgfqpoint{2.999047in}{1.532304in}}%
\pgfpathlineto{\pgfqpoint{2.999241in}{1.956573in}}%
\pgfpathlineto{\pgfqpoint{3.000103in}{1.936008in}}%
\pgfpathlineto{\pgfqpoint{3.000117in}{1.921814in}}%
\pgfpathlineto{\pgfqpoint{3.000130in}{1.938494in}}%
\pgfpathlineto{\pgfqpoint{3.000850in}{1.934172in}}%
\pgfpathlineto{\pgfqpoint{3.001938in}{1.619006in}}%
\pgfpathlineto{\pgfqpoint{3.000949in}{1.935104in}}%
\pgfpathlineto{\pgfqpoint{3.001995in}{1.631784in}}%
\pgfpathlineto{\pgfqpoint{3.002010in}{1.532013in}}%
\pgfpathlineto{\pgfqpoint{3.002759in}{1.532039in}}%
\pgfpathlineto{\pgfqpoint{3.003994in}{1.532634in}}%
\pgfpathlineto{\pgfqpoint{3.004098in}{1.956557in}}%
\pgfpathlineto{\pgfqpoint{3.005009in}{1.922027in}}%
\pgfpathlineto{\pgfqpoint{3.005029in}{1.937819in}}%
\pgfpathlineto{\pgfqpoint{3.005759in}{1.926307in}}%
\pgfpathlineto{\pgfqpoint{3.005907in}{1.935382in}}%
\pgfpathlineto{\pgfqpoint{3.006303in}{1.924243in}}%
\pgfpathlineto{\pgfqpoint{3.006501in}{1.927700in}}%
\pgfpathlineto{\pgfqpoint{3.006998in}{1.532028in}}%
\pgfpathlineto{\pgfqpoint{3.007726in}{1.532052in}}%
\pgfpathlineto{\pgfqpoint{3.008934in}{1.532145in}}%
\pgfpathlineto{\pgfqpoint{3.009028in}{1.955211in}}%
\pgfpathlineto{\pgfqpoint{3.010098in}{1.926331in}}%
\pgfpathlineto{\pgfqpoint{3.010127in}{1.935716in}}%
\pgfpathlineto{\pgfqpoint{3.010346in}{1.924620in}}%
\pgfpathlineto{\pgfqpoint{3.010840in}{1.935262in}}%
\pgfpathlineto{\pgfqpoint{3.011889in}{1.572998in}}%
\pgfpathlineto{\pgfqpoint{3.011902in}{1.532022in}}%
\pgfpathlineto{\pgfqpoint{3.012649in}{1.532041in}}%
\pgfpathlineto{\pgfqpoint{3.013878in}{1.532124in}}%
\pgfpathlineto{\pgfqpoint{3.014111in}{1.959904in}}%
\pgfpathlineto{\pgfqpoint{3.015192in}{1.928919in}}%
\pgfpathlineto{\pgfqpoint{3.015230in}{1.923445in}}%
\pgfpathlineto{\pgfqpoint{3.015243in}{1.936646in}}%
\pgfpathlineto{\pgfqpoint{3.015957in}{1.931231in}}%
\pgfpathlineto{\pgfqpoint{3.015995in}{1.931921in}}%
\pgfpathlineto{\pgfqpoint{3.016423in}{1.917614in}}%
\pgfpathlineto{\pgfqpoint{3.016845in}{1.532014in}}%
\pgfpathlineto{\pgfqpoint{3.017700in}{1.532057in}}%
\pgfpathlineto{\pgfqpoint{3.018823in}{1.532124in}}%
\pgfpathlineto{\pgfqpoint{3.018991in}{1.954542in}}%
\pgfpathlineto{\pgfqpoint{3.020178in}{1.934112in}}%
\pgfpathlineto{\pgfqpoint{3.020193in}{1.926146in}}%
\pgfpathlineto{\pgfqpoint{3.020966in}{1.930359in}}%
\pgfpathlineto{\pgfqpoint{3.020975in}{1.930978in}}%
\pgfpathlineto{\pgfqpoint{3.021364in}{1.926938in}}%
\pgfpathlineto{\pgfqpoint{3.021850in}{1.532032in}}%
\pgfpathlineto{\pgfqpoint{3.022880in}{1.532056in}}%
\pgfpathlineto{\pgfqpoint{3.023771in}{1.532156in}}%
\pgfpathlineto{\pgfqpoint{3.024078in}{1.959061in}}%
\pgfpathlineto{\pgfqpoint{3.024960in}{1.925765in}}%
\pgfpathlineto{\pgfqpoint{3.024964in}{1.923357in}}%
\pgfpathlineto{\pgfqpoint{3.024980in}{1.936350in}}%
\pgfpathlineto{\pgfqpoint{3.025724in}{1.929485in}}%
\pgfpathlineto{\pgfqpoint{3.025734in}{1.931766in}}%
\pgfpathlineto{\pgfqpoint{3.026335in}{1.915856in}}%
\pgfpathlineto{\pgfqpoint{3.026777in}{1.532030in}}%
\pgfpathlineto{\pgfqpoint{3.027502in}{1.532050in}}%
\pgfpathlineto{\pgfqpoint{3.028711in}{1.532100in}}%
\pgfpathlineto{\pgfqpoint{3.028720in}{1.532305in}}%
\pgfpathlineto{\pgfqpoint{3.028815in}{1.955659in}}%
\pgfpathlineto{\pgfqpoint{3.029770in}{1.935959in}}%
\pgfpathlineto{\pgfqpoint{3.029783in}{1.921717in}}%
\pgfpathlineto{\pgfqpoint{3.029830in}{1.938148in}}%
\pgfpathlineto{\pgfqpoint{3.030547in}{1.925630in}}%
\pgfpathlineto{\pgfqpoint{3.030596in}{1.934687in}}%
\pgfpathlineto{\pgfqpoint{3.031239in}{1.911374in}}%
\pgfpathlineto{\pgfqpoint{3.031387in}{1.795675in}}%
\pgfpathlineto{\pgfqpoint{3.031719in}{1.532029in}}%
\pgfpathlineto{\pgfqpoint{3.032342in}{1.532049in}}%
\pgfpathlineto{\pgfqpoint{3.033660in}{1.532120in}}%
\pgfpathlineto{\pgfqpoint{3.033760in}{1.955856in}}%
\pgfpathlineto{\pgfqpoint{3.035162in}{1.934662in}}%
\pgfpathlineto{\pgfqpoint{3.035857in}{1.925657in}}%
\pgfpathlineto{\pgfqpoint{3.035907in}{1.933000in}}%
\pgfpathlineto{\pgfqpoint{3.036302in}{1.824041in}}%
\pgfpathlineto{\pgfqpoint{3.036627in}{1.532015in}}%
\pgfpathlineto{\pgfqpoint{3.037250in}{1.532038in}}%
\pgfpathlineto{\pgfqpoint{3.038605in}{1.532118in}}%
\pgfpathlineto{\pgfqpoint{3.038898in}{1.958265in}}%
\pgfpathlineto{\pgfqpoint{3.040016in}{1.930385in}}%
\pgfpathlineto{\pgfqpoint{3.040024in}{1.924505in}}%
\pgfpathlineto{\pgfqpoint{3.040039in}{1.935447in}}%
\pgfpathlineto{\pgfqpoint{3.040789in}{1.928707in}}%
\pgfpathlineto{\pgfqpoint{3.040809in}{1.930915in}}%
\pgfpathlineto{\pgfqpoint{3.041232in}{1.856793in}}%
\pgfpathlineto{\pgfqpoint{3.041631in}{1.532032in}}%
\pgfpathlineto{\pgfqpoint{3.042217in}{1.532044in}}%
\pgfpathlineto{\pgfqpoint{3.043551in}{1.532130in}}%
\pgfpathlineto{\pgfqpoint{3.043652in}{1.956538in}}%
\pgfpathlineto{\pgfqpoint{3.044929in}{1.928593in}}%
\pgfpathlineto{\pgfqpoint{3.044937in}{1.925457in}}%
\pgfpathlineto{\pgfqpoint{3.044955in}{1.934602in}}%
\pgfpathlineto{\pgfqpoint{3.045653in}{1.931048in}}%
\pgfpathlineto{\pgfqpoint{3.045901in}{1.932884in}}%
\pgfpathlineto{\pgfqpoint{3.046098in}{1.908944in}}%
\pgfpathlineto{\pgfqpoint{3.046518in}{1.532010in}}%
\pgfpathlineto{\pgfqpoint{3.047583in}{1.532062in}}%
\pgfpathlineto{\pgfqpoint{3.048494in}{1.532109in}}%
\pgfpathlineto{\pgfqpoint{3.048536in}{1.722362in}}%
\pgfpathlineto{\pgfqpoint{3.048798in}{1.956087in}}%
\pgfpathlineto{\pgfqpoint{3.049318in}{1.935259in}}%
\pgfpathlineto{\pgfqpoint{3.049330in}{1.916956in}}%
\pgfpathlineto{\pgfqpoint{3.049342in}{1.942834in}}%
\pgfpathlineto{\pgfqpoint{3.050071in}{1.923581in}}%
\pgfpathlineto{\pgfqpoint{3.050337in}{1.936943in}}%
\pgfpathlineto{\pgfqpoint{3.050832in}{1.922692in}}%
\pgfpathlineto{\pgfqpoint{3.050881in}{1.935439in}}%
\pgfpathlineto{\pgfqpoint{3.051129in}{1.843155in}}%
\pgfpathlineto{\pgfqpoint{3.051509in}{1.532028in}}%
\pgfpathlineto{\pgfqpoint{3.052091in}{1.532041in}}%
\pgfpathlineto{\pgfqpoint{3.053444in}{1.532148in}}%
\pgfpathlineto{\pgfqpoint{3.053537in}{1.955742in}}%
\pgfpathlineto{\pgfqpoint{3.054586in}{1.931714in}}%
\pgfpathlineto{\pgfqpoint{3.054882in}{1.924756in}}%
\pgfpathlineto{\pgfqpoint{3.054833in}{1.935486in}}%
\pgfpathlineto{\pgfqpoint{3.055229in}{1.925161in}}%
\pgfpathlineto{\pgfqpoint{3.055278in}{1.935282in}}%
\pgfpathlineto{\pgfqpoint{3.055773in}{1.924688in}}%
\pgfpathlineto{\pgfqpoint{3.055970in}{1.929244in}}%
\pgfpathlineto{\pgfqpoint{3.056452in}{1.532032in}}%
\pgfpathlineto{\pgfqpoint{3.057230in}{1.532049in}}%
\pgfpathlineto{\pgfqpoint{3.058392in}{1.532300in}}%
\pgfpathlineto{\pgfqpoint{3.058586in}{1.955852in}}%
\pgfpathlineto{\pgfqpoint{3.059452in}{1.932656in}}%
\pgfpathlineto{\pgfqpoint{3.059529in}{1.922015in}}%
\pgfpathlineto{\pgfqpoint{3.059479in}{1.938401in}}%
\pgfpathlineto{\pgfqpoint{3.060187in}{1.934583in}}%
\pgfpathlineto{\pgfqpoint{3.060286in}{1.934858in}}%
\pgfpathlineto{\pgfqpoint{3.061344in}{1.562073in}}%
\pgfpathlineto{\pgfqpoint{3.061372in}{1.532027in}}%
\pgfpathlineto{\pgfqpoint{3.062103in}{1.532042in}}%
\pgfpathlineto{\pgfqpoint{3.063334in}{1.532145in}}%
\pgfpathlineto{\pgfqpoint{3.063428in}{1.955255in}}%
\pgfpathlineto{\pgfqpoint{3.064497in}{1.933956in}}%
\pgfpathlineto{\pgfqpoint{3.064683in}{1.924914in}}%
\pgfpathlineto{\pgfqpoint{3.064530in}{1.935482in}}%
\pgfpathlineto{\pgfqpoint{3.065227in}{1.925470in}}%
\pgfpathlineto{\pgfqpoint{3.065623in}{1.934450in}}%
\pgfpathlineto{\pgfqpoint{3.065821in}{1.920051in}}%
\pgfpathlineto{\pgfqpoint{3.065870in}{1.929074in}}%
\pgfpathlineto{\pgfqpoint{3.066325in}{1.532028in}}%
\pgfpathlineto{\pgfqpoint{3.067229in}{1.532056in}}%
\pgfpathlineto{\pgfqpoint{3.068278in}{1.532119in}}%
\pgfpathlineto{\pgfqpoint{3.068318in}{1.738205in}}%
\pgfpathlineto{\pgfqpoint{3.068600in}{1.956362in}}%
\pgfpathlineto{\pgfqpoint{3.069100in}{1.938054in}}%
\pgfpathlineto{\pgfqpoint{3.069106in}{1.945354in}}%
\pgfpathlineto{\pgfqpoint{3.069121in}{1.914872in}}%
\pgfpathlineto{\pgfqpoint{3.069866in}{1.933192in}}%
\pgfpathlineto{\pgfqpoint{3.069874in}{1.936018in}}%
\pgfpathlineto{\pgfqpoint{3.069892in}{1.924211in}}%
\pgfpathlineto{\pgfqpoint{3.070638in}{1.934888in}}%
\pgfpathlineto{\pgfqpoint{3.071034in}{1.692944in}}%
\pgfpathlineto{\pgfqpoint{3.071294in}{1.532028in}}%
\pgfpathlineto{\pgfqpoint{3.071976in}{1.532043in}}%
\pgfpathlineto{\pgfqpoint{3.073226in}{1.532151in}}%
\pgfpathlineto{\pgfqpoint{3.073319in}{1.955746in}}%
\pgfpathlineto{\pgfqpoint{3.074348in}{1.932501in}}%
\pgfpathlineto{\pgfqpoint{3.074645in}{1.924332in}}%
\pgfpathlineto{\pgfqpoint{3.074596in}{1.935647in}}%
\pgfpathlineto{\pgfqpoint{3.075090in}{1.924429in}}%
\pgfpathlineto{\pgfqpoint{3.075486in}{1.934977in}}%
\pgfpathlineto{\pgfqpoint{3.075535in}{1.924259in}}%
\pgfpathlineto{\pgfqpoint{3.075733in}{1.929064in}}%
\pgfpathlineto{\pgfqpoint{3.076233in}{1.532029in}}%
\pgfpathlineto{\pgfqpoint{3.076959in}{1.532052in}}%
\pgfpathlineto{\pgfqpoint{3.078169in}{1.532120in}}%
\pgfpathlineto{\pgfqpoint{3.078458in}{1.957066in}}%
\pgfpathlineto{\pgfqpoint{3.079706in}{1.928248in}}%
\pgfpathlineto{\pgfqpoint{3.079716in}{1.924694in}}%
\pgfpathlineto{\pgfqpoint{3.079734in}{1.935084in}}%
\pgfpathlineto{\pgfqpoint{3.080420in}{1.928098in}}%
\pgfpathlineto{\pgfqpoint{3.080569in}{1.934305in}}%
\pgfpathlineto{\pgfqpoint{3.080618in}{1.924379in}}%
\pgfpathlineto{\pgfqpoint{3.080717in}{1.925598in}}%
\pgfpathlineto{\pgfqpoint{3.081136in}{1.532011in}}%
\pgfpathlineto{\pgfqpoint{3.081904in}{1.532055in}}%
\pgfpathlineto{\pgfqpoint{3.083113in}{1.532110in}}%
\pgfpathlineto{\pgfqpoint{3.083154in}{1.722773in}}%
\pgfpathlineto{\pgfqpoint{3.083414in}{1.957133in}}%
\pgfpathlineto{\pgfqpoint{3.083938in}{1.921994in}}%
\pgfpathlineto{\pgfqpoint{3.083949in}{1.944068in}}%
\pgfpathlineto{\pgfqpoint{3.083965in}{1.916366in}}%
\pgfpathlineto{\pgfqpoint{3.084690in}{1.933955in}}%
\pgfpathlineto{\pgfqpoint{3.084822in}{1.922937in}}%
\pgfpathlineto{\pgfqpoint{3.084773in}{1.937245in}}%
\pgfpathlineto{\pgfqpoint{3.085465in}{1.930188in}}%
\pgfpathlineto{\pgfqpoint{3.085564in}{1.934200in}}%
\pgfpathlineto{\pgfqpoint{3.085713in}{1.896984in}}%
\pgfpathlineto{\pgfqpoint{3.086104in}{1.532027in}}%
\pgfpathlineto{\pgfqpoint{3.086749in}{1.532040in}}%
\pgfpathlineto{\pgfqpoint{3.088063in}{1.532169in}}%
\pgfpathlineto{\pgfqpoint{3.088165in}{1.956384in}}%
\pgfpathlineto{\pgfqpoint{3.089169in}{1.937351in}}%
\pgfpathlineto{\pgfqpoint{3.089187in}{1.922859in}}%
\pgfpathlineto{\pgfqpoint{3.089935in}{1.933540in}}%
\pgfpathlineto{\pgfqpoint{3.090034in}{1.935894in}}%
\pgfpathlineto{\pgfqpoint{3.090825in}{1.665680in}}%
\pgfpathlineto{\pgfqpoint{3.091057in}{1.532027in}}%
\pgfpathlineto{\pgfqpoint{3.091721in}{1.532042in}}%
\pgfpathlineto{\pgfqpoint{3.093005in}{1.532119in}}%
\pgfpathlineto{\pgfqpoint{3.093046in}{1.736972in}}%
\pgfpathlineto{\pgfqpoint{3.093315in}{1.958021in}}%
\pgfpathlineto{\pgfqpoint{3.093825in}{1.920088in}}%
\pgfpathlineto{\pgfqpoint{3.093831in}{1.916675in}}%
\pgfpathlineto{\pgfqpoint{3.093846in}{1.942754in}}%
\pgfpathlineto{\pgfqpoint{3.094572in}{1.924732in}}%
\pgfpathlineto{\pgfqpoint{3.094798in}{1.935563in}}%
\pgfpathlineto{\pgfqpoint{3.094641in}{1.924352in}}%
\pgfpathlineto{\pgfqpoint{3.095342in}{1.933874in}}%
\pgfpathlineto{\pgfqpoint{3.095688in}{1.779437in}}%
\pgfpathlineto{\pgfqpoint{3.096024in}{1.532032in}}%
\pgfpathlineto{\pgfqpoint{3.096608in}{1.532044in}}%
\pgfpathlineto{\pgfqpoint{3.097955in}{1.532236in}}%
\pgfpathlineto{\pgfqpoint{3.098270in}{1.963606in}}%
\pgfpathlineto{\pgfqpoint{3.099040in}{1.941485in}}%
\pgfpathlineto{\pgfqpoint{3.099056in}{1.919008in}}%
\pgfpathlineto{\pgfqpoint{3.099815in}{1.927763in}}%
\pgfpathlineto{\pgfqpoint{3.100009in}{1.933897in}}%
\pgfpathlineto{\pgfqpoint{3.100602in}{1.808209in}}%
\pgfpathlineto{\pgfqpoint{3.100981in}{1.532032in}}%
\pgfpathlineto{\pgfqpoint{3.101519in}{1.532045in}}%
\pgfpathlineto{\pgfqpoint{3.102899in}{1.532157in}}%
\pgfpathlineto{\pgfqpoint{3.102984in}{1.953093in}}%
\pgfpathlineto{\pgfqpoint{3.104087in}{1.927307in}}%
\pgfpathlineto{\pgfqpoint{3.104092in}{1.924754in}}%
\pgfpathlineto{\pgfqpoint{3.104107in}{1.935190in}}%
\pgfpathlineto{\pgfqpoint{3.104846in}{1.931798in}}%
\pgfpathlineto{\pgfqpoint{3.105926in}{1.532032in}}%
\pgfpathlineto{\pgfqpoint{3.107154in}{1.532060in}}%
\pgfpathlineto{\pgfqpoint{3.107845in}{1.532183in}}%
\pgfpathlineto{\pgfqpoint{3.108171in}{1.962485in}}%
\pgfpathlineto{\pgfqpoint{3.108979in}{1.928288in}}%
\pgfpathlineto{\pgfqpoint{3.109017in}{1.922004in}}%
\pgfpathlineto{\pgfqpoint{3.109030in}{1.937818in}}%
\pgfpathlineto{\pgfqpoint{3.109748in}{1.930233in}}%
\pgfpathlineto{\pgfqpoint{3.110853in}{1.532029in}}%
\pgfpathlineto{\pgfqpoint{3.109842in}{1.932319in}}%
\pgfpathlineto{\pgfqpoint{3.111483in}{1.532050in}}%
\pgfpathlineto{\pgfqpoint{3.112789in}{1.532143in}}%
\pgfpathlineto{\pgfqpoint{3.112879in}{1.955533in}}%
\pgfpathlineto{\pgfqpoint{3.113973in}{1.930798in}}%
\pgfpathlineto{\pgfqpoint{3.113984in}{1.923291in}}%
\pgfpathlineto{\pgfqpoint{3.114002in}{1.936636in}}%
\pgfpathlineto{\pgfqpoint{3.114712in}{1.932437in}}%
\pgfpathlineto{\pgfqpoint{3.114959in}{1.934433in}}%
\pgfpathlineto{\pgfqpoint{3.115404in}{1.857887in}}%
\pgfpathlineto{\pgfqpoint{3.115812in}{1.532028in}}%
\pgfpathlineto{\pgfqpoint{3.116399in}{1.532041in}}%
\pgfpathlineto{\pgfqpoint{3.117735in}{1.532156in}}%
\pgfpathlineto{\pgfqpoint{3.117826in}{1.955182in}}%
\pgfpathlineto{\pgfqpoint{3.118860in}{1.926220in}}%
\pgfpathlineto{\pgfqpoint{3.119017in}{1.942572in}}%
\pgfpathlineto{\pgfqpoint{3.119064in}{1.918085in}}%
\pgfpathlineto{\pgfqpoint{3.119459in}{1.942232in}}%
\pgfpathlineto{\pgfqpoint{3.120753in}{1.532027in}}%
\pgfpathlineto{\pgfqpoint{3.121041in}{1.532033in}}%
\pgfpathlineto{\pgfqpoint{3.122680in}{1.532145in}}%
\pgfpathlineto{\pgfqpoint{3.122770in}{1.954337in}}%
\pgfpathlineto{\pgfqpoint{3.123839in}{1.933036in}}%
\pgfpathlineto{\pgfqpoint{3.123849in}{1.923118in}}%
\pgfpathlineto{\pgfqpoint{3.123870in}{1.937096in}}%
\pgfpathlineto{\pgfqpoint{3.124597in}{1.930648in}}%
\pgfpathlineto{\pgfqpoint{3.124696in}{1.927257in}}%
\pgfpathlineto{\pgfqpoint{3.124745in}{1.933723in}}%
\pgfpathlineto{\pgfqpoint{3.124844in}{1.934610in}}%
\pgfpathlineto{\pgfqpoint{3.125388in}{1.732808in}}%
\pgfpathlineto{\pgfqpoint{3.125696in}{1.532028in}}%
\pgfpathlineto{\pgfqpoint{3.126280in}{1.532041in}}%
\pgfpathlineto{\pgfqpoint{3.127626in}{1.532154in}}%
\pgfpathlineto{\pgfqpoint{3.127717in}{1.955289in}}%
\pgfpathlineto{\pgfqpoint{3.128786in}{1.941130in}}%
\pgfpathlineto{\pgfqpoint{3.129173in}{1.918835in}}%
\pgfpathlineto{\pgfqpoint{3.129519in}{1.921357in}}%
\pgfpathlineto{\pgfqpoint{3.129568in}{1.940304in}}%
\pgfpathlineto{\pgfqpoint{3.129618in}{1.919025in}}%
\pgfpathlineto{\pgfqpoint{3.130160in}{1.921705in}}%
\pgfpathlineto{\pgfqpoint{3.130633in}{1.532028in}}%
\pgfpathlineto{\pgfqpoint{3.131260in}{1.532049in}}%
\pgfpathlineto{\pgfqpoint{3.132569in}{1.532119in}}%
\pgfpathlineto{\pgfqpoint{3.132874in}{1.957243in}}%
\pgfpathlineto{\pgfqpoint{3.134143in}{1.935712in}}%
\pgfpathlineto{\pgfqpoint{3.134161in}{1.924487in}}%
\pgfpathlineto{\pgfqpoint{3.134924in}{1.925406in}}%
\pgfpathlineto{\pgfqpoint{3.134974in}{1.934214in}}%
\pgfpathlineto{\pgfqpoint{3.135172in}{1.879537in}}%
\pgfpathlineto{\pgfqpoint{3.135573in}{1.532029in}}%
\pgfpathlineto{\pgfqpoint{3.136195in}{1.532049in}}%
\pgfpathlineto{\pgfqpoint{3.137515in}{1.532132in}}%
\pgfpathlineto{\pgfqpoint{3.137722in}{1.959133in}}%
\pgfpathlineto{\pgfqpoint{3.138808in}{1.935985in}}%
\pgfpathlineto{\pgfqpoint{3.138823in}{1.924234in}}%
\pgfpathlineto{\pgfqpoint{3.139582in}{1.931913in}}%
\pgfpathlineto{\pgfqpoint{3.140483in}{1.532022in}}%
\pgfpathlineto{\pgfqpoint{3.141970in}{1.532071in}}%
\pgfpathlineto{\pgfqpoint{3.142450in}{1.532082in}}%
\pgfpathlineto{\pgfqpoint{3.142460in}{1.532130in}}%
\pgfpathlineto{\pgfqpoint{3.142698in}{1.961385in}}%
\pgfpathlineto{\pgfqpoint{3.143728in}{1.934154in}}%
\pgfpathlineto{\pgfqpoint{3.143738in}{1.921734in}}%
\pgfpathlineto{\pgfqpoint{3.143754in}{1.938136in}}%
\pgfpathlineto{\pgfqpoint{3.144495in}{1.927100in}}%
\pgfpathlineto{\pgfqpoint{3.144512in}{1.932513in}}%
\pgfpathlineto{\pgfqpoint{3.145047in}{1.892376in}}%
\pgfpathlineto{\pgfqpoint{3.145472in}{1.532028in}}%
\pgfpathlineto{\pgfqpoint{3.146103in}{1.532050in}}%
\pgfpathlineto{\pgfqpoint{3.147407in}{1.532147in}}%
\pgfpathlineto{\pgfqpoint{3.147501in}{1.956274in}}%
\pgfpathlineto{\pgfqpoint{3.148574in}{1.928521in}}%
\pgfpathlineto{\pgfqpoint{3.148585in}{1.924005in}}%
\pgfpathlineto{\pgfqpoint{3.148710in}{1.935980in}}%
\pgfpathlineto{\pgfqpoint{3.149301in}{1.924427in}}%
\pgfpathlineto{\pgfqpoint{3.149697in}{1.934925in}}%
\pgfpathlineto{\pgfqpoint{3.149746in}{1.924243in}}%
\pgfpathlineto{\pgfqpoint{3.149944in}{1.928109in}}%
\pgfpathlineto{\pgfqpoint{3.150401in}{1.532027in}}%
\pgfpathlineto{\pgfqpoint{3.151211in}{1.532054in}}%
\pgfpathlineto{\pgfqpoint{3.152350in}{1.532119in}}%
\pgfpathlineto{\pgfqpoint{3.152391in}{1.735154in}}%
\pgfpathlineto{\pgfqpoint{3.152670in}{1.956953in}}%
\pgfpathlineto{\pgfqpoint{3.153176in}{1.944928in}}%
\pgfpathlineto{\pgfqpoint{3.153192in}{1.915531in}}%
\pgfpathlineto{\pgfqpoint{3.153957in}{1.926701in}}%
\pgfpathlineto{\pgfqpoint{3.154626in}{1.924438in}}%
\pgfpathlineto{\pgfqpoint{3.154053in}{1.935642in}}%
\pgfpathlineto{\pgfqpoint{3.154725in}{1.926336in}}%
\pgfpathlineto{\pgfqpoint{3.154774in}{1.930616in}}%
\pgfpathlineto{\pgfqpoint{3.154923in}{1.899527in}}%
\pgfpathlineto{\pgfqpoint{3.155350in}{1.532031in}}%
\pgfpathlineto{\pgfqpoint{3.156111in}{1.532048in}}%
\pgfpathlineto{\pgfqpoint{3.157298in}{1.532151in}}%
\pgfpathlineto{\pgfqpoint{3.157604in}{1.958019in}}%
\pgfpathlineto{\pgfqpoint{3.158508in}{1.937667in}}%
\pgfpathlineto{\pgfqpoint{3.158526in}{1.923229in}}%
\pgfpathlineto{\pgfqpoint{3.159288in}{1.927788in}}%
\pgfpathlineto{\pgfqpoint{3.159300in}{1.931935in}}%
\pgfpathlineto{\pgfqpoint{3.159871in}{1.907402in}}%
\pgfpathlineto{\pgfqpoint{3.160288in}{1.532031in}}%
\pgfpathlineto{\pgfqpoint{3.160986in}{1.532048in}}%
\pgfpathlineto{\pgfqpoint{3.162241in}{1.532119in}}%
\pgfpathlineto{\pgfqpoint{3.162342in}{1.956571in}}%
\pgfpathlineto{\pgfqpoint{3.163769in}{1.933783in}}%
\pgfpathlineto{\pgfqpoint{3.164264in}{1.926304in}}%
\pgfpathlineto{\pgfqpoint{3.164511in}{1.928136in}}%
\pgfpathlineto{\pgfqpoint{3.164659in}{1.932881in}}%
\pgfpathlineto{\pgfqpoint{3.164709in}{1.925816in}}%
\pgfpathlineto{\pgfqpoint{3.164808in}{1.925902in}}%
\pgfpathlineto{\pgfqpoint{3.165210in}{1.532022in}}%
\pgfpathlineto{\pgfqpoint{3.166014in}{1.532056in}}%
\pgfpathlineto{\pgfqpoint{3.167185in}{1.532109in}}%
\pgfpathlineto{\pgfqpoint{3.167225in}{1.707329in}}%
\pgfpathlineto{\pgfqpoint{3.167490in}{1.956645in}}%
\pgfpathlineto{\pgfqpoint{3.168011in}{1.915377in}}%
\pgfpathlineto{\pgfqpoint{3.168027in}{1.944213in}}%
\pgfpathlineto{\pgfqpoint{3.168780in}{1.937028in}}%
\pgfpathlineto{\pgfqpoint{3.169587in}{1.922861in}}%
\pgfpathlineto{\pgfqpoint{3.169636in}{1.935191in}}%
\pgfpathlineto{\pgfqpoint{3.169834in}{1.821589in}}%
\pgfpathlineto{\pgfqpoint{3.170082in}{1.619489in}}%
\pgfpathlineto{\pgfqpoint{3.170132in}{1.633524in}}%
\pgfusepath{stroke}%
\end{pgfscope}%
\begin{pgfscope}%
\pgfsetrectcap%
\pgfsetmiterjoin%
\pgfsetlinewidth{0.803000pt}%
\definecolor{currentstroke}{rgb}{0.000000,0.000000,0.000000}%
\pgfsetstrokecolor{currentstroke}%
\pgfsetdash{}{0pt}%
\pgfpathmoveto{\pgfqpoint{0.573768in}{1.532029in}}%
\pgfpathlineto{\pgfqpoint{0.573768in}{2.438279in}}%
\pgfusepath{stroke}%
\end{pgfscope}%
\begin{pgfscope}%
\pgfsetrectcap%
\pgfsetmiterjoin%
\pgfsetlinewidth{0.803000pt}%
\definecolor{currentstroke}{rgb}{0.000000,0.000000,0.000000}%
\pgfsetstrokecolor{currentstroke}%
\pgfsetdash{}{0pt}%
\pgfpathmoveto{\pgfqpoint{3.293768in}{1.532029in}}%
\pgfpathlineto{\pgfqpoint{3.293768in}{2.438279in}}%
\pgfusepath{stroke}%
\end{pgfscope}%
\begin{pgfscope}%
\pgfsetrectcap%
\pgfsetmiterjoin%
\pgfsetlinewidth{0.803000pt}%
\definecolor{currentstroke}{rgb}{0.000000,0.000000,0.000000}%
\pgfsetstrokecolor{currentstroke}%
\pgfsetdash{}{0pt}%
\pgfpathmoveto{\pgfqpoint{0.573768in}{1.532029in}}%
\pgfpathlineto{\pgfqpoint{3.293768in}{1.532029in}}%
\pgfusepath{stroke}%
\end{pgfscope}%
\begin{pgfscope}%
\pgfsetrectcap%
\pgfsetmiterjoin%
\pgfsetlinewidth{0.803000pt}%
\definecolor{currentstroke}{rgb}{0.000000,0.000000,0.000000}%
\pgfsetstrokecolor{currentstroke}%
\pgfsetdash{}{0pt}%
\pgfpathmoveto{\pgfqpoint{0.573768in}{2.438279in}}%
\pgfpathlineto{\pgfqpoint{3.293768in}{2.438279in}}%
\pgfusepath{stroke}%
\end{pgfscope}%
\begin{pgfscope}%
\definecolor{textcolor}{rgb}{0.000000,0.000000,0.000000}%
\pgfsetstrokecolor{textcolor}%
\pgfsetfillcolor{textcolor}%
\pgftext[x=0.628168in,y=2.374842in,left,top]{\color{textcolor}{\rmfamily\fontsize{8.000000}{9.600000}\selectfont\catcode`\^=\active\def^{\ifmmode\sp\else\^{}\fi}\catcode`\%=\active\def%{\%}(b)}}%
\end{pgfscope}%
\begin{pgfscope}%
\pgfsetbuttcap%
\pgfsetmiterjoin%
\definecolor{currentfill}{rgb}{1.000000,1.000000,1.000000}%
\pgfsetfillcolor{currentfill}%
\pgfsetlinewidth{0.000000pt}%
\definecolor{currentstroke}{rgb}{0.000000,0.000000,0.000000}%
\pgfsetstrokecolor{currentstroke}%
\pgfsetstrokeopacity{0.000000}%
\pgfsetdash{}{0pt}%
\pgfpathmoveto{\pgfqpoint{0.573768in}{0.462654in}}%
\pgfpathlineto{\pgfqpoint{3.293768in}{0.462654in}}%
\pgfpathlineto{\pgfqpoint{3.293768in}{1.368904in}}%
\pgfpathlineto{\pgfqpoint{0.573768in}{1.368904in}}%
\pgfpathlineto{\pgfqpoint{0.573768in}{0.462654in}}%
\pgfpathclose%
\pgfusepath{fill}%
\end{pgfscope}%
\begin{pgfscope}%
\pgfpathrectangle{\pgfqpoint{0.573768in}{0.462654in}}{\pgfqpoint{2.720000in}{0.906250in}}%
\pgfusepath{clip}%
\pgfsetbuttcap%
\pgfsetroundjoin%
\pgfsetlinewidth{0.501875pt}%
\definecolor{currentstroke}{rgb}{0.690196,0.690196,0.690196}%
\pgfsetstrokecolor{currentstroke}%
\pgfsetstrokeopacity{0.250000}%
\pgfsetdash{{1.850000pt}{0.800000pt}}{0.000000pt}%
\pgfpathmoveto{\pgfqpoint{0.697405in}{0.462654in}}%
\pgfpathlineto{\pgfqpoint{0.697405in}{1.368904in}}%
\pgfusepath{stroke}%
\end{pgfscope}%
\begin{pgfscope}%
\pgfsetbuttcap%
\pgfsetroundjoin%
\definecolor{currentfill}{rgb}{0.000000,0.000000,0.000000}%
\pgfsetfillcolor{currentfill}%
\pgfsetlinewidth{0.803000pt}%
\definecolor{currentstroke}{rgb}{0.000000,0.000000,0.000000}%
\pgfsetstrokecolor{currentstroke}%
\pgfsetdash{}{0pt}%
\pgfsys@defobject{currentmarker}{\pgfqpoint{0.000000in}{-0.048611in}}{\pgfqpoint{0.000000in}{0.000000in}}{%
\pgfpathmoveto{\pgfqpoint{0.000000in}{0.000000in}}%
\pgfpathlineto{\pgfqpoint{0.000000in}{-0.048611in}}%
\pgfusepath{stroke,fill}%
}%
\begin{pgfscope}%
\pgfsys@transformshift{0.697405in}{0.462654in}%
\pgfsys@useobject{currentmarker}{}%
\end{pgfscope}%
\end{pgfscope}%
\begin{pgfscope}%
\definecolor{textcolor}{rgb}{0.000000,0.000000,0.000000}%
\pgfsetstrokecolor{textcolor}%
\pgfsetfillcolor{textcolor}%
\pgftext[x=0.697405in,y=0.365432in,,top]{\color{textcolor}{\rmfamily\fontsize{8.000000}{9.600000}\selectfont\catcode`\^=\active\def^{\ifmmode\sp\else\^{}\fi}\catcode`\%=\active\def%{\%}$\mathdefault{0}$}}%
\end{pgfscope}%
\begin{pgfscope}%
\pgfpathrectangle{\pgfqpoint{0.573768in}{0.462654in}}{\pgfqpoint{2.720000in}{0.906250in}}%
\pgfusepath{clip}%
\pgfsetbuttcap%
\pgfsetroundjoin%
\pgfsetlinewidth{0.501875pt}%
\definecolor{currentstroke}{rgb}{0.690196,0.690196,0.690196}%
\pgfsetstrokecolor{currentstroke}%
\pgfsetstrokeopacity{0.250000}%
\pgfsetdash{{1.850000pt}{0.800000pt}}{0.000000pt}%
\pgfpathmoveto{\pgfqpoint{1.191950in}{0.462654in}}%
\pgfpathlineto{\pgfqpoint{1.191950in}{1.368904in}}%
\pgfusepath{stroke}%
\end{pgfscope}%
\begin{pgfscope}%
\pgfsetbuttcap%
\pgfsetroundjoin%
\definecolor{currentfill}{rgb}{0.000000,0.000000,0.000000}%
\pgfsetfillcolor{currentfill}%
\pgfsetlinewidth{0.803000pt}%
\definecolor{currentstroke}{rgb}{0.000000,0.000000,0.000000}%
\pgfsetstrokecolor{currentstroke}%
\pgfsetdash{}{0pt}%
\pgfsys@defobject{currentmarker}{\pgfqpoint{0.000000in}{-0.048611in}}{\pgfqpoint{0.000000in}{0.000000in}}{%
\pgfpathmoveto{\pgfqpoint{0.000000in}{0.000000in}}%
\pgfpathlineto{\pgfqpoint{0.000000in}{-0.048611in}}%
\pgfusepath{stroke,fill}%
}%
\begin{pgfscope}%
\pgfsys@transformshift{1.191950in}{0.462654in}%
\pgfsys@useobject{currentmarker}{}%
\end{pgfscope}%
\end{pgfscope}%
\begin{pgfscope}%
\definecolor{textcolor}{rgb}{0.000000,0.000000,0.000000}%
\pgfsetstrokecolor{textcolor}%
\pgfsetfillcolor{textcolor}%
\pgftext[x=1.191950in,y=0.365432in,,top]{\color{textcolor}{\rmfamily\fontsize{8.000000}{9.600000}\selectfont\catcode`\^=\active\def^{\ifmmode\sp\else\^{}\fi}\catcode`\%=\active\def%{\%}$\mathdefault{1}$}}%
\end{pgfscope}%
\begin{pgfscope}%
\pgfpathrectangle{\pgfqpoint{0.573768in}{0.462654in}}{\pgfqpoint{2.720000in}{0.906250in}}%
\pgfusepath{clip}%
\pgfsetbuttcap%
\pgfsetroundjoin%
\pgfsetlinewidth{0.501875pt}%
\definecolor{currentstroke}{rgb}{0.690196,0.690196,0.690196}%
\pgfsetstrokecolor{currentstroke}%
\pgfsetstrokeopacity{0.250000}%
\pgfsetdash{{1.850000pt}{0.800000pt}}{0.000000pt}%
\pgfpathmoveto{\pgfqpoint{1.686496in}{0.462654in}}%
\pgfpathlineto{\pgfqpoint{1.686496in}{1.368904in}}%
\pgfusepath{stroke}%
\end{pgfscope}%
\begin{pgfscope}%
\pgfsetbuttcap%
\pgfsetroundjoin%
\definecolor{currentfill}{rgb}{0.000000,0.000000,0.000000}%
\pgfsetfillcolor{currentfill}%
\pgfsetlinewidth{0.803000pt}%
\definecolor{currentstroke}{rgb}{0.000000,0.000000,0.000000}%
\pgfsetstrokecolor{currentstroke}%
\pgfsetdash{}{0pt}%
\pgfsys@defobject{currentmarker}{\pgfqpoint{0.000000in}{-0.048611in}}{\pgfqpoint{0.000000in}{0.000000in}}{%
\pgfpathmoveto{\pgfqpoint{0.000000in}{0.000000in}}%
\pgfpathlineto{\pgfqpoint{0.000000in}{-0.048611in}}%
\pgfusepath{stroke,fill}%
}%
\begin{pgfscope}%
\pgfsys@transformshift{1.686496in}{0.462654in}%
\pgfsys@useobject{currentmarker}{}%
\end{pgfscope}%
\end{pgfscope}%
\begin{pgfscope}%
\definecolor{textcolor}{rgb}{0.000000,0.000000,0.000000}%
\pgfsetstrokecolor{textcolor}%
\pgfsetfillcolor{textcolor}%
\pgftext[x=1.686496in,y=0.365432in,,top]{\color{textcolor}{\rmfamily\fontsize{8.000000}{9.600000}\selectfont\catcode`\^=\active\def^{\ifmmode\sp\else\^{}\fi}\catcode`\%=\active\def%{\%}$\mathdefault{2}$}}%
\end{pgfscope}%
\begin{pgfscope}%
\pgfpathrectangle{\pgfqpoint{0.573768in}{0.462654in}}{\pgfqpoint{2.720000in}{0.906250in}}%
\pgfusepath{clip}%
\pgfsetbuttcap%
\pgfsetroundjoin%
\pgfsetlinewidth{0.501875pt}%
\definecolor{currentstroke}{rgb}{0.690196,0.690196,0.690196}%
\pgfsetstrokecolor{currentstroke}%
\pgfsetstrokeopacity{0.250000}%
\pgfsetdash{{1.850000pt}{0.800000pt}}{0.000000pt}%
\pgfpathmoveto{\pgfqpoint{2.181041in}{0.462654in}}%
\pgfpathlineto{\pgfqpoint{2.181041in}{1.368904in}}%
\pgfusepath{stroke}%
\end{pgfscope}%
\begin{pgfscope}%
\pgfsetbuttcap%
\pgfsetroundjoin%
\definecolor{currentfill}{rgb}{0.000000,0.000000,0.000000}%
\pgfsetfillcolor{currentfill}%
\pgfsetlinewidth{0.803000pt}%
\definecolor{currentstroke}{rgb}{0.000000,0.000000,0.000000}%
\pgfsetstrokecolor{currentstroke}%
\pgfsetdash{}{0pt}%
\pgfsys@defobject{currentmarker}{\pgfqpoint{0.000000in}{-0.048611in}}{\pgfqpoint{0.000000in}{0.000000in}}{%
\pgfpathmoveto{\pgfqpoint{0.000000in}{0.000000in}}%
\pgfpathlineto{\pgfqpoint{0.000000in}{-0.048611in}}%
\pgfusepath{stroke,fill}%
}%
\begin{pgfscope}%
\pgfsys@transformshift{2.181041in}{0.462654in}%
\pgfsys@useobject{currentmarker}{}%
\end{pgfscope}%
\end{pgfscope}%
\begin{pgfscope}%
\definecolor{textcolor}{rgb}{0.000000,0.000000,0.000000}%
\pgfsetstrokecolor{textcolor}%
\pgfsetfillcolor{textcolor}%
\pgftext[x=2.181041in,y=0.365432in,,top]{\color{textcolor}{\rmfamily\fontsize{8.000000}{9.600000}\selectfont\catcode`\^=\active\def^{\ifmmode\sp\else\^{}\fi}\catcode`\%=\active\def%{\%}$\mathdefault{3}$}}%
\end{pgfscope}%
\begin{pgfscope}%
\pgfpathrectangle{\pgfqpoint{0.573768in}{0.462654in}}{\pgfqpoint{2.720000in}{0.906250in}}%
\pgfusepath{clip}%
\pgfsetbuttcap%
\pgfsetroundjoin%
\pgfsetlinewidth{0.501875pt}%
\definecolor{currentstroke}{rgb}{0.690196,0.690196,0.690196}%
\pgfsetstrokecolor{currentstroke}%
\pgfsetstrokeopacity{0.250000}%
\pgfsetdash{{1.850000pt}{0.800000pt}}{0.000000pt}%
\pgfpathmoveto{\pgfqpoint{2.675587in}{0.462654in}}%
\pgfpathlineto{\pgfqpoint{2.675587in}{1.368904in}}%
\pgfusepath{stroke}%
\end{pgfscope}%
\begin{pgfscope}%
\pgfsetbuttcap%
\pgfsetroundjoin%
\definecolor{currentfill}{rgb}{0.000000,0.000000,0.000000}%
\pgfsetfillcolor{currentfill}%
\pgfsetlinewidth{0.803000pt}%
\definecolor{currentstroke}{rgb}{0.000000,0.000000,0.000000}%
\pgfsetstrokecolor{currentstroke}%
\pgfsetdash{}{0pt}%
\pgfsys@defobject{currentmarker}{\pgfqpoint{0.000000in}{-0.048611in}}{\pgfqpoint{0.000000in}{0.000000in}}{%
\pgfpathmoveto{\pgfqpoint{0.000000in}{0.000000in}}%
\pgfpathlineto{\pgfqpoint{0.000000in}{-0.048611in}}%
\pgfusepath{stroke,fill}%
}%
\begin{pgfscope}%
\pgfsys@transformshift{2.675587in}{0.462654in}%
\pgfsys@useobject{currentmarker}{}%
\end{pgfscope}%
\end{pgfscope}%
\begin{pgfscope}%
\definecolor{textcolor}{rgb}{0.000000,0.000000,0.000000}%
\pgfsetstrokecolor{textcolor}%
\pgfsetfillcolor{textcolor}%
\pgftext[x=2.675587in,y=0.365432in,,top]{\color{textcolor}{\rmfamily\fontsize{8.000000}{9.600000}\selectfont\catcode`\^=\active\def^{\ifmmode\sp\else\^{}\fi}\catcode`\%=\active\def%{\%}$\mathdefault{4}$}}%
\end{pgfscope}%
\begin{pgfscope}%
\pgfpathrectangle{\pgfqpoint{0.573768in}{0.462654in}}{\pgfqpoint{2.720000in}{0.906250in}}%
\pgfusepath{clip}%
\pgfsetbuttcap%
\pgfsetroundjoin%
\pgfsetlinewidth{0.501875pt}%
\definecolor{currentstroke}{rgb}{0.690196,0.690196,0.690196}%
\pgfsetstrokecolor{currentstroke}%
\pgfsetstrokeopacity{0.250000}%
\pgfsetdash{{1.850000pt}{0.800000pt}}{0.000000pt}%
\pgfpathmoveto{\pgfqpoint{3.170132in}{0.462654in}}%
\pgfpathlineto{\pgfqpoint{3.170132in}{1.368904in}}%
\pgfusepath{stroke}%
\end{pgfscope}%
\begin{pgfscope}%
\pgfsetbuttcap%
\pgfsetroundjoin%
\definecolor{currentfill}{rgb}{0.000000,0.000000,0.000000}%
\pgfsetfillcolor{currentfill}%
\pgfsetlinewidth{0.803000pt}%
\definecolor{currentstroke}{rgb}{0.000000,0.000000,0.000000}%
\pgfsetstrokecolor{currentstroke}%
\pgfsetdash{}{0pt}%
\pgfsys@defobject{currentmarker}{\pgfqpoint{0.000000in}{-0.048611in}}{\pgfqpoint{0.000000in}{0.000000in}}{%
\pgfpathmoveto{\pgfqpoint{0.000000in}{0.000000in}}%
\pgfpathlineto{\pgfqpoint{0.000000in}{-0.048611in}}%
\pgfusepath{stroke,fill}%
}%
\begin{pgfscope}%
\pgfsys@transformshift{3.170132in}{0.462654in}%
\pgfsys@useobject{currentmarker}{}%
\end{pgfscope}%
\end{pgfscope}%
\begin{pgfscope}%
\definecolor{textcolor}{rgb}{0.000000,0.000000,0.000000}%
\pgfsetstrokecolor{textcolor}%
\pgfsetfillcolor{textcolor}%
\pgftext[x=3.170132in,y=0.365432in,,top]{\color{textcolor}{\rmfamily\fontsize{8.000000}{9.600000}\selectfont\catcode`\^=\active\def^{\ifmmode\sp\else\^{}\fi}\catcode`\%=\active\def%{\%}$\mathdefault{5}$}}%
\end{pgfscope}%
\begin{pgfscope}%
\definecolor{textcolor}{rgb}{0.000000,0.000000,0.000000}%
\pgfsetstrokecolor{textcolor}%
\pgfsetfillcolor{textcolor}%
\pgftext[x=1.933768in,y=0.211111in,,top]{\color{textcolor}{\rmfamily\fontsize{8.000000}{9.600000}\selectfont\catcode`\^=\active\def^{\ifmmode\sp\else\^{}\fi}\catcode`\%=\active\def%{\%}Time (ms)}}%
\end{pgfscope}%
\begin{pgfscope}%
\pgfpathrectangle{\pgfqpoint{0.573768in}{0.462654in}}{\pgfqpoint{2.720000in}{0.906250in}}%
\pgfusepath{clip}%
\pgfsetbuttcap%
\pgfsetroundjoin%
\pgfsetlinewidth{0.501875pt}%
\definecolor{currentstroke}{rgb}{0.690196,0.690196,0.690196}%
\pgfsetstrokecolor{currentstroke}%
\pgfsetstrokeopacity{0.250000}%
\pgfsetdash{{1.850000pt}{0.800000pt}}{0.000000pt}%
\pgfpathmoveto{\pgfqpoint{0.573768in}{0.462654in}}%
\pgfpathlineto{\pgfqpoint{3.293768in}{0.462654in}}%
\pgfusepath{stroke}%
\end{pgfscope}%
\begin{pgfscope}%
\pgfsetbuttcap%
\pgfsetroundjoin%
\definecolor{currentfill}{rgb}{0.000000,0.000000,0.000000}%
\pgfsetfillcolor{currentfill}%
\pgfsetlinewidth{0.803000pt}%
\definecolor{currentstroke}{rgb}{0.000000,0.000000,0.000000}%
\pgfsetstrokecolor{currentstroke}%
\pgfsetdash{}{0pt}%
\pgfsys@defobject{currentmarker}{\pgfqpoint{-0.048611in}{0.000000in}}{\pgfqpoint{-0.000000in}{0.000000in}}{%
\pgfpathmoveto{\pgfqpoint{-0.000000in}{0.000000in}}%
\pgfpathlineto{\pgfqpoint{-0.048611in}{0.000000in}}%
\pgfusepath{stroke,fill}%
}%
\begin{pgfscope}%
\pgfsys@transformshift{0.573768in}{0.462654in}%
\pgfsys@useobject{currentmarker}{}%
\end{pgfscope}%
\end{pgfscope}%
\begin{pgfscope}%
\definecolor{textcolor}{rgb}{0.000000,0.000000,0.000000}%
\pgfsetstrokecolor{textcolor}%
\pgfsetfillcolor{textcolor}%
\pgftext[x=0.325695in, y=0.424074in, left, base]{\color{textcolor}{\rmfamily\fontsize{8.000000}{9.600000}\selectfont\catcode`\^=\active\def^{\ifmmode\sp\else\^{}\fi}\catcode`\%=\active\def%{\%}0.0}}%
\end{pgfscope}%
\begin{pgfscope}%
\pgfpathrectangle{\pgfqpoint{0.573768in}{0.462654in}}{\pgfqpoint{2.720000in}{0.906250in}}%
\pgfusepath{clip}%
\pgfsetbuttcap%
\pgfsetroundjoin%
\pgfsetlinewidth{0.501875pt}%
\definecolor{currentstroke}{rgb}{0.690196,0.690196,0.690196}%
\pgfsetstrokecolor{currentstroke}%
\pgfsetstrokeopacity{0.250000}%
\pgfsetdash{{1.850000pt}{0.800000pt}}{0.000000pt}%
\pgfpathmoveto{\pgfqpoint{0.573768in}{0.689217in}}%
\pgfpathlineto{\pgfqpoint{3.293768in}{0.689217in}}%
\pgfusepath{stroke}%
\end{pgfscope}%
\begin{pgfscope}%
\pgfsetbuttcap%
\pgfsetroundjoin%
\definecolor{currentfill}{rgb}{0.000000,0.000000,0.000000}%
\pgfsetfillcolor{currentfill}%
\pgfsetlinewidth{0.803000pt}%
\definecolor{currentstroke}{rgb}{0.000000,0.000000,0.000000}%
\pgfsetstrokecolor{currentstroke}%
\pgfsetdash{}{0pt}%
\pgfsys@defobject{currentmarker}{\pgfqpoint{-0.048611in}{0.000000in}}{\pgfqpoint{-0.000000in}{0.000000in}}{%
\pgfpathmoveto{\pgfqpoint{-0.000000in}{0.000000in}}%
\pgfpathlineto{\pgfqpoint{-0.048611in}{0.000000in}}%
\pgfusepath{stroke,fill}%
}%
\begin{pgfscope}%
\pgfsys@transformshift{0.573768in}{0.689217in}%
\pgfsys@useobject{currentmarker}{}%
\end{pgfscope}%
\end{pgfscope}%
\begin{pgfscope}%
\definecolor{textcolor}{rgb}{0.000000,0.000000,0.000000}%
\pgfsetstrokecolor{textcolor}%
\pgfsetfillcolor{textcolor}%
\pgftext[x=0.325695in, y=0.650637in, left, base]{\color{textcolor}{\rmfamily\fontsize{8.000000}{9.600000}\selectfont\catcode`\^=\active\def^{\ifmmode\sp\else\^{}\fi}\catcode`\%=\active\def%{\%}6.2}}%
\end{pgfscope}%
\begin{pgfscope}%
\pgfpathrectangle{\pgfqpoint{0.573768in}{0.462654in}}{\pgfqpoint{2.720000in}{0.906250in}}%
\pgfusepath{clip}%
\pgfsetbuttcap%
\pgfsetroundjoin%
\pgfsetlinewidth{0.501875pt}%
\definecolor{currentstroke}{rgb}{0.690196,0.690196,0.690196}%
\pgfsetstrokecolor{currentstroke}%
\pgfsetstrokeopacity{0.250000}%
\pgfsetdash{{1.850000pt}{0.800000pt}}{0.000000pt}%
\pgfpathmoveto{\pgfqpoint{0.573768in}{0.915779in}}%
\pgfpathlineto{\pgfqpoint{3.293768in}{0.915779in}}%
\pgfusepath{stroke}%
\end{pgfscope}%
\begin{pgfscope}%
\pgfsetbuttcap%
\pgfsetroundjoin%
\definecolor{currentfill}{rgb}{0.000000,0.000000,0.000000}%
\pgfsetfillcolor{currentfill}%
\pgfsetlinewidth{0.803000pt}%
\definecolor{currentstroke}{rgb}{0.000000,0.000000,0.000000}%
\pgfsetstrokecolor{currentstroke}%
\pgfsetdash{}{0pt}%
\pgfsys@defobject{currentmarker}{\pgfqpoint{-0.048611in}{0.000000in}}{\pgfqpoint{-0.000000in}{0.000000in}}{%
\pgfpathmoveto{\pgfqpoint{-0.000000in}{0.000000in}}%
\pgfpathlineto{\pgfqpoint{-0.048611in}{0.000000in}}%
\pgfusepath{stroke,fill}%
}%
\begin{pgfscope}%
\pgfsys@transformshift{0.573768in}{0.915779in}%
\pgfsys@useobject{currentmarker}{}%
\end{pgfscope}%
\end{pgfscope}%
\begin{pgfscope}%
\definecolor{textcolor}{rgb}{0.000000,0.000000,0.000000}%
\pgfsetstrokecolor{textcolor}%
\pgfsetfillcolor{textcolor}%
\pgftext[x=0.266667in, y=0.877199in, left, base]{\color{textcolor}{\rmfamily\fontsize{8.000000}{9.600000}\selectfont\catcode`\^=\active\def^{\ifmmode\sp\else\^{}\fi}\catcode`\%=\active\def%{\%}12.5}}%
\end{pgfscope}%
\begin{pgfscope}%
\pgfpathrectangle{\pgfqpoint{0.573768in}{0.462654in}}{\pgfqpoint{2.720000in}{0.906250in}}%
\pgfusepath{clip}%
\pgfsetbuttcap%
\pgfsetroundjoin%
\pgfsetlinewidth{0.501875pt}%
\definecolor{currentstroke}{rgb}{0.690196,0.690196,0.690196}%
\pgfsetstrokecolor{currentstroke}%
\pgfsetstrokeopacity{0.250000}%
\pgfsetdash{{1.850000pt}{0.800000pt}}{0.000000pt}%
\pgfpathmoveto{\pgfqpoint{0.573768in}{1.142342in}}%
\pgfpathlineto{\pgfqpoint{3.293768in}{1.142342in}}%
\pgfusepath{stroke}%
\end{pgfscope}%
\begin{pgfscope}%
\pgfsetbuttcap%
\pgfsetroundjoin%
\definecolor{currentfill}{rgb}{0.000000,0.000000,0.000000}%
\pgfsetfillcolor{currentfill}%
\pgfsetlinewidth{0.803000pt}%
\definecolor{currentstroke}{rgb}{0.000000,0.000000,0.000000}%
\pgfsetstrokecolor{currentstroke}%
\pgfsetdash{}{0pt}%
\pgfsys@defobject{currentmarker}{\pgfqpoint{-0.048611in}{0.000000in}}{\pgfqpoint{-0.000000in}{0.000000in}}{%
\pgfpathmoveto{\pgfqpoint{-0.000000in}{0.000000in}}%
\pgfpathlineto{\pgfqpoint{-0.048611in}{0.000000in}}%
\pgfusepath{stroke,fill}%
}%
\begin{pgfscope}%
\pgfsys@transformshift{0.573768in}{1.142342in}%
\pgfsys@useobject{currentmarker}{}%
\end{pgfscope}%
\end{pgfscope}%
\begin{pgfscope}%
\definecolor{textcolor}{rgb}{0.000000,0.000000,0.000000}%
\pgfsetstrokecolor{textcolor}%
\pgfsetfillcolor{textcolor}%
\pgftext[x=0.266667in, y=1.103762in, left, base]{\color{textcolor}{\rmfamily\fontsize{8.000000}{9.600000}\selectfont\catcode`\^=\active\def^{\ifmmode\sp\else\^{}\fi}\catcode`\%=\active\def%{\%}18.8}}%
\end{pgfscope}%
\begin{pgfscope}%
\pgfpathrectangle{\pgfqpoint{0.573768in}{0.462654in}}{\pgfqpoint{2.720000in}{0.906250in}}%
\pgfusepath{clip}%
\pgfsetbuttcap%
\pgfsetroundjoin%
\pgfsetlinewidth{0.501875pt}%
\definecolor{currentstroke}{rgb}{0.690196,0.690196,0.690196}%
\pgfsetstrokecolor{currentstroke}%
\pgfsetstrokeopacity{0.250000}%
\pgfsetdash{{1.850000pt}{0.800000pt}}{0.000000pt}%
\pgfpathmoveto{\pgfqpoint{0.573768in}{1.368904in}}%
\pgfpathlineto{\pgfqpoint{3.293768in}{1.368904in}}%
\pgfusepath{stroke}%
\end{pgfscope}%
\begin{pgfscope}%
\pgfsetbuttcap%
\pgfsetroundjoin%
\definecolor{currentfill}{rgb}{0.000000,0.000000,0.000000}%
\pgfsetfillcolor{currentfill}%
\pgfsetlinewidth{0.803000pt}%
\definecolor{currentstroke}{rgb}{0.000000,0.000000,0.000000}%
\pgfsetstrokecolor{currentstroke}%
\pgfsetdash{}{0pt}%
\pgfsys@defobject{currentmarker}{\pgfqpoint{-0.048611in}{0.000000in}}{\pgfqpoint{-0.000000in}{0.000000in}}{%
\pgfpathmoveto{\pgfqpoint{-0.000000in}{0.000000in}}%
\pgfpathlineto{\pgfqpoint{-0.048611in}{0.000000in}}%
\pgfusepath{stroke,fill}%
}%
\begin{pgfscope}%
\pgfsys@transformshift{0.573768in}{1.368904in}%
\pgfsys@useobject{currentmarker}{}%
\end{pgfscope}%
\end{pgfscope}%
\begin{pgfscope}%
\definecolor{textcolor}{rgb}{0.000000,0.000000,0.000000}%
\pgfsetstrokecolor{textcolor}%
\pgfsetfillcolor{textcolor}%
\pgftext[x=0.266667in, y=1.330324in, left, base]{\color{textcolor}{\rmfamily\fontsize{8.000000}{9.600000}\selectfont\catcode`\^=\active\def^{\ifmmode\sp\else\^{}\fi}\catcode`\%=\active\def%{\%}25.0}}%
\end{pgfscope}%
\begin{pgfscope}%
\definecolor{textcolor}{rgb}{0.000000,0.000000,0.000000}%
\pgfsetstrokecolor{textcolor}%
\pgfsetfillcolor{textcolor}%
\pgftext[x=0.211111in,y=0.915779in,,bottom,rotate=90.000000]{\color{textcolor}{\rmfamily\fontsize{8.000000}{9.600000}\selectfont\catcode`\^=\active\def^{\ifmmode\sp\else\^{}\fi}\catcode`\%=\active\def%{\%}$v_o$ (V)}}%
\end{pgfscope}%
\begin{pgfscope}%
\pgfpathrectangle{\pgfqpoint{0.573768in}{0.462654in}}{\pgfqpoint{2.720000in}{0.906250in}}%
\pgfusepath{clip}%
\pgfsetrectcap%
\pgfsetroundjoin%
\pgfsetlinewidth{1.505625pt}%
\definecolor{currentstroke}{rgb}{0.000000,0.619608,0.450980}%
\pgfsetstrokecolor{currentstroke}%
\pgfsetdash{}{0pt}%
\pgfpathmoveto{\pgfqpoint{0.697405in}{0.462654in}}%
\pgfpathlineto{\pgfqpoint{0.699634in}{0.463334in}}%
\pgfpathlineto{\pgfqpoint{0.703069in}{0.474846in}}%
\pgfpathlineto{\pgfqpoint{0.704500in}{0.475547in}}%
\pgfpathlineto{\pgfqpoint{0.708031in}{0.498563in}}%
\pgfpathlineto{\pgfqpoint{0.709421in}{0.499220in}}%
\pgfpathlineto{\pgfqpoint{0.712266in}{0.533040in}}%
\pgfpathlineto{\pgfqpoint{0.712953in}{0.532942in}}%
\pgfpathlineto{\pgfqpoint{0.714305in}{0.532823in}}%
\pgfpathlineto{\pgfqpoint{0.714352in}{0.533494in}}%
\pgfpathlineto{\pgfqpoint{0.717214in}{0.577111in}}%
\pgfpathlineto{\pgfqpoint{0.717787in}{0.576978in}}%
\pgfpathlineto{\pgfqpoint{0.719232in}{0.576665in}}%
\pgfpathlineto{\pgfqpoint{0.719270in}{0.577133in}}%
\pgfpathlineto{\pgfqpoint{0.722160in}{0.629565in}}%
\pgfpathlineto{\pgfqpoint{0.722608in}{0.629414in}}%
\pgfpathlineto{\pgfqpoint{0.724184in}{0.628915in}}%
\pgfpathlineto{\pgfqpoint{0.724217in}{0.629373in}}%
\pgfpathlineto{\pgfqpoint{0.727108in}{0.688662in}}%
\pgfpathlineto{\pgfqpoint{0.727573in}{0.688449in}}%
\pgfpathlineto{\pgfqpoint{0.729125in}{0.687759in}}%
\pgfpathlineto{\pgfqpoint{0.729173in}{0.688523in}}%
\pgfpathlineto{\pgfqpoint{0.732053in}{0.752764in}}%
\pgfpathlineto{\pgfqpoint{0.732633in}{0.752424in}}%
\pgfpathlineto{\pgfqpoint{0.734071in}{0.751598in}}%
\pgfpathlineto{\pgfqpoint{0.734119in}{0.752338in}}%
\pgfpathlineto{\pgfqpoint{0.736999in}{0.819968in}}%
\pgfpathlineto{\pgfqpoint{0.737381in}{0.819692in}}%
\pgfpathlineto{\pgfqpoint{0.739019in}{0.818544in}}%
\pgfpathlineto{\pgfqpoint{0.739067in}{0.819408in}}%
\pgfpathlineto{\pgfqpoint{0.741942in}{0.888489in}}%
\pgfpathlineto{\pgfqpoint{0.742468in}{0.888037in}}%
\pgfpathlineto{\pgfqpoint{0.743950in}{0.886765in}}%
\pgfpathlineto{\pgfqpoint{0.744005in}{0.887378in}}%
\pgfpathlineto{\pgfqpoint{0.746887in}{0.956353in}}%
\pgfpathlineto{\pgfqpoint{0.747218in}{0.956025in}}%
\pgfpathlineto{\pgfqpoint{0.748907in}{0.954344in}}%
\pgfpathlineto{\pgfqpoint{0.748963in}{0.955144in}}%
\pgfpathlineto{\pgfqpoint{0.751832in}{1.021567in}}%
\pgfpathlineto{\pgfqpoint{0.752172in}{1.021183in}}%
\pgfpathlineto{\pgfqpoint{0.753853in}{1.019289in}}%
\pgfpathlineto{\pgfqpoint{0.753916in}{1.020222in}}%
\pgfpathlineto{\pgfqpoint{0.756776in}{1.082527in}}%
\pgfpathlineto{\pgfqpoint{0.757221in}{1.081970in}}%
\pgfpathlineto{\pgfqpoint{0.758798in}{1.079999in}}%
\pgfpathlineto{\pgfqpoint{0.758862in}{1.080917in}}%
\pgfpathlineto{\pgfqpoint{0.761720in}{1.137683in}}%
\pgfpathlineto{\pgfqpoint{0.762162in}{1.137080in}}%
\pgfpathlineto{\pgfqpoint{0.763745in}{1.134925in}}%
\pgfpathlineto{\pgfqpoint{0.763816in}{1.135896in}}%
\pgfpathlineto{\pgfqpoint{0.766665in}{1.185506in}}%
\pgfpathlineto{\pgfqpoint{0.767038in}{1.184961in}}%
\pgfpathlineto{\pgfqpoint{0.768702in}{1.182562in}}%
\pgfpathlineto{\pgfqpoint{0.768773in}{1.183655in}}%
\pgfpathlineto{\pgfqpoint{0.771607in}{1.224894in}}%
\pgfpathlineto{\pgfqpoint{0.771778in}{1.224631in}}%
\pgfpathlineto{\pgfqpoint{0.773651in}{1.221780in}}%
\pgfpathlineto{\pgfqpoint{0.773741in}{1.222960in}}%
\pgfpathlineto{\pgfqpoint{0.776552in}{1.254735in}}%
\pgfpathlineto{\pgfqpoint{0.776649in}{1.254580in}}%
\pgfpathlineto{\pgfqpoint{0.778603in}{1.251477in}}%
\pgfpathlineto{\pgfqpoint{0.778721in}{1.252715in}}%
\pgfpathlineto{\pgfqpoint{0.781073in}{1.272436in}}%
\pgfpathlineto{\pgfqpoint{0.781493in}{1.274378in}}%
\pgfpathlineto{\pgfqpoint{0.782036in}{1.273491in}}%
\pgfpathlineto{\pgfqpoint{0.783540in}{1.271040in}}%
\pgfpathlineto{\pgfqpoint{0.783703in}{1.272199in}}%
\pgfpathlineto{\pgfqpoint{0.785922in}{1.282972in}}%
\pgfpathlineto{\pgfqpoint{0.786435in}{1.283547in}}%
\pgfpathlineto{\pgfqpoint{0.786754in}{1.283022in}}%
\pgfpathlineto{\pgfqpoint{0.788506in}{1.280147in}}%
\pgfpathlineto{\pgfqpoint{0.788761in}{1.281186in}}%
\pgfpathlineto{\pgfqpoint{0.790246in}{1.283985in}}%
\pgfpathlineto{\pgfqpoint{0.791038in}{1.283094in}}%
\pgfpathlineto{\pgfqpoint{0.793472in}{1.279096in}}%
\pgfpathlineto{\pgfqpoint{0.793848in}{1.279916in}}%
\pgfpathlineto{\pgfqpoint{0.794666in}{1.280418in}}%
\pgfpathlineto{\pgfqpoint{0.794879in}{1.280250in}}%
\pgfpathlineto{\pgfqpoint{0.796909in}{1.276961in}}%
\pgfpathlineto{\pgfqpoint{0.798436in}{1.274491in}}%
\pgfpathlineto{\pgfqpoint{0.798814in}{1.275311in}}%
\pgfpathlineto{\pgfqpoint{0.799659in}{1.275777in}}%
\pgfpathlineto{\pgfqpoint{0.799883in}{1.275575in}}%
\pgfpathlineto{\pgfqpoint{0.803374in}{1.269927in}}%
\pgfpathlineto{\pgfqpoint{0.804255in}{1.271274in}}%
\pgfpathlineto{\pgfqpoint{0.805025in}{1.270755in}}%
\pgfpathlineto{\pgfqpoint{0.808320in}{1.265415in}}%
\pgfpathlineto{\pgfqpoint{0.808695in}{1.266232in}}%
\pgfpathlineto{\pgfqpoint{0.809556in}{1.266735in}}%
\pgfpathlineto{\pgfqpoint{0.809754in}{1.266570in}}%
\pgfpathlineto{\pgfqpoint{0.812481in}{1.262203in}}%
\pgfpathlineto{\pgfqpoint{0.813265in}{1.260963in}}%
\pgfpathlineto{\pgfqpoint{0.813637in}{1.261776in}}%
\pgfpathlineto{\pgfqpoint{0.814534in}{1.262272in}}%
\pgfpathlineto{\pgfqpoint{0.814732in}{1.262094in}}%
\pgfpathlineto{\pgfqpoint{0.817937in}{1.256970in}}%
\pgfpathlineto{\pgfqpoint{0.818214in}{1.256570in}}%
\pgfpathlineto{\pgfqpoint{0.818586in}{1.257385in}}%
\pgfpathlineto{\pgfqpoint{0.819487in}{1.257875in}}%
\pgfpathlineto{\pgfqpoint{0.819636in}{1.257750in}}%
\pgfpathlineto{\pgfqpoint{0.822053in}{1.253935in}}%
\pgfpathlineto{\pgfqpoint{0.823149in}{1.252195in}}%
\pgfpathlineto{\pgfqpoint{0.823545in}{1.252991in}}%
\pgfpathlineto{\pgfqpoint{0.824413in}{1.253497in}}%
\pgfpathlineto{\pgfqpoint{0.824611in}{1.253334in}}%
\pgfpathlineto{\pgfqpoint{0.826798in}{1.249897in}}%
\pgfpathlineto{\pgfqpoint{0.828093in}{1.247872in}}%
\pgfpathlineto{\pgfqpoint{0.828462in}{1.248695in}}%
\pgfpathlineto{\pgfqpoint{0.829398in}{1.249224in}}%
\pgfpathlineto{\pgfqpoint{0.829596in}{1.249042in}}%
\pgfpathlineto{\pgfqpoint{0.832534in}{1.244417in}}%
\pgfpathlineto{\pgfqpoint{0.833050in}{1.243634in}}%
\pgfpathlineto{\pgfqpoint{0.833430in}{1.244443in}}%
\pgfpathlineto{\pgfqpoint{0.834388in}{1.244923in}}%
\pgfpathlineto{\pgfqpoint{0.834579in}{1.244732in}}%
\pgfpathlineto{\pgfqpoint{0.837982in}{1.239415in}}%
\pgfpathlineto{\pgfqpoint{0.838792in}{1.240751in}}%
\pgfpathlineto{\pgfqpoint{0.839602in}{1.240448in}}%
\pgfpathlineto{\pgfqpoint{0.842929in}{1.235255in}}%
\pgfpathlineto{\pgfqpoint{0.843429in}{1.236280in}}%
\pgfpathlineto{\pgfqpoint{0.844274in}{1.236611in}}%
\pgfpathlineto{\pgfqpoint{0.844471in}{1.236424in}}%
\pgfpathlineto{\pgfqpoint{0.847886in}{1.231161in}}%
\pgfpathlineto{\pgfqpoint{0.848556in}{1.232407in}}%
\pgfpathlineto{\pgfqpoint{0.849268in}{1.232483in}}%
\pgfpathlineto{\pgfqpoint{0.849417in}{1.232334in}}%
\pgfpathlineto{\pgfqpoint{0.852832in}{1.227098in}}%
\pgfpathlineto{\pgfqpoint{0.853604in}{1.228405in}}%
\pgfpathlineto{\pgfqpoint{0.854465in}{1.228125in}}%
\pgfpathlineto{\pgfqpoint{0.857778in}{1.223075in}}%
\pgfpathlineto{\pgfqpoint{0.858216in}{1.223998in}}%
\pgfpathlineto{\pgfqpoint{0.859073in}{1.224456in}}%
\pgfpathlineto{\pgfqpoint{0.859254in}{1.224314in}}%
\pgfpathlineto{\pgfqpoint{0.861881in}{1.220334in}}%
\pgfpathlineto{\pgfqpoint{0.862720in}{1.219091in}}%
\pgfpathlineto{\pgfqpoint{0.863077in}{1.219904in}}%
\pgfpathlineto{\pgfqpoint{0.863984in}{1.220545in}}%
\pgfpathlineto{\pgfqpoint{0.864128in}{1.220455in}}%
\pgfpathlineto{\pgfqpoint{0.865823in}{1.217969in}}%
\pgfpathlineto{\pgfqpoint{0.867666in}{1.215186in}}%
\pgfpathlineto{\pgfqpoint{0.868048in}{1.215974in}}%
\pgfpathlineto{\pgfqpoint{0.868926in}{1.216576in}}%
\pgfpathlineto{\pgfqpoint{0.869109in}{1.216456in}}%
\pgfpathlineto{\pgfqpoint{0.870992in}{1.213675in}}%
\pgfpathlineto{\pgfqpoint{0.872600in}{1.211273in}}%
\pgfpathlineto{\pgfqpoint{0.872957in}{1.212092in}}%
\pgfpathlineto{\pgfqpoint{0.873929in}{1.212735in}}%
\pgfpathlineto{\pgfqpoint{0.874099in}{1.212604in}}%
\pgfpathlineto{\pgfqpoint{0.876086in}{1.209656in}}%
\pgfpathlineto{\pgfqpoint{0.877555in}{1.207473in}}%
\pgfpathlineto{\pgfqpoint{0.877911in}{1.208284in}}%
\pgfpathlineto{\pgfqpoint{0.878872in}{1.208929in}}%
\pgfpathlineto{\pgfqpoint{0.879032in}{1.208813in}}%
\pgfpathlineto{\pgfqpoint{0.881042in}{1.205857in}}%
\pgfpathlineto{\pgfqpoint{0.882502in}{1.203692in}}%
\pgfpathlineto{\pgfqpoint{0.882881in}{1.204473in}}%
\pgfpathlineto{\pgfqpoint{0.883832in}{1.205072in}}%
\pgfpathlineto{\pgfqpoint{0.883982in}{1.204961in}}%
\pgfpathlineto{\pgfqpoint{0.886106in}{1.201849in}}%
\pgfpathlineto{\pgfqpoint{0.887426in}{1.199888in}}%
\pgfpathlineto{\pgfqpoint{0.887809in}{1.200675in}}%
\pgfpathlineto{\pgfqpoint{0.888760in}{1.201325in}}%
\pgfpathlineto{\pgfqpoint{0.888910in}{1.201224in}}%
\pgfpathlineto{\pgfqpoint{0.890902in}{1.198338in}}%
\pgfpathlineto{\pgfqpoint{0.892381in}{1.196180in}}%
\pgfpathlineto{\pgfqpoint{0.892734in}{1.196995in}}%
\pgfpathlineto{\pgfqpoint{0.893759in}{1.197651in}}%
\pgfpathlineto{\pgfqpoint{0.893947in}{1.197493in}}%
\pgfpathlineto{\pgfqpoint{0.897139in}{1.192801in}}%
\pgfpathlineto{\pgfqpoint{0.897336in}{1.192544in}}%
\pgfpathlineto{\pgfqpoint{0.897703in}{1.193329in}}%
\pgfpathlineto{\pgfqpoint{0.898642in}{1.194002in}}%
\pgfpathlineto{\pgfqpoint{0.898786in}{1.193915in}}%
\pgfpathlineto{\pgfqpoint{0.900690in}{1.191207in}}%
\pgfpathlineto{\pgfqpoint{0.902271in}{1.188904in}}%
\pgfpathlineto{\pgfqpoint{0.902625in}{1.189714in}}%
\pgfpathlineto{\pgfqpoint{0.903668in}{1.190389in}}%
\pgfpathlineto{\pgfqpoint{0.903816in}{1.190272in}}%
\pgfpathlineto{\pgfqpoint{0.906504in}{1.186382in}}%
\pgfpathlineto{\pgfqpoint{0.907231in}{1.185354in}}%
\pgfpathlineto{\pgfqpoint{0.907597in}{1.186146in}}%
\pgfpathlineto{\pgfqpoint{0.908607in}{1.186792in}}%
\pgfpathlineto{\pgfqpoint{0.908773in}{1.186662in}}%
\pgfpathlineto{\pgfqpoint{0.911035in}{1.183408in}}%
\pgfpathlineto{\pgfqpoint{0.912166in}{1.181787in}}%
\pgfpathlineto{\pgfqpoint{0.912538in}{1.182566in}}%
\pgfpathlineto{\pgfqpoint{0.913511in}{1.183252in}}%
\pgfpathlineto{\pgfqpoint{0.913635in}{1.183176in}}%
\pgfpathlineto{\pgfqpoint{0.915226in}{1.180965in}}%
\pgfpathlineto{\pgfqpoint{0.917112in}{1.178266in}}%
\pgfpathlineto{\pgfqpoint{0.917473in}{1.179050in}}%
\pgfpathlineto{\pgfqpoint{0.918515in}{1.179727in}}%
\pgfpathlineto{\pgfqpoint{0.918700in}{1.179576in}}%
\pgfpathlineto{\pgfqpoint{0.921290in}{1.175869in}}%
\pgfpathlineto{\pgfqpoint{0.922056in}{1.174790in}}%
\pgfpathlineto{\pgfqpoint{0.922423in}{1.175584in}}%
\pgfpathlineto{\pgfqpoint{0.923485in}{1.176258in}}%
\pgfpathlineto{\pgfqpoint{0.923633in}{1.176138in}}%
\pgfpathlineto{\pgfqpoint{0.926390in}{1.172220in}}%
\pgfpathlineto{\pgfqpoint{0.926998in}{1.171375in}}%
\pgfpathlineto{\pgfqpoint{0.927351in}{1.172188in}}%
\pgfpathlineto{\pgfqpoint{0.928392in}{1.172912in}}%
\pgfpathlineto{\pgfqpoint{0.928531in}{1.172814in}}%
\pgfpathlineto{\pgfqpoint{0.930660in}{1.169832in}}%
\pgfpathlineto{\pgfqpoint{0.931953in}{1.168016in}}%
\pgfpathlineto{\pgfqpoint{0.932305in}{1.168824in}}%
\pgfpathlineto{\pgfqpoint{0.933303in}{1.169560in}}%
\pgfpathlineto{\pgfqpoint{0.933443in}{1.169476in}}%
\pgfpathlineto{\pgfqpoint{0.934973in}{1.167385in}}%
\pgfpathlineto{\pgfqpoint{0.936902in}{1.164686in}}%
\pgfpathlineto{\pgfqpoint{0.937251in}{1.165476in}}%
\pgfpathlineto{\pgfqpoint{0.938302in}{1.166190in}}%
\pgfpathlineto{\pgfqpoint{0.938457in}{1.166076in}}%
\pgfpathlineto{\pgfqpoint{0.940511in}{1.163217in}}%
\pgfpathlineto{\pgfqpoint{0.941847in}{1.161364in}}%
\pgfpathlineto{\pgfqpoint{0.942208in}{1.162140in}}%
\pgfpathlineto{\pgfqpoint{0.943191in}{1.162872in}}%
\pgfpathlineto{\pgfqpoint{0.943312in}{1.162805in}}%
\pgfpathlineto{\pgfqpoint{0.944888in}{1.160693in}}%
\pgfpathlineto{\pgfqpoint{0.946779in}{1.158058in}}%
\pgfpathlineto{\pgfqpoint{0.947123in}{1.158857in}}%
\pgfpathlineto{\pgfqpoint{0.948151in}{1.159659in}}%
\pgfpathlineto{\pgfqpoint{0.948259in}{1.159601in}}%
\pgfpathlineto{\pgfqpoint{0.949738in}{1.157641in}}%
\pgfpathlineto{\pgfqpoint{0.951739in}{1.154868in}}%
\pgfpathlineto{\pgfqpoint{0.952101in}{1.155644in}}%
\pgfpathlineto{\pgfqpoint{0.953124in}{1.156370in}}%
\pgfpathlineto{\pgfqpoint{0.953228in}{1.156308in}}%
\pgfpathlineto{\pgfqpoint{0.955031in}{1.153888in}}%
\pgfpathlineto{\pgfqpoint{0.956682in}{1.151636in}}%
\pgfpathlineto{\pgfqpoint{0.957038in}{1.152392in}}%
\pgfpathlineto{\pgfqpoint{0.958118in}{1.153122in}}%
\pgfpathlineto{\pgfqpoint{0.958284in}{1.152995in}}%
\pgfpathlineto{\pgfqpoint{0.960533in}{1.149908in}}%
\pgfpathlineto{\pgfqpoint{0.961619in}{1.148435in}}%
\pgfpathlineto{\pgfqpoint{0.961973in}{1.149224in}}%
\pgfpathlineto{\pgfqpoint{0.963026in}{1.150008in}}%
\pgfpathlineto{\pgfqpoint{0.963125in}{1.149951in}}%
\pgfpathlineto{\pgfqpoint{0.964652in}{1.147942in}}%
\pgfpathlineto{\pgfqpoint{0.966572in}{1.145336in}}%
\pgfpathlineto{\pgfqpoint{0.966918in}{1.146120in}}%
\pgfpathlineto{\pgfqpoint{0.968056in}{1.146849in}}%
\pgfpathlineto{\pgfqpoint{0.968207in}{1.146724in}}%
\pgfpathlineto{\pgfqpoint{0.971508in}{1.142219in}}%
\pgfpathlineto{\pgfqpoint{0.972279in}{1.143608in}}%
\pgfpathlineto{\pgfqpoint{0.973088in}{1.143683in}}%
\pgfpathlineto{\pgfqpoint{0.973187in}{1.143590in}}%
\pgfpathlineto{\pgfqpoint{0.976454in}{1.139166in}}%
\pgfpathlineto{\pgfqpoint{0.977091in}{1.140445in}}%
\pgfpathlineto{\pgfqpoint{0.977983in}{1.140713in}}%
\pgfpathlineto{\pgfqpoint{0.978128in}{1.140586in}}%
\pgfpathlineto{\pgfqpoint{0.981408in}{1.136162in}}%
\pgfpathlineto{\pgfqpoint{0.982155in}{1.137565in}}%
\pgfpathlineto{\pgfqpoint{0.983075in}{1.137581in}}%
\pgfpathlineto{\pgfqpoint{0.983125in}{1.137528in}}%
\pgfpathlineto{\pgfqpoint{0.986346in}{1.133184in}}%
\pgfpathlineto{\pgfqpoint{0.986876in}{1.134244in}}%
\pgfpathlineto{\pgfqpoint{0.987789in}{1.134735in}}%
\pgfpathlineto{\pgfqpoint{0.987934in}{1.134641in}}%
\pgfpathlineto{\pgfqpoint{0.990241in}{1.131583in}}%
\pgfpathlineto{\pgfqpoint{0.991276in}{1.130198in}}%
\pgfpathlineto{\pgfqpoint{0.991637in}{1.130973in}}%
\pgfpathlineto{\pgfqpoint{0.992706in}{1.131823in}}%
\pgfpathlineto{\pgfqpoint{0.992794in}{1.131778in}}%
\pgfpathlineto{\pgfqpoint{0.993977in}{1.130309in}}%
\pgfpathlineto{\pgfqpoint{0.996225in}{1.127293in}}%
\pgfpathlineto{\pgfqpoint{0.996590in}{1.128075in}}%
\pgfpathlineto{\pgfqpoint{0.997703in}{1.128880in}}%
\pgfpathlineto{\pgfqpoint{0.997827in}{1.128799in}}%
\pgfpathlineto{\pgfqpoint{0.999756in}{1.126276in}}%
\pgfpathlineto{\pgfqpoint{1.001167in}{1.124404in}}%
\pgfpathlineto{\pgfqpoint{1.001526in}{1.125188in}}%
\pgfpathlineto{\pgfqpoint{1.002721in}{1.125979in}}%
\pgfpathlineto{\pgfqpoint{1.002853in}{1.125871in}}%
\pgfpathlineto{\pgfqpoint{1.005952in}{1.121764in}}%
\pgfpathlineto{\pgfqpoint{1.006128in}{1.121566in}}%
\pgfpathlineto{\pgfqpoint{1.006479in}{1.122335in}}%
\pgfpathlineto{\pgfqpoint{1.007576in}{1.123169in}}%
\pgfpathlineto{\pgfqpoint{1.007643in}{1.123134in}}%
\pgfpathlineto{\pgfqpoint{1.009162in}{1.121235in}}%
\pgfpathlineto{\pgfqpoint{1.011074in}{1.118726in}}%
\pgfpathlineto{\pgfqpoint{1.011415in}{1.119512in}}%
\pgfpathlineto{\pgfqpoint{1.012526in}{1.120389in}}%
\pgfpathlineto{\pgfqpoint{1.012624in}{1.120336in}}%
\pgfpathlineto{\pgfqpoint{1.014305in}{1.118212in}}%
\pgfpathlineto{\pgfqpoint{1.016006in}{1.115969in}}%
\pgfpathlineto{\pgfqpoint{1.016372in}{1.116742in}}%
\pgfpathlineto{\pgfqpoint{1.017507in}{1.117562in}}%
\pgfpathlineto{\pgfqpoint{1.017641in}{1.117473in}}%
\pgfpathlineto{\pgfqpoint{1.019688in}{1.114827in}}%
\pgfpathlineto{\pgfqpoint{1.020972in}{1.113186in}}%
\pgfpathlineto{\pgfqpoint{1.021325in}{1.113941in}}%
\pgfpathlineto{\pgfqpoint{1.022492in}{1.114751in}}%
\pgfpathlineto{\pgfqpoint{1.022628in}{1.114653in}}%
\pgfpathlineto{\pgfqpoint{1.025153in}{1.111378in}}%
\pgfpathlineto{\pgfqpoint{1.025895in}{1.110426in}}%
\pgfpathlineto{\pgfqpoint{1.026235in}{1.111219in}}%
\pgfpathlineto{\pgfqpoint{1.027376in}{1.112126in}}%
\pgfpathlineto{\pgfqpoint{1.027475in}{1.112072in}}%
\pgfpathlineto{\pgfqpoint{1.029155in}{1.109970in}}%
\pgfpathlineto{\pgfqpoint{1.030861in}{1.107772in}}%
\pgfpathlineto{\pgfqpoint{1.031199in}{1.108560in}}%
\pgfpathlineto{\pgfqpoint{1.032369in}{1.109425in}}%
\pgfpathlineto{\pgfqpoint{1.032419in}{1.109396in}}%
\pgfpathlineto{\pgfqpoint{1.034074in}{1.107344in}}%
\pgfpathlineto{\pgfqpoint{1.035801in}{1.105122in}}%
\pgfpathlineto{\pgfqpoint{1.036157in}{1.105891in}}%
\pgfpathlineto{\pgfqpoint{1.037288in}{1.106771in}}%
\pgfpathlineto{\pgfqpoint{1.037348in}{1.106741in}}%
\pgfpathlineto{\pgfqpoint{1.038888in}{1.104864in}}%
\pgfpathlineto{\pgfqpoint{1.040745in}{1.102498in}}%
\pgfpathlineto{\pgfqpoint{1.041074in}{1.103276in}}%
\pgfpathlineto{\pgfqpoint{1.042254in}{1.104194in}}%
\pgfpathlineto{\pgfqpoint{1.042303in}{1.104167in}}%
\pgfpathlineto{\pgfqpoint{1.043922in}{1.102190in}}%
\pgfpathlineto{\pgfqpoint{1.045698in}{1.099933in}}%
\pgfpathlineto{\pgfqpoint{1.046046in}{1.100687in}}%
\pgfpathlineto{\pgfqpoint{1.047244in}{1.101558in}}%
\pgfpathlineto{\pgfqpoint{1.047377in}{1.101466in}}%
\pgfpathlineto{\pgfqpoint{1.049715in}{1.098503in}}%
\pgfpathlineto{\pgfqpoint{1.050621in}{1.097357in}}%
\pgfpathlineto{\pgfqpoint{1.050962in}{1.098122in}}%
\pgfpathlineto{\pgfqpoint{1.052174in}{1.099044in}}%
\pgfpathlineto{\pgfqpoint{1.052273in}{1.098982in}}%
\pgfpathlineto{\pgfqpoint{1.054277in}{1.096483in}}%
\pgfpathlineto{\pgfqpoint{1.055570in}{1.094853in}}%
\pgfpathlineto{\pgfqpoint{1.055896in}{1.095607in}}%
\pgfpathlineto{\pgfqpoint{1.057122in}{1.096596in}}%
\pgfpathlineto{\pgfqpoint{1.057154in}{1.096579in}}%
\pgfpathlineto{\pgfqpoint{1.058670in}{1.094763in}}%
\pgfpathlineto{\pgfqpoint{1.060531in}{1.092431in}}%
\pgfpathlineto{\pgfqpoint{1.060860in}{1.093210in}}%
\pgfpathlineto{\pgfqpoint{1.062073in}{1.094134in}}%
\pgfpathlineto{\pgfqpoint{1.062145in}{1.094091in}}%
\pgfpathlineto{\pgfqpoint{1.063910in}{1.091930in}}%
\pgfpathlineto{\pgfqpoint{1.065472in}{1.089975in}}%
\pgfpathlineto{\pgfqpoint{1.065811in}{1.090741in}}%
\pgfpathlineto{\pgfqpoint{1.066992in}{1.091685in}}%
\pgfpathlineto{\pgfqpoint{1.068266in}{1.090231in}}%
\pgfpathlineto{\pgfqpoint{1.070426in}{1.087544in}}%
\pgfpathlineto{\pgfqpoint{1.070768in}{1.088283in}}%
\pgfpathlineto{\pgfqpoint{1.072044in}{1.089137in}}%
\pgfpathlineto{\pgfqpoint{1.072154in}{1.089051in}}%
\pgfpathlineto{\pgfqpoint{1.075351in}{1.085078in}}%
\pgfpathlineto{\pgfqpoint{1.076006in}{1.086439in}}%
\pgfpathlineto{\pgfqpoint{1.077003in}{1.086818in}}%
\pgfpathlineto{\pgfqpoint{1.077106in}{1.086739in}}%
\pgfpathlineto{\pgfqpoint{1.080140in}{1.082959in}}%
\pgfpathlineto{\pgfqpoint{1.080316in}{1.082783in}}%
\pgfpathlineto{\pgfqpoint{1.080650in}{1.083545in}}%
\pgfpathlineto{\pgfqpoint{1.081854in}{1.084488in}}%
\pgfpathlineto{\pgfqpoint{1.081908in}{1.084459in}}%
\pgfpathlineto{\pgfqpoint{1.083321in}{1.082795in}}%
\pgfpathlineto{\pgfqpoint{1.085240in}{1.080410in}}%
\pgfpathlineto{\pgfqpoint{1.085587in}{1.081187in}}%
\pgfpathlineto{\pgfqpoint{1.086856in}{1.082140in}}%
\pgfpathlineto{\pgfqpoint{1.086970in}{1.082062in}}%
\pgfpathlineto{\pgfqpoint{1.089466in}{1.078994in}}%
\pgfpathlineto{\pgfqpoint{1.090197in}{1.078114in}}%
\pgfpathlineto{\pgfqpoint{1.090527in}{1.078897in}}%
\pgfpathlineto{\pgfqpoint{1.091773in}{1.079877in}}%
\pgfpathlineto{\pgfqpoint{1.093254in}{1.078167in}}%
\pgfpathlineto{\pgfqpoint{1.095152in}{1.075840in}}%
\pgfpathlineto{\pgfqpoint{1.095504in}{1.076590in}}%
\pgfpathlineto{\pgfqpoint{1.096680in}{1.077545in}}%
\pgfpathlineto{\pgfqpoint{1.096689in}{1.077541in}}%
\pgfpathlineto{\pgfqpoint{1.098010in}{1.076068in}}%
\pgfpathlineto{\pgfqpoint{1.100075in}{1.073540in}}%
\pgfpathlineto{\pgfqpoint{1.100407in}{1.074320in}}%
\pgfpathlineto{\pgfqpoint{1.101700in}{1.075334in}}%
\pgfpathlineto{\pgfqpoint{1.101799in}{1.075271in}}%
\pgfpathlineto{\pgfqpoint{1.104251in}{1.072304in}}%
\pgfpathlineto{\pgfqpoint{1.105034in}{1.071369in}}%
\pgfpathlineto{\pgfqpoint{1.105368in}{1.072144in}}%
\pgfpathlineto{\pgfqpoint{1.106679in}{1.073150in}}%
\pgfpathlineto{\pgfqpoint{1.106729in}{1.073120in}}%
\pgfpathlineto{\pgfqpoint{1.108786in}{1.070663in}}%
\pgfpathlineto{\pgfqpoint{1.109976in}{1.069228in}}%
\pgfpathlineto{\pgfqpoint{1.110305in}{1.070011in}}%
\pgfpathlineto{\pgfqpoint{1.111531in}{1.071050in}}%
\pgfpathlineto{\pgfqpoint{1.112795in}{1.069668in}}%
\pgfpathlineto{\pgfqpoint{1.114927in}{1.067091in}}%
\pgfpathlineto{\pgfqpoint{1.115255in}{1.067861in}}%
\pgfpathlineto{\pgfqpoint{1.116569in}{1.068881in}}%
\pgfpathlineto{\pgfqpoint{1.116667in}{1.068816in}}%
\pgfpathlineto{\pgfqpoint{1.119203in}{1.065774in}}%
\pgfpathlineto{\pgfqpoint{1.119871in}{1.064997in}}%
\pgfpathlineto{\pgfqpoint{1.120196in}{1.065771in}}%
\pgfpathlineto{\pgfqpoint{1.121534in}{1.066823in}}%
\pgfpathlineto{\pgfqpoint{1.121584in}{1.066794in}}%
\pgfpathlineto{\pgfqpoint{1.123618in}{1.064391in}}%
\pgfpathlineto{\pgfqpoint{1.124822in}{1.062968in}}%
\pgfpathlineto{\pgfqpoint{1.125155in}{1.063736in}}%
\pgfpathlineto{\pgfqpoint{1.126381in}{1.064784in}}%
\pgfpathlineto{\pgfqpoint{1.126408in}{1.064773in}}%
\pgfpathlineto{\pgfqpoint{1.127681in}{1.063378in}}%
\pgfpathlineto{\pgfqpoint{1.129748in}{1.060891in}}%
\pgfpathlineto{\pgfqpoint{1.130082in}{1.061659in}}%
\pgfpathlineto{\pgfqpoint{1.131337in}{1.062780in}}%
\pgfpathlineto{\pgfqpoint{1.132514in}{1.061533in}}%
\pgfpathlineto{\pgfqpoint{1.134698in}{1.058917in}}%
\pgfpathlineto{\pgfqpoint{1.135036in}{1.059690in}}%
\pgfpathlineto{\pgfqpoint{1.136360in}{1.060755in}}%
\pgfpathlineto{\pgfqpoint{1.136419in}{1.060721in}}%
\pgfpathlineto{\pgfqpoint{1.138334in}{1.058490in}}%
\pgfpathlineto{\pgfqpoint{1.139640in}{1.056941in}}%
\pgfpathlineto{\pgfqpoint{1.139975in}{1.057719in}}%
\pgfpathlineto{\pgfqpoint{1.141267in}{1.058835in}}%
\pgfpathlineto{\pgfqpoint{1.142688in}{1.057279in}}%
\pgfpathlineto{\pgfqpoint{1.144587in}{1.055017in}}%
\pgfpathlineto{\pgfqpoint{1.144920in}{1.055798in}}%
\pgfpathlineto{\pgfqpoint{1.146259in}{1.056913in}}%
\pgfpathlineto{\pgfqpoint{1.146314in}{1.056882in}}%
\pgfpathlineto{\pgfqpoint{1.147872in}{1.055103in}}%
\pgfpathlineto{\pgfqpoint{1.149531in}{1.053138in}}%
\pgfpathlineto{\pgfqpoint{1.149856in}{1.053923in}}%
\pgfpathlineto{\pgfqpoint{1.151215in}{1.055070in}}%
\pgfpathlineto{\pgfqpoint{1.151303in}{1.055019in}}%
\pgfpathlineto{\pgfqpoint{1.153508in}{1.052452in}}%
\pgfpathlineto{\pgfqpoint{1.154480in}{1.051315in}}%
\pgfpathlineto{\pgfqpoint{1.154816in}{1.052093in}}%
\pgfpathlineto{\pgfqpoint{1.156160in}{1.053224in}}%
\pgfpathlineto{\pgfqpoint{1.157611in}{1.051624in}}%
\pgfpathlineto{\pgfqpoint{1.159428in}{1.049491in}}%
\pgfpathlineto{\pgfqpoint{1.159743in}{1.050279in}}%
\pgfpathlineto{\pgfqpoint{1.161159in}{1.051470in}}%
\pgfpathlineto{\pgfqpoint{1.161209in}{1.051440in}}%
\pgfpathlineto{\pgfqpoint{1.163285in}{1.049048in}}%
\pgfpathlineto{\pgfqpoint{1.164367in}{1.047780in}}%
\pgfpathlineto{\pgfqpoint{1.164687in}{1.048560in}}%
\pgfpathlineto{\pgfqpoint{1.166124in}{1.049757in}}%
\pgfpathlineto{\pgfqpoint{1.168002in}{1.047631in}}%
\pgfpathlineto{\pgfqpoint{1.169320in}{1.046105in}}%
\pgfpathlineto{\pgfqpoint{1.169608in}{1.046844in}}%
\pgfpathlineto{\pgfqpoint{1.171042in}{1.048163in}}%
\pgfpathlineto{\pgfqpoint{1.172503in}{1.046563in}}%
\pgfpathlineto{\pgfqpoint{1.174273in}{1.044509in}}%
\pgfpathlineto{\pgfqpoint{1.174581in}{1.045292in}}%
\pgfpathlineto{\pgfqpoint{1.175907in}{1.046594in}}%
\pgfpathlineto{\pgfqpoint{1.177130in}{1.045346in}}%
\pgfpathlineto{\pgfqpoint{1.179216in}{1.042915in}}%
\pgfpathlineto{\pgfqpoint{1.179531in}{1.043703in}}%
\pgfpathlineto{\pgfqpoint{1.180958in}{1.044974in}}%
\pgfpathlineto{\pgfqpoint{1.182577in}{1.043190in}}%
\pgfpathlineto{\pgfqpoint{1.184161in}{1.041359in}}%
\pgfpathlineto{\pgfqpoint{1.184470in}{1.042148in}}%
\pgfpathlineto{\pgfqpoint{1.185888in}{1.043479in}}%
\pgfpathlineto{\pgfqpoint{1.187385in}{1.041863in}}%
\pgfpathlineto{\pgfqpoint{1.189097in}{1.039877in}}%
\pgfpathlineto{\pgfqpoint{1.189417in}{1.040667in}}%
\pgfpathlineto{\pgfqpoint{1.190762in}{1.042062in}}%
\pgfpathlineto{\pgfqpoint{1.191920in}{1.040918in}}%
\pgfpathlineto{\pgfqpoint{1.194036in}{1.038468in}}%
\pgfpathlineto{\pgfqpoint{1.194344in}{1.039261in}}%
\pgfpathlineto{\pgfqpoint{1.195771in}{1.040719in}}%
\pgfpathlineto{\pgfqpoint{1.197042in}{1.039408in}}%
\pgfpathlineto{\pgfqpoint{1.198984in}{1.037164in}}%
\pgfpathlineto{\pgfqpoint{1.199290in}{1.037945in}}%
\pgfpathlineto{\pgfqpoint{1.200679in}{1.039472in}}%
\pgfpathlineto{\pgfqpoint{1.201878in}{1.038297in}}%
\pgfpathlineto{\pgfqpoint{1.203929in}{1.035935in}}%
\pgfpathlineto{\pgfqpoint{1.204204in}{1.036649in}}%
\pgfpathlineto{\pgfqpoint{1.205703in}{1.038306in}}%
\pgfpathlineto{\pgfqpoint{1.206947in}{1.037037in}}%
\pgfpathlineto{\pgfqpoint{1.208879in}{1.034829in}}%
\pgfpathlineto{\pgfqpoint{1.209168in}{1.035624in}}%
\pgfpathlineto{\pgfqpoint{1.210589in}{1.037331in}}%
\pgfpathlineto{\pgfqpoint{1.211679in}{1.036324in}}%
\pgfpathlineto{\pgfqpoint{1.213823in}{1.033855in}}%
\pgfpathlineto{\pgfqpoint{1.214130in}{1.034649in}}%
\pgfpathlineto{\pgfqpoint{1.215538in}{1.036380in}}%
\pgfpathlineto{\pgfqpoint{1.216506in}{1.035538in}}%
\pgfpathlineto{\pgfqpoint{1.218769in}{1.032955in}}%
\pgfpathlineto{\pgfqpoint{1.219055in}{1.033763in}}%
\pgfpathlineto{\pgfqpoint{1.220626in}{1.035579in}}%
\pgfpathlineto{\pgfqpoint{1.221911in}{1.034264in}}%
\pgfpathlineto{\pgfqpoint{1.223709in}{1.032210in}}%
\pgfpathlineto{\pgfqpoint{1.223973in}{1.032967in}}%
\pgfpathlineto{\pgfqpoint{1.225506in}{1.034923in}}%
\pgfpathlineto{\pgfqpoint{1.226640in}{1.033857in}}%
\pgfpathlineto{\pgfqpoint{1.228659in}{1.031554in}}%
\pgfpathlineto{\pgfqpoint{1.228937in}{1.032331in}}%
\pgfpathlineto{\pgfqpoint{1.230465in}{1.034238in}}%
\pgfpathlineto{\pgfqpoint{1.231589in}{1.033175in}}%
\pgfpathlineto{\pgfqpoint{1.233609in}{1.030882in}}%
\pgfpathlineto{\pgfqpoint{1.233885in}{1.031690in}}%
\pgfpathlineto{\pgfqpoint{1.235451in}{1.033585in}}%
\pgfpathlineto{\pgfqpoint{1.236633in}{1.032422in}}%
\pgfpathlineto{\pgfqpoint{1.238555in}{1.030232in}}%
\pgfpathlineto{\pgfqpoint{1.238838in}{1.031043in}}%
\pgfpathlineto{\pgfqpoint{1.240417in}{1.032923in}}%
\pgfpathlineto{\pgfqpoint{1.241710in}{1.031618in}}%
\pgfpathlineto{\pgfqpoint{1.243492in}{1.029586in}}%
\pgfpathlineto{\pgfqpoint{1.243778in}{1.030391in}}%
\pgfpathlineto{\pgfqpoint{1.245244in}{1.032323in}}%
\pgfpathlineto{\pgfqpoint{1.246283in}{1.031431in}}%
\pgfpathlineto{\pgfqpoint{1.248435in}{1.028978in}}%
\pgfpathlineto{\pgfqpoint{1.248710in}{1.029768in}}%
\pgfpathlineto{\pgfqpoint{1.250247in}{1.031705in}}%
\pgfpathlineto{\pgfqpoint{1.251365in}{1.030660in}}%
\pgfpathlineto{\pgfqpoint{1.253381in}{1.028365in}}%
\pgfpathlineto{\pgfqpoint{1.253665in}{1.029174in}}%
\pgfpathlineto{\pgfqpoint{1.255196in}{1.031084in}}%
\pgfpathlineto{\pgfqpoint{1.256403in}{1.029929in}}%
\pgfpathlineto{\pgfqpoint{1.258328in}{1.027743in}}%
\pgfpathlineto{\pgfqpoint{1.258617in}{1.028556in}}%
\pgfpathlineto{\pgfqpoint{1.260160in}{1.030477in}}%
\pgfpathlineto{\pgfqpoint{1.261274in}{1.029429in}}%
\pgfpathlineto{\pgfqpoint{1.263273in}{1.027162in}}%
\pgfpathlineto{\pgfqpoint{1.263542in}{1.027924in}}%
\pgfpathlineto{\pgfqpoint{1.265113in}{1.029887in}}%
\pgfpathlineto{\pgfqpoint{1.266236in}{1.028820in}}%
\pgfpathlineto{\pgfqpoint{1.268220in}{1.026575in}}%
\pgfpathlineto{\pgfqpoint{1.268494in}{1.027337in}}%
\pgfpathlineto{\pgfqpoint{1.270067in}{1.029297in}}%
\pgfpathlineto{\pgfqpoint{1.271207in}{1.028208in}}%
\pgfpathlineto{\pgfqpoint{1.273165in}{1.026004in}}%
\pgfpathlineto{\pgfqpoint{1.273394in}{1.026684in}}%
\pgfpathlineto{\pgfqpoint{1.274944in}{1.028814in}}%
\pgfpathlineto{\pgfqpoint{1.275978in}{1.027924in}}%
\pgfpathlineto{\pgfqpoint{1.278112in}{1.025512in}}%
\pgfpathlineto{\pgfqpoint{1.278394in}{1.026333in}}%
\pgfpathlineto{\pgfqpoint{1.279958in}{1.028286in}}%
\pgfpathlineto{\pgfqpoint{1.281096in}{1.027207in}}%
\pgfpathlineto{\pgfqpoint{1.283059in}{1.024990in}}%
\pgfpathlineto{\pgfqpoint{1.283347in}{1.025804in}}%
\pgfpathlineto{\pgfqpoint{1.284837in}{1.027753in}}%
\pgfpathlineto{\pgfqpoint{1.285813in}{1.026925in}}%
\pgfpathlineto{\pgfqpoint{1.287999in}{1.024456in}}%
\pgfpathlineto{\pgfqpoint{1.288276in}{1.025259in}}%
\pgfpathlineto{\pgfqpoint{1.289809in}{1.027271in}}%
\pgfpathlineto{\pgfqpoint{1.290927in}{1.026263in}}%
\pgfpathlineto{\pgfqpoint{1.292950in}{1.023975in}}%
\pgfpathlineto{\pgfqpoint{1.293243in}{1.024745in}}%
\pgfpathlineto{\pgfqpoint{1.294854in}{1.026650in}}%
\pgfpathlineto{\pgfqpoint{1.295979in}{1.025542in}}%
\pgfpathlineto{\pgfqpoint{1.297883in}{1.023394in}}%
\pgfpathlineto{\pgfqpoint{1.298172in}{1.024206in}}%
\pgfpathlineto{\pgfqpoint{1.299701in}{1.026175in}}%
\pgfpathlineto{\pgfqpoint{1.300802in}{1.025181in}}%
\pgfpathlineto{\pgfqpoint{1.302843in}{1.022898in}}%
\pgfpathlineto{\pgfqpoint{1.303122in}{1.023709in}}%
\pgfpathlineto{\pgfqpoint{1.304678in}{1.025674in}}%
\pgfpathlineto{\pgfqpoint{1.305858in}{1.024565in}}%
\pgfpathlineto{\pgfqpoint{1.307787in}{1.022397in}}%
\pgfpathlineto{\pgfqpoint{1.308068in}{1.023207in}}%
\pgfpathlineto{\pgfqpoint{1.309738in}{1.025157in}}%
\pgfpathlineto{\pgfqpoint{1.311127in}{1.023723in}}%
\pgfpathlineto{\pgfqpoint{1.312729in}{1.021930in}}%
\pgfpathlineto{\pgfqpoint{1.313002in}{1.022695in}}%
\pgfpathlineto{\pgfqpoint{1.314576in}{1.024724in}}%
\pgfpathlineto{\pgfqpoint{1.315674in}{1.023720in}}%
\pgfpathlineto{\pgfqpoint{1.317668in}{1.021483in}}%
\pgfpathlineto{\pgfqpoint{1.317937in}{1.022262in}}%
\pgfpathlineto{\pgfqpoint{1.319502in}{1.024290in}}%
\pgfpathlineto{\pgfqpoint{1.320650in}{1.023232in}}%
\pgfpathlineto{\pgfqpoint{1.322622in}{1.021026in}}%
\pgfpathlineto{\pgfqpoint{1.322895in}{1.021821in}}%
\pgfpathlineto{\pgfqpoint{1.324466in}{1.023855in}}%
\pgfpathlineto{\pgfqpoint{1.325609in}{1.022800in}}%
\pgfpathlineto{\pgfqpoint{1.327573in}{1.020604in}}%
\pgfpathlineto{\pgfqpoint{1.327853in}{1.021422in}}%
\pgfpathlineto{\pgfqpoint{1.329303in}{1.023426in}}%
\pgfpathlineto{\pgfqpoint{1.330312in}{1.022638in}}%
\pgfpathlineto{\pgfqpoint{1.332519in}{1.020168in}}%
\pgfpathlineto{\pgfqpoint{1.332798in}{1.020984in}}%
\pgfpathlineto{\pgfqpoint{1.334260in}{1.022994in}}%
\pgfpathlineto{\pgfqpoint{1.335249in}{1.022218in}}%
\pgfpathlineto{\pgfqpoint{1.337464in}{1.019742in}}%
\pgfpathlineto{\pgfqpoint{1.337742in}{1.020552in}}%
\pgfpathlineto{\pgfqpoint{1.339244in}{1.022556in}}%
\pgfpathlineto{\pgfqpoint{1.340319in}{1.021639in}}%
\pgfpathlineto{\pgfqpoint{1.342410in}{1.019303in}}%
\pgfpathlineto{\pgfqpoint{1.342687in}{1.020114in}}%
\pgfpathlineto{\pgfqpoint{1.344192in}{1.022135in}}%
\pgfpathlineto{\pgfqpoint{1.345274in}{1.021216in}}%
\pgfpathlineto{\pgfqpoint{1.347361in}{1.018903in}}%
\pgfpathlineto{\pgfqpoint{1.347635in}{1.019708in}}%
\pgfpathlineto{\pgfqpoint{1.349084in}{1.021678in}}%
\pgfpathlineto{\pgfqpoint{1.350071in}{1.020907in}}%
\pgfpathlineto{\pgfqpoint{1.352295in}{1.018413in}}%
\pgfpathlineto{\pgfqpoint{1.352586in}{1.019190in}}%
\pgfpathlineto{\pgfqpoint{1.354215in}{1.021163in}}%
\pgfpathlineto{\pgfqpoint{1.355431in}{1.019972in}}%
\pgfpathlineto{\pgfqpoint{1.357235in}{1.017967in}}%
\pgfpathlineto{\pgfqpoint{1.357496in}{1.018735in}}%
\pgfpathlineto{\pgfqpoint{1.359057in}{1.020812in}}%
\pgfpathlineto{\pgfqpoint{1.360139in}{1.019865in}}%
\pgfpathlineto{\pgfqpoint{1.362201in}{1.017593in}}%
\pgfpathlineto{\pgfqpoint{1.362475in}{1.018392in}}%
\pgfpathlineto{\pgfqpoint{1.363923in}{1.020354in}}%
\pgfpathlineto{\pgfqpoint{1.364871in}{1.019627in}}%
\pgfpathlineto{\pgfqpoint{1.367128in}{1.017110in}}%
\pgfpathlineto{\pgfqpoint{1.367396in}{1.017865in}}%
\pgfpathlineto{\pgfqpoint{1.368977in}{1.019953in}}%
\pgfpathlineto{\pgfqpoint{1.370070in}{1.018980in}}%
\pgfpathlineto{\pgfqpoint{1.372068in}{1.016756in}}%
\pgfpathlineto{\pgfqpoint{1.372331in}{1.017524in}}%
\pgfpathlineto{\pgfqpoint{1.373887in}{1.019606in}}%
\pgfpathlineto{\pgfqpoint{1.374957in}{1.018679in}}%
\pgfpathlineto{\pgfqpoint{1.377023in}{1.016375in}}%
\pgfpathlineto{\pgfqpoint{1.377318in}{1.017181in}}%
\pgfpathlineto{\pgfqpoint{1.378938in}{1.019163in}}%
\pgfpathlineto{\pgfqpoint{1.380100in}{1.018049in}}%
\pgfpathlineto{\pgfqpoint{1.381960in}{1.015994in}}%
\pgfpathlineto{\pgfqpoint{1.382236in}{1.016801in}}%
\pgfpathlineto{\pgfqpoint{1.383816in}{1.018849in}}%
\pgfpathlineto{\pgfqpoint{1.384948in}{1.017819in}}%
\pgfpathlineto{\pgfqpoint{1.386913in}{1.015643in}}%
\pgfpathlineto{\pgfqpoint{1.387180in}{1.016416in}}%
\pgfpathlineto{\pgfqpoint{1.388742in}{1.018504in}}%
\pgfpathlineto{\pgfqpoint{1.389849in}{1.017536in}}%
\pgfpathlineto{\pgfqpoint{1.391856in}{1.015312in}}%
\pgfpathlineto{\pgfqpoint{1.392120in}{1.016069in}}%
\pgfpathlineto{\pgfqpoint{1.393716in}{1.018135in}}%
\pgfpathlineto{\pgfqpoint{1.394888in}{1.017057in}}%
\pgfpathlineto{\pgfqpoint{1.396802in}{1.014943in}}%
\pgfpathlineto{\pgfqpoint{1.397069in}{1.015721in}}%
\pgfpathlineto{\pgfqpoint{1.398675in}{1.017790in}}%
\pgfpathlineto{\pgfqpoint{1.399821in}{1.016736in}}%
\pgfpathlineto{\pgfqpoint{1.401750in}{1.014598in}}%
\pgfpathlineto{\pgfqpoint{1.402034in}{1.015374in}}%
\pgfpathlineto{\pgfqpoint{1.403664in}{1.017355in}}%
\pgfpathlineto{\pgfqpoint{1.404920in}{1.016131in}}%
\pgfpathlineto{\pgfqpoint{1.406697in}{1.014183in}}%
\pgfpathlineto{\pgfqpoint{1.406972in}{1.014988in}}%
\pgfpathlineto{\pgfqpoint{1.408401in}{1.017042in}}%
\pgfpathlineto{\pgfqpoint{1.409391in}{1.016344in}}%
\pgfpathlineto{\pgfqpoint{1.411648in}{1.013854in}}%
\pgfpathlineto{\pgfqpoint{1.411923in}{1.014665in}}%
\pgfpathlineto{\pgfqpoint{1.413410in}{1.016679in}}%
\pgfpathlineto{\pgfqpoint{1.414448in}{1.015842in}}%
\pgfpathlineto{\pgfqpoint{1.416583in}{1.013477in}}%
\pgfpathlineto{\pgfqpoint{1.416844in}{1.014237in}}%
\pgfpathlineto{\pgfqpoint{1.418455in}{1.016334in}}%
\pgfpathlineto{\pgfqpoint{1.419566in}{1.015328in}}%
\pgfpathlineto{\pgfqpoint{1.421531in}{1.013161in}}%
\pgfpathlineto{\pgfqpoint{1.421795in}{1.013923in}}%
\pgfpathlineto{\pgfqpoint{1.423366in}{1.015992in}}%
\pgfpathlineto{\pgfqpoint{1.424484in}{1.015006in}}%
\pgfpathlineto{\pgfqpoint{1.426489in}{1.012816in}}%
\pgfpathlineto{\pgfqpoint{1.426759in}{1.013613in}}%
\pgfpathlineto{\pgfqpoint{1.428362in}{1.015628in}}%
\pgfpathlineto{\pgfqpoint{1.429560in}{1.014505in}}%
\pgfpathlineto{\pgfqpoint{1.431421in}{1.012455in}}%
\pgfpathlineto{\pgfqpoint{1.431697in}{1.013262in}}%
\pgfpathlineto{\pgfqpoint{1.433182in}{1.015317in}}%
\pgfpathlineto{\pgfqpoint{1.434171in}{1.014557in}}%
\pgfpathlineto{\pgfqpoint{1.436370in}{1.012139in}}%
\pgfpathlineto{\pgfqpoint{1.436650in}{1.012946in}}%
\pgfpathlineto{\pgfqpoint{1.438177in}{1.014992in}}%
\pgfpathlineto{\pgfqpoint{1.439236in}{1.014106in}}%
\pgfpathlineto{\pgfqpoint{1.441313in}{1.011824in}}%
\pgfpathlineto{\pgfqpoint{1.441590in}{1.012634in}}%
\pgfpathlineto{\pgfqpoint{1.443127in}{1.014702in}}%
\pgfpathlineto{\pgfqpoint{1.444201in}{1.013798in}}%
\pgfpathlineto{\pgfqpoint{1.446258in}{1.011537in}}%
\pgfpathlineto{\pgfqpoint{1.446532in}{1.012343in}}%
\pgfpathlineto{\pgfqpoint{1.448112in}{1.014409in}}%
\pgfpathlineto{\pgfqpoint{1.449228in}{1.013419in}}%
\pgfpathlineto{\pgfqpoint{1.451207in}{1.011254in}}%
\pgfpathlineto{\pgfqpoint{1.451480in}{1.012057in}}%
\pgfpathlineto{\pgfqpoint{1.452944in}{1.014113in}}%
\pgfpathlineto{\pgfqpoint{1.453933in}{1.013384in}}%
\pgfpathlineto{\pgfqpoint{1.456142in}{1.010945in}}%
\pgfpathlineto{\pgfqpoint{1.456423in}{1.011751in}}%
\pgfpathlineto{\pgfqpoint{1.458053in}{1.013795in}}%
\pgfpathlineto{\pgfqpoint{1.459310in}{1.012597in}}%
\pgfpathlineto{\pgfqpoint{1.461109in}{1.010646in}}%
\pgfpathlineto{\pgfqpoint{1.461382in}{1.011436in}}%
\pgfpathlineto{\pgfqpoint{1.463006in}{1.013475in}}%
\pgfpathlineto{\pgfqpoint{1.464153in}{1.012409in}}%
\pgfpathlineto{\pgfqpoint{1.466038in}{1.010346in}}%
\pgfpathlineto{\pgfqpoint{1.466308in}{1.011137in}}%
\pgfpathlineto{\pgfqpoint{1.467879in}{1.013213in}}%
\pgfpathlineto{\pgfqpoint{1.469019in}{1.012211in}}%
\pgfpathlineto{\pgfqpoint{1.470981in}{1.010057in}}%
\pgfpathlineto{\pgfqpoint{1.471260in}{1.010868in}}%
\pgfpathlineto{\pgfqpoint{1.472784in}{1.012946in}}%
\pgfpathlineto{\pgfqpoint{1.473861in}{1.012064in}}%
\pgfpathlineto{\pgfqpoint{1.475932in}{1.009789in}}%
\pgfpathlineto{\pgfqpoint{1.476214in}{1.010601in}}%
\pgfpathlineto{\pgfqpoint{1.477941in}{1.012579in}}%
\pgfpathlineto{\pgfqpoint{1.479350in}{1.011136in}}%
\pgfpathlineto{\pgfqpoint{1.480875in}{1.009461in}}%
\pgfpathlineto{\pgfqpoint{1.481162in}{1.010273in}}%
\pgfpathlineto{\pgfqpoint{1.482745in}{1.012338in}}%
\pgfpathlineto{\pgfqpoint{1.483814in}{1.011401in}}%
\pgfpathlineto{\pgfqpoint{1.485822in}{1.009209in}}%
\pgfpathlineto{\pgfqpoint{1.486096in}{1.010021in}}%
\pgfpathlineto{\pgfqpoint{1.487711in}{1.012091in}}%
\pgfpathlineto{\pgfqpoint{1.488885in}{1.011013in}}%
\pgfpathlineto{\pgfqpoint{1.490768in}{1.008946in}}%
\pgfpathlineto{\pgfqpoint{1.491051in}{1.009751in}}%
\pgfpathlineto{\pgfqpoint{1.492653in}{1.011793in}}%
\pgfpathlineto{\pgfqpoint{1.493818in}{1.010728in}}%
\pgfpathlineto{\pgfqpoint{1.495714in}{1.008657in}}%
\pgfpathlineto{\pgfqpoint{1.495992in}{1.009469in}}%
\pgfpathlineto{\pgfqpoint{1.497694in}{1.011515in}}%
\pgfpathlineto{\pgfqpoint{1.498970in}{1.010253in}}%
\pgfpathlineto{\pgfqpoint{1.500658in}{1.008412in}}%
\pgfpathlineto{\pgfqpoint{1.500935in}{1.009226in}}%
\pgfpathlineto{\pgfqpoint{1.502525in}{1.011310in}}%
\pgfpathlineto{\pgfqpoint{1.503637in}{1.010329in}}%
\pgfpathlineto{\pgfqpoint{1.505603in}{1.008170in}}%
\pgfpathlineto{\pgfqpoint{1.505891in}{1.008963in}}%
\pgfpathlineto{\pgfqpoint{1.507404in}{1.011011in}}%
\pgfpathlineto{\pgfqpoint{1.508460in}{1.010154in}}%
\pgfpathlineto{\pgfqpoint{1.510542in}{1.007867in}}%
\pgfpathlineto{\pgfqpoint{1.510826in}{1.008687in}}%
\pgfpathlineto{\pgfqpoint{1.512423in}{1.010777in}}%
\pgfpathlineto{\pgfqpoint{1.513540in}{1.009789in}}%
\pgfpathlineto{\pgfqpoint{1.515495in}{1.007655in}}%
\pgfpathlineto{\pgfqpoint{1.515770in}{1.008450in}}%
\pgfpathlineto{\pgfqpoint{1.517330in}{1.010521in}}%
\pgfpathlineto{\pgfqpoint{1.518412in}{1.009600in}}%
\pgfpathlineto{\pgfqpoint{1.520440in}{1.007390in}}%
\pgfpathlineto{\pgfqpoint{1.520700in}{1.008154in}}%
\pgfpathlineto{\pgfqpoint{1.522321in}{1.010288in}}%
\pgfpathlineto{\pgfqpoint{1.523430in}{1.009303in}}%
\pgfpathlineto{\pgfqpoint{1.525385in}{1.007173in}}%
\pgfpathlineto{\pgfqpoint{1.525659in}{1.007979in}}%
\pgfpathlineto{\pgfqpoint{1.527118in}{1.010072in}}%
\pgfpathlineto{\pgfqpoint{1.528106in}{1.009378in}}%
\pgfpathlineto{\pgfqpoint{1.530331in}{1.006950in}}%
\pgfpathlineto{\pgfqpoint{1.530595in}{1.007710in}}%
\pgfpathlineto{\pgfqpoint{1.532186in}{1.009816in}}%
\pgfpathlineto{\pgfqpoint{1.533327in}{1.008813in}}%
\pgfpathlineto{\pgfqpoint{1.535274in}{1.006693in}}%
\pgfpathlineto{\pgfqpoint{1.535537in}{1.007454in}}%
\pgfpathlineto{\pgfqpoint{1.537124in}{1.009600in}}%
\pgfpathlineto{\pgfqpoint{1.538203in}{1.008686in}}%
\pgfpathlineto{\pgfqpoint{1.540217in}{1.006493in}}%
\pgfpathlineto{\pgfqpoint{1.540491in}{1.007291in}}%
\pgfpathlineto{\pgfqpoint{1.542013in}{1.009383in}}%
\pgfpathlineto{\pgfqpoint{1.543078in}{1.008531in}}%
\pgfpathlineto{\pgfqpoint{1.545166in}{1.006246in}}%
\pgfpathlineto{\pgfqpoint{1.545448in}{1.007039in}}%
\pgfpathlineto{\pgfqpoint{1.547016in}{1.009102in}}%
\pgfpathlineto{\pgfqpoint{1.548094in}{1.008178in}}%
\pgfpathlineto{\pgfqpoint{1.550112in}{1.005981in}}%
\pgfpathlineto{\pgfqpoint{1.550391in}{1.006792in}}%
\pgfpathlineto{\pgfqpoint{1.552092in}{1.008858in}}%
\pgfpathlineto{\pgfqpoint{1.553345in}{1.007635in}}%
\pgfpathlineto{\pgfqpoint{1.555058in}{1.005775in}}%
\pgfpathlineto{\pgfqpoint{1.555335in}{1.006588in}}%
\pgfpathlineto{\pgfqpoint{1.556983in}{1.008678in}}%
\pgfpathlineto{\pgfqpoint{1.558178in}{1.007563in}}%
\pgfpathlineto{\pgfqpoint{1.560000in}{1.005568in}}%
\pgfpathlineto{\pgfqpoint{1.560294in}{1.006369in}}%
\pgfpathlineto{\pgfqpoint{1.562041in}{1.008377in}}%
\pgfpathlineto{\pgfqpoint{1.562066in}{1.008366in}}%
\pgfpathlineto{\pgfqpoint{1.563590in}{1.006788in}}%
\pgfpathlineto{\pgfqpoint{1.564948in}{1.005324in}}%
\pgfpathlineto{\pgfqpoint{1.565225in}{1.006137in}}%
\pgfpathlineto{\pgfqpoint{1.566925in}{1.008212in}}%
\pgfpathlineto{\pgfqpoint{1.568221in}{1.006948in}}%
\pgfpathlineto{\pgfqpoint{1.569894in}{1.005134in}}%
\pgfpathlineto{\pgfqpoint{1.570166in}{1.005921in}}%
\pgfpathlineto{\pgfqpoint{1.571630in}{1.008038in}}%
\pgfpathlineto{\pgfqpoint{1.572618in}{1.007349in}}%
\pgfpathlineto{\pgfqpoint{1.574841in}{1.004929in}}%
\pgfpathlineto{\pgfqpoint{1.575119in}{1.005731in}}%
\pgfpathlineto{\pgfqpoint{1.576816in}{1.007789in}}%
\pgfpathlineto{\pgfqpoint{1.578093in}{1.006545in}}%
\pgfpathlineto{\pgfqpoint{1.579785in}{1.004712in}}%
\pgfpathlineto{\pgfqpoint{1.580060in}{1.005527in}}%
\pgfpathlineto{\pgfqpoint{1.581616in}{1.007645in}}%
\pgfpathlineto{\pgfqpoint{1.582668in}{1.006786in}}%
\pgfpathlineto{\pgfqpoint{1.584731in}{1.004547in}}%
\pgfpathlineto{\pgfqpoint{1.585000in}{1.005316in}}%
\pgfpathlineto{\pgfqpoint{1.586606in}{1.007426in}}%
\pgfpathlineto{\pgfqpoint{1.587710in}{1.006462in}}%
\pgfpathlineto{\pgfqpoint{1.589668in}{1.004330in}}%
\pgfpathlineto{\pgfqpoint{1.589901in}{1.005012in}}%
\pgfpathlineto{\pgfqpoint{1.591471in}{1.007260in}}%
\pgfpathlineto{\pgfqpoint{1.592496in}{1.006466in}}%
\pgfpathlineto{\pgfqpoint{1.594616in}{1.004159in}}%
\pgfpathlineto{\pgfqpoint{1.594897in}{1.004965in}}%
\pgfpathlineto{\pgfqpoint{1.596426in}{1.007057in}}%
\pgfpathlineto{\pgfqpoint{1.597482in}{1.006212in}}%
\pgfpathlineto{\pgfqpoint{1.599567in}{1.003952in}}%
\pgfpathlineto{\pgfqpoint{1.599800in}{1.004640in}}%
\pgfpathlineto{\pgfqpoint{1.601413in}{1.006884in}}%
\pgfpathlineto{\pgfqpoint{1.602500in}{1.005978in}}%
\pgfpathlineto{\pgfqpoint{1.604511in}{1.003788in}}%
\pgfpathlineto{\pgfqpoint{1.604802in}{1.004592in}}%
\pgfpathlineto{\pgfqpoint{1.606274in}{1.006648in}}%
\pgfpathlineto{\pgfqpoint{1.607267in}{1.005923in}}%
\pgfpathlineto{\pgfqpoint{1.609449in}{1.003549in}}%
\pgfpathlineto{\pgfqpoint{1.609722in}{1.004347in}}%
\pgfpathlineto{\pgfqpoint{1.611344in}{1.006494in}}%
\pgfpathlineto{\pgfqpoint{1.612501in}{1.005470in}}%
\pgfpathlineto{\pgfqpoint{1.614404in}{1.003403in}}%
\pgfpathlineto{\pgfqpoint{1.614687in}{1.004209in}}%
\pgfpathlineto{\pgfqpoint{1.616369in}{1.006246in}}%
\pgfpathlineto{\pgfqpoint{1.617700in}{1.004952in}}%
\pgfpathlineto{\pgfqpoint{1.619349in}{1.003177in}}%
\pgfpathlineto{\pgfqpoint{1.619579in}{1.003867in}}%
\pgfpathlineto{\pgfqpoint{1.621147in}{1.006128in}}%
\pgfpathlineto{\pgfqpoint{1.622171in}{1.005348in}}%
\pgfpathlineto{\pgfqpoint{1.624297in}{1.003053in}}%
\pgfpathlineto{\pgfqpoint{1.624569in}{1.003846in}}%
\pgfpathlineto{\pgfqpoint{1.626039in}{1.005925in}}%
\pgfpathlineto{\pgfqpoint{1.627028in}{1.005219in}}%
\pgfpathlineto{\pgfqpoint{1.629240in}{1.002817in}}%
\pgfpathlineto{\pgfqpoint{1.629520in}{1.003621in}}%
\pgfpathlineto{\pgfqpoint{1.631267in}{1.005678in}}%
\pgfpathlineto{\pgfqpoint{1.632687in}{1.004250in}}%
\pgfpathlineto{\pgfqpoint{1.634184in}{1.002630in}}%
\pgfpathlineto{\pgfqpoint{1.634468in}{1.003433in}}%
\pgfpathlineto{\pgfqpoint{1.635976in}{1.005528in}}%
\pgfpathlineto{\pgfqpoint{1.636948in}{1.004805in}}%
\pgfpathlineto{\pgfqpoint{1.639129in}{1.002438in}}%
\pgfpathlineto{\pgfqpoint{1.639412in}{1.003245in}}%
\pgfpathlineto{\pgfqpoint{1.641038in}{1.005337in}}%
\pgfpathlineto{\pgfqpoint{1.642149in}{1.004343in}}%
\pgfpathlineto{\pgfqpoint{1.644074in}{1.002253in}}%
\pgfpathlineto{\pgfqpoint{1.644361in}{1.003050in}}%
\pgfpathlineto{\pgfqpoint{1.645847in}{1.005149in}}%
\pgfpathlineto{\pgfqpoint{1.646835in}{1.004435in}}%
\pgfpathlineto{\pgfqpoint{1.649012in}{1.002074in}}%
\pgfpathlineto{\pgfqpoint{1.649244in}{1.002757in}}%
\pgfpathlineto{\pgfqpoint{1.650826in}{1.005031in}}%
\pgfpathlineto{\pgfqpoint{1.651847in}{1.004248in}}%
\pgfpathlineto{\pgfqpoint{1.653972in}{1.001958in}}%
\pgfpathlineto{\pgfqpoint{1.654245in}{1.002760in}}%
\pgfpathlineto{\pgfqpoint{1.655837in}{1.004847in}}%
\pgfpathlineto{\pgfqpoint{1.657002in}{1.003826in}}%
\pgfpathlineto{\pgfqpoint{1.658912in}{1.001755in}}%
\pgfpathlineto{\pgfqpoint{1.659197in}{1.002558in}}%
\pgfpathlineto{\pgfqpoint{1.660694in}{1.004650in}}%
\pgfpathlineto{\pgfqpoint{1.661639in}{1.003970in}}%
\pgfpathlineto{\pgfqpoint{1.663854in}{1.001574in}}%
\pgfpathlineto{\pgfqpoint{1.664118in}{1.002340in}}%
\pgfpathlineto{\pgfqpoint{1.665739in}{1.004507in}}%
\pgfpathlineto{\pgfqpoint{1.666846in}{1.003550in}}%
\pgfpathlineto{\pgfqpoint{1.668798in}{1.001442in}}%
\pgfpathlineto{\pgfqpoint{1.669051in}{1.002187in}}%
\pgfpathlineto{\pgfqpoint{1.670631in}{1.004380in}}%
\pgfpathlineto{\pgfqpoint{1.671741in}{1.003467in}}%
\pgfpathlineto{\pgfqpoint{1.673744in}{1.001299in}}%
\pgfpathlineto{\pgfqpoint{1.674019in}{1.002095in}}%
\pgfpathlineto{\pgfqpoint{1.675600in}{1.004178in}}%
\pgfpathlineto{\pgfqpoint{1.676669in}{1.003274in}}%
\pgfpathlineto{\pgfqpoint{1.678694in}{1.001089in}}%
\pgfpathlineto{\pgfqpoint{1.678956in}{1.001852in}}%
\pgfpathlineto{\pgfqpoint{1.680586in}{1.004044in}}%
\pgfpathlineto{\pgfqpoint{1.681679in}{1.003105in}}%
\pgfpathlineto{\pgfqpoint{1.683635in}{1.000996in}}%
\pgfpathlineto{\pgfqpoint{1.683913in}{1.001802in}}%
\pgfpathlineto{\pgfqpoint{1.685478in}{1.003901in}}%
\pgfpathlineto{\pgfqpoint{1.686544in}{1.003022in}}%
\pgfpathlineto{\pgfqpoint{1.688584in}{1.000813in}}%
\pgfpathlineto{\pgfqpoint{1.688867in}{1.001613in}}%
\pgfpathlineto{\pgfqpoint{1.690450in}{1.003723in}}%
\pgfpathlineto{\pgfqpoint{1.691479in}{1.002871in}}%
\pgfpathlineto{\pgfqpoint{1.693524in}{1.000653in}}%
\pgfpathlineto{\pgfqpoint{1.693806in}{1.001459in}}%
\pgfpathlineto{\pgfqpoint{1.695403in}{1.003580in}}%
\pgfpathlineto{\pgfqpoint{1.696504in}{1.002640in}}%
\pgfpathlineto{\pgfqpoint{1.698482in}{1.000510in}}%
\pgfpathlineto{\pgfqpoint{1.698756in}{1.001311in}}%
\pgfpathlineto{\pgfqpoint{1.700258in}{1.003409in}}%
\pgfpathlineto{\pgfqpoint{1.701290in}{1.002634in}}%
\pgfpathlineto{\pgfqpoint{1.703429in}{1.000331in}}%
\pgfpathlineto{\pgfqpoint{1.703698in}{1.001127in}}%
\pgfpathlineto{\pgfqpoint{1.705234in}{1.003265in}}%
\pgfpathlineto{\pgfqpoint{1.706280in}{1.002453in}}%
\pgfpathlineto{\pgfqpoint{1.708372in}{1.000211in}}%
\pgfpathlineto{\pgfqpoint{1.708601in}{1.000899in}}%
\pgfpathlineto{\pgfqpoint{1.710165in}{1.003136in}}%
\pgfpathlineto{\pgfqpoint{1.711189in}{1.002357in}}%
\pgfpathlineto{\pgfqpoint{1.713312in}{1.000068in}}%
\pgfpathlineto{\pgfqpoint{1.713582in}{1.000865in}}%
\pgfpathlineto{\pgfqpoint{1.715197in}{1.003014in}}%
\pgfpathlineto{\pgfqpoint{1.716288in}{1.002083in}}%
\pgfpathlineto{\pgfqpoint{1.718257in}{0.999963in}}%
\pgfpathlineto{\pgfqpoint{1.718538in}{1.000752in}}%
\pgfpathlineto{\pgfqpoint{1.720117in}{1.002829in}}%
\pgfpathlineto{\pgfqpoint{1.721231in}{1.001880in}}%
\pgfpathlineto{\pgfqpoint{1.723195in}{0.999762in}}%
\pgfpathlineto{\pgfqpoint{1.723455in}{1.000525in}}%
\pgfpathlineto{\pgfqpoint{1.725062in}{1.002732in}}%
\pgfpathlineto{\pgfqpoint{1.726148in}{1.001831in}}%
\pgfpathlineto{\pgfqpoint{1.728153in}{0.999682in}}%
\pgfpathlineto{\pgfqpoint{1.728423in}{1.000480in}}%
\pgfpathlineto{\pgfqpoint{1.729978in}{1.002611in}}%
\pgfpathlineto{\pgfqpoint{1.731032in}{1.001770in}}%
\pgfpathlineto{\pgfqpoint{1.733104in}{0.999550in}}%
\pgfpathlineto{\pgfqpoint{1.733374in}{1.000351in}}%
\pgfpathlineto{\pgfqpoint{1.734931in}{1.002466in}}%
\pgfpathlineto{\pgfqpoint{1.736001in}{1.001601in}}%
\pgfpathlineto{\pgfqpoint{1.738038in}{0.999397in}}%
\pgfpathlineto{\pgfqpoint{1.738332in}{1.000200in}}%
\pgfpathlineto{\pgfqpoint{1.739823in}{1.002271in}}%
\pgfpathlineto{\pgfqpoint{1.740791in}{1.001570in}}%
\pgfpathlineto{\pgfqpoint{1.742980in}{0.999212in}}%
\pgfpathlineto{\pgfqpoint{1.743248in}{0.999979in}}%
\pgfpathlineto{\pgfqpoint{1.744855in}{1.002155in}}%
\pgfpathlineto{\pgfqpoint{1.745921in}{1.001269in}}%
\pgfpathlineto{\pgfqpoint{1.747923in}{0.999103in}}%
\pgfpathlineto{\pgfqpoint{1.748206in}{0.999902in}}%
\pgfpathlineto{\pgfqpoint{1.749670in}{1.001991in}}%
\pgfpathlineto{\pgfqpoint{1.750660in}{1.001308in}}%
\pgfpathlineto{\pgfqpoint{1.752875in}{0.998915in}}%
\pgfpathlineto{\pgfqpoint{1.753161in}{0.999717in}}%
\pgfpathlineto{\pgfqpoint{1.754761in}{1.001823in}}%
\pgfpathlineto{\pgfqpoint{1.755865in}{1.000877in}}%
\pgfpathlineto{\pgfqpoint{1.757816in}{0.998785in}}%
\pgfpathlineto{\pgfqpoint{1.758067in}{0.999525in}}%
\pgfpathlineto{\pgfqpoint{1.759649in}{1.001739in}}%
\pgfpathlineto{\pgfqpoint{1.760712in}{1.000893in}}%
\pgfpathlineto{\pgfqpoint{1.762778in}{0.998686in}}%
\pgfpathlineto{\pgfqpoint{1.763046in}{0.999483in}}%
\pgfpathlineto{\pgfqpoint{1.764723in}{1.001586in}}%
\pgfpathlineto{\pgfqpoint{1.765987in}{1.000403in}}%
\pgfpathlineto{\pgfqpoint{1.767711in}{0.998541in}}%
\pgfpathlineto{\pgfqpoint{1.768003in}{0.999324in}}%
\pgfpathlineto{\pgfqpoint{1.769468in}{1.001406in}}%
\pgfpathlineto{\pgfqpoint{1.770407in}{1.000775in}}%
\pgfpathlineto{\pgfqpoint{1.772649in}{0.998358in}}%
\pgfpathlineto{\pgfqpoint{1.772928in}{0.999154in}}%
\pgfpathlineto{\pgfqpoint{1.774442in}{1.001304in}}%
\pgfpathlineto{\pgfqpoint{1.775431in}{1.000592in}}%
\pgfpathlineto{\pgfqpoint{1.777606in}{0.998262in}}%
\pgfpathlineto{\pgfqpoint{1.777881in}{0.999049in}}%
\pgfpathlineto{\pgfqpoint{1.779538in}{1.001130in}}%
\pgfpathlineto{\pgfqpoint{1.780766in}{0.999996in}}%
\pgfpathlineto{\pgfqpoint{1.782557in}{0.998095in}}%
\pgfpathlineto{\pgfqpoint{1.782825in}{0.998894in}}%
\pgfpathlineto{\pgfqpoint{1.784438in}{1.001039in}}%
\pgfpathlineto{\pgfqpoint{1.785585in}{1.000050in}}%
\pgfpathlineto{\pgfqpoint{1.787489in}{0.998002in}}%
\pgfpathlineto{\pgfqpoint{1.787767in}{0.998804in}}%
\pgfpathlineto{\pgfqpoint{1.789345in}{1.000934in}}%
\pgfpathlineto{\pgfqpoint{1.790423in}{1.000047in}}%
\pgfpathlineto{\pgfqpoint{1.792439in}{0.997885in}}%
\pgfpathlineto{\pgfqpoint{1.792704in}{0.998655in}}%
\pgfpathlineto{\pgfqpoint{1.794328in}{1.000831in}}%
\pgfpathlineto{\pgfqpoint{1.795422in}{0.999901in}}%
\pgfpathlineto{\pgfqpoint{1.797377in}{0.997798in}}%
\pgfpathlineto{\pgfqpoint{1.797656in}{0.998594in}}%
\pgfpathlineto{\pgfqpoint{1.799149in}{1.000726in}}%
\pgfpathlineto{\pgfqpoint{1.800138in}{1.000029in}}%
\pgfpathlineto{\pgfqpoint{1.802327in}{0.997667in}}%
\pgfpathlineto{\pgfqpoint{1.802616in}{0.998450in}}%
\pgfpathlineto{\pgfqpoint{1.804096in}{1.000563in}}%
\pgfpathlineto{\pgfqpoint{1.805024in}{0.999942in}}%
\pgfpathlineto{\pgfqpoint{1.807271in}{0.997540in}}%
\pgfpathlineto{\pgfqpoint{1.807538in}{0.998333in}}%
\pgfpathlineto{\pgfqpoint{1.809098in}{1.000508in}}%
\pgfpathlineto{\pgfqpoint{1.810145in}{0.999695in}}%
\pgfpathlineto{\pgfqpoint{1.812222in}{0.997461in}}%
\pgfpathlineto{\pgfqpoint{1.812507in}{0.998268in}}%
\pgfpathlineto{\pgfqpoint{1.814205in}{1.000330in}}%
\pgfpathlineto{\pgfqpoint{1.815400in}{0.999198in}}%
\pgfpathlineto{\pgfqpoint{1.817160in}{0.997302in}}%
\pgfpathlineto{\pgfqpoint{1.817446in}{0.998116in}}%
\pgfpathlineto{\pgfqpoint{1.818907in}{1.000252in}}%
\pgfpathlineto{\pgfqpoint{1.819896in}{0.999601in}}%
\pgfpathlineto{\pgfqpoint{1.822125in}{0.997227in}}%
\pgfpathlineto{\pgfqpoint{1.822395in}{0.998021in}}%
\pgfpathlineto{\pgfqpoint{1.823960in}{1.000115in}}%
\pgfpathlineto{\pgfqpoint{1.825042in}{0.999226in}}%
\pgfpathlineto{\pgfqpoint{1.827056in}{0.997058in}}%
\pgfpathlineto{\pgfqpoint{1.827344in}{0.997864in}}%
\pgfpathlineto{\pgfqpoint{1.828935in}{0.999968in}}%
\pgfpathlineto{\pgfqpoint{1.830004in}{0.999074in}}%
\pgfpathlineto{\pgfqpoint{1.831998in}{0.996940in}}%
\pgfpathlineto{\pgfqpoint{1.832237in}{0.997650in}}%
\pgfpathlineto{\pgfqpoint{1.833859in}{0.999907in}}%
\pgfpathlineto{\pgfqpoint{1.834992in}{0.998960in}}%
\pgfpathlineto{\pgfqpoint{1.836943in}{0.996865in}}%
\pgfpathlineto{\pgfqpoint{1.837181in}{0.997563in}}%
\pgfpathlineto{\pgfqpoint{1.838813in}{0.999814in}}%
\pgfpathlineto{\pgfqpoint{1.839904in}{0.998908in}}%
\pgfpathlineto{\pgfqpoint{1.841895in}{0.996772in}}%
\pgfpathlineto{\pgfqpoint{1.842176in}{0.997574in}}%
\pgfpathlineto{\pgfqpoint{1.843919in}{0.999642in}}%
\pgfpathlineto{\pgfqpoint{1.845328in}{0.998253in}}%
\pgfpathlineto{\pgfqpoint{1.846834in}{0.996645in}}%
\pgfpathlineto{\pgfqpoint{1.847081in}{0.997376in}}%
\pgfpathlineto{\pgfqpoint{1.848691in}{0.999622in}}%
\pgfpathlineto{\pgfqpoint{1.849729in}{0.998792in}}%
\pgfpathlineto{\pgfqpoint{1.851786in}{0.996583in}}%
\pgfpathlineto{\pgfqpoint{1.852065in}{0.997343in}}%
\pgfpathlineto{\pgfqpoint{1.853704in}{0.999441in}}%
\pgfpathlineto{\pgfqpoint{1.854753in}{0.998526in}}%
\pgfpathlineto{\pgfqpoint{1.856726in}{0.996433in}}%
\pgfpathlineto{\pgfqpoint{1.856992in}{0.997226in}}%
\pgfpathlineto{\pgfqpoint{1.858482in}{0.999415in}}%
\pgfpathlineto{\pgfqpoint{1.859471in}{0.998754in}}%
\pgfpathlineto{\pgfqpoint{1.861679in}{0.996401in}}%
\pgfpathlineto{\pgfqpoint{1.861949in}{0.997195in}}%
\pgfpathlineto{\pgfqpoint{1.863566in}{0.999323in}}%
\pgfpathlineto{\pgfqpoint{1.864659in}{0.998388in}}%
\pgfpathlineto{\pgfqpoint{1.866625in}{0.996300in}}%
\pgfpathlineto{\pgfqpoint{1.866896in}{0.997098in}}%
\pgfpathlineto{\pgfqpoint{1.868440in}{0.999253in}}%
\pgfpathlineto{\pgfqpoint{1.869505in}{0.998429in}}%
\pgfpathlineto{\pgfqpoint{1.871563in}{0.996223in}}%
\pgfpathlineto{\pgfqpoint{1.871837in}{0.997026in}}%
\pgfpathlineto{\pgfqpoint{1.873478in}{0.999176in}}%
\pgfpathlineto{\pgfqpoint{1.874602in}{0.998195in}}%
\pgfpathlineto{\pgfqpoint{1.876507in}{0.996154in}}%
\pgfpathlineto{\pgfqpoint{1.876787in}{0.996961in}}%
\pgfpathlineto{\pgfqpoint{1.878395in}{0.999083in}}%
\pgfpathlineto{\pgfqpoint{1.879501in}{0.998141in}}%
\pgfpathlineto{\pgfqpoint{1.881452in}{0.996050in}}%
\pgfpathlineto{\pgfqpoint{1.881718in}{0.996830in}}%
\pgfpathlineto{\pgfqpoint{1.883333in}{0.999022in}}%
\pgfpathlineto{\pgfqpoint{1.884456in}{0.998082in}}%
\pgfpathlineto{\pgfqpoint{1.886415in}{0.996002in}}%
\pgfpathlineto{\pgfqpoint{1.886683in}{0.996791in}}%
\pgfpathlineto{\pgfqpoint{1.888235in}{0.998898in}}%
\pgfpathlineto{\pgfqpoint{1.889295in}{0.998057in}}%
\pgfpathlineto{\pgfqpoint{1.891361in}{0.995863in}}%
\pgfpathlineto{\pgfqpoint{1.891629in}{0.996661in}}%
\pgfpathlineto{\pgfqpoint{1.893294in}{0.998786in}}%
\pgfpathlineto{\pgfqpoint{1.894507in}{0.997684in}}%
\pgfpathlineto{\pgfqpoint{1.896292in}{0.995775in}}%
\pgfpathlineto{\pgfqpoint{1.896558in}{0.996548in}}%
\pgfpathlineto{\pgfqpoint{1.898192in}{0.998699in}}%
\pgfpathlineto{\pgfqpoint{1.899333in}{0.997707in}}%
\pgfpathlineto{\pgfqpoint{1.901234in}{0.995670in}}%
\pgfpathlineto{\pgfqpoint{1.901510in}{0.996474in}}%
\pgfpathlineto{\pgfqpoint{1.903110in}{0.998624in}}%
\pgfpathlineto{\pgfqpoint{1.904184in}{0.997737in}}%
\pgfpathlineto{\pgfqpoint{1.906180in}{0.995600in}}%
\pgfpathlineto{\pgfqpoint{1.906456in}{0.996397in}}%
\pgfpathlineto{\pgfqpoint{1.908052in}{0.998541in}}%
\pgfpathlineto{\pgfqpoint{1.909160in}{0.997617in}}%
\pgfpathlineto{\pgfqpoint{1.911130in}{0.995515in}}%
\pgfpathlineto{\pgfqpoint{1.911394in}{0.996283in}}%
\pgfpathlineto{\pgfqpoint{1.913023in}{0.998463in}}%
\pgfpathlineto{\pgfqpoint{1.914171in}{0.997480in}}%
\pgfpathlineto{\pgfqpoint{1.916077in}{0.995443in}}%
\pgfpathlineto{\pgfqpoint{1.916364in}{0.996227in}}%
\pgfpathlineto{\pgfqpoint{1.917873in}{0.998310in}}%
\pgfpathlineto{\pgfqpoint{1.918830in}{0.997617in}}%
\pgfpathlineto{\pgfqpoint{1.921013in}{0.995284in}}%
\pgfpathlineto{\pgfqpoint{1.921261in}{0.996019in}}%
\pgfpathlineto{\pgfqpoint{1.922860in}{0.998291in}}%
\pgfpathlineto{\pgfqpoint{1.923894in}{0.997494in}}%
\pgfpathlineto{\pgfqpoint{1.925972in}{0.995285in}}%
\pgfpathlineto{\pgfqpoint{1.926240in}{0.996072in}}%
\pgfpathlineto{\pgfqpoint{1.927747in}{0.998199in}}%
\pgfpathlineto{\pgfqpoint{1.928779in}{0.997445in}}%
\pgfpathlineto{\pgfqpoint{1.930913in}{0.995158in}}%
\pgfpathlineto{\pgfqpoint{1.931193in}{0.995926in}}%
\pgfpathlineto{\pgfqpoint{1.932888in}{0.998036in}}%
\pgfpathlineto{\pgfqpoint{1.934102in}{0.996906in}}%
\pgfpathlineto{\pgfqpoint{1.935859in}{0.995049in}}%
\pgfpathlineto{\pgfqpoint{1.936131in}{0.995846in}}%
\pgfpathlineto{\pgfqpoint{1.937816in}{0.997971in}}%
\pgfpathlineto{\pgfqpoint{1.938944in}{0.996946in}}%
\pgfpathlineto{\pgfqpoint{1.940799in}{0.994965in}}%
\pgfpathlineto{\pgfqpoint{1.941077in}{0.995777in}}%
\pgfpathlineto{\pgfqpoint{1.942685in}{0.997932in}}%
\pgfpathlineto{\pgfqpoint{1.943789in}{0.997008in}}%
\pgfpathlineto{\pgfqpoint{1.945761in}{0.994920in}}%
\pgfpathlineto{\pgfqpoint{1.946032in}{0.995720in}}%
\pgfpathlineto{\pgfqpoint{1.947601in}{0.997826in}}%
\pgfpathlineto{\pgfqpoint{1.948676in}{0.996954in}}%
\pgfpathlineto{\pgfqpoint{1.950692in}{0.994793in}}%
\pgfpathlineto{\pgfqpoint{1.950980in}{0.995591in}}%
\pgfpathlineto{\pgfqpoint{1.952504in}{0.997697in}}%
\pgfpathlineto{\pgfqpoint{1.953478in}{0.996981in}}%
\pgfpathlineto{\pgfqpoint{1.955631in}{0.994677in}}%
\pgfpathlineto{\pgfqpoint{1.955907in}{0.995470in}}%
\pgfpathlineto{\pgfqpoint{1.957352in}{0.997633in}}%
\pgfpathlineto{\pgfqpoint{1.958341in}{0.997023in}}%
\pgfpathlineto{\pgfqpoint{1.960586in}{0.994621in}}%
\pgfpathlineto{\pgfqpoint{1.960878in}{0.995416in}}%
\pgfpathlineto{\pgfqpoint{1.962370in}{0.997503in}}%
\pgfpathlineto{\pgfqpoint{1.963406in}{0.996747in}}%
\pgfpathlineto{\pgfqpoint{1.965523in}{0.994495in}}%
\pgfpathlineto{\pgfqpoint{1.965797in}{0.995299in}}%
\pgfpathlineto{\pgfqpoint{1.967428in}{0.997474in}}%
\pgfpathlineto{\pgfqpoint{1.968519in}{0.996553in}}%
\pgfpathlineto{\pgfqpoint{1.970471in}{0.994466in}}%
\pgfpathlineto{\pgfqpoint{1.970751in}{0.995276in}}%
\pgfpathlineto{\pgfqpoint{1.972294in}{0.997409in}}%
\pgfpathlineto{\pgfqpoint{1.973337in}{0.996613in}}%
\pgfpathlineto{\pgfqpoint{1.975421in}{0.994392in}}%
\pgfpathlineto{\pgfqpoint{1.975692in}{0.995165in}}%
\pgfpathlineto{\pgfqpoint{1.977302in}{0.997317in}}%
\pgfpathlineto{\pgfqpoint{1.978363in}{0.996439in}}%
\pgfpathlineto{\pgfqpoint{1.980359in}{0.994310in}}%
\pgfpathlineto{\pgfqpoint{1.980638in}{0.995118in}}%
\pgfpathlineto{\pgfqpoint{1.982239in}{0.997266in}}%
\pgfpathlineto{\pgfqpoint{1.983350in}{0.996339in}}%
\pgfpathlineto{\pgfqpoint{1.985305in}{0.994250in}}%
\pgfpathlineto{\pgfqpoint{1.985582in}{0.995046in}}%
\pgfpathlineto{\pgfqpoint{1.987117in}{0.997213in}}%
\pgfpathlineto{\pgfqpoint{1.988137in}{0.996461in}}%
\pgfpathlineto{\pgfqpoint{1.990253in}{0.994200in}}%
\pgfpathlineto{\pgfqpoint{1.990532in}{0.995011in}}%
\pgfpathlineto{\pgfqpoint{1.992104in}{0.997159in}}%
\pgfpathlineto{\pgfqpoint{1.993168in}{0.996314in}}%
\pgfpathlineto{\pgfqpoint{1.995206in}{0.994153in}}%
\pgfpathlineto{\pgfqpoint{1.995474in}{0.994940in}}%
\pgfpathlineto{\pgfqpoint{1.996962in}{0.997094in}}%
\pgfpathlineto{\pgfqpoint{1.997950in}{0.996426in}}%
\pgfpathlineto{\pgfqpoint{2.000151in}{0.994089in}}%
\pgfpathlineto{\pgfqpoint{2.000421in}{0.994878in}}%
\pgfpathlineto{\pgfqpoint{2.001865in}{0.997029in}}%
\pgfpathlineto{\pgfqpoint{2.002854in}{0.996420in}}%
\pgfpathlineto{\pgfqpoint{2.005107in}{0.994041in}}%
\pgfpathlineto{\pgfqpoint{2.005374in}{0.994833in}}%
\pgfpathlineto{\pgfqpoint{2.007070in}{0.996927in}}%
\pgfpathlineto{\pgfqpoint{2.008364in}{0.995710in}}%
\pgfpathlineto{\pgfqpoint{2.010051in}{0.993932in}}%
\pgfpathlineto{\pgfqpoint{2.010321in}{0.994730in}}%
\pgfpathlineto{\pgfqpoint{2.011983in}{0.996858in}}%
\pgfpathlineto{\pgfqpoint{2.013170in}{0.995791in}}%
\pgfpathlineto{\pgfqpoint{2.014988in}{0.993869in}}%
\pgfpathlineto{\pgfqpoint{2.015260in}{0.994669in}}%
\pgfpathlineto{\pgfqpoint{2.016866in}{0.996805in}}%
\pgfpathlineto{\pgfqpoint{2.017974in}{0.995877in}}%
\pgfpathlineto{\pgfqpoint{2.019925in}{0.993795in}}%
\pgfpathlineto{\pgfqpoint{2.020178in}{0.994537in}}%
\pgfpathlineto{\pgfqpoint{2.021775in}{0.996755in}}%
\pgfpathlineto{\pgfqpoint{2.022843in}{0.995907in}}%
\pgfpathlineto{\pgfqpoint{2.024871in}{0.993743in}}%
\pgfpathlineto{\pgfqpoint{2.025147in}{0.994540in}}%
\pgfpathlineto{\pgfqpoint{2.026716in}{0.996703in}}%
\pgfpathlineto{\pgfqpoint{2.027765in}{0.995884in}}%
\pgfpathlineto{\pgfqpoint{2.029826in}{0.993701in}}%
\pgfpathlineto{\pgfqpoint{2.030092in}{0.994491in}}%
\pgfpathlineto{\pgfqpoint{2.031673in}{0.996659in}}%
\pgfpathlineto{\pgfqpoint{2.032771in}{0.995777in}}%
\pgfpathlineto{\pgfqpoint{2.034767in}{0.993645in}}%
\pgfpathlineto{\pgfqpoint{2.035057in}{0.994446in}}%
\pgfpathlineto{\pgfqpoint{2.036534in}{0.996523in}}%
\pgfpathlineto{\pgfqpoint{2.037491in}{0.995871in}}%
\pgfpathlineto{\pgfqpoint{2.039714in}{0.993506in}}%
\pgfpathlineto{\pgfqpoint{2.039998in}{0.994301in}}%
\pgfpathlineto{\pgfqpoint{2.041661in}{0.996422in}}%
\pgfpathlineto{\pgfqpoint{2.042824in}{0.995381in}}%
\pgfpathlineto{\pgfqpoint{2.044653in}{0.993453in}}%
\pgfpathlineto{\pgfqpoint{2.044922in}{0.994251in}}%
\pgfpathlineto{\pgfqpoint{2.046437in}{0.996419in}}%
\pgfpathlineto{\pgfqpoint{2.047427in}{0.995727in}}%
\pgfpathlineto{\pgfqpoint{2.049599in}{0.993407in}}%
\pgfpathlineto{\pgfqpoint{2.049874in}{0.994208in}}%
\pgfpathlineto{\pgfqpoint{2.051574in}{0.996361in}}%
\pgfpathlineto{\pgfqpoint{2.052787in}{0.995245in}}%
\pgfpathlineto{\pgfqpoint{2.054549in}{0.993369in}}%
\pgfpathlineto{\pgfqpoint{2.054832in}{0.994175in}}%
\pgfpathlineto{\pgfqpoint{2.056628in}{0.996243in}}%
\pgfpathlineto{\pgfqpoint{2.058260in}{0.994595in}}%
\pgfpathlineto{\pgfqpoint{2.059490in}{0.993298in}}%
\pgfpathlineto{\pgfqpoint{2.059752in}{0.994068in}}%
\pgfpathlineto{\pgfqpoint{2.061335in}{0.996278in}}%
\pgfpathlineto{\pgfqpoint{2.062404in}{0.995441in}}%
\pgfpathlineto{\pgfqpoint{2.064442in}{0.993283in}}%
\pgfpathlineto{\pgfqpoint{2.064714in}{0.994078in}}%
\pgfpathlineto{\pgfqpoint{2.066269in}{0.996206in}}%
\pgfpathlineto{\pgfqpoint{2.067331in}{0.995376in}}%
\pgfpathlineto{\pgfqpoint{2.069397in}{0.993197in}}%
\pgfpathlineto{\pgfqpoint{2.069663in}{0.993992in}}%
\pgfpathlineto{\pgfqpoint{2.071364in}{0.996123in}}%
\pgfpathlineto{\pgfqpoint{2.072648in}{0.994927in}}%
\pgfpathlineto{\pgfqpoint{2.074328in}{0.993141in}}%
\pgfpathlineto{\pgfqpoint{2.074592in}{0.993912in}}%
\pgfpathlineto{\pgfqpoint{2.076205in}{0.996091in}}%
\pgfpathlineto{\pgfqpoint{2.077268in}{0.995221in}}%
\pgfpathlineto{\pgfqpoint{2.079278in}{0.993092in}}%
\pgfpathlineto{\pgfqpoint{2.079554in}{0.993886in}}%
\pgfpathlineto{\pgfqpoint{2.081078in}{0.996018in}}%
\pgfpathlineto{\pgfqpoint{2.082100in}{0.995269in}}%
\pgfpathlineto{\pgfqpoint{2.084234in}{0.993022in}}%
\pgfpathlineto{\pgfqpoint{2.084500in}{0.993814in}}%
\pgfpathlineto{\pgfqpoint{2.085950in}{0.995941in}}%
\pgfpathlineto{\pgfqpoint{2.086939in}{0.995312in}}%
\pgfpathlineto{\pgfqpoint{2.089169in}{0.992945in}}%
\pgfpathlineto{\pgfqpoint{2.089446in}{0.993743in}}%
\pgfpathlineto{\pgfqpoint{2.091140in}{0.995858in}}%
\pgfpathlineto{\pgfqpoint{2.092419in}{0.994669in}}%
\pgfpathlineto{\pgfqpoint{2.094105in}{0.992869in}}%
\pgfpathlineto{\pgfqpoint{2.094388in}{0.993673in}}%
\pgfpathlineto{\pgfqpoint{2.095867in}{0.995810in}}%
\pgfpathlineto{\pgfqpoint{2.096852in}{0.995153in}}%
\pgfpathlineto{\pgfqpoint{2.099053in}{0.992807in}}%
\pgfpathlineto{\pgfqpoint{2.099325in}{0.993596in}}%
\pgfpathlineto{\pgfqpoint{2.100856in}{0.995782in}}%
\pgfpathlineto{\pgfqpoint{2.101880in}{0.995042in}}%
\pgfpathlineto{\pgfqpoint{2.103998in}{0.992786in}}%
\pgfpathlineto{\pgfqpoint{2.104276in}{0.993590in}}%
\pgfpathlineto{\pgfqpoint{2.105873in}{0.995742in}}%
\pgfpathlineto{\pgfqpoint{2.106965in}{0.994848in}}%
\pgfpathlineto{\pgfqpoint{2.108962in}{0.992749in}}%
\pgfpathlineto{\pgfqpoint{2.109229in}{0.993541in}}%
\pgfpathlineto{\pgfqpoint{2.110942in}{0.995646in}}%
\pgfpathlineto{\pgfqpoint{2.112286in}{0.994370in}}%
\pgfpathlineto{\pgfqpoint{2.113895in}{0.992657in}}%
\pgfpathlineto{\pgfqpoint{2.114183in}{0.993429in}}%
\pgfpathlineto{\pgfqpoint{2.115860in}{0.995510in}}%
\pgfpathlineto{\pgfqpoint{2.116981in}{0.994494in}}%
\pgfpathlineto{\pgfqpoint{2.118827in}{0.992534in}}%
\pgfpathlineto{\pgfqpoint{2.119097in}{0.993330in}}%
\pgfpathlineto{\pgfqpoint{2.120699in}{0.995551in}}%
\pgfpathlineto{\pgfqpoint{2.121774in}{0.994692in}}%
\pgfpathlineto{\pgfqpoint{2.123783in}{0.992544in}}%
\pgfpathlineto{\pgfqpoint{2.124081in}{0.993342in}}%
\pgfpathlineto{\pgfqpoint{2.125567in}{0.995436in}}%
\pgfpathlineto{\pgfqpoint{2.126525in}{0.994780in}}%
\pgfpathlineto{\pgfqpoint{2.128714in}{0.992446in}}%
\pgfpathlineto{\pgfqpoint{2.128991in}{0.993247in}}%
\pgfpathlineto{\pgfqpoint{2.130468in}{0.995443in}}%
\pgfpathlineto{\pgfqpoint{2.131457in}{0.994812in}}%
\pgfpathlineto{\pgfqpoint{2.133671in}{0.992454in}}%
\pgfpathlineto{\pgfqpoint{2.133946in}{0.993249in}}%
\pgfpathlineto{\pgfqpoint{2.135558in}{0.995412in}}%
\pgfpathlineto{\pgfqpoint{2.136672in}{0.994484in}}%
\pgfpathlineto{\pgfqpoint{2.138633in}{0.992419in}}%
\pgfpathlineto{\pgfqpoint{2.138902in}{0.993216in}}%
\pgfpathlineto{\pgfqpoint{2.140596in}{0.995326in}}%
\pgfpathlineto{\pgfqpoint{2.141861in}{0.994152in}}%
\pgfpathlineto{\pgfqpoint{2.143579in}{0.992347in}}%
\pgfpathlineto{\pgfqpoint{2.143848in}{0.993141in}}%
\pgfpathlineto{\pgfqpoint{2.145522in}{0.995257in}}%
\pgfpathlineto{\pgfqpoint{2.146697in}{0.994195in}}%
\pgfpathlineto{\pgfqpoint{2.148515in}{0.992279in}}%
\pgfpathlineto{\pgfqpoint{2.148787in}{0.993076in}}%
\pgfpathlineto{\pgfqpoint{2.150253in}{0.995226in}}%
\pgfpathlineto{\pgfqpoint{2.151240in}{0.994593in}}%
\pgfpathlineto{\pgfqpoint{2.153450in}{0.992236in}}%
\pgfpathlineto{\pgfqpoint{2.153729in}{0.993032in}}%
\pgfpathlineto{\pgfqpoint{2.155230in}{0.995181in}}%
\pgfpathlineto{\pgfqpoint{2.156269in}{0.994446in}}%
\pgfpathlineto{\pgfqpoint{2.158415in}{0.992185in}}%
\pgfpathlineto{\pgfqpoint{2.158682in}{0.992980in}}%
\pgfpathlineto{\pgfqpoint{2.160412in}{0.995097in}}%
\pgfpathlineto{\pgfqpoint{2.161737in}{0.993834in}}%
\pgfpathlineto{\pgfqpoint{2.163349in}{0.992123in}}%
\pgfpathlineto{\pgfqpoint{2.163631in}{0.992925in}}%
\pgfpathlineto{\pgfqpoint{2.165369in}{0.995019in}}%
\pgfpathlineto{\pgfqpoint{2.166658in}{0.993784in}}%
\pgfpathlineto{\pgfqpoint{2.168289in}{0.992059in}}%
\pgfpathlineto{\pgfqpoint{2.168562in}{0.992858in}}%
\pgfpathlineto{\pgfqpoint{2.170087in}{0.995041in}}%
\pgfpathlineto{\pgfqpoint{2.171113in}{0.994308in}}%
\pgfpathlineto{\pgfqpoint{2.173235in}{0.992049in}}%
\pgfpathlineto{\pgfqpoint{2.173511in}{0.992853in}}%
\pgfpathlineto{\pgfqpoint{2.175109in}{0.995011in}}%
\pgfpathlineto{\pgfqpoint{2.176159in}{0.994165in}}%
\pgfpathlineto{\pgfqpoint{2.178184in}{0.992016in}}%
\pgfpathlineto{\pgfqpoint{2.178458in}{0.992823in}}%
\pgfpathlineto{\pgfqpoint{2.179969in}{0.994993in}}%
\pgfpathlineto{\pgfqpoint{2.180958in}{0.994312in}}%
\pgfpathlineto{\pgfqpoint{2.183126in}{0.992005in}}%
\pgfpathlineto{\pgfqpoint{2.183399in}{0.992795in}}%
\pgfpathlineto{\pgfqpoint{2.185006in}{0.994947in}}%
\pgfpathlineto{\pgfqpoint{2.186075in}{0.994069in}}%
\pgfpathlineto{\pgfqpoint{2.188071in}{0.991945in}}%
\pgfpathlineto{\pgfqpoint{2.188348in}{0.992752in}}%
\pgfpathlineto{\pgfqpoint{2.190007in}{0.994928in}}%
\pgfpathlineto{\pgfqpoint{2.191152in}{0.993930in}}%
\pgfpathlineto{\pgfqpoint{2.193021in}{0.991950in}}%
\pgfpathlineto{\pgfqpoint{2.193277in}{0.992697in}}%
\pgfpathlineto{\pgfqpoint{2.194918in}{0.994894in}}%
\pgfpathlineto{\pgfqpoint{2.196018in}{0.993966in}}%
\pgfpathlineto{\pgfqpoint{2.197966in}{0.991903in}}%
\pgfpathlineto{\pgfqpoint{2.198237in}{0.992705in}}%
\pgfpathlineto{\pgfqpoint{2.199714in}{0.994877in}}%
\pgfpathlineto{\pgfqpoint{2.200703in}{0.994241in}}%
\pgfpathlineto{\pgfqpoint{2.202907in}{0.991895in}}%
\pgfpathlineto{\pgfqpoint{2.203185in}{0.992694in}}%
\pgfpathlineto{\pgfqpoint{2.204681in}{0.994844in}}%
\pgfpathlineto{\pgfqpoint{2.205670in}{0.994171in}}%
\pgfpathlineto{\pgfqpoint{2.207853in}{0.991850in}}%
\pgfpathlineto{\pgfqpoint{2.208108in}{0.992597in}}%
\pgfpathlineto{\pgfqpoint{2.209691in}{0.994823in}}%
\pgfpathlineto{\pgfqpoint{2.210754in}{0.994005in}}%
\pgfpathlineto{\pgfqpoint{2.212798in}{0.991833in}}%
\pgfpathlineto{\pgfqpoint{2.213077in}{0.992639in}}%
\pgfpathlineto{\pgfqpoint{2.214600in}{0.994791in}}%
\pgfpathlineto{\pgfqpoint{2.215624in}{0.994051in}}%
\pgfpathlineto{\pgfqpoint{2.217744in}{0.991796in}}%
\pgfpathlineto{\pgfqpoint{2.218022in}{0.992604in}}%
\pgfpathlineto{\pgfqpoint{2.219701in}{0.994753in}}%
\pgfpathlineto{\pgfqpoint{2.220961in}{0.993605in}}%
\pgfpathlineto{\pgfqpoint{2.222706in}{0.991773in}}%
\pgfpathlineto{\pgfqpoint{2.222976in}{0.992571in}}%
\pgfpathlineto{\pgfqpoint{2.224622in}{0.994701in}}%
\pgfpathlineto{\pgfqpoint{2.225784in}{0.993683in}}%
\pgfpathlineto{\pgfqpoint{2.227640in}{0.991709in}}%
\pgfpathlineto{\pgfqpoint{2.227920in}{0.992467in}}%
\pgfpathlineto{\pgfqpoint{2.229576in}{0.994581in}}%
\pgfpathlineto{\pgfqpoint{2.230796in}{0.993486in}}%
\pgfpathlineto{\pgfqpoint{2.232577in}{0.991609in}}%
\pgfpathlineto{\pgfqpoint{2.232847in}{0.992405in}}%
\pgfpathlineto{\pgfqpoint{2.234384in}{0.994608in}}%
\pgfpathlineto{\pgfqpoint{2.235407in}{0.993875in}}%
\pgfpathlineto{\pgfqpoint{2.237526in}{0.991625in}}%
\pgfpathlineto{\pgfqpoint{2.237802in}{0.992423in}}%
\pgfpathlineto{\pgfqpoint{2.239408in}{0.994574in}}%
\pgfpathlineto{\pgfqpoint{2.240495in}{0.993677in}}%
\pgfpathlineto{\pgfqpoint{2.242468in}{0.991579in}}%
\pgfpathlineto{\pgfqpoint{2.242749in}{0.992386in}}%
\pgfpathlineto{\pgfqpoint{2.244272in}{0.994561in}}%
\pgfpathlineto{\pgfqpoint{2.245297in}{0.993827in}}%
\pgfpathlineto{\pgfqpoint{2.247420in}{0.991576in}}%
\pgfpathlineto{\pgfqpoint{2.247694in}{0.992381in}}%
\pgfpathlineto{\pgfqpoint{2.249163in}{0.994535in}}%
\pgfpathlineto{\pgfqpoint{2.250152in}{0.993899in}}%
\pgfpathlineto{\pgfqpoint{2.252362in}{0.991549in}}%
\pgfpathlineto{\pgfqpoint{2.252637in}{0.992348in}}%
\pgfpathlineto{\pgfqpoint{2.254245in}{0.994525in}}%
\pgfpathlineto{\pgfqpoint{2.255351in}{0.993618in}}%
\pgfpathlineto{\pgfqpoint{2.257312in}{0.991543in}}%
\pgfpathlineto{\pgfqpoint{2.257546in}{0.992222in}}%
\pgfpathlineto{\pgfqpoint{2.259176in}{0.994465in}}%
\pgfpathlineto{\pgfqpoint{2.260239in}{0.993606in}}%
\pgfpathlineto{\pgfqpoint{2.262253in}{0.991475in}}%
\pgfpathlineto{\pgfqpoint{2.262526in}{0.992276in}}%
\pgfpathlineto{\pgfqpoint{2.264118in}{0.994469in}}%
\pgfpathlineto{\pgfqpoint{2.265200in}{0.993610in}}%
\pgfpathlineto{\pgfqpoint{2.267198in}{0.991489in}}%
\pgfpathlineto{\pgfqpoint{2.267475in}{0.992285in}}%
\pgfpathlineto{\pgfqpoint{2.268985in}{0.994421in}}%
\pgfpathlineto{\pgfqpoint{2.270017in}{0.993681in}}%
\pgfpathlineto{\pgfqpoint{2.272150in}{0.991412in}}%
\pgfpathlineto{\pgfqpoint{2.272438in}{0.992188in}}%
\pgfpathlineto{\pgfqpoint{2.274113in}{0.994291in}}%
\pgfpathlineto{\pgfqpoint{2.275278in}{0.993239in}}%
\pgfpathlineto{\pgfqpoint{2.277086in}{0.991334in}}%
\pgfpathlineto{\pgfqpoint{2.277349in}{0.992116in}}%
\pgfpathlineto{\pgfqpoint{2.278914in}{0.994339in}}%
\pgfpathlineto{\pgfqpoint{2.279972in}{0.993544in}}%
\pgfpathlineto{\pgfqpoint{2.282034in}{0.991354in}}%
\pgfpathlineto{\pgfqpoint{2.282312in}{0.992156in}}%
\pgfpathlineto{\pgfqpoint{2.283851in}{0.994311in}}%
\pgfpathlineto{\pgfqpoint{2.284868in}{0.993565in}}%
\pgfpathlineto{\pgfqpoint{2.286983in}{0.991320in}}%
\pgfpathlineto{\pgfqpoint{2.287245in}{0.992084in}}%
\pgfpathlineto{\pgfqpoint{2.288849in}{0.994282in}}%
\pgfpathlineto{\pgfqpoint{2.289943in}{0.993399in}}%
\pgfpathlineto{\pgfqpoint{2.291930in}{0.991296in}}%
\pgfpathlineto{\pgfqpoint{2.292201in}{0.992088in}}%
\pgfpathlineto{\pgfqpoint{2.293720in}{0.994251in}}%
\pgfpathlineto{\pgfqpoint{2.294748in}{0.993518in}}%
\pgfpathlineto{\pgfqpoint{2.296877in}{0.991254in}}%
\pgfpathlineto{\pgfqpoint{2.297158in}{0.992017in}}%
\pgfpathlineto{\pgfqpoint{2.298825in}{0.994135in}}%
\pgfpathlineto{\pgfqpoint{2.300049in}{0.993030in}}%
\pgfpathlineto{\pgfqpoint{2.301823in}{0.991164in}}%
\pgfpathlineto{\pgfqpoint{2.302095in}{0.991953in}}%
\pgfpathlineto{\pgfqpoint{2.303805in}{0.994090in}}%
\pgfpathlineto{\pgfqpoint{2.305094in}{0.992890in}}%
\pgfpathlineto{\pgfqpoint{2.306769in}{0.991127in}}%
\pgfpathlineto{\pgfqpoint{2.307048in}{0.991920in}}%
\pgfpathlineto{\pgfqpoint{2.308694in}{0.994040in}}%
\pgfpathlineto{\pgfqpoint{2.309846in}{0.993031in}}%
\pgfpathlineto{\pgfqpoint{2.311704in}{0.991060in}}%
\pgfpathlineto{\pgfqpoint{2.311984in}{0.991861in}}%
\pgfpathlineto{\pgfqpoint{2.313450in}{0.994016in}}%
\pgfpathlineto{\pgfqpoint{2.314440in}{0.993387in}}%
\pgfpathlineto{\pgfqpoint{2.316653in}{0.991035in}}%
\pgfpathlineto{\pgfqpoint{2.316930in}{0.991843in}}%
\pgfpathlineto{\pgfqpoint{2.318551in}{0.994014in}}%
\pgfpathlineto{\pgfqpoint{2.319705in}{0.993040in}}%
\pgfpathlineto{\pgfqpoint{2.321602in}{0.991033in}}%
\pgfpathlineto{\pgfqpoint{2.321875in}{0.991831in}}%
\pgfpathlineto{\pgfqpoint{2.323328in}{0.993989in}}%
\pgfpathlineto{\pgfqpoint{2.324317in}{0.993377in}}%
\pgfpathlineto{\pgfqpoint{2.326548in}{0.991003in}}%
\pgfpathlineto{\pgfqpoint{2.326831in}{0.991800in}}%
\pgfpathlineto{\pgfqpoint{2.328336in}{0.993940in}}%
\pgfpathlineto{\pgfqpoint{2.329326in}{0.993254in}}%
\pgfpathlineto{\pgfqpoint{2.331485in}{0.990959in}}%
\pgfpathlineto{\pgfqpoint{2.331718in}{0.991639in}}%
\pgfpathlineto{\pgfqpoint{2.333286in}{0.993948in}}%
\pgfpathlineto{\pgfqpoint{2.334313in}{0.993219in}}%
\pgfpathlineto{\pgfqpoint{2.336435in}{0.990967in}}%
\pgfpathlineto{\pgfqpoint{2.336709in}{0.991759in}}%
\pgfpathlineto{\pgfqpoint{2.338250in}{0.993926in}}%
\pgfpathlineto{\pgfqpoint{2.339301in}{0.993143in}}%
\pgfpathlineto{\pgfqpoint{2.341386in}{0.990932in}}%
\pgfpathlineto{\pgfqpoint{2.341664in}{0.991733in}}%
\pgfpathlineto{\pgfqpoint{2.343365in}{0.993866in}}%
\pgfpathlineto{\pgfqpoint{2.344625in}{0.992701in}}%
\pgfpathlineto{\pgfqpoint{2.346330in}{0.990898in}}%
\pgfpathlineto{\pgfqpoint{2.346609in}{0.991697in}}%
\pgfpathlineto{\pgfqpoint{2.348203in}{0.993840in}}%
\pgfpathlineto{\pgfqpoint{2.349238in}{0.993011in}}%
\pgfpathlineto{\pgfqpoint{2.351276in}{0.990860in}}%
\pgfpathlineto{\pgfqpoint{2.351499in}{0.991528in}}%
\pgfpathlineto{\pgfqpoint{2.353093in}{0.993854in}}%
\pgfpathlineto{\pgfqpoint{2.354163in}{0.993055in}}%
\pgfpathlineto{\pgfqpoint{2.356222in}{0.990866in}}%
\pgfpathlineto{\pgfqpoint{2.356505in}{0.991625in}}%
\pgfpathlineto{\pgfqpoint{2.358168in}{0.993734in}}%
\pgfpathlineto{\pgfqpoint{2.359348in}{0.992678in}}%
\pgfpathlineto{\pgfqpoint{2.361151in}{0.990769in}}%
\pgfpathlineto{\pgfqpoint{2.361428in}{0.991577in}}%
\pgfpathlineto{\pgfqpoint{2.362903in}{0.993775in}}%
\pgfpathlineto{\pgfqpoint{2.363892in}{0.993157in}}%
\pgfpathlineto{\pgfqpoint{2.366118in}{0.990808in}}%
\pgfpathlineto{\pgfqpoint{2.366385in}{0.991595in}}%
\pgfpathlineto{\pgfqpoint{2.367961in}{0.993755in}}%
\pgfpathlineto{\pgfqpoint{2.369027in}{0.992918in}}%
\pgfpathlineto{\pgfqpoint{2.371058in}{0.990760in}}%
\pgfpathlineto{\pgfqpoint{2.371338in}{0.991517in}}%
\pgfpathlineto{\pgfqpoint{2.372994in}{0.993629in}}%
\pgfpathlineto{\pgfqpoint{2.374238in}{0.992512in}}%
\pgfpathlineto{\pgfqpoint{2.375995in}{0.990662in}}%
\pgfpathlineto{\pgfqpoint{2.376269in}{0.991463in}}%
\pgfpathlineto{\pgfqpoint{2.377924in}{0.993647in}}%
\pgfpathlineto{\pgfqpoint{2.379068in}{0.992658in}}%
\pgfpathlineto{\pgfqpoint{2.380961in}{0.990675in}}%
\pgfpathlineto{\pgfqpoint{2.381227in}{0.991466in}}%
\pgfpathlineto{\pgfqpoint{2.382918in}{0.993592in}}%
\pgfpathlineto{\pgfqpoint{2.384173in}{0.992441in}}%
\pgfpathlineto{\pgfqpoint{2.385894in}{0.990625in}}%
\pgfpathlineto{\pgfqpoint{2.386168in}{0.991427in}}%
\pgfpathlineto{\pgfqpoint{2.387656in}{0.993586in}}%
\pgfpathlineto{\pgfqpoint{2.388645in}{0.992931in}}%
\pgfpathlineto{\pgfqpoint{2.390843in}{0.990614in}}%
\pgfpathlineto{\pgfqpoint{2.391112in}{0.991408in}}%
\pgfpathlineto{\pgfqpoint{2.392700in}{0.993563in}}%
\pgfpathlineto{\pgfqpoint{2.393760in}{0.992718in}}%
\pgfpathlineto{\pgfqpoint{2.395785in}{0.990579in}}%
\pgfpathlineto{\pgfqpoint{2.396055in}{0.991362in}}%
\pgfpathlineto{\pgfqpoint{2.397661in}{0.993553in}}%
\pgfpathlineto{\pgfqpoint{2.398751in}{0.992670in}}%
\pgfpathlineto{\pgfqpoint{2.400743in}{0.990587in}}%
\pgfpathlineto{\pgfqpoint{2.401007in}{0.991371in}}%
\pgfpathlineto{\pgfqpoint{2.402479in}{0.993505in}}%
\pgfpathlineto{\pgfqpoint{2.403468in}{0.992859in}}%
\pgfpathlineto{\pgfqpoint{2.405681in}{0.990525in}}%
\pgfpathlineto{\pgfqpoint{2.405950in}{0.991325in}}%
\pgfpathlineto{\pgfqpoint{2.407555in}{0.993498in}}%
\pgfpathlineto{\pgfqpoint{2.408656in}{0.992599in}}%
\pgfpathlineto{\pgfqpoint{2.410620in}{0.990512in}}%
\pgfpathlineto{\pgfqpoint{2.410907in}{0.991313in}}%
\pgfpathlineto{\pgfqpoint{2.412626in}{0.993408in}}%
\pgfpathlineto{\pgfqpoint{2.413924in}{0.992182in}}%
\pgfpathlineto{\pgfqpoint{2.415558in}{0.990455in}}%
\pgfpathlineto{\pgfqpoint{2.415833in}{0.991248in}}%
\pgfpathlineto{\pgfqpoint{2.417326in}{0.993433in}}%
\pgfpathlineto{\pgfqpoint{2.418315in}{0.992784in}}%
\pgfpathlineto{\pgfqpoint{2.420515in}{0.990464in}}%
\pgfpathlineto{\pgfqpoint{2.420788in}{0.991254in}}%
\pgfpathlineto{\pgfqpoint{2.422280in}{0.993406in}}%
\pgfpathlineto{\pgfqpoint{2.423269in}{0.992742in}}%
\pgfpathlineto{\pgfqpoint{2.425470in}{0.990435in}}%
\pgfpathlineto{\pgfqpoint{2.425732in}{0.991217in}}%
\pgfpathlineto{\pgfqpoint{2.427215in}{0.993382in}}%
\pgfpathlineto{\pgfqpoint{2.428204in}{0.992737in}}%
\pgfpathlineto{\pgfqpoint{2.430403in}{0.990410in}}%
\pgfpathlineto{\pgfqpoint{2.430679in}{0.991210in}}%
\pgfpathlineto{\pgfqpoint{2.432177in}{0.993368in}}%
\pgfpathlineto{\pgfqpoint{2.433166in}{0.992701in}}%
\pgfpathlineto{\pgfqpoint{2.435349in}{0.990388in}}%
\pgfpathlineto{\pgfqpoint{2.435629in}{0.991189in}}%
\pgfpathlineto{\pgfqpoint{2.437379in}{0.993287in}}%
\pgfpathlineto{\pgfqpoint{2.438709in}{0.992006in}}%
\pgfpathlineto{\pgfqpoint{2.440289in}{0.990341in}}%
\pgfpathlineto{\pgfqpoint{2.440558in}{0.991127in}}%
\pgfpathlineto{\pgfqpoint{2.442127in}{0.993330in}}%
\pgfpathlineto{\pgfqpoint{2.443185in}{0.992526in}}%
\pgfpathlineto{\pgfqpoint{2.445232in}{0.990353in}}%
\pgfpathlineto{\pgfqpoint{2.445513in}{0.991157in}}%
\pgfpathlineto{\pgfqpoint{2.447044in}{0.993314in}}%
\pgfpathlineto{\pgfqpoint{2.448101in}{0.992534in}}%
\pgfpathlineto{\pgfqpoint{2.450186in}{0.990323in}}%
\pgfpathlineto{\pgfqpoint{2.450477in}{0.991115in}}%
\pgfpathlineto{\pgfqpoint{2.451983in}{0.993238in}}%
\pgfpathlineto{\pgfqpoint{2.452935in}{0.992585in}}%
\pgfpathlineto{\pgfqpoint{2.455123in}{0.990276in}}%
\pgfpathlineto{\pgfqpoint{2.455393in}{0.991068in}}%
\pgfpathlineto{\pgfqpoint{2.456924in}{0.993260in}}%
\pgfpathlineto{\pgfqpoint{2.457950in}{0.992529in}}%
\pgfpathlineto{\pgfqpoint{2.460078in}{0.990279in}}%
\pgfpathlineto{\pgfqpoint{2.460365in}{0.991068in}}%
\pgfpathlineto{\pgfqpoint{2.462051in}{0.993161in}}%
\pgfpathlineto{\pgfqpoint{2.463310in}{0.992000in}}%
\pgfpathlineto{\pgfqpoint{2.465023in}{0.990211in}}%
\pgfpathlineto{\pgfqpoint{2.465293in}{0.991005in}}%
\pgfpathlineto{\pgfqpoint{2.466993in}{0.993166in}}%
\pgfpathlineto{\pgfqpoint{2.468173in}{0.992098in}}%
\pgfpathlineto{\pgfqpoint{2.469967in}{0.990209in}}%
\pgfpathlineto{\pgfqpoint{2.470225in}{0.990966in}}%
\pgfpathlineto{\pgfqpoint{2.471878in}{0.993170in}}%
\pgfpathlineto{\pgfqpoint{2.473017in}{0.992197in}}%
\pgfpathlineto{\pgfqpoint{2.474911in}{0.990198in}}%
\pgfpathlineto{\pgfqpoint{2.475185in}{0.991002in}}%
\pgfpathlineto{\pgfqpoint{2.476883in}{0.993158in}}%
\pgfpathlineto{\pgfqpoint{2.478138in}{0.992010in}}%
\pgfpathlineto{\pgfqpoint{2.479856in}{0.990190in}}%
\pgfpathlineto{\pgfqpoint{2.480139in}{0.990988in}}%
\pgfpathlineto{\pgfqpoint{2.481787in}{0.993114in}}%
\pgfpathlineto{\pgfqpoint{2.482901in}{0.992147in}}%
\pgfpathlineto{\pgfqpoint{2.484802in}{0.990139in}}%
\pgfpathlineto{\pgfqpoint{2.485080in}{0.990940in}}%
\pgfpathlineto{\pgfqpoint{2.486810in}{0.993077in}}%
\pgfpathlineto{\pgfqpoint{2.488079in}{0.991885in}}%
\pgfpathlineto{\pgfqpoint{2.489749in}{0.990129in}}%
\pgfpathlineto{\pgfqpoint{2.490020in}{0.990929in}}%
\pgfpathlineto{\pgfqpoint{2.491496in}{0.993100in}}%
\pgfpathlineto{\pgfqpoint{2.492485in}{0.992467in}}%
\pgfpathlineto{\pgfqpoint{2.494699in}{0.990135in}}%
\pgfpathlineto{\pgfqpoint{2.494968in}{0.990931in}}%
\pgfpathlineto{\pgfqpoint{2.496581in}{0.993099in}}%
\pgfpathlineto{\pgfqpoint{2.497675in}{0.992199in}}%
\pgfpathlineto{\pgfqpoint{2.499643in}{0.990132in}}%
\pgfpathlineto{\pgfqpoint{2.499911in}{0.990921in}}%
\pgfpathlineto{\pgfqpoint{2.501486in}{0.993080in}}%
\pgfpathlineto{\pgfqpoint{2.502573in}{0.992223in}}%
\pgfpathlineto{\pgfqpoint{2.504587in}{0.990101in}}%
\pgfpathlineto{\pgfqpoint{2.504864in}{0.990895in}}%
\pgfpathlineto{\pgfqpoint{2.506513in}{0.993035in}}%
\pgfpathlineto{\pgfqpoint{2.507690in}{0.992006in}}%
\pgfpathlineto{\pgfqpoint{2.509522in}{0.990068in}}%
\pgfpathlineto{\pgfqpoint{2.509800in}{0.990865in}}%
\pgfpathlineto{\pgfqpoint{2.511302in}{0.993030in}}%
\pgfpathlineto{\pgfqpoint{2.512291in}{0.992361in}}%
\pgfpathlineto{\pgfqpoint{2.514475in}{0.990053in}}%
\pgfpathlineto{\pgfqpoint{2.514746in}{0.990846in}}%
\pgfpathlineto{\pgfqpoint{2.516264in}{0.993028in}}%
\pgfpathlineto{\pgfqpoint{2.517293in}{0.992305in}}%
\pgfpathlineto{\pgfqpoint{2.519423in}{0.990047in}}%
\pgfpathlineto{\pgfqpoint{2.519703in}{0.990807in}}%
\pgfpathlineto{\pgfqpoint{2.521372in}{0.992924in}}%
\pgfpathlineto{\pgfqpoint{2.522600in}{0.991816in}}%
\pgfpathlineto{\pgfqpoint{2.524358in}{0.989966in}}%
\pgfpathlineto{\pgfqpoint{2.524627in}{0.990764in}}%
\pgfpathlineto{\pgfqpoint{2.526234in}{0.992984in}}%
\pgfpathlineto{\pgfqpoint{2.527288in}{0.992152in}}%
\pgfpathlineto{\pgfqpoint{2.529312in}{0.990013in}}%
\pgfpathlineto{\pgfqpoint{2.529587in}{0.990807in}}%
\pgfpathlineto{\pgfqpoint{2.531188in}{0.992961in}}%
\pgfpathlineto{\pgfqpoint{2.532288in}{0.992063in}}%
\pgfpathlineto{\pgfqpoint{2.534258in}{0.989986in}}%
\pgfpathlineto{\pgfqpoint{2.534505in}{0.990711in}}%
\pgfpathlineto{\pgfqpoint{2.536155in}{0.992945in}}%
\pgfpathlineto{\pgfqpoint{2.537275in}{0.992006in}}%
\pgfpathlineto{\pgfqpoint{2.539199in}{0.989974in}}%
\pgfpathlineto{\pgfqpoint{2.539475in}{0.990779in}}%
\pgfpathlineto{\pgfqpoint{2.541057in}{0.992949in}}%
\pgfpathlineto{\pgfqpoint{2.542142in}{0.992091in}}%
\pgfpathlineto{\pgfqpoint{2.544151in}{0.989979in}}%
\pgfpathlineto{\pgfqpoint{2.544423in}{0.990776in}}%
\pgfpathlineto{\pgfqpoint{2.546054in}{0.992920in}}%
\pgfpathlineto{\pgfqpoint{2.547226in}{0.991914in}}%
\pgfpathlineto{\pgfqpoint{2.549089in}{0.989944in}}%
\pgfpathlineto{\pgfqpoint{2.549347in}{0.990698in}}%
\pgfpathlineto{\pgfqpoint{2.550979in}{0.992935in}}%
\pgfpathlineto{\pgfqpoint{2.552052in}{0.992064in}}%
\pgfpathlineto{\pgfqpoint{2.554038in}{0.989966in}}%
\pgfpathlineto{\pgfqpoint{2.554312in}{0.990767in}}%
\pgfpathlineto{\pgfqpoint{2.555841in}{0.992928in}}%
\pgfpathlineto{\pgfqpoint{2.556863in}{0.992195in}}%
\pgfpathlineto{\pgfqpoint{2.558980in}{0.989952in}}%
\pgfpathlineto{\pgfqpoint{2.559256in}{0.990748in}}%
\pgfpathlineto{\pgfqpoint{2.560808in}{0.992911in}}%
\pgfpathlineto{\pgfqpoint{2.561873in}{0.992103in}}%
\pgfpathlineto{\pgfqpoint{2.563933in}{0.989934in}}%
\pgfpathlineto{\pgfqpoint{2.564208in}{0.990728in}}%
\pgfpathlineto{\pgfqpoint{2.565675in}{0.992864in}}%
\pgfpathlineto{\pgfqpoint{2.566664in}{0.992227in}}%
\pgfpathlineto{\pgfqpoint{2.568871in}{0.989890in}}%
\pgfpathlineto{\pgfqpoint{2.569146in}{0.990697in}}%
\pgfpathlineto{\pgfqpoint{2.570802in}{0.992881in}}%
\pgfpathlineto{\pgfqpoint{2.571903in}{0.991943in}}%
\pgfpathlineto{\pgfqpoint{2.573816in}{0.989918in}}%
\pgfpathlineto{\pgfqpoint{2.574094in}{0.990718in}}%
\pgfpathlineto{\pgfqpoint{2.575628in}{0.992864in}}%
\pgfpathlineto{\pgfqpoint{2.576687in}{0.992076in}}%
\pgfpathlineto{\pgfqpoint{2.578762in}{0.989880in}}%
\pgfpathlineto{\pgfqpoint{2.579038in}{0.990683in}}%
\pgfpathlineto{\pgfqpoint{2.580659in}{0.992872in}}%
\pgfpathlineto{\pgfqpoint{2.581730in}{0.991998in}}%
\pgfpathlineto{\pgfqpoint{2.583722in}{0.989907in}}%
\pgfpathlineto{\pgfqpoint{2.583991in}{0.990701in}}%
\pgfpathlineto{\pgfqpoint{2.585634in}{0.992829in}}%
\pgfpathlineto{\pgfqpoint{2.586785in}{0.991829in}}%
\pgfpathlineto{\pgfqpoint{2.588659in}{0.989841in}}%
\pgfpathlineto{\pgfqpoint{2.588951in}{0.990614in}}%
\pgfpathlineto{\pgfqpoint{2.590643in}{0.992703in}}%
\pgfpathlineto{\pgfqpoint{2.591982in}{0.991452in}}%
\pgfpathlineto{\pgfqpoint{2.593591in}{0.989753in}}%
\pgfpathlineto{\pgfqpoint{2.593869in}{0.990559in}}%
\pgfpathlineto{\pgfqpoint{2.595464in}{0.992751in}}%
\pgfpathlineto{\pgfqpoint{2.596508in}{0.991932in}}%
\pgfpathlineto{\pgfqpoint{2.598561in}{0.989781in}}%
\pgfpathlineto{\pgfqpoint{2.598827in}{0.990574in}}%
\pgfpathlineto{\pgfqpoint{2.600556in}{0.992703in}}%
\pgfpathlineto{\pgfqpoint{2.601930in}{0.991401in}}%
\pgfpathlineto{\pgfqpoint{2.603489in}{0.989755in}}%
\pgfpathlineto{\pgfqpoint{2.603768in}{0.990556in}}%
\pgfpathlineto{\pgfqpoint{2.605337in}{0.992708in}}%
\pgfpathlineto{\pgfqpoint{2.606407in}{0.991873in}}%
\pgfpathlineto{\pgfqpoint{2.608435in}{0.989725in}}%
\pgfpathlineto{\pgfqpoint{2.608711in}{0.990528in}}%
\pgfpathlineto{\pgfqpoint{2.610336in}{0.992699in}}%
\pgfpathlineto{\pgfqpoint{2.611420in}{0.991802in}}%
\pgfpathlineto{\pgfqpoint{2.613384in}{0.989727in}}%
\pgfpathlineto{\pgfqpoint{2.613654in}{0.990520in}}%
\pgfpathlineto{\pgfqpoint{2.615134in}{0.992702in}}%
\pgfpathlineto{\pgfqpoint{2.616118in}{0.992078in}}%
\pgfpathlineto{\pgfqpoint{2.618333in}{0.989743in}}%
\pgfpathlineto{\pgfqpoint{2.618607in}{0.990531in}}%
\pgfpathlineto{\pgfqpoint{2.620305in}{0.992659in}}%
\pgfpathlineto{\pgfqpoint{2.621591in}{0.991472in}}%
\pgfpathlineto{\pgfqpoint{2.623271in}{0.989703in}}%
\pgfpathlineto{\pgfqpoint{2.623546in}{0.990504in}}%
\pgfpathlineto{\pgfqpoint{2.625132in}{0.992683in}}%
\pgfpathlineto{\pgfqpoint{2.626201in}{0.991842in}}%
\pgfpathlineto{\pgfqpoint{2.628216in}{0.989710in}}%
\pgfpathlineto{\pgfqpoint{2.628490in}{0.990502in}}%
\pgfpathlineto{\pgfqpoint{2.630111in}{0.992679in}}%
\pgfpathlineto{\pgfqpoint{2.631248in}{0.991730in}}%
\pgfpathlineto{\pgfqpoint{2.633162in}{0.989707in}}%
\pgfpathlineto{\pgfqpoint{2.633436in}{0.990500in}}%
\pgfpathlineto{\pgfqpoint{2.634937in}{0.992677in}}%
\pgfpathlineto{\pgfqpoint{2.635927in}{0.992016in}}%
\pgfpathlineto{\pgfqpoint{2.638111in}{0.989707in}}%
\pgfpathlineto{\pgfqpoint{2.638384in}{0.990510in}}%
\pgfpathlineto{\pgfqpoint{2.639855in}{0.992674in}}%
\pgfpathlineto{\pgfqpoint{2.640827in}{0.992063in}}%
\pgfpathlineto{\pgfqpoint{2.643056in}{0.989696in}}%
\pgfpathlineto{\pgfqpoint{2.643341in}{0.990493in}}%
\pgfpathlineto{\pgfqpoint{2.645072in}{0.992610in}}%
\pgfpathlineto{\pgfqpoint{2.646294in}{0.991464in}}%
\pgfpathlineto{\pgfqpoint{2.647995in}{0.989667in}}%
\pgfpathlineto{\pgfqpoint{2.648272in}{0.990476in}}%
\pgfpathlineto{\pgfqpoint{2.649891in}{0.992655in}}%
\pgfpathlineto{\pgfqpoint{2.650982in}{0.991757in}}%
\pgfpathlineto{\pgfqpoint{2.652949in}{0.989680in}}%
\pgfpathlineto{\pgfqpoint{2.653224in}{0.990488in}}%
\pgfpathlineto{\pgfqpoint{2.654834in}{0.992653in}}%
\pgfpathlineto{\pgfqpoint{2.655913in}{0.991773in}}%
\pgfpathlineto{\pgfqpoint{2.657893in}{0.989681in}}%
\pgfpathlineto{\pgfqpoint{2.658168in}{0.990486in}}%
\pgfpathlineto{\pgfqpoint{2.659671in}{0.992656in}}%
\pgfpathlineto{\pgfqpoint{2.660660in}{0.991989in}}%
\pgfpathlineto{\pgfqpoint{2.662838in}{0.989677in}}%
\pgfpathlineto{\pgfqpoint{2.663124in}{0.990478in}}%
\pgfpathlineto{\pgfqpoint{2.664860in}{0.992588in}}%
\pgfpathlineto{\pgfqpoint{2.666202in}{0.991309in}}%
\pgfpathlineto{\pgfqpoint{2.667782in}{0.989649in}}%
\pgfpathlineto{\pgfqpoint{2.668049in}{0.990429in}}%
\pgfpathlineto{\pgfqpoint{2.669668in}{0.992633in}}%
\pgfpathlineto{\pgfqpoint{2.670765in}{0.991742in}}%
\pgfpathlineto{\pgfqpoint{2.672727in}{0.989669in}}%
\pgfpathlineto{\pgfqpoint{2.673002in}{0.990475in}}%
\pgfpathlineto{\pgfqpoint{2.674687in}{0.992621in}}%
\pgfpathlineto{\pgfqpoint{2.675917in}{0.991508in}}%
\pgfpathlineto{\pgfqpoint{2.677675in}{0.989654in}}%
\pgfpathlineto{\pgfqpoint{2.677948in}{0.990430in}}%
\pgfpathlineto{\pgfqpoint{2.679586in}{0.992617in}}%
\pgfpathlineto{\pgfqpoint{2.680706in}{0.991674in}}%
\pgfpathlineto{\pgfqpoint{2.682618in}{0.989650in}}%
\pgfpathlineto{\pgfqpoint{2.682902in}{0.990436in}}%
\pgfpathlineto{\pgfqpoint{2.684562in}{0.992561in}}%
\pgfpathlineto{\pgfqpoint{2.685725in}{0.991533in}}%
\pgfpathlineto{\pgfqpoint{2.687562in}{0.989598in}}%
\pgfpathlineto{\pgfqpoint{2.687826in}{0.990357in}}%
\pgfpathlineto{\pgfqpoint{2.689439in}{0.992577in}}%
\pgfpathlineto{\pgfqpoint{2.690536in}{0.991696in}}%
\pgfpathlineto{\pgfqpoint{2.692505in}{0.989604in}}%
\pgfpathlineto{\pgfqpoint{2.692787in}{0.990403in}}%
\pgfpathlineto{\pgfqpoint{2.694265in}{0.992535in}}%
\pgfpathlineto{\pgfqpoint{2.695254in}{0.991884in}}%
\pgfpathlineto{\pgfqpoint{2.697458in}{0.989548in}}%
\pgfpathlineto{\pgfqpoint{2.697740in}{0.990352in}}%
\pgfpathlineto{\pgfqpoint{2.699487in}{0.992468in}}%
\pgfpathlineto{\pgfqpoint{2.700814in}{0.991200in}}%
\pgfpathlineto{\pgfqpoint{2.702402in}{0.989534in}}%
\pgfpathlineto{\pgfqpoint{2.702671in}{0.990328in}}%
\pgfpathlineto{\pgfqpoint{2.704178in}{0.992518in}}%
\pgfpathlineto{\pgfqpoint{2.705213in}{0.991806in}}%
\pgfpathlineto{\pgfqpoint{2.707344in}{0.989550in}}%
\pgfpathlineto{\pgfqpoint{2.707618in}{0.990343in}}%
\pgfpathlineto{\pgfqpoint{2.709207in}{0.992514in}}%
\pgfpathlineto{\pgfqpoint{2.710293in}{0.991648in}}%
\pgfpathlineto{\pgfqpoint{2.712294in}{0.989538in}}%
\pgfpathlineto{\pgfqpoint{2.712526in}{0.990218in}}%
\pgfpathlineto{\pgfqpoint{2.714137in}{0.992515in}}%
\pgfpathlineto{\pgfqpoint{2.715203in}{0.991696in}}%
\pgfpathlineto{\pgfqpoint{2.717238in}{0.989552in}}%
\pgfpathlineto{\pgfqpoint{2.717511in}{0.990354in}}%
\pgfpathlineto{\pgfqpoint{2.718976in}{0.992503in}}%
\pgfpathlineto{\pgfqpoint{2.719965in}{0.991877in}}%
\pgfpathlineto{\pgfqpoint{2.722184in}{0.989529in}}%
\pgfpathlineto{\pgfqpoint{2.722455in}{0.990323in}}%
\pgfpathlineto{\pgfqpoint{2.724006in}{0.992504in}}%
\pgfpathlineto{\pgfqpoint{2.725068in}{0.991707in}}%
\pgfpathlineto{\pgfqpoint{2.727130in}{0.989529in}}%
\pgfpathlineto{\pgfqpoint{2.727405in}{0.990334in}}%
\pgfpathlineto{\pgfqpoint{2.728864in}{0.992500in}}%
\pgfpathlineto{\pgfqpoint{2.729854in}{0.991890in}}%
\pgfpathlineto{\pgfqpoint{2.732074in}{0.989530in}}%
\pgfpathlineto{\pgfqpoint{2.732361in}{0.990308in}}%
\pgfpathlineto{\pgfqpoint{2.734015in}{0.992423in}}%
\pgfpathlineto{\pgfqpoint{2.735186in}{0.991390in}}%
\pgfpathlineto{\pgfqpoint{2.737017in}{0.989474in}}%
\pgfpathlineto{\pgfqpoint{2.737286in}{0.990273in}}%
\pgfpathlineto{\pgfqpoint{2.738794in}{0.992471in}}%
\pgfpathlineto{\pgfqpoint{2.739783in}{0.991811in}}%
\pgfpathlineto{\pgfqpoint{2.741967in}{0.989500in}}%
\pgfpathlineto{\pgfqpoint{2.742249in}{0.990302in}}%
\pgfpathlineto{\pgfqpoint{2.743707in}{0.992425in}}%
\pgfpathlineto{\pgfqpoint{2.744662in}{0.991831in}}%
\pgfpathlineto{\pgfqpoint{2.746911in}{0.989450in}}%
\pgfpathlineto{\pgfqpoint{2.747192in}{0.990252in}}%
\pgfpathlineto{\pgfqpoint{2.748691in}{0.992415in}}%
\pgfpathlineto{\pgfqpoint{2.749664in}{0.991767in}}%
\pgfpathlineto{\pgfqpoint{2.751857in}{0.989454in}}%
\pgfpathlineto{\pgfqpoint{2.752128in}{0.990252in}}%
\pgfpathlineto{\pgfqpoint{2.753705in}{0.992427in}}%
\pgfpathlineto{\pgfqpoint{2.754754in}{0.991618in}}%
\pgfpathlineto{\pgfqpoint{2.756798in}{0.989457in}}%
\pgfpathlineto{\pgfqpoint{2.757071in}{0.990247in}}%
\pgfpathlineto{\pgfqpoint{2.758602in}{0.992427in}}%
\pgfpathlineto{\pgfqpoint{2.759661in}{0.991659in}}%
\pgfpathlineto{\pgfqpoint{2.761748in}{0.989456in}}%
\pgfpathlineto{\pgfqpoint{2.762020in}{0.990258in}}%
\pgfpathlineto{\pgfqpoint{2.763519in}{0.992432in}}%
\pgfpathlineto{\pgfqpoint{2.764490in}{0.991794in}}%
\pgfpathlineto{\pgfqpoint{2.766693in}{0.989457in}}%
\pgfpathlineto{\pgfqpoint{2.766979in}{0.990251in}}%
\pgfpathlineto{\pgfqpoint{2.768433in}{0.992379in}}%
\pgfpathlineto{\pgfqpoint{2.769396in}{0.991784in}}%
\pgfpathlineto{\pgfqpoint{2.771638in}{0.989423in}}%
\pgfpathlineto{\pgfqpoint{2.771858in}{0.990082in}}%
\pgfpathlineto{\pgfqpoint{2.773449in}{0.992420in}}%
\pgfpathlineto{\pgfqpoint{2.774508in}{0.991645in}}%
\pgfpathlineto{\pgfqpoint{2.776584in}{0.989451in}}%
\pgfpathlineto{\pgfqpoint{2.776858in}{0.990255in}}%
\pgfpathlineto{\pgfqpoint{2.778375in}{0.992421in}}%
\pgfpathlineto{\pgfqpoint{2.779403in}{0.991697in}}%
\pgfpathlineto{\pgfqpoint{2.781538in}{0.989446in}}%
\pgfpathlineto{\pgfqpoint{2.781808in}{0.990239in}}%
\pgfpathlineto{\pgfqpoint{2.783345in}{0.992398in}}%
\pgfpathlineto{\pgfqpoint{2.784389in}{0.991627in}}%
\pgfpathlineto{\pgfqpoint{2.786476in}{0.989417in}}%
\pgfpathlineto{\pgfqpoint{2.786759in}{0.990222in}}%
\pgfpathlineto{\pgfqpoint{2.788500in}{0.992332in}}%
\pgfpathlineto{\pgfqpoint{2.789788in}{0.991107in}}%
\pgfpathlineto{\pgfqpoint{2.791414in}{0.989382in}}%
\pgfpathlineto{\pgfqpoint{2.791697in}{0.990189in}}%
\pgfpathlineto{\pgfqpoint{2.793383in}{0.992341in}}%
\pgfpathlineto{\pgfqpoint{2.794567in}{0.991280in}}%
\pgfpathlineto{\pgfqpoint{2.796367in}{0.989378in}}%
\pgfpathlineto{\pgfqpoint{2.796646in}{0.990176in}}%
\pgfpathlineto{\pgfqpoint{2.798403in}{0.992291in}}%
\pgfpathlineto{\pgfqpoint{2.799729in}{0.991017in}}%
\pgfpathlineto{\pgfqpoint{2.801324in}{0.989362in}}%
\pgfpathlineto{\pgfqpoint{2.801587in}{0.990145in}}%
\pgfpathlineto{\pgfqpoint{2.803073in}{0.992314in}}%
\pgfpathlineto{\pgfqpoint{2.804062in}{0.991670in}}%
\pgfpathlineto{\pgfqpoint{2.806258in}{0.989346in}}%
\pgfpathlineto{\pgfqpoint{2.806537in}{0.990146in}}%
\pgfpathlineto{\pgfqpoint{2.808298in}{0.992238in}}%
\pgfpathlineto{\pgfqpoint{2.809611in}{0.990971in}}%
\pgfpathlineto{\pgfqpoint{2.811195in}{0.989298in}}%
\pgfpathlineto{\pgfqpoint{2.811476in}{0.990106in}}%
\pgfpathlineto{\pgfqpoint{2.812962in}{0.992290in}}%
\pgfpathlineto{\pgfqpoint{2.813951in}{0.991654in}}%
\pgfpathlineto{\pgfqpoint{2.816148in}{0.989336in}}%
\pgfpathlineto{\pgfqpoint{2.816382in}{0.990024in}}%
\pgfpathlineto{\pgfqpoint{2.817999in}{0.992299in}}%
\pgfpathlineto{\pgfqpoint{2.819055in}{0.991477in}}%
\pgfpathlineto{\pgfqpoint{2.821095in}{0.989317in}}%
\pgfpathlineto{\pgfqpoint{2.821383in}{0.990090in}}%
\pgfpathlineto{\pgfqpoint{2.823068in}{0.992189in}}%
\pgfpathlineto{\pgfqpoint{2.824354in}{0.991007in}}%
\pgfpathlineto{\pgfqpoint{2.826032in}{0.989234in}}%
\pgfpathlineto{\pgfqpoint{2.826311in}{0.990040in}}%
\pgfpathlineto{\pgfqpoint{2.827833in}{0.992212in}}%
\pgfpathlineto{\pgfqpoint{2.828859in}{0.991486in}}%
\pgfpathlineto{\pgfqpoint{2.830987in}{0.989244in}}%
\pgfpathlineto{\pgfqpoint{2.831265in}{0.990037in}}%
\pgfpathlineto{\pgfqpoint{2.832867in}{0.992176in}}%
\pgfpathlineto{\pgfqpoint{2.833957in}{0.991286in}}%
\pgfpathlineto{\pgfqpoint{2.835923in}{0.989200in}}%
\pgfpathlineto{\pgfqpoint{2.836204in}{0.989999in}}%
\pgfpathlineto{\pgfqpoint{2.837784in}{0.992165in}}%
\pgfpathlineto{\pgfqpoint{2.838860in}{0.991319in}}%
\pgfpathlineto{\pgfqpoint{2.840873in}{0.989183in}}%
\pgfpathlineto{\pgfqpoint{2.841164in}{0.989966in}}%
\pgfpathlineto{\pgfqpoint{2.842806in}{0.992074in}}%
\pgfpathlineto{\pgfqpoint{2.843923in}{0.991109in}}%
\pgfpathlineto{\pgfqpoint{2.845816in}{0.989108in}}%
\pgfpathlineto{\pgfqpoint{2.846088in}{0.989876in}}%
\pgfpathlineto{\pgfqpoint{2.847697in}{0.992055in}}%
\pgfpathlineto{\pgfqpoint{2.848779in}{0.991179in}}%
\pgfpathlineto{\pgfqpoint{2.850768in}{0.989087in}}%
\pgfpathlineto{\pgfqpoint{2.851040in}{0.989882in}}%
\pgfpathlineto{\pgfqpoint{2.852606in}{0.992060in}}%
\pgfpathlineto{\pgfqpoint{2.853640in}{0.991279in}}%
\pgfpathlineto{\pgfqpoint{2.855712in}{0.989095in}}%
\pgfpathlineto{\pgfqpoint{2.855946in}{0.989785in}}%
\pgfpathlineto{\pgfqpoint{2.857548in}{0.992074in}}%
\pgfpathlineto{\pgfqpoint{2.858597in}{0.991283in}}%
\pgfpathlineto{\pgfqpoint{2.860658in}{0.989106in}}%
\pgfpathlineto{\pgfqpoint{2.860936in}{0.989900in}}%
\pgfpathlineto{\pgfqpoint{2.862446in}{0.992037in}}%
\pgfpathlineto{\pgfqpoint{2.863412in}{0.991374in}}%
\pgfpathlineto{\pgfqpoint{2.865598in}{0.989065in}}%
\pgfpathlineto{\pgfqpoint{2.865862in}{0.989838in}}%
\pgfpathlineto{\pgfqpoint{2.867484in}{0.992064in}}%
\pgfpathlineto{\pgfqpoint{2.868542in}{0.991219in}}%
\pgfpathlineto{\pgfqpoint{2.870561in}{0.989107in}}%
\pgfpathlineto{\pgfqpoint{2.870826in}{0.989894in}}%
\pgfpathlineto{\pgfqpoint{2.872521in}{0.992018in}}%
\pgfpathlineto{\pgfqpoint{2.873786in}{0.990858in}}%
\pgfpathlineto{\pgfqpoint{2.875498in}{0.989068in}}%
\pgfpathlineto{\pgfqpoint{2.875766in}{0.989866in}}%
\pgfpathlineto{\pgfqpoint{2.877384in}{0.992040in}}%
\pgfpathlineto{\pgfqpoint{2.878496in}{0.991124in}}%
\pgfpathlineto{\pgfqpoint{2.880440in}{0.989075in}}%
\pgfpathlineto{\pgfqpoint{2.880713in}{0.989867in}}%
\pgfpathlineto{\pgfqpoint{2.882337in}{0.992027in}}%
\pgfpathlineto{\pgfqpoint{2.883411in}{0.991140in}}%
\pgfpathlineto{\pgfqpoint{2.885381in}{0.989062in}}%
\pgfpathlineto{\pgfqpoint{2.885622in}{0.989770in}}%
\pgfpathlineto{\pgfqpoint{2.887205in}{0.992039in}}%
\pgfpathlineto{\pgfqpoint{2.888252in}{0.991260in}}%
\pgfpathlineto{\pgfqpoint{2.890343in}{0.989071in}}%
\pgfpathlineto{\pgfqpoint{2.890609in}{0.989863in}}%
\pgfpathlineto{\pgfqpoint{2.892330in}{0.991995in}}%
\pgfpathlineto{\pgfqpoint{2.893655in}{0.990753in}}%
\pgfpathlineto{\pgfqpoint{2.895270in}{0.989048in}}%
\pgfpathlineto{\pgfqpoint{2.895508in}{0.989736in}}%
\pgfpathlineto{\pgfqpoint{2.897072in}{0.992011in}}%
\pgfpathlineto{\pgfqpoint{2.898097in}{0.991282in}}%
\pgfpathlineto{\pgfqpoint{2.900219in}{0.989031in}}%
\pgfpathlineto{\pgfqpoint{2.900506in}{0.989827in}}%
\pgfpathlineto{\pgfqpoint{2.902243in}{0.991923in}}%
\pgfpathlineto{\pgfqpoint{2.903553in}{0.990677in}}%
\pgfpathlineto{\pgfqpoint{2.905168in}{0.988987in}}%
\pgfpathlineto{\pgfqpoint{2.905441in}{0.989785in}}%
\pgfpathlineto{\pgfqpoint{2.907121in}{0.991950in}}%
\pgfpathlineto{\pgfqpoint{2.908231in}{0.990976in}}%
\pgfpathlineto{\pgfqpoint{2.910112in}{0.988998in}}%
\pgfpathlineto{\pgfqpoint{2.910367in}{0.989744in}}%
\pgfpathlineto{\pgfqpoint{2.912019in}{0.991967in}}%
\pgfpathlineto{\pgfqpoint{2.913183in}{0.990978in}}%
\pgfpathlineto{\pgfqpoint{2.915055in}{0.989004in}}%
\pgfpathlineto{\pgfqpoint{2.915318in}{0.989768in}}%
\pgfpathlineto{\pgfqpoint{2.916966in}{0.991980in}}%
\pgfpathlineto{\pgfqpoint{2.918118in}{0.991004in}}%
\pgfpathlineto{\pgfqpoint{2.920004in}{0.989020in}}%
\pgfpathlineto{\pgfqpoint{2.920279in}{0.989819in}}%
\pgfpathlineto{\pgfqpoint{2.921815in}{0.991970in}}%
\pgfpathlineto{\pgfqpoint{2.922815in}{0.991246in}}%
\pgfpathlineto{\pgfqpoint{2.924947in}{0.988985in}}%
\pgfpathlineto{\pgfqpoint{2.925240in}{0.989782in}}%
\pgfpathlineto{\pgfqpoint{2.926729in}{0.991923in}}%
\pgfpathlineto{\pgfqpoint{2.927759in}{0.991217in}}%
\pgfpathlineto{\pgfqpoint{2.929888in}{0.988964in}}%
\pgfpathlineto{\pgfqpoint{2.930165in}{0.989769in}}%
\pgfpathlineto{\pgfqpoint{2.931751in}{0.991932in}}%
\pgfpathlineto{\pgfqpoint{2.932810in}{0.991099in}}%
\pgfpathlineto{\pgfqpoint{2.934838in}{0.988950in}}%
\pgfpathlineto{\pgfqpoint{2.935121in}{0.989744in}}%
\pgfpathlineto{\pgfqpoint{2.936637in}{0.991900in}}%
\pgfpathlineto{\pgfqpoint{2.937637in}{0.991202in}}%
\pgfpathlineto{\pgfqpoint{2.939785in}{0.988946in}}%
\pgfpathlineto{\pgfqpoint{2.940051in}{0.989736in}}%
\pgfpathlineto{\pgfqpoint{2.941656in}{0.991930in}}%
\pgfpathlineto{\pgfqpoint{2.942720in}{0.991082in}}%
\pgfpathlineto{\pgfqpoint{2.944730in}{0.988964in}}%
\pgfpathlineto{\pgfqpoint{2.944989in}{0.989718in}}%
\pgfpathlineto{\pgfqpoint{2.946640in}{0.991930in}}%
\pgfpathlineto{\pgfqpoint{2.947793in}{0.990950in}}%
\pgfpathlineto{\pgfqpoint{2.949671in}{0.988970in}}%
\pgfpathlineto{\pgfqpoint{2.949940in}{0.989752in}}%
\pgfpathlineto{\pgfqpoint{2.951503in}{0.991940in}}%
\pgfpathlineto{\pgfqpoint{2.952551in}{0.991152in}}%
\pgfpathlineto{\pgfqpoint{2.954634in}{0.988975in}}%
\pgfpathlineto{\pgfqpoint{2.954900in}{0.989766in}}%
\pgfpathlineto{\pgfqpoint{2.956613in}{0.991902in}}%
\pgfpathlineto{\pgfqpoint{2.957925in}{0.990681in}}%
\pgfpathlineto{\pgfqpoint{2.959561in}{0.988956in}}%
\pgfpathlineto{\pgfqpoint{2.959838in}{0.989747in}}%
\pgfpathlineto{\pgfqpoint{2.961367in}{0.991912in}}%
\pgfpathlineto{\pgfqpoint{2.962390in}{0.991180in}}%
\pgfpathlineto{\pgfqpoint{2.964511in}{0.988938in}}%
\pgfpathlineto{\pgfqpoint{2.964795in}{0.989718in}}%
\pgfpathlineto{\pgfqpoint{2.966415in}{0.991853in}}%
\pgfpathlineto{\pgfqpoint{2.967544in}{0.990903in}}%
\pgfpathlineto{\pgfqpoint{2.969449in}{0.988898in}}%
\pgfpathlineto{\pgfqpoint{2.969724in}{0.989696in}}%
\pgfpathlineto{\pgfqpoint{2.971306in}{0.991886in}}%
\pgfpathlineto{\pgfqpoint{2.972390in}{0.991040in}}%
\pgfpathlineto{\pgfqpoint{2.974415in}{0.988922in}}%
\pgfpathlineto{\pgfqpoint{2.974683in}{0.989715in}}%
\pgfpathlineto{\pgfqpoint{2.976408in}{0.991828in}}%
\pgfpathlineto{\pgfqpoint{2.977778in}{0.990531in}}%
\pgfpathlineto{\pgfqpoint{2.979343in}{0.988878in}}%
\pgfpathlineto{\pgfqpoint{2.979599in}{0.989629in}}%
\pgfpathlineto{\pgfqpoint{2.981227in}{0.991877in}}%
\pgfpathlineto{\pgfqpoint{2.982304in}{0.991014in}}%
\pgfpathlineto{\pgfqpoint{2.984293in}{0.988916in}}%
\pgfpathlineto{\pgfqpoint{2.984571in}{0.989699in}}%
\pgfpathlineto{\pgfqpoint{2.986203in}{0.991849in}}%
\pgfpathlineto{\pgfqpoint{2.987356in}{0.990867in}}%
\pgfpathlineto{\pgfqpoint{2.989235in}{0.988896in}}%
\pgfpathlineto{\pgfqpoint{2.989501in}{0.989680in}}%
\pgfpathlineto{\pgfqpoint{2.991030in}{0.991886in}}%
\pgfpathlineto{\pgfqpoint{2.992057in}{0.991167in}}%
\pgfpathlineto{\pgfqpoint{2.994184in}{0.988922in}}%
\pgfpathlineto{\pgfqpoint{2.994454in}{0.989716in}}%
\pgfpathlineto{\pgfqpoint{2.995957in}{0.991891in}}%
\pgfpathlineto{\pgfqpoint{2.996946in}{0.991230in}}%
\pgfpathlineto{\pgfqpoint{2.999142in}{0.988925in}}%
\pgfpathlineto{\pgfqpoint{2.999408in}{0.989712in}}%
\pgfpathlineto{\pgfqpoint{3.001146in}{0.991840in}}%
\pgfpathlineto{\pgfqpoint{3.002512in}{0.990544in}}%
\pgfpathlineto{\pgfqpoint{3.004083in}{0.988897in}}%
\pgfpathlineto{\pgfqpoint{3.004355in}{0.989693in}}%
\pgfpathlineto{\pgfqpoint{3.005858in}{0.991836in}}%
\pgfpathlineto{\pgfqpoint{3.006889in}{0.991115in}}%
\pgfpathlineto{\pgfqpoint{3.009016in}{0.988867in}}%
\pgfpathlineto{\pgfqpoint{3.009291in}{0.989658in}}%
\pgfpathlineto{\pgfqpoint{3.010890in}{0.991835in}}%
\pgfpathlineto{\pgfqpoint{3.012006in}{0.990934in}}%
\pgfpathlineto{\pgfqpoint{3.013966in}{0.988862in}}%
\pgfpathlineto{\pgfqpoint{3.014250in}{0.989660in}}%
\pgfpathlineto{\pgfqpoint{3.016011in}{0.991742in}}%
\pgfpathlineto{\pgfqpoint{3.017453in}{0.990338in}}%
\pgfpathlineto{\pgfqpoint{3.018913in}{0.988812in}}%
\pgfpathlineto{\pgfqpoint{3.019184in}{0.989608in}}%
\pgfpathlineto{\pgfqpoint{3.020843in}{0.991786in}}%
\pgfpathlineto{\pgfqpoint{3.021941in}{0.990847in}}%
\pgfpathlineto{\pgfqpoint{3.023856in}{0.988820in}}%
\pgfpathlineto{\pgfqpoint{3.024139in}{0.989619in}}%
\pgfpathlineto{\pgfqpoint{3.025868in}{0.991739in}}%
\pgfpathlineto{\pgfqpoint{3.027106in}{0.990580in}}%
\pgfpathlineto{\pgfqpoint{3.028796in}{0.988789in}}%
\pgfpathlineto{\pgfqpoint{3.029081in}{0.989601in}}%
\pgfpathlineto{\pgfqpoint{3.030794in}{0.991747in}}%
\pgfpathlineto{\pgfqpoint{3.032095in}{0.990541in}}%
\pgfpathlineto{\pgfqpoint{3.033747in}{0.988801in}}%
\pgfpathlineto{\pgfqpoint{3.034024in}{0.989608in}}%
\pgfpathlineto{\pgfqpoint{3.035610in}{0.991761in}}%
\pgfpathlineto{\pgfqpoint{3.036671in}{0.990923in}}%
\pgfpathlineto{\pgfqpoint{3.038695in}{0.988786in}}%
\pgfpathlineto{\pgfqpoint{3.038973in}{0.989582in}}%
\pgfpathlineto{\pgfqpoint{3.040602in}{0.991737in}}%
\pgfpathlineto{\pgfqpoint{3.041631in}{0.990890in}}%
\pgfpathlineto{\pgfqpoint{3.043640in}{0.988777in}}%
\pgfpathlineto{\pgfqpoint{3.043911in}{0.989581in}}%
\pgfpathlineto{\pgfqpoint{3.045406in}{0.991763in}}%
\pgfpathlineto{\pgfqpoint{3.046395in}{0.991116in}}%
\pgfpathlineto{\pgfqpoint{3.048582in}{0.988802in}}%
\pgfpathlineto{\pgfqpoint{3.048843in}{0.989559in}}%
\pgfpathlineto{\pgfqpoint{3.050486in}{0.991763in}}%
\pgfpathlineto{\pgfqpoint{3.051646in}{0.990781in}}%
\pgfpathlineto{\pgfqpoint{3.053525in}{0.988798in}}%
\pgfpathlineto{\pgfqpoint{3.053798in}{0.989588in}}%
\pgfpathlineto{\pgfqpoint{3.055426in}{0.991771in}}%
\pgfpathlineto{\pgfqpoint{3.056537in}{0.990848in}}%
\pgfpathlineto{\pgfqpoint{3.058488in}{0.988811in}}%
\pgfpathlineto{\pgfqpoint{3.058756in}{0.989607in}}%
\pgfpathlineto{\pgfqpoint{3.060484in}{0.991732in}}%
\pgfpathlineto{\pgfqpoint{3.061855in}{0.990435in}}%
\pgfpathlineto{\pgfqpoint{3.063416in}{0.988791in}}%
\pgfpathlineto{\pgfqpoint{3.063697in}{0.989598in}}%
\pgfpathlineto{\pgfqpoint{3.065277in}{0.991746in}}%
\pgfpathlineto{\pgfqpoint{3.066349in}{0.990899in}}%
\pgfpathlineto{\pgfqpoint{3.068367in}{0.988775in}}%
\pgfpathlineto{\pgfqpoint{3.068600in}{0.989466in}}%
\pgfpathlineto{\pgfqpoint{3.070193in}{0.991767in}}%
\pgfpathlineto{\pgfqpoint{3.071228in}{0.991008in}}%
\pgfpathlineto{\pgfqpoint{3.073308in}{0.988813in}}%
\pgfpathlineto{\pgfqpoint{3.073567in}{0.989566in}}%
\pgfpathlineto{\pgfqpoint{3.075140in}{0.991760in}}%
\pgfpathlineto{\pgfqpoint{3.076194in}{0.990960in}}%
\pgfpathlineto{\pgfqpoint{3.078258in}{0.988786in}}%
\pgfpathlineto{\pgfqpoint{3.078525in}{0.989561in}}%
\pgfpathlineto{\pgfqpoint{3.080124in}{0.991762in}}%
\pgfpathlineto{\pgfqpoint{3.081178in}{0.990935in}}%
\pgfpathlineto{\pgfqpoint{3.083200in}{0.988802in}}%
\pgfpathlineto{\pgfqpoint{3.083463in}{0.989564in}}%
\pgfpathlineto{\pgfqpoint{3.085070in}{0.991775in}}%
\pgfpathlineto{\pgfqpoint{3.086162in}{0.990902in}}%
\pgfpathlineto{\pgfqpoint{3.088144in}{0.988810in}}%
\pgfpathlineto{\pgfqpoint{3.088417in}{0.989611in}}%
\pgfpathlineto{\pgfqpoint{3.090084in}{0.991799in}}%
\pgfpathlineto{\pgfqpoint{3.091276in}{0.990754in}}%
\pgfpathlineto{\pgfqpoint{3.093094in}{0.988842in}}%
\pgfpathlineto{\pgfqpoint{3.093364in}{0.989612in}}%
\pgfpathlineto{\pgfqpoint{3.094995in}{0.991791in}}%
\pgfpathlineto{\pgfqpoint{3.096113in}{0.990860in}}%
\pgfpathlineto{\pgfqpoint{3.098036in}{0.988825in}}%
\pgfpathlineto{\pgfqpoint{3.098312in}{0.989573in}}%
\pgfpathlineto{\pgfqpoint{3.099978in}{0.991710in}}%
\pgfpathlineto{\pgfqpoint{3.101173in}{0.990650in}}%
\pgfpathlineto{\pgfqpoint{3.102984in}{0.988751in}}%
\pgfpathlineto{\pgfqpoint{3.103253in}{0.989545in}}%
\pgfpathlineto{\pgfqpoint{3.104881in}{0.991743in}}%
\pgfpathlineto{\pgfqpoint{3.105926in}{0.990897in}}%
\pgfpathlineto{\pgfqpoint{3.107929in}{0.988775in}}%
\pgfpathlineto{\pgfqpoint{3.108215in}{0.989554in}}%
\pgfpathlineto{\pgfqpoint{3.109865in}{0.991677in}}%
\pgfpathlineto{\pgfqpoint{3.111038in}{0.990650in}}%
\pgfpathlineto{\pgfqpoint{3.112868in}{0.988726in}}%
\pgfpathlineto{\pgfqpoint{3.113143in}{0.989530in}}%
\pgfpathlineto{\pgfqpoint{3.114761in}{0.991722in}}%
\pgfpathlineto{\pgfqpoint{3.115855in}{0.990832in}}%
\pgfpathlineto{\pgfqpoint{3.117816in}{0.988760in}}%
\pgfpathlineto{\pgfqpoint{3.118086in}{0.989546in}}%
\pgfpathlineto{\pgfqpoint{3.119697in}{0.991733in}}%
\pgfpathlineto{\pgfqpoint{3.120753in}{0.990885in}}%
\pgfpathlineto{\pgfqpoint{3.122770in}{0.988772in}}%
\pgfpathlineto{\pgfqpoint{3.123037in}{0.989556in}}%
\pgfpathlineto{\pgfqpoint{3.124696in}{0.991730in}}%
\pgfpathlineto{\pgfqpoint{3.125884in}{0.990694in}}%
\pgfpathlineto{\pgfqpoint{3.127707in}{0.988770in}}%
\pgfpathlineto{\pgfqpoint{3.127963in}{0.989517in}}%
\pgfpathlineto{\pgfqpoint{3.129568in}{0.991752in}}%
\pgfpathlineto{\pgfqpoint{3.130603in}{0.990954in}}%
\pgfpathlineto{\pgfqpoint{3.132658in}{0.988790in}}%
\pgfpathlineto{\pgfqpoint{3.132924in}{0.989558in}}%
\pgfpathlineto{\pgfqpoint{3.134529in}{0.991743in}}%
\pgfpathlineto{\pgfqpoint{3.135609in}{0.990877in}}%
\pgfpathlineto{\pgfqpoint{3.137602in}{0.988774in}}%
\pgfpathlineto{\pgfqpoint{3.137883in}{0.989577in}}%
\pgfpathlineto{\pgfqpoint{3.139419in}{0.991738in}}%
\pgfpathlineto{\pgfqpoint{3.140415in}{0.991023in}}%
\pgfpathlineto{\pgfqpoint{3.142549in}{0.988768in}}%
\pgfpathlineto{\pgfqpoint{3.142836in}{0.989564in}}%
\pgfpathlineto{\pgfqpoint{3.144459in}{0.991681in}}%
\pgfpathlineto{\pgfqpoint{3.145608in}{0.990701in}}%
\pgfpathlineto{\pgfqpoint{3.147489in}{0.988732in}}%
\pgfpathlineto{\pgfqpoint{3.147758in}{0.989529in}}%
\pgfpathlineto{\pgfqpoint{3.149301in}{0.991732in}}%
\pgfpathlineto{\pgfqpoint{3.150365in}{0.990957in}}%
\pgfpathlineto{\pgfqpoint{3.152439in}{0.988771in}}%
\pgfpathlineto{\pgfqpoint{3.152670in}{0.989445in}}%
\pgfpathlineto{\pgfqpoint{3.154280in}{0.991733in}}%
\pgfpathlineto{\pgfqpoint{3.155350in}{0.990911in}}%
\pgfpathlineto{\pgfqpoint{3.157383in}{0.988762in}}%
\pgfpathlineto{\pgfqpoint{3.157665in}{0.989565in}}%
\pgfpathlineto{\pgfqpoint{3.159362in}{0.991681in}}%
\pgfpathlineto{\pgfqpoint{3.160591in}{0.990556in}}%
\pgfpathlineto{\pgfqpoint{3.162330in}{0.988730in}}%
\pgfpathlineto{\pgfqpoint{3.162599in}{0.989527in}}%
\pgfpathlineto{\pgfqpoint{3.164264in}{0.991722in}}%
\pgfpathlineto{\pgfqpoint{3.165421in}{0.990719in}}%
\pgfpathlineto{\pgfqpoint{3.167274in}{0.988766in}}%
\pgfpathlineto{\pgfqpoint{3.167540in}{0.989537in}}%
\pgfpathlineto{\pgfqpoint{3.169142in}{0.991736in}}%
\pgfpathlineto{\pgfqpoint{3.170132in}{0.990975in}}%
\pgfpathlineto{\pgfqpoint{3.170132in}{0.990975in}}%
\pgfusepath{stroke}%
\end{pgfscope}%
\begin{pgfscope}%
\pgfsetrectcap%
\pgfsetmiterjoin%
\pgfsetlinewidth{0.803000pt}%
\definecolor{currentstroke}{rgb}{0.000000,0.000000,0.000000}%
\pgfsetstrokecolor{currentstroke}%
\pgfsetdash{}{0pt}%
\pgfpathmoveto{\pgfqpoint{0.573768in}{0.462654in}}%
\pgfpathlineto{\pgfqpoint{0.573768in}{1.368904in}}%
\pgfusepath{stroke}%
\end{pgfscope}%
\begin{pgfscope}%
\pgfsetrectcap%
\pgfsetmiterjoin%
\pgfsetlinewidth{0.803000pt}%
\definecolor{currentstroke}{rgb}{0.000000,0.000000,0.000000}%
\pgfsetstrokecolor{currentstroke}%
\pgfsetdash{}{0pt}%
\pgfpathmoveto{\pgfqpoint{3.293768in}{0.462654in}}%
\pgfpathlineto{\pgfqpoint{3.293768in}{1.368904in}}%
\pgfusepath{stroke}%
\end{pgfscope}%
\begin{pgfscope}%
\pgfsetrectcap%
\pgfsetmiterjoin%
\pgfsetlinewidth{0.803000pt}%
\definecolor{currentstroke}{rgb}{0.000000,0.000000,0.000000}%
\pgfsetstrokecolor{currentstroke}%
\pgfsetdash{}{0pt}%
\pgfpathmoveto{\pgfqpoint{0.573768in}{0.462654in}}%
\pgfpathlineto{\pgfqpoint{3.293768in}{0.462654in}}%
\pgfusepath{stroke}%
\end{pgfscope}%
\begin{pgfscope}%
\pgfsetrectcap%
\pgfsetmiterjoin%
\pgfsetlinewidth{0.803000pt}%
\definecolor{currentstroke}{rgb}{0.000000,0.000000,0.000000}%
\pgfsetstrokecolor{currentstroke}%
\pgfsetdash{}{0pt}%
\pgfpathmoveto{\pgfqpoint{0.573768in}{1.368904in}}%
\pgfpathlineto{\pgfqpoint{3.293768in}{1.368904in}}%
\pgfusepath{stroke}%
\end{pgfscope}%
\begin{pgfscope}%
\definecolor{textcolor}{rgb}{0.000000,0.000000,0.000000}%
\pgfsetstrokecolor{textcolor}%
\pgfsetfillcolor{textcolor}%
\pgftext[x=0.628168in,y=1.305467in,left,top]{\color{textcolor}{\rmfamily\fontsize{8.000000}{9.600000}\selectfont\catcode`\^=\active\def^{\ifmmode\sp\else\^{}\fi}\catcode`\%=\active\def%{\%}(c)}}%
\end{pgfscope}%
\begin{pgfscope}%
\pgfsetbuttcap%
\pgfsetmiterjoin%
\definecolor{currentfill}{rgb}{1.000000,1.000000,1.000000}%
\pgfsetfillcolor{currentfill}%
\pgfsetlinewidth{0.000000pt}%
\definecolor{currentstroke}{rgb}{0.000000,0.000000,0.000000}%
\pgfsetstrokecolor{currentstroke}%
\pgfsetstrokeopacity{0.000000}%
\pgfsetdash{}{0pt}%
\pgfpathmoveto{\pgfqpoint{2.197902in}{3.064738in}}%
\pgfpathlineto{\pgfqpoint{3.177102in}{3.064738in}}%
\pgfpathlineto{\pgfqpoint{3.177102in}{3.390988in}}%
\pgfpathlineto{\pgfqpoint{2.197902in}{3.390988in}}%
\pgfpathlineto{\pgfqpoint{2.197902in}{3.064738in}}%
\pgfpathclose%
\pgfusepath{fill}%
\end{pgfscope}%
\begin{pgfscope}%
\pgfpathrectangle{\pgfqpoint{2.197902in}{3.064738in}}{\pgfqpoint{0.979200in}{0.326250in}}%
\pgfusepath{clip}%
\pgfsetbuttcap%
\pgfsetroundjoin%
\pgfsetlinewidth{0.401500pt}%
\definecolor{currentstroke}{rgb}{0.690196,0.690196,0.690196}%
\pgfsetstrokecolor{currentstroke}%
\pgfsetstrokeopacity{0.250000}%
\pgfsetdash{{1.480000pt}{0.640000pt}}{0.000000pt}%
\pgfpathmoveto{\pgfqpoint{2.197902in}{3.064738in}}%
\pgfpathlineto{\pgfqpoint{2.197902in}{3.390988in}}%
\pgfusepath{stroke}%
\end{pgfscope}%
\begin{pgfscope}%
\pgfsetbuttcap%
\pgfsetroundjoin%
\definecolor{currentfill}{rgb}{0.000000,0.000000,0.000000}%
\pgfsetfillcolor{currentfill}%
\pgfsetlinewidth{0.803000pt}%
\definecolor{currentstroke}{rgb}{0.000000,0.000000,0.000000}%
\pgfsetstrokecolor{currentstroke}%
\pgfsetdash{}{0pt}%
\pgfsys@defobject{currentmarker}{\pgfqpoint{0.000000in}{0.000000in}}{\pgfqpoint{0.000000in}{0.027778in}}{%
\pgfpathmoveto{\pgfqpoint{0.000000in}{0.000000in}}%
\pgfpathlineto{\pgfqpoint{0.000000in}{0.027778in}}%
\pgfusepath{stroke,fill}%
}%
\begin{pgfscope}%
\pgfsys@transformshift{2.197902in}{3.064738in}%
\pgfsys@useobject{currentmarker}{}%
\end{pgfscope}%
\end{pgfscope}%
\begin{pgfscope}%
\definecolor{textcolor}{rgb}{0.000000,0.000000,0.000000}%
\pgfsetstrokecolor{textcolor}%
\pgfsetfillcolor{textcolor}%
\pgftext[x=2.197902in,y=3.050849in,,top]{\color{textcolor}{\rmfamily\fontsize{5.000000}{6.000000}\selectfont\catcode`\^=\active\def^{\ifmmode\sp\else\^{}\fi}\catcode`\%=\active\def%{\%}4.980}}%
\end{pgfscope}%
\begin{pgfscope}%
\pgfpathrectangle{\pgfqpoint{2.197902in}{3.064738in}}{\pgfqpoint{0.979200in}{0.326250in}}%
\pgfusepath{clip}%
\pgfsetbuttcap%
\pgfsetroundjoin%
\pgfsetlinewidth{0.401500pt}%
\definecolor{currentstroke}{rgb}{0.690196,0.690196,0.690196}%
\pgfsetstrokecolor{currentstroke}%
\pgfsetstrokeopacity{0.250000}%
\pgfsetdash{{1.480000pt}{0.640000pt}}{0.000000pt}%
\pgfpathmoveto{\pgfqpoint{2.442702in}{3.064738in}}%
\pgfpathlineto{\pgfqpoint{2.442702in}{3.390988in}}%
\pgfusepath{stroke}%
\end{pgfscope}%
\begin{pgfscope}%
\pgfsetbuttcap%
\pgfsetroundjoin%
\definecolor{currentfill}{rgb}{0.000000,0.000000,0.000000}%
\pgfsetfillcolor{currentfill}%
\pgfsetlinewidth{0.803000pt}%
\definecolor{currentstroke}{rgb}{0.000000,0.000000,0.000000}%
\pgfsetstrokecolor{currentstroke}%
\pgfsetdash{}{0pt}%
\pgfsys@defobject{currentmarker}{\pgfqpoint{0.000000in}{0.000000in}}{\pgfqpoint{0.000000in}{0.027778in}}{%
\pgfpathmoveto{\pgfqpoint{0.000000in}{0.000000in}}%
\pgfpathlineto{\pgfqpoint{0.000000in}{0.027778in}}%
\pgfusepath{stroke,fill}%
}%
\begin{pgfscope}%
\pgfsys@transformshift{2.442702in}{3.064738in}%
\pgfsys@useobject{currentmarker}{}%
\end{pgfscope}%
\end{pgfscope}%
\begin{pgfscope}%
\definecolor{textcolor}{rgb}{0.000000,0.000000,0.000000}%
\pgfsetstrokecolor{textcolor}%
\pgfsetfillcolor{textcolor}%
\pgftext[x=2.442702in,y=3.050849in,,top]{\color{textcolor}{\rmfamily\fontsize{5.000000}{6.000000}\selectfont\catcode`\^=\active\def^{\ifmmode\sp\else\^{}\fi}\catcode`\%=\active\def%{\%}4.985}}%
\end{pgfscope}%
\begin{pgfscope}%
\pgfpathrectangle{\pgfqpoint{2.197902in}{3.064738in}}{\pgfqpoint{0.979200in}{0.326250in}}%
\pgfusepath{clip}%
\pgfsetbuttcap%
\pgfsetroundjoin%
\pgfsetlinewidth{0.401500pt}%
\definecolor{currentstroke}{rgb}{0.690196,0.690196,0.690196}%
\pgfsetstrokecolor{currentstroke}%
\pgfsetstrokeopacity{0.250000}%
\pgfsetdash{{1.480000pt}{0.640000pt}}{0.000000pt}%
\pgfpathmoveto{\pgfqpoint{2.687502in}{3.064738in}}%
\pgfpathlineto{\pgfqpoint{2.687502in}{3.390988in}}%
\pgfusepath{stroke}%
\end{pgfscope}%
\begin{pgfscope}%
\pgfsetbuttcap%
\pgfsetroundjoin%
\definecolor{currentfill}{rgb}{0.000000,0.000000,0.000000}%
\pgfsetfillcolor{currentfill}%
\pgfsetlinewidth{0.803000pt}%
\definecolor{currentstroke}{rgb}{0.000000,0.000000,0.000000}%
\pgfsetstrokecolor{currentstroke}%
\pgfsetdash{}{0pt}%
\pgfsys@defobject{currentmarker}{\pgfqpoint{0.000000in}{0.000000in}}{\pgfqpoint{0.000000in}{0.027778in}}{%
\pgfpathmoveto{\pgfqpoint{0.000000in}{0.000000in}}%
\pgfpathlineto{\pgfqpoint{0.000000in}{0.027778in}}%
\pgfusepath{stroke,fill}%
}%
\begin{pgfscope}%
\pgfsys@transformshift{2.687502in}{3.064738in}%
\pgfsys@useobject{currentmarker}{}%
\end{pgfscope}%
\end{pgfscope}%
\begin{pgfscope}%
\definecolor{textcolor}{rgb}{0.000000,0.000000,0.000000}%
\pgfsetstrokecolor{textcolor}%
\pgfsetfillcolor{textcolor}%
\pgftext[x=2.687502in,y=3.050849in,,top]{\color{textcolor}{\rmfamily\fontsize{5.000000}{6.000000}\selectfont\catcode`\^=\active\def^{\ifmmode\sp\else\^{}\fi}\catcode`\%=\active\def%{\%}4.990}}%
\end{pgfscope}%
\begin{pgfscope}%
\pgfpathrectangle{\pgfqpoint{2.197902in}{3.064738in}}{\pgfqpoint{0.979200in}{0.326250in}}%
\pgfusepath{clip}%
\pgfsetbuttcap%
\pgfsetroundjoin%
\pgfsetlinewidth{0.401500pt}%
\definecolor{currentstroke}{rgb}{0.690196,0.690196,0.690196}%
\pgfsetstrokecolor{currentstroke}%
\pgfsetstrokeopacity{0.250000}%
\pgfsetdash{{1.480000pt}{0.640000pt}}{0.000000pt}%
\pgfpathmoveto{\pgfqpoint{2.932302in}{3.064738in}}%
\pgfpathlineto{\pgfqpoint{2.932302in}{3.390988in}}%
\pgfusepath{stroke}%
\end{pgfscope}%
\begin{pgfscope}%
\pgfsetbuttcap%
\pgfsetroundjoin%
\definecolor{currentfill}{rgb}{0.000000,0.000000,0.000000}%
\pgfsetfillcolor{currentfill}%
\pgfsetlinewidth{0.803000pt}%
\definecolor{currentstroke}{rgb}{0.000000,0.000000,0.000000}%
\pgfsetstrokecolor{currentstroke}%
\pgfsetdash{}{0pt}%
\pgfsys@defobject{currentmarker}{\pgfqpoint{0.000000in}{0.000000in}}{\pgfqpoint{0.000000in}{0.027778in}}{%
\pgfpathmoveto{\pgfqpoint{0.000000in}{0.000000in}}%
\pgfpathlineto{\pgfqpoint{0.000000in}{0.027778in}}%
\pgfusepath{stroke,fill}%
}%
\begin{pgfscope}%
\pgfsys@transformshift{2.932302in}{3.064738in}%
\pgfsys@useobject{currentmarker}{}%
\end{pgfscope}%
\end{pgfscope}%
\begin{pgfscope}%
\definecolor{textcolor}{rgb}{0.000000,0.000000,0.000000}%
\pgfsetstrokecolor{textcolor}%
\pgfsetfillcolor{textcolor}%
\pgftext[x=2.932302in,y=3.050849in,,top]{\color{textcolor}{\rmfamily\fontsize{5.000000}{6.000000}\selectfont\catcode`\^=\active\def^{\ifmmode\sp\else\^{}\fi}\catcode`\%=\active\def%{\%}4.995}}%
\end{pgfscope}%
\begin{pgfscope}%
\pgfpathrectangle{\pgfqpoint{2.197902in}{3.064738in}}{\pgfqpoint{0.979200in}{0.326250in}}%
\pgfusepath{clip}%
\pgfsetbuttcap%
\pgfsetroundjoin%
\pgfsetlinewidth{0.401500pt}%
\definecolor{currentstroke}{rgb}{0.690196,0.690196,0.690196}%
\pgfsetstrokecolor{currentstroke}%
\pgfsetstrokeopacity{0.250000}%
\pgfsetdash{{1.480000pt}{0.640000pt}}{0.000000pt}%
\pgfpathmoveto{\pgfqpoint{3.177102in}{3.064738in}}%
\pgfpathlineto{\pgfqpoint{3.177102in}{3.390988in}}%
\pgfusepath{stroke}%
\end{pgfscope}%
\begin{pgfscope}%
\pgfsetbuttcap%
\pgfsetroundjoin%
\definecolor{currentfill}{rgb}{0.000000,0.000000,0.000000}%
\pgfsetfillcolor{currentfill}%
\pgfsetlinewidth{0.803000pt}%
\definecolor{currentstroke}{rgb}{0.000000,0.000000,0.000000}%
\pgfsetstrokecolor{currentstroke}%
\pgfsetdash{}{0pt}%
\pgfsys@defobject{currentmarker}{\pgfqpoint{0.000000in}{0.000000in}}{\pgfqpoint{0.000000in}{0.027778in}}{%
\pgfpathmoveto{\pgfqpoint{0.000000in}{0.000000in}}%
\pgfpathlineto{\pgfqpoint{0.000000in}{0.027778in}}%
\pgfusepath{stroke,fill}%
}%
\begin{pgfscope}%
\pgfsys@transformshift{3.177102in}{3.064738in}%
\pgfsys@useobject{currentmarker}{}%
\end{pgfscope}%
\end{pgfscope}%
\begin{pgfscope}%
\definecolor{textcolor}{rgb}{0.000000,0.000000,0.000000}%
\pgfsetstrokecolor{textcolor}%
\pgfsetfillcolor{textcolor}%
\pgftext[x=3.177102in,y=3.050849in,,top]{\color{textcolor}{\rmfamily\fontsize{5.000000}{6.000000}\selectfont\catcode`\^=\active\def^{\ifmmode\sp\else\^{}\fi}\catcode`\%=\active\def%{\%}5.000}}%
\end{pgfscope}%
\begin{pgfscope}%
\pgfpathrectangle{\pgfqpoint{2.197902in}{3.064738in}}{\pgfqpoint{0.979200in}{0.326250in}}%
\pgfusepath{clip}%
\pgfsetbuttcap%
\pgfsetroundjoin%
\pgfsetlinewidth{0.401500pt}%
\definecolor{currentstroke}{rgb}{0.690196,0.690196,0.690196}%
\pgfsetstrokecolor{currentstroke}%
\pgfsetstrokeopacity{0.250000}%
\pgfsetdash{{1.480000pt}{0.640000pt}}{0.000000pt}%
\pgfpathmoveto{\pgfqpoint{2.197902in}{3.064738in}}%
\pgfpathlineto{\pgfqpoint{3.177102in}{3.064738in}}%
\pgfusepath{stroke}%
\end{pgfscope}%
\begin{pgfscope}%
\pgfsetbuttcap%
\pgfsetroundjoin%
\definecolor{currentfill}{rgb}{0.000000,0.000000,0.000000}%
\pgfsetfillcolor{currentfill}%
\pgfsetlinewidth{0.803000pt}%
\definecolor{currentstroke}{rgb}{0.000000,0.000000,0.000000}%
\pgfsetstrokecolor{currentstroke}%
\pgfsetdash{}{0pt}%
\pgfsys@defobject{currentmarker}{\pgfqpoint{0.000000in}{0.000000in}}{\pgfqpoint{0.027778in}{0.000000in}}{%
\pgfpathmoveto{\pgfqpoint{0.000000in}{0.000000in}}%
\pgfpathlineto{\pgfqpoint{0.027778in}{0.000000in}}%
\pgfusepath{stroke,fill}%
}%
\begin{pgfscope}%
\pgfsys@transformshift{2.197902in}{3.064738in}%
\pgfsys@useobject{currentmarker}{}%
\end{pgfscope}%
\end{pgfscope}%
\begin{pgfscope}%
\definecolor{textcolor}{rgb}{0.000000,0.000000,0.000000}%
\pgfsetstrokecolor{textcolor}%
\pgfsetfillcolor{textcolor}%
\pgftext[x=1.981463in, y=3.040625in, left, base]{\color{textcolor}{\rmfamily\fontsize{5.000000}{6.000000}\selectfont\catcode`\^=\active\def^{\ifmmode\sp\else\^{}\fi}\catcode`\%=\active\def%{\%}-0.50}}%
\end{pgfscope}%
\begin{pgfscope}%
\pgfpathrectangle{\pgfqpoint{2.197902in}{3.064738in}}{\pgfqpoint{0.979200in}{0.326250in}}%
\pgfusepath{clip}%
\pgfsetbuttcap%
\pgfsetroundjoin%
\pgfsetlinewidth{0.401500pt}%
\definecolor{currentstroke}{rgb}{0.690196,0.690196,0.690196}%
\pgfsetstrokecolor{currentstroke}%
\pgfsetstrokeopacity{0.250000}%
\pgfsetdash{{1.480000pt}{0.640000pt}}{0.000000pt}%
\pgfpathmoveto{\pgfqpoint{2.197902in}{3.227863in}}%
\pgfpathlineto{\pgfqpoint{3.177102in}{3.227863in}}%
\pgfusepath{stroke}%
\end{pgfscope}%
\begin{pgfscope}%
\pgfsetbuttcap%
\pgfsetroundjoin%
\definecolor{currentfill}{rgb}{0.000000,0.000000,0.000000}%
\pgfsetfillcolor{currentfill}%
\pgfsetlinewidth{0.803000pt}%
\definecolor{currentstroke}{rgb}{0.000000,0.000000,0.000000}%
\pgfsetstrokecolor{currentstroke}%
\pgfsetdash{}{0pt}%
\pgfsys@defobject{currentmarker}{\pgfqpoint{0.000000in}{0.000000in}}{\pgfqpoint{0.027778in}{0.000000in}}{%
\pgfpathmoveto{\pgfqpoint{0.000000in}{0.000000in}}%
\pgfpathlineto{\pgfqpoint{0.027778in}{0.000000in}}%
\pgfusepath{stroke,fill}%
}%
\begin{pgfscope}%
\pgfsys@transformshift{2.197902in}{3.227863in}%
\pgfsys@useobject{currentmarker}{}%
\end{pgfscope}%
\end{pgfscope}%
\begin{pgfscope}%
\definecolor{textcolor}{rgb}{0.000000,0.000000,0.000000}%
\pgfsetstrokecolor{textcolor}%
\pgfsetfillcolor{textcolor}%
\pgftext[x=2.014257in, y=3.203750in, left, base]{\color{textcolor}{\rmfamily\fontsize{5.000000}{6.000000}\selectfont\catcode`\^=\active\def^{\ifmmode\sp\else\^{}\fi}\catcode`\%=\active\def%{\%}0.73}}%
\end{pgfscope}%
\begin{pgfscope}%
\pgfpathrectangle{\pgfqpoint{2.197902in}{3.064738in}}{\pgfqpoint{0.979200in}{0.326250in}}%
\pgfusepath{clip}%
\pgfsetbuttcap%
\pgfsetroundjoin%
\pgfsetlinewidth{0.401500pt}%
\definecolor{currentstroke}{rgb}{0.690196,0.690196,0.690196}%
\pgfsetstrokecolor{currentstroke}%
\pgfsetstrokeopacity{0.250000}%
\pgfsetdash{{1.480000pt}{0.640000pt}}{0.000000pt}%
\pgfpathmoveto{\pgfqpoint{2.197902in}{3.390988in}}%
\pgfpathlineto{\pgfqpoint{3.177102in}{3.390988in}}%
\pgfusepath{stroke}%
\end{pgfscope}%
\begin{pgfscope}%
\pgfsetbuttcap%
\pgfsetroundjoin%
\definecolor{currentfill}{rgb}{0.000000,0.000000,0.000000}%
\pgfsetfillcolor{currentfill}%
\pgfsetlinewidth{0.803000pt}%
\definecolor{currentstroke}{rgb}{0.000000,0.000000,0.000000}%
\pgfsetstrokecolor{currentstroke}%
\pgfsetdash{}{0pt}%
\pgfsys@defobject{currentmarker}{\pgfqpoint{0.000000in}{0.000000in}}{\pgfqpoint{0.027778in}{0.000000in}}{%
\pgfpathmoveto{\pgfqpoint{0.000000in}{0.000000in}}%
\pgfpathlineto{\pgfqpoint{0.027778in}{0.000000in}}%
\pgfusepath{stroke,fill}%
}%
\begin{pgfscope}%
\pgfsys@transformshift{2.197902in}{3.390988in}%
\pgfsys@useobject{currentmarker}{}%
\end{pgfscope}%
\end{pgfscope}%
\begin{pgfscope}%
\definecolor{textcolor}{rgb}{0.000000,0.000000,0.000000}%
\pgfsetstrokecolor{textcolor}%
\pgfsetfillcolor{textcolor}%
\pgftext[x=2.014257in, y=3.366875in, left, base]{\color{textcolor}{\rmfamily\fontsize{5.000000}{6.000000}\selectfont\catcode`\^=\active\def^{\ifmmode\sp\else\^{}\fi}\catcode`\%=\active\def%{\%}1.95}}%
\end{pgfscope}%
\begin{pgfscope}%
\pgfpathrectangle{\pgfqpoint{2.197902in}{3.064738in}}{\pgfqpoint{0.979200in}{0.326250in}}%
\pgfusepath{clip}%
\pgfsetrectcap%
\pgfsetroundjoin%
\pgfsetlinewidth{1.505625pt}%
\definecolor{currentstroke}{rgb}{0.000000,0.447059,0.698039}%
\pgfsetstrokecolor{currentstroke}%
\pgfsetdash{}{0pt}%
\pgfpathmoveto{\pgfqpoint{2.194249in}{3.131319in}}%
\pgfpathlineto{\pgfqpoint{2.198331in}{3.131319in}}%
\pgfpathlineto{\pgfqpoint{2.198889in}{3.291244in}}%
\pgfpathlineto{\pgfqpoint{2.199306in}{3.232586in}}%
\pgfpathlineto{\pgfqpoint{2.199819in}{3.244645in}}%
\pgfpathlineto{\pgfqpoint{2.200623in}{3.151644in}}%
\pgfpathlineto{\pgfqpoint{2.200793in}{3.177863in}}%
\pgfpathlineto{\pgfqpoint{2.201079in}{3.148342in}}%
\pgfpathlineto{\pgfqpoint{2.201588in}{3.172045in}}%
\pgfpathlineto{\pgfqpoint{2.202553in}{3.140227in}}%
\pgfpathlineto{\pgfqpoint{2.204483in}{3.161965in}}%
\pgfpathlineto{\pgfqpoint{2.208343in}{3.133149in}}%
\pgfpathlineto{\pgfqpoint{2.212915in}{3.163390in}}%
\pgfpathlineto{\pgfqpoint{2.217811in}{3.141094in}}%
\pgfpathlineto{\pgfqpoint{2.222707in}{3.172825in}}%
\pgfpathlineto{\pgfqpoint{2.227603in}{3.150691in}}%
\pgfpathlineto{\pgfqpoint{2.232499in}{3.182404in}}%
\pgfpathlineto{\pgfqpoint{2.237395in}{3.160270in}}%
\pgfpathlineto{\pgfqpoint{2.242291in}{3.191986in}}%
\pgfpathlineto{\pgfqpoint{2.247187in}{3.169863in}}%
\pgfpathlineto{\pgfqpoint{2.252083in}{3.201567in}}%
\pgfpathlineto{\pgfqpoint{2.256979in}{3.179456in}}%
\pgfpathlineto{\pgfqpoint{2.261875in}{3.211148in}}%
\pgfpathlineto{\pgfqpoint{2.266771in}{3.189049in}}%
\pgfpathlineto{\pgfqpoint{2.271667in}{3.220729in}}%
\pgfpathlineto{\pgfqpoint{2.276563in}{3.198643in}}%
\pgfpathlineto{\pgfqpoint{2.281459in}{3.230310in}}%
\pgfpathlineto{\pgfqpoint{2.286355in}{3.208236in}}%
\pgfpathlineto{\pgfqpoint{2.291251in}{3.239891in}}%
\pgfpathlineto{\pgfqpoint{2.296147in}{3.217828in}}%
\pgfpathlineto{\pgfqpoint{2.301043in}{3.249472in}}%
\pgfpathlineto{\pgfqpoint{2.305939in}{3.227421in}}%
\pgfpathlineto{\pgfqpoint{2.310835in}{3.259052in}}%
\pgfpathlineto{\pgfqpoint{2.315731in}{3.237014in}}%
\pgfpathlineto{\pgfqpoint{2.320627in}{3.268633in}}%
\pgfpathlineto{\pgfqpoint{2.325523in}{3.246606in}}%
\pgfpathlineto{\pgfqpoint{2.330419in}{3.278213in}}%
\pgfpathlineto{\pgfqpoint{2.335315in}{3.256199in}}%
\pgfpathlineto{\pgfqpoint{2.340211in}{3.287793in}}%
\pgfpathlineto{\pgfqpoint{2.345107in}{3.265791in}}%
\pgfpathlineto{\pgfqpoint{2.350003in}{3.297374in}}%
\pgfpathlineto{\pgfqpoint{2.354899in}{3.275383in}}%
\pgfpathlineto{\pgfqpoint{2.359795in}{3.306954in}}%
\pgfpathlineto{\pgfqpoint{2.364691in}{3.284975in}}%
\pgfpathlineto{\pgfqpoint{2.369587in}{3.316534in}}%
\pgfpathlineto{\pgfqpoint{2.374483in}{3.294567in}}%
\pgfpathlineto{\pgfqpoint{2.379379in}{3.326114in}}%
\pgfpathlineto{\pgfqpoint{2.384275in}{3.304159in}}%
\pgfpathlineto{\pgfqpoint{2.389171in}{3.335694in}}%
\pgfpathlineto{\pgfqpoint{2.393791in}{3.313481in}}%
\pgfpathlineto{\pgfqpoint{2.393796in}{3.320660in}}%
\pgfpathlineto{\pgfqpoint{2.393846in}{3.270263in}}%
\pgfpathlineto{\pgfqpoint{2.394351in}{3.280011in}}%
\pgfpathlineto{\pgfqpoint{2.395893in}{3.301180in}}%
\pgfpathlineto{\pgfqpoint{2.397786in}{3.256272in}}%
\pgfpathlineto{\pgfqpoint{2.398894in}{3.266254in}}%
\pgfpathlineto{\pgfqpoint{2.400139in}{3.231878in}}%
\pgfpathlineto{\pgfqpoint{2.401415in}{3.131320in}}%
\pgfpathlineto{\pgfqpoint{2.687969in}{3.131319in}}%
\pgfpathlineto{\pgfqpoint{2.688503in}{3.281725in}}%
\pgfpathlineto{\pgfqpoint{2.688965in}{3.241362in}}%
\pgfpathlineto{\pgfqpoint{2.689554in}{3.219000in}}%
\pgfpathlineto{\pgfqpoint{2.690227in}{3.123028in}}%
\pgfpathlineto{\pgfqpoint{2.690656in}{3.202822in}}%
\pgfpathlineto{\pgfqpoint{2.691515in}{3.114656in}}%
\pgfpathlineto{\pgfqpoint{2.693233in}{3.190813in}}%
\pgfpathlineto{\pgfqpoint{2.696669in}{3.104575in}}%
\pgfpathlineto{\pgfqpoint{2.700882in}{3.189196in}}%
\pgfpathlineto{\pgfqpoint{2.705778in}{3.111218in}}%
\pgfpathlineto{\pgfqpoint{2.710674in}{3.198405in}}%
\pgfpathlineto{\pgfqpoint{2.715570in}{3.120809in}}%
\pgfpathlineto{\pgfqpoint{2.720466in}{3.207952in}}%
\pgfpathlineto{\pgfqpoint{2.725362in}{3.130425in}}%
\pgfpathlineto{\pgfqpoint{2.730258in}{3.217517in}}%
\pgfpathlineto{\pgfqpoint{2.735154in}{3.140042in}}%
\pgfpathlineto{\pgfqpoint{2.740050in}{3.227075in}}%
\pgfpathlineto{\pgfqpoint{2.744946in}{3.149659in}}%
\pgfpathlineto{\pgfqpoint{2.749842in}{3.236632in}}%
\pgfpathlineto{\pgfqpoint{2.754738in}{3.159275in}}%
\pgfpathlineto{\pgfqpoint{2.759634in}{3.246190in}}%
\pgfpathlineto{\pgfqpoint{2.764530in}{3.168892in}}%
\pgfpathlineto{\pgfqpoint{2.769426in}{3.255747in}}%
\pgfpathlineto{\pgfqpoint{2.774322in}{3.178508in}}%
\pgfpathlineto{\pgfqpoint{2.779218in}{3.265304in}}%
\pgfpathlineto{\pgfqpoint{2.784114in}{3.188125in}}%
\pgfpathlineto{\pgfqpoint{2.789010in}{3.274861in}}%
\pgfpathlineto{\pgfqpoint{2.793906in}{3.197741in}}%
\pgfpathlineto{\pgfqpoint{2.798802in}{3.284419in}}%
\pgfpathlineto{\pgfqpoint{2.803698in}{3.207357in}}%
\pgfpathlineto{\pgfqpoint{2.808594in}{3.293976in}}%
\pgfpathlineto{\pgfqpoint{2.813490in}{3.216973in}}%
\pgfpathlineto{\pgfqpoint{2.818386in}{3.303533in}}%
\pgfpathlineto{\pgfqpoint{2.823282in}{3.226589in}}%
\pgfpathlineto{\pgfqpoint{2.828178in}{3.313090in}}%
\pgfpathlineto{\pgfqpoint{2.833074in}{3.236205in}}%
\pgfpathlineto{\pgfqpoint{2.837970in}{3.322646in}}%
\pgfpathlineto{\pgfqpoint{2.842866in}{3.245820in}}%
\pgfpathlineto{\pgfqpoint{2.847762in}{3.332203in}}%
\pgfpathlineto{\pgfqpoint{2.852658in}{3.255436in}}%
\pgfpathlineto{\pgfqpoint{2.857554in}{3.341760in}}%
\pgfpathlineto{\pgfqpoint{2.862450in}{3.265051in}}%
\pgfpathlineto{\pgfqpoint{2.867346in}{3.351317in}}%
\pgfpathlineto{\pgfqpoint{2.872242in}{3.274667in}}%
\pgfpathlineto{\pgfqpoint{2.877138in}{3.360873in}}%
\pgfpathlineto{\pgfqpoint{2.880264in}{3.282558in}}%
\pgfpathlineto{\pgfqpoint{2.883391in}{3.366948in}}%
\pgfpathlineto{\pgfqpoint{2.883446in}{3.269564in}}%
\pgfpathlineto{\pgfqpoint{2.883921in}{3.278819in}}%
\pgfpathlineto{\pgfqpoint{2.885370in}{3.299255in}}%
\pgfpathlineto{\pgfqpoint{2.886298in}{3.294389in}}%
\pgfpathlineto{\pgfqpoint{2.887169in}{3.244173in}}%
\pgfpathlineto{\pgfqpoint{2.888228in}{3.289761in}}%
\pgfpathlineto{\pgfqpoint{2.890817in}{3.150662in}}%
\pgfpathlineto{\pgfqpoint{2.892281in}{3.131319in}}%
\pgfpathlineto{\pgfqpoint{3.177102in}{3.131319in}}%
\pgfpathlineto{\pgfqpoint{3.177102in}{3.131319in}}%
\pgfusepath{stroke}%
\end{pgfscope}%
\begin{pgfscope}%
\pgfsetrectcap%
\pgfsetmiterjoin%
\pgfsetlinewidth{0.803000pt}%
\definecolor{currentstroke}{rgb}{0.000000,0.000000,0.000000}%
\pgfsetstrokecolor{currentstroke}%
\pgfsetdash{}{0pt}%
\pgfpathmoveto{\pgfqpoint{2.197902in}{3.064738in}}%
\pgfpathlineto{\pgfqpoint{2.197902in}{3.390988in}}%
\pgfusepath{stroke}%
\end{pgfscope}%
\begin{pgfscope}%
\pgfsetrectcap%
\pgfsetmiterjoin%
\pgfsetlinewidth{0.803000pt}%
\definecolor{currentstroke}{rgb}{0.000000,0.000000,0.000000}%
\pgfsetstrokecolor{currentstroke}%
\pgfsetdash{}{0pt}%
\pgfpathmoveto{\pgfqpoint{3.177102in}{3.064738in}}%
\pgfpathlineto{\pgfqpoint{3.177102in}{3.390988in}}%
\pgfusepath{stroke}%
\end{pgfscope}%
\begin{pgfscope}%
\pgfsetrectcap%
\pgfsetmiterjoin%
\pgfsetlinewidth{0.803000pt}%
\definecolor{currentstroke}{rgb}{0.000000,0.000000,0.000000}%
\pgfsetstrokecolor{currentstroke}%
\pgfsetdash{}{0pt}%
\pgfpathmoveto{\pgfqpoint{2.197902in}{3.064738in}}%
\pgfpathlineto{\pgfqpoint{3.177102in}{3.064738in}}%
\pgfusepath{stroke}%
\end{pgfscope}%
\begin{pgfscope}%
\pgfsetrectcap%
\pgfsetmiterjoin%
\pgfsetlinewidth{0.803000pt}%
\definecolor{currentstroke}{rgb}{0.000000,0.000000,0.000000}%
\pgfsetstrokecolor{currentstroke}%
\pgfsetdash{}{0pt}%
\pgfpathmoveto{\pgfqpoint{2.197902in}{3.390988in}}%
\pgfpathlineto{\pgfqpoint{3.177102in}{3.390988in}}%
\pgfusepath{stroke}%
\end{pgfscope}%
\begin{pgfscope}%
\pgfsetbuttcap%
\pgfsetmiterjoin%
\definecolor{currentfill}{rgb}{1.000000,1.000000,1.000000}%
\pgfsetfillcolor{currentfill}%
\pgfsetlinewidth{0.000000pt}%
\definecolor{currentstroke}{rgb}{0.000000,0.000000,0.000000}%
\pgfsetstrokecolor{currentstroke}%
\pgfsetstrokeopacity{0.000000}%
\pgfsetdash{}{0pt}%
\pgfpathmoveto{\pgfqpoint{2.197902in}{1.995363in}}%
\pgfpathlineto{\pgfqpoint{3.177102in}{1.995363in}}%
\pgfpathlineto{\pgfqpoint{3.177102in}{2.321613in}}%
\pgfpathlineto{\pgfqpoint{2.197902in}{2.321613in}}%
\pgfpathlineto{\pgfqpoint{2.197902in}{1.995363in}}%
\pgfpathclose%
\pgfusepath{fill}%
\end{pgfscope}%
\begin{pgfscope}%
\pgfpathrectangle{\pgfqpoint{2.197902in}{1.995363in}}{\pgfqpoint{0.979200in}{0.326250in}}%
\pgfusepath{clip}%
\pgfsetbuttcap%
\pgfsetroundjoin%
\pgfsetlinewidth{0.401500pt}%
\definecolor{currentstroke}{rgb}{0.690196,0.690196,0.690196}%
\pgfsetstrokecolor{currentstroke}%
\pgfsetstrokeopacity{0.250000}%
\pgfsetdash{{1.480000pt}{0.640000pt}}{0.000000pt}%
\pgfpathmoveto{\pgfqpoint{2.197902in}{1.995363in}}%
\pgfpathlineto{\pgfqpoint{2.197902in}{2.321613in}}%
\pgfusepath{stroke}%
\end{pgfscope}%
\begin{pgfscope}%
\pgfsetbuttcap%
\pgfsetroundjoin%
\definecolor{currentfill}{rgb}{0.000000,0.000000,0.000000}%
\pgfsetfillcolor{currentfill}%
\pgfsetlinewidth{0.803000pt}%
\definecolor{currentstroke}{rgb}{0.000000,0.000000,0.000000}%
\pgfsetstrokecolor{currentstroke}%
\pgfsetdash{}{0pt}%
\pgfsys@defobject{currentmarker}{\pgfqpoint{0.000000in}{0.000000in}}{\pgfqpoint{0.000000in}{0.027778in}}{%
\pgfpathmoveto{\pgfqpoint{0.000000in}{0.000000in}}%
\pgfpathlineto{\pgfqpoint{0.000000in}{0.027778in}}%
\pgfusepath{stroke,fill}%
}%
\begin{pgfscope}%
\pgfsys@transformshift{2.197902in}{1.995363in}%
\pgfsys@useobject{currentmarker}{}%
\end{pgfscope}%
\end{pgfscope}%
\begin{pgfscope}%
\definecolor{textcolor}{rgb}{0.000000,0.000000,0.000000}%
\pgfsetstrokecolor{textcolor}%
\pgfsetfillcolor{textcolor}%
\pgftext[x=2.197902in,y=1.981474in,,top]{\color{textcolor}{\rmfamily\fontsize{5.000000}{6.000000}\selectfont\catcode`\^=\active\def^{\ifmmode\sp\else\^{}\fi}\catcode`\%=\active\def%{\%}4.980}}%
\end{pgfscope}%
\begin{pgfscope}%
\pgfpathrectangle{\pgfqpoint{2.197902in}{1.995363in}}{\pgfqpoint{0.979200in}{0.326250in}}%
\pgfusepath{clip}%
\pgfsetbuttcap%
\pgfsetroundjoin%
\pgfsetlinewidth{0.401500pt}%
\definecolor{currentstroke}{rgb}{0.690196,0.690196,0.690196}%
\pgfsetstrokecolor{currentstroke}%
\pgfsetstrokeopacity{0.250000}%
\pgfsetdash{{1.480000pt}{0.640000pt}}{0.000000pt}%
\pgfpathmoveto{\pgfqpoint{2.442702in}{1.995363in}}%
\pgfpathlineto{\pgfqpoint{2.442702in}{2.321613in}}%
\pgfusepath{stroke}%
\end{pgfscope}%
\begin{pgfscope}%
\pgfsetbuttcap%
\pgfsetroundjoin%
\definecolor{currentfill}{rgb}{0.000000,0.000000,0.000000}%
\pgfsetfillcolor{currentfill}%
\pgfsetlinewidth{0.803000pt}%
\definecolor{currentstroke}{rgb}{0.000000,0.000000,0.000000}%
\pgfsetstrokecolor{currentstroke}%
\pgfsetdash{}{0pt}%
\pgfsys@defobject{currentmarker}{\pgfqpoint{0.000000in}{0.000000in}}{\pgfqpoint{0.000000in}{0.027778in}}{%
\pgfpathmoveto{\pgfqpoint{0.000000in}{0.000000in}}%
\pgfpathlineto{\pgfqpoint{0.000000in}{0.027778in}}%
\pgfusepath{stroke,fill}%
}%
\begin{pgfscope}%
\pgfsys@transformshift{2.442702in}{1.995363in}%
\pgfsys@useobject{currentmarker}{}%
\end{pgfscope}%
\end{pgfscope}%
\begin{pgfscope}%
\definecolor{textcolor}{rgb}{0.000000,0.000000,0.000000}%
\pgfsetstrokecolor{textcolor}%
\pgfsetfillcolor{textcolor}%
\pgftext[x=2.442702in,y=1.981474in,,top]{\color{textcolor}{\rmfamily\fontsize{5.000000}{6.000000}\selectfont\catcode`\^=\active\def^{\ifmmode\sp\else\^{}\fi}\catcode`\%=\active\def%{\%}4.985}}%
\end{pgfscope}%
\begin{pgfscope}%
\pgfpathrectangle{\pgfqpoint{2.197902in}{1.995363in}}{\pgfqpoint{0.979200in}{0.326250in}}%
\pgfusepath{clip}%
\pgfsetbuttcap%
\pgfsetroundjoin%
\pgfsetlinewidth{0.401500pt}%
\definecolor{currentstroke}{rgb}{0.690196,0.690196,0.690196}%
\pgfsetstrokecolor{currentstroke}%
\pgfsetstrokeopacity{0.250000}%
\pgfsetdash{{1.480000pt}{0.640000pt}}{0.000000pt}%
\pgfpathmoveto{\pgfqpoint{2.687502in}{1.995363in}}%
\pgfpathlineto{\pgfqpoint{2.687502in}{2.321613in}}%
\pgfusepath{stroke}%
\end{pgfscope}%
\begin{pgfscope}%
\pgfsetbuttcap%
\pgfsetroundjoin%
\definecolor{currentfill}{rgb}{0.000000,0.000000,0.000000}%
\pgfsetfillcolor{currentfill}%
\pgfsetlinewidth{0.803000pt}%
\definecolor{currentstroke}{rgb}{0.000000,0.000000,0.000000}%
\pgfsetstrokecolor{currentstroke}%
\pgfsetdash{}{0pt}%
\pgfsys@defobject{currentmarker}{\pgfqpoint{0.000000in}{0.000000in}}{\pgfqpoint{0.000000in}{0.027778in}}{%
\pgfpathmoveto{\pgfqpoint{0.000000in}{0.000000in}}%
\pgfpathlineto{\pgfqpoint{0.000000in}{0.027778in}}%
\pgfusepath{stroke,fill}%
}%
\begin{pgfscope}%
\pgfsys@transformshift{2.687502in}{1.995363in}%
\pgfsys@useobject{currentmarker}{}%
\end{pgfscope}%
\end{pgfscope}%
\begin{pgfscope}%
\definecolor{textcolor}{rgb}{0.000000,0.000000,0.000000}%
\pgfsetstrokecolor{textcolor}%
\pgfsetfillcolor{textcolor}%
\pgftext[x=2.687502in,y=1.981474in,,top]{\color{textcolor}{\rmfamily\fontsize{5.000000}{6.000000}\selectfont\catcode`\^=\active\def^{\ifmmode\sp\else\^{}\fi}\catcode`\%=\active\def%{\%}4.990}}%
\end{pgfscope}%
\begin{pgfscope}%
\pgfpathrectangle{\pgfqpoint{2.197902in}{1.995363in}}{\pgfqpoint{0.979200in}{0.326250in}}%
\pgfusepath{clip}%
\pgfsetbuttcap%
\pgfsetroundjoin%
\pgfsetlinewidth{0.401500pt}%
\definecolor{currentstroke}{rgb}{0.690196,0.690196,0.690196}%
\pgfsetstrokecolor{currentstroke}%
\pgfsetstrokeopacity{0.250000}%
\pgfsetdash{{1.480000pt}{0.640000pt}}{0.000000pt}%
\pgfpathmoveto{\pgfqpoint{2.932302in}{1.995363in}}%
\pgfpathlineto{\pgfqpoint{2.932302in}{2.321613in}}%
\pgfusepath{stroke}%
\end{pgfscope}%
\begin{pgfscope}%
\pgfsetbuttcap%
\pgfsetroundjoin%
\definecolor{currentfill}{rgb}{0.000000,0.000000,0.000000}%
\pgfsetfillcolor{currentfill}%
\pgfsetlinewidth{0.803000pt}%
\definecolor{currentstroke}{rgb}{0.000000,0.000000,0.000000}%
\pgfsetstrokecolor{currentstroke}%
\pgfsetdash{}{0pt}%
\pgfsys@defobject{currentmarker}{\pgfqpoint{0.000000in}{0.000000in}}{\pgfqpoint{0.000000in}{0.027778in}}{%
\pgfpathmoveto{\pgfqpoint{0.000000in}{0.000000in}}%
\pgfpathlineto{\pgfqpoint{0.000000in}{0.027778in}}%
\pgfusepath{stroke,fill}%
}%
\begin{pgfscope}%
\pgfsys@transformshift{2.932302in}{1.995363in}%
\pgfsys@useobject{currentmarker}{}%
\end{pgfscope}%
\end{pgfscope}%
\begin{pgfscope}%
\definecolor{textcolor}{rgb}{0.000000,0.000000,0.000000}%
\pgfsetstrokecolor{textcolor}%
\pgfsetfillcolor{textcolor}%
\pgftext[x=2.932302in,y=1.981474in,,top]{\color{textcolor}{\rmfamily\fontsize{5.000000}{6.000000}\selectfont\catcode`\^=\active\def^{\ifmmode\sp\else\^{}\fi}\catcode`\%=\active\def%{\%}4.995}}%
\end{pgfscope}%
\begin{pgfscope}%
\pgfpathrectangle{\pgfqpoint{2.197902in}{1.995363in}}{\pgfqpoint{0.979200in}{0.326250in}}%
\pgfusepath{clip}%
\pgfsetbuttcap%
\pgfsetroundjoin%
\pgfsetlinewidth{0.401500pt}%
\definecolor{currentstroke}{rgb}{0.690196,0.690196,0.690196}%
\pgfsetstrokecolor{currentstroke}%
\pgfsetstrokeopacity{0.250000}%
\pgfsetdash{{1.480000pt}{0.640000pt}}{0.000000pt}%
\pgfpathmoveto{\pgfqpoint{3.177102in}{1.995363in}}%
\pgfpathlineto{\pgfqpoint{3.177102in}{2.321613in}}%
\pgfusepath{stroke}%
\end{pgfscope}%
\begin{pgfscope}%
\pgfsetbuttcap%
\pgfsetroundjoin%
\definecolor{currentfill}{rgb}{0.000000,0.000000,0.000000}%
\pgfsetfillcolor{currentfill}%
\pgfsetlinewidth{0.803000pt}%
\definecolor{currentstroke}{rgb}{0.000000,0.000000,0.000000}%
\pgfsetstrokecolor{currentstroke}%
\pgfsetdash{}{0pt}%
\pgfsys@defobject{currentmarker}{\pgfqpoint{0.000000in}{0.000000in}}{\pgfqpoint{0.000000in}{0.027778in}}{%
\pgfpathmoveto{\pgfqpoint{0.000000in}{0.000000in}}%
\pgfpathlineto{\pgfqpoint{0.000000in}{0.027778in}}%
\pgfusepath{stroke,fill}%
}%
\begin{pgfscope}%
\pgfsys@transformshift{3.177102in}{1.995363in}%
\pgfsys@useobject{currentmarker}{}%
\end{pgfscope}%
\end{pgfscope}%
\begin{pgfscope}%
\definecolor{textcolor}{rgb}{0.000000,0.000000,0.000000}%
\pgfsetstrokecolor{textcolor}%
\pgfsetfillcolor{textcolor}%
\pgftext[x=3.177102in,y=1.981474in,,top]{\color{textcolor}{\rmfamily\fontsize{5.000000}{6.000000}\selectfont\catcode`\^=\active\def^{\ifmmode\sp\else\^{}\fi}\catcode`\%=\active\def%{\%}5.000}}%
\end{pgfscope}%
\begin{pgfscope}%
\pgfpathrectangle{\pgfqpoint{2.197902in}{1.995363in}}{\pgfqpoint{0.979200in}{0.326250in}}%
\pgfusepath{clip}%
\pgfsetbuttcap%
\pgfsetroundjoin%
\pgfsetlinewidth{0.401500pt}%
\definecolor{currentstroke}{rgb}{0.690196,0.690196,0.690196}%
\pgfsetstrokecolor{currentstroke}%
\pgfsetstrokeopacity{0.250000}%
\pgfsetdash{{1.480000pt}{0.640000pt}}{0.000000pt}%
\pgfpathmoveto{\pgfqpoint{2.197902in}{1.995363in}}%
\pgfpathlineto{\pgfqpoint{3.177102in}{1.995363in}}%
\pgfusepath{stroke}%
\end{pgfscope}%
\begin{pgfscope}%
\pgfsetbuttcap%
\pgfsetroundjoin%
\definecolor{currentfill}{rgb}{0.000000,0.000000,0.000000}%
\pgfsetfillcolor{currentfill}%
\pgfsetlinewidth{0.803000pt}%
\definecolor{currentstroke}{rgb}{0.000000,0.000000,0.000000}%
\pgfsetstrokecolor{currentstroke}%
\pgfsetdash{}{0pt}%
\pgfsys@defobject{currentmarker}{\pgfqpoint{0.000000in}{0.000000in}}{\pgfqpoint{0.027778in}{0.000000in}}{%
\pgfpathmoveto{\pgfqpoint{0.000000in}{0.000000in}}%
\pgfpathlineto{\pgfqpoint{0.027778in}{0.000000in}}%
\pgfusepath{stroke,fill}%
}%
\begin{pgfscope}%
\pgfsys@transformshift{2.197902in}{1.995363in}%
\pgfsys@useobject{currentmarker}{}%
\end{pgfscope}%
\end{pgfscope}%
\begin{pgfscope}%
\definecolor{textcolor}{rgb}{0.000000,0.000000,0.000000}%
\pgfsetstrokecolor{textcolor}%
\pgfsetfillcolor{textcolor}%
\pgftext[x=1.966996in, y=1.971250in, left, base]{\color{textcolor}{\rmfamily\fontsize{5.000000}{6.000000}\selectfont\catcode`\^=\active\def^{\ifmmode\sp\else\^{}\fi}\catcode`\%=\active\def%{\%}45.00}}%
\end{pgfscope}%
\begin{pgfscope}%
\pgfpathrectangle{\pgfqpoint{2.197902in}{1.995363in}}{\pgfqpoint{0.979200in}{0.326250in}}%
\pgfusepath{clip}%
\pgfsetbuttcap%
\pgfsetroundjoin%
\pgfsetlinewidth{0.401500pt}%
\definecolor{currentstroke}{rgb}{0.690196,0.690196,0.690196}%
\pgfsetstrokecolor{currentstroke}%
\pgfsetstrokeopacity{0.250000}%
\pgfsetdash{{1.480000pt}{0.640000pt}}{0.000000pt}%
\pgfpathmoveto{\pgfqpoint{2.197902in}{2.158488in}}%
\pgfpathlineto{\pgfqpoint{3.177102in}{2.158488in}}%
\pgfusepath{stroke}%
\end{pgfscope}%
\begin{pgfscope}%
\pgfsetbuttcap%
\pgfsetroundjoin%
\definecolor{currentfill}{rgb}{0.000000,0.000000,0.000000}%
\pgfsetfillcolor{currentfill}%
\pgfsetlinewidth{0.803000pt}%
\definecolor{currentstroke}{rgb}{0.000000,0.000000,0.000000}%
\pgfsetstrokecolor{currentstroke}%
\pgfsetdash{}{0pt}%
\pgfsys@defobject{currentmarker}{\pgfqpoint{0.000000in}{0.000000in}}{\pgfqpoint{0.027778in}{0.000000in}}{%
\pgfpathmoveto{\pgfqpoint{0.000000in}{0.000000in}}%
\pgfpathlineto{\pgfqpoint{0.027778in}{0.000000in}}%
\pgfusepath{stroke,fill}%
}%
\begin{pgfscope}%
\pgfsys@transformshift{2.197902in}{2.158488in}%
\pgfsys@useobject{currentmarker}{}%
\end{pgfscope}%
\end{pgfscope}%
\begin{pgfscope}%
\definecolor{textcolor}{rgb}{0.000000,0.000000,0.000000}%
\pgfsetstrokecolor{textcolor}%
\pgfsetfillcolor{textcolor}%
\pgftext[x=1.966996in, y=2.134375in, left, base]{\color{textcolor}{\rmfamily\fontsize{5.000000}{6.000000}\selectfont\catcode`\^=\active\def^{\ifmmode\sp\else\^{}\fi}\catcode`\%=\active\def%{\%}65.00}}%
\end{pgfscope}%
\begin{pgfscope}%
\pgfpathrectangle{\pgfqpoint{2.197902in}{1.995363in}}{\pgfqpoint{0.979200in}{0.326250in}}%
\pgfusepath{clip}%
\pgfsetbuttcap%
\pgfsetroundjoin%
\pgfsetlinewidth{0.401500pt}%
\definecolor{currentstroke}{rgb}{0.690196,0.690196,0.690196}%
\pgfsetstrokecolor{currentstroke}%
\pgfsetstrokeopacity{0.250000}%
\pgfsetdash{{1.480000pt}{0.640000pt}}{0.000000pt}%
\pgfpathmoveto{\pgfqpoint{2.197902in}{2.321613in}}%
\pgfpathlineto{\pgfqpoint{3.177102in}{2.321613in}}%
\pgfusepath{stroke}%
\end{pgfscope}%
\begin{pgfscope}%
\pgfsetbuttcap%
\pgfsetroundjoin%
\definecolor{currentfill}{rgb}{0.000000,0.000000,0.000000}%
\pgfsetfillcolor{currentfill}%
\pgfsetlinewidth{0.803000pt}%
\definecolor{currentstroke}{rgb}{0.000000,0.000000,0.000000}%
\pgfsetstrokecolor{currentstroke}%
\pgfsetdash{}{0pt}%
\pgfsys@defobject{currentmarker}{\pgfqpoint{0.000000in}{0.000000in}}{\pgfqpoint{0.027778in}{0.000000in}}{%
\pgfpathmoveto{\pgfqpoint{0.000000in}{0.000000in}}%
\pgfpathlineto{\pgfqpoint{0.027778in}{0.000000in}}%
\pgfusepath{stroke,fill}%
}%
\begin{pgfscope}%
\pgfsys@transformshift{2.197902in}{2.321613in}%
\pgfsys@useobject{currentmarker}{}%
\end{pgfscope}%
\end{pgfscope}%
\begin{pgfscope}%
\definecolor{textcolor}{rgb}{0.000000,0.000000,0.000000}%
\pgfsetstrokecolor{textcolor}%
\pgfsetfillcolor{textcolor}%
\pgftext[x=1.966996in, y=2.297500in, left, base]{\color{textcolor}{\rmfamily\fontsize{5.000000}{6.000000}\selectfont\catcode`\^=\active\def^{\ifmmode\sp\else\^{}\fi}\catcode`\%=\active\def%{\%}85.00}}%
\end{pgfscope}%
\begin{pgfscope}%
\pgfpathrectangle{\pgfqpoint{2.197902in}{1.995363in}}{\pgfqpoint{0.979200in}{0.326250in}}%
\pgfusepath{clip}%
\pgfsetrectcap%
\pgfsetroundjoin%
\pgfsetlinewidth{1.505625pt}%
\definecolor{currentstroke}{rgb}{0.835294,0.368627,0.000000}%
\pgfsetstrokecolor{currentstroke}%
\pgfsetdash{}{0pt}%
\pgfpathmoveto{\pgfqpoint{2.400234in}{1.988418in}}%
\pgfpathlineto{\pgfqpoint{2.401415in}{2.164073in}}%
\pgfpathlineto{\pgfqpoint{2.404735in}{2.197210in}}%
\pgfpathlineto{\pgfqpoint{2.405893in}{2.201463in}}%
\pgfpathlineto{\pgfqpoint{2.407565in}{2.160270in}}%
\pgfpathlineto{\pgfqpoint{2.409240in}{2.149930in}}%
\pgfpathlineto{\pgfqpoint{2.410852in}{2.185986in}}%
\pgfpathlineto{\pgfqpoint{2.412399in}{2.145528in}}%
\pgfpathlineto{\pgfqpoint{2.413190in}{2.154457in}}%
\pgfpathlineto{\pgfqpoint{2.414260in}{2.185579in}}%
\pgfpathlineto{\pgfqpoint{2.415639in}{2.144963in}}%
\pgfpathlineto{\pgfqpoint{2.416079in}{2.149899in}}%
\pgfpathlineto{\pgfqpoint{2.417417in}{2.183099in}}%
\pgfpathlineto{\pgfqpoint{2.418767in}{2.145607in}}%
\pgfpathlineto{\pgfqpoint{2.420125in}{2.183802in}}%
\pgfpathlineto{\pgfqpoint{2.420599in}{2.177642in}}%
\pgfpathlineto{\pgfqpoint{2.421814in}{2.146422in}}%
\pgfpathlineto{\pgfqpoint{2.423255in}{2.183315in}}%
\pgfpathlineto{\pgfqpoint{2.423730in}{2.174898in}}%
\pgfpathlineto{\pgfqpoint{2.424686in}{2.147219in}}%
\pgfpathlineto{\pgfqpoint{2.425176in}{2.153560in}}%
\pgfpathlineto{\pgfqpoint{2.426287in}{2.182631in}}%
\pgfpathlineto{\pgfqpoint{2.427812in}{2.147779in}}%
\pgfpathlineto{\pgfqpoint{2.428311in}{2.156663in}}%
\pgfpathlineto{\pgfqpoint{2.429279in}{2.181945in}}%
\pgfpathlineto{\pgfqpoint{2.430812in}{2.148484in}}%
\pgfpathlineto{\pgfqpoint{2.431325in}{2.156772in}}%
\pgfpathlineto{\pgfqpoint{2.432349in}{2.181320in}}%
\pgfpathlineto{\pgfqpoint{2.433848in}{2.149175in}}%
\pgfpathlineto{\pgfqpoint{2.434356in}{2.157285in}}%
\pgfpathlineto{\pgfqpoint{2.435361in}{2.180709in}}%
\pgfpathlineto{\pgfqpoint{2.436889in}{2.149846in}}%
\pgfpathlineto{\pgfqpoint{2.437407in}{2.158228in}}%
\pgfpathlineto{\pgfqpoint{2.438406in}{2.180107in}}%
\pgfpathlineto{\pgfqpoint{2.439920in}{2.150489in}}%
\pgfpathlineto{\pgfqpoint{2.440441in}{2.158706in}}%
\pgfpathlineto{\pgfqpoint{2.441450in}{2.179512in}}%
\pgfpathlineto{\pgfqpoint{2.442976in}{2.151151in}}%
\pgfpathlineto{\pgfqpoint{2.444509in}{2.178893in}}%
\pgfpathlineto{\pgfqpoint{2.446093in}{2.152065in}}%
\pgfpathlineto{\pgfqpoint{2.447658in}{2.177806in}}%
\pgfpathlineto{\pgfqpoint{2.449181in}{2.152939in}}%
\pgfpathlineto{\pgfqpoint{2.450714in}{2.177191in}}%
\pgfpathlineto{\pgfqpoint{2.452243in}{2.153636in}}%
\pgfpathlineto{\pgfqpoint{2.453787in}{2.176468in}}%
\pgfpathlineto{\pgfqpoint{2.455331in}{2.154522in}}%
\pgfpathlineto{\pgfqpoint{2.456899in}{2.175384in}}%
\pgfpathlineto{\pgfqpoint{2.457954in}{2.154929in}}%
\pgfpathlineto{\pgfqpoint{2.458480in}{2.156029in}}%
\pgfpathlineto{\pgfqpoint{2.459589in}{2.176198in}}%
\pgfpathlineto{\pgfqpoint{2.460101in}{2.173271in}}%
\pgfpathlineto{\pgfqpoint{2.461301in}{2.154266in}}%
\pgfpathlineto{\pgfqpoint{2.462899in}{2.176095in}}%
\pgfpathlineto{\pgfqpoint{2.463352in}{2.170337in}}%
\pgfpathlineto{\pgfqpoint{2.464458in}{2.154937in}}%
\pgfpathlineto{\pgfqpoint{2.466049in}{2.175274in}}%
\pgfpathlineto{\pgfqpoint{2.467171in}{2.156699in}}%
\pgfpathlineto{\pgfqpoint{2.467805in}{2.157558in}}%
\pgfpathlineto{\pgfqpoint{2.469100in}{2.175368in}}%
\pgfpathlineto{\pgfqpoint{2.470305in}{2.157428in}}%
\pgfpathlineto{\pgfqpoint{2.470989in}{2.157889in}}%
\pgfpathlineto{\pgfqpoint{2.472207in}{2.175070in}}%
\pgfpathlineto{\pgfqpoint{2.474008in}{2.156816in}}%
\pgfpathlineto{\pgfqpoint{2.475236in}{2.174393in}}%
\pgfpathlineto{\pgfqpoint{2.475903in}{2.170196in}}%
\pgfpathlineto{\pgfqpoint{2.476686in}{2.158139in}}%
\pgfpathlineto{\pgfqpoint{2.477492in}{2.159378in}}%
\pgfpathlineto{\pgfqpoint{2.478897in}{2.173769in}}%
\pgfpathlineto{\pgfqpoint{2.480184in}{2.156886in}}%
\pgfpathlineto{\pgfqpoint{2.480875in}{2.161544in}}%
\pgfpathlineto{\pgfqpoint{2.481725in}{2.173130in}}%
\pgfpathlineto{\pgfqpoint{2.482514in}{2.170267in}}%
\pgfpathlineto{\pgfqpoint{2.483435in}{2.158146in}}%
\pgfpathlineto{\pgfqpoint{2.484142in}{2.159722in}}%
\pgfpathlineto{\pgfqpoint{2.485612in}{2.172898in}}%
\pgfpathlineto{\pgfqpoint{2.486932in}{2.157319in}}%
\pgfpathlineto{\pgfqpoint{2.488369in}{2.172655in}}%
\pgfpathlineto{\pgfqpoint{2.489184in}{2.169732in}}%
\pgfpathlineto{\pgfqpoint{2.489931in}{2.159492in}}%
\pgfpathlineto{\pgfqpoint{2.490890in}{2.160961in}}%
\pgfpathlineto{\pgfqpoint{2.491797in}{2.172077in}}%
\pgfpathlineto{\pgfqpoint{2.492540in}{2.170750in}}%
\pgfpathlineto{\pgfqpoint{2.494108in}{2.159291in}}%
\pgfpathlineto{\pgfqpoint{2.495526in}{2.172936in}}%
\pgfpathlineto{\pgfqpoint{2.496993in}{2.158280in}}%
\pgfpathlineto{\pgfqpoint{2.498609in}{2.172433in}}%
\pgfpathlineto{\pgfqpoint{2.500547in}{2.158480in}}%
\pgfpathlineto{\pgfqpoint{2.502526in}{2.172588in}}%
\pgfpathlineto{\pgfqpoint{2.504524in}{2.159102in}}%
\pgfpathlineto{\pgfqpoint{2.506293in}{2.171957in}}%
\pgfpathlineto{\pgfqpoint{2.507945in}{2.159103in}}%
\pgfpathlineto{\pgfqpoint{2.509601in}{2.172130in}}%
\pgfpathlineto{\pgfqpoint{2.511365in}{2.159187in}}%
\pgfpathlineto{\pgfqpoint{2.513133in}{2.171908in}}%
\pgfpathlineto{\pgfqpoint{2.515046in}{2.159677in}}%
\pgfpathlineto{\pgfqpoint{2.516849in}{2.171350in}}%
\pgfpathlineto{\pgfqpoint{2.518640in}{2.160098in}}%
\pgfpathlineto{\pgfqpoint{2.520384in}{2.171092in}}%
\pgfpathlineto{\pgfqpoint{2.522135in}{2.160252in}}%
\pgfpathlineto{\pgfqpoint{2.523885in}{2.170929in}}%
\pgfpathlineto{\pgfqpoint{2.525684in}{2.160548in}}%
\pgfpathlineto{\pgfqpoint{2.527518in}{2.170391in}}%
\pgfpathlineto{\pgfqpoint{2.528349in}{2.162978in}}%
\pgfpathlineto{\pgfqpoint{2.529420in}{2.161471in}}%
\pgfpathlineto{\pgfqpoint{2.530450in}{2.169800in}}%
\pgfpathlineto{\pgfqpoint{2.531494in}{2.168327in}}%
\pgfpathlineto{\pgfqpoint{2.532477in}{2.160731in}}%
\pgfpathlineto{\pgfqpoint{2.535677in}{2.169536in}}%
\pgfpathlineto{\pgfqpoint{2.537397in}{2.160519in}}%
\pgfpathlineto{\pgfqpoint{2.539765in}{2.168857in}}%
\pgfpathlineto{\pgfqpoint{2.541994in}{2.166067in}}%
\pgfpathlineto{\pgfqpoint{2.544140in}{2.161774in}}%
\pgfpathlineto{\pgfqpoint{2.547176in}{2.170698in}}%
\pgfpathlineto{\pgfqpoint{2.552072in}{2.160632in}}%
\pgfpathlineto{\pgfqpoint{2.556968in}{2.169969in}}%
\pgfpathlineto{\pgfqpoint{2.561864in}{2.162468in}}%
\pgfpathlineto{\pgfqpoint{2.566760in}{2.167173in}}%
\pgfpathlineto{\pgfqpoint{2.576552in}{2.163650in}}%
\pgfpathlineto{\pgfqpoint{2.581448in}{2.168913in}}%
\pgfpathlineto{\pgfqpoint{2.586344in}{2.161064in}}%
\pgfpathlineto{\pgfqpoint{2.591240in}{2.170395in}}%
\pgfpathlineto{\pgfqpoint{2.596136in}{2.160602in}}%
\pgfpathlineto{\pgfqpoint{2.601032in}{2.169487in}}%
\pgfpathlineto{\pgfqpoint{2.605928in}{2.162415in}}%
\pgfpathlineto{\pgfqpoint{2.610824in}{2.166603in}}%
\pgfpathlineto{\pgfqpoint{2.615720in}{2.165551in}}%
\pgfpathlineto{\pgfqpoint{2.620616in}{2.163076in}}%
\pgfpathlineto{\pgfqpoint{2.625512in}{2.168415in}}%
\pgfpathlineto{\pgfqpoint{2.630408in}{2.160512in}}%
\pgfpathlineto{\pgfqpoint{2.635304in}{2.169481in}}%
\pgfpathlineto{\pgfqpoint{2.640200in}{2.159944in}}%
\pgfpathlineto{\pgfqpoint{2.645096in}{2.165797in}}%
\pgfpathlineto{\pgfqpoint{2.649992in}{2.160059in}}%
\pgfpathlineto{\pgfqpoint{2.654888in}{2.100808in}}%
\pgfpathlineto{\pgfqpoint{2.660661in}{1.988418in}}%
\pgfpathmoveto{\pgfqpoint{2.889953in}{1.988418in}}%
\pgfpathlineto{\pgfqpoint{2.890817in}{2.160811in}}%
\pgfpathlineto{\pgfqpoint{2.894171in}{2.195090in}}%
\pgfpathlineto{\pgfqpoint{2.895638in}{2.200943in}}%
\pgfpathlineto{\pgfqpoint{2.898009in}{2.157888in}}%
\pgfpathlineto{\pgfqpoint{2.901173in}{2.150423in}}%
\pgfpathlineto{\pgfqpoint{2.905760in}{2.191993in}}%
\pgfpathlineto{\pgfqpoint{2.910656in}{2.129417in}}%
\pgfpathlineto{\pgfqpoint{2.915552in}{2.201562in}}%
\pgfpathlineto{\pgfqpoint{2.920448in}{2.131887in}}%
\pgfpathlineto{\pgfqpoint{2.924459in}{2.192234in}}%
\pgfpathlineto{\pgfqpoint{2.926527in}{2.160581in}}%
\pgfpathlineto{\pgfqpoint{2.927171in}{2.132411in}}%
\pgfpathlineto{\pgfqpoint{2.929049in}{2.197983in}}%
\pgfpathlineto{\pgfqpoint{2.930924in}{2.137634in}}%
\pgfpathlineto{\pgfqpoint{2.932211in}{2.196593in}}%
\pgfpathlineto{\pgfqpoint{2.933595in}{2.134683in}}%
\pgfpathlineto{\pgfqpoint{2.934085in}{2.145929in}}%
\pgfpathlineto{\pgfqpoint{2.935186in}{2.195528in}}%
\pgfpathlineto{\pgfqpoint{2.936721in}{2.135732in}}%
\pgfpathlineto{\pgfqpoint{2.937227in}{2.151584in}}%
\pgfpathlineto{\pgfqpoint{2.938194in}{2.194290in}}%
\pgfpathlineto{\pgfqpoint{2.939719in}{2.136944in}}%
\pgfpathlineto{\pgfqpoint{2.940234in}{2.151345in}}%
\pgfpathlineto{\pgfqpoint{2.941260in}{2.193129in}}%
\pgfpathlineto{\pgfqpoint{2.942763in}{2.138080in}}%
\pgfpathlineto{\pgfqpoint{2.943275in}{2.152595in}}%
\pgfpathlineto{\pgfqpoint{2.944279in}{2.192042in}}%
\pgfpathlineto{\pgfqpoint{2.945800in}{2.139217in}}%
\pgfpathlineto{\pgfqpoint{2.946321in}{2.153932in}}%
\pgfpathlineto{\pgfqpoint{2.947332in}{2.190920in}}%
\pgfpathlineto{\pgfqpoint{2.948854in}{2.140377in}}%
\pgfpathlineto{\pgfqpoint{2.950390in}{2.189771in}}%
\pgfpathlineto{\pgfqpoint{2.951969in}{2.141974in}}%
\pgfpathlineto{\pgfqpoint{2.953538in}{2.187764in}}%
\pgfpathlineto{\pgfqpoint{2.955058in}{2.143523in}}%
\pgfpathlineto{\pgfqpoint{2.956596in}{2.186601in}}%
\pgfpathlineto{\pgfqpoint{2.958123in}{2.144775in}}%
\pgfpathlineto{\pgfqpoint{2.959671in}{2.185240in}}%
\pgfpathlineto{\pgfqpoint{2.961215in}{2.146403in}}%
\pgfpathlineto{\pgfqpoint{2.962249in}{2.183239in}}%
\pgfpathlineto{\pgfqpoint{2.962792in}{2.183111in}}%
\pgfpathlineto{\pgfqpoint{2.963856in}{2.146686in}}%
\pgfpathlineto{\pgfqpoint{2.964376in}{2.149299in}}%
\pgfpathlineto{\pgfqpoint{2.965503in}{2.185086in}}%
\pgfpathlineto{\pgfqpoint{2.966011in}{2.178883in}}%
\pgfpathlineto{\pgfqpoint{2.967123in}{2.145851in}}%
\pgfpathlineto{\pgfqpoint{2.968700in}{2.184780in}}%
\pgfpathlineto{\pgfqpoint{2.969171in}{2.175780in}}%
\pgfpathlineto{\pgfqpoint{2.970274in}{2.146711in}}%
\pgfpathlineto{\pgfqpoint{2.971876in}{2.183511in}}%
\pgfpathlineto{\pgfqpoint{2.973507in}{2.148940in}}%
\pgfpathlineto{\pgfqpoint{2.974689in}{2.182228in}}%
\pgfpathlineto{\pgfqpoint{2.975335in}{2.176264in}}%
\pgfpathlineto{\pgfqpoint{2.976777in}{2.150543in}}%
\pgfpathlineto{\pgfqpoint{2.977978in}{2.182307in}}%
\pgfpathlineto{\pgfqpoint{2.978553in}{2.175419in}}%
\pgfpathlineto{\pgfqpoint{2.979867in}{2.149914in}}%
\pgfpathlineto{\pgfqpoint{2.981091in}{2.181140in}}%
\pgfpathlineto{\pgfqpoint{2.981765in}{2.174698in}}%
\pgfpathlineto{\pgfqpoint{2.982506in}{2.153699in}}%
\pgfpathlineto{\pgfqpoint{2.983374in}{2.154402in}}%
\pgfpathlineto{\pgfqpoint{2.984748in}{2.180705in}}%
\pgfpathlineto{\pgfqpoint{2.986020in}{2.150778in}}%
\pgfpathlineto{\pgfqpoint{2.986715in}{2.157378in}}%
\pgfpathlineto{\pgfqpoint{2.987497in}{2.177681in}}%
\pgfpathlineto{\pgfqpoint{2.988376in}{2.174879in}}%
\pgfpathlineto{\pgfqpoint{2.989912in}{2.153521in}}%
\pgfpathlineto{\pgfqpoint{2.991282in}{2.180199in}}%
\pgfpathlineto{\pgfqpoint{2.992609in}{2.152133in}}%
\pgfpathlineto{\pgfqpoint{2.993307in}{2.157279in}}%
\pgfpathlineto{\pgfqpoint{2.994205in}{2.177542in}}%
\pgfpathlineto{\pgfqpoint{2.994950in}{2.175179in}}%
\pgfpathlineto{\pgfqpoint{2.996466in}{2.153645in}}%
\pgfpathlineto{\pgfqpoint{2.997905in}{2.179203in}}%
\pgfpathlineto{\pgfqpoint{2.999363in}{2.152430in}}%
\pgfpathlineto{\pgfqpoint{3.001011in}{2.178172in}}%
\pgfpathlineto{\pgfqpoint{3.002958in}{2.152656in}}%
\pgfpathlineto{\pgfqpoint{3.004933in}{2.178453in}}%
\pgfpathlineto{\pgfqpoint{3.006905in}{2.153678in}}%
\pgfpathlineto{\pgfqpoint{3.008644in}{2.177521in}}%
\pgfpathlineto{\pgfqpoint{3.010385in}{2.153890in}}%
\pgfpathlineto{\pgfqpoint{3.012073in}{2.177414in}}%
\pgfpathlineto{\pgfqpoint{3.013841in}{2.154095in}}%
\pgfpathlineto{\pgfqpoint{3.015624in}{2.176902in}}%
\pgfpathlineto{\pgfqpoint{3.017489in}{2.154982in}}%
\pgfpathlineto{\pgfqpoint{3.019245in}{2.176141in}}%
\pgfpathlineto{\pgfqpoint{3.020997in}{2.155344in}}%
\pgfpathlineto{\pgfqpoint{3.022710in}{2.175961in}}%
\pgfpathlineto{\pgfqpoint{3.024568in}{2.155923in}}%
\pgfpathlineto{\pgfqpoint{3.026376in}{2.174925in}}%
\pgfpathlineto{\pgfqpoint{3.028211in}{2.157007in}}%
\pgfpathlineto{\pgfqpoint{3.029192in}{2.172182in}}%
\pgfpathlineto{\pgfqpoint{3.030093in}{2.173242in}}%
\pgfpathlineto{\pgfqpoint{3.030953in}{2.159093in}}%
\pgfpathlineto{\pgfqpoint{3.032115in}{2.160201in}}%
\pgfpathlineto{\pgfqpoint{3.033229in}{2.174983in}}%
\pgfpathlineto{\pgfqpoint{3.036111in}{2.157763in}}%
\pgfpathlineto{\pgfqpoint{3.037358in}{2.173623in}}%
\pgfpathlineto{\pgfqpoint{3.041012in}{2.158753in}}%
\pgfpathlineto{\pgfqpoint{3.043199in}{2.175079in}}%
\pgfpathlineto{\pgfqpoint{3.047036in}{2.157148in}}%
\pgfpathlineto{\pgfqpoint{3.050624in}{2.171261in}}%
\pgfpathlineto{\pgfqpoint{3.054673in}{2.163554in}}%
\pgfpathlineto{\pgfqpoint{3.059489in}{2.164316in}}%
\pgfpathlineto{\pgfqpoint{3.064385in}{2.169932in}}%
\pgfpathlineto{\pgfqpoint{3.069281in}{2.158639in}}%
\pgfpathlineto{\pgfqpoint{3.074177in}{2.174084in}}%
\pgfpathlineto{\pgfqpoint{3.079073in}{2.156299in}}%
\pgfpathlineto{\pgfqpoint{3.083969in}{2.174083in}}%
\pgfpathlineto{\pgfqpoint{3.088865in}{2.158333in}}%
\pgfpathlineto{\pgfqpoint{3.093761in}{2.169941in}}%
\pgfpathlineto{\pgfqpoint{3.098657in}{2.163665in}}%
\pgfpathlineto{\pgfqpoint{3.103553in}{2.163626in}}%
\pgfpathlineto{\pgfqpoint{3.108449in}{2.169653in}}%
\pgfpathlineto{\pgfqpoint{3.113345in}{2.158123in}}%
\pgfpathlineto{\pgfqpoint{3.118241in}{2.173310in}}%
\pgfpathlineto{\pgfqpoint{3.123137in}{2.155954in}}%
\pgfpathlineto{\pgfqpoint{3.128033in}{2.172600in}}%
\pgfpathlineto{\pgfqpoint{3.132929in}{2.157914in}}%
\pgfpathlineto{\pgfqpoint{3.142721in}{2.087996in}}%
\pgfpathlineto{\pgfqpoint{3.149394in}{1.988418in}}%
\pgfpathlineto{\pgfqpoint{3.149394in}{1.988418in}}%
\pgfusepath{stroke}%
\end{pgfscope}%
\begin{pgfscope}%
\pgfsetrectcap%
\pgfsetmiterjoin%
\pgfsetlinewidth{0.803000pt}%
\definecolor{currentstroke}{rgb}{0.000000,0.000000,0.000000}%
\pgfsetstrokecolor{currentstroke}%
\pgfsetdash{}{0pt}%
\pgfpathmoveto{\pgfqpoint{2.197902in}{1.995363in}}%
\pgfpathlineto{\pgfqpoint{2.197902in}{2.321613in}}%
\pgfusepath{stroke}%
\end{pgfscope}%
\begin{pgfscope}%
\pgfsetrectcap%
\pgfsetmiterjoin%
\pgfsetlinewidth{0.803000pt}%
\definecolor{currentstroke}{rgb}{0.000000,0.000000,0.000000}%
\pgfsetstrokecolor{currentstroke}%
\pgfsetdash{}{0pt}%
\pgfpathmoveto{\pgfqpoint{3.177102in}{1.995363in}}%
\pgfpathlineto{\pgfqpoint{3.177102in}{2.321613in}}%
\pgfusepath{stroke}%
\end{pgfscope}%
\begin{pgfscope}%
\pgfsetrectcap%
\pgfsetmiterjoin%
\pgfsetlinewidth{0.803000pt}%
\definecolor{currentstroke}{rgb}{0.000000,0.000000,0.000000}%
\pgfsetstrokecolor{currentstroke}%
\pgfsetdash{}{0pt}%
\pgfpathmoveto{\pgfqpoint{2.197902in}{1.995363in}}%
\pgfpathlineto{\pgfqpoint{3.177102in}{1.995363in}}%
\pgfusepath{stroke}%
\end{pgfscope}%
\begin{pgfscope}%
\pgfsetrectcap%
\pgfsetmiterjoin%
\pgfsetlinewidth{0.803000pt}%
\definecolor{currentstroke}{rgb}{0.000000,0.000000,0.000000}%
\pgfsetstrokecolor{currentstroke}%
\pgfsetdash{}{0pt}%
\pgfpathmoveto{\pgfqpoint{2.197902in}{2.321613in}}%
\pgfpathlineto{\pgfqpoint{3.177102in}{2.321613in}}%
\pgfusepath{stroke}%
\end{pgfscope}%
\begin{pgfscope}%
\pgfsetbuttcap%
\pgfsetmiterjoin%
\definecolor{currentfill}{rgb}{1.000000,1.000000,1.000000}%
\pgfsetfillcolor{currentfill}%
\pgfsetlinewidth{0.000000pt}%
\definecolor{currentstroke}{rgb}{0.000000,0.000000,0.000000}%
\pgfsetstrokecolor{currentstroke}%
\pgfsetstrokeopacity{0.000000}%
\pgfsetdash{}{0pt}%
\pgfpathmoveto{\pgfqpoint{2.197902in}{0.579321in}}%
\pgfpathlineto{\pgfqpoint{3.177102in}{0.579321in}}%
\pgfpathlineto{\pgfqpoint{3.177102in}{0.905571in}}%
\pgfpathlineto{\pgfqpoint{2.197902in}{0.905571in}}%
\pgfpathlineto{\pgfqpoint{2.197902in}{0.579321in}}%
\pgfpathclose%
\pgfusepath{fill}%
\end{pgfscope}%
\begin{pgfscope}%
\pgfpathrectangle{\pgfqpoint{2.197902in}{0.579321in}}{\pgfqpoint{0.979200in}{0.326250in}}%
\pgfusepath{clip}%
\pgfsetbuttcap%
\pgfsetroundjoin%
\pgfsetlinewidth{0.401500pt}%
\definecolor{currentstroke}{rgb}{0.690196,0.690196,0.690196}%
\pgfsetstrokecolor{currentstroke}%
\pgfsetstrokeopacity{0.250000}%
\pgfsetdash{{1.480000pt}{0.640000pt}}{0.000000pt}%
\pgfpathmoveto{\pgfqpoint{2.197902in}{0.579321in}}%
\pgfpathlineto{\pgfqpoint{2.197902in}{0.905571in}}%
\pgfusepath{stroke}%
\end{pgfscope}%
\begin{pgfscope}%
\pgfsetbuttcap%
\pgfsetroundjoin%
\definecolor{currentfill}{rgb}{0.000000,0.000000,0.000000}%
\pgfsetfillcolor{currentfill}%
\pgfsetlinewidth{0.803000pt}%
\definecolor{currentstroke}{rgb}{0.000000,0.000000,0.000000}%
\pgfsetstrokecolor{currentstroke}%
\pgfsetdash{}{0pt}%
\pgfsys@defobject{currentmarker}{\pgfqpoint{0.000000in}{0.000000in}}{\pgfqpoint{0.000000in}{0.027778in}}{%
\pgfpathmoveto{\pgfqpoint{0.000000in}{0.000000in}}%
\pgfpathlineto{\pgfqpoint{0.000000in}{0.027778in}}%
\pgfusepath{stroke,fill}%
}%
\begin{pgfscope}%
\pgfsys@transformshift{2.197902in}{0.579321in}%
\pgfsys@useobject{currentmarker}{}%
\end{pgfscope}%
\end{pgfscope}%
\begin{pgfscope}%
\definecolor{textcolor}{rgb}{0.000000,0.000000,0.000000}%
\pgfsetstrokecolor{textcolor}%
\pgfsetfillcolor{textcolor}%
\pgftext[x=2.197902in,y=0.565432in,,top]{\color{textcolor}{\rmfamily\fontsize{5.000000}{6.000000}\selectfont\catcode`\^=\active\def^{\ifmmode\sp\else\^{}\fi}\catcode`\%=\active\def%{\%}4.980}}%
\end{pgfscope}%
\begin{pgfscope}%
\pgfpathrectangle{\pgfqpoint{2.197902in}{0.579321in}}{\pgfqpoint{0.979200in}{0.326250in}}%
\pgfusepath{clip}%
\pgfsetbuttcap%
\pgfsetroundjoin%
\pgfsetlinewidth{0.401500pt}%
\definecolor{currentstroke}{rgb}{0.690196,0.690196,0.690196}%
\pgfsetstrokecolor{currentstroke}%
\pgfsetstrokeopacity{0.250000}%
\pgfsetdash{{1.480000pt}{0.640000pt}}{0.000000pt}%
\pgfpathmoveto{\pgfqpoint{2.442702in}{0.579321in}}%
\pgfpathlineto{\pgfqpoint{2.442702in}{0.905571in}}%
\pgfusepath{stroke}%
\end{pgfscope}%
\begin{pgfscope}%
\pgfsetbuttcap%
\pgfsetroundjoin%
\definecolor{currentfill}{rgb}{0.000000,0.000000,0.000000}%
\pgfsetfillcolor{currentfill}%
\pgfsetlinewidth{0.803000pt}%
\definecolor{currentstroke}{rgb}{0.000000,0.000000,0.000000}%
\pgfsetstrokecolor{currentstroke}%
\pgfsetdash{}{0pt}%
\pgfsys@defobject{currentmarker}{\pgfqpoint{0.000000in}{0.000000in}}{\pgfqpoint{0.000000in}{0.027778in}}{%
\pgfpathmoveto{\pgfqpoint{0.000000in}{0.000000in}}%
\pgfpathlineto{\pgfqpoint{0.000000in}{0.027778in}}%
\pgfusepath{stroke,fill}%
}%
\begin{pgfscope}%
\pgfsys@transformshift{2.442702in}{0.579321in}%
\pgfsys@useobject{currentmarker}{}%
\end{pgfscope}%
\end{pgfscope}%
\begin{pgfscope}%
\definecolor{textcolor}{rgb}{0.000000,0.000000,0.000000}%
\pgfsetstrokecolor{textcolor}%
\pgfsetfillcolor{textcolor}%
\pgftext[x=2.442702in,y=0.565432in,,top]{\color{textcolor}{\rmfamily\fontsize{5.000000}{6.000000}\selectfont\catcode`\^=\active\def^{\ifmmode\sp\else\^{}\fi}\catcode`\%=\active\def%{\%}4.985}}%
\end{pgfscope}%
\begin{pgfscope}%
\pgfpathrectangle{\pgfqpoint{2.197902in}{0.579321in}}{\pgfqpoint{0.979200in}{0.326250in}}%
\pgfusepath{clip}%
\pgfsetbuttcap%
\pgfsetroundjoin%
\pgfsetlinewidth{0.401500pt}%
\definecolor{currentstroke}{rgb}{0.690196,0.690196,0.690196}%
\pgfsetstrokecolor{currentstroke}%
\pgfsetstrokeopacity{0.250000}%
\pgfsetdash{{1.480000pt}{0.640000pt}}{0.000000pt}%
\pgfpathmoveto{\pgfqpoint{2.687502in}{0.579321in}}%
\pgfpathlineto{\pgfqpoint{2.687502in}{0.905571in}}%
\pgfusepath{stroke}%
\end{pgfscope}%
\begin{pgfscope}%
\pgfsetbuttcap%
\pgfsetroundjoin%
\definecolor{currentfill}{rgb}{0.000000,0.000000,0.000000}%
\pgfsetfillcolor{currentfill}%
\pgfsetlinewidth{0.803000pt}%
\definecolor{currentstroke}{rgb}{0.000000,0.000000,0.000000}%
\pgfsetstrokecolor{currentstroke}%
\pgfsetdash{}{0pt}%
\pgfsys@defobject{currentmarker}{\pgfqpoint{0.000000in}{0.000000in}}{\pgfqpoint{0.000000in}{0.027778in}}{%
\pgfpathmoveto{\pgfqpoint{0.000000in}{0.000000in}}%
\pgfpathlineto{\pgfqpoint{0.000000in}{0.027778in}}%
\pgfusepath{stroke,fill}%
}%
\begin{pgfscope}%
\pgfsys@transformshift{2.687502in}{0.579321in}%
\pgfsys@useobject{currentmarker}{}%
\end{pgfscope}%
\end{pgfscope}%
\begin{pgfscope}%
\definecolor{textcolor}{rgb}{0.000000,0.000000,0.000000}%
\pgfsetstrokecolor{textcolor}%
\pgfsetfillcolor{textcolor}%
\pgftext[x=2.687502in,y=0.565432in,,top]{\color{textcolor}{\rmfamily\fontsize{5.000000}{6.000000}\selectfont\catcode`\^=\active\def^{\ifmmode\sp\else\^{}\fi}\catcode`\%=\active\def%{\%}4.990}}%
\end{pgfscope}%
\begin{pgfscope}%
\pgfpathrectangle{\pgfqpoint{2.197902in}{0.579321in}}{\pgfqpoint{0.979200in}{0.326250in}}%
\pgfusepath{clip}%
\pgfsetbuttcap%
\pgfsetroundjoin%
\pgfsetlinewidth{0.401500pt}%
\definecolor{currentstroke}{rgb}{0.690196,0.690196,0.690196}%
\pgfsetstrokecolor{currentstroke}%
\pgfsetstrokeopacity{0.250000}%
\pgfsetdash{{1.480000pt}{0.640000pt}}{0.000000pt}%
\pgfpathmoveto{\pgfqpoint{2.932302in}{0.579321in}}%
\pgfpathlineto{\pgfqpoint{2.932302in}{0.905571in}}%
\pgfusepath{stroke}%
\end{pgfscope}%
\begin{pgfscope}%
\pgfsetbuttcap%
\pgfsetroundjoin%
\definecolor{currentfill}{rgb}{0.000000,0.000000,0.000000}%
\pgfsetfillcolor{currentfill}%
\pgfsetlinewidth{0.803000pt}%
\definecolor{currentstroke}{rgb}{0.000000,0.000000,0.000000}%
\pgfsetstrokecolor{currentstroke}%
\pgfsetdash{}{0pt}%
\pgfsys@defobject{currentmarker}{\pgfqpoint{0.000000in}{0.000000in}}{\pgfqpoint{0.000000in}{0.027778in}}{%
\pgfpathmoveto{\pgfqpoint{0.000000in}{0.000000in}}%
\pgfpathlineto{\pgfqpoint{0.000000in}{0.027778in}}%
\pgfusepath{stroke,fill}%
}%
\begin{pgfscope}%
\pgfsys@transformshift{2.932302in}{0.579321in}%
\pgfsys@useobject{currentmarker}{}%
\end{pgfscope}%
\end{pgfscope}%
\begin{pgfscope}%
\definecolor{textcolor}{rgb}{0.000000,0.000000,0.000000}%
\pgfsetstrokecolor{textcolor}%
\pgfsetfillcolor{textcolor}%
\pgftext[x=2.932302in,y=0.565432in,,top]{\color{textcolor}{\rmfamily\fontsize{5.000000}{6.000000}\selectfont\catcode`\^=\active\def^{\ifmmode\sp\else\^{}\fi}\catcode`\%=\active\def%{\%}4.995}}%
\end{pgfscope}%
\begin{pgfscope}%
\pgfpathrectangle{\pgfqpoint{2.197902in}{0.579321in}}{\pgfqpoint{0.979200in}{0.326250in}}%
\pgfusepath{clip}%
\pgfsetbuttcap%
\pgfsetroundjoin%
\pgfsetlinewidth{0.401500pt}%
\definecolor{currentstroke}{rgb}{0.690196,0.690196,0.690196}%
\pgfsetstrokecolor{currentstroke}%
\pgfsetstrokeopacity{0.250000}%
\pgfsetdash{{1.480000pt}{0.640000pt}}{0.000000pt}%
\pgfpathmoveto{\pgfqpoint{3.177102in}{0.579321in}}%
\pgfpathlineto{\pgfqpoint{3.177102in}{0.905571in}}%
\pgfusepath{stroke}%
\end{pgfscope}%
\begin{pgfscope}%
\pgfsetbuttcap%
\pgfsetroundjoin%
\definecolor{currentfill}{rgb}{0.000000,0.000000,0.000000}%
\pgfsetfillcolor{currentfill}%
\pgfsetlinewidth{0.803000pt}%
\definecolor{currentstroke}{rgb}{0.000000,0.000000,0.000000}%
\pgfsetstrokecolor{currentstroke}%
\pgfsetdash{}{0pt}%
\pgfsys@defobject{currentmarker}{\pgfqpoint{0.000000in}{0.000000in}}{\pgfqpoint{0.000000in}{0.027778in}}{%
\pgfpathmoveto{\pgfqpoint{0.000000in}{0.000000in}}%
\pgfpathlineto{\pgfqpoint{0.000000in}{0.027778in}}%
\pgfusepath{stroke,fill}%
}%
\begin{pgfscope}%
\pgfsys@transformshift{3.177102in}{0.579321in}%
\pgfsys@useobject{currentmarker}{}%
\end{pgfscope}%
\end{pgfscope}%
\begin{pgfscope}%
\definecolor{textcolor}{rgb}{0.000000,0.000000,0.000000}%
\pgfsetstrokecolor{textcolor}%
\pgfsetfillcolor{textcolor}%
\pgftext[x=3.177102in,y=0.565432in,,top]{\color{textcolor}{\rmfamily\fontsize{5.000000}{6.000000}\selectfont\catcode`\^=\active\def^{\ifmmode\sp\else\^{}\fi}\catcode`\%=\active\def%{\%}5.000}}%
\end{pgfscope}%
\begin{pgfscope}%
\pgfpathrectangle{\pgfqpoint{2.197902in}{0.579321in}}{\pgfqpoint{0.979200in}{0.326250in}}%
\pgfusepath{clip}%
\pgfsetbuttcap%
\pgfsetroundjoin%
\pgfsetlinewidth{0.401500pt}%
\definecolor{currentstroke}{rgb}{0.690196,0.690196,0.690196}%
\pgfsetstrokecolor{currentstroke}%
\pgfsetstrokeopacity{0.250000}%
\pgfsetdash{{1.480000pt}{0.640000pt}}{0.000000pt}%
\pgfpathmoveto{\pgfqpoint{2.197902in}{0.579321in}}%
\pgfpathlineto{\pgfqpoint{3.177102in}{0.579321in}}%
\pgfusepath{stroke}%
\end{pgfscope}%
\begin{pgfscope}%
\pgfsetbuttcap%
\pgfsetroundjoin%
\definecolor{currentfill}{rgb}{0.000000,0.000000,0.000000}%
\pgfsetfillcolor{currentfill}%
\pgfsetlinewidth{0.803000pt}%
\definecolor{currentstroke}{rgb}{0.000000,0.000000,0.000000}%
\pgfsetstrokecolor{currentstroke}%
\pgfsetdash{}{0pt}%
\pgfsys@defobject{currentmarker}{\pgfqpoint{0.000000in}{0.000000in}}{\pgfqpoint{0.027778in}{0.000000in}}{%
\pgfpathmoveto{\pgfqpoint{0.000000in}{0.000000in}}%
\pgfpathlineto{\pgfqpoint{0.027778in}{0.000000in}}%
\pgfusepath{stroke,fill}%
}%
\begin{pgfscope}%
\pgfsys@transformshift{2.197902in}{0.579321in}%
\pgfsys@useobject{currentmarker}{}%
\end{pgfscope}%
\end{pgfscope}%
\begin{pgfscope}%
\definecolor{textcolor}{rgb}{0.000000,0.000000,0.000000}%
\pgfsetstrokecolor{textcolor}%
\pgfsetfillcolor{textcolor}%
\pgftext[x=1.966996in, y=0.555209in, left, base]{\color{textcolor}{\rmfamily\fontsize{5.000000}{6.000000}\selectfont\catcode`\^=\active\def^{\ifmmode\sp\else\^{}\fi}\catcode`\%=\active\def%{\%}14.50}}%
\end{pgfscope}%
\begin{pgfscope}%
\pgfpathrectangle{\pgfqpoint{2.197902in}{0.579321in}}{\pgfqpoint{0.979200in}{0.326250in}}%
\pgfusepath{clip}%
\pgfsetbuttcap%
\pgfsetroundjoin%
\pgfsetlinewidth{0.401500pt}%
\definecolor{currentstroke}{rgb}{0.690196,0.690196,0.690196}%
\pgfsetstrokecolor{currentstroke}%
\pgfsetstrokeopacity{0.250000}%
\pgfsetdash{{1.480000pt}{0.640000pt}}{0.000000pt}%
\pgfpathmoveto{\pgfqpoint{2.197902in}{0.742446in}}%
\pgfpathlineto{\pgfqpoint{3.177102in}{0.742446in}}%
\pgfusepath{stroke}%
\end{pgfscope}%
\begin{pgfscope}%
\pgfsetbuttcap%
\pgfsetroundjoin%
\definecolor{currentfill}{rgb}{0.000000,0.000000,0.000000}%
\pgfsetfillcolor{currentfill}%
\pgfsetlinewidth{0.803000pt}%
\definecolor{currentstroke}{rgb}{0.000000,0.000000,0.000000}%
\pgfsetstrokecolor{currentstroke}%
\pgfsetdash{}{0pt}%
\pgfsys@defobject{currentmarker}{\pgfqpoint{0.000000in}{0.000000in}}{\pgfqpoint{0.027778in}{0.000000in}}{%
\pgfpathmoveto{\pgfqpoint{0.000000in}{0.000000in}}%
\pgfpathlineto{\pgfqpoint{0.027778in}{0.000000in}}%
\pgfusepath{stroke,fill}%
}%
\begin{pgfscope}%
\pgfsys@transformshift{2.197902in}{0.742446in}%
\pgfsys@useobject{currentmarker}{}%
\end{pgfscope}%
\end{pgfscope}%
\begin{pgfscope}%
\definecolor{textcolor}{rgb}{0.000000,0.000000,0.000000}%
\pgfsetstrokecolor{textcolor}%
\pgfsetfillcolor{textcolor}%
\pgftext[x=1.966996in, y=0.718334in, left, base]{\color{textcolor}{\rmfamily\fontsize{5.000000}{6.000000}\selectfont\catcode`\^=\active\def^{\ifmmode\sp\else\^{}\fi}\catcode`\%=\active\def%{\%}14.55}}%
\end{pgfscope}%
\begin{pgfscope}%
\pgfpathrectangle{\pgfqpoint{2.197902in}{0.579321in}}{\pgfqpoint{0.979200in}{0.326250in}}%
\pgfusepath{clip}%
\pgfsetbuttcap%
\pgfsetroundjoin%
\pgfsetlinewidth{0.401500pt}%
\definecolor{currentstroke}{rgb}{0.690196,0.690196,0.690196}%
\pgfsetstrokecolor{currentstroke}%
\pgfsetstrokeopacity{0.250000}%
\pgfsetdash{{1.480000pt}{0.640000pt}}{0.000000pt}%
\pgfpathmoveto{\pgfqpoint{2.197902in}{0.905571in}}%
\pgfpathlineto{\pgfqpoint{3.177102in}{0.905571in}}%
\pgfusepath{stroke}%
\end{pgfscope}%
\begin{pgfscope}%
\pgfsetbuttcap%
\pgfsetroundjoin%
\definecolor{currentfill}{rgb}{0.000000,0.000000,0.000000}%
\pgfsetfillcolor{currentfill}%
\pgfsetlinewidth{0.803000pt}%
\definecolor{currentstroke}{rgb}{0.000000,0.000000,0.000000}%
\pgfsetstrokecolor{currentstroke}%
\pgfsetdash{}{0pt}%
\pgfsys@defobject{currentmarker}{\pgfqpoint{0.000000in}{0.000000in}}{\pgfqpoint{0.027778in}{0.000000in}}{%
\pgfpathmoveto{\pgfqpoint{0.000000in}{0.000000in}}%
\pgfpathlineto{\pgfqpoint{0.027778in}{0.000000in}}%
\pgfusepath{stroke,fill}%
}%
\begin{pgfscope}%
\pgfsys@transformshift{2.197902in}{0.905571in}%
\pgfsys@useobject{currentmarker}{}%
\end{pgfscope}%
\end{pgfscope}%
\begin{pgfscope}%
\definecolor{textcolor}{rgb}{0.000000,0.000000,0.000000}%
\pgfsetstrokecolor{textcolor}%
\pgfsetfillcolor{textcolor}%
\pgftext[x=1.966996in, y=0.881459in, left, base]{\color{textcolor}{\rmfamily\fontsize{5.000000}{6.000000}\selectfont\catcode`\^=\active\def^{\ifmmode\sp\else\^{}\fi}\catcode`\%=\active\def%{\%}14.60}}%
\end{pgfscope}%
\begin{pgfscope}%
\pgfpathrectangle{\pgfqpoint{2.197902in}{0.579321in}}{\pgfqpoint{0.979200in}{0.326250in}}%
\pgfusepath{clip}%
\pgfsetrectcap%
\pgfsetroundjoin%
\pgfsetlinewidth{1.505625pt}%
\definecolor{currentstroke}{rgb}{0.000000,0.619608,0.450980}%
\pgfsetstrokecolor{currentstroke}%
\pgfsetdash{}{0pt}%
\pgfpathmoveto{\pgfqpoint{2.194249in}{0.809077in}}%
\pgfpathlineto{\pgfqpoint{2.402983in}{0.606629in}}%
\pgfpathlineto{\pgfqpoint{2.404735in}{0.607061in}}%
\pgfpathlineto{\pgfqpoint{2.405893in}{0.609043in}}%
\pgfpathlineto{\pgfqpoint{2.433330in}{0.684036in}}%
\pgfpathlineto{\pgfqpoint{2.443967in}{0.709783in}}%
\pgfpathlineto{\pgfqpoint{2.464911in}{0.755338in}}%
\pgfpathlineto{\pgfqpoint{2.488369in}{0.797691in}}%
\pgfpathlineto{\pgfqpoint{2.500547in}{0.816257in}}%
\pgfpathlineto{\pgfqpoint{2.516043in}{0.836199in}}%
\pgfpathlineto{\pgfqpoint{2.529420in}{0.850354in}}%
\pgfpathlineto{\pgfqpoint{2.544140in}{0.862498in}}%
\pgfpathlineto{\pgfqpoint{2.556968in}{0.870144in}}%
\pgfpathlineto{\pgfqpoint{2.566760in}{0.874228in}}%
\pgfpathlineto{\pgfqpoint{2.576552in}{0.876750in}}%
\pgfpathlineto{\pgfqpoint{2.586344in}{0.877701in}}%
\pgfpathlineto{\pgfqpoint{2.596136in}{0.877073in}}%
\pgfpathlineto{\pgfqpoint{2.605928in}{0.874867in}}%
\pgfpathlineto{\pgfqpoint{2.615720in}{0.871089in}}%
\pgfpathlineto{\pgfqpoint{2.625512in}{0.865752in}}%
\pgfpathlineto{\pgfqpoint{2.635304in}{0.858865in}}%
\pgfpathlineto{\pgfqpoint{2.645096in}{0.850452in}}%
\pgfpathlineto{\pgfqpoint{2.893255in}{0.609993in}}%
\pgfpathlineto{\pgfqpoint{2.894171in}{0.610313in}}%
\pgfpathlineto{\pgfqpoint{2.895638in}{0.612608in}}%
\pgfpathlineto{\pgfqpoint{2.910656in}{0.655331in}}%
\pgfpathlineto{\pgfqpoint{2.943747in}{0.734507in}}%
\pgfpathlineto{\pgfqpoint{2.971876in}{0.788996in}}%
\pgfpathlineto{\pgfqpoint{2.979867in}{0.802341in}}%
\pgfpathlineto{\pgfqpoint{3.005851in}{0.837730in}}%
\pgfpathlineto{\pgfqpoint{3.014771in}{0.847333in}}%
\pgfpathlineto{\pgfqpoint{3.030953in}{0.861550in}}%
\pgfpathlineto{\pgfqpoint{3.041012in}{0.868196in}}%
\pgfpathlineto{\pgfqpoint{3.054673in}{0.874537in}}%
\pgfpathlineto{\pgfqpoint{3.064385in}{0.877173in}}%
\pgfpathlineto{\pgfqpoint{3.074177in}{0.878287in}}%
\pgfpathlineto{\pgfqpoint{3.083969in}{0.877862in}}%
\pgfpathlineto{\pgfqpoint{3.093761in}{0.875899in}}%
\pgfpathlineto{\pgfqpoint{3.103553in}{0.872387in}}%
\pgfpathlineto{\pgfqpoint{3.113345in}{0.867307in}}%
\pgfpathlineto{\pgfqpoint{3.123137in}{0.860645in}}%
\pgfpathlineto{\pgfqpoint{3.137825in}{0.847964in}}%
\pgfpathlineto{\pgfqpoint{3.177102in}{0.809690in}}%
\pgfpathlineto{\pgfqpoint{3.177102in}{0.809690in}}%
\pgfusepath{stroke}%
\end{pgfscope}%
\begin{pgfscope}%
\pgfsetrectcap%
\pgfsetmiterjoin%
\pgfsetlinewidth{0.803000pt}%
\definecolor{currentstroke}{rgb}{0.000000,0.000000,0.000000}%
\pgfsetstrokecolor{currentstroke}%
\pgfsetdash{}{0pt}%
\pgfpathmoveto{\pgfqpoint{2.197902in}{0.579321in}}%
\pgfpathlineto{\pgfqpoint{2.197902in}{0.905571in}}%
\pgfusepath{stroke}%
\end{pgfscope}%
\begin{pgfscope}%
\pgfsetrectcap%
\pgfsetmiterjoin%
\pgfsetlinewidth{0.803000pt}%
\definecolor{currentstroke}{rgb}{0.000000,0.000000,0.000000}%
\pgfsetstrokecolor{currentstroke}%
\pgfsetdash{}{0pt}%
\pgfpathmoveto{\pgfqpoint{3.177102in}{0.579321in}}%
\pgfpathlineto{\pgfqpoint{3.177102in}{0.905571in}}%
\pgfusepath{stroke}%
\end{pgfscope}%
\begin{pgfscope}%
\pgfsetrectcap%
\pgfsetmiterjoin%
\pgfsetlinewidth{0.803000pt}%
\definecolor{currentstroke}{rgb}{0.000000,0.000000,0.000000}%
\pgfsetstrokecolor{currentstroke}%
\pgfsetdash{}{0pt}%
\pgfpathmoveto{\pgfqpoint{2.197902in}{0.579321in}}%
\pgfpathlineto{\pgfqpoint{3.177102in}{0.579321in}}%
\pgfusepath{stroke}%
\end{pgfscope}%
\begin{pgfscope}%
\pgfsetrectcap%
\pgfsetmiterjoin%
\pgfsetlinewidth{0.803000pt}%
\definecolor{currentstroke}{rgb}{0.000000,0.000000,0.000000}%
\pgfsetstrokecolor{currentstroke}%
\pgfsetdash{}{0pt}%
\pgfpathmoveto{\pgfqpoint{2.197902in}{0.905571in}}%
\pgfpathlineto{\pgfqpoint{3.177102in}{0.905571in}}%
\pgfusepath{stroke}%
\end{pgfscope}%
\end{pgfpicture}%
\makeatother%
\endgroup%

	\caption{(a) corrente no dreno do transistor, (b) tensão dreno--source do transistor e (c) tensão de saída do circuito Flyback para $K~L1~L2~0.99$ e circuito \textit{clamper} RCD.}
	\label{fig:ex4_3}
\end{figure}

\medskip
\textbf{d)} Projete agora o circuito de \textit{snubber} e apresente novamente as formas de onda atualizadas.

\medskip
\textbf{Resolução}: para o projeto do circuito \textit{snubber}, considerou-se que a capacitância equivalente vista pelo dreno do MOSFET é aproximadamente igual a capacitância de saída do MOSFET $C_{oss}$. Para uma tensão $V_{DS} = 65$ V, a folha de dados do fabricante fornece $C_{oss} \approx 310$ pF~\cite{irfh5007_datasheet}.  Os valores obtidos para o circuito \textit{snubber} foram~\cite{aula6_isolados}: $R_{snb} = 80{,}1~\Omega$, $C_{snb} = 994{,}5$ pF, considerando uma perda $P_{loss,snb} = 0.5P_{loss,clamper} = 0{,}4176$ W.

A Figura~\ref{fig:ex4_4} apresenta as formas de onda solicitadas no enunciado. Pode-se observar a diminuição da amplitude da oscilação na tensão $V_{DS}$. Além disso, também é observada a redução da tensão de saída. como consequência das perdas $P_{loss,snb}$ e $P_{loss,clamper}$.
\begin{figure}[htb!]
	\centering
	%% Creator: Matplotlib, PGF backend
%%
%% To include the figure in your LaTeX document, write
%%   \input{<filename>.pgf}
%%
%% Make sure the required packages are loaded in your preamble
%%   \usepackage{pgf}
%%
%% Also ensure that all the required font packages are loaded; for instance,
%% the lmodern package is sometimes necessary when using math font.
%%   \usepackage{lmodern}
%%
%% Figures using additional raster images can only be included by \input if
%% they are in the same directory as the main LaTeX file. For loading figures
%% from other directories you can use the `import` package
%%   \usepackage{import}
%%
%% and then include the figures with
%%   \import{<path to file>}{<filename>.pgf}
%%
%% Matplotlib used the following preamble
%%   \def\mathdefault#1{#1}
%%   \everymath=\expandafter{\the\everymath\displaystyle}
%%   \IfFileExists{scrextend.sty}{
%%     \usepackage[fontsize=8.000000pt]{scrextend}
%%   }{
%%     \renewcommand{\normalsize}{\fontsize{8.000000}{9.600000}\selectfont}
%%     \normalsize
%%   }
%%   
%%           \usepackage{amsmath}
%%           \usepackage{amssymb}
%%           \usepackage{mathtools}
%%           \usepackage{dsfont}
%%           \providecommand{\mathdefault}[1]{#1}
%%       
%%   \makeatletter\@ifpackageloaded{underscore}{}{\usepackage[strings]{underscore}}\makeatother
%%
\begingroup%
\makeatletter%
\begin{pgfpicture}%
\pgfpathrectangle{\pgfpointorigin}{\pgfqpoint{3.393768in}{3.646235in}}%
\pgfusepath{use as bounding box, clip}%
\begin{pgfscope}%
\pgfsetbuttcap%
\pgfsetmiterjoin%
\definecolor{currentfill}{rgb}{1.000000,1.000000,1.000000}%
\pgfsetfillcolor{currentfill}%
\pgfsetlinewidth{0.000000pt}%
\definecolor{currentstroke}{rgb}{1.000000,1.000000,1.000000}%
\pgfsetstrokecolor{currentstroke}%
\pgfsetdash{}{0pt}%
\pgfpathmoveto{\pgfqpoint{0.000000in}{0.000000in}}%
\pgfpathlineto{\pgfqpoint{3.393768in}{0.000000in}}%
\pgfpathlineto{\pgfqpoint{3.393768in}{3.646235in}}%
\pgfpathlineto{\pgfqpoint{0.000000in}{3.646235in}}%
\pgfpathlineto{\pgfqpoint{0.000000in}{0.000000in}}%
\pgfpathclose%
\pgfusepath{fill}%
\end{pgfscope}%
\begin{pgfscope}%
\pgfsetbuttcap%
\pgfsetmiterjoin%
\definecolor{currentfill}{rgb}{1.000000,1.000000,1.000000}%
\pgfsetfillcolor{currentfill}%
\pgfsetlinewidth{0.000000pt}%
\definecolor{currentstroke}{rgb}{0.000000,0.000000,0.000000}%
\pgfsetstrokecolor{currentstroke}%
\pgfsetstrokeopacity{0.000000}%
\pgfsetdash{}{0pt}%
\pgfpathmoveto{\pgfqpoint{0.573768in}{2.601404in}}%
\pgfpathlineto{\pgfqpoint{3.293768in}{2.601404in}}%
\pgfpathlineto{\pgfqpoint{3.293768in}{3.507654in}}%
\pgfpathlineto{\pgfqpoint{0.573768in}{3.507654in}}%
\pgfpathlineto{\pgfqpoint{0.573768in}{2.601404in}}%
\pgfpathclose%
\pgfusepath{fill}%
\end{pgfscope}%
\begin{pgfscope}%
\pgfpathrectangle{\pgfqpoint{0.573768in}{2.601404in}}{\pgfqpoint{2.720000in}{0.906250in}}%
\pgfusepath{clip}%
\pgfsetbuttcap%
\pgfsetroundjoin%
\pgfsetlinewidth{0.501875pt}%
\definecolor{currentstroke}{rgb}{0.690196,0.690196,0.690196}%
\pgfsetstrokecolor{currentstroke}%
\pgfsetstrokeopacity{0.250000}%
\pgfsetdash{{1.850000pt}{0.800000pt}}{0.000000pt}%
\pgfpathmoveto{\pgfqpoint{0.697405in}{2.601404in}}%
\pgfpathlineto{\pgfqpoint{0.697405in}{3.507654in}}%
\pgfusepath{stroke}%
\end{pgfscope}%
\begin{pgfscope}%
\pgfsetbuttcap%
\pgfsetroundjoin%
\definecolor{currentfill}{rgb}{0.000000,0.000000,0.000000}%
\pgfsetfillcolor{currentfill}%
\pgfsetlinewidth{0.803000pt}%
\definecolor{currentstroke}{rgb}{0.000000,0.000000,0.000000}%
\pgfsetstrokecolor{currentstroke}%
\pgfsetdash{}{0pt}%
\pgfsys@defobject{currentmarker}{\pgfqpoint{0.000000in}{-0.048611in}}{\pgfqpoint{0.000000in}{0.000000in}}{%
\pgfpathmoveto{\pgfqpoint{0.000000in}{0.000000in}}%
\pgfpathlineto{\pgfqpoint{0.000000in}{-0.048611in}}%
\pgfusepath{stroke,fill}%
}%
\begin{pgfscope}%
\pgfsys@transformshift{0.697405in}{2.601404in}%
\pgfsys@useobject{currentmarker}{}%
\end{pgfscope}%
\end{pgfscope}%
\begin{pgfscope}%
\pgfpathrectangle{\pgfqpoint{0.573768in}{2.601404in}}{\pgfqpoint{2.720000in}{0.906250in}}%
\pgfusepath{clip}%
\pgfsetbuttcap%
\pgfsetroundjoin%
\pgfsetlinewidth{0.501875pt}%
\definecolor{currentstroke}{rgb}{0.690196,0.690196,0.690196}%
\pgfsetstrokecolor{currentstroke}%
\pgfsetstrokeopacity{0.250000}%
\pgfsetdash{{1.850000pt}{0.800000pt}}{0.000000pt}%
\pgfpathmoveto{\pgfqpoint{1.191950in}{2.601404in}}%
\pgfpathlineto{\pgfqpoint{1.191950in}{3.507654in}}%
\pgfusepath{stroke}%
\end{pgfscope}%
\begin{pgfscope}%
\pgfsetbuttcap%
\pgfsetroundjoin%
\definecolor{currentfill}{rgb}{0.000000,0.000000,0.000000}%
\pgfsetfillcolor{currentfill}%
\pgfsetlinewidth{0.803000pt}%
\definecolor{currentstroke}{rgb}{0.000000,0.000000,0.000000}%
\pgfsetstrokecolor{currentstroke}%
\pgfsetdash{}{0pt}%
\pgfsys@defobject{currentmarker}{\pgfqpoint{0.000000in}{-0.048611in}}{\pgfqpoint{0.000000in}{0.000000in}}{%
\pgfpathmoveto{\pgfqpoint{0.000000in}{0.000000in}}%
\pgfpathlineto{\pgfqpoint{0.000000in}{-0.048611in}}%
\pgfusepath{stroke,fill}%
}%
\begin{pgfscope}%
\pgfsys@transformshift{1.191950in}{2.601404in}%
\pgfsys@useobject{currentmarker}{}%
\end{pgfscope}%
\end{pgfscope}%
\begin{pgfscope}%
\pgfpathrectangle{\pgfqpoint{0.573768in}{2.601404in}}{\pgfqpoint{2.720000in}{0.906250in}}%
\pgfusepath{clip}%
\pgfsetbuttcap%
\pgfsetroundjoin%
\pgfsetlinewidth{0.501875pt}%
\definecolor{currentstroke}{rgb}{0.690196,0.690196,0.690196}%
\pgfsetstrokecolor{currentstroke}%
\pgfsetstrokeopacity{0.250000}%
\pgfsetdash{{1.850000pt}{0.800000pt}}{0.000000pt}%
\pgfpathmoveto{\pgfqpoint{1.686496in}{2.601404in}}%
\pgfpathlineto{\pgfqpoint{1.686496in}{3.507654in}}%
\pgfusepath{stroke}%
\end{pgfscope}%
\begin{pgfscope}%
\pgfsetbuttcap%
\pgfsetroundjoin%
\definecolor{currentfill}{rgb}{0.000000,0.000000,0.000000}%
\pgfsetfillcolor{currentfill}%
\pgfsetlinewidth{0.803000pt}%
\definecolor{currentstroke}{rgb}{0.000000,0.000000,0.000000}%
\pgfsetstrokecolor{currentstroke}%
\pgfsetdash{}{0pt}%
\pgfsys@defobject{currentmarker}{\pgfqpoint{0.000000in}{-0.048611in}}{\pgfqpoint{0.000000in}{0.000000in}}{%
\pgfpathmoveto{\pgfqpoint{0.000000in}{0.000000in}}%
\pgfpathlineto{\pgfqpoint{0.000000in}{-0.048611in}}%
\pgfusepath{stroke,fill}%
}%
\begin{pgfscope}%
\pgfsys@transformshift{1.686496in}{2.601404in}%
\pgfsys@useobject{currentmarker}{}%
\end{pgfscope}%
\end{pgfscope}%
\begin{pgfscope}%
\pgfpathrectangle{\pgfqpoint{0.573768in}{2.601404in}}{\pgfqpoint{2.720000in}{0.906250in}}%
\pgfusepath{clip}%
\pgfsetbuttcap%
\pgfsetroundjoin%
\pgfsetlinewidth{0.501875pt}%
\definecolor{currentstroke}{rgb}{0.690196,0.690196,0.690196}%
\pgfsetstrokecolor{currentstroke}%
\pgfsetstrokeopacity{0.250000}%
\pgfsetdash{{1.850000pt}{0.800000pt}}{0.000000pt}%
\pgfpathmoveto{\pgfqpoint{2.181041in}{2.601404in}}%
\pgfpathlineto{\pgfqpoint{2.181041in}{3.507654in}}%
\pgfusepath{stroke}%
\end{pgfscope}%
\begin{pgfscope}%
\pgfsetbuttcap%
\pgfsetroundjoin%
\definecolor{currentfill}{rgb}{0.000000,0.000000,0.000000}%
\pgfsetfillcolor{currentfill}%
\pgfsetlinewidth{0.803000pt}%
\definecolor{currentstroke}{rgb}{0.000000,0.000000,0.000000}%
\pgfsetstrokecolor{currentstroke}%
\pgfsetdash{}{0pt}%
\pgfsys@defobject{currentmarker}{\pgfqpoint{0.000000in}{-0.048611in}}{\pgfqpoint{0.000000in}{0.000000in}}{%
\pgfpathmoveto{\pgfqpoint{0.000000in}{0.000000in}}%
\pgfpathlineto{\pgfqpoint{0.000000in}{-0.048611in}}%
\pgfusepath{stroke,fill}%
}%
\begin{pgfscope}%
\pgfsys@transformshift{2.181041in}{2.601404in}%
\pgfsys@useobject{currentmarker}{}%
\end{pgfscope}%
\end{pgfscope}%
\begin{pgfscope}%
\pgfpathrectangle{\pgfqpoint{0.573768in}{2.601404in}}{\pgfqpoint{2.720000in}{0.906250in}}%
\pgfusepath{clip}%
\pgfsetbuttcap%
\pgfsetroundjoin%
\pgfsetlinewidth{0.501875pt}%
\definecolor{currentstroke}{rgb}{0.690196,0.690196,0.690196}%
\pgfsetstrokecolor{currentstroke}%
\pgfsetstrokeopacity{0.250000}%
\pgfsetdash{{1.850000pt}{0.800000pt}}{0.000000pt}%
\pgfpathmoveto{\pgfqpoint{2.675587in}{2.601404in}}%
\pgfpathlineto{\pgfqpoint{2.675587in}{3.507654in}}%
\pgfusepath{stroke}%
\end{pgfscope}%
\begin{pgfscope}%
\pgfsetbuttcap%
\pgfsetroundjoin%
\definecolor{currentfill}{rgb}{0.000000,0.000000,0.000000}%
\pgfsetfillcolor{currentfill}%
\pgfsetlinewidth{0.803000pt}%
\definecolor{currentstroke}{rgb}{0.000000,0.000000,0.000000}%
\pgfsetstrokecolor{currentstroke}%
\pgfsetdash{}{0pt}%
\pgfsys@defobject{currentmarker}{\pgfqpoint{0.000000in}{-0.048611in}}{\pgfqpoint{0.000000in}{0.000000in}}{%
\pgfpathmoveto{\pgfqpoint{0.000000in}{0.000000in}}%
\pgfpathlineto{\pgfqpoint{0.000000in}{-0.048611in}}%
\pgfusepath{stroke,fill}%
}%
\begin{pgfscope}%
\pgfsys@transformshift{2.675587in}{2.601404in}%
\pgfsys@useobject{currentmarker}{}%
\end{pgfscope}%
\end{pgfscope}%
\begin{pgfscope}%
\pgfpathrectangle{\pgfqpoint{0.573768in}{2.601404in}}{\pgfqpoint{2.720000in}{0.906250in}}%
\pgfusepath{clip}%
\pgfsetbuttcap%
\pgfsetroundjoin%
\pgfsetlinewidth{0.501875pt}%
\definecolor{currentstroke}{rgb}{0.690196,0.690196,0.690196}%
\pgfsetstrokecolor{currentstroke}%
\pgfsetstrokeopacity{0.250000}%
\pgfsetdash{{1.850000pt}{0.800000pt}}{0.000000pt}%
\pgfpathmoveto{\pgfqpoint{3.170132in}{2.601404in}}%
\pgfpathlineto{\pgfqpoint{3.170132in}{3.507654in}}%
\pgfusepath{stroke}%
\end{pgfscope}%
\begin{pgfscope}%
\pgfsetbuttcap%
\pgfsetroundjoin%
\definecolor{currentfill}{rgb}{0.000000,0.000000,0.000000}%
\pgfsetfillcolor{currentfill}%
\pgfsetlinewidth{0.803000pt}%
\definecolor{currentstroke}{rgb}{0.000000,0.000000,0.000000}%
\pgfsetstrokecolor{currentstroke}%
\pgfsetdash{}{0pt}%
\pgfsys@defobject{currentmarker}{\pgfqpoint{0.000000in}{-0.048611in}}{\pgfqpoint{0.000000in}{0.000000in}}{%
\pgfpathmoveto{\pgfqpoint{0.000000in}{0.000000in}}%
\pgfpathlineto{\pgfqpoint{0.000000in}{-0.048611in}}%
\pgfusepath{stroke,fill}%
}%
\begin{pgfscope}%
\pgfsys@transformshift{3.170132in}{2.601404in}%
\pgfsys@useobject{currentmarker}{}%
\end{pgfscope}%
\end{pgfscope}%
\begin{pgfscope}%
\pgfpathrectangle{\pgfqpoint{0.573768in}{2.601404in}}{\pgfqpoint{2.720000in}{0.906250in}}%
\pgfusepath{clip}%
\pgfsetbuttcap%
\pgfsetroundjoin%
\pgfsetlinewidth{0.501875pt}%
\definecolor{currentstroke}{rgb}{0.690196,0.690196,0.690196}%
\pgfsetstrokecolor{currentstroke}%
\pgfsetstrokeopacity{0.250000}%
\pgfsetdash{{1.850000pt}{0.800000pt}}{0.000000pt}%
\pgfpathmoveto{\pgfqpoint{0.573768in}{2.644559in}}%
\pgfpathlineto{\pgfqpoint{3.293768in}{2.644559in}}%
\pgfusepath{stroke}%
\end{pgfscope}%
\begin{pgfscope}%
\pgfsetbuttcap%
\pgfsetroundjoin%
\definecolor{currentfill}{rgb}{0.000000,0.000000,0.000000}%
\pgfsetfillcolor{currentfill}%
\pgfsetlinewidth{0.803000pt}%
\definecolor{currentstroke}{rgb}{0.000000,0.000000,0.000000}%
\pgfsetstrokecolor{currentstroke}%
\pgfsetdash{}{0pt}%
\pgfsys@defobject{currentmarker}{\pgfqpoint{-0.048611in}{0.000000in}}{\pgfqpoint{-0.000000in}{0.000000in}}{%
\pgfpathmoveto{\pgfqpoint{-0.000000in}{0.000000in}}%
\pgfpathlineto{\pgfqpoint{-0.048611in}{0.000000in}}%
\pgfusepath{stroke,fill}%
}%
\begin{pgfscope}%
\pgfsys@transformshift{0.573768in}{2.644559in}%
\pgfsys@useobject{currentmarker}{}%
\end{pgfscope}%
\end{pgfscope}%
\begin{pgfscope}%
\definecolor{textcolor}{rgb}{0.000000,0.000000,0.000000}%
\pgfsetstrokecolor{textcolor}%
\pgfsetfillcolor{textcolor}%
\pgftext[x=0.417518in, y=2.605979in, left, base]{\color{textcolor}{\rmfamily\fontsize{8.000000}{9.600000}\selectfont\catcode`\^=\active\def^{\ifmmode\sp\else\^{}\fi}\catcode`\%=\active\def%{\%}$\mathdefault{0}$}}%
\end{pgfscope}%
\begin{pgfscope}%
\pgfpathrectangle{\pgfqpoint{0.573768in}{2.601404in}}{\pgfqpoint{2.720000in}{0.906250in}}%
\pgfusepath{clip}%
\pgfsetbuttcap%
\pgfsetroundjoin%
\pgfsetlinewidth{0.501875pt}%
\definecolor{currentstroke}{rgb}{0.690196,0.690196,0.690196}%
\pgfsetstrokecolor{currentstroke}%
\pgfsetstrokeopacity{0.250000}%
\pgfsetdash{{1.850000pt}{0.800000pt}}{0.000000pt}%
\pgfpathmoveto{\pgfqpoint{0.573768in}{3.076107in}}%
\pgfpathlineto{\pgfqpoint{3.293768in}{3.076107in}}%
\pgfusepath{stroke}%
\end{pgfscope}%
\begin{pgfscope}%
\pgfsetbuttcap%
\pgfsetroundjoin%
\definecolor{currentfill}{rgb}{0.000000,0.000000,0.000000}%
\pgfsetfillcolor{currentfill}%
\pgfsetlinewidth{0.803000pt}%
\definecolor{currentstroke}{rgb}{0.000000,0.000000,0.000000}%
\pgfsetstrokecolor{currentstroke}%
\pgfsetdash{}{0pt}%
\pgfsys@defobject{currentmarker}{\pgfqpoint{-0.048611in}{0.000000in}}{\pgfqpoint{-0.000000in}{0.000000in}}{%
\pgfpathmoveto{\pgfqpoint{-0.000000in}{0.000000in}}%
\pgfpathlineto{\pgfqpoint{-0.048611in}{0.000000in}}%
\pgfusepath{stroke,fill}%
}%
\begin{pgfscope}%
\pgfsys@transformshift{0.573768in}{3.076107in}%
\pgfsys@useobject{currentmarker}{}%
\end{pgfscope}%
\end{pgfscope}%
\begin{pgfscope}%
\definecolor{textcolor}{rgb}{0.000000,0.000000,0.000000}%
\pgfsetstrokecolor{textcolor}%
\pgfsetfillcolor{textcolor}%
\pgftext[x=0.417518in, y=3.037527in, left, base]{\color{textcolor}{\rmfamily\fontsize{8.000000}{9.600000}\selectfont\catcode`\^=\active\def^{\ifmmode\sp\else\^{}\fi}\catcode`\%=\active\def%{\%}$\mathdefault{5}$}}%
\end{pgfscope}%
\begin{pgfscope}%
\pgfpathrectangle{\pgfqpoint{0.573768in}{2.601404in}}{\pgfqpoint{2.720000in}{0.906250in}}%
\pgfusepath{clip}%
\pgfsetbuttcap%
\pgfsetroundjoin%
\pgfsetlinewidth{0.501875pt}%
\definecolor{currentstroke}{rgb}{0.690196,0.690196,0.690196}%
\pgfsetstrokecolor{currentstroke}%
\pgfsetstrokeopacity{0.250000}%
\pgfsetdash{{1.850000pt}{0.800000pt}}{0.000000pt}%
\pgfpathmoveto{\pgfqpoint{0.573768in}{3.507654in}}%
\pgfpathlineto{\pgfqpoint{3.293768in}{3.507654in}}%
\pgfusepath{stroke}%
\end{pgfscope}%
\begin{pgfscope}%
\pgfsetbuttcap%
\pgfsetroundjoin%
\definecolor{currentfill}{rgb}{0.000000,0.000000,0.000000}%
\pgfsetfillcolor{currentfill}%
\pgfsetlinewidth{0.803000pt}%
\definecolor{currentstroke}{rgb}{0.000000,0.000000,0.000000}%
\pgfsetstrokecolor{currentstroke}%
\pgfsetdash{}{0pt}%
\pgfsys@defobject{currentmarker}{\pgfqpoint{-0.048611in}{0.000000in}}{\pgfqpoint{-0.000000in}{0.000000in}}{%
\pgfpathmoveto{\pgfqpoint{-0.000000in}{0.000000in}}%
\pgfpathlineto{\pgfqpoint{-0.048611in}{0.000000in}}%
\pgfusepath{stroke,fill}%
}%
\begin{pgfscope}%
\pgfsys@transformshift{0.573768in}{3.507654in}%
\pgfsys@useobject{currentmarker}{}%
\end{pgfscope}%
\end{pgfscope}%
\begin{pgfscope}%
\definecolor{textcolor}{rgb}{0.000000,0.000000,0.000000}%
\pgfsetstrokecolor{textcolor}%
\pgfsetfillcolor{textcolor}%
\pgftext[x=0.358489in, y=3.469074in, left, base]{\color{textcolor}{\rmfamily\fontsize{8.000000}{9.600000}\selectfont\catcode`\^=\active\def^{\ifmmode\sp\else\^{}\fi}\catcode`\%=\active\def%{\%}$\mathdefault{10}$}}%
\end{pgfscope}%
\begin{pgfscope}%
\definecolor{textcolor}{rgb}{0.000000,0.000000,0.000000}%
\pgfsetstrokecolor{textcolor}%
\pgfsetfillcolor{textcolor}%
\pgftext[x=0.302933in,y=3.054529in,,bottom,rotate=90.000000]{\color{textcolor}{\rmfamily\fontsize{8.000000}{9.600000}\selectfont\catcode`\^=\active\def^{\ifmmode\sp\else\^{}\fi}\catcode`\%=\active\def%{\%}$i_D$ (A)}}%
\end{pgfscope}%
\begin{pgfscope}%
\pgfpathrectangle{\pgfqpoint{0.573768in}{2.601404in}}{\pgfqpoint{2.720000in}{0.906250in}}%
\pgfusepath{clip}%
\pgfsetrectcap%
\pgfsetroundjoin%
\pgfsetlinewidth{1.505625pt}%
\definecolor{currentstroke}{rgb}{0.000000,0.447059,0.698039}%
\pgfsetstrokecolor{currentstroke}%
\pgfsetdash{}{0pt}%
\pgfpathmoveto{\pgfqpoint{0.697405in}{2.644559in}}%
\pgfpathlineto{\pgfqpoint{0.697411in}{2.644559in}}%
\pgfpathlineto{\pgfqpoint{0.697417in}{2.745750in}}%
\pgfpathlineto{\pgfqpoint{0.698159in}{2.691910in}}%
\pgfpathlineto{\pgfqpoint{0.699384in}{2.768830in}}%
\pgfpathlineto{\pgfqpoint{0.699411in}{2.755155in}}%
\pgfpathlineto{\pgfqpoint{0.700204in}{2.644559in}}%
\pgfpathlineto{\pgfqpoint{0.702357in}{2.644563in}}%
\pgfpathlineto{\pgfqpoint{0.702367in}{2.881803in}}%
\pgfpathlineto{\pgfqpoint{0.703092in}{2.821526in}}%
\pgfpathlineto{\pgfqpoint{0.703141in}{2.805483in}}%
\pgfpathlineto{\pgfqpoint{0.703685in}{2.858757in}}%
\pgfpathlineto{\pgfqpoint{0.703734in}{2.842801in}}%
\pgfpathlineto{\pgfqpoint{0.704329in}{2.899141in}}%
\pgfpathlineto{\pgfqpoint{0.704368in}{2.838275in}}%
\pgfpathlineto{\pgfqpoint{0.704379in}{2.845917in}}%
\pgfpathlineto{\pgfqpoint{0.705178in}{2.644559in}}%
\pgfpathlineto{\pgfqpoint{0.707302in}{2.644564in}}%
\pgfpathlineto{\pgfqpoint{0.707319in}{2.984090in}}%
\pgfpathlineto{\pgfqpoint{0.708053in}{2.903721in}}%
\pgfpathlineto{\pgfqpoint{0.708894in}{3.005454in}}%
\pgfpathlineto{\pgfqpoint{0.709274in}{3.029281in}}%
\pgfpathlineto{\pgfqpoint{0.710056in}{2.644559in}}%
\pgfpathlineto{\pgfqpoint{0.712248in}{2.644565in}}%
\pgfpathlineto{\pgfqpoint{0.712265in}{3.090648in}}%
\pgfpathlineto{\pgfqpoint{0.713030in}{3.017961in}}%
\pgfpathlineto{\pgfqpoint{0.713870in}{3.108595in}}%
\pgfpathlineto{\pgfqpoint{0.714220in}{3.130473in}}%
\pgfpathlineto{\pgfqpoint{0.715290in}{2.644559in}}%
\pgfpathlineto{\pgfqpoint{0.717193in}{2.644564in}}%
\pgfpathlineto{\pgfqpoint{0.717215in}{3.186425in}}%
\pgfpathlineto{\pgfqpoint{0.717934in}{3.102235in}}%
\pgfpathlineto{\pgfqpoint{0.718775in}{3.209828in}}%
\pgfpathlineto{\pgfqpoint{0.719165in}{3.234278in}}%
\pgfpathlineto{\pgfqpoint{0.720068in}{2.644559in}}%
\pgfpathlineto{\pgfqpoint{0.722138in}{2.644566in}}%
\pgfpathlineto{\pgfqpoint{0.722161in}{3.265290in}}%
\pgfpathlineto{\pgfqpoint{0.722920in}{3.221710in}}%
\pgfpathlineto{\pgfqpoint{0.722969in}{3.208830in}}%
\pgfpathlineto{\pgfqpoint{0.723513in}{3.258940in}}%
\pgfpathlineto{\pgfqpoint{0.723563in}{3.246099in}}%
\pgfpathlineto{\pgfqpoint{0.724084in}{3.294742in}}%
\pgfpathlineto{\pgfqpoint{0.724146in}{3.224683in}}%
\pgfpathlineto{\pgfqpoint{0.724156in}{3.261913in}}%
\pgfpathlineto{\pgfqpoint{0.724981in}{2.644559in}}%
\pgfpathlineto{\pgfqpoint{0.727084in}{2.644564in}}%
\pgfpathlineto{\pgfqpoint{0.727106in}{3.319723in}}%
\pgfpathlineto{\pgfqpoint{0.727844in}{3.299086in}}%
\pgfpathlineto{\pgfqpoint{0.727893in}{3.260129in}}%
\pgfpathlineto{\pgfqpoint{0.728536in}{3.342436in}}%
\pgfpathlineto{\pgfqpoint{0.728586in}{3.303685in}}%
\pgfpathlineto{\pgfqpoint{0.729019in}{3.372638in}}%
\pgfpathlineto{\pgfqpoint{0.729124in}{3.135771in}}%
\pgfpathlineto{\pgfqpoint{0.729872in}{2.644559in}}%
\pgfpathlineto{\pgfqpoint{0.732029in}{2.644566in}}%
\pgfpathlineto{\pgfqpoint{0.732055in}{3.377356in}}%
\pgfpathlineto{\pgfqpoint{0.732821in}{3.323666in}}%
\pgfpathlineto{\pgfqpoint{0.733860in}{3.397748in}}%
\pgfpathlineto{\pgfqpoint{0.733955in}{3.403755in}}%
\pgfpathlineto{\pgfqpoint{0.734070in}{3.287948in}}%
\pgfpathlineto{\pgfqpoint{0.734820in}{2.644559in}}%
\pgfpathlineto{\pgfqpoint{0.736975in}{2.644567in}}%
\pgfpathlineto{\pgfqpoint{0.736999in}{3.411802in}}%
\pgfpathlineto{\pgfqpoint{0.737725in}{3.345737in}}%
\pgfpathlineto{\pgfqpoint{0.738665in}{3.423399in}}%
\pgfpathlineto{\pgfqpoint{0.739003in}{3.320067in}}%
\pgfpathlineto{\pgfqpoint{0.738947in}{3.441080in}}%
\pgfpathlineto{\pgfqpoint{0.739015in}{3.334471in}}%
\pgfpathlineto{\pgfqpoint{0.739755in}{2.644559in}}%
\pgfpathlineto{\pgfqpoint{0.741920in}{2.644564in}}%
\pgfpathlineto{\pgfqpoint{0.741947in}{3.428402in}}%
\pgfpathlineto{\pgfqpoint{0.742674in}{3.357107in}}%
\pgfpathlineto{\pgfqpoint{0.743613in}{3.433274in}}%
\pgfpathlineto{\pgfqpoint{0.743948in}{3.330336in}}%
\pgfpathlineto{\pgfqpoint{0.743893in}{3.450787in}}%
\pgfpathlineto{\pgfqpoint{0.743961in}{3.362348in}}%
\pgfpathlineto{\pgfqpoint{0.744740in}{2.644559in}}%
\pgfpathlineto{\pgfqpoint{0.746866in}{2.644571in}}%
\pgfpathlineto{\pgfqpoint{0.746888in}{3.421775in}}%
\pgfpathlineto{\pgfqpoint{0.747644in}{3.342863in}}%
\pgfpathlineto{\pgfqpoint{0.748485in}{3.426921in}}%
\pgfpathlineto{\pgfqpoint{0.748892in}{3.335080in}}%
\pgfpathlineto{\pgfqpoint{0.748838in}{3.449015in}}%
\pgfpathlineto{\pgfqpoint{0.748905in}{3.356892in}}%
\pgfpathlineto{\pgfqpoint{0.749657in}{2.644559in}}%
\pgfpathlineto{\pgfqpoint{0.751811in}{2.644564in}}%
\pgfpathlineto{\pgfqpoint{0.751837in}{3.402945in}}%
\pgfpathlineto{\pgfqpoint{0.752574in}{3.322660in}}%
\pgfpathlineto{\pgfqpoint{0.753514in}{3.396455in}}%
\pgfpathlineto{\pgfqpoint{0.753839in}{3.294310in}}%
\pgfpathlineto{\pgfqpoint{0.753784in}{3.413360in}}%
\pgfpathlineto{\pgfqpoint{0.753852in}{3.340664in}}%
\pgfpathlineto{\pgfqpoint{0.754604in}{2.644559in}}%
\pgfpathlineto{\pgfqpoint{0.756757in}{2.644566in}}%
\pgfpathlineto{\pgfqpoint{0.756777in}{3.361654in}}%
\pgfpathlineto{\pgfqpoint{0.757548in}{3.282505in}}%
\pgfpathlineto{\pgfqpoint{0.758587in}{3.356280in}}%
\pgfpathlineto{\pgfqpoint{0.758784in}{3.249666in}}%
\pgfpathlineto{\pgfqpoint{0.758683in}{3.362284in}}%
\pgfpathlineto{\pgfqpoint{0.758797in}{3.296177in}}%
\pgfpathlineto{\pgfqpoint{0.759540in}{2.644559in}}%
\pgfpathlineto{\pgfqpoint{0.761702in}{2.644565in}}%
\pgfpathlineto{\pgfqpoint{0.761721in}{3.323456in}}%
\pgfpathlineto{\pgfqpoint{0.762485in}{3.197592in}}%
\pgfpathlineto{\pgfqpoint{0.763326in}{3.305120in}}%
\pgfpathlineto{\pgfqpoint{0.763674in}{3.326930in}}%
\pgfpathlineto{\pgfqpoint{0.764709in}{2.644559in}}%
\pgfpathlineto{\pgfqpoint{0.766648in}{2.644565in}}%
\pgfpathlineto{\pgfqpoint{0.766666in}{3.254257in}}%
\pgfpathlineto{\pgfqpoint{0.767420in}{3.135168in}}%
\pgfpathlineto{\pgfqpoint{0.768261in}{3.215151in}}%
\pgfpathlineto{\pgfqpoint{0.768620in}{3.237654in}}%
\pgfpathlineto{\pgfqpoint{0.770185in}{2.644559in}}%
\pgfpathlineto{\pgfqpoint{0.771593in}{2.644566in}}%
\pgfpathlineto{\pgfqpoint{0.771612in}{3.169059in}}%
\pgfpathlineto{\pgfqpoint{0.772363in}{3.057886in}}%
\pgfpathlineto{\pgfqpoint{0.773500in}{3.135703in}}%
\pgfpathlineto{\pgfqpoint{0.773620in}{3.028842in}}%
\pgfpathlineto{\pgfqpoint{0.773565in}{3.139800in}}%
\pgfpathlineto{\pgfqpoint{0.773632in}{3.073800in}}%
\pgfpathlineto{\pgfqpoint{0.774396in}{2.644559in}}%
\pgfpathlineto{\pgfqpoint{0.776538in}{2.644566in}}%
\pgfpathlineto{\pgfqpoint{0.776554in}{3.084840in}}%
\pgfpathlineto{\pgfqpoint{0.777274in}{2.964448in}}%
\pgfpathlineto{\pgfqpoint{0.777324in}{2.962773in}}%
\pgfpathlineto{\pgfqpoint{0.777472in}{2.976870in}}%
\pgfpathlineto{\pgfqpoint{0.777521in}{2.975195in}}%
\pgfpathlineto{\pgfqpoint{0.778511in}{3.042100in}}%
\pgfpathlineto{\pgfqpoint{0.778511in}{3.015907in}}%
\pgfpathlineto{\pgfqpoint{0.778538in}{3.023713in}}%
\pgfpathlineto{\pgfqpoint{0.779269in}{2.644559in}}%
\pgfpathlineto{\pgfqpoint{0.781484in}{2.644564in}}%
\pgfpathlineto{\pgfqpoint{0.781499in}{2.984482in}}%
\pgfpathlineto{\pgfqpoint{0.782246in}{2.863524in}}%
\pgfpathlineto{\pgfqpoint{0.782295in}{2.859981in}}%
\pgfpathlineto{\pgfqpoint{0.782641in}{2.888369in}}%
\pgfpathlineto{\pgfqpoint{0.782691in}{2.884831in}}%
\pgfpathlineto{\pgfqpoint{0.783420in}{2.937242in}}%
\pgfpathlineto{\pgfqpoint{0.783483in}{2.921070in}}%
\pgfpathlineto{\pgfqpoint{0.784278in}{2.644559in}}%
\pgfpathlineto{\pgfqpoint{0.786429in}{2.644566in}}%
\pgfpathlineto{\pgfqpoint{0.786445in}{2.877127in}}%
\pgfpathlineto{\pgfqpoint{0.787159in}{2.764699in}}%
\pgfpathlineto{\pgfqpoint{0.787208in}{2.740495in}}%
\pgfpathlineto{\pgfqpoint{0.787752in}{2.801918in}}%
\pgfpathlineto{\pgfqpoint{0.787802in}{2.777832in}}%
\pgfpathlineto{\pgfqpoint{0.788402in}{2.842640in}}%
\pgfpathlineto{\pgfqpoint{0.788466in}{2.733096in}}%
\pgfpathlineto{\pgfqpoint{0.789246in}{2.644559in}}%
\pgfpathlineto{\pgfqpoint{0.791375in}{2.644559in}}%
\pgfpathlineto{\pgfqpoint{0.791451in}{2.615542in}}%
\pgfpathlineto{\pgfqpoint{0.792133in}{2.723205in}}%
\pgfpathlineto{\pgfqpoint{0.792183in}{2.659003in}}%
\pgfpathlineto{\pgfqpoint{0.792826in}{2.766553in}}%
\pgfpathlineto{\pgfqpoint{0.792875in}{2.702645in}}%
\pgfpathlineto{\pgfqpoint{0.793309in}{2.796801in}}%
\pgfpathlineto{\pgfqpoint{0.793434in}{2.644559in}}%
\pgfpathlineto{\pgfqpoint{0.793602in}{2.644559in}}%
\pgfpathlineto{\pgfqpoint{0.796318in}{2.644559in}}%
\pgfpathlineto{\pgfqpoint{0.796330in}{2.810162in}}%
\pgfpathlineto{\pgfqpoint{0.796453in}{2.631019in}}%
\pgfpathlineto{\pgfqpoint{0.797051in}{2.665675in}}%
\pgfpathlineto{\pgfqpoint{0.797892in}{2.754698in}}%
\pgfpathlineto{\pgfqpoint{0.798265in}{2.778115in}}%
\pgfpathlineto{\pgfqpoint{0.798825in}{2.644559in}}%
\pgfpathlineto{\pgfqpoint{0.801264in}{2.644559in}}%
\pgfpathlineto{\pgfqpoint{0.801275in}{2.788532in}}%
\pgfpathlineto{\pgfqpoint{0.801410in}{2.622304in}}%
\pgfpathlineto{\pgfqpoint{0.802018in}{2.658768in}}%
\pgfpathlineto{\pgfqpoint{0.802859in}{2.761881in}}%
\pgfpathlineto{\pgfqpoint{0.803238in}{2.785637in}}%
\pgfpathlineto{\pgfqpoint{0.803751in}{2.644559in}}%
\pgfpathlineto{\pgfqpoint{0.806210in}{2.644559in}}%
\pgfpathlineto{\pgfqpoint{0.806220in}{2.788749in}}%
\pgfpathlineto{\pgfqpoint{0.806350in}{2.627306in}}%
\pgfpathlineto{\pgfqpoint{0.806953in}{2.663102in}}%
\pgfpathlineto{\pgfqpoint{0.807794in}{2.756175in}}%
\pgfpathlineto{\pgfqpoint{0.808184in}{2.780577in}}%
\pgfpathlineto{\pgfqpoint{0.808677in}{2.644559in}}%
\pgfpathlineto{\pgfqpoint{0.811155in}{2.644559in}}%
\pgfpathlineto{\pgfqpoint{0.811166in}{2.788959in}}%
\pgfpathlineto{\pgfqpoint{0.811290in}{2.615421in}}%
\pgfpathlineto{\pgfqpoint{0.811889in}{2.650622in}}%
\pgfpathlineto{\pgfqpoint{0.812730in}{2.767884in}}%
\pgfpathlineto{\pgfqpoint{0.813103in}{2.791197in}}%
\pgfpathlineto{\pgfqpoint{0.813593in}{2.644559in}}%
\pgfpathlineto{\pgfqpoint{0.816101in}{2.644559in}}%
\pgfpathlineto{\pgfqpoint{0.816115in}{2.780513in}}%
\pgfpathlineto{\pgfqpoint{0.816232in}{2.634716in}}%
\pgfpathlineto{\pgfqpoint{0.816829in}{2.669286in}}%
\pgfpathlineto{\pgfqpoint{0.817670in}{2.749383in}}%
\pgfpathlineto{\pgfqpoint{0.818045in}{2.772914in}}%
\pgfpathlineto{\pgfqpoint{0.818618in}{2.644559in}}%
\pgfpathlineto{\pgfqpoint{0.821046in}{2.644559in}}%
\pgfpathlineto{\pgfqpoint{0.821057in}{2.809625in}}%
\pgfpathlineto{\pgfqpoint{0.821217in}{2.639784in}}%
\pgfpathlineto{\pgfqpoint{0.821798in}{2.694528in}}%
\pgfpathlineto{\pgfqpoint{0.821847in}{2.678419in}}%
\pgfpathlineto{\pgfqpoint{0.822391in}{2.731767in}}%
\pgfpathlineto{\pgfqpoint{0.822441in}{2.715745in}}%
\pgfpathlineto{\pgfqpoint{0.822977in}{2.768560in}}%
\pgfpathlineto{\pgfqpoint{0.823087in}{2.676068in}}%
\pgfpathlineto{\pgfqpoint{0.823875in}{2.644559in}}%
\pgfpathlineto{\pgfqpoint{0.825991in}{2.644559in}}%
\pgfpathlineto{\pgfqpoint{0.826001in}{2.808904in}}%
\pgfpathlineto{\pgfqpoint{0.826131in}{2.615185in}}%
\pgfpathlineto{\pgfqpoint{0.826734in}{2.650763in}}%
\pgfpathlineto{\pgfqpoint{0.827575in}{2.771467in}}%
\pgfpathlineto{\pgfqpoint{0.827965in}{2.795927in}}%
\pgfpathlineto{\pgfqpoint{0.828439in}{2.644559in}}%
\pgfpathlineto{\pgfqpoint{0.830937in}{2.644559in}}%
\pgfpathlineto{\pgfqpoint{0.830954in}{2.781081in}}%
\pgfpathlineto{\pgfqpoint{0.831085in}{2.627006in}}%
\pgfpathlineto{\pgfqpoint{0.831696in}{2.663548in}}%
\pgfpathlineto{\pgfqpoint{0.832536in}{2.762086in}}%
\pgfpathlineto{\pgfqpoint{0.832911in}{2.785549in}}%
\pgfpathlineto{\pgfqpoint{0.833396in}{2.644559in}}%
\pgfpathlineto{\pgfqpoint{0.835884in}{2.644697in}}%
\pgfpathlineto{\pgfqpoint{0.835899in}{2.815952in}}%
\pgfpathlineto{\pgfqpoint{0.836013in}{2.619811in}}%
\pgfpathlineto{\pgfqpoint{0.836659in}{2.719198in}}%
\pgfpathlineto{\pgfqpoint{0.836708in}{2.660256in}}%
\pgfpathlineto{\pgfqpoint{0.837351in}{2.762558in}}%
\pgfpathlineto{\pgfqpoint{0.837400in}{2.703886in}}%
\pgfpathlineto{\pgfqpoint{0.837826in}{2.792307in}}%
\pgfpathlineto{\pgfqpoint{0.837942in}{2.644559in}}%
\pgfpathlineto{\pgfqpoint{0.838101in}{2.644559in}}%
\pgfpathlineto{\pgfqpoint{0.840829in}{2.644561in}}%
\pgfpathlineto{\pgfqpoint{0.840844in}{2.806825in}}%
\pgfpathlineto{\pgfqpoint{0.840969in}{2.616129in}}%
\pgfpathlineto{\pgfqpoint{0.841575in}{2.651866in}}%
\pgfpathlineto{\pgfqpoint{0.842415in}{2.775014in}}%
\pgfpathlineto{\pgfqpoint{0.842802in}{2.799201in}}%
\pgfpathlineto{\pgfqpoint{0.843281in}{2.644559in}}%
\pgfpathlineto{\pgfqpoint{0.845775in}{2.644562in}}%
\pgfpathlineto{\pgfqpoint{0.845790in}{2.809361in}}%
\pgfpathlineto{\pgfqpoint{0.845905in}{2.625529in}}%
\pgfpathlineto{\pgfqpoint{0.846505in}{2.659979in}}%
\pgfpathlineto{\pgfqpoint{0.847346in}{2.766568in}}%
\pgfpathlineto{\pgfqpoint{0.847719in}{2.789984in}}%
\pgfpathlineto{\pgfqpoint{0.848212in}{2.644559in}}%
\pgfpathlineto{\pgfqpoint{0.850718in}{2.644559in}}%
\pgfpathlineto{\pgfqpoint{0.850734in}{2.788524in}}%
\pgfpathlineto{\pgfqpoint{0.850855in}{2.635869in}}%
\pgfpathlineto{\pgfqpoint{0.851458in}{2.670880in}}%
\pgfpathlineto{\pgfqpoint{0.852299in}{2.758545in}}%
\pgfpathlineto{\pgfqpoint{0.852693in}{2.783230in}}%
\pgfpathlineto{\pgfqpoint{0.853205in}{2.644559in}}%
\pgfpathlineto{\pgfqpoint{0.855666in}{2.644568in}}%
\pgfpathlineto{\pgfqpoint{0.855680in}{2.813252in}}%
\pgfpathlineto{\pgfqpoint{0.855824in}{2.626191in}}%
\pgfpathlineto{\pgfqpoint{0.856398in}{2.717235in}}%
\pgfpathlineto{\pgfqpoint{0.856448in}{2.663984in}}%
\pgfpathlineto{\pgfqpoint{0.857091in}{2.760629in}}%
\pgfpathlineto{\pgfqpoint{0.857140in}{2.707580in}}%
\pgfpathlineto{\pgfqpoint{0.857638in}{2.794911in}}%
\pgfpathlineto{\pgfqpoint{0.857725in}{2.644559in}}%
\pgfpathlineto{\pgfqpoint{0.857844in}{2.644559in}}%
\pgfpathlineto{\pgfqpoint{0.860611in}{2.644635in}}%
\pgfpathlineto{\pgfqpoint{0.860625in}{2.818822in}}%
\pgfpathlineto{\pgfqpoint{0.860770in}{2.627477in}}%
\pgfpathlineto{\pgfqpoint{0.861345in}{2.718159in}}%
\pgfpathlineto{\pgfqpoint{0.861395in}{2.665355in}}%
\pgfpathlineto{\pgfqpoint{0.862038in}{2.761554in}}%
\pgfpathlineto{\pgfqpoint{0.862087in}{2.708951in}}%
\pgfpathlineto{\pgfqpoint{0.862584in}{2.795750in}}%
\pgfpathlineto{\pgfqpoint{0.862670in}{2.644559in}}%
\pgfpathlineto{\pgfqpoint{0.862788in}{2.644559in}}%
\pgfpathlineto{\pgfqpoint{0.865554in}{2.644559in}}%
\pgfpathlineto{\pgfqpoint{0.865564in}{2.807154in}}%
\pgfpathlineto{\pgfqpoint{0.865709in}{2.638788in}}%
\pgfpathlineto{\pgfqpoint{0.866332in}{2.676103in}}%
\pgfpathlineto{\pgfqpoint{0.867172in}{2.764444in}}%
\pgfpathlineto{\pgfqpoint{0.867529in}{2.786813in}}%
\pgfpathlineto{\pgfqpoint{0.868092in}{2.644559in}}%
\pgfpathlineto{\pgfqpoint{0.870500in}{2.644559in}}%
\pgfpathlineto{\pgfqpoint{0.870517in}{2.792261in}}%
\pgfpathlineto{\pgfqpoint{0.870629in}{2.625048in}}%
\pgfpathlineto{\pgfqpoint{0.871234in}{2.659646in}}%
\pgfpathlineto{\pgfqpoint{0.872075in}{2.777741in}}%
\pgfpathlineto{\pgfqpoint{0.872448in}{2.801082in}}%
\pgfpathlineto{\pgfqpoint{0.872926in}{2.644559in}}%
\pgfpathlineto{\pgfqpoint{0.875446in}{2.644559in}}%
\pgfpathlineto{\pgfqpoint{0.875463in}{2.791618in}}%
\pgfpathlineto{\pgfqpoint{0.875581in}{2.635255in}}%
\pgfpathlineto{\pgfqpoint{0.876193in}{2.670938in}}%
\pgfpathlineto{\pgfqpoint{0.877034in}{2.770768in}}%
\pgfpathlineto{\pgfqpoint{0.877420in}{2.794943in}}%
\pgfpathlineto{\pgfqpoint{0.877905in}{2.644559in}}%
\pgfpathlineto{\pgfqpoint{0.880392in}{2.644559in}}%
\pgfpathlineto{\pgfqpoint{0.880408in}{2.790970in}}%
\pgfpathlineto{\pgfqpoint{0.880522in}{2.631135in}}%
\pgfpathlineto{\pgfqpoint{0.881134in}{2.666551in}}%
\pgfpathlineto{\pgfqpoint{0.881975in}{2.776760in}}%
\pgfpathlineto{\pgfqpoint{0.882365in}{2.801235in}}%
\pgfpathlineto{\pgfqpoint{0.882861in}{2.644559in}}%
\pgfpathlineto{\pgfqpoint{0.885337in}{2.644559in}}%
\pgfpathlineto{\pgfqpoint{0.885353in}{2.789439in}}%
\pgfpathlineto{\pgfqpoint{0.885474in}{2.635199in}}%
\pgfpathlineto{\pgfqpoint{0.886094in}{2.671795in}}%
\pgfpathlineto{\pgfqpoint{0.886935in}{2.775240in}}%
\pgfpathlineto{\pgfqpoint{0.887311in}{2.798790in}}%
\pgfpathlineto{\pgfqpoint{0.887828in}{2.644559in}}%
\pgfpathlineto{\pgfqpoint{0.890282in}{2.644559in}}%
\pgfpathlineto{\pgfqpoint{0.890298in}{2.787806in}}%
\pgfpathlineto{\pgfqpoint{0.890426in}{2.629816in}}%
\pgfpathlineto{\pgfqpoint{0.891053in}{2.667417in}}%
\pgfpathlineto{\pgfqpoint{0.891894in}{2.783045in}}%
\pgfpathlineto{\pgfqpoint{0.892256in}{2.805732in}}%
\pgfpathlineto{\pgfqpoint{0.892756in}{2.644559in}}%
\pgfpathlineto{\pgfqpoint{0.895228in}{2.644559in}}%
\pgfpathlineto{\pgfqpoint{0.895243in}{2.783700in}}%
\pgfpathlineto{\pgfqpoint{0.895986in}{2.683643in}}%
\pgfpathlineto{\pgfqpoint{0.896826in}{2.766877in}}%
\pgfpathlineto{\pgfqpoint{0.897202in}{2.790411in}}%
\pgfpathlineto{\pgfqpoint{0.897805in}{2.644559in}}%
\pgfpathlineto{\pgfqpoint{0.900173in}{2.644559in}}%
\pgfpathlineto{\pgfqpoint{0.900185in}{2.781236in}}%
\pgfpathlineto{\pgfqpoint{0.900322in}{2.641714in}}%
\pgfpathlineto{\pgfqpoint{0.900910in}{2.719002in}}%
\pgfpathlineto{\pgfqpoint{0.900960in}{2.680255in}}%
\pgfpathlineto{\pgfqpoint{0.901603in}{2.762439in}}%
\pgfpathlineto{\pgfqpoint{0.901652in}{2.723808in}}%
\pgfpathlineto{\pgfqpoint{0.902147in}{2.796586in}}%
\pgfpathlineto{\pgfqpoint{0.902239in}{2.644559in}}%
\pgfpathlineto{\pgfqpoint{0.902301in}{2.644559in}}%
\pgfpathlineto{\pgfqpoint{0.905119in}{2.644559in}}%
\pgfpathlineto{\pgfqpoint{0.905132in}{2.784751in}}%
\pgfpathlineto{\pgfqpoint{0.905255in}{2.641304in}}%
\pgfpathlineto{\pgfqpoint{0.905882in}{2.678757in}}%
\pgfpathlineto{\pgfqpoint{0.906723in}{2.774084in}}%
\pgfpathlineto{\pgfqpoint{0.907093in}{2.797250in}}%
\pgfpathlineto{\pgfqpoint{0.907656in}{2.644559in}}%
\pgfpathlineto{\pgfqpoint{0.910064in}{2.644559in}}%
\pgfpathlineto{\pgfqpoint{0.910077in}{2.783825in}}%
\pgfpathlineto{\pgfqpoint{0.910814in}{2.687652in}}%
\pgfpathlineto{\pgfqpoint{0.911655in}{2.764192in}}%
\pgfpathlineto{\pgfqpoint{0.912038in}{2.788221in}}%
\pgfpathlineto{\pgfqpoint{0.912683in}{2.644559in}}%
\pgfpathlineto{\pgfqpoint{0.915010in}{2.644559in}}%
\pgfpathlineto{\pgfqpoint{0.915022in}{2.782406in}}%
\pgfpathlineto{\pgfqpoint{0.915742in}{2.682781in}}%
\pgfpathlineto{\pgfqpoint{0.916583in}{2.767207in}}%
\pgfpathlineto{\pgfqpoint{0.916956in}{2.790641in}}%
\pgfpathlineto{\pgfqpoint{0.917552in}{2.644559in}}%
\pgfpathlineto{\pgfqpoint{0.919955in}{2.644559in}}%
\pgfpathlineto{\pgfqpoint{0.919967in}{2.781481in}}%
\pgfpathlineto{\pgfqpoint{0.920077in}{2.631623in}}%
\pgfpathlineto{\pgfqpoint{0.920695in}{2.667797in}}%
\pgfpathlineto{\pgfqpoint{0.921536in}{2.783559in}}%
\pgfpathlineto{\pgfqpoint{0.921929in}{2.808191in}}%
\pgfpathlineto{\pgfqpoint{0.922406in}{2.644559in}}%
\pgfpathlineto{\pgfqpoint{0.924901in}{2.644559in}}%
\pgfpathlineto{\pgfqpoint{0.924913in}{2.780472in}}%
\pgfpathlineto{\pgfqpoint{0.925029in}{2.629304in}}%
\pgfpathlineto{\pgfqpoint{0.925657in}{2.666705in}}%
\pgfpathlineto{\pgfqpoint{0.926498in}{2.786971in}}%
\pgfpathlineto{\pgfqpoint{0.926874in}{2.810554in}}%
\pgfpathlineto{\pgfqpoint{0.927335in}{2.644559in}}%
\pgfpathlineto{\pgfqpoint{0.929846in}{2.644559in}}%
\pgfpathlineto{\pgfqpoint{0.929858in}{2.778331in}}%
\pgfpathlineto{\pgfqpoint{0.929978in}{2.628041in}}%
\pgfpathlineto{\pgfqpoint{0.930610in}{2.666063in}}%
\pgfpathlineto{\pgfqpoint{0.931451in}{2.788843in}}%
\pgfpathlineto{\pgfqpoint{0.931820in}{2.811922in}}%
\pgfpathlineto{\pgfqpoint{0.932310in}{2.644559in}}%
\pgfpathlineto{\pgfqpoint{0.934792in}{2.644559in}}%
\pgfpathlineto{\pgfqpoint{0.934803in}{2.776250in}}%
\pgfpathlineto{\pgfqpoint{0.934934in}{2.628121in}}%
\pgfpathlineto{\pgfqpoint{0.935526in}{2.734518in}}%
\pgfpathlineto{\pgfqpoint{0.935576in}{2.667328in}}%
\pgfpathlineto{\pgfqpoint{0.936218in}{2.777862in}}%
\pgfpathlineto{\pgfqpoint{0.936268in}{2.710972in}}%
\pgfpathlineto{\pgfqpoint{0.936765in}{2.812068in}}%
\pgfpathlineto{\pgfqpoint{0.936850in}{2.644559in}}%
\pgfpathlineto{\pgfqpoint{0.936994in}{2.644559in}}%
\pgfpathlineto{\pgfqpoint{0.939737in}{2.644559in}}%
\pgfpathlineto{\pgfqpoint{0.939747in}{2.773605in}}%
\pgfpathlineto{\pgfqpoint{0.939878in}{2.634985in}}%
\pgfpathlineto{\pgfqpoint{0.940470in}{2.727682in}}%
\pgfpathlineto{\pgfqpoint{0.940520in}{2.674128in}}%
\pgfpathlineto{\pgfqpoint{0.941163in}{2.771055in}}%
\pgfpathlineto{\pgfqpoint{0.941212in}{2.717743in}}%
\pgfpathlineto{\pgfqpoint{0.941711in}{2.805376in}}%
\pgfpathlineto{\pgfqpoint{0.941796in}{2.644559in}}%
\pgfpathlineto{\pgfqpoint{0.941916in}{2.644559in}}%
\pgfpathlineto{\pgfqpoint{0.944682in}{2.644559in}}%
\pgfpathlineto{\pgfqpoint{0.944692in}{2.773438in}}%
\pgfpathlineto{\pgfqpoint{0.944816in}{2.641325in}}%
\pgfpathlineto{\pgfqpoint{0.945454in}{2.679867in}}%
\pgfpathlineto{\pgfqpoint{0.946295in}{2.776450in}}%
\pgfpathlineto{\pgfqpoint{0.946656in}{2.799089in}}%
\pgfpathlineto{\pgfqpoint{0.947188in}{2.644559in}}%
\pgfpathlineto{\pgfqpoint{0.949628in}{2.644559in}}%
\pgfpathlineto{\pgfqpoint{0.949638in}{2.773007in}}%
\pgfpathlineto{\pgfqpoint{0.949759in}{2.642906in}}%
\pgfpathlineto{\pgfqpoint{0.950397in}{2.681400in}}%
\pgfpathlineto{\pgfqpoint{0.951238in}{2.774484in}}%
\pgfpathlineto{\pgfqpoint{0.951602in}{2.797267in}}%
\pgfpathlineto{\pgfqpoint{0.952157in}{2.644559in}}%
\pgfpathlineto{\pgfqpoint{0.954573in}{2.644559in}}%
\pgfpathlineto{\pgfqpoint{0.954583in}{2.772481in}}%
\pgfpathlineto{\pgfqpoint{0.955332in}{2.687952in}}%
\pgfpathlineto{\pgfqpoint{0.956173in}{2.766641in}}%
\pgfpathlineto{\pgfqpoint{0.956547in}{2.790114in}}%
\pgfpathlineto{\pgfqpoint{0.957174in}{2.644559in}}%
\pgfpathlineto{\pgfqpoint{0.959519in}{2.644559in}}%
\pgfpathlineto{\pgfqpoint{0.959529in}{2.771671in}}%
\pgfpathlineto{\pgfqpoint{0.960258in}{2.695666in}}%
\pgfpathlineto{\pgfqpoint{0.961296in}{2.768829in}}%
\pgfpathlineto{\pgfqpoint{0.961547in}{2.692212in}}%
\pgfpathlineto{\pgfqpoint{0.961493in}{2.781175in}}%
\pgfpathlineto{\pgfqpoint{0.961560in}{2.703858in}}%
\pgfpathlineto{\pgfqpoint{0.962321in}{2.644559in}}%
\pgfpathlineto{\pgfqpoint{0.964464in}{2.644559in}}%
\pgfpathlineto{\pgfqpoint{0.964474in}{2.770599in}}%
\pgfpathlineto{\pgfqpoint{0.965235in}{2.684646in}}%
\pgfpathlineto{\pgfqpoint{0.966075in}{2.771307in}}%
\pgfpathlineto{\pgfqpoint{0.966438in}{2.794060in}}%
\pgfpathlineto{\pgfqpoint{0.967030in}{2.644559in}}%
\pgfpathlineto{\pgfqpoint{0.969410in}{2.644559in}}%
\pgfpathlineto{\pgfqpoint{0.969420in}{2.769144in}}%
\pgfpathlineto{\pgfqpoint{0.969526in}{2.631654in}}%
\pgfpathlineto{\pgfqpoint{0.970151in}{2.668672in}}%
\pgfpathlineto{\pgfqpoint{0.970992in}{2.783481in}}%
\pgfpathlineto{\pgfqpoint{0.971384in}{2.808015in}}%
\pgfpathlineto{\pgfqpoint{0.971850in}{2.644559in}}%
\pgfpathlineto{\pgfqpoint{0.974355in}{2.644559in}}%
\pgfpathlineto{\pgfqpoint{0.974365in}{2.767974in}}%
\pgfpathlineto{\pgfqpoint{0.974483in}{2.632711in}}%
\pgfpathlineto{\pgfqpoint{0.975123in}{2.671420in}}%
\pgfpathlineto{\pgfqpoint{0.975963in}{2.783985in}}%
\pgfpathlineto{\pgfqpoint{0.976329in}{2.806861in}}%
\pgfpathlineto{\pgfqpoint{0.976841in}{2.644559in}}%
\pgfpathlineto{\pgfqpoint{0.979301in}{2.644559in}}%
\pgfpathlineto{\pgfqpoint{0.979311in}{2.766143in}}%
\pgfpathlineto{\pgfqpoint{0.979441in}{2.630044in}}%
\pgfpathlineto{\pgfqpoint{0.980039in}{2.732113in}}%
\pgfpathlineto{\pgfqpoint{0.980089in}{2.669902in}}%
\pgfpathlineto{\pgfqpoint{0.980732in}{2.775460in}}%
\pgfpathlineto{\pgfqpoint{0.980781in}{2.713544in}}%
\pgfpathlineto{\pgfqpoint{0.981226in}{2.806421in}}%
\pgfpathlineto{\pgfqpoint{0.981361in}{2.644559in}}%
\pgfpathlineto{\pgfqpoint{0.981500in}{2.644559in}}%
\pgfpathlineto{\pgfqpoint{0.984246in}{2.644559in}}%
\pgfpathlineto{\pgfqpoint{0.984256in}{2.764344in}}%
\pgfpathlineto{\pgfqpoint{0.984383in}{2.628146in}}%
\pgfpathlineto{\pgfqpoint{0.984982in}{2.733716in}}%
\pgfpathlineto{\pgfqpoint{0.985032in}{2.667930in}}%
\pgfpathlineto{\pgfqpoint{0.985674in}{2.777059in}}%
\pgfpathlineto{\pgfqpoint{0.985724in}{2.711575in}}%
\pgfpathlineto{\pgfqpoint{0.986220in}{2.811171in}}%
\pgfpathlineto{\pgfqpoint{0.986303in}{2.644559in}}%
\pgfpathlineto{\pgfqpoint{0.986445in}{2.644559in}}%
\pgfpathlineto{\pgfqpoint{0.989192in}{2.644559in}}%
\pgfpathlineto{\pgfqpoint{0.989201in}{2.763178in}}%
\pgfpathlineto{\pgfqpoint{0.989328in}{2.630333in}}%
\pgfpathlineto{\pgfqpoint{0.989926in}{2.731011in}}%
\pgfpathlineto{\pgfqpoint{0.989976in}{2.670064in}}%
\pgfpathlineto{\pgfqpoint{0.990618in}{2.774365in}}%
\pgfpathlineto{\pgfqpoint{0.990668in}{2.713699in}}%
\pgfpathlineto{\pgfqpoint{0.991165in}{2.808579in}}%
\pgfpathlineto{\pgfqpoint{0.991253in}{2.644559in}}%
\pgfpathlineto{\pgfqpoint{0.991387in}{2.644559in}}%
\pgfpathlineto{\pgfqpoint{0.994137in}{2.644559in}}%
\pgfpathlineto{\pgfqpoint{0.994146in}{2.764440in}}%
\pgfpathlineto{\pgfqpoint{0.994875in}{2.688670in}}%
\pgfpathlineto{\pgfqpoint{0.995715in}{2.762407in}}%
\pgfpathlineto{\pgfqpoint{0.996111in}{2.787218in}}%
\pgfpathlineto{\pgfqpoint{0.996790in}{2.644559in}}%
\pgfpathlineto{\pgfqpoint{0.999082in}{2.644559in}}%
\pgfpathlineto{\pgfqpoint{0.999091in}{2.763332in}}%
\pgfpathlineto{\pgfqpoint{0.999844in}{2.708460in}}%
\pgfpathlineto{\pgfqpoint{0.999894in}{2.695772in}}%
\pgfpathlineto{\pgfqpoint{1.000438in}{2.745708in}}%
\pgfpathlineto{\pgfqpoint{1.000487in}{2.733086in}}%
\pgfpathlineto{\pgfqpoint{1.001019in}{2.782187in}}%
\pgfpathlineto{\pgfqpoint{1.001111in}{2.683670in}}%
\pgfpathlineto{\pgfqpoint{1.001124in}{2.709398in}}%
\pgfpathlineto{\pgfqpoint{1.001873in}{2.644559in}}%
\pgfpathlineto{\pgfqpoint{1.004028in}{2.644559in}}%
\pgfpathlineto{\pgfqpoint{1.004036in}{2.762014in}}%
\pgfpathlineto{\pgfqpoint{1.004756in}{2.695762in}}%
\pgfpathlineto{\pgfqpoint{1.005972in}{2.777380in}}%
\pgfpathlineto{\pgfqpoint{1.006068in}{2.694308in}}%
\pgfpathlineto{\pgfqpoint{1.006853in}{2.644559in}}%
\pgfpathlineto{\pgfqpoint{1.008973in}{2.644559in}}%
\pgfpathlineto{\pgfqpoint{1.008982in}{2.760625in}}%
\pgfpathlineto{\pgfqpoint{1.009738in}{2.700989in}}%
\pgfpathlineto{\pgfqpoint{1.010947in}{2.776127in}}%
\pgfpathlineto{\pgfqpoint{1.010974in}{2.762750in}}%
\pgfpathlineto{\pgfqpoint{1.011726in}{2.644559in}}%
\pgfpathlineto{\pgfqpoint{1.013919in}{2.644559in}}%
\pgfpathlineto{\pgfqpoint{1.013925in}{2.760332in}}%
\pgfpathlineto{\pgfqpoint{1.014051in}{2.637390in}}%
\pgfpathlineto{\pgfqpoint{1.014651in}{2.722769in}}%
\pgfpathlineto{\pgfqpoint{1.014700in}{2.677131in}}%
\pgfpathlineto{\pgfqpoint{1.015343in}{2.766158in}}%
\pgfpathlineto{\pgfqpoint{1.015392in}{2.720731in}}%
\pgfpathlineto{\pgfqpoint{1.015893in}{2.800584in}}%
\pgfpathlineto{\pgfqpoint{1.015979in}{2.644559in}}%
\pgfpathlineto{\pgfqpoint{1.016086in}{2.644559in}}%
\pgfpathlineto{\pgfqpoint{1.018864in}{2.644559in}}%
\pgfpathlineto{\pgfqpoint{1.018870in}{2.760309in}}%
\pgfpathlineto{\pgfqpoint{1.018994in}{2.642631in}}%
\pgfpathlineto{\pgfqpoint{1.019591in}{2.717053in}}%
\pgfpathlineto{\pgfqpoint{1.019641in}{2.682116in}}%
\pgfpathlineto{\pgfqpoint{1.020284in}{2.760465in}}%
\pgfpathlineto{\pgfqpoint{1.020333in}{2.725694in}}%
\pgfpathlineto{\pgfqpoint{1.020838in}{2.795213in}}%
\pgfpathlineto{\pgfqpoint{1.020924in}{2.644559in}}%
\pgfpathlineto{\pgfqpoint{1.020990in}{2.644559in}}%
\pgfpathlineto{\pgfqpoint{1.023810in}{2.644559in}}%
\pgfpathlineto{\pgfqpoint{1.023816in}{2.760232in}}%
\pgfpathlineto{\pgfqpoint{1.024553in}{2.699931in}}%
\pgfpathlineto{\pgfqpoint{1.024899in}{2.719611in}}%
\pgfpathlineto{\pgfqpoint{1.025784in}{2.775142in}}%
\pgfpathlineto{\pgfqpoint{1.025810in}{2.762447in}}%
\pgfpathlineto{\pgfqpoint{1.026583in}{2.644559in}}%
\pgfpathlineto{\pgfqpoint{1.028755in}{2.644559in}}%
\pgfpathlineto{\pgfqpoint{1.028761in}{2.760113in}}%
\pgfpathlineto{\pgfqpoint{1.029526in}{2.694045in}}%
\pgfpathlineto{\pgfqpoint{1.030466in}{2.765981in}}%
\pgfpathlineto{\pgfqpoint{1.030729in}{2.782510in}}%
\pgfpathlineto{\pgfqpoint{1.031804in}{2.644559in}}%
\pgfpathlineto{\pgfqpoint{1.033701in}{2.644559in}}%
\pgfpathlineto{\pgfqpoint{1.033707in}{2.760022in}}%
\pgfpathlineto{\pgfqpoint{1.034428in}{2.693697in}}%
\pgfpathlineto{\pgfqpoint{1.035467in}{2.767014in}}%
\pgfpathlineto{\pgfqpoint{1.035645in}{2.778203in}}%
\pgfpathlineto{\pgfqpoint{1.035742in}{2.690711in}}%
\pgfpathlineto{\pgfqpoint{1.036526in}{2.644559in}}%
\pgfpathlineto{\pgfqpoint{1.038646in}{2.644559in}}%
\pgfpathlineto{\pgfqpoint{1.038652in}{2.759832in}}%
\pgfpathlineto{\pgfqpoint{1.038758in}{2.636954in}}%
\pgfpathlineto{\pgfqpoint{1.039388in}{2.674599in}}%
\pgfpathlineto{\pgfqpoint{1.040229in}{2.775056in}}%
\pgfpathlineto{\pgfqpoint{1.040620in}{2.799586in}}%
\pgfpathlineto{\pgfqpoint{1.041123in}{2.644559in}}%
\pgfpathlineto{\pgfqpoint{1.043592in}{2.644559in}}%
\pgfpathlineto{\pgfqpoint{1.043598in}{2.759652in}}%
\pgfpathlineto{\pgfqpoint{1.043712in}{2.635778in}}%
\pgfpathlineto{\pgfqpoint{1.044355in}{2.674743in}}%
\pgfpathlineto{\pgfqpoint{1.045196in}{2.777476in}}%
\pgfpathlineto{\pgfqpoint{1.045565in}{2.800614in}}%
\pgfpathlineto{\pgfqpoint{1.046104in}{2.644559in}}%
\pgfpathlineto{\pgfqpoint{1.048537in}{2.644559in}}%
\pgfpathlineto{\pgfqpoint{1.048543in}{2.759413in}}%
\pgfpathlineto{\pgfqpoint{1.048670in}{2.634994in}}%
\pgfpathlineto{\pgfqpoint{1.049274in}{2.723825in}}%
\pgfpathlineto{\pgfqpoint{1.049324in}{2.675149in}}%
\pgfpathlineto{\pgfqpoint{1.049967in}{2.767235in}}%
\pgfpathlineto{\pgfqpoint{1.050016in}{2.718728in}}%
\pgfpathlineto{\pgfqpoint{1.050511in}{2.801341in}}%
\pgfpathlineto{\pgfqpoint{1.050598in}{2.644559in}}%
\pgfpathlineto{\pgfqpoint{1.050705in}{2.644559in}}%
\pgfpathlineto{\pgfqpoint{1.053482in}{2.644559in}}%
\pgfpathlineto{\pgfqpoint{1.053488in}{2.759145in}}%
\pgfpathlineto{\pgfqpoint{1.053621in}{2.633987in}}%
\pgfpathlineto{\pgfqpoint{1.054232in}{2.725475in}}%
\pgfpathlineto{\pgfqpoint{1.054281in}{2.674731in}}%
\pgfpathlineto{\pgfqpoint{1.054924in}{2.768878in}}%
\pgfpathlineto{\pgfqpoint{1.054973in}{2.718317in}}%
\pgfpathlineto{\pgfqpoint{1.055413in}{2.799511in}}%
\pgfpathlineto{\pgfqpoint{1.055545in}{2.644559in}}%
\pgfpathlineto{\pgfqpoint{1.055651in}{2.644559in}}%
\pgfpathlineto{\pgfqpoint{1.058428in}{2.644559in}}%
\pgfpathlineto{\pgfqpoint{1.058434in}{2.758680in}}%
\pgfpathlineto{\pgfqpoint{1.058522in}{2.633450in}}%
\pgfpathlineto{\pgfqpoint{1.059193in}{2.726303in}}%
\pgfpathlineto{\pgfqpoint{1.059242in}{2.675520in}}%
\pgfpathlineto{\pgfqpoint{1.059885in}{2.769702in}}%
\pgfpathlineto{\pgfqpoint{1.059934in}{2.719110in}}%
\pgfpathlineto{\pgfqpoint{1.060366in}{2.799841in}}%
\pgfpathlineto{\pgfqpoint{1.060490in}{2.644559in}}%
\pgfpathlineto{\pgfqpoint{1.060595in}{2.644559in}}%
\pgfpathlineto{\pgfqpoint{1.063373in}{2.644559in}}%
\pgfpathlineto{\pgfqpoint{1.063379in}{2.758453in}}%
\pgfpathlineto{\pgfqpoint{1.063472in}{2.632465in}}%
\pgfpathlineto{\pgfqpoint{1.064145in}{2.727303in}}%
\pgfpathlineto{\pgfqpoint{1.064195in}{2.675175in}}%
\pgfpathlineto{\pgfqpoint{1.064837in}{2.770698in}}%
\pgfpathlineto{\pgfqpoint{1.064887in}{2.718769in}}%
\pgfpathlineto{\pgfqpoint{1.065315in}{2.800612in}}%
\pgfpathlineto{\pgfqpoint{1.065435in}{2.644559in}}%
\pgfpathlineto{\pgfqpoint{1.065588in}{2.644559in}}%
\pgfpathlineto{\pgfqpoint{1.068319in}{2.644559in}}%
\pgfpathlineto{\pgfqpoint{1.068325in}{2.757913in}}%
\pgfpathlineto{\pgfqpoint{1.068426in}{2.631163in}}%
\pgfpathlineto{\pgfqpoint{1.069055in}{2.668733in}}%
\pgfpathlineto{\pgfqpoint{1.069896in}{2.778617in}}%
\pgfpathlineto{\pgfqpoint{1.070267in}{2.801884in}}%
\pgfpathlineto{\pgfqpoint{1.070780in}{2.644559in}}%
\pgfpathlineto{\pgfqpoint{1.073264in}{2.644559in}}%
\pgfpathlineto{\pgfqpoint{1.073270in}{2.757408in}}%
\pgfpathlineto{\pgfqpoint{1.073375in}{2.630187in}}%
\pgfpathlineto{\pgfqpoint{1.074009in}{2.668354in}}%
\pgfpathlineto{\pgfqpoint{1.074850in}{2.779898in}}%
\pgfpathlineto{\pgfqpoint{1.075238in}{2.804203in}}%
\pgfpathlineto{\pgfqpoint{1.075732in}{2.644559in}}%
\pgfpathlineto{\pgfqpoint{1.078210in}{2.644559in}}%
\pgfpathlineto{\pgfqpoint{1.078216in}{2.757165in}}%
\pgfpathlineto{\pgfqpoint{1.078344in}{2.638795in}}%
\pgfpathlineto{\pgfqpoint{1.078954in}{2.718896in}}%
\pgfpathlineto{\pgfqpoint{1.079003in}{2.679478in}}%
\pgfpathlineto{\pgfqpoint{1.079646in}{2.762286in}}%
\pgfpathlineto{\pgfqpoint{1.079696in}{2.723077in}}%
\pgfpathlineto{\pgfqpoint{1.080137in}{2.793069in}}%
\pgfpathlineto{\pgfqpoint{1.080273in}{2.644559in}}%
\pgfpathlineto{\pgfqpoint{1.080378in}{2.644559in}}%
\pgfpathlineto{\pgfqpoint{1.083155in}{2.644559in}}%
\pgfpathlineto{\pgfqpoint{1.083161in}{2.756484in}}%
\pgfpathlineto{\pgfqpoint{1.083287in}{2.642739in}}%
\pgfpathlineto{\pgfqpoint{1.083896in}{2.714233in}}%
\pgfpathlineto{\pgfqpoint{1.083945in}{2.683258in}}%
\pgfpathlineto{\pgfqpoint{1.084588in}{2.757642in}}%
\pgfpathlineto{\pgfqpoint{1.084637in}{2.726838in}}%
\pgfpathlineto{\pgfqpoint{1.085083in}{2.788647in}}%
\pgfpathlineto{\pgfqpoint{1.085218in}{2.644559in}}%
\pgfpathlineto{\pgfqpoint{1.085281in}{2.644559in}}%
\pgfpathlineto{\pgfqpoint{1.088101in}{2.644559in}}%
\pgfpathlineto{\pgfqpoint{1.088107in}{2.756147in}}%
\pgfpathlineto{\pgfqpoint{1.088231in}{2.644362in}}%
\pgfpathlineto{\pgfqpoint{1.088838in}{2.712187in}}%
\pgfpathlineto{\pgfqpoint{1.088887in}{2.684711in}}%
\pgfpathlineto{\pgfqpoint{1.089530in}{2.755603in}}%
\pgfpathlineto{\pgfqpoint{1.089579in}{2.728283in}}%
\pgfpathlineto{\pgfqpoint{1.090074in}{2.789731in}}%
\pgfpathlineto{\pgfqpoint{1.090162in}{2.644559in}}%
\pgfpathlineto{\pgfqpoint{1.090205in}{2.644559in}}%
\pgfpathlineto{\pgfqpoint{1.093046in}{2.644559in}}%
\pgfpathlineto{\pgfqpoint{1.093052in}{2.755836in}}%
\pgfpathlineto{\pgfqpoint{1.093779in}{2.710243in}}%
\pgfpathlineto{\pgfqpoint{1.093829in}{2.685979in}}%
\pgfpathlineto{\pgfqpoint{1.094373in}{2.747464in}}%
\pgfpathlineto{\pgfqpoint{1.094422in}{2.723321in}}%
\pgfpathlineto{\pgfqpoint{1.095020in}{2.788025in}}%
\pgfpathlineto{\pgfqpoint{1.095077in}{2.687000in}}%
\pgfpathlineto{\pgfqpoint{1.095092in}{2.695613in}}%
\pgfpathlineto{\pgfqpoint{1.095846in}{2.644559in}}%
\pgfpathlineto{\pgfqpoint{1.097992in}{2.644559in}}%
\pgfpathlineto{\pgfqpoint{1.097997in}{2.755260in}}%
\pgfpathlineto{\pgfqpoint{1.098717in}{2.705779in}}%
\pgfpathlineto{\pgfqpoint{1.098766in}{2.689238in}}%
\pgfpathlineto{\pgfqpoint{1.099310in}{2.743014in}}%
\pgfpathlineto{\pgfqpoint{1.099360in}{2.726566in}}%
\pgfpathlineto{\pgfqpoint{1.099965in}{2.784095in}}%
\pgfpathlineto{\pgfqpoint{1.100034in}{2.689594in}}%
\pgfpathlineto{\pgfqpoint{1.100818in}{2.644559in}}%
\pgfpathlineto{\pgfqpoint{1.102937in}{2.644559in}}%
\pgfpathlineto{\pgfqpoint{1.102943in}{2.756938in}}%
\pgfpathlineto{\pgfqpoint{1.103070in}{2.641842in}}%
\pgfpathlineto{\pgfqpoint{1.103681in}{2.714477in}}%
\pgfpathlineto{\pgfqpoint{1.103730in}{2.682592in}}%
\pgfpathlineto{\pgfqpoint{1.104373in}{2.757883in}}%
\pgfpathlineto{\pgfqpoint{1.104423in}{2.726175in}}%
\pgfpathlineto{\pgfqpoint{1.104865in}{2.788682in}}%
\pgfpathlineto{\pgfqpoint{1.104999in}{2.644559in}}%
\pgfpathlineto{\pgfqpoint{1.105062in}{2.644559in}}%
\pgfpathlineto{\pgfqpoint{1.107882in}{2.644559in}}%
\pgfpathlineto{\pgfqpoint{1.107888in}{2.756334in}}%
\pgfpathlineto{\pgfqpoint{1.108617in}{2.709110in}}%
\pgfpathlineto{\pgfqpoint{1.108666in}{2.686371in}}%
\pgfpathlineto{\pgfqpoint{1.109210in}{2.746332in}}%
\pgfpathlineto{\pgfqpoint{1.109260in}{2.723712in}}%
\pgfpathlineto{\pgfqpoint{1.109856in}{2.786833in}}%
\pgfpathlineto{\pgfqpoint{1.109914in}{2.687458in}}%
\pgfpathlineto{\pgfqpoint{1.109928in}{2.694733in}}%
\pgfpathlineto{\pgfqpoint{1.110672in}{2.644559in}}%
\pgfpathlineto{\pgfqpoint{1.112828in}{2.644559in}}%
\pgfpathlineto{\pgfqpoint{1.112833in}{2.756494in}}%
\pgfpathlineto{\pgfqpoint{1.113564in}{2.710324in}}%
\pgfpathlineto{\pgfqpoint{1.113614in}{2.685341in}}%
\pgfpathlineto{\pgfqpoint{1.114158in}{2.747542in}}%
\pgfpathlineto{\pgfqpoint{1.114207in}{2.722686in}}%
\pgfpathlineto{\pgfqpoint{1.114802in}{2.787897in}}%
\pgfpathlineto{\pgfqpoint{1.114860in}{2.684891in}}%
\pgfpathlineto{\pgfqpoint{1.114874in}{2.695997in}}%
\pgfpathlineto{\pgfqpoint{1.115641in}{2.644559in}}%
\pgfpathlineto{\pgfqpoint{1.117773in}{2.644559in}}%
\pgfpathlineto{\pgfqpoint{1.117779in}{2.755987in}}%
\pgfpathlineto{\pgfqpoint{1.118501in}{2.705522in}}%
\pgfpathlineto{\pgfqpoint{1.118550in}{2.688713in}}%
\pgfpathlineto{\pgfqpoint{1.119094in}{2.742755in}}%
\pgfpathlineto{\pgfqpoint{1.119144in}{2.726042in}}%
\pgfpathlineto{\pgfqpoint{1.119747in}{2.783694in}}%
\pgfpathlineto{\pgfqpoint{1.119817in}{2.690134in}}%
\pgfpathlineto{\pgfqpoint{1.120600in}{2.644559in}}%
\pgfpathlineto{\pgfqpoint{1.122719in}{2.644559in}}%
\pgfpathlineto{\pgfqpoint{1.122724in}{2.755714in}}%
\pgfpathlineto{\pgfqpoint{1.123487in}{2.684425in}}%
\pgfpathlineto{\pgfqpoint{1.124328in}{2.764388in}}%
\pgfpathlineto{\pgfqpoint{1.124693in}{2.787251in}}%
\pgfpathlineto{\pgfqpoint{1.125325in}{2.644559in}}%
\pgfpathlineto{\pgfqpoint{1.127664in}{2.644559in}}%
\pgfpathlineto{\pgfqpoint{1.127670in}{2.755379in}}%
\pgfpathlineto{\pgfqpoint{1.128420in}{2.684709in}}%
\pgfpathlineto{\pgfqpoint{1.129261in}{2.762119in}}%
\pgfpathlineto{\pgfqpoint{1.129638in}{2.785783in}}%
\pgfpathlineto{\pgfqpoint{1.130286in}{2.644559in}}%
\pgfpathlineto{\pgfqpoint{1.132610in}{2.644559in}}%
\pgfpathlineto{\pgfqpoint{1.132615in}{2.755161in}}%
\pgfpathlineto{\pgfqpoint{1.133360in}{2.685369in}}%
\pgfpathlineto{\pgfqpoint{1.134201in}{2.760699in}}%
\pgfpathlineto{\pgfqpoint{1.134584in}{2.784685in}}%
\pgfpathlineto{\pgfqpoint{1.135267in}{2.644559in}}%
\pgfpathlineto{\pgfqpoint{1.137555in}{2.644559in}}%
\pgfpathlineto{\pgfqpoint{1.137561in}{2.754749in}}%
\pgfpathlineto{\pgfqpoint{1.138300in}{2.687474in}}%
\pgfpathlineto{\pgfqpoint{1.139239in}{2.763729in}}%
\pgfpathlineto{\pgfqpoint{1.139529in}{2.781905in}}%
\pgfpathlineto{\pgfqpoint{1.140401in}{2.644559in}}%
\pgfpathlineto{\pgfqpoint{1.142501in}{2.644559in}}%
\pgfpathlineto{\pgfqpoint{1.142506in}{2.754166in}}%
\pgfpathlineto{\pgfqpoint{1.143266in}{2.702358in}}%
\pgfpathlineto{\pgfqpoint{1.143315in}{2.695300in}}%
\pgfpathlineto{\pgfqpoint{1.143761in}{2.733403in}}%
\pgfpathlineto{\pgfqpoint{1.143810in}{2.726390in}}%
\pgfpathlineto{\pgfqpoint{1.144439in}{2.775985in}}%
\pgfpathlineto{\pgfqpoint{1.144510in}{2.715061in}}%
\pgfpathlineto{\pgfqpoint{1.144519in}{2.746054in}}%
\pgfpathlineto{\pgfqpoint{1.145270in}{2.644559in}}%
\pgfpathlineto{\pgfqpoint{1.147446in}{2.644559in}}%
\pgfpathlineto{\pgfqpoint{1.147452in}{2.754047in}}%
\pgfpathlineto{\pgfqpoint{1.148203in}{2.700927in}}%
\pgfpathlineto{\pgfqpoint{1.148252in}{2.695413in}}%
\pgfpathlineto{\pgfqpoint{1.148599in}{2.725765in}}%
\pgfpathlineto{\pgfqpoint{1.148648in}{2.720284in}}%
\pgfpathlineto{\pgfqpoint{1.149380in}{2.774829in}}%
\pgfpathlineto{\pgfqpoint{1.149448in}{2.759154in}}%
\pgfpathlineto{\pgfqpoint{1.150195in}{2.644559in}}%
\pgfpathlineto{\pgfqpoint{1.152392in}{2.644559in}}%
\pgfpathlineto{\pgfqpoint{1.152397in}{2.753563in}}%
\pgfpathlineto{\pgfqpoint{1.153132in}{2.694733in}}%
\pgfpathlineto{\pgfqpoint{1.154365in}{2.773696in}}%
\pgfpathlineto{\pgfqpoint{1.154366in}{2.749093in}}%
\pgfpathlineto{\pgfqpoint{1.154391in}{2.758669in}}%
\pgfpathlineto{\pgfqpoint{1.155132in}{2.644559in}}%
\pgfpathlineto{\pgfqpoint{1.157337in}{2.644559in}}%
\pgfpathlineto{\pgfqpoint{1.157342in}{2.753416in}}%
\pgfpathlineto{\pgfqpoint{1.158072in}{2.694134in}}%
\pgfpathlineto{\pgfqpoint{1.159285in}{2.772088in}}%
\pgfpathlineto{\pgfqpoint{1.159311in}{2.749115in}}%
\pgfpathlineto{\pgfqpoint{1.159338in}{2.759027in}}%
\pgfpathlineto{\pgfqpoint{1.160058in}{2.644559in}}%
\pgfpathlineto{\pgfqpoint{1.162282in}{2.644559in}}%
\pgfpathlineto{\pgfqpoint{1.162288in}{2.752932in}}%
\pgfpathlineto{\pgfqpoint{1.163031in}{2.696977in}}%
\pgfpathlineto{\pgfqpoint{1.163179in}{2.703781in}}%
\pgfpathlineto{\pgfqpoint{1.164256in}{2.771417in}}%
\pgfpathlineto{\pgfqpoint{1.164257in}{2.748930in}}%
\pgfpathlineto{\pgfqpoint{1.164283in}{2.758947in}}%
\pgfpathlineto{\pgfqpoint{1.165005in}{2.644559in}}%
\pgfpathlineto{\pgfqpoint{1.167228in}{2.644559in}}%
\pgfpathlineto{\pgfqpoint{1.167233in}{2.752680in}}%
\pgfpathlineto{\pgfqpoint{1.167976in}{2.694252in}}%
\pgfpathlineto{\pgfqpoint{1.169202in}{2.773920in}}%
\pgfpathlineto{\pgfqpoint{1.169202in}{2.748755in}}%
\pgfpathlineto{\pgfqpoint{1.169229in}{2.758814in}}%
\pgfpathlineto{\pgfqpoint{1.169962in}{2.644559in}}%
\pgfpathlineto{\pgfqpoint{1.172173in}{2.644559in}}%
\pgfpathlineto{\pgfqpoint{1.172179in}{2.752425in}}%
\pgfpathlineto{\pgfqpoint{1.172920in}{2.696661in}}%
\pgfpathlineto{\pgfqpoint{1.173069in}{2.703379in}}%
\pgfpathlineto{\pgfqpoint{1.174147in}{2.771119in}}%
\pgfpathlineto{\pgfqpoint{1.174147in}{2.748668in}}%
\pgfpathlineto{\pgfqpoint{1.174174in}{2.758704in}}%
\pgfpathlineto{\pgfqpoint{1.174891in}{2.644559in}}%
\pgfpathlineto{\pgfqpoint{1.177119in}{2.644559in}}%
\pgfpathlineto{\pgfqpoint{1.177124in}{2.751915in}}%
\pgfpathlineto{\pgfqpoint{1.177865in}{2.693917in}}%
\pgfpathlineto{\pgfqpoint{1.179093in}{2.773684in}}%
\pgfpathlineto{\pgfqpoint{1.179093in}{2.748536in}}%
\pgfpathlineto{\pgfqpoint{1.179120in}{2.758593in}}%
\pgfpathlineto{\pgfqpoint{1.179840in}{2.644559in}}%
\pgfpathlineto{\pgfqpoint{1.182064in}{2.644559in}}%
\pgfpathlineto{\pgfqpoint{1.182070in}{2.751565in}}%
\pgfpathlineto{\pgfqpoint{1.182825in}{2.694859in}}%
\pgfpathlineto{\pgfqpoint{1.184038in}{2.773342in}}%
\pgfpathlineto{\pgfqpoint{1.184038in}{2.748371in}}%
\pgfpathlineto{\pgfqpoint{1.184065in}{2.758436in}}%
\pgfpathlineto{\pgfqpoint{1.184790in}{2.644559in}}%
\pgfpathlineto{\pgfqpoint{1.187010in}{2.644559in}}%
\pgfpathlineto{\pgfqpoint{1.187015in}{2.751347in}}%
\pgfpathlineto{\pgfqpoint{1.187766in}{2.694196in}}%
\pgfpathlineto{\pgfqpoint{1.188984in}{2.773523in}}%
\pgfpathlineto{\pgfqpoint{1.188984in}{2.748289in}}%
\pgfpathlineto{\pgfqpoint{1.189011in}{2.758342in}}%
\pgfpathlineto{\pgfqpoint{1.189744in}{2.644559in}}%
\pgfpathlineto{\pgfqpoint{1.191955in}{2.644559in}}%
\pgfpathlineto{\pgfqpoint{1.191961in}{2.751209in}}%
\pgfpathlineto{\pgfqpoint{1.192717in}{2.688753in}}%
\pgfpathlineto{\pgfqpoint{1.193656in}{2.761884in}}%
\pgfpathlineto{\pgfqpoint{1.193929in}{2.778990in}}%
\pgfpathlineto{\pgfqpoint{1.194911in}{2.644559in}}%
\pgfpathlineto{\pgfqpoint{1.196901in}{2.644559in}}%
\pgfpathlineto{\pgfqpoint{1.196906in}{2.750921in}}%
\pgfpathlineto{\pgfqpoint{1.197665in}{2.704488in}}%
\pgfpathlineto{\pgfqpoint{1.197714in}{2.690335in}}%
\pgfpathlineto{\pgfqpoint{1.198258in}{2.741747in}}%
\pgfpathlineto{\pgfqpoint{1.198308in}{2.727639in}}%
\pgfpathlineto{\pgfqpoint{1.198838in}{2.778155in}}%
\pgfpathlineto{\pgfqpoint{1.198942in}{2.683986in}}%
\pgfpathlineto{\pgfqpoint{1.199705in}{2.644559in}}%
\pgfpathlineto{\pgfqpoint{1.201846in}{2.644559in}}%
\pgfpathlineto{\pgfqpoint{1.201852in}{2.746575in}}%
\pgfpathlineto{\pgfqpoint{1.201928in}{2.630049in}}%
\pgfpathlineto{\pgfqpoint{1.202593in}{2.721248in}}%
\pgfpathlineto{\pgfqpoint{1.202642in}{2.671260in}}%
\pgfpathlineto{\pgfqpoint{1.203285in}{2.764614in}}%
\pgfpathlineto{\pgfqpoint{1.203334in}{2.714884in}}%
\pgfpathlineto{\pgfqpoint{1.203775in}{2.795302in}}%
\pgfpathlineto{\pgfqpoint{1.203910in}{2.644559in}}%
\pgfpathlineto{\pgfqpoint{1.204006in}{2.644559in}}%
\pgfpathlineto{\pgfqpoint{1.206791in}{2.644559in}}%
\pgfpathlineto{\pgfqpoint{1.206797in}{2.746912in}}%
\pgfpathlineto{\pgfqpoint{1.206873in}{2.628587in}}%
\pgfpathlineto{\pgfqpoint{1.207538in}{2.722124in}}%
\pgfpathlineto{\pgfqpoint{1.207588in}{2.669870in}}%
\pgfpathlineto{\pgfqpoint{1.208231in}{2.765484in}}%
\pgfpathlineto{\pgfqpoint{1.208280in}{2.713500in}}%
\pgfpathlineto{\pgfqpoint{1.208721in}{2.796161in}}%
\pgfpathlineto{\pgfqpoint{1.208852in}{2.644559in}}%
\pgfpathlineto{\pgfqpoint{1.208992in}{2.644559in}}%
\pgfpathlineto{\pgfqpoint{1.211737in}{2.644559in}}%
\pgfpathlineto{\pgfqpoint{1.211743in}{2.746634in}}%
\pgfpathlineto{\pgfqpoint{1.211819in}{2.628917in}}%
\pgfpathlineto{\pgfqpoint{1.212484in}{2.721769in}}%
\pgfpathlineto{\pgfqpoint{1.212533in}{2.670230in}}%
\pgfpathlineto{\pgfqpoint{1.213176in}{2.765131in}}%
\pgfpathlineto{\pgfqpoint{1.213226in}{2.713858in}}%
\pgfpathlineto{\pgfqpoint{1.213666in}{2.795801in}}%
\pgfpathlineto{\pgfqpoint{1.213798in}{2.644559in}}%
\pgfpathlineto{\pgfqpoint{1.213893in}{2.644559in}}%
\pgfpathlineto{\pgfqpoint{1.216682in}{2.644559in}}%
\pgfpathlineto{\pgfqpoint{1.216688in}{2.749905in}}%
\pgfpathlineto{\pgfqpoint{1.216774in}{2.637958in}}%
\pgfpathlineto{\pgfqpoint{1.217445in}{2.713067in}}%
\pgfpathlineto{\pgfqpoint{1.217494in}{2.680477in}}%
\pgfpathlineto{\pgfqpoint{1.218137in}{2.756495in}}%
\pgfpathlineto{\pgfqpoint{1.218187in}{2.724038in}}%
\pgfpathlineto{\pgfqpoint{1.218619in}{2.786725in}}%
\pgfpathlineto{\pgfqpoint{1.218750in}{2.644559in}}%
\pgfpathlineto{\pgfqpoint{1.218832in}{2.644559in}}%
\pgfpathlineto{\pgfqpoint{1.221628in}{2.644559in}}%
\pgfpathlineto{\pgfqpoint{1.221633in}{2.748789in}}%
\pgfpathlineto{\pgfqpoint{1.221711in}{2.625406in}}%
\pgfpathlineto{\pgfqpoint{1.222376in}{2.724842in}}%
\pgfpathlineto{\pgfqpoint{1.222426in}{2.666849in}}%
\pgfpathlineto{\pgfqpoint{1.223068in}{2.768188in}}%
\pgfpathlineto{\pgfqpoint{1.223118in}{2.710492in}}%
\pgfpathlineto{\pgfqpoint{1.223558in}{2.798809in}}%
\pgfpathlineto{\pgfqpoint{1.223684in}{2.644559in}}%
\pgfpathlineto{\pgfqpoint{1.223828in}{2.644559in}}%
\pgfpathlineto{\pgfqpoint{1.226573in}{2.644559in}}%
\pgfpathlineto{\pgfqpoint{1.226579in}{2.750671in}}%
\pgfpathlineto{\pgfqpoint{1.226668in}{2.635210in}}%
\pgfpathlineto{\pgfqpoint{1.227344in}{2.715793in}}%
\pgfpathlineto{\pgfqpoint{1.227393in}{2.678365in}}%
\pgfpathlineto{\pgfqpoint{1.228036in}{2.759211in}}%
\pgfpathlineto{\pgfqpoint{1.228086in}{2.721936in}}%
\pgfpathlineto{\pgfqpoint{1.228514in}{2.789181in}}%
\pgfpathlineto{\pgfqpoint{1.228639in}{2.644559in}}%
\pgfpathlineto{\pgfqpoint{1.228743in}{2.644559in}}%
\pgfpathlineto{\pgfqpoint{1.231518in}{2.644559in}}%
\pgfpathlineto{\pgfqpoint{1.231524in}{2.751271in}}%
\pgfpathlineto{\pgfqpoint{1.231621in}{2.632529in}}%
\pgfpathlineto{\pgfqpoint{1.232254in}{2.670464in}}%
\pgfpathlineto{\pgfqpoint{1.233095in}{2.768648in}}%
\pgfpathlineto{\pgfqpoint{1.233467in}{2.791968in}}%
\pgfpathlineto{\pgfqpoint{1.233999in}{2.644559in}}%
\pgfpathlineto{\pgfqpoint{1.236463in}{2.644559in}}%
\pgfpathlineto{\pgfqpoint{1.236469in}{2.753634in}}%
\pgfpathlineto{\pgfqpoint{1.236589in}{2.620901in}}%
\pgfpathlineto{\pgfqpoint{1.237195in}{2.727393in}}%
\pgfpathlineto{\pgfqpoint{1.237244in}{2.661364in}}%
\pgfpathlineto{\pgfqpoint{1.237887in}{2.770746in}}%
\pgfpathlineto{\pgfqpoint{1.237937in}{2.705001in}}%
\pgfpathlineto{\pgfqpoint{1.238438in}{2.805218in}}%
\pgfpathlineto{\pgfqpoint{1.238523in}{2.644559in}}%
\pgfpathlineto{\pgfqpoint{1.238660in}{2.644559in}}%
\pgfpathlineto{\pgfqpoint{1.241409in}{2.644559in}}%
\pgfpathlineto{\pgfqpoint{1.241415in}{2.752216in}}%
\pgfpathlineto{\pgfqpoint{1.241519in}{2.628039in}}%
\pgfpathlineto{\pgfqpoint{1.242158in}{2.666751in}}%
\pgfpathlineto{\pgfqpoint{1.242998in}{2.773444in}}%
\pgfpathlineto{\pgfqpoint{1.243384in}{2.797578in}}%
\pgfpathlineto{\pgfqpoint{1.243884in}{2.644559in}}%
\pgfpathlineto{\pgfqpoint{1.246355in}{2.644559in}}%
\pgfpathlineto{\pgfqpoint{1.246360in}{2.753413in}}%
\pgfpathlineto{\pgfqpoint{1.246456in}{2.633463in}}%
\pgfpathlineto{\pgfqpoint{1.247087in}{2.671192in}}%
\pgfpathlineto{\pgfqpoint{1.247928in}{2.767043in}}%
\pgfpathlineto{\pgfqpoint{1.248302in}{2.790457in}}%
\pgfpathlineto{\pgfqpoint{1.248846in}{2.644559in}}%
\pgfpathlineto{\pgfqpoint{1.251300in}{2.644559in}}%
\pgfpathlineto{\pgfqpoint{1.251306in}{2.754986in}}%
\pgfpathlineto{\pgfqpoint{1.251411in}{2.627258in}}%
\pgfpathlineto{\pgfqpoint{1.252050in}{2.666079in}}%
\pgfpathlineto{\pgfqpoint{1.252891in}{2.773851in}}%
\pgfpathlineto{\pgfqpoint{1.253274in}{2.797864in}}%
\pgfpathlineto{\pgfqpoint{1.253782in}{2.644559in}}%
\pgfpathlineto{\pgfqpoint{1.256246in}{2.644559in}}%
\pgfpathlineto{\pgfqpoint{1.256251in}{2.753353in}}%
\pgfpathlineto{\pgfqpoint{1.256350in}{2.632410in}}%
\pgfpathlineto{\pgfqpoint{1.256983in}{2.670473in}}%
\pgfpathlineto{\pgfqpoint{1.257824in}{2.768056in}}%
\pgfpathlineto{\pgfqpoint{1.258195in}{2.791309in}}%
\pgfpathlineto{\pgfqpoint{1.258734in}{2.644559in}}%
\pgfpathlineto{\pgfqpoint{1.261191in}{2.644559in}}%
\pgfpathlineto{\pgfqpoint{1.261197in}{2.755325in}}%
\pgfpathlineto{\pgfqpoint{1.261312in}{2.620064in}}%
\pgfpathlineto{\pgfqpoint{1.261963in}{2.660108in}}%
\pgfpathlineto{\pgfqpoint{1.262804in}{2.782048in}}%
\pgfpathlineto{\pgfqpoint{1.263165in}{2.804658in}}%
\pgfpathlineto{\pgfqpoint{1.263643in}{2.644559in}}%
\pgfpathlineto{\pgfqpoint{1.266136in}{2.644559in}}%
\pgfpathlineto{\pgfqpoint{1.266142in}{2.756421in}}%
\pgfpathlineto{\pgfqpoint{1.266260in}{2.619865in}}%
\pgfpathlineto{\pgfqpoint{1.266864in}{2.726973in}}%
\pgfpathlineto{\pgfqpoint{1.266913in}{2.660126in}}%
\pgfpathlineto{\pgfqpoint{1.267556in}{2.770324in}}%
\pgfpathlineto{\pgfqpoint{1.267606in}{2.703764in}}%
\pgfpathlineto{\pgfqpoint{1.268111in}{2.805021in}}%
\pgfpathlineto{\pgfqpoint{1.268194in}{2.644559in}}%
\pgfpathlineto{\pgfqpoint{1.268318in}{2.644559in}}%
\pgfpathlineto{\pgfqpoint{1.271082in}{2.644559in}}%
\pgfpathlineto{\pgfqpoint{1.271088in}{2.755968in}}%
\pgfpathlineto{\pgfqpoint{1.271212in}{2.622749in}}%
\pgfpathlineto{\pgfqpoint{1.271823in}{2.724702in}}%
\pgfpathlineto{\pgfqpoint{1.271872in}{2.663627in}}%
\pgfpathlineto{\pgfqpoint{1.272515in}{2.768064in}}%
\pgfpathlineto{\pgfqpoint{1.272564in}{2.707255in}}%
\pgfpathlineto{\pgfqpoint{1.273009in}{2.799036in}}%
\pgfpathlineto{\pgfqpoint{1.273145in}{2.644559in}}%
\pgfpathlineto{\pgfqpoint{1.273241in}{2.644559in}}%
\pgfpathlineto{\pgfqpoint{1.276027in}{2.644559in}}%
\pgfpathlineto{\pgfqpoint{1.276033in}{2.760362in}}%
\pgfpathlineto{\pgfqpoint{1.276140in}{2.625876in}}%
\pgfpathlineto{\pgfqpoint{1.276780in}{2.664805in}}%
\pgfpathlineto{\pgfqpoint{1.277620in}{2.774406in}}%
\pgfpathlineto{\pgfqpoint{1.278002in}{2.798284in}}%
\pgfpathlineto{\pgfqpoint{1.278512in}{2.644559in}}%
\pgfpathlineto{\pgfqpoint{1.280973in}{2.644559in}}%
\pgfpathlineto{\pgfqpoint{1.280979in}{2.757225in}}%
\pgfpathlineto{\pgfqpoint{1.281101in}{2.624685in}}%
\pgfpathlineto{\pgfqpoint{1.281708in}{2.722258in}}%
\pgfpathlineto{\pgfqpoint{1.281757in}{2.665253in}}%
\pgfpathlineto{\pgfqpoint{1.282400in}{2.765627in}}%
\pgfpathlineto{\pgfqpoint{1.282450in}{2.708874in}}%
\pgfpathlineto{\pgfqpoint{1.282947in}{2.799863in}}%
\pgfpathlineto{\pgfqpoint{1.283030in}{2.644559in}}%
\pgfpathlineto{\pgfqpoint{1.283165in}{2.644559in}}%
\pgfpathlineto{\pgfqpoint{1.285918in}{2.644559in}}%
\pgfpathlineto{\pgfqpoint{1.285923in}{2.761727in}}%
\pgfpathlineto{\pgfqpoint{1.286002in}{2.640295in}}%
\pgfpathlineto{\pgfqpoint{1.286667in}{2.707219in}}%
\pgfpathlineto{\pgfqpoint{1.286716in}{2.681615in}}%
\pgfpathlineto{\pgfqpoint{1.287260in}{2.744455in}}%
\pgfpathlineto{\pgfqpoint{1.287310in}{2.718942in}}%
\pgfpathlineto{\pgfqpoint{1.287848in}{2.781359in}}%
\pgfpathlineto{\pgfqpoint{1.287960in}{2.681474in}}%
\pgfpathlineto{\pgfqpoint{1.288705in}{2.644559in}}%
\pgfpathlineto{\pgfqpoint{1.290863in}{2.644559in}}%
\pgfpathlineto{\pgfqpoint{1.290869in}{2.762492in}}%
\pgfpathlineto{\pgfqpoint{1.290995in}{2.641533in}}%
\pgfpathlineto{\pgfqpoint{1.291606in}{2.705660in}}%
\pgfpathlineto{\pgfqpoint{1.291656in}{2.682290in}}%
\pgfpathlineto{\pgfqpoint{1.292200in}{2.742900in}}%
\pgfpathlineto{\pgfqpoint{1.292249in}{2.719613in}}%
\pgfpathlineto{\pgfqpoint{1.292793in}{2.780138in}}%
\pgfpathlineto{\pgfqpoint{1.292901in}{2.673525in}}%
\pgfpathlineto{\pgfqpoint{1.293665in}{2.644559in}}%
\pgfpathlineto{\pgfqpoint{1.295809in}{2.644559in}}%
\pgfpathlineto{\pgfqpoint{1.295814in}{2.766120in}}%
\pgfpathlineto{\pgfqpoint{1.295915in}{2.634904in}}%
\pgfpathlineto{\pgfqpoint{1.296598in}{2.714022in}}%
\pgfpathlineto{\pgfqpoint{1.296647in}{2.679256in}}%
\pgfpathlineto{\pgfqpoint{1.297290in}{2.757442in}}%
\pgfpathlineto{\pgfqpoint{1.297339in}{2.722826in}}%
\pgfpathlineto{\pgfqpoint{1.297784in}{2.788389in}}%
\pgfpathlineto{\pgfqpoint{1.297870in}{2.644559in}}%
\pgfpathlineto{\pgfqpoint{1.297988in}{2.644559in}}%
\pgfpathlineto{\pgfqpoint{1.300754in}{2.644559in}}%
\pgfpathlineto{\pgfqpoint{1.300760in}{2.764188in}}%
\pgfpathlineto{\pgfqpoint{1.300840in}{2.639467in}}%
\pgfpathlineto{\pgfqpoint{1.301506in}{2.707751in}}%
\pgfpathlineto{\pgfqpoint{1.301555in}{2.680999in}}%
\pgfpathlineto{\pgfqpoint{1.302099in}{2.744984in}}%
\pgfpathlineto{\pgfqpoint{1.302149in}{2.718329in}}%
\pgfpathlineto{\pgfqpoint{1.302686in}{2.781798in}}%
\pgfpathlineto{\pgfqpoint{1.302796in}{2.682638in}}%
\pgfpathlineto{\pgfqpoint{1.303547in}{2.644559in}}%
\pgfpathlineto{\pgfqpoint{1.305700in}{2.644559in}}%
\pgfpathlineto{\pgfqpoint{1.305705in}{2.768395in}}%
\pgfpathlineto{\pgfqpoint{1.305811in}{2.628170in}}%
\pgfpathlineto{\pgfqpoint{1.306499in}{2.721259in}}%
\pgfpathlineto{\pgfqpoint{1.306549in}{2.673158in}}%
\pgfpathlineto{\pgfqpoint{1.307192in}{2.764648in}}%
\pgfpathlineto{\pgfqpoint{1.307241in}{2.716758in}}%
\pgfpathlineto{\pgfqpoint{1.307674in}{2.794904in}}%
\pgfpathlineto{\pgfqpoint{1.307763in}{2.644559in}}%
\pgfpathlineto{\pgfqpoint{1.307924in}{2.644559in}}%
\pgfpathlineto{\pgfqpoint{1.310645in}{2.644559in}}%
\pgfpathlineto{\pgfqpoint{1.310651in}{2.772109in}}%
\pgfpathlineto{\pgfqpoint{1.310727in}{2.619519in}}%
\pgfpathlineto{\pgfqpoint{1.311393in}{2.726799in}}%
\pgfpathlineto{\pgfqpoint{1.311443in}{2.660922in}}%
\pgfpathlineto{\pgfqpoint{1.312085in}{2.770148in}}%
\pgfpathlineto{\pgfqpoint{1.312135in}{2.704562in}}%
\pgfpathlineto{\pgfqpoint{1.312575in}{2.800809in}}%
\pgfpathlineto{\pgfqpoint{1.312703in}{2.644559in}}%
\pgfpathlineto{\pgfqpoint{1.312842in}{2.644559in}}%
\pgfpathlineto{\pgfqpoint{1.315590in}{2.644559in}}%
\pgfpathlineto{\pgfqpoint{1.315597in}{2.754345in}}%
\pgfpathlineto{\pgfqpoint{1.315678in}{2.613428in}}%
\pgfpathlineto{\pgfqpoint{1.316343in}{2.732634in}}%
\pgfpathlineto{\pgfqpoint{1.316393in}{2.655351in}}%
\pgfpathlineto{\pgfqpoint{1.317035in}{2.775938in}}%
\pgfpathlineto{\pgfqpoint{1.317085in}{2.699036in}}%
\pgfpathlineto{\pgfqpoint{1.317523in}{2.806423in}}%
\pgfpathlineto{\pgfqpoint{1.317645in}{2.644559in}}%
\pgfpathlineto{\pgfqpoint{1.317819in}{2.644559in}}%
\pgfpathlineto{\pgfqpoint{1.320536in}{2.644559in}}%
\pgfpathlineto{\pgfqpoint{1.320542in}{2.774094in}}%
\pgfpathlineto{\pgfqpoint{1.320666in}{2.618432in}}%
\pgfpathlineto{\pgfqpoint{1.321371in}{2.733085in}}%
\pgfpathlineto{\pgfqpoint{1.321420in}{2.665280in}}%
\pgfpathlineto{\pgfqpoint{1.322063in}{2.776430in}}%
\pgfpathlineto{\pgfqpoint{1.322113in}{2.708924in}}%
\pgfpathlineto{\pgfqpoint{1.322511in}{2.804410in}}%
\pgfpathlineto{\pgfqpoint{1.322592in}{2.644559in}}%
\pgfpathlineto{\pgfqpoint{1.322803in}{2.644559in}}%
\pgfpathlineto{\pgfqpoint{1.325481in}{2.644559in}}%
\pgfpathlineto{\pgfqpoint{1.325487in}{2.777447in}}%
\pgfpathlineto{\pgfqpoint{1.325597in}{2.624775in}}%
\pgfpathlineto{\pgfqpoint{1.326386in}{2.730404in}}%
\pgfpathlineto{\pgfqpoint{1.326435in}{2.676333in}}%
\pgfpathlineto{\pgfqpoint{1.327078in}{2.773779in}}%
\pgfpathlineto{\pgfqpoint{1.327128in}{2.719946in}}%
\pgfpathlineto{\pgfqpoint{1.327456in}{2.797442in}}%
\pgfpathlineto{\pgfqpoint{1.327541in}{2.644559in}}%
\pgfpathlineto{\pgfqpoint{1.327835in}{2.644559in}}%
\pgfpathlineto{\pgfqpoint{1.330427in}{2.644559in}}%
\pgfpathlineto{\pgfqpoint{1.330432in}{2.777265in}}%
\pgfpathlineto{\pgfqpoint{1.330543in}{2.624469in}}%
\pgfpathlineto{\pgfqpoint{1.331233in}{2.724638in}}%
\pgfpathlineto{\pgfqpoint{1.331283in}{2.669834in}}%
\pgfpathlineto{\pgfqpoint{1.331926in}{2.768012in}}%
\pgfpathlineto{\pgfqpoint{1.331975in}{2.713450in}}%
\pgfpathlineto{\pgfqpoint{1.332402in}{2.797823in}}%
\pgfpathlineto{\pgfqpoint{1.332486in}{2.644559in}}%
\pgfpathlineto{\pgfqpoint{1.332646in}{2.644559in}}%
\pgfpathlineto{\pgfqpoint{1.335375in}{2.644809in}}%
\pgfpathlineto{\pgfqpoint{1.335378in}{2.778776in}}%
\pgfpathlineto{\pgfqpoint{1.335487in}{2.633117in}}%
\pgfpathlineto{\pgfqpoint{1.336123in}{2.671837in}}%
\pgfpathlineto{\pgfqpoint{1.336963in}{2.764800in}}%
\pgfpathlineto{\pgfqpoint{1.337347in}{2.788866in}}%
\pgfpathlineto{\pgfqpoint{1.337895in}{2.644559in}}%
\pgfpathlineto{\pgfqpoint{1.340321in}{2.644789in}}%
\pgfpathlineto{\pgfqpoint{1.340323in}{2.778641in}}%
\pgfpathlineto{\pgfqpoint{1.340427in}{2.634017in}}%
\pgfpathlineto{\pgfqpoint{1.341059in}{2.672216in}}%
\pgfpathlineto{\pgfqpoint{1.341900in}{2.763224in}}%
\pgfpathlineto{\pgfqpoint{1.342293in}{2.787862in}}%
\pgfpathlineto{\pgfqpoint{1.342830in}{2.644559in}}%
\pgfpathlineto{\pgfqpoint{1.345266in}{2.644596in}}%
\pgfpathlineto{\pgfqpoint{1.345268in}{2.771115in}}%
\pgfpathlineto{\pgfqpoint{1.346028in}{2.690671in}}%
\pgfpathlineto{\pgfqpoint{1.347238in}{2.770763in}}%
\pgfpathlineto{\pgfqpoint{1.347238in}{2.757869in}}%
\pgfpathlineto{\pgfqpoint{1.347964in}{2.644559in}}%
\pgfpathlineto{\pgfqpoint{1.350211in}{2.644587in}}%
\pgfpathlineto{\pgfqpoint{1.350214in}{2.768895in}}%
\pgfpathlineto{\pgfqpoint{1.350955in}{2.689986in}}%
\pgfpathlineto{\pgfqpoint{1.352184in}{2.770201in}}%
\pgfpathlineto{\pgfqpoint{1.352184in}{2.757863in}}%
\pgfpathlineto{\pgfqpoint{1.352941in}{2.644559in}}%
\pgfpathlineto{\pgfqpoint{1.355157in}{2.644570in}}%
\pgfpathlineto{\pgfqpoint{1.355159in}{2.759880in}}%
\pgfpathlineto{\pgfqpoint{1.355265in}{2.643682in}}%
\pgfpathlineto{\pgfqpoint{1.355901in}{2.682129in}}%
\pgfpathlineto{\pgfqpoint{1.356841in}{2.759626in}}%
\pgfpathlineto{\pgfqpoint{1.357129in}{2.777695in}}%
\pgfpathlineto{\pgfqpoint{1.357966in}{2.644559in}}%
\pgfpathlineto{\pgfqpoint{1.360102in}{2.644565in}}%
\pgfpathlineto{\pgfqpoint{1.360104in}{2.752778in}}%
\pgfpathlineto{\pgfqpoint{1.360211in}{2.643508in}}%
\pgfpathlineto{\pgfqpoint{1.360897in}{2.703940in}}%
\pgfpathlineto{\pgfqpoint{1.360946in}{2.688220in}}%
\pgfpathlineto{\pgfqpoint{1.361490in}{2.741171in}}%
\pgfpathlineto{\pgfqpoint{1.361540in}{2.725552in}}%
\pgfpathlineto{\pgfqpoint{1.362074in}{2.777813in}}%
\pgfpathlineto{\pgfqpoint{1.362131in}{2.685425in}}%
\pgfpathlineto{\pgfqpoint{1.362145in}{2.687601in}}%
\pgfpathlineto{\pgfqpoint{1.362888in}{2.644559in}}%
\pgfpathlineto{\pgfqpoint{1.365048in}{2.644566in}}%
\pgfpathlineto{\pgfqpoint{1.365050in}{2.753803in}}%
\pgfpathlineto{\pgfqpoint{1.365157in}{2.642873in}}%
\pgfpathlineto{\pgfqpoint{1.365794in}{2.681401in}}%
\pgfpathlineto{\pgfqpoint{1.366733in}{2.760256in}}%
\pgfpathlineto{\pgfqpoint{1.367020in}{2.778226in}}%
\pgfpathlineto{\pgfqpoint{1.367818in}{2.644559in}}%
\pgfpathlineto{\pgfqpoint{1.369993in}{2.644566in}}%
\pgfpathlineto{\pgfqpoint{1.369995in}{2.753497in}}%
\pgfpathlineto{\pgfqpoint{1.370102in}{2.643472in}}%
\pgfpathlineto{\pgfqpoint{1.370787in}{2.703529in}}%
\pgfpathlineto{\pgfqpoint{1.370836in}{2.688127in}}%
\pgfpathlineto{\pgfqpoint{1.371380in}{2.740760in}}%
\pgfpathlineto{\pgfqpoint{1.371430in}{2.725458in}}%
\pgfpathlineto{\pgfqpoint{1.371965in}{2.777452in}}%
\pgfpathlineto{\pgfqpoint{1.372022in}{2.685649in}}%
\pgfpathlineto{\pgfqpoint{1.372036in}{2.687189in}}%
\pgfpathlineto{\pgfqpoint{1.372792in}{2.644559in}}%
\pgfpathlineto{\pgfqpoint{1.374938in}{2.644565in}}%
\pgfpathlineto{\pgfqpoint{1.374941in}{2.752900in}}%
\pgfpathlineto{\pgfqpoint{1.375048in}{2.643026in}}%
\pgfpathlineto{\pgfqpoint{1.375733in}{2.703907in}}%
\pgfpathlineto{\pgfqpoint{1.375783in}{2.687715in}}%
\pgfpathlineto{\pgfqpoint{1.376327in}{2.741136in}}%
\pgfpathlineto{\pgfqpoint{1.376376in}{2.725048in}}%
\pgfpathlineto{\pgfqpoint{1.376911in}{2.777781in}}%
\pgfpathlineto{\pgfqpoint{1.376968in}{2.684226in}}%
\pgfpathlineto{\pgfqpoint{1.376982in}{2.687819in}}%
\pgfpathlineto{\pgfqpoint{1.377736in}{2.644559in}}%
\pgfpathlineto{\pgfqpoint{1.379884in}{2.644563in}}%
\pgfpathlineto{\pgfqpoint{1.379887in}{2.746446in}}%
\pgfpathlineto{\pgfqpoint{1.379997in}{2.642951in}}%
\pgfpathlineto{\pgfqpoint{1.380685in}{2.704183in}}%
\pgfpathlineto{\pgfqpoint{1.380734in}{2.687961in}}%
\pgfpathlineto{\pgfqpoint{1.381278in}{2.741412in}}%
\pgfpathlineto{\pgfqpoint{1.381328in}{2.725295in}}%
\pgfpathlineto{\pgfqpoint{1.381856in}{2.777676in}}%
\pgfpathlineto{\pgfqpoint{1.381913in}{2.683476in}}%
\pgfpathlineto{\pgfqpoint{1.381928in}{2.688026in}}%
\pgfpathlineto{\pgfqpoint{1.382695in}{2.644559in}}%
\pgfpathlineto{\pgfqpoint{1.384829in}{2.644562in}}%
\pgfpathlineto{\pgfqpoint{1.384833in}{2.746442in}}%
\pgfpathlineto{\pgfqpoint{1.384942in}{2.643162in}}%
\pgfpathlineto{\pgfqpoint{1.385629in}{2.703825in}}%
\pgfpathlineto{\pgfqpoint{1.385679in}{2.688126in}}%
\pgfpathlineto{\pgfqpoint{1.386223in}{2.741055in}}%
\pgfpathlineto{\pgfqpoint{1.386272in}{2.725459in}}%
\pgfpathlineto{\pgfqpoint{1.386802in}{2.777368in}}%
\pgfpathlineto{\pgfqpoint{1.386859in}{2.684132in}}%
\pgfpathlineto{\pgfqpoint{1.386873in}{2.687623in}}%
\pgfpathlineto{\pgfqpoint{1.387650in}{2.644559in}}%
\pgfpathlineto{\pgfqpoint{1.389775in}{2.644562in}}%
\pgfpathlineto{\pgfqpoint{1.389778in}{2.746414in}}%
\pgfpathlineto{\pgfqpoint{1.389886in}{2.644163in}}%
\pgfpathlineto{\pgfqpoint{1.390572in}{2.702375in}}%
\pgfpathlineto{\pgfqpoint{1.390621in}{2.688948in}}%
\pgfpathlineto{\pgfqpoint{1.391165in}{2.739609in}}%
\pgfpathlineto{\pgfqpoint{1.391215in}{2.726277in}}%
\pgfpathlineto{\pgfqpoint{1.391747in}{2.776121in}}%
\pgfpathlineto{\pgfqpoint{1.391817in}{2.685850in}}%
\pgfpathlineto{\pgfqpoint{1.392560in}{2.644559in}}%
\pgfpathlineto{\pgfqpoint{1.394720in}{2.644562in}}%
\pgfpathlineto{\pgfqpoint{1.394724in}{2.746481in}}%
\pgfpathlineto{\pgfqpoint{1.394832in}{2.643882in}}%
\pgfpathlineto{\pgfqpoint{1.395518in}{2.702621in}}%
\pgfpathlineto{\pgfqpoint{1.395567in}{2.688697in}}%
\pgfpathlineto{\pgfqpoint{1.396111in}{2.739854in}}%
\pgfpathlineto{\pgfqpoint{1.396161in}{2.726027in}}%
\pgfpathlineto{\pgfqpoint{1.396693in}{2.776321in}}%
\pgfpathlineto{\pgfqpoint{1.396749in}{2.686176in}}%
\pgfpathlineto{\pgfqpoint{1.396763in}{2.686348in}}%
\pgfpathlineto{\pgfqpoint{1.397532in}{2.644559in}}%
\pgfpathlineto{\pgfqpoint{1.399666in}{2.644564in}}%
\pgfpathlineto{\pgfqpoint{1.399668in}{2.750298in}}%
\pgfpathlineto{\pgfqpoint{1.399778in}{2.642500in}}%
\pgfpathlineto{\pgfqpoint{1.400465in}{2.703936in}}%
\pgfpathlineto{\pgfqpoint{1.400515in}{2.687433in}}%
\pgfpathlineto{\pgfqpoint{1.401059in}{2.741164in}}%
\pgfpathlineto{\pgfqpoint{1.401108in}{2.724767in}}%
\pgfpathlineto{\pgfqpoint{1.401638in}{2.777489in}}%
\pgfpathlineto{\pgfqpoint{1.401695in}{2.682789in}}%
\pgfpathlineto{\pgfqpoint{1.401709in}{2.687890in}}%
\pgfpathlineto{\pgfqpoint{1.402468in}{2.644559in}}%
\pgfpathlineto{\pgfqpoint{1.404611in}{2.644562in}}%
\pgfpathlineto{\pgfqpoint{1.404614in}{2.746658in}}%
\pgfpathlineto{\pgfqpoint{1.404722in}{2.643913in}}%
\pgfpathlineto{\pgfqpoint{1.405408in}{2.702223in}}%
\pgfpathlineto{\pgfqpoint{1.405458in}{2.688681in}}%
\pgfpathlineto{\pgfqpoint{1.406002in}{2.739456in}}%
\pgfpathlineto{\pgfqpoint{1.406051in}{2.726011in}}%
\pgfpathlineto{\pgfqpoint{1.406584in}{2.775958in}}%
\pgfpathlineto{\pgfqpoint{1.406654in}{2.686104in}}%
\pgfpathlineto{\pgfqpoint{1.407414in}{2.644559in}}%
\pgfpathlineto{\pgfqpoint{1.409557in}{2.644562in}}%
\pgfpathlineto{\pgfqpoint{1.409560in}{2.746736in}}%
\pgfpathlineto{\pgfqpoint{1.409669in}{2.643223in}}%
\pgfpathlineto{\pgfqpoint{1.410355in}{2.702827in}}%
\pgfpathlineto{\pgfqpoint{1.410404in}{2.688051in}}%
\pgfpathlineto{\pgfqpoint{1.410948in}{2.740058in}}%
\pgfpathlineto{\pgfqpoint{1.410998in}{2.725383in}}%
\pgfpathlineto{\pgfqpoint{1.411529in}{2.776480in}}%
\pgfpathlineto{\pgfqpoint{1.411586in}{2.683928in}}%
\pgfpathlineto{\pgfqpoint{1.411600in}{2.687101in}}%
\pgfpathlineto{\pgfqpoint{1.412368in}{2.644559in}}%
\pgfpathlineto{\pgfqpoint{1.414502in}{2.644562in}}%
\pgfpathlineto{\pgfqpoint{1.414505in}{2.746728in}}%
\pgfpathlineto{\pgfqpoint{1.414613in}{2.643891in}}%
\pgfpathlineto{\pgfqpoint{1.415298in}{2.701876in}}%
\pgfpathlineto{\pgfqpoint{1.415348in}{2.688608in}}%
\pgfpathlineto{\pgfqpoint{1.415892in}{2.739110in}}%
\pgfpathlineto{\pgfqpoint{1.415941in}{2.725937in}}%
\pgfpathlineto{\pgfqpoint{1.416474in}{2.775655in}}%
\pgfpathlineto{\pgfqpoint{1.416545in}{2.685877in}}%
\pgfpathlineto{\pgfqpoint{1.417298in}{2.644559in}}%
\pgfpathlineto{\pgfqpoint{1.419448in}{2.644562in}}%
\pgfpathlineto{\pgfqpoint{1.419451in}{2.746762in}}%
\pgfpathlineto{\pgfqpoint{1.419559in}{2.643198in}}%
\pgfpathlineto{\pgfqpoint{1.420245in}{2.702516in}}%
\pgfpathlineto{\pgfqpoint{1.420295in}{2.687981in}}%
\pgfpathlineto{\pgfqpoint{1.420839in}{2.739747in}}%
\pgfpathlineto{\pgfqpoint{1.420888in}{2.725313in}}%
\pgfpathlineto{\pgfqpoint{1.421420in}{2.776209in}}%
\pgfpathlineto{\pgfqpoint{1.421477in}{2.684044in}}%
\pgfpathlineto{\pgfqpoint{1.421491in}{2.686821in}}%
\pgfpathlineto{\pgfqpoint{1.422259in}{2.644559in}}%
\pgfpathlineto{\pgfqpoint{1.424393in}{2.644562in}}%
\pgfpathlineto{\pgfqpoint{1.424396in}{2.746811in}}%
\pgfpathlineto{\pgfqpoint{1.424505in}{2.643250in}}%
\pgfpathlineto{\pgfqpoint{1.425190in}{2.702261in}}%
\pgfpathlineto{\pgfqpoint{1.425240in}{2.688000in}}%
\pgfpathlineto{\pgfqpoint{1.425784in}{2.739492in}}%
\pgfpathlineto{\pgfqpoint{1.425833in}{2.725331in}}%
\pgfpathlineto{\pgfqpoint{1.426365in}{2.775986in}}%
\pgfpathlineto{\pgfqpoint{1.426422in}{2.684240in}}%
\pgfpathlineto{\pgfqpoint{1.426436in}{2.686594in}}%
\pgfpathlineto{\pgfqpoint{1.427175in}{2.644559in}}%
\pgfpathlineto{\pgfqpoint{1.429338in}{2.644562in}}%
\pgfpathlineto{\pgfqpoint{1.429342in}{2.746797in}}%
\pgfpathlineto{\pgfqpoint{1.429450in}{2.643250in}}%
\pgfpathlineto{\pgfqpoint{1.430135in}{2.702119in}}%
\pgfpathlineto{\pgfqpoint{1.430185in}{2.687977in}}%
\pgfpathlineto{\pgfqpoint{1.430729in}{2.739350in}}%
\pgfpathlineto{\pgfqpoint{1.430778in}{2.725308in}}%
\pgfpathlineto{\pgfqpoint{1.431311in}{2.775867in}}%
\pgfpathlineto{\pgfqpoint{1.431368in}{2.684351in}}%
\pgfpathlineto{\pgfqpoint{1.431382in}{2.686450in}}%
\pgfpathlineto{\pgfqpoint{1.432129in}{2.644559in}}%
\pgfpathlineto{\pgfqpoint{1.434285in}{2.645312in}}%
\pgfpathlineto{\pgfqpoint{1.434287in}{2.746779in}}%
\pgfpathlineto{\pgfqpoint{1.434396in}{2.643232in}}%
\pgfpathlineto{\pgfqpoint{1.435031in}{2.681707in}}%
\pgfpathlineto{\pgfqpoint{1.435970in}{2.757767in}}%
\pgfpathlineto{\pgfqpoint{1.436256in}{2.775699in}}%
\pgfpathlineto{\pgfqpoint{1.437102in}{2.644559in}}%
\pgfpathlineto{\pgfqpoint{1.439230in}{2.645310in}}%
\pgfpathlineto{\pgfqpoint{1.439233in}{2.747040in}}%
\pgfpathlineto{\pgfqpoint{1.439346in}{2.639581in}}%
\pgfpathlineto{\pgfqpoint{1.439987in}{2.678639in}}%
\pgfpathlineto{\pgfqpoint{1.440827in}{2.755801in}}%
\pgfpathlineto{\pgfqpoint{1.441202in}{2.779270in}}%
\pgfpathlineto{\pgfqpoint{1.441845in}{2.644559in}}%
\pgfpathlineto{\pgfqpoint{1.444176in}{2.645275in}}%
\pgfpathlineto{\pgfqpoint{1.444178in}{2.747078in}}%
\pgfpathlineto{\pgfqpoint{1.444292in}{2.638711in}}%
\pgfpathlineto{\pgfqpoint{1.444934in}{2.677851in}}%
\pgfpathlineto{\pgfqpoint{1.445774in}{2.756599in}}%
\pgfpathlineto{\pgfqpoint{1.446147in}{2.779961in}}%
\pgfpathlineto{\pgfqpoint{1.446784in}{2.644559in}}%
\pgfpathlineto{\pgfqpoint{1.449121in}{2.645255in}}%
\pgfpathlineto{\pgfqpoint{1.449124in}{2.747037in}}%
\pgfpathlineto{\pgfqpoint{1.449237in}{2.639077in}}%
\pgfpathlineto{\pgfqpoint{1.449878in}{2.678163in}}%
\pgfpathlineto{\pgfqpoint{1.450719in}{2.756060in}}%
\pgfpathlineto{\pgfqpoint{1.451093in}{2.779488in}}%
\pgfpathlineto{\pgfqpoint{1.451736in}{2.644559in}}%
\pgfpathlineto{\pgfqpoint{1.454066in}{2.645234in}}%
\pgfpathlineto{\pgfqpoint{1.454069in}{2.747192in}}%
\pgfpathlineto{\pgfqpoint{1.454184in}{2.636533in}}%
\pgfpathlineto{\pgfqpoint{1.454826in}{2.675723in}}%
\pgfpathlineto{\pgfqpoint{1.455666in}{2.758550in}}%
\pgfpathlineto{\pgfqpoint{1.456038in}{2.781844in}}%
\pgfpathlineto{\pgfqpoint{1.456644in}{2.644559in}}%
\pgfpathlineto{\pgfqpoint{1.459012in}{2.645208in}}%
\pgfpathlineto{\pgfqpoint{1.459014in}{2.747147in}}%
\pgfpathlineto{\pgfqpoint{1.459129in}{2.636996in}}%
\pgfpathlineto{\pgfqpoint{1.459770in}{2.676124in}}%
\pgfpathlineto{\pgfqpoint{1.460611in}{2.757894in}}%
\pgfpathlineto{\pgfqpoint{1.460984in}{2.781266in}}%
\pgfpathlineto{\pgfqpoint{1.461581in}{2.644559in}}%
\pgfpathlineto{\pgfqpoint{1.463957in}{2.644563in}}%
\pgfpathlineto{\pgfqpoint{1.463960in}{2.746503in}}%
\pgfpathlineto{\pgfqpoint{1.464072in}{2.640016in}}%
\pgfpathlineto{\pgfqpoint{1.464760in}{2.704795in}}%
\pgfpathlineto{\pgfqpoint{1.464810in}{2.685109in}}%
\pgfpathlineto{\pgfqpoint{1.465354in}{2.742016in}}%
\pgfpathlineto{\pgfqpoint{1.465403in}{2.722451in}}%
\pgfpathlineto{\pgfqpoint{1.465929in}{2.778079in}}%
\pgfpathlineto{\pgfqpoint{1.466016in}{2.644559in}}%
\pgfpathlineto{\pgfqpoint{1.468902in}{2.644563in}}%
\pgfpathlineto{\pgfqpoint{1.468905in}{2.746506in}}%
\pgfpathlineto{\pgfqpoint{1.469018in}{2.639448in}}%
\pgfpathlineto{\pgfqpoint{1.469707in}{2.705300in}}%
\pgfpathlineto{\pgfqpoint{1.469756in}{2.684573in}}%
\pgfpathlineto{\pgfqpoint{1.470300in}{2.742518in}}%
\pgfpathlineto{\pgfqpoint{1.470349in}{2.721917in}}%
\pgfpathlineto{\pgfqpoint{1.470874in}{2.778535in}}%
\pgfpathlineto{\pgfqpoint{1.470961in}{2.644559in}}%
\pgfpathlineto{\pgfqpoint{1.470972in}{2.644559in}}%
\pgfpathlineto{\pgfqpoint{1.473848in}{2.645179in}}%
\pgfpathlineto{\pgfqpoint{1.473851in}{2.747399in}}%
\pgfpathlineto{\pgfqpoint{1.473968in}{2.634610in}}%
\pgfpathlineto{\pgfqpoint{1.474610in}{2.673915in}}%
\pgfpathlineto{\pgfqpoint{1.475451in}{2.760179in}}%
\pgfpathlineto{\pgfqpoint{1.475820in}{2.783298in}}%
\pgfpathlineto{\pgfqpoint{1.476403in}{2.644559in}}%
\pgfpathlineto{\pgfqpoint{1.478794in}{2.645081in}}%
\pgfpathlineto{\pgfqpoint{1.478796in}{2.747601in}}%
\pgfpathlineto{\pgfqpoint{1.478915in}{2.633573in}}%
\pgfpathlineto{\pgfqpoint{1.479557in}{2.672961in}}%
\pgfpathlineto{\pgfqpoint{1.480398in}{2.761080in}}%
\pgfpathlineto{\pgfqpoint{1.480765in}{2.784081in}}%
\pgfpathlineto{\pgfqpoint{1.481336in}{2.644559in}}%
\pgfpathlineto{\pgfqpoint{1.483739in}{2.644562in}}%
\pgfpathlineto{\pgfqpoint{1.483742in}{2.746473in}}%
\pgfpathlineto{\pgfqpoint{1.483854in}{2.639768in}}%
\pgfpathlineto{\pgfqpoint{1.484542in}{2.704562in}}%
\pgfpathlineto{\pgfqpoint{1.484591in}{2.684821in}}%
\pgfpathlineto{\pgfqpoint{1.485135in}{2.741782in}}%
\pgfpathlineto{\pgfqpoint{1.485185in}{2.722163in}}%
\pgfpathlineto{\pgfqpoint{1.485711in}{2.777873in}}%
\pgfpathlineto{\pgfqpoint{1.485797in}{2.644559in}}%
\pgfpathlineto{\pgfqpoint{1.488685in}{2.645120in}}%
\pgfpathlineto{\pgfqpoint{1.488687in}{2.747450in}}%
\pgfpathlineto{\pgfqpoint{1.488805in}{2.633803in}}%
\pgfpathlineto{\pgfqpoint{1.489447in}{2.673130in}}%
\pgfpathlineto{\pgfqpoint{1.490288in}{2.760577in}}%
\pgfpathlineto{\pgfqpoint{1.490656in}{2.783648in}}%
\pgfpathlineto{\pgfqpoint{1.491235in}{2.644559in}}%
\pgfpathlineto{\pgfqpoint{1.493629in}{2.644562in}}%
\pgfpathlineto{\pgfqpoint{1.493633in}{2.746670in}}%
\pgfpathlineto{\pgfqpoint{1.493747in}{2.638096in}}%
\pgfpathlineto{\pgfqpoint{1.494435in}{2.706110in}}%
\pgfpathlineto{\pgfqpoint{1.494485in}{2.683280in}}%
\pgfpathlineto{\pgfqpoint{1.495029in}{2.743324in}}%
\pgfpathlineto{\pgfqpoint{1.495078in}{2.720629in}}%
\pgfpathlineto{\pgfqpoint{1.495602in}{2.779234in}}%
\pgfpathlineto{\pgfqpoint{1.495687in}{2.644559in}}%
\pgfpathlineto{\pgfqpoint{1.495710in}{2.644559in}}%
\pgfpathlineto{\pgfqpoint{1.498576in}{2.645102in}}%
\pgfpathlineto{\pgfqpoint{1.498578in}{2.747590in}}%
\pgfpathlineto{\pgfqpoint{1.498698in}{2.632925in}}%
\pgfpathlineto{\pgfqpoint{1.499340in}{2.672311in}}%
\pgfpathlineto{\pgfqpoint{1.500180in}{2.761367in}}%
\pgfpathlineto{\pgfqpoint{1.500547in}{2.784344in}}%
\pgfpathlineto{\pgfqpoint{1.501115in}{2.644559in}}%
\pgfpathlineto{\pgfqpoint{1.503520in}{2.644562in}}%
\pgfpathlineto{\pgfqpoint{1.503524in}{2.746840in}}%
\pgfpathlineto{\pgfqpoint{1.503640in}{2.636892in}}%
\pgfpathlineto{\pgfqpoint{1.504329in}{2.707256in}}%
\pgfpathlineto{\pgfqpoint{1.504378in}{2.682187in}}%
\pgfpathlineto{\pgfqpoint{1.504922in}{2.744466in}}%
\pgfpathlineto{\pgfqpoint{1.504972in}{2.719540in}}%
\pgfpathlineto{\pgfqpoint{1.505493in}{2.780220in}}%
\pgfpathlineto{\pgfqpoint{1.505579in}{2.644559in}}%
\pgfpathlineto{\pgfqpoint{1.505624in}{2.644559in}}%
\pgfpathlineto{\pgfqpoint{1.508466in}{2.645075in}}%
\pgfpathlineto{\pgfqpoint{1.508469in}{2.747753in}}%
\pgfpathlineto{\pgfqpoint{1.508590in}{2.631595in}}%
\pgfpathlineto{\pgfqpoint{1.509232in}{2.671063in}}%
\pgfpathlineto{\pgfqpoint{1.510073in}{2.762621in}}%
\pgfpathlineto{\pgfqpoint{1.510438in}{2.785470in}}%
\pgfpathlineto{\pgfqpoint{1.510991in}{2.644559in}}%
\pgfpathlineto{\pgfqpoint{1.513412in}{2.645105in}}%
\pgfpathlineto{\pgfqpoint{1.513414in}{2.747527in}}%
\pgfpathlineto{\pgfqpoint{1.513534in}{2.632757in}}%
\pgfpathlineto{\pgfqpoint{1.514175in}{2.672100in}}%
\pgfpathlineto{\pgfqpoint{1.515016in}{2.761108in}}%
\pgfpathlineto{\pgfqpoint{1.515384in}{2.784120in}}%
\pgfpathlineto{\pgfqpoint{1.515952in}{2.644559in}}%
\pgfpathlineto{\pgfqpoint{1.518357in}{2.645004in}}%
\pgfpathlineto{\pgfqpoint{1.518360in}{2.747802in}}%
\pgfpathlineto{\pgfqpoint{1.518481in}{2.631407in}}%
\pgfpathlineto{\pgfqpoint{1.519123in}{2.670866in}}%
\pgfpathlineto{\pgfqpoint{1.519964in}{2.762503in}}%
\pgfpathlineto{\pgfqpoint{1.520329in}{2.785358in}}%
\pgfpathlineto{\pgfqpoint{1.520876in}{2.644559in}}%
\pgfpathlineto{\pgfqpoint{1.523303in}{2.645092in}}%
\pgfpathlineto{\pgfqpoint{1.523305in}{2.747631in}}%
\pgfpathlineto{\pgfqpoint{1.523426in}{2.631764in}}%
\pgfpathlineto{\pgfqpoint{1.524068in}{2.671172in}}%
\pgfpathlineto{\pgfqpoint{1.524909in}{2.762034in}}%
\pgfpathlineto{\pgfqpoint{1.525274in}{2.784947in}}%
\pgfpathlineto{\pgfqpoint{1.525843in}{2.644559in}}%
\pgfpathlineto{\pgfqpoint{1.528248in}{2.645068in}}%
\pgfpathlineto{\pgfqpoint{1.528251in}{2.747802in}}%
\pgfpathlineto{\pgfqpoint{1.528373in}{2.630480in}}%
\pgfpathlineto{\pgfqpoint{1.529015in}{2.669981in}}%
\pgfpathlineto{\pgfqpoint{1.529856in}{2.763324in}}%
\pgfpathlineto{\pgfqpoint{1.530220in}{2.786096in}}%
\pgfpathlineto{\pgfqpoint{1.530767in}{2.644559in}}%
\pgfpathlineto{\pgfqpoint{1.533194in}{2.645113in}}%
\pgfpathlineto{\pgfqpoint{1.533196in}{2.747666in}}%
\pgfpathlineto{\pgfqpoint{1.533318in}{2.630724in}}%
\pgfpathlineto{\pgfqpoint{1.533960in}{2.670172in}}%
\pgfpathlineto{\pgfqpoint{1.534801in}{2.762906in}}%
\pgfpathlineto{\pgfqpoint{1.535165in}{2.785745in}}%
\pgfpathlineto{\pgfqpoint{1.535723in}{2.644559in}}%
\pgfpathlineto{\pgfqpoint{1.538139in}{2.645197in}}%
\pgfpathlineto{\pgfqpoint{1.538142in}{2.747484in}}%
\pgfpathlineto{\pgfqpoint{1.538262in}{2.631467in}}%
\pgfpathlineto{\pgfqpoint{1.538904in}{2.670824in}}%
\pgfpathlineto{\pgfqpoint{1.539744in}{2.761898in}}%
\pgfpathlineto{\pgfqpoint{1.540111in}{2.784844in}}%
\pgfpathlineto{\pgfqpoint{1.540666in}{2.644559in}}%
\pgfpathlineto{\pgfqpoint{1.543085in}{2.645038in}}%
\pgfpathlineto{\pgfqpoint{1.543087in}{2.747825in}}%
\pgfpathlineto{\pgfqpoint{1.543210in}{2.630132in}}%
\pgfpathlineto{\pgfqpoint{1.543852in}{2.669629in}}%
\pgfpathlineto{\pgfqpoint{1.544693in}{2.763380in}}%
\pgfpathlineto{\pgfqpoint{1.545056in}{2.786142in}}%
\pgfpathlineto{\pgfqpoint{1.545593in}{2.644559in}}%
\pgfpathlineto{\pgfqpoint{1.548030in}{2.645145in}}%
\pgfpathlineto{\pgfqpoint{1.548033in}{2.747713in}}%
\pgfpathlineto{\pgfqpoint{1.548156in}{2.629565in}}%
\pgfpathlineto{\pgfqpoint{1.548798in}{2.669062in}}%
\pgfpathlineto{\pgfqpoint{1.549639in}{2.763824in}}%
\pgfpathlineto{\pgfqpoint{1.550002in}{2.786564in}}%
\pgfpathlineto{\pgfqpoint{1.550529in}{2.644559in}}%
\pgfpathlineto{\pgfqpoint{1.552976in}{2.645107in}}%
\pgfpathlineto{\pgfqpoint{1.552978in}{2.747760in}}%
\pgfpathlineto{\pgfqpoint{1.553101in}{2.629313in}}%
\pgfpathlineto{\pgfqpoint{1.553744in}{2.668827in}}%
\pgfpathlineto{\pgfqpoint{1.554584in}{2.764007in}}%
\pgfpathlineto{\pgfqpoint{1.554947in}{2.786721in}}%
\pgfpathlineto{\pgfqpoint{1.555488in}{2.644559in}}%
\pgfpathlineto{\pgfqpoint{1.557921in}{2.645250in}}%
\pgfpathlineto{\pgfqpoint{1.557924in}{2.747524in}}%
\pgfpathlineto{\pgfqpoint{1.558046in}{2.629714in}}%
\pgfpathlineto{\pgfqpoint{1.558687in}{2.669108in}}%
\pgfpathlineto{\pgfqpoint{1.559528in}{2.763055in}}%
\pgfpathlineto{\pgfqpoint{1.559893in}{2.785897in}}%
\pgfpathlineto{\pgfqpoint{1.560449in}{2.644559in}}%
\pgfpathlineto{\pgfqpoint{1.562866in}{2.645091in}}%
\pgfpathlineto{\pgfqpoint{1.562869in}{2.747788in}}%
\pgfpathlineto{\pgfqpoint{1.562992in}{2.629513in}}%
\pgfpathlineto{\pgfqpoint{1.563634in}{2.669018in}}%
\pgfpathlineto{\pgfqpoint{1.564475in}{2.763787in}}%
\pgfpathlineto{\pgfqpoint{1.564838in}{2.786513in}}%
\pgfpathlineto{\pgfqpoint{1.565381in}{2.644559in}}%
\pgfpathlineto{\pgfqpoint{1.567812in}{2.645068in}}%
\pgfpathlineto{\pgfqpoint{1.567815in}{2.747936in}}%
\pgfpathlineto{\pgfqpoint{1.567941in}{2.627878in}}%
\pgfpathlineto{\pgfqpoint{1.568535in}{2.709593in}}%
\pgfpathlineto{\pgfqpoint{1.568585in}{2.667578in}}%
\pgfpathlineto{\pgfqpoint{1.569228in}{2.752965in}}%
\pgfpathlineto{\pgfqpoint{1.569277in}{2.711197in}}%
\pgfpathlineto{\pgfqpoint{1.569784in}{2.787771in}}%
\pgfpathlineto{\pgfqpoint{1.569870in}{2.644559in}}%
\pgfpathlineto{\pgfqpoint{1.569965in}{2.644559in}}%
\pgfpathlineto{\pgfqpoint{1.572757in}{2.645262in}}%
\pgfpathlineto{\pgfqpoint{1.572760in}{2.747671in}}%
\pgfpathlineto{\pgfqpoint{1.572884in}{2.628452in}}%
\pgfpathlineto{\pgfqpoint{1.573526in}{2.667951in}}%
\pgfpathlineto{\pgfqpoint{1.574367in}{2.764411in}}%
\pgfpathlineto{\pgfqpoint{1.574729in}{2.787088in}}%
\pgfpathlineto{\pgfqpoint{1.575280in}{2.644559in}}%
\pgfpathlineto{\pgfqpoint{1.577703in}{2.644898in}}%
\pgfpathlineto{\pgfqpoint{1.577705in}{2.748268in}}%
\pgfpathlineto{\pgfqpoint{1.577828in}{2.630647in}}%
\pgfpathlineto{\pgfqpoint{1.578468in}{2.669981in}}%
\pgfpathlineto{\pgfqpoint{1.579309in}{2.761805in}}%
\pgfpathlineto{\pgfqpoint{1.579674in}{2.784694in}}%
\pgfpathlineto{\pgfqpoint{1.580219in}{2.644559in}}%
\pgfpathlineto{\pgfqpoint{1.582648in}{2.645029in}}%
\pgfpathlineto{\pgfqpoint{1.582651in}{2.747982in}}%
\pgfpathlineto{\pgfqpoint{1.582778in}{2.627868in}}%
\pgfpathlineto{\pgfqpoint{1.583372in}{2.709592in}}%
\pgfpathlineto{\pgfqpoint{1.583421in}{2.667580in}}%
\pgfpathlineto{\pgfqpoint{1.584064in}{2.752964in}}%
\pgfpathlineto{\pgfqpoint{1.584114in}{2.711199in}}%
\pgfpathlineto{\pgfqpoint{1.584620in}{2.787756in}}%
\pgfpathlineto{\pgfqpoint{1.584706in}{2.644559in}}%
\pgfpathlineto{\pgfqpoint{1.584803in}{2.644559in}}%
\pgfpathlineto{\pgfqpoint{1.587594in}{2.644886in}}%
\pgfpathlineto{\pgfqpoint{1.587598in}{2.749377in}}%
\pgfpathlineto{\pgfqpoint{1.587727in}{2.617502in}}%
\pgfpathlineto{\pgfqpoint{1.588323in}{2.719891in}}%
\pgfpathlineto{\pgfqpoint{1.588372in}{2.657466in}}%
\pgfpathlineto{\pgfqpoint{1.589015in}{2.763231in}}%
\pgfpathlineto{\pgfqpoint{1.589065in}{2.701117in}}%
\pgfpathlineto{\pgfqpoint{1.589565in}{2.797627in}}%
\pgfpathlineto{\pgfqpoint{1.589646in}{2.644559in}}%
\pgfpathlineto{\pgfqpoint{1.589793in}{2.644559in}}%
\pgfpathlineto{\pgfqpoint{1.592539in}{2.645037in}}%
\pgfpathlineto{\pgfqpoint{1.592542in}{2.748038in}}%
\pgfpathlineto{\pgfqpoint{1.592669in}{2.628018in}}%
\pgfpathlineto{\pgfqpoint{1.593264in}{2.709363in}}%
\pgfpathlineto{\pgfqpoint{1.593313in}{2.667760in}}%
\pgfpathlineto{\pgfqpoint{1.593956in}{2.752734in}}%
\pgfpathlineto{\pgfqpoint{1.594006in}{2.711379in}}%
\pgfpathlineto{\pgfqpoint{1.594511in}{2.787455in}}%
\pgfpathlineto{\pgfqpoint{1.594597in}{2.644559in}}%
\pgfpathlineto{\pgfqpoint{1.594692in}{2.644559in}}%
\pgfpathlineto{\pgfqpoint{1.597485in}{2.644871in}}%
\pgfpathlineto{\pgfqpoint{1.597489in}{2.750159in}}%
\pgfpathlineto{\pgfqpoint{1.597618in}{2.616205in}}%
\pgfpathlineto{\pgfqpoint{1.598213in}{2.720903in}}%
\pgfpathlineto{\pgfqpoint{1.598263in}{2.656143in}}%
\pgfpathlineto{\pgfqpoint{1.598906in}{2.764238in}}%
\pgfpathlineto{\pgfqpoint{1.598955in}{2.699799in}}%
\pgfpathlineto{\pgfqpoint{1.599456in}{2.798657in}}%
\pgfpathlineto{\pgfqpoint{1.599539in}{2.644559in}}%
\pgfpathlineto{\pgfqpoint{1.599667in}{2.644559in}}%
\pgfpathlineto{\pgfqpoint{1.602430in}{2.645198in}}%
\pgfpathlineto{\pgfqpoint{1.602433in}{2.747788in}}%
\pgfpathlineto{\pgfqpoint{1.602560in}{2.627726in}}%
\pgfpathlineto{\pgfqpoint{1.603154in}{2.709372in}}%
\pgfpathlineto{\pgfqpoint{1.603203in}{2.667402in}}%
\pgfpathlineto{\pgfqpoint{1.603846in}{2.752743in}}%
\pgfpathlineto{\pgfqpoint{1.603895in}{2.711021in}}%
\pgfpathlineto{\pgfqpoint{1.604402in}{2.787538in}}%
\pgfpathlineto{\pgfqpoint{1.604488in}{2.644559in}}%
\pgfpathlineto{\pgfqpoint{1.604583in}{2.644559in}}%
\pgfpathlineto{\pgfqpoint{1.607376in}{2.645147in}}%
\pgfpathlineto{\pgfqpoint{1.607378in}{2.747871in}}%
\pgfpathlineto{\pgfqpoint{1.607509in}{2.628350in}}%
\pgfpathlineto{\pgfqpoint{1.608104in}{2.708627in}}%
\pgfpathlineto{\pgfqpoint{1.608153in}{2.668191in}}%
\pgfpathlineto{\pgfqpoint{1.608796in}{2.752000in}}%
\pgfpathlineto{\pgfqpoint{1.608846in}{2.711808in}}%
\pgfpathlineto{\pgfqpoint{1.609347in}{2.786480in}}%
\pgfpathlineto{\pgfqpoint{1.609434in}{2.644559in}}%
\pgfpathlineto{\pgfqpoint{1.609530in}{2.644559in}}%
\pgfpathlineto{\pgfqpoint{1.612321in}{2.645193in}}%
\pgfpathlineto{\pgfqpoint{1.612324in}{2.747903in}}%
\pgfpathlineto{\pgfqpoint{1.612455in}{2.628742in}}%
\pgfpathlineto{\pgfqpoint{1.613050in}{2.708628in}}%
\pgfpathlineto{\pgfqpoint{1.613100in}{2.668630in}}%
\pgfpathlineto{\pgfqpoint{1.613743in}{2.752001in}}%
\pgfpathlineto{\pgfqpoint{1.613792in}{2.712246in}}%
\pgfpathlineto{\pgfqpoint{1.614293in}{2.786422in}}%
\pgfpathlineto{\pgfqpoint{1.614380in}{2.644559in}}%
\pgfpathlineto{\pgfqpoint{1.614476in}{2.644559in}}%
\pgfpathlineto{\pgfqpoint{1.617267in}{2.645143in}}%
\pgfpathlineto{\pgfqpoint{1.617271in}{2.743607in}}%
\pgfpathlineto{\pgfqpoint{1.617400in}{2.625005in}}%
\pgfpathlineto{\pgfqpoint{1.617994in}{2.711840in}}%
\pgfpathlineto{\pgfqpoint{1.618044in}{2.664800in}}%
\pgfpathlineto{\pgfqpoint{1.618687in}{2.755211in}}%
\pgfpathlineto{\pgfqpoint{1.618736in}{2.708419in}}%
\pgfpathlineto{\pgfqpoint{1.619238in}{2.789717in}}%
\pgfpathlineto{\pgfqpoint{1.619325in}{2.644559in}}%
\pgfpathlineto{\pgfqpoint{1.619421in}{2.644559in}}%
\pgfpathlineto{\pgfqpoint{1.622212in}{2.645165in}}%
\pgfpathlineto{\pgfqpoint{1.622215in}{2.747948in}}%
\pgfpathlineto{\pgfqpoint{1.622347in}{2.629120in}}%
\pgfpathlineto{\pgfqpoint{1.622942in}{2.708102in}}%
\pgfpathlineto{\pgfqpoint{1.622992in}{2.669031in}}%
\pgfpathlineto{\pgfqpoint{1.623635in}{2.751477in}}%
\pgfpathlineto{\pgfqpoint{1.623684in}{2.712646in}}%
\pgfpathlineto{\pgfqpoint{1.624184in}{2.785838in}}%
\pgfpathlineto{\pgfqpoint{1.624272in}{2.644559in}}%
\pgfpathlineto{\pgfqpoint{1.624367in}{2.644559in}}%
\pgfpathlineto{\pgfqpoint{1.627157in}{2.645117in}}%
\pgfpathlineto{\pgfqpoint{1.627162in}{2.743748in}}%
\pgfpathlineto{\pgfqpoint{1.627291in}{2.624943in}}%
\pgfpathlineto{\pgfqpoint{1.627885in}{2.711693in}}%
\pgfpathlineto{\pgfqpoint{1.627935in}{2.664730in}}%
\pgfpathlineto{\pgfqpoint{1.628578in}{2.755064in}}%
\pgfpathlineto{\pgfqpoint{1.628627in}{2.708350in}}%
\pgfpathlineto{\pgfqpoint{1.629129in}{2.789581in}}%
\pgfpathlineto{\pgfqpoint{1.629216in}{2.644559in}}%
\pgfpathlineto{\pgfqpoint{1.629311in}{2.644559in}}%
\pgfpathlineto{\pgfqpoint{1.632103in}{2.645149in}}%
\pgfpathlineto{\pgfqpoint{1.632107in}{2.743910in}}%
\pgfpathlineto{\pgfqpoint{1.632236in}{2.624260in}}%
\pgfpathlineto{\pgfqpoint{1.632831in}{2.712594in}}%
\pgfpathlineto{\pgfqpoint{1.632880in}{2.664052in}}%
\pgfpathlineto{\pgfqpoint{1.633523in}{2.755962in}}%
\pgfpathlineto{\pgfqpoint{1.633573in}{2.707674in}}%
\pgfpathlineto{\pgfqpoint{1.634074in}{2.790470in}}%
\pgfpathlineto{\pgfqpoint{1.634160in}{2.644559in}}%
\pgfpathlineto{\pgfqpoint{1.634253in}{2.644559in}}%
\pgfpathlineto{\pgfqpoint{1.637048in}{2.645113in}}%
\pgfpathlineto{\pgfqpoint{1.637053in}{2.745078in}}%
\pgfpathlineto{\pgfqpoint{1.637170in}{2.615536in}}%
\pgfpathlineto{\pgfqpoint{1.637805in}{2.654420in}}%
\pgfpathlineto{\pgfqpoint{1.638646in}{2.775128in}}%
\pgfpathlineto{\pgfqpoint{1.639020in}{2.798526in}}%
\pgfpathlineto{\pgfqpoint{1.639509in}{2.644559in}}%
\pgfpathlineto{\pgfqpoint{1.641994in}{2.644853in}}%
\pgfpathlineto{\pgfqpoint{1.641998in}{2.752014in}}%
\pgfpathlineto{\pgfqpoint{1.642087in}{2.629338in}}%
\pgfpathlineto{\pgfqpoint{1.642745in}{2.708470in}}%
\pgfpathlineto{\pgfqpoint{1.642794in}{2.670809in}}%
\pgfpathlineto{\pgfqpoint{1.643437in}{2.751888in}}%
\pgfpathlineto{\pgfqpoint{1.643486in}{2.714381in}}%
\pgfpathlineto{\pgfqpoint{1.643924in}{2.782409in}}%
\pgfpathlineto{\pgfqpoint{1.644056in}{2.644559in}}%
\pgfpathlineto{\pgfqpoint{1.644152in}{2.644559in}}%
\pgfpathlineto{\pgfqpoint{1.646939in}{2.644934in}}%
\pgfpathlineto{\pgfqpoint{1.646944in}{2.749239in}}%
\pgfpathlineto{\pgfqpoint{1.647050in}{2.618094in}}%
\pgfpathlineto{\pgfqpoint{1.647672in}{2.655609in}}%
\pgfpathlineto{\pgfqpoint{1.648513in}{2.770704in}}%
\pgfpathlineto{\pgfqpoint{1.648885in}{2.793992in}}%
\pgfpathlineto{\pgfqpoint{1.649348in}{2.644559in}}%
\pgfpathlineto{\pgfqpoint{1.651885in}{2.645077in}}%
\pgfpathlineto{\pgfqpoint{1.651889in}{2.745484in}}%
\pgfpathlineto{\pgfqpoint{1.652018in}{2.621511in}}%
\pgfpathlineto{\pgfqpoint{1.652612in}{2.714806in}}%
\pgfpathlineto{\pgfqpoint{1.652661in}{2.661277in}}%
\pgfpathlineto{\pgfqpoint{1.653304in}{2.758164in}}%
\pgfpathlineto{\pgfqpoint{1.653354in}{2.704910in}}%
\pgfpathlineto{\pgfqpoint{1.653856in}{2.792697in}}%
\pgfpathlineto{\pgfqpoint{1.653937in}{2.644559in}}%
\pgfpathlineto{\pgfqpoint{1.654065in}{2.644559in}}%
\pgfpathlineto{\pgfqpoint{1.656830in}{2.644850in}}%
\pgfpathlineto{\pgfqpoint{1.656834in}{2.752815in}}%
\pgfpathlineto{\pgfqpoint{1.656921in}{2.631416in}}%
\pgfpathlineto{\pgfqpoint{1.657577in}{2.706048in}}%
\pgfpathlineto{\pgfqpoint{1.657627in}{2.672519in}}%
\pgfpathlineto{\pgfqpoint{1.658270in}{2.749476in}}%
\pgfpathlineto{\pgfqpoint{1.658319in}{2.716083in}}%
\pgfpathlineto{\pgfqpoint{1.658758in}{2.780121in}}%
\pgfpathlineto{\pgfqpoint{1.658894in}{2.644559in}}%
\pgfpathlineto{\pgfqpoint{1.658973in}{2.644559in}}%
\pgfpathlineto{\pgfqpoint{1.661776in}{2.645210in}}%
\pgfpathlineto{\pgfqpoint{1.661780in}{2.743351in}}%
\pgfpathlineto{\pgfqpoint{1.661899in}{2.616058in}}%
\pgfpathlineto{\pgfqpoint{1.662535in}{2.655025in}}%
\pgfpathlineto{\pgfqpoint{1.663376in}{2.774293in}}%
\pgfpathlineto{\pgfqpoint{1.663747in}{2.797549in}}%
\pgfpathlineto{\pgfqpoint{1.664226in}{2.644559in}}%
\pgfpathlineto{\pgfqpoint{1.666721in}{2.644920in}}%
\pgfpathlineto{\pgfqpoint{1.666725in}{2.750004in}}%
\pgfpathlineto{\pgfqpoint{1.666829in}{2.622694in}}%
\pgfpathlineto{\pgfqpoint{1.667447in}{2.659753in}}%
\pgfpathlineto{\pgfqpoint{1.668288in}{2.765591in}}%
\pgfpathlineto{\pgfqpoint{1.668663in}{2.789114in}}%
\pgfpathlineto{\pgfqpoint{1.669178in}{2.644559in}}%
\pgfpathlineto{\pgfqpoint{1.671666in}{2.644562in}}%
\pgfpathlineto{\pgfqpoint{1.671669in}{2.747485in}}%
\pgfpathlineto{\pgfqpoint{1.671802in}{2.629073in}}%
\pgfpathlineto{\pgfqpoint{1.672445in}{2.668863in}}%
\pgfpathlineto{\pgfqpoint{1.673286in}{2.762797in}}%
\pgfpathlineto{\pgfqpoint{1.673638in}{2.784818in}}%
\pgfpathlineto{\pgfqpoint{1.674180in}{2.644559in}}%
\pgfpathlineto{\pgfqpoint{1.676611in}{2.644563in}}%
\pgfpathlineto{\pgfqpoint{1.676614in}{2.748997in}}%
\pgfpathlineto{\pgfqpoint{1.676747in}{2.628480in}}%
\pgfpathlineto{\pgfqpoint{1.677389in}{2.668208in}}%
\pgfpathlineto{\pgfqpoint{1.678230in}{2.763358in}}%
\pgfpathlineto{\pgfqpoint{1.678584in}{2.785480in}}%
\pgfpathlineto{\pgfqpoint{1.679120in}{2.644559in}}%
\pgfpathlineto{\pgfqpoint{1.681557in}{2.644569in}}%
\pgfpathlineto{\pgfqpoint{1.681559in}{2.758485in}}%
\pgfpathlineto{\pgfqpoint{1.681693in}{2.628868in}}%
\pgfpathlineto{\pgfqpoint{1.682285in}{2.707026in}}%
\pgfpathlineto{\pgfqpoint{1.682335in}{2.668586in}}%
\pgfpathlineto{\pgfqpoint{1.682978in}{2.750403in}}%
\pgfpathlineto{\pgfqpoint{1.683027in}{2.712199in}}%
\pgfpathlineto{\pgfqpoint{1.683529in}{2.784909in}}%
\pgfpathlineto{\pgfqpoint{1.683620in}{2.644559in}}%
\pgfpathlineto{\pgfqpoint{1.683716in}{2.644559in}}%
\pgfpathlineto{\pgfqpoint{1.686502in}{2.644567in}}%
\pgfpathlineto{\pgfqpoint{1.686505in}{2.756686in}}%
\pgfpathlineto{\pgfqpoint{1.686638in}{2.628882in}}%
\pgfpathlineto{\pgfqpoint{1.687280in}{2.668593in}}%
\pgfpathlineto{\pgfqpoint{1.688121in}{2.762827in}}%
\pgfpathlineto{\pgfqpoint{1.688474in}{2.784934in}}%
\pgfpathlineto{\pgfqpoint{1.689008in}{2.644559in}}%
\pgfpathlineto{\pgfqpoint{1.691448in}{2.644571in}}%
\pgfpathlineto{\pgfqpoint{1.691450in}{2.761178in}}%
\pgfpathlineto{\pgfqpoint{1.691584in}{2.629315in}}%
\pgfpathlineto{\pgfqpoint{1.692177in}{2.706592in}}%
\pgfpathlineto{\pgfqpoint{1.692227in}{2.669048in}}%
\pgfpathlineto{\pgfqpoint{1.692870in}{2.749970in}}%
\pgfpathlineto{\pgfqpoint{1.692919in}{2.712661in}}%
\pgfpathlineto{\pgfqpoint{1.693420in}{2.784417in}}%
\pgfpathlineto{\pgfqpoint{1.693511in}{2.644559in}}%
\pgfpathlineto{\pgfqpoint{1.693602in}{2.644559in}}%
\pgfpathlineto{\pgfqpoint{1.696392in}{2.644559in}}%
\pgfpathlineto{\pgfqpoint{1.696398in}{2.748558in}}%
\pgfpathlineto{\pgfqpoint{1.696506in}{2.637834in}}%
\pgfpathlineto{\pgfqpoint{1.697128in}{2.675208in}}%
\pgfpathlineto{\pgfqpoint{1.697969in}{2.750678in}}%
\pgfpathlineto{\pgfqpoint{1.698340in}{2.773980in}}%
\pgfpathlineto{\pgfqpoint{1.698994in}{2.644559in}}%
\pgfpathlineto{\pgfqpoint{1.701339in}{2.644571in}}%
\pgfpathlineto{\pgfqpoint{1.701341in}{2.761165in}}%
\pgfpathlineto{\pgfqpoint{1.701474in}{2.628970in}}%
\pgfpathlineto{\pgfqpoint{1.702067in}{2.706784in}}%
\pgfpathlineto{\pgfqpoint{1.702116in}{2.668652in}}%
\pgfpathlineto{\pgfqpoint{1.702759in}{2.750162in}}%
\pgfpathlineto{\pgfqpoint{1.702809in}{2.712265in}}%
\pgfpathlineto{\pgfqpoint{1.703311in}{2.784702in}}%
\pgfpathlineto{\pgfqpoint{1.703402in}{2.644559in}}%
\pgfpathlineto{\pgfqpoint{1.703494in}{2.644559in}}%
\pgfpathlineto{\pgfqpoint{1.706284in}{2.644579in}}%
\pgfpathlineto{\pgfqpoint{1.706287in}{2.765970in}}%
\pgfpathlineto{\pgfqpoint{1.706419in}{2.628778in}}%
\pgfpathlineto{\pgfqpoint{1.707011in}{2.706558in}}%
\pgfpathlineto{\pgfqpoint{1.707060in}{2.668366in}}%
\pgfpathlineto{\pgfqpoint{1.707703in}{2.749935in}}%
\pgfpathlineto{\pgfqpoint{1.707753in}{2.711980in}}%
\pgfpathlineto{\pgfqpoint{1.708256in}{2.784566in}}%
\pgfpathlineto{\pgfqpoint{1.708347in}{2.644559in}}%
\pgfpathlineto{\pgfqpoint{1.708441in}{2.644559in}}%
\pgfpathlineto{\pgfqpoint{1.711230in}{2.644580in}}%
\pgfpathlineto{\pgfqpoint{1.711232in}{2.766356in}}%
\pgfpathlineto{\pgfqpoint{1.711361in}{2.628366in}}%
\pgfpathlineto{\pgfqpoint{1.712000in}{2.667728in}}%
\pgfpathlineto{\pgfqpoint{1.712841in}{2.762339in}}%
\pgfpathlineto{\pgfqpoint{1.713202in}{2.784945in}}%
\pgfpathlineto{\pgfqpoint{1.713740in}{2.644559in}}%
\pgfpathlineto{\pgfqpoint{1.716175in}{2.644604in}}%
\pgfpathlineto{\pgfqpoint{1.716178in}{2.771909in}}%
\pgfpathlineto{\pgfqpoint{1.716296in}{2.632535in}}%
\pgfpathlineto{\pgfqpoint{1.716926in}{2.670893in}}%
\pgfpathlineto{\pgfqpoint{1.717767in}{2.756682in}}%
\pgfpathlineto{\pgfqpoint{1.718147in}{2.780497in}}%
\pgfpathlineto{\pgfqpoint{1.718714in}{2.644559in}}%
\pgfpathlineto{\pgfqpoint{1.721121in}{2.644583in}}%
\pgfpathlineto{\pgfqpoint{1.721123in}{2.767249in}}%
\pgfpathlineto{\pgfqpoint{1.721252in}{2.628591in}}%
\pgfpathlineto{\pgfqpoint{1.721891in}{2.667963in}}%
\pgfpathlineto{\pgfqpoint{1.722732in}{2.762227in}}%
\pgfpathlineto{\pgfqpoint{1.723093in}{2.784802in}}%
\pgfpathlineto{\pgfqpoint{1.723651in}{2.644559in}}%
\pgfpathlineto{\pgfqpoint{1.726066in}{2.644606in}}%
\pgfpathlineto{\pgfqpoint{1.726069in}{2.772092in}}%
\pgfpathlineto{\pgfqpoint{1.726189in}{2.630504in}}%
\pgfpathlineto{\pgfqpoint{1.726820in}{2.669013in}}%
\pgfpathlineto{\pgfqpoint{1.727661in}{2.758692in}}%
\pgfpathlineto{\pgfqpoint{1.728038in}{2.782312in}}%
\pgfpathlineto{\pgfqpoint{1.728596in}{2.644559in}}%
\pgfpathlineto{\pgfqpoint{1.731012in}{2.644579in}}%
\pgfpathlineto{\pgfqpoint{1.731014in}{2.765942in}}%
\pgfpathlineto{\pgfqpoint{1.731148in}{2.629173in}}%
\pgfpathlineto{\pgfqpoint{1.731740in}{2.706442in}}%
\pgfpathlineto{\pgfqpoint{1.731789in}{2.668834in}}%
\pgfpathlineto{\pgfqpoint{1.732432in}{2.749821in}}%
\pgfpathlineto{\pgfqpoint{1.732482in}{2.712446in}}%
\pgfpathlineto{\pgfqpoint{1.732984in}{2.784345in}}%
\pgfpathlineto{\pgfqpoint{1.733075in}{2.644559in}}%
\pgfpathlineto{\pgfqpoint{1.733168in}{2.644559in}}%
\pgfpathlineto{\pgfqpoint{1.735957in}{2.644736in}}%
\pgfpathlineto{\pgfqpoint{1.735960in}{2.777043in}}%
\pgfpathlineto{\pgfqpoint{1.736060in}{2.643482in}}%
\pgfpathlineto{\pgfqpoint{1.736720in}{2.693477in}}%
\pgfpathlineto{\pgfqpoint{1.736769in}{2.685557in}}%
\pgfpathlineto{\pgfqpoint{1.737214in}{2.724547in}}%
\pgfpathlineto{\pgfqpoint{1.737264in}{2.716623in}}%
\pgfpathlineto{\pgfqpoint{1.737893in}{2.767185in}}%
\pgfpathlineto{\pgfqpoint{1.737963in}{2.712331in}}%
\pgfpathlineto{\pgfqpoint{1.737972in}{2.731333in}}%
\pgfpathlineto{\pgfqpoint{1.738726in}{2.644559in}}%
\pgfpathlineto{\pgfqpoint{1.740903in}{2.644590in}}%
\pgfpathlineto{\pgfqpoint{1.740905in}{2.769469in}}%
\pgfpathlineto{\pgfqpoint{1.741029in}{2.626995in}}%
\pgfpathlineto{\pgfqpoint{1.741665in}{2.666017in}}%
\pgfpathlineto{\pgfqpoint{1.742506in}{2.762979in}}%
\pgfpathlineto{\pgfqpoint{1.742874in}{2.786032in}}%
\pgfpathlineto{\pgfqpoint{1.743417in}{2.644559in}}%
\pgfpathlineto{\pgfqpoint{1.745848in}{2.644852in}}%
\pgfpathlineto{\pgfqpoint{1.745851in}{2.777674in}}%
\pgfpathlineto{\pgfqpoint{1.745959in}{2.634429in}}%
\pgfpathlineto{\pgfqpoint{1.746575in}{2.671263in}}%
\pgfpathlineto{\pgfqpoint{1.747416in}{2.752757in}}%
\pgfpathlineto{\pgfqpoint{1.747791in}{2.776306in}}%
\pgfpathlineto{\pgfqpoint{1.748393in}{2.644559in}}%
\pgfpathlineto{\pgfqpoint{1.750793in}{2.644594in}}%
\pgfpathlineto{\pgfqpoint{1.750796in}{2.770297in}}%
\pgfpathlineto{\pgfqpoint{1.750919in}{2.627939in}}%
\pgfpathlineto{\pgfqpoint{1.751554in}{2.666836in}}%
\pgfpathlineto{\pgfqpoint{1.752395in}{2.761888in}}%
\pgfpathlineto{\pgfqpoint{1.752765in}{2.785083in}}%
\pgfpathlineto{\pgfqpoint{1.753296in}{2.644559in}}%
\pgfpathlineto{\pgfqpoint{1.755739in}{2.644837in}}%
\pgfpathlineto{\pgfqpoint{1.755742in}{2.777532in}}%
\pgfpathlineto{\pgfqpoint{1.756499in}{2.685053in}}%
\pgfpathlineto{\pgfqpoint{1.757711in}{2.766295in}}%
\pgfpathlineto{\pgfqpoint{1.757711in}{2.760021in}}%
\pgfpathlineto{\pgfqpoint{1.758437in}{2.644559in}}%
\pgfpathlineto{\pgfqpoint{1.760684in}{2.644611in}}%
\pgfpathlineto{\pgfqpoint{1.760687in}{2.772692in}}%
\pgfpathlineto{\pgfqpoint{1.760807in}{2.633211in}}%
\pgfpathlineto{\pgfqpoint{1.761438in}{2.671690in}}%
\pgfpathlineto{\pgfqpoint{1.762279in}{2.756091in}}%
\pgfpathlineto{\pgfqpoint{1.762656in}{2.779736in}}%
\pgfpathlineto{\pgfqpoint{1.763213in}{2.644559in}}%
\pgfpathlineto{\pgfqpoint{1.765630in}{2.644834in}}%
\pgfpathlineto{\pgfqpoint{1.765633in}{2.777531in}}%
\pgfpathlineto{\pgfqpoint{1.766381in}{2.684381in}}%
\pgfpathlineto{\pgfqpoint{1.767602in}{2.766397in}}%
\pgfpathlineto{\pgfqpoint{1.768371in}{2.644559in}}%
\pgfpathlineto{\pgfqpoint{1.770572in}{2.644559in}}%
\pgfpathlineto{\pgfqpoint{1.770578in}{2.765439in}}%
\pgfpathlineto{\pgfqpoint{1.770685in}{2.616551in}}%
\pgfpathlineto{\pgfqpoint{1.771351in}{2.720982in}}%
\pgfpathlineto{\pgfqpoint{1.771400in}{2.659753in}}%
\pgfpathlineto{\pgfqpoint{1.772043in}{2.764346in}}%
\pgfpathlineto{\pgfqpoint{1.772092in}{2.703378in}}%
\pgfpathlineto{\pgfqpoint{1.772518in}{2.794062in}}%
\pgfpathlineto{\pgfqpoint{1.772635in}{2.644559in}}%
\pgfpathlineto{\pgfqpoint{1.772776in}{2.644559in}}%
\pgfpathlineto{\pgfqpoint{1.775517in}{2.644559in}}%
\pgfpathlineto{\pgfqpoint{1.775524in}{2.768095in}}%
\pgfpathlineto{\pgfqpoint{1.775638in}{2.614786in}}%
\pgfpathlineto{\pgfqpoint{1.776313in}{2.723202in}}%
\pgfpathlineto{\pgfqpoint{1.776362in}{2.658948in}}%
\pgfpathlineto{\pgfqpoint{1.777005in}{2.766556in}}%
\pgfpathlineto{\pgfqpoint{1.777055in}{2.702584in}}%
\pgfpathlineto{\pgfqpoint{1.777493in}{2.797061in}}%
\pgfpathlineto{\pgfqpoint{1.777573in}{2.644559in}}%
\pgfpathlineto{\pgfqpoint{1.777743in}{2.644559in}}%
\pgfpathlineto{\pgfqpoint{1.780466in}{2.644792in}}%
\pgfpathlineto{\pgfqpoint{1.780469in}{2.777460in}}%
\pgfpathlineto{\pgfqpoint{1.781232in}{2.691155in}}%
\pgfpathlineto{\pgfqpoint{1.781282in}{2.688161in}}%
\pgfpathlineto{\pgfqpoint{1.781529in}{2.709801in}}%
\pgfpathlineto{\pgfqpoint{1.781578in}{2.706798in}}%
\pgfpathlineto{\pgfqpoint{1.782404in}{2.764779in}}%
\pgfpathlineto{\pgfqpoint{1.782461in}{2.748108in}}%
\pgfpathlineto{\pgfqpoint{1.783256in}{2.644559in}}%
\pgfpathlineto{\pgfqpoint{1.785412in}{2.645107in}}%
\pgfpathlineto{\pgfqpoint{1.785415in}{2.777400in}}%
\pgfpathlineto{\pgfqpoint{1.785539in}{2.626320in}}%
\pgfpathlineto{\pgfqpoint{1.786173in}{2.665080in}}%
\pgfpathlineto{\pgfqpoint{1.787014in}{2.762894in}}%
\pgfpathlineto{\pgfqpoint{1.787384in}{2.786086in}}%
\pgfpathlineto{\pgfqpoint{1.787922in}{2.644559in}}%
\pgfpathlineto{\pgfqpoint{1.790357in}{2.644767in}}%
\pgfpathlineto{\pgfqpoint{1.790360in}{2.777250in}}%
\pgfpathlineto{\pgfqpoint{1.790459in}{2.642510in}}%
\pgfpathlineto{\pgfqpoint{1.791117in}{2.694037in}}%
\pgfpathlineto{\pgfqpoint{1.791166in}{2.684346in}}%
\pgfpathlineto{\pgfqpoint{1.791611in}{2.725104in}}%
\pgfpathlineto{\pgfqpoint{1.791661in}{2.715416in}}%
\pgfpathlineto{\pgfqpoint{1.792292in}{2.767834in}}%
\pgfpathlineto{\pgfqpoint{1.792365in}{2.706269in}}%
\pgfpathlineto{\pgfqpoint{1.792374in}{2.736224in}}%
\pgfpathlineto{\pgfqpoint{1.793112in}{2.644559in}}%
\pgfpathlineto{\pgfqpoint{1.795300in}{2.644559in}}%
\pgfpathlineto{\pgfqpoint{1.795305in}{2.776851in}}%
\pgfpathlineto{\pgfqpoint{1.795394in}{2.621339in}}%
\pgfpathlineto{\pgfqpoint{1.796143in}{2.720137in}}%
\pgfpathlineto{\pgfqpoint{1.796192in}{2.668072in}}%
\pgfpathlineto{\pgfqpoint{1.796835in}{2.763533in}}%
\pgfpathlineto{\pgfqpoint{1.796885in}{2.711667in}}%
\pgfpathlineto{\pgfqpoint{1.797228in}{2.788147in}}%
\pgfpathlineto{\pgfqpoint{1.797365in}{2.644559in}}%
\pgfpathlineto{\pgfqpoint{1.797566in}{2.644559in}}%
\pgfpathlineto{\pgfqpoint{1.800245in}{2.644559in}}%
\pgfpathlineto{\pgfqpoint{1.800251in}{2.776128in}}%
\pgfpathlineto{\pgfqpoint{1.800348in}{2.619409in}}%
\pgfpathlineto{\pgfqpoint{1.801102in}{2.722467in}}%
\pgfpathlineto{\pgfqpoint{1.801151in}{2.667323in}}%
\pgfpathlineto{\pgfqpoint{1.801794in}{2.765852in}}%
\pgfpathlineto{\pgfqpoint{1.801843in}{2.710928in}}%
\pgfpathlineto{\pgfqpoint{1.802180in}{2.790036in}}%
\pgfpathlineto{\pgfqpoint{1.802312in}{2.644559in}}%
\pgfpathlineto{\pgfqpoint{1.802555in}{2.644559in}}%
\pgfpathlineto{\pgfqpoint{1.805191in}{2.644559in}}%
\pgfpathlineto{\pgfqpoint{1.805196in}{2.776045in}}%
\pgfpathlineto{\pgfqpoint{1.805293in}{2.619446in}}%
\pgfpathlineto{\pgfqpoint{1.806047in}{2.722404in}}%
\pgfpathlineto{\pgfqpoint{1.806096in}{2.667331in}}%
\pgfpathlineto{\pgfqpoint{1.806739in}{2.765789in}}%
\pgfpathlineto{\pgfqpoint{1.806789in}{2.710936in}}%
\pgfpathlineto{\pgfqpoint{1.807125in}{2.789984in}}%
\pgfpathlineto{\pgfqpoint{1.807257in}{2.644559in}}%
\pgfpathlineto{\pgfqpoint{1.807486in}{2.644559in}}%
\pgfpathlineto{\pgfqpoint{1.810136in}{2.644559in}}%
\pgfpathlineto{\pgfqpoint{1.810142in}{2.776654in}}%
\pgfpathlineto{\pgfqpoint{1.810274in}{2.624117in}}%
\pgfpathlineto{\pgfqpoint{1.810963in}{2.716294in}}%
\pgfpathlineto{\pgfqpoint{1.811013in}{2.669699in}}%
\pgfpathlineto{\pgfqpoint{1.811656in}{2.759703in}}%
\pgfpathlineto{\pgfqpoint{1.811705in}{2.713279in}}%
\pgfpathlineto{\pgfqpoint{1.812111in}{2.788229in}}%
\pgfpathlineto{\pgfqpoint{1.812197in}{2.644559in}}%
\pgfpathlineto{\pgfqpoint{1.812368in}{2.644559in}}%
\pgfpathlineto{\pgfqpoint{1.815085in}{2.645060in}}%
\pgfpathlineto{\pgfqpoint{1.815087in}{2.777410in}}%
\pgfpathlineto{\pgfqpoint{1.815180in}{2.630969in}}%
\pgfpathlineto{\pgfqpoint{1.815834in}{2.704734in}}%
\pgfpathlineto{\pgfqpoint{1.815883in}{2.672079in}}%
\pgfpathlineto{\pgfqpoint{1.816526in}{2.748175in}}%
\pgfpathlineto{\pgfqpoint{1.816576in}{2.715628in}}%
\pgfpathlineto{\pgfqpoint{1.817014in}{2.778762in}}%
\pgfpathlineto{\pgfqpoint{1.817139in}{2.644559in}}%
\pgfpathlineto{\pgfqpoint{1.817252in}{2.644559in}}%
\pgfpathlineto{\pgfqpoint{1.820027in}{2.644559in}}%
\pgfpathlineto{\pgfqpoint{1.820033in}{2.776080in}}%
\pgfpathlineto{\pgfqpoint{1.820128in}{2.619442in}}%
\pgfpathlineto{\pgfqpoint{1.820881in}{2.722240in}}%
\pgfpathlineto{\pgfqpoint{1.820931in}{2.667172in}}%
\pgfpathlineto{\pgfqpoint{1.821574in}{2.765626in}}%
\pgfpathlineto{\pgfqpoint{1.821623in}{2.710776in}}%
\pgfpathlineto{\pgfqpoint{1.821961in}{2.789880in}}%
\pgfpathlineto{\pgfqpoint{1.822094in}{2.644559in}}%
\pgfpathlineto{\pgfqpoint{1.822331in}{2.644559in}}%
\pgfpathlineto{\pgfqpoint{1.824972in}{2.644559in}}%
\pgfpathlineto{\pgfqpoint{1.824978in}{2.775485in}}%
\pgfpathlineto{\pgfqpoint{1.825071in}{2.619715in}}%
\pgfpathlineto{\pgfqpoint{1.825822in}{2.721568in}}%
\pgfpathlineto{\pgfqpoint{1.825871in}{2.666988in}}%
\pgfpathlineto{\pgfqpoint{1.826514in}{2.764957in}}%
\pgfpathlineto{\pgfqpoint{1.826563in}{2.710589in}}%
\pgfpathlineto{\pgfqpoint{1.826904in}{2.789372in}}%
\pgfpathlineto{\pgfqpoint{1.827039in}{2.644559in}}%
\pgfpathlineto{\pgfqpoint{1.827281in}{2.644559in}}%
\pgfpathlineto{\pgfqpoint{1.829918in}{2.644559in}}%
\pgfpathlineto{\pgfqpoint{1.829924in}{2.775781in}}%
\pgfpathlineto{\pgfqpoint{1.830016in}{2.619909in}}%
\pgfpathlineto{\pgfqpoint{1.830766in}{2.721356in}}%
\pgfpathlineto{\pgfqpoint{1.830815in}{2.667054in}}%
\pgfpathlineto{\pgfqpoint{1.831458in}{2.764745in}}%
\pgfpathlineto{\pgfqpoint{1.831507in}{2.710654in}}%
\pgfpathlineto{\pgfqpoint{1.831848in}{2.789210in}}%
\pgfpathlineto{\pgfqpoint{1.831984in}{2.644559in}}%
\pgfpathlineto{\pgfqpoint{1.832215in}{2.644559in}}%
\pgfpathlineto{\pgfqpoint{1.834863in}{2.644559in}}%
\pgfpathlineto{\pgfqpoint{1.834869in}{2.774909in}}%
\pgfpathlineto{\pgfqpoint{1.834969in}{2.619665in}}%
\pgfpathlineto{\pgfqpoint{1.835626in}{2.715868in}}%
\pgfpathlineto{\pgfqpoint{1.835675in}{2.661679in}}%
\pgfpathlineto{\pgfqpoint{1.836318in}{2.759253in}}%
\pgfpathlineto{\pgfqpoint{1.836367in}{2.705284in}}%
\pgfpathlineto{\pgfqpoint{1.836801in}{2.789486in}}%
\pgfpathlineto{\pgfqpoint{1.836930in}{2.644559in}}%
\pgfpathlineto{\pgfqpoint{1.837077in}{2.644559in}}%
\pgfpathlineto{\pgfqpoint{1.839809in}{2.644559in}}%
\pgfpathlineto{\pgfqpoint{1.839815in}{2.773732in}}%
\pgfpathlineto{\pgfqpoint{1.839925in}{2.622305in}}%
\pgfpathlineto{\pgfqpoint{1.840587in}{2.714218in}}%
\pgfpathlineto{\pgfqpoint{1.840637in}{2.665405in}}%
\pgfpathlineto{\pgfqpoint{1.841280in}{2.757610in}}%
\pgfpathlineto{\pgfqpoint{1.841329in}{2.709004in}}%
\pgfpathlineto{\pgfqpoint{1.841754in}{2.787339in}}%
\pgfpathlineto{\pgfqpoint{1.841875in}{2.644559in}}%
\pgfpathlineto{\pgfqpoint{1.842014in}{2.644559in}}%
\pgfpathlineto{\pgfqpoint{1.844754in}{2.644559in}}%
\pgfpathlineto{\pgfqpoint{1.844760in}{2.774205in}}%
\pgfpathlineto{\pgfqpoint{1.844864in}{2.620563in}}%
\pgfpathlineto{\pgfqpoint{1.845522in}{2.715175in}}%
\pgfpathlineto{\pgfqpoint{1.845572in}{2.662955in}}%
\pgfpathlineto{\pgfqpoint{1.846214in}{2.758563in}}%
\pgfpathlineto{\pgfqpoint{1.846264in}{2.706558in}}%
\pgfpathlineto{\pgfqpoint{1.846694in}{2.788625in}}%
\pgfpathlineto{\pgfqpoint{1.846821in}{2.644559in}}%
\pgfpathlineto{\pgfqpoint{1.846962in}{2.644559in}}%
\pgfpathlineto{\pgfqpoint{1.849700in}{2.644559in}}%
\pgfpathlineto{\pgfqpoint{1.849706in}{2.774708in}}%
\pgfpathlineto{\pgfqpoint{1.849800in}{2.619551in}}%
\pgfpathlineto{\pgfqpoint{1.850550in}{2.721682in}}%
\pgfpathlineto{\pgfqpoint{1.850600in}{2.666907in}}%
\pgfpathlineto{\pgfqpoint{1.851242in}{2.765070in}}%
\pgfpathlineto{\pgfqpoint{1.851292in}{2.710509in}}%
\pgfpathlineto{\pgfqpoint{1.851632in}{2.789448in}}%
\pgfpathlineto{\pgfqpoint{1.851766in}{2.644559in}}%
\pgfpathlineto{\pgfqpoint{1.851999in}{2.644559in}}%
\pgfpathlineto{\pgfqpoint{1.854646in}{2.644559in}}%
\pgfpathlineto{\pgfqpoint{1.854652in}{2.752042in}}%
\pgfpathlineto{\pgfqpoint{1.854757in}{2.620746in}}%
\pgfpathlineto{\pgfqpoint{1.855419in}{2.715330in}}%
\pgfpathlineto{\pgfqpoint{1.855468in}{2.663513in}}%
\pgfpathlineto{\pgfqpoint{1.856111in}{2.758716in}}%
\pgfpathlineto{\pgfqpoint{1.856161in}{2.707118in}}%
\pgfpathlineto{\pgfqpoint{1.856588in}{2.788590in}}%
\pgfpathlineto{\pgfqpoint{1.856712in}{2.644559in}}%
\pgfpathlineto{\pgfqpoint{1.856863in}{2.644559in}}%
\pgfpathlineto{\pgfqpoint{1.859591in}{2.644559in}}%
\pgfpathlineto{\pgfqpoint{1.859597in}{2.773462in}}%
\pgfpathlineto{\pgfqpoint{1.859702in}{2.621621in}}%
\pgfpathlineto{\pgfqpoint{1.860362in}{2.714238in}}%
\pgfpathlineto{\pgfqpoint{1.860411in}{2.664228in}}%
\pgfpathlineto{\pgfqpoint{1.861054in}{2.757629in}}%
\pgfpathlineto{\pgfqpoint{1.861104in}{2.707828in}}%
\pgfpathlineto{\pgfqpoint{1.861532in}{2.787585in}}%
\pgfpathlineto{\pgfqpoint{1.861657in}{2.644559in}}%
\pgfpathlineto{\pgfqpoint{1.861811in}{2.644559in}}%
\pgfpathlineto{\pgfqpoint{1.864536in}{2.644559in}}%
\pgfpathlineto{\pgfqpoint{1.864542in}{2.771526in}}%
\pgfpathlineto{\pgfqpoint{1.864651in}{2.616435in}}%
\pgfpathlineto{\pgfqpoint{1.865316in}{2.719945in}}%
\pgfpathlineto{\pgfqpoint{1.865366in}{2.659592in}}%
\pgfpathlineto{\pgfqpoint{1.866009in}{2.763311in}}%
\pgfpathlineto{\pgfqpoint{1.866058in}{2.703217in}}%
\pgfpathlineto{\pgfqpoint{1.866482in}{2.792968in}}%
\pgfpathlineto{\pgfqpoint{1.866599in}{2.644559in}}%
\pgfpathlineto{\pgfqpoint{1.866740in}{2.644559in}}%
\pgfpathlineto{\pgfqpoint{1.869482in}{2.644559in}}%
\pgfpathlineto{\pgfqpoint{1.869488in}{2.771575in}}%
\pgfpathlineto{\pgfqpoint{1.869595in}{2.616973in}}%
\pgfpathlineto{\pgfqpoint{1.870259in}{2.719292in}}%
\pgfpathlineto{\pgfqpoint{1.870309in}{2.659975in}}%
\pgfpathlineto{\pgfqpoint{1.870952in}{2.762661in}}%
\pgfpathlineto{\pgfqpoint{1.871001in}{2.703597in}}%
\pgfpathlineto{\pgfqpoint{1.871427in}{2.792392in}}%
\pgfpathlineto{\pgfqpoint{1.871547in}{2.644559in}}%
\pgfpathlineto{\pgfqpoint{1.871690in}{2.644559in}}%
\pgfpathlineto{\pgfqpoint{1.874427in}{2.644559in}}%
\pgfpathlineto{\pgfqpoint{1.874433in}{2.772410in}}%
\pgfpathlineto{\pgfqpoint{1.874539in}{2.623813in}}%
\pgfpathlineto{\pgfqpoint{1.875200in}{2.712031in}}%
\pgfpathlineto{\pgfqpoint{1.875249in}{2.666455in}}%
\pgfpathlineto{\pgfqpoint{1.875892in}{2.755428in}}%
\pgfpathlineto{\pgfqpoint{1.875942in}{2.710048in}}%
\pgfpathlineto{\pgfqpoint{1.876369in}{2.785341in}}%
\pgfpathlineto{\pgfqpoint{1.876492in}{2.644559in}}%
\pgfpathlineto{\pgfqpoint{1.876640in}{2.644559in}}%
\pgfpathlineto{\pgfqpoint{1.879372in}{2.644559in}}%
\pgfpathlineto{\pgfqpoint{1.879379in}{2.771020in}}%
\pgfpathlineto{\pgfqpoint{1.879488in}{2.616442in}}%
\pgfpathlineto{\pgfqpoint{1.880152in}{2.719878in}}%
\pgfpathlineto{\pgfqpoint{1.880202in}{2.659580in}}%
\pgfpathlineto{\pgfqpoint{1.880845in}{2.763245in}}%
\pgfpathlineto{\pgfqpoint{1.880894in}{2.703204in}}%
\pgfpathlineto{\pgfqpoint{1.881319in}{2.792902in}}%
\pgfpathlineto{\pgfqpoint{1.881435in}{2.644559in}}%
\pgfpathlineto{\pgfqpoint{1.881577in}{2.644559in}}%
\pgfpathlineto{\pgfqpoint{1.884318in}{2.644559in}}%
\pgfpathlineto{\pgfqpoint{1.884324in}{2.770452in}}%
\pgfpathlineto{\pgfqpoint{1.884435in}{2.615990in}}%
\pgfpathlineto{\pgfqpoint{1.885101in}{2.720446in}}%
\pgfpathlineto{\pgfqpoint{1.885150in}{2.659302in}}%
\pgfpathlineto{\pgfqpoint{1.885793in}{2.763810in}}%
\pgfpathlineto{\pgfqpoint{1.885843in}{2.702929in}}%
\pgfpathlineto{\pgfqpoint{1.886265in}{2.793374in}}%
\pgfpathlineto{\pgfqpoint{1.886379in}{2.644559in}}%
\pgfpathlineto{\pgfqpoint{1.886570in}{2.644559in}}%
\pgfpathlineto{\pgfqpoint{1.889263in}{2.644559in}}%
\pgfpathlineto{\pgfqpoint{1.889270in}{2.772308in}}%
\pgfpathlineto{\pgfqpoint{1.889375in}{2.623252in}}%
\pgfpathlineto{\pgfqpoint{1.890035in}{2.712538in}}%
\pgfpathlineto{\pgfqpoint{1.890085in}{2.665858in}}%
\pgfpathlineto{\pgfqpoint{1.890728in}{2.755933in}}%
\pgfpathlineto{\pgfqpoint{1.890777in}{2.709453in}}%
\pgfpathlineto{\pgfqpoint{1.891205in}{2.785864in}}%
\pgfpathlineto{\pgfqpoint{1.891329in}{2.644559in}}%
\pgfpathlineto{\pgfqpoint{1.891475in}{2.644559in}}%
\pgfpathlineto{\pgfqpoint{1.894209in}{2.644559in}}%
\pgfpathlineto{\pgfqpoint{1.894215in}{2.771213in}}%
\pgfpathlineto{\pgfqpoint{1.894325in}{2.627788in}}%
\pgfpathlineto{\pgfqpoint{1.894987in}{2.708207in}}%
\pgfpathlineto{\pgfqpoint{1.895036in}{2.670767in}}%
\pgfpathlineto{\pgfqpoint{1.895679in}{2.751618in}}%
\pgfpathlineto{\pgfqpoint{1.895729in}{2.714346in}}%
\pgfpathlineto{\pgfqpoint{1.896154in}{2.781377in}}%
\pgfpathlineto{\pgfqpoint{1.896274in}{2.644559in}}%
\pgfpathlineto{\pgfqpoint{1.896375in}{2.644559in}}%
\pgfpathlineto{\pgfqpoint{1.899154in}{2.644559in}}%
\pgfpathlineto{\pgfqpoint{1.899160in}{2.771543in}}%
\pgfpathlineto{\pgfqpoint{1.899266in}{2.624535in}}%
\pgfpathlineto{\pgfqpoint{1.899927in}{2.711247in}}%
\pgfpathlineto{\pgfqpoint{1.899977in}{2.667161in}}%
\pgfpathlineto{\pgfqpoint{1.900620in}{2.754646in}}%
\pgfpathlineto{\pgfqpoint{1.900669in}{2.710752in}}%
\pgfpathlineto{\pgfqpoint{1.901097in}{2.784553in}}%
\pgfpathlineto{\pgfqpoint{1.901218in}{2.644559in}}%
\pgfpathlineto{\pgfqpoint{1.901365in}{2.644559in}}%
\pgfpathlineto{\pgfqpoint{1.904100in}{2.644559in}}%
\pgfpathlineto{\pgfqpoint{1.904106in}{2.770996in}}%
\pgfpathlineto{\pgfqpoint{1.904216in}{2.627403in}}%
\pgfpathlineto{\pgfqpoint{1.904877in}{2.708526in}}%
\pgfpathlineto{\pgfqpoint{1.904926in}{2.670336in}}%
\pgfpathlineto{\pgfqpoint{1.905569in}{2.751936in}}%
\pgfpathlineto{\pgfqpoint{1.905619in}{2.713917in}}%
\pgfpathlineto{\pgfqpoint{1.906044in}{2.781714in}}%
\pgfpathlineto{\pgfqpoint{1.906163in}{2.644559in}}%
\pgfpathlineto{\pgfqpoint{1.906312in}{2.644559in}}%
\pgfpathlineto{\pgfqpoint{1.909045in}{2.644559in}}%
\pgfpathlineto{\pgfqpoint{1.909051in}{2.770896in}}%
\pgfpathlineto{\pgfqpoint{1.909162in}{2.627914in}}%
\pgfpathlineto{\pgfqpoint{1.909823in}{2.708041in}}%
\pgfpathlineto{\pgfqpoint{1.909872in}{2.670877in}}%
\pgfpathlineto{\pgfqpoint{1.910515in}{2.751452in}}%
\pgfpathlineto{\pgfqpoint{1.910565in}{2.714456in}}%
\pgfpathlineto{\pgfqpoint{1.910990in}{2.781215in}}%
\pgfpathlineto{\pgfqpoint{1.911111in}{2.644559in}}%
\pgfpathlineto{\pgfqpoint{1.911212in}{2.644559in}}%
\pgfpathlineto{\pgfqpoint{1.913991in}{2.644559in}}%
\pgfpathlineto{\pgfqpoint{1.913997in}{2.770591in}}%
\pgfpathlineto{\pgfqpoint{1.914107in}{2.627471in}}%
\pgfpathlineto{\pgfqpoint{1.914768in}{2.708449in}}%
\pgfpathlineto{\pgfqpoint{1.914818in}{2.670429in}}%
\pgfpathlineto{\pgfqpoint{1.915461in}{2.751859in}}%
\pgfpathlineto{\pgfqpoint{1.915510in}{2.714010in}}%
\pgfpathlineto{\pgfqpoint{1.915936in}{2.781618in}}%
\pgfpathlineto{\pgfqpoint{1.916056in}{2.644559in}}%
\pgfpathlineto{\pgfqpoint{1.916159in}{2.644559in}}%
\pgfpathlineto{\pgfqpoint{1.918936in}{2.644559in}}%
\pgfpathlineto{\pgfqpoint{1.918943in}{2.768822in}}%
\pgfpathlineto{\pgfqpoint{1.919057in}{2.614920in}}%
\pgfpathlineto{\pgfqpoint{1.919728in}{2.721881in}}%
\pgfpathlineto{\pgfqpoint{1.919778in}{2.658779in}}%
\pgfpathlineto{\pgfqpoint{1.920421in}{2.765238in}}%
\pgfpathlineto{\pgfqpoint{1.920470in}{2.702412in}}%
\pgfpathlineto{\pgfqpoint{1.920911in}{2.795916in}}%
\pgfpathlineto{\pgfqpoint{1.920992in}{2.644559in}}%
\pgfpathlineto{\pgfqpoint{1.921188in}{2.644559in}}%
\pgfpathlineto{\pgfqpoint{1.923882in}{2.644559in}}%
\pgfpathlineto{\pgfqpoint{1.923888in}{2.768978in}}%
\pgfpathlineto{\pgfqpoint{1.924003in}{2.614752in}}%
\pgfpathlineto{\pgfqpoint{1.924674in}{2.722090in}}%
\pgfpathlineto{\pgfqpoint{1.924724in}{2.658639in}}%
\pgfpathlineto{\pgfqpoint{1.925367in}{2.765447in}}%
\pgfpathlineto{\pgfqpoint{1.925416in}{2.702273in}}%
\pgfpathlineto{\pgfqpoint{1.925856in}{2.796103in}}%
\pgfpathlineto{\pgfqpoint{1.925937in}{2.644559in}}%
\pgfpathlineto{\pgfqpoint{1.926140in}{2.644559in}}%
\pgfpathlineto{\pgfqpoint{1.928829in}{2.644559in}}%
\pgfpathlineto{\pgfqpoint{1.928834in}{2.750394in}}%
\pgfpathlineto{\pgfqpoint{1.928921in}{2.643491in}}%
\pgfpathlineto{\pgfqpoint{1.929568in}{2.690666in}}%
\pgfpathlineto{\pgfqpoint{1.929618in}{2.683508in}}%
\pgfpathlineto{\pgfqpoint{1.930063in}{2.721711in}}%
\pgfpathlineto{\pgfqpoint{1.930112in}{2.714599in}}%
\pgfpathlineto{\pgfqpoint{1.930755in}{2.765172in}}%
\pgfpathlineto{\pgfqpoint{1.930837in}{2.705392in}}%
\pgfpathlineto{\pgfqpoint{1.930846in}{2.735561in}}%
\pgfpathlineto{\pgfqpoint{1.931602in}{2.644559in}}%
\pgfpathlineto{\pgfqpoint{1.933773in}{2.644559in}}%
\pgfpathlineto{\pgfqpoint{1.933780in}{2.752341in}}%
\pgfpathlineto{\pgfqpoint{1.933871in}{2.626909in}}%
\pgfpathlineto{\pgfqpoint{1.934521in}{2.707496in}}%
\pgfpathlineto{\pgfqpoint{1.934570in}{2.667711in}}%
\pgfpathlineto{\pgfqpoint{1.935213in}{2.750908in}}%
\pgfpathlineto{\pgfqpoint{1.935263in}{2.711289in}}%
\pgfpathlineto{\pgfqpoint{1.935703in}{2.781597in}}%
\pgfpathlineto{\pgfqpoint{1.935838in}{2.644559in}}%
\pgfpathlineto{\pgfqpoint{1.935936in}{2.644559in}}%
\pgfpathlineto{\pgfqpoint{1.938718in}{2.644559in}}%
\pgfpathlineto{\pgfqpoint{1.938724in}{2.768383in}}%
\pgfpathlineto{\pgfqpoint{1.938838in}{2.615100in}}%
\pgfpathlineto{\pgfqpoint{1.939509in}{2.721289in}}%
\pgfpathlineto{\pgfqpoint{1.939558in}{2.658822in}}%
\pgfpathlineto{\pgfqpoint{1.940201in}{2.764648in}}%
\pgfpathlineto{\pgfqpoint{1.940250in}{2.702454in}}%
\pgfpathlineto{\pgfqpoint{1.940693in}{2.795440in}}%
\pgfpathlineto{\pgfqpoint{1.940774in}{2.644559in}}%
\pgfpathlineto{\pgfqpoint{1.940962in}{2.644559in}}%
\pgfpathlineto{\pgfqpoint{1.943663in}{2.644559in}}%
\pgfpathlineto{\pgfqpoint{1.943670in}{2.770771in}}%
\pgfpathlineto{\pgfqpoint{1.943780in}{2.622959in}}%
\pgfpathlineto{\pgfqpoint{1.944441in}{2.712938in}}%
\pgfpathlineto{\pgfqpoint{1.944490in}{2.665893in}}%
\pgfpathlineto{\pgfqpoint{1.945133in}{2.756333in}}%
\pgfpathlineto{\pgfqpoint{1.945182in}{2.709488in}}%
\pgfpathlineto{\pgfqpoint{1.945608in}{2.786100in}}%
\pgfpathlineto{\pgfqpoint{1.945729in}{2.644559in}}%
\pgfpathlineto{\pgfqpoint{1.945874in}{2.644559in}}%
\pgfpathlineto{\pgfqpoint{1.948609in}{2.644559in}}%
\pgfpathlineto{\pgfqpoint{1.948615in}{2.766344in}}%
\pgfpathlineto{\pgfqpoint{1.948739in}{2.614664in}}%
\pgfpathlineto{\pgfqpoint{1.949366in}{2.653046in}}%
\pgfpathlineto{\pgfqpoint{1.950207in}{2.772281in}}%
\pgfpathlineto{\pgfqpoint{1.950584in}{2.795843in}}%
\pgfpathlineto{\pgfqpoint{1.951035in}{2.644559in}}%
\pgfpathlineto{\pgfqpoint{1.953554in}{2.644559in}}%
\pgfpathlineto{\pgfqpoint{1.953561in}{2.769436in}}%
\pgfpathlineto{\pgfqpoint{1.953672in}{2.615935in}}%
\pgfpathlineto{\pgfqpoint{1.954338in}{2.720291in}}%
\pgfpathlineto{\pgfqpoint{1.954387in}{2.659246in}}%
\pgfpathlineto{\pgfqpoint{1.955030in}{2.763655in}}%
\pgfpathlineto{\pgfqpoint{1.955080in}{2.702873in}}%
\pgfpathlineto{\pgfqpoint{1.955502in}{2.793201in}}%
\pgfpathlineto{\pgfqpoint{1.955615in}{2.644559in}}%
\pgfpathlineto{\pgfqpoint{1.955806in}{2.644559in}}%
\pgfpathlineto{\pgfqpoint{1.958500in}{2.644559in}}%
\pgfpathlineto{\pgfqpoint{1.958506in}{2.766662in}}%
\pgfpathlineto{\pgfqpoint{1.958626in}{2.617270in}}%
\pgfpathlineto{\pgfqpoint{1.959251in}{2.655284in}}%
\pgfpathlineto{\pgfqpoint{1.960092in}{2.769323in}}%
\pgfpathlineto{\pgfqpoint{1.960474in}{2.793292in}}%
\pgfpathlineto{\pgfqpoint{1.960921in}{2.644559in}}%
\pgfpathlineto{\pgfqpoint{1.963445in}{2.644559in}}%
\pgfpathlineto{\pgfqpoint{1.963452in}{2.768801in}}%
\pgfpathlineto{\pgfqpoint{1.963564in}{2.616026in}}%
\pgfpathlineto{\pgfqpoint{1.964233in}{2.720243in}}%
\pgfpathlineto{\pgfqpoint{1.964282in}{2.659565in}}%
\pgfpathlineto{\pgfqpoint{1.964925in}{2.763607in}}%
\pgfpathlineto{\pgfqpoint{1.964975in}{2.703191in}}%
\pgfpathlineto{\pgfqpoint{1.965395in}{2.793024in}}%
\pgfpathlineto{\pgfqpoint{1.965506in}{2.644559in}}%
\pgfpathlineto{\pgfqpoint{1.965697in}{2.644559in}}%
\pgfpathlineto{\pgfqpoint{1.968391in}{2.644559in}}%
\pgfpathlineto{\pgfqpoint{1.968397in}{2.768188in}}%
\pgfpathlineto{\pgfqpoint{1.968509in}{2.616447in}}%
\pgfpathlineto{\pgfqpoint{1.969178in}{2.719699in}}%
\pgfpathlineto{\pgfqpoint{1.969227in}{2.659939in}}%
\pgfpathlineto{\pgfqpoint{1.969870in}{2.763065in}}%
\pgfpathlineto{\pgfqpoint{1.969920in}{2.703564in}}%
\pgfpathlineto{\pgfqpoint{1.970340in}{2.792493in}}%
\pgfpathlineto{\pgfqpoint{1.970452in}{2.644559in}}%
\pgfpathlineto{\pgfqpoint{1.970644in}{2.644559in}}%
\pgfpathlineto{\pgfqpoint{1.973336in}{2.644559in}}%
\pgfpathlineto{\pgfqpoint{1.973343in}{2.768314in}}%
\pgfpathlineto{\pgfqpoint{1.973454in}{2.616296in}}%
\pgfpathlineto{\pgfqpoint{1.974120in}{2.719645in}}%
\pgfpathlineto{\pgfqpoint{1.974169in}{2.659575in}}%
\pgfpathlineto{\pgfqpoint{1.974812in}{2.763010in}}%
\pgfpathlineto{\pgfqpoint{1.974862in}{2.703201in}}%
\pgfpathlineto{\pgfqpoint{1.975284in}{2.792549in}}%
\pgfpathlineto{\pgfqpoint{1.975397in}{2.644559in}}%
\pgfpathlineto{\pgfqpoint{1.975589in}{2.644559in}}%
\pgfpathlineto{\pgfqpoint{1.978282in}{2.644559in}}%
\pgfpathlineto{\pgfqpoint{1.978288in}{2.767916in}}%
\pgfpathlineto{\pgfqpoint{1.978402in}{2.615820in}}%
\pgfpathlineto{\pgfqpoint{1.979071in}{2.720412in}}%
\pgfpathlineto{\pgfqpoint{1.979120in}{2.659443in}}%
\pgfpathlineto{\pgfqpoint{1.979763in}{2.763775in}}%
\pgfpathlineto{\pgfqpoint{1.979813in}{2.703071in}}%
\pgfpathlineto{\pgfqpoint{1.980256in}{2.794644in}}%
\pgfpathlineto{\pgfqpoint{1.980340in}{2.644559in}}%
\pgfpathlineto{\pgfqpoint{1.980505in}{2.644559in}}%
\pgfpathlineto{\pgfqpoint{1.983227in}{2.644559in}}%
\pgfpathlineto{\pgfqpoint{1.983234in}{2.766393in}}%
\pgfpathlineto{\pgfqpoint{1.983356in}{2.615562in}}%
\pgfpathlineto{\pgfqpoint{1.983982in}{2.653822in}}%
\pgfpathlineto{\pgfqpoint{1.984823in}{2.771260in}}%
\pgfpathlineto{\pgfqpoint{1.985202in}{2.794953in}}%
\pgfpathlineto{\pgfqpoint{1.985673in}{2.644559in}}%
\pgfpathlineto{\pgfqpoint{1.988173in}{2.644559in}}%
\pgfpathlineto{\pgfqpoint{1.988179in}{2.768080in}}%
\pgfpathlineto{\pgfqpoint{1.988292in}{2.616164in}}%
\pgfpathlineto{\pgfqpoint{1.988961in}{2.720053in}}%
\pgfpathlineto{\pgfqpoint{1.989010in}{2.659716in}}%
\pgfpathlineto{\pgfqpoint{1.989653in}{2.763417in}}%
\pgfpathlineto{\pgfqpoint{1.989703in}{2.703343in}}%
\pgfpathlineto{\pgfqpoint{1.990147in}{2.794362in}}%
\pgfpathlineto{\pgfqpoint{1.990232in}{2.644559in}}%
\pgfpathlineto{\pgfqpoint{1.990424in}{2.644559in}}%
\pgfpathlineto{\pgfqpoint{1.993118in}{2.644559in}}%
\pgfpathlineto{\pgfqpoint{1.993125in}{2.766916in}}%
\pgfpathlineto{\pgfqpoint{1.993242in}{2.615051in}}%
\pgfpathlineto{\pgfqpoint{1.993864in}{2.652808in}}%
\pgfpathlineto{\pgfqpoint{1.994705in}{2.770993in}}%
\pgfpathlineto{\pgfqpoint{1.995093in}{2.795251in}}%
\pgfpathlineto{\pgfqpoint{1.995568in}{2.644559in}}%
\pgfpathlineto{\pgfqpoint{1.998064in}{2.644559in}}%
\pgfpathlineto{\pgfqpoint{1.998070in}{2.766647in}}%
\pgfpathlineto{\pgfqpoint{1.998190in}{2.616834in}}%
\pgfpathlineto{\pgfqpoint{1.998816in}{2.654908in}}%
\pgfpathlineto{\pgfqpoint{1.999656in}{2.769804in}}%
\pgfpathlineto{\pgfqpoint{2.000038in}{2.793707in}}%
\pgfpathlineto{\pgfqpoint{2.000500in}{2.644559in}}%
\pgfpathlineto{\pgfqpoint{2.003009in}{2.644559in}}%
\pgfpathlineto{\pgfqpoint{2.003016in}{2.765481in}}%
\pgfpathlineto{\pgfqpoint{2.003141in}{2.614581in}}%
\pgfpathlineto{\pgfqpoint{2.003770in}{2.653128in}}%
\pgfpathlineto{\pgfqpoint{2.004611in}{2.772451in}}%
\pgfpathlineto{\pgfqpoint{2.004984in}{2.795775in}}%
\pgfpathlineto{\pgfqpoint{2.005450in}{2.644559in}}%
\pgfpathlineto{\pgfqpoint{2.007954in}{2.644559in}}%
\pgfpathlineto{\pgfqpoint{2.007961in}{2.766988in}}%
\pgfpathlineto{\pgfqpoint{2.008078in}{2.614952in}}%
\pgfpathlineto{\pgfqpoint{2.008701in}{2.652745in}}%
\pgfpathlineto{\pgfqpoint{2.009542in}{2.771274in}}%
\pgfpathlineto{\pgfqpoint{2.009929in}{2.795500in}}%
\pgfpathlineto{\pgfqpoint{2.010381in}{2.644559in}}%
\pgfpathlineto{\pgfqpoint{2.012900in}{2.644559in}}%
\pgfpathlineto{\pgfqpoint{2.012906in}{2.766332in}}%
\pgfpathlineto{\pgfqpoint{2.013026in}{2.614741in}}%
\pgfpathlineto{\pgfqpoint{2.013650in}{2.652719in}}%
\pgfpathlineto{\pgfqpoint{2.014491in}{2.771487in}}%
\pgfpathlineto{\pgfqpoint{2.014874in}{2.795482in}}%
\pgfpathlineto{\pgfqpoint{2.015351in}{2.644559in}}%
\pgfpathlineto{\pgfqpoint{2.017845in}{2.644559in}}%
\pgfpathlineto{\pgfqpoint{2.017852in}{2.766258in}}%
\pgfpathlineto{\pgfqpoint{2.017972in}{2.617338in}}%
\pgfpathlineto{\pgfqpoint{2.018597in}{2.655386in}}%
\pgfpathlineto{\pgfqpoint{2.019438in}{2.769210in}}%
\pgfpathlineto{\pgfqpoint{2.019820in}{2.793121in}}%
\pgfpathlineto{\pgfqpoint{2.020290in}{2.644559in}}%
\pgfpathlineto{\pgfqpoint{2.022791in}{2.644559in}}%
\pgfpathlineto{\pgfqpoint{2.022797in}{2.766398in}}%
\pgfpathlineto{\pgfqpoint{2.022914in}{2.616145in}}%
\pgfpathlineto{\pgfqpoint{2.023536in}{2.653800in}}%
\pgfpathlineto{\pgfqpoint{2.024377in}{2.769744in}}%
\pgfpathlineto{\pgfqpoint{2.024765in}{2.794075in}}%
\pgfpathlineto{\pgfqpoint{2.025217in}{2.644559in}}%
\pgfpathlineto{\pgfqpoint{2.027736in}{2.644559in}}%
\pgfpathlineto{\pgfqpoint{2.027743in}{2.766315in}}%
\pgfpathlineto{\pgfqpoint{2.027862in}{2.614904in}}%
\pgfpathlineto{\pgfqpoint{2.028486in}{2.652843in}}%
\pgfpathlineto{\pgfqpoint{2.029327in}{2.771268in}}%
\pgfpathlineto{\pgfqpoint{2.029711in}{2.795303in}}%
\pgfpathlineto{\pgfqpoint{2.030171in}{2.644559in}}%
\pgfpathlineto{\pgfqpoint{2.032682in}{2.644559in}}%
\pgfpathlineto{\pgfqpoint{2.032688in}{2.765512in}}%
\pgfpathlineto{\pgfqpoint{2.032812in}{2.615589in}}%
\pgfpathlineto{\pgfqpoint{2.033440in}{2.653967in}}%
\pgfpathlineto{\pgfqpoint{2.034281in}{2.771231in}}%
\pgfpathlineto{\pgfqpoint{2.034656in}{2.794743in}}%
\pgfpathlineto{\pgfqpoint{2.035125in}{2.644559in}}%
\pgfpathlineto{\pgfqpoint{2.037627in}{2.644559in}}%
\pgfpathlineto{\pgfqpoint{2.037634in}{2.765699in}}%
\pgfpathlineto{\pgfqpoint{2.037755in}{2.617294in}}%
\pgfpathlineto{\pgfqpoint{2.038380in}{2.655383in}}%
\pgfpathlineto{\pgfqpoint{2.039221in}{2.769151in}}%
\pgfpathlineto{\pgfqpoint{2.039602in}{2.792985in}}%
\pgfpathlineto{\pgfqpoint{2.040076in}{2.644559in}}%
\pgfpathlineto{\pgfqpoint{2.042573in}{2.644559in}}%
\pgfpathlineto{\pgfqpoint{2.042579in}{2.766404in}}%
\pgfpathlineto{\pgfqpoint{2.042694in}{2.617676in}}%
\pgfpathlineto{\pgfqpoint{2.043364in}{2.718430in}}%
\pgfpathlineto{\pgfqpoint{2.043413in}{2.661328in}}%
\pgfpathlineto{\pgfqpoint{2.044056in}{2.761798in}}%
\pgfpathlineto{\pgfqpoint{2.044105in}{2.704950in}}%
\pgfpathlineto{\pgfqpoint{2.044547in}{2.792560in}}%
\pgfpathlineto{\pgfqpoint{2.044635in}{2.644559in}}%
\pgfpathlineto{\pgfqpoint{2.044826in}{2.644559in}}%
\pgfpathlineto{\pgfqpoint{2.047518in}{2.644559in}}%
\pgfpathlineto{\pgfqpoint{2.047525in}{2.764434in}}%
\pgfpathlineto{\pgfqpoint{2.047656in}{2.614374in}}%
\pgfpathlineto{\pgfqpoint{2.048289in}{2.653394in}}%
\pgfpathlineto{\pgfqpoint{2.049130in}{2.773187in}}%
\pgfpathlineto{\pgfqpoint{2.049493in}{2.795863in}}%
\pgfpathlineto{\pgfqpoint{2.049993in}{2.644559in}}%
\pgfpathlineto{\pgfqpoint{2.052464in}{2.644559in}}%
\pgfpathlineto{\pgfqpoint{2.052470in}{2.765708in}}%
\pgfpathlineto{\pgfqpoint{2.052592in}{2.616914in}}%
\pgfpathlineto{\pgfqpoint{2.053218in}{2.655076in}}%
\pgfpathlineto{\pgfqpoint{2.054059in}{2.769591in}}%
\pgfpathlineto{\pgfqpoint{2.054438in}{2.793347in}}%
\pgfpathlineto{\pgfqpoint{2.054935in}{2.644559in}}%
\pgfpathlineto{\pgfqpoint{2.057409in}{2.644559in}}%
\pgfpathlineto{\pgfqpoint{2.057416in}{2.764381in}}%
\pgfpathlineto{\pgfqpoint{2.057547in}{2.614385in}}%
\pgfpathlineto{\pgfqpoint{2.058180in}{2.653407in}}%
\pgfpathlineto{\pgfqpoint{2.059021in}{2.773154in}}%
\pgfpathlineto{\pgfqpoint{2.059384in}{2.795824in}}%
\pgfpathlineto{\pgfqpoint{2.059888in}{2.644559in}}%
\pgfpathlineto{\pgfqpoint{2.062355in}{2.644559in}}%
\pgfpathlineto{\pgfqpoint{2.062361in}{2.765398in}}%
\pgfpathlineto{\pgfqpoint{2.062484in}{2.616508in}}%
\pgfpathlineto{\pgfqpoint{2.063111in}{2.654770in}}%
\pgfpathlineto{\pgfqpoint{2.063951in}{2.770064in}}%
\pgfpathlineto{\pgfqpoint{2.064329in}{2.793694in}}%
\pgfpathlineto{\pgfqpoint{2.064790in}{2.644559in}}%
\pgfpathlineto{\pgfqpoint{2.067300in}{2.644559in}}%
\pgfpathlineto{\pgfqpoint{2.067307in}{2.765266in}}%
\pgfpathlineto{\pgfqpoint{2.067412in}{2.637448in}}%
\pgfpathlineto{\pgfqpoint{2.068069in}{2.697142in}}%
\pgfpathlineto{\pgfqpoint{2.068119in}{2.679622in}}%
\pgfpathlineto{\pgfqpoint{2.068663in}{2.734379in}}%
\pgfpathlineto{\pgfqpoint{2.068712in}{2.716950in}}%
\pgfpathlineto{\pgfqpoint{2.069241in}{2.770634in}}%
\pgfpathlineto{\pgfqpoint{2.069342in}{2.677544in}}%
\pgfpathlineto{\pgfqpoint{2.070111in}{2.644559in}}%
\pgfpathlineto{\pgfqpoint{2.072245in}{2.644559in}}%
\pgfpathlineto{\pgfqpoint{2.072252in}{2.765878in}}%
\pgfpathlineto{\pgfqpoint{2.072372in}{2.615379in}}%
\pgfpathlineto{\pgfqpoint{2.072996in}{2.653349in}}%
\pgfpathlineto{\pgfqpoint{2.073837in}{2.770857in}}%
\pgfpathlineto{\pgfqpoint{2.074220in}{2.794834in}}%
\pgfpathlineto{\pgfqpoint{2.074706in}{2.644559in}}%
\pgfpathlineto{\pgfqpoint{2.077191in}{2.644559in}}%
\pgfpathlineto{\pgfqpoint{2.077198in}{2.764350in}}%
\pgfpathlineto{\pgfqpoint{2.077334in}{2.643241in}}%
\pgfpathlineto{\pgfqpoint{2.077922in}{2.689567in}}%
\pgfpathlineto{\pgfqpoint{2.077972in}{2.682428in}}%
\pgfpathlineto{\pgfqpoint{2.078417in}{2.720613in}}%
\pgfpathlineto{\pgfqpoint{2.078466in}{2.713518in}}%
\pgfpathlineto{\pgfqpoint{2.079165in}{2.767588in}}%
\pgfpathlineto{\pgfqpoint{2.079201in}{2.705659in}}%
\pgfpathlineto{\pgfqpoint{2.079209in}{2.734520in}}%
\pgfpathlineto{\pgfqpoint{2.079961in}{2.644559in}}%
\pgfpathlineto{\pgfqpoint{2.082136in}{2.644559in}}%
\pgfpathlineto{\pgfqpoint{2.082143in}{2.764805in}}%
\pgfpathlineto{\pgfqpoint{2.082239in}{2.640532in}}%
\pgfpathlineto{\pgfqpoint{2.082888in}{2.693213in}}%
\pgfpathlineto{\pgfqpoint{2.082937in}{2.681394in}}%
\pgfpathlineto{\pgfqpoint{2.083481in}{2.730460in}}%
\pgfpathlineto{\pgfqpoint{2.083531in}{2.718711in}}%
\pgfpathlineto{\pgfqpoint{2.084068in}{2.767285in}}%
\pgfpathlineto{\pgfqpoint{2.084151in}{2.712629in}}%
\pgfpathlineto{\pgfqpoint{2.084161in}{2.720890in}}%
\pgfpathlineto{\pgfqpoint{2.084933in}{2.644559in}}%
\pgfpathlineto{\pgfqpoint{2.087082in}{2.644559in}}%
\pgfpathlineto{\pgfqpoint{2.087089in}{2.763835in}}%
\pgfpathlineto{\pgfqpoint{2.087836in}{2.685647in}}%
\pgfpathlineto{\pgfqpoint{2.089056in}{2.762601in}}%
\pgfpathlineto{\pgfqpoint{2.089081in}{2.747719in}}%
\pgfpathlineto{\pgfqpoint{2.089857in}{2.644559in}}%
\pgfpathlineto{\pgfqpoint{2.092027in}{2.644559in}}%
\pgfpathlineto{\pgfqpoint{2.092034in}{2.765328in}}%
\pgfpathlineto{\pgfqpoint{2.092156in}{2.617382in}}%
\pgfpathlineto{\pgfqpoint{2.092782in}{2.655527in}}%
\pgfpathlineto{\pgfqpoint{2.093622in}{2.769010in}}%
\pgfpathlineto{\pgfqpoint{2.094002in}{2.792761in}}%
\pgfpathlineto{\pgfqpoint{2.094463in}{2.644559in}}%
\pgfpathlineto{\pgfqpoint{2.096973in}{2.644559in}}%
\pgfpathlineto{\pgfqpoint{2.096979in}{2.764798in}}%
\pgfpathlineto{\pgfqpoint{2.097086in}{2.636463in}}%
\pgfpathlineto{\pgfqpoint{2.097745in}{2.698203in}}%
\pgfpathlineto{\pgfqpoint{2.097794in}{2.678782in}}%
\pgfpathlineto{\pgfqpoint{2.098338in}{2.735435in}}%
\pgfpathlineto{\pgfqpoint{2.098388in}{2.716114in}}%
\pgfpathlineto{\pgfqpoint{2.098915in}{2.771590in}}%
\pgfpathlineto{\pgfqpoint{2.099010in}{2.668944in}}%
\pgfpathlineto{\pgfqpoint{2.099790in}{2.644559in}}%
\pgfpathlineto{\pgfqpoint{2.101918in}{2.644559in}}%
\pgfpathlineto{\pgfqpoint{2.101925in}{2.763745in}}%
\pgfpathlineto{\pgfqpoint{2.102678in}{2.686571in}}%
\pgfpathlineto{\pgfqpoint{2.103893in}{2.761987in}}%
\pgfpathlineto{\pgfqpoint{2.103936in}{2.728967in}}%
\pgfpathlineto{\pgfqpoint{2.104726in}{2.644559in}}%
\pgfpathlineto{\pgfqpoint{2.106864in}{2.644559in}}%
\pgfpathlineto{\pgfqpoint{2.106870in}{2.764857in}}%
\pgfpathlineto{\pgfqpoint{2.106978in}{2.638597in}}%
\pgfpathlineto{\pgfqpoint{2.107635in}{2.696134in}}%
\pgfpathlineto{\pgfqpoint{2.107685in}{2.680977in}}%
\pgfpathlineto{\pgfqpoint{2.108229in}{2.733374in}}%
\pgfpathlineto{\pgfqpoint{2.108278in}{2.718300in}}%
\pgfpathlineto{\pgfqpoint{2.108805in}{2.769558in}}%
\pgfpathlineto{\pgfqpoint{2.108905in}{2.675670in}}%
\pgfpathlineto{\pgfqpoint{2.109647in}{2.644559in}}%
\pgfpathlineto{\pgfqpoint{2.111809in}{2.644559in}}%
\pgfpathlineto{\pgfqpoint{2.111816in}{2.764396in}}%
\pgfpathlineto{\pgfqpoint{2.111908in}{2.641526in}}%
\pgfpathlineto{\pgfqpoint{2.112553in}{2.691743in}}%
\pgfpathlineto{\pgfqpoint{2.112603in}{2.681719in}}%
\pgfpathlineto{\pgfqpoint{2.113048in}{2.722785in}}%
\pgfpathlineto{\pgfqpoint{2.113097in}{2.712814in}}%
\pgfpathlineto{\pgfqpoint{2.113737in}{2.766048in}}%
\pgfpathlineto{\pgfqpoint{2.113823in}{2.709720in}}%
\pgfpathlineto{\pgfqpoint{2.113833in}{2.727547in}}%
\pgfpathlineto{\pgfqpoint{2.114618in}{2.644559in}}%
\pgfpathlineto{\pgfqpoint{2.116755in}{2.644559in}}%
\pgfpathlineto{\pgfqpoint{2.116761in}{2.764815in}}%
\pgfpathlineto{\pgfqpoint{2.116871in}{2.635476in}}%
\pgfpathlineto{\pgfqpoint{2.117533in}{2.699572in}}%
\pgfpathlineto{\pgfqpoint{2.117583in}{2.678203in}}%
\pgfpathlineto{\pgfqpoint{2.118127in}{2.736800in}}%
\pgfpathlineto{\pgfqpoint{2.118176in}{2.715538in}}%
\pgfpathlineto{\pgfqpoint{2.118700in}{2.772749in}}%
\pgfpathlineto{\pgfqpoint{2.118806in}{2.644560in}}%
\pgfpathlineto{\pgfqpoint{2.121700in}{2.644559in}}%
\pgfpathlineto{\pgfqpoint{2.121707in}{2.763490in}}%
\pgfpathlineto{\pgfqpoint{2.122457in}{2.686445in}}%
\pgfpathlineto{\pgfqpoint{2.123674in}{2.761829in}}%
\pgfpathlineto{\pgfqpoint{2.123718in}{2.730521in}}%
\pgfpathlineto{\pgfqpoint{2.124474in}{2.644559in}}%
\pgfpathlineto{\pgfqpoint{2.126645in}{2.644559in}}%
\pgfpathlineto{\pgfqpoint{2.126652in}{2.764333in}}%
\pgfpathlineto{\pgfqpoint{2.126790in}{2.642947in}}%
\pgfpathlineto{\pgfqpoint{2.127379in}{2.689981in}}%
\pgfpathlineto{\pgfqpoint{2.127428in}{2.682221in}}%
\pgfpathlineto{\pgfqpoint{2.127874in}{2.721026in}}%
\pgfpathlineto{\pgfqpoint{2.127923in}{2.713312in}}%
\pgfpathlineto{\pgfqpoint{2.128620in}{2.767867in}}%
\pgfpathlineto{\pgfqpoint{2.128655in}{2.705645in}}%
\pgfpathlineto{\pgfqpoint{2.128664in}{2.734481in}}%
\pgfpathlineto{\pgfqpoint{2.129441in}{2.644559in}}%
\pgfpathlineto{\pgfqpoint{2.131591in}{2.644559in}}%
\pgfpathlineto{\pgfqpoint{2.131598in}{2.763915in}}%
\pgfpathlineto{\pgfqpoint{2.132345in}{2.685436in}}%
\pgfpathlineto{\pgfqpoint{2.133565in}{2.762653in}}%
\pgfpathlineto{\pgfqpoint{2.133589in}{2.747459in}}%
\pgfpathlineto{\pgfqpoint{2.134348in}{2.644559in}}%
\pgfpathlineto{\pgfqpoint{2.136536in}{2.644559in}}%
\pgfpathlineto{\pgfqpoint{2.136543in}{2.764640in}}%
\pgfpathlineto{\pgfqpoint{2.136653in}{2.637993in}}%
\pgfpathlineto{\pgfqpoint{2.137312in}{2.696848in}}%
\pgfpathlineto{\pgfqpoint{2.137361in}{2.680613in}}%
\pgfpathlineto{\pgfqpoint{2.137905in}{2.734087in}}%
\pgfpathlineto{\pgfqpoint{2.137955in}{2.717939in}}%
\pgfpathlineto{\pgfqpoint{2.138480in}{2.770143in}}%
\pgfpathlineto{\pgfqpoint{2.138578in}{2.676498in}}%
\pgfpathlineto{\pgfqpoint{2.139335in}{2.644559in}}%
\pgfpathlineto{\pgfqpoint{2.141482in}{2.644559in}}%
\pgfpathlineto{\pgfqpoint{2.141489in}{2.764075in}}%
\pgfpathlineto{\pgfqpoint{2.141579in}{2.641901in}}%
\pgfpathlineto{\pgfqpoint{2.142224in}{2.691265in}}%
\pgfpathlineto{\pgfqpoint{2.142273in}{2.681845in}}%
\pgfpathlineto{\pgfqpoint{2.142718in}{2.722308in}}%
\pgfpathlineto{\pgfqpoint{2.142768in}{2.712939in}}%
\pgfpathlineto{\pgfqpoint{2.143411in}{2.765766in}}%
\pgfpathlineto{\pgfqpoint{2.143492in}{2.704797in}}%
\pgfpathlineto{\pgfqpoint{2.143501in}{2.734967in}}%
\pgfpathlineto{\pgfqpoint{2.144267in}{2.644559in}}%
\pgfpathlineto{\pgfqpoint{2.146427in}{2.644559in}}%
\pgfpathlineto{\pgfqpoint{2.146434in}{2.764317in}}%
\pgfpathlineto{\pgfqpoint{2.146534in}{2.639630in}}%
\pgfpathlineto{\pgfqpoint{2.147183in}{2.694002in}}%
\pgfpathlineto{\pgfqpoint{2.147233in}{2.680849in}}%
\pgfpathlineto{\pgfqpoint{2.147777in}{2.731246in}}%
\pgfpathlineto{\pgfqpoint{2.147826in}{2.718168in}}%
\pgfpathlineto{\pgfqpoint{2.148361in}{2.767919in}}%
\pgfpathlineto{\pgfqpoint{2.148469in}{2.673827in}}%
\pgfpathlineto{\pgfqpoint{2.149229in}{2.644559in}}%
\pgfpathlineto{\pgfqpoint{2.151373in}{2.644559in}}%
\pgfpathlineto{\pgfqpoint{2.151379in}{2.764183in}}%
\pgfpathlineto{\pgfqpoint{2.151474in}{2.640764in}}%
\pgfpathlineto{\pgfqpoint{2.152121in}{2.692635in}}%
\pgfpathlineto{\pgfqpoint{2.152171in}{2.681351in}}%
\pgfpathlineto{\pgfqpoint{2.152715in}{2.729883in}}%
\pgfpathlineto{\pgfqpoint{2.152764in}{2.718668in}}%
\pgfpathlineto{\pgfqpoint{2.153303in}{2.766794in}}%
\pgfpathlineto{\pgfqpoint{2.153387in}{2.711939in}}%
\pgfpathlineto{\pgfqpoint{2.153397in}{2.722867in}}%
\pgfpathlineto{\pgfqpoint{2.154155in}{2.644559in}}%
\pgfpathlineto{\pgfqpoint{2.156318in}{2.644559in}}%
\pgfpathlineto{\pgfqpoint{2.156325in}{2.762825in}}%
\pgfpathlineto{\pgfqpoint{2.157076in}{2.686323in}}%
\pgfpathlineto{\pgfqpoint{2.158293in}{2.761766in}}%
\pgfpathlineto{\pgfqpoint{2.158336in}{2.729523in}}%
\pgfpathlineto{\pgfqpoint{2.159117in}{2.644559in}}%
\pgfpathlineto{\pgfqpoint{2.161264in}{2.644559in}}%
\pgfpathlineto{\pgfqpoint{2.161270in}{2.764179in}}%
\pgfpathlineto{\pgfqpoint{2.161361in}{2.641805in}}%
\pgfpathlineto{\pgfqpoint{2.162006in}{2.691437in}}%
\pgfpathlineto{\pgfqpoint{2.162056in}{2.681830in}}%
\pgfpathlineto{\pgfqpoint{2.162501in}{2.722480in}}%
\pgfpathlineto{\pgfqpoint{2.162550in}{2.712925in}}%
\pgfpathlineto{\pgfqpoint{2.163193in}{2.765937in}}%
\pgfpathlineto{\pgfqpoint{2.163277in}{2.708833in}}%
\pgfpathlineto{\pgfqpoint{2.163287in}{2.728913in}}%
\pgfpathlineto{\pgfqpoint{2.164045in}{2.644559in}}%
\pgfpathlineto{\pgfqpoint{2.166209in}{2.644559in}}%
\pgfpathlineto{\pgfqpoint{2.166216in}{2.763483in}}%
\pgfpathlineto{\pgfqpoint{2.166349in}{2.643564in}}%
\pgfpathlineto{\pgfqpoint{2.166984in}{2.682413in}}%
\pgfpathlineto{\pgfqpoint{2.168023in}{2.756533in}}%
\pgfpathlineto{\pgfqpoint{2.168184in}{2.766636in}}%
\pgfpathlineto{\pgfqpoint{2.169752in}{2.644559in}}%
\pgfpathlineto{\pgfqpoint{2.171155in}{2.644559in}}%
\pgfpathlineto{\pgfqpoint{2.171161in}{2.764417in}}%
\pgfpathlineto{\pgfqpoint{2.171263in}{2.639142in}}%
\pgfpathlineto{\pgfqpoint{2.171915in}{2.694695in}}%
\pgfpathlineto{\pgfqpoint{2.171964in}{2.680668in}}%
\pgfpathlineto{\pgfqpoint{2.172508in}{2.731938in}}%
\pgfpathlineto{\pgfqpoint{2.172558in}{2.717989in}}%
\pgfpathlineto{\pgfqpoint{2.173091in}{2.768486in}}%
\pgfpathlineto{\pgfqpoint{2.173196in}{2.674603in}}%
\pgfpathlineto{\pgfqpoint{2.173978in}{2.644559in}}%
\pgfpathlineto{\pgfqpoint{2.176100in}{2.644559in}}%
\pgfpathlineto{\pgfqpoint{2.176107in}{2.763769in}}%
\pgfpathlineto{\pgfqpoint{2.176244in}{2.642913in}}%
\pgfpathlineto{\pgfqpoint{2.176832in}{2.689655in}}%
\pgfpathlineto{\pgfqpoint{2.176882in}{2.682099in}}%
\pgfpathlineto{\pgfqpoint{2.177327in}{2.720701in}}%
\pgfpathlineto{\pgfqpoint{2.177376in}{2.713191in}}%
\pgfpathlineto{\pgfqpoint{2.178074in}{2.767616in}}%
\pgfpathlineto{\pgfqpoint{2.178110in}{2.705501in}}%
\pgfpathlineto{\pgfqpoint{2.178119in}{2.734327in}}%
\pgfpathlineto{\pgfqpoint{2.178887in}{2.644559in}}%
\pgfpathlineto{\pgfqpoint{2.181046in}{2.644559in}}%
\pgfpathlineto{\pgfqpoint{2.181052in}{2.763397in}}%
\pgfpathlineto{\pgfqpoint{2.181801in}{2.686367in}}%
\pgfpathlineto{\pgfqpoint{2.183020in}{2.761743in}}%
\pgfpathlineto{\pgfqpoint{2.183063in}{2.731117in}}%
\pgfpathlineto{\pgfqpoint{2.183809in}{2.644559in}}%
\pgfpathlineto{\pgfqpoint{2.185991in}{2.644559in}}%
\pgfpathlineto{\pgfqpoint{2.185998in}{2.763511in}}%
\pgfpathlineto{\pgfqpoint{2.186746in}{2.685619in}}%
\pgfpathlineto{\pgfqpoint{2.187965in}{2.762465in}}%
\pgfpathlineto{\pgfqpoint{2.187990in}{2.747640in}}%
\pgfpathlineto{\pgfqpoint{2.188764in}{2.644559in}}%
\pgfpathlineto{\pgfqpoint{2.190936in}{2.644559in}}%
\pgfpathlineto{\pgfqpoint{2.190943in}{2.763546in}}%
\pgfpathlineto{\pgfqpoint{2.191691in}{2.685445in}}%
\pgfpathlineto{\pgfqpoint{2.192911in}{2.762558in}}%
\pgfpathlineto{\pgfqpoint{2.192935in}{2.747397in}}%
\pgfpathlineto{\pgfqpoint{2.193721in}{2.644559in}}%
\pgfpathlineto{\pgfqpoint{2.195882in}{2.644559in}}%
\pgfpathlineto{\pgfqpoint{2.195889in}{2.763409in}}%
\pgfpathlineto{\pgfqpoint{2.196642in}{2.686291in}}%
\pgfpathlineto{\pgfqpoint{2.197856in}{2.762039in}}%
\pgfpathlineto{\pgfqpoint{2.197899in}{2.725944in}}%
\pgfpathlineto{\pgfqpoint{2.198655in}{2.644559in}}%
\pgfpathlineto{\pgfqpoint{2.200827in}{2.644559in}}%
\pgfpathlineto{\pgfqpoint{2.200834in}{2.763636in}}%
\pgfpathlineto{\pgfqpoint{2.200971in}{2.642984in}}%
\pgfpathlineto{\pgfqpoint{2.201559in}{2.689466in}}%
\pgfpathlineto{\pgfqpoint{2.201608in}{2.682125in}}%
\pgfpathlineto{\pgfqpoint{2.202053in}{2.720512in}}%
\pgfpathlineto{\pgfqpoint{2.202103in}{2.713216in}}%
\pgfpathlineto{\pgfqpoint{2.202802in}{2.767480in}}%
\pgfpathlineto{\pgfqpoint{2.202837in}{2.705470in}}%
\pgfpathlineto{\pgfqpoint{2.202846in}{2.734300in}}%
\pgfpathlineto{\pgfqpoint{2.203630in}{2.644559in}}%
\pgfpathlineto{\pgfqpoint{2.205773in}{2.644559in}}%
\pgfpathlineto{\pgfqpoint{2.205780in}{2.763708in}}%
\pgfpathlineto{\pgfqpoint{2.205874in}{2.640708in}}%
\pgfpathlineto{\pgfqpoint{2.206521in}{2.692412in}}%
\pgfpathlineto{\pgfqpoint{2.206570in}{2.681191in}}%
\pgfpathlineto{\pgfqpoint{2.207114in}{2.729660in}}%
\pgfpathlineto{\pgfqpoint{2.207164in}{2.718507in}}%
\pgfpathlineto{\pgfqpoint{2.207703in}{2.766595in}}%
\pgfpathlineto{\pgfqpoint{2.207787in}{2.711795in}}%
\pgfpathlineto{\pgfqpoint{2.207797in}{2.722739in}}%
\pgfpathlineto{\pgfqpoint{2.208571in}{2.644559in}}%
\pgfpathlineto{\pgfqpoint{2.210718in}{2.644559in}}%
\pgfpathlineto{\pgfqpoint{2.210725in}{2.762264in}}%
\pgfpathlineto{\pgfqpoint{2.210847in}{2.618847in}}%
\pgfpathlineto{\pgfqpoint{2.211469in}{2.656708in}}%
\pgfpathlineto{\pgfqpoint{2.212310in}{2.766603in}}%
\pgfpathlineto{\pgfqpoint{2.212693in}{2.790559in}}%
\pgfpathlineto{\pgfqpoint{2.213192in}{2.644559in}}%
\pgfpathlineto{\pgfqpoint{2.215664in}{2.644559in}}%
\pgfpathlineto{\pgfqpoint{2.215670in}{2.764004in}}%
\pgfpathlineto{\pgfqpoint{2.215776in}{2.637358in}}%
\pgfpathlineto{\pgfqpoint{2.216431in}{2.696756in}}%
\pgfpathlineto{\pgfqpoint{2.216480in}{2.679322in}}%
\pgfpathlineto{\pgfqpoint{2.217024in}{2.733992in}}%
\pgfpathlineto{\pgfqpoint{2.217074in}{2.716649in}}%
\pgfpathlineto{\pgfqpoint{2.217603in}{2.770307in}}%
\pgfpathlineto{\pgfqpoint{2.217705in}{2.677377in}}%
\pgfpathlineto{\pgfqpoint{2.218477in}{2.644559in}}%
\pgfpathlineto{\pgfqpoint{2.220609in}{2.644559in}}%
\pgfpathlineto{\pgfqpoint{2.220616in}{2.763417in}}%
\pgfpathlineto{\pgfqpoint{2.221364in}{2.685408in}}%
\pgfpathlineto{\pgfqpoint{2.222584in}{2.762552in}}%
\pgfpathlineto{\pgfqpoint{2.222607in}{2.747309in}}%
\pgfpathlineto{\pgfqpoint{2.223371in}{2.644559in}}%
\pgfpathlineto{\pgfqpoint{2.225555in}{2.644559in}}%
\pgfpathlineto{\pgfqpoint{2.225561in}{2.762823in}}%
\pgfpathlineto{\pgfqpoint{2.225670in}{2.618337in}}%
\pgfpathlineto{\pgfqpoint{2.226326in}{2.715746in}}%
\pgfpathlineto{\pgfqpoint{2.226375in}{2.660756in}}%
\pgfpathlineto{\pgfqpoint{2.227018in}{2.759129in}}%
\pgfpathlineto{\pgfqpoint{2.227068in}{2.704365in}}%
\pgfpathlineto{\pgfqpoint{2.227496in}{2.789070in}}%
\pgfpathlineto{\pgfqpoint{2.227611in}{2.644559in}}%
\pgfpathlineto{\pgfqpoint{2.227775in}{2.644559in}}%
\pgfpathlineto{\pgfqpoint{2.230500in}{2.644559in}}%
\pgfpathlineto{\pgfqpoint{2.230507in}{2.762780in}}%
\pgfpathlineto{\pgfqpoint{2.230616in}{2.618517in}}%
\pgfpathlineto{\pgfqpoint{2.231272in}{2.715616in}}%
\pgfpathlineto{\pgfqpoint{2.231322in}{2.661010in}}%
\pgfpathlineto{\pgfqpoint{2.231965in}{2.758998in}}%
\pgfpathlineto{\pgfqpoint{2.232014in}{2.704618in}}%
\pgfpathlineto{\pgfqpoint{2.232442in}{2.788905in}}%
\pgfpathlineto{\pgfqpoint{2.232557in}{2.644559in}}%
\pgfpathlineto{\pgfqpoint{2.232716in}{2.644559in}}%
\pgfpathlineto{\pgfqpoint{2.235446in}{2.644559in}}%
\pgfpathlineto{\pgfqpoint{2.235452in}{2.762841in}}%
\pgfpathlineto{\pgfqpoint{2.235562in}{2.618798in}}%
\pgfpathlineto{\pgfqpoint{2.236219in}{2.715394in}}%
\pgfpathlineto{\pgfqpoint{2.236269in}{2.661387in}}%
\pgfpathlineto{\pgfqpoint{2.236912in}{2.758777in}}%
\pgfpathlineto{\pgfqpoint{2.236961in}{2.704994in}}%
\pgfpathlineto{\pgfqpoint{2.237388in}{2.788638in}}%
\pgfpathlineto{\pgfqpoint{2.237502in}{2.644559in}}%
\pgfpathlineto{\pgfqpoint{2.237662in}{2.644559in}}%
\pgfpathlineto{\pgfqpoint{2.240391in}{2.644559in}}%
\pgfpathlineto{\pgfqpoint{2.240397in}{2.762802in}}%
\pgfpathlineto{\pgfqpoint{2.240507in}{2.618589in}}%
\pgfpathlineto{\pgfqpoint{2.241164in}{2.715526in}}%
\pgfpathlineto{\pgfqpoint{2.241213in}{2.661109in}}%
\pgfpathlineto{\pgfqpoint{2.241856in}{2.758909in}}%
\pgfpathlineto{\pgfqpoint{2.241905in}{2.704716in}}%
\pgfpathlineto{\pgfqpoint{2.242333in}{2.788802in}}%
\pgfpathlineto{\pgfqpoint{2.242448in}{2.644559in}}%
\pgfpathlineto{\pgfqpoint{2.242606in}{2.644559in}}%
\pgfpathlineto{\pgfqpoint{2.245336in}{2.644559in}}%
\pgfpathlineto{\pgfqpoint{2.245343in}{2.762779in}}%
\pgfpathlineto{\pgfqpoint{2.245453in}{2.618503in}}%
\pgfpathlineto{\pgfqpoint{2.246109in}{2.715579in}}%
\pgfpathlineto{\pgfqpoint{2.246158in}{2.660992in}}%
\pgfpathlineto{\pgfqpoint{2.246801in}{2.758962in}}%
\pgfpathlineto{\pgfqpoint{2.246850in}{2.704600in}}%
\pgfpathlineto{\pgfqpoint{2.247278in}{2.788869in}}%
\pgfpathlineto{\pgfqpoint{2.247393in}{2.644559in}}%
\pgfpathlineto{\pgfqpoint{2.247552in}{2.644559in}}%
\pgfpathlineto{\pgfqpoint{2.250282in}{2.644559in}}%
\pgfpathlineto{\pgfqpoint{2.250288in}{2.762784in}}%
\pgfpathlineto{\pgfqpoint{2.250393in}{2.617922in}}%
\pgfpathlineto{\pgfqpoint{2.251045in}{2.715645in}}%
\pgfpathlineto{\pgfqpoint{2.251094in}{2.659743in}}%
\pgfpathlineto{\pgfqpoint{2.251737in}{2.759026in}}%
\pgfpathlineto{\pgfqpoint{2.251786in}{2.703352in}}%
\pgfpathlineto{\pgfqpoint{2.252219in}{2.789231in}}%
\pgfpathlineto{\pgfqpoint{2.252338in}{2.644559in}}%
\pgfpathlineto{\pgfqpoint{2.252497in}{2.644559in}}%
\pgfpathlineto{\pgfqpoint{2.255227in}{2.644559in}}%
\pgfpathlineto{\pgfqpoint{2.255234in}{2.762712in}}%
\pgfpathlineto{\pgfqpoint{2.255338in}{2.617886in}}%
\pgfpathlineto{\pgfqpoint{2.255990in}{2.715645in}}%
\pgfpathlineto{\pgfqpoint{2.256039in}{2.659684in}}%
\pgfpathlineto{\pgfqpoint{2.256682in}{2.759027in}}%
\pgfpathlineto{\pgfqpoint{2.256732in}{2.703292in}}%
\pgfpathlineto{\pgfqpoint{2.257164in}{2.789240in}}%
\pgfpathlineto{\pgfqpoint{2.257284in}{2.644559in}}%
\pgfpathlineto{\pgfqpoint{2.257442in}{2.644559in}}%
\pgfpathlineto{\pgfqpoint{2.260173in}{2.644559in}}%
\pgfpathlineto{\pgfqpoint{2.260179in}{2.762683in}}%
\pgfpathlineto{\pgfqpoint{2.260283in}{2.617873in}}%
\pgfpathlineto{\pgfqpoint{2.260935in}{2.715625in}}%
\pgfpathlineto{\pgfqpoint{2.260984in}{2.659637in}}%
\pgfpathlineto{\pgfqpoint{2.261627in}{2.759007in}}%
\pgfpathlineto{\pgfqpoint{2.261677in}{2.703246in}}%
\pgfpathlineto{\pgfqpoint{2.262110in}{2.789234in}}%
\pgfpathlineto{\pgfqpoint{2.262229in}{2.644559in}}%
\pgfpathlineto{\pgfqpoint{2.262388in}{2.644559in}}%
\pgfpathlineto{\pgfqpoint{2.265118in}{2.644559in}}%
\pgfpathlineto{\pgfqpoint{2.265125in}{2.762641in}}%
\pgfpathlineto{\pgfqpoint{2.265229in}{2.617875in}}%
\pgfpathlineto{\pgfqpoint{2.265880in}{2.715585in}}%
\pgfpathlineto{\pgfqpoint{2.265929in}{2.659603in}}%
\pgfpathlineto{\pgfqpoint{2.266572in}{2.758967in}}%
\pgfpathlineto{\pgfqpoint{2.266622in}{2.703212in}}%
\pgfpathlineto{\pgfqpoint{2.267055in}{2.789209in}}%
\pgfpathlineto{\pgfqpoint{2.267175in}{2.644559in}}%
\pgfpathlineto{\pgfqpoint{2.267333in}{2.644559in}}%
\pgfpathlineto{\pgfqpoint{2.270064in}{2.644559in}}%
\pgfpathlineto{\pgfqpoint{2.270070in}{2.762605in}}%
\pgfpathlineto{\pgfqpoint{2.270174in}{2.617877in}}%
\pgfpathlineto{\pgfqpoint{2.270825in}{2.715551in}}%
\pgfpathlineto{\pgfqpoint{2.270874in}{2.659576in}}%
\pgfpathlineto{\pgfqpoint{2.271517in}{2.758933in}}%
\pgfpathlineto{\pgfqpoint{2.271567in}{2.703185in}}%
\pgfpathlineto{\pgfqpoint{2.272000in}{2.789187in}}%
\pgfpathlineto{\pgfqpoint{2.272120in}{2.644559in}}%
\pgfpathlineto{\pgfqpoint{2.272279in}{2.644559in}}%
\pgfpathlineto{\pgfqpoint{2.275009in}{2.644559in}}%
\pgfpathlineto{\pgfqpoint{2.275016in}{2.762566in}}%
\pgfpathlineto{\pgfqpoint{2.275119in}{2.617878in}}%
\pgfpathlineto{\pgfqpoint{2.275770in}{2.715523in}}%
\pgfpathlineto{\pgfqpoint{2.275820in}{2.659553in}}%
\pgfpathlineto{\pgfqpoint{2.276462in}{2.758905in}}%
\pgfpathlineto{\pgfqpoint{2.276512in}{2.703162in}}%
\pgfpathlineto{\pgfqpoint{2.276946in}{2.789168in}}%
\pgfpathlineto{\pgfqpoint{2.277066in}{2.644559in}}%
\pgfpathlineto{\pgfqpoint{2.277224in}{2.644559in}}%
\pgfpathlineto{\pgfqpoint{2.279955in}{2.644559in}}%
\pgfpathlineto{\pgfqpoint{2.279961in}{2.762530in}}%
\pgfpathlineto{\pgfqpoint{2.280064in}{2.618051in}}%
\pgfpathlineto{\pgfqpoint{2.280715in}{2.715301in}}%
\pgfpathlineto{\pgfqpoint{2.280764in}{2.659670in}}%
\pgfpathlineto{\pgfqpoint{2.281407in}{2.758683in}}%
\pgfpathlineto{\pgfqpoint{2.281457in}{2.703278in}}%
\pgfpathlineto{\pgfqpoint{2.281891in}{2.788969in}}%
\pgfpathlineto{\pgfqpoint{2.282011in}{2.644559in}}%
\pgfpathlineto{\pgfqpoint{2.282170in}{2.644559in}}%
\pgfpathlineto{\pgfqpoint{2.284900in}{2.644559in}}%
\pgfpathlineto{\pgfqpoint{2.284907in}{2.762492in}}%
\pgfpathlineto{\pgfqpoint{2.285010in}{2.617858in}}%
\pgfpathlineto{\pgfqpoint{2.285661in}{2.715520in}}%
\pgfpathlineto{\pgfqpoint{2.285711in}{2.659541in}}%
\pgfpathlineto{\pgfqpoint{2.286353in}{2.758902in}}%
\pgfpathlineto{\pgfqpoint{2.286403in}{2.703150in}}%
\pgfpathlineto{\pgfqpoint{2.286837in}{2.789160in}}%
\pgfpathlineto{\pgfqpoint{2.286957in}{2.644559in}}%
\pgfpathlineto{\pgfqpoint{2.287115in}{2.644559in}}%
\pgfpathlineto{\pgfqpoint{2.289846in}{2.644559in}}%
\pgfpathlineto{\pgfqpoint{2.289852in}{2.762461in}}%
\pgfpathlineto{\pgfqpoint{2.289955in}{2.618061in}}%
\pgfpathlineto{\pgfqpoint{2.290605in}{2.715218in}}%
\pgfpathlineto{\pgfqpoint{2.290654in}{2.659600in}}%
\pgfpathlineto{\pgfqpoint{2.291297in}{2.758601in}}%
\pgfpathlineto{\pgfqpoint{2.291347in}{2.703208in}}%
\pgfpathlineto{\pgfqpoint{2.291781in}{2.788919in}}%
\pgfpathlineto{\pgfqpoint{2.291902in}{2.644559in}}%
\pgfpathlineto{\pgfqpoint{2.292061in}{2.644559in}}%
\pgfpathlineto{\pgfqpoint{2.294791in}{2.644559in}}%
\pgfpathlineto{\pgfqpoint{2.294798in}{2.762866in}}%
\pgfpathlineto{\pgfqpoint{2.295545in}{2.685239in}}%
\pgfpathlineto{\pgfqpoint{2.296765in}{2.762403in}}%
\pgfpathlineto{\pgfqpoint{2.296789in}{2.747178in}}%
\pgfpathlineto{\pgfqpoint{2.297584in}{2.644559in}}%
\pgfpathlineto{\pgfqpoint{2.299736in}{2.644559in}}%
\pgfpathlineto{\pgfqpoint{2.299743in}{2.762570in}}%
\pgfpathlineto{\pgfqpoint{2.299846in}{2.638827in}}%
\pgfpathlineto{\pgfqpoint{2.300496in}{2.694534in}}%
\pgfpathlineto{\pgfqpoint{2.300546in}{2.680267in}}%
\pgfpathlineto{\pgfqpoint{2.301090in}{2.731776in}}%
\pgfpathlineto{\pgfqpoint{2.301139in}{2.717589in}}%
\pgfpathlineto{\pgfqpoint{2.301672in}{2.768333in}}%
\pgfpathlineto{\pgfqpoint{2.301778in}{2.674611in}}%
\pgfpathlineto{\pgfqpoint{2.302517in}{2.644559in}}%
\pgfpathlineto{\pgfqpoint{2.304682in}{2.644559in}}%
\pgfpathlineto{\pgfqpoint{2.304689in}{2.762913in}}%
\pgfpathlineto{\pgfqpoint{2.304821in}{2.643787in}}%
\pgfpathlineto{\pgfqpoint{2.305454in}{2.682462in}}%
\pgfpathlineto{\pgfqpoint{2.306492in}{2.755997in}}%
\pgfpathlineto{\pgfqpoint{2.306712in}{2.667492in}}%
\pgfpathlineto{\pgfqpoint{2.306656in}{2.766284in}}%
\pgfpathlineto{\pgfqpoint{2.306725in}{2.696676in}}%
\pgfpathlineto{\pgfqpoint{2.307474in}{2.644559in}}%
\pgfpathlineto{\pgfqpoint{2.309627in}{2.644559in}}%
\pgfpathlineto{\pgfqpoint{2.309634in}{2.761940in}}%
\pgfpathlineto{\pgfqpoint{2.309740in}{2.630389in}}%
\pgfpathlineto{\pgfqpoint{2.310395in}{2.703306in}}%
\pgfpathlineto{\pgfqpoint{2.310444in}{2.672371in}}%
\pgfpathlineto{\pgfqpoint{2.311087in}{2.746723in}}%
\pgfpathlineto{\pgfqpoint{2.311137in}{2.715944in}}%
\pgfpathlineto{\pgfqpoint{2.311567in}{2.776804in}}%
\pgfpathlineto{\pgfqpoint{2.311692in}{2.644559in}}%
\pgfpathlineto{\pgfqpoint{2.311788in}{2.644559in}}%
\pgfpathlineto{\pgfqpoint{2.314573in}{2.644559in}}%
\pgfpathlineto{\pgfqpoint{2.314580in}{2.762771in}}%
\pgfpathlineto{\pgfqpoint{2.314684in}{2.638472in}}%
\pgfpathlineto{\pgfqpoint{2.315335in}{2.694888in}}%
\pgfpathlineto{\pgfqpoint{2.315384in}{2.680053in}}%
\pgfpathlineto{\pgfqpoint{2.315928in}{2.732129in}}%
\pgfpathlineto{\pgfqpoint{2.315978in}{2.717376in}}%
\pgfpathlineto{\pgfqpoint{2.316510in}{2.768613in}}%
\pgfpathlineto{\pgfqpoint{2.316614in}{2.675107in}}%
\pgfpathlineto{\pgfqpoint{2.317366in}{2.644559in}}%
\pgfpathlineto{\pgfqpoint{2.319518in}{2.644559in}}%
\pgfpathlineto{\pgfqpoint{2.319525in}{2.762668in}}%
\pgfpathlineto{\pgfqpoint{2.319620in}{2.640499in}}%
\pgfpathlineto{\pgfqpoint{2.320266in}{2.692389in}}%
\pgfpathlineto{\pgfqpoint{2.320316in}{2.680975in}}%
\pgfpathlineto{\pgfqpoint{2.320860in}{2.729636in}}%
\pgfpathlineto{\pgfqpoint{2.320909in}{2.718291in}}%
\pgfpathlineto{\pgfqpoint{2.321448in}{2.766566in}}%
\pgfpathlineto{\pgfqpoint{2.321532in}{2.711905in}}%
\pgfpathlineto{\pgfqpoint{2.321543in}{2.721737in}}%
\pgfpathlineto{\pgfqpoint{2.322316in}{2.644559in}}%
\pgfpathlineto{\pgfqpoint{2.324464in}{2.644559in}}%
\pgfpathlineto{\pgfqpoint{2.324471in}{2.762427in}}%
\pgfpathlineto{\pgfqpoint{2.324584in}{2.634583in}}%
\pgfpathlineto{\pgfqpoint{2.325199in}{2.671395in}}%
\pgfpathlineto{\pgfqpoint{2.326040in}{2.749759in}}%
\pgfpathlineto{\pgfqpoint{2.326412in}{2.773089in}}%
\pgfpathlineto{\pgfqpoint{2.327019in}{2.644559in}}%
\pgfpathlineto{\pgfqpoint{2.329409in}{2.644559in}}%
\pgfpathlineto{\pgfqpoint{2.329416in}{2.761579in}}%
\pgfpathlineto{\pgfqpoint{2.329528in}{2.621236in}}%
\pgfpathlineto{\pgfqpoint{2.330144in}{2.658122in}}%
\pgfpathlineto{\pgfqpoint{2.330985in}{2.762948in}}%
\pgfpathlineto{\pgfqpoint{2.331357in}{2.786268in}}%
\pgfpathlineto{\pgfqpoint{2.331871in}{2.644559in}}%
\pgfpathlineto{\pgfqpoint{2.334355in}{2.644559in}}%
\pgfpathlineto{\pgfqpoint{2.334361in}{2.762791in}}%
\pgfpathlineto{\pgfqpoint{2.334454in}{2.641360in}}%
\pgfpathlineto{\pgfqpoint{2.335098in}{2.691399in}}%
\pgfpathlineto{\pgfqpoint{2.335147in}{2.681357in}}%
\pgfpathlineto{\pgfqpoint{2.335592in}{2.722441in}}%
\pgfpathlineto{\pgfqpoint{2.335642in}{2.712453in}}%
\pgfpathlineto{\pgfqpoint{2.336282in}{2.765738in}}%
\pgfpathlineto{\pgfqpoint{2.336368in}{2.709686in}}%
\pgfpathlineto{\pgfqpoint{2.336378in}{2.726855in}}%
\pgfpathlineto{\pgfqpoint{2.337159in}{2.644559in}}%
\pgfpathlineto{\pgfqpoint{2.339300in}{2.644559in}}%
\pgfpathlineto{\pgfqpoint{2.339306in}{2.771479in}}%
\pgfpathlineto{\pgfqpoint{2.339429in}{2.631242in}}%
\pgfpathlineto{\pgfqpoint{2.340099in}{2.704254in}}%
\pgfpathlineto{\pgfqpoint{2.340149in}{2.675144in}}%
\pgfpathlineto{\pgfqpoint{2.340792in}{2.747646in}}%
\pgfpathlineto{\pgfqpoint{2.340841in}{2.718743in}}%
\pgfpathlineto{\pgfqpoint{2.341274in}{2.777890in}}%
\pgfpathlineto{\pgfqpoint{2.341358in}{2.644559in}}%
\pgfpathlineto{\pgfqpoint{2.341482in}{2.644559in}}%
\pgfpathlineto{\pgfqpoint{2.344246in}{2.644559in}}%
\pgfpathlineto{\pgfqpoint{2.344252in}{2.762618in}}%
\pgfpathlineto{\pgfqpoint{2.344352in}{2.618203in}}%
\pgfpathlineto{\pgfqpoint{2.345001in}{2.714921in}}%
\pgfpathlineto{\pgfqpoint{2.345050in}{2.659428in}}%
\pgfpathlineto{\pgfqpoint{2.345693in}{2.758305in}}%
\pgfpathlineto{\pgfqpoint{2.345743in}{2.703035in}}%
\pgfpathlineto{\pgfqpoint{2.346179in}{2.788749in}}%
\pgfpathlineto{\pgfqpoint{2.346302in}{2.644559in}}%
\pgfpathlineto{\pgfqpoint{2.346448in}{2.644559in}}%
\pgfpathlineto{\pgfqpoint{2.349191in}{2.644559in}}%
\pgfpathlineto{\pgfqpoint{2.349198in}{2.762243in}}%
\pgfpathlineto{\pgfqpoint{2.349306in}{2.629511in}}%
\pgfpathlineto{\pgfqpoint{2.349963in}{2.704418in}}%
\pgfpathlineto{\pgfqpoint{2.350013in}{2.671797in}}%
\pgfpathlineto{\pgfqpoint{2.350656in}{2.747831in}}%
\pgfpathlineto{\pgfqpoint{2.350705in}{2.715376in}}%
\pgfpathlineto{\pgfqpoint{2.351133in}{2.777752in}}%
\pgfpathlineto{\pgfqpoint{2.351255in}{2.644559in}}%
\pgfpathlineto{\pgfqpoint{2.351357in}{2.644559in}}%
\pgfpathlineto{\pgfqpoint{2.354136in}{2.644559in}}%
\pgfpathlineto{\pgfqpoint{2.354143in}{2.762506in}}%
\pgfpathlineto{\pgfqpoint{2.354255in}{2.634757in}}%
\pgfpathlineto{\pgfqpoint{2.354868in}{2.671282in}}%
\pgfpathlineto{\pgfqpoint{2.355708in}{2.749329in}}%
\pgfpathlineto{\pgfqpoint{2.356083in}{2.772804in}}%
\pgfpathlineto{\pgfqpoint{2.356681in}{2.644559in}}%
\pgfpathlineto{\pgfqpoint{2.359082in}{2.644559in}}%
\pgfpathlineto{\pgfqpoint{2.359089in}{2.761655in}}%
\pgfpathlineto{\pgfqpoint{2.359199in}{2.620415in}}%
\pgfpathlineto{\pgfqpoint{2.359859in}{2.713705in}}%
\pgfpathlineto{\pgfqpoint{2.359909in}{2.663057in}}%
\pgfpathlineto{\pgfqpoint{2.360551in}{2.757088in}}%
\pgfpathlineto{\pgfqpoint{2.360601in}{2.706664in}}%
\pgfpathlineto{\pgfqpoint{2.361026in}{2.786838in}}%
\pgfpathlineto{\pgfqpoint{2.361138in}{2.644559in}}%
\pgfpathlineto{\pgfqpoint{2.361300in}{2.644559in}}%
\pgfpathlineto{\pgfqpoint{2.364027in}{2.644559in}}%
\pgfpathlineto{\pgfqpoint{2.364034in}{2.762499in}}%
\pgfpathlineto{\pgfqpoint{2.364141in}{2.636305in}}%
\pgfpathlineto{\pgfqpoint{2.364797in}{2.697635in}}%
\pgfpathlineto{\pgfqpoint{2.364846in}{2.678370in}}%
\pgfpathlineto{\pgfqpoint{2.365390in}{2.734868in}}%
\pgfpathlineto{\pgfqpoint{2.365440in}{2.715701in}}%
\pgfpathlineto{\pgfqpoint{2.365968in}{2.771101in}}%
\pgfpathlineto{\pgfqpoint{2.366065in}{2.668648in}}%
\pgfpathlineto{\pgfqpoint{2.366847in}{2.644559in}}%
\pgfpathlineto{\pgfqpoint{2.368972in}{2.644559in}}%
\pgfpathlineto{\pgfqpoint{2.368978in}{2.773161in}}%
\pgfpathlineto{\pgfqpoint{2.369103in}{2.627145in}}%
\pgfpathlineto{\pgfqpoint{2.369775in}{2.708578in}}%
\pgfpathlineto{\pgfqpoint{2.369825in}{2.671197in}}%
\pgfpathlineto{\pgfqpoint{2.370468in}{2.751985in}}%
\pgfpathlineto{\pgfqpoint{2.370517in}{2.714781in}}%
\pgfpathlineto{\pgfqpoint{2.370947in}{2.782028in}}%
\pgfpathlineto{\pgfqpoint{2.371034in}{2.644559in}}%
\pgfpathlineto{\pgfqpoint{2.371203in}{2.644559in}}%
\pgfpathlineto{\pgfqpoint{2.373918in}{2.644559in}}%
\pgfpathlineto{\pgfqpoint{2.373925in}{2.762589in}}%
\pgfpathlineto{\pgfqpoint{2.374020in}{2.618707in}}%
\pgfpathlineto{\pgfqpoint{2.374666in}{2.714060in}}%
\pgfpathlineto{\pgfqpoint{2.374715in}{2.659271in}}%
\pgfpathlineto{\pgfqpoint{2.375358in}{2.757447in}}%
\pgfpathlineto{\pgfqpoint{2.375408in}{2.702875in}}%
\pgfpathlineto{\pgfqpoint{2.375848in}{2.788131in}}%
\pgfpathlineto{\pgfqpoint{2.375985in}{2.644559in}}%
\pgfpathlineto{\pgfqpoint{2.376106in}{2.644559in}}%
\pgfpathlineto{\pgfqpoint{2.378863in}{2.644559in}}%
\pgfpathlineto{\pgfqpoint{2.378869in}{2.772091in}}%
\pgfpathlineto{\pgfqpoint{2.378986in}{2.630900in}}%
\pgfpathlineto{\pgfqpoint{2.379653in}{2.703942in}}%
\pgfpathlineto{\pgfqpoint{2.379702in}{2.674154in}}%
\pgfpathlineto{\pgfqpoint{2.380345in}{2.747365in}}%
\pgfpathlineto{\pgfqpoint{2.380395in}{2.717722in}}%
\pgfpathlineto{\pgfqpoint{2.380838in}{2.778269in}}%
\pgfpathlineto{\pgfqpoint{2.380926in}{2.644559in}}%
\pgfpathlineto{\pgfqpoint{2.381040in}{2.644559in}}%
\pgfpathlineto{\pgfqpoint{2.383809in}{2.644559in}}%
\pgfpathlineto{\pgfqpoint{2.383816in}{2.761956in}}%
\pgfpathlineto{\pgfqpoint{2.383922in}{2.630472in}}%
\pgfpathlineto{\pgfqpoint{2.384577in}{2.703244in}}%
\pgfpathlineto{\pgfqpoint{2.384626in}{2.672449in}}%
\pgfpathlineto{\pgfqpoint{2.385269in}{2.746661in}}%
\pgfpathlineto{\pgfqpoint{2.385318in}{2.716022in}}%
\pgfpathlineto{\pgfqpoint{2.385749in}{2.776745in}}%
\pgfpathlineto{\pgfqpoint{2.385874in}{2.644559in}}%
\pgfpathlineto{\pgfqpoint{2.385975in}{2.644559in}}%
\pgfpathlineto{\pgfqpoint{2.388754in}{2.644559in}}%
\pgfpathlineto{\pgfqpoint{2.388761in}{2.770122in}}%
\pgfpathlineto{\pgfqpoint{2.388883in}{2.634271in}}%
\pgfpathlineto{\pgfqpoint{2.389553in}{2.701233in}}%
\pgfpathlineto{\pgfqpoint{2.389603in}{2.678081in}}%
\pgfpathlineto{\pgfqpoint{2.390147in}{2.738437in}}%
\pgfpathlineto{\pgfqpoint{2.390196in}{2.715441in}}%
\pgfpathlineto{\pgfqpoint{2.390729in}{2.774941in}}%
\pgfpathlineto{\pgfqpoint{2.390818in}{2.644559in}}%
\pgfpathlineto{\pgfqpoint{2.390839in}{2.644559in}}%
\pgfpathlineto{\pgfqpoint{2.393700in}{2.644559in}}%
\pgfpathlineto{\pgfqpoint{2.393707in}{2.761810in}}%
\pgfpathlineto{\pgfqpoint{2.393821in}{2.622061in}}%
\pgfpathlineto{\pgfqpoint{2.394437in}{2.659048in}}%
\pgfpathlineto{\pgfqpoint{2.395278in}{2.762334in}}%
\pgfpathlineto{\pgfqpoint{2.395649in}{2.785599in}}%
\pgfpathlineto{\pgfqpoint{2.396172in}{2.644559in}}%
\pgfpathlineto{\pgfqpoint{2.398646in}{2.644559in}}%
\pgfpathlineto{\pgfqpoint{2.398652in}{2.761867in}}%
\pgfpathlineto{\pgfqpoint{2.398766in}{2.622173in}}%
\pgfpathlineto{\pgfqpoint{2.399383in}{2.659180in}}%
\pgfpathlineto{\pgfqpoint{2.400224in}{2.762146in}}%
\pgfpathlineto{\pgfqpoint{2.400595in}{2.785396in}}%
\pgfpathlineto{\pgfqpoint{2.401119in}{2.644559in}}%
\pgfpathlineto{\pgfqpoint{2.403591in}{2.644559in}}%
\pgfpathlineto{\pgfqpoint{2.403598in}{2.762284in}}%
\pgfpathlineto{\pgfqpoint{2.403705in}{2.630322in}}%
\pgfpathlineto{\pgfqpoint{2.404362in}{2.703550in}}%
\pgfpathlineto{\pgfqpoint{2.404412in}{2.672517in}}%
\pgfpathlineto{\pgfqpoint{2.405055in}{2.746966in}}%
\pgfpathlineto{\pgfqpoint{2.405104in}{2.716093in}}%
\pgfpathlineto{\pgfqpoint{2.405532in}{2.776930in}}%
\pgfpathlineto{\pgfqpoint{2.405656in}{2.644559in}}%
\pgfpathlineto{\pgfqpoint{2.405752in}{2.644559in}}%
\pgfpathlineto{\pgfqpoint{2.408539in}{2.645257in}}%
\pgfpathlineto{\pgfqpoint{2.408542in}{2.774079in}}%
\pgfpathlineto{\pgfqpoint{2.408664in}{2.612205in}}%
\pgfpathlineto{\pgfqpoint{2.409285in}{2.649879in}}%
\pgfpathlineto{\pgfqpoint{2.410125in}{2.772525in}}%
\pgfpathlineto{\pgfqpoint{2.410511in}{2.796646in}}%
\pgfpathlineto{\pgfqpoint{2.410963in}{2.644559in}}%
\pgfpathlineto{\pgfqpoint{2.413482in}{2.644559in}}%
\pgfpathlineto{\pgfqpoint{2.413489in}{2.762245in}}%
\pgfpathlineto{\pgfqpoint{2.413595in}{2.630691in}}%
\pgfpathlineto{\pgfqpoint{2.414251in}{2.703074in}}%
\pgfpathlineto{\pgfqpoint{2.414300in}{2.672752in}}%
\pgfpathlineto{\pgfqpoint{2.414943in}{2.746492in}}%
\pgfpathlineto{\pgfqpoint{2.414993in}{2.716325in}}%
\pgfpathlineto{\pgfqpoint{2.415422in}{2.776521in}}%
\pgfpathlineto{\pgfqpoint{2.415547in}{2.644559in}}%
\pgfpathlineto{\pgfqpoint{2.415643in}{2.644559in}}%
\pgfpathlineto{\pgfqpoint{2.418430in}{2.645000in}}%
\pgfpathlineto{\pgfqpoint{2.418433in}{2.774302in}}%
\pgfpathlineto{\pgfqpoint{2.418543in}{2.625629in}}%
\pgfpathlineto{\pgfqpoint{2.419196in}{2.707699in}}%
\pgfpathlineto{\pgfqpoint{2.419245in}{2.667730in}}%
\pgfpathlineto{\pgfqpoint{2.419888in}{2.751105in}}%
\pgfpathlineto{\pgfqpoint{2.419938in}{2.711314in}}%
\pgfpathlineto{\pgfqpoint{2.420368in}{2.781143in}}%
\pgfpathlineto{\pgfqpoint{2.420490in}{2.644559in}}%
\pgfpathlineto{\pgfqpoint{2.420637in}{2.644559in}}%
\pgfpathlineto{\pgfqpoint{2.423373in}{2.644559in}}%
\pgfpathlineto{\pgfqpoint{2.423380in}{2.762379in}}%
\pgfpathlineto{\pgfqpoint{2.423497in}{2.638730in}}%
\pgfpathlineto{\pgfqpoint{2.424115in}{2.675866in}}%
\pgfpathlineto{\pgfqpoint{2.425054in}{2.752274in}}%
\pgfpathlineto{\pgfqpoint{2.425347in}{2.770651in}}%
\pgfpathlineto{\pgfqpoint{2.426158in}{2.644559in}}%
\pgfpathlineto{\pgfqpoint{2.428318in}{2.644559in}}%
\pgfpathlineto{\pgfqpoint{2.428325in}{2.762427in}}%
\pgfpathlineto{\pgfqpoint{2.428438in}{2.634678in}}%
\pgfpathlineto{\pgfqpoint{2.429053in}{2.671461in}}%
\pgfpathlineto{\pgfqpoint{2.429894in}{2.749695in}}%
\pgfpathlineto{\pgfqpoint{2.430266in}{2.773043in}}%
\pgfpathlineto{\pgfqpoint{2.430873in}{2.644559in}}%
\pgfpathlineto{\pgfqpoint{2.433266in}{2.644724in}}%
\pgfpathlineto{\pgfqpoint{2.433269in}{2.773515in}}%
\pgfpathlineto{\pgfqpoint{2.433997in}{2.682664in}}%
\pgfpathlineto{\pgfqpoint{2.435211in}{2.761393in}}%
\pgfpathlineto{\pgfqpoint{2.435238in}{2.738083in}}%
\pgfpathlineto{\pgfqpoint{2.435265in}{2.748193in}}%
\pgfpathlineto{\pgfqpoint{2.435985in}{2.644559in}}%
\pgfpathlineto{\pgfqpoint{2.438209in}{2.644559in}}%
\pgfpathlineto{\pgfqpoint{2.438216in}{2.761847in}}%
\pgfpathlineto{\pgfqpoint{2.438329in}{2.621942in}}%
\pgfpathlineto{\pgfqpoint{2.438944in}{2.658822in}}%
\pgfpathlineto{\pgfqpoint{2.439785in}{2.762269in}}%
\pgfpathlineto{\pgfqpoint{2.440157in}{2.785586in}}%
\pgfpathlineto{\pgfqpoint{2.440659in}{2.644559in}}%
\pgfpathlineto{\pgfqpoint{2.443155in}{2.644559in}}%
\pgfpathlineto{\pgfqpoint{2.443161in}{2.762095in}}%
\pgfpathlineto{\pgfqpoint{2.443275in}{2.621624in}}%
\pgfpathlineto{\pgfqpoint{2.443892in}{2.658593in}}%
\pgfpathlineto{\pgfqpoint{2.444732in}{2.762631in}}%
\pgfpathlineto{\pgfqpoint{2.445104in}{2.785892in}}%
\pgfpathlineto{\pgfqpoint{2.445621in}{2.644559in}}%
\pgfpathlineto{\pgfqpoint{2.448100in}{2.644559in}}%
\pgfpathlineto{\pgfqpoint{2.448107in}{2.762102in}}%
\pgfpathlineto{\pgfqpoint{2.448222in}{2.622770in}}%
\pgfpathlineto{\pgfqpoint{2.448839in}{2.659875in}}%
\pgfpathlineto{\pgfqpoint{2.449680in}{2.761688in}}%
\pgfpathlineto{\pgfqpoint{2.450074in}{2.786401in}}%
\pgfpathlineto{\pgfqpoint{2.450578in}{2.644559in}}%
\pgfpathlineto{\pgfqpoint{2.453046in}{2.644559in}}%
\pgfpathlineto{\pgfqpoint{2.453052in}{2.762199in}}%
\pgfpathlineto{\pgfqpoint{2.453168in}{2.623144in}}%
\pgfpathlineto{\pgfqpoint{2.453786in}{2.660302in}}%
\pgfpathlineto{\pgfqpoint{2.454627in}{2.761354in}}%
\pgfpathlineto{\pgfqpoint{2.455020in}{2.786001in}}%
\pgfpathlineto{\pgfqpoint{2.455518in}{2.644559in}}%
\pgfpathlineto{\pgfqpoint{2.457991in}{2.644559in}}%
\pgfpathlineto{\pgfqpoint{2.457998in}{2.762211in}}%
\pgfpathlineto{\pgfqpoint{2.458111in}{2.624553in}}%
\pgfpathlineto{\pgfqpoint{2.458727in}{2.661440in}}%
\pgfpathlineto{\pgfqpoint{2.459568in}{2.759662in}}%
\pgfpathlineto{\pgfqpoint{2.459940in}{2.782965in}}%
\pgfpathlineto{\pgfqpoint{2.460470in}{2.644559in}}%
\pgfpathlineto{\pgfqpoint{2.462939in}{2.644598in}}%
\pgfpathlineto{\pgfqpoint{2.462941in}{2.767691in}}%
\pgfpathlineto{\pgfqpoint{2.463689in}{2.685605in}}%
\pgfpathlineto{\pgfqpoint{2.464911in}{2.761277in}}%
\pgfpathlineto{\pgfqpoint{2.464955in}{2.731978in}}%
\pgfpathlineto{\pgfqpoint{2.465714in}{2.644559in}}%
\pgfpathlineto{\pgfqpoint{2.467882in}{2.644559in}}%
\pgfpathlineto{\pgfqpoint{2.467889in}{2.761926in}}%
\pgfpathlineto{\pgfqpoint{2.468002in}{2.621874in}}%
\pgfpathlineto{\pgfqpoint{2.468617in}{2.658760in}}%
\pgfpathlineto{\pgfqpoint{2.469458in}{2.762353in}}%
\pgfpathlineto{\pgfqpoint{2.469830in}{2.785664in}}%
\pgfpathlineto{\pgfqpoint{2.470351in}{2.644559in}}%
\pgfpathlineto{\pgfqpoint{2.472827in}{2.644559in}}%
\pgfpathlineto{\pgfqpoint{2.472833in}{2.773247in}}%
\pgfpathlineto{\pgfqpoint{2.472957in}{2.625527in}}%
\pgfpathlineto{\pgfqpoint{2.473629in}{2.710070in}}%
\pgfpathlineto{\pgfqpoint{2.473678in}{2.669509in}}%
\pgfpathlineto{\pgfqpoint{2.474321in}{2.753470in}}%
\pgfpathlineto{\pgfqpoint{2.474370in}{2.713100in}}%
\pgfpathlineto{\pgfqpoint{2.474802in}{2.783603in}}%
\pgfpathlineto{\pgfqpoint{2.474889in}{2.644559in}}%
\pgfpathlineto{\pgfqpoint{2.475058in}{2.644559in}}%
\pgfpathlineto{\pgfqpoint{2.477773in}{2.644559in}}%
\pgfpathlineto{\pgfqpoint{2.477780in}{2.762378in}}%
\pgfpathlineto{\pgfqpoint{2.477896in}{2.638820in}}%
\pgfpathlineto{\pgfqpoint{2.478514in}{2.675944in}}%
\pgfpathlineto{\pgfqpoint{2.479454in}{2.752168in}}%
\pgfpathlineto{\pgfqpoint{2.479747in}{2.770558in}}%
\pgfpathlineto{\pgfqpoint{2.480575in}{2.644559in}}%
\pgfpathlineto{\pgfqpoint{2.482718in}{2.644559in}}%
\pgfpathlineto{\pgfqpoint{2.482724in}{2.771894in}}%
\pgfpathlineto{\pgfqpoint{2.482847in}{2.628769in}}%
\pgfpathlineto{\pgfqpoint{2.483517in}{2.706651in}}%
\pgfpathlineto{\pgfqpoint{2.483566in}{2.672649in}}%
\pgfpathlineto{\pgfqpoint{2.484209in}{2.750031in}}%
\pgfpathlineto{\pgfqpoint{2.484259in}{2.716259in}}%
\pgfpathlineto{\pgfqpoint{2.484693in}{2.780307in}}%
\pgfpathlineto{\pgfqpoint{2.484783in}{2.644559in}}%
\pgfpathlineto{\pgfqpoint{2.484926in}{2.644559in}}%
\pgfpathlineto{\pgfqpoint{2.487667in}{2.644944in}}%
\pgfpathlineto{\pgfqpoint{2.487669in}{2.774165in}}%
\pgfpathlineto{\pgfqpoint{2.487771in}{2.629858in}}%
\pgfpathlineto{\pgfqpoint{2.488417in}{2.702756in}}%
\pgfpathlineto{\pgfqpoint{2.488467in}{2.670751in}}%
\pgfpathlineto{\pgfqpoint{2.489110in}{2.746181in}}%
\pgfpathlineto{\pgfqpoint{2.489159in}{2.714318in}}%
\pgfpathlineto{\pgfqpoint{2.489596in}{2.776709in}}%
\pgfpathlineto{\pgfqpoint{2.489729in}{2.644559in}}%
\pgfpathlineto{\pgfqpoint{2.489812in}{2.644559in}}%
\pgfpathlineto{\pgfqpoint{2.492612in}{2.644846in}}%
\pgfpathlineto{\pgfqpoint{2.492614in}{2.774132in}}%
\pgfpathlineto{\pgfqpoint{2.492748in}{2.635334in}}%
\pgfpathlineto{\pgfqpoint{2.493377in}{2.673749in}}%
\pgfpathlineto{\pgfqpoint{2.494217in}{2.751137in}}%
\pgfpathlineto{\pgfqpoint{2.494584in}{2.774114in}}%
\pgfpathlineto{\pgfqpoint{2.495215in}{2.644559in}}%
\pgfpathlineto{\pgfqpoint{2.497557in}{2.644832in}}%
\pgfpathlineto{\pgfqpoint{2.497560in}{2.774160in}}%
\pgfpathlineto{\pgfqpoint{2.497693in}{2.635304in}}%
\pgfpathlineto{\pgfqpoint{2.498322in}{2.673729in}}%
\pgfpathlineto{\pgfqpoint{2.499163in}{2.751165in}}%
\pgfpathlineto{\pgfqpoint{2.499529in}{2.774139in}}%
\pgfpathlineto{\pgfqpoint{2.500157in}{2.644559in}}%
\pgfpathlineto{\pgfqpoint{2.502503in}{2.644801in}}%
\pgfpathlineto{\pgfqpoint{2.502505in}{2.774037in}}%
\pgfpathlineto{\pgfqpoint{2.502628in}{2.637360in}}%
\pgfpathlineto{\pgfqpoint{2.503248in}{2.674740in}}%
\pgfpathlineto{\pgfqpoint{2.504187in}{2.753854in}}%
\pgfpathlineto{\pgfqpoint{2.504474in}{2.771881in}}%
\pgfpathlineto{\pgfqpoint{2.505256in}{2.644559in}}%
\pgfpathlineto{\pgfqpoint{2.507445in}{2.644559in}}%
\pgfpathlineto{\pgfqpoint{2.507451in}{2.772306in}}%
\pgfpathlineto{\pgfqpoint{2.507566in}{2.632097in}}%
\pgfpathlineto{\pgfqpoint{2.508231in}{2.702531in}}%
\pgfpathlineto{\pgfqpoint{2.508281in}{2.675156in}}%
\pgfpathlineto{\pgfqpoint{2.508924in}{2.745959in}}%
\pgfpathlineto{\pgfqpoint{2.508973in}{2.718718in}}%
\pgfpathlineto{\pgfqpoint{2.509394in}{2.775468in}}%
\pgfpathlineto{\pgfqpoint{2.509511in}{2.644559in}}%
\pgfpathlineto{\pgfqpoint{2.509632in}{2.644559in}}%
\pgfpathlineto{\pgfqpoint{2.512394in}{2.644620in}}%
\pgfpathlineto{\pgfqpoint{2.512396in}{2.770015in}}%
\pgfpathlineto{\pgfqpoint{2.513138in}{2.684488in}}%
\pgfpathlineto{\pgfqpoint{2.514365in}{2.762082in}}%
\pgfpathlineto{\pgfqpoint{2.514390in}{2.747342in}}%
\pgfpathlineto{\pgfqpoint{2.515190in}{2.644559in}}%
\pgfpathlineto{\pgfqpoint{2.517336in}{2.644559in}}%
\pgfpathlineto{\pgfqpoint{2.517342in}{2.769899in}}%
\pgfpathlineto{\pgfqpoint{2.517456in}{2.639562in}}%
\pgfpathlineto{\pgfqpoint{2.518121in}{2.695018in}}%
\pgfpathlineto{\pgfqpoint{2.518170in}{2.682507in}}%
\pgfpathlineto{\pgfqpoint{2.518714in}{2.732242in}}%
\pgfpathlineto{\pgfqpoint{2.518763in}{2.719846in}}%
\pgfpathlineto{\pgfqpoint{2.519284in}{2.768010in}}%
\pgfpathlineto{\pgfqpoint{2.519372in}{2.668742in}}%
\pgfpathlineto{\pgfqpoint{2.520122in}{2.644559in}}%
\pgfpathlineto{\pgfqpoint{2.522285in}{2.644616in}}%
\pgfpathlineto{\pgfqpoint{2.522287in}{2.769638in}}%
\pgfpathlineto{\pgfqpoint{2.523031in}{2.684180in}}%
\pgfpathlineto{\pgfqpoint{2.524256in}{2.762516in}}%
\pgfpathlineto{\pgfqpoint{2.524257in}{2.738183in}}%
\pgfpathlineto{\pgfqpoint{2.524279in}{2.746321in}}%
\pgfpathlineto{\pgfqpoint{2.524989in}{2.644559in}}%
\pgfpathlineto{\pgfqpoint{2.527227in}{2.644559in}}%
\pgfpathlineto{\pgfqpoint{2.527233in}{2.770234in}}%
\pgfpathlineto{\pgfqpoint{2.527352in}{2.636241in}}%
\pgfpathlineto{\pgfqpoint{2.528020in}{2.698878in}}%
\pgfpathlineto{\pgfqpoint{2.528070in}{2.679728in}}%
\pgfpathlineto{\pgfqpoint{2.528614in}{2.736089in}}%
\pgfpathlineto{\pgfqpoint{2.528663in}{2.717080in}}%
\pgfpathlineto{\pgfqpoint{2.529202in}{2.772960in}}%
\pgfpathlineto{\pgfqpoint{2.529289in}{2.644559in}}%
\pgfpathlineto{\pgfqpoint{2.532172in}{2.644559in}}%
\pgfpathlineto{\pgfqpoint{2.532178in}{2.772370in}}%
\pgfpathlineto{\pgfqpoint{2.532303in}{2.628141in}}%
\pgfpathlineto{\pgfqpoint{2.532974in}{2.707494in}}%
\pgfpathlineto{\pgfqpoint{2.533023in}{2.672078in}}%
\pgfpathlineto{\pgfqpoint{2.533666in}{2.750904in}}%
\pgfpathlineto{\pgfqpoint{2.533716in}{2.715658in}}%
\pgfpathlineto{\pgfqpoint{2.534147in}{2.781056in}}%
\pgfpathlineto{\pgfqpoint{2.534233in}{2.644559in}}%
\pgfpathlineto{\pgfqpoint{2.534387in}{2.644559in}}%
\pgfpathlineto{\pgfqpoint{2.537118in}{2.644559in}}%
\pgfpathlineto{\pgfqpoint{2.537124in}{2.772963in}}%
\pgfpathlineto{\pgfqpoint{2.537250in}{2.626621in}}%
\pgfpathlineto{\pgfqpoint{2.537922in}{2.709172in}}%
\pgfpathlineto{\pgfqpoint{2.537971in}{2.670720in}}%
\pgfpathlineto{\pgfqpoint{2.538614in}{2.752576in}}%
\pgfpathlineto{\pgfqpoint{2.538664in}{2.714307in}}%
\pgfpathlineto{\pgfqpoint{2.539093in}{2.782558in}}%
\pgfpathlineto{\pgfqpoint{2.539180in}{2.644559in}}%
\pgfpathlineto{\pgfqpoint{2.539349in}{2.644559in}}%
\pgfpathlineto{\pgfqpoint{2.542063in}{2.644559in}}%
\pgfpathlineto{\pgfqpoint{2.542069in}{2.771954in}}%
\pgfpathlineto{\pgfqpoint{2.542186in}{2.631312in}}%
\pgfpathlineto{\pgfqpoint{2.542852in}{2.703466in}}%
\pgfpathlineto{\pgfqpoint{2.542901in}{2.674496in}}%
\pgfpathlineto{\pgfqpoint{2.543544in}{2.746890in}}%
\pgfpathlineto{\pgfqpoint{2.543594in}{2.718062in}}%
\pgfpathlineto{\pgfqpoint{2.544038in}{2.777864in}}%
\pgfpathlineto{\pgfqpoint{2.544125in}{2.644559in}}%
\pgfpathlineto{\pgfqpoint{2.544242in}{2.644559in}}%
\pgfpathlineto{\pgfqpoint{2.547009in}{2.644559in}}%
\pgfpathlineto{\pgfqpoint{2.547015in}{2.773275in}}%
\pgfpathlineto{\pgfqpoint{2.547141in}{2.623678in}}%
\pgfpathlineto{\pgfqpoint{2.547814in}{2.712150in}}%
\pgfpathlineto{\pgfqpoint{2.547863in}{2.667847in}}%
\pgfpathlineto{\pgfqpoint{2.548506in}{2.755541in}}%
\pgfpathlineto{\pgfqpoint{2.548555in}{2.711445in}}%
\pgfpathlineto{\pgfqpoint{2.548984in}{2.785458in}}%
\pgfpathlineto{\pgfqpoint{2.549073in}{2.644559in}}%
\pgfpathlineto{\pgfqpoint{2.549239in}{2.644559in}}%
\pgfpathlineto{\pgfqpoint{2.551954in}{2.644559in}}%
\pgfpathlineto{\pgfqpoint{2.551960in}{2.771981in}}%
\pgfpathlineto{\pgfqpoint{2.552084in}{2.628635in}}%
\pgfpathlineto{\pgfqpoint{2.552754in}{2.706935in}}%
\pgfpathlineto{\pgfqpoint{2.552804in}{2.672500in}}%
\pgfpathlineto{\pgfqpoint{2.553447in}{2.750347in}}%
\pgfpathlineto{\pgfqpoint{2.553496in}{2.716077in}}%
\pgfpathlineto{\pgfqpoint{2.553929in}{2.780572in}}%
\pgfpathlineto{\pgfqpoint{2.554015in}{2.644559in}}%
\pgfpathlineto{\pgfqpoint{2.554136in}{2.644559in}}%
\pgfpathlineto{\pgfqpoint{2.556903in}{2.644922in}}%
\pgfpathlineto{\pgfqpoint{2.556905in}{2.774269in}}%
\pgfpathlineto{\pgfqpoint{2.557009in}{2.628899in}}%
\pgfpathlineto{\pgfqpoint{2.557658in}{2.703876in}}%
\pgfpathlineto{\pgfqpoint{2.557707in}{2.670143in}}%
\pgfpathlineto{\pgfqpoint{2.558350in}{2.747297in}}%
\pgfpathlineto{\pgfqpoint{2.558399in}{2.713713in}}%
\pgfpathlineto{\pgfqpoint{2.558835in}{2.777693in}}%
\pgfpathlineto{\pgfqpoint{2.558966in}{2.644559in}}%
\pgfpathlineto{\pgfqpoint{2.559068in}{2.644559in}}%
\pgfpathlineto{\pgfqpoint{2.561848in}{2.644577in}}%
\pgfpathlineto{\pgfqpoint{2.561850in}{2.761813in}}%
\pgfpathlineto{\pgfqpoint{2.562591in}{2.685954in}}%
\pgfpathlineto{\pgfqpoint{2.562740in}{2.692642in}}%
\pgfpathlineto{\pgfqpoint{2.563820in}{2.760493in}}%
\pgfpathlineto{\pgfqpoint{2.563820in}{2.738057in}}%
\pgfpathlineto{\pgfqpoint{2.563847in}{2.748137in}}%
\pgfpathlineto{\pgfqpoint{2.564575in}{2.644559in}}%
\pgfpathlineto{\pgfqpoint{2.566794in}{2.645314in}}%
\pgfpathlineto{\pgfqpoint{2.566796in}{2.773903in}}%
\pgfpathlineto{\pgfqpoint{2.566919in}{2.612219in}}%
\pgfpathlineto{\pgfqpoint{2.567540in}{2.649922in}}%
\pgfpathlineto{\pgfqpoint{2.568381in}{2.772572in}}%
\pgfpathlineto{\pgfqpoint{2.568765in}{2.796650in}}%
\pgfpathlineto{\pgfqpoint{2.569210in}{2.644559in}}%
\pgfpathlineto{\pgfqpoint{2.571739in}{2.644621in}}%
\pgfpathlineto{\pgfqpoint{2.571741in}{2.770192in}}%
\pgfpathlineto{\pgfqpoint{2.572487in}{2.685110in}}%
\pgfpathlineto{\pgfqpoint{2.573711in}{2.761677in}}%
\pgfpathlineto{\pgfqpoint{2.573736in}{2.747624in}}%
\pgfpathlineto{\pgfqpoint{2.574504in}{2.644559in}}%
\pgfpathlineto{\pgfqpoint{2.576681in}{2.644559in}}%
\pgfpathlineto{\pgfqpoint{2.576687in}{2.773740in}}%
\pgfpathlineto{\pgfqpoint{2.576811in}{2.612279in}}%
\pgfpathlineto{\pgfqpoint{2.577482in}{2.723081in}}%
\pgfpathlineto{\pgfqpoint{2.577531in}{2.656285in}}%
\pgfpathlineto{\pgfqpoint{2.578174in}{2.766425in}}%
\pgfpathlineto{\pgfqpoint{2.578223in}{2.699932in}}%
\pgfpathlineto{\pgfqpoint{2.578656in}{2.796610in}}%
\pgfpathlineto{\pgfqpoint{2.578737in}{2.644559in}}%
\pgfpathlineto{\pgfqpoint{2.578953in}{2.644559in}}%
\pgfpathlineto{\pgfqpoint{2.581630in}{2.644600in}}%
\pgfpathlineto{\pgfqpoint{2.581632in}{2.767904in}}%
\pgfpathlineto{\pgfqpoint{2.582382in}{2.685295in}}%
\pgfpathlineto{\pgfqpoint{2.583602in}{2.761726in}}%
\pgfpathlineto{\pgfqpoint{2.583644in}{2.724562in}}%
\pgfpathlineto{\pgfqpoint{2.584410in}{2.644559in}}%
\pgfpathlineto{\pgfqpoint{2.586572in}{2.644559in}}%
\pgfpathlineto{\pgfqpoint{2.586578in}{2.773275in}}%
\pgfpathlineto{\pgfqpoint{2.586706in}{2.612763in}}%
\pgfpathlineto{\pgfqpoint{2.587379in}{2.723075in}}%
\pgfpathlineto{\pgfqpoint{2.587428in}{2.657105in}}%
\pgfpathlineto{\pgfqpoint{2.588071in}{2.766419in}}%
\pgfpathlineto{\pgfqpoint{2.588121in}{2.700751in}}%
\pgfpathlineto{\pgfqpoint{2.588547in}{2.796202in}}%
\pgfpathlineto{\pgfqpoint{2.588628in}{2.644559in}}%
\pgfpathlineto{\pgfqpoint{2.588851in}{2.644559in}}%
\pgfpathlineto{\pgfqpoint{2.591521in}{2.644590in}}%
\pgfpathlineto{\pgfqpoint{2.591523in}{2.766167in}}%
\pgfpathlineto{\pgfqpoint{2.592276in}{2.685501in}}%
\pgfpathlineto{\pgfqpoint{2.593493in}{2.761741in}}%
\pgfpathlineto{\pgfqpoint{2.593535in}{2.722835in}}%
\pgfpathlineto{\pgfqpoint{2.594327in}{2.644559in}}%
\pgfpathlineto{\pgfqpoint{2.596467in}{2.644932in}}%
\pgfpathlineto{\pgfqpoint{2.596469in}{2.774258in}}%
\pgfpathlineto{\pgfqpoint{2.596575in}{2.628444in}}%
\pgfpathlineto{\pgfqpoint{2.597224in}{2.704445in}}%
\pgfpathlineto{\pgfqpoint{2.597273in}{2.669871in}}%
\pgfpathlineto{\pgfqpoint{2.597916in}{2.747865in}}%
\pgfpathlineto{\pgfqpoint{2.597965in}{2.713442in}}%
\pgfpathlineto{\pgfqpoint{2.598400in}{2.778186in}}%
\pgfpathlineto{\pgfqpoint{2.598529in}{2.644559in}}%
\pgfpathlineto{\pgfqpoint{2.598622in}{2.644559in}}%
\pgfpathlineto{\pgfqpoint{2.601412in}{2.644568in}}%
\pgfpathlineto{\pgfqpoint{2.601416in}{2.782716in}}%
\pgfpathlineto{\pgfqpoint{2.601522in}{2.636838in}}%
\pgfpathlineto{\pgfqpoint{2.602174in}{2.696286in}}%
\pgfpathlineto{\pgfqpoint{2.602224in}{2.678511in}}%
\pgfpathlineto{\pgfqpoint{2.602768in}{2.733520in}}%
\pgfpathlineto{\pgfqpoint{2.602817in}{2.715840in}}%
\pgfpathlineto{\pgfqpoint{2.603348in}{2.769900in}}%
\pgfpathlineto{\pgfqpoint{2.603447in}{2.668439in}}%
\pgfpathlineto{\pgfqpoint{2.604226in}{2.644559in}}%
\pgfpathlineto{\pgfqpoint{2.606357in}{2.644586in}}%
\pgfpathlineto{\pgfqpoint{2.606360in}{2.764927in}}%
\pgfpathlineto{\pgfqpoint{2.607121in}{2.685986in}}%
\pgfpathlineto{\pgfqpoint{2.608329in}{2.761727in}}%
\pgfpathlineto{\pgfqpoint{2.608371in}{2.724647in}}%
\pgfpathlineto{\pgfqpoint{2.609139in}{2.644559in}}%
\pgfpathlineto{\pgfqpoint{2.611300in}{2.644559in}}%
\pgfpathlineto{\pgfqpoint{2.611306in}{2.773433in}}%
\pgfpathlineto{\pgfqpoint{2.611432in}{2.612649in}}%
\pgfpathlineto{\pgfqpoint{2.612104in}{2.723026in}}%
\pgfpathlineto{\pgfqpoint{2.612153in}{2.656877in}}%
\pgfpathlineto{\pgfqpoint{2.612796in}{2.766371in}}%
\pgfpathlineto{\pgfqpoint{2.612846in}{2.700523in}}%
\pgfpathlineto{\pgfqpoint{2.613274in}{2.796289in}}%
\pgfpathlineto{\pgfqpoint{2.613355in}{2.644559in}}%
\pgfpathlineto{\pgfqpoint{2.613578in}{2.644559in}}%
\pgfpathlineto{\pgfqpoint{2.616248in}{2.644577in}}%
\pgfpathlineto{\pgfqpoint{2.616250in}{2.761707in}}%
\pgfpathlineto{\pgfqpoint{2.617013in}{2.687294in}}%
\pgfpathlineto{\pgfqpoint{2.617161in}{2.693970in}}%
\pgfpathlineto{\pgfqpoint{2.618185in}{2.760948in}}%
\pgfpathlineto{\pgfqpoint{2.618220in}{2.738055in}}%
\pgfpathlineto{\pgfqpoint{2.618247in}{2.748134in}}%
\pgfpathlineto{\pgfqpoint{2.618990in}{2.644559in}}%
\pgfpathlineto{\pgfqpoint{2.621191in}{2.644559in}}%
\pgfpathlineto{\pgfqpoint{2.621197in}{2.771933in}}%
\pgfpathlineto{\pgfqpoint{2.621323in}{2.626451in}}%
\pgfpathlineto{\pgfqpoint{2.621995in}{2.709396in}}%
\pgfpathlineto{\pgfqpoint{2.622044in}{2.670538in}}%
\pgfpathlineto{\pgfqpoint{2.622687in}{2.752799in}}%
\pgfpathlineto{\pgfqpoint{2.622737in}{2.714126in}}%
\pgfpathlineto{\pgfqpoint{2.623165in}{2.782763in}}%
\pgfpathlineto{\pgfqpoint{2.623252in}{2.644559in}}%
\pgfpathlineto{\pgfqpoint{2.623422in}{2.644559in}}%
\pgfpathlineto{\pgfqpoint{2.626139in}{2.644637in}}%
\pgfpathlineto{\pgfqpoint{2.626142in}{2.771092in}}%
\pgfpathlineto{\pgfqpoint{2.626883in}{2.684701in}}%
\pgfpathlineto{\pgfqpoint{2.628111in}{2.761859in}}%
\pgfpathlineto{\pgfqpoint{2.628136in}{2.747415in}}%
\pgfpathlineto{\pgfqpoint{2.628934in}{2.644559in}}%
\pgfpathlineto{\pgfqpoint{2.631085in}{2.644702in}}%
\pgfpathlineto{\pgfqpoint{2.631087in}{2.773061in}}%
\pgfpathlineto{\pgfqpoint{2.631821in}{2.683186in}}%
\pgfpathlineto{\pgfqpoint{2.633056in}{2.762950in}}%
\pgfpathlineto{\pgfqpoint{2.633057in}{2.738057in}}%
\pgfpathlineto{\pgfqpoint{2.633083in}{2.748186in}}%
\pgfpathlineto{\pgfqpoint{2.633805in}{2.644559in}}%
\pgfpathlineto{\pgfqpoint{2.636028in}{2.644559in}}%
\pgfpathlineto{\pgfqpoint{2.636035in}{2.754977in}}%
\pgfpathlineto{\pgfqpoint{2.636131in}{2.642350in}}%
\pgfpathlineto{\pgfqpoint{2.636775in}{2.689982in}}%
\pgfpathlineto{\pgfqpoint{2.636825in}{2.682734in}}%
\pgfpathlineto{\pgfqpoint{2.637270in}{2.721025in}}%
\pgfpathlineto{\pgfqpoint{2.637319in}{2.713829in}}%
\pgfpathlineto{\pgfqpoint{2.637957in}{2.764173in}}%
\pgfpathlineto{\pgfqpoint{2.638041in}{2.709838in}}%
\pgfpathlineto{\pgfqpoint{2.638051in}{2.725768in}}%
\pgfpathlineto{\pgfqpoint{2.638822in}{2.644559in}}%
\pgfpathlineto{\pgfqpoint{2.640975in}{2.644565in}}%
\pgfpathlineto{\pgfqpoint{2.640979in}{2.780784in}}%
\pgfpathlineto{\pgfqpoint{2.641089in}{2.635259in}}%
\pgfpathlineto{\pgfqpoint{2.641745in}{2.698284in}}%
\pgfpathlineto{\pgfqpoint{2.641794in}{2.677394in}}%
\pgfpathlineto{\pgfqpoint{2.642338in}{2.735512in}}%
\pgfpathlineto{\pgfqpoint{2.642388in}{2.714729in}}%
\pgfpathlineto{\pgfqpoint{2.642915in}{2.771661in}}%
\pgfpathlineto{\pgfqpoint{2.643024in}{2.644561in}}%
\pgfpathlineto{\pgfqpoint{2.645919in}{2.644559in}}%
\pgfpathlineto{\pgfqpoint{2.645926in}{2.754240in}}%
\pgfpathlineto{\pgfqpoint{2.646058in}{2.633052in}}%
\pgfpathlineto{\pgfqpoint{2.646687in}{2.671601in}}%
\pgfpathlineto{\pgfqpoint{2.647528in}{2.753215in}}%
\pgfpathlineto{\pgfqpoint{2.647893in}{2.776099in}}%
\pgfpathlineto{\pgfqpoint{2.648493in}{2.644559in}}%
\pgfpathlineto{\pgfqpoint{2.650867in}{2.645069in}}%
\pgfpathlineto{\pgfqpoint{2.650869in}{2.774166in}}%
\pgfpathlineto{\pgfqpoint{2.650977in}{2.622911in}}%
\pgfpathlineto{\pgfqpoint{2.651630in}{2.710341in}}%
\pgfpathlineto{\pgfqpoint{2.651679in}{2.664780in}}%
\pgfpathlineto{\pgfqpoint{2.652322in}{2.753735in}}%
\pgfpathlineto{\pgfqpoint{2.652372in}{2.708377in}}%
\pgfpathlineto{\pgfqpoint{2.652803in}{2.783840in}}%
\pgfpathlineto{\pgfqpoint{2.652930in}{2.644559in}}%
\pgfpathlineto{\pgfqpoint{2.653079in}{2.644559in}}%
\pgfpathlineto{\pgfqpoint{2.655810in}{2.644559in}}%
\pgfpathlineto{\pgfqpoint{2.655816in}{2.755385in}}%
\pgfpathlineto{\pgfqpoint{2.655924in}{2.640067in}}%
\pgfpathlineto{\pgfqpoint{2.656575in}{2.692951in}}%
\pgfpathlineto{\pgfqpoint{2.656625in}{2.681898in}}%
\pgfpathlineto{\pgfqpoint{2.657169in}{2.730195in}}%
\pgfpathlineto{\pgfqpoint{2.657218in}{2.719218in}}%
\pgfpathlineto{\pgfqpoint{2.657748in}{2.766552in}}%
\pgfpathlineto{\pgfqpoint{2.657851in}{2.674421in}}%
\pgfpathlineto{\pgfqpoint{2.658597in}{2.644559in}}%
\pgfpathlineto{\pgfqpoint{2.660758in}{2.645299in}}%
\pgfpathlineto{\pgfqpoint{2.660760in}{2.773839in}}%
\pgfpathlineto{\pgfqpoint{2.660883in}{2.612389in}}%
\pgfpathlineto{\pgfqpoint{2.661503in}{2.650037in}}%
\pgfpathlineto{\pgfqpoint{2.662343in}{2.772348in}}%
\pgfpathlineto{\pgfqpoint{2.662729in}{2.796478in}}%
\pgfpathlineto{\pgfqpoint{2.663204in}{2.644559in}}%
\pgfpathlineto{\pgfqpoint{2.665701in}{2.644559in}}%
\pgfpathlineto{\pgfqpoint{2.665707in}{2.755276in}}%
\pgfpathlineto{\pgfqpoint{2.665806in}{2.641608in}}%
\pgfpathlineto{\pgfqpoint{2.666452in}{2.690873in}}%
\pgfpathlineto{\pgfqpoint{2.666502in}{2.682384in}}%
\pgfpathlineto{\pgfqpoint{2.666947in}{2.721913in}}%
\pgfpathlineto{\pgfqpoint{2.666996in}{2.713481in}}%
\pgfpathlineto{\pgfqpoint{2.667632in}{2.764915in}}%
\pgfpathlineto{\pgfqpoint{2.667714in}{2.711918in}}%
\pgfpathlineto{\pgfqpoint{2.667725in}{2.720082in}}%
\pgfpathlineto{\pgfqpoint{2.668472in}{2.644559in}}%
\pgfpathlineto{\pgfqpoint{2.670648in}{2.644561in}}%
\pgfpathlineto{\pgfqpoint{2.670652in}{2.773828in}}%
\pgfpathlineto{\pgfqpoint{2.670773in}{2.637141in}}%
\pgfpathlineto{\pgfqpoint{2.671392in}{2.674466in}}%
\pgfpathlineto{\pgfqpoint{2.672232in}{2.747508in}}%
\pgfpathlineto{\pgfqpoint{2.672620in}{2.771811in}}%
\pgfpathlineto{\pgfqpoint{2.673273in}{2.644559in}}%
\pgfpathlineto{\pgfqpoint{2.675592in}{2.644559in}}%
\pgfpathlineto{\pgfqpoint{2.675598in}{2.754875in}}%
\pgfpathlineto{\pgfqpoint{2.675691in}{2.643349in}}%
\pgfpathlineto{\pgfqpoint{2.676334in}{2.689024in}}%
\pgfpathlineto{\pgfqpoint{2.676384in}{2.683307in}}%
\pgfpathlineto{\pgfqpoint{2.676829in}{2.720069in}}%
\pgfpathlineto{\pgfqpoint{2.676878in}{2.714399in}}%
\pgfpathlineto{\pgfqpoint{2.677518in}{2.763367in}}%
\pgfpathlineto{\pgfqpoint{2.677601in}{2.704366in}}%
\pgfpathlineto{\pgfqpoint{2.677610in}{2.734461in}}%
\pgfpathlineto{\pgfqpoint{2.678348in}{2.644559in}}%
\pgfpathlineto{\pgfqpoint{2.680539in}{2.644563in}}%
\pgfpathlineto{\pgfqpoint{2.680543in}{2.777868in}}%
\pgfpathlineto{\pgfqpoint{2.680655in}{2.634471in}}%
\pgfpathlineto{\pgfqpoint{2.681266in}{2.670808in}}%
\pgfpathlineto{\pgfqpoint{2.682106in}{2.749058in}}%
\pgfpathlineto{\pgfqpoint{2.682482in}{2.772595in}}%
\pgfpathlineto{\pgfqpoint{2.683090in}{2.644559in}}%
\pgfpathlineto{\pgfqpoint{2.685484in}{2.644559in}}%
\pgfpathlineto{\pgfqpoint{2.685488in}{2.765184in}}%
\pgfpathlineto{\pgfqpoint{2.685616in}{2.631467in}}%
\pgfpathlineto{\pgfqpoint{2.686241in}{2.669550in}}%
\pgfpathlineto{\pgfqpoint{2.687082in}{2.754013in}}%
\pgfpathlineto{\pgfqpoint{2.687456in}{2.777488in}}%
\pgfpathlineto{\pgfqpoint{2.688041in}{2.644559in}}%
\pgfpathlineto{\pgfqpoint{2.690427in}{2.644559in}}%
\pgfpathlineto{\pgfqpoint{2.690435in}{2.756013in}}%
\pgfpathlineto{\pgfqpoint{2.691164in}{2.686938in}}%
\pgfpathlineto{\pgfqpoint{2.691214in}{2.684595in}}%
\pgfpathlineto{\pgfqpoint{2.691461in}{2.705568in}}%
\pgfpathlineto{\pgfqpoint{2.691510in}{2.703247in}}%
\pgfpathlineto{\pgfqpoint{2.692402in}{2.764622in}}%
\pgfpathlineto{\pgfqpoint{2.692429in}{2.748183in}}%
\pgfpathlineto{\pgfqpoint{2.693217in}{2.644559in}}%
\pgfpathlineto{\pgfqpoint{2.695372in}{2.644559in}}%
\pgfpathlineto{\pgfqpoint{2.695379in}{2.762512in}}%
\pgfpathlineto{\pgfqpoint{2.695498in}{2.632879in}}%
\pgfpathlineto{\pgfqpoint{2.696116in}{2.670130in}}%
\pgfpathlineto{\pgfqpoint{2.696957in}{2.751538in}}%
\pgfpathlineto{\pgfqpoint{2.697347in}{2.775986in}}%
\pgfpathlineto{\pgfqpoint{2.697901in}{2.644559in}}%
\pgfpathlineto{\pgfqpoint{2.700319in}{2.644559in}}%
\pgfpathlineto{\pgfqpoint{2.700326in}{2.754223in}}%
\pgfpathlineto{\pgfqpoint{2.700458in}{2.633309in}}%
\pgfpathlineto{\pgfqpoint{2.701087in}{2.671854in}}%
\pgfpathlineto{\pgfqpoint{2.701927in}{2.753043in}}%
\pgfpathlineto{\pgfqpoint{2.702293in}{2.775941in}}%
\pgfpathlineto{\pgfqpoint{2.702898in}{2.644559in}}%
\pgfpathlineto{\pgfqpoint{2.705266in}{2.644590in}}%
\pgfpathlineto{\pgfqpoint{2.705269in}{2.765998in}}%
\pgfpathlineto{\pgfqpoint{2.706022in}{2.685461in}}%
\pgfpathlineto{\pgfqpoint{2.707238in}{2.761793in}}%
\pgfpathlineto{\pgfqpoint{2.707281in}{2.724878in}}%
\pgfpathlineto{\pgfqpoint{2.708063in}{2.644559in}}%
\pgfpathlineto{\pgfqpoint{2.710210in}{2.644559in}}%
\pgfpathlineto{\pgfqpoint{2.710217in}{2.754184in}}%
\pgfpathlineto{\pgfqpoint{2.710348in}{2.631898in}}%
\pgfpathlineto{\pgfqpoint{2.710976in}{2.670399in}}%
\pgfpathlineto{\pgfqpoint{2.711817in}{2.754382in}}%
\pgfpathlineto{\pgfqpoint{2.712184in}{2.777356in}}%
\pgfpathlineto{\pgfqpoint{2.712777in}{2.644559in}}%
\pgfpathlineto{\pgfqpoint{2.715157in}{2.644572in}}%
\pgfpathlineto{\pgfqpoint{2.715159in}{2.758517in}}%
\pgfpathlineto{\pgfqpoint{2.715279in}{2.635675in}}%
\pgfpathlineto{\pgfqpoint{2.715896in}{2.672678in}}%
\pgfpathlineto{\pgfqpoint{2.716737in}{2.748759in}}%
\pgfpathlineto{\pgfqpoint{2.717129in}{2.773387in}}%
\pgfpathlineto{\pgfqpoint{2.717750in}{2.644559in}}%
\pgfpathlineto{\pgfqpoint{2.720099in}{2.644559in}}%
\pgfpathlineto{\pgfqpoint{2.720106in}{2.766745in}}%
\pgfpathlineto{\pgfqpoint{2.720225in}{2.635837in}}%
\pgfpathlineto{\pgfqpoint{2.720893in}{2.699063in}}%
\pgfpathlineto{\pgfqpoint{2.720942in}{2.679286in}}%
\pgfpathlineto{\pgfqpoint{2.721486in}{2.736274in}}%
\pgfpathlineto{\pgfqpoint{2.721536in}{2.716639in}}%
\pgfpathlineto{\pgfqpoint{2.722074in}{2.773155in}}%
\pgfpathlineto{\pgfqpoint{2.722161in}{2.644559in}}%
\pgfpathlineto{\pgfqpoint{2.725048in}{2.644571in}}%
\pgfpathlineto{\pgfqpoint{2.725050in}{2.758057in}}%
\pgfpathlineto{\pgfqpoint{2.725177in}{2.634598in}}%
\pgfpathlineto{\pgfqpoint{2.725798in}{2.672257in}}%
\pgfpathlineto{\pgfqpoint{2.726639in}{2.750630in}}%
\pgfpathlineto{\pgfqpoint{2.727020in}{2.774523in}}%
\pgfpathlineto{\pgfqpoint{2.727609in}{2.644559in}}%
\pgfpathlineto{\pgfqpoint{2.729992in}{2.644559in}}%
\pgfpathlineto{\pgfqpoint{2.729999in}{2.754384in}}%
\pgfpathlineto{\pgfqpoint{2.730131in}{2.632475in}}%
\pgfpathlineto{\pgfqpoint{2.730760in}{2.671048in}}%
\pgfpathlineto{\pgfqpoint{2.731601in}{2.753833in}}%
\pgfpathlineto{\pgfqpoint{2.731965in}{2.776671in}}%
\pgfpathlineto{\pgfqpoint{2.732547in}{2.644559in}}%
\pgfpathlineto{\pgfqpoint{2.734939in}{2.644609in}}%
\pgfpathlineto{\pgfqpoint{2.734941in}{2.768945in}}%
\pgfpathlineto{\pgfqpoint{2.735687in}{2.684273in}}%
\pgfpathlineto{\pgfqpoint{2.736911in}{2.762502in}}%
\pgfpathlineto{\pgfqpoint{2.736911in}{2.738181in}}%
\pgfpathlineto{\pgfqpoint{2.736934in}{2.746278in}}%
\pgfpathlineto{\pgfqpoint{2.737682in}{2.644559in}}%
\pgfpathlineto{\pgfqpoint{2.739883in}{2.644559in}}%
\pgfpathlineto{\pgfqpoint{2.739890in}{2.753844in}}%
\pgfpathlineto{\pgfqpoint{2.740015in}{2.633131in}}%
\pgfpathlineto{\pgfqpoint{2.740639in}{2.671110in}}%
\pgfpathlineto{\pgfqpoint{2.741480in}{2.752386in}}%
\pgfpathlineto{\pgfqpoint{2.741856in}{2.775998in}}%
\pgfpathlineto{\pgfqpoint{2.742457in}{2.644559in}}%
\pgfpathlineto{\pgfqpoint{2.744830in}{2.644564in}}%
\pgfpathlineto{\pgfqpoint{2.744834in}{2.779274in}}%
\pgfpathlineto{\pgfqpoint{2.744943in}{2.635295in}}%
\pgfpathlineto{\pgfqpoint{2.745599in}{2.698207in}}%
\pgfpathlineto{\pgfqpoint{2.745649in}{2.677420in}}%
\pgfpathlineto{\pgfqpoint{2.746193in}{2.735435in}}%
\pgfpathlineto{\pgfqpoint{2.746242in}{2.714755in}}%
\pgfpathlineto{\pgfqpoint{2.746769in}{2.771594in}}%
\pgfpathlineto{\pgfqpoint{2.746879in}{2.644569in}}%
\pgfpathlineto{\pgfqpoint{2.748287in}{2.644559in}}%
\pgfpathlineto{\pgfqpoint{2.749774in}{2.644559in}}%
\pgfpathlineto{\pgfqpoint{2.749780in}{2.754711in}}%
\pgfpathlineto{\pgfqpoint{2.749889in}{2.639370in}}%
\pgfpathlineto{\pgfqpoint{2.750541in}{2.693775in}}%
\pgfpathlineto{\pgfqpoint{2.750591in}{2.681364in}}%
\pgfpathlineto{\pgfqpoint{2.751135in}{2.731017in}}%
\pgfpathlineto{\pgfqpoint{2.751184in}{2.718686in}}%
\pgfpathlineto{\pgfqpoint{2.751713in}{2.767294in}}%
\pgfpathlineto{\pgfqpoint{2.751814in}{2.675363in}}%
\pgfpathlineto{\pgfqpoint{2.752567in}{2.644559in}}%
\pgfpathlineto{\pgfqpoint{2.754721in}{2.644629in}}%
\pgfpathlineto{\pgfqpoint{2.754723in}{2.770573in}}%
\pgfpathlineto{\pgfqpoint{2.755464in}{2.684339in}}%
\pgfpathlineto{\pgfqpoint{2.756693in}{2.762154in}}%
\pgfpathlineto{\pgfqpoint{2.756717in}{2.747241in}}%
\pgfpathlineto{\pgfqpoint{2.757490in}{2.644559in}}%
\pgfpathlineto{\pgfqpoint{2.759663in}{2.644559in}}%
\pgfpathlineto{\pgfqpoint{2.759670in}{2.767591in}}%
\pgfpathlineto{\pgfqpoint{2.759783in}{2.639486in}}%
\pgfpathlineto{\pgfqpoint{2.760447in}{2.694729in}}%
\pgfpathlineto{\pgfqpoint{2.760496in}{2.682340in}}%
\pgfpathlineto{\pgfqpoint{2.761040in}{2.731953in}}%
\pgfpathlineto{\pgfqpoint{2.761090in}{2.719679in}}%
\pgfpathlineto{\pgfqpoint{2.761611in}{2.767749in}}%
\pgfpathlineto{\pgfqpoint{2.761700in}{2.668587in}}%
\pgfpathlineto{\pgfqpoint{2.762450in}{2.644559in}}%
\pgfpathlineto{\pgfqpoint{2.764609in}{2.644559in}}%
\pgfpathlineto{\pgfqpoint{2.764616in}{2.757405in}}%
\pgfpathlineto{\pgfqpoint{2.764721in}{2.640371in}}%
\pgfpathlineto{\pgfqpoint{2.765372in}{2.692549in}}%
\pgfpathlineto{\pgfqpoint{2.765421in}{2.681874in}}%
\pgfpathlineto{\pgfqpoint{2.765965in}{2.729790in}}%
\pgfpathlineto{\pgfqpoint{2.766015in}{2.719197in}}%
\pgfpathlineto{\pgfqpoint{2.766546in}{2.766252in}}%
\pgfpathlineto{\pgfqpoint{2.766651in}{2.675958in}}%
\pgfpathlineto{\pgfqpoint{2.767424in}{2.644559in}}%
\pgfpathlineto{\pgfqpoint{2.769556in}{2.644559in}}%
\pgfpathlineto{\pgfqpoint{2.769562in}{2.754092in}}%
\pgfpathlineto{\pgfqpoint{2.769653in}{2.643260in}}%
\pgfpathlineto{\pgfqpoint{2.770296in}{2.688887in}}%
\pgfpathlineto{\pgfqpoint{2.770345in}{2.682990in}}%
\pgfpathlineto{\pgfqpoint{2.770790in}{2.719933in}}%
\pgfpathlineto{\pgfqpoint{2.770840in}{2.714081in}}%
\pgfpathlineto{\pgfqpoint{2.771483in}{2.763394in}}%
\pgfpathlineto{\pgfqpoint{2.771565in}{2.704281in}}%
\pgfpathlineto{\pgfqpoint{2.771574in}{2.734359in}}%
\pgfpathlineto{\pgfqpoint{2.772316in}{2.644559in}}%
\pgfpathlineto{\pgfqpoint{2.774501in}{2.644559in}}%
\pgfpathlineto{\pgfqpoint{2.774508in}{2.754124in}}%
\pgfpathlineto{\pgfqpoint{2.774634in}{2.634864in}}%
\pgfpathlineto{\pgfqpoint{2.775259in}{2.672890in}}%
\pgfpathlineto{\pgfqpoint{2.776099in}{2.750654in}}%
\pgfpathlineto{\pgfqpoint{2.776474in}{2.774176in}}%
\pgfpathlineto{\pgfqpoint{2.777095in}{2.644559in}}%
\pgfpathlineto{\pgfqpoint{2.779446in}{2.644559in}}%
\pgfpathlineto{\pgfqpoint{2.779453in}{2.754477in}}%
\pgfpathlineto{\pgfqpoint{2.779551in}{2.641509in}}%
\pgfpathlineto{\pgfqpoint{2.780197in}{2.690965in}}%
\pgfpathlineto{\pgfqpoint{2.780246in}{2.682181in}}%
\pgfpathlineto{\pgfqpoint{2.780691in}{2.722006in}}%
\pgfpathlineto{\pgfqpoint{2.780741in}{2.713277in}}%
\pgfpathlineto{\pgfqpoint{2.781377in}{2.765045in}}%
\pgfpathlineto{\pgfqpoint{2.781460in}{2.711887in}}%
\pgfpathlineto{\pgfqpoint{2.781470in}{2.720349in}}%
\pgfpathlineto{\pgfqpoint{2.782211in}{2.644559in}}%
\pgfpathlineto{\pgfqpoint{2.784392in}{2.644559in}}%
\pgfpathlineto{\pgfqpoint{2.784399in}{2.754540in}}%
\pgfpathlineto{\pgfqpoint{2.784503in}{2.638931in}}%
\pgfpathlineto{\pgfqpoint{2.785152in}{2.693820in}}%
\pgfpathlineto{\pgfqpoint{2.785202in}{2.680362in}}%
\pgfpathlineto{\pgfqpoint{2.785746in}{2.731059in}}%
\pgfpathlineto{\pgfqpoint{2.785795in}{2.717686in}}%
\pgfpathlineto{\pgfqpoint{2.786328in}{2.767557in}}%
\pgfpathlineto{\pgfqpoint{2.786432in}{2.676072in}}%
\pgfpathlineto{\pgfqpoint{2.787210in}{2.644559in}}%
\pgfpathlineto{\pgfqpoint{2.789339in}{2.644619in}}%
\pgfpathlineto{\pgfqpoint{2.789341in}{2.769845in}}%
\pgfpathlineto{\pgfqpoint{2.790085in}{2.684157in}}%
\pgfpathlineto{\pgfqpoint{2.791311in}{2.762511in}}%
\pgfpathlineto{\pgfqpoint{2.791311in}{2.738182in}}%
\pgfpathlineto{\pgfqpoint{2.791334in}{2.746440in}}%
\pgfpathlineto{\pgfqpoint{2.792068in}{2.644559in}}%
\pgfpathlineto{\pgfqpoint{2.794282in}{2.644559in}}%
\pgfpathlineto{\pgfqpoint{2.794289in}{2.755794in}}%
\pgfpathlineto{\pgfqpoint{2.794379in}{2.644260in}}%
\pgfpathlineto{\pgfqpoint{2.795021in}{2.687812in}}%
\pgfpathlineto{\pgfqpoint{2.795070in}{2.683766in}}%
\pgfpathlineto{\pgfqpoint{2.795417in}{2.712650in}}%
\pgfpathlineto{\pgfqpoint{2.795466in}{2.708638in}}%
\pgfpathlineto{\pgfqpoint{2.796208in}{2.762324in}}%
\pgfpathlineto{\pgfqpoint{2.796283in}{2.748081in}}%
\pgfpathlineto{\pgfqpoint{2.797052in}{2.644559in}}%
\pgfpathlineto{\pgfqpoint{2.799229in}{2.644560in}}%
\pgfpathlineto{\pgfqpoint{2.799234in}{2.767299in}}%
\pgfpathlineto{\pgfqpoint{2.799358in}{2.634846in}}%
\pgfpathlineto{\pgfqpoint{2.799980in}{2.672530in}}%
\pgfpathlineto{\pgfqpoint{2.800820in}{2.750145in}}%
\pgfpathlineto{\pgfqpoint{2.801202in}{2.774055in}}%
\pgfpathlineto{\pgfqpoint{2.801815in}{2.644559in}}%
\pgfpathlineto{\pgfqpoint{2.804173in}{2.644559in}}%
\pgfpathlineto{\pgfqpoint{2.804180in}{2.756803in}}%
\pgfpathlineto{\pgfqpoint{2.804287in}{2.641203in}}%
\pgfpathlineto{\pgfqpoint{2.804939in}{2.692056in}}%
\pgfpathlineto{\pgfqpoint{2.804988in}{2.683047in}}%
\pgfpathlineto{\pgfqpoint{2.805434in}{2.723094in}}%
\pgfpathlineto{\pgfqpoint{2.805483in}{2.714146in}}%
\pgfpathlineto{\pgfqpoint{2.806112in}{2.765657in}}%
\pgfpathlineto{\pgfqpoint{2.806189in}{2.714635in}}%
\pgfpathlineto{\pgfqpoint{2.806983in}{2.644559in}}%
\pgfpathlineto{\pgfqpoint{2.809120in}{2.644561in}}%
\pgfpathlineto{\pgfqpoint{2.809125in}{2.773335in}}%
\pgfpathlineto{\pgfqpoint{2.809244in}{2.632974in}}%
\pgfpathlineto{\pgfqpoint{2.809862in}{2.670216in}}%
\pgfpathlineto{\pgfqpoint{2.810703in}{2.751458in}}%
\pgfpathlineto{\pgfqpoint{2.811093in}{2.775893in}}%
\pgfpathlineto{\pgfqpoint{2.811677in}{2.644559in}}%
\pgfpathlineto{\pgfqpoint{2.814065in}{2.644559in}}%
\pgfpathlineto{\pgfqpoint{2.814071in}{2.754204in}}%
\pgfpathlineto{\pgfqpoint{2.814166in}{2.642277in}}%
\pgfpathlineto{\pgfqpoint{2.814810in}{2.689996in}}%
\pgfpathlineto{\pgfqpoint{2.814859in}{2.682501in}}%
\pgfpathlineto{\pgfqpoint{2.815304in}{2.721039in}}%
\pgfpathlineto{\pgfqpoint{2.815354in}{2.713595in}}%
\pgfpathlineto{\pgfqpoint{2.815993in}{2.764238in}}%
\pgfpathlineto{\pgfqpoint{2.816077in}{2.709577in}}%
\pgfpathlineto{\pgfqpoint{2.816088in}{2.726157in}}%
\pgfpathlineto{\pgfqpoint{2.816851in}{2.644559in}}%
\pgfpathlineto{\pgfqpoint{2.819011in}{2.644559in}}%
\pgfpathlineto{\pgfqpoint{2.819017in}{2.753979in}}%
\pgfpathlineto{\pgfqpoint{2.819150in}{2.629999in}}%
\pgfpathlineto{\pgfqpoint{2.819778in}{2.668572in}}%
\pgfpathlineto{\pgfqpoint{2.820619in}{2.756053in}}%
\pgfpathlineto{\pgfqpoint{2.820984in}{2.778880in}}%
\pgfpathlineto{\pgfqpoint{2.821556in}{2.644559in}}%
\pgfpathlineto{\pgfqpoint{2.823955in}{2.644559in}}%
\pgfpathlineto{\pgfqpoint{2.823962in}{2.757194in}}%
\pgfpathlineto{\pgfqpoint{2.824070in}{2.641479in}}%
\pgfpathlineto{\pgfqpoint{2.824722in}{2.691851in}}%
\pgfpathlineto{\pgfqpoint{2.824771in}{2.683392in}}%
\pgfpathlineto{\pgfqpoint{2.825216in}{2.722889in}}%
\pgfpathlineto{\pgfqpoint{2.825266in}{2.714491in}}%
\pgfpathlineto{\pgfqpoint{2.825894in}{2.765418in}}%
\pgfpathlineto{\pgfqpoint{2.825971in}{2.714590in}}%
\pgfpathlineto{\pgfqpoint{2.826759in}{2.644559in}}%
\pgfpathlineto{\pgfqpoint{2.828902in}{2.644560in}}%
\pgfpathlineto{\pgfqpoint{2.828906in}{2.772552in}}%
\pgfpathlineto{\pgfqpoint{2.829028in}{2.636946in}}%
\pgfpathlineto{\pgfqpoint{2.829647in}{2.674320in}}%
\pgfpathlineto{\pgfqpoint{2.830488in}{2.747746in}}%
\pgfpathlineto{\pgfqpoint{2.830874in}{2.771997in}}%
\pgfpathlineto{\pgfqpoint{2.831529in}{2.644559in}}%
\pgfpathlineto{\pgfqpoint{2.833846in}{2.644559in}}%
\pgfpathlineto{\pgfqpoint{2.833853in}{2.755344in}}%
\pgfpathlineto{\pgfqpoint{2.833954in}{2.641189in}}%
\pgfpathlineto{\pgfqpoint{2.834601in}{2.691476in}}%
\pgfpathlineto{\pgfqpoint{2.834650in}{2.682191in}}%
\pgfpathlineto{\pgfqpoint{2.835095in}{2.722515in}}%
\pgfpathlineto{\pgfqpoint{2.835145in}{2.713288in}}%
\pgfpathlineto{\pgfqpoint{2.835779in}{2.765429in}}%
\pgfpathlineto{\pgfqpoint{2.835860in}{2.712501in}}%
\pgfpathlineto{\pgfqpoint{2.835871in}{2.715448in}}%
\pgfpathlineto{\pgfqpoint{2.836642in}{2.644559in}}%
\pgfpathlineto{\pgfqpoint{2.838793in}{2.644593in}}%
\pgfpathlineto{\pgfqpoint{2.838798in}{2.788236in}}%
\pgfpathlineto{\pgfqpoint{2.838915in}{2.633652in}}%
\pgfpathlineto{\pgfqpoint{2.839531in}{2.670631in}}%
\pgfpathlineto{\pgfqpoint{2.840372in}{2.750453in}}%
\pgfpathlineto{\pgfqpoint{2.840765in}{2.775141in}}%
\pgfpathlineto{\pgfqpoint{2.841362in}{2.644559in}}%
\pgfpathlineto{\pgfqpoint{2.843737in}{2.644559in}}%
\pgfpathlineto{\pgfqpoint{2.843744in}{2.753922in}}%
\pgfpathlineto{\pgfqpoint{2.843870in}{2.635105in}}%
\pgfpathlineto{\pgfqpoint{2.844494in}{2.673091in}}%
\pgfpathlineto{\pgfqpoint{2.845335in}{2.750492in}}%
\pgfpathlineto{\pgfqpoint{2.845711in}{2.774075in}}%
\pgfpathlineto{\pgfqpoint{2.846330in}{2.644559in}}%
\pgfpathlineto{\pgfqpoint{2.848683in}{2.644559in}}%
\pgfpathlineto{\pgfqpoint{2.848689in}{2.755228in}}%
\pgfpathlineto{\pgfqpoint{2.848794in}{2.638847in}}%
\pgfpathlineto{\pgfqpoint{2.849445in}{2.694125in}}%
\pgfpathlineto{\pgfqpoint{2.849494in}{2.680375in}}%
\pgfpathlineto{\pgfqpoint{2.850038in}{2.731363in}}%
\pgfpathlineto{\pgfqpoint{2.850088in}{2.717701in}}%
\pgfpathlineto{\pgfqpoint{2.850619in}{2.767820in}}%
\pgfpathlineto{\pgfqpoint{2.850723in}{2.676664in}}%
\pgfpathlineto{\pgfqpoint{2.851488in}{2.644559in}}%
\pgfpathlineto{\pgfqpoint{2.853629in}{2.644559in}}%
\pgfpathlineto{\pgfqpoint{2.853635in}{2.754055in}}%
\pgfpathlineto{\pgfqpoint{2.853764in}{2.628300in}}%
\pgfpathlineto{\pgfqpoint{2.854391in}{2.666635in}}%
\pgfpathlineto{\pgfqpoint{2.855232in}{2.757557in}}%
\pgfpathlineto{\pgfqpoint{2.855602in}{2.780736in}}%
\pgfpathlineto{\pgfqpoint{2.856156in}{2.644559in}}%
\pgfpathlineto{\pgfqpoint{2.858572in}{2.644559in}}%
\pgfpathlineto{\pgfqpoint{2.858579in}{2.763025in}}%
\pgfpathlineto{\pgfqpoint{2.858699in}{2.630734in}}%
\pgfpathlineto{\pgfqpoint{2.859317in}{2.668009in}}%
\pgfpathlineto{\pgfqpoint{2.860157in}{2.753628in}}%
\pgfpathlineto{\pgfqpoint{2.860547in}{2.778062in}}%
\pgfpathlineto{\pgfqpoint{2.861098in}{2.644559in}}%
\pgfpathlineto{\pgfqpoint{2.863521in}{2.644573in}}%
\pgfpathlineto{\pgfqpoint{2.863523in}{2.759651in}}%
\pgfpathlineto{\pgfqpoint{2.864269in}{2.683581in}}%
\pgfpathlineto{\pgfqpoint{2.865493in}{2.763155in}}%
\pgfpathlineto{\pgfqpoint{2.865493in}{2.737977in}}%
\pgfpathlineto{\pgfqpoint{2.865520in}{2.748089in}}%
\pgfpathlineto{\pgfqpoint{2.866249in}{2.644559in}}%
\pgfpathlineto{\pgfqpoint{2.868464in}{2.644559in}}%
\pgfpathlineto{\pgfqpoint{2.868471in}{2.755532in}}%
\pgfpathlineto{\pgfqpoint{2.869191in}{2.685811in}}%
\pgfpathlineto{\pgfqpoint{2.869241in}{2.684295in}}%
\pgfpathlineto{\pgfqpoint{2.869389in}{2.698232in}}%
\pgfpathlineto{\pgfqpoint{2.869439in}{2.696729in}}%
\pgfpathlineto{\pgfqpoint{2.870438in}{2.764085in}}%
\pgfpathlineto{\pgfqpoint{2.870438in}{2.738287in}}%
\pgfpathlineto{\pgfqpoint{2.870458in}{2.744255in}}%
\pgfpathlineto{\pgfqpoint{2.871177in}{2.644559in}}%
\pgfpathlineto{\pgfqpoint{2.873411in}{2.644559in}}%
\pgfpathlineto{\pgfqpoint{2.873415in}{2.763989in}}%
\pgfpathlineto{\pgfqpoint{2.873545in}{2.631231in}}%
\pgfpathlineto{\pgfqpoint{2.874169in}{2.669335in}}%
\pgfpathlineto{\pgfqpoint{2.875010in}{2.754260in}}%
\pgfpathlineto{\pgfqpoint{2.875384in}{2.777666in}}%
\pgfpathlineto{\pgfqpoint{2.875955in}{2.644559in}}%
\pgfpathlineto{\pgfqpoint{2.878355in}{2.644559in}}%
\pgfpathlineto{\pgfqpoint{2.878362in}{2.753834in}}%
\pgfpathlineto{\pgfqpoint{2.878481in}{2.632504in}}%
\pgfpathlineto{\pgfqpoint{2.879100in}{2.669834in}}%
\pgfpathlineto{\pgfqpoint{2.879941in}{2.752307in}}%
\pgfpathlineto{\pgfqpoint{2.880329in}{2.776641in}}%
\pgfpathlineto{\pgfqpoint{2.880910in}{2.644559in}}%
\pgfpathlineto{\pgfqpoint{2.883302in}{2.644560in}}%
\pgfpathlineto{\pgfqpoint{2.883306in}{2.769237in}}%
\pgfpathlineto{\pgfqpoint{2.883431in}{2.634791in}}%
\pgfpathlineto{\pgfqpoint{2.884052in}{2.672465in}}%
\pgfpathlineto{\pgfqpoint{2.884893in}{2.750218in}}%
\pgfpathlineto{\pgfqpoint{2.885274in}{2.774130in}}%
\pgfpathlineto{\pgfqpoint{2.885895in}{2.644559in}}%
\pgfpathlineto{\pgfqpoint{2.888246in}{2.644559in}}%
\pgfpathlineto{\pgfqpoint{2.888253in}{2.754753in}}%
\pgfpathlineto{\pgfqpoint{2.888356in}{2.640220in}}%
\pgfpathlineto{\pgfqpoint{2.889005in}{2.692480in}}%
\pgfpathlineto{\pgfqpoint{2.889054in}{2.681535in}}%
\pgfpathlineto{\pgfqpoint{2.889598in}{2.729724in}}%
\pgfpathlineto{\pgfqpoint{2.889648in}{2.718854in}}%
\pgfpathlineto{\pgfqpoint{2.890181in}{2.766299in}}%
\pgfpathlineto{\pgfqpoint{2.890287in}{2.674107in}}%
\pgfpathlineto{\pgfqpoint{2.891045in}{2.644559in}}%
\pgfpathlineto{\pgfqpoint{2.893193in}{2.644565in}}%
\pgfpathlineto{\pgfqpoint{2.893198in}{2.780905in}}%
\pgfpathlineto{\pgfqpoint{2.893306in}{2.635503in}}%
\pgfpathlineto{\pgfqpoint{2.893962in}{2.697938in}}%
\pgfpathlineto{\pgfqpoint{2.894011in}{2.677525in}}%
\pgfpathlineto{\pgfqpoint{2.894555in}{2.735167in}}%
\pgfpathlineto{\pgfqpoint{2.894604in}{2.714859in}}%
\pgfpathlineto{\pgfqpoint{2.895132in}{2.771371in}}%
\pgfpathlineto{\pgfqpoint{2.895242in}{2.644560in}}%
\pgfpathlineto{\pgfqpoint{2.898138in}{2.644559in}}%
\pgfpathlineto{\pgfqpoint{2.898142in}{2.757033in}}%
\pgfpathlineto{\pgfqpoint{2.898273in}{2.626880in}}%
\pgfpathlineto{\pgfqpoint{2.898900in}{2.665165in}}%
\pgfpathlineto{\pgfqpoint{2.899740in}{2.758755in}}%
\pgfpathlineto{\pgfqpoint{2.900111in}{2.781957in}}%
\pgfpathlineto{\pgfqpoint{2.900662in}{2.644559in}}%
\pgfpathlineto{\pgfqpoint{2.903084in}{2.644600in}}%
\pgfpathlineto{\pgfqpoint{2.903087in}{2.767732in}}%
\pgfpathlineto{\pgfqpoint{2.903839in}{2.684248in}}%
\pgfpathlineto{\pgfqpoint{2.905056in}{2.762923in}}%
\pgfpathlineto{\pgfqpoint{2.905057in}{2.738200in}}%
\pgfpathlineto{\pgfqpoint{2.905078in}{2.745522in}}%
\pgfpathlineto{\pgfqpoint{2.905788in}{2.644559in}}%
\pgfpathlineto{\pgfqpoint{2.908029in}{2.644559in}}%
\pgfpathlineto{\pgfqpoint{2.908035in}{2.754239in}}%
\pgfpathlineto{\pgfqpoint{2.908165in}{2.626349in}}%
\pgfpathlineto{\pgfqpoint{2.908791in}{2.664672in}}%
\pgfpathlineto{\pgfqpoint{2.909632in}{2.759394in}}%
\pgfpathlineto{\pgfqpoint{2.910002in}{2.782546in}}%
\pgfpathlineto{\pgfqpoint{2.910533in}{2.644559in}}%
\pgfpathlineto{\pgfqpoint{2.912975in}{2.644579in}}%
\pgfpathlineto{\pgfqpoint{2.912978in}{2.762461in}}%
\pgfpathlineto{\pgfqpoint{2.913740in}{2.686413in}}%
\pgfpathlineto{\pgfqpoint{2.914947in}{2.761342in}}%
\pgfpathlineto{\pgfqpoint{2.914989in}{2.725140in}}%
\pgfpathlineto{\pgfqpoint{2.915738in}{2.644559in}}%
\pgfpathlineto{\pgfqpoint{2.917918in}{2.644559in}}%
\pgfpathlineto{\pgfqpoint{2.917925in}{2.760877in}}%
\pgfpathlineto{\pgfqpoint{2.918016in}{2.618990in}}%
\pgfpathlineto{\pgfqpoint{2.918657in}{2.712766in}}%
\pgfpathlineto{\pgfqpoint{2.918706in}{2.658549in}}%
\pgfpathlineto{\pgfqpoint{2.919349in}{2.756156in}}%
\pgfpathlineto{\pgfqpoint{2.919399in}{2.702150in}}%
\pgfpathlineto{\pgfqpoint{2.919844in}{2.787148in}}%
\pgfpathlineto{\pgfqpoint{2.919984in}{2.644559in}}%
\pgfpathlineto{\pgfqpoint{2.920099in}{2.644559in}}%
\pgfpathlineto{\pgfqpoint{2.922865in}{2.644559in}}%
\pgfpathlineto{\pgfqpoint{2.922872in}{2.754307in}}%
\pgfpathlineto{\pgfqpoint{2.922999in}{2.631329in}}%
\pgfpathlineto{\pgfqpoint{2.923624in}{2.669440in}}%
\pgfpathlineto{\pgfqpoint{2.924465in}{2.754195in}}%
\pgfpathlineto{\pgfqpoint{2.924838in}{2.777609in}}%
\pgfpathlineto{\pgfqpoint{2.925425in}{2.644559in}}%
\pgfpathlineto{\pgfqpoint{2.927810in}{2.644559in}}%
\pgfpathlineto{\pgfqpoint{2.927817in}{2.754360in}}%
\pgfpathlineto{\pgfqpoint{2.928580in}{2.688147in}}%
\pgfpathlineto{\pgfqpoint{2.928630in}{2.687453in}}%
\pgfpathlineto{\pgfqpoint{2.928679in}{2.694361in}}%
\pgfpathlineto{\pgfqpoint{2.928728in}{2.693666in}}%
\pgfpathlineto{\pgfqpoint{2.929751in}{2.761677in}}%
\pgfpathlineto{\pgfqpoint{2.929784in}{2.738068in}}%
\pgfpathlineto{\pgfqpoint{2.929811in}{2.748165in}}%
\pgfpathlineto{\pgfqpoint{2.930557in}{2.644559in}}%
\pgfpathlineto{\pgfqpoint{2.932755in}{2.644559in}}%
\pgfpathlineto{\pgfqpoint{2.932761in}{2.757008in}}%
\pgfpathlineto{\pgfqpoint{2.932892in}{2.629388in}}%
\pgfpathlineto{\pgfqpoint{2.933519in}{2.667736in}}%
\pgfpathlineto{\pgfqpoint{2.934360in}{2.756439in}}%
\pgfpathlineto{\pgfqpoint{2.934729in}{2.779571in}}%
\pgfpathlineto{\pgfqpoint{2.935286in}{2.644559in}}%
\pgfpathlineto{\pgfqpoint{2.937701in}{2.644559in}}%
\pgfpathlineto{\pgfqpoint{2.937707in}{2.755380in}}%
\pgfpathlineto{\pgfqpoint{2.937816in}{2.639643in}}%
\pgfpathlineto{\pgfqpoint{2.938469in}{2.693604in}}%
\pgfpathlineto{\pgfqpoint{2.938518in}{2.681649in}}%
\pgfpathlineto{\pgfqpoint{2.939062in}{2.730846in}}%
\pgfpathlineto{\pgfqpoint{2.939112in}{2.718971in}}%
\pgfpathlineto{\pgfqpoint{2.939640in}{2.767120in}}%
\pgfpathlineto{\pgfqpoint{2.939742in}{2.675240in}}%
\pgfpathlineto{\pgfqpoint{2.940507in}{2.644559in}}%
\pgfpathlineto{\pgfqpoint{2.942647in}{2.644559in}}%
\pgfpathlineto{\pgfqpoint{2.942653in}{2.754175in}}%
\pgfpathlineto{\pgfqpoint{2.942783in}{2.627712in}}%
\pgfpathlineto{\pgfqpoint{2.943410in}{2.666059in}}%
\pgfpathlineto{\pgfqpoint{2.944250in}{2.758115in}}%
\pgfpathlineto{\pgfqpoint{2.944620in}{2.781268in}}%
\pgfpathlineto{\pgfqpoint{2.945168in}{2.644559in}}%
\pgfpathlineto{\pgfqpoint{2.947593in}{2.644561in}}%
\pgfpathlineto{\pgfqpoint{2.947597in}{2.775740in}}%
\pgfpathlineto{\pgfqpoint{2.947715in}{2.633360in}}%
\pgfpathlineto{\pgfqpoint{2.948332in}{2.670447in}}%
\pgfpathlineto{\pgfqpoint{2.949173in}{2.750925in}}%
\pgfpathlineto{\pgfqpoint{2.949565in}{2.775522in}}%
\pgfpathlineto{\pgfqpoint{2.950157in}{2.644559in}}%
\pgfpathlineto{\pgfqpoint{2.952536in}{2.644559in}}%
\pgfpathlineto{\pgfqpoint{2.952543in}{2.762307in}}%
\pgfpathlineto{\pgfqpoint{2.952653in}{2.618704in}}%
\pgfpathlineto{\pgfqpoint{2.953310in}{2.714815in}}%
\pgfpathlineto{\pgfqpoint{2.953360in}{2.661070in}}%
\pgfpathlineto{\pgfqpoint{2.954003in}{2.758194in}}%
\pgfpathlineto{\pgfqpoint{2.954052in}{2.704682in}}%
\pgfpathlineto{\pgfqpoint{2.954479in}{2.788046in}}%
\pgfpathlineto{\pgfqpoint{2.954593in}{2.644559in}}%
\pgfpathlineto{\pgfqpoint{2.954754in}{2.644559in}}%
\pgfpathlineto{\pgfqpoint{2.957484in}{2.644561in}}%
\pgfpathlineto{\pgfqpoint{2.957488in}{2.775568in}}%
\pgfpathlineto{\pgfqpoint{2.957608in}{2.637633in}}%
\pgfpathlineto{\pgfqpoint{2.958227in}{2.674877in}}%
\pgfpathlineto{\pgfqpoint{2.959166in}{2.753157in}}%
\pgfpathlineto{\pgfqpoint{2.959456in}{2.771338in}}%
\pgfpathlineto{\pgfqpoint{2.960232in}{2.644559in}}%
\pgfpathlineto{\pgfqpoint{2.962428in}{2.644559in}}%
\pgfpathlineto{\pgfqpoint{2.962434in}{2.756737in}}%
\pgfpathlineto{\pgfqpoint{2.962568in}{2.633598in}}%
\pgfpathlineto{\pgfqpoint{2.963198in}{2.672222in}}%
\pgfpathlineto{\pgfqpoint{2.964039in}{2.752831in}}%
\pgfpathlineto{\pgfqpoint{2.964402in}{2.775580in}}%
\pgfpathlineto{\pgfqpoint{2.965023in}{2.644559in}}%
\pgfpathlineto{\pgfqpoint{2.967375in}{2.644562in}}%
\pgfpathlineto{\pgfqpoint{2.967379in}{2.777488in}}%
\pgfpathlineto{\pgfqpoint{2.967497in}{2.633486in}}%
\pgfpathlineto{\pgfqpoint{2.968113in}{2.670509in}}%
\pgfpathlineto{\pgfqpoint{2.968954in}{2.750761in}}%
\pgfpathlineto{\pgfqpoint{2.969347in}{2.775419in}}%
\pgfpathlineto{\pgfqpoint{2.969930in}{2.644559in}}%
\pgfpathlineto{\pgfqpoint{2.972318in}{2.644559in}}%
\pgfpathlineto{\pgfqpoint{2.972324in}{2.764124in}}%
\pgfpathlineto{\pgfqpoint{2.972445in}{2.629321in}}%
\pgfpathlineto{\pgfqpoint{2.973064in}{2.666718in}}%
\pgfpathlineto{\pgfqpoint{2.973905in}{2.755082in}}%
\pgfpathlineto{\pgfqpoint{2.974293in}{2.779390in}}%
\pgfpathlineto{\pgfqpoint{2.974828in}{2.644559in}}%
\pgfpathlineto{\pgfqpoint{2.977266in}{2.644563in}}%
\pgfpathlineto{\pgfqpoint{2.977270in}{2.778609in}}%
\pgfpathlineto{\pgfqpoint{2.977385in}{2.633922in}}%
\pgfpathlineto{\pgfqpoint{2.978000in}{2.670714in}}%
\pgfpathlineto{\pgfqpoint{2.978841in}{2.750085in}}%
\pgfpathlineto{\pgfqpoint{2.979213in}{2.773385in}}%
\pgfpathlineto{\pgfqpoint{2.979832in}{2.644559in}}%
\pgfpathlineto{\pgfqpoint{2.982209in}{2.644559in}}%
\pgfpathlineto{\pgfqpoint{2.982216in}{2.760081in}}%
\pgfpathlineto{\pgfqpoint{2.982345in}{2.644539in}}%
\pgfpathlineto{\pgfqpoint{2.982970in}{2.682451in}}%
\pgfpathlineto{\pgfqpoint{2.984184in}{2.764596in}}%
\pgfpathlineto{\pgfqpoint{2.984956in}{2.644559in}}%
\pgfpathlineto{\pgfqpoint{2.987157in}{2.644574in}}%
\pgfpathlineto{\pgfqpoint{2.987159in}{2.760086in}}%
\pgfpathlineto{\pgfqpoint{2.987896in}{2.685741in}}%
\pgfpathlineto{\pgfqpoint{2.987946in}{2.686080in}}%
\pgfpathlineto{\pgfqpoint{2.989129in}{2.760414in}}%
\pgfpathlineto{\pgfqpoint{2.989129in}{2.738047in}}%
\pgfpathlineto{\pgfqpoint{2.989156in}{2.748127in}}%
\pgfpathlineto{\pgfqpoint{2.989874in}{2.644559in}}%
\pgfpathlineto{\pgfqpoint{2.992101in}{2.644559in}}%
\pgfpathlineto{\pgfqpoint{2.992107in}{2.755167in}}%
\pgfpathlineto{\pgfqpoint{2.992213in}{2.637833in}}%
\pgfpathlineto{\pgfqpoint{2.992865in}{2.695098in}}%
\pgfpathlineto{\pgfqpoint{2.992914in}{2.679464in}}%
\pgfpathlineto{\pgfqpoint{2.993458in}{2.732333in}}%
\pgfpathlineto{\pgfqpoint{2.993508in}{2.716793in}}%
\pgfpathlineto{\pgfqpoint{2.994038in}{2.768729in}}%
\pgfpathlineto{\pgfqpoint{2.994138in}{2.668303in}}%
\pgfpathlineto{\pgfqpoint{2.994879in}{2.644559in}}%
\pgfpathlineto{\pgfqpoint{2.997047in}{2.644559in}}%
\pgfpathlineto{\pgfqpoint{2.997052in}{2.758760in}}%
\pgfpathlineto{\pgfqpoint{2.997186in}{2.629876in}}%
\pgfpathlineto{\pgfqpoint{2.997814in}{2.668411in}}%
\pgfpathlineto{\pgfqpoint{2.998655in}{2.756199in}}%
\pgfpathlineto{\pgfqpoint{2.999020in}{2.779059in}}%
\pgfpathlineto{\pgfqpoint{2.999568in}{2.644559in}}%
\pgfpathlineto{\pgfqpoint{3.001992in}{2.644559in}}%
\pgfpathlineto{\pgfqpoint{3.001998in}{2.755387in}}%
\pgfpathlineto{\pgfqpoint{3.002108in}{2.639357in}}%
\pgfpathlineto{\pgfqpoint{3.002761in}{2.693856in}}%
\pgfpathlineto{\pgfqpoint{3.002810in}{2.681443in}}%
\pgfpathlineto{\pgfqpoint{3.003354in}{2.731097in}}%
\pgfpathlineto{\pgfqpoint{3.003404in}{2.718766in}}%
\pgfpathlineto{\pgfqpoint{3.003932in}{2.767331in}}%
\pgfpathlineto{\pgfqpoint{3.004032in}{2.675574in}}%
\pgfpathlineto{\pgfqpoint{3.004793in}{2.644559in}}%
\pgfpathlineto{\pgfqpoint{3.006938in}{2.644559in}}%
\pgfpathlineto{\pgfqpoint{3.006942in}{2.764178in}}%
\pgfpathlineto{\pgfqpoint{3.007076in}{2.630276in}}%
\pgfpathlineto{\pgfqpoint{3.007703in}{2.668741in}}%
\pgfpathlineto{\pgfqpoint{3.008544in}{2.755807in}}%
\pgfpathlineto{\pgfqpoint{3.008911in}{2.778793in}}%
\pgfpathlineto{\pgfqpoint{3.009488in}{2.644559in}}%
\pgfpathlineto{\pgfqpoint{3.011884in}{2.644559in}}%
\pgfpathlineto{\pgfqpoint{3.011888in}{2.762036in}}%
\pgfpathlineto{\pgfqpoint{3.012022in}{2.629674in}}%
\pgfpathlineto{\pgfqpoint{3.012650in}{2.668201in}}%
\pgfpathlineto{\pgfqpoint{3.013491in}{2.756463in}}%
\pgfpathlineto{\pgfqpoint{3.013856in}{2.779375in}}%
\pgfpathlineto{\pgfqpoint{3.014434in}{2.644559in}}%
\pgfpathlineto{\pgfqpoint{3.016828in}{2.644559in}}%
\pgfpathlineto{\pgfqpoint{3.016835in}{2.754557in}}%
\pgfpathlineto{\pgfqpoint{3.017572in}{2.685801in}}%
\pgfpathlineto{\pgfqpoint{3.017720in}{2.692451in}}%
\pgfpathlineto{\pgfqpoint{3.018802in}{2.760379in}}%
\pgfpathlineto{\pgfqpoint{3.018802in}{2.737980in}}%
\pgfpathlineto{\pgfqpoint{3.018829in}{2.748046in}}%
\pgfpathlineto{\pgfqpoint{3.019554in}{2.644559in}}%
\pgfpathlineto{\pgfqpoint{3.021775in}{2.644559in}}%
\pgfpathlineto{\pgfqpoint{3.021779in}{2.765940in}}%
\pgfpathlineto{\pgfqpoint{3.021905in}{2.633636in}}%
\pgfpathlineto{\pgfqpoint{3.022528in}{2.671474in}}%
\pgfpathlineto{\pgfqpoint{3.023369in}{2.751537in}}%
\pgfpathlineto{\pgfqpoint{3.023747in}{2.775263in}}%
\pgfpathlineto{\pgfqpoint{3.024349in}{2.644559in}}%
\pgfpathlineto{\pgfqpoint{3.026718in}{2.644559in}}%
\pgfpathlineto{\pgfqpoint{3.026726in}{2.756030in}}%
\pgfpathlineto{\pgfqpoint{3.026825in}{2.641545in}}%
\pgfpathlineto{\pgfqpoint{3.027472in}{2.690952in}}%
\pgfpathlineto{\pgfqpoint{3.027521in}{2.682383in}}%
\pgfpathlineto{\pgfqpoint{3.027966in}{2.721993in}}%
\pgfpathlineto{\pgfqpoint{3.028016in}{2.713479in}}%
\pgfpathlineto{\pgfqpoint{3.028651in}{2.764964in}}%
\pgfpathlineto{\pgfqpoint{3.028733in}{2.712101in}}%
\pgfpathlineto{\pgfqpoint{3.028743in}{2.718829in}}%
\pgfpathlineto{\pgfqpoint{3.029485in}{2.644559in}}%
\pgfpathlineto{\pgfqpoint{3.031665in}{2.644559in}}%
\pgfpathlineto{\pgfqpoint{3.031672in}{2.754449in}}%
\pgfpathlineto{\pgfqpoint{3.031801in}{2.627212in}}%
\pgfpathlineto{\pgfqpoint{3.032428in}{2.665501in}}%
\pgfpathlineto{\pgfqpoint{3.033268in}{2.758364in}}%
\pgfpathlineto{\pgfqpoint{3.033638in}{2.781530in}}%
\pgfpathlineto{\pgfqpoint{3.034190in}{2.644559in}}%
\pgfpathlineto{\pgfqpoint{3.036608in}{2.644559in}}%
\pgfpathlineto{\pgfqpoint{3.036615in}{2.766379in}}%
\pgfpathlineto{\pgfqpoint{3.036736in}{2.633160in}}%
\pgfpathlineto{\pgfqpoint{3.037404in}{2.701819in}}%
\pgfpathlineto{\pgfqpoint{3.037453in}{2.676735in}}%
\pgfpathlineto{\pgfqpoint{3.037997in}{2.739020in}}%
\pgfpathlineto{\pgfqpoint{3.038047in}{2.714098in}}%
\pgfpathlineto{\pgfqpoint{3.038584in}{2.775762in}}%
\pgfpathlineto{\pgfqpoint{3.038673in}{2.644559in}}%
\pgfpathlineto{\pgfqpoint{3.038695in}{2.644559in}}%
\pgfpathlineto{\pgfqpoint{3.041557in}{2.644559in}}%
\pgfpathlineto{\pgfqpoint{3.041561in}{2.765256in}}%
\pgfpathlineto{\pgfqpoint{3.041689in}{2.632076in}}%
\pgfpathlineto{\pgfqpoint{3.042313in}{2.670079in}}%
\pgfpathlineto{\pgfqpoint{3.043154in}{2.753298in}}%
\pgfpathlineto{\pgfqpoint{3.043529in}{2.776825in}}%
\pgfpathlineto{\pgfqpoint{3.044111in}{2.644559in}}%
\pgfpathlineto{\pgfqpoint{3.046499in}{2.644559in}}%
\pgfpathlineto{\pgfqpoint{3.046506in}{2.764200in}}%
\pgfpathlineto{\pgfqpoint{3.046628in}{2.628623in}}%
\pgfpathlineto{\pgfqpoint{3.047247in}{2.666089in}}%
\pgfpathlineto{\pgfqpoint{3.048088in}{2.755928in}}%
\pgfpathlineto{\pgfqpoint{3.048474in}{2.780158in}}%
\pgfpathlineto{\pgfqpoint{3.049008in}{2.644559in}}%
\pgfpathlineto{\pgfqpoint{3.051448in}{2.644560in}}%
\pgfpathlineto{\pgfqpoint{3.051452in}{2.767402in}}%
\pgfpathlineto{\pgfqpoint{3.051585in}{2.630612in}}%
\pgfpathlineto{\pgfqpoint{3.052212in}{2.669004in}}%
\pgfpathlineto{\pgfqpoint{3.053053in}{2.755383in}}%
\pgfpathlineto{\pgfqpoint{3.053420in}{2.778396in}}%
\pgfpathlineto{\pgfqpoint{3.053988in}{2.644559in}}%
\pgfpathlineto{\pgfqpoint{3.056392in}{2.644559in}}%
\pgfpathlineto{\pgfqpoint{3.056399in}{2.754124in}}%
\pgfpathlineto{\pgfqpoint{3.056528in}{2.627414in}}%
\pgfpathlineto{\pgfqpoint{3.057155in}{2.665755in}}%
\pgfpathlineto{\pgfqpoint{3.057996in}{2.758385in}}%
\pgfpathlineto{\pgfqpoint{3.058365in}{2.781542in}}%
\pgfpathlineto{\pgfqpoint{3.058912in}{2.644559in}}%
\pgfpathlineto{\pgfqpoint{3.061339in}{2.644560in}}%
\pgfpathlineto{\pgfqpoint{3.061343in}{2.771393in}}%
\pgfpathlineto{\pgfqpoint{3.061465in}{2.636358in}}%
\pgfpathlineto{\pgfqpoint{3.062085in}{2.673820in}}%
\pgfpathlineto{\pgfqpoint{3.062926in}{2.748423in}}%
\pgfpathlineto{\pgfqpoint{3.063311in}{2.772578in}}%
\pgfpathlineto{\pgfqpoint{3.063949in}{2.644559in}}%
\pgfpathlineto{\pgfqpoint{3.066282in}{2.644559in}}%
\pgfpathlineto{\pgfqpoint{3.066289in}{2.756479in}}%
\pgfpathlineto{\pgfqpoint{3.066394in}{2.639250in}}%
\pgfpathlineto{\pgfqpoint{3.067044in}{2.693690in}}%
\pgfpathlineto{\pgfqpoint{3.067094in}{2.680763in}}%
\pgfpathlineto{\pgfqpoint{3.067638in}{2.730929in}}%
\pgfpathlineto{\pgfqpoint{3.067687in}{2.718087in}}%
\pgfpathlineto{\pgfqpoint{3.068219in}{2.767392in}}%
\pgfpathlineto{\pgfqpoint{3.068323in}{2.676535in}}%
\pgfpathlineto{\pgfqpoint{3.069067in}{2.644559in}}%
\pgfpathlineto{\pgfqpoint{3.071229in}{2.644560in}}%
\pgfpathlineto{\pgfqpoint{3.071234in}{2.772650in}}%
\pgfpathlineto{\pgfqpoint{3.071357in}{2.635928in}}%
\pgfpathlineto{\pgfqpoint{3.071977in}{2.673434in}}%
\pgfpathlineto{\pgfqpoint{3.072818in}{2.748927in}}%
\pgfpathlineto{\pgfqpoint{3.073202in}{2.773021in}}%
\pgfpathlineto{\pgfqpoint{3.073843in}{2.644559in}}%
\pgfpathlineto{\pgfqpoint{3.076173in}{2.644559in}}%
\pgfpathlineto{\pgfqpoint{3.076179in}{2.758490in}}%
\pgfpathlineto{\pgfqpoint{3.076297in}{2.631072in}}%
\pgfpathlineto{\pgfqpoint{3.076914in}{2.668185in}}%
\pgfpathlineto{\pgfqpoint{3.077755in}{2.753066in}}%
\pgfpathlineto{\pgfqpoint{3.078147in}{2.777661in}}%
\pgfpathlineto{\pgfqpoint{3.078701in}{2.644559in}}%
\pgfpathlineto{\pgfqpoint{3.081119in}{2.644559in}}%
\pgfpathlineto{\pgfqpoint{3.081126in}{2.755543in}}%
\pgfpathlineto{\pgfqpoint{3.081230in}{2.639032in}}%
\pgfpathlineto{\pgfqpoint{3.081880in}{2.693769in}}%
\pgfpathlineto{\pgfqpoint{3.081930in}{2.680520in}}%
\pgfpathlineto{\pgfqpoint{3.082474in}{2.731009in}}%
\pgfpathlineto{\pgfqpoint{3.082523in}{2.717844in}}%
\pgfpathlineto{\pgfqpoint{3.083055in}{2.767484in}}%
\pgfpathlineto{\pgfqpoint{3.083160in}{2.676304in}}%
\pgfpathlineto{\pgfqpoint{3.083938in}{2.644559in}}%
\pgfpathlineto{\pgfqpoint{3.086064in}{2.644559in}}%
\pgfpathlineto{\pgfqpoint{3.086071in}{2.754960in}}%
\pgfpathlineto{\pgfqpoint{3.086176in}{2.638298in}}%
\pgfpathlineto{\pgfqpoint{3.086827in}{2.694649in}}%
\pgfpathlineto{\pgfqpoint{3.086877in}{2.679872in}}%
\pgfpathlineto{\pgfqpoint{3.087421in}{2.731885in}}%
\pgfpathlineto{\pgfqpoint{3.087470in}{2.717199in}}%
\pgfpathlineto{\pgfqpoint{3.088001in}{2.768314in}}%
\pgfpathlineto{\pgfqpoint{3.088101in}{2.668483in}}%
\pgfpathlineto{\pgfqpoint{3.088840in}{2.644559in}}%
\pgfpathlineto{\pgfqpoint{3.091010in}{2.644559in}}%
\pgfpathlineto{\pgfqpoint{3.091017in}{2.755788in}}%
\pgfpathlineto{\pgfqpoint{3.091123in}{2.636650in}}%
\pgfpathlineto{\pgfqpoint{3.091776in}{2.696395in}}%
\pgfpathlineto{\pgfqpoint{3.091826in}{2.678412in}}%
\pgfpathlineto{\pgfqpoint{3.092370in}{2.733625in}}%
\pgfpathlineto{\pgfqpoint{3.092419in}{2.715745in}}%
\pgfpathlineto{\pgfqpoint{3.092949in}{2.769941in}}%
\pgfpathlineto{\pgfqpoint{3.093046in}{2.668218in}}%
\pgfpathlineto{\pgfqpoint{3.093814in}{2.644559in}}%
\pgfpathlineto{\pgfqpoint{3.095956in}{2.644559in}}%
\pgfpathlineto{\pgfqpoint{3.095962in}{2.754215in}}%
\pgfpathlineto{\pgfqpoint{3.096092in}{2.626497in}}%
\pgfpathlineto{\pgfqpoint{3.096718in}{2.664808in}}%
\pgfpathlineto{\pgfqpoint{3.097559in}{2.759092in}}%
\pgfpathlineto{\pgfqpoint{3.097929in}{2.782256in}}%
\pgfpathlineto{\pgfqpoint{3.098476in}{2.644559in}}%
\pgfpathlineto{\pgfqpoint{3.100901in}{2.644559in}}%
\pgfpathlineto{\pgfqpoint{3.100907in}{2.754992in}}%
\pgfpathlineto{\pgfqpoint{3.101010in}{2.640674in}}%
\pgfpathlineto{\pgfqpoint{3.101658in}{2.692089in}}%
\pgfpathlineto{\pgfqpoint{3.101707in}{2.681864in}}%
\pgfpathlineto{\pgfqpoint{3.102152in}{2.723127in}}%
\pgfpathlineto{\pgfqpoint{3.102202in}{2.712963in}}%
\pgfpathlineto{\pgfqpoint{3.102835in}{2.765964in}}%
\pgfpathlineto{\pgfqpoint{3.102916in}{2.714530in}}%
\pgfpathlineto{\pgfqpoint{3.103703in}{2.644559in}}%
\pgfpathlineto{\pgfqpoint{3.105847in}{2.644559in}}%
\pgfpathlineto{\pgfqpoint{3.105852in}{2.760669in}}%
\pgfpathlineto{\pgfqpoint{3.105985in}{2.630584in}}%
\pgfpathlineto{\pgfqpoint{3.106613in}{2.669043in}}%
\pgfpathlineto{\pgfqpoint{3.107454in}{2.755407in}}%
\pgfpathlineto{\pgfqpoint{3.107820in}{2.778362in}}%
\pgfpathlineto{\pgfqpoint{3.108396in}{2.644559in}}%
\pgfpathlineto{\pgfqpoint{3.110792in}{2.644559in}}%
\pgfpathlineto{\pgfqpoint{3.110798in}{2.754870in}}%
\pgfpathlineto{\pgfqpoint{3.110903in}{2.638991in}}%
\pgfpathlineto{\pgfqpoint{3.111553in}{2.693935in}}%
\pgfpathlineto{\pgfqpoint{3.111603in}{2.680486in}}%
\pgfpathlineto{\pgfqpoint{3.112147in}{2.731174in}}%
\pgfpathlineto{\pgfqpoint{3.112196in}{2.717811in}}%
\pgfpathlineto{\pgfqpoint{3.112728in}{2.767650in}}%
\pgfpathlineto{\pgfqpoint{3.112832in}{2.676292in}}%
\pgfpathlineto{\pgfqpoint{3.113593in}{2.644559in}}%
\pgfpathlineto{\pgfqpoint{3.115738in}{2.644560in}}%
\pgfpathlineto{\pgfqpoint{3.115743in}{2.767215in}}%
\pgfpathlineto{\pgfqpoint{3.115869in}{2.633542in}}%
\pgfpathlineto{\pgfqpoint{3.116492in}{2.671380in}}%
\pgfpathlineto{\pgfqpoint{3.117332in}{2.751648in}}%
\pgfpathlineto{\pgfqpoint{3.117711in}{2.775368in}}%
\pgfpathlineto{\pgfqpoint{3.118307in}{2.644559in}}%
\pgfpathlineto{\pgfqpoint{3.120683in}{2.644559in}}%
\pgfpathlineto{\pgfqpoint{3.120690in}{2.754587in}}%
\pgfpathlineto{\pgfqpoint{3.120794in}{2.639064in}}%
\pgfpathlineto{\pgfqpoint{3.121443in}{2.693736in}}%
\pgfpathlineto{\pgfqpoint{3.121493in}{2.680513in}}%
\pgfpathlineto{\pgfqpoint{3.122037in}{2.730975in}}%
\pgfpathlineto{\pgfqpoint{3.122086in}{2.717837in}}%
\pgfpathlineto{\pgfqpoint{3.122619in}{2.767471in}}%
\pgfpathlineto{\pgfqpoint{3.122723in}{2.676007in}}%
\pgfpathlineto{\pgfqpoint{3.123495in}{2.644559in}}%
\pgfpathlineto{\pgfqpoint{3.125629in}{2.644560in}}%
\pgfpathlineto{\pgfqpoint{3.125634in}{2.767383in}}%
\pgfpathlineto{\pgfqpoint{3.125759in}{2.634199in}}%
\pgfpathlineto{\pgfqpoint{3.126381in}{2.671959in}}%
\pgfpathlineto{\pgfqpoint{3.127222in}{2.750892in}}%
\pgfpathlineto{\pgfqpoint{3.127602in}{2.774707in}}%
\pgfpathlineto{\pgfqpoint{3.128200in}{2.644559in}}%
\pgfpathlineto{\pgfqpoint{3.130573in}{2.644559in}}%
\pgfpathlineto{\pgfqpoint{3.130580in}{2.755135in}}%
\pgfpathlineto{\pgfqpoint{3.130685in}{2.638022in}}%
\pgfpathlineto{\pgfqpoint{3.131337in}{2.694957in}}%
\pgfpathlineto{\pgfqpoint{3.131387in}{2.679640in}}%
\pgfpathlineto{\pgfqpoint{3.131931in}{2.732192in}}%
\pgfpathlineto{\pgfqpoint{3.131980in}{2.716968in}}%
\pgfpathlineto{\pgfqpoint{3.132511in}{2.768597in}}%
\pgfpathlineto{\pgfqpoint{3.132610in}{2.668373in}}%
\pgfpathlineto{\pgfqpoint{3.133388in}{2.644559in}}%
\pgfpathlineto{\pgfqpoint{3.135520in}{2.644559in}}%
\pgfpathlineto{\pgfqpoint{3.135526in}{2.754424in}}%
\pgfpathlineto{\pgfqpoint{3.135656in}{2.627548in}}%
\pgfpathlineto{\pgfqpoint{3.136282in}{2.665837in}}%
\pgfpathlineto{\pgfqpoint{3.137123in}{2.758005in}}%
\pgfpathlineto{\pgfqpoint{3.137493in}{2.781176in}}%
\pgfpathlineto{\pgfqpoint{3.138034in}{2.644559in}}%
\pgfpathlineto{\pgfqpoint{3.140465in}{2.644559in}}%
\pgfpathlineto{\pgfqpoint{3.140471in}{2.754638in}}%
\pgfpathlineto{\pgfqpoint{3.140579in}{2.639808in}}%
\pgfpathlineto{\pgfqpoint{3.141230in}{2.693276in}}%
\pgfpathlineto{\pgfqpoint{3.141279in}{2.681647in}}%
\pgfpathlineto{\pgfqpoint{3.141823in}{2.730519in}}%
\pgfpathlineto{\pgfqpoint{3.141873in}{2.718968in}}%
\pgfpathlineto{\pgfqpoint{3.142403in}{2.766871in}}%
\pgfpathlineto{\pgfqpoint{3.142505in}{2.674670in}}%
\pgfpathlineto{\pgfqpoint{3.143254in}{2.644559in}}%
\pgfpathlineto{\pgfqpoint{3.145411in}{2.644560in}}%
\pgfpathlineto{\pgfqpoint{3.145415in}{2.768637in}}%
\pgfpathlineto{\pgfqpoint{3.145542in}{2.633687in}}%
\pgfpathlineto{\pgfqpoint{3.146164in}{2.671496in}}%
\pgfpathlineto{\pgfqpoint{3.147005in}{2.751489in}}%
\pgfpathlineto{\pgfqpoint{3.147384in}{2.775237in}}%
\pgfpathlineto{\pgfqpoint{3.147982in}{2.644559in}}%
\pgfpathlineto{\pgfqpoint{3.150355in}{2.644559in}}%
\pgfpathlineto{\pgfqpoint{3.150362in}{2.754981in}}%
\pgfpathlineto{\pgfqpoint{3.150467in}{2.637921in}}%
\pgfpathlineto{\pgfqpoint{3.151119in}{2.695060in}}%
\pgfpathlineto{\pgfqpoint{3.151168in}{2.679543in}}%
\pgfpathlineto{\pgfqpoint{3.151712in}{2.732295in}}%
\pgfpathlineto{\pgfqpoint{3.151762in}{2.716871in}}%
\pgfpathlineto{\pgfqpoint{3.152293in}{2.768697in}}%
\pgfpathlineto{\pgfqpoint{3.152392in}{2.668354in}}%
\pgfpathlineto{\pgfqpoint{3.153176in}{2.644559in}}%
\pgfpathlineto{\pgfqpoint{3.155301in}{2.644559in}}%
\pgfpathlineto{\pgfqpoint{3.155308in}{2.755253in}}%
\pgfpathlineto{\pgfqpoint{3.155414in}{2.636634in}}%
\pgfpathlineto{\pgfqpoint{3.156067in}{2.696257in}}%
\pgfpathlineto{\pgfqpoint{3.156116in}{2.678375in}}%
\pgfpathlineto{\pgfqpoint{3.156660in}{2.733487in}}%
\pgfpathlineto{\pgfqpoint{3.156710in}{2.715708in}}%
\pgfpathlineto{\pgfqpoint{3.157239in}{2.769815in}}%
\pgfpathlineto{\pgfqpoint{3.157337in}{2.668160in}}%
\pgfpathlineto{\pgfqpoint{3.158092in}{2.644559in}}%
\pgfpathlineto{\pgfqpoint{3.160245in}{2.644559in}}%
\pgfpathlineto{\pgfqpoint{3.160251in}{2.763972in}}%
\pgfpathlineto{\pgfqpoint{3.160375in}{2.626836in}}%
\pgfpathlineto{\pgfqpoint{3.160996in}{2.664463in}}%
\pgfpathlineto{\pgfqpoint{3.161836in}{2.757802in}}%
\pgfpathlineto{\pgfqpoint{3.162220in}{2.781825in}}%
\pgfpathlineto{\pgfqpoint{3.162744in}{2.644559in}}%
\pgfpathlineto{\pgfqpoint{3.165194in}{2.644948in}}%
\pgfpathlineto{\pgfqpoint{3.165196in}{2.772122in}}%
\pgfpathlineto{\pgfqpoint{3.165956in}{2.686835in}}%
\pgfpathlineto{\pgfqpoint{3.166303in}{2.706394in}}%
\pgfpathlineto{\pgfqpoint{3.167165in}{2.760586in}}%
\pgfpathlineto{\pgfqpoint{3.167192in}{2.748006in}}%
\pgfpathlineto{\pgfqpoint{3.167937in}{2.644559in}}%
\pgfpathlineto{\pgfqpoint{3.170132in}{2.644559in}}%
\pgfpathlineto{\pgfqpoint{3.170132in}{2.644559in}}%
\pgfusepath{stroke}%
\end{pgfscope}%
\begin{pgfscope}%
\pgfsetrectcap%
\pgfsetmiterjoin%
\pgfsetlinewidth{0.803000pt}%
\definecolor{currentstroke}{rgb}{0.000000,0.000000,0.000000}%
\pgfsetstrokecolor{currentstroke}%
\pgfsetdash{}{0pt}%
\pgfpathmoveto{\pgfqpoint{0.573768in}{2.601404in}}%
\pgfpathlineto{\pgfqpoint{0.573768in}{3.507654in}}%
\pgfusepath{stroke}%
\end{pgfscope}%
\begin{pgfscope}%
\pgfsetrectcap%
\pgfsetmiterjoin%
\pgfsetlinewidth{0.803000pt}%
\definecolor{currentstroke}{rgb}{0.000000,0.000000,0.000000}%
\pgfsetstrokecolor{currentstroke}%
\pgfsetdash{}{0pt}%
\pgfpathmoveto{\pgfqpoint{3.293768in}{2.601404in}}%
\pgfpathlineto{\pgfqpoint{3.293768in}{3.507654in}}%
\pgfusepath{stroke}%
\end{pgfscope}%
\begin{pgfscope}%
\pgfsetrectcap%
\pgfsetmiterjoin%
\pgfsetlinewidth{0.803000pt}%
\definecolor{currentstroke}{rgb}{0.000000,0.000000,0.000000}%
\pgfsetstrokecolor{currentstroke}%
\pgfsetdash{}{0pt}%
\pgfpathmoveto{\pgfqpoint{0.573768in}{2.601404in}}%
\pgfpathlineto{\pgfqpoint{3.293768in}{2.601404in}}%
\pgfusepath{stroke}%
\end{pgfscope}%
\begin{pgfscope}%
\pgfsetrectcap%
\pgfsetmiterjoin%
\pgfsetlinewidth{0.803000pt}%
\definecolor{currentstroke}{rgb}{0.000000,0.000000,0.000000}%
\pgfsetstrokecolor{currentstroke}%
\pgfsetdash{}{0pt}%
\pgfpathmoveto{\pgfqpoint{0.573768in}{3.507654in}}%
\pgfpathlineto{\pgfqpoint{3.293768in}{3.507654in}}%
\pgfusepath{stroke}%
\end{pgfscope}%
\begin{pgfscope}%
\definecolor{textcolor}{rgb}{0.000000,0.000000,0.000000}%
\pgfsetstrokecolor{textcolor}%
\pgfsetfillcolor{textcolor}%
\pgftext[x=0.628168in,y=3.444217in,left,top]{\color{textcolor}{\rmfamily\fontsize{8.000000}{9.600000}\selectfont\catcode`\^=\active\def^{\ifmmode\sp\else\^{}\fi}\catcode`\%=\active\def%{\%}(a)}}%
\end{pgfscope}%
\begin{pgfscope}%
\pgfsetbuttcap%
\pgfsetmiterjoin%
\definecolor{currentfill}{rgb}{1.000000,1.000000,1.000000}%
\pgfsetfillcolor{currentfill}%
\pgfsetlinewidth{0.000000pt}%
\definecolor{currentstroke}{rgb}{0.000000,0.000000,0.000000}%
\pgfsetstrokecolor{currentstroke}%
\pgfsetstrokeopacity{0.000000}%
\pgfsetdash{}{0pt}%
\pgfpathmoveto{\pgfqpoint{0.573768in}{1.532029in}}%
\pgfpathlineto{\pgfqpoint{3.293768in}{1.532029in}}%
\pgfpathlineto{\pgfqpoint{3.293768in}{2.438279in}}%
\pgfpathlineto{\pgfqpoint{0.573768in}{2.438279in}}%
\pgfpathlineto{\pgfqpoint{0.573768in}{1.532029in}}%
\pgfpathclose%
\pgfusepath{fill}%
\end{pgfscope}%
\begin{pgfscope}%
\pgfpathrectangle{\pgfqpoint{0.573768in}{1.532029in}}{\pgfqpoint{2.720000in}{0.906250in}}%
\pgfusepath{clip}%
\pgfsetbuttcap%
\pgfsetroundjoin%
\pgfsetlinewidth{0.501875pt}%
\definecolor{currentstroke}{rgb}{0.690196,0.690196,0.690196}%
\pgfsetstrokecolor{currentstroke}%
\pgfsetstrokeopacity{0.250000}%
\pgfsetdash{{1.850000pt}{0.800000pt}}{0.000000pt}%
\pgfpathmoveto{\pgfqpoint{0.697405in}{1.532029in}}%
\pgfpathlineto{\pgfqpoint{0.697405in}{2.438279in}}%
\pgfusepath{stroke}%
\end{pgfscope}%
\begin{pgfscope}%
\pgfsetbuttcap%
\pgfsetroundjoin%
\definecolor{currentfill}{rgb}{0.000000,0.000000,0.000000}%
\pgfsetfillcolor{currentfill}%
\pgfsetlinewidth{0.803000pt}%
\definecolor{currentstroke}{rgb}{0.000000,0.000000,0.000000}%
\pgfsetstrokecolor{currentstroke}%
\pgfsetdash{}{0pt}%
\pgfsys@defobject{currentmarker}{\pgfqpoint{0.000000in}{-0.048611in}}{\pgfqpoint{0.000000in}{0.000000in}}{%
\pgfpathmoveto{\pgfqpoint{0.000000in}{0.000000in}}%
\pgfpathlineto{\pgfqpoint{0.000000in}{-0.048611in}}%
\pgfusepath{stroke,fill}%
}%
\begin{pgfscope}%
\pgfsys@transformshift{0.697405in}{1.532029in}%
\pgfsys@useobject{currentmarker}{}%
\end{pgfscope}%
\end{pgfscope}%
\begin{pgfscope}%
\pgfpathrectangle{\pgfqpoint{0.573768in}{1.532029in}}{\pgfqpoint{2.720000in}{0.906250in}}%
\pgfusepath{clip}%
\pgfsetbuttcap%
\pgfsetroundjoin%
\pgfsetlinewidth{0.501875pt}%
\definecolor{currentstroke}{rgb}{0.690196,0.690196,0.690196}%
\pgfsetstrokecolor{currentstroke}%
\pgfsetstrokeopacity{0.250000}%
\pgfsetdash{{1.850000pt}{0.800000pt}}{0.000000pt}%
\pgfpathmoveto{\pgfqpoint{1.191950in}{1.532029in}}%
\pgfpathlineto{\pgfqpoint{1.191950in}{2.438279in}}%
\pgfusepath{stroke}%
\end{pgfscope}%
\begin{pgfscope}%
\pgfsetbuttcap%
\pgfsetroundjoin%
\definecolor{currentfill}{rgb}{0.000000,0.000000,0.000000}%
\pgfsetfillcolor{currentfill}%
\pgfsetlinewidth{0.803000pt}%
\definecolor{currentstroke}{rgb}{0.000000,0.000000,0.000000}%
\pgfsetstrokecolor{currentstroke}%
\pgfsetdash{}{0pt}%
\pgfsys@defobject{currentmarker}{\pgfqpoint{0.000000in}{-0.048611in}}{\pgfqpoint{0.000000in}{0.000000in}}{%
\pgfpathmoveto{\pgfqpoint{0.000000in}{0.000000in}}%
\pgfpathlineto{\pgfqpoint{0.000000in}{-0.048611in}}%
\pgfusepath{stroke,fill}%
}%
\begin{pgfscope}%
\pgfsys@transformshift{1.191950in}{1.532029in}%
\pgfsys@useobject{currentmarker}{}%
\end{pgfscope}%
\end{pgfscope}%
\begin{pgfscope}%
\pgfpathrectangle{\pgfqpoint{0.573768in}{1.532029in}}{\pgfqpoint{2.720000in}{0.906250in}}%
\pgfusepath{clip}%
\pgfsetbuttcap%
\pgfsetroundjoin%
\pgfsetlinewidth{0.501875pt}%
\definecolor{currentstroke}{rgb}{0.690196,0.690196,0.690196}%
\pgfsetstrokecolor{currentstroke}%
\pgfsetstrokeopacity{0.250000}%
\pgfsetdash{{1.850000pt}{0.800000pt}}{0.000000pt}%
\pgfpathmoveto{\pgfqpoint{1.686496in}{1.532029in}}%
\pgfpathlineto{\pgfqpoint{1.686496in}{2.438279in}}%
\pgfusepath{stroke}%
\end{pgfscope}%
\begin{pgfscope}%
\pgfsetbuttcap%
\pgfsetroundjoin%
\definecolor{currentfill}{rgb}{0.000000,0.000000,0.000000}%
\pgfsetfillcolor{currentfill}%
\pgfsetlinewidth{0.803000pt}%
\definecolor{currentstroke}{rgb}{0.000000,0.000000,0.000000}%
\pgfsetstrokecolor{currentstroke}%
\pgfsetdash{}{0pt}%
\pgfsys@defobject{currentmarker}{\pgfqpoint{0.000000in}{-0.048611in}}{\pgfqpoint{0.000000in}{0.000000in}}{%
\pgfpathmoveto{\pgfqpoint{0.000000in}{0.000000in}}%
\pgfpathlineto{\pgfqpoint{0.000000in}{-0.048611in}}%
\pgfusepath{stroke,fill}%
}%
\begin{pgfscope}%
\pgfsys@transformshift{1.686496in}{1.532029in}%
\pgfsys@useobject{currentmarker}{}%
\end{pgfscope}%
\end{pgfscope}%
\begin{pgfscope}%
\pgfpathrectangle{\pgfqpoint{0.573768in}{1.532029in}}{\pgfqpoint{2.720000in}{0.906250in}}%
\pgfusepath{clip}%
\pgfsetbuttcap%
\pgfsetroundjoin%
\pgfsetlinewidth{0.501875pt}%
\definecolor{currentstroke}{rgb}{0.690196,0.690196,0.690196}%
\pgfsetstrokecolor{currentstroke}%
\pgfsetstrokeopacity{0.250000}%
\pgfsetdash{{1.850000pt}{0.800000pt}}{0.000000pt}%
\pgfpathmoveto{\pgfqpoint{2.181041in}{1.532029in}}%
\pgfpathlineto{\pgfqpoint{2.181041in}{2.438279in}}%
\pgfusepath{stroke}%
\end{pgfscope}%
\begin{pgfscope}%
\pgfsetbuttcap%
\pgfsetroundjoin%
\definecolor{currentfill}{rgb}{0.000000,0.000000,0.000000}%
\pgfsetfillcolor{currentfill}%
\pgfsetlinewidth{0.803000pt}%
\definecolor{currentstroke}{rgb}{0.000000,0.000000,0.000000}%
\pgfsetstrokecolor{currentstroke}%
\pgfsetdash{}{0pt}%
\pgfsys@defobject{currentmarker}{\pgfqpoint{0.000000in}{-0.048611in}}{\pgfqpoint{0.000000in}{0.000000in}}{%
\pgfpathmoveto{\pgfqpoint{0.000000in}{0.000000in}}%
\pgfpathlineto{\pgfqpoint{0.000000in}{-0.048611in}}%
\pgfusepath{stroke,fill}%
}%
\begin{pgfscope}%
\pgfsys@transformshift{2.181041in}{1.532029in}%
\pgfsys@useobject{currentmarker}{}%
\end{pgfscope}%
\end{pgfscope}%
\begin{pgfscope}%
\pgfpathrectangle{\pgfqpoint{0.573768in}{1.532029in}}{\pgfqpoint{2.720000in}{0.906250in}}%
\pgfusepath{clip}%
\pgfsetbuttcap%
\pgfsetroundjoin%
\pgfsetlinewidth{0.501875pt}%
\definecolor{currentstroke}{rgb}{0.690196,0.690196,0.690196}%
\pgfsetstrokecolor{currentstroke}%
\pgfsetstrokeopacity{0.250000}%
\pgfsetdash{{1.850000pt}{0.800000pt}}{0.000000pt}%
\pgfpathmoveto{\pgfqpoint{2.675587in}{1.532029in}}%
\pgfpathlineto{\pgfqpoint{2.675587in}{2.438279in}}%
\pgfusepath{stroke}%
\end{pgfscope}%
\begin{pgfscope}%
\pgfsetbuttcap%
\pgfsetroundjoin%
\definecolor{currentfill}{rgb}{0.000000,0.000000,0.000000}%
\pgfsetfillcolor{currentfill}%
\pgfsetlinewidth{0.803000pt}%
\definecolor{currentstroke}{rgb}{0.000000,0.000000,0.000000}%
\pgfsetstrokecolor{currentstroke}%
\pgfsetdash{}{0pt}%
\pgfsys@defobject{currentmarker}{\pgfqpoint{0.000000in}{-0.048611in}}{\pgfqpoint{0.000000in}{0.000000in}}{%
\pgfpathmoveto{\pgfqpoint{0.000000in}{0.000000in}}%
\pgfpathlineto{\pgfqpoint{0.000000in}{-0.048611in}}%
\pgfusepath{stroke,fill}%
}%
\begin{pgfscope}%
\pgfsys@transformshift{2.675587in}{1.532029in}%
\pgfsys@useobject{currentmarker}{}%
\end{pgfscope}%
\end{pgfscope}%
\begin{pgfscope}%
\pgfpathrectangle{\pgfqpoint{0.573768in}{1.532029in}}{\pgfqpoint{2.720000in}{0.906250in}}%
\pgfusepath{clip}%
\pgfsetbuttcap%
\pgfsetroundjoin%
\pgfsetlinewidth{0.501875pt}%
\definecolor{currentstroke}{rgb}{0.690196,0.690196,0.690196}%
\pgfsetstrokecolor{currentstroke}%
\pgfsetstrokeopacity{0.250000}%
\pgfsetdash{{1.850000pt}{0.800000pt}}{0.000000pt}%
\pgfpathmoveto{\pgfqpoint{3.170132in}{1.532029in}}%
\pgfpathlineto{\pgfqpoint{3.170132in}{2.438279in}}%
\pgfusepath{stroke}%
\end{pgfscope}%
\begin{pgfscope}%
\pgfsetbuttcap%
\pgfsetroundjoin%
\definecolor{currentfill}{rgb}{0.000000,0.000000,0.000000}%
\pgfsetfillcolor{currentfill}%
\pgfsetlinewidth{0.803000pt}%
\definecolor{currentstroke}{rgb}{0.000000,0.000000,0.000000}%
\pgfsetstrokecolor{currentstroke}%
\pgfsetdash{}{0pt}%
\pgfsys@defobject{currentmarker}{\pgfqpoint{0.000000in}{-0.048611in}}{\pgfqpoint{0.000000in}{0.000000in}}{%
\pgfpathmoveto{\pgfqpoint{0.000000in}{0.000000in}}%
\pgfpathlineto{\pgfqpoint{0.000000in}{-0.048611in}}%
\pgfusepath{stroke,fill}%
}%
\begin{pgfscope}%
\pgfsys@transformshift{3.170132in}{1.532029in}%
\pgfsys@useobject{currentmarker}{}%
\end{pgfscope}%
\end{pgfscope}%
\begin{pgfscope}%
\pgfpathrectangle{\pgfqpoint{0.573768in}{1.532029in}}{\pgfqpoint{2.720000in}{0.906250in}}%
\pgfusepath{clip}%
\pgfsetbuttcap%
\pgfsetroundjoin%
\pgfsetlinewidth{0.501875pt}%
\definecolor{currentstroke}{rgb}{0.690196,0.690196,0.690196}%
\pgfsetstrokecolor{currentstroke}%
\pgfsetstrokeopacity{0.250000}%
\pgfsetdash{{1.850000pt}{0.800000pt}}{0.000000pt}%
\pgfpathmoveto{\pgfqpoint{0.573768in}{1.532029in}}%
\pgfpathlineto{\pgfqpoint{3.293768in}{1.532029in}}%
\pgfusepath{stroke}%
\end{pgfscope}%
\begin{pgfscope}%
\pgfsetbuttcap%
\pgfsetroundjoin%
\definecolor{currentfill}{rgb}{0.000000,0.000000,0.000000}%
\pgfsetfillcolor{currentfill}%
\pgfsetlinewidth{0.803000pt}%
\definecolor{currentstroke}{rgb}{0.000000,0.000000,0.000000}%
\pgfsetstrokecolor{currentstroke}%
\pgfsetdash{}{0pt}%
\pgfsys@defobject{currentmarker}{\pgfqpoint{-0.048611in}{0.000000in}}{\pgfqpoint{-0.000000in}{0.000000in}}{%
\pgfpathmoveto{\pgfqpoint{-0.000000in}{0.000000in}}%
\pgfpathlineto{\pgfqpoint{-0.048611in}{0.000000in}}%
\pgfusepath{stroke,fill}%
}%
\begin{pgfscope}%
\pgfsys@transformshift{0.573768in}{1.532029in}%
\pgfsys@useobject{currentmarker}{}%
\end{pgfscope}%
\end{pgfscope}%
\begin{pgfscope}%
\definecolor{textcolor}{rgb}{0.000000,0.000000,0.000000}%
\pgfsetstrokecolor{textcolor}%
\pgfsetfillcolor{textcolor}%
\pgftext[x=0.417518in, y=1.493449in, left, base]{\color{textcolor}{\rmfamily\fontsize{8.000000}{9.600000}\selectfont\catcode`\^=\active\def^{\ifmmode\sp\else\^{}\fi}\catcode`\%=\active\def%{\%}$\mathdefault{0}$}}%
\end{pgfscope}%
\begin{pgfscope}%
\pgfpathrectangle{\pgfqpoint{0.573768in}{1.532029in}}{\pgfqpoint{2.720000in}{0.906250in}}%
\pgfusepath{clip}%
\pgfsetbuttcap%
\pgfsetroundjoin%
\pgfsetlinewidth{0.501875pt}%
\definecolor{currentstroke}{rgb}{0.690196,0.690196,0.690196}%
\pgfsetstrokecolor{currentstroke}%
\pgfsetstrokeopacity{0.250000}%
\pgfsetdash{{1.850000pt}{0.800000pt}}{0.000000pt}%
\pgfpathmoveto{\pgfqpoint{0.573768in}{1.834113in}}%
\pgfpathlineto{\pgfqpoint{3.293768in}{1.834113in}}%
\pgfusepath{stroke}%
\end{pgfscope}%
\begin{pgfscope}%
\pgfsetbuttcap%
\pgfsetroundjoin%
\definecolor{currentfill}{rgb}{0.000000,0.000000,0.000000}%
\pgfsetfillcolor{currentfill}%
\pgfsetlinewidth{0.803000pt}%
\definecolor{currentstroke}{rgb}{0.000000,0.000000,0.000000}%
\pgfsetstrokecolor{currentstroke}%
\pgfsetdash{}{0pt}%
\pgfsys@defobject{currentmarker}{\pgfqpoint{-0.048611in}{0.000000in}}{\pgfqpoint{-0.000000in}{0.000000in}}{%
\pgfpathmoveto{\pgfqpoint{-0.000000in}{0.000000in}}%
\pgfpathlineto{\pgfqpoint{-0.048611in}{0.000000in}}%
\pgfusepath{stroke,fill}%
}%
\begin{pgfscope}%
\pgfsys@transformshift{0.573768in}{1.834113in}%
\pgfsys@useobject{currentmarker}{}%
\end{pgfscope}%
\end{pgfscope}%
\begin{pgfscope}%
\definecolor{textcolor}{rgb}{0.000000,0.000000,0.000000}%
\pgfsetstrokecolor{textcolor}%
\pgfsetfillcolor{textcolor}%
\pgftext[x=0.358489in, y=1.795533in, left, base]{\color{textcolor}{\rmfamily\fontsize{8.000000}{9.600000}\selectfont\catcode`\^=\active\def^{\ifmmode\sp\else\^{}\fi}\catcode`\%=\active\def%{\%}$\mathdefault{50}$}}%
\end{pgfscope}%
\begin{pgfscope}%
\pgfpathrectangle{\pgfqpoint{0.573768in}{1.532029in}}{\pgfqpoint{2.720000in}{0.906250in}}%
\pgfusepath{clip}%
\pgfsetbuttcap%
\pgfsetroundjoin%
\pgfsetlinewidth{0.501875pt}%
\definecolor{currentstroke}{rgb}{0.690196,0.690196,0.690196}%
\pgfsetstrokecolor{currentstroke}%
\pgfsetstrokeopacity{0.250000}%
\pgfsetdash{{1.850000pt}{0.800000pt}}{0.000000pt}%
\pgfpathmoveto{\pgfqpoint{0.573768in}{2.136196in}}%
\pgfpathlineto{\pgfqpoint{3.293768in}{2.136196in}}%
\pgfusepath{stroke}%
\end{pgfscope}%
\begin{pgfscope}%
\pgfsetbuttcap%
\pgfsetroundjoin%
\definecolor{currentfill}{rgb}{0.000000,0.000000,0.000000}%
\pgfsetfillcolor{currentfill}%
\pgfsetlinewidth{0.803000pt}%
\definecolor{currentstroke}{rgb}{0.000000,0.000000,0.000000}%
\pgfsetstrokecolor{currentstroke}%
\pgfsetdash{}{0pt}%
\pgfsys@defobject{currentmarker}{\pgfqpoint{-0.048611in}{0.000000in}}{\pgfqpoint{-0.000000in}{0.000000in}}{%
\pgfpathmoveto{\pgfqpoint{-0.000000in}{0.000000in}}%
\pgfpathlineto{\pgfqpoint{-0.048611in}{0.000000in}}%
\pgfusepath{stroke,fill}%
}%
\begin{pgfscope}%
\pgfsys@transformshift{0.573768in}{2.136196in}%
\pgfsys@useobject{currentmarker}{}%
\end{pgfscope}%
\end{pgfscope}%
\begin{pgfscope}%
\definecolor{textcolor}{rgb}{0.000000,0.000000,0.000000}%
\pgfsetstrokecolor{textcolor}%
\pgfsetfillcolor{textcolor}%
\pgftext[x=0.299460in, y=2.097616in, left, base]{\color{textcolor}{\rmfamily\fontsize{8.000000}{9.600000}\selectfont\catcode`\^=\active\def^{\ifmmode\sp\else\^{}\fi}\catcode`\%=\active\def%{\%}$\mathdefault{100}$}}%
\end{pgfscope}%
\begin{pgfscope}%
\pgfpathrectangle{\pgfqpoint{0.573768in}{1.532029in}}{\pgfqpoint{2.720000in}{0.906250in}}%
\pgfusepath{clip}%
\pgfsetbuttcap%
\pgfsetroundjoin%
\pgfsetlinewidth{0.501875pt}%
\definecolor{currentstroke}{rgb}{0.690196,0.690196,0.690196}%
\pgfsetstrokecolor{currentstroke}%
\pgfsetstrokeopacity{0.250000}%
\pgfsetdash{{1.850000pt}{0.800000pt}}{0.000000pt}%
\pgfpathmoveto{\pgfqpoint{0.573768in}{2.438279in}}%
\pgfpathlineto{\pgfqpoint{3.293768in}{2.438279in}}%
\pgfusepath{stroke}%
\end{pgfscope}%
\begin{pgfscope}%
\pgfsetbuttcap%
\pgfsetroundjoin%
\definecolor{currentfill}{rgb}{0.000000,0.000000,0.000000}%
\pgfsetfillcolor{currentfill}%
\pgfsetlinewidth{0.803000pt}%
\definecolor{currentstroke}{rgb}{0.000000,0.000000,0.000000}%
\pgfsetstrokecolor{currentstroke}%
\pgfsetdash{}{0pt}%
\pgfsys@defobject{currentmarker}{\pgfqpoint{-0.048611in}{0.000000in}}{\pgfqpoint{-0.000000in}{0.000000in}}{%
\pgfpathmoveto{\pgfqpoint{-0.000000in}{0.000000in}}%
\pgfpathlineto{\pgfqpoint{-0.048611in}{0.000000in}}%
\pgfusepath{stroke,fill}%
}%
\begin{pgfscope}%
\pgfsys@transformshift{0.573768in}{2.438279in}%
\pgfsys@useobject{currentmarker}{}%
\end{pgfscope}%
\end{pgfscope}%
\begin{pgfscope}%
\definecolor{textcolor}{rgb}{0.000000,0.000000,0.000000}%
\pgfsetstrokecolor{textcolor}%
\pgfsetfillcolor{textcolor}%
\pgftext[x=0.299460in, y=2.399699in, left, base]{\color{textcolor}{\rmfamily\fontsize{8.000000}{9.600000}\selectfont\catcode`\^=\active\def^{\ifmmode\sp\else\^{}\fi}\catcode`\%=\active\def%{\%}$\mathdefault{150}$}}%
\end{pgfscope}%
\begin{pgfscope}%
\definecolor{textcolor}{rgb}{0.000000,0.000000,0.000000}%
\pgfsetstrokecolor{textcolor}%
\pgfsetfillcolor{textcolor}%
\pgftext[x=0.243905in,y=1.985154in,,bottom,rotate=90.000000]{\color{textcolor}{\rmfamily\fontsize{8.000000}{9.600000}\selectfont\catcode`\^=\active\def^{\ifmmode\sp\else\^{}\fi}\catcode`\%=\active\def%{\%}$v_{DS}$ (V)}}%
\end{pgfscope}%
\begin{pgfscope}%
\pgfpathrectangle{\pgfqpoint{0.573768in}{1.532029in}}{\pgfqpoint{2.720000in}{0.906250in}}%
\pgfusepath{clip}%
\pgfsetrectcap%
\pgfsetroundjoin%
\pgfsetlinewidth{1.505625pt}%
\definecolor{currentstroke}{rgb}{0.835294,0.368627,0.000000}%
\pgfsetstrokecolor{currentstroke}%
\pgfsetdash{}{0pt}%
\pgfpathmoveto{\pgfqpoint{0.697405in}{1.532030in}}%
\pgfpathlineto{\pgfqpoint{0.697411in}{1.541696in}}%
\pgfpathlineto{\pgfqpoint{0.697414in}{1.545416in}}%
\pgfpathlineto{\pgfqpoint{0.697454in}{1.532031in}}%
\pgfpathlineto{\pgfqpoint{0.698159in}{1.532046in}}%
\pgfpathlineto{\pgfqpoint{0.699411in}{1.532297in}}%
\pgfpathlineto{\pgfqpoint{0.699491in}{1.776985in}}%
\pgfpathlineto{\pgfqpoint{0.700303in}{1.756390in}}%
\pgfpathlineto{\pgfqpoint{0.700352in}{1.754482in}}%
\pgfpathlineto{\pgfqpoint{0.700748in}{1.756616in}}%
\pgfpathlineto{\pgfqpoint{0.701045in}{1.756013in}}%
\pgfpathlineto{\pgfqpoint{0.701885in}{1.757747in}}%
\pgfpathlineto{\pgfqpoint{0.701935in}{1.757449in}}%
\pgfpathlineto{\pgfqpoint{0.702351in}{1.759313in}}%
\pgfpathlineto{\pgfqpoint{0.702357in}{1.764783in}}%
\pgfpathlineto{\pgfqpoint{0.702474in}{1.532074in}}%
\pgfpathlineto{\pgfqpoint{0.702993in}{1.532090in}}%
\pgfpathlineto{\pgfqpoint{0.704349in}{1.532195in}}%
\pgfpathlineto{\pgfqpoint{0.704422in}{1.805644in}}%
\pgfpathlineto{\pgfqpoint{0.705574in}{1.760701in}}%
\pgfpathlineto{\pgfqpoint{0.706414in}{1.764292in}}%
\pgfpathlineto{\pgfqpoint{0.706464in}{1.763716in}}%
\pgfpathlineto{\pgfqpoint{0.707057in}{1.765581in}}%
\pgfpathlineto{\pgfqpoint{0.707299in}{1.771584in}}%
\pgfpathlineto{\pgfqpoint{0.707302in}{1.773043in}}%
\pgfpathlineto{\pgfqpoint{0.707442in}{1.532111in}}%
\pgfpathlineto{\pgfqpoint{0.707576in}{1.532127in}}%
\pgfpathlineto{\pgfqpoint{0.709277in}{1.532159in}}%
\pgfpathlineto{\pgfqpoint{0.709284in}{1.532178in}}%
\pgfpathlineto{\pgfqpoint{0.709301in}{1.532850in}}%
\pgfpathlineto{\pgfqpoint{0.709361in}{1.845448in}}%
\pgfpathlineto{\pgfqpoint{0.710204in}{1.768960in}}%
\pgfpathlineto{\pgfqpoint{0.711045in}{1.773530in}}%
\pgfpathlineto{\pgfqpoint{0.711094in}{1.773158in}}%
\pgfpathlineto{\pgfqpoint{0.712248in}{1.784624in}}%
\pgfpathlineto{\pgfqpoint{0.712401in}{1.532150in}}%
\pgfpathlineto{\pgfqpoint{0.713030in}{1.532164in}}%
\pgfpathlineto{\pgfqpoint{0.714240in}{1.532342in}}%
\pgfpathlineto{\pgfqpoint{0.714318in}{1.882806in}}%
\pgfpathlineto{\pgfqpoint{0.715438in}{1.783528in}}%
\pgfpathlineto{\pgfqpoint{0.716872in}{1.791166in}}%
\pgfpathlineto{\pgfqpoint{0.716922in}{1.791120in}}%
\pgfpathlineto{\pgfqpoint{0.716971in}{1.791654in}}%
\pgfpathlineto{\pgfqpoint{0.717020in}{1.791645in}}%
\pgfpathlineto{\pgfqpoint{0.717193in}{1.799254in}}%
\pgfpathlineto{\pgfqpoint{0.717373in}{1.532185in}}%
\pgfpathlineto{\pgfqpoint{0.717934in}{1.532197in}}%
\pgfpathlineto{\pgfqpoint{0.719188in}{1.532502in}}%
\pgfpathlineto{\pgfqpoint{0.719250in}{1.916354in}}%
\pgfpathlineto{\pgfqpoint{0.720216in}{1.798209in}}%
\pgfpathlineto{\pgfqpoint{0.721848in}{1.808565in}}%
\pgfpathlineto{\pgfqpoint{0.722133in}{1.810555in}}%
\pgfpathlineto{\pgfqpoint{0.722138in}{1.816734in}}%
\pgfpathlineto{\pgfqpoint{0.722311in}{1.532218in}}%
\pgfpathlineto{\pgfqpoint{0.722771in}{1.532228in}}%
\pgfpathlineto{\pgfqpoint{0.724129in}{1.532414in}}%
\pgfpathlineto{\pgfqpoint{0.724156in}{1.576865in}}%
\pgfpathlineto{\pgfqpoint{0.724207in}{1.938925in}}%
\pgfpathlineto{\pgfqpoint{0.724932in}{1.816769in}}%
\pgfpathlineto{\pgfqpoint{0.724981in}{1.814020in}}%
\pgfpathlineto{\pgfqpoint{0.725624in}{1.820084in}}%
\pgfpathlineto{\pgfqpoint{0.725674in}{1.820102in}}%
\pgfpathlineto{\pgfqpoint{0.726712in}{1.827621in}}%
\pgfpathlineto{\pgfqpoint{0.726762in}{1.827502in}}%
\pgfpathlineto{\pgfqpoint{0.727084in}{1.836455in}}%
\pgfpathlineto{\pgfqpoint{0.727246in}{1.532237in}}%
\pgfpathlineto{\pgfqpoint{0.727844in}{1.532262in}}%
\pgfpathlineto{\pgfqpoint{0.729060in}{1.532292in}}%
\pgfpathlineto{\pgfqpoint{0.729078in}{1.532609in}}%
\pgfpathlineto{\pgfqpoint{0.729159in}{1.968885in}}%
\pgfpathlineto{\pgfqpoint{0.730119in}{1.836327in}}%
\pgfpathlineto{\pgfqpoint{0.731059in}{1.844391in}}%
\pgfpathlineto{\pgfqpoint{0.731603in}{1.848069in}}%
\pgfpathlineto{\pgfqpoint{0.732029in}{1.857961in}}%
\pgfpathlineto{\pgfqpoint{0.732208in}{1.532259in}}%
\pgfpathlineto{\pgfqpoint{0.732772in}{1.532272in}}%
\pgfpathlineto{\pgfqpoint{0.734020in}{1.532484in}}%
\pgfpathlineto{\pgfqpoint{0.734047in}{1.578129in}}%
\pgfpathlineto{\pgfqpoint{0.734102in}{2.001983in}}%
\pgfpathlineto{\pgfqpoint{0.734820in}{1.857964in}}%
\pgfpathlineto{\pgfqpoint{0.734870in}{1.855822in}}%
\pgfpathlineto{\pgfqpoint{0.735463in}{1.862946in}}%
\pgfpathlineto{\pgfqpoint{0.735513in}{1.861860in}}%
\pgfpathlineto{\pgfqpoint{0.736057in}{1.866779in}}%
\pgfpathlineto{\pgfqpoint{0.736106in}{1.866775in}}%
\pgfpathlineto{\pgfqpoint{0.736975in}{1.880342in}}%
\pgfpathlineto{\pgfqpoint{0.737159in}{1.532268in}}%
\pgfpathlineto{\pgfqpoint{0.737725in}{1.532281in}}%
\pgfpathlineto{\pgfqpoint{0.738965in}{1.532502in}}%
\pgfpathlineto{\pgfqpoint{0.738992in}{1.578378in}}%
\pgfpathlineto{\pgfqpoint{0.739051in}{2.042413in}}%
\pgfpathlineto{\pgfqpoint{0.739755in}{1.878323in}}%
\pgfpathlineto{\pgfqpoint{0.740695in}{1.887390in}}%
\pgfpathlineto{\pgfqpoint{0.740744in}{1.886923in}}%
\pgfpathlineto{\pgfqpoint{0.741041in}{1.889897in}}%
\pgfpathlineto{\pgfqpoint{0.741090in}{1.889760in}}%
\pgfpathlineto{\pgfqpoint{0.741920in}{1.903088in}}%
\pgfpathlineto{\pgfqpoint{0.742107in}{1.532273in}}%
\pgfpathlineto{\pgfqpoint{0.742674in}{1.532286in}}%
\pgfpathlineto{\pgfqpoint{0.743911in}{1.532508in}}%
\pgfpathlineto{\pgfqpoint{0.743938in}{1.578645in}}%
\pgfpathlineto{\pgfqpoint{0.743988in}{2.064171in}}%
\pgfpathlineto{\pgfqpoint{0.744691in}{1.899866in}}%
\pgfpathlineto{\pgfqpoint{0.745531in}{1.909620in}}%
\pgfpathlineto{\pgfqpoint{0.745581in}{1.908380in}}%
\pgfpathlineto{\pgfqpoint{0.746075in}{1.913367in}}%
\pgfpathlineto{\pgfqpoint{0.746125in}{1.913157in}}%
\pgfpathlineto{\pgfqpoint{0.746866in}{1.925727in}}%
\pgfpathlineto{\pgfqpoint{0.747014in}{1.532267in}}%
\pgfpathlineto{\pgfqpoint{0.747644in}{1.532280in}}%
\pgfpathlineto{\pgfqpoint{0.748862in}{1.532840in}}%
\pgfpathlineto{\pgfqpoint{0.748939in}{2.102175in}}%
\pgfpathlineto{\pgfqpoint{0.749905in}{1.926256in}}%
\pgfpathlineto{\pgfqpoint{0.751805in}{1.941331in}}%
\pgfpathlineto{\pgfqpoint{0.751811in}{1.947425in}}%
\pgfpathlineto{\pgfqpoint{0.751981in}{1.532262in}}%
\pgfpathlineto{\pgfqpoint{0.752377in}{1.532270in}}%
\pgfpathlineto{\pgfqpoint{0.753802in}{1.532486in}}%
\pgfpathlineto{\pgfqpoint{0.753828in}{1.578327in}}%
\pgfpathlineto{\pgfqpoint{0.753886in}{2.123507in}}%
\pgfpathlineto{\pgfqpoint{0.754604in}{1.945451in}}%
\pgfpathlineto{\pgfqpoint{0.755643in}{1.953903in}}%
\pgfpathlineto{\pgfqpoint{0.755692in}{1.953370in}}%
\pgfpathlineto{\pgfqpoint{0.756088in}{1.956806in}}%
\pgfpathlineto{\pgfqpoint{0.756137in}{1.956589in}}%
\pgfpathlineto{\pgfqpoint{0.756757in}{1.967639in}}%
\pgfpathlineto{\pgfqpoint{0.756916in}{1.532244in}}%
\pgfpathlineto{\pgfqpoint{0.757548in}{1.532257in}}%
\pgfpathlineto{\pgfqpoint{0.758747in}{1.532459in}}%
\pgfpathlineto{\pgfqpoint{0.758774in}{1.577848in}}%
\pgfpathlineto{\pgfqpoint{0.758834in}{2.137205in}}%
\pgfpathlineto{\pgfqpoint{0.759540in}{1.966793in}}%
\pgfpathlineto{\pgfqpoint{0.759589in}{1.964684in}}%
\pgfpathlineto{\pgfqpoint{0.760183in}{1.970048in}}%
\pgfpathlineto{\pgfqpoint{0.761699in}{1.984811in}}%
\pgfpathlineto{\pgfqpoint{0.761702in}{1.986106in}}%
\pgfpathlineto{\pgfqpoint{0.761712in}{1.764796in}}%
\pgfpathlineto{\pgfqpoint{0.761860in}{1.532217in}}%
\pgfpathlineto{\pgfqpoint{0.762485in}{1.532230in}}%
\pgfpathlineto{\pgfqpoint{0.763696in}{1.532542in}}%
\pgfpathlineto{\pgfqpoint{0.763780in}{2.135711in}}%
\pgfpathlineto{\pgfqpoint{0.764907in}{1.985554in}}%
\pgfpathlineto{\pgfqpoint{0.765946in}{1.991931in}}%
\pgfpathlineto{\pgfqpoint{0.765995in}{1.991703in}}%
\pgfpathlineto{\pgfqpoint{0.766292in}{1.993556in}}%
\pgfpathlineto{\pgfqpoint{0.766341in}{1.993532in}}%
\pgfpathlineto{\pgfqpoint{0.766648in}{2.002064in}}%
\pgfpathlineto{\pgfqpoint{0.766804in}{1.532194in}}%
\pgfpathlineto{\pgfqpoint{0.767420in}{1.532207in}}%
\pgfpathlineto{\pgfqpoint{0.768646in}{1.533068in}}%
\pgfpathlineto{\pgfqpoint{0.768724in}{2.128075in}}%
\pgfpathlineto{\pgfqpoint{0.769691in}{2.000651in}}%
\pgfpathlineto{\pgfqpoint{0.769740in}{1.999243in}}%
\pgfpathlineto{\pgfqpoint{0.770334in}{2.003073in}}%
\pgfpathlineto{\pgfqpoint{0.771587in}{2.009519in}}%
\pgfpathlineto{\pgfqpoint{0.771593in}{2.015000in}}%
\pgfpathlineto{\pgfqpoint{0.771750in}{1.532165in}}%
\pgfpathlineto{\pgfqpoint{0.772115in}{1.532173in}}%
\pgfpathlineto{\pgfqpoint{0.773584in}{1.532322in}}%
\pgfpathlineto{\pgfqpoint{0.773610in}{1.574689in}}%
\pgfpathlineto{\pgfqpoint{0.773672in}{2.117517in}}%
\pgfpathlineto{\pgfqpoint{0.774396in}{2.009522in}}%
\pgfpathlineto{\pgfqpoint{0.775237in}{2.014465in}}%
\pgfpathlineto{\pgfqpoint{0.775287in}{2.013809in}}%
\pgfpathlineto{\pgfqpoint{0.775781in}{2.016248in}}%
\pgfpathlineto{\pgfqpoint{0.775831in}{2.015805in}}%
\pgfpathlineto{\pgfqpoint{0.776538in}{2.024937in}}%
\pgfpathlineto{\pgfqpoint{0.776684in}{1.532131in}}%
\pgfpathlineto{\pgfqpoint{0.777324in}{1.532144in}}%
\pgfpathlineto{\pgfqpoint{0.778529in}{1.532263in}}%
\pgfpathlineto{\pgfqpoint{0.778555in}{1.573372in}}%
\pgfpathlineto{\pgfqpoint{0.778616in}{2.099339in}}%
\pgfpathlineto{\pgfqpoint{0.779368in}{2.018406in}}%
\pgfpathlineto{\pgfqpoint{0.780208in}{2.022790in}}%
\pgfpathlineto{\pgfqpoint{0.780258in}{2.021727in}}%
\pgfpathlineto{\pgfqpoint{0.780901in}{2.023673in}}%
\pgfpathlineto{\pgfqpoint{0.780950in}{2.023564in}}%
\pgfpathlineto{\pgfqpoint{0.781198in}{2.024281in}}%
\pgfpathlineto{\pgfqpoint{0.781395in}{2.024302in}}%
\pgfpathlineto{\pgfqpoint{0.781479in}{2.027997in}}%
\pgfpathlineto{\pgfqpoint{0.781484in}{2.031420in}}%
\pgfpathlineto{\pgfqpoint{0.781625in}{1.532094in}}%
\pgfpathlineto{\pgfqpoint{0.781906in}{1.532099in}}%
\pgfpathlineto{\pgfqpoint{0.783483in}{1.532693in}}%
\pgfpathlineto{\pgfqpoint{0.783560in}{2.076096in}}%
\pgfpathlineto{\pgfqpoint{0.784526in}{2.026290in}}%
\pgfpathlineto{\pgfqpoint{0.784822in}{2.025436in}}%
\pgfpathlineto{\pgfqpoint{0.785218in}{2.027129in}}%
\pgfpathlineto{\pgfqpoint{0.785267in}{2.026279in}}%
\pgfpathlineto{\pgfqpoint{0.785762in}{2.027410in}}%
\pgfpathlineto{\pgfqpoint{0.786059in}{2.027122in}}%
\pgfpathlineto{\pgfqpoint{0.786158in}{2.027088in}}%
\pgfpathlineto{\pgfqpoint{0.786429in}{2.034081in}}%
\pgfpathlineto{\pgfqpoint{0.786567in}{1.532052in}}%
\pgfpathlineto{\pgfqpoint{0.787208in}{1.532065in}}%
\pgfpathlineto{\pgfqpoint{0.788423in}{1.532161in}}%
\pgfpathlineto{\pgfqpoint{0.788466in}{1.796007in}}%
\pgfpathlineto{\pgfqpoint{0.788519in}{2.056775in}}%
\pgfpathlineto{\pgfqpoint{0.789246in}{2.026856in}}%
\pgfpathlineto{\pgfqpoint{0.789394in}{2.028016in}}%
\pgfpathlineto{\pgfqpoint{0.789345in}{2.026564in}}%
\pgfpathlineto{\pgfqpoint{0.790037in}{2.027502in}}%
\pgfpathlineto{\pgfqpoint{0.790779in}{2.026068in}}%
\pgfpathlineto{\pgfqpoint{0.790977in}{1.895051in}}%
\pgfpathlineto{\pgfqpoint{0.791388in}{1.532013in}}%
\pgfpathlineto{\pgfqpoint{0.792034in}{1.532051in}}%
\pgfpathlineto{\pgfqpoint{0.793365in}{1.532097in}}%
\pgfpathlineto{\pgfqpoint{0.793384in}{1.546407in}}%
\pgfpathlineto{\pgfqpoint{0.793560in}{2.051851in}}%
\pgfpathlineto{\pgfqpoint{0.794299in}{2.024632in}}%
\pgfpathlineto{\pgfqpoint{0.794349in}{2.027212in}}%
\pgfpathlineto{\pgfqpoint{0.794942in}{2.024923in}}%
\pgfpathlineto{\pgfqpoint{0.795140in}{1.901634in}}%
\pgfpathlineto{\pgfqpoint{0.795585in}{1.547728in}}%
\pgfpathlineto{\pgfqpoint{0.795931in}{1.797926in}}%
\pgfpathlineto{\pgfqpoint{0.796179in}{1.930538in}}%
\pgfpathlineto{\pgfqpoint{0.796401in}{1.532030in}}%
\pgfpathlineto{\pgfqpoint{0.798315in}{1.532122in}}%
\pgfpathlineto{\pgfqpoint{0.798457in}{2.045448in}}%
\pgfpathlineto{\pgfqpoint{0.799715in}{2.023879in}}%
\pgfpathlineto{\pgfqpoint{0.799863in}{2.023047in}}%
\pgfpathlineto{\pgfqpoint{0.800110in}{1.817836in}}%
\pgfpathlineto{\pgfqpoint{0.800457in}{1.548669in}}%
\pgfpathlineto{\pgfqpoint{0.800852in}{1.815288in}}%
\pgfpathlineto{\pgfqpoint{0.801050in}{1.929038in}}%
\pgfpathlineto{\pgfqpoint{0.801297in}{1.532041in}}%
\pgfpathlineto{\pgfqpoint{0.801301in}{1.532059in}}%
\pgfpathlineto{\pgfqpoint{0.802068in}{1.532051in}}%
\pgfpathlineto{\pgfqpoint{0.803260in}{1.532113in}}%
\pgfpathlineto{\pgfqpoint{0.803400in}{2.042224in}}%
\pgfpathlineto{\pgfqpoint{0.804740in}{2.021632in}}%
\pgfpathlineto{\pgfqpoint{0.804789in}{2.021469in}}%
\pgfpathlineto{\pgfqpoint{0.804938in}{1.949916in}}%
\pgfpathlineto{\pgfqpoint{0.805432in}{1.549388in}}%
\pgfpathlineto{\pgfqpoint{0.805778in}{1.800536in}}%
\pgfpathlineto{\pgfqpoint{0.806026in}{1.928456in}}%
\pgfpathlineto{\pgfqpoint{0.806254in}{1.532034in}}%
\pgfpathlineto{\pgfqpoint{0.807992in}{1.532072in}}%
\pgfpathlineto{\pgfqpoint{0.808207in}{1.532135in}}%
\pgfpathlineto{\pgfqpoint{0.808347in}{2.040180in}}%
\pgfpathlineto{\pgfqpoint{0.809419in}{2.017872in}}%
\pgfpathlineto{\pgfqpoint{0.809468in}{2.022817in}}%
\pgfpathlineto{\pgfqpoint{0.809814in}{2.009361in}}%
\pgfpathlineto{\pgfqpoint{0.810358in}{1.549519in}}%
\pgfpathlineto{\pgfqpoint{0.810754in}{1.820364in}}%
\pgfpathlineto{\pgfqpoint{0.810952in}{1.928105in}}%
\pgfpathlineto{\pgfqpoint{0.811187in}{1.532039in}}%
\pgfpathlineto{\pgfqpoint{0.811190in}{1.532066in}}%
\pgfpathlineto{\pgfqpoint{0.811889in}{1.532035in}}%
\pgfpathlineto{\pgfqpoint{0.811939in}{1.532052in}}%
\pgfpathlineto{\pgfqpoint{0.813147in}{1.532093in}}%
\pgfpathlineto{\pgfqpoint{0.813157in}{1.532386in}}%
\pgfpathlineto{\pgfqpoint{0.813403in}{2.042317in}}%
\pgfpathlineto{\pgfqpoint{0.814236in}{2.018581in}}%
\pgfpathlineto{\pgfqpoint{0.814434in}{2.019525in}}%
\pgfpathlineto{\pgfqpoint{0.814484in}{2.017709in}}%
\pgfpathlineto{\pgfqpoint{0.814681in}{2.017880in}}%
\pgfpathlineto{\pgfqpoint{0.814731in}{2.017142in}}%
\pgfpathlineto{\pgfqpoint{0.815028in}{1.740186in}}%
\pgfpathlineto{\pgfqpoint{0.815324in}{1.550171in}}%
\pgfpathlineto{\pgfqpoint{0.815720in}{1.826253in}}%
\pgfpathlineto{\pgfqpoint{0.815918in}{1.927425in}}%
\pgfpathlineto{\pgfqpoint{0.816133in}{1.532049in}}%
\pgfpathlineto{\pgfqpoint{0.816136in}{1.532059in}}%
\pgfpathlineto{\pgfqpoint{0.816928in}{1.532040in}}%
\pgfpathlineto{\pgfqpoint{0.818096in}{1.532111in}}%
\pgfpathlineto{\pgfqpoint{0.818136in}{1.718385in}}%
\pgfpathlineto{\pgfqpoint{0.818239in}{2.037756in}}%
\pgfpathlineto{\pgfqpoint{0.818915in}{2.022820in}}%
\pgfpathlineto{\pgfqpoint{0.820300in}{1.551621in}}%
\pgfpathlineto{\pgfqpoint{0.820646in}{1.792654in}}%
\pgfpathlineto{\pgfqpoint{0.820893in}{1.925771in}}%
\pgfpathlineto{\pgfqpoint{0.821110in}{1.532032in}}%
\pgfpathlineto{\pgfqpoint{0.823042in}{1.532122in}}%
\pgfpathlineto{\pgfqpoint{0.823185in}{2.036768in}}%
\pgfpathlineto{\pgfqpoint{0.824469in}{2.015136in}}%
\pgfpathlineto{\pgfqpoint{0.824617in}{2.014667in}}%
\pgfpathlineto{\pgfqpoint{0.824765in}{1.951551in}}%
\pgfpathlineto{\pgfqpoint{0.825260in}{1.552721in}}%
\pgfpathlineto{\pgfqpoint{0.825606in}{1.793860in}}%
\pgfpathlineto{\pgfqpoint{0.825853in}{1.924470in}}%
\pgfpathlineto{\pgfqpoint{0.826080in}{1.532025in}}%
\pgfpathlineto{\pgfqpoint{0.827772in}{1.532075in}}%
\pgfpathlineto{\pgfqpoint{0.827984in}{1.532097in}}%
\pgfpathlineto{\pgfqpoint{0.827993in}{1.532481in}}%
\pgfpathlineto{\pgfqpoint{0.828268in}{2.037958in}}%
\pgfpathlineto{\pgfqpoint{0.829048in}{2.014944in}}%
\pgfpathlineto{\pgfqpoint{0.829691in}{1.980722in}}%
\pgfpathlineto{\pgfqpoint{0.830185in}{1.554579in}}%
\pgfpathlineto{\pgfqpoint{0.830581in}{1.805596in}}%
\pgfpathlineto{\pgfqpoint{0.830828in}{1.922209in}}%
\pgfpathlineto{\pgfqpoint{0.831031in}{1.532026in}}%
\pgfpathlineto{\pgfqpoint{0.832934in}{1.532134in}}%
\pgfpathlineto{\pgfqpoint{0.833066in}{2.032275in}}%
\pgfpathlineto{\pgfqpoint{0.834137in}{2.014459in}}%
\pgfpathlineto{\pgfqpoint{0.834731in}{1.909469in}}%
\pgfpathlineto{\pgfqpoint{0.835176in}{1.554025in}}%
\pgfpathlineto{\pgfqpoint{0.835522in}{1.785326in}}%
\pgfpathlineto{\pgfqpoint{0.835769in}{1.922918in}}%
\pgfpathlineto{\pgfqpoint{0.836013in}{1.532024in}}%
\pgfpathlineto{\pgfqpoint{0.837648in}{1.532075in}}%
\pgfpathlineto{\pgfqpoint{0.837875in}{1.532099in}}%
\pgfpathlineto{\pgfqpoint{0.837904in}{1.576035in}}%
\pgfpathlineto{\pgfqpoint{0.838151in}{2.031469in}}%
\pgfpathlineto{\pgfqpoint{0.838695in}{2.009372in}}%
\pgfpathlineto{\pgfqpoint{0.838992in}{2.011864in}}%
\pgfpathlineto{\pgfqpoint{0.839437in}{2.010792in}}%
\pgfpathlineto{\pgfqpoint{0.839585in}{1.997415in}}%
\pgfpathlineto{\pgfqpoint{0.840129in}{1.555388in}}%
\pgfpathlineto{\pgfqpoint{0.840525in}{1.820703in}}%
\pgfpathlineto{\pgfqpoint{0.840723in}{1.921221in}}%
\pgfpathlineto{\pgfqpoint{0.840917in}{1.532025in}}%
\pgfpathlineto{\pgfqpoint{0.842514in}{1.532074in}}%
\pgfpathlineto{\pgfqpoint{0.842819in}{1.532094in}}%
\pgfpathlineto{\pgfqpoint{0.842827in}{1.532183in}}%
\pgfpathlineto{\pgfqpoint{0.843068in}{2.037172in}}%
\pgfpathlineto{\pgfqpoint{0.844017in}{2.009030in}}%
\pgfpathlineto{\pgfqpoint{0.844166in}{2.009028in}}%
\pgfpathlineto{\pgfqpoint{0.844413in}{2.008457in}}%
\pgfpathlineto{\pgfqpoint{0.844512in}{2.007243in}}%
\pgfpathlineto{\pgfqpoint{0.844957in}{1.596257in}}%
\pgfpathlineto{\pgfqpoint{0.845105in}{1.555909in}}%
\pgfpathlineto{\pgfqpoint{0.845402in}{1.737871in}}%
\pgfpathlineto{\pgfqpoint{0.845699in}{1.920767in}}%
\pgfpathlineto{\pgfqpoint{0.845905in}{1.532026in}}%
\pgfpathlineto{\pgfqpoint{0.845959in}{1.532026in}}%
\pgfpathlineto{\pgfqpoint{0.847769in}{1.532122in}}%
\pgfpathlineto{\pgfqpoint{0.847910in}{2.028517in}}%
\pgfpathlineto{\pgfqpoint{0.849201in}{2.007435in}}%
\pgfpathlineto{\pgfqpoint{0.849448in}{2.006531in}}%
\pgfpathlineto{\pgfqpoint{0.849646in}{1.881839in}}%
\pgfpathlineto{\pgfqpoint{0.850041in}{1.558720in}}%
\pgfpathlineto{\pgfqpoint{0.850437in}{1.806038in}}%
\pgfpathlineto{\pgfqpoint{0.850674in}{1.918031in}}%
\pgfpathlineto{\pgfqpoint{0.850855in}{1.532028in}}%
\pgfpathlineto{\pgfqpoint{0.850882in}{1.532037in}}%
\pgfpathlineto{\pgfqpoint{0.852705in}{1.532081in}}%
\pgfpathlineto{\pgfqpoint{0.852717in}{1.532157in}}%
\pgfpathlineto{\pgfqpoint{0.852868in}{2.025098in}}%
\pgfpathlineto{\pgfqpoint{0.853882in}{2.006952in}}%
\pgfpathlineto{\pgfqpoint{0.854080in}{2.011440in}}%
\pgfpathlineto{\pgfqpoint{0.854575in}{1.921132in}}%
\pgfpathlineto{\pgfqpoint{0.855020in}{1.560366in}}%
\pgfpathlineto{\pgfqpoint{0.855415in}{1.813798in}}%
\pgfpathlineto{\pgfqpoint{0.855666in}{1.922550in}}%
\pgfpathlineto{\pgfqpoint{0.855824in}{1.532028in}}%
\pgfpathlineto{\pgfqpoint{0.857662in}{1.532142in}}%
\pgfpathlineto{\pgfqpoint{0.857806in}{2.027236in}}%
\pgfpathlineto{\pgfqpoint{0.858894in}{2.008412in}}%
\pgfpathlineto{\pgfqpoint{0.859636in}{1.819877in}}%
\pgfpathlineto{\pgfqpoint{0.859982in}{1.560886in}}%
\pgfpathlineto{\pgfqpoint{0.860377in}{1.808239in}}%
\pgfpathlineto{\pgfqpoint{0.860611in}{1.925376in}}%
\pgfpathlineto{\pgfqpoint{0.860770in}{1.532028in}}%
\pgfpathlineto{\pgfqpoint{0.860795in}{1.532039in}}%
\pgfpathlineto{\pgfqpoint{0.862596in}{1.532083in}}%
\pgfpathlineto{\pgfqpoint{0.862607in}{1.532143in}}%
\pgfpathlineto{\pgfqpoint{0.862750in}{2.026128in}}%
\pgfpathlineto{\pgfqpoint{0.863806in}{2.007092in}}%
\pgfpathlineto{\pgfqpoint{0.864597in}{1.828240in}}%
\pgfpathlineto{\pgfqpoint{0.864993in}{1.559823in}}%
\pgfpathlineto{\pgfqpoint{0.865339in}{1.802608in}}%
\pgfpathlineto{\pgfqpoint{0.865554in}{1.924214in}}%
\pgfpathlineto{\pgfqpoint{0.865709in}{1.532031in}}%
\pgfpathlineto{\pgfqpoint{0.865788in}{1.532037in}}%
\pgfpathlineto{\pgfqpoint{0.867547in}{1.532099in}}%
\pgfpathlineto{\pgfqpoint{0.867556in}{1.532288in}}%
\pgfpathlineto{\pgfqpoint{0.867762in}{2.024885in}}%
\pgfpathlineto{\pgfqpoint{0.868636in}{2.001822in}}%
\pgfpathlineto{\pgfqpoint{0.868834in}{2.002472in}}%
\pgfpathlineto{\pgfqpoint{0.869230in}{2.001144in}}%
\pgfpathlineto{\pgfqpoint{0.869279in}{2.001545in}}%
\pgfpathlineto{\pgfqpoint{0.869378in}{1.992265in}}%
\pgfpathlineto{\pgfqpoint{0.869922in}{1.561855in}}%
\pgfpathlineto{\pgfqpoint{0.870367in}{1.849552in}}%
\pgfpathlineto{\pgfqpoint{0.870500in}{1.919462in}}%
\pgfpathlineto{\pgfqpoint{0.870629in}{1.532027in}}%
\pgfpathlineto{\pgfqpoint{0.872322in}{1.532064in}}%
\pgfpathlineto{\pgfqpoint{0.872495in}{1.532110in}}%
\pgfpathlineto{\pgfqpoint{0.872535in}{1.718290in}}%
\pgfpathlineto{\pgfqpoint{0.872619in}{2.020129in}}%
\pgfpathlineto{\pgfqpoint{0.873322in}{2.005619in}}%
\pgfpathlineto{\pgfqpoint{0.873470in}{1.994476in}}%
\pgfpathlineto{\pgfqpoint{0.873421in}{2.007440in}}%
\pgfpathlineto{\pgfqpoint{0.874113in}{2.000277in}}%
\pgfpathlineto{\pgfqpoint{0.874311in}{2.001937in}}%
\pgfpathlineto{\pgfqpoint{0.874360in}{1.987453in}}%
\pgfpathlineto{\pgfqpoint{0.874904in}{1.561612in}}%
\pgfpathlineto{\pgfqpoint{0.875349in}{1.856512in}}%
\pgfpathlineto{\pgfqpoint{0.875446in}{1.913747in}}%
\pgfpathlineto{\pgfqpoint{0.875581in}{1.532029in}}%
\pgfpathlineto{\pgfqpoint{0.877445in}{1.532173in}}%
\pgfpathlineto{\pgfqpoint{0.877574in}{2.021162in}}%
\pgfpathlineto{\pgfqpoint{0.878597in}{1.996831in}}%
\pgfpathlineto{\pgfqpoint{0.878647in}{2.002278in}}%
\pgfpathlineto{\pgfqpoint{0.879290in}{1.996530in}}%
\pgfpathlineto{\pgfqpoint{0.879537in}{1.796228in}}%
\pgfpathlineto{\pgfqpoint{0.879883in}{1.562105in}}%
\pgfpathlineto{\pgfqpoint{0.880279in}{1.824620in}}%
\pgfpathlineto{\pgfqpoint{0.880392in}{1.904416in}}%
\pgfpathlineto{\pgfqpoint{0.880522in}{1.532028in}}%
\pgfpathlineto{\pgfqpoint{0.880655in}{1.532044in}}%
\pgfpathlineto{\pgfqpoint{0.882321in}{1.532068in}}%
\pgfpathlineto{\pgfqpoint{0.882388in}{1.532132in}}%
\pgfpathlineto{\pgfqpoint{0.882520in}{2.021596in}}%
\pgfpathlineto{\pgfqpoint{0.883717in}{1.998736in}}%
\pgfpathlineto{\pgfqpoint{0.884311in}{1.981585in}}%
\pgfpathlineto{\pgfqpoint{0.884855in}{1.562659in}}%
\pgfpathlineto{\pgfqpoint{0.885250in}{1.827605in}}%
\pgfpathlineto{\pgfqpoint{0.885337in}{1.893895in}}%
\pgfpathlineto{\pgfqpoint{0.885474in}{1.532029in}}%
\pgfpathlineto{\pgfqpoint{0.885607in}{1.532044in}}%
\pgfpathlineto{\pgfqpoint{0.887323in}{1.532084in}}%
\pgfpathlineto{\pgfqpoint{0.887334in}{1.532143in}}%
\pgfpathlineto{\pgfqpoint{0.887470in}{2.020868in}}%
\pgfpathlineto{\pgfqpoint{0.888546in}{1.999133in}}%
\pgfpathlineto{\pgfqpoint{0.889386in}{1.892995in}}%
\pgfpathlineto{\pgfqpoint{0.889831in}{1.563429in}}%
\pgfpathlineto{\pgfqpoint{0.890178in}{1.792252in}}%
\pgfpathlineto{\pgfqpoint{0.890282in}{1.880994in}}%
\pgfpathlineto{\pgfqpoint{0.890426in}{1.532028in}}%
\pgfpathlineto{\pgfqpoint{0.890638in}{1.532032in}}%
\pgfpathlineto{\pgfqpoint{0.892190in}{1.532082in}}%
\pgfpathlineto{\pgfqpoint{0.892266in}{1.532081in}}%
\pgfpathlineto{\pgfqpoint{0.892276in}{1.532111in}}%
\pgfpathlineto{\pgfqpoint{0.892315in}{1.698689in}}%
\pgfpathlineto{\pgfqpoint{0.892402in}{2.017815in}}%
\pgfpathlineto{\pgfqpoint{0.893102in}{1.998565in}}%
\pgfpathlineto{\pgfqpoint{0.893251in}{1.991722in}}%
\pgfpathlineto{\pgfqpoint{0.893201in}{2.000273in}}%
\pgfpathlineto{\pgfqpoint{0.893894in}{1.995304in}}%
\pgfpathlineto{\pgfqpoint{0.893993in}{1.996494in}}%
\pgfpathlineto{\pgfqpoint{0.894042in}{1.993996in}}%
\pgfpathlineto{\pgfqpoint{0.894092in}{1.996911in}}%
\pgfpathlineto{\pgfqpoint{0.894240in}{1.983180in}}%
\pgfpathlineto{\pgfqpoint{0.894784in}{1.563879in}}%
\pgfpathlineto{\pgfqpoint{0.895177in}{1.818874in}}%
\pgfpathlineto{\pgfqpoint{0.895228in}{1.865924in}}%
\pgfpathlineto{\pgfqpoint{0.895363in}{1.532032in}}%
\pgfpathlineto{\pgfqpoint{0.895548in}{1.532043in}}%
\pgfpathlineto{\pgfqpoint{0.897214in}{1.532085in}}%
\pgfpathlineto{\pgfqpoint{0.897225in}{1.532151in}}%
\pgfpathlineto{\pgfqpoint{0.897350in}{2.016937in}}%
\pgfpathlineto{\pgfqpoint{0.898448in}{1.993540in}}%
\pgfpathlineto{\pgfqpoint{0.898497in}{1.996383in}}%
\pgfpathlineto{\pgfqpoint{0.898547in}{1.993247in}}%
\pgfpathlineto{\pgfqpoint{0.899140in}{1.993987in}}%
\pgfpathlineto{\pgfqpoint{0.899338in}{1.869622in}}%
\pgfpathlineto{\pgfqpoint{0.899733in}{1.565470in}}%
\pgfpathlineto{\pgfqpoint{0.900124in}{1.805681in}}%
\pgfpathlineto{\pgfqpoint{0.900173in}{1.853084in}}%
\pgfpathlineto{\pgfqpoint{0.900322in}{1.532033in}}%
\pgfpathlineto{\pgfqpoint{0.900534in}{1.532037in}}%
\pgfpathlineto{\pgfqpoint{0.902168in}{1.532117in}}%
\pgfpathlineto{\pgfqpoint{0.902210in}{1.749223in}}%
\pgfpathlineto{\pgfqpoint{0.902491in}{2.018733in}}%
\pgfpathlineto{\pgfqpoint{0.902962in}{1.995416in}}%
\pgfpathlineto{\pgfqpoint{0.903011in}{1.992129in}}%
\pgfpathlineto{\pgfqpoint{0.903753in}{1.993856in}}%
\pgfpathlineto{\pgfqpoint{0.903852in}{1.994057in}}%
\pgfpathlineto{\pgfqpoint{0.904149in}{1.983317in}}%
\pgfpathlineto{\pgfqpoint{0.905255in}{1.532029in}}%
\pgfpathlineto{\pgfqpoint{0.905832in}{1.532055in}}%
\pgfpathlineto{\pgfqpoint{0.907104in}{1.532084in}}%
\pgfpathlineto{\pgfqpoint{0.907115in}{1.532130in}}%
\pgfpathlineto{\pgfqpoint{0.907247in}{2.015041in}}%
\pgfpathlineto{\pgfqpoint{0.908483in}{1.993003in}}%
\pgfpathlineto{\pgfqpoint{0.909076in}{1.991156in}}%
\pgfpathlineto{\pgfqpoint{0.909422in}{1.686697in}}%
\pgfpathlineto{\pgfqpoint{0.909670in}{1.566629in}}%
\pgfpathlineto{\pgfqpoint{0.910013in}{1.779563in}}%
\pgfpathlineto{\pgfqpoint{0.910064in}{1.831507in}}%
\pgfpathlineto{\pgfqpoint{0.910192in}{1.532032in}}%
\pgfpathlineto{\pgfqpoint{0.910490in}{1.532046in}}%
\pgfpathlineto{\pgfqpoint{0.912050in}{1.532086in}}%
\pgfpathlineto{\pgfqpoint{0.912061in}{1.532144in}}%
\pgfpathlineto{\pgfqpoint{0.912184in}{2.013501in}}%
\pgfpathlineto{\pgfqpoint{0.913277in}{1.989972in}}%
\pgfpathlineto{\pgfqpoint{0.913425in}{1.994096in}}%
\pgfpathlineto{\pgfqpoint{0.913376in}{1.989178in}}%
\pgfpathlineto{\pgfqpoint{0.914019in}{1.990107in}}%
\pgfpathlineto{\pgfqpoint{0.914315in}{1.756363in}}%
\pgfpathlineto{\pgfqpoint{0.915129in}{1.532033in}}%
\pgfpathlineto{\pgfqpoint{0.915010in}{1.819559in}}%
\pgfpathlineto{\pgfqpoint{0.915285in}{1.532036in}}%
\pgfpathlineto{\pgfqpoint{0.917004in}{1.532114in}}%
\pgfpathlineto{\pgfqpoint{0.917045in}{1.722838in}}%
\pgfpathlineto{\pgfqpoint{0.917132in}{2.013297in}}%
\pgfpathlineto{\pgfqpoint{0.917800in}{1.986120in}}%
\pgfpathlineto{\pgfqpoint{0.917849in}{1.994844in}}%
\pgfpathlineto{\pgfqpoint{0.918541in}{1.991655in}}%
\pgfpathlineto{\pgfqpoint{0.918640in}{1.992262in}}%
\pgfpathlineto{\pgfqpoint{0.919085in}{1.947257in}}%
\pgfpathlineto{\pgfqpoint{0.920077in}{1.532029in}}%
\pgfpathlineto{\pgfqpoint{0.920343in}{1.532034in}}%
\pgfpathlineto{\pgfqpoint{0.921951in}{1.532130in}}%
\pgfpathlineto{\pgfqpoint{0.922066in}{2.010679in}}%
\pgfpathlineto{\pgfqpoint{0.923290in}{1.989714in}}%
\pgfpathlineto{\pgfqpoint{0.923389in}{1.989892in}}%
\pgfpathlineto{\pgfqpoint{0.923933in}{1.988669in}}%
\pgfpathlineto{\pgfqpoint{0.924180in}{1.808607in}}%
\pgfpathlineto{\pgfqpoint{0.925029in}{1.532028in}}%
\pgfpathlineto{\pgfqpoint{0.925188in}{1.532031in}}%
\pgfpathlineto{\pgfqpoint{0.926695in}{1.532080in}}%
\pgfpathlineto{\pgfqpoint{0.926895in}{1.532116in}}%
\pgfpathlineto{\pgfqpoint{0.926935in}{1.715808in}}%
\pgfpathlineto{\pgfqpoint{0.927021in}{2.010513in}}%
\pgfpathlineto{\pgfqpoint{0.927681in}{1.985715in}}%
\pgfpathlineto{\pgfqpoint{0.927730in}{1.992615in}}%
\pgfpathlineto{\pgfqpoint{0.927780in}{1.984808in}}%
\pgfpathlineto{\pgfqpoint{0.928472in}{1.988489in}}%
\pgfpathlineto{\pgfqpoint{0.928620in}{1.989933in}}%
\pgfpathlineto{\pgfqpoint{0.928670in}{1.986836in}}%
\pgfpathlineto{\pgfqpoint{0.928868in}{1.987470in}}%
\pgfpathlineto{\pgfqpoint{0.929066in}{1.886924in}}%
\pgfpathlineto{\pgfqpoint{0.929978in}{1.532028in}}%
\pgfpathlineto{\pgfqpoint{0.930163in}{1.532048in}}%
\pgfpathlineto{\pgfqpoint{0.931599in}{1.532064in}}%
\pgfpathlineto{\pgfqpoint{0.931840in}{1.532111in}}%
\pgfpathlineto{\pgfqpoint{0.931867in}{1.576333in}}%
\pgfpathlineto{\pgfqpoint{0.931966in}{2.010707in}}%
\pgfpathlineto{\pgfqpoint{0.932656in}{1.986551in}}%
\pgfpathlineto{\pgfqpoint{0.932953in}{1.989041in}}%
\pgfpathlineto{\pgfqpoint{0.933398in}{1.988421in}}%
\pgfpathlineto{\pgfqpoint{0.933942in}{1.952155in}}%
\pgfpathlineto{\pgfqpoint{0.934934in}{1.532028in}}%
\pgfpathlineto{\pgfqpoint{0.935238in}{1.532051in}}%
\pgfpathlineto{\pgfqpoint{0.936775in}{1.532082in}}%
\pgfpathlineto{\pgfqpoint{0.936785in}{1.532111in}}%
\pgfpathlineto{\pgfqpoint{0.936812in}{1.575871in}}%
\pgfpathlineto{\pgfqpoint{0.936911in}{2.009357in}}%
\pgfpathlineto{\pgfqpoint{0.937607in}{1.988028in}}%
\pgfpathlineto{\pgfqpoint{0.937656in}{1.985717in}}%
\pgfpathlineto{\pgfqpoint{0.937706in}{1.988194in}}%
\pgfpathlineto{\pgfqpoint{0.938398in}{1.986869in}}%
\pgfpathlineto{\pgfqpoint{0.938497in}{1.987060in}}%
\pgfpathlineto{\pgfqpoint{0.938843in}{1.979013in}}%
\pgfpathlineto{\pgfqpoint{0.939878in}{1.532029in}}%
\pgfpathlineto{\pgfqpoint{0.940470in}{1.532056in}}%
\pgfpathlineto{\pgfqpoint{0.941723in}{1.532085in}}%
\pgfpathlineto{\pgfqpoint{0.941734in}{1.532143in}}%
\pgfpathlineto{\pgfqpoint{0.941856in}{2.008445in}}%
\pgfpathlineto{\pgfqpoint{0.942982in}{1.987802in}}%
\pgfpathlineto{\pgfqpoint{0.943922in}{1.889813in}}%
\pgfpathlineto{\pgfqpoint{0.944816in}{1.532031in}}%
\pgfpathlineto{\pgfqpoint{0.945028in}{1.532035in}}%
\pgfpathlineto{\pgfqpoint{0.946680in}{1.532152in}}%
\pgfpathlineto{\pgfqpoint{0.946799in}{2.007436in}}%
\pgfpathlineto{\pgfqpoint{0.947880in}{1.983958in}}%
\pgfpathlineto{\pgfqpoint{0.948028in}{1.987054in}}%
\pgfpathlineto{\pgfqpoint{0.947979in}{1.983182in}}%
\pgfpathlineto{\pgfqpoint{0.948622in}{1.983947in}}%
\pgfpathlineto{\pgfqpoint{0.948671in}{1.984284in}}%
\pgfpathlineto{\pgfqpoint{0.948721in}{1.983868in}}%
\pgfpathlineto{\pgfqpoint{0.949759in}{1.532029in}}%
\pgfpathlineto{\pgfqpoint{0.950842in}{1.532066in}}%
\pgfpathlineto{\pgfqpoint{0.951613in}{1.532085in}}%
\pgfpathlineto{\pgfqpoint{0.951625in}{1.532142in}}%
\pgfpathlineto{\pgfqpoint{0.951747in}{2.007225in}}%
\pgfpathlineto{\pgfqpoint{0.952898in}{1.983753in}}%
\pgfpathlineto{\pgfqpoint{0.953195in}{1.984699in}}%
\pgfpathlineto{\pgfqpoint{0.953591in}{1.983136in}}%
\pgfpathlineto{\pgfqpoint{0.953640in}{1.983542in}}%
\pgfpathlineto{\pgfqpoint{0.953789in}{1.923011in}}%
\pgfpathlineto{\pgfqpoint{0.954699in}{1.532031in}}%
\pgfpathlineto{\pgfqpoint{0.954968in}{1.532037in}}%
\pgfpathlineto{\pgfqpoint{0.956570in}{1.532146in}}%
\pgfpathlineto{\pgfqpoint{0.956696in}{2.007182in}}%
\pgfpathlineto{\pgfqpoint{0.957767in}{1.985064in}}%
\pgfpathlineto{\pgfqpoint{0.958509in}{1.982011in}}%
\pgfpathlineto{\pgfqpoint{0.958558in}{1.983327in}}%
\pgfpathlineto{\pgfqpoint{0.958657in}{1.972936in}}%
\pgfpathlineto{\pgfqpoint{0.959660in}{1.532036in}}%
\pgfpathlineto{\pgfqpoint{0.960258in}{1.532049in}}%
\pgfpathlineto{\pgfqpoint{0.961518in}{1.532214in}}%
\pgfpathlineto{\pgfqpoint{0.961639in}{2.006429in}}%
\pgfpathlineto{\pgfqpoint{0.962667in}{1.982745in}}%
\pgfpathlineto{\pgfqpoint{0.962914in}{1.982952in}}%
\pgfpathlineto{\pgfqpoint{0.963557in}{1.981393in}}%
\pgfpathlineto{\pgfqpoint{0.963755in}{1.865185in}}%
\pgfpathlineto{\pgfqpoint{0.964595in}{1.532033in}}%
\pgfpathlineto{\pgfqpoint{0.964836in}{1.532044in}}%
\pgfpathlineto{\pgfqpoint{0.966459in}{1.532122in}}%
\pgfpathlineto{\pgfqpoint{0.966502in}{1.778707in}}%
\pgfpathlineto{\pgfqpoint{0.966586in}{2.005741in}}%
\pgfpathlineto{\pgfqpoint{0.967278in}{1.980478in}}%
\pgfpathlineto{\pgfqpoint{0.967327in}{1.983177in}}%
\pgfpathlineto{\pgfqpoint{0.967377in}{1.980288in}}%
\pgfpathlineto{\pgfqpoint{0.968020in}{1.981487in}}%
\pgfpathlineto{\pgfqpoint{0.968217in}{1.981980in}}%
\pgfpathlineto{\pgfqpoint{0.968514in}{1.980336in}}%
\pgfpathlineto{\pgfqpoint{0.968712in}{1.860988in}}%
\pgfpathlineto{\pgfqpoint{0.969526in}{1.532029in}}%
\pgfpathlineto{\pgfqpoint{0.969793in}{1.532034in}}%
\pgfpathlineto{\pgfqpoint{0.971406in}{1.532136in}}%
\pgfpathlineto{\pgfqpoint{0.971530in}{2.004758in}}%
\pgfpathlineto{\pgfqpoint{0.972690in}{1.980829in}}%
\pgfpathlineto{\pgfqpoint{0.972839in}{1.981302in}}%
\pgfpathlineto{\pgfqpoint{0.973333in}{1.980034in}}%
\pgfpathlineto{\pgfqpoint{0.973383in}{1.980433in}}%
\pgfpathlineto{\pgfqpoint{0.973482in}{1.979179in}}%
\pgfpathlineto{\pgfqpoint{0.973828in}{1.688386in}}%
\pgfpathlineto{\pgfqpoint{0.974483in}{1.532029in}}%
\pgfpathlineto{\pgfqpoint{0.974355in}{1.732451in}}%
\pgfpathlineto{\pgfqpoint{0.974753in}{1.532034in}}%
\pgfpathlineto{\pgfqpoint{0.976351in}{1.532127in}}%
\pgfpathlineto{\pgfqpoint{0.976481in}{2.004298in}}%
\pgfpathlineto{\pgfqpoint{0.977731in}{1.979946in}}%
\pgfpathlineto{\pgfqpoint{0.977780in}{1.980181in}}%
\pgfpathlineto{\pgfqpoint{0.978275in}{1.979467in}}%
\pgfpathlineto{\pgfqpoint{0.978324in}{1.979466in}}%
\pgfpathlineto{\pgfqpoint{0.978423in}{1.978745in}}%
\pgfpathlineto{\pgfqpoint{0.978671in}{1.801901in}}%
\pgfpathlineto{\pgfqpoint{0.979387in}{1.532028in}}%
\pgfpathlineto{\pgfqpoint{0.979682in}{1.532049in}}%
\pgfpathlineto{\pgfqpoint{0.981284in}{1.532082in}}%
\pgfpathlineto{\pgfqpoint{0.981304in}{1.533921in}}%
\pgfpathlineto{\pgfqpoint{0.981431in}{2.003710in}}%
\pgfpathlineto{\pgfqpoint{0.982322in}{1.975644in}}%
\pgfpathlineto{\pgfqpoint{0.982619in}{1.984259in}}%
\pgfpathlineto{\pgfqpoint{0.982570in}{1.974472in}}%
\pgfpathlineto{\pgfqpoint{0.983064in}{1.982090in}}%
\pgfpathlineto{\pgfqpoint{0.983559in}{1.870510in}}%
\pgfpathlineto{\pgfqpoint{0.984332in}{1.532027in}}%
\pgfpathlineto{\pgfqpoint{0.984656in}{1.532033in}}%
\pgfpathlineto{\pgfqpoint{0.986248in}{1.533026in}}%
\pgfpathlineto{\pgfqpoint{0.986366in}{2.001758in}}%
\pgfpathlineto{\pgfqpoint{0.987284in}{1.979702in}}%
\pgfpathlineto{\pgfqpoint{0.987334in}{1.977392in}}%
\pgfpathlineto{\pgfqpoint{0.988076in}{1.978429in}}%
\pgfpathlineto{\pgfqpoint{0.988372in}{1.967623in}}%
\pgfpathlineto{\pgfqpoint{0.989276in}{1.532028in}}%
\pgfpathlineto{\pgfqpoint{0.989976in}{1.532042in}}%
\pgfpathlineto{\pgfqpoint{0.991186in}{1.532116in}}%
\pgfpathlineto{\pgfqpoint{0.991226in}{1.717342in}}%
\pgfpathlineto{\pgfqpoint{0.991312in}{2.000923in}}%
\pgfpathlineto{\pgfqpoint{0.991980in}{1.973138in}}%
\pgfpathlineto{\pgfqpoint{0.992376in}{1.981126in}}%
\pgfpathlineto{\pgfqpoint{0.992722in}{1.979647in}}%
\pgfpathlineto{\pgfqpoint{0.992821in}{1.979674in}}%
\pgfpathlineto{\pgfqpoint{0.993464in}{1.865591in}}%
\pgfpathlineto{\pgfqpoint{0.994250in}{1.532033in}}%
\pgfpathlineto{\pgfqpoint{0.994547in}{1.532044in}}%
\pgfpathlineto{\pgfqpoint{0.996132in}{1.532123in}}%
\pgfpathlineto{\pgfqpoint{0.996255in}{2.000756in}}%
\pgfpathlineto{\pgfqpoint{0.997680in}{1.976771in}}%
\pgfpathlineto{\pgfqpoint{0.997878in}{1.977081in}}%
\pgfpathlineto{\pgfqpoint{0.997927in}{1.976324in}}%
\pgfpathlineto{\pgfqpoint{0.998125in}{1.976246in}}%
\pgfpathlineto{\pgfqpoint{0.998224in}{1.975849in}}%
\pgfpathlineto{\pgfqpoint{0.998471in}{1.805806in}}%
\pgfpathlineto{\pgfqpoint{0.999179in}{1.532033in}}%
\pgfpathlineto{\pgfqpoint{0.999475in}{1.532043in}}%
\pgfpathlineto{\pgfqpoint{1.001074in}{1.532103in}}%
\pgfpathlineto{\pgfqpoint{1.001083in}{1.532319in}}%
\pgfpathlineto{\pgfqpoint{1.001203in}{1.999947in}}%
\pgfpathlineto{\pgfqpoint{1.002170in}{1.976094in}}%
\pgfpathlineto{\pgfqpoint{1.002318in}{1.977004in}}%
\pgfpathlineto{\pgfqpoint{1.002368in}{1.975493in}}%
\pgfpathlineto{\pgfqpoint{1.002862in}{1.976358in}}%
\pgfpathlineto{\pgfqpoint{1.003209in}{1.970878in}}%
\pgfpathlineto{\pgfqpoint{1.004135in}{1.532035in}}%
\pgfpathlineto{\pgfqpoint{1.005053in}{1.532055in}}%
\pgfpathlineto{\pgfqpoint{1.006025in}{1.532149in}}%
\pgfpathlineto{\pgfqpoint{1.006149in}{1.999304in}}%
\pgfpathlineto{\pgfqpoint{1.007298in}{1.975639in}}%
\pgfpathlineto{\pgfqpoint{1.007397in}{1.975668in}}%
\pgfpathlineto{\pgfqpoint{1.008139in}{1.973837in}}%
\pgfpathlineto{\pgfqpoint{1.008485in}{1.691811in}}%
\pgfpathlineto{\pgfqpoint{1.009096in}{1.532036in}}%
\pgfpathlineto{\pgfqpoint{1.008973in}{1.698906in}}%
\pgfpathlineto{\pgfqpoint{1.009441in}{1.532043in}}%
\pgfpathlineto{\pgfqpoint{1.010974in}{1.532315in}}%
\pgfpathlineto{\pgfqpoint{1.011094in}{1.998315in}}%
\pgfpathlineto{\pgfqpoint{1.012022in}{1.974289in}}%
\pgfpathlineto{\pgfqpoint{1.012319in}{1.975459in}}%
\pgfpathlineto{\pgfqpoint{1.012270in}{1.974104in}}%
\pgfpathlineto{\pgfqpoint{1.012764in}{1.974911in}}%
\pgfpathlineto{\pgfqpoint{1.013110in}{1.969785in}}%
\pgfpathlineto{\pgfqpoint{1.014051in}{1.532030in}}%
\pgfpathlineto{\pgfqpoint{1.014947in}{1.532061in}}%
\pgfpathlineto{\pgfqpoint{1.015905in}{1.532087in}}%
\pgfpathlineto{\pgfqpoint{1.015916in}{1.532151in}}%
\pgfpathlineto{\pgfqpoint{1.016048in}{1.998335in}}%
\pgfpathlineto{\pgfqpoint{1.017099in}{1.973019in}}%
\pgfpathlineto{\pgfqpoint{1.017149in}{1.974929in}}%
\pgfpathlineto{\pgfqpoint{1.017841in}{1.973932in}}%
\pgfpathlineto{\pgfqpoint{1.018088in}{1.960337in}}%
\pgfpathlineto{\pgfqpoint{1.018994in}{1.532031in}}%
\pgfpathlineto{\pgfqpoint{1.019542in}{1.532043in}}%
\pgfpathlineto{\pgfqpoint{1.020861in}{1.532144in}}%
\pgfpathlineto{\pgfqpoint{1.020990in}{1.997547in}}%
\pgfpathlineto{\pgfqpoint{1.022064in}{1.972761in}}%
\pgfpathlineto{\pgfqpoint{1.022113in}{1.974110in}}%
\pgfpathlineto{\pgfqpoint{1.022163in}{1.972561in}}%
\pgfpathlineto{\pgfqpoint{1.022806in}{1.972740in}}%
\pgfpathlineto{\pgfqpoint{1.022954in}{1.972070in}}%
\pgfpathlineto{\pgfqpoint{1.023053in}{1.952859in}}%
\pgfpathlineto{\pgfqpoint{1.023922in}{1.532033in}}%
\pgfpathlineto{\pgfqpoint{1.024454in}{1.532045in}}%
\pgfpathlineto{\pgfqpoint{1.025810in}{1.532304in}}%
\pgfpathlineto{\pgfqpoint{1.025935in}{1.996871in}}%
\pgfpathlineto{\pgfqpoint{1.026880in}{1.973406in}}%
\pgfpathlineto{\pgfqpoint{1.027622in}{1.972065in}}%
\pgfpathlineto{\pgfqpoint{1.027671in}{1.972606in}}%
\pgfpathlineto{\pgfqpoint{1.027968in}{1.964410in}}%
\pgfpathlineto{\pgfqpoint{1.028856in}{1.532035in}}%
\pgfpathlineto{\pgfqpoint{1.029724in}{1.532055in}}%
\pgfpathlineto{\pgfqpoint{1.030751in}{1.532135in}}%
\pgfpathlineto{\pgfqpoint{1.030870in}{1.994572in}}%
\pgfpathlineto{\pgfqpoint{1.032051in}{1.972037in}}%
\pgfpathlineto{\pgfqpoint{1.032150in}{1.972076in}}%
\pgfpathlineto{\pgfqpoint{1.032892in}{1.970187in}}%
\pgfpathlineto{\pgfqpoint{1.033238in}{1.690494in}}%
\pgfpathlineto{\pgfqpoint{1.033829in}{1.532034in}}%
\pgfpathlineto{\pgfqpoint{1.034181in}{1.532042in}}%
\pgfpathlineto{\pgfqpoint{1.035697in}{1.532135in}}%
\pgfpathlineto{\pgfqpoint{1.035811in}{1.992933in}}%
\pgfpathlineto{\pgfqpoint{1.037021in}{1.971649in}}%
\pgfpathlineto{\pgfqpoint{1.037861in}{1.965221in}}%
\pgfpathlineto{\pgfqpoint{1.038758in}{1.532031in}}%
\pgfpathlineto{\pgfqpoint{1.039734in}{1.532060in}}%
\pgfpathlineto{\pgfqpoint{1.040642in}{1.532135in}}%
\pgfpathlineto{\pgfqpoint{1.040774in}{1.994835in}}%
\pgfpathlineto{\pgfqpoint{1.041914in}{1.970424in}}%
\pgfpathlineto{\pgfqpoint{1.042063in}{1.971085in}}%
\pgfpathlineto{\pgfqpoint{1.042557in}{1.969898in}}%
\pgfpathlineto{\pgfqpoint{1.042607in}{1.970344in}}%
\pgfpathlineto{\pgfqpoint{1.042804in}{1.965989in}}%
\pgfpathlineto{\pgfqpoint{1.043712in}{1.532031in}}%
\pgfpathlineto{\pgfqpoint{1.044850in}{1.532056in}}%
\pgfpathlineto{\pgfqpoint{1.045589in}{1.532146in}}%
\pgfpathlineto{\pgfqpoint{1.045810in}{1.994080in}}%
\pgfpathlineto{\pgfqpoint{1.046796in}{1.970318in}}%
\pgfpathlineto{\pgfqpoint{1.047044in}{1.970635in}}%
\pgfpathlineto{\pgfqpoint{1.047736in}{1.968550in}}%
\pgfpathlineto{\pgfqpoint{1.048082in}{1.697140in}}%
\pgfpathlineto{\pgfqpoint{1.048618in}{1.532031in}}%
\pgfpathlineto{\pgfqpoint{1.049027in}{1.532039in}}%
\pgfpathlineto{\pgfqpoint{1.050535in}{1.532156in}}%
\pgfpathlineto{\pgfqpoint{1.050666in}{1.993460in}}%
\pgfpathlineto{\pgfqpoint{1.051691in}{1.968904in}}%
\pgfpathlineto{\pgfqpoint{1.051988in}{1.970002in}}%
\pgfpathlineto{\pgfqpoint{1.051939in}{1.968503in}}%
\pgfpathlineto{\pgfqpoint{1.052433in}{1.969369in}}%
\pgfpathlineto{\pgfqpoint{1.052730in}{1.957294in}}%
\pgfpathlineto{\pgfqpoint{1.053570in}{1.532030in}}%
\pgfpathlineto{\pgfqpoint{1.054330in}{1.532056in}}%
\pgfpathlineto{\pgfqpoint{1.055478in}{1.532122in}}%
\pgfpathlineto{\pgfqpoint{1.055518in}{1.731027in}}%
\pgfpathlineto{\pgfqpoint{1.055799in}{1.997466in}}%
\pgfpathlineto{\pgfqpoint{1.056298in}{1.966992in}}%
\pgfpathlineto{\pgfqpoint{1.056694in}{1.969843in}}%
\pgfpathlineto{\pgfqpoint{1.057040in}{1.969271in}}%
\pgfpathlineto{\pgfqpoint{1.057733in}{1.930461in}}%
\pgfpathlineto{\pgfqpoint{1.058522in}{1.532030in}}%
\pgfpathlineto{\pgfqpoint{1.058995in}{1.532050in}}%
\pgfpathlineto{\pgfqpoint{1.060423in}{1.532122in}}%
\pgfpathlineto{\pgfqpoint{1.060463in}{1.726948in}}%
\pgfpathlineto{\pgfqpoint{1.060644in}{1.995176in}}%
\pgfpathlineto{\pgfqpoint{1.061235in}{1.967938in}}%
\pgfpathlineto{\pgfqpoint{1.061482in}{1.966686in}}%
\pgfpathlineto{\pgfqpoint{1.061432in}{1.969414in}}%
\pgfpathlineto{\pgfqpoint{1.061927in}{1.967123in}}%
\pgfpathlineto{\pgfqpoint{1.061976in}{1.968613in}}%
\pgfpathlineto{\pgfqpoint{1.062570in}{1.966867in}}%
\pgfpathlineto{\pgfqpoint{1.062817in}{1.810964in}}%
\pgfpathlineto{\pgfqpoint{1.063472in}{1.532029in}}%
\pgfpathlineto{\pgfqpoint{1.063848in}{1.532048in}}%
\pgfpathlineto{\pgfqpoint{1.065369in}{1.532124in}}%
\pgfpathlineto{\pgfqpoint{1.065408in}{1.723952in}}%
\pgfpathlineto{\pgfqpoint{1.065588in}{1.993816in}}%
\pgfpathlineto{\pgfqpoint{1.066175in}{1.967135in}}%
\pgfpathlineto{\pgfqpoint{1.066422in}{1.966101in}}%
\pgfpathlineto{\pgfqpoint{1.066372in}{1.968688in}}%
\pgfpathlineto{\pgfqpoint{1.066768in}{1.966631in}}%
\pgfpathlineto{\pgfqpoint{1.066818in}{1.968162in}}%
\pgfpathlineto{\pgfqpoint{1.067510in}{1.966289in}}%
\pgfpathlineto{\pgfqpoint{1.067708in}{1.868667in}}%
\pgfpathlineto{\pgfqpoint{1.068426in}{1.532029in}}%
\pgfpathlineto{\pgfqpoint{1.068808in}{1.532048in}}%
\pgfpathlineto{\pgfqpoint{1.070314in}{1.532122in}}%
\pgfpathlineto{\pgfqpoint{1.070354in}{1.721026in}}%
\pgfpathlineto{\pgfqpoint{1.070443in}{1.990678in}}%
\pgfpathlineto{\pgfqpoint{1.071107in}{1.968327in}}%
\pgfpathlineto{\pgfqpoint{1.071156in}{1.965030in}}%
\pgfpathlineto{\pgfqpoint{1.071898in}{1.967004in}}%
\pgfpathlineto{\pgfqpoint{1.071997in}{1.967240in}}%
\pgfpathlineto{\pgfqpoint{1.072541in}{1.952274in}}%
\pgfpathlineto{\pgfqpoint{1.073375in}{1.532029in}}%
\pgfpathlineto{\pgfqpoint{1.074009in}{1.532042in}}%
\pgfpathlineto{\pgfqpoint{1.075262in}{1.532169in}}%
\pgfpathlineto{\pgfqpoint{1.075489in}{1.993146in}}%
\pgfpathlineto{\pgfqpoint{1.076391in}{1.965128in}}%
\pgfpathlineto{\pgfqpoint{1.076539in}{1.968218in}}%
\pgfpathlineto{\pgfqpoint{1.076490in}{1.964159in}}%
\pgfpathlineto{\pgfqpoint{1.077133in}{1.965205in}}%
\pgfpathlineto{\pgfqpoint{1.077331in}{1.966091in}}%
\pgfpathlineto{\pgfqpoint{1.077380in}{1.964461in}}%
\pgfpathlineto{\pgfqpoint{1.077429in}{1.965539in}}%
\pgfpathlineto{\pgfqpoint{1.077726in}{1.748828in}}%
\pgfpathlineto{\pgfqpoint{1.078293in}{1.532029in}}%
\pgfpathlineto{\pgfqpoint{1.078706in}{1.532037in}}%
\pgfpathlineto{\pgfqpoint{1.080205in}{1.532122in}}%
\pgfpathlineto{\pgfqpoint{1.080245in}{1.729320in}}%
\pgfpathlineto{\pgfqpoint{1.080427in}{1.993956in}}%
\pgfpathlineto{\pgfqpoint{1.080989in}{1.967272in}}%
\pgfpathlineto{\pgfqpoint{1.081038in}{1.963708in}}%
\pgfpathlineto{\pgfqpoint{1.081731in}{1.965261in}}%
\pgfpathlineto{\pgfqpoint{1.081879in}{1.966091in}}%
\pgfpathlineto{\pgfqpoint{1.082374in}{1.963872in}}%
\pgfpathlineto{\pgfqpoint{1.082621in}{1.805591in}}%
\pgfpathlineto{\pgfqpoint{1.083236in}{1.532030in}}%
\pgfpathlineto{\pgfqpoint{1.083648in}{1.532038in}}%
\pgfpathlineto{\pgfqpoint{1.085150in}{1.532122in}}%
\pgfpathlineto{\pgfqpoint{1.085190in}{1.726013in}}%
\pgfpathlineto{\pgfqpoint{1.085371in}{1.992299in}}%
\pgfpathlineto{\pgfqpoint{1.085943in}{1.965970in}}%
\pgfpathlineto{\pgfqpoint{1.086091in}{1.963471in}}%
\pgfpathlineto{\pgfqpoint{1.086042in}{1.966437in}}%
\pgfpathlineto{\pgfqpoint{1.086734in}{1.964860in}}%
\pgfpathlineto{\pgfqpoint{1.086932in}{1.965301in}}%
\pgfpathlineto{\pgfqpoint{1.087328in}{1.963220in}}%
\pgfpathlineto{\pgfqpoint{1.087575in}{1.799435in}}%
\pgfpathlineto{\pgfqpoint{1.088179in}{1.532030in}}%
\pgfpathlineto{\pgfqpoint{1.088590in}{1.532038in}}%
\pgfpathlineto{\pgfqpoint{1.090099in}{1.532166in}}%
\pgfpathlineto{\pgfqpoint{1.090319in}{1.989221in}}%
\pgfpathlineto{\pgfqpoint{1.091222in}{1.964937in}}%
\pgfpathlineto{\pgfqpoint{1.091518in}{1.963338in}}%
\pgfpathlineto{\pgfqpoint{1.091469in}{1.965336in}}%
\pgfpathlineto{\pgfqpoint{1.092062in}{1.963524in}}%
\pgfpathlineto{\pgfqpoint{1.092112in}{1.964007in}}%
\pgfpathlineto{\pgfqpoint{1.092211in}{1.963377in}}%
\pgfpathlineto{\pgfqpoint{1.092260in}{1.963498in}}%
\pgfpathlineto{\pgfqpoint{1.092458in}{1.862668in}}%
\pgfpathlineto{\pgfqpoint{1.093123in}{1.532031in}}%
\pgfpathlineto{\pgfqpoint{1.093532in}{1.532038in}}%
\pgfpathlineto{\pgfqpoint{1.095044in}{1.532154in}}%
\pgfpathlineto{\pgfqpoint{1.095265in}{1.988775in}}%
\pgfpathlineto{\pgfqpoint{1.096192in}{1.963930in}}%
\pgfpathlineto{\pgfqpoint{1.096340in}{1.962932in}}%
\pgfpathlineto{\pgfqpoint{1.096390in}{1.964599in}}%
\pgfpathlineto{\pgfqpoint{1.096983in}{1.963141in}}%
\pgfpathlineto{\pgfqpoint{1.097033in}{1.963342in}}%
\pgfpathlineto{\pgfqpoint{1.097231in}{1.961960in}}%
\pgfpathlineto{\pgfqpoint{1.097527in}{1.744912in}}%
\pgfpathlineto{\pgfqpoint{1.098115in}{1.532032in}}%
\pgfpathlineto{\pgfqpoint{1.098519in}{1.532046in}}%
\pgfpathlineto{\pgfqpoint{1.099987in}{1.532129in}}%
\pgfpathlineto{\pgfqpoint{1.100211in}{1.987556in}}%
\pgfpathlineto{\pgfqpoint{1.101263in}{1.962775in}}%
\pgfpathlineto{\pgfqpoint{1.101313in}{1.963669in}}%
\pgfpathlineto{\pgfqpoint{1.101906in}{1.962618in}}%
\pgfpathlineto{\pgfqpoint{1.102005in}{1.962628in}}%
\pgfpathlineto{\pgfqpoint{1.102154in}{1.961929in}}%
\pgfpathlineto{\pgfqpoint{1.102302in}{1.907661in}}%
\pgfpathlineto{\pgfqpoint{1.103019in}{1.532029in}}%
\pgfpathlineto{\pgfqpoint{1.103483in}{1.532049in}}%
\pgfpathlineto{\pgfqpoint{1.104932in}{1.532123in}}%
\pgfpathlineto{\pgfqpoint{1.104972in}{1.723881in}}%
\pgfpathlineto{\pgfqpoint{1.105153in}{1.989573in}}%
\pgfpathlineto{\pgfqpoint{1.105752in}{1.960921in}}%
\pgfpathlineto{\pgfqpoint{1.105801in}{1.963865in}}%
\pgfpathlineto{\pgfqpoint{1.106493in}{1.963159in}}%
\pgfpathlineto{\pgfqpoint{1.106592in}{1.963201in}}%
\pgfpathlineto{\pgfqpoint{1.107186in}{1.949385in}}%
\pgfpathlineto{\pgfqpoint{1.107959in}{1.532030in}}%
\pgfpathlineto{\pgfqpoint{1.108666in}{1.532045in}}%
\pgfpathlineto{\pgfqpoint{1.109880in}{1.532151in}}%
\pgfpathlineto{\pgfqpoint{1.110102in}{1.986679in}}%
\pgfpathlineto{\pgfqpoint{1.111068in}{1.961095in}}%
\pgfpathlineto{\pgfqpoint{1.111117in}{1.963120in}}%
\pgfpathlineto{\pgfqpoint{1.111167in}{1.961046in}}%
\pgfpathlineto{\pgfqpoint{1.111809in}{1.961685in}}%
\pgfpathlineto{\pgfqpoint{1.111908in}{1.961817in}}%
\pgfpathlineto{\pgfqpoint{1.112057in}{1.960539in}}%
\pgfpathlineto{\pgfqpoint{1.112255in}{1.860486in}}%
\pgfpathlineto{\pgfqpoint{1.112906in}{1.532030in}}%
\pgfpathlineto{\pgfqpoint{1.113367in}{1.532048in}}%
\pgfpathlineto{\pgfqpoint{1.114826in}{1.532159in}}%
\pgfpathlineto{\pgfqpoint{1.115048in}{1.986894in}}%
\pgfpathlineto{\pgfqpoint{1.115987in}{1.960981in}}%
\pgfpathlineto{\pgfqpoint{1.116284in}{1.962263in}}%
\pgfpathlineto{\pgfqpoint{1.116234in}{1.960680in}}%
\pgfpathlineto{\pgfqpoint{1.116729in}{1.961657in}}%
\pgfpathlineto{\pgfqpoint{1.117075in}{1.950341in}}%
\pgfpathlineto{\pgfqpoint{1.117898in}{1.532031in}}%
\pgfpathlineto{\pgfqpoint{1.118600in}{1.532053in}}%
\pgfpathlineto{\pgfqpoint{1.119770in}{1.532132in}}%
\pgfpathlineto{\pgfqpoint{1.119995in}{1.985170in}}%
\pgfpathlineto{\pgfqpoint{1.121045in}{1.960355in}}%
\pgfpathlineto{\pgfqpoint{1.121094in}{1.961524in}}%
\pgfpathlineto{\pgfqpoint{1.121787in}{1.960553in}}%
\pgfpathlineto{\pgfqpoint{1.121984in}{1.959231in}}%
\pgfpathlineto{\pgfqpoint{1.122380in}{1.653279in}}%
\pgfpathlineto{\pgfqpoint{1.122838in}{1.532030in}}%
\pgfpathlineto{\pgfqpoint{1.123339in}{1.532049in}}%
\pgfpathlineto{\pgfqpoint{1.124716in}{1.532151in}}%
\pgfpathlineto{\pgfqpoint{1.124939in}{1.985784in}}%
\pgfpathlineto{\pgfqpoint{1.125869in}{1.960599in}}%
\pgfpathlineto{\pgfqpoint{1.126017in}{1.959582in}}%
\pgfpathlineto{\pgfqpoint{1.126067in}{1.961239in}}%
\pgfpathlineto{\pgfqpoint{1.126660in}{1.959826in}}%
\pgfpathlineto{\pgfqpoint{1.126710in}{1.960007in}}%
\pgfpathlineto{\pgfqpoint{1.126908in}{1.958918in}}%
\pgfpathlineto{\pgfqpoint{1.127105in}{1.857781in}}%
\pgfpathlineto{\pgfqpoint{1.127778in}{1.532031in}}%
\pgfpathlineto{\pgfqpoint{1.128173in}{1.532046in}}%
\pgfpathlineto{\pgfqpoint{1.129661in}{1.532142in}}%
\pgfpathlineto{\pgfqpoint{1.129886in}{1.984631in}}%
\pgfpathlineto{\pgfqpoint{1.130879in}{1.960340in}}%
\pgfpathlineto{\pgfqpoint{1.131126in}{1.960489in}}%
\pgfpathlineto{\pgfqpoint{1.131720in}{1.959032in}}%
\pgfpathlineto{\pgfqpoint{1.131769in}{1.959327in}}%
\pgfpathlineto{\pgfqpoint{1.131868in}{1.958392in}}%
\pgfpathlineto{\pgfqpoint{1.132116in}{1.797742in}}%
\pgfpathlineto{\pgfqpoint{1.132721in}{1.532031in}}%
\pgfpathlineto{\pgfqpoint{1.133113in}{1.532046in}}%
\pgfpathlineto{\pgfqpoint{1.134606in}{1.532135in}}%
\pgfpathlineto{\pgfqpoint{1.134828in}{1.983920in}}%
\pgfpathlineto{\pgfqpoint{1.135860in}{1.958964in}}%
\pgfpathlineto{\pgfqpoint{1.136157in}{1.959749in}}%
\pgfpathlineto{\pgfqpoint{1.136553in}{1.958695in}}%
\pgfpathlineto{\pgfqpoint{1.136602in}{1.959189in}}%
\pgfpathlineto{\pgfqpoint{1.136849in}{1.954041in}}%
\pgfpathlineto{\pgfqpoint{1.137664in}{1.532031in}}%
\pgfpathlineto{\pgfqpoint{1.138695in}{1.532054in}}%
\pgfpathlineto{\pgfqpoint{1.139550in}{1.532119in}}%
\pgfpathlineto{\pgfqpoint{1.139671in}{1.981760in}}%
\pgfpathlineto{\pgfqpoint{1.141143in}{1.958086in}}%
\pgfpathlineto{\pgfqpoint{1.141192in}{1.959363in}}%
\pgfpathlineto{\pgfqpoint{1.141786in}{1.956824in}}%
\pgfpathlineto{\pgfqpoint{1.142592in}{1.532032in}}%
\pgfpathlineto{\pgfqpoint{1.143958in}{1.532066in}}%
\pgfpathlineto{\pgfqpoint{1.144493in}{1.532101in}}%
\pgfpathlineto{\pgfqpoint{1.144501in}{1.532309in}}%
\pgfpathlineto{\pgfqpoint{1.144617in}{1.981053in}}%
\pgfpathlineto{\pgfqpoint{1.145566in}{1.958244in}}%
\pgfpathlineto{\pgfqpoint{1.145715in}{1.958943in}}%
\pgfpathlineto{\pgfqpoint{1.145764in}{1.957696in}}%
\pgfpathlineto{\pgfqpoint{1.146259in}{1.958497in}}%
\pgfpathlineto{\pgfqpoint{1.146753in}{1.951077in}}%
\pgfpathlineto{\pgfqpoint{1.147534in}{1.532032in}}%
\pgfpathlineto{\pgfqpoint{1.148500in}{1.532056in}}%
\pgfpathlineto{\pgfqpoint{1.149439in}{1.532105in}}%
\pgfpathlineto{\pgfqpoint{1.149448in}{1.532434in}}%
\pgfpathlineto{\pgfqpoint{1.149565in}{1.980901in}}%
\pgfpathlineto{\pgfqpoint{1.150492in}{1.957303in}}%
\pgfpathlineto{\pgfqpoint{1.150789in}{1.958573in}}%
\pgfpathlineto{\pgfqpoint{1.150739in}{1.957052in}}%
\pgfpathlineto{\pgfqpoint{1.151234in}{1.958084in}}%
\pgfpathlineto{\pgfqpoint{1.151728in}{1.942942in}}%
\pgfpathlineto{\pgfqpoint{1.152470in}{1.532033in}}%
\pgfpathlineto{\pgfqpoint{1.153181in}{1.532048in}}%
\pgfpathlineto{\pgfqpoint{1.154391in}{1.532195in}}%
\pgfpathlineto{\pgfqpoint{1.154509in}{1.980368in}}%
\pgfpathlineto{\pgfqpoint{1.155528in}{1.957771in}}%
\pgfpathlineto{\pgfqpoint{1.156369in}{1.956709in}}%
\pgfpathlineto{\pgfqpoint{1.156418in}{1.956999in}}%
\pgfpathlineto{\pgfqpoint{1.156616in}{1.955804in}}%
\pgfpathlineto{\pgfqpoint{1.156913in}{1.746639in}}%
\pgfpathlineto{\pgfqpoint{1.157412in}{1.532034in}}%
\pgfpathlineto{\pgfqpoint{1.157874in}{1.532043in}}%
\pgfpathlineto{\pgfqpoint{1.159338in}{1.532304in}}%
\pgfpathlineto{\pgfqpoint{1.159453in}{1.979500in}}%
\pgfpathlineto{\pgfqpoint{1.160404in}{1.956803in}}%
\pgfpathlineto{\pgfqpoint{1.160651in}{1.957384in}}%
\pgfpathlineto{\pgfqpoint{1.160602in}{1.956224in}}%
\pgfpathlineto{\pgfqpoint{1.161096in}{1.956993in}}%
\pgfpathlineto{\pgfqpoint{1.161591in}{1.951944in}}%
\pgfpathlineto{\pgfqpoint{1.162394in}{1.532034in}}%
\pgfpathlineto{\pgfqpoint{1.163526in}{1.532058in}}%
\pgfpathlineto{\pgfqpoint{1.164283in}{1.532304in}}%
\pgfpathlineto{\pgfqpoint{1.164399in}{1.979055in}}%
\pgfpathlineto{\pgfqpoint{1.165351in}{1.956404in}}%
\pgfpathlineto{\pgfqpoint{1.165598in}{1.956902in}}%
\pgfpathlineto{\pgfqpoint{1.165549in}{1.955754in}}%
\pgfpathlineto{\pgfqpoint{1.165945in}{1.956551in}}%
\pgfpathlineto{\pgfqpoint{1.166538in}{1.951608in}}%
\pgfpathlineto{\pgfqpoint{1.167311in}{1.532033in}}%
\pgfpathlineto{\pgfqpoint{1.168470in}{1.532059in}}%
\pgfpathlineto{\pgfqpoint{1.169229in}{1.532308in}}%
\pgfpathlineto{\pgfqpoint{1.169345in}{1.978602in}}%
\pgfpathlineto{\pgfqpoint{1.170259in}{1.955257in}}%
\pgfpathlineto{\pgfqpoint{1.170308in}{1.956564in}}%
\pgfpathlineto{\pgfqpoint{1.170358in}{1.955112in}}%
\pgfpathlineto{\pgfqpoint{1.171001in}{1.955590in}}%
\pgfpathlineto{\pgfqpoint{1.171396in}{1.955132in}}%
\pgfpathlineto{\pgfqpoint{1.171495in}{1.948263in}}%
\pgfpathlineto{\pgfqpoint{1.172282in}{1.532033in}}%
\pgfpathlineto{\pgfqpoint{1.173316in}{1.532056in}}%
\pgfpathlineto{\pgfqpoint{1.174174in}{1.532304in}}%
\pgfpathlineto{\pgfqpoint{1.174290in}{1.977994in}}%
\pgfpathlineto{\pgfqpoint{1.175237in}{1.954963in}}%
\pgfpathlineto{\pgfqpoint{1.175386in}{1.956123in}}%
\pgfpathlineto{\pgfqpoint{1.175336in}{1.954649in}}%
\pgfpathlineto{\pgfqpoint{1.175979in}{1.954970in}}%
\pgfpathlineto{\pgfqpoint{1.176177in}{1.955193in}}%
\pgfpathlineto{\pgfqpoint{1.176325in}{1.954360in}}%
\pgfpathlineto{\pgfqpoint{1.176375in}{1.954546in}}%
\pgfpathlineto{\pgfqpoint{1.176474in}{1.938989in}}%
\pgfpathlineto{\pgfqpoint{1.177200in}{1.532033in}}%
\pgfpathlineto{\pgfqpoint{1.177914in}{1.532048in}}%
\pgfpathlineto{\pgfqpoint{1.179120in}{1.532307in}}%
\pgfpathlineto{\pgfqpoint{1.179236in}{1.977615in}}%
\pgfpathlineto{\pgfqpoint{1.180186in}{1.954795in}}%
\pgfpathlineto{\pgfqpoint{1.180334in}{1.955544in}}%
\pgfpathlineto{\pgfqpoint{1.180384in}{1.954254in}}%
\pgfpathlineto{\pgfqpoint{1.180878in}{1.955093in}}%
\pgfpathlineto{\pgfqpoint{1.181422in}{1.939054in}}%
\pgfpathlineto{\pgfqpoint{1.182154in}{1.532033in}}%
\pgfpathlineto{\pgfqpoint{1.182874in}{1.532049in}}%
\pgfpathlineto{\pgfqpoint{1.184065in}{1.532307in}}%
\pgfpathlineto{\pgfqpoint{1.184181in}{1.977133in}}%
\pgfpathlineto{\pgfqpoint{1.185136in}{1.954588in}}%
\pgfpathlineto{\pgfqpoint{1.186224in}{1.953440in}}%
\pgfpathlineto{\pgfqpoint{1.185284in}{1.954931in}}%
\pgfpathlineto{\pgfqpoint{1.186274in}{1.953639in}}%
\pgfpathlineto{\pgfqpoint{1.186372in}{1.937663in}}%
\pgfpathlineto{\pgfqpoint{1.187097in}{1.532033in}}%
\pgfpathlineto{\pgfqpoint{1.187815in}{1.532049in}}%
\pgfpathlineto{\pgfqpoint{1.189011in}{1.532306in}}%
\pgfpathlineto{\pgfqpoint{1.189127in}{1.976674in}}%
\pgfpathlineto{\pgfqpoint{1.190041in}{1.953341in}}%
\pgfpathlineto{\pgfqpoint{1.190338in}{1.954505in}}%
\pgfpathlineto{\pgfqpoint{1.190783in}{1.954134in}}%
\pgfpathlineto{\pgfqpoint{1.191278in}{1.949928in}}%
\pgfpathlineto{\pgfqpoint{1.192074in}{1.532033in}}%
\pgfpathlineto{\pgfqpoint{1.193310in}{1.532060in}}%
\pgfpathlineto{\pgfqpoint{1.193951in}{1.532130in}}%
\pgfpathlineto{\pgfqpoint{1.194072in}{1.976334in}}%
\pgfpathlineto{\pgfqpoint{1.195356in}{1.953142in}}%
\pgfpathlineto{\pgfqpoint{1.195653in}{1.953978in}}%
\pgfpathlineto{\pgfqpoint{1.196048in}{1.952586in}}%
\pgfpathlineto{\pgfqpoint{1.196098in}{1.953187in}}%
\pgfpathlineto{\pgfqpoint{1.196246in}{1.944985in}}%
\pgfpathlineto{\pgfqpoint{1.196993in}{1.532032in}}%
\pgfpathlineto{\pgfqpoint{1.197912in}{1.532052in}}%
\pgfpathlineto{\pgfqpoint{1.198896in}{1.532121in}}%
\pgfpathlineto{\pgfqpoint{1.199108in}{1.977074in}}%
\pgfpathlineto{\pgfqpoint{1.200249in}{1.953252in}}%
\pgfpathlineto{\pgfqpoint{1.201188in}{1.944794in}}%
\pgfpathlineto{\pgfqpoint{1.201928in}{1.532026in}}%
\pgfpathlineto{\pgfqpoint{1.202889in}{1.532061in}}%
\pgfpathlineto{\pgfqpoint{1.203840in}{1.532113in}}%
\pgfpathlineto{\pgfqpoint{1.203882in}{1.731149in}}%
\pgfpathlineto{\pgfqpoint{1.204056in}{1.976720in}}%
\pgfpathlineto{\pgfqpoint{1.204653in}{1.954300in}}%
\pgfpathlineto{\pgfqpoint{1.204703in}{1.950745in}}%
\pgfpathlineto{\pgfqpoint{1.205395in}{1.952478in}}%
\pgfpathlineto{\pgfqpoint{1.205543in}{1.953203in}}%
\pgfpathlineto{\pgfqpoint{1.206038in}{1.951351in}}%
\pgfpathlineto{\pgfqpoint{1.206087in}{1.951829in}}%
\pgfpathlineto{\pgfqpoint{1.206285in}{1.853725in}}%
\pgfpathlineto{\pgfqpoint{1.206873in}{1.532025in}}%
\pgfpathlineto{\pgfqpoint{1.207390in}{1.532036in}}%
\pgfpathlineto{\pgfqpoint{1.208785in}{1.532109in}}%
\pgfpathlineto{\pgfqpoint{1.208826in}{1.709606in}}%
\pgfpathlineto{\pgfqpoint{1.208910in}{1.974848in}}%
\pgfpathlineto{\pgfqpoint{1.209612in}{1.955448in}}%
\pgfpathlineto{\pgfqpoint{1.209909in}{1.946283in}}%
\pgfpathlineto{\pgfqpoint{1.209859in}{1.958283in}}%
\pgfpathlineto{\pgfqpoint{1.210354in}{1.948467in}}%
\pgfpathlineto{\pgfqpoint{1.210403in}{1.955176in}}%
\pgfpathlineto{\pgfqpoint{1.211096in}{1.944731in}}%
\pgfpathlineto{\pgfqpoint{1.211819in}{1.532026in}}%
\pgfpathlineto{\pgfqpoint{1.212632in}{1.532043in}}%
\pgfpathlineto{\pgfqpoint{1.213731in}{1.532110in}}%
\pgfpathlineto{\pgfqpoint{1.213771in}{1.713443in}}%
\pgfpathlineto{\pgfqpoint{1.213857in}{1.974823in}}%
\pgfpathlineto{\pgfqpoint{1.214510in}{1.952798in}}%
\pgfpathlineto{\pgfqpoint{1.214806in}{1.949895in}}%
\pgfpathlineto{\pgfqpoint{1.214757in}{1.953615in}}%
\pgfpathlineto{\pgfqpoint{1.215251in}{1.950560in}}%
\pgfpathlineto{\pgfqpoint{1.215301in}{1.952669in}}%
\pgfpathlineto{\pgfqpoint{1.215993in}{1.950804in}}%
\pgfpathlineto{\pgfqpoint{1.216290in}{1.752174in}}%
\pgfpathlineto{\pgfqpoint{1.216774in}{1.532030in}}%
\pgfpathlineto{\pgfqpoint{1.217297in}{1.532041in}}%
\pgfpathlineto{\pgfqpoint{1.218678in}{1.532121in}}%
\pgfpathlineto{\pgfqpoint{1.218802in}{1.973829in}}%
\pgfpathlineto{\pgfqpoint{1.220233in}{1.951690in}}%
\pgfpathlineto{\pgfqpoint{1.220975in}{1.944931in}}%
\pgfpathlineto{\pgfqpoint{1.221711in}{1.532025in}}%
\pgfpathlineto{\pgfqpoint{1.222821in}{1.532049in}}%
\pgfpathlineto{\pgfqpoint{1.223630in}{1.532558in}}%
\pgfpathlineto{\pgfqpoint{1.223743in}{1.972660in}}%
\pgfpathlineto{\pgfqpoint{1.224655in}{1.947782in}}%
\pgfpathlineto{\pgfqpoint{1.224704in}{1.954080in}}%
\pgfpathlineto{\pgfqpoint{1.225397in}{1.951025in}}%
\pgfpathlineto{\pgfqpoint{1.225595in}{1.951790in}}%
\pgfpathlineto{\pgfqpoint{1.225941in}{1.942200in}}%
\pgfpathlineto{\pgfqpoint{1.226668in}{1.532029in}}%
\pgfpathlineto{\pgfqpoint{1.227492in}{1.532047in}}%
\pgfpathlineto{\pgfqpoint{1.228569in}{1.532119in}}%
\pgfpathlineto{\pgfqpoint{1.228609in}{1.738597in}}%
\pgfpathlineto{\pgfqpoint{1.228792in}{1.974899in}}%
\pgfpathlineto{\pgfqpoint{1.229383in}{1.951867in}}%
\pgfpathlineto{\pgfqpoint{1.229432in}{1.948767in}}%
\pgfpathlineto{\pgfqpoint{1.229482in}{1.952084in}}%
\pgfpathlineto{\pgfqpoint{1.230174in}{1.950633in}}%
\pgfpathlineto{\pgfqpoint{1.230273in}{1.950957in}}%
\pgfpathlineto{\pgfqpoint{1.230916in}{1.929774in}}%
\pgfpathlineto{\pgfqpoint{1.231621in}{1.532028in}}%
\pgfpathlineto{\pgfqpoint{1.232303in}{1.532053in}}%
\pgfpathlineto{\pgfqpoint{1.233514in}{1.532119in}}%
\pgfpathlineto{\pgfqpoint{1.233554in}{1.726835in}}%
\pgfpathlineto{\pgfqpoint{1.233734in}{1.977267in}}%
\pgfpathlineto{\pgfqpoint{1.234318in}{1.950662in}}%
\pgfpathlineto{\pgfqpoint{1.234467in}{1.948591in}}%
\pgfpathlineto{\pgfqpoint{1.234417in}{1.951389in}}%
\pgfpathlineto{\pgfqpoint{1.235110in}{1.949832in}}%
\pgfpathlineto{\pgfqpoint{1.235307in}{1.950438in}}%
\pgfpathlineto{\pgfqpoint{1.235357in}{1.949292in}}%
\pgfpathlineto{\pgfqpoint{1.235654in}{1.949648in}}%
\pgfpathlineto{\pgfqpoint{1.235802in}{1.947862in}}%
\pgfpathlineto{\pgfqpoint{1.236198in}{1.657800in}}%
\pgfpathlineto{\pgfqpoint{1.236537in}{1.532025in}}%
\pgfpathlineto{\pgfqpoint{1.237145in}{1.532038in}}%
\pgfpathlineto{\pgfqpoint{1.238458in}{1.532108in}}%
\pgfpathlineto{\pgfqpoint{1.238486in}{1.577786in}}%
\pgfpathlineto{\pgfqpoint{1.238579in}{1.972105in}}%
\pgfpathlineto{\pgfqpoint{1.239290in}{1.948150in}}%
\pgfpathlineto{\pgfqpoint{1.239586in}{1.951956in}}%
\pgfpathlineto{\pgfqpoint{1.239537in}{1.947317in}}%
\pgfpathlineto{\pgfqpoint{1.240031in}{1.951030in}}%
\pgfpathlineto{\pgfqpoint{1.240872in}{1.890363in}}%
\pgfpathlineto{\pgfqpoint{1.241519in}{1.532027in}}%
\pgfpathlineto{\pgfqpoint{1.242059in}{1.532039in}}%
\pgfpathlineto{\pgfqpoint{1.243408in}{1.532173in}}%
\pgfpathlineto{\pgfqpoint{1.243634in}{1.977154in}}%
\pgfpathlineto{\pgfqpoint{1.244534in}{1.950273in}}%
\pgfpathlineto{\pgfqpoint{1.244830in}{1.948049in}}%
\pgfpathlineto{\pgfqpoint{1.244781in}{1.950338in}}%
\pgfpathlineto{\pgfqpoint{1.245325in}{1.949642in}}%
\pgfpathlineto{\pgfqpoint{1.245770in}{1.921673in}}%
\pgfpathlineto{\pgfqpoint{1.246456in}{1.532028in}}%
\pgfpathlineto{\pgfqpoint{1.247137in}{1.532052in}}%
\pgfpathlineto{\pgfqpoint{1.248350in}{1.532119in}}%
\pgfpathlineto{\pgfqpoint{1.248391in}{1.729799in}}%
\pgfpathlineto{\pgfqpoint{1.248572in}{1.975122in}}%
\pgfpathlineto{\pgfqpoint{1.249145in}{1.950332in}}%
\pgfpathlineto{\pgfqpoint{1.249194in}{1.947082in}}%
\pgfpathlineto{\pgfqpoint{1.249244in}{1.950390in}}%
\pgfpathlineto{\pgfqpoint{1.249886in}{1.948744in}}%
\pgfpathlineto{\pgfqpoint{1.250035in}{1.949372in}}%
\pgfpathlineto{\pgfqpoint{1.250084in}{1.948049in}}%
\pgfpathlineto{\pgfqpoint{1.250480in}{1.948666in}}%
\pgfpathlineto{\pgfqpoint{1.250678in}{1.939110in}}%
\pgfpathlineto{\pgfqpoint{1.251411in}{1.532027in}}%
\pgfpathlineto{\pgfqpoint{1.252347in}{1.532048in}}%
\pgfpathlineto{\pgfqpoint{1.253298in}{1.532154in}}%
\pgfpathlineto{\pgfqpoint{1.253531in}{1.975462in}}%
\pgfpathlineto{\pgfqpoint{1.254426in}{1.949895in}}%
\pgfpathlineto{\pgfqpoint{1.254822in}{1.947340in}}%
\pgfpathlineto{\pgfqpoint{1.255218in}{1.948848in}}%
\pgfpathlineto{\pgfqpoint{1.255663in}{1.921889in}}%
\pgfpathlineto{\pgfqpoint{1.256350in}{1.532028in}}%
\pgfpathlineto{\pgfqpoint{1.257033in}{1.532053in}}%
\pgfpathlineto{\pgfqpoint{1.258241in}{1.532119in}}%
\pgfpathlineto{\pgfqpoint{1.258281in}{1.726359in}}%
\pgfpathlineto{\pgfqpoint{1.258464in}{1.975209in}}%
\pgfpathlineto{\pgfqpoint{1.259066in}{1.946286in}}%
\pgfpathlineto{\pgfqpoint{1.259363in}{1.949435in}}%
\pgfpathlineto{\pgfqpoint{1.259808in}{1.948873in}}%
\pgfpathlineto{\pgfqpoint{1.260648in}{1.902286in}}%
\pgfpathlineto{\pgfqpoint{1.261209in}{1.532024in}}%
\pgfpathlineto{\pgfqpoint{1.261914in}{1.532054in}}%
\pgfpathlineto{\pgfqpoint{1.263185in}{1.532107in}}%
\pgfpathlineto{\pgfqpoint{1.263213in}{1.577300in}}%
\pgfpathlineto{\pgfqpoint{1.263304in}{1.969858in}}%
\pgfpathlineto{\pgfqpoint{1.263989in}{1.947768in}}%
\pgfpathlineto{\pgfqpoint{1.264237in}{1.949667in}}%
\pgfpathlineto{\pgfqpoint{1.264187in}{1.945557in}}%
\pgfpathlineto{\pgfqpoint{1.264583in}{1.948860in}}%
\pgfpathlineto{\pgfqpoint{1.264632in}{1.946337in}}%
\pgfpathlineto{\pgfqpoint{1.264682in}{1.948887in}}%
\pgfpathlineto{\pgfqpoint{1.265374in}{1.947194in}}%
\pgfpathlineto{\pgfqpoint{1.265522in}{1.937254in}}%
\pgfpathlineto{\pgfqpoint{1.266154in}{1.532024in}}%
\pgfpathlineto{\pgfqpoint{1.267161in}{1.532061in}}%
\pgfpathlineto{\pgfqpoint{1.268130in}{1.532106in}}%
\pgfpathlineto{\pgfqpoint{1.268158in}{1.575122in}}%
\pgfpathlineto{\pgfqpoint{1.268260in}{1.970134in}}%
\pgfpathlineto{\pgfqpoint{1.268926in}{1.941029in}}%
\pgfpathlineto{\pgfqpoint{1.268975in}{1.952691in}}%
\pgfpathlineto{\pgfqpoint{1.269668in}{1.948808in}}%
\pgfpathlineto{\pgfqpoint{1.269767in}{1.949639in}}%
\pgfpathlineto{\pgfqpoint{1.270607in}{1.853945in}}%
\pgfpathlineto{\pgfqpoint{1.271161in}{1.532025in}}%
\pgfpathlineto{\pgfqpoint{1.271724in}{1.532052in}}%
\pgfpathlineto{\pgfqpoint{1.273077in}{1.532111in}}%
\pgfpathlineto{\pgfqpoint{1.273118in}{1.726096in}}%
\pgfpathlineto{\pgfqpoint{1.273290in}{1.969665in}}%
\pgfpathlineto{\pgfqpoint{1.273873in}{1.949646in}}%
\pgfpathlineto{\pgfqpoint{1.273922in}{1.944612in}}%
\pgfpathlineto{\pgfqpoint{1.274615in}{1.945825in}}%
\pgfpathlineto{\pgfqpoint{1.274664in}{1.948057in}}%
\pgfpathlineto{\pgfqpoint{1.274713in}{1.945620in}}%
\pgfpathlineto{\pgfqpoint{1.275356in}{1.945740in}}%
\pgfpathlineto{\pgfqpoint{1.275604in}{1.808875in}}%
\pgfpathlineto{\pgfqpoint{1.276087in}{1.532026in}}%
\pgfpathlineto{\pgfqpoint{1.276631in}{1.532050in}}%
\pgfpathlineto{\pgfqpoint{1.278013in}{1.532082in}}%
\pgfpathlineto{\pgfqpoint{1.278025in}{1.532136in}}%
\pgfpathlineto{\pgfqpoint{1.278251in}{1.971327in}}%
\pgfpathlineto{\pgfqpoint{1.279223in}{1.946622in}}%
\pgfpathlineto{\pgfqpoint{1.279470in}{1.947128in}}%
\pgfpathlineto{\pgfqpoint{1.279421in}{1.945887in}}%
\pgfpathlineto{\pgfqpoint{1.279816in}{1.946721in}}%
\pgfpathlineto{\pgfqpoint{1.280360in}{1.937833in}}%
\pgfpathlineto{\pgfqpoint{1.281049in}{1.532026in}}%
\pgfpathlineto{\pgfqpoint{1.282054in}{1.532047in}}%
\pgfpathlineto{\pgfqpoint{1.282969in}{1.532120in}}%
\pgfpathlineto{\pgfqpoint{1.283153in}{1.967505in}}%
\pgfpathlineto{\pgfqpoint{1.284348in}{1.946272in}}%
\pgfpathlineto{\pgfqpoint{1.285287in}{1.943290in}}%
\pgfpathlineto{\pgfqpoint{1.285930in}{1.541996in}}%
\pgfpathlineto{\pgfqpoint{1.286002in}{1.532030in}}%
\pgfpathlineto{\pgfqpoint{1.286667in}{1.532050in}}%
\pgfpathlineto{\pgfqpoint{1.287914in}{1.532119in}}%
\pgfpathlineto{\pgfqpoint{1.288115in}{1.969559in}}%
\pgfpathlineto{\pgfqpoint{1.289299in}{1.945911in}}%
\pgfpathlineto{\pgfqpoint{1.290238in}{1.941927in}}%
\pgfpathlineto{\pgfqpoint{1.290944in}{1.532030in}}%
\pgfpathlineto{\pgfqpoint{1.292299in}{1.532065in}}%
\pgfpathlineto{\pgfqpoint{1.292859in}{1.532114in}}%
\pgfpathlineto{\pgfqpoint{1.292901in}{1.744356in}}%
\pgfpathlineto{\pgfqpoint{1.292987in}{1.967806in}}%
\pgfpathlineto{\pgfqpoint{1.293665in}{1.942578in}}%
\pgfpathlineto{\pgfqpoint{1.293714in}{1.948196in}}%
\pgfpathlineto{\pgfqpoint{1.294407in}{1.945718in}}%
\pgfpathlineto{\pgfqpoint{1.294605in}{1.946412in}}%
\pgfpathlineto{\pgfqpoint{1.295232in}{1.929176in}}%
\pgfpathlineto{\pgfqpoint{1.295915in}{1.532028in}}%
\pgfpathlineto{\pgfqpoint{1.296746in}{1.532047in}}%
\pgfpathlineto{\pgfqpoint{1.297807in}{1.532150in}}%
\pgfpathlineto{\pgfqpoint{1.298038in}{1.972678in}}%
\pgfpathlineto{\pgfqpoint{1.298970in}{1.944828in}}%
\pgfpathlineto{\pgfqpoint{1.299168in}{1.945951in}}%
\pgfpathlineto{\pgfqpoint{1.299217in}{1.944168in}}%
\pgfpathlineto{\pgfqpoint{1.299613in}{1.945465in}}%
\pgfpathlineto{\pgfqpoint{1.300157in}{1.934821in}}%
\pgfpathlineto{\pgfqpoint{1.300840in}{1.532029in}}%
\pgfpathlineto{\pgfqpoint{1.301803in}{1.532056in}}%
\pgfpathlineto{\pgfqpoint{1.302751in}{1.532124in}}%
\pgfpathlineto{\pgfqpoint{1.302938in}{1.966089in}}%
\pgfpathlineto{\pgfqpoint{1.304091in}{1.944652in}}%
\pgfpathlineto{\pgfqpoint{1.304388in}{1.944788in}}%
\pgfpathlineto{\pgfqpoint{1.304734in}{1.944492in}}%
\pgfpathlineto{\pgfqpoint{1.305080in}{1.941658in}}%
\pgfpathlineto{\pgfqpoint{1.305811in}{1.532027in}}%
\pgfpathlineto{\pgfqpoint{1.307192in}{1.532069in}}%
\pgfpathlineto{\pgfqpoint{1.307687in}{1.532083in}}%
\pgfpathlineto{\pgfqpoint{1.307699in}{1.532167in}}%
\pgfpathlineto{\pgfqpoint{1.307924in}{1.971482in}}%
\pgfpathlineto{\pgfqpoint{1.308786in}{1.943074in}}%
\pgfpathlineto{\pgfqpoint{1.308835in}{1.945825in}}%
\pgfpathlineto{\pgfqpoint{1.308885in}{1.942951in}}%
\pgfpathlineto{\pgfqpoint{1.309528in}{1.944174in}}%
\pgfpathlineto{\pgfqpoint{1.309726in}{1.944581in}}%
\pgfpathlineto{\pgfqpoint{1.310022in}{1.941377in}}%
\pgfpathlineto{\pgfqpoint{1.310418in}{1.658972in}}%
\pgfpathlineto{\pgfqpoint{1.310664in}{1.532023in}}%
\pgfpathlineto{\pgfqpoint{1.311393in}{1.532054in}}%
\pgfpathlineto{\pgfqpoint{1.312639in}{1.532103in}}%
\pgfpathlineto{\pgfqpoint{1.312667in}{1.575500in}}%
\pgfpathlineto{\pgfqpoint{1.312751in}{1.964702in}}%
\pgfpathlineto{\pgfqpoint{1.313446in}{1.949953in}}%
\pgfpathlineto{\pgfqpoint{1.313743in}{1.939280in}}%
\pgfpathlineto{\pgfqpoint{1.314188in}{1.940991in}}%
\pgfpathlineto{\pgfqpoint{1.314238in}{1.947085in}}%
\pgfpathlineto{\pgfqpoint{1.314930in}{1.943312in}}%
\pgfpathlineto{\pgfqpoint{1.315177in}{1.816225in}}%
\pgfpathlineto{\pgfqpoint{1.315609in}{1.532017in}}%
\pgfpathlineto{\pgfqpoint{1.316244in}{1.532055in}}%
\pgfpathlineto{\pgfqpoint{1.317584in}{1.532099in}}%
\pgfpathlineto{\pgfqpoint{1.317592in}{1.532239in}}%
\pgfpathlineto{\pgfqpoint{1.317710in}{1.966317in}}%
\pgfpathlineto{\pgfqpoint{1.318677in}{1.943731in}}%
\pgfpathlineto{\pgfqpoint{1.319864in}{1.942678in}}%
\pgfpathlineto{\pgfqpoint{1.319963in}{1.929828in}}%
\pgfpathlineto{\pgfqpoint{1.320554in}{1.532022in}}%
\pgfpathlineto{\pgfqpoint{1.321470in}{1.532059in}}%
\pgfpathlineto{\pgfqpoint{1.322530in}{1.532101in}}%
\pgfpathlineto{\pgfqpoint{1.322539in}{1.532468in}}%
\pgfpathlineto{\pgfqpoint{1.322652in}{1.965653in}}%
\pgfpathlineto{\pgfqpoint{1.323583in}{1.944246in}}%
\pgfpathlineto{\pgfqpoint{1.323633in}{1.942627in}}%
\pgfpathlineto{\pgfqpoint{1.324375in}{1.943555in}}%
\pgfpathlineto{\pgfqpoint{1.324919in}{1.927284in}}%
\pgfpathlineto{\pgfqpoint{1.325544in}{1.532026in}}%
\pgfpathlineto{\pgfqpoint{1.326386in}{1.532056in}}%
\pgfpathlineto{\pgfqpoint{1.327468in}{1.532081in}}%
\pgfpathlineto{\pgfqpoint{1.327478in}{1.532126in}}%
\pgfpathlineto{\pgfqpoint{1.327705in}{1.967833in}}%
\pgfpathlineto{\pgfqpoint{1.328776in}{1.943325in}}%
\pgfpathlineto{\pgfqpoint{1.328875in}{1.943380in}}%
\pgfpathlineto{\pgfqpoint{1.329617in}{1.942412in}}%
\pgfpathlineto{\pgfqpoint{1.329667in}{1.942516in}}%
\pgfpathlineto{\pgfqpoint{1.329716in}{1.942105in}}%
\pgfpathlineto{\pgfqpoint{1.329766in}{1.942130in}}%
\pgfpathlineto{\pgfqpoint{1.329864in}{1.926326in}}%
\pgfpathlineto{\pgfqpoint{1.330490in}{1.532026in}}%
\pgfpathlineto{\pgfqpoint{1.331332in}{1.532057in}}%
\pgfpathlineto{\pgfqpoint{1.332424in}{1.532124in}}%
\pgfpathlineto{\pgfqpoint{1.332646in}{1.965641in}}%
\pgfpathlineto{\pgfqpoint{1.333746in}{1.942760in}}%
\pgfpathlineto{\pgfqpoint{1.334685in}{1.941869in}}%
\pgfpathlineto{\pgfqpoint{1.334784in}{1.935235in}}%
\pgfpathlineto{\pgfqpoint{1.335487in}{1.532028in}}%
\pgfpathlineto{\pgfqpoint{1.336518in}{1.532051in}}%
\pgfpathlineto{\pgfqpoint{1.337370in}{1.532138in}}%
\pgfpathlineto{\pgfqpoint{1.337600in}{1.968288in}}%
\pgfpathlineto{\pgfqpoint{1.338572in}{1.941953in}}%
\pgfpathlineto{\pgfqpoint{1.338721in}{1.943248in}}%
\pgfpathlineto{\pgfqpoint{1.338671in}{1.941520in}}%
\pgfpathlineto{\pgfqpoint{1.339314in}{1.941843in}}%
\pgfpathlineto{\pgfqpoint{1.339364in}{1.942150in}}%
\pgfpathlineto{\pgfqpoint{1.339710in}{1.939492in}}%
\pgfpathlineto{\pgfqpoint{1.340324in}{1.569572in}}%
\pgfpathlineto{\pgfqpoint{1.340427in}{1.532028in}}%
\pgfpathlineto{\pgfqpoint{1.341059in}{1.532042in}}%
\pgfpathlineto{\pgfqpoint{1.342317in}{1.532154in}}%
\pgfpathlineto{\pgfqpoint{1.342545in}{1.967136in}}%
\pgfpathlineto{\pgfqpoint{1.343437in}{1.943133in}}%
\pgfpathlineto{\pgfqpoint{1.343733in}{1.941057in}}%
\pgfpathlineto{\pgfqpoint{1.344228in}{1.942414in}}%
\pgfpathlineto{\pgfqpoint{1.344722in}{1.918013in}}%
\pgfpathlineto{\pgfqpoint{1.345358in}{1.532032in}}%
\pgfpathlineto{\pgfqpoint{1.346127in}{1.532049in}}%
\pgfpathlineto{\pgfqpoint{1.347263in}{1.532182in}}%
\pgfpathlineto{\pgfqpoint{1.347370in}{1.962784in}}%
\pgfpathlineto{\pgfqpoint{1.348360in}{1.939118in}}%
\pgfpathlineto{\pgfqpoint{1.348409in}{1.944340in}}%
\pgfpathlineto{\pgfqpoint{1.349101in}{1.941868in}}%
\pgfpathlineto{\pgfqpoint{1.349200in}{1.942397in}}%
\pgfpathlineto{\pgfqpoint{1.349645in}{1.929980in}}%
\pgfpathlineto{\pgfqpoint{1.350293in}{1.532032in}}%
\pgfpathlineto{\pgfqpoint{1.351153in}{1.532051in}}%
\pgfpathlineto{\pgfqpoint{1.352207in}{1.532135in}}%
\pgfpathlineto{\pgfqpoint{1.352321in}{1.962724in}}%
\pgfpathlineto{\pgfqpoint{1.353435in}{1.941374in}}%
\pgfpathlineto{\pgfqpoint{1.353683in}{1.941515in}}%
\pgfpathlineto{\pgfqpoint{1.353534in}{1.941289in}}%
\pgfpathlineto{\pgfqpoint{1.354128in}{1.941253in}}%
\pgfpathlineto{\pgfqpoint{1.354523in}{1.940174in}}%
\pgfpathlineto{\pgfqpoint{1.354721in}{1.846766in}}%
\pgfpathlineto{\pgfqpoint{1.355265in}{1.532029in}}%
\pgfpathlineto{\pgfqpoint{1.355852in}{1.532048in}}%
\pgfpathlineto{\pgfqpoint{1.357152in}{1.532136in}}%
\pgfpathlineto{\pgfqpoint{1.357385in}{1.967357in}}%
\pgfpathlineto{\pgfqpoint{1.358362in}{1.940717in}}%
\pgfpathlineto{\pgfqpoint{1.358510in}{1.941884in}}%
\pgfpathlineto{\pgfqpoint{1.358461in}{1.940333in}}%
\pgfpathlineto{\pgfqpoint{1.359104in}{1.940601in}}%
\pgfpathlineto{\pgfqpoint{1.359153in}{1.940865in}}%
\pgfpathlineto{\pgfqpoint{1.359450in}{1.939810in}}%
\pgfpathlineto{\pgfqpoint{1.359549in}{1.925338in}}%
\pgfpathlineto{\pgfqpoint{1.360211in}{1.532029in}}%
\pgfpathlineto{\pgfqpoint{1.360996in}{1.532053in}}%
\pgfpathlineto{\pgfqpoint{1.362098in}{1.532138in}}%
\pgfpathlineto{\pgfqpoint{1.362425in}{1.965044in}}%
\pgfpathlineto{\pgfqpoint{1.363283in}{1.941055in}}%
\pgfpathlineto{\pgfqpoint{1.363531in}{1.941523in}}%
\pgfpathlineto{\pgfqpoint{1.364272in}{1.940056in}}%
\pgfpathlineto{\pgfqpoint{1.364322in}{1.940346in}}%
\pgfpathlineto{\pgfqpoint{1.364371in}{1.939640in}}%
\pgfpathlineto{\pgfqpoint{1.364421in}{1.939836in}}%
\pgfpathlineto{\pgfqpoint{1.364668in}{1.801082in}}%
\pgfpathlineto{\pgfqpoint{1.365157in}{1.532029in}}%
\pgfpathlineto{\pgfqpoint{1.365695in}{1.532040in}}%
\pgfpathlineto{\pgfqpoint{1.367043in}{1.532142in}}%
\pgfpathlineto{\pgfqpoint{1.367274in}{1.966301in}}%
\pgfpathlineto{\pgfqpoint{1.368214in}{1.939774in}}%
\pgfpathlineto{\pgfqpoint{1.368510in}{1.941233in}}%
\pgfpathlineto{\pgfqpoint{1.368560in}{1.939719in}}%
\pgfpathlineto{\pgfqpoint{1.368955in}{1.940737in}}%
\pgfpathlineto{\pgfqpoint{1.369450in}{1.922040in}}%
\pgfpathlineto{\pgfqpoint{1.370102in}{1.532029in}}%
\pgfpathlineto{\pgfqpoint{1.370886in}{1.532053in}}%
\pgfpathlineto{\pgfqpoint{1.371988in}{1.532137in}}%
\pgfpathlineto{\pgfqpoint{1.372316in}{1.964336in}}%
\pgfpathlineto{\pgfqpoint{1.373188in}{1.940403in}}%
\pgfpathlineto{\pgfqpoint{1.374276in}{1.938949in}}%
\pgfpathlineto{\pgfqpoint{1.373435in}{1.940976in}}%
\pgfpathlineto{\pgfqpoint{1.374325in}{1.939072in}}%
\pgfpathlineto{\pgfqpoint{1.374622in}{1.748212in}}%
\pgfpathlineto{\pgfqpoint{1.375048in}{1.532029in}}%
\pgfpathlineto{\pgfqpoint{1.375634in}{1.532048in}}%
\pgfpathlineto{\pgfqpoint{1.376934in}{1.532140in}}%
\pgfpathlineto{\pgfqpoint{1.377165in}{1.965409in}}%
\pgfpathlineto{\pgfqpoint{1.378132in}{1.939234in}}%
\pgfpathlineto{\pgfqpoint{1.378280in}{1.940728in}}%
\pgfpathlineto{\pgfqpoint{1.378231in}{1.938974in}}%
\pgfpathlineto{\pgfqpoint{1.378873in}{1.939495in}}%
\pgfpathlineto{\pgfqpoint{1.379071in}{1.939723in}}%
\pgfpathlineto{\pgfqpoint{1.379220in}{1.938732in}}%
\pgfpathlineto{\pgfqpoint{1.379269in}{1.938687in}}%
\pgfpathlineto{\pgfqpoint{1.379516in}{1.796045in}}%
\pgfpathlineto{\pgfqpoint{1.379997in}{1.532029in}}%
\pgfpathlineto{\pgfqpoint{1.380586in}{1.532049in}}%
\pgfpathlineto{\pgfqpoint{1.381880in}{1.532141in}}%
\pgfpathlineto{\pgfqpoint{1.382109in}{1.963988in}}%
\pgfpathlineto{\pgfqpoint{1.383041in}{1.938599in}}%
\pgfpathlineto{\pgfqpoint{1.383338in}{1.940351in}}%
\pgfpathlineto{\pgfqpoint{1.383783in}{1.939866in}}%
\pgfpathlineto{\pgfqpoint{1.384278in}{1.926990in}}%
\pgfpathlineto{\pgfqpoint{1.384942in}{1.532029in}}%
\pgfpathlineto{\pgfqpoint{1.385827in}{1.532055in}}%
\pgfpathlineto{\pgfqpoint{1.386825in}{1.532139in}}%
\pgfpathlineto{\pgfqpoint{1.387055in}{1.964361in}}%
\pgfpathlineto{\pgfqpoint{1.387996in}{1.938291in}}%
\pgfpathlineto{\pgfqpoint{1.388045in}{1.940257in}}%
\pgfpathlineto{\pgfqpoint{1.388738in}{1.939213in}}%
\pgfpathlineto{\pgfqpoint{1.388836in}{1.939334in}}%
\pgfpathlineto{\pgfqpoint{1.389183in}{1.937583in}}%
\pgfpathlineto{\pgfqpoint{1.389778in}{1.580295in}}%
\pgfpathlineto{\pgfqpoint{1.389886in}{1.532029in}}%
\pgfpathlineto{\pgfqpoint{1.390522in}{1.532043in}}%
\pgfpathlineto{\pgfqpoint{1.391770in}{1.532132in}}%
\pgfpathlineto{\pgfqpoint{1.392100in}{1.963359in}}%
\pgfpathlineto{\pgfqpoint{1.392956in}{1.939049in}}%
\pgfpathlineto{\pgfqpoint{1.393203in}{1.939675in}}%
\pgfpathlineto{\pgfqpoint{1.393153in}{1.938281in}}%
\pgfpathlineto{\pgfqpoint{1.393549in}{1.939288in}}%
\pgfpathlineto{\pgfqpoint{1.394192in}{1.918148in}}%
\pgfpathlineto{\pgfqpoint{1.394832in}{1.532029in}}%
\pgfpathlineto{\pgfqpoint{1.395617in}{1.532053in}}%
\pgfpathlineto{\pgfqpoint{1.396716in}{1.532134in}}%
\pgfpathlineto{\pgfqpoint{1.396947in}{1.963685in}}%
\pgfpathlineto{\pgfqpoint{1.397928in}{1.939355in}}%
\pgfpathlineto{\pgfqpoint{1.398224in}{1.937961in}}%
\pgfpathlineto{\pgfqpoint{1.398175in}{1.939376in}}%
\pgfpathlineto{\pgfqpoint{1.398768in}{1.938025in}}%
\pgfpathlineto{\pgfqpoint{1.398818in}{1.938516in}}%
\pgfpathlineto{\pgfqpoint{1.399016in}{1.937559in}}%
\pgfpathlineto{\pgfqpoint{1.399065in}{1.937565in}}%
\pgfpathlineto{\pgfqpoint{1.399362in}{1.746961in}}%
\pgfpathlineto{\pgfqpoint{1.399778in}{1.532028in}}%
\pgfpathlineto{\pgfqpoint{1.400367in}{1.532048in}}%
\pgfpathlineto{\pgfqpoint{1.401662in}{1.532142in}}%
\pgfpathlineto{\pgfqpoint{1.401892in}{1.963325in}}%
\pgfpathlineto{\pgfqpoint{1.402815in}{1.938053in}}%
\pgfpathlineto{\pgfqpoint{1.403012in}{1.939238in}}%
\pgfpathlineto{\pgfqpoint{1.403062in}{1.937558in}}%
\pgfpathlineto{\pgfqpoint{1.403556in}{1.938724in}}%
\pgfpathlineto{\pgfqpoint{1.404051in}{1.930653in}}%
\pgfpathlineto{\pgfqpoint{1.404722in}{1.532029in}}%
\pgfpathlineto{\pgfqpoint{1.405754in}{1.532052in}}%
\pgfpathlineto{\pgfqpoint{1.406606in}{1.532134in}}%
\pgfpathlineto{\pgfqpoint{1.406837in}{1.962365in}}%
\pgfpathlineto{\pgfqpoint{1.407810in}{1.939050in}}%
\pgfpathlineto{\pgfqpoint{1.407859in}{1.937375in}}%
\pgfpathlineto{\pgfqpoint{1.408601in}{1.938204in}}%
\pgfpathlineto{\pgfqpoint{1.408997in}{1.929942in}}%
\pgfpathlineto{\pgfqpoint{1.409669in}{1.532029in}}%
\pgfpathlineto{\pgfqpoint{1.410800in}{1.532054in}}%
\pgfpathlineto{\pgfqpoint{1.411552in}{1.532139in}}%
\pgfpathlineto{\pgfqpoint{1.411783in}{1.963013in}}%
\pgfpathlineto{\pgfqpoint{1.412714in}{1.937190in}}%
\pgfpathlineto{\pgfqpoint{1.413011in}{1.938695in}}%
\pgfpathlineto{\pgfqpoint{1.412961in}{1.937060in}}%
\pgfpathlineto{\pgfqpoint{1.413456in}{1.938190in}}%
\pgfpathlineto{\pgfqpoint{1.413950in}{1.928870in}}%
\pgfpathlineto{\pgfqpoint{1.414613in}{1.532029in}}%
\pgfpathlineto{\pgfqpoint{1.415645in}{1.532052in}}%
\pgfpathlineto{\pgfqpoint{1.416497in}{1.532134in}}%
\pgfpathlineto{\pgfqpoint{1.416827in}{1.962127in}}%
\pgfpathlineto{\pgfqpoint{1.417694in}{1.938431in}}%
\pgfpathlineto{\pgfqpoint{1.418090in}{1.936962in}}%
\pgfpathlineto{\pgfqpoint{1.418485in}{1.937714in}}%
\pgfpathlineto{\pgfqpoint{1.418881in}{1.932090in}}%
\pgfpathlineto{\pgfqpoint{1.419559in}{1.532029in}}%
\pgfpathlineto{\pgfqpoint{1.420888in}{1.532058in}}%
\pgfpathlineto{\pgfqpoint{1.421443in}{1.532138in}}%
\pgfpathlineto{\pgfqpoint{1.421675in}{1.962557in}}%
\pgfpathlineto{\pgfqpoint{1.422605in}{1.936917in}}%
\pgfpathlineto{\pgfqpoint{1.422902in}{1.938141in}}%
\pgfpathlineto{\pgfqpoint{1.422852in}{1.936447in}}%
\pgfpathlineto{\pgfqpoint{1.423347in}{1.937637in}}%
\pgfpathlineto{\pgfqpoint{1.423841in}{1.929513in}}%
\pgfpathlineto{\pgfqpoint{1.424505in}{1.532029in}}%
\pgfpathlineto{\pgfqpoint{1.425586in}{1.532059in}}%
\pgfpathlineto{\pgfqpoint{1.426389in}{1.532137in}}%
\pgfpathlineto{\pgfqpoint{1.426718in}{1.961334in}}%
\pgfpathlineto{\pgfqpoint{1.427571in}{1.936310in}}%
\pgfpathlineto{\pgfqpoint{1.427620in}{1.937995in}}%
\pgfpathlineto{\pgfqpoint{1.427670in}{1.936078in}}%
\pgfpathlineto{\pgfqpoint{1.428313in}{1.936737in}}%
\pgfpathlineto{\pgfqpoint{1.428708in}{1.936238in}}%
\pgfpathlineto{\pgfqpoint{1.428807in}{1.923324in}}%
\pgfpathlineto{\pgfqpoint{1.429450in}{1.532029in}}%
\pgfpathlineto{\pgfqpoint{1.430333in}{1.532055in}}%
\pgfpathlineto{\pgfqpoint{1.431334in}{1.532137in}}%
\pgfpathlineto{\pgfqpoint{1.431565in}{1.961724in}}%
\pgfpathlineto{\pgfqpoint{1.432525in}{1.936235in}}%
\pgfpathlineto{\pgfqpoint{1.432673in}{1.937724in}}%
\pgfpathlineto{\pgfqpoint{1.432624in}{1.935872in}}%
\pgfpathlineto{\pgfqpoint{1.433267in}{1.936373in}}%
\pgfpathlineto{\pgfqpoint{1.433465in}{1.936690in}}%
\pgfpathlineto{\pgfqpoint{1.433613in}{1.935732in}}%
\pgfpathlineto{\pgfqpoint{1.433662in}{1.935979in}}%
\pgfpathlineto{\pgfqpoint{1.433761in}{1.921190in}}%
\pgfpathlineto{\pgfqpoint{1.434396in}{1.532029in}}%
\pgfpathlineto{\pgfqpoint{1.435229in}{1.532047in}}%
\pgfpathlineto{\pgfqpoint{1.436279in}{1.532136in}}%
\pgfpathlineto{\pgfqpoint{1.436510in}{1.960931in}}%
\pgfpathlineto{\pgfqpoint{1.437448in}{1.935609in}}%
\pgfpathlineto{\pgfqpoint{1.437498in}{1.937542in}}%
\pgfpathlineto{\pgfqpoint{1.437547in}{1.935581in}}%
\pgfpathlineto{\pgfqpoint{1.438190in}{1.936364in}}%
\pgfpathlineto{\pgfqpoint{1.438289in}{1.936554in}}%
\pgfpathlineto{\pgfqpoint{1.438635in}{1.935327in}}%
\pgfpathlineto{\pgfqpoint{1.438882in}{1.796839in}}%
\pgfpathlineto{\pgfqpoint{1.439346in}{1.532028in}}%
\pgfpathlineto{\pgfqpoint{1.439937in}{1.532049in}}%
\pgfpathlineto{\pgfqpoint{1.441225in}{1.532137in}}%
\pgfpathlineto{\pgfqpoint{1.441454in}{1.961051in}}%
\pgfpathlineto{\pgfqpoint{1.442389in}{1.935211in}}%
\pgfpathlineto{\pgfqpoint{1.442439in}{1.937166in}}%
\pgfpathlineto{\pgfqpoint{1.443131in}{1.936419in}}%
\pgfpathlineto{\pgfqpoint{1.443230in}{1.936434in}}%
\pgfpathlineto{\pgfqpoint{1.443626in}{1.928867in}}%
\pgfpathlineto{\pgfqpoint{1.444292in}{1.532027in}}%
\pgfpathlineto{\pgfqpoint{1.445478in}{1.532063in}}%
\pgfpathlineto{\pgfqpoint{1.446170in}{1.532138in}}%
\pgfpathlineto{\pgfqpoint{1.446400in}{1.960680in}}%
\pgfpathlineto{\pgfqpoint{1.447328in}{1.935513in}}%
\pgfpathlineto{\pgfqpoint{1.447625in}{1.936857in}}%
\pgfpathlineto{\pgfqpoint{1.447576in}{1.935155in}}%
\pgfpathlineto{\pgfqpoint{1.448070in}{1.936356in}}%
\pgfpathlineto{\pgfqpoint{1.448614in}{1.914913in}}%
\pgfpathlineto{\pgfqpoint{1.449237in}{1.532027in}}%
\pgfpathlineto{\pgfqpoint{1.450027in}{1.532054in}}%
\pgfpathlineto{\pgfqpoint{1.451116in}{1.532137in}}%
\pgfpathlineto{\pgfqpoint{1.451345in}{1.960432in}}%
\pgfpathlineto{\pgfqpoint{1.452280in}{1.934771in}}%
\pgfpathlineto{\pgfqpoint{1.452329in}{1.936749in}}%
\pgfpathlineto{\pgfqpoint{1.453022in}{1.935730in}}%
\pgfpathlineto{\pgfqpoint{1.453121in}{1.935857in}}%
\pgfpathlineto{\pgfqpoint{1.453516in}{1.929175in}}%
\pgfpathlineto{\pgfqpoint{1.454184in}{1.532027in}}%
\pgfpathlineto{\pgfqpoint{1.455419in}{1.532054in}}%
\pgfpathlineto{\pgfqpoint{1.456062in}{1.532141in}}%
\pgfpathlineto{\pgfqpoint{1.456291in}{1.960582in}}%
\pgfpathlineto{\pgfqpoint{1.457219in}{1.934766in}}%
\pgfpathlineto{\pgfqpoint{1.457516in}{1.936323in}}%
\pgfpathlineto{\pgfqpoint{1.457565in}{1.934737in}}%
\pgfpathlineto{\pgfqpoint{1.457961in}{1.935843in}}%
\pgfpathlineto{\pgfqpoint{1.458505in}{1.916397in}}%
\pgfpathlineto{\pgfqpoint{1.459129in}{1.532027in}}%
\pgfpathlineto{\pgfqpoint{1.459918in}{1.532054in}}%
\pgfpathlineto{\pgfqpoint{1.461007in}{1.532141in}}%
\pgfpathlineto{\pgfqpoint{1.461234in}{1.959697in}}%
\pgfpathlineto{\pgfqpoint{1.462172in}{1.934319in}}%
\pgfpathlineto{\pgfqpoint{1.462469in}{1.936041in}}%
\pgfpathlineto{\pgfqpoint{1.462914in}{1.935577in}}%
\pgfpathlineto{\pgfqpoint{1.463458in}{1.913426in}}%
\pgfpathlineto{\pgfqpoint{1.464072in}{1.532028in}}%
\pgfpathlineto{\pgfqpoint{1.464859in}{1.532053in}}%
\pgfpathlineto{\pgfqpoint{1.465953in}{1.532157in}}%
\pgfpathlineto{\pgfqpoint{1.466183in}{1.961687in}}%
\pgfpathlineto{\pgfqpoint{1.467082in}{1.936531in}}%
\pgfpathlineto{\pgfqpoint{1.467132in}{1.933775in}}%
\pgfpathlineto{\pgfqpoint{1.467874in}{1.935436in}}%
\pgfpathlineto{\pgfqpoint{1.468418in}{1.906413in}}%
\pgfpathlineto{\pgfqpoint{1.469018in}{1.532028in}}%
\pgfpathlineto{\pgfqpoint{1.469756in}{1.532044in}}%
\pgfpathlineto{\pgfqpoint{1.470898in}{1.532136in}}%
\pgfpathlineto{\pgfqpoint{1.471227in}{1.959085in}}%
\pgfpathlineto{\pgfqpoint{1.472079in}{1.933996in}}%
\pgfpathlineto{\pgfqpoint{1.472129in}{1.935748in}}%
\pgfpathlineto{\pgfqpoint{1.472178in}{1.933816in}}%
\pgfpathlineto{\pgfqpoint{1.472821in}{1.934519in}}%
\pgfpathlineto{\pgfqpoint{1.472920in}{1.934742in}}%
\pgfpathlineto{\pgfqpoint{1.473266in}{1.933521in}}%
\pgfpathlineto{\pgfqpoint{1.473513in}{1.796127in}}%
\pgfpathlineto{\pgfqpoint{1.473916in}{1.532026in}}%
\pgfpathlineto{\pgfqpoint{1.474561in}{1.532050in}}%
\pgfpathlineto{\pgfqpoint{1.475832in}{1.532082in}}%
\pgfpathlineto{\pgfqpoint{1.475845in}{1.532165in}}%
\pgfpathlineto{\pgfqpoint{1.476069in}{1.959353in}}%
\pgfpathlineto{\pgfqpoint{1.476963in}{1.938157in}}%
\pgfpathlineto{\pgfqpoint{1.477012in}{1.931165in}}%
\pgfpathlineto{\pgfqpoint{1.477705in}{1.933673in}}%
\pgfpathlineto{\pgfqpoint{1.477853in}{1.935692in}}%
\pgfpathlineto{\pgfqpoint{1.478249in}{1.928399in}}%
\pgfpathlineto{\pgfqpoint{1.478806in}{1.536283in}}%
\pgfpathlineto{\pgfqpoint{1.478863in}{1.532026in}}%
\pgfpathlineto{\pgfqpoint{1.479557in}{1.532040in}}%
\pgfpathlineto{\pgfqpoint{1.480789in}{1.532147in}}%
\pgfpathlineto{\pgfqpoint{1.481120in}{1.959416in}}%
\pgfpathlineto{\pgfqpoint{1.481902in}{1.934874in}}%
\pgfpathlineto{\pgfqpoint{1.482198in}{1.933185in}}%
\pgfpathlineto{\pgfqpoint{1.482149in}{1.935530in}}%
\pgfpathlineto{\pgfqpoint{1.482643in}{1.933475in}}%
\pgfpathlineto{\pgfqpoint{1.482693in}{1.934704in}}%
\pgfpathlineto{\pgfqpoint{1.483187in}{1.931684in}}%
\pgfpathlineto{\pgfqpoint{1.483749in}{1.543508in}}%
\pgfpathlineto{\pgfqpoint{1.483854in}{1.532028in}}%
\pgfpathlineto{\pgfqpoint{1.484492in}{1.532042in}}%
\pgfpathlineto{\pgfqpoint{1.485734in}{1.532135in}}%
\pgfpathlineto{\pgfqpoint{1.486063in}{1.958370in}}%
\pgfpathlineto{\pgfqpoint{1.486898in}{1.933098in}}%
\pgfpathlineto{\pgfqpoint{1.486947in}{1.935110in}}%
\pgfpathlineto{\pgfqpoint{1.487640in}{1.934011in}}%
\pgfpathlineto{\pgfqpoint{1.487738in}{1.934176in}}%
\pgfpathlineto{\pgfqpoint{1.488134in}{1.930449in}}%
\pgfpathlineto{\pgfqpoint{1.488687in}{1.584648in}}%
\pgfpathlineto{\pgfqpoint{1.488753in}{1.532026in}}%
\pgfpathlineto{\pgfqpoint{1.489447in}{1.532040in}}%
\pgfpathlineto{\pgfqpoint{1.490680in}{1.532153in}}%
\pgfpathlineto{\pgfqpoint{1.490910in}{1.958496in}}%
\pgfpathlineto{\pgfqpoint{1.491803in}{1.935033in}}%
\pgfpathlineto{\pgfqpoint{1.492099in}{1.932698in}}%
\pgfpathlineto{\pgfqpoint{1.492545in}{1.932978in}}%
\pgfpathlineto{\pgfqpoint{1.492594in}{1.934378in}}%
\pgfpathlineto{\pgfqpoint{1.493089in}{1.929536in}}%
\pgfpathlineto{\pgfqpoint{1.493747in}{1.532027in}}%
\pgfpathlineto{\pgfqpoint{1.495078in}{1.532057in}}%
\pgfpathlineto{\pgfqpoint{1.495625in}{1.532141in}}%
\pgfpathlineto{\pgfqpoint{1.495848in}{1.957454in}}%
\pgfpathlineto{\pgfqpoint{1.496797in}{1.932661in}}%
\pgfpathlineto{\pgfqpoint{1.496847in}{1.934600in}}%
\pgfpathlineto{\pgfqpoint{1.497539in}{1.933605in}}%
\pgfpathlineto{\pgfqpoint{1.497638in}{1.933729in}}%
\pgfpathlineto{\pgfqpoint{1.498034in}{1.928699in}}%
\pgfpathlineto{\pgfqpoint{1.498588in}{1.536302in}}%
\pgfpathlineto{\pgfqpoint{1.498645in}{1.532026in}}%
\pgfpathlineto{\pgfqpoint{1.499340in}{1.532040in}}%
\pgfpathlineto{\pgfqpoint{1.500571in}{1.532139in}}%
\pgfpathlineto{\pgfqpoint{1.500904in}{1.958762in}}%
\pgfpathlineto{\pgfqpoint{1.501729in}{1.933967in}}%
\pgfpathlineto{\pgfqpoint{1.501878in}{1.932113in}}%
\pgfpathlineto{\pgfqpoint{1.501828in}{1.934644in}}%
\pgfpathlineto{\pgfqpoint{1.502521in}{1.933046in}}%
\pgfpathlineto{\pgfqpoint{1.502718in}{1.933486in}}%
\pgfpathlineto{\pgfqpoint{1.502768in}{1.932405in}}%
\pgfpathlineto{\pgfqpoint{1.502916in}{1.932656in}}%
\pgfpathlineto{\pgfqpoint{1.503015in}{1.919882in}}%
\pgfpathlineto{\pgfqpoint{1.503640in}{1.532027in}}%
\pgfpathlineto{\pgfqpoint{1.504576in}{1.532048in}}%
\pgfpathlineto{\pgfqpoint{1.505516in}{1.532140in}}%
\pgfpathlineto{\pgfqpoint{1.505744in}{1.957344in}}%
\pgfpathlineto{\pgfqpoint{1.506656in}{1.933209in}}%
\pgfpathlineto{\pgfqpoint{1.506904in}{1.932349in}}%
\pgfpathlineto{\pgfqpoint{1.506854in}{1.934111in}}%
\pgfpathlineto{\pgfqpoint{1.507448in}{1.932728in}}%
\pgfpathlineto{\pgfqpoint{1.507645in}{1.933088in}}%
\pgfpathlineto{\pgfqpoint{1.507893in}{1.931845in}}%
\pgfpathlineto{\pgfqpoint{1.508140in}{1.798589in}}%
\pgfpathlineto{\pgfqpoint{1.508538in}{1.532025in}}%
\pgfpathlineto{\pgfqpoint{1.509183in}{1.532051in}}%
\pgfpathlineto{\pgfqpoint{1.510460in}{1.532126in}}%
\pgfpathlineto{\pgfqpoint{1.510790in}{1.957209in}}%
\pgfpathlineto{\pgfqpoint{1.511730in}{1.933324in}}%
\pgfpathlineto{\pgfqpoint{1.511978in}{1.933383in}}%
\pgfpathlineto{\pgfqpoint{1.512571in}{1.932433in}}%
\pgfpathlineto{\pgfqpoint{1.512620in}{1.932657in}}%
\pgfpathlineto{\pgfqpoint{1.512868in}{1.930557in}}%
\pgfpathlineto{\pgfqpoint{1.513414in}{1.585483in}}%
\pgfpathlineto{\pgfqpoint{1.513482in}{1.532026in}}%
\pgfpathlineto{\pgfqpoint{1.514175in}{1.532040in}}%
\pgfpathlineto{\pgfqpoint{1.515407in}{1.532139in}}%
\pgfpathlineto{\pgfqpoint{1.515641in}{1.957862in}}%
\pgfpathlineto{\pgfqpoint{1.516535in}{1.933608in}}%
\pgfpathlineto{\pgfqpoint{1.516832in}{1.931569in}}%
\pgfpathlineto{\pgfqpoint{1.516782in}{1.933951in}}%
\pgfpathlineto{\pgfqpoint{1.517277in}{1.931874in}}%
\pgfpathlineto{\pgfqpoint{1.517326in}{1.933224in}}%
\pgfpathlineto{\pgfqpoint{1.517821in}{1.929394in}}%
\pgfpathlineto{\pgfqpoint{1.518429in}{1.532025in}}%
\pgfpathlineto{\pgfqpoint{1.519865in}{1.532069in}}%
\pgfpathlineto{\pgfqpoint{1.520351in}{1.532125in}}%
\pgfpathlineto{\pgfqpoint{1.520683in}{1.955589in}}%
\pgfpathlineto{\pgfqpoint{1.521615in}{1.932899in}}%
\pgfpathlineto{\pgfqpoint{1.522455in}{1.932028in}}%
\pgfpathlineto{\pgfqpoint{1.522505in}{1.932249in}}%
\pgfpathlineto{\pgfqpoint{1.522752in}{1.931040in}}%
\pgfpathlineto{\pgfqpoint{1.523098in}{1.701458in}}%
\pgfpathlineto{\pgfqpoint{1.523374in}{1.532025in}}%
\pgfpathlineto{\pgfqpoint{1.524068in}{1.532040in}}%
\pgfpathlineto{\pgfqpoint{1.525297in}{1.532128in}}%
\pgfpathlineto{\pgfqpoint{1.525520in}{1.957093in}}%
\pgfpathlineto{\pgfqpoint{1.526535in}{1.932502in}}%
\pgfpathlineto{\pgfqpoint{1.526683in}{1.931812in}}%
\pgfpathlineto{\pgfqpoint{1.526733in}{1.932865in}}%
\pgfpathlineto{\pgfqpoint{1.527326in}{1.931971in}}%
\pgfpathlineto{\pgfqpoint{1.527524in}{1.931993in}}%
\pgfpathlineto{\pgfqpoint{1.527573in}{1.931477in}}%
\pgfpathlineto{\pgfqpoint{1.527623in}{1.931715in}}%
\pgfpathlineto{\pgfqpoint{1.527722in}{1.926441in}}%
\pgfpathlineto{\pgfqpoint{1.528321in}{1.532025in}}%
\pgfpathlineto{\pgfqpoint{1.529658in}{1.532067in}}%
\pgfpathlineto{\pgfqpoint{1.530242in}{1.532118in}}%
\pgfpathlineto{\pgfqpoint{1.530465in}{1.954670in}}%
\pgfpathlineto{\pgfqpoint{1.531573in}{1.932332in}}%
\pgfpathlineto{\pgfqpoint{1.531672in}{1.932367in}}%
\pgfpathlineto{\pgfqpoint{1.532414in}{1.931643in}}%
\pgfpathlineto{\pgfqpoint{1.532463in}{1.931690in}}%
\pgfpathlineto{\pgfqpoint{1.532611in}{1.931011in}}%
\pgfpathlineto{\pgfqpoint{1.532760in}{1.886626in}}%
\pgfpathlineto{\pgfqpoint{1.533266in}{1.532025in}}%
\pgfpathlineto{\pgfqpoint{1.534009in}{1.532053in}}%
\pgfpathlineto{\pgfqpoint{1.535187in}{1.532121in}}%
\pgfpathlineto{\pgfqpoint{1.535414in}{1.955904in}}%
\pgfpathlineto{\pgfqpoint{1.536508in}{1.931663in}}%
\pgfpathlineto{\pgfqpoint{1.536558in}{1.932260in}}%
\pgfpathlineto{\pgfqpoint{1.536607in}{1.931613in}}%
\pgfpathlineto{\pgfqpoint{1.537250in}{1.931618in}}%
\pgfpathlineto{\pgfqpoint{1.537596in}{1.930111in}}%
\pgfpathlineto{\pgfqpoint{1.537992in}{1.666346in}}%
\pgfpathlineto{\pgfqpoint{1.538210in}{1.532025in}}%
\pgfpathlineto{\pgfqpoint{1.539003in}{1.532042in}}%
\pgfpathlineto{\pgfqpoint{1.540133in}{1.532127in}}%
\pgfpathlineto{\pgfqpoint{1.540463in}{1.956327in}}%
\pgfpathlineto{\pgfqpoint{1.541354in}{1.931207in}}%
\pgfpathlineto{\pgfqpoint{1.541651in}{1.932152in}}%
\pgfpathlineto{\pgfqpoint{1.542096in}{1.931795in}}%
\pgfpathlineto{\pgfqpoint{1.542591in}{1.918976in}}%
\pgfpathlineto{\pgfqpoint{1.543158in}{1.532025in}}%
\pgfpathlineto{\pgfqpoint{1.544099in}{1.532058in}}%
\pgfpathlineto{\pgfqpoint{1.545078in}{1.532117in}}%
\pgfpathlineto{\pgfqpoint{1.545296in}{1.955071in}}%
\pgfpathlineto{\pgfqpoint{1.546430in}{1.931418in}}%
\pgfpathlineto{\pgfqpoint{1.546479in}{1.931676in}}%
\pgfpathlineto{\pgfqpoint{1.547171in}{1.931200in}}%
\pgfpathlineto{\pgfqpoint{1.547468in}{1.930359in}}%
\pgfpathlineto{\pgfqpoint{1.547666in}{1.839898in}}%
\pgfpathlineto{\pgfqpoint{1.548104in}{1.532025in}}%
\pgfpathlineto{\pgfqpoint{1.548798in}{1.532039in}}%
\pgfpathlineto{\pgfqpoint{1.550023in}{1.532114in}}%
\pgfpathlineto{\pgfqpoint{1.550207in}{1.952556in}}%
\pgfpathlineto{\pgfqpoint{1.551412in}{1.931360in}}%
\pgfpathlineto{\pgfqpoint{1.552401in}{1.930279in}}%
\pgfpathlineto{\pgfqpoint{1.552549in}{1.885314in}}%
\pgfpathlineto{\pgfqpoint{1.553050in}{1.532024in}}%
\pgfpathlineto{\pgfqpoint{1.553793in}{1.532053in}}%
\pgfpathlineto{\pgfqpoint{1.554968in}{1.532113in}}%
\pgfpathlineto{\pgfqpoint{1.555010in}{1.748535in}}%
\pgfpathlineto{\pgfqpoint{1.555088in}{1.952713in}}%
\pgfpathlineto{\pgfqpoint{1.555785in}{1.928190in}}%
\pgfpathlineto{\pgfqpoint{1.555834in}{1.933775in}}%
\pgfpathlineto{\pgfqpoint{1.556526in}{1.931715in}}%
\pgfpathlineto{\pgfqpoint{1.556675in}{1.929971in}}%
\pgfpathlineto{\pgfqpoint{1.556625in}{1.932203in}}%
\pgfpathlineto{\pgfqpoint{1.557318in}{1.930223in}}%
\pgfpathlineto{\pgfqpoint{1.557417in}{1.924621in}}%
\pgfpathlineto{\pgfqpoint{1.557994in}{1.532025in}}%
\pgfpathlineto{\pgfqpoint{1.559132in}{1.532062in}}%
\pgfpathlineto{\pgfqpoint{1.559914in}{1.532115in}}%
\pgfpathlineto{\pgfqpoint{1.560125in}{1.954865in}}%
\pgfpathlineto{\pgfqpoint{1.561289in}{1.931031in}}%
\pgfpathlineto{\pgfqpoint{1.562328in}{1.929163in}}%
\pgfpathlineto{\pgfqpoint{1.562724in}{1.666931in}}%
\pgfpathlineto{\pgfqpoint{1.562940in}{1.532025in}}%
\pgfpathlineto{\pgfqpoint{1.563733in}{1.532041in}}%
\pgfpathlineto{\pgfqpoint{1.564859in}{1.532114in}}%
\pgfpathlineto{\pgfqpoint{1.565041in}{1.952174in}}%
\pgfpathlineto{\pgfqpoint{1.566271in}{1.930756in}}%
\pgfpathlineto{\pgfqpoint{1.567013in}{1.930392in}}%
\pgfpathlineto{\pgfqpoint{1.567261in}{1.929523in}}%
\pgfpathlineto{\pgfqpoint{1.567508in}{1.796501in}}%
\pgfpathlineto{\pgfqpoint{1.567890in}{1.532024in}}%
\pgfpathlineto{\pgfqpoint{1.568585in}{1.532039in}}%
\pgfpathlineto{\pgfqpoint{1.569804in}{1.532107in}}%
\pgfpathlineto{\pgfqpoint{1.569844in}{1.708007in}}%
\pgfpathlineto{\pgfqpoint{1.570113in}{1.952766in}}%
\pgfpathlineto{\pgfqpoint{1.570621in}{1.931321in}}%
\pgfpathlineto{\pgfqpoint{1.570770in}{1.928798in}}%
\pgfpathlineto{\pgfqpoint{1.570720in}{1.932185in}}%
\pgfpathlineto{\pgfqpoint{1.571412in}{1.930393in}}%
\pgfpathlineto{\pgfqpoint{1.572055in}{1.930449in}}%
\pgfpathlineto{\pgfqpoint{1.572253in}{1.922571in}}%
\pgfpathlineto{\pgfqpoint{1.572833in}{1.532024in}}%
\pgfpathlineto{\pgfqpoint{1.574070in}{1.532065in}}%
\pgfpathlineto{\pgfqpoint{1.574750in}{1.532109in}}%
\pgfpathlineto{\pgfqpoint{1.574790in}{1.720746in}}%
\pgfpathlineto{\pgfqpoint{1.574962in}{1.953296in}}%
\pgfpathlineto{\pgfqpoint{1.575546in}{1.926645in}}%
\pgfpathlineto{\pgfqpoint{1.575595in}{1.933664in}}%
\pgfpathlineto{\pgfqpoint{1.576288in}{1.931230in}}%
\pgfpathlineto{\pgfqpoint{1.576436in}{1.928931in}}%
\pgfpathlineto{\pgfqpoint{1.576386in}{1.931786in}}%
\pgfpathlineto{\pgfqpoint{1.577079in}{1.929598in}}%
\pgfpathlineto{\pgfqpoint{1.577178in}{1.929019in}}%
\pgfpathlineto{\pgfqpoint{1.577705in}{1.590585in}}%
\pgfpathlineto{\pgfqpoint{1.577776in}{1.532025in}}%
\pgfpathlineto{\pgfqpoint{1.578468in}{1.532039in}}%
\pgfpathlineto{\pgfqpoint{1.579696in}{1.532117in}}%
\pgfpathlineto{\pgfqpoint{1.579919in}{1.953878in}}%
\pgfpathlineto{\pgfqpoint{1.581028in}{1.930273in}}%
\pgfpathlineto{\pgfqpoint{1.581127in}{1.930362in}}%
\pgfpathlineto{\pgfqpoint{1.582116in}{1.928481in}}%
\pgfpathlineto{\pgfqpoint{1.582512in}{1.667349in}}%
\pgfpathlineto{\pgfqpoint{1.582726in}{1.532024in}}%
\pgfpathlineto{\pgfqpoint{1.583520in}{1.532041in}}%
\pgfpathlineto{\pgfqpoint{1.584640in}{1.532107in}}%
\pgfpathlineto{\pgfqpoint{1.584680in}{1.707898in}}%
\pgfpathlineto{\pgfqpoint{1.584853in}{1.955384in}}%
\pgfpathlineto{\pgfqpoint{1.585458in}{1.928923in}}%
\pgfpathlineto{\pgfqpoint{1.585755in}{1.931349in}}%
\pgfpathlineto{\pgfqpoint{1.585705in}{1.928717in}}%
\pgfpathlineto{\pgfqpoint{1.586200in}{1.930847in}}%
\pgfpathlineto{\pgfqpoint{1.587189in}{1.879326in}}%
\pgfpathlineto{\pgfqpoint{1.587676in}{1.532022in}}%
\pgfpathlineto{\pgfqpoint{1.588422in}{1.532056in}}%
\pgfpathlineto{\pgfqpoint{1.589583in}{1.532093in}}%
\pgfpathlineto{\pgfqpoint{1.589591in}{1.532179in}}%
\pgfpathlineto{\pgfqpoint{1.589707in}{1.950850in}}%
\pgfpathlineto{\pgfqpoint{1.590689in}{1.929897in}}%
\pgfpathlineto{\pgfqpoint{1.591975in}{1.928775in}}%
\pgfpathlineto{\pgfqpoint{1.592123in}{1.883421in}}%
\pgfpathlineto{\pgfqpoint{1.592618in}{1.532024in}}%
\pgfpathlineto{\pgfqpoint{1.593363in}{1.532054in}}%
\pgfpathlineto{\pgfqpoint{1.594531in}{1.532107in}}%
\pgfpathlineto{\pgfqpoint{1.594571in}{1.707277in}}%
\pgfpathlineto{\pgfqpoint{1.594655in}{1.951468in}}%
\pgfpathlineto{\pgfqpoint{1.595315in}{1.931473in}}%
\pgfpathlineto{\pgfqpoint{1.595364in}{1.927631in}}%
\pgfpathlineto{\pgfqpoint{1.596057in}{1.929473in}}%
\pgfpathlineto{\pgfqpoint{1.596205in}{1.930407in}}%
\pgfpathlineto{\pgfqpoint{1.596699in}{1.928817in}}%
\pgfpathlineto{\pgfqpoint{1.596749in}{1.929601in}}%
\pgfpathlineto{\pgfqpoint{1.596996in}{1.921784in}}%
\pgfpathlineto{\pgfqpoint{1.597567in}{1.532022in}}%
\pgfpathlineto{\pgfqpoint{1.598659in}{1.532046in}}%
\pgfpathlineto{\pgfqpoint{1.599484in}{1.532326in}}%
\pgfpathlineto{\pgfqpoint{1.599592in}{1.950355in}}%
\pgfpathlineto{\pgfqpoint{1.600489in}{1.928161in}}%
\pgfpathlineto{\pgfqpoint{1.600539in}{1.930690in}}%
\pgfpathlineto{\pgfqpoint{1.601231in}{1.929742in}}%
\pgfpathlineto{\pgfqpoint{1.601330in}{1.929897in}}%
\pgfpathlineto{\pgfqpoint{1.601973in}{1.907745in}}%
\pgfpathlineto{\pgfqpoint{1.602509in}{1.532024in}}%
\pgfpathlineto{\pgfqpoint{1.603401in}{1.532043in}}%
\pgfpathlineto{\pgfqpoint{1.604422in}{1.532106in}}%
\pgfpathlineto{\pgfqpoint{1.604462in}{1.704843in}}%
\pgfpathlineto{\pgfqpoint{1.604731in}{1.952724in}}%
\pgfpathlineto{\pgfqpoint{1.605216in}{1.930864in}}%
\pgfpathlineto{\pgfqpoint{1.605266in}{1.927406in}}%
\pgfpathlineto{\pgfqpoint{1.605315in}{1.931045in}}%
\pgfpathlineto{\pgfqpoint{1.605958in}{1.929317in}}%
\pgfpathlineto{\pgfqpoint{1.606106in}{1.929966in}}%
\pgfpathlineto{\pgfqpoint{1.606156in}{1.928542in}}%
\pgfpathlineto{\pgfqpoint{1.606551in}{1.929418in}}%
\pgfpathlineto{\pgfqpoint{1.606898in}{1.919464in}}%
\pgfpathlineto{\pgfqpoint{1.607455in}{1.532024in}}%
\pgfpathlineto{\pgfqpoint{1.608500in}{1.532060in}}%
\pgfpathlineto{\pgfqpoint{1.609367in}{1.532107in}}%
\pgfpathlineto{\pgfqpoint{1.609408in}{1.711140in}}%
\pgfpathlineto{\pgfqpoint{1.609580in}{1.953898in}}%
\pgfpathlineto{\pgfqpoint{1.610154in}{1.926408in}}%
\pgfpathlineto{\pgfqpoint{1.610203in}{1.932161in}}%
\pgfpathlineto{\pgfqpoint{1.610253in}{1.925959in}}%
\pgfpathlineto{\pgfqpoint{1.610945in}{1.928759in}}%
\pgfpathlineto{\pgfqpoint{1.611093in}{1.930338in}}%
\pgfpathlineto{\pgfqpoint{1.611143in}{1.927872in}}%
\pgfpathlineto{\pgfqpoint{1.611687in}{1.928033in}}%
\pgfpathlineto{\pgfqpoint{1.611736in}{1.928453in}}%
\pgfpathlineto{\pgfqpoint{1.611786in}{1.927900in}}%
\pgfpathlineto{\pgfqpoint{1.612132in}{1.710674in}}%
\pgfpathlineto{\pgfqpoint{1.612401in}{1.532024in}}%
\pgfpathlineto{\pgfqpoint{1.613149in}{1.532054in}}%
\pgfpathlineto{\pgfqpoint{1.614313in}{1.532108in}}%
\pgfpathlineto{\pgfqpoint{1.614353in}{1.714857in}}%
\pgfpathlineto{\pgfqpoint{1.614525in}{1.952904in}}%
\pgfpathlineto{\pgfqpoint{1.615137in}{1.931364in}}%
\pgfpathlineto{\pgfqpoint{1.615186in}{1.926213in}}%
\pgfpathlineto{\pgfqpoint{1.615236in}{1.931523in}}%
\pgfpathlineto{\pgfqpoint{1.615879in}{1.928913in}}%
\pgfpathlineto{\pgfqpoint{1.616027in}{1.929956in}}%
\pgfpathlineto{\pgfqpoint{1.616077in}{1.927855in}}%
\pgfpathlineto{\pgfqpoint{1.616571in}{1.929048in}}%
\pgfpathlineto{\pgfqpoint{1.616818in}{1.906861in}}%
\pgfpathlineto{\pgfqpoint{1.617348in}{1.532024in}}%
\pgfpathlineto{\pgfqpoint{1.618242in}{1.532043in}}%
\pgfpathlineto{\pgfqpoint{1.619258in}{1.532107in}}%
\pgfpathlineto{\pgfqpoint{1.619299in}{1.712474in}}%
\pgfpathlineto{\pgfqpoint{1.619470in}{1.953170in}}%
\pgfpathlineto{\pgfqpoint{1.620053in}{1.926603in}}%
\pgfpathlineto{\pgfqpoint{1.620202in}{1.931630in}}%
\pgfpathlineto{\pgfqpoint{1.620152in}{1.925759in}}%
\pgfpathlineto{\pgfqpoint{1.620844in}{1.928709in}}%
\pgfpathlineto{\pgfqpoint{1.620993in}{1.929899in}}%
\pgfpathlineto{\pgfqpoint{1.621042in}{1.927565in}}%
\pgfpathlineto{\pgfqpoint{1.621537in}{1.928893in}}%
\pgfpathlineto{\pgfqpoint{1.621735in}{1.916689in}}%
\pgfpathlineto{\pgfqpoint{1.622293in}{1.532024in}}%
\pgfpathlineto{\pgfqpoint{1.623387in}{1.532048in}}%
\pgfpathlineto{\pgfqpoint{1.624204in}{1.532109in}}%
\pgfpathlineto{\pgfqpoint{1.624245in}{1.720395in}}%
\pgfpathlineto{\pgfqpoint{1.624417in}{1.951300in}}%
\pgfpathlineto{\pgfqpoint{1.624985in}{1.930569in}}%
\pgfpathlineto{\pgfqpoint{1.625133in}{1.925157in}}%
\pgfpathlineto{\pgfqpoint{1.625084in}{1.932027in}}%
\pgfpathlineto{\pgfqpoint{1.625776in}{1.928511in}}%
\pgfpathlineto{\pgfqpoint{1.625974in}{1.929968in}}%
\pgfpathlineto{\pgfqpoint{1.626667in}{1.923866in}}%
\pgfpathlineto{\pgfqpoint{1.627239in}{1.532024in}}%
\pgfpathlineto{\pgfqpoint{1.628429in}{1.532050in}}%
\pgfpathlineto{\pgfqpoint{1.629149in}{1.532107in}}%
\pgfpathlineto{\pgfqpoint{1.629190in}{1.712456in}}%
\pgfpathlineto{\pgfqpoint{1.629361in}{1.952574in}}%
\pgfpathlineto{\pgfqpoint{1.629935in}{1.925942in}}%
\pgfpathlineto{\pgfqpoint{1.629985in}{1.931392in}}%
\pgfpathlineto{\pgfqpoint{1.630034in}{1.925351in}}%
\pgfpathlineto{\pgfqpoint{1.630727in}{1.928210in}}%
\pgfpathlineto{\pgfqpoint{1.630875in}{1.929654in}}%
\pgfpathlineto{\pgfqpoint{1.630924in}{1.927232in}}%
\pgfpathlineto{\pgfqpoint{1.631419in}{1.928573in}}%
\pgfpathlineto{\pgfqpoint{1.631617in}{1.919009in}}%
\pgfpathlineto{\pgfqpoint{1.632185in}{1.532024in}}%
\pgfpathlineto{\pgfqpoint{1.633424in}{1.532065in}}%
\pgfpathlineto{\pgfqpoint{1.634094in}{1.532105in}}%
\pgfpathlineto{\pgfqpoint{1.634134in}{1.700402in}}%
\pgfpathlineto{\pgfqpoint{1.634218in}{1.950014in}}%
\pgfpathlineto{\pgfqpoint{1.634884in}{1.926023in}}%
\pgfpathlineto{\pgfqpoint{1.635032in}{1.931084in}}%
\pgfpathlineto{\pgfqpoint{1.634983in}{1.925290in}}%
\pgfpathlineto{\pgfqpoint{1.635675in}{1.928162in}}%
\pgfpathlineto{\pgfqpoint{1.635824in}{1.929432in}}%
\pgfpathlineto{\pgfqpoint{1.635873in}{1.927101in}}%
\pgfpathlineto{\pgfqpoint{1.636368in}{1.928419in}}%
\pgfpathlineto{\pgfqpoint{1.636565in}{1.919065in}}%
\pgfpathlineto{\pgfqpoint{1.637064in}{1.532023in}}%
\pgfpathlineto{\pgfqpoint{1.638349in}{1.532065in}}%
\pgfpathlineto{\pgfqpoint{1.639039in}{1.532101in}}%
\pgfpathlineto{\pgfqpoint{1.639066in}{1.573396in}}%
\pgfpathlineto{\pgfqpoint{1.639166in}{1.949659in}}%
\pgfpathlineto{\pgfqpoint{1.639854in}{1.933479in}}%
\pgfpathlineto{\pgfqpoint{1.640151in}{1.923530in}}%
\pgfpathlineto{\pgfqpoint{1.640596in}{1.925232in}}%
\pgfpathlineto{\pgfqpoint{1.640646in}{1.931103in}}%
\pgfpathlineto{\pgfqpoint{1.641338in}{1.928006in}}%
\pgfpathlineto{\pgfqpoint{1.641437in}{1.928256in}}%
\pgfpathlineto{\pgfqpoint{1.641536in}{1.913123in}}%
\pgfpathlineto{\pgfqpoint{1.642087in}{1.532027in}}%
\pgfpathlineto{\pgfqpoint{1.642992in}{1.532046in}}%
\pgfpathlineto{\pgfqpoint{1.643987in}{1.532114in}}%
\pgfpathlineto{\pgfqpoint{1.644042in}{1.917931in}}%
\pgfpathlineto{\pgfqpoint{1.644114in}{1.949455in}}%
\pgfpathlineto{\pgfqpoint{1.644152in}{1.905289in}}%
\pgfpathlineto{\pgfqpoint{1.644811in}{1.927025in}}%
\pgfpathlineto{\pgfqpoint{1.645108in}{1.929244in}}%
\pgfpathlineto{\pgfqpoint{1.645058in}{1.926834in}}%
\pgfpathlineto{\pgfqpoint{1.645553in}{1.928788in}}%
\pgfpathlineto{\pgfqpoint{1.646394in}{1.926677in}}%
\pgfpathlineto{\pgfqpoint{1.646410in}{1.926924in}}%
\pgfpathlineto{\pgfqpoint{1.646483in}{1.913354in}}%
\pgfpathlineto{\pgfqpoint{1.647050in}{1.532024in}}%
\pgfpathlineto{\pgfqpoint{1.648068in}{1.532047in}}%
\pgfpathlineto{\pgfqpoint{1.648932in}{1.532112in}}%
\pgfpathlineto{\pgfqpoint{1.648974in}{1.749933in}}%
\pgfpathlineto{\pgfqpoint{1.649051in}{1.949104in}}%
\pgfpathlineto{\pgfqpoint{1.649744in}{1.924916in}}%
\pgfpathlineto{\pgfqpoint{1.649793in}{1.930006in}}%
\pgfpathlineto{\pgfqpoint{1.650486in}{1.929040in}}%
\pgfpathlineto{\pgfqpoint{1.650535in}{1.926584in}}%
\pgfpathlineto{\pgfqpoint{1.650585in}{1.929110in}}%
\pgfpathlineto{\pgfqpoint{1.651277in}{1.927367in}}%
\pgfpathlineto{\pgfqpoint{1.651425in}{1.914344in}}%
\pgfpathlineto{\pgfqpoint{1.651967in}{1.532023in}}%
\pgfpathlineto{\pgfqpoint{1.652958in}{1.532045in}}%
\pgfpathlineto{\pgfqpoint{1.653884in}{1.532504in}}%
\pgfpathlineto{\pgfqpoint{1.653996in}{1.949029in}}%
\pgfpathlineto{\pgfqpoint{1.654907in}{1.928795in}}%
\pgfpathlineto{\pgfqpoint{1.654957in}{1.926692in}}%
\pgfpathlineto{\pgfqpoint{1.655649in}{1.927449in}}%
\pgfpathlineto{\pgfqpoint{1.655797in}{1.928011in}}%
\pgfpathlineto{\pgfqpoint{1.656292in}{1.926376in}}%
\pgfpathlineto{\pgfqpoint{1.656490in}{1.842496in}}%
\pgfpathlineto{\pgfqpoint{1.656921in}{1.532027in}}%
\pgfpathlineto{\pgfqpoint{1.657627in}{1.532042in}}%
\pgfpathlineto{\pgfqpoint{1.658823in}{1.532115in}}%
\pgfpathlineto{\pgfqpoint{1.658945in}{1.948982in}}%
\pgfpathlineto{\pgfqpoint{1.660380in}{1.926031in}}%
\pgfpathlineto{\pgfqpoint{1.660429in}{1.928984in}}%
\pgfpathlineto{\pgfqpoint{1.661122in}{1.927362in}}%
\pgfpathlineto{\pgfqpoint{1.661319in}{1.914542in}}%
\pgfpathlineto{\pgfqpoint{1.661847in}{1.532023in}}%
\pgfpathlineto{\pgfqpoint{1.662832in}{1.532044in}}%
\pgfpathlineto{\pgfqpoint{1.663776in}{1.533150in}}%
\pgfpathlineto{\pgfqpoint{1.663881in}{1.947191in}}%
\pgfpathlineto{\pgfqpoint{1.664770in}{1.931345in}}%
\pgfpathlineto{\pgfqpoint{1.664820in}{1.924528in}}%
\pgfpathlineto{\pgfqpoint{1.665512in}{1.925654in}}%
\pgfpathlineto{\pgfqpoint{1.665562in}{1.929215in}}%
\pgfpathlineto{\pgfqpoint{1.665611in}{1.925587in}}%
\pgfpathlineto{\pgfqpoint{1.666205in}{1.925665in}}%
\pgfpathlineto{\pgfqpoint{1.666452in}{1.786681in}}%
\pgfpathlineto{\pgfqpoint{1.666829in}{1.532025in}}%
\pgfpathlineto{\pgfqpoint{1.667497in}{1.532051in}}%
\pgfpathlineto{\pgfqpoint{1.668714in}{1.532114in}}%
\pgfpathlineto{\pgfqpoint{1.668769in}{1.918083in}}%
\pgfpathlineto{\pgfqpoint{1.668928in}{1.948400in}}%
\pgfpathlineto{\pgfqpoint{1.668878in}{1.904100in}}%
\pgfpathlineto{\pgfqpoint{1.669524in}{1.928938in}}%
\pgfpathlineto{\pgfqpoint{1.669573in}{1.925525in}}%
\pgfpathlineto{\pgfqpoint{1.670266in}{1.926994in}}%
\pgfpathlineto{\pgfqpoint{1.670414in}{1.927946in}}%
\pgfpathlineto{\pgfqpoint{1.670908in}{1.926499in}}%
\pgfpathlineto{\pgfqpoint{1.671007in}{1.926516in}}%
\pgfpathlineto{\pgfqpoint{1.671057in}{1.926652in}}%
\pgfpathlineto{\pgfqpoint{1.671156in}{1.925575in}}%
\pgfpathlineto{\pgfqpoint{1.671452in}{1.744756in}}%
\pgfpathlineto{\pgfqpoint{1.671749in}{1.532024in}}%
\pgfpathlineto{\pgfqpoint{1.672495in}{1.532053in}}%
\pgfpathlineto{\pgfqpoint{1.673659in}{1.532108in}}%
\pgfpathlineto{\pgfqpoint{1.673699in}{1.720001in}}%
\pgfpathlineto{\pgfqpoint{1.673784in}{1.948673in}}%
\pgfpathlineto{\pgfqpoint{1.674463in}{1.929290in}}%
\pgfpathlineto{\pgfqpoint{1.674513in}{1.924971in}}%
\pgfpathlineto{\pgfqpoint{1.675205in}{1.926688in}}%
\pgfpathlineto{\pgfqpoint{1.675353in}{1.927992in}}%
\pgfpathlineto{\pgfqpoint{1.675304in}{1.926274in}}%
\pgfpathlineto{\pgfqpoint{1.675947in}{1.926367in}}%
\pgfpathlineto{\pgfqpoint{1.675996in}{1.926505in}}%
\pgfpathlineto{\pgfqpoint{1.676095in}{1.925496in}}%
\pgfpathlineto{\pgfqpoint{1.676343in}{1.789361in}}%
\pgfpathlineto{\pgfqpoint{1.676693in}{1.532024in}}%
\pgfpathlineto{\pgfqpoint{1.677389in}{1.532039in}}%
\pgfpathlineto{\pgfqpoint{1.678604in}{1.532107in}}%
\pgfpathlineto{\pgfqpoint{1.678644in}{1.711952in}}%
\pgfpathlineto{\pgfqpoint{1.678814in}{1.950259in}}%
\pgfpathlineto{\pgfqpoint{1.679424in}{1.930131in}}%
\pgfpathlineto{\pgfqpoint{1.679474in}{1.923923in}}%
\pgfpathlineto{\pgfqpoint{1.680166in}{1.926331in}}%
\pgfpathlineto{\pgfqpoint{1.680314in}{1.928229in}}%
\pgfpathlineto{\pgfqpoint{1.680265in}{1.925752in}}%
\pgfpathlineto{\pgfqpoint{1.680908in}{1.926109in}}%
\pgfpathlineto{\pgfqpoint{1.680957in}{1.926334in}}%
\pgfpathlineto{\pgfqpoint{1.681056in}{1.924219in}}%
\pgfpathlineto{\pgfqpoint{1.681402in}{1.700511in}}%
\pgfpathlineto{\pgfqpoint{1.681639in}{1.532024in}}%
\pgfpathlineto{\pgfqpoint{1.682384in}{1.532053in}}%
\pgfpathlineto{\pgfqpoint{1.683550in}{1.532109in}}%
\pgfpathlineto{\pgfqpoint{1.683591in}{1.725464in}}%
\pgfpathlineto{\pgfqpoint{1.683765in}{1.947958in}}%
\pgfpathlineto{\pgfqpoint{1.684338in}{1.924500in}}%
\pgfpathlineto{\pgfqpoint{1.684387in}{1.929374in}}%
\pgfpathlineto{\pgfqpoint{1.684436in}{1.924238in}}%
\pgfpathlineto{\pgfqpoint{1.685129in}{1.926485in}}%
\pgfpathlineto{\pgfqpoint{1.685277in}{1.927859in}}%
\pgfpathlineto{\pgfqpoint{1.685327in}{1.925823in}}%
\pgfpathlineto{\pgfqpoint{1.685871in}{1.925907in}}%
\pgfpathlineto{\pgfqpoint{1.685920in}{1.926201in}}%
\pgfpathlineto{\pgfqpoint{1.685970in}{1.925839in}}%
\pgfpathlineto{\pgfqpoint{1.686167in}{1.844251in}}%
\pgfpathlineto{\pgfqpoint{1.686584in}{1.532024in}}%
\pgfpathlineto{\pgfqpoint{1.687330in}{1.532053in}}%
\pgfpathlineto{\pgfqpoint{1.688495in}{1.532108in}}%
\pgfpathlineto{\pgfqpoint{1.688536in}{1.720124in}}%
\pgfpathlineto{\pgfqpoint{1.688706in}{1.948443in}}%
\pgfpathlineto{\pgfqpoint{1.689296in}{1.923896in}}%
\pgfpathlineto{\pgfqpoint{1.689345in}{1.929657in}}%
\pgfpathlineto{\pgfqpoint{1.689395in}{1.923797in}}%
\pgfpathlineto{\pgfqpoint{1.690037in}{1.926822in}}%
\pgfpathlineto{\pgfqpoint{1.690285in}{1.925576in}}%
\pgfpathlineto{\pgfqpoint{1.690235in}{1.927871in}}%
\pgfpathlineto{\pgfqpoint{1.690829in}{1.925751in}}%
\pgfpathlineto{\pgfqpoint{1.690878in}{1.925978in}}%
\pgfpathlineto{\pgfqpoint{1.690928in}{1.925662in}}%
\pgfpathlineto{\pgfqpoint{1.691175in}{1.796373in}}%
\pgfpathlineto{\pgfqpoint{1.691531in}{1.532024in}}%
\pgfpathlineto{\pgfqpoint{1.692227in}{1.532039in}}%
\pgfpathlineto{\pgfqpoint{1.693441in}{1.532110in}}%
\pgfpathlineto{\pgfqpoint{1.693482in}{1.731217in}}%
\pgfpathlineto{\pgfqpoint{1.693567in}{1.947685in}}%
\pgfpathlineto{\pgfqpoint{1.694256in}{1.928587in}}%
\pgfpathlineto{\pgfqpoint{1.694306in}{1.924566in}}%
\pgfpathlineto{\pgfqpoint{1.694998in}{1.926190in}}%
\pgfpathlineto{\pgfqpoint{1.695147in}{1.927384in}}%
\pgfpathlineto{\pgfqpoint{1.695641in}{1.925751in}}%
\pgfpathlineto{\pgfqpoint{1.695740in}{1.925818in}}%
\pgfpathlineto{\pgfqpoint{1.695790in}{1.925945in}}%
\pgfpathlineto{\pgfqpoint{1.695888in}{1.924794in}}%
\pgfpathlineto{\pgfqpoint{1.696185in}{1.744091in}}%
\pgfpathlineto{\pgfqpoint{1.696506in}{1.532029in}}%
\pgfpathlineto{\pgfqpoint{1.697227in}{1.532045in}}%
\pgfpathlineto{\pgfqpoint{1.698392in}{1.532289in}}%
\pgfpathlineto{\pgfqpoint{1.698510in}{1.948066in}}%
\pgfpathlineto{\pgfqpoint{1.699439in}{1.923461in}}%
\pgfpathlineto{\pgfqpoint{1.699834in}{1.928644in}}%
\pgfpathlineto{\pgfqpoint{1.700181in}{1.927724in}}%
\pgfpathlineto{\pgfqpoint{1.700972in}{1.868560in}}%
\pgfpathlineto{\pgfqpoint{1.701420in}{1.532024in}}%
\pgfpathlineto{\pgfqpoint{1.702166in}{1.532053in}}%
\pgfpathlineto{\pgfqpoint{1.703332in}{1.532109in}}%
\pgfpathlineto{\pgfqpoint{1.703373in}{1.727995in}}%
\pgfpathlineto{\pgfqpoint{1.703459in}{1.947455in}}%
\pgfpathlineto{\pgfqpoint{1.704128in}{1.927997in}}%
\pgfpathlineto{\pgfqpoint{1.704177in}{1.924681in}}%
\pgfpathlineto{\pgfqpoint{1.704870in}{1.925881in}}%
\pgfpathlineto{\pgfqpoint{1.705018in}{1.926990in}}%
\pgfpathlineto{\pgfqpoint{1.704969in}{1.925621in}}%
\pgfpathlineto{\pgfqpoint{1.705612in}{1.925686in}}%
\pgfpathlineto{\pgfqpoint{1.705760in}{1.924995in}}%
\pgfpathlineto{\pgfqpoint{1.705908in}{1.877075in}}%
\pgfpathlineto{\pgfqpoint{1.706365in}{1.532024in}}%
\pgfpathlineto{\pgfqpoint{1.707159in}{1.532041in}}%
\pgfpathlineto{\pgfqpoint{1.708277in}{1.532109in}}%
\pgfpathlineto{\pgfqpoint{1.708318in}{1.725522in}}%
\pgfpathlineto{\pgfqpoint{1.708404in}{1.947529in}}%
\pgfpathlineto{\pgfqpoint{1.709072in}{1.928205in}}%
\pgfpathlineto{\pgfqpoint{1.709122in}{1.924214in}}%
\pgfpathlineto{\pgfqpoint{1.709814in}{1.925705in}}%
\pgfpathlineto{\pgfqpoint{1.709963in}{1.926994in}}%
\pgfpathlineto{\pgfqpoint{1.709913in}{1.925369in}}%
\pgfpathlineto{\pgfqpoint{1.710556in}{1.925509in}}%
\pgfpathlineto{\pgfqpoint{1.710606in}{1.925572in}}%
\pgfpathlineto{\pgfqpoint{1.710728in}{1.924724in}}%
\pgfpathlineto{\pgfqpoint{1.710831in}{1.890058in}}%
\pgfpathlineto{\pgfqpoint{1.711309in}{1.532024in}}%
\pgfpathlineto{\pgfqpoint{1.712148in}{1.532055in}}%
\pgfpathlineto{\pgfqpoint{1.713222in}{1.532109in}}%
\pgfpathlineto{\pgfqpoint{1.713263in}{1.723821in}}%
\pgfpathlineto{\pgfqpoint{1.713437in}{1.947249in}}%
\pgfpathlineto{\pgfqpoint{1.714036in}{1.924750in}}%
\pgfpathlineto{\pgfqpoint{1.714185in}{1.928795in}}%
\pgfpathlineto{\pgfqpoint{1.714135in}{1.923412in}}%
\pgfpathlineto{\pgfqpoint{1.714828in}{1.926366in}}%
\pgfpathlineto{\pgfqpoint{1.715075in}{1.927060in}}%
\pgfpathlineto{\pgfqpoint{1.715668in}{1.924575in}}%
\pgfpathlineto{\pgfqpoint{1.715965in}{1.749543in}}%
\pgfpathlineto{\pgfqpoint{1.716296in}{1.532025in}}%
\pgfpathlineto{\pgfqpoint{1.716976in}{1.532052in}}%
\pgfpathlineto{\pgfqpoint{1.718159in}{1.532080in}}%
\pgfpathlineto{\pgfqpoint{1.718171in}{1.532139in}}%
\pgfpathlineto{\pgfqpoint{1.718402in}{1.950261in}}%
\pgfpathlineto{\pgfqpoint{1.719283in}{1.924301in}}%
\pgfpathlineto{\pgfqpoint{1.719580in}{1.927804in}}%
\pgfpathlineto{\pgfqpoint{1.719530in}{1.924231in}}%
\pgfpathlineto{\pgfqpoint{1.720025in}{1.926996in}}%
\pgfpathlineto{\pgfqpoint{1.720767in}{1.860711in}}%
\pgfpathlineto{\pgfqpoint{1.721201in}{1.532024in}}%
\pgfpathlineto{\pgfqpoint{1.721941in}{1.532053in}}%
\pgfpathlineto{\pgfqpoint{1.723113in}{1.532109in}}%
\pgfpathlineto{\pgfqpoint{1.723154in}{1.724479in}}%
\pgfpathlineto{\pgfqpoint{1.723330in}{1.946333in}}%
\pgfpathlineto{\pgfqpoint{1.723921in}{1.922222in}}%
\pgfpathlineto{\pgfqpoint{1.723971in}{1.929492in}}%
\pgfpathlineto{\pgfqpoint{1.724663in}{1.926046in}}%
\pgfpathlineto{\pgfqpoint{1.724811in}{1.924363in}}%
\pgfpathlineto{\pgfqpoint{1.724861in}{1.927246in}}%
\pgfpathlineto{\pgfqpoint{1.725454in}{1.924815in}}%
\pgfpathlineto{\pgfqpoint{1.725504in}{1.925056in}}%
\pgfpathlineto{\pgfqpoint{1.725553in}{1.924848in}}%
\pgfpathlineto{\pgfqpoint{1.725800in}{1.797704in}}%
\pgfpathlineto{\pgfqpoint{1.726189in}{1.532025in}}%
\pgfpathlineto{\pgfqpoint{1.726870in}{1.532052in}}%
\pgfpathlineto{\pgfqpoint{1.728060in}{1.532121in}}%
\pgfpathlineto{\pgfqpoint{1.728293in}{1.950857in}}%
\pgfpathlineto{\pgfqpoint{1.729354in}{1.925719in}}%
\pgfpathlineto{\pgfqpoint{1.729552in}{1.926067in}}%
\pgfpathlineto{\pgfqpoint{1.730293in}{1.925153in}}%
\pgfpathlineto{\pgfqpoint{1.730343in}{1.925263in}}%
\pgfpathlineto{\pgfqpoint{1.730491in}{1.924389in}}%
\pgfpathlineto{\pgfqpoint{1.730689in}{1.838177in}}%
\pgfpathlineto{\pgfqpoint{1.731094in}{1.532024in}}%
\pgfpathlineto{\pgfqpoint{1.731839in}{1.532053in}}%
\pgfpathlineto{\pgfqpoint{1.733004in}{1.532110in}}%
\pgfpathlineto{\pgfqpoint{1.733045in}{1.730703in}}%
\pgfpathlineto{\pgfqpoint{1.733110in}{1.945289in}}%
\pgfpathlineto{\pgfqpoint{1.733817in}{1.925052in}}%
\pgfpathlineto{\pgfqpoint{1.733965in}{1.930472in}}%
\pgfpathlineto{\pgfqpoint{1.734014in}{1.920391in}}%
\pgfpathlineto{\pgfqpoint{1.734608in}{1.926971in}}%
\pgfpathlineto{\pgfqpoint{1.734855in}{1.927489in}}%
\pgfpathlineto{\pgfqpoint{1.735597in}{1.870128in}}%
\pgfpathlineto{\pgfqpoint{1.736084in}{1.532031in}}%
\pgfpathlineto{\pgfqpoint{1.736819in}{1.532047in}}%
\pgfpathlineto{\pgfqpoint{1.737955in}{1.532196in}}%
\pgfpathlineto{\pgfqpoint{1.738073in}{1.946949in}}%
\pgfpathlineto{\pgfqpoint{1.739022in}{1.925280in}}%
\pgfpathlineto{\pgfqpoint{1.739171in}{1.925924in}}%
\pgfpathlineto{\pgfqpoint{1.739220in}{1.925026in}}%
\pgfpathlineto{\pgfqpoint{1.739764in}{1.925220in}}%
\pgfpathlineto{\pgfqpoint{1.739962in}{1.925392in}}%
\pgfpathlineto{\pgfqpoint{1.740358in}{1.924439in}}%
\pgfpathlineto{\pgfqpoint{1.740407in}{1.923881in}}%
\pgfpathlineto{\pgfqpoint{1.740803in}{1.670162in}}%
\pgfpathlineto{\pgfqpoint{1.740977in}{1.532024in}}%
\pgfpathlineto{\pgfqpoint{1.741814in}{1.532055in}}%
\pgfpathlineto{\pgfqpoint{1.742895in}{1.532105in}}%
\pgfpathlineto{\pgfqpoint{1.742934in}{1.704651in}}%
\pgfpathlineto{\pgfqpoint{1.743203in}{1.947707in}}%
\pgfpathlineto{\pgfqpoint{1.743691in}{1.926639in}}%
\pgfpathlineto{\pgfqpoint{1.743741in}{1.923541in}}%
\pgfpathlineto{\pgfqpoint{1.743790in}{1.927001in}}%
\pgfpathlineto{\pgfqpoint{1.744433in}{1.925449in}}%
\pgfpathlineto{\pgfqpoint{1.744680in}{1.925938in}}%
\pgfpathlineto{\pgfqpoint{1.744631in}{1.924597in}}%
\pgfpathlineto{\pgfqpoint{1.744928in}{1.925260in}}%
\pgfpathlineto{\pgfqpoint{1.745027in}{1.925413in}}%
\pgfpathlineto{\pgfqpoint{1.745373in}{1.922596in}}%
\pgfpathlineto{\pgfqpoint{1.745854in}{1.563435in}}%
\pgfpathlineto{\pgfqpoint{1.745959in}{1.532029in}}%
\pgfpathlineto{\pgfqpoint{1.746624in}{1.532047in}}%
\pgfpathlineto{\pgfqpoint{1.747838in}{1.532096in}}%
\pgfpathlineto{\pgfqpoint{1.747847in}{1.532287in}}%
\pgfpathlineto{\pgfqpoint{1.748091in}{1.953018in}}%
\pgfpathlineto{\pgfqpoint{1.748856in}{1.925171in}}%
\pgfpathlineto{\pgfqpoint{1.750240in}{1.924287in}}%
\pgfpathlineto{\pgfqpoint{1.750290in}{1.923790in}}%
\pgfpathlineto{\pgfqpoint{1.750582in}{1.751034in}}%
\pgfpathlineto{\pgfqpoint{1.750867in}{1.532024in}}%
\pgfpathlineto{\pgfqpoint{1.751604in}{1.532053in}}%
\pgfpathlineto{\pgfqpoint{1.752786in}{1.532108in}}%
\pgfpathlineto{\pgfqpoint{1.752827in}{1.721627in}}%
\pgfpathlineto{\pgfqpoint{1.752998in}{1.946089in}}%
\pgfpathlineto{\pgfqpoint{1.753577in}{1.921276in}}%
\pgfpathlineto{\pgfqpoint{1.753626in}{1.928747in}}%
\pgfpathlineto{\pgfqpoint{1.754318in}{1.925561in}}%
\pgfpathlineto{\pgfqpoint{1.754467in}{1.923576in}}%
\pgfpathlineto{\pgfqpoint{1.754516in}{1.926495in}}%
\pgfpathlineto{\pgfqpoint{1.755110in}{1.924241in}}%
\pgfpathlineto{\pgfqpoint{1.755209in}{1.924388in}}%
\pgfpathlineto{\pgfqpoint{1.755357in}{1.887058in}}%
\pgfpathlineto{\pgfqpoint{1.755841in}{1.532030in}}%
\pgfpathlineto{\pgfqpoint{1.756647in}{1.532047in}}%
\pgfpathlineto{\pgfqpoint{1.757733in}{1.532125in}}%
\pgfpathlineto{\pgfqpoint{1.757919in}{1.946310in}}%
\pgfpathlineto{\pgfqpoint{1.758981in}{1.924823in}}%
\pgfpathlineto{\pgfqpoint{1.759030in}{1.925161in}}%
\pgfpathlineto{\pgfqpoint{1.759723in}{1.924798in}}%
\pgfpathlineto{\pgfqpoint{1.760168in}{1.923813in}}%
\pgfpathlineto{\pgfqpoint{1.760366in}{1.838363in}}%
\pgfpathlineto{\pgfqpoint{1.760807in}{1.532025in}}%
\pgfpathlineto{\pgfqpoint{1.761488in}{1.532051in}}%
\pgfpathlineto{\pgfqpoint{1.762668in}{1.532080in}}%
\pgfpathlineto{\pgfqpoint{1.762680in}{1.532146in}}%
\pgfpathlineto{\pgfqpoint{1.762907in}{1.948813in}}%
\pgfpathlineto{\pgfqpoint{1.763778in}{1.922706in}}%
\pgfpathlineto{\pgfqpoint{1.763926in}{1.928334in}}%
\pgfpathlineto{\pgfqpoint{1.763876in}{1.921368in}}%
\pgfpathlineto{\pgfqpoint{1.764569in}{1.924818in}}%
\pgfpathlineto{\pgfqpoint{1.764717in}{1.925907in}}%
\pgfpathlineto{\pgfqpoint{1.765113in}{1.923088in}}%
\pgfpathlineto{\pgfqpoint{1.765212in}{1.905753in}}%
\pgfpathlineto{\pgfqpoint{1.765727in}{1.532030in}}%
\pgfpathlineto{\pgfqpoint{1.766678in}{1.532052in}}%
\pgfpathlineto{\pgfqpoint{1.767624in}{1.532126in}}%
\pgfpathlineto{\pgfqpoint{1.767746in}{1.944943in}}%
\pgfpathlineto{\pgfqpoint{1.768865in}{1.924583in}}%
\pgfpathlineto{\pgfqpoint{1.768915in}{1.924934in}}%
\pgfpathlineto{\pgfqpoint{1.769607in}{1.924592in}}%
\pgfpathlineto{\pgfqpoint{1.770052in}{1.923687in}}%
\pgfpathlineto{\pgfqpoint{1.770201in}{1.878814in}}%
\pgfpathlineto{\pgfqpoint{1.770685in}{1.532024in}}%
\pgfpathlineto{\pgfqpoint{1.771450in}{1.532055in}}%
\pgfpathlineto{\pgfqpoint{1.772567in}{1.532105in}}%
\pgfpathlineto{\pgfqpoint{1.772608in}{1.711859in}}%
\pgfpathlineto{\pgfqpoint{1.772691in}{1.946129in}}%
\pgfpathlineto{\pgfqpoint{1.773367in}{1.920926in}}%
\pgfpathlineto{\pgfqpoint{1.773416in}{1.928177in}}%
\pgfpathlineto{\pgfqpoint{1.774109in}{1.925098in}}%
\pgfpathlineto{\pgfqpoint{1.774257in}{1.923160in}}%
\pgfpathlineto{\pgfqpoint{1.774307in}{1.925990in}}%
\pgfpathlineto{\pgfqpoint{1.774900in}{1.923785in}}%
\pgfpathlineto{\pgfqpoint{1.774999in}{1.923898in}}%
\pgfpathlineto{\pgfqpoint{1.775147in}{1.883912in}}%
\pgfpathlineto{\pgfqpoint{1.775638in}{1.532023in}}%
\pgfpathlineto{\pgfqpoint{1.776461in}{1.532041in}}%
\pgfpathlineto{\pgfqpoint{1.777521in}{1.532495in}}%
\pgfpathlineto{\pgfqpoint{1.777632in}{1.945717in}}%
\pgfpathlineto{\pgfqpoint{1.778534in}{1.924675in}}%
\pgfpathlineto{\pgfqpoint{1.778683in}{1.924191in}}%
\pgfpathlineto{\pgfqpoint{1.778633in}{1.924840in}}%
\pgfpathlineto{\pgfqpoint{1.779326in}{1.924428in}}%
\pgfpathlineto{\pgfqpoint{1.779968in}{1.923222in}}%
\pgfpathlineto{\pgfqpoint{1.780265in}{1.748293in}}%
\pgfpathlineto{\pgfqpoint{1.780595in}{1.532030in}}%
\pgfpathlineto{\pgfqpoint{1.781282in}{1.532047in}}%
\pgfpathlineto{\pgfqpoint{1.782461in}{1.532126in}}%
\pgfpathlineto{\pgfqpoint{1.782582in}{1.945672in}}%
\pgfpathlineto{\pgfqpoint{1.783701in}{1.924382in}}%
\pgfpathlineto{\pgfqpoint{1.784295in}{1.924392in}}%
\pgfpathlineto{\pgfqpoint{1.783800in}{1.924331in}}%
\pgfpathlineto{\pgfqpoint{1.784443in}{1.924205in}}%
\pgfpathlineto{\pgfqpoint{1.784839in}{1.923678in}}%
\pgfpathlineto{\pgfqpoint{1.784937in}{1.921704in}}%
\pgfpathlineto{\pgfqpoint{1.785405in}{1.632493in}}%
\pgfpathlineto{\pgfqpoint{1.785413in}{1.633635in}}%
\pgfpathlineto{\pgfqpoint{1.785487in}{1.532027in}}%
\pgfpathlineto{\pgfqpoint{1.786173in}{1.532041in}}%
\pgfpathlineto{\pgfqpoint{1.787407in}{1.532143in}}%
\pgfpathlineto{\pgfqpoint{1.787639in}{1.952796in}}%
\pgfpathlineto{\pgfqpoint{1.788560in}{1.924154in}}%
\pgfpathlineto{\pgfqpoint{1.788807in}{1.923323in}}%
\pgfpathlineto{\pgfqpoint{1.788758in}{1.925292in}}%
\pgfpathlineto{\pgfqpoint{1.789153in}{1.923649in}}%
\pgfpathlineto{\pgfqpoint{1.789203in}{1.924838in}}%
\pgfpathlineto{\pgfqpoint{1.789846in}{1.923191in}}%
\pgfpathlineto{\pgfqpoint{1.789994in}{1.874755in}}%
\pgfpathlineto{\pgfqpoint{1.790483in}{1.532031in}}%
\pgfpathlineto{\pgfqpoint{1.791216in}{1.532047in}}%
\pgfpathlineto{\pgfqpoint{1.792356in}{1.532284in}}%
\pgfpathlineto{\pgfqpoint{1.792457in}{1.943922in}}%
\pgfpathlineto{\pgfqpoint{1.793408in}{1.923606in}}%
\pgfpathlineto{\pgfqpoint{1.793705in}{1.924784in}}%
\pgfpathlineto{\pgfqpoint{1.794150in}{1.924484in}}%
\pgfpathlineto{\pgfqpoint{1.794892in}{1.903053in}}%
\pgfpathlineto{\pgfqpoint{1.795394in}{1.532025in}}%
\pgfpathlineto{\pgfqpoint{1.796341in}{1.532056in}}%
\pgfpathlineto{\pgfqpoint{1.797296in}{1.532113in}}%
\pgfpathlineto{\pgfqpoint{1.797335in}{1.718895in}}%
\pgfpathlineto{\pgfqpoint{1.797516in}{1.947720in}}%
\pgfpathlineto{\pgfqpoint{1.798112in}{1.925312in}}%
\pgfpathlineto{\pgfqpoint{1.798162in}{1.922483in}}%
\pgfpathlineto{\pgfqpoint{1.798211in}{1.925504in}}%
\pgfpathlineto{\pgfqpoint{1.798854in}{1.924094in}}%
\pgfpathlineto{\pgfqpoint{1.799002in}{1.924603in}}%
\pgfpathlineto{\pgfqpoint{1.799052in}{1.923426in}}%
\pgfpathlineto{\pgfqpoint{1.799448in}{1.924117in}}%
\pgfpathlineto{\pgfqpoint{1.799794in}{1.918473in}}%
\pgfpathlineto{\pgfqpoint{1.800348in}{1.532025in}}%
\pgfpathlineto{\pgfqpoint{1.801646in}{1.532053in}}%
\pgfpathlineto{\pgfqpoint{1.802241in}{1.532114in}}%
\pgfpathlineto{\pgfqpoint{1.802362in}{1.944342in}}%
\pgfpathlineto{\pgfqpoint{1.803801in}{1.923794in}}%
\pgfpathlineto{\pgfqpoint{1.803950in}{1.924477in}}%
\pgfpathlineto{\pgfqpoint{1.804444in}{1.923295in}}%
\pgfpathlineto{\pgfqpoint{1.804494in}{1.923780in}}%
\pgfpathlineto{\pgfqpoint{1.804741in}{1.917904in}}%
\pgfpathlineto{\pgfqpoint{1.805293in}{1.532025in}}%
\pgfpathlineto{\pgfqpoint{1.806591in}{1.532053in}}%
\pgfpathlineto{\pgfqpoint{1.807187in}{1.532114in}}%
\pgfpathlineto{\pgfqpoint{1.807306in}{1.945039in}}%
\pgfpathlineto{\pgfqpoint{1.808716in}{1.924238in}}%
\pgfpathlineto{\pgfqpoint{1.808963in}{1.924837in}}%
\pgfpathlineto{\pgfqpoint{1.809655in}{1.922617in}}%
\pgfpathlineto{\pgfqpoint{1.810001in}{1.708006in}}%
\pgfpathlineto{\pgfqpoint{1.810223in}{1.532026in}}%
\pgfpathlineto{\pgfqpoint{1.811013in}{1.532043in}}%
\pgfpathlineto{\pgfqpoint{1.812134in}{1.532134in}}%
\pgfpathlineto{\pgfqpoint{1.812368in}{1.951219in}}%
\pgfpathlineto{\pgfqpoint{1.813281in}{1.923875in}}%
\pgfpathlineto{\pgfqpoint{1.813430in}{1.922877in}}%
\pgfpathlineto{\pgfqpoint{1.813479in}{1.924691in}}%
\pgfpathlineto{\pgfqpoint{1.814073in}{1.923385in}}%
\pgfpathlineto{\pgfqpoint{1.814270in}{1.923766in}}%
\pgfpathlineto{\pgfqpoint{1.814567in}{1.922711in}}%
\pgfpathlineto{\pgfqpoint{1.814666in}{1.906606in}}%
\pgfpathlineto{\pgfqpoint{1.815180in}{1.532028in}}%
\pgfpathlineto{\pgfqpoint{1.816131in}{1.532054in}}%
\pgfpathlineto{\pgfqpoint{1.817078in}{1.532114in}}%
\pgfpathlineto{\pgfqpoint{1.817170in}{1.941068in}}%
\pgfpathlineto{\pgfqpoint{1.818417in}{1.923673in}}%
\pgfpathlineto{\pgfqpoint{1.819505in}{1.922746in}}%
\pgfpathlineto{\pgfqpoint{1.819604in}{1.910925in}}%
\pgfpathlineto{\pgfqpoint{1.820128in}{1.532025in}}%
\pgfpathlineto{\pgfqpoint{1.821129in}{1.532047in}}%
\pgfpathlineto{\pgfqpoint{1.822023in}{1.532114in}}%
\pgfpathlineto{\pgfqpoint{1.822143in}{1.944776in}}%
\pgfpathlineto{\pgfqpoint{1.823559in}{1.924513in}}%
\pgfpathlineto{\pgfqpoint{1.823856in}{1.922325in}}%
\pgfpathlineto{\pgfqpoint{1.823807in}{1.924649in}}%
\pgfpathlineto{\pgfqpoint{1.824400in}{1.922372in}}%
\pgfpathlineto{\pgfqpoint{1.824449in}{1.922903in}}%
\pgfpathlineto{\pgfqpoint{1.824499in}{1.922082in}}%
\pgfpathlineto{\pgfqpoint{1.824932in}{1.652774in}}%
\pgfpathlineto{\pgfqpoint{1.825071in}{1.532025in}}%
\pgfpathlineto{\pgfqpoint{1.825970in}{1.532044in}}%
\pgfpathlineto{\pgfqpoint{1.826969in}{1.532113in}}%
\pgfpathlineto{\pgfqpoint{1.827092in}{1.944502in}}%
\pgfpathlineto{\pgfqpoint{1.828553in}{1.922118in}}%
\pgfpathlineto{\pgfqpoint{1.828602in}{1.924750in}}%
\pgfpathlineto{\pgfqpoint{1.829295in}{1.923172in}}%
\pgfpathlineto{\pgfqpoint{1.829493in}{1.912035in}}%
\pgfpathlineto{\pgfqpoint{1.830016in}{1.532025in}}%
\pgfpathlineto{\pgfqpoint{1.831062in}{1.532059in}}%
\pgfpathlineto{\pgfqpoint{1.831914in}{1.532113in}}%
\pgfpathlineto{\pgfqpoint{1.832040in}{1.944451in}}%
\pgfpathlineto{\pgfqpoint{1.833518in}{1.923488in}}%
\pgfpathlineto{\pgfqpoint{1.833716in}{1.923929in}}%
\pgfpathlineto{\pgfqpoint{1.834408in}{1.920162in}}%
\pgfpathlineto{\pgfqpoint{1.834872in}{1.561942in}}%
\pgfpathlineto{\pgfqpoint{1.834969in}{1.532025in}}%
\pgfpathlineto{\pgfqpoint{1.835626in}{1.532051in}}%
\pgfpathlineto{\pgfqpoint{1.836859in}{1.532114in}}%
\pgfpathlineto{\pgfqpoint{1.836978in}{1.944308in}}%
\pgfpathlineto{\pgfqpoint{1.838389in}{1.924557in}}%
\pgfpathlineto{\pgfqpoint{1.838439in}{1.922117in}}%
\pgfpathlineto{\pgfqpoint{1.839180in}{1.923124in}}%
\pgfpathlineto{\pgfqpoint{1.839378in}{1.913200in}}%
\pgfpathlineto{\pgfqpoint{1.839925in}{1.532025in}}%
\pgfpathlineto{\pgfqpoint{1.841032in}{1.532049in}}%
\pgfpathlineto{\pgfqpoint{1.841805in}{1.532113in}}%
\pgfpathlineto{\pgfqpoint{1.841931in}{1.944132in}}%
\pgfpathlineto{\pgfqpoint{1.843404in}{1.924266in}}%
\pgfpathlineto{\pgfqpoint{1.844146in}{1.922207in}}%
\pgfpathlineto{\pgfqpoint{1.844195in}{1.922887in}}%
\pgfpathlineto{\pgfqpoint{1.844343in}{1.908127in}}%
\pgfpathlineto{\pgfqpoint{1.844864in}{1.532025in}}%
\pgfpathlineto{\pgfqpoint{1.845819in}{1.532057in}}%
\pgfpathlineto{\pgfqpoint{1.846750in}{1.532114in}}%
\pgfpathlineto{\pgfqpoint{1.846870in}{1.944239in}}%
\pgfpathlineto{\pgfqpoint{1.848319in}{1.923179in}}%
\pgfpathlineto{\pgfqpoint{1.848567in}{1.921899in}}%
\pgfpathlineto{\pgfqpoint{1.848517in}{1.924296in}}%
\pgfpathlineto{\pgfqpoint{1.849111in}{1.922114in}}%
\pgfpathlineto{\pgfqpoint{1.849160in}{1.922373in}}%
\pgfpathlineto{\pgfqpoint{1.849210in}{1.922066in}}%
\pgfpathlineto{\pgfqpoint{1.849457in}{1.797378in}}%
\pgfpathlineto{\pgfqpoint{1.849800in}{1.532025in}}%
\pgfpathlineto{\pgfqpoint{1.850550in}{1.532052in}}%
\pgfpathlineto{\pgfqpoint{1.851696in}{1.532113in}}%
\pgfpathlineto{\pgfqpoint{1.851818in}{1.944079in}}%
\pgfpathlineto{\pgfqpoint{1.853269in}{1.922178in}}%
\pgfpathlineto{\pgfqpoint{1.853318in}{1.923909in}}%
\pgfpathlineto{\pgfqpoint{1.854010in}{1.922546in}}%
\pgfpathlineto{\pgfqpoint{1.854208in}{1.915875in}}%
\pgfpathlineto{\pgfqpoint{1.854757in}{1.532025in}}%
\pgfpathlineto{\pgfqpoint{1.856012in}{1.532064in}}%
\pgfpathlineto{\pgfqpoint{1.856641in}{1.532114in}}%
\pgfpathlineto{\pgfqpoint{1.856759in}{1.943821in}}%
\pgfpathlineto{\pgfqpoint{1.858166in}{1.922923in}}%
\pgfpathlineto{\pgfqpoint{1.858858in}{1.921999in}}%
\pgfpathlineto{\pgfqpoint{1.858364in}{1.923856in}}%
\pgfpathlineto{\pgfqpoint{1.859007in}{1.922379in}}%
\pgfpathlineto{\pgfqpoint{1.859106in}{1.921462in}}%
\pgfpathlineto{\pgfqpoint{1.859304in}{1.833844in}}%
\pgfpathlineto{\pgfqpoint{1.859702in}{1.532025in}}%
\pgfpathlineto{\pgfqpoint{1.860411in}{1.532040in}}%
\pgfpathlineto{\pgfqpoint{1.861587in}{1.532113in}}%
\pgfpathlineto{\pgfqpoint{1.861710in}{1.943988in}}%
\pgfpathlineto{\pgfqpoint{1.863138in}{1.924110in}}%
\pgfpathlineto{\pgfqpoint{1.863879in}{1.921619in}}%
\pgfpathlineto{\pgfqpoint{1.863929in}{1.922937in}}%
\pgfpathlineto{\pgfqpoint{1.864127in}{1.907804in}}%
\pgfpathlineto{\pgfqpoint{1.864651in}{1.532023in}}%
\pgfpathlineto{\pgfqpoint{1.865662in}{1.532046in}}%
\pgfpathlineto{\pgfqpoint{1.866531in}{1.532104in}}%
\pgfpathlineto{\pgfqpoint{1.866572in}{1.712949in}}%
\pgfpathlineto{\pgfqpoint{1.866655in}{1.944043in}}%
\pgfpathlineto{\pgfqpoint{1.867350in}{1.925107in}}%
\pgfpathlineto{\pgfqpoint{1.867647in}{1.920191in}}%
\pgfpathlineto{\pgfqpoint{1.867696in}{1.925203in}}%
\pgfpathlineto{\pgfqpoint{1.868092in}{1.921156in}}%
\pgfpathlineto{\pgfqpoint{1.868141in}{1.924260in}}%
\pgfpathlineto{\pgfqpoint{1.868834in}{1.922581in}}%
\pgfpathlineto{\pgfqpoint{1.869081in}{1.903961in}}%
\pgfpathlineto{\pgfqpoint{1.869595in}{1.532024in}}%
\pgfpathlineto{\pgfqpoint{1.870556in}{1.532058in}}%
\pgfpathlineto{\pgfqpoint{1.871477in}{1.532106in}}%
\pgfpathlineto{\pgfqpoint{1.871518in}{1.723366in}}%
\pgfpathlineto{\pgfqpoint{1.871604in}{1.943017in}}%
\pgfpathlineto{\pgfqpoint{1.872260in}{1.922470in}}%
\pgfpathlineto{\pgfqpoint{1.872408in}{1.924468in}}%
\pgfpathlineto{\pgfqpoint{1.872458in}{1.920683in}}%
\pgfpathlineto{\pgfqpoint{1.872952in}{1.923804in}}%
\pgfpathlineto{\pgfqpoint{1.873002in}{1.921816in}}%
\pgfpathlineto{\pgfqpoint{1.873743in}{1.922700in}}%
\pgfpathlineto{\pgfqpoint{1.873991in}{1.916771in}}%
\pgfpathlineto{\pgfqpoint{1.874539in}{1.532025in}}%
\pgfpathlineto{\pgfqpoint{1.875793in}{1.532063in}}%
\pgfpathlineto{\pgfqpoint{1.876423in}{1.532113in}}%
\pgfpathlineto{\pgfqpoint{1.876463in}{1.723727in}}%
\pgfpathlineto{\pgfqpoint{1.876640in}{1.944795in}}%
\pgfpathlineto{\pgfqpoint{1.877205in}{1.919631in}}%
\pgfpathlineto{\pgfqpoint{1.877255in}{1.925155in}}%
\pgfpathlineto{\pgfqpoint{1.877947in}{1.923206in}}%
\pgfpathlineto{\pgfqpoint{1.878095in}{1.921457in}}%
\pgfpathlineto{\pgfqpoint{1.878046in}{1.923677in}}%
\pgfpathlineto{\pgfqpoint{1.878738in}{1.921987in}}%
\pgfpathlineto{\pgfqpoint{1.878837in}{1.922042in}}%
\pgfpathlineto{\pgfqpoint{1.878887in}{1.920992in}}%
\pgfpathlineto{\pgfqpoint{1.879084in}{1.835978in}}%
\pgfpathlineto{\pgfqpoint{1.879488in}{1.532023in}}%
\pgfpathlineto{\pgfqpoint{1.880202in}{1.532039in}}%
\pgfpathlineto{\pgfqpoint{1.881367in}{1.532104in}}%
\pgfpathlineto{\pgfqpoint{1.881408in}{1.713097in}}%
\pgfpathlineto{\pgfqpoint{1.881577in}{1.944744in}}%
\pgfpathlineto{\pgfqpoint{1.882186in}{1.924290in}}%
\pgfpathlineto{\pgfqpoint{1.882482in}{1.919955in}}%
\pgfpathlineto{\pgfqpoint{1.882433in}{1.924992in}}%
\pgfpathlineto{\pgfqpoint{1.882927in}{1.920933in}}%
\pgfpathlineto{\pgfqpoint{1.882977in}{1.923886in}}%
\pgfpathlineto{\pgfqpoint{1.883669in}{1.922437in}}%
\pgfpathlineto{\pgfqpoint{1.883916in}{1.905407in}}%
\pgfpathlineto{\pgfqpoint{1.884435in}{1.532023in}}%
\pgfpathlineto{\pgfqpoint{1.885398in}{1.532059in}}%
\pgfpathlineto{\pgfqpoint{1.886312in}{1.532103in}}%
\pgfpathlineto{\pgfqpoint{1.886353in}{1.704546in}}%
\pgfpathlineto{\pgfqpoint{1.886521in}{1.945210in}}%
\pgfpathlineto{\pgfqpoint{1.887118in}{1.922937in}}%
\pgfpathlineto{\pgfqpoint{1.887266in}{1.920566in}}%
\pgfpathlineto{\pgfqpoint{1.887316in}{1.924116in}}%
\pgfpathlineto{\pgfqpoint{1.887909in}{1.922146in}}%
\pgfpathlineto{\pgfqpoint{1.888107in}{1.923044in}}%
\pgfpathlineto{\pgfqpoint{1.888156in}{1.921631in}}%
\pgfpathlineto{\pgfqpoint{1.888552in}{1.922441in}}%
\pgfpathlineto{\pgfqpoint{1.888849in}{1.909904in}}%
\pgfpathlineto{\pgfqpoint{1.889375in}{1.532025in}}%
\pgfpathlineto{\pgfqpoint{1.890431in}{1.532059in}}%
\pgfpathlineto{\pgfqpoint{1.891259in}{1.532113in}}%
\pgfpathlineto{\pgfqpoint{1.891300in}{1.726597in}}%
\pgfpathlineto{\pgfqpoint{1.891475in}{1.943640in}}%
\pgfpathlineto{\pgfqpoint{1.892057in}{1.923626in}}%
\pgfpathlineto{\pgfqpoint{1.892354in}{1.920584in}}%
\pgfpathlineto{\pgfqpoint{1.892304in}{1.924003in}}%
\pgfpathlineto{\pgfqpoint{1.892799in}{1.921262in}}%
\pgfpathlineto{\pgfqpoint{1.892848in}{1.923314in}}%
\pgfpathlineto{\pgfqpoint{1.893541in}{1.922155in}}%
\pgfpathlineto{\pgfqpoint{1.893788in}{1.913145in}}%
\pgfpathlineto{\pgfqpoint{1.894325in}{1.532026in}}%
\pgfpathlineto{\pgfqpoint{1.895432in}{1.532050in}}%
\pgfpathlineto{\pgfqpoint{1.896205in}{1.532115in}}%
\pgfpathlineto{\pgfqpoint{1.896425in}{1.944858in}}%
\pgfpathlineto{\pgfqpoint{1.897771in}{1.920862in}}%
\pgfpathlineto{\pgfqpoint{1.897820in}{1.923476in}}%
\pgfpathlineto{\pgfqpoint{1.898513in}{1.922131in}}%
\pgfpathlineto{\pgfqpoint{1.898760in}{1.902523in}}%
\pgfpathlineto{\pgfqpoint{1.899266in}{1.532025in}}%
\pgfpathlineto{\pgfqpoint{1.900224in}{1.532057in}}%
\pgfpathlineto{\pgfqpoint{1.901150in}{1.532113in}}%
\pgfpathlineto{\pgfqpoint{1.901190in}{1.719176in}}%
\pgfpathlineto{\pgfqpoint{1.901365in}{1.945514in}}%
\pgfpathlineto{\pgfqpoint{1.901942in}{1.918689in}}%
\pgfpathlineto{\pgfqpoint{1.901991in}{1.925265in}}%
\pgfpathlineto{\pgfqpoint{1.902683in}{1.922878in}}%
\pgfpathlineto{\pgfqpoint{1.902832in}{1.920832in}}%
\pgfpathlineto{\pgfqpoint{1.902782in}{1.923459in}}%
\pgfpathlineto{\pgfqpoint{1.903475in}{1.921537in}}%
\pgfpathlineto{\pgfqpoint{1.903574in}{1.921650in}}%
\pgfpathlineto{\pgfqpoint{1.903623in}{1.920428in}}%
\pgfpathlineto{\pgfqpoint{1.903821in}{1.831983in}}%
\pgfpathlineto{\pgfqpoint{1.904216in}{1.532026in}}%
\pgfpathlineto{\pgfqpoint{1.904926in}{1.532041in}}%
\pgfpathlineto{\pgfqpoint{1.906096in}{1.532114in}}%
\pgfpathlineto{\pgfqpoint{1.906150in}{1.916854in}}%
\pgfpathlineto{\pgfqpoint{1.906312in}{1.945291in}}%
\pgfpathlineto{\pgfqpoint{1.906262in}{1.896030in}}%
\pgfpathlineto{\pgfqpoint{1.906886in}{1.923409in}}%
\pgfpathlineto{\pgfqpoint{1.907182in}{1.920370in}}%
\pgfpathlineto{\pgfqpoint{1.907133in}{1.923705in}}%
\pgfpathlineto{\pgfqpoint{1.907627in}{1.921036in}}%
\pgfpathlineto{\pgfqpoint{1.907677in}{1.923063in}}%
\pgfpathlineto{\pgfqpoint{1.908369in}{1.921908in}}%
\pgfpathlineto{\pgfqpoint{1.908616in}{1.914957in}}%
\pgfpathlineto{\pgfqpoint{1.909162in}{1.532026in}}%
\pgfpathlineto{\pgfqpoint{1.910367in}{1.532053in}}%
\pgfpathlineto{\pgfqpoint{1.911042in}{1.532116in}}%
\pgfpathlineto{\pgfqpoint{1.911261in}{1.943609in}}%
\pgfpathlineto{\pgfqpoint{1.912567in}{1.921748in}}%
\pgfpathlineto{\pgfqpoint{1.912716in}{1.922706in}}%
\pgfpathlineto{\pgfqpoint{1.913210in}{1.921309in}}%
\pgfpathlineto{\pgfqpoint{1.913309in}{1.921369in}}%
\pgfpathlineto{\pgfqpoint{1.913358in}{1.921566in}}%
\pgfpathlineto{\pgfqpoint{1.913457in}{1.921051in}}%
\pgfpathlineto{\pgfqpoint{1.913507in}{1.921107in}}%
\pgfpathlineto{\pgfqpoint{1.913705in}{1.838146in}}%
\pgfpathlineto{\pgfqpoint{1.914107in}{1.532026in}}%
\pgfpathlineto{\pgfqpoint{1.914867in}{1.532052in}}%
\pgfpathlineto{\pgfqpoint{1.915987in}{1.532114in}}%
\pgfpathlineto{\pgfqpoint{1.916209in}{1.944901in}}%
\pgfpathlineto{\pgfqpoint{1.917579in}{1.922682in}}%
\pgfpathlineto{\pgfqpoint{1.917629in}{1.921124in}}%
\pgfpathlineto{\pgfqpoint{1.917678in}{1.922703in}}%
\pgfpathlineto{\pgfqpoint{1.918371in}{1.921393in}}%
\pgfpathlineto{\pgfqpoint{1.918519in}{1.911820in}}%
\pgfpathlineto{\pgfqpoint{1.919057in}{1.532023in}}%
\pgfpathlineto{\pgfqpoint{1.920174in}{1.532048in}}%
\pgfpathlineto{\pgfqpoint{1.920939in}{1.532678in}}%
\pgfpathlineto{\pgfqpoint{1.921053in}{1.942759in}}%
\pgfpathlineto{\pgfqpoint{1.921936in}{1.920719in}}%
\pgfpathlineto{\pgfqpoint{1.921985in}{1.923200in}}%
\pgfpathlineto{\pgfqpoint{1.922034in}{1.920460in}}%
\pgfpathlineto{\pgfqpoint{1.922677in}{1.921729in}}%
\pgfpathlineto{\pgfqpoint{1.922875in}{1.922289in}}%
\pgfpathlineto{\pgfqpoint{1.923419in}{1.920508in}}%
\pgfpathlineto{\pgfqpoint{1.923716in}{1.743049in}}%
\pgfpathlineto{\pgfqpoint{1.924003in}{1.532023in}}%
\pgfpathlineto{\pgfqpoint{1.924724in}{1.532039in}}%
\pgfpathlineto{\pgfqpoint{1.925885in}{1.532621in}}%
\pgfpathlineto{\pgfqpoint{1.925999in}{1.943063in}}%
\pgfpathlineto{\pgfqpoint{1.926880in}{1.919921in}}%
\pgfpathlineto{\pgfqpoint{1.926929in}{1.923880in}}%
\pgfpathlineto{\pgfqpoint{1.926979in}{1.919672in}}%
\pgfpathlineto{\pgfqpoint{1.927622in}{1.921702in}}%
\pgfpathlineto{\pgfqpoint{1.928314in}{1.920399in}}%
\pgfpathlineto{\pgfqpoint{1.927820in}{1.922507in}}%
\pgfpathlineto{\pgfqpoint{1.928364in}{1.920600in}}%
\pgfpathlineto{\pgfqpoint{1.928660in}{1.745219in}}%
\pgfpathlineto{\pgfqpoint{1.928921in}{1.532029in}}%
\pgfpathlineto{\pgfqpoint{1.929667in}{1.532048in}}%
\pgfpathlineto{\pgfqpoint{1.930820in}{1.532095in}}%
\pgfpathlineto{\pgfqpoint{1.930829in}{1.532286in}}%
\pgfpathlineto{\pgfqpoint{1.930948in}{1.942767in}}%
\pgfpathlineto{\pgfqpoint{1.931849in}{1.920002in}}%
\pgfpathlineto{\pgfqpoint{1.932146in}{1.923972in}}%
\pgfpathlineto{\pgfqpoint{1.932097in}{1.919535in}}%
\pgfpathlineto{\pgfqpoint{1.932591in}{1.923062in}}%
\pgfpathlineto{\pgfqpoint{1.933481in}{1.842224in}}%
\pgfpathlineto{\pgfqpoint{1.933871in}{1.532026in}}%
\pgfpathlineto{\pgfqpoint{1.934620in}{1.532051in}}%
\pgfpathlineto{\pgfqpoint{1.935768in}{1.532113in}}%
\pgfpathlineto{\pgfqpoint{1.935809in}{1.722592in}}%
\pgfpathlineto{\pgfqpoint{1.935986in}{1.943984in}}%
\pgfpathlineto{\pgfqpoint{1.936577in}{1.923477in}}%
\pgfpathlineto{\pgfqpoint{1.936874in}{1.919841in}}%
\pgfpathlineto{\pgfqpoint{1.937319in}{1.920517in}}%
\pgfpathlineto{\pgfqpoint{1.937368in}{1.922743in}}%
\pgfpathlineto{\pgfqpoint{1.938060in}{1.921362in}}%
\pgfpathlineto{\pgfqpoint{1.938308in}{1.911594in}}%
\pgfpathlineto{\pgfqpoint{1.938838in}{1.532023in}}%
\pgfpathlineto{\pgfqpoint{1.939904in}{1.532061in}}%
\pgfpathlineto{\pgfqpoint{1.940712in}{1.532100in}}%
\pgfpathlineto{\pgfqpoint{1.940739in}{1.573477in}}%
\pgfpathlineto{\pgfqpoint{1.940824in}{1.941092in}}%
\pgfpathlineto{\pgfqpoint{1.941504in}{1.924937in}}%
\pgfpathlineto{\pgfqpoint{1.941653in}{1.915629in}}%
\pgfpathlineto{\pgfqpoint{1.941603in}{1.927364in}}%
\pgfpathlineto{\pgfqpoint{1.942296in}{1.921388in}}%
\pgfpathlineto{\pgfqpoint{1.942543in}{1.919241in}}%
\pgfpathlineto{\pgfqpoint{1.942493in}{1.923873in}}%
\pgfpathlineto{\pgfqpoint{1.942988in}{1.919738in}}%
\pgfpathlineto{\pgfqpoint{1.943037in}{1.922142in}}%
\pgfpathlineto{\pgfqpoint{1.943235in}{1.911590in}}%
\pgfpathlineto{\pgfqpoint{1.943780in}{1.532025in}}%
\pgfpathlineto{\pgfqpoint{1.945182in}{1.532056in}}%
\pgfpathlineto{\pgfqpoint{1.945659in}{1.532113in}}%
\pgfpathlineto{\pgfqpoint{1.945700in}{1.726021in}}%
\pgfpathlineto{\pgfqpoint{1.945874in}{1.942863in}}%
\pgfpathlineto{\pgfqpoint{1.946439in}{1.922629in}}%
\pgfpathlineto{\pgfqpoint{1.946736in}{1.920010in}}%
\pgfpathlineto{\pgfqpoint{1.946687in}{1.922867in}}%
\pgfpathlineto{\pgfqpoint{1.947181in}{1.920594in}}%
\pgfpathlineto{\pgfqpoint{1.947231in}{1.922338in}}%
\pgfpathlineto{\pgfqpoint{1.947923in}{1.921298in}}%
\pgfpathlineto{\pgfqpoint{1.948170in}{1.918459in}}%
\pgfpathlineto{\pgfqpoint{1.948616in}{1.588621in}}%
\pgfpathlineto{\pgfqpoint{1.948687in}{1.532022in}}%
\pgfpathlineto{\pgfqpoint{1.949366in}{1.532036in}}%
\pgfpathlineto{\pgfqpoint{1.950610in}{1.532228in}}%
\pgfpathlineto{\pgfqpoint{1.950724in}{1.942337in}}%
\pgfpathlineto{\pgfqpoint{1.951678in}{1.921423in}}%
\pgfpathlineto{\pgfqpoint{1.951777in}{1.921521in}}%
\pgfpathlineto{\pgfqpoint{1.953062in}{1.920346in}}%
\pgfpathlineto{\pgfqpoint{1.953211in}{1.877280in}}%
\pgfpathlineto{\pgfqpoint{1.953672in}{1.532023in}}%
\pgfpathlineto{\pgfqpoint{1.954486in}{1.532041in}}%
\pgfpathlineto{\pgfqpoint{1.955549in}{1.532102in}}%
\pgfpathlineto{\pgfqpoint{1.955589in}{1.703542in}}%
\pgfpathlineto{\pgfqpoint{1.955671in}{1.942621in}}%
\pgfpathlineto{\pgfqpoint{1.956346in}{1.922628in}}%
\pgfpathlineto{\pgfqpoint{1.956494in}{1.919296in}}%
\pgfpathlineto{\pgfqpoint{1.956445in}{1.923349in}}%
\pgfpathlineto{\pgfqpoint{1.957088in}{1.921704in}}%
\pgfpathlineto{\pgfqpoint{1.957335in}{1.922124in}}%
\pgfpathlineto{\pgfqpoint{1.958077in}{1.913636in}}%
\pgfpathlineto{\pgfqpoint{1.958626in}{1.532023in}}%
\pgfpathlineto{\pgfqpoint{1.959943in}{1.532052in}}%
\pgfpathlineto{\pgfqpoint{1.960494in}{1.532101in}}%
\pgfpathlineto{\pgfqpoint{1.960521in}{1.574272in}}%
\pgfpathlineto{\pgfqpoint{1.960609in}{1.941574in}}%
\pgfpathlineto{\pgfqpoint{1.961317in}{1.920326in}}%
\pgfpathlineto{\pgfqpoint{1.961465in}{1.922656in}}%
\pgfpathlineto{\pgfqpoint{1.961416in}{1.919800in}}%
\pgfpathlineto{\pgfqpoint{1.962108in}{1.921319in}}%
\pgfpathlineto{\pgfqpoint{1.962355in}{1.921747in}}%
\pgfpathlineto{\pgfqpoint{1.962306in}{1.920656in}}%
\pgfpathlineto{\pgfqpoint{1.962702in}{1.921263in}}%
\pgfpathlineto{\pgfqpoint{1.963048in}{1.906686in}}%
\pgfpathlineto{\pgfqpoint{1.963564in}{1.532023in}}%
\pgfpathlineto{\pgfqpoint{1.964529in}{1.532059in}}%
\pgfpathlineto{\pgfqpoint{1.965440in}{1.532102in}}%
\pgfpathlineto{\pgfqpoint{1.965480in}{1.703302in}}%
\pgfpathlineto{\pgfqpoint{1.965648in}{1.942696in}}%
\pgfpathlineto{\pgfqpoint{1.966235in}{1.922515in}}%
\pgfpathlineto{\pgfqpoint{1.966284in}{1.919377in}}%
\pgfpathlineto{\pgfqpoint{1.966333in}{1.922962in}}%
\pgfpathlineto{\pgfqpoint{1.966976in}{1.921387in}}%
\pgfpathlineto{\pgfqpoint{1.967075in}{1.920784in}}%
\pgfpathlineto{\pgfqpoint{1.967125in}{1.921811in}}%
\pgfpathlineto{\pgfqpoint{1.967224in}{1.921876in}}%
\pgfpathlineto{\pgfqpoint{1.968015in}{1.897673in}}%
\pgfpathlineto{\pgfqpoint{1.968509in}{1.532023in}}%
\pgfpathlineto{\pgfqpoint{1.969425in}{1.532044in}}%
\pgfpathlineto{\pgfqpoint{1.970385in}{1.532103in}}%
\pgfpathlineto{\pgfqpoint{1.970426in}{1.707931in}}%
\pgfpathlineto{\pgfqpoint{1.970594in}{1.944183in}}%
\pgfpathlineto{\pgfqpoint{1.971184in}{1.923910in}}%
\pgfpathlineto{\pgfqpoint{1.971234in}{1.918289in}}%
\pgfpathlineto{\pgfqpoint{1.971926in}{1.920624in}}%
\pgfpathlineto{\pgfqpoint{1.972074in}{1.922267in}}%
\pgfpathlineto{\pgfqpoint{1.972025in}{1.920055in}}%
\pgfpathlineto{\pgfqpoint{1.972668in}{1.920384in}}%
\pgfpathlineto{\pgfqpoint{1.972717in}{1.920695in}}%
\pgfpathlineto{\pgfqpoint{1.972816in}{1.920028in}}%
\pgfpathlineto{\pgfqpoint{1.972866in}{1.920323in}}%
\pgfpathlineto{\pgfqpoint{1.973113in}{1.794795in}}%
\pgfpathlineto{\pgfqpoint{1.973454in}{1.532023in}}%
\pgfpathlineto{\pgfqpoint{1.974169in}{1.532039in}}%
\pgfpathlineto{\pgfqpoint{1.975331in}{1.532103in}}%
\pgfpathlineto{\pgfqpoint{1.975371in}{1.707086in}}%
\pgfpathlineto{\pgfqpoint{1.975540in}{1.944178in}}%
\pgfpathlineto{\pgfqpoint{1.976138in}{1.923981in}}%
\pgfpathlineto{\pgfqpoint{1.976188in}{1.918175in}}%
\pgfpathlineto{\pgfqpoint{1.976880in}{1.920446in}}%
\pgfpathlineto{\pgfqpoint{1.977029in}{1.922241in}}%
\pgfpathlineto{\pgfqpoint{1.976979in}{1.919915in}}%
\pgfpathlineto{\pgfqpoint{1.977622in}{1.920337in}}%
\pgfpathlineto{\pgfqpoint{1.977671in}{1.920572in}}%
\pgfpathlineto{\pgfqpoint{1.977770in}{1.919883in}}%
\pgfpathlineto{\pgfqpoint{1.977820in}{1.920179in}}%
\pgfpathlineto{\pgfqpoint{1.978117in}{1.749439in}}%
\pgfpathlineto{\pgfqpoint{1.978402in}{1.532023in}}%
\pgfpathlineto{\pgfqpoint{1.979170in}{1.532055in}}%
\pgfpathlineto{\pgfqpoint{1.980276in}{1.532102in}}%
\pgfpathlineto{\pgfqpoint{1.980304in}{1.576235in}}%
\pgfpathlineto{\pgfqpoint{1.980396in}{1.941993in}}%
\pgfpathlineto{\pgfqpoint{1.981101in}{1.915790in}}%
\pgfpathlineto{\pgfqpoint{1.981150in}{1.925413in}}%
\pgfpathlineto{\pgfqpoint{1.981842in}{1.922456in}}%
\pgfpathlineto{\pgfqpoint{1.981991in}{1.919023in}}%
\pgfpathlineto{\pgfqpoint{1.981941in}{1.923057in}}%
\pgfpathlineto{\pgfqpoint{1.982634in}{1.920487in}}%
\pgfpathlineto{\pgfqpoint{1.982733in}{1.920672in}}%
\pgfpathlineto{\pgfqpoint{1.982782in}{1.917663in}}%
\pgfpathlineto{\pgfqpoint{1.983128in}{1.700168in}}%
\pgfpathlineto{\pgfqpoint{1.983304in}{1.532023in}}%
\pgfpathlineto{\pgfqpoint{1.984131in}{1.532055in}}%
\pgfpathlineto{\pgfqpoint{1.985220in}{1.532097in}}%
\pgfpathlineto{\pgfqpoint{1.985230in}{1.532420in}}%
\pgfpathlineto{\pgfqpoint{1.985341in}{1.941958in}}%
\pgfpathlineto{\pgfqpoint{1.986217in}{1.921553in}}%
\pgfpathlineto{\pgfqpoint{1.986267in}{1.920234in}}%
\pgfpathlineto{\pgfqpoint{1.986316in}{1.921572in}}%
\pgfpathlineto{\pgfqpoint{1.986959in}{1.920867in}}%
\pgfpathlineto{\pgfqpoint{1.987108in}{1.921119in}}%
\pgfpathlineto{\pgfqpoint{1.987503in}{1.920346in}}%
\pgfpathlineto{\pgfqpoint{1.987553in}{1.920566in}}%
\pgfpathlineto{\pgfqpoint{1.987751in}{1.913515in}}%
\pgfpathlineto{\pgfqpoint{1.988292in}{1.532023in}}%
\pgfpathlineto{\pgfqpoint{1.989554in}{1.532066in}}%
\pgfpathlineto{\pgfqpoint{1.990167in}{1.532103in}}%
\pgfpathlineto{\pgfqpoint{1.990206in}{1.696848in}}%
\pgfpathlineto{\pgfqpoint{1.990289in}{1.942076in}}%
\pgfpathlineto{\pgfqpoint{1.990950in}{1.923640in}}%
\pgfpathlineto{\pgfqpoint{1.990999in}{1.918088in}}%
\pgfpathlineto{\pgfqpoint{1.991692in}{1.920251in}}%
\pgfpathlineto{\pgfqpoint{1.991840in}{1.921992in}}%
\pgfpathlineto{\pgfqpoint{1.991791in}{1.919756in}}%
\pgfpathlineto{\pgfqpoint{1.992434in}{1.920194in}}%
\pgfpathlineto{\pgfqpoint{1.992532in}{1.920402in}}%
\pgfpathlineto{\pgfqpoint{1.992582in}{1.919807in}}%
\pgfpathlineto{\pgfqpoint{1.992631in}{1.920257in}}%
\pgfpathlineto{\pgfqpoint{1.992780in}{1.879793in}}%
\pgfpathlineto{\pgfqpoint{1.993242in}{1.532023in}}%
\pgfpathlineto{\pgfqpoint{1.994062in}{1.532041in}}%
\pgfpathlineto{\pgfqpoint{1.995121in}{1.532558in}}%
\pgfpathlineto{\pgfqpoint{1.995231in}{1.941674in}}%
\pgfpathlineto{\pgfqpoint{1.996112in}{1.919293in}}%
\pgfpathlineto{\pgfqpoint{1.996161in}{1.922137in}}%
\pgfpathlineto{\pgfqpoint{1.996853in}{1.921139in}}%
\pgfpathlineto{\pgfqpoint{1.996952in}{1.921323in}}%
\pgfpathlineto{\pgfqpoint{1.997694in}{1.894876in}}%
\pgfpathlineto{\pgfqpoint{1.998190in}{1.532023in}}%
\pgfpathlineto{\pgfqpoint{1.999063in}{1.532057in}}%
\pgfpathlineto{\pgfqpoint{2.000057in}{1.532100in}}%
\pgfpathlineto{\pgfqpoint{2.000084in}{1.572763in}}%
\pgfpathlineto{\pgfqpoint{2.000186in}{1.941953in}}%
\pgfpathlineto{\pgfqpoint{2.000889in}{1.927211in}}%
\pgfpathlineto{\pgfqpoint{2.000939in}{1.915869in}}%
\pgfpathlineto{\pgfqpoint{2.001631in}{1.917856in}}%
\pgfpathlineto{\pgfqpoint{2.001681in}{1.923854in}}%
\pgfpathlineto{\pgfqpoint{2.001730in}{1.917826in}}%
\pgfpathlineto{\pgfqpoint{2.002373in}{1.920432in}}%
\pgfpathlineto{\pgfqpoint{2.002472in}{1.920986in}}%
\pgfpathlineto{\pgfqpoint{2.002620in}{1.905536in}}%
\pgfpathlineto{\pgfqpoint{2.003089in}{1.532022in}}%
\pgfpathlineto{\pgfqpoint{2.004116in}{1.532060in}}%
\pgfpathlineto{\pgfqpoint{2.005002in}{1.532095in}}%
\pgfpathlineto{\pgfqpoint{2.005010in}{1.532214in}}%
\pgfpathlineto{\pgfqpoint{2.005124in}{1.941629in}}%
\pgfpathlineto{\pgfqpoint{2.006093in}{1.920594in}}%
\pgfpathlineto{\pgfqpoint{2.006588in}{1.920750in}}%
\pgfpathlineto{\pgfqpoint{2.006885in}{1.920585in}}%
\pgfpathlineto{\pgfqpoint{2.007379in}{1.920056in}}%
\pgfpathlineto{\pgfqpoint{2.007478in}{1.919579in}}%
\pgfpathlineto{\pgfqpoint{2.007676in}{1.838844in}}%
\pgfpathlineto{\pgfqpoint{2.008078in}{1.532023in}}%
\pgfpathlineto{\pgfqpoint{2.008800in}{1.532039in}}%
\pgfpathlineto{\pgfqpoint{2.009957in}{1.532523in}}%
\pgfpathlineto{\pgfqpoint{2.010069in}{1.941649in}}%
\pgfpathlineto{\pgfqpoint{2.010975in}{1.921075in}}%
\pgfpathlineto{\pgfqpoint{2.011272in}{1.920134in}}%
\pgfpathlineto{\pgfqpoint{2.011321in}{1.921128in}}%
\pgfpathlineto{\pgfqpoint{2.011717in}{1.920266in}}%
\pgfpathlineto{\pgfqpoint{2.011766in}{1.920861in}}%
\pgfpathlineto{\pgfqpoint{2.012409in}{1.919616in}}%
\pgfpathlineto{\pgfqpoint{2.012508in}{1.907106in}}%
\pgfpathlineto{\pgfqpoint{2.013026in}{1.532023in}}%
\pgfpathlineto{\pgfqpoint{2.014046in}{1.532045in}}%
\pgfpathlineto{\pgfqpoint{2.014902in}{1.532422in}}%
\pgfpathlineto{\pgfqpoint{2.015014in}{1.941577in}}%
\pgfpathlineto{\pgfqpoint{2.015895in}{1.920244in}}%
\pgfpathlineto{\pgfqpoint{2.016043in}{1.920899in}}%
\pgfpathlineto{\pgfqpoint{2.015994in}{1.920181in}}%
\pgfpathlineto{\pgfqpoint{2.016686in}{1.920484in}}%
\pgfpathlineto{\pgfqpoint{2.016834in}{1.920584in}}%
\pgfpathlineto{\pgfqpoint{2.017230in}{1.919934in}}%
\pgfpathlineto{\pgfqpoint{2.017279in}{1.919979in}}%
\pgfpathlineto{\pgfqpoint{2.017378in}{1.919401in}}%
\pgfpathlineto{\pgfqpoint{2.017626in}{1.794416in}}%
\pgfpathlineto{\pgfqpoint{2.017920in}{1.532023in}}%
\pgfpathlineto{\pgfqpoint{2.018696in}{1.532039in}}%
\pgfpathlineto{\pgfqpoint{2.019849in}{1.533059in}}%
\pgfpathlineto{\pgfqpoint{2.019954in}{1.940337in}}%
\pgfpathlineto{\pgfqpoint{2.020834in}{1.922860in}}%
\pgfpathlineto{\pgfqpoint{2.021131in}{1.918689in}}%
\pgfpathlineto{\pgfqpoint{2.021576in}{1.919300in}}%
\pgfpathlineto{\pgfqpoint{2.021626in}{1.921660in}}%
\pgfpathlineto{\pgfqpoint{2.022318in}{1.919551in}}%
\pgfpathlineto{\pgfqpoint{2.022565in}{1.802018in}}%
\pgfpathlineto{\pgfqpoint{2.022914in}{1.532023in}}%
\pgfpathlineto{\pgfqpoint{2.023635in}{1.532039in}}%
\pgfpathlineto{\pgfqpoint{2.024794in}{1.533159in}}%
\pgfpathlineto{\pgfqpoint{2.024903in}{1.941008in}}%
\pgfpathlineto{\pgfqpoint{2.025761in}{1.920347in}}%
\pgfpathlineto{\pgfqpoint{2.025860in}{1.920603in}}%
\pgfpathlineto{\pgfqpoint{2.027245in}{1.919399in}}%
\pgfpathlineto{\pgfqpoint{2.027344in}{1.907409in}}%
\pgfpathlineto{\pgfqpoint{2.027862in}{1.532023in}}%
\pgfpathlineto{\pgfqpoint{2.028882in}{1.532045in}}%
\pgfpathlineto{\pgfqpoint{2.029739in}{1.532451in}}%
\pgfpathlineto{\pgfqpoint{2.029850in}{1.941396in}}%
\pgfpathlineto{\pgfqpoint{2.030765in}{1.920471in}}%
\pgfpathlineto{\pgfqpoint{2.030913in}{1.919871in}}%
\pgfpathlineto{\pgfqpoint{2.030963in}{1.920900in}}%
\pgfpathlineto{\pgfqpoint{2.031556in}{1.920214in}}%
\pgfpathlineto{\pgfqpoint{2.031754in}{1.920396in}}%
\pgfpathlineto{\pgfqpoint{2.031803in}{1.919966in}}%
\pgfpathlineto{\pgfqpoint{2.032100in}{1.919814in}}%
\pgfpathlineto{\pgfqpoint{2.032248in}{1.916776in}}%
\pgfpathlineto{\pgfqpoint{2.032688in}{1.597821in}}%
\pgfpathlineto{\pgfqpoint{2.032761in}{1.532023in}}%
\pgfpathlineto{\pgfqpoint{2.033440in}{1.532037in}}%
\pgfpathlineto{\pgfqpoint{2.034683in}{1.532258in}}%
\pgfpathlineto{\pgfqpoint{2.034798in}{1.941301in}}%
\pgfpathlineto{\pgfqpoint{2.035719in}{1.920078in}}%
\pgfpathlineto{\pgfqpoint{2.036016in}{1.920553in}}%
\pgfpathlineto{\pgfqpoint{2.036461in}{1.920410in}}%
\pgfpathlineto{\pgfqpoint{2.037203in}{1.915045in}}%
\pgfpathlineto{\pgfqpoint{2.037640in}{1.547772in}}%
\pgfpathlineto{\pgfqpoint{2.037703in}{1.532023in}}%
\pgfpathlineto{\pgfqpoint{2.038380in}{1.532037in}}%
\pgfpathlineto{\pgfqpoint{2.039630in}{1.532901in}}%
\pgfpathlineto{\pgfqpoint{2.039734in}{1.939324in}}%
\pgfpathlineto{\pgfqpoint{2.040620in}{1.925483in}}%
\pgfpathlineto{\pgfqpoint{2.041016in}{1.916435in}}%
\pgfpathlineto{\pgfqpoint{2.041362in}{1.917782in}}%
\pgfpathlineto{\pgfqpoint{2.041411in}{1.922789in}}%
\pgfpathlineto{\pgfqpoint{2.042104in}{1.919421in}}%
\pgfpathlineto{\pgfqpoint{2.042401in}{1.760789in}}%
\pgfpathlineto{\pgfqpoint{2.042694in}{1.532023in}}%
\pgfpathlineto{\pgfqpoint{2.043463in}{1.532054in}}%
\pgfpathlineto{\pgfqpoint{2.044567in}{1.532105in}}%
\pgfpathlineto{\pgfqpoint{2.044608in}{1.710906in}}%
\pgfpathlineto{\pgfqpoint{2.044777in}{1.942215in}}%
\pgfpathlineto{\pgfqpoint{2.045369in}{1.919193in}}%
\pgfpathlineto{\pgfqpoint{2.045517in}{1.922994in}}%
\pgfpathlineto{\pgfqpoint{2.045567in}{1.917421in}}%
\pgfpathlineto{\pgfqpoint{2.046160in}{1.920684in}}%
\pgfpathlineto{\pgfqpoint{2.046407in}{1.921270in}}%
\pgfpathlineto{\pgfqpoint{2.047100in}{1.915811in}}%
\pgfpathlineto{\pgfqpoint{2.047602in}{1.532022in}}%
\pgfpathlineto{\pgfqpoint{2.048982in}{1.532052in}}%
\pgfpathlineto{\pgfqpoint{2.049518in}{1.532192in}}%
\pgfpathlineto{\pgfqpoint{2.049634in}{1.940916in}}%
\pgfpathlineto{\pgfqpoint{2.050635in}{1.920144in}}%
\pgfpathlineto{\pgfqpoint{2.051575in}{1.919971in}}%
\pgfpathlineto{\pgfqpoint{2.052020in}{1.918377in}}%
\pgfpathlineto{\pgfqpoint{2.052413in}{1.673075in}}%
\pgfpathlineto{\pgfqpoint{2.052540in}{1.532023in}}%
\pgfpathlineto{\pgfqpoint{2.053416in}{1.532041in}}%
\pgfpathlineto{\pgfqpoint{2.054466in}{1.532658in}}%
\pgfpathlineto{\pgfqpoint{2.054583in}{1.941410in}}%
\pgfpathlineto{\pgfqpoint{2.055479in}{1.917961in}}%
\pgfpathlineto{\pgfqpoint{2.055627in}{1.923679in}}%
\pgfpathlineto{\pgfqpoint{2.055578in}{1.916528in}}%
\pgfpathlineto{\pgfqpoint{2.056270in}{1.920162in}}%
\pgfpathlineto{\pgfqpoint{2.056517in}{1.921274in}}%
\pgfpathlineto{\pgfqpoint{2.056468in}{1.918518in}}%
\pgfpathlineto{\pgfqpoint{2.056863in}{1.920223in}}%
\pgfpathlineto{\pgfqpoint{2.057061in}{1.888487in}}%
\pgfpathlineto{\pgfqpoint{2.057493in}{1.532022in}}%
\pgfpathlineto{\pgfqpoint{2.058378in}{1.532041in}}%
\pgfpathlineto{\pgfqpoint{2.059409in}{1.532192in}}%
\pgfpathlineto{\pgfqpoint{2.059525in}{1.940906in}}%
\pgfpathlineto{\pgfqpoint{2.060482in}{1.920005in}}%
\pgfpathlineto{\pgfqpoint{2.061471in}{1.919854in}}%
\pgfpathlineto{\pgfqpoint{2.061866in}{1.919065in}}%
\pgfpathlineto{\pgfqpoint{2.061965in}{1.906566in}}%
\pgfpathlineto{\pgfqpoint{2.062432in}{1.532023in}}%
\pgfpathlineto{\pgfqpoint{2.063506in}{1.532046in}}%
\pgfpathlineto{\pgfqpoint{2.064357in}{1.532516in}}%
\pgfpathlineto{\pgfqpoint{2.064469in}{1.941005in}}%
\pgfpathlineto{\pgfqpoint{2.065334in}{1.919053in}}%
\pgfpathlineto{\pgfqpoint{2.065384in}{1.920720in}}%
\pgfpathlineto{\pgfqpoint{2.066076in}{1.920249in}}%
\pgfpathlineto{\pgfqpoint{2.066175in}{1.920300in}}%
\pgfpathlineto{\pgfqpoint{2.066917in}{1.905176in}}%
\pgfpathlineto{\pgfqpoint{2.067412in}{1.532027in}}%
\pgfpathlineto{\pgfqpoint{2.068415in}{1.532049in}}%
\pgfpathlineto{\pgfqpoint{2.069297in}{1.532125in}}%
\pgfpathlineto{\pgfqpoint{2.069519in}{1.943926in}}%
\pgfpathlineto{\pgfqpoint{2.070556in}{1.920173in}}%
\pgfpathlineto{\pgfqpoint{2.071149in}{1.919716in}}%
\pgfpathlineto{\pgfqpoint{2.070655in}{1.920215in}}%
\pgfpathlineto{\pgfqpoint{2.071347in}{1.919776in}}%
\pgfpathlineto{\pgfqpoint{2.071792in}{1.918589in}}%
\pgfpathlineto{\pgfqpoint{2.072089in}{1.747945in}}%
\pgfpathlineto{\pgfqpoint{2.072320in}{1.532023in}}%
\pgfpathlineto{\pgfqpoint{2.073145in}{1.532055in}}%
\pgfpathlineto{\pgfqpoint{2.074239in}{1.532098in}}%
\pgfpathlineto{\pgfqpoint{2.074248in}{1.532495in}}%
\pgfpathlineto{\pgfqpoint{2.074359in}{1.940545in}}%
\pgfpathlineto{\pgfqpoint{2.075250in}{1.919541in}}%
\pgfpathlineto{\pgfqpoint{2.075398in}{1.920686in}}%
\pgfpathlineto{\pgfqpoint{2.075349in}{1.919136in}}%
\pgfpathlineto{\pgfqpoint{2.076041in}{1.919943in}}%
\pgfpathlineto{\pgfqpoint{2.076733in}{1.918841in}}%
\pgfpathlineto{\pgfqpoint{2.076981in}{1.793999in}}%
\pgfpathlineto{\pgfqpoint{2.077283in}{1.532029in}}%
\pgfpathlineto{\pgfqpoint{2.078071in}{1.532045in}}%
\pgfpathlineto{\pgfqpoint{2.079192in}{1.532273in}}%
\pgfpathlineto{\pgfqpoint{2.079306in}{1.940954in}}%
\pgfpathlineto{\pgfqpoint{2.080208in}{1.920798in}}%
\pgfpathlineto{\pgfqpoint{2.080505in}{1.918728in}}%
\pgfpathlineto{\pgfqpoint{2.080456in}{1.920904in}}%
\pgfpathlineto{\pgfqpoint{2.080950in}{1.919092in}}%
\pgfpathlineto{\pgfqpoint{2.081000in}{1.920422in}}%
\pgfpathlineto{\pgfqpoint{2.081692in}{1.918514in}}%
\pgfpathlineto{\pgfqpoint{2.082085in}{1.675876in}}%
\pgfpathlineto{\pgfqpoint{2.082239in}{1.532028in}}%
\pgfpathlineto{\pgfqpoint{2.083086in}{1.532051in}}%
\pgfpathlineto{\pgfqpoint{2.084132in}{1.532109in}}%
\pgfpathlineto{\pgfqpoint{2.084174in}{1.745381in}}%
\pgfpathlineto{\pgfqpoint{2.084250in}{1.940502in}}%
\pgfpathlineto{\pgfqpoint{2.084933in}{1.917007in}}%
\pgfpathlineto{\pgfqpoint{2.085230in}{1.922170in}}%
\pgfpathlineto{\pgfqpoint{2.085675in}{1.921324in}}%
\pgfpathlineto{\pgfqpoint{2.085724in}{1.918311in}}%
\pgfpathlineto{\pgfqpoint{2.086466in}{1.919711in}}%
\pgfpathlineto{\pgfqpoint{2.086713in}{1.899940in}}%
\pgfpathlineto{\pgfqpoint{2.087211in}{1.532030in}}%
\pgfpathlineto{\pgfqpoint{2.088183in}{1.532053in}}%
\pgfpathlineto{\pgfqpoint{2.089081in}{1.532151in}}%
\pgfpathlineto{\pgfqpoint{2.089189in}{1.939624in}}%
\pgfpathlineto{\pgfqpoint{2.090253in}{1.919490in}}%
\pgfpathlineto{\pgfqpoint{2.090649in}{1.919920in}}%
\pgfpathlineto{\pgfqpoint{2.090995in}{1.919731in}}%
\pgfpathlineto{\pgfqpoint{2.091588in}{1.917803in}}%
\pgfpathlineto{\pgfqpoint{2.092031in}{1.649280in}}%
\pgfpathlineto{\pgfqpoint{2.092104in}{1.532023in}}%
\pgfpathlineto{\pgfqpoint{2.092782in}{1.532037in}}%
\pgfpathlineto{\pgfqpoint{2.094030in}{1.532842in}}%
\pgfpathlineto{\pgfqpoint{2.094133in}{1.938790in}}%
\pgfpathlineto{\pgfqpoint{2.095007in}{1.925195in}}%
\pgfpathlineto{\pgfqpoint{2.095304in}{1.915465in}}%
\pgfpathlineto{\pgfqpoint{2.095749in}{1.916950in}}%
\pgfpathlineto{\pgfqpoint{2.095799in}{1.922443in}}%
\pgfpathlineto{\pgfqpoint{2.096491in}{1.919052in}}%
\pgfpathlineto{\pgfqpoint{2.096689in}{1.851260in}}%
\pgfpathlineto{\pgfqpoint{2.097086in}{1.532027in}}%
\pgfpathlineto{\pgfqpoint{2.097844in}{1.532050in}}%
\pgfpathlineto{\pgfqpoint{2.098969in}{1.532113in}}%
\pgfpathlineto{\pgfqpoint{2.099086in}{1.940368in}}%
\pgfpathlineto{\pgfqpoint{2.100483in}{1.920904in}}%
\pgfpathlineto{\pgfqpoint{2.100878in}{1.918657in}}%
\pgfpathlineto{\pgfqpoint{2.101274in}{1.919750in}}%
\pgfpathlineto{\pgfqpoint{2.101521in}{1.911085in}}%
\pgfpathlineto{\pgfqpoint{2.102050in}{1.532030in}}%
\pgfpathlineto{\pgfqpoint{2.103173in}{1.532055in}}%
\pgfpathlineto{\pgfqpoint{2.103918in}{1.532186in}}%
\pgfpathlineto{\pgfqpoint{2.104037in}{1.940519in}}%
\pgfpathlineto{\pgfqpoint{2.105023in}{1.919250in}}%
\pgfpathlineto{\pgfqpoint{2.105072in}{1.919974in}}%
\pgfpathlineto{\pgfqpoint{2.105122in}{1.919229in}}%
\pgfpathlineto{\pgfqpoint{2.105765in}{1.919523in}}%
\pgfpathlineto{\pgfqpoint{2.105863in}{1.919584in}}%
\pgfpathlineto{\pgfqpoint{2.106407in}{1.918420in}}%
\pgfpathlineto{\pgfqpoint{2.106655in}{1.792067in}}%
\pgfpathlineto{\pgfqpoint{2.106978in}{1.532028in}}%
\pgfpathlineto{\pgfqpoint{2.107734in}{1.532050in}}%
\pgfpathlineto{\pgfqpoint{2.108860in}{1.532118in}}%
\pgfpathlineto{\pgfqpoint{2.109086in}{1.944366in}}%
\pgfpathlineto{\pgfqpoint{2.110142in}{1.919322in}}%
\pgfpathlineto{\pgfqpoint{2.110191in}{1.919773in}}%
\pgfpathlineto{\pgfqpoint{2.110884in}{1.919412in}}%
\pgfpathlineto{\pgfqpoint{2.111378in}{1.917127in}}%
\pgfpathlineto{\pgfqpoint{2.111814in}{1.634182in}}%
\pgfpathlineto{\pgfqpoint{2.111908in}{1.532028in}}%
\pgfpathlineto{\pgfqpoint{2.112553in}{1.532046in}}%
\pgfpathlineto{\pgfqpoint{2.113804in}{1.532104in}}%
\pgfpathlineto{\pgfqpoint{2.113845in}{1.719233in}}%
\pgfpathlineto{\pgfqpoint{2.114014in}{1.939724in}}%
\pgfpathlineto{\pgfqpoint{2.114618in}{1.924524in}}%
\pgfpathlineto{\pgfqpoint{2.114668in}{1.914233in}}%
\pgfpathlineto{\pgfqpoint{2.115360in}{1.918837in}}%
\pgfpathlineto{\pgfqpoint{2.115509in}{1.921552in}}%
\pgfpathlineto{\pgfqpoint{2.115558in}{1.917476in}}%
\pgfpathlineto{\pgfqpoint{2.116102in}{1.918359in}}%
\pgfpathlineto{\pgfqpoint{2.116151in}{1.919270in}}%
\pgfpathlineto{\pgfqpoint{2.116349in}{1.910758in}}%
\pgfpathlineto{\pgfqpoint{2.116871in}{1.532027in}}%
\pgfpathlineto{\pgfqpoint{2.118226in}{1.532064in}}%
\pgfpathlineto{\pgfqpoint{2.118750in}{1.532113in}}%
\pgfpathlineto{\pgfqpoint{2.118870in}{1.940339in}}%
\pgfpathlineto{\pgfqpoint{2.120298in}{1.920273in}}%
\pgfpathlineto{\pgfqpoint{2.120447in}{1.918036in}}%
\pgfpathlineto{\pgfqpoint{2.120397in}{1.920852in}}%
\pgfpathlineto{\pgfqpoint{2.121090in}{1.918877in}}%
\pgfpathlineto{\pgfqpoint{2.121188in}{1.919032in}}%
\pgfpathlineto{\pgfqpoint{2.121287in}{1.914649in}}%
\pgfpathlineto{\pgfqpoint{2.121831in}{1.532030in}}%
\pgfpathlineto{\pgfqpoint{2.123199in}{1.532061in}}%
\pgfpathlineto{\pgfqpoint{2.123700in}{1.532202in}}%
\pgfpathlineto{\pgfqpoint{2.123820in}{1.939396in}}%
\pgfpathlineto{\pgfqpoint{2.124770in}{1.920065in}}%
\pgfpathlineto{\pgfqpoint{2.124820in}{1.918778in}}%
\pgfpathlineto{\pgfqpoint{2.125562in}{1.919553in}}%
\pgfpathlineto{\pgfqpoint{2.126254in}{1.908376in}}%
\pgfpathlineto{\pgfqpoint{2.126739in}{1.532029in}}%
\pgfpathlineto{\pgfqpoint{2.127824in}{1.532052in}}%
\pgfpathlineto{\pgfqpoint{2.128647in}{1.532273in}}%
\pgfpathlineto{\pgfqpoint{2.128761in}{1.940499in}}%
\pgfpathlineto{\pgfqpoint{2.129689in}{1.919832in}}%
\pgfpathlineto{\pgfqpoint{2.129837in}{1.918152in}}%
\pgfpathlineto{\pgfqpoint{2.129887in}{1.920531in}}%
\pgfpathlineto{\pgfqpoint{2.130480in}{1.919155in}}%
\pgfpathlineto{\pgfqpoint{2.130678in}{1.919660in}}%
\pgfpathlineto{\pgfqpoint{2.130727in}{1.918679in}}%
\pgfpathlineto{\pgfqpoint{2.131024in}{1.918915in}}%
\pgfpathlineto{\pgfqpoint{2.131192in}{1.911762in}}%
\pgfpathlineto{\pgfqpoint{2.131720in}{1.532030in}}%
\pgfpathlineto{\pgfqpoint{2.133037in}{1.532059in}}%
\pgfpathlineto{\pgfqpoint{2.133589in}{1.532143in}}%
\pgfpathlineto{\pgfqpoint{2.133700in}{1.939603in}}%
\pgfpathlineto{\pgfqpoint{2.134744in}{1.919238in}}%
\pgfpathlineto{\pgfqpoint{2.134892in}{1.919519in}}%
\pgfpathlineto{\pgfqpoint{2.134843in}{1.919167in}}%
\pgfpathlineto{\pgfqpoint{2.135485in}{1.919183in}}%
\pgfpathlineto{\pgfqpoint{2.135782in}{1.919071in}}%
\pgfpathlineto{\pgfqpoint{2.135980in}{1.918664in}}%
\pgfpathlineto{\pgfqpoint{2.136079in}{1.918084in}}%
\pgfpathlineto{\pgfqpoint{2.136277in}{1.833204in}}%
\pgfpathlineto{\pgfqpoint{2.136653in}{1.532028in}}%
\pgfpathlineto{\pgfqpoint{2.137411in}{1.532050in}}%
\pgfpathlineto{\pgfqpoint{2.138533in}{1.532121in}}%
\pgfpathlineto{\pgfqpoint{2.138756in}{1.943433in}}%
\pgfpathlineto{\pgfqpoint{2.139830in}{1.919091in}}%
\pgfpathlineto{\pgfqpoint{2.139879in}{1.919520in}}%
\pgfpathlineto{\pgfqpoint{2.139928in}{1.919059in}}%
\pgfpathlineto{\pgfqpoint{2.140571in}{1.919084in}}%
\pgfpathlineto{\pgfqpoint{2.141016in}{1.918177in}}%
\pgfpathlineto{\pgfqpoint{2.141214in}{1.840469in}}%
\pgfpathlineto{\pgfqpoint{2.141579in}{1.532028in}}%
\pgfpathlineto{\pgfqpoint{2.142372in}{1.532045in}}%
\pgfpathlineto{\pgfqpoint{2.143483in}{1.532285in}}%
\pgfpathlineto{\pgfqpoint{2.143586in}{1.939292in}}%
\pgfpathlineto{\pgfqpoint{2.144515in}{1.921374in}}%
\pgfpathlineto{\pgfqpoint{2.144564in}{1.916549in}}%
\pgfpathlineto{\pgfqpoint{2.144614in}{1.922026in}}%
\pgfpathlineto{\pgfqpoint{2.145306in}{1.919416in}}%
\pgfpathlineto{\pgfqpoint{2.145504in}{1.920153in}}%
\pgfpathlineto{\pgfqpoint{2.146048in}{1.905078in}}%
\pgfpathlineto{\pgfqpoint{2.146534in}{1.532028in}}%
\pgfpathlineto{\pgfqpoint{2.147579in}{1.532056in}}%
\pgfpathlineto{\pgfqpoint{2.148423in}{1.532112in}}%
\pgfpathlineto{\pgfqpoint{2.148652in}{1.944419in}}%
\pgfpathlineto{\pgfqpoint{2.149773in}{1.919180in}}%
\pgfpathlineto{\pgfqpoint{2.150367in}{1.919127in}}%
\pgfpathlineto{\pgfqpoint{2.149872in}{1.919112in}}%
\pgfpathlineto{\pgfqpoint{2.150515in}{1.918918in}}%
\pgfpathlineto{\pgfqpoint{2.150812in}{1.918570in}}%
\pgfpathlineto{\pgfqpoint{2.150861in}{1.918324in}}%
\pgfpathlineto{\pgfqpoint{2.150911in}{1.918074in}}%
\pgfpathlineto{\pgfqpoint{2.151109in}{1.837699in}}%
\pgfpathlineto{\pgfqpoint{2.151474in}{1.532028in}}%
\pgfpathlineto{\pgfqpoint{2.152270in}{1.532045in}}%
\pgfpathlineto{\pgfqpoint{2.153368in}{1.532108in}}%
\pgfpathlineto{\pgfqpoint{2.153410in}{1.736808in}}%
\pgfpathlineto{\pgfqpoint{2.153470in}{1.938229in}}%
\pgfpathlineto{\pgfqpoint{2.154155in}{1.915415in}}%
\pgfpathlineto{\pgfqpoint{2.154452in}{1.924500in}}%
\pgfpathlineto{\pgfqpoint{2.154402in}{1.913725in}}%
\pgfpathlineto{\pgfqpoint{2.154897in}{1.922498in}}%
\pgfpathlineto{\pgfqpoint{2.154946in}{1.916330in}}%
\pgfpathlineto{\pgfqpoint{2.155688in}{1.919938in}}%
\pgfpathlineto{\pgfqpoint{2.155985in}{1.885347in}}%
\pgfpathlineto{\pgfqpoint{2.156450in}{1.532030in}}%
\pgfpathlineto{\pgfqpoint{2.157323in}{1.532050in}}%
\pgfpathlineto{\pgfqpoint{2.158318in}{1.532193in}}%
\pgfpathlineto{\pgfqpoint{2.158438in}{1.940237in}}%
\pgfpathlineto{\pgfqpoint{2.159414in}{1.918869in}}%
\pgfpathlineto{\pgfqpoint{2.159711in}{1.919521in}}%
\pgfpathlineto{\pgfqpoint{2.159661in}{1.918729in}}%
\pgfpathlineto{\pgfqpoint{2.160156in}{1.919285in}}%
\pgfpathlineto{\pgfqpoint{2.160848in}{1.914139in}}%
\pgfpathlineto{\pgfqpoint{2.161280in}{1.536155in}}%
\pgfpathlineto{\pgfqpoint{2.161361in}{1.532028in}}%
\pgfpathlineto{\pgfqpoint{2.162006in}{1.532046in}}%
\pgfpathlineto{\pgfqpoint{2.163258in}{1.532103in}}%
\pgfpathlineto{\pgfqpoint{2.163299in}{1.714363in}}%
\pgfpathlineto{\pgfqpoint{2.163383in}{1.940143in}}%
\pgfpathlineto{\pgfqpoint{2.164045in}{1.915591in}}%
\pgfpathlineto{\pgfqpoint{2.164094in}{1.921589in}}%
\pgfpathlineto{\pgfqpoint{2.164787in}{1.920399in}}%
\pgfpathlineto{\pgfqpoint{2.164836in}{1.917609in}}%
\pgfpathlineto{\pgfqpoint{2.164886in}{1.920520in}}%
\pgfpathlineto{\pgfqpoint{2.165578in}{1.918835in}}%
\pgfpathlineto{\pgfqpoint{2.165677in}{1.918902in}}%
\pgfpathlineto{\pgfqpoint{2.165825in}{1.908683in}}%
\pgfpathlineto{\pgfqpoint{2.166298in}{1.532029in}}%
\pgfpathlineto{\pgfqpoint{2.167429in}{1.532056in}}%
\pgfpathlineto{\pgfqpoint{2.168202in}{1.532096in}}%
\pgfpathlineto{\pgfqpoint{2.168211in}{1.532344in}}%
\pgfpathlineto{\pgfqpoint{2.168329in}{1.940109in}}%
\pgfpathlineto{\pgfqpoint{2.169208in}{1.915781in}}%
\pgfpathlineto{\pgfqpoint{2.169357in}{1.923495in}}%
\pgfpathlineto{\pgfqpoint{2.169307in}{1.914359in}}%
\pgfpathlineto{\pgfqpoint{2.170000in}{1.918839in}}%
\pgfpathlineto{\pgfqpoint{2.170148in}{1.920676in}}%
\pgfpathlineto{\pgfqpoint{2.170198in}{1.917053in}}%
\pgfpathlineto{\pgfqpoint{2.170692in}{1.918852in}}%
\pgfpathlineto{\pgfqpoint{2.170890in}{1.838554in}}%
\pgfpathlineto{\pgfqpoint{2.171263in}{1.532028in}}%
\pgfpathlineto{\pgfqpoint{2.172063in}{1.532045in}}%
\pgfpathlineto{\pgfqpoint{2.173151in}{1.532115in}}%
\pgfpathlineto{\pgfqpoint{2.173378in}{1.943897in}}%
\pgfpathlineto{\pgfqpoint{2.174473in}{1.919155in}}%
\pgfpathlineto{\pgfqpoint{2.175313in}{1.918645in}}%
\pgfpathlineto{\pgfqpoint{2.175363in}{1.918709in}}%
\pgfpathlineto{\pgfqpoint{2.175610in}{1.918068in}}%
\pgfpathlineto{\pgfqpoint{2.175709in}{1.909198in}}%
\pgfpathlineto{\pgfqpoint{2.176193in}{1.532029in}}%
\pgfpathlineto{\pgfqpoint{2.177376in}{1.532054in}}%
\pgfpathlineto{\pgfqpoint{2.178101in}{1.532272in}}%
\pgfpathlineto{\pgfqpoint{2.178202in}{1.938782in}}%
\pgfpathlineto{\pgfqpoint{2.179134in}{1.919256in}}%
\pgfpathlineto{\pgfqpoint{2.179282in}{1.918305in}}%
\pgfpathlineto{\pgfqpoint{2.179233in}{1.919575in}}%
\pgfpathlineto{\pgfqpoint{2.179925in}{1.918850in}}%
\pgfpathlineto{\pgfqpoint{2.180519in}{1.918063in}}%
\pgfpathlineto{\pgfqpoint{2.180024in}{1.919008in}}%
\pgfpathlineto{\pgfqpoint{2.180568in}{1.918127in}}%
\pgfpathlineto{\pgfqpoint{2.180667in}{1.905519in}}%
\pgfpathlineto{\pgfqpoint{2.181176in}{1.532030in}}%
\pgfpathlineto{\pgfqpoint{2.182197in}{1.532053in}}%
\pgfpathlineto{\pgfqpoint{2.183046in}{1.532210in}}%
\pgfpathlineto{\pgfqpoint{2.183167in}{1.939235in}}%
\pgfpathlineto{\pgfqpoint{2.184106in}{1.918759in}}%
\pgfpathlineto{\pgfqpoint{2.184254in}{1.919552in}}%
\pgfpathlineto{\pgfqpoint{2.184304in}{1.918287in}}%
\pgfpathlineto{\pgfqpoint{2.184897in}{1.918981in}}%
\pgfpathlineto{\pgfqpoint{2.185540in}{1.917532in}}%
\pgfpathlineto{\pgfqpoint{2.185738in}{1.831985in}}%
\pgfpathlineto{\pgfqpoint{2.186121in}{1.532030in}}%
\pgfpathlineto{\pgfqpoint{2.186894in}{1.532048in}}%
\pgfpathlineto{\pgfqpoint{2.187990in}{1.532152in}}%
\pgfpathlineto{\pgfqpoint{2.188099in}{1.938733in}}%
\pgfpathlineto{\pgfqpoint{2.189160in}{1.918655in}}%
\pgfpathlineto{\pgfqpoint{2.189209in}{1.919104in}}%
\pgfpathlineto{\pgfqpoint{2.189902in}{1.918798in}}%
\pgfpathlineto{\pgfqpoint{2.190495in}{1.917435in}}%
\pgfpathlineto{\pgfqpoint{2.190792in}{1.747985in}}%
\pgfpathlineto{\pgfqpoint{2.191066in}{1.532030in}}%
\pgfpathlineto{\pgfqpoint{2.191839in}{1.532048in}}%
\pgfpathlineto{\pgfqpoint{2.192935in}{1.532143in}}%
\pgfpathlineto{\pgfqpoint{2.193046in}{1.939136in}}%
\pgfpathlineto{\pgfqpoint{2.194117in}{1.918890in}}%
\pgfpathlineto{\pgfqpoint{2.195402in}{1.917856in}}%
\pgfpathlineto{\pgfqpoint{2.195501in}{1.906618in}}%
\pgfpathlineto{\pgfqpoint{2.196015in}{1.532030in}}%
\pgfpathlineto{\pgfqpoint{2.197038in}{1.532053in}}%
\pgfpathlineto{\pgfqpoint{2.197881in}{1.532163in}}%
\pgfpathlineto{\pgfqpoint{2.197999in}{1.939273in}}%
\pgfpathlineto{\pgfqpoint{2.199001in}{1.918705in}}%
\pgfpathlineto{\pgfqpoint{2.199743in}{1.918824in}}%
\pgfpathlineto{\pgfqpoint{2.199842in}{1.918728in}}%
\pgfpathlineto{\pgfqpoint{2.200386in}{1.917417in}}%
\pgfpathlineto{\pgfqpoint{2.200633in}{1.787562in}}%
\pgfpathlineto{\pgfqpoint{2.200920in}{1.532029in}}%
\pgfpathlineto{\pgfqpoint{2.201707in}{1.532045in}}%
\pgfpathlineto{\pgfqpoint{2.202829in}{1.532272in}}%
\pgfpathlineto{\pgfqpoint{2.202943in}{1.939885in}}%
\pgfpathlineto{\pgfqpoint{2.203878in}{1.918746in}}%
\pgfpathlineto{\pgfqpoint{2.204026in}{1.917553in}}%
\pgfpathlineto{\pgfqpoint{2.204076in}{1.920068in}}%
\pgfpathlineto{\pgfqpoint{2.204570in}{1.917989in}}%
\pgfpathlineto{\pgfqpoint{2.204620in}{1.919306in}}%
\pgfpathlineto{\pgfqpoint{2.205262in}{1.917768in}}%
\pgfpathlineto{\pgfqpoint{2.205312in}{1.917953in}}%
\pgfpathlineto{\pgfqpoint{2.205510in}{1.840734in}}%
\pgfpathlineto{\pgfqpoint{2.205874in}{1.532028in}}%
\pgfpathlineto{\pgfqpoint{2.206669in}{1.532045in}}%
\pgfpathlineto{\pgfqpoint{2.207768in}{1.532107in}}%
\pgfpathlineto{\pgfqpoint{2.207810in}{1.736726in}}%
\pgfpathlineto{\pgfqpoint{2.207888in}{1.939473in}}%
\pgfpathlineto{\pgfqpoint{2.208571in}{1.914927in}}%
\pgfpathlineto{\pgfqpoint{2.208620in}{1.921863in}}%
\pgfpathlineto{\pgfqpoint{2.209313in}{1.919788in}}%
\pgfpathlineto{\pgfqpoint{2.209461in}{1.917305in}}%
\pgfpathlineto{\pgfqpoint{2.209412in}{1.920213in}}%
\pgfpathlineto{\pgfqpoint{2.210104in}{1.918285in}}%
\pgfpathlineto{\pgfqpoint{2.210203in}{1.918430in}}%
\pgfpathlineto{\pgfqpoint{2.210302in}{1.916230in}}%
\pgfpathlineto{\pgfqpoint{2.210723in}{1.634404in}}%
\pgfpathlineto{\pgfqpoint{2.210795in}{1.532023in}}%
\pgfpathlineto{\pgfqpoint{2.211469in}{1.532037in}}%
\pgfpathlineto{\pgfqpoint{2.212713in}{1.532103in}}%
\pgfpathlineto{\pgfqpoint{2.212752in}{1.701239in}}%
\pgfpathlineto{\pgfqpoint{2.212921in}{1.940713in}}%
\pgfpathlineto{\pgfqpoint{2.213525in}{1.919608in}}%
\pgfpathlineto{\pgfqpoint{2.213673in}{1.916720in}}%
\pgfpathlineto{\pgfqpoint{2.213624in}{1.920489in}}%
\pgfpathlineto{\pgfqpoint{2.214267in}{1.919158in}}%
\pgfpathlineto{\pgfqpoint{2.214514in}{1.919399in}}%
\pgfpathlineto{\pgfqpoint{2.215108in}{1.917784in}}%
\pgfpathlineto{\pgfqpoint{2.215157in}{1.917994in}}%
\pgfpathlineto{\pgfqpoint{2.215207in}{1.917598in}}%
\pgfpathlineto{\pgfqpoint{2.215404in}{1.837717in}}%
\pgfpathlineto{\pgfqpoint{2.215776in}{1.532027in}}%
\pgfpathlineto{\pgfqpoint{2.216579in}{1.532045in}}%
\pgfpathlineto{\pgfqpoint{2.217661in}{1.532125in}}%
\pgfpathlineto{\pgfqpoint{2.217880in}{1.941841in}}%
\pgfpathlineto{\pgfqpoint{2.218922in}{1.918722in}}%
\pgfpathlineto{\pgfqpoint{2.219714in}{1.918410in}}%
\pgfpathlineto{\pgfqpoint{2.219120in}{1.918934in}}%
\pgfpathlineto{\pgfqpoint{2.219763in}{1.918432in}}%
\pgfpathlineto{\pgfqpoint{2.220109in}{1.917864in}}%
\pgfpathlineto{\pgfqpoint{2.220208in}{1.911366in}}%
\pgfpathlineto{\pgfqpoint{2.220739in}{1.532030in}}%
\pgfpathlineto{\pgfqpoint{2.222056in}{1.532059in}}%
\pgfpathlineto{\pgfqpoint{2.222607in}{1.532142in}}%
\pgfpathlineto{\pgfqpoint{2.222721in}{1.939229in}}%
\pgfpathlineto{\pgfqpoint{2.223766in}{1.918648in}}%
\pgfpathlineto{\pgfqpoint{2.223865in}{1.918775in}}%
\pgfpathlineto{\pgfqpoint{2.224706in}{1.918357in}}%
\pgfpathlineto{\pgfqpoint{2.224755in}{1.918417in}}%
\pgfpathlineto{\pgfqpoint{2.225052in}{1.917796in}}%
\pgfpathlineto{\pgfqpoint{2.225151in}{1.912158in}}%
\pgfpathlineto{\pgfqpoint{2.225670in}{1.532024in}}%
\pgfpathlineto{\pgfqpoint{2.227018in}{1.532066in}}%
\pgfpathlineto{\pgfqpoint{2.227551in}{1.532120in}}%
\pgfpathlineto{\pgfqpoint{2.227775in}{1.942640in}}%
\pgfpathlineto{\pgfqpoint{2.228853in}{1.918731in}}%
\pgfpathlineto{\pgfqpoint{2.229693in}{1.918252in}}%
\pgfpathlineto{\pgfqpoint{2.229743in}{1.918314in}}%
\pgfpathlineto{\pgfqpoint{2.230039in}{1.917437in}}%
\pgfpathlineto{\pgfqpoint{2.230188in}{1.873471in}}%
\pgfpathlineto{\pgfqpoint{2.230616in}{1.532024in}}%
\pgfpathlineto{\pgfqpoint{2.231470in}{1.532055in}}%
\pgfpathlineto{\pgfqpoint{2.232497in}{1.532121in}}%
\pgfpathlineto{\pgfqpoint{2.232716in}{1.940619in}}%
\pgfpathlineto{\pgfqpoint{2.233780in}{1.918557in}}%
\pgfpathlineto{\pgfqpoint{2.234473in}{1.918511in}}%
\pgfpathlineto{\pgfqpoint{2.233879in}{1.918441in}}%
\pgfpathlineto{\pgfqpoint{2.234572in}{1.918395in}}%
\pgfpathlineto{\pgfqpoint{2.235017in}{1.916807in}}%
\pgfpathlineto{\pgfqpoint{2.235363in}{1.704014in}}%
\pgfpathlineto{\pgfqpoint{2.235562in}{1.532024in}}%
\pgfpathlineto{\pgfqpoint{2.236368in}{1.532042in}}%
\pgfpathlineto{\pgfqpoint{2.237442in}{1.532121in}}%
\pgfpathlineto{\pgfqpoint{2.237662in}{1.941559in}}%
\pgfpathlineto{\pgfqpoint{2.238724in}{1.918642in}}%
\pgfpathlineto{\pgfqpoint{2.239664in}{1.918154in}}%
\pgfpathlineto{\pgfqpoint{2.239961in}{1.916894in}}%
\pgfpathlineto{\pgfqpoint{2.240307in}{1.705186in}}%
\pgfpathlineto{\pgfqpoint{2.240507in}{1.532024in}}%
\pgfpathlineto{\pgfqpoint{2.241312in}{1.532042in}}%
\pgfpathlineto{\pgfqpoint{2.242388in}{1.532121in}}%
\pgfpathlineto{\pgfqpoint{2.242606in}{1.941052in}}%
\pgfpathlineto{\pgfqpoint{2.243669in}{1.918459in}}%
\pgfpathlineto{\pgfqpoint{2.243916in}{1.918649in}}%
\pgfpathlineto{\pgfqpoint{2.243768in}{1.918348in}}%
\pgfpathlineto{\pgfqpoint{2.244411in}{1.918243in}}%
\pgfpathlineto{\pgfqpoint{2.244460in}{1.918318in}}%
\pgfpathlineto{\pgfqpoint{2.244757in}{1.917823in}}%
\pgfpathlineto{\pgfqpoint{2.244807in}{1.917811in}}%
\pgfpathlineto{\pgfqpoint{2.244906in}{1.916864in}}%
\pgfpathlineto{\pgfqpoint{2.245252in}{1.705880in}}%
\pgfpathlineto{\pgfqpoint{2.245453in}{1.532024in}}%
\pgfpathlineto{\pgfqpoint{2.246257in}{1.532042in}}%
\pgfpathlineto{\pgfqpoint{2.247333in}{1.532121in}}%
\pgfpathlineto{\pgfqpoint{2.247552in}{1.941294in}}%
\pgfpathlineto{\pgfqpoint{2.248599in}{1.918396in}}%
\pgfpathlineto{\pgfqpoint{2.248747in}{1.918631in}}%
\pgfpathlineto{\pgfqpoint{2.248698in}{1.918308in}}%
\pgfpathlineto{\pgfqpoint{2.249341in}{1.918234in}}%
\pgfpathlineto{\pgfqpoint{2.249539in}{1.918169in}}%
\pgfpathlineto{\pgfqpoint{2.249687in}{1.917830in}}%
\pgfpathlineto{\pgfqpoint{2.249736in}{1.917817in}}%
\pgfpathlineto{\pgfqpoint{2.249835in}{1.917212in}}%
\pgfpathlineto{\pgfqpoint{2.250083in}{1.793525in}}%
\pgfpathlineto{\pgfqpoint{2.250393in}{1.532024in}}%
\pgfpathlineto{\pgfqpoint{2.251144in}{1.532053in}}%
\pgfpathlineto{\pgfqpoint{2.252278in}{1.532120in}}%
\pgfpathlineto{\pgfqpoint{2.252497in}{1.941185in}}%
\pgfpathlineto{\pgfqpoint{2.253596in}{1.918610in}}%
\pgfpathlineto{\pgfqpoint{2.254436in}{1.918071in}}%
\pgfpathlineto{\pgfqpoint{2.254486in}{1.918157in}}%
\pgfpathlineto{\pgfqpoint{2.254783in}{1.917134in}}%
\pgfpathlineto{\pgfqpoint{2.255030in}{1.792547in}}%
\pgfpathlineto{\pgfqpoint{2.255338in}{1.532024in}}%
\pgfpathlineto{\pgfqpoint{2.256089in}{1.532053in}}%
\pgfpathlineto{\pgfqpoint{2.257224in}{1.532120in}}%
\pgfpathlineto{\pgfqpoint{2.257442in}{1.941193in}}%
\pgfpathlineto{\pgfqpoint{2.258541in}{1.918576in}}%
\pgfpathlineto{\pgfqpoint{2.259381in}{1.918038in}}%
\pgfpathlineto{\pgfqpoint{2.259431in}{1.918124in}}%
\pgfpathlineto{\pgfqpoint{2.259728in}{1.917111in}}%
\pgfpathlineto{\pgfqpoint{2.259975in}{1.793047in}}%
\pgfpathlineto{\pgfqpoint{2.260283in}{1.532024in}}%
\pgfpathlineto{\pgfqpoint{2.261034in}{1.532053in}}%
\pgfpathlineto{\pgfqpoint{2.262169in}{1.532120in}}%
\pgfpathlineto{\pgfqpoint{2.262388in}{1.941125in}}%
\pgfpathlineto{\pgfqpoint{2.263487in}{1.918543in}}%
\pgfpathlineto{\pgfqpoint{2.264328in}{1.918006in}}%
\pgfpathlineto{\pgfqpoint{2.264377in}{1.918091in}}%
\pgfpathlineto{\pgfqpoint{2.264674in}{1.917071in}}%
\pgfpathlineto{\pgfqpoint{2.264921in}{1.792617in}}%
\pgfpathlineto{\pgfqpoint{2.265229in}{1.532024in}}%
\pgfpathlineto{\pgfqpoint{2.265979in}{1.532053in}}%
\pgfpathlineto{\pgfqpoint{2.267115in}{1.532120in}}%
\pgfpathlineto{\pgfqpoint{2.267333in}{1.941098in}}%
\pgfpathlineto{\pgfqpoint{2.268433in}{1.918505in}}%
\pgfpathlineto{\pgfqpoint{2.269274in}{1.917978in}}%
\pgfpathlineto{\pgfqpoint{2.269323in}{1.918057in}}%
\pgfpathlineto{\pgfqpoint{2.269570in}{1.917444in}}%
\pgfpathlineto{\pgfqpoint{2.269669in}{1.910605in}}%
\pgfpathlineto{\pgfqpoint{2.270174in}{1.532024in}}%
\pgfpathlineto{\pgfqpoint{2.271418in}{1.532064in}}%
\pgfpathlineto{\pgfqpoint{2.272060in}{1.532120in}}%
\pgfpathlineto{\pgfqpoint{2.272279in}{1.941050in}}%
\pgfpathlineto{\pgfqpoint{2.273330in}{1.918164in}}%
\pgfpathlineto{\pgfqpoint{2.273478in}{1.918481in}}%
\pgfpathlineto{\pgfqpoint{2.273429in}{1.918111in}}%
\pgfpathlineto{\pgfqpoint{2.274072in}{1.918098in}}%
\pgfpathlineto{\pgfqpoint{2.274517in}{1.917406in}}%
\pgfpathlineto{\pgfqpoint{2.274616in}{1.910402in}}%
\pgfpathlineto{\pgfqpoint{2.275119in}{1.532024in}}%
\pgfpathlineto{\pgfqpoint{2.276364in}{1.532064in}}%
\pgfpathlineto{\pgfqpoint{2.277006in}{1.532120in}}%
\pgfpathlineto{\pgfqpoint{2.277224in}{1.941018in}}%
\pgfpathlineto{\pgfqpoint{2.278277in}{1.918142in}}%
\pgfpathlineto{\pgfqpoint{2.278425in}{1.918452in}}%
\pgfpathlineto{\pgfqpoint{2.278376in}{1.918081in}}%
\pgfpathlineto{\pgfqpoint{2.279018in}{1.918061in}}%
\pgfpathlineto{\pgfqpoint{2.279464in}{1.917368in}}%
\pgfpathlineto{\pgfqpoint{2.279562in}{1.910209in}}%
\pgfpathlineto{\pgfqpoint{2.280064in}{1.532024in}}%
\pgfpathlineto{\pgfqpoint{2.281308in}{1.532064in}}%
\pgfpathlineto{\pgfqpoint{2.281951in}{1.532121in}}%
\pgfpathlineto{\pgfqpoint{2.282170in}{1.940904in}}%
\pgfpathlineto{\pgfqpoint{2.283226in}{1.918207in}}%
\pgfpathlineto{\pgfqpoint{2.283819in}{1.918251in}}%
\pgfpathlineto{\pgfqpoint{2.283325in}{1.918097in}}%
\pgfpathlineto{\pgfqpoint{2.283968in}{1.917990in}}%
\pgfpathlineto{\pgfqpoint{2.284017in}{1.918069in}}%
\pgfpathlineto{\pgfqpoint{2.284413in}{1.917296in}}%
\pgfpathlineto{\pgfqpoint{2.284462in}{1.916893in}}%
\pgfpathlineto{\pgfqpoint{2.284709in}{1.788568in}}%
\pgfpathlineto{\pgfqpoint{2.285010in}{1.532024in}}%
\pgfpathlineto{\pgfqpoint{2.285760in}{1.532053in}}%
\pgfpathlineto{\pgfqpoint{2.286897in}{1.532120in}}%
\pgfpathlineto{\pgfqpoint{2.287115in}{1.941012in}}%
\pgfpathlineto{\pgfqpoint{2.288166in}{1.918070in}}%
\pgfpathlineto{\pgfqpoint{2.288315in}{1.918388in}}%
\pgfpathlineto{\pgfqpoint{2.288265in}{1.918019in}}%
\pgfpathlineto{\pgfqpoint{2.288908in}{1.918006in}}%
\pgfpathlineto{\pgfqpoint{2.289353in}{1.917318in}}%
\pgfpathlineto{\pgfqpoint{2.289452in}{1.910537in}}%
\pgfpathlineto{\pgfqpoint{2.289955in}{1.532024in}}%
\pgfpathlineto{\pgfqpoint{2.291198in}{1.532064in}}%
\pgfpathlineto{\pgfqpoint{2.291842in}{1.532122in}}%
\pgfpathlineto{\pgfqpoint{2.292061in}{1.940786in}}%
\pgfpathlineto{\pgfqpoint{2.293130in}{1.918356in}}%
\pgfpathlineto{\pgfqpoint{2.293971in}{1.917831in}}%
\pgfpathlineto{\pgfqpoint{2.294020in}{1.917932in}}%
\pgfpathlineto{\pgfqpoint{2.294317in}{1.917145in}}%
\pgfpathlineto{\pgfqpoint{2.294416in}{1.906049in}}%
\pgfpathlineto{\pgfqpoint{2.294921in}{1.532030in}}%
\pgfpathlineto{\pgfqpoint{2.295990in}{1.532055in}}%
\pgfpathlineto{\pgfqpoint{2.296789in}{1.532142in}}%
\pgfpathlineto{\pgfqpoint{2.296901in}{1.938671in}}%
\pgfpathlineto{\pgfqpoint{2.297980in}{1.918345in}}%
\pgfpathlineto{\pgfqpoint{2.298722in}{1.917979in}}%
\pgfpathlineto{\pgfqpoint{2.298771in}{1.918070in}}%
\pgfpathlineto{\pgfqpoint{2.299315in}{1.916186in}}%
\pgfpathlineto{\pgfqpoint{2.299740in}{1.652365in}}%
\pgfpathlineto{\pgfqpoint{2.299846in}{1.532028in}}%
\pgfpathlineto{\pgfqpoint{2.300496in}{1.532047in}}%
\pgfpathlineto{\pgfqpoint{2.301732in}{1.532115in}}%
\pgfpathlineto{\pgfqpoint{2.301959in}{1.942798in}}%
\pgfpathlineto{\pgfqpoint{2.303061in}{1.918243in}}%
\pgfpathlineto{\pgfqpoint{2.303160in}{1.918287in}}%
\pgfpathlineto{\pgfqpoint{2.303902in}{1.917774in}}%
\pgfpathlineto{\pgfqpoint{2.303951in}{1.917815in}}%
\pgfpathlineto{\pgfqpoint{2.304199in}{1.917170in}}%
\pgfpathlineto{\pgfqpoint{2.304298in}{1.908539in}}%
\pgfpathlineto{\pgfqpoint{2.304770in}{1.532029in}}%
\pgfpathlineto{\pgfqpoint{2.305948in}{1.532054in}}%
\pgfpathlineto{\pgfqpoint{2.306684in}{1.532419in}}%
\pgfpathlineto{\pgfqpoint{2.306784in}{1.937835in}}%
\pgfpathlineto{\pgfqpoint{2.307672in}{1.918556in}}%
\pgfpathlineto{\pgfqpoint{2.307969in}{1.917282in}}%
\pgfpathlineto{\pgfqpoint{2.307919in}{1.918910in}}%
\pgfpathlineto{\pgfqpoint{2.308414in}{1.917554in}}%
\pgfpathlineto{\pgfqpoint{2.308463in}{1.918460in}}%
\pgfpathlineto{\pgfqpoint{2.309106in}{1.917250in}}%
\pgfpathlineto{\pgfqpoint{2.309156in}{1.917285in}}%
\pgfpathlineto{\pgfqpoint{2.309255in}{1.906157in}}%
\pgfpathlineto{\pgfqpoint{2.309740in}{1.532026in}}%
\pgfpathlineto{\pgfqpoint{2.310790in}{1.532058in}}%
\pgfpathlineto{\pgfqpoint{2.311623in}{1.532112in}}%
\pgfpathlineto{\pgfqpoint{2.311663in}{1.722561in}}%
\pgfpathlineto{\pgfqpoint{2.311837in}{1.939497in}}%
\pgfpathlineto{\pgfqpoint{2.312435in}{1.917897in}}%
\pgfpathlineto{\pgfqpoint{2.312682in}{1.916622in}}%
\pgfpathlineto{\pgfqpoint{2.312633in}{1.919493in}}%
\pgfpathlineto{\pgfqpoint{2.313127in}{1.917209in}}%
\pgfpathlineto{\pgfqpoint{2.313177in}{1.918758in}}%
\pgfpathlineto{\pgfqpoint{2.313869in}{1.917996in}}%
\pgfpathlineto{\pgfqpoint{2.314166in}{1.913466in}}%
\pgfpathlineto{\pgfqpoint{2.314580in}{1.600122in}}%
\pgfpathlineto{\pgfqpoint{2.314684in}{1.532028in}}%
\pgfpathlineto{\pgfqpoint{2.315335in}{1.532047in}}%
\pgfpathlineto{\pgfqpoint{2.316569in}{1.532117in}}%
\pgfpathlineto{\pgfqpoint{2.316794in}{1.942359in}}%
\pgfpathlineto{\pgfqpoint{2.317860in}{1.917997in}}%
\pgfpathlineto{\pgfqpoint{2.318108in}{1.918199in}}%
\pgfpathlineto{\pgfqpoint{2.317959in}{1.917875in}}%
\pgfpathlineto{\pgfqpoint{2.318553in}{1.917983in}}%
\pgfpathlineto{\pgfqpoint{2.319097in}{1.915988in}}%
\pgfpathlineto{\pgfqpoint{2.319514in}{1.659414in}}%
\pgfpathlineto{\pgfqpoint{2.319518in}{1.663562in}}%
\pgfpathlineto{\pgfqpoint{2.319620in}{1.532028in}}%
\pgfpathlineto{\pgfqpoint{2.320316in}{1.532043in}}%
\pgfpathlineto{\pgfqpoint{2.321513in}{1.532108in}}%
\pgfpathlineto{\pgfqpoint{2.321555in}{1.740316in}}%
\pgfpathlineto{\pgfqpoint{2.321615in}{1.936832in}}%
\pgfpathlineto{\pgfqpoint{2.322316in}{1.914587in}}%
\pgfpathlineto{\pgfqpoint{2.322464in}{1.922967in}}%
\pgfpathlineto{\pgfqpoint{2.322415in}{1.912796in}}%
\pgfpathlineto{\pgfqpoint{2.323107in}{1.917997in}}%
\pgfpathlineto{\pgfqpoint{2.323255in}{1.919892in}}%
\pgfpathlineto{\pgfqpoint{2.323305in}{1.915922in}}%
\pgfpathlineto{\pgfqpoint{2.323799in}{1.918750in}}%
\pgfpathlineto{\pgfqpoint{2.324145in}{1.882393in}}%
\pgfpathlineto{\pgfqpoint{2.324584in}{1.532027in}}%
\pgfpathlineto{\pgfqpoint{2.325447in}{1.532053in}}%
\pgfpathlineto{\pgfqpoint{2.326459in}{1.532112in}}%
\pgfpathlineto{\pgfqpoint{2.326581in}{1.938958in}}%
\pgfpathlineto{\pgfqpoint{2.328008in}{1.918066in}}%
\pgfpathlineto{\pgfqpoint{2.328255in}{1.916649in}}%
\pgfpathlineto{\pgfqpoint{2.328206in}{1.919348in}}%
\pgfpathlineto{\pgfqpoint{2.328799in}{1.917096in}}%
\pgfpathlineto{\pgfqpoint{2.328849in}{1.917490in}}%
\pgfpathlineto{\pgfqpoint{2.328997in}{1.916291in}}%
\pgfpathlineto{\pgfqpoint{2.329416in}{1.600814in}}%
\pgfpathlineto{\pgfqpoint{2.329528in}{1.532024in}}%
\pgfpathlineto{\pgfqpoint{2.330144in}{1.532038in}}%
\pgfpathlineto{\pgfqpoint{2.331406in}{1.532120in}}%
\pgfpathlineto{\pgfqpoint{2.331629in}{1.941729in}}%
\pgfpathlineto{\pgfqpoint{2.332683in}{1.917846in}}%
\pgfpathlineto{\pgfqpoint{2.332831in}{1.918146in}}%
\pgfpathlineto{\pgfqpoint{2.332781in}{1.917782in}}%
\pgfpathlineto{\pgfqpoint{2.333424in}{1.917752in}}%
\pgfpathlineto{\pgfqpoint{2.333869in}{1.917039in}}%
\pgfpathlineto{\pgfqpoint{2.333968in}{1.908877in}}%
\pgfpathlineto{\pgfqpoint{2.334454in}{1.532028in}}%
\pgfpathlineto{\pgfqpoint{2.335642in}{1.532054in}}%
\pgfpathlineto{\pgfqpoint{2.336349in}{1.532104in}}%
\pgfpathlineto{\pgfqpoint{2.336390in}{1.720681in}}%
\pgfpathlineto{\pgfqpoint{2.336562in}{1.937129in}}%
\pgfpathlineto{\pgfqpoint{2.337159in}{1.915561in}}%
\pgfpathlineto{\pgfqpoint{2.337456in}{1.921303in}}%
\pgfpathlineto{\pgfqpoint{2.337407in}{1.914489in}}%
\pgfpathlineto{\pgfqpoint{2.337901in}{1.920048in}}%
\pgfpathlineto{\pgfqpoint{2.337951in}{1.916160in}}%
\pgfpathlineto{\pgfqpoint{2.338693in}{1.918234in}}%
\pgfpathlineto{\pgfqpoint{2.338950in}{1.898818in}}%
\pgfpathlineto{\pgfqpoint{2.339429in}{1.532024in}}%
\pgfpathlineto{\pgfqpoint{2.340396in}{1.532058in}}%
\pgfpathlineto{\pgfqpoint{2.341297in}{1.532120in}}%
\pgfpathlineto{\pgfqpoint{2.341531in}{1.943884in}}%
\pgfpathlineto{\pgfqpoint{2.342562in}{1.918062in}}%
\pgfpathlineto{\pgfqpoint{2.342809in}{1.918307in}}%
\pgfpathlineto{\pgfqpoint{2.343551in}{1.917415in}}%
\pgfpathlineto{\pgfqpoint{2.343601in}{1.917567in}}%
\pgfpathlineto{\pgfqpoint{2.343798in}{1.916780in}}%
\pgfpathlineto{\pgfqpoint{2.343996in}{1.833223in}}%
\pgfpathlineto{\pgfqpoint{2.344352in}{1.532024in}}%
\pgfpathlineto{\pgfqpoint{2.345149in}{1.532041in}}%
\pgfpathlineto{\pgfqpoint{2.346242in}{1.532124in}}%
\pgfpathlineto{\pgfqpoint{2.346448in}{1.937926in}}%
\pgfpathlineto{\pgfqpoint{2.347533in}{1.917913in}}%
\pgfpathlineto{\pgfqpoint{2.347632in}{1.917781in}}%
\pgfpathlineto{\pgfqpoint{2.348325in}{1.917739in}}%
\pgfpathlineto{\pgfqpoint{2.348770in}{1.916179in}}%
\pgfpathlineto{\pgfqpoint{2.349116in}{1.704594in}}%
\pgfpathlineto{\pgfqpoint{2.349306in}{1.532026in}}%
\pgfpathlineto{\pgfqpoint{2.350112in}{1.532043in}}%
\pgfpathlineto{\pgfqpoint{2.351187in}{1.532113in}}%
\pgfpathlineto{\pgfqpoint{2.351226in}{1.718018in}}%
\pgfpathlineto{\pgfqpoint{2.351406in}{1.940877in}}%
\pgfpathlineto{\pgfqpoint{2.352006in}{1.919632in}}%
\pgfpathlineto{\pgfqpoint{2.352055in}{1.915800in}}%
\pgfpathlineto{\pgfqpoint{2.352748in}{1.917666in}}%
\pgfpathlineto{\pgfqpoint{2.352896in}{1.918605in}}%
\pgfpathlineto{\pgfqpoint{2.352946in}{1.917061in}}%
\pgfpathlineto{\pgfqpoint{2.353440in}{1.917876in}}%
\pgfpathlineto{\pgfqpoint{2.353737in}{1.912226in}}%
\pgfpathlineto{\pgfqpoint{2.354145in}{1.567742in}}%
\pgfpathlineto{\pgfqpoint{2.354255in}{1.532027in}}%
\pgfpathlineto{\pgfqpoint{2.354917in}{1.532048in}}%
\pgfpathlineto{\pgfqpoint{2.356132in}{1.532113in}}%
\pgfpathlineto{\pgfqpoint{2.356234in}{1.936051in}}%
\pgfpathlineto{\pgfqpoint{2.357670in}{1.917491in}}%
\pgfpathlineto{\pgfqpoint{2.357868in}{1.918678in}}%
\pgfpathlineto{\pgfqpoint{2.357917in}{1.916895in}}%
\pgfpathlineto{\pgfqpoint{2.358313in}{1.918100in}}%
\pgfpathlineto{\pgfqpoint{2.358709in}{1.906217in}}%
\pgfpathlineto{\pgfqpoint{2.359199in}{1.532024in}}%
\pgfpathlineto{\pgfqpoint{2.360354in}{1.532062in}}%
\pgfpathlineto{\pgfqpoint{2.361078in}{1.532119in}}%
\pgfpathlineto{\pgfqpoint{2.361300in}{1.942574in}}%
\pgfpathlineto{\pgfqpoint{2.362359in}{1.917697in}}%
\pgfpathlineto{\pgfqpoint{2.362507in}{1.917985in}}%
\pgfpathlineto{\pgfqpoint{2.362457in}{1.917653in}}%
\pgfpathlineto{\pgfqpoint{2.363100in}{1.917622in}}%
\pgfpathlineto{\pgfqpoint{2.363545in}{1.916898in}}%
\pgfpathlineto{\pgfqpoint{2.363644in}{1.908419in}}%
\pgfpathlineto{\pgfqpoint{2.364141in}{1.532027in}}%
\pgfpathlineto{\pgfqpoint{2.365292in}{1.532059in}}%
\pgfpathlineto{\pgfqpoint{2.366023in}{1.532113in}}%
\pgfpathlineto{\pgfqpoint{2.366124in}{1.936489in}}%
\pgfpathlineto{\pgfqpoint{2.367539in}{1.918866in}}%
\pgfpathlineto{\pgfqpoint{2.367836in}{1.916359in}}%
\pgfpathlineto{\pgfqpoint{2.367787in}{1.919125in}}%
\pgfpathlineto{\pgfqpoint{2.368380in}{1.916707in}}%
\pgfpathlineto{\pgfqpoint{2.368430in}{1.917461in}}%
\pgfpathlineto{\pgfqpoint{2.368578in}{1.913707in}}%
\pgfpathlineto{\pgfqpoint{2.369103in}{1.532025in}}%
\pgfpathlineto{\pgfqpoint{2.370468in}{1.532065in}}%
\pgfpathlineto{\pgfqpoint{2.370971in}{1.532136in}}%
\pgfpathlineto{\pgfqpoint{2.371203in}{1.946735in}}%
\pgfpathlineto{\pgfqpoint{2.372089in}{1.918698in}}%
\pgfpathlineto{\pgfqpoint{2.372139in}{1.916606in}}%
\pgfpathlineto{\pgfqpoint{2.372188in}{1.918900in}}%
\pgfpathlineto{\pgfqpoint{2.372881in}{1.917797in}}%
\pgfpathlineto{\pgfqpoint{2.372980in}{1.918004in}}%
\pgfpathlineto{\pgfqpoint{2.373524in}{1.910610in}}%
\pgfpathlineto{\pgfqpoint{2.374020in}{1.532024in}}%
\pgfpathlineto{\pgfqpoint{2.375309in}{1.532053in}}%
\pgfpathlineto{\pgfqpoint{2.375914in}{1.532113in}}%
\pgfpathlineto{\pgfqpoint{2.376016in}{1.936321in}}%
\pgfpathlineto{\pgfqpoint{2.377413in}{1.921295in}}%
\pgfpathlineto{\pgfqpoint{2.377463in}{1.914440in}}%
\pgfpathlineto{\pgfqpoint{2.378205in}{1.918456in}}%
\pgfpathlineto{\pgfqpoint{2.378304in}{1.918510in}}%
\pgfpathlineto{\pgfqpoint{2.378551in}{1.880835in}}%
\pgfpathlineto{\pgfqpoint{2.378986in}{1.532026in}}%
\pgfpathlineto{\pgfqpoint{2.379851in}{1.532053in}}%
\pgfpathlineto{\pgfqpoint{2.380862in}{1.532150in}}%
\pgfpathlineto{\pgfqpoint{2.381089in}{1.943500in}}%
\pgfpathlineto{\pgfqpoint{2.381960in}{1.915679in}}%
\pgfpathlineto{\pgfqpoint{2.382257in}{1.919389in}}%
\pgfpathlineto{\pgfqpoint{2.382702in}{1.918707in}}%
\pgfpathlineto{\pgfqpoint{2.383493in}{1.879331in}}%
\pgfpathlineto{\pgfqpoint{2.383922in}{1.532026in}}%
\pgfpathlineto{\pgfqpoint{2.384774in}{1.532053in}}%
\pgfpathlineto{\pgfqpoint{2.385805in}{1.532112in}}%
\pgfpathlineto{\pgfqpoint{2.385845in}{1.722802in}}%
\pgfpathlineto{\pgfqpoint{2.386123in}{1.938699in}}%
\pgfpathlineto{\pgfqpoint{2.386583in}{1.918136in}}%
\pgfpathlineto{\pgfqpoint{2.386781in}{1.915707in}}%
\pgfpathlineto{\pgfqpoint{2.386830in}{1.919638in}}%
\pgfpathlineto{\pgfqpoint{2.387226in}{1.916484in}}%
\pgfpathlineto{\pgfqpoint{2.387275in}{1.918952in}}%
\pgfpathlineto{\pgfqpoint{2.387967in}{1.917624in}}%
\pgfpathlineto{\pgfqpoint{2.388066in}{1.917778in}}%
\pgfpathlineto{\pgfqpoint{2.388363in}{1.910566in}}%
\pgfpathlineto{\pgfqpoint{2.388770in}{1.543599in}}%
\pgfpathlineto{\pgfqpoint{2.388883in}{1.532025in}}%
\pgfpathlineto{\pgfqpoint{2.389504in}{1.532039in}}%
\pgfpathlineto{\pgfqpoint{2.390753in}{1.532150in}}%
\pgfpathlineto{\pgfqpoint{2.390975in}{1.939697in}}%
\pgfpathlineto{\pgfqpoint{2.391850in}{1.914713in}}%
\pgfpathlineto{\pgfqpoint{2.392147in}{1.921091in}}%
\pgfpathlineto{\pgfqpoint{2.392098in}{1.914495in}}%
\pgfpathlineto{\pgfqpoint{2.392592in}{1.919741in}}%
\pgfpathlineto{\pgfqpoint{2.393482in}{1.813008in}}%
\pgfpathlineto{\pgfqpoint{2.393821in}{1.532024in}}%
\pgfpathlineto{\pgfqpoint{2.394585in}{1.532053in}}%
\pgfpathlineto{\pgfqpoint{2.395697in}{1.532123in}}%
\pgfpathlineto{\pgfqpoint{2.395918in}{1.940573in}}%
\pgfpathlineto{\pgfqpoint{2.396994in}{1.917763in}}%
\pgfpathlineto{\pgfqpoint{2.398230in}{1.915575in}}%
\pgfpathlineto{\pgfqpoint{2.398650in}{1.653528in}}%
\pgfpathlineto{\pgfqpoint{2.398766in}{1.532024in}}%
\pgfpathlineto{\pgfqpoint{2.399383in}{1.532038in}}%
\pgfpathlineto{\pgfqpoint{2.400642in}{1.532123in}}%
\pgfpathlineto{\pgfqpoint{2.400862in}{1.940028in}}%
\pgfpathlineto{\pgfqpoint{2.401933in}{1.917657in}}%
\pgfpathlineto{\pgfqpoint{2.402032in}{1.917506in}}%
\pgfpathlineto{\pgfqpoint{2.402626in}{1.917622in}}%
\pgfpathlineto{\pgfqpoint{2.403170in}{1.915941in}}%
\pgfpathlineto{\pgfqpoint{2.403516in}{1.704806in}}%
\pgfpathlineto{\pgfqpoint{2.403705in}{1.532026in}}%
\pgfpathlineto{\pgfqpoint{2.404511in}{1.532043in}}%
\pgfpathlineto{\pgfqpoint{2.405587in}{1.532116in}}%
\pgfpathlineto{\pgfqpoint{2.405802in}{1.938344in}}%
\pgfpathlineto{\pgfqpoint{2.407086in}{1.917861in}}%
\pgfpathlineto{\pgfqpoint{2.407234in}{1.916609in}}%
\pgfpathlineto{\pgfqpoint{2.407284in}{1.918618in}}%
\pgfpathlineto{\pgfqpoint{2.407877in}{1.916996in}}%
\pgfpathlineto{\pgfqpoint{2.407976in}{1.917167in}}%
\pgfpathlineto{\pgfqpoint{2.408025in}{1.916596in}}%
\pgfpathlineto{\pgfqpoint{2.408075in}{1.917015in}}%
\pgfpathlineto{\pgfqpoint{2.408174in}{1.904727in}}%
\pgfpathlineto{\pgfqpoint{2.408664in}{1.532022in}}%
\pgfpathlineto{\pgfqpoint{2.409680in}{1.532044in}}%
\pgfpathlineto{\pgfqpoint{2.410536in}{1.532179in}}%
\pgfpathlineto{\pgfqpoint{2.410652in}{1.938212in}}%
\pgfpathlineto{\pgfqpoint{2.411640in}{1.917560in}}%
\pgfpathlineto{\pgfqpoint{2.412481in}{1.917500in}}%
\pgfpathlineto{\pgfqpoint{2.412530in}{1.917451in}}%
\pgfpathlineto{\pgfqpoint{2.413025in}{1.916562in}}%
\pgfpathlineto{\pgfqpoint{2.413173in}{1.875686in}}%
\pgfpathlineto{\pgfqpoint{2.413595in}{1.532026in}}%
\pgfpathlineto{\pgfqpoint{2.414449in}{1.532053in}}%
\pgfpathlineto{\pgfqpoint{2.415478in}{1.532112in}}%
\pgfpathlineto{\pgfqpoint{2.415518in}{1.722387in}}%
\pgfpathlineto{\pgfqpoint{2.415692in}{1.939072in}}%
\pgfpathlineto{\pgfqpoint{2.416282in}{1.915987in}}%
\pgfpathlineto{\pgfqpoint{2.416431in}{1.920129in}}%
\pgfpathlineto{\pgfqpoint{2.416381in}{1.914926in}}%
\pgfpathlineto{\pgfqpoint{2.417074in}{1.917696in}}%
\pgfpathlineto{\pgfqpoint{2.417321in}{1.918528in}}%
\pgfpathlineto{\pgfqpoint{2.417272in}{1.916503in}}%
\pgfpathlineto{\pgfqpoint{2.417667in}{1.917895in}}%
\pgfpathlineto{\pgfqpoint{2.418063in}{1.904753in}}%
\pgfpathlineto{\pgfqpoint{2.418543in}{1.532025in}}%
\pgfpathlineto{\pgfqpoint{2.419641in}{1.532049in}}%
\pgfpathlineto{\pgfqpoint{2.420423in}{1.532112in}}%
\pgfpathlineto{\pgfqpoint{2.420462in}{1.714477in}}%
\pgfpathlineto{\pgfqpoint{2.420637in}{1.941190in}}%
\pgfpathlineto{\pgfqpoint{2.421234in}{1.916518in}}%
\pgfpathlineto{\pgfqpoint{2.421530in}{1.918904in}}%
\pgfpathlineto{\pgfqpoint{2.421481in}{1.916165in}}%
\pgfpathlineto{\pgfqpoint{2.421976in}{1.918415in}}%
\pgfpathlineto{\pgfqpoint{2.422025in}{1.916848in}}%
\pgfpathlineto{\pgfqpoint{2.422767in}{1.917370in}}%
\pgfpathlineto{\pgfqpoint{2.423014in}{1.901718in}}%
\pgfpathlineto{\pgfqpoint{2.423497in}{1.532028in}}%
\pgfpathlineto{\pgfqpoint{2.424510in}{1.532050in}}%
\pgfpathlineto{\pgfqpoint{2.425369in}{1.532119in}}%
\pgfpathlineto{\pgfqpoint{2.425608in}{1.942841in}}%
\pgfpathlineto{\pgfqpoint{2.426603in}{1.917004in}}%
\pgfpathlineto{\pgfqpoint{2.426900in}{1.918029in}}%
\pgfpathlineto{\pgfqpoint{2.427345in}{1.917710in}}%
\pgfpathlineto{\pgfqpoint{2.427938in}{1.908220in}}%
\pgfpathlineto{\pgfqpoint{2.428438in}{1.532027in}}%
\pgfpathlineto{\pgfqpoint{2.429597in}{1.532059in}}%
\pgfpathlineto{\pgfqpoint{2.430314in}{1.532113in}}%
\pgfpathlineto{\pgfqpoint{2.430418in}{1.936831in}}%
\pgfpathlineto{\pgfqpoint{2.431842in}{1.920179in}}%
\pgfpathlineto{\pgfqpoint{2.431892in}{1.914881in}}%
\pgfpathlineto{\pgfqpoint{2.432633in}{1.917781in}}%
\pgfpathlineto{\pgfqpoint{2.432732in}{1.917897in}}%
\pgfpathlineto{\pgfqpoint{2.432930in}{1.892476in}}%
\pgfpathlineto{\pgfqpoint{2.433359in}{1.532030in}}%
\pgfpathlineto{\pgfqpoint{2.434343in}{1.532051in}}%
\pgfpathlineto{\pgfqpoint{2.435265in}{1.532280in}}%
\pgfpathlineto{\pgfqpoint{2.435369in}{1.936222in}}%
\pgfpathlineto{\pgfqpoint{2.436292in}{1.914361in}}%
\pgfpathlineto{\pgfqpoint{2.436589in}{1.920636in}}%
\pgfpathlineto{\pgfqpoint{2.437034in}{1.919431in}}%
\pgfpathlineto{\pgfqpoint{2.438023in}{1.789306in}}%
\pgfpathlineto{\pgfqpoint{2.438329in}{1.532024in}}%
\pgfpathlineto{\pgfqpoint{2.439093in}{1.532053in}}%
\pgfpathlineto{\pgfqpoint{2.440206in}{1.532124in}}%
\pgfpathlineto{\pgfqpoint{2.440426in}{1.941670in}}%
\pgfpathlineto{\pgfqpoint{2.441468in}{1.917724in}}%
\pgfpathlineto{\pgfqpoint{2.442061in}{1.917290in}}%
\pgfpathlineto{\pgfqpoint{2.441567in}{1.917739in}}%
\pgfpathlineto{\pgfqpoint{2.442259in}{1.917352in}}%
\pgfpathlineto{\pgfqpoint{2.442753in}{1.912967in}}%
\pgfpathlineto{\pgfqpoint{2.443163in}{1.578094in}}%
\pgfpathlineto{\pgfqpoint{2.443275in}{1.532024in}}%
\pgfpathlineto{\pgfqpoint{2.443892in}{1.532038in}}%
\pgfpathlineto{\pgfqpoint{2.445151in}{1.532120in}}%
\pgfpathlineto{\pgfqpoint{2.445373in}{1.939462in}}%
\pgfpathlineto{\pgfqpoint{2.446435in}{1.917447in}}%
\pgfpathlineto{\pgfqpoint{2.447029in}{1.917492in}}%
\pgfpathlineto{\pgfqpoint{2.446534in}{1.917340in}}%
\pgfpathlineto{\pgfqpoint{2.447177in}{1.917234in}}%
\pgfpathlineto{\pgfqpoint{2.447226in}{1.917307in}}%
\pgfpathlineto{\pgfqpoint{2.447622in}{1.916529in}}%
\pgfpathlineto{\pgfqpoint{2.447671in}{1.916084in}}%
\pgfpathlineto{\pgfqpoint{2.447968in}{1.747978in}}%
\pgfpathlineto{\pgfqpoint{2.448222in}{1.532024in}}%
\pgfpathlineto{\pgfqpoint{2.448988in}{1.532053in}}%
\pgfpathlineto{\pgfqpoint{2.450097in}{1.532128in}}%
\pgfpathlineto{\pgfqpoint{2.450432in}{1.942983in}}%
\pgfpathlineto{\pgfqpoint{2.451277in}{1.917492in}}%
\pgfpathlineto{\pgfqpoint{2.451524in}{1.918169in}}%
\pgfpathlineto{\pgfqpoint{2.451475in}{1.916701in}}%
\pgfpathlineto{\pgfqpoint{2.451969in}{1.917799in}}%
\pgfpathlineto{\pgfqpoint{2.452662in}{1.909699in}}%
\pgfpathlineto{\pgfqpoint{2.453168in}{1.532024in}}%
\pgfpathlineto{\pgfqpoint{2.454379in}{1.532051in}}%
\pgfpathlineto{\pgfqpoint{2.455043in}{1.532129in}}%
\pgfpathlineto{\pgfqpoint{2.455278in}{1.941242in}}%
\pgfpathlineto{\pgfqpoint{2.456245in}{1.917880in}}%
\pgfpathlineto{\pgfqpoint{2.456542in}{1.916807in}}%
\pgfpathlineto{\pgfqpoint{2.456492in}{1.918044in}}%
\pgfpathlineto{\pgfqpoint{2.457086in}{1.916990in}}%
\pgfpathlineto{\pgfqpoint{2.457135in}{1.917348in}}%
\pgfpathlineto{\pgfqpoint{2.457531in}{1.916276in}}%
\pgfpathlineto{\pgfqpoint{2.457630in}{1.902832in}}%
\pgfpathlineto{\pgfqpoint{2.458111in}{1.532025in}}%
\pgfpathlineto{\pgfqpoint{2.459122in}{1.532047in}}%
\pgfpathlineto{\pgfqpoint{2.459987in}{1.532112in}}%
\pgfpathlineto{\pgfqpoint{2.460089in}{1.936162in}}%
\pgfpathlineto{\pgfqpoint{2.461548in}{1.915733in}}%
\pgfpathlineto{\pgfqpoint{2.461845in}{1.919083in}}%
\pgfpathlineto{\pgfqpoint{2.461894in}{1.915658in}}%
\pgfpathlineto{\pgfqpoint{2.462290in}{1.918058in}}%
\pgfpathlineto{\pgfqpoint{2.462636in}{1.875585in}}%
\pgfpathlineto{\pgfqpoint{2.463067in}{1.532030in}}%
\pgfpathlineto{\pgfqpoint{2.463886in}{1.532048in}}%
\pgfpathlineto{\pgfqpoint{2.464937in}{1.532230in}}%
\pgfpathlineto{\pgfqpoint{2.465042in}{1.936419in}}%
\pgfpathlineto{\pgfqpoint{2.466010in}{1.917947in}}%
\pgfpathlineto{\pgfqpoint{2.466060in}{1.917997in}}%
\pgfpathlineto{\pgfqpoint{2.466208in}{1.914699in}}%
\pgfpathlineto{\pgfqpoint{2.466258in}{1.920094in}}%
\pgfpathlineto{\pgfqpoint{2.466851in}{1.916739in}}%
\pgfpathlineto{\pgfqpoint{2.467049in}{1.918147in}}%
\pgfpathlineto{\pgfqpoint{2.467445in}{1.915564in}}%
\pgfpathlineto{\pgfqpoint{2.467642in}{1.830448in}}%
\pgfpathlineto{\pgfqpoint{2.468002in}{1.532024in}}%
\pgfpathlineto{\pgfqpoint{2.468766in}{1.532053in}}%
\pgfpathlineto{\pgfqpoint{2.469879in}{1.532123in}}%
\pgfpathlineto{\pgfqpoint{2.470099in}{1.941498in}}%
\pgfpathlineto{\pgfqpoint{2.471132in}{1.917625in}}%
\pgfpathlineto{\pgfqpoint{2.471182in}{1.917113in}}%
\pgfpathlineto{\pgfqpoint{2.471924in}{1.917258in}}%
\pgfpathlineto{\pgfqpoint{2.472418in}{1.914859in}}%
\pgfpathlineto{\pgfqpoint{2.472831in}{1.661836in}}%
\pgfpathlineto{\pgfqpoint{2.472957in}{1.532025in}}%
\pgfpathlineto{\pgfqpoint{2.473579in}{1.532039in}}%
\pgfpathlineto{\pgfqpoint{2.474825in}{1.532137in}}%
\pgfpathlineto{\pgfqpoint{2.475157in}{1.943648in}}%
\pgfpathlineto{\pgfqpoint{2.475975in}{1.917057in}}%
\pgfpathlineto{\pgfqpoint{2.476173in}{1.918427in}}%
\pgfpathlineto{\pgfqpoint{2.476222in}{1.916267in}}%
\pgfpathlineto{\pgfqpoint{2.476618in}{1.917973in}}%
\pgfpathlineto{\pgfqpoint{2.477459in}{1.879741in}}%
\pgfpathlineto{\pgfqpoint{2.477896in}{1.532028in}}%
\pgfpathlineto{\pgfqpoint{2.478762in}{1.532052in}}%
\pgfpathlineto{\pgfqpoint{2.479769in}{1.532118in}}%
\pgfpathlineto{\pgfqpoint{2.479994in}{1.941132in}}%
\pgfpathlineto{\pgfqpoint{2.481020in}{1.916866in}}%
\pgfpathlineto{\pgfqpoint{2.481316in}{1.917762in}}%
\pgfpathlineto{\pgfqpoint{2.481762in}{1.917454in}}%
\pgfpathlineto{\pgfqpoint{2.482355in}{1.904294in}}%
\pgfpathlineto{\pgfqpoint{2.482847in}{1.532024in}}%
\pgfpathlineto{\pgfqpoint{2.483863in}{1.532046in}}%
\pgfpathlineto{\pgfqpoint{2.484713in}{1.532107in}}%
\pgfpathlineto{\pgfqpoint{2.484754in}{1.721170in}}%
\pgfpathlineto{\pgfqpoint{2.484926in}{1.937017in}}%
\pgfpathlineto{\pgfqpoint{2.485530in}{1.915088in}}%
\pgfpathlineto{\pgfqpoint{2.485827in}{1.920642in}}%
\pgfpathlineto{\pgfqpoint{2.485778in}{1.913871in}}%
\pgfpathlineto{\pgfqpoint{2.486272in}{1.919396in}}%
\pgfpathlineto{\pgfqpoint{2.486322in}{1.915584in}}%
\pgfpathlineto{\pgfqpoint{2.487064in}{1.917611in}}%
\pgfpathlineto{\pgfqpoint{2.487311in}{1.900825in}}%
\pgfpathlineto{\pgfqpoint{2.487771in}{1.532026in}}%
\pgfpathlineto{\pgfqpoint{2.488813in}{1.532057in}}%
\pgfpathlineto{\pgfqpoint{2.489659in}{1.532112in}}%
\pgfpathlineto{\pgfqpoint{2.489782in}{1.938203in}}%
\pgfpathlineto{\pgfqpoint{2.491204in}{1.917263in}}%
\pgfpathlineto{\pgfqpoint{2.491452in}{1.915708in}}%
\pgfpathlineto{\pgfqpoint{2.491402in}{1.918915in}}%
\pgfpathlineto{\pgfqpoint{2.491996in}{1.916315in}}%
\pgfpathlineto{\pgfqpoint{2.492045in}{1.916910in}}%
\pgfpathlineto{\pgfqpoint{2.492144in}{1.916002in}}%
\pgfpathlineto{\pgfqpoint{2.492194in}{1.916330in}}%
\pgfpathlineto{\pgfqpoint{2.492605in}{1.665946in}}%
\pgfpathlineto{\pgfqpoint{2.492613in}{1.666788in}}%
\pgfpathlineto{\pgfqpoint{2.492696in}{1.532028in}}%
\pgfpathlineto{\pgfqpoint{2.493377in}{1.532042in}}%
\pgfpathlineto{\pgfqpoint{2.494605in}{1.532110in}}%
\pgfpathlineto{\pgfqpoint{2.494721in}{1.937583in}}%
\pgfpathlineto{\pgfqpoint{2.496155in}{1.917016in}}%
\pgfpathlineto{\pgfqpoint{2.496303in}{1.917760in}}%
\pgfpathlineto{\pgfqpoint{2.496798in}{1.916738in}}%
\pgfpathlineto{\pgfqpoint{2.496897in}{1.916770in}}%
\pgfpathlineto{\pgfqpoint{2.496996in}{1.916757in}}%
\pgfpathlineto{\pgfqpoint{2.497045in}{1.916413in}}%
\pgfpathlineto{\pgfqpoint{2.497095in}{1.916489in}}%
\pgfpathlineto{\pgfqpoint{2.497193in}{1.904353in}}%
\pgfpathlineto{\pgfqpoint{2.497642in}{1.532028in}}%
\pgfpathlineto{\pgfqpoint{2.498718in}{1.532051in}}%
\pgfpathlineto{\pgfqpoint{2.499550in}{1.532110in}}%
\pgfpathlineto{\pgfqpoint{2.499666in}{1.937610in}}%
\pgfpathlineto{\pgfqpoint{2.501097in}{1.917568in}}%
\pgfpathlineto{\pgfqpoint{2.501938in}{1.916536in}}%
\pgfpathlineto{\pgfqpoint{2.501344in}{1.917618in}}%
\pgfpathlineto{\pgfqpoint{2.501987in}{1.916625in}}%
\pgfpathlineto{\pgfqpoint{2.502086in}{1.915377in}}%
\pgfpathlineto{\pgfqpoint{2.502432in}{1.704446in}}%
\pgfpathlineto{\pgfqpoint{2.502628in}{1.532028in}}%
\pgfpathlineto{\pgfqpoint{2.503445in}{1.532046in}}%
\pgfpathlineto{\pgfqpoint{2.504501in}{1.532271in}}%
\pgfpathlineto{\pgfqpoint{2.504617in}{1.938335in}}%
\pgfpathlineto{\pgfqpoint{2.505553in}{1.918290in}}%
\pgfpathlineto{\pgfqpoint{2.505850in}{1.916314in}}%
\pgfpathlineto{\pgfqpoint{2.506295in}{1.916606in}}%
\pgfpathlineto{\pgfqpoint{2.506344in}{1.917787in}}%
\pgfpathlineto{\pgfqpoint{2.506987in}{1.916228in}}%
\pgfpathlineto{\pgfqpoint{2.507086in}{1.902753in}}%
\pgfpathlineto{\pgfqpoint{2.507566in}{1.532027in}}%
\pgfpathlineto{\pgfqpoint{2.508578in}{1.532049in}}%
\pgfpathlineto{\pgfqpoint{2.509441in}{1.532112in}}%
\pgfpathlineto{\pgfqpoint{2.509563in}{1.937951in}}%
\pgfpathlineto{\pgfqpoint{2.510976in}{1.917767in}}%
\pgfpathlineto{\pgfqpoint{2.511124in}{1.916384in}}%
\pgfpathlineto{\pgfqpoint{2.511075in}{1.918100in}}%
\pgfpathlineto{\pgfqpoint{2.511767in}{1.916758in}}%
\pgfpathlineto{\pgfqpoint{2.511916in}{1.916063in}}%
\pgfpathlineto{\pgfqpoint{2.511866in}{1.916798in}}%
\pgfpathlineto{\pgfqpoint{2.511965in}{1.916279in}}%
\pgfpathlineto{\pgfqpoint{2.512212in}{1.789823in}}%
\pgfpathlineto{\pgfqpoint{2.512519in}{1.532030in}}%
\pgfpathlineto{\pgfqpoint{2.513286in}{1.532048in}}%
\pgfpathlineto{\pgfqpoint{2.514390in}{1.532160in}}%
\pgfpathlineto{\pgfqpoint{2.514497in}{1.936827in}}%
\pgfpathlineto{\pgfqpoint{2.515536in}{1.917079in}}%
\pgfpathlineto{\pgfqpoint{2.515832in}{1.917537in}}%
\pgfpathlineto{\pgfqpoint{2.515783in}{1.916943in}}%
\pgfpathlineto{\pgfqpoint{2.516278in}{1.917324in}}%
\pgfpathlineto{\pgfqpoint{2.516920in}{1.915324in}}%
\pgfpathlineto{\pgfqpoint{2.517299in}{1.682244in}}%
\pgfpathlineto{\pgfqpoint{2.517456in}{1.532026in}}%
\pgfpathlineto{\pgfqpoint{2.518318in}{1.532053in}}%
\pgfpathlineto{\pgfqpoint{2.519332in}{1.532112in}}%
\pgfpathlineto{\pgfqpoint{2.519372in}{1.724926in}}%
\pgfpathlineto{\pgfqpoint{2.519545in}{1.938042in}}%
\pgfpathlineto{\pgfqpoint{2.520122in}{1.918885in}}%
\pgfpathlineto{\pgfqpoint{2.520171in}{1.915804in}}%
\pgfpathlineto{\pgfqpoint{2.520863in}{1.916434in}}%
\pgfpathlineto{\pgfqpoint{2.520913in}{1.918081in}}%
\pgfpathlineto{\pgfqpoint{2.520962in}{1.916425in}}%
\pgfpathlineto{\pgfqpoint{2.521605in}{1.916889in}}%
\pgfpathlineto{\pgfqpoint{2.521704in}{1.916942in}}%
\pgfpathlineto{\pgfqpoint{2.521853in}{1.915714in}}%
\pgfpathlineto{\pgfqpoint{2.522100in}{1.792964in}}%
\pgfpathlineto{\pgfqpoint{2.522411in}{1.532030in}}%
\pgfpathlineto{\pgfqpoint{2.523179in}{1.532048in}}%
\pgfpathlineto{\pgfqpoint{2.524279in}{1.532130in}}%
\pgfpathlineto{\pgfqpoint{2.524399in}{1.937979in}}%
\pgfpathlineto{\pgfqpoint{2.525483in}{1.917380in}}%
\pgfpathlineto{\pgfqpoint{2.526324in}{1.916964in}}%
\pgfpathlineto{\pgfqpoint{2.526373in}{1.917062in}}%
\pgfpathlineto{\pgfqpoint{2.526769in}{1.916152in}}%
\pgfpathlineto{\pgfqpoint{2.526868in}{1.902467in}}%
\pgfpathlineto{\pgfqpoint{2.527352in}{1.532026in}}%
\pgfpathlineto{\pgfqpoint{2.528366in}{1.532048in}}%
\pgfpathlineto{\pgfqpoint{2.529225in}{1.532136in}}%
\pgfpathlineto{\pgfqpoint{2.529458in}{1.945754in}}%
\pgfpathlineto{\pgfqpoint{2.530352in}{1.918031in}}%
\pgfpathlineto{\pgfqpoint{2.530401in}{1.916048in}}%
\pgfpathlineto{\pgfqpoint{2.530451in}{1.918352in}}%
\pgfpathlineto{\pgfqpoint{2.531143in}{1.917148in}}%
\pgfpathlineto{\pgfqpoint{2.531242in}{1.917398in}}%
\pgfpathlineto{\pgfqpoint{2.531786in}{1.910398in}}%
\pgfpathlineto{\pgfqpoint{2.532303in}{1.532026in}}%
\pgfpathlineto{\pgfqpoint{2.533617in}{1.532055in}}%
\pgfpathlineto{\pgfqpoint{2.534171in}{1.532136in}}%
\pgfpathlineto{\pgfqpoint{2.534387in}{1.939975in}}%
\pgfpathlineto{\pgfqpoint{2.535319in}{1.917409in}}%
\pgfpathlineto{\pgfqpoint{2.535468in}{1.916370in}}%
\pgfpathlineto{\pgfqpoint{2.535517in}{1.917951in}}%
\pgfpathlineto{\pgfqpoint{2.536110in}{1.916924in}}%
\pgfpathlineto{\pgfqpoint{2.536308in}{1.917213in}}%
\pgfpathlineto{\pgfqpoint{2.536605in}{1.916392in}}%
\pgfpathlineto{\pgfqpoint{2.536704in}{1.915088in}}%
\pgfpathlineto{\pgfqpoint{2.537082in}{1.682896in}}%
\pgfpathlineto{\pgfqpoint{2.537250in}{1.532025in}}%
\pgfpathlineto{\pgfqpoint{2.538070in}{1.532043in}}%
\pgfpathlineto{\pgfqpoint{2.539116in}{1.532137in}}%
\pgfpathlineto{\pgfqpoint{2.539349in}{1.944622in}}%
\pgfpathlineto{\pgfqpoint{2.540252in}{1.917902in}}%
\pgfpathlineto{\pgfqpoint{2.540401in}{1.916020in}}%
\pgfpathlineto{\pgfqpoint{2.540351in}{1.918288in}}%
\pgfpathlineto{\pgfqpoint{2.541044in}{1.917055in}}%
\pgfpathlineto{\pgfqpoint{2.541241in}{1.917345in}}%
\pgfpathlineto{\pgfqpoint{2.541637in}{1.915481in}}%
\pgfpathlineto{\pgfqpoint{2.541884in}{1.787606in}}%
\pgfpathlineto{\pgfqpoint{2.542186in}{1.532026in}}%
\pgfpathlineto{\pgfqpoint{2.542951in}{1.532051in}}%
\pgfpathlineto{\pgfqpoint{2.544062in}{1.532147in}}%
\pgfpathlineto{\pgfqpoint{2.544291in}{1.941827in}}%
\pgfpathlineto{\pgfqpoint{2.545161in}{1.915982in}}%
\pgfpathlineto{\pgfqpoint{2.545458in}{1.918947in}}%
\pgfpathlineto{\pgfqpoint{2.545409in}{1.915261in}}%
\pgfpathlineto{\pgfqpoint{2.545903in}{1.918172in}}%
\pgfpathlineto{\pgfqpoint{2.546744in}{1.851611in}}%
\pgfpathlineto{\pgfqpoint{2.547141in}{1.532025in}}%
\pgfpathlineto{\pgfqpoint{2.547913in}{1.532053in}}%
\pgfpathlineto{\pgfqpoint{2.548995in}{1.532078in}}%
\pgfpathlineto{\pgfqpoint{2.549007in}{1.532134in}}%
\pgfpathlineto{\pgfqpoint{2.549338in}{1.939348in}}%
\pgfpathlineto{\pgfqpoint{2.550141in}{1.918423in}}%
\pgfpathlineto{\pgfqpoint{2.550438in}{1.915739in}}%
\pgfpathlineto{\pgfqpoint{2.550487in}{1.918445in}}%
\pgfpathlineto{\pgfqpoint{2.550883in}{1.916162in}}%
\pgfpathlineto{\pgfqpoint{2.550932in}{1.917817in}}%
\pgfpathlineto{\pgfqpoint{2.551526in}{1.915781in}}%
\pgfpathlineto{\pgfqpoint{2.551823in}{1.751529in}}%
\pgfpathlineto{\pgfqpoint{2.552084in}{1.532026in}}%
\pgfpathlineto{\pgfqpoint{2.552853in}{1.532052in}}%
\pgfpathlineto{\pgfqpoint{2.553953in}{1.532136in}}%
\pgfpathlineto{\pgfqpoint{2.554186in}{1.943261in}}%
\pgfpathlineto{\pgfqpoint{2.555104in}{1.917345in}}%
\pgfpathlineto{\pgfqpoint{2.555252in}{1.916080in}}%
\pgfpathlineto{\pgfqpoint{2.555302in}{1.918084in}}%
\pgfpathlineto{\pgfqpoint{2.555895in}{1.916783in}}%
\pgfpathlineto{\pgfqpoint{2.556093in}{1.917202in}}%
\pgfpathlineto{\pgfqpoint{2.556439in}{1.916235in}}%
\pgfpathlineto{\pgfqpoint{2.556538in}{1.904597in}}%
\pgfpathlineto{\pgfqpoint{2.557009in}{1.532026in}}%
\pgfpathlineto{\pgfqpoint{2.558053in}{1.532057in}}%
\pgfpathlineto{\pgfqpoint{2.558896in}{1.532112in}}%
\pgfpathlineto{\pgfqpoint{2.558936in}{1.728776in}}%
\pgfpathlineto{\pgfqpoint{2.559216in}{1.935477in}}%
\pgfpathlineto{\pgfqpoint{2.559705in}{1.914719in}}%
\pgfpathlineto{\pgfqpoint{2.560002in}{1.920170in}}%
\pgfpathlineto{\pgfqpoint{2.559952in}{1.913968in}}%
\pgfpathlineto{\pgfqpoint{2.560447in}{1.919027in}}%
\pgfpathlineto{\pgfqpoint{2.560496in}{1.915418in}}%
\pgfpathlineto{\pgfqpoint{2.561238in}{1.917308in}}%
\pgfpathlineto{\pgfqpoint{2.561485in}{1.905119in}}%
\pgfpathlineto{\pgfqpoint{2.561972in}{1.532030in}}%
\pgfpathlineto{\pgfqpoint{2.563036in}{1.532054in}}%
\pgfpathlineto{\pgfqpoint{2.563847in}{1.532279in}}%
\pgfpathlineto{\pgfqpoint{2.563951in}{1.935955in}}%
\pgfpathlineto{\pgfqpoint{2.564890in}{1.912977in}}%
\pgfpathlineto{\pgfqpoint{2.564940in}{1.920059in}}%
\pgfpathlineto{\pgfqpoint{2.565632in}{1.918775in}}%
\pgfpathlineto{\pgfqpoint{2.566572in}{1.818672in}}%
\pgfpathlineto{\pgfqpoint{2.566919in}{1.532022in}}%
\pgfpathlineto{\pgfqpoint{2.567688in}{1.532055in}}%
\pgfpathlineto{\pgfqpoint{2.568783in}{1.532092in}}%
\pgfpathlineto{\pgfqpoint{2.568791in}{1.532178in}}%
\pgfpathlineto{\pgfqpoint{2.568907in}{1.937764in}}%
\pgfpathlineto{\pgfqpoint{2.569897in}{1.917026in}}%
\pgfpathlineto{\pgfqpoint{2.570886in}{1.916840in}}%
\pgfpathlineto{\pgfqpoint{2.571281in}{1.916062in}}%
\pgfpathlineto{\pgfqpoint{2.571373in}{1.905583in}}%
\pgfpathlineto{\pgfqpoint{2.571866in}{1.532030in}}%
\pgfpathlineto{\pgfqpoint{2.572981in}{1.532054in}}%
\pgfpathlineto{\pgfqpoint{2.573736in}{1.532178in}}%
\pgfpathlineto{\pgfqpoint{2.573855in}{1.937872in}}%
\pgfpathlineto{\pgfqpoint{2.574850in}{1.917111in}}%
\pgfpathlineto{\pgfqpoint{2.576235in}{1.915993in}}%
\pgfpathlineto{\pgfqpoint{2.576383in}{1.872734in}}%
\pgfpathlineto{\pgfqpoint{2.576811in}{1.532022in}}%
\pgfpathlineto{\pgfqpoint{2.577679in}{1.532058in}}%
\pgfpathlineto{\pgfqpoint{2.578674in}{1.532092in}}%
\pgfpathlineto{\pgfqpoint{2.578682in}{1.532177in}}%
\pgfpathlineto{\pgfqpoint{2.578798in}{1.937772in}}%
\pgfpathlineto{\pgfqpoint{2.579813in}{1.917014in}}%
\pgfpathlineto{\pgfqpoint{2.581000in}{1.916571in}}%
\pgfpathlineto{\pgfqpoint{2.581225in}{1.914727in}}%
\pgfpathlineto{\pgfqpoint{2.581468in}{1.774310in}}%
\pgfpathlineto{\pgfqpoint{2.581759in}{1.532030in}}%
\pgfpathlineto{\pgfqpoint{2.582530in}{1.532048in}}%
\pgfpathlineto{\pgfqpoint{2.583626in}{1.532156in}}%
\pgfpathlineto{\pgfqpoint{2.583731in}{1.935891in}}%
\pgfpathlineto{\pgfqpoint{2.584756in}{1.917124in}}%
\pgfpathlineto{\pgfqpoint{2.585448in}{1.916889in}}%
\pgfpathlineto{\pgfqpoint{2.584855in}{1.917220in}}%
\pgfpathlineto{\pgfqpoint{2.585547in}{1.916916in}}%
\pgfpathlineto{\pgfqpoint{2.586140in}{1.915727in}}%
\pgfpathlineto{\pgfqpoint{2.586388in}{1.792784in}}%
\pgfpathlineto{\pgfqpoint{2.586654in}{1.532022in}}%
\pgfpathlineto{\pgfqpoint{2.587478in}{1.532056in}}%
\pgfpathlineto{\pgfqpoint{2.588565in}{1.532092in}}%
\pgfpathlineto{\pgfqpoint{2.588572in}{1.532177in}}%
\pgfpathlineto{\pgfqpoint{2.588689in}{1.937726in}}%
\pgfpathlineto{\pgfqpoint{2.589667in}{1.917002in}}%
\pgfpathlineto{\pgfqpoint{2.591052in}{1.916083in}}%
\pgfpathlineto{\pgfqpoint{2.591139in}{1.910043in}}%
\pgfpathlineto{\pgfqpoint{2.591652in}{1.532030in}}%
\pgfpathlineto{\pgfqpoint{2.592969in}{1.532059in}}%
\pgfpathlineto{\pgfqpoint{2.593517in}{1.532145in}}%
\pgfpathlineto{\pgfqpoint{2.593624in}{1.936433in}}%
\pgfpathlineto{\pgfqpoint{2.594673in}{1.916931in}}%
\pgfpathlineto{\pgfqpoint{2.595415in}{1.916968in}}%
\pgfpathlineto{\pgfqpoint{2.595514in}{1.916887in}}%
\pgfpathlineto{\pgfqpoint{2.596008in}{1.915979in}}%
\pgfpathlineto{\pgfqpoint{2.596107in}{1.902298in}}%
\pgfpathlineto{\pgfqpoint{2.596575in}{1.532026in}}%
\pgfpathlineto{\pgfqpoint{2.597619in}{1.532057in}}%
\pgfpathlineto{\pgfqpoint{2.598459in}{1.532112in}}%
\pgfpathlineto{\pgfqpoint{2.598514in}{1.908945in}}%
\pgfpathlineto{\pgfqpoint{2.598671in}{1.937263in}}%
\pgfpathlineto{\pgfqpoint{2.598622in}{1.894287in}}%
\pgfpathlineto{\pgfqpoint{2.599257in}{1.918578in}}%
\pgfpathlineto{\pgfqpoint{2.599307in}{1.915346in}}%
\pgfpathlineto{\pgfqpoint{2.599999in}{1.916550in}}%
\pgfpathlineto{\pgfqpoint{2.600147in}{1.917640in}}%
\pgfpathlineto{\pgfqpoint{2.600098in}{1.916304in}}%
\pgfpathlineto{\pgfqpoint{2.600741in}{1.916479in}}%
\pgfpathlineto{\pgfqpoint{2.600889in}{1.916089in}}%
\pgfpathlineto{\pgfqpoint{2.600939in}{1.916354in}}%
\pgfpathlineto{\pgfqpoint{2.601038in}{1.907991in}}%
\pgfpathlineto{\pgfqpoint{2.601522in}{1.532027in}}%
\pgfpathlineto{\pgfqpoint{2.602718in}{1.532053in}}%
\pgfpathlineto{\pgfqpoint{2.603405in}{1.532113in}}%
\pgfpathlineto{\pgfqpoint{2.603521in}{1.937444in}}%
\pgfpathlineto{\pgfqpoint{2.604918in}{1.916128in}}%
\pgfpathlineto{\pgfqpoint{2.604968in}{1.917917in}}%
\pgfpathlineto{\pgfqpoint{2.605017in}{1.916115in}}%
\pgfpathlineto{\pgfqpoint{2.605660in}{1.916674in}}%
\pgfpathlineto{\pgfqpoint{2.605759in}{1.916761in}}%
\pgfpathlineto{\pgfqpoint{2.605957in}{1.912979in}}%
\pgfpathlineto{\pgfqpoint{2.606360in}{1.647422in}}%
\pgfpathlineto{\pgfqpoint{2.606493in}{1.532030in}}%
\pgfpathlineto{\pgfqpoint{2.607121in}{1.532044in}}%
\pgfpathlineto{\pgfqpoint{2.608354in}{1.532159in}}%
\pgfpathlineto{\pgfqpoint{2.608461in}{1.936570in}}%
\pgfpathlineto{\pgfqpoint{2.609485in}{1.917060in}}%
\pgfpathlineto{\pgfqpoint{2.610178in}{1.916773in}}%
\pgfpathlineto{\pgfqpoint{2.609683in}{1.917259in}}%
\pgfpathlineto{\pgfqpoint{2.610276in}{1.916827in}}%
\pgfpathlineto{\pgfqpoint{2.610820in}{1.916142in}}%
\pgfpathlineto{\pgfqpoint{2.610919in}{1.909032in}}%
\pgfpathlineto{\pgfqpoint{2.611432in}{1.532022in}}%
\pgfpathlineto{\pgfqpoint{2.612648in}{1.532049in}}%
\pgfpathlineto{\pgfqpoint{2.613300in}{1.532178in}}%
\pgfpathlineto{\pgfqpoint{2.613416in}{1.937589in}}%
\pgfpathlineto{\pgfqpoint{2.614410in}{1.916945in}}%
\pgfpathlineto{\pgfqpoint{2.615647in}{1.916450in}}%
\pgfpathlineto{\pgfqpoint{2.615795in}{1.915937in}}%
\pgfpathlineto{\pgfqpoint{2.615894in}{1.901508in}}%
\pgfpathlineto{\pgfqpoint{2.616385in}{1.532030in}}%
\pgfpathlineto{\pgfqpoint{2.617408in}{1.532053in}}%
\pgfpathlineto{\pgfqpoint{2.618247in}{1.532279in}}%
\pgfpathlineto{\pgfqpoint{2.618351in}{1.935901in}}%
\pgfpathlineto{\pgfqpoint{2.619267in}{1.921530in}}%
\pgfpathlineto{\pgfqpoint{2.619317in}{1.912446in}}%
\pgfpathlineto{\pgfqpoint{2.620009in}{1.916069in}}%
\pgfpathlineto{\pgfqpoint{2.620157in}{1.918682in}}%
\pgfpathlineto{\pgfqpoint{2.620108in}{1.915133in}}%
\pgfpathlineto{\pgfqpoint{2.620751in}{1.915392in}}%
\pgfpathlineto{\pgfqpoint{2.620998in}{1.799636in}}%
\pgfpathlineto{\pgfqpoint{2.621323in}{1.532025in}}%
\pgfpathlineto{\pgfqpoint{2.622094in}{1.532052in}}%
\pgfpathlineto{\pgfqpoint{2.623189in}{1.532137in}}%
\pgfpathlineto{\pgfqpoint{2.623422in}{1.944840in}}%
\pgfpathlineto{\pgfqpoint{2.624331in}{1.917276in}}%
\pgfpathlineto{\pgfqpoint{2.624479in}{1.915868in}}%
\pgfpathlineto{\pgfqpoint{2.624529in}{1.917933in}}%
\pgfpathlineto{\pgfqpoint{2.625122in}{1.916658in}}%
\pgfpathlineto{\pgfqpoint{2.625320in}{1.917071in}}%
\pgfpathlineto{\pgfqpoint{2.625369in}{1.916214in}}%
\pgfpathlineto{\pgfqpoint{2.625666in}{1.916160in}}%
\pgfpathlineto{\pgfqpoint{2.625765in}{1.907647in}}%
\pgfpathlineto{\pgfqpoint{2.626264in}{1.532030in}}%
\pgfpathlineto{\pgfqpoint{2.627427in}{1.532057in}}%
\pgfpathlineto{\pgfqpoint{2.628136in}{1.532164in}}%
\pgfpathlineto{\pgfqpoint{2.628232in}{1.935655in}}%
\pgfpathlineto{\pgfqpoint{2.629281in}{1.916893in}}%
\pgfpathlineto{\pgfqpoint{2.630616in}{1.915997in}}%
\pgfpathlineto{\pgfqpoint{2.630715in}{1.906172in}}%
\pgfpathlineto{\pgfqpoint{2.631179in}{1.532030in}}%
\pgfpathlineto{\pgfqpoint{2.632316in}{1.532055in}}%
\pgfpathlineto{\pgfqpoint{2.633083in}{1.532282in}}%
\pgfpathlineto{\pgfqpoint{2.633188in}{1.935627in}}%
\pgfpathlineto{\pgfqpoint{2.634102in}{1.914949in}}%
\pgfpathlineto{\pgfqpoint{2.634300in}{1.921066in}}%
\pgfpathlineto{\pgfqpoint{2.634349in}{1.912497in}}%
\pgfpathlineto{\pgfqpoint{2.634844in}{1.919393in}}%
\pgfpathlineto{\pgfqpoint{2.635882in}{1.764123in}}%
\pgfpathlineto{\pgfqpoint{2.636131in}{1.532028in}}%
\pgfpathlineto{\pgfqpoint{2.636924in}{1.532045in}}%
\pgfpathlineto{\pgfqpoint{2.638022in}{1.532105in}}%
\pgfpathlineto{\pgfqpoint{2.638063in}{1.723409in}}%
\pgfpathlineto{\pgfqpoint{2.638147in}{1.937434in}}%
\pgfpathlineto{\pgfqpoint{2.638822in}{1.914902in}}%
\pgfpathlineto{\pgfqpoint{2.638971in}{1.919168in}}%
\pgfpathlineto{\pgfqpoint{2.638921in}{1.914312in}}%
\pgfpathlineto{\pgfqpoint{2.639614in}{1.916730in}}%
\pgfpathlineto{\pgfqpoint{2.639762in}{1.917847in}}%
\pgfpathlineto{\pgfqpoint{2.639811in}{1.915883in}}%
\pgfpathlineto{\pgfqpoint{2.640306in}{1.917025in}}%
\pgfpathlineto{\pgfqpoint{2.640603in}{1.907042in}}%
\pgfpathlineto{\pgfqpoint{2.641089in}{1.532027in}}%
\pgfpathlineto{\pgfqpoint{2.642338in}{1.532061in}}%
\pgfpathlineto{\pgfqpoint{2.642969in}{1.532112in}}%
\pgfpathlineto{\pgfqpoint{2.643088in}{1.937675in}}%
\pgfpathlineto{\pgfqpoint{2.644526in}{1.917157in}}%
\pgfpathlineto{\pgfqpoint{2.644674in}{1.915553in}}%
\pgfpathlineto{\pgfqpoint{2.644724in}{1.918205in}}%
\pgfpathlineto{\pgfqpoint{2.645317in}{1.916099in}}%
\pgfpathlineto{\pgfqpoint{2.645416in}{1.916326in}}%
\pgfpathlineto{\pgfqpoint{2.645515in}{1.915406in}}%
\pgfpathlineto{\pgfqpoint{2.645924in}{1.622524in}}%
\pgfpathlineto{\pgfqpoint{2.646004in}{1.532026in}}%
\pgfpathlineto{\pgfqpoint{2.646687in}{1.532040in}}%
\pgfpathlineto{\pgfqpoint{2.647916in}{1.532136in}}%
\pgfpathlineto{\pgfqpoint{2.648149in}{1.944227in}}%
\pgfpathlineto{\pgfqpoint{2.649060in}{1.917174in}}%
\pgfpathlineto{\pgfqpoint{2.649209in}{1.915898in}}%
\pgfpathlineto{\pgfqpoint{2.649258in}{1.917816in}}%
\pgfpathlineto{\pgfqpoint{2.649851in}{1.916605in}}%
\pgfpathlineto{\pgfqpoint{2.650049in}{1.916985in}}%
\pgfpathlineto{\pgfqpoint{2.650099in}{1.916190in}}%
\pgfpathlineto{\pgfqpoint{2.650395in}{1.916072in}}%
\pgfpathlineto{\pgfqpoint{2.650494in}{1.906602in}}%
\pgfpathlineto{\pgfqpoint{2.650977in}{1.532025in}}%
\pgfpathlineto{\pgfqpoint{2.652075in}{1.532049in}}%
\pgfpathlineto{\pgfqpoint{2.652859in}{1.532112in}}%
\pgfpathlineto{\pgfqpoint{2.652949in}{1.933868in}}%
\pgfpathlineto{\pgfqpoint{2.654430in}{1.918138in}}%
\pgfpathlineto{\pgfqpoint{2.654578in}{1.914281in}}%
\pgfpathlineto{\pgfqpoint{2.654529in}{1.919343in}}%
\pgfpathlineto{\pgfqpoint{2.655221in}{1.916153in}}%
\pgfpathlineto{\pgfqpoint{2.655320in}{1.916702in}}%
\pgfpathlineto{\pgfqpoint{2.655419in}{1.913643in}}%
\pgfpathlineto{\pgfqpoint{2.655824in}{1.543596in}}%
\pgfpathlineto{\pgfqpoint{2.655924in}{1.532028in}}%
\pgfpathlineto{\pgfqpoint{2.656575in}{1.532047in}}%
\pgfpathlineto{\pgfqpoint{2.657805in}{1.532116in}}%
\pgfpathlineto{\pgfqpoint{2.658034in}{1.942499in}}%
\pgfpathlineto{\pgfqpoint{2.659141in}{1.916974in}}%
\pgfpathlineto{\pgfqpoint{2.659240in}{1.917040in}}%
\pgfpathlineto{\pgfqpoint{2.659982in}{1.916536in}}%
\pgfpathlineto{\pgfqpoint{2.660131in}{1.916443in}}%
\pgfpathlineto{\pgfqpoint{2.660279in}{1.915929in}}%
\pgfpathlineto{\pgfqpoint{2.660328in}{1.915516in}}%
\pgfpathlineto{\pgfqpoint{2.660576in}{1.788410in}}%
\pgfpathlineto{\pgfqpoint{2.660883in}{1.532022in}}%
\pgfpathlineto{\pgfqpoint{2.661651in}{1.532055in}}%
\pgfpathlineto{\pgfqpoint{2.662747in}{1.532092in}}%
\pgfpathlineto{\pgfqpoint{2.662754in}{1.532181in}}%
\pgfpathlineto{\pgfqpoint{2.662856in}{1.936288in}}%
\pgfpathlineto{\pgfqpoint{2.663847in}{1.916810in}}%
\pgfpathlineto{\pgfqpoint{2.664787in}{1.916698in}}%
\pgfpathlineto{\pgfqpoint{2.665281in}{1.915443in}}%
\pgfpathlineto{\pgfqpoint{2.665503in}{1.805547in}}%
\pgfpathlineto{\pgfqpoint{2.665806in}{1.532028in}}%
\pgfpathlineto{\pgfqpoint{2.666601in}{1.532045in}}%
\pgfpathlineto{\pgfqpoint{2.667695in}{1.532108in}}%
\pgfpathlineto{\pgfqpoint{2.667737in}{1.745550in}}%
\pgfpathlineto{\pgfqpoint{2.667813in}{1.937433in}}%
\pgfpathlineto{\pgfqpoint{2.668521in}{1.913869in}}%
\pgfpathlineto{\pgfqpoint{2.668818in}{1.919042in}}%
\pgfpathlineto{\pgfqpoint{2.669263in}{1.918276in}}%
\pgfpathlineto{\pgfqpoint{2.669313in}{1.915395in}}%
\pgfpathlineto{\pgfqpoint{2.670054in}{1.916696in}}%
\pgfpathlineto{\pgfqpoint{2.670302in}{1.899061in}}%
\pgfpathlineto{\pgfqpoint{2.670773in}{1.532027in}}%
\pgfpathlineto{\pgfqpoint{2.671787in}{1.532049in}}%
\pgfpathlineto{\pgfqpoint{2.672643in}{1.532132in}}%
\pgfpathlineto{\pgfqpoint{2.672878in}{1.944877in}}%
\pgfpathlineto{\pgfqpoint{2.673810in}{1.915619in}}%
\pgfpathlineto{\pgfqpoint{2.673859in}{1.917845in}}%
\pgfpathlineto{\pgfqpoint{2.674552in}{1.917054in}}%
\pgfpathlineto{\pgfqpoint{2.674651in}{1.917137in}}%
\pgfpathlineto{\pgfqpoint{2.675244in}{1.897588in}}%
\pgfpathlineto{\pgfqpoint{2.675691in}{1.532028in}}%
\pgfpathlineto{\pgfqpoint{2.676680in}{1.532050in}}%
\pgfpathlineto{\pgfqpoint{2.677592in}{1.532283in}}%
\pgfpathlineto{\pgfqpoint{2.677682in}{1.935255in}}%
\pgfpathlineto{\pgfqpoint{2.678645in}{1.917511in}}%
\pgfpathlineto{\pgfqpoint{2.678694in}{1.915996in}}%
\pgfpathlineto{\pgfqpoint{2.678744in}{1.917595in}}%
\pgfpathlineto{\pgfqpoint{2.679436in}{1.916903in}}%
\pgfpathlineto{\pgfqpoint{2.679535in}{1.917016in}}%
\pgfpathlineto{\pgfqpoint{2.680178in}{1.904547in}}%
\pgfpathlineto{\pgfqpoint{2.680655in}{1.532027in}}%
\pgfpathlineto{\pgfqpoint{2.681760in}{1.532051in}}%
\pgfpathlineto{\pgfqpoint{2.682532in}{1.532112in}}%
\pgfpathlineto{\pgfqpoint{2.682653in}{1.937841in}}%
\pgfpathlineto{\pgfqpoint{2.684080in}{1.918725in}}%
\pgfpathlineto{\pgfqpoint{2.684129in}{1.915295in}}%
\pgfpathlineto{\pgfqpoint{2.684871in}{1.917120in}}%
\pgfpathlineto{\pgfqpoint{2.685168in}{1.885248in}}%
\pgfpathlineto{\pgfqpoint{2.685616in}{1.532026in}}%
\pgfpathlineto{\pgfqpoint{2.686538in}{1.532046in}}%
\pgfpathlineto{\pgfqpoint{2.687480in}{1.532136in}}%
\pgfpathlineto{\pgfqpoint{2.687712in}{1.944605in}}%
\pgfpathlineto{\pgfqpoint{2.688611in}{1.917515in}}%
\pgfpathlineto{\pgfqpoint{2.688759in}{1.915743in}}%
\pgfpathlineto{\pgfqpoint{2.688710in}{1.917873in}}%
\pgfpathlineto{\pgfqpoint{2.689402in}{1.916716in}}%
\pgfpathlineto{\pgfqpoint{2.689501in}{1.916960in}}%
\pgfpathlineto{\pgfqpoint{2.690045in}{1.909414in}}%
\pgfpathlineto{\pgfqpoint{2.690523in}{1.532029in}}%
\pgfpathlineto{\pgfqpoint{2.691807in}{1.532056in}}%
\pgfpathlineto{\pgfqpoint{2.692429in}{1.532380in}}%
\pgfpathlineto{\pgfqpoint{2.692529in}{1.935732in}}%
\pgfpathlineto{\pgfqpoint{2.693415in}{1.916373in}}%
\pgfpathlineto{\pgfqpoint{2.693563in}{1.917621in}}%
\pgfpathlineto{\pgfqpoint{2.693514in}{1.915973in}}%
\pgfpathlineto{\pgfqpoint{2.694206in}{1.916843in}}%
\pgfpathlineto{\pgfqpoint{2.694948in}{1.915371in}}%
\pgfpathlineto{\pgfqpoint{2.695195in}{1.791627in}}%
\pgfpathlineto{\pgfqpoint{2.695498in}{1.532025in}}%
\pgfpathlineto{\pgfqpoint{2.696265in}{1.532052in}}%
\pgfpathlineto{\pgfqpoint{2.697359in}{1.532079in}}%
\pgfpathlineto{\pgfqpoint{2.697371in}{1.532149in}}%
\pgfpathlineto{\pgfqpoint{2.697600in}{1.938659in}}%
\pgfpathlineto{\pgfqpoint{2.698481in}{1.914435in}}%
\pgfpathlineto{\pgfqpoint{2.698778in}{1.920258in}}%
\pgfpathlineto{\pgfqpoint{2.698728in}{1.913307in}}%
\pgfpathlineto{\pgfqpoint{2.699223in}{1.918856in}}%
\pgfpathlineto{\pgfqpoint{2.700113in}{1.811183in}}%
\pgfpathlineto{\pgfqpoint{2.700404in}{1.532026in}}%
\pgfpathlineto{\pgfqpoint{2.701235in}{1.532052in}}%
\pgfpathlineto{\pgfqpoint{2.702316in}{1.532136in}}%
\pgfpathlineto{\pgfqpoint{2.702549in}{1.943032in}}%
\pgfpathlineto{\pgfqpoint{2.703470in}{1.917047in}}%
\pgfpathlineto{\pgfqpoint{2.703618in}{1.915740in}}%
\pgfpathlineto{\pgfqpoint{2.703668in}{1.917769in}}%
\pgfpathlineto{\pgfqpoint{2.704261in}{1.916467in}}%
\pgfpathlineto{\pgfqpoint{2.704459in}{1.916886in}}%
\pgfpathlineto{\pgfqpoint{2.704805in}{1.915921in}}%
\pgfpathlineto{\pgfqpoint{2.704904in}{1.904472in}}%
\pgfpathlineto{\pgfqpoint{2.705398in}{1.532030in}}%
\pgfpathlineto{\pgfqpoint{2.706467in}{1.532055in}}%
\pgfpathlineto{\pgfqpoint{2.707263in}{1.532160in}}%
\pgfpathlineto{\pgfqpoint{2.707367in}{1.936200in}}%
\pgfpathlineto{\pgfqpoint{2.708410in}{1.916644in}}%
\pgfpathlineto{\pgfqpoint{2.708558in}{1.916996in}}%
\pgfpathlineto{\pgfqpoint{2.708509in}{1.916554in}}%
\pgfpathlineto{\pgfqpoint{2.709151in}{1.916620in}}%
\pgfpathlineto{\pgfqpoint{2.709448in}{1.916557in}}%
\pgfpathlineto{\pgfqpoint{2.709646in}{1.916180in}}%
\pgfpathlineto{\pgfqpoint{2.709794in}{1.915245in}}%
\pgfpathlineto{\pgfqpoint{2.710141in}{1.706949in}}%
\pgfpathlineto{\pgfqpoint{2.710294in}{1.532026in}}%
\pgfpathlineto{\pgfqpoint{2.711174in}{1.532044in}}%
\pgfpathlineto{\pgfqpoint{2.712207in}{1.532136in}}%
\pgfpathlineto{\pgfqpoint{2.712439in}{1.946053in}}%
\pgfpathlineto{\pgfqpoint{2.713319in}{1.917951in}}%
\pgfpathlineto{\pgfqpoint{2.713368in}{1.915477in}}%
\pgfpathlineto{\pgfqpoint{2.714061in}{1.916474in}}%
\pgfpathlineto{\pgfqpoint{2.714209in}{1.917081in}}%
\pgfpathlineto{\pgfqpoint{2.714704in}{1.915436in}}%
\pgfpathlineto{\pgfqpoint{2.714802in}{1.901833in}}%
\pgfpathlineto{\pgfqpoint{2.715279in}{1.532028in}}%
\pgfpathlineto{\pgfqpoint{2.716291in}{1.532050in}}%
\pgfpathlineto{\pgfqpoint{2.717156in}{1.532271in}}%
\pgfpathlineto{\pgfqpoint{2.717255in}{1.935490in}}%
\pgfpathlineto{\pgfqpoint{2.718195in}{1.916356in}}%
\pgfpathlineto{\pgfqpoint{2.718492in}{1.917253in}}%
\pgfpathlineto{\pgfqpoint{2.718443in}{1.916248in}}%
\pgfpathlineto{\pgfqpoint{2.718937in}{1.917007in}}%
\pgfpathlineto{\pgfqpoint{2.719728in}{1.907547in}}%
\pgfpathlineto{\pgfqpoint{2.720225in}{1.532025in}}%
\pgfpathlineto{\pgfqpoint{2.721387in}{1.532061in}}%
\pgfpathlineto{\pgfqpoint{2.722098in}{1.532136in}}%
\pgfpathlineto{\pgfqpoint{2.722330in}{1.945110in}}%
\pgfpathlineto{\pgfqpoint{2.723250in}{1.916725in}}%
\pgfpathlineto{\pgfqpoint{2.723399in}{1.915726in}}%
\pgfpathlineto{\pgfqpoint{2.723448in}{1.917771in}}%
\pgfpathlineto{\pgfqpoint{2.723943in}{1.916044in}}%
\pgfpathlineto{\pgfqpoint{2.723992in}{1.917093in}}%
\pgfpathlineto{\pgfqpoint{2.724635in}{1.914668in}}%
\pgfpathlineto{\pgfqpoint{2.724981in}{1.704830in}}%
\pgfpathlineto{\pgfqpoint{2.725177in}{1.532028in}}%
\pgfpathlineto{\pgfqpoint{2.725996in}{1.532046in}}%
\pgfpathlineto{\pgfqpoint{2.727041in}{1.532108in}}%
\pgfpathlineto{\pgfqpoint{2.727083in}{1.746745in}}%
\pgfpathlineto{\pgfqpoint{2.727142in}{1.935346in}}%
\pgfpathlineto{\pgfqpoint{2.727856in}{1.918578in}}%
\pgfpathlineto{\pgfqpoint{2.728153in}{1.912842in}}%
\pgfpathlineto{\pgfqpoint{2.728104in}{1.920775in}}%
\pgfpathlineto{\pgfqpoint{2.728598in}{1.914362in}}%
\pgfpathlineto{\pgfqpoint{2.728648in}{1.918629in}}%
\pgfpathlineto{\pgfqpoint{2.729340in}{1.917155in}}%
\pgfpathlineto{\pgfqpoint{2.729686in}{1.880135in}}%
\pgfpathlineto{\pgfqpoint{2.730077in}{1.532026in}}%
\pgfpathlineto{\pgfqpoint{2.731007in}{1.532054in}}%
\pgfpathlineto{\pgfqpoint{2.731989in}{1.532136in}}%
\pgfpathlineto{\pgfqpoint{2.732221in}{1.945772in}}%
\pgfpathlineto{\pgfqpoint{2.733111in}{1.916547in}}%
\pgfpathlineto{\pgfqpoint{2.733358in}{1.915849in}}%
\pgfpathlineto{\pgfqpoint{2.733309in}{1.917575in}}%
\pgfpathlineto{\pgfqpoint{2.733705in}{1.916162in}}%
\pgfpathlineto{\pgfqpoint{2.733754in}{1.917204in}}%
\pgfpathlineto{\pgfqpoint{2.734446in}{1.915907in}}%
\pgfpathlineto{\pgfqpoint{2.734545in}{1.912081in}}%
\pgfpathlineto{\pgfqpoint{2.734951in}{1.543467in}}%
\pgfpathlineto{\pgfqpoint{2.735066in}{1.532030in}}%
\pgfpathlineto{\pgfqpoint{2.735687in}{1.532043in}}%
\pgfpathlineto{\pgfqpoint{2.736934in}{1.532129in}}%
\pgfpathlineto{\pgfqpoint{2.737126in}{1.934829in}}%
\pgfpathlineto{\pgfqpoint{2.738177in}{1.916658in}}%
\pgfpathlineto{\pgfqpoint{2.738770in}{1.916733in}}%
\pgfpathlineto{\pgfqpoint{2.738276in}{1.916583in}}%
\pgfpathlineto{\pgfqpoint{2.738919in}{1.916520in}}%
\pgfpathlineto{\pgfqpoint{2.739215in}{1.916355in}}%
\pgfpathlineto{\pgfqpoint{2.739364in}{1.915969in}}%
\pgfpathlineto{\pgfqpoint{2.739463in}{1.915308in}}%
\pgfpathlineto{\pgfqpoint{2.739759in}{1.747316in}}%
\pgfpathlineto{\pgfqpoint{2.740015in}{1.532026in}}%
\pgfpathlineto{\pgfqpoint{2.740787in}{1.532052in}}%
\pgfpathlineto{\pgfqpoint{2.741880in}{1.532136in}}%
\pgfpathlineto{\pgfqpoint{2.742112in}{1.944683in}}%
\pgfpathlineto{\pgfqpoint{2.743034in}{1.915661in}}%
\pgfpathlineto{\pgfqpoint{2.743083in}{1.917564in}}%
\pgfpathlineto{\pgfqpoint{2.743776in}{1.916840in}}%
\pgfpathlineto{\pgfqpoint{2.743875in}{1.916925in}}%
\pgfpathlineto{\pgfqpoint{2.744468in}{1.904894in}}%
\pgfpathlineto{\pgfqpoint{2.744943in}{1.532027in}}%
\pgfpathlineto{\pgfqpoint{2.746044in}{1.532051in}}%
\pgfpathlineto{\pgfqpoint{2.746823in}{1.532112in}}%
\pgfpathlineto{\pgfqpoint{2.746926in}{1.935371in}}%
\pgfpathlineto{\pgfqpoint{2.748336in}{1.919371in}}%
\pgfpathlineto{\pgfqpoint{2.748386in}{1.914101in}}%
\pgfpathlineto{\pgfqpoint{2.749128in}{1.917052in}}%
\pgfpathlineto{\pgfqpoint{2.749227in}{1.917109in}}%
\pgfpathlineto{\pgfqpoint{2.749474in}{1.877413in}}%
\pgfpathlineto{\pgfqpoint{2.749889in}{1.532027in}}%
\pgfpathlineto{\pgfqpoint{2.750789in}{1.532047in}}%
\pgfpathlineto{\pgfqpoint{2.751769in}{1.532119in}}%
\pgfpathlineto{\pgfqpoint{2.751996in}{1.942142in}}%
\pgfpathlineto{\pgfqpoint{2.753062in}{1.916548in}}%
\pgfpathlineto{\pgfqpoint{2.753210in}{1.916871in}}%
\pgfpathlineto{\pgfqpoint{2.753161in}{1.916479in}}%
\pgfpathlineto{\pgfqpoint{2.753804in}{1.916457in}}%
\pgfpathlineto{\pgfqpoint{2.754249in}{1.915737in}}%
\pgfpathlineto{\pgfqpoint{2.754348in}{1.907119in}}%
\pgfpathlineto{\pgfqpoint{2.754845in}{1.532030in}}%
\pgfpathlineto{\pgfqpoint{2.756008in}{1.532057in}}%
\pgfpathlineto{\pgfqpoint{2.756717in}{1.532155in}}%
\pgfpathlineto{\pgfqpoint{2.756823in}{1.936498in}}%
\pgfpathlineto{\pgfqpoint{2.757885in}{1.916396in}}%
\pgfpathlineto{\pgfqpoint{2.757935in}{1.916891in}}%
\pgfpathlineto{\pgfqpoint{2.758627in}{1.916650in}}%
\pgfpathlineto{\pgfqpoint{2.759269in}{1.913645in}}%
\pgfpathlineto{\pgfqpoint{2.759624in}{1.688465in}}%
\pgfpathlineto{\pgfqpoint{2.759783in}{1.532026in}}%
\pgfpathlineto{\pgfqpoint{2.760645in}{1.532053in}}%
\pgfpathlineto{\pgfqpoint{2.761659in}{1.532112in}}%
\pgfpathlineto{\pgfqpoint{2.761700in}{1.724856in}}%
\pgfpathlineto{\pgfqpoint{2.761786in}{1.937453in}}%
\pgfpathlineto{\pgfqpoint{2.762450in}{1.918254in}}%
\pgfpathlineto{\pgfqpoint{2.762499in}{1.915125in}}%
\pgfpathlineto{\pgfqpoint{2.763191in}{1.916089in}}%
\pgfpathlineto{\pgfqpoint{2.763241in}{1.917368in}}%
\pgfpathlineto{\pgfqpoint{2.763290in}{1.915941in}}%
\pgfpathlineto{\pgfqpoint{2.763933in}{1.916235in}}%
\pgfpathlineto{\pgfqpoint{2.764032in}{1.916296in}}%
\pgfpathlineto{\pgfqpoint{2.764180in}{1.915213in}}%
\pgfpathlineto{\pgfqpoint{2.764378in}{1.833339in}}%
\pgfpathlineto{\pgfqpoint{2.764721in}{1.532027in}}%
\pgfpathlineto{\pgfqpoint{2.765520in}{1.532045in}}%
\pgfpathlineto{\pgfqpoint{2.766606in}{1.532121in}}%
\pgfpathlineto{\pgfqpoint{2.766830in}{1.941093in}}%
\pgfpathlineto{\pgfqpoint{2.767869in}{1.916463in}}%
\pgfpathlineto{\pgfqpoint{2.768166in}{1.916843in}}%
\pgfpathlineto{\pgfqpoint{2.768611in}{1.916618in}}%
\pgfpathlineto{\pgfqpoint{2.769155in}{1.913281in}}%
\pgfpathlineto{\pgfqpoint{2.769561in}{1.610658in}}%
\pgfpathlineto{\pgfqpoint{2.769653in}{1.532028in}}%
\pgfpathlineto{\pgfqpoint{2.770296in}{1.532046in}}%
\pgfpathlineto{\pgfqpoint{2.771547in}{1.532094in}}%
\pgfpathlineto{\pgfqpoint{2.771556in}{1.532283in}}%
\pgfpathlineto{\pgfqpoint{2.771659in}{1.936656in}}%
\pgfpathlineto{\pgfqpoint{2.772563in}{1.918167in}}%
\pgfpathlineto{\pgfqpoint{2.772612in}{1.915299in}}%
\pgfpathlineto{\pgfqpoint{2.773305in}{1.916102in}}%
\pgfpathlineto{\pgfqpoint{2.773354in}{1.917224in}}%
\pgfpathlineto{\pgfqpoint{2.774047in}{1.915631in}}%
\pgfpathlineto{\pgfqpoint{2.774145in}{1.903675in}}%
\pgfpathlineto{\pgfqpoint{2.774634in}{1.532026in}}%
\pgfpathlineto{\pgfqpoint{2.775654in}{1.532049in}}%
\pgfpathlineto{\pgfqpoint{2.776498in}{1.532131in}}%
\pgfpathlineto{\pgfqpoint{2.776732in}{1.944831in}}%
\pgfpathlineto{\pgfqpoint{2.777663in}{1.915481in}}%
\pgfpathlineto{\pgfqpoint{2.777712in}{1.917540in}}%
\pgfpathlineto{\pgfqpoint{2.778404in}{1.916898in}}%
\pgfpathlineto{\pgfqpoint{2.778503in}{1.916932in}}%
\pgfpathlineto{\pgfqpoint{2.779097in}{1.899987in}}%
\pgfpathlineto{\pgfqpoint{2.779551in}{1.532028in}}%
\pgfpathlineto{\pgfqpoint{2.780592in}{1.532055in}}%
\pgfpathlineto{\pgfqpoint{2.781441in}{1.532108in}}%
\pgfpathlineto{\pgfqpoint{2.781483in}{1.744556in}}%
\pgfpathlineto{\pgfqpoint{2.781543in}{1.935141in}}%
\pgfpathlineto{\pgfqpoint{2.782260in}{1.919993in}}%
\pgfpathlineto{\pgfqpoint{2.782557in}{1.911783in}}%
\pgfpathlineto{\pgfqpoint{2.782508in}{1.921391in}}%
\pgfpathlineto{\pgfqpoint{2.783002in}{1.913641in}}%
\pgfpathlineto{\pgfqpoint{2.783052in}{1.919228in}}%
\pgfpathlineto{\pgfqpoint{2.783744in}{1.917091in}}%
\pgfpathlineto{\pgfqpoint{2.783843in}{1.917221in}}%
\pgfpathlineto{\pgfqpoint{2.784090in}{1.879227in}}%
\pgfpathlineto{\pgfqpoint{2.784503in}{1.532027in}}%
\pgfpathlineto{\pgfqpoint{2.785400in}{1.532047in}}%
\pgfpathlineto{\pgfqpoint{2.786388in}{1.532121in}}%
\pgfpathlineto{\pgfqpoint{2.786613in}{1.941944in}}%
\pgfpathlineto{\pgfqpoint{2.787655in}{1.916541in}}%
\pgfpathlineto{\pgfqpoint{2.788397in}{1.916584in}}%
\pgfpathlineto{\pgfqpoint{2.788496in}{1.916447in}}%
\pgfpathlineto{\pgfqpoint{2.788891in}{1.915491in}}%
\pgfpathlineto{\pgfqpoint{2.789040in}{1.873341in}}%
\pgfpathlineto{\pgfqpoint{2.789465in}{1.532030in}}%
\pgfpathlineto{\pgfqpoint{2.790332in}{1.532050in}}%
\pgfpathlineto{\pgfqpoint{2.791334in}{1.532130in}}%
\pgfpathlineto{\pgfqpoint{2.791436in}{1.935775in}}%
\pgfpathlineto{\pgfqpoint{2.792563in}{1.916482in}}%
\pgfpathlineto{\pgfqpoint{2.792859in}{1.916748in}}%
\pgfpathlineto{\pgfqpoint{2.793305in}{1.916569in}}%
\pgfpathlineto{\pgfqpoint{2.793849in}{1.915487in}}%
\pgfpathlineto{\pgfqpoint{2.794046in}{1.834771in}}%
\pgfpathlineto{\pgfqpoint{2.794379in}{1.532029in}}%
\pgfpathlineto{\pgfqpoint{2.795219in}{1.532050in}}%
\pgfpathlineto{\pgfqpoint{2.796283in}{1.532282in}}%
\pgfpathlineto{\pgfqpoint{2.796388in}{1.935441in}}%
\pgfpathlineto{\pgfqpoint{2.797299in}{1.914570in}}%
\pgfpathlineto{\pgfqpoint{2.797596in}{1.920752in}}%
\pgfpathlineto{\pgfqpoint{2.797547in}{1.912199in}}%
\pgfpathlineto{\pgfqpoint{2.798041in}{1.919106in}}%
\pgfpathlineto{\pgfqpoint{2.799080in}{1.768908in}}%
\pgfpathlineto{\pgfqpoint{2.799358in}{1.532026in}}%
\pgfpathlineto{\pgfqpoint{2.800128in}{1.532051in}}%
\pgfpathlineto{\pgfqpoint{2.801226in}{1.532150in}}%
\pgfpathlineto{\pgfqpoint{2.801454in}{1.942016in}}%
\pgfpathlineto{\pgfqpoint{2.802309in}{1.915060in}}%
\pgfpathlineto{\pgfqpoint{2.802358in}{1.918080in}}%
\pgfpathlineto{\pgfqpoint{2.803051in}{1.916738in}}%
\pgfpathlineto{\pgfqpoint{2.803149in}{1.917035in}}%
\pgfpathlineto{\pgfqpoint{2.803792in}{1.908780in}}%
\pgfpathlineto{\pgfqpoint{2.804287in}{1.532028in}}%
\pgfpathlineto{\pgfqpoint{2.805631in}{1.532062in}}%
\pgfpathlineto{\pgfqpoint{2.806169in}{1.532113in}}%
\pgfpathlineto{\pgfqpoint{2.806366in}{1.937812in}}%
\pgfpathlineto{\pgfqpoint{2.807527in}{1.916481in}}%
\pgfpathlineto{\pgfqpoint{2.807576in}{1.916725in}}%
\pgfpathlineto{\pgfqpoint{2.808269in}{1.916415in}}%
\pgfpathlineto{\pgfqpoint{2.808714in}{1.914655in}}%
\pgfpathlineto{\pgfqpoint{2.809114in}{1.669001in}}%
\pgfpathlineto{\pgfqpoint{2.809120in}{1.674185in}}%
\pgfpathlineto{\pgfqpoint{2.809244in}{1.532026in}}%
\pgfpathlineto{\pgfqpoint{2.809862in}{1.532040in}}%
\pgfpathlineto{\pgfqpoint{2.811116in}{1.532132in}}%
\pgfpathlineto{\pgfqpoint{2.811350in}{1.944389in}}%
\pgfpathlineto{\pgfqpoint{2.812280in}{1.915411in}}%
\pgfpathlineto{\pgfqpoint{2.812329in}{1.917609in}}%
\pgfpathlineto{\pgfqpoint{2.813022in}{1.916740in}}%
\pgfpathlineto{\pgfqpoint{2.813121in}{1.916861in}}%
\pgfpathlineto{\pgfqpoint{2.813714in}{1.900534in}}%
\pgfpathlineto{\pgfqpoint{2.814166in}{1.532028in}}%
\pgfpathlineto{\pgfqpoint{2.815206in}{1.532055in}}%
\pgfpathlineto{\pgfqpoint{2.816058in}{1.532104in}}%
\pgfpathlineto{\pgfqpoint{2.816100in}{1.722009in}}%
\pgfpathlineto{\pgfqpoint{2.816149in}{1.933707in}}%
\pgfpathlineto{\pgfqpoint{2.816851in}{1.910532in}}%
\pgfpathlineto{\pgfqpoint{2.817148in}{1.922586in}}%
\pgfpathlineto{\pgfqpoint{2.817593in}{1.920398in}}%
\pgfpathlineto{\pgfqpoint{2.817642in}{1.913023in}}%
\pgfpathlineto{\pgfqpoint{2.818384in}{1.917349in}}%
\pgfpathlineto{\pgfqpoint{2.818483in}{1.917379in}}%
\pgfpathlineto{\pgfqpoint{2.818730in}{1.867456in}}%
\pgfpathlineto{\pgfqpoint{2.819096in}{1.532025in}}%
\pgfpathlineto{\pgfqpoint{2.819976in}{1.532043in}}%
\pgfpathlineto{\pgfqpoint{2.821008in}{1.532160in}}%
\pgfpathlineto{\pgfqpoint{2.821233in}{1.937964in}}%
\pgfpathlineto{\pgfqpoint{2.822091in}{1.913392in}}%
\pgfpathlineto{\pgfqpoint{2.822141in}{1.920054in}}%
\pgfpathlineto{\pgfqpoint{2.822190in}{1.913190in}}%
\pgfpathlineto{\pgfqpoint{2.822833in}{1.916596in}}%
\pgfpathlineto{\pgfqpoint{2.823031in}{1.917797in}}%
\pgfpathlineto{\pgfqpoint{2.823575in}{1.908350in}}%
\pgfpathlineto{\pgfqpoint{2.824070in}{1.532028in}}%
\pgfpathlineto{\pgfqpoint{2.825414in}{1.532062in}}%
\pgfpathlineto{\pgfqpoint{2.825950in}{1.532113in}}%
\pgfpathlineto{\pgfqpoint{2.826178in}{1.940984in}}%
\pgfpathlineto{\pgfqpoint{2.827303in}{1.916426in}}%
\pgfpathlineto{\pgfqpoint{2.827352in}{1.916739in}}%
\pgfpathlineto{\pgfqpoint{2.828045in}{1.916356in}}%
\pgfpathlineto{\pgfqpoint{2.828490in}{1.914996in}}%
\pgfpathlineto{\pgfqpoint{2.828836in}{1.706808in}}%
\pgfpathlineto{\pgfqpoint{2.829028in}{1.532027in}}%
\pgfpathlineto{\pgfqpoint{2.829845in}{1.532045in}}%
\pgfpathlineto{\pgfqpoint{2.830898in}{1.532134in}}%
\pgfpathlineto{\pgfqpoint{2.831132in}{1.941875in}}%
\pgfpathlineto{\pgfqpoint{2.832066in}{1.915322in}}%
\pgfpathlineto{\pgfqpoint{2.832115in}{1.917631in}}%
\pgfpathlineto{\pgfqpoint{2.832808in}{1.916766in}}%
\pgfpathlineto{\pgfqpoint{2.832906in}{1.916876in}}%
\pgfpathlineto{\pgfqpoint{2.833500in}{1.898397in}}%
\pgfpathlineto{\pgfqpoint{2.833954in}{1.532028in}}%
\pgfpathlineto{\pgfqpoint{2.834947in}{1.532049in}}%
\pgfpathlineto{\pgfqpoint{2.835841in}{1.532111in}}%
\pgfpathlineto{\pgfqpoint{2.835943in}{1.935527in}}%
\pgfpathlineto{\pgfqpoint{2.837384in}{1.915618in}}%
\pgfpathlineto{\pgfqpoint{2.837433in}{1.917331in}}%
\pgfpathlineto{\pgfqpoint{2.838125in}{1.916500in}}%
\pgfpathlineto{\pgfqpoint{2.838422in}{1.907364in}}%
\pgfpathlineto{\pgfqpoint{2.838915in}{1.532026in}}%
\pgfpathlineto{\pgfqpoint{2.840174in}{1.532062in}}%
\pgfpathlineto{\pgfqpoint{2.840790in}{1.532149in}}%
\pgfpathlineto{\pgfqpoint{2.841017in}{1.942176in}}%
\pgfpathlineto{\pgfqpoint{2.841888in}{1.915264in}}%
\pgfpathlineto{\pgfqpoint{2.842184in}{1.918266in}}%
\pgfpathlineto{\pgfqpoint{2.842135in}{1.914854in}}%
\pgfpathlineto{\pgfqpoint{2.842629in}{1.917540in}}%
\pgfpathlineto{\pgfqpoint{2.843470in}{1.855051in}}%
\pgfpathlineto{\pgfqpoint{2.843870in}{1.532026in}}%
\pgfpathlineto{\pgfqpoint{2.844692in}{1.532045in}}%
\pgfpathlineto{\pgfqpoint{2.845735in}{1.532151in}}%
\pgfpathlineto{\pgfqpoint{2.845964in}{1.940597in}}%
\pgfpathlineto{\pgfqpoint{2.846855in}{1.917283in}}%
\pgfpathlineto{\pgfqpoint{2.847152in}{1.915395in}}%
\pgfpathlineto{\pgfqpoint{2.847103in}{1.917709in}}%
\pgfpathlineto{\pgfqpoint{2.847597in}{1.915751in}}%
\pgfpathlineto{\pgfqpoint{2.847647in}{1.917043in}}%
\pgfpathlineto{\pgfqpoint{2.848240in}{1.915501in}}%
\pgfpathlineto{\pgfqpoint{2.848438in}{1.841672in}}%
\pgfpathlineto{\pgfqpoint{2.848794in}{1.532027in}}%
\pgfpathlineto{\pgfqpoint{2.849593in}{1.532045in}}%
\pgfpathlineto{\pgfqpoint{2.850679in}{1.532123in}}%
\pgfpathlineto{\pgfqpoint{2.850901in}{1.938185in}}%
\pgfpathlineto{\pgfqpoint{2.851933in}{1.916624in}}%
\pgfpathlineto{\pgfqpoint{2.852032in}{1.916778in}}%
\pgfpathlineto{\pgfqpoint{2.852774in}{1.916318in}}%
\pgfpathlineto{\pgfqpoint{2.853120in}{1.915841in}}%
\pgfpathlineto{\pgfqpoint{2.853219in}{1.914689in}}%
\pgfpathlineto{\pgfqpoint{2.853594in}{1.685347in}}%
\pgfpathlineto{\pgfqpoint{2.853713in}{1.532025in}}%
\pgfpathlineto{\pgfqpoint{2.854589in}{1.532043in}}%
\pgfpathlineto{\pgfqpoint{2.855625in}{1.532134in}}%
\pgfpathlineto{\pgfqpoint{2.855961in}{1.939642in}}%
\pgfpathlineto{\pgfqpoint{2.856764in}{1.914253in}}%
\pgfpathlineto{\pgfqpoint{2.856813in}{1.919072in}}%
\pgfpathlineto{\pgfqpoint{2.856863in}{1.914024in}}%
\pgfpathlineto{\pgfqpoint{2.857506in}{1.916452in}}%
\pgfpathlineto{\pgfqpoint{2.857704in}{1.917345in}}%
\pgfpathlineto{\pgfqpoint{2.858198in}{1.908691in}}%
\pgfpathlineto{\pgfqpoint{2.858699in}{1.532024in}}%
\pgfpathlineto{\pgfqpoint{2.859959in}{1.532064in}}%
\pgfpathlineto{\pgfqpoint{2.860570in}{1.532123in}}%
\pgfpathlineto{\pgfqpoint{2.860905in}{1.941275in}}%
\pgfpathlineto{\pgfqpoint{2.861795in}{1.916032in}}%
\pgfpathlineto{\pgfqpoint{2.862092in}{1.917019in}}%
\pgfpathlineto{\pgfqpoint{2.862537in}{1.916702in}}%
\pgfpathlineto{\pgfqpoint{2.863131in}{1.911305in}}%
\pgfpathlineto{\pgfqpoint{2.863533in}{1.543063in}}%
\pgfpathlineto{\pgfqpoint{2.863624in}{1.532030in}}%
\pgfpathlineto{\pgfqpoint{2.864269in}{1.532044in}}%
\pgfpathlineto{\pgfqpoint{2.865520in}{1.532282in}}%
\pgfpathlineto{\pgfqpoint{2.865623in}{1.936563in}}%
\pgfpathlineto{\pgfqpoint{2.866546in}{1.918003in}}%
\pgfpathlineto{\pgfqpoint{2.866595in}{1.915083in}}%
\pgfpathlineto{\pgfqpoint{2.867287in}{1.916122in}}%
\pgfpathlineto{\pgfqpoint{2.867436in}{1.917016in}}%
\pgfpathlineto{\pgfqpoint{2.867930in}{1.915558in}}%
\pgfpathlineto{\pgfqpoint{2.867980in}{1.915915in}}%
\pgfpathlineto{\pgfqpoint{2.868079in}{1.910184in}}%
\pgfpathlineto{\pgfqpoint{2.868473in}{1.560932in}}%
\pgfpathlineto{\pgfqpoint{2.868556in}{1.532029in}}%
\pgfpathlineto{\pgfqpoint{2.869241in}{1.532043in}}%
\pgfpathlineto{\pgfqpoint{2.870458in}{1.532102in}}%
\pgfpathlineto{\pgfqpoint{2.870499in}{1.709924in}}%
\pgfpathlineto{\pgfqpoint{2.870582in}{1.937530in}}%
\pgfpathlineto{\pgfqpoint{2.871276in}{1.915358in}}%
\pgfpathlineto{\pgfqpoint{2.871424in}{1.921504in}}%
\pgfpathlineto{\pgfqpoint{2.871473in}{1.911387in}}%
\pgfpathlineto{\pgfqpoint{2.872067in}{1.917614in}}%
\pgfpathlineto{\pgfqpoint{2.872364in}{1.914482in}}%
\pgfpathlineto{\pgfqpoint{2.872314in}{1.918482in}}%
\pgfpathlineto{\pgfqpoint{2.872908in}{1.915141in}}%
\pgfpathlineto{\pgfqpoint{2.872957in}{1.915775in}}%
\pgfpathlineto{\pgfqpoint{2.873007in}{1.915154in}}%
\pgfpathlineto{\pgfqpoint{2.873415in}{1.620410in}}%
\pgfpathlineto{\pgfqpoint{2.873493in}{1.532026in}}%
\pgfpathlineto{\pgfqpoint{2.874169in}{1.532039in}}%
\pgfpathlineto{\pgfqpoint{2.875407in}{1.532136in}}%
\pgfpathlineto{\pgfqpoint{2.875640in}{1.944679in}}%
\pgfpathlineto{\pgfqpoint{2.876532in}{1.917473in}}%
\pgfpathlineto{\pgfqpoint{2.876581in}{1.915316in}}%
\pgfpathlineto{\pgfqpoint{2.876631in}{1.917647in}}%
\pgfpathlineto{\pgfqpoint{2.877323in}{1.916558in}}%
\pgfpathlineto{\pgfqpoint{2.877422in}{1.916759in}}%
\pgfpathlineto{\pgfqpoint{2.877966in}{1.911032in}}%
\pgfpathlineto{\pgfqpoint{2.878367in}{1.548861in}}%
\pgfpathlineto{\pgfqpoint{2.878481in}{1.532027in}}%
\pgfpathlineto{\pgfqpoint{2.879100in}{1.532041in}}%
\pgfpathlineto{\pgfqpoint{2.880351in}{1.532116in}}%
\pgfpathlineto{\pgfqpoint{2.880447in}{1.935034in}}%
\pgfpathlineto{\pgfqpoint{2.881701in}{1.916360in}}%
\pgfpathlineto{\pgfqpoint{2.881751in}{1.916663in}}%
\pgfpathlineto{\pgfqpoint{2.882443in}{1.916346in}}%
\pgfpathlineto{\pgfqpoint{2.882888in}{1.915047in}}%
\pgfpathlineto{\pgfqpoint{2.883185in}{1.744747in}}%
\pgfpathlineto{\pgfqpoint{2.883431in}{1.532026in}}%
\pgfpathlineto{\pgfqpoint{2.884201in}{1.532051in}}%
\pgfpathlineto{\pgfqpoint{2.885299in}{1.532150in}}%
\pgfpathlineto{\pgfqpoint{2.885525in}{1.942047in}}%
\pgfpathlineto{\pgfqpoint{2.886414in}{1.913919in}}%
\pgfpathlineto{\pgfqpoint{2.886463in}{1.918722in}}%
\pgfpathlineto{\pgfqpoint{2.887156in}{1.917122in}}%
\pgfpathlineto{\pgfqpoint{2.887255in}{1.917388in}}%
\pgfpathlineto{\pgfqpoint{2.887947in}{1.876476in}}%
\pgfpathlineto{\pgfqpoint{2.888356in}{1.532028in}}%
\pgfpathlineto{\pgfqpoint{2.889252in}{1.532047in}}%
\pgfpathlineto{\pgfqpoint{2.890242in}{1.532115in}}%
\pgfpathlineto{\pgfqpoint{2.890455in}{1.938242in}}%
\pgfpathlineto{\pgfqpoint{2.891539in}{1.916566in}}%
\pgfpathlineto{\pgfqpoint{2.892776in}{1.915050in}}%
\pgfpathlineto{\pgfqpoint{2.893072in}{1.745556in}}%
\pgfpathlineto{\pgfqpoint{2.893306in}{1.532027in}}%
\pgfpathlineto{\pgfqpoint{2.894110in}{1.532044in}}%
\pgfpathlineto{\pgfqpoint{2.895187in}{1.532113in}}%
\pgfpathlineto{\pgfqpoint{2.895307in}{1.936880in}}%
\pgfpathlineto{\pgfqpoint{2.896736in}{1.915797in}}%
\pgfpathlineto{\pgfqpoint{2.896786in}{1.917151in}}%
\pgfpathlineto{\pgfqpoint{2.897478in}{1.916242in}}%
\pgfpathlineto{\pgfqpoint{2.897726in}{1.914913in}}%
\pgfpathlineto{\pgfqpoint{2.898022in}{1.746065in}}%
\pgfpathlineto{\pgfqpoint{2.898222in}{1.532024in}}%
\pgfpathlineto{\pgfqpoint{2.899048in}{1.532053in}}%
\pgfpathlineto{\pgfqpoint{2.900133in}{1.532123in}}%
\pgfpathlineto{\pgfqpoint{2.900370in}{1.942662in}}%
\pgfpathlineto{\pgfqpoint{2.901365in}{1.915944in}}%
\pgfpathlineto{\pgfqpoint{2.901662in}{1.916933in}}%
\pgfpathlineto{\pgfqpoint{2.902107in}{1.916630in}}%
\pgfpathlineto{\pgfqpoint{2.902701in}{1.909671in}}%
\pgfpathlineto{\pgfqpoint{2.903215in}{1.532030in}}%
\pgfpathlineto{\pgfqpoint{2.904531in}{1.532059in}}%
\pgfpathlineto{\pgfqpoint{2.905078in}{1.532115in}}%
\pgfpathlineto{\pgfqpoint{2.905180in}{1.935344in}}%
\pgfpathlineto{\pgfqpoint{2.906579in}{1.919021in}}%
\pgfpathlineto{\pgfqpoint{2.906628in}{1.914447in}}%
\pgfpathlineto{\pgfqpoint{2.907370in}{1.917114in}}%
\pgfpathlineto{\pgfqpoint{2.907716in}{1.884838in}}%
\pgfpathlineto{\pgfqpoint{2.908114in}{1.532024in}}%
\pgfpathlineto{\pgfqpoint{2.909039in}{1.532056in}}%
\pgfpathlineto{\pgfqpoint{2.910024in}{1.532118in}}%
\pgfpathlineto{\pgfqpoint{2.910258in}{1.942795in}}%
\pgfpathlineto{\pgfqpoint{2.911317in}{1.916562in}}%
\pgfpathlineto{\pgfqpoint{2.912009in}{1.916138in}}%
\pgfpathlineto{\pgfqpoint{2.911515in}{1.916834in}}%
\pgfpathlineto{\pgfqpoint{2.912108in}{1.916191in}}%
\pgfpathlineto{\pgfqpoint{2.912306in}{1.916191in}}%
\pgfpathlineto{\pgfqpoint{2.912454in}{1.915648in}}%
\pgfpathlineto{\pgfqpoint{2.912504in}{1.915632in}}%
\pgfpathlineto{\pgfqpoint{2.912603in}{1.906868in}}%
\pgfpathlineto{\pgfqpoint{2.913112in}{1.532030in}}%
\pgfpathlineto{\pgfqpoint{2.914284in}{1.532057in}}%
\pgfpathlineto{\pgfqpoint{2.914972in}{1.532160in}}%
\pgfpathlineto{\pgfqpoint{2.915069in}{1.935153in}}%
\pgfpathlineto{\pgfqpoint{2.916084in}{1.916616in}}%
\pgfpathlineto{\pgfqpoint{2.916133in}{1.916289in}}%
\pgfpathlineto{\pgfqpoint{2.916183in}{1.916634in}}%
\pgfpathlineto{\pgfqpoint{2.916875in}{1.916424in}}%
\pgfpathlineto{\pgfqpoint{2.917518in}{1.914136in}}%
\pgfpathlineto{\pgfqpoint{2.917923in}{1.642420in}}%
\pgfpathlineto{\pgfqpoint{2.918016in}{1.532024in}}%
\pgfpathlineto{\pgfqpoint{2.918657in}{1.532049in}}%
\pgfpathlineto{\pgfqpoint{2.919914in}{1.532112in}}%
\pgfpathlineto{\pgfqpoint{2.920034in}{1.937412in}}%
\pgfpathlineto{\pgfqpoint{2.921456in}{1.914823in}}%
\pgfpathlineto{\pgfqpoint{2.921505in}{1.918189in}}%
\pgfpathlineto{\pgfqpoint{2.922198in}{1.916246in}}%
\pgfpathlineto{\pgfqpoint{2.922297in}{1.916489in}}%
\pgfpathlineto{\pgfqpoint{2.922495in}{1.907319in}}%
\pgfpathlineto{\pgfqpoint{2.922999in}{1.532026in}}%
\pgfpathlineto{\pgfqpoint{2.924267in}{1.532063in}}%
\pgfpathlineto{\pgfqpoint{2.924862in}{1.532136in}}%
\pgfpathlineto{\pgfqpoint{2.925094in}{1.944690in}}%
\pgfpathlineto{\pgfqpoint{2.925987in}{1.917325in}}%
\pgfpathlineto{\pgfqpoint{2.926037in}{1.915324in}}%
\pgfpathlineto{\pgfqpoint{2.926086in}{1.917599in}}%
\pgfpathlineto{\pgfqpoint{2.926779in}{1.916449in}}%
\pgfpathlineto{\pgfqpoint{2.926878in}{1.916681in}}%
\pgfpathlineto{\pgfqpoint{2.927422in}{1.911164in}}%
\pgfpathlineto{\pgfqpoint{2.927821in}{1.549055in}}%
\pgfpathlineto{\pgfqpoint{2.927898in}{1.532030in}}%
\pgfpathlineto{\pgfqpoint{2.928580in}{1.532044in}}%
\pgfpathlineto{\pgfqpoint{2.929811in}{1.532279in}}%
\pgfpathlineto{\pgfqpoint{2.929900in}{1.934851in}}%
\pgfpathlineto{\pgfqpoint{2.930853in}{1.916772in}}%
\pgfpathlineto{\pgfqpoint{2.931002in}{1.915819in}}%
\pgfpathlineto{\pgfqpoint{2.930952in}{1.917068in}}%
\pgfpathlineto{\pgfqpoint{2.931645in}{1.916370in}}%
\pgfpathlineto{\pgfqpoint{2.931743in}{1.916522in}}%
\pgfpathlineto{\pgfqpoint{2.932337in}{1.914973in}}%
\pgfpathlineto{\pgfqpoint{2.932584in}{1.787824in}}%
\pgfpathlineto{\pgfqpoint{2.932841in}{1.532024in}}%
\pgfpathlineto{\pgfqpoint{2.933667in}{1.532053in}}%
\pgfpathlineto{\pgfqpoint{2.934751in}{1.532121in}}%
\pgfpathlineto{\pgfqpoint{2.934987in}{1.943043in}}%
\pgfpathlineto{\pgfqpoint{2.935974in}{1.915956in}}%
\pgfpathlineto{\pgfqpoint{2.936271in}{1.916819in}}%
\pgfpathlineto{\pgfqpoint{2.936716in}{1.916552in}}%
\pgfpathlineto{\pgfqpoint{2.937309in}{1.912019in}}%
\pgfpathlineto{\pgfqpoint{2.937708in}{1.575003in}}%
\pgfpathlineto{\pgfqpoint{2.937816in}{1.532028in}}%
\pgfpathlineto{\pgfqpoint{2.938469in}{1.532047in}}%
\pgfpathlineto{\pgfqpoint{2.939696in}{1.532118in}}%
\pgfpathlineto{\pgfqpoint{2.939786in}{1.933315in}}%
\pgfpathlineto{\pgfqpoint{2.941002in}{1.916509in}}%
\pgfpathlineto{\pgfqpoint{2.942238in}{1.914714in}}%
\pgfpathlineto{\pgfqpoint{2.942613in}{1.686025in}}%
\pgfpathlineto{\pgfqpoint{2.942731in}{1.532024in}}%
\pgfpathlineto{\pgfqpoint{2.943607in}{1.532043in}}%
\pgfpathlineto{\pgfqpoint{2.944643in}{1.532127in}}%
\pgfpathlineto{\pgfqpoint{2.944879in}{1.943078in}}%
\pgfpathlineto{\pgfqpoint{2.945844in}{1.917270in}}%
\pgfpathlineto{\pgfqpoint{2.945893in}{1.915716in}}%
\pgfpathlineto{\pgfqpoint{2.946635in}{1.916630in}}%
\pgfpathlineto{\pgfqpoint{2.947228in}{1.905927in}}%
\pgfpathlineto{\pgfqpoint{2.947715in}{1.532026in}}%
\pgfpathlineto{\pgfqpoint{2.948827in}{1.532051in}}%
\pgfpathlineto{\pgfqpoint{2.949589in}{1.532131in}}%
\pgfpathlineto{\pgfqpoint{2.949822in}{1.940638in}}%
\pgfpathlineto{\pgfqpoint{2.950736in}{1.917005in}}%
\pgfpathlineto{\pgfqpoint{2.950884in}{1.915511in}}%
\pgfpathlineto{\pgfqpoint{2.950835in}{1.917353in}}%
\pgfpathlineto{\pgfqpoint{2.951527in}{1.916315in}}%
\pgfpathlineto{\pgfqpoint{2.951626in}{1.916526in}}%
\pgfpathlineto{\pgfqpoint{2.952121in}{1.914733in}}%
\pgfpathlineto{\pgfqpoint{2.952417in}{1.746163in}}%
\pgfpathlineto{\pgfqpoint{2.952653in}{1.532024in}}%
\pgfpathlineto{\pgfqpoint{2.953459in}{1.532041in}}%
\pgfpathlineto{\pgfqpoint{2.954532in}{1.532113in}}%
\pgfpathlineto{\pgfqpoint{2.954754in}{1.938980in}}%
\pgfpathlineto{\pgfqpoint{2.955862in}{1.916391in}}%
\pgfpathlineto{\pgfqpoint{2.956604in}{1.916254in}}%
\pgfpathlineto{\pgfqpoint{2.956653in}{1.916165in}}%
\pgfpathlineto{\pgfqpoint{2.956999in}{1.915575in}}%
\pgfpathlineto{\pgfqpoint{2.957098in}{1.910593in}}%
\pgfpathlineto{\pgfqpoint{2.957499in}{1.537362in}}%
\pgfpathlineto{\pgfqpoint{2.957608in}{1.532027in}}%
\pgfpathlineto{\pgfqpoint{2.958227in}{1.532041in}}%
\pgfpathlineto{\pgfqpoint{2.959479in}{1.532129in}}%
\pgfpathlineto{\pgfqpoint{2.959815in}{1.943240in}}%
\pgfpathlineto{\pgfqpoint{2.960677in}{1.916850in}}%
\pgfpathlineto{\pgfqpoint{2.960974in}{1.915627in}}%
\pgfpathlineto{\pgfqpoint{2.960924in}{1.917221in}}%
\pgfpathlineto{\pgfqpoint{2.961518in}{1.915958in}}%
\pgfpathlineto{\pgfqpoint{2.961567in}{1.916322in}}%
\pgfpathlineto{\pgfqpoint{2.962012in}{1.914916in}}%
\pgfpathlineto{\pgfqpoint{2.962309in}{1.745930in}}%
\pgfpathlineto{\pgfqpoint{2.962516in}{1.532025in}}%
\pgfpathlineto{\pgfqpoint{2.963346in}{1.532053in}}%
\pgfpathlineto{\pgfqpoint{2.964425in}{1.532136in}}%
\pgfpathlineto{\pgfqpoint{2.964657in}{1.942024in}}%
\pgfpathlineto{\pgfqpoint{2.965547in}{1.917626in}}%
\pgfpathlineto{\pgfqpoint{2.965596in}{1.915380in}}%
\pgfpathlineto{\pgfqpoint{2.966288in}{1.915845in}}%
\pgfpathlineto{\pgfqpoint{2.966338in}{1.916833in}}%
\pgfpathlineto{\pgfqpoint{2.966931in}{1.915140in}}%
\pgfpathlineto{\pgfqpoint{2.967080in}{1.873095in}}%
\pgfpathlineto{\pgfqpoint{2.967497in}{1.532026in}}%
\pgfpathlineto{\pgfqpoint{2.968361in}{1.532053in}}%
\pgfpathlineto{\pgfqpoint{2.969371in}{1.532152in}}%
\pgfpathlineto{\pgfqpoint{2.969600in}{1.940068in}}%
\pgfpathlineto{\pgfqpoint{2.970491in}{1.917359in}}%
\pgfpathlineto{\pgfqpoint{2.970788in}{1.915361in}}%
\pgfpathlineto{\pgfqpoint{2.970837in}{1.917390in}}%
\pgfpathlineto{\pgfqpoint{2.971233in}{1.915663in}}%
\pgfpathlineto{\pgfqpoint{2.971282in}{1.916899in}}%
\pgfpathlineto{\pgfqpoint{2.971924in}{1.913752in}}%
\pgfpathlineto{\pgfqpoint{2.972241in}{1.715468in}}%
\pgfpathlineto{\pgfqpoint{2.972445in}{1.532024in}}%
\pgfpathlineto{\pgfqpoint{2.973262in}{1.532042in}}%
\pgfpathlineto{\pgfqpoint{2.974314in}{1.532113in}}%
\pgfpathlineto{\pgfqpoint{2.974534in}{1.938892in}}%
\pgfpathlineto{\pgfqpoint{2.975639in}{1.916395in}}%
\pgfpathlineto{\pgfqpoint{2.976430in}{1.916186in}}%
\pgfpathlineto{\pgfqpoint{2.976479in}{1.916110in}}%
\pgfpathlineto{\pgfqpoint{2.976826in}{1.915262in}}%
\pgfpathlineto{\pgfqpoint{2.976974in}{1.870989in}}%
\pgfpathlineto{\pgfqpoint{2.977385in}{1.532026in}}%
\pgfpathlineto{\pgfqpoint{2.978248in}{1.532053in}}%
\pgfpathlineto{\pgfqpoint{2.979259in}{1.532112in}}%
\pgfpathlineto{\pgfqpoint{2.979469in}{1.935166in}}%
\pgfpathlineto{\pgfqpoint{2.980819in}{1.915767in}}%
\pgfpathlineto{\pgfqpoint{2.980968in}{1.917363in}}%
\pgfpathlineto{\pgfqpoint{2.980918in}{1.915383in}}%
\pgfpathlineto{\pgfqpoint{2.981561in}{1.915791in}}%
\pgfpathlineto{\pgfqpoint{2.981660in}{1.915956in}}%
\pgfpathlineto{\pgfqpoint{2.981709in}{1.915360in}}%
\pgfpathlineto{\pgfqpoint{2.981759in}{1.915795in}}%
\pgfpathlineto{\pgfqpoint{2.981858in}{1.904125in}}%
\pgfpathlineto{\pgfqpoint{2.982321in}{1.532030in}}%
\pgfpathlineto{\pgfqpoint{2.983415in}{1.532053in}}%
\pgfpathlineto{\pgfqpoint{2.984206in}{1.532122in}}%
\pgfpathlineto{\pgfqpoint{2.984425in}{1.940283in}}%
\pgfpathlineto{\pgfqpoint{2.985450in}{1.916398in}}%
\pgfpathlineto{\pgfqpoint{2.985549in}{1.916514in}}%
\pgfpathlineto{\pgfqpoint{2.986736in}{1.914973in}}%
\pgfpathlineto{\pgfqpoint{2.987033in}{1.747702in}}%
\pgfpathlineto{\pgfqpoint{2.987279in}{1.532030in}}%
\pgfpathlineto{\pgfqpoint{2.988094in}{1.532048in}}%
\pgfpathlineto{\pgfqpoint{2.989156in}{1.532279in}}%
\pgfpathlineto{\pgfqpoint{2.989259in}{1.936198in}}%
\pgfpathlineto{\pgfqpoint{2.990170in}{1.916195in}}%
\pgfpathlineto{\pgfqpoint{2.990319in}{1.918901in}}%
\pgfpathlineto{\pgfqpoint{2.990368in}{1.913716in}}%
\pgfpathlineto{\pgfqpoint{2.990863in}{1.917868in}}%
\pgfpathlineto{\pgfqpoint{2.991852in}{1.846552in}}%
\pgfpathlineto{\pgfqpoint{2.992213in}{1.532027in}}%
\pgfpathlineto{\pgfqpoint{2.993013in}{1.532044in}}%
\pgfpathlineto{\pgfqpoint{2.994096in}{1.532113in}}%
\pgfpathlineto{\pgfqpoint{2.994213in}{1.936888in}}%
\pgfpathlineto{\pgfqpoint{2.995620in}{1.915788in}}%
\pgfpathlineto{\pgfqpoint{2.995670in}{1.917155in}}%
\pgfpathlineto{\pgfqpoint{2.995719in}{1.915569in}}%
\pgfpathlineto{\pgfqpoint{2.996362in}{1.915925in}}%
\pgfpathlineto{\pgfqpoint{2.996609in}{1.914994in}}%
\pgfpathlineto{\pgfqpoint{2.996758in}{1.872349in}}%
\pgfpathlineto{\pgfqpoint{2.997132in}{1.532025in}}%
\pgfpathlineto{\pgfqpoint{2.998012in}{1.532043in}}%
\pgfpathlineto{\pgfqpoint{2.999045in}{1.532160in}}%
\pgfpathlineto{\pgfqpoint{2.999262in}{1.938015in}}%
\pgfpathlineto{\pgfqpoint{3.000112in}{1.920396in}}%
\pgfpathlineto{\pgfqpoint{3.000409in}{1.913347in}}%
\pgfpathlineto{\pgfqpoint{3.000854in}{1.914374in}}%
\pgfpathlineto{\pgfqpoint{3.000903in}{1.918329in}}%
\pgfpathlineto{\pgfqpoint{3.001596in}{1.914783in}}%
\pgfpathlineto{\pgfqpoint{3.002001in}{1.561703in}}%
\pgfpathlineto{\pgfqpoint{3.002108in}{1.532027in}}%
\pgfpathlineto{\pgfqpoint{3.002761in}{1.532047in}}%
\pgfpathlineto{\pgfqpoint{3.003988in}{1.532120in}}%
\pgfpathlineto{\pgfqpoint{3.004211in}{1.940148in}}%
\pgfpathlineto{\pgfqpoint{3.005287in}{1.916450in}}%
\pgfpathlineto{\pgfqpoint{3.006523in}{1.914901in}}%
\pgfpathlineto{\pgfqpoint{3.006820in}{1.744740in}}%
\pgfpathlineto{\pgfqpoint{3.007021in}{1.532025in}}%
\pgfpathlineto{\pgfqpoint{3.007852in}{1.532053in}}%
\pgfpathlineto{\pgfqpoint{3.008934in}{1.532136in}}%
\pgfpathlineto{\pgfqpoint{3.009266in}{1.942766in}}%
\pgfpathlineto{\pgfqpoint{3.010085in}{1.915131in}}%
\pgfpathlineto{\pgfqpoint{3.010134in}{1.917299in}}%
\pgfpathlineto{\pgfqpoint{3.010827in}{1.916721in}}%
\pgfpathlineto{\pgfqpoint{3.010925in}{1.916727in}}%
\pgfpathlineto{\pgfqpoint{3.011519in}{1.906363in}}%
\pgfpathlineto{\pgfqpoint{3.011968in}{1.532025in}}%
\pgfpathlineto{\pgfqpoint{3.013144in}{1.532050in}}%
\pgfpathlineto{\pgfqpoint{3.013881in}{1.532161in}}%
\pgfpathlineto{\pgfqpoint{3.014103in}{1.935353in}}%
\pgfpathlineto{\pgfqpoint{3.014956in}{1.921381in}}%
\pgfpathlineto{\pgfqpoint{3.015006in}{1.912093in}}%
\pgfpathlineto{\pgfqpoint{3.015698in}{1.914711in}}%
\pgfpathlineto{\pgfqpoint{3.015747in}{1.918273in}}%
\pgfpathlineto{\pgfqpoint{3.016440in}{1.913930in}}%
\pgfpathlineto{\pgfqpoint{3.016954in}{1.532030in}}%
\pgfpathlineto{\pgfqpoint{3.018363in}{1.532061in}}%
\pgfpathlineto{\pgfqpoint{3.018829in}{1.532278in}}%
\pgfpathlineto{\pgfqpoint{3.018933in}{1.935194in}}%
\pgfpathlineto{\pgfqpoint{3.019861in}{1.913037in}}%
\pgfpathlineto{\pgfqpoint{3.020158in}{1.919351in}}%
\pgfpathlineto{\pgfqpoint{3.020603in}{1.918210in}}%
\pgfpathlineto{\pgfqpoint{3.021592in}{1.795475in}}%
\pgfpathlineto{\pgfqpoint{3.021905in}{1.532026in}}%
\pgfpathlineto{\pgfqpoint{3.022676in}{1.532051in}}%
\pgfpathlineto{\pgfqpoint{3.023770in}{1.532133in}}%
\pgfpathlineto{\pgfqpoint{3.024004in}{1.944455in}}%
\pgfpathlineto{\pgfqpoint{3.024901in}{1.917263in}}%
\pgfpathlineto{\pgfqpoint{3.024950in}{1.915210in}}%
\pgfpathlineto{\pgfqpoint{3.024999in}{1.917467in}}%
\pgfpathlineto{\pgfqpoint{3.025692in}{1.916379in}}%
\pgfpathlineto{\pgfqpoint{3.025791in}{1.916585in}}%
\pgfpathlineto{\pgfqpoint{3.026335in}{1.910404in}}%
\pgfpathlineto{\pgfqpoint{3.026825in}{1.532028in}}%
\pgfpathlineto{\pgfqpoint{3.028263in}{1.532064in}}%
\pgfpathlineto{\pgfqpoint{3.028714in}{1.532109in}}%
\pgfpathlineto{\pgfqpoint{3.028756in}{1.750475in}}%
\pgfpathlineto{\pgfqpoint{3.028815in}{1.934934in}}%
\pgfpathlineto{\pgfqpoint{3.029534in}{1.917663in}}%
\pgfpathlineto{\pgfqpoint{3.029683in}{1.912402in}}%
\pgfpathlineto{\pgfqpoint{3.029732in}{1.920019in}}%
\pgfpathlineto{\pgfqpoint{3.030326in}{1.915824in}}%
\pgfpathlineto{\pgfqpoint{3.030523in}{1.917794in}}%
\pgfpathlineto{\pgfqpoint{3.030573in}{1.914759in}}%
\pgfpathlineto{\pgfqpoint{3.030968in}{1.916945in}}%
\pgfpathlineto{\pgfqpoint{3.031364in}{1.879990in}}%
\pgfpathlineto{\pgfqpoint{3.031750in}{1.532024in}}%
\pgfpathlineto{\pgfqpoint{3.032675in}{1.532056in}}%
\pgfpathlineto{\pgfqpoint{3.033661in}{1.532123in}}%
\pgfpathlineto{\pgfqpoint{3.033897in}{1.943366in}}%
\pgfpathlineto{\pgfqpoint{3.034902in}{1.915732in}}%
\pgfpathlineto{\pgfqpoint{3.034952in}{1.916899in}}%
\pgfpathlineto{\pgfqpoint{3.035644in}{1.916333in}}%
\pgfpathlineto{\pgfqpoint{3.035743in}{1.916365in}}%
\pgfpathlineto{\pgfqpoint{3.036237in}{1.907356in}}%
\pgfpathlineto{\pgfqpoint{3.036736in}{1.532025in}}%
\pgfpathlineto{\pgfqpoint{3.037948in}{1.532052in}}%
\pgfpathlineto{\pgfqpoint{3.038607in}{1.532139in}}%
\pgfpathlineto{\pgfqpoint{3.038811in}{1.936687in}}%
\pgfpathlineto{\pgfqpoint{3.039740in}{1.914371in}}%
\pgfpathlineto{\pgfqpoint{3.039790in}{1.918526in}}%
\pgfpathlineto{\pgfqpoint{3.039839in}{1.914222in}}%
\pgfpathlineto{\pgfqpoint{3.040482in}{1.916305in}}%
\pgfpathlineto{\pgfqpoint{3.040680in}{1.917024in}}%
\pgfpathlineto{\pgfqpoint{3.041175in}{1.910303in}}%
\pgfpathlineto{\pgfqpoint{3.041689in}{1.532026in}}%
\pgfpathlineto{\pgfqpoint{3.043055in}{1.532065in}}%
\pgfpathlineto{\pgfqpoint{3.043552in}{1.532136in}}%
\pgfpathlineto{\pgfqpoint{3.043785in}{1.944160in}}%
\pgfpathlineto{\pgfqpoint{3.044693in}{1.916736in}}%
\pgfpathlineto{\pgfqpoint{3.044841in}{1.915303in}}%
\pgfpathlineto{\pgfqpoint{3.044890in}{1.917291in}}%
\pgfpathlineto{\pgfqpoint{3.045484in}{1.916104in}}%
\pgfpathlineto{\pgfqpoint{3.045781in}{1.916341in}}%
\pgfpathlineto{\pgfqpoint{3.046028in}{1.915607in}}%
\pgfpathlineto{\pgfqpoint{3.046127in}{1.908494in}}%
\pgfpathlineto{\pgfqpoint{3.046628in}{1.532024in}}%
\pgfpathlineto{\pgfqpoint{3.047840in}{1.532051in}}%
\pgfpathlineto{\pgfqpoint{3.048495in}{1.532109in}}%
\pgfpathlineto{\pgfqpoint{3.048537in}{1.737609in}}%
\pgfpathlineto{\pgfqpoint{3.048598in}{1.934780in}}%
\pgfpathlineto{\pgfqpoint{3.049289in}{1.918205in}}%
\pgfpathlineto{\pgfqpoint{3.049487in}{1.910343in}}%
\pgfpathlineto{\pgfqpoint{3.049537in}{1.922419in}}%
\pgfpathlineto{\pgfqpoint{3.049932in}{1.912560in}}%
\pgfpathlineto{\pgfqpoint{3.049982in}{1.920175in}}%
\pgfpathlineto{\pgfqpoint{3.050674in}{1.916604in}}%
\pgfpathlineto{\pgfqpoint{3.050773in}{1.917340in}}%
\pgfpathlineto{\pgfqpoint{3.051119in}{1.892326in}}%
\pgfpathlineto{\pgfqpoint{3.051585in}{1.532025in}}%
\pgfpathlineto{\pgfqpoint{3.052509in}{1.532046in}}%
\pgfpathlineto{\pgfqpoint{3.053443in}{1.532136in}}%
\pgfpathlineto{\pgfqpoint{3.053676in}{1.944259in}}%
\pgfpathlineto{\pgfqpoint{3.054598in}{1.916176in}}%
\pgfpathlineto{\pgfqpoint{3.054845in}{1.915380in}}%
\pgfpathlineto{\pgfqpoint{3.054796in}{1.917312in}}%
\pgfpathlineto{\pgfqpoint{3.055192in}{1.915720in}}%
\pgfpathlineto{\pgfqpoint{3.055241in}{1.916893in}}%
\pgfpathlineto{\pgfqpoint{3.055933in}{1.915409in}}%
\pgfpathlineto{\pgfqpoint{3.056032in}{1.905209in}}%
\pgfpathlineto{\pgfqpoint{3.056477in}{1.532024in}}%
\pgfpathlineto{\pgfqpoint{3.057600in}{1.532060in}}%
\pgfpathlineto{\pgfqpoint{3.058388in}{1.532125in}}%
\pgfpathlineto{\pgfqpoint{3.058602in}{1.938433in}}%
\pgfpathlineto{\pgfqpoint{3.059604in}{1.916845in}}%
\pgfpathlineto{\pgfqpoint{3.059653in}{1.915751in}}%
\pgfpathlineto{\pgfqpoint{3.059703in}{1.916905in}}%
\pgfpathlineto{\pgfqpoint{3.060395in}{1.916283in}}%
\pgfpathlineto{\pgfqpoint{3.060494in}{1.916329in}}%
\pgfpathlineto{\pgfqpoint{3.060939in}{1.913555in}}%
\pgfpathlineto{\pgfqpoint{3.061343in}{1.615975in}}%
\pgfpathlineto{\pgfqpoint{3.061465in}{1.532027in}}%
\pgfpathlineto{\pgfqpoint{3.062085in}{1.532040in}}%
\pgfpathlineto{\pgfqpoint{3.063334in}{1.532138in}}%
\pgfpathlineto{\pgfqpoint{3.063567in}{1.944219in}}%
\pgfpathlineto{\pgfqpoint{3.064466in}{1.916921in}}%
\pgfpathlineto{\pgfqpoint{3.064614in}{1.915243in}}%
\pgfpathlineto{\pgfqpoint{3.064565in}{1.917367in}}%
\pgfpathlineto{\pgfqpoint{3.065257in}{1.916177in}}%
\pgfpathlineto{\pgfqpoint{3.065455in}{1.916499in}}%
\pgfpathlineto{\pgfqpoint{3.065850in}{1.914884in}}%
\pgfpathlineto{\pgfqpoint{3.066048in}{1.834845in}}%
\pgfpathlineto{\pgfqpoint{3.066394in}{1.532027in}}%
\pgfpathlineto{\pgfqpoint{3.067193in}{1.532045in}}%
\pgfpathlineto{\pgfqpoint{3.068279in}{1.532123in}}%
\pgfpathlineto{\pgfqpoint{3.068368in}{1.934024in}}%
\pgfpathlineto{\pgfqpoint{3.069562in}{1.916574in}}%
\pgfpathlineto{\pgfqpoint{3.070304in}{1.916119in}}%
\pgfpathlineto{\pgfqpoint{3.070353in}{1.916205in}}%
\pgfpathlineto{\pgfqpoint{3.070798in}{1.915192in}}%
\pgfpathlineto{\pgfqpoint{3.070996in}{1.833689in}}%
\pgfpathlineto{\pgfqpoint{3.071357in}{1.532027in}}%
\pgfpathlineto{\pgfqpoint{3.072125in}{1.532051in}}%
\pgfpathlineto{\pgfqpoint{3.073225in}{1.532141in}}%
\pgfpathlineto{\pgfqpoint{3.073457in}{1.943892in}}%
\pgfpathlineto{\pgfqpoint{3.074337in}{1.917139in}}%
\pgfpathlineto{\pgfqpoint{3.074634in}{1.915374in}}%
\pgfpathlineto{\pgfqpoint{3.074684in}{1.917192in}}%
\pgfpathlineto{\pgfqpoint{3.075079in}{1.915640in}}%
\pgfpathlineto{\pgfqpoint{3.075129in}{1.916748in}}%
\pgfpathlineto{\pgfqpoint{3.075722in}{1.915269in}}%
\pgfpathlineto{\pgfqpoint{3.075821in}{1.902999in}}%
\pgfpathlineto{\pgfqpoint{3.076297in}{1.532025in}}%
\pgfpathlineto{\pgfqpoint{3.077359in}{1.532059in}}%
\pgfpathlineto{\pgfqpoint{3.078159in}{1.532079in}}%
\pgfpathlineto{\pgfqpoint{3.078171in}{1.532146in}}%
\pgfpathlineto{\pgfqpoint{3.078372in}{1.936105in}}%
\pgfpathlineto{\pgfqpoint{3.079295in}{1.915259in}}%
\pgfpathlineto{\pgfqpoint{3.079591in}{1.918332in}}%
\pgfpathlineto{\pgfqpoint{3.079542in}{1.914207in}}%
\pgfpathlineto{\pgfqpoint{3.080036in}{1.917469in}}%
\pgfpathlineto{\pgfqpoint{3.080828in}{1.873259in}}%
\pgfpathlineto{\pgfqpoint{3.081230in}{1.532027in}}%
\pgfpathlineto{\pgfqpoint{3.082078in}{1.532051in}}%
\pgfpathlineto{\pgfqpoint{3.083115in}{1.532122in}}%
\pgfpathlineto{\pgfqpoint{3.083340in}{1.940882in}}%
\pgfpathlineto{\pgfqpoint{3.084383in}{1.916272in}}%
\pgfpathlineto{\pgfqpoint{3.084482in}{1.916460in}}%
\pgfpathlineto{\pgfqpoint{3.085620in}{1.915218in}}%
\pgfpathlineto{\pgfqpoint{3.085768in}{1.875457in}}%
\pgfpathlineto{\pgfqpoint{3.086176in}{1.532027in}}%
\pgfpathlineto{\pgfqpoint{3.087075in}{1.532047in}}%
\pgfpathlineto{\pgfqpoint{3.088060in}{1.532113in}}%
\pgfpathlineto{\pgfqpoint{3.088161in}{1.935173in}}%
\pgfpathlineto{\pgfqpoint{3.089532in}{1.918013in}}%
\pgfpathlineto{\pgfqpoint{3.089581in}{1.914726in}}%
\pgfpathlineto{\pgfqpoint{3.090323in}{1.916476in}}%
\pgfpathlineto{\pgfqpoint{3.090669in}{1.898963in}}%
\pgfpathlineto{\pgfqpoint{3.091123in}{1.532027in}}%
\pgfpathlineto{\pgfqpoint{3.092122in}{1.532049in}}%
\pgfpathlineto{\pgfqpoint{3.093005in}{1.532112in}}%
\pgfpathlineto{\pgfqpoint{3.093124in}{1.937105in}}%
\pgfpathlineto{\pgfqpoint{3.094556in}{1.917436in}}%
\pgfpathlineto{\pgfqpoint{3.094853in}{1.915102in}}%
\pgfpathlineto{\pgfqpoint{3.095347in}{1.916491in}}%
\pgfpathlineto{\pgfqpoint{3.095595in}{1.907274in}}%
\pgfpathlineto{\pgfqpoint{3.096092in}{1.532024in}}%
\pgfpathlineto{\pgfqpoint{3.097213in}{1.532049in}}%
\pgfpathlineto{\pgfqpoint{3.097951in}{1.532119in}}%
\pgfpathlineto{\pgfqpoint{3.098124in}{1.938797in}}%
\pgfpathlineto{\pgfqpoint{3.099218in}{1.916251in}}%
\pgfpathlineto{\pgfqpoint{3.099712in}{1.916314in}}%
\pgfpathlineto{\pgfqpoint{3.099960in}{1.916159in}}%
\pgfpathlineto{\pgfqpoint{3.100454in}{1.915239in}}%
\pgfpathlineto{\pgfqpoint{3.100602in}{1.874233in}}%
\pgfpathlineto{\pgfqpoint{3.101010in}{1.532028in}}%
\pgfpathlineto{\pgfqpoint{3.101905in}{1.532047in}}%
\pgfpathlineto{\pgfqpoint{3.102896in}{1.532113in}}%
\pgfpathlineto{\pgfqpoint{3.103123in}{1.940750in}}%
\pgfpathlineto{\pgfqpoint{3.104247in}{1.916147in}}%
\pgfpathlineto{\pgfqpoint{3.104297in}{1.916478in}}%
\pgfpathlineto{\pgfqpoint{3.104989in}{1.916103in}}%
\pgfpathlineto{\pgfqpoint{3.105434in}{1.914872in}}%
\pgfpathlineto{\pgfqpoint{3.105731in}{1.746223in}}%
\pgfpathlineto{\pgfqpoint{3.105931in}{1.532025in}}%
\pgfpathlineto{\pgfqpoint{3.106761in}{1.532053in}}%
\pgfpathlineto{\pgfqpoint{3.107843in}{1.532136in}}%
\pgfpathlineto{\pgfqpoint{3.108175in}{1.942315in}}%
\pgfpathlineto{\pgfqpoint{3.108962in}{1.917553in}}%
\pgfpathlineto{\pgfqpoint{3.109011in}{1.914941in}}%
\pgfpathlineto{\pgfqpoint{3.109704in}{1.916060in}}%
\pgfpathlineto{\pgfqpoint{3.109852in}{1.916655in}}%
\pgfpathlineto{\pgfqpoint{3.110347in}{1.914930in}}%
\pgfpathlineto{\pgfqpoint{3.110495in}{1.874767in}}%
\pgfpathlineto{\pgfqpoint{3.110903in}{1.532027in}}%
\pgfpathlineto{\pgfqpoint{3.111800in}{1.532047in}}%
\pgfpathlineto{\pgfqpoint{3.112788in}{1.532122in}}%
\pgfpathlineto{\pgfqpoint{3.112979in}{1.936188in}}%
\pgfpathlineto{\pgfqpoint{3.114038in}{1.916412in}}%
\pgfpathlineto{\pgfqpoint{3.114137in}{1.916523in}}%
\pgfpathlineto{\pgfqpoint{3.114879in}{1.916068in}}%
\pgfpathlineto{\pgfqpoint{3.115027in}{1.916037in}}%
\pgfpathlineto{\pgfqpoint{3.115175in}{1.915659in}}%
\pgfpathlineto{\pgfqpoint{3.115225in}{1.915610in}}%
\pgfpathlineto{\pgfqpoint{3.115324in}{1.914800in}}%
\pgfpathlineto{\pgfqpoint{3.115621in}{1.745845in}}%
\pgfpathlineto{\pgfqpoint{3.115869in}{1.532026in}}%
\pgfpathlineto{\pgfqpoint{3.116640in}{1.532051in}}%
\pgfpathlineto{\pgfqpoint{3.117734in}{1.532133in}}%
\pgfpathlineto{\pgfqpoint{3.117968in}{1.944704in}}%
\pgfpathlineto{\pgfqpoint{3.118899in}{1.915091in}}%
\pgfpathlineto{\pgfqpoint{3.118948in}{1.917264in}}%
\pgfpathlineto{\pgfqpoint{3.119640in}{1.916524in}}%
\pgfpathlineto{\pgfqpoint{3.119739in}{1.916593in}}%
\pgfpathlineto{\pgfqpoint{3.120333in}{1.901763in}}%
\pgfpathlineto{\pgfqpoint{3.120794in}{1.532027in}}%
\pgfpathlineto{\pgfqpoint{3.121839in}{1.532056in}}%
\pgfpathlineto{\pgfqpoint{3.122679in}{1.532121in}}%
\pgfpathlineto{\pgfqpoint{3.122768in}{1.933836in}}%
\pgfpathlineto{\pgfqpoint{3.123941in}{1.916503in}}%
\pgfpathlineto{\pgfqpoint{3.124682in}{1.916103in}}%
\pgfpathlineto{\pgfqpoint{3.124732in}{1.916172in}}%
\pgfpathlineto{\pgfqpoint{3.125226in}{1.914035in}}%
\pgfpathlineto{\pgfqpoint{3.125623in}{1.670121in}}%
\pgfpathlineto{\pgfqpoint{3.125629in}{1.675228in}}%
\pgfpathlineto{\pgfqpoint{3.125759in}{1.532026in}}%
\pgfpathlineto{\pgfqpoint{3.126381in}{1.532040in}}%
\pgfpathlineto{\pgfqpoint{3.127625in}{1.532132in}}%
\pgfpathlineto{\pgfqpoint{3.127859in}{1.944352in}}%
\pgfpathlineto{\pgfqpoint{3.128791in}{1.915104in}}%
\pgfpathlineto{\pgfqpoint{3.128841in}{1.917228in}}%
\pgfpathlineto{\pgfqpoint{3.129533in}{1.916510in}}%
\pgfpathlineto{\pgfqpoint{3.129632in}{1.916572in}}%
\pgfpathlineto{\pgfqpoint{3.130226in}{1.900946in}}%
\pgfpathlineto{\pgfqpoint{3.130685in}{1.532027in}}%
\pgfpathlineto{\pgfqpoint{3.131733in}{1.532056in}}%
\pgfpathlineto{\pgfqpoint{3.132569in}{1.532113in}}%
\pgfpathlineto{\pgfqpoint{3.132670in}{1.934867in}}%
\pgfpathlineto{\pgfqpoint{3.134080in}{1.916268in}}%
\pgfpathlineto{\pgfqpoint{3.134228in}{1.917769in}}%
\pgfpathlineto{\pgfqpoint{3.134278in}{1.914668in}}%
\pgfpathlineto{\pgfqpoint{3.134772in}{1.916887in}}%
\pgfpathlineto{\pgfqpoint{3.135217in}{1.881114in}}%
\pgfpathlineto{\pgfqpoint{3.135604in}{1.532024in}}%
\pgfpathlineto{\pgfqpoint{3.136529in}{1.532055in}}%
\pgfpathlineto{\pgfqpoint{3.137515in}{1.532126in}}%
\pgfpathlineto{\pgfqpoint{3.137752in}{1.943535in}}%
\pgfpathlineto{\pgfqpoint{3.138730in}{1.916989in}}%
\pgfpathlineto{\pgfqpoint{3.138780in}{1.915493in}}%
\pgfpathlineto{\pgfqpoint{3.139522in}{1.916296in}}%
\pgfpathlineto{\pgfqpoint{3.139621in}{1.916334in}}%
\pgfpathlineto{\pgfqpoint{3.140066in}{1.913460in}}%
\pgfpathlineto{\pgfqpoint{3.140470in}{1.612894in}}%
\pgfpathlineto{\pgfqpoint{3.140579in}{1.532028in}}%
\pgfpathlineto{\pgfqpoint{3.141230in}{1.532047in}}%
\pgfpathlineto{\pgfqpoint{3.142460in}{1.532116in}}%
\pgfpathlineto{\pgfqpoint{3.142550in}{1.933574in}}%
\pgfpathlineto{\pgfqpoint{3.143749in}{1.916091in}}%
\pgfpathlineto{\pgfqpoint{3.144144in}{1.916380in}}%
\pgfpathlineto{\pgfqpoint{3.144491in}{1.916178in}}%
\pgfpathlineto{\pgfqpoint{3.144985in}{1.914982in}}%
\pgfpathlineto{\pgfqpoint{3.145232in}{1.793176in}}%
\pgfpathlineto{\pgfqpoint{3.145542in}{1.532026in}}%
\pgfpathlineto{\pgfqpoint{3.146312in}{1.532051in}}%
\pgfpathlineto{\pgfqpoint{3.147407in}{1.532133in}}%
\pgfpathlineto{\pgfqpoint{3.147641in}{1.944388in}}%
\pgfpathlineto{\pgfqpoint{3.148575in}{1.915070in}}%
\pgfpathlineto{\pgfqpoint{3.148625in}{1.917281in}}%
\pgfpathlineto{\pgfqpoint{3.149317in}{1.916444in}}%
\pgfpathlineto{\pgfqpoint{3.149416in}{1.916550in}}%
\pgfpathlineto{\pgfqpoint{3.150009in}{1.900169in}}%
\pgfpathlineto{\pgfqpoint{3.150467in}{1.532027in}}%
\pgfpathlineto{\pgfqpoint{3.151515in}{1.532056in}}%
\pgfpathlineto{\pgfqpoint{3.152350in}{1.532113in}}%
\pgfpathlineto{\pgfqpoint{3.152451in}{1.934847in}}%
\pgfpathlineto{\pgfqpoint{3.153868in}{1.916506in}}%
\pgfpathlineto{\pgfqpoint{3.154115in}{1.917832in}}%
\pgfpathlineto{\pgfqpoint{3.154066in}{1.914605in}}%
\pgfpathlineto{\pgfqpoint{3.154461in}{1.916999in}}%
\pgfpathlineto{\pgfqpoint{3.155005in}{1.877939in}}%
\pgfpathlineto{\pgfqpoint{3.155414in}{1.532027in}}%
\pgfpathlineto{\pgfqpoint{3.156314in}{1.532046in}}%
\pgfpathlineto{\pgfqpoint{3.157296in}{1.532112in}}%
\pgfpathlineto{\pgfqpoint{3.157414in}{1.937014in}}%
\pgfpathlineto{\pgfqpoint{3.158834in}{1.915005in}}%
\pgfpathlineto{\pgfqpoint{3.158883in}{1.917312in}}%
\pgfpathlineto{\pgfqpoint{3.159575in}{1.916300in}}%
\pgfpathlineto{\pgfqpoint{3.159872in}{1.908061in}}%
\pgfpathlineto{\pgfqpoint{3.160375in}{1.532023in}}%
\pgfpathlineto{\pgfqpoint{3.161737in}{1.532067in}}%
\pgfpathlineto{\pgfqpoint{3.162240in}{1.532102in}}%
\pgfpathlineto{\pgfqpoint{3.162279in}{1.696547in}}%
\pgfpathlineto{\pgfqpoint{3.162547in}{1.939087in}}%
\pgfpathlineto{\pgfqpoint{3.163041in}{1.917785in}}%
\pgfpathlineto{\pgfqpoint{3.163090in}{1.914256in}}%
\pgfpathlineto{\pgfqpoint{3.163140in}{1.918031in}}%
\pgfpathlineto{\pgfqpoint{3.163783in}{1.916278in}}%
\pgfpathlineto{\pgfqpoint{3.163931in}{1.916943in}}%
\pgfpathlineto{\pgfqpoint{3.163981in}{1.915474in}}%
\pgfpathlineto{\pgfqpoint{3.164475in}{1.916355in}}%
\pgfpathlineto{\pgfqpoint{3.164821in}{1.908559in}}%
\pgfpathlineto{\pgfqpoint{3.165329in}{1.532030in}}%
\pgfpathlineto{\pgfqpoint{3.166550in}{1.532057in}}%
\pgfpathlineto{\pgfqpoint{3.167192in}{1.532279in}}%
\pgfpathlineto{\pgfqpoint{3.167296in}{1.935958in}}%
\pgfpathlineto{\pgfqpoint{3.168234in}{1.916880in}}%
\pgfpathlineto{\pgfqpoint{3.168283in}{1.916773in}}%
\pgfpathlineto{\pgfqpoint{3.168432in}{1.913599in}}%
\pgfpathlineto{\pgfqpoint{3.168481in}{1.918940in}}%
\pgfpathlineto{\pgfqpoint{3.169075in}{1.915601in}}%
\pgfpathlineto{\pgfqpoint{3.169272in}{1.917030in}}%
\pgfpathlineto{\pgfqpoint{3.169668in}{1.914872in}}%
\pgfpathlineto{\pgfqpoint{3.169717in}{1.915584in}}%
\pgfpathlineto{\pgfqpoint{3.170014in}{1.750507in}}%
\pgfpathlineto{\pgfqpoint{3.170132in}{1.670918in}}%
\pgfpathlineto{\pgfqpoint{3.170132in}{1.670918in}}%
\pgfusepath{stroke}%
\end{pgfscope}%
\begin{pgfscope}%
\pgfsetrectcap%
\pgfsetmiterjoin%
\pgfsetlinewidth{0.803000pt}%
\definecolor{currentstroke}{rgb}{0.000000,0.000000,0.000000}%
\pgfsetstrokecolor{currentstroke}%
\pgfsetdash{}{0pt}%
\pgfpathmoveto{\pgfqpoint{0.573768in}{1.532029in}}%
\pgfpathlineto{\pgfqpoint{0.573768in}{2.438279in}}%
\pgfusepath{stroke}%
\end{pgfscope}%
\begin{pgfscope}%
\pgfsetrectcap%
\pgfsetmiterjoin%
\pgfsetlinewidth{0.803000pt}%
\definecolor{currentstroke}{rgb}{0.000000,0.000000,0.000000}%
\pgfsetstrokecolor{currentstroke}%
\pgfsetdash{}{0pt}%
\pgfpathmoveto{\pgfqpoint{3.293768in}{1.532029in}}%
\pgfpathlineto{\pgfqpoint{3.293768in}{2.438279in}}%
\pgfusepath{stroke}%
\end{pgfscope}%
\begin{pgfscope}%
\pgfsetrectcap%
\pgfsetmiterjoin%
\pgfsetlinewidth{0.803000pt}%
\definecolor{currentstroke}{rgb}{0.000000,0.000000,0.000000}%
\pgfsetstrokecolor{currentstroke}%
\pgfsetdash{}{0pt}%
\pgfpathmoveto{\pgfqpoint{0.573768in}{1.532029in}}%
\pgfpathlineto{\pgfqpoint{3.293768in}{1.532029in}}%
\pgfusepath{stroke}%
\end{pgfscope}%
\begin{pgfscope}%
\pgfsetrectcap%
\pgfsetmiterjoin%
\pgfsetlinewidth{0.803000pt}%
\definecolor{currentstroke}{rgb}{0.000000,0.000000,0.000000}%
\pgfsetstrokecolor{currentstroke}%
\pgfsetdash{}{0pt}%
\pgfpathmoveto{\pgfqpoint{0.573768in}{2.438279in}}%
\pgfpathlineto{\pgfqpoint{3.293768in}{2.438279in}}%
\pgfusepath{stroke}%
\end{pgfscope}%
\begin{pgfscope}%
\definecolor{textcolor}{rgb}{0.000000,0.000000,0.000000}%
\pgfsetstrokecolor{textcolor}%
\pgfsetfillcolor{textcolor}%
\pgftext[x=0.628168in,y=2.374842in,left,top]{\color{textcolor}{\rmfamily\fontsize{8.000000}{9.600000}\selectfont\catcode`\^=\active\def^{\ifmmode\sp\else\^{}\fi}\catcode`\%=\active\def%{\%}(b)}}%
\end{pgfscope}%
\begin{pgfscope}%
\pgfsetbuttcap%
\pgfsetmiterjoin%
\definecolor{currentfill}{rgb}{1.000000,1.000000,1.000000}%
\pgfsetfillcolor{currentfill}%
\pgfsetlinewidth{0.000000pt}%
\definecolor{currentstroke}{rgb}{0.000000,0.000000,0.000000}%
\pgfsetstrokecolor{currentstroke}%
\pgfsetstrokeopacity{0.000000}%
\pgfsetdash{}{0pt}%
\pgfpathmoveto{\pgfqpoint{0.573768in}{0.462654in}}%
\pgfpathlineto{\pgfqpoint{3.293768in}{0.462654in}}%
\pgfpathlineto{\pgfqpoint{3.293768in}{1.368904in}}%
\pgfpathlineto{\pgfqpoint{0.573768in}{1.368904in}}%
\pgfpathlineto{\pgfqpoint{0.573768in}{0.462654in}}%
\pgfpathclose%
\pgfusepath{fill}%
\end{pgfscope}%
\begin{pgfscope}%
\pgfpathrectangle{\pgfqpoint{0.573768in}{0.462654in}}{\pgfqpoint{2.720000in}{0.906250in}}%
\pgfusepath{clip}%
\pgfsetbuttcap%
\pgfsetroundjoin%
\pgfsetlinewidth{0.501875pt}%
\definecolor{currentstroke}{rgb}{0.690196,0.690196,0.690196}%
\pgfsetstrokecolor{currentstroke}%
\pgfsetstrokeopacity{0.250000}%
\pgfsetdash{{1.850000pt}{0.800000pt}}{0.000000pt}%
\pgfpathmoveto{\pgfqpoint{0.697405in}{0.462654in}}%
\pgfpathlineto{\pgfqpoint{0.697405in}{1.368904in}}%
\pgfusepath{stroke}%
\end{pgfscope}%
\begin{pgfscope}%
\pgfsetbuttcap%
\pgfsetroundjoin%
\definecolor{currentfill}{rgb}{0.000000,0.000000,0.000000}%
\pgfsetfillcolor{currentfill}%
\pgfsetlinewidth{0.803000pt}%
\definecolor{currentstroke}{rgb}{0.000000,0.000000,0.000000}%
\pgfsetstrokecolor{currentstroke}%
\pgfsetdash{}{0pt}%
\pgfsys@defobject{currentmarker}{\pgfqpoint{0.000000in}{-0.048611in}}{\pgfqpoint{0.000000in}{0.000000in}}{%
\pgfpathmoveto{\pgfqpoint{0.000000in}{0.000000in}}%
\pgfpathlineto{\pgfqpoint{0.000000in}{-0.048611in}}%
\pgfusepath{stroke,fill}%
}%
\begin{pgfscope}%
\pgfsys@transformshift{0.697405in}{0.462654in}%
\pgfsys@useobject{currentmarker}{}%
\end{pgfscope}%
\end{pgfscope}%
\begin{pgfscope}%
\definecolor{textcolor}{rgb}{0.000000,0.000000,0.000000}%
\pgfsetstrokecolor{textcolor}%
\pgfsetfillcolor{textcolor}%
\pgftext[x=0.697405in,y=0.365432in,,top]{\color{textcolor}{\rmfamily\fontsize{8.000000}{9.600000}\selectfont\catcode`\^=\active\def^{\ifmmode\sp\else\^{}\fi}\catcode`\%=\active\def%{\%}$\mathdefault{0}$}}%
\end{pgfscope}%
\begin{pgfscope}%
\pgfpathrectangle{\pgfqpoint{0.573768in}{0.462654in}}{\pgfqpoint{2.720000in}{0.906250in}}%
\pgfusepath{clip}%
\pgfsetbuttcap%
\pgfsetroundjoin%
\pgfsetlinewidth{0.501875pt}%
\definecolor{currentstroke}{rgb}{0.690196,0.690196,0.690196}%
\pgfsetstrokecolor{currentstroke}%
\pgfsetstrokeopacity{0.250000}%
\pgfsetdash{{1.850000pt}{0.800000pt}}{0.000000pt}%
\pgfpathmoveto{\pgfqpoint{1.191950in}{0.462654in}}%
\pgfpathlineto{\pgfqpoint{1.191950in}{1.368904in}}%
\pgfusepath{stroke}%
\end{pgfscope}%
\begin{pgfscope}%
\pgfsetbuttcap%
\pgfsetroundjoin%
\definecolor{currentfill}{rgb}{0.000000,0.000000,0.000000}%
\pgfsetfillcolor{currentfill}%
\pgfsetlinewidth{0.803000pt}%
\definecolor{currentstroke}{rgb}{0.000000,0.000000,0.000000}%
\pgfsetstrokecolor{currentstroke}%
\pgfsetdash{}{0pt}%
\pgfsys@defobject{currentmarker}{\pgfqpoint{0.000000in}{-0.048611in}}{\pgfqpoint{0.000000in}{0.000000in}}{%
\pgfpathmoveto{\pgfqpoint{0.000000in}{0.000000in}}%
\pgfpathlineto{\pgfqpoint{0.000000in}{-0.048611in}}%
\pgfusepath{stroke,fill}%
}%
\begin{pgfscope}%
\pgfsys@transformshift{1.191950in}{0.462654in}%
\pgfsys@useobject{currentmarker}{}%
\end{pgfscope}%
\end{pgfscope}%
\begin{pgfscope}%
\definecolor{textcolor}{rgb}{0.000000,0.000000,0.000000}%
\pgfsetstrokecolor{textcolor}%
\pgfsetfillcolor{textcolor}%
\pgftext[x=1.191950in,y=0.365432in,,top]{\color{textcolor}{\rmfamily\fontsize{8.000000}{9.600000}\selectfont\catcode`\^=\active\def^{\ifmmode\sp\else\^{}\fi}\catcode`\%=\active\def%{\%}$\mathdefault{1}$}}%
\end{pgfscope}%
\begin{pgfscope}%
\pgfpathrectangle{\pgfqpoint{0.573768in}{0.462654in}}{\pgfqpoint{2.720000in}{0.906250in}}%
\pgfusepath{clip}%
\pgfsetbuttcap%
\pgfsetroundjoin%
\pgfsetlinewidth{0.501875pt}%
\definecolor{currentstroke}{rgb}{0.690196,0.690196,0.690196}%
\pgfsetstrokecolor{currentstroke}%
\pgfsetstrokeopacity{0.250000}%
\pgfsetdash{{1.850000pt}{0.800000pt}}{0.000000pt}%
\pgfpathmoveto{\pgfqpoint{1.686496in}{0.462654in}}%
\pgfpathlineto{\pgfqpoint{1.686496in}{1.368904in}}%
\pgfusepath{stroke}%
\end{pgfscope}%
\begin{pgfscope}%
\pgfsetbuttcap%
\pgfsetroundjoin%
\definecolor{currentfill}{rgb}{0.000000,0.000000,0.000000}%
\pgfsetfillcolor{currentfill}%
\pgfsetlinewidth{0.803000pt}%
\definecolor{currentstroke}{rgb}{0.000000,0.000000,0.000000}%
\pgfsetstrokecolor{currentstroke}%
\pgfsetdash{}{0pt}%
\pgfsys@defobject{currentmarker}{\pgfqpoint{0.000000in}{-0.048611in}}{\pgfqpoint{0.000000in}{0.000000in}}{%
\pgfpathmoveto{\pgfqpoint{0.000000in}{0.000000in}}%
\pgfpathlineto{\pgfqpoint{0.000000in}{-0.048611in}}%
\pgfusepath{stroke,fill}%
}%
\begin{pgfscope}%
\pgfsys@transformshift{1.686496in}{0.462654in}%
\pgfsys@useobject{currentmarker}{}%
\end{pgfscope}%
\end{pgfscope}%
\begin{pgfscope}%
\definecolor{textcolor}{rgb}{0.000000,0.000000,0.000000}%
\pgfsetstrokecolor{textcolor}%
\pgfsetfillcolor{textcolor}%
\pgftext[x=1.686496in,y=0.365432in,,top]{\color{textcolor}{\rmfamily\fontsize{8.000000}{9.600000}\selectfont\catcode`\^=\active\def^{\ifmmode\sp\else\^{}\fi}\catcode`\%=\active\def%{\%}$\mathdefault{2}$}}%
\end{pgfscope}%
\begin{pgfscope}%
\pgfpathrectangle{\pgfqpoint{0.573768in}{0.462654in}}{\pgfqpoint{2.720000in}{0.906250in}}%
\pgfusepath{clip}%
\pgfsetbuttcap%
\pgfsetroundjoin%
\pgfsetlinewidth{0.501875pt}%
\definecolor{currentstroke}{rgb}{0.690196,0.690196,0.690196}%
\pgfsetstrokecolor{currentstroke}%
\pgfsetstrokeopacity{0.250000}%
\pgfsetdash{{1.850000pt}{0.800000pt}}{0.000000pt}%
\pgfpathmoveto{\pgfqpoint{2.181041in}{0.462654in}}%
\pgfpathlineto{\pgfqpoint{2.181041in}{1.368904in}}%
\pgfusepath{stroke}%
\end{pgfscope}%
\begin{pgfscope}%
\pgfsetbuttcap%
\pgfsetroundjoin%
\definecolor{currentfill}{rgb}{0.000000,0.000000,0.000000}%
\pgfsetfillcolor{currentfill}%
\pgfsetlinewidth{0.803000pt}%
\definecolor{currentstroke}{rgb}{0.000000,0.000000,0.000000}%
\pgfsetstrokecolor{currentstroke}%
\pgfsetdash{}{0pt}%
\pgfsys@defobject{currentmarker}{\pgfqpoint{0.000000in}{-0.048611in}}{\pgfqpoint{0.000000in}{0.000000in}}{%
\pgfpathmoveto{\pgfqpoint{0.000000in}{0.000000in}}%
\pgfpathlineto{\pgfqpoint{0.000000in}{-0.048611in}}%
\pgfusepath{stroke,fill}%
}%
\begin{pgfscope}%
\pgfsys@transformshift{2.181041in}{0.462654in}%
\pgfsys@useobject{currentmarker}{}%
\end{pgfscope}%
\end{pgfscope}%
\begin{pgfscope}%
\definecolor{textcolor}{rgb}{0.000000,0.000000,0.000000}%
\pgfsetstrokecolor{textcolor}%
\pgfsetfillcolor{textcolor}%
\pgftext[x=2.181041in,y=0.365432in,,top]{\color{textcolor}{\rmfamily\fontsize{8.000000}{9.600000}\selectfont\catcode`\^=\active\def^{\ifmmode\sp\else\^{}\fi}\catcode`\%=\active\def%{\%}$\mathdefault{3}$}}%
\end{pgfscope}%
\begin{pgfscope}%
\pgfpathrectangle{\pgfqpoint{0.573768in}{0.462654in}}{\pgfqpoint{2.720000in}{0.906250in}}%
\pgfusepath{clip}%
\pgfsetbuttcap%
\pgfsetroundjoin%
\pgfsetlinewidth{0.501875pt}%
\definecolor{currentstroke}{rgb}{0.690196,0.690196,0.690196}%
\pgfsetstrokecolor{currentstroke}%
\pgfsetstrokeopacity{0.250000}%
\pgfsetdash{{1.850000pt}{0.800000pt}}{0.000000pt}%
\pgfpathmoveto{\pgfqpoint{2.675587in}{0.462654in}}%
\pgfpathlineto{\pgfqpoint{2.675587in}{1.368904in}}%
\pgfusepath{stroke}%
\end{pgfscope}%
\begin{pgfscope}%
\pgfsetbuttcap%
\pgfsetroundjoin%
\definecolor{currentfill}{rgb}{0.000000,0.000000,0.000000}%
\pgfsetfillcolor{currentfill}%
\pgfsetlinewidth{0.803000pt}%
\definecolor{currentstroke}{rgb}{0.000000,0.000000,0.000000}%
\pgfsetstrokecolor{currentstroke}%
\pgfsetdash{}{0pt}%
\pgfsys@defobject{currentmarker}{\pgfqpoint{0.000000in}{-0.048611in}}{\pgfqpoint{0.000000in}{0.000000in}}{%
\pgfpathmoveto{\pgfqpoint{0.000000in}{0.000000in}}%
\pgfpathlineto{\pgfqpoint{0.000000in}{-0.048611in}}%
\pgfusepath{stroke,fill}%
}%
\begin{pgfscope}%
\pgfsys@transformshift{2.675587in}{0.462654in}%
\pgfsys@useobject{currentmarker}{}%
\end{pgfscope}%
\end{pgfscope}%
\begin{pgfscope}%
\definecolor{textcolor}{rgb}{0.000000,0.000000,0.000000}%
\pgfsetstrokecolor{textcolor}%
\pgfsetfillcolor{textcolor}%
\pgftext[x=2.675587in,y=0.365432in,,top]{\color{textcolor}{\rmfamily\fontsize{8.000000}{9.600000}\selectfont\catcode`\^=\active\def^{\ifmmode\sp\else\^{}\fi}\catcode`\%=\active\def%{\%}$\mathdefault{4}$}}%
\end{pgfscope}%
\begin{pgfscope}%
\pgfpathrectangle{\pgfqpoint{0.573768in}{0.462654in}}{\pgfqpoint{2.720000in}{0.906250in}}%
\pgfusepath{clip}%
\pgfsetbuttcap%
\pgfsetroundjoin%
\pgfsetlinewidth{0.501875pt}%
\definecolor{currentstroke}{rgb}{0.690196,0.690196,0.690196}%
\pgfsetstrokecolor{currentstroke}%
\pgfsetstrokeopacity{0.250000}%
\pgfsetdash{{1.850000pt}{0.800000pt}}{0.000000pt}%
\pgfpathmoveto{\pgfqpoint{3.170132in}{0.462654in}}%
\pgfpathlineto{\pgfqpoint{3.170132in}{1.368904in}}%
\pgfusepath{stroke}%
\end{pgfscope}%
\begin{pgfscope}%
\pgfsetbuttcap%
\pgfsetroundjoin%
\definecolor{currentfill}{rgb}{0.000000,0.000000,0.000000}%
\pgfsetfillcolor{currentfill}%
\pgfsetlinewidth{0.803000pt}%
\definecolor{currentstroke}{rgb}{0.000000,0.000000,0.000000}%
\pgfsetstrokecolor{currentstroke}%
\pgfsetdash{}{0pt}%
\pgfsys@defobject{currentmarker}{\pgfqpoint{0.000000in}{-0.048611in}}{\pgfqpoint{0.000000in}{0.000000in}}{%
\pgfpathmoveto{\pgfqpoint{0.000000in}{0.000000in}}%
\pgfpathlineto{\pgfqpoint{0.000000in}{-0.048611in}}%
\pgfusepath{stroke,fill}%
}%
\begin{pgfscope}%
\pgfsys@transformshift{3.170132in}{0.462654in}%
\pgfsys@useobject{currentmarker}{}%
\end{pgfscope}%
\end{pgfscope}%
\begin{pgfscope}%
\definecolor{textcolor}{rgb}{0.000000,0.000000,0.000000}%
\pgfsetstrokecolor{textcolor}%
\pgfsetfillcolor{textcolor}%
\pgftext[x=3.170132in,y=0.365432in,,top]{\color{textcolor}{\rmfamily\fontsize{8.000000}{9.600000}\selectfont\catcode`\^=\active\def^{\ifmmode\sp\else\^{}\fi}\catcode`\%=\active\def%{\%}$\mathdefault{5}$}}%
\end{pgfscope}%
\begin{pgfscope}%
\definecolor{textcolor}{rgb}{0.000000,0.000000,0.000000}%
\pgfsetstrokecolor{textcolor}%
\pgfsetfillcolor{textcolor}%
\pgftext[x=1.933768in,y=0.211111in,,top]{\color{textcolor}{\rmfamily\fontsize{8.000000}{9.600000}\selectfont\catcode`\^=\active\def^{\ifmmode\sp\else\^{}\fi}\catcode`\%=\active\def%{\%}Time (ms)}}%
\end{pgfscope}%
\begin{pgfscope}%
\pgfpathrectangle{\pgfqpoint{0.573768in}{0.462654in}}{\pgfqpoint{2.720000in}{0.906250in}}%
\pgfusepath{clip}%
\pgfsetbuttcap%
\pgfsetroundjoin%
\pgfsetlinewidth{0.501875pt}%
\definecolor{currentstroke}{rgb}{0.690196,0.690196,0.690196}%
\pgfsetstrokecolor{currentstroke}%
\pgfsetstrokeopacity{0.250000}%
\pgfsetdash{{1.850000pt}{0.800000pt}}{0.000000pt}%
\pgfpathmoveto{\pgfqpoint{0.573768in}{0.462654in}}%
\pgfpathlineto{\pgfqpoint{3.293768in}{0.462654in}}%
\pgfusepath{stroke}%
\end{pgfscope}%
\begin{pgfscope}%
\pgfsetbuttcap%
\pgfsetroundjoin%
\definecolor{currentfill}{rgb}{0.000000,0.000000,0.000000}%
\pgfsetfillcolor{currentfill}%
\pgfsetlinewidth{0.803000pt}%
\definecolor{currentstroke}{rgb}{0.000000,0.000000,0.000000}%
\pgfsetstrokecolor{currentstroke}%
\pgfsetdash{}{0pt}%
\pgfsys@defobject{currentmarker}{\pgfqpoint{-0.048611in}{0.000000in}}{\pgfqpoint{-0.000000in}{0.000000in}}{%
\pgfpathmoveto{\pgfqpoint{-0.000000in}{0.000000in}}%
\pgfpathlineto{\pgfqpoint{-0.048611in}{0.000000in}}%
\pgfusepath{stroke,fill}%
}%
\begin{pgfscope}%
\pgfsys@transformshift{0.573768in}{0.462654in}%
\pgfsys@useobject{currentmarker}{}%
\end{pgfscope}%
\end{pgfscope}%
\begin{pgfscope}%
\definecolor{textcolor}{rgb}{0.000000,0.000000,0.000000}%
\pgfsetstrokecolor{textcolor}%
\pgfsetfillcolor{textcolor}%
\pgftext[x=0.325695in, y=0.424074in, left, base]{\color{textcolor}{\rmfamily\fontsize{8.000000}{9.600000}\selectfont\catcode`\^=\active\def^{\ifmmode\sp\else\^{}\fi}\catcode`\%=\active\def%{\%}0.0}}%
\end{pgfscope}%
\begin{pgfscope}%
\pgfpathrectangle{\pgfqpoint{0.573768in}{0.462654in}}{\pgfqpoint{2.720000in}{0.906250in}}%
\pgfusepath{clip}%
\pgfsetbuttcap%
\pgfsetroundjoin%
\pgfsetlinewidth{0.501875pt}%
\definecolor{currentstroke}{rgb}{0.690196,0.690196,0.690196}%
\pgfsetstrokecolor{currentstroke}%
\pgfsetstrokeopacity{0.250000}%
\pgfsetdash{{1.850000pt}{0.800000pt}}{0.000000pt}%
\pgfpathmoveto{\pgfqpoint{0.573768in}{0.689217in}}%
\pgfpathlineto{\pgfqpoint{3.293768in}{0.689217in}}%
\pgfusepath{stroke}%
\end{pgfscope}%
\begin{pgfscope}%
\pgfsetbuttcap%
\pgfsetroundjoin%
\definecolor{currentfill}{rgb}{0.000000,0.000000,0.000000}%
\pgfsetfillcolor{currentfill}%
\pgfsetlinewidth{0.803000pt}%
\definecolor{currentstroke}{rgb}{0.000000,0.000000,0.000000}%
\pgfsetstrokecolor{currentstroke}%
\pgfsetdash{}{0pt}%
\pgfsys@defobject{currentmarker}{\pgfqpoint{-0.048611in}{0.000000in}}{\pgfqpoint{-0.000000in}{0.000000in}}{%
\pgfpathmoveto{\pgfqpoint{-0.000000in}{0.000000in}}%
\pgfpathlineto{\pgfqpoint{-0.048611in}{0.000000in}}%
\pgfusepath{stroke,fill}%
}%
\begin{pgfscope}%
\pgfsys@transformshift{0.573768in}{0.689217in}%
\pgfsys@useobject{currentmarker}{}%
\end{pgfscope}%
\end{pgfscope}%
\begin{pgfscope}%
\definecolor{textcolor}{rgb}{0.000000,0.000000,0.000000}%
\pgfsetstrokecolor{textcolor}%
\pgfsetfillcolor{textcolor}%
\pgftext[x=0.325695in, y=0.650637in, left, base]{\color{textcolor}{\rmfamily\fontsize{8.000000}{9.600000}\selectfont\catcode`\^=\active\def^{\ifmmode\sp\else\^{}\fi}\catcode`\%=\active\def%{\%}6.2}}%
\end{pgfscope}%
\begin{pgfscope}%
\pgfpathrectangle{\pgfqpoint{0.573768in}{0.462654in}}{\pgfqpoint{2.720000in}{0.906250in}}%
\pgfusepath{clip}%
\pgfsetbuttcap%
\pgfsetroundjoin%
\pgfsetlinewidth{0.501875pt}%
\definecolor{currentstroke}{rgb}{0.690196,0.690196,0.690196}%
\pgfsetstrokecolor{currentstroke}%
\pgfsetstrokeopacity{0.250000}%
\pgfsetdash{{1.850000pt}{0.800000pt}}{0.000000pt}%
\pgfpathmoveto{\pgfqpoint{0.573768in}{0.915779in}}%
\pgfpathlineto{\pgfqpoint{3.293768in}{0.915779in}}%
\pgfusepath{stroke}%
\end{pgfscope}%
\begin{pgfscope}%
\pgfsetbuttcap%
\pgfsetroundjoin%
\definecolor{currentfill}{rgb}{0.000000,0.000000,0.000000}%
\pgfsetfillcolor{currentfill}%
\pgfsetlinewidth{0.803000pt}%
\definecolor{currentstroke}{rgb}{0.000000,0.000000,0.000000}%
\pgfsetstrokecolor{currentstroke}%
\pgfsetdash{}{0pt}%
\pgfsys@defobject{currentmarker}{\pgfqpoint{-0.048611in}{0.000000in}}{\pgfqpoint{-0.000000in}{0.000000in}}{%
\pgfpathmoveto{\pgfqpoint{-0.000000in}{0.000000in}}%
\pgfpathlineto{\pgfqpoint{-0.048611in}{0.000000in}}%
\pgfusepath{stroke,fill}%
}%
\begin{pgfscope}%
\pgfsys@transformshift{0.573768in}{0.915779in}%
\pgfsys@useobject{currentmarker}{}%
\end{pgfscope}%
\end{pgfscope}%
\begin{pgfscope}%
\definecolor{textcolor}{rgb}{0.000000,0.000000,0.000000}%
\pgfsetstrokecolor{textcolor}%
\pgfsetfillcolor{textcolor}%
\pgftext[x=0.266667in, y=0.877199in, left, base]{\color{textcolor}{\rmfamily\fontsize{8.000000}{9.600000}\selectfont\catcode`\^=\active\def^{\ifmmode\sp\else\^{}\fi}\catcode`\%=\active\def%{\%}12.5}}%
\end{pgfscope}%
\begin{pgfscope}%
\pgfpathrectangle{\pgfqpoint{0.573768in}{0.462654in}}{\pgfqpoint{2.720000in}{0.906250in}}%
\pgfusepath{clip}%
\pgfsetbuttcap%
\pgfsetroundjoin%
\pgfsetlinewidth{0.501875pt}%
\definecolor{currentstroke}{rgb}{0.690196,0.690196,0.690196}%
\pgfsetstrokecolor{currentstroke}%
\pgfsetstrokeopacity{0.250000}%
\pgfsetdash{{1.850000pt}{0.800000pt}}{0.000000pt}%
\pgfpathmoveto{\pgfqpoint{0.573768in}{1.142342in}}%
\pgfpathlineto{\pgfqpoint{3.293768in}{1.142342in}}%
\pgfusepath{stroke}%
\end{pgfscope}%
\begin{pgfscope}%
\pgfsetbuttcap%
\pgfsetroundjoin%
\definecolor{currentfill}{rgb}{0.000000,0.000000,0.000000}%
\pgfsetfillcolor{currentfill}%
\pgfsetlinewidth{0.803000pt}%
\definecolor{currentstroke}{rgb}{0.000000,0.000000,0.000000}%
\pgfsetstrokecolor{currentstroke}%
\pgfsetdash{}{0pt}%
\pgfsys@defobject{currentmarker}{\pgfqpoint{-0.048611in}{0.000000in}}{\pgfqpoint{-0.000000in}{0.000000in}}{%
\pgfpathmoveto{\pgfqpoint{-0.000000in}{0.000000in}}%
\pgfpathlineto{\pgfqpoint{-0.048611in}{0.000000in}}%
\pgfusepath{stroke,fill}%
}%
\begin{pgfscope}%
\pgfsys@transformshift{0.573768in}{1.142342in}%
\pgfsys@useobject{currentmarker}{}%
\end{pgfscope}%
\end{pgfscope}%
\begin{pgfscope}%
\definecolor{textcolor}{rgb}{0.000000,0.000000,0.000000}%
\pgfsetstrokecolor{textcolor}%
\pgfsetfillcolor{textcolor}%
\pgftext[x=0.266667in, y=1.103762in, left, base]{\color{textcolor}{\rmfamily\fontsize{8.000000}{9.600000}\selectfont\catcode`\^=\active\def^{\ifmmode\sp\else\^{}\fi}\catcode`\%=\active\def%{\%}18.8}}%
\end{pgfscope}%
\begin{pgfscope}%
\pgfpathrectangle{\pgfqpoint{0.573768in}{0.462654in}}{\pgfqpoint{2.720000in}{0.906250in}}%
\pgfusepath{clip}%
\pgfsetbuttcap%
\pgfsetroundjoin%
\pgfsetlinewidth{0.501875pt}%
\definecolor{currentstroke}{rgb}{0.690196,0.690196,0.690196}%
\pgfsetstrokecolor{currentstroke}%
\pgfsetstrokeopacity{0.250000}%
\pgfsetdash{{1.850000pt}{0.800000pt}}{0.000000pt}%
\pgfpathmoveto{\pgfqpoint{0.573768in}{1.368904in}}%
\pgfpathlineto{\pgfqpoint{3.293768in}{1.368904in}}%
\pgfusepath{stroke}%
\end{pgfscope}%
\begin{pgfscope}%
\pgfsetbuttcap%
\pgfsetroundjoin%
\definecolor{currentfill}{rgb}{0.000000,0.000000,0.000000}%
\pgfsetfillcolor{currentfill}%
\pgfsetlinewidth{0.803000pt}%
\definecolor{currentstroke}{rgb}{0.000000,0.000000,0.000000}%
\pgfsetstrokecolor{currentstroke}%
\pgfsetdash{}{0pt}%
\pgfsys@defobject{currentmarker}{\pgfqpoint{-0.048611in}{0.000000in}}{\pgfqpoint{-0.000000in}{0.000000in}}{%
\pgfpathmoveto{\pgfqpoint{-0.000000in}{0.000000in}}%
\pgfpathlineto{\pgfqpoint{-0.048611in}{0.000000in}}%
\pgfusepath{stroke,fill}%
}%
\begin{pgfscope}%
\pgfsys@transformshift{0.573768in}{1.368904in}%
\pgfsys@useobject{currentmarker}{}%
\end{pgfscope}%
\end{pgfscope}%
\begin{pgfscope}%
\definecolor{textcolor}{rgb}{0.000000,0.000000,0.000000}%
\pgfsetstrokecolor{textcolor}%
\pgfsetfillcolor{textcolor}%
\pgftext[x=0.266667in, y=1.330324in, left, base]{\color{textcolor}{\rmfamily\fontsize{8.000000}{9.600000}\selectfont\catcode`\^=\active\def^{\ifmmode\sp\else\^{}\fi}\catcode`\%=\active\def%{\%}25.0}}%
\end{pgfscope}%
\begin{pgfscope}%
\definecolor{textcolor}{rgb}{0.000000,0.000000,0.000000}%
\pgfsetstrokecolor{textcolor}%
\pgfsetfillcolor{textcolor}%
\pgftext[x=0.211111in,y=0.915779in,,bottom,rotate=90.000000]{\color{textcolor}{\rmfamily\fontsize{8.000000}{9.600000}\selectfont\catcode`\^=\active\def^{\ifmmode\sp\else\^{}\fi}\catcode`\%=\active\def%{\%}$v_o$ (V)}}%
\end{pgfscope}%
\begin{pgfscope}%
\pgfpathrectangle{\pgfqpoint{0.573768in}{0.462654in}}{\pgfqpoint{2.720000in}{0.906250in}}%
\pgfusepath{clip}%
\pgfsetrectcap%
\pgfsetroundjoin%
\pgfsetlinewidth{1.505625pt}%
\definecolor{currentstroke}{rgb}{0.000000,0.619608,0.450980}%
\pgfsetstrokecolor{currentstroke}%
\pgfsetdash{}{0pt}%
\pgfpathmoveto{\pgfqpoint{0.697405in}{0.462654in}}%
\pgfpathlineto{\pgfqpoint{0.699632in}{0.463298in}}%
\pgfpathlineto{\pgfqpoint{0.703092in}{0.474818in}}%
\pgfpathlineto{\pgfqpoint{0.704510in}{0.475520in}}%
\pgfpathlineto{\pgfqpoint{0.708053in}{0.498453in}}%
\pgfpathlineto{\pgfqpoint{0.709413in}{0.499074in}}%
\pgfpathlineto{\pgfqpoint{0.712265in}{0.533010in}}%
\pgfpathlineto{\pgfqpoint{0.712832in}{0.532929in}}%
\pgfpathlineto{\pgfqpoint{0.714306in}{0.532809in}}%
\pgfpathlineto{\pgfqpoint{0.714352in}{0.533456in}}%
\pgfpathlineto{\pgfqpoint{0.717215in}{0.577056in}}%
\pgfpathlineto{\pgfqpoint{0.717835in}{0.576913in}}%
\pgfpathlineto{\pgfqpoint{0.719233in}{0.576611in}}%
\pgfpathlineto{\pgfqpoint{0.719274in}{0.577147in}}%
\pgfpathlineto{\pgfqpoint{0.722161in}{0.629489in}}%
\pgfpathlineto{\pgfqpoint{0.722673in}{0.629316in}}%
\pgfpathlineto{\pgfqpoint{0.724178in}{0.628826in}}%
\pgfpathlineto{\pgfqpoint{0.724227in}{0.629502in}}%
\pgfpathlineto{\pgfqpoint{0.727106in}{0.688576in}}%
\pgfpathlineto{\pgfqpoint{0.727646in}{0.688329in}}%
\pgfpathlineto{\pgfqpoint{0.729124in}{0.687674in}}%
\pgfpathlineto{\pgfqpoint{0.729173in}{0.688414in}}%
\pgfpathlineto{\pgfqpoint{0.732055in}{0.752644in}}%
\pgfpathlineto{\pgfqpoint{0.732574in}{0.752339in}}%
\pgfpathlineto{\pgfqpoint{0.734070in}{0.751475in}}%
\pgfpathlineto{\pgfqpoint{0.734124in}{0.752279in}}%
\pgfpathlineto{\pgfqpoint{0.736999in}{0.819788in}}%
\pgfpathlineto{\pgfqpoint{0.737429in}{0.819478in}}%
\pgfpathlineto{\pgfqpoint{0.739015in}{0.818342in}}%
\pgfpathlineto{\pgfqpoint{0.739063in}{0.818956in}}%
\pgfpathlineto{\pgfqpoint{0.741944in}{0.888183in}}%
\pgfpathlineto{\pgfqpoint{0.742327in}{0.887852in}}%
\pgfpathlineto{\pgfqpoint{0.743961in}{0.886453in}}%
\pgfpathlineto{\pgfqpoint{0.744021in}{0.887405in}}%
\pgfpathlineto{\pgfqpoint{0.746888in}{0.955935in}}%
\pgfpathlineto{\pgfqpoint{0.747298in}{0.955527in}}%
\pgfpathlineto{\pgfqpoint{0.748905in}{0.953928in}}%
\pgfpathlineto{\pgfqpoint{0.748968in}{0.954869in}}%
\pgfpathlineto{\pgfqpoint{0.751834in}{1.021183in}}%
\pgfpathlineto{\pgfqpoint{0.752327in}{1.020626in}}%
\pgfpathlineto{\pgfqpoint{0.753852in}{1.018908in}}%
\pgfpathlineto{\pgfqpoint{0.753911in}{1.019727in}}%
\pgfpathlineto{\pgfqpoint{0.756776in}{1.082171in}}%
\pgfpathlineto{\pgfqpoint{0.757202in}{1.081638in}}%
\pgfpathlineto{\pgfqpoint{0.758797in}{1.079645in}}%
\pgfpathlineto{\pgfqpoint{0.758858in}{1.080484in}}%
\pgfpathlineto{\pgfqpoint{0.761720in}{1.137355in}}%
\pgfpathlineto{\pgfqpoint{0.762139in}{1.136784in}}%
\pgfpathlineto{\pgfqpoint{0.763742in}{1.134602in}}%
\pgfpathlineto{\pgfqpoint{0.763820in}{1.135656in}}%
\pgfpathlineto{\pgfqpoint{0.766664in}{1.185211in}}%
\pgfpathlineto{\pgfqpoint{0.767032in}{1.184675in}}%
\pgfpathlineto{\pgfqpoint{0.768701in}{1.182269in}}%
\pgfpathlineto{\pgfqpoint{0.768781in}{1.183447in}}%
\pgfpathlineto{\pgfqpoint{0.771608in}{1.224579in}}%
\pgfpathlineto{\pgfqpoint{0.771725in}{1.224399in}}%
\pgfpathlineto{\pgfqpoint{0.773647in}{1.221464in}}%
\pgfpathlineto{\pgfqpoint{0.773744in}{1.222674in}}%
\pgfpathlineto{\pgfqpoint{0.776554in}{1.254462in}}%
\pgfpathlineto{\pgfqpoint{0.776631in}{1.254338in}}%
\pgfpathlineto{\pgfqpoint{0.778593in}{1.251218in}}%
\pgfpathlineto{\pgfqpoint{0.778714in}{1.252393in}}%
\pgfpathlineto{\pgfqpoint{0.781346in}{1.273556in}}%
\pgfpathlineto{\pgfqpoint{0.781494in}{1.274176in}}%
\pgfpathlineto{\pgfqpoint{0.782048in}{1.273270in}}%
\pgfpathlineto{\pgfqpoint{0.783538in}{1.270839in}}%
\pgfpathlineto{\pgfqpoint{0.783704in}{1.271968in}}%
\pgfpathlineto{\pgfqpoint{0.785811in}{1.282545in}}%
\pgfpathlineto{\pgfqpoint{0.786434in}{1.283351in}}%
\pgfpathlineto{\pgfqpoint{0.786780in}{1.282781in}}%
\pgfpathlineto{\pgfqpoint{0.788509in}{1.279948in}}%
\pgfpathlineto{\pgfqpoint{0.788745in}{1.280885in}}%
\pgfpathlineto{\pgfqpoint{0.790235in}{1.283771in}}%
\pgfpathlineto{\pgfqpoint{0.791026in}{1.282907in}}%
\pgfpathlineto{\pgfqpoint{0.793508in}{1.278850in}}%
\pgfpathlineto{\pgfqpoint{0.793904in}{1.279531in}}%
\pgfpathlineto{\pgfqpoint{0.794695in}{1.279680in}}%
\pgfpathlineto{\pgfqpoint{0.794843in}{1.279519in}}%
\pgfpathlineto{\pgfqpoint{0.799566in}{1.274416in}}%
\pgfpathlineto{\pgfqpoint{0.801001in}{1.272199in}}%
\pgfpathlineto{\pgfqpoint{0.803400in}{1.268310in}}%
\pgfpathlineto{\pgfqpoint{0.803850in}{1.268999in}}%
\pgfpathlineto{\pgfqpoint{0.804641in}{1.268932in}}%
\pgfpathlineto{\pgfqpoint{0.804740in}{1.268806in}}%
\pgfpathlineto{\pgfqpoint{0.808347in}{1.262987in}}%
\pgfpathlineto{\pgfqpoint{0.809023in}{1.263828in}}%
\pgfpathlineto{\pgfqpoint{0.809814in}{1.263304in}}%
\pgfpathlineto{\pgfqpoint{0.813285in}{1.257725in}}%
\pgfpathlineto{\pgfqpoint{0.813692in}{1.258366in}}%
\pgfpathlineto{\pgfqpoint{0.814484in}{1.258450in}}%
\pgfpathlineto{\pgfqpoint{0.814632in}{1.258282in}}%
\pgfpathlineto{\pgfqpoint{0.818239in}{1.252557in}}%
\pgfpathlineto{\pgfqpoint{0.819014in}{1.253479in}}%
\pgfpathlineto{\pgfqpoint{0.819805in}{1.252830in}}%
\pgfpathlineto{\pgfqpoint{0.823185in}{1.247495in}}%
\pgfpathlineto{\pgfqpoint{0.823578in}{1.248169in}}%
\pgfpathlineto{\pgfqpoint{0.824370in}{1.248323in}}%
\pgfpathlineto{\pgfqpoint{0.824518in}{1.248170in}}%
\pgfpathlineto{\pgfqpoint{0.828126in}{1.242516in}}%
\pgfpathlineto{\pgfqpoint{0.829147in}{1.243461in}}%
\pgfpathlineto{\pgfqpoint{0.830086in}{1.242270in}}%
\pgfpathlineto{\pgfqpoint{0.833066in}{1.237620in}}%
\pgfpathlineto{\pgfqpoint{0.833445in}{1.238287in}}%
\pgfpathlineto{\pgfqpoint{0.834286in}{1.238522in}}%
\pgfpathlineto{\pgfqpoint{0.834434in}{1.238374in}}%
\pgfpathlineto{\pgfqpoint{0.838026in}{1.232852in}}%
\pgfpathlineto{\pgfqpoint{0.838942in}{1.233915in}}%
\pgfpathlineto{\pgfqpoint{0.839734in}{1.233151in}}%
\pgfpathlineto{\pgfqpoint{0.842969in}{1.228190in}}%
\pgfpathlineto{\pgfqpoint{0.843362in}{1.228886in}}%
\pgfpathlineto{\pgfqpoint{0.844215in}{1.229132in}}%
\pgfpathlineto{\pgfqpoint{0.844364in}{1.228986in}}%
\pgfpathlineto{\pgfqpoint{0.847896in}{1.223594in}}%
\pgfpathlineto{\pgfqpoint{0.848756in}{1.224755in}}%
\pgfpathlineto{\pgfqpoint{0.849547in}{1.224192in}}%
\pgfpathlineto{\pgfqpoint{0.852841in}{1.219154in}}%
\pgfpathlineto{\pgfqpoint{0.853205in}{1.219855in}}%
\pgfpathlineto{\pgfqpoint{0.854080in}{1.220281in}}%
\pgfpathlineto{\pgfqpoint{0.854278in}{1.220111in}}%
\pgfpathlineto{\pgfqpoint{0.857806in}{1.214830in}}%
\pgfpathlineto{\pgfqpoint{0.858696in}{1.216056in}}%
\pgfpathlineto{\pgfqpoint{0.859487in}{1.215495in}}%
\pgfpathlineto{\pgfqpoint{0.862727in}{1.210593in}}%
\pgfpathlineto{\pgfqpoint{0.863084in}{1.211317in}}%
\pgfpathlineto{\pgfqpoint{0.864004in}{1.211843in}}%
\pgfpathlineto{\pgfqpoint{0.864152in}{1.211729in}}%
\pgfpathlineto{\pgfqpoint{0.866381in}{1.208426in}}%
\pgfpathlineto{\pgfqpoint{0.867688in}{1.206509in}}%
\pgfpathlineto{\pgfqpoint{0.868043in}{1.207245in}}%
\pgfpathlineto{\pgfqpoint{0.868982in}{1.207782in}}%
\pgfpathlineto{\pgfqpoint{0.869131in}{1.207667in}}%
\pgfpathlineto{\pgfqpoint{0.871531in}{1.204121in}}%
\pgfpathlineto{\pgfqpoint{0.872619in}{1.202515in}}%
\pgfpathlineto{\pgfqpoint{0.872976in}{1.203284in}}%
\pgfpathlineto{\pgfqpoint{0.873915in}{1.203912in}}%
\pgfpathlineto{\pgfqpoint{0.874113in}{1.203766in}}%
\pgfpathlineto{\pgfqpoint{0.876638in}{1.200046in}}%
\pgfpathlineto{\pgfqpoint{0.877574in}{1.198692in}}%
\pgfpathlineto{\pgfqpoint{0.877905in}{1.199425in}}%
\pgfpathlineto{\pgfqpoint{0.878894in}{1.200120in}}%
\pgfpathlineto{\pgfqpoint{0.879043in}{1.200018in}}%
\pgfpathlineto{\pgfqpoint{0.881084in}{1.197057in}}%
\pgfpathlineto{\pgfqpoint{0.882501in}{1.194973in}}%
\pgfpathlineto{\pgfqpoint{0.882861in}{1.195779in}}%
\pgfpathlineto{\pgfqpoint{0.883866in}{1.196493in}}%
\pgfpathlineto{\pgfqpoint{0.884014in}{1.196390in}}%
\pgfpathlineto{\pgfqpoint{0.886094in}{1.193385in}}%
\pgfpathlineto{\pgfqpoint{0.887470in}{1.191415in}}%
\pgfpathlineto{\pgfqpoint{0.887810in}{1.192196in}}%
\pgfpathlineto{\pgfqpoint{0.888842in}{1.192940in}}%
\pgfpathlineto{\pgfqpoint{0.888941in}{1.192875in}}%
\pgfpathlineto{\pgfqpoint{0.890665in}{1.190435in}}%
\pgfpathlineto{\pgfqpoint{0.892402in}{1.187918in}}%
\pgfpathlineto{\pgfqpoint{0.892707in}{1.188657in}}%
\pgfpathlineto{\pgfqpoint{0.893795in}{1.189551in}}%
\pgfpathlineto{\pgfqpoint{0.893943in}{1.189451in}}%
\pgfpathlineto{\pgfqpoint{0.895986in}{1.186534in}}%
\pgfpathlineto{\pgfqpoint{0.897350in}{1.184576in}}%
\pgfpathlineto{\pgfqpoint{0.897656in}{1.185318in}}%
\pgfpathlineto{\pgfqpoint{0.898744in}{1.186243in}}%
\pgfpathlineto{\pgfqpoint{0.898893in}{1.186148in}}%
\pgfpathlineto{\pgfqpoint{0.900811in}{1.183434in}}%
\pgfpathlineto{\pgfqpoint{0.902301in}{1.181328in}}%
\pgfpathlineto{\pgfqpoint{0.902629in}{1.182119in}}%
\pgfpathlineto{\pgfqpoint{0.903704in}{1.183018in}}%
\pgfpathlineto{\pgfqpoint{0.903852in}{1.182924in}}%
\pgfpathlineto{\pgfqpoint{0.905734in}{1.180278in}}%
\pgfpathlineto{\pgfqpoint{0.907229in}{1.178138in}}%
\pgfpathlineto{\pgfqpoint{0.907564in}{1.178966in}}%
\pgfpathlineto{\pgfqpoint{0.908680in}{1.179899in}}%
\pgfpathlineto{\pgfqpoint{0.908779in}{1.179838in}}%
\pgfpathlineto{\pgfqpoint{0.910431in}{1.177554in}}%
\pgfpathlineto{\pgfqpoint{0.912184in}{1.175058in}}%
\pgfpathlineto{\pgfqpoint{0.912486in}{1.175811in}}%
\pgfpathlineto{\pgfqpoint{0.913623in}{1.176828in}}%
\pgfpathlineto{\pgfqpoint{0.913722in}{1.176771in}}%
\pgfpathlineto{\pgfqpoint{0.915311in}{1.174600in}}%
\pgfpathlineto{\pgfqpoint{0.917132in}{1.172023in}}%
\pgfpathlineto{\pgfqpoint{0.917453in}{1.172821in}}%
\pgfpathlineto{\pgfqpoint{0.918591in}{1.173809in}}%
\pgfpathlineto{\pgfqpoint{0.918690in}{1.173750in}}%
\pgfpathlineto{\pgfqpoint{0.920315in}{1.171530in}}%
\pgfpathlineto{\pgfqpoint{0.922066in}{1.169051in}}%
\pgfpathlineto{\pgfqpoint{0.922356in}{1.169795in}}%
\pgfpathlineto{\pgfqpoint{0.923537in}{1.170896in}}%
\pgfpathlineto{\pgfqpoint{0.923636in}{1.170839in}}%
\pgfpathlineto{\pgfqpoint{0.925215in}{1.168699in}}%
\pgfpathlineto{\pgfqpoint{0.927021in}{1.166157in}}%
\pgfpathlineto{\pgfqpoint{0.927335in}{1.166948in}}%
\pgfpathlineto{\pgfqpoint{0.928472in}{1.168012in}}%
\pgfpathlineto{\pgfqpoint{0.928620in}{1.167926in}}%
\pgfpathlineto{\pgfqpoint{0.930363in}{1.165544in}}%
\pgfpathlineto{\pgfqpoint{0.931966in}{1.163311in}}%
\pgfpathlineto{\pgfqpoint{0.932261in}{1.164077in}}%
\pgfpathlineto{\pgfqpoint{0.933448in}{1.165193in}}%
\pgfpathlineto{\pgfqpoint{0.933547in}{1.165137in}}%
\pgfpathlineto{\pgfqpoint{0.935118in}{1.163025in}}%
\pgfpathlineto{\pgfqpoint{0.936911in}{1.160526in}}%
\pgfpathlineto{\pgfqpoint{0.937211in}{1.161309in}}%
\pgfpathlineto{\pgfqpoint{0.938398in}{1.162430in}}%
\pgfpathlineto{\pgfqpoint{0.938497in}{1.162376in}}%
\pgfpathlineto{\pgfqpoint{0.940035in}{1.160322in}}%
\pgfpathlineto{\pgfqpoint{0.941856in}{1.157798in}}%
\pgfpathlineto{\pgfqpoint{0.942141in}{1.158539in}}%
\pgfpathlineto{\pgfqpoint{0.943328in}{1.159715in}}%
\pgfpathlineto{\pgfqpoint{0.943427in}{1.159666in}}%
\pgfpathlineto{\pgfqpoint{0.944868in}{1.157773in}}%
\pgfpathlineto{\pgfqpoint{0.946799in}{1.155102in}}%
\pgfpathlineto{\pgfqpoint{0.947089in}{1.155857in}}%
\pgfpathlineto{\pgfqpoint{0.948325in}{1.157025in}}%
\pgfpathlineto{\pgfqpoint{0.948375in}{1.156998in}}%
\pgfpathlineto{\pgfqpoint{0.949786in}{1.155160in}}%
\pgfpathlineto{\pgfqpoint{0.951747in}{1.152464in}}%
\pgfpathlineto{\pgfqpoint{0.952058in}{1.153272in}}%
\pgfpathlineto{\pgfqpoint{0.953245in}{1.154410in}}%
\pgfpathlineto{\pgfqpoint{0.953343in}{1.154359in}}%
\pgfpathlineto{\pgfqpoint{0.954857in}{1.152371in}}%
\pgfpathlineto{\pgfqpoint{0.956677in}{1.149852in}}%
\pgfpathlineto{\pgfqpoint{0.956976in}{1.150604in}}%
\pgfpathlineto{\pgfqpoint{0.958212in}{1.151815in}}%
\pgfpathlineto{\pgfqpoint{0.958262in}{1.151790in}}%
\pgfpathlineto{\pgfqpoint{0.959609in}{1.150069in}}%
\pgfpathlineto{\pgfqpoint{0.961624in}{1.147291in}}%
\pgfpathlineto{\pgfqpoint{0.961925in}{1.148064in}}%
\pgfpathlineto{\pgfqpoint{0.963161in}{1.149278in}}%
\pgfpathlineto{\pgfqpoint{0.963211in}{1.149254in}}%
\pgfpathlineto{\pgfqpoint{0.964543in}{1.147563in}}%
\pgfpathlineto{\pgfqpoint{0.966565in}{1.144782in}}%
\pgfpathlineto{\pgfqpoint{0.966882in}{1.145586in}}%
\pgfpathlineto{\pgfqpoint{0.968118in}{1.146786in}}%
\pgfpathlineto{\pgfqpoint{0.968168in}{1.146761in}}%
\pgfpathlineto{\pgfqpoint{0.969526in}{1.145037in}}%
\pgfpathlineto{\pgfqpoint{0.971514in}{1.142316in}}%
\pgfpathlineto{\pgfqpoint{0.971800in}{1.143050in}}%
\pgfpathlineto{\pgfqpoint{0.973037in}{1.144337in}}%
\pgfpathlineto{\pgfqpoint{0.973136in}{1.144290in}}%
\pgfpathlineto{\pgfqpoint{0.974563in}{1.142462in}}%
\pgfpathlineto{\pgfqpoint{0.976464in}{1.139881in}}%
\pgfpathlineto{\pgfqpoint{0.976749in}{1.140630in}}%
\pgfpathlineto{\pgfqpoint{0.978028in}{1.141906in}}%
\pgfpathlineto{\pgfqpoint{0.978077in}{1.141881in}}%
\pgfpathlineto{\pgfqpoint{0.979466in}{1.140122in}}%
\pgfpathlineto{\pgfqpoint{0.981412in}{1.137484in}}%
\pgfpathlineto{\pgfqpoint{0.981729in}{1.138310in}}%
\pgfpathlineto{\pgfqpoint{0.982965in}{1.139526in}}%
\pgfpathlineto{\pgfqpoint{0.983015in}{1.139504in}}%
\pgfpathlineto{\pgfqpoint{0.984359in}{1.137825in}}%
\pgfpathlineto{\pgfqpoint{0.986366in}{1.135134in}}%
\pgfpathlineto{\pgfqpoint{0.986641in}{1.135869in}}%
\pgfpathlineto{\pgfqpoint{0.987927in}{1.137171in}}%
\pgfpathlineto{\pgfqpoint{0.987977in}{1.137146in}}%
\pgfpathlineto{\pgfqpoint{0.989354in}{1.135421in}}%
\pgfpathlineto{\pgfqpoint{0.991312in}{1.132808in}}%
\pgfpathlineto{\pgfqpoint{0.991585in}{1.133529in}}%
\pgfpathlineto{\pgfqpoint{0.992871in}{1.134846in}}%
\pgfpathlineto{\pgfqpoint{0.992920in}{1.134822in}}%
\pgfpathlineto{\pgfqpoint{0.994250in}{1.133172in}}%
\pgfpathlineto{\pgfqpoint{0.996237in}{1.130500in}}%
\pgfpathlineto{\pgfqpoint{0.996543in}{1.131295in}}%
\pgfpathlineto{\pgfqpoint{0.997828in}{1.132585in}}%
\pgfpathlineto{\pgfqpoint{0.997878in}{1.132560in}}%
\pgfpathlineto{\pgfqpoint{0.999258in}{1.130841in}}%
\pgfpathlineto{\pgfqpoint{1.001187in}{1.128256in}}%
\pgfpathlineto{\pgfqpoint{1.001478in}{1.129005in}}%
\pgfpathlineto{\pgfqpoint{1.002763in}{1.130335in}}%
\pgfpathlineto{\pgfqpoint{1.002813in}{1.130313in}}%
\pgfpathlineto{\pgfqpoint{1.004135in}{1.128691in}}%
\pgfpathlineto{\pgfqpoint{1.006134in}{1.126027in}}%
\pgfpathlineto{\pgfqpoint{1.006414in}{1.126772in}}%
\pgfpathlineto{\pgfqpoint{1.007743in}{1.128120in}}%
\pgfpathlineto{\pgfqpoint{1.009000in}{1.126608in}}%
\pgfpathlineto{\pgfqpoint{1.011077in}{1.123839in}}%
\pgfpathlineto{\pgfqpoint{1.011379in}{1.124617in}}%
\pgfpathlineto{\pgfqpoint{1.012665in}{1.125940in}}%
\pgfpathlineto{\pgfqpoint{1.012715in}{1.125918in}}%
\pgfpathlineto{\pgfqpoint{1.014051in}{1.124287in}}%
\pgfpathlineto{\pgfqpoint{1.016026in}{1.121673in}}%
\pgfpathlineto{\pgfqpoint{1.016342in}{1.122498in}}%
\pgfpathlineto{\pgfqpoint{1.017643in}{1.123771in}}%
\pgfpathlineto{\pgfqpoint{1.017693in}{1.123746in}}%
\pgfpathlineto{\pgfqpoint{1.019152in}{1.121933in}}%
\pgfpathlineto{\pgfqpoint{1.020970in}{1.119534in}}%
\pgfpathlineto{\pgfqpoint{1.021282in}{1.120347in}}%
\pgfpathlineto{\pgfqpoint{1.022558in}{1.121656in}}%
\pgfpathlineto{\pgfqpoint{1.022608in}{1.121635in}}%
\pgfpathlineto{\pgfqpoint{1.023922in}{1.120049in}}%
\pgfpathlineto{\pgfqpoint{1.025916in}{1.117427in}}%
\pgfpathlineto{\pgfqpoint{1.026190in}{1.118152in}}%
\pgfpathlineto{\pgfqpoint{1.027523in}{1.119557in}}%
\pgfpathlineto{\pgfqpoint{1.027572in}{1.119535in}}%
\pgfpathlineto{\pgfqpoint{1.028934in}{1.117881in}}%
\pgfpathlineto{\pgfqpoint{1.030870in}{1.115365in}}%
\pgfpathlineto{\pgfqpoint{1.031161in}{1.116149in}}%
\pgfpathlineto{\pgfqpoint{1.032497in}{1.117485in}}%
\pgfpathlineto{\pgfqpoint{1.033777in}{1.115963in}}%
\pgfpathlineto{\pgfqpoint{1.035811in}{1.113314in}}%
\pgfpathlineto{\pgfqpoint{1.036090in}{1.114065in}}%
\pgfpathlineto{\pgfqpoint{1.037416in}{1.115463in}}%
\pgfpathlineto{\pgfqpoint{1.037466in}{1.115441in}}%
\pgfpathlineto{\pgfqpoint{1.038808in}{1.113828in}}%
\pgfpathlineto{\pgfqpoint{1.040753in}{1.111294in}}%
\pgfpathlineto{\pgfqpoint{1.041063in}{1.112108in}}%
\pgfpathlineto{\pgfqpoint{1.042359in}{1.113447in}}%
\pgfpathlineto{\pgfqpoint{1.042409in}{1.113426in}}%
\pgfpathlineto{\pgfqpoint{1.043712in}{1.111874in}}%
\pgfpathlineto{\pgfqpoint{1.045700in}{1.109295in}}%
\pgfpathlineto{\pgfqpoint{1.046006in}{1.110106in}}%
\pgfpathlineto{\pgfqpoint{1.047340in}{1.111447in}}%
\pgfpathlineto{\pgfqpoint{1.048589in}{1.109987in}}%
\pgfpathlineto{\pgfqpoint{1.050643in}{1.107324in}}%
\pgfpathlineto{\pgfqpoint{1.050952in}{1.108129in}}%
\pgfpathlineto{\pgfqpoint{1.052285in}{1.109479in}}%
\pgfpathlineto{\pgfqpoint{1.052334in}{1.109455in}}%
\pgfpathlineto{\pgfqpoint{1.053753in}{1.107742in}}%
\pgfpathlineto{\pgfqpoint{1.055582in}{1.105370in}}%
\pgfpathlineto{\pgfqpoint{1.055908in}{1.106182in}}%
\pgfpathlineto{\pgfqpoint{1.057238in}{1.107515in}}%
\pgfpathlineto{\pgfqpoint{1.058495in}{1.106056in}}%
\pgfpathlineto{\pgfqpoint{1.060529in}{1.103434in}}%
\pgfpathlineto{\pgfqpoint{1.060845in}{1.104244in}}%
\pgfpathlineto{\pgfqpoint{1.062174in}{1.105608in}}%
\pgfpathlineto{\pgfqpoint{1.062224in}{1.105586in}}%
\pgfpathlineto{\pgfqpoint{1.063603in}{1.103943in}}%
\pgfpathlineto{\pgfqpoint{1.065475in}{1.101536in}}%
\pgfpathlineto{\pgfqpoint{1.065794in}{1.102359in}}%
\pgfpathlineto{\pgfqpoint{1.067114in}{1.103724in}}%
\pgfpathlineto{\pgfqpoint{1.067164in}{1.103703in}}%
\pgfpathlineto{\pgfqpoint{1.068502in}{1.102126in}}%
\pgfpathlineto{\pgfqpoint{1.070420in}{1.099664in}}%
\pgfpathlineto{\pgfqpoint{1.070742in}{1.100490in}}%
\pgfpathlineto{\pgfqpoint{1.072096in}{1.101842in}}%
\pgfpathlineto{\pgfqpoint{1.073375in}{1.100360in}}%
\pgfpathlineto{\pgfqpoint{1.075368in}{1.097814in}}%
\pgfpathlineto{\pgfqpoint{1.075675in}{1.098604in}}%
\pgfpathlineto{\pgfqpoint{1.077034in}{1.099993in}}%
\pgfpathlineto{\pgfqpoint{1.078293in}{1.098550in}}%
\pgfpathlineto{\pgfqpoint{1.080310in}{1.095976in}}%
\pgfpathlineto{\pgfqpoint{1.080633in}{1.096795in}}%
\pgfpathlineto{\pgfqpoint{1.081978in}{1.098171in}}%
\pgfpathlineto{\pgfqpoint{1.083206in}{1.096778in}}%
\pgfpathlineto{\pgfqpoint{1.085256in}{1.094172in}}%
\pgfpathlineto{\pgfqpoint{1.085571in}{1.094978in}}%
\pgfpathlineto{\pgfqpoint{1.086932in}{1.096372in}}%
\pgfpathlineto{\pgfqpoint{1.088179in}{1.094953in}}%
\pgfpathlineto{\pgfqpoint{1.090205in}{1.092393in}}%
\pgfpathlineto{\pgfqpoint{1.090515in}{1.093207in}}%
\pgfpathlineto{\pgfqpoint{1.091864in}{1.094613in}}%
\pgfpathlineto{\pgfqpoint{1.091914in}{1.094592in}}%
\pgfpathlineto{\pgfqpoint{1.093280in}{1.092998in}}%
\pgfpathlineto{\pgfqpoint{1.095152in}{1.090640in}}%
\pgfpathlineto{\pgfqpoint{1.095456in}{1.091440in}}%
\pgfpathlineto{\pgfqpoint{1.096835in}{1.092854in}}%
\pgfpathlineto{\pgfqpoint{1.098115in}{1.091396in}}%
\pgfpathlineto{\pgfqpoint{1.100101in}{1.088912in}}%
\pgfpathlineto{\pgfqpoint{1.100406in}{1.089723in}}%
\pgfpathlineto{\pgfqpoint{1.101758in}{1.091141in}}%
\pgfpathlineto{\pgfqpoint{1.101807in}{1.091121in}}%
\pgfpathlineto{\pgfqpoint{1.103176in}{1.089536in}}%
\pgfpathlineto{\pgfqpoint{1.105037in}{1.087193in}}%
\pgfpathlineto{\pgfqpoint{1.105358in}{1.088008in}}%
\pgfpathlineto{\pgfqpoint{1.106691in}{1.089429in}}%
\pgfpathlineto{\pgfqpoint{1.107829in}{1.088199in}}%
\pgfpathlineto{\pgfqpoint{1.109989in}{1.085503in}}%
\pgfpathlineto{\pgfqpoint{1.110298in}{1.086319in}}%
\pgfpathlineto{\pgfqpoint{1.111661in}{1.087740in}}%
\pgfpathlineto{\pgfqpoint{1.112877in}{1.086393in}}%
\pgfpathlineto{\pgfqpoint{1.114933in}{1.083828in}}%
\pgfpathlineto{\pgfqpoint{1.115238in}{1.084629in}}%
\pgfpathlineto{\pgfqpoint{1.116630in}{1.086065in}}%
\pgfpathlineto{\pgfqpoint{1.117898in}{1.084640in}}%
\pgfpathlineto{\pgfqpoint{1.119883in}{1.082186in}}%
\pgfpathlineto{\pgfqpoint{1.120180in}{1.082978in}}%
\pgfpathlineto{\pgfqpoint{1.121589in}{1.084421in}}%
\pgfpathlineto{\pgfqpoint{1.122918in}{1.082913in}}%
\pgfpathlineto{\pgfqpoint{1.124825in}{1.080552in}}%
\pgfpathlineto{\pgfqpoint{1.125131in}{1.081357in}}%
\pgfpathlineto{\pgfqpoint{1.126512in}{1.082804in}}%
\pgfpathlineto{\pgfqpoint{1.127752in}{1.081435in}}%
\pgfpathlineto{\pgfqpoint{1.129772in}{1.078946in}}%
\pgfpathlineto{\pgfqpoint{1.130075in}{1.079745in}}%
\pgfpathlineto{\pgfqpoint{1.131473in}{1.081193in}}%
\pgfpathlineto{\pgfqpoint{1.132746in}{1.079775in}}%
\pgfpathlineto{\pgfqpoint{1.134717in}{1.077353in}}%
\pgfpathlineto{\pgfqpoint{1.135020in}{1.078157in}}%
\pgfpathlineto{\pgfqpoint{1.136404in}{1.079619in}}%
\pgfpathlineto{\pgfqpoint{1.137636in}{1.078271in}}%
\pgfpathlineto{\pgfqpoint{1.139653in}{1.075784in}}%
\pgfpathlineto{\pgfqpoint{1.139956in}{1.076588in}}%
\pgfpathlineto{\pgfqpoint{1.141390in}{1.078069in}}%
\pgfpathlineto{\pgfqpoint{1.142724in}{1.076569in}}%
\pgfpathlineto{\pgfqpoint{1.144603in}{1.074257in}}%
\pgfpathlineto{\pgfqpoint{1.144874in}{1.074978in}}%
\pgfpathlineto{\pgfqpoint{1.146308in}{1.076542in}}%
\pgfpathlineto{\pgfqpoint{1.147560in}{1.075171in}}%
\pgfpathlineto{\pgfqpoint{1.149548in}{1.072730in}}%
\pgfpathlineto{\pgfqpoint{1.149849in}{1.073517in}}%
\pgfpathlineto{\pgfqpoint{1.151283in}{1.075001in}}%
\pgfpathlineto{\pgfqpoint{1.152627in}{1.073499in}}%
\pgfpathlineto{\pgfqpoint{1.154495in}{1.071219in}}%
\pgfpathlineto{\pgfqpoint{1.154786in}{1.072000in}}%
\pgfpathlineto{\pgfqpoint{1.156220in}{1.073513in}}%
\pgfpathlineto{\pgfqpoint{1.157516in}{1.072082in}}%
\pgfpathlineto{\pgfqpoint{1.159440in}{1.069733in}}%
\pgfpathlineto{\pgfqpoint{1.159712in}{1.070458in}}%
\pgfpathlineto{\pgfqpoint{1.161146in}{1.072037in}}%
\pgfpathlineto{\pgfqpoint{1.162363in}{1.070726in}}%
\pgfpathlineto{\pgfqpoint{1.164385in}{1.068263in}}%
\pgfpathlineto{\pgfqpoint{1.164659in}{1.068993in}}%
\pgfpathlineto{\pgfqpoint{1.166093in}{1.070570in}}%
\pgfpathlineto{\pgfqpoint{1.167337in}{1.069231in}}%
\pgfpathlineto{\pgfqpoint{1.169331in}{1.066808in}}%
\pgfpathlineto{\pgfqpoint{1.169616in}{1.067569in}}%
\pgfpathlineto{\pgfqpoint{1.171050in}{1.069119in}}%
\pgfpathlineto{\pgfqpoint{1.172307in}{1.067761in}}%
\pgfpathlineto{\pgfqpoint{1.174276in}{1.065373in}}%
\pgfpathlineto{\pgfqpoint{1.174545in}{1.066091in}}%
\pgfpathlineto{\pgfqpoint{1.175979in}{1.067692in}}%
\pgfpathlineto{\pgfqpoint{1.177171in}{1.066432in}}%
\pgfpathlineto{\pgfqpoint{1.179222in}{1.063951in}}%
\pgfpathlineto{\pgfqpoint{1.179493in}{1.064677in}}%
\pgfpathlineto{\pgfqpoint{1.180928in}{1.066278in}}%
\pgfpathlineto{\pgfqpoint{1.182127in}{1.065012in}}%
\pgfpathlineto{\pgfqpoint{1.184167in}{1.062550in}}%
\pgfpathlineto{\pgfqpoint{1.184444in}{1.063288in}}%
\pgfpathlineto{\pgfqpoint{1.185878in}{1.064878in}}%
\pgfpathlineto{\pgfqpoint{1.187097in}{1.063589in}}%
\pgfpathlineto{\pgfqpoint{1.189113in}{1.061163in}}%
\pgfpathlineto{\pgfqpoint{1.189398in}{1.061924in}}%
\pgfpathlineto{\pgfqpoint{1.190832in}{1.063494in}}%
\pgfpathlineto{\pgfqpoint{1.192074in}{1.062177in}}%
\pgfpathlineto{\pgfqpoint{1.194053in}{1.059796in}}%
\pgfpathlineto{\pgfqpoint{1.194317in}{1.060493in}}%
\pgfpathlineto{\pgfqpoint{1.195751in}{1.062144in}}%
\pgfpathlineto{\pgfqpoint{1.196923in}{1.060940in}}%
\pgfpathlineto{\pgfqpoint{1.199009in}{1.058457in}}%
\pgfpathlineto{\pgfqpoint{1.199308in}{1.059257in}}%
\pgfpathlineto{\pgfqpoint{1.200743in}{1.060776in}}%
\pgfpathlineto{\pgfqpoint{1.202032in}{1.059395in}}%
\pgfpathlineto{\pgfqpoint{1.203944in}{1.057109in}}%
\pgfpathlineto{\pgfqpoint{1.204245in}{1.057900in}}%
\pgfpathlineto{\pgfqpoint{1.205642in}{1.059472in}}%
\pgfpathlineto{\pgfqpoint{1.206799in}{1.058298in}}%
\pgfpathlineto{\pgfqpoint{1.208910in}{1.055813in}}%
\pgfpathlineto{\pgfqpoint{1.209190in}{1.056562in}}%
\pgfpathlineto{\pgfqpoint{1.210601in}{1.058131in}}%
\pgfpathlineto{\pgfqpoint{1.211790in}{1.056910in}}%
\pgfpathlineto{\pgfqpoint{1.213835in}{1.054470in}}%
\pgfpathlineto{\pgfqpoint{1.214146in}{1.055272in}}%
\pgfpathlineto{\pgfqpoint{1.215598in}{1.056804in}}%
\pgfpathlineto{\pgfqpoint{1.216904in}{1.055405in}}%
\pgfpathlineto{\pgfqpoint{1.218781in}{1.053182in}}%
\pgfpathlineto{\pgfqpoint{1.219057in}{1.053923in}}%
\pgfpathlineto{\pgfqpoint{1.220481in}{1.055562in}}%
\pgfpathlineto{\pgfqpoint{1.221635in}{1.054403in}}%
\pgfpathlineto{\pgfqpoint{1.223743in}{1.051927in}}%
\pgfpathlineto{\pgfqpoint{1.224012in}{1.052662in}}%
\pgfpathlineto{\pgfqpoint{1.225446in}{1.054280in}}%
\pgfpathlineto{\pgfqpoint{1.226640in}{1.053062in}}%
\pgfpathlineto{\pgfqpoint{1.228675in}{1.050659in}}%
\pgfpathlineto{\pgfqpoint{1.228981in}{1.051460in}}%
\pgfpathlineto{\pgfqpoint{1.230372in}{1.053027in}}%
\pgfpathlineto{\pgfqpoint{1.231518in}{1.051888in}}%
\pgfpathlineto{\pgfqpoint{1.233618in}{1.049411in}}%
\pgfpathlineto{\pgfqpoint{1.233930in}{1.050209in}}%
\pgfpathlineto{\pgfqpoint{1.235357in}{1.051759in}}%
\pgfpathlineto{\pgfqpoint{1.236589in}{1.050488in}}%
\pgfpathlineto{\pgfqpoint{1.238562in}{1.048166in}}%
\pgfpathlineto{\pgfqpoint{1.238845in}{1.048915in}}%
\pgfpathlineto{\pgfqpoint{1.240279in}{1.050555in}}%
\pgfpathlineto{\pgfqpoint{1.241446in}{1.049387in}}%
\pgfpathlineto{\pgfqpoint{1.243510in}{1.046959in}}%
\pgfpathlineto{\pgfqpoint{1.243822in}{1.047745in}}%
\pgfpathlineto{\pgfqpoint{1.245226in}{1.049303in}}%
\pgfpathlineto{\pgfqpoint{1.246402in}{1.048125in}}%
\pgfpathlineto{\pgfqpoint{1.248454in}{1.045730in}}%
\pgfpathlineto{\pgfqpoint{1.248754in}{1.046497in}}%
\pgfpathlineto{\pgfqpoint{1.250134in}{1.048110in}}%
\pgfpathlineto{\pgfqpoint{1.251222in}{1.047080in}}%
\pgfpathlineto{\pgfqpoint{1.253403in}{1.044539in}}%
\pgfpathlineto{\pgfqpoint{1.253679in}{1.045241in}}%
\pgfpathlineto{\pgfqpoint{1.255168in}{1.046870in}}%
\pgfpathlineto{\pgfqpoint{1.256454in}{1.045527in}}%
\pgfpathlineto{\pgfqpoint{1.258345in}{1.043332in}}%
\pgfpathlineto{\pgfqpoint{1.258652in}{1.044110in}}%
\pgfpathlineto{\pgfqpoint{1.260055in}{1.045709in}}%
\pgfpathlineto{\pgfqpoint{1.261199in}{1.044594in}}%
\pgfpathlineto{\pgfqpoint{1.263287in}{1.042163in}}%
\pgfpathlineto{\pgfqpoint{1.263544in}{1.042852in}}%
\pgfpathlineto{\pgfqpoint{1.265028in}{1.044572in}}%
\pgfpathlineto{\pgfqpoint{1.266208in}{1.043391in}}%
\pgfpathlineto{\pgfqpoint{1.268243in}{1.041028in}}%
\pgfpathlineto{\pgfqpoint{1.268530in}{1.041796in}}%
\pgfpathlineto{\pgfqpoint{1.269965in}{1.043428in}}%
\pgfpathlineto{\pgfqpoint{1.271131in}{1.042278in}}%
\pgfpathlineto{\pgfqpoint{1.273180in}{1.039897in}}%
\pgfpathlineto{\pgfqpoint{1.273479in}{1.040676in}}%
\pgfpathlineto{\pgfqpoint{1.274911in}{1.042299in}}%
\pgfpathlineto{\pgfqpoint{1.276087in}{1.041140in}}%
\pgfpathlineto{\pgfqpoint{1.278136in}{1.038784in}}%
\pgfpathlineto{\pgfqpoint{1.278438in}{1.039574in}}%
\pgfpathlineto{\pgfqpoint{1.279915in}{1.041139in}}%
\pgfpathlineto{\pgfqpoint{1.281233in}{1.039766in}}%
\pgfpathlineto{\pgfqpoint{1.283079in}{1.037659in}}%
\pgfpathlineto{\pgfqpoint{1.283360in}{1.038430in}}%
\pgfpathlineto{\pgfqpoint{1.284842in}{1.040071in}}%
\pgfpathlineto{\pgfqpoint{1.286106in}{1.038782in}}%
\pgfpathlineto{\pgfqpoint{1.288025in}{1.036578in}}%
\pgfpathlineto{\pgfqpoint{1.288316in}{1.037353in}}%
\pgfpathlineto{\pgfqpoint{1.289793in}{1.038967in}}%
\pgfpathlineto{\pgfqpoint{1.291048in}{1.037688in}}%
\pgfpathlineto{\pgfqpoint{1.292964in}{1.035486in}}%
\pgfpathlineto{\pgfqpoint{1.293265in}{1.036291in}}%
\pgfpathlineto{\pgfqpoint{1.294753in}{1.037906in}}%
\pgfpathlineto{\pgfqpoint{1.296045in}{1.036578in}}%
\pgfpathlineto{\pgfqpoint{1.297914in}{1.034433in}}%
\pgfpathlineto{\pgfqpoint{1.298220in}{1.035209in}}%
\pgfpathlineto{\pgfqpoint{1.299662in}{1.036798in}}%
\pgfpathlineto{\pgfqpoint{1.300867in}{1.035604in}}%
\pgfpathlineto{\pgfqpoint{1.302859in}{1.033347in}}%
\pgfpathlineto{\pgfqpoint{1.303132in}{1.034098in}}%
\pgfpathlineto{\pgfqpoint{1.304585in}{1.035788in}}%
\pgfpathlineto{\pgfqpoint{1.305739in}{1.034679in}}%
\pgfpathlineto{\pgfqpoint{1.307801in}{1.032315in}}%
\pgfpathlineto{\pgfqpoint{1.308114in}{1.033108in}}%
\pgfpathlineto{\pgfqpoint{1.309577in}{1.034699in}}%
\pgfpathlineto{\pgfqpoint{1.310832in}{1.033435in}}%
\pgfpathlineto{\pgfqpoint{1.312751in}{1.031262in}}%
\pgfpathlineto{\pgfqpoint{1.313002in}{1.031952in}}%
\pgfpathlineto{\pgfqpoint{1.314485in}{1.033706in}}%
\pgfpathlineto{\pgfqpoint{1.315651in}{1.032584in}}%
\pgfpathlineto{\pgfqpoint{1.317693in}{1.030253in}}%
\pgfpathlineto{\pgfqpoint{1.317985in}{1.031023in}}%
\pgfpathlineto{\pgfqpoint{1.319469in}{1.032667in}}%
\pgfpathlineto{\pgfqpoint{1.320718in}{1.031417in}}%
\pgfpathlineto{\pgfqpoint{1.322638in}{1.029233in}}%
\pgfpathlineto{\pgfqpoint{1.322940in}{1.030036in}}%
\pgfpathlineto{\pgfqpoint{1.324424in}{1.031663in}}%
\pgfpathlineto{\pgfqpoint{1.325702in}{1.030378in}}%
\pgfpathlineto{\pgfqpoint{1.327592in}{1.028250in}}%
\pgfpathlineto{\pgfqpoint{1.327887in}{1.029023in}}%
\pgfpathlineto{\pgfqpoint{1.329320in}{1.030653in}}%
\pgfpathlineto{\pgfqpoint{1.330468in}{1.029562in}}%
\pgfpathlineto{\pgfqpoint{1.332538in}{1.027248in}}%
\pgfpathlineto{\pgfqpoint{1.332823in}{1.028016in}}%
\pgfpathlineto{\pgfqpoint{1.334290in}{1.029673in}}%
\pgfpathlineto{\pgfqpoint{1.335487in}{1.028510in}}%
\pgfpathlineto{\pgfqpoint{1.337479in}{1.026266in}}%
\pgfpathlineto{\pgfqpoint{1.337785in}{1.027048in}}%
\pgfpathlineto{\pgfqpoint{1.339265in}{1.028656in}}%
\pgfpathlineto{\pgfqpoint{1.340530in}{1.027392in}}%
\pgfpathlineto{\pgfqpoint{1.342423in}{1.025273in}}%
\pgfpathlineto{\pgfqpoint{1.342693in}{1.025980in}}%
\pgfpathlineto{\pgfqpoint{1.344178in}{1.027695in}}%
\pgfpathlineto{\pgfqpoint{1.345383in}{1.026533in}}%
\pgfpathlineto{\pgfqpoint{1.347357in}{1.024304in}}%
\pgfpathlineto{\pgfqpoint{1.347618in}{1.024997in}}%
\pgfpathlineto{\pgfqpoint{1.349101in}{1.026775in}}%
\pgfpathlineto{\pgfqpoint{1.350246in}{1.025704in}}%
\pgfpathlineto{\pgfqpoint{1.352321in}{1.023389in}}%
\pgfpathlineto{\pgfqpoint{1.352595in}{1.024127in}}%
\pgfpathlineto{\pgfqpoint{1.354078in}{1.025823in}}%
\pgfpathlineto{\pgfqpoint{1.355290in}{1.024652in}}%
\pgfpathlineto{\pgfqpoint{1.357262in}{1.022448in}}%
\pgfpathlineto{\pgfqpoint{1.357565in}{1.023220in}}%
\pgfpathlineto{\pgfqpoint{1.359054in}{1.024843in}}%
\pgfpathlineto{\pgfqpoint{1.360316in}{1.023592in}}%
\pgfpathlineto{\pgfqpoint{1.362208in}{1.021503in}}%
\pgfpathlineto{\pgfqpoint{1.362497in}{1.022271in}}%
\pgfpathlineto{\pgfqpoint{1.363976in}{1.023941in}}%
\pgfpathlineto{\pgfqpoint{1.365182in}{1.022778in}}%
\pgfpathlineto{\pgfqpoint{1.367152in}{1.020586in}}%
\pgfpathlineto{\pgfqpoint{1.367450in}{1.021349in}}%
\pgfpathlineto{\pgfqpoint{1.368857in}{1.023017in}}%
\pgfpathlineto{\pgfqpoint{1.369941in}{1.022056in}}%
\pgfpathlineto{\pgfqpoint{1.372099in}{1.019670in}}%
\pgfpathlineto{\pgfqpoint{1.372387in}{1.020435in}}%
\pgfpathlineto{\pgfqpoint{1.373880in}{1.022107in}}%
\pgfpathlineto{\pgfqpoint{1.375126in}{1.020895in}}%
\pgfpathlineto{\pgfqpoint{1.377043in}{1.018772in}}%
\pgfpathlineto{\pgfqpoint{1.377341in}{1.019537in}}%
\pgfpathlineto{\pgfqpoint{1.378775in}{1.021208in}}%
\pgfpathlineto{\pgfqpoint{1.379912in}{1.020169in}}%
\pgfpathlineto{\pgfqpoint{1.381988in}{1.017877in}}%
\pgfpathlineto{\pgfqpoint{1.382281in}{1.018643in}}%
\pgfpathlineto{\pgfqpoint{1.383783in}{1.020312in}}%
\pgfpathlineto{\pgfqpoint{1.385047in}{1.019077in}}%
\pgfpathlineto{\pgfqpoint{1.386934in}{1.016998in}}%
\pgfpathlineto{\pgfqpoint{1.387234in}{1.017776in}}%
\pgfpathlineto{\pgfqpoint{1.388639in}{1.019449in}}%
\pgfpathlineto{\pgfqpoint{1.389723in}{1.018502in}}%
\pgfpathlineto{\pgfqpoint{1.391881in}{1.016131in}}%
\pgfpathlineto{\pgfqpoint{1.392168in}{1.016888in}}%
\pgfpathlineto{\pgfqpoint{1.393697in}{1.018565in}}%
\pgfpathlineto{\pgfqpoint{1.395016in}{1.017257in}}%
\pgfpathlineto{\pgfqpoint{1.396826in}{1.015272in}}%
\pgfpathlineto{\pgfqpoint{1.397122in}{1.016038in}}%
\pgfpathlineto{\pgfqpoint{1.398521in}{1.017727in}}%
\pgfpathlineto{\pgfqpoint{1.399560in}{1.016847in}}%
\pgfpathlineto{\pgfqpoint{1.401770in}{1.014417in}}%
\pgfpathlineto{\pgfqpoint{1.402040in}{1.015123in}}%
\pgfpathlineto{\pgfqpoint{1.403556in}{1.016862in}}%
\pgfpathlineto{\pgfqpoint{1.404774in}{1.015700in}}%
\pgfpathlineto{\pgfqpoint{1.406717in}{1.013579in}}%
\pgfpathlineto{\pgfqpoint{1.407006in}{1.014337in}}%
\pgfpathlineto{\pgfqpoint{1.408502in}{1.016031in}}%
\pgfpathlineto{\pgfqpoint{1.409721in}{1.014871in}}%
\pgfpathlineto{\pgfqpoint{1.411661in}{1.012748in}}%
\pgfpathlineto{\pgfqpoint{1.411953in}{1.013504in}}%
\pgfpathlineto{\pgfqpoint{1.413456in}{1.015193in}}%
\pgfpathlineto{\pgfqpoint{1.414691in}{1.014011in}}%
\pgfpathlineto{\pgfqpoint{1.416608in}{1.011928in}}%
\pgfpathlineto{\pgfqpoint{1.416898in}{1.012691in}}%
\pgfpathlineto{\pgfqpoint{1.418386in}{1.014388in}}%
\pgfpathlineto{\pgfqpoint{1.419585in}{1.013263in}}%
\pgfpathlineto{\pgfqpoint{1.421552in}{1.011116in}}%
\pgfpathlineto{\pgfqpoint{1.421844in}{1.011871in}}%
\pgfpathlineto{\pgfqpoint{1.423347in}{1.013565in}}%
\pgfpathlineto{\pgfqpoint{1.424583in}{1.012389in}}%
\pgfpathlineto{\pgfqpoint{1.426498in}{1.010312in}}%
\pgfpathlineto{\pgfqpoint{1.426767in}{1.011022in}}%
\pgfpathlineto{\pgfqpoint{1.428313in}{1.012767in}}%
\pgfpathlineto{\pgfqpoint{1.429581in}{1.011540in}}%
\pgfpathlineto{\pgfqpoint{1.431443in}{1.009520in}}%
\pgfpathlineto{\pgfqpoint{1.431736in}{1.010278in}}%
\pgfpathlineto{\pgfqpoint{1.433267in}{1.011969in}}%
\pgfpathlineto{\pgfqpoint{1.434553in}{1.010719in}}%
\pgfpathlineto{\pgfqpoint{1.436389in}{1.008735in}}%
\pgfpathlineto{\pgfqpoint{1.436659in}{1.009444in}}%
\pgfpathlineto{\pgfqpoint{1.438190in}{1.011197in}}%
\pgfpathlineto{\pgfqpoint{1.439426in}{1.010025in}}%
\pgfpathlineto{\pgfqpoint{1.441334in}{1.007959in}}%
\pgfpathlineto{\pgfqpoint{1.441632in}{1.008733in}}%
\pgfpathlineto{\pgfqpoint{1.443082in}{1.010434in}}%
\pgfpathlineto{\pgfqpoint{1.444216in}{1.009426in}}%
\pgfpathlineto{\pgfqpoint{1.446279in}{1.007192in}}%
\pgfpathlineto{\pgfqpoint{1.446570in}{1.007949in}}%
\pgfpathlineto{\pgfqpoint{1.448070in}{1.009662in}}%
\pgfpathlineto{\pgfqpoint{1.449265in}{1.008550in}}%
\pgfpathlineto{\pgfqpoint{1.451225in}{1.006434in}}%
\pgfpathlineto{\pgfqpoint{1.451515in}{1.007190in}}%
\pgfpathlineto{\pgfqpoint{1.453022in}{1.008906in}}%
\pgfpathlineto{\pgfqpoint{1.454237in}{1.007770in}}%
\pgfpathlineto{\pgfqpoint{1.456169in}{1.005684in}}%
\pgfpathlineto{\pgfqpoint{1.456440in}{1.006386in}}%
\pgfpathlineto{\pgfqpoint{1.457961in}{1.008156in}}%
\pgfpathlineto{\pgfqpoint{1.459157in}{1.007050in}}%
\pgfpathlineto{\pgfqpoint{1.461115in}{1.004942in}}%
\pgfpathlineto{\pgfqpoint{1.461405in}{1.005698in}}%
\pgfpathlineto{\pgfqpoint{1.462914in}{1.007422in}}%
\pgfpathlineto{\pgfqpoint{1.464125in}{1.006295in}}%
\pgfpathlineto{\pgfqpoint{1.466057in}{1.004205in}}%
\pgfpathlineto{\pgfqpoint{1.466331in}{1.004896in}}%
\pgfpathlineto{\pgfqpoint{1.467874in}{1.006665in}}%
\pgfpathlineto{\pgfqpoint{1.469124in}{1.005487in}}%
\pgfpathlineto{\pgfqpoint{1.470984in}{1.003461in}}%
\pgfpathlineto{\pgfqpoint{1.471276in}{1.004202in}}%
\pgfpathlineto{\pgfqpoint{1.472821in}{1.005973in}}%
\pgfpathlineto{\pgfqpoint{1.474074in}{1.004793in}}%
\pgfpathlineto{\pgfqpoint{1.475946in}{1.002775in}}%
\pgfpathlineto{\pgfqpoint{1.476217in}{1.003473in}}%
\pgfpathlineto{\pgfqpoint{1.477754in}{1.005266in}}%
\pgfpathlineto{\pgfqpoint{1.478968in}{1.004144in}}%
\pgfpathlineto{\pgfqpoint{1.480894in}{1.002081in}}%
\pgfpathlineto{\pgfqpoint{1.481170in}{1.002796in}}%
\pgfpathlineto{\pgfqpoint{1.482693in}{1.004567in}}%
\pgfpathlineto{\pgfqpoint{1.483882in}{1.003479in}}%
\pgfpathlineto{\pgfqpoint{1.485843in}{1.001386in}}%
\pgfpathlineto{\pgfqpoint{1.486112in}{1.002094in}}%
\pgfpathlineto{\pgfqpoint{1.487640in}{1.003879in}}%
\pgfpathlineto{\pgfqpoint{1.488833in}{1.002789in}}%
\pgfpathlineto{\pgfqpoint{1.490784in}{1.000700in}}%
\pgfpathlineto{\pgfqpoint{1.491058in}{1.001410in}}%
\pgfpathlineto{\pgfqpoint{1.492594in}{1.003194in}}%
\pgfpathlineto{\pgfqpoint{1.493800in}{1.002084in}}%
\pgfpathlineto{\pgfqpoint{1.495733in}{1.000024in}}%
\pgfpathlineto{\pgfqpoint{1.496018in}{1.000770in}}%
\pgfpathlineto{\pgfqpoint{1.497490in}{1.002532in}}%
\pgfpathlineto{\pgfqpoint{1.498619in}{1.001555in}}%
\pgfpathlineto{\pgfqpoint{1.500676in}{0.999357in}}%
\pgfpathlineto{\pgfqpoint{1.500953in}{1.000072in}}%
\pgfpathlineto{\pgfqpoint{1.502471in}{1.001853in}}%
\pgfpathlineto{\pgfqpoint{1.503667in}{1.000772in}}%
\pgfpathlineto{\pgfqpoint{1.505624in}{0.998693in}}%
\pgfpathlineto{\pgfqpoint{1.505892in}{0.999397in}}%
\pgfpathlineto{\pgfqpoint{1.507448in}{1.001189in}}%
\pgfpathlineto{\pgfqpoint{1.508696in}{1.000028in}}%
\pgfpathlineto{\pgfqpoint{1.510575in}{0.998048in}}%
\pgfpathlineto{\pgfqpoint{1.510856in}{0.998797in}}%
\pgfpathlineto{\pgfqpoint{1.512324in}{1.000561in}}%
\pgfpathlineto{\pgfqpoint{1.513442in}{0.999605in}}%
\pgfpathlineto{\pgfqpoint{1.515513in}{0.997399in}}%
\pgfpathlineto{\pgfqpoint{1.515789in}{0.998112in}}%
\pgfpathlineto{\pgfqpoint{1.517326in}{0.999896in}}%
\pgfpathlineto{\pgfqpoint{1.518534in}{0.998795in}}%
\pgfpathlineto{\pgfqpoint{1.520441in}{0.996744in}}%
\pgfpathlineto{\pgfqpoint{1.520740in}{0.997507in}}%
\pgfpathlineto{\pgfqpoint{1.522208in}{0.999292in}}%
\pgfpathlineto{\pgfqpoint{1.523326in}{0.998347in}}%
\pgfpathlineto{\pgfqpoint{1.525406in}{0.996136in}}%
\pgfpathlineto{\pgfqpoint{1.525705in}{0.996906in}}%
\pgfpathlineto{\pgfqpoint{1.527227in}{0.998633in}}%
\pgfpathlineto{\pgfqpoint{1.528453in}{0.997508in}}%
\pgfpathlineto{\pgfqpoint{1.530333in}{0.995492in}}%
\pgfpathlineto{\pgfqpoint{1.530636in}{0.996269in}}%
\pgfpathlineto{\pgfqpoint{1.532117in}{0.998042in}}%
\pgfpathlineto{\pgfqpoint{1.533238in}{0.997083in}}%
\pgfpathlineto{\pgfqpoint{1.535301in}{0.994907in}}%
\pgfpathlineto{\pgfqpoint{1.535594in}{0.995672in}}%
\pgfpathlineto{\pgfqpoint{1.537052in}{0.997420in}}%
\pgfpathlineto{\pgfqpoint{1.538170in}{0.996477in}}%
\pgfpathlineto{\pgfqpoint{1.540247in}{0.994296in}}%
\pgfpathlineto{\pgfqpoint{1.540533in}{0.995051in}}%
\pgfpathlineto{\pgfqpoint{1.541997in}{0.996817in}}%
\pgfpathlineto{\pgfqpoint{1.543116in}{0.995876in}}%
\pgfpathlineto{\pgfqpoint{1.545190in}{0.993693in}}%
\pgfpathlineto{\pgfqpoint{1.545479in}{0.994456in}}%
\pgfpathlineto{\pgfqpoint{1.546974in}{0.996218in}}%
\pgfpathlineto{\pgfqpoint{1.548131in}{0.995208in}}%
\pgfpathlineto{\pgfqpoint{1.550132in}{0.993100in}}%
\pgfpathlineto{\pgfqpoint{1.550415in}{0.993865in}}%
\pgfpathlineto{\pgfqpoint{1.551906in}{0.995655in}}%
\pgfpathlineto{\pgfqpoint{1.553050in}{0.994675in}}%
\pgfpathlineto{\pgfqpoint{1.555069in}{0.992522in}}%
\pgfpathlineto{\pgfqpoint{1.555340in}{0.993232in}}%
\pgfpathlineto{\pgfqpoint{1.556873in}{0.995084in}}%
\pgfpathlineto{\pgfqpoint{1.558022in}{0.994084in}}%
\pgfpathlineto{\pgfqpoint{1.560025in}{0.991973in}}%
\pgfpathlineto{\pgfqpoint{1.560309in}{0.992719in}}%
\pgfpathlineto{\pgfqpoint{1.561833in}{0.994492in}}%
\pgfpathlineto{\pgfqpoint{1.563045in}{0.993408in}}%
\pgfpathlineto{\pgfqpoint{1.564968in}{0.991387in}}%
\pgfpathlineto{\pgfqpoint{1.565257in}{0.992166in}}%
\pgfpathlineto{\pgfqpoint{1.566766in}{0.993950in}}%
\pgfpathlineto{\pgfqpoint{1.567941in}{0.992923in}}%
\pgfpathlineto{\pgfqpoint{1.569904in}{0.990833in}}%
\pgfpathlineto{\pgfqpoint{1.570214in}{0.991617in}}%
\pgfpathlineto{\pgfqpoint{1.571759in}{0.993361in}}%
\pgfpathlineto{\pgfqpoint{1.573016in}{0.992207in}}%
\pgfpathlineto{\pgfqpoint{1.574852in}{0.990272in}}%
\pgfpathlineto{\pgfqpoint{1.575136in}{0.991016in}}%
\pgfpathlineto{\pgfqpoint{1.576634in}{0.992835in}}%
\pgfpathlineto{\pgfqpoint{1.577748in}{0.991900in}}%
\pgfpathlineto{\pgfqpoint{1.579810in}{0.989749in}}%
\pgfpathlineto{\pgfqpoint{1.580097in}{0.990502in}}%
\pgfpathlineto{\pgfqpoint{1.581622in}{0.992267in}}%
\pgfpathlineto{\pgfqpoint{1.582830in}{0.991192in}}%
\pgfpathlineto{\pgfqpoint{1.584741in}{0.989169in}}%
\pgfpathlineto{\pgfqpoint{1.585046in}{0.989949in}}%
\pgfpathlineto{\pgfqpoint{1.586595in}{0.991718in}}%
\pgfpathlineto{\pgfqpoint{1.587859in}{0.990567in}}%
\pgfpathlineto{\pgfqpoint{1.589693in}{0.988647in}}%
\pgfpathlineto{\pgfqpoint{1.589969in}{0.989396in}}%
\pgfpathlineto{\pgfqpoint{1.591480in}{0.991217in}}%
\pgfpathlineto{\pgfqpoint{1.592618in}{0.990251in}}%
\pgfpathlineto{\pgfqpoint{1.594655in}{0.988138in}}%
\pgfpathlineto{\pgfqpoint{1.594938in}{0.988890in}}%
\pgfpathlineto{\pgfqpoint{1.596452in}{0.990669in}}%
\pgfpathlineto{\pgfqpoint{1.597618in}{0.989656in}}%
\pgfpathlineto{\pgfqpoint{1.599576in}{0.987589in}}%
\pgfpathlineto{\pgfqpoint{1.599846in}{0.988295in}}%
\pgfpathlineto{\pgfqpoint{1.601379in}{0.990156in}}%
\pgfpathlineto{\pgfqpoint{1.602536in}{0.989166in}}%
\pgfpathlineto{\pgfqpoint{1.604546in}{0.987086in}}%
\pgfpathlineto{\pgfqpoint{1.604826in}{0.987829in}}%
\pgfpathlineto{\pgfqpoint{1.606354in}{0.989616in}}%
\pgfpathlineto{\pgfqpoint{1.607535in}{0.988582in}}%
\pgfpathlineto{\pgfqpoint{1.609468in}{0.986549in}}%
\pgfpathlineto{\pgfqpoint{1.609769in}{0.987320in}}%
\pgfpathlineto{\pgfqpoint{1.611291in}{0.989103in}}%
\pgfpathlineto{\pgfqpoint{1.612481in}{0.988065in}}%
\pgfpathlineto{\pgfqpoint{1.614415in}{0.986038in}}%
\pgfpathlineto{\pgfqpoint{1.614698in}{0.986771in}}%
\pgfpathlineto{\pgfqpoint{1.616175in}{0.988618in}}%
\pgfpathlineto{\pgfqpoint{1.617278in}{0.987736in}}%
\pgfpathlineto{\pgfqpoint{1.619360in}{0.985552in}}%
\pgfpathlineto{\pgfqpoint{1.619660in}{0.986325in}}%
\pgfpathlineto{\pgfqpoint{1.621191in}{0.988113in}}%
\pgfpathlineto{\pgfqpoint{1.622398in}{0.987053in}}%
\pgfpathlineto{\pgfqpoint{1.624307in}{0.985061in}}%
\pgfpathlineto{\pgfqpoint{1.624592in}{0.985808in}}%
\pgfpathlineto{\pgfqpoint{1.626123in}{0.987643in}}%
\pgfpathlineto{\pgfqpoint{1.627291in}{0.986642in}}%
\pgfpathlineto{\pgfqpoint{1.629250in}{0.984588in}}%
\pgfpathlineto{\pgfqpoint{1.629549in}{0.985359in}}%
\pgfpathlineto{\pgfqpoint{1.631073in}{0.987156in}}%
\pgfpathlineto{\pgfqpoint{1.632263in}{0.986127in}}%
\pgfpathlineto{\pgfqpoint{1.634196in}{0.984099in}}%
\pgfpathlineto{\pgfqpoint{1.634495in}{0.984871in}}%
\pgfpathlineto{\pgfqpoint{1.636021in}{0.986674in}}%
\pgfpathlineto{\pgfqpoint{1.637197in}{0.985662in}}%
\pgfpathlineto{\pgfqpoint{1.639150in}{0.983631in}}%
\pgfpathlineto{\pgfqpoint{1.639410in}{0.984334in}}%
\pgfpathlineto{\pgfqpoint{1.640942in}{0.986215in}}%
\pgfpathlineto{\pgfqpoint{1.642087in}{0.985257in}}%
\pgfpathlineto{\pgfqpoint{1.644091in}{0.983165in}}%
\pgfpathlineto{\pgfqpoint{1.644385in}{0.983944in}}%
\pgfpathlineto{\pgfqpoint{1.645899in}{0.985759in}}%
\pgfpathlineto{\pgfqpoint{1.647050in}{0.984788in}}%
\pgfpathlineto{\pgfqpoint{1.649033in}{0.982716in}}%
\pgfpathlineto{\pgfqpoint{1.649299in}{0.983417in}}%
\pgfpathlineto{\pgfqpoint{1.650832in}{0.985309in}}%
\pgfpathlineto{\pgfqpoint{1.651967in}{0.984368in}}%
\pgfpathlineto{\pgfqpoint{1.653983in}{0.982265in}}%
\pgfpathlineto{\pgfqpoint{1.654264in}{0.983009in}}%
\pgfpathlineto{\pgfqpoint{1.655797in}{0.984847in}}%
\pgfpathlineto{\pgfqpoint{1.656947in}{0.983874in}}%
\pgfpathlineto{\pgfqpoint{1.658924in}{0.981812in}}%
\pgfpathlineto{\pgfqpoint{1.659199in}{0.982535in}}%
\pgfpathlineto{\pgfqpoint{1.660726in}{0.984414in}}%
\pgfpathlineto{\pgfqpoint{1.661874in}{0.983462in}}%
\pgfpathlineto{\pgfqpoint{1.663881in}{0.981385in}}%
\pgfpathlineto{\pgfqpoint{1.664127in}{0.982048in}}%
\pgfpathlineto{\pgfqpoint{1.665661in}{0.983963in}}%
\pgfpathlineto{\pgfqpoint{1.666774in}{0.983056in}}%
\pgfpathlineto{\pgfqpoint{1.668818in}{0.980933in}}%
\pgfpathlineto{\pgfqpoint{1.669103in}{0.981691in}}%
\pgfpathlineto{\pgfqpoint{1.670661in}{0.983532in}}%
\pgfpathlineto{\pgfqpoint{1.671853in}{0.982497in}}%
\pgfpathlineto{\pgfqpoint{1.673761in}{0.980514in}}%
\pgfpathlineto{\pgfqpoint{1.674053in}{0.981279in}}%
\pgfpathlineto{\pgfqpoint{1.675601in}{0.983105in}}%
\pgfpathlineto{\pgfqpoint{1.676798in}{0.982070in}}%
\pgfpathlineto{\pgfqpoint{1.678705in}{0.980087in}}%
\pgfpathlineto{\pgfqpoint{1.679002in}{0.980857in}}%
\pgfpathlineto{\pgfqpoint{1.680512in}{0.982684in}}%
\pgfpathlineto{\pgfqpoint{1.681639in}{0.981756in}}%
\pgfpathlineto{\pgfqpoint{1.683654in}{0.979673in}}%
\pgfpathlineto{\pgfqpoint{1.683941in}{0.980436in}}%
\pgfpathlineto{\pgfqpoint{1.685524in}{0.982266in}}%
\pgfpathlineto{\pgfqpoint{1.686770in}{0.981157in}}%
\pgfpathlineto{\pgfqpoint{1.688598in}{0.979263in}}%
\pgfpathlineto{\pgfqpoint{1.688893in}{0.980041in}}%
\pgfpathlineto{\pgfqpoint{1.690433in}{0.981867in}}%
\pgfpathlineto{\pgfqpoint{1.691610in}{0.980864in}}%
\pgfpathlineto{\pgfqpoint{1.693545in}{0.978863in}}%
\pgfpathlineto{\pgfqpoint{1.693828in}{0.979619in}}%
\pgfpathlineto{\pgfqpoint{1.695394in}{0.981470in}}%
\pgfpathlineto{\pgfqpoint{1.696611in}{0.980411in}}%
\pgfpathlineto{\pgfqpoint{1.698495in}{0.978467in}}%
\pgfpathlineto{\pgfqpoint{1.698772in}{0.979203in}}%
\pgfpathlineto{\pgfqpoint{1.700329in}{0.981056in}}%
\pgfpathlineto{\pgfqpoint{1.701525in}{0.980028in}}%
\pgfpathlineto{\pgfqpoint{1.703436in}{0.978056in}}%
\pgfpathlineto{\pgfqpoint{1.703722in}{0.978820in}}%
\pgfpathlineto{\pgfqpoint{1.705265in}{0.980672in}}%
\pgfpathlineto{\pgfqpoint{1.706445in}{0.979675in}}%
\pgfpathlineto{\pgfqpoint{1.708381in}{0.977676in}}%
\pgfpathlineto{\pgfqpoint{1.708668in}{0.978435in}}%
\pgfpathlineto{\pgfqpoint{1.710210in}{0.980282in}}%
\pgfpathlineto{\pgfqpoint{1.711387in}{0.979288in}}%
\pgfpathlineto{\pgfqpoint{1.713327in}{0.977288in}}%
\pgfpathlineto{\pgfqpoint{1.713609in}{0.978037in}}%
\pgfpathlineto{\pgfqpoint{1.715174in}{0.979898in}}%
\pgfpathlineto{\pgfqpoint{1.716376in}{0.978865in}}%
\pgfpathlineto{\pgfqpoint{1.718276in}{0.976917in}}%
\pgfpathlineto{\pgfqpoint{1.718551in}{0.977625in}}%
\pgfpathlineto{\pgfqpoint{1.720124in}{0.979484in}}%
\pgfpathlineto{\pgfqpoint{1.721331in}{0.978440in}}%
\pgfpathlineto{\pgfqpoint{1.723219in}{0.976517in}}%
\pgfpathlineto{\pgfqpoint{1.723505in}{0.977287in}}%
\pgfpathlineto{\pgfqpoint{1.725059in}{0.979144in}}%
\pgfpathlineto{\pgfqpoint{1.726242in}{0.978141in}}%
\pgfpathlineto{\pgfqpoint{1.728174in}{0.976173in}}%
\pgfpathlineto{\pgfqpoint{1.728468in}{0.976923in}}%
\pgfpathlineto{\pgfqpoint{1.729997in}{0.978719in}}%
\pgfpathlineto{\pgfqpoint{1.731173in}{0.977725in}}%
\pgfpathlineto{\pgfqpoint{1.733095in}{0.975744in}}%
\pgfpathlineto{\pgfqpoint{1.733366in}{0.976463in}}%
\pgfpathlineto{\pgfqpoint{1.734954in}{0.978400in}}%
\pgfpathlineto{\pgfqpoint{1.736137in}{0.977398in}}%
\pgfpathlineto{\pgfqpoint{1.738058in}{0.975425in}}%
\pgfpathlineto{\pgfqpoint{1.738339in}{0.976171in}}%
\pgfpathlineto{\pgfqpoint{1.739863in}{0.978023in}}%
\pgfpathlineto{\pgfqpoint{1.741005in}{0.977090in}}%
\pgfpathlineto{\pgfqpoint{1.742994in}{0.975036in}}%
\pgfpathlineto{\pgfqpoint{1.743294in}{0.975796in}}%
\pgfpathlineto{\pgfqpoint{1.744829in}{0.977640in}}%
\pgfpathlineto{\pgfqpoint{1.745983in}{0.976683in}}%
\pgfpathlineto{\pgfqpoint{1.747939in}{0.974673in}}%
\pgfpathlineto{\pgfqpoint{1.748250in}{0.975436in}}%
\pgfpathlineto{\pgfqpoint{1.749746in}{0.977237in}}%
\pgfpathlineto{\pgfqpoint{1.750867in}{0.976336in}}%
\pgfpathlineto{\pgfqpoint{1.752873in}{0.974274in}}%
\pgfpathlineto{\pgfqpoint{1.753146in}{0.974982in}}%
\pgfpathlineto{\pgfqpoint{1.754714in}{0.976932in}}%
\pgfpathlineto{\pgfqpoint{1.755865in}{0.975988in}}%
\pgfpathlineto{\pgfqpoint{1.757840in}{0.973975in}}%
\pgfpathlineto{\pgfqpoint{1.758123in}{0.974731in}}%
\pgfpathlineto{\pgfqpoint{1.759673in}{0.976574in}}%
\pgfpathlineto{\pgfqpoint{1.760860in}{0.975575in}}%
\pgfpathlineto{\pgfqpoint{1.762784in}{0.973612in}}%
\pgfpathlineto{\pgfqpoint{1.763056in}{0.974316in}}%
\pgfpathlineto{\pgfqpoint{1.764618in}{0.976214in}}%
\pgfpathlineto{\pgfqpoint{1.765779in}{0.975250in}}%
\pgfpathlineto{\pgfqpoint{1.767729in}{0.973273in}}%
\pgfpathlineto{\pgfqpoint{1.767999in}{0.974006in}}%
\pgfpathlineto{\pgfqpoint{1.769558in}{0.975894in}}%
\pgfpathlineto{\pgfqpoint{1.770736in}{0.974912in}}%
\pgfpathlineto{\pgfqpoint{1.772668in}{0.972929in}}%
\pgfpathlineto{\pgfqpoint{1.772958in}{0.973678in}}%
\pgfpathlineto{\pgfqpoint{1.774455in}{0.975561in}}%
\pgfpathlineto{\pgfqpoint{1.775562in}{0.974709in}}%
\pgfpathlineto{\pgfqpoint{1.777620in}{0.972608in}}%
\pgfpathlineto{\pgfqpoint{1.777891in}{0.973332in}}%
\pgfpathlineto{\pgfqpoint{1.779424in}{0.975224in}}%
\pgfpathlineto{\pgfqpoint{1.780544in}{0.974329in}}%
\pgfpathlineto{\pgfqpoint{1.782565in}{0.972268in}}%
\pgfpathlineto{\pgfqpoint{1.782841in}{0.972999in}}%
\pgfpathlineto{\pgfqpoint{1.784393in}{0.974890in}}%
\pgfpathlineto{\pgfqpoint{1.785565in}{0.973924in}}%
\pgfpathlineto{\pgfqpoint{1.787511in}{0.971939in}}%
\pgfpathlineto{\pgfqpoint{1.787788in}{0.972613in}}%
\pgfpathlineto{\pgfqpoint{1.789351in}{0.974489in}}%
\pgfpathlineto{\pgfqpoint{1.790536in}{0.973497in}}%
\pgfpathlineto{\pgfqpoint{1.792446in}{0.971563in}}%
\pgfpathlineto{\pgfqpoint{1.792716in}{0.972294in}}%
\pgfpathlineto{\pgfqpoint{1.794299in}{0.974226in}}%
\pgfpathlineto{\pgfqpoint{1.795473in}{0.973248in}}%
\pgfpathlineto{\pgfqpoint{1.797400in}{0.971282in}}%
\pgfpathlineto{\pgfqpoint{1.797695in}{0.972043in}}%
\pgfpathlineto{\pgfqpoint{1.799200in}{0.973897in}}%
\pgfpathlineto{\pgfqpoint{1.800320in}{0.973018in}}%
\pgfpathlineto{\pgfqpoint{1.802342in}{0.970963in}}%
\pgfpathlineto{\pgfqpoint{1.802614in}{0.971692in}}%
\pgfpathlineto{\pgfqpoint{1.804197in}{0.973608in}}%
\pgfpathlineto{\pgfqpoint{1.805398in}{0.972600in}}%
\pgfpathlineto{\pgfqpoint{1.807287in}{0.970670in}}%
\pgfpathlineto{\pgfqpoint{1.807578in}{0.971429in}}%
\pgfpathlineto{\pgfqpoint{1.809161in}{0.973298in}}%
\pgfpathlineto{\pgfqpoint{1.810380in}{0.972256in}}%
\pgfpathlineto{\pgfqpoint{1.812241in}{0.970371in}}%
\pgfpathlineto{\pgfqpoint{1.812517in}{0.971049in}}%
\pgfpathlineto{\pgfqpoint{1.814122in}{0.972919in}}%
\pgfpathlineto{\pgfqpoint{1.815391in}{0.971809in}}%
\pgfpathlineto{\pgfqpoint{1.817170in}{0.970013in}}%
\pgfpathlineto{\pgfqpoint{1.817441in}{0.970735in}}%
\pgfpathlineto{\pgfqpoint{1.819010in}{0.972678in}}%
\pgfpathlineto{\pgfqpoint{1.820153in}{0.971753in}}%
\pgfpathlineto{\pgfqpoint{1.822124in}{0.969744in}}%
\pgfpathlineto{\pgfqpoint{1.822422in}{0.970524in}}%
\pgfpathlineto{\pgfqpoint{1.824004in}{0.972379in}}%
\pgfpathlineto{\pgfqpoint{1.825229in}{0.971330in}}%
\pgfpathlineto{\pgfqpoint{1.827071in}{0.969461in}}%
\pgfpathlineto{\pgfqpoint{1.827358in}{0.970224in}}%
\pgfpathlineto{\pgfqpoint{1.828899in}{0.972105in}}%
\pgfpathlineto{\pgfqpoint{1.830041in}{0.971186in}}%
\pgfpathlineto{\pgfqpoint{1.832018in}{0.969179in}}%
\pgfpathlineto{\pgfqpoint{1.832302in}{0.969935in}}%
\pgfpathlineto{\pgfqpoint{1.833864in}{0.971823in}}%
\pgfpathlineto{\pgfqpoint{1.835046in}{0.970847in}}%
\pgfpathlineto{\pgfqpoint{1.836960in}{0.968901in}}%
\pgfpathlineto{\pgfqpoint{1.837252in}{0.969666in}}%
\pgfpathlineto{\pgfqpoint{1.838785in}{0.971545in}}%
\pgfpathlineto{\pgfqpoint{1.839925in}{0.970635in}}%
\pgfpathlineto{\pgfqpoint{1.841908in}{0.968625in}}%
\pgfpathlineto{\pgfqpoint{1.842188in}{0.969371in}}%
\pgfpathlineto{\pgfqpoint{1.843750in}{0.971274in}}%
\pgfpathlineto{\pgfqpoint{1.844914in}{0.970324in}}%
\pgfpathlineto{\pgfqpoint{1.846851in}{0.968356in}}%
\pgfpathlineto{\pgfqpoint{1.847133in}{0.969094in}}%
\pgfpathlineto{\pgfqpoint{1.848715in}{0.970995in}}%
\pgfpathlineto{\pgfqpoint{1.849905in}{0.970002in}}%
\pgfpathlineto{\pgfqpoint{1.851797in}{0.968085in}}%
\pgfpathlineto{\pgfqpoint{1.852089in}{0.968854in}}%
\pgfpathlineto{\pgfqpoint{1.853664in}{0.970732in}}%
\pgfpathlineto{\pgfqpoint{1.854861in}{0.969730in}}%
\pgfpathlineto{\pgfqpoint{1.856742in}{0.967823in}}%
\pgfpathlineto{\pgfqpoint{1.857029in}{0.968578in}}%
\pgfpathlineto{\pgfqpoint{1.858562in}{0.970475in}}%
\pgfpathlineto{\pgfqpoint{1.859702in}{0.969575in}}%
\pgfpathlineto{\pgfqpoint{1.861688in}{0.967563in}}%
\pgfpathlineto{\pgfqpoint{1.861958in}{0.968275in}}%
\pgfpathlineto{\pgfqpoint{1.863533in}{0.970211in}}%
\pgfpathlineto{\pgfqpoint{1.864703in}{0.969258in}}%
\pgfpathlineto{\pgfqpoint{1.866632in}{0.967303in}}%
\pgfpathlineto{\pgfqpoint{1.866923in}{0.968064in}}%
\pgfpathlineto{\pgfqpoint{1.868488in}{0.969947in}}%
\pgfpathlineto{\pgfqpoint{1.869673in}{0.968970in}}%
\pgfpathlineto{\pgfqpoint{1.871581in}{0.967047in}}%
\pgfpathlineto{\pgfqpoint{1.871870in}{0.967818in}}%
\pgfpathlineto{\pgfqpoint{1.873397in}{0.969706in}}%
\pgfpathlineto{\pgfqpoint{1.874512in}{0.968837in}}%
\pgfpathlineto{\pgfqpoint{1.876527in}{0.966801in}}%
\pgfpathlineto{\pgfqpoint{1.876806in}{0.967532in}}%
\pgfpathlineto{\pgfqpoint{1.878392in}{0.969437in}}%
\pgfpathlineto{\pgfqpoint{1.879592in}{0.968435in}}%
\pgfpathlineto{\pgfqpoint{1.881469in}{0.966542in}}%
\pgfpathlineto{\pgfqpoint{1.881758in}{0.967300in}}%
\pgfpathlineto{\pgfqpoint{1.883323in}{0.969192in}}%
\pgfpathlineto{\pgfqpoint{1.884487in}{0.968242in}}%
\pgfpathlineto{\pgfqpoint{1.886412in}{0.966292in}}%
\pgfpathlineto{\pgfqpoint{1.886708in}{0.967051in}}%
\pgfpathlineto{\pgfqpoint{1.888255in}{0.968934in}}%
\pgfpathlineto{\pgfqpoint{1.889425in}{0.967990in}}%
\pgfpathlineto{\pgfqpoint{1.891364in}{0.966044in}}%
\pgfpathlineto{\pgfqpoint{1.891644in}{0.966787in}}%
\pgfpathlineto{\pgfqpoint{1.893194in}{0.968698in}}%
\pgfpathlineto{\pgfqpoint{1.894349in}{0.967778in}}%
\pgfpathlineto{\pgfqpoint{1.896310in}{0.965808in}}%
\pgfpathlineto{\pgfqpoint{1.896598in}{0.966560in}}%
\pgfpathlineto{\pgfqpoint{1.898166in}{0.968443in}}%
\pgfpathlineto{\pgfqpoint{1.899343in}{0.967477in}}%
\pgfpathlineto{\pgfqpoint{1.901253in}{0.965561in}}%
\pgfpathlineto{\pgfqpoint{1.901548in}{0.966323in}}%
\pgfpathlineto{\pgfqpoint{1.903129in}{0.968195in}}%
\pgfpathlineto{\pgfqpoint{1.904344in}{0.967176in}}%
\pgfpathlineto{\pgfqpoint{1.906199in}{0.965323in}}%
\pgfpathlineto{\pgfqpoint{1.906481in}{0.966054in}}%
\pgfpathlineto{\pgfqpoint{1.908072in}{0.967958in}}%
\pgfpathlineto{\pgfqpoint{1.909290in}{0.966938in}}%
\pgfpathlineto{\pgfqpoint{1.911147in}{0.965090in}}%
\pgfpathlineto{\pgfqpoint{1.911435in}{0.965853in}}%
\pgfpathlineto{\pgfqpoint{1.913012in}{0.967739in}}%
\pgfpathlineto{\pgfqpoint{1.914210in}{0.966747in}}%
\pgfpathlineto{\pgfqpoint{1.916092in}{0.964864in}}%
\pgfpathlineto{\pgfqpoint{1.916385in}{0.965624in}}%
\pgfpathlineto{\pgfqpoint{1.917876in}{0.967507in}}%
\pgfpathlineto{\pgfqpoint{1.918981in}{0.966687in}}%
\pgfpathlineto{\pgfqpoint{1.921039in}{0.964628in}}%
\pgfpathlineto{\pgfqpoint{1.921306in}{0.965350in}}%
\pgfpathlineto{\pgfqpoint{1.922875in}{0.967286in}}%
\pgfpathlineto{\pgfqpoint{1.924030in}{0.966363in}}%
\pgfpathlineto{\pgfqpoint{1.925985in}{0.964399in}}%
\pgfpathlineto{\pgfqpoint{1.926237in}{0.965069in}}%
\pgfpathlineto{\pgfqpoint{1.927820in}{0.967048in}}%
\pgfpathlineto{\pgfqpoint{1.928975in}{0.966128in}}%
\pgfpathlineto{\pgfqpoint{1.930932in}{0.964168in}}%
\pgfpathlineto{\pgfqpoint{1.931211in}{0.964918in}}%
\pgfpathlineto{\pgfqpoint{1.932789in}{0.966819in}}%
\pgfpathlineto{\pgfqpoint{1.933976in}{0.965844in}}%
\pgfpathlineto{\pgfqpoint{1.935873in}{0.963941in}}%
\pgfpathlineto{\pgfqpoint{1.936152in}{0.964673in}}%
\pgfpathlineto{\pgfqpoint{1.937714in}{0.966599in}}%
\pgfpathlineto{\pgfqpoint{1.938866in}{0.965683in}}%
\pgfpathlineto{\pgfqpoint{1.940824in}{0.963730in}}%
\pgfpathlineto{\pgfqpoint{1.941109in}{0.964496in}}%
\pgfpathlineto{\pgfqpoint{1.942691in}{0.966376in}}%
\pgfpathlineto{\pgfqpoint{1.943908in}{0.965363in}}%
\pgfpathlineto{\pgfqpoint{1.945763in}{0.963504in}}%
\pgfpathlineto{\pgfqpoint{1.946054in}{0.964273in}}%
\pgfpathlineto{\pgfqpoint{1.947626in}{0.966173in}}%
\pgfpathlineto{\pgfqpoint{1.948818in}{0.965200in}}%
\pgfpathlineto{\pgfqpoint{1.950710in}{0.963307in}}%
\pgfpathlineto{\pgfqpoint{1.950985in}{0.964041in}}%
\pgfpathlineto{\pgfqpoint{1.952568in}{0.965958in}}%
\pgfpathlineto{\pgfqpoint{1.953750in}{0.964992in}}%
\pgfpathlineto{\pgfqpoint{1.955648in}{0.963080in}}%
\pgfpathlineto{\pgfqpoint{1.955946in}{0.963837in}}%
\pgfpathlineto{\pgfqpoint{1.957434in}{0.965739in}}%
\pgfpathlineto{\pgfqpoint{1.958528in}{0.964948in}}%
\pgfpathlineto{\pgfqpoint{1.960609in}{0.962882in}}%
\pgfpathlineto{\pgfqpoint{1.960872in}{0.963592in}}%
\pgfpathlineto{\pgfqpoint{1.962454in}{0.965529in}}%
\pgfpathlineto{\pgfqpoint{1.963617in}{0.964591in}}%
\pgfpathlineto{\pgfqpoint{1.965539in}{0.962656in}}%
\pgfpathlineto{\pgfqpoint{1.965840in}{0.963426in}}%
\pgfpathlineto{\pgfqpoint{1.967372in}{0.965314in}}%
\pgfpathlineto{\pgfqpoint{1.968509in}{0.964432in}}%
\pgfpathlineto{\pgfqpoint{1.970486in}{0.962450in}}%
\pgfpathlineto{\pgfqpoint{1.970766in}{0.963175in}}%
\pgfpathlineto{\pgfqpoint{1.972272in}{0.965116in}}%
\pgfpathlineto{\pgfqpoint{1.973368in}{0.964320in}}%
\pgfpathlineto{\pgfqpoint{1.975431in}{0.962251in}}%
\pgfpathlineto{\pgfqpoint{1.975729in}{0.963023in}}%
\pgfpathlineto{\pgfqpoint{1.977276in}{0.964916in}}%
\pgfpathlineto{\pgfqpoint{1.978429in}{0.964007in}}%
\pgfpathlineto{\pgfqpoint{1.980376in}{0.962051in}}%
\pgfpathlineto{\pgfqpoint{1.980664in}{0.962792in}}%
\pgfpathlineto{\pgfqpoint{1.982238in}{0.964708in}}%
\pgfpathlineto{\pgfqpoint{1.983408in}{0.963766in}}%
\pgfpathlineto{\pgfqpoint{1.985329in}{0.961851in}}%
\pgfpathlineto{\pgfqpoint{1.985575in}{0.962510in}}%
\pgfpathlineto{\pgfqpoint{1.987157in}{0.964517in}}%
\pgfpathlineto{\pgfqpoint{1.988292in}{0.963636in}}%
\pgfpathlineto{\pgfqpoint{1.990289in}{0.961669in}}%
\pgfpathlineto{\pgfqpoint{1.990560in}{0.962385in}}%
\pgfpathlineto{\pgfqpoint{1.992137in}{0.964295in}}%
\pgfpathlineto{\pgfqpoint{1.993322in}{0.963334in}}%
\pgfpathlineto{\pgfqpoint{1.995219in}{0.961447in}}%
\pgfpathlineto{\pgfqpoint{1.995469in}{0.962118in}}%
\pgfpathlineto{\pgfqpoint{1.997051in}{0.964117in}}%
\pgfpathlineto{\pgfqpoint{1.998190in}{0.963229in}}%
\pgfpathlineto{\pgfqpoint{2.000169in}{0.961258in}}%
\pgfpathlineto{\pgfqpoint{2.000450in}{0.962012in}}%
\pgfpathlineto{\pgfqpoint{2.002027in}{0.963924in}}%
\pgfpathlineto{\pgfqpoint{2.003193in}{0.962984in}}%
\pgfpathlineto{\pgfqpoint{2.005111in}{0.961073in}}%
\pgfpathlineto{\pgfqpoint{2.005401in}{0.961845in}}%
\pgfpathlineto{\pgfqpoint{2.006984in}{0.963735in}}%
\pgfpathlineto{\pgfqpoint{2.008184in}{0.962749in}}%
\pgfpathlineto{\pgfqpoint{2.010056in}{0.960880in}}%
\pgfpathlineto{\pgfqpoint{2.010332in}{0.961610in}}%
\pgfpathlineto{\pgfqpoint{2.011915in}{0.963547in}}%
\pgfpathlineto{\pgfqpoint{2.013079in}{0.962614in}}%
\pgfpathlineto{\pgfqpoint{2.015001in}{0.960697in}}%
\pgfpathlineto{\pgfqpoint{2.015252in}{0.961365in}}%
\pgfpathlineto{\pgfqpoint{2.016834in}{0.963363in}}%
\pgfpathlineto{\pgfqpoint{2.017972in}{0.962478in}}%
\pgfpathlineto{\pgfqpoint{2.019954in}{0.960514in}}%
\pgfpathlineto{\pgfqpoint{2.020192in}{0.961155in}}%
\pgfpathlineto{\pgfqpoint{2.021774in}{0.963179in}}%
\pgfpathlineto{\pgfqpoint{2.022914in}{0.962302in}}%
\pgfpathlineto{\pgfqpoint{2.024903in}{0.960342in}}%
\pgfpathlineto{\pgfqpoint{2.025168in}{0.961055in}}%
\pgfpathlineto{\pgfqpoint{2.026750in}{0.962992in}}%
\pgfpathlineto{\pgfqpoint{2.027915in}{0.962061in}}%
\pgfpathlineto{\pgfqpoint{2.029837in}{0.960147in}}%
\pgfpathlineto{\pgfqpoint{2.030122in}{0.960903in}}%
\pgfpathlineto{\pgfqpoint{2.031704in}{0.962816in}}%
\pgfpathlineto{\pgfqpoint{2.032892in}{0.961854in}}%
\pgfpathlineto{\pgfqpoint{2.034785in}{0.959975in}}%
\pgfpathlineto{\pgfqpoint{2.035027in}{0.960626in}}%
\pgfpathlineto{\pgfqpoint{2.036609in}{0.962650in}}%
\pgfpathlineto{\pgfqpoint{2.037731in}{0.961793in}}%
\pgfpathlineto{\pgfqpoint{2.039734in}{0.959805in}}%
\pgfpathlineto{\pgfqpoint{2.039977in}{0.960461in}}%
\pgfpathlineto{\pgfqpoint{2.041560in}{0.962475in}}%
\pgfpathlineto{\pgfqpoint{2.042694in}{0.961603in}}%
\pgfpathlineto{\pgfqpoint{2.044668in}{0.959633in}}%
\pgfpathlineto{\pgfqpoint{2.044959in}{0.960392in}}%
\pgfpathlineto{\pgfqpoint{2.046506in}{0.962315in}}%
\pgfpathlineto{\pgfqpoint{2.047632in}{0.961451in}}%
\pgfpathlineto{\pgfqpoint{2.049620in}{0.959480in}}%
\pgfpathlineto{\pgfqpoint{2.049899in}{0.960229in}}%
\pgfpathlineto{\pgfqpoint{2.051476in}{0.962154in}}%
\pgfpathlineto{\pgfqpoint{2.052645in}{0.961222in}}%
\pgfpathlineto{\pgfqpoint{2.054568in}{0.959312in}}%
\pgfpathlineto{\pgfqpoint{2.054836in}{0.960026in}}%
\pgfpathlineto{\pgfqpoint{2.056418in}{0.961985in}}%
\pgfpathlineto{\pgfqpoint{2.057572in}{0.961074in}}%
\pgfpathlineto{\pgfqpoint{2.059511in}{0.959151in}}%
\pgfpathlineto{\pgfqpoint{2.059796in}{0.959913in}}%
\pgfpathlineto{\pgfqpoint{2.061372in}{0.961825in}}%
\pgfpathlineto{\pgfqpoint{2.062537in}{0.960894in}}%
\pgfpathlineto{\pgfqpoint{2.064456in}{0.958985in}}%
\pgfpathlineto{\pgfqpoint{2.064741in}{0.959738in}}%
\pgfpathlineto{\pgfqpoint{2.066324in}{0.961658in}}%
\pgfpathlineto{\pgfqpoint{2.067516in}{0.960696in}}%
\pgfpathlineto{\pgfqpoint{2.069407in}{0.958836in}}%
\pgfpathlineto{\pgfqpoint{2.069697in}{0.959585in}}%
\pgfpathlineto{\pgfqpoint{2.071248in}{0.961471in}}%
\pgfpathlineto{\pgfqpoint{2.072399in}{0.960568in}}%
\pgfpathlineto{\pgfqpoint{2.074347in}{0.958648in}}%
\pgfpathlineto{\pgfqpoint{2.074607in}{0.959354in}}%
\pgfpathlineto{\pgfqpoint{2.076189in}{0.961343in}}%
\pgfpathlineto{\pgfqpoint{2.077334in}{0.960455in}}%
\pgfpathlineto{\pgfqpoint{2.079293in}{0.958510in}}%
\pgfpathlineto{\pgfqpoint{2.079565in}{0.959233in}}%
\pgfpathlineto{\pgfqpoint{2.081148in}{0.961181in}}%
\pgfpathlineto{\pgfqpoint{2.082316in}{0.960255in}}%
\pgfpathlineto{\pgfqpoint{2.084232in}{0.958350in}}%
\pgfpathlineto{\pgfqpoint{2.084488in}{0.959027in}}%
\pgfpathlineto{\pgfqpoint{2.086070in}{0.961049in}}%
\pgfpathlineto{\pgfqpoint{2.087211in}{0.960177in}}%
\pgfpathlineto{\pgfqpoint{2.089189in}{0.958230in}}%
\pgfpathlineto{\pgfqpoint{2.089469in}{0.958980in}}%
\pgfpathlineto{\pgfqpoint{2.091044in}{0.960891in}}%
\pgfpathlineto{\pgfqpoint{2.092209in}{0.959964in}}%
\pgfpathlineto{\pgfqpoint{2.094133in}{0.958063in}}%
\pgfpathlineto{\pgfqpoint{2.094414in}{0.958820in}}%
\pgfpathlineto{\pgfqpoint{2.095997in}{0.960745in}}%
\pgfpathlineto{\pgfqpoint{2.097191in}{0.959788in}}%
\pgfpathlineto{\pgfqpoint{2.099069in}{0.957921in}}%
\pgfpathlineto{\pgfqpoint{2.099345in}{0.958649in}}%
\pgfpathlineto{\pgfqpoint{2.100928in}{0.960612in}}%
\pgfpathlineto{\pgfqpoint{2.102077in}{0.959713in}}%
\pgfpathlineto{\pgfqpoint{2.104022in}{0.957791in}}%
\pgfpathlineto{\pgfqpoint{2.104310in}{0.958562in}}%
\pgfpathlineto{\pgfqpoint{2.105863in}{0.960468in}}%
\pgfpathlineto{\pgfqpoint{2.107002in}{0.959586in}}%
\pgfpathlineto{\pgfqpoint{2.108972in}{0.957651in}}%
\pgfpathlineto{\pgfqpoint{2.109264in}{0.958395in}}%
\pgfpathlineto{\pgfqpoint{2.110834in}{0.960284in}}%
\pgfpathlineto{\pgfqpoint{2.112013in}{0.959341in}}%
\pgfpathlineto{\pgfqpoint{2.113907in}{0.957485in}}%
\pgfpathlineto{\pgfqpoint{2.114189in}{0.958242in}}%
\pgfpathlineto{\pgfqpoint{2.115756in}{0.960190in}}%
\pgfpathlineto{\pgfqpoint{2.116896in}{0.959309in}}%
\pgfpathlineto{\pgfqpoint{2.118851in}{0.957369in}}%
\pgfpathlineto{\pgfqpoint{2.119111in}{0.958057in}}%
\pgfpathlineto{\pgfqpoint{2.120694in}{0.960066in}}%
\pgfpathlineto{\pgfqpoint{2.121831in}{0.959197in}}%
\pgfpathlineto{\pgfqpoint{2.123804in}{0.957249in}}%
\pgfpathlineto{\pgfqpoint{2.124087in}{0.958009in}}%
\pgfpathlineto{\pgfqpoint{2.125660in}{0.959927in}}%
\pgfpathlineto{\pgfqpoint{2.126816in}{0.959016in}}%
\pgfpathlineto{\pgfqpoint{2.128748in}{0.957104in}}%
\pgfpathlineto{\pgfqpoint{2.128996in}{0.957768in}}%
\pgfpathlineto{\pgfqpoint{2.130579in}{0.959790in}}%
\pgfpathlineto{\pgfqpoint{2.131720in}{0.958921in}}%
\pgfpathlineto{\pgfqpoint{2.133700in}{0.956984in}}%
\pgfpathlineto{\pgfqpoint{2.133952in}{0.957664in}}%
\pgfpathlineto{\pgfqpoint{2.135535in}{0.959652in}}%
\pgfpathlineto{\pgfqpoint{2.136677in}{0.958771in}}%
\pgfpathlineto{\pgfqpoint{2.138644in}{0.956845in}}%
\pgfpathlineto{\pgfqpoint{2.138932in}{0.957587in}}%
\pgfpathlineto{\pgfqpoint{2.140522in}{0.959487in}}%
\pgfpathlineto{\pgfqpoint{2.141737in}{0.958498in}}%
\pgfpathlineto{\pgfqpoint{2.143573in}{0.956680in}}%
\pgfpathlineto{\pgfqpoint{2.143822in}{0.957343in}}%
\pgfpathlineto{\pgfqpoint{2.145405in}{0.959391in}}%
\pgfpathlineto{\pgfqpoint{2.146506in}{0.958574in}}%
\pgfpathlineto{\pgfqpoint{2.148536in}{0.956585in}}%
\pgfpathlineto{\pgfqpoint{2.148829in}{0.957326in}}%
\pgfpathlineto{\pgfqpoint{2.150416in}{0.959211in}}%
\pgfpathlineto{\pgfqpoint{2.151632in}{0.958219in}}%
\pgfpathlineto{\pgfqpoint{2.153457in}{0.956412in}}%
\pgfpathlineto{\pgfqpoint{2.153718in}{0.957101in}}%
\pgfpathlineto{\pgfqpoint{2.155342in}{0.959146in}}%
\pgfpathlineto{\pgfqpoint{2.156502in}{0.958231in}}%
\pgfpathlineto{\pgfqpoint{2.158422in}{0.956336in}}%
\pgfpathlineto{\pgfqpoint{2.158703in}{0.957086in}}%
\pgfpathlineto{\pgfqpoint{2.160255in}{0.959014in}}%
\pgfpathlineto{\pgfqpoint{2.161386in}{0.958153in}}%
\pgfpathlineto{\pgfqpoint{2.163360in}{0.956198in}}%
\pgfpathlineto{\pgfqpoint{2.163638in}{0.956931in}}%
\pgfpathlineto{\pgfqpoint{2.165232in}{0.958902in}}%
\pgfpathlineto{\pgfqpoint{2.166403in}{0.957979in}}%
\pgfpathlineto{\pgfqpoint{2.168313in}{0.956098in}}%
\pgfpathlineto{\pgfqpoint{2.168582in}{0.956819in}}%
\pgfpathlineto{\pgfqpoint{2.170148in}{0.958780in}}%
\pgfpathlineto{\pgfqpoint{2.171287in}{0.957910in}}%
\pgfpathlineto{\pgfqpoint{2.173263in}{0.955978in}}%
\pgfpathlineto{\pgfqpoint{2.173557in}{0.956728in}}%
\pgfpathlineto{\pgfqpoint{2.175116in}{0.958618in}}%
\pgfpathlineto{\pgfqpoint{2.176270in}{0.957715in}}%
\pgfpathlineto{\pgfqpoint{2.178191in}{0.955820in}}%
\pgfpathlineto{\pgfqpoint{2.178442in}{0.956492in}}%
\pgfpathlineto{\pgfqpoint{2.180024in}{0.958539in}}%
\pgfpathlineto{\pgfqpoint{2.181152in}{0.957698in}}%
\pgfpathlineto{\pgfqpoint{2.183150in}{0.955728in}}%
\pgfpathlineto{\pgfqpoint{2.183397in}{0.956396in}}%
\pgfpathlineto{\pgfqpoint{2.184996in}{0.958416in}}%
\pgfpathlineto{\pgfqpoint{2.186148in}{0.957527in}}%
\pgfpathlineto{\pgfqpoint{2.188099in}{0.955617in}}%
\pgfpathlineto{\pgfqpoint{2.188368in}{0.956343in}}%
\pgfpathlineto{\pgfqpoint{2.189951in}{0.958298in}}%
\pgfpathlineto{\pgfqpoint{2.191119in}{0.957384in}}%
\pgfpathlineto{\pgfqpoint{2.193031in}{0.955491in}}%
\pgfpathlineto{\pgfqpoint{2.193325in}{0.956258in}}%
\pgfpathlineto{\pgfqpoint{2.194908in}{0.958183in}}%
\pgfpathlineto{\pgfqpoint{2.196094in}{0.957240in}}%
\pgfpathlineto{\pgfqpoint{2.197984in}{0.955384in}}%
\pgfpathlineto{\pgfqpoint{2.198266in}{0.956141in}}%
\pgfpathlineto{\pgfqpoint{2.199842in}{0.958080in}}%
\pgfpathlineto{\pgfqpoint{2.200997in}{0.957180in}}%
\pgfpathlineto{\pgfqpoint{2.202929in}{0.955274in}}%
\pgfpathlineto{\pgfqpoint{2.203185in}{0.955956in}}%
\pgfpathlineto{\pgfqpoint{2.204768in}{0.957965in}}%
\pgfpathlineto{\pgfqpoint{2.205899in}{0.957108in}}%
\pgfpathlineto{\pgfqpoint{2.207869in}{0.955162in}}%
\pgfpathlineto{\pgfqpoint{2.208126in}{0.955843in}}%
\pgfpathlineto{\pgfqpoint{2.209708in}{0.957874in}}%
\pgfpathlineto{\pgfqpoint{2.210847in}{0.957018in}}%
\pgfpathlineto{\pgfqpoint{2.212812in}{0.955069in}}%
\pgfpathlineto{\pgfqpoint{2.213111in}{0.955830in}}%
\pgfpathlineto{\pgfqpoint{2.214663in}{0.957748in}}%
\pgfpathlineto{\pgfqpoint{2.215800in}{0.956882in}}%
\pgfpathlineto{\pgfqpoint{2.217771in}{0.954966in}}%
\pgfpathlineto{\pgfqpoint{2.218054in}{0.955706in}}%
\pgfpathlineto{\pgfqpoint{2.219615in}{0.957630in}}%
\pgfpathlineto{\pgfqpoint{2.220766in}{0.956745in}}%
\pgfpathlineto{\pgfqpoint{2.222706in}{0.954840in}}%
\pgfpathlineto{\pgfqpoint{2.222975in}{0.955560in}}%
\pgfpathlineto{\pgfqpoint{2.224558in}{0.957549in}}%
\pgfpathlineto{\pgfqpoint{2.225694in}{0.956681in}}%
\pgfpathlineto{\pgfqpoint{2.227663in}{0.954762in}}%
\pgfpathlineto{\pgfqpoint{2.227952in}{0.955506in}}%
\pgfpathlineto{\pgfqpoint{2.229495in}{0.957416in}}%
\pgfpathlineto{\pgfqpoint{2.230640in}{0.956548in}}%
\pgfpathlineto{\pgfqpoint{2.232586in}{0.954623in}}%
\pgfpathlineto{\pgfqpoint{2.232887in}{0.955390in}}%
\pgfpathlineto{\pgfqpoint{2.234423in}{0.957324in}}%
\pgfpathlineto{\pgfqpoint{2.235533in}{0.956508in}}%
\pgfpathlineto{\pgfqpoint{2.237554in}{0.954545in}}%
\pgfpathlineto{\pgfqpoint{2.237836in}{0.955285in}}%
\pgfpathlineto{\pgfqpoint{2.239417in}{0.957208in}}%
\pgfpathlineto{\pgfqpoint{2.240582in}{0.956291in}}%
\pgfpathlineto{\pgfqpoint{2.242477in}{0.954416in}}%
\pgfpathlineto{\pgfqpoint{2.242780in}{0.955180in}}%
\pgfpathlineto{\pgfqpoint{2.244312in}{0.957111in}}%
\pgfpathlineto{\pgfqpoint{2.245423in}{0.956296in}}%
\pgfpathlineto{\pgfqpoint{2.247444in}{0.954336in}}%
\pgfpathlineto{\pgfqpoint{2.247729in}{0.955083in}}%
\pgfpathlineto{\pgfqpoint{2.249291in}{0.957005in}}%
\pgfpathlineto{\pgfqpoint{2.250443in}{0.956118in}}%
\pgfpathlineto{\pgfqpoint{2.252389in}{0.954234in}}%
\pgfpathlineto{\pgfqpoint{2.252673in}{0.954980in}}%
\pgfpathlineto{\pgfqpoint{2.254239in}{0.956905in}}%
\pgfpathlineto{\pgfqpoint{2.255388in}{0.956019in}}%
\pgfpathlineto{\pgfqpoint{2.257335in}{0.954135in}}%
\pgfpathlineto{\pgfqpoint{2.257618in}{0.954880in}}%
\pgfpathlineto{\pgfqpoint{2.259184in}{0.956806in}}%
\pgfpathlineto{\pgfqpoint{2.260333in}{0.955921in}}%
\pgfpathlineto{\pgfqpoint{2.262280in}{0.954037in}}%
\pgfpathlineto{\pgfqpoint{2.262564in}{0.954783in}}%
\pgfpathlineto{\pgfqpoint{2.264130in}{0.956709in}}%
\pgfpathlineto{\pgfqpoint{2.265279in}{0.955824in}}%
\pgfpathlineto{\pgfqpoint{2.267226in}{0.953941in}}%
\pgfpathlineto{\pgfqpoint{2.267510in}{0.954687in}}%
\pgfpathlineto{\pgfqpoint{2.269076in}{0.956613in}}%
\pgfpathlineto{\pgfqpoint{2.270224in}{0.955729in}}%
\pgfpathlineto{\pgfqpoint{2.272171in}{0.953846in}}%
\pgfpathlineto{\pgfqpoint{2.272455in}{0.954593in}}%
\pgfpathlineto{\pgfqpoint{2.274022in}{0.956518in}}%
\pgfpathlineto{\pgfqpoint{2.275169in}{0.955635in}}%
\pgfpathlineto{\pgfqpoint{2.277117in}{0.953752in}}%
\pgfpathlineto{\pgfqpoint{2.277401in}{0.954500in}}%
\pgfpathlineto{\pgfqpoint{2.278969in}{0.956425in}}%
\pgfpathlineto{\pgfqpoint{2.280114in}{0.955542in}}%
\pgfpathlineto{\pgfqpoint{2.282062in}{0.953660in}}%
\pgfpathlineto{\pgfqpoint{2.282346in}{0.954405in}}%
\pgfpathlineto{\pgfqpoint{2.283918in}{0.956333in}}%
\pgfpathlineto{\pgfqpoint{2.285086in}{0.955425in}}%
\pgfpathlineto{\pgfqpoint{2.287008in}{0.953568in}}%
\pgfpathlineto{\pgfqpoint{2.287292in}{0.954315in}}%
\pgfpathlineto{\pgfqpoint{2.288859in}{0.956241in}}%
\pgfpathlineto{\pgfqpoint{2.290005in}{0.955360in}}%
\pgfpathlineto{\pgfqpoint{2.291953in}{0.953478in}}%
\pgfpathlineto{\pgfqpoint{2.292235in}{0.954221in}}%
\pgfpathlineto{\pgfqpoint{2.293773in}{0.956155in}}%
\pgfpathlineto{\pgfqpoint{2.294897in}{0.955325in}}%
\pgfpathlineto{\pgfqpoint{2.296886in}{0.953372in}}%
\pgfpathlineto{\pgfqpoint{2.297139in}{0.954045in}}%
\pgfpathlineto{\pgfqpoint{2.298722in}{0.956085in}}%
\pgfpathlineto{\pgfqpoint{2.299846in}{0.955253in}}%
\pgfpathlineto{\pgfqpoint{2.301845in}{0.953310in}}%
\pgfpathlineto{\pgfqpoint{2.302136in}{0.954054in}}%
\pgfpathlineto{\pgfqpoint{2.303704in}{0.955961in}}%
\pgfpathlineto{\pgfqpoint{2.304874in}{0.955048in}}%
\pgfpathlineto{\pgfqpoint{2.306772in}{0.953183in}}%
\pgfpathlineto{\pgfqpoint{2.307029in}{0.953869in}}%
\pgfpathlineto{\pgfqpoint{2.308612in}{0.955917in}}%
\pgfpathlineto{\pgfqpoint{2.309740in}{0.955084in}}%
\pgfpathlineto{\pgfqpoint{2.311727in}{0.953131in}}%
\pgfpathlineto{\pgfqpoint{2.312016in}{0.953894in}}%
\pgfpathlineto{\pgfqpoint{2.313573in}{0.955839in}}%
\pgfpathlineto{\pgfqpoint{2.314708in}{0.954986in}}%
\pgfpathlineto{\pgfqpoint{2.316681in}{0.953071in}}%
\pgfpathlineto{\pgfqpoint{2.316971in}{0.953815in}}%
\pgfpathlineto{\pgfqpoint{2.318553in}{0.955722in}}%
\pgfpathlineto{\pgfqpoint{2.319751in}{0.954771in}}%
\pgfpathlineto{\pgfqpoint{2.321615in}{0.952958in}}%
\pgfpathlineto{\pgfqpoint{2.321875in}{0.953665in}}%
\pgfpathlineto{\pgfqpoint{2.323503in}{0.955693in}}%
\pgfpathlineto{\pgfqpoint{2.324689in}{0.954755in}}%
\pgfpathlineto{\pgfqpoint{2.326560in}{0.952914in}}%
\pgfpathlineto{\pgfqpoint{2.326821in}{0.953598in}}%
\pgfpathlineto{\pgfqpoint{2.328404in}{0.955623in}}%
\pgfpathlineto{\pgfqpoint{2.329528in}{0.954785in}}%
\pgfpathlineto{\pgfqpoint{2.331518in}{0.952856in}}%
\pgfpathlineto{\pgfqpoint{2.331807in}{0.953604in}}%
\pgfpathlineto{\pgfqpoint{2.333375in}{0.955513in}}%
\pgfpathlineto{\pgfqpoint{2.334532in}{0.954617in}}%
\pgfpathlineto{\pgfqpoint{2.336454in}{0.952751in}}%
\pgfpathlineto{\pgfqpoint{2.336705in}{0.953437in}}%
\pgfpathlineto{\pgfqpoint{2.338297in}{0.955487in}}%
\pgfpathlineto{\pgfqpoint{2.339429in}{0.954643in}}%
\pgfpathlineto{\pgfqpoint{2.341408in}{0.952715in}}%
\pgfpathlineto{\pgfqpoint{2.341706in}{0.953461in}}%
\pgfpathlineto{\pgfqpoint{2.343304in}{0.955341in}}%
\pgfpathlineto{\pgfqpoint{2.344535in}{0.954337in}}%
\pgfpathlineto{\pgfqpoint{2.346331in}{0.952578in}}%
\pgfpathlineto{\pgfqpoint{2.346623in}{0.953347in}}%
\pgfpathlineto{\pgfqpoint{2.348226in}{0.955317in}}%
\pgfpathlineto{\pgfqpoint{2.349411in}{0.954383in}}%
\pgfpathlineto{\pgfqpoint{2.351291in}{0.952541in}}%
\pgfpathlineto{\pgfqpoint{2.351583in}{0.953297in}}%
\pgfpathlineto{\pgfqpoint{2.353193in}{0.955231in}}%
\pgfpathlineto{\pgfqpoint{2.354413in}{0.954245in}}%
\pgfpathlineto{\pgfqpoint{2.356234in}{0.952470in}}%
\pgfpathlineto{\pgfqpoint{2.356517in}{0.953232in}}%
\pgfpathlineto{\pgfqpoint{2.358115in}{0.955198in}}%
\pgfpathlineto{\pgfqpoint{2.359304in}{0.954263in}}%
\pgfpathlineto{\pgfqpoint{2.361189in}{0.952430in}}%
\pgfpathlineto{\pgfqpoint{2.361481in}{0.953178in}}%
\pgfpathlineto{\pgfqpoint{2.363051in}{0.955089in}}%
\pgfpathlineto{\pgfqpoint{2.364219in}{0.954181in}}%
\pgfpathlineto{\pgfqpoint{2.366124in}{0.952326in}}%
\pgfpathlineto{\pgfqpoint{2.366402in}{0.953079in}}%
\pgfpathlineto{\pgfqpoint{2.367985in}{0.955066in}}%
\pgfpathlineto{\pgfqpoint{2.369130in}{0.954197in}}%
\pgfpathlineto{\pgfqpoint{2.371074in}{0.952288in}}%
\pgfpathlineto{\pgfqpoint{2.371351in}{0.952956in}}%
\pgfpathlineto{\pgfqpoint{2.372980in}{0.954907in}}%
\pgfpathlineto{\pgfqpoint{2.374203in}{0.953909in}}%
\pgfpathlineto{\pgfqpoint{2.376003in}{0.952158in}}%
\pgfpathlineto{\pgfqpoint{2.376254in}{0.952837in}}%
\pgfpathlineto{\pgfqpoint{2.377858in}{0.954929in}}%
\pgfpathlineto{\pgfqpoint{2.378986in}{0.954095in}}%
\pgfpathlineto{\pgfqpoint{2.380964in}{0.952156in}}%
\pgfpathlineto{\pgfqpoint{2.381266in}{0.952909in}}%
\pgfpathlineto{\pgfqpoint{2.382850in}{0.954810in}}%
\pgfpathlineto{\pgfqpoint{2.384052in}{0.953853in}}%
\pgfpathlineto{\pgfqpoint{2.385894in}{0.952047in}}%
\pgfpathlineto{\pgfqpoint{2.386173in}{0.952772in}}%
\pgfpathlineto{\pgfqpoint{2.387770in}{0.954788in}}%
\pgfpathlineto{\pgfqpoint{2.388910in}{0.953922in}}%
\pgfpathlineto{\pgfqpoint{2.390856in}{0.952027in}}%
\pgfpathlineto{\pgfqpoint{2.391123in}{0.952732in}}%
\pgfpathlineto{\pgfqpoint{2.392741in}{0.954720in}}%
\pgfpathlineto{\pgfqpoint{2.393926in}{0.953781in}}%
\pgfpathlineto{\pgfqpoint{2.395809in}{0.951968in}}%
\pgfpathlineto{\pgfqpoint{2.396093in}{0.952713in}}%
\pgfpathlineto{\pgfqpoint{2.397637in}{0.954647in}}%
\pgfpathlineto{\pgfqpoint{2.398766in}{0.953811in}}%
\pgfpathlineto{\pgfqpoint{2.400754in}{0.951898in}}%
\pgfpathlineto{\pgfqpoint{2.401033in}{0.952637in}}%
\pgfpathlineto{\pgfqpoint{2.402576in}{0.954582in}}%
\pgfpathlineto{\pgfqpoint{2.403705in}{0.953751in}}%
\pgfpathlineto{\pgfqpoint{2.405691in}{0.951820in}}%
\pgfpathlineto{\pgfqpoint{2.405975in}{0.952576in}}%
\pgfpathlineto{\pgfqpoint{2.407531in}{0.954543in}}%
\pgfpathlineto{\pgfqpoint{2.408664in}{0.953704in}}%
\pgfpathlineto{\pgfqpoint{2.410638in}{0.951778in}}%
\pgfpathlineto{\pgfqpoint{2.410910in}{0.952507in}}%
\pgfpathlineto{\pgfqpoint{2.412481in}{0.954487in}}%
\pgfpathlineto{\pgfqpoint{2.413620in}{0.953633in}}%
\pgfpathlineto{\pgfqpoint{2.415581in}{0.951713in}}%
\pgfpathlineto{\pgfqpoint{2.415867in}{0.952468in}}%
\pgfpathlineto{\pgfqpoint{2.417469in}{0.954429in}}%
\pgfpathlineto{\pgfqpoint{2.418646in}{0.953505in}}%
\pgfpathlineto{\pgfqpoint{2.420525in}{0.951667in}}%
\pgfpathlineto{\pgfqpoint{2.420821in}{0.952429in}}%
\pgfpathlineto{\pgfqpoint{2.422421in}{0.954358in}}%
\pgfpathlineto{\pgfqpoint{2.423629in}{0.953394in}}%
\pgfpathlineto{\pgfqpoint{2.425482in}{0.951612in}}%
\pgfpathlineto{\pgfqpoint{2.425771in}{0.952354in}}%
\pgfpathlineto{\pgfqpoint{2.427394in}{0.954265in}}%
\pgfpathlineto{\pgfqpoint{2.428650in}{0.953229in}}%
\pgfpathlineto{\pgfqpoint{2.430404in}{0.951512in}}%
\pgfpathlineto{\pgfqpoint{2.430669in}{0.952218in}}%
\pgfpathlineto{\pgfqpoint{2.432287in}{0.954269in}}%
\pgfpathlineto{\pgfqpoint{2.433436in}{0.953389in}}%
\pgfpathlineto{\pgfqpoint{2.435369in}{0.951502in}}%
\pgfpathlineto{\pgfqpoint{2.435654in}{0.952265in}}%
\pgfpathlineto{\pgfqpoint{2.437232in}{0.954199in}}%
\pgfpathlineto{\pgfqpoint{2.438409in}{0.953286in}}%
\pgfpathlineto{\pgfqpoint{2.440316in}{0.951437in}}%
\pgfpathlineto{\pgfqpoint{2.440602in}{0.952176in}}%
\pgfpathlineto{\pgfqpoint{2.442160in}{0.954111in}}%
\pgfpathlineto{\pgfqpoint{2.443303in}{0.953251in}}%
\pgfpathlineto{\pgfqpoint{2.445240in}{0.951348in}}%
\pgfpathlineto{\pgfqpoint{2.445541in}{0.952114in}}%
\pgfpathlineto{\pgfqpoint{2.447078in}{0.954065in}}%
\pgfpathlineto{\pgfqpoint{2.448197in}{0.953254in}}%
\pgfpathlineto{\pgfqpoint{2.450207in}{0.951309in}}%
\pgfpathlineto{\pgfqpoint{2.450482in}{0.952010in}}%
\pgfpathlineto{\pgfqpoint{2.452118in}{0.953964in}}%
\pgfpathlineto{\pgfqpoint{2.453354in}{0.952952in}}%
\pgfpathlineto{\pgfqpoint{2.455129in}{0.951211in}}%
\pgfpathlineto{\pgfqpoint{2.455418in}{0.951932in}}%
\pgfpathlineto{\pgfqpoint{2.457036in}{0.953914in}}%
\pgfpathlineto{\pgfqpoint{2.458244in}{0.952953in}}%
\pgfpathlineto{\pgfqpoint{2.460089in}{0.951164in}}%
\pgfpathlineto{\pgfqpoint{2.460364in}{0.951908in}}%
\pgfpathlineto{\pgfqpoint{2.461944in}{0.953906in}}%
\pgfpathlineto{\pgfqpoint{2.463067in}{0.953068in}}%
\pgfpathlineto{\pgfqpoint{2.465042in}{0.951142in}}%
\pgfpathlineto{\pgfqpoint{2.465284in}{0.951791in}}%
\pgfpathlineto{\pgfqpoint{2.466851in}{0.953840in}}%
\pgfpathlineto{\pgfqpoint{2.467949in}{0.953059in}}%
\pgfpathlineto{\pgfqpoint{2.469988in}{0.951080in}}%
\pgfpathlineto{\pgfqpoint{2.470276in}{0.951824in}}%
\pgfpathlineto{\pgfqpoint{2.471825in}{0.953759in}}%
\pgfpathlineto{\pgfqpoint{2.472957in}{0.952919in}}%
\pgfpathlineto{\pgfqpoint{2.474929in}{0.951011in}}%
\pgfpathlineto{\pgfqpoint{2.475207in}{0.951708in}}%
\pgfpathlineto{\pgfqpoint{2.476816in}{0.953670in}}%
\pgfpathlineto{\pgfqpoint{2.478002in}{0.952732in}}%
\pgfpathlineto{\pgfqpoint{2.479858in}{0.950918in}}%
\pgfpathlineto{\pgfqpoint{2.480150in}{0.951661in}}%
\pgfpathlineto{\pgfqpoint{2.481762in}{0.953634in}}%
\pgfpathlineto{\pgfqpoint{2.482952in}{0.952693in}}%
\pgfpathlineto{\pgfqpoint{2.484817in}{0.950886in}}%
\pgfpathlineto{\pgfqpoint{2.485094in}{0.951632in}}%
\pgfpathlineto{\pgfqpoint{2.486717in}{0.953618in}}%
\pgfpathlineto{\pgfqpoint{2.487900in}{0.952679in}}%
\pgfpathlineto{\pgfqpoint{2.489761in}{0.950858in}}%
\pgfpathlineto{\pgfqpoint{2.490026in}{0.951558in}}%
\pgfpathlineto{\pgfqpoint{2.491650in}{0.953578in}}%
\pgfpathlineto{\pgfqpoint{2.492827in}{0.952656in}}%
\pgfpathlineto{\pgfqpoint{2.494707in}{0.950821in}}%
\pgfpathlineto{\pgfqpoint{2.494968in}{0.951519in}}%
\pgfpathlineto{\pgfqpoint{2.496551in}{0.953549in}}%
\pgfpathlineto{\pgfqpoint{2.497669in}{0.952727in}}%
\pgfpathlineto{\pgfqpoint{2.499653in}{0.950785in}}%
\pgfpathlineto{\pgfqpoint{2.499910in}{0.951471in}}%
\pgfpathlineto{\pgfqpoint{2.501493in}{0.953510in}}%
\pgfpathlineto{\pgfqpoint{2.502628in}{0.952675in}}%
\pgfpathlineto{\pgfqpoint{2.504602in}{0.950746in}}%
\pgfpathlineto{\pgfqpoint{2.504893in}{0.951516in}}%
\pgfpathlineto{\pgfqpoint{2.506493in}{0.953450in}}%
\pgfpathlineto{\pgfqpoint{2.507699in}{0.952491in}}%
\pgfpathlineto{\pgfqpoint{2.509542in}{0.950690in}}%
\pgfpathlineto{\pgfqpoint{2.509806in}{0.951389in}}%
\pgfpathlineto{\pgfqpoint{2.511421in}{0.953420in}}%
\pgfpathlineto{\pgfqpoint{2.512571in}{0.952536in}}%
\pgfpathlineto{\pgfqpoint{2.514497in}{0.950669in}}%
\pgfpathlineto{\pgfqpoint{2.514751in}{0.951353in}}%
\pgfpathlineto{\pgfqpoint{2.516376in}{0.953367in}}%
\pgfpathlineto{\pgfqpoint{2.517562in}{0.952435in}}%
\pgfpathlineto{\pgfqpoint{2.519436in}{0.950608in}}%
\pgfpathlineto{\pgfqpoint{2.519715in}{0.951351in}}%
\pgfpathlineto{\pgfqpoint{2.521259in}{0.953337in}}%
\pgfpathlineto{\pgfqpoint{2.522359in}{0.952556in}}%
\pgfpathlineto{\pgfqpoint{2.524382in}{0.950582in}}%
\pgfpathlineto{\pgfqpoint{2.524649in}{0.951289in}}%
\pgfpathlineto{\pgfqpoint{2.526225in}{0.953289in}}%
\pgfpathlineto{\pgfqpoint{2.527352in}{0.952454in}}%
\pgfpathlineto{\pgfqpoint{2.529329in}{0.950524in}}%
\pgfpathlineto{\pgfqpoint{2.529606in}{0.951195in}}%
\pgfpathlineto{\pgfqpoint{2.531242in}{0.953161in}}%
\pgfpathlineto{\pgfqpoint{2.532461in}{0.952173in}}%
\pgfpathlineto{\pgfqpoint{2.534256in}{0.950422in}}%
\pgfpathlineto{\pgfqpoint{2.534551in}{0.951180in}}%
\pgfpathlineto{\pgfqpoint{2.536160in}{0.953152in}}%
\pgfpathlineto{\pgfqpoint{2.537355in}{0.952210in}}%
\pgfpathlineto{\pgfqpoint{2.539220in}{0.950398in}}%
\pgfpathlineto{\pgfqpoint{2.539497in}{0.951077in}}%
\pgfpathlineto{\pgfqpoint{2.541142in}{0.953040in}}%
\pgfpathlineto{\pgfqpoint{2.542398in}{0.952008in}}%
\pgfpathlineto{\pgfqpoint{2.544147in}{0.950300in}}%
\pgfpathlineto{\pgfqpoint{2.544440in}{0.951031in}}%
\pgfpathlineto{\pgfqpoint{2.546052in}{0.953009in}}%
\pgfpathlineto{\pgfqpoint{2.547247in}{0.952068in}}%
\pgfpathlineto{\pgfqpoint{2.549095in}{0.950263in}}%
\pgfpathlineto{\pgfqpoint{2.549387in}{0.951019in}}%
\pgfpathlineto{\pgfqpoint{2.550982in}{0.952995in}}%
\pgfpathlineto{\pgfqpoint{2.552137in}{0.952108in}}%
\pgfpathlineto{\pgfqpoint{2.554059in}{0.950244in}}%
\pgfpathlineto{\pgfqpoint{2.554334in}{0.950932in}}%
\pgfpathlineto{\pgfqpoint{2.555945in}{0.952903in}}%
\pgfpathlineto{\pgfqpoint{2.557139in}{0.951961in}}%
\pgfpathlineto{\pgfqpoint{2.558986in}{0.950165in}}%
\pgfpathlineto{\pgfqpoint{2.559266in}{0.950908in}}%
\pgfpathlineto{\pgfqpoint{2.560892in}{0.952922in}}%
\pgfpathlineto{\pgfqpoint{2.562078in}{0.951990in}}%
\pgfpathlineto{\pgfqpoint{2.563951in}{0.950169in}}%
\pgfpathlineto{\pgfqpoint{2.564230in}{0.950919in}}%
\pgfpathlineto{\pgfqpoint{2.565830in}{0.952875in}}%
\pgfpathlineto{\pgfqpoint{2.566999in}{0.951963in}}%
\pgfpathlineto{\pgfqpoint{2.568893in}{0.950121in}}%
\pgfpathlineto{\pgfqpoint{2.569171in}{0.950866in}}%
\pgfpathlineto{\pgfqpoint{2.570737in}{0.952844in}}%
\pgfpathlineto{\pgfqpoint{2.571866in}{0.952009in}}%
\pgfpathlineto{\pgfqpoint{2.573839in}{0.950088in}}%
\pgfpathlineto{\pgfqpoint{2.574118in}{0.950834in}}%
\pgfpathlineto{\pgfqpoint{2.575691in}{0.952803in}}%
\pgfpathlineto{\pgfqpoint{2.576838in}{0.951941in}}%
\pgfpathlineto{\pgfqpoint{2.578784in}{0.950046in}}%
\pgfpathlineto{\pgfqpoint{2.579069in}{0.950808in}}%
\pgfpathlineto{\pgfqpoint{2.580654in}{0.952768in}}%
\pgfpathlineto{\pgfqpoint{2.581812in}{0.951882in}}%
\pgfpathlineto{\pgfqpoint{2.583731in}{0.950017in}}%
\pgfpathlineto{\pgfqpoint{2.584018in}{0.950783in}}%
\pgfpathlineto{\pgfqpoint{2.585596in}{0.952729in}}%
\pgfpathlineto{\pgfqpoint{2.586758in}{0.951840in}}%
\pgfpathlineto{\pgfqpoint{2.588675in}{0.949975in}}%
\pgfpathlineto{\pgfqpoint{2.588949in}{0.950709in}}%
\pgfpathlineto{\pgfqpoint{2.590557in}{0.952697in}}%
\pgfpathlineto{\pgfqpoint{2.591732in}{0.951783in}}%
\pgfpathlineto{\pgfqpoint{2.593624in}{0.949950in}}%
\pgfpathlineto{\pgfqpoint{2.593889in}{0.950661in}}%
\pgfpathlineto{\pgfqpoint{2.595464in}{0.952658in}}%
\pgfpathlineto{\pgfqpoint{2.596598in}{0.951817in}}%
\pgfpathlineto{\pgfqpoint{2.598563in}{0.949903in}}%
\pgfpathlineto{\pgfqpoint{2.598846in}{0.950660in}}%
\pgfpathlineto{\pgfqpoint{2.600444in}{0.952636in}}%
\pgfpathlineto{\pgfqpoint{2.601626in}{0.951721in}}%
\pgfpathlineto{\pgfqpoint{2.603506in}{0.949885in}}%
\pgfpathlineto{\pgfqpoint{2.603781in}{0.950613in}}%
\pgfpathlineto{\pgfqpoint{2.605363in}{0.952615in}}%
\pgfpathlineto{\pgfqpoint{2.606519in}{0.951749in}}%
\pgfpathlineto{\pgfqpoint{2.608461in}{0.949868in}}%
\pgfpathlineto{\pgfqpoint{2.608734in}{0.950600in}}%
\pgfpathlineto{\pgfqpoint{2.610326in}{0.952573in}}%
\pgfpathlineto{\pgfqpoint{2.611484in}{0.951686in}}%
\pgfpathlineto{\pgfqpoint{2.613402in}{0.949822in}}%
\pgfpathlineto{\pgfqpoint{2.613680in}{0.950569in}}%
\pgfpathlineto{\pgfqpoint{2.615251in}{0.952550in}}%
\pgfpathlineto{\pgfqpoint{2.616385in}{0.951708in}}%
\pgfpathlineto{\pgfqpoint{2.618351in}{0.949796in}}%
\pgfpathlineto{\pgfqpoint{2.618635in}{0.950560in}}%
\pgfpathlineto{\pgfqpoint{2.620256in}{0.952503in}}%
\pgfpathlineto{\pgfqpoint{2.621481in}{0.951516in}}%
\pgfpathlineto{\pgfqpoint{2.623292in}{0.949753in}}%
\pgfpathlineto{\pgfqpoint{2.623570in}{0.950429in}}%
\pgfpathlineto{\pgfqpoint{2.625172in}{0.952406in}}%
\pgfpathlineto{\pgfqpoint{2.626344in}{0.951498in}}%
\pgfpathlineto{\pgfqpoint{2.628220in}{0.949674in}}%
\pgfpathlineto{\pgfqpoint{2.628493in}{0.950405in}}%
\pgfpathlineto{\pgfqpoint{2.630072in}{0.952441in}}%
\pgfpathlineto{\pgfqpoint{2.631205in}{0.951613in}}%
\pgfpathlineto{\pgfqpoint{2.633188in}{0.949684in}}%
\pgfpathlineto{\pgfqpoint{2.633432in}{0.950338in}}%
\pgfpathlineto{\pgfqpoint{2.635041in}{0.952395in}}%
\pgfpathlineto{\pgfqpoint{2.636182in}{0.951538in}}%
\pgfpathlineto{\pgfqpoint{2.638125in}{0.949644in}}%
\pgfpathlineto{\pgfqpoint{2.638415in}{0.950414in}}%
\pgfpathlineto{\pgfqpoint{2.640009in}{0.952380in}}%
\pgfpathlineto{\pgfqpoint{2.641167in}{0.951488in}}%
\pgfpathlineto{\pgfqpoint{2.643069in}{0.949630in}}%
\pgfpathlineto{\pgfqpoint{2.643339in}{0.950341in}}%
\pgfpathlineto{\pgfqpoint{2.644922in}{0.952360in}}%
\pgfpathlineto{\pgfqpoint{2.646058in}{0.951519in}}%
\pgfpathlineto{\pgfqpoint{2.648021in}{0.949609in}}%
\pgfpathlineto{\pgfqpoint{2.648297in}{0.950285in}}%
\pgfpathlineto{\pgfqpoint{2.649901in}{0.952256in}}%
\pgfpathlineto{\pgfqpoint{2.651081in}{0.951336in}}%
\pgfpathlineto{\pgfqpoint{2.652949in}{0.949526in}}%
\pgfpathlineto{\pgfqpoint{2.653228in}{0.950275in}}%
\pgfpathlineto{\pgfqpoint{2.654825in}{0.952302in}}%
\pgfpathlineto{\pgfqpoint{2.655973in}{0.951440in}}%
\pgfpathlineto{\pgfqpoint{2.657917in}{0.949558in}}%
\pgfpathlineto{\pgfqpoint{2.658215in}{0.950308in}}%
\pgfpathlineto{\pgfqpoint{2.659784in}{0.952213in}}%
\pgfpathlineto{\pgfqpoint{2.660962in}{0.951302in}}%
\pgfpathlineto{\pgfqpoint{2.662845in}{0.949476in}}%
\pgfpathlineto{\pgfqpoint{2.663116in}{0.950204in}}%
\pgfpathlineto{\pgfqpoint{2.664688in}{0.952233in}}%
\pgfpathlineto{\pgfqpoint{2.665806in}{0.951425in}}%
\pgfpathlineto{\pgfqpoint{2.667796in}{0.949478in}}%
\pgfpathlineto{\pgfqpoint{2.668076in}{0.950214in}}%
\pgfpathlineto{\pgfqpoint{2.669659in}{0.952211in}}%
\pgfpathlineto{\pgfqpoint{2.670800in}{0.951358in}}%
\pgfpathlineto{\pgfqpoint{2.672749in}{0.949463in}}%
\pgfpathlineto{\pgfqpoint{2.673026in}{0.950135in}}%
\pgfpathlineto{\pgfqpoint{2.674651in}{0.952098in}}%
\pgfpathlineto{\pgfqpoint{2.675874in}{0.951116in}}%
\pgfpathlineto{\pgfqpoint{2.677682in}{0.949375in}}%
\pgfpathlineto{\pgfqpoint{2.677952in}{0.950109in}}%
\pgfpathlineto{\pgfqpoint{2.679535in}{0.952141in}}%
\pgfpathlineto{\pgfqpoint{2.680655in}{0.951322in}}%
\pgfpathlineto{\pgfqpoint{2.682633in}{0.949387in}}%
\pgfpathlineto{\pgfqpoint{2.682900in}{0.950084in}}%
\pgfpathlineto{\pgfqpoint{2.684475in}{0.952113in}}%
\pgfpathlineto{\pgfqpoint{2.685593in}{0.951302in}}%
\pgfpathlineto{\pgfqpoint{2.687584in}{0.949364in}}%
\pgfpathlineto{\pgfqpoint{2.687861in}{0.950038in}}%
\pgfpathlineto{\pgfqpoint{2.689501in}{0.952004in}}%
\pgfpathlineto{\pgfqpoint{2.690734in}{0.951003in}}%
\pgfpathlineto{\pgfqpoint{2.692518in}{0.949281in}}%
\pgfpathlineto{\pgfqpoint{2.692772in}{0.949972in}}%
\pgfpathlineto{\pgfqpoint{2.694404in}{0.952048in}}%
\pgfpathlineto{\pgfqpoint{2.695578in}{0.951146in}}%
\pgfpathlineto{\pgfqpoint{2.697476in}{0.949305in}}%
\pgfpathlineto{\pgfqpoint{2.697744in}{0.950012in}}%
\pgfpathlineto{\pgfqpoint{2.699371in}{0.952010in}}%
\pgfpathlineto{\pgfqpoint{2.700562in}{0.951069in}}%
\pgfpathlineto{\pgfqpoint{2.702422in}{0.949270in}}%
\pgfpathlineto{\pgfqpoint{2.702698in}{0.949957in}}%
\pgfpathlineto{\pgfqpoint{2.704311in}{0.951931in}}%
\pgfpathlineto{\pgfqpoint{2.705503in}{0.950993in}}%
\pgfpathlineto{\pgfqpoint{2.707355in}{0.949205in}}%
\pgfpathlineto{\pgfqpoint{2.707618in}{0.949917in}}%
\pgfpathlineto{\pgfqpoint{2.709201in}{0.951958in}}%
\pgfpathlineto{\pgfqpoint{2.710324in}{0.951141in}}%
\pgfpathlineto{\pgfqpoint{2.712310in}{0.949201in}}%
\pgfpathlineto{\pgfqpoint{2.712588in}{0.949870in}}%
\pgfpathlineto{\pgfqpoint{2.714209in}{0.951847in}}%
\pgfpathlineto{\pgfqpoint{2.715411in}{0.950894in}}%
\pgfpathlineto{\pgfqpoint{2.717244in}{0.949126in}}%
\pgfpathlineto{\pgfqpoint{2.717503in}{0.949829in}}%
\pgfpathlineto{\pgfqpoint{2.719086in}{0.951888in}}%
\pgfpathlineto{\pgfqpoint{2.720200in}{0.951085in}}%
\pgfpathlineto{\pgfqpoint{2.722201in}{0.949134in}}%
\pgfpathlineto{\pgfqpoint{2.722479in}{0.949807in}}%
\pgfpathlineto{\pgfqpoint{2.724091in}{0.951785in}}%
\pgfpathlineto{\pgfqpoint{2.725282in}{0.950852in}}%
\pgfpathlineto{\pgfqpoint{2.727130in}{0.949062in}}%
\pgfpathlineto{\pgfqpoint{2.727411in}{0.949816in}}%
\pgfpathlineto{\pgfqpoint{2.728994in}{0.951839in}}%
\pgfpathlineto{\pgfqpoint{2.730108in}{0.951025in}}%
\pgfpathlineto{\pgfqpoint{2.732093in}{0.949086in}}%
\pgfpathlineto{\pgfqpoint{2.732403in}{0.949835in}}%
\pgfpathlineto{\pgfqpoint{2.733902in}{0.951733in}}%
\pgfpathlineto{\pgfqpoint{2.735014in}{0.950953in}}%
\pgfpathlineto{\pgfqpoint{2.737022in}{0.949006in}}%
\pgfpathlineto{\pgfqpoint{2.737303in}{0.949759in}}%
\pgfpathlineto{\pgfqpoint{2.738869in}{0.951773in}}%
\pgfpathlineto{\pgfqpoint{2.739991in}{0.950963in}}%
\pgfpathlineto{\pgfqpoint{2.741985in}{0.949020in}}%
\pgfpathlineto{\pgfqpoint{2.742289in}{0.949762in}}%
\pgfpathlineto{\pgfqpoint{2.743776in}{0.951676in}}%
\pgfpathlineto{\pgfqpoint{2.744870in}{0.950934in}}%
\pgfpathlineto{\pgfqpoint{2.746912in}{0.948950in}}%
\pgfpathlineto{\pgfqpoint{2.747175in}{0.949657in}}%
\pgfpathlineto{\pgfqpoint{2.748781in}{0.951726in}}%
\pgfpathlineto{\pgfqpoint{2.749913in}{0.950889in}}%
\pgfpathlineto{\pgfqpoint{2.751881in}{0.948985in}}%
\pgfpathlineto{\pgfqpoint{2.752176in}{0.949731in}}%
\pgfpathlineto{\pgfqpoint{2.753754in}{0.951647in}}%
\pgfpathlineto{\pgfqpoint{2.754925in}{0.950742in}}%
\pgfpathlineto{\pgfqpoint{2.756810in}{0.948916in}}%
\pgfpathlineto{\pgfqpoint{2.757094in}{0.949676in}}%
\pgfpathlineto{\pgfqpoint{2.758677in}{0.951667in}}%
\pgfpathlineto{\pgfqpoint{2.759811in}{0.950820in}}%
\pgfpathlineto{\pgfqpoint{2.761763in}{0.948919in}}%
\pgfpathlineto{\pgfqpoint{2.762046in}{0.949673in}}%
\pgfpathlineto{\pgfqpoint{2.763636in}{0.951654in}}%
\pgfpathlineto{\pgfqpoint{2.764798in}{0.950770in}}%
\pgfpathlineto{\pgfqpoint{2.766717in}{0.948922in}}%
\pgfpathlineto{\pgfqpoint{2.767010in}{0.949670in}}%
\pgfpathlineto{\pgfqpoint{2.768611in}{0.951592in}}%
\pgfpathlineto{\pgfqpoint{2.769811in}{0.950639in}}%
\pgfpathlineto{\pgfqpoint{2.771645in}{0.948855in}}%
\pgfpathlineto{\pgfqpoint{2.771920in}{0.949585in}}%
\pgfpathlineto{\pgfqpoint{2.773503in}{0.951607in}}%
\pgfpathlineto{\pgfqpoint{2.774634in}{0.950774in}}%
\pgfpathlineto{\pgfqpoint{2.776604in}{0.948859in}}%
\pgfpathlineto{\pgfqpoint{2.776880in}{0.949532in}}%
\pgfpathlineto{\pgfqpoint{2.778503in}{0.951502in}}%
\pgfpathlineto{\pgfqpoint{2.779708in}{0.950545in}}%
\pgfpathlineto{\pgfqpoint{2.781530in}{0.948784in}}%
\pgfpathlineto{\pgfqpoint{2.781767in}{0.949425in}}%
\pgfpathlineto{\pgfqpoint{2.783398in}{0.951570in}}%
\pgfpathlineto{\pgfqpoint{2.784527in}{0.950739in}}%
\pgfpathlineto{\pgfqpoint{2.786498in}{0.948831in}}%
\pgfpathlineto{\pgfqpoint{2.786789in}{0.949569in}}%
\pgfpathlineto{\pgfqpoint{2.788347in}{0.951499in}}%
\pgfpathlineto{\pgfqpoint{2.789492in}{0.950643in}}%
\pgfpathlineto{\pgfqpoint{2.791436in}{0.948774in}}%
\pgfpathlineto{\pgfqpoint{2.791691in}{0.949473in}}%
\pgfpathlineto{\pgfqpoint{2.793305in}{0.951519in}}%
\pgfpathlineto{\pgfqpoint{2.794458in}{0.950645in}}%
\pgfpathlineto{\pgfqpoint{2.796388in}{0.948773in}}%
\pgfpathlineto{\pgfqpoint{2.796630in}{0.949425in}}%
\pgfpathlineto{\pgfqpoint{2.798239in}{0.951493in}}%
\pgfpathlineto{\pgfqpoint{2.799385in}{0.950639in}}%
\pgfpathlineto{\pgfqpoint{2.801328in}{0.948753in}}%
\pgfpathlineto{\pgfqpoint{2.801602in}{0.949442in}}%
\pgfpathlineto{\pgfqpoint{2.803199in}{0.951429in}}%
\pgfpathlineto{\pgfqpoint{2.804362in}{0.950541in}}%
\pgfpathlineto{\pgfqpoint{2.806259in}{0.948697in}}%
\pgfpathlineto{\pgfqpoint{2.806551in}{0.949461in}}%
\pgfpathlineto{\pgfqpoint{2.808120in}{0.951450in}}%
\pgfpathlineto{\pgfqpoint{2.809244in}{0.950626in}}%
\pgfpathlineto{\pgfqpoint{2.811222in}{0.948706in}}%
\pgfpathlineto{\pgfqpoint{2.811498in}{0.949380in}}%
\pgfpathlineto{\pgfqpoint{2.813121in}{0.951351in}}%
\pgfpathlineto{\pgfqpoint{2.814323in}{0.950398in}}%
\pgfpathlineto{\pgfqpoint{2.816149in}{0.948636in}}%
\pgfpathlineto{\pgfqpoint{2.816422in}{0.949372in}}%
\pgfpathlineto{\pgfqpoint{2.818038in}{0.951418in}}%
\pgfpathlineto{\pgfqpoint{2.819175in}{0.950559in}}%
\pgfpathlineto{\pgfqpoint{2.821112in}{0.948679in}}%
\pgfpathlineto{\pgfqpoint{2.821374in}{0.949374in}}%
\pgfpathlineto{\pgfqpoint{2.822982in}{0.951391in}}%
\pgfpathlineto{\pgfqpoint{2.824145in}{0.950503in}}%
\pgfpathlineto{\pgfqpoint{2.826063in}{0.948663in}}%
\pgfpathlineto{\pgfqpoint{2.826352in}{0.949403in}}%
\pgfpathlineto{\pgfqpoint{2.827896in}{0.951346in}}%
\pgfpathlineto{\pgfqpoint{2.829028in}{0.950521in}}%
\pgfpathlineto{\pgfqpoint{2.831004in}{0.948625in}}%
\pgfpathlineto{\pgfqpoint{2.831280in}{0.949328in}}%
\pgfpathlineto{\pgfqpoint{2.832906in}{0.951297in}}%
\pgfpathlineto{\pgfqpoint{2.834110in}{0.950339in}}%
\pgfpathlineto{\pgfqpoint{2.835932in}{0.948571in}}%
\pgfpathlineto{\pgfqpoint{2.836197in}{0.949277in}}%
\pgfpathlineto{\pgfqpoint{2.837829in}{0.951343in}}%
\pgfpathlineto{\pgfqpoint{2.838995in}{0.950448in}}%
\pgfpathlineto{\pgfqpoint{2.840891in}{0.948605in}}%
\pgfpathlineto{\pgfqpoint{2.841166in}{0.949291in}}%
\pgfpathlineto{\pgfqpoint{2.842778in}{0.951279in}}%
\pgfpathlineto{\pgfqpoint{2.843949in}{0.950370in}}%
\pgfpathlineto{\pgfqpoint{2.845837in}{0.948563in}}%
\pgfpathlineto{\pgfqpoint{2.846112in}{0.949273in}}%
\pgfpathlineto{\pgfqpoint{2.847746in}{0.951261in}}%
\pgfpathlineto{\pgfqpoint{2.848950in}{0.950301in}}%
\pgfpathlineto{\pgfqpoint{2.850767in}{0.948530in}}%
\pgfpathlineto{\pgfqpoint{2.851064in}{0.949292in}}%
\pgfpathlineto{\pgfqpoint{2.852626in}{0.951267in}}%
\pgfpathlineto{\pgfqpoint{2.853764in}{0.950431in}}%
\pgfpathlineto{\pgfqpoint{2.855716in}{0.948526in}}%
\pgfpathlineto{\pgfqpoint{2.856011in}{0.949275in}}%
\pgfpathlineto{\pgfqpoint{2.857605in}{0.951246in}}%
\pgfpathlineto{\pgfqpoint{2.858751in}{0.950372in}}%
\pgfpathlineto{\pgfqpoint{2.860681in}{0.948526in}}%
\pgfpathlineto{\pgfqpoint{2.860947in}{0.949216in}}%
\pgfpathlineto{\pgfqpoint{2.862587in}{0.951202in}}%
\pgfpathlineto{\pgfqpoint{2.863806in}{0.950222in}}%
\pgfpathlineto{\pgfqpoint{2.865609in}{0.948470in}}%
\pgfpathlineto{\pgfqpoint{2.865854in}{0.949123in}}%
\pgfpathlineto{\pgfqpoint{2.867485in}{0.951221in}}%
\pgfpathlineto{\pgfqpoint{2.868634in}{0.950355in}}%
\pgfpathlineto{\pgfqpoint{2.870559in}{0.948477in}}%
\pgfpathlineto{\pgfqpoint{2.870839in}{0.949208in}}%
\pgfpathlineto{\pgfqpoint{2.872413in}{0.951214in}}%
\pgfpathlineto{\pgfqpoint{2.873545in}{0.950385in}}%
\pgfpathlineto{\pgfqpoint{2.875511in}{0.948476in}}%
\pgfpathlineto{\pgfqpoint{2.875788in}{0.949147in}}%
\pgfpathlineto{\pgfqpoint{2.877422in}{0.951120in}}%
\pgfpathlineto{\pgfqpoint{2.878640in}{0.950143in}}%
\pgfpathlineto{\pgfqpoint{2.880437in}{0.948402in}}%
\pgfpathlineto{\pgfqpoint{2.880718in}{0.949156in}}%
\pgfpathlineto{\pgfqpoint{2.882295in}{0.951187in}}%
\pgfpathlineto{\pgfqpoint{2.883407in}{0.950384in}}%
\pgfpathlineto{\pgfqpoint{2.885400in}{0.948444in}}%
\pgfpathlineto{\pgfqpoint{2.885701in}{0.949199in}}%
\pgfpathlineto{\pgfqpoint{2.887304in}{0.951122in}}%
\pgfpathlineto{\pgfqpoint{2.888512in}{0.950164in}}%
\pgfpathlineto{\pgfqpoint{2.890332in}{0.948394in}}%
\pgfpathlineto{\pgfqpoint{2.890622in}{0.949147in}}%
\pgfpathlineto{\pgfqpoint{2.892232in}{0.951134in}}%
\pgfpathlineto{\pgfqpoint{2.893411in}{0.950219in}}%
\pgfpathlineto{\pgfqpoint{2.895287in}{0.948396in}}%
\pgfpathlineto{\pgfqpoint{2.895556in}{0.949113in}}%
\pgfpathlineto{\pgfqpoint{2.897132in}{0.951146in}}%
\pgfpathlineto{\pgfqpoint{2.898249in}{0.950340in}}%
\pgfpathlineto{\pgfqpoint{2.900244in}{0.948414in}}%
\pgfpathlineto{\pgfqpoint{2.900544in}{0.949160in}}%
\pgfpathlineto{\pgfqpoint{2.902058in}{0.951070in}}%
\pgfpathlineto{\pgfqpoint{2.903164in}{0.950287in}}%
\pgfpathlineto{\pgfqpoint{2.905167in}{0.948343in}}%
\pgfpathlineto{\pgfqpoint{2.905430in}{0.949048in}}%
\pgfpathlineto{\pgfqpoint{2.907074in}{0.951122in}}%
\pgfpathlineto{\pgfqpoint{2.908270in}{0.950190in}}%
\pgfpathlineto{\pgfqpoint{2.910136in}{0.948391in}}%
\pgfpathlineto{\pgfqpoint{2.910435in}{0.949137in}}%
\pgfpathlineto{\pgfqpoint{2.911960in}{0.951046in}}%
\pgfpathlineto{\pgfqpoint{2.913088in}{0.950232in}}%
\pgfpathlineto{\pgfqpoint{2.915056in}{0.948318in}}%
\pgfpathlineto{\pgfqpoint{2.915342in}{0.949077in}}%
\pgfpathlineto{\pgfqpoint{2.916925in}{0.951090in}}%
\pgfpathlineto{\pgfqpoint{2.918069in}{0.950243in}}%
\pgfpathlineto{\pgfqpoint{2.920014in}{0.948343in}}%
\pgfpathlineto{\pgfqpoint{2.920269in}{0.949008in}}%
\pgfpathlineto{\pgfqpoint{2.921852in}{0.951076in}}%
\pgfpathlineto{\pgfqpoint{2.922948in}{0.950298in}}%
\pgfpathlineto{\pgfqpoint{2.924965in}{0.948337in}}%
\pgfpathlineto{\pgfqpoint{2.925243in}{0.949009in}}%
\pgfpathlineto{\pgfqpoint{2.926878in}{0.950984in}}%
\pgfpathlineto{\pgfqpoint{2.928108in}{0.949995in}}%
\pgfpathlineto{\pgfqpoint{2.929900in}{0.948273in}}%
\pgfpathlineto{\pgfqpoint{2.930161in}{0.948983in}}%
\pgfpathlineto{\pgfqpoint{2.931743in}{0.951044in}}%
\pgfpathlineto{\pgfqpoint{2.932868in}{0.950235in}}%
\pgfpathlineto{\pgfqpoint{2.934862in}{0.948306in}}%
\pgfpathlineto{\pgfqpoint{2.935160in}{0.949046in}}%
\pgfpathlineto{\pgfqpoint{2.936667in}{0.950962in}}%
\pgfpathlineto{\pgfqpoint{2.937754in}{0.950210in}}%
\pgfpathlineto{\pgfqpoint{2.939786in}{0.948239in}}%
\pgfpathlineto{\pgfqpoint{2.940073in}{0.948999in}}%
\pgfpathlineto{\pgfqpoint{2.941645in}{0.951000in}}%
\pgfpathlineto{\pgfqpoint{2.942783in}{0.950166in}}%
\pgfpathlineto{\pgfqpoint{2.944752in}{0.948262in}}%
\pgfpathlineto{\pgfqpoint{2.945050in}{0.948999in}}%
\pgfpathlineto{\pgfqpoint{2.946635in}{0.950918in}}%
\pgfpathlineto{\pgfqpoint{2.947821in}{0.949995in}}%
\pgfpathlineto{\pgfqpoint{2.949674in}{0.948191in}}%
\pgfpathlineto{\pgfqpoint{2.949962in}{0.948913in}}%
\pgfpathlineto{\pgfqpoint{2.951577in}{0.950914in}}%
\pgfpathlineto{\pgfqpoint{2.952758in}{0.949999in}}%
\pgfpathlineto{\pgfqpoint{2.954623in}{0.948181in}}%
\pgfpathlineto{\pgfqpoint{2.954922in}{0.948941in}}%
\pgfpathlineto{\pgfqpoint{2.956505in}{0.950911in}}%
\pgfpathlineto{\pgfqpoint{2.957661in}{0.950037in}}%
\pgfpathlineto{\pgfqpoint{2.959588in}{0.948185in}}%
\pgfpathlineto{\pgfqpoint{2.959864in}{0.948883in}}%
\pgfpathlineto{\pgfqpoint{2.961518in}{0.950847in}}%
\pgfpathlineto{\pgfqpoint{2.962779in}{0.949810in}}%
\pgfpathlineto{\pgfqpoint{2.964530in}{0.948142in}}%
\pgfpathlineto{\pgfqpoint{2.964805in}{0.948843in}}%
\pgfpathlineto{\pgfqpoint{2.966437in}{0.950824in}}%
\pgfpathlineto{\pgfqpoint{2.967655in}{0.949852in}}%
\pgfpathlineto{\pgfqpoint{2.969456in}{0.948099in}}%
\pgfpathlineto{\pgfqpoint{2.969748in}{0.948830in}}%
\pgfpathlineto{\pgfqpoint{2.971381in}{0.950819in}}%
\pgfpathlineto{\pgfqpoint{2.972604in}{0.949844in}}%
\pgfpathlineto{\pgfqpoint{2.974405in}{0.948091in}}%
\pgfpathlineto{\pgfqpoint{2.974705in}{0.948857in}}%
\pgfpathlineto{\pgfqpoint{2.976282in}{0.950822in}}%
\pgfpathlineto{\pgfqpoint{2.977438in}{0.949952in}}%
\pgfpathlineto{\pgfqpoint{2.979363in}{0.948093in}}%
\pgfpathlineto{\pgfqpoint{2.979645in}{0.948854in}}%
\pgfpathlineto{\pgfqpoint{2.981264in}{0.950843in}}%
\pgfpathlineto{\pgfqpoint{2.982450in}{0.949915in}}%
\pgfpathlineto{\pgfqpoint{2.984316in}{0.948117in}}%
\pgfpathlineto{\pgfqpoint{2.984603in}{0.948861in}}%
\pgfpathlineto{\pgfqpoint{2.986192in}{0.950803in}}%
\pgfpathlineto{\pgfqpoint{2.987359in}{0.949903in}}%
\pgfpathlineto{\pgfqpoint{2.989259in}{0.948078in}}%
\pgfpathlineto{\pgfqpoint{2.989527in}{0.948808in}}%
\pgfpathlineto{\pgfqpoint{2.991110in}{0.950815in}}%
\pgfpathlineto{\pgfqpoint{2.992238in}{0.949983in}}%
\pgfpathlineto{\pgfqpoint{2.994196in}{0.948075in}}%
\pgfpathlineto{\pgfqpoint{2.994483in}{0.948832in}}%
\pgfpathlineto{\pgfqpoint{2.996065in}{0.950816in}}%
\pgfpathlineto{\pgfqpoint{2.997211in}{0.949959in}}%
\pgfpathlineto{\pgfqpoint{2.999146in}{0.948085in}}%
\pgfpathlineto{\pgfqpoint{2.999432in}{0.948838in}}%
\pgfpathlineto{\pgfqpoint{3.001052in}{0.950802in}}%
\pgfpathlineto{\pgfqpoint{3.002263in}{0.949840in}}%
\pgfpathlineto{\pgfqpoint{3.004099in}{0.948084in}}%
\pgfpathlineto{\pgfqpoint{3.004385in}{0.948824in}}%
\pgfpathlineto{\pgfqpoint{3.005930in}{0.950772in}}%
\pgfpathlineto{\pgfqpoint{3.007052in}{0.949960in}}%
\pgfpathlineto{\pgfqpoint{3.009038in}{0.948046in}}%
\pgfpathlineto{\pgfqpoint{3.009316in}{0.948737in}}%
\pgfpathlineto{\pgfqpoint{3.010925in}{0.950716in}}%
\pgfpathlineto{\pgfqpoint{3.012100in}{0.949806in}}%
\pgfpathlineto{\pgfqpoint{3.013966in}{0.948001in}}%
\pgfpathlineto{\pgfqpoint{3.014251in}{0.948759in}}%
\pgfpathlineto{\pgfqpoint{3.015846in}{0.950773in}}%
\pgfpathlineto{\pgfqpoint{3.016981in}{0.949925in}}%
\pgfpathlineto{\pgfqpoint{3.018933in}{0.948034in}}%
\pgfpathlineto{\pgfqpoint{3.019170in}{0.948674in}}%
\pgfpathlineto{\pgfqpoint{3.020752in}{0.950757in}}%
\pgfpathlineto{\pgfqpoint{3.021854in}{0.949980in}}%
\pgfpathlineto{\pgfqpoint{3.023876in}{0.948017in}}%
\pgfpathlineto{\pgfqpoint{3.024153in}{0.948690in}}%
\pgfpathlineto{\pgfqpoint{3.025791in}{0.950664in}}%
\pgfpathlineto{\pgfqpoint{3.027008in}{0.949687in}}%
\pgfpathlineto{\pgfqpoint{3.028803in}{0.947956in}}%
\pgfpathlineto{\pgfqpoint{3.029047in}{0.948616in}}%
\pgfpathlineto{\pgfqpoint{3.030672in}{0.950746in}}%
\pgfpathlineto{\pgfqpoint{3.031801in}{0.949918in}}%
\pgfpathlineto{\pgfqpoint{3.033771in}{0.948012in}}%
\pgfpathlineto{\pgfqpoint{3.034074in}{0.948759in}}%
\pgfpathlineto{\pgfqpoint{3.035595in}{0.950660in}}%
\pgfpathlineto{\pgfqpoint{3.036711in}{0.949860in}}%
\pgfpathlineto{\pgfqpoint{3.038695in}{0.947943in}}%
\pgfpathlineto{\pgfqpoint{3.038967in}{0.948665in}}%
\pgfpathlineto{\pgfqpoint{3.040581in}{0.950708in}}%
\pgfpathlineto{\pgfqpoint{3.041741in}{0.949830in}}%
\pgfpathlineto{\pgfqpoint{3.043657in}{0.947971in}}%
\pgfpathlineto{\pgfqpoint{3.043933in}{0.948646in}}%
\pgfpathlineto{\pgfqpoint{3.045533in}{0.950628in}}%
\pgfpathlineto{\pgfqpoint{3.046708in}{0.949727in}}%
\pgfpathlineto{\pgfqpoint{3.048585in}{0.947913in}}%
\pgfpathlineto{\pgfqpoint{3.048846in}{0.948620in}}%
\pgfpathlineto{\pgfqpoint{3.050476in}{0.950701in}}%
\pgfpathlineto{\pgfqpoint{3.051636in}{0.949822in}}%
\pgfpathlineto{\pgfqpoint{3.053547in}{0.947965in}}%
\pgfpathlineto{\pgfqpoint{3.053824in}{0.948641in}}%
\pgfpathlineto{\pgfqpoint{3.055439in}{0.950623in}}%
\pgfpathlineto{\pgfqpoint{3.056634in}{0.949689in}}%
\pgfpathlineto{\pgfqpoint{3.058473in}{0.947901in}}%
\pgfpathlineto{\pgfqpoint{3.058764in}{0.948650in}}%
\pgfpathlineto{\pgfqpoint{3.060346in}{0.950650in}}%
\pgfpathlineto{\pgfqpoint{3.061492in}{0.949799in}}%
\pgfpathlineto{\pgfqpoint{3.063438in}{0.947913in}}%
\pgfpathlineto{\pgfqpoint{3.063716in}{0.948587in}}%
\pgfpathlineto{\pgfqpoint{3.065356in}{0.950565in}}%
\pgfpathlineto{\pgfqpoint{3.066576in}{0.949584in}}%
\pgfpathlineto{\pgfqpoint{3.068368in}{0.947857in}}%
\pgfpathlineto{\pgfqpoint{3.068629in}{0.948554in}}%
\pgfpathlineto{\pgfqpoint{3.070254in}{0.950623in}}%
\pgfpathlineto{\pgfqpoint{3.071409in}{0.949748in}}%
\pgfpathlineto{\pgfqpoint{3.073328in}{0.947886in}}%
\pgfpathlineto{\pgfqpoint{3.073605in}{0.948563in}}%
\pgfpathlineto{\pgfqpoint{3.075228in}{0.950548in}}%
\pgfpathlineto{\pgfqpoint{3.076430in}{0.949601in}}%
\pgfpathlineto{\pgfqpoint{3.078258in}{0.947838in}}%
\pgfpathlineto{\pgfqpoint{3.078530in}{0.948562in}}%
\pgfpathlineto{\pgfqpoint{3.080135in}{0.950605in}}%
\pgfpathlineto{\pgfqpoint{3.081281in}{0.949749in}}%
\pgfpathlineto{\pgfqpoint{3.083226in}{0.947876in}}%
\pgfpathlineto{\pgfqpoint{3.083516in}{0.948617in}}%
\pgfpathlineto{\pgfqpoint{3.085125in}{0.950553in}}%
\pgfpathlineto{\pgfqpoint{3.086333in}{0.949595in}}%
\pgfpathlineto{\pgfqpoint{3.088161in}{0.947840in}}%
\pgfpathlineto{\pgfqpoint{3.088394in}{0.948478in}}%
\pgfpathlineto{\pgfqpoint{3.090026in}{0.950600in}}%
\pgfpathlineto{\pgfqpoint{3.091175in}{0.949744in}}%
\pgfpathlineto{\pgfqpoint{3.093105in}{0.947863in}}%
\pgfpathlineto{\pgfqpoint{3.093369in}{0.948556in}}%
\pgfpathlineto{\pgfqpoint{3.094952in}{0.950602in}}%
\pgfpathlineto{\pgfqpoint{3.096068in}{0.949797in}}%
\pgfpathlineto{\pgfqpoint{3.098052in}{0.947868in}}%
\pgfpathlineto{\pgfqpoint{3.098335in}{0.948611in}}%
\pgfpathlineto{\pgfqpoint{3.099910in}{0.950583in}}%
\pgfpathlineto{\pgfqpoint{3.101060in}{0.949725in}}%
\pgfpathlineto{\pgfqpoint{3.103009in}{0.947859in}}%
\pgfpathlineto{\pgfqpoint{3.103298in}{0.948600in}}%
\pgfpathlineto{\pgfqpoint{3.104841in}{0.950547in}}%
\pgfpathlineto{\pgfqpoint{3.105962in}{0.949739in}}%
\pgfpathlineto{\pgfqpoint{3.107948in}{0.947831in}}%
\pgfpathlineto{\pgfqpoint{3.108225in}{0.948525in}}%
\pgfpathlineto{\pgfqpoint{3.109852in}{0.950501in}}%
\pgfpathlineto{\pgfqpoint{3.111059in}{0.949543in}}%
\pgfpathlineto{\pgfqpoint{3.112877in}{0.947784in}}%
\pgfpathlineto{\pgfqpoint{3.113149in}{0.948505in}}%
\pgfpathlineto{\pgfqpoint{3.114730in}{0.950549in}}%
\pgfpathlineto{\pgfqpoint{3.115845in}{0.949745in}}%
\pgfpathlineto{\pgfqpoint{3.117839in}{0.947809in}}%
\pgfpathlineto{\pgfqpoint{3.118116in}{0.948481in}}%
\pgfpathlineto{\pgfqpoint{3.119739in}{0.950460in}}%
\pgfpathlineto{\pgfqpoint{3.120950in}{0.949503in}}%
\pgfpathlineto{\pgfqpoint{3.122768in}{0.947749in}}%
\pgfpathlineto{\pgfqpoint{3.123050in}{0.948500in}}%
\pgfpathlineto{\pgfqpoint{3.124633in}{0.950514in}}%
\pgfpathlineto{\pgfqpoint{3.125759in}{0.949687in}}%
\pgfpathlineto{\pgfqpoint{3.127731in}{0.947775in}}%
\pgfpathlineto{\pgfqpoint{3.128008in}{0.948448in}}%
\pgfpathlineto{\pgfqpoint{3.129632in}{0.950426in}}%
\pgfpathlineto{\pgfqpoint{3.130842in}{0.949468in}}%
\pgfpathlineto{\pgfqpoint{3.132658in}{0.947718in}}%
\pgfpathlineto{\pgfqpoint{3.132900in}{0.948373in}}%
\pgfpathlineto{\pgfqpoint{3.134525in}{0.950506in}}%
\pgfpathlineto{\pgfqpoint{3.135656in}{0.949678in}}%
\pgfpathlineto{\pgfqpoint{3.137625in}{0.947772in}}%
\pgfpathlineto{\pgfqpoint{3.137928in}{0.948513in}}%
\pgfpathlineto{\pgfqpoint{3.139423in}{0.950419in}}%
\pgfpathlineto{\pgfqpoint{3.140517in}{0.949669in}}%
\pgfpathlineto{\pgfqpoint{3.142550in}{0.947702in}}%
\pgfpathlineto{\pgfqpoint{3.142815in}{0.948406in}}%
\pgfpathlineto{\pgfqpoint{3.144441in}{0.950466in}}%
\pgfpathlineto{\pgfqpoint{3.145594in}{0.949589in}}%
\pgfpathlineto{\pgfqpoint{3.147512in}{0.947730in}}%
\pgfpathlineto{\pgfqpoint{3.147789in}{0.948403in}}%
\pgfpathlineto{\pgfqpoint{3.149416in}{0.950382in}}%
\pgfpathlineto{\pgfqpoint{3.150624in}{0.949424in}}%
\pgfpathlineto{\pgfqpoint{3.152440in}{0.947675in}}%
\pgfpathlineto{\pgfqpoint{3.152693in}{0.948359in}}%
\pgfpathlineto{\pgfqpoint{3.154313in}{0.950463in}}%
\pgfpathlineto{\pgfqpoint{3.155465in}{0.949608in}}%
\pgfpathlineto{\pgfqpoint{3.157395in}{0.947724in}}%
\pgfpathlineto{\pgfqpoint{3.157647in}{0.948377in}}%
\pgfpathlineto{\pgfqpoint{3.159279in}{0.950455in}}%
\pgfpathlineto{\pgfqpoint{3.160455in}{0.949556in}}%
\pgfpathlineto{\pgfqpoint{3.162338in}{0.947716in}}%
\pgfpathlineto{\pgfqpoint{3.162637in}{0.948471in}}%
\pgfpathlineto{\pgfqpoint{3.164178in}{0.950432in}}%
\pgfpathlineto{\pgfqpoint{3.165304in}{0.949626in}}%
\pgfpathlineto{\pgfqpoint{3.167296in}{0.947710in}}%
\pgfpathlineto{\pgfqpoint{3.167541in}{0.948380in}}%
\pgfpathlineto{\pgfqpoint{3.169173in}{0.950439in}}%
\pgfpathlineto{\pgfqpoint{3.170132in}{0.949749in}}%
\pgfpathlineto{\pgfqpoint{3.170132in}{0.949749in}}%
\pgfusepath{stroke}%
\end{pgfscope}%
\begin{pgfscope}%
\pgfsetrectcap%
\pgfsetmiterjoin%
\pgfsetlinewidth{0.803000pt}%
\definecolor{currentstroke}{rgb}{0.000000,0.000000,0.000000}%
\pgfsetstrokecolor{currentstroke}%
\pgfsetdash{}{0pt}%
\pgfpathmoveto{\pgfqpoint{0.573768in}{0.462654in}}%
\pgfpathlineto{\pgfqpoint{0.573768in}{1.368904in}}%
\pgfusepath{stroke}%
\end{pgfscope}%
\begin{pgfscope}%
\pgfsetrectcap%
\pgfsetmiterjoin%
\pgfsetlinewidth{0.803000pt}%
\definecolor{currentstroke}{rgb}{0.000000,0.000000,0.000000}%
\pgfsetstrokecolor{currentstroke}%
\pgfsetdash{}{0pt}%
\pgfpathmoveto{\pgfqpoint{3.293768in}{0.462654in}}%
\pgfpathlineto{\pgfqpoint{3.293768in}{1.368904in}}%
\pgfusepath{stroke}%
\end{pgfscope}%
\begin{pgfscope}%
\pgfsetrectcap%
\pgfsetmiterjoin%
\pgfsetlinewidth{0.803000pt}%
\definecolor{currentstroke}{rgb}{0.000000,0.000000,0.000000}%
\pgfsetstrokecolor{currentstroke}%
\pgfsetdash{}{0pt}%
\pgfpathmoveto{\pgfqpoint{0.573768in}{0.462654in}}%
\pgfpathlineto{\pgfqpoint{3.293768in}{0.462654in}}%
\pgfusepath{stroke}%
\end{pgfscope}%
\begin{pgfscope}%
\pgfsetrectcap%
\pgfsetmiterjoin%
\pgfsetlinewidth{0.803000pt}%
\definecolor{currentstroke}{rgb}{0.000000,0.000000,0.000000}%
\pgfsetstrokecolor{currentstroke}%
\pgfsetdash{}{0pt}%
\pgfpathmoveto{\pgfqpoint{0.573768in}{1.368904in}}%
\pgfpathlineto{\pgfqpoint{3.293768in}{1.368904in}}%
\pgfusepath{stroke}%
\end{pgfscope}%
\begin{pgfscope}%
\definecolor{textcolor}{rgb}{0.000000,0.000000,0.000000}%
\pgfsetstrokecolor{textcolor}%
\pgfsetfillcolor{textcolor}%
\pgftext[x=0.628168in,y=1.305467in,left,top]{\color{textcolor}{\rmfamily\fontsize{8.000000}{9.600000}\selectfont\catcode`\^=\active\def^{\ifmmode\sp\else\^{}\fi}\catcode`\%=\active\def%{\%}(c)}}%
\end{pgfscope}%
\begin{pgfscope}%
\pgfsetbuttcap%
\pgfsetmiterjoin%
\definecolor{currentfill}{rgb}{1.000000,1.000000,1.000000}%
\pgfsetfillcolor{currentfill}%
\pgfsetlinewidth{0.000000pt}%
\definecolor{currentstroke}{rgb}{0.000000,0.000000,0.000000}%
\pgfsetstrokecolor{currentstroke}%
\pgfsetstrokeopacity{0.000000}%
\pgfsetdash{}{0pt}%
\pgfpathmoveto{\pgfqpoint{2.197902in}{3.064738in}}%
\pgfpathlineto{\pgfqpoint{3.177102in}{3.064738in}}%
\pgfpathlineto{\pgfqpoint{3.177102in}{3.390988in}}%
\pgfpathlineto{\pgfqpoint{2.197902in}{3.390988in}}%
\pgfpathlineto{\pgfqpoint{2.197902in}{3.064738in}}%
\pgfpathclose%
\pgfusepath{fill}%
\end{pgfscope}%
\begin{pgfscope}%
\pgfpathrectangle{\pgfqpoint{2.197902in}{3.064738in}}{\pgfqpoint{0.979200in}{0.326250in}}%
\pgfusepath{clip}%
\pgfsetbuttcap%
\pgfsetroundjoin%
\pgfsetlinewidth{0.401500pt}%
\definecolor{currentstroke}{rgb}{0.690196,0.690196,0.690196}%
\pgfsetstrokecolor{currentstroke}%
\pgfsetstrokeopacity{0.250000}%
\pgfsetdash{{1.480000pt}{0.640000pt}}{0.000000pt}%
\pgfpathmoveto{\pgfqpoint{2.197902in}{3.064738in}}%
\pgfpathlineto{\pgfqpoint{2.197902in}{3.390988in}}%
\pgfusepath{stroke}%
\end{pgfscope}%
\begin{pgfscope}%
\pgfsetbuttcap%
\pgfsetroundjoin%
\definecolor{currentfill}{rgb}{0.000000,0.000000,0.000000}%
\pgfsetfillcolor{currentfill}%
\pgfsetlinewidth{0.803000pt}%
\definecolor{currentstroke}{rgb}{0.000000,0.000000,0.000000}%
\pgfsetstrokecolor{currentstroke}%
\pgfsetdash{}{0pt}%
\pgfsys@defobject{currentmarker}{\pgfqpoint{0.000000in}{0.000000in}}{\pgfqpoint{0.000000in}{0.027778in}}{%
\pgfpathmoveto{\pgfqpoint{0.000000in}{0.000000in}}%
\pgfpathlineto{\pgfqpoint{0.000000in}{0.027778in}}%
\pgfusepath{stroke,fill}%
}%
\begin{pgfscope}%
\pgfsys@transformshift{2.197902in}{3.064738in}%
\pgfsys@useobject{currentmarker}{}%
\end{pgfscope}%
\end{pgfscope}%
\begin{pgfscope}%
\definecolor{textcolor}{rgb}{0.000000,0.000000,0.000000}%
\pgfsetstrokecolor{textcolor}%
\pgfsetfillcolor{textcolor}%
\pgftext[x=2.197902in,y=3.050849in,,top]{\color{textcolor}{\rmfamily\fontsize{5.000000}{6.000000}\selectfont\catcode`\^=\active\def^{\ifmmode\sp\else\^{}\fi}\catcode`\%=\active\def%{\%}4.980}}%
\end{pgfscope}%
\begin{pgfscope}%
\pgfpathrectangle{\pgfqpoint{2.197902in}{3.064738in}}{\pgfqpoint{0.979200in}{0.326250in}}%
\pgfusepath{clip}%
\pgfsetbuttcap%
\pgfsetroundjoin%
\pgfsetlinewidth{0.401500pt}%
\definecolor{currentstroke}{rgb}{0.690196,0.690196,0.690196}%
\pgfsetstrokecolor{currentstroke}%
\pgfsetstrokeopacity{0.250000}%
\pgfsetdash{{1.480000pt}{0.640000pt}}{0.000000pt}%
\pgfpathmoveto{\pgfqpoint{2.442702in}{3.064738in}}%
\pgfpathlineto{\pgfqpoint{2.442702in}{3.390988in}}%
\pgfusepath{stroke}%
\end{pgfscope}%
\begin{pgfscope}%
\pgfsetbuttcap%
\pgfsetroundjoin%
\definecolor{currentfill}{rgb}{0.000000,0.000000,0.000000}%
\pgfsetfillcolor{currentfill}%
\pgfsetlinewidth{0.803000pt}%
\definecolor{currentstroke}{rgb}{0.000000,0.000000,0.000000}%
\pgfsetstrokecolor{currentstroke}%
\pgfsetdash{}{0pt}%
\pgfsys@defobject{currentmarker}{\pgfqpoint{0.000000in}{0.000000in}}{\pgfqpoint{0.000000in}{0.027778in}}{%
\pgfpathmoveto{\pgfqpoint{0.000000in}{0.000000in}}%
\pgfpathlineto{\pgfqpoint{0.000000in}{0.027778in}}%
\pgfusepath{stroke,fill}%
}%
\begin{pgfscope}%
\pgfsys@transformshift{2.442702in}{3.064738in}%
\pgfsys@useobject{currentmarker}{}%
\end{pgfscope}%
\end{pgfscope}%
\begin{pgfscope}%
\definecolor{textcolor}{rgb}{0.000000,0.000000,0.000000}%
\pgfsetstrokecolor{textcolor}%
\pgfsetfillcolor{textcolor}%
\pgftext[x=2.442702in,y=3.050849in,,top]{\color{textcolor}{\rmfamily\fontsize{5.000000}{6.000000}\selectfont\catcode`\^=\active\def^{\ifmmode\sp\else\^{}\fi}\catcode`\%=\active\def%{\%}4.985}}%
\end{pgfscope}%
\begin{pgfscope}%
\pgfpathrectangle{\pgfqpoint{2.197902in}{3.064738in}}{\pgfqpoint{0.979200in}{0.326250in}}%
\pgfusepath{clip}%
\pgfsetbuttcap%
\pgfsetroundjoin%
\pgfsetlinewidth{0.401500pt}%
\definecolor{currentstroke}{rgb}{0.690196,0.690196,0.690196}%
\pgfsetstrokecolor{currentstroke}%
\pgfsetstrokeopacity{0.250000}%
\pgfsetdash{{1.480000pt}{0.640000pt}}{0.000000pt}%
\pgfpathmoveto{\pgfqpoint{2.687502in}{3.064738in}}%
\pgfpathlineto{\pgfqpoint{2.687502in}{3.390988in}}%
\pgfusepath{stroke}%
\end{pgfscope}%
\begin{pgfscope}%
\pgfsetbuttcap%
\pgfsetroundjoin%
\definecolor{currentfill}{rgb}{0.000000,0.000000,0.000000}%
\pgfsetfillcolor{currentfill}%
\pgfsetlinewidth{0.803000pt}%
\definecolor{currentstroke}{rgb}{0.000000,0.000000,0.000000}%
\pgfsetstrokecolor{currentstroke}%
\pgfsetdash{}{0pt}%
\pgfsys@defobject{currentmarker}{\pgfqpoint{0.000000in}{0.000000in}}{\pgfqpoint{0.000000in}{0.027778in}}{%
\pgfpathmoveto{\pgfqpoint{0.000000in}{0.000000in}}%
\pgfpathlineto{\pgfqpoint{0.000000in}{0.027778in}}%
\pgfusepath{stroke,fill}%
}%
\begin{pgfscope}%
\pgfsys@transformshift{2.687502in}{3.064738in}%
\pgfsys@useobject{currentmarker}{}%
\end{pgfscope}%
\end{pgfscope}%
\begin{pgfscope}%
\definecolor{textcolor}{rgb}{0.000000,0.000000,0.000000}%
\pgfsetstrokecolor{textcolor}%
\pgfsetfillcolor{textcolor}%
\pgftext[x=2.687502in,y=3.050849in,,top]{\color{textcolor}{\rmfamily\fontsize{5.000000}{6.000000}\selectfont\catcode`\^=\active\def^{\ifmmode\sp\else\^{}\fi}\catcode`\%=\active\def%{\%}4.990}}%
\end{pgfscope}%
\begin{pgfscope}%
\pgfpathrectangle{\pgfqpoint{2.197902in}{3.064738in}}{\pgfqpoint{0.979200in}{0.326250in}}%
\pgfusepath{clip}%
\pgfsetbuttcap%
\pgfsetroundjoin%
\pgfsetlinewidth{0.401500pt}%
\definecolor{currentstroke}{rgb}{0.690196,0.690196,0.690196}%
\pgfsetstrokecolor{currentstroke}%
\pgfsetstrokeopacity{0.250000}%
\pgfsetdash{{1.480000pt}{0.640000pt}}{0.000000pt}%
\pgfpathmoveto{\pgfqpoint{2.932302in}{3.064738in}}%
\pgfpathlineto{\pgfqpoint{2.932302in}{3.390988in}}%
\pgfusepath{stroke}%
\end{pgfscope}%
\begin{pgfscope}%
\pgfsetbuttcap%
\pgfsetroundjoin%
\definecolor{currentfill}{rgb}{0.000000,0.000000,0.000000}%
\pgfsetfillcolor{currentfill}%
\pgfsetlinewidth{0.803000pt}%
\definecolor{currentstroke}{rgb}{0.000000,0.000000,0.000000}%
\pgfsetstrokecolor{currentstroke}%
\pgfsetdash{}{0pt}%
\pgfsys@defobject{currentmarker}{\pgfqpoint{0.000000in}{0.000000in}}{\pgfqpoint{0.000000in}{0.027778in}}{%
\pgfpathmoveto{\pgfqpoint{0.000000in}{0.000000in}}%
\pgfpathlineto{\pgfqpoint{0.000000in}{0.027778in}}%
\pgfusepath{stroke,fill}%
}%
\begin{pgfscope}%
\pgfsys@transformshift{2.932302in}{3.064738in}%
\pgfsys@useobject{currentmarker}{}%
\end{pgfscope}%
\end{pgfscope}%
\begin{pgfscope}%
\definecolor{textcolor}{rgb}{0.000000,0.000000,0.000000}%
\pgfsetstrokecolor{textcolor}%
\pgfsetfillcolor{textcolor}%
\pgftext[x=2.932302in,y=3.050849in,,top]{\color{textcolor}{\rmfamily\fontsize{5.000000}{6.000000}\selectfont\catcode`\^=\active\def^{\ifmmode\sp\else\^{}\fi}\catcode`\%=\active\def%{\%}4.995}}%
\end{pgfscope}%
\begin{pgfscope}%
\pgfpathrectangle{\pgfqpoint{2.197902in}{3.064738in}}{\pgfqpoint{0.979200in}{0.326250in}}%
\pgfusepath{clip}%
\pgfsetbuttcap%
\pgfsetroundjoin%
\pgfsetlinewidth{0.401500pt}%
\definecolor{currentstroke}{rgb}{0.690196,0.690196,0.690196}%
\pgfsetstrokecolor{currentstroke}%
\pgfsetstrokeopacity{0.250000}%
\pgfsetdash{{1.480000pt}{0.640000pt}}{0.000000pt}%
\pgfpathmoveto{\pgfqpoint{3.177102in}{3.064738in}}%
\pgfpathlineto{\pgfqpoint{3.177102in}{3.390988in}}%
\pgfusepath{stroke}%
\end{pgfscope}%
\begin{pgfscope}%
\pgfsetbuttcap%
\pgfsetroundjoin%
\definecolor{currentfill}{rgb}{0.000000,0.000000,0.000000}%
\pgfsetfillcolor{currentfill}%
\pgfsetlinewidth{0.803000pt}%
\definecolor{currentstroke}{rgb}{0.000000,0.000000,0.000000}%
\pgfsetstrokecolor{currentstroke}%
\pgfsetdash{}{0pt}%
\pgfsys@defobject{currentmarker}{\pgfqpoint{0.000000in}{0.000000in}}{\pgfqpoint{0.000000in}{0.027778in}}{%
\pgfpathmoveto{\pgfqpoint{0.000000in}{0.000000in}}%
\pgfpathlineto{\pgfqpoint{0.000000in}{0.027778in}}%
\pgfusepath{stroke,fill}%
}%
\begin{pgfscope}%
\pgfsys@transformshift{3.177102in}{3.064738in}%
\pgfsys@useobject{currentmarker}{}%
\end{pgfscope}%
\end{pgfscope}%
\begin{pgfscope}%
\definecolor{textcolor}{rgb}{0.000000,0.000000,0.000000}%
\pgfsetstrokecolor{textcolor}%
\pgfsetfillcolor{textcolor}%
\pgftext[x=3.177102in,y=3.050849in,,top]{\color{textcolor}{\rmfamily\fontsize{5.000000}{6.000000}\selectfont\catcode`\^=\active\def^{\ifmmode\sp\else\^{}\fi}\catcode`\%=\active\def%{\%}5.000}}%
\end{pgfscope}%
\begin{pgfscope}%
\pgfpathrectangle{\pgfqpoint{2.197902in}{3.064738in}}{\pgfqpoint{0.979200in}{0.326250in}}%
\pgfusepath{clip}%
\pgfsetbuttcap%
\pgfsetroundjoin%
\pgfsetlinewidth{0.401500pt}%
\definecolor{currentstroke}{rgb}{0.690196,0.690196,0.690196}%
\pgfsetstrokecolor{currentstroke}%
\pgfsetstrokeopacity{0.250000}%
\pgfsetdash{{1.480000pt}{0.640000pt}}{0.000000pt}%
\pgfpathmoveto{\pgfqpoint{2.197902in}{3.064738in}}%
\pgfpathlineto{\pgfqpoint{3.177102in}{3.064738in}}%
\pgfusepath{stroke}%
\end{pgfscope}%
\begin{pgfscope}%
\pgfsetbuttcap%
\pgfsetroundjoin%
\definecolor{currentfill}{rgb}{0.000000,0.000000,0.000000}%
\pgfsetfillcolor{currentfill}%
\pgfsetlinewidth{0.803000pt}%
\definecolor{currentstroke}{rgb}{0.000000,0.000000,0.000000}%
\pgfsetstrokecolor{currentstroke}%
\pgfsetdash{}{0pt}%
\pgfsys@defobject{currentmarker}{\pgfqpoint{0.000000in}{0.000000in}}{\pgfqpoint{0.027778in}{0.000000in}}{%
\pgfpathmoveto{\pgfqpoint{0.000000in}{0.000000in}}%
\pgfpathlineto{\pgfqpoint{0.027778in}{0.000000in}}%
\pgfusepath{stroke,fill}%
}%
\begin{pgfscope}%
\pgfsys@transformshift{2.197902in}{3.064738in}%
\pgfsys@useobject{currentmarker}{}%
\end{pgfscope}%
\end{pgfscope}%
\begin{pgfscope}%
\definecolor{textcolor}{rgb}{0.000000,0.000000,0.000000}%
\pgfsetstrokecolor{textcolor}%
\pgfsetfillcolor{textcolor}%
\pgftext[x=1.981463in, y=3.040625in, left, base]{\color{textcolor}{\rmfamily\fontsize{5.000000}{6.000000}\selectfont\catcode`\^=\active\def^{\ifmmode\sp\else\^{}\fi}\catcode`\%=\active\def%{\%}-0.50}}%
\end{pgfscope}%
\begin{pgfscope}%
\pgfpathrectangle{\pgfqpoint{2.197902in}{3.064738in}}{\pgfqpoint{0.979200in}{0.326250in}}%
\pgfusepath{clip}%
\pgfsetbuttcap%
\pgfsetroundjoin%
\pgfsetlinewidth{0.401500pt}%
\definecolor{currentstroke}{rgb}{0.690196,0.690196,0.690196}%
\pgfsetstrokecolor{currentstroke}%
\pgfsetstrokeopacity{0.250000}%
\pgfsetdash{{1.480000pt}{0.640000pt}}{0.000000pt}%
\pgfpathmoveto{\pgfqpoint{2.197902in}{3.227863in}}%
\pgfpathlineto{\pgfqpoint{3.177102in}{3.227863in}}%
\pgfusepath{stroke}%
\end{pgfscope}%
\begin{pgfscope}%
\pgfsetbuttcap%
\pgfsetroundjoin%
\definecolor{currentfill}{rgb}{0.000000,0.000000,0.000000}%
\pgfsetfillcolor{currentfill}%
\pgfsetlinewidth{0.803000pt}%
\definecolor{currentstroke}{rgb}{0.000000,0.000000,0.000000}%
\pgfsetstrokecolor{currentstroke}%
\pgfsetdash{}{0pt}%
\pgfsys@defobject{currentmarker}{\pgfqpoint{0.000000in}{0.000000in}}{\pgfqpoint{0.027778in}{0.000000in}}{%
\pgfpathmoveto{\pgfqpoint{0.000000in}{0.000000in}}%
\pgfpathlineto{\pgfqpoint{0.027778in}{0.000000in}}%
\pgfusepath{stroke,fill}%
}%
\begin{pgfscope}%
\pgfsys@transformshift{2.197902in}{3.227863in}%
\pgfsys@useobject{currentmarker}{}%
\end{pgfscope}%
\end{pgfscope}%
\begin{pgfscope}%
\definecolor{textcolor}{rgb}{0.000000,0.000000,0.000000}%
\pgfsetstrokecolor{textcolor}%
\pgfsetfillcolor{textcolor}%
\pgftext[x=2.014257in, y=3.203750in, left, base]{\color{textcolor}{\rmfamily\fontsize{5.000000}{6.000000}\selectfont\catcode`\^=\active\def^{\ifmmode\sp\else\^{}\fi}\catcode`\%=\active\def%{\%}0.73}}%
\end{pgfscope}%
\begin{pgfscope}%
\pgfpathrectangle{\pgfqpoint{2.197902in}{3.064738in}}{\pgfqpoint{0.979200in}{0.326250in}}%
\pgfusepath{clip}%
\pgfsetbuttcap%
\pgfsetroundjoin%
\pgfsetlinewidth{0.401500pt}%
\definecolor{currentstroke}{rgb}{0.690196,0.690196,0.690196}%
\pgfsetstrokecolor{currentstroke}%
\pgfsetstrokeopacity{0.250000}%
\pgfsetdash{{1.480000pt}{0.640000pt}}{0.000000pt}%
\pgfpathmoveto{\pgfqpoint{2.197902in}{3.390988in}}%
\pgfpathlineto{\pgfqpoint{3.177102in}{3.390988in}}%
\pgfusepath{stroke}%
\end{pgfscope}%
\begin{pgfscope}%
\pgfsetbuttcap%
\pgfsetroundjoin%
\definecolor{currentfill}{rgb}{0.000000,0.000000,0.000000}%
\pgfsetfillcolor{currentfill}%
\pgfsetlinewidth{0.803000pt}%
\definecolor{currentstroke}{rgb}{0.000000,0.000000,0.000000}%
\pgfsetstrokecolor{currentstroke}%
\pgfsetdash{}{0pt}%
\pgfsys@defobject{currentmarker}{\pgfqpoint{0.000000in}{0.000000in}}{\pgfqpoint{0.027778in}{0.000000in}}{%
\pgfpathmoveto{\pgfqpoint{0.000000in}{0.000000in}}%
\pgfpathlineto{\pgfqpoint{0.027778in}{0.000000in}}%
\pgfusepath{stroke,fill}%
}%
\begin{pgfscope}%
\pgfsys@transformshift{2.197902in}{3.390988in}%
\pgfsys@useobject{currentmarker}{}%
\end{pgfscope}%
\end{pgfscope}%
\begin{pgfscope}%
\definecolor{textcolor}{rgb}{0.000000,0.000000,0.000000}%
\pgfsetstrokecolor{textcolor}%
\pgfsetfillcolor{textcolor}%
\pgftext[x=2.014257in, y=3.366875in, left, base]{\color{textcolor}{\rmfamily\fontsize{5.000000}{6.000000}\selectfont\catcode`\^=\active\def^{\ifmmode\sp\else\^{}\fi}\catcode`\%=\active\def%{\%}1.95}}%
\end{pgfscope}%
\begin{pgfscope}%
\pgfpathrectangle{\pgfqpoint{2.197902in}{3.064738in}}{\pgfqpoint{0.979200in}{0.326250in}}%
\pgfusepath{clip}%
\pgfsetrectcap%
\pgfsetroundjoin%
\pgfsetlinewidth{1.505625pt}%
\definecolor{currentstroke}{rgb}{0.000000,0.447059,0.698039}%
\pgfsetstrokecolor{currentstroke}%
\pgfsetdash{}{0pt}%
\pgfpathmoveto{\pgfqpoint{2.194319in}{3.131319in}}%
\pgfpathlineto{\pgfqpoint{2.198259in}{3.131319in}}%
\pgfpathlineto{\pgfqpoint{2.198667in}{3.154131in}}%
\pgfpathlineto{\pgfqpoint{2.198918in}{3.315556in}}%
\pgfpathlineto{\pgfqpoint{2.199211in}{3.294250in}}%
\pgfpathlineto{\pgfqpoint{2.200530in}{3.151423in}}%
\pgfpathlineto{\pgfqpoint{2.200706in}{3.216368in}}%
\pgfpathlineto{\pgfqpoint{2.200879in}{3.147548in}}%
\pgfpathlineto{\pgfqpoint{2.201225in}{3.206676in}}%
\pgfpathlineto{\pgfqpoint{2.201916in}{3.133740in}}%
\pgfpathlineto{\pgfqpoint{2.203297in}{3.184566in}}%
\pgfpathlineto{\pgfqpoint{2.206061in}{3.107849in}}%
\pgfpathlineto{\pgfqpoint{2.208795in}{3.167650in}}%
\pgfpathlineto{\pgfqpoint{2.211196in}{3.103975in}}%
\pgfpathlineto{\pgfqpoint{2.213835in}{3.168725in}}%
\pgfpathlineto{\pgfqpoint{2.216399in}{3.107231in}}%
\pgfpathlineto{\pgfqpoint{2.219092in}{3.172949in}}%
\pgfpathlineto{\pgfqpoint{2.221652in}{3.112008in}}%
\pgfpathlineto{\pgfqpoint{2.224288in}{3.177751in}}%
\pgfpathlineto{\pgfqpoint{2.226861in}{3.117065in}}%
\pgfpathlineto{\pgfqpoint{2.229511in}{3.182806in}}%
\pgfpathlineto{\pgfqpoint{2.232129in}{3.122252in}}%
\pgfpathlineto{\pgfqpoint{2.234832in}{3.187960in}}%
\pgfpathlineto{\pgfqpoint{2.237590in}{3.127639in}}%
\pgfpathlineto{\pgfqpoint{2.240589in}{3.193549in}}%
\pgfpathlineto{\pgfqpoint{2.243949in}{3.133905in}}%
\pgfpathlineto{\pgfqpoint{2.248108in}{3.200875in}}%
\pgfpathlineto{\pgfqpoint{2.253004in}{3.142800in}}%
\pgfpathlineto{\pgfqpoint{2.257900in}{3.210435in}}%
\pgfpathlineto{\pgfqpoint{2.262796in}{3.152414in}}%
\pgfpathlineto{\pgfqpoint{2.267692in}{3.219996in}}%
\pgfpathlineto{\pgfqpoint{2.272588in}{3.162028in}}%
\pgfpathlineto{\pgfqpoint{2.277484in}{3.229556in}}%
\pgfpathlineto{\pgfqpoint{2.282380in}{3.171641in}}%
\pgfpathlineto{\pgfqpoint{2.287276in}{3.239117in}}%
\pgfpathlineto{\pgfqpoint{2.292172in}{3.181255in}}%
\pgfpathlineto{\pgfqpoint{2.297068in}{3.248677in}}%
\pgfpathlineto{\pgfqpoint{2.301964in}{3.190869in}}%
\pgfpathlineto{\pgfqpoint{2.306860in}{3.258237in}}%
\pgfpathlineto{\pgfqpoint{2.311756in}{3.200482in}}%
\pgfpathlineto{\pgfqpoint{2.316652in}{3.267797in}}%
\pgfpathlineto{\pgfqpoint{2.321548in}{3.210095in}}%
\pgfpathlineto{\pgfqpoint{2.326444in}{3.277357in}}%
\pgfpathlineto{\pgfqpoint{2.331340in}{3.219708in}}%
\pgfpathlineto{\pgfqpoint{2.336236in}{3.286917in}}%
\pgfpathlineto{\pgfqpoint{2.341132in}{3.229321in}}%
\pgfpathlineto{\pgfqpoint{2.346028in}{3.296477in}}%
\pgfpathlineto{\pgfqpoint{2.350924in}{3.238934in}}%
\pgfpathlineto{\pgfqpoint{2.355820in}{3.306037in}}%
\pgfpathlineto{\pgfqpoint{2.360716in}{3.248547in}}%
\pgfpathlineto{\pgfqpoint{2.365612in}{3.315597in}}%
\pgfpathlineto{\pgfqpoint{2.370508in}{3.258160in}}%
\pgfpathlineto{\pgfqpoint{2.375404in}{3.325156in}}%
\pgfpathlineto{\pgfqpoint{2.380300in}{3.267772in}}%
\pgfpathlineto{\pgfqpoint{2.385196in}{3.334716in}}%
\pgfpathlineto{\pgfqpoint{2.389493in}{3.276801in}}%
\pgfpathlineto{\pgfqpoint{2.393791in}{3.343100in}}%
\pgfpathlineto{\pgfqpoint{2.393846in}{3.255349in}}%
\pgfpathlineto{\pgfqpoint{2.394319in}{3.264558in}}%
\pgfpathlineto{\pgfqpoint{2.395761in}{3.284944in}}%
\pgfpathlineto{\pgfqpoint{2.396690in}{3.281362in}}%
\pgfpathlineto{\pgfqpoint{2.397557in}{3.224930in}}%
\pgfpathlineto{\pgfqpoint{2.398527in}{3.266316in}}%
\pgfpathlineto{\pgfqpoint{2.399661in}{3.162508in}}%
\pgfpathlineto{\pgfqpoint{2.401005in}{3.203989in}}%
\pgfpathlineto{\pgfqpoint{2.402209in}{3.131319in}}%
\pgfpathlineto{\pgfqpoint{2.688192in}{3.131919in}}%
\pgfpathlineto{\pgfqpoint{2.688433in}{3.328131in}}%
\pgfpathlineto{\pgfqpoint{2.689187in}{3.253039in}}%
\pgfpathlineto{\pgfqpoint{2.689690in}{3.262330in}}%
\pgfpathlineto{\pgfqpoint{2.690604in}{3.179144in}}%
\pgfpathlineto{\pgfqpoint{2.691972in}{3.162057in}}%
\pgfpathlineto{\pgfqpoint{2.693796in}{3.145169in}}%
\pgfpathlineto{\pgfqpoint{2.699168in}{3.133911in}}%
\pgfpathlineto{\pgfqpoint{2.701619in}{3.137658in}}%
\pgfpathlineto{\pgfqpoint{2.704246in}{3.135847in}}%
\pgfpathlineto{\pgfqpoint{2.706851in}{3.141294in}}%
\pgfpathlineto{\pgfqpoint{2.709508in}{3.140306in}}%
\pgfpathlineto{\pgfqpoint{2.712085in}{3.146088in}}%
\pgfpathlineto{\pgfqpoint{2.714703in}{3.145240in}}%
\pgfpathlineto{\pgfqpoint{2.717297in}{3.151121in}}%
\pgfpathlineto{\pgfqpoint{2.719938in}{3.150328in}}%
\pgfpathlineto{\pgfqpoint{2.722577in}{3.156274in}}%
\pgfpathlineto{\pgfqpoint{2.725291in}{3.155559in}}%
\pgfpathlineto{\pgfqpoint{2.728106in}{3.161685in}}%
\pgfpathlineto{\pgfqpoint{2.731194in}{3.161335in}}%
\pgfpathlineto{\pgfqpoint{2.734752in}{3.168191in}}%
\pgfpathlineto{\pgfqpoint{2.739230in}{3.169203in}}%
\pgfpathlineto{\pgfqpoint{2.744126in}{3.177370in}}%
\pgfpathlineto{\pgfqpoint{2.749022in}{3.178790in}}%
\pgfpathlineto{\pgfqpoint{2.753918in}{3.186958in}}%
\pgfpathlineto{\pgfqpoint{2.758814in}{3.188376in}}%
\pgfpathlineto{\pgfqpoint{2.763710in}{3.196545in}}%
\pgfpathlineto{\pgfqpoint{2.768606in}{3.197963in}}%
\pgfpathlineto{\pgfqpoint{2.773502in}{3.206133in}}%
\pgfpathlineto{\pgfqpoint{2.778398in}{3.207549in}}%
\pgfpathlineto{\pgfqpoint{2.783294in}{3.215720in}}%
\pgfpathlineto{\pgfqpoint{2.788190in}{3.217136in}}%
\pgfpathlineto{\pgfqpoint{2.793086in}{3.225308in}}%
\pgfpathlineto{\pgfqpoint{2.797982in}{3.226722in}}%
\pgfpathlineto{\pgfqpoint{2.802878in}{3.234895in}}%
\pgfpathlineto{\pgfqpoint{2.807774in}{3.236309in}}%
\pgfpathlineto{\pgfqpoint{2.812670in}{3.244482in}}%
\pgfpathlineto{\pgfqpoint{2.817566in}{3.245895in}}%
\pgfpathlineto{\pgfqpoint{2.822462in}{3.254069in}}%
\pgfpathlineto{\pgfqpoint{2.827358in}{3.255481in}}%
\pgfpathlineto{\pgfqpoint{2.832254in}{3.263656in}}%
\pgfpathlineto{\pgfqpoint{2.837150in}{3.265067in}}%
\pgfpathlineto{\pgfqpoint{2.842046in}{3.273243in}}%
\pgfpathlineto{\pgfqpoint{2.846942in}{3.274653in}}%
\pgfpathlineto{\pgfqpoint{2.851838in}{3.282829in}}%
\pgfpathlineto{\pgfqpoint{2.856734in}{3.284238in}}%
\pgfpathlineto{\pgfqpoint{2.861630in}{3.292416in}}%
\pgfpathlineto{\pgfqpoint{2.866526in}{3.293824in}}%
\pgfpathlineto{\pgfqpoint{2.871422in}{3.302002in}}%
\pgfpathlineto{\pgfqpoint{2.876318in}{3.303410in}}%
\pgfpathlineto{\pgfqpoint{2.879855in}{3.310258in}}%
\pgfpathlineto{\pgfqpoint{2.883391in}{3.310333in}}%
\pgfpathlineto{\pgfqpoint{2.883443in}{3.255480in}}%
\pgfpathlineto{\pgfqpoint{2.883875in}{3.263866in}}%
\pgfpathlineto{\pgfqpoint{2.886068in}{3.290923in}}%
\pgfpathlineto{\pgfqpoint{2.886928in}{3.223454in}}%
\pgfpathlineto{\pgfqpoint{2.887810in}{3.269397in}}%
\pgfpathlineto{\pgfqpoint{2.888861in}{3.172423in}}%
\pgfpathlineto{\pgfqpoint{2.890119in}{3.205992in}}%
\pgfpathlineto{\pgfqpoint{2.891392in}{3.131319in}}%
\pgfpathlineto{\pgfqpoint{3.177102in}{3.131319in}}%
\pgfpathlineto{\pgfqpoint{3.177102in}{3.131319in}}%
\pgfusepath{stroke}%
\end{pgfscope}%
\begin{pgfscope}%
\pgfsetrectcap%
\pgfsetmiterjoin%
\pgfsetlinewidth{0.803000pt}%
\definecolor{currentstroke}{rgb}{0.000000,0.000000,0.000000}%
\pgfsetstrokecolor{currentstroke}%
\pgfsetdash{}{0pt}%
\pgfpathmoveto{\pgfqpoint{2.197902in}{3.064738in}}%
\pgfpathlineto{\pgfqpoint{2.197902in}{3.390988in}}%
\pgfusepath{stroke}%
\end{pgfscope}%
\begin{pgfscope}%
\pgfsetrectcap%
\pgfsetmiterjoin%
\pgfsetlinewidth{0.803000pt}%
\definecolor{currentstroke}{rgb}{0.000000,0.000000,0.000000}%
\pgfsetstrokecolor{currentstroke}%
\pgfsetdash{}{0pt}%
\pgfpathmoveto{\pgfqpoint{3.177102in}{3.064738in}}%
\pgfpathlineto{\pgfqpoint{3.177102in}{3.390988in}}%
\pgfusepath{stroke}%
\end{pgfscope}%
\begin{pgfscope}%
\pgfsetrectcap%
\pgfsetmiterjoin%
\pgfsetlinewidth{0.803000pt}%
\definecolor{currentstroke}{rgb}{0.000000,0.000000,0.000000}%
\pgfsetstrokecolor{currentstroke}%
\pgfsetdash{}{0pt}%
\pgfpathmoveto{\pgfqpoint{2.197902in}{3.064738in}}%
\pgfpathlineto{\pgfqpoint{3.177102in}{3.064738in}}%
\pgfusepath{stroke}%
\end{pgfscope}%
\begin{pgfscope}%
\pgfsetrectcap%
\pgfsetmiterjoin%
\pgfsetlinewidth{0.803000pt}%
\definecolor{currentstroke}{rgb}{0.000000,0.000000,0.000000}%
\pgfsetstrokecolor{currentstroke}%
\pgfsetdash{}{0pt}%
\pgfpathmoveto{\pgfqpoint{2.197902in}{3.390988in}}%
\pgfpathlineto{\pgfqpoint{3.177102in}{3.390988in}}%
\pgfusepath{stroke}%
\end{pgfscope}%
\begin{pgfscope}%
\pgfsetbuttcap%
\pgfsetmiterjoin%
\definecolor{currentfill}{rgb}{1.000000,1.000000,1.000000}%
\pgfsetfillcolor{currentfill}%
\pgfsetlinewidth{0.000000pt}%
\definecolor{currentstroke}{rgb}{0.000000,0.000000,0.000000}%
\pgfsetstrokecolor{currentstroke}%
\pgfsetstrokeopacity{0.000000}%
\pgfsetdash{}{0pt}%
\pgfpathmoveto{\pgfqpoint{2.197902in}{1.995363in}}%
\pgfpathlineto{\pgfqpoint{3.177102in}{1.995363in}}%
\pgfpathlineto{\pgfqpoint{3.177102in}{2.321613in}}%
\pgfpathlineto{\pgfqpoint{2.197902in}{2.321613in}}%
\pgfpathlineto{\pgfqpoint{2.197902in}{1.995363in}}%
\pgfpathclose%
\pgfusepath{fill}%
\end{pgfscope}%
\begin{pgfscope}%
\pgfpathrectangle{\pgfqpoint{2.197902in}{1.995363in}}{\pgfqpoint{0.979200in}{0.326250in}}%
\pgfusepath{clip}%
\pgfsetbuttcap%
\pgfsetroundjoin%
\pgfsetlinewidth{0.401500pt}%
\definecolor{currentstroke}{rgb}{0.690196,0.690196,0.690196}%
\pgfsetstrokecolor{currentstroke}%
\pgfsetstrokeopacity{0.250000}%
\pgfsetdash{{1.480000pt}{0.640000pt}}{0.000000pt}%
\pgfpathmoveto{\pgfqpoint{2.197902in}{1.995363in}}%
\pgfpathlineto{\pgfqpoint{2.197902in}{2.321613in}}%
\pgfusepath{stroke}%
\end{pgfscope}%
\begin{pgfscope}%
\pgfsetbuttcap%
\pgfsetroundjoin%
\definecolor{currentfill}{rgb}{0.000000,0.000000,0.000000}%
\pgfsetfillcolor{currentfill}%
\pgfsetlinewidth{0.803000pt}%
\definecolor{currentstroke}{rgb}{0.000000,0.000000,0.000000}%
\pgfsetstrokecolor{currentstroke}%
\pgfsetdash{}{0pt}%
\pgfsys@defobject{currentmarker}{\pgfqpoint{0.000000in}{0.000000in}}{\pgfqpoint{0.000000in}{0.027778in}}{%
\pgfpathmoveto{\pgfqpoint{0.000000in}{0.000000in}}%
\pgfpathlineto{\pgfqpoint{0.000000in}{0.027778in}}%
\pgfusepath{stroke,fill}%
}%
\begin{pgfscope}%
\pgfsys@transformshift{2.197902in}{1.995363in}%
\pgfsys@useobject{currentmarker}{}%
\end{pgfscope}%
\end{pgfscope}%
\begin{pgfscope}%
\definecolor{textcolor}{rgb}{0.000000,0.000000,0.000000}%
\pgfsetstrokecolor{textcolor}%
\pgfsetfillcolor{textcolor}%
\pgftext[x=2.197902in,y=1.981474in,,top]{\color{textcolor}{\rmfamily\fontsize{5.000000}{6.000000}\selectfont\catcode`\^=\active\def^{\ifmmode\sp\else\^{}\fi}\catcode`\%=\active\def%{\%}4.980}}%
\end{pgfscope}%
\begin{pgfscope}%
\pgfpathrectangle{\pgfqpoint{2.197902in}{1.995363in}}{\pgfqpoint{0.979200in}{0.326250in}}%
\pgfusepath{clip}%
\pgfsetbuttcap%
\pgfsetroundjoin%
\pgfsetlinewidth{0.401500pt}%
\definecolor{currentstroke}{rgb}{0.690196,0.690196,0.690196}%
\pgfsetstrokecolor{currentstroke}%
\pgfsetstrokeopacity{0.250000}%
\pgfsetdash{{1.480000pt}{0.640000pt}}{0.000000pt}%
\pgfpathmoveto{\pgfqpoint{2.442702in}{1.995363in}}%
\pgfpathlineto{\pgfqpoint{2.442702in}{2.321613in}}%
\pgfusepath{stroke}%
\end{pgfscope}%
\begin{pgfscope}%
\pgfsetbuttcap%
\pgfsetroundjoin%
\definecolor{currentfill}{rgb}{0.000000,0.000000,0.000000}%
\pgfsetfillcolor{currentfill}%
\pgfsetlinewidth{0.803000pt}%
\definecolor{currentstroke}{rgb}{0.000000,0.000000,0.000000}%
\pgfsetstrokecolor{currentstroke}%
\pgfsetdash{}{0pt}%
\pgfsys@defobject{currentmarker}{\pgfqpoint{0.000000in}{0.000000in}}{\pgfqpoint{0.000000in}{0.027778in}}{%
\pgfpathmoveto{\pgfqpoint{0.000000in}{0.000000in}}%
\pgfpathlineto{\pgfqpoint{0.000000in}{0.027778in}}%
\pgfusepath{stroke,fill}%
}%
\begin{pgfscope}%
\pgfsys@transformshift{2.442702in}{1.995363in}%
\pgfsys@useobject{currentmarker}{}%
\end{pgfscope}%
\end{pgfscope}%
\begin{pgfscope}%
\definecolor{textcolor}{rgb}{0.000000,0.000000,0.000000}%
\pgfsetstrokecolor{textcolor}%
\pgfsetfillcolor{textcolor}%
\pgftext[x=2.442702in,y=1.981474in,,top]{\color{textcolor}{\rmfamily\fontsize{5.000000}{6.000000}\selectfont\catcode`\^=\active\def^{\ifmmode\sp\else\^{}\fi}\catcode`\%=\active\def%{\%}4.985}}%
\end{pgfscope}%
\begin{pgfscope}%
\pgfpathrectangle{\pgfqpoint{2.197902in}{1.995363in}}{\pgfqpoint{0.979200in}{0.326250in}}%
\pgfusepath{clip}%
\pgfsetbuttcap%
\pgfsetroundjoin%
\pgfsetlinewidth{0.401500pt}%
\definecolor{currentstroke}{rgb}{0.690196,0.690196,0.690196}%
\pgfsetstrokecolor{currentstroke}%
\pgfsetstrokeopacity{0.250000}%
\pgfsetdash{{1.480000pt}{0.640000pt}}{0.000000pt}%
\pgfpathmoveto{\pgfqpoint{2.687502in}{1.995363in}}%
\pgfpathlineto{\pgfqpoint{2.687502in}{2.321613in}}%
\pgfusepath{stroke}%
\end{pgfscope}%
\begin{pgfscope}%
\pgfsetbuttcap%
\pgfsetroundjoin%
\definecolor{currentfill}{rgb}{0.000000,0.000000,0.000000}%
\pgfsetfillcolor{currentfill}%
\pgfsetlinewidth{0.803000pt}%
\definecolor{currentstroke}{rgb}{0.000000,0.000000,0.000000}%
\pgfsetstrokecolor{currentstroke}%
\pgfsetdash{}{0pt}%
\pgfsys@defobject{currentmarker}{\pgfqpoint{0.000000in}{0.000000in}}{\pgfqpoint{0.000000in}{0.027778in}}{%
\pgfpathmoveto{\pgfqpoint{0.000000in}{0.000000in}}%
\pgfpathlineto{\pgfqpoint{0.000000in}{0.027778in}}%
\pgfusepath{stroke,fill}%
}%
\begin{pgfscope}%
\pgfsys@transformshift{2.687502in}{1.995363in}%
\pgfsys@useobject{currentmarker}{}%
\end{pgfscope}%
\end{pgfscope}%
\begin{pgfscope}%
\definecolor{textcolor}{rgb}{0.000000,0.000000,0.000000}%
\pgfsetstrokecolor{textcolor}%
\pgfsetfillcolor{textcolor}%
\pgftext[x=2.687502in,y=1.981474in,,top]{\color{textcolor}{\rmfamily\fontsize{5.000000}{6.000000}\selectfont\catcode`\^=\active\def^{\ifmmode\sp\else\^{}\fi}\catcode`\%=\active\def%{\%}4.990}}%
\end{pgfscope}%
\begin{pgfscope}%
\pgfpathrectangle{\pgfqpoint{2.197902in}{1.995363in}}{\pgfqpoint{0.979200in}{0.326250in}}%
\pgfusepath{clip}%
\pgfsetbuttcap%
\pgfsetroundjoin%
\pgfsetlinewidth{0.401500pt}%
\definecolor{currentstroke}{rgb}{0.690196,0.690196,0.690196}%
\pgfsetstrokecolor{currentstroke}%
\pgfsetstrokeopacity{0.250000}%
\pgfsetdash{{1.480000pt}{0.640000pt}}{0.000000pt}%
\pgfpathmoveto{\pgfqpoint{2.932302in}{1.995363in}}%
\pgfpathlineto{\pgfqpoint{2.932302in}{2.321613in}}%
\pgfusepath{stroke}%
\end{pgfscope}%
\begin{pgfscope}%
\pgfsetbuttcap%
\pgfsetroundjoin%
\definecolor{currentfill}{rgb}{0.000000,0.000000,0.000000}%
\pgfsetfillcolor{currentfill}%
\pgfsetlinewidth{0.803000pt}%
\definecolor{currentstroke}{rgb}{0.000000,0.000000,0.000000}%
\pgfsetstrokecolor{currentstroke}%
\pgfsetdash{}{0pt}%
\pgfsys@defobject{currentmarker}{\pgfqpoint{0.000000in}{0.000000in}}{\pgfqpoint{0.000000in}{0.027778in}}{%
\pgfpathmoveto{\pgfqpoint{0.000000in}{0.000000in}}%
\pgfpathlineto{\pgfqpoint{0.000000in}{0.027778in}}%
\pgfusepath{stroke,fill}%
}%
\begin{pgfscope}%
\pgfsys@transformshift{2.932302in}{1.995363in}%
\pgfsys@useobject{currentmarker}{}%
\end{pgfscope}%
\end{pgfscope}%
\begin{pgfscope}%
\definecolor{textcolor}{rgb}{0.000000,0.000000,0.000000}%
\pgfsetstrokecolor{textcolor}%
\pgfsetfillcolor{textcolor}%
\pgftext[x=2.932302in,y=1.981474in,,top]{\color{textcolor}{\rmfamily\fontsize{5.000000}{6.000000}\selectfont\catcode`\^=\active\def^{\ifmmode\sp\else\^{}\fi}\catcode`\%=\active\def%{\%}4.995}}%
\end{pgfscope}%
\begin{pgfscope}%
\pgfpathrectangle{\pgfqpoint{2.197902in}{1.995363in}}{\pgfqpoint{0.979200in}{0.326250in}}%
\pgfusepath{clip}%
\pgfsetbuttcap%
\pgfsetroundjoin%
\pgfsetlinewidth{0.401500pt}%
\definecolor{currentstroke}{rgb}{0.690196,0.690196,0.690196}%
\pgfsetstrokecolor{currentstroke}%
\pgfsetstrokeopacity{0.250000}%
\pgfsetdash{{1.480000pt}{0.640000pt}}{0.000000pt}%
\pgfpathmoveto{\pgfqpoint{3.177102in}{1.995363in}}%
\pgfpathlineto{\pgfqpoint{3.177102in}{2.321613in}}%
\pgfusepath{stroke}%
\end{pgfscope}%
\begin{pgfscope}%
\pgfsetbuttcap%
\pgfsetroundjoin%
\definecolor{currentfill}{rgb}{0.000000,0.000000,0.000000}%
\pgfsetfillcolor{currentfill}%
\pgfsetlinewidth{0.803000pt}%
\definecolor{currentstroke}{rgb}{0.000000,0.000000,0.000000}%
\pgfsetstrokecolor{currentstroke}%
\pgfsetdash{}{0pt}%
\pgfsys@defobject{currentmarker}{\pgfqpoint{0.000000in}{0.000000in}}{\pgfqpoint{0.000000in}{0.027778in}}{%
\pgfpathmoveto{\pgfqpoint{0.000000in}{0.000000in}}%
\pgfpathlineto{\pgfqpoint{0.000000in}{0.027778in}}%
\pgfusepath{stroke,fill}%
}%
\begin{pgfscope}%
\pgfsys@transformshift{3.177102in}{1.995363in}%
\pgfsys@useobject{currentmarker}{}%
\end{pgfscope}%
\end{pgfscope}%
\begin{pgfscope}%
\definecolor{textcolor}{rgb}{0.000000,0.000000,0.000000}%
\pgfsetstrokecolor{textcolor}%
\pgfsetfillcolor{textcolor}%
\pgftext[x=3.177102in,y=1.981474in,,top]{\color{textcolor}{\rmfamily\fontsize{5.000000}{6.000000}\selectfont\catcode`\^=\active\def^{\ifmmode\sp\else\^{}\fi}\catcode`\%=\active\def%{\%}5.000}}%
\end{pgfscope}%
\begin{pgfscope}%
\pgfpathrectangle{\pgfqpoint{2.197902in}{1.995363in}}{\pgfqpoint{0.979200in}{0.326250in}}%
\pgfusepath{clip}%
\pgfsetbuttcap%
\pgfsetroundjoin%
\pgfsetlinewidth{0.401500pt}%
\definecolor{currentstroke}{rgb}{0.690196,0.690196,0.690196}%
\pgfsetstrokecolor{currentstroke}%
\pgfsetstrokeopacity{0.250000}%
\pgfsetdash{{1.480000pt}{0.640000pt}}{0.000000pt}%
\pgfpathmoveto{\pgfqpoint{2.197902in}{1.995363in}}%
\pgfpathlineto{\pgfqpoint{3.177102in}{1.995363in}}%
\pgfusepath{stroke}%
\end{pgfscope}%
\begin{pgfscope}%
\pgfsetbuttcap%
\pgfsetroundjoin%
\definecolor{currentfill}{rgb}{0.000000,0.000000,0.000000}%
\pgfsetfillcolor{currentfill}%
\pgfsetlinewidth{0.803000pt}%
\definecolor{currentstroke}{rgb}{0.000000,0.000000,0.000000}%
\pgfsetstrokecolor{currentstroke}%
\pgfsetdash{}{0pt}%
\pgfsys@defobject{currentmarker}{\pgfqpoint{0.000000in}{0.000000in}}{\pgfqpoint{0.027778in}{0.000000in}}{%
\pgfpathmoveto{\pgfqpoint{0.000000in}{0.000000in}}%
\pgfpathlineto{\pgfqpoint{0.027778in}{0.000000in}}%
\pgfusepath{stroke,fill}%
}%
\begin{pgfscope}%
\pgfsys@transformshift{2.197902in}{1.995363in}%
\pgfsys@useobject{currentmarker}{}%
\end{pgfscope}%
\end{pgfscope}%
\begin{pgfscope}%
\definecolor{textcolor}{rgb}{0.000000,0.000000,0.000000}%
\pgfsetstrokecolor{textcolor}%
\pgfsetfillcolor{textcolor}%
\pgftext[x=1.966996in, y=1.971250in, left, base]{\color{textcolor}{\rmfamily\fontsize{5.000000}{6.000000}\selectfont\catcode`\^=\active\def^{\ifmmode\sp\else\^{}\fi}\catcode`\%=\active\def%{\%}45.00}}%
\end{pgfscope}%
\begin{pgfscope}%
\pgfpathrectangle{\pgfqpoint{2.197902in}{1.995363in}}{\pgfqpoint{0.979200in}{0.326250in}}%
\pgfusepath{clip}%
\pgfsetbuttcap%
\pgfsetroundjoin%
\pgfsetlinewidth{0.401500pt}%
\definecolor{currentstroke}{rgb}{0.690196,0.690196,0.690196}%
\pgfsetstrokecolor{currentstroke}%
\pgfsetstrokeopacity{0.250000}%
\pgfsetdash{{1.480000pt}{0.640000pt}}{0.000000pt}%
\pgfpathmoveto{\pgfqpoint{2.197902in}{2.158488in}}%
\pgfpathlineto{\pgfqpoint{3.177102in}{2.158488in}}%
\pgfusepath{stroke}%
\end{pgfscope}%
\begin{pgfscope}%
\pgfsetbuttcap%
\pgfsetroundjoin%
\definecolor{currentfill}{rgb}{0.000000,0.000000,0.000000}%
\pgfsetfillcolor{currentfill}%
\pgfsetlinewidth{0.803000pt}%
\definecolor{currentstroke}{rgb}{0.000000,0.000000,0.000000}%
\pgfsetstrokecolor{currentstroke}%
\pgfsetdash{}{0pt}%
\pgfsys@defobject{currentmarker}{\pgfqpoint{0.000000in}{0.000000in}}{\pgfqpoint{0.027778in}{0.000000in}}{%
\pgfpathmoveto{\pgfqpoint{0.000000in}{0.000000in}}%
\pgfpathlineto{\pgfqpoint{0.027778in}{0.000000in}}%
\pgfusepath{stroke,fill}%
}%
\begin{pgfscope}%
\pgfsys@transformshift{2.197902in}{2.158488in}%
\pgfsys@useobject{currentmarker}{}%
\end{pgfscope}%
\end{pgfscope}%
\begin{pgfscope}%
\definecolor{textcolor}{rgb}{0.000000,0.000000,0.000000}%
\pgfsetstrokecolor{textcolor}%
\pgfsetfillcolor{textcolor}%
\pgftext[x=1.966996in, y=2.134375in, left, base]{\color{textcolor}{\rmfamily\fontsize{5.000000}{6.000000}\selectfont\catcode`\^=\active\def^{\ifmmode\sp\else\^{}\fi}\catcode`\%=\active\def%{\%}65.00}}%
\end{pgfscope}%
\begin{pgfscope}%
\pgfpathrectangle{\pgfqpoint{2.197902in}{1.995363in}}{\pgfqpoint{0.979200in}{0.326250in}}%
\pgfusepath{clip}%
\pgfsetbuttcap%
\pgfsetroundjoin%
\pgfsetlinewidth{0.401500pt}%
\definecolor{currentstroke}{rgb}{0.690196,0.690196,0.690196}%
\pgfsetstrokecolor{currentstroke}%
\pgfsetstrokeopacity{0.250000}%
\pgfsetdash{{1.480000pt}{0.640000pt}}{0.000000pt}%
\pgfpathmoveto{\pgfqpoint{2.197902in}{2.321613in}}%
\pgfpathlineto{\pgfqpoint{3.177102in}{2.321613in}}%
\pgfusepath{stroke}%
\end{pgfscope}%
\begin{pgfscope}%
\pgfsetbuttcap%
\pgfsetroundjoin%
\definecolor{currentfill}{rgb}{0.000000,0.000000,0.000000}%
\pgfsetfillcolor{currentfill}%
\pgfsetlinewidth{0.803000pt}%
\definecolor{currentstroke}{rgb}{0.000000,0.000000,0.000000}%
\pgfsetstrokecolor{currentstroke}%
\pgfsetdash{}{0pt}%
\pgfsys@defobject{currentmarker}{\pgfqpoint{0.000000in}{0.000000in}}{\pgfqpoint{0.027778in}{0.000000in}}{%
\pgfpathmoveto{\pgfqpoint{0.000000in}{0.000000in}}%
\pgfpathlineto{\pgfqpoint{0.027778in}{0.000000in}}%
\pgfusepath{stroke,fill}%
}%
\begin{pgfscope}%
\pgfsys@transformshift{2.197902in}{2.321613in}%
\pgfsys@useobject{currentmarker}{}%
\end{pgfscope}%
\end{pgfscope}%
\begin{pgfscope}%
\definecolor{textcolor}{rgb}{0.000000,0.000000,0.000000}%
\pgfsetstrokecolor{textcolor}%
\pgfsetfillcolor{textcolor}%
\pgftext[x=1.966996in, y=2.297500in, left, base]{\color{textcolor}{\rmfamily\fontsize{5.000000}{6.000000}\selectfont\catcode`\^=\active\def^{\ifmmode\sp\else\^{}\fi}\catcode`\%=\active\def%{\%}85.00}}%
\end{pgfscope}%
\begin{pgfscope}%
\pgfpathrectangle{\pgfqpoint{2.197902in}{1.995363in}}{\pgfqpoint{0.979200in}{0.326250in}}%
\pgfusepath{clip}%
\pgfsetrectcap%
\pgfsetroundjoin%
\pgfsetlinewidth{1.505625pt}%
\definecolor{currentstroke}{rgb}{0.835294,0.368627,0.000000}%
\pgfsetstrokecolor{currentstroke}%
\pgfsetdash{}{0pt}%
\pgfpathmoveto{\pgfqpoint{2.400321in}{1.988418in}}%
\pgfpathlineto{\pgfqpoint{2.401005in}{2.131345in}}%
\pgfpathlineto{\pgfqpoint{2.403972in}{2.153797in}}%
\pgfpathlineto{\pgfqpoint{2.405533in}{2.165745in}}%
\pgfpathlineto{\pgfqpoint{2.407852in}{2.175458in}}%
\pgfpathlineto{\pgfqpoint{2.411453in}{2.123641in}}%
\pgfpathlineto{\pgfqpoint{2.416349in}{2.176639in}}%
\pgfpathlineto{\pgfqpoint{2.421245in}{2.113491in}}%
\pgfpathlineto{\pgfqpoint{2.426141in}{2.177859in}}%
\pgfpathlineto{\pgfqpoint{2.430410in}{2.121300in}}%
\pgfpathlineto{\pgfqpoint{2.433968in}{2.167875in}}%
\pgfpathlineto{\pgfqpoint{2.435117in}{2.133331in}}%
\pgfpathlineto{\pgfqpoint{2.435988in}{2.132852in}}%
\pgfpathlineto{\pgfqpoint{2.437330in}{2.159890in}}%
\pgfpathlineto{\pgfqpoint{2.438575in}{2.138522in}}%
\pgfpathlineto{\pgfqpoint{2.439273in}{2.139759in}}%
\pgfpathlineto{\pgfqpoint{2.440127in}{2.151066in}}%
\pgfpathlineto{\pgfqpoint{2.441188in}{2.151095in}}%
\pgfpathlineto{\pgfqpoint{2.442301in}{2.141779in}}%
\pgfpathlineto{\pgfqpoint{2.445724in}{2.150036in}}%
\pgfpathlineto{\pgfqpoint{2.447589in}{2.142802in}}%
\pgfpathlineto{\pgfqpoint{2.451317in}{2.149603in}}%
\pgfpathlineto{\pgfqpoint{2.455496in}{2.145069in}}%
\pgfpathlineto{\pgfqpoint{2.460392in}{2.147229in}}%
\pgfpathlineto{\pgfqpoint{2.465288in}{2.147457in}}%
\pgfpathlineto{\pgfqpoint{2.470184in}{2.145236in}}%
\pgfpathlineto{\pgfqpoint{2.475080in}{2.149101in}}%
\pgfpathlineto{\pgfqpoint{2.479976in}{2.144338in}}%
\pgfpathlineto{\pgfqpoint{2.484872in}{2.149434in}}%
\pgfpathlineto{\pgfqpoint{2.489768in}{2.144757in}}%
\pgfpathlineto{\pgfqpoint{2.494664in}{2.148576in}}%
\pgfpathlineto{\pgfqpoint{2.499560in}{2.146076in}}%
\pgfpathlineto{\pgfqpoint{2.509352in}{2.147550in}}%
\pgfpathlineto{\pgfqpoint{2.514248in}{2.145933in}}%
\pgfpathlineto{\pgfqpoint{2.519144in}{2.148503in}}%
\pgfpathlineto{\pgfqpoint{2.524040in}{2.145426in}}%
\pgfpathlineto{\pgfqpoint{2.528936in}{2.148615in}}%
\pgfpathlineto{\pgfqpoint{2.533832in}{2.145734in}}%
\pgfpathlineto{\pgfqpoint{2.538728in}{2.147998in}}%
\pgfpathlineto{\pgfqpoint{2.543624in}{2.146567in}}%
\pgfpathlineto{\pgfqpoint{2.553416in}{2.147448in}}%
\pgfpathlineto{\pgfqpoint{2.558312in}{2.146293in}}%
\pgfpathlineto{\pgfqpoint{2.563208in}{2.147965in}}%
\pgfpathlineto{\pgfqpoint{2.568104in}{2.145981in}}%
\pgfpathlineto{\pgfqpoint{2.573000in}{2.147941in}}%
\pgfpathlineto{\pgfqpoint{2.577896in}{2.146167in}}%
\pgfpathlineto{\pgfqpoint{2.582792in}{2.147470in}}%
\pgfpathlineto{\pgfqpoint{2.587688in}{2.146648in}}%
\pgfpathlineto{\pgfqpoint{2.597480in}{2.147119in}}%
\pgfpathlineto{\pgfqpoint{2.602376in}{2.146279in}}%
\pgfpathlineto{\pgfqpoint{2.607272in}{2.147328in}}%
\pgfpathlineto{\pgfqpoint{2.612168in}{2.146018in}}%
\pgfpathlineto{\pgfqpoint{2.617064in}{2.147171in}}%
\pgfpathlineto{\pgfqpoint{2.621960in}{2.146033in}}%
\pgfpathlineto{\pgfqpoint{2.626856in}{2.146705in}}%
\pgfpathlineto{\pgfqpoint{2.636648in}{2.146047in}}%
\pgfpathlineto{\pgfqpoint{2.641544in}{2.146078in}}%
\pgfpathlineto{\pgfqpoint{2.646440in}{2.144941in}}%
\pgfpathlineto{\pgfqpoint{2.651336in}{2.136646in}}%
\pgfpathlineto{\pgfqpoint{2.656232in}{2.111031in}}%
\pgfpathlineto{\pgfqpoint{2.661128in}{2.073860in}}%
\pgfpathlineto{\pgfqpoint{2.669860in}{1.988418in}}%
\pgfpathmoveto{\pgfqpoint{2.889782in}{1.988418in}}%
\pgfpathlineto{\pgfqpoint{2.890119in}{2.056782in}}%
\pgfpathlineto{\pgfqpoint{2.891392in}{2.141057in}}%
\pgfpathlineto{\pgfqpoint{2.894913in}{2.167703in}}%
\pgfpathlineto{\pgfqpoint{2.896285in}{2.173636in}}%
\pgfpathlineto{\pgfqpoint{2.897974in}{2.161810in}}%
\pgfpathlineto{\pgfqpoint{2.900025in}{2.135461in}}%
\pgfpathlineto{\pgfqpoint{2.903548in}{2.153229in}}%
\pgfpathlineto{\pgfqpoint{2.907336in}{2.145271in}}%
\pgfpathlineto{\pgfqpoint{2.910832in}{2.142310in}}%
\pgfpathlineto{\pgfqpoint{2.915728in}{2.153472in}}%
\pgfpathlineto{\pgfqpoint{2.920624in}{2.137790in}}%
\pgfpathlineto{\pgfqpoint{2.925520in}{2.155778in}}%
\pgfpathlineto{\pgfqpoint{2.930416in}{2.138071in}}%
\pgfpathlineto{\pgfqpoint{2.935312in}{2.153532in}}%
\pgfpathlineto{\pgfqpoint{2.940208in}{2.142135in}}%
\pgfpathlineto{\pgfqpoint{2.945104in}{2.148634in}}%
\pgfpathlineto{\pgfqpoint{2.950000in}{2.147482in}}%
\pgfpathlineto{\pgfqpoint{2.954896in}{2.143814in}}%
\pgfpathlineto{\pgfqpoint{2.959792in}{2.151505in}}%
\pgfpathlineto{\pgfqpoint{2.964688in}{2.141248in}}%
\pgfpathlineto{\pgfqpoint{2.969584in}{2.152680in}}%
\pgfpathlineto{\pgfqpoint{2.974480in}{2.141688in}}%
\pgfpathlineto{\pgfqpoint{2.979376in}{2.151042in}}%
\pgfpathlineto{\pgfqpoint{2.984272in}{2.144390in}}%
\pgfpathlineto{\pgfqpoint{2.989168in}{2.147879in}}%
\pgfpathlineto{\pgfqpoint{2.994064in}{2.147735in}}%
\pgfpathlineto{\pgfqpoint{2.998960in}{2.144914in}}%
\pgfpathlineto{\pgfqpoint{3.003856in}{2.150117in}}%
\pgfpathlineto{\pgfqpoint{3.008752in}{2.143451in}}%
\pgfpathlineto{\pgfqpoint{3.013648in}{2.150661in}}%
\pgfpathlineto{\pgfqpoint{3.018544in}{2.143869in}}%
\pgfpathlineto{\pgfqpoint{3.023440in}{2.149478in}}%
\pgfpathlineto{\pgfqpoint{3.028336in}{2.145628in}}%
\pgfpathlineto{\pgfqpoint{3.033232in}{2.147427in}}%
\pgfpathlineto{\pgfqpoint{3.038128in}{2.147685in}}%
\pgfpathlineto{\pgfqpoint{3.043024in}{2.145588in}}%
\pgfpathlineto{\pgfqpoint{3.047920in}{2.149056in}}%
\pgfpathlineto{\pgfqpoint{3.052816in}{2.144737in}}%
\pgfpathlineto{\pgfqpoint{3.057712in}{2.149245in}}%
\pgfpathlineto{\pgfqpoint{3.062608in}{2.145053in}}%
\pgfpathlineto{\pgfqpoint{3.067504in}{2.148375in}}%
\pgfpathlineto{\pgfqpoint{3.072400in}{2.146153in}}%
\pgfpathlineto{\pgfqpoint{3.082192in}{2.147361in}}%
\pgfpathlineto{\pgfqpoint{3.087088in}{2.145833in}}%
\pgfpathlineto{\pgfqpoint{3.091984in}{2.148082in}}%
\pgfpathlineto{\pgfqpoint{3.096880in}{2.145284in}}%
\pgfpathlineto{\pgfqpoint{3.101776in}{2.148033in}}%
\pgfpathlineto{\pgfqpoint{3.106672in}{2.145433in}}%
\pgfpathlineto{\pgfqpoint{3.111568in}{2.147311in}}%
\pgfpathlineto{\pgfqpoint{3.116464in}{2.145995in}}%
\pgfpathlineto{\pgfqpoint{3.126256in}{2.146488in}}%
\pgfpathlineto{\pgfqpoint{3.131152in}{2.145169in}}%
\pgfpathlineto{\pgfqpoint{3.136048in}{2.146131in}}%
\pgfpathlineto{\pgfqpoint{3.140944in}{2.134574in}}%
\pgfpathlineto{\pgfqpoint{3.145840in}{2.109553in}}%
\pgfpathlineto{\pgfqpoint{3.150736in}{2.071514in}}%
\pgfpathlineto{\pgfqpoint{3.159226in}{1.988418in}}%
\pgfpathlineto{\pgfqpoint{3.159226in}{1.988418in}}%
\pgfusepath{stroke}%
\end{pgfscope}%
\begin{pgfscope}%
\pgfsetrectcap%
\pgfsetmiterjoin%
\pgfsetlinewidth{0.803000pt}%
\definecolor{currentstroke}{rgb}{0.000000,0.000000,0.000000}%
\pgfsetstrokecolor{currentstroke}%
\pgfsetdash{}{0pt}%
\pgfpathmoveto{\pgfqpoint{2.197902in}{1.995363in}}%
\pgfpathlineto{\pgfqpoint{2.197902in}{2.321613in}}%
\pgfusepath{stroke}%
\end{pgfscope}%
\begin{pgfscope}%
\pgfsetrectcap%
\pgfsetmiterjoin%
\pgfsetlinewidth{0.803000pt}%
\definecolor{currentstroke}{rgb}{0.000000,0.000000,0.000000}%
\pgfsetstrokecolor{currentstroke}%
\pgfsetdash{}{0pt}%
\pgfpathmoveto{\pgfqpoint{3.177102in}{1.995363in}}%
\pgfpathlineto{\pgfqpoint{3.177102in}{2.321613in}}%
\pgfusepath{stroke}%
\end{pgfscope}%
\begin{pgfscope}%
\pgfsetrectcap%
\pgfsetmiterjoin%
\pgfsetlinewidth{0.803000pt}%
\definecolor{currentstroke}{rgb}{0.000000,0.000000,0.000000}%
\pgfsetstrokecolor{currentstroke}%
\pgfsetdash{}{0pt}%
\pgfpathmoveto{\pgfqpoint{2.197902in}{1.995363in}}%
\pgfpathlineto{\pgfqpoint{3.177102in}{1.995363in}}%
\pgfusepath{stroke}%
\end{pgfscope}%
\begin{pgfscope}%
\pgfsetrectcap%
\pgfsetmiterjoin%
\pgfsetlinewidth{0.803000pt}%
\definecolor{currentstroke}{rgb}{0.000000,0.000000,0.000000}%
\pgfsetstrokecolor{currentstroke}%
\pgfsetdash{}{0pt}%
\pgfpathmoveto{\pgfqpoint{2.197902in}{2.321613in}}%
\pgfpathlineto{\pgfqpoint{3.177102in}{2.321613in}}%
\pgfusepath{stroke}%
\end{pgfscope}%
\begin{pgfscope}%
\pgfsetbuttcap%
\pgfsetmiterjoin%
\definecolor{currentfill}{rgb}{1.000000,1.000000,1.000000}%
\pgfsetfillcolor{currentfill}%
\pgfsetlinewidth{0.000000pt}%
\definecolor{currentstroke}{rgb}{0.000000,0.000000,0.000000}%
\pgfsetstrokecolor{currentstroke}%
\pgfsetstrokeopacity{0.000000}%
\pgfsetdash{}{0pt}%
\pgfpathmoveto{\pgfqpoint{2.197902in}{0.579321in}}%
\pgfpathlineto{\pgfqpoint{3.177102in}{0.579321in}}%
\pgfpathlineto{\pgfqpoint{3.177102in}{0.905571in}}%
\pgfpathlineto{\pgfqpoint{2.197902in}{0.905571in}}%
\pgfpathlineto{\pgfqpoint{2.197902in}{0.579321in}}%
\pgfpathclose%
\pgfusepath{fill}%
\end{pgfscope}%
\begin{pgfscope}%
\pgfpathrectangle{\pgfqpoint{2.197902in}{0.579321in}}{\pgfqpoint{0.979200in}{0.326250in}}%
\pgfusepath{clip}%
\pgfsetbuttcap%
\pgfsetroundjoin%
\pgfsetlinewidth{0.401500pt}%
\definecolor{currentstroke}{rgb}{0.690196,0.690196,0.690196}%
\pgfsetstrokecolor{currentstroke}%
\pgfsetstrokeopacity{0.250000}%
\pgfsetdash{{1.480000pt}{0.640000pt}}{0.000000pt}%
\pgfpathmoveto{\pgfqpoint{2.197902in}{0.579321in}}%
\pgfpathlineto{\pgfqpoint{2.197902in}{0.905571in}}%
\pgfusepath{stroke}%
\end{pgfscope}%
\begin{pgfscope}%
\pgfsetbuttcap%
\pgfsetroundjoin%
\definecolor{currentfill}{rgb}{0.000000,0.000000,0.000000}%
\pgfsetfillcolor{currentfill}%
\pgfsetlinewidth{0.803000pt}%
\definecolor{currentstroke}{rgb}{0.000000,0.000000,0.000000}%
\pgfsetstrokecolor{currentstroke}%
\pgfsetdash{}{0pt}%
\pgfsys@defobject{currentmarker}{\pgfqpoint{0.000000in}{0.000000in}}{\pgfqpoint{0.000000in}{0.027778in}}{%
\pgfpathmoveto{\pgfqpoint{0.000000in}{0.000000in}}%
\pgfpathlineto{\pgfqpoint{0.000000in}{0.027778in}}%
\pgfusepath{stroke,fill}%
}%
\begin{pgfscope}%
\pgfsys@transformshift{2.197902in}{0.579321in}%
\pgfsys@useobject{currentmarker}{}%
\end{pgfscope}%
\end{pgfscope}%
\begin{pgfscope}%
\definecolor{textcolor}{rgb}{0.000000,0.000000,0.000000}%
\pgfsetstrokecolor{textcolor}%
\pgfsetfillcolor{textcolor}%
\pgftext[x=2.197902in,y=0.565432in,,top]{\color{textcolor}{\rmfamily\fontsize{5.000000}{6.000000}\selectfont\catcode`\^=\active\def^{\ifmmode\sp\else\^{}\fi}\catcode`\%=\active\def%{\%}4.980}}%
\end{pgfscope}%
\begin{pgfscope}%
\pgfpathrectangle{\pgfqpoint{2.197902in}{0.579321in}}{\pgfqpoint{0.979200in}{0.326250in}}%
\pgfusepath{clip}%
\pgfsetbuttcap%
\pgfsetroundjoin%
\pgfsetlinewidth{0.401500pt}%
\definecolor{currentstroke}{rgb}{0.690196,0.690196,0.690196}%
\pgfsetstrokecolor{currentstroke}%
\pgfsetstrokeopacity{0.250000}%
\pgfsetdash{{1.480000pt}{0.640000pt}}{0.000000pt}%
\pgfpathmoveto{\pgfqpoint{2.442702in}{0.579321in}}%
\pgfpathlineto{\pgfqpoint{2.442702in}{0.905571in}}%
\pgfusepath{stroke}%
\end{pgfscope}%
\begin{pgfscope}%
\pgfsetbuttcap%
\pgfsetroundjoin%
\definecolor{currentfill}{rgb}{0.000000,0.000000,0.000000}%
\pgfsetfillcolor{currentfill}%
\pgfsetlinewidth{0.803000pt}%
\definecolor{currentstroke}{rgb}{0.000000,0.000000,0.000000}%
\pgfsetstrokecolor{currentstroke}%
\pgfsetdash{}{0pt}%
\pgfsys@defobject{currentmarker}{\pgfqpoint{0.000000in}{0.000000in}}{\pgfqpoint{0.000000in}{0.027778in}}{%
\pgfpathmoveto{\pgfqpoint{0.000000in}{0.000000in}}%
\pgfpathlineto{\pgfqpoint{0.000000in}{0.027778in}}%
\pgfusepath{stroke,fill}%
}%
\begin{pgfscope}%
\pgfsys@transformshift{2.442702in}{0.579321in}%
\pgfsys@useobject{currentmarker}{}%
\end{pgfscope}%
\end{pgfscope}%
\begin{pgfscope}%
\definecolor{textcolor}{rgb}{0.000000,0.000000,0.000000}%
\pgfsetstrokecolor{textcolor}%
\pgfsetfillcolor{textcolor}%
\pgftext[x=2.442702in,y=0.565432in,,top]{\color{textcolor}{\rmfamily\fontsize{5.000000}{6.000000}\selectfont\catcode`\^=\active\def^{\ifmmode\sp\else\^{}\fi}\catcode`\%=\active\def%{\%}4.985}}%
\end{pgfscope}%
\begin{pgfscope}%
\pgfpathrectangle{\pgfqpoint{2.197902in}{0.579321in}}{\pgfqpoint{0.979200in}{0.326250in}}%
\pgfusepath{clip}%
\pgfsetbuttcap%
\pgfsetroundjoin%
\pgfsetlinewidth{0.401500pt}%
\definecolor{currentstroke}{rgb}{0.690196,0.690196,0.690196}%
\pgfsetstrokecolor{currentstroke}%
\pgfsetstrokeopacity{0.250000}%
\pgfsetdash{{1.480000pt}{0.640000pt}}{0.000000pt}%
\pgfpathmoveto{\pgfqpoint{2.687502in}{0.579321in}}%
\pgfpathlineto{\pgfqpoint{2.687502in}{0.905571in}}%
\pgfusepath{stroke}%
\end{pgfscope}%
\begin{pgfscope}%
\pgfsetbuttcap%
\pgfsetroundjoin%
\definecolor{currentfill}{rgb}{0.000000,0.000000,0.000000}%
\pgfsetfillcolor{currentfill}%
\pgfsetlinewidth{0.803000pt}%
\definecolor{currentstroke}{rgb}{0.000000,0.000000,0.000000}%
\pgfsetstrokecolor{currentstroke}%
\pgfsetdash{}{0pt}%
\pgfsys@defobject{currentmarker}{\pgfqpoint{0.000000in}{0.000000in}}{\pgfqpoint{0.000000in}{0.027778in}}{%
\pgfpathmoveto{\pgfqpoint{0.000000in}{0.000000in}}%
\pgfpathlineto{\pgfqpoint{0.000000in}{0.027778in}}%
\pgfusepath{stroke,fill}%
}%
\begin{pgfscope}%
\pgfsys@transformshift{2.687502in}{0.579321in}%
\pgfsys@useobject{currentmarker}{}%
\end{pgfscope}%
\end{pgfscope}%
\begin{pgfscope}%
\definecolor{textcolor}{rgb}{0.000000,0.000000,0.000000}%
\pgfsetstrokecolor{textcolor}%
\pgfsetfillcolor{textcolor}%
\pgftext[x=2.687502in,y=0.565432in,,top]{\color{textcolor}{\rmfamily\fontsize{5.000000}{6.000000}\selectfont\catcode`\^=\active\def^{\ifmmode\sp\else\^{}\fi}\catcode`\%=\active\def%{\%}4.990}}%
\end{pgfscope}%
\begin{pgfscope}%
\pgfpathrectangle{\pgfqpoint{2.197902in}{0.579321in}}{\pgfqpoint{0.979200in}{0.326250in}}%
\pgfusepath{clip}%
\pgfsetbuttcap%
\pgfsetroundjoin%
\pgfsetlinewidth{0.401500pt}%
\definecolor{currentstroke}{rgb}{0.690196,0.690196,0.690196}%
\pgfsetstrokecolor{currentstroke}%
\pgfsetstrokeopacity{0.250000}%
\pgfsetdash{{1.480000pt}{0.640000pt}}{0.000000pt}%
\pgfpathmoveto{\pgfqpoint{2.932302in}{0.579321in}}%
\pgfpathlineto{\pgfqpoint{2.932302in}{0.905571in}}%
\pgfusepath{stroke}%
\end{pgfscope}%
\begin{pgfscope}%
\pgfsetbuttcap%
\pgfsetroundjoin%
\definecolor{currentfill}{rgb}{0.000000,0.000000,0.000000}%
\pgfsetfillcolor{currentfill}%
\pgfsetlinewidth{0.803000pt}%
\definecolor{currentstroke}{rgb}{0.000000,0.000000,0.000000}%
\pgfsetstrokecolor{currentstroke}%
\pgfsetdash{}{0pt}%
\pgfsys@defobject{currentmarker}{\pgfqpoint{0.000000in}{0.000000in}}{\pgfqpoint{0.000000in}{0.027778in}}{%
\pgfpathmoveto{\pgfqpoint{0.000000in}{0.000000in}}%
\pgfpathlineto{\pgfqpoint{0.000000in}{0.027778in}}%
\pgfusepath{stroke,fill}%
}%
\begin{pgfscope}%
\pgfsys@transformshift{2.932302in}{0.579321in}%
\pgfsys@useobject{currentmarker}{}%
\end{pgfscope}%
\end{pgfscope}%
\begin{pgfscope}%
\definecolor{textcolor}{rgb}{0.000000,0.000000,0.000000}%
\pgfsetstrokecolor{textcolor}%
\pgfsetfillcolor{textcolor}%
\pgftext[x=2.932302in,y=0.565432in,,top]{\color{textcolor}{\rmfamily\fontsize{5.000000}{6.000000}\selectfont\catcode`\^=\active\def^{\ifmmode\sp\else\^{}\fi}\catcode`\%=\active\def%{\%}4.995}}%
\end{pgfscope}%
\begin{pgfscope}%
\pgfpathrectangle{\pgfqpoint{2.197902in}{0.579321in}}{\pgfqpoint{0.979200in}{0.326250in}}%
\pgfusepath{clip}%
\pgfsetbuttcap%
\pgfsetroundjoin%
\pgfsetlinewidth{0.401500pt}%
\definecolor{currentstroke}{rgb}{0.690196,0.690196,0.690196}%
\pgfsetstrokecolor{currentstroke}%
\pgfsetstrokeopacity{0.250000}%
\pgfsetdash{{1.480000pt}{0.640000pt}}{0.000000pt}%
\pgfpathmoveto{\pgfqpoint{3.177102in}{0.579321in}}%
\pgfpathlineto{\pgfqpoint{3.177102in}{0.905571in}}%
\pgfusepath{stroke}%
\end{pgfscope}%
\begin{pgfscope}%
\pgfsetbuttcap%
\pgfsetroundjoin%
\definecolor{currentfill}{rgb}{0.000000,0.000000,0.000000}%
\pgfsetfillcolor{currentfill}%
\pgfsetlinewidth{0.803000pt}%
\definecolor{currentstroke}{rgb}{0.000000,0.000000,0.000000}%
\pgfsetstrokecolor{currentstroke}%
\pgfsetdash{}{0pt}%
\pgfsys@defobject{currentmarker}{\pgfqpoint{0.000000in}{0.000000in}}{\pgfqpoint{0.000000in}{0.027778in}}{%
\pgfpathmoveto{\pgfqpoint{0.000000in}{0.000000in}}%
\pgfpathlineto{\pgfqpoint{0.000000in}{0.027778in}}%
\pgfusepath{stroke,fill}%
}%
\begin{pgfscope}%
\pgfsys@transformshift{3.177102in}{0.579321in}%
\pgfsys@useobject{currentmarker}{}%
\end{pgfscope}%
\end{pgfscope}%
\begin{pgfscope}%
\definecolor{textcolor}{rgb}{0.000000,0.000000,0.000000}%
\pgfsetstrokecolor{textcolor}%
\pgfsetfillcolor{textcolor}%
\pgftext[x=3.177102in,y=0.565432in,,top]{\color{textcolor}{\rmfamily\fontsize{5.000000}{6.000000}\selectfont\catcode`\^=\active\def^{\ifmmode\sp\else\^{}\fi}\catcode`\%=\active\def%{\%}5.000}}%
\end{pgfscope}%
\begin{pgfscope}%
\pgfpathrectangle{\pgfqpoint{2.197902in}{0.579321in}}{\pgfqpoint{0.979200in}{0.326250in}}%
\pgfusepath{clip}%
\pgfsetbuttcap%
\pgfsetroundjoin%
\pgfsetlinewidth{0.401500pt}%
\definecolor{currentstroke}{rgb}{0.690196,0.690196,0.690196}%
\pgfsetstrokecolor{currentstroke}%
\pgfsetstrokeopacity{0.250000}%
\pgfsetdash{{1.480000pt}{0.640000pt}}{0.000000pt}%
\pgfpathmoveto{\pgfqpoint{2.197902in}{0.579321in}}%
\pgfpathlineto{\pgfqpoint{3.177102in}{0.579321in}}%
\pgfusepath{stroke}%
\end{pgfscope}%
\begin{pgfscope}%
\pgfsetbuttcap%
\pgfsetroundjoin%
\definecolor{currentfill}{rgb}{0.000000,0.000000,0.000000}%
\pgfsetfillcolor{currentfill}%
\pgfsetlinewidth{0.803000pt}%
\definecolor{currentstroke}{rgb}{0.000000,0.000000,0.000000}%
\pgfsetstrokecolor{currentstroke}%
\pgfsetdash{}{0pt}%
\pgfsys@defobject{currentmarker}{\pgfqpoint{0.000000in}{0.000000in}}{\pgfqpoint{0.027778in}{0.000000in}}{%
\pgfpathmoveto{\pgfqpoint{0.000000in}{0.000000in}}%
\pgfpathlineto{\pgfqpoint{0.027778in}{0.000000in}}%
\pgfusepath{stroke,fill}%
}%
\begin{pgfscope}%
\pgfsys@transformshift{2.197902in}{0.579321in}%
\pgfsys@useobject{currentmarker}{}%
\end{pgfscope}%
\end{pgfscope}%
\begin{pgfscope}%
\definecolor{textcolor}{rgb}{0.000000,0.000000,0.000000}%
\pgfsetstrokecolor{textcolor}%
\pgfsetfillcolor{textcolor}%
\pgftext[x=1.966996in, y=0.555209in, left, base]{\color{textcolor}{\rmfamily\fontsize{5.000000}{6.000000}\selectfont\catcode`\^=\active\def^{\ifmmode\sp\else\^{}\fi}\catcode`\%=\active\def%{\%}13.37}}%
\end{pgfscope}%
\begin{pgfscope}%
\pgfpathrectangle{\pgfqpoint{2.197902in}{0.579321in}}{\pgfqpoint{0.979200in}{0.326250in}}%
\pgfusepath{clip}%
\pgfsetbuttcap%
\pgfsetroundjoin%
\pgfsetlinewidth{0.401500pt}%
\definecolor{currentstroke}{rgb}{0.690196,0.690196,0.690196}%
\pgfsetstrokecolor{currentstroke}%
\pgfsetstrokeopacity{0.250000}%
\pgfsetdash{{1.480000pt}{0.640000pt}}{0.000000pt}%
\pgfpathmoveto{\pgfqpoint{2.197902in}{0.742446in}}%
\pgfpathlineto{\pgfqpoint{3.177102in}{0.742446in}}%
\pgfusepath{stroke}%
\end{pgfscope}%
\begin{pgfscope}%
\pgfsetbuttcap%
\pgfsetroundjoin%
\definecolor{currentfill}{rgb}{0.000000,0.000000,0.000000}%
\pgfsetfillcolor{currentfill}%
\pgfsetlinewidth{0.803000pt}%
\definecolor{currentstroke}{rgb}{0.000000,0.000000,0.000000}%
\pgfsetstrokecolor{currentstroke}%
\pgfsetdash{}{0pt}%
\pgfsys@defobject{currentmarker}{\pgfqpoint{0.000000in}{0.000000in}}{\pgfqpoint{0.027778in}{0.000000in}}{%
\pgfpathmoveto{\pgfqpoint{0.000000in}{0.000000in}}%
\pgfpathlineto{\pgfqpoint{0.027778in}{0.000000in}}%
\pgfusepath{stroke,fill}%
}%
\begin{pgfscope}%
\pgfsys@transformshift{2.197902in}{0.742446in}%
\pgfsys@useobject{currentmarker}{}%
\end{pgfscope}%
\end{pgfscope}%
\begin{pgfscope}%
\definecolor{textcolor}{rgb}{0.000000,0.000000,0.000000}%
\pgfsetstrokecolor{textcolor}%
\pgfsetfillcolor{textcolor}%
\pgftext[x=1.966996in, y=0.718334in, left, base]{\color{textcolor}{\rmfamily\fontsize{5.000000}{6.000000}\selectfont\catcode`\^=\active\def^{\ifmmode\sp\else\^{}\fi}\catcode`\%=\active\def%{\%}13.42}}%
\end{pgfscope}%
\begin{pgfscope}%
\pgfpathrectangle{\pgfqpoint{2.197902in}{0.579321in}}{\pgfqpoint{0.979200in}{0.326250in}}%
\pgfusepath{clip}%
\pgfsetbuttcap%
\pgfsetroundjoin%
\pgfsetlinewidth{0.401500pt}%
\definecolor{currentstroke}{rgb}{0.690196,0.690196,0.690196}%
\pgfsetstrokecolor{currentstroke}%
\pgfsetstrokeopacity{0.250000}%
\pgfsetdash{{1.480000pt}{0.640000pt}}{0.000000pt}%
\pgfpathmoveto{\pgfqpoint{2.197902in}{0.905571in}}%
\pgfpathlineto{\pgfqpoint{3.177102in}{0.905571in}}%
\pgfusepath{stroke}%
\end{pgfscope}%
\begin{pgfscope}%
\pgfsetbuttcap%
\pgfsetroundjoin%
\definecolor{currentfill}{rgb}{0.000000,0.000000,0.000000}%
\pgfsetfillcolor{currentfill}%
\pgfsetlinewidth{0.803000pt}%
\definecolor{currentstroke}{rgb}{0.000000,0.000000,0.000000}%
\pgfsetstrokecolor{currentstroke}%
\pgfsetdash{}{0pt}%
\pgfsys@defobject{currentmarker}{\pgfqpoint{0.000000in}{0.000000in}}{\pgfqpoint{0.027778in}{0.000000in}}{%
\pgfpathmoveto{\pgfqpoint{0.000000in}{0.000000in}}%
\pgfpathlineto{\pgfqpoint{0.027778in}{0.000000in}}%
\pgfusepath{stroke,fill}%
}%
\begin{pgfscope}%
\pgfsys@transformshift{2.197902in}{0.905571in}%
\pgfsys@useobject{currentmarker}{}%
\end{pgfscope}%
\end{pgfscope}%
\begin{pgfscope}%
\definecolor{textcolor}{rgb}{0.000000,0.000000,0.000000}%
\pgfsetstrokecolor{textcolor}%
\pgfsetfillcolor{textcolor}%
\pgftext[x=1.966996in, y=0.881459in, left, base]{\color{textcolor}{\rmfamily\fontsize{5.000000}{6.000000}\selectfont\catcode`\^=\active\def^{\ifmmode\sp\else\^{}\fi}\catcode`\%=\active\def%{\%}13.46}}%
\end{pgfscope}%
\begin{pgfscope}%
\pgfpathrectangle{\pgfqpoint{2.197902in}{0.579321in}}{\pgfqpoint{0.979200in}{0.326250in}}%
\pgfusepath{clip}%
\pgfsetrectcap%
\pgfsetroundjoin%
\pgfsetlinewidth{1.505625pt}%
\definecolor{currentstroke}{rgb}{0.000000,0.619608,0.450980}%
\pgfsetstrokecolor{currentstroke}%
\pgfsetdash{}{0pt}%
\pgfpathmoveto{\pgfqpoint{2.194319in}{0.815098in}}%
\pgfpathlineto{\pgfqpoint{2.403972in}{0.608887in}}%
\pgfpathlineto{\pgfqpoint{2.405533in}{0.608272in}}%
\pgfpathlineto{\pgfqpoint{2.407852in}{0.610681in}}%
\pgfpathlineto{\pgfqpoint{2.411453in}{0.618830in}}%
\pgfpathlineto{\pgfqpoint{2.439273in}{0.693570in}}%
\pgfpathlineto{\pgfqpoint{2.460392in}{0.741980in}}%
\pgfpathlineto{\pgfqpoint{2.475080in}{0.771305in}}%
\pgfpathlineto{\pgfqpoint{2.489768in}{0.797108in}}%
\pgfpathlineto{\pgfqpoint{2.504456in}{0.819290in}}%
\pgfpathlineto{\pgfqpoint{2.519144in}{0.837888in}}%
\pgfpathlineto{\pgfqpoint{2.528936in}{0.848312in}}%
\pgfpathlineto{\pgfqpoint{2.538728in}{0.857150in}}%
\pgfpathlineto{\pgfqpoint{2.548520in}{0.864400in}}%
\pgfpathlineto{\pgfqpoint{2.558312in}{0.870057in}}%
\pgfpathlineto{\pgfqpoint{2.568104in}{0.874123in}}%
\pgfpathlineto{\pgfqpoint{2.577896in}{0.876596in}}%
\pgfpathlineto{\pgfqpoint{2.587688in}{0.877481in}}%
\pgfpathlineto{\pgfqpoint{2.597480in}{0.876781in}}%
\pgfpathlineto{\pgfqpoint{2.607272in}{0.874496in}}%
\pgfpathlineto{\pgfqpoint{2.617064in}{0.870630in}}%
\pgfpathlineto{\pgfqpoint{2.626856in}{0.865183in}}%
\pgfpathlineto{\pgfqpoint{2.636648in}{0.858156in}}%
\pgfpathlineto{\pgfqpoint{2.646440in}{0.849557in}}%
\pgfpathlineto{\pgfqpoint{2.739230in}{0.758132in}}%
\pgfpathlineto{\pgfqpoint{2.893540in}{0.606692in}}%
\pgfpathlineto{\pgfqpoint{2.894913in}{0.606509in}}%
\pgfpathlineto{\pgfqpoint{2.896285in}{0.607752in}}%
\pgfpathlineto{\pgfqpoint{2.900025in}{0.617373in}}%
\pgfpathlineto{\pgfqpoint{2.920624in}{0.674108in}}%
\pgfpathlineto{\pgfqpoint{2.940208in}{0.721856in}}%
\pgfpathlineto{\pgfqpoint{2.954896in}{0.753550in}}%
\pgfpathlineto{\pgfqpoint{2.969584in}{0.781576in}}%
\pgfpathlineto{\pgfqpoint{2.984272in}{0.806148in}}%
\pgfpathlineto{\pgfqpoint{2.998960in}{0.827106in}}%
\pgfpathlineto{\pgfqpoint{3.013648in}{0.844425in}}%
\pgfpathlineto{\pgfqpoint{3.023440in}{0.854027in}}%
\pgfpathlineto{\pgfqpoint{3.033232in}{0.862043in}}%
\pgfpathlineto{\pgfqpoint{3.043024in}{0.868466in}}%
\pgfpathlineto{\pgfqpoint{3.052816in}{0.873290in}}%
\pgfpathlineto{\pgfqpoint{3.062608in}{0.876519in}}%
\pgfpathlineto{\pgfqpoint{3.072400in}{0.878155in}}%
\pgfpathlineto{\pgfqpoint{3.082192in}{0.878205in}}%
\pgfpathlineto{\pgfqpoint{3.091984in}{0.876674in}}%
\pgfpathlineto{\pgfqpoint{3.101776in}{0.873564in}}%
\pgfpathlineto{\pgfqpoint{3.111568in}{0.868873in}}%
\pgfpathlineto{\pgfqpoint{3.121360in}{0.862602in}}%
\pgfpathlineto{\pgfqpoint{3.131152in}{0.854752in}}%
\pgfpathlineto{\pgfqpoint{3.150736in}{0.835807in}}%
\pgfpathlineto{\pgfqpoint{3.177102in}{0.809802in}}%
\pgfpathlineto{\pgfqpoint{3.177102in}{0.809802in}}%
\pgfusepath{stroke}%
\end{pgfscope}%
\begin{pgfscope}%
\pgfsetrectcap%
\pgfsetmiterjoin%
\pgfsetlinewidth{0.803000pt}%
\definecolor{currentstroke}{rgb}{0.000000,0.000000,0.000000}%
\pgfsetstrokecolor{currentstroke}%
\pgfsetdash{}{0pt}%
\pgfpathmoveto{\pgfqpoint{2.197902in}{0.579321in}}%
\pgfpathlineto{\pgfqpoint{2.197902in}{0.905571in}}%
\pgfusepath{stroke}%
\end{pgfscope}%
\begin{pgfscope}%
\pgfsetrectcap%
\pgfsetmiterjoin%
\pgfsetlinewidth{0.803000pt}%
\definecolor{currentstroke}{rgb}{0.000000,0.000000,0.000000}%
\pgfsetstrokecolor{currentstroke}%
\pgfsetdash{}{0pt}%
\pgfpathmoveto{\pgfqpoint{3.177102in}{0.579321in}}%
\pgfpathlineto{\pgfqpoint{3.177102in}{0.905571in}}%
\pgfusepath{stroke}%
\end{pgfscope}%
\begin{pgfscope}%
\pgfsetrectcap%
\pgfsetmiterjoin%
\pgfsetlinewidth{0.803000pt}%
\definecolor{currentstroke}{rgb}{0.000000,0.000000,0.000000}%
\pgfsetstrokecolor{currentstroke}%
\pgfsetdash{}{0pt}%
\pgfpathmoveto{\pgfqpoint{2.197902in}{0.579321in}}%
\pgfpathlineto{\pgfqpoint{3.177102in}{0.579321in}}%
\pgfusepath{stroke}%
\end{pgfscope}%
\begin{pgfscope}%
\pgfsetrectcap%
\pgfsetmiterjoin%
\pgfsetlinewidth{0.803000pt}%
\definecolor{currentstroke}{rgb}{0.000000,0.000000,0.000000}%
\pgfsetstrokecolor{currentstroke}%
\pgfsetdash{}{0pt}%
\pgfpathmoveto{\pgfqpoint{2.197902in}{0.905571in}}%
\pgfpathlineto{\pgfqpoint{3.177102in}{0.905571in}}%
\pgfusepath{stroke}%
\end{pgfscope}%
\end{pgfpicture}%
\makeatother%
\endgroup%

	\caption{(a) corrente no dreno do transistor, (b) tensão dreno--source do transistor e (c) tensão de saída do circuito Flyback para $K~L1~L2~0.99$, circuito \textit{clamper} RCD e circuito \textit{snubber}.}
	\label{fig:ex4_4}
\end{figure}